\documentclass[12pt]{report}

% ---------------------------------------------------------------
% 字型與中文排版（和 C++ 演算法解題 1、2 相同）
% ---------------------------------------------------------------
% 作者在 macOS 上用 Times New Roman／Songti TC／PingFang TC；
% 其他系統（例如 Linux）找不到這些字型時，改用 TeX Gyre Termes 與
% Noto Serif／Sans CJK TC，同一份原始碼兩邊都能編。
\usepackage{fontspec}
\usepackage{xeCJK}
\IfFontExistsTF{Times New Roman}
    {\setmainfont{Times New Roman}}
    {\setmainfont{TeX Gyre Termes}}
\IfFontExistsTF{Songti TC}
    {\setCJKmainfont{Songti TC}}
    {\setCJKmainfont{Noto Serif CJK TC}}
\IfFontExistsTF{PingFang TC}
    {\setCJKsansfont{PingFang TC}}
    {\setCJKsansfont{Noto Sans CJK TC}}
% 等寬字用 Latin Modern Mono（只在 TeX Live 樹裡，用檔名讓 kpathsea 找）。
\setmonofont{lmmono10-regular}[
    Extension = .otf,
    AutoFakeBold = 2.5,
]

% ---------------------------------------------------------------
% 本書排版套件（packages/ 目錄）
% ---------------------------------------------------------------
\usepackage{sleek}
\usepackage{sleek-code}
\usepackage{sleek-diagrams}
\usepackage{sleek-recap}
\usepackage{sleek-aki}
\usepackage{sleek-pseudo}   % CLRS 風格的虛擬碼框（algpseudocode）
\usepackage{sleek-math}     % 定義、定理、證明（amsthm）
% 注意：sleek-theorems.sty 不載入，原因見 sleek-math.sty 開頭。
% 注意：sleek-listings.sty 不載入，程式碼一律用 sleek-code.sty 的 minted 環境。

% 章節標題中文化：「第 N 章」
\usepackage{titlesec}
\titleformat{\chapter}{\normalfont\huge\bfseries}{第~\thechapter~章}{1em}{}
\titlespacing*{\chapter}{0pt}{-20pt}{20pt}
% 頁首寫「第 N 章 章名」，不用英文的 CHAPTER。
\renewcommand{\chaptermark}[1]{\markboth{第~\thechapter~章\quad #1}{}}

% 行內程式碼（\texttt）不能斷行，允許 TeX 多拉開一點字距。
\setlength{\emergencystretch}{3em}

\hypersetup{
    pdftitle={資料結構與演算法：C 語言、虛擬碼與研究所考試},
    pdfauthor={林子祐},
    pdfsubject={註解 Annotated 演算法系列 第 4 本},
}

\title{資料結構與演算法\\[0.3em]{\LARGE C 語言、虛擬碼與研究所考試}\\[1.2em]{\large 註解 Annotated 演算法系列}}
\author{林子祐}
\date{\today}

\begin{document}

\maketitle
\romantableofcontents

% ---------------------------------------------------------------
% 各章內容放在 "chapters/<單元>/chapter.tex"。章號一開始就定好，
% 還沒寫的章先註解掉，已寫的章前面用 \setcounter{chapter}{N - 1} 定章號。
% （Makefile 編單章時會自己從目錄名稱算章號。）
% ---------------------------------------------------------------

% 第一部：基礎
\chapter{演算法分析與漸近符號}
\label{ch:analysis}

% 本章的圖共用的樣式
\tikzset{
    sortedcell/.style={cell, fill=coreFill, draw=coreStroke},
    growth/.style={thick, smooth},
    glabel/.style={font=\scriptsize, anchor=west},
}
% 插入排序追蹤圖的一列：\insrow{y}{左邊的說明}{數值串列}{已排好的格數}{剛插入的格子}
\newcommand{\insrow}[5]{%
    \node[anchor=east, font=\small] at (0.3, #1) {#2};
    \foreach \v [count=\k] in {#3} {
        \ifnum\k=#5\relax
            \node[hlcell] at (\k*0.9, #1) {\v};
        \else\ifnum\k>#4\relax
            \node[cell] at (\k*0.9, #1) {\v};
        \else
            \node[sortedcell] at (\k*0.9, #1) {\v};
        \fi\fi
    }
}

系辦公室收到一批暑期實習的申請，每份申請寫著一個九位數的學號。 有些同學怕沒
送到，同一份申請送了兩次，助理想先知道這批申請裡有沒有重複的學號。 第一種做
法是拿每一份去和其他每一份比對學號；第二種做法是準備一張很大的表，學號每出現一
次就在表上對應的格子打勾，打勾前發現格子已經勾過，就是重複。 申請從 $1000$ 份變
成 $8000$ 份時，兩種做法各要多做幾倍的工作？ 往下讀之前，先估計一下。

這一章講怎麼回答這種問題。 我們不用碼錶量秒數，改成數一個演算法做了多少個基本
步驟，把步數寫成輸入大小 $n$ 的函數，再用\textbf{漸近符號}（asymptotic notation）
$O$、$\Omega$、$\Theta$ 只留下 $n$ 很大時決定快慢的那一項。 中間會用插入排序當例
子，示範三件之後每一章都會用到的事：逐行算出執行次數、用迴圈不變量說明程式為什麼
對、分別討論最好、最壞和平均的情況。

這一章用到《C 程式設計基礎》第~9~章的指標參數（用指標把計數帶回呼叫的地方）和
第~19~章的 \texttt{malloc}、\texttt{calloc}、\texttt{free}，用到時會先在框裡複習。
數學上只需要高中的指數、對數和數列求和；本書用到的求和公式整理在附錄~A。

\section{演算法和它的成本}
\label{sec:ana-cost}

\textbf{演算法}（algorithm）是一串定義清楚的步驟，把\textbf{輸入}（input）變成
\textbf{輸出}（output）。 「定義清楚」的意思是每一步都沒有歧義，任何人照著做都
得到一樣的結果；而且對每一個合法的輸入，它都會在有限步之後停下來，給出正確的輸
出。 一個\textbf{問題}（problem）規定了輸入和輸出的關係，例如「給 $n$ 個學號，回
答有沒有兩個相同」；把具體的一批學號交給它，叫做這個問題的一個\textbf{實例}
（instance）。

\subsection{數步驟，不量秒數}

同一支程式在不同電腦、不同編譯器上跑，秒數可以差好幾倍；同一台電腦上，秒數也會
受到其他程式影響。 我們想要的是一個只和演算法本身有關的量，所以改成數步驟。 數
法建立在一個簡化的電腦模型上，叫\textbf{隨機存取機}（random-access machine，簡稱
RAM 模型）：

\begin{itemize}
\item 一個\textbf{基本運算}花固定的時間：整數的加減乘除和取餘數、比較兩個數、
指定一個變數、讀寫陣列的一格、呼叫函式和返回。
\item 讀寫記憶體的任何一個位置都花一樣的時間，不管位置在哪裡。
\item 一次只做一個運算，沒有平行。
\item 一個整數放得進一個「字組」（word），大小足以存下 $n$ 這個量級的數。 不能把
任意大的數塞進一個字組，假裝運算還是固定時間；第~\ref{sec:ana-input-size}~節會
看到這件事為什麼重要。
\end{itemize}

這個模型和真正的電腦不完全一樣（真正的電腦有快取，讀相鄰的記憶體比較快），但用它
算出來的步數，在 $n$ 變大時的變化趨勢通常和實際的秒數一致。 這本書說「時間」
或「執行時間」（running time），指的都是在這個模型下的步數。

步數會隨輸入變多而變多，所以要先說好怎麼量輸入的大小。 對大部分問題，\textbf{輸
入大小}（input size）就是資料的個數：排序 $n$ 個數，大小是 $n$；圖的問題要用兩個
數，點數 $n$ 和邊數 $m$（第~21~章）。 有些問題的輸入只有一兩個整數，這時要改用
位元數來量，第~\ref{sec:ana-input-size}~節再講。

\subsection{例子：找重複的學號}

章首的兩種做法寫成虛擬碼如下。 學號都在 $[\mathit{base}, \mathit{base} + M)$ 這
個範圍裡，第二種做法用一個長度 $M$ 的陣列 $\mathit{seen}$ 打勾。

\begin{pseudo}[label=alg:dup-pairs]{兩兩比較}
\Function{Has-Dup-Pairs}{$A, n$}
    \For{$i \gets 1$ \To $n - 1$}
        \For{$j \gets i + 1$ \To $n$}
            \If{$A[i] = A[j]$}
                \State \Return \TRUE
            \EndIf
        \EndFor
    \EndFor
    \State \Return \FALSE
\EndFunction
\end{pseudo}

\begin{pseudo}[label=alg:dup-marks]{在表上打勾}
\Function{Has-Dup-Marks}{$A, n, \mathit{base}, M$}
    \State 建立陣列 $\mathit{seen}[0..M-1]$，每一格都設成 \FALSE
    \For{$i \gets 1$ \To $n$}
        \State $k \gets A[i] - \mathit{base}$
        \If{$\mathit{seen}[k]$}
            \State \Return \TRUE
        \EndIf
        \State $\mathit{seen}[k] \gets \TRUE$
    \EndFor
    \State \Return \FALSE
\EndFunction
\end{pseudo}

沒有重複的時候，第一種做法把每一對都比過一次。 $i = 1$ 時比 $n - 1$ 次，$i = 2$
時比 $n - 2$ 次……最後 $i = n - 1$ 時比 $1$ 次，總共
\[
(n - 1) + (n - 2) + \cdots + 1 = \frac{n(n-1)}{2}
\]
次。 第二種做法每個學號查一次表，共 $n$ 次，但開始前要把 $M$ 格都設成
\FALSE，這也要 $M$ 步。

下面的 C 程式把兩種做法各寫成一個函式，用指標參數把「比較了幾次」「查了幾次表」
帶回 \texttt{main}。 學號用 \texttt{BASE + i * 7919 \% RANGE} 產生，$7919$ 和
$10^6$ 沒有公因數，所以 $n \le 10^6$ 時產生的學號都不同，這是兩種做法最花時間的
情況（找到重複就可以提早結束）。

\begin{cbookrecap}[第~9~章 9.2 節、第~19~章 19.3 節]
函式只能回傳一個值。 想多帶回一個數，就多傳一個指標參數，在函式裡透過
\texttt{*cmp} 改它指到的變數，例如下面的 \texttt{has\_dup\_pairs(id, n, \&a)}。

\texttt{malloc(n * sizeof *id)} 向系統要一塊能放 $n$ 個 \texttt{int} 的記憶體，
要不到時回傳 \texttt{NULL}。 \texttt{calloc(count, size)} 一樣要記憶體，但會把整
塊設成 $0$。 用完要 \texttt{free}，同一塊不能 \texttt{free} 兩次。 本書所有的程式
都用 \texttt{-fsanitize=address} 檢查過，忘了 \texttt{free} 或是讀寫超出範圍都會
被抓出來。
\end{cbookrecap}

\begin{example}[label=ex:dup-check]{c}{(dup\_check.c) 找重複學號的兩種做法}
#include <stdio.h>
#include <stdlib.h>

#define BASE 115000000          // 學號的最小值
#define RANGE 1000000           // 學號落在 [BASE, BASE + RANGE)

// 方法一：兩兩比較。*cmp 記錄比較了幾次。
int has_dup_pairs(const int *id, int n, long long *cmp)
{
    *cmp = 0;
    for (int i = 0; i < n; i++) {
        for (int j = i + 1; j < n; j++) {
            (*cmp)++;
            if (id[i] == id[j]) {
                return 1;
            }
        }
    }
    return 0;
}

// 方法二：用一個長度 RANGE 的陣列做記號。*look 記錄查了幾次記號。
int has_dup_marks(const int *id, int n, long long *look)
{
    char *seen = calloc(RANGE, 1);   // 全部設成 0，代表都還沒看過
    if (seen == NULL) {
        return -1;
    }
    int found = 0;
    *look = 0;
    for (int i = 0; i < n && !found; i++) {
        (*look)++;
        int k = id[i] - BASE;
        if (seen[k]) {
            found = 1;
        }
        seen[k] = 1;
    }
    free(seen);
    return found;
}

int main(void)
{
    printf("n pairs marks result\n");
    for (int n = 1000; n <= 8000; n *= 2) {
        int *id = malloc(n * sizeof *id);
        if (id == NULL) {
            return 1;
        }
        for (int i = 0; i < n; i++) {   // 產生 n 個互不相同的學號
            id[i] = BASE + (int)((i * 7919LL) % RANGE);
        }
        long long a, b;
        int r1 = has_dup_pairs(id, n, &a);
        int r2 = has_dup_marks(id, n, &b);
        printf("%d %lld %lld %s\n", n, a, b,
               (r1 == 0 && r2 == 0) ? "no-dup" : "dup");
        free(id);
    }
    return 0;
}
\end{example}
\iosample{}{n pairs marks result \\ 1000 499500 1000 no-dup \\ 2000 1999000 2000 no-dup \\ 4000 7998000 4000 no-dup \\ 8000 31996000 8000 no-dup}

$n$ 每變成兩倍，兩兩比較的次數大約變成四倍，查表的次數剛好變成兩倍。 從 $1000$ 份
到 $8000$ 份，第一種做法多做了大約 $64$ 倍的工作，第二種只多了 $8$ 倍。 章首的問
題答案就是這兩個數。

第二種做法的代價在別的地方：不管 $n$ 多小，它都要先準備 $M = 10^6$ 格的表，花
$M$ 個位元組的記憶體，\texttt{calloc} 把它們設成 $0$ 也要時間，所以它的總步數比
較像 $n + M$。 學號如果是任意的 $64$ 位元整數，$M$ 會大到根本開不出這張表。 這時
要改用第~12~章的雜湊，或是先排序（第~14~章）再比相鄰的兩個。 一個演算法好不
好，要同時看時間、記憶體，以及它對輸入做了什麼假設。

\section{本書的虛擬碼和 C 程式}
\label{sec:ana-conventions}

這本書每個演算法都會給兩種寫法：虛擬碼和 C 程式。 虛擬碼的寫法和多數演算法教科
書、考卷上的寫法相近，用來講道理、算次數、手算追蹤；C 程式用來實際執行，驗證書上
說的結果。 兩種寫法有幾個地方不同，下表整理在一起，之後各章不再重複。 表裡只對照寫法，$i$ 的值兩邊相同；陣列編號差 $1$ 怎麼換，見表後的說明。

\begin{center}
\begin{tabular}{p{0.15\linewidth}p{0.33\linewidth}p{0.42\linewidth}}
\toprule
項目 & 虛擬碼 & C \\
\midrule
陣列編號 & 從 $1$ 開始：$A[1..n]$ & 從 $0$ 開始：\texttt{a[0..n-1]} \\
指定 & $x \gets x + 1$ & \texttt{x = x + 1} \\
比較相等 & $x = y$、不等 $x \ne y$ & \texttt{x == y}、\texttt{x != y} \\
計數迴圈 & \textbf{for} $i \gets 1$ \textbf{to} $n$（包含 $n$）&
\texttt{for (int i = 1; i <= n; i++)} \\
往下數 & \textbf{for} $i \gets n$ \textbf{downto} $1$（包含 $1$）&
\texttt{for (int i = n; i >= 1; i--)} \\
且、或 & \textbf{and}、\textbf{or}，左邊決定結果時右邊不算 & \texttt{\&\&}、\texttt{||}，
一樣不算右邊 \\
空指標 & \NIL & \texttt{NULL} \\
物件的欄位 & $x.\mathit{key}$，$x$ 是指向物件的指標 & \texttt{x->key} \\
程序名稱 & \proc{Insertion-Sort}$(A, n)$，名稱用小型大寫 & \texttt{insertion\_sort(a, n)} \\
回傳 & \textbf{return} $x$ & \texttt{return x;} \\
註解 & 寫在行尾，以 // 開頭 & \texttt{//} 或 \texttt{/* */} \\
區塊 & 用縮排表示，不寫 \textbf{end} & 用大括號 \\
\bottomrule
\end{tabular}
\end{center}

另外幾個約定：

\begin{itemize}
\item 虛擬碼的變數都是區域變數；陣列和物件傳進程序時，傳的是指標，程序裡改了內
容，呼叫的人看得到。 C 也一樣，陣列參數其實是指標。
\item 虛擬碼裡的數學運算不會溢位，除法要寫明是 $\lfloor a / b \rfloor$ 還是
$\lceil a / b \rceil$。 C 程式要自己注意溢位，會超過 \texttt{int} 的計數一律用
\texttt{long long}。
\item 虛擬碼可以寫一句中文代替一段細節，例如「建立陣列 $\mathit{seen}[0..M-1]$，
每一格都設成 \FALSE」，只要那一句的成本清楚。
\item 好幾個字母的變數名稱用斜體整個寫出，例如 $\mathit{key}$、$\mathit{seen}$，指的是一
個變數，不是幾個字母相乘。
\item $\lg n$ 表示 $\log_2 n$，$\ln n$ 表示自然對數。
\end{itemize}

把虛擬碼翻成 C 時最容易錯的地方是編號差 $1$。 這本書的做法是：C 程式的陣列一律從
$0$ 開始，虛擬碼的 $A[i]$ 對應 \texttt{a[i - 1]}，翻譯時整段一起改，不要一行一行
加減 $1$。 例如虛擬碼的 \textbf{for} $i \gets 2$ \textbf{to} $n$ 處理第 $2$ 到第
$n$ 個元素，翻成 C 是 \texttt{for (i = 1; i < n; i++)}，迴圈裡的 $A[i]$ 直接寫
\texttt{a[i]}。

\section{插入排序：逐行計算成本}
\label{sec:ana-insertion}

\textbf{排序}（sorting）的問題是：輸入 $n$ 個數 $a_1, a_2, \ldots, a_n$，輸出它們由
小到大的排列。 第~14~章和第~15~章會講十種左右的排序方法；這一節先用最簡單的一種，示
範怎麼逐行算出執行次數。

\textbf{插入排序}（insertion sort）的想法是讓陣列左邊有一段已經排好。 一開始這段
只有第一個數。 每一輪把這段右邊的下一個數拿出來，叫它 $\mathit{key}$，從這段的最
右邊往左看，比 $\mathit{key}$ 大的數都往右挪一格，直到遇到不比它大的數，或是走到
最左邊，再把 $\mathit{key}$ 放進空出來的位置。 這樣排好的一段就多了一個數。

\begin{pseudo}[label=alg:insertion]{插入排序}
\Procedure{Insertion-Sort}{$A, n$}
    \For{$i \gets 2$ \To $n$}
        \State $\mathit{key} \gets A[i]$
        \State $j \gets i - 1$
        \While{$j \ge 1$ \AND $A[j] > \mathit{key}$}
            \State $A[j + 1] \gets A[j]$ \Comment{往右挪一格}
            \State $j \gets j - 1$
        \EndWhile
        \State $A[j + 1] \gets \mathit{key}$
    \EndFor
\EndProcedure
\end{pseudo}

\subsection{手算追蹤}

用 $A = \langle 27, 12, 35, 6, 19, 8 \rangle$ 跑一次。 下圖每一列是外層迴圈做完一
輪之後的陣列，藍底是已經排好的一段，橘框是這一輪插進去的 $\mathit{key}$。

\begin{center}
\begin{tikzpicture}
\insrow{0}{開始}{27,12,35,6,19,8}{1}{0}
\insrow{-1.1}{$i = 2$，$\mathit{key} = 12$}{12,27,35,6,19,8}{2}{1}
\insrow{-2.2}{$i = 3$，$\mathit{key} = 35$}{12,27,35,6,19,8}{3}{3}
\insrow{-3.3}{$i = 4$，$\mathit{key} = 6$}{6,12,27,35,19,8}{4}{1}
\insrow{-4.4}{$i = 5$，$\mathit{key} = 19$}{6,12,19,27,35,8}{5}{3}
\insrow{-5.5}{$i = 6$，$\mathit{key} = 8$}{6,8,12,19,27,35}{6}{2}
\foreach \k in {1,...,6} { \node[idxlbl] at (\k*0.9, -6.2) {\k}; }
\end{tikzpicture}
\diagcaption{插入排序每一輪之後的陣列。 藍底是已排好的前段，橘框是這一輪的
$\mathit{key}$ 最後放的位置。}
\label{fig:ana-insertion-trace}
\end{center}

每一輪的細節列成下表。 「比較」是第 5 行 $A[j] > \mathit{key}$ 實際比了幾次（$j$
已經變成 $0$ 時，因為 \textbf{and} 左邊已經不成立，右邊不算），「挪動」是第 6 行
執行的次數。

\begin{center}
\begin{tabular}{cccccl}
\toprule
$i$ & $\mathit{key}$ & 比較 & 挪動 & 停下的原因 & 這一輪之後 \\
\midrule
2 & 12 & 1 & 1 & $j = 0$ & 12 27 35 6 19 8 \\
3 & 35 & 1 & 0 & $A[2] = 27 \le 35$ & 12 27 35 6 19 8 \\
4 & 6 & 3 & 3 & $j = 0$ & 6 12 27 35 19 8 \\
5 & 19 & 3 & 2 & $A[2] = 12 \le 19$ & 6 12 19 27 35 8 \\
6 & 8 & 5 & 4 & $A[1] = 6 \le 8$ & 6 8 12 19 27 35 \\
\midrule
合計 & & 13 & 10 & & \\
\bottomrule
\end{tabular}
\end{center}

\begin{examnote}
考卷上的追蹤題，通常要求寫出「每一輪之後」的陣列，像上表最右欄那樣一輪一列。 先
看清楚題目要的是每一輪之後，還是每一次交換、挪動之後，兩者的列數不同。 題目問
比較次數時，記得 $j$ 走到 $0$ 那一次不算比較。

有的課本把 $\mathit{key}$ 先放進 $A[0]$ 當\textbf{哨兵}（sentinel），迴圈條件只寫
$A[j] > \mathit{key}$，省掉 $j \ge 1$ 的檢查。 這個版本裡 $j$ 走到 $0$ 時會和哨兵比
一次，所以第 $i$ 輪最多比 $i$ 次；上面的例子照這個版本數，比較次數是 $15$，不是
$13$。 題目給了虛擬碼時，照題目的版本數。
\end{examnote}

\subsection{C 程式}

下面的程式從輸入讀進好幾組陣列，每組排序之後印出結果，以及比較和挪動的次數。
C 的陣列從 $0$ 開始，所以外層從 \texttt{i = 1} 開始，內層的條件 $j \ge 1$ 變成
\texttt{j >= 0}。 為了數比較次數，把 \textbf{and} 的兩半拆開：\texttt{j >= 0}
寫在 \texttt{while} 的條件，\texttt{a[j] <= key} 寫在迴圈裡面，比之前先加一。

\begin{example}[label=ex:insertion-count]{c}{(insertion\_count.c) 插入排序並計算比較和挪動的次數}
#include <stdio.h>
#include <stdlib.h>

// 插入排序，a[0..n-1] 由小到大。*cmp 記錄 a[j] > key 比較了幾次，
// *shift 記錄往右搬了幾次。
void insertion_sort(int *a, int n, long long *cmp, long long *shift)
{
    *cmp = 0;
    *shift = 0;
    for (int i = 1; i < n; i++) {
        int key = a[i];
        int j = i - 1;
        while (j >= 0) {
            (*cmp)++;
            if (a[j] <= key) {
                break;
            }
            a[j + 1] = a[j];
            (*shift)++;
            j--;
        }
        a[j + 1] = key;
    }
}

int main(void)
{
    int t;
    if (scanf("%d", &t) != 1) {
        return 1;
    }
    while (t--) {
        int n;
        if (scanf("%d", &n) != 1 || n <= 0) {
            return 1;
        }
        int *a = malloc(n * sizeof *a);
        if (a == NULL) {
            return 1;
        }
        for (int i = 0; i < n; i++) {
            if (scanf("%d", &a[i]) != 1) {
                free(a);
                return 1;
            }
        }
        long long cmp, shift;
        insertion_sort(a, n, &cmp, &shift);
        for (int i = 0; i < n; i++) {
            printf("%d%c", a[i], i + 1 < n ? ' ' : '\n');
        }
        printf("comparisons = %lld, shifts = %lld\n", cmp, shift);
        free(a);
    }
    return 0;
}
\end{example}
\iosample{3 \\ 6 \\ 4 9 13 21 30 38 \\ 6 \\ 38 30 21 13 9 4 \\ 6 \\ 27 12 35 6 19 8}{4 9 13 21 30 38 \\ comparisons = 5, shifts = 0 \\ 4 9 13 21 30 38 \\ comparisons = 15, shifts = 15 \\ 6 8 12 19 27 35 \\ comparisons = 13, shifts = 10}

三組輸入都是 $6$ 個數，比較次數卻從 $5$ 到 $15$。 已經排好的陣列每一輪只比一次就
停；完全相反的陣列每一輪都要一路挪到最左邊。 執行時間不只和 $n$ 有關，也和輸入
長什麼樣子有關，這一點在第~\ref{sec:ana-cases}~節會再整理。

\subsection{逐行計算}
\label{sec:ana-line-cost}

現在對一般的 $n$ 算出虛擬碼~\ref{alg:insertion} 的執行時間。 假設第 $k$ 行執行一
次花 $c_k$ 的時間（$c_k$ 是常數，第 2 行到第 8 行各自不同），再數每一行執行幾次。
外層的 $i$ 跑 $2$ 到 $n$。 對固定的 $i$，令 $t_i$ 是第 5 行 \textbf{while} 的條件
被檢查的次數：挪了幾次就檢查幾次，再加上最後讓迴圈停下來的那一次，所以挪動
$t_i - 1$ 次。

\begin{center}
\begin{tabular}{clcc}
\toprule
行 & 內容 & 每次的成本 & 執行次數 \\
\midrule
2 & \textbf{for} $i \gets 2$ \textbf{to} $n$ & $c_2$ & $n$ \\
3 & $\mathit{key} \gets A[i]$ & $c_3$ & $n - 1$ \\
4 & $j \gets i - 1$ & $c_4$ & $n - 1$ \\
5 & \textbf{while} $j \ge 1$ \textbf{and} $A[j] > \mathit{key}$ & $c_5$ &
$\sum_{i=2}^{n} t_i$ \\
6 & $A[j + 1] \gets A[j]$ & $c_6$ & $\sum_{i=2}^{n} (t_i - 1)$ \\
7 & $j \gets j - 1$ & $c_7$ & $\sum_{i=2}^{n} (t_i - 1)$ \\
8 & $A[j + 1] \gets \mathit{key}$ & $c_8$ & $n - 1$ \\
\bottomrule
\end{tabular}
\end{center}

第 2 行執行 $n$ 次：$i = 2, 3, \ldots, n$ 各檢查一次，再加上 $i = n + 1$ 時發現超過
而離開的那一次。 把每一行的成本乘上次數加起來，得到
\[
T(n) = c_2 n + (c_3 + c_4 + c_8)(n - 1) + c_5 \sum_{i=2}^{n} t_i
+ (c_6 + c_7) \sum_{i=2}^{n} (t_i - 1).
\]

\textbf{最好的情況}：陣列本來就排好了。 每一輪第一次檢查就發現 $A[i-1] \le
\mathit{key}$，$t_i = 1$，挪動 $0$ 次。
\[
T(n) = c_2 n + (c_3 + c_4 + c_5 + c_8)(n - 1) = (c_2 + c_3 + c_4 + c_5 + c_8)\, n
- (c_3 + c_4 + c_5 + c_8).
\]
這是 $an + b$ 的形式，$a$、$b$ 是由各個 $c_k$ 決定的常數，執行時間是 $n$ 的一次
函數。

\textbf{最壞的情況}：陣列由大到小排。 第 $i$ 輪的 $\mathit{key}$ 比前面每一個數都
小，要挪 $i - 1$ 次，最後在 $j = 0$ 時停下，$t_i = i$。 用
\[
\sum_{i=2}^{n} i = \frac{n(n+1)}{2} - 1, \qquad
\sum_{i=2}^{n} (i - 1) = \frac{n(n-1)}{2}
\]
代進去：
\[
T(n) = \left(\frac{c_5 + c_6 + c_7}{2}\right) n^2
+ \left(c_2 + c_3 + c_4 + c_8 + \frac{c_5 - c_6 - c_7}{2}\right) n
- (c_3 + c_4 + c_5 + c_8).
\]
這是 $an^2 + bn + c$ 的形式，執行時間是 $n$ 的二次函數。

這些 $c_k$ 是多少，要看機器和編譯器，算不出來，也不重要。 $n$ 很大時，$an^2$ 這一
項會遠遠超過其他項，係數 $a$ 只影響「大多少倍」，不影響「$n$ 變兩倍時時間變幾
倍」。 下一節介紹的漸近符號，就是把這種「只看最高次項、不看係數」的說法寫得精確。

\subsection{挪動次數和逆序數}

在陣列裡，如果 $i < j$ 而 $A[i] > A[j]$，就說 $(i, j)$ 是一個\textbf{逆序對}
（inversion）。 插入排序每挪動一次，剛好消掉一個逆序對：$A[j]$ 往右越過
$\mathit{key}$，原本 $A[j]$ 在 $\mathit{key}$ 左邊而且比較大，挪完之後兩者的順序對
了，而其他數對的前後關係都沒有變。 排好的陣列沒有逆序對，所以
\[
\text{插入排序的挪動次數} = \text{輸入陣列的逆序對個數}.
\]
例如 $\langle 27, 12, 35, 6, 19, 8 \rangle$ 的逆序對有 $(27,12)$、$(27,6)$、
$(27,19)$、$(27,8)$、$(12,6)$、$(12,8)$、$(35,6)$、$(35,19)$、$(35,8)$、$(19,8)$，
共 $10$ 個（這裡直接寫數值，不寫位置），和程式印出的 \texttt{shifts = 10} 一樣。
$n$ 個數最多有 $\binom{n}{2} = n(n-1)/2$ 個逆序對，發生在完全相反的時候。

\section{迴圈不變量}
\label{sec:ana-invariant}

算完執行時間，還要說明插入排序真的會把陣列排好。 測幾組輸入只能說明那幾組是對
的；要說明每一組都對，常用的工具是\textbf{迴圈不變量}（loop invariant）：一個和迴
圈變數有關的敘述，每一輪開始時都成立。 證明分成三步。

\begin{enumerate}
\item \textbf{初始}：第一輪開始之前，敘述成立。
\item \textbf{維持}：如果某一輪開始時敘述成立，那麼這一輪做完、下一輪開始時，敘述
仍然成立。
\item \textbf{結束}：迴圈停下來時，把敘述和「讓迴圈停下來的條件」合在一起，可以推
出我們要的結論。
\end{enumerate}

前兩步和數學歸納法一樣，第三步是歸納法沒有的：歸納法證明對所有正整數都成立，迴
圈卻會停，停的時候要把不變量用起來。

\textbf{插入排序的不變量}：外層迴圈每一輪開始時（$i$ 是這一輪的值），$A[1..i-1]$
裡的數就是原本放在 $A[1..i-1]$ 的那些數，而且已經由小到大排好。

\begin{itemize}
\item \textbf{初始}：第一輪 $i = 2$，$A[1..1]$ 只有一個數，就是原本的 $A[1]$，一
個數一定排好了。
\item \textbf{維持}：這一輪把 $\mathit{key} = A[i]$ 存起來，內層迴圈把
$A[i-1], A[i-2], \ldots$ 之中比 $\mathit{key}$ 大的數依序往右挪一格，停在第一個
不比 $\mathit{key}$ 大的 $A[j]$（或 $j = 0$）。 因為 $A[1..i-1]$ 原本排好了，被
挪的是最右邊連續的一段，也就是所有比 $\mathit{key}$ 大的數；它們挪動後順序不
變。 $\mathit{key}$ 放進 $A[j+1]$ 之後，左邊的數都不大於它，右邊的數都大於它，
所以 $A[1..i]$ 排好了，而且裡面的數正是原本 $A[1..i]$ 的那些數。 下一輪的 $i$
加一，不變量成立。（嚴格的寫法要對內層迴圈另外寫一個不變量，這裡用文字帶過。）
\item \textbf{結束}：外層迴圈在 $i = n + 1$ 時停下。 代進不變量：$A[1..n]$ 是原本
的 $n$ 個數，而且排好了，這正是排序要的結果。
\end{itemize}

\begin{examnote}
用迴圈不變量證明時，三步都要寫出來，而且不變量要寫成和迴圈變數有關的具體敘述
（「$A[1..i-1]$ 排好了」），不能只寫「陣列越來越有序」。 結束那一步要明白寫出迴
圈停下時變數的值（這裡是 $i = n + 1$），再代進不變量。
\end{examnote}

不變量要選得剛好：太弱推不出結論，太強維持不住。 例如只寫「$A[1..i-1]$ 排好了」，
少了「是原本的那些數」，那麼一個把整個陣列都改成 $0$ 的程式也滿足它。
觀念題第~\ref{ana-q-lower-bound}~題的二分搜尋、第~14~章快速排序的分割（partition）都會再用到這個工具。

\section{漸近符號}
\label{sec:ana-notation}

第~\ref{sec:ana-line-cost}~節算出插入排序最壞的情況是 $an^2 + bn + c$。 $n = 10^4$
時，$an^2$ 比 $bn$ 大了大約 $10^4 \cdot a/b$ 倍，只要 $a/b$ 不是非常小，低次項
幾乎沒有影響。 所以我們只想說「它是 $n^2$ 這一級的」。 這一節把這句話定義清楚。

這一節的函數都定義在正整數 $n$ 上，而且 $n$ 夠大以後取值不是負的（執行時間不會是
負的）。

\subsection{$\Theta$、$O$、$\Omega$}

\begin{definition}[$\Theta$]
\label{def:theta}
$\Theta(g(n))$ 是所有滿足下面條件的函數 $f(n)$ 組成的集合：存在正的常數 $c_1$、
$c_2$ 和 $n_0$，使得所有 $n \ge n_0$ 都有
\[
0 \le c_1 g(n) \le f(n) \le c_2 g(n).
\]
\end{definition}

白話是：$n$ 夠大以後，$f(n)$ 被夾在 $g(n)$ 的兩個常數倍之間。 $n_0$ 以前怎麼樣都
不管。 $f(n) \in \Theta(g(n))$ 習慣寫成 $f(n) = \Theta(g(n))$，這裡的等號意思是
「屬於」，不能左右對調：寫 $\Theta(g(n)) = f(n)$ 沒有意義。

\begin{center}
\begin{tikzpicture}[x=0.85cm, y=0.7cm]
\draw[-{Stealth[length=2mm]}] (0,0) -- (8.6,0) node[right, font=\small] {$n$};
\draw[-{Stealth[length=2mm]}] (0,0) -- (0,8.2);
\draw[thick, coreStroke, domain=0:8, samples=40] plot (\x, {1.4*(0.08*\x*\x+0.5)})
    node[right, font=\small] {$c_2 g(n)$};
\draw[thick, coreStroke, dashed, domain=0:8, samples=40] plot (\x, {0.6*(0.08*\x*\x+0.5)})
    node[right, font=\small] {$c_1 g(n)$};
\draw[thick, memStroke] plot[smooth] coordinates {(0.2,2.0) (1,1.6) (2,0.95) (3,1.3)
    (4,1.9) (5,2.3) (6,3.6) (7,4.2) (8,5.6)} node[right, font=\small] {$f(n)$};
\draw[densely dotted] (2,0) -- (2,7.6);
\node[below, font=\small] at (2,0) {$n_0$};
\end{tikzpicture}
\diagcaption{$f(n) = \Theta(g(n))$：在 $n_0$ 右邊，$f(n)$ 一直落在 $c_1 g(n)$ 和
$c_2 g(n)$ 之間。 $n_0$ 左邊的部分不影響結論。}
\label{fig:ana-theta}
\end{center}

只要上界或只要下界時，分別用 $O$ 和 $\Omega$。

\begin{definition}[$O$ 和 $\Omega$]
\label{def:big-o}
$O(g(n))$ 是所有滿足下面條件的 $f(n)$：存在正的常數 $c$ 和 $n_0$，使得所有
$n \ge n_0$ 都有 $0 \le f(n) \le c\, g(n)$。

$\Omega(g(n))$ 是所有滿足下面條件的 $f(n)$：存在正的常數 $c$ 和 $n_0$，使得所有
$n \ge n_0$ 都有 $0 \le c\, g(n) \le f(n)$。
\end{definition}

直接從定義可以得到：

\begin{theorem}
\label{thm:theta-iff}
$f(n) = \Theta(g(n))$ 的充分必要條件是 $f(n) = O(g(n))$ 而且 $f(n) = \Omega(g(n))$。
\end{theorem}

\begin{proof}
如果 $f(n) = \Theta(g(n))$，定義裡的 $c_2$、$n_0$ 就是 $O$ 要的常數，$c_1$、$n_0$
就是 $\Omega$ 要的常數。 反過來，假設 $O$ 的常數是 $c$、$n_1$，$\Omega$ 的常數是
$c'$、$n_2$，取 $c_1 = c'$、$c_2 = c$、$n_0 = \max(n_1, n_2)$，$n \ge n_0$ 時兩個
不等式同時成立。
\end{proof}

\subsection{用定義證明}

\textbf{例 1}：證明 $4n^2 + 7n - 30 = \Theta(n^2)$。

上界：$n \ge 1$ 時 $7n \le 7n^2$，而 $-30 < 0$，所以
$4n^2 + 7n - 30 \le 4n^2 + 7n^2 = 11n^2$。 取 $c_2 = 11$。

下界：想要 $4n^2 + 7n - 30 \ge 3n^2$，也就是 $n^2 + 7n \ge 30$。 $n = 3$ 時左邊是
$30$，剛好成立，$n$ 更大時左邊更大，所以 $n \ge 3$ 時成立。 取 $c_1 = 3$。

兩邊都要成立，取 $n_0 = 3$。 另外 $n \ge 3$ 時 $3n^2 \ge 0$，定義裡的 $0 \le$ 也滿
足。 所以 $c_1 = 3$、$c_2 = 11$、$n_0 = 3$ 符合定義~\ref{def:theta}。

常數的選法不只一組。 下界也可以取 $c_1 = 1$，這時需要 $3n^2 + 7n \ge 30$，
$n = 2$ 時左邊只有 $26$，要到 $n \ge 3$ 才成立。 證明只要找到一組，並且說清楚每個不等式從哪一個 $n$ 開始成立。

\begin{examnote}
用定義證明 $\Theta$、$O$、$\Omega$ 時，答案裡要明白寫出 $c_1$、$c_2$（或 $c$）
和 $n_0$ 的數值，並驗證「所有 $n \ge n_0$」都成立，不能只代一兩個 $n$ 檢查。 常
用的手法是把負的項丟掉得到上界，把低次項併進最高次項的一部分得到下界。
\end{examnote}

\textbf{例 2}：證明 $n^3 \ne O(n^2)$。

用反證法。 假設有正的 $c$、$n_0$，使得所有 $n \ge n_0$ 都有 $n^3 \le c\,n^2$，兩
邊除以 $n^2$ 得到 $n \le c$。 但取 $n = \max(n_0, \lceil c \rceil + 1)$，這個 $n$
大於等於 $n_0$，卻比 $c$ 大，矛盾。

要證明「不是 $O$」，就要說明不管常數怎麼選都不行，所以通常用反證法：假設常數存
在，再找一個夠大的 $n$ 推出矛盾。

\textbf{例 3}：多項式只看最高次。 設
$p(n) = a_d n^d + a_{d-1} n^{d-1} + \cdots + a_0$，$a_d > 0$。 令
$A = |a_{d-1}| + \cdots + |a_0|$。 $n \ge 1$ 時每個 $|a_k| n^k \le |a_k| n^{d-1}$，所以
\[
a_d n^d - A n^{d-1} \le p(n) \le a_d n^d + A n^{d-1}.
\]
右邊不超過 $(a_d + A) n^d$；左邊在 $n \ge 2A / a_d$ 時至少是 $\frac{a_d}{2} n^d$。
所以 $p(n) = \Theta(n^d)$。 插入排序最壞情況的 $an^2 + bn + c$（$a > 0$）是
$\Theta(n^2)$，最好情況的 $an + b$ 是 $\Theta(n)$。

\subsection{$O$ 只是上界}

$O$ 只說「不會比這個大太多」，沒有說「差不多這麼大」。 $5n = O(n^2)$、
$5n = O(n^3)$ 都對，因為 $n \ge 5$ 時 $5n \le n^2 \le n^3$。 所以「插入排序的執行
時間是 $O(n^3)$」是一句對的話，只是不夠精確。

考卷上的選擇題常常利用這一點。 讀到 $O$ 時，先問自己：這句話說的是上界，題目給
的函數有沒有超過這個上界？ 不要把 $O$ 讀成「剛好是」。 想說「剛好是這一級」時用
$\Theta$。

\subsection{$o$ 和 $\omega$}

$O$ 允許 $f$ 和 $g$ 同一級，例如 $2n^2 = O(n^2)$。 想說「$f$ 嚴格比 $g$ 小一級」時，
用小寫的 $o$。

\begin{definition}[$o$ 和 $\omega$]
\label{def:little-o}
$o(g(n))$ 是所有滿足下面條件的 $f(n)$：\textbf{對每一個}正的常數 $c$，都存在
$n_0$，使得所有 $n \ge n_0$ 都有 $0 \le f(n) < c\, g(n)$。

$\omega(g(n))$ 是所有滿足下面條件的 $f(n)$：\textbf{對每一個}正的常數 $c$，都存在
$n_0$，使得所有 $n \ge n_0$ 都有 $0 \le c\, g(n) < f(n)$。
\end{definition}

和 $O$ 的差別在「存在一個 $c$」換成了「每一個 $c$」：$c$ 再小，$n$ 夠大以後
$f(n)$ 都比 $c\,g(n)$ 小。 例如 $2n = o(n^2)$：給定 $c$，取 $n_0 > 2/c$，$n \ge n_0$
時 $2n < c\,n^2$。 但 $2n^2 \ne o(n^2)$：取 $c = 1$，$2n^2 < n^2$ 永遠不成立。

用極限判斷常常比較快。

\begin{theorem}
\label{thm:limit}
設 $g(n) > 0$，而且極限 $L = \lim_{n \to \infty} f(n)/g(n)$ 存在（可以是 $\infty$）。
\begin{enumerate}
\item $L = 0$：$f(n) = o(g(n))$，因此 $f(n) = O(g(n))$，但 $f(n) \ne \Omega(g(n))$。
\item $0 < L < \infty$：$f(n) = \Theta(g(n))$。
\item $L = \infty$：$f(n) = \omega(g(n))$，因此 $f(n) = \Omega(g(n))$，但 $f(n) \ne O(g(n))$。
\end{enumerate}
\end{theorem}

第 2 點的理由：$n$ 夠大以後 $f(n)/g(n)$ 落在 $[L/2, 2L]$ 之間，取 $c_1 = L/2$、
$c_2 = 2L$。 第 1 點：極限是 $0$ 的意思是，對每一個 $c > 0$，$n$ 夠大以後
$f(n)/g(n) < c$，這正是 $o$ 的定義；取 $c = 1$ 就得到 $O$。 如果又有
$f(n) \ge c' g(n)$（$c' > 0$），比值會一直不小於 $c'$，不可能趨近 $0$，所以
$f(n) \ne \Omega(g(n))$。 第 3 點的理由相同，大小對調。

要注意這個定理只有一個方向：極限不存在時，什麼都不能說，要回到定義。 例如
\[
f(n) = \begin{cases} n^2, & n \text{ 是偶數} \\ n, & n \text{ 是奇數} \end{cases}
\]
和 $g(n) = n^2$ 相比，$f(n)/g(n)$ 在 $1$ 和 $1/n$ 之間跳，極限不存在。 但 $f(n) \le
n^2$ 對所有 $n$ 成立，所以 $f(n) = O(n^2)$；奇數的 $n$ 讓 $f(n) = n$，所以 $f(n) \ne
\Omega(n^2)$。

\subsection{性質}

下面幾條性質直接從定義推得，之後會常常用到。 $f$、$g$、$h$ 都是夠大以後非負的函數。

\begin{itemize}
\item \textbf{遞移}：$f = O(g)$ 而且 $g = O(h)$，則 $f = O(h)$。 $\Theta$、$\Omega$、
$o$、$\omega$ 也一樣。
\item \textbf{自反}：$f = \Theta(f)$，也就是 $f = O(f)$、$f = \Omega(f)$。
\item \textbf{對稱}：$f = \Theta(g)$ 的充分必要條件是 $g = \Theta(f)$。
\item \textbf{轉置}：$f = O(g)$ 的充分必要條件是 $g = \Omega(f)$；$f = o(g)$ 的充分
必要條件是 $g = \omega(f)$。
\item \textbf{加法}：$f + g = \Theta(\max(f, g))$。 例如兩段程式接在一起，總時間
由比較慢的那一段決定。
\end{itemize}

這些性質很像數字的大小關係（$O$ 像 $\le$，$\Theta$ 像 $=$，$o$ 像 $<$），但有一點
不一樣：任兩個數一定比得出大小，兩個函數卻不一定。 例如 $g(n) = n^2$，而
$f(n)$ 在偶數時是 $n$、奇數時是 $n^3$。 偶數的 $n$ 讓 $f$ 遠小於 $g$，奇數的 $n$
讓 $f$ 遠大於 $g$，所以 $f \ne O(g)$，$f \ne \Omega(g)$。

\subsection{等式裡的漸近符號}

漸近符號出現在式子的一部分時，代表「某一個屬於這個集合的函數」。 例如
\[
T(n) = 2n^2 + \Theta(n)
\]
是說 $T(n) = 2n^2 + h(n)$，其中 $h(n)$ 是某個 $\Theta(n)$ 的函數。 這種寫法可以保
留關心的部分（最高次的係數 $2$），把不關心的細節收起來。 第~2~章寫遞迴式
$T(n) = 2T(n/2) + \Theta(n)$ 也是這個意思。

式子兩邊都有漸近符號時，例如 $2n^2 + \Theta(n) = \Theta(n^2)$，意思是：左邊不管
$\Theta(n)$ 換成哪一個函數，右邊都找得到一個 $\Theta(n^2)$ 的函數讓等式成立。 這仍
然是單向的，右邊比左邊粗。

\section{常見函數的成長速度}
\label{sec:ana-functions}

分析演算法時，最常出現的函數是對數、多項式、指數和階乘。 這一節整理它們之間的
大小關係，最後給一張由慢到快的表。

\subsection{對數}

$\lg n = \log_2 n$。 換底公式 $\log_a n = \log_b n / \log_b a$ 說明不同底的對數只差
一個常數倍，所以在 $O$、$\Theta$ 裡面寫 $\log n$ 不用寫底：$\log_{10} n =
\Theta(\lg n) = \Theta(\ln n)$。

底數在指數上就不能省了。 $4^n = 2^{2n} = (2^n)^2$，比 $2^n$ 大了 $2^n$ 倍，所以
$4^n \ne O(2^n)$。 另外幾個常用的換算：
\[
2^{\lg n} = n, \qquad 4^{\lg n} = n^2, \qquad a^{\log_b n} = n^{\log_b a}.
\]
最後一個等式兩邊取 $\log_b$ 就看得出來，第~2~章的主定理會一直用到它。

\subsection{對數、多項式、指數}

三條常用的結論，$a$、$b$、$\varepsilon$ 是常數：

\begin{enumerate}
\item 任何次方的對數都比任何正次方的多項式慢：$(\lg n)^a = o(n^{\varepsilon})$，
$a > 0$，$\varepsilon > 0$。
\item 任何多項式都比任何底大於 $1$ 的指數慢：$n^a = o(b^n)$，$b > 1$。
\item 多項式之間只看次方：$a < b$ 時 $n^a = o(n^b)$。
\end{enumerate}

第 1 點可以用定理~\ref{thm:limit}：令 $n = 2^m$，$(\lg n)^a / n^{\varepsilon} =
m^a / 2^{\varepsilon m}$，這是第 2 點的情形（以 $m$ 為變數，底是 $2^{\varepsilon}$），
所以極限是 $0$。 第 2 點可以對 $n^a / b^n$ 用羅必達法則 $\lceil a \rceil$ 次。

這幾條結論的數量感要記住。 $(\lg n)^3$ 在 $n = 10^6$ 時大約是 $7900$，比
$\sqrt{n} = 1000$ 還大；可是 $n = 2^{100}$ 時，$(\lg n)^3 = 10^6$，$\sqrt n = 2^{50}
\approx 10^{15}$。 漸近的比較說的是 $n$ 夠大以後，「夠大」可能非常大。

\subsection{階乘}

$n!$ 比任何 $b^n$ 都快，又比 $n^n$ 慢。 比較精確的估計是\textbf{斯特林公式}
（Stirling's approximation）：
\[
n! = \sqrt{2\pi n}\left(\frac{n}{e}\right)^n \left(1 + \Theta\!\left(\frac{1}{n}\right)\right).
\]
排序的下界（第~15~章）要用到 $\lg(n!)$。 它其實不需要斯特林公式也能估：
\[
n! = 1 \cdot 2 \cdots n \le n^n,
\qquad
n! \ge \left\lceil \tfrac{n}{2} \right\rceil \cdot \left(\left\lceil \tfrac{n}{2} \right\rceil + 1\right) \cdots n
\ge \left(\frac{n}{2}\right)^{n/2}.
\]
第二個式子只留下後半段的乘數，每個都至少 $n/2$，個數至少 $n/2$。 兩邊取 $\lg$：
\[
\frac{n}{2} \lg \frac{n}{2} \le \lg(n!) \le n \lg n,
\]
左邊在 $n \ge 4$ 時至少是 $\frac{n}{4}\lg n$（因為 $\lg(n/2) = \lg n - 1 \ge
\frac{1}{2}\lg n$）。 所以 $\lg(n!) = \Theta(n \lg n)$。

\subsection{級數}

逐行分析時最常出現下面幾種和。

\begin{itemize}
\item \textbf{等差}：$\displaystyle\sum_{i=1}^{n} i = \frac{n(n+1)}{2} = \Theta(n^2)$。 一
般地，$\displaystyle\sum_{i=1}^{n} i^k = \Theta(n^{k+1})$，$k \ge 0$ 是常數。
\item \textbf{等比}：$r \ne 1$ 時
$\displaystyle\sum_{i=0}^{m} r^i = \frac{r^{m+1} - 1}{r - 1}$。 $0 < r < 1$ 時，不管
$m$ 多大，和都小於 $\frac{1}{1-r}$，是 $\Theta(1)$；$r > 1$ 時，和是
$\Theta(r^m)$，由最後一項決定。
\item \textbf{調和}：$H_n = \displaystyle\sum_{i=1}^{n} \frac{1}{i}$。 拿 $1/x$ 的
積分上下夾，得到 $\ln(n+1) \le H_n \le \ln n + 1$，所以 $H_n = \Theta(\lg n)$。
\end{itemize}

等比級數的兩個結論要特別記住：公比小於 $1$ 時，總和和第一項同一級；公比大於 $1$ 時，
總和和最後一項同一級。 下一節分析「每次減半」「每次加倍」的迴圈，都是用這兩句話。

\subsection{由慢到快}

圖~\ref{fig:ana-growth} 畫出 $n \le 20$ 時幾個常見函數的大小，下表再列出 $n$ 更大時
的值。 在一般的電腦上，一秒大約能做 $10^8$ 到 $10^9$ 個基本運算。 對照下表，
$n = 10^6$ 時 $n \lg n$ 的演算法一瞬間就跑完，$n^2$ 要十幾分鐘到幾個小時，指數的演
算法連 $n = 100$ 都跑不完。 演算法課關心的，主要就是能不能從表的下半部換到上半部。

\begin{center}
\begin{tikzpicture}[x=0.42cm, y=0.075cm]
\begin{scope}
\clip (0,0) rectangle (21,62);
\draw[growth, black!60, domain=1:20, samples=60] plot (\x, {ln(\x)/ln(2)});
\draw[growth, coreStroke, domain=1:20, samples=60] plot (\x, {\x});
\draw[growth, greenStroke, domain=1:20, samples=60] plot (\x, {\x*ln(\x)/ln(2)});
\draw[growth, memStroke, domain=1:20, samples=60] plot (\x, {\x*\x});
\draw[growth, optStroke, domain=1:6, samples=60] plot (\x, {pow(2,\x)});
\end{scope}
\draw[-{Stealth[length=2mm]}] (0,0) -- (21.5,0) node[right, font=\small] {$n$};
\draw[-{Stealth[length=2mm]}] (0,0) -- (0,64);
\foreach \t in {5,10,15,20} { \draw (\t,0) -- (\t,-1) node[below, font=\scriptsize] {\t}; }
\foreach \t in {20,40,60} { \draw (0,\t) -- (-0.3,\t) node[left, font=\scriptsize] {\t}; }
\node[glabel, text=black!60] at (20.1,4.3) {$\lg n$};
\node[glabel, text=coreStroke] at (20.1,20) {$n$};
\node[glabel, text=greenStroke] at (15.4,60) {$n \lg n$};
\node[glabel, text=memStroke] at (7.9,60) {$n^2$};
\node[glabel, text=optStroke] at (6.0,60) {$2^n$};
\end{tikzpicture}
\diagcaption{$n \le 20$ 時幾個常見函數的大小（縱軸只畫到 $60$）。 $2^n$ 在
$n = 6$ 左右就超出圖外，$n^2$ 在 $n = 8$ 左右，$n \lg n$ 在 $n = 15$ 左右。}
\label{fig:ana-growth}
\end{center}

\begin{center}
\begin{tabular}{lrrrl}
\toprule
函數 & $n = 10$ & $n = 10^3$ & $n = 10^6$ & 常見的來源 \\
\midrule
$1$ & 1 & 1 & 1 & 陣列取值、堆疊的推入 \\
$\lg n$ & 3.3 & 10.0 & 19.9 & 二分搜尋、平衡樹的一次操作 \\
$\sqrt{n}$ & 3.2 & 31.6 & 1000 & 試除法判斷質數 \\
$n$ & 10 & 1000 & $10^6$ & 掃過陣列一次 \\
$n \lg n$ & 33 & 9966 & $2.0 \times 10^7$ & 合併排序、堆積排序 \\
$n^2$ & 100 & $10^6$ & $10^{12}$ & 兩兩比較、插入排序最壞情況 \\
$n^3$ & 1000 & $10^9$ & $10^{18}$ & 三層迴圈、Floyd--Warshall \\
$2^n$ & 1024 & $1.1 \times 10^{301}$ & 約 $10^{301030}$ & 枚舉所有子集合 \\
$n!$ & $3.6 \times 10^6$ & 約 $10^{2568}$ & 約 $10^{5565709}$ & 枚舉所有排列 \\
\bottomrule
\end{tabular}
\end{center}

\section{分析程式片段}
\label{sec:ana-fragments}

考題和作業常常給一段迴圈，問某一行執行幾次，或問時間複雜度。 做法是從最內層往外
寫成和，再化簡。 下面五段迴圈各自示範一種情況，都數 \texttt{c++} 執行的次數。

\begin{code}{c}{五段迴圈}
// (a) 內層上界跟著外層變
for (int i = 0; i < n; i++)
    for (int j = 0; j < i; j++)
        c++;

// (b) 每次加倍
for (long long i = 1; i < n; i *= 2)
    c++;

// (c) 內層的步長跟著外層變
for (int i = 1; i <= n; i++)
    for (int j = i; j <= n; j += i)
        c++;

// (d) 內層跑到根號
for (int i = 1; i <= n; i++)
    for (long long j = 1; j * j <= i; j++)
        c++;

// (e) 外層每次減半，內層跑 i 次
for (int i = n; i > 0; i /= 2)
    for (int j = 0; j < i; j++)
        c++;
\end{code}

\textbf{(a)} 外層的 $i$ 固定時，內層跑 $i$ 次。 總次數
$\sum_{i=0}^{n-1} i = \frac{n(n-1)}{2} = \Theta(n^2)$。 這和第~\ref{sec:ana-cost}~節的兩
兩比較是同一個和。

\textbf{(b)} $i$ 依序是 $1, 2, 4, \ldots$，只要 $2^k < n$ 就會再跑一次。 執行次數是
滿足 $2^k < n$ 的 $k \ge 0$ 的個數，$n \ge 2$ 時等於 $\lceil \lg n \rceil$，是
$\Theta(\lg n)$。 每次乘以任何大於 $1$ 的常數（例如乘 $3$）都一樣是 $\Theta(\lg n)$，
只差對數的底。

\textbf{(c)} 外層是 $i$ 時，$j$ 是 $i, 2i, 3i, \ldots$，不超過 $n$ 的有 $\lfloor n/i
\rfloor$ 個。 總次數
\[
\sum_{i=1}^{n} \left\lfloor \frac{n}{i} \right\rfloor.
\]
因為 $\frac{n}{i} - 1 < \lfloor \frac{n}{i} \rfloor \le \frac{n}{i}$，總和夾在
$n H_n - n$ 和 $n H_n$ 之間，而 $H_n = \Theta(\lg n)$，所以是 $\Theta(n \lg n)$。 兩
層迴圈不一定是 $n^2$，要看內層實際跑幾次。

\textbf{(d)} 外層是 $i$ 時內層跑 $\lfloor \sqrt{i} \rfloor$ 次。 上界：每一項都不超過
$\sqrt{n}$，總和不超過 $n\sqrt{n}$。 下界：只看 $i \ge n/2$ 的那一半，每一項至少
$\lfloor \sqrt{n/2} \rfloor$，大約有 $n/2$ 項，總和至少是 $n^{3/2}$ 的常數倍。 所以
是 $\Theta(n^{3/2})$。 精確一點的估計是 $\frac{2}{3} n^{3/2}$，來自積分
$\int_0^n \sqrt{x}\,dx$。

「只看後半」是求下界的常用手法：前半的項丟掉，後半的每一項都用最小的那一項代替。
上一節估 $\lg(n!)$ 也是這樣做的。

\textbf{(e)} 外層的 $i$ 依序是 $n, \lfloor n/2 \rfloor, \lfloor n/4 \rfloor, \ldots,
1$，內層跑 $i$ 次，總次數不超過
\[
n + \frac{n}{2} + \frac{n}{4} + \cdots < 2n.
\]
又至少是 $n$（第一輪），所以是 $\Theta(n)$。 這是公比 $1/2$ 的等比級數，和第一項同
一級。 外層雖然跑了 $\Theta(\lg n)$ 輪，但不能直接寫成 $n \lg n$：每一輪的工作量在
減半。

\subsection{實際數一次}

把五段迴圈各包成一個函式，對幾個 $n$ 實際數出次數。 推出公式之後跑一次這種短程式，
很容易抓到推導裡的錯。

\begin{example}[label=ex:loop-counts]{c}{(loop\_counts.c) 五段迴圈的執行次數}
#include <stdio.h>

// 五段迴圈，各自回傳最內層執行的次數。
long long frag_a(int n)
{
    long long c = 0;
    for (int i = 0; i < n; i++)
        for (int j = 0; j < i; j++)
            c++;
    return c;
}

long long frag_b(int n)
{
    long long c = 0;
    for (long long i = 1; i < n; i *= 2)
        c++;
    return c;
}

long long frag_c(int n)
{
    long long c = 0;
    for (int i = 1; i <= n; i++)
        for (int j = i; j <= n; j += i)
            c++;
    return c;
}

long long frag_d(int n)
{
    long long c = 0;
    for (int i = 1; i <= n; i++)
        for (long long j = 1; j * j <= i; j++)
            c++;
    return c;
}

long long frag_e(int n)
{
    long long c = 0;
    for (int i = n; i > 0; i /= 2)
        for (int j = 0; j < i; j++)
            c++;
    return c;
}

int main(void)
{
    int ns[] = {16, 100, 1000, 10000};
    printf("n a b c d e\n");
    for (int k = 0; k < 4; k++) {
        int n = ns[k];
        printf("%d %lld %lld %lld %lld %lld\n", n, frag_a(n),
               frag_b(n), frag_c(n), frag_d(n), frag_e(n));
    }
    return 0;
}
\end{example}
\iosample{}{n a b c d e \\ 16 120 4 50 38 31 \\ 100 4950 7 482 625 197 \\ 1000 499500 10 7069 20615 1994 \\ 10000 49995000 14 93668 661750 19995}

和公式對照（$n = 10^4$）：(a) $n(n-1)/2 = 49{,}995{,}000$，完全一樣；(b)
$\lceil \lg 10^4 \rceil = 14$；(c) 落在 $nH_n - n \approx 87{,}876$ 和
$nH_n \approx 97{,}876$ 之間；(d) 接近 $\frac{2}{3} n^{3/2} \approx 666{,}667$；(e) 小於
$2n = 20{,}000$。 $n$ 從 $10^3$ 變成 $10^4$ 時，(a) 變成 $100$ 倍，(c) 大約 $13$
倍，(d) 大約 $32$ 倍（$10^{1.5} \approx 31.6$），(e) 大約 $10$ 倍，各自符合
$n^2$、$n \lg n$、$n^{1.5}$、$n$ 的成長方式。

\begin{examnote}
問「某一行執行幾次」時，先寫出和（例如 $\sum_{i=1}^{n} \lfloor n/i \rfloor$），再化
成封閉形式或給出上下界，最後才寫 $\Theta$。 寫完後代一個小的 $n$（例如 $n = 4$）
手數一次，和公式比對，可以抓到大部分差 $1$ 的錯誤，例如 $<$ 和 $\le$ 搞混。
\end{examnote}

\section{最好、最壞和平均的情況}
\label{sec:ana-cases}

插入排序的例子說明，同樣大小的輸入，步數可以差很多。 所以談執行時間時，要說清楚
是哪一種：

\begin{itemize}
\item \textbf{最壞情況}（worst case）：所有大小為 $n$ 的輸入裡，步數的最大值。 它
是一個保證：不管輸入是什麼，都不會比這個慢。 這本書沒有特別說明時，執行時間指的
都是最壞情況。
\item \textbf{最好情況}（best case）：步數的最小值。 用處不大，因為輸入通常不會剛
好是最好的那一種，但常出現在選擇題。
\item \textbf{平均情況}（average case）：先假設輸入怎麼分布，再算步數的期望值。
排序常用的假設是：$n$ 個不同的數，$n!$ 種排列出現的機會都一樣。
\end{itemize}

\subsection{插入排序的平均情況}

在「每種排列機會一樣」的假設下，插入排序平均挪動幾次？ 挪動次數等於逆序對的個數，
所以要算逆序對個數的期望值。

任取兩個位置 $i < j$。 把一個排列的第 $i$、$j$ 個數對調，得到另一個排列；這樣兩兩
配成一對，每一對裡剛好一個在 $(i, j)$ 是逆序。 所以 $(i, j)$ 是逆序的排列剛好佔一
半，機率是 $1/2$。 一共有 $\binom{n}{2}$ 對位置，期望值可以一對一對加起來（期望值
的線性），得到
\[
\text{平均挪動次數} = \binom{n}{2} \cdot \frac{1}{2} = \frac{n(n-1)}{4}.
\]
這剛好是最壞情況的一半，仍然是 $\Theta(n^2)$。

$n$ 很小時，可以把 $n!$ 種排列全部列出來，直接驗證這個公式。 下面的程式用遞迴產
生 $0, 1, \ldots, n-1$ 的所有排列（《C 程式設計基礎》第~15~章回溯法的寫法），對每
一個排列跑一次插入排序，加總挪動次數。

\begin{example}[label=ex:avg-shift]{c}{(avg\_shift.c) 所有排列的平均挪動次數}
#include <stdio.h>

#define MAXN 8

int perm[MAXN];
int used[MAXN];
long long total;        // 所有排列的挪動次數總和
long long count;        // 排列個數

// 插入排序挪動的次數，等於 a[0..n-1] 裡 i < j 且 a[i] > a[j] 的對數。
// 這裡直接把插入排序跑一次來數，不用這個性質。
long long shifts(const int *src, int n)
{
    int a[MAXN];
    long long s = 0;
    for (int i = 0; i < n; i++)
        a[i] = src[i];
    for (int i = 1; i < n; i++) {
        int key = a[i];
        int j = i - 1;
        while (j >= 0 && a[j] > key) {
            a[j + 1] = a[j];
            s++;
            j--;
        }
        a[j + 1] = key;
    }
    return s;
}

// 依序填 perm[pos..n-1]，填滿時算一次。
void gen(int pos, int n)
{
    if (pos == n) {
        total += shifts(perm, n);
        count++;
        return;
    }
    for (int v = 0; v < n; v++) {
        if (!used[v]) {
            used[v] = 1;
            perm[pos] = v;
            gen(pos + 1, n);
            used[v] = 0;
        }
    }
}

int main(void)
{
    printf("n n! total average n(n-1)/4\n");
    for (int n = 1; n <= MAXN; n++) {
        total = 0;
        count = 0;
        gen(0, n);
        printf("%d %lld %lld %.2f %.2f\n", n, count, total,
               (double)total / count, n * (n - 1) / 4.0);
    }
    return 0;
}
\end{example}
\iosample{}{n n! total average n(n-1)/4 \\ 1 1 0 0.00 0.00 \\ 2 2 1 0.50 0.50 \\ 3 6 9 1.50 1.50 \\ 4 24 72 3.00 3.00 \\ 5 120 600 5.00 5.00 \\ 6 720 5400 7.50 7.50 \\ 7 5040 52920 10.50 10.50 \\ 8 40320 564480 14.00 14.00}

八個 $n$ 都和公式一致。 這種「小範圍全部列舉」的程式證明不了公式，但公式推錯時
通常馬上會露出來。

\subsection{「插入排序是 $\Theta(n^2)$」對不對}

這句話常在考卷的是非題出現，答案要看它指的是什麼。

\begin{itemize}
\item 「插入排序最壞情況的執行時間是 $\Theta(n^2)$」：對。
\item 「插入排序的執行時間是 $O(n^2)$」：對。 每一種輸入都不超過最壞情況，所以
$O(n^2)$ 是對所有輸入都成立的上界。
\item 「插入排序的執行時間是 $\Omega(n)$」：對。 每一種輸入至少要把每個數看一次。
\item 「插入排序的執行時間是 $\Theta(n^2)$」，沒有說哪一種情況：不對。 排好的輸入
只要 $\Theta(n)$，所以不是每一種輸入都至少 $n^2$ 的常數倍。
\item 「插入排序的執行時間是 $\Omega(n^2)$」：同上，不對。
\end{itemize}

整理成一句：對所有輸入成立的上界，看最壞情況；對所有輸入成立的下界，看最好情況。
「演算法 X 是 $\Theta(f(n))$」只有在最好和最壞情況同一級時才能不加說明。

\section{輸入大小要用位元數來算的時候}
\label{sec:ana-input-size}

前面的例子輸入是 $n$ 個數，大小就是 $n$。 如果輸入只有一個整數 $x$，問 $x$ 是不是
質數呢？ 試除法從 $2$ 開始試，試到 $\sqrt{x}$ 為止，步數大約是 $\sqrt{x}$。 這看起
來比線性還快，但輸入的大小並不是 $x$。 把 $x$ 寫下來只需要 $b = \lfloor \lg x
\rfloor + 1$ 個位元（或 $d$ 個十進位位數），輸入大小應該是 $b$。 用 $b$ 來寫，
步數是 $\sqrt{x} \approx 2^{b/2}$，是 $b$ 的指數函數。

\begin{example}[label=ex:trial-div]{c}{(trial\_div.c) 試除法的步數和位數}
#include <stdio.h>

// 試除法判斷質數。*steps 記錄做了幾次除法（取餘數）。
int is_prime(long long x, long long *steps)
{
    *steps = 0;
    if (x < 2) {
        return 0;
    }
    for (long long d = 2; d * d <= x; d++) {
        (*steps)++;
        if (x % d == 0) {
            return 0;
        }
    }
    return 1;
}

int main(void)
{
    long long xs[] = {101, 10007, 1000003, 100000007, 10000000019LL,
                      1000000000039LL};
    printf("x digits prime steps\n");
    for (int k = 0; k < 6; k++) {
        long long steps;
        int p = is_prime(xs[k], &steps);
        int digits = 0;
        for (long long t = xs[k]; t > 0; t /= 10)
            digits++;
        printf("%lld %d %s %lld\n", xs[k], digits,
               p ? "yes" : "no", steps);
    }
    return 0;
}
\end{example}
\iosample{}{x digits prime steps \\ 101 3 yes 9 \\ 10007 5 yes 99 \\ 1000003 7 yes 999 \\ 100000007 9 yes 9999 \\ 10000000019 11 yes 99999 \\ 1000000000039 13 yes 999999}

六個數都是質數，所以每次都要試到最後。 位數每多 $2$，步數就變成 $10$ 倍：輸入只
長了一點，工作量卻乘上一個固定倍數，這是指數成長的樣子。 $30$ 位數的質數，試除法
要試大約 $10^{15}$ 次。

像這樣執行時間是輸入\textbf{數值}的多項式、卻是輸入\textbf{位元數}的指數的演算
法，叫做\textbf{擬多項式時間}（pseudo-polynomial time）演算法。 第~18~章的 0/1 背
包動態規劃要 $\Theta(nW)$ 的時間，$W$ 是背包容量的數值，也是擬多項式；第~28~章談
NP 完備時，會說明為什麼這個區別很重要。

RAM 模型假設一個字組放得下 $n$ 這個量級的數，所以排序 $n$ 個 \texttt{int} 時，比
較兩個數算一步是合理的。 但如果數字本身有幾千位（例如密碼學裡的大整數），加法和
乘法就不能算一步，要另外數位元運算的次數。

\section{本章常見錯誤}
\label{sec:ana-pitfalls}

\begin{center}
\begin{tabular}{>{\raggedright\arraybackslash}p{0.24\linewidth}>{\raggedright\arraybackslash}p{0.50\linewidth}p{0.16\linewidth}}
\toprule
錯誤 & 例子 & 見 \\
\midrule
把 $O$ 當成「剛好是」 & 以為 $n = O(n^2)$ 是錯的；以為「$O(n^2)$ 的演算法一定比
$O(n)$ 的慢」 & 第~\ref{sec:ana-notation}~節 \\
沒說是哪一種情況 & 「插入排序是 $\Theta(n^2)$」不加說明 &
第~\ref{sec:ana-cases}~節 \\
證明只代幾個 $n$ & 用定義證明 $\Theta$ 時沒有寫出 $c_1$、$c_2$、$n_0$，或只驗證
$n = 1, 2, 3$ & 第~\ref{sec:ana-notation}~節 \\
兩層迴圈就寫 $n^2$ & 內層步長是 $i$ 的迴圈（$n \lg n$），外層減半的迴圈（$n$）&
第~\ref{sec:ana-fragments}~節 \\
指數的底當成常數丟掉 & 以為 $4^n = O(2^n)$；以為 $2^{2\lg n} = \Theta(n)$ &
第~\ref{sec:ana-functions}~節 \\
極限不存在就下結論 & 用 $\lim f/g$ 判斷，卻沒檢查極限存在 &
第~\ref{sec:ana-notation}~節 \\
輸入大小量錯 & 說試除法是 $O(\sqrt{x})$ 所以是「多項式時間」 &
第~\ref{sec:ana-input-size}~節 \\
C 和虛擬碼的編號差 $1$ & 虛擬碼的 $j \ge 1$ 翻成 \texttt{j >= 1}，少處理了
\texttt{a[0]} & 第~\ref{sec:ana-conventions}~節 \\
\bottomrule
\end{tabular}
\end{center}

\section{觀念題}
\label{sec:ana-exercises}

題型和研究所考試、期中期末考的筆試相近，有選擇、計算、追蹤和證明。 解答在本章最後。

\begin{enumerate}
\item （選擇，可能不只一個答案）下列哪些敘述正確？
\begin{enumerate}
\item $100 n \lg n = O(n^2)$
\item $n^2 = O(100 n \lg n)$
\item $2^{n+3} = \Theta(2^n)$
\item $3^n = O(2^n)$
\item $\lg(n^5) = \Theta(\lg n)$
\end{enumerate}

\item （計算）求下面的程式片段中 \texttt{x++} 執行的確切次數，以 $n$ 表示，並寫出
$\Theta$。
\begin{code}{c}{觀念題 2}
for (int i = 1; i <= n; i++)
    for (int j = 1; j <= i; j++)
        for (int k = j; k <= i; k++)
            x++;
\end{code}

\item （計算）求下面的程式片段的時間複雜度，以 $\Theta$ 表示，並說明理由。
\begin{code}{c}{觀念題 3}
for (int i = n; i >= 1; i /= 3)
    for (long long j = 1; j <= (long long)i * i; j++)
        x++;
\end{code}

\item （計算）$n$ 是 $2$ 的次方。 求 \texttt{x++} 執行的確切次數。
\begin{code}{c}{觀念題 4}
for (int i = 1; i <= n; i *= 2)
    for (int j = 1; j <= i; j++)
        x++;
\end{code}

\item （排序）把下列函數依成長速度由慢到快排列。 同一級（彼此是 $\Theta$）的用
等號連起來。
\[
n^{1.5},\quad n^{\lg 3},\quad (\lg n)^{\lg n},\quad 2^{\lg n},\quad \lg(n!),\quad
4^{\lg n},\quad n^{1/\lg n},\quad \sqrt{n}\,\lg n,\quad n \lg n,\quad (1.5)^n
\]

\item （證明）用定義證明 $5n^2 - 12n + 4 = \Theta(n^2)$，寫出 $c_1$、$c_2$、$n_0$。

\item （證明或反例）$f$、$g$、$h$ 都是正的函數。 下列敘述對的請證明，錯的請給反例。
\begin{enumerate}
\item 若 $f(n) = O(g(n))$，則 $f(n)^2 = O(g(n)^2)$。
\item 若 $f(n) = O(h(n))$ 且 $g(n) = O(h(n))$，則 $f(n) \cdot g(n) = O(h(n)^2)$。
\item 若 $f(n) = O(g(n))$，則 $\lim_{n \to \infty} f(n)/g(n)$ 存在。
\item 若 $f(n) = O(g(n))$，則 $f(n) - g(n) = O(1)$。
\end{enumerate}

\item （證明）用定義~\ref{def:little-o} 或定理~\ref{thm:limit} 證明
$n \lg n = o(n^{1.1})$，並說明為什麼 $n^{1.1} \ne O(n \lg n)$。

\item \label{ana-q-trace}（追蹤）用插入排序（虛擬碼~\ref{alg:insertion}）排序
$\langle 50, 20, 40, 10, 70, 30, 60 \rangle$。 寫出外層迴圈每一輪之後的陣列，以及
$A[j] > \mathit{key}$ 的總比較次數和總挪動次數。

\item （計算）插入排序排一個長度 $n$、有 $k$ 個逆序對的陣列。 證明比較次數
（只數 $A[j] > \mathit{key}$）等於 $k + (n - 1) - z$，其中 $z$ 是「$\mathit{key}$
被放到第 $1$ 格」的輪數。 用第~\ref{ana-q-trace}~題的結果驗證這個公式。

\item \label{ana-q-lower-bound}（證明）下面的程序在已經由小到大排好的 $A[1..n]$ 裡，回傳第一個大於等於
$x$ 的位置；全部都小於 $x$ 時回傳 $n + 1$。 寫出 \textbf{while} 迴圈的不變量，用
它證明程序正確，並說明迴圈一定會停。
\begin{pseudo}[unbreakable]{第一個大於等於 $x$ 的位置}
\Function{Lower-Bound}{$A, n, x$}
    \State $\mathit{lo} \gets 1$
    \State $\mathit{hi} \gets n + 1$
    \While{$\mathit{lo} < \mathit{hi}$}
        \State $\mathit{mid} \gets \lfloor (\mathit{lo} + \mathit{hi}) / 2 \rfloor$
        \If{$A[\mathit{mid}] < x$}
            \State $\mathit{lo} \gets \mathit{mid} + 1$
        \Else
            \State $\mathit{hi} \gets \mathit{mid}$
        \EndIf
    \EndWhile
    \State \Return $\mathit{lo}$
\EndFunction
\end{pseudo}

\item （選擇）在 $n$ 個互不相同的數裡從左到右找 $x$，找到就停。 已知 $x$ 一定在
裡面，而且在每個位置的機率都是 $1/n$。 平均比較幾次？
(A) $n/2$ \quad (B) $(n+1)/2$ \quad (C) $n$ \quad (D) $\lg n$ \quad (E) $(n-1)/2$

\item （是非）下面的函式用迴圈算 $1 + 2 + \cdots + x$，迴圈跑 $x$ 次。 「因為執行
時間是 $\Theta(x)$，所以它是多項式時間的演算法。」 這句話對嗎？ 說明理由。
\begin{code}{c}{觀念題 13}
long long sum_to(long long x)
{
    long long s = 0;
    for (long long i = 1; i <= x; i++)
        s += i;
    return s;
}
\end{code}

\item （證明）證明 $f(n) + g(n) = \Theta(\max(f(n), g(n)))$ 不一定能推廣到
「$n$ 個函數」：給一個例子，$f_1, f_2, \ldots$ 都是 $\Theta(1)$，但
$\sum_{i=1}^{n} f_i(n)$ 不是 $\Theta(1)$。 再說明這和第~\ref{sec:ana-fragments}~節
(e) 不能寫成「$\Theta(\lg n)$ 輪，每輪 $O(n)$，所以 $\Theta(n \lg n)$」有什麼關係。
\end{enumerate}

\section{歷屆試題}
\label{sec:ana-pastexam}

以下是研究所入學考試的原題，題目照原卷（語言、記號、配分照錄）。 詳解在本章最後。

\begin{pastexam}[label={px:ana-1}]{臺大資工所 115 學年度 第 2 題}
(4 points) For positive functions $f(n)$ and $g(n)$ on $\mathbb{N}$. If $f(n)=O(g(n))$
and $f(n)=\Omega(1)$, how many of the following statements are true?
\begin{itemize}
\item $g(n)=\Omega(f(n))$
\item $f(n)\cdot\log(1+f(n))=O(g(n)\cdot\log(1+g(n)))$
\item $2^{f(n)}=O(2^{g(n)})$
\item $(f(n))^k=O((g(n))^k)$ for any $0<k<1$
\end{itemize}
(A) 0 \qquad (B) 1 \qquad (C) 2 \qquad (D) 3 \qquad (E) 4
\end{pastexam}

\begin{pastexam}[label={px:ana-2}]{臺科大資工所 113 學年度 第 I-1 題}
I.\quad (35\%) Multiple choice questions.

(5\%) Order the following functions by growth rate in increasing order.\\
$N$, $\sqrt{N}$, $N^{1.5}$, $N\log N$, $N^2$, $N\log(\log N)$, $N\log^2 N$,
$N\log(N^2)$, $2^N$, 2048\\
Which is the 5th function in the ordered sequence?\\
(a) $N$,\quad (b) $\sqrt{N}$,\quad (c) $N^{1.5}$,\quad (d) $N\log N$,\quad (e) $N^2$,\\
(f) $N\log(\log N)$,\quad (g) $N\log^2 N$,\quad (h) $N\log(N^2)$,\quad (i) $2^N$,\quad
(j) 2048
\end{pastexam}

\begin{pastexam}[label={px:ana-3}]{臺科大電子所甲組 110 學年度 第 4 題}
Suppose we have three functions:
\begin{minted}[fontsize=\small]{text}
foo(int N)
    do something in A steps
    for i equals 1 to N
        do something in B steps
        for j equals 1 to N
            do something in C steps

goo(int N)
    do something in A steps
    for i equals 1 to N
        do something in B step
        for j equals 1 to i
            do something in C steps

poo(int N)
    do something in A steps
    for i equals 1 to N
        do something in B steps
        for j equals 1 to 250
            do something in C steps
\end{minted}
\begin{enumerate}[label=(\alph*)]
\item (6\%) For each of the three methods above (foo(), goo(), and poo()), what is the
exact runtime for each method (in steps) in terms of A, B, C, and N? You should not count
loop related operations.
\item (4\%) For N = 100, which one is the fastest algorithm? Why?
\item (6\%) Assuming that A, B, and C are constants, what is the asymptotic runtime of
each method in terms of N?
\item (4\%) As N grows infinitely large, which one is the fastest algorithm?
\end{enumerate}
\end{pastexam}

\begin{pastexam}[label={px:ana-4}]{中山資工所 112 學年度 第 2 題}
The pseudo code for the insertion sort is given as follows.
\begin{minted}[fontsize=\small]{text}
Read input elements into a[ ]; //stored in a[1], a[2], ..., a[n]
for i = 2 to n do
{
   j = i;
   while (j >= 2) and  (a[j] < a[j-1]) do
   {
      Swap(a[j], a[j-1]); // exchange a[j] and a[j-1]
      j = j-1;
   }
}
\end{minted}
\begin{enumerate}[label=(\alph*)]
\item (5\%) Suppose that the input elements are 3, 9, 6, 5, 8, 2. How many times
(counts) are needed for the execution of \texttt{Swap( )}?
\item (5\%) Suppose that the input is a permutation of $\{1,2,3,\ldots,n\}$. What
permutation will make \texttt{Swap( )} be executed the least times? How many times?
\item (5\%) Suppose that the input is a permutation of $\{1,2,3,\ldots,n\}$. What
permutation will make \texttt{Swap( )} be executed the most times? How many times?
\end{enumerate}
\end{pastexam}

\begin{pastexam}[label={px:ana-5}]{政大資科所 110 學年度 第 1 題}
(10\%) Prove or disprove the following statements ($n$ is the size of the input array):
\begin{enumerate}[label=\alph*.]
\item (5\%) The time complexity of any sorting algorithm is $O(n^2)$.
\item (5\%) The best-case time complexity of insertion sort is $\theta(n\log n)$.
\end{enumerate}
\end{pastexam}

\section{實作練習}
\label{sec:ana-practice}

下面是 Aki-Judge 上和本章相關的題目。 這一章的重點是分析，能寫成程式題的不多，
其他練習請以觀念題為主。

\begin{center}
\begin{tabular}{llc}
\toprule
小節 & 題目 & 難度 \\
\midrule
分析程式片段 & \akiref{count-nested-loop-executions}{計算巢狀迴圈執行次數} & 暖身 \\
插入排序 & \akiref{insertion-sort-function}{插入排序函式} & 基本 \\
\bottomrule
\end{tabular}
\end{center}

「插入排序函式」只要交排序函式，不寫 \texttt{main}。 交之前先用
範例~\ref{ex:insertion-count} 的三組輸入自己測一次，特別是已經排好和完全相反的兩組。

\section{讀完你應該會}
\label{sec:ana-checklist}

\begin{itemize}
\item 看到一段虛擬碼或 C 程式，能逐行寫出執行次數，化成封閉形式或上下界，再寫成
$\Theta$。
\item 能用定義證明 $f(n) = \Theta(g(n))$，寫出 $c_1$、$c_2$、$n_0$；能用反證法說明
$f(n) \ne O(g(n))$；能用極限判斷，並知道極限不存在時要回到定義。
\item 能把對數、多項式、指數、階乘混在一起的一串函數依成長速度排好，包括
$a^{\log_b n} = n^{\log_b a}$ 和 $\lg(n!) = \Theta(n \lg n)$ 這類換算。
\item 能用迴圈不變量的三步（初始、維持、結束）證明一個迴圈的正確性。
\item 能分別說出插入排序最好、最壞、平均情況的次數，並判斷「某演算法是
$\Theta(f(n))$」這類敘述在哪一種意義下成立。
\item 輸入是一個整數時，能用位元數判斷一個演算法是不是多項式時間。
\end{itemize}

\textbf{延伸閱讀}（書末會整理完整的參考資料）

\begin{itemize}
\item Thomas H. Cormen 等，\textit{Introduction to Algorithms}，第 4 版，第 2、3
章與附錄 A：插入排序的逐行分析、迴圈不變量、漸近符號的完整定義；CLRS 的附錄 A 整理了求
和的技巧（本書的求和工具見附錄~A）。
\item Steven S. Skiena，\textit{The Algorithm Design Manual}，第 3 版，第 2 章：
用較少的符號講 RAM 模型和成長速度，例子偏向實務。
\item Antti Laaksonen，\textit{Competitive Programmer's Handbook}，第 2 章：幾頁
就講完常見的複雜度和「資料量對應可接受的複雜度」，適合快速複習。
\end{itemize}

\addcontentsline{toc}{section}{觀念題解答}
\section*{觀念題解答}

\begin{enumerate}
\item (a)、(c)、(e) 正確。 (a)：$\lg n \le n$，所以 $100n\lg n \le 100n^2$。
(b)：$n^2 / (100 n \lg n) = n / (100 \lg n) \to \infty$，不是 $O$。 (c)：
$2^{n+3} = 8 \cdot 2^n$，常數倍。 (d)：$3^n / 2^n = 1.5^n \to \infty$。 (e)：
$\lg(n^5) = 5\lg n$。

\item 固定 $i$、$j$ 時，最內層跑 $i - j + 1$ 次。 固定 $i$ 時，
$\sum_{j=1}^{i} (i - j + 1) = 1 + 2 + \cdots + i = \frac{i(i+1)}{2}$。 總次數
\[
\sum_{i=1}^{n} \frac{i(i+1)}{2}
= \frac{1}{2}\left(\frac{n(n+1)(2n+1)}{6} + \frac{n(n+1)}{2}\right)
= \frac{n(n+1)(n+2)}{6}.
\]
是 $\Theta(n^3)$。 驗算 $n = 2$：$i = 1$ 時 $1$ 次，$i = 2$ 時 $j = 1$ 有 $2$ 次、
$j = 2$ 有 $1$ 次，共 $4 = \frac{2 \cdot 3 \cdot 4}{6}$。

\item 外層的 $i$ 是 $n, \lfloor n/3 \rfloor, \lfloor n/9 \rfloor, \ldots$，內層跑
$i^2$ 次。 總和不超過 $n^2 + n^2/9 + n^2/81 + \cdots = n^2 \cdot \frac{1}{1 - 1/9} =
\frac{9}{8}n^2$，又至少是第一輪的 $n^2$，所以是 $\Theta(n^2)$。 公比 $1/9 < 1$，總
和和第一項同一級。

\item 設 $n = 2^m$。 $i = 1, 2, 4, \ldots, 2^m$，內層跑 $i$ 次，總次數
$1 + 2 + \cdots + 2^m = 2^{m+1} - 1 = 2n - 1$，是 $\Theta(n)$。

\item 先化簡：$2^{\lg n} = n$，$4^{\lg n} = n^2$，$n^{1/\lg n} = 2^{\lg n / \lg n} = 2$，
$(\lg n)^{\lg n} = 2^{\lg n \cdot \lg\lg n} = n^{\lg \lg n}$，$\lg(n!) = \Theta(n \lg n)$，
$n^{\lg 3} \approx n^{1.585}$。 由慢到快：
\[
n^{1/\lg n} < \sqrt{n}\,\lg n < 2^{\lg n} < \lg(n!) = n\lg n < n^{1.5} < n^{\lg 3}
< 4^{\lg n} < (\lg n)^{\lg n} < (1.5)^n.
\]
$(\lg n)^{\lg n} = n^{\lg\lg n}$ 的次方 $\lg \lg n$ 會無限變大，所以比任何多項式都快；
和 $(1.5)^n$ 比，取 $\lg$ 後是 $\lg n \cdot \lg \lg n$ 對 $n \lg 1.5$，後者大。

\item 上界：$n \ge 1$ 時 $-12n + 4 \le -8 < 0$，所以 $5n^2 - 12n + 4 \le 5n^2$，取
$c_2 = 5$。 下界：要 $5n^2 - 12n + 4 \ge 2n^2$，即 $3n^2 - 12n + 4 \ge 0$。 這個二次
式的較大根是 $\frac{12 + \sqrt{96}}{6} \approx 3.63$，所以 $n \ge 4$ 時成立（$n = 4$：
$48 - 48 + 4 = 4 \ge 0$）。 取 $c_1 = 2$、$c_2 = 5$、$n_0 = 4$，此時 $2n^2 \ge 0$ 也
成立。

\item
\begin{enumerate}
\item 對。 $n \ge n_0$ 時 $0 \le f(n) \le c\,g(n)$，兩邊平方得
$f(n)^2 \le c^2 g(n)^2$，取常數 $c^2$。
\item 對。 $f \le c_1 h$（$n \ge n_1$），$g \le c_2 h$（$n \ge n_2$），相乘得
$fg \le c_1 c_2 h^2$，$n \ge \max(n_1, n_2)$。
\item 錯。 $g(n) = n$，$f(n)$ 在偶數時是 $n$、奇數時是 $1$。 $f(n) \le g(n)$，所以
$f = O(g)$；但 $f/g$ 在 $1$ 和 $1/n$ 之間跳，沒有極限。
\item 錯。 $f(n) = 2n$，$g(n) = n$。 $f = O(g)$（取 $c = 2$），但 $f(n) - g(n) = n$
不是 $O(1)$。 $O$ 允許差一個常數倍，不是差一個常數。
\end{enumerate}

\item $\dfrac{n \lg n}{n^{1.1}} = \dfrac{\lg n}{n^{0.1}}$。 用羅必達法則（把 $\lg n$
寫成 $\ln n / \ln 2$）：$\dfrac{1/(n \ln 2)}{0.1\, n^{-0.9}} = \dfrac{10}{\ln 2 \cdot
n^{0.1}} \to 0$。 由定理~\ref{thm:limit}，$n \lg n = o(n^{1.1})$。 同一個定理的第 1 點
也說明 $n \lg n \ne \Omega(n^{1.1})$，由轉置性質，$n^{1.1} \ne O(n \lg n)$。 注意
$n = 2^{60}$ 時 $\lg n = 60$，$n^{0.1} = 64$，兩者才差不多大；$n^{1.1}$ 要在非常
大的 $n$ 才明顯超過 $n \lg n$。

\item 每一輪之後的陣列：
\begin{center}
\begin{tabular}{cccl}
\toprule
$i$ & 比較 & 挪動 & 這一輪之後 \\
\midrule
2 & 1 & 1 & 20 50 40 10 70 30 60 \\
3 & 2 & 1 & 20 40 50 10 70 30 60 \\
4 & 3 & 3 & 10 20 40 50 70 30 60 \\
5 & 1 & 0 & 10 20 40 50 70 30 60 \\
6 & 4 & 3 & 10 20 30 40 50 70 60 \\
7 & 2 & 1 & 10 20 30 40 50 60 70 \\
\midrule
合計 & 13 & 9 & \\
\bottomrule
\end{tabular}
\end{center}
驗證的方法：把範例~\ref{ex:insertion-count} 的輸入改成三行
\texttt{1}、\texttt{7}、\texttt{50 20 40 10 70 30 60}，程式印出
\texttt{comparisons = 13, shifts = 9}。

\item 第 $i$ 輪挪動 $s_i$ 次。 如果這一輪因為某個 $A[j] \le \mathit{key}$ 而停下，
比較了 $s_i + 1$ 次；如果因為 $j = 0$ 停下（$\mathit{key}$ 放到第 $1$ 格），最後那
次沒有比較，比較了 $s_i$ 次。 把 $n - 1$ 輪加起來：比較次數 $= \sum s_i + (n - 1) - z
= k + (n - 1) - z$，因為挪動總次數等於逆序對個數 $k$。 第~\ref{ana-q-trace}~題
$k = 9$、$n = 7$，
$\mathit{key}$ 被放到第 $1$ 格的是 $i = 2$（20）和 $i = 4$（10），$z = 2$，
$9 + 6 - 2 = 13$，和表上一樣。

\item 不變量：每次檢查 \textbf{while} 條件時，$1 \le \mathit{lo} \le \mathit{hi} \le
n + 1$，而且 $A[1..\mathit{lo}-1]$ 的每個數都小於 $x$，$A[\mathit{hi}..n]$ 的每個數都
大於等於 $x$。

初始：$\mathit{lo} = 1$、$\mathit{hi} = n + 1$，兩段都是空的，成立。

維持：$\mathit{lo} < \mathit{hi}$ 時 $\mathit{lo} \le \mathit{mid} < \mathit{hi}$，所以
$A[\mathit{mid}]$ 存在。 若 $A[\mathit{mid}] < x$，因為陣列排好了，
$A[\mathit{lo}..\mathit{mid}]$ 都小於 $x$，設 $\mathit{lo} \gets \mathit{mid} + 1$ 後第
一段仍然都小於 $x$，而且 $\mathit{lo} \le \mathit{hi}$。 否則 $A[\mathit{mid}] \ge x$，
$A[\mathit{mid}..n]$ 都大於等於 $x$，設 $\mathit{hi} \gets \mathit{mid}$ 後第二段仍然
成立，而且 $\mathit{lo} \le \mathit{mid} = \mathit{hi}$。

結束：迴圈停在 $\mathit{lo} = \mathit{hi}$。 這時 $A[1..\mathit{lo}-1]$ 都小於 $x$、
$A[\mathit{lo}..n]$ 都大於等於 $x$，所以 $\mathit{lo}$ 正是第一個大於等於 $x$ 的位置
（$\mathit{lo} = n + 1$ 表示全部都小於 $x$）。

會停：每一輪 $\mathit{hi} - \mathit{lo}$ 都嚴格變小（$\mathit{lo}$ 變成
$\mathit{mid} + 1 > \mathit{lo}$，或 $\mathit{hi}$ 變成 $\mathit{mid} < \mathit{hi}$），
而它是非負整數，所以有限輪後一定變成 $0$。 每一輪大約減半，共 $\Theta(\lg n)$ 輪。

\item (B)。 $x$ 在第 $k$ 個位置時比較 $k$ 次，平均
$\frac{1}{n}\sum_{k=1}^{n} k = \frac{n+1}{2}$。

\item 不對。 輸入是一個整數 $x$，大小是它的位元數 $b \approx \lg x$。 執行時間
$\Theta(x) = \Theta(2^b)$，是輸入大小的指數，這是擬多項式時間。 （用公式
$x(x+1)/2$ 可以在 $O(1)$ 次算術運算內算完，$x$ 很大時要算大整數的乘法，時間是
$b$ 的多項式。）

\item 取 $f_i(n) = 1$。 每一個都是 $\Theta(1)$，但 $\sum_{i=1}^{n} f_i(n) = n$ 不是
$\Theta(1)$。 「加法」性質只對固定個數的函數成立；項數跟著 $n$ 變多時，要真的把和
算出來。 (e) 的情況相反：輪數 $\Theta(\lg n)$，每一輪的上界 $O(n)$ 都對，所以
$O(n \lg n)$ 是對的上界，但每一輪的實際工作量在減半，把和算出來只有 $\Theta(n)$。
「輪數乘上每輪的上界」只能得到上界，要得到 $\Theta$ 必須算和。
\end{enumerate}

\addcontentsline{toc}{section}{歷屆試題詳解}
\section*{歷屆試題詳解}

\pastsol{px:ana-1}{臺大資工所 115 學年度 第 2 題}
答案是 (D)：第 1、2、4 句對，第 3 句錯。

先把已知條件照定義~\ref{def:big-o} 寫出來。 存在 $c > 0$、$n_1$，使 $n \ge n_1$ 時
$f(n) \le c\,g(n)$；存在 $\delta > 0$、$n_2$，使 $n \ge n_2$ 時 $f(n) \ge \delta$。 把
$c$ 調大不會讓不等式變錯，所以不妨設 $c \ge 1$。 以下都假設 $n \ge \max(n_1, n_2)$，
$\log$ 的底是任何大於 $1$ 的固定數。

\begin{itemize}
\item 第 1 句對。 這是第~\ref{sec:ana-notation}~節的轉置性質：$f(n) \le c\,g(n)$ 就是
$\frac{1}{c} f(n) \le g(n)$，取常數 $1/c$ 得到 $g(n) = \Omega(f(n))$。
\item 第 2 句對。 先估 $\log(1 + f)$。 因為 $c \ge 1$，
$1 + f \le 1 + c\,g \le c\,(1 + g)$，所以 $\log(1 + f) \le \log c + \log(1 + g)$。 要
把常數 $\log c$ 併進 $\log(1 + g)$，需要 $\log(1 + g)$ 不會趨近 $0$。 這裡用到
$f = \Omega(1)$：$g \ge f / c \ge \delta / c$，所以
$\log(1 + g) \ge \log(1 + \delta / c) = \lambda$，$\lambda$ 是一個正的常數。 於是
$\log c \le \frac{\log c}{\lambda} \log(1 + g)$，
\[
\log(1 + f) \le \Bigl(1 + \frac{\log c}{\lambda}\Bigr) \log(1 + g).
\]
再和 $f \le c\,g$ 相乘（兩邊都是正的），得到
$f \log(1 + f) \le c\,\bigl(1 + \frac{\log c}{\lambda}\bigr)\, g \log(1 + g)$，括號外
是常數。
\item 第 3 句錯。 反例：$f(n) = 2n$、$g(n) = n$。 $f = O(g)$（取 $c = 2$），
$f = \Omega(1)$，但 $2^{f(n)} / 2^{g(n)} = 2^{n} \to \infty$，由定理~\ref{thm:limit}，
$2^{2n} \ne O(2^n)$。 第~\ref{sec:ana-functions}~節說過，指數上的常數倍不能丟。
\item 第 4 句對。 $k > 0$ 時 $x^k$ 在正數上是遞增函數，$f \le c\,g$ 兩邊取 $k$ 次方得
$f^k \le c^k g^k$，取常數 $c^k$。 這一步其實對任何 $k > 0$ 都成立。
\end{itemize}

\pastsol{px:ana-2}{臺科大資工所 113 學年度 第 I-1 題}
答案是 (d) $N \log N$。

先化簡：$N \log(N^2) = 2N \log N$，$N \log^2 N = N (\log N)^2$。 $\log$ 的底不影響比
較（第~\ref{sec:ana-functions}~節）。 由小到大是
\[
2048 < \sqrt{N} < N < N \log(\log N) < N \log N < N \log(N^2) < N \log^2 N < N^{1.5}
< N^2 < 2^N .
\]
相鄰兩個逐一說明（除了第 5、6 個，每一對的比值極限都是 $0$，用定理~\ref{thm:limit}）：
\begin{itemize}
\item $2048$ 是常數，$\sqrt{N} \to \infty$；$\sqrt{N} / N = 1/\sqrt{N} \to 0$。
\item $N / (N \log\log N) = 1 / \log\log N \to 0$。
\item $N \log\log N / (N \log N) = \log\log N / \log N$：令 $m = \log N$，這是
$\log m / m \to 0$。
\item $N \log N$ 和 $N \log(N^2) = 2N \log N$ 只差 $2$ 倍，是同一個 $\Theta$；但
$N \ge 2$ 時 $N \log N$ 嚴格比較小。
\item $2N \log N / (N \log^2 N) = 2 / \log N \to 0$。
\item $N \log^2 N / N^{1.5} = (\log N)^2 / N^{0.5} \to 0$（第~\ref{sec:ana-functions}~節
第 1 點：任何次方的對數都比任何正次方的多項式慢）。
\item $N^{1.5} / N^2 \to 0$；$N^2 / 2^N \to 0$（第 2 點）。
\end{itemize}
第 5 個是 $N \log N$。 題目要的是「排好的序列裡的第 5 個」，第 5、6 個同一級時，照
實際大小排，$N \log N$ 在前。 （代 $N = 2^{64}$ 實際算出十個值再排序，順序和上面相
同。）

\pastsol{px:ana-3}{臺科大電子所甲組 110 學年度 第 4 題}
題目說不計迴圈本身的操作，所以只數三種 do something 各做幾次，再乘上各自的步數，
和第~\ref{sec:ana-line-cost}~節的逐行計算一樣。 原題 goo 第 4 行的 ``B step'' 應該是
``B steps''，不影響答案。

(a)
\begin{itemize}
\item foo：A 做 $1$ 次；外層 $i$ 從 $1$ 到 $N$，B 做 $N$ 次；內層每輪跑 $N$ 次，C 做
$N \cdot N$ 次。 $T_{\mathrm{foo}}(N) = A + NB + N^2 C$。
\item goo：內層 $j$ 從 $1$ 到 $i$，跑 $i$ 次，C 共做
$\sum_{i=1}^{N} i = N(N+1)/2$ 次。 $T_{\mathrm{goo}}(N) = A + NB + N(N+1)C/2$。
\item poo：內層固定跑 $250$ 次，C 共做 $250N$ 次。
$T_{\mathrm{poo}}(N) = A + NB + 250NC$。
\end{itemize}

(b) 代 $N = 100$：
\[
T_{\mathrm{foo}} = A + 100B + 10000C,\quad
T_{\mathrm{goo}} = A + 100B + 5050C,\quad
T_{\mathrm{poo}} = A + 100B + 25000C .
\]
三者的 A、B 部分相同，只比 C 的次數：$5050 < 10000 < 25000$，所以 goo 最快（假設
$C > 0$；$C = 0$ 時三者一樣）。

(c) A、B、C 是正的常數。 foo 的最高次項是 $CN^2$，是 $\Theta(N^2)$（第~\ref{sec:ana-notation}~節
例 3 的多項式只看最高次）。 goo 是 $\frac{C}{2}N^2 + (B + \frac{C}{2})N + A$，也是
$\Theta(N^2)$。 poo 是 $(B + 250C)N + A$，是 $\Theta(N)$。

(d) poo。 它是 $\Theta(N)$，另外兩個是 $\Theta(N^2)$，$N$ 夠大以後一定是 poo 最快。
交叉點也算得出來：poo 的 C 比 foo 少，等於 $250N < N^2$，也就是 $N > 250$；比 goo
少，等於 $250N < N(N+1)/2$，也就是 $N > 499$。 所以 $N \ge 500$ 時 poo 最快，
$N = 499$ 時 poo 和 goo 的 C 都做 $124750$ 次。 (b) 和 (d) 的答案不同，因為漸近的
比較只管 $N$ 夠大的情形，$N = 100$ 時常數 $250$ 還佔上風。 （用小程式實際數過
$N = 100$、$499$、$500$ 時三段程式裡 C 執行的次數，和公式一致。）

\pastsol{px:ana-4}{中山資工所 112 學年度 第 2 題}
題目的版本用相鄰交換：每一輪把 $a[j]$ 和左邊比它大的鄰居交換，一路換到左邊不比它
大為止。 一次交換讓兩個相鄰而且逆序的數變成順序對，其他數對的前後關係不變，所以剛好
消掉一個逆序對。 這和第~\ref{sec:ana-insertion}~節「挪動次數等於逆序對個數」是同一
個論證：\texttt{Swap} 執行的次數等於輸入的逆序對個數。

(a) $9$ 次。 每一輪之後的陣列：
\begin{center}
\begin{tabular}{cccl}
\toprule
$i$ & 這一輪的 $a[i]$ & 交換次數 & 這一輪之後 \\
\midrule
2 & 9 & 0 & 3 9 6 5 8 2 \\
3 & 6 & 1 & 3 6 9 5 8 2 \\
4 & 5 & 2 & 3 5 6 9 8 2 \\
5 & 8 & 1 & 3 5 6 8 9 2 \\
6 & 2 & 5 & 2 3 5 6 8 9 \\
\midrule
合計 & & 9 & \\
\bottomrule
\end{tabular}
\end{center}
用逆序對驗算：$(3,2)$、$(9,6)$、$(9,5)$、$(9,8)$、$(9,2)$、$(6,5)$、$(6,2)$、
$(5,2)$、$(8,2)$，共 $9$ 個。

(b) $1, 2, \ldots, n$（已經排好），$0$ 次。 交換次數等於逆序對個數，不會是負的；等於
$0$ 表示沒有任何逆序對，也就是已經由小到大排好，而 $\{1, \ldots, n\}$ 排好的排列只有
這一個。

(c) $n, n-1, \ldots, 1$（完全相反），$n(n-1)/2$ 次。 $n$ 個位置共有 $\binom{n}{2}$
對，每一對最多貢獻一個逆序對，所以交換次數最多 $n(n-1)/2$；要每一對都是逆序，排列
只能是完全相反的那一個。 直接數也一樣：第 $i$ 輪的 $a[i]$ 比左邊每個數都小，要換
$i - 1$ 次，$\sum_{i=2}^{n} (i - 1) = n(n-1)/2$。 （用程式對 $n \le 7$ 的所有排列檢
查過，最少和最多都只有這兩個排列達到。）

\pastsol{px:ana-5}{政大資科所 110 學年度 第 1 題}
兩句都是錯的。 原題 b 的小寫 $\theta$ 指的就是 $\Theta$。

a.\ 錯。 這句話說的是「每一個排序演算法」，要否定它，舉出一個不是 $O(n^2)$ 的排序
演算法就夠了。 例子：先用遞迴產生 $n$ 個位置的全部 $n!$ 種排列（範例~\ref{ex:avg-shift}
的 \texttt{gen} 就是這樣產生排列），每一種都花 $n - 1$ 次比較檢查照這個順序讀出來是
不是由小到大，最後輸出排好的那一種。 排好的順序一定在這 $n!$ 種裡，所以輸出正確；
不管輸入是什麼，比較次數都是 $n! \cdot (n - 1)$。 而
\[
\frac{n! \cdot (n-1)}{n^2} = \frac{(n-1)! \cdot (n-1)}{n} \ge \frac{(n-1)!}{2} \to \infty
\qquad (n \ge 2),
\]
由定理~\ref{thm:limit}，$n!\,(n-1) \ne O(n^2)$。 這個演算法很笨，但它是一個正確的排
序演算法，所以原句不成立。 如果把原句讀成「排序這個問題可以在 $O(n^2)$ 時間內解決」，
那是對的（插入排序最壞是 $\Theta(n^2)$，第~\ref{sec:ana-cases}~節）；作答時把這個差
別寫出來。

b.\ 錯。 插入排序的最好情況是 $\Theta(n)$。 輸入已經排好時，每一輪第一次比較就停，
共 $n - 1$ 次比較、$0$ 次挪動，第~\ref{sec:ana-line-cost}~節算出執行時間是 $an + b$
的形式，$a > 0$；任何輸入都至少要把 $n$ 個數各看一次，最好情況也不會更少。 而
$n / (n \lg n) = 1 / \lg n \to 0$，由定理~\ref{thm:limit}，$n \ne \Omega(n \lg n)$，
所以最好情況的時間不是 $\Theta(n \log n)$。

\chapter{遞迴式}
\label{ch:recurrences}

% 本章的圖共用的樣式
\tikzset{
    rnode/.style={draw=black!55, rounded corners=2pt, inner sep=2.5pt,
                  font=\small, minimum width=0.95cm, align=center},
    rleaf/.style={draw=black!35, rounded corners=2pt, inner sep=2pt,
                  font=\scriptsize, minimum width=0.8cm, align=center},
    rsum/.style={font=\small, anchor=west},
    redge/.style={draw=black!60},
}

會計室要核對一份很多天的營業額報表。 主任的做法是：報表只有一天，就直接把那天
的數字當成總和；超過一天，就把報表從中間分成前後兩份，交給兩位組員各自算出總
和，組員也用同樣的方法再往下分。 兩份總和交回來以後，主任不放心，自己把整份報
表從頭到尾再加一次，確認兩份加起來沒有錯。 每一層的人都這樣做。 報表從
$1000$ 天變成 $8000$ 天，全部的人加起來，要多做幾倍的加法？ 往下讀之前，先估計
一下。

第~1~章分析的是迴圈，執行次數可以直接寫成一個和。 遞迴的程式不一樣：規模
$n$ 的工作量，是用「規模比較小的同一件事」來描述的。 寫出來是一個
\textbf{遞迴式}（recurrence），例如
\[
T(n) = 2T(n/2) + n .
\]
這一章先講怎麼從程式寫出遞迴式，再介紹幾種解法：代換法、遞迴樹、主定理，以及處
理 $T(n-1)$ 這類形式的展開和特徵方程式，最後簡單介紹生成函數和 Akra--Bazzi
方法。 之後的合併排序、快速排序（第~14~章）、分治（第~17~章）和很多樹的操作，
都要先寫出遞迴式才能分析。

這一章用到《C 程式設計基礎》第~15~章的遞迴，用到時會先在框裡複習。 數學上用到
第~1~章的漸近符號、等比級數和 $a^{\log_b n} = n^{\log_b a}$，以及數學歸納法（附錄~A
第~\ref{sec:mth-induction}~節整理了幾種寫法）。

\section{從程式寫出遞迴式}
\label{sec:rec-write}

遞迴式是一個等式（或不等式），把函數在 $n$ 的值用它在比較小的輸入上的值表示出
來，再加上最小的幾個 $n$ 的值，叫做\textbf{基礎情況}（base case）。 例如
\[
T(n) = \begin{cases} 0, & n = 1, \\ 2T(n/2) + n, & n \ge 2. \end{cases}
\]
光有遞迴式，還看不出 $T(n)$ 有多大。 \textbf{解}遞迴式的意思是找出 $T(n)$ 的封閉
形式（例如 $n \lg n$），或至少找出它的 $\Theta$。

\subsection{寫遞迴式要回答的四件事}

\begin{enumerate}
\item \textbf{規模}：用哪個量當 $n$。 處理陣列的一段 $A[p..r]$ 時，通常取
$n = r - p + 1$。
\item \textbf{基礎情況}：$n$ 多小時直接回答、不再呼叫，要花多少時間。 通常是
$\Theta(1)$。
\item \textbf{遞迴呼叫}：呼叫了幾次，每一次的規模是多少。 每一次呼叫寫一項
$T(\text{規模})$；同樣的規模呼叫兩次，就寫 $2T(\cdot)$。
\item \textbf{呼叫以外的工作}：切分、迴圈、合併答案，記成 $f(n)$。
\end{enumerate}

四件事合起來是 $T(n) = (\text{各次呼叫的 } T \text{ 加起來}) + f(n)$。 最容易漏掉的
是第 3 點的「次數」：同一個規模呼叫了兩次，卻只寫一個 $T$。

\begin{cbookrecap}[第~15~章 15.2、15.4、15.6 節]
遞迴函式要有基礎情況，而且每次呼叫都要往基礎情況靠近，否則不會停。 每一次呼叫
都有自己的一份參數和區域變數，放在呼叫堆疊上，所以下面的 \texttt{audit\_sum} 裡
的 \texttt{p}、\texttt{q}、\texttt{r}、\texttt{s} 在每一層都是不同的變數，不會互
相蓋掉。

處理陣列的一段時，用兩個整數 \texttt{p}、\texttt{r} 表示「從 \texttt{a[p]} 到
\texttt{a[r]}」，取中間 \texttt{q = p + (r - p) / 2} 分成兩半，這樣寫不會像
\texttt{(p + r) / 2} 那樣在 \texttt{p + r} 超過 \texttt{int} 時溢位。 想數函式被
呼叫幾次，最簡單的方法是用全域變數當計數器，每次進入函式就加一。
\end{cbookrecap}

\subsection{例子：遞迴版的二分搜尋}

在排好的 $A[p..r]$ 裡找 $x$，找到就回傳位置，找不到回傳 \NIL。

\begin{pseudo}[label=alg:rec-search]{遞迴版的二分搜尋}
\Function{Rec-Search}{$A, p, r, x$}
    \If{$p > r$}
        \State \Return \NIL \Comment{空的一段}
    \EndIf
    \State $q \gets \lfloor (p + r) / 2 \rfloor$
    \If{$A[q] = x$}
        \State \Return $q$
    \ElsIf{$A[q] < x$}
        \State \Return \Call{Rec-Search}{$A, q + 1, r, x$}
    \Else
        \State \Return \Call{Rec-Search}{$A, p, q - 1, x$}
    \EndIf
\EndFunction
\end{pseudo}

令 $n = r - p + 1$。 左邊那段 $A[p..q-1]$ 有 $q - p = \lceil n/2 \rceil - 1$ 個數，
右邊那段 $A[q+1..r]$ 有 $r - q = \lfloor n/2 \rfloor$ 個數，比較大的是
$\lfloor n/2 \rfloor$。 每一次呼叫只會往下呼叫\textbf{一次}，呼叫以外只做常數個
步驟。 分析最壞情況時，假設每一步都沒找到，而且都走到比較大的那一段，得到
\[
W(n) = W(\lfloor n/2 \rfloor) + c \quad (n \ge 1), \qquad W(0) = c,
\]
$c$ 是一次呼叫除了遞迴以外的成本。 觀念題第 12 題會證明這時總共呼叫
$\lfloor \lg n \rfloor + 2$ 次，所以 $W(n) = \Theta(\lg n)$。

\subsection{例子：層層核對的報表加總}

章首的做法寫成虛擬碼如下。 \proc{Audit-Sum} 回傳 $A[p..r]$ 的總和；兩半交給下一
層算，回來之後把整段重新加一次，核對兩半的小計。

\begin{pseudo}[label=alg:rec-audit]{層層核對的加總}
\Function{Audit-Sum}{$A, p, r$}
    \If{$p = r$}
        \State \Return $A[p]$
    \EndIf
    \State $q \gets \lfloor (p + r) / 2 \rfloor$
    \State $s_1 \gets$ \Call{Audit-Sum}{$A, p, q$}
    \State $s_2 \gets$ \Call{Audit-Sum}{$A, q + 1, r$}
    \State $s \gets 0$
    \For{$i \gets p$ \To $r$}
        \State $s \gets s + A[i]$ \Comment{整段重算一次}
    \EndFor
    \If{$s \ne s_1 + s_2$}
        \State 回報「$A[p..r]$ 的小計不符」
    \EndIf
    \State \Return $s$
\EndFunction
\end{pseudo}

照四件事來寫。 規模是 $n = r - p + 1$。 基礎情況是 $n = 1$，不做加法。 遞迴呼叫兩
次：左半 $A[p..q]$ 有 $q - p + 1 = \lceil n/2 \rceil$ 天，右半有
$\lfloor n/2 \rfloor$ 天。 呼叫以外，最後的 \textbf{for} 迴圈做 $n$ 次加法。 所以加法的次
數 $T(n)$ 滿足
\[
T(n) = \begin{cases}
0, & n = 1, \\
T(\lceil n/2 \rceil) + T(\lfloor n/2 \rfloor) + n, & n \ge 2.
\end{cases}
\]
被呼叫的次數 $C(n)$ 也可以用同樣的方法寫：每次呼叫自己算一次，再加上兩邊的呼叫，
\[
C(1) = 1, \qquad C(n) = C(\lceil n/2 \rceil) + C(\lfloor n/2 \rfloor) + 1 \quad (n \ge 2).
\]
整個程序的執行時間是 $\Theta(T(n) + C(n))$：每一次加法花常數時間，每一次呼叫除了
迴圈以外也只花常數時間。

\subsection{手算追蹤}

$n = 5$ 時把每一次呼叫畫成一個節點，節點裡寫它處理的那一段，下面寫它自己做了幾
次加法（不含下一層的）。 這種圖叫\textbf{遞迴樹}（recursion tree），第~\ref{sec:rec-tree}~節
會把它當成解遞迴式的工具。

\begin{center}
\begin{tikzpicture}[level distance=1.45cm]
\node[rnode] (a) at (0,0) {$A[1..5]$\\[-1pt]{\scriptsize 加 5 次}};
\node[rnode] (b) at (-3,-1.6) {$A[1..3]$\\[-1pt]{\scriptsize 加 3 次}};
\node[rnode] (c) at (3,-1.6) {$A[4..5]$\\[-1pt]{\scriptsize 加 2 次}};
\node[rnode] (d) at (-4.4,-3.2) {$A[1..2]$\\[-1pt]{\scriptsize 加 2 次}};
\node[rnode] (e) at (-1.6,-3.2) {$A[3..3]$\\[-1pt]{\scriptsize 0}};
\node[rnode] (f) at (1.8,-3.2) {$A[4..4]$\\[-1pt]{\scriptsize 0}};
\node[rnode] (g) at (4.2,-3.2) {$A[5..5]$\\[-1pt]{\scriptsize 0}};
\node[rnode] (h) at (-5.4,-4.8) {$A[1..1]$\\[-1pt]{\scriptsize 0}};
\node[rnode] (i) at (-3.4,-4.8) {$A[2..2]$\\[-1pt]{\scriptsize 0}};
\foreach \u/\v in {a/b, a/c, b/d, b/e, c/f, c/g, d/h, d/i} { \draw[redge] (\u) -- (\v); }
\node[rsum] at (6.0,0) {第 0 層：5};
\node[rsum] at (6.0,-1.6) {第 1 層：$3 + 2 = 5$};
\node[rsum] at (6.0,-3.2) {第 2 層：2};
\node[rsum] at (6.0,-4.8) {第 3 層：0};
\end{tikzpicture}
\diagcaption{\proc{Audit-Sum} 處理 5 天的遞迴樹。 共 9 次呼叫、$5 + 5 + 2 = 12$ 次
加法。 第 0、1 層各加 5 次；第 2 層只有 $A[1..2]$ 要加，其他都是基礎情況。}
\label{fig:rec-audit5}
\end{center}

考卷上寫追蹤時，也可以把遞迴式一個一個代進去，由小到大列表：
\begin{center}
\begin{tabular}{cccl}
\toprule
$n$ & $\lceil n/2 \rceil$、$\lfloor n/2 \rfloor$ & $T(n)$ & 計算 \\
\midrule
1 & -- & 0 & 基礎情況 \\
2 & 1、1 & 2 & $0 + 0 + 2$ \\
3 & 2、1 & 5 & $T(2) + T(1) + 3 = 2 + 0 + 3$ \\
4 & 2、2 & 8 & $2 + 2 + 4$ \\
5 & 3、2 & 12 & $T(3) + T(2) + 5 = 5 + 2 + 5$ \\
\bottomrule
\end{tabular}
\end{center}

\subsection{C 程式}

下面的程式把 \proc{Audit-Sum} 翻成 C，用兩個全域變數數呼叫次數和加法次數，再印出
$n \lfloor \lg n \rfloor$ 和 $n \lceil \lg n \rceil$ 來比較（第~\ref{sec:rec-subst}~節
會證明 $T(n)$ 一定夾在這兩個數中間）。

\begin{example}[label=ex:rec-audit]{c}{(audit\_sum.c) 層層核對的加總，數呼叫和加法的次數}
#include <stdio.h>
#include <stdlib.h>

long long calls;    // audit_sum 被呼叫的次數
long long adds;     // 核對時重算總和做的加法次數

// 回傳 a[p..r] 的總和。兩半各自交給下一層算，回來後把整段重算一次，
// 核對兩半的小計加起來對不對。
long long audit_sum(const int *a, int p, int r)
{
    calls++;
    if (p == r) {
        return a[p];
    }
    int q = p + (r - p) / 2;        // 左半 a[p..q] 有 ceil(n/2) 天
    long long s1 = audit_sum(a, p, q);
    long long s2 = audit_sum(a, q + 1, r);
    long long s = 0;
    for (int i = p; i <= r; i++) {  // 重算整段：n 次加法
        s += a[i];
        adds++;
    }
    if (s != s1 + s2) {
        printf("mismatch in [%d, %d]\n", p, r);
    }
    return s;
}

// floor(lg n) 和 ceil(lg n)，n >= 1
int floor_lg(int n)
{
    int k = 0;
    while ((2LL << k) <= n) {
        k++;
    }
    return k;
}

int ceil_lg(int n)
{
    int k = 0;
    while ((1LL << k) < n) {
        k++;
    }
    return k;
}

int main(void)
{
    int ns[] = {1, 2, 3, 4, 5, 6, 7, 8, 1000, 2000, 4000, 8000};
    printf("n calls adds n*floor(lg n) n*ceil(lg n)\n");
    for (int t = 0; t < 12; t++) {
        int n = ns[t];
        int *a = malloc(n * sizeof *a);
        if (a == NULL) {
            return 1;
        }
        for (int i = 0; i < n; i++) {
            a[i] = 500 + (i * 37) % 101;  // 第 i 天的營業額（自己產生）
        }
        calls = 0;
        adds = 0;
        audit_sum(a, 0, n - 1);
        printf("%d %lld %lld %lld %lld\n", n, calls, adds,
               (long long)n * floor_lg(n), (long long)n * ceil_lg(n));
        free(a);
    }
    return 0;
}
\end{example}
\iosample{}{n calls adds n*floor(lg n) n*ceil(lg n) \\ 1 1 0 0 0 \\ 2 3 2 2 2 \\ 3 5 5 3 6 \\ 4 7 8 8 8 \\ 5 9 12 10 15 \\ 6 11 16 12 18 \\ 7 13 20 14 21 \\ 8 15 24 24 24 \\ 1000 1999 9976 9000 10000 \\ 2000 3999 21952 20000 22000 \\ 4000 7999 47904 44000 48000 \\ 8000 15999 103808 96000 104000}

幾件事可以從表上讀出來。 呼叫次數每一列都是 $2n - 1$。 加法次數在 $n$ 是 $2$ 的
次方時剛好等於 $n \lg n$（$n = 2, 4, 8$），其他時候夾在兩欄中間。 章首的問題：
$1000$ 天要 $9976$ 次加法，$8000$ 天要 $103808$ 次，大約 $10.4$ 倍。 天數變成 $8$
倍，每天被加的次數（大約是層數 $\lg n$）也從大約 $10$ 變成大約 $13$，所以是
$8 \times 13/10 \approx 10.4$ 倍，比 $8$ 倍多一些。

\subsection{地板、天花板和基礎情況}
\label{sec:rec-floors}

$n$ 是 $2$ 的次方時，每次都剛好對半分，遞迴式變成 $T(n) = 2T(n/2) + n$。 令
$n = 2^k$，兩邊除以 $2^k$：
\[
\frac{T(2^k)}{2^k} = \frac{T(2^{k-1})}{2^{k-1}} + 1
= \frac{T(2^{k-2})}{2^{k-2}} + 2 = \cdots = \frac{T(1)}{1} + k = k,
\]
所以 $T(n) = n \lg n$，和程式印出的一樣。 其他的 $n$ 有地板、天花板，算起來麻煩很
多。 下面的定理說明，在常見的情況下，只算 $b$ 的次方就夠了。

\begin{theorem}[只看 $b$ 的次方就夠]
\label{thm:rec-smooth}
設 $b \ge 2$ 是整數，$T(n)$ 和 $g(n)$ 都是非負的非遞減函數，而且存在常數 $d$ 讓
$n$ 夠大時 $g(bn) \le d\, g(n)$。 如果 $n$ 是 $b$ 的次方時 $T(n) = \Theta(g(n))$，那麼
對所有的 $n$ 都有 $T(n) = \Theta(g(n))$。
\end{theorem}

\begin{proof}
設 $n$ 是 $b$ 的次方而且 $n \ge n_0$ 時 $c_1 g(n) \le T(n) \le c_2 g(n)$，並且
$m \ge n_0$ 時 $g(bm) \le d\, g(m)$。 對 $n \ge b n_0$，取 $k$ 讓
$b^k \le n < b^{k+1}$，這時 $b^k > n/b \ge n_0$。 用 $T$、$g$ 非遞減：
\[
T(n) \le T(b^{k+1}) \le c_2\, g(b^{k+1}) \le c_2 d\, g(b^k) \le c_2 d\, g(n),
\]
\[
T(n) \ge T(b^k) \ge c_1 g(b^k) \ge \frac{c_1}{d}\, g(b^{k+1}) \ge \frac{c_1}{d}\, g(n).
\]
取常數 $c_1/d$、$c_2 d$ 和 $b n_0$，就符合 $\Theta$ 的定義。
\end{proof}

用在報表加總：$T(n)$ 是非遞減的（$n$ 加一時，兩半之中剛好有一半多一天，對 $n$ 做
歸納就得到 $T(n + 1) > T(n)$），$g(n) = n \lg n$ 滿足 $g(2n) = 2n(\lg n + 1) \le
4 n \lg n$（$n \ge 2$）。 所以從 $T(2^k) = 2^k k$ 就能推到所有的 $n$ 都有
$T(n) = \Theta(n \lg n)$。

條件 $g(bn) \le d\, g(n)$ 的意思是「$n$ 變成常數倍時，$g$ 也只變成常數倍」。
$n^k$、$n^k \lg n$、$\lg n$ 都滿足；$2^n$ 不滿足，因為 $2^{2n} / 2^n = 2^n$ 不是常數。
所以解出來是指數的遞迴式，不能只看特殊的 $n$ 就下結論。

基礎情況也常常省略不寫，只寫 $T(n) = 2T(n/2) + \Theta(n)$，意思是「$n$ 小於某個
常數時 $T(n) = \Theta(1)$」。 在本章的遞迴式裡，把基礎情況的常數換掉，只會讓答案
差一個常數倍（第~\ref{sec:rec-master}~節的引理~\ref{lem:rec-unroll} 會看到，它只影
響葉子的總成本），不會改變 $\Theta$。

\section{代換法}
\label{sec:rec-subst}

\textbf{代換法}（substitution method）分兩步：先猜出答案的形式，再用數學歸納法證
明，順便找出常數。 猜的方法可以是算幾個小的值、畫遞迴樹、或和看過的遞迴式比
較。 這一節的例子都已經有答案可以猜，重點放在第二步。

\subsection{例子：報表加總的上下界}

從上一節的表可以猜 $n \lfloor \lg n \rfloor \le T(n) \le n \lceil \lg n \rceil$。
證明前先準備地板、天花板的兩個等式。

\begin{lemma}
\label{lem:rec-lg-half}
$n \ge 2$ 時，$\lceil \lg \lceil n/2 \rceil \rceil = \lceil \lg n \rceil - 1$，而且
$\lfloor \lg \lfloor n/2 \rfloor \rfloor = \lfloor \lg n \rfloor - 1$。
\end{lemma}

\begin{proof}
令 $k = \lceil \lg n \rceil \ge 1$，則 $2^{k-1} < n \le 2^k$，所以
$2^{k-2} < n/2 \le \lceil n/2 \rceil \le 2^{k-1}$（$2^{k-1}$ 是整數，不小於
$n/2$，也就不小於 $\lceil n/2 \rceil$）。 落在 $(2^{k-2}, 2^{k-1}]$ 的數取
$\lceil \lg \cdot \rceil$ 都是 $k - 1$。 第二個等式類似：令
$k = \lfloor \lg n \rfloor \ge 1$，則 $2^k \le n < 2^{k+1}$，所以
$2^{k-1} \le \lfloor n/2 \rfloor \le n/2 < 2^k$。
\end{proof}

\begin{theorem}
\label{thm:rec-audit}
報表加總的加法次數對所有 $n \ge 1$ 滿足
$n \lfloor \lg n \rfloor \le T(n) \le n \lceil \lg n \rceil$，所以
$T(n) = \Theta(n \lg n)$。
\end{theorem}

\begin{proof}
對 $n$ 做強歸納法：假設比 $n$ 小的值都成立，證明 $n$ 也成立。

$n = 1$：$T(1) = 0$，兩邊也都是 $0$。

$n \ge 2$：$\lceil n/2 \rceil$、$\lfloor n/2 \rfloor$ 都在 $1$ 和 $n - 1$ 之間，可以
用歸納假設。 上界：因為 $\lfloor n/2 \rfloor \le \lceil n/2 \rceil$，兩個
$\lceil \lg \cdot \rceil$ 都不超過 $\lceil \lg \lceil n/2 \rceil \rceil =
\lceil \lg n \rceil - 1$，所以
\[
T(n) \le \left\lceil \tfrac{n}{2} \right\rceil (\lceil \lg n \rceil - 1)
+ \left\lfloor \tfrac{n}{2} \right\rfloor (\lceil \lg n \rceil - 1) + n
= n(\lceil \lg n \rceil - 1) + n = n \lceil \lg n \rceil.
\]
下界：兩個 $\lfloor \lg \cdot \rfloor$ 都至少是
$\lfloor \lg \lfloor n/2 \rfloor \rfloor = \lfloor \lg n \rfloor - 1$，同樣推導得到
$T(n) \ge n \lfloor \lg n \rfloor$。

最後，$n \ge 2$ 時 $\lfloor \lg n \rfloor \ge \frac{1}{2} \lg n$、$\lceil \lg n \rceil
\le 2 \lg n$，取 $c_1 = 1/2$、$c_2 = 2$、$n_0 = 2$，得到 $T(n) = \Theta(n \lg n)$。
\end{proof}

這個證明沒有忽略任何地板或天花板，而且得到的上下界和程式印出的兩欄完全一樣。 能
這樣做是因為猜的形式 $n \lceil \lg n \rceil$ 剛好和天花板配合；猜不出這麼合適的形
式時，就先用第~\ref{sec:rec-floors}~節的方法只看 $2$ 的次方。

\begin{examnote}
用代換法時，先寫出要證明的不等式，常數寫成符號（例如 $T(n) \le c\, n \lg n$），歸
納步驟要推出\textbf{同一個} $c$ 的同一個式子，最後再說明 $c$ 和 $n_0$ 可以取多少、
基礎情況成不成立。 推出「$\le cn + n$」這種多一項的式子時，證明還沒有完成。
\end{examnote}

\subsection{常見的錯誤：常數在歸納中變大}

下面是一個錯誤的「證明」，聲稱報表加總的 $T(n) = O(n)$。

\begin{quote}
假設比 $n$ 小的值都有 $T(k) \le ck$。 則
$T(n) \le c \lceil n/2 \rceil + c \lfloor n/2 \rfloor + n = cn + n = (c + 1) n =
O(n)$。
\end{quote}

錯在最後一個等號。 歸納假設用的是常數 $c$，要證明的也必須是 $T(n) \le cn$；推出來
的卻是 $(c + 1)n$。 每往上一層，「常數」就多 $1$，$\lg n$ 層以後是 $c + \lg n$，已
經不是常數了。 定理~\ref{thm:rec-audit} 的下界 $n \lfloor \lg n \rfloor$ 也說明
$T(n) \ne O(n)$。 $O(n)$ 這個寫法藏起了常數，在歸納法裡要把常數寫出來，才看得到
它有沒有變。

\subsection{減去低次項}

有時候猜的答案是對的，歸納卻差一點點過不去。 以呼叫次數
$C(n) = C(\lceil n/2 \rceil) + C(\lfloor n/2 \rfloor) + 1$、$C(1) = 1$ 為例，程式
印出 $2n - 1$，猜 $C(n) \le cn$：
\[
C(n) \le c \lceil n/2 \rceil + c \lfloor n/2 \rfloor + 1 = cn + 1.
\]
多了 $1$，不管 $c$ 取多大都推不出 $C(n) \le cn$。 這時該做的是把猜的式子\textbf{減去一個低次項}，改成
$C(n) \le cn - d$，不必改猜更大的 $n \lg n$ 或 $n^2$：
\[
C(n) \le (c \lceil n/2 \rceil - d) + (c \lfloor n/2 \rfloor - d) + 1 = cn - 2d + 1
\le cn - d,
\]
最後一步只要 $d \ge 1$。 基礎情況要 $C(1) = 1 \le c - d$。 取 $c = 2$、$d = 1$，得到
$C(n) \le 2n - 1$。 把上面的 $\le$ 都換成 $=$ 再做一次，就得到 $C(n) = 2n - 1$，和
程式一致。

減去一項之後，要證明的式子變強了，可是歸納假設也一起變強，兩個呼叫各自少了 $d$，
合起來少 $2d$，剛好吸收掉多出來的 $1$。 這是代換法最常用的技巧。

\subsection{換變數}

$T(n) = 2T(\sqrt{n}) + 1$、$T(2) = 1$，$n$ 取 $2, 4, 16, 256, \ldots$（這樣 $\sqrt{n}$
一直是整數）。 開根號不好處理，可以換成指數：令 $m = \lg n$，$S(m) = T(2^m)$。
$\sqrt{n} = 2^{m/2}$，所以
\[
S(m) = 2S(m/2) + 1, \qquad S(1) = 1.
\]
這和上一小段的呼叫次數是同一種遞迴式，$m$ 是 $2$ 的次方時 $S(m) = 2m - 1$。 換
回來：$T(n) = 2 \lg n - 1 = \Theta(\lg n)$。 驗算：$T(4) = 2T(2) + 1 = 3$，
$T(16) = 2T(4) + 1 = 7$，都等於 $2 \lg n - 1$。 把規模開根號的資料結構會出現這種
遞迴式。

\section{遞迴樹}
\label{sec:rec-tree}

代換法需要先猜。 \textbf{遞迴樹}（recursion tree）是一個好用的猜法，很多時候也可
以直接算出答案：每一次呼叫畫成一個節點，節點裡寫這次呼叫\textbf{自己}做的工作
（不含下面的呼叫），一層一層加起來，最後把所有層加總。

遞迴樹的層從根算起，根是第 $0$ 層，第 $j$ 層的節點離根 $j$ 步，這樣第 $j$ 層的節點數
剛好是 $a^j$ 這種形式。 第~7~章講一般的樹時，根改叫第 $1$ 層、高度數節點；兩種數法
差 $1$，作答時寫清楚用的是哪一種就好。

\subsection{畫樹、算每層、加總}

以 $T(n) = 3T(n/2) + n$ 為例，$n$ 是 $2$ 的次方，$T(1) = \Theta(1)$。 第~17~章的
大數乘法就是這個形狀：一個規模 $n$ 的問題拆成三個規模 $n/2$ 的問題，另外做線性的
工作。

\begin{center}
\begin{tikzpicture}
\node[rnode] (r) at (0,0) {$n$};
\foreach \i/\x in {1/-4, 2/0, 3/4} {
    \node[rnode] (a\i) at (\x,-1.4) {$n/2$};
    \draw[redge] (r) -- (a\i);
    \foreach \j/\dx in {1/-1.25, 2/0, 3/1.25} {
        \node[rnode, minimum width=0.8cm] (b\i\j) at (\x+\dx,-2.8) {$n/4$};
        \draw[redge] (a\i) -- (b\i\j);
    }
}
\foreach \x in {-5.25,-4,-2.75,-1.25,0,1.25,2.75,4,5.25} {
    \node[font=\small] at (\x,-3.6) {$\vdots$};
}
\foreach \x in {-5.25,-4,-1.25,0,2.75,5.25} {
    \node[rleaf] at (\x,-4.5) {$\Theta(1)$};
}
\node[font=\small] at (-2.6,-4.5) {$\cdots$};
\node[font=\small] at (1.4,-4.5) {$\cdots$};
\node[font=\small] at (4.0,-4.5) {$\cdots$};
\draw[decorate, decoration={brace, amplitude=5pt}] (-6.3,-4.5) -- (-6.3,0)
    node[midway, left=6pt, font=\small, align=right] {$\lg n$\\層};
\node[rsum] at (6.3,0) {$n$};
\node[rsum] at (6.3,-1.4) {$\tfrac{3}{2}\,n$};
\node[rsum] at (6.3,-2.8) {$\left(\tfrac{3}{2}\right)^{2} n$};
\node[rsum] at (6.3,-3.6) {$\vdots$};
\node[rsum] at (6.3,-4.5) {$\Theta(n^{\lg 3})$};
\end{tikzpicture}
\diagcaption{$T(n) = 3T(n/2) + n$ 的遞迴樹。 右邊是每一層的總和，一層比一層大
$3/2$ 倍；最下面一層是 $3^{\lg n} = n^{\lg 3}$ 個葉子。}
\label{fig:rec-tree-3half}
\end{center}

一層一層數：

\begin{itemize}
\item 第 $j$ 層有 $3^j$ 個節點，每個的規模是 $n/2^j$，自己的工作是 $n/2^j$，這一層
合計 $3^j \cdot n/2^j = (3/2)^j n$。
\item 規模在 $j = \lg n$ 時變成 $1$，所以第 $0$ 到第 $\lg n - 1$ 層是內部節點，第
$\lg n$ 層是葉子，共有 $3^{\lg n} = n^{\lg 3}$ 個，每個 $\Theta(1)$。
\end{itemize}

加起來（等比級數，公比 $3/2$）：
\[
T(n) = n \sum_{j=0}^{\lg n - 1} \left(\frac{3}{2}\right)^{j} + \Theta(n^{\lg 3})
= n \cdot \frac{(3/2)^{\lg n} - 1}{3/2 - 1} + \Theta(n^{\lg 3})
= 2n^{\lg 3} - 2n + \Theta(n^{\lg 3}) = \Theta(n^{\lg 3}).
\]
這裡用了 $(3/2)^{\lg n} = n^{\lg (3/2)} = n^{\lg 3 - 1}$。 $\lg 3 \approx 1.585$，所以
這個演算法比 $n^2$ 好。 這個例子每一步都是精確的等式，不需要再驗證；如果畫樹時有
近似（例如忽略了地板），習慣上再用代換法確認一次。

\subsection{三種形狀}

上面的例子每一層是上一層的 $3/2$ 倍。 一般來說，相鄰兩層的總和之比 $r$ 決定了答
案，這正是第~1~章等比級數的兩句話：公比小於 $1$，總和和第一項同一級；公比大於
$1$，總和和最後一項同一級。

\begin{center}
\begin{tabular}{llll}
\toprule
每層的變化 & 誰最大 & 答案 & 例子 \\
\midrule
越來越小（$r < 1$） & 根 & $\Theta(\text{根的工作})$ &
$2T(n/2) + n^2 = \Theta(n^2)$ \\
每層一樣（$r = 1$） & 平手 & $\Theta(\text{每層} \times \text{層數})$ &
$2T(n/2) + n = \Theta(n \lg n)$ \\
越來越大（$r > 1$） & 葉子 & $\Theta(\text{葉子的個數})$ &
$3T(n/2) + n = \Theta(n^{\lg 3})$ \\
\bottomrule
\end{tabular}
\end{center}

第一列的例子：每層的總和是 $n^2, n^2/2, n^2/4, \ldots$，全部加起來不超過 $2n^2$；
葉子只有 $n$ 個。 所以是 $\Theta(n^2)$，大部分的工作在根做完。 下一節的主定理就
是把這張表寫成精確的定理。

\begin{examnote}
用遞迴樹作答時，寫出第 $j$ 層的節點數、每個節點的規模、這一層的總和、總層數，以
及葉子的個數和總成本，再寫加總。 葉子那一層最容易漏掉，層數也容易差一。 葉子的個數
$n^{\log_b a}$ 可以拿來檢查答案：答案不可能比它小。 例如把 $3T(n/2) + n$ 答成
$\Theta(n \lg n)$，一比 $n^{\lg 3} \approx n^{1.585}$ 就知道錯了。
\end{examnote}

\subsection{不平衡的樹}

兩個子問題不一樣大時，遞迴樹的左右兩邊深度不同。 看
\[
T(n) = T(n/4) + T(3n/4) + n, \qquad T(n) = \Theta(1) \ (n < 4).
\]

\begin{center}
\begin{tikzpicture}
\node[rnode] (r) at (0,0) {$n$};
\node[rnode] (a) at (-3,-1.4) {$n/4$};
\node[rnode] (b) at (3,-1.4) {$3n/4$};
\node[rnode] (c) at (-4.4,-2.8) {$n/16$};
\node[rnode] (d) at (-1.6,-2.8) {$3n/16$};
\node[rnode] (e) at (1.6,-2.8) {$3n/16$};
\node[rnode] (f) at (4.4,-2.8) {$9n/16$};
\foreach \u/\v in {r/a, r/b, a/c, a/d, b/e, b/f} { \draw[redge] (\u) -- (\v); }
\node[font=\small] at (-4.4,-3.6) {$\vdots$};
\node[font=\small] at (4.4,-3.6) {$\vdots$};
\node[rleaf] (l) at (-4.4,-4.5) {$\Theta(1)$};
\node[rleaf] (rr) at (4.4,-5.4) {$\Theta(1)$};
\node[font=\scriptsize, align=center] at (-1.9,-4.5) {最淺的葉子\\深度約 $\log_4 n$};
\node[font=\scriptsize, align=center] at (1.9,-5.4) {最深的葉子\\深度約 $\log_{4/3} n$};
\node[rsum] at (6.0,0) {$n$};
\node[rsum] at (6.0,-1.4) {$n/4 + 3n/4 = n$};
\node[rsum] at (6.0,-2.8) {$n$};
\node[rsum] at (6.0,-3.9) {不滿的層 $\le n$};
\end{tikzpicture}
\diagcaption{$T(n) = T(n/4) + T(3n/4) + n$ 的遞迴樹。 每一層的規模加起來是 $n$；
往左走規模縮得快，往右走縮得慢，所以最下面幾層是不滿的。}
\label{fig:rec-tree-unbalanced}
\end{center}

同一層的規模加起來是 $n$（每個節點分成的兩塊加起來就是它自己），所以每一層滿的層
成本剛好是 $n$，不滿的層不超過 $n$。 一直往左走，規模每層變成 $1/4$，大約
$\log_4 n$ 層就到葉子；一直往右走，每層只變成 $3/4$，要大約 $\log_{4/3} n$ 層。

\begin{itemize}
\item \textbf{上界}：最多 $\log_{4/3} n + 1$ 層，每層的內部成本不超過 $n$；葉子的
規模都至少 $1$（規模 $s \ge 4$ 的節點分成 $s/4 \ge 1$ 和 $3s/4 \ge 3$），規模加起
來不超過 $n$，所以葉子最多 $n$ 個。 合計 $O(n \lg n)$。
\item \textbf{下界}：前 $\lfloor \log_4 n \rfloor$ 層都是滿的，每層成本 $n$，合計
$\Omega(n \lg n)$。
\end{itemize}

所以 $T(n) = \Theta(n \lg n)$。 這裡忽略了「規模不是整數」等細節，用代換法驗證上
界：猜 $T(n) \le d\, n \lg n$（$n \ge 2$）。 $n \ge 8$ 時兩個子問題的規模都至少 $2$，
用歸納假設：
\[
T(n) \le d\,\frac{n}{4} \lg \frac{n}{4} + d\,\frac{3n}{4} \lg \frac{3n}{4} + n
= d\, n \lg n - d\, n \left(\frac{1}{4} \lg 4 + \frac{3}{4} \lg \frac{4}{3}\right) + n .
\]
括號裡的數 $h = \frac{1}{2} + \frac{3}{4}\lg\frac{4}{3} \approx 0.811$ 是正的常數，
只要 $d \ge 1/h$（大約 $1.23$），減掉的 $d h n$ 就蓋過加上的 $n$，得到
$T(n) \le d\, n \lg n$。 $2 \le n < 8$ 時最多往下三層就到基礎情況，$T(n)$ 不超過某個常數
$M$，而 $n \lg n \ge 2$，所以 $d$ 再取大一點（$d \ge M/2$）就都成立。

兩塊的比例換成任何固定的 $\alpha : (1 - \alpha)$（$0 < \alpha < 1$），結論都一樣是
$\Theta(n \lg n)$，只是 $\alpha$ 越偏，常數越大。 第~14~章分析快速排序時會用到這
個結論。

\section{主定理}
\label{sec:rec-master}

很多分治演算法的遞迴式長成
\[
T(n) = aT(n/b) + f(n),
\]
意思是：規模 $n$ 的問題拆成 $a$ 個規模 $n/b$ 的子問題，呼叫以外另外花 $f(n)$。 這
裡 $a \ge 1$、$b > 1$ 是常數。 遞迴樹的第 $j$ 層有 $a^j$ 個節點，每個的成本是
$f(n/b^j)$；葉子在第 $\log_b n$ 層（深度 $\log_b n$），共有
\[
a^{\log_b n} = n^{\log_b a}
\]
個。 上一節的表是在比較根的成本 $f(n)$ 和葉子的總數 $n^{\log_b a}$，主定理把這個比
較說清楚。

\subsection{定理}

\begin{theorem}[主定理]
\label{thm:rec-master}
設 $a \ge 1$、$b > 1$ 是常數，$f(n)$ 是正的函數，$T(n)$ 定義在 $b$ 的次方
$n = 1, b, b^2, \ldots$ 上：$T(1)$ 是正的常數，$n \ge b$ 時
$T(n) = aT(n/b) + f(n)$。
\begin{enumerate}
\item 若存在常數 $\varepsilon > 0$ 使 $f(n) = O(n^{\log_b a - \varepsilon})$，則
$T(n) = \Theta(n^{\log_b a})$。
\item 若 $f(n) = \Theta(n^{\log_b a})$，則 $T(n) = \Theta(n^{\log_b a} \lg n)$。
\item 若存在常數 $\varepsilon > 0$ 使 $f(n) = \Omega(n^{\log_b a + \varepsilon})$，
而且存在常數 $c < 1$，使得 $n$ 夠大時
\[
a f(n/b) \le c f(n) \qquad \text{（正則條件）},
\]
則 $T(n) = \Theta(f(n))$。
\end{enumerate}
\end{theorem}

三種情況對應上一節的三列：情況 1 是葉子贏，情況 2 平手，情況 3 根贏。 有兩個地方
要注意。

第一，情況 1 和情況 3 要求 $f(n)$ 和 $n^{\log_b a}$ 差了\textbf{至少一個}
$n^{\varepsilon}$ 的因子，叫做「多項式地小（大）」。 只是「比較小」不夠。 例如和
$n^2$ 比，$n \lg n$ 是多項式地小（$n \lg n = O(n^{1.5})$），但 $n^2 / \lg n$ 只差了
$\lg n$，任何 $n^{\varepsilon}$ 最後都比 $\lg n$ 大，所以不算。

第二，情況 3 的\textbf{正則條件}（regularity condition）說第 1 層的總和
$a f(n/b)$ 不超過根的 $c$ 倍，這樣每一層都至少縮成上一層的 $c$ 倍，總和是遞減的
等比級數。 $f(n) = n^k$ 而且 $k > \log_b a$ 時它一定成立：
$a (n/b)^k = (a/b^k) n^k$，而 $a/b^k = b^{\log_b a - k} < 1$，取 $c = a/b^k$。

情況 2 常用下面這個推廣，它把 $f(n)$ 多一個 $\lg^k n$ 因子的情形也包進來。

\begin{theorem}[推廣的情況 2]
\label{thm:rec-master2}
同定理~\ref{thm:rec-master} 的假設。 若存在常數 $k \ge 0$ 使
$f(n) = \Theta(n^{\log_b a} \lg^k n)$，則 $T(n) = \Theta(n^{\log_b a} \lg^{k+1} n)$。
\end{theorem}

$k = 0$ 就是原本的情況 2。 定理只對 $b$ 的次方敘述；遞迴式寫成
$T(\lfloor n/b \rfloor)$ 或 $T(\lceil n/b \rceil)$ 時結論也一樣成立，證明比較繁
瑣，本書不做。 如果知道 $T$ 是非遞減的，也可以直接用定理~\ref{thm:rec-smooth} 推到
所有的 $n$。 證明放在本節最後的框裡。

\subsection{怎麼用}

\begin{enumerate}
\item 確認是 $aT(n/b) + f(n)$ 的形式：$a$、$b$ 是常數，$a \ge 1$，$b > 1$。
\item 算出 $\log_b a$。
\item 比較 $f(n)$ 和 $n^{\log_b a}$：差到多項式倍（情況 1 或 3），還是同一級（情況
2，或推廣的情況 2）。
\item 情況 3 還要驗證正則條件。
\end{enumerate}

\begin{center}
\begin{tabular}{lcccl}
\toprule
遞迴式 & $a$、$b$ & $n^{\log_b a}$ & 情況 & 答案 \\
\midrule
$5T(n/2) + n^2$ & 5、2 & $n^{2.32\ldots}$ & 1（$\varepsilon = 0.3$） & $\Theta(n^{\lg 5})$ \\
$27T(n/3) + n^3$ & 27、3 & $n^3$ & 2 & $\Theta(n^3 \lg n)$ \\
$3T(n/9) + n$ & 3、9 & $n^{1/2}$ & 3（$\varepsilon = 1/2$，$c = 1/3$） & $\Theta(n)$ \\
$T(n/2) + 1$ & 1、2 & $n^0 = 1$ & 2 & $\Theta(\lg n)$ \\
$4T(n/16) + \sqrt{n}$ & 4、16 & $n^{1/2}$ & 2 & $\Theta(\sqrt{n} \lg n)$ \\
$3T(n/3) + n \lg n$ & 3、3 & $n$ & 推廣的 2（$k = 1$） & $\Theta(n \lg^2 n)$ \\
\bottomrule
\end{tabular}
\end{center}

幾列的說明。 第 3 列：$f(n) = n = \Omega(n^{1/2 + 1/2})$，正則條件
$3 \cdot (n/9) = n/3 \le \frac{1}{3} n$，取 $c = 1/3$。 第 4 列是二分搜尋，$a = 1$
時 $n^{\log_b a} = 1$，$f(n) = 1$ 和它同一級。 最後一列：$n \lg n$ 比 $n$ 大，但沒
有大到 $n^{1+\varepsilon}$（$\lg n = o(n^{\varepsilon})$），原本的三種情況都不適
用；它是推廣情況 2 的 $k = 1$。

\begin{examnote}
用主定理作答時，依序寫出 $a$、$b$、$\log_b a$、$f(n)$、屬於哪一種情況。 情況 1、3
寫出一個可用的 $\varepsilon$，情況 3 再寫出正則條件的 $c$ 和驗證過程。 只寫「由主
定理得 $\Theta(n^2)$」，看不出你有沒有檢查條件。
\end{examnote}

\subsection{不能用的時候}

\textbf{$f(n)$ 和 $n^{\log_b a}$ 差不到多項式倍。} 例如
$T(n) = 4T(n/2) + n^2 / \lg n$。 $n^{\log_2 4} = n^2$，$f(n)$ 比 $n^2$ 小，但只小了
$\lg n$，不是情況 1；推廣的情況 2 要 $k \ge 0$，這裡等於 $k = -1$，也不行。 直接用
遞迴樹算（$n = 2^m$）：第 $j$ 層的總和是
\[
4^j \cdot \frac{(n/2^j)^2}{\lg (n/2^j)} = \frac{n^2}{m - j},
\]
把 $j = 0, \ldots, m - 1$ 加起來是 $n^2 (1/m + \cdots + 1/1) = n^2 H_m$，葉子是
$4^m = n^2$ 個。 $H_m = \Theta(\lg m)$，所以
$T(n) = \Theta(n^2 \lg \lg n)$，比 $n^2$ 大，比 $n^2 \lg n$ 小。

\textbf{不是 $aT(n/b) + f(n)$ 的形式。} 規模是減的，例如 $T(n) = T(n-1) + n$（見
第~\ref{sec:rec-other}~節）；子問題不一樣大，例如 $T(n/4) + T(3n/4) + n$（見
第~\ref{sec:rec-tree}~節和第~\ref{sec:rec-gf}~節）；$a$ 不是常數，例如
$T(n) = n\,T(n/2) + 1$；$a < 1$，例如 $T(n) = \frac{1}{2} T(n/2) + n$，「呼叫半
次」不是一個合理的遞迴程式。

\textbf{情況 3 的正則條件不成立。} $f(n)$ 是多項式時一定成立，但 $f(n)$ 忽大忽小
時可能不成立，這時定理不給結論，要回到遞迴樹自己算。

\subsection{實際算一次}

下面的程式對 $n = 2^k$ 直接用遞迴式一層一層算出 $T(n)$（$T(1) = 1$），再除以猜的
答案。 比值趨近一個正的常數，表示猜的 $\Theta$ 對；比值趨近 $0$ 或越來越大，表示
猜錯了。 四個遞迴式和印出的比值是：
A：$3T(n/2) + n$，印 $T(n)/n^{\lg 3}$；
B：$2T(n/2) + n$，印 $T(n)/(n \lg n)$；
C：$2T(n/2) + n^2$，印 $T(n)/n^2$；
D：$4T(n/2) + n^2/\lg n$，D1 印 $T(n)/(n^2 \lg n)$，D2 印 $T(n)/(n^2 \lg \lg n)$。

\begin{example}[label=ex:rec-ratio]{c}{(recur\_ratio.c) 用比值檢查遞迴式的解}
#include <stdio.h>
#include <math.h>

// n = 2^k 時，用遞迴式本身一層一層算出 T(n)，T(1) = 1。
//   A: T(n) = 3T(n/2) + n            除以 n^(lg 3) = 3^k
//   B: T(n) = 2T(n/2) + n            除以 n lg n
//   C: T(n) = 2T(n/2) + n^2          除以 n^2
//   D: T(n) = 4T(n/2) + n^2 / lg n   D1 除以 n^2 lg n，D2 除以 n^2 lg lg n
int main(void)
{
    double ta = 1, tb = 1, tc = 1, td = 1;  // T(2^0)
    double n = 1, pow3 = 1;                  // 2^k、3^k
    printf("k A B C D1 D2\n");
    for (int k = 1; k <= 64; k++) {
        n *= 2;
        pow3 *= 3;
        ta = 3 * ta + n;
        tb = 2 * tb + n;
        tc = 2 * tc + n * n;
        td = 4 * td + n * n / k;              // lg n = k
        if (k == 2 || k == 4 || k == 8 || k == 16 || k == 32 || k == 64) {
            printf("%d %.4f %.4f %.4f %.4f %.4f\n", k, ta / pow3,
                   tb / (n * k), tc / (n * n), td / (n * n * k),
                   td / (n * n * log2(k)));
        }
    }
    return 0;
}
\end{example}
\iosample{}{k A B C D1 D2 \\ 2 2.1111 1.5000 1.7500 1.2500 2.5000 \\ 4 2.6049 1.2500 1.9375 0.7708 1.5417 \\ 8 2.9220 1.1250 1.9961 0.4647 1.2393 \\ 16 2.9970 1.0625 2.0000 0.2738 1.0952 \\ 32 3.0000 1.0312 2.0000 0.1581 1.0117 \\ 64 3.0000 1.0156 2.0000 0.0897 0.9573}

A、B、C 都很快穩定下來：A 的精確解是 $3n^{\lg 3} - 2n$，比值是
$3 - 2(2/3)^k$；B 是 $n \lg n + n$，比值 $1 + 1/k$；C 是 $2n^2 - n$。 D1 一路往 $0$
掉，說明 $n^2 \lg n$ 猜得太大，把 $4T(n/2) + n^2/\lg n$ 硬套成情況 2 會得到錯的答
案。 D2 的比值是 $(1 + H_k)/\lg k$，會慢慢趨近 $\ln 2 \approx 0.69$；$k$ 到 $64$
（$n = 2^{64}$）都還沒停下來。 牽涉到 $\lg \lg n$ 的比值收斂得非常慢，這種程式只能
幫忙檢查，不能代替證明。

\begin{whybox}[主定理的證明（$n$ 是 $b$ 的次方）]
令 $n = b^m$，$w = \log_b a$，所以 $a^m = n^w$，$a = b^w$。 $f$ 是正的，而且每個
$O$、$\Omega$、$\Theta$ 的不等式只在 $n \ge n_0$ 時成立；$n_0$ 以下只有有限多個
$b$ 的次方，把常數調大（或調小）就能讓不等式對所有出現的 $b$ 的次方都成立。 下面
直接這樣用。

\begin{lemma}
\label{lem:rec-unroll}
$T(n) = T(1)\, n^{w} + g(n)$，其中 $g(n) = \sum_{j=0}^{m-1} a^j f(n/b^j)$。
\end{lemma}

\begin{proof}
反覆代入：$T(n) = f(n) + aT(n/b) = f(n) + af(n/b) + a^2 T(n/b^2) = \cdots =
\sum_{j=0}^{m-1} a^j f(n/b^j) + a^m T(1)$，而 $a^m = n^w$。
\end{proof}

$T(1)\, n^w$ 是葉子的總成本，$g(n)$ 是內部節點的總成本，第 $j$ 項是第 $j$ 層。 剩
下的工作是估計 $g(n)$。

\textbf{情況 1。} 設 $f(x) \le c\, x^{w - \varepsilon}$。 因為
$a^j / (b^j)^{w - \varepsilon} = (b^{\varepsilon})^j$，
\[
g(n) \le c \sum_{j=0}^{m-1} a^j \left(\frac{n}{b^j}\right)^{w - \varepsilon}
= c\, n^{w - \varepsilon} \sum_{j=0}^{m-1} (b^{\varepsilon})^j
= c\, n^{w - \varepsilon} \cdot \frac{b^{\varepsilon m} - 1}{b^{\varepsilon} - 1}
< \frac{c}{b^{\varepsilon} - 1}\, n^{w},
\]
最後一步用了 $b^{\varepsilon m} = n^{\varepsilon}$。 所以
$T(1) n^w \le T(n) \le (T(1) + \frac{c}{b^\varepsilon - 1}) n^w$，$T(n) = \Theta(n^w)$。

\textbf{情況 2（含推廣，$k \ge 0$）。} $g(n)$ 裡 $f$ 的參數 $n/b^j$ 都至少是 $b$，
$\lg$ 是正的。 設 $c_1 x^w \lg^k x \le f(x) \le c_2 x^w \lg^k x$。 因為
$a^j (n/b^j)^w = n^w$、$\lg(n/b^j) = (m - j)\lg b$，第 $j$ 層夾在
$c_1 n^w (\lg b)^k (m-j)^k$ 和 $c_2 n^w (\lg b)^k (m-j)^k$ 之間。 而
\[
\left(\frac{m}{2}\right)^{k+1} \le \sum_{i=1}^{m} i^k \le m^{k+1}
\]
（左邊只留 $i \ge m/2$ 的項，至少 $m/2$ 項，每項至少 $(m/2)^k$）。 $m = \lg n / \lg b$，
所以 $g(n) = \Theta(n^w \lg^{k+1} n)$。 $n \ge b$ 時 $\lg n \ge \lg b > 0$，葉子的
$n^w$ 被它蓋過，$T(n) = \Theta(n^w \lg^{k+1} n)$。

\textbf{情況 3。} 先假設正則條件對所有 $x \ge b$ 都成立。 用 $j$ 次得到
$a^j f(n/b^j) \le c^j f(n)$，所以
\[
f(n) \le g(n) \le f(n) \sum_{j=0}^{\infty} c^j = \frac{f(n)}{1 - c},
\]
左邊是 $j = 0$ 那一項。 $g(n) = \Theta(f(n))$；又
$f(n) = \Omega(n^{w + \varepsilon})$，所以葉子的 $n^w = O(f(n))$，$T(n) =
\Theta(f(n))$。 如果正則條件只在 $x \ge n_1$ 時成立，參數小於 $n_1$ 的那幾層用不
上它；這樣的層最多 $\log_b n_1 + 1$ 個（常數），每層不超過
$a^m \cdot \max_{x < n_1} f(x) = O(n^w)$，合計 $O(n^w) = O(f(n))$，結論不變。
\end{whybox}

\section{其他常見的形式}
\label{sec:rec-other}

\subsection{\texorpdfstring{每次減一：$T(n) = T(n-1) + f(n)$}{每次減一：T(n) = T(n-1) + f(n)}}

每次只把規模減 $1$ 的遞迴，例如走訪串列、選擇排序、快速排序的最壞情況（第~14~章），
遞迴樹是一條鏈，直接展開：
\[
T(n) = f(n) + T(n-1) = f(n) + f(n-1) + T(n-2) = \cdots = T(0) + \sum_{i=1}^{n} f(i).
\]
問題變成第~1~章的求和。

\begin{center}
\begin{tabular}{lll}
\toprule
$f(n)$ & $\sum_{i=1}^{n} f(i)$ & $T(n)$ \\
\midrule
$1$ & $n$ & $\Theta(n)$ \\
$n$ & $n(n+1)/2$ & $\Theta(n^2)$ \\
$n^k$（$k \ge 0$ 常數） & & $\Theta(n^{k+1})$ \\
$\lg n$ & $\lg(n!)$ & $\Theta(n \lg n)$ \\
$1/n$ & $H_n$ & $\Theta(\lg n)$ \\
\bottomrule
\end{tabular}
\end{center}

如果每次呼叫自己\textbf{兩次}，形式變成 $T(n) = 2T(n-1) + f(n)$，就完全不同了。
《C 程式設計基礎》第~15~章 15.8 節的河內塔，搬 $n$ 個盤子的步數是
$H(n) = 2H(n-1) + 1$、$H(0) = 0$。 兩邊加 $1$：
$H(n) + 1 = 2\bigl(H(n-1) + 1\bigr)$，所以 $H(n) + 1 = 2^n (H(0) + 1)$，
$H(n) = 2^n - 1$。 規模\textbf{減}一個常數、呼叫次數又大於一，答案是指數；規模
\textbf{除以}一個常數，答案是多項式。 看到遞迴式時，先分清楚是哪一種。

\subsection{\texorpdfstring{開根號：$T(n) = T(\sqrt{n}) + 1$}{開根號：T(n) = T(√n) + 1}}

規模每次開根號，$k$ 次之後是 $n^{1/2^k}$。 設 $n \le 2$ 時停下：
\[
n^{1/2^k} \le 2 \iff \frac{\lg n}{2^k} \le 1 \iff k \ge \lg \lg n,
\]
所以 $T(n) = \Theta(\lg \lg n)$。 也可以用上一節的換變數：$m = \lg n$，
$S(m) = S(m/2) + 1 = \Theta(\lg m)$。 $\lg \lg n$ 長得非常慢，$n = 2^{64}$ 時只有
$6$。

\subsection{線性遞迴與特徵方程式}

像 $a_n = a_{n-1} + 2a_{n-2}$ 這種「每一項是前面固定幾項的常數倍相加」的遞迴式叫做
\textbf{常係數線性遞迴}（linear recurrence with constant coefficients）。 往前看兩項
的叫二階。 它們有一套標準的解法。

\begin{theorem}[二階齊次線性遞迴]
\label{thm:rec-linear}
設 $c_1$、$c_2$ 是常數，$c_2 \ne 0$，數列滿足 $a_n = c_1 a_{n-1} + c_2 a_{n-2}$
（$n \ge 2$）。 方程式 $x^2 = c_1 x + c_2$ 叫做它的\textbf{特徵方程式}
（characteristic equation）。
\begin{enumerate}
\item 特徵方程式有兩個不同的根 $r_1$、$r_2$ 時，$a_n = A r_1^n + B r_2^n$。
\item 特徵方程式有重根 $r$ 時，$a_n = (A + Bn) r^n$。
\end{enumerate}
常數 $A$、$B$ 由 $a_0$、$a_1$ 決定。
\end{theorem}

\begin{proof}
先說明公式裡的數列都滿足遞迴式。 $c_2 \ne 0$，所以 $0$ 不是根。 $r^n$ 滿足遞迴式
等於 $r^n = c_1 r^{n-1} + c_2 r^{n-2}$，兩邊除以 $r^{n-2}$ 就是 $r$ 是特徵方程式的
根。 兩個滿足遞迴式的數列，它們的常數倍相加也滿足（遞迴式兩邊都是線性的）。 重根
時 $r = c_1/2$、$c_2 = -r^2$，檢查 $n r^n$：
\[
c_1 (n-1) r^{n-1} + c_2 (n-2) r^{n-2} = r^{n}\bigl(2(n-1) - (n-2)\bigr) = n r^n .
\]

再說明一定選得到 $A$、$B$ 符合 $a_0$、$a_1$。 不同根：$A + B = a_0$、
$A r_1 + B r_2 = a_1$，係數的行列式 $r_2 - r_1 \ne 0$，有解。 重根：$A = a_0$、
$(A + B) r = a_1$，$r \ne 0$，有解。

最後，滿足同一個二階遞迴式、而且前兩項相同的兩個數列，用歸納法可知每一項都相同。
所以 $a_n$ 就是公式給的數列。
\end{proof}

根是複數時公式仍然成立，只是要用複數計算；本書的例子都是實數根。 《C 程式設計基礎》第~15~章
的費氏數列也可以這樣解：$x^2 = x + 1$ 的根是 $\varphi = (1 + \sqrt{5})/2$ 和
$(1 - \sqrt{5})/2$，第二個根的絕對值小於 $1$，所以費氏數是
$\Theta(\varphi^n)$。

\textbf{例 1（兩個不同的根）。} 一條寬 $1$、長 $n$ 的走道要鋪地磚。 磚有兩種長度：
長 $1$ 的只有灰色，長 $2$ 的有紅、藍兩色。 鋪法有幾種？ 看最右邊那一塊：是灰磚，剩
下長 $n - 1$；是紅磚或藍磚，剩下長 $n - 2$。 所以
\[
a_n = a_{n-1} + 2a_{n-2}, \qquad a_0 = 1,\ a_1 = 1
\]
（長 $0$ 只有「什麼都不鋪」一種）。 特徵方程式 $x^2 = x + 2$，根是 $2$ 和 $-1$，
$a_n = A \cdot 2^n + B(-1)^n$。 代 $n = 0, 1$：$A + B = 1$、$2A - B = 1$，解得
$A = 2/3$、$B = 1/3$。
\[
a_n = \frac{2^{n+1} + (-1)^n}{3}.
\]
驗算：$a_2 = a_1 + 2a_0 = 3 = (8 + 1)/3$，$a_3 = 3 + 2 = 5 = (16 - 1)/3$。

\textbf{例 2（重根）。} $a_n = 6a_{n-1} - 9a_{n-2}$，$a_0 = 1$、$a_1 = 6$。 特徵方
程式 $x^2 - 6x + 9 = (x - 3)^2$，重根 $3$，$a_n = (A + Bn) 3^n$。 $a_0 = A = 1$，
$a_1 = 3(1 + B) = 6$，$B = 1$，所以 $a_n = (n + 1) 3^n$。 驗算 $a_2 = 36 - 9 = 27 =
3 \cdot 9$。 重根時如果寫成 $A \cdot 3^n + B \cdot 3^n$，那其實只有一個常數，配不出
兩個初始值。

\subsection{非齊次的線性遞迴}

多了一項和 $a$ 無關的 $g(n)$：
\[
a_n = c_1 a_{n-1} + c_2 a_{n-2} + g(n).
\]
如果已經找到\textbf{任何一個}滿足它的數列 $p_n$（叫\textbf{特解}），那麼所有的解都
是 $p_n$ 加上齊次遞迴式（把 $g(n)$ 拿掉）的解。 理由：兩個解相減，$g(n)$ 消掉了，
差滿足齊次遞迴式。 所以解法是三步：解齊次部分（定理~\ref{thm:rec-linear}）、猜一
個特解、\textbf{最後}才用初始值求 $A$、$B$。

特解照 $g(n)$ 的樣子猜：

\begin{center}
\begin{tabular}{ll}
\toprule
$g(n)$ & 特解的猜法 \\
\midrule
常數 $d$ & 常數 $p$ \\
$k$ 次多項式 & $k$ 次多項式（係數待定） \\
$d \cdot s^n$ & $p \cdot s^n$ \\
\bottomrule
\end{tabular}
\end{center}

例外：如果 $1$（前兩列）或 $s$（第三列）剛好是特徵方程式的根，猜的式子會被齊次部
分吃掉，要再乘上 $n$（重根乘 $n^2$）。 觀念題第 11 題 (b) 就是這種情形。

\textbf{例 3。} 回到鋪地磚。 照定義直接寫遞迴函式，最右邊是長 $2$ 的磚時，紅、藍
各呼叫一次，所以呼叫次數是
\[
C(n) = C(n-1) + 2C(n-2) + 1 \ (n \ge 2), \qquad C(0) = C(1) = 1 .
\]
齊次部分和例 1 一樣，是 $A \cdot 2^n + B(-1)^n$。 $g(n) = 1$，$1$ 不是根（$1 \ne 1
+ 2$），猜常數 $p$：$p = p + 2p + 1$，$p = -1/2$。 所以
$C(n) = A \cdot 2^n + B(-1)^n - 1/2$。 現在才代初始值：$A + B - 1/2 = 1$、
$2A - B - 1/2 = 1$，得 $A = 1$、$B = 1/2$：
\[
C(n) = 2^n + \frac{(-1)^n - 1}{2},
\]
$n$ 是偶數時剛好是 $2^n$，奇數時是 $2^n - 1$。

\subsection{實際數一次：鋪地磚}

\begin{example}[label=ex:rec-tiles]{c}{(tiles.c) 鋪地磚的方法數和呼叫次數}
#include <stdio.h>

long long calls;    // tile 被呼叫的次數

// 長 n 的走道有幾種鋪法。磚有兩種長度：長 1 的只有灰色，
// 長 2 的有紅、藍兩色。看走道最右邊那一塊是哪一種。
long long tile(int n)
{
    calls++;
    if (n <= 1) {
        return 1;                   // 長 0：什麼都不鋪；長 1：一塊灰磚
    }
    long long ways = tile(n - 1);   // 最右邊是灰磚
    for (int color = 0; color < 2; color++) {
        ways += tile(n - 2);        // 最右邊是紅磚或藍磚
    }
    return ways;
}

int main(void)
{
    int ns[] = {0, 1, 2, 3, 4, 5, 6, 7, 8, 12, 16, 20};
    printf("n ways closed calls closed\n");
    for (int t = 0; t < 12; t++) {
        int n = ns[t];
        long long sign = (n % 2 == 0) ? 1 : -1;     // (-1)^n
        calls = 0;
        long long w = tile(n);
        printf("%d %lld %lld %lld %lld\n", n, w, ((2LL << n) + sign) / 3,
               calls, (1LL << n) + (sign - 1) / 2);
    }
    return 0;
}
\end{example}
\iosample{}{n ways closed calls closed \\ 0 1 1 1 1 \\ 1 1 1 1 1 \\ 2 3 3 4 4 \\ 3 5 5 7 7 \\ 4 11 11 16 16 \\ 5 21 21 31 31 \\ 6 43 43 64 64 \\ 7 85 85 127 127 \\ 8 171 171 256 256 \\ 12 2731 2731 4096 4096 \\ 16 43691 43691 65536 65536 \\ 20 699051 699051 1048576 1048576}

第二、三欄是方法數和例 1 的公式，第四、五欄是呼叫次數和例 3 的公式，每一列都相同。
這支程式的時間是 $\Theta(2^n)$：同一個 \texttt{tile(k)} 被重算了很多次。 用《C 程式
設計基礎》第~15~章 15.9 節的記憶化，或從小到大填一個陣列，呼叫次數可以降到
$\Theta(n)$。

\section{生成函數與 Akra--Bazzi 方法}
\label{sec:rec-gf}

這一節介紹兩個更一般的工具，只給結論和一個例子。

\subsection{生成函數}

數列 $a_0, a_1, a_2, \ldots$ 的\textbf{生成函數}（generating function）是冪級數
$A(x) = \sum_{n \ge 0} a_n x^n$。 遞迴式乘上 $x^n$ 再對 $n$ 加總，會變成 $A(x)$ 的
一個方程式。 用鋪地磚的 $a_n = a_{n-1} + 2a_{n-2}$（$a_0 = a_1 = 1$）示範：
\[
A(x) - x A(x) - 2x^2 A(x) = a_0 + (a_1 - a_0) x = 1 .
\]
（左邊 $x^n$ 的係數，$n \ge 2$ 時是 $a_n - a_{n-1} - 2a_{n-2} = 0$；$n = 0, 1$ 時剩
下右邊兩項。） 所以
\[
A(x) = \frac{1}{1 - x - 2x^2} = \frac{1}{(1 - 2x)(1 + x)}
= \frac{2/3}{1 - 2x} + \frac{1/3}{1 + x}.
\]
再用 $\frac{1}{1 - rx} = \sum_n r^n x^n$ 展開，$x^n$ 的係數是
$\frac{2}{3} 2^n + \frac{1}{3}(-1)^n$，和例 1 一樣。 分母
$1 - c_1 x - c_2 x^2$ 拆成 $(1 - r_1 x)(1 - r_2 x)$ 時，$r_1$、$r_2$ 正是特徵方程式
的根，所以對線性遞迴，兩種方法是同一件事。 生成函數的好處是還能處理計數問題裡很
多別的形式，例如第~7~章二元樹的個數（卡特蘭數）。

\subsection{Akra--Bazzi 方法}

子問題大小不同的遞迴式，例如
$T(n) = T(n/4) + T(3n/4) + n$，主定理用不上。 Akra 和 Bazzi 在 1998 年給了一個一般的
公式。

\begin{theorem}[Akra--Bazzi，只給結論]
\label{thm:rec-akra}
設 $T(n) = f(n) + \sum_{i=1}^{k} a_i T(n/b_i)$，$a_i > 0$、$b_i > 1$ 都是常數，
$f(n)$ 是非負函數，而且 $n$ 變成常數倍時 $f(n)$ 只變成常數倍（例如 $n^k$、
$n^k \lg n$ 可以，$2^n$ 不行）。 令 $p$ 是滿足 $\sum_{i=1}^{k} a_i / b_i^{\,p} = 1$
的實數（只有一個），則
\[
T(n) = \Theta\!\left(n^{p} \left(1 + \int_{1}^{n} \frac{f(u)}{u^{p+1}}\, du\right)\right).
\]
\end{theorem}

\textbf{例。} $T(n) = T(n/4) + T(3n/4) + n$：$\frac{1}{4^p} + \frac{1}{(4/3)^p} = 1$，
$p = 1$ 時 $\frac{1}{4} + \frac{3}{4} = 1$，所以 $p = 1$。
$\int_1^n \frac{u}{u^2}\, du = \ln n$，所以 $T(n) = \Theta(n(1 + \ln n)) =
\Theta(n \lg n)$，和第~\ref{sec:rec-tree}~節用遞迴樹得到的一樣。 如果改成
$T(n/2) + T(n/3) + n$，$\frac{1}{2} + \frac{1}{3} < 1$，所以 $p < 1$，積分是
$\Theta(n^{1-p})$，答案是 $\Theta(n)$：每一層比上一層少，根的工作最大。

這個公式對地板、天花板也不敏感（把 $n/b_i$ 換成 $\lfloor n/b_i \rfloor$ 結論不變）。
$k = 1$ 時它包含主定理的大部分情形。

\section{本章常見錯誤}
\label{sec:rec-pitfalls}

\begin{center}
\begin{tabular}{p{0.24\linewidth}p{0.50\linewidth}p{0.16\linewidth}}
\toprule
錯誤 & 例子 & 見 \\
\midrule
漏掉呼叫的次數 & 呼叫兩次 $T(n/2)$ 卻寫成 $T(n/2) + n$；兩種顏色各呼叫一次卻只寫
一個 $C(n-2)$ & 第~\ref{sec:rec-write}~節 \\
歸納時常數變大 & 從 $T(k) \le ck$ 推出 $T(n) \le cn + n$，就說「$= O(n)$」 &
第~\ref{sec:rec-subst}~節 \\
猜對了卻證不出來就放棄 & 差一個常數時改猜更大的答案，而不是減去低次項 &
第~\ref{sec:rec-subst}~節 \\
遞迴樹漏了葉子 & $3T(n/2) + n$ 只算內部節點，或層數寫錯 &
第~\ref{sec:rec-tree}~節 \\
不平衡的樹用一條路徑代表 & $T(n/4) + T(3n/4) + n$ 只看最短路徑就說是
$\Theta(n \log_4 n)$ 而沒有說明上界 & 第~\ref{sec:rec-tree}~節 \\
主定理只比大小 & 把 $n^2/\lg n$ 當成 $O(n^{2 - \varepsilon})$；把 $n \lg n$ 當成
$\Omega(n^{1+\varepsilon})$ & 第~\ref{sec:rec-master}~節 \\
情況 3 沒驗正則條件 & 只寫「$f$ 比較大，所以 $\Theta(f(n))$」 &
第~\ref{sec:rec-master}~節 \\
$\log_b a$ 算錯 & 把 $\log_4 2$ 算成 $2$；把 $\log_2 8$ 和 $\log_8 2$ 搞混 &
第~\ref{sec:rec-master}~節 \\
減法的遞迴套主定理 & $T(n) = 2T(n-1) + 1$ 說成 $\Theta(n)$ &
第~\ref{sec:rec-other}~節 \\
特徵方程式重根寫錯 & 重根 $r$ 寫成 $A r^n + B r^n$ & 第~\ref{sec:rec-other}~節 \\
非齊次先代初始值 & 齊次解先用 $a_0$、$a_1$ 求出 $A$、$B$，再加上特解 &
第~\ref{sec:rec-other}~節 \\
\bottomrule
\end{tabular}
\end{center}

\section{觀念題}
\label{sec:rec-exercises}

題型和研究所考試、期中期末考的筆試相近，有選擇、計算、追蹤和證明。 沒有特別說明
時，$T(n)$ 在 $n$ 小於某個常數時是 $\Theta(1)$。 解答在本章最後。

\begin{enumerate}
\item （選擇，可能不只一個答案）下列哪些遞迴式可以\textbf{直接}套用主定理（含推廣
的情況 2）？
\begin{enumerate}
\item $T(n) = 4T(n/2) + n \lg n$
\item $T(n) = 8T(n/2) + n^3 / \lg n$
\item $T(n) = T(n/2) + 2^n$
\item $T(n) = 0.5\, T(n/2) + n$
\item $T(n) = 2T(n-1) + 1$
\item $T(n) = 6T(n/3) + n^2 \lg n$
\end{enumerate}

\item （計算）用主定理解下列遞迴式。 寫出 $a$、$b$、$\log_b a$、情況，情況 3 要驗
證正則條件。
\begin{enumerate}
\item $T(n) = 7T(n/3) + n^2$
\item $T(n) = 16T(n/2) + n^4$
\item $T(n) = 10T(n/3) + n^2$
\item $T(n) = 2T(n/8) + \sqrt[3]{n}$
\item $T(n) = 3T(n/9) + \sqrt{n} \lg n$
\end{enumerate}

\item （追蹤）\texttt{x} 是全域變數，一開始是 $0$。
\begin{code}{c}{觀念題 3}
void g(int n)
{
    if (n < 4) {
        x++;
        return;
    }
    for (int i = 0; i < n; i += 2)
        x++;
    for (int k = 0; k < 5; k++)
        g(n / 4);
}
\end{code}
(a) 呼叫 \texttt{g(64)} 之後，\texttt{x} 是多少？ \texttt{g} 一共被呼叫幾次（含
\texttt{g(64)} 本身）？ (b) 寫出 \texttt{x++} 執行次數 $X(n)$ 的遞迴式。 (c) 求
$X(n)$ 的 $\Theta$。

\item （證明）$T(1) = 1$，$n \ge 2$ 時 $T(n) = 2T(\lfloor n/2 \rfloor) + n^2$。 用代
換法證明對所有 $n \ge 1$ 都有 $n^2 \le T(n) \le 2n^2$。

\item （證明）$n$ 是 $3$ 的次方，$T(1) = 1$，$T(n) = 3T(n/3) + 1$。 (a) 說明為什麼
用歸納法證明 $T(n) \le cn$ 會失敗。 (b) 減去一個低次項，證明 $T(n) \le cn - d$，求
出可用的 $c$、$d$，並說明 $T(n) = (3n - 1)/2$。

\item （找錯）下面的「證明」錯在哪裡？ $T(n) = T(n-1) + n$ 的正確答案是什麼？
\begin{quote}
設 $T(0) = 1$。 假設 $T(n-1) \le c(n - 1)$，則 $T(n) \le c(n-1) + n \le cn + n =
O(n)$。 所以 $T(n) = O(n)$。
\end{quote}

\item （計算）$T(n) = 2T(n/5) + T(n/3) + n$。 畫出遞迴樹的前三層，寫出每一層的總
和，猜出 $\Theta$，再用代換法證明上界。

\item （計算）用換變數解 $T(n) = 4T(\sqrt{n}) + \lg^2 n$，寫成 $\Theta$。

\item （計算）$T(1) = 0$，$T(n) = T(n-1) + n \lg n$。 求 $T(n)$ 的 $\Theta$，上界和
下界都要說明。

\item （計算）解下列遞迴式，寫出封閉形式。
\begin{enumerate}
\item $a_n = 5a_{n-1} - 6a_{n-2}$，$a_0 = 0$，$a_1 = 1$。
\item $a_n = 4a_{n-1} - 4a_{n-2}$，$a_0 = 2$，$a_1 = 2$。
\end{enumerate}

\item （計算）解下列遞迴式，寫出封閉形式。
\begin{enumerate}
\item $T(0) = 0$，$T(n) = 3T(n-1) + 2^n$。
\item $T(0) = 0$，$T(n) = 2T(n-1) + 2^n$。
\end{enumerate}

\item （證明）$R(1) = 0$，$n \ge 2$ 時 $R(n) = R(\lfloor n/2 \rfloor) + 1$。 證明
$R(n) = \lfloor \lg n \rfloor$。 再用它說明：虛擬碼~\ref{alg:rec-search} 在 $n \ge 1$
個數裡找一個比全部都大的 $x$ 時，一共呼叫 $\lfloor \lg n \rfloor + 2$ 次。

\item （是非）$T(1) = 1$，$n \ge 2$ 時 $T(n) = 2T(\lfloor n/2 \rfloor) + 1$。 已知
$n$ 是 $2$ 的次方時 $T(n) = 2n - 1$。 下列敘述對的打 O，錯的打 X，並說明理由。
\begin{enumerate}
\item 對所有 $n \ge 1$，$T(n) = 2n - 1$。
\item 對所有 $n \ge 1$，$T(n) \le 2n - 1$。
\item 對所有 $n \ge 1$，$T(n) \ge n$。
\item 對所有 $n$，$T(n) = \Theta(n)$。
\end{enumerate}

\item （計算）河內塔的步數 $a_0 = 0$、$a_n = 2a_{n-1} + 1$。 用生成函數求 $a_n$。

\item （計算）用 Akra--Bazzi 方法解 $T(n) = 2T(n/4) + T(n/2) + n \lg n$。
\end{enumerate}

\section{實作練習}
\label{sec:rec-practice}

下面是 Aki-Judge 上和本章相關的題目。 這一章以手算和證明為主，能寫成程式題的不
多，練習請以觀念題為主。

\begin{center}
\begin{tabular}{llc}
\toprule
小節 & 題目 & 難度 \\
\midrule
每次減一 & \akiref{evaluate-a-linear-recurrence}{計算線性遞迴式的值} & 暖身 \\
從程式寫出遞迴式 & \akiref{counting-calls-between-two-recursive-functions}{互相呼叫的呼叫次數} & 基本 \\
\bottomrule
\end{tabular}
\end{center}

「計算線性遞迴式的值」可以照遞迴式用迴圈算，也可以先用第~\ref{sec:rec-other}~節的
方法求出封閉形式，再和程式的答案比對。 「互相呼叫的呼叫次數」裡兩個函式互相呼叫，
用範例~\ref{ex:rec-audit} 的全域計數器寫法；交之前先對小的 $n$ 手算呼叫樹，和程式
的計數比對。

\section{讀完你應該會}
\label{sec:rec-checklist}

\begin{itemize}
\item 看到一個遞迴程式，能寫出它的執行時間或呼叫次數的遞迴式，包括基礎情況、呼叫
次數和每次的規模。
\item 能用代換法證明遞迴式的上下界，寫出常數；歸納過不去時會減去低次項；會用換變
數處理 $\sqrt{n}$。
\item 能畫出遞迴樹，寫出每一層的總和、層數、葉子數，判斷是根、每層平手還是葉子決
定答案；能處理兩個子問題不一樣大的情形。
\item 能用主定理解 $aT(n/b) + f(n)$，寫出情況、$\varepsilon$ 和正則條件，並認出不
能用的遞迴式。
\item 能解 $T(n-1) + f(n)$ 型的遞迴式，以及用特徵方程式解二階常係數線性遞迴（不同
根、重根、非齊次）。
\end{itemize}

\textbf{延伸閱讀}（書末會整理完整的參考資料）

\begin{itemize}
\item Thomas H. Cormen 等，\textit{Introduction to Algorithms}，第 4 版，4.3--4.7 節：
代換法、遞迴樹、主定理（證明包含連續版本）和 Akra--Bazzi 方法；這一版的主定理情
況 2 已經是本章定理~\ref{thm:rec-master2} 的形式。
\item Steven S. Skiena，\textit{The Algorithm Design Manual}，第 3 版，5.3--5.4 節：
用很短的篇幅講分治的遞迴式和主定理三種情況的直覺。
\item Ronald L. Graham、Donald E. Knuth、Oren Patashnik，\textit{Concrete
Mathematics}，第 2 版，第 1 章與第 7 章：前者用幾個經典問題示範怎麼找出遞迴式的封閉
形式，後者完整講生成函數，適合想把第~\ref{sec:rec-gf}~節讀深的人。
\end{itemize}

\addcontentsline{toc}{section}{觀念題解答}
\section*{觀念題解答}

\begin{enumerate}
\item (a)、(c)、(f)。 (a)：$n^{\log_2 4} = n^2$，$n \lg n = O(n^{1.5})$，情況 1，
$\Theta(n^2)$。 (b)：$n^{\log_2 8} = n^3$，$f(n) = n^3 / \lg n$ 只小了 $\lg n$，不是
情況 1，也不是推廣的情況 2（$k = -1$）；用遞迴樹可算出 $\Theta(n^3 \lg \lg n)$。
(c)：$n^{\log_2 1} = 1$，$2^n = \Omega(n^{1})$，正則條件 $2^{n/2} \le \frac{1}{2} 2^n$
在 $n \ge 2$ 時成立，情況 3，$\Theta(2^n)$。 (d)：$a = 0.5 < 1$，不符合定理的假設
（答案碰巧是 $\Theta(n)$，但不是主定理給的）。 (e)：規模是減 $1$，不是除以 $b$。
(f)：$\log_3 6 \approx 1.63$，$n^2 \lg n = \Omega(n^{1.63 + 0.3})$，正則條件
$6 (n/3)^2 \lg (n/3) = \frac{2}{3} n^2 \lg(n/3) \le \frac{2}{3} n^2 \lg n$，$c = 2/3$，
情況 3，$\Theta(n^2 \lg n)$。

\item
\begin{enumerate}
\item $a = 7$、$b = 3$，$\log_3 7 \approx 1.77$。 $n^2 = \Omega(n^{1.77 + 0.2})$，
正則條件 $7(n/3)^2 = \frac{7}{9} n^2$，$c = 7/9$。 情況 3，$\Theta(n^2)$。
\item $a = 16$、$b = 2$，$\log_2 16 = 4$。 $f(n) = n^4$，情況 2，$\Theta(n^4 \lg n)$。
\item $a = 10$、$b = 3$，$\log_3 10 \approx 2.096$。 $n^2 = O(n^{2.096 - 0.09})$，情況
1，$\Theta(n^{\log_3 10})$。 注意 $\log_3 9 = 2$，多一個子問題就讓答案超過 $n^2$。
\item $a = 2$、$b = 8$，$\log_8 2 = 1/3$。 $f(n) = n^{1/3}$，情況 2，
$\Theta(\sqrt[3]{n} \lg n)$。
\item $a = 3$、$b = 9$，$\log_9 3 = 1/2$。 $f(n) = n^{1/2} \lg n$，推廣的情況 2，
$k = 1$，$\Theta(\sqrt{n} \lg^2 n)$。
\end{enumerate}

\item (a) 由下往上算。 $n < 4$ 時 \texttt{x++} 一次。 \texttt{g(4)}：迴圈
$i = 0, 2$ 兩次，加上 $5$ 次 \texttt{g(1)}，共 $2 + 5 = 7$。 \texttt{g(16)}：
$8 + 5 \cdot 7 = 43$。 \texttt{g(64)}：$32 + 5 \cdot 43 = 247$，所以 \texttt{x} 是
$247$。 呼叫次數 $1 + 5 + 25 + 125 = 156$。 (b) $X(n) = 1$（$n < 4$），
$X(n) = 5X(\lfloor n/4 \rfloor) + \lceil n/2 \rceil$（$n \ge 4$）；迴圈
$i = 0, 2, 4, \ldots$ 小於 $n$ 的有 $\lceil n/2 \rceil$ 個。 (c) $n$ 是 $4$ 的次方
時，$a = 5$、$b = 4$，$\log_4 5 \approx 1.16$，$f(n) = n/2 = O(n^{1.16 - 0.1})$，
情況 1，$\Theta(n^{\log_4 5})$。 $X$ 是非遞減的，$n^{\log_4 5}$ 滿足
定理~\ref{thm:rec-smooth} 的條件，所以對所有 $n$ 都是 $\Theta(n^{\log_4 5})$。

\item 對 $n$ 做強歸納。 $n = 1$：$T(1) = 1$，$1 \le 1 \le 2$。 $n \ge 2$：
$1 \le \lfloor n/2 \rfloor < n$，用歸納假設。 上界：
$T(n) \le 2 \cdot 2\lfloor n/2 \rfloor^2 + n^2 \le 4 (n/2)^2 + n^2 = 2n^2$。 下界：
$T(\lfloor n/2 \rfloor) \ge 0$，所以 $T(n) \ge n^2$。 因此 $T(n) = \Theta(n^2)$，
$c_1 = 1$、$c_2 = 2$、$n_0 = 1$。

\item (a) 假設 $T(n/3) \le c\,n/3$，得到 $T(n) \le 3 \cdot c\,n/3 + 1 = cn + 1$，比要
證的 $cn$ 多 $1$，不管 $c$ 多大都推不出來。 (b) 假設 $T(n/3) \le c\,n/3 - d$：
$T(n) \le 3(c\,n/3 - d) + 1 = cn - 3d + 1 \le cn - d$，只要 $d \ge 1/2$。 基礎情況
$T(1) = 1 \le c - d$。 取 $d = 1/2$、$c = 3/2$，得 $T(n) \le (3n - 1)/2$。 把每個
$\le$ 換成 $=$ 再做一次：$T(1) = 1 = (3 - 1)/2$，
$3 \cdot \frac{3(n/3) - 1}{2} + 1 = \frac{3n - 1}{2}$，所以 $T(n) = (3n - 1)/2$。

\item 錯在「$cn + n = O(n)$」。 歸納要證明的是 $T(n) \le cn$，常數必須和假設的
$c$ 相同；推出來的是 $(c + 1) n$，歸納沒有完成。 每減一，常數就多 $1$，$n$ 步之後
「常數」是 $c + n$。 正確的答案是展開：
$T(n) = T(0) + \sum_{i=1}^{n} i = 1 + \frac{n(n+1)}{2} = \Theta(n^2)$。

\item 第 0 層：$n$。 第 1 層：兩個 $n/5$、一個 $n/3$，總和
$(\frac{2}{5} + \frac{1}{3}) n = \frac{11}{15} n$。 第 2 層：每個節點又分成總和是自
己 $\frac{11}{15}$ 倍的三塊，所以總和是 $(\frac{11}{15})^2 n = \frac{121}{225} n$。
每層是上一層的 $\frac{11}{15} < 1$ 倍，總和不超過
$n \sum_j (\frac{11}{15})^j = \frac{15}{4} n$，猜 $\Theta(n)$。 下界：
$T(n) \ge n$。 上界用代換法：假設比 $n$ 小的值都有 $T(k) \le ck$，則
$T(n) \le 2c\,\frac{n}{5} + c\,\frac{n}{3} + n = \frac{11}{15} cn + n \le cn$，只要
$c \ge \frac{15}{4}$；$c$ 再取大一點讓有限個基礎情況也成立。 所以
$T(n) = \Theta(n)$。

\item 令 $m = \lg n$，$S(m) = T(2^m)$。 $\sqrt{n} = 2^{m/2}$，$\lg^2 n = m^2$，所以
$S(m) = 4S(m/2) + m^2$。 $\log_2 4 = 2$，情況 2，$S(m) = \Theta(m^2 \lg m)$。 換回
來：$T(n) = \Theta(\lg^2 n \cdot \lg \lg n)$。

\item 展開得 $T(n) = \sum_{i=2}^{n} i \lg i$。 上界：每一項不超過 $n \lg n$，共不到
$n$ 項，$T(n) \le n^2 \lg n$。 下界：只留 $i \ge n/2$ 的項，至少 $n/2$ 項，每項至少
$\frac{n}{2} \lg \frac{n}{2}$，所以 $T(n) \ge \frac{n^2}{4}(\lg n - 1)$；$n \ge 4$ 時
$\lg n - 1 \ge \frac{1}{2} \lg n$，$T(n) \ge \frac{1}{8} n^2 \lg n$。 所以
$T(n) = \Theta(n^2 \lg n)$。

\item
\begin{enumerate}
\item 特徵方程式 $x^2 - 5x + 6 = (x - 2)(x - 3)$，$a_n = A \cdot 2^n + B \cdot 3^n$。
$A + B = 0$、$2A + 3B = 1$，得 $B = 1$、$A = -1$，$a_n = 3^n - 2^n$。 驗算
$a_2 = 5 - 0 = 5 = 9 - 4$。
\item $x^2 - 4x + 4 = (x - 2)^2$，重根 $2$，$a_n = (A + Bn) 2^n$。 $A = 2$，
$(2 + B) \cdot 2 = 2$，$B = -1$，$a_n = (2 - n) 2^n$。 驗算
$a_2 = 8 - 8 = 0 = (2 - 2) \cdot 4$。
\end{enumerate}

\item
\begin{enumerate}
\item 齊次部分 $A \cdot 3^n$。 $2$ 不是根，猜 $p \cdot 2^n$：
$p \cdot 2^n = 3p \cdot 2^{n-1} + 2^n$，兩邊除以 $2^{n-1}$：$2p = 3p + 2$，$p = -2$。
$T(n) = A \cdot 3^n - 2^{n+1}$，$T(0) = A - 2 = 0$，$A = 2$。
$T(n) = 2 \cdot 3^n - 2^{n+1}$。 驗算 $T(1) = 0 + 2 = 2 = 6 - 4$。
\item 齊次部分 $A \cdot 2^n$，$2$ 是根，猜 $p \cdot 2^n$ 會被吃掉（代進去得 $0 = 2^n$），
改猜 $p\, n 2^n$：$p n 2^n = 2p(n-1) 2^{n-1} + 2^n$，除以 $2^n$：$pn = p(n - 1) + 1$，
$p = 1$。 $T(n) = A \cdot 2^n + n 2^n$，$T(0) = A = 0$，$T(n) = n 2^n$。 驗算
$T(1) = 0 + 2 = 2$，$T(2) = 4 + 4 = 8 = 2 \cdot 4$。
\end{enumerate}

\item 對 $n$ 做強歸納。 $R(1) = 0 = \lfloor \lg 1 \rfloor$。 $n \ge 2$ 時
$1 \le \lfloor n/2 \rfloor < n$，由歸納假設和引理~\ref{lem:rec-lg-half}：
$R(n) = \lfloor \lg \lfloor n/2 \rfloor \rfloor + 1 = (\lfloor \lg n \rfloor - 1) + 1
= \lfloor \lg n \rfloor$。

二分搜尋：$x$ 比全部都大，每次 $A[q] < x$，往右邊 $A[q+1..r]$ 找，規模從 $n$ 變成
$\lfloor n/2 \rfloor$，直到規模 $0$ 的那次呼叫回傳 \NIL。 令 $W(n)$ 是呼叫次數，
$W(0) = 1$，$W(n) = W(\lfloor n/2 \rfloor) + 1$（$n \ge 1$），所以 $W(1) = 2$。
$n \ge 1$ 時，$W(n) - 2$ 和 $R(n)$ 滿足同一個遞迴式、同一個基礎情況
（$W(1) - 2 = 0$），所以 $W(n) = \lfloor \lg n \rfloor + 2$。

\item 先求出精確解：$T(n) = 2^{\lfloor \lg n \rfloor + 1} - 1$。 （歸納：$n \ge 2$ 時
$T(n) = 2(2^{\lfloor \lg n \rfloor} - 1) + 1$，用了引理~\ref{lem:rec-lg-half}。）
令 $k = \lfloor \lg n \rfloor$，$2^k \le n < 2^{k+1}$。
\begin{enumerate}
\item X。 $T(3) = 2T(1) + 1 = 3 \ne 5$。 $n$ 是 $2$ 的次方時的公式不一定對其他 $n$ 成
立。
\item O。 $2^{k+1} \le 2n$，所以 $T(n) \le 2n - 1$。
\item O。 $2^{k+1} > n$，而且是整數，所以 $2^{k+1} \ge n + 1$，$T(n) \ge n$。
\item O。 由 (b)、(c)。 也可以用定理~\ref{thm:rec-smooth}：$T$ 非遞減，$g(n) = n$。
\end{enumerate}

\item 令 $A(x) = \sum_{n \ge 0} a_n x^n$。 $n \ge 1$ 時 $a_n - 2a_{n-1} = 1$，$a_0 = 0$，
所以 $A(x) - 2xA(x) = \sum_{n \ge 1} x^n = \frac{x}{1 - x}$。
\[
A(x) = \frac{x}{(1 - x)(1 - 2x)} = \frac{1}{1 - 2x} - \frac{1}{1 - x}
\]
（通分驗算：$(1 - x) - (1 - 2x) = x$）。 $x^n$ 的係數是 $2^n - 1$，所以
$a_n = 2^n - 1$。

\item $a_1 = 2$、$b_1 = 4$、$a_2 = 1$、$b_2 = 2$。 找 $p$：
$\frac{2}{4^p} + \frac{1}{2^p} = 1$。 令 $y = 2^{-p}$，$2y^2 + y - 1 = 0$，
$(2y - 1)(y + 1) = 0$，$y = 1/2$，$p = 1$。 積分
$\int_1^n \frac{u \lg u}{u^2}\, du = \int_1^n \frac{\lg u}{u}\, du =
\frac{(\ln n)^2}{2 \ln 2} = \Theta(\lg^2 n)$。 所以
$T(n) = \Theta(n (1 + \lg^2 n)) = \Theta(n \lg^2 n)$。 （每一層的規模加起來是
$\frac{2}{4} + \frac{1}{2} = 1$ 倍，和第~\ref{sec:rec-tree}~節的「每層一樣」同一型，
只是每層的成本帶著 $\lg$。）
\end{enumerate}

\chapter{平攤分析}
\label{ch:amortized}

% 本章的圖共用的樣式
\tikzset{
    ctrcell/.style={cell, minimum width=0.75cm, minimum height=0.6cm},
    flipcell/.style={hlcell, minimum width=0.75cm, minimum height=0.6cm},
    costwrite/.style={draw=coreStroke, fill=coreFill},
    costcopy/.style={draw=memStroke, fill=memFill},
}
% 計數器追蹤圖的一列：\ctrrow{y}{左邊的說明}{位元/是否翻轉, ...}{右邊的說明}
\newcommand{\ctrrow}[4]{%
    \node[anchor=east, font=\small] at (0.3, #1) {#2};
    \foreach \b/\h [count=\k] in {#3} {
        \ifnum\h=1\relax
            \node[flipcell] at (\k*0.85, #1) {\b};
        \else
            \node[ctrcell] at (\k*0.85, #1) {\b};
        \fi
    }
    \node[anchor=west, font=\small] at (3.9, #1) {#4};
}

班代用一本筆記本記班費，每收一筆寫一行。 本子寫滿了，他就買一本行數是舊本子兩倍的
新本子，把舊帳一行一行抄過去，再接著記。 抄帳那天特別累，例如記第 1025 筆之前，
要先抄 1024 行。 這學期一共收了大約一千筆，平均下來每收一筆，他總共要寫幾行？ 會
不會因為偶爾要抄一大本，平均起來就變得很多？ 往下讀之前，先猜一個數。

上一章分析的是「一次呼叫」的成本。 這一章換一個角度：資料結構上的操作通常一次做
很多個，其中少數幾次很貴，大部分很便宜。 只看最貴的那一次再乘上次數，得到的上界
往往太鬆。 \textbf{平攤分析}（amortized analysis）改成證明「任何一串 $n$ 次操作的總
成本」的上界，再把它平均分給每一次操作。 這裡的平均不涉及機率：它對最壞的那一串操
作也成立。

這一章介紹三種做法：聚合法、記帳法、位能法，都用同兩個例子（可以一次拿掉很多個的
堆疊、二進位計數器）示範，最後用它們分析動態陣列，回答章首的問題。 後面的章會再
用到這些工具：第~8~章費氏堆積的各項成本、第~12~章雜湊表的重新配置、第~13~章互斥
集合、第~27~章 KMP 演算法的線性時間，都是平攤的結論。

這一章用到《C 程式設計基礎》第~14~章的結構指標、第~19~章的 \texttt{realloc} 和第~20~章
用陣列寫的堆疊，用到時會先在框裡複習。

\section{聚合法}
\label{sec:amo-aggregate}

\subsection{可以一次拿掉很多個的堆疊}

堆疊（stack）是後放進去的先拿出來的資料結構，第~6~章會完整介紹。 這一節只需要三
個操作：

\begin{itemize}
\item \proc{Push}$(S, x)$：把 $x$ 放到堆疊 $S$ 的最上面。
\item \proc{Pop}$(S)$：拿掉最上面的元素並回傳它。 呼叫前要確定 $S$ 不是空的。
\item \proc{Multipop}$(S, k)$：一次拿掉最上面的 $k$ 個；不到 $k$ 個就全部拿掉。
\end{itemize}

前兩個操作都花常數時間，記成成本 $1$。 \proc{Multipop} 的寫法如下，堆疊目前有 $s$
個元素時，迴圈做 $\min(s, k)$ 次，成本記成 $\min(s, k)$。

\begin{pseudo}[label=alg:multipop]{一次拿掉 $k$ 個}
\Procedure{Multipop}{$S, k$}
    \While{\textbf{not} \Call{Stack-Empty}{$S$} \AND $k > 0$}
        \State \Call{Pop}{$S$}
        \State $k \gets k - 1$
    \EndWhile
\EndProcedure
\end{pseudo}

現在從空的堆疊開始，做 $n$ 次操作，每次是三種之一。 總成本是多少？

照第~1~章的方法，先找每一次操作的最壞情況。 堆疊最多有 $n$ 個元素，一次
\proc{Multipop} 最多拿掉 $n$ 個，成本 $O(n)$；$n$ 次操作，總成本 $O(n^2)$。 這個上
界是對的，但太鬆了。 要讓某一次 \proc{Multipop} 拿掉很多個，前面就得先花很多次
\proc{Push} 把它們放進去，而放進去的元素只能被拿掉一次。

把這句話寫清楚：每個元素被 \proc{Pop}（包括 \proc{Multipop} 裡面的 \proc{Pop}）拿
掉之前，一定被 \proc{Push} 放進去過，而且只會被拿掉一次。 所以整串操作裡，
\proc{Pop} 實際執行的總次數不超過 \proc{Push} 的次數，而後者不超過 $n$。 總成本
\[
\text{\proc{Push} 的次數} + \text{\proc{Pop} 實際執行的次數} \le n + n = 2n.
\]
任何 $n$ 次操作的總成本都是 $O(n)$，平均每次 $O(1)$。

這種做法叫\textbf{聚合法}（aggregate analysis）：不去估每一次操作，直接對整串操作
算出總成本的上界 $T(n)$，再說每次操作的平攤成本是 $T(n)/n$。 聚合法給所有種類的操
作同一個平攤成本，這裡三種操作的平攤成本都是 $2n/n = 2$。

\begin{definition}[平攤成本]
\label{def:amortized}
一個資料結構從初始狀態開始，執行一串操作，第 $i$ 次操作的實際成本是 $c_i$。 如果
給每一次操作指定一個數 $\hat{c}_i$，使得對\textbf{每一串}操作、每一個 $n$ 都有
\[
\sum_{i=1}^{n} c_i \le \sum_{i=1}^{n} \hat{c}_i,
\]
就說 $\hat{c}_i$ 是第 $i$ 次操作的\textbf{平攤成本}（amortized cost）。
\end{definition}

定義裡的「每一串」很重要：平攤成本的上界對最壞的操作順序也成立，所以它仍然是一種
最壞情況的保證，只是保證的對象是總和。 它和第~1~章的平均情況不同，平均情況要先假
設輸入的機率分布，算的是期望值；平攤分析不用任何機率。

它也不保證單一次操作很快。 上面的堆疊裡，某一次 \proc{Multipop} 還是可能花 $n - 1$
的時間。 如果應用要求每一次回應都很快（例如控制機器的程式），只有平攤上界是不夠的。

\begin{cbookrecap}[第~14~章 14.5 節、第~20~章 20.8 節]
\texttt{Stack s = \{malloc(n * sizeof(int)), 0\};} 用大括號依欄位的順序初始化一個
結構。 函式要改呼叫端的結構時傳結構的位址，在函式裡用 \texttt{s->size} 存取欄位，
它是 \texttt{(*s).size} 的簡寫。 第~20~章用同樣的寫法把陣列和堆疊頂端的位置包成一
個結構，這一章的堆疊沿用那個做法。 字串比較用 \texttt{strcmp}，相等時回傳 $0$
（第~11~章）。
\end{cbookrecap}

下面的程式讀進一串操作，印出每一次的實際成本和堆疊的大小。 輸入只有
\texttt{push} 和 \texttt{multipop} 兩種，\proc{Pop} 就是 \texttt{multipop 1}。

\begin{example}[label=ex:multipop]{c}{(multipop.c) 每次操作的實際成本}
#include <stdio.h>
#include <stdlib.h>
#include <string.h>

// 用陣列做的堆疊。n 次操作最多 push n 次，容量開 n 格就夠。
typedef struct {
    int *data;
    int size;
} Stack;

// 一般的 pop，實際成本 1。呼叫前要確定堆疊不是空的。
int pop(Stack *s)
{
    s->size--;
    return s->data[s->size];
}

// 一次 pop 掉 k 個，堆疊不夠 k 個就全部 pop 掉。回傳這次 pop 了幾個。
int multipop(Stack *s, int k)
{
    int popped = 0;
    while (s->size > 0 && k > 0) {
        pop(s);
        popped++;
        k--;
    }
    return popped;
}

int main(void)
{
    int n;
    if (scanf("%d", &n) != 1 || n <= 0) {
        return 1;
    }
    Stack s = {malloc(n * sizeof(int)), 0};
    if (s.data == NULL) {
        return 1;
    }
    long long total = 0;
    int pushes = 0;
    char op[16];
    for (int i = 1; i <= n; i++) {
        int x, cost;
        if (scanf("%15s %d", op, &x) != 2) {
            free(s.data);
            return 1;
        }
        if (strcmp(op, "push") == 0) {
            s.data[s.size++] = x;
            pushes++;
            cost = 1;
        } else {                       // multipop x
            cost = multipop(&s, x);
        }
        total += cost;
        printf("%d %s %d: cost %d, size %d\n", i, op, x, cost, s.size);
    }
    printf("total %lld, pushes %d, 2n = %d\n", total, pushes, 2 * n);
    free(s.data);
    return 0;
}
\end{example}
\iosample{10 \\ push 5 \\ push 8 \\ push 3 \\ multipop 2 \\ push 7 \\ push 1 \\ push 4 \\ push 6 \\ multipop 10 \\ push 2}{1 push 5: cost 1, size 1 \\ 2 push 8: cost 1, size 2 \\ 3 push 3: cost 1, size 3 \\ 4 multipop 2: cost 2, size 1 \\ 5 push 7: cost 1, size 2 \\ 6 push 1: cost 1, size 3 \\ 7 push 4: cost 1, size 4 \\ 8 push 6: cost 1, size 5 \\ 9 multipop 10: cost 5, size 0 \\ 10 push 2: cost 1, size 1 \\ total 15, pushes 8, 2n = 20}

第 9 次操作一次花了 $5$，但在它之前已經有 $5$ 次便宜的 \proc{Push} 把這 $5$ 個元素
放進去。 總成本 $15$ 剛好是 \proc{Push} 的 $8$ 次加上 \proc{Pop} 的 $7$ 次，在
$2n = 20$ 以內。

\subsection{二進位計數器}
\label{sec:amo-counter}

第二個例子是一個 $k$ 位元的二進位計數器，用陣列 $A[1..k]$ 存，$A[1]$ 是最低位。 計
數器表示的數是 $\sum_{i=1}^{k} A[i] \cdot 2^{i-1}$，一開始每一位都是 $0$。 唯一的
操作是加一：

\begin{pseudo}[label=alg:increment]{計數器加一}
\Procedure{Increment}{$A, k$}
    \State $i \gets 1$
    \While{$i \le k$ \AND $A[i] = 1$}
        \State $A[i] \gets 0$ \Comment{尾端連續的 1 變成 0}
        \State $i \gets i + 1$
    \EndWhile
    \If{$i \le k$}
        \State $A[i] \gets 1$ \Comment{第一個 0 變成 1}
    \EndIf
\EndProcedure
\end{pseudo}

全部都是 $1$ 時，$i$ 會走到 $k + 1$，每一位都變成 $0$，計數器回到 $0$，這叫\textbf{溢
位}（overflow）。 一次 \proc{Increment} 的成本是它翻轉（$0$ 變 $1$ 或 $1$ 變 $0$）
的位元個數。 下圖是 $4$ 位元的計數器從 $0$ 加到 $8$ 的過程，左邊是最高位
$A[4]$，橘框是這一次翻轉的位元。

\begin{center}
\begin{tikzpicture}[yscale=0.75]
\node[font=\small\bfseries, anchor=west] at (3.9, 0.9) {成本　累計};
\foreach \i [count=\k] in {4,3,2,1} { \node[idxlbl] at (\k*0.85, 0.9) {$A[\i]$}; }
\ctrrow{0}{$0$}{0/0,0/0,0/0,0/0}{}
\ctrrow{-1}{$1$}{0/0,0/0,0/0,1/1}{$1$\qquad $1$}
\ctrrow{-2}{$2$}{0/0,0/0,1/1,0/1}{$2$\qquad $3$}
\ctrrow{-3}{$3$}{0/0,0/0,1/0,1/1}{$1$\qquad $4$}
\ctrrow{-4}{$4$}{0/0,1/1,0/1,0/1}{$3$\qquad $7$}
\ctrrow{-5}{$5$}{0/0,1/0,0/0,1/1}{$1$\qquad $8$}
\ctrrow{-6}{$6$}{0/0,1/0,1/1,0/1}{$2$\qquad $10$}
\ctrrow{-7}{$7$}{0/0,1/0,1/0,1/1}{$1$\qquad $11$}
\ctrrow{-8}{$8$}{1/1,0/1,0/1,0/1}{$4$\qquad $15$}
\end{tikzpicture}
\diagcaption{$4$ 位元計數器從 $0$ 加到 $8$。 每一列左邊是計數器的值，橘框是這次
\proc{Increment} 翻轉的位元，右邊是這次的成本和到目前為止的總成本。}
\label{fig:amo-counter}
\end{center}

一次 \proc{Increment} 最壞要翻轉全部 $k$ 位（例如從 $0111\cdots1$ 變成
$1000\cdots0$），所以 $n$ 次的總成本是 $O(nk)$。 但從圖上看得出來，$A[1]$ 每次都翻，
$A[2]$ 每兩次翻一次，$A[3]$ 每四次翻一次，越高的位越少動。 照每一位分開數，得到更
好的上界。

\begin{theorem}
\label{thm:counter-aggregate}
$k$ 位元計數器從 $0$ 開始執行 $n$ 次 \proc{Increment}，翻轉的位元總數是
$\sum_{i=1}^{k} \lfloor n / 2^{i-1} \rfloor$，小於 $2n$。 這個結論不管中間有沒有溢位
都成立。
\end{theorem}

\begin{proof}
第 $m$ 次 \proc{Increment} 把計數器從 $(m - 1) \bmod 2^k$ 變成 $m \bmod 2^k$。 它翻
轉的是最低位開始的一段：如果 $m$ 的二進位最低的 $j$ 位都是 $0$、第 $j + 1$ 位是
$1$，翻轉的就是 $A[1..j+1]$（$m$ 是 $2^k$ 的倍數時是溢位，翻轉全部 $k$ 位）。 換句
話說，$A[i]$ 在第 $m$ 次被翻轉，若且唯若 $2^{i-1}$ 整除 $m$。 $1$ 到 $n$ 之間
$2^{i-1}$ 的倍數有 $\lfloor n / 2^{i-1} \rfloor$ 個，所以
\[
\sum_{i=1}^{k} \left\lfloor \frac{n}{2^{i-1}} \right\rfloor
< \sum_{j=0}^{\infty} \frac{n}{2^{j}} = 2n.
\qedhere
\]
\end{proof}

所以 $n$ 次 \proc{Increment} 的總成本是 $O(n)$，每次的平攤成本是 $O(1)$，和位元數
$k$ 無關。 證明裡用的是第~1~章的等比級數：公比 $1/2$，總和和第一項同一級。

下面的程式實際數一次。 第一部分印出圖~\ref{fig:amo-counter} 的過程，第二部分用
$32$ 位元的計數器加到 $10^6$，第三部分印出 $n = 1000$ 時最低的 $10$ 位各翻了幾次。

\begin{example}[label=ex:counter-flips]{c}{(counter\_flips.c) 計數器加一的翻轉次數}
#include <stdio.h>
#include <stdlib.h>

#define K 32                    // 計數器的位元數

// 把 bit[0..K-1] 表示的數加一，bit[0] 是最低位。回傳翻轉了幾個位元。
int increment(int *bit)
{
    int i = 0, flips = 0;
    while (i < K && bit[i] == 1) {     // 尾端連續的 1 都變成 0
        bit[i] = 0;
        flips++;
        i++;
    }
    if (i < K) {                       // 第一個 0 變成 1（溢位時沒有這一步）
        bit[i] = 1;
        flips++;
    }
    return flips;
}

void clear(int *bit)
{
    for (int i = 0; i < K; i++) {
        bit[i] = 0;
    }
}

int main(void)
{
    int *bit = calloc(K, sizeof *bit);   // 一開始是 0
    if (bit == NULL) {
        return 1;
    }
    // 第一部分：前 8 次加一，印出低 4 位
    printf("step bits flips total\n");
    long long total = 0;
    for (int step = 1; step <= 8; step++) {
        int f = increment(bit);
        total += f;
        printf("%d %d%d%d%d %d %lld\n", step, bit[3], bit[2], bit[1], bit[0],
               f, total);
    }

    // 第二部分：從 0 開始加 n 次，總翻轉次數和 2n 比較
    int ns[] = {16, 1000, 1000000};
    printf("n total 2n\n");
    for (int k = 0; k < 3; k++) {
        clear(bit);
        total = 0;
        for (int step = 0; step < ns[k]; step++) {
            total += increment(bit);
        }
        printf("%d %lld %d\n", ns[k], total, 2 * ns[k]);
    }

    // 第三部分：n = 1000 時每一位各翻了幾次
    clear(bit);
    long long per_bit[K] = {0};
    for (int step = 0; step < 1000; step++) {
        int f = increment(bit);
        for (int i = 0; i < f; i++) {   // 沒有溢位時，翻的正是第 0 到第 f-1 位
            per_bit[i]++;
        }
    }
    printf("flips of bit 0..9:");
    for (int i = 0; i < 10; i++) {
        printf(" %lld", per_bit[i]);
    }
    printf("\n");
    free(bit);
    return 0;
}
\end{example}
\iosample{}{step bits flips total \\ 1 0001 1 1 \\ 2 0010 2 3 \\ 3 0011 1 4 \\ 4 0100 3 7 \\ 5 0101 1 8 \\ 6 0110 2 10 \\ 7 0111 1 11 \\ 8 1000 4 15 \\ n total 2n \\ 16 31 32 \\ 1000 1994 2000 \\ 1000000 1999993 2000000 \\ flips of bit 0..9: 1000 500 250 125 62 31 15 7 3 1}

C 程式的陣列從 $0$ 開始，\texttt{bit[j]} 對應虛擬碼的 $A[j + 1]$，所以它翻了
$\lfloor n / 2^j \rfloor$ 次。 第三部分的 $1000, 500, 250, \ldots, 1$ 正是
$\lfloor 1000 / 2^j \rfloor$，加起來 $1994$。 總數都比 $2n$ 少一點點：$2000 - 1994
= 6$，而 $1000$ 的二進位 $1111101000$ 剛好有 $6$ 個 $1$。 這不是巧合，觀念題第 11
題請你證明。

\begin{examnote}
題目問「平攤成本」或「平攤時間複雜度」時，答案要針對\textbf{任意一串} $n$ 次操作的
總成本論證，再除以 $n$。 不要寫「平均來說」「通常」「大部分的操作」這類機率的說
法，也不要只舉一串特定的操作算給人看：舉例只能說明那一串，不能說明所有的串。
\end{examnote}

\section{記帳法}
\label{sec:amo-accounting}

聚合法給每一種操作同一個平攤成本。 \textbf{記帳法}（accounting method）可以給不同
的操作不同的平攤成本，想法像預付：便宜的操作多收一點錢，多的部分存起來，以後貴的操
作來了，用存款付。

具體地說，給每一種操作訂一個收費 $\hat{c}_i$。 第 $i$ 次操作收費 $\hat{c}_i$、實際
花 $c_i$；$\hat{c}_i > c_i$ 時把差額存起來，存款通常想成「放在資料結構裡某些元素
身上」；$\hat{c}_i < c_i$ 時從存款提出差額。 前 $j$ 次操作之後的存款是
\[
\sum_{i=1}^{j} \hat{c}_i - \sum_{i=1}^{j} c_i.
\]
只要能證明\textbf{存款永遠不會變成負的}，也就是對所有 $j$ 這個式子都 $\ge 0$，收費
$\hat{c}_i$ 就符合定義~\ref{def:amortized}，是平攤成本。 記帳法的證明就是說明這一
件事。

\subsection{堆疊}

收費訂成：\proc{Push} 收 $2$，\proc{Pop} 和 \proc{Multipop} 收 $0$。 \proc{Push}
實際花 $1$，多收的 $1$ 存在剛放進去的那個元素身上。 之後這個元素被拿掉時（不管是
\proc{Pop} 還是 \proc{Multipop} 裡的一次），用它身上的 $1$ 付這次的成本。
\proc{Multipop} 拿掉幾個元素，就有幾個元素各自付 $1$，所以收 $0$ 也付得出來。

存款不會變負：堆疊裡每個元素身上都剛好有 $1$，存款等於堆疊的大小，大小不會是負
的。 所以 $n$ 次操作的總成本不超過總收費，最多 $2n$。

用上一節程式的那 $10$ 次操作追蹤一次：

\begin{center}
\begin{tabular}{clccc}
\toprule
$i$ & 操作 & 實際成本 $c_i$ & 收費 $\hat{c}_i$ & 之後的存款 \\
\midrule
1 & \texttt{push 5} & 1 & 2 & 1 \\
2 & \texttt{push 8} & 1 & 2 & 2 \\
3 & \texttt{push 3} & 1 & 2 & 3 \\
4 & \texttt{multipop 2} & 2 & 0 & 1 \\
5 & \texttt{push 7} & 1 & 2 & 2 \\
6 & \texttt{push 1} & 1 & 2 & 3 \\
7 & \texttt{push 4} & 1 & 2 & 4 \\
8 & \texttt{push 6} & 1 & 2 & 5 \\
9 & \texttt{multipop 10} & 5 & 0 & 0 \\
10 & \texttt{push 2} & 1 & 2 & 1 \\
\midrule
合計 & & 15 & 16 & \\
\bottomrule
\end{tabular}
\end{center}

最右欄每一列都等於那時候堆疊的大小，最小是第 9 次之後的 $0$，沒有變成負的。 總收費
$16$ 減去總成本 $15$，剩下的 $1$ 就是最後堆疊裡那個元素身上的存款。

\subsection{計數器}

\proc{Increment} 收 $2$。 把一個位元從 $0$ 設成 $1$ 時，用掉 $1$ 付這次翻轉，另外
$1$ 存在這個位元身上；之後這個位元從 $1$ 變回 $0$ 時，用它身上的存款付。

每次 \proc{Increment} 的 \textbf{while} 迴圈把一些 $1$ 變成 $0$，這些都由位元身上的
存款付；迴圈後面最多把一個位元從 $0$ 設成 $1$，由這次收的 $2$ 付，還剩下 $1$ 存起
來。 所以收 $2$ 夠用。 存款等於計數器裡 $1$ 的個數，永遠不是負的。 $n$ 次
\proc{Increment} 的總成本最多 $2n$，和定理~\ref{thm:counter-aggregate} 一樣。

記帳法的關鍵是決定存款放在哪裡。 放在「之後要付錢的那個東西」身上最方便：堆疊的存
款放在以後要被拿掉的元素上，計數器的存款放在以後要變回 $0$ 的位元上。 這樣存款夠
不夠用，只要看「每個要付錢的東西身上是不是都有錢」。

\section{位能法}
\label{sec:amo-potential}

記帳法的存款分散在各個元素身上。 \textbf{位能法}（potential method）把存款的總額寫
成資料結構狀態的一個函數，直接用式子計算。

設資料結構一開始的狀態是 $D_0$，第 $i$ 次操作之後的狀態是 $D_i$。 選一個\textbf{位
能函數}（potential function）$\Phi$，它把每一個狀態對應到一個實數 $\Phi(D_i)$。 定
義第 $i$ 次操作的平攤成本是
\[
\hat{c}_i = c_i + \Phi(D_i) - \Phi(D_{i-1}),
\]
實際成本加上這次操作讓位能增加了多少。 位能增加，表示這次多付了錢存起來；位能減少，
表示這次用了存款。

\begin{theorem}
\label{thm:potential}
對任何位能函數 $\Phi$ 和任何一串操作，
\[
\sum_{i=1}^{n} c_i = \sum_{i=1}^{n} \hat{c}_i - \Phi(D_n) + \Phi(D_0).
\]
因此，如果對所有 $i$ 都有 $\Phi(D_i) \ge \Phi(D_0)$，$\hat{c}_i$ 就是平攤成本。
\end{theorem}

\begin{proof}
把 $\hat{c}_i$ 的定義加起來，位能的部分前後相消：
\[
\sum_{i=1}^{n} \hat{c}_i = \sum_{i=1}^{n} \bigl( c_i + \Phi(D_i) - \Phi(D_{i-1}) \bigr)
= \sum_{i=1}^{n} c_i + \Phi(D_n) - \Phi(D_0).
\]
移項就是第一個式子。 如果 $\Phi(D_n) \ge \Phi(D_0)$，右邊減掉的
$\Phi(D_n) - \Phi(D_0)$ 不是負的，所以 $\sum c_i \le \sum \hat{c}_i$。 要求對所有
$i$ 成立，是因為定義~\ref{def:amortized} 要求對每一個 $n$ 都成立。
\end{proof}

實際使用時，通常把位能函數選成 $\Phi(D_0) = 0$、而且永遠 $\Phi(D_i) \ge 0$，這樣條
件自動滿足。 剩下的工作就是對每一種操作算出 $\hat{c}_i$ 的上界。

\subsection{堆疊}

令 $\Phi$ 是堆疊裡的元素個數。 空堆疊的 $\Phi(D_0) = 0$，而且元素個數不會是負的。

\begin{itemize}
\item \proc{Push}：實際成本 $1$，位能加 $1$，$\hat{c}_i = 1 + 1 = 2$。
\item \proc{Pop}：實際成本 $1$，位能減 $1$，$\hat{c}_i = 1 - 1 = 0$。
\item \proc{Multipop}$(S, k)$：設堆疊原本有 $s$ 個，拿掉 $k' = \min(s, k)$ 個。 實
際成本 $k'$，位能減 $k'$，$\hat{c}_i = k' - k' = 0$。
\end{itemize}

三種操作的平攤成本都不超過 $2$，由定理~\ref{thm:potential}，$n$ 次操作的總成本最
多 $2n$。 位能就是記帳法裡的存款總額：每個元素身上 $1$，總共是元素個數。

\subsection{計數器}

令 $\Phi(D_i) = b_i$，第 $i$ 次操作之後計數器裡 $1$ 的個數。 從 $0$ 開始時
$b_0 = 0$，而且 $b_i \ge 0$。

設第 $i$ 次 \proc{Increment} 的 \textbf{while} 迴圈把 $t_i$ 個 $1$ 變成 $0$。 沒有
溢位時，迴圈後面再把一個 $0$ 變成 $1$，所以
\[
c_i = t_i + 1, \qquad b_i = b_{i-1} - t_i + 1, \qquad
\hat{c}_i = (t_i + 1) + (1 - t_i) = 2.
\]
溢位時 $t_i = k$，全部變成 $0$，$c_i = k$、$b_i = 0 = b_{i-1} - k$，$\hat{c}_i = k
- k = 0$。 兩種情況都有 $\hat{c}_i \le 2$，所以總成本最多 $2n$。

照圖~\ref{fig:amo-counter} 的前 $8$ 次逐列算：

\begin{center}
\begin{tabular}{cccccc}
\toprule
$i$ & 之後的值 & $c_i$ & $\Phi(D_i) = b_i$ & $\Phi(D_i) - \Phi(D_{i-1})$ & $\hat{c}_i$ \\
\midrule
1 & 0001 & 1 & 1 & $+1$ & 2 \\
2 & 0010 & 2 & 1 & $0$ & 2 \\
3 & 0011 & 1 & 2 & $+1$ & 2 \\
4 & 0100 & 3 & 1 & $-1$ & 2 \\
5 & 0101 & 1 & 2 & $+1$ & 2 \\
6 & 0110 & 2 & 2 & $0$ & 2 \\
7 & 0111 & 1 & 3 & $+1$ & 2 \\
8 & 1000 & 4 & 1 & $-2$ & 2 \\
\midrule
合計 & & 15 & & & 16 \\
\bottomrule
\end{tabular}
\end{center}

用定理~\ref{thm:potential} 驗算：$\sum c_i = 16 - \Phi(D_8) + \Phi(D_0) = 16 - 1 + 0
= 15$。

計數器不從 $0$ 開始時，$\Phi(D_0) = b_0$ 不是 $0$。 這時不能直接說總成本不超過
$2n$，要用定理~\ref{thm:potential} 的等式：總成本 $= \sum \hat{c}_i - b_n + b_0 \le
2n + b_0$。 一開始有幾個 $1$，總成本最多就多出幾。 觀念題第 2 題有一個例子。

\subsection{三種方法的比較}

\begin{center}
\begin{tabular}{p{0.13\linewidth}p{0.40\linewidth}p{0.36\linewidth}}
\toprule
方法 & 做法 & 證明的重點 \\
\midrule
聚合法 & 直接算任意 $n$ 次操作的總成本上界 $T(n)$，平攤成本是 $T(n)/n$ &
從整體數，例如「每個元素最多被拿掉一次」「$A[i]$ 翻 $\lfloor n/2^{i-1} \rfloor$ 次」 \\
記帳法 & 每種操作訂收費，多收的存在元素身上 & 存款永遠不是負的 \\
位能法 & 選 $\Phi$，$\hat{c}_i = c_i + \Delta\Phi$ & $\Phi(D_0) = 0$、$\Phi \ge 0$，
再逐種操作算 $\hat{c}_i$ \\
\bottomrule
\end{tabular}
\end{center}

三種方法得到的結論一樣，差別在寫法。 聚合法最直接，但只能給所有操作同一個平攤成
本。 記帳法容易想，適合用文字說明。 位能法最適合寫成嚴謹的證明，操作種類多、情況
多的時候（例如下一節會縮小的動態陣列），逐種情況代公式比追蹤每個元素身上的錢清楚。
記帳法的存款總額就是一個位能函數，所以一個記帳法的論證通常可以直接改寫成位能法。

\begin{examnote}
用位能法作答時依序寫出：(1) $\Phi$ 的定義；(2) $\Phi(D_0) = 0$ 而且 $\Phi(D_i) \ge
0$ 的理由；(3) 每一種操作的 $c_i$、$\Phi$ 的變化、$\hat{c}_i$，有特殊情況（溢位、擴
充、縮小）的要分開算；(4) 結論「任意 $n$ 次操作的總成本 $\le \sum \hat{c}_i$」。
用記帳法時，寫清楚每種操作收多少、存款放在哪裡、為什麼不會變負。
\end{examnote}

\section{動態陣列：加倍與縮小}
\label{sec:amo-dynarray}

陣列可以用 $O(1)$ 時間讀寫任何一格，但大小要事先決定。 事先不知道會放多少個元素時，
可以用\textbf{動態陣列}（dynamic array）：記住目前放了幾個（$\mathit{num}$）和目前
的容量（$\mathit{cap}$），放滿了就配置一塊更大的陣列，把舊的元素搬過去，釋放舊的。
章首的筆記本就是這樣做的。

這一節的動態陣列只在尾端加入和刪除，也就是一個用陣列做的堆疊。 成本這樣算：寫入或
刪除一個元素算 $1$，換容量時每搬一個元素算 $1$。 配置和釋放記憶體當作常數時間，不
算進去。

\subsection{只會變大的情形}

\begin{pseudo}[label=alg:table-push]{換容量和在尾端加入}
\Procedure{Resize}{$T, m$}
    \State 配置一個 $m$ 格的新陣列 $B$（$m = 0$ 時 $B \gets \NIL$）
    \For{$i \gets 1$ \To $T.\mathit{num}$}
        \State $B[i] \gets T.\mathit{data}[i]$ \Comment{每搬一個成本 1}
    \EndFor
    \State 釋放舊的陣列
    \State $T.\mathit{data} \gets B$
    \State $T.\mathit{cap} \gets m$
\EndProcedure
\Statex
\Procedure{Table-Push}{$T, x$}
    \If{$T.\mathit{num} = T.\mathit{cap}$}
        \If{$T.\mathit{cap} = 0$}
            \State \Call{Resize}{$T, 1$}
        \Else
            \State \Call{Resize}{$T, 2 \cdot T.\mathit{cap}$}
        \EndIf
    \EndIf
    \State $T.\mathit{num} \gets T.\mathit{num} + 1$
    \State $T.\mathit{data}[T.\mathit{num}] \gets x$
\EndProcedure
\end{pseudo}

從空的陣列（$\mathit{num} = \mathit{cap} = 0$）開始連續 \proc{Table-Push}。 第 $i$
次放入時，如果前面已經有 $i - 1$ 個而且陣列滿了，要先搬 $i - 1$ 個。 容量依序是
$1, 2, 4, 8, \ldots$，所以會搬的是第 $2, 3, 5, 9, 17, \ldots$ 次，也就是 $i - 1$ 是
$2$ 的次方的那幾次。 下圖是前 $17$ 次的實際成本。

\begin{center}
\begin{tikzpicture}[x=0.62cm, y=0.24cm]
\foreach \c [count=\i] in {1,2,3,1,5,1,1,1,9,1,1,1,1,1,1,1,17} {
    \draw[costwrite] (\i-0.3,0) rectangle (\i+0.3,1);
    \ifnum\c>1\relax
        \draw[costcopy] (\i-0.3,1) rectangle (\i+0.3,\c);
    \fi
    \node[font=\scriptsize, text=black!60] at (\i,-0.9) {\i};
}
\draw (0.4,0) -- (17.8,0);
\draw[-{Stealth[length=2mm]}] (0.4,0) -- (0.4,18.5);
\foreach \t in {1,3,5,9,17} { \draw (0.4,\t) -- (0.25,\t) node[left, font=\scriptsize] {\t}; }
\draw[thick, dashed, optStroke] (0.4,3) -- (17.8,3) node[right, font=\small] {平攤 $3$};
\node[font=\small, anchor=west] at (18.2,-0.9) {$i$};
\draw[costwrite] (2,15) rectangle (2.5,16); \node[anchor=west, font=\small] at (2.6,15.5) {寫入};
\draw[costcopy] (2,12.5) rectangle (2.5,13.5); \node[anchor=west, font=\small] at (2.6,13) {搬動};
\end{tikzpicture}
\diagcaption{連續 \proc{Table-Push} 時第 $i$ 次的實際成本。 藍色是寫入自己的 $1$，
橘色是搬動舊元素的成本。 偶爾一次很高，但都在一長串成本 $1$ 之後。}
\label{fig:amo-push-cost}
\end{center}

\textbf{聚合法}：$n$ 次放入的寫入成本是 $n$。 搬動發生在容量是 $1, 2, 4, \ldots$ 而
且已滿的時候，每次搬的個數等於當時的容量，而這些容量都小於 $n$。 它們是 $2$ 的次方，
最大的一個小於 $n$，由等比級數，總和小於最大那一項的兩倍，也就小於 $2n$。 總成本小
於 $n + 2n = 3n$，平攤成本是 $3$。 精確地說，$n \ge 2$ 時搬動的總數是
$1 + 2 + \cdots + 2^{p} = 2^{p+1} - 1$，其中 $2^p$ 是小於 $n$ 的最大的 $2$ 的次方。

\textbf{記帳法}：每次 \proc{Table-Push} 收 $3$。 $1$ 付寫入自己；$1$ 存在自己身上，
留著下次擴充時搬自己；最後 $1$ 存在一個「上次擴充時已經被搬過、身上沒錢的」舊元素
身上。 為什麼夠：容量剛從 $m$ 變成 $2m$ 時（$m \ge 1$），陣列裡有 $m + 1$ 個元素。
到下一次擴充之前，第 $m + 1$ 到第 $2m$ 個元素各存了 $2$，共 $2m$，剛好夠下一次擴
充時搬 $2m$ 個元素。

\textbf{位能法}：令 $\Phi = 2 \cdot \mathit{num} - \mathit{cap}$。 一開始
$\Phi = 0$。 只會變大時，除了一開始，陣列至少半滿（擴充後 $\mathit{num} =
\mathit{cap}/2 + 1$），所以 $\Phi \ge 0$。
\begin{itemize}
\item 不用擴充：$c_i = 1$，$\mathit{num}$ 加 $1$，$\Phi$ 加 $2$，$\hat{c}_i = 3$。
\item 擴充，原本 $\mathit{num} = \mathit{cap} = m \ge 1$：$c_i = m + 1$。 擴充前
$\Phi = 2m - m = m$，擴充後 $\Phi = 2(m + 1) - 2m = 2$，$\hat{c}_i = (m + 1) + 2 - m
= 3$。
\item 第一次放入（$\mathit{cap} = 0$）：$c_i = 1$，$\Phi$ 從 $0$ 變成 $1$，
$\hat{c}_i = 2$。
\end{itemize}
擴充前的位能 $m$ 正好付掉搬 $m$ 個元素的成本，這就是記帳法裡存下來的錢。

章首的問題答案是：平均每記一筆寫不到 $3$ 行。 後面的程式會印出，記 $1000$ 筆時平均
大約 $2.02$ 行。

\subsubsection{為什麼要乘以倍數，不能加固定的量}

如果每次擴充只多 $d$ 格（容量 $d, 2d, 3d, \ldots$），搬動的個數是
$d + 2d + \cdots$，$n$ 次放入大約擴充 $n/d$ 次，搬動總數大約
\[
d \left( 1 + 2 + \cdots + \frac{n}{d} \right) \approx \frac{n^2}{2d} = \Theta(n^2),
\]
$d$ 是常數時平攤成本是 $\Theta(n)$，不是常數。 這是等差級數，和最後一項的平方同一
級。 乘以一個固定的倍數 $r > 1$ 時，容量是等比級數，搬動總數和最後一項同一級，平攤
成本才是常數。 $r$ 不一定要是 $2$：$r$ 越大，搬動越少，但陣列可能有更大的比例是空
的（剛擴充完只用了大約 $1/r$）。 觀念題第 3、4 題分別算這兩種情形。

\subsection{也會變小：為什麼縮到 1/4 才減半}

元素被刪掉很多以後，我們希望把多餘的記憶體還回去，讓 $\mathit{cap}$ 和
$\mathit{num}$ 維持在同一級。 很自然的想法是：放滿就加倍，少於一半就減半。 這樣
做會出問題。

設容量是 $m$（$m$ 是 $2$ 的次方），剛好放滿 $m$ 個。 接著重複下面四次操作：
\begin{enumerate}
\item 加入：滿了，擴充成 $2m$，搬 $m$ 個。 現在有 $m + 1$ 個。
\item 刪除：剩 $m$ 個，剛好一半，不縮小。
\item 刪除：剩 $m - 1$ 個，少於一半，縮成 $m$，搬 $m - 1$ 個。
\item 加入：$m$ 個，剛好放滿，回到一開始的樣子。
\end{enumerate}
每四次操作就搬了大約 $2m$ 個元素。 $m$ 和 $n$ 同一級時，平攤成本是 $\Theta(n)$。
問題出在擴充完的陣列剛好半滿、縮小完的陣列剛好全滿，兩個門檻緊貼在一起，來回一兩
次操作就會再換一次容量。

解法是把縮小的門檻放低，讓剛換完容量的陣列離兩個門檻都有一段距離：

\begin{pseudo}[label=alg:table-pop]{在尾端刪除}
\Function{Table-Pop}{$T$}
    \State $x \gets T.\mathit{data}[T.\mathit{num}]$
    \State $T.\mathit{num} \gets T.\mathit{num} - 1$
    \If{$T.\mathit{num} = 0$}
        \State \Call{Resize}{$T, 0$} \Comment{空了就全部還回去}
    \ElsIf{$4 \cdot T.\mathit{num} < T.\mathit{cap}$}
        \State \Call{Resize}{$T, T.\mathit{cap} / 2$}
    \EndIf
    \State \Return $x$
\EndFunction
\end{pseudo}

少於四分之一才減半，減半之後大約半滿；放滿才加倍，加倍之後也大約半滿。 從半滿到下
一次換容量，至少要再做容量的四分之一那麼多次操作，換容量的成本就分攤得掉。 下面用位
能法把這件事證明清楚。 容量一開始是 $0$，之後只會是 $1$ 或乘 $2$、除以 $2$、變回
$0$，所以永遠是 $0$ 或 $2$ 的次方。

\begin{theorem}
\label{thm:dyn-array}
動態陣列從空的開始，用 \proc{Table-Push} 和 \proc{Table-Pop} 做任意 $n$ 次操作，實
際總成本不超過 $3n$。 而且任何時候只要 $\mathit{num} \ge 1$，就有 $\mathit{cap} \le
4 \cdot \mathit{num}$。
\end{theorem}

\begin{proof}
令 $\Phi = |2 \cdot \mathit{num} - \mathit{cap}|$。 一開始 $\Phi = 0$，而且絕對值不
會是負的。 證明每一種情況都有 $\hat{c}_i \le 3$。 以下 $c$ 是操作之前的容量。

\textbf{加入、不擴充}：$c_i = 1$。 $2 \cdot \mathit{num}$ 加 $2$，絕對值最多變大 $2$，
$\hat{c}_i \le 3$。

\textbf{加入、擴充}：$c = 0$ 時 $c_i = 1$，$\Phi$ 從 $0$ 變成 $|2 - 1| = 1$，
$\hat{c}_i = 2$。 $c \ge 1$ 時原本 $\mathit{num} = c$，$c_i = c + 1$，$\Phi$ 從
$|2c - c| = c$ 變成 $|2(c + 1) - 2c| = 2$，$\hat{c}_i = (c + 1) + 2 - c = 3$。

以下是刪除，設刪除之後剩 $r$ 個（$r = \mathit{num} - 1$）。

\textbf{刪除之後變空}（$r = 0$）：$c_i = 1$，不用搬，$\Phi$ 變成 $0$，
$\hat{c}_i = 1 - \Phi_{\text{之前}} \le 1$。

\textbf{刪除、不縮小}（$r \ge 1$ 且 $4r \ge c$）：$c_i = 1$，$2 \cdot \mathit{num}$
減 $2$，絕對值最多變大 $2$，$\hat{c}_i \le 3$。

\textbf{刪除、縮小}（$r \ge 1$ 且 $4r < c$）：因為 $c > 4r \ge 4$ 而 $c$ 是 $2$ 的
次方，$c \ge 8$。 實際成本是刪除的 $1$ 加上搬 $r$ 個，$c_i = 1 + r$。 縮小之後
$2r < c/2$，所以 $\Phi_{\text{之後}} = c/2 - 2r$。 刪除之前有 $r + 1$ 個，
$2(r + 1) < c/2 + 2 \le c$，所以 $\Phi_{\text{之前}} = c - 2r - 2$。
\[
\hat{c}_i = (1 + r) + \left(\frac{c}{2} - 2r\right) - (c - 2r - 2)
= 3 + r - \frac{c}{2} < 3 + \frac{c}{4} - \frac{c}{2} < 3.
\]

每一種情況都有 $\hat{c}_i \le 3$，由定理~\ref{thm:potential}，總成本不超過 $3n$。

記憶體的部分對操作次數做歸納。 一開始 $\mathit{num} = 0$，不用檢查。 加入不擴充時
$\mathit{num}$ 變大，不等式仍然成立；擴充時 $\mathit{cap}$ 變成 $2c$、$\mathit{num}$
變成 $c + 1$（或從 $0$ 變成容量 $1$、$1$ 個），成立。 刪除不縮小時，$4r \ge c$ 正是要
的不等式。 刪除而縮小時，由歸納假設，刪除之前 $4(r + 1) \ge c$，所以 $4r \ge c - 4
\ge c/2$（因為 $c \ge 8$），縮小後的容量 $c/2$ 不超過 $4r$。
\end{proof}

用 $9$ 次加入、$7$ 次刪除追蹤一次。 表裡的 $\mathit{num}$、$\mathit{cap}$、$\Phi$ 都
是操作之後的值，一開始三個都是 $0$。

\begin{center}
\begin{tabular}{clccccc}
\toprule
$i$ & 操作 & $\mathit{num}$ & $\mathit{cap}$ & $c_i$ & $\Phi$ & $\hat{c}_i$ \\
\midrule
1 & 加入 & 1 & 1 & 1 & 1 & 2 \\
2 & 加入 & 2 & 2 & 2 & 2 & 3 \\
3 & 加入 & 3 & 4 & 3 & 2 & 3 \\
4 & 加入 & 4 & 4 & 1 & 4 & 3 \\
5 & 加入 & 5 & 8 & 5 & 2 & 3 \\
6--8 & 加入 & 6, 7, 8 & 8 & 各 1 & 4, 6, 8 & 各 3 \\
9 & 加入 & 9 & 16 & 9 & 2 & 3 \\
10 & 刪除 & 8 & 16 & 1 & 0 & $-1$ \\
11--14 & 刪除 & 7, 6, 5, 4 & 16 & 各 1 & 2, 4, 6, 8 & 各 3 \\
15 & 刪除 & 3 & 8 & 4 & 2 & $-2$ \\
16 & 刪除 & 2 & 8 & 1 & 4 & 3 \\
\midrule
合計 & & & & 34 & & 38 \\
\bottomrule
\end{tabular}
\end{center}

第 15 次刪除後剩 $3$ 個，$4 \cdot 3 < 16$，縮成 $8$，搬 $3$ 個，實際成本 $4$。 它的
平攤成本是負的，因為它用掉了第 11 到 14 次刪除存下來的位能。 驗算：
$\sum c_i = 38 - \Phi(D_{16}) + \Phi(D_0) = 38 - 4 + 0 = 34$。

\begin{examnote}
動態陣列的追蹤題，每一列寫出操作之後的 $\mathit{num}$、$\mathit{cap}$ 和這次的實際
成本，換容量的那一列另外註明搬了幾個。 先確認題目的規則：容量從 $0$ 還是 $1$ 開始、
縮小的條件是「少於」還是「不超過」四分之一、判斷是在刪除之前還是之後，這些細節會改
變換容量的時間點。
\end{examnote}

\subsection{C 實作}

\begin{cbookrecap}[第~19~章 19.3 節]
\texttt{realloc(p, size)} 把 \texttt{p} 指的那塊記憶體改成 \texttt{size} 個位元
組，回傳新的位址。 它可能在原地變大或變小，也可能另外找一塊、把舊內容複製過去再釋
放舊的。 失敗時回傳 \texttt{NULL}，而且\textbf{原本那塊仍然有效}，所以要先用暫存指標
接：\texttt{tmp = realloc(p, ...)}，確定不是 \texttt{NULL} 才寫 \texttt{p = tmp}。
直接寫 \texttt{p = realloc(p, ...)}，失敗時唯一的位址被蓋掉，就漏掉了。
\texttt{realloc} 傳入大小 $0$ 的行為由各個系統自己決定，下面的程式在容量變成 $0$
時改用 \texttt{free}。
\end{cbookrecap}

程式把動態陣列包成一個結構，縮小的門檻寫成參數 \texttt{shrink\_den}：
\texttt{num * shrink\_den < cap} 時減半。 \texttt{shrink\_den} 是 $4$ 就是
虛擬碼~\ref{alg:table-pop}，是 $2$ 就是會來回換容量的做法。 成本照前面的算法記在
\texttt{cost} 裡。 \texttt{realloc} 有時能原地變大，不用真的搬，但最壞情況要複製全部
的元素，所以成本照搬 $\mathit{num}$ 個來算。

\begin{example}[label=ex:dyn-array]{c}{(dyn\_array.c) 會加倍和減半的動態陣列}
#include <stdio.h>
#include <stdlib.h>

// 會自己變大變小的陣列（拿來當堆疊用：只在尾端加入和刪除）。
typedef struct {
    int *data;
    int num;            // 目前放了幾個
    int cap;            // 目前的容量
    int shrink_den;     // num * shrink_den < cap 時把容量減半
    long long cost;     // 實際成本：寫入或刪除一個算 1，搬一個也算 1
    int resizes;        // 換了幾次容量
} DynArray;

void da_init(DynArray *a, int shrink_den)
{
    a->data = NULL;
    a->num = 0;
    a->cap = 0;
    a->shrink_den = shrink_den;
    a->cost = 0;
    a->resizes = 0;
}

// 把容量改成 new_cap，成功回傳 1。搬動 num 個元素的成本算進 cost。
int da_resize(DynArray *a, int new_cap)
{
    if (new_cap == 0) {
        free(a->data);
        a->data = NULL;
    } else {
        int *tmp = realloc(a->data, new_cap * sizeof *tmp);
        if (tmp == NULL) {              // 失敗時 a->data 仍然有效，不要蓋掉
            return 0;
        }
        a->data = tmp;
    }
    a->cost += a->num;
    a->cap = new_cap;
    a->resizes++;
    return 1;
}

int da_push(DynArray *a, int x)
{
    if (a->num == a->cap) {             // 滿了：加倍（容量 0 時變成 1）
        if (!da_resize(a, a->cap == 0 ? 1 : 2 * a->cap)) {
            return 0;
        }
    }
    a->data[a->num++] = x;
    a->cost += 1;
    return 1;
}

// 刪除並回傳最後一個元素。呼叫前要確定 num > 0。
int da_pop(DynArray *a)
{
    int x = a->data[--a->num];
    a->cost += 1;
    if (a->num == 0) {                  // 空了就整塊還回去
        da_resize(a, 0);
    } else if (a->num * a->shrink_den < a->cap) {
        da_resize(a, a->cap / 2);      // 縮小失敗也沒關係，舊的那塊還能用
    }
    return x;
}

int main(void)
{
    // 第一部分：從空的開始 push n 次
    int ns[] = {10, 100, 1000, 1025, 1000000};
    printf("n cost cost/n cap\n");
    for (int k = 0; k < 5; k++) {
        DynArray a;
        da_init(&a, 4);
        for (int i = 0; i < ns[k]; i++) {
            if (!da_push(&a, i)) {
                free(a.data);
                return 1;
            }
        }
        printf("%d %lld %.3f %d\n", ns[k], a.cost, (double)a.cost / ns[k], a.cap);
        free(a.data);
    }

    // 第二部分：先放 1024 個，再重複 1000 次「push, pop, pop, push」
    printf("rule cost resizes\n");
    int dens[] = {2, 4};
    for (int k = 0; k < 2; k++) {
        DynArray a;
        da_init(&a, dens[k]);
        for (int i = 0; i < 1024; i++) {
            if (!da_push(&a, i)) {
                free(a.data);
                return 1;
            }
        }
        a.cost = 0;                     // 只算後面這 4000 次操作
        a.resizes = 0;
        int ok = 1;
        for (int r = 0; r < 1000 && ok; r++) {
            ok = da_push(&a, r);
            da_pop(&a);
            da_pop(&a);
            ok = ok && da_push(&a, r);
        }
        if (!ok) {
            free(a.data);
            return 1;
        }
        printf("1/%d %lld %d\n", dens[k], a.cost, a.resizes);
        while (a.num > 0) {             // 全部 pop 掉，最後整塊還回去
            da_pop(&a);
        }
        printf("after popping all: num %d, cap %d\n", a.num, a.cap);
    }
    return 0;
}
\end{example}
\iosample{}{n cost cost/n cap \\ 10 25 2.500 16 \\ 100 227 2.270 128 \\ 1000 2023 2.023 1024 \\ 1025 3072 2.997 2048 \\ 1000000 2048575 2.049 1048576 \\ rule cost resizes \\ 1/2 2051000 2000 \\ after popping all: num 0, cap 0 \\ 1/4 5024 1 \\ after popping all: num 0, cap 0}

第一部分：$n = 1000$ 時搬了 $1 + 2 + \cdots + 512 = 1023$ 個，總成本 $2023$，平均
每次 $2.023$，這就是章首的答案。 $n = 1025$ 是最不划算的時候，剛好在第 $1025$ 次多
搬了 $1024$ 個，平均 $2.997$，仍然小於 $3$。

第二部分是同一串 $4000$ 次操作，只差在縮小的門檻。 門檻是一半時，每一輪「加入、刪
除、刪除、加入」都擴充一次、縮小一次，各搬大約 $1024$ 個，每輪成本
$1025 + 1 + 1024 + 1 = 2051$，$1000$ 輪是 $2{,}051{,}000$。 門檻是四分之一時，只有
第一次加入擴充成 $2048$（搬 $1024$ 個），之後元素在 $1023$ 到 $1025$ 個之間變動，
離縮小的門檻 $512$ 還很遠，總成本 $4000 + 1024 = 5024$。 兩者差了大約 $400$ 倍。

\section{本章常見錯誤}
\label{sec:amo-pitfalls}

\begin{center}
\begin{tabular}{p{0.26\linewidth}p{0.48\linewidth}p{0.16\linewidth}}
\toprule
錯誤 & 例子 & 見 \\
\midrule
把平攤當成平均情況 & 說「平均來說 \proc{Multipop} 很快」，或去假設輸入的分布 &
第~\ref{sec:amo-aggregate}~節 \\
以為平攤 $O(1)$ 就是每次 $O(1)$ & 說動態陣列的每一次加入都是常數時間；擴充那一次是
$\Theta(n)$ & 第~\ref{sec:amo-aggregate}、\ref{sec:amo-dynarray}~節 \\
每次的最壞情況乘以 $n$ 就當答案 & 說 $n$ 次 \proc{Increment} 是 $\Theta(nk)$；
$O(nk)$ 是對的上界，但不緊 & 第~\ref{sec:amo-counter}~節 \\
記帳法沒有證明存款不會變負 & 只寫「\proc{Push} 收 $2$」，沒有說明每個元素身上都有
$1$ & 第~\ref{sec:amo-accounting}~節 \\
位能函數沒檢查 $\Phi(D_i) \ge \Phi(D_0)$ & 選了會變負的 $\Phi$，或計數器不從 $0$ 開始
時忘了加 $b_0$ & 第~\ref{sec:amo-potential}~節 \\
特殊情況漏算 & 只算了不擴充的加入，沒算擴充、溢位、縮小那一次 &
第~\ref{sec:amo-potential}、\ref{sec:amo-dynarray}~節 \\
擴充時加固定的量 & 每次多 $100$ 格，以為平攤仍是 $O(1)$ & 第~\ref{sec:amo-dynarray}~節 \\
縮小的門檻設在一半 & 剛好在邊界來回加入刪除，每次都換容量 &
第~\ref{sec:amo-dynarray}~節 \\
\texttt{p = realloc(p, ...)} & 失敗時原本的位址被 \texttt{NULL} 蓋掉，記憶體漏掉 &
第~\ref{sec:amo-dynarray}~節 \\
\bottomrule
\end{tabular}
\end{center}

\section{觀念題}
\label{sec:amo-exercises}

題型和研究所考試、期中期末考的筆試相近，有選擇、計算、追蹤和證明。 解答在本章最後。

\begin{enumerate}
\item （選擇，可能不只一個答案）下列哪些敘述正確？
\begin{enumerate}
\item 某個操作的平攤成本是 $O(1)$，表示每一次執行這個操作都只花 $O(1)$ 時間。
\item 平攤分析要先假設輸入的機率分布，算出來的是期望值。
\item 用記帳法證明了每種操作的收費都不超過 $3$ 而且存款不會變負，則從初始狀態開始
的任何 $n$ 次操作，實際總成本不超過 $3n$。
\item 用位能法時，只要 $\Phi(D_0) = 0$，$\Phi$ 在中間的某些狀態取負值也沒關係。
\item 記帳法裡，某一次操作的平攤成本可以比它的實際成本小。
\end{enumerate}

\item （追蹤）一個 $6$ 位元的計數器目前是 $001011$。 連續做 $6$ 次 \proc{Increment}，
寫出每一次之後的值和翻轉的位元數，求總數。 再用位能 $\Phi = $ 計數器裡 $1$ 的個數，
驗證定理~\ref{thm:potential} 的等式。

\item （計算）動態陣列一開始容量 $1$、沒有元素，放滿時把容量變成原來的 $3$ 倍（不
縮小）。 成本照第~\ref{sec:amo-dynarray}~節的算法。
\begin{enumerate}
\item 連續放入 $82$ 個元素，總成本是多少？ 最後的容量是多少？
\item 證明連續放入 $n$ 個的總成本小於 $2.5n$。
\item 和加倍比，這個做法在空間上的缺點是什麼？
\end{enumerate}

\item （計算）動態陣列一開始容量 $4$，放滿時容量加 $4$（$4, 8, 12, \ldots$）。 連續
放入 $40$ 個元素總共搬動幾次？ 對 $n = 4m$ 寫出搬動總數的公式，並說明每次放入的
平攤成本是 $\Theta(n)$。

\item （證明）三進位計數器的每一位是 $0$、$1$ 或 $2$。 加一的做法是：從最低位開始，
把連續的 $2$ 都改成 $0$，再把遇到的第一個不是 $2$ 的位加一。 成本是改動的位數。
\begin{enumerate}
\item 用聚合法證明從 $0$ 開始做 $n$ 次加一，改動的總位數小於 $1.5n$。
\item 找一個位能函數，使每次加一的平攤成本都不超過 $3/2$，並驗證。
\end{enumerate}

\item （證明）一個堆疊保持「由下到上遞增」。 \proc{Insert}$(S, x)$ 先把頂端比 $x$
大的元素一個一個 \proc{Pop} 掉，再 \proc{Push}$(x)$，成本是 \proc{Pop} 的次數加 $1$。
\begin{enumerate}
\item 從空堆疊開始做 $n$ 次 \proc{Insert}，證明總成本不超過 $2n$。
\item 舉一串 $n$ 次 \proc{Insert}，其中某一次的成本是 $n$。
\end{enumerate}

\item （追蹤）用虛擬碼~\ref{alg:table-push}、\ref{alg:table-pop} 的動態陣列，目前
$\mathit{num} = 5$、$\mathit{cap} = 8$。 依序做 $4$ 次刪除、$5$ 次加入。 每次寫出操
作之後的 $\mathit{num}$、$\mathit{cap}$、實際成本 $c_i$、$\Phi = |2 \cdot \mathit{num}
- \mathit{cap}|$ 和 $\hat{c}_i$，並驗證定理~\ref{thm:potential} 的等式。

\item （選擇，可能不只一個答案）下列動態陣列的規則中，哪些能讓「從空的開始任意 $n$
次加入、刪除」的總成本是 $O(n)$？ 搬動的成本照本章的算法。
\begin{enumerate}
\item 放滿時加倍；刪除後少於一半時減半。
\item 放滿時加倍；刪除後少於四分之一時減半。
\item 放滿時加倍；刪除後少於三分之一時減半。
\item 放滿時加 $100$ 格；永遠不縮小。
\item 放滿時變成 $3$ 倍；刪除後少於九分之一時變成三分之一。
\end{enumerate}

\item （證明）一個集合用陣列存。 \proc{Insert} 放在已用部分的尾端，成本 $1$（假設
陣列夠大）。 \proc{Delete} 給定位置，只把那一格標成「已刪除」，成本 $1$；標完之後，
如果已刪除的格數超過已用格數的一半，就做一次\textbf{整理}：從頭掃過已用的 $u$ 格，
把沒刪除的往前搬，成本 $u$，整理完沒有已刪除的格子。 用位能
$\Phi = 2 \times (\text{已刪除的格數})$ 證明每種操作的平攤成本都不超過 $3$。

\item （找錯）有人分析可以一次拿掉很多個的堆疊時，選位能 $\Phi = -(\text{堆疊的大
小})$，得到 \proc{Push} 的平攤成本是 $1 + (-1) = 0$，於是說「$n$ 次 \proc{Push} 的
總成本是 $0$」。 錯在哪裡？

\item （證明）$k$ 位元計數器從 $0$ 開始做 $n$ 次 \proc{Increment}，$n < 2^k$。 證明翻轉
的位元總數剛好是 $2n - b(n)$，其中 $b(n)$ 是 $n$ 的二進位裡 $1$ 的個數。 用
範例~\ref{ex:counter-flips} 的輸出驗證 $n = 16$ 和 $n = 10^6$。
\end{enumerate}

\section{實作練習}
\label{sec:amo-practice}

這一章的重點是證明，Aki-Judge 上直接對應的題目只有一題，其他練習請以觀念題為主。

\begin{center}
\begin{tabular}{llc}
\toprule
小節 & 題目 & 難度 \\
\midrule
動態陣列 & \akiref{growable-array-stack}{可自動擴充的陣列堆疊} & 基本 \\
\bottomrule
\end{tabular}
\end{center}

這一題要交標頭檔和實作檔（《C 程式設計基礎》第~20~章的分檔寫法），只要求放滿時加倍，
不要求縮小。 寫完後可以自己加上縮小的規則，用範例~\ref{ex:dyn-array} 第二部分的操
作順序測試，確認門檻是四分之一時不會來回換容量。

\section{讀完你應該會}
\label{sec:amo-checklist}

\begin{itemize}
\item 說明平攤成本和最壞情況、平均情況的差別，並判斷「平攤 $O(1)$」這類敘述在什麼
意義下成立。
\item 用聚合法算出 \proc{Multipop}、二進位計數器這類操作序列的總成本上界。
\item 用記帳法訂出每種操作的收費，說明存款放在哪裡、為什麼不會變負。
\item 選一個位能函數，檢查 $\Phi(D_0) = 0$、$\Phi \ge 0$，逐種情況算出平攤成本。
\item 手算動態陣列一串加入、刪除後的 $\mathit{num}$、$\mathit{cap}$ 和每次的成本，說
明為什麼要乘以倍數擴充、為什麼縮小的門檻要比一半低。
\item 用 \texttt{realloc} 正確地寫出會變大變小的陣列：用暫存指標接回傳值，容量變成
$0$ 時改用 \texttt{free}。
\end{itemize}

\textbf{延伸閱讀}（書末會整理完整的參考資料）

\begin{itemize}
\item Thomas H. Cormen 等，\textit{Introduction to Algorithms}，第 4 版，第 16 章：
三種方法的完整講法，動態表格同時考慮擴充和縮小。
\item Robert E. Tarjan，``Amortized Computational Complexity''，\textit{SIAM Journal
on Algebraic and Discrete Methods}，1985：位能法的經典論文，篇幅不長。
\item Steven S. Skiena，\textit{The Algorithm Design Manual}，第 3 版，第 3.1 節：
從實作的角度講動態陣列。
\item Antti Laaksonen，\textit{Competitive Programmer's Handbook}，第 8 章：平攤分析
在雙指標、單調堆疊這類解題技巧上的用法。
\end{itemize}

\addcontentsline{toc}{section}{觀念題解答}
\section*{觀念題解答}

\begin{enumerate}
\item (c)、(e) 正確。 (a)：平攤 $O(1)$ 只保證總和，單一次可能很貴，例如動態陣列擴
充的那一次。 (b)：平攤分析不用機率。 (c)：存款不變負就是 $\sum_{i \le j} c_i \le
\sum_{i \le j} \hat{c}_i \le 3j$。 (d)：定理~\ref{thm:potential} 要求每個 $i$ 都有
$\Phi(D_i) \ge \Phi(D_0)$；某個 $D_j$ 的位能是負的，前 $j$ 次的
$\sum \hat{c}_i$ 就可能小於實際成本。 (e)：例如 \proc{Multipop} 收 $0$，實際成本可
以很大，差額由存款付。

\item 每一次之後的值和翻轉的位元數：
\begin{center}
\begin{tabular}{ccccc}
\toprule
次 & 之後的值 & 翻轉 & $\Phi$ & $\hat{c}_i$ \\
\midrule
開始 & 001011 & & 3 & \\
1 & 001100 & 3 & 2 & 2 \\
2 & 001101 & 1 & 3 & 2 \\
3 & 001110 & 2 & 3 & 2 \\
4 & 001111 & 1 & 4 & 2 \\
5 & 010000 & 5 & 1 & 2 \\
6 & 010001 & 1 & 2 & 2 \\
\midrule
合計 & & 13 & & 12 \\
\bottomrule
\end{tabular}
\end{center}
沒有溢位，每次 $\hat{c}_i = 2$。 $\sum c_i = \sum \hat{c}_i - \Phi(D_6) + \Phi(D_0)
= 12 - 2 + 3 = 13$，和直接數的一樣。 總數比 $2 \cdot 6 = 12$ 多 $1$，因為一開始的
$3$ 個 $1$ 帶著位能。

\item
\begin{enumerate}
\item 容量依序是 $1, 3, 9, 27, 81, 243$。 放第 $2, 4, 10, 28, 82$ 個時陣列是滿的，
分別搬 $1, 3, 9, 27, 81$ 個，共 $121$ 個。 總成本 $82 + 121 = 203$，最後容量 $243$。
\item 搬動的個數是 $3$ 的次方，而且都小於 $n$（擴充時已有的元素個數小於這次放入後
的 $n$）。 設最大的一項是 $3^p < n$，總和
$1 + 3 + \cdots + 3^p = \frac{3^{p+1} - 1}{2} < \frac{3}{2} \cdot 3^p < 1.5n$。 加上
寫入的 $n$，總成本小於 $2.5n$。
\item 剛擴充完時，容量是元素個數的大約 $3$ 倍，大約三分之二的空間沒有用到；加倍時最
多大約一半沒用到。 倍數越大，搬得越少，浪費的空間越多。
\end{enumerate}

\item 容量是 $4j$ 而且放滿時，放下一個要搬 $4j$ 個。 放 $40$ 個時，在第
$5, 9, \ldots, 37$ 個擴充，搬 $4 + 8 + \cdots + 36 = 4(1 + 2 + \cdots + 9) = 180$
次。 一般地，$n = 4m$ 時搬動總數是
\[
4 \bigl(1 + 2 + \cdots + (m - 1)\bigr) = 2m(m - 1) = \frac{n(n - 4)}{8},
\]
是 $\Theta(n^2)$。 總成本是 $n + \frac{n(n-4)}{8} = \Theta(n^2)$，除以 $n$，平攤成本是
$\Theta(n)$。 加固定的量時，搬動的個數是等差級數，總和和項數的平方同一級。

\item
\begin{enumerate}
\item 第 $m$ 次加一改動最低位開始的一段：從最低位起連續的 $2$，加上它們上面那一位。
所以從低位數來第 $j + 1$ 位（$j \ge 0$）在第 $m$ 次被改動，若且唯若 $m - 1$ 最低的
$j$ 位全是 $2$，也就是 $m - 1$ 除以 $3^j$ 的餘數是 $3^j - 1$，也就是 $3^j$ 整除 $m$。
$1$ 到 $n$ 之間這樣的 $m$ 有 $\lfloor n/3^j \rfloor$ 個。 總數
\[
\sum_{j \ge 0} \left\lfloor \frac{n}{3^j} \right\rfloor
< n \sum_{j \ge 0} \frac{1}{3^j} = \frac{3}{2} n.
\]
\item 令 $\Phi = \frac{1}{2} \times (\text{值是 } 1 \text{ 的位數}) + (\text{值是 } 2
\text{ 的位數})$。 一開始全是 $0$，$\Phi = 0$，而且 $\Phi \ge 0$。 一次加一把 $t$
個 $2$ 改成 $0$（位能減 $t$），再把一個 $0$ 改成 $1$ 或把一個 $1$ 改成 $2$（位能都加
$1/2$）。 實際成本 $t + 1$，$\hat{c}_i = (t + 1) - t + \frac{1}{2} = \frac{3}{2}$。
全部是 $2$ 而溢位時，成本是位數 $t$，位能減 $t$，$\hat{c}_i = 0$。
\end{enumerate}

\item
\begin{enumerate}
\item 每個元素最多被 \proc{Push} 一次、\proc{Pop} 一次。 $n$ 次 \proc{Insert} 共
\proc{Push} $n$ 次，所以 \proc{Pop} 的總次數不超過 $n$，總成本不超過 $2n$。 用位能
法寫：$\Phi = $ 堆疊大小，一次 \proc{Insert} pop 掉 $p$ 個，$c_i = p + 1$，位能變化
$1 - p$，$\hat{c}_i = 2$。
\item 依序插入 $2, 3, \ldots, n$，再插入 $1$。 前 $n - 1$ 次都不用 pop；最後一次要
把 $n - 1$ 個全部 pop 掉，成本 $(n - 1) + 1 = n$。 找「左邊第一個比它小的元素」這類問題用的單調堆疊就是這個結構。
\end{enumerate}

\item 一開始 $\Phi = |10 - 8| = 2$。
\begin{center}
\begin{tabular}{clccccc}
\toprule
$i$ & 操作 & $\mathit{num}$ & $\mathit{cap}$ & $c_i$ & $\Phi$ & $\hat{c}_i$ \\
\midrule
1 & 刪除 & 4 & 8 & 1 & 0 & $-1$ \\
2 & 刪除 & 3 & 8 & 1 & 2 & 3 \\
3 & 刪除 & 2 & 8 & 1 & 4 & 3 \\
4 & 刪除 & 1 & 4 & 2 & 2 & 0 \\
5 & 加入 & 2 & 4 & 1 & 0 & $-1$ \\
6 & 加入 & 3 & 4 & 1 & 2 & 3 \\
7 & 加入 & 4 & 4 & 1 & 4 & 3 \\
8 & 加入 & 5 & 8 & 5 & 2 & 3 \\
9 & 加入 & 6 & 8 & 1 & 4 & 3 \\
\midrule
合計 & & & & 14 & & 16 \\
\bottomrule
\end{tabular}
\end{center}
第 4 次刪除後剩 $1$ 個，$4 \cdot 1 < 8$，縮成 $4$，搬 $1$ 個；第 8 次加入時陣列是滿的
（$4$ 個），擴充成 $8$，搬 $4$ 個。 驗算：$16 - \Phi(D_9) + \Phi(D_0) = 16 - 4 + 2
= 14$。

\item (b)、(c)、(e)。 (a) 是本章說明過的來回換容量，平攤 $\Theta(n)$。 (d) 只放入也要
搬 $100 + 200 + \cdots$，是 $\Theta(n^2/100) = \Theta(n^2)$。 (b)、(c)、(e) 的共同點
是：剛換完容量時，元素個數離兩個門檻都差了容量的一個固定比例。 (c) 加倍後半滿，要再
放入容量的一半才會滿、刪除容量的六分之一才會低於三分之一；減半後大約三分之二滿，
要放入或刪除大約三分之一的容量才會再換容量。 (e) 變成 $3$ 倍後是三
分之一滿，縮小後也大約是三分之一滿，離「滿」和「九分之一」都差了容量的一個固定比
例。 每次換容量搬動的個數不超過當時的容量，而在下一次換容量之前，至少要做容量的固
定比例那麼多次操作，所以搬動的成本可以平攤成每次 $O(1)$。

\item 設已用 $u$ 格、其中已刪除 $d$ 格。 一開始 $d = 0$，$\Phi = 0$，而且
$\Phi \ge 0$。
\begin{itemize}
\item \proc{Insert}：$c_i = 1$，$d$ 不變，$\hat{c}_i = 1$。
\item \proc{Delete}、不整理：$c_i = 1$，$d$ 加 $1$，$\hat{c}_i = 1 + 2 = 3$。
\item \proc{Delete}、要整理：標完之後 $d > u/2$，也就是 $2d > u$。 實際成本 $1 + u$。
操作前 $\Phi = 2(d - 1)$，整理後 $\Phi = 0$。
$\hat{c}_i = (1 + u) + 0 - 2(d - 1) = 3 + u - 2d < 3$。
\end{itemize}
由定理~\ref{thm:potential}，任意 $n$ 次操作的總成本不超過 $3n$。 如果每次刪除都立刻
往前搬，一次刪除就要 $\Theta(u)$；等到刪除的格子夠多才一起整理，前面每次刪除多付的
$2$ 剛好付得起整理。 第~12~章雜湊表的「刪除標記」會用到同樣的想法。

\item 定理~\ref{thm:potential} 要求 $\Phi(D_i) \ge \Phi(D_0)$。 這裡 $\Phi(D_0) = 0$，
放了 $i$ 個之後 $\Phi(D_i) = -i < 0$，條件不成立，所以 $\hat{c}_i$ 加起來不是實際成本
的上界。 用等式算：$n$ 次 \proc{Push} 的 $\sum c_i = \sum \hat{c}_i - \Phi(D_n) +
\Phi(D_0) = 0 - (-n) + 0 = n$，實際總成本是 $n$。 位能是預先存的錢，不能借了不還。

\item 令 $\Phi = $ 計數器裡 $1$ 的個數。 $n < 2^k$，所以 $n$ 次都沒有溢位，由
第~\ref{sec:amo-potential}~節的計算每次 $\hat{c}_i$ 剛好是 $2$。 由
定理~\ref{thm:potential}，
\[
\sum_{i=1}^{n} c_i = 2n - \Phi(D_n) + \Phi(D_0) = 2n - b(n) + 0,
\]
因為做完 $n$ 次後計數器的值就是 $n$。 驗證：$16 = 10000_2$，$b = 1$，$32 - 1 = 31$；
$10^6 = 11110100001001000000_2$，$b = 7$，$2 \times 10^6 - 7 = 1{,}999{,}993$，都和
程式的輸出一樣。
\end{enumerate}


% 第二部：資料結構
\chapter{陣列、矩陣與多項式}
\label{ch:arrays}

% 本章的圖共用的樣式（名稱都以 arr- 開頭，避免和其他章衝突）
\tikzset{
    arr-r1/.style={draw=coreStroke, fill=coreFill, minimum width=0.95cm,
        minimum height=0.55cm, inner sep=1pt, font=\scriptsize},
    arr-r2/.style={arr-r1, draw=memStroke, fill=memFill},
    arr-r3/.style={arr-r1, draw=optStroke, fill=optFill},
    arr-tri/.style={draw=coreStroke, fill=coreFill, minimum width=0.7cm,
        minimum height=0.7cm, inner sep=1pt, font=\small},
    arr-zero/.style={draw=black!25, text=black!35, minimum width=0.7cm,
        minimum height=0.7cm, inner sep=1pt, font=\small},
}

期末考借了一間大教室，座位排成 $9$ 排，每排 $7$ 個位子。 監考老師拿到的座位名單
是一長串學號，從第 $1$ 排最左邊的位子開始，一排排完再換下一排。 坐在第 $6$ 排左
邊數來第 $4$ 個位子的同學，是名單上的第幾個？ 如果名單改成先把每排最左邊那個位子
從前排到後排列完，再列左邊數來第二個位子，依此類推，同一個同學又變成第幾個？ 不
用把名單數一遍，能不能直接算出來？ 往下讀之前，先算算看。

這一章講的就是這種計算。 記憶體是一長條編了號的格子，二維的表格、三維的方塊放進
去都要先「攤平」，攤平的規則決定了每一格的位置怎麼算。 我們先推導一般陣列的位址
公式，再看幾種形狀特別的矩陣（三角、對稱、帶狀）怎麼只存必要的部分、同時仍然能用
公式在 $O(1)$ 時間找到每一格。 接著是大部分元素都是 $0$ 的稀疏矩陣：只存非零的
項，轉置和相乘要怎麼做才快。 最後用同樣的眼光看多項式：存全部的係數，或只存非零
的項，相加、相乘和求值各要花多少時間。

第~\ref{sec:ana-cost}~節的 RAM 模型假設讀寫任何一格記憶體都花固定時間，這一章說明陣列
為什麼能做到這件事。 第~\ref{sec:ll-poly}~節會把這一章的多項式和稀疏矩陣改用串列存，那裡處理
「項數一直變動」的情況，這一章只用陣列。 用到的 C 語法有二維陣列、陣列參數、指標
算術和結構，用到時會先在框裡複習。

\section{陣列與位址計算}
\label{sec:arr-address}

陣列（array）的元素在記憶體裡一個緊接著一個，而且每個元素占一樣多的空間。 這兩
個性質讓人不用從頭數，只靠一次乘法和一次加法就知道某個元素在哪裡。 這一節把這
個「在哪裡」寫成公式，從一維開始，推到二維、三維和一般的 $d$ 維。

\subsection{一維陣列}

記憶體的每個位元組（byte）有一個編號，叫做\textbf{位址}（address）。 一個一維陣
列寫成 $A[l..u]$，表示下標從 $l$ 到 $u$，共有 $u - l + 1$ 個元素。 $l$ 叫
\textbf{下界}（lower bound），$u$ 叫\textbf{上界}（upper bound）。 C 的陣列下界一
定是 $0$，本書的虛擬碼通常從 $1$ 開始，但考題和一些程式語言允許任意的下界，例如
$A[-5..20]$。

設每個元素占 $w$ 個位元組，第一個元素 $A[l]$ 的位址是 $\alpha$，叫做
\textbf{起始位址}（base address）。 $A[i]$ 前面有 $A[l], A[l+1], \ldots, A[i-1]$
共 $i - l$ 個元素，每個占 $w$，所以
\begin{equation}
\mathrm{loc}(A[i]) = \alpha + (i - l)\,w .
\label{eq:arr-1d}
\end{equation}
這個元素占用的位址是從 $\mathrm{loc}(A[i])$ 到 $\mathrm{loc}(A[i]) + w - 1$。 算
$\mathrm{loc}(A[i])$ 只要一次減法、一次乘法、一次加法，和 $i$ 是多少、陣列多長都
沒有關係，這就是陣列可以在 $O(1)$ 時間讀寫任何一格的原因。

\textbf{例}。 陣列 $T[-10..30]$ 記錄每個整數氣溫（攝氏）出現過幾天，每個元素是
$4$ 位元組的整數，$T[-10]$ 的位址是十六進位的 \texttt{1A00}。 共有
$30 - (-10) + 1 = 41$ 個元素。 $T[17]$ 前面有 $17 - (-10) = 27$ 個元素，位移是
$27 \times 4 = 108$，十六進位是 \texttt{6C}，所以
$\mathrm{loc}(T[17]) = \texttt{1A00} + \texttt{6C} = \texttt{1A6C}$。 最後一個元素
$T[30]$ 的位址是 $\texttt{1A00} + 40 \times 4 = \texttt{1A00} + \texttt{A0} =
\texttt{1AA0}$，整個陣列占用 \texttt{1A00} 到 \texttt{1AA3}，共 $164$ 個位元組。

位址是十六進位時，可以先把位移用十進位算好，再換成十六進位相加；也可以把起始位址
換成十進位，算完再換回去。 兩種做法都行，重點是不要把十六進位的 \texttt{1A00} 當
成十進位的 $1A00$ 去加。

\subsection{二維陣列：列優先與行優先}

二維陣列 $A[l_1..u_1,\ l_2..u_2]$ 有 $m = u_1 - l_1 + 1$ 列（row）、
$n = u_2 - l_2 + 1$ 行（column，也叫欄）。 記憶體只有一維，要先決定這 $mn$ 個元素
排成一條時的順序。 常用的有兩種：

\begin{itemize}
\item \textbf{列優先}（row-major，也叫以列為主）：先放第 $l_1$ 列的 $n$ 個元素，
再放下一列，依此類推。 C、C++ 的陣列用這種順序。
\item \textbf{行優先}（column-major，以行為主）：先放第 $l_2$ 行的 $m$ 個元素，再
放下一行。 Fortran、MATLAB 的陣列用這種順序。
\end{itemize}

下圖是 $A[1..3, 1..4]$ 的兩種排法，每一格寫的是它的下標「列, 行」，同一列用同一
種顏色。 格子下面的數字是它前面有幾個元素。

\begin{center}
\begin{tikzpicture}
\node[arr-r1] at (0.95,-0.62) {1,1};
\node[arr-r1] at (1.90,-0.62) {1,2};
\node[arr-r1] at (2.85,-0.62) {1,3};
\node[arr-r1] at (3.80,-0.62) {1,4};
\node[arr-r2] at (0.95,-1.24) {2,1};
\node[arr-r2] at (1.90,-1.24) {2,2};
\node[arr-r2] at (2.85,-1.24) {2,3};
\node[arr-r2] at (3.80,-1.24) {2,4};
\node[arr-r3] at (0.95,-1.86) {3,1};
\node[arr-r3] at (1.90,-1.86) {3,2};
\node[arr-r3] at (2.85,-1.86) {3,3};
\node[arr-r3] at (3.80,-1.86) {3,4};
\node[idxlbl, anchor=east] at (0.50,-0.62) {列 1};
\node[idxlbl, anchor=east] at (0.50,-1.24) {列 2};
\node[idxlbl, anchor=east] at (0.50,-1.86) {列 3};
\node[idxlbl] at (0.95,-0.15) {行 1};
\node[idxlbl] at (1.90,-0.15) {行 2};
\node[idxlbl] at (2.85,-0.15) {行 3};
\node[idxlbl] at (3.80,-0.15) {行 4};
\node[anchor=east, font=\small] at (0.5,-3.0) {列優先};
\node[arr-r1] at (1.075,-3.0) {1,1};
\node[idxlbl] at (1.075,-3.5) {0};
\node[arr-r1] at (2.025,-3.0) {1,2};
\node[idxlbl] at (2.025,-3.5) {1};
\node[arr-r1] at (2.975,-3.0) {1,3};
\node[idxlbl] at (2.975,-3.5) {2};
\node[arr-r1] at (3.925,-3.0) {1,4};
\node[idxlbl] at (3.925,-3.5) {3};
\node[arr-r2] at (4.875,-3.0) {2,1};
\node[idxlbl] at (4.875,-3.5) {4};
\node[arr-r2] at (5.825,-3.0) {2,2};
\node[idxlbl] at (5.825,-3.5) {5};
\node[arr-r2] at (6.775,-3.0) {2,3};
\node[idxlbl] at (6.775,-3.5) {6};
\node[arr-r2] at (7.725,-3.0) {2,4};
\node[idxlbl] at (7.725,-3.5) {7};
\node[arr-r3] at (8.675,-3.0) {3,1};
\node[idxlbl] at (8.675,-3.5) {8};
\node[arr-r3] at (9.625,-3.0) {3,2};
\node[idxlbl] at (9.625,-3.5) {9};
\node[arr-r3] at (10.575,-3.0) {3,3};
\node[idxlbl] at (10.575,-3.5) {10};
\node[arr-r3] at (11.525,-3.0) {3,4};
\node[idxlbl] at (11.525,-3.5) {11};
\node[anchor=east, font=\small] at (0.5,-4.6) {行優先};
\node[arr-r1] at (1.075,-4.6) {1,1};
\node[idxlbl] at (1.075,-5.1) {0};
\node[arr-r2] at (2.025,-4.6) {2,1};
\node[idxlbl] at (2.025,-5.1) {1};
\node[arr-r3] at (2.975,-4.6) {3,1};
\node[idxlbl] at (2.975,-5.1) {2};
\node[arr-r1] at (3.925,-4.6) {1,2};
\node[idxlbl] at (3.925,-5.1) {3};
\node[arr-r2] at (4.875,-4.6) {2,2};
\node[idxlbl] at (4.875,-5.1) {4};
\node[arr-r3] at (5.825,-4.6) {3,2};
\node[idxlbl] at (5.825,-5.1) {5};
\node[arr-r1] at (6.775,-4.6) {1,3};
\node[idxlbl] at (6.775,-5.1) {6};
\node[arr-r2] at (7.725,-4.6) {2,3};
\node[idxlbl] at (7.725,-5.1) {7};
\node[arr-r3] at (8.675,-4.6) {3,3};
\node[idxlbl] at (8.675,-5.1) {8};
\node[arr-r1] at (9.625,-4.6) {1,4};
\node[idxlbl] at (9.625,-5.1) {9};
\node[arr-r2] at (10.575,-4.6) {2,4};
\node[idxlbl] at (10.575,-5.1) {10};
\node[arr-r3] at (11.525,-4.6) {3,4};
\node[idxlbl] at (11.525,-5.1) {11};
\end{tikzpicture}
\diagcaption{$3 \times 4$ 的陣列用列優先和行優先排成一維。 列優先時同一列的元素相
鄰；行優先時同一行的元素相鄰，同一列的元素每隔 $3$ 格出現一次。}
\label{fig:arr-order}
\end{center}

\textbf{列優先的公式}。 $A[i, j]$ 前面有哪些元素？ 第 $l_1$ 列到第 $i - 1$ 列全部
排在它前面，共 $i - l_1$ 列，每列 $n$ 個；第 $i$ 列裡面，$A[i, l_2]$ 到
$A[i, j - 1]$ 也在它前面，共 $j - l_2$ 個。 所以
\begin{equation}
\mathrm{loc}(A[i, j]) = \alpha + \bigl[(i - l_1)\,n + (j - l_2)\bigr]\,w ,
\qquad n = u_2 - l_2 + 1 .
\label{eq:arr-row}
\end{equation}

\textbf{行優先的公式}。 把列和行的角色對調：前面有 $j - l_2$ 整行，每行 $m$ 個，
再加上第 $j$ 行裡排在它上面的 $i - l_1$ 個：
\begin{equation}
\mathrm{loc}(A[i, j]) = \alpha + \bigl[(j - l_2)\,m + (i - l_1)\bigr]\,w ,
\qquad m = u_1 - l_1 + 1 .
\label{eq:arr-col}
\end{equation}

兩個公式都只用到「另一個維度」的大小：列優先要乘的是\textbf{行數} $n$（一列有幾
個），行優先要乘的是\textbf{列數} $m$。 列優先的公式裡完全沒有 $m$，所以陣列有幾
列不影響位址，只影響陣列的總大小。 這也是 C 的二維陣列參數可以省略第一維、不能省
略第二維的原因。

\textbf{例 1}。 $A[-2..5,\ 3..9]$，每個元素 $4$ 位元組，$\alpha = 2000$（也就是
$A[-2, 3]$ 的位址）。 求 $A[1, 7]$ 的位址。

先算大小：$m = 5 - (-2) + 1 = 8$ 列，$n = 9 - 3 + 1 = 7$ 行。

\begin{itemize}
\item 列優先：前面有 $(1 - (-2)) \times 7 + (7 - 3) = 21 + 4 = 25$ 個元素，位址是
$2000 + 25 \times 4 = 2100$。
\item 行優先：前面有 $(7 - 3) \times 8 + (1 - (-2)) = 32 + 3 = 35$ 個元素，位址是
$2000 + 35 \times 4 = 2140$。
\end{itemize}

\textbf{例 2}（反推）。 $B[1..30,\ 1..n]$ 以列優先存放，每個元素 $2$ 位元組，$n$
未知。 已知 $\mathrm{loc}(B[3, 5]) = 1130$、$\mathrm{loc}(B[6, 2]) = 1316$，求 $n$ 和
$B[1, 1]$ 的位址。

把兩個位址都寫成公式~(\ref{eq:arr-row}) 再相減，$\alpha$ 就消掉了：
\[
1316 - 1130 = \bigl[(6 - 3)\,n + (2 - 5)\bigr] \times 2
\quad\Longrightarrow\quad 3n - 3 = 93
\quad\Longrightarrow\quad n = 32 .
\]
再代回任一個：$1130 = \alpha + [(3 - 1) \times 32 + (5 - 1)] \times 2 = \alpha + 136$，
所以 $\alpha = 994$。 驗算 $B[6, 2]$：$994 + (5 \times 32 + 1) \times 2 = 994 + 322
= 1316$，相符。 注意 $B[6, 2]$ 的行號比 $B[3, 5]$ 小，$j - l_2$ 那一項相減後是負的，
不要直接寫成 $3n + 3$。

章首的座位問題就是 $S[1..9, 1..7]$，名單上的「第幾個」等於前面的人數加 $1$。 一排
一排列是列優先：$(6 - 1) \times 7 + (4 - 1) + 1 = 39$，第 $39$ 個。 改成先列每排最
左邊的位子，是行優先：$(4 - 1) \times 9 + (6 - 1) + 1 = 33$，第 $33$ 個。

\begin{examnote}
位址題先把題目給的條件抄成一行：上下界 $l_1..u_1$、$l_2..u_2$（下標從 $0$ 開始、
從 $1$ 開始，還是有負的）、每格大小 $w$、起始位址 $\alpha$ 是哪一格的位址、列優先
還是行優先。 每格大小的單位也要看清楚：題目說「每個元素占 $4$ 個位元組」和「每個元
素占 $1$ 個字組（word）、位址以字組為單位」，乘的數不同。 起始位址是十六進位時，
答案也寫成十六進位。 另外，中國大陸的教材「行」「列」的意思和台灣相反（「行」是
row），讀到簡體的題目時，以英文的 row-major、column-major 為準。
\end{examnote}

\subsection{三維與更多維}

三維陣列 $A[l_1..u_1,\ l_2..u_2,\ l_3..u_3]$ 的三個維度大小記作
$d_k = u_k - l_k + 1$（$k = 1, 2, 3$）。 列優先推廣成「最後一個下標變得最快」：先固
定 $i_1 = l_1$、$i_2 = l_2$，讓 $i_3$ 從 $l_3$ 走到 $u_3$；再讓 $i_2$ 加一，依此類
推。 可以把它想成 $d_1$ 張 $d_2 \times d_3$ 的表格疊起來，一張一張用列優先放。
$A[i_1, i_2, i_3]$ 前面有：

\begin{itemize}
\item 前面 $i_1 - l_1$ 張完整的表格，每張 $d_2 d_3$ 個；
\item 同一張表格裡，前面 $i_2 - l_2$ 整列，每列 $d_3$ 個；
\item 同一列裡，前面 $i_3 - l_3$ 個。
\end{itemize}

所以
\begin{equation}
\mathrm{loc}(A[i_1, i_2, i_3]) = \alpha + \bigl[(i_1 - l_1)\,d_2 d_3 + (i_2 - l_2)\,d_3
+ (i_3 - l_3)\bigr]\,w .
\label{eq:arr-3d-row}
\end{equation}

行優先是「第一個下標變得最快」，把下標的順序倒過來：
\begin{equation}
\mathrm{loc}(A[i_1, i_2, i_3]) = \alpha + \bigl[(i_3 - l_3)\,d_1 d_2 + (i_2 - l_2)\,d_1
+ (i_1 - l_1)\bigr]\,w .
\label{eq:arr-3d-col}
\end{equation}

\textbf{例 3}。 $B[1..4,\ -1..3,\ 0..5]$，每個元素 $8$ 位元組，$\alpha = 5000$。 求
$B[3, 2, 4]$ 的位址。 三個維度的大小是 $d_1 = 4$、$d_2 = 5$、$d_3 = 6$，下標減掉下
界後是 $(3 - 1, 2 - (-1), 4 - 0) = (2, 3, 4)$。

\begin{itemize}
\item 列優先：$2 \times (5 \times 6) + 3 \times 6 + 4 = 60 + 18 + 4 = 82$，位址
$5000 + 82 \times 8 = 5656$。
\item 行優先：$4 \times (4 \times 5) + 3 \times 4 + 2 = 80 + 12 + 2 = 94$，位址
$5000 + 94 \times 8 = 5752$。
\end{itemize}

\textbf{$d$ 維的情形}。 $A[l_1..u_1, \ldots, l_d..u_d]$，記 $e_k = i_k - l_k$（第 $k$
個下標減掉下界）。 列優先時，$A[i_1, \ldots, i_d]$ 前面的元素有幾個？ 一個元素排在
它前面，若且唯若兩者的下標從左邊比起，第一個不同的位置上前者比較小（就像字典裡的
順序）。 把這些元素依「第一個不同的位置 $k$」分類：前 $k - 1$ 個下標和
$A[i_1, \ldots, i_d]$ 相同，第 $k$ 個下標比 $i_k$ 小（有 $e_k$ 種），後面的下標隨
便（有 $d_{k+1} \cdots d_d$ 種）。 加起來，前面共有
\begin{equation}
\mathrm{off} = \sum_{k=1}^{d} e_k \prod_{t=k+1}^{d} d_t
= \Bigl(\cdots\bigl((e_1\,d_2 + e_2)\,d_3 + e_3\bigr)\cdots\Bigr)\,d_d + e_d
\label{eq:arr-dd}
\end{equation}
個元素，位址是 $\alpha + \mathrm{off} \cdot w$。 右邊那種一層包一層的寫法，就是
第~\ref{sec:arr-poly}~節求多項式值用的 Horner 法：從左到右，每一步把目前的值乘上
下一個維度的大小，再加上下一個 $e_k$。 照這個寫法算，只要 $d - 1$ 次乘法。

\begin{pseudo}[label=alg:arr-offset]{列優先的位移（$d$ 維）}
\Function{Row-Major-Offset}{$d, \mathit{lo}, \mathit{hi}, \mathit{idx}$}
    \State $\mathit{off} \gets 0$
    \For{$k \gets 1$ \To $d$}
        \State $\mathit{off} \gets \mathit{off} \cdot (\mathit{hi}[k] - \mathit{lo}[k] + 1)
        + (\mathit{idx}[k] - \mathit{lo}[k])$
    \EndFor
    \State \Return $\mathit{off}$
\EndFunction
\end{pseudo}

第 4 行第一次執行時 $\mathit{off} = 0$，乘法其實可以省掉，所以實際需要的乘法是
$d - 1$ 次。 行優先只要把 $k$ 的順序改成 $d$ \textbf{downto} $1$。

公式~(\ref{eq:arr-dd}) 也告訴我們每個下標加 $1$ 時位址變多少，這個量叫做那一維的
\textbf{步長}（stride）。 列優先時第 $k$ 維的步長是 $w \cdot d_{k+1} \cdots d_d$，最
後一維的步長是 $w$；行優先時第 $k$ 維的步長是 $w \cdot d_1 \cdots d_{k-1}$，第一維
的步長是 $w$。 例 3 的 $B$ 用列優先存放時，三個維度的步長依序是 $240$、$48$、$8$
位元組。

\subsection{C 的陣列}

C 的多維陣列一律是列優先、下界一律是 $0$。 宣告 \texttt{int a[R][C]} 之後，
\texttt{a[i][j]} 的位址是 \texttt{a} 的位址加上 $(i \cdot C + j) \cdot
\texttt{sizeof(int)}$ 個位元組，就是公式~(\ref{eq:arr-row}) 取 $l_1 = l_2 = 0$、
$n = C$。 編譯器算這個位址時要知道 $C$，所以二維陣列當參數傳進函式時，第二維的大
小一定要寫出來。

\begin{cbookrecap}[第~6~章 6.11 節、第~8~章 8.6 節、第~9~章 9.6、9.7 節、第~14~章 14.4 節]
二維陣列在記憶體裡是先放完第 $0$ 列再放第 $1$ 列的一長條，\texttt{grid[i][j]} 的
位移是 \texttt{i * C + j}。 指標加上整數 $k$，移動的是 $k$ 個「它指的型別」的大小，
不是 $k$ 個位元組；兩個指向同一個陣列的指標相減，得到的是相差幾個元素。 想算位元組
的差，先把兩個指標都轉成 \texttt{char *} 再相減，結果的型別是
\texttt{ptrdiff\_t}，用 \texttt{\%td} 印。 \texttt{a[i][j]} 會展開成
\texttt{*(*(a + i) + j)}。 多維陣列傳進函式時，第一維可以省略，後面的維度一定要寫
清楚，例如 \texttt{void f(int n, int a[][8])}。
\end{cbookrecap}

C 不能宣告下界不是 $0$ 的陣列。 需要 $A[l..u]$ 的時候，配置 $u - l + 1$ 格，存取時
自己把下標減掉 $l$：寫一個函式或巨集回傳 \texttt{buf[i - l]}。 有人會先算一個
\texttt{p = buf - l}，之後直接寫 \texttt{p[i]}。 $l > 0$ 時 \texttt{buf - l} 指到陣
列開頭之前，C 標準規定這樣的指標運算是未定義行為，即使之後再加回來也一樣，本書不這
樣寫。

下面的程式把公式~(\ref{eq:arr-dd}) 寫成兩個函式，算出例 1、例 3 和一維例子的位
址；最後拿一個真正的 C 三維陣列，比較公式算出的位移和編譯器實際的位移。

\begin{example}[label=ex:arr-addr]{c}{(addr\_check.c) 公式算的位址和實際的位址}
#include <stdio.h>
#include <stddef.h>

// d 維陣列，第 k 個下標的範圍是 lo[k]..hi[k]。
// 回傳 idx 這一格前面有幾個元素。列優先：最後一個下標變得最快。
long row_major(int d, const int *lo, const int *hi, const int *idx)
{
    long off = 0;
    for (int k = 0; k < d; k++)
        off = off * (hi[k] - lo[k] + 1) + (idx[k] - lo[k]);
    return off;
}

// 行優先：第一個下標變得最快，等於把下標倒過來做列優先。
long col_major(int d, const int *lo, const int *hi, const int *idx)
{
    long off = 0;
    for (int k = d - 1; k >= 0; k--)
        off = off * (hi[k] - lo[k] + 1) + (idx[k] - lo[k]);
    return off;
}

static double c[4][5][6];       // C 的三維陣列，下標從 0 開始

int main(void)
{
    // 一維：T[-10..30]，每格 4 bytes，T[-10] 在 0x1A00
    int lo1[] = {-10}, hi1[] = {30}, t[] = {17};
    printf("T[17]: 0x%lX\n", 0x1A00 + 4 * row_major(1, lo1, hi1, t));

    // 二維：A[-2..5, 3..9]，每格 4 bytes，A[-2,3] 在 2000
    int lo2[] = {-2, 3}, hi2[] = {5, 9}, a[] = {1, 7};
    printf("A[1,7]: row-major %ld, column-major %ld\n",
           2000 + 4 * row_major(2, lo2, hi2, a),
           2000 + 4 * col_major(2, lo2, hi2, a));

    // 三維：B[1..4, -1..3, 0..5]，每格 8 bytes，B[1,-1,0] 在 5000
    int lo3[] = {1, -1, 0}, hi3[] = {4, 3, 5}, b[] = {3, 2, 4};
    printf("B[3,2,4]: row-major %ld, column-major %ld\n",
           5000 + 8 * row_major(3, lo3, hi3, b),
           5000 + 8 * col_major(3, lo3, hi3, b));

    // 和編譯器實際的位址比較：公式算的位移 vs. 指標相減
    int loc[] = {0, 0, 0}, hic[] = {3, 4, 5};
    int probe[][3] = {{0, 0, 1}, {0, 1, 0}, {1, 0, 0},
                      {2, 3, 4}, {3, 4, 5}};
    for (int p = 0; p < 5; p++) {
        int *x = probe[p];
        long f = (long)sizeof(double) * row_major(3, loc, hic, x);
        char *base = (char *)&c[0][0][0];
        ptrdiff_t real = (char *)&c[x[0]][x[1]][x[2]] - base;
        printf("c[%d][%d][%d]: formula %ld, actual %td\n",
               x[0], x[1], x[2], f, real);
    }
    return 0;
}
\end{example}
\iosample{}{T[17]: 0x1A6C \\ A[1,7]: row-major 2100, column-major 2140 \\ B[3,2,4]: row-major 5656, column-major 5752 \\ c[0][0][1]: formula 8, actual 8 \\ c[0][1][0]: formula 48, actual 48 \\ c[1][0][0]: formula 240, actual 240 \\ c[2][3][4]: formula 656, actual 656 \\ c[3][4][5]: formula 952, actual 952}

最後五行左邊是公式~(\ref{eq:arr-dd}) 乘上 \texttt{sizeof(double)}（$8$），右邊是
編譯器實際放的位置，五組都一樣。 \texttt{c[0][0][1]}、\texttt{c[0][1][0]}、
\texttt{c[1][0][0]} 三行正好是三個維度的步長：$8$、$6 \times 8 = 48$、
$5 \times 6 \times 8 = 240$。

\textbf{指標陣列}。 另一種存二維資料的方法是先配置一個指標陣列
\texttt{int *row[R]}，每一列各自用 \texttt{malloc} 配置，\texttt{row[i]} 指向第 $i$
列（《C 程式設計基礎》第~9~章 9.8 節）。 存取 \texttt{row[i][j]} 時先讀出
\texttt{row[i]} 這個指標，再往後數 $j$ 格；不用乘上行數，但多讀了一次記憶體。 它的
好處是每一列可以不一樣長，下一節的三角矩陣就可以一列配置 $i$ 格。 壞處是各列不一
定相鄰，不能用一個公式從頭算到任何一格，整個陣列也不能一次複製。

\textbf{走訪的順序}。 第~1~章提過，真正的電腦有快取，讀相鄰的位址比跳著讀快。 C 的
陣列是列優先，外層迴圈跑列、內層迴圈跑行時，連續兩次存取的位址相差 $w$；兩層對調
之後，連續兩次存取相差一整列 $n \cdot w$ 個位元組。 在 RAM 模型裡兩種寫法的步數一
樣，實際上 $n$ 很大時後者可能慢上好幾倍。 寫雙層迴圈時，讓最內層的迴圈跑最後一個
下標。

\section{特殊矩陣的壓縮存放}
\label{sec:arr-special}

有些 $n \times n$ 矩陣有一大片位置一定是 $0$，或是一半的元素和另一半相同。 照上
一節的方法存，要用 $n^2$ 格，其中很多格的內容早就知道了。 這一節把真正需要的元素
依某個順序放進一個一維陣列 $B$，再推導「$A[i, j]$ 放在 $B$ 的第幾格」的公式。 有
了公式，讀寫一格仍然是 $O(1)$ 時間，只是省下了空間。

這一節的矩陣下標都從 $1$ 開始，$B$ 也從 $1$ 開始，「$A[i, j]$ 是 $B$ 的第 $k$ 格」
的位址是 $\alpha + (k - 1)\,w$；下標從 $0$ 開始的版本，最後另外整理。

\begin{definition}[特殊矩陣]
\label{def:arr-special}
$A$ 是 $n \times n$ 的矩陣。
\begin{itemize}
\item $i < j$ 時 $A[i, j] = 0$，叫做\textbf{下三角矩陣}（lower triangular matrix）：
非零元素只在對角線和它的左下方。
\item $i > j$ 時 $A[i, j] = 0$，叫做\textbf{上三角矩陣}（upper triangular matrix）。
\item 對所有 $i, j$ 都有 $A[i, j] = A[j, i]$，叫做\textbf{對稱矩陣}（symmetric
matrix）。
\item $|i - j| > 1$ 時 $A[i, j] = 0$，叫做\textbf{三對角矩陣}（tridiagonal matrix）；
更一般地，$|i - j| > a$ 時 $A[i, j] = 0$，叫做半頻寬 $a$ 的\textbf{帶狀矩陣}
（band matrix）。
\end{itemize}
\end{definition}

\subsection{下三角矩陣}

下三角矩陣第 $i$ 列需要存的是 $A[i, 1], \ldots, A[i, i]$，共 $i$ 個，全部加起來
$1 + 2 + \cdots + n = n(n + 1)/2$ 個，大約是 $n^2$ 的一半。

\textbf{列優先}。 $A[i, j]$（$i \ge j$）前面有第 $1$ 到第 $i - 1$ 列的全部，共
$1 + 2 + \cdots + (i - 1) = i(i - 1)/2$ 個，再加上第 $i$ 列的 $A[i, 1]$ 到
$A[i, j - 1]$ 共 $j - 1$ 個。 它是第
\begin{equation}
k = \frac{i(i - 1)}{2} + j
\label{eq:arr-lower-row}
\end{equation}
格。 這個公式和 $n$ 無關：矩陣再大，前幾列的位置都不變。

\textbf{行優先}。 第 $c$ 行要存的是 $A[c, c], \ldots, A[n, c]$，共 $n - c + 1$ 個。
$A[i, j]$ 前面有第 $1$ 到第 $j - 1$ 行的全部，共
\[
\sum_{c=1}^{j-1} (n - c + 1) = (j - 1)(n + 1) - \frac{(j - 1)j}{2}
= \frac{(j - 1)(2n - j + 2)}{2}
\]
個，再加上第 $j$ 行裡排在它上面的 $A[j, j], \ldots, A[i - 1, j]$，共 $i - j$ 個。
它是第
\begin{equation}
k = \frac{(j - 1)(2n - j + 2)}{2} + (i - j + 1)
\label{eq:arr-lower-col}
\end{equation}
格。 這個公式就和 $n$ 有關了。

下圖是 $n = 5$ 時兩種順序的編號，灰色的 $0$ 不存。

\begin{center}
\begin{tikzpicture}
\node[arr-tri] at (0.70,-0.70) {1};
\node[arr-zero] at (1.40,-0.70) {0};
\node[arr-zero] at (2.10,-0.70) {0};
\node[arr-zero] at (2.80,-0.70) {0};
\node[arr-zero] at (3.50,-0.70) {0};
\node[arr-tri] at (0.70,-1.40) {2};
\node[arr-tri] at (1.40,-1.40) {3};
\node[arr-zero] at (2.10,-1.40) {0};
\node[arr-zero] at (2.80,-1.40) {0};
\node[arr-zero] at (3.50,-1.40) {0};
\node[arr-tri] at (0.70,-2.10) {4};
\node[arr-tri] at (1.40,-2.10) {5};
\node[arr-tri] at (2.10,-2.10) {6};
\node[arr-zero] at (2.80,-2.10) {0};
\node[arr-zero] at (3.50,-2.10) {0};
\node[arr-tri] at (0.70,-2.80) {7};
\node[arr-tri] at (1.40,-2.80) {8};
\node[arr-tri] at (2.10,-2.80) {9};
\node[arr-tri] at (2.80,-2.80) {10};
\node[arr-zero] at (3.50,-2.80) {0};
\node[arr-tri] at (0.70,-3.50) {11};
\node[arr-tri] at (1.40,-3.50) {12};
\node[arr-tri] at (2.10,-3.50) {13};
\node[arr-tri] at (2.80,-3.50) {14};
\node[arr-tri] at (3.50,-3.50) {15};
\node[font=\small] at (2.10,-4.25) {列優先};
\node[arr-tri] at (5.30,-0.70) {1};
\node[arr-zero] at (6.00,-0.70) {0};
\node[arr-zero] at (6.70,-0.70) {0};
\node[arr-zero] at (7.40,-0.70) {0};
\node[arr-zero] at (8.10,-0.70) {0};
\node[arr-tri] at (5.30,-1.40) {2};
\node[arr-tri] at (6.00,-1.40) {6};
\node[arr-zero] at (6.70,-1.40) {0};
\node[arr-zero] at (7.40,-1.40) {0};
\node[arr-zero] at (8.10,-1.40) {0};
\node[arr-tri] at (5.30,-2.10) {3};
\node[arr-tri] at (6.00,-2.10) {7};
\node[arr-tri] at (6.70,-2.10) {10};
\node[arr-zero] at (7.40,-2.10) {0};
\node[arr-zero] at (8.10,-2.10) {0};
\node[arr-tri] at (5.30,-2.80) {4};
\node[arr-tri] at (6.00,-2.80) {8};
\node[arr-tri] at (6.70,-2.80) {11};
\node[arr-tri] at (7.40,-2.80) {13};
\node[arr-zero] at (8.10,-2.80) {0};
\node[arr-tri] at (5.30,-3.50) {5};
\node[arr-tri] at (6.00,-3.50) {9};
\node[arr-tri] at (6.70,-3.50) {12};
\node[arr-tri] at (7.40,-3.50) {14};
\node[arr-tri] at (8.10,-3.50) {15};
\node[font=\small] at (6.70,-4.25) {行優先};
\end{tikzpicture}
\diagcaption{$5 \times 5$ 下三角矩陣的非零格，依列優先（左）和行優先（右）放進
$B[1..15]$ 時的編號。}
\label{fig:arr-tri}
\end{center}

\textbf{例 4}。 $6 \times 6$ 的下三角矩陣，每個元素 $4$ 位元組，$B[1]$ 的位址是
$3000$。 求 $A[5, 3]$ 的位址。

\begin{itemize}
\item 列優先：$k = 5 \times 4 / 2 + 3 = 13$，位址 $3000 + 12 \times 4 = 3048$。
\item 行優先：$k = 2 \times (12 - 3 + 2)/2 + (5 - 3 + 1) = 11 + 3 = 14$，位址
$3000 + 13 \times 4 = 3052$。
\end{itemize}

可以用圖驗算：$n = 5$ 時列優先的 $A[5, 3]$ 是 $13$，和 $n$ 無關，所以 $n = 6$ 時也
是 $13$；行優先時第 $1$、$2$ 行在 $n = 6$ 時各有 $6$、$5$ 個，$A[5, 3]$ 前面還有
$A[3, 3]$、$A[4, 3]$，所以是 $6 + 5 + 2 + 1 = 14$。

\subsection{上三角矩陣}

上三角矩陣不必重新推導。 $A$ 是上三角時，它的轉置 $A^{\mathsf T}$（$A^{\mathsf
T}[j, i] = A[i, j]$）是下三角。 用列優先存 $A$，等於一列一列存 $A^{\mathsf T}$ 的
行，也就是用行優先存 $A^{\mathsf T}$。 所以「上三角、列優先」的公式就是「下三角、
行優先」的公式把 $i$、$j$ 對調；另外兩種也一樣。 下表把四種情形整理在一起。

\begin{center}
\renewcommand{\arraystretch}{1.6}
\begin{tabular}{lll}
\toprule
 & 列優先 & 行優先 \\
\midrule
下三角（$i \ge j$）& $\dfrac{i(i - 1)}{2} + j$ & $\dfrac{(j - 1)(2n - j + 2)}{2} + (i - j + 1)$ \\
上三角（$i \le j$）& $\dfrac{(i - 1)(2n - i + 2)}{2} + (j - i + 1)$ & $\dfrac{j(j - 1)}{2} + i$ \\
\bottomrule
\end{tabular}
\end{center}

這些公式不必硬背，推導的步驟都一樣：先問「它前面有幾整列（或幾整行），每一列各
有幾個」，再加上同一列裡排在它前面的個數。 每一列個數固定是 $1, 2, 3, \ldots$ 的那兩種情形，公式和
$n$ 無關；每一列個數是 $n, n - 1, \ldots$ 的那兩種，公式裡有 $n$。

\textbf{下標從 $0$ 開始}。 C 的程式裡 $i$、$j$、$k$ 都從 $0$ 開始。 下三角、列優先時，
第 $0$ 到第 $i - 1$ 列共有 $1 + 2 + \cdots + i = i(i + 1)/2$ 個，所以 $A[i][j]$
放在 \texttt{b[i * (i + 1) / 2 + j]}。 不要把 $1$ 起算的公式~(\ref{eq:arr-lower-row})
直接搬過來再減 $1$，代入 $i$ 的值也要跟著換。

\begin{examnote}
特殊矩陣的題目先確認四件事：上三角還是下三角、列優先還是行優先、下標和陣列 $B$ 從
$0$ 還是從 $1$ 開始、對角線有沒有存（距離表這類對角線全是 $0$ 的矩陣常常不存）。
推出公式之後，代 $A[1, 1]$（或 $A[0][0]$）和這一列的最後一格驗算一次，看是不是
$B$ 的第一格和預期的位置。
\end{examnote}

\subsection{對稱矩陣：距離表}

對稱矩陣只要存下三角（含對角線），讀 $A[i, j]$ 時如果 $i < j$，就改讀 $A[j, i]$。
空間從 $n^2$ 降到 $n(n + 1)/2$。

校園裡幾個地點兩兩之間的步行距離是一個對稱矩陣，而且對角線全是 $0$（自己到自己），
所以連對角線都不用存，只存 $i > j$ 的 $n(n - 1)/2$ 格。 下標從 $0$ 開始時，第 $0$ 列
沒有要存的格子，第 $1$ 到第 $i - 1$ 列各存 $1, \ldots, i - 1$ 格，所以 $A[i][j]$
（$i > j$）放在 \texttt{d[i * (i - 1) / 2 + j]}。

\begin{cbookrecap}[第~8~章 8.6 節、第~9~章 9.3 節、第~19~章 19.3 節]
一維陣列當參數傳進函式時傳的是位址，\texttt{const int *d} 表示函式只讀不改。
\texttt{malloc(size * sizeof *d)} 配置 \texttt{size} 個 \texttt{int}，要不到時回傳
\texttt{NULL}，用完要 \texttt{free}。
\end{cbookrecap}

\begin{example}[label=ex:arr-tri-dist]{c}{(tri\_dist.c) 三角矩陣存距離表}
#include <stdio.h>
#include <stdlib.h>

// 對稱、對角線是 0 的距離表，只存 i > j 的格子（下標從 0 開始）。
// (i, j) 和 (j, i) 都對應到 d[] 的同一格。
static long tri_index(int i, int j)
{
    if (i < j) {                // 對稱：換成下三角的那一格
        int tmp = i; i = j; j = tmp;
    }
    return (long)i * (i - 1) / 2 + j;
}

int dist_get(const int *d, int i, int j)
{
    return i == j ? 0 : d[tri_index(i, j)];
}

int main(void)
{
    int n = 5;                  // 0 lib, 1 gym, 2 cafe, 3 eng, 4 dorm
    int pairs[][3] = {          // 地點、地點、距離（公尺），順序任意
        {0, 1, 450}, {2, 0, 300}, {1, 2, 380}, {0, 3, 520},
        {3, 1, 610}, {2, 3, 270}, {4, 0, 800}, {1, 4, 350},
        {4, 2, 540}, {3, 4, 700},
    };
    long size = (long)n * (n - 1) / 2;
    int *d = malloc(size * sizeof *d);
    if (d == NULL)
        return 1;
    for (int p = 0; p < 10; p++)
        d[tri_index(pairs[p][0], pairs[p][1])] = pairs[p][2];

    printf("packed:");
    for (long k = 0; k < size; k++)
        printf(" %d", d[k]);
    printf("\n");
    for (int i = 0; i < n; i++) {         // 還原成完整的表
        for (int j = 0; j < n; j++)
            printf(j > 0 ? " %d" : "%d", dist_get(d, i, j));
        printf("\n");
    }
    for (long m = 50; m <= 5000; m *= 10)
        printf("n = %ld: full %ld, packed %ld\n",
               m, m * m, m * (m - 1) / 2);
    free(d);
    return 0;
}
\end{example}
\iosample{}{packed: 450 300 380 520 610 270 800 350 540 700 \\ 0 450 300 520 800 \\ 450 0 380 610 350 \\ 300 380 0 270 540 \\ 520 610 270 0 700 \\ 800 350 540 700 0 \\ n = 50: full 2500, packed 1225 \\ n = 500: full 250000, packed 124750 \\ n = 5000: full 25000000, packed 12497500}

\texttt{pairs} 裡的地點順序是亂的，\texttt{tri\_index} 會先把 $(i, j)$ 換成 $i > j$
的那一格。 第一行印出的 \texttt{d[0..9]} 依序是 $(1,0)$、$(2,0)$、$(2,1)$、$(3,0)$、
\ldots、$(4,3)$ 的距離，正是下三角、列優先的順序；接著還原出完整的 $5 \times 5$ 表，
它對稱、對角線是 $0$。 最後三行是空間的比較，$n$ 很大時省下一半多一點。

\subsection{三對角與帶狀矩陣}

三對角矩陣的第 $1$ 列和第 $n$ 列各有 $2$ 個可能非零的元素，其他列各 $3$ 個，共
$3n - 2$ 個。 用列優先放進 $B$：第 $1$ 列占 $B[1..2]$；第 $i$ 列（$i \ge 2$）前面
有 $2 + 3(i - 2) = 3i - 4$ 個，這一列的三個元素是 $A[i, i - 1]$、$A[i, i]$、
$A[i, i + 1]$，$A[i, j]$ 是其中第 $j - i + 2$ 個。 所以
\[
k = (3i - 4) + (j - i + 2) = 2i + j - 2 .
\]
代 $i = 1$ 檢查：$A[1, 1]$、$A[1, 2]$ 得到 $1$、$2$，第 $1$ 列也適用。

半頻寬 $a$ 更大的帶狀矩陣，前 $a$ 列和最後 $a$ 列的長度都不一樣，緊密存放的公式會
分段，不好用。 實務上常改成每一列都留 $2a + 1$ 格：開一個 $n \times (2a + 1)$ 的陣
列 $E$，把 $A[i, j]$ 放在 $E[i,\ j - i + a + 1]$，也就是第 $i$ 列只存從 $A[i, i - a]$
到 $A[i, i + a]$ 這一段。 前 $a$ 列左邊、最後 $a$ 列右邊會有幾格用不到（共
$a(a + 1)$ 格），換來的是 $E$ 是普通的矩形陣列，用列優先攤平時第 $k = (i - 1)(2a + 1)
+ (j - i + a + 1)$ 格，沒有例外。 下圖是 $n = 6$、$a = 1$ 的兩種存法，右邊打叉的兩
格是浪費的。

\begin{center}
\begin{tikzpicture}
\node[arr-tri] at (0.70,-0.70) {1};
\node[arr-tri] at (1.40,-0.70) {2};
\node[arr-zero] at (2.10,-0.70) {0};
\node[arr-zero] at (2.80,-0.70) {0};
\node[arr-zero] at (3.50,-0.70) {0};
\node[arr-zero] at (4.20,-0.70) {0};
\node[arr-tri] at (0.70,-1.40) {3};
\node[arr-tri] at (1.40,-1.40) {4};
\node[arr-tri] at (2.10,-1.40) {5};
\node[arr-zero] at (2.80,-1.40) {0};
\node[arr-zero] at (3.50,-1.40) {0};
\node[arr-zero] at (4.20,-1.40) {0};
\node[arr-zero] at (0.70,-2.10) {0};
\node[arr-tri] at (1.40,-2.10) {6};
\node[arr-tri] at (2.10,-2.10) {7};
\node[arr-tri] at (2.80,-2.10) {8};
\node[arr-zero] at (3.50,-2.10) {0};
\node[arr-zero] at (4.20,-2.10) {0};
\node[arr-zero] at (0.70,-2.80) {0};
\node[arr-zero] at (1.40,-2.80) {0};
\node[arr-tri] at (2.10,-2.80) {9};
\node[arr-tri] at (2.80,-2.80) {10};
\node[arr-tri] at (3.50,-2.80) {11};
\node[arr-zero] at (4.20,-2.80) {0};
\node[arr-zero] at (0.70,-3.50) {0};
\node[arr-zero] at (1.40,-3.50) {0};
\node[arr-zero] at (2.10,-3.50) {0};
\node[arr-tri] at (2.80,-3.50) {12};
\node[arr-tri] at (3.50,-3.50) {13};
\node[arr-tri] at (4.20,-3.50) {14};
\node[arr-zero] at (0.70,-4.20) {0};
\node[arr-zero] at (1.40,-4.20) {0};
\node[arr-zero] at (2.10,-4.20) {0};
\node[arr-zero] at (2.80,-4.20) {0};
\node[arr-tri] at (3.50,-4.20) {15};
\node[arr-tri] at (4.20,-4.20) {16};
\node[font=\small] at (2.45,-4.95) {$A$ 的非零格依列優先編號};
\node[arr-zero, minimum width=1.05cm] at (7.25,-0.70) {$\times$};
\node[arr-tri, minimum width=1.05cm] at (8.30,-0.70) {$a_{11}$};
\node[arr-tri, minimum width=1.05cm] at (9.35,-0.70) {$a_{12}$};
\node[arr-tri, minimum width=1.05cm] at (7.25,-1.40) {$a_{21}$};
\node[arr-tri, minimum width=1.05cm] at (8.30,-1.40) {$a_{22}$};
\node[arr-tri, minimum width=1.05cm] at (9.35,-1.40) {$a_{23}$};
\node[arr-tri, minimum width=1.05cm] at (7.25,-2.10) {$a_{32}$};
\node[arr-tri, minimum width=1.05cm] at (8.30,-2.10) {$a_{33}$};
\node[arr-tri, minimum width=1.05cm] at (9.35,-2.10) {$a_{34}$};
\node[arr-tri, minimum width=1.05cm] at (7.25,-2.80) {$a_{43}$};
\node[arr-tri, minimum width=1.05cm] at (8.30,-2.80) {$a_{44}$};
\node[arr-tri, minimum width=1.05cm] at (9.35,-2.80) {$a_{45}$};
\node[arr-tri, minimum width=1.05cm] at (7.25,-3.50) {$a_{54}$};
\node[arr-tri, minimum width=1.05cm] at (8.30,-3.50) {$a_{55}$};
\node[arr-tri, minimum width=1.05cm] at (9.35,-3.50) {$a_{56}$};
\node[arr-tri, minimum width=1.05cm] at (7.25,-4.20) {$a_{65}$};
\node[arr-tri, minimum width=1.05cm] at (8.30,-4.20) {$a_{66}$};
\node[arr-zero, minimum width=1.05cm] at (9.35,-4.20) {$\times$};
\node[idxlbl] at (7.25,-0.2) {1};
\node[idxlbl] at (8.30,-0.2) {2};
\node[idxlbl] at (9.35,-0.2) {3};
\node[idxlbl, anchor=east] at (6.60,-0.70) {1};
\node[idxlbl, anchor=east] at (6.60,-1.40) {2};
\node[idxlbl, anchor=east] at (6.60,-2.10) {3};
\node[idxlbl, anchor=east] at (6.60,-2.80) {4};
\node[idxlbl, anchor=east] at (6.60,-3.50) {5};
\node[idxlbl, anchor=east] at (6.60,-4.20) {6};
\node[font=\small] at (8.30,-4.95) {每列 $3$ 格的 $E[1..6, 1..3]$};
\end{tikzpicture}
\diagcaption{三對角矩陣的兩種存法。 左：只存非零格，第 $k = 2i + j - 2$ 格。 右：每
列留 $3$ 格，$A[i, j]$ 放在 $E[i, j - i + 2]$，多用了 $2$ 格。}
\label{fig:arr-band}
\end{center}

$a = 1$ 時右邊的攤平公式是 $(i - 1) \times 3 + (j - i + 2) = 2i + j - 1$，比緊密存放
多 $1$，差的就是第 $1$ 列左邊浪費的那一格。

\section{稀疏矩陣}
\label{sec:arr-sparse}

一門課有 $6$ 位同學、$7$ 次實驗課，助教只在同學上台示範時加分，表上大部分是 $0$：
\[
A = \begin{pmatrix}
0 & 0 & 3 & 0 & 0 & 1 & 0 \\
2 & 0 & 0 & 0 & 0 & 0 & 0 \\
0 & 0 & 5 & 0 & 0 & 0 & 4 \\
0 & 1 & 0 & 0 & 0 & 0 & 0 \\
0 & 0 & 2 & 0 & 0 & 3 & 0 \\
1 & 0 & 0 & 0 & 0 & 0 & 0
\end{pmatrix}
\]
$42$ 格裡只有 $9$ 格不是 $0$。 非零元素只占很小一部分的矩陣叫做\textbf{稀疏矩陣}
（sparse matrix）。 「很小」沒有固定的門檻，看的是只存非零元素划不划算：$1000
\times 1000$ 的矩陣只有幾千個非零元素時，存滿 $10^6$ 格顯然浪費；特殊矩陣的零集中
在固定的位置，可以用上一節的公式，稀疏矩陣的零散在各處，要另外記下每個非零元素的
位置。

\subsection{三元組表示法}

每個非零元素記成一個\textbf{三元組}（triple）$(\text{列}, \text{行}, \text{值})$，
整個矩陣是一個三元組陣列，再加上列數、行數和非零項的個數 $t$。 三元組依列號由小到
大排，列號相同時依行號由小到大，也就是列優先的順序。 這一節的列號、行號和三元組陣
列都從 $0$ 開始，和 Aki-Judge 的題目一致；虛擬碼也照寫 $[0..t-1]$，C 程式的輸出才
能直接和手算對照。

上面的加分表寫成三元組（$m = 6$ 列、$n = 7$ 行、$t = 9$）：

\begin{center}
\begin{tabular}{c|ccccccccc}
\toprule
位置 $p$ & 0 & 1 & 2 & 3 & 4 & 5 & 6 & 7 & 8 \\
\midrule
列 & 0 & 0 & 1 & 2 & 2 & 3 & 4 & 4 & 5 \\
行 & 2 & 5 & 0 & 2 & 6 & 1 & 2 & 5 & 0 \\
值 & 3 & 1 & 2 & 5 & 4 & 1 & 2 & 3 & 1 \\
\bottomrule
\end{tabular}
\end{center}

虛擬碼把一個稀疏矩陣寫成物件 $A$，欄位有 $A.\mathit{rows}$、$A.\mathit{cols}$、
$A.t$ 和三元組陣列 $A.T[0..t-1]$，每個三元組有 $\mathit{row}$、$\mathit{col}$、
$\mathit{val}$ 三個欄位。

\textbf{空間}。 列號、行號、值各占一個整數時，三元組表示法用 $3t$ 個整數（加上記
列數、行數、$t$ 的 $3$ 個），存滿要 $mn$ 個。 $3t + 3 < mn$ 時才划算，加分表的
$3 \times 9 + 3 = 30 < 42$，省得不多；$1000 \times 1000$ 有 $5000$ 個非零元素時，是
$15003$ 對 $10^6$。

\textbf{讀一格}。 三元組依（列, 行）排好，讀 $A[i, j]$ 可以在三元組陣列上做二分搜
尋（第~2~章用過它的遞迴版，第~17~章細講），$O(\lg t)$；找不到就是 $0$。 \textbf{寫一格}比較麻煩：把一個 $0$
改成非零，要插入一個三元組，後面的全部往後挪；把非零改成 $0$ 要刪除，後面的往前
挪。 最壞 $\Theta(t)$。 矩陣常常要改時，第~\ref{sec:ll-poly}~節的正交串列比較適合；這一節
處理的是建好之後只讀的矩陣。

\subsection{轉置}

$m \times n$ 矩陣 $A$ 的\textbf{轉置}（transpose）$A^{\mathsf T}$ 是 $n \times m$
的矩陣，$A^{\mathsf T}[j, i] = A[i, j]$。 存滿的矩陣轉置要 $\Theta(mn)$。 三元組表
示法下，把每個三元組的列和行對調只要 $\Theta(t)$，但結果的順序亂了：對調後的列號
是原來的行號，原本依列排好的三元組，現在依行排才對。 加分表對調後依序是
$(2,0)$、$(5,0)$、$(0,1)$……，不是列優先的順序。 後面的相乘、二分搜尋都需要排好的
三元組，所以轉置要連順序一起處理。

最直接的方法：結果的第 $0$ 列是 $A$ 的第 $0$ 行，就把 $A$ 的三元組從頭掃一遍，挑
出行號是 $0$ 的；再掃一遍挑行號是 $1$ 的，依此類推。

\begin{pseudo}[label=alg:arr-simple-transpose]{簡單轉置}
\Function{Simple-Transpose}{$A$}
    \State $B.\mathit{rows} \gets A.\mathit{cols}$；$B.\mathit{cols} \gets A.\mathit{rows}$；$B.t \gets A.t$
    \State $q \gets 0$
    \For{$c \gets 0$ \To $A.\mathit{cols} - 1$}
        \For{$p \gets 0$ \To $A.t - 1$}
            \If{$A.T[p].\mathit{col} = c$}
                \State $B.T[q] \gets (c,\ A.T[p].\mathit{row},\ A.T[p].\mathit{val})$
                \State $q \gets q + 1$
            \EndIf
        \EndFor
    \EndFor
    \State \Return $B$
\EndFunction
\end{pseudo}

結果依列排好，因為第 4 行的 $c$ 由小到大；同一列裡依行排好，因為第 5 行由前往後
掃，行號相同的三元組在 $A$ 裡依列號由小到大出現，對調後就是依行號由小到大。 第 6
行的比較執行 $A.\mathit{cols} \cdot t$ 次，加分表是 $7 \times 9 = 63$ 次。 時間
$\Theta(n \cdot t)$。 非零元素少時還好，但 $t$ 接近 $mn$ 時是 $\Theta(mn^2)$，比存
滿的矩陣直接轉置還慢。

\subsection{快速轉置}

簡單轉置浪費在「每一行都掃一遍」。 換個想法：如果事先知道 $A$ 的第 $c$ 行有幾個
非零元素，就知道結果的第 $c$ 列從 $B$ 的第幾格開始放，接著只要把 $A$ 掃一遍，每個
三元組直接放到它該去的地方。 這個方法叫\textbf{快速轉置}（fast transpose）。

用兩個長度 $n$ 的陣列：$\mathit{cnt}[c]$ 是 $A$ 的第 $c$ 行有幾個非零元素，
$\mathit{start}[c]$ 是結果第 $c$ 列下一個要放的位置。 第 $c$ 列前面是第 $0$ 到第
$c - 1$ 列，所以一開始 $\mathit{start}[c] = \mathit{cnt}[0] + \cdots +
\mathit{cnt}[c - 1]$。 這種「前面幾格的和」叫\textbf{前綴和}（prefix sum），不用每
一格重新加一次：$\mathit{start}[c] = \mathit{start}[c - 1] + \mathit{cnt}[c - 1]$，
由左往右一格一格算，全部只要 $n - 1$ 次加法。

\begin{pseudo}[label=alg:arr-fast-transpose]{快速轉置}
\Function{Fast-Transpose}{$A$}
    \State $B.\mathit{rows} \gets A.\mathit{cols}$；$B.\mathit{cols} \gets A.\mathit{rows}$；$B.t \gets A.t$
    \State 建立陣列 $\mathit{cnt}[0..A.\mathit{cols} - 1]$，每一格設成 $0$
    \For{$p \gets 0$ \To $A.t - 1$}
        \State $\mathit{cnt}[A.T[p].\mathit{col}] \gets \mathit{cnt}[A.T[p].\mathit{col}] + 1$
    \EndFor
    \State $\mathit{start}[0] \gets 0$
    \For{$c \gets 1$ \To $A.\mathit{cols} - 1$}
        \State $\mathit{start}[c] \gets \mathit{start}[c - 1] + \mathit{cnt}[c - 1]$
    \EndFor
    \For{$p \gets 0$ \To $A.t - 1$}
        \State $c \gets A.T[p].\mathit{col}$
        \State $B.T[\mathit{start}[c]] \gets (c,\ A.T[p].\mathit{row},\ A.T[p].\mathit{val})$
        \State $\mathit{start}[c] \gets \mathit{start}[c] + 1$
    \EndFor
    \State \Return $B$
\EndFunction
\end{pseudo}

\textbf{手算追蹤}。 用加分表跑一次。 第 4、5 行數完之後，第 7、8 行做前綴和：

\begin{center}
\begin{tabular}{c|ccccccc}
\toprule
行 $c$ & 0 & 1 & 2 & 3 & 4 & 5 & 6 \\
\midrule
$\mathit{cnt}[c]$ & 2 & 1 & 3 & 0 & 0 & 2 & 1 \\
$\mathit{start}[c]$ & 0 & 2 & 3 & 6 & 6 & 6 & 8 \\
\bottomrule
\end{tabular}
\end{center}

第 $3$、$4$ 行沒有非零元素，$\mathit{start}[3] = \mathit{start}[4] =
\mathit{start}[5] = 6$：結果的第 $3$、$4$ 列是空的，第 $5$ 列從 $B[6]$ 開始。 接著第
9 到 12 行依序處理每個三元組，每一步放進 $B[\mathit{start}[c]]$，再把
$\mathit{start}[c]$ 加一：

\begin{center}
\begin{tabular}{cccccl}
\toprule
$p$ & $A.T[p]$ & $c$ & 放進 & 之後的 $\mathit{start}[c]$ & 放進去的三元組 \\
\midrule
0 & $(0,2,3)$ & 2 & $B[3]$ & 4 & $(2,0,3)$ \\
1 & $(0,5,1)$ & 5 & $B[6]$ & 7 & $(5,0,1)$ \\
2 & $(1,0,2)$ & 0 & $B[0]$ & 1 & $(0,1,2)$ \\
3 & $(2,2,5)$ & 2 & $B[4]$ & 5 & $(2,2,5)$ \\
4 & $(2,6,4)$ & 6 & $B[8]$ & 9 & $(6,2,4)$ \\
5 & $(3,1,1)$ & 1 & $B[2]$ & 3 & $(1,3,1)$ \\
6 & $(4,2,2)$ & 2 & $B[5]$ & 6 & $(2,4,2)$ \\
7 & $(4,5,3)$ & 5 & $B[7]$ & 8 & $(5,4,3)$ \\
8 & $(5,0,1)$ & 0 & $B[1]$ & 2 & $(0,5,1)$ \\
\bottomrule
\end{tabular}
\end{center}

最後的 $B$ 是
$(0,1,2)$、$(0,5,1)$、$(1,3,1)$、$(2,0,3)$、$(2,2,5)$、$(2,4,2)$、$(5,0,1)$、$(5,4,3)$、$(6,2,4)$，
依（列, 行）排好。 做完之後每個 $\mathit{start}[c]$ 都變成原本的
$\mathit{start}[c + 1]$（例如 $\mathit{start}[2]$ 從 $3$ 變成 $6$），因為第 $c$ 列
剛好放滿 $\mathit{cnt}[c]$ 個。

\begin{theorem}
\label{thm:arr-fast-transpose}
$A$ 的三元組依（列, 行）排好時，\proc{Fast-Transpose} 回傳的 $B$ 是
$A^{\mathsf T}$，三元組也依（列, 行）排好。 時間是 $\Theta(n + t)$，其中 $n =
A.\mathit{cols}$，額外空間 $\Theta(n)$。
\end{theorem}

\begin{proof}
第 4、5 行做完時，$\mathit{cnt}[c]$ 是 $A$ 的第 $c$ 行非零元素的個數。 第 7、8 行
做完時，$\mathit{start}[c] = \sum_{c' < c} \mathit{cnt}[c']$，也就是行號小於 $c$ 的
三元組個數。 依（列, 行）排好的 $A^{\mathsf T}$ 裡，列號是 $c$ 的三元組（$A$ 的第
$c$ 行）應該占 $B[\mathit{start}[c] .. \mathit{start}[c] + \mathit{cnt}[c] - 1]$ 這
一段，各段互不重疊、剛好鋪滿 $B[0..t-1]$。

第 9--12 行的迴圈每遇到一個行號是 $c$ 的三元組，就放進這一段裡下一個空位，所以
第 $c$ 段依序收到 $A$ 的第 $c$ 行的 $\mathit{cnt}[c]$ 個元素，不會超出這一段。 它們
被處理的順序是在 $A$ 裡出現的順序；$A$ 依列排好，所以這些元素的列號遞增，放進 $B$
之後就是同一列裡行號遞增。 因此 $B$ 依（列, 行）排好，而且每個 $(i, j, v)$ 都變成
了 $(j, i, v)$，$B = A^{\mathsf T}$。

時間：第 3 行和第 7、8 行各 $\Theta(n)$，第 4、5 行和第 9--12 行各 $\Theta(t)$，共
$\Theta(n + t)$。 額外的陣列是 $\mathit{cnt}$ 和 $\mathit{start}$，各 $n$ 格。
\end{proof}

證明用到了「$A$ 依列排好」。 如果 $A$ 的三元組順序是亂的，每一列的範圍仍然對，但同
一列裡的順序就跟著 $A$ 亂掉（觀念題第 12 題）。 這個「先數、再算起點、再一個一個放」
的方法，和第~15~章的計數排序是同一件事：快速轉置就是用行號當鍵，對三元組做一次穩
定的計數排序。

下表比較三種轉置。 $t$ 遠小於 $mn$ 時快速轉置最好；$t$ 接近 $mn$ 時，快速轉置和直
接轉置都是 $\Theta(mn)$，簡單轉置變成 $\Theta(mn^2)$。

\begin{center}
\begin{tabular}{lll}
\toprule
方法 & 時間 & 額外空間 \\
\midrule
存滿的矩陣直接轉置 & $\Theta(mn)$ & 結果的 $mn$ 格 \\
簡單轉置（虛擬碼~\ref{alg:arr-simple-transpose}）& $\Theta(n \cdot t)$ & 結果的 $t$ 個三元組 \\
快速轉置（虛擬碼~\ref{alg:arr-fast-transpose}）& $\Theta(n + t)$ & 結果的 $t$ 個三元組，加 $2n$ 格 \\
\bottomrule
\end{tabular}
\end{center}

快速轉置的 C 程式和下一小節的相乘寫在同一支程式裡（範例~\ref{ex:arr-sparse}），它
的輸出前半就是上面兩張追蹤表。

\subsection{相乘}

$A$ 是 $m \times n$、$B$ 是 $n \times r$，乘積 $C = AB$ 是 $m \times r$，
\[
C[i, j] = \sum_{k=0}^{n-1} A[i, k] \cdot B[k, j] .
\]
存滿的矩陣照定義算要 $mnr$ 次乘法（第~17~章的 Strassen 方法可以更少）。 兩個都稀疏
時，大部分的乘積 $A[i, k] \cdot B[k, j]$ 都有一邊是 $0$，不該去算。

$C[i, j]$ 是「$A$ 的第 $i$ 列」和「$B$ 的第 $j$ 行」的內積，只有兩邊在同一個 $k$ 都
非零時才有貢獻。 $A$ 的第 $i$ 列在三元組陣列裡是連續的一段，依行號 $k$ 排好；可惜
$B$ 的第 $j$ 行散在各處。 所以先把 $B$ 快速轉置成 $B^{\mathsf T}$，$B$ 的第 $j$ 行就
變成 $B^{\mathsf T}$ 的第 $j$ 列，也是連續的一段、依 $k$ 排好。 兩段都依 $k$ 排好，
就能像合併兩個排好的序列一樣，用兩個位置同步往前走，找出 $k$ 相同的配對。

要找到每一列在三元組陣列裡的範圍，先算每一列的起點 $\mathit{rs}[i]$，做法和快速轉
置的 $\mathit{start}$ 一樣，只是多留一格 $\mathit{rs}[m] = t$ 當結尾：第 $i$ 列是
$T[\mathit{rs}[i] .. \mathit{rs}[i + 1] - 1]$，空的列 $\mathit{rs}[i] =
\mathit{rs}[i + 1]$。 算起點的 \proc{Row-Starts} 寫在相乘的虛擬碼後面（虛擬
碼~\ref{alg:arr-row-starts}）。

\begin{pseudo}[label=alg:arr-sparse-mul, unbreakable]{稀疏矩陣相乘}
\Function{Sparse-Multiply}{$A, B$} \Comment{$A$ 是 $m \times n$，$B$ 是 $n \times r$}
    \State $B' \gets \Call{Fast-Transpose}{B}$
    \State $\mathit{ra} \gets \Call{Row-Starts}{A}$；$\mathit{rb} \gets \Call{Row-Starts}{B'}$
    \State $C.\mathit{rows} \gets A.\mathit{rows}$；$C.\mathit{cols} \gets B.\mathit{cols}$；$C.t \gets 0$
    \For{$i \gets 0$ \To $A.\mathit{rows} - 1$}
        \For{$j \gets 0$ \To $B.\mathit{cols} - 1$}
            \State $p \gets \mathit{ra}[i]$；$q \gets \mathit{rb}[j]$；$s \gets 0$
            \While{$p < \mathit{ra}[i + 1]$ \AND $q < \mathit{rb}[j + 1]$}
                \If{$A.T[p].\mathit{col} < B'.T[q].\mathit{col}$}
                    \State $p \gets p + 1$
                \ElsIf{$A.T[p].\mathit{col} > B'.T[q].\mathit{col}$}
                    \State $q \gets q + 1$
                \Else
                    \State $s \gets s + A.T[p].\mathit{val} \cdot B'.T[q].\mathit{val}$
                    \State $p \gets p + 1$；$q \gets q + 1$
                \EndIf
            \EndWhile
            \If{$s \ne 0$}
                \State $C.T[C.t] \gets (i, j, s)$；$C.t \gets C.t + 1$
            \EndIf
        \EndFor
    \EndFor
    \State \Return $C$
\EndFunction
\end{pseudo}

$B'$ 的第 $j$ 列的三元組是 $(j, k, B[k, j])$，所以第 9、11 行比的 $B'.T[q].\mathit{col}$
是 $B$ 原本的列號 $k$。 $i$、$j$ 都由小到大，$C$ 的三元組自然依（列, 行）排好。 第
16 行只存非零的和：兩個非零的乘積可能剛好互相抵消，$C[i, j]$ 是 $0$ 就不存。

\begin{pseudo}[label=alg:arr-row-starts, unbreakable]{每一列的起點}
\Function{Row-Starts}{$M$}
    \State 建立陣列 $\mathit{rs}[0..M.\mathit{rows}]$，每一格設成 $0$
    \For{$p \gets 0$ \To $M.t - 1$}
        \State $\mathit{rs}[M.T[p].\mathit{row} + 1] \gets \mathit{rs}[M.T[p].\mathit{row} + 1] + 1$
    \EndFor
    \For{$i \gets 1$ \To $M.\mathit{rows}$}
        \State $\mathit{rs}[i] \gets \mathit{rs}[i] + \mathit{rs}[i - 1]$
    \EndFor
    \State \Return $\mathit{rs}$
\EndFunction
\end{pseudo}

第 4 行把第 $i$ 列的個數記在 $\mathit{rs}[i + 1]$，往後挪了一格，第 5、6 行的前綴
和做完，$\mathit{rs}[i]$ 就是第 $0$ 到 $i - 1$ 列的總個數。

\textbf{手算追蹤}。 用一組小的例子：
\[
A = \begin{pmatrix} 2 & 0 & 0 & 1 \\ 0 & 0 & 3 & 0 \\ 0 & -1 & 0 & 0 \end{pmatrix},
\qquad
B = \begin{pmatrix} 0 & 4 & 1 \\ 1 & 0 & 0 \\ 0 & 0 & -2 \\ 3 & 0 & -2 \end{pmatrix}.
\]
$A$ 的三元組是 $(0,0,2)$、$(0,3,1)$、$(1,2,3)$、$(2,1,-1)$。 $B$ 轉置之後，$B'$ 的三元組
是 $(0,1,1)$、$(0,3,3)$、$(1,0,4)$、$(2,0,1)$、$(2,2,-2)$、$(2,3,-2)$。 下表把每一列寫成
「$k$:值」，列出九對 $(i, j)$ 各自配對到的 $k$：

\begin{center}
\begin{tabular}{cllll}
\toprule
$(i, j)$ & $A$ 的第 $i$ 列 & $B'$ 的第 $j$ 列 & 配對 & $s$ \\
\midrule
$(0, 0)$ & 0:2, 3:1 & 1:1, 3:3 & $k = 3$：$1 \cdot 3$ & $3$ \\
$(0, 1)$ & 0:2, 3:1 & 0:4 & $k = 0$：$2 \cdot 4$ & $8$ \\
$(0, 2)$ & 0:2, 3:1 & 0:1, 2:$-2$, 3:$-2$ & $k = 0$：$2 \cdot 1$；$k = 3$：$1 \cdot (-2)$ & $0$，不存 \\
$(1, 0)$ & 2:3 & 1:1, 3:3 & 無 & $0$ \\
$(1, 1)$ & 2:3 & 0:4 & 無 & $0$ \\
$(1, 2)$ & 2:3 & 0:1, 2:$-2$, 3:$-2$ & $k = 2$：$3 \cdot (-2)$ & $-6$ \\
$(2, 0)$ & 1:$-1$ & 1:1, 3:3 & $k = 1$：$(-1) \cdot 1$ & $-1$ \\
$(2, 1)$ & 1:$-1$ & 0:4 & 無 & $0$ \\
$(2, 2)$ & 1:$-1$ & 0:1, 2:$-2$, 3:$-2$ & 無 & $0$ \\
\bottomrule
\end{tabular}
\end{center}

結果 $C$ 有 $4$ 項：$(0,0,3)$、$(0,1,8)$、$(1,2,-6)$、$(2,0,-1)$。 總共做了 $6$ 次乘法，存
滿的矩陣要 $3 \times 4 \times 3 = 36$ 次。 $(0, 2)$ 那一格做了兩次乘法，加起來是
$0$，沒有存進 $C$。

\textbf{成本}。 設 $A$ 有 $t_A$ 個非零元素，第 $i$ 列有 $a_i$ 個；$B$ 有 $t_B$ 個，
第 $j$ 行有 $b_j$ 個。 處理一對 $(i, j)$ 時，\textbf{while} 迴圈每一輪至少讓 $p$ 或
$q$ 前進一格，最多 $a_i + b_j$ 輪。 全部加起來
\[
\sum_{i=0}^{m-1} \sum_{j=0}^{r-1} (a_i + b_j) = r \cdot t_A + m \cdot t_B .
\]
再加上兩層 \textbf{for} 本身的 $mr$ 次、轉置的 $\Theta(r + t_B)$ 和算起點的
$\Theta(m + r + t_A + t_B)$，總時間 $O(mr + r \cdot t_A + m \cdot t_B)$。 存滿的矩陣
是 $\Theta(mnr)$；每列、每行只有少數幾個非零元素時，差距非常大。 C 程式在第 $i$ 列
全是 $0$ 時直接跳過內層迴圈，對很多空列的矩陣能再省一些。

另一種常用的做法不轉置 $B$，改成一列一列算 $C$：$C$ 的第 $i$ 列等於「$A$ 的第 $i$
列的每個非零元素 $A[i, k]$ 乘上 $B$ 的第 $k$ 列」加起來。 用一個長度 $r$ 的陣列累
加，這一列做完後把非零的位置收集起來。 它不會為了配對而掃過沒有貢獻的元素，工作量
和非零乘積的個數成正比，另外再加上和 $m$、$r$ 有關的固定成本。 代價是累加陣列要清
空，收集出來的位置也要排序。 本章的實作練習有一題就是這個做法。

\begin{cbookrecap}[第~9~章 9.3 節、第~14~章 14.3、14.5、14.6 節、第~19~章 19.3 節]
函式收 \texttt{const Sparse *A}，表示只讀不改；要把結果交回呼叫端，就傳
\texttt{Sparse *C} 進去，在函式裡設定 \texttt{C->t}、\texttt{C->a}。
\texttt{(Triple)\{c, r, v\}} 是 C99 起的\textbf{複合字面值}，臨時造出一個
\texttt{Triple} 結構的值，效果和分三行設定 \texttt{row}、\texttt{col}、\texttt{val}
一樣。 \texttt{calloc(n, size)} 配置 $n$ 格並全部設成 $0$，適合當計數陣列。
\texttt{free(NULL)} 什麼都不做，所以配置失敗時可以把幾個指標一起 \texttt{free}。
\end{cbookrecap}

下面的程式先對加分表做快速轉置，印出 $\mathit{cnt}$、$\mathit{start}$ 和每一步放進
$B$ 的哪一格；再算上面的 $C = AB$。 C 的陣列從 $0$ 開始，和虛擬碼、追蹤表完全一樣，
輸出可以逐行對照。

\begin{example}[label=ex:arr-sparse]{c}{(sparse.c) 稀疏矩陣的快速轉置與相乘}
#include <stdio.h>
#include <stdlib.h>

typedef struct {
    int row, col, val;
} Triple;

typedef struct {
    int rows, cols, t;          // 列數、行數、非零項個數
    Triple *a;                  // a[0..t-1]，依（列, 行）由小到大
} Sparse;

// 先印大小，再每行印 5 個三元組。
void print_terms(const char *title, const Sparse *m)
{
    printf("%s: %d x %d, %d terms\n", title, m->rows, m->cols, m->t);
    for (int p = 0; p < m->t; p++)
        printf("(%d,%d,%d)%s", m->a[p].row, m->a[p].col, m->a[p].val,
               (p % 5 == 4 || p == m->t - 1) ? "\n" : " ");
}

// 快速轉置。B->a 由這個函式配置；trace 不是 0 時印出每一步。
// 假設 A->cols >= 1。成功回傳 1。
int fast_transpose(const Sparse *A, Sparse *B, int trace)
{
    int n = A->cols;
    int *cnt = calloc(n, sizeof *cnt);
    int *start = malloc(n * sizeof *start);
    B->a = malloc((A->t > 0 ? A->t : 1) * sizeof *B->a);
    if (cnt == NULL || start == NULL || B->a == NULL) {
        free(cnt); free(start); free(B->a);
        return 0;
    }
    B->rows = A->cols;
    B->cols = A->rows;
    B->t = A->t;
    for (int p = 0; p < A->t; p++)
        cnt[A->a[p].col]++;
    start[0] = 0;
    for (int c = 1; c < n; c++)
        start[c] = start[c - 1] + cnt[c - 1];
    if (trace) {
        printf("cnt:");
        for (int c = 0; c < n; c++) printf(" %d", cnt[c]);
        printf("\nstart:");
        for (int c = 0; c < n; c++) printf(" %d", start[c]);
        printf("\n");
    }
    for (int p = 0; p < A->t; p++) {
        int c = A->a[p].col;
        int q = start[c]++;
        B->a[q] = (Triple){c, A->a[p].row, A->a[p].val};
        if (trace)
            printf("p=%d -> B[%d], start[%d]=%d\n",
                   p, q, c, start[c]);
    }
    free(cnt);
    free(start);
    return 1;
}

// rs[i] 是第 i 列第一項的位置，rs[rows] = t。由呼叫者 free。
int *row_starts(const Sparse *M)
{
    int *rs = calloc(M->rows + 1, sizeof *rs);
    if (rs == NULL)
        return NULL;
    for (int p = 0; p < M->t; p++)
        rs[M->a[p].row + 1]++;
    for (int i = 1; i <= M->rows; i++)
        rs[i] += rs[i - 1];
    return rs;
}

// C = A * B。C->a 由這個函式配置。*mul 記錄乘法次數。成功回傳 1。
int sparse_mul(const Sparse *A, const Sparse *B, Sparse *C, long *mul)
{
    Sparse Bt;
    if (A->cols != B->rows || !fast_transpose(B, &Bt, 0))
        return 0;
    int *ra = row_starts(A), *rb = row_starts(&Bt);
    long cap = (long)A->rows * B->cols;     // C 最多這麼多項
    C->a = malloc((cap > 0 ? cap : 1) * sizeof *C->a);
    if (ra == NULL || rb == NULL || C->a == NULL) {
        free(ra); free(rb); free(Bt.a); free(C->a);
        return 0;
    }
    C->rows = A->rows;
    C->cols = B->cols;
    C->t = 0;
    *mul = 0;
    for (int i = 0; i < A->rows; i++) {
        if (ra[i] == ra[i + 1])
            continue;                       // A 的第 i 列全是 0
        for (int j = 0; j < Bt.rows; j++) {
            int p = ra[i], q = rb[j];
            long sum = 0;
            while (p < ra[i + 1] && q < rb[j + 1]) {
                if (A->a[p].col < Bt.a[q].col) {
                    p++;
                } else if (A->a[p].col > Bt.a[q].col) {
                    q++;
                } else {
                    sum += (long)A->a[p].val * Bt.a[q].val;
                    (*mul)++;
                    p++;
                    q++;
                }
            }
            if (sum != 0)                   // 相消成 0 的不存
                C->a[C->t++] = (Triple){i, j, (int)sum};
        }
    }
    free(ra);
    free(rb);
    free(Bt.a);
    return 1;
}

int main(void)
{
    // 6 位同學、7 次實驗課的加分表，只記非零的 9 格
    Triple g[] = {{0, 2, 3}, {0, 5, 1}, {1, 0, 2}, {2, 2, 5},
                  {2, 6, 4}, {3, 1, 1}, {4, 2, 2}, {4, 5, 3},
                  {5, 0, 1}};
    Sparse G = {6, 7, 9, g}, Gt;
    if (!fast_transpose(&G, &Gt, 1))
        return 1;
    print_terms("Gt", &Gt);
    free(Gt.a);

    Triple ta[] = {{0, 0, 2}, {0, 3, 1}, {1, 2, 3}, {2, 1, -1}};
    Triple tb[] = {{0, 1, 4}, {0, 2, 1}, {1, 0, 1},
                   {2, 2, -2}, {3, 0, 3}, {3, 2, -2}};
    Sparse A = {3, 4, 4, ta}, B = {4, 3, 6, tb}, C;
    long mul;
    if (!sparse_mul(&A, &B, &C, &mul))
        return 1;
    print_terms("C", &C);
    printf("multiplications = %ld (dense: %d)\n",
           mul, A.rows * A.cols * B.cols);
    free(C.a);
    return 0;
}
\end{example}
\iosample{}{cnt: 2 1 3 0 0 2 1 \\ start: 0 2 3 6 6 6 8 \\ p=0 -> B[3], start[2]=4 \\ p=1 -> B[6], start[5]=7 \\ p=2 -> B[0], start[0]=1 \\ p=3 -> B[4], start[2]=5 \\ p=4 -> B[8], start[6]=9 \\ p=5 -> B[2], start[1]=3 \\ p=6 -> B[5], start[2]=6 \\ p=7 -> B[7], start[5]=8 \\ p=8 -> B[1], start[0]=2 \\ Gt: 7 x 6, 9 terms \\ (0,1,2) (0,5,1) (1,3,1) (2,0,3) (2,2,5) \\ (2,4,2) (5,0,1) (5,4,3) (6,2,4) \\ C: 3 x 3, 4 terms \\ (0,0,3) (0,1,8) (1,2,-6) (2,0,-1) \\ multiplications = 6 (dense: 36)}

\texttt{int q = start[c]++;} 先取出 $\mathit{start}[c]$ 的值放進 \texttt{q}，再把
\texttt{start[c]} 加一，對應虛擬碼~\ref{alg:arr-fast-transpose} 的第 11、12 行。 快
速轉置的三個迴圈合計只跑了 $9 + 6 + 9 = 24$ 輪，簡單轉置要比較 $63$ 次。
\texttt{sparse\_mul} 結束前把轉置用的 \texttt{Bt.a} 和兩個起點陣列都釋放掉，只留下
\texttt{C->a} 給 \texttt{main}。 \texttt{C->a} 一開始配置 $mr$ 格，這是 $C$ 最多可能
的項數；在意空間時，可以先算一遍有幾項再配置，或是像第~3~章的動態陣列那樣不夠時
加倍。

\section{多項式}
\label{sec:arr-poly}

單變數多項式 $p(x) = c_d x^d + c_{d-1} x^{d-1} + \cdots + c_1 x + c_0$ 由係數決定。
$c_d \ne 0$ 時 $d$ 是它的\textbf{次數}（degree）。 這一節比較兩種用陣列存多項式的
方法，再看相加、相乘和求值各要多少時間。

\subsection{兩種表示法}

\textbf{係數陣列}。 開一個長度 $d + 1$ 的陣列 $c[0..d]$，$c[i]$ 是 $x^i$ 的係數，
係數是 $0$ 的也存。 下標就是次方，要讀寫某一次的係數是 $O(1)$。 例如
$p(x) = 2x^4 - 3x^3 + 5x - 6$ 存成 $c[0..4] = (-6, 5, 0, -3, 2)$。

\textbf{非零項陣列}。 只存係數不是 $0$ 的項，每一項是一個（係數, 次方）的配對，依次
方由大到小排。 $p(x)$ 存成 $(2, 4), (-3, 3), (5, 1), (-6, 0)$。 有 $t$ 個非零項時用
$2t$ 格。

兩者的空間是 $d + 1$ 對 $2t$。 $x^{1000} + 1$ 用係數陣列要 $1001$ 格，非零項陣列只
要 $4$ 格；反過來，$3x^5 - x^4 + 2x^3 + 7x^2 - x + 4$ 每一項都不是 $0$，係數陣列用
$6$ 格，非零項陣列要 $12$ 格。 和稀疏矩陣一樣，非零項陣列適合「次數高、項數少」的
多項式。

很多個多項式要一起存時，非零項陣列還可以全部放進同一個大陣列，每個多項式只記住自
己從第幾格到第幾格。 這樣不用為每個多項式各配置一塊，但某個多項式變長時，後面的多
項式就得整個往後搬；刪掉一個多項式則會在中間留下空洞。 第~5~章用串列存多項式，就
是為了避開這兩個問題。

\begin{examnote}
表示法的約定先看清楚再動手。 係數陣列的第一個數是常數項還是最高次項，題目各有不
同（本章和大部分 Aki-Judge 的題目是由低次到高次，但也有由高次寫到常數項的）。 稀疏
矩陣的三元組陣列，有些課本（例如 Horowitz 等人的 \textit{Fundamentals of Data
Structures in C}）把第 $0$ 格拿來存（列數, 行數, 非零個數），真正的三元組從第 $1$
格開始；這時快速轉置的起點陣列也從 $1$ 開始算，追蹤表上的每個位置都比本
章多 $1$。
\end{examnote}

\subsection{相加}

係數陣列相加就是對應的格子相加：$r[i] = a[i] + b[i]$，短的那個多項式缺的高次項當成
$0$。 時間 $\Theta(\max(d_a, d_b) + 1)$。 和的最高次係數可能抵消成 $0$（例如
$x^2 + 1$ 加 $-x^2 + x$），要知道真正的次數，就從最高次往下找第一個不是 $0$ 的。

非零項陣列相加和第~\ref{sec:ll-poly}~節的 \proc{Poly-Add}（虛擬碼~\ref{alg:ll-poly-add}）是同一個演算法：兩個位置 $i$、$j$
各指著一個多項式目前的項，次方大的先抄到結果；次方相同就把係數加起來，和不是 $0$
才抄，兩邊一起前進；有一邊抄完了，另一邊剩下的照抄。 那裡的節點換成陣列的格子，
「接到結果後面」換成「寫進結果的下一格」，其他完全一樣。 設兩個多項式各有 $t_a$、
$t_b$ 項，結果最多 $t_a + t_b$ 項，所以陣列版可以事先配置這麼多格；比較次方最多
$t_a + t_b - 1$ 次，時間 $\Theta(t_a + t_b)$。

\subsection{相乘}

係數陣列相乘：$a(x)$ 的 $x^i$ 項乘上 $b(x)$ 的 $x^j$ 項會貢獻到 $x^{i+j}$，所以
\[
r[k] = \sum_{i + j = k} a[i] \cdot b[j], \qquad 0 \le k \le d_a + d_b .
\]

\begin{pseudo}[label=alg:arr-poly-mul]{係數陣列相乘}
\Function{Poly-Mul}{$a, d_a, b, d_b$}
    \State 建立陣列 $r[0..d_a + d_b]$，每一格設成 $0$
    \For{$i \gets 0$ \To $d_a$}
        \For{$j \gets 0$ \To $d_b$}
            \State $r[i + j] \gets r[i + j] + a[i] \cdot b[j]$
        \EndFor
    \EndFor
    \State \Return $r$
\EndFunction
\end{pseudo}

第 5 行執行 $(d_a + 1)(d_b + 1)$ 次，時間 $\Theta(d_a d_b)$（$d_a, d_b \ge 1$ 時）。
兩個次數都是 $d$ 時是 $\Theta(d^2)$。 第~17~章的 Karatsuba 方法用分治把它降到
$\Theta(d^{\lg 3})$，那裡的大數乘法就是把整數的每一位當成多項式的係數。

非零項陣列相乘：$a$ 的每一項乘上 $b$ 的每一項，共 $t_a t_b$ 個乘積，難處在把次方相
同的乘積合併起來。 三種做法：

\begin{itemize}
\item 把「$a$ 的一項乘上整個 $b$」得到的 $t_b$ 項（次方仍由大到小）用上一小節的相加
加進結果，重複 $t_a$ 次。 第 $i$ 次相加時結果最多 $(i - 1)t_b$ 項，這次相加是
$O(i \cdot t_b)$，合計 $O(t_a^2 t_b)$。 第~5~章的串列版也是這樣分析。
\item 把 $t_a t_b$ 個乘積全部算出來，依次方排序（第~14~章），再把相鄰次方相同的合
併。 時間 $O(t_a t_b \lg(t_a t_b))$。
\item 結果的次數 $d_a + d_b$ 不大時，直接用一個係數陣列累加，再把非零的挑出來。 時間
$O(t_a t_b + d_a + d_b)$。
\end{itemize}

兩個多項式都很稀疏、次數又很高時（例如次數 $10^6$、各 $50$ 項），係數陣列的乘法要
做約 $10^{12}$ 次，非零項陣列只要 $2500$ 個乘積，這時一定用非零項陣列。

\subsection{Horner 法求值}

要算 $p(x_0)$，直接照定義一項一項算：第 $i$ 項的 $x_0^i$ 如果每次重新乘，要 $i - 1$
次乘法，全部加起來是 $\Theta(d^2)$ 次。 一邊加一邊維護 $x_0$ 的次方（每一項之後乘一
次 $x_0$），可以降到大約 $2d$ 次。 \textbf{Horner 法}（Horner's rule）更省：把多項式
寫成一層包一層的形式
\[
p(x) = \Bigl(\cdots\bigl((c_d\,x + c_{d-1})\,x + c_{d-2}\bigr)\,x + \cdots + c_1\Bigr)\,x + c_0 ,
\]
從最裡面往外算，每一層一次乘法、一次加法。

\begin{pseudo}[label=alg:arr-horner]{Horner 法，$c[0..d]$ 是係數}
\Function{Horner}{$c, d, x$}
    \State $s \gets c[d]$
    \For{$i \gets d - 1$ \DownTo $0$}
        \State $s \gets s \cdot x + c[i]$
    \EndFor
    \State \Return $s$
\EndFunction
\end{pseudo}

\textbf{手算追蹤}。 $p(x) = 2x^4 - 3x^3 + 5x - 6$ 在 $x = 3$：

\begin{center}
\begin{tabular}{cccl}
\toprule
$i$ & $c[i]$ & $s \cdot 3 + c[i]$ & 這一步之後的 $s$ \\
\midrule
（第 2 行）& $c[4] = 2$ & & $2$ \\
3 & $-3$ & $2 \cdot 3 - 3$ & $3$ \\
2 & $0$ & $3 \cdot 3 + 0$ & $9$ \\
1 & $5$ & $9 \cdot 3 + 5$ & $32$ \\
0 & $-6$ & $32 \cdot 3 - 6$ & $90$ \\
\bottomrule
\end{tabular}
\end{center}

$p(3) = 2 \cdot 81 - 3 \cdot 27 + 15 - 6 = 90$，相符。 一共 $4$ 次乘法；$0$ 係數那一
項也要做一次，Horner 法對每個次方都做一步。

\begin{theorem}
\label{thm:arr-horner}
\proc{Horner}$(c, d, x)$ 回傳 $\sum_{k=0}^{d} c[k]\,x^k$，用 $d$ 次乘法和 $d$ 次加法。
\end{theorem}

\begin{proof}
對 $0 \le i \le d$，定義只看 $c[i..d]$ 的多項式
\[
q_i(x) = c[i] + c[i+1]\,x + \cdots + c[d]\,x^{d - i} = \sum_{k=i}^{d} c[k]\,x^{k - i} .
\]
$q_d(x) = c[d]$，$q_0$ 就是 $p$，而且把 $q_{i+1}$ 的每一項乘上 $x$、再補上常數項
$c[i]$，就得到 $q_i$：$q_i(x) = q_{i+1}(x) \cdot x + c[i]$。

不變量：第 3 行的迴圈要處理 $i$ 之前（$i = d - 1, d - 2, \ldots, 0$），$s =
q_{i+1}(x)$。 初始：$i = d - 1$ 時 $s = c[d] = q_d(x)$。 維持：第 4 行把 $s$ 換成
$q_{i+1}(x) \cdot x + c[i] = q_i(x)$，下一個要處理的是 $i - 1$，不變量成立。 結束：迴
圈做完 $i = 0$ 之後 $s = q_0(x) = p(x)$。 迴圈跑 $d$ 輪，每輪一次乘法、一次加法。
\end{proof}

\textbf{Horner 法和除法}。 上面的追蹤表中間出現的 $2, 3, 9, 32$ 也有用：它們正好是
$p(x)$ 除以 $(x - 3)$ 的商的係數，最後的 $90$ 是餘數，
\[
2x^4 - 3x^3 + 5x - 6 = (x - 3)(2x^3 + 3x^2 + 9x + 32) + 90 .
\]
把右邊展開就能驗證。 這就是高中學過的綜合除法，也說明了為什麼 $p(x_0)$ 等於 $p(x)$
除以 $(x - x_0)$ 的餘數。 一般的證明留作觀念題第 15 題。

非零項陣列也能用 Horner 法：相鄰兩項的次方差 $g$ 不是 $1$ 時，那一步改成乘上
$x^{g}$，用第~17~章的快速冪算 $x^g$ 只要 $O(\lg g)$ 次乘法。 $t$ 項、次數 $d$ 的多
項式總共 $O(t \lg d)$ 次乘法。

下面的程式用係數陣列做相加、相乘和 Horner 法求值，順便印出綜合除法的商；係數一律
由低次印到高次。 非零項陣列的相加和第~5~章的串列版只差在「寫進下一格」，留給實作
練習。

\begin{example}[label=ex:arr-poly-dense]{c}{(poly\_dense.c) 係數陣列的多項式}
#include <stdio.h>

// 係數陣列：c[i] 是 x^i 的係數，c[0..d]。由低次到高次印出來。
void print_poly(const char *name, const long long *c, int d)
{
    printf("%s:", name);
    for (int i = 0; i <= d; i++)
        printf(" %lld", c[i]);
    printf("\n");
}

// r[0..max(da,db)] = a + b
void poly_add(const long long *a, int da, const long long *b, int db,
              long long *r)
{
    int d = da > db ? da : db;
    for (int i = 0; i <= d; i++)
        r[i] = (i <= da ? a[i] : 0) + (i <= db ? b[i] : 0);
}

// r[0..da+db] = a * b，每一對 (i, j) 都乘一次。
void poly_mul(const long long *a, int da, const long long *b, int db,
              long long *r)
{
    for (int k = 0; k <= da + db; k++)
        r[k] = 0;
    for (int i = 0; i <= da; i++)
        for (int j = 0; j <= db; j++)
            r[i + j] += a[i] * b[j];
}

// Horner 法求 p(x)。q 不是 NULL 時，順便存下每一步的值：
// q[d-1..0] 是 p(x) 除以 (x - x0) 的商。*mul 記錄乘法次數。
long long horner(const long long *c, int d, long long x, long long *q,
                 int *mul)
{
    long long s = c[d];
    *mul = 0;
    for (int i = d - 1; i >= 0; i--) {
        if (q != NULL)
            q[i] = s;
        s = s * x + c[i];
        (*mul)++;
    }
    return s;
}

int main(void)
{
    long long p[] = {-6, 5, 0, -3, 2};      // 2x^4 - 3x^3 + 5x - 6
    long long g[] = {-1, 4, 1};             // x^2 + 4x - 1
    long long sum[5], prod[7], quot[4];
    int mul;
    print_poly("p", p, 4);
    print_poly("g", g, 2);
    poly_add(p, 4, g, 2, sum);
    print_poly("p + g", sum, 4);
    poly_mul(p, 4, g, 2, prod);
    print_poly("p * g", prod, 6);
    long long v = horner(p, 4, 3, quot, &mul);
    printf("p(3) = %lld, multiplications = %d\n", v, mul);
    print_poly("p / (x - 3)", quot, 3);
    printf("remainder = %lld\n", v);
    return 0;
}
\end{example}
\iosample{}{p: -6 5 0 -3 2 \\ g: -1 4 1 \\ p + g: -7 9 1 -3 2 \\ p * g: 6 -29 14 8 -14 5 2 \\ p(3) = 90, multiplications = 4 \\ p / (x - 3): 32 9 3 2 \\ remainder = 90}

\texttt{p * g} 的 $7$ 個係數由 \texttt{poly\_mul} 的雙層迴圈做了 $5 \times 3 = 15$
次乘法得到。 \texttt{horner} 在第 $i$ 步把當下的 $s$ 存進 \texttt{q[i]}，所以
\texttt{q[0..3]} 是商 $2x^3 + 3x^2 + 9x + 32$ 由低次到高次的係數。

\section{本章常見錯誤}
\label{sec:arr-pitfalls}

\begin{center}
\begin{tabular}{p{0.24\linewidth}p{0.50\linewidth}p{0.16\linewidth}}
\toprule
錯誤 & 例子 & 見 \\
\midrule
忘了減下界 & $A[-2..5, 3..9]$ 的 $A[1, 7]$ 直接用 $1 \times 7 + 7$ 算位移 &
第~\ref{sec:arr-address}~節 \\
乘錯維度 & 列優先乘了列數 $m$，或行優先乘了行數 $n$ & 第~\ref{sec:arr-address}~節 \\
位址單位搞混 & 每格 $4$ 位元組卻只加了位移的元素個數；十六進位的位址用十進位相加 &
第~\ref{sec:arr-address}~節 \\
反推時正負號錯 & 兩個位址相減時，行號比較小的那一項寫成正的 &
第~\ref{sec:arr-address}~節 \\
C 用指標假裝下界 & \texttt{int *p = buf - l;} 再用 \texttt{p[i]}：指到陣列外，
未定義行為 & 第~\ref{sec:arr-address}~節 \\
三角矩陣的起點差 $1$ & 把從 $1$ 起算的 $i(i-1)/2 + j$ 用在下標從 $0$ 開始的 C 程式 &
第~\ref{sec:arr-special}~節 \\
上三角套用下三角公式 & 沒有先把 $i$、$j$ 對調；行優先時忘了公式裡有 $n$ &
第~\ref{sec:arr-special}~節 \\
轉置只對調列和行 & 結果的三元組沒有依（列, 行）排好，之後的二分搜尋、相乘出錯 &
第~\ref{sec:arr-sparse}~節 \\
起點陣列漏了空行 & 只為出現過的行算 $\mathit{start}$，行數是 $0$ 的那幾行沒有填 &
第~\ref{sec:arr-sparse}~節 \\
放完沒有更新起點 & 快速轉置放進 $B[\mathit{start}[c]]$ 後忘了加一，同一行的元素互相
覆蓋 & 第~\ref{sec:arr-sparse}~節 \\
存了 $0$ 的項 & 稀疏相乘或多項式相加的結果裡留下抵消成 $0$ 的項 &
第~\ref{sec:arr-sparse}、\ref{sec:arr-poly}~節 \\
係數順序讀反 & 題目由高次寫到常數項，程式卻當成 $c[0]$ 是常數項 &
第~\ref{sec:arr-poly}~節 \\
Horner 法從常數項開始 & $s \gets c[0]$，再依序乘 $x$ 加 $c[1], c[2], \ldots$，算出來
的是係數倒過來的多項式 & 第~\ref{sec:arr-poly}~節 \\
\bottomrule
\end{tabular}
\end{center}

\section{觀念題}
\label{sec:arr-exercises}

題型和研究所考試、期中期末考的筆試相近，有選擇、計算、追蹤和證明。 除非題目另外
說明，位址都用十進位，每格大小的單位是位元組。 解答在本章最後。

\begin{enumerate}
\item （選擇，可能不只一個答案）$A[1..m, 1..n]$，每個元素 $w$ 位元組。 下列哪些敘
述正確？
\begin{enumerate}
\item 列優先存放時，$A[i, j]$ 和 $A[i, j + 1]$（$j < n$）的位址相差 $w$。
\item 行優先存放時，$A[i, j]$ 和 $A[i + 1, j]$（$i < m$）的位址相差 $n \cdot w$。
\item 用列優先存 $A$，和用行優先存 $A^{\mathsf T}$（$n \times m$），每個元素在記憶體
裡的位置完全一樣。
\item 列優先存放時，只要知道 $\alpha$、$w$、$i$、$j$，不用知道 $n$ 就能算出
$A[i, j]$ 的位址。
\item C 的 \texttt{int a[5][8]} 裡，\texttt{a[2][3]} 在 \texttt{a[0][0]} 之後第 $19$
個 \texttt{int} 的位置。
\end{enumerate}

\item （計算）陣列 $H[-20..44]$，每個元素 $6$ 位元組，$H[-20]$ 的位址是十六進位的
\texttt{2F40}。 求元素個數、$H[0]$ 和 $H[31]$ 的位址，以及整個陣列占用的最後一個位
元組的位址（都用十六進位）。

\item （計算）$P[0..11, -4..15]$，每個元素 $4$ 位元組，$P[0, -4]$ 的位址是 $4000$。
\begin{enumerate}
\item 分別用列優先和行優先，求 $P[7, 3]$ 的位址。
\item 列優先時，位址 $4400$ 是哪一個元素？ 行優先時呢？
\end{enumerate}

\item （計算）$Q[1..m, 1..8]$ 以行優先存放，每個元素 $4$ 位元組，$m$ 未知。 已知
$Q[2, 3]$ 的位址是 $1276$，$Q[5, 7]$ 的位址是 $1432$。 求 $m$、$Q[1, 1]$ 的位址，
以及最後一個元素 $Q[m, 8]$ 的位址。

\item （證明）$A[l_1..u_1, l_2..u_2]$ 以列優先存放，$n = u_2 - l_2 + 1$。 證明位移
$\mathrm{off}(i, j) = (i - l_1)\,n + (j - l_2)$ 把 $mn$ 個下標一對一地對應到
$0, 1, \ldots, mn - 1$，並寫出反過來的公式：給 $\mathrm{off}$，求 $(i, j)$。

\item （計算）下面的程式片段把一個 C 的二維陣列清成 $0$。
\begin{code}{c}{觀念題 6}
double a[400][300];
for (int j = 0; j < 300; j++)
    for (int i = 0; i < 400; i++)
        a[i][j] = 0;
\end{code}
\begin{enumerate}
\item 內層迴圈相鄰兩次寫入的位址相差幾個位元組？
\item 從 \texttt{a[399][j]} 換到 \texttt{a[0][j + 1]} 時，位址差多少？
\item 把兩層迴圈對調（外層 \texttt{i}、內層 \texttt{j}）之後，相鄰兩次寫入的位址相差
多少？ 換列的時候呢？
\end{enumerate}

\item （計算）C 程式用一維陣列 \texttt{b[]} 存 $8 \times 8$ 下三角矩陣的下三角（含對
角線），列優先，下標全部從 $0$ 開始。 \texttt{b} 要開幾格？ $A[6][2]$ 放在
\texttt{b} 的第幾格？ 如果 $A$ 是對稱矩陣，$A[2][6]$ 要讀哪一格？

\item （計算）$10 \times 10$ 的上三角矩陣，下標從 $1$ 開始，把上三角（含對角線）放進
$B[1..55]$，每個元素 $8$ 位元組，$B[1]$ 的位址是 $600$。 分別用列優先和行優先，求
$A[4, 7]$ 在 $B$ 的第幾格和它的位址。

\item （推導）下三角矩陣以列優先放進 $B[1..n(n+1)/2]$（下標從 $1$ 開始）。 給 $k$，
求它是 $A$ 的哪一格 $(i, j)$。 寫出 $i$ 的公式並說明理由，再算 $k = 23$。

\item （計算）$50 \times 50$ 的三對角矩陣，下標從 $1$ 開始。
\begin{enumerate}
\item 非零元素最多幾個？
\item 依列優先緊密放進 $B[1..]$ 時，$A[20, 21]$ 和 $A[50, 49]$ 各是第幾格？
\item 改用每列 $3$ 格的 $E[1..50, 1..3]$（$A[i, j]$ 放在 $E[i, j - i + 2]$），再用列
優先攤平，同樣兩個元素各是第幾格？ $E$ 共用幾格？
\end{enumerate}

\item （追蹤）$5 \times 6$ 稀疏矩陣的三元組依序是
$(0,1,7)$、$(0,4,2)$、$(1,1,-3)$、$(2,0,6)$、$(2,5,1)$、$(3,3,9)$、$(4,0,-4)$、$(4,4,8)$。 用
\proc{Fast-Transpose}（虛擬碼~\ref{alg:arr-fast-transpose}）轉置。 寫出
$\mathit{cnt}$、一開始的 $\mathit{start}$、每一步放進 $B$ 的哪一格，以及最後的 $B$。

\item （說明）\proc{Fast-Transpose} 的輸入如果沒有依（列, 行）排好，例如 $3 \times 3$
矩陣的三元組依序是 $(2,0,6)$、$(0,0,1)$、$(1,2,4)$、$(0,2,5)$，輸出是什麼？ 它是不是
$A^{\mathsf T}$？ 是不是依（列, 行）排好？ \proc{Simple-Transpose} 會不會有同樣的
問題？

\item （追蹤）用 \proc{Sparse-Multiply}（虛擬碼~\ref{alg:arr-sparse-mul}）計算 $C = AB$，
$A$ 是 $2 \times 3$，三元組 $(0,0,1)$、$(0,2,-2)$、$(1,1,3)$；$B$ 是 $3 \times 2$，三元組
$(0,0,2)$、$(0,1,1)$、$(1,1,-1)$、$(2,0,1)$、$(2,1,4)$。 寫出 $B'$、$\mathit{ra}$、
$\mathit{rb}$、$C$ 的三元組，以及總共做了幾次乘法。 哪一對 $(i, j)$ 做了乘法、卻沒
有存進 $C$？

\item （計算）用 Horner 法計算 $p(x) = x^5 - 4x^3 + 2x^2 - 7$ 在 $x = -2$ 的值，列出每
一步的 $s$，並寫出 $p(x)$ 除以 $(x + 2)$ 的商和餘數。

\item （證明）$p(x) = \sum_{k=0}^{d} c_k x^k$，$d \ge 1$。 對一個數 $a$，令 $s_d =
c_d$，$s_i = s_{i+1} \cdot a + c_i$（$i = d - 1, \ldots, 0$），也就是 \proc{Horner}
$(c, d, a)$ 每一步的 $s$。 證明
\[
p(x) = (x - a)\bigl(s_d\,x^{d-1} + s_{d-1}\,x^{d-2} + \cdots + s_1\bigr) + s_0 ,
\]
並由此說明 $p(a) = s_0$。

\item （計算）$f(x) = x^{1000000} + 3x^{500000} - 1$，$g(x)$ 的次數是 $10^6$，有
$40$ 個非零項。
\begin{enumerate}
\item 用係數陣列和非零項陣列存 $f$，各要幾格？
\item 用兩種表示法計算 $f + g$ 和 $f \cdot g$，各大約要做多少步（只要量級）？
\item $h(x)$ 的次數是 $30$，$31$ 個係數都不是 $0$。 存 $h$ 用哪一種表示法比較省？
\end{enumerate}
\end{enumerate}

\section{實作練習}
\label{sec:arr-practice}

下面是 Aki-Judge 上和本章相關的題目。 有幾題只要交函式和標頭檔，由評測程式讀輸入、
呼叫你的函式。 每一題的列號、行號從 $0$ 還是 $1$ 開始，係數由低次還是高次排，題目
都寫得很清楚，動手前先對照本章的約定。

\begin{center}
\begin{tabular}{llc}
\toprule
小節 & 題目 & 難度 \\
\midrule
位址計算 & \akiref{locating-cells-in-a-shared-pool}{共用記憶體區塊的元素定位} & 進階 \\
稀疏矩陣 & \akiref{transpose-a-sparse-matrix}{稀疏矩陣轉置} & 暖身 \\
稀疏矩陣 & \akiref{sparse-matrix-set-and-get}{稀疏矩陣設值與查詢} & 基本 \\
稀疏矩陣 & \akiref{multiply-two-sparse-matrices}{稀疏矩陣相乘} & 基本 \\
稀疏矩陣 & \akiref{dot-products-of-sparse-rows}{稀疏列的內積} & 基本 \\
稀疏矩陣 & \akiref{row-by-row-sparse-matrix-product}{逐列累加的稀疏矩陣乘積} & 進階 \\
多項式 & \akiref{evaluate-a-polynomial-at-a-point}{多項式在某點的值} & 暖身 \\
多項式 & \akiref{add-two-polynomials}{兩個多項式相加} & 暖身 \\
多項式 & \akiref{dense-polynomial-set-eval-and-degree}{稠密多項式設值、求值與次數} & 基本 \\
多項式 & \akiref{sparse-polynomial-terms}{稀疏多項式的項} & 基本 \\
多項式 & \akiref{polynomial-long-division}{多項式長除法} & 基本 \\
多項式 & \akiref{polynomial-composition}{多項式代入多項式} & 進階 \\
\bottomrule
\end{tabular}
\end{center}

幾點提示：

\begin{itemize}
\item 「共用記憶體區塊的元素定位」是把位址公式反過來用：給一個指標，先找出它落在哪
個矩陣，再用觀念題第 5 題的反函數求列和行。
\item 「稀疏矩陣轉置」規定的輸出順序不是排好的順序，照題目做就好。 寫完之後，可以
自己改成輸出依（列, 行）排好的結果，用快速轉置再寫一次。
\item 「稀疏列的內積」用一個指標陣列記每一列的開頭，作用和 \proc{Row-Starts} 的
$\mathit{rs}$ 一樣；兩列的內積就是 \proc{Sparse-Multiply} 第 8--15 行的合併。
\item 「逐列累加的稀疏矩陣乘積」用的是第~\ref{sec:arr-sparse}~節最後提到的另一種乘
法。 資料量大，每一列都把長度 $r$ 的累加陣列整個清空會太慢，要想辦法只清用過的位置。
\item 「多項式長除法」的係數由最高次寫到常數項，和本章相反。 除式的最高次係數是
$1$，每一步的做法和綜合除法類似。
\item 「多項式代入多項式」把 Horner 法裡的「乘上 $x$」換成「乘上 $g(x)$」，每一步
是一次係數陣列的乘法。
\end{itemize}

\section{讀完你應該會}
\label{sec:arr-checklist}

\begin{itemize}
\item 給上下界、每格大小和起始位址，能用列優先和行優先算出一維、二維、三維陣列任
一格的位址；給兩格的位址，能反推未知的維度大小或起始位址。
\item 能推導上、下三角矩陣在列優先、行優先、下標從 $0$ 或 $1$ 開始時的對應公式，
並用對稱性從一種推出另一種；能寫出三對角矩陣和帶狀矩陣的存法。
\item 能把稀疏矩陣寫成三元組，手算快速轉置的 $\mathit{cnt}$、$\mathit{start}$ 和每一
步的位置，並說明簡單轉置、快速轉置的時間。
\item 能用「先轉置、再逐對合併」計算兩個稀疏矩陣的乘積，並算出它的時間上界。
\item 能依次數和項數選擇多項式的表示法，說出相加、相乘的成本；能用 Horner 法求值，
並從過程讀出除以 $(x - a)$ 的商和餘數。
\end{itemize}

\textbf{延伸閱讀}（書末會整理完整的參考資料）

\begin{itemize}
\item Thomas H. Cormen 等，\textit{Introduction to Algorithms}，第 4 版，第 10.1 節：
陣列、列優先與行優先、一個陣列和多個陣列的存法，還提到把矩陣切成小塊存放的做法。
第 30 章：多項式的係數表示和點值表示，以及用快速傅立葉轉換在 $\Theta(n \lg n)$ 時
間做多項式乘法。
\item Donald E. Knuth，\textit{The Art of Computer Programming}，第 1 冊，第 2.2.6 節：
循序配置的位址公式、三角矩陣的存法和稀疏矩陣，是這些題材比較早的系統整理。
\item Fred G. Gustavson，``Two Fast Algorithms for Sparse Matrices: Multiplication and
Permuted Transposition''，\textit{ACM Transactions on Mathematical Software}，1978：
用計數算起點的轉置，以及逐列累加的稀疏矩陣乘法，適合想知道實際的稀疏矩陣程式庫怎
麼做的人。
\item 《C 程式設計基礎》第~6~章 6.11 節、第~9~章 9.6 到 9.8 節：C 的二維陣列在記憶
體裡的排法、指標算術和指標陣列，C 語法的部分寫得比本章詳細。
\end{itemize}

\addcontentsline{toc}{section}{觀念題解答}
\section*{觀念題解答}

\begin{enumerate}
\item (a)、(c)、(e) 正確。 (a)：同一列相鄰，位移差 $1$。 (b)：行優先時同一行相鄰，
$A[i, j]$ 和 $A[i + 1, j]$ 相差 $w$；相差 $m \cdot w$ 的是 $A[i, j]$ 和
$A[i, j + 1]$。 (c)：$A^{\mathsf T}$ 的第 $i$ 行就是 $A$ 的第 $i$ 列，行優先存
$A^{\mathsf T}$ 時 $A^{\mathsf T}[j, i]$ 的位移是 $(i - 1)\,n + (j - 1)$，和列優先存
$A$ 時 $A[i, j]$ 的位移相同。 (d)：公式~(\ref{eq:arr-row}) 要乘 $n$。 (e)：
$2 \times 8 + 3 = 19$。

\item 元素個數 $44 - (-20) + 1 = 65$。 $H[0]$：位移 $20 \times 6 = 120 = \texttt{78}$，
位址 $\texttt{2F40} + \texttt{78} = \texttt{2FB8}$。 $H[31]$：位移 $51 \times 6 = 306
= \texttt{132}$，位址 $\texttt{3072}$。 最後一個元素 $H[44]$ 的位移 $64 \times 6 = 384
= \texttt{180}$，從 $\texttt{30C0}$ 開始占 $6$ 個位元組，最後一個位元組是
$\texttt{30C5}$（整個陣列共 $390$ 個位元組）。

\item $m = 12$，$n = 20$。
\begin{enumerate}
\item 列優先：$(7 - 0) \times 20 + (3 - (-4)) = 147$，位址 $4000 + 147 \times 4 = 4588$。
行優先：$(3 - (-4)) \times 12 + (7 - 0) = 91$，位址 $4000 + 91 \times 4 = 4364$。
\item 位移 $(4400 - 4000)/4 = 100$。 列優先：$100 = 5 \times 20 + 0$，是 $P[5, -4]$。
行優先：$100 = 8 \times 12 + 4$，行是 $-4 + 8 = 4$、列是 $0 + 4 = 4$，是 $P[4, 4]$。
\end{enumerate}

\item 行優先的位移是 $(j - 1)\,m + (i - 1)$。 兩式相減：$1432 - 1276 = 156 =
[(7 - 3)\,m + (5 - 2)] \times 4$，$4m + 3 = 39$，$m = 9$。 代回：$1276 = \alpha + [(3 - 1)
\times 9 + 1] \times 4 = \alpha + 76$，$\alpha = 1200$。 $Q[9, 8]$ 的位移是
$7 \times 9 + 8 = 71$，位址 $1200 + 284 = 1484$。

\item 令 $e_1 = i - l_1 \in \{0, \ldots, m - 1\}$、$e_2 = j - l_2 \in \{0, \ldots, n - 1\}$，
$\mathrm{off} = e_1 n + e_2$。 範圍：最小 $0$，最大 $(m - 1)n + (n - 1) = mn - 1$。 一對
一：若 $e_1 n + e_2 = e_1' n + e_2'$，則 $(e_1 - e_1')\,n = e_2' - e_2$，右邊的絕對值
小於 $n$，所以兩邊都是 $0$，$e_1 = e_1'$、$e_2 = e_2'$。 $mn$ 個下標對應到 $mn$ 個不
同的值，又都落在 $0..mn - 1$，所以剛好一對一對應。 反過來的公式就是除法的商和餘數：
$i = l_1 + \lfloor \mathrm{off} / n \rfloor$，$j = l_2 + (\mathrm{off} \bmod n)$。

\item (a) 內層是 \texttt{i} 加一，跳過一整列：$300 \times 8 = 2400$ 位元組。 (b)
$(0 - 399) \times 2400 + 8 = -957592$，往回跳了將近整個陣列。 (c) 對調後內層是
\texttt{j} 加一，相差 $8$ 位元組；換列時從 \texttt{a[i][299]} 到
\texttt{a[i + 1][0]}，位移 $(i + 1) \times 300 - (i \times 300 + 299) = 1$，也只差 $8$。
整個陣列從頭到尾依序寫一遍。

\item 共 $8 \times 9 / 2 = 36$ 格。 $A[6][2]$：第 $0$ 到第 $5$ 列共 $1 + 2 + \cdots + 6
= 21$ 格，所以是 \texttt{b[21 + 2] = b[23]}（公式 $i(i + 1)/2 + j$）。 對稱時
$A[2][6] = A[6][2]$，也讀 \texttt{b[23]}。

\item 列優先：$k = (4 - 1)(20 - 4 + 2)/2 + (7 - 4 + 1) = 27 + 4 = 31$，位址
$600 + 30 \times 8 = 840$。 驗算：第 $1$、$2$、$3$ 列各有 $10$、$9$、$8$ 個，共 $27$；
第 $4$ 列從 $A[4, 4]$ 數到 $A[4, 7]$ 是第 $4$ 個。 行優先：$k = 7 \times 6 / 2 + 4 =
25$，位址 $600 + 24 \times 8 = 792$。 驗算：第 $1$ 到 $6$ 行各有 $1$ 到 $6$ 個，共
$21$；第 $7$ 行從 $A[1, 7]$ 數到 $A[4, 7]$ 是第 $4$ 個。

\item 第 $i$ 列占 $B$ 的第 $i(i - 1)/2 + 1$ 到第 $i(i + 1)/2$ 格。 所以 $i$ 是滿足
$i(i + 1)/2 \ge k$ 的最小正整數。 $i^2 + i - 2k \ge 0$ 的正根是
$(\sqrt{8k + 1} - 1)/2$，因此
\[
i = \left\lceil \frac{\sqrt{8k + 1} - 1}{2} \right\rceil ,\qquad
j = k - \frac{i(i - 1)}{2} .
\]
$k = 23$：$\sqrt{185} \approx 13.6$，$(13.6 - 1)/2 \approx 6.3$，$i = 7$；$j = 23 - 21
= 2$，是 $A[7, 2]$。 驗算 $7 \times 6/2 + 2 = 23$。

\item (a) $3 \times 50 - 2 = 148$。 (b) $k = 2i + j - 2$：$A[20, 21]$ 是 $40 + 21 - 2 =
59$，$A[50, 49]$ 是 $100 + 49 - 2 = 147$。 (c) $k = (i - 1) \times 3 + (j - i + 2)$：
$A[20, 21]$ 是 $57 + 3 = 60$，$A[50, 49]$ 是 $147 + 1 = 148$；$E$ 共 $150$ 格，比
(a) 多 $2$ 格（第 $1$ 列左邊、第 $50$ 列右邊各一格）。

\item $\mathit{cnt} = (2, 2, 0, 1, 2, 1)$，$\mathit{start} = (0, 2, 4, 4, 5, 7)$。 每一步：
\begin{center}
\begin{tabular}{cccc}
\toprule
$p$ & 三元組 & 放進 & 之後的 $\mathit{start}[c]$ \\
\midrule
0 & $(0,1,7)$ & $B[2]$ & $\mathit{start}[1] = 3$ \\
1 & $(0,4,2)$ & $B[5]$ & $\mathit{start}[4] = 6$ \\
2 & $(1,1,-3)$ & $B[3]$ & $\mathit{start}[1] = 4$ \\
3 & $(2,0,6)$ & $B[0]$ & $\mathit{start}[0] = 1$ \\
4 & $(2,5,1)$ & $B[7]$ & $\mathit{start}[5] = 8$ \\
5 & $(3,3,9)$ & $B[4]$ & $\mathit{start}[3] = 5$ \\
6 & $(4,0,-4)$ & $B[1]$ & $\mathit{start}[0] = 2$ \\
7 & $(4,4,8)$ & $B[6]$ & $\mathit{start}[4] = 7$ \\
\bottomrule
\end{tabular}
\end{center}
$B$（$6 \times 5$）：$(0,2,6)$、$(0,4,-4)$、$(1,0,7)$、$(1,1,-3)$、$(3,3,9)$、$(4,0,2)$、$(4,4,8)$、$(5,2,1)$。 第 $2$ 行沒有元素，$\mathit{start}[2] = \mathit{start}[3] = 4$，結果的第
$2$ 列是空的。 把範例~\ref{ex:arr-sparse} 的 \texttt{g[]} 和大小換成這一組，
程式印出的追蹤和上表相同。

\item $\mathit{cnt} = (2, 0, 2)$，$\mathit{start} = (0, 2, 2)$。 依序放：$(2,0,6)$ 進
$B[0]$ 成為 $(0,2,6)$；$(0,0,1)$ 進 $B[1]$ 成為 $(0,0,1)$；$(1,2,4)$ 進 $B[2]$ 成為
$(2,1,4)$；$(0,2,5)$ 進 $B[3]$ 成為 $(2,0,5)$。 輸出 $(0,2,6)$、$(0,0,1)$、$(2,1,4)$、$(2,0,5)$。 每個元素都對調了，而且列號已經分好組，所以它代表的矩陣是
$A^{\mathsf T}$；但同一列裡的行號是 $2, 0$ 和 $1, 0$，沒有排好。 \proc{Simple-Transpose}
也一樣：它對每個 $c$ 由前往後掃 $A$，同一列裡的順序跟著 $A$ 的順序。 兩者都只保證
「列」排好，「同一列依行排好」要靠輸入依列排好。

\item $B'$（$2 \times 3$）：$(0,0,2)$、$(0,2,1)$、$(1,0,1)$、$(1,1,-1)$、$(1,2,4)$。
$\mathit{ra} = (0, 2, 3)$，$\mathit{rb} = (0, 2, 5)$。 四對：
\begin{itemize}
\item $(0, 0)$：$A$ 第 $0$ 列 0:1、2:$-2$，$B'$ 第 $0$ 列 0:2、2:1，配對 $k = 0, 2$，
$1 \cdot 2 + (-2) \cdot 1 = 0$，不存。
\item $(0, 1)$：$B'$ 第 $1$ 列 0:1、1:$-1$、2:4，配對 $k = 0, 2$，$1 - 8 = -7$。
\item $(1, 0)$：$A$ 第 $1$ 列 1:3，$B'$ 第 $0$ 列沒有 $k = 1$，$0$。
\item $(1, 1)$：配對 $k = 1$，$3 \cdot (-1) = -3$。
\end{itemize}
$C = (0,1,-7), (1,1,-3)$，共 $5$ 次乘法（存滿要 $2 \times 3 \times 2 = 12$ 次）。
$(0, 0)$ 做了兩次乘法，和是 $0$，沒有存。

\item 係數由高到低是 $1, 0, -4, 2, 0, -7$。 $s$ 依序是 $1$；$1 \cdot (-2) + 0 = -2$；
$(-2)(-2) - 4 = 0$；$0 \cdot (-2) + 2 = 2$；$2 \cdot (-2) + 0 = -4$；$(-4)(-2) - 7 = 1$。
$p(-2) = 1$（驗算：$-32 + 32 + 8 - 7 = 1$）。 除以 $(x + 2) = (x - (-2))$ 的商是
$x^4 - 2x^3 + 0x^2 + 2x - 4$，餘數 $1$。 共 $5$ 次乘法。

\item 展開右邊的乘積：
\[
(x - a)\sum_{i=1}^{d} s_i\,x^{i-1} = \sum_{i=1}^{d} s_i\,x^{i} - \sum_{i=0}^{d-1} a\,s_{i+1}\,x^{i} .
\]
比較各次的係數：$x^d$ 的係數是 $s_d = c_d$；$1 \le i \le d - 1$ 時，$x^i$ 的係數是
$s_i - a\,s_{i+1} = c_i$（由 $s_i$ 的定義）；常數項是 $-a\,s_1$，再加上 $s_0 = a\,s_1
+ c_0$，得到 $c_0$。 所以右邊等於 $\sum c_k x^k = p(x)$。 代入 $x = a$，第一項是 $0$，
得到 $p(a) = s_0$。 這也證明了餘式定理：$p(x)$ 除以 $(x - a)$ 的餘數是 $p(a)$。

\item (a) 係數陣列 $10^6 + 1$ 格；非零項陣列 $3$ 項、$6$ 格。 (b) $f + g$：係數陣列
約 $10^6$ 步，非零項陣列約 $3 + 40$ 步。 $f \cdot g$：係數陣列約 $(10^6)^2 = 10^{12}$
次乘法；非零項陣列 $3 \times 40 = 120$ 個乘積，用「一項乘整個多項式再相加」合併，相加
的總步數不超過 $40 + 80 + 120$ 的量級，全部在幾百步以內。 (c) 係數陣列 $31$ 格，非零
項陣列要 $62$ 格，用係數陣列。
\end{enumerate}

\chapter{鏈結串列}
\label{ch:lists}

% 本章的圖共用的樣式：一個節點畫成「資料｜next」兩格。
\tikzset{
    lkey/.style={draw=black!60, minimum width=0.85cm, minimum height=0.65cm,
        font=\ttfamily\small, inner sep=1pt},
    lptr/.style={draw=black!60, minimum width=0.40cm, minimum height=0.65cm, inner sep=0pt},
    lhl/.style={lkey, draw=memStroke, thick, fill=memFill},
    lhd/.style={lkey, fill=black!8},
    lname/.style={font=\small\ttfamily},
    lnk/.style={-{Stealth[length=1.8mm]}, thick, draw=black!70},
    lnew/.style={lnk, draw=memStroke},
    lold/.style={lnk, densely dashed, draw=black!30},
    stepno/.style={circle, draw=optStroke, fill=optFill, inner sep=0.5pt,
        font=\scriptsize\bfseries, text=optStroke},
}
% \lnode[樣式]{名稱}{座標}{內容}：資料格叫 <名稱>，next 格叫 <名稱>p。
\newcommand{\lnode}[4][lkey]{%
    \node[#1] (#2) at (#3) {#4};
    \node[lptr, anchor=west] (#2p) at (#2.east) {};
}
% \lnil{名稱}：next 是 NIL，在 next 格畫一條斜線。
\newcommand{\lnil}[1]{\draw[black!60] (#1p.south west) -- (#1p.north east);}
% \llink[樣式]{名稱}{路徑}：從 <名稱> 的 next 格畫出指標，例如 \llink{a}{-- (b)}。
\newcommand{\llink}[3][lnk]{%
    \fill[black!70] (#2p.center) circle (1.1pt);
    \draw[#1] (#2p.center) #3;
}

社團的迎新活動開放報名，總召把報名的同學依學號由小到大寫在一張紙上。 每天都有
人新報名，也有人取消。 新報名的人要插進中間某個位置，名單如果寫在一整排連續的格子
裡，插入一個人就得把後面所有的人往後抄一格。 換一種記法：每個人旁邊只註明「下一位
是誰」，插入時只要改一兩個「下一位」。 名單有 $n$ 個人時，兩種記法插入一個人各要
改幾個地方？ 第二種記法有沒有什麼事變得比較慢？ 往下讀之前，先想一想。

第二種記法就是\textbf{鏈結串列}（linked list）。 這一章先講單向串列的搜尋、插入和
刪除，再講幾種變形：帶頭節點的串列、環狀串列、雙向串列。 接著用串列表示多項式和稀
疏矩陣，說明怎麼自己管理一池節點，再整理反轉、合併、找中間這三個最常考的操作，最後
用快慢指標判斷串列有沒有環、找出環從哪裡開始。 每個操作都會給虛擬碼、C 程式、手算
追蹤和成本。

這一章用到《C 程式設計基礎》第~14~章的結構指標（\texttt{p->next}）、第~19~章的
\texttt{malloc}、\texttt{free} 和串列入門。 那一章已經寫過頭端插入、指標的指標和
刪除；這一章從資料結構的角度重新整理，重點放在成本、正確性和考卷上的追蹤，用到時
會先在框裡複習。

\section{單向鏈結串列}
\label{sec:ll-singly}

\subsection{節點與串列}

鏈結串列由\textbf{節點}（node）組成。 單向串列（singly linked list）的每個節點有
兩個欄位：$\mathit{key}$ 存資料，$\mathit{next}$ 指向下一個節點；最後一個節點的
$\mathit{next}$ 是 \NIL。 串列本身只記一個指標 $L.\mathit{head}$，指向第一個節點；
空串列的 $L.\mathit{head} = \NIL$。 實際的節點常常還有別的資料欄位（姓名、成績），
虛擬碼只寫 $\mathit{key}$ 和指標，其他欄位跟著節點一起搬。

\begin{center}
\begin{tikzpicture}
\node[lname] (hd) at (0,0) {head};
\lnode{a}{1.4,0}{115}
\lnode{b}{3.4,0}{160}
\lnode{c}{5.4,0}{207}
\lnode{d}{7.4,0}{342}
\draw[lnk] (hd) -- (a);
\llink{a}{-- (b)}
\llink{b}{-- (c)}
\llink{c}{-- (d)}
\lnil{d}
\end{tikzpicture}
\diagcaption{有四個節點的單向串列。 每個節點左格是 $\mathit{key}$，右格是
$\mathit{next}$；斜線表示 \NIL。}
\label{fig:ll-singly}
\end{center}

節點在記憶體裡不必相鄰，每個節點各自用 \texttt{malloc} 要來，彼此只靠
$\mathit{next}$ 串起來。 這帶來兩個後果。 好處是：只要手上有某個節點 $p$，在它後面
接一個新節點只要改兩個指標，和串列多長無關。 壞處是：想拿第 $k$ 個節點，只能從
$\mathit{head}$ 開始一個一個往下走，沒有辦法像陣列那樣用位址公式直接算出來（第~4~
章）。 下表把兩者並排。

\begin{center}
\begin{tabular}{lll}
\toprule
操作（$n$ 個元素） & 陣列（依序存放） & 單向串列（只記 $\mathit{head}$） \\
\midrule
取第 $k$ 個 & $\Theta(1)$ & $\Theta(k)$ \\
找某個 $\mathit{key}$（沒有排序） & $\Theta(n)$ & $\Theta(n)$ \\
找某個 $\mathit{key}$（已經排序） & $\Theta(\lg n)$，二分搜尋 & $\Theta(n)$ \\
在最前面插入 & $\Theta(n)$，全部往後挪 & $\Theta(1)$ \\
已知節點 $p$，在它後面插入或刪除 & $\Theta(n)$ & $\Theta(1)$ \\
額外記憶體 & 沒有（但要先開好容量） & 每個節點一個指標 \\
\bottomrule
\end{tabular}
\end{center}

排好的串列不能二分搜尋：二分搜尋要直接跳到中間那一格，串列走到中間就要
$\Theta(n)$ 步。 所以「串列比陣列好」或「陣列比串列好」都不對，要看哪些操作比較常
做。 第~6~章的堆疊與佇列、第~12~章雜湊的鏈結法、第~21~章圖的鄰接串列，都是挑串列
擅長的操作來用。

\subsection{搜尋、插入、刪除}

\textbf{搜尋}從 $\mathit{head}$ 往下走，遇到 $\mathit{key}$ 相同的節點就停。

\begin{pseudo}[label=alg:ll-search]{在串列裡找 $k$}
\Function{List-Search}{$L, k$}
    \State $x \gets L.\mathit{head}$
    \While{$x \ne \NIL$ \AND $x.\mathit{key} \ne k$}
        \State $x \gets x.\mathit{next}$
    \EndWhile
    \State \Return $x$ \Comment{找不到時是 \NIL}
\EndFunction
\end{pseudo}

最壞情況（$k$ 不在串列裡）要走過全部 $n$ 個節點，$\Theta(n)$。 第 3 行的
\textbf{and} 先檢查 $x \ne \NIL$，$x$ 是 \NIL 時不會去讀 $x.\mathit{key}$；C 的
\texttt{\&\&} 也是這樣，兩個條件的順序不能對調。

\textbf{插入}要先說好插在哪裡。 單向串列的節點只知道下一個是誰，不知道上一個是誰，
所以插入和刪除都以「前一個節點」$p$ 為準：把新節點 $x$ 接在 $p$ 後面，或拿掉 $p$
後面那一個。 $p = \NIL$ 代表「在最前面」。

\begin{pseudo}[label=alg:ll-insert]{把節點 $x$ 接在 $p$ 後面}
\Procedure{List-Insert-After}{$L, p, x$}
    \If{$p = \NIL$} \Comment{接在最前面}
        \State $x.\mathit{next} \gets L.\mathit{head}$
        \State $L.\mathit{head} \gets x$
    \Else
        \State $x.\mathit{next} \gets p.\mathit{next}$
        \State $p.\mathit{next} \gets x$
    \EndIf
\EndProcedure
\end{pseudo}

\begin{pseudo}[label=alg:ll-delete]{拿掉 $p$ 後面的節點}
\Function{List-Delete-After}{$L, p$}
    \If{$p = \NIL$} \Comment{拿掉第一個}
        \State $x \gets L.\mathit{head}$
        \State $L.\mathit{head} \gets x.\mathit{next}$
    \Else
        \State $x \gets p.\mathit{next}$
        \State $p.\mathit{next} \gets x.\mathit{next}$
    \EndIf
    \State \Return $x$ \Comment{由呼叫的人決定要不要釋放}
\EndFunction
\end{pseudo}

兩個程序都是 $\Theta(1)$：不管串列多長，都只改兩個指標。 前提是被拿掉的節點存在
（$p = \NIL$ 時串列不是空的，否則 $p.\mathit{next} \ne \NIL$）。 實際上要刪除某個
$\mathit{key}$ 時，還得先花 $\Theta(n)$ 找到它的前一個節點，所以「串列刪除是
$\Theta(1)$」只在已經拿到前一個節點時成立。

插入的兩行順序不能對調。 下圖把 $\mathit{key} = 180$ 的節點 $x$ 接在 $160$ 後面。
如果先做 $p.\mathit{next} \gets x$，$207$ 那個節點就再也沒有人指著，接下來的
$x.\mathit{next} \gets p.\mathit{next}$ 讀到的是 $x$ 自己，$x$ 指向自己，後半段整條
遺失。

\begin{center}
\begin{tikzpicture}
% 插入前
\node[anchor=east, font=\small] at (-0.2,0) {插入前};
\lnode{a}{0.6,0}{115}
\lnode[lhl]{b}{2.6,0}{160}
\lnode{c}{4.6,0}{207}
\lnode{d}{6.6,0}{342}
\llink{a}{-- (b)}
\llink{b}{-- (c)}
\llink{c}{-- (d)}
\lnil{d}
\node[font=\scriptsize, text=memStroke] at (2.8,0.6) {$p$};
% 第一步
\node[anchor=east, font=\small] at (-0.2,-1.9) {第 6 行};
\lnode{a2}{0.6,-1.6}{115}
\lnode[lhl]{b2}{2.6,-1.6}{160}
\lnode{c2}{4.6,-1.6}{207}
\lnode{d2}{6.6,-1.6}{342}
\lnode{x2}{3.6,-2.5}{180}
\llink{a2}{-- (b2)}
\llink{b2}{-- (c2)}
\llink{c2}{-- (d2)}
\lnil{d2}
\llink[lnew]{x2}{-- (c2.south)}
\node[font=\scriptsize, text=black!60] at (2.6,-2.5) {$x$};
% 第二步
\node[anchor=east, font=\small] at (-0.2,-4.1) {第 7 行};
\lnode{a3}{0.6,-3.8}{115}
\lnode[lhl]{b3}{2.6,-3.8}{160}
\lnode{c3}{4.6,-3.8}{207}
\lnode{d3}{6.6,-3.8}{342}
\lnode{x3}{3.6,-4.7}{180}
\llink{a3}{-- (b3)}
\fill[black!70] (b3p.center) circle (1.1pt);
\draw[lold] (b3p.center) -- (c3);
\draw[lnew] (b3p.center) -- (x3.north west);
\llink{x3}{-- (c3.south)}
\llink{c3}{-- (d3)}
\lnil{d3}
\end{tikzpicture}
\diagcaption{在 $p$（$160$）後面插入 $x$（$180$）。 橘色是這一步新設的指標，灰色虛線
是被蓋掉的舊指標。 先讓 $x$ 接住後半段，再讓 $p$ 指向 $x$。}
\label{fig:ll-insert}
\end{center}

刪除也有一個順序：先把要刪的節點記在變數 $x$ 裡，接好前後，最後才釋放 $x$。 如果
先釋放，$x.\mathit{next}$ 就讀不到了。

\begin{examnote}
追蹤指標的題目，一步畫一張圖：先畫操作前的串列，每執行一行指定，就把那一個指標的
新箭頭畫上、舊箭頭劃掉。 題目問「執行完後串列是什麼」時，從 $\mathit{head}$ 沿著
箭頭重新走一次再寫答案，不要憑印象寫。 有些節點會變成沒有人指著，答案裡要寫出它們
已經不在串列上。
\end{examnote}

\subsection{依序插入與指標的指標}

報名名單要依學號排好。 插入新的學號 $k$ 時，從頭走到第一個 $\mathit{key} \ge k$
的節點 $q$，把新節點接在 $q$ 的前一個節點 $p$ 後面。 單向串列沒辦法從 $q$ 回頭找
$p$，所以走的時候讓 $p$ 跟在 $q$ 後面一格。

\begin{pseudo}[label=alg:ll-sorted-insert]{依 $\mathit{key}$ 由小到大插入}
\Procedure{Sorted-Insert}{$L, x$}
    \State $p \gets \NIL$ \Comment{$p$ 永遠跟在 $q$ 後面一格}
    \State $q \gets L.\mathit{head}$
    \While{$q \ne \NIL$ \AND $q.\mathit{key} < x.\mathit{key}$}
        \State $p \gets q$
        \State $q \gets q.\mathit{next}$
    \EndWhile
    \State \Call{List-Insert-After}{$L, p, x$}
\EndProcedure
\end{pseudo}

為什麼插完之後串列還是排好的？ 用第~1~章的迴圈不變量：每次檢查 \textbf{while} 的
條件時，從 $\mathit{head}$ 到 $p$（含 $p$）的每個節點都滿足 $\mathit{key} <
x.\mathit{key}$，而且 $q$ 是 $p$ 的下一個節點（$p = \NIL$ 時 $q$ 是第一個節點）。
一開始 $p = \NIL$，前一段是空的，成立；每一輪 $q.\mathit{key} < x.\mathit{key}$ 才
往前走，$p$ 移到 $q$ 之後仍然成立。 迴圈停下時，$q = \NIL$ 或 $q.\mathit{key} \ge
x.\mathit{key}$。 原本的串列排好了，$p$ 之前的都不大於 $p$，所以把 $x$ 放在 $p$ 和
$q$ 中間，左邊都比 $x$ 小、右邊都不比 $x$ 小，串列仍然排好。 最壞情況走完整條串
列，$\Theta(n)$。

虛擬碼裡有兩個地方要分開處理「在最前面」：$p = \NIL$ 的初值，以及
\proc{List-Insert-After} 的第一個分支。 C 有一個寫法可以把它們合成一種：不記「前一
個節點」，改記「指向 $q$ 的那個指標欄位在哪裡」。 $q$ 是第一個節點時，指著它的是
$\mathit{head}$ 這個變數；否則是前一個節點的 \texttt{next} 欄位。 兩者的型別都是
\texttt{Node *}，所以它們的位址都是 \texttt{Node **}，可以用同一個變數
\texttt{link} 記。

\begin{cbookrecap}[第~14~章 14.5 節、第~19~章 19.9、19.10 節]
\texttt{p->next} 是 \texttt{(*p).next} 的簡寫：先沿著指標 \texttt{p} 找到結構，再
取它的 \texttt{next} 欄位。

C 的參數都是傳值。 函式要改呼叫端的 \texttt{Node *head}，必須收 \texttt{Node **}，
呼叫時傳 \texttt{\&head}，在函式裡寫 \texttt{*head = ...}。 同樣地，
\texttt{\&p->next} 是節點 \texttt{p} 的 \texttt{next} 欄位的位址；拿到它，就能改
「\texttt{p} 的下一個是誰」。 第~19~章用這個寫法把「刪第一個」和「刪中間」寫成同一
段程式，這一節用同樣的寫法做依序插入。
\end{cbookrecap}

\begin{center}
\begin{tikzpicture}
\node[lname, draw=black!60, minimum width=1.0cm, minimum height=0.65cm] (hd) at (0,0) {};
\node[lname, above=1pt] at (hd.north) {head};
\lnode{a}{1.8,0}{115}
\lnode{c}{3.8,0}{207}
\lnode{d}{5.8,0}{342}
\fill[black!70] (hd.center) circle (1.1pt);
\draw[lnk] (hd.center) -- (a);
\llink{a}{-- (c)}
\llink{c}{-- (d)}
\lnil{d}
\node[lname, text=memStroke] (lk1) at (0,-1.4) {link（開始）};
\draw[lnew] (lk1) -- (hd.south);
\node[lname, text=memStroke] (lk2) at (2.4,-1.4) {link（停下）};
\draw[lnew] (lk2) -- (ap.south);
\end{tikzpicture}
\diagcaption{插入 $160$ 時 \texttt{link} 的位置。 它一開始指向 \texttt{head} 這個
變數；$115 < 160$，往前一格，改指向 $115$ 的 \texttt{next} 欄位；那個欄位指著
$207 \ge 160$，停下。 新節點就接在 \texttt{*link} 這個位置。}
\label{fig:ll-link}
\end{center}

\texttt{link} 停下時，\texttt{*link} 就是虛擬碼的 $q$。 插入只要兩行：
\texttt{p->next = *link;} 讓新節點接住 $q$，\texttt{*link = p;} 讓原本指著 $q$ 的
那個指標（不管是 \texttt{head} 還是某個 \texttt{next}）改指新節點。 刪除也一樣：
\texttt{*link = dead->next;}。 下面的程式用它管理報名名單。 為了圖好畫，學號只寫
末三碼。

\begin{example}[label=ex:ll-signup]{c}{(signup.c) 依學號排好的報名名單}
#include <stdio.h>
#include <stdlib.h>

typedef struct Node {
    int id;                 // 學號（這裡只用末三碼）
    struct Node *next;
} Node;

// 回傳「指向第一個 id >= x 的節點」的那個指標欄位的位址，
// 可能是 head 本身，也可能是某個節點的 next 欄位。
Node **find_link(Node **head, int x)
{
    Node **link = head;
    while (*link != NULL && (*link)->id < x)
        link = &(*link)->next;
    return link;
}

// 依學號由小到大插入。成功回傳 1，已經在名單上回傳 0，
// 要不到記憶體回傳 -1。
int add(Node **head, int x)
{
    Node **link = find_link(head, x);
    if (*link != NULL && (*link)->id == x)
        return 0;
    Node *p = malloc(sizeof *p);
    if (p == NULL)
        return -1;
    p->id = x;
    p->next = *link;        // 先接住後面那一段
    *link = p;              // 再讓前面指向新節點
    return 1;
}

// 刪除學號 x。成功回傳 1，找不到回傳 0。
int drop(Node **head, int x)
{
    Node **link = find_link(head, x);
    if (*link == NULL || (*link)->id != x)
        return 0;
    Node *dead = *link;
    *link = dead->next;     // 前面直接跳過 dead
    free(dead);
    return 1;
}

void print_list(const Node *p)
{
    printf("list:");
    for (; p != NULL; p = p->next)
        printf(" %d", p->id);
    printf("\n");
}

int main(void)
{
    Node *head = NULL;
    int join[] = {207, 115, 342, 160, 115, 101};
    int quit[] = {115, 300, 342};
    for (int i = 0; i < 6; i++) {
        int r = add(&head, join[i]);
        if (r < 0)
            break;
        printf("add %d: %s\n", join[i], r ? "ok" : "already in");
    }
    print_list(head);
    for (int i = 0; i < 3; i++)
        printf("drop %d: %s\n", quit[i],
               drop(&head, quit[i]) ? "ok" : "not found");
    print_list(head);
    while (head != NULL)    // 全部釋放：每次刪掉第一個
        drop(&head, head->id);
    return 0;
}
\end{example}
\iosample{}{add 207: ok \\ add 115: ok \\ add 342: ok \\ add 160: ok \\ add 115: already in \\ add 101: ok \\ list: 101 115 160 207 342 \\ drop 115: ok \\ drop 300: not found \\ drop 342: ok \\ list: 101 160 207}

這支程式有幾點可以對照前面的說明：

\begin{itemize}
\item 加入 $115$ 時它插在最前面，$342$ 插在最後面，$160$ 插在中間，三種情況用的是
同一段程式。 刪除 $115$ 時它是第一個節點，$342$ 是最後一個，也一樣。
\item \texttt{find\_link} 只用 \texttt{<} 往前走，停下時 \texttt{*link} 是第一個
\texttt{id >= x} 的節點，所以只要再比一次 \texttt{==} 就知道 \texttt{x} 在不在名
單上。 重複報名、取消不存在的學號都在這裡擋掉。
\item 程式最後一直刪第一個節點，直到名單變空，所有節點都釋放了。 \texttt{drop} 先
把 \texttt{dead->next} 接好才 \texttt{free(dead)}；順序反過來，就是讀了已經釋放的記
憶體，AddressSanitizer（《C 程式設計基礎》第~19~章 19.5 節）會直接報錯。
\end{itemize}

下表是加入 $160$ 時 \texttt{find\_link} 的過程（這時名單是 $115, 207, 342$）。

\begin{center}
\begin{tabular}{clll}
\toprule
檢查 & \texttt{link} 指向 & \texttt{*link} 指向的節點 & 結果 \\
\midrule
1 & \texttt{head} & $115$ & $115 < 160$，往前 \\
2 & $115$ 的 \texttt{next} & $207$ & $207 \ge 160$，停下 \\
\bottomrule
\end{tabular}
\end{center}

插入之後 $115$ 的 \texttt{next} 指向新節點 $160$，$160$ 的 \texttt{next} 指向
$207$。 每次操作最多走過整條名單，所以 $n$ 個人的名單每次報名或取消都是 $O(n)$。

\section{帶頭節點、環狀與雙向串列}
\label{sec:ll-variants}

上一節的程式裡，「在最前面」一直是要另外想的情況。 這一節的三種變形各自解決單向串
列的一個不方便：頭節點去掉特殊情況，環狀串列讓頭尾都能 $\Theta(1)$ 操作，雙向串列
讓節點知道前一個是誰。

\subsection{頭節點}

\textbf{頭節點}（header node，也叫 dummy node）是放在串列最前面、不存資料的節點。
串列改成記住頭節點 $h$，真正的第一個元素是 $h.\mathit{next}$；空串列時
$h.\mathit{next} = \NIL$，但 $h$ 本身一直存在。 有了它，每個資料節點前面都有一個
節點，\proc{List-Insert-After} 和 \proc{List-Delete-After} 的 $p$ 永遠不是 \NIL，
虛擬碼~\ref{alg:ll-insert} 第 2--5 行的特殊情況可以整段刪掉，只剩兩行指定。

\begin{center}
\begin{tikzpicture}
\node[lname] (h) at (-0.3,0) {h};
\lnode[lhd]{z}{0.8,0}{}
\lnode{a}{2.8,0}{12}
\lnode{b}{4.8,0}{31}
\lnil{b}
\draw[lnk] (h) -- (z);
\llink{z}{-- (a)}
\llink{a}{-- (b)}
\node[lname] (h2) at (7.4,0) {h};
\lnode[lhd]{e}{8.5,0}{}
\lnil{e}
\draw[lnk] (h2) -- (e);
\node[font=\small] at (3.0,-0.8) {兩個元素};
\node[font=\small] at (8.8,-0.8) {空串列};
\end{tikzpicture}
\diagcaption{帶頭節點的串列。 灰底的頭節點不存資料；空串列也有頭節點。}
\label{fig:ll-header}
\end{center}

代價是多一個節點的空間，而且走訪時要從 $h.\mathit{next}$ 開始。 C 程式常把頭節點
宣告成函式裡的區域變數 \texttt{Node head}，用 \texttt{\&head} 當「最後一個節點」開始
往後接，最後回傳 \texttt{head.next}；這樣接第一個節點時不用判斷串列是不是空的。
第~\ref{sec:ll-poly}~節的多項式加法和第~\ref{sec:ll-ops}~節的合併都用這個寫法。 它
和上一節的 \texttt{Node **link} 解決的是同一件事，選一個自己寫得順的就好。

\subsection{環狀串列}

把最後一個節點的 $\mathit{next}$ 接回第一個節點，就是\textbf{環狀串列}（circular
list）。 環上沒有 \NIL，任何一個節點往下走都會繞回自己。

環狀串列通常記的是\textbf{最後一個}節點 $R.\mathit{last}$。 第一
個節點就是 $R.\mathit{last}.\mathit{next}$，一步就拿得到；而最後一個節點正是「最前
面」的前一個節點，所以在最前面和在最後面插入，都是接在 $R.\mathit{last}$ 後面，差
別只在插完之後 $R.\mathit{last}$ 要不要移到新節點。

\begin{pseudo}[label=alg:ll-ring-push]{在環狀串列的最後面插入}
\Procedure{Ring-Push-Back}{$R, x$}
    \If{$R.\mathit{last} = \NIL$} \Comment{空的環：自己接自己}
        \State $x.\mathit{next} \gets x$
    \Else
        \State $x.\mathit{next} \gets R.\mathit{last}.\mathit{next}$
        \State $R.\mathit{last}.\mathit{next} \gets x$
    \EndIf
    \State $R.\mathit{last} \gets x$
\EndProcedure
\end{pseudo}

把第 7 行改成「只有原本是空的環時才做」，就是在最前面插入。 兩者都是
$\Theta(1)$。 如果只記第一個節點，在最前面插入時要改最後一個節點的
$\mathit{next}$，得先繞一圈找到它，反而變成 $\Theta(n)$。

走訪環狀串列不能用「走到 \NIL 為止」當停止條件，那會永遠走不完。 正確的寫法是從
第一個節點出發，處理完最後一個（$x = R.\mathit{last}$）就停，或是用
\textbf{do-while}：先處理再前進，前進後回到起點就停。

\begin{center}
\begin{tikzpicture}[lkey/.append style={minimum width=1.1cm}]
\lnode{a}{0,0}{Ming}
\lnode{b}{2.4,0}{Hua}
\lnode{c}{4.8,0}{Ting}
\llink{a}{-- (b)}
\llink{b}{-- (c)}
\fill[black!70] (cp.center) circle (1.1pt);
\draw[lnk, rounded corners=6pt] (cp.center) -- ++(0.5,0) -- ++(0,-0.75) -- ($(a.south)+(0,-0.42)$) -- (a.south);
\node[lname, text=memStroke] (lst) at (5.05,1.1) {last};
\draw[lnew] (lst) -- (c.north);
\end{tikzpicture}
\diagcaption{三個人的值日環。 \texttt{last} 指向 Ting，今天值日的是
\texttt{last->next}，也就是 Ming。}
\label{fig:ll-ring}
\end{center}

下面的例子是實驗室的值日輪替：每天由環上的下一位值日；新加入的人排在這一輪的最後
（也就是下一個要值日的人的前一位）；有人離開就從環上拿掉。 「過一天」只要把
\texttt{last} 往前移一格，整個環就轉了一格，不用搬任何節點。

\begin{example}[label=ex:ll-duty]{c}{(duty\_ring.c) 用環狀串列排值日}
#include <stdio.h>
#include <stdlib.h>
#include <string.h>

typedef struct Member {
    char name[16];
    struct Member *next;
} Member;

// 環只記 last：last 是這一輪的最後一位，last->next 是今天值日的人。
// 新成員接在 last 後面，再讓 last 指向他，也就是排到這一輪的最後。
int join(Member **last, const char *name)
{
    Member *p = malloc(sizeof *p);
    if (p == NULL)
        return 0;
    snprintf(p->name, sizeof p->name, "%s", name);
    if (*last == NULL) {
        p->next = p;                // 只有一個人：自己接自己
    } else {
        p->next = (*last)->next;
        (*last)->next = p;
    }
    *last = p;
    return 1;
}

// 名字是 name 的人退出。找到回傳 1，找不到回傳 0。
int leave(Member **last, const char *name)
{
    if (*last == NULL)
        return 0;
    Member *prev = *last;
    Member *cur = prev->next;       // 從今天值日的人開始繞一圈
    do {
        if (strcmp(cur->name, name) == 0) {
            if (cur == prev)            // 環上只剩他一個
                *last = NULL;
            else {
                prev->next = cur->next;
                if (cur == *last)       // 刪掉的是最後一位
                    *last = prev;
            }
            free(cur);
            return 1;
        }
        prev = cur;
        cur = cur->next;
    } while (prev != *last);
    return 0;
}

int main(void)
{
    Member *last = NULL;
    char op, name[16];
    int day = 0;
    while (scanf(" %c", &op) == 1) {
        if (op == 'N') {            // 過一天：環轉一格
            if (last == NULL)
                printf("day %d: nobody\n", ++day);
            else {
                printf("day %d: %s\n", ++day, last->next->name);
                last = last->next;
            }
        } else if (scanf("%15s", name) == 1) {
            if (op == 'J' && !join(&last, name))
                break;
            if (op == 'L' && !leave(&last, name))
                printf("%s is not on duty list\n", name);
        }
    }
    while (last != NULL)            // 全部釋放
        leave(&last, last->next->name);
    return 0;
}
\end{example}
\iosample{J Ming \\ J Hua \\ J Ting \\ N \\ N \\ J Wei \\ N \\ L Hua \\ N \\ N \\ L Ting \\ L Ming \\ L Ting \\ N \\ N \\ L Wei \\ N}{day 1: Ming \\ day 2: Hua \\ day 3: Ting \\ day 4: Ming \\ day 5: Wei \\ Ting is not on duty list \\ day 6: Wei \\ day 7: Wei \\ day 8: nobody}

\texttt{J} 是加入，\texttt{L} 是離開，\texttt{N} 是過一天。 Wei 是在第 2 天之後加入
的，那時 \texttt{last} 是 Hua，下一個輪到 Ting。 Wei 接在 Hua 後面、Ting 前面，輪替
順序變成 Ting、Ming、Hua、Wei。 Hua 在第 3 天之後就離開了，所以第 4 天是 Ming，第 5
天輪到 Wei。 環上只剩 Wei 時，他的 \texttt{next} 指向自己，每天都是他。

\texttt{leave} 有三種情況要分開想：刪掉之後環變空了、刪掉的是 \texttt{last}（
\texttt{last} 要退回前一個）、其他。 它的迴圈用 \textbf{do-while}，停止條件是「剛檢
查完的是 \texttt{last}」，繞完一圈就停。 程式結束前，每次讓今天值日的人
（\texttt{last->next}）離開，\texttt{leave} 第一個就找到他，一次 $\Theta(1)$，直到環
變空。 第~\ref{sec:ll-avail}~節會看到，如果節點是還給自己管理的節點池，整個環可以一
次 $\Theta(1)$ 還回去。

\subsection{雙向串列}

單向串列拿到一個節點 $x$，沒有辦法 $\Theta(1)$ 把它刪掉，因為改的是前一個節點的
$\mathit{next}$，而前一個節點只能從頭找。 \textbf{雙向串列}（doubly linked list）
的每個節點多一個 $\mathit{prev}$ 指向前一個節點，兩個方向都能走，刪除任何已知的節
點都是 $\Theta(1)$。

雙向串列最好寫的形式是\textbf{帶頭節點的環狀雙向串列}：頭節點 $h$ 的
$\mathit{next}$ 是第一個元素，$\mathit{prev}$ 是最後一個元素；空串列時
$h.\mathit{next} = h.\mathit{prev} = h$。 這樣每個資料節點的前後都一定有節點（可能
是 $h$），插入、刪除完全沒有特殊情況。

\begin{center}
\begin{tikzpicture}[
    dn/.style={draw=black!60, minimum width=0.85cm, minimum height=0.65cm, font=\ttfamily\small},
]
\node[dn, fill=black!8] (h) at (0,0) {h};
\node[dn] (a) at (2.2,0) {12};
\node[dn] (b) at (4.4,0) {31};
\node[dn] (c) at (6.6,0) {47};
\foreach \u/\v in {h/a, a/b, b/c} {
    \draw[lnk] ([yshift=4pt]\u.east) -- ([yshift=4pt]\v.west);
    \draw[lnk] ([yshift=-4pt]\v.west) -- ([yshift=-4pt]\u.east);
}
\draw[lnk, rounded corners=6pt] (c.north) -- ++(0,0.55) -| (h.north);
\draw[lnk, rounded corners=6pt] ([xshift=-6pt]h.south) -- ++(0,-0.55) -| ([xshift=-6pt]c.south);
\node[font=\scriptsize, text=black!60] at (3.3,1.12) {最後一個的 next 是 h};
\node[font=\scriptsize, text=black!60] at (3.3,-1.12) {h 的 prev 是最後一個};
\end{tikzpicture}
\diagcaption{帶頭節點的環狀雙向串列。 上排箭頭是 $\mathit{next}$，下排是
$\mathit{prev}$。}
\label{fig:ll-doubly}
\end{center}

\begin{pseudo}[label=alg:ll-dinsert]{雙向串列：把 $x$ 接在 $p$ 後面}
\Procedure{DList-Insert-After}{$p, x$}
    \State $x.\mathit{prev} \gets p$
    \State $x.\mathit{next} \gets p.\mathit{next}$
    \State $p.\mathit{next}.\mathit{prev} \gets x$
    \State $p.\mathit{next} \gets x$
\EndProcedure
\end{pseudo}

\begin{pseudo}[label=alg:ll-dremove]{雙向串列：拿掉節點 $x$}
\Procedure{DList-Remove}{$x$}
    \State $x.\mathit{prev}.\mathit{next} \gets x.\mathit{next}$
    \State $x.\mathit{next}.\mathit{prev} \gets x.\mathit{prev}$
\EndProcedure
\end{pseudo}

插入要改四個指標：新節點自己的兩個、前一個的 $\mathit{next}$、後一個的
$\mathit{prev}$。 順序的限制只有一條：第 5 行會蓋掉 $p.\mathit{next}$，而第 3、4 行
都要用到舊的 $p.\mathit{next}$，所以這兩行要在第 5 行之前做。 第 4 行若改寫成
$x.\mathit{next}.\mathit{prev} \gets x$，限制就換成「要在第 3 行之後」。 在頭節點
$h$ 後面插入是放到最前面，在 $h.\mathit{prev}$ 後面插入是放到
最後面。 刪除只改兩個指標，被刪的 $x$ 自己的指標不用動。

\begin{center}
\begin{tikzpicture}[
    dn/.style={draw=black!60, minimum width=0.85cm, minimum height=0.65cm, font=\ttfamily\small},
]
\node[dn] (p) at (0,0) {31};
\node[dn] (q) at (4.4,0) {47};
\node[dn, draw=memStroke, thick, fill=memFill] (x) at (2.2,-2.0) {40};
\draw[lold] ([yshift=4pt]p.east) -- ([yshift=4pt]q.west);
\draw[lold] ([yshift=-4pt]q.west) -- ([yshift=-4pt]p.east);
\draw[lnew] (x.west) to[bend left=25] node[stepno, pos=0.5, left=2pt] {2} (p.south west);
\draw[lnew] (x.north east) -- node[stepno, pos=0.5, left=2pt] {3} (q.south west);
\draw[lnew] (q.south east) to[bend left=25] node[stepno, pos=0.5, right=2pt] {4} (x.east);
\draw[lnew] (p.south east) -- node[stepno, pos=0.5, right=2pt] {5} (x.north west);
\node[font=\scriptsize, text=memStroke] at (-0.6,0.5) {$p$};
\end{tikzpicture}
\diagcaption{在 $p$（$31$）後面插入 $x$（$40$）。 圓圈裡的數字是虛擬碼~\ref{alg:ll-dinsert}
的行號，灰色虛線是被取代的兩個舊指標。}
\label{fig:ll-dinsert}
\end{center}

翻成 C 時，節點多一個 \texttt{struct DNode *prev} 欄位，四行指定一行對一行照抄（例如
\texttt{x->prev = p;}）；空串列由 \texttt{h->next = h->prev = h} 開始。
\proc{DList-Remove} 不負責釋放，拿下來的節點可以 \texttt{free}，也可以接到別的地方，
例如再接到 $h$ 後面，就是「移到最前面」。

\subsection{幾種串列的比較}

下表整理本節三種結構和單向串列的成本。 「單向」只記 $\mathit{head}$，「單向＋tail」
另外記最後一個節點，「單向環狀」只記 $\mathit{last}$，「雙向環狀」帶頭節點。 「已知
節點 $x$」是說手上已經有指向 $x$ 的指標，不包含找到它的時間。

\begin{center}
{\small
\begin{tabular}{lcccc}
\toprule
操作 & 單向 & 單向＋tail & 單向環狀 & 雙向環狀 \\
\midrule
在最前面插入 & $\Theta(1)$ & $\Theta(1)$ & $\Theta(1)$ & $\Theta(1)$ \\
在最後面插入 & $\Theta(n)$ & $\Theta(1)$ & $\Theta(1)$ & $\Theta(1)$ \\
刪除第一個 & $\Theta(1)$ & $\Theta(1)$ & $\Theta(1)$ & $\Theta(1)$ \\
刪除最後一個 & $\Theta(n)$ & $\Theta(n)$ & $\Theta(n)$ & $\Theta(1)$ \\
刪除已知節點 $x$ & $\Theta(n)$ & $\Theta(n)$ & $\Theta(n)$ & $\Theta(1)$ \\
每個節點的指標數 & 1 & 1 & 1 & 2 \\
\bottomrule
\end{tabular}}
\end{center}

「刪除最後一個」在單向串列裡都是 $\Theta(n)$，即使記了 \texttt{tail} 也一樣：刪掉
之後 \texttt{tail} 要改成倒數第二個，而倒數第二個只能從頭找。 第~6~章用串列做佇列
時，會選「在最後面插入、刪除第一個」這兩個在單向串列上都便宜的操作。

\begin{examnote}
寫串列的程式題時，寫完後拿四種情況各檢查一次：空串列、只有一個節點、操作的是第一
個節點、操作的是最後一個節點。 大部分錯誤（讀到 \NIL 的欄位、\texttt{head} 或
\texttt{last} 沒有更新）都出在這四種情況。 有頭節點或用 \texttt{Node **} 時，可以
在答案裡說明為什麼這些情況不必另外處理。
\end{examnote}

\section{用串列表示多項式與稀疏矩陣}
\label{sec:ll-poly}

第~4~章用陣列存多項式和稀疏矩陣：多項式可以存全部的係數，也可以只存非零項；稀疏
矩陣用三元組 $(\text{列}, \text{行}, \text{值})$ 的陣列。 只存非零項省空間，但項數
會變的時候（相加後有些項相消、相乘後項數變多），陣列要預先開多大很難說，插入一項還
得挪動後面的項。 改用串列，每一項是一個節點，要多少就配置多少。

\subsection{多項式}

多項式的每一個非零項存成一個節點，欄位是係數 $\mathit{coef}$、次方 $\mathit{exp}$
和 $\mathit{next}$，節點依次方由大到小串起來。 例如
\[
A(x) = 5x^{12} - 3x^{7} + x^{2} + 8, \qquad B(x) = 2x^{9} + 3x^{7} - x^{2} + 4x
\]
存成下圖的兩條串列，每個節點寫成「係數, 次方」。

\begin{center}
\begin{tikzpicture}[lkey/.append style={minimum width=1.05cm}]
\node[lname] (ha) at (-0.4,0) {A};
\lnode{a1}{0.8,0}{5,12}
\lnode{a2}{3.0,0}{-3,7}
\lnode{a3}{5.2,0}{1,2}
\lnode{a4}{7.4,0}{8,0}
\draw[lnk] (ha) -- (a1);
\llink{a1}{-- (a2)}
\llink{a2}{-- (a3)}
\llink{a3}{-- (a4)}
\lnil{a4}
\node[lname] (hb) at (-0.4,-1.1) {B};
\lnode{b1}{0.8,-1.1}{2,9}
\lnode{b2}{3.0,-1.1}{3,7}
\lnode{b3}{5.2,-1.1}{-1,2}
\lnode{b4}{7.4,-1.1}{4,1}
\draw[lnk] (hb) -- (b1);
\llink{b1}{-- (b2)}
\llink{b2}{-- (b3)}
\llink{b3}{-- (b4)}
\lnil{b4}
\end{tikzpicture}
\diagcaption{多項式 $A$、$B$ 的串列，依次方由大到小。 係數是 $0$ 的項不存。}
\label{fig:ll-poly}
\end{center}

相加的做法和合併兩條排好的串列一樣：$a$、$b$ 各指著一條串列目前的項，次方大的那一
項先抄到結果；次方一樣就把係數加起來，和不是 $0$ 才抄，兩邊一起前進。 有一邊走完
了，另一邊剩下的項全部照抄。

\begin{pseudo}[label=alg:ll-poly-add]{多項式相加}
\Function{Poly-Add}{$a, b$} \Comment{$a$、$b$ 指向兩個多項式的第一項}
    \State $h \gets$ 新的頭節點；$t \gets h$ \Comment{$t$ 是結果的最後一項}
    \While{$a \ne \NIL$ \AND $b \ne \NIL$}
        \If{$a.\mathit{exp} > b.\mathit{exp}$}
            \State $t \gets \Call{Attach}{t, a.\mathit{coef}, a.\mathit{exp}}$
            \State $a \gets a.\mathit{next}$
        \ElsIf{$a.\mathit{exp} < b.\mathit{exp}$}
            \State $t \gets \Call{Attach}{t, b.\mathit{coef}, b.\mathit{exp}}$
            \State $b \gets b.\mathit{next}$
        \Else
            \State $s \gets a.\mathit{coef} + b.\mathit{coef}$
            \If{$s \ne 0$}
                \State $t \gets \Call{Attach}{t, s, a.\mathit{exp}}$
            \EndIf
            \State $a \gets a.\mathit{next}$；$b \gets b.\mathit{next}$
        \EndIf
    \EndWhile
    \State 把 $a$ 或 $b$ 剩下的項依序 \proc{Attach} 到 $t$ 後面
    \State \Return $h.\mathit{next}$
\EndFunction
\end{pseudo}

\proc{Attach}$(t, c, e)$ 配置一個新節點、接在 $t$ 後面，並回傳這個新節點（新的最
後一項）。 結果用頭節點 $h$ 起頭，第一次 \proc{Attach} 就不用另外處理。

\textbf{成本}。 設 $A$ 有 $m$ 項、$B$ 有 $n$ 項。 迴圈每一輪比較一次次方，至少讓
$a$、$b$ 其中一個前進一格，所以迴圈最多 $m + n$ 輪；更仔細地說，最後一個被處理的
項不會在迴圈裡被比較，比較次數最多 $m + n - 1$。 剩下的項照抄，總時間
$\Theta(m + n)$，和多項式的次數無關。 如果用長度 $d + 1$ 的係數陣列（$d$ 是次方），
相加要 $\Theta(d)$；$A(x) = x^{1000} + 1$ 這種項數少、次方高的多項式，串列快得多。

\begin{example}[label=ex:ll-poly]{c}{(poly\_add.c) 多項式相加}
#include <stdio.h>
#include <stdlib.h>

typedef struct Term {
    int coef;               // 係數，不會是 0
    int exp;                // 次方，由大到小排
    struct Term *next;
} Term;

// 在 tail 後面接上一個新項，回傳新的尾巴；要不到記憶體回傳 NULL。
Term *attach(Term *tail, int coef, int exp)
{
    Term *t = malloc(sizeof *t);
    if (t == NULL)
        return NULL;
    t->coef = coef;
    t->exp = exp;
    t->next = NULL;
    tail->next = t;
    return t;
}

// 從陣列 c[0..m-1]、e[0..m-1] 建一個多項式。
Term *build(const int *c, const int *e, int m)
{
    Term head = {0, 0, NULL};       // 頭節點：只用它的 next
    Term *t = &head;
    for (int i = 0; i < m && t != NULL; i++)
        t = attach(t, c[i], e[i]);
    return head.next;
}

// 回傳 a + b，a、b 不變。*cmp 記錄比較了幾次次方。
Term *poly_add(const Term *a, const Term *b, int *cmp)
{
    Term head = {0, 0, NULL};
    Term *t = &head;
    *cmp = 0;
    while (a != NULL && b != NULL && t != NULL) {
        (*cmp)++;
        if (a->exp > b->exp) {
            t = attach(t, a->coef, a->exp);
            a = a->next;
        } else if (a->exp < b->exp) {
            t = attach(t, b->coef, b->exp);
            b = b->next;
        } else {
            if (a->coef + b->coef != 0)     // 係數相消就不留這一項
                t = attach(t, a->coef + b->coef, a->exp);
            a = a->next;
            b = b->next;
        }
    }
    for (; a != NULL && t != NULL; a = a->next)     // 剩下的照抄
        t = attach(t, a->coef, a->exp);
    for (; b != NULL && t != NULL; b = b->next)
        t = attach(t, b->coef, b->exp);
    return head.next;
}

// 每一項印成 (係數,次方)，然後釋放整條串列。
void print_and_free(const char *name, Term *p)
{
    printf("%s:", name);
    while (p != NULL) {
        Term *next = p->next;
        printf(" (%d,%d)", p->coef, p->exp);
        free(p);
        p = next;
    }
    printf("\n");
}

int main(void)
{
    int ca[] = {5, -3, 1, 8}, ea[] = {12, 7, 2, 0};
    int cb[] = {2, 3, -1, 4}, eb[] = {9, 7, 2, 1};
    Term *a = build(ca, ea, 4);
    Term *b = build(cb, eb, 4);
    int cmp;
    Term *c = poly_add(a, b, &cmp);
    print_and_free("A", a);
    print_and_free("B", b);
    print_and_free("A+B", c);
    printf("exponent comparisons = %d\n", cmp);
    return 0;
}
\end{example}
\iosample{}{A: (5,12) (-3,7) (1,2) (8,0) \\ B: (2,9) (3,7) (-1,2) (4,1) \\ A+B: (5,12) (2,9) (4,1) (8,0) \\ exponent comparisons = 5}

$A$、$B$ 的係數和次方寫在 \texttt{main} 的陣列裡，次方由大到小。 輸出把每一項印成
（係數, 次方），結果是 $A + B = 5x^{12} + 2x^{9} + 4x + 8$。

逐輪追蹤（「抄」是把那一項接到結果的最後面）：

\begin{center}
\begin{tabular}{cccl}
\toprule
比較 & $a$ 目前的項 & $b$ 目前的項 & 動作 \\
\midrule
1 & $5x^{12}$ & $2x^{9}$ & $12 > 9$，抄 $5x^{12}$ \\
2 & $-3x^{7}$ & $2x^{9}$ & $7 < 9$，抄 $2x^{9}$ \\
3 & $-3x^{7}$ & $3x^{7}$ & 係數和 $0$，不抄，兩邊都前進 \\
4 & $x^{2}$ & $-x^{2}$ & 係數和 $0$，不抄，兩邊都前進 \\
5 & $8$ & $4x$ & $0 < 1$，抄 $4x$；$b$ 走完了 \\
-- & $8$ & -- & 迴圈結束，抄 $8$ \\
\bottomrule
\end{tabular}
\end{center}

$m = n = 4$，比較了 $5$ 次，少於上限 $m + n - 1 = 7$，因為有兩次是一次消掉兩項。

多項式乘法可以拿 $A$ 的每一項去乘 $B$ 的每一項，得到 $mn$ 個乘積項，再一項一項加
進結果串列。 結果最多 $mn$ 項，每加一項要在結果裡找位置，最壞情況要
$O((mn)^2)$；改成每次把「$A$ 的一項乘上整個 $B$」得到的一條排好的串列，用
\proc{Poly-Add} 加進結果。 第 $i$ 次相加時結果最多 $(i-1)n$ 項，這次相加是 $O(in)$，
$m$ 次合計 $O(m^2 n)$。

\subsection{稀疏矩陣：正交串列}

稀疏矩陣如果常常要改（把某一格改成 $0$、或把 $0$ 改成非零），三元組陣列每次都要挪
動，這時可以改用\textbf{正交串列}（orthogonal list，也叫十字串列）。 這一小節只講
結構和成本，不寫程式。

每個非零元素是一個節點，有五個欄位：列號 $\mathit{row}$、行號 $\mathit{col}$、值
$\mathit{val}$，以及兩個指標 $\mathit{right}$（同一列的下一個非零元素）和
$\mathit{down}$（同一行的下一個非零元素）。 每一列有一個列頭節點，把這一列的非零元
素依行號串成一條環狀串列；每一行也有一個行頭節點，把這一行的非零元素依列號串成環。
一個節點同時在一條「橫的」串列和一條「直的」串列上，所以叫做正交。 下圖左邊是矩陣
$M$，右邊是它的正交串列。

\begin{center}
$M = \begin{pmatrix}
0 & 0 & 3 & 0 \\
8 & 0 & 0 & 0 \\
0 & 0 & 0 & 0 \\
0 & 6 & 0 & 2
\end{pmatrix}$\hspace{1.2cm}%
\begin{tikzpicture}[baseline=(current bounding box.center),
    hd/.style={draw=black!60, fill=black!8, minimum width=0.8cm, minimum height=0.6cm, font=\small},
    en/.style={draw=coreStroke, fill=coreFill, minimum width=1.25cm, minimum height=0.6cm, font=\scriptsize},
    lnk/.style={-{Stealth[length=1.6mm]}, thick, draw=black!70},
]
\foreach \j in {1,...,4} { \node[hd] (C\j) at (\j*1.9,0) {C\j}; }
\foreach \i in {1,...,4} { \node[hd] (R\i) at (0,-\i*0.95) {R\i}; }
\node[en] (e13) at (3*1.9,-1*0.95) {1,3 \quad 3};
\node[en] (e21) at (1*1.9,-2*0.95) {2,1 \quad 8};
\node[en] (e42) at (2*1.9,-4*0.95) {4,2 \quad 6};
\node[en] (e44) at (4*1.9,-4*0.95) {4,4 \quad 2};
% right 指標（同一列）
\draw[lnk] (R1) -- (e13);
\draw[lnk] (R2) -- (e21);
\draw[lnk] (R4) -- (e42);
\draw[lnk] (e42) -- (e44);
% down 指標（同一行）
\draw[lnk, draw=memStroke] (C1) -- (e21);
\draw[lnk, draw=memStroke] (C2) -- (e42);
\draw[lnk, draw=memStroke] (C3) -- (e13);
\draw[lnk, draw=memStroke] (C4) -- (e44);
\end{tikzpicture}
\diagcaption{$M$ 的正交串列。 黑色箭頭是 $\mathit{right}$，橘色箭頭是
$\mathit{down}$；每個節點寫「列, 行　值」。 每一列、每一行的最後一個節點指回自己的
頭節點，圖上沒有畫出這些指標。 第 3 列沒有非零元素，R3 的 $\mathit{right}$ 指向自
己。}
\label{fig:ll-ortho}
\end{center}

幾個成本：

\begin{itemize}
\item \textbf{空間}：$r \times c$ 的矩陣有 $t$ 個非零元素時，共 $t$ 個元素節點和
$r + c$ 個頭節點，$\Theta(t + r + c)$。 三元組陣列是 $\Theta(t)$，正交串列多了頭節
點，每個節點也比三元組多兩個指標。
\item \textbf{取 $M[i][j]$}：沿第 $i$ 列走，$O(\text{第 } i \text{ 列的非零個數})$。
\item \textbf{設定 $M[i][j]$}（新增或刪除一個節點）：在第 $i$ 列找到插入位置的前一個節
點，在第 $j$ 行也找一次，各改一個指標。 時間是兩條串列的長度，不用挪動其他元素。
\item \textbf{轉置}：只要把每個節點的 $\mathit{row}$ 和 $\mathit{col}$、
$\mathit{right}$ 和 $\mathit{down}$ 對調，列頭和行頭對調，$\Theta(t + r + c)$。
\end{itemize}

矩陣只讀不改時，三元組陣列比較省空間，也比較好寫；元素常常變動時，正交串列省下挪動
的成本。 第~21~章圖的鄰接串列，可以看成只保留「橫的」那一半的正交串列。

\section{可用空間串列}
\label{sec:ll-avail}

到目前為止，每個節點都用 \texttt{malloc} 要、用 \texttt{free} 還。 有些情況不想或
不能這樣做：沒有作業系統的小型裝置可能沒有 \texttt{malloc}；節點很小、配置很頻繁時，
每次都呼叫 \texttt{malloc} 的額外成本比節點本身還大；有些考題也會直接給兩個陣列
$\mathit{key}[\,]$、$\mathit{next}[\,]$，要你用陣列模擬串列。 這時可以一次準備好一池
節點，自己管理哪些在用、哪些空著。

空著的節點也串成一條串列，叫\textbf{可用空間串列}（available-space list，或 free
list），開頭記在 $\mathit{avail}$。 要節點就從這條串列的最前面拿一個，還節點就放回
最前面，兩者都是 $\Theta(1)$。 可用空間串列其實就是一個堆疊（第~6~章）。

\begin{pseudo}[label=alg:ll-getnode, unbreakable]{從可用空間串列拿一個節點}
\Function{Get-Node}{}
    \If{$\mathit{avail} = \NIL$}
        \State \Error{沒有可用的節點}
    \EndIf
    \State $x \gets \mathit{avail}$
    \State $\mathit{avail} \gets x.\mathit{next}$
    \State \Return $x$
\EndFunction
\end{pseudo}

\begin{pseudo}[label=alg:ll-retnode]{把節點 $x$ 還回去}
\Procedure{Ret-Node}{$x$}
    \State $x.\mathit{next} \gets \mathit{avail}$
    \State $\mathit{avail} \gets x$
\EndProcedure
\end{pseudo}

一開始所有節點都空著，把它們依序串起來就是最初的可用空間串列。 用陣列實作時，
「指標」改成陣列的編號，\NIL 用一個不會是編號的值（例如 $-1$）表示；
$x.\mathit{next}$ 寫成 $\mathit{next}[x]$。 這種寫法也叫\textbf{游標}（cursor）實
作。 一般的串列操作照寫，只是 \texttt{malloc} 換成 \proc{Get-Node}、\texttt{free}
換成 \proc{Ret-Node}；刪除時一樣要找到前一個節點，插入一樣先接後面再改前面。

下面用一個只有 $5$ 格的節點池追蹤一次。 一開始 $\mathit{avail} = 0$，
$\mathit{next}[\,] = \langle 1, 2, 3, 4, -1 \rangle$，串列是空的（$\mathit{head} =
-1$）。 「插入」都插在最前面；「$40@0$」表示 $40$ 放在第 $0$ 格。

\begin{center}
\begin{tabular}{llll}
\toprule
操作 & 用到的格子 & 串列（從 $\mathit{head}$ 走） & 可用空間串列 \\
\midrule
插入 40 & 拿第 0 格 & $40@0$ & 1, 2, 3, 4 \\
插入 25 & 拿第 1 格 & $25@1, 40@0$ & 2, 3, 4 \\
插入 63 & 拿第 2 格 & $63@2, 25@1, 40@0$ & 3, 4 \\
刪除 25 & 還第 1 格 & $63@2, 40@0$ & 1, 3, 4 \\
插入 18 & 拿第 1 格 & $18@1, 63@2, 40@0$ & 3, 4 \\
插入 77 & 拿第 3 格 & $77@3, 18@1, 63@2, 40@0$ & 4 \\
插入 52 & 拿第 4 格 & $52@4, 77@3, 18@1, 63@2, 40@0$ & （空） \\
插入 11 & 失敗 & 不變 & （空） \\
刪除 18 & 還第 1 格 & $52@4, 77@3, 63@2, 40@0$ & 1 \\
\bottomrule
\end{tabular}
\end{center}

\begin{center}
最後的陣列是\quad
\begin{tabular}{cccccc}
\toprule
編號 & 0 & 1 & 2 & 3 & 4 \\
\midrule
$\mathit{key}$ & 40 & 18 & 63 & 77 & 52 \\
$\mathit{next}$ & $-1$ & $-1$ & 0 & 2 & 3 \\
\bottomrule
\end{tabular}
\qquad $\mathit{head} = 4$，$\mathit{avail} = 1$。
\end{center}

幾件要看清楚的事：

\begin{itemize}
\item 刪除 $25$ 把第 $1$ 格還回去，它被放在可用空間串列的\textbf{最前面}，所以下一
次插入 $18$ 拿到的又是第 $1$ 格。 剛還的節點最先被重用，和堆疊一樣後進先出。
\item 五格都用完之後，插入 $11$ 失敗。 這是自己管理節點池的代價：容量要事先決定。
\item 第 $1$ 格的 $\mathit{key}$ 還是 $18$。 還回去的節點不會被清掉，只是不在串列上
了；判斷一格是不是在用，要看它在哪一條串列上，不能看裡面的值。
\end{itemize}

\begin{examnote}
用陣列模擬串列的追蹤題，每一步都寫出 $\mathit{head}$、$\mathit{avail}$ 和有變動的
$\mathit{next}[\,]$。 最後答案除了陣列內容，最好再從 $\mathit{head}$ 和
$\mathit{avail}$ 各走一次，寫出兩條串列，例如「串列：$4, 3, 2, 0$；可用：$1$」。
這兩條串列加起來要剛好是每一格各一次，可以拿來驗算。
\end{examnote}

\subsection{一次還一整條環狀串列}

要把一條有 $k$ 個節點的單向串列整條還回去，得先走到它的最後一個節點，才能把它接到
$\mathit{avail}$ 前面，$\Theta(k)$。 環狀串列記的正是最後一個節點，所以可以
$\Theta(1)$ 做完：

\enlargethispage{\baselineskip}% 讓整個虛擬碼框留在同一頁
\begin{pseudo}[label=alg:ll-ring-erase]{把整個環還給可用空間串列}
\Procedure{Ring-Erase}{$R$}
    \If{$R.\mathit{last} \ne \NIL$}
        \State $\mathit{first} \gets R.\mathit{last}.\mathit{next}$
        \State $R.\mathit{last}.\mathit{next} \gets \mathit{avail}$
        \State $\mathit{avail} \gets \mathit{first}$
        \State $R.\mathit{last} \gets \NIL$
    \EndIf
\EndProcedure
\end{pseudo}

第 4 行把環在 $\mathit{last}$ 後面剪開、接上原本的可用空間串列，第 5 行讓
$\mathit{avail}$ 從環的第一個節點開始。 做完之後，原本環上的 $k$ 個節點排在可用空間
串列的最前面，順序和在環上一樣。 和第~3~章的平攤分析放在一起看：建立一個 $k$ 個
節點的環要 $k$ 次 \proc{Get-Node}，拆掉只要 $\Theta(1)$，一連串操作的總時間只和配
置過的節點數有關。

\section{常見操作：反轉、合併、找中間}
\label{sec:ll-ops}

這一節的三個操作都只改 $\mathit{next}$，不配置新節點，額外空間 $\Theta(1)$。 它們
常常一起出現：第~14~章串列版的合併排序，就是「找中間切成兩半、各自排好、再合併」。
這一節用同一個情境示範：兩位助教各自登記了一部分同學的小考成績，名單都依學號排好；
要把它們合成一份，再切成前後兩半分到兩間教室，後半那間要從學號大的開始點名。

\subsection{反轉}

\textbf{反轉}把串列 $x_1 \to x_2 \to \cdots \to x_n$ 改成 $x_n \to \cdots \to x_1$。
做法是從頭一個一個拆下來，放到另一條串列 $d$ 的最前面。 拆之前要先記住下一個節點，
因為改掉 $h.\mathit{next}$ 之後就找不到它了。

\begin{pseudo}[label=alg:ll-reverse, unbreakable]{就地反轉}
\Function{List-Reverse}{$h$} \Comment{$h$ 指向第一個節點}
    \State $d \gets \NIL$ \Comment{已經反轉好的一段}
    \While{$h \ne \NIL$}
        \State $r \gets h.\mathit{next}$ \Comment{先記住還沒處理的部分}
        \State $h.\mathit{next} \gets d$
        \State $d \gets h$
        \State $h \gets r$
    \EndWhile
    \State \Return $d$
\EndFunction
\end{pseudo}

用 $112 \to 115 \to 117 \to 120$ 追蹤，每一列是 \textbf{while} 做完一輪之後的狀態：

\begin{center}
\begin{tabular}{clll}
\toprule
輪 & 這一輪搬的節點 & $d$ 那一段 & $h$ 那一段 \\
\midrule
開始 & -- & （空） & $112 \to 115 \to 117 \to 120$ \\
1 & $112$ & $112$ & $115 \to 117 \to 120$ \\
2 & $115$ & $115 \to 112$ & $117 \to 120$ \\
3 & $117$ & $117 \to 115 \to 112$ & $120$ \\
4 & $120$ & $120 \to 117 \to 115 \to 112$ & （空） \\
\bottomrule
\end{tabular}
\end{center}

\begin{theorem}
\label{thm:ll-reverse}
\proc{List-Reverse} 回傳反轉後的串列，時間 $\Theta(n)$，額外空間 $\Theta(1)$。
\end{theorem}

\begin{proof}
設原本的串列是 $x_1 \to \cdots \to x_n$。 不變量：第 $i$ 輪開始時（$1 \le i \le
n + 1$），$d$ 是 $x_{i-1} \to x_{i-2} \to \cdots \to x_1 \to \NIL$，$h$ 是 $x_i \to
x_{i+1} \to \cdots \to x_n \to \NIL$，而且 $x_i$ 以後的 $\mathit{next}$ 都沒有被改過
（$i = n + 1$ 時 $h = \NIL$）。

初始：$i = 1$，$d = \NIL$ 是空的一段，$h$ 是整條原本的串列，成立。

維持：第 $i$ 輪 $h = x_i \ne \NIL$。 第 4 行記下 $r = x_{i+1}$（$i = n$ 時是
\NIL）；第 5 行讓 $x_i$ 指向 $d$，$d$ 這一段變成 $x_i \to x_{i-1} \to \cdots \to
x_1$；第 6、7 行讓 $d = x_i$、$h = x_{i+1}$。 第 5 行只改了 $x_i$ 的
$\mathit{next}$，$x_{i+1}$ 以後沒有動，所以 $h$ 仍然是 $x_{i+1} \to \cdots \to x_n$。
下一輪開始時不變量對 $i + 1$ 成立。

結束：迴圈在 $h = \NIL$ 時停，也就是 $i = n + 1$。 這時 $d$ 是 $x_n \to \cdots \to
x_1 \to \NIL$，正是反轉後的串列。

迴圈跑 $n$ 輪，每輪 $\Theta(1)$；只用了 $d$、$r$ 兩個額外的指標。
\end{proof}

反轉也可以寫成遞迴：先反轉 $x_2$ 以後的部分，再把 $x_1$ 接到它的最後面。 遞迴的深
度是 $n$，要 $\Theta(n)$ 的堆疊空間，串列很長時可能用光呼叫堆疊（《C 程式設計基礎》
第~19~章 19.7 節）。 觀念題第 7 題會證明遞迴版的正確性。

\subsection{合併兩條排好的串列}

兩條依 $\mathit{key}$ 由小到大排好的串列，要合成一條排好的串列。 和多項式相加一
樣，兩邊各一個指標，每次把比較小的那個節點接到結果的最後面；不同的是這裡直接搬原本
的節點，不配置新的。

\begin{pseudo}[label=alg:ll-merge]{合併兩條排好的串列}
\Function{Merge-Lists}{$a, b$}
    \State $h \gets$ 頭節點；$t \gets h$
    \While{$a \ne \NIL$ \AND $b \ne \NIL$}
        \If{$a.\mathit{key} \le b.\mathit{key}$}
            \State $t.\mathit{next} \gets a$；$a \gets a.\mathit{next}$
        \Else
            \State $t.\mathit{next} \gets b$；$b \gets b.\mathit{next}$
        \EndIf
        \State $t \gets t.\mathit{next}$
    \EndWhile
    \If{$a \ne \NIL$}
        \State $t.\mathit{next} \gets a$ \Comment{剩下的一整段直接接上}
    \Else
        \State $t.\mathit{next} \gets b$
    \EndIf
    \State \Return $h.\mathit{next}$
\EndFunction
\end{pseudo}

兩點說明。 第一，第 4 行用 $\le$：$\mathit{key}$ 相同時先拿 $a$ 的，所以同樣
$\mathit{key}$ 的節點保持「$a$ 的在前、$b$ 的在後」。 這叫做\textbf{穩定}
（stable），第~14~章的合併排序靠這一點才是穩定的排序。 第二，某一邊用完時，另一邊
剩下的部分本來就排好了，而且都不比已經接上的節點小，整段接上只要 $\Theta(1)$，不用
一個一個搬。 這是串列比陣列方便的地方。

設兩條串列各有 $m$、$n$ 個節點。 每比較一次，就有一個節點接到結果上，而迴圈在其中
一邊用完時就停，所以比較次數介於 $\min(m, n)$ 和 $m + n - 1$ 之間（觀念題第 6 題會證
明兩個界都達得到）。 剩下的一段只花 $\Theta(1)$ 接上，所以總時間也介於
$\Theta(\min(m, n))$ 和 $\Theta(m + n)$ 之間，一般寫成 $O(m + n)$。 陣列版的合併
必須把剩下的元素一個一個抄過去，永遠是 $\Theta(m + n)$。

\subsection{找中間：快慢指標}

要把串列切成前後兩半，得先知道中間在哪裡。 直接的做法是先走一次數出長度 $n$，再走
$\lceil n/2 \rceil$ 步。 另一種做法只走一次：兩個指標一起從頭出發，\textbf{慢指
標} $s$ 每次走一步，\textbf{快指標} $f$ 每次走兩步。 $f$ 走到尾巴時，$s$ 剛好在中
間。

\begin{pseudo}[label=alg:ll-split]{切成前後兩半，前半有 $\lceil n/2 \rceil$ 個節點}
\Function{Split-Half}{$h$}
    \If{$h = \NIL$}
        \State \Return \NIL
    \EndIf
    \State $s \gets h$；$f \gets h.\mathit{next}$
    \While{$f \ne \NIL$ \AND $f.\mathit{next} \ne \NIL$}
        \State $s \gets s.\mathit{next}$
        \State $f \gets f.\mathit{next}.\mathit{next}$
    \EndWhile
    \State $b \gets s.\mathit{next}$；$s.\mathit{next} \gets \NIL$
    \State \Return $b$ \Comment{後半的開頭；前半仍從 $h$ 開始}
\EndFunction
\end{pseudo}

\begin{theorem}
\label{thm:ll-split}
設 $h$ 是 $x_1 \to \cdots \to x_n$，$n \ge 1$。 \proc{Split-Half} 的迴圈跑
$\lfloor (n-1)/2 \rfloor$ 輪，結束時 $s = x_{\lceil n/2 \rceil}$，所以前半是
$x_1, \ldots, x_{\lceil n/2 \rceil}$。 $n \ge 2$ 時 $1 \le \lceil n/2 \rceil < n$，
前後兩半都不是空的。
\end{theorem}

\begin{proof}
做完 $k$ 輪之後，$s = x_{1+k}$，$f = x_{2+2k}$（$2 + 2k > n$ 時 $f = \NIL$）。
這對 $k = 0$ 成立，每一輪 $s$ 前進一格、$f$ 前進兩格，所以一直成立。 迴圈繼續的條
件是 $x_{2+2k}$ 和 $x_{3+2k}$ 都存在，也就是 $2k + 3 \le n$。 所以迴圈跑的輪數
$K$ 是滿足 $2k + 3 \le n$ 的非負整數 $k$ 的個數，$K = \max(0, \lfloor (n-3)/2 \rfloor
+ 1) = \lfloor (n-1)/2 \rfloor$。

$n = 2m + 1$ 時 $K = m$，$s = x_{m+1} = x_{\lceil n/2 \rceil}$；$n = 2m$ 時
$K = m - 1$，$s = x_m = x_{n/2}$。 兩種情況都是 $x_{\lceil n/2 \rceil}$。
\end{proof}

$f$ 從 $h.\mathit{next}$ 出發而不是從 $h$ 出發，是為了偶數長度時讓 $s$ 停在兩個中
間節點的\textbf{前一個}。 $n = 2$ 時這樣切成 $1 + 1$；如果 $f$ 從 $h$ 出發，會切成
$2 + 0$，合併排序用它的話會一直切不開（觀念題第 11 題）。

快慢指標和「先數長度」的做法都走了大約 $1.5n$ 個指標：$f$ 走 $n$ 步、$s$ 走 $n/2$
步。 它的好處在不需要先知道長度。 同樣的技巧也能判斷串列有沒有環：有環的話快指標會在環
上追上慢指標，沒有環的話快指標會先走到 \NIL。 第~\ref{sec:ll-cycle}~節會把這件事講完整。

\subsection{例子：兩位助教的成績名單}

下面的程式把三個操作串起來：合併兩位助教的名單，切成兩半，再把後半反轉。 節點除了
學號還有成績；為了版面只印學號，成績跟著節點一起搬，不用另外處理。 \texttt{merge} 和 \texttt{split\_half} 是虛擬
碼~\ref{alg:ll-merge}、\ref{alg:ll-split} 的直接翻譯，\texttt{merge} 的頭節點是區域
變數。

\begin{example}[label=ex:ll-ops]{c}{(list\_ops.c) 合併、切半、反轉}
#include <stdio.h>

typedef struct Node {
    int id, score;          // 學號、小考成績
    struct Node *next;
} Node;

// 把陣列 a[0..n-1] 裡的節點依序串起來，回傳第一個。
Node *link_all(Node *a, int n)
{
    for (int i = 0; i + 1 < n; i++)
        a[i].next = &a[i + 1];
    a[n - 1].next = NULL;
    return &a[0];
}

// 合併兩條依學號排好的串列，不配置新節點。*cmp 記錄比較次數。
Node *merge(Node *a, Node *b, int *cmp)
{
    Node head = {0, 0, NULL};
    Node *t = &head;
    *cmp = 0;
    while (a != NULL && b != NULL) {
        (*cmp)++;
        if (a->id <= b->id) {
            t->next = a;
            a = a->next;
        } else {
            t->next = b;
            b = b->next;
        }
        t = t->next;
    }
    t->next = (a != NULL) ? a : b;  // 剩下的一整段直接接上
    return head.next;
}

// 切成前後兩半，前半有 ceil(n/2) 個節點。回傳後半的開頭。
Node *split_half(Node *h)
{
    if (h == NULL)
        return NULL;
    Node *slow = h, *fast = h->next;
    while (fast != NULL && fast->next != NULL) {
        slow = slow->next;          // 走一步
        fast = fast->next->next;    // 走兩步
    }
    Node *back = slow->next;
    slow->next = NULL;              // slow 是前半的最後一個
    return back;
}

// 就地反轉，回傳新的開頭。
Node *reverse(Node *h)
{
    Node *done = NULL;              // 已經反轉好的一段
    while (h != NULL) {
        Node *rest = h->next;       // 先記住還沒處理的部分
        h->next = done;
        done = h;
        h = rest;
    }
    return done;
}

void print_list(const char *name, const Node *p)   // 只印學號
{
    printf("%s:", name);
    for (; p != NULL; p = p->next)
        printf(" %d", p->id);
    printf("\n");
}

int main(void)
{
    // 節點直接放在陣列裡；三個操作都只改 next，不需要 malloc。
    Node ta1[] = {{103, 72, NULL}, {108, 90, NULL}, {112, 65, NULL},
                  {120, 88, NULL}};
    Node ta2[] = {{101, 80, NULL}, {109, 77, NULL}, {110, 93, NULL},
                  {115, 58, NULL}, {117, 84, NULL}};
    Node *a = link_all(ta1, 4), *b = link_all(ta2, 5);
    print_list("TA1", a);
    print_list("TA2", b);
    int cmp;
    Node *all = merge(a, b, &cmp);
    print_list("merged", all);
    printf("comparisons = %d\n", cmp);
    Node *back = split_half(all);
    print_list("room A", all);
    print_list("room B", back);
    print_list("room B reversed", reverse(back));
    return 0;
}
\end{example}
\iosample{}{TA1: 103 108 112 120 \\ TA2: 101 109 110 115 117 \\ merged: 101 103 108 109 110 112 115 117 120 \\ comparisons = 8 \\ room A: 101 103 108 109 110 \\ room B: 112 115 117 120 \\ room B reversed: 120 117 115 112}

$m + n = 9$，比較了 $8 = m + n - 1$ 次，是最多的情況：最大的 $120$ 在助教一的名單，
第二大的 $117$ 在助教二的名單，一直到只剩 $120$ 才停。 合併後有 $9$ 個節點，\texttt{split\_half} 讓前半
有 $\lceil 9/2 \rceil = 5$ 個。 節點都放在 \texttt{main} 的兩個陣列裡，三個操作
只改 \texttt{next}，沒有配置也沒有釋放任何節點。

\section{判斷串列有沒有環}
\label{sec:ll-cycle}

前面的串列都假設最後一個節點的 $\mathit{next}$ 是 \NIL。 程式有錯時這件事不一定成
立：觀念題第 7 題 (b) 拿掉一行之後，$1$ 和 $2$ 互相指著；兩條串列接錯地方，也會讓某
個節點指回前面。 這時「走到 \NIL 為止」的迴圈永遠停不下來。 這一節的問題是：只給
$\mathit{head}$，怎麼判斷從它出發會不會走進一個環？ 有環的話，環從哪一個節點開始、
有幾個節點？ 而且希望只用 $\Theta(1)$ 的額外空間，也不改動串列。

\subsection{$\rho$ 形的串列}

從 $h$ 出發一直沿著 $\mathit{next}$ 走，依序經過的節點記成
\[
x_0 = h, \qquad x_{i+1} = x_i.\mathit{next}.
\]
沒有環時，有 $n$ 個節點的串列走 $n$ 步就到 \NIL。 有環時永遠走不完，但節點只有有限
個，總會走到一個之前走過的節點。 令 $\mu$ 是最小的足標，使得 $x_\mu$ 之後還會再出現；
再令 $\lambda$ 是最小的正整數，使得 $x_{\mu+\lambda} = x_\mu$。 $x_\mu$ 叫\textbf{環的
起點}，$\lambda$ 叫\textbf{環長}；$x_0, \ldots, x_{\mu-1}$ 是環前面的\textbf{尾巴}，
$\mu$ 是尾巴的長度（整條串列就是一個環時 $\mu = 0$）。 畫出來像希臘字母 $\rho$：先
一段直線，再一個圈。 下圖的 $\mu = 3$、$\lambda = 5$。

\begin{center}
\begin{tikzpicture}
\node[lname] (hd) at (-1.15,0) {head};
\foreach \k/\v in {0/12, 1/25, 2/31, 3/47, 4/50, 5/66, 6/73, 7/88} {
    \lnode{c\k}{\k*1.75,0}{\v}
    \node[font=\scriptsize, text=black!60] at (\k*1.75+0.2,0.55) {$x_{\k}$};
}
\draw[lnk] (hd) -- (c0);
\foreach \k [evaluate=\k as \j using int(\k+1)] in {0,...,6} {
    \llink{c\k}{-- (c\j)}
}
\fill[black!70] (c7p.center) circle (1.1pt);
\draw[lnk, rounded corners=6pt] (c7p.center) -- ++(0.5,0) -- ++(0,-0.75) -| (c3.south);
\draw[decorate, decoration={brace, amplitude=4pt}, black!60]
    (-0.45,0.85) -- node[above=3pt, font=\small] {尾巴 $\mu = 3$} (3.95,0.85);
\draw[decorate, decoration={brace, amplitude=4pt}, black!60]
    (4.8,0.85) -- node[above=3pt, font=\small] {環 $\lambda = 5$} (12.9,0.85);
\end{tikzpicture}
\diagcaption{一條有環的串列。 $88$ 的 $\mathit{next}$ 指回 $47$，所以環的起點是
$x_3 = 47$。 從 head 走下去的節點序列是 $12, 25, 31, 47, 50, 66, 73, 88, 47, 50,
\ldots$。}
\label{fig:ll-rho}
\end{center}

\begin{lemma}
\label{lem:ll-rho}
設從 $h$ 出發會走進環。 則 $x_0, \ldots, x_{\mu+\lambda-1}$ 兩兩不同，從 $h$ 走得到
的節點恰好是這 $n = \mu + \lambda$ 個。 而且對 $0 \le i < j$，
\[
x_i = x_j \iff i \ge \mu \text{ 而且 } \lambda \text{ 整除 } j - i .
\]
\end{lemma}

\begin{proof}
先證三件事。

(a) 對所有 $t \ge \mu$，$x_{t+\lambda} = x_t$。 對 $t$ 歸納：$t = \mu$ 是 $\lambda$ 的
定義；若 $x_{t+\lambda} = x_t$，兩邊各走一步得到 $x_{t+1+\lambda} = x_{t+1}$。 重複
用 (a) 就得到：$i \ge \mu$ 而且 $\lambda$ 整除 $j - i$ 時，$x_j = x_i$。

(b) 若 $x_i = x_j$、$i < j$，則 $i \ge \mu$：$x_i$ 之後又出現了，由 $\mu$ 的最小性，
$i \ge \mu$。

(c) 若 $\mu \le i < j$ 而且 $x_i = x_j$，則 $\lambda$ 整除 $j - i$。 寫成 $j - i = q\lambda
+ r$，$0 \le r < \lambda$。 由 (a)，$x_j = x_{i+r}$，所以 $x_i = x_{i+r}$。 取 $t \ge 0$
使 $i + t = \mu + s\lambda$（$s$ 是某個非負整數），兩邊各再走 $t$ 步，由 (a) 得
$x_\mu = x_{i+t} = x_{i+r+t} = x_{\mu+r}$。 如果 $r > 0$，這和 $\lambda$ 是最小的矛
盾，所以 $r = 0$。

(a)、(b)、(c) 合起來就是要證的等價式。 若 $x_0, \ldots, x_{\mu+\lambda-1}$ 裡有兩個
相同的 $x_i = x_j$（$i < j$），則 $\lambda$ 整除 $j - i$，但 $0 < j - i < \lambda$，
矛盾，所以它們兩兩不同。 由 (a)，之後的節點都在 $x_\mu, \ldots, x_{\mu+\lambda-1}$
之中重複，所以走得到的節點恰好是這 $\mu + \lambda$ 個。
\end{proof}

引理說的是：尾巴上的節點各只出現一次；環上的節點每隔 $\lambda$ 步出現一次。

\subsection{Floyd 的快慢指標}

用一個集合記住走過的節點，就能發現第一個重複的節點，但要 $\Theta(n)$ 的空間。
\textbf{Floyd 的判環法}（也叫龜兔賽跑）只用兩個指標：慢指標 $s$ 一次走一步，快指標
$f$ 一次走兩步，兩個都從 $h$ 出發。 沒有環時，$f$ 先走到 \NIL；有環時，兩個都進了環
之後，$f$ 每一輪比 $s$ 多走一步，一定會在環上追上 $s$。 追上之後，再把一個指標放回
$h$，兩個都改成一次一步，下一次碰頭的地方就是環的起點。

\begin{pseudo}[label=alg:ll-floyd]{Floyd 判環：相遇點、環的起點、環長}
\Function{Floyd-Meet}{$h$} \Comment{有環時回傳相遇的節點，沒有環回傳 \NIL}
    \State $s \gets h$；$f \gets h$
    \While{$f \ne \NIL$ \AND $f.\mathit{next} \ne \NIL$}
        \State $s \gets s.\mathit{next}$
        \State $f \gets f.\mathit{next}.\mathit{next}$
        \If{$s = f$}
            \State \Return $s$
        \EndIf
    \EndWhile
    \State \Return \NIL
\EndFunction
\Statex
\Function{Cycle-Start}{$h$}
    \State $m \gets \Call{Floyd-Meet}{h}$
    \If{$m = \NIL$}
        \State \Return \NIL \Comment{沒有環}
    \EndIf
    \State $s \gets h$；$f \gets m$ \Comment{一個從頭、一個從相遇點}
    \While{$s \ne f$}
        \State $s \gets s.\mathit{next}$；$f \gets f.\mathit{next}$ \Comment{都走一步}
    \EndWhile
    \State \Return $s$
\EndFunction
\Statex
\Function{Cycle-Length}{$c$} \Comment{$c$ 是環上的節點}
    \State $\lambda \gets 1$；$t \gets c.\mathit{next}$
    \While{$t \ne c$}
        \State $t \gets t.\mathit{next}$；$\lambda \gets \lambda + 1$
    \EndWhile
    \State \Return $\lambda$
\EndFunction
\end{pseudo}

\proc{Floyd-Meet} 的迴圈條件和 \proc{Split-Half}（虛擬碼~\ref{alg:ll-split}）一樣：
$f$ 和 $f.\mathit{next}$ 都存在，才能走兩步。 差別在兩個指標都從 $h$ 出發，而且每走
完一輪就比一次。

\begin{theorem}
\label{thm:ll-floyd}
\begin{enumerate}
\item 沒有環時（$n$ 個節點），\proc{Floyd-Meet} 的迴圈跑 $\lfloor n/2 \rfloor$ 輪後
停下，回傳 \NIL。
\item 有環時，\proc{Floyd-Meet} 一定回傳一個節點。 迴圈跑的輪數 $k^*$ 是
$\lambda$ 的倍數裡、不小於 $\max(\mu, 1)$ 的最小那一個，而且 $k^* \le n = \mu +
\lambda$。
\item 有環時，\proc{Cycle-Start} 的第二個迴圈恰好跑 $\mu$ 輪，回傳環的起點
$x_\mu$。
\end{enumerate}
\end{theorem}

\begin{proof}
先看 \proc{Floyd-Meet}：做完 $k$ 輪之後 $s = x_k$、$f = x_{2k}$。 $k = 0$ 時成立，每
一輪 $s$ 前進一格、$f$ 前進兩格，所以一直成立。

(1) 沒有環時，$x_0, \ldots, x_{n-1}$ 兩兩不同，$x_n = \NIL$。 做完 $k$ 輪後還要再
做一輪，條件是 $x_{2k}$ 和 $x_{2k+1}$ 都存在，也就是 $2k + 2 \le n$。 滿足這個條件
的 $k \ge 0$ 有 $\lfloor n/2 \rfloor$ 個，所以迴圈跑 $\lfloor n/2 \rfloor$ 輪。 中
間每次比較的是 $x_k$ 和 $x_{2k}$（$k \ge 1$），足標不同，節點也不同，第 7 行不會執行，
最後回傳 \NIL。

(2) 有環時串列上沒有 \NIL，迴圈只會從第 7 行離開。 第 $k$ 輪（$k \ge 1$）結束時
$s = f$，就是 $x_k = x_{2k}$；由引理~\ref{lem:ll-rho}（取 $i = k$、$j = 2k$），這等
價於 $k \ge \mu$ 而且 $\lambda$ 整除 $k$。 所以迴圈在第一個滿足這兩個條件的 $k \ge 1$
停下，就是定理說的 $k^*$。 $\mu = 0$ 時 $k^* = \lambda = n$。 $\mu \ge 1$ 時，
$\mu, \mu + 1, \ldots, \mu + \lambda - 1$ 這 $\lambda$ 個連續整數裡恰好有一個是
$\lambda$ 的倍數，所以 $k^* \le \mu + \lambda - 1 < n$。

(3) 令 $m = x_{k^*}$。 \proc{Cycle-Start} 的第二個迴圈做完 $i$ 輪後 $s = x_i$、
$f = x_{k^* + i}$。 $i < \mu$ 時，$x_i$ 在尾巴上，由引理~\ref{lem:ll-rho} 它和其他
足標的節點都不同，所以 $s \ne f$，迴圈繼續。 $i = \mu$ 時，兩個足標 $\mu$ 和 $k^* +
\mu$ 都 $\ge \mu$，相差 $k^*$，是 $\lambda$ 的倍數，所以 $s = f = x_\mu$，迴圈停下。
\end{proof}

證明用的是代數，下面兩段把同樣的事用「追趕」的方式說一次，考試寫說明時比較順手。

\textbf{為什麼一定會相遇，不會跳過去。} 慢指標在第 $\mu$ 輪走到環的起點，這時快指標
已經在環上。 沿著環的方向看，快指標要再追 $g$ 步才碰到慢指標（$0 \le g < \lambda$）。
之後每一輪慢指標走 $1$ 步、快指標走 $2$ 步，$g$ 恰好減少 $1$，所以 $g$ 會一路減到
$0$，中間不會從 $1$ 直接變成負的。 也就是說，再過 $g$ 輪就相遇（$g = 0$ 表示第 $\mu$
輪就相遇了；$\mu = 0$ 時兩個指標一開始就在一起，要等繞完第一圈）。

\textbf{為什麼從頭和從相遇點一起走，會在起點碰頭。} 相遇點是從 $h$ 走 $k^*$ 步到的，
$k^*$ 是 $\lambda$ 的倍數。 從相遇點再走 $\mu$ 步，等於從 $h$ 走 $k^* + \mu$ 步：先走
$\mu$ 步到起點，再繞 $k^* / \lambda$ 整圈，還是停在起點。 從 $h$ 出發的那個指標走
$\mu$ 步也剛好到起點。 兩個指標同一輪到達起點；在那之前，從頭出發的指標還在尾巴上，
另一個在環上，不可能碰頭。 注意相遇點本身\textbf{不一定}是起點：下面的例子相遇在
$66$，起點是 $47$。

\textbf{成本。} \proc{Floyd-Meet} 最多 $\lfloor n/2 \rfloor$ 或 $k^* \le n$ 輪，
\proc{Cycle-Start} 的第二個迴圈 $\mu$ 輪，\proc{Cycle-Length} 從起點出發繞一圈
$\lambda$ 輪，每輪 $\Theta(1)$，總時間 $O(n)$。 反過來，沒有環時跑
$\lfloor n/2 \rfloor$ 輪，有環時 $k^* \ge \max(\mu, \lambda) \ge n/2$，所以也是
$\Omega(n)$，合起來 $\Theta(n)$。 額外空間只有幾個指標，$\Theta(1)$，而且完全沒有
改動串列。

\textbf{手算追蹤。} 用圖~\ref{fig:ll-rho} 的串列，第 $k$ 列是 \proc{Floyd-Meet} 做完
$k$ 輪後兩個指標指著的節點：

\begin{center}
\begin{tabular}{cccc}
\toprule
輪 $k$ & $s = x_k$ & $f = x_{2k}$ & 相等？ \\
\midrule
0 & $12$ & $12$ & （還沒開始比）\\
1 & $25$ & $31$ & 否 \\
2 & $31$ & $50$ & 否 \\
3 & $47$ & $73$ & 否 \\
4 & $50$ & $47$ & 否 \\
5 & $66$ & $66$ & 相遇 \\
\bottomrule
\end{tabular}
\end{center}

$\mu = 3$、$\lambda = 5$，$5$ 的倍數裡不小於 $3$ 的最小是 $5$，所以第 $5$ 輪相遇，
和定理一致。 第 $3$ 輪慢指標剛進環時，快指標在 $73$，沿著環再走 $2$ 步才碰到慢指標
（$73 \to 88 \to 47$），所以再過 $2$ 輪相遇。 接著 \proc{Cycle-Start}：

\begin{center}
\begin{tabular}{ccc}
\toprule
第二個迴圈的輪 & $s$（從 head 出發）& $f$（從相遇點出發）\\
\midrule
0 & $12$ & $66$ \\
1 & $25$ & $73$ \\
2 & $31$ & $88$ \\
3 & $47$ & $47$ \\
\bottomrule
\end{tabular}
\end{center}

跑了 $3 = \mu$ 輪，回傳 $47$。 最後從 $47$ 繞一圈 $47 \to 50 \to 66 \to 73 \to 88 \to
47$，環長 $5$。

\begin{examnote}
判環的題目，先寫出你用的版本：兩個指標從哪裡出發、每輪先走再比還是先比再走。 有些
課本讓 $f$ 從 $h.\mathit{next}$ 出發、先比再走，這時第一階段停下時，$s$ 走了 $k$ 步，
$f$ 走了 $2k + 1$ 步，比兩倍多一步，所以第二階段要讓 $f$ 從相遇點的\textbf{下一個}節點出發，才會在起點碰頭；
照本書的寫法從相遇點出發，兩個指標會永遠碰不到（環長是 $1$ 時除外）。 追蹤時一輪一列，寫出
$s$、$f$；最後分開寫「相遇點」「環的起點」「環長」三個答案。
\end{examnote}

\subsection{和其他做法比較}

\begin{center}
\begin{tabular}{>{\raggedright\arraybackslash}p{0.34\linewidth}lll}
\toprule
做法 & 時間 & 額外空間 & 改動串列 \\
\midrule
用集合記住走過的節點（雜湊表，第~12~章）& 期望 $\Theta(n)$ & $\Theta(n)$ & 不會 \\
每個節點多一個「走過了」旗標 & $\Theta(n)$ & 每個節點一個欄位 & 會，用完要清掉 \\
Floyd 快慢指標 & $\Theta(n)$ & $\Theta(1)$ & 不會 \\
\bottomrule
\end{tabular}
\end{center}

前兩種做法一走到重複的節點就知道它是環的起點，比較直接；Floyd 的方法用多走幾輪換
到常數空間，而且串列可以是唯讀的（例如別人的資料結構，或很多程式同時在讀）。
「先數長度」行不通：有環時數不完。

同一個方法不只能用在串列上。 只要有一個「下一個是誰」的規則，例如陣列
$\mathit{next}[\,]$ 表示的串列，或亂數產生器 $x_{i+1} = (a x_i + c) \bmod M$，從
任何起點一直往下走，最後都會進入一個環，Floyd 的方法可以在不記錄整個序列的情況下
找出尾巴長度和週期。

\subsection{C 程式}

下面的程式把圖~\ref{fig:ll-rho} 的八個節點放在陣列裡，改變最後一個節點的
\texttt{next}，做出四種串列：沒有環、$\rho$ 形、整條是一個環、最後一個節點指向自
己。 每一種印出相遇點、第一階段跑了幾輪、環的起點、尾巴長度和環長。

\begin{example}[label=ex:ll-cycle]{c}{(cycle.c) Floyd 判環}
#include <stdio.h>

typedef struct Node {
    int key;
    struct Node *next;
} Node;

// 第一階段：有環時回傳相遇的節點，沒有環回傳 NULL。*rounds 記錄跑了幾輪。
Node *floyd_meet(Node *h, int *rounds)
{
    Node *slow = h, *fast = h;
    *rounds = 0;
    while (fast != NULL && fast->next != NULL) {
        slow = slow->next;          // 走一步
        fast = fast->next->next;    // 走兩步
        (*rounds)++;
        if (slow == fast)
            return slow;
    }
    return NULL;
}

// 第二階段：一個從頭、一個從相遇點，每次各走一步，碰頭的地方是環的起點。
Node *cycle_start(Node *h, Node *meet, int *tail)
{
    *tail = 0;
    while (h != meet) {
        h = h->next;
        meet = meet->next;
        (*tail)++;
    }
    return h;
}

int cycle_length(const Node *c)     // c 在環上：繞一圈數節點
{
    int len = 1;
    for (const Node *t = c->next; t != c; t = t->next)
        len++;
    return len;
}

void report(const char *name, Node *h)
{
    int rounds, tail;
    Node *meet = floyd_meet(h, &rounds);
    printf("%s: ", name);
    if (meet == NULL) {
        printf("no cycle (%d rounds)\n", rounds);
        return;
    }
    Node *start = cycle_start(h, meet, &tail);
    printf("meet %d after %d rounds, ", meet->key, rounds);
    printf("start %d, tail %d, length %d\n", start->key, tail,
           cycle_length(start));
}

int main(void)
{
    int keys[] = {12, 25, 31, 47, 50, 66, 73, 88};
    Node a[8];
    for (int i = 0; i < 8; i++) {
        a[i].key = keys[i];
        a[i].next = (i + 1 < 8) ? &a[i + 1] : NULL;
    }
    report("chain", &a[0]);
    a[7].next = &a[3];      // 88 指回 47
    report("rho", &a[0]);
    a[7].next = &a[0];      // 88 指回 12：整條是一個環
    report("ring", &a[0]);
    a[7].next = &a[7];      // 88 指向自己
    report("self", &a[0]);
    return 0;
}
\end{example}
\iosample{}{chain: no cycle (4 rounds) \\ rho: meet 66 after 5 rounds, start 47, tail 3, length 5 \\ ring: meet 12 after 8 rounds, start 12, tail 0, length 8 \\ self: meet 88 after 7 rounds, start 88, tail 7, length 1}

四行都和定理~\ref{thm:ll-floyd} 對得上：沒有環的 $8$ 個節點跑 $\lfloor 8/2 \rfloor =
4$ 輪；整條是環時 $\mu = 0$，要等 $k^* = \lambda = 8$ 輪，兩個指標都回到 $12$ 才相
遇；$88$ 指向自己時 $\mu = 7$、$\lambda = 1$，第 $7$ 輪慢指標一走到 $88$ 就和早已停
在那裡打轉的快指標相遇。 第三種情況相遇點就是起點，第二種不是，所以第二階段不能省。

\section{本章常見錯誤}
\label{sec:ll-pitfalls}

\begin{center}
\begin{tabular}{p{0.24\linewidth}p{0.50\linewidth}p{0.16\linewidth}}
\toprule
錯誤 & 例子 & 見 \\
\midrule
改指標的順序錯 & 插入時先寫 \texttt{p->next = x} 再寫 \texttt{x->next = p->next}，
$x$ 指向自己，後半段遺失 & 第~\ref{sec:ll-singly}~節 \\
釋放之後還在用 & \texttt{free(p); p = p->next;}；刪除時先 \texttt{free} 才接前後 &
第~\ref{sec:ll-singly}~節 \\
改到參數的複本 & 函式收 \texttt{Node *head}，在裡面寫 \texttt{head = ...}，呼叫端的
\texttt{head} 沒變 & 第~\ref{sec:ll-singly}~節 \\
忘了第一個節點 & 刪除或插入在最前面時沒有更新 \texttt{head} &
第~\ref{sec:ll-singly}~節 \\
環狀串列用 \NIL 當停止條件 & \texttt{while (p != NULL)} 走環，永遠停不下來 &
第~\ref{sec:ll-variants}~節 \\
雙向串列只改一個方向 & 刪除時只改 \texttt{x->prev->next}，後一個節點的
\texttt{prev} 還指著 $x$ & 第~\ref{sec:ll-variants}~節 \\
係數相消的項留在結果裡 & 多項式相加得到「$0x^7$」這種節點 &
第~\ref{sec:ll-poly}~節 \\
以為還回去的節點被清空 & 用 $\mathit{key}$ 判斷一格有沒有在用；還回去之後還有別的
指標指著它 & 第~\ref{sec:ll-avail}~節 \\
反轉時沒先記住下一個 & 先改 \texttt{h->next = d}，再寫 \texttt{h = h->next}，
回到了 $d$ & 第~\ref{sec:ll-ops}~節 \\
快慢指標的起點選錯 & 偶數長度時切成「全部＋空的」，合併排序無窮遞迴 &
第~\ref{sec:ll-ops}~節 \\
把相遇點當成環的起點 & 判環時第一階段停下的節點直接當答案；$f$ 從 $h.\mathit{next}$
出發，第二階段卻從相遇點出發 & 第~\ref{sec:ll-cycle}~節 \\
\bottomrule
\end{tabular}
\end{center}

\section{觀念題}
\label{sec:ll-exercises}

題型和研究所考試、期中期末考的筆試相近，有選擇、計算、追蹤、證明和寫程式。 解答在
本章最後。

\begin{enumerate}
\item （選擇，可能不只一個答案）一條有 $n$ 個節點的單向串列，只記
$L.\mathit{head}$，沒有頭節點、也不記最後一個節點。 下列哪些操作在最壞情況下是
$\Theta(1)$？
\begin{enumerate}
\item 在最前面插入一個節點
\item 刪除最後一個節點
\item 已知指向節點 $p$ 的指標，在 $p$ 後面插入一個節點
\item 已知指向節點 $p$ 的指標，在 $p$ 前面插入一個節點，而且串列上原有的每個節點
仍然存著原本的資料
\item 求串列的長度
\end{enumerate}

\item （追蹤）對單向串列 $L$ 執行下面的程序。
\begin{pseudo}{觀念題 2}
\Procedure{Mystery}{$L$}
    \State $p \gets L.\mathit{head}$
    \While{$p \ne \NIL$ \AND $p.\mathit{next} \ne \NIL$}
        \State $q \gets p.\mathit{next}$
        \State $p.\mathit{next} \gets q.\mathit{next}$
        \State $q.\mathit{next} \gets L.\mathit{head}$
        \State $L.\mathit{head} \gets q$
        \State $p \gets p.\mathit{next}$
    \EndWhile
\EndProcedure
\end{pseudo}
\begin{enumerate}
\item $L$ 是 $1 \to 2 \to 3 \to 4 \to 5 \to 6 \to 7$。 寫出 \textbf{while} 每一輪之後
的串列。
\item 用一句話描述這個程序對一般的串列做了什麼，並寫出 $n$ 個節點時迴圈跑幾輪。
\end{enumerate}

\item （選擇，可能不只一個答案）在雙向串列裡把節點 $x$ 接在 $p$ 後面（$p$ 後面原本
有節點），要執行下面四個指定：
\[
\text{(1)}\ x.\mathit{next} \gets p.\mathit{next} \quad
\text{(2)}\ x.\mathit{prev} \gets p \quad
\text{(3)}\ p.\mathit{next}.\mathit{prev} \gets x \quad
\text{(4)}\ p.\mathit{next} \gets x
\]
下列哪些執行順序是正確的？
(A) 1, 2, 3, 4 \quad (B) 4, 1, 2, 3 \quad (C) 2, 1, 4, 3 \quad (D) 3, 1, 2, 4 \quad
(E) 1, 4, 2, 3

\item （追蹤）用陣列模擬串列，節點池有 $6$ 格。 目前的狀態如下，「--」表示那一格的
$\mathit{key}$ 沒有意義。 $\mathit{head} = 2$，$\mathit{avail} = 3$。
\begin{center}
\begin{tabular}{ccccccc}
\toprule
編號 & 0 & 1 & 2 & 3 & 4 & 5 \\
\midrule
$\mathit{key}$ & 31 & -- & 17 & -- & -- & 8 \\
$\mathit{next}$ & 5 & 4 & 0 & 1 & $-1$ & $-1$ \\
\bottomrule
\end{tabular}
\end{center}
依序執行：刪除 $31$、在最前面插入 $44$、在最前面插入 $66$、刪除 $8$。 插入用
\proc{Get-Node}，刪除時找到前一個節點接好後用 \proc{Ret-Node} 把那一格還回去。 寫出最後的
$\mathit{key}$、$\mathit{next}$、$\mathit{head}$、$\mathit{avail}$，以及從
$\mathit{head}$ 和 $\mathit{avail}$ 出發的兩條串列。

\item （追蹤）$3 \times 4$ 的矩陣 $N$ 只有四個非零元素：$N[1][2] = 5$、
$N[2][4] = -2$、$N[3][1] = 9$、$N[3][3] = 4$。 用第~\ref{sec:ll-poly}~節的正交串列
存（每列、每行各一個頭節點，串列都是環狀的）。
\begin{enumerate}
\item 一共有幾個節點？
\item 從第 $3$ 列的頭節點沿 $\mathit{right}$ 走一圈，依序經過哪些節點？
\item 把 $N[3][3]$ 改成 $0$（刪掉那個節點），要改哪兩個指標，各改成指向誰？
\item 把 $N[2][2]$ 設成 $7$（新增節點 $y$），在第 $2$ 列和第 $2$ 行各要改哪些指標？
\end{enumerate}

\item （證明）兩條排好的串列各有 $m$、$n$ 個節點（$m, n \ge 1$），用
\proc{Merge-Lists}（虛擬碼~\ref{alg:ll-merge}）合併。 證明 $\mathit{key}$ 的比較次
數 $C$ 滿足 $\min(m, n) \le C \le m + n - 1$，並對任意的 $m$、$n$ 各給一組輸入，分
別讓 $C = \min(m, n)$ 和 $C = m + n - 1$。

\item （證明）下面是遞迴版的反轉，回傳反轉後的第一個節點。
\begin{pseudo}{觀念題 7}
\Function{Rev-Rec}{$h$}
    \If{$h = \NIL$ \OR $h.\mathit{next} = \NIL$}
        \State \Return $h$
    \EndIf
    \State $r \gets \Call{Rev-Rec}{h.\mathit{next}}$
    \State $h.\mathit{next}.\mathit{next} \gets h$
    \State $h.\mathit{next} \gets \NIL$
    \State \Return $r$
\EndFunction
\end{pseudo}
\begin{enumerate}
\item 對節點個數 $n$ 用數學歸納法證明它是對的。
\item 如果拿掉第 6 行，對 $1 \to 2 \to 3$ 執行之後會發生什麼事？
\item 它的額外空間是多少？
\end{enumerate}

\item （寫程式）用第~\ref{sec:ll-singly}~節的 \texttt{Node} 和 \texttt{Node **}
寫法，寫出
\begin{center}
\texttt{void remove\_all(Node **head, int k)}
\end{center}
刪除串列裡所有 \texttt{id} 等於 \texttt{k} 的節點並釋放它們。 要求 $O(n)$ 時間，
不能另外寫「刪第一個節點」的特殊情況。

\item （找錯）下面兩個函式各有什麼錯？ 說明怎麼改。
\begin{code}{c}{觀念題 9}
// (a) 釋放整條串列
void free_all(Node *head)
{
    for (Node *p = head; p != NULL; p = p->next)
        free(p);
}

// (b) 刪除第一個節點
void drop_first(Node *head)
{
    Node *p = head;
    head = head->next;
    free(p);
}
\end{code}

\item （設計）值日環用單向環狀串列表示，只記最後一個節點 $R.\mathit{last}$。 現在
要把它拆成兩組：從今天值日的人一直到節點 $p$（含 $p$）為第一組，$p$ 之後到
$R.\mathit{last}$ 為第二組，兩組各自成為一個環，組內的輪替順序不變。 已知 $p$ 在環上
而且 $p \ne R.\mathit{last}$。 寫出 $\Theta(1)$ 的程序 \proc{Ring-Split}$(R, p)$，
回傳兩個環各自的最後一個節點。 如果環改成只記第一個節點，還能 $\Theta(1)$ 做到嗎？

\item （計算）把 \proc{Split-Half}（虛擬碼~\ref{alg:ll-split}）第 4 行的
$f \gets h.\mathit{next}$ 改成 $f \gets h$，其他不變。
\begin{enumerate}
\item $n = 1, 2, \ldots, 6$ 時，迴圈結束時 $s$ 是第幾個節點？ 寫出一般的公式。
\item 用改過的版本切半，再拿去做合併排序（切成兩半、各自遞迴排序、再合併），$n = 2$
時會發生什麼事？
\end{enumerate}

\item （計算）一條串列有 $10$ 個節點 $x_0, x_1, \ldots, x_9$，$x_0$ 是 head，
$x_i.\mathit{next} = x_{i+1}$（$i \le 8$），$x_9.\mathit{next} = x_6$。
\begin{enumerate}
\item 寫出這條串列的 $\mu$、$\lambda$。 \proc{Floyd-Meet}（虛擬碼~\ref{alg:ll-floyd}）
跑幾輪、在哪個節點相遇？ 列出每一輪之後的 $s$、$f$。
\item \proc{Cycle-Start} 的第二個迴圈跑幾輪？
\item 把 $x_9.\mathit{next}$ 改成 $x_0$，(a)、(b) 的答案變成什麼？
\item 節點數固定是 $10$，$x_9.\mathit{next}$ 可以指向 $x_0, \ldots, x_9$ 任何一個。
哪一種讓 \proc{Floyd-Meet} 跑最多輪、哪一種最少？ 各幾輪？
\end{enumerate}

\item （設計）兩條沒有環的單向串列，開頭分別是 $a$、$b$，長度 $m$、$n$。 它們可能從
某個節點 $X$ 開始共用後半段（$X$ 以後的節點兩條串列都有，兩條分支在 $X$ 合成一條，像英文字母 Y）。
寫出一個 $O(m + n)$ 時間、$O(1)$ 額外空間的程序，回傳第一個共用的節點 $X$，沒有共
用的節點時回傳 \NIL。 可以暫時改動指標，但結束前要改回來。 （提示：用第~\ref{sec:ll-cycle}~節的
\proc{Cycle-Start}。）
\end{enumerate}

\section{歷屆試題}
\label{sec:ll-pastexam}

以下是研究所入學考試的原題，題目照原卷（語言、記號、配分照錄）。 詳解在本章最後。

\begin{pastexam}[label={px:ll-1}]{臺大資工所 110 學年度 第 VII 題}
（原卷為單選題，每題答對 5 分，答錯倒扣 2.5 分。）

There is a linked-list with $n$ elements, which are $x_1,x_2,\ldots,x_n$. Let $f(n)$ be
the time complexity of inserting a new element $x_{n+1}$ to the tail of the list, and
$g(n)$ be that of deleting the last element $x_n$. If there is only one pointer pointed
to the head of the list, we know to access the list is time-consuming, and functions
$f(n)$ and $g(n)$ are denoted as $f_1(n)$ and $g_1(n)$ in this case. To improve it, we
may also add another pointer pointed to the tail of it, so the functions $f(n)$ and
$g(n)$ become $f_2(n)$ and $g_2(n)$. Which of the following option is false if $n$
becomes large. \underline{\makebox[5em]{(13)}}

(A) $\frac{f_1(n)}{g_1(n)}=1$ \qquad (B) $\frac{f_2(n)}{g_2(n)}=n$ \qquad
(C) $\frac{f_1(n)}{f_2(n)}=n$ \qquad (D) $\frac{g_1(n)}{g_2(n)}=1$
\end{pastexam}

\begin{pastexam}[label={px:ll-2}]{成大資訊聯招 114 學年度 第 2 題}
{[3\%]} Given a node structure as shown below:
\begin{minted}[fontsize=\small]{c}
typedef struct Node {
    int data;
    struct Node* next;
} Node;
\end{minted}
What will the following function do when used on a singly linked list?
\begin{minted}[fontsize=\small]{c}
void unknownFunction(Node* node) {
    if (node == NULL || node->next == NULL) {
        return;
    }
    node->data = node->next->data;
    Node* temp = node->next;
    node->next = node->next->next;
    free(temp);
}
\end{minted}
\begin{enumerate}[label=(\Alph*)]
\item The function will delete the specified node, including the last node in the list.
\item The function will delete the specified node only if it is a middle node or any node
except the last one.
\item The function will produce a segmentation fault when applied to the last node in
the list because it attempts to access node->next->data.
\item The function will remove the node but corrupt the list structure if the deleted
node has duplicate values in the list.
\item None of the above is correct.
\end{enumerate}
\end{pastexam}

\begin{pastexam}[label={px:ll-3}]{中央資工所 113 學年度 第 8 題}
（原卷為複選題，每題 5 分，全部答對才給分，答錯倒扣一分。）

Consider the following function that takes reference to head of a Doubly Linked List as
parameter. Assume that a node of doubly linked list has previous pointer as prev and next
pointer as next.
\begin{minted}[fontsize=\small]{c}
void fun(struct node **head_ref)
{
    struct node *temp = NULL;
    struct node *current = *head_ref;

    while (current != NULL)
    {
        temp = current->prev;
        current->prev = current->next;
        current->next = temp;
        current = current->prev;
    }

    if(temp != NULL )
        *head_ref = temp->prev;
}
\end{minted}
Assume that reference of head of a doubly linked list L is passed to above function. What
is the resulting linked list after the function call?
\begin{enumerate}[label=(\Alph*)]
\item \texttt{2 <--> 1 <--> 4 <--> 3 <--> 6 <-->5} when L is
\texttt{2 <--> 1 <--> 4 <--> 3 <--> 6 <-->5}
\item \texttt{5 <--> 4 <--> 3 <--> 2 <--> 1 <-->6} when L is
\texttt{5 <-->4 <--> 3 <-->2 <--> 6<-->1}
\item \texttt{6 <--> 5 <--> 4 <--> 3 <--> 2 <--> 1} when L is
\texttt{1 <--> 2 <--> 3 <--> 4 <--> 5 <-->6}
\item \texttt{6 <--> 5 <--> 4 <--> 3 <--> 1 <--> 2} when L is
\texttt{1 <--> 2 <--> 3 <--> 4 <--> 5 <-->6}
\end{enumerate}
\end{pastexam}

\begin{pastexam}[label={px:ll-4}]{臺科大資工所 112 學年度 第 2 題}
(10\%) Consider the function \textbf{func} as shown below, where the two lists list1 and
list2 are both ordered in ascending order.
\begin{minted}[fontsize=\small]{c}
struct item
{
   int value;
   struct item *next;
};

int func(struct item *list1, struct item *list2)
{
   struct item *ptr1, *ptr2;
   int count = 0;
   ptr1 = list1;
   ptr2 = list2;
   while ((ptr1 != NULL)&&(ptr2 != NULL))
   {
      if (ptr1->value < ptr2->value)
         ptr1 = ptr1 ->next;
      else if (ptr1->value > ptr2->value)
         ptr2 = ptr2->next;
      else
      {
         ptr1 = ptr1->next;
         ptr2 = ptr2->next;
         count++;
      }
   }
   return count;
}
\end{minted}
If $m$ and $n$ are the number of items in list1 and list2 respectively, what is the
worst-case time complexity of func(list1, list2)?
\end{pastexam}

\begin{pastexam}[label={px:ll-5}]{政大資科所 108 學年度 第 9 題}
(15\%) Design an algorithm for the following task.

\quad Input: a linked list.\\
\quad Task: determine whether the linked list contains a cycle.\\
\quad For example, the following linked list contains a cycle.
\begin{center}
\begin{tikzpicture}[>={Stealth[length=2.2mm]},thick,b/.style={draw,minimum size=6mm}]
\foreach \i in {1,...,5} \node[b] (n\i) at (1.2*\i,0) {};
\foreach \i/\j in {1/2,2/3,3/4,4/5} \draw[->] (n\i) -- (n\j);
\draw[->] (n5.south) to[out=-130,in=-50] (n1.south);
\end{tikzpicture}
\end{center}
\end{pastexam}

\section{實作練習}
\label{sec:ll-practice}

下面是 Aki-Judge 上和本章相關的題目。 大部分題目只要交函式和標頭檔，不寫
\texttt{main}；評測程式會檢查你有沒有真的用節點和指標，有幾題還規定不能配置新節點，
只能重接原本的 \texttt{next}。

\begin{center}
\begin{tabular}{llc}
\toprule
小節 & 題目 & 難度 \\
\midrule
單向串列 & \akiref{build-and-print-a-list}{建立並列印串列} & 暖身 \\
單向串列 & \akiref{insert-and-delete-at-a-position}{依位置插入與刪除} & 基本 \\
雙向串列 & \akiref{doubly-linked-list}{雙向鏈結串列} & 基本 \\
雙向串列 & \akiref{move-to-front-in-a-doubly-linked-list}{雙向鏈結串列的移到最前端} & 基本 \\
環狀串列 & \akiref{cutting-a-bead-ring-by-color}{依顏色拆開珠串環} & 進階 \\
用串列表示資料 & \akiref{bead-runs-pack-and-locate}{串珠壓縮與定位} & 基本 \\
可用空間 & \akiref{best-fit-memory-pool}{最佳適配記憶體池} & 進階 \\
反轉 & \akiref{reverse-a-linked-list-in-place-iterative}{就地反轉鏈結串列（迴圈）} & 暖身 \\
反轉 & \akiref{reverse-a-linked-list-in-place-recursive}{遞迴反轉鏈結串列} & 基本 \\
合併 & \akiref{merge-two-sorted-linked-lists}{合併兩條已排序的鏈結串列} & 基本 \\
合併 & \akiref{merge-chains-drop-shared-values}{合併兩條鏈並剔除共同值} & 進階 \\
重接指標 & \akiref{riffle-two-chains-by-counts}{依張數交錯洗牌兩條鏈} & 進階 \\
重接指標 & \akiref{follow-the-card}{照卡片指示抽牌} & 進階 \\
重接指標 & \akiref{threading-replies-under-posts}{把回覆接到貼文底下} & 進階 \\
綜合 & \akiref{clean-up-a-sign-up-list}{整理報名名單} & 進階 \\
\bottomrule
\end{tabular}
\end{center}

「依張數交錯洗牌兩條鏈」和「合併兩條鏈並剔除共同值」的鏈是用指標陣列表示的，
\texttt{next} 欄位換成了另一個陣列，但想法和本章一樣。 「整理報名名單」的資料量
大，要在串列上做 $O(n \lg n)$ 的排序，可以先讀第~14~章串列版的合併排序再回來寫。

\section{讀完你應該會}
\label{sec:ll-checklist}

\begin{itemize}
\item 能畫出單向、環狀、雙向串列插入和刪除時每一步的指標變化，並說出各種操作在不同
串列上的成本，以及和陣列比較的取捨。
\item 能用頭節點或 \texttt{Node **} 寫出沒有特殊情況的插入與刪除，並且釋放所有節
點，通過 AddressSanitizer 的檢查。
\item 能用串列做多項式相加並說明比較次數的上限，能畫出稀疏矩陣的正交串列並說出它
的空間和操作成本。
\item 能用陣列模擬串列和可用空間串列，逐步追蹤 $\mathit{head}$、$\mathit{avail}$ 和
$\mathit{next}[\,]$。
\item 能寫出反轉、合併、快慢指標切半，並用迴圈不變量證明反轉的正確性、算出切半時
前半的長度。
\item 能用 Floyd 的快慢指標判斷串列有沒有環、找出環的起點和環長，說明快慢指標為什麼一
定相遇、為什麼從頭和從相遇點一起走會在起點碰頭，並算出每個階段跑幾輪。
\end{itemize}

\textbf{延伸閱讀}（書末會整理完整的參考資料）

\begin{itemize}
\item Thomas H. Cormen 等，\textit{Introduction to Algorithms}，第 4 版，第 10.2 節：
雙向串列的搜尋、插入、刪除，以及用哨兵（sentinel）去掉邊界情況，對應本章的頭節點。
\item Donald E. Knuth，\textit{The Art of Computer Programming}，第 1 冊，第 2.2.3 到
2.2.6 節：鏈結配置、環狀串列、雙向串列和正交串列的原始講法，可用空間串列也在這
裡。 寫法比較老，適合想看這些結構來源的人。
\item Donald E. Knuth，\textit{The Art of Computer Programming}，第 2 冊，第 3.1 節習題
6：Floyd 判環法最早的出處，放在「亂數產生器的週期」的脈絡下，正好是第~\ref{sec:ll-cycle}~節最後說的推廣。
\item Steven S. Skiena，\textit{The Algorithm Design Manual}，第 3 版，第 3.1 節：
連續存放和鏈結存放的比較，偏重實務上的取捨。
\item 《C 程式設計基礎》第~19~章：串列的不變式、頭端插入、指標的指標、串接和有序
合併，C 語法的部分寫得比本章詳細。
\end{itemize}

\addcontentsline{toc}{section}{觀念題解答}
\section*{觀念題解答}

\begin{enumerate}
\item (a)、(c)。 (a) 和 (c) 都只改兩個指標。 (b) 要把倒數第二個節點的
$\mathit{next}$ 設成 \NIL，只能從頭走過去找，$\Theta(n)$。 (d) 要改 $p$ 前一個節點
的 $\mathit{next}$，一樣要從頭找。 （如果允許搬資料，可以在 $p$ 後面插入新節點，
再把 $p$ 和新節點的資料對調，$\Theta(1)$；但這樣 $p$ 這個節點存的資料變了，其他指
著 $p$ 的指標會看到不同的東西，題目已經排除這種做法。） (e) 要數過每個節點，
$\Theta(n)$。

\item
\begin{enumerate}
\item 每一輪把 $p$ 後面那個節點 $q$ 拿下來放到最前面，$p$ 再往前一格。
\begin{center}
\begin{tabular}{clc}
\toprule
輪 & 這一輪之後的串列 & 下一輪的 $p$ \\
\midrule
1 & $2 \to 1 \to 3 \to 4 \to 5 \to 6 \to 7$ & 3 \\
2 & $4 \to 2 \to 1 \to 3 \to 5 \to 6 \to 7$ & 5 \\
3 & $6 \to 4 \to 2 \to 1 \to 3 \to 5 \to 7$ & 7 \\
\bottomrule
\end{tabular}
\end{center}
第 3 輪之後 $p = 7$，$p.\mathit{next} = \NIL$，迴圈停。
\item 原本在偶數位置的節點，依相反的順序搬到最前面；奇數位置的節點保持原本的順序
接在後面。 $p$ 依序是原本的第 $1, 3, 5, \ldots$ 個節點，每輪處理一個偶數位置，迴
圈跑 $\lfloor n/2 \rfloor$ 輪，總時間 $\Theta(n)$。
\end{enumerate}

\item (A)、(D)。 (3) 和 (1) 都讀了 $p.\mathit{next}$，要的是原本的後一個節點，所以
都要在 (4) 之前；(2) 放在哪裡都可以。 (B)：(4) 先做，(1) 讓 $x.\mathit{next} = x$。
(C)：(3) 在 (4) 之後，$p.\mathit{next}$ 已經是 $x$，結果是 $x.\mathit{prev} = x$，
原本後一個節點的 $\mathit{prev}$ 沒有改。 (E)：同 (C) 的問題。

\item 四步之後：
\begin{center}
\begin{tabular}{ccccccc}
\toprule
編號 & 0 & 1 & 2 & 3 & 4 & 5 \\
\midrule
$\mathit{key}$ & 44 & -- & 17 & 66 & -- & 8 \\
$\mathit{next}$ & 2 & 4 & $-1$ & 0 & $-1$ & 1 \\
\bottomrule
\end{tabular}
\end{center}
$\mathit{head} = 3$，$\mathit{avail} = 5$。 串列：$66, 44, 17$（格 $3, 0, 2$）；可用：
格 $5, 1, 4$。 過程：刪除 $31$（格 $0$）時 $\mathit{next}[2] \gets 5$，格 $0$ 還回去，
$\mathit{avail} = 0 \to 3 \to 1 \to 4$；插入 $44$ 拿到格 $0$；插入 $66$ 拿到格 $3$；
刪除 $8$（格 $5$）時 $\mathit{next}[2] \gets -1$，格 $5$ 放到可用空間串列的最前面。
格 $5$ 的 $\mathit{key}$ 還是 $8$，但它已經不在串列上。

\item
\begin{enumerate}
\item $4$ 個元素節點，加上 $3$ 個列頭、$4$ 個行頭，共 $11$ 個。
\item 列頭 R3 $\to (3,1)\,9 \to (3,3)\,4 \to$ R3。
\item $(3,3)$ 是第 $3$ 列的最後一個，也是第 $3$ 行唯一的元素。 第 $3$ 列：
$(3,1).\mathit{right} \gets \text{R3}$。 第 $3$ 行：$\text{C3}.\mathit{down} \gets
\text{C3}$（第 $3$ 行變成空的）。
\item 第 $2$ 列原本只有 $(2,4)$，$y$ 的行號 $2 < 4$，要放在它前面：
$y.\mathit{right} \gets (2,4)$，$\text{R2}.\mathit{right} \gets y$。 第 $2$ 行原本只
有 $(1,2)$，$y$ 的列號 $2 > 1$，放在它後面：$y.\mathit{down} \gets \text{C2}$，
$(1,2).\mathit{down} \gets y$。
\end{enumerate}

\item 每次比較之後，恰好有一個節點接到結果上，而且以後不再參與比較。 迴圈在 $a$ 或
$b$ 變成 \NIL 時停。

下界：迴圈停下時，某一條串列的節點已經全部經由比較接上了，所以 $C \ge \min(m, n)$。

上界：迴圈停下時，另一條串列至少還剩一個節點沒有比較就整段接上，經由比較接上的節點
最多 $m + n - 1$ 個，所以 $C \le m + n - 1$。

達到下界：$a$ 是 $1, 2, \ldots, m$，$b$ 是 $m + 1, \ldots, m + n$，而且 $m \le n$。
每次比較都接上 $a$ 的節點，$m$ 次後 $a$ 用完，$C = m = \min(m, n)$（$m > n$ 時把兩條
對調）。

達到上界：讓兩條串列的最大值分別是全部的最大和第二大，例如 $a$ 是
$1, 2, \ldots, m - 1, m + n$，$b$ 是 $m, m + 1, \ldots, m + n - 1$。 $b$ 的最後一個
節點 $m + n - 1$ 接上之前，$a$ 一直還有 $m + n$ 沒接，所以前 $m + n - 1$ 個節點都是
經由比較接上的，$C = m + n - 1$。

\item
\begin{enumerate}
\item 要證明的敘述：對 $n$ 個節點的串列 $x_1 \to \cdots \to x_n$（$n = 0$ 時 $h =
\NIL$），\proc{Rev-Rec}$(h)$ 回傳 $x_n$，執行完之後串列是 $x_n \to \cdots \to x_1 \to
\NIL$。

基底：$n = 0$ 或 $n = 1$，第 2、3 行直接回傳 $h$，空串列和一個節點的串列反轉後不變。

歸納：設 $n \ge 2$，而且敘述對 $n - 1$ 個節點成立。 第 4 行對 $x_2 \to \cdots \to
x_n$ 遞迴，由歸納假設，回傳 $r = x_n$，而且 $x_n \to \cdots \to x_2 \to \NIL$。 遞迴
只改了 $x_2, \ldots, x_n$ 的 $\mathit{next}$，所以 $x_1.\mathit{next}$ 仍然是
$x_2$。 第 5 行讓 $x_2.\mathit{next} = x_1$，第 6 行讓 $x_1.\mathit{next} = \NIL$，
串列變成 $x_n \to \cdots \to x_2 \to x_1 \to \NIL$，回傳 $x_n$。

\item 這裡第 6 行指的是 $h.\mathit{next} \gets \NIL$。 拿掉之後，最外層呼叫結束時
$x_1.\mathit{next}$ 仍然是 $x_2$，而 $x_2.\mathit{next} = x_1$，串列變成 $3 \to 2
\to 1 \to 2 \to 1 \to \cdots$，$1$ 和 $2$ 互相指著，形成一個環。 走訪它的程式會停不
下來。 （內層呼叫留下的 $x_i.\mathit{next} = x_{i+1}$ 不會出問題，因為上一層的第 5
行會把它改成 $x_{i-1}$；只有最外層的 $x_1$ 沒有上一層來改。）

\item $n \ge 1$ 時呼叫最多同時疊了 $n$ 層（$x_1$ 到 $x_n$ 各一層），每層 $\Theta(1)$，
額外空間 $\Theta(n)$；迴圈版只要 $\Theta(1)$。
\end{enumerate}

\item 用 \texttt{link} 指著「目前要檢查的節點的那個指標欄位」。 刪掉時
\texttt{link} 不動（它現在指著下一個節點），沒刪才往前一格。
\begin{code}{c}{觀念題 8 的解答}
void remove_all(Node **head, int k)
{
    Node **link = head;
    while (*link != NULL) {
        if ((*link)->id == k) {
            Node *dead = *link;
            *link = dead->next;     // 跳過 dead
            free(dead);
        } else {
            link = &(*link)->next;
        }
    }
}
\end{code}
每個節點檢查一次，$O(n)$。 第一個節點要刪時，\texttt{link} 就是 \texttt{head}，
\texttt{*link = dead->next} 直接更新了呼叫端的 \texttt{head}，不需要特殊情況。 連續
好幾個節點都要刪時，\texttt{link} 一直不動，也是對的。

\item (a) \texttt{free(p)} 之後，迴圈的 \texttt{p = p->next} 讀了已經釋放的記憶體。
要先記住下一個再釋放：迴圈裡寫 \texttt{Node *next = p->next; free(p); p = next;}，
和範例~\ref{ex:ll-poly} 的 \texttt{print\_and\_free} 一樣。

(b) 兩個錯。 第一，\texttt{head} 是參數的複本，改它不會改到呼叫端；函式回來之後，
呼叫端的 \texttt{head} 還指著已經釋放的節點。 第二，空串列時 \texttt{head->next}
讀了 \NIL 的欄位。 改成收 \texttt{Node **head}：先檢查 \texttt{*head != NULL}，再
寫 \texttt{Node *p = *head; *head = p->next; free(p);}，呼叫時傳
\texttt{\&head}。 另一種改法是回傳新的開頭，呼叫端寫
\texttt{head = drop\_first(head);}。

\item 第一組的第一個人是原本的 $f = R.\mathit{last}.\mathit{next}$，最後一個人是 $p$；
第二組的第一個人是 $g = p.\mathit{next}$，最後一個人是原本的 $R.\mathit{last}$。 拆開
之後，$p$ 要接回 $f$，$R.\mathit{last}$ 要接到 $g$。
\begin{pseudo}{觀念題 10 的解答}
\Function{Ring-Split}{$R, p$}
    \State $f \gets R.\mathit{last}.\mathit{next}$ \Comment{今天值日的人}
    \State $g \gets p.\mathit{next}$ \Comment{第二組的第一個人}
    \State $p.\mathit{next} \gets f$
    \State $R.\mathit{last}.\mathit{next} \gets g$
    \State \Return $(p, R.\mathit{last})$
\EndFunction
\end{pseudo}
只讀兩個、改兩個指標，$\Theta(1)$。 $p$ 就是 $f$ 的時候（第一組只有一個人），第 4 行讓
$p$ 指向自己，仍然正確。 只記第一個節點時，第 5 行要改的是最後一個節點的
$\mathit{next}$，得繞一圈才找得到，變成 $\Theta(n)$。

\item
\begin{enumerate}
\item 做完 $k$ 輪後 $s = x_{1+k}$、$f = x_{1+2k}$。 迴圈繼續的條件是 $x_{1+2k}$ 和
$x_{2+2k}$ 都存在，即 $2k + 2 \le n$，所以跑 $\lfloor n/2 \rfloor$ 輪，結束時
$s = x_{\lfloor n/2 \rfloor + 1}$。 $n = 1, \ldots, 6$ 依序是第 $1, 2, 2, 3, 3, 4$
個。 $n$ 是奇數時和原來一樣停在正中間；$n$ 是偶數時停在兩個中間節點的後一個。
\item $n = 2$ 時 $s = x_2$，在 $s$ 後面切開，前半是 $x_1, x_2$，後半是空的。 合併排序
對前半遞迴，又是同一條長度 $2$ 的串列，永遠切不小，無窮遞迴（實際執行會用光呼叫堆
疊）。 所以給合併排序用的切法，偶數長度時一定要停在前一個中間節點，也就是原本的
$f \gets h.\mathit{next}$。
\end{enumerate}

\item
\begin{enumerate}
\item 尾巴是 $x_0, \ldots, x_5$，$\mu = 6$；環是 $x_6 \to x_7 \to x_8 \to x_9 \to x_6$，
$\lambda = 4$。 做完 $k$ 輪後 $s = x_k$、$f = x_{2k}$；足標 $t \ge 6$ 時，$x_t$ 就是
$x_{6 + ((t - 6) \bmod 4)}$。
\begin{center}
\begin{tabular}{ccccccccc}
\toprule
輪 $k$ & 1 & 2 & 3 & 4 & 5 & 6 & 7 & 8 \\
\midrule
$s$ & $x_1$ & $x_2$ & $x_3$ & $x_4$ & $x_5$ & $x_6$ & $x_7$ & $x_8$ \\
$f$ & $x_2$ & $x_4$ & $x_6$ & $x_8$ & $x_6$ & $x_8$ & $x_6$ & $x_8$ \\
\bottomrule
\end{tabular}
\end{center}
第 $8$ 輪相遇在 $x_8$。 用定理~\ref{thm:ll-floyd}：$4$ 的倍數裡不小於 $6$ 的最小是
$8$，一致。
\item $\mu = 6$ 輪：$s$ 從 $x_0$、$f$ 從 $x_8$ 出發，$6$ 輪後 $s = x_6$，$f = x_{14} =
x_6$。
\item $\mu = 0$、$\lambda = 10$。 \proc{Floyd-Meet} 要等 $k^* = 10$ 輪，相遇在
$x_{10} = x_0$；第二個迴圈一開始 $s = f = x_0$，跑 $0$ 輪，起點是 $x_0$。
\item 指向 $x_\mu$ 時尾巴長 $\mu$、環長 $10 - \mu$，$k^*$ 是 $10 - \mu$ 的倍數裡不小於
$\max(\mu, 1)$ 的最小那一個：
\begin{center}
\begin{tabular}{ccccccccccc}
\toprule
$\mu$ & 0 & 1 & 2 & 3 & 4 & 5 & 6 & 7 & 8 & 9 \\
\midrule
$k^*$ & 10 & 9 & 8 & 7 & 6 & 5 & 8 & 9 & 8 & 9 \\
\bottomrule
\end{tabular}
\end{center}
最多 $10$ 輪（指向 $x_0$，整條是一個環），最少 $5$ 輪（指向 $x_5$）。 一般來說
$k^*$ 是 $\lambda$ 的正倍數，又不小於 $\mu$，所以 $k^* \ge \max(\mu, \lambda) \ge
\lceil n/2 \rceil$；$\mu = \lfloor n/2 \rfloor$ 時 $\lambda = \lceil n/2 \rceil \ge \mu$，
$k^* = \lambda$，下界達到。 上界 $k^* \le n$ 只在 $\mu = 0$ 時達到。
\end{enumerate}

\item 想法：把 $a$ 那條串列接成一個環，再從 $b$ 出發找環的起點。
\begin{pseudo}{觀念題 13 的解答}
\Function{First-Common}{$a, b$}
    \If{$a = \NIL$ \OR $b = \NIL$}
        \State \Return \NIL
    \EndIf
    \State $t \gets a$
    \While{$t.\mathit{next} \ne \NIL$} \Comment{找 $a$ 的最後一個節點}
        \State $t \gets t.\mathit{next}$
    \EndWhile
    \State $t.\mathit{next} \gets a$ \Comment{暫時接成環}
    \State $X \gets \Call{Cycle-Start}{b}$
    \State $t.\mathit{next} \gets \NIL$ \Comment{改回來}
    \State \Return $X$
\EndFunction
\end{pseudo}
正確性：接成環之後，環上的節點恰好是 $a$ 那條串列的全部節點。 如果兩條串列從 $X$
開始共用，從 $b$ 出發會先經過 $b$ 自己的節點（不在 $a$ 上，也就不在環上），第一個
走到的環上節點就是 $X$；之後一直在環上打轉。 所以從 $b$ 看，這是一條尾巴是 $b$ 獨有
的那一段、環的起點是 $X$ 的 $\rho$ 形串列，由定理~\ref{thm:ll-floyd}，
\proc{Cycle-Start} 回傳 $X$。 如果沒有共用的節點，從 $b$ 出發永遠碰不到環，走到
\NIL，\proc{Cycle-Start} 回傳 \NIL。 第 9 行在回傳之前把 $t.\mathit{next}$ 改回
\NIL，兩條串列恢復原狀。

時間：找 $t$ 走 $m - 1$ 步；從 $b$ 走得到的節點最多 $m + n$ 個，\proc{Cycle-Start}
是 $O(m + n)$。 合計 $O(m + n)$，額外空間只有幾個指標。

不改動指標的做法：先各走一次求出 $m$、$n$ 和兩條的最後一個節點；最後一個節點不同
就沒有共用。 相同的話，讓長的那條先走 $|m - n|$ 步，再兩條一起一次一步，第一次指到
同一個節點就是 $X$（共用的後半段等長，所以兩個指標會同時到 $X$）。 一樣是
$O(m + n)$ 時間、$O(1)$ 空間。
\end{enumerate}

\addcontentsline{toc}{section}{歷屆試題詳解}
\section*{歷屆試題詳解}

\pastsol{px:ll-1}{臺大資工所 110 學年度 第 VII 題}
答案是 (B)。

題目沒有說是哪一種串列，但它說刪掉最後一個元素要花時間、而且只多一個 tail 指標，
這裡照單向串列來解（最後說明為什麼不是雙向）。 選項寫的「$= 1$」「$= n$」不能照字面
當成比值，應該讀成成長的階數：「$= 1$」是分子和分母同階，「$= n$」是分子比分母多一
個 $n$ 的因子。

四個函數（第~\ref{sec:ll-variants}~節的比較表就是這四格）：
\begin{itemize}
\item $f_1(n) = \Theta(n)$：只有 head，要把 $x_{n+1}$ 接到 $x_n$ 後面，得先從 head 走
$n - 1$ 步找到 $x_n$。
\item $g_1(n) = \Theta(n)$：刪 $x_n$ 要把 $x_{n-1}.\mathit{next}$ 設成 \NIL，得先從
head 走到 $x_{n-1}$。
\item $f_2(n) = \Theta(1)$：$\mathit{tail}.\mathit{next} \gets x_{n+1}$，
$\mathit{tail} \gets x_{n+1}$。
\item $g_2(n) = \Theta(n)$：刪掉 $x_n$ 之後 tail 要改成指向 $x_{n-1}$，單向串列不能從
$x_n$ 往回走，還是得從 head 找 $x_{n-1}$。 多一個 tail 指標對刪除最後一個沒有幫助。
\end{itemize}

逐項檢查：
\begin{itemize}
\item (A) $f_1 / g_1 = \Theta(n) / \Theta(n)$，同階，對。
\item (B) $f_2 / g_2 = \Theta(1) / \Theta(n) = \Theta(1/n)$，不是 $n$，\textbf{錯}。 加了 tail 之後
變快的只有插入，刪除沒有變。
\item (C) $f_1 / f_2 = \Theta(n) / \Theta(1) = \Theta(n)$，對。
\item (D) $g_1 / g_2 = \Theta(n) / \Theta(n)$，同階，對。
\end{itemize}

如果是雙向串列，$g_2$ 會變成 $\Theta(1)$（$x_n.\mathit{prev}$ 就是 $x_{n-1}$），這
時 (D) 的比值是 $n$，也錯了，單選題會有兩個答案。 所以題意是單向串列。

\pastsol{px:ll-2}{成大資訊聯招 114 學年度 第 2 題}
答案是 (B)。

函式做的事：把下一個節點的資料抄到 \texttt{node} 裡，再把下一個節點從串列上拿掉、
釋放。 被釋放的其實是下一個節點，但串列上少掉的是 \texttt{node} 原本的那筆資料，看
起來就像把 \texttt{node} 刪掉了。 觀念題第 1 題 (d) 的解答用同一招在 $p$ 前面插入，
這裡反過來用在刪除。 例如串列 $3 \to 8 \to 5$：
\begin{itemize}
\item 對存 $8$ 的節點呼叫：它的 \texttt{data} 變成 $5$，\texttt{next} 改成原本 $5$
那個節點的 \texttt{next}（\texttt{NULL}），再釋放原本存 $5$ 的節點。 結果是 $3 \to 5$，
正是「刪掉 $8$」。
\item 對最後一個節點（存 $5$）呼叫：\texttt{node->next == NULL}，直接返回，串列不變。
\end{itemize}

逐項檢查：
\begin{itemize}
\item (A) 錯。 對最後一個節點，函式第一個 \texttt{if} 就返回了，什麼都沒刪。 這也不
是寫法的疏忽：刪掉最後一個節點要把它前一個節點的 \texttt{next} 設成 \texttt{NULL}，
函式只拿到 \texttt{node}，沒辦法找到前一個。
\item (B) 對。 只要 \texttt{node} 不是最後一個，串列上就少了它的資料。
\item (C) 錯。 函式一開始就檢查 \texttt{node->next == NULL}，最後一個節點在這裡就
返回了，後面的 \texttt{node->next->data} 根本沒有執行，不會 segmentation fault。
\item (D) 錯。 函式沒有用資料值去找任何東西，串列裡有沒有重複的值都一樣；它改的兩
個地方（\texttt{node->data}、\texttt{node->next}）都接得好好的，結構沒有壞。
\item (E) 錯，因為 (B) 對。
\end{itemize}

補充：這個技巧在考卷上常出現，但實務上要小心。 真正被 \texttt{free} 的是下一個節點，
如果程式別的地方還留著指向它的指標（例如串列另外記了 tail，而 \texttt{node} 正好是
倒數第二個），那個指標就指到已經釋放的記憶體；指向 \texttt{node} 的指標也會發現它
存的資料變了。

\pastsol{px:ll-3}{中央資工所 113 學年度 第 8 題}
答案是 (C)。

\texttt{fun} 對每個節點把 \texttt{prev}、\texttt{next} 兩個指標對調，然後讓
\texttt{current} 走到對調後的 \texttt{prev}，也就是原本的下一個節點，所以每個節點都
會被處理一次。 每個節點兩個方向都反過來，整條串列就反轉了。 迴圈結束時
\texttt{temp} 是最後處理的節點 $x_n$ 原本的 \texttt{prev}，也就是 $x_{n-1}$；
$x_{n-1}$ 對調之後 \texttt{prev} 是原本的 \texttt{next}，也就是 $x_n$，所以
\texttt{*head\_ref = temp->prev} 讓新的開頭是 $x_n$。

用 L $= 1 \leftrightarrow 2 \leftrightarrow 3 \leftrightarrow 4 \leftrightarrow 5
\leftrightarrow 6$ 追蹤，每一列是迴圈的一輪（「之後」是對調後的值）：
\begin{center}
\begin{tabular}{ccccc}
\toprule
\texttt{current} & \texttt{temp} & 之後的 \texttt{prev} & 之後的 \texttt{next} & 下一個 \texttt{current} \\
\midrule
1 & \NIL & 2 & \NIL & 2 \\
2 & 1 & 3 & 1 & 3 \\
3 & 2 & 4 & 2 & 4 \\
4 & 3 & 5 & 3 & 5 \\
5 & 4 & 6 & 4 & 6 \\
6 & 5 & \NIL & 5 & \NIL \\
\bottomrule
\end{tabular}
\end{center}
迴圈結束時 \texttt{temp} 是節點 $5$，\texttt{temp->prev} 是 $6$，新的開頭是 $6$。
從 $6$ 沿著 \texttt{next} 走：$6, 5, 4, 3, 2, 1$；沿著 \texttt{prev} 從 $1$ 往回走也
是同樣的順序，兩個方向都接對了。

逐項檢查：
\begin{itemize}
\item (A) 錯。 \texttt{fun} 會把任何串列反轉，$2, 1, 4, 3, 6, 5$ 反轉是
$5, 6, 3, 4, 1, 2$，不是原樣。
\item (B) 錯。 $5, 4, 3, 2, 6, 1$ 反轉是 $1, 6, 2, 3, 4, 5$。
\item (C) 對，上面的追蹤。
\item (D) 錯。 反轉後 $1$ 在最後，不會是 $\ldots 1, 2$。
\end{itemize}

邊界情況：只有一個節點時，迴圈跑一輪，\texttt{temp} 是 \texttt{NULL}，開頭不變，正
確；空串列時迴圈不跑，也正確。 時間 $\Theta(n)$，額外空間 $\Theta(1)$。 和
第~\ref{sec:ll-ops}~節單向串列的 \proc{List-Reverse} 比較：雙向串列每個節點本來就
有兩個方向的指標，對調就好，不用另外記「已經反轉好的一段」。

\pastsol{px:ll-4}{臺科大資工所 112 學年度 第 2 題}
答案：最壞情況 $\Theta(m + n)$（寫成 $O(m + n)$ 也可以）。

\textbf{上界。} 迴圈每一輪做一次（或兩次）比較，每一輪至少讓 \texttt{ptr1}、
\texttt{ptr2} 其中一個前進一格。 \texttt{ptr1} 最多前進 $m$ 次，\texttt{ptr2} 最多
$n$ 次。 設只動 \texttt{ptr1} 的有 $a$ 輪、只動 \texttt{ptr2} 的有 $b$ 輪、兩個都動
的（值相等）有 $c$ 輪，則 $a + c \le m$、$b + c \le n$，輪數是 $a + b + c$。 迴圈在
其中一個變成 \texttt{NULL} 時停：
\begin{itemize}
\item 若最後一輪是值相等、兩個同時變成 \texttt{NULL}，則 $c \ge 1$，輪數
$a + b + c \le (m - c) + (n - c) + c = m + n - c \le m + n - 1$。
\item 否則停下時只有一個走完，例如 \texttt{ptr1} 走完、\texttt{ptr2} 還沒，
$b + c \le n - 1$，輪數 $\le m + n - 1$。 反過來一樣。
\end{itemize}
所以最多 $m + n - 1$ 輪，每輪 $\Theta(1)$，時間 $O(m + n)$。

\textbf{達得到。} 讓兩條串列沒有相同的值，而且全部裡最大和第二大的值分在兩條，例如
list1 是 $1, 2, \ldots, m - 1, m + n$，list2 是 $m, m + 1, \ldots, m + n - 1$。
\texttt{ptr1} 先走過 $1, \ldots, m - 1$（$m - 1$ 輪），停在 $m + n$；之後
\texttt{ptr2} 走過 list2 全部 $n$ 個才變成 \texttt{NULL}。 共 $m + n - 1$ 輪，所以最
壞情況是 $\Theta(m + n)$。 （這和觀念題第 6 題讓 \proc{Merge-Lists} 比較 $m + n - 1$
次的輸入是同一組。）

函式算的是兩條串列共同的值有幾個（兩條各自沒有重複值時，就是交集的大小），走法和
第~\ref{sec:ll-ops}~節的 \proc{Merge-Lists} 一樣，只是不接節點。 最好的情況是一條
串列的值全部比另一條小，$\min(m, n)$ 輪就停，$\Theta(\min(m, n))$。

\pastsol{px:ll-5}{政大資科所 108 學年度 第 9 題}
約定：輸入是串列第一個節點的指標 $h$，每個節點有 $\mathit{next}$；輸出 \TRUE（有環）
或 \FALSE。 用 Floyd 的快慢指標（第~\ref{sec:ll-cycle}~節）：

\begin{pseudo}{政大 108 第 9 題的解答}
\Function{Has-Cycle}{$h$}
    \State $s \gets h$；$f \gets h$ \Comment{慢指標、快指標}
    \While{$f \ne \NIL$ \AND $f.\mathit{next} \ne \NIL$}
        \State $s \gets s.\mathit{next}$ \Comment{走一步}
        \State $f \gets f.\mathit{next}.\mathit{next}$ \Comment{走兩步}
        \If{$s = f$}
            \State \Return \TRUE
        \EndIf
    \EndWhile
    \State \Return \FALSE
\EndFunction
\end{pseudo}

\textbf{正確性。} 做完 $k$ 輪後，$s$ 是從 $h$ 走 $k$ 步到的節點，$f$ 是走 $2k$ 步到
的節點。
\begin{itemize}
\item 沒有環：串列有 $n$ 個節點，$f$ 一直在 $s$ 前面（$2k > k$），兩者是不同的節點，
第 7 行不會執行；$f$ 每輪前進兩格，最多 $\lfloor n/2 \rfloor$ 輪就使 $f$ 或
$f.\mathit{next}$ 變成 \NIL，迴圈停下，回傳 \FALSE。
\item 有環：串列上沒有 \NIL，迴圈只會從第 7 行離開。 設尾巴長 $\mu$、環長 $\lambda$。
第 $\mu$ 輪 $s$ 走到環上，這時 $f$ 早已在環上，沿著環的方向還要追 $g$ 步才碰到 $s$
（$0 \le g < \lambda$）。 之後每輪 $f$ 比 $s$ 多走一步，$g$ 恰好減 $1$，不會跳過
$0$，所以再 $g$ 輪內相遇（$\mu = 0$ 時第 $\lambda$ 輪相遇），回傳 \TRUE。 完整的證
明是定理~\ref{thm:ll-floyd}。
\end{itemize}

\textbf{複雜度。} 沒有環時跑 $\lfloor n/2 \rfloor$ 輪；有環時跑的輪數不超過
$\mu + \lambda = n$。 每輪 $\Theta(1)$，時間 $O(n)$；只用兩個指標，額外空間
$\Theta(1)$，也不改動串列。 題目圖上的串列是整個環（$\mu = 0$、$\lambda = 5$），第
$5$ 輪兩個指標在第一個節點相遇。

另一種寫法是用雜湊表（第~12~章）記錄走過的節點，走到記錄過的節點就有環，走到 \NIL
就沒有；時間一樣是期望 $\Theta(n)$，但要 $\Theta(n)$ 的空間。 題目沒有限制空間時這
也是正確答案，寫 Floyd 的方法比較好。 如果題目還要求環的起點，接著做
\proc{Cycle-Start}（虛擬碼~\ref{alg:ll-floyd}）。

\chapter{堆疊與佇列}
\label{ch:stackqueue}

% 本章的圖共用的樣式與巨集（名稱都以 sq 開頭，避免和其他章衝突）
\tikzset{
    sqold/.style={cell, text=black!35},
    sqfill/.style={cell, fill=coreFill, draw=coreStroke},
    sqnode/.style={draw=black!65, minimum width=0.9cm, minimum height=0.7cm, font=\ttfamily},
}
\newcommand{\sqpow}{\text{\textasciicircum}}

用文書軟體打字時，按一次「復原」會取消最後做的那個動作，再按一次，取消的是再前面
那一個。 系辦公室的共用印表機剛好相反，誰先送出文件，就先印誰的。 這兩件事都要記
住一串還沒處理完的東西，差別只在要從哪一頭拿。 如果要你設計記錄的方式，讓「放進
一個」和「拿出一個」都只花固定的時間，你會怎麼存？ 往下讀之前，先想一想。

復原用的是\textbf{堆疊}（stack）：最後放進去的最先拿出來。 印表機用的是
\textbf{佇列}（queue）：最先放進去的最先拿出來。 兩者都只開放兩三個操作，正因為
限制多，每個操作都能做到 $\Theta(1)$。 這一章先用陣列和鏈結串列各實作一次，再看
它們最常見的幾個用途：運算式的轉換與求值、把遞迴改寫成迴圈，以及好幾個堆疊共用
一塊記憶體。 中間會穿插考卷上常見的手算：堆疊能排出哪些順序、環狀佇列的
$\mathit{front}$、$\mathit{rear}$ 怎麼變、中序轉後序的逐步堆疊表。

這一章用到《C 程式設計基礎》第~14~章的結構指標、第~15~章遞迴與呼叫堆疊、第~19~章
的 \texttt{malloc} 與串列，以及第~20~章用 \texttt{struct} 包裝的堆疊；用到時先在
框裡複習。 鏈結串列本身的操作在本書第~5~章。

\section{堆疊：陣列與串列兩種實作}
\label{sec:sq-stack}

\textbf{堆疊}是一串資料，加入和取出都只能在同一端。 這一端叫\textbf{頂端}（top），
另一端叫底部（bottom）。 它支援的操作：

\begin{itemize}
\item \proc{Push}$(S, x)$：把 $x$ 放到頂端。
\item \proc{Pop}$(S)$：取出頂端的元素並回傳。
\item \proc{Top}$(S)$：回傳頂端的元素，不取出（有些書叫 peek）。
\item \proc{Stack-Empty}$(S)$：堆疊是不是空的。
\end{itemize}

最後放進去的元素最先被取出，所以堆疊也叫\textbf{後進先出}（last in, first out，
LIFO）的結構。 對空的堆疊做 \proc{Pop} 叫\textbf{下溢}（underflow）；容量有限的
堆疊已經滿了還做 \proc{Push}，叫\textbf{溢位}（overflow）。 兩者都是錯誤，程式要
檢查，不能假設不會發生。

上面只規定了有哪些操作、每個操作做什麼，沒有規定資料怎麼存。 這種只描述操作的規
格叫\textbf{抽象資料型態}（abstract data type，ADT）。 同一個 ADT 可以有好幾種實
作，用的人只看操作，不必知道裡面怎麼存。 下面給兩種實作。

\subsection{用陣列存}

用一個陣列 $S.\mathit{item}[1..S.\mathit{cap}]$ 存元素，再用一個整數
$S.\mathit{top}$ 記目前有幾個元素。 $S.\mathit{item}[1]$ 是底部，
$S.\mathit{item}[S.\mathit{top}]$ 是頂端；$S.\mathit{top} = 0$ 表示堆疊是空的。

\begin{pseudo}[label=alg:sq-array-stack]{陣列堆疊}
\Function{Stack-Empty}{$S$}
    \State \Return $S.\mathit{top} = 0$
\EndFunction
\Statex
\Procedure{Push}{$S, x$}
    \If{$S.\mathit{top} = S.\mathit{cap}$}
        \State \Error{overflow}
    \EndIf
    \State $S.\mathit{top} \gets S.\mathit{top} + 1$
    \State $S.\mathit{item}[S.\mathit{top}] \gets x$
\EndProcedure
\Statex
\Function{Pop}{$S$}
    \If{\Call{Stack-Empty}{$S$}}
        \State \Error{underflow}
    \EndIf
    \State $S.\mathit{top} \gets S.\mathit{top} - 1$
    \State \Return $S.\mathit{item}[S.\mathit{top} + 1]$
\EndFunction
\end{pseudo}

\proc{Top}$(S)$ 和 \proc{Pop} 一樣先檢查是不是空的，再回傳
$S.\mathit{item}[S.\mathit{top}]$，但不改 $S.\mathit{top}$。 每個操作都只做幾個
固定的步驟，是 $\Theta(1)$。

\proc{Pop} 只把 $S.\mathit{top}$ 減一，被取出的值還留在陣列裡，只是不再屬於堆疊。
下表用容量 $3$ 的陣列走一串操作，灰字是這種「還在陣列裡、但已經不在堆疊裡」的值。

\begin{center}
\begin{tabular}{llcccc}
\toprule
操作 & 回傳 & $S.\mathit{item}[1]$ & $[2]$ & $[3]$ & $S.\mathit{top}$ \\
\midrule
\proc{Push}$(S, 4)$ & & 4 & & & 1 \\
\proc{Push}$(S, 7)$ & & 4 & 7 & & 2 \\
\proc{Push}$(S, 1)$ & & 4 & 7 & 1 & 3 \\
\proc{Pop}$(S)$ & 1 & 4 & 7 & \textcolor{black!40}{1} & 2 \\
\proc{Push}$(S, 6)$ & & 4 & 7 & 6 & 3 \\
\proc{Push}$(S, 2)$ & overflow & 4 & 7 & 6 & 3 \\
\proc{Pop}$(S)$ & 6 & 4 & 7 & \textcolor{black!40}{6} & 2 \\
\proc{Pop}$(S)$ & 7 & 4 & \textcolor{black!40}{7} & \textcolor{black!40}{6} & 1 \\
\proc{Pop}$(S)$ & 4 & \textcolor{black!40}{4} & \textcolor{black!40}{7} & \textcolor{black!40}{6} & 0 \\
\proc{Pop}$(S)$ & underflow & \textcolor{black!40}{4} & \textcolor{black!40}{7} & \textcolor{black!40}{6} & 0 \\
\bottomrule
\end{tabular}
\end{center}

容量固定是陣列堆疊的缺點：開太大浪費記憶體，開太小會溢位。 改用「滿了就換一個兩
倍大的陣列」可以避免溢位，單次 \proc{Push} 最壞要搬 $\Theta(n)$ 個元素，但平均每
次仍然是 $\Theta(1)$，第~3~章用平攤分析證明這件事。

\subsection{用串列存}

另一種做法是把每個元素放進一個節點，節點用 $\mathit{next}$ 指向它下面的那一個，
$S.\mathit{top}$ 指向頂端的節點，空堆疊時是 \NIL。 \proc{Push} 和 \proc{Pop} 都
只動串列的頭。

\begin{pseudo}[label=alg:sq-list-stack]{串列堆疊}
\Procedure{Push-List}{$S, x$}
    \State 配置一個新節點 $p$
    \State $p.\mathit{key} \gets x$
    \State $p.\mathit{next} \gets S.\mathit{top}$
    \State $S.\mathit{top} \gets p$
\EndProcedure
\Statex
\Function{Pop-List}{$S$}
    \If{$S.\mathit{top} = \NIL$}
        \State \Error{underflow}
    \EndIf
    \State $p \gets S.\mathit{top}$
    \State $x \gets p.\mathit{key}$
    \State $S.\mathit{top} \gets p.\mathit{next}$
    \State 釋放 $p$
    \State \Return $x$
\EndFunction
\end{pseudo}

為什麼頂端要放在串列的頭，不放在尾巴？ 單向串列要刪掉最後一個節點，得先找到它
前面那一個，這要從頭走過去，花 $\Theta(n)$。 放在頭的話，加入和刪除都只改一兩個
指標。

\begin{center}
\begin{tikzpicture}
% 左：陣列
\node[font=\small] at (1.8, 1.2) {陣列，$\mathit{cap} = 5$};
\foreach \v [count=\k] in {4,7,6} { \node[sqfill] at (\k*0.9-0.9, 0) {\v}; }
\foreach \k in {4,5} { \node[cell] at (\k*0.9-0.9, 0) {}; }
\foreach \k in {1,...,5} { \node[idxlbl] at (\k*0.9-0.9, -0.65) {\k}; }
\node[font=\small] (tp) at (1.8, -1.4) {$\mathit{top} = 3$};
\draw[diagarrow] (tp.north) -- (1.8, -0.8);
% 右：串列
\node[font=\small] at (9.6, 1.2) {串列};
\node[font=\small] (top) at (6.6, 0) {$\mathit{top}$};
\node[sqnode, fill=coreFill, draw=coreStroke] (a) at (8.0, 0) {6};
\node[sqnode] (b) at (9.6, 0) {7};
\node[sqnode] (c) at (11.2, 0) {4};
\node[font=\small] (nil) at (12.6, 0) {\NIL};
\draw[diagarrow] (top) -- (a);
\draw[diagarrow] (a) -- (b);
\draw[diagarrow] (b) -- (c);
\draw[diagarrow] (c) -- (nil);
\end{tikzpicture}
\diagcaption{同一個堆疊（由底到頂是 $4, 7, 6$）的兩種存法。 陣列的頂端在
$\mathit{item}[\mathit{top}]$；串列的頂端是第一個節點，箭頭指向它下面的元素。}
\label{fig:sq-two-stacks}
\end{center}

兩種做法比較如下。 四個操作在兩種做法裡都是 $\Theta(1)$，差別在記憶體。

\begin{center}
\begin{tabular}{p{0.22\linewidth}p{0.34\linewidth}p{0.34\linewidth}}
\toprule
 & 陣列 & 串列 \\
\midrule
容量 & 事先決定，會溢位（除非換大陣列）& 只受記憶體限制 \\
每個元素的空間 & 只有資料本身 & 資料再加一個指標 \\
每次 \proc{Push} & 寫一格 & 呼叫一次 \texttt{malloc} \\
沒用到的空間 & 陣列後段空著 & 沒有 \\
\bottomrule
\end{tabular}
\end{center}

\subsection{C 程式}

\begin{cbookrecap}[第~14~章 14.5 節、第~19~章 19.3 節、第~20~章 20.2 節]
第~20~章把「陣列加上 \texttt{top}」包進一個 \texttt{struct}，操作函式都收一個指
向它的指標，例如 \texttt{stack\_push(\&s, x)}；陣列版的 push 就是
\texttt{s->data[s->top++] = x}，pop 是 \texttt{return s->data[--s->top]}。 透過結
構指標取欄位用箭頭：\texttt{s->top} 就是 \texttt{(*s).top}。

節點用 \texttt{malloc(sizeof *p)} 配置，\texttt{sizeof *p} 會跟著 \texttt{p} 的型
別走，型別改了也不用改這一行。 取出節點時先把要的值和 \texttt{next} 存起來，再
\texttt{free}，順序反過來就是讀已經釋放的記憶體。
\end{cbookrecap}

陣列版在《C 程式設計基礎》第~20~章寫過了，這裡寫串列版，並用上面那張表的同一串操
作測試，最後多 push 一個 $5$。 每一行印出這次操作的結果、堆疊的元素個數和頂端。

\begin{example}[label=ex:sq-list-stack]{c}{(list\_stack.c) 串列堆疊}
#include <stdio.h>
#include <stdlib.h>

typedef struct Node {
    int key;
    struct Node *next;      // 指向下面那一個
} Node;

typedef struct {
    Node *top;              // 空堆疊時是 NULL
    int size;
} Stack;

int push(Stack *s, int x)   // 只有 malloc 失敗才回傳 0
{
    Node *p = malloc(sizeof *p);
    if (p == NULL) {
        return 0;
    }
    p->key = x;
    p->next = s->top;
    s->top = p;
    s->size++;
    return 1;
}

int pop(Stack *s, int *x)   // 空的回傳 0
{
    if (s->top == NULL) {
        return 0;
    }
    Node *p = s->top;
    *x = p->key;            // 先把要的東西存起來，再 free
    s->top = p->next;
    free(p);
    s->size--;
    return 1;
}

int main(void)
{
    Stack s = {NULL, 0};
    int ops[] = {4, 7, 1, 0, 6, 2, 0, 0, 0, 0, 0, 5};  // 正數 x 是 push x，0 是 pop
    int m = sizeof ops / sizeof ops[0];
    for (int i = 0; i < m; i++) {
        int x;
        if (ops[i] > 0) {
            if (!push(&s, ops[i])) {
                return 1;
            }
            printf("push %d:", ops[i]);
        } else if (pop(&s, &x)) {
            printf("pop -> %d:", x);
        } else {
            printf("pop -> underflow:");
        }
        if (s.top != NULL) {
            printf(" size %d, top %d\n", s.size, s.top->key);
        } else {
            printf(" empty\n");
        }
    }
    int x;
    while (pop(&s, &x)) {   // 結束前把剩下的節點都釋放
    }
    return 0;
}
\end{example}
\iosample{}{push 4: size 1, top 4 \\ push 7: size 2, top 7 \\ push 1: size 3, top 1 \\ pop -> 1: size 2, top 7 \\ push 6: size 3, top 6 \\ push 2: size 4, top 2 \\ pop -> 2: size 3, top 6 \\ pop -> 6: size 2, top 7 \\ pop -> 7: size 1, top 4 \\ pop -> 4: empty \\ pop -> underflow: empty \\ push 5: size 1, top 5}

到 \texttt{push 2} 之前都和陣列版的表一樣。 串列沒有容量限制，$2$ 放進去了，所以
之後 pop 出來的是 $2, 6, 7, 4$，陣列版是 $6, 7, 4$ 再下溢。 溢位不能只印一句警告
就繼續：呼叫的人以為 $2$ 在堆疊裡，後面的結果就全錯了。 範例的函式用回傳值告訴呼
叫的人成功與否，由呼叫的人決定怎麼處理。 程式結束前堆疊裡還有 $5$，最後的迴圈把
剩下的節點都 pop 掉，每個 \texttt{malloc} 都有對應的 \texttt{free}。

C 的陣列從 $0$ 開始，陣列版的 \texttt{top} 同時是「元素個數」和「下一個空位的索
引」，\texttt{data[top - 1]} 才是頂端。 虛擬碼的 $S.\mathit{item}[S.\mathit{top}]$
是頂端，因為虛擬碼從 $1$ 開始；兩邊的 $\mathit{top}$ 都是元素個數，翻譯時整段一起
改，不要逐行加減 $1$。

\subsection{應用：括號配對}
\label{sec:sq-brackets}

一串括號是\textbf{配對正確}的，意思是每個右括號都和它左邊最近、還沒配對的左括號
同一種。 「最近、還沒配對」正是堆疊頂端的意思：讀到左括號就 push，讀到右括號就
pop 一個出來比對。

\begin{pseudo}[label=alg:sq-brackets]{括號配對}
\Function{Brackets-Match}{$s, n$}
    \State 建立空堆疊 $S$
    \For{$i \gets 1$ \To $n$}
        \If{$s[i]$ 是左括號}
            \State \Call{Push}{$S, s[i]$}
        \ElsIf{$s[i]$ 是右括號}
            \If{\Call{Stack-Empty}{$S$} \OR \Call{Pop}{$S$} 和 $s[i]$ 不是同一種}
                \State \Return \FALSE
            \EndIf
        \EndIf
    \EndFor
    \State \Return \Call{Stack-Empty}{$S$}
\EndFunction
\end{pseudo}

有三種錯誤會被抓到：右括號來的時候堆疊是空的（右括號多了）、種類不合（例如
\texttt{[(])} 讀到 \texttt{]} 時頂端是 \texttt{(}），以及讀完之後堆疊不是空的（左
括號多了）。 最後一行很容易漏掉，漏掉之後 \texttt{((} 會被當成正確。 每個字元最多
push 一次、pop 一次，時間 $\Theta(n)$，堆疊最多放 $n$ 個。

\subsection{哪些輸出順序做得到}
\label{sec:sq-perm}

考卷上很常見的一種題目：$1, 2, \ldots, n$ 依序進入一個堆疊，中間可以隨時 pop，
pop 出來的數依序寫下，問哪些順序寫得出來。

例如 $n = 5$，順序 $3, 2, 5, 4, 1$ 做得到：push 1、push 2、push 3、pop（3）、
pop（2）、push 4、push 5、pop（5）、pop（4）、pop（1）。 順序 $4, 2, 3, 1, 5$ 做
不到：要先拿到 $4$，得先 push $1, 2, 3, 4$，pop 出 $4$ 之後頂端是 $3$，$2$ 被壓在
$3$ 下面，拿不到。

判斷的方法就是照著做一次：依序看要輸出的數，還沒 push 進來就一直 push，直到它進
來；它在頂端就 pop，不在頂端就表示做不到。 這個過程沒有選擇的餘地：要的數在頂端
時如果改去 push 別的數，它就被壓住了，而壓住它的數比它大、還沒輸出，之後一定得先
pop 掉，順序就不對了。 所以一個做得到的輸出順序，對應唯一的一串 push、pop。

不用模擬也能判斷，看有沒有下面這種「大、小、中」的樣子。

\begin{theorem}
\label{thm:sq-312}
$1, 2, \ldots, n$ 依序進入堆疊。 輸出順序 $p_1, p_2, \ldots, p_n$ 做得到的充分必
要條件是：不存在 $i < j < k$ 使得 $p_j < p_k < p_i$。
\end{theorem}

\begin{proof}
必要：假設有這樣的 $i < j < k$。 $p_i$ 被 pop 的時候，所有比 $p_i$ 小的數都已經
push 過了，其中 $p_j$、$p_k$ 都還沒輸出，所以兩個都在堆疊裡。 $p_j < p_k$，$p_j$
比較早 push，在 $p_k$ 下面，必須等 $p_k$ 出來之後才能出來，和 $j < k$ 矛盾。

充分：假設沒有這種三元組，照上面的方法模擬。 唯一會失敗的情況是：輪到 $p_j$ 時，
它已經在堆疊裡，但不在頂端，頂端是某個 $x > p_j$（堆疊由底到頂遞增，因為是依序
push 的）。 $x$ 還沒輸出，所以 $x = p_k$，$k > j$。 處理 $p_j$ 時不會再 push
（$p_j$ 已經在堆疊裡），所以 $x$ 是在處理更早的某個輸出 $p_i$（$i < j$）時 push
的；那一步 push 到 $p_i$ 為止，所以 $x \le p_i$。 $p_i$ 已經輸出、$x$ 還在堆疊裡，
所以 $x < p_i$。 這樣 $p_j < p_k < p_i$，$i < j < k$，和假設矛盾。
\end{proof}

用這個定理檢查 $4, 2, 3, 1, 5$：$4$ 先、$2$ 在中間、$3$ 在後，$2 < 3 < 4$，是
「大、小、中」，做不到。

能做出幾種順序？ $n = 3$ 的六種排列裡，只有 $3, 1, 2$ 做不到，它本身就是「大、
小、中」。 $n = 4$ 的 $24$ 種裡有 $10$ 種做不到：
\[
1423,\ 2413,\ 3124,\ 3142,\ 3412,\ 4123,\ 4132,\ 4213,\ 4231,\ 4312,
\]
每一種都含「大、小、中」，例如 $2413$ 的 $4, 1, 3$。 其餘 $14$ 種做得到。 把
$n = 1, 2, \ldots$ 做得到的個數列出來是 $1, 2, 5, 14, 42, 132, \ldots$，這是
\textbf{卡特蘭數}（Catalan number）
\[
C_n = \frac{1}{n+1}\binom{2n}{n}.
\]

\begin{theorem}
\label{thm:sq-catalan}
$1, 2, \ldots, n$ 依序進入堆疊，做得到的輸出順序恰好有 $C_n$ 種。
\end{theorem}

\begin{whybox}[卡特蘭數]
前面說過，做得到的輸出順序和「合法的 push、pop 序列」一一對應。 一串 push、pop
是合法的，條件是：共 $n$ 個 push、$n$ 個 pop，而且每一個前綴裡 pop 的個數都不超過
push 的個數（不會對空堆疊 pop）。 所以只要數合法的序列。

$n$ 個 push、$n$ 個 pop 任意排，共 $\binom{2n}{n}$ 種。 數不合法的：不合法的序列
一定有第一個「pop 比 push 多一個」的位置。 把這個位置\textbf{之後}的每一步對調
（push 換成 pop，pop 換成 push）。 那個位置之前 push 比 pop 少一個，之後原本 push、
pop 各剩 $n - t$ 和 $n - t - 1$ 個（$t$ 是前面 push 的個數），對調後總共變成
$n - 1$ 個 push、$n + 1$ 個 pop。 反過來，任何 $n - 1$ 個 push、$n + 1$ 個 pop 的
序列，pop 總數比較多，一定有第一個「pop 比 push 多一個」的位置，把之後對調回去，
就得到一個不合法的原序列。 兩邊一一對應，所以不合法的有 $\binom{2n}{n-1}$ 種。

合法的個數是
\[
\binom{2n}{n} - \binom{2n}{n-1}
= \binom{2n}{n} - \frac{n}{n+1}\binom{2n}{n}
= \frac{1}{n+1}\binom{2n}{n}.
\]
這裡用了 $\binom{2n}{n-1} = \binom{2n}{n} \cdot \frac{n}{n+1}$，把兩個組合數寫成階
乘就看得出來。
\end{whybox}

另一個角度：看 $1$ 是第幾個輸出的。 如果 $1$ 是第 $k + 1$ 個，那麼在它之前輸出的
恰好是 $2, \ldots, k + 1$（它們都在 $1$ 上面進出），之後輸出的是
$k + 2, \ldots, n$。 兩段各自是一個小一點的堆疊問題，所以
$C_n = \sum_{k=0}^{n-1} C_k\, C_{n-1-k}$，$C_0 = 1$。 這個遞迴式在第~7~章數二元樹
的形狀時會再出現。

\begin{examnote}
判斷一個輸出順序做不做得到，答案寫「做不到」時，要指出卡住的那一步：「輸出 $4$
之後頂端是 $3$，$2$ 在 $3$ 下面」，或指出一組 $p_j < p_k < p_i$。 寫「做得到」時，
寫出完整的 push、pop 序列。 只寫是或否，閱卷的人看不出你是推的還是猜的。
\end{examnote}

\section{佇列：環狀陣列與串列}
\label{sec:sq-queue}

\textbf{佇列}也是一串資料，但從一端加入、從另一端取出。 加入的一端叫\textbf{尾端}
（rear），取出的一端叫\textbf{前端}（front）。 操作：

\begin{itemize}
\item \proc{Enqueue}$(Q, x)$：把 $x$ 接在尾端。
\item \proc{Dequeue}$(Q)$：取出前端的元素並回傳。
\item \proc{Front}$(Q)$：回傳前端的元素，不取出。
\item \proc{Queue-Empty}$(Q)$：佇列是不是空的。
\end{itemize}

最先加入的最先取出，所以佇列是\textbf{先進先出}（first in, first out，FIFO）的結
構。 排隊、印表機的工作、第~21~章廣度優先搜尋裡「等著被處理的點」，都是佇列。

\subsection{直接用陣列的問題}

最直接的想法：用 $Q[1..n]$ 存，元素放在 $Q[\mathit{front}..\mathit{rear}]$。 加
入時 $\mathit{rear}$ 加一，取出時 $\mathit{front}$ 加一。 問題是兩個索引都只往右
走，取出的元素留下的空位再也用不到。 做了 $n$ 次加入、$n - 1$ 次取出之後，佇列只
剩一個元素，$\mathit{rear}$ 卻已經走到 $n$，不能再加入了。 這種「陣列還有空位，
卻說滿了」的情況，有些書叫假溢位。

一種修法是每次取出後把剩下的元素都往前搬一格，但這讓 \proc{Dequeue} 變成
$\Theta(n)$。 比較好的修法是讓索引走到最後一格之後，繞回第一格。

\subsection{環狀陣列}
\label{sec:sq-circular}

把陣列 $Q[0..n-1]$ 想成頭尾相接的一圈：$Q[n-1]$ 的下一格是 $Q[0]$。 往後走一格就
是索引加一再對 $n$ 取餘數，$i \gets (i + 1) \bmod n$。 這一節的陣列從 $0$ 開始，取
餘數比較方便，C 程式也一樣，翻譯時不用改編號。

用兩個索引：
\begin{itemize}
\item $\mathit{front}$：前端元素的位置；
\item $\mathit{rear}$：下一個要放的位置（尾端元素的下一格）。
\end{itemize}
佇列的元素依序在 $Q[\mathit{front}], Q[\mathit{front} + 1], \ldots$，一路到
$Q[\mathit{rear} - 1]$，中間可能繞過陣列的結尾。

\textbf{空和滿分不出來。} 一開始 $\mathit{front} = \mathit{rear} = 0$，佇列是空
的，所以「$\mathit{front} = \mathit{rear}$」表示空。 但如果一直加入，$n$ 次之後
$\mathit{rear}$ 繞了一圈回到 $0$，又是 $\mathit{front} = \mathit{rear}$，這時佇列
是滿的。 只看這兩個索引，分不出空和滿。 這不是寫法的問題：固定 $\mathit{front}$
時，$\mathit{rear}$ 只有 $n$ 種可能，佇列的元素個數卻有 $0, 1, \ldots, n$ 共
$n + 1$ 種，$n$ 個值不夠表示 $n + 1$ 種狀態（觀念題第~\ref{q:sq-pigeon}~題）。 常見的解法有三種：

\begin{enumerate}
\item \textbf{留一格不用}：最多只放 $n - 1$ 個。 空是 $\mathit{front} =
\mathit{rear}$，滿是 $(\mathit{rear} + 1) \bmod n = \mathit{front}$，也就是尾端再
往後一格就撞到前端。
\item \textbf{多記一個個數} $\mathit{count}$：空是 $\mathit{count} = 0$，滿是
$\mathit{count} = n$，$n$ 格都能用。 這時 $\mathit{rear}$ 可以不存，由
$(\mathit{front} + \mathit{count}) \bmod n$ 算出來。
\item \textbf{多記一個旗標} $\mathit{tag}$：最近一次操作是加入就設成 $1$，是取出就
設成 $0$。 $\mathit{front} = \mathit{rear}$ 時，$\mathit{tag} = 1$ 表示滿（剛才
加入才變相等的），$\mathit{tag} = 0$ 表示空。
\end{enumerate}

下面的虛擬碼用第一種。 用第一種時，佇列裡的元素個數是
\[
(\mathit{rear} - \mathit{front} + n) \bmod n.
\]
$\mathit{rear} \ge \mathit{front}$ 時就是 $\mathit{rear} - \mathit{front}$；
$\mathit{rear} < \mathit{front}$ 表示繞過了結尾，要補上 $n$。

\begin{pseudo}[label=alg:sq-circular]{環狀佇列（留一格不用）}
\Procedure{Enqueue}{$Q, x$}
    \If{$(Q.\mathit{rear} + 1) \bmod Q.n = Q.\mathit{front}$}
        \State \Error{overflow}
    \EndIf
    \State $Q.\mathit{item}[Q.\mathit{rear}] \gets x$
    \State $Q.\mathit{rear} \gets (Q.\mathit{rear} + 1) \bmod Q.n$
\EndProcedure
\Statex
\Function{Dequeue}{$Q$}
    \If{$Q.\mathit{front} = Q.\mathit{rear}$}
        \State \Error{underflow}
    \EndIf
    \State $x \gets Q.\mathit{item}[Q.\mathit{front}]$
    \State $Q.\mathit{front} \gets (Q.\mathit{front} + 1) \bmod Q.n$
    \State \Return $x$
\EndFunction
\end{pseudo}

\subsection{手算追蹤}

用 $n = 5$（最多放 $4$ 個），從 $\mathit{front} = \mathit{rear} = 0$ 開始。 下表
「E 31」表示 \proc{Enqueue}$(Q, 31)$，「D」表示 \proc{Dequeue}$(Q)$。 陣列那五欄
寫的是實際存在陣列裡的值，灰字是已經取出、但還沒被蓋掉的值。

\begin{center}
\begin{tabular}{llccccccc}
\toprule
操作 & 結果 & $Q[0]$ & $Q[1]$ & $Q[2]$ & $Q[3]$ & $Q[4]$ & $\mathit{front}$ & $\mathit{rear}$ \\
\midrule
E 31 & & 31 & & & & & 0 & 1 \\
E 42 & & 31 & 42 & & & & 0 & 2 \\
E 53 & & 31 & 42 & 53 & & & 0 & 3 \\
D & 取出 31 & \textcolor{black!40}{31} & 42 & 53 & & & 1 & 3 \\
E 64 & & \textcolor{black!40}{31} & 42 & 53 & 64 & & 1 & 4 \\
E 75 & & \textcolor{black!40}{31} & 42 & 53 & 64 & 75 & 1 & 0 \\
E 86 & 滿了 & \textcolor{black!40}{31} & 42 & 53 & 64 & 75 & 1 & 0 \\
D & 取出 42 & \textcolor{black!40}{31} & \textcolor{black!40}{42} & 53 & 64 & 75 & 2 & 0 \\
D & 取出 53 & \textcolor{black!40}{31} & \textcolor{black!40}{42} & \textcolor{black!40}{53} & 64 & 75 & 3 & 0 \\
E 97 & & 97 & \textcolor{black!40}{42} & \textcolor{black!40}{53} & 64 & 75 & 3 & 1 \\
D & 取出 64 & 97 & \textcolor{black!40}{42} & \textcolor{black!40}{53} & \textcolor{black!40}{64} & 75 & 4 & 1 \\
D & 取出 75 & 97 & \textcolor{black!40}{42} & \textcolor{black!40}{53} & \textcolor{black!40}{64} & \textcolor{black!40}{75} & 0 & 1 \\
\bottomrule
\end{tabular}
\end{center}

\proc{Enqueue}$(Q, 75)$ 把 $75$ 放進 $Q[4]$，$\mathit{rear}$ 從 $4$ 繞回 $0$。 接
著加入 $86$ 時，$(0 + 1) \bmod 5 = 1 = \mathit{front}$，判定為滿，雖然 $Q[0]$ 其實
空著，這就是留下不用的那一格。 下圖是這個時候的樣子。

\begin{center}
\begin{tikzpicture}
\foreach \v/\k in {{}/0, 42/1, 53/2, 64/3, 75/4} {
    \pgfmathsetmacro{\ang}{90 - 72*\k}
    \ifnum\k=0
        \node[cell] (q\k) at (\ang:1.7) {};
    \else
        \node[sqfill] (q\k) at (\ang:1.7) {\v};
    \fi
    \node[idxlbl] at (\ang:2.65) {$Q[\k]$};
}
\draw[diagarrow, draw=black!40] (60:0.75) arc (60:-240:0.75);
\node[font=\small, text=black!55] at (0,0) {方向};
\node[font=\small, text=memStroke] (f) at (3.9, 0.1) {$\mathit{front} = 1$};
\draw[diagarrow, draw=memStroke] (f.west) -- (q1.east);
\node[font=\small, text=optStroke] (r) at (-3.6, 1.6) {$\mathit{rear} = 0$};
\draw[diagarrow, draw=optStroke] (r.east) -- (q0.west);
\end{tikzpicture}
\diagcaption{$n = 5$ 的環狀佇列，存了 $42, 53, 64, 75$ 四個元素（藍底）。
$\mathit{rear}$ 指著下一個要放的 $Q[0]$，再往後一格就是 $\mathit{front}$，所以判
定為滿。 元素沿順時針方向排列。}
\label{fig:sq-ring}
\end{center}

\begin{examnote}
環狀佇列的題目，先看題目怎麼定義 $\mathit{front}$ 和 $\mathit{rear}$。 有的書讓
$\mathit{front}$ 指向前端元素的\textbf{前一格}、$\mathit{rear}$ 指向尾端元素本身；
這時空和滿的條件、元素個數的公式都要跟著改。 作答時先寫一句「依題意，
$\mathit{front}$ 指向……」，再照那個定義算，每做一個操作寫一列
$\mathit{front}$、$\mathit{rear}$ 和陣列內容。
\end{examnote}

\subsection{C 程式}

C 的 \texttt{\%} 在左邊是負數時，結果也是負的（例如 \texttt{-1 \% 5} 是
\texttt{-1}），所以算個數時寫 \texttt{(rear - front + n) \% n}，先加 $n$ 再取餘
數。 下面是虛擬碼~\ref{alg:sq-circular} 翻成的 C 函式，陣列用 \texttt{malloc}
配置 \texttt{n} 格。 把上面那張表的操作依序丟給它們，每一步的 \texttt{front}、
\texttt{rear} 都和表上一樣；例如做完 \texttt{E 75} 時 $\mathit{rear} = 0 <
\mathit{front} = 1$，\texttt{q\_size} 算出 $(0 - 1 + 5) \bmod 5 = 4$。

\begin{code}[label=code:sq-circ]{c}{環狀佇列（留一格不用）}
// q[0..n-1] 當成一圈。front 是前端元素的位置，rear 是下一個要放的位置。
// 留一格不用，最多放 n - 1 個。
typedef struct {
    int *q;
    int n, front, rear;
} Queue;

int q_size(const Queue *Q)
{
    return (Q->rear - Q->front + Q->n) % Q->n;  // 先加 n，避免負數
}

int enqueue(Queue *Q, int x)                    // 滿了回傳 0
{
    if ((Q->rear + 1) % Q->n == Q->front) {
        return 0;
    }
    Q->q[Q->rear] = x;
    Q->rear = (Q->rear + 1) % Q->n;
    return 1;
}

int dequeue(Queue *Q, int *x)                   // 空的回傳 0
{
    if (Q->front == Q->rear) {
        return 0;
    }
    *x = Q->q[Q->front];
    Q->front = (Q->front + 1) % Q->n;
    return 1;
}
\end{code}

\subsection{用串列存}

用串列存佇列，要記兩個指標：$Q.\mathit{front}$ 指向第一個節點，
$Q.\mathit{rear}$ 指向最後一個節點。 加入接在 $\mathit{rear}$ 後面，取出拿
$\mathit{front}$。 方向不能反過來：在單向串列的尾巴刪除要先找到前一個節點，是
$\Theta(n)$；在尾巴加入只要改 $\mathit{rear}.\mathit{next}$，是 $\Theta(1)$。

\begin{pseudo}[label=alg:sq-list-queue]{串列佇列}
\Procedure{Enqueue-List}{$Q, x$}
    \State 配置一個新節點 $p$
    \State $p.\mathit{key} \gets x$
    \State $p.\mathit{next} \gets \NIL$
    \If{$Q.\mathit{rear} = \NIL$} \Comment{原本是空的}
        \State $Q.\mathit{front} \gets p$
    \Else
        \State $Q.\mathit{rear}.\mathit{next} \gets p$
    \EndIf
    \State $Q.\mathit{rear} \gets p$
\EndProcedure
\Statex
\Function{Dequeue-List}{$Q$}
    \If{$Q.\mathit{front} = \NIL$}
        \State \Error{underflow}
    \EndIf
    \State $p \gets Q.\mathit{front}$
    \State $x \gets p.\mathit{key}$
    \State $Q.\mathit{front} \gets p.\mathit{next}$
    \If{$Q.\mathit{front} = \NIL$} \Comment{取走的是最後一個}
        \State $Q.\mathit{rear} \gets \NIL$
    \EndIf
    \State 釋放 $p$
    \State \Return $x$
\EndFunction
\end{pseudo}

兩個 \textbf{if} 處理的都是空佇列的邊界：加入第一個節點時 $\mathit{front}$ 也要指
向它；取走最後一個節點時 $\mathit{rear}$ 也要清成 \NIL，不然 $\mathit{rear}$ 會指
向已經釋放的節點，下一次加入就寫進釋放過的記憶體。

\textbf{例題：印表機排程。} 系辦的印表機一次印一份文件，每頁要 $4$ 秒，大家送出
的文件依送出的先後排隊。 輸入每份文件送出的時間和頁數（已經依時間排好），算出每
份何時開始印、何時印完、等了多久。 印表機印完一份時，把這段時間內送來的文件都加
入佇列，再從前端取出下一份；佇列是空的就閒著，等下一份送來。

\begin{example}[label=ex:sq-printer]{c}{(printer\_queue.c) 用串列佇列模擬印表機}
#include <stdio.h>
#include <stdlib.h>

#define MAXJOBS 100
#define SEC_PER_PAGE 4      // 印一頁要 4 秒

typedef struct QNode {
    int job;                // 文件的編號
    struct QNode *next;
} QNode;

typedef struct {
    QNode *front, *rear;    // 從 front 拿，往 rear 後面接
    int size;
} Queue;

int enqueue(Queue *Q, int job)
{
    QNode *p = malloc(sizeof *p);
    if (p == NULL) {
        return 0;
    }
    p->job = job;
    p->next = NULL;
    if (Q->rear == NULL) {  // 原本是空的：新節點同時是頭和尾
        Q->front = p;
    } else {
        Q->rear->next = p;
    }
    Q->rear = p;
    Q->size++;
    return 1;
}

int dequeue(Queue *Q)       // 呼叫前要確定佇列不是空的
{
    QNode *p = Q->front;
    int job = p->job;
    Q->front = p->next;
    if (Q->front == NULL) { // 拿走最後一個：尾巴也要清掉
        Q->rear = NULL;
    }
    free(p);
    Q->size--;
    return job;
}

int main(void)
{
    int n, arrive[MAXJOBS], pages[MAXJOBS];
    if (scanf("%d", &n) != 1 || n < 1 || n > MAXJOBS) {
        return 1;
    }
    for (int i = 0; i < n; i++) {   // 依送出的時間排好
        if (scanf("%d %d", &arrive[i], &pages[i]) != 2) {
            return 1;
        }
    }
    Queue Q = {NULL, NULL, 0};
    int next = 0, now = 0, longest = 0, total_wait = 0;
    printf("job arrive start finish wait\n");
    while (next < n || Q.size > 0) {
        if (Q.size == 0 && now < arrive[next]) {
            now = arrive[next];     // 印表機閒著，等下一份送來
        }
        while (next < n && arrive[next] <= now) {
            if (!enqueue(&Q, next++)) {
                return 1;
            }
        }
        if (Q.size > longest) {
            longest = Q.size;
        }
        int j = dequeue(&Q);
        int finish = now + pages[j] * SEC_PER_PAGE;
        printf("%d %d %d %d %d\n", j + 1, arrive[j], now, finish, now - arrive[j]);
        total_wait += now - arrive[j];
        now = finish;
    }
    printf("average wait = %.2f, longest queue = %d\n",
           (double)total_wait / n, longest);
    return 0;
}
\end{example}
\iosample{6 \\ 0 5 \\ 3 2 \\ 6 10 \\ 9 1 \\ 50 3 \\ 58 2}{job arrive start finish wait \\ 1 0 0 20 0 \\ 2 3 20 28 17 \\ 3 6 28 68 22 \\ 4 9 68 72 59 \\ 5 50 72 84 22 \\ 6 58 84 92 26 \\ average wait = 24.33, longest queue = 3}

手算對一次：第 1 份在 $0$ 秒開始，$5$ 頁印到 $20$ 秒。 這段時間第 2、3、4 份陸續
送來，$20$ 秒時佇列是 $\langle 2, 3, 4 \rangle$，這是排得最長的時候，所以
\texttt{longest queue = 3}。 取出第 2 份印到 $28$，再取出第 3 份，$10$ 頁印到
$68$。 第 4 份只有 $1$ 頁，卻排在 $10$ 頁的那份後面，等了 $59$ 秒。

先來先服務很公平，但不一定讓平均等待最短。 如果每次挑頁數最少的先印，第 4 份就不
用等那麼久；這時要的是「每次取出最小的」，那是第~8~章的優先佇列，第~19~章會證明
這樣排能讓平均等待時間最小。

程式結束前佇列一定是空的（迴圈的條件），所以每個 \texttt{malloc} 出來的節點都在
\texttt{dequeue} 裡釋放了。 佇列不保證空的程式，結束前要像範例~\ref{ex:sq-list-stack}
最後的迴圈那樣，把剩下的節點都取出來。

\subsection{雙端佇列}
\label{sec:sq-deque}

\textbf{雙端佇列}（double-ended queue，常簡稱 deque）兩端都可以加入、也都可以取出，
共四個操作：前端加入、前端取出、尾端加入、尾端取出。 只用一端，它就是堆疊；一端
加、另一端取，它就是佇列。

環狀陣列可以直接改成雙端佇列。 尾端加入、前端取出和佇列一樣；前端加入是
$\mathit{front}$ 先往回退一格再放，尾端取出是 $\mathit{rear}$ 先往回退一格再拿。
往回退一格是 $(i - 1 + n) \bmod n$，加 $n$ 的理由和前面一樣。 例如前端加入：
先檢查滿了沒有，再做 $\mathit{front} \gets (\mathit{front} - 1 + n) \bmod n$、
$Q[\mathit{front}] \gets x$；尾端取出：先檢查空了沒有，再做
$\mathit{rear} \gets (\mathit{rear} - 1 + n) \bmod n$，回傳 $Q[\mathit{rear}]$。

用串列做的話要用雙向串列：尾端取出需要找到倒數第二個節點，單向串列做不到
$\Theta(1)$。

考卷上還常出現兩種受限的雙端佇列。 \textbf{輸入受限}（input-restricted）的只能
從一端加入、兩端都能取出；\textbf{輸出受限}（output-restricted）的兩端都能加入、
只能從一端取出。 題目通常像第~\ref{sec:sq-perm}~節那樣，讓 $1, 2, \ldots, n$ 依
序進入，問哪些輸出順序做得到（觀念題第~\ref{q:sq-deque}~題）。 解法也一樣：照著要的輸出順序走，
每一步看要的數在不在可以取出的那一端。

\section{運算式：中序、前序、後序}
\label{sec:sq-expr}

\subsection{三種寫法}

平常寫的 $a + b$ 把運算子放在兩個運算元中間，叫\textbf{中序}（infix）。 把運算
子放在前面的 $+\;a\;b$ 叫\textbf{前序}（prefix，也叫波蘭表示法），放在後面的
$a\;b\;+$ 叫\textbf{後序}（postfix，也叫逆波蘭表示法）。

中序要靠\textbf{優先順序}（precedence）和\textbf{結合性}（associativity）才能讀出
唯一的意思：$a - b * c$ 先算乘法，$a - b - c$ 是 $(a - b) - c$，從左往右結合；
次方通常從右往左結合，$a \sqpow b \sqpow c$ 是 $a^{(b^c)}$。 想改變順序就要加括
號。 前序和後序完全不需要括號，也不需要優先順序：運算子的位置本身就決定了它作用
在誰身上。

手算轉換最穩的方法是\textbf{加滿括號}：先依優先順序和結合性，把每個運算加上一對
括號，再把每個運算子搬到它那對括號的右括號（後序）或左括號（前序）的位置，最後
把括號全部擦掉。 例如 $a * b + c / d$：

\begin{center}
\begin{tabular}{ll}
\toprule
步驟 & 結果 \\
\midrule
加滿括號 & $((a * b) + (c / d))$ \\
運算子移到右括號，擦掉括號（後序）& $a\;b\;*\;c\;d\;/\;+$ \\
運算子移到左括號，擦掉括號（前序）& $+\;*\;a\;b\;/\;c\;d$ \\
\bottomrule
\end{tabular}
\end{center}

三種寫法裡，運算元出現的順序都一樣（$a, b, c, d$），只有運算子的位置不同。 這是
檢查答案的好方法：轉完之後，運算元的順序變了就一定錯了。 運算式和二元樹的關係
（中序、前序、後序剛好是同一棵樹的三種走訪）在第~7~章再講。

\subsection{後序求值}
\label{sec:sq-postfix-eval}

後序由左往右讀就能算：遇到運算元就 push；遇到運算子就 pop 兩個數出來算，結果
push 回去。 讀完時堆疊裡剩下的一個數就是答案。 要注意 pop 的順序：\textbf{先 pop
出來的是右邊的運算元}。 $8\;3\;-$ 是 $8 - 3$：讀到 $-$ 時頂端是 $3$，先 pop 出
$3$ 當右邊，再 pop 出 $8$ 當左邊。

\begin{pseudo}[label=alg:sq-postfix-eval]{後序求值}
\Function{Postfix-Eval}{$T, m$} \Comment{$T[1..m]$ 是記號}
    \State 建立空堆疊 $S$
    \For{$i \gets 1$ \To $m$}
        \If{$T[i]$ 是運算元}
            \State \Call{Push}{$S, T[i]$}
        \Else
            \State $b \gets$ \Call{Pop}{$S$} \Comment{右邊的運算元}
            \State $a \gets$ \Call{Pop}{$S$} \Comment{左邊的運算元}
            \State \Call{Push}{$S$, $a$ $T[i]$ $b$ 的值}
        \EndIf
    \EndFor
    \State \Return \Call{Pop}{$S$}
\EndFunction
\end{pseudo}

用 $8\;2\;5\;*\;+\;1\;3\;2\;*\;+\;4\;-\;/$ 走一次，堆疊由底到頂寫：

\begin{center}
\begin{tabular}{clll}
\toprule
記號 & 動作 & 堆疊（底 $\to$ 頂）\\
\midrule
8, 2, 5 & 依序 push & 8 2 5 \\
$*$ & pop 5、2，push $2 * 5 = 10$ & 8 10 \\
$+$ & pop 10、8，push $8 + 10 = 18$ & 18 \\
1, 3, 2 & 依序 push & 18 1 3 2 \\
$*$ & pop 2、3，push $3 * 2 = 6$ & 18 1 6 \\
$+$ & pop 6、1，push $1 + 6 = 7$ & 18 7 \\
4 & push & 18 7 4 \\
$-$ & pop 4、7，push $7 - 4 = 3$ & 18 3 \\
$/$ & pop 3、18，push $18 / 3 = 6$ & 6 \\
\bottomrule
\end{tabular}
\end{center}

答案是 $6$，對應的中序是 $(8 + 2 * 5) / (1 + 3 * 2 - 4)$。 每個記號只處理一次，
每次 $\Theta(1)$，總時間 $\Theta(m)$。

\subsection{中序轉後序}
\label{sec:sq-infix-postfix}

後序好算，但人寫的是中序，所以要先轉換。 轉換時運算元的順序不變，直接輸出；運算
子則要等一下，因為讀到 $a - b$ 的 $-$ 時，還不知道右邊是不是 $b * c$ 這種要先算
的東西。 等待中的運算子放在堆疊裡。 規則如下，由左往右讀每個記號：

\begin{enumerate}
\item 運算元：直接輸出。
\item 左括號：push。 它像一道牆，牆下面的運算子要等右括號出現才處理。
\item 右括號：一直 pop 並輸出，直到遇到左括號；左括號 pop 掉，不輸出。
\item 運算子 $\mathit{op}$：只要頂端是運算子、而且它\textbf{應該比 $\mathit{op}$
先算}，就 pop 並輸出；最後把 $\mathit{op}$ push。 「頂端的 $t$ 應該比
$\mathit{op}$ 先算」的意思是：$t$ 的優先順序比較高，或兩者一樣高而且是由左往右
結合。
\item 讀完之後，把堆疊裡剩下的運算子依序 pop 並輸出。
\end{enumerate}

本書用的優先順序（越上面越先算）：

\begin{center}
\begin{tabular}{ccc}
\toprule
運算子 & 優先順序 & 結合性 \\
\midrule
\sqpow & 3 & 由右往左 \\
$*$、$/$ & 2 & 由左往右 \\
$+$、$-$ & 1 & 由左往右 \\
\bottomrule
\end{tabular}
\end{center}

第 4 條規則裡的「一樣高」決定了結合性。 讀 $a - b - c$ 的第二個 $-$ 時，頂端是第
一個 $-$，一樣高、由左往右結合，所以先輸出第一個 $-$，得到 $a\;b\;-\;c\;-$，也就是
$(a - b) - c$。 讀 $a \sqpow b \sqpow c$ 的第二個 \sqpow{} 時，一樣高但由右往左結
合，不 pop，直接 push，最後得到 $a\;b\;c\;\sqpow\;\sqpow$，也就是
$a \sqpow (b \sqpow c)$。

\begin{pseudo}[label=alg:sq-infix-postfix]{中序轉後序}
\Procedure{Infix-To-Postfix}{$T, m$}
    \State 建立空堆疊 $S$
    \For{$i \gets 1$ \To $m$}
        \If{$T[i]$ 是運算元}
            \State 輸出 $T[i]$
        \ElsIf{$T[i]$ 是左括號}
            \State \Call{Push}{$S, T[i]$}
        \ElsIf{$T[i]$ 是右括號}
            \While{\Call{Top}{$S$} 不是左括號}
                \State 輸出 \Call{Pop}{$S$}
            \EndWhile
            \State \Call{Pop}{$S$} \Comment{丟掉左括號}
        \Else \Comment{$T[i]$ 是運算子}
            \While{$S$ 不是空的 \AND \Call{Top}{$S$} 是運算子 \AND \Call{Top}{$S$} 應該比 $T[i]$ 先算}
                \State 輸出 \Call{Pop}{$S$}
            \EndWhile
            \State \Call{Push}{$S, T[i]$}
        \EndIf
    \EndFor
    \While{$S$ 不是空的}
        \State 輸出 \Call{Pop}{$S$}
    \EndWhile
\EndProcedure
\end{pseudo}

每個運算子最多 push 一次、pop 一次，每個括號也一樣，所以總時間是 $\Theta(m)$。 內
層的 \textbf{while} 有時一次 pop 好幾個，但全部加起來不會超過 push 的次數，這是
第~3~章聚合法的典型例子。

\textbf{手算追蹤}：轉換 \texttt{A - B * (C + D \sqpow{} E \sqpow{} F) / G + H}。
每讀一個記號寫一列，堆疊由底到頂寫。

\begin{center}
\small
\begin{tabular}{clll}
\toprule
記號 & 動作 & 堆疊（底 $\to$ 頂）& 輸出 \\
\midrule
A & 輸出 & & A \\
$-$ & push & $-$ & A \\
B & 輸出 & $-$ & A B \\
$*$ & $*$ 比 $-$ 高，push & $-\;*$ & A B \\
( & push & $-\;*\;($ & A B \\
C & 輸出 & $-\;*\;($ & A B C \\
$+$ & 頂端是 (，push & $-\;*\;(\;+$ & A B C \\
D & 輸出 & $-\;*\;(\;+$ & A B C D \\
\sqpow & \sqpow{} 比 $+$ 高，push & $-\;*\;(\;+\;\sqpow$ & A B C D \\
E & 輸出 & $-\;*\;(\;+\;\sqpow$ & A B C D E \\
\sqpow & 一樣高但由右往左，push & $-\;*\;(\;+\;\sqpow\;\sqpow$ & A B C D E \\
F & 輸出 & $-\;*\;(\;+\;\sqpow\;\sqpow$ & A B C D E F \\
) & pop \sqpow{} \sqpow{} $+$，丟掉 ( & $-\;*$ & A B C D E F \sqpow{} \sqpow{} $+$ \\
$/$ & pop $*$（一樣高、由左往右），push & $-\;/$ & A B C D E F \sqpow{} \sqpow{} $+$ $*$ \\
G & 輸出 & $-\;/$ & \ldots{} $+$ $*$ G \\
$+$ & pop $/$、$-$，push & $+$ & \ldots{} $+$ $*$ G $/$ $-$ \\
H & 輸出 & $+$ & \ldots{} $+$ $*$ G $/$ $-$ H \\
結束 & pop $+$ & & \ldots{} $+$ $*$ G $/$ $-$ H $+$ \\
\bottomrule
\end{tabular}
\end{center}

結果是 \texttt{A B C D E F \sqpow{} \sqpow{} + * G / - H +}。 檢查：運算元的順序
\texttt{A} 到 \texttt{H} 沒變；\texttt{D E F \sqpow{} \sqpow} 表示
$D^{(E^F)}$，次方由右往左結合；$*$ 在 $/$ 前面，表示 $(B * (\ldots)) / G$，乘除
由左往右結合。

\begin{examnote}
中序轉後序的追蹤題，表格固定寫四欄：讀到的記號、做了什麼、堆疊（寫明由底到頂或由
頂到底）、目前為止的輸出。 遇到同優先順序時，在「動作」欄寫出你依據的結合性。 最
後一列別忘了把堆疊清空。 寫完用「運算元順序沒變」檢查一次。
\end{examnote}

很多教科書和考卷把第 4 條規則寫成兩張數字表：運算子在堆疊裡時的優先順序
（in-stack priority，isp），和剛讀進來時的優先順序（incoming priority，icp）。 規
則變成「頂端的 isp 大於等於讀進來的 icp 就 pop」。 下面這組數字和本書的規則等價：

\begin{center}
\begin{tabular}{ccccc}
\toprule
 & ( & \sqpow & $*$ $/$ & $+$ $-$ \\
\midrule
isp（在堆疊裡）& 0 & 3 & 2 & 1 \\
icp（剛讀進來）& 5 & 4 & 2 & 1 \\
\bottomrule
\end{tabular}
\end{center}

左括號的 icp 最高，所以一定會被 push；進了堆疊之後 isp 最低，所以不會被任何運算子
pop 掉。 \sqpow{} 的 icp 比 isp 大，所以新讀到的 \sqpow{} 不會把堆疊裡的 \sqpow{}
pop 出來，這就是由右往左結合。 $*$、$/$、$+$、$-$ 的 isp 和 icp 相等，一樣高時會
pop，這就是由左往右結合。

\subsection{C 程式}

下面的程式每讀一行中序運算式，就印出後序，再用後序算出值。 為了專心在堆疊上，運
算元只有一位數（\texttt{0} 到 \texttt{9}），空白略過，而且假設輸入的運算式都是合
法的。 運算子有 \texttt{+ - * /} 和 \texttt{\sqpow}（次方），除法和 C 一樣捨去小
數。 兩個堆疊都開成和一行一樣長，不會溢位。

\begingroup\catcode`\^=12
\begin{example}[label=ex:sq-infix]{c}{(infix\_postfix.c) 中序轉後序並求值}
#include <stdio.h>
#include <string.h>
#include <ctype.h>

#define MAXLEN 256

int prec(char op)               // 優先順序，數字越大越先算
{
    switch (op) {
    case '^': return 3;
    case '*': case '/': return 2;
    case '+': case '-': return 1;
    default:  return 0;
    }
}

// 中序 in 轉成後序 out。運算元是一位數，空白略過。
void to_postfix(const char *in, char *out)
{
    char st[MAXLEN];            // 運算子堆疊
    int top = 0, k = 0;
    for (int i = 0; in[i] != '\0'; i++) {
        char c = in[i];
        if (isdigit((unsigned char)c)) {
            out[k++] = c;       // 運算元直接輸出
        } else if (c == '(') {
            st[top++] = c;
        } else if (c == ')') {
            while (st[top - 1] != '(') {
                out[k++] = st[--top];
            }
            top--;              // 丟掉 '('
        } else if (strchr("+-*/^", c) != NULL) {
            // 頂端的運算子應該先算的話，先輸出它；^ 由右往左結合，一樣高時不 pop
            while (top > 0 && st[top - 1] != '(' &&
                   (prec(st[top - 1]) > prec(c) ||
                    (prec(st[top - 1]) == prec(c) && c != '^'))) {
                out[k++] = st[--top];
            }
            st[top++] = c;
        }
    }
    while (top > 0) {
        out[k++] = st[--top];
    }
    out[k] = '\0';
}

long long eval_postfix(const char *p)
{
    long long st[MAXLEN];
    int top = 0;
    for (int i = 0; p[i] != '\0'; i++) {
        if (isdigit((unsigned char)p[i])) {
            st[top++] = p[i] - '0';
            continue;
        }
        long long b = st[--top];    // 先拿出來的是右邊的運算元
        long long a = st[--top];
        long long r = 1;
        switch (p[i]) {
        case '+': r = a + b; break;
        case '-': r = a - b; break;
        case '*': r = a * b; break;
        case '/': r = a / b; break;
        case '^': for (long long j = 0; j < b; j++) r *= a; break;
        }
        st[top++] = r;
    }
    return st[0];
}

int main(void)
{
    char line[MAXLEN], post[MAXLEN];
    while (fgets(line, sizeof line, stdin) != NULL) {
        line[strcspn(line, "\n")] = '\0';
        to_postfix(line, post);
        for (int i = 0; post[i] != '\0'; i++) {
            printf("%c ", post[i]);
        }
        printf("= %lld\n", eval_postfix(post));
    }
    return 0;
}
\end{example}
\iosample{7 - 4 * (1 + 2 ^ 3 ^ 2) / 9 + 5 \\ (8 + 2 * 5) / (1 + 3 * 2 - 4) \\ 9 - 5 - 3 \\ 2 ^ 3 ^ 2 \\ (2 ^ 3) ^ 2 \\ 9 / 4 * 3}{7 4 1 2 3 2 ^ ^ + * 9 / - 5 + = -216 \\ 8 2 5 * + 1 3 2 * + 4 - / = 6 \\ 9 5 - 3 - = 1 \\ 2 3 2 ^ ^ = 512 \\ 2 3 ^ 2 ^ = 64 \\ 9 4 / 3 * = 6}
\endgroup

第一行就是追蹤表的運算式，把 \texttt{A} 到 \texttt{H} 換成 $7, 4, 1, 2, 3, 2, 9, 5$，
後序的形狀和手算的一樣。 $2 \sqpow 3 \sqpow 2 = 2^9 = 512$，加了括號的
$(2 \sqpow 3) \sqpow 2 = 8^2 = 64$，結合性不同，答案差很多。 \texttt{9 - 5 - 3}
由左往右是 $1$，錯當成由右往左會得到 $9 - 2 = 7$；\texttt{9 / 4 * 3} 由左往右是
$2 * 3 = 6$，由右往左會得到 $9 / 12 = 0$。

要支援多位數，就先把一行字切成一個一個記號（數字、運算子、括號），堆疊裡放記號
而不是字元，演算法不變。 程式也沒有處理一元負號（例如 \texttt{-3 + 5} 開頭的
\texttt{-}）：它和減號是同一個字元，要看前一個記號是什麼來分辨，通常在切記號時改
成另一個符號再處理。

\subsection{前序}
\label{sec:sq-prefix}

前序求值有兩種做法。 第一種是\textbf{由右往左}讀，規則和後序一樣，只是 pop 的順序
反過來：\textbf{先 pop 出來的是左邊的運算元}。 例如 $-\;8\;3$ 由右往左讀，先 push
$3$、再 push $8$，讀到 $-$ 時先 pop 出 $8$（左）、再 pop 出 $3$（右），算 $8 - 3$。

第二種是遞迴：前序運算式不是一個數，就是「一個運算子，後面接兩個前序運算式」。
寫一個函式，讀到數字就回傳它；讀到運算子，就遞迴讀左邊的運算式、再遞迴讀右邊的，
算完回傳。 這個遞迴用的就是系統的呼叫堆疊，下一節說明它和第一種做法的關係。

中序轉前序，手算用加滿括號最不容易錯。 用堆疊的做法要由右往左讀，並且注意結合性
跟著反過來，容易出錯，這裡不展開；第~7~章用運算式樹可以一次得到三種寫法。

\section{遞迴與堆疊：把遞迴改成迴圈}
\label{sec:sq-recursion}

\subsection{呼叫堆疊}

\begin{cbookrecap}[第~15~章 15.4 節、15.13 節]
每呼叫一次函式，系統就在\textbf{呼叫堆疊}（call stack）上放一個 frame，裡面有這次
呼叫的參數、區域變數，以及返回之後要接著執行的位置。 函式返回時，它的 frame 被
pop 掉，回到呼叫它的那一層。 遞迴每深一層就多一個 frame，所以遞迴深度太大會用完
呼叫堆疊（stack overflow）。 第~15~章 15.13 節用一個指標陣列當堆疊，把「反向印出
串列」改成了迴圈。
\end{cbookrecap}

呼叫堆疊就是一個堆疊：最後呼叫的函式最先返回。 遞迴之所以能「暫停目前的工作，先
去做子問題，做完再回來」，就是因為暫停的工作都存在呼叫堆疊裡。 所以任何遞迴都能
改寫成迴圈加上一個自己管理的堆疊。 自己管理的好處是：深度只受自己配置的記憶體限
制，而且堆疊裡只放真正需要的東西。

這也說明了遞迴的\textbf{空間}要怎麼算：呼叫堆疊最深時有幾層，每層多大。 遞迴深度
$d$、每層 $\Theta(1)$ 的遞迴，額外空間是 $\Theta(d)$，不是 $\Theta(1)$。 例如遞迴
版的二分搜尋，深度 $\Theta(\lg n)$，用了 $\Theta(\lg n)$ 的空間；迴圈版只要
$\Theta(1)$。

\subsection{尾遞迴：不需要堆疊}

如果遞迴呼叫是函式的最後一個動作，回來之後什麼都不用做，叫\textbf{尾遞迴}（tail
recursion）。 這時目前這一層的資料之後用不到，不用存，直接把參數換成新的值、跳回
開頭就好。 例如輾轉相除法：

\begin{pseudo}[label=alg:sq-gcd]{最大公因數：遞迴和迴圈}
\Function{Gcd}{$a, b$}
    \If{$b = 0$}
        \State \Return $a$
    \EndIf
    \State \Return \Call{Gcd}{$b, a \bmod b$}
\EndFunction
\Statex
\Function{Gcd-Iter}{$a, b$}
    \While{$b \ne 0$}
        \State $r \gets a \bmod b$
        \State $a \gets b$
        \State $b \gets r$
    \EndWhile
    \State \Return $a$
\EndFunction
\end{pseudo}

\proc{Gcd}$(a, b)$ 的最後一步是回傳 \proc{Gcd}$(b, a \bmod b)$ 的結果，所以迴圈
版只要把 $(a, b)$ 換成 $(b, a \bmod b)$。 注意要用暫存的 $r$：直接寫
$a \gets b$、$b \gets a \bmod b$，第二行用到的已經是新的 $a$。

\subsection{一般的遞迴：堆疊裡放待辦的工作}

遞迴呼叫之後還有事要做時，就需要堆疊了。 以河內塔為例：把 $k$ 個盤子從柱子
$a$ 搬到柱子 $c$、以 $b$ 暫放，遞迴的做法分三步：先把上面 $k - 1$ 個從 $a$ 搬到
$b$，再把最大的從 $a$ 搬到 $c$，最後把那 $k - 1$ 個從 $b$ 搬到 $c$。

\begin{pseudo}[label=alg:sq-hanoi]{河內塔：遞迴版}
\Procedure{Hanoi}{$k, a, c, b$} \Comment{$k$ 個盤子從 $a$ 搬到 $c$，$b$ 暫放}
    \If{$k = 0$}
        \State \Return
    \EndIf
    \State \Call{Hanoi}{$k - 1, a, b, c$}
    \State 印出「$a \to c$」
    \State \Call{Hanoi}{$k - 1, b, c, a$}
\EndProcedure
\end{pseudo}

改成迴圈的想法：堆疊裡放「還沒做的工作」，每個工作是一組 $(k, a, c, b)$，意思是
「把 $k$ 個盤子從 $a$ 搬到 $c$，$b$ 暫放」。 搬一個盤子也是一個工作，$k = 1$。
每次 pop 出一個工作：$k = 1$ 就直接搬；否則把它拆成上面三步，放回堆疊。 因為堆疊
後進先出，三步要\textbf{倒過來} push，最先要做的最後放，下一輪才會先拿到它。

\begin{pseudo}[label=alg:sq-hanoi-iter]{河內塔：用堆疊的迴圈版}
\Procedure{Hanoi-Iter}{$n, a, c, b$}
    \State 建立空堆疊 $S$
    \State \Call{Push}{$S, (n, a, c, b)$}
    \While{$S$ 不是空的}
        \State $(k, x, z, y) \gets$ \Call{Pop}{$S$}
        \If{$k = 1$}
            \State 印出「$x \to z$」
        \Else
            \State \Call{Push}{$S, (k - 1, y, z, x)$} \Comment{第三步，最後做}
            \State \Call{Push}{$S, (1, x, z, y)$} \Comment{第二步：搬最大的}
            \State \Call{Push}{$S, (k - 1, x, y, z)$} \Comment{第一步，最先做}
        \EndIf
    \EndWhile
\EndProcedure
\end{pseudo}

\textbf{手算追蹤}：$n = 3$，從 A 搬到 C，B 暫放。 堆疊由底到頂寫，
$(k, x, z, y)$ 簡寫成 $k{:}x{\to}z$（暫放的柱子可以從另外兩根推出來）。

\begin{center}
\small
\begin{tabular}{clll}
\toprule
輪 & pop 出的工作 & 做的事 & 之後的堆疊（底 $\to$ 頂）\\
\midrule
1 & $3{:}A{\to}C$ & 拆成三個 & $2{:}B{\to}C$\quad $1{:}A{\to}C$\quad $2{:}A{\to}B$ \\
2 & $2{:}A{\to}B$ & 拆成三個 & $2{:}B{\to}C$\quad $1{:}A{\to}C$\quad $1{:}C{\to}B$\quad $1{:}A{\to}B$\quad $1{:}A{\to}C$ \\
3 & $1{:}A{\to}C$ & 搬 A$\to$C & $2{:}B{\to}C$\quad $1{:}A{\to}C$\quad $1{:}C{\to}B$\quad $1{:}A{\to}B$ \\
4 & $1{:}A{\to}B$ & 搬 A$\to$B & $2{:}B{\to}C$\quad $1{:}A{\to}C$\quad $1{:}C{\to}B$ \\
5 & $1{:}C{\to}B$ & 搬 C$\to$B & $2{:}B{\to}C$\quad $1{:}A{\to}C$ \\
6 & $1{:}A{\to}C$ & 搬 A$\to$C & $2{:}B{\to}C$ \\
7 & $2{:}B{\to}C$ & 拆成三個 & $1{:}A{\to}C$\quad $1{:}B{\to}C$\quad $1{:}B{\to}A$ \\
8 & $1{:}B{\to}A$ & 搬 B$\to$A & $1{:}A{\to}C$\quad $1{:}B{\to}C$ \\
9 & $1{:}B{\to}C$ & 搬 B$\to$C & $1{:}A{\to}C$ \\
10 & $1{:}A{\to}C$ & 搬 A$\to$C & （空）\\
\bottomrule
\end{tabular}
\end{center}

七次搬動依序是 A$\to$C、A$\to$B、C$\to$B、A$\to$C、B$\to$A、B$\to$C、A$\to$C，和
遞迴版一樣。 第 2 輪之後堆疊最高，有 $5$ 個工作。

下面的程式把兩個版本都跑一次，比對 $n = 1$ 到 $10$ 的搬動順序是否完全相同，並記錄
迴圈版堆疊最高時有幾個工作。 $n = 3$ 時把搬動印出來。

\begin{example}[label=ex:sq-hanoi]{c}{(hanoi\_stack.c) 河內塔的遞迴版和堆疊版}
#include <stdio.h>
#include <stdlib.h>
#include <string.h>

#define MAXN 10

char rec[1 << MAXN][3], itr[1 << MAXN][3];  // 每一步存成 "AC" 這樣的字串
int nrec, nitr;

void hanoi_rec(int n, char a, char c, char b)   // n 個盤子從 a 搬到 c，b 暫放
{
    if (n == 0) {
        return;
    }
    hanoi_rec(n - 1, a, b, c);
    rec[nrec][0] = a;
    rec[nrec++][1] = c;
    hanoi_rec(n - 1, b, c, a);
}

typedef struct {
    int n;
    char a, c, b;
} Task;                     // 待辦的工作：n 個盤子從 a 搬到 c，b 暫放

int hanoi_iter(int n, char a, char c, char b)   // 回傳堆疊最高時有幾個工作
{
    Task *st = malloc(2 * n * sizeof *st);
    if (st == NULL) {
        return -1;
    }
    int top = 0, deepest = 0;
    st[top++] = (Task){n, a, c, b};
    while (top > 0) {
        if (top > deepest) {
            deepest = top;
        }
        Task t = st[--top];
        if (t.n == 1) {
            itr[nitr][0] = t.a;
            itr[nitr++][1] = t.c;
        } else {            // 三步倒過來 push，最先要做的最後放
            st[top++] = (Task){t.n - 1, t.b, t.c, t.a};
            st[top++] = (Task){1, t.a, t.c, t.b};
            st[top++] = (Task){t.n - 1, t.a, t.b, t.c};
        }
    }
    free(st);
    return deepest;
}

int main(void)
{
    printf("n moves same deepest\n");
    for (int n = 1; n <= MAXN; n++) {
        nrec = nitr = 0;
        hanoi_rec(n, 'A', 'C', 'B');
        int deepest = hanoi_iter(n, 'A', 'C', 'B');
        int same = (nrec == nitr);
        for (int i = 0; same && i < nrec; i++) {
            same = (strcmp(rec[i], itr[i]) == 0);
        }
        printf("%d %d %s %d\n", n, nitr, same ? "yes" : "no", deepest);
        if (n == 3) {
            for (int i = 0; i < nitr; i++) {
                printf("%s%c", itr[i], i + 1 < nitr ? ' ' : '\n');
            }
        }
    }
    return 0;
}
\end{example}
\iosample{}{n moves same deepest \\ 1 1 yes 1 \\ 2 3 yes 3 \\ 3 7 yes 5 \\ AC AB CB AC BA BC AC \\ 4 15 yes 7 \\ 5 31 yes 9 \\ 6 63 yes 11 \\ 7 127 yes 13 \\ 8 255 yes 15 \\ 9 511 yes 17 \\ 10 1023 yes 19}

搬動次數是 $2^n - 1$，每個 $n$ 的順序都和遞迴版相同。 堆疊最高是 $2n - 1$ 個工作
（觀念題第~\ref{q:sq-hanoi-depth}~題證明），所以程式配置 $2n$ 格就夠。 遞迴版的呼叫堆疊最深是 $n + 1$
層（$n, n-1, \ldots, 1, 0$），兩者都是 $\Theta(n)$ 的空間。

改寫的一般步驟可以整理成三句：堆疊裡放「之後還要做的事」，需要的參數都放進去；
一個工作拆成幾個子工作時，倒過來 push；遞迴的終止條件變成「pop 出來的工作可以直
接做完」。 第~7~章二元樹走訪的迴圈版、第~14~章的快速排序、第~21~章的深度優先
搜尋，都可以用這個方法把遞迴改成迴圈。

\section{多個堆疊共用一個陣列}
\label{sec:sq-multi}

一個程式需要好幾個堆疊，而且只能用一塊固定大小的記憶體時，要怎麼分？ 平均切成幾
段最簡單，但可能其中一個堆疊滿了，別的堆疊還空著很多。

\subsection{兩個堆疊}

兩個堆疊可以做到「只要總數不超過陣列大小就不會溢位」：一個從陣列的左端往右長，另
一個從右端往左長，兩個頂端在中間相遇時才是真的滿了。
用 $A[1..n]$ 的話，一號堆疊佔 $A[1..t_1]$，二號堆疊佔 $A[t_2..n]$，一開始
$t_1 = 0$、$t_2 = n + 1$。 一號 push 是 $t_1 \gets t_1 + 1$ 再放進 $A[t_1]$，二號
push 是 $t_2 \gets t_2 - 1$ 再放進 $A[t_2]$；兩者共用同一個滿的條件
$t_1 + 1 = t_2$，也就是中間一格都不剩。 兩個堆疊的大小此消彼長時（例如一個只在
另一個縮小時才長大），這樣分最省。

\subsection{$m$ 個堆疊}

三個以上就沒辦法都從端點長了。 常見的做法是先把 $A[1..n]$ 平均分成 $m$ 段，用
$b[i]$ 記第 $i$ 段的邊界、$t[i]$ 記第 $i$ 個堆疊的頂端：第 $i$ 個堆疊佔
$A[b[i]+1..t[i]]$，空的時候 $t[i] = b[i]$。 一開始
\[
b[i] = t[i] = \left\lfloor \frac{(i - 1)\, n}{m} \right\rfloor, \quad
i = 1, \ldots, m, \qquad b[m+1] = n.
\]
第 $i$ 個堆疊滿了，是指 $t[i] = b[i + 1]$，頂端碰到下一段的邊界。

\begin{center}
\begin{tikzpicture}
% 兩個堆疊
\node[font=\small, anchor=east] at (-0.6, 0) {兩個};
\foreach \k in {1,...,10} {
    \ifnum\k<4
        \node[sqfill] at (\k*0.9-0.9, 0) {};
    \else\ifnum\k>6
        \node[hlcell] at (\k*0.9-0.9, 0) {};
    \else
        \node[cell] at (\k*0.9-0.9, 0) {};
    \fi\fi
}
\draw[diagarrow, draw=coreStroke] (2.3, 0.65) -- (3.3, 0.65);
\draw[diagarrow, draw=memStroke] (5.1, 0.65) -- (4.1, 0.65);
\node[font=\scriptsize] at (1.0, 0.65) {一號往右長};
\node[font=\scriptsize] at (6.6, 0.65) {二號往左長};
% m 個堆疊
\node[font=\small, anchor=east] at (-0.6, -1.6) {$m = 3$};
\foreach \k in {1,...,12} {
    \pgfmathtruncatemacro{\seg}{(\k-1)/4}
    \pgfmathtruncatemacro{\pos}{mod(\k-1,4)}
    \ifnum\seg=0 \ifnum\pos<2 \node[sqfill] at (\k*0.9-0.9, -1.6) {}; \else \node[cell] at (\k*0.9-0.9, -1.6) {}; \fi \fi
    \ifnum\seg=1 \node[hlcell] at (\k*0.9-0.9, -1.6) {}; \fi
    \ifnum\seg=2 \ifnum\pos<1 \node[sqfill] at (\k*0.9-0.9, -1.6) {}; \else \node[cell] at (\k*0.9-0.9, -1.6) {}; \fi \fi
}
\foreach \x/\l in {-0.45/b[1], 3.15/b[2], 6.75/b[3], 10.35/b[4]} {
    \draw[thick] (\x, -2.15) -- (\x, -1.05);
    \node[idxlbl] at (\x, -2.4) {$\l$};
}
\node[idxlbl] at (0.9, -1.0) {$t[1]$};
\node[idxlbl] at (6.3, -1.0) {$t[2]$};
\node[idxlbl] at (7.2, -1.0) {$t[3]$};
\end{tikzpicture}
\diagcaption{上：兩個堆疊從兩端往中間長，中間的空格誰都能用。 下：三個堆疊各分到
$4$ 格，第 2 個（橘色）已經滿了，但第 1、3 個還有空格。}
\label{fig:sq-multi}
\end{center}

圖中第 2 個堆疊滿了，再 push 進第 2 個時，陣列其實還有 $5$ 格空著。 這時可以把別
的堆疊挪一挪：先往右找，找到第一個還沒滿的堆疊 $j > i$，把第 $i + 1$ 到第 $j$ 個
堆疊的內容整段往右移一格（它們的 $b$、$t$ 都加一），第 $i$ 個堆疊就多出一格；右
邊都滿了，就往左找，把第 $j$ 到第 $i$ 個堆疊往左移一格。 兩邊都找不到，才是真的
滿了。

這個做法保證只要陣列還有空位就能 push，但一次 push 最壞要搬 $\Theta(n)$ 個元素。
如果某個堆疊一直長，幾乎每次 push 都要搬。 所以它比較像是觀念：記憶體很緊、堆疊
個數固定時的折衷。 記憶體夠用時，每個堆疊各自用一個會自動加倍的陣列（第~3~章），
或乾脆用串列，每個操作都是 $\Theta(1)$（平攤或最壞），也不用管邊界。

\section{本章常見錯誤}
\label{sec:sq-pitfalls}

\begin{center}
\begin{tabular}{p{0.24\linewidth}p{0.50\linewidth}p{0.16\linewidth}}
\toprule
錯誤 & 例子 & 見 \\
\midrule
沒檢查空和滿 & 對空堆疊 pop，讀到 \texttt{data[-1]}；溢位只印警告就繼續，後面全錯 &
第~\ref{sec:sq-stack}~節 \\
\texttt{top} 差一 & C 裡 \texttt{top} 是個數，頂端在 \texttt{data[top - 1]}，卻讀
\texttt{data[top]} & 第~\ref{sec:sq-stack}~節 \\
括號配對漏了最後一步 & 讀完沒檢查堆疊是否為空，\texttt{((} 被當成正確 &
第~\ref{sec:sq-brackets}~節 \\
環狀佇列的空滿混淆 & 只用 $\mathit{front} = \mathit{rear}$ 同時當空和滿；忘了題目
的 $\mathit{front}$ 定義和自己習慣的不同 & 第~\ref{sec:sq-circular}~節 \\
負數取餘數 & C 寫 \texttt{(rear - front) \% n}，$\mathit{rear} < \mathit{front}$ 時
得到負數 & 第~\ref{sec:sq-circular}~節 \\
串列佇列沒顧邊界 & 取走最後一個節點時沒把 $\mathit{rear}$ 清成 \NIL &
第~\ref{sec:sq-queue}~節 \\
運算元順序弄反 & 後序 $8\;3\;-$ 算成 $3 - 8$；前序由右往左時，把先 pop 的當成右邊 &
第~\ref{sec:sq-postfix-eval}、\ref{sec:sq-prefix}~節 \\
結合性弄錯 & $a \sqpow b \sqpow c$ 轉成 $a\;b\;\sqpow\;c\;\sqpow$；$a - b - c$ 轉成
$a\;b\;c\;-\;-$ & 第~\ref{sec:sq-infix-postfix}~節 \\
改寫遞迴時順序弄反 & 子工作照原來的順序 push，結果最後的那一步最先做 &
第~\ref{sec:sq-recursion}~節 \\
\bottomrule
\end{tabular}
\end{center}

\section{觀念題}
\label{sec:sq-exercises}

題型和研究所考試、期中期末考的筆試相近，有選擇、計算、追蹤和證明。 解答在本章最後。

\begin{enumerate}
\item （選擇，可能不只一個答案）$1, 2, 3, 4, 5, 6$ 依序進入一個堆疊，中間可以隨時
pop。 下列哪些輸出順序\textbf{做不到}？
\begin{enumerate}
\item $3, 2, 1, 6, 5, 4$
\item $4, 5, 3, 6, 2, 1$
\item $2, 5, 3, 4, 1, 6$
\item $6, 1, 2, 3, 4, 5$
\item $1, 4, 3, 6, 5, 2$
\end{enumerate}

\item （計算）$1, 2, 3, 4, 5$ 依序進入一個堆疊。
\begin{enumerate}
\item 做得到的輸出順序中，最後一個是 $1$ 的有幾種？
\item 做得到的輸出順序中，第一個是 $3$ 的有幾種？ 全部列出來。
\end{enumerate}

\item （追蹤）用虛擬碼~\ref{alg:sq-circular}（留一格不用），$n = 6$，一開始
$\mathit{front} = \mathit{rear} = 0$。 依序做：E 10、E 20、E 30、E 40、D、D、E 50、
E 60、E 70、E 80、D、E 90。 寫出每個操作之後的 $\mathit{front}$、
$\mathit{rear}$，哪個操作失敗？ 最後陣列 $Q[0..5]$ 實際存的值是什麼？

\item （計算）某課本的環狀佇列用 $Q[0..n-1]$，規定 $\mathit{front}$ 指向前端元素的
\textbf{前一格}，$\mathit{rear}$ 指向尾端元素本身，一開始
$\mathit{front} = \mathit{rear} = 0$，一樣留一格不用。
\begin{enumerate}
\item 寫出這個定義下空、滿的條件和元素個數的公式。
\item $n = 10$、$\mathit{front} = 7$、$\mathit{rear} = 3$ 時，佇列有幾個元素？ 前
端元素在哪一格？
\end{enumerate}

\item \label{q:sq-pigeon}（證明）一個用 $Q[0..n-1]$ 的環狀佇列，只存兩個 $0$ 到 $n - 1$ 之間的整數
$\mathit{front}$、$\mathit{rear}$，兩者的意義可以自己定，但要滿足：(i) 佇列有幾個
元素，完全由 $(\mathit{front}, \mathit{rear})$ 決定（陣列裡殘留的舊值和真的元素分
不出來，不能拿來判斷）；(ii) \proc{Enqueue} 不改 $\mathit{front}$。 證明這個佇列最
多只能放 $n - 1$ 個元素。

\item （追蹤）用第~\ref{sec:sq-infix-postfix}~節的規則，把
\texttt{A * (B - C) / D \sqpow{} E \sqpow{} F + G * H} 轉成後序，寫出逐步的堆疊表。
再用加滿括號的方法寫出它的前序。

\item （計算）求值（\sqpow{} 是次方）：
\begin{enumerate}
\item 後序 $6\;3\;-\;4\;2\;\sqpow\;*\;8\;2\;/\;-$
\item 前序 $-\;*\;+\;4\;2\;3\;/\;8\;\sqpow\;2\;2$
\end{enumerate}

\item （證明）一串記號只含運算元和二元運算子。 令 $d_k$ 是前 $k$ 個記號裡「運算元
個數減去運算子個數」。 證明：這串記號能被虛擬碼~\ref{alg:sq-postfix-eval} 算完（過
程中從不對空堆疊 pop，結束時堆疊恰好剩一個數）的充分必要條件是：對所有
$1 \le k \le m$ 都有 $d_k \ge 1$，而且 $d_m = 1$。 再用它判斷 \texttt{a b c - + d}
和 \texttt{a b + - c d} 是不是合法的後序。

\item \label{q:sq-deque}（選擇）$1, 2, 3, 4$ 依序進入一個\textbf{輸入受限}的雙端佇列（只能從尾端加入，
兩端都能取出，中間可以隨時取出）。 下列哪些輸出順序\textbf{做不到}？
(A) $4, 1, 3, 2$ \quad (B) $4, 2, 1, 3$ \quad (C) $1, 4, 2, 3$ \quad
(D) $4, 2, 3, 1$ \quad (E) $3, 1, 2, 4$

\item （設計）用一個\textbf{環狀單向串列}存佇列，而且只記一個指標
$\mathit{rear}$ 指向尾端節點（尾端節點的 $\mathit{next}$ 指回前端節點）。 寫出
\proc{Enqueue} 和 \proc{Dequeue} 的虛擬碼，兩者都要 $\Theta(1)$，並處理佇列為空和
只剩一個節點的情況。 為什麼只記 $\mathit{front}$ 不行？

\item \label{q:sq-hanoi-depth}（證明）證明虛擬碼~\ref{alg:sq-hanoi-iter} 對 $n \ge 1$ 執行時，堆疊裡最多同
時有 $2n - 1$ 個工作，而且這個數會達到。 （提示：對 $k$ 做歸納，證明一個更強的敘
述：堆疊下面已經有 $s$ 個工作、頂端是一個 $k$ 的工作時，從這時起到這個工作「做完」
（堆疊回到只剩下面那 $s$ 個）為止，每次 \textbf{while} 開始時堆疊最多有
$s + 2k - 1$ 個工作，而且下面那 $s$ 個沒有被動過。）

\item \label{q:sq-fib}（計算）考慮下面的函式。
\begin{code}{c}{觀念題的函式 f}
int f(int n)
{
    if (n <= 1)
        return 1;
    return f(n - 1) + f(n - 2);
}
\end{code}
\begin{enumerate}
\item 呼叫 \texttt{f(6)}，總共呼叫 \texttt{f} 幾次（含 \texttt{f(6)} 本身）？
\item 呼叫堆疊最深時有幾層 \texttt{f} 的 frame？ 額外空間是 $n$ 的什麼函數？
\item 這個函式是不是尾遞迴？ 說明理由。
\end{enumerate}

\item （計算）$m = 3$ 個堆疊共用 $A[1..12]$，依第~\ref{sec:sq-multi}~節的方式平均
分段。
\begin{enumerate}
\item 寫出一開始的 $b[1..4]$ 和 $t[1..3]$。
\item 依序 push 進第 $2$ 個堆疊 $5$ 次。 前 $4$ 次之後 $t[2]$ 是多少？ 第 $5$ 次時
發生什麼事，用「先往右找」的規則挪動之後，$b$、$t$ 變成多少？
\end{enumerate}
\end{enumerate}

\section{實作練習}
\label{sec:sq-practice}

下面是 Aki-Judge 上和本章相關的題目。 「實作 ADT」類的題目只要交標頭檔和函式，不
寫 \texttt{main}；交之前先用本章的範例自己測一次，特別是空的、滿的、繞回陣列開頭
的情況。

\begin{center}
\begin{tabular}{llc}
\toprule
小節 & 題目 & 難度 \\
\midrule
堆疊的實作 & \akiref{growable-array-stack}{可自動擴充的陣列堆疊} & 基本 \\
 & \akiref{linked-node-stack}{鏈結節點堆疊} & 基本 \\
 & \akiref{container-yard-stacks}{貨櫃堆疊場} & 基本 \\
括號配對 & \akiref{bracket-matching-check}{括號配對檢查} & 基本 \\
佇列的實作 & \akiref{ring-buffer-queue}{環狀緩衝區佇列} & 基本 \\
 & \akiref{sushi-conveyor-belt}{迴轉壽司輸送帶} & 基本 \\
雙端佇列 & \akiref{double-ended-queue}{雙端佇列} & 基本 \\
 & \akiref{queue-line-manager}{排隊名單調度} & 進階 \\
運算式 & \akiref{evaluate-postfix-expression}{後序運算式求值} & 暖身 \\
 & \akiref{convert-infix-to-postfix}{中序運算式轉後序} & 基本 \\
 & \akiref{postfix-string-calculator}{後序字串計算機} & 基本 \\
 & \akiref{prefix-calculator-lazy-conditional}{附條件運算的前序算式} & 進階 \\
單調堆疊、單調佇列 & \akiref{next-larger-value-to-the-right}{右側第一個較大的數} & 進階 \\
 & \akiref{maximum-of-every-window}{每個視窗的最大值} & 進階 \\
\bottomrule
\end{tabular}
\end{center}

幾題的提示：

\begin{itemize}
\item 「可自動擴充的陣列堆疊」要在陣列滿的時候把容量加倍，加倍為什麼平均是
$\Theta(1)$，見第~3~章。
\item 「迴轉壽司輸送帶」滿了還要放時，要先取出最舊的一盤，再放新的；查第 $i$ 盤
就是從 $\mathit{front}$ 往後數 $i - 1$ 格，記得取餘數。
\item 「排隊名單調度」要從隊伍中間拉出任何一個人，用雙向串列，另外開一個陣列從編
號直接找到節點（第~5~章）。
\item 「附條件運算的前序算式」照第~\ref{sec:sq-prefix}~節的遞迴寫法最自然；\texttt{?}
只計算被選中的那一邊，所以沒被選中的那一邊只要檢查語法、不能算值。
\item 最後兩題是本章正文沒有講的延伸。 「右側第一個較大的數」：由左往右掃，堆疊
裡放還沒找到答案的位置，它們的值由底到頂不增；新的數比頂端大，頂端的答案就是它，
pop 掉再比下一個。 「每個視窗的最大值」：雙端佇列裡放「還可能成為最大值」的位置，
值由前到後遞減，新的數從尾端進來前，先從尾端 pop 掉比它小的；前端的位置滑出視窗
就從前端 pop。 兩題每個位置都只進出一次，總時間 $\Theta(n)$，這也是第~3~章聚合法
的例子。
\end{itemize}

\section{讀完你應該會}
\label{sec:sq-checklist}

\begin{itemize}
\item 用陣列和串列各寫出堆疊、佇列的所有操作，處理空、滿、只剩一個元素的邊界，並
說出兩種做法在空間上的差別。
\item 給一個輸出順序，判斷 $1, \ldots, n$ 經過堆疊（或受限的雙端佇列）能不能得到，
寫出 push、pop 序列或指出卡住的地方；知道做得到的順序有 $C_n$ 種。
\item 環狀佇列依題目給的 $\mathit{front}$、$\mathit{rear}$ 定義，寫出空、滿的條件
和元素個數，逐步追蹤每個操作。
\item 把中序運算式轉成前序、後序（加滿括號或用堆疊），寫出逐步的堆疊表，並正確處
理括號和結合性；能用堆疊算出後序、前序的值。
\item 把一個遞迴函式改寫成迴圈加上自己的堆疊，並算出遞迴的深度和額外空間。
\end{itemize}

\textbf{延伸閱讀}（書末會整理完整的參考資料）

\begin{itemize}
\item Thomas H. Cormen 等，\textit{Introduction to Algorithms}，第 4 版，10.1--10.2
節：用陣列和串列實作堆疊、佇列的標準寫法，以及雙向串列。
\item Steven S. Skiena，\textit{The Algorithm Design Manual}，第 3 版，3.2 節：從
「什麼時候取出的順序重要」的角度談堆疊和佇列。
\item Antti Laaksonen，\textit{Competitive Programmer's Handbook}，8.2、8.3 節：用堆
疊找每個位置最近的較小元素、用雙端佇列求滑動視窗的最小值，就是上面最後兩題的做法，
並用平攤分析說明為什麼是線性時間。
\item 《C 程式設計基礎》第~20~章：把堆疊包成標頭檔加實作檔，換掉實作而不改用的人
的程式，是本章 ADT 觀念在 C 裡的寫法。
\end{itemize}

\addcontentsline{toc}{section}{觀念題解答}
\section*{觀念題解答}

\begin{enumerate}
\item (c)、(d) 做不到。 (c)：輸出 $5$ 時 $1$ 到 $5$ 都 push 過了，堆疊由底到頂是
$1, 3, 4$，頂端是 $4$，下一個要 $3$ 卻在 $4$ 下面；也可以直接看 $5, 3, 4$ 是
「大、小、中」。 (d)：輸出 $6$ 之後堆疊是 $1, 2, 3, 4, 5$，頂端是 $5$，拿不到
$1$。 其他三個都做得到，例如 (b)：push 1、2、3、4，pop 4，push 5，pop 5，pop 3，
push 6，pop 6，pop 2，pop 1。

\item
\begin{enumerate}
\item $14$ 種。 $1$ 最後輸出，表示它一進去就待在底部直到最後，$2, 3, 4, 5$ 在它
上面自己進出，是 $n = 4$ 的堆疊問題，共 $C_4 = 14$ 種。
\item $9$ 種。 先 push 1、2、3 再 pop 3，這時堆疊是 $1, 2$，還有 $4, 5$ 沒進來。 下
一個輸出是 $2$、$4$ 或 $5$：
\begin{itemize}
\item 下一個是 $2$：剩下堆疊 $1$，還有 $4, 5$，得到 32145、32154、32415、32451、
32541，共 $5$ 種。
\item 下一個是 $4$：push 4 再 pop，堆疊 $1, 2$，還有 $5$，得到 34215、34251、
34521，共 $3$ 種。
\item 下一個是 $5$：push 4、5 再 pop 5，之後只能依序 pop，得到 35421，$1$ 種。
\end{itemize}
\end{enumerate}

\item 每個操作之後的 $(\mathit{front}, \mathit{rear})$：E 10 $(0,1)$、E 20 $(0,2)$、
E 30 $(0,3)$、E 40 $(0,4)$、D 取出 10 $(1,4)$、D 取出 20 $(2,4)$、E 50 $(2,5)$、
E 60 $(2,0)$、E 70 $(2,1)$、E 80 失敗（$(1 + 1) \bmod 6 = 2 = \mathit{front}$，滿
了）、D 取出 30 $(3,1)$、E 90 $(3,2)$。 最後陣列 $Q[0..5] = 70, 90, 30, 40, 50, 60$，
其中 $Q[2] = 30$ 已經取出，佇列由前到後是 $40, 50, 60, 70, 90$，共
$(2 - 3 + 6) \bmod 6 = 5$ 個，剛好是滿的。

\item
\begin{enumerate}
\item 前端元素在 $(\mathit{front} + 1) \bmod n$。 加入是先
$\mathit{rear} \gets (\mathit{rear} + 1) \bmod n$ 再放；取出是先
$\mathit{front} \gets (\mathit{front} + 1) \bmod n$ 再拿。 空：
$\mathit{front} = \mathit{rear}$。 滿：$(\mathit{rear} + 1) \bmod n =
\mathit{front}$。 個數：$(\mathit{rear} - \mathit{front} + n) \bmod n$。 公式和本
書的定義長得一樣，因為兩個索引都往前挪了一格，差不變。
\item $(3 - 7 + 10) \bmod 10 = 6$ 個。 元素依序在 $Q[8]$、$Q[9]$、$Q[0]$、
$Q[1]$、$Q[2]$、$Q[3]$。
\end{enumerate}

\item 設這個佇列最多能放 $c$ 個元素。 從空佇列開始，連續做 $c$ 次
\proc{Enqueue}，中間經過 $c + 1$ 個狀態，元素個數依序是 $0, 1, \ldots, c$。 由
(ii)，這 $c + 1$ 個狀態的 $\mathit{front}$ 都一樣。 由 (i)，元素個數不同的兩個狀態，
$(\mathit{front}, \mathit{rear})$ 不能相同，所以這 $c + 1$ 個狀態的 $\mathit{rear}$
兩兩不同。 $\mathit{rear}$ 只有 $0, 1, \ldots, n - 1$ 共 $n$ 種值，所以
$c + 1 \le n$，$c \le n - 1$。

這個論證不管 $\mathit{front}$、$\mathit{rear}$ 怎麼定義都成立：只存兩個索引，
就一定要空下一格。 想用滿 $n$ 格，就要多存一點
資訊，例如第~\ref{sec:sq-circular}~節的個數或旗標，這時 (i) 不成立了。

\item 每讀一個記號的堆疊（底 $\to$ 頂）依序是：
A：空；$*$：$*$；(：$*\;($；B：同上；$-$：$*\;(\;-$；C：同上；
)：pop $-$，剩 $*$；$/$：$*$ 和 $/$ 一樣高、由左往右，pop $*$ 再 push，剩 $/$；
D：同上；\sqpow：$/\;\sqpow$；E：同上；\sqpow：一樣高但由右往左，不 pop，
$/\;\sqpow\;\sqpow$；F：同上；$+$：比三者都低，pop \sqpow{} \sqpow{} $/$，剩 $+$；
G：同上；$*$：比 $+$ 高，$+\;*$；H：同上；結束：pop $*$ $+$。 輸出依序累積成
\texttt{A B C - * D E F \sqpow{} \sqpow{} / G H * +}。 作答時照
第~\ref{sec:sq-infix-postfix}~節的四欄表寫，一個記號一列。 前序：加滿括號是
$(((A * (B - C)) / (D \sqpow (E \sqpow F))) + (G * H))$，把運算子移到左括號，得到
\texttt{+ / * A - B C \sqpow{} D \sqpow{} E F * G H}。 兩者的運算元順序都是 $A$ 到
$H$。


\item
\begin{enumerate}
\item 堆疊依序：$6$；$6, 3$；$-$ 得 $3$；$3, 4, 2$；\sqpow{} 得 $3, 16$；$*$ 得
$48$；$48, 8, 2$；$/$ 得 $48, 4$；$-$ 得 $44$。 答案 $44$，即
$(6 - 3) \cdot 4^2 - 8 / 2$。
\item 由右往左讀，先 pop 的是左邊：push 2、2；\sqpow{} pop 出 $2$（左）、$2$（右）
得 $4$；push 8；$/$ 得 $8 / 4 = 2$；push 3、2、4；$+$ 得 $4 + 2 = 6$，堆疊由底到頂
是 $2, 3, 6$；$*$ 得 $6 * 3 = 18$；$-$ 得 $18 - 2 = 16$。 答案 $16$，即
$(4 + 2) \cdot 3 - 8 / 2^2$。
\end{enumerate}

\item 先用歸納法證明：只要前 $k$ 個記號都處理成功，處理完之後堆疊裡有 $d_k$ 個數。
$k = 0$ 時堆疊是空的，$d_0 = 0$。 處理第 $k$ 個記號：運算元 push 一個，$d$ 加一；
運算子 pop 兩個、push 一個，堆疊少一個，$d$ 減一。 兩邊一致。

充分：若所有 $d_k \ge 1$。 第 $k$ 個記號是運算子時，$d_k = d_{k-1} - 1 \ge 1$，所
以處理前堆疊有 $d_{k-1} \ge 2$ 個數，pop 兩次不會下溢。 結束時堆疊有 $d_m = 1$ 個
數。

必要：若第一個 $d_k < 1$ 出現在某個 $k$。 $d_0 = 0$，$k = 1$ 時第一個記號是運算子，
對空堆疊 pop；$k > 1$ 時 $d_{k-1} \ge 1$、$d_k = d_{k-1} - 1 < 1$，所以
$d_{k-1} = 1$，處理第 $k$ 個運算子時堆疊只有一個數，第二次 pop 下溢。 若所有
$d_k \ge 1$ 但 $d_m \ne 1$，結束時堆疊不是恰好一個數。

判斷：\texttt{a b c - + d} 的 $d$ 依序是 $1, 2, 3, 2, 1, 2$，結尾 $d_m = 2 \ne 1$，
不合法（最後多一個運算元）。 \texttt{a b + - c d} 的 $d$ 依序是 $1, 2, 1, 0, \ldots$，
第 4 個記號 $d = 0$，不合法（$-$ 只有一個運算元可用）。

\item (B)、(D) 做不到。 要先輸出 $4$，得先讓 $1, 2, 3, 4$ 都進來，雙端佇列由前到後
是 $1, 2, 3, 4$。 取出 $4$（尾端）之後剩 $1, 2, 3$，只能取兩端的 $1$ 或 $3$，所以
第二個是 $2$ 的 (B)、(D) 都不行。 (A)：取 4（尾）、1（前）、3（尾）、2。 (C)：
$1$ 進來馬上取出，之後 $2, 3, 4$ 都進來，取 4（尾）、2（前）、3。 (E)：$1, 2, 3$
進來，取 3（尾）、1（前）、2，再讓 $4$ 進來取出。 (C) 和 (E) 用普通堆疊做不到
（$1, 4, 2, 3$ 和 $3, 1, 2$ 都含「大、小、中」），這也說明雙端佇列比堆疊有彈性。

\item 前端節點是 $\mathit{rear}.\mathit{next}$。
\begin{pseudo}{環狀串列的佇列}
\Procedure{Enqueue-Ring}{$Q, x$}
    \State 配置新節點 $p$，$p.\mathit{key} \gets x$
    \If{$Q.\mathit{rear} = \NIL$}
        \State $p.\mathit{next} \gets p$ \Comment{只有一個節點，自己接自己}
    \Else
        \State $p.\mathit{next} \gets Q.\mathit{rear}.\mathit{next}$ \Comment{新節點接到前端}
        \State $Q.\mathit{rear}.\mathit{next} \gets p$
    \EndIf
    \State $Q.\mathit{rear} \gets p$
\EndProcedure
\Statex
\Function{Dequeue-Ring}{$Q$}
    \If{$Q.\mathit{rear} = \NIL$}
        \State \Error{underflow}
    \EndIf
    \State $p \gets Q.\mathit{rear}.\mathit{next}$ \Comment{前端節點}
    \If{$p = Q.\mathit{rear}$}
        \State $Q.\mathit{rear} \gets \NIL$ \Comment{取走的是唯一的節點}
    \Else
        \State $Q.\mathit{rear}.\mathit{next} \gets p.\mathit{next}$
    \EndIf
    \State $x \gets p.\mathit{key}$，釋放 $p$
    \State \Return $x$
\EndFunction
\end{pseudo}
只記 $\mathit{front}$ 的話，加入時要找到尾端節點，得從前端繞一圈，是
$\Theta(n)$。 記 $\mathit{rear}$ 時，前端只差一步，兩端都是 $\Theta(1)$。

\item 用歸納法證明提示裡的敘述 $P(k)$：堆疊下面有 $s$ 個工作、頂端是一個 $k$ 的工作
時，到這個工作做完為止，每次 \textbf{while} 開始時堆疊最多有 $s + 2k - 1$ 個工作，
下面那 $s$ 個沒有被動過，而且最大值 $s + 2k - 1$ 會達到。

$k = 1$：只有一輪，開始時有 $s + 1 = s + 2 \cdot 1 - 1$ 個，pop 出來直接搬，沒有動到
下面。

$k \ge 2$：第一輪開始時有 $s + 1$ 個。 pop 出這個工作、push 三個之後，由底到頂是
「下面 $s$ 個、$(k - 1)$、$(1)$、$(k - 1)$」，共 $s + 3$ 個。 接著依序做三個工作，
每個都套用歸納假設：
\begin{itemize}
\item 頂端的 $(k - 1)$，下面有 $s + 2$ 個：最多 $(s + 2) + 2(k - 1) - 1 = s + 2k - 1$
個，而且會達到；做完時剩 $s + 2$ 個，下面沒被動過。
\item 接著的 $(1)$，下面有 $s + 1$ 個：最多 $s + 2$ 個。
\item 最後的 $(k - 1)$，下面有 $s$ 個：最多 $s + 2k - 3$ 個。
\end{itemize}
第一輪的 $s + 1$ 也不超過 $s + 2k - 1$。 所以整個過程最多 $s + 2k - 1$ 個，而且第一
段會達到；三段都不動自己下面的工作，所以原本那 $s$ 個沒被動過。

取 $s = 0$、$k = n$，就得到最多 $2n - 1$ 個，而且會達到。 達到的時刻是一路拆最上面
的工作 $(n), (n - 1), \ldots, (2)$，每拆一次多兩個，從 $1$ 個變成
$1 + 2(n - 1) = 2n - 1$ 個，和範例~\ref{ex:sq-hanoi} 印出的 \texttt{deepest} 一致。
同一個歸納也說明了迴圈版的搬動順序和遞迴版相同：兩者都是先做完第一個 $(k-1)$，
再搬最大的，再做第二個 $(k-1)$。

\item
\begin{enumerate}
\item 令 $c(n)$ 是呼叫 \texttt{f(n)} 的總次數，$c(0) = c(1) = 1$，
$c(n) = 1 + c(n - 1) + c(n - 2)$。 依序 $c(2) = 3$、$c(3) = 5$、$c(4) = 9$、
$c(5) = 15$、$c(6) = 25$。 共 $25$ 次。
\item 最深的一條是 \texttt{f(6)}$\to$\texttt{f(5)}$\to \cdots \to$\texttt{f(1)}，
共 $6$ 層。 一般是 $n$ 層（$n \ge 1$），額外空間 $\Theta(n)$。 雖然總呼叫次數是指
數級，同一時間在呼叫堆疊上的只有一條路徑。
\item 不是。 遞迴呼叫 \texttt{f(n - 1)} 回來之後，還要再呼叫 \texttt{f(n - 2)}、做
加法，最後一個動作是加法，不是遞迴呼叫。
\end{enumerate}

\item
\begin{enumerate}
\item $b[i] = \lfloor (i-1) \cdot 12 / 3 \rfloor$：$b[1] = 0$、$b[2] = 4$、
$b[3] = 8$、$b[4] = 12$；$t[1] = 0$、$t[2] = 4$、$t[3] = 8$。
\item 前 $4$ 次放在 $A[5..8]$，$t[2] = 8 = b[3]$，第 $2$ 個堆疊滿了。 第 $5$ 次時
往右找：第 $3$ 個堆疊 $t[3] = 8 < b[4] = 12$，還沒滿。 把第 $3$ 個堆疊（目前是空
的，沒有元素要搬）往右移一格：$b[3] = 9$、$t[3] = 9$。 第 $2$ 個堆疊多出一格，放進
$A[9]$，$t[2] = 9$。 最後 $b = (0, 4, 9, 12)$，$t = (0, 9, 9)$。
\end{enumerate}
\end{enumerate}

\chapter{樹與二元樹}
\label{ch:trees}

% 本章的圖共用的樣式與巨集（名稱都以 bt 開頭，避免和其他章衝突）
\tikzset{
    btn/.style={circle, draw=black!65, minimum size=6.5mm, inner sep=0pt, font=\small},
    bthl/.style={btn, fill=coreFill, draw=coreStroke},
    btop/.style={btn, fill=memFill, draw=memStroke},
    btthread/.style={-{Stealth[length=1.6mm]}, dashed, thick, draw=optStroke},
    bttree/.style={
        level distance=10mm,
        edge from parent/.style={draw=black!60, thick},
    },
    btmain/.style={bttree,
        level 1/.style={sibling distance=36mm},
        level 2/.style={sibling distance=18mm},
        level 3/.style={sibling distance=11mm},
        level 4/.style={sibling distance=9mm}},
}
% 本章大部分例子用的那棵樹（圖 7.1），節點名稱就是字母本身
\newcommand{\btmaintree}{%
\node[btn] (A) {A}
    child {node[btn] (B) {B}
        child {node[btn] (D) {D}}
        child {node[btn] (E) {E}
            child {node[btn] (G) {G}}
            child[missing]}}
    child {node[btn] (C) {C}
        child[missing]
        child {node[btn] (F) {F}
            child {node[btn] (H) {H}
                child[missing]
                child {node[btn] (J) {J}}}
            child {node[btn] (I) {I}}}};}

程式設計社的幹部名單要做成一張組織圖：社長底下有兩位副社長，每位副社長底下再各帶
幾個組長，組長底下還有組員。 這張圖不是一條直線，名單從社長開始往下分岔。 如果要
把它存進程式，之後還要「依職位由高到低列出所有人」或「從某個組長算起，他底下總共
有幾個人」，你會怎麼存？ 往下讀之前，先想一想。

前兩章的串列、堆疊和佇列都是\textbf{線性}的：每個元素最多只有一個「下一個」。 組
織圖、資料夾、比賽的對戰表、運算式的結構都是一層一層往下分岔的，這種結構叫
\textbf{樹}（tree）。 這一章先把名詞和幾個常考的性質整理清楚，再講二元樹的兩種存
法、四種走訪，以及考卷上最常見的幾種手算：由走訪序列把樹畫回來、在樹上加引線、數
有幾種不同形狀的二元樹、把運算式畫成樹。

這一章用到《C 程式設計基礎》第~14~章的結構指標、第~15~章的遞迴，以及第~19~章的
\texttt{malloc} 和二元樹入門（三種遍歷、由前序和中序重建）；用到時先在框裡複習。
本書第~6~章的堆疊、佇列和卡特蘭數，這一章會直接拿來用。

\section{名詞與基本性質}
\label{sec:bt-terms}

\subsection{樹的名詞}

一棵\textbf{樹}由有限個\textbf{節點}（node）組成。 不是空的樹有一個特別的節點叫
\textbf{根}（root），其餘的節點分成 $m \ge 0$ 組，每一組自己又是一棵樹，叫做根的
\textbf{子樹}（subtree）。 這個定義是遞迴的：樹是「一個根加上幾棵比較小的樹」。

根和每棵子樹的根之間連一條\textbf{邊}（edge）。 邊的上端叫\textbf{父節點}
（parent），下端叫\textbf{子節點}或\textbf{孩子}（child）；同一個父節點的孩子互為
\textbf{兄弟}（sibling）。 從根往下走到 $x$ 的路上經過的節點（含根，不含 $x$）是
$x$ 的\textbf{祖先}（ancestor），反過來 $x$ 是它們的\textbf{子孫}（descendant）。
除了根，每個節點恰好有一個父節點，所以 $n$ 個節點的樹恰好有 $n - 1$ 條邊。

一個節點的孩子個數叫它的\textbf{分支度}（degree）；整棵樹裡最大的分支度叫這棵樹的
分支度。 分支度為 $0$ 的節點是\textbf{葉節點}（leaf，也叫終端節點），其他的是
\textbf{內部節點}（internal node）。 好幾棵互不相交的樹放在一起叫\textbf{森林}
（forest）；把一棵樹的根拿掉，剩下的子樹就是一個森林。

孩子之間有先後順序（第一個孩子、第二個孩子……）的樹叫\textbf{有序樹}（ordered
tree）。 資料夾裡的項目、組織圖上由左到右的組別都有順序，本章的樹預設都是有序樹。

\subsection{深度、層與高度}

這三個詞各本書的定義不太一樣，本書固定用下面這一套：

\begin{itemize}
\item 節點 $x$ 的\textbf{深度}（depth）是從根走到 $x$ 經過的邊數，根的深度是 $0$。
\item 深度為 $d$ 的節點在第 $d + 1$ \textbf{層}（level），根在第 $1$ 層。
\item 樹的\textbf{高度}（height）是所有節點深度的最大值，也就是從根到最深的葉節點
經過的邊數。 只有一個節點的樹高度是 $0$，空樹的高度定為 $-1$。
\end{itemize}

所以高度為 $h$ 的樹共有 $h + 1$ 層。 有些書把「高度」定義成層數（只有根的樹高度是
$1$），有些書讓根在第 $0$ 層；同一個公式在不同定義下會差 $1$。

\begin{examnote}
題目問到高度、深度、層數，或「深度為 $k$ 的二元樹最多有幾個節點」這類公式時，先寫
一句「依題意，根在第 $1$ 層」或「這裡的高度是邊數」，再代公式。 題目沒有定義時，
寫出你用的定義，答案才不會因為差 $1$ 被當成錯的。 下面的「滿」「完整」也一樣，各
本書的用詞不同，先寫出定義。
\end{examnote}

\subsection{二元樹}

\begin{definition}[二元樹]
\label{def:bt-binary-tree}
一棵\textbf{二元樹}（binary tree）要嘛是空的，要嘛由一個根、一棵\textbf{左子樹}和
一棵\textbf{右子樹}組成，左右子樹各自也是二元樹（可以是空的）。
\end{definition}

二元樹和「分支度不超過 $2$ 的有序樹」不一樣：二元樹的孩子分左右，只有一個孩子時也
要說清楚是左孩子還是右孩子。 例如兩個節點的二元樹有兩種（根的左邊或右邊接一個節
點），兩個節點的有序樹只有一種。 二元樹也可以是空的，樹的定義通常不允許空樹。

圖~\ref{fig:bt-main} 是本章大部分例子用的二元樹，共 $10$ 個節點。 $A$ 是根；$B$、
$C$ 是 $A$ 的孩子；$C$ 只有右孩子 $F$；$D$、$G$、$J$、$I$ 是葉節點；$H$ 的深度是
$3$，在第 $4$ 層；這棵樹的高度是 $4$（$A \to C \to F \to H \to J$）。 $E$ 的祖先是
$B$、$A$，$F$ 的子孫是 $H$、$I$、$J$。

\begin{center}
\begin{tikzpicture}[btmain]
\btmaintree
\foreach \d/\y in {0/0, 1/-1, 2/-2, 3/-3, 4/-4} {
    \node[idxlbl, anchor=east] at (-4.6, \y) {深度 \d};
}
\end{tikzpicture}
\diagcaption{本章的例子樹 $T$。 分支度為 $2$ 的有 $A$、$B$、$F$，分支度為 $1$ 的有
$C$、$E$、$H$，其餘四個是葉節點。}
\label{fig:bt-main}
\end{center}

在正文裡畫樹很佔地方，所以本章也用一種文字寫法：$x(L, R)$ 表示根是 $x$、左子樹是
$L$、右子樹是 $R$，空的子樹寫 $-$，葉節點只寫名字。 圖~\ref{fig:bt-main} 的樹寫成
\[
A\bigl(B(D, E(G, -)),\; C(-, F(H(-, J), I))\bigr).
\]

\subsection{節點數與高度}

二元樹第 $1$ 層只有根，每往下一層，節點數最多變成兩倍，所以：

\begin{theorem}
\label{thm:bt-count}
\begin{enumerate}
\item 二元樹的第 $i$ 層最多有 $2^{i-1}$ 個節點（深度 $d$ 最多 $2^d$ 個）。
\item 高度為 $h$ 的二元樹最少有 $h + 1$ 個節點，最多有 $2^{h+1} - 1$ 個。
\item $n \ge 1$ 個節點的二元樹，高度 $h$ 滿足
$\lfloor \lg n \rfloor \le h \le n - 1$。
\end{enumerate}
\end{theorem}

\begin{proof}
(1) 對層數歸納：第 $1$ 層只有根，$2^0 = 1$。 第 $i$ 層的每個節點最多兩個孩子，所以
第 $i + 1$ 層最多 $2 \cdot 2^{i-1} = 2^i$ 個。

(2) 高度 $h$ 表示有一個深度 $h$ 的節點，它和它的 $h$ 個祖先就有 $h + 1$ 個節點。
最多的情形是每層都滿：$\sum_{d=0}^{h} 2^d = 2^{h+1} - 1$。

(3) 由 (2)，$n \le 2^{h+1} - 1 < 2^{h+1}$，取對數得 $\lg n < h + 1$。 $h$ 是整數，
所以 $h \ge \lfloor \lg n \rfloor$（若 $h < \lfloor \lg n \rfloor$，則
$h + 1 \le \lfloor \lg n \rfloor \le \lg n$，矛盾）。 上界由 (2) 的
$n \ge h + 1$ 得到。
\end{proof}

高度下界常寫成 $\lceil \lg(n + 1) \rceil - 1$，和 $\lfloor \lg n \rfloor$ 對所有
$n \ge 1$ 都相等。 第~9、10~章的搜尋樹每個操作都花 $\Theta(h)$，(3) 告訴我們 $h$ 最
好是 $\Theta(\lg n)$、最壞是 $\Theta(n)$，平衡樹要做的就是把 $h$ 壓在下界附近。

下一個性質是考卷上最常出現的一個等式。 令 $n_0$、$n_1$、$n_2$ 分別是分支度為
$0$、$1$、$2$ 的節點數。

\begin{theorem}
\label{thm:bt-n0n2}
不是空的二元樹都有 $n_0 = n_2 + 1$。
\end{theorem}

\begin{proof}
用兩種方式數邊。 每個節點都屬於三類之一，所以節點數 $n = n_0 + n_1 + n_2$，邊數是
$n - 1$。 另一方面，每條邊都是某個節點連到它的孩子，分支度 $1$ 的節點貢獻一條，分
支度 $2$ 的貢獻兩條，所以邊數也等於 $n_1 + 2 n_2$。 兩式相等：
$n_0 + n_1 + n_2 - 1 = n_1 + 2 n_2$，整理得 $n_0 = n_2 + 1$。
\end{proof}

圖~\ref{fig:bt-main} 的樹有 $n_2 = 3$、$n_0 = 4$，符合定理。 注意 $n_1$ 被消掉了：
分支度為 $1$ 的節點不管有幾個，都不影響葉節點的個數。 同樣的數邊法用在分支度最多
為 $k$ 的樹，令 $n_i$ 是分支度為 $i$ 的節點數，由
$\sum_i n_i - 1 = \sum_i i\, n_i$ 得到
\[
n_0 = 1 + \sum_{i=2}^{k} (i - 1)\, n_i = 1 + n_2 + 2 n_3 + \cdots + (k - 1) n_k .
\]

\subsection{滿二元樹、完整二元樹、嚴格二元樹}

有三種形狀特別常用，名稱在各本書差異很大，本書用的名稱和定義如下（英文名稱在不同書
裡也不一致，所以一併列出常見的叫法）：

\begin{center}
\begin{tabular}{p{0.2\linewidth}p{0.44\linewidth}p{0.26\linewidth}}
\toprule
本書用的名稱 & 定義 & 其他書的叫法 \\
\midrule
滿二元樹 & 每一層都是滿的；高度 $h$ 時恰有 $2^{h+1} - 1$ 個節點 & perfect binary tree；
有些書叫 full binary tree \\
完整二元樹 & 除了最後一層都是滿的，最後一層的節點都靠左邊 & complete binary tree；
也有人譯作完全二元樹 \\
嚴格二元樹 & 每個節點的分支度都是 $0$ 或 $2$ & CLRS 叫 full binary tree；也叫
proper binary tree \\
\bottomrule
\end{tabular}
\end{center}

\begin{center}
\begin{tikzpicture}[bttree, level 1/.style={sibling distance=18mm}, level 2/.style={sibling distance=9mm}]
\node[btn] {1}
    child {node[btn] {2} child {node[btn] {4}} child {node[btn] {5}}}
    child {node[btn] {3} child {node[btn] {6}} child {node[btn] {7}}};
\node at (0, -2.8) {\small 滿二元樹};
\begin{scope}[xshift=5cm]
\node[btn] {1}
    child {node[btn] {2} child {node[btn] {4}} child {node[btn] {5}}}
    child {node[btn] {3} child {node[btn] {6}} child[missing]};
\node at (0, -2.8) {\small 完整二元樹};
\end{scope}
\begin{scope}[xshift=10cm]
\node[btn] {}
    child {node[btn] {}}
    child {node[btn] {} child {node[btn] {}} child {node[btn] {}}};
\node at (0, -2.8) {\small 嚴格二元樹};
\end{scope}
\end{tikzpicture}
\diagcaption{三種特別的二元樹。 前兩棵的編號是由上到下、由左到右的順序，完整二元樹
的節點恰好是滿二元樹的編號 $1$ 到 $n$。}
\label{fig:bt-shapes}
\end{center}

完整二元樹的另一個說法是：把滿二元樹的節點由上到下、同一層由左到右編號 $1, 2,
\ldots$，$n$ 個節點的完整二元樹就是編號 $1$ 到 $n$ 的那些節點。 滿二元樹是完整二元
樹，反過來不一定。 嚴格二元樹由定理~\ref{thm:bt-n0n2} 有 $n_0 = n_2 + 1$、
$n_1 = 0$，所以節點數 $n = 2 n_2 + 1$ 一定是奇數。

\begin{theorem}
\label{thm:bt-complete-height}
$n$ 個節點的完整二元樹，高度是 $\lfloor \lg n \rfloor$；編號 $i$ 的節點深度是
$\lfloor \lg i \rfloor$。
\end{theorem}

\begin{proof}
深度 $d$ 那一層滿的時候，編號是 $2^d$ 到 $2^{d+1} - 1$（前 $d$ 層共 $2^d - 1$ 個節
點）。 所以編號 $i$ 的節點深度 $d$ 滿足 $2^d \le i < 2^{d+1}$，即
$d = \lfloor \lg i \rfloor$。 最深的節點是編號 $n$，高度是
$\lfloor \lg n \rfloor$。
\end{proof}

\section{表示法}
\label{sec:bt-repr}

\subsection{用陣列存}
\label{sec:bt-array}

把每個節點放在「它在滿二元樹裡的編號」那一格：根放 $A[1]$，編號 $i$ 的節點左孩子放
$A[2i]$、右孩子放 $A[2i + 1]$。 這樣父子關係都能用算的，不用存指標：

\begin{center}
\begin{tabular}{lll}
\toprule
 & 從 $1$ 開始（虛擬碼） & 從 $0$ 開始（C） \\
\midrule
父節點 & $\lfloor i / 2 \rfloor$（$i > 1$）& \texttt{(i - 1) / 2}（\texttt{i > 0}）\\
左孩子 & $2i$ & \texttt{2 * i + 1} \\
右孩子 & $2i + 1$ & \texttt{2 * i + 2} \\
\bottomrule
\end{tabular}
\end{center}

為什麼左孩子是 $2i$？ 由定理~\ref{thm:bt-complete-height} 的證明，深度 $d$ 那一層
的編號從 $2^d$ 開始。 編號 $i$ 是那一層的第 $i - 2^d$ 個（從 $0$ 數），它前面的
$i - 2^d$ 個節點在下一層各有兩個孩子，所以它的左孩子是下一層的第 $2(i - 2^d)$ 個，
編號 $2^{d+1} + 2(i - 2^d) = 2i$。 C 的版本是把每個編號減一：$2(i + 1) - 1 = 2i + 1$。

完整二元樹用這種存法剛好填滿 $A[1..n]$，一格都不浪費，第~8~章的堆積就是這樣存的。
不完整的樹就會有空格。 圖~\ref{fig:bt-main} 的樹存起來是：

\begin{center}
\setlength{\tabcolsep}{4pt}
\begin{tabular}{l*{15}{c}c}
\toprule
$i$ & 1 & 2 & 3 & 4 & 5 & 6 & 7 & 8 & 9 & 10 & 11 & 12 & 13 & 14 & 15 & $\cdots$ 29 \\
\midrule
$A[i]$ & A & B & C & D & E & & F & & & G & & & & H & I & $\cdots$ J \\
\bottomrule
\end{tabular}
\end{center}

$J$ 是 $H$（編號 $14$）的右孩子，編號 $29$，所以這棵 $10$ 個節點的樹要用 $29$ 格，
$16$ 到 $28$ 全空著。 最糟的是歪斜的樹（每個節點都只有右孩子）：$n$ 個節點的編號是
$1, 3, 7, \ldots, 2^n - 1$，要 $2^n - 1$ 格。 一般來說，高度 $h$ 的樹要準備
$2^{h+1} - 1$ 格。 另外，插入或刪除中間的節點時，整棵子樹的編號都要改，得搬很多
格。 所以陣列存法只適合完整（或接近完整）、形狀不常變的樹。

\subsection{用指標存}
\label{sec:bt-linked}

一般的二元樹用節點加指標存：每個節點有 $\mathit{key}$、$\mathit{left}$、
$\mathit{right}$，沒有的孩子是 \NIL；整棵樹用 $T.\mathit{root}$ 指向根，空樹時
$T.\mathit{root} = \NIL$。 需要從節點往上走時（例如第~9~章找後繼），再多存一個
$\mathit{parent}$。 這種存法每個節點的空間固定，插入、刪除只改幾個指標。

$n$ 個節點共有 $2n$ 個孩子指標，其中真的指向節點的恰好是邊數 $n - 1$ 個，所以有
\[
2n - (n - 1) = n + 1
\]
個指標是 \NIL，比一半還多。 第~\ref{sec:bt-thread}~節的引線二元樹就是要把這些空
指標拿來用。

\subsection{一般的樹：左孩子右兄弟}
\label{sec:bt-lcrs}

孩子個數不固定的樹（例如資料夾）怎麼存？ 直覺的做法是每個節點準備 $k$ 個孩子指標，
$k$ 是樹的分支度。 這很浪費：$n$ 個節點共 $nk$ 個指標，用到的只有 $n - 1$ 個，空指
標有 $n(k - 1) + 1$ 個；$k = 10$ 時九成以上的指標是空的。 而且 $k$ 要事先決定，有一
個資料夾特別大就全部跟著變大。

\textbf{左孩子右兄弟}（left-child, right-sibling）表示法只用兩個指標：
\begin{itemize}
\item $\mathit{child}$ 指向第一個（最左邊的）孩子；
\item $\mathit{sibling}$ 指向自己右邊的下一個兄弟。
\end{itemize}
一個節點的所有孩子從 $\mathit{child}$ 開始，沿著 $\mathit{sibling}$ 串成一條串列。
空指標只有 $n + 1$ 個，和二元樹一樣。

把 $\mathit{child}$ 看成左指標、$\mathit{sibling}$ 看成右指標，每一棵有序樹就變成一
棵二元樹，而且根沒有右孩子（根沒有兄弟）。 手算時用三個步驟：
\begin{enumerate}
\item 同一個父節點的孩子，由左到右用線連起來（兄弟相連）；
\item 每個節點只保留到第一個孩子的那條線，到其他孩子的線擦掉；
\item 原本往下的線當左指標、兄弟之間的線當右指標，順時針轉 $45^\circ$ 擺好。
\end{enumerate}

\begin{center}
\begin{tikzpicture}[bttree]
\begin{scope}[level 1/.style={sibling distance=20mm}, level 2/.style={sibling distance=9mm}]
\node[btn] (R) {R}
    child {node[btn] (P) {P} child {node[btn] (T) {T}} child {node[btn] (U) {U}}}
    child {node[btn] (Q) {Q}}
    child {node[btn] (S) {S} child {node[btn] (V) {V}} child {node[btn] (W) {W}} child {node[btn] (X) {X}}};
\draw[btthread, -] (P) -- (Q);
\draw[btthread, -] (Q) -- (S);
\draw[btthread, -] (T) -- (U);
\draw[btthread, -] (V) -- (W);
\draw[btthread, -] (W) -- (X);
\node at (0, -2.9) {\small 一般的樹（虛線是兄弟相連）};
\end{scope}
\begin{scope}[xshift=8.2cm, level 1/.style={sibling distance=20mm}, level 2/.style={sibling distance=14mm}, level 3/.style={sibling distance=10mm}, level 4/.style={sibling distance=9mm}, level distance=8.5mm]
\node[btn] {R}
    child {node[btn] {P}
        child {node[btn] {T} child[missing] child {node[btn] {U}}}
        child {node[btn] {Q}
            child[missing]
            child {node[btn] {S}
                child {node[btn] {V} child[missing] child {node[btn] {W} child[missing] child {node[btn] {X}}}}
                child[missing]}}}
    child[missing];
\node at (0, -5.9) {\small 對應的二元樹};
\end{scope}
\end{tikzpicture}
\diagcaption{左：一棵有序樹，$R$ 有三個孩子 $P$、$Q$、$S$。 右：左孩子右兄弟轉成的
二元樹，左指標是「第一個孩子」，右指標是「下一個兄弟」。 根 $R$ 沒有右孩子。}
\label{fig:bt-lcrs}
\end{center}

反過來，二元樹轉回一般的樹：每個節點 $x$ 的左孩子是 $x$ 的第一個孩子，從左孩子沿著
右指標一路走下去經過的節點，全都是 $x$ 的孩子。

\textbf{森林也可以轉成一棵二元樹}。 森林 $F = (T_1, T_2, \ldots, T_m)$ 對應的二元樹
$B(F)$ 遞迴定義如下：$m = 0$ 時是空樹；否則根是 $T_1$ 的根，左子樹是
$B(\text{$T_1$ 的根的子樹們})$，右子樹是 $B(T_2, \ldots, T_m)$。 手算的做法是先把每
棵樹各自轉好，再把 $T_2$ 的根接成 $T_1$ 的根的右孩子、$T_3$ 的根接成 $T_2$ 的根的右
孩子……，也就是把各棵樹的根當成兄弟。 轉回來的時候，從二元樹的根沿著右指標走，經過
的每個節點都是一棵樹的根。

\begin{center}
\begin{tikzpicture}[bttree]
\begin{scope}[level 1/.style={sibling distance=9mm}]
\node[btn] {A} child {node[btn] {B}} child {node[btn] {C}} child {node[btn] {D}};
\end{scope}
\begin{scope}[xshift=2.2cm]
\node[btn] {E} child {node[btn] {F}};
\end{scope}
\begin{scope}[xshift=4.0cm, level 1/.style={sibling distance=10mm}, level 2/.style={sibling distance=9mm}]
\node[btn] {G} child {node[btn] {H} child {node[btn] {I}} child[missing]} child {node[btn] {J}};
\end{scope}
\node at (2.0, -2.7) {\small 三棵樹的森林};
\begin{scope}[xshift=9.5cm, level 1/.style={sibling distance=30mm}, level 2/.style={sibling distance=14mm}, level 3/.style={sibling distance=10mm}, level distance=8.5mm]
\node[btn] {A}
    child {node[btn] {B} child[missing] child {node[btn] {C} child[missing] child {node[btn] {D}}}}
    child {node[btn] {E}
        child {node[btn] {F}}
        child {node[btn] {G}
            child {node[btn] {H} child {node[btn] {I}} child {node[btn] {J}}}
            child[missing]}};
\node at (0, -4.2) {\small 對應的二元樹};
\end{scope}
\end{tikzpicture}
\diagcaption{森林轉成二元樹。 三棵樹的根 $A$、$E$、$G$ 沿著右指標串在一起；$H$ 的
左孩子 $I$ 是它的第一個孩子，$H$ 的右孩子 $J$ 是它的下一個兄弟。}
\label{fig:bt-forest}
\end{center}

\textbf{走訪的對應}（第~\ref{sec:bt-walk}~節定義前序、中序、後序）。 森林的
\textbf{前序}是：先拜訪 $T_1$ 的根，再用前序走 $T_1$ 的根的子樹們，最後用前序走
$T_2, \ldots, T_m$；\textbf{後序}是：先用後序走 $T_1$ 的根的子樹們，再拜訪 $T_1$ 的
根，最後用後序走 $T_2, \ldots, T_m$。 對照 $B(F)$ 的定義，根是 $T_1$ 的根，左子樹是
$T_1$ 的子樹們，右子樹是其餘的樹，由歸納法立刻得到：
\begin{itemize}
\item 森林的前序 $=$ $B(F)$ 的前序；
\item 森林的後序 $=$ $B(F)$ 的\textbf{中序}。
\end{itemize}
以圖~\ref{fig:bt-forest} 為例，森林的前序是 $A\,B\,C\,D\,E\,F\,G\,H\,I\,J$，後序是
$B\,C\,D\,A\,F\,E\,I\,H\,J\,G$，正好是右邊二元樹的前序和中序。 有些書把森林的後序叫
做森林的「中序」，就是因為這個對應。 二元樹的後序則沒有自然的對應：它會把 $T_1$
的根排在 $T_2, \ldots, T_m$ 所有節點的後面，在原來的森林裡不是任何一種常用的順序。

\textbf{只存父節點}。 如果只需要從節點往上走（例如找路徑、判斷兩個節點是不是在同
一棵樹），可以每個節點只存父節點，根的父節點記成 $0$ 或 \NIL；陣列
$\mathit{parent}[1..n]$ 就能存一整個森林。 第~13~章的互斥集合就是這樣存的。 缺點是
找孩子要掃過所有節點。

\section{走訪}
\label{sec:bt-walk}

\textbf{走訪}（traversal）是把每個節點都拜訪恰好一次。 串列只有一種自然的順序，二
元樹就有好幾種。 最常用的三種都是遞迴的，差別只在「拜訪根」排在走左子樹、右子樹的
前面、中間還是後面：

\begin{itemize}
\item \textbf{前序}（preorder）：根、左子樹、右子樹。
\item \textbf{中序}（inorder）：左子樹、根、右子樹。
\item \textbf{後序}（postorder）：左子樹、右子樹、根。
\end{itemize}

左子樹一律在右子樹前面；把左右對調的三種（例如「根、右、左」）偶爾也會用到，但考
卷上說的前序、中序、後序就是上面這三種。

\begin{pseudo}[label=alg:bt-walks]{三種走訪}
\Procedure{Preorder}{$x$}
    \If{$x \ne \NIL$}
        \State 拜訪 $x$
        \State \Call{Preorder}{$x.\mathit{left}$}
        \State \Call{Preorder}{$x.\mathit{right}$}
    \EndIf
\EndProcedure
\Statex
\Procedure{Inorder}{$x$}
    \If{$x \ne \NIL$}
        \State \Call{Inorder}{$x.\mathit{left}$}
        \State 拜訪 $x$
        \State \Call{Inorder}{$x.\mathit{right}$}
    \EndIf
\EndProcedure
\Statex
\Procedure{Postorder}{$x$}
    \If{$x \ne \NIL$}
        \State \Call{Postorder}{$x.\mathit{left}$}
        \State \Call{Postorder}{$x.\mathit{right}$}
        \State 拜訪 $x$
    \EndIf
\EndProcedure
\end{pseudo}

圖~\ref{fig:bt-main} 的樹：
\begin{center}
\begin{tabular}{ll}
\toprule
前序 & $A\;B\;D\;E\;G\;C\;F\;H\;J\;I$ \\
中序 & $D\;B\;G\;E\;A\;C\;H\;J\;F\;I$ \\
後序 & $D\;G\;E\;B\;J\;H\;I\;F\;C\;A$ \\
\bottomrule
\end{tabular}
\end{center}

\textbf{手算的方法。} 想像從根的左上方出發，貼著樹的外緣逆時針繞一圈（每個節點都
繞過去，空的子樹也當成一小段路繞過）。 每個節點會被經過三次：從它的左邊經過、從它
的下方經過（走完左子樹、要往右子樹去的時候）、從它的右邊經過。 前序是第一次經過時
寫下它，中序是第二次，後序是第三次。 這個方法的好處是三種順序可以一次寫出來，而且
不用在腦子裡記遞迴到哪一層。

另外有兩個檢查答案的方法：前序的第一個一定是根，後序的最後一個一定是根；中序裡根
的左邊恰好是左子樹的節點。 第~\ref{sec:bt-rebuild}~節由走訪序列重建樹，用的就是這
兩件事。

\textbf{時間與空間。} 對有 $n$ 個節點的樹呼叫一次 \proc{Inorder}，每個節點會被呼叫
一次，每個 \NIL 指標也會被呼叫一次（立刻返回）。 由第~\ref{sec:bt-linked}~節，
\NIL 指標有 $n + 1$ 個，所以總共呼叫 $2n + 1$ 次，每次除了遞迴呼叫之外只做常數個步
驟，時間是 $\Theta(n)$。 另外兩種一樣。 遞迴最深時，呼叫堆疊上是從根到最深節點的
一條路徑，再加一個 \NIL，所以額外空間是 $\Theta(h)$：平衡的樹是 $\Theta(\lg n)$，歪
斜的樹是 $\Theta(n)$。

\subsection{C 程式：社團組織圖}

\begin{cbookrecap}[第~14~章 14.5 節、第~15~章 15.10 節、第~19~章 19.14--19.17 節]
節點的結構裡有兩個指向同型結構的指標：\texttt{struct Node *left, *right;}，所以要先
寫 \texttt{typedef struct Node \{ ... \} Node;} 的完整標籤，結構裡面才能用
\texttt{struct Node}。 第~15~章說過，資料的定義是遞迴的（樹要嘛是 \texttt{NULL}，
要嘛是一個節點接上兩棵樹），函式就照著寫：先處理 \texttt{NULL}，再處理
\texttt{t->left} 和 \texttt{t->right}。 第~19~章 19.16 節的三種遍歷只差
\texttt{printf} 放在兩個遞迴呼叫的前、中、後；19.17 節用前序和中序重建過二元樹。

釋放整棵樹要用後序：先釋放兩棵子樹，最後才 \texttt{free(t)}。 先 \texttt{free(t)}
的話，之後再讀 \texttt{t->left} 就是讀已經釋放的記憶體。
\end{cbookrecap}

程式設計社的幹部用一個大寫字母代表，組織圖存成二元樹。 樹用\textbf{延伸前序}輸入：
照前序寫出每個節點，空的子樹也寫成一個 \texttt{.}，這樣一個序列就能唯一決定一棵樹
（第~\ref{sec:bt-rebuild}~節說明）。 圖~\ref{fig:bt-main} 的樹寫成
\texttt{ABD..EG...C.FH.J..I..}，$10$ 個字母加 $11$ 個點，點的個數正好是 \NIL 指標的
個數 $n + 1$。 程式印出三種走訪、高度，以及分支度 $0$、$1$、$2$ 的節點數，用來驗證定
理~\ref{thm:bt-n0n2}。

\begin{example}[label=ex:bt-club]{c}{(club\_tree.c) 社團組織圖的三種走訪}
#include <stdio.h>
#include <stdlib.h>

typedef struct Node {
    char key;                   // 幹部代號，一個大寫字母
    struct Node *left, *right;
} Node;

// 從延伸前序建樹：字母是節點，'.' 是空的子樹，*pos 是讀到的位置
Node *build(const char *s, int *pos)
{
    char c = s[*pos];
    if (c == '\0') {            // 字串不夠長：當成空的
        return NULL;
    }
    (*pos)++;
    if (c == '.') {
        return NULL;
    }
    Node *p = malloc(sizeof *p);
    if (p == NULL) {
        exit(1);
    }
    p->key = c;
    p->left = build(s, pos);    // 先讀完整個左子樹
    p->right = build(s, pos);
    return p;
}

// when 是 0、1、2：在走左子樹之前、之間、之後印出根（前序、中序、後序）
void walk(const Node *t, int when)
{
    if (t == NULL) {
        return;
    }
    if (when == 0) {
        printf(" %c", t->key);
    }
    walk(t->left, when);
    if (when == 1) {
        printf(" %c", t->key);
    }
    walk(t->right, when);
    if (when == 2) {
        printf(" %c", t->key);
    }
}

int height(const Node *t)       // 邊數；空樹是 -1
{
    if (t == NULL) {
        return -1;
    }
    int hl = height(t->left), hr = height(t->right);
    return 1 + (hl > hr ? hl : hr);
}

void count_degree(const Node *t, int cnt[3])    // cnt[d]：分支度 d 的個數
{
    if (t != NULL) {
        cnt[(t->left != NULL) + (t->right != NULL)]++;
        count_degree(t->left, cnt);
        count_degree(t->right, cnt);
    }
}

void free_tree(Node *t)         // 後序：孩子都釋放了才釋放自己
{
    if (t != NULL) {
        free_tree(t->left);
        free_tree(t->right);
        free(t);
    }
}

int main(void)
{
    char s[128];
    while (scanf("%127s", s) == 1) {
        int pos = 0, cnt[3] = {0, 0, 0};
        Node *root = build(s, &pos);
        const char *name[3] = {"pre", "in", "post"};
        printf("%s\n", s);
        for (int w = 0; w < 3; w++) {
            printf("%s:", name[w]);
            walk(root, w);
            printf("\n");
        }
        count_degree(root, cnt);
        printf("height %d, n0 %d, n1 %d, n2 %d\n", height(root), cnt[0], cnt[1], cnt[2]);
        free_tree(root);
    }
    return 0;
}
\end{example}
\iosample{ABD..EG...C.FH.J..I.. \\ A.B.C..}{ABD..EG...C.FH.J..I.. \\ pre: A B D E G C F H J I \\ in: D B G E A C H J F I \\ post: D G E B J H I F C A \\ height 4, n0 4, n1 3, n2 3 \\ A.B.C.. \\ pre: A B C \\ in: A B C \\ post: C B A \\ height 2, n0 1, n1 2, n2 0}

三種走訪寫成同一個函式 \texttt{walk}，參數 \texttt{when} 決定 \texttt{printf}
放在兩個遞迴呼叫的前、中、後。 \texttt{height} 要先知道兩棵子樹的高度才能算出自己
的，是後序；\texttt{free\_tree} 也是後序。 高度用邊數，所以空樹回傳 $-1$，葉節點得
到 $1 + \max(-1, -1) = 0$；題目把高度定義成層數時，空樹改回傳 $0$ 就好。 第二棵樹
$A(-, B(-, C))$ 是只往右長的鏈，前序和中序相同，高度是 $2$。 層序和迴圈版的走訪，下
面兩小節先講虛擬碼，C 的寫法留給實作練習的前幾題。

\subsection{層序走訪}
\label{sec:bt-level}

\textbf{層序}（level order）是一層一層、同一層由左到右拜訪。 層序沒有遞迴的寫法，
要用佇列（第~6~章）：先放進根；每次從佇列拿出一個節點，拜訪它，再把它的左、右孩子
依序放進佇列。

\begin{pseudo}[label=alg:bt-level]{層序走訪}
\Procedure{Level-Order}{$T$}
    \If{$T.\mathit{root} = \NIL$}
        \State \Return
    \EndIf
    \State 建立空佇列 $Q$
    \State \Call{Enqueue}{$Q, T.\mathit{root}$}
    \While{$Q$ 不是空的}
        \State $x \gets$ \Call{Dequeue}{$Q$}
        \State 拜訪 $x$
        \If{$x.\mathit{left} \ne \NIL$}
            \State \Call{Enqueue}{$Q, x.\mathit{left}$}
        \EndIf
        \If{$x.\mathit{right} \ne \NIL$}
            \State \Call{Enqueue}{$Q, x.\mathit{right}$}
        \EndIf
    \EndWhile
\EndProcedure
\end{pseudo}

用圖~\ref{fig:bt-main} 的樹走一次，每一列是一次 \proc{Dequeue}：

\begin{center}
\begin{tabular}{ccl}
\toprule
拿出並拜訪 & 放進佇列 & 之後的佇列（前 $\to$ 後）\\
\midrule
$A$ & $B$, $C$ & $B\;C$ \\
$B$ & $D$, $E$ & $C\;D\;E$ \\
$C$ & $F$ & $D\;E\;F$ \\
$D$ & & $E\;F$ \\
$E$ & $G$ & $F\;G$ \\
$F$ & $H$, $I$ & $G\;H\;I$ \\
$G$ & & $H\;I$ \\
$H$ & $J$ & $I\;J$ \\
$I$ & & $J$ \\
$J$ & & （空）\\
\bottomrule
\end{tabular}
\end{center}

為什麼這樣會一層一層出來？ 佇列裡的節點，深度由前到後不會減少，而且最多只差一層：
拿出深度 $d$ 的節點時放進去的是深度 $d + 1$ 的節點，排在所有已經在佇列裡的節點後
面。 同一層裡，左邊節點的孩子先進佇列，所以也先出來。 每個節點進出佇列各一次，時
間 $\Theta(n)$；佇列最長時大約是最寬的那一層，完整二元樹的最後一層可以有大約 $n/2$
個節點。 第~21~章的廣度優先搜尋就是把層序推廣到圖上。

\subsection{不用遞迴的走訪}
\label{sec:bt-iter}

第~6~章 6.4 節說過，任何遞迴都能改寫成迴圈加上自己管理的堆疊。 樹很深的時候（例如
十萬個節點排成一條鏈），遞迴版會用完系統的呼叫堆疊，這時就需要迴圈版。

\textbf{中序。} 遞迴版進到一個節點時，要先把左子樹整個做完，才輪到它自己，所以「它
自己」要先存起來。 迴圈版從根開始一路往左走，路上的節點都 push 進堆疊；走到底之
後，pop 出來的那個就是下一個要拜訪的（它的左子樹已經做完了），拜訪它，再轉到它的
右子樹，重複同樣的事。

\begin{pseudo}[label=alg:bt-inorder-iter]{中序走訪的迴圈版}
\Procedure{Inorder-Iter}{$T$}
    \State 建立空堆疊 $S$
    \State $x \gets T.\mathit{root}$
    \While{$x \ne \NIL$ \OR $S$ 不是空的}
        \While{$x \ne \NIL$} \Comment{一路往左，路上的節點先存著}
            \State \Call{Push}{$S, x$}
            \State $x \gets x.\mathit{left}$
        \EndWhile
        \State $x \gets$ \Call{Pop}{$S$} \Comment{它的左子樹已經做完}
        \State 拜訪 $x$
        \State $x \gets x.\mathit{right}$
    \EndWhile
\EndProcedure
\end{pseudo}

用圖~\ref{fig:bt-main} 的樹追蹤，每一列是一次拜訪，「推入」是拜訪前第 5--7 行
push 的節點：

\begin{center}
\begin{tabular}{ccl}
\toprule
拜訪 & 拜訪前推入 & 拜訪後的堆疊（底 $\to$ 頂）\\
\midrule
$D$ & $A, B, D$ & $A\;B$ \\
$B$ & & $A$ \\
$G$ & $E, G$ & $A\;E$ \\
$E$ & & $A$ \\
$A$ & & （空）\\
$C$ & $C$ & （空）\\
$H$ & $F, H$ & $F$ \\
$J$ & $J$ & $F$ \\
$F$ & & （空）\\
$I$ & $I$ & （空）\\
\bottomrule
\end{tabular}
\end{center}

拜訪的順序就是中序。 注意 push 的順序是 $A, B, D, E, G, C, F, H, J, I$，正好是前
序：這不是巧合，第~\ref{sec:bt-count}~節會用到這件事。 堆疊裡的節點永遠在同一條從
根往下的路徑上，所以最多 $h + 1$ 個（觀念題第~\ref{q:bt-stack-height}~題）。

\textbf{前序。} 前序比較簡單：pop 出一個節點就拜訪它，再把右孩子、左孩子依序 push。
左孩子後放，所以先出來，整個左子樹會在右孩子 pop 出來之前做完。 這就是第~6~章「把
子工作倒過來 push」的規則。

\textbf{後序。} 後序最麻煩：一個節點要等左右子樹都做完才能拜訪，而它從堆疊頂端被
看到的時候，可能是剛做完左子樹（還要去右子樹），也可能是剛做完右子樹（該拜訪它了）。
分辨的方法是記住上一個拜訪的節點 $\mathit{last}$：頂端節點的右孩子存在而且不是
$\mathit{last}$，表示右子樹還沒做；否則就拜訪它、pop 掉。

\begin{pseudo}[label=alg:bt-postorder-iter]{後序走訪的迴圈版}
\Procedure{Postorder-Iter}{$T$}
    \State 建立空堆疊 $S$
    \State $x \gets T.\mathit{root}$，$\mathit{last} \gets \NIL$
    \While{$x \ne \NIL$ \OR $S$ 不是空的}
        \While{$x \ne \NIL$}
            \State \Call{Push}{$S, x$}
            \State $x \gets x.\mathit{left}$
        \EndWhile
        \State $p \gets$ \Call{Top}{$S$} \Comment{先看，不拿}
        \If{$p.\mathit{right} \ne \NIL$ \AND $p.\mathit{right} \ne \mathit{last}$}
            \State $x \gets p.\mathit{right}$ \Comment{右子樹還沒做}
        \Else
            \State 拜訪 $p$
            \State $\mathit{last} \gets p$
            \State \Call{Pop}{$S$}
        \EndIf
    \EndWhile
\EndProcedure
\end{pseudo}

第 9 行為什麼可以只比較右孩子？ 頂端節點 $p$ 被第二次看到時，剛做完的一定是它的
某棵子樹，而子樹做完時最後拜訪的就是子樹的根（後序最後才拜訪根）。 所以剛做完右子
樹時 $\mathit{last} = p.\mathit{right}$；剛做完左子樹時，$\mathit{last}$ 是
$p.\mathit{left}$（或在 $p$ 沒有左孩子時是更早的節點），不會等於
$p.\mathit{right}$。

\section{由走訪序列重建}
\label{sec:bt-rebuild}

只給一種走訪序列，畫不出唯一的樹：前序 $A\,B$ 可以是 $A(B, -)$，也可以是
$A(-, B)$。 考卷上常給兩種序列，要你把樹畫回來。 這一節假設節點的名稱互不相同。

\subsection{前序加中序}

前序的第一個是根。 在中序裡找到根，它左邊的是左子樹的節點、右邊的是右子樹的節點，
個數也就知道了。 前序在根之後，先是整個左子樹、再是整個右子樹，所以照左子樹的個數
就能把前序切成兩段。 兩邊各自又是「前序加中序」，遞迴下去。

\begin{theorem}
\label{thm:bt-pre-in}
節點名稱互不相同時，前序和中序可以唯一決定一棵二元樹。 後序和中序也可以。
\end{theorem}

\begin{proof}
對節點數 $n$ 做強歸納。 $n = 0$ 只有空樹。 $n \ge 1$ 時，設兩棵樹 $T$、$T'$ 的前序
都是 $p_1 \cdots p_n$、中序都是 $q_1 \cdots q_n$。 兩棵樹的根都是 $p_1$。 設 $p_1 =
q_k$；名稱互不相同，$k$ 是唯一的。 中序是「左子樹、根、右子樹」，所以兩棵樹的左子
樹都由 $q_1 \cdots q_{k-1}$ 組成，中序就是這一串；右子樹的中序都是
$q_{k+1} \cdots q_n$。 左子樹有 $k - 1$ 個節點，前序是「根、左子樹、右子樹」，所以
兩棵樹的左子樹前序都是 $p_2 \cdots p_k$，右子樹前序都是 $p_{k+1} \cdots p_n$。 兩
棵左子樹的前序、中序都相同，節點數小於 $n$，由歸納假設是同一棵；右子樹也一樣。 所
以 $T = T'$。 後序加中序的證明相同，只是根改成後序的最後一個。
\end{proof}

這個證明本身就是演算法。 用一個表 $\mathit{pos}[c]$ 記住名稱 $c$ 在中序裡的位置，找
根就是 $\Theta(1)$：

\begin{pseudo}[label=alg:bt-pre-in]{由前序和中序建樹}
\Function{Build-Pre-In}{$i, j, m$} \Comment{$\mathit{pre}[i..i{+}m{-}1]$ 和 $\mathit{in}[j..j{+}m{-}1]$ 是同一棵子樹}
    \If{$m = 0$}
        \State \Return \NIL
    \EndIf
    \State $k \gets \mathit{pos}[\mathit{pre}[i]]$ \Comment{根在中序的位置}
    \State $\ell \gets k - j$ \Comment{左子樹的節點數}
    \State 配置新節點 $x$，$x.\mathit{key} \gets \mathit{pre}[i]$
    \State $x.\mathit{left} \gets$ \Call{Build-Pre-In}{$i + 1, j, \ell$}
    \State $x.\mathit{right} \gets$ \Call{Build-Pre-In}{$i + 1 + \ell, k + 1, m - 1 - \ell$}
    \State \Return $x$
\EndFunction
\end{pseudo}

每個節點建一次，每次 $\Theta(1)$，總時間 $\Theta(n)$。 不用 $\mathit{pos}$ 表、每次
在中序裡從頭找根的話，最壞（歪斜的樹）是 $\Theta(n^2)$。

\textbf{手算追蹤。} 前序 $A\,B\,D\,E\,G\,C\,F\,H\,J\,I$、中序
$D\,B\,G\,E\,A\,C\,H\,J\,F\,I$。 每一列處理一棵子樹：

\begin{center}
\begin{tabular}{llcll}
\toprule
子樹的前序 & 子樹的中序 & 根 & 左子樹（中序）& 右子樹（中序）\\
\midrule
$ABDEGCFHJI$ & $DBGE\,|\,A\,|\,CHJFI$ & $A$ & $DBGE$ & $CHJFI$ \\
$BDEG$ & $D\,|\,B\,|\,GE$ & $B$ & $D$ & $GE$ \\
$EG$ & $G\,|\,E$ & $E$ & $G$ & 空 \\
$CFHJI$ & $C\,|\,HJFI$ & $C$ & 空 & $HJFI$ \\
$FHJI$ & $HJ\,|\,F\,|\,I$ & $F$ & $HJ$ & $I$ \\
$HJ$ & $H\,|\,J$ & $H$ & 空 & $J$ \\
\bottomrule
\end{tabular}
\end{center}

只剩一個節點的子樹（$D$、$G$、$I$、$J$）不必再列。 由表就能畫出圖~\ref{fig:bt-main}
的樹。 切前序時要數個數：$A$ 的左子樹有 $4$ 個節點，所以前序 $A$ 後面的 $BDEG$ 是
左子樹，剩下的 $CFHJI$ 是右子樹。

\begin{examnote}
重建的題目，每一步寫出「根是誰、左子樹有哪些節點、右子樹有哪些節點」，像上面的表一
樣一棵子樹一列，最後再畫樹。 畫完之後用另一種走訪（例如後序）走一次自己畫的樹，和
題目給的序列核對；題目問「這樣的樹存不存在」時，矛盾出現在哪一步就寫哪一步。
\end{examnote}

\textbf{後序加中序}的做法對稱：根是後序的最後一個，左子樹有 $\ell$ 個節點時，後序
的前 $\ell$ 個是左子樹、接下來到倒數第二個是右子樹。 \textbf{層序加中序}也能唯一決定
樹：層序裡第一個出現的是根；切開中序之後，左子樹的層序就是整個層序裡屬於左子樹的那
些節點，順序不變（觀念題第~\ref{q:bt-level-in}~題）。

\subsection{前序加後序不夠}

前序和後序都不告訴你「哪些節點在左邊」。 最小的反例是兩個節點：
$A(B, -)$ 和 $A(-, B)$ 的前序都是 $A\,B$、後序都是 $B\,A$。 一般來說，問題出在只有
一個孩子的節點：前序、後序看得出它有一個孩子，看不出是左還是右。

用前序和後序還是可以找出每個節點有哪些孩子。 設某棵子樹的前序是 $x\,y \cdots$，則
$y$ 是 $x$ 的第一棵子樹的根；在後序裡，$y$ 的子樹以 $y$ 結尾，由 $y$ 的位置可以算出
這棵子樹有幾個節點。 如果它就是 $x$ 底下全部的節點，$x$ 只有一個孩子，左右無法決定；
否則 $x$ 有兩個孩子，$y$ 在左邊，另一棵子樹的根是前序裡接在 $y$ 的子樹後面的那一個。
所以：

\begin{theorem}
\label{thm:bt-pre-post}
節點名稱互不相同。 給定前序和後序，若至少有一棵二元樹符合，則符合的樹恰好有 $2^m$
棵，$m$ 是分支度為 $1$ 的節點個數（對每一棵符合的樹，$m$ 都一樣）。 特別地，嚴格二
元樹可以由前序和後序唯一決定。
\end{theorem}

證明留作觀念題第~\ref{q:bt-pre-post}~題。 例如前序 $A\,B\,C$、後序 $C\,B\,A$：$A$、
$B$ 都只有一個孩子，共 $2^2 = 4$ 棵，分別是 $B$ 接在 $A$ 的左或右、$C$ 接在 $B$ 的
左或右。

\textbf{延伸前序}也能唯一決定樹，因為它把每個節點的分支度和左右都寫進去了：讀到字
母就是一個節點，接著讀它的左子樹、再讀右子樹；讀到 \texttt{.} 就是空的子樹。 範
例~\ref{ex:bt-club} 的 \texttt{build} 就是照這個定義寫的。

\subsection{C 程式}

下面的程式讀入多組前序和中序，建好樹之後用後序印出來。 兩個序
列可能互相矛盾，例如前序 $A\,B\,C\,D$、中序 $B\,D\,A\,C$ 根本沒有對應的樹，程式要
能發現。 發現的方法很簡單：每次找到的根，在中序裡的位置必須落在目前這一段裡面。

\begin{example}[label=ex:bt-rebuild]{c}{(rebuild.c) 由兩種走訪序列重建}
#include <stdio.h>
#include <stdlib.h>
#include <string.h>

typedef struct Node {
    char key;
    struct Node *left, *right;
} Node;

static int pos[256];            // pos[c]：c 在中序裡的位置，不在就是 -1
static int ok;                  // 序列互相矛盾時設成 0

Node *new_node(char c)
{
    Node *p = malloc(sizeof *p);
    if (p == NULL) {
        exit(1);
    }
    p->key = c;
    p->left = p->right = NULL;
    return p;
}

// 前序 pre[pl..pl+n-1]、中序 in[il..il+n-1] 描述同一棵子樹
Node *from_pre_in(const char *pre, int pl, int il, int n)
{
    if (n == 0 || !ok) {
        return NULL;
    }
    int k = pos[(unsigned char)pre[pl]];    // 根在中序的位置
    if (k < il || k >= il + n) {            // 根不在這一段：矛盾
        ok = 0;
        return NULL;
    }
    int nl = k - il;                        // 左子樹的節點數
    Node *p = new_node(pre[pl]);
    p->left = from_pre_in(pre, pl + 1, il, nl);
    p->right = from_pre_in(pre, pl + 1 + nl, k + 1, n - 1 - nl);
    return p;
}

void postorder(const Node *t)
{
    if (t != NULL) {
        postorder(t->left);
        postorder(t->right);
        putchar(t->key);
    }
}

void free_tree(Node *t)
{
    if (t != NULL) {
        free_tree(t->left);
        free_tree(t->right);
        free(t);
    }
}

int main(void)
{
    char pre[64], in[64];
    while (scanf("%63s %63s", pre, in) == 2) {   // 每行：前序 中序
        int n = strlen(in);
        ok = (int)strlen(pre) == n;
        memset(pos, -1, sizeof pos);
        for (int i = 0; i < n; i++) {
            unsigned char c = in[i];
            if (pos[c] != -1) {
                ok = 0;                     // 中序裡有重複的字母
            }
            pos[c] = i;
        }
        Node *root = from_pre_in(pre, 0, 0, n);
        printf("pre %s, in %s -> ", pre, in);
        if (ok) {
            printf("post ");
            postorder(root);
            printf("\n");
        } else {
            printf("no such tree\n");
        }
        free_tree(root);
    }
    return 0;
}
\end{example}
\iosample{ABDEGCFHJI DBGEACHJFI \\ ABCD BADC \\ ABCD BDAC \\ ABC ABD}{pre ABDEGCFHJI, in DBGEACHJFI -> post DGEBJHIFCA \\ pre ABCD, in BADC -> post BDCA \\ pre ABCD, in BDAC -> no such tree \\ pre ABC, in ABD -> no such tree}

為什麼只要檢查「根落在這一段」就夠了？ 對子樹大小歸納：一段長度為 $1$ 時，根落在
這一段表示前序和中序的這個字母相同；長度更長時，左右兩段各自通過檢查，由歸納假設
建出來的左右子樹的中序分別是那兩段，接上根之後整棵樹的中序就是整段。 所以程式沒有
回報矛盾時，建出來的樹的前序和中序一定和輸入相同。 第三組在處理 $B$ 的右子樹時，前
序剩下 $C$、中序剩下 $D$，$C$ 不在這一段，所以回報 \texttt{no such tree}。 第四組
的前序裡有 $C$，中序裡卻沒有，\texttt{pos} 查到 $-1$，一樣不在範圍內。 後序加中序
的版本只要把根改成後序的最後一個，切法照上面說的改。

\section{引線二元樹}
\label{sec:bt-thread}

第~\ref{sec:bt-linked}~節算過，$n$ 個節點的二元樹有 $n + 1$ 個 \NIL 指標。 另一方
面，中序走訪的迴圈版要一個堆疊，遞迴版要呼叫堆疊，都要 $\Theta(h)$ 的額外空間。
\textbf{引線二元樹}（threaded binary tree）把這些空的指標拿來指向中序的前後節點，
走訪時就不需要堆疊了。

\subsection{中序引線}

\textbf{中序引線}的規則：
\begin{itemize}
\item 節點的左指標原本是 \NIL 時，改成指向它在中序裡的\textbf{前一個}節點（前驅，
predecessor）；
\item 右指標原本是 \NIL 時，改成指向中序裡的\textbf{下一個}節點（後繼，successor）。
\end{itemize}
這種指標叫\textbf{引線}（thread）。 改過之後，一個指標可能指向孩子，也可能是引線，
所以每個節點再多兩個旗標 $\mathit{lthread}$、$\mathit{rthread}$：為真表示對應的指
標是引線。 中序第一個節點沒有前驅、最後一個節點沒有後繼，這兩條引線本書維持
\NIL（旗標仍然設成真）。

\begin{center}
\begin{tikzpicture}[btmain]
\btmaintree
\draw[btthread] (D) to[out=-60, in=-100] (B);
\draw[btthread] (G) to[out=150, in=-45] (B);
\draw[btthread] (G) to[out=-30, in=-90] (E);
\draw[btthread] (E) to[out=-10, in=-20, looseness=1.4] (A);
\draw[btthread] (C) to[out=-170, in=-60] (A);
\draw[btthread] (H) to[out=150, in=-120] (C);
\draw[btthread] (J) to[out=170, in=-90] (H);
\draw[btthread] (J) to[out=-10, in=-100] (F);
\draw[btthread] (I) to[out=-140, in=-60] (F);
\node[font=\scriptsize, text=optStroke] (n1) at (-4.0, -2.75) {\NIL};
\draw[btthread] (D) -- (n1);
\node[font=\scriptsize, text=optStroke] (n2) at (4.2, -3.75) {\NIL};
\draw[btthread] (I) -- (n2);
\end{tikzpicture}
\diagcaption{圖~\ref{fig:bt-main} 的樹加上中序引線（虛線）。 中序是
$D\,B\,G\,E\,A\,C\,H\,J\,F\,I$，例如 $G$ 的左引線指向前驅 $B$、右引線指向後繼 $E$。
$11$ 個空指標裡，$9$ 個變成引線，中序的第一個 $D$ 的左邊和最後一個 $I$ 的右邊仍是
\NIL。}
\label{fig:bt-threads}
\end{center}

\textbf{找後繼。} 節點 $x$ 的中序後繼分兩種情況：
\begin{itemize}
\item $x$ 的右指標是引線：引線指的就是後繼。
\item $x$ 有右子樹：後繼是右子樹裡中序最前面的節點，也就是從 $x.\mathit{right}$ 一
路往左走到底（走到左指標是引線的那個節點）。
\end{itemize}
沒有引線時，第一種情況要往上找祖先（第~9~章用 $\mathit{parent}$ 指標做這件事），
有了引線就一步到位。

\begin{pseudo}[label=alg:bt-insucc]{中序引線樹：找後繼、走訪}
\Function{In-Succ}{$x$}
    \If{$x.\mathit{rthread}$}
        \State \Return $x.\mathit{right}$
    \EndIf
    \State $y \gets x.\mathit{right}$
    \While{\textbf{not} $y.\mathit{lthread}$} \Comment{右子樹裡最左邊的節點}
        \State $y \gets y.\mathit{left}$
    \EndWhile
    \State \Return $y$
\EndFunction
\Statex
\Procedure{Threaded-Inorder}{$T$}
    \State $x \gets T.\mathit{root}$
    \If{$x = \NIL$}
        \State \Return
    \EndIf
    \While{\textbf{not} $x.\mathit{lthread}$} \Comment{中序第一個：一路往左}
        \State $x \gets x.\mathit{left}$
    \EndWhile
    \While{$x \ne \NIL$}
        \State 拜訪 $x$
        \State $x \gets$ \Call{In-Succ}{$x$}
    \EndWhile
\EndProcedure
\end{pseudo}

找前驅 \proc{In-Pred} 完全對稱：左指標是引線就直接回傳，否則從
$x.\mathit{left}$ 一路往右走到底。

\proc{Threaded-Inorder} 只用幾個變數，不用堆疊。 單次 \proc{In-Succ} 最壞要走
$\Theta(h)$ 步，但「一路往左」走過的都是左孩子邊，每條左孩子邊在整趟走訪中只走一次，
所以總時間仍是 $\Theta(n)$，這是第~3~章聚合法的例子。

\textbf{追蹤。} 在圖~\ref{fig:bt-threads} 上找幾個後繼：
\begin{itemize}
\item $\proc{In-Succ}(G)$：$G$ 的右邊是引線，直接得到 $E$。
\item $\proc{In-Succ}(A)$：$A$ 有右子樹，從 $C$ 開始往左走；$C$ 的左指標是引線，
所以答案是 $C$。
\item $\proc{In-Succ}(C)$：從 $F$ 往左走到 $H$；$H$ 的左指標是引線，答案是 $H$。
\item $\proc{In-Succ}(I)$：右邊是引線，值是 \NIL，$I$ 是中序最後一個。
\end{itemize}

\subsection{建立引線與插入節點}

把一棵普通的二元樹加上引線，只要做一次中序走訪，用 $\mathit{prev}$ 記住上一個拜訪
的節點：拜訪 $x$ 時，$x$ 的左邊空著就指向 $\mathit{prev}$，$\mathit{prev}$ 的右邊空
著就指向 $x$。 中序最後一個節點的右邊，等走訪結束後再補上旗標。

\begin{pseudo}[label=alg:bt-make-threads]{加上中序引線}
\Procedure{Make-Threads}{$x$} \Comment{$\mathit{prev}$ 一開始是 \NIL}
    \If{$x = \NIL$}
        \State \Return
    \EndIf
    \State \Call{Make-Threads}{$x.\mathit{left}$}
    \If{$x.\mathit{left} = \NIL$}
        \State $x.\mathit{left} \gets \mathit{prev}$，$x.\mathit{lthread} \gets \TRUE$
    \EndIf
    \If{$\mathit{prev} \ne \NIL$ \AND $\mathit{prev}.\mathit{right} = \NIL$}
        \State $\mathit{prev}.\mathit{right} \gets x$，$\mathit{prev}.\mathit{rthread} \gets \TRUE$
    \EndIf
    \State $\mathit{prev} \gets x$
    \State \Call{Make-Threads}{$x.\mathit{right}$}
\EndProcedure
\end{pseudo}

這裡有個細節：第 6 行改的是 $x.\mathit{left}$，但 $x$ 的左子樹在第 4 行已經做完了，
之後不會再沿著它走；第 8 行改的是 $\mathit{prev}$ 的右指標，那時 $\mathit{prev}$ 的
右子樹（是空的）也已經「做完了」。 所以走訪過程中不會誤把引線當成孩子。

\textbf{插入。} 把新節點 $r$ 接成 $s$ 的右孩子，$s$ 原本的右邊（不管是子樹還是引
線）改接到 $r$ 的右邊。 在中序裡，$r$ 剛好插在 $s$ 和 $s$ 原本的後繼之間。

\begin{pseudo}[label=alg:bt-insert-right]{在引線樹插入右孩子}
\Procedure{Insert-Right}{$s, r$}
    \State $r.\mathit{right} \gets s.\mathit{right}$，$r.\mathit{rthread} \gets s.\mathit{rthread}$
    \State $r.\mathit{left} \gets s$，$r.\mathit{lthread} \gets \TRUE$ \Comment{$r$ 的前驅是 $s$}
    \State $s.\mathit{right} \gets r$，$s.\mathit{rthread} \gets \FALSE$
    \If{\textbf{not} $r.\mathit{rthread}$} \Comment{$r$ 下面接著 $s$ 原本的右子樹}
        \State $t \gets$ \Call{In-Succ}{$r$}
        \State $t.\mathit{left} \gets r$ \Comment{$t$ 原本的左引線指向 $s$}
    \EndIf
\EndProcedure
\end{pseudo}

分兩種情況看。 \textbf{$s$ 的右邊原本是引線}（例如在圖~\ref{fig:bt-threads} 把 $L$
接成 $G$ 的右孩子）：$L$ 繼承 $G$ 的右引線（指向 $E$），左引線指向 $G$，不用改別的
節點。 \textbf{$s$ 原本有右子樹}（例如把 $K$ 接成 $C$ 的右孩子）：$F$ 的子樹整個掛到
$K$ 的右邊。 這時 $K$ 的後繼是 $F$ 的子樹裡最左邊的 $H$，$H$ 的左引線原本指向 $C$，
現在 $C$ 和 $H$ 中間多了 $K$，所以第 6--7 行把它改成指向 $K$。 這一步最容易漏掉。

\begin{examnote}
引線樹的作圖題，實線畫孩子、虛線畫引線，每條引線都要畫到它指的節點，中序第一個的
左引線和最後一個的右引線也要畫出來（指向 \NIL 或題目規定的頭節點）。 畫完之後，數
一數引線是不是 $n + 1$ 條，再沿著引線和「往左走到底」走一次，看是不是中序。
\end{examnote}

\textbf{頭節點。} 很多書的引線樹多放一個\textbf{頭節點}：頭節點的左指標指向根（是孩
子，不是引線），右指標指向自己；中序第一個節點的左引線和最後一個節點的右引線都指向
頭節點，而不是 \NIL。 這樣迴圈不用特別判斷 \NIL：從頭節點出發一直找後繼，回到頭節點
就結束；空樹是頭節點的左指標也指向自己（左旗標設為引線）。 兩種寫法的差別只在兩個端
點，題目用哪一種就照哪一種畫。

\subsection{C 程式}

下面的程式建好圖~\ref{fig:bt-main} 的樹、加上引線，印出每條引線，接著用後繼走一次中
序，再做上面那兩次插入。 最後釋放時也沿著後繼走：先找好下一個再
\texttt{free}，因為 \proc{In-Succ} 只會往後走到中序後面的節點，已經釋放的節點不會
再被讀到。

\begin{example}[label=ex:bt-thread]{c}{(threaded.c) 中序引線二元樹}
#include <stdio.h>
#include <stdlib.h>

typedef struct TNode {
    char key;
    struct TNode *left, *right;
    int lthread, rthread;       // 1：這個指標是引線，不是孩子
} TNode;

TNode *new_node(char c)
{
    TNode *p = malloc(sizeof *p);
    if (p == NULL) {
        exit(1);
    }
    p->key = c;
    p->left = p->right = NULL;
    p->lthread = p->rthread = 0;
    return p;
}

TNode *build(const char *s, int *pos)   // 延伸前序，'.' 是空的子樹
{
    char c = s[*pos];
    if (c == '\0') {
        return NULL;
    }
    (*pos)++;
    if (c == '.') {
        return NULL;
    }
    TNode *p = new_node(c);
    p->left = build(s, pos);
    p->right = build(s, pos);
    return p;
}

// 中序走一次，把空的指標改成引線。*prev 是中序裡的前一個節點。
void make_threads(TNode *p, TNode **prev)
{
    if (p == NULL) {
        return;
    }
    make_threads(p->left, prev);
    if (p->left == NULL) {          // 左邊空著：指向前驅
        p->lthread = 1;
        p->left = *prev;
    }
    if (*prev != NULL && (*prev)->right == NULL) {
        (*prev)->rthread = 1;       // 前一個的右邊空著：指向我
        (*prev)->right = p;
    }
    *prev = p;
    make_threads(p->right, prev);
}

TNode *in_succ(const TNode *x)
{
    if (x->rthread) {
        return x->right;            // 引線直接就是答案（可能是 NULL）
    }
    TNode *y = x->right;            // 右子樹最左邊的節點
    while (!y->lthread) {
        y = y->left;
    }
    return y;
}

// 把新節點 r 接成 s 的右孩子，s 原本的右子樹改接在 r 的右邊
void insert_right(TNode *s, TNode *r)
{
    r->right = s->right;
    r->rthread = s->rthread;
    r->left = s;
    r->lthread = 1;
    s->right = r;
    s->rthread = 0;
    if (!r->rthread) {              // r 下面有原本的右子樹
        TNode *t = in_succ(r);      // 它最左邊的節點，左引線原本指向 s
        t->left = r;
    }
}

void print_inorder(TNode *root)     // 不用堆疊，也不用遞迴
{
    TNode *x = root;
    while (!x->lthread) {
        x = x->left;
    }
    printf("inorder:");
    for (; x != NULL; x = in_succ(x)) {
        printf(" %c", x->key);
    }
    printf("\n");
}

int main(void)
{
    char s[] = "ABD..EG...C.FH.J..I..";
    int pos = 0;
    TNode *root = build(s, &pos), *prev = NULL;
    make_threads(root, &prev);
    prev->rthread = 1;              // 中序最後一個節點：右邊是 NULL 引線
    TNode *first = root;
    while (!first->lthread) {
        first = first->left;
    }
    printf("threads:");
    for (TNode *x = first; x != NULL; x = in_succ(x)) {
        if (x->lthread && x->left != NULL) {
            printf(" %c.L=%c", x->key, x->left->key);
        }
        if (x->rthread && x->right != NULL) {
            printf(" %c.R=%c", x->key, x->right->key);
        }
    }
    printf("\n");
    print_inorder(root);
    TNode *c = root->right, *g = root->left->right->left;
    insert_right(c, new_node('K'));     // C 原本有右孩子 F
    insert_right(g, new_node('L'));     // G 的右邊原本是引線
    print_inorder(root);
    for (TNode *x = first; x != NULL; ) {   // 依中序逐一釋放
        TNode *next = in_succ(x);       // 先找好下一個再 free
        free(x);
        x = next;
    }
    return 0;
}
\end{example}
\iosample{}{threads: D.R=B G.L=B G.R=E E.R=A C.L=A H.L=C J.L=H J.R=F I.L=F \\ inorder: D B G E A C H J F I \\ inorder: D B G L E A C K H J F I}

輸出第一行的 $9$ 條引線和圖~\ref{fig:bt-threads} 的虛線一一對應。 插入 $K$、$L$ 之
後，中序變成 $D\,B\,G\,L\,E\,A\,C\,K\,H\,J\,F\,I$，$L$ 插在 $G$ 和 $E$ 之間，$K$ 插
在 $C$ 和 $H$ 之間。 如果拿掉 \texttt{insert\_right} 最後那個 \texttt{if}，$H$ 的左
引線還指向 $C$，用後繼往前走不會出錯，用 \proc{In-Pred} 往回走時從 $H$ 會直接跳回 $C$，把 $K$ 漏掉。

引線樹的代價是每個節點多兩個旗標，插入、刪除也比較麻煩（刪除要同時修好前驅、後繼
的引線）。 它適合「建好之後主要用來走訪」、又不想用堆疊的情況。 在很多實際的程式
裡，同樣的需求改用 $\mathit{parent}$ 指標來做，第~9~章的二元搜尋樹就是這樣。

\section{二元樹的個數}
\label{sec:bt-count}

$n$ 個節點的二元樹有幾種不同的形狀？ $n = 3$ 時有下面五種。

\begin{center}
\begin{tikzpicture}[bttree, level distance=8mm, level 1/.style={sibling distance=10mm}, level 2/.style={sibling distance=8mm}]
\node[btn] {} child {node[btn] {} child {node[btn] {}} child[missing]} child[missing];
\begin{scope}[xshift=2.6cm]
\node[btn] {} child {node[btn] {} child[missing] child {node[btn] {}}} child[missing];
\end{scope}
\begin{scope}[xshift=5.2cm]
\node[btn] {} child {node[btn] {}} child {node[btn] {}};
\end{scope}
\begin{scope}[xshift=7.8cm]
\node[btn] {} child[missing] child {node[btn] {} child {node[btn] {}} child[missing]};
\end{scope}
\begin{scope}[xshift=10.4cm]
\node[btn] {} child[missing] child {node[btn] {} child[missing] child {node[btn] {}}};
\end{scope}
\end{tikzpicture}
\diagcaption{$3$ 個節點的二元樹共 $5$ 種形狀。 左邊兩種的左子樹有 $2$ 個節點，中間
那種左右各 $1$ 個，右邊兩種的右子樹有 $2$ 個節點。}
\label{fig:bt-five}
\end{center}

\subsection{遞迴式}

令 $b_n$ 是 $n$ 個節點的二元樹的形狀數。 空樹只有一種，$b_0 = 1$。 $n \ge 1$ 時，
根佔掉一個節點，剩下 $n - 1$ 個分給左右子樹：左子樹 $k$ 個、右子樹 $n - 1 - k$ 個，
$k = 0, 1, \ldots, n - 1$。 左子樹有 $b_k$ 種形狀、右子樹有 $b_{n-1-k}$ 種，任意搭
配，所以
\[
b_n = \sum_{k=0}^{n-1} b_k\, b_{n-1-k}, \qquad b_0 = 1 .
\]
算出來是 $b_1 = 1$、$b_2 = 2$、$b_3 = 2 + 1 + 2 = 5$、$b_4 = 14$、$b_5 = 42$。 這正
是第~6~章出現過的卡特蘭數：

\begin{theorem}
\label{thm:bt-catalan}
$n$ 個節點的二元樹共有 $C_n = \dfrac{1}{n+1}\dbinom{2n}{n}$ 種不同的形狀。
\end{theorem}

\begin{proof}
第~6~章 6.1 節證明了：$1, \ldots, n$ 依序進入堆疊時，做得到的輸出順序有 $C_n$ 種（第~6~
章的定理~6.2），而且這個個數滿足同一個遞迴式 $C_n = \sum_{k=0}^{n-1} C_k\,
C_{n-1-k}$、$C_0 = 1$。 $b_n$ 和 $C_n$ 的起始值和遞迴式都相同，對 $n$ 歸納就得到
$b_n = C_n$。
\end{proof}

第~2~章提過的生成函數也能直接解這個遞迴式：令 $B(x) = \sum_n b_n x^n$，遞迴式等於
$B(x) = 1 + x\,B(x)^2$，解二次方程式再展開，得到同一個公式。 本書不展開這個做法。

\subsection{前序固定時，中序就是堆疊的輸出}

二元樹和堆疊的輸出順序之間，其實有一個直接的對應。 把一棵樹的節點依前序編號
$1, 2, \ldots, n$，再看它的中序。 第~\ref{sec:bt-iter}~節的 \proc{Inorder-Iter} 追
蹤時已經看到：push 的順序是前序，pop（拜訪）的順序是中序。 push 的順序一定是前序，
因為 \proc{Inorder-Iter} 第一次走到一個節點時就 push 它，而它走到節點的順序是根、
整個左子樹、整個右子樹（右子樹要等根 pop 出來才會進去）。 所以\textbf{中序一定是
$1, \ldots, n$ 經過一個堆疊能排出的順序}。 反過來，不同的樹有不同的中序（定
理~\ref{thm:bt-pre-in}），所以「前序是 $1, \ldots, n$ 的樹」對應到的中序兩兩不同，
而這兩邊的個數都是 $C_n$，這個對應是一對一而且映成的。 結論：

\begin{theorem}
\label{thm:bt-pre-stack}
給定前序 $p_1 \cdots p_n$（名稱互不相同），序列 $q_1 \cdots q_n$ 是某一棵前序為
$p_1 \cdots p_n$ 的二元樹的中序，充分必要條件是：$p_1, \ldots, p_n$ 依序 push 進堆
疊、中間隨時可以 pop，pop 出來的順序能夠是 $q_1 \cdots q_n$。 這樣的樹恰好有一棵。
\end{theorem}

這給了一個不用畫樹就能判斷「前序和中序有沒有矛盾」的方法：把前序換成 $1, \ldots,
n$，用第~6~章的定理~6.1 檢查中序裡有沒有「大、小、中」的樣子。 例如前序
$A\,B\,C\,D$、中序 $B\,D\,A\,C$ 換成編號是 $2, 4, 1, 3$，其中 $4, 1, 3$ 是「大、小、
中」，所以沒有這樣的樹，和範例~\ref{ex:bt-rebuild} 的第三組一致。

\subsection{其他用同一個數算的東西}

下面幾種東西的個數也都是卡特蘭數，因為它們都能和二元樹一一對應：

\begin{itemize}
\item \textbf{$n$ 個節點的森林}有 $C_n$ 種：第~\ref{sec:bt-lcrs}~節的轉換是一一對
應。
\item \textbf{$n$ 個節點的有序樹}有 $C_{n-1}$ 種：轉成二元樹後根沒有右孩子，根的
左子樹可以是任何 $n - 1$ 個節點的二元樹。
\item \textbf{有 $n$ 個內部節點的嚴格二元樹}有 $C_n$ 種：把任何一棵 $n$ 個節點的
二元樹的每個 \NIL 都換成一個葉節點（共 $n + 1$ 個），就得到一棵嚴格二元樹，原來的
節點都變成內部節點；反過來把葉節點全部拿掉。 這樣加上葉節點的樹叫\textbf{延伸二元
樹}（extended binary tree），新加的葉節點叫外部節點。
\item \textbf{$n + 1$ 個數連乘、加滿括號}的方法有 $C_n$ 種：每種加法對應一棵有 $n$ 個
運算子（內部節點）、$n + 1$ 個運算元（葉節點）的運算式樹，正是上一項的嚴格二元樹。
\end{itemize}

延伸二元樹的前序，把內部節點寫成 \texttt{1}、外部節點寫成 \texttt{0}，就是範
例~\ref{ex:bt-club} 的延伸前序把字母換成 \texttt{1}、點換成 \texttt{0}。 一串
\texttt{0}、\texttt{1} 是不是某棵樹的延伸前序，可以由左往右數「還要讀幾棵子樹」：一
開始是 $1$；讀到 \texttt{1} 用掉一棵、多出兩棵，淨加一；讀到 \texttt{0} 減一。 合法
的條件是這個數在最後一個字元之前都不是 $0$，最後恰好是 $0$，和第~6~章判斷後序是否合
法的條件是同一類。 用電腦枚舉所有 $n$ 個 \texttt{1}、$n + 1$ 個 \texttt{0} 的字串，數
出合法的個數，就能驗證 $b_n = C_n$。 $C_n$ 成長得很快，大約是
$4^n / (n^{3/2}\sqrt{\pi})$：$C_{10} = 16796$，$C_{20}$ 已經超過六十億。

\section{運算式樹}
\label{sec:bt-expr}

只有二元運算子的運算式可以畫成一棵二元樹：葉節點是運算元，內部節點是運算子，運算
子的左右子樹是它的兩個運算元（各自又是一個運算式）。 這叫\textbf{運算式樹}
（expression tree）。 圖~\ref{fig:bt-expr} 是 $(6 + 2) * 5 - 8 / (3 - 1) \text{\textasciicircum} 2$ 的
運算式樹，\textasciicircum{} 是次方。

\begin{center}
\begin{tikzpicture}[bttree, level 1/.style={sibling distance=40mm}, level 2/.style={sibling distance=20mm}, level 3/.style={sibling distance=12mm}, level 4/.style={sibling distance=9mm}]
\node[btop] {$-$}
    child {node[btop] {$*$}
        child {node[btop] {$+$} child {node[btn] {6}} child {node[btn] {2}}}
        child {node[btn] {5}}}
    child {node[btop] {$/$}
        child {node[btn] {8}}
        child {node[btop] {\textasciicircum}
            child {node[btop] {$-$} child {node[btn] {3}} child {node[btn] {1}}}
            child {node[btn] {2}}}};
\end{tikzpicture}
\diagcaption{$(6 + 2) * 5 - 8 / (3 - 1) \text{\textasciicircum} 2$ 的運算式樹。 橘色是運算子（內部節
點），白色是運算元（葉節點）。 越晚算的運算越靠近根，根是最後算的減法。}
\label{fig:bt-expr}
\end{center}

\subsection{三種走訪就是三種寫法}

第~6~章的前序、中序、後序運算式，就是這棵樹的三種走訪：

\begin{center}
\begin{tabular}{ll}
\toprule
前序走訪（前序式）& $-\;*\;+\;6\;2\;5\;/\;8\;\text{\textasciicircum}\;-\;3\;1\;2$ \\
中序走訪（中序式，要加括號）& $(((6 + 2) * 5) - (8 / ((3 - 1) \text{\textasciicircum} 2)))$ \\
後序走訪（後序式）& $6\;2\;+\;5\;*\;8\;3\;1\;-\;2\;\text{\textasciicircum}\;/\;-$ \\
\bottomrule
\end{tabular}
\end{center}

前序和後序不需要括號：樹的形狀完全由記號的順序決定。 中序就不行了，$6 + 2 * 5$ 可
能是 $(6 + 2) * 5$ 也可能是 $6 + (2 * 5)$，兩棵樹不同，中序走訪卻一樣。 所以從樹印出
中序時要加括號。 最簡單的是每個運算都加（上表那樣），比較好看的做法是只加必要的：
子樹的根是運算子，而且
\begin{itemize}
\item 它的優先順序比父節點低，或
\item 優先順序一樣高，但它在「結合方向的另一邊」：父節點由左往右結合（$+$、$-$、$*$、
$/$）而它是右孩子，或父節點由右往左結合（\textasciicircum）而它是左孩子，
\end{itemize}
就要加括號。 照這個規則，圖~\ref{fig:bt-expr} 印成
$(6 + 2) * 5 - 8 / (3 - 1) \text{\textasciicircum} 2$。 第二條規則對 $+$、$*$ 也加括號，所以
$a + (b + c)$ 會保留括號：數學上它等於 $a + b + c$，但它們是兩棵不同的樹。

注意這裡又看到了定理~\ref{thm:bt-n0n2}：只有二元運算子時，運算式樹是嚴格二元樹，
運算元（葉節點）恰好比運算子多一個。 圖~\ref{fig:bt-expr} 有 $6$ 個運算子、$7$ 個
運算元。

另外，第~\ref{sec:bt-rebuild}~節說前序加後序決定不了一般的二元樹，但一個後序式就能
決定運算式樹，因為看記號就知道每個節點有幾個孩子：運算元是葉節點，二元運算子有兩個
孩子。 這和延伸前序的道理一樣，分支度的資訊已經寫在序列裡了。

\subsection{由後序式建樹與求值}

由左往右讀後序式，堆疊裡放的是「已經建好的子樹」：讀到運算元，建一個葉節點 push；
讀到運算子，pop 出兩棵子樹，\textbf{先 pop 的接在右邊}，接到新的運算子節點下面再
push 回去。 讀完時堆疊裡剩下的唯一一棵就是整棵樹。 這和第~6~章的後序求值是同一個
迴圈，只是堆疊裡放的從「數」換成「樹」。

\begin{pseudo}[label=alg:bt-postfix-tree]{由後序式建運算式樹}
\Function{Postfix-To-Tree}{$T, m$} \Comment{$T[1..m]$ 是記號}
    \State 建立空堆疊 $S$
    \For{$i \gets 1$ \To $m$}
        \State 配置新節點 $x$，$x.\mathit{key} \gets T[i]$
        \If{$T[i]$ 是運算子}
            \State $x.\mathit{right} \gets$ \Call{Pop}{$S$}
            \State $x.\mathit{left} \gets$ \Call{Pop}{$S$}
        \Else
            \State $x.\mathit{left} \gets \NIL$，$x.\mathit{right} \gets \NIL$
        \EndIf
        \State \Call{Push}{$S, x$}
    \EndFor
    \State \Return \Call{Pop}{$S$}
\EndFunction
\end{pseudo}

用 $6\;2\;+\;5\;*\;8\;3\;1\;-\;2\;\text{\textasciicircum}\;/\;-$ 追蹤，堆疊裡的每棵子樹用加滿括號的中
序寫：

\begin{center}
\begin{tabular}{cl}
\toprule
記號 & 之後的堆疊（底 $\to$ 頂，每一項是一棵子樹）\\
\midrule
6, 2 & $6$ \quad $2$ \\
$+$ & $(6 + 2)$ \\
5 & $(6 + 2)$ \quad $5$ \\
$*$ & $((6 + 2) * 5)$ \\
8, 3, 1 & $((6 + 2) * 5)$ \quad $8$ \quad $3$ \quad $1$ \\
$-$ & $((6 + 2) * 5)$ \quad $8$ \quad $(3 - 1)$ \\
2 & $((6 + 2) * 5)$ \quad $8$ \quad $(3 - 1)$ \quad $2$ \\
\textasciicircum & $((6 + 2) * 5)$ \quad $8$ \quad $((3 - 1) \text{\textasciicircum} 2)$ \\
$/$ & $((6 + 2) * 5)$ \quad $(8 / ((3 - 1) \text{\textasciicircum} 2))$ \\
$-$ & $(((6 + 2) * 5) - (8 / ((3 - 1) \text{\textasciicircum} 2)))$ \\
\bottomrule
\end{tabular}
\end{center}

\textbf{求值}就是後序走訪：先算出左右子樹的值，再用根的運算子合起來，葉節點的值是
它自己。 由前序式建樹可以用遞迴（第~6~章 6.3 節前序求值的第二種做法，把「算值」換
成「建節點」）；由中序式建樹，可以先用第~6~章的方法轉成後序，再用上面的迴圈。

\subsection{C 程式}

程式一行讀一個後序式，記號用空白隔開，建好樹之後印出前序式、只加必要括號的中序和值。 不合法的後序式（運算元不夠，或最後剩下不只一棵樹）印出一行訊息。

\begingroup\catcode`\^=12
\begin{example}[label=ex:bt-expr]{c}{(expr\_tree.c) 運算式樹}
#include <stdio.h>
#include <stdlib.h>
#include <string.h>

typedef struct ENode {
    char op;                    // 運算子；運算元時是 0
    long long val;              // 運算元的值
    struct ENode *left, *right;
} ENode;

int prec(char op)
{
    return op == '^' ? 3 : (op == '*' || op == '/') ? 2 : 1;
}

void free_tree(ENode *t)
{
    if (t != NULL) {
        free_tree(t->left);
        free_tree(t->right);
        free(t);
    }
}

// 由後序建樹：運算元 push，運算子 pop 兩棵子樹接起來再 push。
// 不是合法的後序時回傳 NULL。
ENode *build(char *line)
{
    ENode *st[64];
    int top = 0, ok = 1;
    for (char *tok = strtok(line, " \n"); tok != NULL && ok; tok = strtok(NULL, " \n")) {
        ENode *p = malloc(sizeof *p);
        if (p == NULL) {
            exit(1);
        }
        p->left = p->right = NULL;
        if (strchr("+-*/^", tok[0]) != NULL && tok[1] == '\0') {
            p->op = tok[0];
            if (top < 2) {
                ok = 0;             // 運算元不夠
            } else {
                p->right = st[--top];   // 先 pop 出來的是右邊
                p->left = st[--top];
            }
        } else {
            p->op = 0;
            p->val = atoll(tok);
        }
        if (top == 64) {
            ok = 0;
        }
        if (ok) {
            st[top++] = p;
        } else {
            free(p);
        }
    }
    if (ok && top == 1) {
        return st[0];
    }
    while (top > 0) {               // 不合法：把建到一半的子樹都釋放
        free_tree(st[--top]);
    }
    return NULL;
}

void prefix(const ENode *t)
{
    if (t->op) {
        printf(" %c", t->op);
        prefix(t->left);
        prefix(t->right);
    } else {
        printf(" %lld", t->val);
    }
}

// 只加必要的括號：子樹的運算子比較低，或一樣高但結合方向不對
void infix(const ENode *t)
{
    if (!t->op) {
        printf("%lld", t->val);
        return;
    }
    const ENode *l = t->left, *r = t->right;
    int right_assoc = (t->op == '^');
    int pl = l->op && (prec(l->op) < prec(t->op) ||
                       (prec(l->op) == prec(t->op) && right_assoc));
    int pr = r->op && (prec(r->op) < prec(t->op) ||
                       (prec(r->op) == prec(t->op) && !right_assoc));
    if (pl) {
        printf("(");
    }
    infix(l);
    if (pl) {
        printf(")");
    }
    printf(" %c ", t->op);
    if (pr) {
        printf("(");
    }
    infix(r);
    if (pr) {
        printf(")");
    }
}

long long eval(const ENode *t)      // 後序：先算兩邊，再算自己
{
    if (!t->op) {
        return t->val;
    }
    long long a = eval(t->left), b = eval(t->right), r = 1;
    switch (t->op) {
    case '+': return a + b;
    case '-': return a - b;
    case '*': return a * b;
    case '/': return b != 0 ? a / b : 0;    // 範例的資料不會除以 0
    default:
        for (long long i = 0; i < b; i++) {
            r *= a;
        }
        return r;
    }
}

int main(void)
{
    char line[256];
    while (fgets(line, sizeof line, stdin) != NULL) {
        line[strcspn(line, "\n")] = '\0';
        printf("postfix: %s\n", line);
        ENode *root = build(line);
        if (root == NULL) {
            printf("not a valid postfix expression\n");
            continue;
        }
        printf("prefix:");
        prefix(root);
        printf("\ninfix: ");
        infix(root);
        printf("\nvalue: %lld\n", eval(root));
        free_tree(root);
    }
    return 0;
}
\end{example}
\iosample{6 2 + 5 * 8 3 1 - 2 ^ / - \\ 2 3 2 ^ ^ \\ 2 3 ^ 2 ^ \\ 9 4 2 - - \\ 7 + 1}{postfix: 6 2 + 5 * 8 3 1 - 2 ^ / - \\ prefix: - * + 6 2 5 / 8 ^ - 3 1 2 \\ infix: (6 + 2) * 5 - 8 / (3 - 1) ^ 2 \\ value: 38 \\ postfix: 2 3 2 ^ ^ \\ prefix: ^ 2 ^ 3 2 \\ infix: 2 ^ 3 ^ 2 \\ value: 512 \\ postfix: 2 3 ^ 2 ^ \\ prefix: ^ ^ 2 3 2 \\ infix: (2 ^ 3) ^ 2 \\ value: 64 \\ postfix: 9 4 2 - - \\ prefix: - 9 - 4 2 \\ infix: 9 - (4 - 2) \\ value: 7 \\ postfix: 7 + 1 \\ not a valid postfix expression}
\endgroup

第二、三行對照了結合方向：$2 \text{\textasciicircum} 3 \text{\textasciicircum} 2$ 由右往左結合，
樹的右子樹是 $3 \text{\textasciicircum} 2$，印的時候不用括號；左子樹是次方時就要加括號，
值也不同（$512$ 和 $64$）。 第四行的減法由左往右結合，右子樹是減法就要加括號。 最
後一行的 \texttt{7 + 1} 讀到 \texttt{+} 時堆疊裡只有一棵子樹，所以不合法。
\texttt{build} 遇到不合法的輸入時，要把已經建好、還在堆疊裡的子樹都釋放掉，不然
AddressSanitizer 會回報記憶體洩漏。

\section{本章常見錯誤}
\label{sec:bt-pitfalls}

\begin{center}
\begin{tabular}{p{0.22\linewidth}p{0.52\linewidth}p{0.16\linewidth}}
\toprule
錯誤 & 例子 & 見 \\
\midrule
高度差一 & 高度用邊數算，卻套「層數」的公式 $2^h - 1$；空樹和葉節點的高度都寫成 $0$ &
第~\ref{sec:bt-terms}~節 \\
名稱混用 & 把「每個節點 $0$ 或 $2$ 個孩子」當成完整二元樹；沒寫出自己用的定義 &
第~\ref{sec:bt-terms}~節 \\
陣列索引沒換 & C 從 $0$ 開始存，左孩子卻寫 \texttt{2 * i}、父節點寫 \texttt{i / 2} &
第~\ref{sec:bt-array}~節 \\
森林轉換弄反 & 第一個孩子接到右指標、兄弟接到左指標；以為森林的後序對應二元樹的後序 &
第~\ref{sec:bt-lcrs}~節 \\
釋放順序錯 & 先 \texttt{free(t)} 再讀 \texttt{t->left}；只 \texttt{free} 根就以為整棵
樹都釋放了 & 第~\ref{sec:bt-walk}~節 \\
後序迴圈版沒記上一個 & 只看頂端節點有沒有右孩子，右子樹做完回來又走進去一次，變成
無窮迴圈 & 第~\ref{sec:bt-iter}~節 \\
重建時切錯段 & 用中序裡根的位置直接切前序；前序要照左子樹的\textbf{個數}切 &
第~\ref{sec:bt-rebuild}~節 \\
以為前序加後序就夠 & 只有一個孩子的節點，左右看不出來，答案只畫一棵 &
第~\ref{sec:bt-rebuild}~節 \\
引線插入漏一步 & 插在有右子樹的節點下面，沒把後繼的左引線改成指向新節點 &
第~\ref{sec:bt-thread}~節 \\
運算式樹左右顛倒 & 由後序建樹時，先 pop 出來的接在左邊，$8\;3\;-$ 變成 $3 - 8$ &
第~\ref{sec:bt-expr}~節 \\
\bottomrule
\end{tabular}
\end{center}

\section{觀念題}
\label{sec:bt-exercises}

題型和研究所考試、期中期末考的筆試相近，有選擇、計算、追蹤和證明。 解答在本章最後。
題目裡的高度、深度、層數都照第~\ref{sec:bt-terms}~節的定義。

\begin{enumerate}
\item （選擇，可能不只一個答案）下列敘述哪些正確？
\begin{enumerate}
\item 有 $30$ 個葉節點的二元樹，恰好有 $29$ 個分支度為 $2$ 的節點。
\item 有 $100$ 個節點的完整二元樹，高度是 $6$。
\item 高度為 $5$ 的二元樹最多有 $63$ 個節點。
\item 有 $n$ 個節點的二元樹用左右孩子指標存，共有 $n - 1$ 個 \NIL 指標。
\item 節點名稱互不相同時，層序加前序可以唯一決定一棵二元樹。
\end{enumerate}

\item （計算）一棵樹的分支度是 $4$，其中分支度為 $1$、$2$、$3$、$4$ 的節點分別有
$6$、$4$、$3$、$2$ 個。
\begin{enumerate}
\item 這棵樹有幾個葉節點？ 共有幾個節點？
\item 每個節點準備 $4$ 個孩子指標來存，共有幾個空指標？ 改用左孩子右兄弟呢？
\end{enumerate}

\item （計算）一棵 $200$ 個節點的完整二元樹存在 $A[1..200]$。
\begin{enumerate}
\item $A[75]$ 的父節點、左孩子、右孩子在哪一格？ 它的深度是多少？
\item 這棵樹的高度是多少？ 有幾個葉節點？ 有幾個節點恰好只有一個孩子？
\item 改用 C 的 \texttt{a[0..199]} 存，\texttt{a[75]} 的父節點和兩個孩子在哪一格？
\item 一棵完整的三元樹（每個內部節點最多三個孩子，由上到下、由左到右編號）存在
$A[1..n]$。 寫出編號 $i$ 的節點的三個孩子和父節點的編號，並說明理由。
\end{enumerate}

\item \label{q:bt-convert}（追蹤）一棵有序樹寫成 $K(L(P, Q), M, N(R(T), S))$，括號裡是由左到右的
孩子（這一題的括號是一般樹的孩子，不分左右）。
\begin{enumerate}
\item 用左孩子右兄弟轉成二元樹，用本章的 $x(L, R)$ 寫法寫出來。
\item 寫出這棵二元樹的前序和中序，並和原來那棵樹的前序、後序比較。
\end{enumerate}

\item （追蹤）一棵二元樹的前序和中序如下。 畫出這棵樹，寫出它的後序和層序。
\[
\text{前序：}M\,K\,H\,C\,A\,F\,J\,B\,D\,L\,E \qquad
\text{中序：}C\,A\,H\,K\,J\,F\,B\,M\,D\,E\,L
\]

\item （計算）一棵二元樹的前序是 $A\,B\,C\,D\,E$，後序是 $C\,E\,D\,B\,A$。
\begin{enumerate}
\item 符合的二元樹有幾棵？ 全部寫出來。
\item 如果另外知道中序是 $C\,B\,E\,D\,A$，是哪一棵？
\end{enumerate}

\item \label{q:bt-pre-post}（證明）證明定理~\ref{thm:bt-pre-post}：節點名稱互不相同，給
定前序和後序，若至少有一棵二元樹符合，則符合的樹恰好有 $2^m$ 棵，$m$ 是分支度為 $1$
的節點個數。

\item \label{q:bt-level-in}（證明）證明：節點名稱互不相同時，層序和中序可以唯一決定
一棵二元樹。

\item \label{q:bt-stack-height}（證明）證明虛擬碼~\ref{alg:bt-inorder-iter} 執行時，
堆疊裡最多同時有 $h + 1$ 個節點（$h$ 是樹高），並舉一棵會達到 $h + 1$ 的樹。

\item （追蹤）用第~5~題的樹。
\begin{enumerate}
\item 加上中序引線，列出每一條引線（例如「$A$ 的左引線指向 $C$」）。 共有幾條？
\item 用 \proc{In-Succ} 找 $K$、$A$、$B$ 的後繼，用 \proc{In-Pred} 找 $M$ 的前驅。
\item 用 \proc{Insert-Right} 把新節點 $N$ 接成 $H$ 的右孩子，寫出改了哪些指標和旗
標。
\item 接著把新節點 $P$ 接成 $D$ 的右孩子，寫出改了哪些指標和旗標，以及最後的中序。
\end{enumerate}

\item （計算）
\begin{enumerate}
\item $6$ 個節點的二元樹有幾種形狀？
\item 其中根的左子樹恰好有 $2$ 個節點的有幾種？
\item 其中高度是 $5$ 的有幾種？
\item 前序是 $A\,B\,C\,D\,E$ 的二元樹有幾棵？ 其中有沒有中序是 $C\,B\,A\,E\,D$ 的？
有沒有中序是 $C\,A\,B\,E\,D$ 的？ 有的話畫出來。
\end{enumerate}

\item （追蹤）運算式 $(a + b * c) / (d - e) - f * g$。
\begin{enumerate}
\item 畫出運算式樹，寫出前序式和後序式。
\item 用虛擬碼~\ref{alg:bt-postfix-tree} 由後序式建樹，寫出每讀一個記號之後的堆疊
（每棵子樹用加滿括號的中序寫）。
\item $a = 4$、$b = 3$、$c = 2$、$d = 7$、$e = 5$、$f = 1$、$g = 3$ 時，用後序走訪
求值。
\end{enumerate}

\end{enumerate}

\section{實作練習}
\label{sec:bt-practice}

下面是 Aki-Judge 上和本章相關的題目。 很多題只要交一個 \texttt{.c} 檔，實作題目規
定的函式，不寫 \texttt{main}；評測程式會在同一次執行裡呼叫好幾次，所以函式不能依賴
上一次呼叫留下的全域變數，用完的樹也要讓評測程式能逐一 \texttt{free}。

\begin{center}
\begin{tabular}{llc}
\toprule
小節 & 題目 & 難度 \\
\midrule
走訪 & \akiref{build-tree-and-traverse}{建立二元樹並進行各種遍歷} & 基本 \\
 & \akiref{height-and-size}{二元樹的高度與節點數} & 暖身 \\
 & \akiref{level-order-traversal}{二元樹的層序遍歷} & 基本 \\
 & \akiref{inorder-without-recursion}{不用遞迴的中序遍歷} & 基本 \\
 & \akiref{parcel-sorter-gate-tree}{包裹分揀閘門樹} & 基本 \\
 & \akiref{heavier-subtree-first}{重的子樹先走} & 基本 \\
 & \akiref{smallest-overloaded-subtrees}{最小超載子樹} & 進階 \\
 & \akiref{trimmed-copy-symmetric-level-widths}{二元樹的截斷複製、對稱子樹與每層寬度} & 進階 \\
 & \akiref{patch-a-binary-tree}{用修補樹改寫二元樹} & 進階 \\
 & \akiref{lighter-side-first-to-list}{輕的先串：把二元樹改接成雙向串列} & 進階 \\
 & \akiref{deal-deck-into-binary-tree}{輪流發牌建二元樹} & 進階 \\
重建 & \akiref{regrow-tree-from-preorder-record}{照前序紀錄把盆栽種回來} & 基本 \\
 & \akiref{rebuild-from-inorder-depths}{由中序深度還原二元樹} & 進階 \\
一般的樹 & \akiref{folder-tree-paths-sizes-moves}{資料夾樹：路徑、容量與搬移} & 基本 \\
 & \akiref{folder-tree-from-one-line}{一行字串還原資料夾樹} & 進階 \\
 & \akiref{bonsai-pruning-match}{盆栽修剪配對} & 進階 \\
運算式樹 & \akiref{gate-tree-settle-time}{邏輯閘樹的訊號穩定時間} & 進階 \\
\bottomrule
\end{tabular}
\end{center}

幾題的提示：

\begin{itemize}
\item 前四題共用同一個二元樹模組，樹用第~\ref{sec:bt-array}~節的陣列存法給（從 $0$
開始，空格是 $-1$），先轉成指標存法再做。 題目規定空樹的高度是 $0$，和本書的 $-1$
不同。 「不用遞迴的中序遍歷」有一條很長的鏈，要用第~\ref{sec:bt-iter}~節的迴圈版。
\item 「重的子樹先走」「最小超載子樹」先用一次後序算好每棵子樹的總和。 「輕的先串」
是中序的變形，遞迴函式回傳子串列的頭和尾最方便。
\item 「照前序紀錄把盆栽種回來」的紀錄和延伸前序一樣能唯一決定樹，照範
例~\ref{ex:bt-club} 的 \texttt{build} 寫。 「由中序深度還原二元樹」：一段中序裡深度
最小的就是這段的根，照第~\ref{sec:bt-rebuild}~節的方式切開。
\item 「一行字串還原資料夾樹」是左孩子右兄弟；「資料夾樹：路徑、容量與搬移」只給父
節點（第~\ref{sec:bt-lcrs}~節最後）。 「邏輯閘樹的訊號穩定時間」是運算式樹的推廣，
每個閘有很多輸入，求值仍然是後序。
\end{itemize}

\section{讀完你應該會}
\label{sec:bt-checklist}

\begin{itemize}
\item 依指定的定義算出深度、層數、高度，用 $n_0 = n_2 + 1$ 和它對 $k$ 元樹的推廣算
節點數，說出滿、完整、嚴格二元樹的差別。
\item 用陣列存完整二元樹，從 $1$ 或 $0$ 開始都能寫出父節點、孩子的編號；把一般的樹
和森林轉成二元樹，再轉回來。
\item 手寫一棵樹的前序、中序、後序、層序，寫出遞迴版和迴圈版的程式，說出各自用的額
外空間。
\item 由前序（或後序、層序）加中序畫出唯一的樹並寫出每一步；說明前序加後序為什麼不
夠，並算出符合的樹有幾棵。
\item 畫出中序引線，用引線找後繼、前驅，並在引線樹上插入節點。
\item 用卡特蘭數算二元樹、森林、加括號方法的個數；把運算式畫成運算式樹，由後序式建
樹並求值。
\end{itemize}

\textbf{延伸閱讀}（書末會整理完整的參考資料）

\begin{itemize}
\item Thomas H. Cormen 等，\textit{Introduction to Algorithms}，第 4 版，10.3 節（用
指標表示有根樹、左孩子右兄弟）、12.1 節（中序走訪與它的時間）、附錄 B.5 節（樹的名
詞；它的 full、complete 和本書的用詞不同，對照著讀正好練習「先看定義」）。
\item Donald E. Knuth，\textit{The Art of Computer Programming}，Vol.~1，第 3 版，
2.3 節：樹的走訪、引線二元樹、森林和二元樹的對應，本章這幾節的經典出處，寫法嚴謹，
適合想看原始證明的人。
\item Ronald L. Graham、Donald E. Knuth、Oren Patashnik，\textit{Concrete
Mathematics}，第 2 版，7.5 節：用生成函數解出卡特蘭數，就是本章略過的那個做法。
\item Antti Laaksonen，\textit{Competitive Programmer's Handbook}，第 14 章：一般的
樹的走訪、直徑和每個節點的子樹資訊，偏實作。
\item 《C 程式設計基礎》第~19~章 19.14--19.18 節：二元搜尋樹的插入、遍歷、重建與集
合運算，是本書第~9~章的預習。
\end{itemize}

\addcontentsline{toc}{section}{觀念題解答}
\section*{觀念題解答}

\begin{enumerate}
\item (a)、(b)、(c) 正確。 (a)：$n_2 = n_0 - 1 = 29$。 (b)：$\lfloor \lg 100 \rfloor
= 6$（$64 \le 100 < 128$）。 (c)：$2^{5+1} - 1 = 63$。 (d) 錯，是 $n + 1$ 個：$2n$ 個
指標裡只有 $n - 1$ 個指向節點。 (e) 錯：$A(B, -)$ 和 $A(-, B)$ 的層序、前序都是
$A\,B$。

\item
\begin{enumerate}
\item $n_0 = 1 + 1 \cdot n_2 + 2 \cdot n_3 + 3 \cdot n_4 = 1 + 4 + 6 + 6 = 17$。 共
$17 + 6 + 4 + 3 + 2 = 32$ 個節點。 驗算：邊數 $31 = 1 \cdot 6 + 2 \cdot 4 + 3 \cdot 3
+ 4 \cdot 2$。
\item 每個節點 $4$ 個指標：$4 \cdot 32 - 31 = 97$ 個空指標，也就是 $n(k - 1) + 1$。
左孩子右兄弟：$2 \cdot 32 - 31 = 33 = n + 1$ 個。
\end{enumerate}

\item
\begin{enumerate}
\item 父節點 $A[37]$（$\lfloor 75/2 \rfloor$），左孩子 $A[150]$，右孩子 $A[151]$，兩
個都 $\le 200$，存在。 深度 $\lfloor \lg 75 \rfloor = 6$（$64 \le 75 < 128$），在第
$7$ 層。
\item 高度 $\lfloor \lg 200 \rfloor = 7$。 有孩子的節點是 $2i \le 200$ 的那些，即
$i = 1, \ldots, 100$，所以葉節點是 $A[101..200]$，共 $100$ 個。 $A[100]$ 的左孩子是
$A[200]$，右孩子 $A[201]$ 不存在，只有它恰好一個孩子。 驗算：$n_2 = 99$，$n_0 = 100
= n_2 + 1$。
\item 父節點 \texttt{a[37]}（\texttt{(75 - 1) / 2}），孩子 \texttt{a[151]}、
\texttt{a[152]}。 \texttt{a[75]} 是第 $76$ 個節點，和 (a) 不是同一個節點。
\item 依編號順序，每個節點依序「生」三個孩子：節點 $1$ 的孩子是 $2, 3, 4$，節點
$2$ 的是 $5, 6, 7$……。 編號 $1$ 到 $i - 1$ 的節點一共生了 $3(i - 1)$ 個孩子，佔
掉編號 $2$ 到 $3(i - 1) + 1$，所以節點 $i$ 的孩子是 $3i - 1$、$3i$、$3i + 1$。 反過
來，編號 $j \ge 2$ 的父節點是滿足 $3i - 1 \le j \le 3i + 1$ 的 $i$，即
$i = \lfloor (j + 1)/3 \rfloor$。 一般的 $k$ 元樹：孩子是 $k(i - 1) + 2$ 到
$ki + 1$，父節點是 $\lfloor (j - 2)/k \rfloor + 1$。
\end{enumerate}

\item
\begin{enumerate}
\item $K(L(P(-, Q),\, M(-, N(R(T, S), -))),\, -)$。 $L$ 的左孩子 $P$ 是它的第一個孩
子，$P$ 的右孩子 $Q$ 是 $P$ 的弟弟；$L$ 的右邊依序串著 $M$、$N$；$R$ 的左孩子 $T$ 是
$R$ 的孩子，$R$ 的右孩子 $S$ 是 $R$ 的弟弟。
\item 二元樹的前序 $K\,L\,P\,Q\,M\,N\,R\,T\,S$，等於原樹的前序。 二元樹的中序
$P\,Q\,L\,M\,T\,R\,S\,N\,K$，等於原樹的後序（$P, Q$ 做完才輪到 $L$；$T$ 做完輪到
$R$，接著 $S$，再來 $N$，最後是根 $K$）。
\end{enumerate}

\item 根是 $M$，中序切成 $CAHKJFB \,|\, M \,|\, DEL$。 左邊 $7$ 個，前序 $KHCAFJB$：
根 $K$，切成 $CAH \,|\, K \,|\, JFB$；$HCA$ 的根 $H$，左邊 $CA$（根 $C$，$A$ 在右）；
$FJB$ 的根 $F$，左 $J$、右 $B$。 右邊 $DLE$：根 $D$，左邊空，$LE$ 的根 $L$，左孩子
$E$。 樹是
\[
M\bigl(K(H(C(-, A), -), F(J, B)),\; D(-, L(E, -))\bigr).
\]
後序 $A\,C\,H\,J\,B\,F\,K\,E\,L\,D\,M$，層序 $M\,K\,D\,H\,F\,L\,C\,J\,B\,E\,A$。

\item
\begin{enumerate}
\item 根 $A$；前序第二個 $B$ 在後序是第 $4$ 個，所以 $B$ 的子樹有 $4$ 個節點，就是
$A$ 底下全部，$A$ 只有一個孩子。 $B$ 的子樹前序 $BCDE$、後序 $CEDB$：$C$ 在後序是第
$1$ 個，子樹只有 $1$ 個節點，不是全部，所以 $B$ 有兩個孩子，左 $C$、右是前序接著的
$D$。 $D$ 的子樹前序 $DE$、後序 $ED$：只有一個孩子 $E$。 $A$、$D$ 各有左右兩種選
擇，共 $2^2 = 4$ 棵：
$A(B(C, D(E, -)), -)$、$A(B(C, D(-, E)), -)$、$A(-, B(C, D(E, -)))$、
$A(-, B(C, D(-, E)))$。
\item 四棵的中序依序是 $CBEDA$、$CBDEA$、$ACBED$、$ACBDE$，所以是第一棵
$A(B(C, D(E, -)), -)$。
\end{enumerate}

\item 對節點數 $n$ 做強歸納，同時證明「每一棵符合的樹，分支度為 $1$ 的節點都是同一
批」，所以 $m$ 有定義。 設前序 $p_1 \cdots p_n$、後序 $q_1 \cdots q_n$，至少有一棵樹
符合。 $n \le 1$ 時恰好一棵，$m = 0$。

$n \ge 2$：任何符合的樹，根都是 $p_1$，而且根至少有一個孩子。 根的「第一棵子樹」（有
左子樹就是左子樹，否則是右子樹）在前序裡緊接在根後面，所以它的根是 $p_2$；在後序裡
它排在最前面，而且以自己的根結尾，所以若 $p_2 = q_s$，這棵子樹恰好有 $s$ 個節點，前
序是 $p_2 \cdots p_{s+1}$、後序是 $q_1 \cdots q_s$。 $s$ 只由兩個序列決定，和是哪一
棵樹無關。
\begin{itemize}
\item $s = n - 1$：根只有一個孩子，那棵子樹可以接在左邊或右邊。 子問題的符合樹由歸
納假設有 $2^{m'}$ 棵（$m'$ 是子樹裡分支度 $1$ 的節點數，對每棵都一樣），乘上左右兩種
選擇，共 $2^{m'+1}$ 棵，而 $m = m' + 1$（根本身分支度 $1$）。
\item $s < n - 1$：根還有其他子孫，所以有兩個孩子，第一棵子樹必定在左邊。 右子樹的
前序是 $p_{s+2} \cdots p_n$、後序是 $q_{s+1} \cdots q_{n-1}$。 左右兩個子問題各自有
$2^{m_L}$、$2^{m_R}$ 棵，任意搭配，共 $2^{m_L + m_R} = 2^m$ 棵（根的分支度是 $2$，不
算）。
\end{itemize}
兩種情況下，分支度為 $1$ 的節點都由序列決定，歸納完成。 嚴格二元樹 $m = 0$，所以唯
一。

\item 對節點數 $n$ 做強歸納，$n = 0$ 只有空樹。 設層序 $L$、中序 $I$。 根是唯一深度
$0$ 的節點，所以是 $L$ 的第一個 $r$。 名稱互不相同，$r$ 在 $I$ 的位置唯一，左邊的是
左子樹的節點集合 $X$、右邊的是右子樹的 $Y$，左子樹的中序就是 $I$ 裡 $r$ 左邊那一段。
接著說明左子樹的層序就是 $L$ 裡屬於 $X$ 的那些項，順序不變。 層序是先依深度、同一深
度再由左到右排；左子樹裡每個節點的深度都比在整棵樹裡少 $1$，同一深度的節點左右順序
也不變，所以把 $L$ 限制到 $X$ 上，正好是左子樹的層序。 右子樹同理。 兩棵子樹的層序和
中序都由 $L$、$I$ 決定，節點數小於 $n$，由歸納假設唯一，所以整棵樹唯一。

\item 證明下面的不變量：第 4 行每次檢查條件時，堆疊由底到頂 $s_1, \ldots, s_t$ 滿足
每個 $s_{i+1}$ 都在 $s_i$ 的左子樹裡，而且 $x \ne \NIL$、$t \ge 1$ 時，$x$ 在 $s_t$
的左子樹裡。 一開始堆疊是空的，成立。 內層迴圈 push $x$ 之後令
$x \gets x.\mathit{left}$：被 push 的 $x$ 原本就在舊頂端的左子樹裡，新的 $x$ 在它的
左子樹裡，維持住。 第 8--10 行 pop 出 $s_t$、令 $x \gets s_t.\mathit{right}$：新的
$x$ 在 $s_t$ 的子樹裡，而 $s_t$ 在 $s_{t-1}$ 的左子樹裡，所以 $x$ 也在新頂端
$s_{t-1}$ 的左子樹裡，維持住。

由不變量，$s_1, s_2, \ldots, s_t$ 一個是下一個的祖先，全都在從 $s_1$ 往下到 $s_t$ 的
一條路徑上，一條路徑最多 $h + 1$ 個節點，所以 $t \le h + 1$。 只往左長的鏈
（$h + 1$ 個節點）會把整條鏈都 push 進去，達到 $h + 1$。

\item 樹是第~5~題的
\[
M\bigl(K(H(C(-, A), -), F(J, B)),\; D(-, L(E, -))\bigr),
\]
中序 $C\,A\,H\,K\,J\,F\,B\,M\,D\,E\,L$。
\begin{enumerate}
\item 空指標有 $n + 1 = 12$ 個，全都變成引線：$C$ 的左引線（\NIL）；$A$ 左 $\to C$、
右 $\to H$；$H$ 右 $\to K$；$J$ 左 $\to K$、右 $\to F$；$B$ 左 $\to F$、右 $\to M$；
$D$ 左 $\to M$；$E$ 左 $\to D$、右 $\to L$；$L$ 的右引線（\NIL）。 指向節點的有
$10$ 條。
\item $K$ 有右子樹，從 $F$ 往左走到 $J$，$J$ 的左邊是引線，後繼是 $J$。 $A$、$B$ 的
右邊是引線，後繼分別是 $H$、$M$。 $M$ 有左子樹，從 $K$ 往右走：$K \to F \to B$，
$B$ 的右邊是引線，前驅是 $B$。
\item $H$ 的右邊原本是引線（指向 $K$）。 $N.\mathit{right} \gets K$、
$N.\mathit{rthread} \gets \TRUE$；$N.\mathit{left} \gets H$、
$N.\mathit{lthread} \gets \TRUE$；$H.\mathit{right} \gets N$、
$H.\mathit{rthread} \gets \FALSE$。 $N$ 的右邊是引線，不用改別的節點。
\item $D$ 原本有右子樹 $L(E, -)$。 $P.\mathit{right} \gets L$、
$P.\mathit{rthread} \gets \FALSE$；$P.\mathit{left} \gets D$、
$P.\mathit{lthread} \gets \TRUE$；$D.\mathit{right} \gets P$。 $P$ 的後繼是 $L$ 的子
樹最左邊的 $E$，$E$ 的左引線原本指向 $D$，改成指向 $P$。 最後的中序是
$C\,A\,H\,N\,K\,J\,F\,B\,M\,D\,P\,E\,L$。
\end{enumerate}

\item
\begin{enumerate}
\item $C_6 = 132$。
\item 左子樹 $2$ 個、右子樹 $3$ 個：$C_2 \cdot C_3 = 2 \cdot 5 = 10$。
\item 高度 $5$ 表示 $6$ 個節點排成一條路徑，除了根，每個節點是父節點的左孩子或右孩
子，$2^5 = 32$ 種。
\item $C_5 = 42$ 棵。 把 $A, \ldots, E$ 編成 $1, \ldots, 5$：$C\,B\,A\,E\,D$ 是
$3, 2, 1, 5, 4$，沒有「大、小、中」，所以有，而且恰好一棵：根 $A$，左邊 $CB$（$B$
的左孩子 $C$），右邊 $ED$（$D$ 的左孩子 $E$），即 $A(B(C, -), D(E, -))$。
$C\,A\,B\,E\,D$ 是 $3, 1, 2, 5, 4$，含 $3, 1, 2$，所以沒有。
\end{enumerate}

\item
\begin{enumerate}
\item 根是最後算的減法，左子樹是 $(a + b * c) / (d - e)$，右子樹是 $f * g$，寫成
$-\bigl(/(+(a, *(b, c)), -(d, e)),\; *(f, g)\bigr)$。 前序式
$-\;/\;+\;a\;*\;b\;c\;-\;d\;e\;*\;f\;g$，後序式
$a\;b\;c\;*\;+\;d\;e\;-\;/\;f\;g\;*\;-$。
\item 依序：$a$；$a$ $b$；$a$ $b$ $c$；$*$ 之後 $a$ $(b * c)$；$+$ 之後
$(a + (b * c))$；$d$、$e$ 之後 $(a + (b * c))$ $d$ $e$；$-$ 之後 $(a + (b * c))$
$(d - e)$；$/$ 之後 $((a + (b * c)) / (d - e))$；$f$、$g$ 之後再多兩項 $f$ $g$；$*$
之後 $((a + (b * c)) / (d - e))$ $(f * g)$；$-$ 之後剩下一棵
$(((a + (b * c)) / (d - e)) - (f * g))$。
\item $b * c = 6$，$a + 6 = 10$，$d - e = 2$，$10 / 2 = 5$，$f * g = 3$，$5 - 3 = 2$。
\end{enumerate}

\end{enumerate}

\chapter{堆積與優先佇列}
\label{ch:heaps}

% 本章的圖共用的樣式與巨集（名稱都以 hp 開頭，避免和其他章衝突）
\tikzset{
    hpn/.style={circle, draw=black!65, minimum size=0.66cm, inner sep=0pt, font=\small},
    hpa/.style={hpn, draw=memStroke, thick, fill=memFill},
    hpb/.style={hpn, draw=optStroke, thick, fill=optFill},
    hpg/.style={hpn, draw=black!30, text=black!40},
    hpe/.style={hpn, draw=black!40, densely dashed},
    hpedge/.style={draw=black!55},
    hpidx/.style={font=\scriptsize, text=black!50},
    hpcn/.style={cell, minimum width=0.8cm, minimum height=0.75cm},
    hpca/.style={hpcn, draw=memStroke, thick, fill=memFill},
    hpcb/.style={hpcn, draw=optStroke, thick, fill=optFill},
    hpcg/.style={hpcn, text=black!35},
}
% \hpheap{根的座標}{最底層的寬度}{值/樣式,值/樣式,...}{節點名稱前綴}
% 依陣列順序（從 1 開始）畫一棵完整二元樹。樣式是 n、a、b、g、e 之一（e 是空的根）。
% \hpshowidxtrue 之後畫的樹，每個節點右上角會標出它在陣列裡的位置。
\newif\ifhpshowidx
\newcommand{\hpheap}[4]{%
    \foreach \hpv/\hps [count=\hpi] in {#3} {%
        \pgfmathtruncatemacro{\hplev}{floor(ln(\hpi)/ln(2) + 0.0001)}%
        \pgfmathsetmacro{\hpx}{(\hpi - 2^\hplev + 0.5) * #2 / 2^\hplev - #2 / 2}%
        \node[hp\hps] (#4\hpi) at ($(#1) + (\hpx, -0.95*\hplev)$) {\hpv};%
        \ifhpshowidx
            \node[hpidx, anchor=south west] at (#4\hpi.north east) {\hpi};%
        \fi
        \ifnum\hpi>1
            \pgfmathtruncatemacro{\hpp}{floor(\hpi/2)}%
            \draw[hpedge] (#4\hpp) -- (#4\hpi);%
        \fi
    }%
}
% \hparray{左端的座標}{值/樣式,...}：一排格子，下面標位置（從 1 開始）。
\newcommand{\hparray}[2]{%
    \foreach \hpv/\hps [count=\hpi] in {#2} {%
        \node[hpc\hps] at ($(#1) + (0.8*\hpi, 0)$) {\hpv};%
        \node[idxlbl] at ($(#1) + (0.8*\hpi, -0.6)$) {\hpi};%
    }%
}

醫院急診室的護理站旁邊有一塊白板，寫著正在等候的病人和他們的檢傷級數：1 級最緊急，
5 級最輕微。 每次有診間空出來，護理師就叫級數最小的人；級數一樣時，先到的先看。
病人隨時會來，等候中的人病情也可能變重，要往前排。 白板上的名單如果照到達順序寫，
每次叫號都要從頭看到尾；如果一直保持由急到緩排好，每來一個人又要找位置插進去，後面
的人全部往下挪。 有沒有一種寫法，讓「來一個人」和「叫一個人」都不必看過整份名單？
往下讀之前，先想一想。

這種「每次只要最大或最小的那一個」的需求叫\textbf{優先佇列}（priority queue）。 第~6~章
的佇列依進入的先後取出，優先佇列則依每個元素的\textbf{鍵}（key）大小取出。 最常用的實
作是\textbf{二元堆積}（binary heap）：把元素放在一個陣列裡，看成一棵形狀固定的二元樹，
只要求「父節點不比子節點小」，不要求整個排好。 這個較弱的要求讓插入和取出最大值都
是 $O(\lg n)$，從一堆亂數建出堆積只要 $\Theta(n)$。

本章前四節講二元堆積和它的變化：陣列存法與上浮、下沉，由下而上建堆積為什麼是線性
時間，帶索引的堆積怎麼支援 decrease-key，以及每個節點有 $d$ 個子節點的 $d$ 元堆積。
第~\ref{sec:hp-mergeable}~節講考試常見的變形：同時取最小和最大的三種結構（最小最大堆積、
deap、對稱最小最大堆積），以及可以快速合併的左傾樹和二項堆積，都以手算追蹤為主。 接著用
第~3~章的位能法，整理費氏堆積各操作的平攤成本。 章末收了幾題研究所入學考試的原題，詳解
在本章最後。

這一章用到第~7~章的樹的名詞（深度、高度、完整二元樹）和第~3~章的平攤分析；C 程式用到
《C 程式設計基礎》第~14~章的結構指標、第~15~章的遞迴和第~19~章的 \texttt{malloc}、
\texttt{realloc}，用到時先在框裡複習。 順帶一提，《C 程式設計基礎》第~19~章 19.2 節
說的 heap 是 \texttt{malloc} 配置記憶體的區域，和這一章的堆積只是同名，兩者沒有關係。

\section{二元堆積}
\label{sec:hp-binary}

\subsection{完整二元樹與陣列存法}

堆積用的是第~7~章的\textbf{完整二元樹}（complete binary tree，CLRS 第 6 章叫 nearly
complete binary tree）：除了最後一層，每一層都填滿，最後一層的節點靠左排。 第~\ref{sec:bt-array}~節說過，把完整二元樹一層一層、
每層由左到右編號 $1, 2, \ldots, n$，就能剛好填滿陣列 $A[1..n]$，不需要任何指標，父子的
位置都用算的：
\[
\proc{Parent}(i) = \lfloor i / 2 \rfloor, \qquad
\proc{Left}(i) = 2i, \qquad
\proc{Right}(i) = 2i + 1.
\]
$\proc{Left}(i) > n$ 表示 $i$ 沒有左子節點，也就是 $i$ 是葉節點；$\proc{Right}(i) > n$ 表示沒
有右子節點。

C 的陣列從 $0$ 開始。 照第~1~章的約定整段改寫，不逐個加減 $1$：位置 $i$ 的父節點是
\texttt{(i - 1) / 2}，子節點是 \texttt{2 * i + 1} 和 \texttt{2 * i + 2}。 這兩組公式是
最常寫錯的地方，寫程式時先拿 $i = 0$、$1$、$2$ 驗一次。

完整二元樹的形狀只由 $n$ 決定，這一章會一直用到下面兩個性質：

\begin{itemize}
\item \textbf{高度}是 $\lfloor \lg n \rfloor + 1$（高度數節點），位置 $i$ 的\textbf{深度}是
$\lfloor \lg i \rfloor$，見定理~\ref{thm:bt-complete-height}。
\item \textbf{葉節點}是位置 $\lfloor n/2 \rfloor + 1$ 到 $n$，共 $\lceil n/2 \rceil$ 個。 因為
$i$ 有子節點的條件是 $2i \le n$，也就是 $i \le \lfloor n/2 \rfloor$。
\end{itemize}

\subsection{堆積性質}

\begin{definition}[最大堆積、最小堆積]
\label{def:hp-heap}
存在陣列 $A[1..n]$ 的完整二元樹，如果每個 $i = 2, 3, \ldots, n$ 都滿足
$A[\lfloor i/2 \rfloor] \ge A[i]$，就叫\textbf{最大堆積}（max-heap）；如果每個 $i$ 都滿足
$A[\lfloor i/2 \rfloor] \le A[i]$，就叫\textbf{最小堆積}（min-heap）。 這個條件叫
\textbf{堆積性質}（heap property）。
\end{definition}

堆積性質只管父子之間。 沿著父節點往上走，值不會變小，所以最大堆積裡，每個節點都是它
那棵子樹的最大值，整個堆積的最大值在 $A[1]$。 兄弟之間、不同子樹之間則沒有任何規定：
堆積不是排好的陣列。 反過來，由大到小排好的陣列一定是最大堆積。

下圖是一個有 $8$ 個元素的最大堆積，右上角的小字是它在陣列裡的位置。

\begin{center}
\begin{tikzpicture}
\hpshowidxtrue
\hpheap{0,0}{6.4}{71/n,62/n,53/n,40/n,47/n,18/n,26/n,35/n}{t}
\hpshowidxfalse
\hparray{4.2,-1.4}{71/n,62/n,53/n,40/n,47/n,18/n,26/n,35/n}
\node[font=\small, anchor=west] at (4.6,-0.55) {陣列 $A[1..8]$};
\end{tikzpicture}
\diagcaption{一個最大堆積的兩種樣子。 左邊是樹，右邊是實際存的陣列。 $A[4] = 40$ 的父節
點是 $A[2] = 62$，子節點是 $A[8] = 35$；$A[5]$ 到 $A[8]$ 是葉節點。}
\label{fig:hp-heap-array}
\end{center}

\subsection{上浮與下沉}

堆積的所有操作都建立在兩個修補動作上。 假設整個陣列只有一個位置 $i$ 可能違反堆積性
質：

\begin{itemize}
\item $A[i]$ 可能比父節點大：把它和父節點交換，再看新位置和新的父節點，一路往上，直到
不比父節點大或是到了根。 這叫\textbf{上浮}（sift-up）。
\item $A[i]$ 可能比子節點小（$i$ 的兩棵子樹本身都是堆積）：把它和\textbf{較大的}子節點交
換，一路往下，直到不比任何子節點小或是到了葉節點。 這叫\textbf{下沉}（sift-down，CLRS
叫 \textsc{Max-Heapify}）。
\end{itemize}

下沉一定要和較大的子節點換。 換上來的那個數要當另一個子節點的父節點，只有較大的那個
才不會比另一邊小。

陣列 $A$ 帶一個欄位 $A.\mathit{size}$，記目前堆積裡有幾個元素；$A[A.\mathit{size}+1]$
以後的格子不屬於堆積。

\begin{pseudo}[label=alg:hp-sift]{上浮與下沉（最大堆積）}
\Procedure{Sift-Up}{$A, i$}
    \While{$i > 1$ \AND $A[\lfloor i/2 \rfloor] < A[i]$}
        \State 交換 $A[i]$ 和 $A[\lfloor i/2 \rfloor]$
        \State $i \gets \lfloor i/2 \rfloor$
    \EndWhile
\EndProcedure
\Statex
\Procedure{Sift-Down}{$A, i$}
    \While{$2i \le A.\mathit{size}$}
        \State $j \gets 2i$ \Comment{左子節點}
        \If{$j + 1 \le A.\mathit{size}$ \AND $A[j + 1] > A[j]$}
            \State $j \gets j + 1$ \Comment{右子節點比較大}
        \EndIf
        \If{$A[i] \ge A[j]$}
            \State \Return
        \EndIf
        \State 交換 $A[i]$ 和 $A[j]$
        \State $i \gets j$
    \EndWhile
\EndProcedure
\end{pseudo}

\proc{Sift-Up} 每一輪往上一層，最多 $\lfloor \lg n \rfloor$ 輪（最深的節點深度是
$\lfloor \lg n \rfloor$，也就是高度減 $1$）；\proc{Sift-Down} 每一輪往下一層，從位置 $i$
開始最多走「$i$ 的高度減 $1$」輪（$i$ 的高度是以 $i$ 為根的子樹的層數），每輪比較兩次
（第 8 行和第 10 行）。
兩者都是 $O(\lg n)$。

為什麼下沉之後整棵子樹都是堆積？ 一輪交換之前，$A[i]$ 的兩棵子樹都是堆積；換完之後
$A[i]$ 是原本三個數裡最大的，所以 $i$ 和兩個子節點之間沒問題，沒被換到的那棵子樹也沒
動過。 唯一可能出問題的是被換下去的那個數，它現在在位置 $j$，而 $j$ 的兩棵子樹仍然是
堆積，這正是下一輪的前提。 上浮的道理一樣：被換下來的父節點原本就不比它的其他子孫小，
換下來以後仍然是那棵子樹的最大值。

\subsection{插入與取出最大值}

\textbf{插入}：把新元素放到陣列最後一格，堆積多一格，形狀仍是完整二元樹；唯一可能違反
的地方是新元素和它的父節點，所以做一次上浮。

\textbf{取出最大值}：最大值在 $A[1]$。 把最後一個元素搬到 $A[1]$，堆積少一格，形狀仍是
完整二元樹；唯一可能違反的地方是新的根，它的兩棵子樹都沒動過，所以做一次下沉。

\begin{pseudo}[label=alg:hp-insert-extract, unbreakable]{插入與取出最大值}
\Procedure{Heap-Insert}{$A, x$} \Comment{假設陣列還有空間}
    \State $A.\mathit{size} \gets A.\mathit{size} + 1$
    \State $A[A.\mathit{size}] \gets x$
    \State \Call{Sift-Up}{$A, A.\mathit{size}$}
\EndProcedure
\Statex
\Function{Heap-Extract-Max}{$A$}
    \If{$A.\mathit{size} = 0$}
        \State \Error{下溢}
    \EndIf
    \State $\mathit{max} \gets A[1]$
    \State $A[1] \gets A[A.\mathit{size}]$
    \State $A.\mathit{size} \gets A.\mathit{size} - 1$
    \State \Call{Sift-Down}{$A, 1$}
    \State \Return $\mathit{max}$
\EndFunction
\end{pseudo}

兩個操作都是 $O(\lg n)$。 只看最大值不取出（\proc{Heap-Maximum}）直接回傳 $A[1]$，是
$\Theta(1)$。 最小堆積把所有的 $<$、$>$ 對調即可，之後提到「最小堆積的上浮、下沉」都是
指這個對調過的版本。

\subsection{手算追蹤}
\label{sec:hp-trace}

把 $35, 62, 18, 47, 71, 26, 53, 40$ 依序插入一個空的最大堆積，再取出兩次最大值。 下表
每一列是一個操作做完之後的陣列，「交換」是上浮或下沉交換了幾次。

\begin{center}
\begin{tabular}{lc*{8}{c}}
\toprule
操作 & 交換 & 1 & 2 & 3 & 4 & 5 & 6 & 7 & 8 \\
\midrule
插入 35 & 0 & 35 \\
插入 62 & 1 & 62 & 35 \\
插入 18 & 0 & 62 & 35 & 18 \\
插入 47 & 1 & 62 & 47 & 18 & 35 \\
插入 71 & 2 & 71 & 62 & 18 & 35 & 47 \\
插入 26 & 1 & 71 & 62 & 26 & 35 & 47 & 18 \\
插入 53 & 1 & 71 & 62 & 53 & 35 & 47 & 18 & 26 \\
插入 40 & 1 & 71 & 62 & 53 & 40 & 47 & 18 & 26 & 35 \\
\midrule
取出 71 & 2 & 62 & 47 & 53 & 40 & 35 & 18 & 26 \\
取出 62 & 1 & 53 & 47 & 26 & 40 & 35 & 18 \\
\bottomrule
\end{tabular}
\end{center}

插入 $71$ 時，它先放在 $A[5]$。 插入之前 $A[1] = 62$、$A[2] = 47$，所以 $71$ 先和
$A[2] = 47$ 交換，再和 $A[1] = 62$ 交換，共兩次；$47$ 被換到 $A[5]$，$62$ 被換到 $A[2]$。
手算時最容易出錯的是拿錯一列的陣列，每一步都要從上一列抄起。

第一次取出最大值的過程畫成樹如下：最後一個元素 $35$ 搬到根，和較大的子節點 $62$ 交換，
再和較大的子節點 $47$ 交換，到了沒有子節點的 $A[5]$ 停下。

\begin{center}
\begin{tikzpicture}
\hpheap{0,0}{4.6}{35/a,62/n,53/n,40/n,47/n,18/n,26/n}{p}
\hpheap{5.4,0}{4.6}{62/n,35/a,53/n,40/n,47/n,18/n,26/n}{q}
\hpheap{10.8,0}{4.6}{62/n,47/n,53/n,40/n,35/a,18/n,26/n}{r}
\node[font=\small] at (0,-3.0) {(a) $35$ 搬到根};
\node[font=\small] at (5.4,-3.0) {(b) 和 $62$ 交換};
\node[font=\small] at (10.8,-3.0) {(c) 和 $47$ 交換，停};
\end{tikzpicture}
\diagcaption{取出 $71$ 之後的下沉。 橘色是正在往下走的 $35$。 每一步比較它和兩個子節點，
和較大的那個交換。}
\label{fig:hp-extract}
\end{center}

\begin{examnote}
堆積的追蹤題通常要求「每插入（刪除）一個元素後畫出堆積」。 樹和陣列可以只畫一種，
但要看清楚題目要哪一種；畫樹時節點的左右位置要對，因為左右決定了陣列的順序。 兩個
子節點相等時，下沉換哪一邊會影響之後的形狀，本書一律換左邊（虛擬碼~\ref{alg:hp-sift} 第 8 行用 $>$），
CLRS 的 \textsc{Max-Heapify} 也是先看左子節點、右邊要嚴格比較大才換右邊，結果一樣。
作答時寫一句你用的規則。
\end{examnote}

\subsection{C 程式}

\begin{cbookrecap}[第~14~章 14.5 節、第~19~章 19.3 節]
\texttt{h->a} 是 \texttt{(*h).a} 的簡寫：\texttt{h} 是指向結構的指標，箭號取出它指到的
那個結構的欄位。 函式要改呼叫者的結構，就傳結構的位址進去。

\texttt{realloc(p, n)} 把 \texttt{p} 指的區塊改成 \texttt{n} 個位元組，內容保留，位址可能
換掉。 要不到記憶體時回傳 \texttt{NULL}，原本的區塊還在，所以要先用另一個指標接，確定
不是 \texttt{NULL} 才蓋回去，不能直接寫 \texttt{p = realloc(p, n)}。 \texttt{realloc}
的第一個參數是 \texttt{NULL} 時，效果和 \texttt{malloc} 一樣。
\end{cbookrecap}

下面的程式從輸入讀操作：\texttt{I x} 插入 $x$，\texttt{D} 取出最大值。 每個操作之後印
出整個陣列；取出時先在括號裡印出取出的數。 陣列滿了就用 \texttt{realloc} 把容量加倍，
第~3~章證明過這樣每次插入的平攤成本仍是 $O(1)$ 加上上浮的成本。

\begin{example}[label=ex:hp-max-heap]{c}{(max\_heap.c) 陣列實作的最大堆積}
#include <stdio.h>
#include <stdlib.h>

typedef struct {
    int *a;       // a[0..size-1] 是堆積，a[0] 是根
    int size;     // 目前有幾個元素
    int cap;      // a 的容量
} MaxHeap;

static void swap(int *x, int *y)
{
    int t = *x;
    *x = *y;
    *y = t;
}

// 位置 i 的值可能比父節點大：一路往上換。
void sift_up(MaxHeap *h, int i)
{
    while (i > 0 && h->a[(i - 1) / 2] < h->a[i]) {
        swap(&h->a[(i - 1) / 2], &h->a[i]);
        i = (i - 1) / 2;
    }
}

// 位置 i 的值可能比子節點小：和較大的子節點交換，一路往下。
void sift_down(MaxHeap *h, int i)
{
    for (;;) {
        int j = 2 * i + 1;              // 左子節點
        if (j >= h->size) {
            return;                     // i 是葉節點
        }
        if (j + 1 < h->size && h->a[j + 1] > h->a[j]) {
            j++;                        // 右子節點比較大
        }
        if (h->a[i] >= h->a[j]) {
            return;
        }
        swap(&h->a[i], &h->a[j]);
        i = j;
    }
}

int heap_insert(MaxHeap *h, int x)
{
    if (h->size == h->cap) {            // 滿了：容量加倍（第 3 章）
        int ncap = h->cap == 0 ? 4 : 2 * h->cap;
        int *na = realloc(h->a, ncap * sizeof *na);
        if (na == NULL) {
            return 0;
        }
        h->a = na;
        h->cap = ncap;
    }
    h->a[h->size] = x;
    h->size++;
    sift_up(h, h->size - 1);
    return 1;
}

// 呼叫前要確定堆積不是空的。
int heap_extract_max(MaxHeap *h)
{
    int top = h->a[0];
    h->size--;
    h->a[0] = h->a[h->size];            // 最後一個元素搬到根
    sift_down(h, 0);
    return top;
}

void print_heap(const MaxHeap *h)
{
    for (int i = 0; i < h->size; i++) {
        printf("%d%c", h->a[i], i + 1 < h->size ? ' ' : '\n');
    }
    if (h->size == 0) {
        printf("empty\n");
    }
}

int main(void)
{
    MaxHeap h = {NULL, 0, 0};
    char op;
    while (scanf(" %c", &op) == 1) {
        if (op == 'I') {
            int x;
            if (scanf("%d", &x) != 1 || !heap_insert(&h, x)) {
                break;
            }
        } else if (op == 'D') {
            if (h.size == 0) {
                printf("underflow\n");
                continue;
            }
            printf("(%d) ", heap_extract_max(&h));
        } else {
            break;
        }
        print_heap(&h);
    }
    free(h.a);
    return 0;
}
\end{example}
\iosample{I 35 \\ I 62 \\ I 18 \\ I 47 \\ I 71 \\ I 26 \\ I 53 \\ I 40 \\ D \\ D}{35 \\ 62 35 \\ 62 35 18 \\ 62 47 18 35 \\ 71 62 18 35 47 \\ 71 62 26 35 47 18 \\ 71 62 53 35 47 18 26 \\ 71 62 53 40 47 18 26 35 \\ (71) 62 47 53 40 35 18 26 \\ (62) 53 47 26 40 35 18}

輸出的每一列和上一小節的表一樣。 \texttt{heap\_extract\_max} 在只剩一個元素時也對：
\texttt{size} 變成 $0$，\texttt{a[0] = a[0]}，\texttt{sift\_down} 發現 \texttt{j >= size}
立刻返回。

\section{由下而上建堆積}
\label{sec:hp-build}

手上已經有 $n$ 個數，要把它們變成堆積。 最直接的做法是一個一個插入，每次 $O(\lg n)$，
共 $O(n \lg n)$。 這個上界是緊的：輸入由小到大時，每個新來的數都比堆積裡所有的數大，
要一路浮到根，第 $i$ 個數交換 $\lfloor \lg i \rfloor$ 次，總共
$\sum_{i=1}^{n} \lfloor \lg i \rfloor = \Theta(n \lg n)$ 次（觀念題第~\ref{hp-q-sum}~題）。

另一種做法是把 $n$ 個數直接放進 $A[1..n]$，再由下往上修補。 葉節點自己就是只有一個節點
的堆積，不用修；從最後一個有子節點的位置 $\lfloor n/2 \rfloor$ 開始，倒著對每個位置做
下沉。 輪到位置 $i$ 時，比它大的位置都已經修好了，所以 $i$ 的兩棵子樹都是堆積，正好
符合下沉的前提。

\begin{pseudo}[label=alg:hp-build]{由下而上建最大堆積}
\Procedure{Build-Max-Heap}{$A, n$}
    \State $A.\mathit{size} \gets n$
    \For{$i \gets \lfloor n/2 \rfloor$ \DownTo $1$}
        \State \Call{Sift-Down}{$A, i$}
    \EndFor
\EndProcedure
\end{pseudo}

\textbf{正確性}用迴圈不變量（第~\ref{sec:ana-invariant}~節）：第 3 行的迴圈每一輪開始時，位置
$i + 1, i + 2, \ldots, n$ 的每一個節點都是一個最大堆積的根。 初始時 $i = \lfloor n/2
\rfloor$，後面的位置都是葉節點，成立。 維持的論證就是上一段的「兩棵子樹都是堆積」。 結
束時 $i = 0$，代進去就是位置 $1$ 是最大堆積的根，也就是整個陣列是最大堆積。 （觀念題
第~\ref{hp-q-build}~題要你把三步完整寫出來。）

\subsection{手算追蹤}

用 $A = \langle 14, 37, 52, 9, 61, 28, 45, 70, 23, 66 \rangle$，$n = 10$，$i$ 從 $5$ 做到
$1$。 每一列是那一次 \proc{Sift-Down} 做完之後的陣列。

\begin{center}
\begin{tabular}{lc*{10}{c}}
\toprule
& 交換 & 1 & 2 & 3 & 4 & 5 & 6 & 7 & 8 & 9 & 10 \\
\midrule
開始 & & 14 & 37 & 52 & 9 & 61 & 28 & 45 & 70 & 23 & 66 \\
$i = 5$ & 1 & 14 & 37 & 52 & 9 & 66 & 28 & 45 & 70 & 23 & 61 \\
$i = 4$ & 1 & 14 & 37 & 52 & 70 & 66 & 28 & 45 & 9 & 23 & 61 \\
$i = 3$ & 0 & 14 & 37 & 52 & 70 & 66 & 28 & 45 & 9 & 23 & 61 \\
$i = 2$ & 1 & 14 & 70 & 52 & 37 & 66 & 28 & 45 & 9 & 23 & 61 \\
$i = 1$ & 3 & 70 & 66 & 52 & 37 & 61 & 28 & 45 & 9 & 23 & 14 \\
\bottomrule
\end{tabular}
\end{center}

$i = 5$ 只有左子節點 $A[10] = 66$；$i = 3$ 的 $52$ 已經比 $28$、$45$ 大，不用動；
$i = 1$ 的 $14$ 一路換到葉節點：先和 $70$ 換，再和 $A[4]$、$A[5]$ 中較大的 $66$ 換，最後和
$A[10] = 61$ 換。 下圖左邊是開始時的樹，右邊是建好的堆積。

\begin{center}
\begin{tikzpicture}
\hpheap{0,0}{6.0}{14/a,37/n,52/n,9/n,61/n,28/n,45/n,70/n,23/n,66/n}{s}
\hpheap{7.4,0}{6.0}{70/n,66/n,52/n,37/n,61/n,28/n,45/n,9/n,23/n,14/a}{e}
\node[font=\small] at (0,-3.6) {開始};
\node[font=\small] at (7.4,-3.6) {建好之後};
\end{tikzpicture}
\diagcaption{由下而上建堆積。 原本在根的 $14$（橘色）最後落到葉節點。}
\label{fig:hp-build}
\end{center}

同一組數如果一個一個插入，得到的是 $\langle 70, 66, 45, 52, 61, 28, 37, 9, 23, 14
\rangle$，交換 $10$ 次；由下而上建只交換 $6$ 次，結果也不一樣。 兩個都是合法的最大堆
積，考題問哪一種方法，就要照那一種做。

\begin{examnote}
題目說「建立堆積」（build a heap）而沒有指定方法時，先看題目有沒有「依序插入」之類的
字。 有就逐一插入；沒有的話，通常指的是由下而上的方法，CLRS 6.3 節的
\textsc{Build-Max-Heap} 就是這一種。 作答時先寫明你用哪一種，再寫追蹤。 由下而上的方法要從 $\lfloor n/2 \rfloor$ 開始倒著做，從 $1$ 往下做是錯的。
\end{examnote}

\subsection{為什麼是 $O(n)$}

令 $h(x)$ 是節點 $x$ 的高度，也就是以 $x$ 為根的子樹的層數，葉節點的高度是 $1$。 從 $x$
往下走到最遠的葉節點要 $h(x) - 1$ 步，所以對位置 $i$ 做 \proc{Sift-Down}，交換次數不超過
$h(i) - 1$，比較次數不超過 $2(h(i) - 1) + 2$。 所以只要證明「所有節點的 $h(x) - 1$ 加
起來是 $O(n)$」。 直覺是：一半的節點是葉節點，高度 $1$，一步都不用走；四分之一高度 $2$，
最多走一步；八分之一高度 $3$，最多走兩步……高度大的節點很少。

\begin{theorem}
\label{thm:hp-build}
有 $n$ 個節點的完整二元樹，$\sum_x \bigl(h(x) - 1\bigr) < 2n$，加總跑過所有節點。 因此
\proc{Build-Max-Heap} 的交換次數小於 $2n$，執行時間是 $\Theta(n)$。
\end{theorem}

\begin{proof}
先算\textbf{滿二元樹}（第~7~章的 perfect binary tree，每一層都填滿）。 設高度是
$k \ge 1$，共 $N = 2^k - 1$ 個節點，深度 $d$ 的節點高度是 $k - d$。 換一個方向數：對每個
$\ell = 2, 3, \ldots, k$，數「高度至少是 $\ell$」的節點有幾個，再把這些個數加起來，每個
高度 $h$ 的節點剛好被數了 $h - 1$ 次（$\ell = 2, \ldots, h$ 各一次）。 高度至少 $\ell$ 的
節點就是深度不超過 $k - \ell$ 的節點，也就是最上面 $k - \ell + 1$ 層，有
$2^{k - \ell + 1} - 1$ 個。 所以
\[
\sum_{x} \bigl(h(x) - 1\bigr)
= \sum_{\ell=2}^{k} \left(2^{k - \ell + 1} - 1\right)
= (2^k - 2) - (k - 1)
= N - k.
\]
一般的 $n$，令 $k = \lfloor \lg n \rfloor + 1$，這是這棵完整二元樹的高度。 它的節點是高度
$k$ 的滿二元樹的位置 $1..n$，每個節點的子樹只會比滿二元樹裡少，高度不會比較大；而且每一
項 $h(x) - 1$ 都不是負的，少掉的節點只會讓總和變小。 所以總和不超過
$2^k - 1 - k < 2^k = 2^{\lfloor \lg n \rfloor + 1} \le 2n$。

交換次數不超過這個總和，所以小於 $2n$；只有 $\lfloor n/2 \rfloor$ 個非葉節點要下沉，比
較次數不超過 $2\sum_x (h(x) - 1) + 2\lfloor n/2 \rfloor < 5n$。 另一方面每個元素至少要看
一次，所以是 $\Theta(n)$。
\end{proof}

其實 $\sum_x (h(x) - 1)$ 恰好是 $n - b(n)$，$b(n)$ 是 $n$ 的二進位表示中 $1$ 的個數（這裡
不證）；換成高度和就是 $2n - b(n)$。 用邊數算高度時每個節點的高度都少 $1$，同一件事
會寫成「高度和是 $n - b(n)$」。 要注意這個論證和「一共 $\lfloor n/2 \rfloor$ 次下沉、每次
$O(\lg n)$，所以 $O(n \lg n)$」並不矛盾：那個估法給的是對的上界，只是不夠緊，因為大部
分的下沉都從接近葉節點的地方開始。

把上浮和下沉寫成 C 實際數一下：輸入由小到大時，由下而上的方法在 $n = 10^6$ 只交換 $999{,}988$ 次，比
$n - b(n) = 999{,}993$ 還少一點；逐一插入則交換了大約 $1.8 \times 10^7$ 次，接近
$n \lg n - 2n$，是前者的 $18$ 倍。

\subsection{堆積排序}

建好最大堆積以後，把根（最大值）和最後一格交換，堆積縮小一格，再對根做一次下沉；重
複 $n - 1$ 次，陣列就由小到大排好了。 這就是\textbf{堆積排序}（heapsort）：原地排序，
最壞情況 $O(n \lg n)$。 它的比較次數、不穩定的原因，以及和其他排序的比較放在第~\ref{sec:so-heap}~節。

\section{優先佇列與帶索引的堆積}
\label{sec:hp-pq}

\textbf{最小優先佇列}（min-priority queue）是一個抽象資料型態（第~6~章），維護一群各自帶
有鍵的元素，支援下面的操作：

\begin{itemize}
\item \proc{Insert}$(Q, x, k)$：加入元素 $x$，鍵是 $k$。
\item \proc{Minimum}$(Q)$：回傳鍵最小的元素，不取出。
\item \proc{Extract-Min}$(Q)$：取出並回傳鍵最小的元素。
\item \proc{Decrease-Key}$(Q, x, k)$：把元素 $x$ 的鍵改成比較小的 $k$。
\item \proc{Delete}$(Q, x)$：把元素 $x$ 從佇列中拿掉。
\end{itemize}

最大優先佇列把「最小」換成「最大」，\proc{Decrease-Key} 換成 \proc{Increase-Key}。 第~25~章
的 Dijkstra 演算法和第~24~章的 Prim 演算法都用最小優先佇列，而且大量使用
\proc{Decrease-Key}，所以這一節用最小堆積來寫。

幾種實作的成本比較如下（$n$ 是佇列裡的元素個數；\proc{Decrease-Key} 假設已經知道 $x$
存在哪裡）。

\begin{center}
\begin{tabular}{lcccc}
\toprule
實作 & \proc{Insert} & \proc{Minimum} & \proc{Extract-Min} & \proc{Decrease-Key} \\
\midrule
未排序的陣列 & $\Theta(1)$ & $\Theta(n)$ & $\Theta(n)$ & $\Theta(1)$ \\
由大到小排好的陣列 & $\Theta(n)$ & $\Theta(1)$ & $\Theta(1)$ & $\Theta(n)$ \\
排好的鏈結串列 & $\Theta(n)$ & $\Theta(1)$ & $\Theta(1)$ & $\Theta(n)$ \\
二元堆積 & $O(\lg n)$ & $\Theta(1)$ & $O(\lg n)$ & $O(\lg n)$ \\
\bottomrule
\end{tabular}
\end{center}

排好的陣列由大到小放，最小值在最後一格，取出時不用挪動其他元素。 前三種總有一個操作
要 $\Theta(n)$；堆積把每個操作都壓到 $O(\lg n)$，做 $n$ 次插入、$n$ 次取出的總成本從
$\Theta(n^2)$ 降到 $O(n \lg n)$。

\subsection{decrease-key 要先知道元素在哪裡}

堆積不支援搜尋：在最小堆積裡找鍵 $k$，只能跳過根比 $k$ 大的子樹，其他的都得往下看，
最壞要看 $\Theta(n)$ 格（例如 $k$ 比所有的鍵都大）。 所以
\proc{Decrease-Key}$(Q, x, k)$ 的 $x$ 不能只是一個鍵，必須能直接找到它在陣列的位置。 常
見的做法是給每個元素一個編號 $v = 1, 2, \ldots, m$（例如圖的點、病人的掛號序號），陣列
$H.A$ 裡存的是編號，另外開兩個以編號為索引的陣列：

\begin{itemize}
\item $H.\mathit{key}[v]$：元素 $v$ 目前的鍵。
\item $H.\mathit{pos}[v]$：元素 $v$ 目前在 $H.A$ 的哪一格；不在堆積裡時是 $0$。
\end{itemize}

這叫\textbf{帶索引的堆積}（indexed heap）。 每次交換 $H.A$ 的兩格，兩個元素的
$\mathit{pos}$ 都要跟著改，所以把交換寫成一個程序。 鍵變小時，元素只可能比父節點小，
做一次上浮；刪除任意一個元素時，用最後一個元素補進它的位置，補進來的數可能比新的父節
點小，也可能比子節點大，兩個方向都要試。

\begin{pseudo}[label=alg:hp-indexed]{帶索引的最小堆積}
\Procedure{Exchange}{$H, i, j$}
    \State 交換 $H.A[i]$ 和 $H.A[j]$
    \State $H.\mathit{pos}[H.A[i]] \gets i$
    \State $H.\mathit{pos}[H.A[j]] \gets j$
\EndProcedure
\Statex
\Procedure{Decrease-Key}{$H, v, k$}
    \If{$k > H.\mathit{key}[v]$}
        \State \Error{新的鍵比原來大}
    \EndIf
    \State $H.\mathit{key}[v] \gets k$
    \State $i \gets H.\mathit{pos}[v]$
    \While{$i > 1$ \AND $H.\mathit{key}[H.A[\lfloor i/2 \rfloor]] > H.\mathit{key}[H.A[i]]$}
        \State \Call{Exchange}{$H, i, \lfloor i/2 \rfloor$}
        \State $i \gets \lfloor i/2 \rfloor$
    \EndWhile
\EndProcedure
\Statex
\Procedure{Heap-Delete}{$H, v$}
    \State $i \gets H.\mathit{pos}[v]$
    \State \Call{Exchange}{$H, i, H.\mathit{size}$}
    \State $H.\mathit{size} \gets H.\mathit{size} - 1$
    \State $H.\mathit{pos}[v] \gets 0$
    \If{$i \le H.\mathit{size}$}
        \State \Call{Min-Sift-Up}{$H, i$}
        \State \Call{Min-Sift-Down}{$H, i$}
    \EndIf
\EndProcedure
\end{pseudo}

\proc{Min-Sift-Up}、\proc{Min-Sift-Down} 是虛擬碼~\ref{alg:hp-sift} 把比較反過來、比較的
對象改成 $H.\mathit{key}[H.A[\cdot]]$、交換改成 \proc{Exchange} 的版本。 第 19、20 行兩
個都呼叫，但最多只有一個會真的動：補進來的數如果比父節點小，往上浮，浮上去的路徑上原
本的數都比它大，不會再需要下沉；反之亦然。 \proc{Extract-Min} 就是
\proc{Heap-Delete}$(H, H.A[1])$。 所有操作都是 $O(\lg n)$，多花的空間是兩個長度 $m$ 的
陣列。

有些實作不做 \proc{Decrease-Key}，鍵變小時直接再插入一份新的，舊的那份留在堆積裡，取出
時發現是過期的就丟掉。 這樣寫比較短，代價是堆積裡最多可能有和操作次數一樣多的元素。
寫 Dijkstra 演算法（第~25~章）時兩種都常見。

\subsection{例子：急診室的候診名單}

回到章首的急診室。 病人依到達順序編號 $1, 2, 3, \ldots$，鍵是（級數，編號）這一對：級
數小的優先，同級時編號小（先到）的優先。 只用級數當鍵是不夠的，堆積不會保持相同鍵的
先後順序（下沉時同一層的元素會被換來換去），所以把到達順序也放進鍵裡。

下面的程式處理一個晚上的紀錄。 \texttt{A} 是新病人到達，\texttt{T} 是叫號看診，
\texttt{U 3 1} 是 $3$ 號病人病情變化、改成 $1$ 級，\texttt{L 2} 是 $2$ 號病人沒等到就離
開。 看診和離開都是把一個病人從堆積拿掉，程式用同一個函式 \texttt{er\_remove}，也就是
虛擬碼的 \proc{Heap-Delete}。 每個操作之後印出堆積陣列，\texttt{6(2)} 表示 $6$ 號病人、
$2$ 級。

\begin{example}[label=ex:hp-er-queue]{c}{(er\_queue.c) 用帶索引的最小堆積管理候診名單}
#include <stdio.h>
#include <stdlib.h>

// 帶索引的最小堆積。病人編號 1, 2, 3, ... 依到達順序給。
// 比較的鍵是 (級數, 編號)：級數小的先看，同級的先到先看。
typedef struct {
    int *heap;    // heap[0..size-1]：放病人編號
    int *pos;     // pos[v]：病人 v 在 heap 的哪一格
    int *level;   // level[v]：病人 v 的檢傷級數（1 最緊急）
    int size;
} ERQueue;

// 病人 u 是否應該排在 v 前面
static int before(const ERQueue *q, int u, int v)
{
    return q->level[u] < q->level[v] ||
           (q->level[u] == q->level[v] && u < v);
}

// 交換 heap[i]、heap[j]，兩人的 pos 也要跟著改
static void swap_at(ERQueue *q, int i, int j)
{
    int u = q->heap[i], v = q->heap[j];
    q->heap[i] = v;
    q->heap[j] = u;
    q->pos[v] = i;
    q->pos[u] = j;
}

static void sift_up(ERQueue *q, int i)
{
    while (i > 0 && before(q, q->heap[i], q->heap[(i - 1) / 2])) {
        swap_at(q, i, (i - 1) / 2);
        i = (i - 1) / 2;
    }
}

static void sift_down(ERQueue *q, int i)
{
    for (;;) {
        int j = 2 * i + 1;
        if (j >= q->size) {
            return;
        }
        if (j + 1 < q->size &&
            before(q, q->heap[j + 1], q->heap[j])) {
            j++;
        }
        if (!before(q, q->heap[j], q->heap[i])) {
            return;
        }
        swap_at(q, i, j);
        i = j;
    }
}

void er_arrive(ERQueue *q, int v, int lv)
{
    q->level[v] = lv;
    q->heap[q->size] = v;
    q->pos[v] = q->size;
    q->size++;
    sift_up(q, q->size - 1);
}

// 看診或離開：把病人 v 拿掉，最後一個補進他的位置，再往上或往下調整。
void er_remove(ERQueue *q, int v)
{
    int i = q->pos[v];
    q->size--;
    if (i != q->size) {
        swap_at(q, i, q->size);
        sift_up(q, i);               // 兩個方向最多只有一個會動
        sift_down(q, i);
    }
}

void er_upgrade(ERQueue *q, int v, int lv)   // decrease-key
{
    q->level[v] = lv;
    sift_up(q, q->pos[v]);
}

int main(void)
{
    // 一個晚上的紀錄。A x：新病人到達，x 級；T：叫號看診；
    // U x y：x 號病人改成 y 級；L x：x 號病人離開。
    struct { char op; int x, y; } ops[] = {
        {'A', 4, 0}, {'A', 4, 0}, {'A', 5, 0}, {'A', 3, 0},
        {'A', 4, 0}, {'A', 2, 0}, {'A', 3, 0}, {'L', 2, 0},
        {'U', 3, 1}, {'A', 3, 0}, {'T', 0, 0}, {'T', 0, 0}};
    int m = sizeof ops / sizeof ops[0];      // 病人最多 m 個
    size_t sz = (m + 1) * sizeof(int);       // 編號從 1 用到 m
    ERQueue q = {malloc(sz), malloc(sz), malloc(sz), 0};
    if (q.heap == NULL || q.pos == NULL || q.level == NULL) {
        free(q.heap);
        free(q.pos);
        free(q.level);
        return 1;
    }
    int next = 1;                            // 下一個病人的編號
    for (int k = 0; k < m; k++) {
        int x = ops[k].x, y = ops[k].y;
        if (ops[k].op == 'A') {
            er_arrive(&q, next, x);
            printf("arrive %d(%d):", next, x);
            next++;
        } else if (ops[k].op == 'T') {       // 看診時一定有人
            int v = q.heap[0];
            er_remove(&q, v);
            printf("treat %d(%d):", v, q.level[v]);
        } else if (ops[k].op == 'U') {
            er_upgrade(&q, x, y);
            printf("upgrade %d(%d):", x, y);
        } else {
            er_remove(&q, x);
            printf("leave %d:", x);
        }
        for (int i = 0; i < q.size; i++) {
            printf(" %d(%d)", q.heap[i], q.level[q.heap[i]]);
        }
        printf("\n");
    }
    free(q.heap);
    free(q.pos);
    free(q.level);
    return 0;
}
\end{example}
\iosample{}{arrive 1(4): 1(4) \\ arrive 2(4): 1(4) 2(4) \\ arrive 3(5): 1(4) 2(4) 3(5) \\ arrive 4(3): 4(3) 1(4) 3(5) 2(4) \\ arrive 5(4): 4(3) 1(4) 3(5) 2(4) 5(4) \\ arrive 6(2): 6(2) 1(4) 4(3) 2(4) 5(4) 3(5) \\ arrive 7(3): 6(2) 1(4) 4(3) 2(4) 5(4) 3(5) 7(3) \\ leave 2: 6(2) 7(3) 4(3) 1(4) 5(4) 3(5) \\ upgrade 3(1): 3(1) 7(3) 6(2) 1(4) 5(4) 4(3) \\ arrive 8(3): 3(1) 7(3) 6(2) 1(4) 5(4) 4(3) 8(3) \\ treat 3(1): 6(2) 7(3) 4(3) 1(4) 5(4) 8(3) \\ treat 6(2): 4(3) 7(3) 8(3) 1(4) 5(4)}

\subsection{手算追蹤}

上面的輸出整理成考卷的寫法，一個操作一列，記陣列的內容（位置從 $1$ 開始）。

\begin{center}
\begin{tabular}{l*{7}{c}}
\toprule
操作 & 1 & 2 & 3 & 4 & 5 & 6 & 7 \\
\midrule
到達 1(4)、2(4)、3(5) & 1(4) & 2(4) & 3(5) \\
到達 4(3) & 4(3) & 1(4) & 3(5) & 2(4) \\
到達 5(4) & 4(3) & 1(4) & 3(5) & 2(4) & 5(4) \\
到達 6(2) & 6(2) & 1(4) & 4(3) & 2(4) & 5(4) & 3(5) \\
到達 7(3) & 6(2) & 1(4) & 4(3) & 2(4) & 5(4) & 3(5) & 7(3) \\
離開 2 & 6(2) & 7(3) & 4(3) & 1(4) & 5(4) & 3(5) \\
3 號改成 1 級 & 3(1) & 7(3) & 6(2) & 1(4) & 5(4) & 4(3) \\
到達 8(3) & 3(1) & 7(3) & 6(2) & 1(4) & 5(4) & 4(3) & 8(3) \\
看診 3 & 6(2) & 7(3) & 4(3) & 1(4) & 5(4) & 8(3) \\
看診 6 & 4(3) & 7(3) & 8(3) & 1(4) & 5(4) \\
\bottomrule
\end{tabular}
\end{center}

幾個值得看的步驟：

\begin{itemize}
\item \textbf{到達 6(2)}：放在第 6 格，父節點第 3 格是 3(5)，交換；到了第 3 格，父節點
4(3) 也比較後面，再交換，6 號一路到根。
\item \textbf{離開 2}：2 號在第 4 格。 最後一格的 7(3) 補進第 4 格，它比父節點 1(4) 優
先，所以要\textbf{上浮}，和 1(4) 交換。 只寫下沉的程式在這裡會出錯。
\item \textbf{3 號改成 1 級}：3 號在第 6 格，鍵變小以後上浮兩次到根，這就是
\proc{Decrease-Key}。 程式用 \texttt{pos[3]} 直接找到第 6 格，不用搜尋。
\item \textbf{看診 3}：最後一格的 8(3) 搬到根，兩個子節點 7(3)、6(2) 中 6(2) 優先，交換；
到了第 3 格，唯一的子節點 4(3) 和它同級但編號較小，所以再交換一次。
\end{itemize}

\subsection{其他常見用法}

\begin{itemize}
\item \textbf{合併 $k$ 個排好的數列}：最小堆積裡放每個數列目前的第一個數，每次取出最小
的輸出，再把它所屬數列的下一個數放進去。 堆積最多 $k$ 個元素，$N$ 個數各進出一次，
共 $O(N \lg k)$。 第~15~章的外部排序會用到同樣的想法。
\item \textbf{找前 $k$ 大}：用一個最多放 $k$ 個元素的最小堆積，掃過每個數，堆積不滿就放
進去；滿了而新的數比堆積的最小值大，就把最小值換成新的數。 最後堆積的根就是第 $k$
大，總時間 $O(n \lg k)$，只要 $O(k)$ 的額外空間。
\item \textbf{取出後立刻放入}：取出最小值後馬上要放一個新的數時，不必先取出（下沉一次）
再插入（上浮一次）。 直接把根改成新的數，做一次下沉就好。 這叫 replace 操作，上面
「找前 $k$ 大」的換掉最小值就是它。
\end{itemize}

\section{$d$ 元堆積}
\label{sec:hp-dary}

把「每個節點最多兩個子節點」改成最多 $d$ 個，就是 \textbf{$d$ 元堆積}（$d$-ary heap）。
仍然用陣列存，一層一層由左到右編號。 從 $1$ 開始編號時，位置 $i$ 的子節點和父節點是
\[
\text{子節點：} d(i - 1) + 2,\; d(i - 1) + 3,\; \ldots,\; di + 1, \qquad
\text{父節點：} \left\lfloor \frac{i - 2}{d} \right\rfloor + 1.
\]
從 $0$ 開始編號（C 程式）時比較整齊：子節點是 $di + 1$ 到 $di + d$，父節點是
$\lfloor (i - 1)/d \rfloor$。 $d = 2$ 時兩組公式都會回到上一節的樣子。 最後一個有子節點
的位置是 $n$ 的父節點，從 $1$ 編號時是 $\lfloor (n - 2)/d \rfloor + 1$。

高度 $h \ge 1$ 的完整 $d$ 元樹，前 $h - 1$ 層有
$1 + d + \cdots + d^{h-2} = (d^{h-1} - 1)/(d - 1)$ 個節點，第 $h$ 層至少一個、最多
$d^{h-1}$ 個，所以
\[
\frac{d^{h-1} - 1}{d - 1} + 1 \le n \le \frac{d^h - 1}{d - 1}.
\]
右邊的不等式是 $d^h \ge (d - 1)n + 1$；左邊乘開是 $(d - 1)n + 1 \ge d^{h-1} + d - 1 >
d^{h-1}$。 兩者合起來，$h - 1 < \log_d\bigl((d - 1)n + 1\bigr) \le h$，所以
\[
h = \left\lceil \log_d \bigl((d - 1)n + 1\bigr) \right\rceil = \Theta(\log_d n).
\]
$d = 2$ 時就是定理~\ref{thm:bt-count} 的 $\lceil \lg(n + 1) \rceil$，和完整二元樹的
$\lfloor \lg n \rfloor + 1$ 相等。

操作的成本跟著改變：

\begin{itemize}
\item \textbf{上浮}（插入、decrease-key）：每層比較一次，$O(\log_d n)$。 $d$ 越大，樹
越矮，越快。
\item \textbf{下沉}（取出最小值）：每層要在 $d$ 個子節點裡找最小的，再和自己比，$d$
次比較，共 $O(d \log_d n)$。 $d$ 越大，這一項越慢。
\end{itemize}

所以 $d$ 的選擇看兩類操作的比例。 Dijkstra 演算法在 $n$ 個點、$m$ 條邊的圖上，大約做
$n$ 次取出最小值和 $m$ 次 decrease-key，總共 $O(n d \log_d n + m \log_d n)$。 取
$d \approx \max(2, m/n)$，兩項平衡，邊很多的圖會比二元堆積快（第~25~章）。 只看取出最
小值時，$d \log_d n = (d / \ln d) \ln n$，而 $d / \ln d$ 在整數裡以 $d = 3$ 最小，所以三
元堆積的比較次數比二元堆積少一點。

\subsection{三元堆積的手算追蹤}

用 $\langle 44, 17, 62, 28, 9, 51, 33, 70, 12, 38 \rangle$ 建三元最小堆積，再取出最小值、
插入 $5$。 $n = 10$，最後一個有子節點的位置是 $\lfloor 8/3 \rfloor + 1 = 3$，$i$ 從 $3$
做到 $1$。 位置 $1$ 的子節點是 $2, 3, 4$；位置 $2$ 的是 $5, 6, 7$；位置 $3$ 的是
$8, 9, 10$。

\begin{center}
\begin{tabular}{lc*{10}{c}}
\toprule
& 交換 & 1 & 2 & 3 & 4 & 5 & 6 & 7 & 8 & 9 & 10 \\
\midrule
開始 & & 44 & 17 & 62 & 28 & 9 & 51 & 33 & 70 & 12 & 38 \\
$i = 3$ & 1 & 44 & 17 & 12 & 28 & 9 & 51 & 33 & 70 & 62 & 38 \\
$i = 2$ & 1 & 44 & 9 & 12 & 28 & 17 & 51 & 33 & 70 & 62 & 38 \\
$i = 1$ & 2 & 9 & 17 & 12 & 28 & 44 & 51 & 33 & 70 & 62 & 38 \\
\midrule
取出 9 & 1 & 12 & 17 & 38 & 28 & 44 & 51 & 33 & 70 & 62 \\
插入 5 & 2 & 5 & 17 & 12 & 28 & 44 & 51 & 33 & 70 & 62 & 38 \\
\bottomrule
\end{tabular}
\end{center}

$i = 1$ 時 $44$ 先和 $\{9, 12, 28\}$ 中的 $9$ 交換，到了位置 $2$，再和 $\{17, 51, 33\}$
中的 $17$ 交換，到了沒有子節點的位置 $5$。 取出 $9$ 時，最後一個 $38$ 搬到根，和
$\{17, 12, 28\}$ 中的 $12$ 交換，到了位置 $3$，子節點只剩 $70$、$62$，都比 $38$ 大，停。
插入 $5$ 放在位置 $10$，父節點是 $\lfloor 8/3 \rfloor + 1 = 3$，交換，再和根交換。

\begin{center}
\begin{tikzpicture}[level distance=1.0cm,
    level 1/.style={sibling distance=3.0cm},
    level 2/.style={sibling distance=0.85cm},
    every node/.style={hpn}, edge from parent/.style={hpedge, draw}]
\node {9}
    child {node {17} child {node {44}} child {node {51}} child {node {33}}}
    child {node {12} child {node {70}} child {node {62}} child {node {38}}}
    child {node {28}};
\end{tikzpicture}
\diagcaption{建好的三元最小堆積。 每個節點最多三個子節點，$10$ 個元素的高度是 $3$，二
元堆積則是 $4$。}
\label{fig:hp-ternary}
\end{center}

實際的執行時間還和快取有關：$d$ 個子節點在陣列裡是相鄰的，一次讀進快取，樹又比較矮，所
以實務上常用 $d = 4$ 左右。

\section{雙端優先佇列、左傾樹與二項堆積}
\label{sec:hp-mergeable}

二元堆積有兩個做不到的地方。 第一，最小堆積找最大值要看所有的葉節點，$\Theta(n)$；想同時
取最小和最大，就要另外想辦法。 第二，兩個堆積合併成一個，只能把兩個陣列接起來重建，
$\Theta(n)$。 這一節的前三種結構（最小最大堆積、deap、對稱最小最大堆積）處理第一件，後兩
種（左傾樹、二項堆積）處理第二件。 每一種都以手算追蹤為主，考卷上最常見的題目就是「插入
（刪除）之後畫出結果」。

\subsection{最小最大堆積}

\textbf{雙端優先佇列}（double-ended priority queue）同時支援取出最小值和取出最大值。 例如
急診室同時想知道最緊急和最不緊急的病人，後者可以轉到一般門診。 \textbf{最小最大堆積}
（min-max heap）用一個陣列做到這件事。

\begin{definition}[最小最大堆積]
\label{def:hp-minmax}
最小最大堆積是存在陣列裡的完整二元樹，各層輪流是\textbf{最小層}和\textbf{最大層}：根所在
的深度 $0$ 是最小層，深度 $1$ 是最大層，深度 $2$ 又是最小層，依此類推。 最小層的節點
不大於它的所有子孫，最大層的節點不小於它的所有子孫。
\end{definition}

所以最小值在根，最大值是根的兩個子節點中較大的那個（只有一個元素時就是根）。 位置 $i$
的深度是 $\lfloor \lg i \rfloor$，偶數就在最小層。 和二元堆積不同的是，比較的對象常常
是\textbf{祖父}（位置 $\lfloor i/4 \rfloor$）和\textbf{孫子}，因為隔一層才是同一種層。

\textbf{插入}：新元素 $x$ 放在位置 $j = n + 1$，先和父節點比一次，決定它屬於「小的那一
邊」還是「大的那一邊」，之後只沿著同一種層的祖父往上走。

\begin{pseudo}[label=alg:hp-minmax-insert]{最小最大堆積的插入}
\Procedure{MinMax-Insert}{$A, x$}
    \State $A.\mathit{size} \gets A.\mathit{size} + 1$
    \State $j \gets A.\mathit{size}$
    \If{$j = 1$}
        \State $A[1] \gets x$
    \ElsIf{$j$ 在最小層}
        \If{$x > A[\lfloor j/2 \rfloor]$}
            \State $A[j] \gets A[\lfloor j/2 \rfloor]$ \Comment{父節點在最大層}
            \State \Call{Push-Up-Max}{$A, \lfloor j/2 \rfloor, x$}
        \Else
            \State \Call{Push-Up-Min}{$A, j, x$}
        \EndIf
    \Else \Comment{$j$ 在最大層}
        \If{$x < A[\lfloor j/2 \rfloor]$}
            \State $A[j] \gets A[\lfloor j/2 \rfloor]$ \Comment{父節點在最小層}
            \State \Call{Push-Up-Min}{$A, \lfloor j/2 \rfloor, x$}
        \Else
            \State \Call{Push-Up-Max}{$A, j, x$}
        \EndIf
    \EndIf
\EndProcedure
\Statex
\Procedure{Push-Up-Min}{$A, i, x$} \Comment{$i$ 在最小層，$x$ 要放在 $i$ 或它的祖先}
    \While{$i \ge 4$ \AND $x < A[\lfloor i/4 \rfloor]$}
        \State $A[i] \gets A[\lfloor i/4 \rfloor]$
        \State $i \gets \lfloor i/4 \rfloor$
    \EndWhile
    \State $A[i] \gets x$
\EndProcedure
\end{pseudo}

\proc{Push-Up-Max} 把第 19 行的 $<$ 改成 $>$。 這裡用「挖洞」的寫法：祖父往下挪，$x$ 最
後才放進去。 每一步往上兩層，所以插入只要 $O(\lg n)$，比較次數大約是 $\frac{1}{2}\lg n
+ 1$。

\textbf{取出最小值}：最小值在根。 把最後一個元素 $x$ 拿出來，堆積少一格，然後替 $x$ 在
根底下找位置。 根在最小層，要放的是「$x$ 和根的所有子孫中最小的」，而那個最小的一定在
根的子節點或孫節點裡（再往下的都被它們的祖先管住了）。

\begin{pseudo}[label=alg:hp-minmax-delete]{最小最大堆積的取出最小值}
\Function{MinMax-Extract-Min}{$A$} \Comment{假設 $A.\mathit{size} \ge 1$}
    \State $\mathit{min} \gets A[1]$
    \State $x \gets A[A.\mathit{size}]$
    \State $A.\mathit{size} \gets A.\mathit{size} - 1$
    \State $i \gets 1$
    \While{$2i \le A.\mathit{size}$}
        \State $k \gets$ $i$ 的子節點和孫節點之中，值最小的位置
        \If{$x \le A[k]$}
            \State \Break
        \EndIf
        \State $A[i] \gets A[k]$
        \If{$k$ 是 $i$ 的子節點}
            \State $i \gets k$
            \State \Break
        \EndIf
        \If{$x > A[\lfloor k/2 \rfloor]$} \Comment{$k$ 的父節點在最大層}
            \State 交換 $x$ 和 $A[\lfloor k/2 \rfloor]$
        \EndIf
        \State $i \gets k$
    \EndWhile
    \If{$A.\mathit{size} \ge 1$}
        \State $A[i] \gets x$
    \EndIf
    \State \Return $\mathit{min}$
\EndFunction
\end{pseudo}

每一輪往下兩層，看最多六個位置，所以是 $O(\lg n)$。 第 11 行處理 $k$ 是子節點的情況。
$k$ 在最大層，它的子孫都不大於 $A[k]$；而 $A[k]$ 又是所有候選（包括 $k$ 的子節點）裡最小
的，所以 $k$ 如果有子孫，值都等於 $A[k]$。 這時 $x > A[k]$ 不小於 $k$ 的任何子孫，放在最
大層的 $k$ 是合法的，搬到 $i$ 的 $A[k]$ 也不大於 $i$ 的任何子孫，可以結束。 第 15 行的交
換是因為 $x$ 要往孫子那裡走，路上會經過最大層的 $A[\lfloor k/2 \rfloor]$，$x$ 比它大的
話就不能待在它底下，兩者對調。 取出最大值的做法對稱：最大值在位置 $2$ 或 $3$，從那裡開始，把「最小」都換成「最大」。

\textbf{手算追蹤}。 下圖 (a) 是一個有 $10$ 個元素的最小最大堆積，陣列是
\[
\langle 11, 46, 72, 31, 23, 58, 38, 36, 37, 34 \rangle.
\]
依序做三個操作：插入 $8$、插入 $77$、取出最小值。

\begin{center}
\begin{tikzpicture}
\foreach \y/\t in {0/min, -0.95/max, -1.9/min, -2.85/max}
    \node[hpidx, anchor=east] at (-3.3,\y) {\t};
\hpheap{0,0}{6.0}{11/n,46/n,72/n,31/n,23/n,58/n,38/n,36/n,37/n,34/n}{m}
\hpheap{7.2,0}{6.0}{8/a,46/n,72/n,31/n,11/a,58/n,38/n,36/n,37/n,34/n,23/a}{n}
\node[font=\small] at (0,-3.65) {(a) 開始};
\node[font=\small] at (7.2,-3.65) {(b) 插入 $8$};
\end{tikzpicture}

\begin{tikzpicture}
\foreach \y/\t in {0/min, -0.95/max, -1.9/min, -2.85/max}
    \node[hpidx, anchor=east] at (-3.3,\y) {\t};
\hpheap{0,0}{6.0}{8/n,46/n,77/a,31/n,11/n,58/n,38/n,36/n,37/n,34/n,23/n,72/a}{o}
\hpheap{7.2,0}{6.0}{11/a,72/a,77/n,31/n,23/a,58/n,38/n,36/n,37/n,34/n,46/a}{p}
\node[font=\small] at (0,-3.65) {(c) 插入 $77$};
\node[font=\small] at (7.2,-3.65) {(d) 取出最小值 $8$};
\end{tikzpicture}
\diagcaption{最小最大堆積的三個操作。 左邊標出每一層是最小層還是最大層，橘色是這一步改
變的位置。}
\label{fig:hp-minmax}
\end{center}

\begin{itemize}
\item \textbf{插入 $8$}：放在位置 $11$，深度 $3$，是最大層。 父節點 $A[5] = 23$ 在最小
層，$8 < 23$，所以 $23$ 挪到位置 $11$，$8$ 改走最小層：從位置 $5$ 看祖父 $A[1] = 11$，
$8 < 11$，$11$ 挪到位置 $5$，$8$ 到了根。
\item \textbf{插入 $77$}：放在位置 $12$，最大層。 父節點 $A[6] = 58$，$77 > 58$，不用和
父節點換，直接走最大層：祖父 $A[3] = 72$，$77 > 72$，$72$ 挪到位置 $12$，$77$ 放在位
置 $3$。
\item \textbf{取出最小值}：取出 $8$，最後一個元素 $x = 72$ 拿出來。 根的子孫中，子節點
$46, 77$ 和孫節點 $31, 11, 58, 38$ 裡最小的是 $A[5] = 11$（孫節點），搬到根。 $11$ 的父
節點 $A[2] = 46$ 在最大層，$x = 72 > 46$，對調，$46$ 繼續往下找位置，$72$ 留在
$A[2]$。 到了位置 $5$，它的子節點 $34, 23$ 裡最小的是 $A[11] = 23$（子節點），比 $46$
小，$23$ 搬到位置 $5$，$46$ 放進位置 $11$，結束。
\end{itemize}

最後的陣列是 $\langle 11, 72, 77, 31, 23, 58, 38, 36, 37, 34, 46 \rangle$。 可以逐層檢查：
最大層的 $72$ 不小於它底下的 $31, 23, 36, 37, 34, 46$，最小層的 $23$ 不大於 $34, 46$。

\subsection{雙端堆積（deap）}
\label{sec:hp-deap}

最小最大堆積把小的和大的放在交錯的層裡。 \textbf{deap}（double-ended heap，中文常叫雙端
堆積）改成左右分開：根不放元素，根的左子樹是一個最小堆積，右子樹是一個最大堆積，兩邊再用
「對應點」互相約束，保證左邊每個元素都不大於右邊對應位置的元素。

仍然用陣列存。 $A[1]$ 是空的根，元素放在 $A[2..\mathit{last}]$，共 $\mathit{last} - 1$ 個。 深度
$d \ge 1$ 的位置是 $2^d$ 到 $2^{d+1} - 1$，前一半屬於最小堆積，後一半屬於最大堆積：位置 $i$
在最大堆積，若且唯若 $i \ge 2^d + 2^{d-1}$，其中 $d = \lfloor \lg i \rfloor$。 兩個堆積的根分別是
$A[2]$ 和 $A[3]$；除了這兩個根，每個位置的父節點仍然是 $\lfloor i/2 \rfloor$，而且和自己在同
一邊。

\begin{definition}[deap]
\label{def:hp-deap}
deap 是存在陣列 $A[1..\mathit{last}]$ 的完整二元樹，根 $A[1]$ 不放元素，而且滿足：
\begin{enumerate}
\item 根的左子樹是最小堆積，右子樹是最大堆積（只有一個元素時右子樹是空的）。
\item 對最小堆積的每個位置 $i$，令 $j = i + 2^{\lfloor \lg i \rfloor - 1}$，這是最大堆積裡同一層、
同一個相對位置的\textbf{對應點}；如果 $j > \mathit{last}$，改用它的父節點 $\lfloor j/2 \rfloor$。
那麼 $A[i] \le A[j]$。
\end{enumerate}
\end{definition}

改用父節點是因為最後一層從左邊填起，最小堆積那一邊的最後一層可能比右邊多出幾個節點，它們
在右邊還沒有對應點，就和那個位置的父節點比。 反過來，最大堆積的位置 $j$ 的對應點是
$i = j - 2^{\lfloor \lg j \rfloor - 1}$，這個 $i$ 一定存在。

最小值在 $A[2]$，最大值在 $A[3]$（只有一個元素時兩者都是 $A[2]$）。 $A[2]$ 是最小堆積的最
小值；最大堆積的任何位置 $j$，對應點 $i$ 都存在而且 $A[i] \le A[j]$，所以 $A[j]$ 也不小於
$A[2]$。 $A[3]$ 是最大值的理由對稱。

\textbf{插入}。 新元素 $x$ 先放在新的最後一格 $p$。 如果 $p$ 在最小堆積，看它的對應點 $j$：
$x > A[j]$ 表示 $x$ 屬於大的那一邊，把 $A[j]$ 搬到 $p$，$x$ 從 $j$ 開始在最大堆積裡上浮；否則
$x$ 從 $p$ 開始在最小堆積裡上浮。 $p$ 在最大堆積時反過來，和對應點比，$x$ 比較小就換到最小
堆積那一邊。 上浮用最小最大堆積那種「挖洞」的寫法：\proc{Min-Up}$(A, i, x)$ 從位置 $i$ 沿著父
節點往上，父節點比 $x$ 大就把父節點往下挪，最多走到最小堆積的根 $A[2]$；\proc{Max-Up} 對稱，
最多走到 $A[3]$。 新的最後一格 $p$ 在最小堆積時，它的對應點一定超出範圍，改用的
$\lfloor j/2 \rfloor$ 正好是 $p$ 的父節點的對應點，所以搬到 $p$ 的 $A[j]$ 不小於 $p$ 的父節點，
最小堆積的性質不會被破壞。

\textbf{取出最小值}。 最小值在 $A[2]$。 把最後一格的元素 $x$ 拿出來，$\mathit{last}$ 減 $1$。
接著讓 $A[2]$ 這個洞沿著最小堆積往下走：每次把較小的子節點搬上來，直到洞到了沒有子節點的
位置 $i$。 最後把 $x$ 放進位置 $i$，做法和插入時一樣，先和 $i$ 的對應點比。 和二元堆積的
下沉不同，洞不在中途停下來，一路走到最底：這樣 $x$ 只需要在最後決定去哪一邊，再往上浮一次。

\textbf{取出最大值}。 對稱：最大值在 $A[3]$（只有一個元素時就是 $A[2]$），洞沿著最大堆積較
大的子節點走到底，到了位置 $j$ 再放 $x$。 這時有一個細節：如果對應點 $i$ 有子節點，這些子節
點在最大堆積那邊沒有對應點，依定義都要和 $j$ 比，所以 $x$ 要和「$i$ 的子節點中較大的那個」比，
不能只和 $A[i]$ 比。

\begin{pseudo}[label=alg:hp-deap]{deap 的插入與取出最小值}
\Procedure{Deap-Insert}{$A, x$}
    \State $A.\mathit{last} \gets A.\mathit{last} + 1$
    \If{$A.\mathit{last}$ 在最大堆積}
        \State \Call{Deap-Put-Max}{$A, A.\mathit{last}, x$}
    \Else
        \State \Call{Deap-Put-Min}{$A, A.\mathit{last}, x$}
    \EndIf
\EndProcedure
\Statex
\Procedure{Deap-Put-Min}{$A, i, x$} \Comment{位置 $i$ 在最小堆積，目前是空的}
    \State $j \gets i + 2^{\lfloor \lg i \rfloor - 1}$ \Comment{對應點}
    \If{$j > A.\mathit{last}$}
        \State $j \gets \lfloor j/2 \rfloor$ \Comment{最大堆積是空的時 $j = 1$}
    \EndIf
    \If{$j \ge 3$ \AND $x > A[j]$}
        \State $A[i] \gets A[j]$
        \State \Call{Max-Up}{$A, j, x$}
    \Else
        \State \Call{Min-Up}{$A, i, x$}
    \EndIf
\EndProcedure
\Statex
\Procedure{Deap-Put-Max}{$A, j, x$} \Comment{位置 $j$ 在最大堆積，目前是空的}
    \State $i \gets j - 2^{\lfloor \lg j \rfloor - 1}$ \Comment{對應點}
    \If{$2i \le A.\mathit{last}$} \Comment{$i$ 有子節點：改看較大的子節點}
        \State $i \gets 2i$
        \If{$i + 1 \le A.\mathit{last}$ \AND $A[i + 1] > A[i]$}
            \State $i \gets i + 1$
        \EndIf
    \EndIf
    \If{$x < A[i]$}
        \State $A[j] \gets A[i]$
        \State \Call{Min-Up}{$A, i, x$}
    \Else
        \State \Call{Max-Up}{$A, j, x$}
    \EndIf
\EndProcedure
\Statex
\Function{Deap-Delete-Min}{$A$} \Comment{假設 $A.\mathit{last} \ge 2$}
    \State $\mathit{min} \gets A[2]$
    \State $x \gets A[A.\mathit{last}]$
    \State $A.\mathit{last} \gets A.\mathit{last} - 1$
    \If{$A.\mathit{last} \ge 2$}
        \State $i \gets 2$
        \While{$2i \le A.\mathit{last}$} \Comment{洞沿著較小的子節點走到底}
            \State $c \gets 2i$
            \If{$c + 1 \le A.\mathit{last}$ \AND $A[c + 1] < A[c]$}
                \State $c \gets c + 1$
            \EndIf
            \State $A[i] \gets A[c]$
            \State $i \gets c$
        \EndWhile
        \State \Call{Deap-Put-Min}{$A, i, x$}
    \EndIf
    \State \Return $\mathit{min}$
\EndFunction
\end{pseudo}

\proc{Deap-Delete-Max} 把 \proc{Deap-Delete-Min} 的 $A[2]$ 換成 $A[3]$（只有一個元素時直接
取 $A[2]$）、「較小的子節點」換成「較大的子節點」，最後呼叫 \proc{Deap-Put-Max}。 樹高是
$\lfloor \lg \mathit{last} \rfloor + 1$，洞往下走和上浮都只沿一條路徑，所以三個操作都是
$O(\lg n)$。

\textbf{手算追蹤}。 把 $35, 69, 29, 33, 65, 52, 96, 61, 24, 14, 55$ 依序插入空的 deap，再取出最
小值、取出最大值。 下表每一列是操作之後的 $A[2..12]$，粗體的位置屬於最大堆積。

\begin{center}
\small
\begin{tabular}{l*{11}{c}}
\toprule
操作 & 2 & \textbf{3} & 4 & 5 & \textbf{6} & \textbf{7} & 8 & 9 & 10 & 11 & \textbf{12} \\
\midrule
插入 35 & 35 \\
插入 69 & 35 & 69 \\
插入 29 & 29 & 69 & 35 \\
插入 33 & 29 & 69 & 35 & 33 \\
插入 65 & 29 & 69 & 35 & 33 & 65 \\
插入 52 & 29 & 69 & 35 & 33 & 65 & 52 \\
插入 96 & 29 & 96 & 35 & 33 & 69 & 52 & 65 \\
插入 61 & 29 & 96 & 35 & 33 & 69 & 52 & 65 & 61 \\
插入 24 & 24 & 96 & 35 & 29 & 69 & 52 & 65 & 61 & 33 \\
插入 14 & 14 & 96 & 35 & 24 & 69 & 52 & 65 & 61 & 33 & 29 \\
插入 55 & 14 & 96 & 35 & 24 & 69 & 52 & 55 & 61 & 33 & 29 & 65 \\
\midrule
取出最小 14 & 24 & 96 & 35 & 29 & 69 & 65 & 55 & 61 & 33 & 52 \\
取出最大 96 & 24 & 69 & 35 & 29 & 61 & 65 & 55 & 52 & 33 \\
\bottomrule
\end{tabular}
\end{center}

\begin{itemize}
\item \textbf{插入 $29$}：放在位置 $4$（最小堆積）。 對應點 $4 + 2 = 6$ 超出範圍，改用父節點
$3$；$29 \le A[3] = 69$，留在最小堆積上浮，$A[2] = 35$ 比它大，$35$ 往下挪，$29$ 到位置 $2$。
\item \textbf{插入 $96$}：放在位置 $8$（深度 $3$ 的 $8$ 到 $11$ 屬於最小堆積）。 對應點 $12$ 超出
範圍，改用 $6$；$96 > A[6] = 65$，所以 $65$ 搬到位置 $8$，$96$ 從位置 $6$ 在最大堆積上浮，父節
點 $A[3] = 69$ 比它小，$69$ 往下挪，$96$ 到位置 $3$。
\item \textbf{插入 $55$}：放在位置 $12$（最大堆積）。 對應點 $12 - 4 = 8$，它沒有子節點；
$55 < A[8] = 65$，所以 $65$ 搬到位置 $12$，$55$ 從位置 $8$ 在最小堆積上浮，父節點 $A[4] = 35$
比它小，停。
\item \textbf{取出最小值 $14$}：最後一格的 $65$ 拿出來。 洞從位置 $2$ 往下：子節點 $35, 24$ 中
較小的 $24$ 上來，洞到位置 $5$；子節點 $33, 29$ 中較小的 $29$ 上來，洞到位置 $11$，沒有子
節點了。 位置 $11$ 的對應點 $15$ 超出範圍，改用 $7$；$65 > A[7] = 52$，所以 $52$ 搬到位置
$11$，$65$ 從位置 $7$ 上浮，父節點 $96$ 比它大，停在 $7$。
\item \textbf{取出最大值 $96$}：最後一格的 $52$ 拿出來。 洞從位置 $3$ 往下：子節點 $69, 65$ 中
較大的 $69$ 上來，洞到位置 $6$，沒有子節點了。 對應點 $6 - 2 = 4$ 有子節點 $55, 61$，取較大的
$61$（位置 $9$）；$52 < 61$，所以 $61$ 搬到位置 $6$，$52$ 從位置 $9$ 在最小堆積上浮，父節點
$35$ 比它小，停。
\end{itemize}

最後一步如果只和 $A[4] = 35$ 比，$52$ 會留在位置 $6$；位置 $9$ 的 $61$ 對應點 $13$ 超出範圍，改
用 $6$，就違反了 $61 \le 52$。 這就是 \proc{Deap-Put-Max} 要看子節點的原因。

\begin{center}
\begin{tikzpicture}
\hpshowidxtrue
\hpheap{0,0}{6.0}{/e,14/n,96/n,35/n,24/n,69/n,52/n,55/n,61/n,33/n,29/n,65/n}{da}
\hpshowidxfalse
\node[hpidx] at (-1.5,-3.35) {最小堆積};
\node[hpidx] at (1.5,-3.35) {最大堆積};
\hpheap{7.2,0}{6.0}{/e,24/a,96/n,35/n,29/a,69/n,65/a,55/n,61/n,33/n,52/a}{db}
\node[font=\small] at (0,-3.9) {(a) 插入 $11$ 個之後};
\node[font=\small] at (7.2,-3.9) {(b) 取出最小值 $14$};
\end{tikzpicture}

\begin{tikzpicture}
\hpheap{0,0}{6.0}{/e,24/n,69/a,35/n,29/n,61/a,65/n,55/n,52/a,33/n}{dc}
\node[font=\small] at (0,-3.9) {(c) 取出最大值 $96$};
\end{tikzpicture}
\diagcaption{deap 的三個狀態。 虛線圓是不放元素的根，左子樹是最小堆積，右子樹是最大堆積；
(a) 右上角的小字是陣列位置，橘色是上一個狀態之後改變的位置。}
\label{fig:hp-deap}
\end{center}

\begin{examnote}
deap 的題目先確認三件事。 一、根是空的，元素從 $A[2]$ 開始；有的書把 $A[2]$ 叫做「位置 1」，
對應點的公式要跟著平移，照題目給的圖和編號。 二、對應點超出範圍時改用父節點。 三、取出最
大值時，要和「對應點的較大子節點」比（對應點沒有子節點才和對應點本身比）。 例如依序插入
$10, 30, 40, 20$ 得到 $A[2..5] = \langle 10, 40, 30, 20 \rangle$，取出最大值 $40$ 時，$20$ 要和
$A[4] = 30$ 比，結果是 $\langle 10, 30, 20 \rangle$；只和 $A[2] = 10$ 比會得到
$\langle 10, 20, 30 \rangle$，位置 $4$ 的 $30$ 比它的對應點 $A[3] = 20$ 大，不是 deap。
層數照全書約定從根算第 $1$ 層，所以 $A[2]$、$A[3]$ 在第 $2$ 層。
\end{examnote}

\subsection{對稱最小最大堆積}
\label{sec:hp-smmh}

\textbf{對稱最小最大堆積}（symmetric min-max heap，簡稱 SMMH）也是根不放元素的完整二元樹，
但規則只看兄弟和祖父，不用算對應點，所以程式比 deap 短。

\begin{definition}[對稱最小最大堆積]
\label{def:hp-smmh}
存在陣列 $A[1..\mathit{last}]$ 的完整二元樹，根 $A[1]$ 不放元素，其他每個節點放一個元素，而且：
\begin{enumerate}
\item[P1.] 每個有右兄弟的節點，元素不大於右兄弟：偶數位置 $x$ 有 $A[x] \le A[x + 1]$。
\item[P2.] 每個有祖父 $g$ 的節點 $x$，$A[2g] \le A[x]$：祖父的左子節點不大於它。
\item[P3.] 每個有祖父 $g$ 的節點 $x$，$A[2g + 1] \ge A[x]$：祖父的右子節點不小於它。
\end{enumerate}
\end{definition}

換個方式讀：任何一個節點 $y$，把 $y$ 的子樹扣掉 $y$ 自己，剩下的元素裡最小的在 $y$ 的左子節
點，最大的在 $y$ 的右子節點。 （P1 管左右子節點本身，P2、P3 管下面兩層，再往下一層一層套
用就得到整棵子樹。） 拿根當 $y$，最小值在 $A[2]$，最大值在 $A[3]$。

\textbf{插入}。 新元素 $x$ 放在新的最後一格 $c$。 如果 $c$ 是右子節點而 $x$ 比左兄弟小，會違反
P1，所以先把左兄弟搬到 $c$，$x$ 改看左兄弟的位置。 接著只要 $c$ 有祖父 $g$（$c \ge 4$），就和
$A[2g]$、$A[2g + 1]$ 比：比 $A[2g]$ 小就把它搬下來、$x$ 往上到 $2g$；比 $A[2g + 1]$ 大就把它
搬下來、$x$ 往上到 $2g + 1$；介於兩者之間就停。 每一步往上一層。

\textbf{取出最小值}。 最小值在 $A[2]$。 把最後一格的元素 $y$ 拿出來，替 $y$ 從洞 $e = 2$ 往下
找位置。 洞 $e$ 一定是左子節點；扣掉它以後，$e$ 和右兄弟 $e + 1$ 底下（不含 $e + 1$ 本身）最
小的元素，依上面的讀法一定是 $A[2e]$ 或 $A[2e + 2]$，也就是 $e$ 和 $e + 1$ 的左子節點。 取兩者
中較小的 $k$：$y \le A[k]$ 就把 $y$ 放在 $e$；否則 $A[k]$ 上來，洞移到 $k$。 洞移到新位置之後，
$y$ 將來會放在這個洞或它底下，要先檢查 P1：$y$ 比新洞的右兄弟大，就和右兄弟的元素對調，
帶著比較小的那個繼續往下。

\begin{pseudo}[label=alg:hp-smmh]{對稱最小最大堆積的插入與取出最小值}
\Procedure{SMMH-Insert}{$A, x$}
    \State $A.\mathit{last} \gets A.\mathit{last} + 1$
    \State $c \gets A.\mathit{last}$
    \If{$c$ 是奇數 \AND $x < A[c - 1]$} \Comment{P1：比左兄弟小}
        \State $A[c] \gets A[c - 1]$
        \State $c \gets c - 1$
    \EndIf
    \While{$c \ge 4$}
        \State $g \gets \lfloor c/4 \rfloor$ \Comment{祖父}
        \If{$x < A[2g]$} \Comment{P2}
            \State $A[c] \gets A[2g]$
            \State $c \gets 2g$
        \ElsIf{$x > A[2g + 1]$} \Comment{P3}
            \State $A[c] \gets A[2g + 1]$
            \State $c \gets 2g + 1$
        \Else
            \State \Break
        \EndIf
    \EndWhile
    \State $A[c] \gets x$
\EndProcedure
\Statex
\Function{SMMH-Delete-Min}{$A$} \Comment{假設 $A.\mathit{last} \ge 2$}
    \State $\mathit{min} \gets A[2]$
    \State $y \gets A[A.\mathit{last}]$
    \State $A.\mathit{last} \gets A.\mathit{last} - 1$
    \If{$A.\mathit{last} \ge 2$}
        \State $e \gets 2$ \Comment{洞，一定是左子節點}
        \While{$2e \le A.\mathit{last}$}
            \State $k \gets 2e$
            \If{$2e + 2 \le A.\mathit{last}$ \AND $A[2e + 2] < A[2e]$}
                \State $k \gets 2e + 2$
            \EndIf
            \If{$y \le A[k]$}
                \State \Break
            \EndIf
            \State $A[e] \gets A[k]$
            \State $e \gets k$
            \If{$e + 1 \le A.\mathit{last}$ \AND $y > A[e + 1]$} \Comment{P1}
                \State 交換 $y$ 和 $A[e + 1]$
            \EndIf
        \EndWhile
        \State $A[e] \gets y$
    \EndIf
    \State \Return $\mathit{min}$
\EndFunction
\end{pseudo}

\proc{SMMH-Delete-Max} 對稱：最大值在 $A[3]$（只有一個元素時是 $A[2]$），洞 $e$ 一定是右子
節點，候選是 $e - 1$ 和 $e$ 的子節點中最大的（右子節點存在就是右子節點，否則是左子節點），
洞移下去之後檢查 $y$ 有沒有比新洞的左兄弟小。 插入每層比兩次，取出每層最多比三次（兩個候
選、$y$、兄弟），都是 $O(\lg n)$。

\textbf{手算追蹤}。 用和 deap 相同的輸入 $35, 69, 29, 33, 65, 52, 96, 61, 24, 14, 55$，再取出最
小值和最大值。

\begin{center}
\small
\begin{tabular}{l*{11}{c}}
\toprule
操作 & 2 & 3 & 4 & 5 & 6 & 7 & 8 & 9 & 10 & 11 & 12 \\
\midrule
插入 35 & 35 \\
插入 69 & 35 & 69 \\
插入 29 & 29 & 69 & 35 \\
插入 33 & 29 & 69 & 33 & 35 \\
插入 65 & 29 & 69 & 33 & 35 & 65 \\
插入 52 & 29 & 69 & 33 & 35 & 52 & 65 \\
插入 96 & 29 & 96 & 33 & 69 & 52 & 65 & 35 \\
插入 61 & 29 & 96 & 33 & 69 & 52 & 65 & 35 & 61 \\
插入 24 & 24 & 96 & 29 & 69 & 52 & 65 & 35 & 61 & 33 \\
插入 14 & 14 & 96 & 24 & 69 & 52 & 65 & 35 & 61 & 29 & 33 \\
插入 55 & 14 & 96 & 24 & 69 & 52 & 65 & 35 & 61 & 29 & 33 & 55 \\
\midrule
取出最小 14 & 24 & 96 & 29 & 69 & 52 & 65 & 35 & 61 & 33 & 55 \\
取出最大 96 & 24 & 69 & 29 & 61 & 52 & 65 & 35 & 55 & 33 \\
\bottomrule
\end{tabular}
\end{center}

\begin{itemize}
\item \textbf{插入 $33$}：位置 $5$ 是右子節點，比左兄弟 $35$ 小，$35$ 搬到位置 $5$，$33$ 改看位
置 $4$。 祖父是根，它的兩個子節點是 $29$ 和 $69$，$29 \le 33 \le 69$，停在位置 $4$。
\item \textbf{插入 $96$}：位置 $8$ 的祖父是 $2$，看 $A[4] = 33$、$A[5] = 35$；$96 > 35$，$35$ 搬
下來，$96$ 到位置 $5$。 位置 $5$ 的祖父是根，看 $A[2] = 29$、$A[3] = 69$；$96 > 69$，$69$ 搬下
來，$96$ 到位置 $3$。
\item \textbf{插入 $14$}：位置 $11$ 是右子節點，比左兄弟 $33$ 小，先換到位置 $10$。 再往上兩
步：比 $A[4] = 29$ 小，$29$ 搬到位置 $10$；比 $A[2] = 24$ 小，$24$ 搬到位置 $4$，$14$ 到根的
左子節點。
\item \textbf{取出最小值 $14$}：最後一格的 $55$ 拿出來，洞在位置 $2$。 候選 $A[4] = 24$、
$A[6] = 52$，$24$ 較小而且比 $55$ 小，搬上來，洞到位置 $4$；右兄弟 $A[5] = 69$ 不比 $55$ 小，不
用對調。 候選 $A[8] = 35$、$A[10] = 29$，$29$ 上來，洞到位置 $10$；右兄弟 $A[11] = 33 < 55$，
對調，$55$ 放進位置 $11$，手上換成 $33$。 位置 $10$ 沒有子節點，$33$ 放進去。
\item \textbf{取出最大值 $96$}：最後一格的 $55$ 拿出來，洞在位置 $3$。 候選是位置 $2$ 的右子節
點 $A[5] = 69$ 和位置 $3$ 的右子節點 $A[7] = 35$，$69$ 上來，洞到位置 $5$；左兄弟 $A[4] = 29$
不比 $55$ 大，不用對調。 候選是 $A[9] = 61$ 和 $A[10] = 33$（位置 $5$ 只有左子節點），$61$
上來，洞到位置 $9$，沒有子節點，$55$ 放進去。
\end{itemize}

\begin{center}
\begin{tikzpicture}
\hpshowidxtrue
\hpheap{0,0}{6.0}{/e,14/n,96/n,24/n,69/n,52/n,65/n,35/n,61/n,29/n,33/n,55/n}{sa}
\hpshowidxfalse
\hpheap{7.2,0}{6.0}{/e,24/a,96/n,29/a,69/n,52/n,65/n,35/n,61/n,33/a,55/a}{sb}
\node[font=\small] at (0,-3.65) {(a) 插入 $11$ 個之後};
\node[font=\small] at (7.2,-3.65) {(b) 取出最小值 $14$};
\end{tikzpicture}

\begin{tikzpicture}
\hpheap{0,0}{6.0}{/e,24/n,69/a,29/n,61/a,52/n,65/n,35/n,55/a,33/n}{sc}
\node[font=\small] at (0,-3.65) {(c) 取出最大值 $96$};
\end{tikzpicture}
\diagcaption{對稱最小最大堆積的三個狀態，輸入和圖~\ref{fig:hp-deap} 相同。 每一對兄弟都是
左小右大，而且夾住它們底下的所有元素。}
\label{fig:hp-smmh}
\end{center}

同樣的輸入，deap 和 SMMH 得到的陣列不一樣，兩個都是對的；考卷問哪一種，就照哪一種的規則
做。

\begin{examnote}
SMMH 的根也是空的，「第 $k$ 層」照全書約定從根算第 $1$ 層，最小值、最大值在第 $2$ 層。 插入
時先處理兄弟（P1），再一層一層和「祖父的兩個子節點」比，不要和根的兩個子節點比：例如
$A[2..10] = \langle 5, 50, 10, 40, 20, 30, 15, 35, 45 \rangle$ 的 $45$ 雖然不大於 $A[3] = 50$，
卻大於它祖父的右子節點 $A[5] = 40$，違反 P3。 取出時洞每往下一層，都要檢查一次兄弟。
\end{examnote}

\subsection{幾種雙端優先佇列的比較}
\label{sec:hp-depq}

還有一種常見的結構叫\textbf{區間堆積}（interval heap）：每個節點放兩個元素 $a \le b$，看成區間
$[a, b]$（最後一個節點可以只放一個），每個節點的區間都包含子節點的區間；所以根的 $a$ 是最小
值，$b$ 是最大值。 把 SMMH 去掉根、每一對兄弟合成一個節點，剛好就是區間堆積：P1 說每一對
左小右大，P2、P3 說底下的元素都夾在這一對之間。 例如圖~\ref{fig:hp-smmh} (a) 兩兩配對成
$[14, 96]$；$[24, 69]$、$[52, 65]$；$[35, 61]$、$[29, 33]$、$[55]$（分號隔開各層），每個區間
都包含它子節點的區間。 它的插入和取出和 SMMH
幾乎一樣，只是樹高少一層，這裡不再展開。

\begin{center}
\small
\begin{tabular}{>{\raggedright\arraybackslash}p{0.16\linewidth}>{\raggedright\arraybackslash}p{0.16\linewidth}>{\raggedright\arraybackslash}p{0.28\linewidth}>{\raggedright\arraybackslash}p{0.29\linewidth}}
\toprule
結構 & 最小、最大 & 插入 & 取出最小值 \\
\midrule
最小最大堆積 & $A[1]$；$A[2]$ 或 $A[3]$ & 和父節點比一次決定哪一種層，再沿祖父往上 &
看子節點和孫節點（最多六個），每次下兩層 \\
deap & $A[2]$；$A[3]$ & 和對應點比一次決定哪一邊，再在那一邊上浮 &
洞沿最小堆積走到底，再和對應點比 \\
SMMH & $A[2]$；$A[3]$ & 先和兄弟比，再和祖父的兩個子節點比 &
洞每層看兩個候選，再和兄弟比 \\
區間堆積 & 根的兩端 & 和 SMMH 相同，一個節點放兩個元素 & 和 SMMH 相同 \\
\bottomrule
\end{tabular}
\end{center}

四種的插入、取出最小值、取出最大值都是 $O(\lg n)$，找最小值和最大值都是 $\Theta(1)$。 deap 和
SMMH 多用一格空的根；SMMH 的規則最少，是最好寫的一種。

\subsection{左傾樹}

左傾樹的目的是\textbf{合併}（meld）：兩個優先佇列併成一個，只花 $O(\lg n)$。 它不再是完
整二元樹，改用節點和指標存，形狀可以很不平衡，但\textbf{右邊}一定很短。

對每個節點 $x$，令 $s(x)$ 是從 $x$ 往下走到最近的空位（\NIL）要經過幾個節點：
$s(\NIL) = 0$，$s(x) = 1 + \min\bigl(s(x.\mathit{left}), s(x.\mathit{right})\bigr)$。 葉節點的
$s$ 是 $1$；只有一個子節點的節點，$s$ 也是 $1$。 實作時每個節點用一個欄位 $x.s$ 存
$s(x)$，$s(\NIL)$ 直接當成 $0$。

\begin{definition}[左傾樹]
\label{def:hp-leftist}
二元樹的每個節點 $x$ 都滿足 $s(x.\mathit{left}) \ge s(x.\mathit{right})$，叫\textbf{左傾樹}
（leftist tree）。 如果同時滿足最小堆積的性質（父節點不大於子節點），叫\textbf{最小左傾
樹}。
\end{definition}

左傾的條件讓 $s(x) = 1 + s(x.\mathit{right})$：一路往右走就是最短的那條路。

\begin{theorem}
\label{thm:hp-leftist}
左傾樹中，以 $x$ 為根的子樹至少有 $2^{s(x)} - 1$ 個節點，從 $x$ 一路往右走到 \NIL 剛好經
過 $s(x)$ 個節點。 因此 $n$ 個節點的左傾樹，最右邊的路徑最多有 $\lg(n + 1)$ 個節點。
\end{theorem}

\begin{proof}
由 $s$ 的定義，從 $x$ 往下走不到 $s(x)$ 步不會遇到 \NIL，所以 $x$ 底下的前 $s(x)$ 層（深
度 $0$ 到 $s(x) - 1$，相對於 $x$）都是滿的，至少 $1 + 2 + \cdots + 2^{s(x) - 1} = 2^{s(x)}
- 1$ 個節點。 第二句從 $s(x) = 1 + s(x.\mathit{right})$ 往下歸納即可。 把兩者合起來，
$n \ge 2^{s(\text{根})} - 1$，所以右邊的路徑長 $s(\text{根}) \le \lg(n + 1)$。
\end{proof}

\textbf{合併}兩棵最小左傾樹 $a$、$b$：根比較小的那棵當新的根，把它的右子樹和另一棵\textbf{遞
迴合併}，結果接回右邊；回來時如果右邊的 $s$ 比左邊大，左右對調，再更新 $s$。

\begin{pseudo}[label=alg:hp-meld]{合併兩棵最小左傾樹}
\Function{Meld}{$a, b$}
    \If{$a = \NIL$}
        \State \Return $b$
    \EndIf
    \If{$b = \NIL$}
        \State \Return $a$
    \EndIf
    \If{$b.\mathit{key} < a.\mathit{key}$}
        \State 交換 $a$ 和 $b$ \Comment{讓 $a$ 是根比較小的那棵}
    \EndIf
    \State $a.\mathit{right} \gets$ \Call{Meld}{$a.\mathit{right}, b$}
    \If{$s(a.\mathit{left}) < s(a.\mathit{right})$}
        \State 交換 $a.\mathit{left}$ 和 $a.\mathit{right}$
    \EndIf
    \State $a.s \gets s(a.\mathit{right}) + 1$
    \State \Return $a$
\EndFunction
\end{pseudo}

每一層遞迴都沿著 $a$ 或 $b$ 的右邊路徑往下走一步，所以遞迴的次數不超過兩條右邊路徑的長度
和，由定理~\ref{thm:hp-leftist} 是 $O(\lg n_a + \lg n_b) = O(\lg n)$。 有了合併，其他操作
都很簡單：\textbf{插入}是和只有一個節點的樹合併；\textbf{取出最小值}是拿掉根，把左右兩棵
子樹合併。 兩者都是 $O(\lg n)$。

\textbf{手算追蹤}。 下圖把 $A$、$B$ 兩棵最小左傾樹合併，節點右邊的小字是 $s$。 考卷上
的做法分兩步：先把兩條右邊路徑依鍵的大小合成一條，掛在右邊（圖 (c)）；再從下往上，遇
到左邊的 $s$ 比右邊小就對調左右，同時更新 $s$（圖 (d)）。 兩步的結果和虛擬碼的遞迴完
全一樣，只是把「往下」和「回來」分開畫。

\begin{center}
\tikzset{lt/.style={hpn, label={[hpidx, inner sep=1pt]east:#1}},
    ltlev/.style={level distance=0.95cm, edge from parent/.style={hpedge, draw},
        level 1/.style={sibling distance=2.0cm},
        level 2/.style={sibling distance=1.0cm},
        level 3/.style={sibling distance=0.75cm}}}
\begin{tikzpicture}[ltlev]
\node[lt=2] {12}
    child {node[lt=1] {28}
        child {node[lt=2] {33} child {node[lt=1] {61}} child {node[lt=1] {65}}}
        child[missing]}
    child {node[lt=1] {37}};
\node[font=\small] at (0,-3.5) {(a) $A$};
\end{tikzpicture}\hspace{1.2cm}%
\begin{tikzpicture}[ltlev]
\node[lt=2] {15}
    child {node[lt=2] {31} child {node[lt=1] {42}} child {node[lt=1] {49}}}
    child {node[lt=1] {23}};
\node[font=\small] at (0,-3.5) {(b) $B$};
\end{tikzpicture}

\medskip
\begin{tikzpicture}[ltlev, level 1/.style={sibling distance=3.0cm},
    level 2/.style={sibling distance=1.5cm}]
\node[hpb] {12}
    child {node[lt=1] {28}
        child {node[lt=2] {33} child {node[lt=1] {61}} child {node[lt=1] {65}}}
        child[missing]}
    child {node[hpa] {15}
        child {node[lt=2] {31} child {node[lt=1] {42}} child {node[lt=1] {49}}}
        child {node[hpb] {23} child[missing] child {node[hpa] {37}}}};
\node[font=\small] at (0,-3.6) {(c) 右邊路徑合成一條};
\end{tikzpicture}\hspace{0.6cm}%
\begin{tikzpicture}[ltlev, level 1/.style={sibling distance=3.0cm},
    level 2/.style={sibling distance=1.5cm}]
\node[lt=2] {12}
    child {node[lt=2] {15}
        child {node[lt=2] {31} child {node[lt=1] {42}} child {node[lt=1] {49}}}
        child {node[lt=1] {23} child {node[lt=1] {37}} child[missing]}}
    child {node[lt=1] {28}
        child {node[lt=2] {33} child {node[lt=1] {61}} child {node[lt=1] {65}}}
        child[missing]};
\node[font=\small] at (0,-3.6) {(d) 對調之後};
\end{tikzpicture}
\diagcaption{合併兩棵最小左傾樹。 (c) 中有顏色的是合成之後的右邊路徑 $12, 15, 23, 37$：
粉紅色是左邊的 $s$ 比右邊小、回來時要對調的節點，橘色是不用對調的。}
\label{fig:hp-leftist}
\end{center}

合併的過程：$A$ 的右邊路徑是 $12, 37$，$B$ 的是 $15, 23$。 依大小合成
$12 \to 15 \to 23 \to 37$，每個節點保留自己原本的左子樹（圖 (c)）。 接著由下往上：
$37$ 是葉節點，$s = 1$；$23$ 的左邊是空的（$s = 0$）、右邊 $s = 1$，對調，$s(23) = 1$；
$15$ 的左邊 $s(31) = 2$ 不小於右邊 $s(23) = 1$，不用對調，$s(15) = 2$；$12$ 的左邊
$s(28) = 1$ 小於右邊 $s(15) = 2$，對調，$s(12) = 1 + s(28) = 2$。

接著取出最小值 $12$：把它的左右子樹（根是 $15$ 和 $28$）合併。 右邊路徑 $15, 23$ 和
$28$ 合成 $15 \to 23 \to 28$；由下往上，$23$ 左右的 $s$ 都是 $1$，不用對調，$s(23) =
2$；$15$ 左右都是 $2$，也不用對調。 結果是 $15$ 的左邊 $31(42, 49)$、右邊 $23$，$23$ 的
左邊 $37$、右邊 $28(33(61, 65), \NIL)$。

\begin{examnote}
$s$ 的定義各書不同。 有的書叫 shortest$(x)$，和本書一樣 $s(\NIL) = 0$；有的書叫 null path
length（npl），規定 $\mathrm{npl}(\NIL) = -1$、葉節點是 $0$，每個值都比本書小 $1$。 左傾的
條件和合併的做法不受影響，但算數值時要照題目的定義。 另外有一種\textbf{重量左傾樹}
（weight-biased leftist tree），比較的是子樹的節點個數而不是 $s$，見觀念題第~\ref{hp-q-wblt}~題。
\end{examnote}

\subsection{二項堆積}

二項堆積用另一種方式做到快速合併：把元素分成好幾棵形狀固定的樹，合併時像二進位加法一
樣處理同樣大小的樹。

\textbf{二項樹}（binomial tree）$B_k$ 遞迴定義：$B_0$ 是一個節點；$B_k$ 由兩棵 $B_{k-1}$
組成，其中一棵的根變成另一棵的根的最左邊的子節點。

\begin{center}
\tikzset{bt/.style={circle, draw=black!65, fill=black!12, minimum size=0.28cm, inner sep=0pt},
    btlev/.style={level distance=0.7cm, edge from parent/.style={hpedge, draw}}}
\begin{tikzpicture}[btlev]
\node[bt] {};
\node[font=\small] at (0,-2.6) {$B_0$};
\end{tikzpicture}\hspace{1.0cm}%
\begin{tikzpicture}[btlev]
\node[bt] {} child {node[bt] {}};
\node[font=\small] at (0,-2.6) {$B_1$};
\end{tikzpicture}\hspace{1.0cm}%
\begin{tikzpicture}[btlev, level 1/.style={sibling distance=0.8cm}]
\node[bt] {} child {node[bt] {} child {node[bt] {}}} child {node[bt] {}};
\node[font=\small] at (0,-2.6) {$B_2$};
\end{tikzpicture}\hspace{1.0cm}%
\begin{tikzpicture}[btlev, level 1/.style={sibling distance=1.0cm},
    level 2/.style={sibling distance=0.6cm}]
\node[bt] {}
    child {node[bt] {} child {node[bt] {} child {node[bt] {}}} child {node[bt] {}}}
    child {node[bt] {} child {node[bt] {}}}
    child {node[bt] {}};
\node[font=\small] at (0,-2.6) {$B_3$};
\end{tikzpicture}
\diagcaption{二項樹 $B_0$ 到 $B_3$。 $B_3$ 的根有三個子節點，由左到右是 $B_2$、$B_1$、
$B_0$ 的根。}
\label{fig:hp-binomial-trees}
\end{center}

從定義用歸納法可以得到 $B_k$ 的性質：

\begin{enumerate}
\item 有 $2^k$ 個節點，高度是 $k + 1$（最深的節點深度是 $k$）。
\item 深度 $d$ 的節點剛好有 $\binom{k}{d}$ 個。 （$B_k$ 深度 $d$ 的節點，來自一棵 $B_{k-1}$
的深度 $d$ 和另一棵 $B_{k-1}$ 的深度 $d - 1$，所以個數滿足巴斯卡公式
$\binom{k}{d} = \binom{k-1}{d} + \binom{k-1}{d-1}$，這也是「二項」的由來。）
\item 根的子節點數（叫\textbf{階數}，degree）是 $k$，比其他節點都多；根的子節點由左到右
是 $B_{k-1}, B_{k-2}, \ldots, B_0$ 的根。
\end{enumerate}

\textbf{二項堆積}（binomial heap）是一群二項樹，每棵都滿足最小堆積的性質，而且\textbf{同一個
階數最多一棵}。 $n$ 個元素的二項堆積有哪幾棵樹是固定的：把 $n$ 寫成二進位，第 $k$ 位是
$1$ 就有一棵 $B_k$。 例如 $n = 13 = 1101_2$，就是 $B_3$、$B_2$、$B_0$ 各一棵。 所以樹最多
$\lfloor \lg n \rfloor + 1$ 棵，每棵的階數不超過 $\lfloor \lg n \rfloor$，高度不超過
$\lfloor \lg n \rfloor + 1$。 各樹的
根用串列串起來，叫\textbf{根串列}（root list），依階數由小到大排。

每個節點的子節點個數不固定，所以用第~\ref{sec:bt-lcrs}~節的左孩子右兄弟存：
$x.\mathit{child}$ 指向最左邊的子節點，$x.\mathit{sibling}$ 指向右邊的下一個兄弟，另外
$x.p$ 指向父節點，$x.\mathit{degree}$ 記子節點個數。 根串列就是用各棵樹的根的
$\mathit{sibling}$ 串起來的。

基本的動作是把兩棵同階的樹連成一棵：根比較大的那棵掛到另一棵的根底下，成為最左邊的子
節點。 一棵 $B_k$ 和一棵 $B_k$ 連成 $B_{k+1}$，堆積性質仍然成立，$O(1)$。

\begin{pseudo}[label=alg:hp-binomial-link]{連結兩棵同階的二項樹}
\Function{Binomial-Link}{$y, z$} \Comment{$y$、$z$ 是同階的根}
    \If{$z.\mathit{key} < y.\mathit{key}$}
        \State 交換 $y$ 和 $z$
    \EndIf
    \State $z.p \gets y$
    \State $z.\mathit{sibling} \gets y.\mathit{child}$ \Comment{$z$ 變成 $y$ 最左邊的子節點}
    \State $y.\mathit{child} \gets z$
    \State $y.\mathit{degree} \gets y.\mathit{degree} + 1$
    \State \Return $y$
\EndFunction
\end{pseudo}

其他操作都從\textbf{合併}來：

\begin{itemize}
\item \textbf{合併}兩個二項堆積：從階數 $0$ 開始往上，像二進位加法。 每個階數 $k$ 可能有
$0$ 到 $3$ 棵樹（兩個堆積各一棵，加上前一位的進位）。 $0$ 或 $1$ 棵就留下；$2$ 棵就連成
一棵 $B_{k+1}$ 進位；$3$ 棵就留下一棵，另外兩棵連成進位。 根串列長 $O(\lg n)$，共
$O(\lg n)$。
\item \textbf{插入}：和只有一棵 $B_0$ 的堆積合併。 最壞要連結 $O(\lg n)$ 次（$n = 2^k - 1$
時一路進位），但這和第~3~章的二進位計數器一樣：連結的次數等於 $n$ 的二進位尾端連續的
$1$ 的個數，平攤只有 $O(1)$。 換個方式看，從空的堆積連續插入 $n$ 次，一共產生 $n$ 棵
$B_0$，每次連結讓樹少一棵，最後剩 $b(n)$ 棵，所以總共剛好連結 $n - b(n)$ 次。
\item \textbf{找最小值}：最小值一定是某棵樹的根，掃過根串列，$O(\lg n)$（另外記一個指向
最小根的指標，就是 $O(1)$）。
\item \textbf{取出最小值}：找到最小的根，從根串列拿掉；它的子節點由左到右是
$B_{k-1}, \ldots, B_0$，反過來排就是一個合法的二項堆積，再和剩下的堆積合併。
$O(\lg n)$。
\item \textbf{減少鍵值}：在那棵樹裡像二元堆積一樣往上換，樹高 $O(\lg n)$。 \textbf{刪除}
任意元素：把它的鍵減成 $-\infty$，再取出最小值。
\end{itemize}

\textbf{手算追蹤}。 把 $27, 13, 41, 8, 35, 19, 50$ 依序插入空的二項堆積，再取出最小值。
下表的「之後」用 $r[\,\cdots]$ 表示根是 $r$、方括號裡由左到右是它的子樹。

\begin{center}
\begin{tabular}{lccl}
\toprule
插入 & 插入前的 $n$ & 連結次數 & 之後的根串列 \\
\midrule
27 & $0 = 000_2$ & 0 & $B_0{:}\,27$ \\
13 & $1 = 001_2$ & 1 & $B_1{:}\,13[27]$ \\
41 & $2 = 010_2$ & 0 & $B_0{:}\,41$，$B_1{:}\,13[27]$ \\
8 & $3 = 011_2$ & 2 & $B_2{:}\,8[13[27],\,41]$ \\
35 & $4 = 100_2$ & 0 & $B_0{:}\,35$，$B_2{:}\,8[13[27],\,41]$ \\
19 & $5 = 101_2$ & 1 & $B_1{:}\,19[35]$，$B_2{:}\,8[13[27],\,41]$ \\
50 & $6 = 110_2$ & 0 & $B_0{:}\,50$，$B_1{:}\,19[35]$，$B_2{:}\,8[13[27],\,41]$ \\
\bottomrule
\end{tabular}
\end{center}

插入 $8$ 時，$B_0{:}\,8$ 和 $B_0{:}\,41$ 連成 $B_1{:}\,8[41]$，進位後和 $B_1{:}\,13[27]$ 連
成 $B_2$，根是較小的 $8$，$13$ 那棵掛在最左邊。 連結次數總共 $4 = 7 - b(7)$。

取出最小值 $8$：它的子樹 $B_1{:}\,13[27]$、$B_0{:}\,41$ 自成一個堆積，和剩下的
$\{B_0{:}\,50,\ B_1{:}\,19[35]\}$ 合併。 階數 $0$ 有 $50$、$41$ 兩棵，連成 $41[50]$ 進位；
階數 $1$ 有 $19[35]$、$13[27]$ 和進位的 $41[50]$ 三棵，留下進位的那棵，另外兩棵連成
$13[19[35],\,27]$ 進位到階數 $2$。 結果是 $6 = 110_2$ 個元素，兩棵樹。

\begin{center}
\tikzset{btlev/.style={level distance=0.9cm, edge from parent/.style={hpedge, draw},
    level 1/.style={sibling distance=1.0cm}}}
\begin{tikzpicture}[btlev]
\node[hpn] (r0) at (0,0) {50};
\node[hpn] (r1) at (1.3,0) {19} child {node[hpn] {35}};
\node[hpn] (r2) at (3.3,0) {8} child {node[hpn] {13} child {node[hpn] {27}}}
    child {node[hpn] {41}};
\draw[densely dashed, black!50] (r0) -- (r1) -- (r2);
\node[font=\small] at (1.8,-2.6) {(a) 插入 $7$ 個之後};
\end{tikzpicture}\hspace{1.4cm}%
\begin{tikzpicture}[btlev]
\node[hpn] (s1) at (0,0) {41} child {node[hpn] {50}};
\node[hpn] (s2) at (2.0,0) {13} child {node[hpn] {19} child {node[hpn] {35}}}
    child {node[hpn] {27}};
\draw[densely dashed, black!50] (s1) -- (s2);
\node[font=\small] at (1.2,-2.6) {(b) 取出 $8$ 之後};
\end{tikzpicture}
\diagcaption{二項堆積。 虛線是根串列，依階數由小到大。}
\label{fig:hp-binomial-heap}
\end{center}

\begin{examnote}
合併時某個階數有三棵樹，留哪一棵、連哪兩棵，各書的規定不同，結果都是合法的二項堆積，
但形狀不一樣。 作答時寫明你的規則（本書留下進位來的那棵）。 另外有的書把二項堆積寫成
「延遲合併」的版本：插入和合併只把根串列接起來，$O(1)$，等到取出最小值時才把同階的樹
兩兩連結。 如果題目的圖裡有兩棵同階的樹，就是這個版本，取出最小值時要把所有同階的樹
連到每個階數只剩一棵為止。 下一節的費氏堆積就是從延遲合併的版本再往前走一步。
\end{examnote}

\section{費氏堆積的平攤成本}
\label{sec:hp-fib}

前面的堆積，\proc{Decrease-Key} 都是 $O(\lg n)$。 Dijkstra 演算法在邊很多的圖上會做
$\Theta(m)$ 次 \proc{Decrease-Key}，卻只做 $n$ 次 \proc{Extract-Min}；如果
\proc{Decrease-Key} 能變成 $O(1)$，總時間就能從 $O(m \lg n)$ 降到 $O(m + n \lg n)$。
\textbf{費氏堆積}（Fibonacci heap）做到了這件事，不過是\textbf{平攤}的 $O(1)$。 這一節只
講它怎麼做、各操作的平攤成本是多少，以及怎麼用第~3~章的位能法得到這些結果，不做逐步
追蹤。

\subsection{結構與操作}

費氏堆積也是一群滿足最小堆積性質的樹，根用環狀雙向串列串起來，另外有一個指標指向最小
的根。 和二項堆積相比，它更懶：

\begin{itemize}
\item \textbf{插入}：新節點直接加進根串列，必要時更新最小指標。 不做任何連結。
\item \textbf{合併}：兩個根串列接起來，最小指標取兩者較小的。
\item \textbf{取出最小值}：拿掉最小的根，它的子節點都加進根串列；然後\textbf{整理}
（consolidate）根串列，把同階數的樹兩兩連結，直到每個階數只剩一棵，順便找出新的最
小值。 插入、合併時沒做的連結，都留到這一步才做。
\item \textbf{減少鍵值}：鍵變小以後如果比父節點小，就把這個節點連同它的子樹\textbf{剪
下來}，放進根串列。 為了不讓樹被剪得太稀疏，每個節點有一個標記：非根的節點第一次失去
子節點時做標記；已經有標記的節點再失去一個子節點，自己也被剪下來，再往上檢查它的父節
點，這叫\textbf{連鎖剪下}（cascading cut）。
\item \textbf{刪除}：減少鍵值到 $-\infty$，再取出最小值。
\end{itemize}

連鎖剪下保證一個階數是 $k$ 的節點，子樹至少有 $F_{k+2} \ge \phi^k$ 個節點，$F$ 是費氏數、
$\phi = (1 + \sqrt{5})/2$（名字就是這麼來的）。 所以 $n$ 個節點的費氏堆積，最大的階數
$D(n) \le \log_\phi n = O(\lg n)$。 這個結論這裡不證明，只拿來用。

\subsection{平攤成本}

\begin{center}
\begin{tabular}{lccc}
\toprule
操作 & 二元堆積（最壞） & 二項堆積（最壞） & 費氏堆積（平攤） \\
\midrule
\proc{Insert} & $O(\lg n)$ & $O(\lg n)$ & $O(1)$ \\
\proc{Minimum} & $\Theta(1)$ & $O(\lg n)$ & $O(1)$ \\
\proc{Extract-Min} & $O(\lg n)$ & $O(\lg n)$ & $O(\lg n)$ \\
合併 & $\Theta(n)$ & $O(\lg n)$ & $O(1)$ \\
\proc{Decrease-Key} & $O(\lg n)$ & $O(\lg n)$ & $O(1)$ \\
\proc{Delete} & $O(\lg n)$ & $O(\lg n)$ & $O(\lg n)$ \\
\bottomrule
\end{tabular}
\end{center}

二元堆積的合併是把兩個陣列接起來再由下而上建一次。 二項堆積的 \proc{Insert} 平攤是
$O(1)$（上一節），\proc{Minimum} 多記一個指標也是 $O(1)$。

表上費氏堆積那一欄是平攤成本：任何一串操作，實際的總成本不超過各操作平攤成本的和，但
單獨一次操作可能很貴。 例如連續插入 $n$ 個元素之後，根串列有 $n$ 棵只有一個節點的樹，
接下來那一次 \proc{Extract-Min} 要整理 $n - 1$ 棵樹，實際成本是 $\Theta(n)$。 需要每次
操作都有保證的場合（例如即時系統），平攤的界不夠用。 另外它的常數大、指標多，實務上
不一定比二元堆積或 $d$ 元堆積快；它的價值主要在理論上的時間界，第~24、25~章會看到
Prim 和 Dijkstra 用它得到 $O(m + n \lg n)$。

\begin{whybox}[用位能法算費氏堆積的平攤成本]
設 $t(H)$ 是根串列裡樹的棵數，$m(H)$ 是有標記的節點數，取位能
\[
\Phi(H) = t(H) + 2\,m(H).
\]
空堆積 $\Phi = 0$，而且 $\Phi$ 永遠不是負的，所以對所有 $i$ 都有 $\Phi(D_i) \ge
\Phi(D_0)$，由定理~\ref{thm:potential}，平攤成本的總和是實際成本總和的上界。 下面把位能的一個
單位定成「夠付一個常數時間的工作」，所以 $O(1)$ 的實際工作可以用一個單位的位能抵銷。

\textbf{插入}：實際 $O(1)$；多一棵樹，$\Delta\Phi = 1$。 平攤 $O(1)$。 \textbf{合併}：實際
$O(1)$；樹和標記都只是加在一起，$\Delta\Phi = 0$。 平攤 $O(1)$。

\textbf{取出最小值}：最小的根有最多 $D(n)$ 個子節點，加進根串列後最多 $t(H) + D(n) - 1$
棵樹要整理，每次連結少一棵，實際成本 $O(t(H) + D(n))$。 整理完每個階數最多一棵，最
多剩 $D(n) + 1$ 棵；標記沒有增加。 所以
\[
\hat{c} \le O(t(H) + D(n)) + \bigl(D(n) + 1 + 2m(H)\bigr) - \bigl(t(H) + 2m(H)\bigr)
= O(D(n)) = O(\lg n).
\]
$t(H)$ 那一項被位能的下降抵銷了，這就是「前面插入時多付的錢」。

\textbf{減少鍵值}：設連鎖剪下總共剪了 $c$ 次（包括第一次剪下 $x$），實際成本 $O(c)$。
根串列多了 $c$ 棵樹；前 $c - 1$ 次是剪下有標記的節點，標記少了 $c - 1$ 個，最後最多多
一個新標記。 所以
\[
\Delta\Phi \le c + 2\bigl(-(c - 1) + 1\bigr) = 4 - c,
\qquad \hat{c} \le O(c) + 4 - c = O(1).
\]
標記的係數取 $2$ 就是為了這裡：每剪掉一個有標記的節點，一個單位付它變成新樹的位能，另
一個單位付剪下的工作。

\textbf{刪除}：減少鍵值 $O(1)$ 加上取出最小值 $O(\lg n)$。
\end{whybox}

\begin{examnote}
費氏堆積的題目如果問「哪些操作是 $O(1)$」，作答時要寫清楚是平攤還是最壞：
\proc{Extract-Min} 的平攤是 $O(\lg n)$，單獨一次最壞可以是 $\Theta(n)$；
\proc{Decrease-Key} 的平攤是 $O(1)$，單獨一次的連鎖剪下可能很長。 只寫「費氏堆積的插入
是 $O(1)$」而不說平攤或實際，在證明題裡會被認為不完整。
\end{examnote}

\section{本章常見錯誤}
\label{sec:hp-pitfalls}

\begin{center}
\begin{tabular}{>{\raggedright\arraybackslash}p{0.24\linewidth}>{\raggedright\arraybackslash}p{0.50\linewidth}>{\raggedright\arraybackslash}p{0.16\linewidth}}
\toprule
錯誤 & 例子 & 見 \\
\midrule
下沉時換錯子節點 & 和第一個比自己大的子節點換，而不是和較大的那個換，換上來的數比另一
個子節點小 & 第~\ref{sec:hp-binary}~節 \\
C 的位置公式用錯 & 從 $0$ 開始卻寫父節點 \texttt{i / 2}、子節點 \texttt{2 * i}，
\texttt{a[0]} 的子節點變成自己 & 第~\ref{sec:hp-binary}~節 \\
以為堆積是排好的 & 以為最大堆積的最小值在 $A[n]$；以為同一層由左到右遞增 &
第~\ref{sec:hp-binary}~節 \\
建堆積的方向反了 & \proc{Build-Max-Heap} 的 $i$ 從 $1$ 往上做，子樹還不是堆積就下沉 &
第~\ref{sec:hp-build}~節 \\
把建堆積算成 $\Theta(n \lg n)$ & 把「$n/2$ 次下沉、每次 $O(\lg n)$」當成答案；那只是上界，不是緊的界 &
第~\ref{sec:hp-build}~節 \\
刪除任意元素只做下沉 & 補進來的最後一個元素比新的父節點小，應該上浮 &
第~\ref{sec:hp-pq}~節 \\
帶索引的堆積忘了更新位置 & 交換兩格時只改了陣列，沒改 $\mathit{pos}$，之後的
decrease-key 找錯格 & 第~\ref{sec:hp-pq}~節 \\
用堆積排同鍵的先後 & 以為同級的病人會照到達順序出來，沒有把到達順序放進鍵 &
第~\ref{sec:hp-pq}~節 \\
$d$ 元堆積公式混用編號 & 從 $1$ 編號卻寫子節點 $di + 1$ & 第~\ref{sec:hp-dary}~節 \\
最小最大堆積只看父節點 & 插入時和父節點比完以後，繼續和父節點比，而不是和祖父比 &
第~\ref{sec:hp-mergeable}~節 \\
deap 取出最大值只看對應點 & 對應點有子節點時沒有改看較大的子節點，那個子節點變得比它的
對應位置大 & 第~\ref{sec:hp-deap}~節 \\
SMMH 和根的子節點比 & 插入時應該和「祖父的兩個子節點」比；或是忘了先和左兄弟比 &
第~\ref{sec:hp-smmh}~節 \\
左傾樹依鍵對調左右 & 回來時比較的是左右子節點的鍵，應該比較 $s$；忘了更新 $s$ &
第~\ref{sec:hp-mergeable}~節 \\
二項堆積連結錯樹 & 把不同階的樹連在一起；連結後子節點的順序放錯 &
第~\ref{sec:hp-mergeable}~節 \\
平攤和最壞混為一談 & 說「費氏堆積的 \proc{Extract-Min} 最壞是 $O(\lg n)$」 &
第~\ref{sec:hp-fib}~節 \\
\bottomrule
\end{tabular}
\end{center}

\section{觀念題}
\label{sec:hp-exercises}

題型和研究所考試、期中期末考的筆試相近，有選擇、計算、追蹤和證明。 陣列一律從 $1$ 開始
編號，下沉時兩個子節點相等就換左邊。 解答在本章最後。

\begin{enumerate}
\item （選擇，可能不只一個答案）下列哪些陣列是最大堆積？
\begin{enumerate}
\item $\langle 90, 70, 80, 20, 60, 75, 10 \rangle$
\item $\langle 90, 80, 70, 75, 60, 85, 50 \rangle$
\item $\langle 50, 50, 40, 50, 20 \rangle$
\item $\langle 100, 30, 90, 40, 20 \rangle$
\item 任何一個由大到小排好的陣列
\end{enumerate}

\item （追蹤）把 $24, 11, 39, 7, 52, 30, 18, 45$ 依序插入一個空的\textbf{最小}堆積，寫出每次
插入後的陣列；再連續取出兩次最小值，寫出每次之後的陣列。

\item （追蹤）$A = \langle 33, 18, 56, 9, 41, 27, 62, 15, 48 \rangle$。
\begin{enumerate}
\item 用 \proc{Build-Max-Heap} 建成最大堆積，寫出每次 \proc{Sift-Down} 之後的陣列和總交
換次數。
\item 改成把這 $9$ 個數依序插入空的最大堆積，最後的陣列是什麼？
\end{enumerate}

\item （證明）$n$ 個元素互不相同的最大堆積。
\begin{enumerate}
\item 證明第 $k$ 大的元素深度不超過 $k - 1$。
\item 由 (a) 推出第 $2$ 大的元素一定在位置 $2$ 或 $3$。
\item 證明最小的元素一定是葉節點，並寫出它可能在的位置範圍。
\end{enumerate}

\item\label{hp-q-build} （證明）用迴圈不變量證明 \proc{Build-Max-Heap}（虛擬碼~\ref{alg:hp-build}）的正確性，
寫出初始、維持、結束三步。 可以直接引用「若 $i$ 的兩棵子樹都是最大堆積，\proc{Sift-Down}
$(A, i)$ 之後以 $i$ 為根的子樹是最大堆積」。

\item\label{hp-q-sum} （計算）把 $1, 2, \ldots, n$ 依序插入空的最大堆積，總交換次數是
$S(n) = \sum_{i=1}^{n} \lfloor \lg i \rfloor$。
\begin{enumerate}
\item 求 $S(10)$。
\item 證明 $S(n) = \Theta(n \lg n)$。
\end{enumerate}

\item （選擇，可能不只一個答案）下列關於優先佇列的敘述，哪些正確？
\begin{enumerate}
\item 用未排序的陣列實作，\proc{Insert} 是 $\Theta(1)$，\proc{Extract-Min} 是 $\Theta(n)$。
\item 在二元最小堆積裡找最大值只要 $O(\lg n)$。
\item 二元堆積的 \proc{Decrease-Key} 如果只給要改的鍵值、不給位置，最壞要 $\Theta(n)$ 才
找得到它。
\item 兩個各有 $n$ 個元素的二元堆積可以在 $O(n)$ 時間合併成一個。
\item 二元最小堆積中，同一層的元素由左到右遞增。
\end{enumerate}

\item （計算）$d$ 元堆積，從 $1$ 開始編號。
\begin{enumerate}
\item $d = 4$ 時，位置 $100$ 的父節點是哪一格？ 那一格的子節點是哪幾格？
\item $d = 4$、$n = 1000$ 時，樹的高度是多少？
\item 寫出取出最小值最多要比較幾次（以 $d$ 和高度 $h$ 表示），說明理由。
\item 一個應用會做 $10^6$ 次 decrease-key 和 $10^3$ 次取出最小值，堆積裡大約有 $10^3$ 個
元素。 $d = 2$ 和 $d = 32$ 哪一個比較好？ 用比較次數的上界估計。
\end{enumerate}

\item （追蹤）下面是一個最小最大堆積：$\langle 10, 65, 70, 20, 40, 30, 60, 25, 45, 50 \rangle$。
\begin{enumerate}
\item 畫出這棵樹，標出最小層和最大層，並檢查它確實滿足定義。
\item 插入 $15$，寫出之後的陣列。
\item 接著取出最大值，寫出之後的陣列。
\item 再取出最小值，寫出之後的陣列。
\end{enumerate}

\item\label{hp-q-deap} （追蹤）下面是一個 deap（根 $A[1]$ 是空的），
$A[2..10] = \langle 13, 68, 27, 16, 61, 44, 40, 28, 34 \rangle$。
\begin{enumerate}
\item 畫出這棵樹。 對最小堆積的位置 $4$ 到 $10$，寫出各自要和哪個位置比（對應點或它的父節點），
並檢查定義~\ref{def:hp-deap} 成立。
\item 依序插入 $94$、$21$，寫出每次插入後的陣列。
\item 接著取出最大值，寫出之後的陣列。 如果取出時只和對應點本身比（不看它的子節點），會得
到什麼陣列？ 哪裡違反定義？
\item 從 (c) 正確的結果再取出最小值，寫出之後的陣列。
\end{enumerate}

\item\label{hp-q-smmh} （選擇、追蹤）對稱最小最大堆積。
\begin{enumerate}
\item 下列陣列都是 $A[2..]$（根 $A[1]$ 是空的），哪些是 SMMH？ 不是的，指出違反 P1、P2、P3
的哪一條。 (A) $\langle 5, 50, 10, 40, 20, 30 \rangle$ \quad (B) $\langle 5, 50, 10, 40, 30, 20
\rangle$ \quad (C) $\langle 5, 50, 10, 40, 20, 30, 15, 35, 25 \rangle$ \quad (D) $\langle 5, 50,
10, 40, 20, 30, 15, 35, 45 \rangle$
\item 把 $47, 18, 83, 60, 25, 91, 12, 70, 35$ 依序插入空的 SMMH，寫出每次插入後的陣列。
\item 接著取出最大值，再取出最小值，寫出每次之後的陣列。
\item 把 (b) 的結果每一對兄弟合成一個節點，寫出得到的區間堆積，並檢查每個節點的區間都包
含子節點的區間。
\end{enumerate}

\item （追蹤）合併下面兩棵最小左傾樹，畫出合併後的樹並標出每個節點的 $s$；再取出最小
值，畫出結果。

\begin{center}
\tikzset{ltq/.style={level distance=0.9cm, edge from parent/.style={hpedge, draw},
    level 1/.style={sibling distance=1.8cm}, level 2/.style={sibling distance=0.9cm},
    every node/.style={hpn}}}
\begin{tikzpicture}[ltq]
\node {23}
    child {node {26} child {node {32} child {node {62}} child {node {60}}} child[missing]}
    child {node {61}};
\node[draw=none, font=\small] at (0,-3.3) {$A$};
\end{tikzpicture}\hspace{1.5cm}%
\begin{tikzpicture}[ltq]
\node {24}
    child {node {27} child {node {38}} child[missing]}
    child {node {39} child {node {57}} child[missing]};
\node[draw=none, font=\small] at (0,-2.4) {$B$};
\end{tikzpicture}
\end{center}

\item\label{hp-q-wblt} （證明）\textbf{重量左傾樹}：令 $w(x)$ 是以 $x$ 為根的子樹的節點數（$w(\NIL) = 0$），
規定每個節點都有 $w(x.\mathit{left}) \ge w(x.\mathit{right})$。 證明從根一路往右走到 \NIL，
經過的節點數不超過 $\lg(n + 1)$。

\item （計算、追蹤）二項堆積。
\begin{enumerate}
\item $n = 13$ 的二項堆積有幾棵樹？ 各是哪一階？
\item $B_5$ 有幾個節點在深度 $2$？
\item 把 $36, 14, 52, 29, 8, 61, 23, 47, 17, 40$ 依序插入空的二項堆積，寫出最後的各棵樹，
以及總共連結了幾次。
\item 接著取出最小值，寫出結果。
\end{enumerate}

\item （選擇，可能不只一個答案；計算）
\begin{enumerate}
\item 下列關於費氏堆積的敘述，哪些正確？ (A) \proc{Extract-Min} 最壞情況是 $O(\lg n)$。
(B) 從空的堆積插入 $n$ 個元素，接著的第一次 \proc{Extract-Min} 實際成本是 $\Theta(n)$。
(C) 某一次 \proc{Decrease-Key} 的實際成本可能超過常數。 (D) 位能 $t(H) + 2m(H)$ 在任何
時候都不小於 $0$。
\item Dijkstra 演算法在 $n$ 個點、$m$ 條邊的圖上做 $n$ 次 \proc{Insert}、$n$ 次
\proc{Extract-Min}、最多 $m$ 次 \proc{Decrease-Key}。 分別寫出用二元堆積和費氏堆積的總
時間；$m = \Theta(n^2)$ 時兩者各是多少？ $m = \Theta(n)$ 時呢？
\end{enumerate}
\end{enumerate}

\section{歷屆試題}
\label{sec:hp-pastexam}

以下是研究所入學考試的原題，題目照原卷（語言、記號、配分照錄）。 詳解在本章最後。

\begin{pastexam}[label={px:hp-1}]{中正資工所甲組 114 學年度 第 3 題}
(5\%) (Choose \underline{one} best answer) Consider the integer array \{52, 21, 40, 33, 88, 57, 46, 71\}. What will the array look like after heapifying it to a Max Heap? Note that we use the bottom-up method that runs in $O(n)$ time.
\begin{enumerate}[label=(\alph*)]
\item \{88, 71, 57, 52, 21, 40, 46, 33\}
\item \{88, 71, 57, 52, 46, 40, 33, 21\}
\item \{21, 33, 40, 46, 52, 57, 71, 88\}
\item \{88, 71, 57, 52, 33, 40, 46, 21\}
\item None of the above.
\end{enumerate}
\end{pastexam}

\begin{pastexam}[label={px:hp-2}]{臺科大資工所 110 學年度 第 6 題（節錄 (d)(e) 小題）}
Answer the following questions. (28\%)
\begin{enumerate}[label=(\alph*), start=4]
\item Draw the binary min heap that results from inserting 12, 9, 14, 16, 4, 17, 6, 8, 2 in that order into an initially empty binary heap. Draw the final tree. (10\%)
\item Remove the first two smallest values from the heap created in (d) and reconstruct the min heap. Draw the resulting tree. (5\%)
\end{enumerate}
\end{pastexam}

\begin{pastexam}[label={px:hp-3}]{成大資訊聯招 113 學年度 第 2 題}
\textbf{[1\%, 1\%, 3\%]} Considering the data pairs with keys in the given order: 20, 5, 10, 18, 4, 22, 11, 32, 21, as inputs to create a tree.
\begin{enumerate}[label=\roman*)]
\item If you create a \underline{max heap} for them, what is the level of the node containing 5?
\begin{enumerate}[label=(\Alph*)]
\item 1
\item 2
\item 3
\item 4
\item None of the above
\end{enumerate}
\item If you create a \underline{min heap} for them, which of the following statements is true?
\begin{enumerate}[label=(\Alph*)]
\item The node containing 11 is in the left subtree of the root.
\item The node containing 18 is in the right subtree of the root.
\item The node containing 20 is the root.
\item The node containing 21 is at level 3.
\item None of the above.
\end{enumerate}
\item If you create a \underline{symmetric min-max heap} for them, which of the following statements is true?
\begin{enumerate}[label=(\Alph*)]
\item To delete the node containing 5, the node containing 21 is used to replace the deleted node.
\item There are two nodes at level 4.
\item The leftmost grandchild of the root is 5.
\item The keys of the nodes in the right subtree of the root are 32, 10, and 20.
\item None of the above.
\end{enumerate}
\end{enumerate}
\end{pastexam}

\begin{pastexam}[label={px:hp-4}]{臺大資工所 113 學年度 第 2 題}
（原卷說明：本題為複選題。）

(5 points) A Maximum Binary Tree (abbreviated as MaxBT) is defined as a binary tree that adheres to the max-heap property without the necessity of following the complete binary tree property. In a MaxBT, each node $p$ consists of a value ($p$.data) and two links ($p$.left and $p$.right) pointing to the left and right children, respectively. A link to nil indicates the absence of further children.

The provided pseudocode, although incomplete, implements a recursive function for merging two MaxBTs into a single MaxBT. This function takes the root nodes of the two MaxBTs (\emph{root1} and \emph{root2}) as input and returns the root node of the merged MaxBT.
\begin{center}
\small
\begin{tabular}{rl}
1: & function MergeMaxBinaryTrees(\emph{root1}, \emph{root2}) \\
2: & \quad if \emph{root1} = nil then return \underline{\makebox[1.6cm]{(a)}} end if \\
3: & \quad if \emph{root2} = nil then return \underline{\makebox[1.6cm]{(b)}} end if \\
4: & \quad if \emph{root1}.data $>$ \emph{root2}.data then \\
5: & \quad\quad \underline{\makebox[1.6cm]{(c)}} $\leftarrow$ MergeMaxBinaryTrees(\underline{\makebox[1.6cm]{(d)}}, \underline{\makebox[1.6cm]{(e)}}) \\
6: & \quad\quad return \underline{\makebox[1.6cm]{(f)}} \\
7: & \quad else \\
8: & \quad\quad \underline{\makebox[1.6cm]{(g)}} $\leftarrow$ MergeMaxBinaryTrees(\underline{\makebox[1.6cm]{(h)}}, \underline{\makebox[1.6cm]{(i)}}) \\
9: & \quad\quad return \underline{\makebox[1.6cm]{(j)}} \\
10: & \quad end if \\
11: & end function
\end{tabular}
\end{center}
Assuming all the blanks in the above function have been appropriately filled to ensure its proper functioning, please select the correct description(s) below.
\begin{enumerate}[label=\Alph*.]
\item (b) and (f) are the same.
\item (c), (d), and (e) are all distinct.
\item In the cases of (c) and (g), each has only one correct answer.
\item Exchanging the filled content of (h) and (i) leads to a correct implementation.
\item The function works correctly even when the two arguments passed into it refer to the same tree.
\end{enumerate}
\end{pastexam}

\begin{pastexam}[label={px:hp-5}]{成大電通所 110 學年度 第 2 題（節錄 (b)(c) 小題）}
(\textbf{20 pts}) Heap
\begin{enumerate}[label=\alph*., start=2]
\item (4 \textbf{pts}) Describe the properties of a binomial min heap (i.e., four conditions).
\item (4 \textbf{pts}) Please show the resulting figure after you delete the node 1 in the following binomial min heap

\begin{center}
\begin{tikzpicture}[x=0.9cm,y=-0.9cm,
  v/.style={circle,draw,thick,inner sep=0pt,minimum size=6mm,font=\small},
  a/.style={-{Latex[length=2.2mm]},thick}]
\node[anchor=east] (hd) at (2.9,0.35) {Head[H]};
\node[v] (r14) at (4.3,0.35) {14};
\node[v] (r1) at (6.65,0.35) {1};
\node[v] (r6) at (11.75,0.35) {6};
\node[v] (a10) at (5.3,1.35) {10};
\node[v] (a23) at (6.65,1.35) {23};
\node[v] (a20) at (5.3,2.45) {20};
\node[v] (b10) at (9.5,1.35) {10};
\node[v] (b7) at (10.6,1.35) {7};
\node[v] (b15) at (11.75,1.35) {15};
\node[v] (b13) at (8.4,2.45) {13};
\node[v] (b12) at (9.5,2.45) {12};
\node[v] (b19) at (10.6,2.45) {19};
\node[v] (b30) at (8.4,3.45) {30};
\draw[a] (hd) -- (r14);
\draw[a] (r14) -- (r1);
\draw[a] (r1) -- (r6);
\draw[a] (r1) -- (a10);
\draw[a] (r1) -- (a23);
\draw[a] (a10) -- (a20);
\draw[a] (r6) -- (b10);
\draw[a] (r6) -- (b7);
\draw[a] (r6) -- (b15);
\draw[a] (b10) -- (b13);
\draw[a] (b10) -- (b12);
\draw[a] (b7) -- (b19);
\draw[a] (b13) -- (b30);
\end{tikzpicture}
\end{center}
\end{enumerate}
\end{pastexam}

\section{實作練習}
\label{sec:hp-practice}

下面是 Aki-Judge 上和本章相關的題目。 前幾題只要交標頭檔和函式，不寫 \texttt{main}；交
之前先用本章的範例和觀念題的資料自己測一次，特別是空的堆積、只有一個元素、所有元素相
等的情況。

\begin{center}
\begin{tabular}{llc}
\toprule
小節 & 題目 & 難度 \\
\midrule
二元堆積 & \akiref{check-min-heap}{檢查陣列是否為最小堆積} & 暖身 \\
二元堆積 & \akiref{min-heap-operations}{最小堆積的基本操作} & 基本 \\
由下而上建堆積 & \akiref{bottom-up-max-heap}{由下而上建立最大堆積} & 基本 \\
由下而上建堆積 & \akiref{heapsort-an-array}{以堆積排序整數陣列} & 基本 \\
優先佇列 & \akiref{kth-largest-element}{第 K 大的元素} & 基本 \\
優先佇列 & \akiref{merge-sorted-lists}{合併多個已排序數列} & 基本 \\
優先佇列 & \akiref{priority-queue-with-decrease}{可降低數值的優先佇列} & 進階 \\
優先佇列 & \akiref{halving-string-queue}{對半拆字串佇列} & 進階 \\
$d$ 元堆積 & \akiref{ternary-min-heap-with-replace}{三叉最小堆積與替換操作} & 進階 \\
\bottomrule
\end{tabular}
\end{center}

幾題的提示：

\begin{itemize}
\item 「由下而上建立最大堆積」的題目把這個方法叫做 Floyd 的方法（R. W. Floyd 在 1964 年
提出），就是第~\ref{sec:hp-build}~節的 \proc{Build-Max-Heap}，C 從 \texttt{n / 2 - 1} 倒
著做到 \texttt{0}。
\item 「以堆積排序整數陣列」先讀第~\ref{sec:hp-build}~節最後一小段，完整的虛擬碼和分析在第~\ref{sec:so-heap}~節。 常見的錯是每次把根
換到尾端以後，下沉時仍然用原本的大小 $n$，把已經排好的尾端又拉回堆積。
\item 「第 K 大的元素」和「合併多個已排序數列」照第~\ref{sec:hp-pq}~節「其他常見用法」的
做法，時間分別是 $O(n \lg k)$ 和 $O(N \lg k)$。
\item 「可降低數值的優先佇列」就是範例~\ref{ex:hp-er-queue} 的帶索引堆積：開一個位置陣
列，每次交換都更新。
\item 「對半拆字串佇列」每次取出字典序最小的片段再切成兩半放回去。 片段不必複製字串，記
住它在原字串的起點和長度就好，比較時用 \texttt{memcmp} 比共同的長度，再比長度。
\item 「三叉最小堆積與替換操作」照第~\ref{sec:hp-dary}~節的公式；題目規定幾個子節點一樣
小時取位置最小的，\texttt{R x} 就是第~\ref{sec:hp-pq}~節說的 replace，只做一次下沉，和
「先取出再插入」得到的陣列不一樣。
\end{itemize}

\section{讀完你應該會}
\label{sec:hp-checklist}

\begin{itemize}
\item 給一串數，畫出依序插入二元堆積（最大或最小）的每一步，以及取出最大（最小）值時下
沉的每一步，樹和陣列都會寫；從 $0$ 或從 $1$ 編號的位置公式都能正確使用。
\item 用由下而上的方法建堆積，寫出逐步陣列，並證明它的時間是 $\Theta(n)$。
\item 寫出帶索引的堆積，說明 decrease-key 和刪除任意元素為什麼需要位置陣列、刪除時為什麼
上浮和下沉都要試；知道相同的鍵要靠加入順序來分先後。
\item 寫出 $d$ 元堆積的位置公式和高度，說明 $d$ 變大時哪些操作變快、哪些變慢。
\item 在最小最大堆積、deap、對稱最小最大堆積上做插入、取出最小值和最大值，每一步都畫得出
來；寫出 deap 的對應點公式，說出 SMMH 的 P1--P3，知道區間堆積和 SMMH 的關係。
\item 合併兩棵左傾樹並更新 $s$；在二項堆積上做插入（含進位）、合併和取出最小值，每一步都
畫得出來。
\item 說出費氏堆積各操作的平攤成本，用位能 $t(H) + 2m(H)$ 說明插入、取出最小值、減少鍵
值的平攤成本從哪裡來，並分清楚平攤和最壞。
\end{itemize}

\textbf{延伸閱讀}（書末會整理完整的參考資料）

\begin{itemize}
\item Thomas H. Cormen 等，\textit{Introduction to Algorithms}，第 4 版，第 6 章：二元
堆積、建堆積的線性時間分析、堆積排序和優先佇列，章末問題有 $d$ 元堆積。 費氏堆積在第
4 版只在第五部的導言提到結論，完整的內容和分析在第 3 版第 19 章；第 2 版第 19 章是二項
堆積。
\item Ellis Horowitz、Sartaj Sahni、Susan Anderson-Freed，\textit{Fundamentals of Data
Structures in C}，第 9 章：最小最大堆積、左傾樹、二項堆積、費氏堆積都有 C 程式，各版的
節次稍有不同；二項堆積用的是延遲合併的版本。
\item Sartaj Sahni，``Double-Ended Priority Queues''，收在 Dinesh P. Mehta、Sartaj Sahni 編，
\textit{Handbook of Data Structures and Applications}（Chapman \& Hall/CRC，2005）第 8 章：
對稱最小最大堆積、區間堆積、最小最大堆積和 deap 放在一起比較。 deap 最早見於 Svante
Carlsson，``The Deap---A Double-Ended Heap to Implement Double-Ended Priority Queues''，
\textit{Information Processing Letters} 26 (1987)；SMMH 見於 A. Arvind、C. Pandu Rangan，
``Symmetric Min-Max Heap: A Simpler Data Structure for Double-Ended Priority Queue''，
\textit{Information Processing Letters} 69 (1999)。
\item Mark Allen Weiss，\textit{Data Structures and Algorithm Analysis in C}，第 6 章：
$d$ 元堆積、左傾堆積（用 npl 的定義）、斜堆積和二項佇列，取向偏實作。
\item 《C 程式設計基礎》第~19~章：用指標和 \texttt{malloc} 建二元樹的基本功，寫左傾樹前
可以先複習。
\end{itemize}

\addcontentsline{toc}{section}{觀念題解答}
\section*{觀念題解答}

\begin{enumerate}
\item (a)、(c)、(e)。 (a)：$90 \ge 70, 80$；$70 \ge 20, 60$；$80 \ge 75, 10$。 (b)：$A[6] =
85$ 比父節點 $A[3] = 70$ 大。 (c)：相等是允許的。 (d)：$A[4] = 40$ 比父節點 $A[2] = 30$
大。 (e)：$i > \lfloor i/2 \rfloor$，排在後面的不比前面的大。

\item 每列是操作之後的陣列，括號是交換次數。
\begin{center}
\begin{tabular}{lcl}
\toprule
操作 & 交換 & 陣列 \\
\midrule
插入 24 & 0 & 24 \\
插入 11 & 1 & 11 24 \\
插入 39 & 0 & 11 24 39 \\
插入 7 & 2 & 7 11 39 24 \\
插入 52 & 0 & 7 11 39 24 52 \\
插入 30 & 1 & 7 11 30 24 52 39 \\
插入 18 & 1 & 7 11 18 24 52 39 30 \\
插入 45 & 0 & 7 11 18 24 52 39 30 45 \\
取出 7 & 2 & 11 24 18 45 52 39 30 \\
取出 11 & 1 & 18 24 30 45 52 39 \\
\bottomrule
\end{tabular}
\end{center}
取出 $7$：$45$ 搬到根，和較小的子節點 $11$ 交換，再和 $24$ 交換。 取出 $11$：$30$ 搬到
根，和 $18$ 交換，到了位置 $3$，子節點 $39$ 比它大，停。

\item
\begin{enumerate}
\item $n = 9$，$i$ 從 $4$ 做到 $1$。
\begin{center}
\begin{tabular}{lcl}
\toprule
& 交換 & 陣列 \\
\midrule
$i = 4$ & 1 & 33 18 56 48 41 27 62 15 9 \\
$i = 3$ & 1 & 33 18 62 48 41 27 56 15 9 \\
$i = 2$ & 1 & 33 48 62 18 41 27 56 15 9 \\
$i = 1$ & 2 & 62 48 56 18 41 27 33 15 9 \\
\bottomrule
\end{tabular}
\end{center}
共交換 $5$ 次。 $i = 2$ 時 $18$ 和較大的子節點 $48$ 交換後，到了位置 $4$，子節點 $15$、
$9$ 都比它小，停。
\item 依序插入的結果是 $\langle 62, 48, 56, 41, 18, 27, 33, 9, 15 \rangle$，和 (a) 不同。
\end{enumerate}

\item
\begin{enumerate}
\item 深度 $d$ 的節點有 $d$ 個祖先，由堆積性質，它們都比這個節點大（元素互不相同）。
如果這個節點是第 $k$ 大，比它大的只有 $k - 1$ 個，所以 $d \le k - 1$。
\item 第 $2$ 大的深度不超過 $1$，又不是根（根是第 $1$ 大），所以在深度 $1$，也就是位置
$2$ 或 $3$。
\item 如果最小的元素有子節點，子節點比它小，矛盾，所以它是葉節點，在位置
$\lfloor n/2 \rfloor + 1$ 到 $n$ 之間。 每個葉節點都有可能：由大到小排好的陣列是最大堆積，
最小值在位置 $n$；把它和任何一個葉節點位置 $\ell$ 的數對調，位置 $\ell$ 沒有子節點，位置
$n$ 的父節點 $\lfloor n/2 \rfloor < \ell$ 比換過去的數大，仍是合法的堆積。
\end{enumerate}

\item 不變量：第 3 行的迴圈每一輪開始時，位置 $i + 1, \ldots, n$ 的每個節點都是一個最大
堆積的根。

初始：$i = \lfloor n/2 \rfloor$，位置 $\lfloor n/2 \rfloor + 1$ 到 $n$ 都是葉節點，一個節點
是最大堆積。

維持：$i$ 的子節點 $2i$、$2i + 1$ 都大於 $i$，由不變量它們都是最大堆積的根（不存在就不
用管），所以 \proc{Sift-Down}$(A, i)$ 之後以 $i$ 為根的子樹是最大堆積，裡面每個節點也都是
最大堆積的根。 \proc{Sift-Down} 只改動這棵子樹裡的位置，其他位置的子樹沒有被動到。 所以位
置 $i, i + 1, \ldots, n$ 都是最大堆積的根；下一輪 $i$ 減 $1$，不變量成立。

結束：迴圈在 $i = 0$ 停下，位置 $1, \ldots, n$ 都是最大堆積的根，特別是位置 $1$，也就是
整個陣列是最大堆積。

\item
\begin{enumerate}
\item $\lfloor \lg i \rfloor$ 在 $i = 1$ 到 $10$ 是 $0, 1, 1, 2, 2, 2, 2, 3, 3, 3$，和是
$19$。
\item 上界：每一項不超過 $\lg n$，$S(n) \le n \lg n$。 下界：$i \ge n/2$ 的項有至少
$n/2$ 項，每一項 $\lfloor \lg i \rfloor \ge \lg(n/2) - 1 = \lg n - 2$，所以
$S(n) \ge \frac{n}{2}(\lg n - 2)$，$n \ge 16$ 時至少 $\frac{n}{4}\lg n$。 所以
$S(n) = \Theta(n \lg n)$。
\end{enumerate}

\item (a)、(c)、(d)。 (b)：最大值在某個葉節點，要看 $\lceil n/2 \rceil$ 個葉節點，
$\Theta(n)$。 (c)：堆積只保證父子的大小，不能像二元搜尋那樣排除一半。 (d)：把兩個陣列接
起來，用 \proc{Build-Max-Heap}（或最小堆積版）建一次，$O(n)$。 (e)：同一層之間沒有任何
規定。

\item
\begin{enumerate}
\item 父節點 $\lfloor (100 - 2)/4 \rfloor + 1 = 25$。 位置 $25$ 的子節點是
$4 \cdot 24 + 2 = 98$ 到 $4 \cdot 25 + 1 = 101$，也就是 $98, 99, 100, 101$。
\item $h = \lceil \log_4(3 \cdot 1000 + 1) \rceil = \lceil \log_4 3001 \rceil = 6$
（$4^5 = 1024 < 3001 \le 4096 = 4^6$）。 驗算：高度 $5$ 最多 $(4^5 - 1)/3 = 341$ 個節點，
高度 $6$ 最多 $(4^6 - 1)/3 = 1365$ 個，$341 < 1000 \le 1365$。
\item 每往下一層，在最多 $d$ 個子節點裡找最小的要比 $d - 1$ 次，再和自己比 $1$ 次，共
$d$ 次。 從根（深度 $0$）往下最多走到深度 $h - 1$，也就是 $h - 1$ 層，所以最多
$d(h - 1)$ 次。
\item $n \approx 10^3$。 上浮和下沉最多走「高度減 $1$」層。 $d = 2$：高度
$\lfloor \lg 1000 \rfloor + 1 = 10$，最多走 $9$ 層，decrease-key 共
$10^6 \times 9 = 9 \times 10^6$，取出共 $10^3 \times 2 \times 9 \approx 1.8 \times 10^4$。
$d = 32$：高度 $\lceil \log_{32} 31001 \rceil = 3$（$32^2 = 1024 < 31001 \le 32^3$），最多
走 $2$ 層，decrease-key 共 $2 \times 10^6$，取出共
$10^3 \times 32 \times 2 = 6.4 \times 10^4$。 總數 $d = 32$ 大約是 $2.1 \times 10^6$，
比 $d = 2$ 的大約 $9 \times 10^6$ 少很多，選 $d = 32$。 decrease-key 遠多於取出時，$d$ 要取大。
\end{enumerate}

\item
\begin{enumerate}
\item 深度 $0$（最小層）：$10$；深度 $1$（最大層）：$65, 70$；深度 $2$（最小層）：$20,
40, 30, 60$；深度 $3$（最大層）：$25, 45, 50$。 檢查：$10$ 不大於所有元素；$65$ 不小於
$20, 40, 25, 45, 50$；$70$ 不小於 $30, 60$；$20$ 不大於 $25, 45$；$40$ 不大於 $50$。
\item $15$ 放在位置 $11$（最大層），父節點 $A[5] = 40$ 在最小層，$15 < 40$，$40$ 挪到位置
$11$；從位置 $5$ 沿最小層往上，祖父 $A[1] = 10$ 不比 $15$ 大，停。 陣列
$\langle 10, 65, 70, 20, 15, 30, 60, 25, 45, 50, 40 \rangle$。
\item 最大值是 $A[2]$、$A[3]$ 中較大的 $70$。 拿出最後一個 $x = 40$，從位置 $3$ 開始：子
節點 $30, 60$（沒有孫節點）中最大的是 $A[7] = 60$，是子節點，$60$ 搬到位置 $3$，$40$ 放
進位置 $7$。 陣列 $\langle 10, 65, 60, 20, 15, 30, 40, 25, 45, 50 \rangle$。
\item 取出 $10$，拿出 $x = 50$。 根的子節點 $65, 60$、孫節點 $20, 15, 30, 40$ 中最小的是
$A[5] = 15$（孫節點），搬到根；它的父節點 $A[2] = 65$ 不比 $50$ 小，不用對調。 位置 $5$
沒有子節點了（$n = 9$），$50$ 放進位置 $5$。 陣列
$\langle 15, 65, 60, 20, 50, 30, 40, 25, 45 \rangle$。
\end{enumerate}

\item
\begin{enumerate}
\item 位置 $2$ 到 $10$ 依序是 $13,\allowbreak 68,\allowbreak 27,\allowbreak 16,\allowbreak 61,\allowbreak 44,\allowbreak 40,\allowbreak 28,\allowbreak 34$；最小堆積的根 $13$ 有子節點 $27$、$16$，$27$ 的子節點是 $40$、$28$，$16$ 的子節點是 $34$；
最大堆積的根 $68$ 有子節點 $61$、$44$。 位置 $4$、$5$ 在深度 $2$，對應點是 $4 + 2 = 6$、
$5 + 2 = 7$；位置 $8, 9, 10$ 在深度 $3$，對應點 $12, 13, 14$ 都超過 $10$，改用父節點 $6, 6, 7$。
檢查：$27 \le 61$、$16 \le 44$、$40 \le 61$、$28 \le 61$、$34 \le 44$，兩邊各自也是堆積，所以是
deap。
\item 插入 $94$：放在位置 $11$（最小堆積），對應點 $15$ 改用 $7$；$94 > A[7] = 44$，$44$ 搬到位
置 $11$，$94$ 從位置 $7$ 上浮，比 $A[3] = 68$ 大，$68$ 挪到位置 $7$，$94$ 到位置 $3$：
$\langle 13, 94, 27, 16, 61, 68, 40, 28, 34, 44 \rangle$。 插入 $21$：放在位置 $12$（最大堆積），
對應點 $12 - 4 = 8$ 沒有子節點；$21 < A[8] = 40$，$40$ 搬到位置 $12$，$21$ 從位置 $8$ 上浮，
比 $A[4] = 27$ 小，$27$ 挪到位置 $8$，再和 $A[2] = 13$ 比，停在位置 $4$：
$\langle 13, 94, 21, 16, 61, 68, 27, 28, 34, 44, 40 \rangle$。
\item 取出 $94$，最後一格的 $40$ 拿出來。 洞從位置 $3$ 往下，子節點 $61, 68$ 中較大的 $68$ 上
來，洞到位置 $7$，沒有子節點（$14 > 11$）。 對應點 $7 - 2 = 5$ 有子節點 $34, 44$，取較大的
$44$（位置 $11$）；$40 < 44$，$44$ 搬到位置 $7$，$40$ 從位置 $11$ 上浮，父節點 $A[5] = 16$ 比它
小，停：$\langle 13, 68, 21, 16, 61, 44, 27, 28, 34, 40 \rangle$。 只和 $A[5] = 16$ 比的話，
$40$ 會留在位置 $7$，得到 $\langle 13, 68, 21, 16, 61, 40, 27, 28, 34, 44 \rangle$；位置 $11$ 的
$44$ 對應點 $15$ 超出範圍，改用 $7$，$44 > 40$，違反定義的第 2 條。
\item 取出 $13$，最後一格的 $40$ 拿出來。 洞從位置 $2$ 往下：$21, 16$ 中 $16$ 上來，洞到位置
$5$；它只剩子節點 $34$，$34$ 上來，洞到位置 $10$。 對應點 $14$ 改用 $7$；$40 \le A[7] = 44$，留
在最小堆積，父節點 $34$ 比它小，停：$\langle 16, 68, 21, 34, 61, 44, 27, 28, 40 \rangle$。
\end{enumerate}

\item
\begin{enumerate}
\item (A)、(C) 是。 (B) 位置 $6$、$7$ 是兄弟，$30 > 20$，違反 P1。 (D) 位置 $10$ 的祖父是 $2$，
$45 > A[5] = 40$，違反 P3（雖然 $45 \le A[3] = 50$）。
\item 每列是插入後的 $A[2..]$。
\begin{center}
\begin{tabular}{ll}
\toprule
插入 & 陣列 \\
\midrule
47 & 47 \\
18 & 18 47 \\
83 & 18 83 47 \\
60 & 18 83 47 60 \\
25 & 18 83 47 60 25 \\
91 & 18 91 47 60 25 83 \\
12 & 12 91 18 60 25 83 47 \\
70 & 12 91 18 70 25 83 47 60 \\
35 & 12 91 18 70 25 83 47 60 35 \\
\bottomrule
\end{tabular}
\end{center}
插入 $18$ 時位置 $3$ 是右子節點，比左兄弟 $47$ 小，對調。 插入 $83$、$91$ 時都比祖父（根）的右
子節點大，換上去。 插入 $12$：位置 $8$ 的祖父是 $2$，$12 < A[4] = 47$，$47$ 下來；再和根的左
子節點 $18$ 比，$18$ 下來，$12$ 到位置 $2$。 插入 $70$：位置 $9$ 比左兄弟 $47$ 大；祖父 $2$ 的右
子節點 $A[5] = 60 < 70$，$60$ 下來；再看根的兩個子節點 $12, 91$，停在位置 $5$。
\item 取出 $91$：最後的 $35$ 拿出來，洞在位置 $3$，候選 $A[5] = 70$、$A[7] = 83$，$83$ 上來，洞到
位置 $7$，左兄弟 $A[6] = 25$ 不比 $35$ 大，沒有子節點，$35$ 放進去：
$\langle 12, 83, 18, 70, 25, 35, 47, 60 \rangle$。 取出 $12$：最後的 $60$ 拿出來，洞在位置 $2$，
候選 $A[4] = 18$、$A[6] = 25$，$18$ 上來，洞到位置 $4$，右兄弟 $A[5] = 70$ 不比 $60$ 小；候選只
剩 $A[8] = 47$，$47 < 60$，上來，洞到位置 $8$，沒有兄弟也沒有子節點，$60$ 放進去：
$\langle 18, 83, 47, 70, 25, 35, 60 \rangle$。
\item 兄弟配對：$(2, 3)$、$(4, 5)$、$(6, 7)$、$(8, 9)$、$(10)$，得到根 $[12, 91]$，它的子節點
$[18, 70]$、$[25, 83]$，$[18, 70]$ 的子節點 $[47, 60]$ 和 $[35]$。 檢查：$[18, 70]$、$[25, 83]$
都在 $[12, 91]$ 裡；$[47, 60]$ 和 $35$ 都在 $[18, 70]$ 裡。
\end{enumerate}

\item 右邊路徑 $A$：$23, 61$；$B$：$24, 39$。 合成 $23 \to 24 \to 39 \to 61$。 由下往上：
$61$ 的 $s = 1$；$39$ 左邊 $57$（$s = 1$）、右邊 $61$（$s = 1$），不對調，$s = 2$；$24$
左邊 $27$（$s = 1$）小於右邊 $39$（$s = 2$），對調，$s = 2$；$23$ 左邊 $26$（$s = 1$）
小於右邊 $24$（$s = 2$），對調，$s = 2$。 結果用「根(左, 右)」寫成
\[
23\bigl(24(39(57, 61),\ 27(38, -)),\ 26(32(62, 60), -)\bigr),
\]
各節點的 $s$：$23{:}2$、$24{:}2$、$39{:}2$、$27{:}1$、$26{:}1$、$32{:}2$，葉節點都是 $1$。

取出 $23$：合併 $24(\ldots)$ 和 $26(\ldots)$。 右邊路徑 $24, 27$ 和 $26$ 合成
$24 \to 26 \to 27$。 $27$ 的 $s = 1$；$26$ 左邊 $32$（$s = 2$）不小於右邊 $27$（$s = 1$），
$s(26) = 2$；$24$ 左邊 $39$（$s = 2$）不小於右邊 $26$（$s = 2$），$s(24) = 3$。 結果
\[
24\bigl(39(57, 61),\ 26(32(62, 60),\ 27(38, -))\bigr).
\]

\item 對節點 $x$，令 $r(x)$ 是從 $x$ 一路往右走到 \NIL 經過的節點數，證明
$r(x) \le \lg(w(x) + 1)$。 對子樹大小歸納。 $x = \NIL$ 時兩邊都是 $0$。 否則
$w(x) = 1 + w(\mathit{left}) + w(\mathit{right})$，而 $w(\mathit{left}) \ge
w(\mathit{right})$，所以 $w(\mathit{right}) \le (w(x) - 1)/2$。 由歸納假設，
\[
r(x) = 1 + r(\mathit{right}) \le 1 + \lg\bigl(w(\mathit{right}) + 1\bigr)
\le 1 + \lg\frac{w(x) + 1}{2} = \lg\bigl(w(x) + 1\bigr).
\]
取 $x$ 是根，$w = n$，得證。 和定理~\ref{thm:hp-leftist} 一樣，合併只沿右邊路徑走，所以
重量左傾樹的合併也是 $O(\lg n)$。

\item
\begin{enumerate}
\item $13 = 1101_2$，三棵：$B_3$、$B_2$、$B_0$。
\item $\binom{5}{2} = 10$ 個。
\item 最後 $n = 10 = 1010_2$，兩棵：$B_1{:}\,17[40]$ 和
$B_3{:}\,8[14[29[52],\,36],\ 23[47],\ 61]$。 連結次數是 $10 - b(10) = 10 - 2 = 8$（逐次是
$0, 1, 0, 2, 0, 1, 0, 3, 0, 1$）。
\item 取出 $8$，子樹 $14[29[52], 36]$（$B_2$）、$23[47]$（$B_1$）、$61$（$B_0$）和剩下
的 $17[40]$（$B_1$）合併。 階數 $0$ 只有 $61$；階數 $1$ 有 $17[40]$、$23[47]$，連成
$17[23[47],\,40]$ 進位；階數 $2$ 有 $14[\ldots]$ 和進位，連成
$14[17[23[47],\,40],\ 29[52],\ 36]$。 結果 $n = 9 = 1001_2$：$B_0{:}\,61$ 和這棵 $B_3$。
\end{enumerate}

\item
\begin{enumerate}
\item (B)、(C)、(D)。 (A) 錯：$O(\lg n)$ 是平攤，(B) 就是一次 $\Theta(n)$ 的例子。 (C)：
連鎖剪下可能剪很多次。 (D)：棵數和標記數都不是負的。
\item 二元堆積：$O((n + m) \lg n)$。 費氏堆積：$n$ 次 \proc{Insert} 和 $m$ 次
\proc{Decrease-Key} 各 $O(1)$，$n$ 次 \proc{Extract-Min} 各 $O(\lg n)$，共
$O(m + n \lg n)$。 $m = \Theta(n^2)$：$O(n^2 \lg n)$ 對 $O(n^2)$，費氏堆積快一個
$\lg n$。 $m = \Theta(n)$：兩者都是 $O(n \lg n)$。
\end{enumerate}
\end{enumerate}

\addcontentsline{toc}{section}{歷屆試題詳解}
\section*{歷屆試題詳解}

\pastsol{px:hp-1}{中正資工所甲組 114 學年度 第 3 題}
答案是 (a)。 由下而上建堆積照虛擬碼~\ref{alg:hp-build}：$n = 8$，$i$ 從 $\lfloor 8/2 \rfloor = 4$
做到 $1$，每次對位置 $i$ 做 \proc{Sift-Down}。

\begin{center}
\begin{tabular}{lc*{8}{c}}
\toprule
& 交換 & 1 & 2 & 3 & 4 & 5 & 6 & 7 & 8 \\
\midrule
開始 & & 52 & 21 & 40 & 33 & 88 & 57 & 46 & 71 \\
$i = 4$ & 1 & 52 & 21 & 40 & 71 & 88 & 57 & 46 & 33 \\
$i = 3$ & 1 & 52 & 21 & 57 & 71 & 88 & 40 & 46 & 33 \\
$i = 2$ & 1 & 52 & 88 & 57 & 71 & 21 & 40 & 46 & 33 \\
$i = 1$ & 2 & 88 & 71 & 57 & 52 & 21 & 40 & 46 & 33 \\
\bottomrule
\end{tabular}
\end{center}

$i = 4$：$33$ 只有左子節點 $A[8] = 71$，交換。 $i = 3$：$40$ 和較大的子節點 $57$ 交換。 $i = 2$：
$21$ 和較大的子節點 $88$ 交換，到了位置 $5$，沒有子節點。 $i = 1$：$52$ 先和 $88$ 交換，到了位
置 $2$，子節點 $71, 21$ 中 $71$ 比它大，再交換，到了位置 $4$，子節點 $33$ 比它小，停。 結果
$\{88, 71, 57, 52, 21, 40, 46, 33\}$，就是 (a)。

其他選項：(b) 是由大到小排好的陣列，它也是最大堆積，但不是這個方法的結果。 (c) 是由小到大
排好，根是最小值，不是最大堆積。 (d) 也是合法的最大堆積（逐層檢查父子都成立），但位置 $5$
和 $8$ 的 $33$、$21$ 和由下而上的結果對調了。 (a) 正確，所以 (e) 不對。 題目問的是特定方法
的結果，「是不是合法的最大堆積」不夠，要真的做一次。

\pastsol{px:hp-2}{臺科大資工所 110 學年度 第 6 題（節錄 (d)(e) 小題）}
(d) 每次插入放在最後一格再上浮（第~\ref{sec:hp-trace}~節，最小堆積版）。 每列是插入之後的陣
列，位置從 $1$ 開始。

\begin{center}
\begin{tabular}{lc*{9}{c}}
\toprule
插入 & 交換 & 1 & 2 & 3 & 4 & 5 & 6 & 7 & 8 & 9 \\
\midrule
12 & 0 & 12 \\
9 & 1 & 9 & 12 \\
14 & 0 & 9 & 12 & 14 \\
16 & 0 & 9 & 12 & 14 & 16 \\
4 & 2 & 4 & 9 & 14 & 16 & 12 \\
17 & 0 & 4 & 9 & 14 & 16 & 12 & 17 \\
6 & 1 & 4 & 9 & 6 & 16 & 12 & 17 & 14 \\
8 & 2 & 4 & 8 & 6 & 9 & 12 & 17 & 14 & 16 \\
2 & 3 & 2 & 4 & 6 & 8 & 12 & 17 & 14 & 16 & 9 \\
\bottomrule
\end{tabular}
\end{center}

插入 $4$：放在位置 $5$，父節點 $A[2] = 12$ 比它大，交換；再和根 $9$ 交換。 插入 $8$：放在位置 $8$，
和父節點 $16$、再和 $9$ 交換，父節點 $4$ 比它小，停。 插入 $2$：位置 $9$ 往上的父節點依序是
$A[4] = 9$、$A[2] = 8$、$A[1] = 4$，都比它大，交換三次，$2$ 到根。 最後的樹是下圖 (a)。

(e) 取出 $2$：最後一格的 $9$ 搬到根，和較小的子節點 $4$ 交換，到了位置 $2$，子節點 $8, 12$ 中
$8$ 比它小，交換，到了位置 $4$，子節點只剩 $16$，停：
$\langle 4, 8, 6, 9, 12, 17, 14, 16 \rangle$。 取出 $4$：最後一格的 $16$ 搬到根，和較小的子節點
$6$ 交換，到了位置 $3$，子節點 $17, 14$ 中 $14$ 比它小，交換，到了位置 $7$，停：
$\langle 6, 8, 14, 9, 12, 17, 16 \rangle$，是下圖 (b)。

\begin{center}
\begin{tikzpicture}
\hpheap{0,0}{6.0}{2/n,4/n,6/n,8/n,12/n,17/n,14/n,16/n,9/n}{xa}
\hpheap{7.4,0}{4.6}{6/n,8/n,14/a,9/n,12/n,17/n,16/a}{xb}
\node[font=\small] at (0,-3.65) {(a) (d) 的答案};
\node[font=\small] at (7.4,-3.65) {(b) (e) 的答案};
\end{tikzpicture}
\end{center}

\pastsol{px:hp-3}{成大資訊聯招 113 學年度 第 2 題}
題目說 ``in the given order''，三小題都解讀成依序一個一個插入。 層數照題目的慣例（也是本書的
約定）從根算第 $1$ 層。 答案：i) (D)；ii) (E)；iii) (A)。

i) 依序插入最大堆積，每次之後的陣列：
\begin{center}
\small
\begin{tabular}{l*{9}{c}}
\toprule
插入 & 1 & 2 & 3 & 4 & 5 & 6 & 7 & 8 & 9 \\
\midrule
20 & 20 \\
5 & 20 & 5 \\
10 & 20 & 5 & 10 \\
18 & 20 & 18 & 10 & 5 \\
4 & 20 & 18 & 10 & 5 & 4 \\
22 & 22 & 18 & 20 & 5 & 4 & 10 \\
11 & 22 & 18 & 20 & 5 & 4 & 10 & 11 \\
32 & 32 & 22 & 20 & 18 & 4 & 10 & 11 & 5 \\
21 & 32 & 22 & 20 & 21 & 4 & 10 & 11 & 5 & 18 \\
\bottomrule
\end{tabular}
\end{center}
插入 $32$ 時它從位置 $8$ 一路換到根，把 $5$ 帶到位置 $8$。 位置 $8$ 的深度是
$\lfloor \lg 8 \rfloor = 3$，在第 $4$ 層，選 (D)。 （如果題目的層數從 $0$ 算，答案會變成 (C)；成大
的考卷沿用從 $1$ 算的定義。）

ii) 依序插入最小堆積：
\begin{center}
\small
\begin{tabular}{l*{9}{c}}
\toprule
插入 & 1 & 2 & 3 & 4 & 5 & 6 & 7 & 8 & 9 \\
\midrule
20 & 20 \\
5 & 5 & 20 \\
10 & 5 & 20 & 10 \\
18 & 5 & 18 & 10 & 20 \\
4 & 4 & 5 & 10 & 20 & 18 \\
22 & 4 & 5 & 10 & 20 & 18 & 22 \\
11 & 4 & 5 & 10 & 20 & 18 & 22 & 11 \\
32 & 4 & 5 & 10 & 20 & 18 & 22 & 11 & 32 \\
21 & 4 & 5 & 10 & 20 & 18 & 22 & 11 & 32 & 21 \\
\bottomrule
\end{tabular}
\end{center}
結果是 $\langle 4,\allowbreak 5,\allowbreak 10,\allowbreak 20,\allowbreak 18,\allowbreak 22,\allowbreak 11,\allowbreak 32,\allowbreak 21 \rangle$。 根的左子樹是位置 $2,\allowbreak 4,\allowbreak 5,\allowbreak 8,\allowbreak 9$，也就是 $5,\allowbreak 20,\allowbreak 18,\allowbreak 32,\allowbreak 21$；右子樹是位
置 $3, 6, 7$，也就是 $10, 22, 11$。 (A) $11$ 在右子樹，錯。 (B) $18$ 在左子樹，錯。 (C) 根是
$4$，錯。 (D) $21$ 在位置 $9$，第 $4$ 層，錯。 所以選 (E)。

iii) 照第~\ref{sec:hp-smmh}~節的 \proc{SMMH-Insert}，下表是每次之後的 $A[2..]$：
\begin{center}
\small
\begin{tabular}{l*{9}{c}}
\toprule
插入 & 2 & 3 & 4 & 5 & 6 & 7 & 8 & 9 & 10 \\
\midrule
20 & 20 \\
5 & 5 & 20 \\
10 & 5 & 20 & 10 \\
18 & 5 & 20 & 10 & 18 \\
4 & 4 & 20 & 10 & 18 & 5 \\
22 & 4 & 22 & 10 & 18 & 5 & 20 \\
11 & 4 & 22 & 10 & 18 & 5 & 20 & 11 \\
32 & 4 & 32 & 10 & 22 & 5 & 20 & 11 & 18 \\
21 & 4 & 32 & 10 & 22 & 5 & 20 & 11 & 18 & 21 \\
\bottomrule
\end{tabular}
\end{center}
插入 $5$：位置 $3$ 是右子節點，比左兄弟 $20$ 小，對調。 插入 $4$：位置 $6$ 的祖父是根，
$4 < A[2] = 5$，$5$ 下來。 插入 $22$：比 $A[3] = 20$ 大，$20$ 下來。 插入 $32$：位置 $9$ 的祖父是
$2$，$32 > A[5] = 18$，$18$ 下來；位置 $5$ 的祖父是根，$32 > A[3] = 22$，$22$ 下來，$32$ 到位置
$3$。 插入 $21$：位置 $10$ 的祖父是 $2$，$A[4] = 10 \le 21 \le A[5] = 22$，不動。 最後的樹如下
圖 (c)：空的根，第 $2$ 層 $4, 32$，第 $3$ 層 $10, 22, 5, 20$，第 $4$ 層 $11, 18, 21$。

(B) 第 $4$ 層有三個節點，錯。 (C) 根的孫節點由左到右是 $10, 22, 5, 20$，最左邊是 $10$，錯。
(D) 根的右子樹是位置 $3, 6, 7$，鍵是 $32, 5, 20$，不是 $32, 10, 20$，錯。 (A)：本書只定義了
取出最小值和最大值，題目沒有說刪除任意節點怎麼做；照取出最小值的做法類推，刪掉位置 $6$
的 $5$ 之後，拿最後一個節點的 $21$ 來補這個洞（洞沒有子節點；$21$ 比右兄弟 $A[7] = 20$ 大，
兩者再對調，最後是 $A[6] = 20$、$A[7] = 21$）。 「用 $21$ 那個節點來補」這件事成立，所以 (A)
對，選 (A)。 這也是唯一可能的答案：(B)、(C)、(D) 都確定是錯的。

\begin{center}
\begin{tikzpicture}
\hpheap{0,0}{5.6}{32/n,22/n,20/n,21/n,4/n,10/n,11/n,5/a,18/n}{ya}
\hpheap{6.6,0}{5.6}{4/n,5/n,10/n,20/n,18/n,22/n,11/n,32/n,21/n}{yb}
\node[font=\small] at (0,-3.65) {(a) 最大堆積，$5$ 在第 $4$ 層};
\node[font=\small] at (6.6,-3.65) {(b) 最小堆積};
\end{tikzpicture}

\begin{tikzpicture}
\hpheap{0,0}{6.0}{/e,4/n,32/n,10/n,22/n,5/a,20/n,11/n,18/n,21/n}{yc}
\node[font=\small] at (0,-3.65) {(c) 對稱最小最大堆積};
\end{tikzpicture}
\end{center}

\pastsol{px:hp-4}{臺大資工所 113 學年度 第 2 題}
答案是 A、D。 先把函式補完。 兩棵樹都不空時，根比較大的那個當新的根（$>$ 不成立時選
\emph{root2}，相等也可以），把它的某一棵子樹和另一整棵樹遞迴合併，結果接回同一邊。 這和
虛擬碼~\ref{alg:hp-meld} 左傾樹的 \proc{Meld} 是同一個架構，只是少了比較 $s$、對調左右那兩步。
一種正確的填法：

\begin{pseudo}[unbreakable]{臺大 113 第 2 題的一種填法}
\Function{MergeMaxBinaryTrees}{$\mathit{root1}, \mathit{root2}$}
    \If{$\mathit{root1} = \NIL$} \Return $\mathit{root2}$ \Comment{(a)} \EndIf
    \If{$\mathit{root2} = \NIL$} \Return $\mathit{root1}$ \Comment{(b)} \EndIf
    \If{$\mathit{root1}.\mathit{data} > \mathit{root2}.\mathit{data}$}
        \State $\mathit{root1}.\mathit{right} \gets$ \Call{MergeMaxBinaryTrees}{$\mathit{root1}.\mathit{right}, \mathit{root2}$} \Comment{(c)(d)(e)}
        \State \Return $\mathit{root1}$ \Comment{(f)}
    \Else
        \State $\mathit{root2}.\mathit{right} \gets$ \Call{MergeMaxBinaryTrees}{$\mathit{root1}, \mathit{root2}.\mathit{right}$} \Comment{(g)(h)(i)}
        \State \Return $\mathit{root2}$ \Comment{(j)}
    \EndIf
\EndFunction
\end{pseudo}

正確性對兩棵樹的節點總數歸納：遞迴呼叫的兩棵樹節點總數少了一個（新的根不在裡面），由歸納假
設回傳一棵 MaxBT，它的節點都來自原本的兩棵樹，所以都不大於新的根；接回去以後新的根仍然不小
於兩棵子樹的所有節點，整棵是 MaxBT。

逐個選項：
\begin{itemize}
\item[A.] 對。 (b) 是 \emph{root2} 為空時回傳 \emph{root1}；(f) 是 \emph{root1} 當新的根時回傳
\emph{root1}。 兩格都必須是 \emph{root1}。
\item[B.] 錯。 (c) 是被改寫的那個欄位，它原本指的子樹必須當作遞迴的一個參數傳進去，否則那棵子
樹就從樹上掉了。 所以 (c) 一定和 (d)、(e) 其中一個相同（例如上面的 (c) 和 (d) 都是
\emph{root1}.right）。
\item[C.] 錯。 MaxBT 不要求形狀，(c) 填 \emph{root1}.left（同時 (d) 或 (e) 改成
\emph{root1}.left）也一樣正確；(g) 同理。 每一格都至少有兩種正確答案。
\item[D.] 對。 上面的正確性論證沒有用到兩個參數的順序：交換 (h)、(i) 之後，遞迴呼叫仍然是合併同
樣兩棵樹，回傳的仍然是一棵包含它們全部節點的 MaxBT。 相等的情況也沒問題，相等時選誰當根都
滿足最大堆積性質。
\item[E.] 錯。 兩個參數是同一棵樹時，資料相等，走 else，把 \emph{root2}.right 改成「這棵樹和它
自己的右子樹合併」的結果，同一批節點被接到兩個地方。 最簡單的例子是只有一個節點 $r$：
\proc{MergeMaxBinaryTrees}$(r, r)$ 走 else，$r.\mathit{right} \gets$
\proc{MergeMaxBinaryTrees}$(r, \NIL) = r$，$r$ 的右子節點變成自己，不再是一棵樹。
\end{itemize}

補充：沒有像左傾樹那樣控制右邊路徑的長度，這個合併最壞要沿著很長的路徑遞迴，時間是
$\Theta(n)$，而左傾樹保證 $O(\lg n)$（定理~\ref{thm:hp-leftist}）。

\pastsol{px:hp-5}{成大電通所 110 學年度 第 2 題（節錄 (b)(c) 小題）}
(b) 照第~\ref{sec:hp-mergeable}~節的定義，二項最小堆積是一群樹，滿足：
\begin{enumerate}
\item 每棵樹都是二項樹 $B_k$（$B_0$ 是一個節點，$B_k$ 是兩棵 $B_{k-1}$ 把一棵的根接成另一棵根
的最左邊的子節點）。
\item 每棵樹都滿足最小堆積性質：每個節點的鍵不小於它的父節點。
\item 每個階數 $k$ 最多只有一棵 $B_k$；所以 $n$ 個元素時有哪幾棵樹由 $n$ 的二進位決定。
\item 各棵樹的根用根串列串起來，依階數由小到大排列（\texttt{Head[H]} 指向階數最小的那棵）。
\end{enumerate}

(c) 圖中的根串列是 $B_0{:}\,14$、$B_2{:}\,1[10[20],\,23]$、$B_3{:}\,6[10[13[30],\,12],\ 7[19],\
15]$，共 $13 = 1101_2$ 個元素。 照第~\ref{sec:hp-mergeable}~節的取出最小值：
\begin{enumerate}
\item 最小的根是 $1$，從根串列拿掉，剩下 $\{B_0{:}\,14,\ B_3{:}\,6[\ldots]\}$。
\item $1$ 的子樹由左到右是 $B_1{:}\,10[20]$、$B_0{:}\,23$，反過來排，成為另一個二項堆積：
$B_0{:}\,23$、$B_1{:}\,10[20]$。
\item 兩個堆積合併。 階數 $0$：$14$、$23$ 兩棵，連成 $14[23]$ 進位。 階數 $1$：$10[20]$ 和進位的
$14[23]$ 兩棵，根 $10 < 14$，$14$ 那棵掛到 $10$ 的最左邊，得到 $B_2{:}\,10[14[23],\,20]$ 進位。
階數 $2$：只有進位的那棵，留下。 階數 $3$：$6[\ldots]$，留下。
\end{enumerate}
結果有 $12 = 1100_2$ 個元素，兩棵樹：

\begin{center}
\begin{tikzpicture}[x=0.9cm,y=-0.9cm,
  v/.style={circle,draw,thick,inner sep=0pt,minimum size=6mm,font=\small},
  a/.style={-{Latex[length=2.2mm]},thick}]
\node[anchor=east] (hd) at (2.9,0.35) {Head[H]};
\node[v] (r10) at (4.6,0.35) {10};
\node[v] (c14) at (3.8,1.35) {14};
\node[v] (c20) at (5.4,1.35) {20};
\node[v] (c23) at (3.8,2.45) {23};
\node[v] (r6) at (9.9,0.35) {6};
\node[v] (b10) at (7.7,1.35) {10};
\node[v] (b7) at (8.8,1.35) {7};
\node[v] (b15) at (9.9,1.35) {15};
\node[v] (b13) at (6.6,2.45) {13};
\node[v] (b12) at (7.7,2.45) {12};
\node[v] (b19) at (8.8,2.45) {19};
\node[v] (b30) at (6.6,3.45) {30};
\draw[a] (hd) -- (r10);
\draw[a] (r10) -- (r6);
\draw[a] (r10) -- (c14);
\draw[a] (r10) -- (c20);
\draw[a] (c14) -- (c23);
\draw[a] (r6) -- (b10);
\draw[a] (r6) -- (b7);
\draw[a] (r6) -- (b15);
\draw[a] (b10) -- (b13);
\draw[a] (b10) -- (b12);
\draw[a] (b7) -- (b19);
\draw[a] (b13) -- (b30);
\end{tikzpicture}
\end{center}

也就是 $B_2{:}\,10[14[23],\,20]$ 和原本的 $B_3$。 圖中兩棵樹各有一個 $10$，鍵可以重複，不影響
做法。

\chapter{二元搜尋樹}
\label{ch:bst}

% 本章的圖共用的樣式（名稱都以 b 開頭，避免和其他章衝突）
\tikzset{
    bnode/.style={circle, draw=black!65, minimum size=7mm, inner sep=0pt, font=\small},
    bnodes/.style={bnode, minimum size=6mm, font=\footnotesize},
    bhl/.style={draw=memStroke, thick, fill=memFill},
    bnew/.style={draw=optStroke, thick, fill=optFill},
    bpath/.style={edge from parent/.style={draw=memStroke, very thick}},
    bst/.style={level distance=10mm, every child node/.style={bnode},
        edge from parent/.style={draw=black!55, thick},
        level 1/.style={sibling distance=48mm},
        level 2/.style={sibling distance=24mm},
        level 3/.style={sibling distance=12mm},
        level 4/.style={sibling distance=7mm}},
    bsts/.style={bst, level distance=9mm, every child node/.style={bnodes},
        level 1/.style={sibling distance=34mm},
        level 2/.style={sibling distance=17mm},
        level 3/.style={sibling distance=8.5mm},
        level 4/.style={sibling distance=6mm}},
    bsize/.style={font=\scriptsize, text=coreStroke, inner sep=1pt},
}

健身房的櫃台用會員編號管理會員，一整天都在做三件事：刷卡時查這個編號是不是有效的
會員，有人入會就加一筆，有人退會就刪一筆。 第~5~章排好的串列，查詢和加入都得從頭
一個一個走；排好的陣列可以二分搜尋，加入一筆卻要把後面的資料全部往後搬。 有沒有一
種存法，讓三件事都只要走大約 $\lg n$ 步？ 提示：回想二分搜尋，每一步其實只看了一個
元素，看完之後就決定往左半邊還是右半邊走。 往下讀之前，先想一想。

把二分搜尋「先看中間、再往左或往右」的結構直接用節點和指標存下來，就是\textbf{二元
搜尋樹}（binary search tree，簡稱 BST）。 每個節點的左邊放比它小的、右邊放比它大
的；查詢從根往下走，每一步往左或往右，插入就是查到空位時接上一個新節點，不用搬動任
何東西。 這一章先講搜尋、最小／最大和前驅／後繼，再講插入和刪除（刪除有三種情況，
每種都畫出來）。 所有操作的成本都和樹高成正比，所以第~\ref{sec:bst-height}~節討論
樹高：最壞會長成一條鏈，隨機的插入順序則讓樹高維持在 $O(\lg n)$ 左右。 最後在每個
節點多記一個「子樹大小」，用它在 $O(h)$ 時間內找第 $k$ 小、算名次。

這一章用到《C 程式設計基礎》第~14~章的結構指標、第~15~章的遞迴、第~19~章的
\texttt{malloc}、指標的指標和 BST 入門（那裡寫過插入和三種走訪，沒有寫刪除），用到
時會先在框裡複習。 樹的名詞和走訪在本書第~7~章，第~5~章的 \texttt{Node **} 寫法這
一章會再用一次。

\section{搜尋、最小／最大、前驅與後繼}
\label{sec:bst-query}

\subsection{定義}

二元搜尋樹是一棵二元樹，每個節點 $x$ 有三個欄位：$x.\mathit{key}$ 是鍵值，
$x.\mathit{left}$、$x.\mathit{right}$ 指向左、右子節點（沒有就是 \NIL）。 樹本身只
記一個指標 $T.\mathit{root}$。 實際的節點還會帶著其他資料（會員姓名、到期日），和
第~5~章一樣，虛擬碼只寫鍵值和指標。 這些節點必須滿足下面的性質。

\begin{definition}[二元搜尋樹性質]
\label{def:bst}
對樹上的每個節點 $x$：$x$ 的左子樹裡\textbf{每一個}節點 $y$ 都有 $y.\mathit{key} <
x.\mathit{key}$，右子樹裡\textbf{每一個}節點 $y$ 都有 $y.\mathit{key} >
x.\mathit{key}$。
\end{definition}

條件是對整棵子樹，只看父子兩個節點不夠。 《C 程式設計基礎》第~19~章有一個反例：
$5$ 的右子節點是 $8$，$8$ 的左子節點是 $4$，每一對父子都沒有排錯，但 $4$ 在 $5$ 的右
子樹裡卻比 $5$ 小。

這一章假設鍵值都不相同，重複的鍵值在插入時忽略，和《C 程式設計基礎》的約定一樣。
真的要存重複的鍵值，常見的做法是在節點裡多記一個「出現次數」，或是把鍵值換成
「鍵值＋序號」這種不會重複的組合（第~\ref{sec:bst-rank}~節的例子會用到）。 有些課
本允許相等的鍵值放在右邊，定義裡的 $>$ 改成 $\ge$；考卷上遇到時，照題目的規定。

\textbf{深度}和\textbf{高度}照第~7~章的定義：深度數邊，根的深度是 $0$，節點的深度是
它到根的路徑上有幾條邊；高度數節點，樹高 $h$ 是樹的層數，也就是最大深度加 $1$。 只有
一個節點的樹高度是 $1$，空樹的高度是 $0$。 CLRS 的高度數邊，算出來會比這裡少 $1$。

下面是本章前兩節一直用的例子：依序加入 $52,\allowbreak 27,\allowbreak 74,\allowbreak 15,\allowbreak 40,\allowbreak 63,\allowbreak 88,\allowbreak 33,\allowbreak 46,\allowbreak 58,\allowbreak 70,\allowbreak 95,\allowbreak 36$ 這 $13$ 位會員（編號只寫末兩碼），得到的樹。 第~\ref{sec:bst-modify}~節會一
步一步插入；這裡先看結果。

\begin{center}
\begin{tikzpicture}[bst]
\node[bnode] {52}
  child {node {27}
    child {node {15}}
    child {node {40}
      child {node {33}
        child[missing]
        child {node {36}}}
      child {node {46}}}}
  child {node {74}
    child {node {63}
      child {node {58}}
      child {node {70}}}
    child {node {88}
      child[missing]
      child {node {95}}}};
\end{tikzpicture}
\diagcaption{$13$ 個會員編號組成的二元搜尋樹 $T$，高度 $5$（最深的是 $36$，深度 $4$）。 根
$52$ 的左子樹都比 $52$ 小，右子樹都比 $52$ 大；每一棵子樹也都是這樣。}
\label{fig:bst-T}
\end{center}

對一棵二元搜尋樹做第~7~章的\textbf{中序走訪}（先左子樹、再自己、再右子樹），會由
小到大印出所有鍵值。 在圖~\ref{fig:bst-T} 上是 $15,\allowbreak 27,\allowbreak 33,\allowbreak 36,\allowbreak 40,\allowbreak 46,\allowbreak 52,\allowbreak 58,\allowbreak 63,\allowbreak 70,\allowbreak 74,\allowbreak 88,\allowbreak 95$。

\begin{theorem}
\label{thm:bst-inorder}
對二元搜尋樹做中序走訪，印出的鍵值由小到大排列。 $n$ 個節點時走訪的時間是
$\Theta(n)$。
\end{theorem}

\begin{proof}
對子樹的節點數做歸納。 空樹什麼都不印，成立。 設根是 $x$，左、右子樹的節點數都比整
棵少，由歸納假設，走訪左子樹印出的是左子樹的鍵值，由小到大；右子樹也是。 由
定義~\ref{def:bst}，左子樹的鍵值都比 $x.\mathit{key}$ 小，右子樹的都比它大，所以
「左子樹的遞增序列、$x.\mathit{key}$、右子樹的遞增序列」接起來仍然遞增。 時間：每
個節點印一次，每個 \NIL 子樹各花常數時間，而 $n$ 個節點的二元樹恰好有 $n + 1$ 個
\NIL 子樹（第~7~章），合計 $\Theta(n)$。
\end{proof}

反過來也成立：一棵二元樹的中序走訪如果是遞增的，它就是二元搜尋樹。 所以「檢查一棵
樹是不是 BST」可以只做一次中序走訪，看每個鍵值是否比前一個大。

\subsection{搜尋}

二元搜尋樹和二分搜尋的關係，看下面這張圖最清楚。 對排好的陣列 $15,\allowbreak 27,\allowbreak 33,\allowbreak 40,\allowbreak 52,\allowbreak 63,\allowbreak 74$ 做二分搜尋，第一步一定看中間的 $40$；比 $40$ 小就到左半邊 $15, 27, 33$，
看它的中間 $27$；比 $40$ 大就看右半邊的中間 $63$。 把「每一步可能看的元素」畫出來，
就是一棵二元搜尋樹。

\begin{center}
\begin{tikzpicture}
\foreach \v [count=\k] in {15, 27, 33, 40, 52, 63, 74} {
    \node[cell] (a\k) at (\k*0.9, 0) {\v};
}
\node[idxlbl, anchor=east] at (0.3, 0) {陣列};
\begin{scope}[xshift=10.6cm, yshift=0.6cm, bsts]
\node[bnodes] {40}
  child {node {27} child {node {15}} child {node {33}}}
  child {node {63} child {node {52}} child {node {74}}};
\end{scope}
\draw[diagarrow] (7.0, 0) -- (8.2, 0);
\end{tikzpicture}
\diagcaption{二分搜尋每一步看的元素排成一棵樹。 陣列可以用位址算出中間那一格；
樹把「中間是誰」直接存成指標，所以插入新元素時不必搬動其他元素。}
\label{fig:bst-binsearch}
\end{center}

搜尋鍵值 $k$ 時，從根開始和 $x.\mathit{key}$ 比較：相等就找到了；$k$ 比較小，答案只
可能在左子樹，往左走；比較大就往右走。 走到 \NIL 表示 $k$ 不在樹裡。

\begin{pseudo}[label=alg:bst-search]{在以 $x$ 為根的子樹裡找 $k$（遞迴）}
\Function{Tree-Search}{$x, k$}
    \If{$x = \NIL$ \OR $k = x.\mathit{key}$}
        \State \Return $x$
    \EndIf
    \If{$k < x.\mathit{key}$}
        \State \Return \Call{Tree-Search}{$x.\mathit{left}, k$}
    \Else
        \State \Return \Call{Tree-Search}{$x.\mathit{right}, k$}
    \EndIf
\EndFunction
\end{pseudo}

呼叫的方式是 \proc{Tree-Search}$(T.\mathit{root}, k)$。 每一層遞迴只往下走一步，而且
遞迴呼叫是函式的最後一個動作，可以直接改成迴圈，省掉呼叫堆疊（第~6~章）：

\begin{pseudo}[label=alg:bst-search-iter]{在以 $x$ 為根的子樹裡找 $k$（迴圈）}
\Function{Iterative-Tree-Search}{$x, k$}
    \While{$x \ne \NIL$ \AND $k \ne x.\mathit{key}$}
        \If{$k < x.\mathit{key}$}
            \State $x \gets x.\mathit{left}$
        \Else
            \State $x \gets x.\mathit{right}$
        \EndIf
    \EndWhile
    \State \Return $x$
\EndFunction
\end{pseudo}

兩個版本走過的節點完全一樣，都在一條從根往下的路徑上，每層最多一個，所以最多比較
$h$ 個節點，時間 $O(h)$，$h$ 是樹高。 下表是在圖~\ref{fig:bst-T} 上搜尋 $46$ 和 $60$ 的過程。

\begin{center}
\begin{tabular}{cccl}
\toprule
搜尋 & 第幾次比較 & $x.\mathit{key}$ & 結果 \\
\midrule
$46$ & 1 & 52 & $46 < 52$，往左 \\
     & 2 & 27 & $46 > 27$，往右 \\
     & 3 & 40 & $46 > 40$，往右 \\
     & 4 & 46 & 相等，找到 \\
\midrule
$60$ & 1 & 52 & $60 > 52$，往右 \\
     & 2 & 74 & $60 < 74$，往左 \\
     & 3 & 63 & $60 < 63$，往左 \\
     & 4 & 58 & $60 > 58$，往右，$x = \NIL$，找不到 \\
\bottomrule
\end{tabular}
\end{center}

找不到 $60$ 時，最後停在 $58$ 的右邊。 這個空位剛好就是插入 $60$ 時新節點要接的地
方，第~\ref{sec:bst-modify}~節的插入就是利用這件事。

\subsection{最小與最大}

最小的鍵值在哪裡？ 從根一直往左走，走到沒有左子節點為止。 理由：如果 $x$ 有左子節
點，左子樹裡的鍵值都比 $x$ 小，最小值一定不是 $x$，也不在 $x$ 的右子樹，只能在左子
樹；如果 $x$ 沒有左子節點，子樹裡其他的鍵值都在右邊，都比 $x$ 大，所以 $x$ 就是最
小值。 最大值左右對調。

\begin{pseudo}[label=alg:bst-minmax]{以 $x$ 為根的子樹裡的最小值、最大值（$x \ne \NIL$）}
\Function{Tree-Minimum}{$x$}
    \While{$x.\mathit{left} \ne \NIL$}
        \State $x \gets x.\mathit{left}$
    \EndWhile
    \State \Return $x$
\EndFunction
\Function{Tree-Maximum}{$x$}
    \While{$x.\mathit{right} \ne \NIL$}
        \State $x \gets x.\mathit{right}$
    \EndWhile
    \State \Return $x$
\EndFunction
\end{pseudo}

圖~\ref{fig:bst-T} 的最小值是 $52 \to 27 \to 15$，最大值是 $52 \to 74 \to 88 \to
95$。 注意最大值 $95$ 剛好是葉節點，但這不是必然的：最小值所在的節點一定沒有左子節
點，卻可能有右子節點。 例如刪掉 $15$ 之後，最小值變成 $27$，它有右子樹。 時間一樣
是 $O(h)$。

\subsection{後繼與前驅}

節點 $x$ 的\textbf{後繼}（successor）是中序走訪裡排在 $x$ 下一個的節點，也就是鍵值
比 $x.\mathit{key}$ 大的節點裡最小的那個；\textbf{前驅}（predecessor）是排在前一個
的節點。 $x$ 是最大值時沒有後繼，回傳 \NIL。 刪除（下一節）和第 $k$ 小的應用都需要
它。

找後繼分兩種情況。

\begin{enumerate}
\item \textbf{$x$ 有右子樹。} 右子樹的鍵值都比 $x$ 大；樹上其他比 $x$ 大的鍵值，如果
在 $x$ 的某個祖先 $y$ 那裡，那是因為 $x$ 在 $y$ 的左子樹裡，這時 $x$ 的右子樹也整棵
在 $y$ 的左子樹裡，都比 $y$ 小。 所以後繼是右子樹的最小值。 圖~\ref{fig:bst-T} 裡
$27$ 的後繼是右子樹 $\{40, 33, 36, 46\}$ 的最小值 $33$。
\item \textbf{$x$ 沒有右子樹。} 後繼一定是 $x$ 的祖先。 從 $x$ 往上走，只要目前的節
點是父節點的右子節點，父節點就比它小，繼續往上；第一次「從左邊走上來」時，那個父
節點就是後繼。 圖~\ref{fig:bst-T} 裡 $46$ 是 $40$ 的右子節點，$40$ 是 $27$ 的右子節
點，$27$ 是 $52$ 的左子節點，所以 $46$ 的後繼是 $52$。 走到根都沒有從左邊上來，表
示 $x$ 是最大值。
\end{enumerate}

\begin{center}
\begin{tikzpicture}[bst]
\node[bnode, bhl] {52}
  child[bpath] {node[bhl] {27}
    child[black!55] {node {15}}
    child[bpath] {node[bhl] {40}
      child {node {33}
        child[missing]
        child {node {36}}}
      child[bpath] {node[bnew] {46}}}}
  child {node {74}
    child {node {63}
      child {node {58}}
      child {node {70}}}
    child {node {88}
      child[missing]
      child {node {95}}}};
\node[font=\scriptsize, text=memStroke, anchor=west] at (0.5, 0) {後繼};
\end{tikzpicture}
\diagcaption{$46$ 沒有右子樹，往上走：$46 \to 40$、$40 \to 27$ 都是從右邊上去；
$27 \to 52$ 是從左邊上去，所以 $46$ 的後繼是 $52$。}
\label{fig:bst-succ}
\end{center}

第二種情況要能從節點往上走，所以節點要多一個欄位 $x.p$ 指向父節點（根的 $p$ 是
\NIL）。 有 $p$ 的寫法如下：

\begin{pseudo}[label=alg:bst-succ]{節點 $x$ 的後繼（節點有父指標 $p$）}
\Function{Tree-Successor}{$x$}
    \If{$x.\mathit{right} \ne \NIL$}
        \State \Return \Call{Tree-Minimum}{$x.\mathit{right}$}
    \EndIf
    \State $y \gets x.p$
    \While{$y \ne \NIL$ \AND $x = y.\mathit{right}$} \Comment{從右邊上來，繼續}
        \State $x \gets y$
        \State $y \gets y.p$
    \EndWhile
    \State \Return $y$
\EndFunction
\end{pseudo}

不想多存父指標，也可以從根往下找。 從根往下走向 $x.\mathit{key}$ 的路上，每次往左
轉，那個節點就比 $x.\mathit{key}$ 大，是後繼的候選；越晚往左轉的候選越小。 所以只要
記住「最後一次往左轉的節點」。 走到 $x$ 之後，如果 $x$ 有右子樹，就照同樣的規則繼
續走進右子樹：右子樹的鍵值都比 $x.\mathit{key}$ 大，先往右一步，接著一路往左，每一
步都更新候選，最後停在右子樹的最小值。 兩種情況因此可以寫成同一個迴圈：

\begin{pseudo}[label=alg:bst-succ-key]{比 $k$ 大的最小鍵值（不用父指標）}
\Function{Successor-Of-Key}{$T, k$}
    \State $s \gets \NIL$ \Comment{目前看到比 $k$ 大的最小節點}
    \State $x \gets T.\mathit{root}$
    \While{$x \ne \NIL$}
        \If{$k < x.\mathit{key}$}
            \State $s \gets x$
            \State $x \gets x.\mathit{left}$
        \Else
            \State $x \gets x.\mathit{right}$
        \EndIf
    \EndWhile
    \State \Return $s$
\EndFunction
\end{pseudo}

這個版本還有一個好處：$k$ 不必在樹裡。 例如 $k = 60$ 時它回傳 $63$，回答的是「編號
$60$ 之後的下一位會員是誰」。 下表是 $k = 46$ 和 $k = 27$ 的過程。

\begin{center}
\begin{tabular}{ccccc}
\toprule
$k$ & $x.\mathit{key}$ & 比較 & 動作 & $s$ \\
\midrule
$46$ & 52 & $46 < 52$ & $s \gets 52$，往左 & 52 \\
     & 27 & $46 \ge 27$ & 往右 & 52 \\
     & 40 & $46 \ge 40$ & 往右 & 52 \\
     & 46 & $46 \ge 46$ & 往右，$x = \NIL$ & 52 \\
\midrule
$27$ & 52 & $27 < 52$ & $s \gets 52$，往左 & 52 \\
     & 27 & $27 \ge 27$ & 往右 & 52 \\
     & 40 & $27 < 40$ & $s \gets 40$，往左 & 40 \\
     & 33 & $27 < 33$ & $s \gets 33$，往左，$x = \NIL$ & 33 \\
\bottomrule
\end{tabular}
\end{center}

兩個版本都只沿著一條路徑往上或往下走，時間 $O(h)$。 前驅把左右對調：有左子樹就是左
子樹的最大值，否則往上找第一個「從右邊走上來」的祖先；不用父指標的版本則改成記住最後
一次往右轉的節點。 圖~\ref{fig:bst-T} 裡 $58$ 沒有左子樹，往上 $58 \to 63 \to 74$
都是從左邊上去，$74 \to 52$ 是從右邊上去，所以 $58$ 的前驅是 $52$。

\begin{examnote}
追蹤搜尋、後繼、前驅時，把走過的節點依序寫成一列，例如「$52 \to 27 \to 40 \to 46$，
找到」，每一步旁邊註明比較的結果。 求後繼時先寫「$x$ 有沒有右子樹」，再寫用哪一種情
況。 題目給的樹如果沒有說節點有父指標，就用從根往下的寫法，不要假設可以往上走。
\end{examnote}

\section{插入、刪除}
\label{sec:bst-modify}

\subsection{插入}

插入新節點 $z$ 的方法就是先搜尋 $z.\mathit{key}$：搜尋失敗時停下的那個 \NIL，就是
$z$ 的位置。 走的時候讓 $y$ 跟在 $x$ 後面一層，$x$ 變成 \NIL 時，$y$ 就是 $z$ 的父節
點，再看 $z.\mathit{key}$ 比 $y.\mathit{key}$ 小還是大，決定接在左邊還是右邊。 這和
第~5~章依序插入串列時讓 $p$ 跟在 $q$ 後面是同一個做法。

\begin{pseudo}[label=alg:bst-insert]{插入節點 $z$（$z.\mathit{left} = z.\mathit{right} = \NIL$，鍵值不在樹裡）}
\Procedure{Tree-Insert}{$T, z$}
    \State $y \gets \NIL$ \Comment{$y$ 跟在 $x$ 後面一層}
    \State $x \gets T.\mathit{root}$
    \While{$x \ne \NIL$}
        \State $y \gets x$
        \If{$z.\mathit{key} < x.\mathit{key}$}
            \State $x \gets x.\mathit{left}$
        \Else
            \State $x \gets x.\mathit{right}$
        \EndIf
    \EndWhile
    \If{$y = \NIL$}
        \State $T.\mathit{root} \gets z$ \Comment{樹原本是空的}
    \ElsIf{$z.\mathit{key} < y.\mathit{key}$}
        \State $y.\mathit{left} \gets z$
    \Else
        \State $y.\mathit{right} \gets z$
    \EndIf
\EndProcedure
\end{pseudo}

新節點永遠接成葉節點，原本的節點一個都不動，所以插入之後樹仍然滿足
定義~\ref{def:bst}：$z$ 走過的每個節點，$z$ 都放在正確的那一邊。 走過的節點最多
$h$ 個（$h$ 是插入前的樹高），時間 $O(h)$。 節點有父指標時，在第 10 行之前加一行 $z.p \gets y$。

第~\ref{sec:bst-query}~節的例子就是對空樹依序插入 $52,\allowbreak 27,\allowbreak 74,\allowbreak 15,\allowbreak 40,\allowbreak 63,\allowbreak 88,\allowbreak 33,\allowbreak 46,\allowbreak 58,\allowbreak 70,\allowbreak 95,\allowbreak 36$。 考卷上常要求寫出每一次插入的比較過程，下表一列寫一個鍵值。

\begin{center}
\begin{tabular}{cll}
\toprule
插入 & 比較過的節點 & 接在 \\
\midrule
52 & （空樹） & 成為根 \\
27 & 52 & 52 的左邊 \\
74 & 52 & 52 的右邊 \\
15 & 52, 27 & 27 的左邊 \\
40 & 52, 27 & 27 的右邊 \\
63 & 52, 74 & 74 的左邊 \\
88 & 52, 74 & 74 的右邊 \\
33 & 52, 27, 40 & 40 的左邊 \\
46 & 52, 27, 40 & 40 的右邊 \\
58 & 52, 74, 63 & 63 的左邊 \\
70 & 52, 74, 63 & 63 的右邊 \\
95 & 52, 74, 88 & 88 的右邊 \\
36 & 52, 27, 40, 33 & 33 的右邊 \\
\bottomrule
\end{tabular}
\end{center}

一個鍵值插入時比較了幾個節點，就等於它的深度。 把這一欄加起來，就是建這棵樹總共比
較的次數：$0 + 1 + 1 + 2 + 2 + 2 + 2 + 3 + 3 + 3 + 3 + 3 + 4 = 29$。 樹的形狀完全由
插入順序決定。 同樣這 $13$ 個編號，如果由小到大插入，每個新節點都接在前一個的右
邊，樹會變成一條鏈，第~\ref{sec:bst-height}~節再回來看這件事。

\subsection{刪除的三種情況}

刪除節點 $z$ 比插入麻煩，因為 $z$ 可能在樹的中間，拿掉它之後兩棵子樹要有人接手。 依
$z$ 有幾個子節點分成三種情況。

\begin{enumerate}
\item \textbf{$z$ 沒有子節點（葉節點）。} 直接拿掉，父節點指向 $z$ 的那個指標改成
\NIL。
\item \textbf{$z$ 只有一個子節點。} 讓這個子節點連同它的子樹頂上來，接在 $z$ 原本的
位置。 這樣做仍然是 BST：那棵子樹原本就在 $z$ 的父節點的同一邊。
\item \textbf{$z$ 有兩個子節點。} 找 $z$ 的後繼 $s$，也就是右子樹的最小值，用 $s$ 取
代 $z$。 $s$ 比左子樹的每個鍵值大（因為它比 $z$ 大），又比右子樹裡其他的鍵值小（因為
它是右子樹的最小值），放在 $z$ 的位置剛好。 $s$ 是右子樹的最小值，所以\textbf{沒有左
子節點}，把它從原位拿掉時屬於情況 1 或 2，不會再遇到兩個子節點。
\end{enumerate}

情況 3 也可以用前驅（左子樹的最大值）取代 $z$，理由左右對調，前驅一定沒有右子節
點。 兩種都是對的，但結果的樹不一樣；本書沒有特別說明時都用後繼。

下面用圖~\ref{fig:bst-T} 依序刪除 $46$、$88$、$27$、$74$，四種狀況各遇到一次。

\begin{enumerate}
\item 刪除 $46$：葉節點，$40$ 的右子節點改成 \NIL。
\item 刪除 $88$：只有右子節點 $95$，$95$ 頂上來成為 $74$ 的右子節點。
\item 刪除 $27$：兩個子節點。 後繼是右子樹 $\{40, 33, 36\}$ 的最小值 $33$，它在右子
樹的深處，而且有右子節點 $36$。 先把 $33$ 從原位拿掉（情況 2：$36$ 頂上來成為
$40$ 的左子節點），再讓 $33$ 坐到 $27$ 的位置，接手 $27$ 的左子樹 $15$ 和右子樹
$40$。
\item 刪除 $74$：兩個子節點。 這時 $74$ 的右子樹只剩 $95$，後繼 $95$ 就是 $74$ 的右子
節點本身。 $95$ 沒有左子節點，直接坐到 $74$ 的位置，接手左子樹 $63$，右邊是它原本的
右子樹（空的）。
\end{enumerate}

\begin{center}
\begin{minipage}{0.49\linewidth}\centering
\begin{tikzpicture}[bsts]
\node[bnodes] {52}
  child {node {27}
    child {node {15}}
    child {node {40}
      child {node {33}
        child[missing]
        child {node {36}}}
      child[missing]}}
  child {node {74}
    child {node {63}
      child {node {58}}
      child {node {70}}}
    child {node[bhl] {88}
      child[missing]
      child {node {95}}}};
\end{tikzpicture}

{\footnotesize (a) 刪除 $46$ 之後；下一步刪 $88$}
\end{minipage}\hfill
\begin{minipage}{0.49\linewidth}\centering
\begin{tikzpicture}[bsts]
\node[bnodes] {52}
  child {node[bhl] {27}
    child {node {15}}
    child {node {40}
      child {node[bnew] {33}
        child[missing]
        child {node {36}}}
      child[missing]}}
  child {node {74}
    child {node {63}
      child {node {58}}
      child {node {70}}}
    child {node {95}}};
\end{tikzpicture}

{\footnotesize (b) 刪除 $88$ 之後；下一步刪 $27$，後繼是 $33$}
\end{minipage}

\medskip
\begin{minipage}{0.49\linewidth}\centering
\begin{tikzpicture}[bsts]
\node[bnodes] {52}
  child {node {33}
    child {node {15}}
    child {node {40}
      child {node {36}}
      child[missing]}}
  child {node[bhl] {74}
    child {node {63}
      child {node {58}}
      child {node {70}}}
    child {node[bnew] {95}}};
\end{tikzpicture}

{\footnotesize (c) 刪除 $27$ 之後；下一步刪 $74$，後繼是 $95$}
\end{minipage}\hfill
\begin{minipage}{0.49\linewidth}\centering
\begin{tikzpicture}[bsts]
\node[bnodes] {52}
  child {node {33}
    child {node {15}}
    child {node {40}
      child {node {36}}
      child[missing]}}
  child {node[bnew] {95}
    child {node {63}
      child {node {58}}
      child {node {70}}}
    child[missing]};
\end{tikzpicture}

{\footnotesize (d) 刪除 $74$ 之後}
\end{minipage}
\diagcaption{從圖~\ref{fig:bst-T} 依序刪除 $46$、$88$、$27$、$74$。 橘色是下一步要刪
的節點，粉紅色是用來取代它的後繼。 (b) 到 (c) 時，$33$ 先離開原位（$36$ 頂上去），
再坐到 $27$ 的位置；(c) 到 (d) 時，後繼 $95$ 就是 $74$ 的右子節點。}
\label{fig:bst-delete}
\end{center}

\subsection{刪除的虛擬碼}

情況 3 有兩種常見的寫法。 第一種把後繼 $s$ 的鍵值（連同其他資料）複製到 $z$ 裡，再到
右子樹裡刪掉 $s$；節點 $z$ 留在原地，換掉的是它存的東西。 第二種不複製資料，改接指
標，把節點 $s$ 本身搬到 $z$ 的位置。 兩種寫法刪除之後，每個位置上的鍵值完全一樣，畫
出來的樹相同。 差別在節點：如果程式的其他地方存著指向某個節點的指標（例如「目前選取
的會員」），複製的寫法會讓 $s$ 的資料搬到另一個節點去，那些指標就指錯了。

虛擬碼用第一種，寫成遞迴：函式回傳刪除之後這棵子樹的根，呼叫的人把它接回去。 這
和《C 程式設計基礎》第~19~章遞迴插入的寫法一樣，呼叫方式是 $T.\mathit{root} \gets
\proc{Tree-Delete}(T.\mathit{root}, k)$。

\begin{pseudo}[label=alg:bst-delete]{刪除鍵值 $k$，回傳子樹新的根}
\Function{Tree-Delete}{$x, k$}
    \If{$x = \NIL$}
        \State \Return \NIL \Comment{$k$ 不在樹裡}
    \EndIf
    \If{$k < x.\mathit{key}$}
        \State $x.\mathit{left} \gets \Call{Tree-Delete}{x.\mathit{left}, k}$
    \ElsIf{$k > x.\mathit{key}$}
        \State $x.\mathit{right} \gets \Call{Tree-Delete}{x.\mathit{right}, k}$
    \ElsIf{$x.\mathit{left} = \NIL$} \Comment{情況 1、2：沒有左子節點}
        \State \Return $x.\mathit{right}$
    \ElsIf{$x.\mathit{right} = \NIL$} \Comment{情況 2：只有左子節點}
        \State \Return $x.\mathit{left}$
    \Else \Comment{情況 3：兩個子節點}
        \State $s \gets \Call{Tree-Minimum}{x.\mathit{right}}$
        \State $x.\mathit{key} \gets s.\mathit{key}$ \Comment{連同其他資料一起複製}
        \State $x.\mathit{right} \gets \Call{Tree-Delete}{x.\mathit{right}, s.\mathit{key}}$
    \EndIf
    \State \Return $x$
\EndFunction
\end{pseudo}

第 8 行同時處理了情況 1 和「只有右子節點」的情況 2：沒有左子節點時，右子樹（可能是
空的）頂上來。 第 15 行到右子樹裡刪 $s.\mathit{key}$，那一次遞迴會停在 $s$，而 $s$ 沒
有左子節點，走第 9 行，不會再進入情況 3。 被拿掉的節點（第 9、11 行回傳之後的 $x$）
由呼叫的程式釋放。

時間：第 4 到 7 行沿著一條路徑往下找 $k$；找到之後，情況 3 再從 $x$ 往下走到 $s$，
第 15 行的遞迴走的也是這一段路。 全部都在從根往下的一條路徑上，時間 $O(h)$。

\begin{cbookrecap}[第~14~章、第~15~章、第~19~章 19.10、19.11、19.15 節]
\texttt{x->left} 是 \texttt{(*x).left}。 遞迴函式回傳「更新後的子樹根」時，呼叫
端一定要寫 \texttt{x->left = f(x->left, ...)} 把它接回去；只呼叫不接，子樹的改變
會遺失。

\texttt{Node **link} 存的是「某個 \texttt{Node *} 變數的位址」，可能是根指標
\texttt{root}，也可能是某個節點的 \texttt{left} 或 \texttt{right} 欄位。 寫
\texttt{*link = q} 就把那個指標改成指向 \texttt{q}，不必知道它是根還是哪個節點的
哪一邊。 第~19~章用它寫串列的刪除，本書第~5~章用它寫依序插入；這一節把同樣的寫法
用在樹上。
\end{cbookrecap}

\subsection{C 實作}

C 版用第二種寫法（接指標，不複製資料），而且不用遞迴，也不存父指標。 關鍵是
\texttt{find\_link}：從 \texttt{root} 往下找 \texttt{key}，回傳「指向那個節點的指標
欄位」的位址。 找不到時，回傳的欄位是 \texttt{NULL}，正好是插入時要接的空位，所以插
入和刪除共用它。 刪除情況 3 時，再用一個 \texttt{slink} 從 \texttt{\&z->right} 一路往
左找到「指向後繼 \texttt{s} 的欄位」，先 \texttt{*slink = s->right} 把 $s$ 拿下來，
再讓 $s$ 接手 $z$ 的兩棵子樹。

程式依序插入那 $13$ 個編號，查幾個編號在不在、下一位是誰（\proc{Successor-Of-Key}
的 C 版），再依序刪除 $46$、$88$、$27$、$74$。 \texttt{show} 印出樹的形狀，格式是
「鍵值(左子樹,右子樹)」，空的子樹印 \texttt{-}，葉節點只印鍵值，可以對照
圖~\ref{fig:bst-T} 和圖~\ref{fig:bst-delete}。

\begin{example}[label=ex:bst-member]{c}{(member\_tree.c) 會員編號的二元搜尋樹}
#include <stdio.h>
#include <stdlib.h>

typedef struct Node {
    int key;                    // 會員編號（這裡只用末兩碼）
    struct Node *left, *right;
} Node;

// 回傳「指向 key 所在節點」的那個指標欄位的位址。key 不在樹裡時，
// 那個欄位是 NULL，正好是 key 應該插入的空位。
Node **find_link(Node **root, int key)
{
    Node **link = root;
    while (*link != NULL && (*link)->key != key)
        link = key < (*link)->key ? &(*link)->left : &(*link)->right;
    return link;
}

// 成功回傳 1，已經在樹裡回傳 0，要不到記憶體回傳 -1。
int insert(Node **root, int key)
{
    Node **link = find_link(root, key);
    if (*link != NULL)
        return 0;
    Node *z = malloc(sizeof *z);
    if (z == NULL)
        return -1;
    z->key = key;
    z->left = z->right = NULL;
    *link = z;                  // 接在找到的空位上
    return 1;
}

// 比 key 大的最小鍵值所在的節點；沒有就回傳 NULL。
const Node *successor(const Node *x, int key)
{
    const Node *best = NULL;
    while (x != NULL) {
        if (key < x->key) {
            best = x;           // 往左轉：x 是目前的候選
            x = x->left;
        } else {
            x = x->right;
        }
    }
    return best;
}

// 刪除 key。成功回傳 1，找不到回傳 0。
int erase(Node **root, int key)
{
    Node **link = find_link(root, key);
    Node *z = *link;
    if (z == NULL)
        return 0;
    if (z->left == NULL) {              // 沒有左子：右子樹頂上來
        *link = z->right;
    } else if (z->right == NULL) {      // 只有左子：左子樹頂上來
        *link = z->left;
    } else {                            // 兩個子節點：用後繼 s 取代 z
        Node **slink = &z->right;
        while ((*slink)->left != NULL)
            slink = &(*slink)->left;
        Node *s = *slink;
        *slink = s->right;              // 先把 s 從原位拿掉
        s->left = z->left;              // 再讓 s 接手 z 的兩棵子樹
        s->right = z->right;
        *link = s;
    }
    free(z);
    return 1;
}

// 印出樹的形狀：鍵值(左子樹,右子樹)；空的子樹印 -，葉節點只印鍵值。
void show(const Node *x)
{
    if (x == NULL) {
        putchar('-');
        return;
    }
    printf("%d", x->key);
    if (x->left != NULL || x->right != NULL) {
        putchar('(');
        show(x->left);
        putchar(',');
        show(x->right);
        putchar(')');
    }
}

void free_tree(Node *x)         // 後序：先釋放兩棵子樹，再釋放自己
{
    if (x == NULL)
        return;
    free_tree(x->left);
    free_tree(x->right);
    free(x);
}

int main(void)
{
    Node *root = NULL;
    int join[] = {52, 27, 74, 15, 40, 63, 88, 33, 46, 58, 70, 95, 36};
    int ask[] = {46, 27, 60, 95}, quit[] = {46, 88, 27, 74};
    for (int i = 0; i < 13; i++)
        if (insert(&root, join[i]) < 0)
            break;
    show(root);
    printf("\n");
    for (int i = 0; i < 4; i++) {
        const Node *s = successor(root, ask[i]);
        printf("%d %s, next ", ask[i],
               *find_link(&root, ask[i]) != NULL ? "found" : "missing");
        if (s != NULL)
            printf("%d\n", s->key);
        else
            printf("none\n");
    }
    for (int i = 0; i < 4; i++) {
        erase(&root, quit[i]);
        printf("erase %d: ", quit[i]);
        show(root);
        printf("\n");
    }
    free_tree(root);
    return 0;
}
\end{example}
\iosample{}{52(27(15,40(33(-,36),46)),74(63(58,70),88(-,95))) \\ 46 found, next 52 \\ 27 found, next 33 \\ 60 missing, next 63 \\ 95 found, next none \\ erase 46: 52(27(15,40(33(-,36),-)),74(63(58,70),88(-,95))) \\ erase 88: 52(27(15,40(33(-,36),-)),74(63(58,70),95)) \\ erase 27: 52(33(15,40(36,-)),74(63(58,70),95)) \\ erase 74: 52(33(15,40(36,-)),95(63(58,70),-))}

程式有三個地方要說明：

\begin{itemize}
\item \texttt{erase} 的前兩個分支對應虛擬碼的第 8 到 11 行，\texttt{*link} 可能是
\texttt{root} 本身，所以刪根的時候也不用另外處理。
\item 情況 3 的兩個子情況（後繼在深處、後繼就是右子節點）用的是同一段程式。 刪除
$74$ 時，\texttt{slink} 一開始就是 \texttt{\&z->right}，迴圈一次都沒跑；
\texttt{*slink = s->right} 讓 \texttt{z->right} 變成 \texttt{s} 的右子樹（空的），下
一行 \texttt{s->right = z->right} 讀到的就是這個新值，結果仍然正確。 可以對照圖
~\ref{fig:bst-delete} 的 (c)、(d) 自己走一次。
\item 釋放整棵樹要用後序，先釋放兩棵子樹，最後才釋放自己；用前序的話，
\texttt{free(x)} 之後還要讀 \texttt{x->left}。 \texttt{show} 和 \texttt{free\_tree}
的遞迴深度等於樹高，樹長成一條鏈時會到 $n$；搜尋、插入、刪除都用迴圈，沒有這個問
題。
\end{itemize}

\begin{examnote}
刪除題先判斷節點有幾個子節點，寫出是哪一種情況；兩個子節點時寫出用的是後繼還是前
驅，以及它是誰。 題目沒有指定時，選一種並寫明，整題都用同一種。 答案畫完整棵樹，不
要只畫改到的那一塊，最後用中序走訪檢查一次是否仍然遞增。
\end{examnote}

\section{高度：最壞情況與隨機插入}
\label{sec:bst-height}

前兩節的每個操作都沿著一條從根往下的路徑走，時間都是 $O(h)$。 所以 BST 快不快，全看
樹高 $h$。

\subsection{最好與最壞}

$n$ 個節點的二元樹，高度最少是 $\lceil \lg(n + 1) \rceil$，最多是 $n$（第~7~章）。
兩個極端都會出現在 BST 上。 鍵值由小到大插入時，每個新鍵值都比樹上所有的大，從根一
路往右走到底，接在最右邊，樹變成一條往右的鏈，高度 $n$；由大到小插入是往左的
鏈。 這時 BST 和排好的串列一樣，搜尋要 $\Theta(n)$，依序插入 $n$ 個鍵值總共比較
\[
0 + 1 + 2 + \cdots + (n - 1) = \frac{n(n-1)}{2}
\]
次，是 $\Theta(n^2)$。 反過來，如果每次都先插入中間的鍵值（像圖~\ref{fig:bst-binsearch}
那樣），可以讓每一層除了最後一層都填滿，高度是最矮的 $\lceil \lg(n + 1) \rceil$，
也就是 $\lfloor \lg n \rfloor + 1$。

\begin{center}
\begin{minipage}{0.45\linewidth}\centering
\begin{tikzpicture}[bsts, level distance=8mm,
    level 1/.style={sibling distance=10mm}, level 2/.style={sibling distance=10mm},
    level 3/.style={sibling distance=10mm}, level 4/.style={sibling distance=10mm}]
\node[bnodes] {15}
  child[missing]
  child {node {27}
    child[missing]
    child {node {33}
      child[missing]
      child {node {40}
        child[missing]
        child {node {52}}}}};
\end{tikzpicture}

{\footnotesize (a) 依 $15,\allowbreak 27,\allowbreak 33,\allowbreak 40,\allowbreak 52$ 插入，高度 $5$}
\end{minipage}\hfill
\begin{minipage}{0.45\linewidth}\centering
\begin{tikzpicture}[bsts, level 1/.style={sibling distance=20mm},
    level 2/.style={sibling distance=10mm}]
\node[bnodes] {33}
  child {node {27} child {node {15}} child[missing]}
  child {node {40} child[missing] child {node {52}}};
\end{tikzpicture}

{\footnotesize (b) 依 $33,\allowbreak 27,\allowbreak 40,\allowbreak 15,\allowbreak 52$ 插入，高度 $3$}
\end{minipage}
\diagcaption{同一組鍵值，插入順序不同，樹高可以差很多。}
\label{fig:bst-shapes}
\end{center}

實際的資料常常接近排好的順序：會員編號、訂單編號、時間戳記都是一直往上加的。 這正
是最壞的情況。 所以真正用在系統裡的搜尋樹，會在插入和刪除時順便調整形狀，把高度壓在
$O(\lg n)$，這是第~10~章的 AVL 樹和紅黑樹。

\subsection{隨機插入的期望高度}

如果插入順序是隨機的呢？ 說清楚一點：$n$ 個相異鍵值，$n!$ 種插入順序每一種的機率
都一樣，依這個順序插入空樹，得到的樹叫做\textbf{隨機建立的二元搜尋樹}（randomly
built binary search tree）。 結論如下，證明不在本書的範圍，延伸閱讀有出處。

\begin{theorem}
\label{thm:bst-random-height}
$n$ 個相異鍵值隨機建立的二元搜尋樹，期望高度是 $O(\lg n)$。 更精確地說，期望高度
約為 $4.311 \ln n$，大約是 $2.99 \lg n$。
\end{theorem}

為什麼會這樣？ 先看直覺。 第一個插入的鍵值當根，它在所有鍵值裡的名次是隨機的，
所以根把其餘的鍵值切成兩邊，大小是隨機的。 切得太偏（例如一邊不到四分之一）的機率大約
只有一半；每往下一層，子樹的鍵值又被隨機切一次。 一個鍵值每經過一次「切得還
算平均」的節點，它所在的子樹就縮小到原本的四分之三以下，這種事最多發生
$\log_{4/3} n$ 次，所以路徑的長度大約是 $\lg n$ 的常數倍。 這和第~14~章隨機選基準的
快速排序是同一件事：根就是第一個基準，左右子樹就是切開之後的兩堆。 如果快速排序每次
都拿子陣列裡「在插入順序中最早出現」的元素當基準，它比較過的鍵值對，和依同一個順序建
BST 時比較過的鍵值對完全相同。

高度要考慮最深的那一個節點，證明比較長。 平均深度（全部節點深度的平均）就簡單得
多，下面的框給出完整的計算，順便說明「插入順序」和「誰是誰的祖先」的關係，觀念題會
用到。

\begin{whybox}[平均深度大約是 $2 \ln n$]
把鍵值依大小編號成 $1, 2, \ldots, n$（名次）。 先看一個和機率無關的性質。

\textbf{引理。} 依某個順序把相異鍵值插入空樹。 對兩個鍵值 $i \ne j$，$i$ 是 $j$ 的
祖先，若且唯若在 $i$ 和 $j$ 之間（含兩端）的所有鍵值裡，$i$ 是最早插入的。

\textbf{證明。} 設 $I$ 是 $i$ 和 $j$ 之間（含兩端）的鍵值集合，$c$ 是 $I$ 裡最早插入
的。 在 $c$ 插入之前，樹上的節點都不在 $I$ 裡，它們每一個不是比 $I$ 的全部都小，就
是比 $I$ 的全部都大，所以 $I$ 裡的任何鍵值和它們比較，結果都一樣，往同一邊走。
因此 $I$ 的鍵值插入時都會走到 $c$ 所在的那個位置；$c$ 先到，就成為其他人的祖先。
如果 $c = i$，$i$ 就是 $j$ 的祖先。 如果 $c = j$，$j$ 是 $i$ 的祖先，$i$ 就不可能是
$j$ 的祖先。 如果 $c$ 在 $i$、$j$ 中間，$i$ 和 $j$ 到了 $c$ 會分往左右兩邊，誰也不是
誰的祖先。 \qed

隨機插入時，$I$ 有 $|i - j| + 1$ 個鍵值，每一個最早插入的機率都一樣，所以
\[
\Pr[\,i \text{ 是 } j \text{ 的祖先}\,] = \frac{1}{|i - j| + 1}.
\]
$j$ 的深度等於它的祖先個數。 用第~1~章計算插入排序平均情況時的方法（期望值的線性），
\begin{align*}
\mathrm{E}[\,j \text{ 的深度}\,]
&= \sum_{i \ne j} \frac{1}{|i - j| + 1}
= \left(\frac12 + \cdots + \frac1j\right) + \left(\frac12 + \cdots + \frac{1}{n-j+1}\right) \\
&= H_j + H_{n-j+1} - 2,
\end{align*}
其中 $H_m = 1 + \frac12 + \cdots + \frac1m$ 是第~1~章的調和數，$H_m \le \ln m + 1$。 所以每個
鍵值的期望深度都不超過 $2 \ln n$。 對 $j$ 加總，用 $\sum_{j=1}^{n} H_j = (n+1)H_n - n$，
得到全部節點的期望深度總和是 $2(n+1)H_n - 4n$，平均深度是
\[
2\left(1 + \frac1n\right) H_n - 4 \approx 2 \ln n \approx 1.39 \lg n.
\]
\end{whybox}

拿實際的數字對照：$n = 1000$ 時公式給出平均深度約 $11.0$，$n = 10^5$ 時約 $20.2$；
$n$ 變成 $100$ 倍，平均深度只多了大約 $9$。 由小到大插入時，平均深度是 $(n-1)/2$，
$n = 1000$ 時已經是 $499.5$。 用亂數洗牌建樹量出來的平均高度，$n = 1000$ 時大約
$22$、$n = 10^5$ 時大約 $41$，是 $\lg n$ 的 $2.2$ 到 $2.5$ 倍；定理說的 $2.99$ 倍是 $n$ 趨
近無限大時的比值，$n$ 不大時還有一個成長得很慢的負項把它壓低。

最後要提醒兩件事。 第一，定理只保證「插入順序隨機」時的期望值，資料本身的順序我們
通常控制不了，遇到接近排好的資料一樣會退化。 第二，定理只管插入；插入和刪除交錯很
多次之後，樹的形狀就不再是「隨機建立」的分布，沒有這麼簡單的結論。 要在任何順序下
都保證 $O(\lg n)$，就要用第~10~章的平衡樹。

\section{加上子樹大小：第 $k$ 小、排名}
\label{sec:bst-rank}

\subsection{子樹大小}

馬拉松的終點一直有人跑進來，記錄組想隨時回答兩種問題：「目前第 $k$ 快的成績是幾分
鐘？」和「完成時間 $199$ 分鐘的人目前排第幾？」 成績會不斷加入，也可能因為違規被取
消。 只用前兩節的 BST，第一個問題要做中序走訪數到第 $k$ 個，至少走過 $k$ 個節點；第二
個問題要數有幾個比它小，最壞 $\Theta(n)$。

做法是在每個節點多記一個欄位 $x.\mathit{size}$：以 $x$ 為根的子樹有幾個節點（含
$x$）。 約定 \NIL 的 $\mathit{size}$ 是 $0$，那麼對每個節點
\[
x.\mathit{size} = x.\mathit{left}.\mathit{size} + x.\mathit{right}.\mathit{size} + 1.
\]
在 $x$ 自己這棵子樹裡，比 $x$ 小的恰好是左子樹的那些節點，所以 $x$ 在子樹裡排第
$x.\mathit{left}.\mathit{size} + 1$ 小。 這個數字是下面兩個操作的關鍵。

例子：依序加入 $212,\allowbreak 187,\allowbreak 245,\allowbreak 166,\allowbreak 199,\allowbreak 230,\allowbreak 271,\allowbreak 158,\allowbreak 193,\allowbreak 238$ 這 $10$ 個完成時間
（分鐘，假設都不相同；實際上可以用「時間＋號碼布」當鍵值避免重複）。 圖中每個節點
右邊的藍色小字是 $\mathit{size}$。

\begin{center}
\begin{tikzpicture}[bst, bnode/.append style={minimum size=8.5mm},
    level 1/.style={sibling distance=56mm},
    level 2/.style={sibling distance=28mm}, level 3/.style={sibling distance=14mm}]
\node[bnode, label={[bsize]0:10}] {212}
  child {node[label={[bsize]0:5}] {187}
    child {node[label={[bsize]0:2}] {166}
      child {node[label={[bsize]0:1}] {158}}
      child[missing]}
    child {node[label={[bsize]0:2}] {199}
      child {node[label={[bsize]0:1}] {193}}
      child[missing]}}
  child {node[label={[bsize]0:4}] {245}
    child {node[label={[bsize]0:2}] {230}
      child[missing]
      child {node[label={[bsize]0:1}] {238}}}
    child {node[label={[bsize]0:1}] {271}}};
\end{tikzpicture}
\diagcaption{記了子樹大小的 BST。 由小到大是 $158,\allowbreak 166,\allowbreak 187,\allowbreak 193,\allowbreak 199,\allowbreak 212,\allowbreak 230,\allowbreak 238,\allowbreak 245,\allowbreak 271$；根 $212$ 的左子樹有 $5$ 個節點，所以 $212$ 排第 $6$。}
\label{fig:bst-size}
\end{center}

\subsection{第 $k$ 小}

從根開始，$r = x.\mathit{left}.\mathit{size} + 1$ 是 $x$ 在目前這棵子樹裡的名次。
$k = r$ 就是 $x$；$k < r$ 時答案在左子樹，仍然是第 $k$ 小；$k > r$ 時答案在右子樹，
左子樹和 $x$ 共 $r$ 個都比它小，所以在右子樹裡找第 $k - r$ 小。

\begin{pseudo}[label=alg:bst-select]{第 $k$ 小的節點（$k$ 從 $1$ 開始）}
\Function{Tree-Select}{$T, k$}
    \State $x \gets T.\mathit{root}$
    \While{$x \ne \NIL$}
        \State $r \gets x.\mathit{left}.\mathit{size} + 1$ \Comment{$x$ 在這棵子樹裡的名次}
        \If{$k = r$}
            \State \Return $x$
        \ElsIf{$k < r$}
            \State $x \gets x.\mathit{left}$
        \Else
            \State $k \gets k - r$
            \State $x \gets x.\mathit{right}$
        \EndIf
    \EndWhile
    \State \Return \NIL \Comment{$k < 1$ 或 $k$ 大於節點數}
\EndFunction
\end{pseudo}

在圖~\ref{fig:bst-size} 上找第 $7$ 小：

\begin{center}
\begin{tabular}{ccccl}
\toprule
步 & $x.\mathit{key}$ & $k$ & $r$ & 動作 \\
\midrule
1 & 212 & 7 & $5 + 1 = 6$ & $k > r$，$k \gets 7 - 6 = 1$，往右 \\
2 & 245 & 1 & $2 + 1 = 3$ & $k < r$，往左 \\
3 & 230 & 1 & $0 + 1 = 1$ & $k = r$，回傳 $230$ \\
\bottomrule
\end{tabular}
\end{center}

\subsection{排名}

反過來，給一個鍵值 $k$，問樹裡有幾個鍵值 $\le k$。 $k$ 在樹裡時，這就是 $k$ 的名
次。 從根往下搜尋 $k$，每次往右走時，$x$ 和它的左子樹都 $\le k$，把
$x.\mathit{left}.\mathit{size} + 1$ 加進答案；往左走時，$x$ 和右子樹都比 $k$ 大，不
加。

\begin{pseudo}[label=alg:bst-rank]{樹裡有幾個鍵值 $\le k$}
\Function{Tree-Rank}{$T, k$}
    \State $r \gets 0$
    \State $x \gets T.\mathit{root}$
    \While{$x \ne \NIL$}
        \If{$k < x.\mathit{key}$}
            \State $x \gets x.\mathit{left}$
        \Else \Comment{$x$ 和左子樹都 $\le k$}
            \State $r \gets r + x.\mathit{left}.\mathit{size} + 1$
            \State $x \gets x.\mathit{right}$
        \EndIf
    \EndWhile
    \State \Return $r$
\EndFunction
\end{pseudo}

這個版本找到 $k$ 之後沒有停，而是繼續往右子樹走；右子樹的鍵值都比 $k$ 大，之後只會
往左，不會再加。 這樣寫的好處是 $k$ 不在樹裡也能用。 下表是 $k = 199$ 和 $k = 240$
（不在樹裡）的過程。

\begin{center}
\begin{tabular}{ccccl}
\toprule
$k$ & $x.\mathit{key}$ & 比較 & $r$ & 說明 \\
\midrule
$199$ & 212 & $199 < 212$ & 0 & 往左 \\
      & 187 & $199 \ge 187$ & $0 + 2 + 1 = 3$ & 加上 $158, 166, 187$，往右 \\
      & 199 & $199 \ge 199$ & $3 + 1 + 1 = 5$ & 加上 $193, 199$，往右，$x = \NIL$ \\
\midrule
$240$ & 212 & $240 \ge 212$ & $0 + 5 + 1 = 6$ & 往右 \\
      & 245 & $240 < 245$ & 6 & 往左 \\
      & 230 & $240 \ge 230$ & $6 + 0 + 1 = 7$ & 往右 \\
      & 238 & $240 \ge 238$ & $7 + 0 + 1 = 8$ & 往右，$x = \NIL$ \\
\bottomrule
\end{tabular}
\end{center}

$199$ 排第 $5$，比 $240$ 分鐘快的有 $8$ 人。 兩個操作都沿著一條路徑往下走，時間
$O(h)$。 有了 \proc{Tree-Rank}，「完成時間在 $[a, b]$ 之間的有幾人」也是 $O(h)$：
用 $\le b$ 的個數減掉 $< a$ 的個數，觀念題 8 會要你寫出來。

\subsection{插入和刪除時維護 size}

$\mathit{size}$ 只有在插入和刪除時會變，而且只有\textbf{路徑上}的節點會變。

\begin{itemize}
\item \textbf{插入}：新節點的 $\mathit{size} = 1$；從根到新節點的父節點，路上每一棵
子樹都多了一個節點，各加 $1$。 在 \proc{Tree-Insert} 的第 5 行之後加上
$x.\mathit{size} \gets x.\mathit{size} + 1$ 就可以，前提是鍵值確定不在樹裡；如果可
能重複，要先搜尋一次，確定要插入了才加，否則路上的 $\mathit{size}$ 會被多加。
\item \textbf{刪除}：從根到被刪除的節點 $z$（不含 $z$），路上每一棵子樹都少一個，各減
$1$。 情況 3 時後繼 $s$ 離開原位，從 $z.\mathit{right}$ 到 $s$ 的父節點，路上的子樹也
都少了 $s$，各減 $1$；$s$ 坐到 $z$ 的位置後，它的子樹就是 $z$ 原本的子樹少掉 $z$，所
以 $s.\mathit{size} \gets z.\mathit{size} - 1$。
\end{itemize}

以圖~\ref{fig:bst-size} 取消 $187$ 為例（兩個子節點）。 從根往下找 $187$，經過的
$212$ 從 $10$ 減成 $9$。 後繼是右子樹 $\{199, 193\}$ 的最小值 $193$，從 $199$ 往左才
找到，經過的 $199$ 從 $2$ 減成 $1$。 $193$ 坐到 $187$ 的位置，$\mathit{size} = 5 - 1
= 4$。 其他節點的子樹沒有變，$\mathit{size}$ 不動。

路徑上最多 $h$ 個節點，所以多記 $\mathit{size}$ 不會讓插入和刪除變慢，仍然是
$O(h)$。 這個想法可以推廣：一個欄位只要能從「自己的資料」和「兩個子節點的同一個欄
位」算出來（例如子樹的鍵值總和、子樹的最大值），插入和刪除時就只有路徑上的節點要重
算，都能在 $O(h)$ 內維護。 反例是「節點在整棵樹裡的名次」：插入一個很小的鍵值，所有
節點的名次都加 $1$，要改 $\Theta(n)$ 個節點。 第~10~章的平衡樹會用旋轉調整形狀，到時
候還要確認旋轉時這些欄位也能在 $O(1)$ 內重算。

下面的程式依序加入那 $10$ 個成績，查第 $1, 6, 7, 11$ 快的成績，以及 $199$、$240$、
$150$ 分鐘的名次（$\le$ 它的有幾人）。 插入走兩趟：第一趟確定鍵值不在樹裡，第二趟才沿
路把 $\mathit{size}$ 加 $1$。 刪除和範例~\ref{ex:bst-member} 的 \texttt{erase} 一樣，
再照上面的規則沿路減 $1$，留給讀者練習。

\begin{example}[label=ex:bst-rank]{c}{(rank\_tree.c) 第 $k$ 快的成績與名次}
#include <stdio.h>
#include <stdlib.h>

typedef struct Node {
    int key;                    // 完成時間（分鐘），假設都不相同
    int size;                   // 以這個節點為根的子樹有幾個節點
    struct Node *left, *right;
} Node;

int size(const Node *x)
{
    return x == NULL ? 0 : x->size;
}

// 成功回傳 1，已經在樹裡回傳 0，要不到記憶體回傳 -1。
int insert(Node **root, int key)
{
    Node **link = root;                 // 第一趟：確定 key 不在樹裡
    while (*link != NULL && (*link)->key != key)
        link = key < (*link)->key ? &(*link)->left : &(*link)->right;
    if (*link != NULL)
        return 0;
    Node *z = malloc(sizeof *z);
    if (z == NULL)
        return -1;
    for (link = root; *link != NULL; ) { // 第二趟：經過的子樹都多一個
        (*link)->size++;
        link = key < (*link)->key ? &(*link)->left : &(*link)->right;
    }
    z->key = key;
    z->size = 1;
    z->left = z->right = NULL;
    *link = z;
    return 1;
}

// 第 k 小的節點（k 從 1 開始）；k 超出範圍回傳 NULL。
const Node *select_kth(const Node *x, int k)
{
    while (x != NULL) {
        int r = size(x->left) + 1;      // x 在自己這棵子樹裡排第 r
        if (k == r)
            return x;
        if (k < r) {
            x = x->left;
        } else {
            k -= r;
            x = x->right;
        }
    }
    return NULL;
}

// 樹裡有幾個鍵值 <= key。key 在樹裡時，就是它的名次。
int rank_of(const Node *x, int key)
{
    int r = 0;
    while (x != NULL) {
        if (key < x->key) {
            x = x->left;
        } else {
            r += size(x->left) + 1;     // x 和它的左子樹都 <= key
            x = x->right;
        }
    }
    return r;
}

void free_tree(Node *x)
{
    if (x == NULL)
        return;
    free_tree(x->left);
    free_tree(x->right);
    free(x);
}

int main(void)
{
    Node *root = NULL;
    int finish[] = {212, 187, 245, 166, 199, 230, 271, 158, 193, 238};
    int ks[] = {1, 6, 7, 11}, keys[] = {199, 240, 150};
    for (int i = 0; i < 10; i++)
        if (insert(&root, finish[i]) < 0)
            break;
    for (int i = 0; i < 4; i++) {
        const Node *x = select_kth(root, ks[i]);
        if (x != NULL)
            printf("kth %d: %d\n", ks[i], x->key);
        else
            printf("kth %d: none\n", ks[i]);
    }
    for (int i = 0; i < 3; i++)
        printf("rank %d: %d\n", keys[i], rank_of(root, keys[i]));
    free_tree(root);
    return 0;
}
\end{example}
\iosample{}{kth 1: 158 \\ kth 6: 212 \\ kth 7: 230 \\ kth 11: none \\ rank 199: 5 \\ rank 240: 8 \\ rank 150: 0}

\texttt{select\_kth} 和 \texttt{rank\_of} 是虛擬碼~\ref{alg:bst-select}、
\ref{alg:bst-rank} 逐行翻譯，NIL 的 $\mathit{size}$ 由函式 \texttt{size} 回傳 $0$。
$150$ 比所有成績都快，名次查詢回答 $0$，表示它插入的話會排第 $1$。

\begin{examnote}
第 $k$ 小和排名的追蹤，一步寫一列：目前的節點、$k$（或累加的 $r$）、左子樹的
$\mathit{size}$、往哪裡走。 題目只給樹、沒給 $\mathit{size}$ 時，先在每個節點旁邊標
上 $\mathit{size}$（由葉節點往上算），檢查根的 $\mathit{size}$ 等於節點總數，再開始
追蹤。
\end{examnote}

\section{本章常見錯誤}
\label{sec:bst-pitfalls}

\begin{center}
\begin{tabular}{p{0.24\linewidth}p{0.50\linewidth}p{0.16\linewidth}}
\toprule
錯誤 & 例子 & 見 \\
\midrule
只比較父子兩個節點 & 以為每個節點比左子節點大、比右子節點小就是 BST；$4$ 在 $5$
的右子樹裡卻比 $5$ 小 & 第~\ref{sec:bst-query}~節 \\
高度、深度差 1 & 題目的高度數邊（CLRS），卻套用本書數節點的公式，或反過來；把深度
和層數混用 &
第~\ref{sec:bst-query}~節 \\
以為最小值一定是葉節點 & 最小值沒有左子節點，但可能有右子樹；刪除最小值時忘了把右
子樹接回去 & 第~\ref{sec:bst-query}~節 \\
後繼只看右子樹 & 節點沒有右子樹時回答「沒有後繼」，其實後繼在祖先裡 &
第~\ref{sec:bst-query}~節 \\
插入時改到參數的複本 & \texttt{void insert(Node *root, ...)} 裡寫 \texttt{root =
malloc(...)}，呼叫端的樹沒有變 & 第~\ref{sec:bst-modify}~節 \\
遞迴刪除沒有接回子樹 & 呼叫 \texttt{erase(x->left, k)} 卻沒有寫 \texttt{x->left =}，
刪掉的節點還被指著 & 第~\ref{sec:bst-modify}~節 \\
情況 3 刪錯鍵值 & 把後繼的鍵值複製上來之後，到右子樹裡刪的還是原本的 $k$，後繼留了兩
份 & 第~\ref{sec:bst-modify}~節 \\
情況 3 找錯替代節點 & 用右子節點（而不是右子樹的最小值）取代 $z$，右子節點的左子樹
沒地方放 & 第~\ref{sec:bst-modify}~節 \\
用前序釋放 & 先 \texttt{free(x)} 再讀 \texttt{x->left} & 第~\ref{sec:bst-modify}~節 \\
把期望高度當成保證 & 以為 BST 的操作一定是 $O(\lg n)$；依序插入時是 $\Theta(n)$ &
第~\ref{sec:bst-height}~節 \\
重複鍵值時 size 多加 & 插入前沒先確定鍵值不在樹裡，沿路加了 $\mathit{size}$ 才發現重複
& 第~\ref{sec:bst-rank}~節 \\
\bottomrule
\end{tabular}
\end{center}

\section{觀念題}
\label{sec:bst-exercises}

題型和研究所考試、期中期末考的筆試相近，有選擇、計算、追蹤、證明和寫程式。 解答在
本章最後。 除非題目另外說明，鍵值都不相同，高度以節點數計算（第~7~章的定義），刪除
兩個子節點的節點時用後繼取代。

\begin{enumerate}
\item （選擇，可能不只一個答案）$T$ 是有 $n \ge 2$ 個節點的二元搜尋樹。 下列哪些敘述
一定成立？
\begin{enumerate}
\item 最小鍵值所在的節點沒有左子節點。
\item 最小鍵值所在的節點是葉節點。
\item 如果根有右子節點，根的後繼在根的右子樹裡。
\item 前序走訪的第一個鍵值是最小值。
\item 高度為 $h$ 的 BST 最少有 $h$ 個節點，最多有 $2^h - 1$ 個節點。
\end{enumerate}

\item （追蹤）依序把 $61,\allowbreak 34,\allowbreak 80,\allowbreak 12,\allowbreak 47,\allowbreak 70,\allowbreak 93,\allowbreak 25,\allowbreak 40,\allowbreak 55,\allowbreak 88,\allowbreak 43$ 插入空的二元
搜尋樹。
\begin{enumerate}
\item 畫出這棵樹，寫出它的高度和後序走訪。
\item 從 (a) 的樹刪除 $34$，分別用後繼和前驅取代，畫出兩種結果。
\item 從 (b) 用後繼的那棵樹開始，依序刪除 $80$ 和 $61$，畫出每一步之後的樹。
\end{enumerate}

\item （計算）在第 2 題 (a) 的樹上用 \proc{Iterative-Tree-Search}（虛擬碼
~\ref{alg:bst-search-iter}）搜尋整數 $k$。
\begin{enumerate}
\item 搜尋 $50$ 時依序和哪些節點比較？ 如果接著插入 $50$，它會成為哪個節點的哪一邊
子節點？
\item 哪些整數 $k$ 會在恰好和 $4$ 個節點比較之後失敗（回傳 \NIL）？ 用區間寫出答案。
\end{enumerate}

\item （計算）
\begin{enumerate}
\item 依序插入 $50,\allowbreak 30,\allowbreak 70,\allowbreak 20,\allowbreak 40,\allowbreak 60,\allowbreak 45$ 得到一棵 BST。 這 $7$ 個鍵值的 $7!$ 種
插入順序裡，有幾種會得到形狀和鍵值位置都完全相同的樹？
\item 設 $T$ 的根是 $x$，左、右子樹分別有 $a$、$b$ 個節點，能產生它們的插入順序各有
$f_L$、$f_R$ 種。 寫出能產生 $T$ 的插入順序個數，並說明理由。
\item $1, 2, \ldots, n$ 的 $n!$ 種插入順序裡，有幾種得到的樹高度是 $n$？
\end{enumerate}

\item （追蹤）在第 2 題 (a) 的樹上，每個節點加上 $\mathit{size}$。
\begin{enumerate}
\item 寫出每個節點的 $\mathit{size}$。
\item 用表格追蹤 \proc{Tree-Select}$(T, 5)$（虛擬碼~\ref{alg:bst-select}），每一步寫
出目前的節點、$k$、$r$。
\item 用表格追蹤 \proc{Tree-Rank}$(T, 50)$（虛擬碼~\ref{alg:bst-rank}）。
\end{enumerate}

\item （證明）$a < b$ 是 BST 裡的兩個鍵值，而且樹裡沒有任何鍵值介於 $a$ 和 $b$ 之間
（也就是說，排好之後它們相鄰）。 證明 $a$ 和 $b$ 所在的節點，一定有一個是另一個的祖
先。 （這裡的樹不一定只經過插入，也可能做過刪除。）

\item （證明）證明 \proc{Tree-Delete}$(x, k)$（虛擬碼~\ref{alg:bst-delete}）回傳的
樹仍然是二元搜尋樹，而且它的鍵值集合恰好是原本的集合拿掉 $k$（$k$ 不在樹裡時不
變）。 提示：對子樹的節點數做歸納。

\item （設計）$T$ 是記了 $\mathit{size}$ 的 BST，鍵值是整數。
\begin{enumerate}
\item 寫出 \proc{Tree-Floor}$(T, k)$：回傳樹裡 $\le k$ 的最大鍵值所在的節點，沒有就回
傳 \NIL。 這一小題用不到 $\mathit{size}$。
\item 利用 \proc{Tree-Rank}，在 $O(h)$ 時間內算出樹裡有幾個鍵值落在 $[a, b]$ 之間
（$a \le b$）。
\item 如果鍵值是實數，(b) 要怎麼改？
\end{enumerate}

\item （找錯）下面三個函式各有什麼錯？ 說明會發生什麼事，並改正。 \texttt{Node} 的
定義和範例~\ref{ex:bst-member} 相同。
\begin{code}{c}{觀念題 9}
// (a) 插入
void insert(Node *root, int key)
{
    if (root == NULL) {
        root = malloc(sizeof *root);
        root->key = key;
        root->left = root->right = NULL;
    } else if (key < root->key) {
        insert(root->left, key);
    } else if (key > root->key) {
        insert(root->right, key);
    }
}

// (b) 刪除，回傳子樹新的根
Node *del(Node *x, int k)
{
    if (x == NULL)
        return NULL;
    if (k < x->key) {
        x->left = del(x->left, k);
    } else if (k > x->key) {
        x->right = del(x->right, k);
    } else if (x->left == NULL || x->right == NULL) {
        Node *child = x->left != NULL ? x->left : x->right;
        free(x);
        return child;
    } else {
        Node *s = x->right;
        while (s->left != NULL)
            s = s->left;
        x->key = s->key;
        x->right = del(x->right, k);
    }
    return x;
}

// (c) 釋放整棵樹
void free_tree(Node *x)
{
    if (x == NULL)
        return;
    free(x);
    free_tree(x->left);
    free_tree(x->right);
}
\end{code}

\item （計算）隨機建立的 BST：$n$ 個鍵值的每一種插入順序機率都一樣。
\begin{enumerate}
\item $n = 3$，鍵值 $1, 2, 3$。 列出 $6$ 種順序各自得到的樹高，算出期望高度。
\item 同上，算出鍵值 $2$ 的期望深度，並和第~\ref{sec:bst-height}~節框裡的公式
$H_j + H_{n-j+1} - 2$ 比對。
\item 用同一個公式算 $n = 4$ 時鍵值 $1$ 的期望深度。
\end{enumerate}

\end{enumerate}

\section{實作練習}
\label{sec:bst-practice}

下面是 Aki-Judge 上和本章相關的題目。 「二元搜尋樹的插入與搜尋」和「二元搜尋樹的刪
除」要交函式和標頭檔，不寫 \texttt{main}；評測程式會呼叫你的插入、搜尋、刪除和中序
走訪。 「書架編號樹」只交函式，題目規定了節點的欄位名稱，重複的編號要累加本數，還要
回報那個節點的深度（根的深度是 $0$）。

\begin{center}
\begin{tabular}{llc}
\toprule
小節 & 題目 & 難度 \\
\midrule
搜尋、插入 & \akiref{bst-insert-and-search}{二元搜尋樹的插入與搜尋} & 基本 \\
搜尋、插入 & \akiref{shelf-index-tree}{書架編號樹} & 基本 \\
刪除 & \akiref{bst-deleting-keys}{二元搜尋樹的刪除} & 基本 \\
排名 & \akiref{count-keys-in-a-range-of-a-bst}{二元搜尋樹的區間計數} & 基本 \\
\bottomrule
\end{tabular}
\end{center}

「書架編號樹」的編號可能依遞增順序出現，樹會長成一條鏈，深度接近操作數；這時用遞迴
插入，遞迴深度也會跟著變大，用第~\ref{sec:bst-modify}~節的迴圈寫法比較安全。 「二元
搜尋樹的區間計數」資料量小，直接走訪就能過；也可以練習只走進可能有答案的子樹，或照觀
念題 8 加上 $\mathit{size}$。 本章的題目不多，觀念題請一定要做。

\section{讀完你應該會}
\label{sec:bst-checklist}

\begin{itemize}
\item 能依插入順序畫出 BST，寫出每次插入比較過的節點，並說出每個操作的時間是
$O(h)$。
\item 能在給定的樹上追蹤搜尋、最小／最大、後繼、前驅，分清楚「有右子樹」和「往上
找祖先」兩種情況，也會不用父指標從根往下找。
\item 能判斷刪除屬於哪一種情況，用後繼或前驅畫出刪除之後的樹，並用 \texttt{Node **}
寫出沒有特殊情況、不會漏放記憶體的 C 版插入和刪除。
\item 能說明 BST 最壞高度 $n$ 的成因，說出隨機插入的期望高度 $O(\lg n)$，並用祖
先的引理算出平均深度。
\item 能在節點加上 $\mathit{size}$，追蹤第 $k$ 小和排名，說明插入、刪除時哪些節點的
$\mathit{size}$ 要改。
\end{itemize}

\textbf{延伸閱讀}（書末會整理完整的參考資料）

\begin{itemize}
\item Thomas H. Cormen 等，\textit{Introduction to Algorithms}，第 4 版，第 12 章：
BST 的查詢、插入和刪除，刪除用另一種接指標的寫法並且維護父指標，可以和本章對照。
第 17.1 節是加上子樹大小的順序統計樹，建在紅黑樹上。 隨機建立的 BST 期望高度
$O(\lg n)$ 的完整證明在第 3 版第 12.4 節。
\item Donald E. Knuth，\textit{The Art of Computer Programming}，第 3 冊，第 6.2.2
節：二元樹搜尋的原始分析，包括隨機插入時平均比較次數的精確公式，以及刪除對樹形的影
響。
\item Steven S. Skiena，\textit{The Algorithm Design Manual}，第 3 版，第 3.4 節：
BST 的實作與「BST 到底好不好」的討論，偏重實務。
\item 《C 程式設計基礎》第~19~章：BST 的不變式、遞迴插入、三種走訪和由走訪重建，C
語法寫得比本章詳細。
\item 本書第~10~章（AVL 樹、紅黑樹）保證高度 $O(\lg n)$；第~18~章的最佳二元搜尋樹
討論已知查詢頻率時，怎麼排出平均比較次數最少的樹。
\end{itemize}

\addcontentsline{toc}{section}{觀念題解答}
\section*{觀念題解答}

\begin{enumerate}
\item (a)、(c)、(e)。 (a)：有左子節點就有更小的鍵值。 (b) 不對：最小值可以有右子
樹，例如依序插入 $1, 2$。 (c)：根有右子樹，後繼是右子樹的最小值。 (d) 不對：前序的
第一個是根，根不一定最小。 (e)：高度 $h$ 表示有 $h$ 層，從根到第 $h$ 層的某個節點
有一條 $h$ 個節點的路徑，所以至少 $h$ 個；第 $i$ 層最多 $2^{i-1}$ 個節點，
$i = 1, \ldots, h$ 加起來最多 $2^h - 1$ 個。

\item
\begin{enumerate}
\item 樹如下，高度 $5$（$61 \to 34 \to 47 \to 40 \to 43$，$5$ 個節點）。 後序走訪：$25,\allowbreak 12,\allowbreak 43,\allowbreak 40,\allowbreak 55,\allowbreak 47,\allowbreak 34,\allowbreak 70,\allowbreak 88,\allowbreak 93,\allowbreak 80,\allowbreak 61$。
\begin{center}
\begin{tikzpicture}[bsts]
\node[bnodes] {61}
  child {node {34}
    child {node {12} child[missing] child {node {25}}}
    child {node {47}
      child {node {40} child[missing] child {node {43}}}
      child {node {55}}}}
  child {node {80}
    child {node {70}}
    child {node {93} child {node {88}} child[missing]}};
\end{tikzpicture}
\end{center}

\item 用後繼：$34$ 的後繼是右子樹 $\{47, 40, 43, 55\}$ 的最小值 $40$。 $40$ 只有右子
節點 $43$，$43$ 頂上來成為 $47$ 的左子節點，$40$ 坐到 $34$ 的位置。 用前驅：前驅是
左子樹 $\{12, 25\}$ 的最大值 $25$，它是葉節點，直接拿下來坐到 $34$ 的位置，$12$ 的右
子節點變成 \NIL。 寫成「鍵值(左,右)」：
\begin{center}
\begin{tabular}{ll}
後繼 & $61(40(12(-,25),47(43,55)),80(70,93(88,-)))$ \\
前驅 & $61(25(12,47(40(-,43),55)),80(70,93(88,-)))$
\end{tabular}
\end{center}

\item 刪除 $80$：兩個子節點，後繼是 $93$ 子樹的最小值 $88$（葉節點），得到
$61(40(12(-,25),47(43,55)),88(70,93))$。 刪除 $61$：後繼是右子樹 $\{88, 70, 93\}$ 的
最小值 $70$（葉節點），得到 $70(40(12(-,25),47(43,55)),88(-,93))$。
\end{enumerate}

\item
\begin{enumerate}
\item $61 \to 34 \to 47 \to 55$，$50 < 55$ 往左是 \NIL，失敗。 插入 $50$ 會成為 $55$
的左子節點。
\item 失敗時比較的節點數等於「最後一個節點的深度 $+ 1$」，所以要找深度 $3$ 的節點底下
的空位。 深度 $3$ 的節點是 $25, 40, 55, 88$（$43$ 在深度 $4$）。 每個空位對應一個區
間：$25$ 的左、右是 $(12, 25)$、$(25, 34)$；$40$ 的左是 $(34, 40)$（右邊有 $43$）；
$55$ 的左、右是 $(47, 55)$、$(55, 61)$；$88$ 的左、右是 $(80, 88)$、$(88, 93)$。
整數答案：$13 \sim 24$、$26 \sim 33$、$35 \sim 39$、$48 \sim 54$、$56 \sim 60$、
$81 \sim 87$、$89 \sim 92$。
\end{enumerate}

\item
\begin{enumerate}
\item 樹是 $50(30(20,40(-,45)),70(60,-))$。 用 (b) 的公式：$50$ 必須第一個插入；左子
樹 $\{30, 20, 40, 45\}$ 和右子樹 $\{70, 60\}$ 的插入順序可以任意交錯。 左子樹裡 $30$
最先，接著 $\{20\}$ 和 $\{40, 45\}$ 交錯，$40$ 要在 $45$ 前面：$\binom{3}{1} \cdot 1
\cdot 1 = 3$ 種。 右子樹只能 $70, 60$：$1$ 種。 合計 $\binom{6}{4} \cdot 3 \cdot 1 =
45$ 種。
\item $\binom{a+b}{a} f_L f_R$。 根必須最先插入，否則先插入的那個鍵值會成為根。 之後
左子樹的鍵值插入時一定走到根的左邊，只和左子樹的其他鍵值互相影響；右子樹也一樣。 所
以一個插入順序產生 $T$，若且唯若根在最前面，而且把它拆成「左子樹的鍵值」和「右子樹
的鍵值」兩個子序列時，各自產生左、右子樹。 左右兩個子序列各有 $f_L$、$f_R$ 種，在
後面 $a + b$ 個位置裡，選 $a$ 個位置放左子樹的鍵值，有 $\binom{a+b}{a}$ 種選法。
\item 高度 $n$ 表示 $n$ 層各一個節點，每個節點最多一個子節點，樹是一條鏈。 這時每一
步插入的都必須是剩下的鍵值裡最小或最大的那個，否則剩下的鍵值會分到它的兩邊。 前
$n - 1$ 步各有 $2$ 種選擇，最後一步只剩一個，共 $2^{n-1}$ 種。 例如 $n = 7$ 時是 $64$ 種。
\end{enumerate}

\item
\begin{enumerate}
\item 葉節點 $25,\allowbreak 43,\allowbreak 55,\allowbreak 70,\allowbreak 88$ 都是 $1$；$12$：$2$；$40$：$2$；$47$：$4$；$34$：
$7$；$93$：$2$；$80$：$4$；$61$：$12$。
\item
\begin{tabular}[t]{ccccl}
\toprule
步 & $x.\mathit{key}$ & $k$ & $r$ & 動作 \\
\midrule
1 & 61 & 5 & $7 + 1 = 8$ & $k < r$，往左 \\
2 & 34 & 5 & $2 + 1 = 3$ & $k > r$，$k \gets 2$，往右 \\
3 & 47 & 2 & $2 + 1 = 3$ & $k < r$，往左 \\
4 & 40 & 2 & $0 + 1 = 1$ & $k > r$，$k \gets 1$，往右 \\
5 & 43 & 1 & $0 + 1 = 1$ & $k = r$，回傳 $43$ \\
\bottomrule
\end{tabular}

第 $5$ 小是 $43$（由小到大：$12,\allowbreak 25,\allowbreak 34,\allowbreak 40,\allowbreak 43,\allowbreak \ldots$）。
\item
\begin{tabular}[t]{cccl}
\toprule
$x.\mathit{key}$ & 比較 & $r$ & 動作 \\
\midrule
61 & $50 < 61$ & 0 & 往左 \\
34 & $50 \ge 34$ & $0 + 2 + 1 = 3$ & 往右 \\
47 & $50 \ge 47$ & $3 + 2 + 1 = 6$ & 往右 \\
55 & $50 < 55$ & 6 & 往左，$x = \NIL$ \\
\bottomrule
\end{tabular}

回傳 $6$：$\le 50$ 的有 $12,\allowbreak 25,\allowbreak 34,\allowbreak 40,\allowbreak 43,\allowbreak 47$。
\end{enumerate}

\item 設 $c$ 是 $a$ 和 $b$ 的最低共同祖先（同時是兩者祖先的節點裡最深的那個；節點也
算自己的祖先）。 如果 $c$ 就是 $a$ 或 $b$，結論成立。 否則 $a$ 和 $b$ 不可能在 $c$
的同一棵子樹裡（不然 $c$ 的那個子節點也是共同祖先，比 $c$ 深），所以一個在左子樹、一
個在右子樹。 由 BST 性質，左邊的比 $c.\mathit{key}$ 小，右邊的比 $c.\mathit{key}$
大，而 $a < b$，所以 $a < c.\mathit{key} < b$，和「沒有鍵值介於 $a$、$b$ 之間」矛
盾。 證明只用到 BST 性質，和樹是怎麼來的無關，所以做過刪除也成立。

\item 對 $x$ 這棵子樹的節點數 $m$ 做歸納，證明：如果 $x$ 是一棵 BST，
\proc{Tree-Delete}$(x, k)$ 回傳一棵 BST，鍵值集合是原本的集合拿掉 $k$。

$m = 0$：$x = \NIL$，回傳 \NIL，成立。

$m \ge 1$，分情況。 (i) $k < x.\mathit{key}$：由歸納假設，第 5 行之後左子樹是 BST，
鍵值是原本的左子樹拿掉 $k$，仍然都比 $x.\mathit{key}$ 小；右子樹沒動；所以整棵仍是
BST，鍵值少了 $k$（$k$ 不在樹裡時不變）。 (ii) $k > x.\mathit{key}$：對稱。
(iii) $k = x.\mathit{key}$ 而且至少一個子樹是空的（第 8 到 11 行）：回傳另一棵子樹，
它本來就是 BST，鍵值恰好是原本的集合拿掉 $k$。 (iv) 兩個子樹都不空：$s$ 是右子樹的
最小值。 第 14 行之後 $x.\mathit{key} = s.\mathit{key}$；由歸納假設，第 15 行之後右
子樹是 BST，鍵值是原本的右子樹拿掉 $s.\mathit{key}$，剩下的都比 $s.\mathit{key}$
大。 左子樹的鍵值都比原本的 $k$ 小，而 $k < s.\mathit{key}$，所以也都比新的
$x.\mathit{key}$ 小。 因此 $x$ 仍滿足 BST 性質。 鍵值集合：$k$ 被 $s.\mathit{key}$
蓋掉，右子樹少了一份 $s.\mathit{key}$，合起來恰好少了 $k$。

\item
\begin{enumerate}
\item 和 \proc{Successor-Of-Key} 對稱，往右走時記下候選：
\begin{pseudo}{觀念題 8 的解答}
\Function{Tree-Floor}{$T, k$}
    \State $f \gets \NIL$
    \State $x \gets T.\mathit{root}$
    \While{$x \ne \NIL$}
        \If{$x.\mathit{key} = k$}
            \State \Return $x$
        \ElsIf{$x.\mathit{key} < k$}
            \State $f \gets x$ \Comment{$x$ 是目前 $\le k$ 的最大候選}
            \State $x \gets x.\mathit{right}$
        \Else
            \State $x \gets x.\mathit{left}$
        \EndIf
    \EndWhile
    \State \Return $f$
\EndFunction
\end{pseudo}
只走一條路徑，$O(h)$。
\item 答案是 $\proc{Tree-Rank}(T, b) - \proc{Tree-Rank}(T, a - 1)$：前者是 $\le b$ 的
個數，後者是 $\le a - 1$，也就是 $< a$ 的個數（鍵值是整數）。 兩次各 $O(h)$。
\item 實數沒有「$a - 1$」可以用。 改成 $\proc{Tree-Rank}(T, b) - \proc{Tree-Rank}(T,
a)$，如果 $a$ 在樹裡（用 \proc{Tree-Search} 查，$O(h)$），再加 $1$。 也可以另外寫一個
「$< k$ 的個數」，把 \proc{Tree-Rank} 第 5 行的 $k < x.\mathit{key}$ 改成 $k \le
x.\mathit{key}$。
\end{enumerate}

\item
(a) \texttt{root} 是參數的複本。 樹是空的時，\texttt{malloc} 來的節點只存在函式裡，
回到呼叫端時樹還是空的，而且那塊記憶體再也沒人指著（洩漏）。 遞迴到
\texttt{root->left == NULL} 時也一樣，新節點沒有接上。 改法：收 \texttt{Node **}，寫
\texttt{*root = ...}；或回傳子樹的根，呼叫端寫 \texttt{root->left = insert(root->left,
key)}。

(b) 兩個子節點時，把後繼的鍵值複製到 \texttt{x} 之後，到右子樹刪的還是 \texttt{k}。
\texttt{k} 比右子樹的每個鍵值都小，什麼也刪不到，結果是後繼的鍵值出現兩次，\texttt{k}
卻消失了（而且沒有任何節點被釋放）。 改成 \texttt{x->right = del(x->right, x->key);}
（此時 \texttt{x->key} 就是後繼的鍵值）。 其他部分是對的：只有一個子節點或沒有子節點
時，\texttt{child} 是那個子節點或 \texttt{NULL}。

(c) \texttt{free(x)} 之後又讀 \texttt{x->left}、\texttt{x->right}，是讀已經釋放的記憶
體。 改成後序：先遞迴釋放兩棵子樹，最後才 \texttt{free(x)}。

\item
\begin{enumerate}
\item $2$ 先插入的兩種順序（$2, 1, 3$ 和 $2, 3, 1$）得到高度 $2$；其他四種（$1$ 或 $3$
先插入）都是一條鏈，高度 $3$。 期望高度 $(2 \cdot 2 + 4 \cdot 3)/6 = 8/3$。
\item $2$ 最先插入時深度 $0$（$2$ 種）。 $1, 2, 3$ 和 $3, 2, 1$ 時深度 $1$；$1, 3, 2$
和 $3, 1, 2$ 時深度 $2$。 期望深度 $(0 + 0 + 1 + 1 + 2 + 2)/6 = 1$。 公式：$j = 2$，
$n = 3$，$H_2 + H_2 - 2 = \frac32 + \frac32 - 2 = 1$，相同。
\item $j = 1$，$n = 4$：$H_1 + H_4 - 2 = 1 + \frac{25}{12} - 2 = \frac{13}{12}$。
\end{enumerate}

\end{enumerate}

\chapter{平衡搜尋樹：AVL 與紅黑樹}
\label{ch:balanced}

% 本章的圖共用的樣式
\definecolor{rbRed}{HTML}{C8102E}
\tikzset{
    tn/.style={circle, draw=black!65, minimum size=0.56cm, inner sep=0pt, font=\small},
    tnA/.style={draw=memStroke, very thick, fill=memFill},      % 失衡的節點
    tnN/.style={draw=coreStroke, thick, fill=coreFill},         % 剛插入的節點
    tnB/.style={tn, draw=black, fill=black!85, text=white},     % 紅黑樹的黑節點
    tnR/.style={tn, draw=rbRed, fill=rbRed, text=white},        % 紅黑樹的紅節點
    tnG/.style={tn, draw=black!50, fill=black!25},              % 紅黑都可以
    tnil/.style={rectangle, fill=black!85, minimum size=0.17cm, inner sep=0pt},
    te/.style={draw=black!60},
    bfl/.style={anchor=south west, font=\scriptsize, text=black!60, inner sep=0.5pt},
    annl/.style={font=\small\bfseries, inner sep=0.8pt, text=memStroke},
    plab/.style={font=\small},
    tarr/.style={-{Stealth[length=2mm]}, thick, draw=black!70},
}
% \subtri[高]{x}{y}{名稱}：頂點在 (x, y)、往下畫的子樹三角形
\newcommand{\subtri}[4][0.9]{%
    \draw[black!60, fill=black!4] (#2,#3) -- ({#2-0.4},{#3-#1}) -- ({#2+0.4},{#3-#1}) -- cycle;
    \node[font=\scriptsize] at ({#2},{#3-0.66*#1}) {#4};
}

學校健身房把會員資料存在一棵二元搜尋樹裡，鍵是會員編號。 編號是開卡時依序發的，
所以每一位新會員的編號，都比樹裡所有的編號大。 三年後累積了一萬名會員，櫃台要查
一位剛入會的會員，從根往下要比幾次？ 如果每次插入之後，允許你花一點工夫把樹整理
一下，你希望它整理成什麼樣子？ 每次整理願意花多少時間？ 往下讀之前，先想一想。

第~9~章的二元搜尋樹，搜尋、插入、刪除都要 $O(h)$ 的時間，$h$ 是樹高。 鍵依序插入
時，樹會長成一條鏈，$h = n$，上面那個問題的答案是一萬次。 這一章講兩種在每次插入
和刪除之後，只沿著一條路徑做局部調整，就能把高度維持在 $O(\lg n)$ 的樹：\textbf{AVL 樹}
用「左右子樹高度差不超過 $1$」當條件；\textbf{紅黑樹}用節點的顏色當條件，條件比較
寬，調整的次數比較少。 兩種樹的調整都靠同一個動作，叫\textbf{旋轉}。

這一章用到《C 程式設計基礎》第~14~章的結構指標、第~15~章的遞迴和第~19~章的
\texttt{malloc}、\texttt{free}；樹的名詞（根、子樹、葉節點、高度）見第~7~章，二元搜
尋樹的插入和刪除見第~9~章。

本章的\textbf{高度}和全書一樣用第~7~章的定義：從根往下最長的一條路徑上的
\textbf{節點數}，也就是層數，空樹的高度是 $0$，只有一個節點的樹高度是 $1$；深度仍然數
邊，根的深度是 $0$。 有些書（例如 CLRS）數的是邊數，葉節點的高度是 $0$、空
樹是 $-1$，兩種數法只差 $1$。 本章用到的平衡因子是兩邊高度相減，不受影響；但
「高度 $h$ 的 AVL 樹最少幾個節點」這類題目，答案會跟著差一格，考卷有給定義時照題目
的定義算。

\section{為什麼要平衡、旋轉}
\label{sec:bal-rotate}

\subsection{樹會變成一條鏈}

把 $11, 22, 33, 44, 55$ 依序插入空的二元搜尋樹，每個新鍵都比樹裡的鍵大，一路往右
走到底，接在最右邊，得到圖~\ref{fig:bal-chain}(a) 的一條鏈。 找 $55$ 要比 $5$ 次，
和在串列裡從頭找一樣。 同樣的五個鍵，如果樹長成 (b) 的樣子，最多比 $3$ 次。

\begin{center}
\begin{tikzpicture}[x=0.55cm,y=0.7cm,baseline=(current bounding box.north)]
\node[tn] (chain11) at (0,0) {11};
\node[tn] (chain22) at (1,-1) {22};
\node[tn] (chain33) at (2,-2) {33};
\node[tn] (chain44) at (3,-3) {44};
\node[tn] (chain55) at (4,-4) {55};
\draw[te] (chain11) -- (chain22);
\draw[te] (chain22) -- (chain33);
\draw[te] (chain33) -- (chain44);
\draw[te] (chain44) -- (chain55);
\node[plab] at (2.0,-4.75) {(a) 依序插入的二元搜尋樹};
\end{tikzpicture}\hspace{1.2cm}\begin{tikzpicture}[x=0.55cm,y=0.7cm,baseline=(current bounding box.north)]
\node[tn] (chainavl11) at (0,-1) {11};
\node[tn] (chainavl22) at (1,0) {22};
\node[tn] (chainavl33) at (2,-2) {33};
\node[tn] (chainavl44) at (3,-1) {44};
\node[tn] (chainavl55) at (4,-2) {55};
\draw[te] (chainavl22) -- (chainavl11);
\draw[te] (chainavl22) -- (chainavl44);
\draw[te] (chainavl44) -- (chainavl33);
\draw[te] (chainavl44) -- (chainavl55);
\node[plab] at (2.0,-4.75) {(b) 同樣五個鍵的 AVL 樹};
\end{tikzpicture}
\diagcaption{(a) 鍵依序插入，二元搜尋樹退化成一條鏈，高度 $5$。 (b) 同樣的五個鍵
插入本章的 AVL 樹，高度 $3$。}
\label{fig:bal-chain}
\end{center}

第~7~章證明過，$n$ 個節點的二元樹高度至少是 $\lceil \lg(n+1) \rceil$。 第~9~章說過，鍵以隨機的順
序插入時，期望的高度是 $O(\lg n)$。 但實際的資料常常不隨機：會員編號、訂單編號、
時間戳記都是遞增的，這時二元搜尋樹的每個操作都變成 $\Theta(n)$。

一種搜尋樹叫做\textbf{平衡}的（balanced），是指它有一個條件，保證任何滿足條件、有
$n$ 個節點的樹，高度都是 $O(\lg n)$，而且插入和刪除之後可以在 $O(\lg n)$ 的時間內
把樹調整回滿足條件的狀態。 每次插入後都把樹重建成最矮的樣子不划算，重建要
$\Theta(n)$。 我們要的是只動插入點到根這條路徑附近的少數指標。

\subsection{旋轉}

\textbf{旋轉}（rotation）只改兩、三個指標，就讓一個子節點升上來、它的父節點降下
去。 圖~\ref{fig:bal-rotation} 左邊的樹，$y$ 的左子節點是 $x$，$T_1$、$T_2$、$T_3$
是三棵子樹（可以是空的）。 對 $y$ 做\textbf{右旋}（right rotation），$x$ 升上來當
這一段的根，$y$ 變成 $x$ 的右子節點，原本夾在中間的 $T_2$ 改接到 $y$ 的左邊。 反
過來，對右邊那棵樹的 $x$ 做\textbf{左旋}（left rotation），就回到左邊的樣子。

\begin{center}
\begin{tikzpicture}[x=0.9cm, y=0.9cm]
\node[tn] (y) at (0,0) {$y$};
\node[tn] (x) at (-1,-1) {$x$};
\draw[te] (y) -- (x); \draw[te] (y) -- (0.8,-1.0); \draw[te] (x) -- (-1.6,-1.95); \draw[te] (x) -- (-0.4,-1.95);
\subtri{-1.6}{-1.95}{$T_1$}
\subtri{-0.4}{-1.95}{$T_2$}
\subtri{0.8}{-1.0}{$T_3$}
\draw[tarr] (1.9,-0.7) -- node[above, font=\small] {右旋 $y$} (3.6,-0.7);
\draw[tarr] (3.6,-1.3) -- node[below, font=\small] {左旋 $x$} (1.9,-1.3);
\begin{scope}[xshift=5.6cm]
\node[tn] (x2) at (0,0) {$x$};
\node[tn] (y2) at (1,-1) {$y$};
\draw[te] (x2) -- (y2); \draw[te] (x2) -- (-0.8,-1.0); \draw[te] (y2) -- (0.4,-1.95); \draw[te] (y2) -- (1.6,-1.95);
\subtri{-0.8}{-1.0}{$T_1$}
\subtri{0.4}{-1.95}{$T_2$}
\subtri{1.6}{-1.95}{$T_3$}
\end{scope}
\end{tikzpicture}
\diagcaption{右旋和左旋。 兩邊的中序都是 $T_1, x, T_2, y, T_3$。 右旋後 $T_1$ 上升
一層，$T_3$ 下降一層，$T_2$ 的深度不變。}
\label{fig:bal-rotation}
\end{center}

旋轉不會破壞二元搜尋樹的性質。 旋轉前，$T_1$ 的鍵都小於 $x$，$T_2$ 的鍵在 $x$ 和
$y$ 之間，$T_3$ 的鍵都大於 $y$；旋轉後每棵子樹接的位置，仍然在同樣兩個鍵的中間，
所以中序走訪的順序 $T_1, x, T_2, y, T_3$ 完全沒變。 改的指標只有三個：$y$ 的左指
標、$x$ 的右指標，以及原本指向 $y$ 的那個指標（$y$ 的父節點的某個子指標，或是
根）。 所以旋轉是 $\Theta(1)$。

\begin{pseudo}[label=alg:bal-rotate-right]{右旋，回傳這一段新的根}
\Function{Rotate-Right}{$y$}
    \State $x \gets y.\mathit{left}$
    \State $y.\mathit{left} \gets x.\mathit{right}$ \Comment{$T_2$ 改接到 $y$ 底下}
    \State $x.\mathit{right} \gets y$
    \State \Return $x$
\EndFunction
\end{pseudo}

\proc{Rotate-Left}$(x)$ 是把上面的 $\mathit{left}$ 和 $\mathit{right}$ 對調。 這裡的
寫法不用父指標：呼叫的人負責把原本指向 $y$ 的指標改成回傳的 $x$，例如
$t.\mathit{left} \gets \proc{Rotate-Right}(t.\mathit{left})$。 AVL 樹的程式用這種寫
法；紅黑樹的程式要往上找父節點和叔叔，第~\ref{sec:bal-rb}~節會改成帶父指標的版本。

記旋轉時有一個簡單的檢查：升上去的節點（右旋時的 $x$）原本靠內側的那棵子樹
（$T_2$），會換到降下來的節點（$y$）底下。 其他子樹都跟著原本的父節點走。

\section{AVL 樹}
\label{sec:bal-avl}

\subsection{平衡因子}

\begin{definition}[AVL 樹]
\label{def:bal-avl}
節點 $x$ 的\textbf{平衡因子}（balance factor）是
\[
\mathrm{BF}(x) = h(x.\mathit{left}) - h(x.\mathit{right}),
\]
$h$ 是子樹的高度，空樹的高度是 $0$。 一棵二元搜尋樹如果\textbf{每一個}節點的平衡
因子都是 $-1$、$0$ 或 $+1$，就叫做 \textbf{AVL 樹}。
\end{definition}

這個名字來自 1962 年提出它的兩位作者 Adelson-Velsky 和 Landis。 平衡因子是正的表
示左邊比較高。 圖~\ref{fig:bal-bf} 的每個節點右上方標了它的平衡因子。 (b) 的根是
$+1$，看起來沒問題，但 $20$ 的左子樹高度 $2$、右子樹是空的，平衡因子 $+2$，所以整
棵不是 AVL 樹。 條件要對每個節點都檢查，不是只看根。

\begin{center}
\begin{tikzpicture}[x=0.6cm,y=0.85cm,baseline=(current bounding box.north)]
\node[tn] (bfgood10) at (0,-2) {10};
\node[tn] (bfgood20) at (1,-1) {20};
\node[tn] (bfgood30) at (2,-2) {30};
\node[tn] (bfgood40) at (3,0) {40};
\node[tn] (bfgood60) at (4,-1) {60};
\node[tn] (bfgood70) at (5,-2) {70};
\draw[te] (bfgood40) -- (bfgood20);
\draw[te] (bfgood20) -- (bfgood10);
\draw[te] (bfgood20) -- (bfgood30);
\draw[te] (bfgood40) -- (bfgood60);
\draw[te] (bfgood60) -- (bfgood70);
\node[bfl] at (bfgood10.north east) {$0$};
\node[bfl] at (bfgood20.north east) {$0$};
\node[bfl] at (bfgood30.north east) {$0$};
\node[bfl] at (bfgood40.north east) {$0$};
\node[bfl] at (bfgood60.north east) {$-1$};
\node[bfl] at (bfgood70.north east) {$0$};
\node[plab] at (2.5,-3.75) {(a) AVL 樹};
\end{tikzpicture}\hspace{1.5cm}\begin{tikzpicture}[x=0.6cm,y=0.85cm,baseline=(current bounding box.north)]
\node[tn] (bfbad5) at (0,-3) {5};
\node[tn] (bfbad10) at (1,-2) {10};
\node[tn,tnA] (bfbad20) at (2,-1) {20};
\node[tn] (bfbad50) at (3,0) {50};
\node[tn] (bfbad60) at (4,-2) {60};
\node[tn] (bfbad70) at (5,-1) {70};
\draw[te] (bfbad50) -- (bfbad20);
\draw[te] (bfbad20) -- (bfbad10);
\draw[te] (bfbad10) -- (bfbad5);
\draw[te] (bfbad50) -- (bfbad70);
\draw[te] (bfbad70) -- (bfbad60);
\node[bfl] at (bfbad5.north east) {$0$};
\node[bfl] at (bfbad10.north east) {$+1$};
\node[bfl] at (bfbad20.north east) {$+2$};
\node[bfl] at (bfbad50.north east) {$+1$};
\node[bfl] at (bfbad60.north east) {$0$};
\node[bfl] at (bfbad70.north east) {$+1$};
\node[plab] at (2.5,-3.75) {(b) 不是 AVL 樹};
\end{tikzpicture}
\diagcaption{節點右上方的數字是平衡因子。 (b) 的根平衡，但 $20$ 不平衡（橘色）。}
\label{fig:bal-bf}
\end{center}

\subsection{高度上界}
\label{sec:bal-avl-height}

AVL 的條件只限制每個節點左右差一層，看起來很寬鬆。 要說明它真的保證
$h = O(\lg n)$，反過來問：高度 $h$ 的 AVL 樹\textbf{最少}有幾個節點？ 記成 $N(h)$。

$N(0) = 0$，$N(1) = 1$，$N(2) = 2$（根加一個子節點）。 $h \ge 2$ 時，高度 $h$ 的樹
的根有一棵子樹高度是 $h - 1$；根的平衡因子不超過 $1$，所以另一棵子樹高度至少是
$h - 2$。 兩棵子樹也都是 AVL 樹。 高度 $h$ 的 AVL 樹裡面一定有一棵高度 $h - 1$ 的
子樹，所以 $N(h) > N(h-1)$，$N$ 是遞增的；要節點最少，另一棵子樹就取高度剛好
$h - 2$、而且節點最少的樣子。 因此
\[
N(h) = N(h-1) + N(h-2) + 1, \qquad h \ge 2.
\]

\begin{center}
\begin{tabular}{c*{11}{c}}
\toprule
$h$ & 0 & 1 & 2 & 3 & 4 & 5 & 6 & 7 & 8 & 9 & 10 \\
\midrule
$N(h)$ & 0 & 1 & 2 & 4 & 7 & 12 & 20 & 33 & 54 & 88 & 143 \\
$F_{h+2}$ & 1 & 2 & 3 & 5 & 8 & 13 & 21 & 34 & 55 & 89 & 144 \\
\bottomrule
\end{tabular}
\end{center}

達到 $N(h)$ 的樹，每個非葉節點的兩棵子樹都差一層，叫做\textbf{費氏樹}（Fibonacci
tree）。 圖~\ref{fig:bal-fib} 是 $h = 1$ 到 $5$ 的樣子（左右對調也可以）。

\begin{center}
\begin{tikzpicture}[x=0.42cm,y=0.62cm,baseline=(current bounding box.north)]
\node[tn] (fb1x1) at (0,0) {};
\node[plab] at (0.0,-4.75) {$h = 1$};
\end{tikzpicture}\hspace{0.5cm}\begin{tikzpicture}[x=0.42cm,y=0.62cm,baseline=(current bounding box.north)]
\node[tn] (fb2x1) at (0,-1) {};
\node[tn] (fb2x2) at (1,0) {};
\draw[te] (fb2x2) -- (fb2x1);
\node[plab] at (0.5,-4.75) {$h = 2$};
\end{tikzpicture}\hspace{0.5cm}\begin{tikzpicture}[x=0.42cm,y=0.62cm,baseline=(current bounding box.north)]
\node[tn] (fb3x1) at (0,-2) {};
\node[tn] (fb3x2) at (1,-1) {};
\node[tn] (fb3x3) at (2,0) {};
\node[tn] (fb3x4) at (3,-1) {};
\draw[te] (fb3x3) -- (fb3x2);
\draw[te] (fb3x2) -- (fb3x1);
\draw[te] (fb3x3) -- (fb3x4);
\node[plab] at (1.5,-4.75) {$h = 3$};
\end{tikzpicture}\hspace{0.5cm}\begin{tikzpicture}[x=0.42cm,y=0.62cm,baseline=(current bounding box.north)]
\node[tn] (fb4x1) at (0,-3) {};
\node[tn] (fb4x2) at (1,-2) {};
\node[tn] (fb4x3) at (2,-1) {};
\node[tn] (fb4x4) at (3,-2) {};
\node[tn] (fb4x5) at (4,0) {};
\node[tn] (fb4x6) at (5,-2) {};
\node[tn] (fb4x7) at (6,-1) {};
\draw[te] (fb4x5) -- (fb4x3);
\draw[te] (fb4x3) -- (fb4x2);
\draw[te] (fb4x2) -- (fb4x1);
\draw[te] (fb4x3) -- (fb4x4);
\draw[te] (fb4x5) -- (fb4x7);
\draw[te] (fb4x7) -- (fb4x6);
\node[plab] at (3.0,-4.75) {$h = 4$};
\end{tikzpicture}\hspace{0.5cm}\begin{tikzpicture}[x=0.42cm,y=0.62cm,baseline=(current bounding box.north)]
\node[tn] (fb5x1) at (0,-4) {};
\node[tn] (fb5x2) at (1,-3) {};
\node[tn] (fb5x3) at (2,-2) {};
\node[tn] (fb5x4) at (3,-3) {};
\node[tn] (fb5x5) at (4,-1) {};
\node[tn] (fb5x6) at (5,-3) {};
\node[tn] (fb5x7) at (6,-2) {};
\node[tn] (fb5x8) at (7,0) {};
\node[tn] (fb5x9) at (8,-3) {};
\node[tn] (fb5x10) at (9,-2) {};
\node[tn] (fb5x11) at (10,-1) {};
\node[tn] (fb5x12) at (11,-2) {};
\draw[te] (fb5x8) -- (fb5x5);
\draw[te] (fb5x5) -- (fb5x3);
\draw[te] (fb5x3) -- (fb5x2);
\draw[te] (fb5x2) -- (fb5x1);
\draw[te] (fb5x3) -- (fb5x4);
\draw[te] (fb5x5) -- (fb5x7);
\draw[te] (fb5x7) -- (fb5x6);
\draw[te] (fb5x8) -- (fb5x11);
\draw[te] (fb5x11) -- (fb5x10);
\draw[te] (fb5x10) -- (fb5x9);
\draw[te] (fb5x11) -- (fb5x12);
\node[plab] at (5.5,-4.75) {$h = 5$};
\end{tikzpicture}
\diagcaption{節點最少的 AVL 樹，節點數依序是 $1, 2, 4, 7, 12$。 每一棵都是根加上
前兩棵。}
\label{fig:bal-fib}
\end{center}

先用一個粗略的估計看出它是對數：$N(h) = N(h-1) + N(h-2) + 1 > 2N(h-2)$，每往下兩
層節點數至少減半，從 $N(1) = 1$、$N(2) = 2$ 往上推，得到 $N(h) \ge 2^{(h-1)/2}$。 所
以 $n \ge 2^{(h-1)/2}$，$h \le 2\lg n + 1$。 精確的常數要用費氏數。 令 $F_1 = F_2 =
1$、$F_k = F_{k-1} + F_{k-2}$，$\varphi = (1 + \sqrt5)/2 \approx 1.618$ 是
$\varphi^2 = \varphi + 1$ 的正根。

\begin{theorem}
\label{thm:bal-avl-height}
(1) $N(h) = F_{h+2} - 1$。 (2) 有 $n$ 個節點的 AVL 樹，高度
\[
h \le \log_\varphi(n + 1) = \frac{\lg(n+1)}{\lg \varphi} < 1.4405 \lg(n+1).
\]
\end{theorem}

\begin{proof}
(1) 對 $h$ 歸納。 $N(0) + 1 = 1 = F_2$，$N(1) + 1 = 2 = F_3$。 $h \ge 2$ 時，遞迴式
兩邊加 $1$：
\[
N(h) + 1 = \bigl(N(h-1) + 1\bigr) + \bigl(N(h-2) + 1\bigr) = F_{h+1} + F_h = F_{h+2}.
\]

(2) 先證明 $F_k \ge \varphi^{k-2}$（$k \ge 1$）。 $F_1 = 1 > \varphi^{-1}$，
$F_2 = 1 = \varphi^0$。 $k \ge 3$ 時，由歸納假設
\[
F_k = F_{k-1} + F_{k-2} \ge \varphi^{k-3} + \varphi^{k-4} = \varphi^{k-4}(\varphi + 1)
= \varphi^{k-4} \varphi^2 = \varphi^{k-2}.
\]
高度 $h$ 的 AVL 樹有 $n \ge N(h)$ 個節點，所以
$n + 1 \ge N(h) + 1 = F_{h+2} \ge \varphi^h$，兩邊取 $\log_\varphi$ 得到
$h \le \log_\varphi(n+1)$。 $1 / \lg \varphi \approx 1.44042$。
\end{proof}

和最矮的可能 $\lceil \lg(n+1) \rceil$ 比，AVL 樹最壞只高了大約 $44\%$。 一百萬個節
點的二元樹至少高 $20$；AVL 樹最高 $28$（觀念題第 13 題）。

考卷上常問「$n$ 個節點的 AVL 樹最高可能是多少」。 做法是把 $N(h)$ 的表列到超過 $n$
為止，答案是滿足 $N(h) \le n$ 的最大 $h$。 例如 $n = 100$：$N(9) = 88 \le 100 <
143 = N(10)$，最高是 $9$；最矮是 $\lceil \lg 101 \rceil = 7$。 定理的上界
$\log_\varphi 101 \approx 9.59$ 只是估計，要算確切的值還是查表。

\subsection{插入：四種失衡和旋轉}
\label{sec:bal-avl-insert}

插入分三步。

\begin{enumerate}
\item 照第~9~章的方法，把新鍵插成葉節點。
\item 從新節點的父節點往根走，更新路上每個節點的高度和平衡因子。 只有這條路徑上的
節點，子樹的內容變了，其他節點的高度不會變。
\item 在這條路上找\textbf{第一個}（離新節點最近的）平衡因子變成 $\pm 2$ 的節點，
叫它 $A$。 從 $A$ 往新節點走兩步，依照這兩步的方向分成四種情況，各用一種旋轉修
好。 沒有任何節點變成 $\pm 2$，就不用旋轉。
\end{enumerate}

\begin{center}
\begin{tabular}{cll}
\toprule
類型 & 新節點在哪裡 & 修正 \\
\midrule
LL & $A$ 的左子節點的左子樹 & 對 $A$ 右旋 \\
RR & $A$ 的右子節點的右子樹 & 對 $A$ 左旋 \\
LR & $A$ 的左子節點的右子樹 & 先對 $A$ 的左子節點左旋，再對 $A$ 右旋 \\
RL & $A$ 的右子節點的左子樹 & 先對 $A$ 的右子節點右旋，再對 $A$ 左旋 \\
\bottomrule
\end{tabular}
\end{center}

LL、RR 叫\textbf{單旋轉}，LR、RL 叫\textbf{雙旋轉}。 RR、RL 是 LL、LR 左右對調，下
面只分析 LL 和 LR。

\textbf{LL}。 設 $B$ 是 $A$ 的左子節點。 插入後 $A$ 是 $+2$，$B$ 的平衡因子只能是
$+1$：如果是 $0$，表示 $B$ 的高度沒有因為這次插入而增加，$A$ 就不會失衡；如果是
$\pm 2$，最近的失衡節點就是 $B$ 而不是 $A$。 所以插入前 $B$ 是 $0$，兩棵子樹
$T_1$、$T_2$ 同高，設為 $k$；$A$ 插入前是 $+1$，右子樹 $T_3$ 的高度是 $k$。 新節
點讓 $T_1$ 變成 $k + 1$。 對 $A$ 右旋之後，$A$ 的兩棵子樹是 $T_2$、$T_3$，都是
$k$；$B$ 的兩棵子樹是 $T_1$ 和 $A$，都是 $k + 1$。 兩個節點都變成 $0$，這一段的
高度是 $k + 2$，和插入\textbf{前}一樣。

\textbf{LR}。 新節點在 $B$ 的右子樹裡，設 $C$ 是 $B$ 的右子節點，$C$ 的兩棵子樹是
$T_2$、$T_3$。 同樣的理由，插入前 $B$、$C$ 都是 $0$，$T_1$ 高 $k$，$T_2$、$T_3$ 高
$k - 1$，$T_4$（$A$ 的右子樹）高 $k$。 新節點讓 $T_2$ 或 $T_3$ 其中一棵變成 $k$。
這時對 $A$ 右旋沒有用：$T_2$、$T_3$ 會跟著 $C$ 一起掛到 $A$ 底下，左邊變成右邊太
高。 正確的做法是讓 $C$ 升兩層：先對 $B$ 左旋（變成 LL 的樣子），再對 $A$ 右旋。
結果 $C$ 是根，左邊是 $B$（子樹 $T_1$、$T_2$），右邊是 $A$（子樹 $T_3$、$T_4$），
$B$、$A$ 的高度都是 $k + 1$，$C$ 是 $k + 2$，也和插入前一樣。

\begin{center}
\begin{tikzpicture}[x=0.8cm, y=0.82cm]
% LL 之前
\node[tn, tnA] (A) at (0,0) {$A$};
\node[tn] (B) at (-1.1,-1) {$B$};
\draw[te] (A) -- (B); \draw[te] (A) -- (1.0,-1.0); \draw[te] (B) -- (-1.7,-2); \draw[te] (B) -- (-0.5,-2);
\subtri[1.25]{-1.7}{-2}{$T_1$}
\subtri{-0.5}{-2}{$T_2$}
\subtri{1.0}{-1}{$T_3$}
\node[bfl] at (A.north east) {$+2$}; \node[bfl] at (B.north east) {$+1$};
\node[font=\scriptsize] at (-1.7,-3.5) {$k+1$};
\node[font=\scriptsize] at (-0.5,-3.15) {$k$};
\node[font=\scriptsize] at (1.0,-2.15) {$k$};
\fill[coreStroke] (-1.7,-3.15) circle (2.5pt);
\draw[tarr] (1.9,-1.2) -- node[above, font=\small] {右旋 $A$} (3.4,-1.2);
\begin{scope}[xshift=4.9cm]
\node[tn] (B2) at (0,0) {$B$};
\node[tn] (A2) at (1.1,-1) {$A$};
\draw[te] (B2) -- (A2); \draw[te] (B2) -- (-1.0,-1.0); \draw[te] (A2) -- (0.5,-2); \draw[te] (A2) -- (1.7,-2);
\subtri[1.25]{-1.0}{-1}{$T_1$}
\subtri{0.5}{-2}{$T_2$}
\subtri{1.7}{-2}{$T_3$}
\node[bfl] at (A2.north east) {$0$}; \node[bfl] at (B2.north east) {$0$};
\fill[coreStroke] (-1.0,-2.15) circle (2.5pt);
\end{scope}
\end{tikzpicture}\\*[0.6em]
\begin{tikzpicture}[x=0.8cm, y=0.82cm]
% LR 之前
\node[tn, tnA] (A) at (0,0) {$A$};
\node[tn] (B) at (-1.3,-1) {$B$};
\node[tn] (C) at (-0.6,-2) {$C$};
\draw[te] (A) -- (B); \draw[te] (A) -- (1.0,-1.0); \draw[te] (B) -- (-2.0,-2); \draw[te] (B) -- (C);
\draw[te] (C) -- (-1.1,-3); \draw[te] (C) -- (-0.1,-3);
\subtri[1.6]{-2.0}{-2}{$T_1$}
\subtri[0.9]{-1.1}{-3}{$T_2$}
\subtri[0.9]{-0.1}{-3}{$T_3$}
\subtri[1.6]{1.0}{-1}{$T_4$}
\node[bfl] at (A.north east) {$+2$}; \node[bfl] at (B.north east) {$-1$};
\node[font=\scriptsize] at (-2.0,-3.85) {$k$};
\node[font=\scriptsize] at (1.0,-2.85) {$k$};
\node[font=\scriptsize, align=center] at (-0.6,-4.35) {其中一棵 $k$，\\ 另一棵 $k-1$};
\draw[tarr] (1.9,-1.4) -- node[above, font=\small] {左旋 $B$} node[below, font=\small] {再右旋 $A$} (3.6,-1.4);
\begin{scope}[xshift=5.6cm]
\node[tn] (C2) at (0,0) {$C$};
\node[tn] (B2) at (-1.1,-1) {$B$};
\node[tn] (A2) at (1.1,-1) {$A$};
\draw[te] (C2) -- (B2); \draw[te] (C2) -- (A2);
\draw[te] (B2) -- (-1.6,-2); \draw[te] (B2) -- (-0.6,-2); \draw[te] (A2) -- (0.6,-2); \draw[te] (A2) -- (1.6,-2);
\subtri[1.6]{-1.6}{-2}{$T_1$}
\subtri[0.9]{-0.6}{-2}{$T_2$}
\subtri[0.9]{0.6}{-2}{$T_3$}
\subtri[1.6]{1.6}{-2}{$T_4$}
\node[bfl] at (C2.north east) {$0$};
\end{scope}
\end{tikzpicture}
\diagcaption{上：LL，藍點是新節點，落在 $T_1$。 下：LR，新節點落在 $T_2$ 或
$T_3$（$k = 0$ 時 $C$ 自己就是新節點）。 兩種修完之後，這一段的高度都回到插入
前的 $k + 2$。}
\label{fig:bal-ll-lr}
\end{center}

兩種情況修完，這一段的高度都和插入前一樣，所以 $A$ 上面的祖先，高度和平衡因子都
回到插入前的樣子，不會再失衡。 這證明了：

\begin{theorem}
\label{thm:bal-avl-insert}
AVL 樹插入一個鍵，最多做一次單旋轉或一次雙旋轉（也就是最多兩次旋轉）。 整個插入
的時間是 $O(\lg n)$。
\end{theorem}

雙旋轉之後 $A$、$B$ 的平衡因子是多少，要看新節點落在 $T_2$ 還是 $T_3$，觀念題第 5
題會整理。 程式不用記這張表：每個節點存高度，旋轉後重算就好。

\begin{pseudo}[label=alg:bal-avl-insert]{AVL 樹的插入，回傳新的根}
\Function{AVL-Insert}{$t, k$}
    \If{$t = \NIL$}
        \State \Return 新節點（$\mathit{key} = k$、$\mathit{height} = 1$、兩個子指標都是 \NIL）
    \EndIf
    \If{$k < t.\mathit{key}$}
        \State $t.\mathit{left} \gets \Call{AVL-Insert}{t.\mathit{left}, k}$
    \ElsIf{$k > t.\mathit{key}$}
        \State $t.\mathit{right} \gets \Call{AVL-Insert}{t.\mathit{right}, k}$
    \Else
        \State \Return $t$ \Comment{已經在樹裡}
    \EndIf
    \State $t.\mathit{height} \gets 1 + \max(h(t.\mathit{left}), h(t.\mathit{right}))$
    \State \Return \Call{Rebalance}{$t$}
\EndFunction
\end{pseudo}

\begin{pseudo}[label=alg:bal-rebalance]{檢查 $t$ 是否失衡，失衡就旋轉}
\Function{Rebalance}{$t$}
    \If{$\mathrm{BF}(t) = 2$}
        \If{$\mathrm{BF}(t.\mathit{left}) < 0$} \Comment{LR：先轉成 LL}
            \State $t.\mathit{left} \gets \Call{Rotate-Left}{t.\mathit{left}}$
        \EndIf
        \State \Return \Call{Rotate-Right}{$t$}
    \ElsIf{$\mathrm{BF}(t) = -2$}
        \If{$\mathrm{BF}(t.\mathit{right}) > 0$} \Comment{RL：先轉成 RR}
            \State $t.\mathit{right} \gets \Call{Rotate-Right}{t.\mathit{right}}$
        \EndIf
        \State \Return \Call{Rotate-Left}{$t$}
    \EndIf
    \State \Return $t$
\EndFunction
\end{pseudo}

$h(\NIL) = 0$。 這裡的 \proc{Rotate-Right}、\proc{Rotate-Left} 是虛擬碼~\ref{alg:bal-rotate-right}
加上最後兩步：先重算降下去的節點的高度，再重算升上來的節點（它的高度要用到前者）。

遞迴版本在回程時對路徑上的每個節點都呼叫 \proc{Rebalance}，第一個失衡的節點就是
最深的那個，也就是上面說的 $A$。 修完之後高度回到插入前，上面的節點自然不會再旋
轉。 \proc{Rebalance} 第 3 行用的是嚴格的 $< 0$：插入時 $A$ 的子節點不會是 $0$，
但刪除時會，那時必須做單旋轉（第~\ref{sec:bal-avl-delete}~節）。

\subsection{手算追蹤：依序登錄會員卡}
\label{sec:bal-avl-trace}

健身房把十張紙本會員卡登錄進系統，依收到卡片的順序插入 AVL 樹，編號是
\[
52,\ 41,\ 16,\ 34,\ 89,\ 94,\ 64,\ 46,\ 56,\ 58 .
\]
下面每張圖的節點右上方是平衡因子，藍底是剛插入的節點，橘框是第一個失衡的節點
$A$。 沒有引起旋轉的插入合併成一張。

\begin{center}
\begin{tikzpicture}[x=0.5cm,y=0.72cm,baseline=(current bounding box.north)]
\node[tn] (aia41) at (0,-1) {41};
\node[tn] (aia52) at (1,0) {52};
\draw[te] (aia52) -- (aia41);
\node[bfl] at (aia41.north east) {$0$};
\node[bfl] at (aia52.north east) {$+1$};
\node[plab] at (0.5,-2.75) {(a) 插入 52、41};
\end{tikzpicture}\hspace{0.8cm}\begin{tikzpicture}[x=0.5cm,y=0.72cm,baseline=(current bounding box.north)]
\node[tn,tnN] (aib16) at (0,-2) {16};
\node[tn] (aib41) at (1,-1) {41};
\node[tn,tnA] (aib52) at (2,0) {52};
\draw[te] (aib52) -- (aib41);
\draw[te] (aib41) -- (aib16);
\node[bfl] at (aib16.north east) {$0$};
\node[bfl] at (aib41.north east) {$+1$};
\node[bfl] at (aib52.north east) {$+2$};
\node[plab] at (1.0,-2.75) {(b) 插入 16：LL};
\end{tikzpicture}\hspace{0.8cm}\begin{tikzpicture}[x=0.5cm,y=0.72cm,baseline=(current bounding box.north)]
\node[tn] (aic16) at (0,-1) {16};
\node[tn] (aic41) at (1,0) {41};
\node[tn] (aic52) at (2,-1) {52};
\draw[te] (aic41) -- (aic16);
\draw[te] (aic41) -- (aic52);
\node[bfl] at (aic16.north east) {$0$};
\node[bfl] at (aic41.north east) {$0$};
\node[bfl] at (aic52.north east) {$0$};
\node[plab] at (1.0,-2.75) {(c) 右旋之後};
\end{tikzpicture}\\*[0.6em]
\begin{tikzpicture}[x=0.5cm,y=0.72cm,baseline=(current bounding box.north)]
\node[tn] (aid16) at (0,-1) {16};
\node[tn] (aid34) at (1,-2) {34};
\node[tn] (aid41) at (2,0) {41};
\node[tn] (aid52) at (3,-1) {52};
\node[tn] (aid89) at (4,-2) {89};
\draw[te] (aid41) -- (aid16);
\draw[te] (aid16) -- (aid34);
\draw[te] (aid41) -- (aid52);
\draw[te] (aid52) -- (aid89);
\node[bfl] at (aid16.north east) {$-1$};
\node[bfl] at (aid34.north east) {$0$};
\node[bfl] at (aid41.north east) {$0$};
\node[bfl] at (aid52.north east) {$-1$};
\node[bfl] at (aid89.north east) {$0$};
\node[plab] at (2.0,-3.75) {(d) 插入 34、89};
\end{tikzpicture}\hspace{0.6cm}\begin{tikzpicture}[x=0.5cm,y=0.72cm,baseline=(current bounding box.north)]
\node[tn] (aie16) at (0,-1) {16};
\node[tn] (aie34) at (1,-2) {34};
\node[tn] (aie41) at (2,0) {41};
\node[tn,tnA] (aie52) at (3,-1) {52};
\node[tn] (aie89) at (4,-2) {89};
\node[tn,tnN] (aie94) at (5,-3) {94};
\draw[te] (aie41) -- (aie16);
\draw[te] (aie16) -- (aie34);
\draw[te] (aie41) -- (aie52);
\draw[te] (aie52) -- (aie89);
\draw[te] (aie89) -- (aie94);
\node[bfl] at (aie16.north east) {$-1$};
\node[bfl] at (aie34.north east) {$0$};
\node[bfl] at (aie41.north east) {$-1$};
\node[bfl] at (aie52.north east) {$-2$};
\node[bfl] at (aie89.north east) {$-1$};
\node[bfl] at (aie94.north east) {$0$};
\node[plab] at (2.5,-3.75) {(e) 插入 94：RR};
\end{tikzpicture}\hspace{0.6cm}\begin{tikzpicture}[x=0.5cm,y=0.72cm,baseline=(current bounding box.north)]
\node[tn] (aif16) at (0,-1) {16};
\node[tn] (aif34) at (1,-2) {34};
\node[tn] (aif41) at (2,0) {41};
\node[tn] (aif52) at (3,-2) {52};
\node[tn] (aif89) at (4,-1) {89};
\node[tn] (aif94) at (5,-2) {94};
\draw[te] (aif41) -- (aif16);
\draw[te] (aif16) -- (aif34);
\draw[te] (aif41) -- (aif89);
\draw[te] (aif89) -- (aif52);
\draw[te] (aif89) -- (aif94);
\node[bfl] at (aif16.north east) {$-1$};
\node[bfl] at (aif34.north east) {$0$};
\node[bfl] at (aif41.north east) {$0$};
\node[bfl] at (aif52.north east) {$0$};
\node[bfl] at (aif89.north east) {$0$};
\node[bfl] at (aif94.north east) {$0$};
\node[plab] at (2.5,-3.75) {(f) 左旋之後};
\end{tikzpicture}
\diagcaption{插入 $52$ 到 $94$。 (b) 從 $16$ 往上，$41$ 變 $+1$，$52$ 變 $+2$，兩步
都往左，是 LL。 (e) $89$ 變 $-1$，$52$ 變 $-2$，兩步都往右，是 RR。 根 $41$ 是 $-1$，沒有失衡。}
\label{fig:bal-avl-trace1}
\end{center}

插入 $64$、$46$ 都沒有失衡：$64$ 讓 $52$ 變 $-1$、$89$ 變 $+1$、$41$ 變 $-1$；
$46$ 讓 $52$ 回到 $0$，$52$ 的高度沒變，再往上的節點都不受影響。

\begin{center}
\begin{tikzpicture}[x=0.5cm,y=0.72cm,baseline=(current bounding box.north)]
\node[tn] (aig16) at (0,-1) {16};
\node[tn] (aig34) at (1,-2) {34};
\node[tn] (aig41) at (2,0) {41};
\node[tn] (aig46) at (3,-3) {46};
\node[tn] (aig52) at (4,-2) {52};
\node[tn] (aig64) at (5,-3) {64};
\node[tn] (aig89) at (6,-1) {89};
\node[tn] (aig94) at (7,-2) {94};
\draw[te] (aig41) -- (aig16);
\draw[te] (aig16) -- (aig34);
\draw[te] (aig41) -- (aig89);
\draw[te] (aig89) -- (aig52);
\draw[te] (aig52) -- (aig46);
\draw[te] (aig52) -- (aig64);
\draw[te] (aig89) -- (aig94);
\node[bfl] at (aig16.north east) {$-1$};
\node[bfl] at (aig34.north east) {$0$};
\node[bfl] at (aig41.north east) {$-1$};
\node[bfl] at (aig46.north east) {$0$};
\node[bfl] at (aig52.north east) {$0$};
\node[bfl] at (aig64.north east) {$0$};
\node[bfl] at (aig89.north east) {$+1$};
\node[bfl] at (aig94.north east) {$0$};
\node[plab] at (3.5,-4.75) {(g) 插入 64、46};
\end{tikzpicture}\hspace{0.4cm}\begin{tikzpicture}[x=0.5cm,y=0.72cm,baseline=(current bounding box.north)]
\node[tn] (aih16) at (0,-1) {16};
\node[tn] (aih34) at (1,-2) {34};
\node[tn] (aih41) at (2,0) {41};
\node[tn] (aih46) at (3,-3) {46};
\node[tn] (aih52) at (4,-2) {52};
\node[tn,tnN] (aih56) at (5,-4) {56};
\node[tn] (aih64) at (6,-3) {64};
\node[tn,tnA] (aih89) at (7,-1) {89};
\node[tn] (aih94) at (8,-2) {94};
\draw[te] (aih41) -- (aih16);
\draw[te] (aih16) -- (aih34);
\draw[te] (aih41) -- (aih89);
\draw[te] (aih89) -- (aih52);
\draw[te] (aih52) -- (aih46);
\draw[te] (aih52) -- (aih64);
\draw[te] (aih64) -- (aih56);
\draw[te] (aih89) -- (aih94);
\node[bfl] at (aih16.north east) {$-1$};
\node[bfl] at (aih34.north east) {$0$};
\node[bfl] at (aih41.north east) {$-2$};
\node[bfl] at (aih46.north east) {$0$};
\node[bfl] at (aih52.north east) {$-1$};
\node[bfl] at (aih56.north east) {$0$};
\node[bfl] at (aih64.north east) {$+1$};
\node[bfl] at (aih89.north east) {$+2$};
\node[bfl] at (aih94.north east) {$0$};
\node[plab] at (4.0,-4.75) {(h) 插入 56：LR};
\end{tikzpicture}\hspace{0.4cm}\begin{tikzpicture}[x=0.5cm,y=0.72cm,baseline=(current bounding box.north)]
\node[tn] (aii16) at (0,-1) {16};
\node[tn] (aii34) at (1,-2) {34};
\node[tn] (aii41) at (2,0) {41};
\node[tn] (aii46) at (3,-3) {46};
\node[tn] (aii52) at (4,-2) {52};
\node[tn] (aii56) at (5,-3) {56};
\node[tn] (aii64) at (6,-1) {64};
\node[tn] (aii89) at (7,-2) {89};
\node[tn] (aii94) at (8,-3) {94};
\draw[te] (aii41) -- (aii16);
\draw[te] (aii16) -- (aii34);
\draw[te] (aii41) -- (aii64);
\draw[te] (aii64) -- (aii52);
\draw[te] (aii52) -- (aii46);
\draw[te] (aii52) -- (aii56);
\draw[te] (aii64) -- (aii89);
\draw[te] (aii89) -- (aii94);
\node[bfl] at (aii16.north east) {$-1$};
\node[bfl] at (aii34.north east) {$0$};
\node[bfl] at (aii41.north east) {$-1$};
\node[bfl] at (aii46.north east) {$0$};
\node[bfl] at (aii52.north east) {$0$};
\node[bfl] at (aii56.north east) {$0$};
\node[bfl] at (aii64.north east) {$0$};
\node[bfl] at (aii89.north east) {$-1$};
\node[bfl] at (aii94.north east) {$0$};
\node[plab] at (4.0,-4.75) {(i) 雙旋轉之後};
\end{tikzpicture}
\diagcaption{插入 $56$。 從 $56$ 往上：$64$ 變 $+1$，$52$ 變 $-1$，$89$ 變 $+2$。 從
$89$ 往 $56$ 先左（$52$）再右（$64$），是 LR：先對 $52$ 左旋，再對 $89$ 右旋，
$64$ 升到這一段的根。 $64$ 原本的左子節點 $56$ 換到 $52$ 的右邊。}
\label{fig:bal-avl-trace2}
\end{center}

\begin{center}
\begin{tikzpicture}[x=0.5cm,y=0.72cm,baseline=(current bounding box.north)]
\node[tn] (aij16) at (0,-1) {16};
\node[tn] (aij34) at (1,-2) {34};
\node[tn,tnA] (aij41) at (2,0) {41};
\node[tn] (aij46) at (3,-3) {46};
\node[tn] (aij52) at (4,-2) {52};
\node[tn] (aij56) at (5,-3) {56};
\node[tn,tnN] (aij58) at (6,-4) {58};
\node[tn] (aij64) at (7,-1) {64};
\node[tn] (aij89) at (8,-2) {89};
\node[tn] (aij94) at (9,-3) {94};
\draw[te] (aij41) -- (aij16);
\draw[te] (aij16) -- (aij34);
\draw[te] (aij41) -- (aij64);
\draw[te] (aij64) -- (aij52);
\draw[te] (aij52) -- (aij46);
\draw[te] (aij52) -- (aij56);
\draw[te] (aij56) -- (aij58);
\draw[te] (aij64) -- (aij89);
\draw[te] (aij89) -- (aij94);
\node[bfl] at (aij16.north east) {$-1$};
\node[bfl] at (aij34.north east) {$0$};
\node[bfl] at (aij41.north east) {$-2$};
\node[bfl] at (aij46.north east) {$0$};
\node[bfl] at (aij52.north east) {$-1$};
\node[bfl] at (aij56.north east) {$-1$};
\node[bfl] at (aij58.north east) {$0$};
\node[bfl] at (aij64.north east) {$+1$};
\node[bfl] at (aij89.north east) {$-1$};
\node[bfl] at (aij94.north east) {$0$};
\node[plab] at (4.5,-4.75) {(j) 插入 58：RL};
\end{tikzpicture}\hspace{1.0cm}\begin{tikzpicture}[x=0.5cm,y=0.72cm,baseline=(current bounding box.north)]
\node[tn] (aik16) at (0,-2) {16};
\node[tn] (aik34) at (1,-3) {34};
\node[tn] (aik41) at (2,-1) {41};
\node[tn] (aik46) at (3,-2) {46};
\node[tn] (aik52) at (4,0) {52};
\node[tn] (aik56) at (5,-2) {56};
\node[tn] (aik58) at (6,-3) {58};
\node[tn] (aik64) at (7,-1) {64};
\node[tn] (aik89) at (8,-2) {89};
\node[tn] (aik94) at (9,-3) {94};
\draw[te] (aik52) -- (aik41);
\draw[te] (aik41) -- (aik16);
\draw[te] (aik16) -- (aik34);
\draw[te] (aik41) -- (aik46);
\draw[te] (aik52) -- (aik64);
\draw[te] (aik64) -- (aik56);
\draw[te] (aik56) -- (aik58);
\draw[te] (aik64) -- (aik89);
\draw[te] (aik89) -- (aik94);
\node[bfl] at (aik16.north east) {$-1$};
\node[bfl] at (aik34.north east) {$0$};
\node[bfl] at (aik41.north east) {$+1$};
\node[bfl] at (aik46.north east) {$0$};
\node[bfl] at (aik52.north east) {$0$};
\node[bfl] at (aik56.north east) {$-1$};
\node[bfl] at (aik58.north east) {$0$};
\node[bfl] at (aik64.north east) {$0$};
\node[bfl] at (aik89.north east) {$-1$};
\node[bfl] at (aik94.north east) {$0$};
\node[plab] at (4.5,-4.75) {(k) 雙旋轉之後};
\end{tikzpicture}
\diagcaption{插入 $58$。 $56$ 變 $-1$，$52$ 變 $-1$，$64$ 變 $+1$，根 $41$ 變
$-2$。 從 $41$ 先右（$64$）再左（$52$），是 RL：先對 $64$ 右旋，再對 $41$ 左旋，
$52$ 升到根。 $52$ 原本的左子樹 $46$ 接到 $41$ 的右邊，右子樹 $56$ 接到 $64$ 的
左邊。}
\label{fig:bal-avl-trace3}
\end{center}

整理成表：

\begin{center}
\begin{tabular}{cccl}
\toprule
插入 & 失衡節點 $A$ & 類型 & 修完後這一段的根 \\
\midrule
16 & 52 & LL & 41 \\
94 & 52 & RR & 89 \\
56 & 89 & LR & 64 \\
58 & 41 & RL & 52 \\
\bottomrule
\end{tabular}
\end{center}

圖~\ref{fig:bal-avl-trace3}(k) 的最後一步最容易畫錯。 雙旋轉的結果可以不畫中間那一
步，直接寫：$A$、$B$、$C$ 三個節點依鍵的大小排好，中間的 $C$ 當根，小的當左子節點，
大的當右子節點；原本掛在這三個節點底下的四棵子樹，依中序順序從左到右接上去。 這裡
三個節點是 $41 < 52 < 64$，四棵子樹依序是 $\{16, 34\}$、$\{46\}$、$\{56, 58\}$、
$\{89, 94\}$。 單旋轉也可以這樣想，只是三個節點是 $A$、$B$ 和 $B$ 往同一邊的子節點。

\begin{examnote}
AVL 的追蹤題，每次旋轉都要畫出旋轉後的整棵樹，並寫出失衡的節點和類型（LL、RR、
LR、RL）。 沒有旋轉的插入可以幾個合併成一張。 判斷類型時，從\textbf{最靠近新節
點}的失衡節點往下看兩步，不是從根看。 最後檢查一次：中序是不是由小到大，每個節點
的平衡因子是不是都在 $-1$ 到 $1$。
\end{examnote}

\subsection{C 程式}
\label{sec:bal-avl-c}

C 程式照虛擬碼~\ref{alg:bal-avl-insert}、\ref{alg:bal-rebalance} 寫。 每個節點多存
一個 \texttt{height}，平衡因子要用時再算。 遞迴函式回傳這棵子樹新的根，呼叫的地方
寫 \texttt{t->left = avl\_insert(t->left, key, rot)}，旋轉換了根也接得回去。
\texttt{Rotation} 結構只是用來把旋轉的類型和位置帶回 \texttt{main} 印出來。

\begin{cbookrecap}[第~14~章 14.5 節、第~15~章、第~19~章 19.10、19.15、19.17 節]
\texttt{t->left} 是 \texttt{(*t).left} 的簡寫，\texttt{t} 是指向結構的指標。 遞迴
函式回傳 \texttt{Node *}，呼叫端把回傳值接回原本的指標，是串列和樹最常見的寫法：
不用傳指標的指標，空樹的情況（\texttt{t == NULL}）也不用另外處理。 釋放整棵樹要
先釋放兩棵子樹，最後才釋放自己（後序），順序反過來就會讀到已經釋放的記憶體。
\end{cbookrecap}

\begin{example}[label=ex:bal-avl-insert]{c}{(avl\_insert.c) AVL 樹的插入}
#include <stdio.h>
#include <stdlib.h>

typedef struct Node {
    int key;
    int height;                 // 這棵子樹的高度，只有一個節點時是 1
    struct Node *left, *right;
} Node;

typedef struct {
    const char *type;           // "LL"、"RR"、"LR"、"RL"；沒有旋轉時是 NULL
    int at;                     // 失衡節點的鍵
} Rotation;

static int height(const Node *t)
{
    return t == NULL ? 0 : t->height;
}

static void update(Node *t)
{
    int hl = height(t->left), hr = height(t->right);
    t->height = 1 + (hl > hr ? hl : hr);
}

static int balance(const Node *t)       // 平衡因子：左高減右高
{
    return height(t->left) - height(t->right);
}

static Node *rotate_right(Node *a)      // a 的左子節點 b 轉上來
{
    Node *b = a->left;
    a->left = b->right;
    b->right = a;
    update(a);                          // a 現在在 b 下面，先更新 a
    update(b);
    return b;
}

static Node *rotate_left(Node *a)       // a 的右子節點 b 轉上來
{
    Node *b = a->right;
    a->right = b->left;
    b->left = a;
    update(a);
    update(b);
    return b;
}

// 檢查 t 是否失衡；失衡就旋轉，回傳這棵子樹新的根。
static Node *rebalance(Node *t, Rotation *rot)
{
    int b = balance(t);
    if (b == 2) {                               // 左邊太高
        rot->at = t->key;
        if (balance(t->left) < 0) {
            rot->type = "LR";
            t->left = rotate_left(t->left);
        } else {
            rot->type = "LL";
        }
        return rotate_right(t);
    }
    if (b == -2) {                              // 右邊太高
        rot->at = t->key;
        if (balance(t->right) > 0) {
            rot->type = "RL";
            t->right = rotate_right(t->right);
        } else {
            rot->type = "RR";
        }
        return rotate_left(t);
    }
    return t;
}

// 把 key 插進以 t 為根的 AVL 樹，回傳新的根。key 已經在樹裡就不插。
Node *avl_insert(Node *t, int key, Rotation *rot)
{
    if (t == NULL) {
        Node *x = malloc(sizeof *x);
        if (x == NULL) {
            exit(1);
        }
        x->key = key;
        x->height = 1;
        x->left = x->right = NULL;
        return x;
    }
    if (key < t->key) {
        t->left = avl_insert(t->left, key, rot);
    } else if (key > t->key) {
        t->right = avl_insert(t->right, key, rot);
    } else {
        return t;
    }
    update(t);
    return rebalance(t, rot);
}

// 用「根(左子樹,右子樹)」的格式印出，空的子樹印成 .
void print_tree(const Node *t)
{
    if (t == NULL) {
        putchar('.');
        return;
    }
    printf("%d", t->key);
    if (t->left != NULL || t->right != NULL) {
        putchar('(');
        print_tree(t->left);
        putchar(',');
        print_tree(t->right);
        putchar(')');
    }
}

void free_tree(Node *t)
{
    if (t != NULL) {
        free_tree(t->left);
        free_tree(t->right);
        free(t);
    }
}

int main(void)
{
    int keys[] = {52, 41, 16, 34, 89, 94, 64, 46, 56, 58};
    int n = sizeof keys / sizeof keys[0];
    Node *root = NULL;
    for (int i = 0; i < n; i++) {
        Rotation rot = {NULL, 0};
        root = avl_insert(root, keys[i], &rot);
        printf("insert %d: ", keys[i]);
        if (rot.type != NULL) {
            printf("%s at %d, ", rot.type, rot.at);
        }
        print_tree(root);
        putchar('\n');
    }
    printf("height = %d\n", height(root));
    free_tree(root);

    // 會員編號由小到大插入：1, 2, ..., m
    for (int m = 1000; m <= 1000000; m *= 10) {
        Rotation rot = {NULL, 0};
        root = NULL;
        for (int k = 1; k <= m; k++) {
            root = avl_insert(root, k, &rot);
        }
        int lg = 0;                     // lg = ceil(lg(m + 1))，m 個節點的最小高度
        while ((1L << lg) < m + 1) {
            lg++;
        }
        printf("sorted m = %d: height = %d, minimum = %d\n", m, height(root), lg);
        free_tree(root);
    }
    return 0;
}
\end{example}
\iosample{}{insert 52: 52 \\ insert 41: 52(41,.) \\ insert 16: LL at 52, 41(16,52) \\ insert 34: 41(16(.,34),52) \\ insert 89: 41(16(.,34),52(.,89)) \\ insert 94: RR at 52, 41(16(.,34),89(52,94)) \\ insert 64: 41(16(.,34),89(52(.,64),94)) \\ insert 46: 41(16(.,34),89(52(46,64),94)) \\ insert 56: LR at 89, 41(16(.,34),64(52(46,56),89(.,94))) \\ insert 58: RL at 41, 52(41(16(.,34),46),64(56(.,58),89(.,94))) \\ height = 4 \\ sorted m = 1000: height = 10, minimum = 10 \\ sorted m = 10000: height = 14, minimum = 14 \\ sorted m = 100000: height = 17, minimum = 17 \\ sorted m = 1000000: height = 20, minimum = 20}

前十行和第~\ref{sec:bal-avl-trace}~節的圖一致。 後四行把 $1, 2, \ldots, m$ 依序插
入：普通的二元搜尋樹會長成高度 $m$ 的鏈，AVL 樹的高度剛好是最小可能值
$\lceil \lg(m+1) \rceil$。 依序插入時，需要旋轉的都是 RR，最後得到的樹非常整齊；一般的輸入不會這麼好，但高度不會超過定理~\ref{thm:bal-avl-height} 的上界。

\texttt{rotate\_right} 裡先 \texttt{update(a)} 再 \texttt{update(b)}：旋轉後 \texttt{a}
是 \texttt{b} 的子節點，\texttt{b} 的高度要用到 \texttt{a} 新的高度。 這兩行對調，
\texttt{b} 會用舊的高度算，之後的平衡因子全部算錯，而且不會當掉，很難查。

\subsection{刪除}
\label{sec:bal-avl-delete}

刪除也分三步。

\begin{enumerate}
\item 照第~9~章的方法刪除。 要刪的節點有兩個子節點時，用右子樹最小的節點（中序的
\textbf{後繼}）的鍵取代它，再去右子樹刪掉那個後繼。 真正被拿掉的節點最多只有一個
子節點。
\item 從被拿掉的節點的父節點往根走，更新高度和平衡因子。
\item 遇到平衡因子變成 $\pm 2$ 的節點 $A$，就旋轉。 \textbf{旋轉之後要繼續往上檢查}。
\end{enumerate}

刪除和插入不一樣的地方在第 3 步。 假設 $A$ 變成 $+2$（右邊少了一層），$B$ 是 $A$
的左子節點，也就是比較高的那一邊。 插入時 $B$ 一定是 $\pm 1$，刪除時 $B$ 可能是
$0$，因為 $B$ 這邊沒有動過。 三種情況：

\begin{center}
\begin{tabular}{cclc}
\toprule
$\mathrm{BF}(B)$ & 有些書的名稱 & 修正 & 這一段的高度 \\
\midrule
$+1$ & R1 & 對 $A$ 右旋（同 LL） & 少 $1$，要繼續往上 \\
$0$ & R0 & 對 $A$ 右旋（同 LL） & 不變，可以停 \\
$-1$ & R$-1$ & 先左旋 $B$，再右旋 $A$（同 LR） & 少 $1$，要繼續往上 \\
\bottomrule
\end{tabular}
\end{center}

名稱裡的 R 表示刪除發生在 $A$ 的右子樹，數字是 $\mathrm{BF}(B)$。 $A$ 變成 $-2$ 的
情況左右對調，叫 L0、L1、L$-1$，修正是 RR 或 RL。

$\mathrm{BF}(B) = 0$ 為什麼只能單旋轉、而且高度不變？ 設刪除後 $A$ 的右子樹高
$k$，$B$ 的兩棵子樹都高 $k + 1$，$A$ 原本高 $k + 3$。 對 $A$ 右旋後，$A$ 的子樹是
$B$ 原本的右子樹（$k + 1$）和 $A$ 的右子樹（$k$），$A$ 高 $k + 2$、平衡因子 $+1$；
$B$ 的子樹是 $k + 1$ 和 $k + 2$，平衡因子 $-1$，高度 $k + 3$，和刪除前一樣。 如果
這時誤用雙旋轉，$C$（$B$ 的右子節點，高 $k + 1$）升上來，$B$ 降成 $C$ 的左子節點：
$B$ 的左子樹（$k + 1$）留在 $B$ 底下，右邊換成 $C$ 原本的左子樹，只有 $k$ 或 $k - 1$
高；是 $k - 1$ 時，$B$ 的平衡因子變成 $+2$，又失衡了。 這就是 \proc{Rebalance} 第 3 行要用
$< 0$ 的原因：$\mathrm{BF}(B) = 0$ 時走單旋轉。

$\mathrm{BF}(B) = \pm 1$ 時，計算和插入的 LL、LR 一樣，修完的高度是 $k + 2$，比刪除
前的 $k + 3$ 少一層，$A$ 的父節點可能因此失衡。 所以刪除一個鍵，可能從底下一路旋轉
到根，最多 $O(\lg n)$ 次旋轉。 每次旋轉是 $\Theta(1)$，路徑長是 $O(\lg n)$，總時間
仍然是 $O(\lg n)$。

\begin{pseudo}[label=alg:bal-avl-delete]{AVL 樹的刪除，回傳新的根}
\Function{AVL-Delete}{$t, k$}
    \If{$t = \NIL$}
        \State \Return \NIL \Comment{$k$ 不在樹裡}
    \EndIf
    \If{$k < t.\mathit{key}$}
        \State $t.\mathit{left} \gets \Call{AVL-Delete}{t.\mathit{left}, k}$
    \ElsIf{$k > t.\mathit{key}$}
        \State $t.\mathit{right} \gets \Call{AVL-Delete}{t.\mathit{right}, k}$
    \ElsIf{$t.\mathit{left} = \NIL$ \OR $t.\mathit{right} = \NIL$}
        \State $c \gets$ $t$ 不是 \NIL 的那個子節點（兩個都是 \NIL 時 $c = \NIL$）
        \State 釋放 $t$
        \State \Return $c$
    \Else
        \State $s \gets t.\mathit{right}$ 一路往左走到底的節點 \Comment{後繼}
        \State $t.\mathit{key} \gets s.\mathit{key}$
        \State $t.\mathit{right} \gets \Call{AVL-Delete}{t.\mathit{right}, s.\mathit{key}}$
    \EndIf
    \State $t.\mathit{height} \gets 1 + \max(h(t.\mathit{left}), h(t.\mathit{right}))$
    \State \Return \Call{Rebalance}{$t$}
\EndFunction
\end{pseudo}

第 14 行只複製鍵；節點如果還存了別的資料（例如會員姓名），也要一起複製。 和插入
不同，這裡每一層回程的 \proc{Rebalance} 都可能真的旋轉。

\subsection{刪除的手算追蹤}
\label{sec:bal-avl-deltrace}

\textbf{例 1}。 從圖~\ref{fig:bal-avl-trace3}(k) 的樹開始，依序刪除
$41, 16, 34, 46$。 兩個子節點的節點用後繼取代。

\begin{center}
\begin{tikzpicture}[x=0.5cm,y=0.72cm,baseline=(current bounding box.north)]
\node[tn] (adb16) at (0,-2) {16};
\node[tn] (adb34) at (1,-3) {34};
\node[tn,tnA] (adb46) at (2,-1) {46};
\node[tn] (adb52) at (3,0) {52};
\node[tn] (adb56) at (4,-2) {56};
\node[tn] (adb58) at (5,-3) {58};
\node[tn] (adb64) at (6,-1) {64};
\node[tn] (adb89) at (7,-2) {89};
\node[tn] (adb94) at (8,-3) {94};
\draw[te] (adb52) -- (adb46);
\draw[te] (adb46) -- (adb16);
\draw[te] (adb16) -- (adb34);
\draw[te] (adb52) -- (adb64);
\draw[te] (adb64) -- (adb56);
\draw[te] (adb56) -- (adb58);
\draw[te] (adb64) -- (adb89);
\draw[te] (adb89) -- (adb94);
\node[bfl] at (adb16.north east) {$-1$};
\node[bfl] at (adb34.north east) {$0$};
\node[bfl] at (adb46.north east) {$+2$};
\node[bfl] at (adb52.north east) {$0$};
\node[bfl] at (adb56.north east) {$-1$};
\node[bfl] at (adb58.north east) {$0$};
\node[bfl] at (adb64.north east) {$0$};
\node[bfl] at (adb89.north east) {$-1$};
\node[bfl] at (adb94.north east) {$0$};
\node[plab] at (4.0,-3.75) {(a) 刪 41：46 補上，46 失衡};
\end{tikzpicture}\hspace{0.5cm}\begin{tikzpicture}[x=0.5cm,y=0.72cm,baseline=(current bounding box.north)]
\node[tn] (adc16) at (0,-2) {16};
\node[tn] (adc34) at (1,-1) {34};
\node[tn] (adc46) at (2,-2) {46};
\node[tn] (adc52) at (3,0) {52};
\node[tn] (adc56) at (4,-2) {56};
\node[tn] (adc58) at (5,-3) {58};
\node[tn] (adc64) at (6,-1) {64};
\node[tn] (adc89) at (7,-2) {89};
\node[tn] (adc94) at (8,-3) {94};
\draw[te] (adc52) -- (adc34);
\draw[te] (adc34) -- (adc16);
\draw[te] (adc34) -- (adc46);
\draw[te] (adc52) -- (adc64);
\draw[te] (adc64) -- (adc56);
\draw[te] (adc56) -- (adc58);
\draw[te] (adc64) -- (adc89);
\draw[te] (adc89) -- (adc94);
\node[bfl] at (adc16.north east) {$0$};
\node[bfl] at (adc34.north east) {$0$};
\node[bfl] at (adc46.north east) {$0$};
\node[bfl] at (adc52.north east) {$-1$};
\node[bfl] at (adc56.north east) {$-1$};
\node[bfl] at (adc58.north east) {$0$};
\node[bfl] at (adc64.north east) {$0$};
\node[bfl] at (adc89.north east) {$-1$};
\node[bfl] at (adc94.north east) {$0$};
\node[plab] at (4.0,-3.75) {(b) 雙旋轉（LR）之後};
\end{tikzpicture}\hspace{0.5cm}\begin{tikzpicture}[x=0.5cm,y=0.72cm,baseline=(current bounding box.north)]
\node[tn] (ade46) at (0,-1) {46};
\node[tn,tnA] (ade52) at (1,0) {52};
\node[tn] (ade56) at (2,-2) {56};
\node[tn] (ade58) at (3,-3) {58};
\node[tn] (ade64) at (4,-1) {64};
\node[tn] (ade89) at (5,-2) {89};
\node[tn] (ade94) at (6,-3) {94};
\draw[te] (ade52) -- (ade46);
\draw[te] (ade52) -- (ade64);
\draw[te] (ade64) -- (ade56);
\draw[te] (ade56) -- (ade58);
\draw[te] (ade64) -- (ade89);
\draw[te] (ade89) -- (ade94);
\node[bfl] at (ade46.north east) {$0$};
\node[bfl] at (ade52.north east) {$-2$};
\node[bfl] at (ade56.north east) {$-1$};
\node[bfl] at (ade58.north east) {$0$};
\node[bfl] at (ade64.north east) {$0$};
\node[bfl] at (ade89.north east) {$-1$};
\node[bfl] at (ade94.north east) {$0$};
\node[plab] at (3.0,-3.75) {(c) 刪 16、34：52 失衡};
\end{tikzpicture}\\*[0.6em]
\begin{tikzpicture}[x=0.5cm,y=0.72cm,baseline=(current bounding box.north)]
\node[tn] (adf46) at (0,-2) {46};
\node[tn] (adf52) at (1,-1) {52};
\node[tn] (adf56) at (2,-2) {56};
\node[tn] (adf58) at (3,-3) {58};
\node[tn] (adf64) at (4,0) {64};
\node[tn] (adf89) at (5,-1) {89};
\node[tn] (adf94) at (6,-2) {94};
\draw[te] (adf64) -- (adf52);
\draw[te] (adf52) -- (adf46);
\draw[te] (adf52) -- (adf56);
\draw[te] (adf56) -- (adf58);
\draw[te] (adf64) -- (adf89);
\draw[te] (adf89) -- (adf94);
\node[bfl] at (adf46.north east) {$0$};
\node[bfl] at (adf52.north east) {$-1$};
\node[bfl] at (adf56.north east) {$-1$};
\node[bfl] at (adf58.north east) {$0$};
\node[bfl] at (adf64.north east) {$+1$};
\node[bfl] at (adf89.north east) {$-1$};
\node[bfl] at (adf94.north east) {$0$};
\node[plab] at (3.0,-3.75) {(d) 左旋（RR）之後};
\end{tikzpicture}\hspace{0.8cm}\begin{tikzpicture}[x=0.5cm,y=0.72cm,baseline=(current bounding box.north)]
\node[tn,tnA] (adg52) at (0,-1) {52};
\node[tn] (adg56) at (1,-2) {56};
\node[tn] (adg58) at (2,-3) {58};
\node[tn] (adg64) at (3,0) {64};
\node[tn] (adg89) at (4,-1) {89};
\node[tn] (adg94) at (5,-2) {94};
\draw[te] (adg64) -- (adg52);
\draw[te] (adg52) -- (adg56);
\draw[te] (adg56) -- (adg58);
\draw[te] (adg64) -- (adg89);
\draw[te] (adg89) -- (adg94);
\node[bfl] at (adg52.north east) {$-2$};
\node[bfl] at (adg56.north east) {$-1$};
\node[bfl] at (adg58.north east) {$0$};
\node[bfl] at (adg64.north east) {$+1$};
\node[bfl] at (adg89.north east) {$-1$};
\node[bfl] at (adg94.north east) {$0$};
\node[plab] at (2.5,-3.75) {(e) 刪 46：52 失衡};
\end{tikzpicture}\hspace{0.8cm}\begin{tikzpicture}[x=0.5cm,y=0.72cm,baseline=(current bounding box.north)]
\node[tn] (adh52) at (0,-2) {52};
\node[tn] (adh56) at (1,-1) {56};
\node[tn] (adh58) at (2,-2) {58};
\node[tn] (adh64) at (3,0) {64};
\node[tn] (adh89) at (4,-1) {89};
\node[tn] (adh94) at (5,-2) {94};
\draw[te] (adh64) -- (adh56);
\draw[te] (adh56) -- (adh52);
\draw[te] (adh56) -- (adh58);
\draw[te] (adh64) -- (adh89);
\draw[te] (adh89) -- (adh94);
\node[bfl] at (adh52.north east) {$0$};
\node[bfl] at (adh56.north east) {$0$};
\node[bfl] at (adh58.north east) {$0$};
\node[bfl] at (adh64.north east) {$0$};
\node[bfl] at (adh89.north east) {$-1$};
\node[bfl] at (adh94.north east) {$0$};
\node[plab] at (2.5,-3.75) {(f) 左旋（RR）之後};
\end{tikzpicture}
\diagcaption{依序刪除 $41, 16, 34, 46$，開始的樹是圖~\ref{fig:bal-avl-trace3}(k)。 橘框是失衡
的節點。 刪 $16$ 沒有旋轉，沒有另外畫。}
\label{fig:bal-avl-del}
\end{center}

\begin{itemize}
\item 刪 $41$：$41$ 有兩個子節點，後繼是右子樹最小的 $46$。 $46$ 取代 $41$，再刪掉原
本的 $46$（葉節點）。 這個節點（現在是 $46$）左高 $2$、右高 $0$，平衡因子 $+2$；左子
節點 $16$ 是 $-1$，是 R$-1$，做 LR：$34$ 升上來，(b)。 這一段高度從 $3$ 變 $2$，往上
到根 $52$：左高 $2$、右高 $3$，$-1$，沒事。
\item 刪 $16$：葉節點，$34$ 變 $-1$，高度不變，結束。
\item 刪 $34$：$34$ 只有一個子節點，$46$ 接上來。 $52$ 左高 $1$、右高 $3$，變 $-2$；右
子節點 $64$ 是 $0$，是 L0，對 $52$ 左旋（RR），(d)。 這一段高度還是 $4$，可以停。
\item 刪 $46$：$52$ 左高 $0$、右高 $2$，又是 $-2$；右子節點 $56$ 是 $-1$，是 L1，對
$52$ 左旋，(f)。 這一段變矮，往上檢查根 $64$：左高 $2$、右高 $2$，平衡。
\end{itemize}

\textbf{例 2}。 刪除可能一路旋轉上去。 圖~\ref{fig:bal-avl-cascade}(a) 是一棵高度
$5$、只有 $12$ 個節點的 AVL 樹，也就是一棵費氏樹（$N(5) = 12$）。 刪除 $70$：

\begin{center}
\begin{tikzpicture}[x=0.47cm,y=0.72cm,baseline=(current bounding box.north)]
\node[tn] (aca5) at (0,-3) {5};
\node[tn] (aca12) at (1,-2) {12};
\node[tn] (aca18) at (2,-1) {18};
\node[tn] (aca23) at (3,-4) {23};
\node[tn] (aca27) at (4,-3) {27};
\node[tn] (aca31) at (5,-2) {31};
\node[tn] (aca36) at (6,-3) {36};
\node[tn] (aca44) at (7,0) {44};
\node[tn] (aca50) at (8,-3) {50};
\node[tn] (aca57) at (9,-2) {57};
\node[tn] (aca63) at (10,-1) {63};
\node[tn] (aca70) at (11,-2) {70};
\draw[te] (aca44) -- (aca18);
\draw[te] (aca18) -- (aca12);
\draw[te] (aca12) -- (aca5);
\draw[te] (aca18) -- (aca31);
\draw[te] (aca31) -- (aca27);
\draw[te] (aca27) -- (aca23);
\draw[te] (aca31) -- (aca36);
\draw[te] (aca44) -- (aca63);
\draw[te] (aca63) -- (aca57);
\draw[te] (aca57) -- (aca50);
\draw[te] (aca63) -- (aca70);
\node[bfl] at (aca5.north east) {$0$};
\node[bfl] at (aca12.north east) {$+1$};
\node[bfl] at (aca18.north east) {$-1$};
\node[bfl] at (aca23.north east) {$0$};
\node[bfl] at (aca27.north east) {$+1$};
\node[bfl] at (aca31.north east) {$+1$};
\node[bfl] at (aca36.north east) {$0$};
\node[bfl] at (aca44.north east) {$+1$};
\node[bfl] at (aca50.north east) {$0$};
\node[bfl] at (aca57.north east) {$+1$};
\node[bfl] at (aca63.north east) {$+1$};
\node[bfl] at (aca70.north east) {$0$};
\node[plab] at (5.5,-4.75) {(a) 開始：$h = 5$};
\end{tikzpicture}\hspace{0.6cm}\begin{tikzpicture}[x=0.47cm,y=0.72cm,baseline=(current bounding box.north)]
\node[tn] (acb5) at (0,-3) {5};
\node[tn] (acb12) at (1,-2) {12};
\node[tn] (acb18) at (2,-1) {18};
\node[tn] (acb23) at (3,-4) {23};
\node[tn] (acb27) at (4,-3) {27};
\node[tn] (acb31) at (5,-2) {31};
\node[tn] (acb36) at (6,-3) {36};
\node[tn] (acb44) at (7,0) {44};
\node[tn] (acb50) at (8,-3) {50};
\node[tn] (acb57) at (9,-2) {57};
\node[tn,tnA] (acb63) at (10,-1) {63};
\draw[te] (acb44) -- (acb18);
\draw[te] (acb18) -- (acb12);
\draw[te] (acb12) -- (acb5);
\draw[te] (acb18) -- (acb31);
\draw[te] (acb31) -- (acb27);
\draw[te] (acb27) -- (acb23);
\draw[te] (acb31) -- (acb36);
\draw[te] (acb44) -- (acb63);
\draw[te] (acb63) -- (acb57);
\draw[te] (acb57) -- (acb50);
\node[bfl] at (acb5.north east) {$0$};
\node[bfl] at (acb12.north east) {$+1$};
\node[bfl] at (acb18.north east) {$-1$};
\node[bfl] at (acb23.north east) {$0$};
\node[bfl] at (acb27.north east) {$+1$};
\node[bfl] at (acb31.north east) {$+1$};
\node[bfl] at (acb36.north east) {$0$};
\node[bfl] at (acb44.north east) {$+1$};
\node[bfl] at (acb50.north east) {$0$};
\node[bfl] at (acb57.north east) {$+1$};
\node[bfl] at (acb63.north east) {$+2$};
\node[plab] at (5.0,-4.75) {(b) 刪 70：63 失衡};
\end{tikzpicture}\\*[0.6em]
\begin{tikzpicture}[x=0.47cm,y=0.72cm,baseline=(current bounding box.north)]
\node[tn] (acc5) at (0,-3) {5};
\node[tn] (acc12) at (1,-2) {12};
\node[tn] (acc18) at (2,-1) {18};
\node[tn] (acc23) at (3,-4) {23};
\node[tn] (acc27) at (4,-3) {27};
\node[tn] (acc31) at (5,-2) {31};
\node[tn] (acc36) at (6,-3) {36};
\node[tn,tnA] (acc44) at (7,0) {44};
\node[tn] (acc50) at (8,-2) {50};
\node[tn] (acc57) at (9,-1) {57};
\node[tn] (acc63) at (10,-2) {63};
\draw[te] (acc44) -- (acc18);
\draw[te] (acc18) -- (acc12);
\draw[te] (acc12) -- (acc5);
\draw[te] (acc18) -- (acc31);
\draw[te] (acc31) -- (acc27);
\draw[te] (acc27) -- (acc23);
\draw[te] (acc31) -- (acc36);
\draw[te] (acc44) -- (acc57);
\draw[te] (acc57) -- (acc50);
\draw[te] (acc57) -- (acc63);
\node[bfl] at (acc5.north east) {$0$};
\node[bfl] at (acc12.north east) {$+1$};
\node[bfl] at (acc18.north east) {$-1$};
\node[bfl] at (acc23.north east) {$0$};
\node[bfl] at (acc27.north east) {$+1$};
\node[bfl] at (acc31.north east) {$+1$};
\node[bfl] at (acc36.north east) {$0$};
\node[bfl] at (acc44.north east) {$+2$};
\node[bfl] at (acc50.north east) {$0$};
\node[bfl] at (acc57.north east) {$0$};
\node[bfl] at (acc63.north east) {$0$};
\node[plab] at (5.0,-4.75) {(c) 右旋 63：換 44 失衡};
\end{tikzpicture}\hspace{0.6cm}\begin{tikzpicture}[x=0.47cm,y=0.72cm,baseline=(current bounding box.north)]
\node[tn] (acd5) at (0,-3) {5};
\node[tn] (acd12) at (1,-2) {12};
\node[tn] (acd18) at (2,-1) {18};
\node[tn] (acd23) at (3,-3) {23};
\node[tn] (acd27) at (4,-2) {27};
\node[tn] (acd31) at (5,0) {31};
\node[tn] (acd36) at (6,-2) {36};
\node[tn] (acd44) at (7,-1) {44};
\node[tn] (acd50) at (8,-3) {50};
\node[tn] (acd57) at (9,-2) {57};
\node[tn] (acd63) at (10,-3) {63};
\draw[te] (acd31) -- (acd18);
\draw[te] (acd18) -- (acd12);
\draw[te] (acd12) -- (acd5);
\draw[te] (acd18) -- (acd27);
\draw[te] (acd27) -- (acd23);
\draw[te] (acd31) -- (acd44);
\draw[te] (acd44) -- (acd36);
\draw[te] (acd44) -- (acd57);
\draw[te] (acd57) -- (acd50);
\draw[te] (acd57) -- (acd63);
\node[bfl] at (acd5.north east) {$0$};
\node[bfl] at (acd12.north east) {$+1$};
\node[bfl] at (acd18.north east) {$0$};
\node[bfl] at (acd23.north east) {$0$};
\node[bfl] at (acd27.north east) {$+1$};
\node[bfl] at (acd31.north east) {$0$};
\node[bfl] at (acd36.north east) {$0$};
\node[bfl] at (acd44.north east) {$-1$};
\node[bfl] at (acd50.north east) {$0$};
\node[bfl] at (acd57.north east) {$0$};
\node[bfl] at (acd63.north east) {$0$};
\node[plab] at (5.0,-4.75) {(d) 雙旋轉（LR）之後};
\end{tikzpicture}
\diagcaption{在費氏樹裡刪除 $70$，先在 $63$ 做一次單旋轉，這一段變矮，又讓根 $44$
失衡，再做一次雙旋轉。}
\label{fig:bal-avl-cascade}
\end{center}

\begin{itemize}
\item $63$ 左高 $2$、右高 $0$，$+2$；左子節點 $57$ 是 $+1$，是 R1，對 $63$ 右旋。
這一段的高度從 $3$ 變成 $2$。
\item 往上到根 $44$：左高 $4$、右高 $2$，$+2$；左子節點 $18$ 是 $-1$，是 R$-1$，做
LR。 $18 < 31 < 44$，$31$ 升到根；$31$ 原本的子樹 $\{23, 27\}$ 接到 $18$ 的右邊，
$\{36\}$ 接到 $44$ 的左邊。
\end{itemize}

一次刪除做了三次旋轉（一次單旋轉、一次雙旋轉）。 更大的費氏樹可以讓每一層都旋轉
一次，所以 AVL 刪除的旋轉次數最壞是 $\Theta(\lg n)$，不像插入最多兩次。

\begin{examnote}
刪除有兩個子節點的節點時，用後繼（右子樹最小）或前驅（左子樹最大）取代都是對的，
但畫出來的樹不一樣。 題目沒指定時，在答案第一行寫明用哪一種，整題都用同一種。 每
次旋轉之後要繼續往上檢查到根，不能像插入那樣轉完就停。
\end{examnote}

\section{紅黑樹}
\label{sec:bal-rb}

\subsection{五條性質}

紅黑樹把「平衡」放寬成另一種條件：每個節點多存一個位元的顏色，用顏色限制路徑長度。
為了說得清楚，把每個空的子指標都想成一個\textbf{外部節點}（第~7~章延伸二元樹的外部節點），叫
NIL（圖上畫成小黑方塊）；原本存鍵的節點叫\textbf{內部節點}。 C 程式裡 NIL 就是 \texttt{NULL}，不用真
的配置。

\begin{definition}[紅黑樹]
\label{def:bal-rb}
一棵二元搜尋樹，每個節點是紅色或黑色，而且滿足：
\begin{enumerate}
\item 每個節點不是紅色就是黑色。
\item 根是黑色。
\item 每個 NIL 都是黑色。
\item 紅色節點的兩個子節點都是黑色。 也就是說，不會有兩個紅色節點相連。
\item 對每個節點 $x$，從 $x$ 往下走到任何一個 NIL 的每一條路徑上，黑色節點的個數都
一樣。
\end{enumerate}
\end{definition}

第 5 條說的個數叫 $x$ 的\textbf{黑高}（black-height），記成 $\mathrm{bh}(x)$：從 $x$
往下走到任一個 NIL，路上黑色節點的個數，\textbf{不算 $x$ 自己，算 NIL}。 整棵樹的黑
高是根的黑高。 有的書不算 NIL、或把 $x$ 自己算進去，數字會差 $1$；考卷有定義時照題目。

\begin{center}
\begin{tikzpicture}[x=0.42cm,y=0.85cm,baseline=(current bounding box.north)]
\node[tnil] (rbdefNIL) at (0,-3) {};
\node[tnB] (rbdef11) at (1,-2) {11};
\node[tnil] (rbdefnil2) at (2,-3) {};
\node[tnR] (rbdef19) at (3,-1) {19};
\node[tnil] (rbdefnil4) at (4,-3) {};
\node[tnB] (rbdef27) at (5,-2) {27};
\node[tnil] (rbdefnil6) at (6,-4) {};
\node[tnR] (rbdef31) at (7,-3) {31};
\node[tnil] (rbdefnil8) at (8,-4) {};
\node[tnB] (rbdef36) at (9,0) {36};
\node[tnil] (rbdefnil10) at (10,-3) {};
\node[tnR] (rbdef44) at (11,-2) {44};
\node[tnil] (rbdefnil12) at (12,-3) {};
\node[tnB] (rbdef58) at (13,-1) {58};
\node[tnil] (rbdefnil14) at (14,-3) {};
\node[tnR] (rbdef73) at (15,-2) {73};
\node[tnil] (rbdefnil16) at (16,-3) {};
\draw[te] (rbdef36) -- (rbdef19);
\draw[te] (rbdef19) -- (rbdef11);
\draw[te] (rbdef11) -- (rbdefNIL);
\draw[te] (rbdef11) -- (rbdefnil2);
\draw[te] (rbdef19) -- (rbdef27);
\draw[te] (rbdef27) -- (rbdefnil4);
\draw[te] (rbdef27) -- (rbdef31);
\draw[te] (rbdef31) -- (rbdefnil6);
\draw[te] (rbdef31) -- (rbdefnil8);
\draw[te] (rbdef36) -- (rbdef58);
\draw[te] (rbdef58) -- (rbdef44);
\draw[te] (rbdef44) -- (rbdefnil10);
\draw[te] (rbdef44) -- (rbdefnil12);
\draw[te] (rbdef58) -- (rbdef73);
\draw[te] (rbdef73) -- (rbdefnil14);
\draw[te] (rbdef73) -- (rbdefnil16);
\end{tikzpicture}
\diagcaption{一棵紅黑樹，小方塊是 NIL。 根 $36$ 的黑高是 $2$：例如
$36 \to 19 \to 11 \to \text{NIL}$ 經過 $11$、NIL 兩個黑色；$36 \to 58 \to 44 \to
\text{NIL}$ 經過 $58$、NIL。 $19$ 和 $58$ 的黑高都是 $1$。}
\label{fig:bal-rbdef}
\end{center}

之後的圖不再畫 NIL，但要記得它們在那裡：例如 $27$ 只有一個紅色的右子節點 $31$，它
的左邊是一個 NIL，$27 \to \text{NIL}$ 和 $27 \to 31 \to \text{NIL}$ 都只有一個黑色
（NIL），所以 $27$ 是合法的。 如果 $31$ 是黑的，右邊那條就有兩個黑色，違反第 5 條。

\subsection{高度上界}

\begin{lemma}
\label{lem:bal-rb-count}
以 $x$ 為根的子樹至少有 $2^{\mathrm{bh}(x)} - 1$ 個內部節點。
\end{lemma}

\begin{proof}
對 $x$ 的高度歸納。 $x$ 是 NIL 時，$\mathrm{bh}(x) = 0$，子樹有 $0 = 2^0 - 1$ 個內部
節點。 $x$ 是內部節點時，它的子節點 $c$ 如果是紅的，$\mathrm{bh}(c) = \mathrm{bh}(x)$；
如果是黑的，$\mathrm{bh}(c) = \mathrm{bh}(x) - 1$。 不管哪種，兩個子節點的黑高都至
少 $\mathrm{bh}(x) - 1$，而且子樹比較矮，由歸納假設各至少有 $2^{\mathrm{bh}(x)-1} - 1$
個內部節點。 加上 $x$ 自己：
\[
2\bigl(2^{\mathrm{bh}(x)-1} - 1\bigr) + 1 = 2^{\mathrm{bh}(x)} - 1. \qedhere
\]
\end{proof}

\begin{theorem}
\label{thm:bal-rb-height}
有 $n$ 個內部節點的紅黑樹，高度 $h \le 2 \lg(n + 1)$。
\end{theorem}

\begin{proof}
取一條最長的路徑，從根經過 $h$ 個內部節點 $v_1 = \text{根}, v_2, \ldots, v_h$ 走到一
個 NIL。 看 $v_2, \ldots, v_h$ 加上最後的 NIL，共 $h$ 個節點。 由第 4 條，紅色節點
的下一個一定是黑色，所以每個紅色節點都可以和它下一個節點配成一對，各對互不重疊；
沒配到的都是黑色。 因此這 $h$ 個節點裡黑色至少佔一半，$\mathrm{bh}(\text{根}) \ge
h/2$。 由引理~\ref{lem:bal-rb-count}，$n \ge 2^{h/2} - 1$，移項取對數得到
$h \le 2\lg(n+1)$。
\end{proof}

證明用到的事實是：任何一條路徑上，紅色不會比黑色多。 所以同一個節點往下的最長路
徑，不會超過最短路徑的兩倍（觀念題第 9 題）。 這比 AVL 的條件寬，$2$ 對上 $1.44$，
換來的是插入、刪除時要調整的地方比較少。

搜尋和二元搜尋樹完全一樣，不看顏色，$O(\lg n)$。 要處理的是插入和刪除：照二元搜尋
樹的方法改完之後，顏色可能違反性質，要用改色和旋轉修回來。

\subsection{帶父指標的旋轉}

紅黑樹修正時要往上看父節點、祖父節點和叔叔（父節點的兄弟），所以節點多存一個父指
標 $p$，旋轉時三組父子關係都要雙向更新。 下面是左旋，右旋把 $\mathit{left}$、
$\mathit{right}$ 對調。

\begin{pseudo}[label=alg:bal-left-rotate]{帶父指標的左旋：$a$ 的右子節點 $b$ 升上來}
\Procedure{Left-Rotate}{$T, a$}
    \State $b \gets a.\mathit{right}$
    \State $a.\mathit{right} \gets b.\mathit{left}$ \Comment{$b$ 靠內側的子樹給 $a$}
    \If{$b.\mathit{left} \ne \NIL$}
        \State $b.\mathit{left}.p \gets a$
    \EndIf
    \State $b.p \gets a.p$ \Comment{$b$ 接到 $a$ 原本的位置}
    \If{$a.p = \NIL$}
        \State $T.\mathit{root} \gets b$
    \ElsIf{$a = a.p.\mathit{left}$}
        \State $a.p.\mathit{left} \gets b$
    \Else
        \State $a.p.\mathit{right} \gets b$
    \EndIf
    \State $b.\mathit{left} \gets a$
    \State $a.p \gets b$
\EndProcedure
\end{pseudo}

\subsection{插入}
\label{sec:bal-rb-insert}

先照二元搜尋樹插入新節點 $z$，\textbf{塗成紅色}。 紅色節點不改變任何路徑上的黑色個
數，所以第 5 條不會壞；第 1、3 條也不會壞。 可能壞的只有兩條：
\begin{itemize}
\item 第 2 條：$z$ 是根（原本是空樹）。 把它改成黑色就好。
\item 第 4 條：$z$ 的父節點 $p$ 也是紅的。 這時 $p$ 不是根（根是黑的），所以祖父
$g$ 存在，而且 $g$ 是黑的（樹原本合法，$p$ 紅則 $g$ 黑）。
\end{itemize}

修第 4 條時，關鍵是叔叔 $u$（$g$ 的另一個子節點）的顏色。 下面假設 $p$ 是 $g$ 的
左子節點，右邊的情況左右對調。

\begin{itemize}
\item \textbf{情況 1：$u$ 是紅的。} 把 $p$、$u$ 改黑，$g$ 改紅。 經過 $g$ 的每條路徑，
少了 $g$ 一個黑色、多了 $p$ 或 $u$ 一個黑色，個數不變。 $z$ 和 $p$ 不再紅紅相連，但
$g$ 變紅之後，它的父節點可能是紅的。 把 $z$ 移到 $g$，往上兩層繼續檢查。
\item \textbf{情況 2：$u$ 是黑的，$z$ 是 $p$ 的右子節點（內側）。} 把 $z$ 移到 $p$，
對新的 $z$ 左旋。 兩個紅色節點換了上下，紅紅相連還在，但現在是外側，變成情況 3。
這次旋轉只動兩個紅色節點，路徑上的黑色個數不變。
\item \textbf{情況 3：$u$ 是黑的，$z$ 是 $p$ 的左子節點（外側）。} $p$ 改黑，$g$ 改
紅，對 $g$ 右旋。 $p$ 升上來，黑色的 $p$ 取代了原本黑色的 $g$ 的位置，左右兩邊的黑
色個數都和原本一樣，而且 $p$ 是黑的，上面不會有紅紅相連，修正結束。
\end{itemize}

迴圈結束時（$p$ 是黑的，或 $z$ 走到了根），最後把根塗黑。 情況 1 只有在 $g$ 變紅時
才可能讓根變紅，這一步正是處理它。

\begin{center}
\begin{tikzpicture}[x=0.75cm, y=0.8cm]
% 情況 1
\node[font=\small\bfseries, anchor=west] at (-2.4,0.7) {情況 1};
\node[tnB] (g) at (0,0) {$g$};
\node[tnR] (p) at (-1.2,-1) {$p$};
\node[tnR] (u) at (1.2,-1) {$u$};
\node[tnR] (z) at (-1.8,-2) {$z$};
\draw[te] (g) -- (p); \draw[te] (g) -- (u); \draw[te] (p) -- (z); \draw[te] (p) -- (-0.6,-2.0);
\subtri[0.7]{-0.6}{-2.0}{3}
\draw[tarr] (2.0,-1) -- node[above, font=\scriptsize] {改色} (3.3,-1);
\begin{scope}[xshift=4.4cm]
\node[tnR] (g) at (0,0) {$g$};
\node[tnB] (p) at (-1.2,-1) {$p$};
\node[tnB] (u) at (1.2,-1) {$u$};
\node[tnR] (z) at (-1.8,-2) {};
\draw[te] (g) -- (p); \draw[te] (g) -- (u); \draw[te] (p) -- (z); \draw[te] (p) -- (-0.6,-2.0);
\subtri[0.7]{-0.6}{-2.0}{3}
\node[annl, anchor=south east] at (g.north west) {新的 $z$};
\end{scope}
% 情況 2 → 3
\begin{scope}[yshift=-3.6cm]
\node[font=\small\bfseries, anchor=west] at (-2.4,0.7) {情況 2};
\node[tnB] (g) at (0,0) {$g$};
\node[tnR] (p) at (-1.2,-1) {$p$};
\node[tnB] (u) at (1.2,-1) {$u$};
\node[tnR] (z) at (-0.6,-2) {$z$};
\draw[te] (g) -- (p); \draw[te] (g) -- (u); \draw[te] (p) -- (-1.8,-2.0); \draw[te] (p) -- (z);
\draw[te] (z) -- (-1.0,-3.0); \draw[te] (z) -- (-0.2,-3.0);
\subtri[0.7]{-1.8}{-2.0}{1}
\subtri[0.7]{-1.0}{-3.0}{2}
\subtri[0.7]{-0.2}{-3.0}{3}
\draw[tarr] (2.0,-1) -- node[above, font=\scriptsize] {左旋 $p$} (3.3,-1);
\begin{scope}[xshift=4.4cm]
\node[tnB] (g) at (0,0) {$g$};
\node[tnR] (z) at (-1.2,-1) {$z'$};
\node[tnB] (u) at (1.2,-1) {$u$};
\node[tnR] (p) at (-1.8,-2) {$p'$};
\draw[te] (g) -- (z); \draw[te] (g) -- (u); \draw[te] (z) -- (p); \draw[te] (z) -- (-0.6,-2.0);
\draw[te] (p) -- (-2.2,-3.0); \draw[te] (p) -- (-1.4,-3.0);
\subtri[0.7]{-2.2}{-3.0}{1}
\subtri[0.7]{-1.4}{-3.0}{2}
\subtri[0.7]{-0.6}{-2.0}{3}
\end{scope}
\end{scope}
% 情況 3
\begin{scope}[xshift=9.6cm]
\node[font=\small\bfseries, anchor=west] at (-2.4,0.7) {情況 3};
\node[tnB] (g) at (0,0) {$g$};
\node[tnR] (p) at (-1.2,-1) {$p$};
\node[tnB] (u) at (1.2,-1) {$u$};
\node[tnR] (z) at (-1.8,-2) {$z$};
\draw[te] (g) -- (p); \draw[te] (g) -- (u); \draw[te] (p) -- (z); \draw[te] (p) -- (-0.6,-2.0);
\subtri[0.7]{-0.6}{-2.0}{3}
\draw[tarr] (0,-2.9) -- node[right, font=\scriptsize] {右旋 $g$、改色} (0,-3.9);
\begin{scope}[yshift=-4.4cm]
\node[tnB] (p) at (0,0) {$p$};
\node[tnR] (z) at (-1.2,-1) {$z$};
\node[tnR] (g) at (1.2,-1) {$g$};
\node[tnB] (u) at (1.8,-2) {$u$};
\draw[te] (p) -- (z); \draw[te] (p) -- (g); \draw[te] (g) -- (0.6,-2.0); \draw[te] (g) -- (u);
\subtri[0.7]{0.6}{-2.0}{3}
\end{scope}
\end{scope}
\end{tikzpicture}
\diagcaption{插入的三種情況（$p$ 是 $g$ 的左子節點）。 三角形是黑高相同的子樹，
$z$ 本身的子樹沒有畫。 情況 1 的 $z$ 也可以是 $p$ 的右子節點，做法一樣。 情況 2
左旋之後，$z'$（原本的 $z$）和 $p'$（原本的 $p$）換了上下，再接情況 3：這時的
「$z$」是 $p'$，「$p$」是 $z'$。}
\label{fig:bal-rb-insert-cases}
\end{center}

\proc{RB-Insert}$(T, k)$ 和第~9~章的插入相同：從根往下找到位置，新節點 $z$ 的父指標
設成找到的父節點，兩個子指標是 \NIL，顏色是紅色；接好之後呼叫
\proc{RB-Insert-Fixup}$(T, z)$。

\begin{pseudo}[label=alg:bal-rb-fixup]{插入後的修正}
\Procedure{RB-Insert-Fixup}{$T, z$}
    \While{$z.p$ 是紅的}
        \State $p \gets z.p$
        \State $g \gets p.p$
        \If{$p = g.\mathit{left}$}
            \State $u \gets g.\mathit{right}$
            \If{$u$ 是紅的} \Comment{情況 1}
                \State $p$、$u$ 改黑，$g$ 改紅
                \State $z \gets g$
            \Else
                \If{$z = p.\mathit{right}$} \Comment{情況 2}
                    \State $z \gets p$
                    \State \Call{Left-Rotate}{$T, z$}
                    \State $p \gets z.p$
                \EndIf
                \State $p$ 改黑，$g$ 改紅 \Comment{情況 3}
                \State \Call{Right-Rotate}{$T, g$}
            \EndIf
        \Else
            \State 和第 6 到 16 行相同，$\mathit{left}$、$\mathit{right}$ 對調，\proc{Left-Rotate}、\proc{Right-Rotate} 對調
        \EndIf
    \EndWhile
    \State $T.\mathit{root}$ 改黑
\EndProcedure
\end{pseudo}

NIL 一律當成黑色：$z$ 是根時 $z.p = \NIL$，第 2 行的條件不成立；叔叔是 NIL 時，它
是黑的，走情況 2 或 3。

\textbf{成本。} 情況 1 讓 $z$ 往上兩層，最多做 $O(\lg n)$ 次，每次只改顏色。 情況 2
最多一次，接著情況 3 一次，做完迴圈就結束。 所以插入一共最多\textbf{兩次旋轉}，時間
$O(\lg n)$。

\subsection{插入的完整追蹤}
\label{sec:bal-rb-trace}

依序插入
\[
85,\ 28,\ 24,\ 87,\ 75,\ 50,\ 69,\ 46,\ 52,\ 35 .
\]
圖上深色是黑節點，紅底是紅節點，「$z$」標的是目前要檢查的節點。 每一步的說明寫
在圖下；「情況 3」之類的字，指的是\textbf{接下來}要做的事。

\begin{center}
\begin{tikzpicture}[x=0.5cm,y=0.7cm,baseline=(current bounding box.north)]
\node[tnB] (ria85) at (0,0) {85};
\node[plab] at (0.0,-2.75) {(a) 插入 85，根改黑};
\end{tikzpicture}\hspace{0.5cm}\begin{tikzpicture}[x=0.5cm,y=0.7cm,baseline=(current bounding box.north)]
\node[tnR] (rib24) at (0,-2) {24};
\node[tnR] (rib28) at (1,-1) {28};
\node[tnB] (rib85) at (2,0) {85};
\draw[te] (rib85) -- (rib28);
\draw[te] (rib28) -- (rib24);
\node[annl, anchor=south east] at (rib24.north west) {$z$};
\node[plab] at (1.0,-2.75) {(b) 插入 28、24：情況 3};
\end{tikzpicture}\hspace{0.5cm}\begin{tikzpicture}[x=0.5cm,y=0.7cm,baseline=(current bounding box.north)]
\node[tnR] (ric24) at (0,-1) {24};
\node[tnB] (ric28) at (1,0) {28};
\node[tnR] (ric85) at (2,-1) {85};
\draw[te] (ric28) -- (ric24);
\draw[te] (ric28) -- (ric85);
\node[plab] at (1.0,-2.75) {(c) 右旋 85 之後};
\end{tikzpicture}\\*[0.6em]
\begin{tikzpicture}[x=0.5cm,y=0.7cm,baseline=(current bounding box.north)]
\node[tnR] (rid24) at (0,-1) {24};
\node[tnB] (rid28) at (1,0) {28};
\node[tnR] (rid85) at (2,-1) {85};
\node[tnR] (rid87) at (3,-2) {87};
\draw[te] (rid28) -- (rid24);
\draw[te] (rid28) -- (rid85);
\draw[te] (rid85) -- (rid87);
\node[annl, anchor=south east] at (rid87.north west) {$z$};
\node[plab] at (1.5,-2.75) {(d) 插入 87：情況 1};
\end{tikzpicture}\hspace{0.6cm}\begin{tikzpicture}[x=0.5cm,y=0.7cm,baseline=(current bounding box.north)]
\node[tnB] (rie24) at (0,-1) {24};
\node[tnB] (rie28) at (1,0) {28};
\node[tnB] (rie85) at (2,-1) {85};
\node[tnR] (rie87) at (3,-2) {87};
\draw[te] (rie28) -- (rie24);
\draw[te] (rie28) -- (rie85);
\draw[te] (rie85) -- (rie87);
\node[plab] at (1.5,-2.75) {(e) 改色、根改回黑};
\end{tikzpicture}
\diagcaption{插入 $85$ 到 $87$。}
\label{fig:bal-rb-trace1}
\end{center}

\begin{itemize}
\item $85$：樹是空的，$z$ 是根，改黑。 $28$：父節點 $85$ 是黑的，不用修。
\item $24$：父 $28$ 紅，叔叔是 $85$ 的右邊，NIL，黑的。 $24$ 是 $28$ 的左子節點、
$28$ 是 $85$ 的左子節點，在外側，情況 3：$28$ 改黑、$85$ 改紅，右旋 $85$。
\item $87$：父 $85$ 紅，叔叔 $24$ 也紅，情況 1：$85$、$24$ 改黑，$28$ 改紅。 $z$ 移到
$28$，它是根，迴圈結束，根改回黑。 這是黑高增加的唯一方式：情況 1 一路改色到根，根
被改紅再改黑，整棵樹的每條路徑同時多一個黑色。
\end{itemize}

\begin{center}
\begin{tikzpicture}[x=0.5cm,y=0.7cm,baseline=(current bounding box.north)]
\node[tnB] (rif24) at (0,-1) {24};
\node[tnB] (rif28) at (1,0) {28};
\node[tnR] (rif50) at (2,-3) {50};
\node[tnR] (rif75) at (3,-2) {75};
\node[tnB] (rif85) at (4,-1) {85};
\node[tnR] (rif87) at (5,-2) {87};
\draw[te] (rif28) -- (rif24);
\draw[te] (rif28) -- (rif85);
\draw[te] (rif85) -- (rif75);
\draw[te] (rif75) -- (rif50);
\draw[te] (rif85) -- (rif87);
\node[annl, anchor=south east] at (rif50.north west) {$z$};
\node[plab] at (2.5,-4.75) {(f) 插入 75、50：情況 1};
\end{tikzpicture}\hspace{0.4cm}\begin{tikzpicture}[x=0.5cm,y=0.7cm,baseline=(current bounding box.north)]
\node[tnB] (rig24) at (0,-1) {24};
\node[tnB] (rig28) at (1,0) {28};
\node[tnR] (rig50) at (2,-3) {50};
\node[tnB] (rig75) at (3,-2) {75};
\node[tnR] (rig85) at (4,-1) {85};
\node[tnB] (rig87) at (5,-2) {87};
\draw[te] (rig28) -- (rig24);
\draw[te] (rig28) -- (rig85);
\draw[te] (rig85) -- (rig75);
\draw[te] (rig75) -- (rig50);
\draw[te] (rig85) -- (rig87);
\node[plab] at (2.5,-4.75) {(g) 改色之後};
\end{tikzpicture}\hspace{0.4cm}\begin{tikzpicture}[x=0.5cm,y=0.7cm,baseline=(current bounding box.north)]
\node[tnB] (rih24) at (0,-1) {24};
\node[tnB] (rih28) at (1,0) {28};
\node[tnR] (rih50) at (2,-3) {50};
\node[tnR] (rih69) at (3,-4) {69};
\node[tnB] (rih75) at (4,-2) {75};
\node[tnR] (rih85) at (5,-1) {85};
\node[tnB] (rih87) at (6,-2) {87};
\draw[te] (rih28) -- (rih24);
\draw[te] (rih28) -- (rih85);
\draw[te] (rih85) -- (rih75);
\draw[te] (rih75) -- (rih50);
\draw[te] (rih50) -- (rih69);
\draw[te] (rih85) -- (rih87);
\node[annl, anchor=south east] at (rih69.north west) {$z$};
\node[plab] at (3.0,-4.75) {(h) 插入 69：情況 2};
\end{tikzpicture}\\*[0.6em]
\begin{tikzpicture}[x=0.5cm,y=0.7cm,baseline=(current bounding box.north)]
\node[tnB] (rii24) at (0,-1) {24};
\node[tnB] (rii28) at (1,0) {28};
\node[tnR] (rii50) at (2,-4) {50};
\node[tnR] (rii69) at (3,-3) {69};
\node[tnB] (rii75) at (4,-2) {75};
\node[tnR] (rii85) at (5,-1) {85};
\node[tnB] (rii87) at (6,-2) {87};
\draw[te] (rii28) -- (rii24);
\draw[te] (rii28) -- (rii85);
\draw[te] (rii85) -- (rii75);
\draw[te] (rii75) -- (rii69);
\draw[te] (rii69) -- (rii50);
\draw[te] (rii85) -- (rii87);
\node[annl, anchor=south east] at (rii50.north west) {$z$};
\node[plab] at (3.0,-4.75) {(i) 左旋 50：情況 3};
\end{tikzpicture}\hspace{0.4cm}\begin{tikzpicture}[x=0.5cm,y=0.7cm,baseline=(current bounding box.north)]
\node[tnB] (rij24) at (0,-1) {24};
\node[tnB] (rij28) at (1,0) {28};
\node[tnR] (rij50) at (2,-3) {50};
\node[tnB] (rij69) at (3,-2) {69};
\node[tnR] (rij75) at (4,-3) {75};
\node[tnR] (rij85) at (5,-1) {85};
\node[tnB] (rij87) at (6,-2) {87};
\draw[te] (rij28) -- (rij24);
\draw[te] (rij28) -- (rij85);
\draw[te] (rij85) -- (rij69);
\draw[te] (rij69) -- (rij50);
\draw[te] (rij69) -- (rij75);
\draw[te] (rij85) -- (rij87);
\node[plab] at (3.0,-4.75) {(j) 右旋 75 之後};
\end{tikzpicture}
\diagcaption{插入 $75$、$50$、$69$。}
\label{fig:bal-rb-trace2}
\end{center}

\begin{itemize}
\item $75$：父 $85$ 黑，不用修。
\item $50$：父 $75$ 紅，叔叔 $87$ 紅，情況 1：$75$、$87$ 改黑，$85$ 改紅。 $z$ 移到
$85$，它的父節點 $28$ 是黑的，結束。
\item $69$：父 $50$ 紅；$50$ 是 $75$ 的左子節點，叔叔是 $75$ 的右邊（NIL），黑的。
$69$ 是 $50$ 的右子節點，在內側，情況 2：$z$ 移到 $50$，左旋 $50$，$69$ 升上來，
(i)。 現在 $z = 50$ 是 $69$ 的左子節點，在外側，情況 3：$69$ 改黑、$75$ 改紅，右
旋 $75$，(j)。
\end{itemize}

\begin{center}
\begin{tikzpicture}[x=0.5cm,y=0.7cm,baseline=(current bounding box.north)]
\node[tnB] (rik24) at (0,-1) {24};
\node[tnB] (rik28) at (1,0) {28};
\node[tnR] (rik46) at (2,-4) {46};
\node[tnR] (rik50) at (3,-3) {50};
\node[tnB] (rik69) at (4,-2) {69};
\node[tnR] (rik75) at (5,-3) {75};
\node[tnR] (rik85) at (6,-1) {85};
\node[tnB] (rik87) at (7,-2) {87};
\draw[te] (rik28) -- (rik24);
\draw[te] (rik28) -- (rik85);
\draw[te] (rik85) -- (rik69);
\draw[te] (rik69) -- (rik50);
\draw[te] (rik50) -- (rik46);
\draw[te] (rik69) -- (rik75);
\draw[te] (rik85) -- (rik87);
\node[annl, anchor=south east] at (rik46.north west) {$z$};
\node[plab] at (3.5,-4.75) {(k) 插入 46：情況 1};
\end{tikzpicture}\hspace{0.4cm}\begin{tikzpicture}[x=0.5cm,y=0.7cm,baseline=(current bounding box.north)]
\node[tnB] (ril24) at (0,-1) {24};
\node[tnB] (ril28) at (1,0) {28};
\node[tnR] (ril46) at (2,-4) {46};
\node[tnB] (ril50) at (3,-3) {50};
\node[tnR] (ril69) at (4,-2) {69};
\node[tnB] (ril75) at (5,-3) {75};
\node[tnR] (ril85) at (6,-1) {85};
\node[tnB] (ril87) at (7,-2) {87};
\draw[te] (ril28) -- (ril24);
\draw[te] (ril28) -- (ril85);
\draw[te] (ril85) -- (ril69);
\draw[te] (ril69) -- (ril50);
\draw[te] (ril50) -- (ril46);
\draw[te] (ril69) -- (ril75);
\draw[te] (ril85) -- (ril87);
\node[annl, anchor=south east] at (ril69.north west) {$z$};
\node[plab] at (3.5,-4.75) {(l) 改色後 $z = 69$：情況 2};
\end{tikzpicture}\\*[0.6em]
\begin{tikzpicture}[x=0.5cm,y=0.7cm,baseline=(current bounding box.north)]
\node[tnB] (rim24) at (0,-1) {24};
\node[tnB] (rim28) at (1,0) {28};
\node[tnR] (rim46) at (2,-3) {46};
\node[tnB] (rim50) at (3,-2) {50};
\node[tnR] (rim69) at (4,-1) {69};
\node[tnB] (rim75) at (5,-3) {75};
\node[tnR] (rim85) at (6,-2) {85};
\node[tnB] (rim87) at (7,-3) {87};
\draw[te] (rim28) -- (rim24);
\draw[te] (rim28) -- (rim69);
\draw[te] (rim69) -- (rim50);
\draw[te] (rim50) -- (rim46);
\draw[te] (rim69) -- (rim85);
\draw[te] (rim85) -- (rim75);
\draw[te] (rim85) -- (rim87);
\node[annl, anchor=south east] at (rim85.north west) {$z$};
\node[plab] at (3.5,-4.75) {(m) 右旋 85：情況 3};
\end{tikzpicture}\hspace{0.4cm}\begin{tikzpicture}[x=0.5cm,y=0.7cm,baseline=(current bounding box.north)]
\node[tnB] (rin24) at (0,-2) {24};
\node[tnR] (rin28) at (1,-1) {28};
\node[tnR] (rin46) at (2,-3) {46};
\node[tnB] (rin50) at (3,-2) {50};
\node[tnB] (rin69) at (4,0) {69};
\node[tnB] (rin75) at (5,-2) {75};
\node[tnR] (rin85) at (6,-1) {85};
\node[tnB] (rin87) at (7,-2) {87};
\draw[te] (rin69) -- (rin28);
\draw[te] (rin28) -- (rin24);
\draw[te] (rin28) -- (rin50);
\draw[te] (rin50) -- (rin46);
\draw[te] (rin69) -- (rin85);
\draw[te] (rin85) -- (rin75);
\draw[te] (rin85) -- (rin87);
\node[plab] at (3.5,-4.75) {(n) 左旋 28 之後};
\end{tikzpicture}
\diagcaption{插入 $46$：情況 1 之後接著情況 2、情況 3。}
\label{fig:bal-rb-trace3}
\end{center}

\begin{itemize}
\item $46$：父 $50$ 紅，叔叔 $75$ 紅，情況 1：$50$、$75$ 改黑，$69$ 改紅，$z$ 移到
$69$，(l)。 $69$ 的父節點 $85$ 是紅的，還要修。 這次 $p = 85$ 是祖父 $28$ 的
\textbf{右}子節點，用左右對調的版本：叔叔是 $24$，黑的；$z = 69$ 是 $85$ 的左子節
點，在內側，情況 2：$z$ 移到 $85$，右旋 $85$，(m)。 接著情況 3：$69$ 改黑、$28$
改紅，左旋 $28$，$69$ 成為新的根，(n)。
\end{itemize}

\begin{center}
\begin{tikzpicture}[x=0.5cm,y=0.7cm,baseline=(current bounding box.north)]
\node[tnB] (rio24) at (0,-2) {24};
\node[tnR] (rio28) at (1,-1) {28};
\node[tnR] (rio35) at (2,-4) {35};
\node[tnR] (rio46) at (3,-3) {46};
\node[tnB] (rio50) at (4,-2) {50};
\node[tnR] (rio52) at (5,-3) {52};
\node[tnB] (rio69) at (6,0) {69};
\node[tnB] (rio75) at (7,-2) {75};
\node[tnR] (rio85) at (8,-1) {85};
\node[tnB] (rio87) at (9,-2) {87};
\draw[te] (rio69) -- (rio28);
\draw[te] (rio28) -- (rio24);
\draw[te] (rio28) -- (rio50);
\draw[te] (rio50) -- (rio46);
\draw[te] (rio46) -- (rio35);
\draw[te] (rio50) -- (rio52);
\draw[te] (rio69) -- (rio85);
\draw[te] (rio85) -- (rio75);
\draw[te] (rio85) -- (rio87);
\node[annl, anchor=south east] at (rio35.north west) {$z$};
\node[plab] at (4.5,-4.75) {(o) 插入 52、35：情況 1};
\end{tikzpicture}\hspace{0.5cm}\begin{tikzpicture}[x=0.5cm,y=0.7cm,baseline=(current bounding box.north)]
\node[tnB] (rip24) at (0,-2) {24};
\node[tnR] (rip28) at (1,-1) {28};
\node[tnR] (rip35) at (2,-4) {35};
\node[tnB] (rip46) at (3,-3) {46};
\node[tnR] (rip50) at (4,-2) {50};
\node[tnB] (rip52) at (5,-3) {52};
\node[tnB] (rip69) at (6,0) {69};
\node[tnB] (rip75) at (7,-2) {75};
\node[tnR] (rip85) at (8,-1) {85};
\node[tnB] (rip87) at (9,-2) {87};
\draw[te] (rip69) -- (rip28);
\draw[te] (rip28) -- (rip24);
\draw[te] (rip28) -- (rip50);
\draw[te] (rip50) -- (rip46);
\draw[te] (rip46) -- (rip35);
\draw[te] (rip50) -- (rip52);
\draw[te] (rip69) -- (rip85);
\draw[te] (rip85) -- (rip75);
\draw[te] (rip85) -- (rip87);
\node[annl, anchor=south east] at (rip50.north west) {$z$};
\node[plab] at (4.5,-4.75) {(p) 改色後 $z = 50$：情況 1};
\end{tikzpicture}\\*[0.6em]
\begin{tikzpicture}[x=0.5cm,y=0.7cm,baseline=(current bounding box.north)]
\node[tnB] (riq24) at (0,-2) {24};
\node[tnB] (riq28) at (1,-1) {28};
\node[tnR] (riq35) at (2,-4) {35};
\node[tnB] (riq46) at (3,-3) {46};
\node[tnR] (riq50) at (4,-2) {50};
\node[tnB] (riq52) at (5,-3) {52};
\node[tnB] (riq69) at (6,0) {69};
\node[tnB] (riq75) at (7,-2) {75};
\node[tnB] (riq85) at (8,-1) {85};
\node[tnB] (riq87) at (9,-2) {87};
\draw[te] (riq69) -- (riq28);
\draw[te] (riq28) -- (riq24);
\draw[te] (riq28) -- (riq50);
\draw[te] (riq50) -- (riq46);
\draw[te] (riq46) -- (riq35);
\draw[te] (riq50) -- (riq52);
\draw[te] (riq69) -- (riq85);
\draw[te] (riq85) -- (riq75);
\draw[te] (riq85) -- (riq87);
\node[plab] at (4.5,-4.75) {(q) 改色、根改回黑};
\end{tikzpicture}
\diagcaption{插入 $52$、$35$：情況 1 連續兩次，改色一路到根。}
\label{fig:bal-rb-trace4}
\end{center}

\begin{itemize}
\item $52$：父 $50$ 黑，不用修。
\item $35$：父 $46$ 紅，叔叔 $52$ 紅，情況 1：$46$、$52$ 改黑，$50$ 改紅，$z$ 移到
$50$，(p)。 $50$ 的父節點 $28$ 紅，叔叔 $85$ 也紅，再一次情況 1：$28$、$85$ 改黑，
$69$ 改紅，$z$ 移到根 $69$。 最後根改回黑，(q)。
\end{itemize}

\begin{center}
\begin{tabular}{cll}
\toprule
插入 & 情況（L／R：父節點是祖父的左／右子節點） & 旋轉次數 \\
\midrule
85 & 根改黑 & 0 \\
28、75、52 & 不用修 & 0 \\
24 & 3L & 1 \\
87 & 1R，根改黑 & 0 \\
50 & 1L & 0 \\
69 & 2L、3L & 2 \\
46 & 1L、2R、3R & 2 \\
35 & 1L、1L，根改黑 & 0 \\
\bottomrule
\end{tabular}
\end{center}

最後的樹黑高是 $3$，高度 $5$。 同樣十個鍵放進 AVL 樹，高度是 $4$；紅黑樹允許比較
不整齊的形狀。

\begin{examnote}
紅黑樹的追蹤題，每一步先寫出 $z$、父節點、叔叔的顏色和用了哪個情況，再畫樹；情況
1 之後 $z$ 往上移了兩層，要用新的 $z$ 重新判斷。 考卷通常只能用一種顏色的筆，紅色
節點畫成雙圈或在旁邊標 R，並在答案開頭說明記號。
\end{examnote}

\subsection{C 程式}
\label{sec:bal-rb-c}

C 程式用 \texttt{NULL} 當 NIL，\texttt{is\_red} 把 \texttt{NULL} 當成黑色，這樣叔叔
是 NIL 的情況不用另外寫。 \texttt{replace\_child} 處理「把 $x$ 原本的位置交給另一個
節點」，左旋和右旋都用它。 \texttt{black\_count} 從下往上檢查第 4、5 條，順便算出
黑高。

\begin{example}[label=ex:bal-rb-insert]{c}{(rb\_insert.c) 紅黑樹的插入}
#include <stdio.h>
#include <stdlib.h>

typedef enum { RED, BLACK } Color;

typedef struct RBNode {
    int key;
    Color color;
    struct RBNode *left, *right, *parent;
} RBNode;

typedef struct {
    RBNode *root;
} RBTree;

// NULL 當成 NIL 葉子，NIL 是黑的。
static int is_red(const RBNode *x)
{
    return x != NULL && x->color == RED;
}

// 把 child 接到 x 原本在父節點底下的位置（x 是根時改 t->root）。
static void replace_child(RBTree *t, RBNode *x, RBNode *child)
{
    child->parent = x->parent;
    if (x->parent == NULL) {
        t->root = child;
    } else if (x == x->parent->left) {
        x->parent->left = child;
    } else {
        x->parent->right = child;
    }
}

static void rotate_left(RBTree *t, RBNode *x)   // x 的右子節點 y 轉上來
{
    RBNode *y = x->right;
    x->right = y->left;
    if (y->left != NULL) {
        y->left->parent = x;
    }
    replace_child(t, x, y);
    y->left = x;
    x->parent = y;
}

static void rotate_right(RBTree *t, RBNode *x)  // x 的左子節點 y 轉上來
{
    RBNode *y = x->left;
    x->left = y->right;
    if (y->right != NULL) {
        y->right->parent = x;
    }
    replace_child(t, x, y);
    y->right = x;
    x->parent = y;
}

static int reported = 0;        // 這次插入印了幾項

static void report(const char *what, char side)
{
    reported++;
    printf(" %s", what);
    if (side != '\0') {
        putchar(side);
    }
}

// 修正紅紅相連。印出用到的情況：1、2、3 後面的 L／R 表示父節點是祖父的
// 左／右子節點。
static void insert_fixup(RBTree *t, RBNode *z)
{
    while (is_red(z->parent)) {
        RBNode *p = z->parent, *g = p->parent;  // p 是紅的，所以不是根，g 存在
        int left = (p == g->left);
        RBNode *u = left ? g->right : g->left;  // 叔叔
        char side = left ? 'L' : 'R';
        if (is_red(u)) {                        // 情況 1：叔叔是紅的，改色
            report("1", side);
            p->color = BLACK;
            u->color = BLACK;
            g->color = RED;
            z = g;
            continue;
        }
        if (left && z == p->right) {            // 情況 2：z 在內側，先轉成外側
            report("2", side);
            z = p;
            rotate_left(t, z);
        } else if (!left && z == p->left) {
            report("2", side);
            z = p;
            rotate_right(t, z);
        }
        report("3", side);                      // 情況 3：z 在外側
        z->parent->color = BLACK;
        g->color = RED;
        if (left) {
            rotate_right(t, g);
        } else {
            rotate_left(t, g);
        }
    }
    if (t->root->color == RED) {
        report("root", '\0');
        t->root->color = BLACK;
    }
}

// 插入 key；key 已經在樹裡就不插。
void rb_insert(RBTree *t, int key)
{
    RBNode *y = NULL, *x = t->root;
    while (x != NULL) {                         // 和二元搜尋樹一樣找位置
        y = x;
        if (key == x->key) {
            return;
        }
        x = key < x->key ? x->left : x->right;
    }
    RBNode *z = malloc(sizeof *z);
    if (z == NULL) {
        exit(1);
    }
    z->key = key;
    z->color = RED;                             // 新節點一律是紅的
    z->left = z->right = NULL;
    z->parent = y;
    if (y == NULL) {
        t->root = z;
    } else if (key < y->key) {
        y->left = z;
    } else {
        y->right = z;
    }
    insert_fixup(t, z);
}

// 從 x 往下走到 NIL，路上的黑色節點數（含 x 和 NIL）。
// 有紅紅相連，或左右兩邊數出來不一樣時，回傳 -1。
static int black_count(const RBNode *x)
{
    if (x == NULL) {
        return 1;                               // NIL 是黑的
    }
    if (is_red(x) && (is_red(x->left) || is_red(x->right))) {
        return -1;
    }
    int a = black_count(x->left), b = black_count(x->right);
    if (a < 0 || a != b) {
        return -1;
    }
    return a + (x->color == BLACK);
}

void print_tree(const RBNode *x)
{
    if (x == NULL) {
        putchar('.');
        return;
    }
    printf("%d%c", x->key, x->color == RED ? 'R' : 'B');
    if (x->left != NULL || x->right != NULL) {
        putchar('(');
        print_tree(x->left);
        putchar(',');
        print_tree(x->right);
        putchar(')');
    }
}

void free_tree(RBNode *x)
{
    if (x != NULL) {
        free_tree(x->left);
        free_tree(x->right);
        free(x);
    }
}

int main(void)
{
    int keys[] = {85, 28, 24, 87, 75, 50, 69, 46, 52, 35};
    int n = sizeof keys / sizeof keys[0];
    RBTree t = {NULL};
    for (int i = 0; i < n; i++) {
        printf("insert %d:", keys[i]);
        reported = 0;
        rb_insert(&t, keys[i]);
        if (reported == 0) {
            printf(" -");
        }
        putchar('\n');
    }
    print_tree(t.root);
    putchar('\n');
    printf("black-height = %d\n", black_count(t.root) - 1);
    free_tree(t.root);
    return 0;
}
\end{example}
\iosample{}{insert 85: root \\ insert 28: - \\ insert 24: 3L \\ insert 87: 1R root \\ insert 75: - \\ insert 50: 1L \\ insert 69: 2L 3L \\ insert 46: 1L 2R 3R \\ insert 52: - \\ insert 35: 1L 1L root \\ 69B(28B(24B,50R(46B(35R,.),52B)),85B(75B,87B)) \\ black-height = 3}

輸出和上面的追蹤一致。 把 \texttt{main} 改成依序插入 $1, 2, \ldots, 10^6$，再算樹高，紅黑樹的高度
是 $37$，比 AVL 樹的 $20$（範例~\ref{ex:bal-avl-insert}）高，但仍然在定
理~\ref{thm:bal-rb-height} 的 $2 \lg(10^6+1) \approx 39.9$ 之內。

\subsection{刪除}
\label{sec:bal-rb-delete}

刪除比插入多幾種情況，但想法一樣：先照二元搜尋樹刪，再看壞了哪一條性質。

\textbf{第一步：決定真正拿掉的節點。} 要刪的節點 $z$ 有兩個子節點時，把後繼 $y$
（右子樹最小的節點）的鍵複製到 $z$，改成拿掉 $y$。 $y$ 沒有左子節點，所以真正被拿
掉的節點最多只有一個不是 NIL 的子節點。 下面說的「被拿掉的節點」都是指這一個，它
的顏色才是要看的顏色；$z$ 留在原位，顏色不變。 被拿掉的節點的子節點（可能是 NIL）
叫 $x$，$x$ 接到它原本的位置；$x$ 新的父節點叫 $p$。

\textbf{第二步：看被拿掉的節點的顏色。}
\begin{itemize}
\item 它是紅的：什麼都不用做。 拿掉紅色節點不改變黑色個數，它的父節點和 $x$ 都是
黑的，接起來也不會紅紅相連。
\item 它是黑的、$x$ 是紅的：把 $x$ 改黑，補回少掉的那個黑色，結束。
\item 它是黑的、$x$ 也是黑的（包括 $x$ 是 NIL）：經過 $x$ 的路徑都少了一個黑色，違
反第 5 條。 這是要修正的情況。
\end{itemize}

修正時，想像 $x$ 身上多背了一個黑色（叫它\textbf{雙重黑}），只要能把這個多出來的黑
色放到某個紅色節點上（紅改黑），或是一路推到根（直接丟掉，所有路徑一起少一個），就
修好了。 設 $x$ 是 $p$ 的左子節點，$s$ 是 $x$ 的兄弟（$p$ 的右子節點）。 $s$ 不會是
NIL：經過 $s$ 的路徑比經過 $x$ 的多一個黑色，至少有兩個黑色，所以 $s$ 是內部節點。
$s$ 靠近 $x$ 的子節點叫 $n$（$s$ 的左子節點），遠的叫 $f$（$s$ 的右子節點）。

\begin{itemize}
\item \textbf{情況 1：$s$ 是紅的。} 這時 $p$、$n$、$f$ 都是黑的。 $s$ 改黑、$p$ 改
紅，對 $p$ 左旋。 黑色個數不變，$x$ 的新兄弟是原本的 $n$，黑的，變成情況 2、3、4
之一。
\item \textbf{情況 2：$s$ 是黑的，$n$、$f$ 都是黑的。} $s$ 改紅，經過 $s$ 的路徑也少
一個黑色，和 $x$ 這邊扯平，但整棵 $p$ 的子樹少了一個黑色。 把雙重黑交給 $p$：$x$
移到 $p$。 如果 $p$ 是紅的（例如剛做完情況 1），迴圈結束，最後把它改黑就補回來
了；如果 $p$ 是黑的，往上繼續。
\item \textbf{情況 3：$s$ 是黑的，$n$ 是紅的，$f$ 是黑的。} $n$ 改黑、$s$ 改紅，對
$s$ 右旋。 $n$ 升上來當 $x$ 的新兄弟，它的遠側子節點是紅的 $s$，變成情況 4。
\item \textbf{情況 4：$s$ 是黑的，$f$ 是紅的。} $s$ 改成 $p$ 的顏色，$p$ 改黑，$f$
改黑，對 $p$ 左旋。 $x$ 這邊多了一個黑色的 $p$，抵掉雙重黑；$s$ 那邊少了 $s$ 這個
黑色，由 $f$ 改黑補回來；最上面的節點 $s$ 用的是原本 $p$ 的顏色，對更上面沒有影
響。 修正結束。
\end{itemize}

$x$ 是右子節點時，以上全部左右對調。

\begin{center}
\begin{tikzpicture}[x=0.64cm, y=0.78cm]
% 情況 1
\node[font=\small\bfseries, anchor=west] at (-2.2,0.7) {情況 1};
\node[tnB] (p) at (0,0) {$p$};
\node[tnB] (x) at (-1.2,-1) {$x$};
\node[tnR] (s) at (1.2,-1) {$s$};
\node[tnB] (n) at (0.6,-2) {$n$};
\node[tnB] (f) at (1.8,-2) {$f$};
\draw[te] (p) -- (x); \draw[te] (p) -- (s); \draw[te] (s) -- (n); \draw[te] (s) -- (f);
\draw[tarr] (2.6,-1) -- node[above, font=\scriptsize] {左旋 $p$} (4.0,-1);
\begin{scope}[xshift=5.6cm]
\node[tnB] (s) at (0,0) {$s$};
\node[tnR] (p) at (-1.2,-1) {$p$};
\node[tnB] (f) at (1.2,-1) {$f$};
\node[tnB] (x) at (-1.8,-2) {$x$};
\node[tnB] (n) at (-0.6,-2) {$n$};
\draw[te] (s) -- (p); \draw[te] (s) -- (f); \draw[te] (p) -- (x); \draw[te] (p) -- (n);
\node[annl, anchor=north] at (n.south) {新的 $s$};
\end{scope}
% 情況 2
\begin{scope}[xshift=8.3cm]
\node[font=\small\bfseries, anchor=west] at (-2.2,0.7) {情況 2};
\node[tnG] (p) at (0,0) {$p$};
\node[tnB] (x) at (-1.2,-1) {$x$};
\node[tnB] (s) at (1.2,-1) {$s$};
\node[tnB] (n) at (0.6,-2) {$n$};
\node[tnB] (f) at (1.8,-2) {$f$};
\draw[te] (p) -- (x); \draw[te] (p) -- (s); \draw[te] (s) -- (n); \draw[te] (s) -- (f);
\draw[tarr] (2.6,-1) -- node[above, font=\scriptsize] {改色} (4.0,-1);
\begin{scope}[xshift=4.0cm]
\node[tnG] (p) at (0,0) {$p$};
\node[tnB] (x) at (-1.2,-1) {};
\node[tnR] (s) at (1.2,-1) {$s$};
\node[tnB] (n) at (0.6,-2) {$n$};
\node[tnB] (f) at (1.8,-2) {$f$};
\draw[te] (p) -- (x); \draw[te] (p) -- (s); \draw[te] (s) -- (n); \draw[te] (s) -- (f);
\node[annl, anchor=south east] at (p.north west) {新的 $x$};
\end{scope}
\end{scope}
% 情況 3
\begin{scope}[yshift=-3.6cm]
\node[font=\small\bfseries, anchor=west] at (-2.2,0.7) {情況 3};
\node[tnG] (p) at (0,0) {$p$};
\node[tnB] (x) at (-1.2,-1) {$x$};
\node[tnB] (s) at (1.2,-1) {$s$};
\node[tnR] (n) at (0.6,-2) {$n$};
\node[tnB] (f) at (1.8,-2) {$f$};
\draw[te] (p) -- (x); \draw[te] (p) -- (s); \draw[te] (s) -- (n); \draw[te] (s) -- (f);
\draw[tarr] (2.6,-1) -- node[above, font=\scriptsize] {右旋 $s$} (4.0,-1);
\begin{scope}[xshift=5.0cm]
\node[tnG] (p) at (0,0) {$p$};
\node[tnB] (x) at (-1.2,-1) {$x$};
\node[tnB] (n) at (1.2,-1) {$n$};
\node[tnR] (s) at (1.8,-2) {$s$};
\node[tnB] (f) at (2.4,-3) {$f$};
\draw[te] (p) -- (x); \draw[te] (p) -- (n); \draw[te] (n) -- (s); \draw[te] (s) -- (f);
\node[annl, anchor=south west] at (n.north east) {新的 $s$};
\end{scope}
\end{scope}
% 情況 4
\begin{scope}[xshift=8.3cm, yshift=-3.6cm]
\node[font=\small\bfseries, anchor=west] at (-2.2,0.7) {情況 4};
\node[tnG] (p) at (0,0) {$p$};
\node[tnB] (x) at (-1.2,-1) {$x$};
\node[tnB] (s) at (1.2,-1) {$s$};
\node[tnG] (n) at (0.6,-2) {$n$};
\node[tnR] (f) at (1.8,-2) {$f$};
\draw[te] (p) -- (x); \draw[te] (p) -- (s); \draw[te] (s) -- (n); \draw[te] (s) -- (f);
\draw[tarr] (2.6,-1) -- node[above, font=\scriptsize] {左旋 $p$} (4.0,-1);
\begin{scope}[xshift=5.6cm]
\node[tnG] (s) at (0,0) {$s$};
\node[tnB] (p) at (-1.2,-1) {$p$};
\node[tnB] (f) at (1.2,-1) {$f$};
\node[tnB] (x) at (-1.8,-2) {$x$};
\node[tnG] (n) at (-0.6,-2) {$n$};
\draw[te] (s) -- (p); \draw[te] (s) -- (f); \draw[te] (p) -- (x); \draw[te] (p) -- (n);
\end{scope}
\end{scope}
\end{tikzpicture}
\diagcaption{刪除修正的四種情況（$x$ 是 $p$ 的左子節點）。 $x$ 是雙重黑；灰色表示
紅黑都可以，情況 4 右邊的 $s$ 用的是左邊 $p$ 的顏色。 $n$、$f$ 底下的子樹和 $x$ 的
子樹都沒有畫，它們跟著原本的父節點走，$n$ 的子樹在情況 3 會分給 $s$ 和 $n$，照
旋轉的規則（圖~\ref{fig:bal-rotation}）接。}
\label{fig:bal-rb-delete-cases}
\end{center}

\begin{pseudo}[label=alg:bal-rb-delete]{紅黑樹的刪除}
\Procedure{RB-Delete}{$T, z$}
    \If{$z.\mathit{left} \ne \NIL$ \AND $z.\mathit{right} \ne \NIL$}
        \State $y \gets z.\mathit{right}$ 一路往左走到底的節點
        \State $z.\mathit{key} \gets y.\mathit{key}$
        \State $z \gets y$ \Comment{改成拿掉後繼}
    \EndIf
    \State $x \gets z$ 不是 \NIL 的那個子節點（兩個都是 \NIL 時 $x = \NIL$）
    \State $p \gets z.p$
    \State 把 $x$ 接到 $z$ 原本的位置（$x \ne \NIL$ 時設 $x.p \gets p$；$p = \NIL$ 時 $T.\mathit{root} \gets x$）
    \If{$z$ 是黑的}
        \State \Call{RB-Delete-Fixup}{$T, x, p$}
    \EndIf
    \State 釋放 $z$
\EndProcedure
\end{pseudo}

\begin{pseudo}[label=alg:bal-rb-delfix]{刪除後的修正，$x$ 可能是 \NIL，所以另外傳父節點 $p$}
\Procedure{RB-Delete-Fixup}{$T, x, p$}
    \While{$x \ne T.\mathit{root}$ \AND $x$ 是黑的}
        \If{$x = p.\mathit{left}$}
            \State $s \gets p.\mathit{right}$
            \If{$s$ 是紅的} \Comment{情況 1}
                \State $s$ 改黑，$p$ 改紅
                \State \Call{Left-Rotate}{$T, p$}
                \State $s \gets p.\mathit{right}$
            \EndIf
            \If{$s.\mathit{left}$、$s.\mathit{right}$ 都是黑的} \Comment{情況 2}
                \State $s$ 改紅
                \State $x \gets p$
                \State $p \gets x.p$
            \Else
                \If{$s.\mathit{right}$ 是黑的} \Comment{情況 3}
                    \State $s.\mathit{left}$ 改黑，$s$ 改紅
                    \State \Call{Right-Rotate}{$T, s$}
                    \State $s \gets p.\mathit{right}$
                \EndIf
                \State $s$ 改成 $p$ 的顏色，$p$ 改黑，$s.\mathit{right}$ 改黑 \Comment{情況 4}
                \State \Call{Left-Rotate}{$T, p$}
                \State $x \gets T.\mathit{root}$
            \EndIf
        \Else
            \State 和第 4 到 20 行相同，$\mathit{left}$、$\mathit{right}$ 對調，\proc{Left-Rotate}、\proc{Right-Rotate} 對調
        \EndIf
    \EndWhile
    \If{$x \ne \NIL$}
        \State $x$ 改黑
    \EndIf
\EndProcedure
\end{pseudo}

\proc{RB-Delete} 只要被拿掉的節點是黑的就呼叫修正。 如果 $x$ 是紅的，第 2 行的迴圈
一次都不做，直接跳到最後把 $x$ 改黑，就是第二步的第二種情況。 NIL 一律當成黑的；
$x$ 是 NIL 時，第 3 行的 $x = p.\mathit{left}$ 看的是 $p$ 的左邊是不是 NIL，因為兄
弟 $s$ 一定不是 NIL，不會兩邊都是 NIL。

\textbf{成本。} 情況 2 讓 $x$ 往上一層，最多 $O(\lg n)$ 次，只改顏色。 其他情況的轉
移是：情況 1 之後是 2、3 或 4，而且情況 1 之後如果是情況 2，$p$ 是紅的，迴圈馬上結
束；情況 3 之後一定是情況 4；情況 4 做完就結束。 所以刪除一共最多\textbf{三次旋
轉}（1、3、4 各一次），時間 $O(\lg n)$。 這是紅黑樹和 AVL 樹的一個差別：AVL 刪除
可能每一層都要旋轉。

\subsection{刪除的完整例子}
\label{sec:bal-rb-deltrace}

從第~\ref{sec:bal-rb-trace}~節最後的樹刪除 $24$。 這個例子依序用到情況 1、3、4。

\begin{center}
\begin{tikzpicture}[x=0.41cm,y=0.7cm,baseline=(current bounding box.north)]
\node[tnB] (rda24) at (0,-2) {24};
\node[tnB] (rda28) at (1,-1) {28};
\node[tnR] (rda35) at (2,-4) {35};
\node[tnB] (rda46) at (3,-3) {46};
\node[tnR] (rda50) at (4,-2) {50};
\node[tnB] (rda52) at (5,-3) {52};
\node[tnB] (rda69) at (6,0) {69};
\node[tnB] (rda75) at (7,-2) {75};
\node[tnB] (rda85) at (8,-1) {85};
\node[tnB] (rda87) at (9,-2) {87};
\draw[te] (rda69) -- (rda28);
\draw[te] (rda28) -- (rda24);
\draw[te] (rda28) -- (rda50);
\draw[te] (rda50) -- (rda46);
\draw[te] (rda46) -- (rda35);
\draw[te] (rda50) -- (rda52);
\draw[te] (rda69) -- (rda85);
\draw[te] (rda85) -- (rda75);
\draw[te] (rda85) -- (rda87);
\node[annl, anchor=south east] at (rda24.north west) {刪除};
\node[plab] at (4.5,-4.75) {(a) 刪除 24 之前};
\end{tikzpicture}\hspace{0.4cm}\begin{tikzpicture}[x=0.41cm,y=0.7cm,baseline=(current bounding box.north)]
\node[tnil] (rdbNIL) at (0,-2) {};
\node[tnB] (rdb28) at (1,-1) {28};
\node[tnR] (rdb35) at (2,-4) {35};
\node[tnB] (rdb46) at (3,-3) {46};
\node[tnR] (rdb50) at (4,-2) {50};
\node[tnB] (rdb52) at (5,-3) {52};
\node[tnB] (rdb69) at (6,0) {69};
\node[tnB] (rdb75) at (7,-2) {75};
\node[tnB] (rdb85) at (8,-1) {85};
\node[tnB] (rdb87) at (9,-2) {87};
\draw[te] (rdb69) -- (rdb28);
\draw[te] (rdb28) -- (rdbNIL);
\draw[te] (rdb28) -- (rdb50);
\draw[te] (rdb50) -- (rdb46);
\draw[te] (rdb46) -- (rdb35);
\draw[te] (rdb50) -- (rdb52);
\draw[te] (rdb69) -- (rdb85);
\draw[te] (rdb85) -- (rdb75);
\draw[te] (rdb85) -- (rdb87);
\node[annl, anchor=south east] at (rdbNIL.north west) {$x$};
\node[annl, anchor=south east] at (rdb50.north west) {$s$};
\node[plab] at (4.5,-4.75) {(b) 拿掉 24：情況 1};
\end{tikzpicture}\hspace{0.4cm}\begin{tikzpicture}[x=0.41cm,y=0.7cm,baseline=(current bounding box.north)]
\node[tnil] (rdcNIL) at (0,-3) {};
\node[tnR] (rdc28) at (1,-2) {28};
\node[tnR] (rdc35) at (2,-4) {35};
\node[tnB] (rdc46) at (3,-3) {46};
\node[tnB] (rdc50) at (4,-1) {50};
\node[tnB] (rdc52) at (5,-2) {52};
\node[tnB] (rdc69) at (6,0) {69};
\node[tnB] (rdc75) at (7,-2) {75};
\node[tnB] (rdc85) at (8,-1) {85};
\node[tnB] (rdc87) at (9,-2) {87};
\draw[te] (rdc69) -- (rdc50);
\draw[te] (rdc50) -- (rdc28);
\draw[te] (rdc28) -- (rdcNIL);
\draw[te] (rdc28) -- (rdc46);
\draw[te] (rdc46) -- (rdc35);
\draw[te] (rdc50) -- (rdc52);
\draw[te] (rdc69) -- (rdc85);
\draw[te] (rdc85) -- (rdc75);
\draw[te] (rdc85) -- (rdc87);
\node[annl, anchor=south east] at (rdcNIL.north west) {$x$};
\node[annl, anchor=south east] at (rdc46.north west) {$s$};
\node[plab] at (4.5,-4.75) {(c) 情況 1 之後：情況 3};
\end{tikzpicture}\\*[0.6em]
\begin{tikzpicture}[x=0.41cm,y=0.7cm,baseline=(current bounding box.north)]
\node[tnil] (rddNIL) at (0,-3) {};
\node[tnR] (rdd28) at (1,-2) {28};
\node[tnB] (rdd35) at (2,-3) {35};
\node[tnR] (rdd46) at (3,-4) {46};
\node[tnB] (rdd50) at (4,-1) {50};
\node[tnB] (rdd52) at (5,-2) {52};
\node[tnB] (rdd69) at (6,0) {69};
\node[tnB] (rdd75) at (7,-2) {75};
\node[tnB] (rdd85) at (8,-1) {85};
\node[tnB] (rdd87) at (9,-2) {87};
\draw[te] (rdd69) -- (rdd50);
\draw[te] (rdd50) -- (rdd28);
\draw[te] (rdd28) -- (rddNIL);
\draw[te] (rdd28) -- (rdd35);
\draw[te] (rdd35) -- (rdd46);
\draw[te] (rdd50) -- (rdd52);
\draw[te] (rdd69) -- (rdd85);
\draw[te] (rdd85) -- (rdd75);
\draw[te] (rdd85) -- (rdd87);
\node[annl, anchor=south east] at (rddNIL.north west) {$x$};
\node[annl, anchor=south east] at (rdd35.north west) {$s$};
\node[plab] at (4.5,-4.75) {(d) 情況 3 之後：情況 4};
\end{tikzpicture}\hspace{0.6cm}\begin{tikzpicture}[x=0.41cm,y=0.7cm,baseline=(current bounding box.north)]
\node[tnB] (rde28) at (0,-3) {28};
\node[tnR] (rde35) at (1,-2) {35};
\node[tnB] (rde46) at (2,-3) {46};
\node[tnB] (rde50) at (3,-1) {50};
\node[tnB] (rde52) at (4,-2) {52};
\node[tnB] (rde69) at (5,0) {69};
\node[tnB] (rde75) at (6,-2) {75};
\node[tnB] (rde85) at (7,-1) {85};
\node[tnB] (rde87) at (8,-2) {87};
\draw[te] (rde69) -- (rde50);
\draw[te] (rde50) -- (rde35);
\draw[te] (rde35) -- (rde28);
\draw[te] (rde35) -- (rde46);
\draw[te] (rde50) -- (rde52);
\draw[te] (rde69) -- (rde85);
\draw[te] (rde85) -- (rde75);
\draw[te] (rde85) -- (rde87);
\node[plab] at (4.0,-4.75) {(e) 情況 4 之後};
\end{tikzpicture}
\diagcaption{刪除 $24$。 小方塊是 $x$ 所在的 NIL，$s$ 是它的兄弟。}
\label{fig:bal-rb-del}
\end{center}

\begin{enumerate}
\item $24$ 是葉節點，被拿掉的就是它自己，黑色。 $x$ 是它留下的 NIL，黑的，所以 $x$ 是
雙重黑，$p = 28$。 $x$ 是 $p$ 的左子節點，兄弟 $s = 50$ 是紅的：\textbf{情況 1}。
$50$ 改黑、$28$ 改紅，左旋 $28$，(c)。 這一步之後 $x$ 還是 $28$ 的左子節點，新的兄
弟是原本 $50$ 的左子節點 $46$。
\item $s = 46$ 是黑的，近的子節點 $n = 35$ 是紅的，遠的子節點 $f$ 是 NIL（黑）：
\textbf{情況 3}。 $35$ 改黑、$46$ 改紅，右旋 $46$，(d)。 新的兄弟是 $35$。
\item $s = 35$ 是黑的，遠的子節點 $f = 46$ 是紅的：\textbf{情況 4}。 $35$ 改成 $p =
28$ 的顏色（紅），$28$ 改黑，$46$ 改黑，左旋 $28$，(e)。 $x$ 設成根，迴圈結束。
\end{enumerate}

檢查第 5 條：從根 $69$ 往下，經過 $50 \to 35 \to 28 \to \text{NIL}$ 的黑色是 $50$、
$28$、NIL，經過 $85 \to 75 \to \text{NIL}$ 的是 $85$、$75$、NIL，都是 $3$ 個，黑高
沒變。 這次刪除用了三次旋轉，正好是上限。

\section{紅黑樹和 2-3-4 樹}
\label{sec:bal-234}

把紅黑樹裡每個黑色節點和它的紅色子節點合併成一個大節點。 例如圖~\ref{fig:bal-rb-del}(a)
的 $28$ 和它的紅色子節點 $50$ 合成一個存 $\{28, 50\}$ 的節點，$46$ 和紅色的 $35$ 合成
$\{35, 46\}$，其他黑色節點各自一個。 合併之後，每個大節點存 $1$ 到 $3$ 個鍵、有 $2$ 到 $4$ 個子
節點，而且由第 5 條，所有的 NIL 都在同一層。 這正是第~11~章要講的 \textbf{2-3-4
樹}，紅黑樹是用二元樹來存它的一種方法，樹的黑高就是 2-3-4 樹的高度。

從這個角度看，插入的情況 1 對應 2-3-4 樹把滿了的節點分裂、中間的鍵往上送；情況 2、
3 只是在同一個大節點內部換一種二元的排法。 刪除的情況 2 對應兩個節點合併，情況 3、
4 對應向兄弟借一個鍵，情況 1 只是換一種二元的排法。 詳細的對應放在第~\ref{sec:mw-rb}~節，讀完 2-3-4 樹的分裂、合併和
借之後再回來看這一章的情況表，會比較好記。

\section{AVL 樹和紅黑樹的比較}
\label{sec:bal-compare}

\begin{center}
\begin{tabular}{p{0.24\linewidth}p{0.33\linewidth}p{0.33\linewidth}}
\toprule
 & AVL 樹 & 紅黑樹 \\
\midrule
平衡條件 & 每個節點左右高度差 $\le 1$ & 五條顏色性質 \\
高度上界 & $\log_\varphi(n+1) < 1.4405\lg(n+1)$ & $2\lg(n+1)$ \\
每個節點多存 & 高度（或 $2$ 位元的平衡因子） & $1$ 位元的顏色，加上父指標 \\
搜尋 & $O(\lg n)$，樹比較矮 & $O(\lg n)$ \\
插入的旋轉 & 最多 $2$ 次 & 最多 $2$ 次 \\
刪除的旋轉 & 最壞 $\Theta(\lg n)$ 次 & 最多 $3$ 次 \\
依序插入 $10^6$ 個鍵的高度 & $20$ & $37$ \\
\bottomrule
\end{tabular}
\end{center}

兩種樹的每個操作都是 $O(\lg n)$，差別在常數。 AVL 樹比較矮，搜尋多的時候稍微快一
點；紅黑樹的條件比較寬，插入和刪除時改結構的次數比較少，改顏色只動一個位元。 很多
程式庫的有序容器用紅黑樹，例如 Java 的 \texttt{TreeMap}、Linux 核心的 rbtree，以及
多數 C++ 標準函式庫的 \texttt{std::map}。 搜尋遠多於更新的資料（例如建好之後很少改
的索引）用 AVL 樹也很合適。

另外要記住，「AVL 樹比紅黑樹矮」說的是最壞情況的上界。 同樣 $n$ 個節點，某一棵 AVL
樹可能比某一棵紅黑樹高（觀念題第 1 題）。

\section{本章常見錯誤}
\label{sec:bal-pitfalls}

\begin{center}
\begin{tabular}{p{0.24\linewidth}p{0.50\linewidth}p{0.16\linewidth}}
\toprule
錯誤 & 例子 & 見 \\
\midrule
從根找失衡節點 & 插入後從根往下看，把根當成 $A$；應該取離新節點最近的那個 &
第~\ref{sec:bal-avl-insert}~節 \\
LR、RL 只轉一次 & 對 $A$ 直接右旋，$C$ 的子樹跟著掛到 $A$ 底下，換成另一邊太高 &
第~\ref{sec:bal-avl-insert}~節 \\
刪除時 $\mathrm{BF}(B) = 0$ 用雙旋轉 & \proc{Rebalance} 的條件寫成 $\le 0$；應該做
單旋轉 & 第~\ref{sec:bal-avl-delete}~節 \\
刪除轉完一次就停 & 以為刪除和插入一樣最多轉一次，沒有繼續往上檢查 &
第~\ref{sec:bal-avl-deltrace}~節 \\
旋轉後高度更新順序錯 & 先更新升上來的節點，它用到降下去的節點的舊高度 &
第~\ref{sec:bal-avl-c}~節 \\
沒接回旋轉後的根 & C 程式寫 \texttt{rotate\_right(t)}，沒有寫
\texttt{t = rotate\_right(t)}；紅黑樹旋轉時忘了改父指標或 \texttt{root} &
第~\ref{sec:bal-avl-c}~節、\newline 第~\ref{sec:bal-rb-c}~節 \\
高度的定義混用 & 「高度 $h$ 最少幾個節點」用數邊的定義算，答案差一格 &
第~\ref{sec:bal-avl-height}~節 \\
新節點塗黑 & 插入時新節點用黑色，第 5 條一定壞 & 第~\ref{sec:bal-rb-insert}~節 \\
只看父節點不看叔叔 & 父節點紅就直接旋轉；叔叔是紅的時候應該改色 &
第~\ref{sec:bal-rb-insert}~節 \\
情況 1 之後沒有往上 & 改完色就停，沒有檢查 $g$ 和它的父節點 &
第~\ref{sec:bal-rb-trace}~節 \\
看錯被拿掉的節點 & 有兩個子節點時，看 $z$ 的顏色決定要不要修正；應該看後繼的顏色 &
第~\ref{sec:bal-rb-delete}~節 \\
黑高算錯 & 把 $x$ 自己算進去，或漏了 NIL & 第~\ref{sec:bal-rb}~節 \\
\bottomrule
\end{tabular}
\end{center}

\section{觀念題}
\label{sec:bal-exercises}

題型和研究所考試、期中期末考的筆試相近，有選擇、計算、追蹤、證明和設計。 除非題目
另外說明，高度照本章的定義（數節點），刪除有兩個子節點的節點時用後繼取代。 解答在
本章最後。

\begin{enumerate}
\item （選擇，可能不只一個答案）下列哪些敘述正確？
\begin{enumerate}
\item AVL 樹插入一個鍵，最多做一次單旋轉或一次雙旋轉。
\item AVL 樹刪除一個鍵，最多做一次單旋轉或一次雙旋轉。
\item 紅黑樹裡，從根到任一個 NIL 的路徑上，紅色節點的個數不超過黑色內部節點的個數。
\item 每一棵紅黑樹都滿足 AVL 樹的平衡條件。
\item 同樣 $n$ 個鍵，AVL 樹的高度一定不超過紅黑樹的高度。
\end{enumerate}

\item （計算）
\begin{enumerate}
\item 高度 $7$ 的 AVL 樹最少有幾個節點？ 最多有幾個？
\item 有 $60$ 個節點的 AVL 樹，高度最大和最小各是多少？
\item 如果改用數邊的定義（葉節點高度 $0$），「高度 $7$ 的 AVL 樹」最少和最多各有幾個節點？
\end{enumerate}

\item （追蹤）依序把 $72, 98, 90, 41, 82, 80, 64, 70, 36$ 插入空的 AVL 樹。 寫出每次
旋轉的失衡節點和類型，並畫出每次旋轉之後的樹。

\item （追蹤）對下面的 AVL 樹依序刪除 $47, 71, 93, 12$。 畫出每次刪除之後的樹，有旋
轉的話寫出失衡的節點和類型（用 R0、R1、R$-1$、L0、L1、L$-1$ 的名稱）。
\begin{center}
\begin{tikzpicture}[x=0.5cm,y=0.8cm,baseline=(current bounding box.north)]
\node[tn] (x4t12) at (0,-3) {12};
\node[tn] (x4t18) at (1,-2) {18};
\node[tn] (x4t26) at (2,-1) {26};
\node[tn] (x4t31) at (3,-3) {31};
\node[tn] (x4t35) at (4,-2) {35};
\node[tn] (x4t40) at (5,-3) {40};
\node[tn] (x4t47) at (6,0) {47};
\node[tn] (x4t62) at (7,-2) {62};
\node[tn] (x4t71) at (8,-1) {71};
\node[tn] (x4t84) at (9,-2) {84};
\node[tn] (x4t93) at (10,-3) {93};
\draw[te] (x4t47) -- (x4t26);
\draw[te] (x4t26) -- (x4t18);
\draw[te] (x4t18) -- (x4t12);
\draw[te] (x4t26) -- (x4t35);
\draw[te] (x4t35) -- (x4t31);
\draw[te] (x4t35) -- (x4t40);
\draw[te] (x4t47) -- (x4t71);
\draw[te] (x4t71) -- (x4t62);
\draw[te] (x4t71) -- (x4t84);
\draw[te] (x4t84) -- (x4t93);
\end{tikzpicture}
\end{center}

\item （計算）AVL 樹插入時發生 LR，符號同圖~\ref{fig:bal-ll-lr}。
\begin{enumerate}
\item 分成三種情形：新節點在 $T_2$、新節點在 $T_3$、$C$ 本身就是新節點。 各寫出雙旋
轉之後 $A$、$B$、$C$ 的平衡因子。
\item 如果每個節點不存高度、只存平衡因子，雙旋轉之後要怎麼更新三個節點的平衡因子？
用旋轉前 $\mathrm{BF}(C)$ 的值寫出規則。
\end{enumerate}

\item （證明）AVL 樹插入新節點後，從它的父節點往上更新平衡因子。 證明：如果某個節點
$v$ 的平衡因子更新後變成 $0$，$v$ 的所有祖先的平衡因子都和插入前一樣，可以停止更新。
再說明更新後變成 $\pm 1$ 時為什麼不能停。

\item （設計）給一棵節點個數最少的 AVL 樹和一個要刪除的鍵，使刪除後發生
$\mathrm{BF}(B) = 0$ 的單旋轉（R0 或 L0）。 畫出刪除前後的樹，並說明為什麼節點不能更少。

\item （計算）一棵紅黑樹根的黑高是 $k$（$k \ge 1$）。
\begin{enumerate}
\item 它最少有幾個內部節點？ 最多有幾個？ 證明你的答案，並說明怎樣的樹可以達到。
\item $k = 2$ 時，最高可能是多少？
\end{enumerate}

\item （證明）證明在紅黑樹裡，對任何節點 $x$，從 $x$ 往下走到 NIL 的最長路徑，邊數不
超過最短路徑的兩倍。

\item （選擇）下面四棵樹（深色是黑節點，紅底是紅節點，NIL 沒有畫），哪些是合法的紅黑
樹？ 不合法的違反了哪一條？
\begin{center}
\begin{tikzpicture}[x=0.55cm,y=0.7cm,baseline=(current bounding box.north)]
\node[tnB] (q10a10) at (0,-2) {10};
\node[tnR] (q10a15) at (1,-1) {15};
\node[tnB] (q10a20) at (2,-2) {20};
\node[tnB] (q10a30) at (3,0) {30};
\node[tnB] (q10a40) at (4,-1) {40};
\draw[te] (q10a30) -- (q10a15);
\draw[te] (q10a15) -- (q10a10);
\draw[te] (q10a15) -- (q10a20);
\draw[te] (q10a30) -- (q10a40);
\node[plab] at (2.0,-2.75) {(A)};
\end{tikzpicture}\hspace{0.6cm}\begin{tikzpicture}[x=0.55cm,y=0.7cm,baseline=(current bounding box.north)]
\node[tnR] (q10b10) at (0,-2) {10};
\node[tnB] (q10b15) at (1,-1) {15};
\node[tnB] (q10b30) at (2,0) {30};
\node[tnR] (q10b40) at (3,-1) {40};
\draw[te] (q10b30) -- (q10b15);
\draw[te] (q10b15) -- (q10b10);
\draw[te] (q10b30) -- (q10b40);
\node[plab] at (1.5,-2.75) {(B)};
\end{tikzpicture}\hspace{0.6cm}\begin{tikzpicture}[x=0.55cm,y=0.7cm,baseline=(current bounding box.north)]
\node[tnR] (q10c10) at (0,-2) {10};
\node[tnR] (q10c15) at (1,-1) {15};
\node[tnB] (q10c30) at (2,0) {30};
\node[tnR] (q10c40) at (3,-1) {40};
\draw[te] (q10c30) -- (q10c15);
\draw[te] (q10c15) -- (q10c10);
\draw[te] (q10c30) -- (q10c40);
\node[plab] at (1.5,-2.75) {(C)};
\end{tikzpicture}\hspace{0.6cm}\begin{tikzpicture}[x=0.55cm,y=0.7cm,baseline=(current bounding box.north)]
\node[tnB] (q10d15) at (0,-1) {15};
\node[tnB] (q10d30) at (1,0) {30};
\node[tnB] (q10d35) at (2,-2) {35};
\node[tnR] (q10d40) at (3,-1) {40};
\node[tnB] (q10d50) at (4,-2) {50};
\draw[te] (q10d30) -- (q10d15);
\draw[te] (q10d30) -- (q10d40);
\draw[te] (q10d40) -- (q10d35);
\draw[te] (q10d40) -- (q10d50);
\node[plab] at (2.0,-2.75) {(D)};
\end{tikzpicture}
\end{center}

\item （追蹤）
\begin{enumerate}
\item 依序把 $62, 90, 31, 13, 25, 45, 43, 49$ 插入空的紅黑樹。 寫出每次插入用到的情況，
並畫出插入 $13$、$25$、$43$、$49$ 之後的樹。
\item 接著從 (a) 最後的樹依序刪除 $31$、$62$，畫出每次刪除之後的樹，寫出用到的情況。
\end{enumerate}

\item （簡答）
\begin{enumerate}
\item 紅黑樹插入時，如果把新節點塗成黑色，哪一條性質一定會壞（樹原本不是空的）？
為什麼這比塗成紅色難修？
\item 情況 1 改色之後，為什麼下一輪要檢查的是 $g$，而不是 $p$？
\end{enumerate}

\item （計算）有 $10^6$ 個節點。 (a) AVL 樹的高度最大是多少？ 用 $N(h)$ 和費氏數算出
確切的值。 (b) 由定理~\ref{thm:bal-rb-height}，紅黑樹的高度不超過多少？ (c) 任何二元
樹的高度至少是多少？

\end{enumerate}

\section{實作練習}
\label{sec:bal-practice}

下面是 Aki-Judge 上和本章相關的題目，都是 AVL 樹。 五題用同一組介面：交
\texttt{avl.h}、\texttt{avl.c} 兩個檔案，實作初始化、插入、搜尋、刪除、中序輸出和
高度六個函式，題目要求每次插入、刪除之後樹都維持平衡，測資會在操作之間查詢樹高。 題目的高度和本章的定義
相同（空樹是 $0$）。 範例~\ref{ex:bal-avl-insert} 的插入可以直接用，刪除照虛擬
碼~\ref{alg:bal-avl-delete} 自己寫。 紅黑樹目前沒有對應的題目，請用觀念題第 10、11
題練習。

\begin{center}
\begin{tabular}{llc}
\toprule
小節 & 題目 & 難度 \\
\midrule
插入 & \akiref{avl-insert-and-search}{AVL 樹：插入與搜尋} & 基本 \\
插入 & \akiref{avl-sorted-traversal}{AVL 樹：中序走訪輸出} & 基本 \\
高度上界 & \akiref{avl-height-after-sorted-inserts}{AVL 樹：循序插入後的高度} & 基本 \\
刪除 & \akiref{avl-deletion}{AVL 樹：刪除} & 進階 \\
綜合 & \akiref{avl-long-mixed-workloads}{AVL 樹：長串混合操作} & 進階 \\
\bottomrule
\end{tabular}
\end{center}

寫刪除時，先用第~\ref{sec:bal-avl-deltrace}~節的兩個例子自己測：每刪一個鍵就印出整棵
樹，和圖~\ref{fig:bal-avl-del}、圖~\ref{fig:bal-avl-cascade} 比對。 「刪除」那題要
求刪除不存在的鍵時樹完全不變，虛擬碼~\ref{alg:bal-avl-delete} 的第 2、3 行已經處理。

\section{讀完你應該會}
\label{sec:bal-checklist}

\begin{itemize}
\item 能畫出右旋和左旋前後的樹，說出哪棵子樹換了父節點，並說明旋轉為什麼不改變中序。
\item 能用 $N(h) = N(h-1) + N(h-2) + 1$ 和費氏數，算出 $n$ 個節點的 AVL 樹最高和最矮
的高度，並證明 $h < 1.4405 \lg(n+1)$。
\item 能對一串鍵做 AVL 樹的插入和刪除，找出失衡的節點、判斷 LL／RR／LR／RL（刪除時
R0、R1、R$-1$ 等），畫出每次旋轉之後的樹。
\item 能寫出 AVL 樹插入的 C 程式，用遞迴回傳新的根，旋轉後依正確順序更新高度。
\item 能說出紅黑樹的五條性質、算黑高，並證明 $h \le 2\lg(n+1)$。
\item 能對一串鍵做紅黑樹的插入，每一步寫出情況和改色、旋轉；能說明刪除的四種情況，並
追蹤一次完整的刪除。
\end{itemize}

\textbf{延伸閱讀}（書末會整理完整的參考資料）

\begin{itemize}
\item Thomas H. Cormen 等，\textit{Introduction to Algorithms}，第 4 版，第 13 章：紅
黑樹的性質、旋轉、插入和刪除，用哨兵節點代替 NIL，證明寫得完整；章末問題有 AVL 樹。
\item Donald E. Knuth，\textit{The Art of Computer Programming}，第 3 冊，第 6.2.3 節：
平衡樹的原始講法，包括 AVL 樹高度的精確分析和只存平衡因子的插入程式。
\item Robert Sedgewick、Kevin Wayne，\textit{Algorithms}，第 4 版，第 3.3 節：從 2-3
樹出發講紅黑樹，用的是只允許左傾紅色連結的版本，規則和本章不同，可以和第~11~章一起
讀。
\item Steven S. Skiena，\textit{The Algorithm Design Manual}，第 3 版，第 3.4 節：平
衡搜尋樹在實務上的用法，什麼時候該用、什麼時候用雜湊（第~12~章）就夠了。
\end{itemize}

\addcontentsline{toc}{section}{觀念題解答}
\section*{觀念題解答}

\begin{enumerate}
\item (a)、(c)。
(a) 對，定理~\ref{thm:bal-avl-insert}。 (b) 錯，圖~\ref{fig:bal-avl-cascade} 刪一個鍵
做了一次單旋轉加一次雙旋轉，更大的費氏樹每一層都可能旋轉。 (c) 對：根是黑的，所以
路徑上每個紅色節點的父節點也在路徑上，而且是黑色的內部節點；不同的紅色節點父節點不
同，所以紅色不比黑色內部節點多。 (d) 錯，下面這棵是合法的紅黑樹（每條路徑到 NIL 都
有 $2$ 個黑色，不算根），但根的左子樹高 $3$、右子樹高 $1$，平衡因子 $+2$。
\begin{center}
\begin{tikzpicture}[x=0.55cm,y=0.8cm,baseline=(current bounding box.north)]
\node[tnR] (xf5) at (0,-3) {5};
\node[tnB] (xf10) at (1,-2) {10};
\node[tnR] (xf20) at (2,-1) {20};
\node[tnB] (xf30) at (3,-2) {30};
\node[tnB] (xf40) at (4,0) {40};
\node[tnB] (xf50) at (5,-1) {50};
\draw[te] (xf40) -- (xf20);
\draw[te] (xf20) -- (xf10);
\draw[te] (xf10) -- (xf5);
\draw[te] (xf20) -- (xf30);
\draw[te] (xf40) -- (xf50);
\end{tikzpicture}
\end{center}
(e) 錯。 $7$ 個節點的費氏樹（圖~\ref{fig:bal-fib} 的 $h = 4$）是 AVL 樹，高度 $4$；
同樣 $7$ 個鍵排成滿二元樹、全部塗黑，是高度 $3$ 的紅黑樹。

\item (a) 最少 $N(7) = 33$，最多 $2^7 - 1 = 127$。 (b) $N(8) = 54 \le 60 < 88 = N(9)$，
最大 $8$；最小 $\lceil \lg 61 \rceil = 6$。 (c) 數邊的高度 $7$ 等於數節點的高度
$8$：最少 $N(8) = 54$，最多 $2^8 - 1 = 255$。

\item 旋轉四次：
\begin{center}
\begin{tabular}{cccl}
\toprule
插入 & 失衡節點 & 類型 & 說明 \\
\midrule
90 & 72 & RL & $72 < 90 < 98$，$90$ 升到根 \\
80 & 90 & LR & 從 $90$ 往下先左（$72$）再右（$82$），$82$ 升到根 \\
70 & 41 & RR & $41$ 右邊是 $64$、$64$ 右邊是 $70$，對 $41$ 左旋 \\
36 & 72 & LL & $72$ 左邊是 $64$、$64$ 左邊是 $41$，對 $72$ 右旋 \\
\bottomrule
\end{tabular}
\end{center}
插入 $41$、$82$、$64$ 沒有失衡。 每次旋轉後的樹（右上方是平衡因子）：
\begin{center}
\begin{tikzpicture}[x=0.48cm,y=0.75cm,baseline=(current bounding box.north)]
\node[tn] (xaa72) at (0,-1) {72};
\node[tn] (xaa90) at (1,0) {90};
\node[tn] (xaa98) at (2,-1) {98};
\draw[te] (xaa90) -- (xaa72);
\draw[te] (xaa90) -- (xaa98);
\node[bfl] at (xaa72.north east) {$0$};
\node[bfl] at (xaa90.north east) {$0$};
\node[bfl] at (xaa98.north east) {$0$};
\node[plab] at (1.0,-3.75) {(a) 插入 90 後（RL）};
\end{tikzpicture}\hspace{0.5cm}\begin{tikzpicture}[x=0.48cm,y=0.75cm,baseline=(current bounding box.north)]
\node[tn] (xab41) at (0,-2) {41};
\node[tn] (xab72) at (1,-1) {72};
\node[tn] (xab80) at (2,-2) {80};
\node[tn] (xab82) at (3,0) {82};
\node[tn] (xab90) at (4,-1) {90};
\node[tn] (xab98) at (5,-2) {98};
\draw[te] (xab82) -- (xab72);
\draw[te] (xab72) -- (xab41);
\draw[te] (xab72) -- (xab80);
\draw[te] (xab82) -- (xab90);
\draw[te] (xab90) -- (xab98);
\node[bfl] at (xab41.north east) {$0$};
\node[bfl] at (xab72.north east) {$0$};
\node[bfl] at (xab80.north east) {$0$};
\node[bfl] at (xab82.north east) {$0$};
\node[bfl] at (xab90.north east) {$-1$};
\node[bfl] at (xab98.north east) {$0$};
\node[plab] at (2.5,-3.75) {(b) 插入 80 後（LR）};
\end{tikzpicture}\hspace{0.5cm}\begin{tikzpicture}[x=0.48cm,y=0.75cm,baseline=(current bounding box.north)]
\node[tn] (xac41) at (0,-3) {41};
\node[tn] (xac64) at (1,-2) {64};
\node[tn] (xac70) at (2,-3) {70};
\node[tn] (xac72) at (3,-1) {72};
\node[tn] (xac80) at (4,-2) {80};
\node[tn] (xac82) at (5,0) {82};
\node[tn] (xac90) at (6,-1) {90};
\node[tn] (xac98) at (7,-2) {98};
\draw[te] (xac82) -- (xac72);
\draw[te] (xac72) -- (xac64);
\draw[te] (xac64) -- (xac41);
\draw[te] (xac64) -- (xac70);
\draw[te] (xac72) -- (xac80);
\draw[te] (xac82) -- (xac90);
\draw[te] (xac90) -- (xac98);
\node[bfl] at (xac41.north east) {$0$};
\node[bfl] at (xac64.north east) {$0$};
\node[bfl] at (xac70.north east) {$0$};
\node[bfl] at (xac72.north east) {$+1$};
\node[bfl] at (xac80.north east) {$0$};
\node[bfl] at (xac82.north east) {$+1$};
\node[bfl] at (xac90.north east) {$-1$};
\node[bfl] at (xac98.north east) {$0$};
\node[plab] at (3.5,-3.75) {(c) 插入 70 後（RR）};
\end{tikzpicture}\\*[0.6em]
\begin{tikzpicture}[x=0.48cm,y=0.75cm,baseline=(current bounding box.north)]
\node[tn] (xad36) at (0,-3) {36};
\node[tn] (xad41) at (1,-2) {41};
\node[tn] (xad64) at (2,-1) {64};
\node[tn] (xad70) at (3,-3) {70};
\node[tn] (xad72) at (4,-2) {72};
\node[tn] (xad80) at (5,-3) {80};
\node[tn] (xad82) at (6,0) {82};
\node[tn] (xad90) at (7,-1) {90};
\node[tn] (xad98) at (8,-2) {98};
\draw[te] (xad82) -- (xad64);
\draw[te] (xad64) -- (xad41);
\draw[te] (xad41) -- (xad36);
\draw[te] (xad64) -- (xad72);
\draw[te] (xad72) -- (xad70);
\draw[te] (xad72) -- (xad80);
\draw[te] (xad82) -- (xad90);
\draw[te] (xad90) -- (xad98);
\node[bfl] at (xad36.north east) {$0$};
\node[bfl] at (xad41.north east) {$+1$};
\node[bfl] at (xad64.north east) {$0$};
\node[bfl] at (xad70.north east) {$0$};
\node[bfl] at (xad72.north east) {$0$};
\node[bfl] at (xad80.north east) {$0$};
\node[bfl] at (xad82.north east) {$+1$};
\node[bfl] at (xad90.north east) {$-1$};
\node[bfl] at (xad98.north east) {$0$};
\node[plab] at (4.0,-3.75) {(d) 插入 36 後（LL）};
\end{tikzpicture}
\end{center}
插入 $36$ 時，從 $36$ 往上，$41$ 變 $+1$、$64$ 變 $+1$、$72$ 變 $+2$，根 $82$ 也變成
$+2$（左高 $4$、右高 $2$）。 要修的是最近的 $72$，不是根：修好 $72$ 之後，這一段回到插入前
的高度 $3$，根也回到 $+1$，不用再轉。 如果從根開始修，會轉錯地方。

\item
\begin{center}
\begin{tikzpicture}[x=0.46cm,y=0.75cm,baseline=(current bounding box.north)]
\node[tn] (xba12) at (0,-3) {12};
\node[tn] (xba18) at (1,-2) {18};
\node[tn] (xba26) at (2,-1) {26};
\node[tn] (xba31) at (3,-3) {31};
\node[tn] (xba35) at (4,-2) {35};
\node[tn] (xba40) at (5,-3) {40};
\node[tn] (xba62) at (6,0) {62};
\node[tn] (xba71) at (7,-2) {71};
\node[tn] (xba84) at (8,-1) {84};
\node[tn] (xba93) at (9,-2) {93};
\draw[te] (xba62) -- (xba26);
\draw[te] (xba26) -- (xba18);
\draw[te] (xba18) -- (xba12);
\draw[te] (xba26) -- (xba35);
\draw[te] (xba35) -- (xba31);
\draw[te] (xba35) -- (xba40);
\draw[te] (xba62) -- (xba84);
\draw[te] (xba84) -- (xba71);
\draw[te] (xba84) -- (xba93);
\node[bfl] at (xba12.north east) {$0$};
\node[bfl] at (xba18.north east) {$+1$};
\node[bfl] at (xba26.north east) {$0$};
\node[bfl] at (xba31.north east) {$0$};
\node[bfl] at (xba35.north east) {$0$};
\node[bfl] at (xba40.north east) {$0$};
\node[bfl] at (xba62.north east) {$+1$};
\node[bfl] at (xba71.north east) {$0$};
\node[bfl] at (xba84.north east) {$0$};
\node[bfl] at (xba93.north east) {$0$};
\node[plab] at (4.5,-3.75) {(a) 刪 47};
\end{tikzpicture}\hspace{0.4cm}\begin{tikzpicture}[x=0.46cm,y=0.75cm,baseline=(current bounding box.north)]
\node[tn] (xbb12) at (0,-3) {12};
\node[tn] (xbb18) at (1,-2) {18};
\node[tn] (xbb26) at (2,-1) {26};
\node[tn] (xbb31) at (3,-3) {31};
\node[tn] (xbb35) at (4,-2) {35};
\node[tn] (xbb40) at (5,-3) {40};
\node[tn] (xbb62) at (6,0) {62};
\node[tn] (xbb84) at (7,-1) {84};
\node[tn] (xbb93) at (8,-2) {93};
\draw[te] (xbb62) -- (xbb26);
\draw[te] (xbb26) -- (xbb18);
\draw[te] (xbb18) -- (xbb12);
\draw[te] (xbb26) -- (xbb35);
\draw[te] (xbb35) -- (xbb31);
\draw[te] (xbb35) -- (xbb40);
\draw[te] (xbb62) -- (xbb84);
\draw[te] (xbb84) -- (xbb93);
\node[bfl] at (xbb12.north east) {$0$};
\node[bfl] at (xbb18.north east) {$+1$};
\node[bfl] at (xbb26.north east) {$0$};
\node[bfl] at (xbb31.north east) {$0$};
\node[bfl] at (xbb35.north east) {$0$};
\node[bfl] at (xbb40.north east) {$0$};
\node[bfl] at (xbb62.north east) {$+1$};
\node[bfl] at (xbb84.north east) {$-1$};
\node[bfl] at (xbb93.north east) {$0$};
\node[plab] at (4.0,-3.75) {(b) 刪 71};
\end{tikzpicture}\hspace{0.4cm}\begin{tikzpicture}[x=0.46cm,y=0.75cm,baseline=(current bounding box.north)]
\node[tn] (xbc12) at (0,-2) {12};
\node[tn] (xbc18) at (1,-1) {18};
\node[tn] (xbc26) at (2,0) {26};
\node[tn] (xbc31) at (3,-3) {31};
\node[tn] (xbc35) at (4,-2) {35};
\node[tn] (xbc40) at (5,-3) {40};
\node[tn] (xbc62) at (6,-1) {62};
\node[tn] (xbc84) at (7,-2) {84};
\draw[te] (xbc26) -- (xbc18);
\draw[te] (xbc18) -- (xbc12);
\draw[te] (xbc26) -- (xbc62);
\draw[te] (xbc62) -- (xbc35);
\draw[te] (xbc35) -- (xbc31);
\draw[te] (xbc35) -- (xbc40);
\draw[te] (xbc62) -- (xbc84);
\node[bfl] at (xbc12.north east) {$0$};
\node[bfl] at (xbc18.north east) {$+1$};
\node[bfl] at (xbc26.north east) {$-1$};
\node[bfl] at (xbc31.north east) {$0$};
\node[bfl] at (xbc35.north east) {$0$};
\node[bfl] at (xbc40.north east) {$0$};
\node[bfl] at (xbc62.north east) {$+1$};
\node[bfl] at (xbc84.north east) {$0$};
\node[plab] at (3.5,-3.75) {(c) 刪 93};
\end{tikzpicture}\hspace{0.4cm}\begin{tikzpicture}[x=0.46cm,y=0.75cm,baseline=(current bounding box.north)]
\node[tn] (xbd18) at (0,-2) {18};
\node[tn] (xbd26) at (1,-1) {26};
\node[tn] (xbd31) at (2,-2) {31};
\node[tn] (xbd35) at (3,0) {35};
\node[tn] (xbd40) at (4,-2) {40};
\node[tn] (xbd62) at (5,-1) {62};
\node[tn] (xbd84) at (6,-2) {84};
\draw[te] (xbd35) -- (xbd26);
\draw[te] (xbd26) -- (xbd18);
\draw[te] (xbd26) -- (xbd31);
\draw[te] (xbd35) -- (xbd62);
\draw[te] (xbd62) -- (xbd40);
\draw[te] (xbd62) -- (xbd84);
\node[bfl] at (xbd18.north east) {$0$};
\node[bfl] at (xbd26.north east) {$0$};
\node[bfl] at (xbd31.north east) {$0$};
\node[bfl] at (xbd35.north east) {$0$};
\node[bfl] at (xbd40.north east) {$0$};
\node[bfl] at (xbd62.north east) {$0$};
\node[bfl] at (xbd84.north east) {$0$};
\node[plab] at (3.0,-3.75) {(d) 刪 12};
\end{tikzpicture}
\end{center}
\begin{itemize}
\item 刪 $47$：後繼是 $62$，$62$ 取代根，再刪掉原本的 $62$。 $71$ 左邊空了、右邊 $84$
高 $2$，平衡因子 $-2$；$84$ 是 $-1$，刪除發生在左邊，是 L1，對 $71$ 左旋。 根 $62$ 左高
$3$、右高 $2$，平衡。
\item 刪 $71$：$84$ 變 $-1$，高度不變，不用旋轉。
\item 刪 $93$：$84$ 變成葉節點，根 $62$ 左高 $3$、右高 $1$，$+2$；左子節點 $26$ 是 $0$，
是 R0，對 $62$ 右旋。 高度還是 $4$。
\item 刪 $12$：根 $26$ 左高 $1$、右高 $3$，$-2$；右子節點 $62$ 是 $+1$，是 L$-1$，做 RL，
$35$ 升到根，$31$ 接到 $26$ 的右邊，$40$ 接到 $62$ 的左邊。
\end{itemize}

\item (a) 設 $T_1$、$T_4$ 高 $k$。
\begin{center}
\begin{tabular}{lccc}
\toprule
情形 & $\mathrm{BF}(B)$ & $\mathrm{BF}(A)$ & $\mathrm{BF}(C)$ \\
\midrule
新節點在 $T_2$（$T_2$ 高 $k$、$T_3$ 高 $k - 1$） & $0$ & $-1$ & $0$ \\
新節點在 $T_3$（$T_2$ 高 $k - 1$、$T_3$ 高 $k$） & $+1$ & $0$ & $0$ \\
$C$ 是新節點（$k = 0$） & $0$ & $0$ & $0$ \\
\bottomrule
\end{tabular}
\end{center}
旋轉後 $B$ 的子樹是 $T_1$、$T_2$，$A$ 的子樹是 $T_3$、$T_4$，直接相減即得。 $B$、$A$
都高 $k + 1$，所以 $C$ 是 $0$。

(b) 旋轉前 $\mathrm{BF}(C) = +1$ 表示新節點在 $T_2$，$-1$ 表示在 $T_3$，$0$ 表示 $C$
是新節點。 所以先記下舊的 $\mathrm{BF}(C)$，旋轉後：$+1 \Rightarrow (B, A) = (0, -1)$；
$-1 \Rightarrow (B, A) = (+1, 0)$；$0 \Rightarrow (0, 0)$；$C$ 一律設成 $0$。 RL 的規則
左右對調。

\item $v$ 的平衡因子從插入前的值變成 $0$，插入前只可能是 $\pm 1$（插入只讓一邊高度最多
加 $1$），而且新節點加在原本比較矮的那一邊，讓兩邊一樣高。 $v$ 的高度是
$1 + \max(h_\ell, h_r)$，比較高的那一邊沒有變，所以 $v$ 的高度不變。 $v$ 的祖先只透過
子樹的高度受影響，$v$ 這棵子樹的高度沒變，祖先的高度和平衡因子都不變，可以停止。

如果更新後變成 $\pm 1$，插入前一定是 $0$，新節點讓其中一邊加高 $1$，$v$ 的高度也加 $1$，
父節點的平衡因子會改變，所以要繼續往上。 （變成 $\pm 2$ 時旋轉，旋轉後高度回到插入前，
也可以停。）

\item 例如下面左邊的樹刪除 $33$：$26$ 左高 $2$、右高 $0$，左子節點 $14$ 是 $0$，是 R0，
對 $26$ 右旋，得到右邊的樹，高度仍然是 $3$。
\begin{center}
\begin{tikzpicture}[x=0.55cm,y=0.85cm,baseline=(current bounding box.north)]
\node[tn] (xg19) at (0,-2) {9};
\node[tn] (xg114) at (1,-1) {14};
\node[tn] (xg119) at (2,-2) {19};
\node[tn] (xg126) at (3,0) {26};
\node[tn] (xg133) at (4,-1) {33};
\draw[te] (xg126) -- (xg114);
\draw[te] (xg114) -- (xg19);
\draw[te] (xg114) -- (xg119);
\draw[te] (xg126) -- (xg133);
\node[bfl] at (xg19.north east) {$0$};
\node[bfl] at (xg114.north east) {$0$};
\node[bfl] at (xg119.north east) {$0$};
\node[bfl] at (xg126.north east) {$+1$};
\node[bfl] at (xg133.north east) {$0$};
\node[plab] at (2.0,-2.75) {(a) 刪除前};
\end{tikzpicture}\hspace{1.2cm}\begin{tikzpicture}[x=0.55cm,y=0.85cm,baseline=(current bounding box.north)]
\node[tn] (xg29) at (0,-1) {9};
\node[tn] (xg214) at (1,0) {14};
\node[tn] (xg219) at (2,-2) {19};
\node[tn] (xg226) at (3,-1) {26};
\draw[te] (xg214) -- (xg29);
\draw[te] (xg214) -- (xg226);
\draw[te] (xg226) -- (xg219);
\node[bfl] at (xg29.north east) {$0$};
\node[bfl] at (xg214.north east) {$-1$};
\node[bfl] at (xg219.north east) {$0$};
\node[bfl] at (xg226.north east) {$+1$};
\node[plab] at (1.5,-2.75) {(b) 刪 33、右旋 26 之後};
\end{tikzpicture}
\end{center}
節點不能更少：$B$ 是 $A$ 比較高的子節點，刪除後 $A$ 是 $\pm 2$，所以 $B$ 至少高 $2$；
$\mathrm{BF}(B) = 0$ 表示 $B$ 的兩棵子樹一樣高，都至少高 $1$，$B$ 這邊至少 $3$ 個節點。
$A$ 的另一邊是因為刪除才變矮的，刪除前至少有 $1$ 個節點。 加上 $A$ 自己，至少 $5$ 個節點。

\item (a) 最少 $2^k - 1$，最多 $4^k - 1$。 下界就是引理~\ref{lem:bal-rb-count}，高度 $k$
的滿二元樹全部塗黑可以達到。 上界：根到任一個 NIL 的路徑上，黑色內部節點有 $k$ 個（根
和不算 NIL 的 $k - 1$ 個），由第 1 題 (c)，紅色最多也是 $k$ 個，所以高度不超過 $2k$，內部
節點不超過 $2^{2k} - 1 = 4^k - 1$。 高度 $2k$ 的滿二元樹，第 $1, 3, \ldots, 2k-1$ 層塗黑、
第 $2, 4, \ldots, 2k$ 層塗紅，可以達到。 (b) $2k = 4$。

\item 設從 $x$ 往下到 NIL 的某條路徑有 $L$ 條邊，也就是經過 $x$ 以下的 $L$ 個節點（含最後
的 NIL）。 由第 5 條，這 $L$ 個節點裡黑色的恰好有 $\mathrm{bh}(x)$ 個，所以任何路徑
$L \ge \mathrm{bh}(x)$，最短路徑也是。 另一方面，這 $L$ 個節點裡每個紅色節點的子節點也在
路徑上，而且是黑色；不同的紅色節點子節點不同，所以紅色最多 $\mathrm{bh}(x)$ 個，
$L \le 2\,\mathrm{bh}(x)$。 因此最長 $\le 2\,\mathrm{bh}(x) \le 2 \times$ 最短。

\item (A) 和 (D) 合法。 (B) 違反第 5 條：$30 \to 15 \to \text{NIL}$ 有 $2$ 個黑色
（$15$、NIL），$30 \to 40 \to \text{NIL}$ 只有 $1$ 個。 (C) 違反第 4 條：$15$ 和 $10$ 都
是紅的而且相連（黑色個數倒是都一樣）。 (D) 的根到 NIL 每條路徑都是 $2$ 個黑色，紅色的
$40$ 兩個子節點都是黑的。

\item (a)
\begin{center}
\begin{tabular}{cp{0.78\linewidth}}
\toprule
插入 & 情況 \\
\midrule
62 & 根改黑 \\
90、31 & 不用修 \\
13 & 1L（叔叔 $90$ 紅），$z$ 到根 $62$，根改黑 \\
25 & 2L（$25$ 在 $13$ 的內側），左旋 $13$；3L，右旋 $31$ \\
45 & 1R（父 $31$ 是 $25$ 的右子節點，叔叔 $13$ 紅），$z = 25$，父 $62$ 黑，結束 \\
43 & 2R（$43$ 在 $45$ 的內側），右旋 $45$；3R，左旋 $31$ \\
49 & 1R（叔叔 $31$ 紅），$z = 43$；父 $25$ 紅、叔叔 $90$ 黑、$43$ 在內側：2L，左旋 $25$；
3L，右旋 $62$ \\
\bottomrule
\end{tabular}
\end{center}
\begin{center}
\begin{tikzpicture}[x=0.48cm,y=0.75cm,baseline=(current bounding box.north)]
\node[tnR] (xca13) at (0,-2) {13};
\node[tnB] (xca31) at (1,-1) {31};
\node[tnB] (xca62) at (2,0) {62};
\node[tnB] (xca90) at (3,-1) {90};
\draw[te] (xca62) -- (xca31);
\draw[te] (xca31) -- (xca13);
\draw[te] (xca62) -- (xca90);
\node[plab] at (1.5,-3.75) {(a) 插入 13 之後};
\end{tikzpicture}\hspace{0.5cm}\begin{tikzpicture}[x=0.48cm,y=0.75cm,baseline=(current bounding box.north)]
\node[tnR] (xcb13) at (0,-2) {13};
\node[tnB] (xcb25) at (1,-1) {25};
\node[tnR] (xcb31) at (2,-2) {31};
\node[tnB] (xcb62) at (3,0) {62};
\node[tnB] (xcb90) at (4,-1) {90};
\draw[te] (xcb62) -- (xcb25);
\draw[te] (xcb25) -- (xcb13);
\draw[te] (xcb25) -- (xcb31);
\draw[te] (xcb62) -- (xcb90);
\node[plab] at (2.0,-3.75) {(b) 插入 25 之後};
\end{tikzpicture}\\*[0.6em]
\begin{tikzpicture}[x=0.48cm,y=0.75cm,baseline=(current bounding box.north)]
\node[tnB] (xcc13) at (0,-2) {13};
\node[tnR] (xcc25) at (1,-1) {25};
\node[tnR] (xcc31) at (2,-3) {31};
\node[tnB] (xcc43) at (3,-2) {43};
\node[tnR] (xcc45) at (4,-3) {45};
\node[tnB] (xcc62) at (5,0) {62};
\node[tnB] (xcc90) at (6,-1) {90};
\draw[te] (xcc62) -- (xcc25);
\draw[te] (xcc25) -- (xcc13);
\draw[te] (xcc25) -- (xcc43);
\draw[te] (xcc43) -- (xcc31);
\draw[te] (xcc43) -- (xcc45);
\draw[te] (xcc62) -- (xcc90);
\node[plab] at (3.0,-3.75) {(c) 插入 43 之後};
\end{tikzpicture}\hspace{0.5cm}\begin{tikzpicture}[x=0.48cm,y=0.75cm,baseline=(current bounding box.north)]
\node[tnB] (xcd13) at (0,-2) {13};
\node[tnR] (xcd25) at (1,-1) {25};
\node[tnB] (xcd31) at (2,-2) {31};
\node[tnB] (xcd43) at (3,0) {43};
\node[tnB] (xcd45) at (4,-2) {45};
\node[tnR] (xcd49) at (5,-3) {49};
\node[tnR] (xcd62) at (6,-1) {62};
\node[tnB] (xcd90) at (7,-2) {90};
\draw[te] (xcd43) -- (xcd25);
\draw[te] (xcd25) -- (xcd13);
\draw[te] (xcd25) -- (xcd31);
\draw[te] (xcd43) -- (xcd62);
\draw[te] (xcd62) -- (xcd45);
\draw[te] (xcd45) -- (xcd49);
\draw[te] (xcd62) -- (xcd90);
\node[plab] at (3.5,-3.75) {(d) 插入 49 之後};
\end{tikzpicture}
\end{center}

(b) 刪 $31$：黑色葉節點，$x$ 是 NIL，$p = 25$（紅），兄弟 $s = 13$ 黑、兩個子節點都是 NIL：
情況 2（左右對調的版本），$13$ 改紅，$x$ 移到 $25$。 $25$ 是紅的，迴圈結束，改黑。

刪 $62$：有兩個子節點，後繼是 $90$，$90$ 的鍵複製到 $62$ 的位置，改成拿掉原本的 $90$（黑色
葉節點）。 $x$ 是 NIL，是 $p$（現在存 $90$）的右子節點，兄弟 $s = 45$ 黑；用對調的版本，近的
子節點是 $45$ 的右子節點 $49$（紅），遠的是左邊的 NIL（黑）：情況 3，$49$ 改黑、$45$ 改紅，
左旋 $45$。 新的兄弟 $49$，遠的子節點 $45$ 紅：情況 4，$49$ 取 $p$ 的顏色（紅），$p$ 改黑，
$45$ 改黑，右旋 $p$。
\begin{center}
\begin{tikzpicture}[x=0.5cm,y=0.7cm,baseline=(current bounding box.north)]
\node[tnR] (xda13) at (0,-2) {13};
\node[tnB] (xda25) at (1,-1) {25};
\node[tnB] (xda43) at (2,0) {43};
\node[tnB] (xda45) at (3,-2) {45};
\node[tnR] (xda49) at (4,-3) {49};
\node[tnR] (xda62) at (5,-1) {62};
\node[tnB] (xda90) at (6,-2) {90};
\draw[te] (xda43) -- (xda25);
\draw[te] (xda25) -- (xda13);
\draw[te] (xda43) -- (xda62);
\draw[te] (xda62) -- (xda45);
\draw[te] (xda45) -- (xda49);
\draw[te] (xda62) -- (xda90);
\node[plab] at (3.0,-3.75) {(a) 刪 31 之後};
\end{tikzpicture}\hspace{1.0cm}\begin{tikzpicture}[x=0.5cm,y=0.7cm,baseline=(current bounding box.north)]
\node[tnR] (xdb13) at (0,-2) {13};
\node[tnB] (xdb25) at (1,-1) {25};
\node[tnB] (xdb43) at (2,0) {43};
\node[tnB] (xdb45) at (3,-2) {45};
\node[tnR] (xdb49) at (4,-1) {49};
\node[tnB] (xdb90) at (5,-2) {90};
\draw[te] (xdb43) -- (xdb25);
\draw[te] (xdb25) -- (xdb13);
\draw[te] (xdb43) -- (xdb49);
\draw[te] (xdb49) -- (xdb45);
\draw[te] (xdb49) -- (xdb90);
\node[plab] at (2.5,-3.75) {(b) 再刪 62 之後};
\end{tikzpicture}
\end{center}

\item (a) 第 5 條。 新的黑色節點讓經過它的路徑多一個黑色，其他路徑沒有，要修就得在別的
路徑上各補一個黑色，或在這條路徑上去掉一個，牽動的範圍可能是整棵樹。 塗成紅色時，黑色
個數完全不變，只可能和父節點紅紅相連，問題只在 $z$ 和 $p$ 兩個節點之間，可以在附近修。

(b) 改色之後 $p$ 是黑的，$z$ 和 $p$ 不再是兩個相連的紅色；$u$ 也是黑的。 唯一可能出新問題的
是 $g$：它從黑變紅，如果 $g$ 的父節點也是紅的，就是新的紅紅相連。 所以下一輪的 $z$ 是
$g$。 如果檢查 $p$，$p$ 是黑的，迴圈會直接結束，漏掉 $g$ 的問題。

\item (a) 要找滿足 $N(h) = F_{h+2} - 1 \le 10^6$ 的最大 $h$。 $F_{30} = 832040$、
$F_{31} = 1346269$，$N(28) = 832039 \le 10^6 < 1346268 = N(29)$，最大 $28$。 （定理的
估計 $1.4404 \times \lg(10^6 + 1) \approx 28.7$，一致。） (b) $2\lg(10^6 + 1) \approx 39.9$，
不超過 $39$。 (c) $\lceil \lg(10^6 + 1) \rceil = 20$。 依序插入 $10^6$ 個鍵，實際高度是 $20$
和 $37$（第~\ref{sec:bal-rb-c}~節）。

\end{enumerate}

\chapter{多路搜尋樹：2-3 樹、2-3-4 樹、B 樹與 \texorpdfstring{B$^{+}$}{B+} 樹}
\label{ch:multiway}

% 本章的圖共用的樣式。樹的圖由模擬程式依照本章的演算法產生，每一步都和演算法一致。
\definecolor{mwRed}{HTML}{C8102E}
\tikzset{
    mwk/.style={rectangle, draw=black!65, minimum width=0.5cm, minimum height=0.42cm,
        inner sep=0pt, outer sep=0pt, font=\small},
    mwkN/.style={fill=coreFill},                          % 剛插入的鍵
    mwkM/.style={fill=greenFill},                         % 剛移動的鍵（上升、下降、借來）
    mwkB/.style={fill=memFill},                           % 鍵太多或太少的節點
    mwbad/.style={draw=memStroke, very thick},
    mwempty/.style={draw=memStroke, very thick, dashed, fill=memFill},
    mwe/.style={draw=black!60},
    mwlink/.style={-{Stealth[length=1.6mm]}, draw=black!60},
    mwlab/.style={font=\footnotesize},
    mwB/.style={circle, draw=black, fill=black!85, text=white, minimum size=0.52cm,
        inner sep=0pt, font=\small},
    mwR/.style={circle, draw=mwRed, fill=mwRed, text=white, minimum size=0.52cm,
        inner sep=0pt, font=\small},
    mwgrp/.style={rounded corners=2.5mm, fill=black!8, draw=black!25, inner sep=1.5pt},
    mwsub/.style={font=\scriptsize, text=black!70},
}

學校圖書館的館藏目錄有幾百萬筆，存在硬碟上。 從硬碟讀一次資料，比在記憶體裡比較
兩個數慢上萬倍，而且每次讀進來的是一整塊，例如 4096 個位元組，不是一筆。 第~10~章
的紅黑樹每個節點只放一個鍵，查一本書要從根往下走二十幾層，每走一層就讀一次硬碟。
如果每次讀進來的那一整塊可以放幾百個鍵，樹應該長成什麼樣子？ 查一本書要讀幾次硬
碟？ 往下讀之前，先想一想。

這一章的樹，每個節點放好幾個鍵、有好幾個子節點，而且所有的葉節點（下面常簡稱葉子）
都在同一層。 節點太滿就分裂，太空就向兄弟借一個鍵或和兄弟合併，樹只會從根長高、從
根變矮，所以葉子永遠同一層，高度永遠是 $O(\log n)$。 第~\ref{sec:mw-23}~節
先用最小的兩種，2-3 樹和 2-3-4 樹，把分裂、借、合併三個動作練熟，並說明 2-3-4 樹和
第~10~章紅黑樹的對應；第~\ref{sec:mw-btree}~節推廣成任意階數的 B 樹，證明高度的
上下界，寫出搜尋與插入的 C 程式，插入和刪除各做一次完整的追蹤；
第~\ref{sec:mw-bplus}~節講資料庫和檔案系統實際用的 B$^{+}$ 樹，以及怎麼估計讀幾次
磁碟。

這一章用到《C 程式設計基礎》第~14~章的結構和結構指標、第~19~章的
\texttt{malloc}、\texttt{free}。 樹的名詞見第~7~章，二元搜尋樹見第~9~章，紅黑樹見
第~10~章。 高度照全書的定義數節點：只有根的樹高度是 $1$，根在第 $1$ 層。

\section{2-3 樹與 2-3-4 樹}
\label{sec:mw-23}

\subsection{多路搜尋樹}

二元搜尋樹的節點放一個鍵，把鍵的範圍切成兩段，各交給一棵子樹。 把它推廣：一個節點
放 $k$ 個由小到大的鍵 $a_1 < a_2 < \cdots < a_k$，就把範圍切成 $k + 1$ 段，接
$k + 1$ 棵子樹 $T_0, T_1, \ldots, T_k$。 $T_0$ 的鍵都小於 $a_1$，$T_i$ 的鍵都在
$a_i$ 和 $a_{i+1}$ 之間，$T_k$ 的鍵都大於 $a_k$（圖~\ref{fig:mw-node}）。 每個節點
最多 $m$ 棵子樹的這種樹叫做 \textbf{$m$ 路搜尋樹}（$m$-way search tree）。 有 $d$ 棵
子樹（$d - 1$ 個鍵）的節點叫 \textbf{$d$-節點}，例如 2-節點就是二元搜尋樹的節點，
3-節點放兩個鍵、接三棵子樹。

\begin{center}
\begin{tikzpicture}[mwk/.append style={minimum width=0.62cm}]
\node[mwk] at (0.31,0) {$a_1$};
\node[mwk] at (0.93,0) {$a_2$};
\draw[mwe] (0,-0.21) -- (-1.4,-1.0);
\draw[mwe] (0.62,-0.21) -- (0.62,-1.0);
\draw[mwe] (1.24,-0.21) -- (2.6,-1.0);
\draw[black!60, fill=black!4] (-1.4,-1.0) -- (-2.0,-2.0) -- (-0.8,-2.0) -- cycle;
\draw[black!60, fill=black!4] (0.62,-1.0) -- (0.02,-2.0) -- (1.22,-2.0) -- cycle;
\draw[black!60, fill=black!4] (2.6,-1.0) -- (2.0,-2.0) -- (3.2,-2.0) -- cycle;
\node[mwsub] at (-1.4,-1.65) {$T_0$};
\node[mwsub] at (0.62,-1.65) {$T_1$};
\node[mwsub] at (2.6,-1.65) {$T_2$};
\node[mwsub] at (-1.4,-2.35) {$< a_1$};
\node[mwsub] at (0.62,-2.35) {$a_1 \sim a_2$ 之間};
\node[mwsub] at (2.6,-2.35) {$> a_2$};
\end{tikzpicture}
\diagcaption{一個 3-節點：兩個鍵把範圍切成三段，三棵子樹各管一段。 圖上的子指標從兩
個鍵之間的縫往下畫，本章的樹都這樣畫。}
\label{fig:mw-node}
\end{center}

搜尋 $k$ 時，在目前的節點裡從左到右找第一個 $\ge k$ 的鍵 $a_i$。 $a_i = k$ 就找到
了；否則 $k$ 落在 $a_{i-1}$ 和 $a_i$ 之間，往 $T_{i-1}$ 走（$k$ 比所有鍵都大時往
$T_k$ 走）。 走到空的子樹就是找不到。 每個節點要比較最多 $m - 1$ 次，但往下走的層數
少了。

只規定「每個節點最多 $m$ 棵子樹」不夠，多路搜尋樹一樣可能長成一條鏈。 這一節的兩種
樹再加上兩條規定：每個節點要夠滿，而且所有葉節點在同一層。

\subsection{2-3 樹}

\begin{definition}[2-3 樹]
\label{def:mw-23}
一棵多路搜尋樹，滿足：
\begin{enumerate}
\item 每個內部節點是 2-節點（$1$ 個鍵、$2$ 個子節點）或 3-節點（$2$ 個鍵、$3$ 個子
節點）。
\item 每個葉節點放 $1$ 或 $2$ 個鍵。
\item 所有的葉節點在同一層。
\end{enumerate}
\end{definition}

高度 $h$ 的 2-3 樹，鍵最少時每個節點都是 2-節點，有 $2^h - 1$ 個鍵；最多時每個節點都
是 3-節點，有 $2(1 + 3 + \cdots + 3^{h-1}) = 3^h - 1$ 個鍵。 所以 $n$ 個鍵的 2-3 樹
$\log_3(n + 1) \le h \le \lg(n + 1)$。

\subsection{插入：分裂}
\label{sec:mw-23-insert}

插入 $k$ 時，先照搜尋的路線走到一片葉子，把 $k$ 依大小放進去。 葉子原本只有一個鍵，
放進去變成兩個，結束。 葉子原本有兩個鍵，放進去就有三個，超過上限，叫做\textbf{溢位}
（overflow）。 這時把它\textbf{分裂}（split）：
\begin{itemize}
\item 最小的鍵留在原來的節點，最大的鍵放進一個新節點；
\item 中間的鍵往上放進父節點，夾在這兩個節點中間，當它們的分隔鍵。
\end{itemize}
父節點多了一個鍵和一個子節點，也可能溢位，就用同樣的方法往上分裂。 根溢位時，中間
的鍵自己成為新的根，樹高加 $1$。 樹只在根長高，每一片葉子到根的距離同時加 $1$，所以
葉子永遠在同一層。

圖~\ref{fig:mw-23-ins} 把 $47, 82, 15, 63, 29, 91, 54, 38, 70, 21$ 依序插入空的 2-3
樹。 藍底是剛插入的鍵，綠底是剛移動（上升、下降、借來）的鍵，橘框是鍵太多或太少的
節點；本章其他的圖也用同樣的顏色。 沒有分裂的插入合併成一張。

\begin{figure}[tbp]
\centering
\begin{tikzpicture}[baseline=(current bounding box.north), mwk/.append style={minimum width=0.5cm, minimum height=0.42cm}]
\node[mwk,mwkN] at (0.25,0) {47};
\node[mwk,mwkN] at (0.75,0) {82};
\node[mwlab] at (0.5,-1.28) {(a) 插入 47、82};
\end{tikzpicture}\hspace{0.55cm}\begin{tikzpicture}[baseline=(current bounding box.north), mwk/.append style={minimum width=0.5cm, minimum height=0.42cm}]
\node[mwk,mwkN] at (0.25,0) {15};
\node[mwk,mwkB] at (0.75,0) {47};
\node[mwk,mwkB] at (1.25,0) {82};
\draw[mwbad] (0,-0.21) rectangle (1.5,0.21);
\node[mwlab] at (0.75,-1.28) {(b) 插入 15：溢位};
\end{tikzpicture}\hspace{0.55cm}\begin{tikzpicture}[baseline=(current bounding box.north), mwk/.append style={minimum width=0.5cm, minimum height=0.42cm}]
\draw[mwe] (0.375,-0.21) -- (0.25,-0.59);
\draw[mwe] (0.875,-0.21) -- (1,-0.59);
\node[mwk] at (0.25,-0.8) {15};
\node[mwk] at (1,-0.8) {82};
\node[mwk,mwkM] at (0.625,0) {47};
\node[mwlab] at (0.625,-1.28) {(c) 分裂，47 上升};
\end{tikzpicture}\hspace{0.55cm}\begin{tikzpicture}[baseline=(current bounding box.north), mwk/.append style={minimum width=0.5cm, minimum height=0.42cm}]
\draw[mwe] (0.875,-0.21) -- (0.5,-0.59);
\draw[mwe] (1.375,-0.21) -- (1.75,-0.59);
\node[mwk] at (0.25,-0.8) {15};
\node[mwk,mwkN] at (0.75,-0.8) {29};
\node[mwk,mwkN] at (1.5,-0.8) {63};
\node[mwk] at (2,-0.8) {82};
\node[mwk] at (1.125,0) {47};
\node[mwlab] at (1.125,-1.28) {(d) 插入 63、29};
\end{tikzpicture}\\*[0.45em]
\begin{tikzpicture}[baseline=(current bounding box.north), mwk/.append style={minimum width=0.5cm, minimum height=0.42cm}]
\draw[mwe] (1,-0.21) -- (0.5,-0.59);
\draw[mwe] (1.5,-0.21) -- (2,-0.59);
\node[mwk] at (0.25,-0.8) {15};
\node[mwk] at (0.75,-0.8) {29};
\node[mwk,mwkB] at (1.5,-0.8) {63};
\node[mwk,mwkB] at (2,-0.8) {82};
\node[mwk,mwkN] at (2.5,-0.8) {91};
\draw[mwbad] (1.25,-1.01) rectangle (2.75,-0.59);
\node[mwk] at (1.25,0) {47};
\node[mwlab] at (1.375,-1.28) {(e) 插入 91：溢位};
\end{tikzpicture}\hspace{0.55cm}\begin{tikzpicture}[baseline=(current bounding box.north), mwk/.append style={minimum width=0.5cm, minimum height=0.42cm}]
\draw[mwe] (0.875,-0.21) -- (0.5,-0.59);
\draw[mwe] (1.375,-0.21) -- (1.5,-0.59);
\draw[mwe] (1.875,-0.21) -- (2.25,-0.59);
\node[mwk] at (0.25,-0.8) {15};
\node[mwk] at (0.75,-0.8) {29};
\node[mwk] at (1.5,-0.8) {63};
\node[mwk] at (2.25,-0.8) {91};
\node[mwk] at (1.125,0) {47};
\node[mwk,mwkM] at (1.625,0) {82};
\node[mwlab] at (1.25,-1.28) {(f) 分裂，82 上升};
\end{tikzpicture}\hspace{0.55cm}\begin{tikzpicture}[baseline=(current bounding box.north), mwk/.append style={minimum width=0.5cm, minimum height=0.42cm}]
\draw[mwe] (1.5,-0.21) -- (0.75,-0.59);
\draw[mwe] (2,-0.21) -- (2.25,-0.59);
\draw[mwe] (2.5,-0.21) -- (3.25,-0.59);
\node[mwk,mwkB] at (0.25,-0.8) {15};
\node[mwk,mwkB] at (0.75,-0.8) {29};
\node[mwk,mwkN] at (1.25,-0.8) {38};
\draw[mwbad] (0,-1.01) rectangle (1.5,-0.59);
\node[mwk,mwkN] at (2,-0.8) {54};
\node[mwk] at (2.5,-0.8) {63};
\node[mwk] at (3.25,-0.8) {91};
\node[mwk] at (1.75,0) {47};
\node[mwk] at (2.25,0) {82};
\node[mwlab] at (1.75,-1.28) {(g) 插入 54；插入 38 溢位};
\end{tikzpicture}\\*[0.45em]
\begin{tikzpicture}[baseline=(current bounding box.north), mwk/.append style={minimum width=0.5cm, minimum height=0.42cm}]
\draw[mwe] (0.875,-0.21) -- (0.25,-0.59);
\draw[mwe] (1.375,-0.21) -- (1,-0.59);
\draw[mwe] (1.875,-0.21) -- (2,-0.59);
\draw[mwe] (2.375,-0.21) -- (3,-0.59);
\node[mwk] at (0.25,-0.8) {15};
\node[mwk] at (1,-0.8) {38};
\node[mwk] at (1.75,-0.8) {54};
\node[mwk] at (2.25,-0.8) {63};
\node[mwk] at (3,-0.8) {91};
\node[mwk,mwkM] at (1.125,0) {29};
\node[mwk,mwkB] at (1.625,0) {47};
\node[mwk,mwkB] at (2.125,0) {82};
\draw[mwbad] (0.875,-0.21) rectangle (2.375,0.21);
\node[mwlab] at (1.625,-2.08) {(h) 分裂，29 上升};
\end{tikzpicture}\hspace{0.55cm}\begin{tikzpicture}[baseline=(current bounding box.north), mwk/.append style={minimum width=0.5cm, minimum height=0.42cm}]
\draw[mwe] (0.375,-1.01) -- (0.25,-1.39);
\draw[mwe] (0.875,-1.01) -- (1,-1.39);
\draw[mwe] (2.25,-1.01) -- (2,-1.39);
\draw[mwe] (2.75,-1.01) -- (3,-1.39);
\draw[mwe] (1.312,-0.21) -- (0.625,-0.59);
\draw[mwe] (1.812,-0.21) -- (2.5,-0.59);
\node[mwk] at (0.25,-1.6) {15};
\node[mwk] at (1,-1.6) {38};
\node[mwk] at (0.625,-0.8) {29};
\node[mwk] at (1.75,-1.6) {54};
\node[mwk] at (2.25,-1.6) {63};
\node[mwk] at (3,-1.6) {91};
\node[mwk] at (2.5,-0.8) {82};
\node[mwk,mwkM] at (1.562,0) {47};
\node[mwlab] at (1.625,-2.08) {(i) 根分裂，47 上升};
\end{tikzpicture}\hspace{0.55cm}\begin{tikzpicture}[baseline=(current bounding box.north), mwk/.append style={minimum width=0.5cm, minimum height=0.42cm}]
\draw[mwe] (0.375,-1.01) -- (0.25,-1.39);
\draw[mwe] (0.875,-1.01) -- (1,-1.39);
\draw[mwe] (2.625,-1.01) -- (2.25,-1.39);
\draw[mwe] (3.125,-1.01) -- (3.5,-1.39);
\draw[mwe] (1.5,-0.21) -- (0.625,-0.59);
\draw[mwe] (2,-0.21) -- (2.875,-0.59);
\node[mwk] at (0.25,-1.6) {15};
\node[mwk] at (1,-1.6) {38};
\node[mwk] at (0.625,-0.8) {29};
\node[mwk,mwkB] at (1.75,-1.6) {54};
\node[mwk,mwkB] at (2.25,-1.6) {63};
\node[mwk,mwkN] at (2.75,-1.6) {70};
\draw[mwbad] (1.5,-1.81) rectangle (3,-1.39);
\node[mwk] at (3.5,-1.6) {91};
\node[mwk] at (2.875,-0.8) {82};
\node[mwk] at (1.75,0) {47};
\node[mwlab] at (1.875,-2.08) {(j) 插入 70：溢位};
\end{tikzpicture}\\*[0.45em]
\begin{tikzpicture}[baseline=(current bounding box.north), mwk/.append style={minimum width=0.5cm, minimum height=0.42cm}]
\draw[mwe] (0.375,-1.01) -- (0.25,-1.39);
\draw[mwe] (0.875,-1.01) -- (1,-1.39);
\draw[mwe] (2,-1.01) -- (1.75,-1.39);
\draw[mwe] (2.5,-1.01) -- (2.5,-1.39);
\draw[mwe] (3,-1.01) -- (3.25,-1.39);
\draw[mwe] (1.312,-0.21) -- (0.625,-0.59);
\draw[mwe] (1.812,-0.21) -- (2.5,-0.59);
\node[mwk] at (0.25,-1.6) {15};
\node[mwk] at (1,-1.6) {38};
\node[mwk] at (0.625,-0.8) {29};
\node[mwk] at (1.75,-1.6) {54};
\node[mwk] at (2.5,-1.6) {70};
\node[mwk] at (3.25,-1.6) {91};
\node[mwk,mwkM] at (2.25,-0.8) {63};
\node[mwk] at (2.75,-0.8) {82};
\node[mwk] at (1.562,0) {47};
\node[mwlab] at (1.75,-2.08) {(k) 分裂，63 上升};
\end{tikzpicture}\hspace{0.55cm}\begin{tikzpicture}[baseline=(current bounding box.north), mwk/.append style={minimum width=0.5cm, minimum height=0.42cm}]
\draw[mwe] (0.75,-1.01) -- (0.5,-1.39);
\draw[mwe] (1.25,-1.01) -- (1.5,-1.39);
\draw[mwe] (2.5,-1.01) -- (2.25,-1.39);
\draw[mwe] (3,-1.01) -- (3,-1.39);
\draw[mwe] (3.5,-1.01) -- (3.75,-1.39);
\draw[mwe] (1.75,-0.21) -- (1,-0.59);
\draw[mwe] (2.25,-0.21) -- (3,-0.59);
\node[mwk] at (0.25,-1.6) {15};
\node[mwk,mwkN] at (0.75,-1.6) {21};
\node[mwk] at (1.5,-1.6) {38};
\node[mwk] at (1,-0.8) {29};
\node[mwk] at (2.25,-1.6) {54};
\node[mwk] at (3,-1.6) {70};
\node[mwk] at (3.75,-1.6) {91};
\node[mwk] at (2.75,-0.8) {63};
\node[mwk] at (3.25,-0.8) {82};
\node[mwk] at (2,0) {47};
\node[mwlab] at (2,-2.08) {(l) 插入 21};
\end{tikzpicture}
\diagcaption{依序插入 $47, 82, 15, 63, 29, 91, 54, 38, 70, 21$。 (g) 的 $38$ 讓葉子
$\{15, 29\}$ 溢位，(h) 分裂後 $29$ 上升，根變成三個鍵，又溢位；(i) 根分裂，$47$ 成為
新的根，高度從 $2$ 變成 $3$。}
\label{fig:mw-23-ins}
\end{figure}

(h) 到 (i) 是\textbf{連鎖分裂}：一次插入可能讓從葉子到根路上的每個節點都分裂，但每層
最多分裂一次，所以插入的時間是 $O(h) = O(\lg n)$。

\subsection{刪除：借與合併}
\label{sec:mw-23-delete}

刪除要先處理一件事：被刪的鍵如果在內部節點，不能直接拿掉，否則它兩邊的子樹會少一個
分隔鍵。 做法和第~9~章的二元搜尋樹一樣，用\textbf{後繼}（右邊子樹裡最小的鍵）取代
它，再改成刪除那個後繼。 從右邊的子樹一路往最左邊走，最後停在一片葉子上，後繼就是
那片葉子最左邊的鍵。 所以不管刪哪個鍵，真正拿掉鍵的地方都是一片葉子。

葉子原本有兩個鍵，拿掉一個還剩一個，結束。 原本只有一個鍵，拿掉就空了，叫做
\textbf{不足}（underflow），要從旁邊補：

\begin{itemize}
\item \textbf{借}：相鄰的兄弟是 3-節點，有多的鍵。 父節點裡夾在兩人中間的分隔鍵
下來放進空的節點，兄弟靠近這一邊的鍵上去補分隔鍵的位置。 如果是內
部節點在借，兄弟靠近這一邊的那棵子樹也跟著搬過來。 父節點的鍵數不變，結束。
\item \textbf{合併}：相鄰的兄弟都是 2-節點，沒有多的鍵可借。 把空的節點、父節點裡的
分隔鍵、兄弟三者合成一個 3-節點。 父節點少了一個鍵，可能變空，就在父節點用同樣的
方法往上處理。 根變空時把它拿掉，它唯一的子節點成為新的根，樹高減 $1$。
\end{itemize}

借只看相鄰的兄弟，也就是和不足的節點共用一個分隔鍵的那一個或兩個。 兩個兄
弟都可能時，本書固定的順序是：先看右兄弟能不能借，再看左兄弟；都不能借就合併，有右
兄弟就和右兄弟合併，沒有才和左兄弟合併。 這個順序只是一種約定，換一種順序得到的樹
也合法，但會不一樣。

圖~\ref{fig:mw-23-del} 從圖~\ref{fig:mw-23-ins}(l) 的樹依序刪除 $21, 38, 63, 47$。

\begin{center}
\begin{tikzpicture}[baseline=(current bounding box.north), mwk/.append style={minimum width=0.5cm, minimum height=0.42cm}]
\draw[mwe] (0.375,-1.01) -- (0.25,-1.39);
\draw[mwe] (0.875,-1.01) -- (1,-1.39);
\draw[mwe] (2,-1.01) -- (1.75,-1.39);
\draw[mwe] (2.5,-1.01) -- (2.5,-1.39);
\draw[mwe] (3,-1.01) -- (3.25,-1.39);
\draw[mwe] (1.312,-0.21) -- (0.625,-0.59);
\draw[mwe] (1.812,-0.21) -- (2.5,-0.59);
\node[mwk] at (0.25,-1.6) {15};
\node[mwk] at (1,-1.6) {38};
\node[mwk] at (0.625,-0.8) {29};
\node[mwk] at (1.75,-1.6) {54};
\node[mwk] at (2.5,-1.6) {70};
\node[mwk] at (3.25,-1.6) {91};
\node[mwk] at (2.25,-0.8) {63};
\node[mwk] at (2.75,-0.8) {82};
\node[mwk] at (1.562,0) {47};
\node[mwlab] at (1.75,-2.08) {(a) 刪除 21};
\end{tikzpicture}\hspace{0.55cm}\begin{tikzpicture}[baseline=(current bounding box.north), mwk/.append style={minimum width=0.5cm, minimum height=0.42cm}]
\draw[mwe] (0.375,-1.01) -- (0.25,-1.39);
\draw[mwe] (0.875,-1.01) -- (1,-1.39);
\draw[mwe] (2,-1.01) -- (1.75,-1.39);
\draw[mwe] (2.5,-1.01) -- (2.5,-1.39);
\draw[mwe] (3,-1.01) -- (3.25,-1.39);
\draw[mwe] (1.312,-0.21) -- (0.625,-0.59);
\draw[mwe] (1.812,-0.21) -- (2.5,-0.59);
\node[mwk] at (0.25,-1.6) {15};
\draw[mwempty] (0.75,-1.81) rectangle (1.25,-1.39);
\node[mwk] at (0.625,-0.8) {29};
\node[mwk] at (1.75,-1.6) {54};
\node[mwk] at (2.5,-1.6) {70};
\node[mwk] at (3.25,-1.6) {91};
\node[mwk] at (2.25,-0.8) {63};
\node[mwk] at (2.75,-0.8) {82};
\node[mwk] at (1.562,0) {47};
\node[mwlab] at (1.75,-2.08) {(b) 刪除 38：葉子變空};
\end{tikzpicture}\hspace{0.55cm}\begin{tikzpicture}[baseline=(current bounding box.north), mwk/.append style={minimum width=0.5cm, minimum height=0.42cm}]
\draw[mwe] (0.5,-1.01) -- (0.5,-1.39);
\draw[mwe] (1.75,-1.01) -- (1.5,-1.39);
\draw[mwe] (2.25,-1.01) -- (2.25,-1.39);
\draw[mwe] (2.75,-1.01) -- (3,-1.39);
\draw[mwe] (1.125,-0.21) -- (0.5,-0.59);
\draw[mwe] (1.625,-0.21) -- (2.25,-0.59);
\node[mwk] at (0.25,-1.6) {15};
\node[mwk,mwkM] at (0.75,-1.6) {29};
\draw[mwempty] (0.25,-1.01) rectangle (0.75,-0.59);
\node[mwk] at (1.5,-1.6) {54};
\node[mwk] at (2.25,-1.6) {70};
\node[mwk] at (3,-1.6) {91};
\node[mwk] at (2,-0.8) {63};
\node[mwk] at (2.5,-0.8) {82};
\node[mwk] at (1.375,0) {47};
\node[mwlab] at (1.625,-2.08) {(c) 和左兄弟合併};
\end{tikzpicture}\hspace{0.55cm}\begin{tikzpicture}[baseline=(current bounding box.north), mwk/.append style={minimum width=0.5cm, minimum height=0.42cm}]
\draw[mwe] (0.75,-1.01) -- (0.5,-1.39);
\draw[mwe] (1.25,-1.01) -- (1.5,-1.39);
\draw[mwe] (2.375,-1.01) -- (2.25,-1.39);
\draw[mwe] (2.875,-1.01) -- (3,-1.39);
\draw[mwe] (1.562,-0.21) -- (1,-0.59);
\draw[mwe] (2.062,-0.21) -- (2.625,-0.59);
\node[mwk] at (0.25,-1.6) {15};
\node[mwk] at (0.75,-1.6) {29};
\node[mwk] at (1.5,-1.6) {54};
\node[mwk,mwkM] at (1,-0.8) {47};
\node[mwk] at (2.25,-1.6) {70};
\node[mwk] at (3,-1.6) {91};
\node[mwk] at (2.625,-0.8) {82};
\node[mwk,mwkM] at (1.812,0) {63};
\node[mwlab] at (1.625,-2.08) {(d) 向右兄弟借};
\end{tikzpicture}\\*[0.45em]
\begin{tikzpicture}[baseline=(current bounding box.north), mwk/.append style={minimum width=0.5cm, minimum height=0.42cm}]
\draw[mwe] (0.75,-1.01) -- (0.5,-1.39);
\draw[mwe] (1.25,-1.01) -- (1.5,-1.39);
\draw[mwe] (2.375,-1.01) -- (2.25,-1.39);
\draw[mwe] (2.875,-1.01) -- (3,-1.39);
\draw[mwe] (1.562,-0.21) -- (1,-0.59);
\draw[mwe] (2.062,-0.21) -- (2.625,-0.59);
\node[mwk] at (0.25,-1.6) {15};
\node[mwk] at (0.75,-1.6) {29};
\node[mwk] at (1.5,-1.6) {54};
\node[mwk] at (1,-0.8) {47};
\draw[mwempty] (2,-1.81) rectangle (2.5,-1.39);
\node[mwk] at (3,-1.6) {91};
\node[mwk] at (2.625,-0.8) {82};
\node[mwk,mwkM] at (1.812,0) {70};
\node[mwlab] at (1.625,-2.08) {(e) 刪除 63：換成後繼 70};
\end{tikzpicture}\hspace{0.55cm}\begin{tikzpicture}[baseline=(current bounding box.north), mwk/.append style={minimum width=0.5cm, minimum height=0.42cm}]
\draw[mwe] (0.75,-1.01) -- (0.5,-1.39);
\draw[mwe] (1.25,-1.01) -- (1.5,-1.39);
\draw[mwe] (2.5,-1.01) -- (2.5,-1.39);
\draw[mwe] (1.5,-0.21) -- (1,-0.59);
\draw[mwe] (2,-0.21) -- (2.5,-0.59);
\node[mwk] at (0.25,-1.6) {15};
\node[mwk] at (0.75,-1.6) {29};
\node[mwk] at (1.5,-1.6) {54};
\node[mwk] at (1,-0.8) {47};
\node[mwk,mwkM] at (2.25,-1.6) {82};
\node[mwk] at (2.75,-1.6) {91};
\draw[mwempty] (2.25,-1.01) rectangle (2.75,-0.59);
\node[mwk] at (1.75,0) {70};
\node[mwlab] at (1.5,-2.08) {(f) 和右兄弟合併};
\end{tikzpicture}\hspace{0.55cm}\begin{tikzpicture}[baseline=(current bounding box.north), mwk/.append style={minimum width=0.5cm, minimum height=0.42cm}]
\draw[mwe] (1,-1.01) -- (0.5,-1.39);
\draw[mwe] (1.5,-1.01) -- (1.5,-1.39);
\draw[mwe] (2,-1.01) -- (2.5,-1.39);
\draw[mwe] (1.5,-0.21) -- (1.5,-0.59);
\node[mwk] at (0.25,-1.6) {15};
\node[mwk] at (0.75,-1.6) {29};
\node[mwk] at (1.5,-1.6) {54};
\node[mwk] at (2.25,-1.6) {82};
\node[mwk] at (2.75,-1.6) {91};
\node[mwk] at (1.25,-0.8) {47};
\node[mwk,mwkM] at (1.75,-0.8) {70};
\draw[mwempty] (1.25,-0.21) rectangle (1.75,0.21);
\node[mwlab] at (1.5,-2.08) {(g) 和左兄弟合併};
\end{tikzpicture}\hspace{0.55cm}\begin{tikzpicture}[baseline=(current bounding box.north), mwk/.append style={minimum width=0.5cm, minimum height=0.42cm}]
\draw[mwe] (1,-0.21) -- (0.5,-0.59);
\draw[mwe] (1.5,-0.21) -- (1.5,-0.59);
\draw[mwe] (2,-0.21) -- (2.5,-0.59);
\node[mwk] at (0.25,-0.8) {15};
\node[mwk] at (0.75,-0.8) {29};
\node[mwk] at (1.5,-0.8) {54};
\node[mwk] at (2.25,-0.8) {82};
\node[mwk] at (2.75,-0.8) {91};
\node[mwk] at (1.25,0) {47};
\node[mwk] at (1.75,0) {70};
\node[mwlab] at (1.5,-2.08) {(h) 拿掉空的根};
\end{tikzpicture}\\*[0.45em]
\begin{tikzpicture}[baseline=(current bounding box.north), mwk/.append style={minimum width=0.5cm, minimum height=0.42cm}]
\draw[mwe] (1,-0.21) -- (0.5,-0.59);
\draw[mwe] (1.5,-0.21) -- (1.5,-0.59);
\draw[mwe] (2,-0.21) -- (2.5,-0.59);
\node[mwk] at (0.25,-0.8) {15};
\node[mwk] at (0.75,-0.8) {29};
\draw[mwempty] (1.25,-1.01) rectangle (1.75,-0.59);
\node[mwk] at (2.25,-0.8) {82};
\node[mwk] at (2.75,-0.8) {91};
\node[mwk,mwkM] at (1.25,0) {54};
\node[mwk] at (1.75,0) {70};
\node[mwlab] at (1.5,-1.28) {(i) 刪除 47：換成後繼 54};
\end{tikzpicture}\hspace{0.55cm}\begin{tikzpicture}[baseline=(current bounding box.north), mwk/.append style={minimum width=0.5cm, minimum height=0.42cm}]
\draw[mwe] (0.875,-0.21) -- (0.5,-0.59);
\draw[mwe] (1.375,-0.21) -- (1.5,-0.59);
\draw[mwe] (1.875,-0.21) -- (2.25,-0.59);
\node[mwk] at (0.25,-0.8) {15};
\node[mwk] at (0.75,-0.8) {29};
\node[mwk,mwkM] at (1.5,-0.8) {70};
\node[mwk] at (2.25,-0.8) {91};
\node[mwk] at (1.125,0) {54};
\node[mwk,mwkM] at (1.625,0) {82};
\node[mwlab] at (1.25,-1.28) {(j) 向右兄弟借};
\end{tikzpicture}
\diagcaption{依序刪除 $21, 38, 63, 47$。 虛線框是變空的節點。}
\label{fig:mw-23-del}
\end{center}

刪除 $21$ 之後葉子還剩 $15$，結束。 其他三次：
\begin{itemize}
\item 刪除 $38$：葉子變空。 它是父節點 $\{29\}$ 的最後一個子節點，沒有右兄弟；左兄
弟 $\{15\}$ 只有一個鍵，不能借。 和左兄弟合併：$15$、分隔鍵 $29$ 合成 $\{15, 29\}$。
父節點 $\{29\}$ 因此變空，它的右兄弟 $\{63, 82\}$ 有兩個鍵，可以借：根的 $47$ 下來，
$63$ 上去當根，$63$ 原本左邊的子樹 $\{54\}$ 改接到 $\{47\}$ 的右邊。
\item 刪除 $63$：它是根，後繼是右子樹最左的葉子裡的 $70$。 $70$ 移上去取代 $63$，葉
子 $\{70\}$ 變空。 它沒有左兄弟，右兄弟 $\{91\}$ 不能借，合併成 $\{82, 91\}$；父節點
變空，左兄弟 $\{47\}$ 也不能借，再合併成 $\{47, 70\}$；根變空，拿掉。
\item 刪除 $47$：後繼 $54$ 上去，葉子變空。 右兄弟 $\{82, 91\}$ 可以借：$70$ 下來，
$82$ 上去。
\end{itemize}

刪除和插入一樣，每層最多處理一次，時間 $O(\lg n)$。 合併只發生在兄弟都不能借的時
候，借一次之後就結束，不會再往上。

\subsection{2-3-4 樹}
\label{sec:mw-234}

\textbf{2-3-4 樹}（2-3-4 tree）把定義~\ref{def:mw-23} 放寬一點：內部節點可以是
2-節點、3-節點或 4-節點（$3$ 個鍵、$4$ 個子節點），葉子放 $1$ 到 $3$ 個鍵，所有葉子
同一層。 和 2-3 樹一樣算，$n$ 個鍵的 2-3-4 樹高度滿足
$\log_4(n + 1) \le h \le \lg(n + 1)$。

插入可以照 2-3 樹的做法，葉子放到 $4$ 個鍵時分裂。 但 $4$ 個鍵沒有正中間的那一個，
要先約定往上送第 $2$ 個還是第 $3$ 個。 2-3-4 樹更常用的是\textbf{由上而下}
（top-down）的插入：從根往下走的路上，遇到 4-節點就先把它分裂，中間的鍵
（$3$ 個鍵的第 $2$ 個）送進父節點。 父節點在上一步已經處理過，不會是 4-節點，放得
下。 走到葉子時，葉子一定不是 4-節點，直接放進去就好，不用再往回走。

虛擬碼裡 $T.\mathit{root}$ 是樹根，$x.c[i]$ 是節點 $x$ 由左往右數的第 $i$ 個子節點；
$x.c[i]$ 有 $3$ 個鍵就是 4-節點。 第~\ref{sec:mw-btree}~節寫 B 樹的虛擬碼時，會把節點的欄位
完整列出來。

\begin{pseudo}[label=alg:mw-234-insert]{2-3-4 樹由上而下的插入（樹不是空的）}
\Procedure{Insert-234}{$T, k$}
    \If{$T.\mathit{root}$ 有 $3$ 個鍵}
        \State 根的中間鍵自己成為新的根，左右兩個鍵各成一個 2-節點
    \EndIf
    \State $x \gets T.\mathit{root}$
    \While{$x$ 不是葉節點}
        \State $i \gets$ $k$ 應該往下走的子樹編號
        \If{$x.c[i]$ 有 $3$ 個鍵}
            \State $x.c[i]$ 的中間鍵插進 $x$，$x.c[i]$ 分成兩個 2-節點
            \If{$k >$ 剛插進 $x$ 的那個鍵}
                \State $i \gets i + 1$
            \EndIf
        \EndIf
        \State $x \gets x.c[i]$
    \EndWhile
    \State 把 $k$ 依大小放進葉節點 $x$
\EndProcedure
\end{pseudo}

第 3 行是唯一讓樹長高的地方。 第 8 行分裂之後，原本的 $x.c[i]$ 變成兩個節點，要看
$k$ 在上來的鍵的哪一邊，決定走左半還是右半（第 9、10 行）。
圖~\ref{fig:mw-234-ins} 把 $36, 72, 18, 55, 90, 27, 64, 81, 45, 9, 50, 99$ 依序插入
空的 2-3-4 樹。

\begin{center}
\begin{tikzpicture}[baseline=(current bounding box.north), mwk/.append style={minimum width=0.5cm, minimum height=0.42cm}]
\node[mwk,mwkN] at (0.25,0) {18};
\node[mwk,mwkN] at (0.75,0) {36};
\node[mwk,mwkN] at (1.25,0) {72};
\node[mwlab] at (0.75,-1.28) {(a) 插入 36、72、18};
\end{tikzpicture}\hspace{0.55cm}\begin{tikzpicture}[baseline=(current bounding box.north), mwk/.append style={minimum width=0.5cm, minimum height=0.42cm}]
\draw[mwe] (0.5,-0.21) -- (0.25,-0.59);
\draw[mwe] (1,-0.21) -- (1.25,-0.59);
\node[mwk] at (0.25,-0.8) {18};
\node[mwk,mwkN] at (1,-0.8) {55};
\node[mwk] at (1.5,-0.8) {72};
\node[mwk,mwkM] at (0.75,0) {36};
\node[mwlab] at (0.875,-1.28) {(b) 根先分裂，再插入 55};
\end{tikzpicture}\hspace{0.55cm}\begin{tikzpicture}[baseline=(current bounding box.north), mwk/.append style={minimum width=0.5cm, minimum height=0.42cm}]
\draw[mwe] (1,-0.21) -- (0.5,-0.59);
\draw[mwe] (1.5,-0.21) -- (2,-0.59);
\node[mwk] at (0.25,-0.8) {18};
\node[mwk,mwkN] at (0.75,-0.8) {27};
\node[mwk] at (1.5,-0.8) {55};
\node[mwk] at (2,-0.8) {72};
\node[mwk,mwkN] at (2.5,-0.8) {90};
\node[mwk] at (1.25,0) {36};
\node[mwlab] at (1.375,-1.28) {(c) 插入 90、27};
\end{tikzpicture}\\*[0.45em]
\begin{tikzpicture}[baseline=(current bounding box.north), mwk/.append style={minimum width=0.5cm, minimum height=0.42cm}]
\draw[mwe] (1.125,-0.21) -- (0.5,-0.59);
\draw[mwe] (1.625,-0.21) -- (1.75,-0.59);
\draw[mwe] (2.125,-0.21) -- (2.75,-0.59);
\node[mwk] at (0.25,-0.8) {18};
\node[mwk] at (0.75,-0.8) {27};
\node[mwk] at (1.5,-0.8) {55};
\node[mwk,mwkN] at (2,-0.8) {64};
\node[mwk] at (2.75,-0.8) {90};
\node[mwk] at (1.375,0) {36};
\node[mwk,mwkM] at (1.875,0) {72};
\node[mwlab] at (1.5,-1.28) {(d) 先分裂，再插入 64};
\end{tikzpicture}\hspace{0.55cm}\begin{tikzpicture}[baseline=(current bounding box.north), mwk/.append style={minimum width=0.5cm, minimum height=0.42cm}]
\draw[mwe] (1.875,-0.21) -- (0.75,-0.59);
\draw[mwe] (2.375,-0.21) -- (2.5,-0.59);
\draw[mwe] (2.875,-0.21) -- (4,-0.59);
\node[mwk,mwkN] at (0.25,-0.8) {9};
\node[mwk] at (0.75,-0.8) {18};
\node[mwk] at (1.25,-0.8) {27};
\node[mwk,mwkN] at (2,-0.8) {45};
\node[mwk] at (2.5,-0.8) {55};
\node[mwk] at (3,-0.8) {64};
\node[mwk,mwkN] at (3.75,-0.8) {81};
\node[mwk] at (4.25,-0.8) {90};
\node[mwk] at (2.125,0) {36};
\node[mwk] at (2.625,0) {72};
\node[mwlab] at (2.25,-1.28) {(e) 插入 81、45、9};
\end{tikzpicture}\hspace{0.55cm}\begin{tikzpicture}[baseline=(current bounding box.north), mwk/.append style={minimum width=0.5cm, minimum height=0.42cm}]
\draw[mwe] (1.75,-0.21) -- (0.75,-0.59);
\draw[mwe] (2.25,-0.21) -- (2.25,-0.59);
\draw[mwe] (2.75,-0.21) -- (3.25,-0.59);
\draw[mwe] (3.25,-0.21) -- (4.25,-0.59);
\node[mwk] at (0.25,-0.8) {9};
\node[mwk] at (0.75,-0.8) {18};
\node[mwk] at (1.25,-0.8) {27};
\node[mwk] at (2,-0.8) {45};
\node[mwk,mwkN] at (2.5,-0.8) {50};
\node[mwk] at (3.25,-0.8) {64};
\node[mwk] at (4,-0.8) {81};
\node[mwk] at (4.5,-0.8) {90};
\node[mwk] at (2,0) {36};
\node[mwk,mwkM] at (2.5,0) {55};
\node[mwk] at (3,0) {72};
\node[mwlab] at (2.375,-1.28) {(f) 先分裂，再插入 50};
\end{tikzpicture}\\*[0.45em]
\begin{tikzpicture}[baseline=(current bounding box.north), mwk/.append style={minimum width=0.5cm, minimum height=0.42cm}]
\draw[mwe] (1.25,-1.01) -- (0.75,-1.39);
\draw[mwe] (1.75,-1.01) -- (2.25,-1.39);
\draw[mwe] (3.5,-1.01) -- (3.25,-1.39);
\draw[mwe] (4,-1.01) -- (4.25,-1.39);
\draw[mwe] (2.375,-0.21) -- (1.5,-0.59);
\draw[mwe] (2.875,-0.21) -- (3.75,-0.59);
\node[mwk] at (0.25,-1.6) {9};
\node[mwk] at (0.75,-1.6) {18};
\node[mwk] at (1.25,-1.6) {27};
\node[mwk] at (2,-1.6) {45};
\node[mwk] at (2.5,-1.6) {50};
\node[mwk] at (1.5,-0.8) {36};
\node[mwk] at (3.25,-1.6) {64};
\node[mwk] at (4,-1.6) {81};
\node[mwk] at (4.5,-1.6) {90};
\node[mwk] at (3.75,-0.8) {72};
\node[mwk,mwkM] at (2.625,0) {55};
\node[mwlab] at (2.375,-2.08) {(g) 插入 99 之前，根先分裂};
\end{tikzpicture}\hspace{0.55cm}\begin{tikzpicture}[baseline=(current bounding box.north), mwk/.append style={minimum width=0.5cm, minimum height=0.42cm}]
\draw[mwe] (1.25,-1.01) -- (0.75,-1.39);
\draw[mwe] (1.75,-1.01) -- (2.25,-1.39);
\draw[mwe] (3.625,-1.01) -- (3.25,-1.39);
\draw[mwe] (4.125,-1.01) -- (4.5,-1.39);
\draw[mwe] (2.438,-0.21) -- (1.5,-0.59);
\draw[mwe] (2.938,-0.21) -- (3.875,-0.59);
\node[mwk] at (0.25,-1.6) {9};
\node[mwk] at (0.75,-1.6) {18};
\node[mwk] at (1.25,-1.6) {27};
\node[mwk] at (2,-1.6) {45};
\node[mwk] at (2.5,-1.6) {50};
\node[mwk] at (1.5,-0.8) {36};
\node[mwk] at (3.25,-1.6) {64};
\node[mwk] at (4,-1.6) {81};
\node[mwk] at (4.5,-1.6) {90};
\node[mwk,mwkN] at (5,-1.6) {99};
\node[mwk] at (3.875,-0.8) {72};
\node[mwk] at (2.688,0) {55};
\node[mwlab] at (2.625,-2.08) {(h) 插入 99};
\end{tikzpicture}
\diagcaption{2-3-4 樹由上而下的插入。 (b) 插入 $55$ 時根 $\{18, 36, 72\}$ 是 4-節點，先
分裂（$36$ 上升）再往下走。 (d)、(f) 分裂的是路上的子節點。 (g) 插入 $99$ 時，要去的
葉子 $\{81, 90\}$ 明明放得下，但路上的根是 4-節點，照規則還是先分裂，高度變成 $3$。}
\label{fig:mw-234-ins}
\end{center}

(g) 顯示了兩種做法的差別。 同樣十二個鍵改用由下而上的做法（溢位時送第 $2$ 個鍵，即
第~\ref{sec:mw-insert}~節的規則），最後得到的是高度 $2$ 的
$\{36, 64, 81\}$，下面接 $\{9, 18, 27\}$、$\{45, 50, 55\}$、$\{72\}$、$\{90, 99\}$。
兩棵樹都是合法的 2-3-4 樹。 由上而下的做法可能比較早分裂，換來的好處是只走一趟，不
用記住回去的路。

刪除和 2-3 樹的規則完全一樣，只是 3-節點和 4-節點都有多的鍵可以借。
圖~\ref{fig:mw-234-del} 從圖~\ref{fig:mw-234-ins}(h) 依序刪除 $99, 64, 72$。

\begin{center}
\begin{tikzpicture}[baseline=(current bounding box.north), mwk/.append style={minimum width=0.5cm, minimum height=0.42cm}]
\draw[mwe] (1.25,-1.01) -- (0.75,-1.39);
\draw[mwe] (1.75,-1.01) -- (2.25,-1.39);
\draw[mwe] (3.5,-1.01) -- (3.25,-1.39);
\draw[mwe] (4,-1.01) -- (4.25,-1.39);
\draw[mwe] (2.375,-0.21) -- (1.5,-0.59);
\draw[mwe] (2.875,-0.21) -- (3.75,-0.59);
\node[mwk] at (0.25,-1.6) {9};
\node[mwk] at (0.75,-1.6) {18};
\node[mwk] at (1.25,-1.6) {27};
\node[mwk] at (2,-1.6) {45};
\node[mwk] at (2.5,-1.6) {50};
\node[mwk] at (1.5,-0.8) {36};
\node[mwk] at (3.25,-1.6) {64};
\node[mwk] at (4,-1.6) {81};
\node[mwk] at (4.5,-1.6) {90};
\node[mwk] at (3.75,-0.8) {72};
\node[mwk] at (2.625,0) {55};
\node[mwlab] at (2.375,-2.08) {(a) 刪除 99};
\end{tikzpicture}\hspace{0.55cm}\begin{tikzpicture}[baseline=(current bounding box.north), mwk/.append style={minimum width=0.5cm, minimum height=0.42cm}]
\draw[mwe] (1.25,-1.01) -- (0.75,-1.39);
\draw[mwe] (1.75,-1.01) -- (2.25,-1.39);
\draw[mwe] (3.375,-1.01) -- (3.25,-1.39);
\draw[mwe] (3.875,-1.01) -- (4,-1.39);
\draw[mwe] (2.312,-0.21) -- (1.5,-0.59);
\draw[mwe] (2.812,-0.21) -- (3.625,-0.59);
\node[mwk] at (0.25,-1.6) {9};
\node[mwk] at (0.75,-1.6) {18};
\node[mwk] at (1.25,-1.6) {27};
\node[mwk] at (2,-1.6) {45};
\node[mwk] at (2.5,-1.6) {50};
\node[mwk] at (1.5,-0.8) {36};
\node[mwk,mwkM] at (3.25,-1.6) {72};
\node[mwk] at (4,-1.6) {90};
\node[mwk,mwkM] at (3.625,-0.8) {81};
\node[mwk] at (2.562,0) {55};
\node[mwlab] at (2.125,-2.08) {(b) 刪除 64：向右兄弟借};
\end{tikzpicture}\hspace{0.55cm}\begin{tikzpicture}[baseline=(current bounding box.north), mwk/.append style={minimum width=0.5cm, minimum height=0.42cm}]
\draw[mwe] (1.25,-1.01) -- (0.75,-1.39);
\draw[mwe] (1.75,-1.01) -- (2.25,-1.39);
\draw[mwe] (3.5,-1.01) -- (3.5,-1.39);
\draw[mwe] (2.25,-0.21) -- (1.5,-0.59);
\draw[mwe] (2.75,-0.21) -- (3.5,-0.59);
\node[mwk] at (0.25,-1.6) {9};
\node[mwk] at (0.75,-1.6) {18};
\node[mwk] at (1.25,-1.6) {27};
\node[mwk] at (2,-1.6) {45};
\node[mwk] at (2.5,-1.6) {50};
\node[mwk] at (1.5,-0.8) {36};
\node[mwk,mwkM] at (3.25,-1.6) {81};
\node[mwk] at (3.75,-1.6) {90};
\draw[mwempty] (3.25,-1.01) rectangle (3.75,-0.59);
\node[mwk] at (2.5,0) {55};
\node[mwlab] at (2,-2.08) {(c) 刪除 72：和右兄弟合併};
\end{tikzpicture}\\*[0.45em]
\begin{tikzpicture}[baseline=(current bounding box.north), mwk/.append style={minimum width=0.5cm, minimum height=0.42cm}]
\draw[mwe] (1.625,-0.21) -- (0.75,-0.59);
\draw[mwe] (2.125,-0.21) -- (2.25,-0.59);
\draw[mwe] (2.625,-0.21) -- (3.5,-0.59);
\node[mwk] at (0.25,-0.8) {9};
\node[mwk] at (0.75,-0.8) {18};
\node[mwk] at (1.25,-0.8) {27};
\node[mwk] at (2,-0.8) {45};
\node[mwk] at (2.5,-0.8) {50};
\node[mwk] at (3.25,-0.8) {81};
\node[mwk] at (3.75,-0.8) {90};
\node[mwk] at (1.875,0) {36};
\node[mwk] at (2.375,0) {55};
\node[mwlab] at (2,-1.28) {(d) 再合併，拿掉空的根};
\end{tikzpicture}
\diagcaption{2-3-4 樹的刪除。 (b) 刪除 $64$，葉子變空，右兄弟 $\{81, 90\}$ 可以借：
$72$ 下來，$81$ 上去。 (c) 刪除 $72$，右兄弟 $\{90\}$ 只有一個鍵，合併成
$\{81, 90\}$，父節點變空。 (d) 父節點只有左兄弟 $\{36\}$，不能借，合併成
$\{36, 55\}$，根變空後拿掉。}
\label{fig:mw-234-del}
\end{center}

\subsection{2-3-4 樹和紅黑樹}
\label{sec:mw-rb}

第~\ref{sec:bal-234}~節提到，把紅黑樹的每個黑色節點和它的紅色子節點合成一個大節點，就得到一棵
2-3-4 樹。 反過來，每棵 2-3-4 樹都可以這樣存成紅黑樹：
\begin{itemize}
\item 2-節點 $\{a\}$：一個黑色節點 $a$。
\item 3-節點 $\{a, b\}$：黑色的 $a$ 加上紅色的右子節點 $b$；或者黑色的 $b$ 加上紅色
的左子節點 $a$，兩種都可以。
\item 4-節點 $\{a, b, c\}$：黑色的 $b$，左右兩個子節點 $a$、$c$ 都是紅色。
\end{itemize}
原本 2-3-4 樹的子樹，依照鍵的範圍接到這幾個二元節點底下。 圖~\ref{fig:mw-234-rb}
是圖~\ref{fig:mw-234-ins}(h) 的 2-3-4 樹和它對應的一棵紅黑樹，灰底框起來的就是同一
個 2-3-4 節點。

\begin{center}
\begin{tikzpicture}[baseline=(current bounding box.north), mwk/.append style={minimum width=0.5cm, minimum height=0.42cm}]
\draw[mwe] (1.275,-1.11) -- (0.75,-1.59);
\draw[mwe] (1.775,-1.11) -- (2.3,-1.59);
\draw[mwe] (3.75,-1.11) -- (3.35,-1.59);
\draw[mwe] (4.25,-1.11) -- (4.65,-1.59);
\draw[mwe] (2.512,-0.21) -- (1.525,-0.69);
\draw[mwe] (3.012,-0.21) -- (4,-0.69);
\node[mwk] at (0.25,-1.8) {9};
\node[mwk] at (0.75,-1.8) {18};
\node[mwk] at (1.25,-1.8) {27};
\node[mwk] at (2.05,-1.8) {45};
\node[mwk] at (2.55,-1.8) {50};
\node[mwk] at (1.525,-0.9) {36};
\node[mwk] at (3.35,-1.8) {64};
\node[mwk] at (4.15,-1.8) {81};
\node[mwk] at (4.65,-1.8) {90};
\node[mwk] at (5.15,-1.8) {99};
\node[mwk] at (4,-0.9) {72};
\node[mwk] at (2.762,0) {55};
\node[mwlab] at (2.7,-2.28) {(a) 2-3-4 樹};
\end{tikzpicture}\hspace{1.0cm}\begin{tikzpicture}[baseline=(current bounding box.north)]
\node[mwR] (rb9) at (0.00,-2.46) {9};
\node[mwB] (rb18) at (0.58,-1.64) {18};
\node[mwR] (rb27) at (1.16,-2.46) {27};
\node[mwB] (rb36) at (1.74,-0.82) {36};
\node[mwB] (rb45) at (2.32,-1.64) {45};
\node[mwR] (rb50) at (2.90,-2.46) {50};
\node[mwB] (rb55) at (3.48,0.00) {55};
\node[mwB] (rb64) at (4.06,-1.64) {64};
\node[mwB] (rb72) at (4.64,-0.82) {72};
\node[mwR] (rb81) at (5.22,-2.46) {81};
\node[mwB] (rb90) at (5.80,-1.64) {90};
\node[mwR] (rb99) at (6.38,-2.46) {99};
\draw[mwe] (rb55) -- (rb36);
\draw[mwe] (rb36) -- (rb18);
\draw[mwe] (rb18) -- (rb9);
\draw[mwe] (rb18) -- (rb27);
\draw[mwe] (rb36) -- (rb45);
\draw[mwe] (rb45) -- (rb50);
\draw[mwe] (rb55) -- (rb72);
\draw[mwe] (rb72) -- (rb64);
\draw[mwe] (rb72) -- (rb90);
\draw[mwe] (rb90) -- (rb81);
\draw[mwe] (rb90) -- (rb99);
\begin{scope}[on background layer]
\node[mwgrp, fit=(rb9) (rb18) (rb27)] {};
\node[mwgrp, fit=(rb45) (rb50)] {};
\node[mwgrp, fit=(rb81) (rb90) (rb99)] {};
\end{scope}
\node[mwlab] at (3.19,-3.21) {(b) 對應的紅黑樹};
\end{tikzpicture}
\diagcaption{同一棵樹的兩種畫法。 4-節點 $\{9, 18, 27\}$ 變成黑色的 $18$ 帶兩個紅色
子節點；3-節點 $\{45, 50\}$ 這裡選了黑色的 $45$ 帶紅色的右子節點 $50$。 紅黑樹的黑
高是 $3$，等於 2-3-4 樹的高度。}
\label{fig:mw-234-rb}
\end{center}

檢查第~10~章的五條性質：紅色節點只出現在黑色節點底下、和它同一組，所以不會有兩個紅
色相連；從根往下每經過 2-3-4 樹的一層，剛好經過一個黑色節點，所有葉子同層，所以每
條路徑的黑色個數一樣，而且（照第~10~章不算自己、算 NIL 的數法）根的黑高就是 2-3-4
樹的高度。 一棵紅黑樹只對應一棵 2-3-4 樹；一棵 2-3-4 樹有幾個 3-節點，就有 $2$ 的
那麼多次方種紅黑樹。

這個對應也給了第~10~章高度上界的另一個證明：2-3-4 樹的高度不超過 $\lg(n+1)$，每一層
在紅黑樹裡最多佔兩層（黑色加紅色），所以紅黑樹的高度不超過 $2\lg(n+1)$。

兩邊的操作也一一對應。 下表的情況編號是第~10~章的編號。

\begin{center}
\begin{tabular}{p{0.40\linewidth}p{0.50\linewidth}}
\toprule
紅黑樹 & 2-3-4 樹 \\
\midrule
新節點塗紅 & 新鍵放進葉節點，不增加層數 \\
插入情況 1：叔叔是紅的，$p$、$u$ 改黑，$g$ 改紅 & $g$、$p$、$u$ 是一個 4-節點，加進
新鍵後分裂：中間的 $g$ 變紅，就是往上併進父節點，$p$、$u$ 各成一個節點 \\
插入情況 2、3：叔叔是黑的，旋轉 & 新鍵放進一個 3-節點，變成 4-節點；旋轉只是把三個
鍵換成「中間黑、兩邊紅」的排法 \\
刪除情況 2：兄弟和它的子節點都是黑的，兄弟改紅 & 兄弟是 2-節點，不能借，合併：
父節點的鍵下來（$p$ 的雙重黑往上傳，就是父節點少了一個鍵） \\
刪除情況 3、4：兄弟有紅色的子節點 & 兄弟是 3- 或 4-節點，向它借一個鍵 \\
刪除情況 1：兄弟是紅的 & 父節點是 3-節點，換一種二元的排法，2-3-4 樹沒有變 \\
\bottomrule
\end{tabular}
\end{center}

要注意，紅黑樹插入時遇到 4-節點，往上送的是原本三個鍵中間的那一個（$g$），新鍵留在
分出去的其中一邊。 這和本節由上而下插入時往上送的鍵相同，但紅黑樹是插入之後才往上
修，所以兩邊畫出來的樹不一定每一步都一樣。

\section{B 樹}
\label{sec:mw-btree}

\subsection{定義}

2-3 樹和 2-3-4 樹只是子節點數上限不同。 把上限寫成參數 $m$，就是 B 樹。

\begin{definition}[$m$ 階 B 樹]
\label{def:mw-btree}
設 $m \ge 3$。 一棵\textbf{$m$ 階 B 樹}（B-tree of order $m$）是一棵 $m$ 路搜尋樹，
滿足：
\begin{enumerate}
\item 每個節點最多 $m$ 個子節點，也就是最多 $m - 1$ 個鍵。
\item 根如果不是葉節點，至少有 $2$ 個子節點。
\item 根以外的內部節點，至少有 $\lceil m/2 \rceil$ 個子節點。
\item 根以外的每個節點（包括葉節點），至少有 $\lceil m/2 \rceil - 1$ 個鍵。
\item 有 $k$ 個鍵的內部節點剛好有 $k + 1$ 個子節點。
\item 所有的葉節點在同一層。
\end{enumerate}
\end{definition}

第 3 條和第 5 條合起來等於第 4 條對內部節點的要求；第 4 條另外規定了葉節點。 根是例外，
只要至少一個鍵。

\begin{center}
\begin{tabular}{ccccc}
\toprule
階數 $m$ & 非根節點的鍵數 & 非根內部節點的子節點數 & 根的鍵數 & 別名 \\
\midrule
$3$ & $1 \sim 2$ & $2 \sim 3$ & $1 \sim 2$ & 2-3 樹 \\
$4$ & $1 \sim 3$ & $2 \sim 4$ & $1 \sim 3$ & 2-3-4 樹 \\
$5$ & $2 \sim 4$ & $3 \sim 5$ & $1 \sim 4$ & \\
$m$ & $\lceil m/2 \rceil - 1 \sim m - 1$ & $\lceil m/2 \rceil \sim m$ & $1 \sim m - 1$ & \\
\bottomrule
\end{tabular}
\end{center}

有些書在葉子底下再畫一層方塊，叫\textbf{外部節點}或\textbf{失敗節點}（failure
node），代表搜尋失敗時會走到的地方，第 6 條就改說「所有外部節點在同一層」。 有 $n$
個鍵的 B 樹剛好有 $n + 1$ 個外部節點（觀念題第 9 題），因為 $n$ 個鍵把數線切成
$n + 1$ 段。 本章的圖不畫外部節點，高度也不算這一層。

\begin{examnote}
「階數」有兩種常見的定義，作答前先看清楚題目用哪一種。 本書（以及多數資料結構課本）
的 $m$ 階是\textbf{最多 $m$ 個子節點}。 CLRS 用\textbf{最小分支度}（minimum degree）
$t$：非根節點有 $t - 1$ 到 $2t - 1$ 個鍵、$t$ 到 $2t$ 個子節點，換成本書的說法就是
$m = 2t$ 階。 所以 CLRS 的 $t = 2$ 是 2-3-4 樹（$4$ 階），而奇數階（例如 $3$ 階、
$5$ 階）在 CLRS 的定義裡寫不出來。 另外，B 樹最早的論文說的「$k$ 階」是每個非根節點
有 $k$ 到 $2k$ 個鍵，對應本書的 $m = 2k + 1$。
\end{examnote}

\subsection{高度}
\label{sec:mw-height}

\begin{theorem}
\label{thm:mw-height}
有 $n \ge 1$ 個鍵、高度 $h$ 的 $m$ 階 B 樹，令 $c = \lceil m/2 \rceil$，則
\[
2c^{h-1} - 1 \;\le\; n \;\le\; m^h - 1,
\qquad\text{也就是}\qquad
\log_m(n + 1) \;\le\; h \;\le\; 1 + \log_c \frac{n + 1}{2}.
\]
\end{theorem}

\begin{proof}
\textbf{上界 $n \le m^h - 1$。} 每個節點最多 $m$ 個子節點，第 $i$ 層最多 $m^{i-1}$
個節點，每個節點最多 $m - 1$ 個鍵，所以
\[
n \le (m - 1)(1 + m + m^2 + \cdots + m^{h-1}) = (m - 1) \cdot \frac{m^h - 1}{m - 1}
= m^h - 1.
\]

\textbf{下界 $n \ge 2c^{h-1} - 1$。} 數每一層最少有幾個節點。 第 $1$ 層是根，
$1$ 個。 $h \ge 2$ 時根至少有 $2$ 個子節點，第 $2$ 層至少 $2$ 個節點；從第 $2$ 層起，
每個節點都是非根節點，第 $h$ 層以外的都是內部節點，各至少有 $c$ 個子節點，所以第
$i$ 層（$2 \le i \le h$）至少有 $2c^{i-2}$ 個節點。 根至少 $1$ 個鍵，其他節點至少
$c - 1$ 個鍵：
\[
n \;\ge\; 1 + (c - 1) \sum_{i=2}^{h} 2c^{i-2}
\;=\; 1 + 2(c - 1) \cdot \frac{c^{h-1} - 1}{c - 1}
\;=\; 2c^{h-1} - 1.
\]
$h = 1$ 時右邊是 $1$，也成立。 兩個不等式分別取對數，得到 $h$ 的範圍。
\end{proof}

下界可以達到：每個非根節點剛好 $c - 1$ 個鍵、根剛好 $1$ 個鍵的樹就是。
圖~\ref{fig:mw-b-ins}(o) 那棵 $5$ 階 B 樹有 $17$ 個鍵、高度 $3$，正好是
$2 \cdot 3^2 - 1 = 17$。 上界也可以達到：每個節點都放滿 $m - 1$ 個鍵。

$m$ 是常數時，兩邊都是 $\Theta(\log n)$。 但底數差很多：$m = 200$ 時，高度下界
$\log_m(n+1)$ 的底數是 $200$，上界 $1 + \log_c \frac{n+1}{2}$ 的底數是 $c = 100$。
$\log_{100} n = \lg n / \lg 100$，所以最壞的高度也只有平衡二元樹的 $1/6.6$ 左右
（$\lg 100 \approx 6.6$）。 B 樹
的好處就在這裡。

\begin{examnote}
本書的高度數節點，CLRS 數邊，CLRS 的高度比本書少 $1$。 例如定理~\ref{thm:mw-height}
的上界，用 CLRS 的高度 $h' = h - 1$ 和 $m = 2t$ 寫，就是 CLRS 的
$h' \le \log_t \frac{n+1}{2}$，同一件事。 計算題要先看題目的高度怎麼定義（「只有根的樹
高度是 $0$ 還是 $1$」），答案差 $1$；題目沒說時，寫出你用的定義再算。
\end{examnote}

考卷上的高度題，用定理的兩個不等式直接算，不要用對數的近似值，邊界很容易差一格。

\textbf{例 1.} $5$ 階 B 樹有 $1000$ 個鍵，高度可能是多少？ 這裡 $c = 3$。 最矮：要
$5^h - 1 \ge 1000$，$5^4 = 625$ 不夠，$5^5 = 3125$ 夠，所以 $h \ge 5$。 最高：要
$2 \cdot 3^{h-1} - 1 \le 1000$，也就是 $3^{h-1} \le 500.5$；$3^5 = 243$ 可以，
$3^6 = 729$ 不行，所以 $h - 1 \le 5$，$h \le 6$。 答案是 $5$ 或 $6$。

\textbf{例 2.} 高度 $4$ 的 $5$ 階 B 樹最少、最多有幾個鍵？ 最少 $2 \cdot 3^3 - 1 = 53$，
最多 $5^4 - 1 = 624$。


\subsection{搜尋}

搜尋和多路搜尋樹一樣。 虛擬碼的節點 $x$ 有 $x.n$ 個鍵
$x.\mathit{key}[1..x.n]$，內部節點有 $x.n + 1$ 個子指標 $x.c[1..x.n+1]$，
$x.\mathit{leaf}$ 表示是不是葉子。

\begin{pseudo}[label=alg:mw-search, unbreakable]{在以 $x$ 為根的子樹裡找 $k$，回傳節點和位置}
\Function{B-Tree-Search}{$x, k$}
    \State $i \gets 1$
    \While{$i \le x.n$ \AND $k > x.\mathit{key}[i]$}
        \State $i \gets i + 1$
    \EndWhile
    \If{$i \le x.n$ \AND $k = x.\mathit{key}[i]$}
        \State \Return $(x, i)$
    \EndIf
    \If{$x.\mathit{leaf}$}
        \State \Return \NIL
    \EndIf
    \State 從磁碟讀入節點 $x.c[i]$
    \State \Return \Call{B-Tree-Search}{$x.c[i], k$}
\EndFunction
\end{pseudo}

走過的節點最多 $h$ 個，第 9 行讀磁碟最多 $h = O(\log_c n)$ 次。 第 3 行在節點裡循序找，
CPU 時間 $O(m \log_c n)$；改用二分搜尋是 $O(\lg m \cdot \log_c n) = O(\lg n)$。 磁碟上
的 B 樹，瓶頸是讀取次數。

\subsection{插入}
\label{sec:mw-insert}

插入照 2-3 樹的做法：放進葉子，溢位就分裂，一路往上。 節點溢位時有 $m$ 個鍵
$k_1 < k_2 < \cdots < k_m$，令 $p = \lceil m/2 \rceil$：
\begin{itemize}
\item $k_1, \ldots, k_{p-1}$ 留在原節點，共 $p - 1 = \lceil m/2 \rceil - 1$ 個；
\item $k_p$ 往上送進父節點；
\item $k_{p+1}, \ldots, k_m$ 放進新節點，共 $m - p = \lfloor m/2 \rfloor$ 個。
\end{itemize}
兩邊都至少有 $\lceil m/2 \rceil - 1$ 個鍵，符合定義的第 4 條（觀念題第 8 題）。 $m$ 是
奇數時 $k_p$ 正好在中間；$m$ 是偶數時左邊比右邊少一個。

虛擬碼用遞迴寫：\proc{Insert-Down} 把 $k$ 插進 $x$ 的子樹，如果 $x$ 因此分裂，就把
「上升的鍵」和「分出來的右半節點」交回給呼叫它的父節點。 節點的陣列要多留一格，讓它
可以暫時放 $m$ 個鍵。 下面依序是最上層的 \proc{B-Tree-Insert}、分裂一個節點的
\proc{Split}，最後是 \proc{Insert-Down}。

\begin{pseudo}[label=alg:mw-insert]{B 樹的插入}
\Procedure{B-Tree-Insert}{$T, k$}
    \If{$T.\mathit{root} = \NIL$}
        \State $T.\mathit{root} \gets$ 只放 $k$ 的新葉節點
        \State \Return
    \EndIf
    \State $(u, r) \gets$ \Call{Insert-Down}{$T.\mathit{root}, k$}
    \If{$r \ne \NIL$} \Comment{根分裂了}
        \State $s \gets$ 新節點，$s.n \gets 1$，$s.\mathit{key}[1] \gets u$
        \State $s.c[1] \gets T.\mathit{root}$，$s.c[2] \gets r$
        \State $T.\mathit{root} \gets s$
    \EndIf
\EndProcedure
\end{pseudo}

\begin{pseudo}[label=alg:mw-split]{分裂有 $m$ 個鍵的節點 $x$}
\Function{Split}{$x$}
    \State $p \gets \lceil m/2 \rceil$
    \State $r \gets$ 新節點，$r.\mathit{leaf} \gets x.\mathit{leaf}$
    \State $r.\mathit{key}[1..m-p] \gets x.\mathit{key}[p+1..m]$
    \If{$x$ 不是葉子}
        \State $r.c[1..m-p+1] \gets x.c[p+1..m+1]$
    \EndIf
    \State $r.n \gets m - p$，$x.n \gets p - 1$
    \State \Return $(x.\mathit{key}[p], r)$
\EndFunction
\end{pseudo}

\begin{pseudo}[label=alg:mw-insdown, unbreakable]{把 $k$ 插進 $x$ 的子樹；$x$ 分裂時回傳上升的鍵和新的右半節點}
\Function{Insert-Down}{$x, k$}
    \State $i \gets 1$
    \While{$i \le x.n$ \AND $k > x.\mathit{key}[i]$}
        \State $i \gets i + 1$
    \EndWhile
    \If{$x.\mathit{leaf}$}
        \State $x.\mathit{key}[i..x.n]$ 往後移一格，$x.\mathit{key}[i] \gets k$
        \State $x.n \gets x.n + 1$
    \Else
        \State $(u, r) \gets$ \Call{Insert-Down}{$x.c[i], k$}
        \If{$r = \NIL$}
            \State \Return $(\NIL, \NIL)$
        \EndIf
        \State $x.\mathit{key}[i..x.n]$、$x.c[i+1..x.n+1]$ 往後移一格
        \State $x.\mathit{key}[i] \gets u$，$x.c[i+1] \gets r$，$x.n \gets x.n + 1$
    \EndIf
    \If{$x.n < m$}
        \State \Return $(\NIL, \NIL)$
    \EndIf
    \State \Return \Call{Split}{$x$}
\EndFunction
\end{pseudo}

\proc{Insert-Down} 第 3 行找到 $k$ 該去的位置；葉節點在第 6 行放進去，內部節點在第
9 行往下遞迴，下面分裂時在第 12、13 行把上升的鍵 $u$ 和右半節點 $r$ 接進來。 第 14
行檢查溢位，溢位就在第 16 行分裂。 \proc{Split} 的第 5、6 行：內部節點分裂時，右邊的 $m - p$
個鍵帶著周圍的 $m - p + 1$ 個子指標搬走，左邊留下 $p - 1$ 個鍵和 $p$ 個子
指標，子指標合計 $m + 1$ 個。

\textbf{成本。} 往下走 $h$ 層，回程每層最多分裂一次，每次分裂搬 $O(m)$ 個鍵和指標，
讀寫磁碟 $O(h)$ 次，CPU 時間 $O(mh) = O(m \log_c n)$。 分裂其實很少：每次分裂多出一
個節點，而每個非根節點至少有 $c - 1$ 個鍵，所以從空樹插入 $n$ 個鍵，總共的分裂次數
不超過節點數，約 $n / (c - 1)$。 用第~3~章的說法，每次插入平攤只有 $O(1/m)$ 次分裂。

\subsection{插入的完整追蹤}
\label{sec:mw-b-trace}

把 $40, 57, 65, 56, 29, 12, 53, 23, 50, 14, 28, 16, 81, 48, 96, 72, 47$ 依序插入空
的 $5$ 階 B 樹。 每個節點最多 $4$ 個鍵，放到第 $5$ 個就溢位，分裂時第 $3$ 個鍵
（$p = 3$）上升，左右各留 $2$ 個。

\begin{figure}[tbp]
\centering
\begin{tikzpicture}[baseline=(current bounding box.north), mwk/.append style={minimum width=0.46cm, minimum height=0.42cm}]
\node[mwk,mwkN] at (0.23,0) {40};
\node[mwk,mwkN] at (0.69,0) {56};
\node[mwk,mwkN] at (1.15,0) {57};
\node[mwk,mwkN] at (1.61,0) {65};
\node[mwlab] at (0.92,-1.28) {(a) 插入 40、57、65、56};
\end{tikzpicture}\hspace{0.55cm}\begin{tikzpicture}[baseline=(current bounding box.north), mwk/.append style={minimum width=0.46cm, minimum height=0.42cm}]
\node[mwk,mwkN] at (0.23,0) {29};
\node[mwk,mwkB] at (0.69,0) {40};
\node[mwk,mwkB] at (1.15,0) {56};
\node[mwk,mwkB] at (1.61,0) {57};
\node[mwk,mwkB] at (2.07,0) {65};
\draw[mwbad] (0,-0.21) rectangle (2.3,0.21);
\node[mwlab] at (1.15,-1.28) {(b) 插入 29：溢位};
\end{tikzpicture}\hspace{0.55cm}\begin{tikzpicture}[baseline=(current bounding box.north), mwk/.append style={minimum width=0.46cm, minimum height=0.42cm}]
\draw[mwe] (0.84,-0.21) -- (0.46,-0.59);
\draw[mwe] (1.3,-0.21) -- (1.68,-0.59);
\node[mwk] at (0.23,-0.8) {29};
\node[mwk] at (0.69,-0.8) {40};
\node[mwk] at (1.45,-0.8) {57};
\node[mwk] at (1.91,-0.8) {65};
\node[mwk,mwkM] at (1.07,0) {56};
\node[mwlab] at (1.07,-1.28) {(c) 分裂，56 上升成新根};
\end{tikzpicture}\\*[0.45em]
\begin{tikzpicture}[baseline=(current bounding box.north), mwk/.append style={minimum width=0.46cm, minimum height=0.42cm}]
\draw[mwe] (1.53,-0.21) -- (0.92,-0.59);
\draw[mwe] (1.99,-0.21) -- (2.6,-0.59);
\node[mwk,mwkN] at (0.23,-0.8) {12};
\node[mwk] at (0.69,-0.8) {29};
\node[mwk] at (1.15,-0.8) {40};
\node[mwk,mwkN] at (1.61,-0.8) {53};
\node[mwk] at (2.37,-0.8) {57};
\node[mwk] at (2.83,-0.8) {65};
\node[mwk] at (1.76,0) {56};
\node[mwlab] at (1.53,-1.28) {(d) 插入 12、53};
\end{tikzpicture}\hspace{0.55cm}\begin{tikzpicture}[baseline=(current bounding box.north), mwk/.append style={minimum width=0.46cm, minimum height=0.42cm}]
\draw[mwe] (1.875,-0.21) -- (1.15,-0.59);
\draw[mwe] (2.335,-0.21) -- (3.06,-0.59);
\node[mwk,mwkB] at (0.23,-0.8) {12};
\node[mwk,mwkN] at (0.69,-0.8) {23};
\node[mwk,mwkB] at (1.15,-0.8) {29};
\node[mwk,mwkB] at (1.61,-0.8) {40};
\node[mwk,mwkB] at (2.07,-0.8) {53};
\draw[mwbad] (0,-1.01) rectangle (2.3,-0.59);
\node[mwk] at (2.83,-0.8) {57};
\node[mwk] at (3.29,-0.8) {65};
\node[mwk] at (2.105,0) {56};
\node[mwlab] at (1.76,-1.28) {(e) 插入 23：溢位};
\end{tikzpicture}\hspace{0.55cm}\begin{tikzpicture}[baseline=(current bounding box.north), mwk/.append style={minimum width=0.46cm, minimum height=0.42cm}]
\draw[mwe] (1.22,-0.21) -- (0.46,-0.59);
\draw[mwe] (1.68,-0.21) -- (1.68,-0.59);
\draw[mwe] (2.14,-0.21) -- (2.9,-0.59);
\node[mwk] at (0.23,-0.8) {12};
\node[mwk] at (0.69,-0.8) {23};
\node[mwk] at (1.45,-0.8) {40};
\node[mwk] at (1.91,-0.8) {53};
\node[mwk] at (2.67,-0.8) {57};
\node[mwk] at (3.13,-0.8) {65};
\node[mwk,mwkM] at (1.45,0) {29};
\node[mwk] at (1.91,0) {56};
\node[mwlab] at (1.68,-1.28) {(f) 分裂，29 上升};
\end{tikzpicture}\\*[0.45em]
\begin{tikzpicture}[baseline=(current bounding box.north), mwk/.append style={minimum width=0.46cm, minimum height=0.42cm}]
\draw[mwe] (2.14,-0.21) -- (0.92,-0.59);
\draw[mwe] (2.6,-0.21) -- (2.83,-0.59);
\draw[mwe] (3.06,-0.21) -- (4.28,-0.59);
\node[mwk] at (0.23,-0.8) {12};
\node[mwk,mwkN] at (0.69,-0.8) {14};
\node[mwk] at (1.15,-0.8) {23};
\node[mwk,mwkN] at (1.61,-0.8) {28};
\node[mwk] at (2.37,-0.8) {40};
\node[mwk,mwkN] at (2.83,-0.8) {50};
\node[mwk] at (3.29,-0.8) {53};
\node[mwk] at (4.05,-0.8) {57};
\node[mwk] at (4.51,-0.8) {65};
\node[mwk] at (2.37,0) {29};
\node[mwk] at (2.83,0) {56};
\node[mwlab] at (2.37,-1.28) {(g) 插入 50、14、28};
\end{tikzpicture}\hspace{0.55cm}\begin{tikzpicture}[baseline=(current bounding box.north), mwk/.append style={minimum width=0.46cm, minimum height=0.42cm}]
\draw[mwe] (2.485,-0.21) -- (1.15,-0.59);
\draw[mwe] (2.945,-0.21) -- (3.29,-0.59);
\draw[mwe] (3.405,-0.21) -- (4.74,-0.59);
\node[mwk,mwkB] at (0.23,-0.8) {12};
\node[mwk,mwkB] at (0.69,-0.8) {14};
\node[mwk,mwkN] at (1.15,-0.8) {16};
\node[mwk,mwkB] at (1.61,-0.8) {23};
\node[mwk,mwkB] at (2.07,-0.8) {28};
\draw[mwbad] (0,-1.01) rectangle (2.3,-0.59);
\node[mwk] at (2.83,-0.8) {40};
\node[mwk] at (3.29,-0.8) {50};
\node[mwk] at (3.75,-0.8) {53};
\node[mwk] at (4.51,-0.8) {57};
\node[mwk] at (4.97,-0.8) {65};
\node[mwk] at (2.715,0) {29};
\node[mwk] at (3.175,0) {56};
\node[mwlab] at (2.6,-1.28) {(h) 插入 16：溢位};
\end{tikzpicture}\\*[0.45em]
\begin{tikzpicture}[baseline=(current bounding box.north), mwk/.append style={minimum width=0.46cm, minimum height=0.42cm}]
\draw[mwe] (1.83,-0.21) -- (0.46,-0.59);
\draw[mwe] (2.29,-0.21) -- (1.68,-0.59);
\draw[mwe] (2.75,-0.21) -- (3.13,-0.59);
\draw[mwe] (3.21,-0.21) -- (4.58,-0.59);
\node[mwk] at (0.23,-0.8) {12};
\node[mwk] at (0.69,-0.8) {14};
\node[mwk] at (1.45,-0.8) {23};
\node[mwk] at (1.91,-0.8) {28};
\node[mwk] at (2.67,-0.8) {40};
\node[mwk] at (3.13,-0.8) {50};
\node[mwk] at (3.59,-0.8) {53};
\node[mwk] at (4.35,-0.8) {57};
\node[mwk] at (4.81,-0.8) {65};
\node[mwk,mwkM] at (2.06,0) {16};
\node[mwk] at (2.52,0) {29};
\node[mwk] at (2.98,0) {56};
\node[mwlab] at (2.52,-1.28) {(i) 分裂，16 上升};
\end{tikzpicture}\hspace{0.55cm}\begin{tikzpicture}[baseline=(current bounding box.north), mwk/.append style={minimum width=0.46cm, minimum height=0.42cm}]
\draw[mwe] (2.29,-0.21) -- (0.46,-0.59);
\draw[mwe] (2.75,-0.21) -- (1.68,-0.59);
\draw[mwe] (3.21,-0.21) -- (3.36,-0.59);
\draw[mwe] (3.67,-0.21) -- (5.5,-0.59);
\node[mwk] at (0.23,-0.8) {12};
\node[mwk] at (0.69,-0.8) {14};
\node[mwk] at (1.45,-0.8) {23};
\node[mwk] at (1.91,-0.8) {28};
\node[mwk] at (2.67,-0.8) {40};
\node[mwk,mwkN] at (3.13,-0.8) {48};
\node[mwk] at (3.59,-0.8) {50};
\node[mwk] at (4.05,-0.8) {53};
\node[mwk] at (4.81,-0.8) {57};
\node[mwk] at (5.27,-0.8) {65};
\node[mwk,mwkN] at (5.73,-0.8) {81};
\node[mwk,mwkN] at (6.19,-0.8) {96};
\node[mwk] at (2.52,0) {16};
\node[mwk] at (2.98,0) {29};
\node[mwk] at (3.44,0) {56};
\node[mwlab] at (3.21,-1.28) {(j) 插入 81、48、96};
\end{tikzpicture}\\*[0.45em]
\begin{tikzpicture}[baseline=(current bounding box.north), mwk/.append style={minimum width=0.46cm, minimum height=0.42cm}]
\draw[mwe] (2.405,-0.21) -- (0.46,-0.59);
\draw[mwe] (2.865,-0.21) -- (1.68,-0.59);
\draw[mwe] (3.325,-0.21) -- (3.36,-0.59);
\draw[mwe] (3.785,-0.21) -- (5.73,-0.59);
\node[mwk] at (0.23,-0.8) {12};
\node[mwk] at (0.69,-0.8) {14};
\node[mwk] at (1.45,-0.8) {23};
\node[mwk] at (1.91,-0.8) {28};
\node[mwk] at (2.67,-0.8) {40};
\node[mwk] at (3.13,-0.8) {48};
\node[mwk] at (3.59,-0.8) {50};
\node[mwk] at (4.05,-0.8) {53};
\node[mwk,mwkB] at (4.81,-0.8) {57};
\node[mwk,mwkB] at (5.27,-0.8) {65};
\node[mwk,mwkN] at (5.73,-0.8) {72};
\node[mwk,mwkB] at (6.19,-0.8) {81};
\node[mwk,mwkB] at (6.65,-0.8) {96};
\draw[mwbad] (4.58,-1.01) rectangle (6.88,-0.59);
\node[mwk] at (2.635,0) {16};
\node[mwk] at (3.095,0) {29};
\node[mwk] at (3.555,0) {56};
\node[mwlab] at (3.44,-1.28) {(k) 插入 72：溢位};
\end{tikzpicture}\hspace{0.55cm}\begin{tikzpicture}[baseline=(current bounding box.north), mwk/.append style={minimum width=0.46cm, minimum height=0.42cm}]
\draw[mwe] (2.44,-0.21) -- (0.46,-0.59);
\draw[mwe] (2.9,-0.21) -- (1.68,-0.59);
\draw[mwe] (3.36,-0.21) -- (3.36,-0.59);
\draw[mwe] (3.82,-0.21) -- (5.04,-0.59);
\draw[mwe] (4.28,-0.21) -- (6.26,-0.59);
\node[mwk] at (0.23,-0.8) {12};
\node[mwk] at (0.69,-0.8) {14};
\node[mwk] at (1.45,-0.8) {23};
\node[mwk] at (1.91,-0.8) {28};
\node[mwk] at (2.67,-0.8) {40};
\node[mwk] at (3.13,-0.8) {48};
\node[mwk] at (3.59,-0.8) {50};
\node[mwk] at (4.05,-0.8) {53};
\node[mwk] at (4.81,-0.8) {57};
\node[mwk] at (5.27,-0.8) {65};
\node[mwk] at (6.03,-0.8) {81};
\node[mwk] at (6.49,-0.8) {96};
\node[mwk] at (2.67,0) {16};
\node[mwk] at (3.13,0) {29};
\node[mwk] at (3.59,0) {56};
\node[mwk,mwkM] at (4.05,0) {72};
\node[mwlab] at (3.36,-1.28) {(l) 分裂，72 上升};
\end{tikzpicture}\\*[0.45em]
\begin{tikzpicture}[baseline=(current bounding box.north), mwk/.append style={minimum width=0.46cm, minimum height=0.42cm}]
\draw[mwe] (2.67,-0.21) -- (0.46,-0.59);
\draw[mwe] (3.13,-0.21) -- (1.68,-0.59);
\draw[mwe] (3.59,-0.21) -- (3.59,-0.59);
\draw[mwe] (4.05,-0.21) -- (5.5,-0.59);
\draw[mwe] (4.51,-0.21) -- (6.72,-0.59);
\node[mwk] at (0.23,-0.8) {12};
\node[mwk] at (0.69,-0.8) {14};
\node[mwk] at (1.45,-0.8) {23};
\node[mwk] at (1.91,-0.8) {28};
\node[mwk,mwkB] at (2.67,-0.8) {40};
\node[mwk,mwkN] at (3.13,-0.8) {47};
\node[mwk,mwkB] at (3.59,-0.8) {48};
\node[mwk,mwkB] at (4.05,-0.8) {50};
\node[mwk,mwkB] at (4.51,-0.8) {53};
\draw[mwbad] (2.44,-1.01) rectangle (4.74,-0.59);
\node[mwk] at (5.27,-0.8) {57};
\node[mwk] at (5.73,-0.8) {65};
\node[mwk] at (6.49,-0.8) {81};
\node[mwk] at (6.95,-0.8) {96};
\node[mwk] at (2.9,0) {16};
\node[mwk] at (3.36,0) {29};
\node[mwk] at (3.82,0) {56};
\node[mwk] at (4.28,0) {72};
\node[mwlab] at (3.59,-1.28) {(m) 插入 47：溢位};
\end{tikzpicture}\hspace{0.55cm}\begin{tikzpicture}[baseline=(current bounding box.north), mwk/.append style={minimum width=0.46cm, minimum height=0.42cm}]
\draw[mwe] (2.36,-0.21) -- (0.46,-0.59);
\draw[mwe] (2.82,-0.21) -- (1.68,-0.59);
\draw[mwe] (3.28,-0.21) -- (2.9,-0.59);
\draw[mwe] (3.74,-0.21) -- (4.12,-0.59);
\draw[mwe] (4.2,-0.21) -- (5.34,-0.59);
\draw[mwe] (4.66,-0.21) -- (6.56,-0.59);
\node[mwk] at (0.23,-0.8) {12};
\node[mwk] at (0.69,-0.8) {14};
\node[mwk] at (1.45,-0.8) {23};
\node[mwk] at (1.91,-0.8) {28};
\node[mwk] at (2.67,-0.8) {40};
\node[mwk] at (3.13,-0.8) {47};
\node[mwk] at (3.89,-0.8) {50};
\node[mwk] at (4.35,-0.8) {53};
\node[mwk] at (5.11,-0.8) {57};
\node[mwk] at (5.57,-0.8) {65};
\node[mwk] at (6.33,-0.8) {81};
\node[mwk] at (6.79,-0.8) {96};
\node[mwk,mwkB] at (2.59,0) {16};
\node[mwk,mwkB] at (3.05,0) {29};
\node[mwk,mwkM] at (3.51,0) {48};
\node[mwk,mwkB] at (3.97,0) {56};
\node[mwk,mwkB] at (4.43,0) {72};
\draw[mwbad] (2.36,-0.21) rectangle (4.66,0.21);
\node[mwlab] at (3.51,-1.28) {(n) 分裂，48 上升，父節點溢位};
\end{tikzpicture}\\*[0.45em]
\begin{tikzpicture}[baseline=(current bounding box.north), mwk/.append style={minimum width=0.46cm, minimum height=0.42cm}]
\draw[mwe] (1.22,-1.01) -- (0.46,-1.39);
\draw[mwe] (1.68,-1.01) -- (1.68,-1.39);
\draw[mwe] (2.14,-1.01) -- (2.9,-1.39);
\draw[mwe] (4.88,-1.01) -- (4.12,-1.39);
\draw[mwe] (5.34,-1.01) -- (5.34,-1.39);
\draw[mwe] (5.8,-1.01) -- (6.56,-1.39);
\draw[mwe] (3.28,-0.21) -- (1.68,-0.59);
\draw[mwe] (3.74,-0.21) -- (5.34,-0.59);
\node[mwk] at (0.23,-1.6) {12};
\node[mwk] at (0.69,-1.6) {14};
\node[mwk] at (1.45,-1.6) {23};
\node[mwk] at (1.91,-1.6) {28};
\node[mwk] at (2.67,-1.6) {40};
\node[mwk] at (3.13,-1.6) {47};
\node[mwk] at (1.45,-0.8) {16};
\node[mwk] at (1.91,-0.8) {29};
\node[mwk] at (3.89,-1.6) {50};
\node[mwk] at (4.35,-1.6) {53};
\node[mwk] at (5.11,-1.6) {57};
\node[mwk] at (5.57,-1.6) {65};
\node[mwk] at (6.33,-1.6) {81};
\node[mwk] at (6.79,-1.6) {96};
\node[mwk] at (5.11,-0.8) {56};
\node[mwk] at (5.57,-0.8) {72};
\node[mwk,mwkM] at (3.51,0) {48};
\node[mwlab] at (3.51,-2.08) {(o) 根分裂，48 上升成新根};
\end{tikzpicture}
\diagcaption{$5$ 階 B 樹的插入。 每次溢位畫一張，分裂之後再畫一張；沒有分裂的插入合
併成一張。 (m) 到 (o) 是連鎖分裂：葉子分裂後 $48$ 上升，根變成 $5$ 個鍵，再分裂一
次，$48$ 又上升成為新的根。}
\label{fig:mw-b-ins}
\end{figure}

有分裂的插入是 $29, 23, 16, 72, 47$ 這五次，一共分裂六次：插入 $47$ 是連鎖分裂，兩次
上升的都是 $48$，它先從葉子升到根，剛好又是根裡的第 $3$ 個鍵。

\begin{examnote}
B 樹的追蹤題，每次分裂（或借、合併）都要畫出整棵樹，並寫出哪個節點溢位或不足、哪個
鍵上升或下降。 題目沒指定的規則，在答案開頭寫一句你用的約定：偶數階溢位時上升第
$m/2$ 個還是第 $m/2 + 1$ 個鍵；刪除內部節點的鍵用後繼還是前驅；先向哪一邊的兄弟借、
和哪一邊合併。 畫完檢查三件事：鍵數都在範圍內、葉子都同層、由左到
右讀出來是排好的。
\end{examnote}

\subsection{C 程式}
\label{sec:mw-c}

C 程式照虛擬碼~\ref{alg:mw-insert} 到 \ref{alg:mw-insdown} 寫，階數用 \texttt{M} 固定成
$5$，陣列從 $0$ 開始。 \texttt{key} 開 \texttt{M} 格，比上限多一格，給溢位時暫放。
\texttt{insert\_down} 回傳分裂出來的右半節點（沒有分裂時是 \texttt{NULL}），上升的
鍵透過指標參數 \texttt{up} 帶回；遞迴呼叫時把 \texttt{\&k} 當成 \texttt{up} 傳下去，
子節點分裂時 \texttt{k} 就換成上升的鍵，接下來插進 \texttt{x} 的正是它。
\texttt{main} 跑第~\ref{sec:mw-b-trace}~節的例子，印出有分裂的插入和最後每一層的節
點，再找三個鍵。

\begin{cbookrecap}[第~14~章 14.2、14.5 節，第~19~章 19.3 節]
結構裡可以放陣列，例如 \texttt{int key[M]}，整個陣列跟著結構一起配置，
\texttt{malloc(sizeof *x)} 一次就要到整個節點。 \texttt{struct Node *child[M + 1]}
是「指標的陣列」，每一格指向一個子節點。 \texttt{x->key[i]} 是 \texttt{(*x).key[i]}
的簡寫。 釋放一棵樹要先釋放所有子節點，最後才釋放自己。
\end{cbookrecap}

\begin{example}[label=ex:mw-btree]{c}{(btree\_insert.c) B 樹的搜尋與插入}
#include <stdio.h>
#include <stdlib.h>

#define M 5                         // 階數：每個節點最多 M 個子節點

typedef struct Node {
    int n, leaf;                    // 鍵的個數；是不是葉子
    int key[M];                     // 最多 M-1 個，溢位時暫時 M 個
    struct Node *child[M + 1];
} Node;

static int splits;                  // 統計分裂的次數

static Node *new_node(int leaf)
{
    Node *x = malloc(sizeof *x);
    if (x == NULL) {
        exit(1);
    }
    x->n = 0;
    x->leaf = leaf;
    return x;
}

static int find_pos(const Node *x, int k)      // 第一個 >= k 的位置
{
    int i = 0;
    while (i < x->n && k > x->key[i]) {
        i++;
    }
    return i;
}

// 回傳 k 所在的節點（找不到是 NULL），*visits 是讀了幾個節點。
Node *bt_search(Node *x, int k, int *visits)
{
    *visits = 0;
    while (x != NULL) {
        (*visits)++;
        int i = find_pos(x, k);
        if (i < x->n && x->key[i] == k) {
            return x;
        }
        x = x->leaf ? NULL : x->child[i];
    }
    return NULL;
}

// x 有 M 個鍵：前 p-1 個留下，第 p 個上升（放進 *up），其餘搬到新節點。
static Node *split(Node *x, int *up)
{
    int p = (M + 1) / 2;            // ceil(M/2)
    Node *r = new_node(x->leaf);
    r->n = M - p;
    for (int j = 0; j < r->n; j++) {
        r->key[j] = x->key[p + j];
    }
    for (int j = 0; !x->leaf && j <= r->n; j++) {
        r->child[j] = x->child[p + j];
    }
    *up = x->key[p - 1];
    x->n = p - 1;
    splits++;
    return r;
}

// 把 k 插進 x 的子樹。x 分裂時回傳右半節點，否則回傳 NULL。
static Node *insert_down(Node *x, int k, int *up)
{
    int i = find_pos(x, k);
    if (i < x->n && x->key[i] == k) {
        return NULL;                // 已經有了，不插
    }
    Node *r = NULL;
    if (!x->leaf) {
        r = insert_down(x->child[i], k, &k);   // 子節點分裂時，k 換成上升的鍵
        if (r == NULL) {
            return NULL;
        }
    }
    for (int j = x->n; j > i; j--) {            // 騰出位置 i
        x->key[j] = x->key[j - 1];
        if (!x->leaf) {
            x->child[j + 1] = x->child[j];
        }
    }
    x->key[i] = k;
    if (!x->leaf) {
        x->child[i + 1] = r;
    }
    x->n++;
    return x->n < M ? NULL : split(x, up);
}

Node *bt_insert(Node *root, int k)  // 回傳新的根
{
    if (root == NULL) {
        root = new_node(1);
        root->key[0] = k;
        root->n = 1;
        return root;
    }
    int up;
    Node *r = insert_down(root, k, &up);
    if (r == NULL) {
        return root;
    }
    Node *s = new_node(0);          // 根分裂，長出新的根
    s->n = 1;
    s->key[0] = up;
    s->child[0] = root;
    s->child[1] = r;
    return s;
}

int height(const Node *x)           // 數節點：只有根時是 1
{
    return x == NULL ? 0 : 1 + (x->leaf ? 0 : height(x->child[0]));
}

static void print_level(const Node *x, int level)   // 根是第 1 層
{
    if (level > 1) {
        for (int j = 0; j <= x->n; j++) {
            print_level(x->child[j], level - 1);
        }
        return;
    }
    printf(" [");
    for (int j = 0; j < x->n; j++) {
        printf(j == 0 ? "%d" : " %d", x->key[j]);
    }
    putchar(']');
}

void free_tree(Node *x)
{
    for (int j = 0; !x->leaf && j <= x->n; j++) {
        free_tree(x->child[j]);
    }
    free(x);
}

int main(void)
{
    int keys[] = {40, 57, 65, 56, 29, 12, 53, 23, 50,
                  14, 28, 16, 81, 48, 96, 72, 47};
    Node *root = NULL;
    for (int i = 0; i < 17; i++) {
        int before = splits;
        root = bt_insert(root, keys[i]);
        if (splits > before) {
            printf("insert %d: %d split(s), height %d\n",
                   keys[i], splits - before, height(root));
        }
    }
    for (int lv = 1; lv <= height(root); lv++) {
        printf("level %d:", lv);
        print_level(root, lv);
        putchar('\n');
    }
    int probe[] = {48, 50, 44};
    for (int i = 0; i < 3; i++) {
        int v;
        Node *x = bt_search(root, probe[i], &v);
        printf("search %d: %s, %d node(s) read\n",
               probe[i], x != NULL ? "found" : "not found", v);
    }
    free_tree(root);
    return 0;
}
\end{example}
\iosample{}{insert 29: 1 split(s), height 2 \\ insert 23: 1 split(s), height 2 \\ insert 16: 1 split(s), height 2 \\ insert 72: 1 split(s), height 2 \\ insert 47: 2 split(s), height 3 \\ level 1: [48] \\ level 2: [16 29] [56 72] \\ level 3: [12 14] [23 28] [40 47] [50 53] [57 65] [81 96] \\ search 48: found, 1 node(s) read \\ search 50: found, 3 node(s) read \\ search 44: not found, 3 node(s) read}

輸出和圖~\ref{fig:mw-b-ins} 一致：五次有分裂的插入，最後一次分裂兩次。 找 $48$ 只讀
了根；找 $50$ 和找不到的 $44$ 都走到葉子，讀了 $3$ 個節點，等於樹高。
\texttt{split} 裡 \texttt{p = (M + 1) / 2} 就是 $\lceil m/2 \rceil$。 左半留下
\texttt{key[0..p-2]}，上升的是 \texttt{key[p-1]}，右半從 \texttt{key[p]} 開始；0 起
算和虛擬碼的 1 起算差一格，這裡最容易寫錯。 把 \texttt{M} 改成 $3$ 或 $4$，同一支程
式就是 2-3 樹或（由下而上的）2-3-4 樹。

\subsection{刪除}
\label{sec:mw-delete}

刪除照 2-3 樹的做法，只是「不足」的門檻從 $0$ 個鍵變成少於 $\lceil m/2 \rceil - 1$
個。 令 $\ell = \lceil m/2 \rceil - 1$，是非根節點最少的鍵數。
\begin{enumerate}
\item 要刪的鍵在內部節點：換成後繼，改刪葉子裡的後繼。
\item 從葉子拿掉鍵。 剩下的鍵 $\ge \ell$，或這片葉子就是根，結束。
\item 節點不足（剩 $\ell - 1$ 個鍵）時：右兄弟有多於 $\ell$ 個鍵，向它借；否則左兄弟有多於
$\ell$ 個鍵，向它借；都沒有就合併，有右兄弟和右兄弟合併，否則和左兄弟合併。
\item 合併讓父節點少一個鍵；父節點不足時回到第 3 步處理父節點。 根變成 $0$ 個鍵時，拿
掉它，高度減 $1$。
\end{enumerate}

\proc{B-Tree-Delete}$(T, k)$ 呼叫 \proc{Delete-Down}$(T.\mathit{root}, k)$；回來之後如果根
沒有鍵了，根是葉子就讓樹變空，否則把根換成它唯一的子節點 $T.\mathit{root}.c[1]$，高度
減 $1$。 不足的情況在父節點那一層處理，所以 \proc{Delete-Down} 處理的是 $x$ 的子節點。

\begin{pseudo}[label=alg:mw-deldown]{從 $x$ 的子樹刪除 $k$；$x$ 自己不足時交給父節點處理}
\Procedure{Delete-Down}{$x, k$}
    \State $i \gets 1$
    \While{$i \le x.n$ \AND $k > x.\mathit{key}[i]$}
        \State $i \gets i + 1$
    \EndWhile
    \If{$x.\mathit{leaf}$}
        \If{$i \le x.n$ \AND $k = x.\mathit{key}[i]$}
            \State 刪掉 $x.\mathit{key}[i]$，後面的鍵往前移一格，$x.n \gets x.n - 1$
        \EndIf
        \State \Return
    \EndIf
    \If{$i \le x.n$ \AND $k = x.\mathit{key}[i]$} \Comment{$k$ 在內部節點}
        \State $y \gets x.c[i+1]$，再一路往最左邊的子節點走到葉節點
        \State $x.\mathit{key}[i] \gets y.\mathit{key}[1]$ \Comment{換成後繼}
        \State $k \gets y.\mathit{key}[1]$
        \State $i \gets i + 1$ \Comment{改到右子樹刪後繼}
    \EndIf
    \State \Call{Delete-Down}{$x.c[i], k$}
    \If{$x.c[i].n < \lceil m/2 \rceil - 1$}
        \State \Call{Fix-Child}{$x, i$}
    \EndIf
\EndProcedure
\end{pseudo}

\begin{pseudo}[label=alg:mw-fix, unbreakable]{$x.c[i]$ 的鍵太少：借或合併}
\Procedure{Fix-Child}{$x, i$}
    \State $\ell \gets \lceil m/2 \rceil - 1$
    \If{$i \le x.n$ \AND $x.c[i+1].n > \ell$} \Comment{向右兄弟借}
        \State $x.\mathit{key}[i]$ 移到 $x.c[i]$ 的最後面
        \State $x.c[i+1]$ 最小的鍵移上來，當新的 $x.\mathit{key}[i]$
        \State $x.c[i+1]$ 最左的子指標（如果有）移到 $x.c[i]$ 的最右邊
    \ElsIf{$i > 1$ \AND $x.c[i-1].n > \ell$} \Comment{向左兄弟借}
        \State 第 4--6 行左右對調：$x.\mathit{key}[i-1]$ 下來，$x.c[i-1]$ 最大的鍵上去
    \Else \Comment{合併}
        \If{$i > x.n$}
            \State $i \gets i - 1$ \Comment{沒有右兄弟，改和左兄弟合併}
        \EndIf
        \State $x.\mathit{key}[i]$ 和 $x.c[i+1]$ 的鍵、子指標依序接到 $x.c[i]$ 後面
        \State 從 $x$ 刪掉 $x.\mathit{key}[i]$ 和 $x.c[i+1]$，$x.n \gets x.n - 1$
        \State 釋放原本的 $x.c[i+1]$
    \EndIf
\EndProcedure
\end{pseudo}

\proc{Delete-Down} 第 11 行換成後繼之後，第 14 行到右子樹刪掉它；後繼是右子樹最小的
鍵，往下的每一層都走最左邊，最後在第 7 行從葉節點拿掉。 每一層遞迴回來時，第 15 行檢
查剛剛處理的子節點有沒有不足。 \proc{Fix-Child} 第 10、11 行處理最右邊的子節點：它
沒有右兄弟，把 $i$ 減 $1$ 之後，合併的是左兄弟 $x.c[i]$ 和它自己 $x.c[i+1]$。

\begin{whybox}[借和合併之後仍然是 B 樹]
不足的節點有 $\ell - 1$ 個鍵（$\ell = \lceil m/2 \rceil - 1$）。 \textbf{借}：兄弟原本多於 $\ell$
個鍵，借出一個還有至少 $\ell$ 個；不足的節點拿到一個，變成 $\ell$ 個。 分隔鍵換成兄弟最靠
近的鍵，大小關係不變：例如向右兄弟借，下來的 $x.\mathit{key}[i]$ 比右兄弟所有的鍵小、
比不足的節點所有的鍵大，放在最後面；上去的是右兄弟最小的鍵，比留在右兄弟的都小。
搬過去的子樹，裡面的鍵也剛好落在新的範圍裡。 \textbf{合併}：兄弟剛好 $\ell$ 個鍵，合起
來有 $(\ell - 1) + 1 + \ell = 2\ell = 2\lceil m/2 \rceil - 2 \le (m + 1) - 2 = m - 1$ 個，沒有超
過上限。 借和合併都只在同一層搬鍵，不改變任何葉子的深度；唯一改變高度的是拿掉空的
根，所有葉子一起上升一層。 所以第 6 條「葉子同層」一直成立。
\end{whybox}

刪除的成本和插入一樣：往下 $h$ 層，回程每層最多一次借或合併，每次 $O(m)$，讀寫磁碟
$O(h)$ 次。 CLRS 用的是另一種寫法：往下走的時候，先確定要進去的子節點至少有 $\ell + 1$
個鍵（不夠就先借或合併），這樣到了葉子拿掉一個鍵也不會不足，不必回頭。 兩種做法得到
的樹可能不同；本書的追蹤都用上面這種由下而上的做法。

\subsection{刪除的完整追蹤}
\label{sec:mw-b-deltrace}

從圖~\ref{fig:mw-b-ins}(o) 的樹出發，再插入 $20, 33, 60$（都沒有分裂），得到
圖~\ref{fig:mw-b-del}(a)。 接著依序刪除 $12, 96, 48$。 $5$ 階 B 樹的非根節點至少要
$2$ 個鍵。

\begin{center}
\begin{tikzpicture}[baseline=(current bounding box.north), mwk/.append style={minimum width=0.46cm, minimum height=0.42cm}]
\draw[mwe] (1.565,-1.01) -- (0.46,-1.39);
\draw[mwe] (2.025,-1.01) -- (1.91,-1.39);
\draw[mwe] (2.485,-1.01) -- (3.59,-1.39);
\draw[mwe] (6.03,-1.01) -- (5.04,-1.39);
\draw[mwe] (6.49,-1.01) -- (6.49,-1.39);
\draw[mwe] (6.95,-1.01) -- (7.94,-1.39);
\draw[mwe] (4.027,-0.21) -- (2.025,-0.59);
\draw[mwe] (4.487,-0.21) -- (6.49,-0.59);
\node[mwk] at (0.23,-1.6) {12};
\node[mwk] at (0.69,-1.6) {14};
\node[mwk,mwkN] at (1.45,-1.6) {20};
\node[mwk] at (1.91,-1.6) {23};
\node[mwk] at (2.37,-1.6) {28};
\node[mwk,mwkN] at (3.13,-1.6) {33};
\node[mwk] at (3.59,-1.6) {40};
\node[mwk] at (4.05,-1.6) {47};
\node[mwk] at (1.795,-0.8) {16};
\node[mwk] at (2.255,-0.8) {29};
\node[mwk] at (4.81,-1.6) {50};
\node[mwk] at (5.27,-1.6) {53};
\node[mwk] at (6.03,-1.6) {57};
\node[mwk,mwkN] at (6.49,-1.6) {60};
\node[mwk] at (6.95,-1.6) {65};
\node[mwk] at (7.71,-1.6) {81};
\node[mwk] at (8.17,-1.6) {96};
\node[mwk] at (6.26,-0.8) {56};
\node[mwk] at (6.72,-0.8) {72};
\node[mwk] at (4.258,0) {48};
\node[mwlab] at (4.2,-2.08) {(a) 開始的樹};
\end{tikzpicture}\\*[0.45em]
\begin{tikzpicture}[baseline=(current bounding box.north), mwk/.append style={minimum width=0.46cm, minimum height=0.42cm}]
\draw[mwe] (1.22,-1.01) -- (0.23,-1.39);
\draw[mwe] (1.68,-1.01) -- (1.45,-1.39);
\draw[mwe] (2.14,-1.01) -- (3.13,-1.39);
\draw[mwe] (5.57,-1.01) -- (4.58,-1.39);
\draw[mwe] (6.03,-1.01) -- (6.03,-1.39);
\draw[mwe] (6.49,-1.01) -- (7.48,-1.39);
\draw[mwe] (3.625,-0.21) -- (1.68,-0.59);
\draw[mwe] (4.085,-0.21) -- (6.03,-0.59);
\node[mwk,mwkB] at (0.23,-1.6) {14};
\draw[mwbad] (0,-1.81) rectangle (0.46,-1.39);
\node[mwk] at (0.99,-1.6) {20};
\node[mwk] at (1.45,-1.6) {23};
\node[mwk] at (1.91,-1.6) {28};
\node[mwk] at (2.67,-1.6) {33};
\node[mwk] at (3.13,-1.6) {40};
\node[mwk] at (3.59,-1.6) {47};
\node[mwk] at (1.45,-0.8) {16};
\node[mwk] at (1.91,-0.8) {29};
\node[mwk] at (4.35,-1.6) {50};
\node[mwk] at (4.81,-1.6) {53};
\node[mwk] at (5.57,-1.6) {57};
\node[mwk] at (6.03,-1.6) {60};
\node[mwk] at (6.49,-1.6) {65};
\node[mwk] at (7.25,-1.6) {81};
\node[mwk] at (7.71,-1.6) {96};
\node[mwk] at (5.8,-0.8) {56};
\node[mwk] at (6.26,-0.8) {72};
\node[mwk] at (3.855,0) {48};
\node[mwlab] at (3.97,-2.08) {(b) 刪除 12：葉子不足};
\end{tikzpicture}\\*[0.45em]
\begin{tikzpicture}[baseline=(current bounding box.north), mwk/.append style={minimum width=0.46cm, minimum height=0.42cm}]
\draw[mwe] (1.335,-1.01) -- (0.46,-1.39);
\draw[mwe] (1.795,-1.01) -- (1.68,-1.39);
\draw[mwe] (2.255,-1.01) -- (3.13,-1.39);
\draw[mwe] (5.57,-1.01) -- (4.58,-1.39);
\draw[mwe] (6.03,-1.01) -- (6.03,-1.39);
\draw[mwe] (6.49,-1.01) -- (7.48,-1.39);
\draw[mwe] (3.682,-0.21) -- (1.795,-0.59);
\draw[mwe] (4.143,-0.21) -- (6.03,-0.59);
\node[mwk] at (0.23,-1.6) {14};
\node[mwk,mwkM] at (0.69,-1.6) {16};
\node[mwk] at (1.45,-1.6) {23};
\node[mwk] at (1.91,-1.6) {28};
\node[mwk] at (2.67,-1.6) {33};
\node[mwk] at (3.13,-1.6) {40};
\node[mwk] at (3.59,-1.6) {47};
\node[mwk,mwkM] at (1.565,-0.8) {20};
\node[mwk] at (2.025,-0.8) {29};
\node[mwk] at (4.35,-1.6) {50};
\node[mwk] at (4.81,-1.6) {53};
\node[mwk] at (5.57,-1.6) {57};
\node[mwk] at (6.03,-1.6) {60};
\node[mwk] at (6.49,-1.6) {65};
\node[mwk] at (7.25,-1.6) {81};
\node[mwk] at (7.71,-1.6) {96};
\node[mwk] at (5.8,-0.8) {56};
\node[mwk] at (6.26,-0.8) {72};
\node[mwk] at (3.912,0) {48};
\node[mwlab] at (3.97,-2.08) {(c) 向右兄弟借：16 下來，20 上去};
\end{tikzpicture}\\*[0.45em]
\begin{tikzpicture}[baseline=(current bounding box.north), mwk/.append style={minimum width=0.46cm, minimum height=0.42cm}]
\draw[mwe] (1.335,-1.01) -- (0.46,-1.39);
\draw[mwe] (1.795,-1.01) -- (1.68,-1.39);
\draw[mwe] (2.255,-1.01) -- (3.13,-1.39);
\draw[mwe] (5.455,-1.01) -- (4.58,-1.39);
\draw[mwe] (5.915,-1.01) -- (6.03,-1.39);
\draw[mwe] (6.375,-1.01) -- (7.25,-1.39);
\draw[mwe] (3.625,-0.21) -- (1.795,-0.59);
\draw[mwe] (4.085,-0.21) -- (5.915,-0.59);
\node[mwk] at (0.23,-1.6) {14};
\node[mwk] at (0.69,-1.6) {16};
\node[mwk] at (1.45,-1.6) {23};
\node[mwk] at (1.91,-1.6) {28};
\node[mwk] at (2.67,-1.6) {33};
\node[mwk] at (3.13,-1.6) {40};
\node[mwk] at (3.59,-1.6) {47};
\node[mwk] at (1.565,-0.8) {20};
\node[mwk] at (2.025,-0.8) {29};
\node[mwk] at (4.35,-1.6) {50};
\node[mwk] at (4.81,-1.6) {53};
\node[mwk] at (5.57,-1.6) {57};
\node[mwk] at (6.03,-1.6) {60};
\node[mwk] at (6.49,-1.6) {65};
\node[mwk,mwkB] at (7.25,-1.6) {81};
\draw[mwbad] (7.02,-1.81) rectangle (7.48,-1.39);
\node[mwk] at (5.685,-0.8) {56};
\node[mwk] at (6.145,-0.8) {72};
\node[mwk] at (3.855,0) {48};
\node[mwlab] at (3.74,-2.08) {(d) 刪除 96：葉子不足};
\end{tikzpicture}\hspace{0.55cm}\begin{tikzpicture}[baseline=(current bounding box.north), mwk/.append style={minimum width=0.46cm, minimum height=0.42cm}]
\draw[mwe] (1.335,-1.01) -- (0.46,-1.39);
\draw[mwe] (1.795,-1.01) -- (1.68,-1.39);
\draw[mwe] (2.255,-1.01) -- (3.13,-1.39);
\draw[mwe] (5.34,-1.01) -- (4.58,-1.39);
\draw[mwe] (5.8,-1.01) -- (5.8,-1.39);
\draw[mwe] (6.26,-1.01) -- (7.02,-1.39);
\draw[mwe] (3.567,-0.21) -- (1.795,-0.59);
\draw[mwe] (4.027,-0.21) -- (5.8,-0.59);
\node[mwk] at (0.23,-1.6) {14};
\node[mwk] at (0.69,-1.6) {16};
\node[mwk] at (1.45,-1.6) {23};
\node[mwk] at (1.91,-1.6) {28};
\node[mwk] at (2.67,-1.6) {33};
\node[mwk] at (3.13,-1.6) {40};
\node[mwk] at (3.59,-1.6) {47};
\node[mwk] at (1.565,-0.8) {20};
\node[mwk] at (2.025,-0.8) {29};
\node[mwk] at (4.35,-1.6) {50};
\node[mwk] at (4.81,-1.6) {53};
\node[mwk] at (5.57,-1.6) {57};
\node[mwk] at (6.03,-1.6) {60};
\node[mwk,mwkM] at (6.79,-1.6) {72};
\node[mwk] at (7.25,-1.6) {81};
\node[mwk] at (5.57,-0.8) {56};
\node[mwk,mwkM] at (6.03,-0.8) {65};
\node[mwk] at (3.797,0) {48};
\node[mwlab] at (3.74,-2.08) {(e) 向左兄弟借：72 下來，65 上去};
\end{tikzpicture}\\*[0.45em]
\begin{tikzpicture}[baseline=(current bounding box.north), mwk/.append style={minimum width=0.46cm, minimum height=0.42cm}]
\draw[mwe] (1.335,-1.01) -- (0.46,-1.39);
\draw[mwe] (1.795,-1.01) -- (1.68,-1.39);
\draw[mwe] (2.255,-1.01) -- (3.13,-1.39);
\draw[mwe] (4.995,-1.01) -- (4.35,-1.39);
\draw[mwe] (5.455,-1.01) -- (5.34,-1.39);
\draw[mwe] (5.915,-1.01) -- (6.56,-1.39);
\draw[mwe] (3.395,-0.21) -- (1.795,-0.59);
\draw[mwe] (3.855,-0.21) -- (5.455,-0.59);
\node[mwk] at (0.23,-1.6) {14};
\node[mwk] at (0.69,-1.6) {16};
\node[mwk] at (1.45,-1.6) {23};
\node[mwk] at (1.91,-1.6) {28};
\node[mwk] at (2.67,-1.6) {33};
\node[mwk] at (3.13,-1.6) {40};
\node[mwk] at (3.59,-1.6) {47};
\node[mwk] at (1.565,-0.8) {20};
\node[mwk] at (2.025,-0.8) {29};
\node[mwk,mwkB] at (4.35,-1.6) {53};
\draw[mwbad] (4.12,-1.81) rectangle (4.58,-1.39);
\node[mwk] at (5.11,-1.6) {57};
\node[mwk] at (5.57,-1.6) {60};
\node[mwk] at (6.33,-1.6) {72};
\node[mwk] at (6.79,-1.6) {81};
\node[mwk] at (5.225,-0.8) {56};
\node[mwk] at (5.685,-0.8) {65};
\node[mwk,mwkM] at (3.625,0) {50};
\node[mwlab] at (3.51,-2.08) {(f) 刪除 48：換成後繼 50，葉子不足};
\end{tikzpicture}\hspace{0.55cm}\begin{tikzpicture}[baseline=(current bounding box.north), mwk/.append style={minimum width=0.46cm, minimum height=0.42cm}]
\draw[mwe] (1.335,-1.01) -- (0.46,-1.39);
\draw[mwe] (1.795,-1.01) -- (1.68,-1.39);
\draw[mwe] (2.255,-1.01) -- (3.13,-1.39);
\draw[mwe] (5.65,-1.01) -- (5.04,-1.39);
\draw[mwe] (6.11,-1.01) -- (6.72,-1.39);
\draw[mwe] (3.607,-0.21) -- (1.795,-0.59);
\draw[mwe] (4.067,-0.21) -- (5.88,-0.59);
\node[mwk] at (0.23,-1.6) {14};
\node[mwk] at (0.69,-1.6) {16};
\node[mwk] at (1.45,-1.6) {23};
\node[mwk] at (1.91,-1.6) {28};
\node[mwk] at (2.67,-1.6) {33};
\node[mwk] at (3.13,-1.6) {40};
\node[mwk] at (3.59,-1.6) {47};
\node[mwk] at (1.565,-0.8) {20};
\node[mwk] at (2.025,-0.8) {29};
\node[mwk] at (4.35,-1.6) {53};
\node[mwk,mwkM] at (4.81,-1.6) {56};
\node[mwk] at (5.27,-1.6) {57};
\node[mwk] at (5.73,-1.6) {60};
\node[mwk] at (6.49,-1.6) {72};
\node[mwk] at (6.95,-1.6) {81};
\node[mwk,mwkB] at (5.88,-0.8) {65};
\draw[mwbad] (5.65,-1.01) rectangle (6.11,-0.59);
\node[mwk] at (3.837,0) {50};
\node[mwlab] at (3.59,-2.08) {(g) 和右兄弟合併，56 下來，父節點不足};
\end{tikzpicture}\\*[0.45em]
\begin{tikzpicture}[baseline=(current bounding box.north), mwk/.append style={minimum width=0.46cm, minimum height=0.42cm}]
\draw[mwe] (2.67,-1.01) -- (0.46,-1.39);
\draw[mwe] (3.13,-1.01) -- (1.68,-1.39);
\draw[mwe] (3.59,-1.01) -- (3.13,-1.39);
\draw[mwe] (4.05,-1.01) -- (5.04,-1.39);
\draw[mwe] (4.51,-1.01) -- (6.72,-1.39);
\draw[mwe] (3.59,-0.21) -- (3.59,-0.59);
\node[mwk] at (0.23,-1.6) {14};
\node[mwk] at (0.69,-1.6) {16};
\node[mwk] at (1.45,-1.6) {23};
\node[mwk] at (1.91,-1.6) {28};
\node[mwk] at (2.67,-1.6) {33};
\node[mwk] at (3.13,-1.6) {40};
\node[mwk] at (3.59,-1.6) {47};
\node[mwk] at (4.35,-1.6) {53};
\node[mwk] at (4.81,-1.6) {56};
\node[mwk] at (5.27,-1.6) {57};
\node[mwk] at (5.73,-1.6) {60};
\node[mwk] at (6.49,-1.6) {72};
\node[mwk] at (6.95,-1.6) {81};
\node[mwk] at (2.9,-0.8) {20};
\node[mwk] at (3.36,-0.8) {29};
\node[mwk,mwkM] at (3.82,-0.8) {50};
\node[mwk] at (4.28,-0.8) {65};
\draw[mwempty] (3.36,-0.21) rectangle (3.82,0.21);
\node[mwlab] at (3.59,-2.08) {(h) 和左兄弟合併，50 下來，根變空};
\end{tikzpicture}\hspace{0.55cm}\begin{tikzpicture}[baseline=(current bounding box.north), mwk/.append style={minimum width=0.46cm, minimum height=0.42cm}]
\draw[mwe] (2.67,-0.21) -- (0.46,-0.59);
\draw[mwe] (3.13,-0.21) -- (1.68,-0.59);
\draw[mwe] (3.59,-0.21) -- (3.13,-0.59);
\draw[mwe] (4.05,-0.21) -- (5.04,-0.59);
\draw[mwe] (4.51,-0.21) -- (6.72,-0.59);
\node[mwk] at (0.23,-0.8) {14};
\node[mwk] at (0.69,-0.8) {16};
\node[mwk] at (1.45,-0.8) {23};
\node[mwk] at (1.91,-0.8) {28};
\node[mwk] at (2.67,-0.8) {33};
\node[mwk] at (3.13,-0.8) {40};
\node[mwk] at (3.59,-0.8) {47};
\node[mwk] at (4.35,-0.8) {53};
\node[mwk] at (4.81,-0.8) {56};
\node[mwk] at (5.27,-0.8) {57};
\node[mwk] at (5.73,-0.8) {60};
\node[mwk] at (6.49,-0.8) {72};
\node[mwk] at (6.95,-0.8) {81};
\node[mwk] at (2.9,0) {20};
\node[mwk] at (3.36,0) {29};
\node[mwk] at (3.82,0) {50};
\node[mwk] at (4.28,0) {65};
\node[mwlab] at (3.59,-2.08) {(i) 拿掉空的根};
\end{tikzpicture}
\diagcaption{$5$ 階 B 樹的刪除。 (h) 之後拿掉空的根，高度從 $3$ 變成 $2$。}
\label{fig:mw-b-del}
\end{center}

\begin{itemize}
\item 刪除 $12$：葉子剩 $\{14\}$，不足。 右兄弟 $\{20, 23, 28\}$ 有 $3$ 個鍵，可以借：
分隔鍵 $16$ 下來，$20$ 上去。
\item 刪除 $96$：葉子剩 $\{81\}$，不足。 它是最右邊的子節點，沒有右兄弟；左兄弟
$\{57, 60, 65\}$ 可以借：$72$ 下來，$65$ 上去。
\item 刪除 $48$：$48$ 在根，後繼是右子樹最左的葉子 $\{50, 53\}$ 裡的 $50$。 $50$ 上去
取代 $48$，葉子剩 $\{53\}$，不足。 右兄弟 $\{57, 60\}$ 只有 $2$ 個鍵，不能借；它是第
一個子節點，沒有左兄弟。 和右兄弟合併：$53$、分隔鍵 $56$、$57, 60$ 合成
$\{53, 56, 57, 60\}$。 父節點剩 $\{65\}$，不足；它是根的最後一個子節點，左兄弟
$\{20, 29\}$ 只有 $2$ 個鍵，合併成 $\{20, 29, 50, 65\}$，根變空，拿掉。
\end{itemize}

最後一步合併出來的 $\{20, 29, 50, 65\}$ 有 $4$ 個鍵，剛好是上限 $m - 1$。 這就是上面
「為什麼一定對」框裡 $2\ell \le m - 1$ 等號成立的情況：$m$ 是奇數時合併的結果可能剛好滿。

\section{\texorpdfstring{B$^{+}$}{B+} 樹與磁碟存取}
\label{sec:mw-bplus}

\subsection{磁碟存取的成本}
\label{sec:mw-disk}

硬碟和固態硬碟都以\textbf{區塊}（block，也叫頁 page）為單位讀寫，常見的大小是 4 KiB
到 16 KiB。 讀一個區塊的時間，傳統硬碟約幾毫秒，固態硬碟約幾十到一百多微秒；記憶
體裡比較兩個整數只要幾奈秒。 所以磁碟上的資料結構，成本幾乎只看\textbf{讀寫了幾個區
塊}。 B 樹的做法是讓一個節點剛好是一個區塊：區塊放得下幾個鍵，就把階數 $m$
開到多大，樹高就是查一筆資料要讀的區塊數。

\textbf{例 3.} 區塊 $4096$ 位元組，鍵 $8$ 位元組，子指標（區塊編號）$8$ 位元組，每個
鍵還要帶一個指向資料記錄的指標 $8$ 位元組。 一個 B 樹節點有 $m - 1$ 個鍵和記錄指標、
$m$ 個子指標，要
\[
16(m - 1) + 8m \le 4096 \;\Longrightarrow\; 24m \le 4112 \;\Longrightarrow\; m \le 171.
\]
取 $m = 171$，$c = 86$。 存 $10^8$ 筆資料時，由定理~\ref{thm:mw-height}：
$171^3 = 5{,}000{,}211 < 10^8 + 1$，所以 $h \ge 4$；$86^3 = 636{,}056 \le 5 \times 10^7$、
$86^4 = 54{,}700{,}816 > 5 \times 10^7$，所以 $h - 1 \le 3$，$h \le 4$。 高度剛好是
$4$，查任何一筆最多讀 $4$ 個區塊；根常駐在記憶體的話只要讀 $3$ 個。 同樣 $10^8$ 筆
存在紅黑樹裡，高度可能到 $2\lg(10^8 + 1) \approx 53$。

\subsection{\texorpdfstring{B$^{+}$}{B+} 樹}

資料庫和檔案系統實際用的大多是 B 樹的變形 \textbf{B$^{+}$ 樹}（B$^{+}$-tree）。 和 B
樹的差別：
\begin{itemize}
\item \textbf{所有的鍵和資料記錄（或記錄指標）都在葉子。} 內部節點只放\textbf{分隔鍵}，
用來決定往哪棵子樹走，是葉子裡某些鍵的複本。
\item \textbf{葉子由左到右串成一條串列}，每片葉子多一個指標指向下一片。
\item 內部節點的規則和 B 樹一樣（最多 $m$ 個子節點、所有葉子同層）；葉子能放幾個鍵另
外由區塊大小決定，可以和 $m - 1$ 不同。
\end{itemize}

本書的約定是：分隔鍵等於它右邊那棵子樹裡最小的鍵。 搜尋 $k$ 時，在內部節點遇到
$k \ge$ 分隔鍵就往右走，所以即使在內部節點看到和 $k$ 相等的鍵，也要繼續往下，搜尋一
定走到葉子才結束。 圖~\ref{fig:mw-bplus} 是一棵內部節點 $4$ 階、葉子最多放 $3$ 個鍵的
B$^{+}$ 樹。

\begin{center}
\begin{tikzpicture}[baseline=(current bounding box.north), mwk/.append style={minimum width=0.5cm, minimum height=0.42cm}]
\draw[mwe] (1.125,-1.06) -- (0.5,-1.49);
\draw[mwe] (1.625,-1.06) -- (2.25,-1.49);
\draw[mwe] (5.25,-1.06) -- (4,-1.49);
\draw[mwe] (5.75,-1.06) -- (5.75,-1.49);
\draw[mwe] (6.25,-1.06) -- (7.5,-1.49);
\draw[mwe] (3.312,-0.21) -- (1.375,-0.64);
\draw[mwe] (3.812,-0.21) -- (5.75,-0.64);
\node[mwk] at (0.25,-1.7) {12};
\node[mwk] at (0.75,-1.7) {18};
\node[mwk] at (1.75,-1.7) {25};
\node[mwk] at (2.25,-1.7) {31};
\node[mwk] at (2.75,-1.7) {40};
\node[mwk] at (1.375,-0.85) {25};
\node[mwk] at (3.75,-1.7) {47};
\node[mwk] at (4.25,-1.7) {53};
\node[mwk] at (5.25,-1.7) {60};
\node[mwk] at (5.75,-1.7) {66};
\node[mwk] at (6.25,-1.7) {74};
\node[mwk] at (7.25,-1.7) {82};
\node[mwk] at (7.75,-1.7) {95};
\node[mwk] at (5.5,-0.85) {60};
\node[mwk] at (6,-0.85) {82};
\node[mwk] at (3.562,0) {47};
\draw[mwlink] (1,-1.7) -- (1.5,-1.7);
\draw[mwlink] (3,-1.7) -- (3.5,-1.7);
\draw[mwlink] (4.5,-1.7) -- (5,-1.7);
\draw[mwlink] (6.5,-1.7) -- (7,-1.7);
\node[mwlab] at (4,-2.18) {};
\end{tikzpicture}
\diagcaption{一棵 B$^{+}$ 樹。 根的 $47$ 和內部節點的 $25, 60, 82$ 都是分隔鍵，它們
在葉子裡也都有。 葉子之間的箭頭把所有的鍵串成由小到大的一條串列。}
\label{fig:mw-bplus}
\end{center}

\textbf{範圍查詢}是 B$^{+}$ 樹最主要的好處。 要找 $30 \le k \le 70$ 的所有資料，先像
搜尋一樣走到 $30$ 該在的葉子 $\{25, 31, 40\}$，接著沿著葉子的串列往右讀，讀到超過
$70$ 為止：$31, 40, 47, 53, 60, 66$。 讀的區塊數是路上的 $h - 1$ 個內部節點，加上涵蓋
這個範圍的葉子數，這裡是 $2 + 3 = 5$ 個。
B 樹做同樣的事要做一次中序走訪，上上下下會重複讀內部節點。

\textbf{插入}時，葉子溢位就分裂，右半的第一個鍵\textbf{複製}一份上去當分隔鍵，它自己
仍然留在葉子裡；內部節點溢位時，和 B 樹一樣把中間的鍵\textbf{移}上去，不留在原節點。
例如在圖~\ref{fig:mw-bplus} 插入 $35$，葉子變成 $\{25, 31, 35, 40\}$，溢位，分成
$\{25, 31\}$ 和 $\{35, 40\}$，$35$ 複製上去，父節點變成 $\{25, 35\}$，葉子的串列也要
接上新的葉子。 \textbf{刪除}時從葉子拿掉鍵，不足就和兄弟葉子借或合併；內部節點裡
的分隔鍵即使在葉子裡已經被刪掉，仍然能正確地分隔左右，可以先留著。

\begin{center}
\begin{tabular}{p{0.26\linewidth}p{0.32\linewidth}p{0.32\linewidth}}
\toprule
 & B 樹 & B$^{+}$ 樹 \\
\midrule
鍵出現幾次 & 每個鍵一次，可能在任何一層 & 每個鍵都在葉子；分隔鍵在內部節點另有複本 \\
資料記錄放在 & 每個鍵旁邊 & 只在葉子 \\
搜尋在哪裡結束 & 找到就停，可能在內部節點 & 一定走到葉子 \\
內部節點的分支數 & 鍵要帶記錄指標，分支較少 & 只放分隔鍵和子指標，分支較多，樹較矮 \\
範圍查詢 & 中序走訪 & 找到起點後沿葉子串列往右讀 \\
分裂時上去的鍵 & 移上去 & 葉子：複製上去；內部節點：移上去 \\
\bottomrule
\end{tabular}
\end{center}

第四列可以用例 3 的數字算：B$^{+}$ 樹的內部節點不放記錄指標，$8(m - 1) + 8m \le 4096$，
$m$ 可以到 $256$，比 B 樹的 $171$ 大。 有些課本還會提到 \textbf{B$^{*}$ 樹}：節點溢位
時先把鍵分給兄弟，兄弟也滿了才把兩個節點分成三個，讓非根節點至少 $2/3$ 滿。 想法和
B 樹相同，本書不細講。

\section{本章常見錯誤}
\label{sec:mw-pitfalls}

\begin{center}
\begin{tabular}{p{0.24\linewidth}p{0.50\linewidth}p{0.16\linewidth}}
\toprule
錯誤 & 例子 & 見 \\
\midrule
階數的定義混用 & 把 CLRS 的 $t$ 當成 $m$，$t = 3$ 的 B 樹當成 $3$ 階（其實是 $6$ 階） &
第~\ref{sec:mw-btree}~節 \\
最少鍵數算錯 & $5$ 階 B 樹的非根節點寫成至少 $\lfloor 5/2 \rfloor = 2$ 個子節點；應該是
$\lceil 5/2 \rceil = 3$ 個子節點、$2$ 個鍵 & 定義~\ref{def:mw-btree} \\
高度差一格 & 用 CLRS 數邊的高度代入本書的公式，或反過來 & 第~\ref{sec:mw-height}~節 \\
分裂時中間鍵留在原節點 & B 樹的分裂把 $k_p$ 同時留在左半又送上去；只有 B$^{+}$ 樹的葉子
分裂才複製 & 第~\ref{sec:mw-insert}、\ref{sec:mw-bplus}~節 \\
直接刪內部節點的鍵 & 從內部節點拿掉鍵，兩棵子樹少了分隔鍵；應該換成後繼（或前驅） &
第~\ref{sec:mw-23-delete}~節 \\
借鍵時跳過父節點 & 兄弟的鍵直接搬到不足的節點；應該是分隔鍵下來、兄弟的鍵上去 &
第~\ref{sec:mw-23-delete}~節 \\
內部節點借鍵時忘了搬子樹 & 鍵借過來了，兄弟那一邊的子指標沒有跟著搬，子節點數和鍵數
對不上 & 圖~\ref{fig:mw-23-del}(d) \\
合併後忘了往上 & 父節點少了一個鍵也可能不足，要繼續處理；根變空要拿掉 &
第~\ref{sec:mw-delete}~節 \\
由上而下、由下而上混用 & 2-3-4 樹題目指定由上而下，卻等到葉子溢位才分裂 &
第~\ref{sec:mw-234}~節 \\
\bottomrule
\end{tabular}
\end{center}

\section{觀念題}
\label{sec:mw-exercises}

題型和研究所考試、期中期末考的筆試相近，有選擇、計算、追蹤和證明。 除非題目另外說
明，階數 $m$ 是子節點數的上限，高度照本章的定義（數節點），插入用由下而上的分裂（上升
第 $\lceil m/2 \rceil$ 個鍵），刪除內部節點的鍵用後繼取代，借和合併的順序照
第~\ref{sec:mw-delete}~節。 解答在本章最後。

\begin{enumerate}
\item （選擇，可能不只一個答案）下列哪些敘述正確？
\begin{enumerate}
\item $5$ 階 B 樹裡，根以外的每個節點至少有 $2$ 個鍵。
\item B 樹的插入可能讓樹高減少。
\item B 樹的刪除如果發生了借，就不會再往上影響父節點以上的節點。
\item CLRS 最小分支度 $t = 2$ 的 B 樹就是 2-3-4 樹。
\item 在 B$^{+}$ 樹裡搜尋，只要在內部節點看到要找的鍵就可以停。
\end{enumerate}

\item （計算）
\begin{enumerate}
\item $7$ 階 B 樹有 $2000$ 個鍵，高度最小和最大各是多少？
\item 高度 $3$ 的 $7$ 階 B 樹最少、最多有幾個鍵？
\item 如果改用 CLRS 數邊的高度（只有根的樹高度 $0$），(b) 的「高度 $3$」最少、最多有
幾個鍵？
\end{enumerate}

\item （追蹤）依序把 $64, 18, 93, 40, 75, 27, 52, 86, 11, 33$ 插入空的 2-3 樹，畫出每
次分裂之後的樹。

\item （追蹤）接著從第 3 題最後的樹依序刪除 $33, 64, 93, 75$，畫出每次借或合併之後的
樹，寫出用的是借還是合併、向哪一邊。

\item （追蹤）依序把 $31, 74, 58, 12, 87, 45, 66, 20, 93, 38, 51, 79, 26, 62, 70, 84$
插入空的 $5$ 階 B 樹，畫出每次分裂之後的樹。 再從最後的樹依序刪除
$87, 93, 31, 45, 20, 12$。

\item （追蹤）
\begin{enumerate}
\item 用由上而下的做法，依序把 $5, 10, 15, \ldots, 50$（十個鍵）插入空的 2-3-4 樹，畫
出最後的樹。
\item 畫出一棵和 (a) 的結果對應的紅黑樹。
\item 和 (a) 的結果對應的紅黑樹一共有幾棵？
\end{enumerate}

\item （計算）把 $10, 20, 30, 40$ 依序插入空的 $4$ 階 B 樹。
\begin{enumerate}
\item 用本章由下而上的規則（溢位時上升第 $\lceil m/2 \rceil = 2$ 個鍵）。
\item 改成溢位時上升第 $3$ 個鍵。
\item 改用第~\ref{sec:mw-234}~節由上而下的做法。
\end{enumerate}
三種做法各得到什麼樹？ 都是合法的 2-3-4 樹嗎？

\item （證明）對所有 $m \ge 3$：
\begin{enumerate}
\item 證明第~\ref{sec:mw-insert}~節的分裂之後，左右兩個節點都至少有
$\lceil m/2 \rceil - 1$ 個鍵。
\item 如果把「非根節點至少 $\lceil m/2 \rceil - 1$ 個鍵」改成至少 $\lceil m/2 \rceil$
個鍵，分裂還可能合法嗎？ 舉 $m = 4$ 說明。
\end{enumerate}

\item （證明）
\begin{enumerate}
\item 證明有 $n$ 個鍵的 B 樹（在葉子底下畫出外部節點）剛好有 $n + 1$ 個外部節點。
\item 用 (a) 重新證明定理~\ref{thm:mw-height} 的下界 $n \ge 2c^{h-1} - 1$。
\end{enumerate}

\item （計算）磁碟區塊 $8192$ 位元組，鍵 $16$ 位元組，子指標和記錄指標都是 $8$ 位元
組。 用 B$^{+}$ 樹存 $10^9$ 筆資料。
\begin{enumerate}
\item 內部節點放 $m - 1$ 個鍵和 $m$ 個子指標，$m$ 最大是多少？
\item 葉子放 $L$ 組（鍵，記錄指標），再加一個指向下一片葉子的指標，$L$ 最大是多少？
\item 假設除了根以外，每個內部節點至少 $\lceil m/2 \rceil$ 個子節點，每片葉子至少
$\lceil L/2 \rceil$ 個鍵（葉子不只一片）。 這棵樹的高度（包括葉子那一層）最少和最多
是多少？
\end{enumerate}

\end{enumerate}

\section{實作練習}
\label{sec:mw-practice}

Aki-Judge 目前還沒有和本章對應的題目，請用觀念題和範例~\ref{ex:mw-btree} 練習；寫刪
除時，可以照虛擬碼~\ref{alg:mw-deldown}、\ref{alg:mw-fix} 加進範例的程式，用
圖~\ref{fig:mw-b-del} 的例子檢查每一步。

\section{讀完你應該會}
\label{sec:mw-checklist}

\begin{itemize}
\item 能說出 $m$ 階 B 樹的定義和各種節點的鍵數範圍，換算 CLRS 的 $t$（$m = 2t$）；能證明 $\log_m(n+1) \le h \le 1 + \log_{\lceil m/2 \rceil}\frac{n+1}{2}$，並用
它算出給定鍵數的最小、最大高度，或給定高度的最少、最多鍵數。
\item 能對 2-3 樹、2-3-4 樹和一般的 $m$ 階 B 樹做插入和刪除的追蹤，每次分裂、借、合併
之後畫出整棵樹。
\item 能寫出 B 樹搜尋和插入的 C 程式，分裂時正確地分配鍵和子指標。
\item 能把 2-3-4 樹畫成紅黑樹、把紅黑樹畫成 2-3-4 樹，說出紅黑樹插入和刪除的各種情況
對應分裂、合併還是借。
\item 能說明 B$^{+}$ 樹和 B 樹的差別，用區塊大小算出階數和要讀幾次磁碟。
\end{itemize}

\textbf{延伸閱讀}（書末會整理完整的參考資料）

\begin{itemize}
\item Thomas H. Cormen 等，\textit{Introduction to Algorithms}，第 4 版，第 18 章：用
最小分支度 $t$ 定義 B 樹，插入和刪除都是由上而下一趟完成，並分析讀寫磁碟的次數。 讀
的時候記得換算 $m = 2t$，高度數邊。
\item Donald E. Knuth，\textit{The Art of Computer Programming}，第 3 冊，第 6.2.4 節：
多路搜尋樹、B 樹及其變形（包括 B$^{*}$ 樹）的原始講法和分析。
\item Robert Sedgewick、Kevin Wayne，\textit{Algorithms}，第 4 版，第 3.3 節：從 2-3
樹講到左傾紅黑樹，和本章第~\ref{sec:mw-rb}~節的對應互相參照。
\item Douglas Comer，``The Ubiquitous B-Tree''，\textit{ACM Computing Surveys}
11(2)，1979：B 樹和 B$^{+}$ 樹在資料庫、檔案系統裡的用法，篇幅短。
\end{itemize}

\addcontentsline{toc}{section}{觀念題解答}
\section*{觀念題解答}

解答裡用文字寫樹：$[a\ b]$ 是一個節點，後面括號裡依序是它的子樹，例如
$[47]([29],[63\ 82])$ 是根 $47$、左子節點 $\{29\}$、右子節點 $\{63, 82\}$。

\begin{enumerate}
\item (a)、(c)、(d)。
(a) 對：$\lceil 5/2 \rceil - 1 = 2$。 (b) 錯：插入只會讓根分裂，高度不變或加 $1$。
(c) 對：借之後父節點的鍵數不變，父節點以上都不受影響；只有合併會往上傳。
(d) 對：$t = 2$ 時非根節點有 $1$ 到 $3$ 個鍵、$2$ 到 $4$ 個子節點，就是 $m = 2t = 4$
階。 (e) 錯：內部節點的鍵只是分隔鍵，資料記錄在葉子，要走到葉子才算找到；而且刪除
之後分隔鍵可能還留著，對應的資料已經不在了。

\item $m = 7$，$c = \lceil 7/2 \rceil = 4$。
(a) 最小：$7^3 - 1 = 342 < 2000$，$7^4 - 1 = 2400 \ge 2000$，所以 $h \ge 4$。 最大：
$2 \cdot 4^{h-1} - 1 \le 2000$，即 $4^{h-1} \le 1000.5$；$4^4 = 256$ 可以，$4^5 = 1024$
不行，$h \le 5$。 高度是 $4$ 或 $5$。
(b) 最少 $2 \cdot 4^2 - 1 = 31$ 個，最多 $7^3 - 1 = 342$ 個。
(c) CLRS 的高度 $3$ 是本書的高度 $4$：最少 $2 \cdot 4^3 - 1 = 127$ 個，最多
$7^4 - 1 = 2400$ 個。

\item 有分裂的插入共四次（「溢位」後面是分裂前的節點）：
\begin{itemize}
\item 插入 $93$：根 $[18\ 64\ 93]$ 溢位，$64$ 上升成新根：$[64]([18],\allowbreak[93])$。
\item 插入 $27$：$[18\ 27\ 40]$ 溢位，$27$ 上升：$[27\ 64]([18],\allowbreak[40],\allowbreak[75\ 93])$。
\item 插入 $86$：$[75\ 86\ 93]$ 溢位，$86$ 上升，根 $[27\ 64\ 86]$ 溢位，$64$ 上升成新根：
$[64]([27]([18],\allowbreak[40\ 52]),[86]([75],\allowbreak[93]))$。
\item 插入 $33$：$[33\ 40\ 52]$ 溢位，$40$ 上升：
$[64]([27\ 40]([11\ 18],\allowbreak[33],\allowbreak[52]),[86]([75],\allowbreak[93]))$，高度 $3$。
\end{itemize}

\item
\begin{itemize}
\item 刪除 $33$：葉子變空；右兄弟 $[52]$ 不能借，左兄弟 $[11\ 18]$ 可以，向左兄弟借：
$27$ 下來，$18$ 上去：$[64]([18\ 40]([11],\allowbreak[27],\allowbreak[52]),[86]([75],\allowbreak[93]))$。
\item 刪除 $64$：後繼 $75$ 上去，葉子 $[75]$ 變空；沒有左兄弟，右兄弟 $[93]$ 不能借，和右
兄弟合併成 $[86\ 93]$；父節點變空，它的左兄弟 $[18\ 40]$ 可以借：$75$ 下來，$40$ 上去，
$[52]$ 這棵子樹跟著搬過來：$[40]([18]([11],\allowbreak[27]),[75]([52],\allowbreak[86\ 93]))$。
\item 刪除 $93$：葉子剩 $[86]$，結束。
\item 刪除 $75$：後繼 $86$ 上去，葉子變空；沒有右兄弟，左兄弟 $[52]$ 不能借，和左兄弟
合併成 $[52\ 86]$；父節點變空，左兄弟 $[18]$ 不能借，合併成 $[18\ 40]$；根變空拿掉：
$[18\ 40]([11],\allowbreak[27],\allowbreak[52\ 86])$，高度 $2$。
\end{itemize}

\item 插入時有三次分裂：
\begin{itemize}
\item 插入 $87$：根 $[12\ 31\ 58\ 74\ 87]$ 溢位，$58$ 上升：$[58]([12\ 31],\allowbreak[74\ 87])$。
\item 插入 $38$：$[12\ 20\ 31\ 38\ 45]$ 溢位，$31$ 上升：
$[31\ 58]([12\ 20],\allowbreak[38\ 45],\allowbreak[66\ 74\ 87\ 93])$。
\item 插入 $79$：$[66\ 74\ 79\ 87\ 93]$ 溢位，$79$ 上升：
$[31\ 58\ 79]([12\ 20],\allowbreak[38\ 45\ 51],\allowbreak[66\ 74],\allowbreak[87\ 93])$。
\end{itemize}
插完是 $[31\ 58\ 79]([12\ 20\ 26],\allowbreak[38\ 45\ 51],\allowbreak[62\ 66\ 70\ 74],$
$[84\ 87\ 93])$。 刪除：
\begin{itemize}
\item $87$：葉子剩 $[84\ 93]$，結束。
\item $93$：剩 $[84]$，不足；沒有右兄弟，左兄弟 $[62\ 66\ 70\ 74]$ 可以借：$79$ 下來，
$74$ 上去：$[31\ 58\ 74]([12\ 20\ 26],\allowbreak[38\ 45\ 51],\allowbreak[62\ 66\ 70],\allowbreak[79\ 84])$。
\item $31$：後繼 $38$ 上去，葉子剩 $[45\ 51]$，結束。
\item $45$：剩 $[51]$，不足；右兄弟 $[62\ 66\ 70]$ 可以借（左兄弟也可以，但本書先看右
邊）：$58$ 下來，$62$ 上去：$[38\ 62\ 74]([12\ 20\ 26],\allowbreak[51\ 58],\allowbreak[66\ 70],\allowbreak[79\ 84])$。
\item $20$：葉子剩 $[12\ 26]$，結束。
\item $12$：剩 $[26]$，不足；右兄弟 $[51\ 58]$ 不能借，沒有左兄弟，和右兄弟合併：
$[62\ 74]([26\ 38\ 51\ 58],\allowbreak[66\ 70],\allowbreak[79\ 84])$。 根還有兩個鍵，高度仍是 $2$。
\end{itemize}

\item (a) 分裂發生在插入 $20$（根 $[5\ 10\ 15]$，$10$ 上升）、$30$（$[15\ 20\ 25]$，$20$
上升）、$40$（$[25\ 30\ 35]$，$30$ 上升）、$45$（根 $[10\ 20\ 30]$ 已滿，先分裂，$20$ 上升成新根；要去的葉子 $[35\ 40]$ 不滿，直接放）和 $50$（$[35\ 40\ 45]$，$40$
上升）之前。 最後是 $[20]([10]([5],\allowbreak[15]),[30\ 40]([25],\allowbreak[35],\allowbreak[45\ 50]))$，高度 $3$。
(b) 下面是其中一棵（3-節點都用「黑色帶紅色右子節點」）：
\begin{center}
\begin{tikzpicture}[baseline=(current bounding box.north)]
\node[mwB] (rb5) at (0.00,-1.64) {5};
\node[mwB] (rb10) at (0.58,-0.82) {10};
\node[mwB] (rb15) at (1.16,-1.64) {15};
\node[mwB] (rb20) at (1.74,0.00) {20};
\node[mwB] (rb25) at (2.32,-1.64) {25};
\node[mwB] (rb30) at (2.90,-0.82) {30};
\node[mwB] (rb35) at (3.48,-2.46) {35};
\node[mwR] (rb40) at (4.06,-1.64) {40};
\node[mwB] (rb45) at (4.64,-2.46) {45};
\node[mwR] (rb50) at (5.22,-3.28) {50};
\draw[mwe] (rb20) -- (rb10);
\draw[mwe] (rb10) -- (rb5);
\draw[mwe] (rb10) -- (rb15);
\draw[mwe] (rb20) -- (rb30);
\draw[mwe] (rb30) -- (rb25);
\draw[mwe] (rb30) -- (rb40);
\draw[mwe] (rb40) -- (rb35);
\draw[mwe] (rb40) -- (rb45);
\draw[mwe] (rb45) -- (rb50);
\begin{scope}[on background layer]
\node[mwgrp, fit=(rb30) (rb40)] {};
\node[mwgrp, fit=(rb45) (rb50)] {};
\end{scope}
\node[mwlab] at (2.61,-4.03) {};
\end{tikzpicture}
\end{center}
(c) 有兩個 3-節點 $[30\ 40]$、$[45\ 50]$，每個有兩種畫法，共 $2^2 = 4$ 棵。

\item (a) $[10\ 20\ 30\ 40]$ 溢位，上升第 $2$ 個：$[20]([10],\allowbreak[30\ 40])$。
(b) 上升第 $3$ 個：$[30]([10\ 20],\allowbreak[40])$。
(c) 插入 $40$ 時根 $[10\ 20\ 30]$ 是 4-節點，先分裂成 $[20]([10],\allowbreak[30])$，再放進 $40$：
$[20]([10],\allowbreak[30\ 40])$，和 (a) 相同。 三棵都是合法的 2-3-4 樹，所以作答時一定要寫出用
的規則。

\item 令 $c = \lceil m/2 \rceil$。
(a) 左邊留 $p - 1 = c - 1$ 個，剛好是下限。 右邊有 $m - c = \lfloor m/2 \rfloor$ 個，而
$\lceil m/2 \rceil - \lfloor m/2 \rfloor \le 1$，所以 $\lfloor m/2 \rfloor \ge c - 1$。
(b) 不可能。 溢位的節點有 $m$ 個鍵，一個上升後剩 $m - 1$ 個分給兩邊，兩邊都要
$\ge c$ 就需要 $m - 1 \ge 2c \ge m$，不可能。 例如 $m = 4$：$4$ 個鍵上升一個，剩 $3$
個，不可能兩邊都至少 $2$ 個。 所以下限一定要寫成 $\lceil m/2 \rceil - 1$。

\item (a) 設樹裡（不算外部節點）有 $N$ 個節點。 把外部節點也算進去，共 $N + E$ 個節
點，除了根以外每個節點都有一條連到父節點的邊，邊數是 $N + E - 1$。 另一方面，有
$k$ 個鍵的節點恰有 $k + 1$ 個子節點（葉子的子節點就是外部節點），邊數也等於
$\sum (k_x + 1) = n + N$。 兩式相等得到 $E = n + 1$。
(b) 外部節點都在第 $h + 1$ 層。 照定理的證明，第 $i$ 層（$2 \le i \le h + 1$）至少
$2c^{i-2}$ 個節點（第 $h$ 層的葉子也至少有 $c$ 個外部子節點），所以
$n + 1 = E \ge 2c^{h-1}$。

\item (a) $16(m - 1) + 8m \le 8192$，$24m \le 8208$，$m = 342$。
(b) $24L + 8 \le 8192$，$L = 341$。
(c) 葉子的片數：每片最多 $341$ 個鍵，至少要 $\lceil 10^9 / 341 \rceil = 2{,}932{,}552$ 片；
每片至少 $171$ 個鍵，最多 $\lfloor 10^9 / 171 \rfloor = 5{,}847{,}953$ 片。 高度 $H$（葉
子在第 $H$ 層）時，葉子最多 $342^{H-1}$ 片，最少 $2 \cdot 171^{H-2}$ 片。
最矮：$342^2 = 116{,}964$ 不夠，$342^3 = 40{,}001{,}688$ 夠，$H \ge 4$。
最高：$2 \cdot 171^{H-2} \le 5{,}847{,}953$，$171^2 = 29{,}241$ 可以，$171^3 = 5{,}000{,}211$
不行，$H \le 4$。 所以高度一定是 $4$，查一筆資料讀 $4$ 個區塊。
\end{enumerate}

\chapter{雜湊}
\label{ch:hash}

% 本章的圖共用的樣式（名稱都以 h 開頭，避免和其他章衝突）
\tikzset{
    hcell/.style={draw=black!55, minimum width=0.92cm, minimum height=0.62cm,
        inner sep=0pt, font=\ttfamily\footnotesize},
    hlnew/.style={draw=memStroke, thick, fill=memFill},
    hlprobe/.style={fill=optFill},
    hldead/.style={draw=optStroke, thick, fill=optFill},
    hbucket/.style={draw=black!55, minimum width=0.8cm, minimum height=0.6cm,
        inner sep=0pt, font=\ttfamily\footnotesize},
    hnode/.style={draw=black!55, minimum height=0.6cm, inner sep=2pt,
        font=\ttfamily\scriptsize},
    hdir/.style={draw=black!55, minimum width=0.9cm, minimum height=0.5cm,
        inner sep=0pt, font=\ttfamily\footnotesize},
    hpage/.style={draw=coreStroke, fill=coreFill, minimum width=1.9cm,
        minimum height=0.5cm, inner sep=1pt, font=\footnotesize},
    hbit/.style={draw=black!55, minimum width=0.8cm, minimum height=0.55cm,
        inner sep=0pt, font=\ttfamily\footnotesize},
    hbzero/.style={draw=optStroke, thick, fill=optStroke, text=white},
}
% 刪除標記在虛擬碼裡的寫法
\newcommand{\hashDEL}{\textsc{deleted}}

學生餐廳的刷卡機一到中午就忙不過來。 每張學生證有一個九位數的卡號，機器要在一瞬間
查出這張卡還剩多少錢。 卡號有十億種可能，開一張以卡號當索引的表太浪費；把卡號排好
再二分搜尋，或是存進第~10~章的平衡樹，每次也要走二三十步。 有沒有辦法讓每次查詢平
均只看一兩個位置，而且這個次數幾乎和全校有幾張卡無關？ 往下讀之前，先想一想。

第~1~章找重複學號時用過一個辦法：開一個以學號為索引的陣列，在上面打勾。 它快，是
因為從學號可以直接算出要看哪一格；它浪費，是因為格子數要和學號的範圍一樣多。
\textbf{雜湊}（hashing）保留「從鍵直接算出位置」的做法，但只開 $m$ 格，$m$ 和實際的
資料量差不多，再用一個\textbf{雜湊函數}（hash function）把很大的鍵空間壓到
$0, 1, \ldots, m-1$。 代價是兩個不同的鍵可能算出同一格，叫做\textbf{碰撞}
（collision）。 這一章依序講怎麼選雜湊函數、碰撞了以後怎麼存（鏈結法與開放定址）、平
均要看幾格（期望值的推導），最後講表滿了怎麼辦，以及可延伸雜湊和通用雜湊兩個觀念。 最後一節的布隆過濾器不存鍵本身，只回答
「在不在」，用一點誤判換來很大的空間節省。

雜湊表只支援插入、刪除、搜尋這三個操作，平均每次 $O(1)$。 它不維持任何順序，找最
小值、找後繼、依序列出都得把整張表掃一遍；需要這些操作時，還是要用第~9、10~章的搜
尋樹。 這一章用到《C 程式設計基礎》第~14~章的結構、第~19~章的 \texttt{malloc}、
\texttt{calloc} 和串列，以及本書第~5~章「指標的指標」的刪除寫法和第~3~章的平攤分析。

\section{直接定址表與雜湊函數}
\label{sec:hash-func}

\subsection{直接定址表}

先看鍵的範圍很小的情形。 設所有可能的鍵組成集合 $U = \{0, 1, \ldots, u - 1\}$，叫做
\textbf{鍵空間}（universe）。 \textbf{直接定址表}（direct-address table）就是一個陣
列 $T[0..u-1]$：鍵 $k$ 的資料放在 $T[k]$，沒有這個鍵時 $T[k] = \NIL$。 插入、刪除、
搜尋都只碰一格，最壞情況也是 $O(1)$。 第~1~章的 \proc{Has-Dup-Marks} 就是這種表，只
是每格只存一個「有沒有看過」。

這裡的陣列從 $0$ 開始編號，和本書虛擬碼的慣例（從 $1$ 開始）不同。 雜湊表的格子號
碼是雜湊函數算出來的餘數，從 $0$ 開始最自然，所以這一章的表一律寫成
$T[0..m-1]$，C 程式也一樣，不用換算。

直接定址表的問題是空間。 它要 $\Theta(u)$ 格，和實際存了幾個鍵無關。 九位數的卡號
$u = 10^9$，就算每格只存一個位元組也要 1 GB，而全校的卡也許只有三萬張。

\subsection{雜湊表、碰撞與載入因子}

\textbf{雜湊表}（hash table）只開 $m$ 格，$T[0..m-1]$，鍵 $k$ 放在第 $h(k)$ 格附
近。 這裡
\[
h : U \to \{0, 1, \ldots, m - 1\}
\]
是雜湊函數，$h(k)$ 叫做 $k$ 的\textbf{雜湊值}（hash value）或\textbf{起始位置}
（home slot）。 存了 $n$ 個鍵時，
\[
\alpha = \frac{n}{m}
\]
叫做\textbf{載入因子}（load factor），也就是平均每格分到幾個鍵。 之後所有的成本都用
$\alpha$ 表示。

$u > m$ 時一定有會碰撞的鍵：把整個鍵空間的 $u$ 個鍵都算一次雜湊值，由鴿籠原理，至少有
一格分到 $\lceil u/m \rceil$ 個鍵。 這還說明了一件更糟的事：不管 $h$ 怎麼選，只要 $u \ge nm$，
那一格就分到至少 $\lceil u/m \rceil \ge n$ 個鍵，也就是有 $n$ 個鍵的雜湊值全部相同。 輸入剛好是這 $n$ 個鍵時，任何雜湊表都會退化成在一
堆碰撞的鍵裡逐一比對，最壞情況是 $\Theta(n)$。 所以雜湊的分析一律談期望值，
而且要先說清楚機率從哪裡來：可能是假設鍵的雜湊值像亂數一樣均勻（第~\ref{sec:hash-chain}、
\ref{sec:hash-analysis}~節），也可能是從一族函數裡隨機挑一個（第~\ref{sec:hash-dynamic}~節
的通用雜湊）。

實用的雜湊函數要算得快，而且鍵的每一部分都會影響結果，這樣只差一兩位數的鍵（學號、
卡號常常這樣）才會被分開。 下面四種是考題最常出現的整數鍵雜湊函數。

\subsection{除法}

\textbf{除法}（division method）取餘數：
\[
h(k) = k \bmod m.
\]
算得最快，但 $m$ 不能亂選。 $m = 10^p$ 時 $h(k)$ 就是 $k$ 的最後 $p$ 位數字，前面的位
數完全沒用到；如果卡號的最後兩位是系所代碼，$m = 100$ 會讓同一系的卡全部擠在同一格。
$m = 2^p$ 時同理，只用到最低的 $p$ 個位元。 另一個陷阱是公因數。

\begin{lemma}
\label{lem:hash-gcd}
如果每個鍵都是 $d$ 的倍數，令 $g = \gcd(d, m)$，則 $k \bmod m$ 一定是 $g$ 的倍數，
所以除法最多只用到 $m/g$ 格。
\end{lemma}

\begin{proof}
寫成 $k = qm + r$，$0 \le r < m$。 $g$ 整除 $k$（因為 $g \mid d$、$d \mid k$），也整除
$m$，所以整除 $r = k - qm$。 $0$ 到 $m - 1$ 之間 $g$ 的倍數只有 $m/g$ 個。
\end{proof}

例如記憶體位址通常是 $4$ 的倍數。 鍵 $8, 20, 36, 44, 52, 64, 76, 88$ 用 $m = 12$ 時，
餘數依序是 $8, 8, 0, 8, 4, 4, 4, 4$，只用到 $0, 4, 8$ 三格；改用 $m = 13$，餘數是
$8, 7, 10, 5, 0, 12, 11, 10$，分散開了。 所以除法的 $m$ 通常選質數，而且不要
太靠近 $2$ 或 $10$ 的次方。

\subsection{乘法}

\textbf{乘法}（multiplication method）先選一個常數 $0 < A < 1$，取 $kA$ 的小數部分，
再放大到 $0$ 到 $m-1$：
\[
h(k) = \bigl\lfloor m \,( kA \bmod 1 ) \bigr\rfloor,
\]
$kA \bmod 1 = kA - \lfloor kA \rfloor$ 表示小數部分。 常用的 $A$ 是
$(\sqrt{5} - 1)/2 \approx 0.6180339887$，這個值讓連續的鍵分得很開。 例如 $m = 1000$：
\[
\begin{aligned}
5347 A &= 3304.62773\ldots, & h(5347) &= \lfloor 627.73\ldots \rfloor = 627, \\
5348 A &= 3305.24577\ldots, & h(5348) &= \lfloor 245.77\ldots \rfloor = 245.
\end{aligned}
\]
只差 $1$ 的兩個鍵落在相距很遠的兩格。 乘法對 $m$ 沒有要求，通常取 $m = 2^p$。

程式裡不用浮點數。 設一個字組有 $w$ 個位元，取整數 $s = \lfloor A \cdot 2^w \rfloor$，
$m = 2^p$。 $s \approx A \cdot 2^w$，所以 $ks \approx kA \cdot 2^w$；對 $2^w$ 取餘數，等於丟掉
$kA$ 的整數部分，剩下的低 $w$ 位元大約是 $kA$ 的小數部分乘上 $2^w$。 這個 $w$ 位元的數，
最高的 $p$ 個位元就是它除以 $2^{w-p}$ 的商，大約是小數部分乘上 $2^p = m$ 再取整數，也就是
$h(k)$：
\[
h(k) = \bigl( (k \cdot s) \bmod 2^w \bigr) \gg (w - p).
\]
在 C 裡用 \texttt{unsigned} 的乘法，溢位自動就是 $\bmod 2^w$，再右移，一共兩個運算。
用 $w = 16$、$p = 4$（$m = 16$）算 $k = 3051$：$s = \lfloor 0.6180339887 \cdot 65536
\rfloor = 40503$，$3051 \cdot 40503 = 123574653$，除以 $2^{16}$ 的餘數是
$39293 = (1001\,1001\,0111\,1101)_2$，最高 $4$ 位是 $(1001)_2 = 9$。 用實數算也一樣：
$3051A = 1885.6217\ldots$，$\lfloor 16 \times 0.6217\ldots \rfloor = 9$。

\subsection{中間平方}

\textbf{中間平方法}（mid-square method）把鍵平方，取中間的幾位。 例如四位數的鍵、
$m = 100$：$k^2$ 補成八位數，取第 $4$、$5$ 位。
\[
5347^2 = 28590409 = 285\,\underline{90}\,409, \qquad h(5347) = 90.
\]
中間的位數是由 $k$ 的每一位一起乘出來的，所以 $k$ 的每一位都會影響結果；最後幾位只
和 $k$ 的最後幾位有關，前面幾位主要由 $k$ 的前幾位決定，都不適合。 二進位的版本取
$k^2$ 中間的 $p$ 個位元，$m = 2^p$。

\subsection{摺疊}

\textbf{摺疊法}（folding）把很長的鍵切成幾段，每段的位數和 $m - 1$ 相同，再加起來，
最後只留下和 $m - 1$ 一樣多的位數（多出來的進位丟掉）。 例如九位數的卡號
$417961974$、$m = 1000$，切成 $417 \mid 961 \mid 974$：

\begin{itemize}
\item \textbf{移位摺疊}（shift folding）：直接相加，$417 + 961 + 974 = 2352$，取後三
位，$h = 352$。
\item \textbf{邊界摺疊}（boundary folding）：像摺紙一樣，每隔一段反過來，
$417 + 169 + 974 = 1560$，$h = 560$。
\end{itemize}

摺疊適合位數比 $m$ 長很多的鍵（身分證號、電話號碼）。 實務上也常先摺疊成一個整數，
再用除法或乘法。

\subsection{字串鍵}

字串 $c_0 c_1 \cdots c_{\ell-1}$ 先看成一個 $B$ 進位的數 $c_0 B^{\ell-1} + c_1
B^{\ell-2} + \cdots + c_{\ell-1}$，再對 $m$ 取餘數。 用 Horner 法（第~4~章算多項式的
方法）從左往右算，每一步都取餘數：$h \gets (h \cdot B + c_i) \bmod m$，數字就不會變大。
$B$ 選一個比字元種類多的數（例如 $131$），$m$ 選質數。 只把各字元的碼加起來不行：
\texttt{"listen"} 和 \texttt{"silent"} 用的字母一樣，和也一樣。 第~27~章的
Rabin--Karp 字串比對會用到同一個式子。

\section{鏈結法}
\label{sec:hash-chain}

處理碰撞最直接的方法是讓每一格掛一條串列：第 $j$ 格存一條串列的開頭，所有
$h(k) = j$ 的鍵都放在這條串列上。 這叫做\textbf{鏈結法}（chaining，也叫
separate chaining）。

\begin{pseudo}[label=alg:hash-chain]{鏈結法的三個操作}
\Procedure{Chain-Insert}{$T, x$}
    \State 把節點 $x$ 插在串列 $T[h(x.\mathit{key})]$ 的最前面
\EndProcedure
\Function{Chain-Search}{$T, k$}
    \State $p \gets T[h(k)]$
    \While{$p \ne \NIL$ \AND $p.\mathit{key} \ne k$}
        \State $p \gets p.\mathit{next}$
    \EndWhile
    \State \Return $p$
\EndFunction
\Procedure{Chain-Delete}{$T, x$}
    \State 把節點 $x$ 從串列 $T[h(x.\mathit{key})]$ 拿掉
\EndProcedure
\end{pseudo}

插入是 $O(1)$：插在最前面不用走到尾端。 這假設 $x$ 的鍵還不在表裡；要檢查重複，就
得先做一次搜尋。 搜尋的時間和那條串列的長度成正比。 刪除時如果用的是第~5~章的雙向
串列，給了節點 $x$ 就能 $O(1)$ 拿掉；單向串列要先找到 $x$ 的前一個節點，時間和搜尋
一樣。 下面的 C 程式用單向串列，刪除時用「指標的指標」一邊找一邊記住要改的那個
\texttt{next}。

\subsection{C 實作}
\label{sec:hash-chain-c}

章首的刷卡機用鏈結法存「卡號 $\to$ 餘額」。 為了印得下，示範只開 $m = 7$ 格，放
$8$ 張卡，$\alpha = 8/7$。 鏈結法的 $\alpha$ 可以大於 $1$，串列只是平均長一點。

\begin{cbookrecap}[第~14~章 14.3、14.5 節，第~19~章 19.3、19.9、19.10 節]
\texttt{ChainTable t = \{calloc(7, sizeof(Entry *)), 7, 0\};} 用大括號依欄位的順序
初始化結構。 \texttt{calloc} 要一塊能放 $7$ 個指標的記憶體，而且全部設成 $0$；在常
見的系統上這就是 \texttt{NULL}，所以每條串列一開始都是空的。（標準沒有保證空指標的
位元全是 $0$；要完全不靠這一點，可以改用 \texttt{malloc} 再逐格設成 \texttt{NULL}。）

\texttt{Entry **pp} 存的是「某個 \texttt{Entry *} 變數的位址」：一開始是
\texttt{\&t->bucket[b]}，往後走一步變成 \texttt{\&(*pp)->next}。 找到要刪的節點時
寫 \texttt{*pp = dead->next}，不管它是串列的第一個還是中間的，都是同一行。 這是本書
第~\ref{sec:ll-singly}~節的寫法。
\end{cbookrecap}

\begin{example}[label=ex:hash-chain]{c}{(chain\_table.c) 鏈結法的雜湊表}
#include <stdio.h>
#include <stdlib.h>

// 鏈結法的雜湊表：卡號 -> 餘額。新的鍵插在串列的最前面。
typedef struct Entry {
    long long key;
    int value;
    struct Entry *next;
} Entry;

typedef struct {
    Entry **bucket;     // bucket[0..m-1] 是 m 條串列的開頭
    int m, n;           // 格數、鍵數
} ChainTable;

// 找 key，沒找到回傳 NULL。*cmp 記錄比較了幾個鍵。
Entry *ct_find(const ChainTable *t, long long key, int *cmp)
{
    *cmp = 0;
    Entry *p = t->bucket[key % t->m];      // h(k) = k mod m
    for (; p != NULL; p = p->next) {
        (*cmp)++;
        if (p->key == key) {
            return p;
        }
    }
    return NULL;
}

// 已經有 key 就更新值，否則插在串列最前面。記憶體不夠時回傳 0。
int ct_put(ChainTable *t, long long key, int value)
{
    int cmp;
    Entry *p = ct_find(t, key, &cmp);
    if (p == NULL) {
        p = malloc(sizeof *p);
        if (p == NULL) {
            return 0;
        }
        p->key = key;
        p->next = t->bucket[key % t->m];
        t->bucket[key % t->m] = p;
        t->n++;
    }
    p->value = value;
    return 1;
}

// 刪除 key，成功回傳 1。用指標的指標，刪第一個節點不用另外處理。
int ct_delete(ChainTable *t, long long key)
{
    Entry **pp = &t->bucket[key % t->m];
    while (*pp != NULL && (*pp)->key != key) {
        pp = &(*pp)->next;
    }
    if (*pp == NULL) {
        return 0;
    }
    Entry *dead = *pp;
    *pp = dead->next;
    free(dead);
    t->n--;
    return 1;
}

void print_bucket(const ChainTable *t, int b)
{
    printf("T[%d]:", b);
    for (Entry *p = t->bucket[b]; p != NULL; p = p->next) {
        const char *arrow = (p == t->bucket[b]) ? "" : " ->";
        printf("%s %lld:%d", arrow, p->key, p->value);
    }
    printf("\n");
}

int main(void)
{
    const long long card[] = {417961974, 416402873, 417721181,
        413819238, 419364875, 410986908, 411465496, 412266090};
    const int balance[] = {350, 120, 80, 500, 45, 260, 15, 900};
    const long long ask[] = {417961974, 410986908, 414000002};
    ChainTable t = {calloc(7, sizeof(Entry *)), 7, 0};
    if (t.bucket == NULL) {
        return 1;
    }
    for (int i = 0; i < 8; i++) {
        if (!ct_put(&t, card[i], balance[i])) {
            printf("out of memory\n");
            break;
        }
    }
    for (int b = 0; b < t.m; b++) {
        print_bucket(&t, b);
    }
    for (int i = 0; i < 3; i++) {
        int cmp;
        Entry *p = ct_find(&t, ask[i], &cmp);
        int v = (p != NULL) ? p->value : -1;
        printf("get %lld: %d, cmp = %d\n", ask[i], v, cmp);
    }
    ct_put(&t, 417961974, 330);     // 已經有這張卡：只改餘額
    int first = ct_delete(&t, 417721181);
    int second = ct_delete(&t, 417721181);
    printf("delete twice: %d %d, n = %d\n", first, second, t.n);
    print_bucket(&t, 3);
    for (int b = 0; b < t.m; b++) { // 釋放：一直刪串列的第一個
        while (t.bucket[b] != NULL) {
            ct_delete(&t, t.bucket[b]->key);
        }
    }
    free(t.bucket);
    return 0;
}
\end{example}
\iosample{}{T[0]: 413819238:500 \\ T[1]: 411465496:15 \\ T[2]: \\ T[3]: 410986908:260 -> 417721181:80 -> 417961974:350 \\ T[4]: \\ T[5]: 412266090:900 -> 416402873:120 \\ T[6]: 419364875:45 \\ get 417961974: 350, cmp = 3 \\ get 410986908: 260, cmp = 1 \\ get 414000002: -1, cmp = 3 \\ delete twice: 1 0, n = 7 \\ T[3]: 410986908:260 -> 417961974:330}

前七行是放完八張卡的表，\texttt{T[3]} 那一行就是第 $3$ 格的串列。 第 $3$ 格的三張卡依放入的先後是
$417961974$、$417721181$、$410986908$，每次都插在最前面，所以順序反過來。 找
$417961974$ 要比 $3$ 次，找最後放進去的 $410986908$ 只要 $1$ 次；找不在表裡的
$414000002$ 時，它的 $h$ 也是 $3$，整條串列都要比過，比了 $3$ 次，印出 $-1$。 同
一張卡刪兩次，第二次回傳 $0$；最後一行是刪除、改餘額之後的第 $3$ 格。 程式結束前
用 \texttt{ct\_delete} 逐一刪掉每條串列的第一個節點，把記憶體還回去。 第 $2$、$4$ 格
是空的：$8$ 個鍵放進 $7$ 格，不會剛好一格一個。


\subsection{載入因子與期望長度}
\label{sec:hash-chain-analysis}

最壞的情況是 $n$ 個鍵全部落在同一格，搜尋要 $\Theta(n)$。 要說明鏈結法「平均」很
快，得先假設雜湊值怎麼分布。

\begin{definition}[簡單均勻雜湊]
\label{def:hash-suh}
每個鍵的雜湊值都等可能是 $0, 1, \ldots, m-1$ 之中的任何一個，而且不同鍵的雜湊值互
相獨立。 這個假設叫做\textbf{簡單均勻雜湊}（simple uniform hashing）。
\end{definition}

這是一個關於「雜湊函數加上輸入」的假設，不是證出來的性質。 它的意思是：雜湊函數把
實際的輸入打散得像亂數一樣。 下面的推導要用到兩個工具。 事件 $E$ 的\textbf{指示變
數}（indicator variable）$I\{E\}$ 在 $E$ 發生時是 $1$，否則是 $0$，所以
$\mathrm{E}[I\{E\}] = 1 \cdot \Pr[E] + 0 \cdot (1 - \Pr[E]) = \Pr[E]$。 期望值是線性
的：$\mathrm{E}[X + Y] = \mathrm{E}[X] + \mathrm{E}[Y]$，不需要 $X$、$Y$ 獨立（第~1~章
算插入排序的平均挪動次數用過）。 把一個要數的量寫成很多個指示變數的和，再一項一項
取期望值，是這一章最常用的手法。

\begin{theorem}
\label{thm:hash-chain-len}
在簡單均勻雜湊下，$n$ 個鍵放進 $m$ 格：
\begin{enumerate}
\item 每一條串列的期望長度都是 $\alpha = n/m$。
\item 搜尋一個不在表裡的鍵，期望比較 $\alpha$ 個鍵。
\item 搜尋一個在表裡的鍵（$n$ 個鍵被搜尋的機會一樣，新的鍵插在最前面），期望比較
\[
1 + \frac{n - 1}{2m} = 1 + \frac{\alpha}{2} - \frac{\alpha}{2n}
\]
個鍵。
\end{enumerate}
\end{theorem}

\begin{proof}
把鍵依放入的順序叫做 $x_1, x_2, \ldots, x_n$。

(1) 固定第 $j$ 格。 第 $j$ 條串列的長度是 $L_j = \sum_{i=1}^{n} I\{h(x_i) = j\}$，每
一項的期望值是 $1/m$，所以 $\mathrm{E}[L_j] = n/m$。

(2) 要找的鍵 $k$ 不在表裡，搜尋會把第 $h(k)$ 條串列整條比完。 被比到的鍵個數是
$\sum_{i=1}^{n} I\{h(x_i) = h(k)\}$。 $h(k)$ 和 $h(x_i)$ 獨立而且都均勻，所以
\[
\Pr[h(x_i) = h(k)] = \sum_{j=0}^{m-1} \Pr[h(x_i) = j] \Pr[h(k) = j]
= m \cdot \frac{1}{m^2} = \frac{1}{m},
\]
期望值是 $n/m = \alpha$。

(3) 新的鍵插在最前面，所以找 $x_i$ 時，比到的是 $x_i$ 自己，加上排在它前面、也就是
比它晚放進來而且落在同一格的鍵。 比較次數是
\[
C_i = 1 + \sum_{j = i+1}^{n} I\{h(x_j) = h(x_i)\}, \qquad
\mathrm{E}[C_i] = 1 + \frac{n - i}{m},
\]
這裡用了 (2) 的同一個計算：兩個不同的鍵落在同一格的機率是 $1/m$。 每個鍵被找的機
會是 $1/n$，所以
\[
\frac{1}{n} \sum_{i=1}^{n} \left(1 + \frac{n - i}{m}\right)
= 1 + \frac{1}{nm} \sum_{i=1}^{n} (n - i)
= 1 + \frac{1}{nm} \cdot \frac{n(n-1)}{2}
= 1 + \frac{n-1}{2m}. \qedhere
\]
\end{proof}

再加上算 $h(k)$ 和走到那一格的 $O(1)$，成功和失敗的搜尋都是期望 $\Theta(1 +
\alpha)$。 只要表的格數和鍵數成正比（$n = O(m)$，例如規定 $\alpha \le 1$），$\alpha$
就是常數，三個操作的期望時間都是 $O(1)$。 這就是章首問題的答案：查詢的次數只和
$\alpha$ 有關，和全校有幾張卡無關。

另外，(1) 說的是「隨便指定一格」的期望長度；「某個已經在表裡的鍵所在的那條串列」期望
長度是 $1 + (n-1)/m$，比 $\alpha$ 大，因為長的串列裝了比較多鍵。 期望長度是
$\alpha$ 也不代表每條都差不多長：$\alpha = 1$ 時大約有 $(1 - 1/m)^n \approx e^{-1}
\approx 37\%$ 的格子是空的（觀念題第 8 題）。

\section{開放定址}
\label{sec:hash-open}

鏈結法每個鍵要多一個指標，每次插入要 \texttt{malloc} 一個節點。 \textbf{開放定址}
（open addressing）把所有鍵直接放在表 $T[0..m-1]$ 裡，每格最多一個鍵，所以
$n \le m$、$\alpha \le 1$。 碰撞時照一個固定的順序去試別的格子，叫做\textbf{探測}
（probe）。 每個鍵 $k$ 的探測順序是
\[
h(k, 0),\ h(k, 1),\ \ldots,\ h(k, m - 1),
\]
叫做 $k$ 的\textbf{探測序列}（probe sequence），$i$ 是第幾次探測，從 $0$ 開始。
插入時沿著探測序列找第一個空格；搜尋時沿著同一個序列找，遇到 $k$ 就是找到了，遇到
空格就表示 $k$ 不在表裡：如果 $k$ 在表裡，插入它時一定會停在這個空格或更前面。 為
了讓表不滿時一定插得進去，探測序列最好是 $0, 1, \ldots, m-1$ 的一個排列，每一格都
會被試到。

\begin{pseudo}[label=alg:hash-open]{開放定址的插入與搜尋（還沒有刪除）}
\Function{Open-Insert}{$T, k$}
    \For{$i \gets 0$ \To $m - 1$}
        \State $j \gets h(k, i)$
        \If{$T[j] = \NIL$}
            \State $T[j] \gets k$
            \State \Return $j$
        \EndIf
    \EndFor
    \State \Error{表已滿}
\EndFunction
\Function{Open-Search}{$T, k$}
    \For{$i \gets 0$ \To $m - 1$}
        \State $j \gets h(k, i)$
        \If{$T[j] = k$}
            \State \Return $j$
        \ElsIf{$T[j] = \NIL$}
            \State \Return \NIL
        \EndIf
    \EndFor
    \State \Return \NIL
\EndFunction
\end{pseudo}

一次插入或搜尋的成本就是\textbf{探測次數}：看了幾格，包括讓它停下來的那一格。 下面
三種探測法的差別只在 $h(k, i)$ 怎麼算。 三種都用同一組例子：停車場發的停車證號碼，
$m = 13$，$h(k) = k \bmod 13$，依序插入
\[
5347,\ 3051,\ 6296,\ 1682,\ 2014,\ 8170,\ 2299,\ 5906,\ 7193.
\]
它們的 $h(k)$ 依序是 $4, 9, 4, 5, 12, 6, 11, 4, 4$。 放完之後
$\alpha = 9/13 \approx 0.69$。

\subsection{線性探測}

\textbf{線性探測}（linear probing）從起始位置開始一格一格往後試，到尾端繞回 $0$：
\[
h(k, i) = (h(k) + i) \bmod m.
\]
圖~\ref{fig:hash-linear} 每一列是插入一個鍵之後的表。 橘框是這次放進去的格子，粉紅底是這次探測時看過、
但已經有鍵的格子。

\begin{figure}[tbp]
\centering
\begin{tikzpicture}
\foreach \j in {0,...,12} { \node[idxlbl] at ({0.6+\j*0.92}, 0.55) {\j}; }
\node[anchor=east, font=\footnotesize] at (0.2,0.00) {放 5347，1 次};
\node[hcell] at (0.60,0.00) {};
\node[hcell] at (1.52,0.00) {};
\node[hcell] at (2.44,0.00) {};
\node[hcell] at (3.36,0.00) {};
\node[hcell, hlnew] at (4.28,0.00) {5347};
\node[hcell] at (5.20,0.00) {};
\node[hcell] at (6.12,0.00) {};
\node[hcell] at (7.04,0.00) {};
\node[hcell] at (7.96,0.00) {};
\node[hcell] at (8.88,0.00) {};
\node[hcell] at (9.80,0.00) {};
\node[hcell] at (10.72,0.00) {};
\node[hcell] at (11.64,0.00) {};
\node[anchor=east, font=\footnotesize] at (0.2,-0.78) {放 3051，1 次};
\node[hcell] at (0.60,-0.78) {};
\node[hcell] at (1.52,-0.78) {};
\node[hcell] at (2.44,-0.78) {};
\node[hcell] at (3.36,-0.78) {};
\node[hcell] at (4.28,-0.78) {5347};
\node[hcell] at (5.20,-0.78) {};
\node[hcell] at (6.12,-0.78) {};
\node[hcell] at (7.04,-0.78) {};
\node[hcell] at (7.96,-0.78) {};
\node[hcell, hlnew] at (8.88,-0.78) {3051};
\node[hcell] at (9.80,-0.78) {};
\node[hcell] at (10.72,-0.78) {};
\node[hcell] at (11.64,-0.78) {};
\node[anchor=east, font=\footnotesize] at (0.2,-1.56) {放 6296，2 次};
\node[hcell] at (0.60,-1.56) {};
\node[hcell] at (1.52,-1.56) {};
\node[hcell] at (2.44,-1.56) {};
\node[hcell] at (3.36,-1.56) {};
\node[hcell, hlprobe] at (4.28,-1.56) {5347};
\node[hcell, hlnew] at (5.20,-1.56) {6296};
\node[hcell] at (6.12,-1.56) {};
\node[hcell] at (7.04,-1.56) {};
\node[hcell] at (7.96,-1.56) {};
\node[hcell] at (8.88,-1.56) {3051};
\node[hcell] at (9.80,-1.56) {};
\node[hcell] at (10.72,-1.56) {};
\node[hcell] at (11.64,-1.56) {};
\node[anchor=east, font=\footnotesize] at (0.2,-2.34) {放 1682，2 次};
\node[hcell] at (0.60,-2.34) {};
\node[hcell] at (1.52,-2.34) {};
\node[hcell] at (2.44,-2.34) {};
\node[hcell] at (3.36,-2.34) {};
\node[hcell] at (4.28,-2.34) {5347};
\node[hcell, hlprobe] at (5.20,-2.34) {6296};
\node[hcell, hlnew] at (6.12,-2.34) {1682};
\node[hcell] at (7.04,-2.34) {};
\node[hcell] at (7.96,-2.34) {};
\node[hcell] at (8.88,-2.34) {3051};
\node[hcell] at (9.80,-2.34) {};
\node[hcell] at (10.72,-2.34) {};
\node[hcell] at (11.64,-2.34) {};
\node[anchor=east, font=\footnotesize] at (0.2,-3.12) {放 2014，1 次};
\node[hcell] at (0.60,-3.12) {};
\node[hcell] at (1.52,-3.12) {};
\node[hcell] at (2.44,-3.12) {};
\node[hcell] at (3.36,-3.12) {};
\node[hcell] at (4.28,-3.12) {5347};
\node[hcell] at (5.20,-3.12) {6296};
\node[hcell] at (6.12,-3.12) {1682};
\node[hcell] at (7.04,-3.12) {};
\node[hcell] at (7.96,-3.12) {};
\node[hcell] at (8.88,-3.12) {3051};
\node[hcell] at (9.80,-3.12) {};
\node[hcell] at (10.72,-3.12) {};
\node[hcell, hlnew] at (11.64,-3.12) {2014};
\node[anchor=east, font=\footnotesize] at (0.2,-3.90) {放 8170，2 次};
\node[hcell] at (0.60,-3.90) {};
\node[hcell] at (1.52,-3.90) {};
\node[hcell] at (2.44,-3.90) {};
\node[hcell] at (3.36,-3.90) {};
\node[hcell] at (4.28,-3.90) {5347};
\node[hcell] at (5.20,-3.90) {6296};
\node[hcell, hlprobe] at (6.12,-3.90) {1682};
\node[hcell, hlnew] at (7.04,-3.90) {8170};
\node[hcell] at (7.96,-3.90) {};
\node[hcell] at (8.88,-3.90) {3051};
\node[hcell] at (9.80,-3.90) {};
\node[hcell] at (10.72,-3.90) {};
\node[hcell] at (11.64,-3.90) {2014};
\node[anchor=east, font=\footnotesize] at (0.2,-4.68) {放 2299，1 次};
\node[hcell] at (0.60,-4.68) {};
\node[hcell] at (1.52,-4.68) {};
\node[hcell] at (2.44,-4.68) {};
\node[hcell] at (3.36,-4.68) {};
\node[hcell] at (4.28,-4.68) {5347};
\node[hcell] at (5.20,-4.68) {6296};
\node[hcell] at (6.12,-4.68) {1682};
\node[hcell] at (7.04,-4.68) {8170};
\node[hcell] at (7.96,-4.68) {};
\node[hcell] at (8.88,-4.68) {3051};
\node[hcell] at (9.80,-4.68) {};
\node[hcell, hlnew] at (10.72,-4.68) {2299};
\node[hcell] at (11.64,-4.68) {2014};
\node[anchor=east, font=\footnotesize] at (0.2,-5.46) {放 5906，5 次};
\node[hcell] at (0.60,-5.46) {};
\node[hcell] at (1.52,-5.46) {};
\node[hcell] at (2.44,-5.46) {};
\node[hcell] at (3.36,-5.46) {};
\node[hcell, hlprobe] at (4.28,-5.46) {5347};
\node[hcell, hlprobe] at (5.20,-5.46) {6296};
\node[hcell, hlprobe] at (6.12,-5.46) {1682};
\node[hcell, hlprobe] at (7.04,-5.46) {8170};
\node[hcell, hlnew] at (7.96,-5.46) {5906};
\node[hcell] at (8.88,-5.46) {3051};
\node[hcell] at (9.80,-5.46) {};
\node[hcell] at (10.72,-5.46) {2299};
\node[hcell] at (11.64,-5.46) {2014};
\node[anchor=east, font=\footnotesize] at (0.2,-6.24) {放 7193，7 次};
\node[hcell] at (0.60,-6.24) {};
\node[hcell] at (1.52,-6.24) {};
\node[hcell] at (2.44,-6.24) {};
\node[hcell] at (3.36,-6.24) {};
\node[hcell, hlprobe] at (4.28,-6.24) {5347};
\node[hcell, hlprobe] at (5.20,-6.24) {6296};
\node[hcell, hlprobe] at (6.12,-6.24) {1682};
\node[hcell, hlprobe] at (7.04,-6.24) {8170};
\node[hcell, hlprobe] at (7.96,-6.24) {5906};
\node[hcell, hlprobe] at (8.88,-6.24) {3051};
\node[hcell, hlnew] at (9.80,-6.24) {7193};
\node[hcell] at (10.72,-6.24) {2299};
\node[hcell] at (11.64,-6.24) {2014};
\end{tikzpicture}
\diagcaption{線性探測依序插入九張停車證（$m = 13$）。 左邊是插入的鍵和探測次數。}
\label{fig:hash-linear}
\end{figure}

前七個鍵都只試了一兩次。 第 $8$ 個鍵 $5906$ 的起始位置是 $4$，這時第 $4$ 到 $7$ 格已
經連成一整段，它要一路試到第 $8$ 格，共 $5$ 次。 第 $9$ 個鍵 $7193$ 也從 $4$ 開始，
這段又多了第 $8$ 格，再往後還接上本來就有鍵的第 $9$ 格，它試了 $7$ 次。 九個鍵一共
探測 $22$ 次，平均每個 $22/9 \approx 2.44$ 次。 之後搜尋這九個鍵，每個鍵的探測次數
和它當初插入時一樣，所以成功搜尋的平均也是 $2.44$ 次。

失敗的搜尋呢？ 如果要找的鍵起始位置是 $0$ 到 $3$，第一格就是空的，$1$ 次就停；起始
位置是 $4$，要把第 $4$ 到 $12$ 格看完、再看第 $0$ 格，共 $10$ 次；起始位置是 $12$，
看第 $12$、$0$ 格，共 $2$ 次。 假設起始位置是 $0$ 到 $12$ 的機會一樣，平均是
\[
\frac{4 \cdot 1 + (10 + 9 + 8 + 7 + 6 + 5 + 4 + 3 + 2)}{13} = \frac{58}{13} \approx 4.46.
\]

這種連成一段、越來越長的現象叫做\textbf{主要群聚}（primary clustering）。 原因很簡
單：一段連續的 $L$ 格後面接著一個空格，起始位置落在這 $L$ 格或那個空格本身的鍵都會
被放進那個空格，機率大約是 $(L + 1)/m$；段越長，下一個鍵越可能接在它後面，讓它更長。
兩段之間只隔一格時，那一格一填上，兩段就合成一段，上面第 $4$ 到 $12$ 格就是這樣接起
來的。

\subsection{二次探測}

\textbf{二次探測}（quadratic probing）讓跳的距離越來越大：
\[
h(k, i) = (h(k) + i^2) \bmod m,
\]
從起始位置依序跳 $0, 1, 4, 9, 16, \ldots$ 格（都取 $\bmod m$）。 起始位置不同的鍵，
就算碰上同一格，下一步也會跳到不同的地方，不會跟著同一段一路往後走，所以沒有主要群
聚。 但起始位置\textbf{相同}的鍵，探測序列整條一模一樣，例子裡 $6296$、$5906$、
$7193$ 都從 $4$ 開始，依序試 $4, 5, 8, 0$，後一個總是把前一個試過的路再走一遍。 這
叫做\textbf{次要群聚}（secondary clustering）。

二次探測有一個新問題：$i^2 \bmod m$ 不一定走得遍所有的格子。 例如 $m = 8$ 時
$i^2 \bmod 8$ 只有 $0, 1, 4$ 三種值，從起始位置出發只到得了三格，表還有空位也插不
進去（觀念題第 6 題）。 $m$ 是質數時，至少前一半是不重複的。

\begin{theorem}
\label{thm:hash-quad}
$m$ 是奇質數時，二次探測的前 $(m+1)/2$ 個位置
$h(k, 0), h(k, 1), \ldots, h(k, (m-1)/2)$ 兩兩不同。 因此只要表裡的鍵少於
$(m+1)/2$ 個（例如 $\alpha < 1/2$），插入一定成功。
\end{theorem}

\begin{proof}
假設 $0 \le i < j \le (m-1)/2$ 而 $h(k, i) = h(k, j)$。 兩邊減掉 $h(k)$，得到
$i^2 \equiv j^2 \pmod m$，也就是 $m$ 整除 $j^2 - i^2 = (j - i)(j + i)$。 $m$ 是質數，
所以 $m$ 整除 $j - i$ 或 $j + i$。 但 $1 \le j - i < m$，而且
$1 \le j + i \le m - 2 < m$，兩個都不可能被 $m$ 整除，矛盾。 前 $(m+1)/2$ 個位置兩
兩不同；表裡的鍵少於 $(m+1)/2$ 個時，這些位置裡至少有一個是空的，插入會停在那裡。
\end{proof}

這個界是緊的：$i$ 和 $m - i$ 給出同一個 $i^2 \bmod m$，所以後一半只是把前一半倒過來
再走一次，二次探測最多只看得到 $(m+1)/2$ 個不同的格子。 有兩種常見的改法讓它走遍整
張表：$m = 2^p$ 時改跳三角數 $h(k) + i(i+1)/2$（觀念題第 7 題證明它走得遍）；或是正
負交替地跳 $h(k) + 1^2, h(k) - 1^2, h(k) + 2^2, h(k) - 2^2, \ldots$，可以證明 $m$ 是
除以 $4$ 餘 $3$ 的質數時，這樣會走遍全部 $m$ 格。

\subsection{雙重雜湊}

\textbf{雙重雜湊}（double hashing）用第二個雜湊函數決定步長：
\[
h(k, i) = (h_1(k) + i \cdot h_2(k)) \bmod m.
\]
兩個鍵就算 $h_1$ 相同，$h_2$ 通常不同，探測序列在第二步就分開了，所以主要和次要群
聚都沒有。 $h_2(k)$ 不能是 $0$（會原地不動），而且要和 $m$ 互質。

\begin{theorem}
\label{thm:hash-double}
如果 $\gcd(h_2(k), m) = 1$，雙重雜湊的探測序列 $h(k, 0), \ldots, h(k, m-1)$ 是
$0, 1, \ldots, m-1$ 的一個排列。
\end{theorem}

\begin{proof}
假設 $0 \le i < j \le m - 1$ 而 $h(k, i) = h(k, j)$，則 $m$ 整除 $(j - i)\,h_2(k)$。
$m$ 和 $h_2(k)$ 互質，所以 $m$ 整除 $j - i$，但 $1 \le j - i < m$，矛盾。 $m$ 個位置
兩兩不同，就是全部的 $m$ 格。
\end{proof}

常見的兩種取法：$m$ 取質數，$h_2(k)$ 取 $1$ 到 $m - 1$ 之間的值；或是 $m = 2^p$，
$h_2(k)$ 一律取奇數。 例子裡 $m = 13$，取 $h_2(k) = 1 + (k \bmod 7)$，值在 $1$ 到
$7$ 之間，都和 $13$ 互質。

下表把三種探測法放在一起。 每一格是那個鍵插入時依序試過的格子，最後一個數就是它放
進去的位置。

\begin{center}
\begin{tabular}{rcclll}
\toprule
鍵 & $h_1$ & $h_2$ & 線性探測 & 二次探測 & 雙重雜湊 \\
\midrule
5347 & 4 & 7 & 4 & 4 & 4 \\
3051 & 9 & 7 & 9 & 9 & 9 \\
6296 & 4 & 4 & 4 5 & 4 5 & 4 8 \\
1682 & 5 & 3 & 5 6 & 5 6 & 5 \\
2014 & 12 & 6 & 12 & 12 & 12 \\
8170 & 6 & 2 & 6 7 & 6 7 & 6 \\
2299 & 11 & 4 & 11 & 11 & 11 \\
5906 & 4 & 6 & 4 5 6 7 8 & 4 5 8 & 4 10 \\
7193 & 4 & 5 & 4 5 6 7 8 9 10 & 4 5 8 0 & 4 9 1 \\
\midrule
探測總次數 & & & 22 & 17 & 13 \\
\bottomrule
\end{tabular}
\end{center}

前七個鍵三種方法幾乎一樣，差別出在最後兩個起始位置是 $4$ 的鍵。 線性探測被第 $4$
到 $10$ 格那一整段拖住；二次探測跳開了，但 $5906$、$7193$ 都重走 $4, 5, 8$；雙重雜
湊的 $7193$ 步長是 $5$，第二步就跳到第 $9$ 格，和 $6296$（步長 $4$）、$5906$（步長
$6$）走的路都不一樣。 三種方法最後的表在範例~\ref{ex:hash-open} 的輸出裡。


用「一共有幾種不同的探測序列」來比較也很清楚。 線性探測和二次探測的整條序列只由
起始位置決定，各只有 $m$ 種；雙重雜湊由 $(h_1(k), h_2(k))$ 決定，$m$ 是質數時有
$m(m-1)$ 種。 第~\ref{sec:hash-analysis}~節的理想模型允許全部 $m!$ 種排列，雙重雜湊
是三種方法裡最接近的。

\begin{examnote}
探測的追蹤題，每個鍵寫一列：$h(k)$（雙重雜湊加寫 $h_2(k)$）、依序試過的格子、最後
放的位置，最後畫出整張表。 先看題目的 $i$ 從 $0$ 還是從 $1$ 開始數，以及「探測次數」
包不包括第一次看起始位置；本章的寫法是 $i$ 從 $0$ 開始、看過幾格就算幾次。 二次探
測要寫出 $i^2 \bmod m$ 的值再加，不要心算跳過，繞回 $0$ 的地方最容易錯。
\end{examnote}

\subsection{刪除：刪除標記}
\label{sec:hash-tomb}

開放定址的刪除不能把格子改回空的。 在線性探測的表裡刪掉 $1682$（第 $6$ 格），如果
把第 $6$ 格清空，再找 $5906$：從第 $4$ 格開始，看到第 $6$ 格是空的就停下，回答「不在
表裡」。 但 $5906$ 明明在第 $8$ 格，它當初是越過第 $6$ 格才放進去的。 清空一格會把
經過它的探測序列從中間切斷。

做法是把刪掉的格子標成第三種狀態 \hashDEL，叫做\textbf{刪除標記}，也叫
墓碑（tombstone）。 搜尋時把它當成有鍵的格子，繼續往後找；插入時把它當成
空格，可以放新的鍵。 \proc{Open-Search} 一個字都不用改：\hashDEL{} 既不等於 $k$ 也不
是 \NIL，虛擬碼~\ref{alg:hash-open} 的第 11--14 行自然會跳過它。 要改的是插入和
刪除。

\begin{pseudo}[label=alg:hash-tomb]{有刪除標記的插入與刪除}
\Function{Mark-Insert}{$T, k$}
    \State $a \gets \NIL$ \Comment{第一個可以放的格子}
    \For{$i \gets 0$ \To $m - 1$}
        \State $j \gets h(k, i)$
        \If{$T[j] = k$}
            \State \Return $j$ \Comment{已經在表裡}
        \EndIf
        \If{($T[j] = \NIL$ \OR $T[j] = \hashDEL$) \AND $a = \NIL$}
            \State $a \gets j$
        \EndIf
        \If{$T[j] = \NIL$}
            \State \Break
        \EndIf
    \EndFor
    \If{$a = \NIL$}
        \State \Error{表已滿}
    \EndIf
    \State $T[a] \gets k$
    \State \Return $a$
\EndFunction
\Procedure{Mark-Delete}{$T, k$}
    \State $j \gets$ \Call{Open-Search}{$T, k$}
    \If{$j \ne \NIL$}
        \State $T[j] \gets \hashDEL$
    \EndIf
\EndProcedure
\end{pseudo}

\proc{Mark-Insert} 遇到第一個刪除標記時先記在 $a$（第 7、8 行），但不能馬上放：$k$ 可
能就在更後面，直接放會讓表裡有兩個 $k$。 所以它繼續往後找，直到遇到空格（第 9、10
行）或找到 $k$（第 5、6 行）才停，最後把 $k$ 放在第一個可以放的格子。 回到線性探測
的例子，刪除 $1682$ 之後做三件事，見圖~\ref{fig:hash-tomb}。

\begin{figure}[tbp]
\centering
\begin{tikzpicture}
\foreach \j in {0,...,12} { \node[idxlbl] at ({0.6+\j*0.92}, 0.55) {\j}; }
\node[anchor=east, font=\footnotesize] at (0.2,0.00) {刪除 1682};
\node[hcell] at (0.60,0.00) {};
\node[hcell] at (1.52,0.00) {};
\node[hcell] at (2.44,0.00) {};
\node[hcell] at (3.36,0.00) {};
\node[hcell] at (4.28,0.00) {5347};
\node[hcell, hlprobe] at (5.20,0.00) {6296};
\node[hcell, hldead] at (6.12,0.00) {$\times$};
\node[hcell] at (7.04,0.00) {8170};
\node[hcell] at (7.96,0.00) {5906};
\node[hcell] at (8.88,0.00) {3051};
\node[hcell] at (9.80,0.00) {7193};
\node[hcell] at (10.72,0.00) {2299};
\node[hcell] at (11.64,0.00) {2014};
\node[anchor=east, font=\footnotesize] at (0.2,-0.78) {找 5906，5 次};
\node[hcell] at (0.60,-0.78) {};
\node[hcell] at (1.52,-0.78) {};
\node[hcell] at (2.44,-0.78) {};
\node[hcell] at (3.36,-0.78) {};
\node[hcell, hlprobe] at (4.28,-0.78) {5347};
\node[hcell, hlprobe] at (5.20,-0.78) {6296};
\node[hcell, hlprobe] at (6.12,-0.78) {$\times$};
\node[hcell, hlprobe] at (7.04,-0.78) {8170};
\node[hcell, hlprobe] at (7.96,-0.78) {5906};
\node[hcell] at (8.88,-0.78) {3051};
\node[hcell] at (9.80,-0.78) {7193};
\node[hcell] at (10.72,-0.78) {2299};
\node[hcell] at (11.64,-0.78) {2014};
\node[anchor=east, font=\footnotesize] at (0.2,-1.56) {找 1682，9 次};
\node[hcell, hlprobe] at (0.60,-1.56) {};
\node[hcell] at (1.52,-1.56) {};
\node[hcell] at (2.44,-1.56) {};
\node[hcell] at (3.36,-1.56) {};
\node[hcell] at (4.28,-1.56) {5347};
\node[hcell, hlprobe] at (5.20,-1.56) {6296};
\node[hcell, hlprobe] at (6.12,-1.56) {$\times$};
\node[hcell, hlprobe] at (7.04,-1.56) {8170};
\node[hcell, hlprobe] at (7.96,-1.56) {5906};
\node[hcell, hlprobe] at (8.88,-1.56) {3051};
\node[hcell, hlprobe] at (9.80,-1.56) {7193};
\node[hcell, hlprobe] at (10.72,-1.56) {2299};
\node[hcell, hlprobe] at (11.64,-1.56) {2014};
\node[anchor=east, font=\footnotesize] at (0.2,-2.34) {放 4880，9 次};
\node[hcell, hlprobe] at (0.60,-2.34) {};
\node[hcell] at (1.52,-2.34) {};
\node[hcell] at (2.44,-2.34) {};
\node[hcell] at (3.36,-2.34) {};
\node[hcell] at (4.28,-2.34) {5347};
\node[hcell, hlprobe] at (5.20,-2.34) {6296};
\node[hcell, hlnew] at (6.12,-2.34) {4880};
\node[hcell, hlprobe] at (7.04,-2.34) {8170};
\node[hcell, hlprobe] at (7.96,-2.34) {5906};
\node[hcell, hlprobe] at (8.88,-2.34) {3051};
\node[hcell, hlprobe] at (9.80,-2.34) {7193};
\node[hcell, hlprobe] at (10.72,-2.34) {2299};
\node[hcell, hlprobe] at (11.64,-2.34) {2014};
\end{tikzpicture}
\diagcaption{刪除標記（$\times$）。 找 $5906$ 時越過第 $6$ 格的標記，在第 $8$ 格找到；
找已經刪掉的 $1682$ 要一路看到第 $0$ 格的空格才停；放 $4880$（$h = 5$）時記住第 $6$
格，確定後面沒有 $4880$ 之後才放回第 $6$ 格。}
\label{fig:hash-tomb}
\end{figure}

\begin{examnote}
刪除標記的題目，先看清楚題目規定插入時遇到標記怎麼做。 本章（以及多數教科書）的規
則是「繼續找到空格或找到 $k$ 為止，再放進第一個標記或空格」；有些題目為了簡化，直
接放進第一個遇到的標記。 兩種規則放的位置可能一樣，探測次數卻不同，作答時把採用的
規則寫一句在答案前面。
\end{examnote}

刪除標記讓刪除變得正確，但它不會自己消失。 插入、刪除做久了，表裡可能只剩少數幾個鍵，
卻有很多格是刪除標記；失敗的搜尋要一直走到真正的空格才停，速度由「鍵加上標記」佔的
比例決定，而不是 $\alpha$。 常見的處理是記下標記的個數 $d$，當 $(n + d)/m$ 超過某個
門檻時，把所有的鍵重新插入一張乾淨的表（第~\ref{sec:hash-dynamic}~節的再雜湊）。
刪除和放進標記格的插入都不會改變 $n + d$，只有放進真正空格的插入會讓它加一；只要
重建後的表夠空（例如 $(n + d)/m$ 不超過門檻的一半），下一次重建之前至少要再插入
$\Theta(m)$ 個鍵，重建的 $O(m)$ 平攤到這些插入上，每次 $O(1)$。 第~3~章觀念題「刪除
時先只做標記、標記多了再一起整理」用的是同一類論證。

\subsection{C 實作}
\label{sec:hash-open-c}

下面的程式把三種探測法寫在同一個函式 \texttt{probe} 裡，用 \texttt{mode} 選擇；每格
有 \texttt{EMPTY}、\texttt{FULL}、\texttt{DELETED} 三種狀態。 \texttt{search} 就
是虛擬碼的 \proc{Open-Search} 加上 \proc{Mark-Insert} 要的「第一個可以放的格子」，
\texttt{operate} 的插入、搜尋、刪除都先呼叫它，再依結果放鍵或立標記。 \texttt{main} 把這一節的九張停車證用三種方法各插一次，
線性探測再做一次刪除、搜尋和插入。 線性探測的每個操作印一行：\texttt{i}、\texttt{f}、
\texttt{d} 是插入、搜尋、刪除，接著是試過的格子、結果的格子（\texttt{-1} 表示沒找
到）和探測次數，可以和圖~\ref{fig:hash-linear}、圖~\ref{fig:hash-tomb} 逐格對照；另
外兩種方法只印最後的表，和上一小節的三列表對照。

\begin{example}[label=ex:hash-open]{c}{(open\_addr.c) 開放定址：三種探測法與刪除標記}
#include <stdio.h>
#include <stdlib.h>

// 開放定址的雜湊表（鍵是非負整數），三種探測法，刪除用標記。
enum { EMPTY, FULL, DELETED };            // 每一格的三種狀態
enum { LINEAR, QUADRATIC, DOUBLE };       // 探測法
static int trace = 1;

typedef struct {
    int *key;
    unsigned char *state;
    int m, mode;
} OpenTable;

// 第 i 次探測的格子，i 從 0 開始。h2(k) = 1 + k mod 7。
int probe(const OpenTable *t, int key, long long i)
{
    long long step = i;                           // 線性探測
    if (t->mode == QUADRATIC) {
        step = i * i;
    } else if (t->mode == DOUBLE) {
        step = i * (1 + key % 7);
    }
    return (int)((key % t->m + step) % t->m);
}

// 沿探測序列找 key，回傳它所在的格子，沒找到回傳 -1。
// *avail 設成第一個可以放 key 的格子（第一個標記，
// 或停下來的空格），沒有就是 -1。trace 是 1 時印出看過的格子。
int search(const OpenTable *t, int key, int *avail, int *probes)
{
    *avail = -1;
    for (int i = 0; i < t->m; i++) {
        int j = probe(t, key, i);
        *probes = i + 1;
        if (trace) {
            printf(" %d", j);
        }
        if (t->state[j] == FULL && t->key[j] == key) {
            return j;
        }
        if (t->state[j] != FULL && *avail < 0) {
            *avail = j;
        }
        if (t->state[j] == EMPTY) {
            return -1;
        }
    }
    return -1;
}

// op 是 'i' 插入、'f' 搜尋、'd' 刪除，每個操作都先搜尋。
void operate(OpenTable *t, char op, int key)
{
    int avail, probes;
    if (trace) {
        printf("%c %d:", op, key);
    }
    int j = search(t, key, &avail, &probes);
    if (op == 'i' && j < 0 && avail >= 0) {
        t->key[avail] = key;
        t->state[avail] = FULL;
        j = avail;
    } else if (op == 'd' && j >= 0) {
        t->state[j] = DELETED;      // 立標記，不能改回 EMPTY
    }
    if (trace) {
        printf(" -> %d, probes = %d\n", j, probes);
    }
}

void print_table(const OpenTable *t)
{
    printf("table:");
    for (int j = 0; j < t->m; j++) {
        if (t->state[j] == FULL) {
            printf(" %d", t->key[j]);
        } else {
            printf(" %s", t->state[j] == EMPTY ? "-" : "D");
        }
    }
    printf("\n");
}

int main(void)
{
    const int permit[] = {5347, 3051, 6296, 1682, 2014,
                          8170, 2299, 5906, 7193};
    const char *name[] = {"linear", "quadratic", "double"};
    for (int mode = LINEAR; mode <= DOUBLE; mode++) {
        OpenTable t = {malloc(13 * sizeof(int)), calloc(13, 1),
                       13, mode};    // calloc：全部是 EMPTY
        if (t.key != NULL && t.state != NULL) {
            trace = (mode == LINEAR);   // 只印線性探測的過程
            printf("== %s\n", name[mode]);
            for (int i = 0; i < 9; i++) {
                operate(&t, 'i', permit[i]);
            }
            print_table(&t);
            if (mode == LINEAR) {       // 刪除、搜尋、再插入
                operate(&t, 'd', 1682);
                operate(&t, 'f', 5906);
                operate(&t, 'f', 1682);
                operate(&t, 'i', 4880);
                print_table(&t);
            }
        }
        free(t.key);
        free(t.state);
    }
    return 0;
}
\end{example}
\iosample{}{== linear \\ i 5347: 4 -> 4, probes = 1 \\ i 3051: 9 -> 9, probes = 1 \\ i 6296: 4 5 -> 5, probes = 2 \\ i 1682: 5 6 -> 6, probes = 2 \\ i 2014: 12 -> 12, probes = 1 \\ i 8170: 6 7 -> 7, probes = 2 \\ i 2299: 11 -> 11, probes = 1 \\ i 5906: 4 5 6 7 8 -> 8, probes = 5 \\ i 7193: 4 5 6 7 8 9 10 -> 10, probes = 7 \\ table: - - - - 5347 6296 1682 8170 5906 3051 7193 2299 2014 \\ d 1682: 5 6 -> 6, probes = 2 \\ f 5906: 4 5 6 7 8 -> 8, probes = 5 \\ f 1682: 5 6 7 8 9 10 11 12 0 -> -1, probes = 9 \\ i 4880: 5 6 7 8 9 10 11 12 0 -> 6, probes = 9 \\ table: - - - - 5347 6296 4880 8170 5906 3051 7193 2299 2014 \\ == quadratic \\ table: 7193 - - - 5347 6296 1682 8170 5906 3051 - 2299 2014 \\ == double \\ table: - 7193 - - 5347 1682 8170 - 6296 3051 5906 2299 2014}

輸出的 \texttt{table} 一行由第 $0$ 格印到第 $12$ 格，\texttt{-} 是空格。 刪除後的
\texttt{table} 裡看不到 \texttt{D}，是因為最後插入的 $4880$ 剛好放進了第 $6$ 格的刪
除標記。 \texttt{probe} 用 \texttt{long long} 算 \texttt{i * i} 和
\texttt{i * (1 + key \% 7)}，$m$ 很大時也不會溢位。

\section{期望探測次數}
\label{sec:hash-analysis}

第~\ref{sec:hash-chain-analysis}~節算出鏈結法的期望比較次數。 這一節對開放定址做同
樣的事。 先要一個和簡單均勻雜湊對應的假設，它講的是整條探測序列。

\begin{definition}[均勻雜湊]
\label{def:hash-uh}
每個鍵的探測序列 $\langle h(k, 0), \ldots, h(k, m-1) \rangle$ 等可能是
$0, 1, \ldots, m-1$ 的 $m!$ 種排列之中的任何一種，而且不同鍵的探測序列互相獨立。
這個假設叫做\textbf{均勻雜湊}（uniform hashing）。
\end{definition}

上一節的三種方法都達不到這個假設（線性和二次探測只有 $m$ 種序列，雙重雜湊大約
$m^2$ 種）；它是一個理想模型，實際的方法離它多遠，最後用程式量。 這一節假設表裡沒有
刪除標記。

\subsection{失敗的搜尋與插入}

\begin{theorem}
\label{thm:hash-unsucc}
在均勻雜湊下，開放定址的表有 $m$ 格、$n$ 個鍵，$n < m$。 搜尋一個不在表裡的鍵，期
望探測次數是
\[
\frac{m + 1}{m - n + 1} < \frac{1}{1 - \alpha}.
\]
\end{theorem}

\begin{proof}
表裡哪 $n$ 格有鍵是由其他鍵的探測序列決定的；要找的鍵 $k$ 的探測序列和它們獨立，所
以可以先把有鍵的 $n$ 格固定下來，叫它們「滿格」，其餘 $m - n$ 格叫「空格」。 $k$ 不
在表裡，搜尋會沿著 $k$ 的探測序列一直看，直到第一個空格才停。 所以
\[
\text{探測次數} = 1 + (\text{排在第一個空格前面的滿格個數}).
\]
對每一個滿格 $s$，令 $I_s$ 是「$s$ 在 $k$ 的探測序列裡排在所有空格前面」的指示變數。
$k$ 的探測序列是均勻隨機的排列，所以 $s$ 和 $m - n$ 個空格這 $m - n + 1$ 格之間的先
後順序也是均勻的，每一格排在最前面的機率都一樣，都是 $1/(m - n + 1)$。 因此
$\mathrm{E}[I_s] = 1/(m - n + 1)$，再用期望值的線性：
\[
\mathrm{E}[\text{探測次數}] = 1 + \sum_{s} \mathrm{E}[I_s]
= 1 + \frac{n}{m - n + 1} = \frac{m + 1}{m - n + 1}.
\]
最後 $m - n + 1 > m - n$，所以這個值小於
$1 + \dfrac{n}{m - n} = \dfrac{m}{m - n} = \dfrac{1}{1 - \alpha}$。
\end{proof}

\begin{corollary}
\label{cor:hash-insert}
在均勻雜湊下，把一個新的鍵插入有 $n$ 個鍵的表，期望探測次數是
$(m + 1)/(m - n + 1)$，小於 $1/(1 - \alpha)$。
\end{corollary}

插入就是做一次失敗的搜尋，再把鍵放進停下來的那個空格，探測序列完全一樣。

$1/(1 - \alpha)$ 有一個好記的意思：每看一格，是空格的機率大約是 $1 - \alpha$，像擲一
枚正面機率 $1 - \alpha$ 的硬幣，期望要擲 $1/(1 - \alpha)$ 次才出現第一個正面。 半滿
時要 $2$ 次，$90\%$ 滿時要 $10$ 次。 表越滿，成本漲得越快。

\subsection{成功的搜尋}

\begin{theorem}
\label{thm:hash-succ}
在均勻雜湊下，表有 $m$ 格、$n$ 個鍵，$n < m$，表裡的每個鍵被搜尋的機會一樣。 成功
搜尋的期望探測次數是
\[
\frac{1}{n} \sum_{i=0}^{n-1} \frac{m+1}{m-i+1}
\;\le\; \frac{1}{\alpha} \ln \frac{1}{1 - \alpha}.
\]
\end{theorem}

\begin{proof}
把鍵依插入的順序編號。 第 $i + 1$ 個鍵插入時表裡有 $i$ 個鍵。 因為沒有刪除，格子只
會由空變滿，搜尋這個鍵時走的路和當初插入它時完全一樣：前面那些格子當時是滿的，現在
還是滿的，它停的那一格現在放的就是它。 所以搜尋它的探測次數等於插入它的探測次數，
由系理~\ref{cor:hash-insert}，期望值是 $(m + 1)/(m - i + 1)$。 每個鍵被搜尋的機率
是 $1/n$，把 $i = 0$ 到 $n - 1$ 平均就是左邊的式子。

接著估上界。 $(m + 1)(m - i) \le m(m - i + 1)$ 展開後是 $-i \le 0$，成立，所以每一項
$\frac{m+1}{m-i+1} \le \frac{m}{m-i}$。 於是
\[
\frac{1}{n} \sum_{i=0}^{n-1} \frac{m+1}{m-i+1}
\le \frac{m}{n} \sum_{i=0}^{n-1} \frac{1}{m - i}
= \frac{1}{\alpha} \sum_{j = m-n+1}^{m} \frac{1}{j}.
\]
$1/x$ 是遞減的，在區間 $[j - 1, j]$ 上 $1/x \ge 1/j$，所以
$\frac{1}{j} \le \int_{j-1}^{j} \frac{dx}{x}$。 把 $j = m - n + 1$ 到 $m$ 加起來：
\[
\sum_{j=m-n+1}^{m} \frac{1}{j} \le \int_{m-n}^{m} \frac{dx}{x}
= \ln \frac{m}{m - n} = \ln \frac{1}{1 - \alpha}. \qedhere
\]
\end{proof}

左邊的式子可以寫成 $\frac{m+1}{n}(H_{m+1} - H_{m-n+1})$，$H_t$ 是第~1~章的調和數。
用第~\ref{sec:hash-open}~節的例子驗算：$m = 13$、$n = 9$，失敗搜尋是
$14/5 = 2.8$，成功搜尋是 $\frac{14}{9}(H_{14} - H_5) \approx 1.51$。 線性探測在同一張
表上實際是 $4.46$ 和 $2.44$，雙重雜湊的成功搜尋是 $13/9 \approx 1.44$。 九個鍵太少，
這些數只能看出方向；大量的資料要用程式量。

\subsection{線性探測的公式}

線性探測有主要群聚，均勻雜湊的公式不適用。 Knuth 在 \textit{The Art of Computer
Programming} 第 3 冊 6.4 節分析了線性探測（假設起始位置 $h(k)$ 均勻而且獨立），$m$
很大、$\alpha$ 固定時，期望探測次數趨近於
\[
\text{成功：} \frac{1}{2}\left(1 + \frac{1}{1 - \alpha}\right), \qquad
\text{失敗：} \frac{1}{2}\left(1 + \frac{1}{(1 - \alpha)^2}\right).
\]
這裡只引用結果，推導要用到比本書多的組合計數，見延伸閱讀。 雙重雜湊雖然只有大約
$m^2$ 種探測序列，已經證明它在 $m$ 很大時的期望探測次數趨近於均勻雜湊的公式（Lueker
與 Molodowitch，1993）。 幾個 $\alpha$ 的公式值，下一小節的程式會和實測值一起印出來。

$\alpha = 0.5$ 時，線性探測的失敗搜尋是 $2.5$ 次，均勻雜湊是 $2$ 次，差不多；過了 $0.7$ 左右，線性探測的失敗搜尋就比均勻雜湊
貴很多，$0.9$ 時差了五倍。 所以開放定址的表通常把 $\alpha$ 控制在 $0.5$ 到 $0.75$
以下，超過就擴表。 反過來說，線性探測每次看的是相鄰的格子，在有快取的真實電腦上一
次讀進來好幾格，$\alpha$ 不高時常常比雙重雜湊快。 RAM 模型（第~1~章）數不出這種
差別。

\subsection{實測}
\label{sec:hash-measure}

下面的程式開一張 $m = 100003$ 格（質數）的表，用自己寫的線性同餘亂數產生鍵，載入因
子依序加到 $0.10$、$0.25$、$0.50$、$0.75$、$0.90$、$0.95$。 每到一個載入因子，就把
表裡每個鍵搜尋一次，算成功搜尋的平均；再搜尋 $100000$ 個一定不在表裡的鍵，算失敗搜
尋的平均。 放進表裡的鍵都是偶數、拿來做失敗搜尋的都是奇數，這樣不用另外檢查。 線性
探測和雙重雜湊（$h_2(k) = 1 + (k \bmod (m - 2))$）用同一串鍵，旁邊印出對應的公式值。

亂數用固定的種子，所以每次執行的輸出都一樣，可以逐字對照。 C 標準函式庫的
\texttt{rand} 在不同的系統上演算法不同，同一個種子會得到不同的數，這裡不用它。 線性
同餘法每一步算 $x \gets (ax + c) \bmod 2^{64}$，$a$、$c$ 是兩個選好的常數；
\texttt{unsigned long long} 的乘法和加法溢位時自動就是 $\bmod 2^{64}$，這是 C 標準規
定的行為，不是未定義行為。 低位元的週期很短，所以取最高的 $31$ 位元當鍵。

\begin{example}[label=ex:hash-probe-count]{c}{(probe\_count.c) 載入因子和平均探測次數}
#include <stdio.h>
#include <stdlib.h>
#include <math.h>

#define M 100003            // 格數（質數）
#define TRIALS 100000       // 每個載入因子做幾次失敗搜尋

// 線性同餘法的亂數產生器。種子固定，每次執行的結果都一樣。
static unsigned long long seed;

int next_key(void)
{
    seed = seed * 6364136223846793005ULL + 1442695040888963407ULL;
    return (int)(seed >> 33);           // 取最高的 31 位元
}

// 雙重雜湊的步長 h2(k) = 1 + k mod (M - 2)；線性探測的步長是 1。
int step_of(int key, int dbl)
{
    return dbl ? 1 + key % (M - 2) : 1;
}

// 找 key，回傳探測次數，*slot 是停下來的格子。-1 代表空格。
int search(const int *table, int key, int dbl, int *slot)
{
    int step = step_of(key, dbl);
    for (long long i = 0; i < M; i++) {
        *slot = (int)((key % M + i * step) % M);
        if (table[*slot] == key || table[*slot] == -1) {
            return (int)i + 1;
        }
    }
    return M;                       // 表滿了（這裡不會發生）
}

int main(void)
{
    const double alphas[] = {0.10, 0.25, 0.50, 0.75, 0.90, 0.95};
    double res[6][4];
    int *table = malloc(M * sizeof *table);
    int *keys = malloc(M * sizeof *keys);
    if (table == NULL || keys == NULL) {
        free(table);
        free(keys);
        return 1;
    }
    for (int dbl = 0; dbl <= 1; dbl++) {
        seed = 20261119ULL;             // 兩種方法用同一串鍵
        for (int j = 0; j < M; j++) {
            table[j] = -1;
        }
        int n = 0, slot;
        for (int a = 0; a < 6; a++) {
            while (n < (int)(alphas[a] * M)) {   // 插入偶數鍵
                int key = next_key() & ~1;
                search(table, key, dbl, &slot);
                if (table[slot] == -1) {   // 停在空格：不在表裡
                    table[slot] = key;
                    keys[n++] = key;
                }
            }
            long long hit = 0, miss = 0;
            for (int i = 0; i < n; i++) {    // 成功：每個鍵找一次
                hit += search(table, keys[i], dbl, &slot);
            }
            for (int t = 0; t < TRIALS; t++) {   // 失敗：奇數鍵
                miss += search(table, next_key() | 1, dbl, &slot);
            }
            res[a][2 * dbl] = (double)hit / n;
            res[a][2 * dbl + 1] = (double)miss / TRIALS;
        }
    }
    printf("        ------ linear probing ------   ");
    printf("------ double hashing ------\n");
    printf("alpha   hit   Knuth   miss   Knuth     ");
    printf("hit  uniform   miss uniform\n");
    for (int a = 0; a < 6; a++) {
        double x = alphas[a];
        printf("%.2f %6.2f %7.2f %6.2f %7.2f  ", x,
               res[a][0], 0.5 * (1 + 1 / (1 - x)),
               res[a][1], 0.5 * (1 + 1 / ((1 - x) * (1 - x))));
        printf("%6.2f %7.2f %6.2f %7.2f\n",
               res[a][2], log(1 / (1 - x)) / x,
               res[a][3], 1 / (1 - x));
    }
    free(table);
    free(keys);
    return 0;
}
\end{example}
\iosample{}{\ \ \ \ \ \ \ \ ------\ linear\ probing\ ------\ \ \ ------\ double\ hashing\ ------ \\ alpha\ \ \ hit\ \ \ Knuth\ \ \ miss\ \ \ Knuth\ \ \ \ \ hit\ \ uniform\ \ \ miss\ uniform \\ 0.10\ \ \ 1.05\ \ \ \ 1.06\ \ \ 1.12\ \ \ \ 1.12\ \ \ \ 1.05\ \ \ \ 1.05\ \ \ 1.11\ \ \ \ 1.11 \\ 0.25\ \ \ 1.17\ \ \ \ 1.17\ \ \ 1.39\ \ \ \ 1.39\ \ \ \ 1.16\ \ \ \ 1.15\ \ \ 1.33\ \ \ \ 1.33 \\ 0.50\ \ \ 1.50\ \ \ \ 1.50\ \ \ 2.49\ \ \ \ 2.50\ \ \ \ 1.41\ \ \ \ 1.39\ \ \ 1.94\ \ \ \ 2.00 \\ 0.75\ \ \ 2.50\ \ \ \ 2.50\ \ \ 8.27\ \ \ \ 8.50\ \ \ \ 1.91\ \ \ \ 1.85\ \ \ 3.80\ \ \ \ 4.00 \\ 0.90\ \ \ 5.35\ \ \ \ 5.50\ \ 48.83\ \ \ 50.50\ \ \ \ 2.66\ \ \ \ 2.56\ \ \ 9.49\ \ \ 10.00 \\ 0.95\ \ 10.00\ \ \ 10.50\ 193.24\ \ 200.50\ \ \ \ 3.28\ \ \ \ 3.15\ \ 19.29\ \ \ 20.00}

\texttt{hit} 是成功搜尋、\texttt{miss} 是失敗搜尋。 線性探測的實測值和 Knuth 的公
式相當接近，$\alpha = 0.95$ 時失敗搜尋要將近 $200$ 次，公式是 $200.5$；雙重雜湊緊跟著
均勻雜湊的公式。

失敗搜尋的實測都比公式低一些，原因不在定理~\ref{thm:hash-unsucc} 的上界：$m$ 這麼大
時，精確值 $(m+1)/(m-n+1)$ 和 $1/(1 - \alpha)$ 差不到 $0.01$。 線性探測的群聚長短起伏
很大，這裡只量了一張表；換一個種子，$\alpha = 0.95$ 的失敗搜尋從一百八十幾到兩百二十
幾都可能。 雙重雜湊偏低，是因為這裡的鍵都小於 $2^{31}$。 寫成 $k = qM + r$，
$h_1(k) = r$，而 $k \bmod (M - 2) = (2q + r) \bmod (M - 2)$ 也由同一個商 $q$ 決定，所以
$h_1$、$h_2$ 並不獨立，和均勻雜湊的假設有一點距離。 把鍵換成 $62$ 位元的亂數再量，實測
和公式就幾乎一樣。

\subsection{鏈結法和開放定址怎麼選}

鏈結法的 $\alpha$ 可以大於 $1$，變大時慢慢變慢；刪除直接從串列拿掉；代價是每個鍵多
一個指標、每次插入配置一個節點，走串列時每個節點可能在記憶體的不同地方。 開放定址只
有一個陣列，鍵小的時候省很多空間，線性探測讀的是相鄰的格子，對快取友善；代價是
$\alpha$ 接近 $1$ 時急遽變慢，刪除要用標記，標記多了要重建。 所以鍵數難預估、刪除頻
繁、每筆資料很大時用鏈結法；鍵數大致已知、資料小、以查詢為主時用開放定址。

\section{再雜湊、可延伸雜湊與通用雜湊}
\label{sec:hash-dynamic}

\subsection{再雜湊}

前面的分析都假設 $\alpha$ 是常數。 鍵一直加進來，$\alpha$ 就一直變大，鏈結法的串列
越來越長，開放定址最後會滿。 做法和第~\ref{sec:amo-dynarray}~節的動態陣列一樣：$\alpha$ 超過門檻
（例如開放定址 $0.5$、鏈結法 $1$）時，配置一張大約兩倍大的新表，把舊表的鍵逐一插入
新表，再釋放舊表。 這叫\textbf{再雜湊}（rehashing）。

搬的時候不能照格子號碼直接複製，因為 $h$ 和 $m$ 有關，$m$ 變了，每個鍵的位置都要重
算。 把第~\ref{sec:hash-open}~節線性探測的表（$m = 13$，九個鍵）搬到 $m' = 29$（大
於 $2 \times 13$ 的最小質數），依舊表的格子順序逐一插入：

\begin{center}
\begin{tabular}{lccccccccc}
\toprule
鍵 & 5347 & 6296 & 1682 & 8170 & 5906 & 3051 & 7193 & 2299 & 2014 \\
\midrule
舊位置（$m = 13$）& 4 & 5 & 6 & 7 & 8 & 9 & 10 & 11 & 12 \\
$k \bmod 29$ & 11 & 3 & 0 & 21 & 19 & 6 & 1 & 8 & 13 \\
\bottomrule
\end{tabular}
\end{center}

九個鍵在新表裡沒有任何碰撞，每個都只探測 $1$ 次，舊表裡那一長段群聚完全散開了，
$\alpha$ 降到 $9/29 \approx 0.31$。 一次再雜湊要 $\Theta(n + m)$，但和動態陣列一樣，
表變成兩倍之後要再插入大約那麼多個鍵才會再擴表，把搬動的成本平攤到這些插入上，每次
插入的平攤成本仍然是期望 $O(1)$。 再雜湊也用來清掉刪除標記（第~\ref{sec:hash-tomb}~節），
這時新表的大小可以不變。

要注意「再雜湊」這個詞在一些資料結構課本裡有另一個意思：碰撞時改用第二個、第三個
雜湊函數 $h_2, h_3, \ldots$ 算下一個位置（雙重雜湊是它的一個特例）。 考題寫「rehashing」
時，先從上下文判斷是哪一種。

\subsection{可延伸雜湊}

再雜湊一次要搬動全部的鍵。 資料存在磁碟上時，這代價太大：搬一次要把整個檔案讀寫一
遍。 \textbf{可延伸雜湊}（extendible hashing）讓表一次只長一點點，每次只拆開一個桶。

資料放在固定容量的\textbf{桶}（bucket，在磁碟上就是一頁）裡。 另外有一個
\textbf{目錄}（directory），是一個有 $2^d$ 個指標的陣列，$d$ 叫做\textbf{全域深度}
（global depth）。 每個鍵先算出一個夠長的雜湊值（下面用 $5$ 個位元），取\textbf{最
前面的 $d$ 個位元}當目錄的索引，順著指標找到桶。 幾個相鄰的目錄項可以指向同一個桶；
每個桶記一個\textbf{局部深度}（local depth）$d'$，表示這個桶裡的鍵前 $d'$ 個位元都
相同，有 $2^{d - d'}$ 個目錄項指向它。

插入時如果桶滿了，就把它拆成兩個桶，局部深度各加一，依第 $d' + 1$ 個位元把鍵分到兩
邊：

\begin{itemize}
\item $d' < d$：有不只一個目錄項指向這個桶，把其中一半改指新的桶就好，目錄不變。
\item $d' = d$：只有一個目錄項指向它，目錄先加倍（$d$ 加一，新的第 $x$ 項指向舊的第
$\lfloor x/2 \rfloor$ 項指的桶，也就是去掉最後一個位元），再照上一種情形拆。
\end{itemize}

拆完如果新的桶還是滿的（所有鍵的下一個位元都一樣），就再拆一次。 下面用容量 $2$ 的桶
示範。 七個鍵的雜湊值如下：

\begin{center}
\begin{tabular}{lccccccc}
\toprule
鍵 & A & B & C & D & E & F & G \\
\midrule
$h(k)$ & \texttt{00110} & \texttt{10011} & \texttt{01101} & \texttt{11010} &
\texttt{01001} & \texttt{11100} & \texttt{01110} \\
\bottomrule
\end{tabular}
\end{center}

放 E = \texttt{0}1001 時第一個桶滿了，而且 $d' = d = 1$：目錄先加倍成 $d = 2$
（\texttt{00}、\texttt{01} 都先指向原來的桶），再依第 $2$ 個位元拆，A = 0\texttt{0}110
留下，C、E（0\texttt{1}\ldots）到新的桶，見圖~\ref{fig:hash-extendible}(b)；B、D 的桶沒有動，局部深度仍是 $1$。
放 F = \texttt{11}100 時 B、D 的桶滿了，但 $d' = 1 < d = 2$，只拆桶、把目錄項
\texttt{11} 改指新的桶，見 (c)。 放 G = \texttt{01}110 時 C、E 的桶滿了，$d' = d = 2$，
目錄再加倍成 $d = 3$，依第 $3$ 個位元把 E = 01\texttt{0}01 和 C、G（01\texttt{1}\ldots）
分開，見 (d)；其他三個桶不動，各有兩個目錄項指著。

\begin{figure}[htbp]
\centering
\begin{tikzpicture}[x=1cm, y=0.52cm]
% (a) d = 1，放完 A、B、C、D
\begin{scope}
\node[font=\small] at (1.2, 1.1) {(a) $d = 1$};
\node[hdir] (a0) at (0, 0) {0};
\node[hdir] (a1) at (0, -1) {1};
\node[hpage] (ap) at (2.4, 0) {A C \quad $d'{=}1$};
\node[hpage] (aq) at (2.4, -1) {B D \quad $d'{=}1$};
\draw[diagarrow] (a0) -- (ap);
\draw[diagarrow] (a1) -- (aq);
\end{scope}
% (b) 放 E：目錄加倍
\begin{scope}[xshift=6.6cm]
\node[font=\small] at (1.2, 1.1) {(b) 放 E，目錄加倍，$d = 2$};
\foreach \b [count=\i from 0] in {00,01,10,11} { \node[hdir] (b\i) at (0, -\i) {\b}; }
\node[hpage] (bp) at (2.4, 0) {A \quad $d'{=}2$};
\node[hpage] (bq) at (2.4, -1) {C E \quad $d'{=}2$};
\node[hpage] (br) at (2.4, -2.5) {B D \quad $d'{=}1$};
\draw[diagarrow] (b0) -- (bp);
\draw[diagarrow] (b1) -- (bq);
\draw[diagarrow] (b2) -- (br);
\draw[diagarrow] (b3) -- (br);
\end{scope}
% (c) 放 F：拆桶，目錄不變
\begin{scope}[yshift=-2.9cm]
\node[font=\small] at (1.2, 1.1) {(c) 放 F，只拆桶，$d = 2$};
\foreach \b [count=\i from 0] in {00,01,10,11} { \node[hdir] (c\i) at (0, -\i) {\b}; }
\node[hpage] (cp) at (2.4, 0) {A \quad $d'{=}2$};
\node[hpage] (cq) at (2.4, -1) {C E \quad $d'{=}2$};
\node[hpage] (cr) at (2.4, -2) {B \quad $d'{=}2$};
\node[hpage] (cs) at (2.4, -3) {D F \quad $d'{=}2$};
\foreach \i/\t in {0/cp, 1/cq, 2/cr, 3/cs} { \draw[diagarrow] (c\i) -- (\t); }
\end{scope}
% (d) 放 G：目錄再加倍
\begin{scope}[xshift=6.6cm, yshift=-2.9cm]
\node[font=\small] at (1.2, 1.1) {(d) 放 G，目錄加倍，$d = 3$};
\foreach \b [count=\i from 0] in {000,001,010,011,100,101,110,111} {
    \node[hdir] (d\i) at (0, -\i) {\b};
}
\node[hpage] (dp) at (2.6, -0.5) {A \quad $d'{=}2$};
\node[hpage] (dq) at (2.6, -2) {E \quad $d'{=}3$};
\node[hpage] (dr) at (2.6, -3) {C G \quad $d'{=}3$};
\node[hpage] (ds) at (2.6, -4.5) {B \quad $d'{=}2$};
\node[hpage] (dt) at (2.6, -6.5) {D F \quad $d'{=}2$};
\foreach \i/\t in {0/dp, 1/dp, 2/dq, 3/dr, 4/ds, 5/ds, 6/dt, 7/dt} {
    \draw[diagarrow] (d\i.east) -- (\t.west);
}
\end{scope}
\end{tikzpicture}
\diagcaption{可延伸雜湊，桶的容量是 $2$，用雜湊值最前面的 $d$ 個位元查目錄。}
\label{fig:hash-extendible}
\end{figure}

每次只拆一個桶、只搬那個桶裡的鍵；加倍的只是目錄，目錄只存指標，常常整個放得進記憶
體，搜尋一個鍵只要讀一個桶。 代價是只要有一個桶一直被拆，目錄就得一直加倍。

\begin{examnote}
可延伸雜湊的題目，先確認用的是雜湊值的最前面幾個位元還是最後面幾個位元（本章用最前
面，也有很多課本用最後面），以及桶的容量。 每一步畫出目錄（每一項的位元字串）和每個
桶的局部深度、內容；目錄加倍和只拆桶要分清楚，判斷的依據是 $d'$ 和 $d$ 相不相等。
\end{examnote}

\subsection{通用雜湊}

任何一個固定的 $h$ 都有一批鍵全部落在同一格（第~\ref{sec:hash-func}~節）。 有人知
道伺服器用什麼雜湊函數，就能故意送一堆碰撞的鍵讓它變慢，簡單均勻雜湊的假設不成立。 \textbf{通用雜湊}（universal hashing）的想法是：先準備
一整族函數，程式開始執行時才從裡面隨機挑一個。 輸入是事先決定的，不知道會挑到哪一
個，所以沒有一組輸入對大部分的函數都很糟。 這時的期望值是對「挑函數」這件事取的，
不需要對輸入做任何假設。

\begin{definition}[通用雜湊族]
\label{def:hash-universal}
一族函數 $\mathcal{H}$，每個都把 $U$ 對應到 $\{0, \ldots, m-1\}$。 如果對任兩個不
同的鍵 $x \ne y$，滿足 $h(x) = h(y)$ 的函數 $h$ 不超過 $|\mathcal{H}|/m$ 個，就說
$\mathcal{H}$ 是\textbf{通用的}（universal）。 換句話說，從 $\mathcal{H}$ 隨機挑一
個 $h$，任兩個固定的不同鍵碰撞的機率不超過 $1/m$。
\end{definition}

\begin{theorem}
\label{thm:hash-universal}
用鏈結法，$h$ 從通用雜湊族裡隨機挑選。 對任何一組事先固定的 $n$ 個鍵，以及任何一個
鍵 $k$，第 $h(k)$ 條串列裡除了 $k$ 本身以外的鍵，期望個數不超過 $\alpha$。 所以搜尋
$k$ 期望比較的次數，$k$ 不在表裡時不超過 $\alpha$，在表裡時不超過 $1 + \alpha$。
\end{theorem}

\begin{proof}
對表裡每個不等於 $k$ 的鍵 $x$，令 $I_x = I\{h(x) = h(k)\}$。 由定義
$\mathrm{E}[I_x] = \Pr[h(x) = h(k)] \le 1/m$。 這樣的 $x$ 最多 $n$ 個，加起來期望值不
超過 $n/m = \alpha$。 搜尋 $k$ 比較的就是這條串列上的鍵；$k$ 在表裡時再加上它自己。
\end{proof}

和定理~\ref{thm:hash-chain-len} 比，結論的形式一樣，前提卻不同：這裡沒有假設鍵的分
布，隨機性完全來自程式自己挑函數。 一個常用的通用雜湊族是：取一個比所有鍵都大的質
數 $p$，
\[
h_{a,b}(k) = \bigl( (ak + b) \bmod p \bigr) \bmod m, \qquad
1 \le a \le p - 1,\ 0 \le b \le p - 1.
\]
每組 $(a, b)$ 是一個函數，一共 $p(p-1)$ 個。 它是通用的（Carter 與 Wegman，1979；
CLRS 第 4 版第 11.3.4 節有完整證明，這裡不證）。 例如 $p = 10007$、$m = 10$，執行時隨機挑到
$a = 37$、$b = 58$，則 $h(2026) = (75020 \bmod 10007) \bmod 10 = 4971 \bmod 10 = 1$。
下次執行挑到別的 $(a, b)$，同一個鍵就落在別的格子。

\section{布隆過濾器}
\label{sec:hash-bloom}

刷卡機還有一件事要做：擋下停權的卡（掛失、畢業、被盜刷）。 停權名單在學校的伺服器
上，有二十萬張；每刷一次卡就連線去問，中午排隊的人會等不及。 把名單整份放進刷卡機
也不行：用前面的雜湊表存二十萬個九位數卡號，光是鍵就要 $1.6$ MB，刷卡機只剩幾百 KB
的記憶體。 不過絕大多數刷卡的人都不在名單上，刷卡機其實只需要一個很小的東西，能很
快地說出「這張一定沒有停權」；少數說不準的，再去問伺服器。

\textbf{布隆過濾器}（Bloom filter，Burton Bloom 在 1970 年提出）做的就是這件事。 它只
回答「$x$ 在不在集合 $S$ 裡」，不存鍵本身，所以很省空間；代價是偶爾會把不在 $S$ 裡
的 $x$ 說成「可能在」。 反過來的錯誤不會發生：它說「不在」，就一定不在。

\subsection{結構與操作}

布隆過濾器是一個有 $m$ 個位元的陣列 $B[0..m-1]$，一開始全部是 $0$，再加上 $k$ 個雜湊
函數 $h_1, h_2, \ldots, h_k$，每個都把鍵對應到 $\{0, 1, \ldots, m - 1\}$。
\begin{itemize}
\item \textbf{插入} $x$：把 $B[h_1(x)], B[h_2(x)], \ldots, B[h_k(x)]$ 這 $k$ 個位元都設成
$1$（原本已經是 $1$ 的就不變）。
\item \textbf{查詢} $x$：看同樣的 $k$ 個位元。 只要有一個是 $0$，回答「不在」；全部是
$1$，回答「可能在」。
\end{itemize}

\begin{pseudo}[label=alg:hash-bloom]{布隆過濾器的插入與查詢}
\Procedure{Bloom-Insert}{$B, x$}
    \For{$i \gets 1$ \To $k$}
        \State $B[h_i(x)] \gets 1$
    \EndFor
\EndProcedure
\Function{Bloom-Query}{$B, x$}
    \For{$i \gets 1$ \To $k$}
        \If{$B[h_i(x)] = 0$}
            \State \Return \FALSE \Comment{一定不在}
        \EndIf
    \EndFor
    \State \Return \TRUE \Comment{可能在}
\EndFunction
\end{pseudo}

兩個操作都只算 $k$ 個雜湊值、看 $k$ 個位元，時間是 $O(k)$，和集合裡有幾個鍵無關。 空
間是 $m$ 個位元，也和鍵有多長無關：九位數的卡號和一百個字元的網址，佔的空間一樣。

\begin{theorem}
\label{thm:hash-bloom-noneg}
布隆過濾器不會有\textbf{偽陰性}（false negative）：如果 $x$ 已經插入過，
$\proc{Bloom-Query}(x)$ 一定回傳 \TRUE。
\end{theorem}

\begin{proof}
插入 $x$ 時，第 $3$ 行把 $B[h_1(x)], \ldots, B[h_k(x)]$ 都設成 $1$。 之後的操作只有插
入和查詢，插入只會把位元設成 $1$，查詢不改任何位元，所以這 $k$ 個位元一直是 $1$。 查
詢 $x$ 時第 $6$ 行的條件永遠不成立，迴圈跑完，回傳 \TRUE。
\end{proof}

會出錯的是另一個方向。 $x$ 沒有插入過，但它的 $k$ 個位元剛好都被別的鍵設成了 $1$，查
詢就回答「可能在」，這叫\textbf{偽陽性}（false positive）。 所以「可能在」只是可能：刷
卡機遇到「可能在」，再連線問伺服器確認；遇到「不在」就直接放行。 只要偽陽性很少，大部
分的刷卡都不必連線。

\subsection{手算例子}
\label{sec:hash-bloom-trace}

取 $m = 13$、$k = 3$。 三個雜湊函數借用第~\ref{sec:hash-open}~節雙重雜湊的式子：
\[
h_i(x) = \bigl( (x \bmod 13) + (i - 1)\,(1 + x \bmod 11) \bigr) \bmod 13, \qquad i = 1, 2, 3,
\]
也就是從 $x \bmod 13$ 出發，每次跳 $1 + (x \bmod 11)$ 格。 實務上常這樣用兩個雜湊值組出
$k$ 個，不必真的寫 $k$ 個不同的函數。 依序插入四張停權卡 $2037$、$5512$、$6180$、
$7423$，再查詢三張卡，見圖~\ref{fig:hash-bloom}。

\begin{figure}[htbp]
\centering
\begin{tikzpicture}
\foreach \j in {0,...,12} { \node[idxlbl] at ({0.6+\j*0.8}, 0.5) {\j}; }
\node[anchor=east, font=\footnotesize] at (0.2,0.00) {開始};
\node[hbit] at (0.60,0.00) {0};
\node[hbit] at (1.40,0.00) {0};
\node[hbit] at (2.20,0.00) {0};
\node[hbit] at (3.00,0.00) {0};
\node[hbit] at (3.80,0.00) {0};
\node[hbit] at (4.60,0.00) {0};
\node[hbit] at (5.40,0.00) {0};
\node[hbit] at (6.20,0.00) {0};
\node[hbit] at (7.00,0.00) {0};
\node[hbit] at (7.80,0.00) {0};
\node[hbit] at (8.60,0.00) {0};
\node[hbit] at (9.40,0.00) {0};
\node[hbit] at (10.20,0.00) {0};
\node[anchor=east, font=\footnotesize] at (0.2,-0.70) {插入 2037：9, 12, 2};
\node[hbit] at (0.60,-0.70) {0};
\node[hbit] at (1.40,-0.70) {0};
\node[hbit, hlnew] at (2.20,-0.70) {1};
\node[hbit] at (3.00,-0.70) {0};
\node[hbit] at (3.80,-0.70) {0};
\node[hbit] at (4.60,-0.70) {0};
\node[hbit] at (5.40,-0.70) {0};
\node[hbit] at (6.20,-0.70) {0};
\node[hbit] at (7.00,-0.70) {0};
\node[hbit, hlnew] at (7.80,-0.70) {1};
\node[hbit] at (8.60,-0.70) {0};
\node[hbit] at (9.40,-0.70) {0};
\node[hbit, hlnew] at (10.20,-0.70) {1};
\node[anchor=east, font=\footnotesize] at (0.2,-1.40) {插入 5512：0, 2, 4};
\node[hbit, hlnew] at (0.60,-1.40) {1};
\node[hbit] at (1.40,-1.40) {0};
\node[hbit, hlprobe] at (2.20,-1.40) {1};
\node[hbit] at (3.00,-1.40) {0};
\node[hbit, hlnew] at (3.80,-1.40) {1};
\node[hbit] at (4.60,-1.40) {0};
\node[hbit] at (5.40,-1.40) {0};
\node[hbit] at (6.20,-1.40) {0};
\node[hbit] at (7.00,-1.40) {0};
\node[hbit] at (7.80,-1.40) {1};
\node[hbit] at (8.60,-1.40) {0};
\node[hbit] at (9.40,-1.40) {0};
\node[hbit] at (10.20,-1.40) {1};
\node[anchor=east, font=\footnotesize] at (0.2,-2.10) {插入 6180：5, 2, 12};
\node[hbit] at (0.60,-2.10) {1};
\node[hbit] at (1.40,-2.10) {0};
\node[hbit, hlprobe] at (2.20,-2.10) {1};
\node[hbit] at (3.00,-2.10) {0};
\node[hbit] at (3.80,-2.10) {1};
\node[hbit, hlnew] at (4.60,-2.10) {1};
\node[hbit] at (5.40,-2.10) {0};
\node[hbit] at (6.20,-2.10) {0};
\node[hbit] at (7.00,-2.10) {0};
\node[hbit] at (7.80,-2.10) {1};
\node[hbit] at (8.60,-2.10) {0};
\node[hbit] at (9.40,-2.10) {0};
\node[hbit, hlprobe] at (10.20,-2.10) {1};
\node[anchor=east, font=\footnotesize] at (0.2,-2.80) {插入 7423：0, 10, 7};
\node[hbit, hlprobe] at (0.60,-2.80) {1};
\node[hbit] at (1.40,-2.80) {0};
\node[hbit] at (2.20,-2.80) {1};
\node[hbit] at (3.00,-2.80) {0};
\node[hbit] at (3.80,-2.80) {1};
\node[hbit] at (4.60,-2.80) {1};
\node[hbit] at (5.40,-2.80) {0};
\node[hbit, hlnew] at (6.20,-2.80) {1};
\node[hbit] at (7.00,-2.80) {0};
\node[hbit] at (7.80,-2.80) {1};
\node[hbit, hlnew] at (8.60,-2.80) {1};
\node[hbit] at (9.40,-2.80) {0};
\node[hbit] at (10.20,-2.80) {1};
\node[anchor=east, font=\footnotesize] at (0.2,-3.75) {查詢 6180：5, 2, 12};
\node[hbit] at (0.60,-3.75) {1};
\node[hbit] at (1.40,-3.75) {0};
\node[hbit, hlprobe] at (2.20,-3.75) {1};
\node[hbit] at (3.00,-3.75) {0};
\node[hbit] at (3.80,-3.75) {1};
\node[hbit, hlprobe] at (4.60,-3.75) {1};
\node[hbit] at (5.40,-3.75) {0};
\node[hbit] at (6.20,-3.75) {1};
\node[hbit] at (7.00,-3.75) {0};
\node[hbit] at (7.80,-3.75) {1};
\node[hbit] at (8.60,-3.75) {1};
\node[hbit] at (9.40,-3.75) {0};
\node[hbit, hlprobe] at (10.20,-3.75) {1};
\node[anchor=east, font=\footnotesize] at (0.2,-4.45) {查詢 3141：8, 2, 9};
\node[hbit] at (0.60,-4.45) {1};
\node[hbit] at (1.40,-4.45) {0};
\node[hbit] at (2.20,-4.45) {1};
\node[hbit] at (3.00,-4.45) {0};
\node[hbit] at (3.80,-4.45) {1};
\node[hbit] at (4.60,-4.45) {1};
\node[hbit] at (5.40,-4.45) {0};
\node[hbit] at (6.20,-4.45) {1};
\node[hbit, hbzero] at (7.00,-4.45) {0};
\node[hbit] at (7.80,-4.45) {1};
\node[hbit] at (8.60,-4.45) {1};
\node[hbit] at (9.40,-4.45) {0};
\node[hbit] at (10.20,-4.45) {1};
\node[anchor=east, font=\footnotesize] at (0.2,-5.15) {查詢 4126：5, 7, 9};
\node[hbit] at (0.60,-5.15) {1};
\node[hbit] at (1.40,-5.15) {0};
\node[hbit] at (2.20,-5.15) {1};
\node[hbit] at (3.00,-5.15) {0};
\node[hbit] at (3.80,-5.15) {1};
\node[hbit, hlprobe] at (4.60,-5.15) {1};
\node[hbit] at (5.40,-5.15) {0};
\node[hbit, hlprobe] at (6.20,-5.15) {1};
\node[hbit] at (7.00,-5.15) {0};
\node[hbit, hlprobe] at (7.80,-5.15) {1};
\node[hbit] at (8.60,-5.15) {1};
\node[hbit] at (9.40,-5.15) {0};
\node[hbit] at (10.20,-5.15) {1};

\end{tikzpicture}
\diagcaption{布隆過濾器（$m = 13$，$k = 3$）。 左邊是操作和算出的三個位置。 插入時橘框
是這次由 $0$ 變 $1$ 的位元，粉紅底是本來就是 $1$ 的；查詢時粉紅底是看過、是 $1$ 的位元，
深色底是看到的 $0$，看到它就停。}
\label{fig:hash-bloom}
\end{figure}

以 $2037$ 為例：$2037 = 13 \times 156 + 9$，$2037 = 11 \times 185 + 2$，所以從第 $9$ 格出
發、每次跳 $3$ 格，三個位置是 $9$、$12$、$15 \bmod 13 = 2$。 其他三張卡依同樣的算法得到
$\{0, 2, 4\}$、$\{5, 2, 12\}$、$\{0, 10, 7\}$。 四張卡一共設了 $12$ 次位元，但第 $0$、
$2$、$12$ 格被設了不只一次，最後只有 $8$ 個位元是 $1$。

三個查詢：
\begin{itemize}
\item $6180$ 插入過，三個位元 $5, 2, 12$ 都是 $1$，回答「可能在」。 這是對的。
\item $3141$：$3141 \bmod 13 = 8$，$1 + 3141 \bmod 11 = 8$，位置是 $8, 2, 9$。 第 $8$ 格
是 $0$，回答「不在」。 看到第一個 $0$ 就可以停，不必看完三格。
\item $4126$：$4126 \bmod 13 = 5$，$1 + 4126 \bmod 11 = 2$，位置是 $5, 7, 9$，分別被
$6180$、$7423$、$2037$ 設成 $1$，回答「可能在」。 但 $4126$ 從來沒有插入過，這是一次偽
陽性。
\end{itemize}

這個例子的過濾器已經有 $8/13$ 的位元是 $1$，太滿了。 把 $1000$ 到 $9999$ 之間沒插入過
的卡號都查一次，大約有五分之一會被誤判成「可能在」。 下一小節會算出這個比例怎麼由 $m$、
$n$、$k$ 決定。

\subsection{偽陽性率}
\label{sec:hash-bloom-fp}

和第~\ref{sec:hash-chain-analysis}~節一樣，先假設雜湊值像亂數：每個 $h_i(x)$ 都等可能
是 $0$ 到 $m - 1$ 的任何一格，而且不同的 $i$、不同的鍵之間互相獨立。 插入了 $n$ 個鍵之
後，查詢一個不在集合裡的鍵 $y$，回答「可能在」的機率是多少？

先看一個固定的位元 $B[j]$。 $n$ 個鍵一共算了 $kn$ 個雜湊值，每一個落在 $j$ 以外的機率
是 $1 - 1/m$，而且互相獨立，所以
\[
\Pr[B[j] = 0] = \left(1 - \frac{1}{m}\right)^{kn} \approx e^{-kn/m}.
\]
近似的根據是引理~\ref{lem:mth-exp}：把 $(1 - 1/m)^{kn}$ 寫成
$\bigl((1 - 1/m)^m\bigr)^{kn/m}$，括號裡的數夾在 $e^{-1}(1 - 1/m)$ 和 $e^{-1}$ 之間，$m$
很大時就是 $e^{-1}$。 記 $p = e^{-kn/m}$，表示大約有 $p$ 比例的位元還是 $0$。

$y$ 的 $k$ 個位置也是均勻亂選的，每一個看到 $1$ 的機率大約是 $1 - p$。 $k$ 個都看到 $1$
才會誤判，把它們當成互相獨立，得到偽陽性率
\begin{equation}
\label{eq:hash-bloom-fp}
\varepsilon \approx (1 - p)^k = \bigl(1 - e^{-kn/m}\bigr)^k.
\end{equation}
嚴格說，各個位元是不是 $0$ 並不完全獨立（一個位元是 $1$，別的位元是 $1$ 的機會就小一
點），所以這是近似值；$m$ 很大時，$1$ 的比例會非常接近它的期望值，這個近似相當準。
第~\ref{sec:hash-bloom-c}~節的程式會拿實測值來比。

用手算的例子驗證：$m = 13$、$n = 4$、$k = 3$，式~\eqref{eq:hash-bloom-fp} 給
$(1 - e^{-12/13})^3 \approx 0.22$，和上一小節「大約五分之一」相近。

\paragraph{$k$ 要取多少？}
$m$、$n$ 固定時，$k$ 太小，每次只檢查少數幾個位元，容易全部碰巧是 $1$；$k$ 太大，每個
鍵設很多位元，陣列很快就被 $1$ 塞滿。 中間有一個最好的 $k$。

\begin{theorem}
\label{thm:hash-bloom-k}
把 $k$ 當成正實數，$\bigl(1 - e^{-kn/m}\bigr)^k$ 在
\[
k^{*} = \frac{m}{n} \ln 2 \approx 0.693\,\frac{m}{n}
\]
時最小。 這時 $p = e^{-k^{*} n/m} = \frac12$，也就是大約一半的位元是 $0$，偽陽性率是
\[
\left(\frac12\right)^{k^{*}} = 2^{-(m/n)\ln 2} \approx 0.6185^{\,m/n}.
\]
\end{theorem}

\begin{whybox}[定理~\ref{thm:hash-bloom-k} 的證明]
取對數比較好處理：$\ln \varepsilon = k \ln(1 - e^{-kn/m})$。 改用 $p = e^{-kn/m}$ 當變
數，$k$ 從 $0$ 變到 $\infty$ 時 $p$ 從 $1$ 遞減到 $0$，而且 $k = -\frac{m}{n} \ln p$，代進去
\[
\ln \varepsilon = -\frac{m}{n}\,\ln p \cdot \ln(1 - p).
\]
$m/n$ 是固定的正數，所以要讓 $\varepsilon$ 最小，就是要讓 $\varphi(p) = \ln p \cdot \ln(1 - p)$
在 $0 < p < 1$ 上最大（兩個對數都是負的，$\varphi > 0$）。 微分：
\[
\varphi'(p) = \frac{\ln(1 - p)}{p} - \frac{\ln p}{1 - p}
= \frac{\chi(p)}{p(1 - p)}, \qquad
\chi(p) = (1 - p)\ln(1 - p) - p \ln p .
\]
分母是正的，只要看 $\chi$ 的正負。 $\chi(\frac12) = 0$，$\chi(1 - p) = -\chi(p)$，而且
$p \to 0^{+}$ 時 $p \ln p \to 0$，所以 $\chi(p) \to 0$。 再微分兩次：
\[
\chi'(p) = -\ln(1 - p) - 1 - \ln p - 1, \qquad
\chi''(p) = -\frac{1}{p} + \frac{1}{1 - p} = \frac{2p - 1}{p(1 - p)}.
\]
$0 < p < \frac12$ 時 $\chi'' < 0$，$\chi$ 是凹函數，圖形在兩端點 $(0, 0)$、
$(\frac12, 0)$ 的連線之上，所以 $\chi(p) > 0$；由 $\chi(1 - p) = -\chi(p)$，
$\frac12 < p < 1$ 時 $\chi(p) < 0$。 於是 $\varphi$ 在 $p < \frac12$ 遞增、$p > \frac12$
遞減，最大值在 $p = \frac12$。 $e^{-kn/m} = \frac12$ 解出 $k = \frac{m}{n}\ln 2$；代回
去，$\varepsilon = (1 - \frac12)^{k^{*}} = 2^{-k^{*}}$，而 $2^{-\ln 2} = e^{-(\ln 2)^2}
\approx 0.6185$。
\end{whybox}

$k$ 必須是整數。 由證明，$\varepsilon$ 在 $k^{*}$ 左邊遞減、右邊遞增，所以只要比較
$\lfloor k^{*} \rfloor$ 和 $\lceil k^{*} \rceil$ 兩個值，取比較小的那個。 下表是每個鍵分
到 $m/n$ 個位元時的結果。

\begin{center}
\begin{tabular}{cccc}
\toprule
每個鍵的位元數 $m/n$ & $k^{*} = \frac{m}{n}\ln 2$ & 最好的整數 $k$ & 偽陽性率 \\
\midrule
$4$ & $2.77$ & $3$ & $14.7\%$ \\
$8$ & $5.55$ & $6$ & $2.16\%$ \\
$10$ & $6.93$ & $7$ & $0.82\%$ \\
$16$ & $11.09$ & $11$ & $0.046\%$ \\
\bottomrule
\end{tabular}
\end{center}

反過來，想要偽陽性率是 $\varepsilon$，令 $2^{-k^{*}} = \varepsilon$，得到
$k^{*} = \lg(1/\varepsilon)$，每個鍵需要
\[
\frac{m}{n} = \frac{k^{*}}{\ln 2} = \frac{\lg(1/\varepsilon)}{\ln 2} \approx 1.44 \lg \frac{1}{\varepsilon}
\]
個位元。 $\varepsilon = 1\%$ 時大約是 $9.6$ 個位元、$k = 7$。 回到章首的刷卡機：二十萬張
停權卡、每張 $10$ 個位元，一共兩百萬個位元，大約 $244$ KB，放得進刷卡機；沒有停權的
卡，每刷一百次平均不到一次要連線確認。 注意這個數字和鍵有多長完全無關。

\subsection{不能刪除}

卡片解除停權時，能不能把它的 $k$ 個位元清成 $0$？ 不行。 那些位元可能也屬於別的鍵。
在圖~\ref{fig:hash-bloom} 的過濾器裡「刪除」$6180$，把第 $5$、$2$、$12$ 格清成 $0$；
再查 $2037$，它的位置是 $9, 12, 2$，第 $12$ 格已經是 $0$，回答「不在」。 $2037$ 明明
還在名單上，這是一次偽陰性，定理~\ref{thm:hash-bloom-noneg} 的保證就沒了。 這和開放
定址不能把刪掉的格子改回空的（第~\ref{sec:hash-tomb}~節）是同一類問題：一個位置被好幾
個鍵共用，清掉它會影響別人。

需要刪除時可以用\textbf{計數型布隆過濾器}（counting Bloom filter）：每個位置從一個位元
改成一個小計數器，插入時 $k$ 個計數器各加一，刪除時各減一，查詢時看 $k$ 個計數器是不是
都大於 $0$，代價是空間變成好幾倍。 名單變動不頻繁時，另一個做法是定期用最新的名單重建
整個過濾器，這和第~\ref{sec:hash-tomb}~節清理刪除標記的想法一樣。

\begin{examnote}
布隆過濾器的題目常見三種：(1) 給位元陣列和雜湊函數，判斷插入會不會改到位元、查詢回答
什麼，逐一算出 $k$ 個位置再看位元，有一個 $0$ 就是「不在」；(2) 問最佳的雜湊函數個數，
答 $\frac{m}{n}\ln 2$，題目的選項常寫成「$3\ln 2$」這種形式，先認清 $m$（位元數）和 $n$
（鍵數）；(3) 問偽陽性率，寫出 $\bigl(1 - (1 - \frac1m)^{kn}\bigr)^k \approx
\bigl(1 - e^{-kn/m}\bigr)^k$ 兩種寫法都可以。 題目說「刪除某個鍵就把它的位元清成 $0$」，
幾乎都是要你指出這樣會造成偽陰性。
\end{examnote}

\subsection{C 實作與實測}
\label{sec:hash-bloom-c}

下面的程式先重做圖~\ref{fig:hash-bloom} 的例子，再做一個實驗：$m = 10007$ 個位元、
$n = 1000$ 個鍵（每個鍵大約 $10$ 個位元），$k$ 從 $1$ 試到 $10$，每個 $k$ 都查詢
$100000$ 個不在集合裡的鍵，算出誤判的比例，旁邊印出式~\eqref{eq:hash-bloom-fp} 的值。
雜湊函數用和手算例子相同的寫法，$m = 13$ 時 $m - 2 = 11$，正好是上面的式子。 鍵用
第~\ref{sec:hash-measure}~節的線性同餘亂數產生，插入的都是偶數、查詢的都是奇數，所以查
詢的鍵一定不在集合裡。

為了好讀，程式用一個 \texttt{unsigned char} 存一個位元，浪費了八倍的空間。 真正的實作
把八個位元塞進一個位元組，用《C 程式設計基礎》第~12~章的位元運算讀寫：第 $j$ 個位元
在 \texttt{bit[j / 8]} 的第 \texttt{j \% 8} 位。

\begin{example}[label=ex:hash-bloom]{c}{(bloom\_filter.c) 布隆過濾器：手算例子與偽陽性率實測}
#include <stdio.h>
#include <stdlib.h>
#include <math.h>

// Bloom filter：m 個位元，k 個雜湊函數。
// 第 i 個函數（i 從 0 開始）用雙重雜湊的式子：
//   g_i(x) = (x mod m + i * (1 + x mod (m - 2))) mod m
typedef struct {
    unsigned char *bit;     // 一個位元用一個 unsigned char 存
    int m, k;
} Bloom;

int bloom_init(Bloom *b, int m, int k)
{
    b->bit = calloc(m, 1);              // 全部是 0
    b->m = m;
    b->k = k;
    return b->bit != NULL;
}

int g(const Bloom *b, long long x, int i)
{
    long long step = 1 + x % (b->m - 2);
    return (int)((x % b->m + i * step) % b->m);
}

void bloom_insert(Bloom *b, long long x)
{
    for (int i = 0; i < b->k; i++) {
        b->bit[g(b, x, i)] = 1;
    }
}

// 回傳 1 表示「可能在」，0 表示「一定不在」。
int bloom_query(const Bloom *b, long long x)
{
    for (int i = 0; i < b->k; i++) {
        if (b->bit[g(b, x, i)] == 0) {
            return 0;
        }
    }
    return 1;
}

// 和 probe_count.c 相同的線性同餘亂數，種子固定。
static unsigned long long seed;

long long next_key(void)
{
    seed = seed * 6364136223846793005ULL + 1442695040888963407ULL;
    return (long long)(seed >> 33);
}

int main(void)
{
    // 手算的例子：m = 13，k = 3，四個停權卡號。
    const long long banned[] = {2037, 5512, 6180, 7423};
    const long long ask[] = {6180, 3141, 4126};
    Bloom b;
    if (!bloom_init(&b, 13, 3)) {
        return 1;
    }
    for (int i = 0; i < 4; i++) {
        bloom_insert(&b, banned[i]);
    }
    printf("bits:");
    for (int j = 0; j < b.m; j++) {
        printf(" %d", b.bit[j]);
    }
    printf("\n");
    for (int i = 0; i < 3; i++) {
        printf("%lld:", ask[i]);
        for (int t = 0; t < b.k; t++) {
            printf(" %d", g(&b, ask[i], t));
        }
        printf(" -> %s\n", bloom_query(&b, ask[i]) ? "maybe" : "no");
    }
    free(b.bit);

    // 實測：m = 10007，放 n = 1000 個偶數鍵，再問 100000 個奇數鍵。
    const int m = 10007, n = 1000, trials = 100000;
    printf("m = %d, n = %d\n", m, n);
    printf(" k  measured  formula\n");
    for (int k = 1; k <= 10; k++) {
        if (!bloom_init(&b, m, k)) {
            return 1;
        }
        seed = 20261119ULL;             // 每個 k 用同一串鍵
        for (int i = 0; i < n; i++) {
            bloom_insert(&b, next_key() & ~1LL);
        }
        int fp = 0;                     // 奇數鍵一定不在集合裡
        for (int t = 0; t < trials; t++) {
            fp += bloom_query(&b, next_key() | 1);
        }
        double p = pow(1 - exp(-(double)k * n / m), k);
        printf("%2d  %8.4f  %7.4f\n", k, (double)fp / trials, p);
        free(b.bit);
    }
    return 0;
}
\end{example}
\iosample{}{bits: 1 0 1 0 1 1 0 1 0 1 1 0 1 \\ 6180: 5 2 12 -> maybe \\ 3141: 8 2 9 -> no \\ 4126: 5 7 9 -> maybe \\ m = 10007, n = 1000 \\ \ k\ \ measured\ \ formula \\ \ 1\ \ \ \ 0.0959\ \ \ 0.0951 \\ \ 2\ \ \ \ 0.0331\ \ \ 0.0328 \\ \ 3\ \ \ \ 0.0185\ \ \ 0.0174 \\ \ 4\ \ \ \ 0.0120\ \ \ 0.0118 \\ \ 5\ \ \ \ 0.0101\ \ \ 0.0094 \\ \ 6\ \ \ \ 0.0089\ \ \ 0.0084 \\ \ 7\ \ \ \ 0.0086\ \ \ 0.0082 \\ \ 8\ \ \ \ 0.0086\ \ \ 0.0084 \\ \ 9\ \ \ \ 0.0096\ \ \ 0.0091 \\ 10\ \ \ \ 0.0109\ \ \ 0.0101}

前四行就是圖~\ref{fig:hash-bloom}：位元陣列，和三個查詢各自算出的位置與回答。 實測的
部分，$m/n \approx 10$，由定理~\ref{thm:hash-bloom-k}，最好的 $k$ 是 $10.007 \ln 2
\approx 6.94$，取 $k = 7$，公式值 $0.82\%$，實測 $0.86\%$；$k = 1$ 時偽陽性率將近
$10\%$，$k = 10$ 又回升到 $1\%$ 以上，和上面「$k$ 太小、太大都不好」的分析一致。 實測都
比公式高一點，原因是公式把各個位元當成互相獨立，又假設 $k$ 個位置完全亂選；這裡的 $k$
個位置由兩個雜湊值組出來，也不是完全獨立。 差距在 $10\%$ 以內，拿來決定 $m$ 和 $k$ 已
經夠用。

最後把布隆過濾器和雜湊表放在一起比。 雜湊表存了鍵本身，能回答「在不在」而且不會錯，
能刪除，也能把鍵列出來；布隆過濾器每個鍵只用幾個位元，和鍵的長度無關，但只能回答
「一定不在」或「可能在」，不能刪除，也列不出存了哪些鍵。 適合它的場合是：集合很大、
大部分查詢的答案是「不在」、偶爾誤判只是多花一點時間（多問一次伺服器、多讀一次磁碟），
不會造成錯誤的結果。

\section{本章常見錯誤}
\label{sec:hash-pitfalls}

\begin{center}
\small
\begin{tabular}{p{0.24\linewidth}p{0.50\linewidth}p{0.16\linewidth}}
\toprule
錯誤 & 例子 & 見 \\
\midrule
除法的 $m$ 沒選好 & $m = 100$ 只用到最後兩位；鍵都是 $4$ 的倍數時用 $m = 12$ &
第~\ref{sec:hash-func}~節 \\
負數取餘數 & \texttt{-7 \% 13} 是 $-7$；要寫 \texttt{((k \% m) + m) \% m} &
第~\ref{sec:hash-func}~節 \\
探測次數少算一次 & 只數碰撞，沒算最後放進去或停下來的那一格 &
第~\ref{sec:hash-open}~節 \\
二次探測當成走得遍 & $m$ 不是質數，或 $\alpha > 1/2$，以為一定插得進去 &
第~\ref{sec:hash-open}~節 \\
$h_2(k)$ 選錯 & $h_2(k) = k \bmod 7$ 可能是 $0$，原地不動；$m = 12$、$h_2(k) = 4$
不互質，只走得到三格 & 第~\ref{sec:hash-open}~節 \\
刪除時清成空格 & 之後經過這一格才放進去的鍵都找不到了 & 第~\ref{sec:hash-tomb}~節 \\
插入遇到刪除標記就放 & 沒有確認後面有沒有同一個鍵，表裡出現兩個 $k$ &
第~\ref{sec:hash-tomb}~節 \\
公式張冠李戴 & 線性探測套用 $1/(1-\alpha)$；成功和失敗的公式對調 &
第~\ref{sec:hash-analysis}~節 \\
布隆過濾器清位元 & 刪除時把鍵的 $k$ 個位元清成 $0$；這些位元可能也屬於別的鍵，之後查那個鍵會
答「不在」 & 第~\ref{sec:hash-bloom}~節 \\
「可能在」當成「在」 & 布隆過濾器只保證「不在」一定對；回答「可能在」還要再確認 &
第~\ref{sec:hash-bloom}~節 \\
\bottomrule
\end{tabular}
\end{center}

\section{觀念題}
\label{sec:hash-exercises}

題型和研究所考試、期中期末考的筆試相近，有選擇、計算、追蹤和證明。 除非題目另外說
明，探測次數照本章的數法（看過幾格就算幾次，包括最後一格），$i$ 從 $0$ 開始。 解答
在本章最後。

\begin{enumerate}
\item （選擇，可能不只一個答案）下列哪些敘述正確？
\begin{enumerate}[nosep]
\item 鏈結法的載入因子可以大於 $1$。
\item 線性探測裡，起始位置不同的兩個鍵，探測序列不會經過同一格。
\item 二次探測沒有主要群聚，但起始位置相同的鍵，探測序列完全相同。
\item $m = 16$ 的雙重雜湊，只要 $h_2(k)$ 一定是奇數，每個鍵的探測序列都會走遍 $16$ 格。
\item 開放定址刪除一個鍵時，把那一格改回空的，之後的搜尋仍然正確。
\item 在均勻雜湊下，$\alpha = 0.8$ 的表做一次失敗搜尋，期望探測次數小於 $5$。
\end{enumerate}

\item （計算）鍵 $k = 6083$。
\begin{enumerate}[nosep]
\item 用除法，$m = 97$，求 $h(k)$。
\item 用乘法，$m = 1000$，$A = 0.6180339887$，求 $h(k)$。
\item 用中間平方法，$m = 100$：$k^2$ 補成八位數，取第 $4$、$5$ 位。
\item 另一個鍵 $526803147$ 切成三位數一段，$m = 1000$，分別用移位摺疊和邊界摺疊求
$h$。
\end{enumerate}

\item （追蹤）$m = 11$，$h(k) = k \bmod 11$，用線性探測依序插入
$63, 85, 75, 54, 96, 24, 33$。
\begin{enumerate}[nosep]
\item 寫出每個鍵試過的格子和最後的表。
\item 成功搜尋的平均探測次數是多少？ 假設起始位置是 $0$ 到 $10$ 的機會一樣，失敗搜
尋的平均探測次數是多少？
\item 接著刪除 $54$（用刪除標記），再依序搜尋 $33$、搜尋 $96$、插入 $44$。 寫出每
個操作的探測次數，以及 $44$ 放在哪一格。
\end{enumerate}

\item （追蹤）用二次探測依序插入 $5, 16, 27, 38, 49, 60, 71$，$m = 11$，
$h(k, i) = (k \bmod 11 + i^2) \bmod 11$。 寫出每個鍵試過的格子。 最後一個鍵插得進去嗎？ 這時表裡
還有幾個空格？ 用定理~\ref{thm:hash-quad} 解釋。

\item （追蹤）$m = 11$，雙重雜湊 $h_1(k) = k \bmod 11$、
$h_2(k) = 1 + (\lfloor k / 11 \rfloor \bmod 10)$，依序插入第 3 題的七個鍵。 寫出每個
鍵的 $h_2$、試過的格子和最後的表，並和第 3 題比較探測的總次數。

\item （計算）$m = 8$，$h(k, i) = (h(k) + i^2) \bmod 8$。
\begin{enumerate}[nosep]
\item 列出 $i = 0, 1, \ldots, 7$ 的 $i^2 \bmod 8$。 起始位置是 $0$ 的鍵，最多試得到
幾個不同的格子？
\item 舉一個例子：表裡還有空格，插入卻失敗。
\item 改成 $h(k, i) = (h(k) + i(i+1)/2) \bmod 8$，重做 (a)。
\end{enumerate}

\item （證明）$m = 2^p$，$p \ge 1$。 證明 $h(k, i) = (h(k) + i(i+1)/2) \bmod m$，
$i = 0, 1, \ldots, m - 1$ 兩兩不同。（提示：$j - i$ 和 $i + j + 1$ 一個是奇數、一個是
偶數。）

\item （計算）鏈結法，$m = 100$ 格、$n = 100$ 個鍵，簡單均勻雜湊。
\begin{enumerate}[nosep]
\item 空的格子期望有幾個？（提示：對每一格寫一個「這一格是空的」的指示變數。）
\item 成功搜尋和失敗搜尋的期望比較次數各是多少？
\item 平均每條串列長 $1$，為什麼還有大約三分之一的格子是空的？
\end{enumerate}

\item （計算）希望失敗搜尋的期望探測次數不超過 $3$。 用均勻雜湊的 $1/(1-\alpha)$ 和
線性探測的 Knuth 公式，各求 $\alpha$ 最大可以是多少。 $m = 1000$ 時，兩種情形各大
約可以放幾個鍵？ 和範例~\ref{ex:hash-probe-count} 的輸出比對。

\item （追蹤）接著圖~\ref{fig:hash-extendible}(d) 的狀態，依序插入
H（\texttt{00011}）和 I（\texttt{00101}）。 每插入一個，說明有沒有拆桶、目錄有沒有加
倍，並畫出最後的目錄和桶。 最後有幾個目錄項指向 D、F 那個桶？

\item （計算）$U = \{a, b, c, d\}$，$m = 2$，$\mathcal{H} = \{h_1, h_2, h_3\}$：
\begin{center}
\begin{tabular}{ccccc}
\toprule
 & $a$ & $b$ & $c$ & $d$ \\
\midrule
$h_1$ & 0 & 0 & 1 & 1 \\
$h_2$ & 0 & 1 & 0 & 1 \\
$h_3$ & 1 & 0 & 0 & 1 \\
\bottomrule
\end{tabular}
\end{center}
\begin{enumerate}[nosep]
\item $\mathcal{H}$ 是通用的嗎？
\item 把 $h_3$ 換成 $h_3'(a, b, c, d) = (0, 0, 0, 1)$，還是通用的嗎？
\end{enumerate}

\item \label{q:hash-bloom-trace}（追蹤）布隆過濾器 $m = 11$、$k = 2$，$h_1(x) = x \bmod 11$，
$h_2(x) = \lfloor x / 11 \rfloor \bmod 11$，依序插入 $25, 47, 63, 90$。
\begin{enumerate}[nosep]
\item 寫出每個鍵的兩個位置和最後的位元陣列 $B[0..10]$。
\item 查詢 $36$、$58$、$80$，各回答什麼？ 哪些是偽陽性？
\item 想刪除 $63$，把它的兩個位元清成 $0$。 之後查詢哪一個插入過的鍵會答錯？
\item 改成計數型布隆過濾器，寫出插入四個鍵之後的計數器 $C[0..10]$，以及刪除 $63$
（兩個計數器各減一）之後的計數器。 這時查詢 $90$、$58$ 各回答什麼？
\end{enumerate}

\item （計算）布隆過濾器要存 $n = 100$ 個鍵，用 $m = 1200$ 個位元。 用
式~\eqref{eq:hash-bloom-fp} 估計。
\begin{enumerate}[nosep]
\item 最好的整數 $k$ 是多少？ 這時偽陽性率大約多少？ 大約多少比例的位元是 $0$？
\item 想讓偽陽性率降到 $0.1\%$ 以下，$n = 100$ 時大約要多少個位元？ $k$ 取多少？
\item 有人說「$k$ 越大，查詢時檢查的位元越多，誤判一定越少」。 用 (a) 的數字說明這句話
哪裡不對。
\end{enumerate}

\end{enumerate}

\section{歷屆試題}
\label{sec:hash-pastexam}

以下是研究所入學考試的原題，題目照原卷（語言、記號、配分照錄）。 詳解在本章最後。

\begin{pastexam}[label={px:hash-1}]{成大資訊聯招 110 學年度 第 3 題}
(10\%) Assume that you have to store the four strings: ``abcdef'', ``bcdefa'', ``cdefab'', and ``defabc'', in the hash table by using the hashing technique.
\begin{enumerate}[label=(\arabic*),itemsep=0pt,topsep=2pt]
\item (5\%) Consider the hash function: the index for a specific string will be equal to the sum of the ASCII values of the characters modulo 11. Please determine how many collisions if we use this hash function.
\item (5\%) Consider the hash function: the index for a specific string will be equal to sum of ASCII values of characters multiplied by their respective order in the string after which it is modulo with 11. Please determine how many collisions if we use this hash function.
\end{enumerate}
Note: The ASCII values of a, b, c, d, e, and f are 97, 98, 99, 100, 101, and 102 respectively.
\end{pastexam}

\begin{pastexam}[label={px:hash-2}]{政大資科所 115 學年度 第 13 題}
Consider a hash table of size $M=20$ (slots indexed from 0 to 19). The primary hash function is defined as: $h_1(k)=k \mathit{\ mod\ } 20$. The secondary hash function is defined as: $h_2(k)=(k \mathit{\ mod\ } 7)+1$. Assume the hash table is initially empty before each part of the question. We will insert the following sequence of keys S in the order given:
\[ S=[8,38,3,5,28,18,65,83,25] \]
(a) (5\%) Assuming we use double hashing for collision resolution, determine the final structure of the hash table after all keys in S have been inserted. Note: The probe sequence is given by:
\[ h(k,i)=\big(h_1(k)+i\cdot h_2(k)\big) \mathit{\ mod\ } M \]
where $i$ is the probe number.

(b) (5\%) If we search for 48, what is the total number of probes made during the search?

(c) (10\%) Suppose we are given $n$ keys, $m$ hash table slots, and two uniform hash functions $h_a$ and $h_b$. Further suppose hashing scheme uses $h_a$ for the odd keys and $h_b$ for the even keys. What is the expected number of keys in a slot?
\end{pastexam}

\begin{pastexam}[label={px:hash-3}]{中央資工所 115 學年度 第 12 題}
（多選題，每一個選項單獨計分，每題 5 分）
Given an empty 11-bucket hash table with hash function h(k)=k\%11. Below is the resulting hash table after inserting 8 items.
\begin{center}
\begin{tabular}{|l|l|}
\hline
0 & \\ \hline
1 & \\ \hline
2 & 24 \\ \hline
3 & 57 \\ \hline
4 & 37 \\ \hline
5 & 27 \\ \hline
6 & 16 \\ \hline
7 & 17 \\ \hline
8 & 15 \\ \hline
9 & 38 \\ \hline
10 & \\ \hline
\end{tabular}
\end{center}
Which of the following statements is/are true?
\begin{enumerate}[label=\Alph*.]
\item 27 37 16 17 24 15 57 38 is a possible input sequence when quadratic probing is adopted.
\item 37 27 16 24 15 57 38 17 is a possible input sequence when quadratic probing is adopted.
\item 37 27 16 24 57 17 15 38 is a possible input sequence when liner probing is adopted.
\item 27 16 37 38 57 15 17 24 is a possible input sequence when liner probing is adopted.
\end{enumerate}
\end{pastexam}

\begin{pastexam}[label={px:hash-4}]{臺大資工所 115 學年度 第 15 題}
(4 points) Most hash table implementations (such as Java's \texttt{HashMap}) resize the underlying array when the load factor $\alpha$ exceeds a certain threshold (typically 0.75). Why is this resizing operation necessary for a hash table that uses separate chaining?
\begin{enumerate}[label=(\Alph*),leftmargin=2.2em,itemsep=1pt]
\item \textbf{Saturation:} To prevent the table from reaching a load factor of 1.0, at which point inserting a new element would be impossible.
\item \textbf{Ordering:} To re-sort the keys alphabetically during the transfer, enabling the use of binary search for future lookups.
\item \textbf{Complexity:} To prevent the average length of the linked lists (collision chains) from growing linearly with the number of elements, thereby maintaining $O(1)$ average access time.
\item \textbf{Uniqueness:} To redistribute the keys across a larger range so that all existing collisions are permanently eliminated.
\item \textbf{Memory contiguity:} To ensure that all elements in the linked lists are stored in adjacent memory addresses, minimizing CPU cache misses and memory fragmentation.
\end{enumerate}
\end{pastexam}

\begin{pastexam}[label={px:hash-5}]{成大資訊聯招 113 學年度 第 5 題}
\textbf{[3\%, 2\%]} Given the following bloom filter.
\begin{center}
\small
\begin{tabular}{|*{15}{c|}}
\hline
[0] & [1] & [2] & [3] & [4] & [5] & [6] & [7] & [8] & [9] & [10] & [11] & [12] & [13] & [14] \\ \hline
1 & 1 & 0 & 1 & 1 & 1 & 1 & 1 & 0 & 1 & 0 & 0 & 0 & 1 & 0 \\ \hline
\end{tabular}
\end{center}
\begin{enumerate}[label=\roman*)]
\item Assume that the bloom filter has 3 hash functions $f_1(k)$, $f_2(k)$, and $f_3(k)$:
\[ f_1(k)=(3k) \bmod m,\qquad f_2(k)=(2k+1) \bmod m,\qquad f_3(k)=k^2 \bmod m, \]
where $k$ represents the key and $m$ is the size of bit array for the bloom filter. Which of the following statements is true?
\begin{enumerate}[label=(\Alph*)]
\item When inserting $k$=12, the bits in the bloom filter are not changed.
\item To delete $k$=2, the bits at indices 4, 5, and 6 are changed to 0.
\item When querying the bloom filter for $k$=15, the answer is ``yes.''
\item When querying the bloom filter for $k$=8, the answer is ``yes.''
\item None of the above.
\end{enumerate}
\item Considering that the elements in the bloom filter is 5. To minimize the false positive probability, what is the ideal optimum number of hash functions?
\begin{enumerate}[label=(\Alph*)]
\item 15
\item 5
\item 3*ln(2)
\item 5*ln(2)
\item None of the above
\end{enumerate}
\end{enumerate}
\end{pastexam}


\section{實作練習}
\label{sec:hash-practice}

下面是 Aki-Judge 上和本章相關的題目。

\begin{center}
\begin{tabular}{llc}
\toprule
小節 & 題目 & 難度 \\
\midrule
鏈結法 & \akiref{integer-hash-set-with-chaining}{用鏈結法實作整數集合} & 基本 \\
鏈結法 & \akiref{integer-hash-map-with-chaining}{用鏈結法實作整數對映} & 基本 \\
開放定址 & \akiref{open-addressing-integer-set}{線性探測的整數集合} & 基本 \\
刪除標記 & \akiref{linear-probing-with-deletion-marks}{線性探測雜湊表與刪除標記} & 進階 \\
再雜湊 & \akiref{parking-plate-registry-with-move-to-front-chains}{停車場車牌登記：前移串鏈與擴表} & 進階 \\
應用 & \akiref{number-of-distinct-values}{不同數值的個數} & 暖身 \\
應用 & \akiref{most-frequent-value}{出現次數最多的數} & 暖身 \\
應用 & \akiref{first-value-appearing-once}{第一個只出現一次的數} & 暖身 \\
\bottomrule
\end{tabular}
\end{center}

前三題只交函式，位置都規定用
\texttt{((key \% M) + M) \% M} 算，鍵可能是負數。 「線性探測雜湊表與刪除標記」的
規則和第~\ref{sec:hash-tomb}~節相同，可以先用圖~\ref{fig:hash-tomb} 的例子對一次。 最後三題是計數：把範例~\ref{ex:hash-chain} 的「餘額」改成「出現次數」就能做。

\section{讀完你應該會}
\label{sec:hash-checklist}

\begin{itemize}
\item 能用除法、乘法、中間平方、摺疊算出雜湊值，並說明除法的 $m$ 為什麼要選質數。
\item 能寫出鏈結法的三個操作，並用指示變數推導成功與失敗搜尋的期望比較次數。
\item 能對一串鍵做線性探測、二次探測、雙重雜湊的追蹤，寫出每個鍵試過的格子和最後的
表，並指出主要群聚和次要群聚。
\item 能說明刪除為什麼要用標記，正確地處理標記上的搜尋和插入。
\item 能在均勻雜湊下推導失敗搜尋 $\frac{m+1}{m-n+1} < \frac{1}{1-\alpha}$ 和成功
搜尋 $\le \frac{1}{\alpha}\ln\frac{1}{1-\alpha}$，記得線性探測的 Knuth 公式，並用它
們算出該把 $\alpha$ 控制在多少。
\item 能追蹤可延伸雜湊的插入，並說明通用雜湊的定義和它保證了什麼。
\item 能做布隆過濾器的插入與查詢，說明它為什麼只有偽陽性、為什麼不能直接刪除，推導偽陽性
率 $(1 - e^{-kn/m})^k$ 和最好的 $k = \frac{m}{n}\ln 2$。
\end{itemize}

\textbf{延伸閱讀}（書末會整理完整的參考資料）

\begin{itemize}
\item Thomas H. Cormen 等，\textit{Introduction to Algorithms}，第 4 版，第 11 章：
兩種方法的分析、通用雜湊族的證明，以及不用刪除標記的線性探測刪除。
\item Donald E. Knuth，\textit{The Art of Computer Programming}，第 3 冊，第 2 版，
第 6.4 節：雜湊的原始整理，包括線性探測公式的完整推導和各種探測法的比較。
\item Ronald Fagin、Jürg Nievergelt、Nicholas Pippenger、H. Raymond Strong，
``Extendible Hashing---A Fast Access Method for Dynamic Files''，\textit{ACM
Transactions on Database Systems} 4(3)，1979：可延伸雜湊的原始論文。
\item J. Lawrence Carter、Mark N. Wegman，``Universal Classes of Hash Functions''，
\textit{Journal of Computer and System Sciences} 18(2)，1979：通用雜湊的原始論文。
\item Steven S. Skiena，\textit{The Algorithm Design Manual}，第 3 版，第 3.7 節與
第 6 章：雜湊在實務上的用法，以及用機率看雜湊的角度。
\item Burton H. Bloom，``Space/Time Trade-offs in Hash Coding with Allowable Errors''，
\textit{Communications of the ACM} 13(7)，1970：布隆過濾器的原始論文。
\item Andrei Broder、Michael Mitzenmacher，``Network Applications of Bloom Filters: A
Survey''，\textit{Internet Mathematics} 1(4)，2004：偽陽性率的推導、計數型布隆過濾器和
各種應用。
\end{itemize}

\addcontentsline{toc}{section}{觀念題解答}
\section*{觀念題解答}

\begin{enumerate}[itemsep=0.15em]
\item (a)、(c)、(d)、(f)。 (b)：遇到同一段群聚就會走同樣的格子，這正是主要群聚。
(d)：奇數和 $16$ 互質（定理~\ref{thm:hash-double}）。 (e)：會切斷經過這一格的探測序
列。 (f)：精確值 $\frac{m+1}{m-n+1} < \frac{1}{1-0.8} = 5$。

\item (a) $6083 = 97 \times 62 + 69$，$h = 69$。 (b) $6083A = 3759.50075\ldots$，小數部
分乘以 $1000$ 是 $500.75\ldots$，$h = 500$。 (c) $6083^2 = 37002889$，第 $4$、$5$ 位
是 $0$、$2$，$h = 2$。 (d) 移位：$526 + 803 + 147 = 1476$，$h = 476$；
邊界：$526 + 308 + 147 = 981$，$h = 981$。

\item (a) 起始位置依序是 $8, 8, 9, 10, 8, 2, 0$，試過的格子依序是：$63$：8；
$85$：8, 9；$75$：9, 10；$54$：10, 0；$96$：8, 9, 10, 0, 1；$24$：2；$33$：0, 1, 2, 3。
最後的表（第 $0$ 到 $10$ 格）：$54, 96, 24, 33, -, -, -, -, 63, 85, 75$。 (b) 成功：$(1+2+2+2+5+1+4)/7 = 17/7 \approx 2.43$。 失敗：起始位置 $0$ 到 $3$ 分別要
$5, 4, 3, 2$ 次，$4$ 到 $7$ 各 $1$ 次，$8, 9, 10$ 分別要 $8, 7, 6$ 次，合計 $39$，平均
$39/11 \approx 3.55$。 (c) 刪除 $54$：試 $10, 0$，第 $0$ 格改成標記。 搜尋 $33$：$0$（標記，繼續）、$1, 2,
3$，找到，$4$ 次。 搜尋 $96$：$8, 9, 10, 0$（標記）、$1$，找到，$5$ 次。 插入 $44$
（$h = 0$）：第 $0$ 格是標記，記下來，繼續看 $1, 2, 3$，第 $4$ 格是空的，停；$44$
放進第 $0$ 格，共探測 $5$ 次。

\item 七個鍵的起始位置都是 $5$。 $i^2 \bmod 11$ 依序是 $0, 1, 4, 9, 5, 3, 3, 5, 9, 4,
1$，所以探測序列是 $5, 6, 9, 3, 10, 8, 8, 10, 3, 9, 6$。 前六個鍵依序放在
$5, 6, 9, 3, 10, 8$。 $71$ 整條序列都是滿的，插入失敗，但還有 $5$ 個空格
（第 $0, 1, 2, 4, 7$ 格）。 由定理~\ref{thm:hash-quad}，
二次探測只看得到 $(11+1)/2 = 6$ 個不同的格子，這六格都被起始位置相同的鍵佔了。

\item $h_2$ 依序是 $6, 8, 7, 5, 9, 3, 4$。 試過的格子：$63$：8；$85$：8, 5；
$75$：9；$54$：10；$96$：8, 6；$24$：2；$33$：0。
最後的表：$33, -, 24, -, -, 85, 96, -, 63, 75, 54$。 共探測 $9$ 次，線性探測是 $17$ 次。

\item (a) $0, 1, 4, 1, 0, 1, 4, 1$，只有 $0, 1, 4$ 三種，起始位置 $0$ 的鍵只試得到第
$0, 1, 4$ 格。 (b) 依序插入三個起始位置都是 $0$ 的鍵（例如 $0, 8, 16$），佔滿第 $0,
1, 4$ 格；再插入 $24$，八次探測都落在這三格，失敗，但表裡還有 $5$ 個空格。 (c)
$i(i+1)/2$ 依序是 $0, 1, 3, 6, 10, 15, 21, 28$，除以 $8$ 的餘數是
$0, 1, 3, 6, 2, 7, 5, 4$，八格都走到了。

\item 假設 $0 \le i < j \le m - 1$ 而兩個位置相同，則 $2^p$ 整除
$\frac{j(j+1)}{2} - \frac{i(i+1)}{2} = \frac{(j - i)(j + i + 1)}{2}$，所以 $2^{p+1}$
整除 $(j - i)(j + i + 1)$。 兩個因數的和 $2j + 1$ 是奇數，所以一奇一偶，$2^{p+1}$
必須整除偶數的那個。 但
$1 \le j - i \le m - 1 < 2^{p+1}$，$1 \le j + i + 1 \le 2m - 2 < 2^{p+1}$，兩個都是比
$2^{p+1}$ 小的正數，不可能被它整除，矛盾。

\item (a) 第 $j$ 格是空的，表示 $100$ 個鍵都沒有落在它，機率是 $(1 - 1/100)^{100}$。
對 $100$ 格的指示變數取期望值再相加，得到 $100 \times 0.99^{100} \approx 36.6$ 個。
(b) 成功：$1 + 99/200 = 1.495$。 失敗：$\alpha = 1$。 (c) 有的串列長 $2$、$3$ 甚至更
長，平均要是 $1$，就要有對應數量的空串列。

\item 均勻雜湊：$\frac{1}{1-\alpha} \le 3$，$\alpha \le 2/3$，約可放 $666$ 個。 線性探
測：$\frac12\bigl(1 + \frac{1}{(1-\alpha)^2}\bigr) \le 3$，$(1-\alpha)^2 \ge 1/5$，
$\alpha \le 1 - 1/\sqrt5 \approx 0.553$，約可放 $552$ 個。 範例~\ref{ex:hash-probe-count}
的實測（線性 $\alpha = 0.5$ 時 $2.49$；雙重 $0.75$ 時 $3.80$）和這兩個界相符。

\item H = \texttt{000}11 查目錄 \texttt{000}，指向 A 的桶，還有空位，放進去：\{A, H\}，
沒有拆桶。 I = \texttt{001}01 查目錄 \texttt{001}，也指向這個桶，滿了；局部深度
$2 < 3$，目錄不變，只拆桶。 依第 $3$ 個位元分：H = 00\texttt{0}11 到 \texttt{000}，
A = 00\texttt{1}10 和 I = 00\texttt{1}01 到 \texttt{001}，兩個桶的局部深度都變成 $3$，
其他桶和目錄不變。 D、F 的桶局部深度 $2$，有 $2^{3-2} = 2$ 個
目錄項指向它。

\item (a) 是。 $|\mathcal{H}|/m = 1.5$，每一對鍵最多只能被一個函數撞。 $h_1$ 撞
$\{a, b\}$、$\{c, d\}$，$h_2$ 撞 $\{a, c\}$、$\{b, d\}$，$h_3$ 撞 $\{a, d\}$、
$\{b, c\}$，六對各一次。 (b) 不是。 $h_3'$ 撞 $\{a, b\}$、$\{a, c\}$、$\{b, c\}$，
$\{a, b\}$ 被 $h_1$、$h_3'$ 撞了 $2$ 次。

\item (a) 位置：$25$：$3, 2$；$47$：$3, 4$；$63$：$8, 5$；$90$：$2, 8$。 最後第
$2, 3, 4, 5, 8$ 格是 $1$：
\[
B[0..10] = 0, 0, 1, 1, 1, 1, 0, 0, 1, 0, 0.
\]
(b) $36$：$36 \bmod 11 = 3$，$\lfloor 36/11 \rfloor = 3$，兩個位置都是第 $3$ 格，是
$1$，回答「可能在」，偽陽性。 $58$：位置 $3, 5$，都是 $1$，「可能在」，偽陽性。 $80$：
位置 $3, 7$，第 $7$ 格是 $0$，「不在」，正確。（$36$ 的兩個位置重合，等於只檢查一個位
元，這是 $h_1$、$h_2$ 不夠獨立的後果。）
(c) 清掉第 $8$、$5$ 格之後，$90$ 的位置是 $2, 8$，第 $8$ 格變成 $0$，回答「不在」，但
$90$ 還在集合裡：偽陰性。 $25$、$47$ 的位置沒有被清掉，仍然答對。
(d) 插入四個鍵之後，第 $2$ 格由 $25$、$90$ 各加一，第 $3$ 格由 $25$、$47$，第 $8$ 格由
$63$、$90$，第 $4$、$5$ 格各一次；刪除 $63$ 時第 $8$、$5$ 格各減一：
\[
\begin{aligned}
\text{插入後：}\ C[0..10] &= 0, 0, 2, 2, 1, 1, 0, 0, 2, 0, 0, \\
\text{刪除 $63$ 後：}\ C[0..10] &= 0, 0, 2, 2, 1, 0, 0, 0, 1, 0, 0.
\end{aligned}
\]
查 $90$：$C[2] = 2$、$C[8] = 1$ 都大於 $0$，「可能
在」，正確。 查 $58$：$C[5] = 0$，「不在」；刪除 $63$ 之後，原本的偽陽性也消失了。

\item (a) $m/n = 12$，$k^{*} = 12 \ln 2 \approx 8.32$。 比較兩個整數：$k = 8$ 時
$(1 - e^{-8/12})^8 = (1 - 0.5134)^8 \approx 0.00314$；$k = 9$ 時
$(1 - e^{-9/12})^9 = (1 - 0.4724)^9 \approx 0.00317$。 取 $k = 8$，偽陽性率約
$0.31\%$，這時 $0$ 位元的比例約 $e^{-8/12} \approx 51\%$，接近定理~\ref{thm:hash-bloom-k}
說的一半。
(b) 用最佳的 $k$，$\varepsilon = 2^{-k^{*}} \le 0.001$，$k^{*} \ge \lg 1000 \approx 9.97$，
每個鍵要 $k^{*}/\ln 2 = \ln 1000 / (\ln 2)^2 \approx 14.4$ 個位元，$n = 100$ 時大約
$1438$ 個位元，$k$ 取 $10$。 驗算：$m = 1438$、$k = 10$，$(1 - e^{-1000/1438})^{10}
\approx 0.000999$，剛好在 $0.1\%$ 以下。
(c) $m$、$n$ 固定時，$k$ 越大，被設成 $1$ 的位元越多。 $k = 8$ 時約 $49\%$ 的位元是 $1$，
$k = 20$ 時 $1 - e^{-20/12} \approx 81\%$ 是 $1$，偽陽性率變成 $0.81^{20} \approx 1.5\%$，
是 $k = 8$ 時 $0.31\%$ 的將近五倍。 檢查的位元多了，但每個位元是 $1$ 的機率也變大，兩
者要取平衡，最好的 $k$ 讓一半的位元是 $0$。
\end{enumerate}

\addcontentsline{toc}{section}{歷屆試題詳解}
\section*{歷屆試題詳解}

\pastsol{px:hash-1}{成大資訊聯招 110 學年度 第 3 題}
題目沒有說「幾次碰撞」怎麼數。 這裡照最常見的算法：依題目給的順序把字串放進表裡，放進去時
起始位置已經有別的字串，就算一次碰撞。 另一種數法是數「落在同一格的字串有幾對」，答案寫在
括號裡。

(1) 四個字串都是 a 到 f 各用一次，只是順序不同，字元碼的和都是
\[
97 + 98 + 99 + 100 + 101 + 102 = 597, \qquad 597 = 11 \times 54 + 3,
\]
四個字串的位置都是 $3$。 第一個放進去沒有碰撞，後面三個都撞上，共 $3$ 次碰撞（落在同一格
的有 $\binom{4}{2} = 6$ 對）。 這就是第~\ref{sec:hash-func}~節說的：只把字元碼加起來，用
同一組字母的字串一定撞在一起。

(2) 第 $i$ 個字元乘上 $i$（$i = 1, \ldots, 6$）再加起來：
\begin{center}
\small
\begin{tabular}{lll}
\toprule
字串 & 加權和 & 除以 $11$ 的餘數 \\
\midrule
abcdef & $1 \cdot 97 + 2 \cdot 98 + 3 \cdot 99 + 4 \cdot 100 + 5 \cdot 101 + 6 \cdot 102 = 2107$ & $2107 - 11 \times 191 = 6$ \\
bcdefa & $1 \cdot 98 + 2 \cdot 99 + 3 \cdot 100 + 4 \cdot 101 + 5 \cdot 102 + 6 \cdot 97 = 2092$ & $2092 - 11 \times 190 = 2$ \\
cdefab & $1 \cdot 99 + 2 \cdot 100 + 3 \cdot 101 + 4 \cdot 102 + 5 \cdot 97 + 6 \cdot 98 = 2083$ & $2083 - 11 \times 189 = 4$ \\
defabc & $1 \cdot 100 + 2 \cdot 101 + 3 \cdot 102 + 4 \cdot 97 + 5 \cdot 98 + 6 \cdot 99 = 2080$ & $2080 - 11 \times 189 = 1$ \\
\bottomrule
\end{tabular}
\end{center}
四個位置 $6, 2, 4, 1$ 都不同，$0$ 次碰撞（$0$ 對）。 乘上位置以後，字元的順序會影響結果，
所以能分開這幾個字串；第~\ref{sec:hash-func}~節用 $B$ 進位（Horner 法）是同樣的想法，而且
權重更分散。

\pastsol{px:hash-2}{政大資科所 115 學年度 第 13 題}
(a) 每個鍵算出 $h_1(k) = k \bmod 20$、$h_2(k) = (k \bmod 7) + 1$，再依序試
$h_1, h_1 + h_2, h_1 + 2h_2, \ldots$（都取 $\bmod 20$）。 $i$ 從 $0$ 開始。
\begin{center}
\small
\begin{tabular}{rcccl}
\toprule
鍵 $k$ & $h_1(k)$ & $k \bmod 7$ & $h_2(k)$ & 依序試過的格子（最後一個是放進去的位置） \\
\midrule
8 & 8 & 1 & 2 & 8 \\
38 & 18 & 3 & 4 & 18 \\
3 & 3 & 3 & 4 & 3 \\
5 & 5 & 5 & 6 & 5 \\
28 & 8 & 0 & 1 & 8（有 8），9 \\
18 & 18 & 4 & 5 & 18（有 38），3（有 3），8（有 8），13 \\
65 & 5 & 2 & 3 & 5（有 5），8（有 8），11 \\
83 & 3 & 6 & 7 & 3（有 3），10 \\
25 & 5 & 4 & 5 & 5（有 5），10（有 83），15 \\
\bottomrule
\end{tabular}
\end{center}
最後的表（沒列出的格子是空的）：
\begin{center}
\small
\begin{tabular}{*{9}{c}}
\toprule
3 & 5 & 8 & 9 & 10 & 11 & 13 & 15 & 18 \\
\midrule
3 & 5 & 8 & 28 & 83 & 65 & 18 & 25 & 38 \\
\bottomrule
\end{tabular}
\end{center}
附帶一提，$M = 20$ 不是質數，而 $h_2(k)$ 可能是 $2, 4, 5, 6$，和 $20$ 不互質；例如 $18$ 的
步長 $5$，探測序列只會經過 $18, 3, 8, 13$ 四格。 由定理~\ref{thm:hash-double}，這種鍵的探測
序列走不遍整張表，表還有空格也可能插不進去。 這一題剛好沒有發生。

(b) $48 \bmod 20 = 8$，$48 \bmod 7 = 6$，$h_2(48) = 7$。 依序看第 $8$ 格（是 $8$，不是 $48$）、
第 $15$ 格（是 $25$）、第 $22 \bmod 20 = 2$ 格（空格），停下來，$48$ 不在表裡。 共
$3$ 次探測（照本章的數法，包括最後看到空格的那一次；如果只數「和鍵比較」的次數，是 $2$ 次）。

(c) 答案是 $n/m$。 這裡把「uniform hash function」解讀成：每個鍵的雜湊值等可能是 $m$ 格中
的任何一格（第~\ref{sec:hash-chain-analysis}~節的簡單均勻雜湊）。 設 $n$ 個鍵裡有 $n_o$ 個
奇數、$n_e$ 個偶數，$n_o + n_e = n$。 固定第 $j$ 格，落在這一格的鍵數是
\[
X_j = \sum_{\text{奇數鍵 } x} I\{h_a(x) = j\} + \sum_{\text{偶數鍵 } x} I\{h_b(x) = j\}.
\]
每個奇數鍵用 $h_a$，落在第 $j$ 格的機率是 $1/m$；每個偶數鍵用 $h_b$，機率也是 $1/m$。 由期
望值的線性，
\[
\mathrm{E}[X_j] = \frac{n_o}{m} + \frac{n_e}{m} = \frac{n}{m} = \alpha .
\]
期望值的線性不需要獨立，所以 $h_a$、$h_b$ 之間有沒有關係都不影響答案；奇偶各用哪個函數也
不影響，只要每個函數本身是均勻的。

\pastsol{px:hash-3}{中央資工所 115 學年度 第 12 題}
題目的 ``liner probing'' 是 linear probing 的筆誤。 題目沒有寫出二次探測的公式，這裡照
第~\ref{sec:hash-open}~節的 $h(k, i) = (h(k) + i^2) \bmod 11$，$i$ 從 $0$ 開始，答案是
A、B、C。 如果改用該節最後提到的正負交替跳法（$h, h + 1^2, h - 1^2, h + 2^2, \ldots$），
A、B 都不成立：$15$ 從第 $4$ 格出發依序試 $4, 5, 3$，在 A、B 裡 $15$ 插入時 $57$ 還沒來，
第 $3$ 格是空的，$15$ 會放在第 $3$ 格，和表不符，答案只剩 C。 作答時先寫明用的是哪一種。

八個鍵的起始位置是：
\[
24 \to 2,\ 57 \to 2,\ 37 \to 4,\ 15 \to 4,\ 27 \to 5,\ 16 \to 5,\ 38 \to 5,\ 17 \to 6 .
\]
判斷方法就是照每個選項的順序真的插一次，看最後的表是不是和題目一樣。 二次探測從 $h$ 出
發依序試 $h, h + 1, h + 4, h + 9, \ldots$。 下面每個鍵寫出試過的格子，最後一個是放進去的
位置。
\begin{itemize}
\item A（二次探測）：$27$：5；$37$：4；$16$：5, 6；$17$：6, 7；$24$：2；$15$：4, 5, 8；
$57$：2, 3；$38$：5, 6, 9。 和題目的表完全相同，正確。
\item B（二次探測）：$37$：4；$27$：5；$16$：5, 6；$24$：2；$15$：4, 5, 8；$57$：2, 3；
$38$：5, 6, 9；$17$：6, 7。 也完全相同，正確。（$17$ 最後才放，第 $7$ 格一直空著，所以
它還是落在第 $7$ 格。）
\item C（線性探測）：$37$：4；$27$：5；$16$：5, 6；$24$：2；$57$：2, 3；$17$：6, 7；
$15$：4, 5, 6, 7, 8；$38$：5, 6, 7, 8, 9。 完全相同，正確。
\item D（線性探測）：$27$：5；$16$：5, 6；$37$：4；$38$：5, 6, 7，放在第 $7$ 格。 但題目的
表第 $7$ 格是 $17$，第 $9$ 格才是 $38$，已經不同，錯誤。 原因是 $38$ 比 $17$ 先插入，搶走了
第 $7$ 格；繼續做下去，最後的表是 $2$：57、$3$：24、$7$：38、$9$：17，和題目不同。
\end{itemize}
也可以只用一個必要條件先刪選項：線性探測裡，鍵 $k$ 放在第 $j$ 格，那麼從 $h(k)$ 到
$j - 1$ 的每一格在 $k$ 插入時都已經有鍵。 D 裡 $38$ 放在第 $9$ 格需要第 $7$ 格先有鍵，也就是
$17$ 要比 $38$ 先插入，但 D 的順序是 $38$ 在前，所以不可能。

\pastsol{px:hash-4}{臺大資工所 115 學年度 第 15 題}
答案是 (C)。
\begin{itemize}
\item (A) 錯。 鏈結法的載入因子可以超過 $1$（第~\ref{sec:hash-chain-c}~節的例子就是
$\alpha = 8/7$），串列只是變長，插入永遠做得到。 「$\alpha = 1$ 就插不進去」是開放定址的情形。
\item (B) 錯。 雜湊表不維持任何順序，再雜湊時每個鍵用新的 $m$ 重算位置
（第~\ref{sec:hash-dynamic}~節），不會排序，之後的搜尋也不用二分搜尋。
\item (C) 對。 在簡單均勻雜湊下，串列的期望長度是 $\alpha = n/m$，搜尋的期望時間是
$\Theta(1 + \alpha)$（定理~\ref{thm:hash-chain-len}）。 $m$ 固定不變時，$\alpha$ 隨 $n$ 線
性成長，搜尋就變成 $\Theta(n)$；$\alpha$ 超過門檻就把表加倍，$\alpha$ 一直維持在常數以下，
搜尋的期望時間是 $O(1)$。 搬動的成本平攤到插入上，每次仍是 $O(1)$。
\item (D) 錯。 表變大只讓碰撞的期望次數變少，不能消除碰撞：任何固定的雜湊函數，鍵空間比
表大時一定有鍵撞在一起（第~\ref{sec:hash-func}~節的鴿籠原理），擴表之後再插入的鍵也一樣會
碰撞。
\item (E) 錯。 擴表配置的是新的串列開頭陣列，串列的節點還是一個一個配置的，不保證放在
相鄰的位置；減少快取失誤也不是擴表的目的。
\end{itemize}

\pastsol{px:hash-5}{成大資訊聯招 113 學年度 第 5 題}
位元陣列有 $15$ 格，所以 $m = 15$。

i) 答案是 (C)。 逐一算三個位置 $f_1(k) = 3k \bmod 15$、$f_2(k) = (2k + 1) \bmod 15$、
$f_3(k) = k^2 \bmod 15$：
\begin{center}
\small
\begin{tabular}{cccccl}
\toprule
$k$ & $f_1$ & $f_2$ & $f_3$ & 三格的位元 & 結論 \\
\midrule
12 & $36 \bmod 15 = 6$ & $25 \bmod 15 = 10$ & $144 \bmod 15 = 9$ & $1, 0, 1$ & 插入會把第 $10$ 格改成 $1$ \\
2 & $6$ & $5$ & $4$ & $1, 1, 1$ & 位置確實是 $4, 5, 6$ \\
15 & $45 \bmod 15 = 0$ & $31 \bmod 15 = 1$ & $225 \bmod 15 = 0$ & $1, 1, 1$ & 查詢回答 yes \\
8 & $24 \bmod 15 = 9$ & $17 \bmod 15 = 2$ & $64 \bmod 15 = 4$ & $1, 0, 1$ & 查詢回答 no \\
\bottomrule
\end{tabular}
\end{center}
\begin{itemize}
\item (A) 錯。 第 $10$ 格原本是 $0$，插入 $12$ 會把它改成 $1$。
\item (B) 錯。 $2$ 的三個位置雖然是 $4, 5, 6$，但布隆過濾器不能用清成 $0$ 的方式刪除：這三
個位元可能也屬於別的鍵，清掉之後查那些鍵會答「不在」，變成偽陰性
（第~\ref{sec:hash-bloom}~節）。
\item (C) 對。 $15$ 的位置是第 $0$、$1$、$0$ 格（$f_1$ 和 $f_3$ 剛好相同），都是 $1$，回答
yes。 至於 $15$ 是不是真的插入過，從過濾器看不出來，可能是偽陽性。
\item (D) 錯。 $8$ 的第二個位置第 $2$ 格是 $0$，回答 no。
\item (E) 錯，因為 (C) 正確。
\end{itemize}

ii) 答案是 (C)。 由定理~\ref{thm:hash-bloom-k}，最佳的雜湊函數個數是
$k^{*} = \frac{m}{n}\ln 2$，這裡 $m = 15$ 個位元、$n = 5$ 個元素，$k^{*} = 3 \ln 2 \approx
2.08$。 (A) 是 $m$，(B) 是 $n$，(D) 把 $m/n$ 錯寫成 $n$。 實際要取整數時比較 $k = 2$ 和
$k = 3$：$(1 - e^{-2/3})^2 \approx 0.237$，$(1 - e^{-1})^3 \approx 0.253$，取 $k = 2$。

\chapter{互斥集合}
\label{ch:disjoint}

% 本章的圖共用的樣式。 森林圖由模擬程式依虛擬碼逐步執行後產生，箭頭由子節點指向父節點。
\tikzset{
    dn/.style={circle, draw=black!65, minimum size=0.56cm, inner sep=0pt, font=\small},
    dnroot/.style={draw=coreStroke, very thick, fill=coreFill},   % 根（集合的代表）
    dnhl/.style={draw=memStroke, very thick, fill=memFill},       % 這一步父指標有改的節點
    dptr/.style={-{Stealth[length=1.6mm]}, draw=black!65, shorten >=0.5pt},
    plab/.style={font=\small},
    lnode/.style={draw=black!65, minimum width=0.62cm, minimum height=0.56cm, inner sep=1pt, font=\small},
    lset/.style={draw=coreStroke, thick, fill=coreFill, minimum height=0.56cm, inner sep=2pt, font=\small},
    lptr/.style={-{Stealth[length=1.6mm]}, draw=black!70},
    lrep/.style={-{Stealth[length=1.4mm]}, draw=coreStroke!80, densely dashed},
}

系上的電腦教室有幾十台電腦，網管這學期要重新拉網路線。 他手上有一張表，一行寫一
條線接了哪兩台電腦，照著表一條一條接。 接線的同時，助教隨時會跑來問：「第 7 台和第
23 台現在連得到嗎？」 線接上以後不會再拆，大家只在乎誰和誰在同一個網段。 如果每
接一條線，就把整間教室走一遍，替每台電腦重新標上網段編號，接幾百條線就要走幾百
遍。 有沒有辦法讓每次接線、每次回答都只做很少的工作？ 往下讀之前，先想想你會怎麼
記錄「誰和誰同一組」。

這一章的資料結構叫\textbf{互斥集合}（disjoint sets），也常叫做 union--find。 它只
支援三件事：建立只有一個元素的集合、問一個元素在哪個集合、把兩個集合合併。 功能很
少，卻是許多圖論演算法的零件，第~24~章的 Kruskal 最小生成樹演算法就靠它判斷一條邊
會不會形成迴路。 這一章先用串列表示集合，看「把短的接到長的後面」這個簡單的規則能省
下多少工作；再改用第~7~章說過的「只存父節點」的森林，配上依大小或依秩合併、路徑壓
縮兩個技巧，最後說明這兩個技巧合起來，每個操作的平攤成本幾乎是常數。

這一章用到《C 程式設計基礎》第~14~章的結構、第~15~章的遞迴和第~19~章的
\texttt{malloc}、\texttt{free}；平攤的意思見第~3~章，森林和父節點陣列見第~7~章。

本章的\textbf{高度}照第~7~章的定義數節點：只有一個節點的樹高度是 $1$，空樹是 $0$；
\textbf{深度}是從根往下走的邊數，根的深度是 $0$。 第~\ref{sec:ds-rank}~節會出現另
一個量叫\textbf{秩}，它是另外定義的，和高度差一，不要混在一起。

\section{操作與串列表示}
\label{sec:ds-list}

\subsection{三個操作}

互斥集合維護一群\textbf{兩兩沒有共同元素}的集合 $S_1, S_2, \ldots, S_k$。 每個集合選
一個成員當\textbf{代表}（representative），用代表來指稱整個集合。 代表是誰可以由實作
決定，但只要集合沒有被改過，同一個集合問兩次，必須得到同一個代表。 支援的操作有三
個：

\begin{itemize}
\item \proc{Make-Set}$(x)$：建立只含 $x$ 的新集合，$x$ 就是它的代表。 $x$ 原本不能在
任何集合裡。
\item \proc{Find-Set}$(x)$：回傳 $x$ 所在集合的代表。
\item \proc{Union}$(x, y)$：把 $x$ 所在的集合和 $y$ 所在的集合合成一個，原來的兩個集
合不再存在。 新集合的代表從兩個舊代表裡選一個。 $x$、$y$ 已經在同一個集合時什麼都
不做。
\end{itemize}

問「$x$ 和 $y$ 在不在同一組」，就是比較 $\proc{Find-Set}(x) = \proc{Find-Set}(y)$。
互斥集合不支援把一個集合拆開，也不支援把元素移出集合；第~\ref{sec:ds-app}~節會看到
遇到「拆線」時怎麼繞過這個限制。

分析時用兩個參數：$n$ 是 \proc{Make-Set} 的次數，也就是元素的總數；$m$ 是所有操作
的總數，三種都算。 每個元素都要先做一次 \proc{Make-Set}，所以 $m \ge n$。 每做一次
真正的合併，集合就少一個，所以真正的合併最多 $n - 1$ 次。

章首的教室用這三個操作來寫：每台電腦做一次 \proc{Make-Set}，每接一條線 $(u, v)$ 做一次
\proc{Union}$(u, v)$；助教問 $u$、$v$ 連不連得到，就回答 \proc{Find-Set}$(u) =
\proc{Find-Set}(v)$ 是否成立。 如果所有的線一次給齊，也可以用第~21~章的廣度或深度優
先搜尋，一次把每台電腦標上組別；互斥集合的好處在於線是一條一條來的，每來一條線只做一
次合併，中間隨時可以問。

\subsection{串列表示}

最直接的做法是把每個集合存成一條鏈結串列（第~5~章）。 每個集合有一個\textbf{集合物
件}，記錄串列的頭 $\mathit{head}$、尾 $\mathit{tail}$ 和元素個數 $\mathit{size}$；串
列裡每個元素 $x$ 除了往下一個的指標 $x.\mathit{next}$，還多一個指回集合物件的指標
$x.\mathit{set}$。 代表就是串列的第一個元素（圖~\ref{fig:ds-list}）。

\begin{center}
\begin{tikzpicture}[x=1cm, y=1cm]
\node[lset] (sa) at (1.4, 1.25) {$\mathit{head}\ \ \mathit{tail}\ \ \mathit{size}=3$};
\node[lnode] (a4) at (0, 0) {4};
\node[lnode] (a9) at (1.4, 0) {9};
\node[lnode] (a2) at (2.8, 0) {2};
\node[font=\small] (an) at (3.9, 0) {\NIL};
\draw[lptr] (a4) -- (a9);
\draw[lptr] (a9) -- (a2);
\draw[lptr] (a2) -- (an);
\foreach \v in {a4, a9, a2} { \draw[lrep] (\v.north) -- (\v.north |- sa.south); }
\node[plab] at (1.4, -0.85) {集合 $\{4, 9, 2\}$，代表是 $4$};
\begin{scope}[xshift=6.2cm]
\node[lset] (sb) at (0.7, 1.25) {$\mathit{head}\ \ \mathit{tail}\ \ \mathit{size}=2$};
\node[lnode] (b6) at (0, 0) {6};
\node[lnode] (b1) at (1.4, 0) {1};
\node[font=\small] (bn) at (2.5, 0) {\NIL};
\draw[lptr] (b6) -- (b1);
\draw[lptr] (b1) -- (bn);
\foreach \v in {b6, b1} { \draw[lrep] (\v.north) -- (\v.north |- sb.south); }
\node[plab] at (0.7, -0.85) {集合 $\{6, 1\}$，代表是 $6$};
\end{scope}
\end{tikzpicture}
\diagcaption{用串列表示兩個集合。 實線是 $\mathit{next}$，虛線是每個元素指回集合物件
的 $\mathit{set}$。 \proc{Find-Set}$(2)$ 沿虛線找到集合物件，再讀 $\mathit{head}$，
得到 $4$。}
\label{fig:ds-list}
\end{center}

\begin{pseudo}[label=alg:ds-list]{串列表示，加權合併}
\Procedure{Make-Set}{$x$}
    \State 建立集合物件 $S$；$S.\mathit{head} \gets x$，$S.\mathit{tail} \gets x$，$S.\mathit{size} \gets 1$
    \State $x.\mathit{set} \gets S$；$x.\mathit{next} \gets \NIL$
\EndProcedure
\Function{Find-Set}{$x$}
    \State \Return $x.\mathit{set}.\mathit{head}$
\EndFunction
\Procedure{Union}{$x, y$}
    \State $S \gets x.\mathit{set}$
    \State $T \gets y.\mathit{set}$
    \If{$S = T$}
        \State \Return
    \EndIf
    \If{$S.\mathit{size} < T.\mathit{size}$} \Comment{加權：讓 $T$ 是比較短的那條}
        \State 交換 $S$ 和 $T$
    \EndIf
    \State $S.\mathit{tail}.\mathit{next} \gets T.\mathit{head}$
    \State $S.\mathit{tail} \gets T.\mathit{tail}$
    \State $z \gets T.\mathit{head}$
    \While{$z \ne \NIL$} \Comment{$T$ 的每個元素改指 $S$}
        \State $z.\mathit{set} \gets S$
        \State $z \gets z.\mathit{next}$
    \EndWhile
    \State $S.\mathit{size} \gets S.\mathit{size} + T.\mathit{size}$
    \State 釋放集合物件 $T$
\EndProcedure
\end{pseudo}

\proc{Make-Set} 和 \proc{Find-Set} 都是 $O(1)$。 \proc{Union} 把 $T$ 整條接到 $S$ 的
尾巴，接上去只要改兩個指標（第 13--14 行），貴的是第 16--18 行：$T$ 的每個元素都要把
$\mathit{set}$ 改成 $S$，花 $\Theta(T.\mathit{size})$ 的時間。

先看拿掉第 11--12 行、永遠把 $y$ 的串列接到 $x$ 後面的版本。 教室裡新到一台電腦就
接到已經連好的大群組上，寫成操作是：\proc{Make-Set} $x_1, \ldots, x_n$，再依序做
\[
\proc{Union}(x_2, x_1),\ \proc{Union}(x_3, x_1),\ \ldots,\ \proc{Union}(x_n, x_1).
\]
第 $i$ 次合併時，$x_1$ 所在的集合已經有 $i$ 個元素，整條被接到只有一個元素的新集合
後面，要改 $i$ 個 $\mathit{set}$ 指標。 總共改了 $1 + 2 + \cdots + (n - 1) =
n(n-1)/2$ 次，而操作數只有 $m = 2n - 1$。 平均每個操作要 $\Theta(n)$ 的時間。

毛病出在每次都把長的接到短的後面。 第 11--12 行的\textbf{加權合併}（weighted-union
heuristic）規則是：比較兩條串列的長度，把短的接到長的後面，只改短的那一條的指標。
一樣長時接哪一條都可以，本章約定接 $y$ 的那一條。 同樣那串操作，每次只要改 $1$ 個
指標。

\begin{theorem}
\label{thm:ds-weighted-list}
串列表示使用加權合併時，$m$ 個操作（其中 $n$ 個是 \proc{Make-Set}）的總時間是
$O(m + n \lg n)$。
\end{theorem}

\begin{proof}
每個操作除了改 $\mathit{set}$ 指標以外，都只花常數時間，合計 $O(m)$。 剩下的是改指
標的總次數。 固定一個元素 $x$，看它的 $\mathit{set}$ 指標被改的時候：$x$ 一定在比較
短（或一樣長）的那一條裡，合併後的集合至少是 $x$ 原本所在集合的兩倍大。 $x$ 的集合
一開始只有 $1$ 個元素，被改過 $k$ 次之後至少有 $2^k$ 個元素；集合不會超過 $n$ 個元
素，所以 $2^k \le n$，$k \le \lg n$。 $n$ 個元素加起來最多改 $n \lg n$ 次，總時間
$O(m + n \lg n)$。
\end{proof}

這個上界在最壞情況是緊的。 設 $n = 2^k$，先兩兩合併成 $n/2$ 個大小 $2$ 的集合，再兩
兩合併成 $n/4$ 個大小 $4$ 的集合……一共 $k$ 輪。 每一輪的每次合併都是兩條一樣長的
串列，有一半的元素要改指標，每一輪改 $n/2$ 次，總共 $(n/2)\lg n$ 次。

\section{森林表示}
\label{sec:ds-forest}

串列表示讓 \proc{Find-Set} 很快，代價是合併要改一整條的指標。 反過來想：合併時只改
一個指標，\proc{Find-Set} 多走幾步。 做法是把每個集合存成一棵有根樹，每個節點只記
它的父節點 $p[x]$，根就是代表，根的父節點設成自己。 第~7~章說過，一個陣列
$p[1..n]$ 就能存下整個森林。

\begin{pseudo}[label=alg:ds-naive]{森林表示，不用任何技巧}
\Procedure{Make-Set}{$x$}
    \State $p[x] \gets x$
\EndProcedure
\Function{Find-Set}{$x$}
    \While{$p[x] \ne x$}
        \State $x \gets p[x]$
    \EndWhile
    \State \Return $x$
\EndFunction
\Procedure{Union}{$x, y$}
    \State $r_x \gets \Call{Find-Set}{x}$
    \State $r_y \gets \Call{Find-Set}{y}$
    \If{$r_x \ne r_y$}
        \State $p[r_y] \gets r_x$ \Comment{$y$ 的根接到 $x$ 的根底下}
    \EndIf
\EndProcedure
\end{pseudo}

第 11 行接的是兩個\textbf{根}。 寫成 $p[y] \gets x$ 是常見的錯：$y$ 原本的根和它其他
的子孫會留在原來的樹裡，兩個集合沒有真的合起來。

\proc{Find-Set} 從 $x$ 走到根，時間和 $x$ 的深度成正比；\proc{Union} 除了兩次
\proc{Find-Set} 只有常數時間。 問題是樹可能長得很高。 把第~\ref{sec:ds-list}~節那串
操作搬過來：\proc{Union}$(x_{i+1}, x_1)$ 每次都把 $x_1$ 所在的樹的根接到新節點底下，
樹長成一條 $x_1 \to x_2 \to \cdots \to x_n$ 的鏈，第 $i$ 次合併光是
\proc{Find-Set}$(x_1)$ 就要經過 $i$ 個節點，總時間又是 $\Theta(n^2)$。 下面兩個技巧
分別從「合併時怎麼接」和「找的時候順便整理」下手。

\subsection{依大小合併}
\label{sec:ds-size}

合併時讓節點比較少的那棵樹的根，接到節點比較多的那棵樹的根底下，叫做\textbf{依大小
合併}（union by size）。 每個根多記一個 $\mathit{size}$，是這棵樹的節點數；不是根的
節點的 $\mathit{size}$ 不再使用。

\begin{pseudo}[label=alg:ds-size]{依大小合併（\proc{Make-Set}$(x)$ 另外設 $\mathit{size}[x] \gets 1$）}
\Procedure{Union}{$x, y$}
    \State $r_x \gets \Call{Find-Set}{x}$
    \State $r_y \gets \Call{Find-Set}{y}$
    \If{$r_x = r_y$}
        \State \Return
    \EndIf
    \If{$\mathit{size}[r_x] < \mathit{size}[r_y]$}
        \State 交換 $r_x$ 和 $r_y$ \Comment{讓 $r_x$ 是節點比較多的根}
    \EndIf
    \State $p[r_y] \gets r_x$
    \State $\mathit{size}[r_x] \gets \mathit{size}[r_x] + \mathit{size}[r_y]$
\EndProcedure
\end{pseudo}

第 4--5 行不能省：$x$、$y$ 已經同一組時如果照樣執行第 9 行，$\mathit{size}$ 會被重
複累加。 兩棵樹一樣大時第 6 行不成立，$y$ 的根接到 $x$ 的根底下，和
虛擬碼~\ref{alg:ds-naive} 的方向相同。 第~\ref{sec:ds-cost}~節會證明，依大小合併的
樹，高度不超過 $\lfloor \lg n \rfloor + 1$。

\subsection{依秩合併}
\label{sec:ds-rank}

另一種接法是比較兩棵樹的\textbf{秩}（rank）。 秩是每個節點記的一個非負整數，規則只
有三條：

\begin{itemize}
\item \proc{Make-Set}$(x)$ 時 $\mathit{rank}[x] = 0$。 只有一個節點的樹，根的秩是
$0$（這棵樹的高度是 $1$）。
\item 合併兩個秩不同的根，秩小的接到秩大的底下，兩個秩都不變。
\item 合併兩個秩相同的根，$y$ 的根接到 $x$ 的根底下，留下來的根秩加一。
\end{itemize}

\begin{pseudo}[label=alg:ds-rank]{依秩合併（\proc{Make-Set}$(x)$ 另外設 $\mathit{rank}[x] \gets 0$）}
\Procedure{Union}{$x, y$}
    \State $r_x \gets \Call{Find-Set}{x}$
    \State $r_y \gets \Call{Find-Set}{y}$
    \If{$r_x = r_y$}
        \State \Return
    \EndIf
    \If{$\mathit{rank}[r_x] < \mathit{rank}[r_y]$}
        \State 交換 $r_x$ 和 $r_y$
    \EndIf
    \State $p[r_y] \gets r_x$
    \If{$\mathit{rank}[r_x] = \mathit{rank}[r_y]$}
        \State $\mathit{rank}[r_x] \gets \mathit{rank}[r_x] + 1$
    \EndIf
\EndProcedure
\end{pseudo}

秩和高度的關係要分兩種情況說。 不做路徑壓縮時，節點 $x$ 的秩\textbf{恰好}等於以 $x$
為根的子樹裡，從 $x$ 往下最長的路徑有幾條邊（下一節的路徑壓縮會讓樹變矮，
但不改任何秩）。 所以一般來說，秩是這個邊數的上界。 換成本書數節點的高度：
\[
\text{以 $x$ 為根的子樹的高度} \le \mathit{rank}[x] + 1,
\]
不做路徑壓縮時等號成立。 秩只在根上會變，一個節點變成別人的子節點以後，它的秩就固
定了。

\begin{examnote}
CLRS 和許多考卷用「邊數」定義高度，單一節點的高度是 $0$，所以會看到「秩是高度的上
界」「不壓縮時秩等於高度」這種寫法；換成本書數節點的高度，要寫成高度 $\le$ 秩 $+ 1$。
也有書把依秩合併叫做\textbf{依高度合併}（union by height），內容相同。 作答時先寫出
自己用的定義，題目問「秩」就照秩的三條規則算，不要拿實際的樹高代替。
\end{examnote}

依大小和依秩多數時候做出一樣的選擇，但不一定。 例如一棵是 $5$ 個節點、高度 $2$ 的
星形（根的秩 $1$），另一棵是 $4$ 個節點、高度 $3$ 的樹（根的秩 $2$）：依大小合併會
讓星形的根留下來，依秩合併會讓另一棵的根留下來。 觀念題第~\ref{q:ds-compare}~題就
是這種情形。

\subsection{路徑壓縮}
\label{sec:ds-compress}

\proc{Find-Set}$(x)$ 已經走過從 $x$ 到根的整條路徑，順便把路上每個節點的父指標都直接
改成根，下次再找這些節點，一步就到。 這叫\textbf{路徑壓縮}（path compression）。 它
只改父指標，不改任何秩和 $\mathit{size}$。

\begin{pseudo}[label=alg:ds-find-rec]{路徑壓縮（遞迴）}
\Function{Find-Set}{$x$}
    \If{$p[x] \ne x$}
        \State $p[x] \gets \Call{Find-Set}{p[x]}$
    \EndIf
    \State \Return $p[x]$
\EndFunction
\end{pseudo}

遞迴版一路往上找到根，回來的時候（第 3 行）把根的編號一層一層填回去。 它其實走了兩
趟：去的時候找根，回來的時候改指標。 寫成迴圈就把兩趟分開：

\begin{pseudo}[label=alg:ds-find-iter]{路徑壓縮（迴圈）}
\Function{Find-Set-Iter}{$x$}
    \State $r \gets x$
    \While{$p[r] \ne r$} \Comment{第一趟：找根}
        \State $r \gets p[r]$
    \EndWhile
    \While{$p[x] \ne r$} \Comment{第二趟：路上的節點都改接到根}
        \State $\mathit{next} \gets p[x]$
        \State $p[x] \gets r$
        \State $x \gets \mathit{next}$
    \EndWhile
    \State \Return $r$
\EndFunction
\end{pseudo}

第 6 行要先記下 $p[x]$，第 7 行改掉之後就找不到原本的下一個節點了。 兩個版本對同一
個森林做同一次 \proc{Find-Set}，得到的父節點陣列完全一樣，第~\ref{sec:ds-c}~節的 C
程式兩種都用到。 遞迴版寫起來短，但遞迴的深度等於路徑的長度；如果不搭配依大小或依
秩合併，樹可以是一條長度 $n$ 的鏈，$n$ 很大時會用光呼叫堆疊。

\subsection{路徑減半與路徑分割}
\label{sec:ds-onepass}

路徑壓縮要走兩趟。 另外有兩種只走一趟的變形，不必先找到根：
\begin{itemize}
\item \textbf{路徑減半}（path halving）：路上每隔一個節點，把它的父指標改成祖父，再
跳到（新的）父節點。
\item \textbf{路徑分割}（path splitting）：路上\textbf{每個}節點都把父指標改成祖父，再
跳到原本的父節點。
\end{itemize}

\begin{pseudo}[label=alg:ds-find-onepass]{只走一趟的兩種變形}
\Function{Find-Halving}{$x$}
    \While{$p[x] \ne x$}
        \State $p[x] \gets p[p[x]]$
        \State $x \gets p[x]$
    \EndWhile
    \State \Return $x$
\EndFunction
\Function{Find-Splitting}{$x$}
    \While{$p[x] \ne x$}
        \State $\mathit{next} \gets p[x]$
        \State $p[x] \gets p[\mathit{next}]$
        \State $x \gets \mathit{next}$
    \EndWhile
    \State \Return $x$
\EndFunction
\end{pseudo}

圖~\ref{fig:ds-onepass} 是同一條路徑 $1 \to 2 \to \cdots \to 6$ 做
\proc{Find-Set}$(1)$ 的三種結果，橘色是父指標有改的節點。 路徑壓縮讓整條路徑都掛到根
底下；減半只改了 $1$ 和 $3$，路徑長度大約減半；分割改了 $1$ 到 $4$，原來的路徑拆成
$1 \to 3 \to 5 \to 6$ 和 $2 \to 4 \to 6$ 兩條。 一次 \proc{Find-Set} 之後，減半和分
割的樹都比壓縮的高，但 Tarjan 和 van Leeuwen 證明，搭配依秩（或依大小）合併時，三種
做法的總時間是同一級（第~\ref{sec:ds-cost}~節）。

\begin{center}
\begin{tikzpicture}[x=0.8cm,y=0.85cm,baseline=(current bounding box.north)]
\node[dn] (pa1) at (0,-5) {1};
\node[dn] (pa2) at (0,-4) {2};
\node[dn] (pa3) at (0,-3) {3};
\node[dn] (pa4) at (0,-2) {4};
\node[dn] (pa5) at (0,-1) {5};
\node[dn,dnroot] (pa6) at (0,0) {6};
\draw[dptr] (pa1) -- (pa2);
\draw[dptr] (pa2) -- (pa3);
\draw[dptr] (pa3) -- (pa4);
\draw[dptr] (pa4) -- (pa5);
\draw[dptr] (pa5) -- (pa6);
\node[plab, anchor=north] at (0,-5.45) {(a) 原來的路徑};
\end{tikzpicture}\hspace{0.8cm}%
\begin{tikzpicture}[x=0.8cm,y=0.85cm,baseline=(current bounding box.north)]
\node[dn,dnhl] (pb1) at (0,-1) {1};
\node[dn,dnhl] (pb2) at (1,-1) {2};
\node[dn,dnhl] (pb3) at (2,-1) {3};
\node[dn,dnhl] (pb4) at (3,-1) {4};
\node[dn] (pb5) at (4,-1) {5};
\node[dn,dnroot] (pb6) at (2,0) {6};
\draw[dptr] (pb1) -- (pb6);
\draw[dptr] (pb2) -- (pb6);
\draw[dptr] (pb3) -- (pb6);
\draw[dptr] (pb4) -- (pb6);
\draw[dptr] (pb5) -- (pb6);
\node[plab, anchor=north] at (2,-5.45) {(b) 路徑壓縮};
\end{tikzpicture}\hspace{0.8cm}%
\begin{tikzpicture}[x=0.8cm,y=0.85cm,baseline=(current bounding box.north)]
\node[dn,dnhl] (pc1) at (0,-3) {1};
\node[dn] (pc2) at (1,-3) {2};
\node[dn,dnhl] (pc3) at (0.5,-2) {3};
\node[dn] (pc4) at (2,-2) {4};
\node[dn] (pc5) at (1.25,-1) {5};
\node[dn,dnroot] (pc6) at (1.25,0) {6};
\draw[dptr] (pc1) -- (pc3);
\draw[dptr] (pc2) -- (pc3);
\draw[dptr] (pc3) -- (pc5);
\draw[dptr] (pc4) -- (pc5);
\draw[dptr] (pc5) -- (pc6);
\node[plab, anchor=north] at (1.25,-5.45) {(c) 路徑減半};
\end{tikzpicture}\hspace{0.8cm}%
\begin{tikzpicture}[x=0.8cm,y=0.85cm,baseline=(current bounding box.north)]
\node[dn,dnhl] (pd1) at (1,-3) {1};
\node[dn,dnhl] (pd2) at (0,-2) {2};
\node[dn,dnhl] (pd3) at (1,-2) {3};
\node[dn,dnhl] (pd4) at (0,-1) {4};
\node[dn] (pd5) at (1,-1) {5};
\node[dn,dnroot] (pd6) at (0.5,0) {6};
\draw[dptr] (pd1) -- (pd3);
\draw[dptr] (pd2) -- (pd4);
\draw[dptr] (pd3) -- (pd5);
\draw[dptr] (pd4) -- (pd6);
\draw[dptr] (pd5) -- (pd6);
\node[plab, anchor=north] at (0.5,-5.45) {(d) 路徑分割};
\end{tikzpicture}
\diagcaption{對同一條路徑做 \proc{Find-Set}$(1)$。 藍框是根，橘色是父指標有改的節點，
箭頭由子節點指向父節點。 父節點陣列 $p[1..6]$ 依序是 (b) $6,6,6,6,6,6$；(c)
$3,3,5,5,6,6$；(d) $3,4,5,6,6,6$。}
\label{fig:ds-onepass}
\end{center}

\subsection{手算追蹤：實驗室網路分組}
\label{sec:ds-trace}

實驗室有 $10$ 台電腦，編號 $1$ 到 $10$，一開始各自一組。 網管依序做了下面 $11$ 個操
作，用\textbf{依大小合併}（虛擬碼~\ref{alg:ds-size}）加\textbf{路徑壓縮}（虛擬
碼~\ref{alg:ds-find-rec}），一樣大時 $y$ 的根接到 $x$ 的根底下。 \proc{Union} 裡的兩
次 \proc{Find-Set} 也會壓縮路徑，這一點常被忽略。 下表每一列是做完那個操作之後的
$p[1..10]$，粗體是這一步改過的格子；圖~\ref{fig:ds-trace} 畫出其中四個時間點的森林。

\begin{center}
\small
\setlength{\tabcolsep}{4.2pt}
\begin{tabular}{rl*{10}{c}}
\toprule
步 & 操作 & $p[1]$ & $p[2]$ & $p[3]$ & $p[4]$ & $p[5]$ & $p[6]$ & $p[7]$ & $p[8]$ & $p[9]$ & $p[10]$ \\
\midrule
0 & 開始 & 1 & 2 & 3 & 4 & 5 & 6 & 7 & 8 & 9 & 10 \\
1 & \proc{Union}$(1, 2)$ & 1 & \textbf{1} & 3 & 4 & 5 & 6 & 7 & 8 & 9 & 10 \\
2 & \proc{Union}$(3, 4)$ & 1 & 1 & 3 & \textbf{3} & 5 & 6 & 7 & 8 & 9 & 10 \\
3 & \proc{Union}$(5, 6)$ & 1 & 1 & 3 & 3 & 5 & \textbf{5} & 7 & 8 & 9 & 10 \\
4 & \proc{Union}$(7, 8)$ & 1 & 1 & 3 & 3 & 5 & 5 & 7 & \textbf{7} & 9 & 10 \\
5 & \proc{Union}$(2, 4)$ & 1 & 1 & \textbf{1} & 3 & 5 & 5 & 7 & 7 & 9 & 10 \\
6 & \proc{Union}$(6, 8)$ & 1 & 1 & 1 & 3 & 5 & 5 & \textbf{5} & 7 & 9 & 10 \\
7 & \proc{Union}$(9, 10)$ & 1 & 1 & 1 & 3 & 5 & 5 & 5 & 7 & 9 & \textbf{9} \\
8 & \proc{Union}$(6, 2)$ & \textbf{5} & 1 & 1 & 3 & 5 & 5 & 5 & 7 & 9 & 9 \\
9 & \proc{Find-Set}$(4)$ & 5 & 1 & \textbf{5} & \textbf{5} & 5 & 5 & 5 & 7 & 9 & 9 \\
10 & \proc{Union}$(10, 8)$ & 5 & 1 & 5 & 5 & 5 & 5 & 5 & \textbf{5} & \textbf{5} & 9 \\
11 & \proc{Find-Set}$(10)$ & 5 & 1 & 5 & 5 & 5 & 5 & 5 & 5 & 5 & \textbf{5} \\
\bottomrule
\end{tabular}
\end{center}

\begin{center}
\begin{tikzpicture}[x=0.8cm,y=0.85cm,baseline=(current bounding box.north)]
\node[dn,dnroot] (tr6x1) at (0.5,0) {1};
\node[dn] (tr6x2) at (0,-1) {2};
\node[dn] (tr6x3) at (1,-1) {3};
\node[dn] (tr6x4) at (1,-2) {4};
\node[dn,dnroot] (tr6x5) at (2.85,0) {5};
\node[dn] (tr6x6) at (2.35,-1) {6};
\node[dn,dnhl] (tr6x7) at (3.35,-1) {7};
\node[dn] (tr6x8) at (3.35,-2) {8};
\node[dn,dnroot] (tr6x9) at (4.7,0) {9};
\node[dn,dnroot] (tr6x10) at (6.05,0) {10};
\draw[dptr] (tr6x2) -- (tr6x1);
\draw[dptr] (tr6x3) -- (tr6x1);
\draw[dptr] (tr6x4) -- (tr6x3);
\draw[dptr] (tr6x6) -- (tr6x5);
\draw[dptr] (tr6x7) -- (tr6x5);
\draw[dptr] (tr6x8) -- (tr6x7);
\node[plab, anchor=north] at (3.02,-3.45) {(a) 第 6 步之後};
\end{tikzpicture}\hspace{1.0cm}%
\begin{tikzpicture}[x=0.8cm,y=0.85cm,baseline=(current bounding box.north)]
\node[dn,dnhl] (tr8x1) at (0.5,-1) {1};
\node[dn] (tr8x2) at (0,-2) {2};
\node[dn] (tr8x3) at (1,-2) {3};
\node[dn] (tr8x4) at (1,-3) {4};
\node[dn,dnroot] (tr8x5) at (1.75,0) {5};
\node[dn] (tr8x6) at (2,-1) {6};
\node[dn] (tr8x7) at (3,-1) {7};
\node[dn] (tr8x8) at (3,-2) {8};
\node[dn,dnroot] (tr8x9) at (4.35,0) {9};
\node[dn] (tr8x10) at (4.35,-1) {10};
\draw[dptr] (tr8x1) -- (tr8x5);
\draw[dptr] (tr8x2) -- (tr8x1);
\draw[dptr] (tr8x3) -- (tr8x1);
\draw[dptr] (tr8x4) -- (tr8x3);
\draw[dptr] (tr8x6) -- (tr8x5);
\draw[dptr] (tr8x7) -- (tr8x5);
\draw[dptr] (tr8x8) -- (tr8x7);
\draw[dptr] (tr8x10) -- (tr8x9);
\node[plab, anchor=north] at (2.17,-3.45) {(b) 第 8 步之後};
\end{tikzpicture}

\vspace{0.5em}
\begin{tikzpicture}[x=0.8cm,y=0.85cm,baseline=(current bounding box.north)]
\node[dn] (tr9x1) at (0,-1) {1};
\node[dn] (tr9x2) at (0,-2) {2};
\node[dn,dnhl] (tr9x3) at (1,-1) {3};
\node[dn,dnhl] (tr9x4) at (2,-1) {4};
\node[dn,dnroot] (tr9x5) at (2,0) {5};
\node[dn] (tr9x6) at (3,-1) {6};
\node[dn] (tr9x7) at (4,-1) {7};
\node[dn] (tr9x8) at (4,-2) {8};
\node[dn,dnroot] (tr9x9) at (5.35,0) {9};
\node[dn] (tr9x10) at (5.35,-1) {10};
\draw[dptr] (tr9x1) -- (tr9x5);
\draw[dptr] (tr9x2) -- (tr9x1);
\draw[dptr] (tr9x3) -- (tr9x5);
\draw[dptr] (tr9x4) -- (tr9x5);
\draw[dptr] (tr9x6) -- (tr9x5);
\draw[dptr] (tr9x7) -- (tr9x5);
\draw[dptr] (tr9x8) -- (tr9x7);
\draw[dptr] (tr9x10) -- (tr9x9);
\node[plab, anchor=north] at (2.67,-2.45) {(c) 第 9 步之後};
\end{tikzpicture}\hspace{0.8cm}%
\begin{tikzpicture}[x=0.8cm,y=0.85cm,baseline=(current bounding box.north)]
\node[dn] (tr11x1) at (0,-1) {1};
\node[dn] (tr11x2) at (0,-2) {2};
\node[dn] (tr11x3) at (1,-1) {3};
\node[dn] (tr11x4) at (2,-1) {4};
\node[dn,dnroot] (tr11x5) at (3.5,0) {5};
\node[dn] (tr11x6) at (3,-1) {6};
\node[dn] (tr11x7) at (4,-1) {7};
\node[dn] (tr11x8) at (5,-1) {8};
\node[dn] (tr11x9) at (6,-1) {9};
\node[dn,dnhl] (tr11x10) at (7,-1) {10};
\draw[dptr] (tr11x1) -- (tr11x5);
\draw[dptr] (tr11x2) -- (tr11x1);
\draw[dptr] (tr11x3) -- (tr11x5);
\draw[dptr] (tr11x4) -- (tr11x5);
\draw[dptr] (tr11x6) -- (tr11x5);
\draw[dptr] (tr11x7) -- (tr11x5);
\draw[dptr] (tr11x8) -- (tr11x5);
\draw[dptr] (tr11x9) -- (tr11x5);
\draw[dptr] (tr11x10) -- (tr11x5);
\node[plab, anchor=north] at (3.5,-2.45) {(d) 第 11 步之後};
\end{tikzpicture}
\diagcaption{上表幾個時間點的森林。 藍框是根，橘色是這一步父指標有改的節點，箭頭由子
節點指向父節點。}
\label{fig:ds-trace}
\end{center}

第 1--4 步都是兩個單點合併，一樣大，$y$ 的根接到 $x$ 底下。 其他幾步要注意的地方：

\begin{itemize}
\item 第 5 步：\proc{Find-Set}$(2) = 1$，\proc{Find-Set}$(4) = 3$，兩棵都是 $2$ 個節點，
$3$ 接到 $1$ 底下，$\mathit{size}[1] = 4$。 注意改的是 $p[3]$，不是 $p[4]$。
\item 第 8 步：兩個根 $5$ 和 $1$ 都是 $4$ 個節點，$1$ 接到 $5$ 底下。 這時 $4$ 的深度
是 $3$，路徑 $4 \to 3 \to 1 \to 5$ 是整個森林最長的，圖 (b) 的樹高度是 $4$。
\item 第 9 步：\proc{Find-Set}$(4)$ 走過 $4, 3, 1, 5$，把 $4$ 和 $3$ 改接到 $5$；$1$ 本
來就是 $5$ 的子節點，不用改。 $2$ 不在這條路徑上，仍然掛在 $1$ 底下。
\item 第 10 步：\proc{Find-Set}$(10) = 9$，\proc{Find-Set}$(8)$ 走過 $8, 7, 5$，順便把
$8$ 改接到 $5$。 $\mathit{size}[9] = 2 < \mathit{size}[5] = 8$，兩個根交換，$9$ 接到
$5$ 底下。 這一步改了兩格：$p[8]$ 來自壓縮，$p[9]$ 來自合併。
\end{itemize}

\begin{examnote}
追蹤題通常會規定：(1) 一樣大（或秩相同）時誰當根；(2) \proc{Union}$(x, y)$ 裡的
\proc{Find-Set} 要不要壓縮；(3) 陣列的格式。 有些課本讓根的格子存\textbf{負的節點數}
（例如 $p[1] = -4$ 表示 $1$ 是根、這棵樹有 $4$ 個節點），不另外開 $\mathit{size}$
陣列；照這種寫法，上表第 6 步之後是 $-4, 1, 1, 3, -4, 5, 5, 7, -1, -1$。 這一系的課本
把依大小合併叫\textbf{加權規則}（weighting rule），把路徑壓縮叫\textbf{塌陷規則}
（collapsing rule）。 題目沒有規定時，在答案開頭寫清楚自己的約定，再開始畫。
\end{examnote}

\section{成本}
\label{sec:ds-cost}

\subsection{依大小或依秩合併：高度是 $O(\lg n)$}

\begin{theorem}
\label{thm:ds-size-height}
只用依大小合併（有沒有路徑壓縮都可以），任何時候，一棵有 $s$ 個節點的樹，高度不超
過 $\lfloor \lg s \rfloor + 1$。 所以 \proc{Find-Set} 最壞是 $O(\lg n)$。
\end{theorem}

\begin{proof}
先看沒有路徑壓縮的情形。 固定一個節點 $x$，它的深度只有在一種情況會增加：$x$ 所在
的樹的根 $r_y$ 被接到另一個根 $r_x$ 底下，這時深度加一。 依大小合併保證
$\mathit{size}[r_x] \ge \mathit{size}[r_y]$，所以合併後 $x$ 所在的樹的節點數
$\mathit{size}[r_x] + \mathit{size}[r_y] \ge 2\,\mathit{size}[r_y]$，至少是原來的兩
倍。 $x$ 一開始在只有 $1$ 個節點的樹裡；深度變成 $d$ 時，它所在的樹至少有 $2^d$ 個
節點。 現在這棵樹有 $s$ 個節點，所以 $2^d \le s$，$d \le \lfloor \lg s \rfloor$。
樹的高度是最大深度加一，不超過 $\lfloor \lg s \rfloor + 1$。

路徑壓縮只會把節點改接到更上面的祖先，任何節點的深度都不會因此變大，每棵樹的節點數
也不變，所以上界仍然成立。
\end{proof}

依秩合併的上界要先整理秩的性質。

\begin{lemma}
\label{lem:ds-rank}
使用依秩合併（有沒有路徑壓縮都可以），任何時候：
\begin{enumerate}
\item 若 $x$ 不是根，$\mathit{rank}[x] < \mathit{rank}[p[x]]$。
\item 秩為 $r$ 的根，它的樹至少有 $2^r$ 個節點。 所以每個秩都不超過
$\lfloor \lg n \rfloor$。
\item 以 $x$ 為根的子樹，從 $x$ 往下最長的路徑不超過 $\mathit{rank}[x]$ 條邊；樹的高
度不超過根的秩加一，也就不超過 $\lfloor \lg n \rfloor + 1$。
\end{enumerate}
\end{lemma}

\begin{whybox}[引理~\ref{lem:ds-rank} 的證明]
(1) $x$ 變成子節點的那一刻（虛擬碼~\ref{alg:ds-rank} 第 8 行），$\mathit{rank}[x] \le
\mathit{rank}[p[x]]$，相等時第 10 行把父節點的秩加一，所以嚴格小於。 之後 $x$ 不再是
根，秩不再變；它的父節點如果還是根，秩只會變大；路徑壓縮把 $p[x]$ 換成更上面的祖先，
由同樣的道理，祖先的秩更大。 不等式一直成立。

(2) 對操作的次數做歸納。 一開始每個根的秩是 $0$，樹有 $1 = 2^0$ 個節點。 合併時，
如果兩個根的秩不同，留下來的根秩不變、節點只會變多；如果兩個根的秩都是 $r$，各至少
$2^r$ 個節點，合併後的根秩是 $r + 1$，節點至少 $2^{r+1}$ 個。 路徑壓縮不改秩，也不
改每棵樹的節點數。 一棵樹最多 $n$ 個節點，$2^r \le n$，得到 $r \le \lfloor \lg n
\rfloor$。 不是根的節點，秩停在它最後一次當根時的值，也不超過這個上界。

(3) 從 $x$ 往下走，由 (1)，每走一條邊秩至少減一，而秩不小於 $0$，所以最多走
$\mathit{rank}[x]$ 條邊。 樹的高度（數節點）是最長路徑的邊數加一，不超過
$\mathit{rank}[\text{根}] + 1 \le \lfloor \lg n \rfloor + 1$。
\end{whybox}

這個上界做得到。 $n = 8$ 時依序做 \proc{Union}$(1, 2)$、$(3, 4)$、$(5, 6)$、
$(7, 8)$、$(1, 3)$、$(5, 7)$、$(1, 5)$，每次都是兩棵一樣的樹合併：根 $1$ 的秩是 $3$，高度 $4 = \lg 8 + 1$，$p[1..8] = 1, 1,
1, 3, 1, 5, 5, 7$。 依大小合併在這串操作上做出同一棵樹。 這時每次
\proc{Find-Set}$(8)$ 都要經過 $4$ 個節點（不壓縮的話），一般來說，只用其中一種合併規
則，$m$ 個操作的總時間是 $\Theta(m \lg n)$。


\subsection{加上路徑壓縮：$O(m\,\alpha(n))$}

兩個技巧一起用，總時間的上界降到幾乎是線性。 這裡只給結論，證明要用到精細的平攤分
析（位能法，第~3~章），超出本書的範圍。

\begin{theorem}[Tarjan 1975；CLRS 第 4 版 19.4 節]
\label{thm:ds-alpha}
森林表示同時使用依秩合併和路徑壓縮時，$m$ 個操作（其中 $n$ 個是 \proc{Make-Set}）
最壞情況的總時間是 $O(m\,\alpha(n))$。
\end{theorem}

$\alpha(n)$ 是一個成長極慢的函數，是 Ackermann 函數的某種反函數。 本書用 CLRS 的版
本：對 $k \ge 0$、$j \ge 1$，令
\[
A_0(j) = j + 1, \qquad A_k(j) = \underbrace{A_{k-1}(A_{k-1}(\cdots A_{k-1}}_{j + 1
\text{ 次}}(j)\cdots)) \quad (k \ge 1),
\]
再定義 $\alpha(n) = \min\{k : A_k(1) \ge n\}$。 算幾個值：$A_1(j) = 2j + 1$，
$A_2(j) = 2^{j+1}(j + 1) - 1$，所以 $A_0(1) = 2$、$A_1(1) = 3$、$A_2(1) = 7$、$A_3(1)
= A_2(A_2(1)) = A_2(7) = 2047$，而 $A_4(1) = A_3(2047)$ 已經遠大於 $10^{80}$。 下表把
$\alpha(n)$ 和 $\lg n$、迭代對數 $\lg^* n$（$n$ 要連續取幾次 $\lg$ 才會不超過 $1$）並
排：

\begin{center}
\begin{tabular}{rrrr}
\toprule
$n$ & $\lg n$ & $\lg^* n$ & $\alpha(n)$ \\
\midrule
$10$ & $3.3$ & $3$ & $3$ \\
$10^3$ & $10.0$ & $4$ & $3$ \\
$10^6$ & $19.9$ & $5$ & $4$ \\
$10^{9}$ & $29.9$ & $5$ & $4$ \\
$10^{18}$ & $59.8$ & $5$ & $4$ \\
$2^{65536}$ & $65536$ & $5$ & $4$ \\
\bottomrule
\end{tabular}
\end{center}

任何實際能存進電腦的 $n$，$\alpha(n) \le 4$。 所以定理~\ref{thm:ds-alpha} 實際上就是
「平攤每個操作常數時間」，但在數學上 $\alpha(n)$ 會無限變大，不能寫成 $O(m)$。
Tarjan 原本的結果用的是兩個參數的 $\alpha(m, n)$，CLRS 第 4 版改用單參數的
$\alpha(n)$；對任何實際的 $m$、$n$，兩者都不超過 $4$。 更早，Hopcroft 和 Ullman（1973）證明了 $O(m \lg^* n)$ 的上界，也很
接近線性。

幾個相關的結果，一樣只列結論：
\begin{itemize}
\item 依大小合併換掉依秩合併，或路徑壓縮換成路徑減半、路徑分割，上界一樣是
$O(m\,\alpha(n))$（Tarjan 與 van Leeuwen 1984）。
\item 只用路徑壓縮、不用任何合併規則時，$n$ 個 \proc{Make-Set}、最多 $n - 1$ 個
\proc{Union} 和 $f$ 個 \proc{Find-Set} 的最壞總時間是
$\Theta\bigl(n + f \cdot (1 + \log_{2 + f/n} n)\bigr)$（同一篇論文）。
\item 在相當一般的計算模型（cell probe 模型）下，任何互斥集合的資料結構都要 $\Omega(m\,\alpha(n))$ 的
時間（Fredman 與 Saks 1989），所以這個上界已經無法再改進。
\end{itemize}

最後要分清楚\textbf{平攤}和\textbf{單次最壞}。 $O(m\,\alpha(n))$ 是 $m$ 個操作的總
和，平均分給每個操作才是 $O(\alpha(n))$。 單獨一次 \proc{Find-Set} 仍然可能走過一
條長度 $\lfloor \lg n \rfloor + 1$ 的路徑（引理~\ref{lem:ds-rank}），例如上面
$8$ 個節點的例子第一次找 $8$ 的時候。 各種做法整理如下。

\begin{center}
\small
\begin{tabular}{p{0.30\linewidth}p{0.24\linewidth}p{0.36\linewidth}}
\toprule
做法 & 單次最壞 & $m$ 個操作的總時間 \\
\midrule
串列，不加權 & \proc{Find} $O(1)$，\proc{Union} $O(n)$ & $\Theta(n^2)$ 做得到（$m = 2n - 1$） \\
串列，加權合併 & \proc{Find} $O(1)$，\proc{Union} $O(n)$ & $O(m + n \lg n)$ \\
森林，不用技巧 & $O(n)$ & $O(mn)$ \\
依大小或依秩合併 & $O(\lg n)$ & $O(m \lg n)$ \\
只用路徑壓縮 & $O(n)$ & 見上面第二點 \\
依大小或依秩＋路徑壓縮（或減半、分割） & $O(\lg n)$ & $O(m\,\alpha(n))$ \\
\bottomrule
\end{tabular}
\end{center}

\section{C 實作}
\label{sec:ds-c}

C 程式用陣列存森林：\texttt{parent[x]} 是父節點，\texttt{size[x]} 只在根上有意
義。 元素從 $0$ 編號；題目或輸入從 $1$ 編號時，讀進來先減一，印出去再加一，整段一起
換，不在演算法裡面加減。

\begin{cbookrecap}[第~14~章 14.5 節、第~15~章 15.4 節、第~19~章 19.3 節]
\texttt{DSU *d} 是指向結構的指標，\texttt{d->parent} 等於 \texttt{(*d).parent}。 結
構裡的兩個陣列用 \texttt{malloc} 要記憶體，所以結構本身只存兩個指標；用完要各
\texttt{free} 一次。 第一個 \texttt{malloc} 成功、第二個失敗時，第一塊也要放掉，
\texttt{free(NULL)} 什麼都不做，所以可以直接兩個都 \texttt{free}。

遞迴函式 \texttt{find\_rec} 每呼叫一層就在呼叫堆疊上多一個框，路徑多長就疊幾層。
\end{cbookrecap}

下面的程式實作依大小合併，以及遞迴和迴圈兩種路徑壓縮：\texttt{dsu\_union} 裡用迴
圈版，單獨的 \texttt{F} 操作用遞迴版。 兩個版本都把整條路徑接到根，結果一樣。 程式
照第~\ref{sec:ds-trace}~節的 $11$ 個操作執行，每一步印出 \texttt{parent} 陣列。

\begin{example}[label=ex:ds-trace]{c}{(dsu\_trace.c) 依大小合併與路徑壓縮，印出每一步的父節點陣列}
#include <stdio.h>
#include <stdlib.h>

typedef struct {
    int n;          // 元素是 0..n-1
    int *parent;    // parent[x] == x 表示 x 是根
    int *size;      // 只有根的 size 有意義：這棵樹有幾個節點
    int sets;       // 目前有幾個集合
} DSU;

int dsu_init(DSU *d, int n)
{
    d->parent = malloc(n * sizeof *d->parent);
    d->size = malloc(n * sizeof *d->size);
    if (d->parent == NULL || d->size == NULL) {
        free(d->parent);
        free(d->size);
        return 0;
    }
    for (int x = 0; x < n; x++) {    // n 次 Make-Set
        d->parent[x] = x;
        d->size[x] = 1;
    }
    d->n = n;
    d->sets = n;
    return 1;
}

void dsu_free(DSU *d)
{
    free(d->parent);
    free(d->size);
}

// 路徑壓縮，遞迴版：回來的路上把每個節點直接接到根。
int find_rec(int *parent, int x)
{
    if (parent[x] != x) {
        parent[x] = find_rec(parent, parent[x]);
    }
    return parent[x];
}

// 路徑壓縮，迴圈版：第一趟找根，第二趟把路上的節點接到根。
int find_iter(int *parent, int x)
{
    int root = x;
    while (parent[root] != root) {
        root = parent[root];
    }
    while (parent[x] != root) {
        int next = parent[x];
        parent[x] = root;
        x = next;
    }
    return root;
}

// 依大小合併，一樣大時 y 的根接到 x 的根底下。
int dsu_union(DSU *d, int x, int y)
{
    int rx = find_iter(d->parent, x);
    int ry = find_iter(d->parent, y);
    if (rx == ry) {
        return 0;
    }
    if (d->size[rx] < d->size[ry]) {   // 讓 rx 是比較大的那棵
        int t = rx;
        rx = ry;
        ry = t;
    }
    d->parent[ry] = rx;
    d->size[rx] += d->size[ry];
    d->sets--;
    return 1;
}

// 印出 parent 陣列。 程式從 0 編號，印的時候加 1。
void print_parent(const DSU *d)
{
    for (int x = 0; x < d->n; x++) {
        printf(x > 0 ? " %d" : "%d", d->parent[x] + 1);
    }
    printf("\n");
}

int main(void)
{
    // 手算追蹤的 11 個操作，電腦從 1 編號
    struct { char op; int a, b; } ops[] = {
        {'U', 1, 2}, {'U', 3, 4}, {'U', 5, 6}, {'U', 7, 8},
        {'U', 2, 4}, {'U', 6, 8}, {'U', 9, 10}, {'U', 6, 2},
        {'F', 4, 0}, {'U', 10, 8}, {'F', 10, 0},
    };
    DSU d;
    if (!dsu_init(&d, 10)) {
        return 1;
    }
    for (int i = 0; i < 11; i++) {
        int a = ops[i].a - 1, b = ops[i].b - 1;
        if (ops[i].op == 'U') {
            dsu_union(&d, a, b);
            printf("U %d %d: ", a + 1, b + 1);
        } else {
            int r = find_rec(d.parent, a);
            printf("F %d -> %d: ", a + 1, r + 1);
        }
        print_parent(&d);
    }
    printf("sets = %d\n", d.sets);
    dsu_free(&d);
    return 0;
}
\end{example}
\iosample{}{U 1 2: 1 1 3 4 5 6 7 8 9 10 \\ U 3 4: 1 1 3 3 5 6 7 8 9 10 \\ U 5 6: 1 1 3 3 5 5 7 8 9 10 \\ U 7 8: 1 1 3 3 5 5 7 7 9 10 \\ U 2 4: 1 1 1 3 5 5 7 7 9 10 \\ U 6 8: 1 1 1 3 5 5 5 7 9 10 \\ U 9 10: 1 1 1 3 5 5 5 7 9 9 \\ U 6 2: 5 1 1 3 5 5 5 7 9 9 \\ F 4 -> 5: 5 1 5 5 5 5 5 7 9 9 \\ U 10 8: 5 1 5 5 5 5 5 5 5 9 \\ F 10 -> 5: 5 1 5 5 5 5 5 5 5 5 \\ sets = 1}

每一列都和第~\ref{sec:ds-trace}~節的表相同。 \texttt{dsu\_union} 先 find 再比大小，
所以合併時的兩次 find 也會壓縮路徑，第 10 步的 \texttt{p[8]} 就是這樣改的。 只問
「有幾組」時不必掃描整個陣列：每次真正合併讓 \texttt{sets} 減一就好。

\subsection{應用：拔線之後還剩幾個網段}
\label{sec:ds-app}

互斥集合只會合併，不會拆開。 如果實驗室要做斷線演練：先接好 $m$ 條線，再照一張清單
依序拔掉其中 $k$ 條，每拔一條就要回報還剩幾個網段、最大的網段有幾台電腦，該怎麼辦？

清單是事先就知道的，所以可以\textbf{倒過來做}。 拔完全部 $k$ 條之後剩下的線，是
「從來沒被拔過」的那些；先把它們接上，得到最後的狀態。 接著從清單的最後一條往前，
一條一條接回去：接回第 $i$ 條之前的狀態，正好是「依序拔完第 $1$ 到第 $i$ 條」之後
的狀態。 倒著做只有合併，互斥集合就用得上了。 最大的網段也一樣：倒著做時集合只會
變大，每次合併後和目前的最大值比一下就好。

\begin{example}[label=ex:ds-lab-cut]{c}{(lab\_cut.c) 依序拔線，倒過來用互斥集合回答}
#include <stdio.h>
#include <stdlib.h>

typedef struct {
    int *parent;
    int *size;
    int sets;       // 目前有幾個網段
    int largest;    // 目前最大的網段有幾台
} DSU;

typedef struct {
    int u, v;       // 兩端的電腦（程式裡從 0 編號）
    int cut;        // 1：這條線會被拔掉
} Cable;

typedef struct {
    int sets, largest;
} State;

int find(int *parent, int x)        // 迴圈版的路徑壓縮
{
    int root = x;
    while (parent[root] != root) {
        root = parent[root];
    }
    while (parent[x] != root) {
        int next = parent[x];
        parent[x] = root;
        x = next;
    }
    return root;
}

void unite(DSU *d, int x, int y)    // 依大小合併
{
    int rx = find(d->parent, x);
    int ry = find(d->parent, y);
    if (rx == ry) {
        return;
    }
    if (d->size[rx] < d->size[ry]) {
        int t = rx;
        rx = ry;
        ry = t;
    }
    d->parent[ry] = rx;
    d->size[rx] += d->size[ry];
    d->sets--;
    if (d->size[rx] > d->largest) {     // 最大值只會變大
        d->largest = d->size[rx];
    }
}

int main(void)
{
    int n, m, k = 0;
    if (scanf("%d %d", &n, &m) != 2 || n <= 0 || m < 0) {
        return 1;
    }
    Cable *c = malloc((m + 1) * sizeof *c);       // 線 1..m
    int *order = malloc((m + 1) * sizeof *order); // 拔的順序
    State *ans = malloc((m + 1) * sizeof *ans);
    DSU d;
    d.parent = malloc(n * sizeof *d.parent);
    d.size = malloc(n * sizeof *d.size);
    int ok = c != NULL && order != NULL && ans != NULL
             && d.parent != NULL && d.size != NULL;
    for (int e = 1; ok && e <= m; e++) {
        ok = scanf("%d %d", &c[e].u, &c[e].v) == 2;
        c[e].u--;
        c[e].v--;
        c[e].cut = 0;
    }
    ok = ok && scanf("%d", &k) == 1 && k >= 0 && k <= m;
    for (int i = 1; ok && i <= k; i++) {
        ok = scanf("%d", &order[i]) == 1
             && order[i] >= 1 && order[i] <= m;
        if (ok) {
            c[order[i]].cut = 1;
        }
    }
    if (ok) {
        for (int x = 0; x < n; x++) {
            d.parent[x] = x;
            d.size[x] = 1;
        }
        d.sets = n;
        d.largest = 1;
        for (int e = 1; e <= m; e++) {      // 先接從沒拔的線
            if (!c[e].cut) {
                unite(&d, c[e].u, c[e].v);
            }
        }
        for (int i = k; i >= 1; i--) {      // 倒著把線接回去
            ans[i].sets = d.sets;       // 拔完第 i 條之後
            ans[i].largest = d.largest;
            unite(&d, c[order[i]].u, c[order[i]].v);
        }
        printf("cut ends groups largest\n");
        printf("- - %d %d\n", d.sets, d.largest);   // 還沒拔線
        for (int i = 1; i <= k; i++) {
            Cable *e = &c[order[i]];
            printf("%d %d-%d %d %d\n", order[i], e->u + 1, e->v + 1,
                   ans[i].sets, ans[i].largest);
        }
    }
    free(c);
    free(order);
    free(ans);
    free(d.parent);
    free(d.size);
    return ok ? 0 : 1;
}
\end{example}
\iosample{8 9 \\ 1 2 \\ 2 3 \\ 3 1 \\ 3 4 \\ 5 6 \\ 6 7 \\ 4 5 \\ 7 8 \\ 6 8 \\ 5 \\ 2 7 9 6 1}{cut ends groups largest \\ - - 1 8 \\ 2 2-3 1 8 \\ 7 4-5 2 4 \\ 9 6-8 2 4 \\ 6 6-7 3 4 \\ 1 1-2 4 3}

輸入第一列是電腦數和線數，接著 $m$ 條線（依序是第 $1$ 到第 $m$ 條），再來是要拔幾
條、依序拔哪幾條；假設每條線最多拔一次。 輸出第二列是還沒拔線時的狀態。 拔第 $2$ 條
（2--3）時，$1, 2, 3$ 之間還有另外兩條線，網段數不變；拔第 $7$ 條（4--5）把教室切成
兩半；第 $9$ 條（6--8）拔掉後 $6, 7, 8$ 仍然繞得過去。 每個問題只多花一次合併，總時
間是 $O((m + n)\,\alpha(n))$。 如果拔線的順序要邊做邊決定（線上的版本），這招就不能
用了，需要更複雜的資料結構，不在本書的範圍。

\section{本章常見錯誤}
\label{sec:ds-pitfalls}

\begin{center}
\begin{tabular}{p{0.24\linewidth}p{0.50\linewidth}p{0.16\linewidth}}
\toprule
錯誤 & 例子 & 見 \\
\midrule
接的不是根 & \proc{Union} 寫成 $p[y] \gets x$；$y$ 原來的根和其他節點留在舊的樹裡 &
第~\ref{sec:ds-forest}~節 \\
比父節點不比代表 & 判斷同一組時比 $p[a] = p[b]$，兩個節點在同一棵樹的不同深度就答錯 &
第~\ref{sec:ds-list}~節 \\
大小重複累加 & $x$、$y$ 已經同一組還照樣 $\mathit{size}[r_x] \gets \mathit{size}[r_x]
+ \mathit{size}[r_y]$；或把大小加在非根節點上 & 第~\ref{sec:ds-size}~節 \\
把秩當成高度 & 路徑壓縮後去更新秩；用實際樹高決定誰當根；忘了秩相等時加一，或秩不等也
加一 & 第~\ref{sec:ds-rank}~節 \\
高度差一 & 說「秩 $r$ 的樹高度是 $r$」，沒有換成本書數節點的定義（高度 $\le r + 1$） &
第~\ref{sec:ds-rank}~節 \\
\proc{Union} 裡的 find 不壓縮 & 手算時只在單獨的 \proc{Find-Set} 才壓縮，漏改了合併時
經過的節點 & 第~\ref{sec:ds-trace}~節 \\
平攤當成單次 & 說「每次 \proc{Find-Set} 最壞是 $O(\alpha(n))$」 & 第~\ref{sec:ds-cost}~節 \\
遞迴太深 & 沒有合併規則時用遞迴的 find，鏈長 $10^6$ 就用光堆疊 & 第~\ref{sec:ds-compress}~節 \\
以為可以拆開 & 想用互斥集合直接處理刪除線路；應該倒過來做，或改用別的方法 &
第~\ref{sec:ds-app}~節 \\
\bottomrule
\end{tabular}
\end{center}

\section{觀念題}
\label{sec:ds-exercises}

題型和研究所考試、期中期末考的筆試相近，有選擇、計算、追蹤和證明。 沒有特別說明時，
\proc{Union}$(x, y)$ 在兩個根一樣大（或秩相同）時讓 $y$ 的根接到 $x$ 的根底下，
\proc{Union} 裡的 \proc{Find-Set} 也照題目規定的方式壓縮；高度數節點。 解答在本章最
後。

\begin{enumerate}
\item （選擇，可能不只一個答案）下列哪些敘述正確？
\begin{enumerate}
\item 依秩合併、不做路徑壓縮時，每個節點的秩等於從它往下最長路徑的邊數。
\item 做路徑壓縮時，要把路徑上節點的秩改成新的值，否則之後的合併會選錯根。
\item 只用依大小合併時，$100$ 個元素的森林裡，最高的樹高度不超過 $7$。
\item 依秩合併加路徑壓縮時，每一次 \proc{Find-Set} 最壞都是 $O(\alpha(n))$。
\item 串列表示用加權合併時，\proc{Find-Set} 是 $O(1)$。
\end{enumerate}

\item （計算）用串列表示，\proc{Make-Set} 了 $1$ 到 $8$ 之後依序做
\proc{Union}$(1, 2)$、$(3, 1)$、$(4, 1)$、$(5, 6)$、$(5, 7)$、$(8, 5)$、$(8, 4)$。
分別在 (a) 一律把 $y$ 的串列接到 $x$ 的後面、(b) 加權合併（一樣長時接 $y$ 的）兩種
做法下，求每次合併改了幾個 $\mathit{set}$ 指標和總數，以及最後的代表。

\item \label{q:ds-rank-trace}（追蹤）$9$ 個元素，用依秩合併（虛擬碼~\ref{alg:ds-rank}）
加路徑壓縮，依序做 \proc{Union}$(1, 2)$、$(3, 4)$、$(4, 2)$、$(5, 6)$、$(7, 5)$、
$(8, 9)$、$(9, 6)$、$(4, 7)$，最後 \proc{Find-Set}$(6)$。 寫出每一步之後的 $p[1..9]$
和 $\mathit{rank}[1..9]$。 最後一步之後，根的秩和樹的高度各是多少？ 這和
第~\ref{sec:ds-rank}~節說的關係一致嗎？

\item （追蹤）用「根存負的節點數」的陣列（第~\ref{sec:ds-trace}~節作答提醒），目前
\[
p[1..8] = -3,\ 1,\ 1,\ -4,\ 4,\ 5,\ 4,\ -1.
\]
依序做 \proc{Union}$(3, 6)$、\proc{Find-Set}$(2)$、\proc{Union}$(8, 2)$，用加權規則和塌
陷規則（依大小合併、路徑壓縮）。 寫出每一步之後的陣列。

\item \label{q:ds-compare}（比較）$9$ 個元素，依序做 \proc{Union}$(2, 1)$、$(3, 1)$、$(4, 3)$、
$(2, 9)$、$(6, 5)$、$(7, 8)$、$(6, 8)$、$(5, 4)$，最後 \proc{Find-Set}$(8)$。
分別用下列五種做法，寫出最後的 $p[1..9]$、最高的樹的高度，以及最後的
\proc{Find-Set}$(8)$ 經過幾個節點。
(a) 不用任何技巧（虛擬碼~\ref{alg:ds-naive}）；(b) 依大小合併，不壓縮；(c) 依秩合
併，不壓縮；(d) 依大小合併＋路徑壓縮；(e) 依秩合併＋路徑壓縮。
再說明 (b) 和 (c) 在第 8 個操作為什麼選了不同的根。

\item （證明）只用依大小合併、不壓縮。 對合併的次數做歸納，證明：高度是 $h$ 的樹至
少有 $2^{h-1}$ 個節點。 這和定理~\ref{thm:ds-size-height} 的證明有什麼不同？

\item （證明）依秩合併（有沒有路徑壓縮都可以），證明：在任何時候，秩恰好是 $r$ 的節
點最多 $n / 2^r$ 個。（提示：把每個秩 $r$ 的節點，對應到它變成秩 $r$ 那一刻樹裡的
節點。）

\item （追蹤）森林的父節點是 $p[2] = 1$、$p[3] = 1$、$p[4] = 3$、$p[5] = 3$、
$p[6] = 5$、$p[7] = 5$、$p[8] = 7$、$p[9] = 7$，$1$ 是根。 分別用路徑壓縮、路徑減
半、路徑分割做一次 \proc{Find-Set}$(9)$，寫出之後的 $p[1..9]$ 和樹的高度。

\item （計算）
\begin{enumerate}
\item 求 $\alpha(5)$、$\alpha(100)$、$\alpha(5000)$、$\alpha(10^{30})$。
\item 求 $\lg^* 16$、$\lg^* 65536$、$\lg^* 2^{65536}$。
\item 「依秩合併＋路徑壓縮的總時間是 $O(m\,\alpha(n))$，而 $\alpha(n) \le 4$，所以總
時間是 $O(m)$。」 這句話哪裡不對？
\end{enumerate}

\item （追蹤）範例~\ref{ex:ds-lab-cut} 的線路不變，把拔線的清單改成 \texttt{4} 條：
\texttt{7 3 2 9}。 (a) 程式倒著接回去的順序是什麼？ (b) 寫出程式的輸出。 (c) 為什
麼互斥集合不能直接照拔線的順序處理？
\end{enumerate}

\section{實作練習}
\label{sec:ds-practice}

下面是 Aki-Judge 上和本章相關的題目。 題目裡的「不相交集合」「併查集」就是本章的互斥
集合。 交之前先用第~\ref{sec:ds-trace}~節的操作序列和範例~\ref{ex:ds-trace} 的輸出自己
對一次。

\begin{center}
\begin{tabular}{llc}
\toprule
小節 & 題目 & 難度 \\
\midrule
應用 & \akiref{counting-connected-components}{計算連通分量數} & 暖身 \\
森林表示 & \akiref{disjoint-set-union-data-structure}{不相交集合資料結構} & 基本 \\
 & \akiref{merges-and-connectivity-queries}{合併與連通查詢} & 基本 \\
依大小合併 & \akiref{size-of-the-largest-component}{最大連通分量的大小} & 基本 \\
 & \akiref{device-group-size}{裝置群組大小} & 進階 \\
綜合 & \akiref{power-grid-regions}{配電網路分區} & 進階 \\
\bottomrule
\end{tabular}
\end{center}

幾題的提示：

\begin{itemize}
\item 「計算連通分量數」：每次真正的合併讓連通元件數減一，或是全部做完以後數有幾個
$i$ 滿足 $\proc{Find-Set}(i) = i$。 不要用 $n - m$，重複的邊和自環不會減少分量。
\item 「不相交集合資料結構」只交標頭檔和函式。 \texttt{init} 可能出現不只一次，每次
都要把整個結構重設；\texttt{count} 就是範例~\ref{ex:ds-trace} 的 \texttt{sets}。
\item 「合併與連通查詢」的操作多達 $5 \times 10^5$ 個，兩個技巧都要用。 比較的是兩個
\proc{Find-Set} 的結果，不是兩個父節點。
\item 「最大連通分量的大小」照第~\ref{sec:ds-size}~節，大小只記在根上，已經同一組時不
能再加。
\item 「裝置群組大小」的編號可以大到 $10^9$，不能直接當陣列索引。 先把所有出現過的
編號排序、去掉重複，再用二分搜尋把編號換成 $0, 1, 2, \ldots$（排序見第~14~章，二分搜尋可以用
第~1~章觀念題的 \proc{Lower-Bound}），之後就是一般的依大小合併。 $a = b$ 時也要回答所在
群組的大小。
\item 「配電網路分區」可以拆電纜，互斥集合不能拆。 站數最多 $30$，每次查詢時把目前所
有的電纜重新做一次合併就夠快；這是「資料小、但有刪除」時的直接做法。 電纜總長要另外
加總，每條只算一次。
\end{itemize}

\section{讀完你應該會}
\label{sec:ds-checklist}

\begin{itemize}
\item 說出互斥集合的三個操作和參數 $n$、$m$ 的意思，用它寫出「依序接線、隨時查詢」
這類程序。
\item 用串列表示做合併，數出改了幾個代表指標，並證明加權合併的總時間是 $O(m + n \lg
n)$。
\item 給一串操作和規定（依大小或依秩、平手規則、要不要壓縮），寫出每一步之後的父節點
陣列（含根存負數的寫法）和秩陣列，畫出森林。
\item 分清楚秩、高度、深度：單一節點的秩是 $0$、高度是 $1$；高度 $\le$ 秩 $+ 1$；路徑壓
縮不改秩。
\item 證明依大小或依秩合併的樹高度是 $O(\lg n)$，說出加上路徑壓縮後 $O(m\,\alpha(n))$
的結論與出處，並分清楚平攤和單次最壞。
\item 寫出 C 的陣列版互斥集合（遞迴和迴圈的路徑壓縮），知道遇到刪除時可以把操作倒過
來處理。
\end{itemize}

\textbf{延伸閱讀}（書末會整理完整的參考資料）

\begin{itemize}
\item Thomas H. Cormen 等，\textit{Introduction to Algorithms}，第 4 版，第 19 章：
串列與森林兩種表示、依秩合併與路徑壓縮，19.4 節用位能法完整證明 $O(m\,\alpha(n))$；
章末的附註整理了這個問題的研究歷史。
\item Robert E. Tarjan，``Efficiency of a Good But Not Linear Set Union Algorithm''，
\textit{Journal of the ACM} 22(2)，1975：第一次證明 $\alpha$ 的上界。
\item Robert E. Tarjan、Jan van Leeuwen，``Worst-Case Analysis of Set Union
Algorithms''，\textit{Journal of the ACM} 31(2)，1984：路徑減半、路徑分割等變形，以
及只用路徑壓縮時的上界。
\item Ellis Horowitz、Sartaj Sahni、Susan Anderson-Freed，\textit{Fundamentals of Data
Structures in C}，第 5 章：根存負數的陣列寫法、加權規則和塌陷規則，附 C 程式。
\end{itemize}

\addcontentsline{toc}{section}{觀念題解答}
\section*{觀念題解答}

\begin{enumerate}
\item (a)、(c)、(e) 正確。 (a)：不壓縮時，合併兩個秩 $r$ 的根正好讓最長路徑多一條邊，
秩不同時最長路徑不變，歸納可得。 (b)：路徑壓縮不改秩，秩本來就只是上界，不需要準確。
(c)：$\lfloor \lg 100 \rfloor + 1 = 6 + 1 = 7$。 (d)：$O(\alpha(n))$ 是平攤，單次最壞
可以是 $\Theta(\lg n)$。 (e)：讀 $x.\mathit{set}.\mathit{head}$ 就好。

\item 每次合併改的指標數：
\begin{center}
\begin{tabular}{lccccccc|c}
\toprule
合併 & $(1,2)$ & $(3,1)$ & $(4,1)$ & $(5,6)$ & $(5,7)$ & $(8,5)$ & $(8,4)$ & 合計 \\
\midrule
(a) 一律接 $y$ & 1 & 2 & 3 & 1 & 1 & 3 & 4 & 15 \\
(b) 加權 & 1 & 1 & 1 & 1 & 1 & 1 & 4 & 10 \\
\bottomrule
\end{tabular}
\end{center}
(a) 的最後一條串列是 $8, 5, 6, 7, 4, 3, 1, 2$，代表 $8$。 (b) 的 $(3, 1)$、$(4, 1)$、
$(8, 5)$ 都是把單一元素接到長串列後面；最後一次兩條都是 $4$ 個，接 $y$ 的，最後的串列
是 $5, 6, 7, 8, 1, 2, 3, 4$，代表 $5$。

\item 每列是做完那一步之後的狀態：
\begin{center}
\small
\begin{tabular}{ll}
\toprule
操作 & $p[1..9]$ \qquad $\mathit{rank}[1..9]$ \\
\midrule
\proc{Union}$(1, 2)$ & $1,1,3,4,5,6,7,8,9$ \qquad $1,0,0,0,0,0,0,0,0$ \\
\proc{Union}$(3, 4)$ & $1,1,3,3,5,6,7,8,9$ \qquad $1,0,1,0,0,0,0,0,0$ \\
\proc{Union}$(4, 2)$ & $3,1,3,3,5,6,7,8,9$ \qquad $1,0,2,0,0,0,0,0,0$ \\
\proc{Union}$(5, 6)$ & $3,1,3,3,5,5,7,8,9$ \qquad $1,0,2,0,1,0,0,0,0$ \\
\proc{Union}$(7, 5)$ & $3,1,3,3,5,5,5,8,9$ \qquad $1,0,2,0,1,0,0,0,0$ \\
\proc{Union}$(8, 9)$ & $3,1,3,3,5,5,5,8,8$ \qquad $1,0,2,0,1,0,0,1,0$ \\
\proc{Union}$(9, 6)$ & $3,1,3,3,8,5,5,8,8$ \qquad $1,0,2,0,1,0,0,2,0$ \\
\proc{Union}$(4, 7)$ & $3,1,3,3,8,5,8,3,8$ \qquad $1,0,3,0,1,0,0,2,0$ \\
\proc{Find-Set}$(6)$ & $3,1,3,3,3,3,8,3,8$ \qquad $1,0,3,0,1,0,0,2,0$ \\
\bottomrule
\end{tabular}
\end{center}
幾個關鍵：$(4, 2)$ 的兩個根 $3$、$1$ 秩都是 $1$，$1$ 接到 $3$ 底下，$3$ 的秩變 $2$。
$(7, 5)$ 是秩 $0$ 的 $7$ 接到秩 $1$ 的 $5$ 底下，秩都不變。 $(9, 6)$ 的兩個根 $8$、$5$
秩都是 $1$，$5$ 接到 $8$ 底下。 $(4, 7)$ 先 \proc{Find-Set}$(7)$，路徑 $7 \to 5 \to 8$
把 $7$ 改接到 $8$；兩個根 $3$、$8$ 秩都是 $2$，$8$ 接到 $3$ 底下，$3$ 的秩變 $3$。 這時
$6 \to 5 \to 8 \to 3$ 是最長的路徑，高度 $4 = 3 + 1$。 最後 \proc{Find-Set}$(6)$ 把 $6$、
$5$ 改接到 $3$，最長的路徑剩 $7 \to 8 \to 3$（或 $9 \to 8 \to 3$、$2 \to 1 \to 3$），
高度變成 $3$，根的秩仍然是 $3$。 高度 $3 \le 3 + 1$，符合「高度 $\le$ 秩 $+ 1$」；壓縮
以後等號不再成立。

\item 每一步之後的陣列：
\begin{center}
\begin{tabular}{ll}
\toprule
操作 & $p[1..8]$ \\
\midrule
開始 & $-3,\ 1,\ 1,\ -4,\ 4,\ 5,\ 4,\ -1$ \\
\proc{Union}$(3, 6)$ & $4,\ 1,\ 1,\ -7,\ 4,\ 4,\ 4,\ -1$ \\
\proc{Find-Set}$(2)$ & $4,\ 4,\ 1,\ -7,\ 4,\ 4,\ 4,\ -1$ \\
\proc{Union}$(8, 2)$ & $4,\ 4,\ 1,\ -8,\ 4,\ 4,\ 4,\ 4$ \\
\bottomrule
\end{tabular}
\end{center}
\proc{Union}$(3, 6)$：\proc{Find-Set}$(3) = 1$（$3$ 個節點），\proc{Find-Set}$(6)$ 走過
$6, 5, 4$，把 $6$ 改接到 $4$；$1$ 的樹比較小，接到 $4$ 底下，$p[4] = -(3 + 4) = -7$。
\proc{Find-Set}$(2)$ 走過 $2, 1, 4$，$2$ 改接到 $4$。 最後 $8$ 只有一個節點，接到 $4$
底下。 注意 $p[3] = 1$ 沒變：$3$ 從來不在被壓縮的路徑上。

\item 五種做法的結果：
\begin{center}
\small
\begin{tabular}{lccc}
\toprule
做法 & 最後的 $p[1..9]$ & 最高的樹高度 & \proc{Find-Set}$(8)$ 經過 \\
\midrule
(a) 不用技巧 & $2,3,4,6,6,6,6,7,4$ & $5$ & $3$（$8, 7, 6$）\\
(b) 依大小 & $2,2,2,2,6,2,6,7,2$ & $4$ & $4$（$8, 7, 6, 2$）\\
(c) 依秩 & $2,6,2,2,6,6,6,7,2$ & $3$ & $3$（$8, 7, 6$）\\
(d) 依大小＋壓縮 & $2,2,2,2,6,2,2,2,2$ & $3$ & $4$ \\
(e) 依秩＋壓縮 & $2,6,2,2,6,6,6,6,2$ & $3$ & $3$ \\
\bottomrule
\end{tabular}
\end{center}
(a) 每次都把 $y$ 的根接到 $x$ 的根底下：前三次做出鏈 $1 \to 2 \to 3 \to 4$，
\proc{Union}$(2, 9)$ 讓 $9$ 接到 $4$，最後 \proc{Union}$(5, 4)$ 讓 $4$ 接到 $6$，最長的
路徑是 $1 \to 2 \to 3 \to 4 \to 6$。 (b)、(c) 前七個操作的結果相同：$2$ 是 $5$ 個節點、
高度 $2$ 的星形（秩 $1$），$6$ 是 $4$ 個節點、高度 $3$ 的樹（秩 $2$）。 第 8 個操作依大小
選節點多的 $2$ 當根，高的樹被接到下面，高度變 $4$；依秩選秩大的 $6$ 當根，高度維持 $3$。
(d)、(e) 在前八個操作都沒有可以壓縮的路徑（每次 find 的節點都在根或根的子節點），所以
和 (b)、(c) 一樣，最後的 \proc{Find-Set}$(8)$ 才壓縮：(d) 把 $8$ 和 $7$ 接到 $2$，(e) 把
$8$ 接到 $6$。 兩個合併規則都保證 $O(\lg n)$ 的高度，但單看一串操作，誰比較矮不一定。

\item 歸納的對象是「目前森林裡的每棵樹都滿足：高度 $h$ 的樹至少 $2^{h-1}$ 個節點」。
一開始每棵樹只有一個節點，$h = 1$，$2^0 = 1$，成立。 合併兩棵樹 $T_x$（$s_x$ 個節點、
高度 $h_x$）和 $T_y$（$s_y$ 個節點、高度 $h_y$），設 $s_x \ge s_y$，$T_y$ 的根接到
$T_x$ 的根底下，新的高度是 $h = \max(h_x, h_y + 1)$。 若 $h = h_x$，節點數 $s_x + s_y
> s_x \ge 2^{h_x - 1}$。 若 $h = h_y + 1$，節點數 $s_x + s_y \ge 2 s_y \ge 2 \cdot
2^{h_y - 1} = 2^{h - 1}$。 其他樹不變，所以性質一直成立；$s \ge 2^{h-1}$ 等於 $h \le \lg
s + 1$。 定理~\ref{thm:ds-size-height} 是追蹤\textbf{單一節點}的深度什麼時候增加，這裡
是對\textbf{整棵樹}做歸納；兩種都常見，結論相同。

\item 對每個秩是 $r$ 的節點 $x$，看它的秩變成 $r$ 的那一刻：那時 $x$ 是根，由
引理~\ref{lem:ds-rank} (2)，它的樹至少有 $2^r$ 個節點，把這些節點記成 $N(x)$。 只要證明
不同的 $x$ 得到的 $N(x)$ 兩兩不相交，就有（秩 $r$ 的節點數）$\cdot\, 2^r \le n$。

假設 $z$ 同時在 $N(x)$ 和 $N(y)$ 裡，$x \ne y$，而且 $x$ 比 $y$ 先變成秩 $r$。 集合只會
合併、不會拆開，所以 $x$ 變成秩 $r$ 之後，$z$ 和 $x$ 一直在同一棵樹裡。 到了 $y$ 變成秩
$r$ 的那一刻，$z$ 在 $y$ 的樹裡，所以 $x$ 也在 $y$ 的樹裡，而 $y$ 是根、$x \ne y$，$x$
是 $y$ 的子孫。 由引理~\ref{lem:ds-rank} (1)，沿著 $x$ 到 $y$ 的路徑秩嚴格變大，
$\mathit{rank}[x] < \mathit{rank}[y] = r$；但秩不會變小，$\mathit{rank}[x] \ge r$，矛盾。
所以秩 $r$ 的節點最多 $n / 2^r$ 個。 （這也再一次說明秩不超過 $\lg n$。）

\item 路徑是 $9 \to 7 \to 5 \to 3 \to 1$。
\begin{center}
\begin{tabular}{lcc}
\toprule
做法 & 之後的 $p[1..9]$ & 高度 \\
\midrule
路徑壓縮 & $1,1,1,3,1,5,1,7,1$ & $3$ \\
路徑減半 & $1,1,1,3,1,5,5,7,5$ & $4$ \\
路徑分割 & $1,1,1,3,1,5,3,7,5$ & $4$ \\
\bottomrule
\end{tabular}
\end{center}
壓縮：$9, 7, 5$ 都改接到 $1$（$3$ 本來就是）；最長的路徑變成 $4 \to 3 \to 1$ 這類，高度
$3$。 減半：$9$ 改指祖父 $5$，跳到 $5$；$5$ 改指祖父 $1$，跳到 $1$ 結束；$7$ 沒有改。
分割：$9 \to 5$、$7 \to 3$、$5 \to 1$；$3$ 的父節點就是根，改指祖父還是根，不變。
減半之後還有 $8 \to 7 \to 5 \to 1$，分割之後還有 $8 \to 7 \to 3 \to 1$，都是 $4$ 個節點的
路徑，高度 $4$。

\item (a) $A_2(1) = 7 \ge 5$ 且 $A_1(1) = 3 < 5$，所以 $\alpha(5) = 2$。 $A_2(1) = 7 < 100 \le
A_3(1) = 2047$，$\alpha(100) = 3$。 $5000 > 2047$，$\alpha(5000) = 4$。 $10^{30} < A_4(1)$，
$\alpha(10^{30}) = 4$。

(b) $16 \to 4 \to 2 \to 1$，$\lg^* 16 = 3$；$65536 \to 16 \to 4 \to 2 \to 1$，
$\lg^* 65536 = 4$；$\lg^* 2^{65536} = 5$。

(c) $\alpha(n) \le 4$ 只對「實際會出現的 $n$」成立。 $O$ 描述的是 $n \to \infty$ 的行為，
$\alpha(n)$ 會無限變大，所以數學上不能寫成 $O(m)$；而且 Fredman 和 Saks 證明了
$\Omega(m\,\alpha(n))$ 的下界，總時間確實不是 $O(m)$。 在工程上說「幾乎線性」是可以的。

\item (a) 倒著接回去的順序是第 $9$、$2$、$3$、$7$ 條。 程式先接從沒被拔的第 $1, 4, 5,
6, 8$ 條：$\{1, 2\}$、$\{3, 4\}$、$\{5, 6, 7, 8\}$。
(b) 輸出：
\begin{center}
\ttfamily
\begin{tabular}{l}
cut ends groups largest \\
- - 1 8 \\
7 4-5 2 4 \\
3 3-1 2 4 \\
2 2-3 3 4 \\
9 6-8 3 4 \\
\end{tabular}
\end{center}
拔掉 4--5 把教室分成 $\{1, 2, 3, 4\}$ 和 $\{5, 6, 7, 8\}$；再拔 3--1，$1, 2, 3$ 之間還有
1--2、2--3，不變；拔 2--3 之後 $\{1, 2\}$ 和 $\{3, 4\}$ 分開；拔 6--8 不影響 $5$ 到 $8$。
(c) 拔線會讓一個集合分成兩個，互斥集合沒有拆開的操作；父節點陣列也沒有記錄是哪幾條線
把樹連起來的，拔掉一條線時無從知道該拆哪裡。
\end{enumerate}


% 第三部：排序與選擇
\chapter{排序}
\label{ch:sorting}

% 本章的圖共用的樣式與巨集（名稱都以 so 開頭，避免和其他章衝突）
\tikzset{
    socell/.style={draw=black!55, minimum width=0.62cm, minimum height=0.5cm, inner sep=0pt, font=\small},
    socellg/.style={socell, draw=black!30, text=black!45},
    soarr/.style={-{Stealth[length=1.6mm]}, draw=black!55},
    son/.style={circle, draw=black!65, minimum size=0.62cm, inner sep=0pt, font=\small},
    soa/.style={son, draw=memStroke, thick, fill=memFill},
    sog/.style={son, draw=black!30, text=black!40, dashed},
    soedgen/.style={draw=black!55},
    soedgea/.style={draw=black!55},
    soedgeg/.style={draw=black!25, dashed},
    soreg/.style={draw=black!60, minimum height=0.62cm, inner sep=0pt, font=\small},
}
% 已經排好、不再參與的格子（灰色字）
\newcommand{\sofin}[1]{\textcolor{black!45}{#1}}
% \soheap{根的座標}{最底層的寬度}{值/樣式,值/樣式,...}{節點名稱前綴}
% 依陣列順序（從 1 開始）畫一棵完整二元樹。樣式是 n（一般）、a（標示）、g（已排好）。
\newcommand{\soheap}[4]{%
    \foreach \sov/\sos [count=\soi] in {#3} {%
        \pgfmathtruncatemacro{\solev}{floor(ln(\soi)/ln(2) + 0.0001)}%
        \pgfmathsetmacro{\sox}{(\soi - 2^\solev + 0.5) * #2 / 2^\solev - #2 / 2}%
        \node[so\sos] (#4\soi) at ($(#1) + (\sox, -0.95*\solev)$) {\sov};%
        \ifnum\soi>1
            \pgfmathtruncatemacro{\sop}{floor(\soi/2)}%
            \draw[soedge\sos] (#4\sop) -- (#4\soi);%
        \fi
    }%
}

助教收齊了一班的期中考卷，考卷已經照座號疊好。 老師要她依分數由低到高重新排一
次，還交代同分的考卷要維持座號的順序，登記成績時比較好找。 她把考卷在桌上攤成
一列，想到兩種做法：第一種是從左到右，每次拿起下一張，往左邊已經排好的那幾張裡
插進適當的位置；第二種是每次從還沒排的考卷裡找出分數最低的一張，和還沒排的第一
張對調位置。 兩種做法都能把考卷依分數排好，但只有一種一定做得到老師的第二個要
求。 是哪一種？ 往下讀之前，先想一想。

第~1~章用插入排序示範了怎麼逐行算執行次數。 這一章把常見的\textbf{比較排序}
（comparison sort）放在一起看：它們只靠「比較兩個鍵值誰大」來決定順序。 第
\ref{sec:so-terms}~節先定義穩定、原地這些用來比較排序方法的名詞。 第
\ref{sec:so-simple}~節比較氣泡、選擇、插入三種 $\Theta(n^2)$ 的排序，第
\ref{sec:so-shell}~節的 Shell 排序是插入排序的改良。 第
\ref{sec:so-merge}--\ref{sec:so-heap}~節是三種 $O(n \lg n)$ 的方法：合併排序、快速排
序、堆積排序，其中快速排序的平均情況會用指示變數完整證明。 最後一節整理成一
張比較表，並用同一組資料把每一種排序逐趟跑一次，這是考卷上很常見的題型。

比較排序在最壞情況至少要比 $\Omega(n \lg n)$ 次，這個下界和不靠比較、可以做到線性時
間的計數排序、基數排序，以及資料放不進記憶體時的外部排序，都在第~15~章。 「只找第
$k$ 小，不必全部排好」的選擇問題在第~16~章。

這一章用到第~1~章的插入排序、迴圈不變量和逆序對，第~2~章的遞迴式，第~5~章串列的
合併與切半，第~8~章的堆積，以及第~12~章的指示變數。 C 程式用到《C 程式設計基礎》第~9~章的函式指標、第
~15~章的遞迴、第~18~章的 \texttt{qsort} 和第~19~章的 \texttt{malloc}，用到時先在框
裡複習。

\section{名詞：穩定、原地、比較次數}
\label{sec:so-terms}

\textbf{排序}的輸入是 $n$ 筆\textbf{紀錄}（record），每一筆有一個\textbf{鍵}
（key），輸出是這 $n$ 筆紀錄的一個排列，使得鍵由小到大（不遞減）。 紀錄除了鍵還帶著別的
資料，例如考卷除了分數還有座號和姓名，排序時跟著鍵一起搬；紀錄很大時，常改成排序指向紀錄
的指標，每次只搬一個指標。 這一章的虛擬碼只寫鍵。

\begin{definition}[穩定的排序]
\label{def:so-stable}
如果輸入裡任兩筆鍵相同的紀錄 $x$、$y$，$x$ 在 $y$ 前面，排序之後 $x$ 仍然在 $y$ 前
面，就說這個排序方法是\textbf{穩定}的（stable）。
\end{definition}

鍵會重複時，穩定的排序保留了輸入原本的次序：章首的考卷照座號疊好，用穩定的方法依分數排，
同分的考卷就自動依座號排。 一般地說，「先依次要的鍵排，再用穩定的方法依主要的鍵排」，結
果就是依（主要鍵，次要鍵）排好，第~15~章的基數排序就是這樣做的。 任何排序都可以改成穩定
的：把鍵換成（鍵，原本的位置），鍵相同時比位置，代價是每筆紀錄多存一個整數。

除了輸入的陣列以外，只用 $O(1)$ 個元素大小的額外空間的排序，叫\textbf{原地}（in place）
排序。 遞迴用到的呼叫堆疊，本書另外寫出來；有些課本把 $O(\lg n)$ 的堆疊也算在原地裡，遇
到時看題目怎麼定義。

比較排序只透過比較兩個鍵來得到順序的資訊，所以\textbf{比較次數}是衡量它們最常用的量；紀
錄很大時，\textbf{移動次數}（交換或挪動）也要算。 第~1~章的\textbf{逆序對}個數 $I$ 可以
看成陣列亂的程度，從 $0$（排好）到 $n(n-1)/2$（完全相反），$n!$ 種排列平均是
$n(n-1)/4$。 執行時間會隨 $I$ 變小而變小的排序叫\textbf{自適應}（adaptive）的排序，插入
排序的 $\Theta(n + I)$ 就是典型的例子。

\section{氣泡、選擇、插入}
\label{sec:so-simple}

這三種排序都是兩層迴圈，最壞情況都是 $\Theta(n^2)$，考試卻常常考它們的差別：比較和
交換各幾次、穩不穩定、已經排好時會不會比較快、跑完第 $k$ 趟陣列長什麼樣子。 插入
排序在第~1~章已經逐行分析過（第~1~章虛擬碼~\ref{alg:insertion}），這一節補上另外兩種，最後放在一
起比較。

\subsection{氣泡排序}

\textbf{氣泡排序}（bubble sort）一趟從左到右比較相鄰的兩個數，左邊比較大就交換。
一趟走完，最大的數會被一路換到最右邊，像氣泡浮上來一樣。 第二趟不必再看最後一格，
以此類推。 某一趟完全沒有交換時，陣列已經排好，可以提早結束。

\begin{pseudo}[label=alg:so-bubble]{氣泡排序（沒有交換就提早結束）}
\Procedure{Bubble-Sort}{$A, n$}
    \For{$i \gets 1$ \To $n - 1$}
        \State $\mathit{swapped} \gets \FALSE$
        \For{$j \gets 1$ \To $n - i$}
            \If{$A[j] > A[j + 1]$}
                \State 交換 $A[j]$ 和 $A[j + 1]$
                \State $\mathit{swapped} \gets \TRUE$
            \EndIf
        \EndFor
        \If{$\mathit{swapped} = \FALSE$}
            \State \Return
        \EndIf
    \EndFor
\EndProcedure
\end{pseudo}

第 $i$ 趟比較 $n - i$ 次，結束時 $A[n-i+1..n]$ 是最大的 $i$ 個數，已經排好。

\begin{itemize}
\item \textbf{比較}：輸入已經排好時，第一趟比 $n - 1$ 次、沒有交換，第 8、9 行就結束，
$\Theta(n)$。 最壞情況做滿 $n - 1$ 趟，共 $n(n-1)/2$ 次。 沒有第 3、7、8、9 行（不提早結
束）的版本不管輸入如何都比 $n(n-1)/2$ 次。
\item \textbf{交換}：每次交換相鄰的兩個數，剛好消掉一個逆序對，和第~1~章插入排序的挪動
一樣，所以交換次數等於逆序對的個數 $I$。
\item \textbf{趟數}：每一趟裡，左邊有比它大的數的那些數，都恰好被這一趟往右帶的最大值越過
一次、往左移一格。 所以趟數是 $1$ 加上「左邊比它大的數的個數」的最大值，最多 $n - 1$ 趟；
一個很小的數放在最右邊，就要走滿 $n - 1$ 趟。
\item \textbf{穩定}：只有 $A[j] > A[j+1]$ 時才交換，相等的兩個數不會越過對方。 把第 5 行
改成 $\ge$ 就不穩定了。
\end{itemize}

\subsection{選擇排序}

\textbf{選擇排序}（selection sort）第 $i$ 趟從 $A[i..n]$ 裡找出最小的數，和 $A[i]$ 交
換。 做完第 $i$ 趟，$A[1..i]$ 就是最小的 $i$ 個數，而且排好了。

\begin{pseudo}[label=alg:so-selection]{選擇排序}
\Procedure{Selection-Sort}{$A, n$}
    \For{$i \gets 1$ \To $n - 1$}
        \State $m \gets i$ \Comment{目前看到的最小值的位置}
        \For{$j \gets i + 1$ \To $n$}
            \If{$A[j] < A[m]$}
                \State $m \gets j$
            \EndIf
        \EndFor
        \State 交換 $A[i]$ 和 $A[m]$
    \EndFor
\EndProcedure
\end{pseudo}

\begin{itemize}
\item \textbf{比較}：第 $i$ 趟比 $n - i$ 次，不管輸入長什麼樣子都一樣，共
$n(n-1)/2$ 次。 已經排好的輸入也不會比較快，它不是自適應的。
\item \textbf{交換}：每趟一次，共 $n - 1$ 次（$m = i$ 時等於沒動）。 這是本章所有
排序裡寫入最少的，紀錄很大、搬動很貴時有用。
\item \textbf{不穩定}：遠距離的交換會讓相同的鍵越過對方。 $\langle 5_a, 5_b, 2
\rangle$（下標只是用來分辨兩個 $5$）第一趟把 $2$ 和 $5_a$ 交換，得到
$\langle 2, 5_b, 5_a \rangle$，之後不再變動，兩個 $5$ 的先後反了。 章首的第二種做
法就是選擇排序，所以它做不到老師的第二個要求；第一種做法是插入排序，做得到。
\end{itemize}

\subsection{三種排序的比較}

插入排序（第~1~章虛擬碼~\ref{alg:insertion}）只在 $A[j] > \mathit{key}$ 時挪動，相等的數不會越過
$\mathit{key}$，所以是穩定的；挪動次數等於 $I$，比較次數介於 $I$ 和 $I + n - 1$ 之間
（第~1~章觀念題第 10 題）。 下表 $I$ 是逆序對個數，「平均」是 $n!$ 種排列機會一樣時的期望值。

\begin{center}
\begin{tabular}{llll}
\toprule
& 氣泡（提早結束） & 選擇 & 插入 \\
\midrule
比較：最好 & $n - 1$ & $n(n-1)/2$ & $n - 1$ \\
比較：最壞 & $n(n-1)/2$ & $n(n-1)/2$ & $n(n-1)/2$ \\
交換或挪動 & $I$ 次交換 & $n - 1$ 次交換 & $I$ 次挪動 \\
交換或挪動：平均 & $n(n-1)/4$ & $n - 1$ & $n(n-1)/4$ \\
第 $k$ 趟之後 & 最後 $k$ 格就定位 & 最前 $k$ 格就定位 & 前 $k + 1$ 格排好 \\
穩定 & 是 & 否 & 是 \\
自適應 & 是 & 否 & 是 \\
\bottomrule
\end{tabular}
\end{center}

「就定位」是說那一格已經是最後的答案，之後不會再變。 插入排序前段雖然排好了，後
面的數還可能插進來，所以只是「排好」。 這一點是逐趟追蹤題用來判斷「這是哪一種排
序跑了幾趟」的依據。

三者裡實務上最常用的是插入排序：氣泡排序的交換和插入排序的挪動一樣多，但一次交換要三次
指定；$n$ 很小時插入排序甚至比 $O(n \lg n)$ 的方法快。 C 程式在範例~\ref{ex:so-stability}。

\section{Shell 排序}
\label{sec:so-shell}

插入排序慢在每次只能挪一格：一個本來在最右邊、其實很小的數，要和前面每一個數各比
一次、挪一次。 D. L. Shell 在 1959 年提出的改良是：先讓距離遠的數互相比較，把它們
大致放到正確的區域，最後再做一次普通的插入排序收尾。 那時候陣列已經幾乎排好，插入
排序只要做很少的挪動。

\subsection{$h$ 排序和增量序列}

\begin{definition}[$h$ 排序好]
\label{def:so-hsorted}
如果每個 $i$（$1 \le i \le n - h$）都有 $A[i] \le A[i + h]$，就說陣列是 \textbf{$h$ 排
序好}的（$h$-sorted）。
\end{definition}

$h$ 排序好的意思是：把位置 $1, 1 + h, 1 + 2h, \ldots$ 抽出來是排好的，位置 $2, 2 + h,
\ldots$ 抽出來也是排好的……一共 $h$ 個子序列各自排好。 $1$ 排序好就是整個陣列排好。
要把陣列變成 $h$ 排序好，只要把插入排序裡的「往左一格」改成「往左 $h$ 格」。

\textbf{Shell 排序}（Shellsort）取一串遞減的\textbf{增量}（increment，也叫 gap）
$h_1 > h_2 > \cdots > h_t = 1$，依序對每個 $h_s$ 做一次 $h_s$ 排序。 最後一個增量一定
要是 $1$，這一趟就是普通的插入排序，所以不管前面的增量怎麼選，結果一定是排好的；前面
幾趟只影響速度。

\begin{pseudo}[label=alg:so-shell]{Shell 排序，增量是 $H[1..t]$，$H[t] = 1$}
\Procedure{Shell-Sort}{$A, n, H, t$}
    \For{$s \gets 1$ \To $t$}
        \State $h \gets H[s]$
        \For{$i \gets h + 1$ \To $n$} \Comment{間隔 $h$ 的插入排序}
            \State $\mathit{key} \gets A[i]$
            \State $j \gets i - h$
            \While{$j \ge 1$ \AND $A[j] > \mathit{key}$}
                \State $A[j + h] \gets A[j]$
                \State $j \gets j - h$
            \EndWhile
            \State $A[j + h] \gets \mathit{key}$
        \EndFor
    \EndFor
\EndProcedure
\end{pseudo}

第 4 行的 $i$ 由小到大走，$h$ 個子序列是交錯著處理的：$i = h + 1$ 處理第一個子序列，
$i = h + 2$ 處理第二個……效果和把每個子序列分開做插入排序一樣，因為不同子序列的數不
會互相比較。

\subsection{逐趟追蹤}

用 Shell 原本的增量 $\lfloor n/2 \rfloor, \lfloor n/4 \rfloor, \ldots, 1$ 排 $A = \langle
58, 23, 91, 7, 46, 14, 80, 35, 62, 19 \rangle$，$n = 10$，增量是 $5, 2, 1$。 下表每一
列是一趟做完之後的陣列，「比較」是第 7 行 $A[j] > \mathit{key}$ 比了幾次，「挪動」是第 8 行
執行幾次。

\begin{center}
\begin{tabular}{lcc*{10}{c}}
\toprule
& 比較 & 挪動 & 1 & 2 & 3 & 4 & 5 & 6 & 7 & 8 & 9 & 10 \\
\midrule
開始 & & & 58 & 23 & 91 & 7 & 46 & 14 & 80 & 35 & 62 & 19 \\
$h = 5$ & 5 & 3 & 14 & 23 & 35 & 7 & 19 & 58 & 80 & 91 & 62 & 46 \\
$h = 2$ & 12 & 5 & 14 & 7 & 19 & 23 & 35 & 46 & 62 & 58 & 80 & 91 \\
$h = 1$ & 10 & 2 & 7 & 14 & 19 & 23 & 35 & 46 & 58 & 62 & 80 & 91 \\
\midrule
合計 & 27 & 10 & \multicolumn{10}{l}{} \\
\bottomrule
\end{tabular}
\end{center}

$h = 5$ 的五個子序列是 $(58, 14)$、$(23, 80)$、$(91, 35)$、$(7, 62)$、$(46, 19)$，三組要
對調。 $h = 2$ 把奇數位置 $\langle 14, 35, 19, 80, 62 \rangle$ 和偶數位置 $\langle 23, 7, 58,
91, 46 \rangle$ 各自排好再交錯放回去。 最後只剩 $(14, 7)$、$(62, 58)$ 兩個逆序對，$h = 1$ 只
挪兩次。

同一個陣列直接做插入排序，要比較 31 次、挪動 24 次。 $n = 10$ 時差別還不
大，$n$ 變大以後的差距見本節最後。

Shell 排序不穩定。 $\langle 3, 2_a, 2_b, 1 \rangle$ 用增量 $2, 1$：$h = 2$ 時位置 $(1, 3)$ 是
$(3, 2_b)$，對調；位置 $(2, 4)$ 是 $(2_a, 1)$，對調，得到 $\langle 2_b, 1, 3, 2_a
\rangle$。 最後一趟插入排序不會讓兩個 $2$ 互相越過，結果是 $\langle 1, 2_b, 2_a, 3
\rangle$。 間隔 $h$ 的移動會跳過中間的數，相同的鍵就可能越過對方。

\subsection{增量怎麼選}

為什麼前面幾趟的成果不會被後面的趟數破壞？ 有一個不容易證明、但很重要的性質（證明
見延伸閱讀的 Knuth）：一個 $h$ 排序好的陣列再做 $k$ 排序之後，仍然是 $h$ 排序好的。 所
以每一趟都是在前面的基礎上再整理，逆序對一趟比一趟少。

增量的選法決定了最壞情況。 下表是幾種常見的序列，都由大到小用，用到 $1$ 為止。

\begin{center}
\begin{tabular}{lll}
\toprule
序列 & 前幾項 & 最壞情況 \\
\midrule
Shell（1959），每次減半 & $\lfloor n/2 \rfloor, \lfloor n/4 \rfloor, \ldots, 1$ & $\Theta(n^2)$ \\
Hibbard（1963），$2^k - 1$ & $1, 3, 7, 15, 31, \ldots$ & $\Theta(n^{3/2})$ \\
Knuth，$(3^k - 1)/2$ & $1, 4, 13, 40, 121, \ldots$ & $O(n^{3/2})$ \\
Sedgewick（1986），$4^k + 3 \cdot 2^{k-1} + 1$ & $1, 8, 23, 77, 281, \ldots$ & $O(n^{4/3})$ \\
Pratt（1971），$2^p 3^q$ & $1, 2, 3, 4, 6, 8, 9, 12, \ldots$ & $\Theta(n \lg^2 n)$ \\
\bottomrule
\end{tabular}
\end{center}

Pratt 的序列漸近最好，但項數有 $\Theta(\lg^2 n)$ 個，實際上反而比較慢。 實務上常用
Knuth 的序列，或是 Ciura 用實驗找出來的 $1, 4, 10, 23, 57, 132, 301, 701$。 Shell 排序的
平均情況，對大部分序列都還沒有完整的理論結果。

Shell 原本的減半序列為什麼最壞是 $\Theta(n^2)$？ $n$ 是 $2$ 的次方時，除了最後的 $1$，
每個增量都是偶數，所以奇數位置的數和偶數位置的數，在最後一趟之前從來不會互相比較。
如果最小的一半全部在奇數位置、最大的一半全部在偶數位置，前面那些趟數完全沒有幫忙。 例
如 $n = 4096$，奇數位置依序放 $1..2048$、偶數位置依序放 $2049..4096$：前面 $11$ 個偶數增量
只做了比較、什麼都沒挪，最後一趟要挪 $2048 \cdot 2047 / 2 = 2{,}096{,}128$ 次，比直接做插入
排序還多了前面 $11$ 趟的比較。 Knuth 的序列裡有奇數，奇偶位置在前面幾趟就混在一起了，
不會有這個問題。

實際量一下。 用固定種子的虛擬亂數產生 $n = 16000$ 個數，插入排序比了大約 $6.4 \times
10^7$ 次，減半序列大約 $4.9 \times 10^5$ 次，Knuth 的序列大約 $4.2 \times 10^5$ 次，少了一百
倍以上；$n$ 每變成 $4$ 倍，插入排序的次數變成大約 $16$ 倍，兩種 Shell 排序大約 $5$ 到 $6$
倍。 上面那個 $n = 4096$ 的壞輸入，減半序列比了 $2{,}141{,}185$ 次，比插入排序的
$2{,}100{,}223$ 次還多，Knuth 的序列只比了 $45{,}532$ 次。 Shell 排序的 C 程式只是把插入排
序的 $1$ 換成 $h$，寫在第~\ref{sec:so-compare}~節的範例~\ref{ex:so-stability} 裡。

\section{合併排序}
\label{sec:so-merge}

前面幾種排序都是「每次把一個數放到對的地方」。 \textbf{合併排序}（merge sort）換一個
方向：把陣列切成兩半，兩半各自排好，再把兩段排好的序列\textbf{合併}（merge）成一段。
各自排好的那兩半，用同樣的方法遞迴處理；只剩一個數時本來就排好了。 這是\textbf{分
治法}（divide and conquer）最典型的例子，第~17~章會再看到很多。

\subsection{合併兩段排好的子陣列}

第~5~章第~\ref{sec:ll-ops}~節合併過兩條排好的串列（虛擬碼~\ref{alg:ll-merge}）。 陣列版的想法一樣，兩邊
各一個指標，每次取比較小的那個；不同的是陣列不能直接把節點接過去，所以先把
$A[p..r]$ 抄到暫存陣列 $B$，再一個一個寫回 $A$。

\begin{pseudo}[label=alg:so-merge]{合併 $A[p..q]$ 和 $A[q+1..r]$，$B$ 是暫存陣列}
\Procedure{Merge}{$A, B, p, q, r$}
    \For{$k \gets p$ \To $r$}
        \State $B[k] \gets A[k]$
    \EndFor
    \State $i \gets p$；$j \gets q + 1$
    \For{$k \gets p$ \To $r$}
        \If{$i \le q$ \AND ($j > r$ \OR $B[i] \le B[j]$)}
            \State $A[k] \gets B[i]$；$i \gets i + 1$
        \Else
            \State $A[k] \gets B[j]$；$j \gets j + 1$
        \EndIf
    \EndFor
\EndProcedure
\end{pseudo}

第 6 行的條件是「左邊還有，而且右邊已經用完或左邊的比較小」。 \textbf{and}、
\textbf{or} 照第~1~章的約定，左邊決定結果時右邊不算，所以 $B[i] \le B[j]$ 只在兩邊都
還有的時候才比較，不會讀到範圍外。 兩點要注意：

\begin{itemize}
\item 用 $\le$ 而不是 $<$：鍵相同時先拿左邊的，左半的紀錄本來就在右半前面，所以合併
排序是穩定的。 改成 $<$ 就變成先拿右邊的，不穩定。
\item 兩段長度是 $m_1$、$m_2$ 時，比較次數介於 $\min(m_1, m_2)$ 和 $m_1 + m_2 - 1$ 之
間，但寫入一定是 $m_1 + m_2$ 次，抄到 $B$ 也是，所以時間是 $\Theta(m_1 + m_2)$。
\end{itemize}

\subsection{遞迴版}

\begin{pseudo}[label=alg:so-mergesort]{合併排序（遞迴）}
\Procedure{Merge-Sort}{$A, B, p, r$}
    \If{$p < r$}
        \State $q \gets \lfloor (p + r)/2 \rfloor$
        \State \Call{Merge-Sort}{$A, B, p, q$}
        \State \Call{Merge-Sort}{$A, B, q + 1, r$}
        \State \Call{Merge}{$A, B, p, q, r$}
    \EndIf
\EndProcedure
\end{pseudo}

呼叫 \proc{Merge-Sort}$(A, B, 1, n)$ 排整個陣列。 長度 $n = r - p + 1$ 的子陣列，左半
有 $\lceil n/2 \rceil$ 個，右半有 $\lfloor n/2 \rfloor$ 個。

用 $A = \langle 52, 27, 8, 66, 41, 13 \rangle$ 跑一次。 下圖上半是一路切下去的樣子，
下半是一路合併上來的結果；灰色的是已經只剩一個、不再切的數，直接往下帶。

\begin{center}
\begin{tikzpicture}
\node[socell] at (-1.550,0.000) {52};
\node[socell] at (-0.930,0.000) {27};
\node[socell] at (-0.310,0.000) {8};
\node[socell] at (0.310,0.000) {66};
\node[socell] at (0.930,0.000) {41};
\node[socell] at (1.550,0.000) {13};
\node[socell] at (-1.710,-0.800) {52};
\node[socell] at (-1.090,-0.800) {27};
\node[socell] at (-0.470,-0.800) {8};
\node[socell] at (0.470,-0.800) {66};
\node[socell] at (1.090,-0.800) {41};
\node[socell] at (1.710,-0.800) {13};
\node[socell] at (-2.030,-1.600) {52};
\node[socell] at (-1.410,-1.600) {27};
\node[socell] at (-0.470,-1.600) {8};
\node[socell] at (0.470,-1.600) {66};
\node[socell] at (1.090,-1.600) {41};
\node[socell] at (2.030,-1.600) {13};
\node[socell] at (-2.350,-2.400) {52};
\node[socell] at (-1.410,-2.400) {27};
\node[socellg] at (-0.470,-2.400) {8};
\node[socell] at (0.470,-2.400) {66};
\node[socell] at (1.410,-2.400) {41};
\node[socellg] at (2.350,-2.400) {13};
\node[socell] at (-2.030,-3.200) {27};
\node[socell] at (-1.410,-3.200) {52};
\node[socell] at (-0.470,-3.200) {8};
\node[socell] at (0.470,-3.200) {41};
\node[socell] at (1.090,-3.200) {66};
\node[socell] at (2.030,-3.200) {13};
\node[socell] at (-1.710,-4.000) {8};
\node[socell] at (-1.090,-4.000) {27};
\node[socell] at (-0.470,-4.000) {52};
\node[socell] at (0.470,-4.000) {13};
\node[socell] at (1.090,-4.000) {41};
\node[socell] at (1.710,-4.000) {66};
\node[socell] at (-1.550,-4.800) {8};
\node[socell] at (-0.930,-4.800) {13};
\node[socell] at (-0.310,-4.800) {27};
\node[socell] at (0.310,-4.800) {41};
\node[socell] at (0.930,-4.800) {52};
\node[socell] at (1.550,-4.800) {66};
\draw[soarr] (0.000,-0.300) -- (-1.090,-0.500);
\draw[soarr, draw=memStroke] (-1.090,-4.300) -- (0.000,-4.500);
\draw[soarr] (0.000,-0.300) -- (1.090,-0.500);
\draw[soarr, draw=memStroke] (1.090,-4.300) -- (0.000,-4.500);
\draw[soarr] (-1.090,-1.100) -- (-1.720,-1.300);
\draw[soarr, draw=memStroke] (-1.720,-3.500) -- (-1.090,-3.700);
\draw[soarr] (-1.090,-1.100) -- (-0.470,-1.300);
\draw[soarr, draw=memStroke] (-0.470,-3.500) -- (-1.090,-3.700);
\draw[soarr] (-1.720,-1.900) -- (-2.350,-2.100);
\draw[soarr, draw=memStroke] (-2.350,-2.700) -- (-1.720,-2.900);
\draw[soarr] (-1.720,-1.900) -- (-1.410,-2.100);
\draw[soarr, draw=memStroke] (-1.410,-2.700) -- (-1.720,-2.900);
\draw[soarr] (1.090,-1.100) -- (0.780,-1.300);
\draw[soarr, draw=memStroke] (0.780,-3.500) -- (1.090,-3.700);
\draw[soarr] (1.090,-1.100) -- (2.030,-1.300);
\draw[soarr, draw=memStroke] (2.030,-3.500) -- (1.090,-3.700);
\draw[soarr] (0.780,-1.900) -- (0.470,-2.100);
\draw[soarr, draw=memStroke] (0.470,-2.700) -- (0.780,-2.900);
\draw[soarr] (0.780,-1.900) -- (1.410,-2.100);
\draw[soarr, draw=memStroke] (1.410,-2.700) -- (0.780,-2.900);
\draw[densely dotted, black!35] (-0.470,-1.900) -- (-0.470,-2.100);
\draw[densely dotted, black!35] (2.350,-1.900) -- (2.350,-2.100);
\node[font=\small, anchor=east] at (-2.6,0) {切};
\node[font=\small, anchor=east] at (-2.6,-4.8) {合併};
\end{tikzpicture}
\diagcaption{遞迴版合併排序。 黑色箭頭是切開，橘色箭頭是合併，最下面一列是結果。}
\label{fig:so-merge-tree}
\end{center}

考卷上常用的寫法是「每合併一次寫一列」。 遞迴版合併的先後順序是深度優先的：左半全部
做完，才輪到右半。 粗體是這一次合併的範圍。

\begin{center}
\begin{tabular}{l*{6}{c}}
\toprule
合併的範圍 & 1 & 2 & 3 & 4 & 5 & 6 \\
\midrule
開始 & 52 & 27 & 8 & 66 & 41 & 13 \\
$A[1..1] + A[2..2]$ & \textbf{27} & \textbf{52} & 8 & 66 & 41 & 13 \\
$A[1..2] + A[3..3]$ & \textbf{8} & \textbf{27} & \textbf{52} & 66 & 41 & 13 \\
$A[4..4] + A[5..5]$ & 8 & 27 & 52 & \textbf{41} & \textbf{66} & 13 \\
$A[4..5] + A[6..6]$ & 8 & 27 & 52 & \textbf{13} & \textbf{41} & \textbf{66} \\
$A[1..3] + A[4..6]$ & \textbf{8} & \textbf{13} & \textbf{27} & \textbf{41} & \textbf{52} & \textbf{66} \\
\bottomrule
\end{tabular}
\end{center}

\subsection{分析}

設 $C(n)$ 是長度 $n$ 時最多比較幾次。 兩邊遞迴，再加上一次合併：
\[
C(1) = 0, \qquad C(n) \le C(\lceil n/2 \rceil) + C(\lfloor n/2 \rfloor) + n - 1 \quad (n \ge 2).
\]
這和第~2~章的報表加總 $T(n) = T(\lceil n/2 \rceil) + T(\lfloor n/2 \rfloor) + n$ 是同一
個形狀，只是每次少 $1$。 對 $n$ 歸納馬上得到 $C(n) \le T(n)$，而第~2~章定理~\ref{thm:rec-audit} 證明
了 $T(n) \le n \lceil \lg n \rceil$，所以
\[
C(n) \le n \lceil \lg n \rceil .
\]
確切的最壞值是 $n \lceil \lg n \rceil - 2^{\lceil \lg n \rceil} + 1$（觀念題會驗證 $n$
是 $2$ 的次方的情形）。 最好的情況，例如輸入已經排好、$n$ 是 $2$ 的次方，每次合併只比
左半的個數次，共 $\frac{n}{2} \lg n$ 次，仍然是 $\Theta(n \lg n)$。 寫入的次數更固定：遞
迴樹每一層合併的數加起來不超過 $n$ 個，最上面 $\lfloor \lg n \rfloor$ 層每層剛好 $n$ 個，每個數都要抄到 $B$ 再寫回來。 所以合併排序不
管輸入如何，時間都是 $\Theta(n \lg n)$。

空間方面，$B$ 要 $n$ 格，遞迴深度是 $\lceil \lg n \rceil$，合計 $\Theta(n)$ 的額外空間。
這是合併排序最大的缺點。

\subsection{由下而上的版本}

不用遞迴也可以：先把相鄰的一個和一個合併成長度 $2$ 的段，再把相鄰的兩段合併成長度
$4$，依此類推。 每一趟段的長度加倍，$\lceil \lg n \rceil$ 趟就排好了。 最後一段可能不足
長度，合併時右邊界要取 $\min$；右邊根本沒有段時，這一段留著不動。

\begin{pseudo}[label=alg:so-merge-bu]{合併排序（由下而上）}
\Procedure{Merge-Sort-BU}{$A, B, n$}
    \State $w \gets 1$ \Comment{這一趟每一段的長度}
    \While{$w < n$}
        \State $p \gets 1$
        \While{$p + w \le n$} \Comment{右邊還有一段才合併}
            \State \Call{Merge}{$A, B, p, p + w - 1, \min(p + 2w - 1, n)$}
            \State $p \gets p + 2w$
        \EndWhile
        \State $w \gets 2w$
    \EndWhile
\EndProcedure
\end{pseudo}

同一個陣列用由下而上的版本，一趟寫一列：

\begin{center}
\begin{tabular}{l*{6}{c}}
\toprule
這一趟 & 1 & 2 & 3 & 4 & 5 & 6 \\
\midrule
開始 & 52 & 27 & 8 & 66 & 41 & 13 \\
$w = 1$ & 27 & 52 & 8 & 66 & 13 & 41 \\
$w = 2$ & 8 & 27 & 52 & 66 & 13 & 41 \\
$w = 4$ & 8 & 13 & 27 & 41 & 52 & 66 \\
\bottomrule
\end{tabular}
\end{center}

兩個版本最後的結果一樣，中間的陣列卻不同。 遞迴版把 $6$ 個切成 $3 + 3$，再切成
$2 + 1$；由下而上的版本先兩兩合併成 $2 + 2 + 2$，第二趟合併成 $4 + 2$。 只有 $n$ 是
$2$ 的次方時，兩者合併的段才完全相同（只是先後次序不同）。 比較次數也可能不一樣，
這個例子遞迴版比了 9 次，由下而上比了 10 次。

由下而上的版本沒有遞迴，很適合用在串列和外部排序（第~15~章第~\ref{sec:lin-external}~節）：每一趟都是從頭
到尾依序讀、依序寫。

\subsection{串列版}

第~5~章的 \proc{Split-Half}（虛擬碼~\ref{alg:ll-split}）把串列切成前後兩半，\proc{Merge-Lists}
（虛擬碼~\ref{alg:ll-merge}）合併兩條排好的串列，兩個都只改 $\mathit{next}$。 把它們接起來就
是串列的合併排序：

\begin{pseudo}[label=alg:so-merge-list]{串列的合併排序，回傳排好的串列的開頭}
\Function{Merge-Sort-List}{$h$}
    \If{$h = \NIL$ \OR $h.\mathit{next} = \NIL$}
        \State \Return $h$
    \EndIf
    \State $b \gets$ \Call{Split-Half}{$h$} \Comment{前半從 $h$ 開始，後半從 $b$ 開始}
    \State \Return \proc{Merge-Lists}(\proc{Merge-Sort-List}($h$), \proc{Merge-Sort-List}($b$))
\EndFunction
\end{pseudo}

串列版不需要暫存陣列，額外空間只有 $\Theta(\lg n)$ 的遞迴堆疊。 \proc{Split-Half} 讓前
半在前，\proc{Merge-Lists} 遇到相同的鍵先拿前半的，所以也是穩定的。 第~5~章定理~\ref{thm:ll-split} 證明過
\proc{Split-Half} 在 $n \ge 2$ 時兩半都不是空的，遞迴一定會停。 串列不能隨機存取，快
速排序、堆積排序都不好做，合併排序是在串列上排序的標準做法。

實務上常見的改良：子陣列很短時改用插入排序；合併前先檢查 $A[q] \le A[q + 1]$，成立就不
必合併，排好的輸入只要 $\Theta(n)$。

\subsection{C 程式}

\begin{cbookrecap}[第~15~章 15.6 節、第~19~章 19.3 節]
處理陣列的一段時，C 程式習慣用兩個下標 \texttt{lo}、\texttt{hi} 表示範圍，遞迴函式
只處理這一段，陣列本身用指標傳進去。 中間點寫成 \texttt{lo + (hi - lo) / 2} 而不是
\texttt{(lo + hi) / 2}，兩個很大的下標相加可能超過 \texttt{int} 的範圍。

暫存陣列 \texttt{b} 只配置一次，整個排序共用，用完 \texttt{free}。 不要在每次合併時各
\texttt{malloc} 一次，那會多出 $\Theta(n)$ 次配置。 下面的程式用一次 \texttt{malloc} 要
$3n$ 格，再用指標分成原始資料、\texttt{a}、\texttt{b} 三段。 \texttt{memcpy(dst,
src, k)} 把 \texttt{src} 開始的 \texttt{k} 個位元組抄到 \texttt{dst}，兩塊不能重疊，
宣告在 \texttt{<string.h>}。
\end{cbookrecap}

下面的程式讀入一個陣列，用兩個版本各排一次：遞迴版每合併一次就印出整個陣列和合併的
範圍，由下而上的版本每一趟印一次，最後印出兩者的比較次數。 C 的下標從 $0$ 開始，所以
\texttt{a[0..1]} 是虛擬碼的 $A[1..2]$。 比較次數用逗號運算子數：\texttt{(cmp\_count++,
b[i] <= b[j])} 先加一，再比較，和第 6 行一樣只在兩邊都還有時才算。

\begin{example}[label=ex:so-merge]{c}{(merge\_sort.c) 遞迴與由下而上的合併排序}
#include <stdio.h>
#include <stdlib.h>
#include <string.h>

int n;                                  // 陣列長度
long long cmp_count;                    // b[i] <= b[j] 比較了幾次

void print_array(const int *a)
{
    for (int i = 0; i < n; i++) {
        printf(" %d", a[i]);
    }
    printf("\n");
}

// 合併 a[lo..mid] 和 a[mid+1..hi]，b 是暫存陣列。
void merge(int *a, int *b, int lo, int mid, int hi)
{
    memcpy(b + lo, a + lo, (hi - lo + 1) * sizeof *a);
    int i = lo, j = mid + 1;
    for (int k = lo; k <= hi; k++) {
        if (i <= mid && (j > hi || (cmp_count++, b[i] <= b[j]))) {
            a[k] = b[i++];              // 相等時拿左邊的：穩定
        } else {
            a[k] = b[j++];
        }
    }
}

// 遞迴版：每合併一次就印出範圍和整個陣列。
void merge_sort_rec(int *a, int *b, int lo, int hi)
{
    if (lo < hi) {
        int mid = lo + (hi - lo) / 2;
        merge_sort_rec(a, b, lo, mid);
        merge_sort_rec(a, b, mid + 1, hi);
        merge(a, b, lo, mid, hi);
        printf("a[%d..%d]", lo, hi);
        print_array(a);
    }
}

// 由下而上版：寬度 w = 1, 2, 4, ... 的段兩兩合併，每一趟印一次。
void merge_sort_bu(int *a, int *b)
{
    for (int w = 1; w < n; w *= 2) {
        for (int lo = 0; lo + w < n; lo += 2 * w) {
            int hi = lo + 2 * w - 1 < n - 1 ? lo + 2 * w - 1 : n - 1;
            merge(a, b, lo, lo + w - 1, hi);
        }
        printf("w=%d", w);
        print_array(a);
    }
}

int main(void)
{
    if (scanf("%d", &n) != 1 || n <= 0) {
        return 1;
    }
    int *src = malloc(3 * n * sizeof *src);     // 原始資料、a、b 共用一塊
    if (src == NULL) {
        return 1;
    }
    int *a = src + n, *b = src + 2 * n;
    for (int i = 0; i < n; i++) {
        if (scanf("%d", &src[i]) != 1) {
            free(src);
            return 1;
        }
    }
    memcpy(a, src, n * sizeof *a);
    merge_sort_rec(a, b, 0, n - 1);
    long long c1 = cmp_count;
    cmp_count = 0;
    memcpy(a, src, n * sizeof *a);
    merge_sort_bu(a, b);
    printf("comparisons: %lld %lld\n", c1, cmp_count);
    free(src);
    return 0;
}
\end{example}
\iosample{6 \\ 52 27 8 66 41 13}{a[0..1] 27 52 8 66 41 13 \\ a[0..2] 8 27 52 66 41 13 \\ a[3..4] 8 27 52 41 66 13 \\ a[3..5] 8 27 52 13 41 66 \\ a[0..5] 8 13 27 41 52 66 \\ w=1 27 52 8 66 13 41 \\ w=2 8 27 52 66 13 41 \\ w=4 8 13 27 41 52 66 \\ comparisons: 9 10}

輸出和前面兩張表一樣。

\section{快速排序}
\label{sec:so-quick}

合併排序的工作在遞迴\textbf{之後}：兩半排好了，再花力氣合併。 \textbf{快速排序}
（quicksort，C. A. R. Hoare 在 1961 年發表）把工作放在遞迴\textbf{之前}：先選一個數
當\textbf{基準}（pivot），把比它小的放左邊、比它大的放右邊，這一步叫\textbf{分割}
（partition）。 左右兩邊各自遞迴排好以後，整段自然就排好了，不需要合併。 分割可以原
地做，所以快速排序不需要暫存陣列。

分割有兩種常見的寫法，Lomuto 的和 Hoare 的。 兩種寫法回傳的東西不一樣，遞迴的範圍也
不一樣，是考試和寫程式時最容易混淆的地方，這一節兩種都講，並用同一個陣列追蹤。

\subsection{Lomuto 分割}

Lomuto 分割取最後一個數 $x = A[r]$ 當基準，從左到右掃一次。 掃的過程中維持一個分界
$i$：$A[p..i-1]$ 都不大於 $x$，$A[i..j-1]$ 都大於 $x$。 看到 $A[j] \le x$ 時，把它和
$A[i]$ 交換，分界往右一格。 最後把基準換到分界上。

\begin{pseudo}[label=alg:so-lomuto]{Lomuto 分割，回傳基準最後的位置}
\Function{Lomuto-Partition}{$A, p, r$}
    \State $x \gets A[r]$ \Comment{基準}
    \State $i \gets p$
    \For{$j \gets p$ \To $r - 1$}
        \If{$A[j] \le x$}
            \State 交換 $A[i]$ 和 $A[j]$
            \State $i \gets i + 1$
        \EndIf
    \EndFor
    \State 交換 $A[i]$ 和 $A[r]$
    \State \Return $i$
\EndFunction
\end{pseudo}

\begin{center}
\begin{tikzpicture}
\node[soreg, fill=coreFill, minimum width=2.4cm, anchor=west] at (0,0) {$\le x$};
\node[soreg, fill=memFill, minimum width=2.4cm, anchor=west] at (2.4,0) {$> x$};
\node[soreg, minimum width=2.6cm, anchor=west] at (4.8,0) {還沒看};
\node[soreg, fill=optFill, minimum width=0.7cm, anchor=west] at (7.4,0) {$x$};
\node[idxlbl] at (0.2,-0.55) {$p$};
\node[idxlbl] at (2.6,-0.55) {$i$};
\node[idxlbl] at (5.0,-0.55) {$j$};
\node[idxlbl] at (7.75,-0.55) {$r$};
\end{tikzpicture}
\diagcaption{Lomuto 分割的第 4 行每一輪開始時，$A[p..r]$ 分成四段。 中間兩段可能是空的。}
\label{fig:so-lomuto}
\end{center}

用第~1~章第~\ref{sec:ana-invariant}~節的迴圈不變量證明它是對的。 \textbf{不變量}：第 4 行每一輪開始時，
$A[p..i-1]$ 的每個數都 $\le x$，$A[i..j-1]$ 的每個數都 $> x$，而且 $A[r] = x$。

\begin{itemize}
\item \textbf{初始}：$i = j = p$，兩段都是空的；$A[r]$ 還沒動過。
\item \textbf{維持}：如果 $A[j] > x$，什麼都不做，$j$ 加一以後第二段多了這個數，仍然
成立。 如果 $A[j] \le x$，交換 $A[i]$ 和 $A[j]$：$A[i]$ 換成一個 $\le x$ 的數，第一段
多一格；原本的 $A[i]$ 是第二段的第一個數（$> x$），被換到 $A[j]$，正好是第二段新的最
後一格。 （第二段是空的時候 $i = j$，自己和自己交換。） $i$、$j$ 各加一，仍然成立。 $j$ 只
跑到 $r - 1$，$A[r]$ 不會被換到。
\item \textbf{結束}：$j = r$。 這時 $A[p..i-1] \le x$，$A[i..r-1] > x$。 第 8 行把基準換
到 $A[i]$，原本的 $A[i]$（$> x$，或者 $i = r$ 時就是基準自己）換到 $A[r]$，得到
$A[p..i-1] \le A[i] = x < A[i+1..r]$。
\end{itemize}

第 5 行一共比較 $r - p$ 次，所以分割長度 $m$ 的子陣列剛好比 $m - 1$ 次，時間
$\Theta(m)$。 回傳的 $q$ 是基準最後的位置，基準已經就定位，兩邊遞迴時都不必再含它：

\begin{pseudo}[label=alg:so-quick]{快速排序（Lomuto 分割）}
\Procedure{Quicksort}{$A, p, r$}
    \If{$p < r$}
        \State $q \gets$ \Call{Lomuto-Partition}{$A, p, r$}
        \State \Call{Quicksort}{$A, p, q - 1$}
        \State \Call{Quicksort}{$A, q + 1, r$}
    \EndIf
\EndProcedure
\end{pseudo}

\subsection{Lomuto 的追蹤}

用 $A = \langle 47, 82, 15, 63, 29, 91, 36, 54, 40 \rangle$。 第一次分割的基準是 $x = 40$，
下表每一列是第 4 行做完一輪之後的陣列。

\begin{center}
\small
\begin{tabular}{ccc*{9}{c}}
\toprule
$j$ & $A[j] \le 40$？ & 之後的 $i$ & 1 & 2 & 3 & 4 & 5 & 6 & 7 & 8 & 9 \\
\midrule
開始 & & 1 & 47 & 82 & 15 & 63 & 29 & 91 & 36 & 54 & 40 \\
1 & 否 & 1 & 47 & 82 & 15 & 63 & 29 & 91 & 36 & 54 & 40 \\
2 & 否 & 1 & 47 & 82 & 15 & 63 & 29 & 91 & 36 & 54 & 40 \\
3 & 是，交換 $A[1]$、$A[3]$ & 2 & 15 & 82 & 47 & 63 & 29 & 91 & 36 & 54 & 40 \\
4 & 否 & 2 & 15 & 82 & 47 & 63 & 29 & 91 & 36 & 54 & 40 \\
5 & 是，交換 $A[2]$、$A[5]$ & 3 & 15 & 29 & 47 & 63 & 82 & 91 & 36 & 54 & 40 \\
6 & 否 & 3 & 15 & 29 & 47 & 63 & 82 & 91 & 36 & 54 & 40 \\
7 & 是，交換 $A[3]$、$A[7]$ & 4 & 15 & 29 & 36 & 63 & 82 & 91 & 47 & 54 & 40 \\
8 & 否 & 4 & 15 & 29 & 36 & 63 & 82 & 91 & 47 & 54 & 40 \\
結束 & 交換 $A[4]$、$A[9]$ & 4 & 15 & 29 & 36 & \textbf{40} & 82 & 91 & 47 & 54 & 63 \\
\bottomrule
\end{tabular}
\end{center}

整個排序每一次分割（範圍至少有兩個數的呼叫）寫一列，依呼叫的先後排列，粗體是這次歸
位的基準。

\begin{center}
\begin{tabular}{ccc*{9}{c}}
\toprule
範圍 & $x$ & $q$ & 1 & 2 & 3 & 4 & 5 & 6 & 7 & 8 & 9 \\
\midrule
開始 & & & 47 & 82 & 15 & 63 & 29 & 91 & 36 & 54 & 40 \\
$[1..9]$ & 40 & 4 & 15 & 29 & 36 & \textbf{40} & 82 & 91 & 47 & 54 & 63 \\
$[1..3]$ & 36 & 3 & 15 & 29 & \textbf{36} & 40 & 82 & 91 & 47 & 54 & 63 \\
$[1..2]$ & 29 & 2 & 15 & \textbf{29} & 36 & 40 & 82 & 91 & 47 & 54 & 63 \\
$[5..9]$ & 63 & 7 & 15 & 29 & 36 & 40 & 47 & 54 & \textbf{63} & 91 & 82 \\
$[5..6]$ & 54 & 6 & 15 & 29 & 36 & 40 & 47 & \textbf{54} & 63 & 91 & 82 \\
$[8..9]$ & 82 & 8 & 15 & 29 & 36 & 40 & 47 & 54 & 63 & \textbf{82} & 91 \\
\bottomrule
\end{tabular}
\end{center}

範圍 $[4..3]$、$[5..4]$ 這些空的呼叫和只有一個數的呼叫直接返回，表上不寫。 一共比較了
$8 + 2 + 1 + 4 + 1 + 1 = 17$ 次。

\subsection{Hoare 分割}

Hoare 原本的寫法是兩個指標從兩端往中間走：$i$ 從左邊找一個 $\ge x$ 的數，$j$ 從右邊找
一個 $\le x$ 的數，兩個都找到而且還沒交錯，就把它們交換，繼續走；交錯了就停。 下面的
版本取第一個數 $A[p]$ 當基準。

\begin{pseudo}[label=alg:so-hoare]{Hoare 分割，回傳分界 $j$}
\Function{Hoare-Partition}{$A, p, r$}
    \State $x \gets A[p]$ \Comment{基準}
    \State $i \gets p$；$j \gets r$
    \While{\TRUE}
        \While{$A[i] < x$}
            \State $i \gets i + 1$
        \EndWhile
        \While{$A[j] > x$}
            \State $j \gets j - 1$
        \EndWhile
        \If{$i \ge j$}
            \State \Return $j$
        \EndIf
        \State 交換 $A[i]$ 和 $A[j]$
        \State $i \gets i + 1$；$j \gets j - 1$
    \EndWhile
\EndFunction
\end{pseudo}

和 Lomuto 最大的不同在回傳值。 Hoare 回傳的 $j$ 只保證
\[
A[p..j] \text{ 的每個數} \;\le\; x \;\le\; A[j+1..r] \text{ 的每個數},
\]
基準 $x$ 本身不一定在 $A[j]$，可能在左邊，也可能在右邊。 所以兩邊遞迴的範圍是
$[p..j]$ 和 $[j+1..r]$，$A[j]$ 要包含在左邊：

\begin{pseudo}[label=alg:so-quick-hoare]{快速排序（Hoare 分割）}
\Procedure{Quicksort-H}{$A, p, r$}
    \If{$p < r$}
        \State $j \gets$ \Call{Hoare-Partition}{$A, p, r$}
        \State \Call{Quicksort-H}{$A, p, j$}
        \State \Call{Quicksort-H}{$A, j + 1, r$}
    \EndIf
\EndProcedure
\end{pseudo}

這個版本的第 5、7 行迴圈沒有檢查邊界，第 9 行的 $\ge$ 也不能改成 $>$。 它對的理由寫在
下面的框裡。

\begin{whybox}[Hoare 分割為什麼不會越界、一定會變小]
設 $p < r$。 要說明三件事：$i$、$j$ 不會走出 $[p..r]$；回傳的 $j$ 滿足 $p \le j < r$；
回傳時 $A[p..j] \le x \le A[j+1..r]$。

\textbf{第一輪}：$A[p] = x$ 不小於 $x$，第 5 行的迴圈在 $i = p$ 就停。 $j$ 往左走，最
晚在 $A[p] = x$ 停下，所以 $p \le j \le r$。 如果 $j = p$，第 10 行回傳 $p < r$。 否則
交換 $A[p]$ 和 $A[j]$，之後 $A[p] \le x$、$A[j] = x$。

\textbf{之後每一輪開始時}（剛交換完、$i$ 和 $j$ 各走了一步），以下成立：$A[p..i-1]$ 都
$\le x$，$A[j+1..r]$ 都 $\ge x$，而且上一次交換留下了兩個「擋板」：位置 $i - 1$ 的數
$\le x$，位置 $j + 1$ 的數 $\ge x$。 第 5 行的 $i$ 最晚停在擋板 $j + 1$，第 7 行的 $j$
最晚停在擋板 $i - 1$，所以都不會越界。 $j$ 在第一次交換之後至少減過一次，以後只會變
小，所以回傳時 $j \le r - 1$；$j$ 停的位置不會小於 $p$（$A[p] \le x$ 擋著）。

\textbf{停下的時候}：$i$ 走過的數都 $< x$，$j$ 走過的數都 $> x$，加上原本的兩段，得到
$A[p..i-1] \le x$、$A[j+1..r] \ge x$。 $i \ge j$ 時只有兩種情形：$i = j + 1$，這時
$A[p..j] = A[p..i-1] \le x$；或 $i = j$，這時 $A[i] \ge x$（$i$ 停下的原因）而且
$A[j] \le x$，所以 $A[j] = x$，$A[p..j] \le x$ 仍然成立。 （$i > j + 1$ 不可能：那表示
$i$ 走過了 $A[j+1]$，它應該 $< x$，卻又 $\ge x$。）

$p \le j < r$ 保證兩邊的範圍都比 $[p..r]$ 小，遞迴一定會停。 如果基準改取 $A[r]$，第一
輪的 $j$ 會停在 $r$，可能回傳 $r$，右邊是空的、左邊和原來一樣大，就會無窮遞迴。
\end{whybox}

\subsection{Hoare 的追蹤}

同一個陣列 $\langle 47, 82, 15, 63, 29, 91, 36, 54, 40 \rangle$，基準 $x = 47$。 第一次分割
每交換一次寫一列，粗體是被交換的兩格。

\begin{center}
\begin{tabular}{l*{9}{c}}
\toprule
步驟 & 1 & 2 & 3 & 4 & 5 & 6 & 7 & 8 & 9 \\
\midrule
開始，$x = 47$ & 47 & 82 & 15 & 63 & 29 & 91 & 36 & 54 & 40 \\
$i = 1$、$j = 9$，交換 & \textbf{40} & 82 & 15 & 63 & 29 & 91 & 36 & 54 & \textbf{47} \\
$i = 2$、$j = 7$，交換 & 40 & \textbf{36} & 15 & 63 & 29 & 91 & \textbf{82} & 54 & 47 \\
$i = 4$、$j = 5$，交換 & 40 & 36 & 15 & \textbf{29} & \textbf{63} & 91 & 82 & 54 & 47 \\
$i = 5$、$j = 4$，回傳 4 & 40 & 36 & 15 & 29 & 63 & 91 & 82 & 54 & 47 \\
\bottomrule
\end{tabular}
\end{center}

第一次交換時 $i = 1$ 停在基準自己，$j = 9$ 停在 $40$。 最後 $i$ 停在 $A[5] = 63$，$j$
停在 $A[4] = 29$，已經交錯，回傳 $4$。 基準 $47$ 這時在 $A[9]$，不在分界上。 整個排序
的每一次分割：

\begin{center}
\begin{tabular}{ccc*{9}{c}}
\toprule
範圍 & $x$ & $j$ & 1 & 2 & 3 & 4 & 5 & 6 & 7 & 8 & 9 \\
\midrule
開始 & & & 47 & 82 & 15 & 63 & 29 & 91 & 36 & 54 & 40 \\
$[1..9]$ & 47 & 4 & 40 & 36 & 15 & 29 & 63 & 91 & 82 & 54 & 47 \\
$[1..4]$ & 40 & 3 & 29 & 36 & 15 & 40 & 63 & 91 & 82 & 54 & 47 \\
$[1..3]$ & 29 & 1 & 15 & 36 & 29 & 40 & 63 & 91 & 82 & 54 & 47 \\
$[2..3]$ & 36 & 2 & 15 & 29 & 36 & 40 & 63 & 91 & 82 & 54 & 47 \\
$[5..9]$ & 63 & 6 & 15 & 29 & 36 & 40 & 47 & 54 & 82 & 91 & 63 \\
$[5..6]$ & 47 & 5 & 15 & 29 & 36 & 40 & 47 & 54 & 82 & 91 & 63 \\
$[7..9]$ & 82 & 7 & 15 & 29 & 36 & 40 & 47 & 54 & 63 & 91 & 82 \\
$[8..9]$ & 91 & 8 & 15 & 29 & 36 & 40 & 47 & 54 & 63 & 82 & 91 \\
\bottomrule
\end{tabular}
\end{center}

和 Lomuto 的表對照：Lomuto 做了 6 次分割，每次都讓一個基準就定位；Hoare 做了 8 次，
基準常常留在子陣列裡，要到範圍縮得很小才就定位。 Hoare 的長處是交換少：它只交換
「左邊太大、右邊太小」的一對，Lomuto 則是每看到一個 $\le x$ 的數就交換一次。 這個例子
Lomuto 執行了 $15$ 次交換（其中 $7$ 次是自己和自己換），Hoare 只有 $10$ 次；換成
$n = 2000$ 的隨機排列，Hoare 的交換次數平均大約只有 Lomuto 的四成。

\begin{examnote}
兩種分割的下標規則不一樣，不要混用。 Lomuto 回傳的 $q$ 是基準最後的位置，$A[q]$ 已經就
定位，遞迴範圍是 $[p..q-1]$ 和 $[q+1..r]$。 Hoare 回傳的 $j$ 只是分界，遞迴範圍是
$[p..j]$ 和 $[j+1..r]$；寫成 $[p..j-1]$ 會漏掉 $A[j]$，結果可能沒排好。 Hoare 的基準取
$A[r]$ 時可能回傳 $r$ 而無窮遞迴，$\langle 1, 2 \rangle$ 就是例子。 另外有些課本取
$A[p]$ 當基準，兩邊往中間找完以後再把 $A[p]$ 換到 $A[j]$，這種版本的 $j$ 就是基準的位
置，遞迴範圍和 Lomuto 一樣。 作答前先看清楚題目用哪一種寫法、陣列從 $0$ 還是從 $1$ 開
始，再寫追蹤。
\end{examnote}

\subsection{最壞和最好的情況}

一次分割的時間和子陣列長度成正比，所以總時間取決於每次切得平不平均。

\textbf{最壞的情況}：每次都切成 $0$ 和 $m - 1$ 個。 這時遞迴式是 $T(n) = T(n-1) +
\Theta(n)$，照第~2~章第~\ref{sec:rec-other}~節的展開，$T(n) = \Theta(n^2)$。 用 Lomuto 分割、基準取最後一
個數時，已經排好的輸入每次的基準都是最大值，比較次數剛好是 $(n-1) + (n-2) + \cdots +
1 = n(n-1)/2$；由大到小的輸入每次的基準都是最小值，一樣是 $n(n-1)/2$。 這兩種輸入在實
務上很常見，所以「固定取最後一個」不是好的選法。

\textbf{最好的情況}：每次都切成兩半，$T(n) = 2T(n/2) + \Theta(n) = \Theta(n \lg n)$。 其實
不必剛好一半：第~2~章第~\ref{sec:rec-tree}~節的不平衡遞迴樹說明，只要每次都至少切出固定的比例（例如
$1:9$），遞迴樹的高度就是 $O(\lg n)$，每一層的工作不超過 $n$，總時間仍然是
$\Theta(n \lg n)$，只是常數變大。 快速排序怕的是「幾乎每次都切得極偏」，偶爾切偏一次
沒有關係。

\subsection{隨機選基準}

要避免固定的壞輸入，可以每次從子陣列裡隨機選一個數當基準。 這時沒有任何一個輸入固定
是壞的，執行時間是隨機的，我們算它的期望值。

\begin{pseudo}[label=alg:so-rand]{隨機化快速排序}
\Function{Randomized-Partition}{$A, p, r$}
    \State $k \gets$ \Call{Random}{$p, r$} \Comment{從 $p..r$ 均勻選一個整數}
    \State 交換 $A[k]$ 和 $A[r]$
    \State \Return \Call{Lomuto-Partition}{$A, p, r$}
\EndFunction
\Statex
\Procedure{Randomized-Quicksort}{$A, p, r$}
    \If{$p < r$}
        \State $q \gets$ \Call{Randomized-Partition}{$A, p, r$}
        \State \Call{Randomized-Quicksort}{$A, p, q - 1$}
        \State \Call{Randomized-Quicksort}{$A, q + 1, r$}
    \EndIf
\EndProcedure
\end{pseudo}

以下假設鍵互不相同。 把它們由小到大叫做 $z_1 < z_2 < \cdots < z_n$，並令
$Z_{ij} = \{z_i, z_{i+1}, \ldots, z_j\}$。 令 $X$ 是整個排序裡 \proc{Lomuto-Partition} 第 5 行
比較的總次數。

\begin{lemma}
\label{lem:so-once}
每一次比較都是「這次分割的基準」和「範圍裡的另一個數」比；任兩個數在整個排序中最多
比較一次。
\end{lemma}

\begin{proof}
第 5 行比較的是 $A[j]$ 和 $x$，$x$ 是這次的基準。 分割結束後基準在 $A[q]$，而兩個遞迴
呼叫的範圍是 $[p..q-1]$ 和 $[q+1..r]$，都不含 $q$，所以這個數以後不會再出現在任何一次
分割裡。 兩個數 $a$、$b$ 比較時，其中一個是基準；比完之後那個數就不再參加，所以它們
不會再比第二次。
\end{proof}

\begin{lemma}
\label{lem:so-pair}
$i < j$ 時，$z_i$ 和 $z_j$ 會被比較，若且唯若 $Z_{ij}$ 裡第一個被選為基準的是 $z_i$ 或
$z_j$。
\end{lemma}

\begin{proof}
先說明：在 $Z_{ij}$ 的任何一個數被選為基準之前，$Z_{ij}$ 的數全部在同一次呼叫的範圍
裡。 一開始全部在 $[1..n]$。 一次分割如果選了 $Z_{ij}$ 以外的基準 $y$，$y$ 不是比
$z_i$ 小就是比 $z_j$ 大，$Z_{ij}$ 的數和 $y$ 比的結果都一樣，所以全部分到同一邊。

考慮 $Z_{ij}$ 裡第一個被選為基準的數 $z_k$，這時 $z_i$、$z_j$ 都在這次的範圍裡。 若
$k = i$ 或 $k = j$，這次分割會拿基準和範圍裡每一個數比，其中就包括另一個。 若
$i < k < j$，分割之後 $z_i$ 到左邊、$z_j$ 到右邊，從此在不同的子陣列裡，不會再相比；
在這之前，它們都還沒當過基準，而由引理~\ref{lem:so-once}，每次比較都有一方是基準，所以
之前也沒有比過。
\end{proof}

\begin{theorem}
\label{thm:so-quick-avg}
鍵互不相同時，對任何輸入，\proc{Randomized-Quicksort} 的比較次數期望值是
\[
\mathrm{E}[X] = \sum_{i=1}^{n-1} \sum_{j=i+1}^{n} \frac{2}{j - i + 1} = 2(n+1)H_n - 4n
\;\le\; 2n \ln n,
\]
其中 $H_n = 1 + \frac12 + \cdots + \frac1n$。 期望執行時間是 $O(n \lg n)$。
\end{theorem}

\begin{proof}
對每一對 $i < j$，令指示變數（第~12~章第~\ref{sec:hash-chain-analysis}~節）$X_{ij}$ 在 $z_i$、$z_j$ 被比較時為 $1$，否則為 $0$。
由引理~\ref{lem:so-once}，每一對最多比較一次，所以 $X = \sum_{i<j} X_{ij}$。 用期望值的線
性（第~1~章第~\ref{sec:ana-cases}~節算插入排序平均情況時用過），$\mathrm{E}[X] = \sum_{i<j} \Pr[z_i
\text{ 和 } z_j \text{ 被比較}]$。

考慮 $Z_{ij}$ 裡第一次有數被選為基準的那次分割。 那時 $Z_{ij}$ 的 $j - i + 1$ 個數都在
這次的範圍裡，而基準是在範圍裡均勻選出的；在「選中的數屬於 $Z_{ij}$」的條件下，這
$j - i + 1$ 個數被選中的機會都一樣。 由引理~\ref{lem:so-pair}，
\[
\Pr[z_i \text{ 和 } z_j \text{ 被比較}] = \Pr[z_i \text{ 或 } z_j \text{ 先被選中}] = \frac{2}{j - i + 1}.
\]
加總時令 $d = j - i$，距離為 $d$ 的數對有 $n - d$ 對；再令 $m = d + 1$：
\[
\mathrm{E}[X] = \sum_{d=1}^{n-1} \frac{2(n - d)}{d + 1}
= \sum_{m=2}^{n} \frac{2(n + 1 - m)}{m}
= 2(n+1)(H_n - 1) - 2(n - 1)
= 2(n+1)H_n - 4n .
\]
用第~1~章的 $H_n \le \ln n + 1$，得到 $\mathrm{E}[X] \le 2(n+1)\ln n - 2n + 2 \le 2n\ln n$
（最後一步用 $\ln n \le n - 1$）。

執行時間：每一次 $p < r$ 的呼叫會讓一個數當基準，而由引理~\ref{lem:so-once} 這個數不會再
當基準，所以這種呼叫最多 $n$ 次，每次除了第 5 行的比較以外只花常數時間；$p \ge r$ 的呼
叫都是它們的子呼叫，最多 $2n + 1$ 次，每次常數時間。 所以總時間是 $O(n + X)$，期望值是
$O(n + n \lg n) = O(n \lg n)$。
\end{proof}

$2n \ln n \approx 1.39\, n \lg n$，只比合併排序最壞情況的 $n \lg n$ 多大約四成；而快速
排序的內層迴圈很簡單，又不用暫存陣列，所以實際上常常比較快。

兩點補充。 第一，這個期望值是對演算法自己的隨機選擇取的，對\textbf{每一個}輸入都成立。
如果基準固定取最後一個，改成假設「輸入的 $n!$ 種排列機會一樣」，平均比較次數也是同一個
數：這時最後一個數在子陣列裡的名次是均勻的，而且 Lomuto 分割之後兩邊各自仍是隨機排列
（這裡不證）。 所以「快速排序的平均情況是 $\Theta(n \lg n)$」有這兩種意義，前者不怕特
定的壞輸入。

第二，第~9~章第~\ref{sec:bst-height}~節的框算過：$n$ 個鍵隨機順序建成的二元搜尋樹，所有節點的期望深度總
和也是 $2(n+1)H_n - 4n$。 兩者相同不是巧合：依某個順序選基準的快速排序，比較過的數對
正好是依同一個順序建 BST 時「祖先和子孫」的數對，引理~\ref{lem:so-pair} 和第~9~章那個
引理說的是同一件事。

\subsection{三數取中、相同的鍵、堆疊深度}

\textbf{三數取中}（median-of-three）：取 $A[p]$、$A[\lfloor (p+r)/2 \rfloor]$、$A[r]$ 三個
數的中位數當基準，換到 $A[r]$ 再用 Lomuto 分割。 多花三次比較，但排好和由大到小的輸入
會被切得很平均。 它沒有改變最壞情況：仍然有特別構造的輸入讓每次都切得很偏，所以最壞仍
是 $\Theta(n^2)$。 實務上的排序函式常常把三數取中、小範圍改用插入排序合在一起用。

\textbf{相同的鍵}：所有鍵都相同時，Lomuto 分割的第 5 行每次都成立，基準被換到最右邊，
切成 $m - 1$ 和 $0$，又是 $\Theta(n^2)$。 隨機選基準、三數取中都救不了，因為不管選誰都
一樣。 Hoare 分割遇到等於 $x$ 的數，$i$ 和 $j$ 都會停下來交換，兩個指標在中間相遇，切得
很平均。 另一種做法是\textbf{三路分割}：分成 $< x$、$= x$、$> x$ 三段，只對頭尾兩段遞
迴，鍵全部相同時一趟就結束。

\textbf{堆疊深度}：最壞情況下遞迴深度是 $n$，$n$ 很大時會用光呼叫堆疊。 常見的改法是分割
之後只對\textbf{比較短}的一邊遞迴，比較長的一邊改用迴圈處理（把 $p$ 或 $r$ 改掉，回到迴圈開
頭再分割一次）。 每往下遞迴一層，範圍的長度至少減半，所以同時在堆疊上的呼叫不超過
$\lfloor \lg n \rfloor + 1$ 層，總時間不變。

快速排序不穩定。 $\langle 2_a, 2_b, 1 \rangle$ 第一次分割的基準是 $1$，$2_a$、$2_b$ 都比
它大，最後一步把 $1$ 和 $2_a$ 交換，得到 $\langle 1, 2_b, 2_a \rangle$。 Lomuto 和 Hoare
都會做遠距離的交換。

\subsection{實際數一次}

兩種分割的 C 程式在第~\ref{sec:so-compare}~節的範例~\ref{ex:so-stability}。 把它們加上三
數取中和隨機基準，對 $n = 2000$ 的四種輸入各排一次，數鍵的比較次數（三數取中的那三次、
Hoare 第 5、7 行的每一次都算），結果如下。 作為對照，定理~\ref{thm:so-quick-avg} 的期望值
$2(n+1)H_n - 4n$ 在 $n = 2000$ 時大約是 $24{,}730$。

\begin{center}
\begin{tabular}{lrrrr}
\toprule
輸入 & Lomuto，取最後 & Lomuto，三數取中 & Lomuto，隨機 & Hoare，取第一個 \\
\midrule
排好 & 1{,}999{,}000 & 21{,}033 & 23{,}429 & 2{,}002{,}998 \\
由大到小 & 1{,}999{,}000 & 36{,}967 & 24{,}156 & 2{,}003{,}998 \\
隨機 & 25{,}079 & 24{,}252 & 24{,}796 & 34{,}332 \\
全部相同 & 1{,}999{,}000 & 2{,}004{,}997 & 1{,}999{,}000 & 25{,}726 \\
\bottomrule
\end{tabular}
\end{center}

\begin{itemize}
\item 基準取最後一個，在排好和由大到小的輸入上比較次數正好是 $n(n-1)/2 = 1{,}999{,}000$；
Hoare 基準取第一個，在這兩種輸入上一樣退化。
\item 三數取中和隨機基準在前三種輸入都只比了兩萬到四萬次，和期望值同一級。
\item 全部相同時，三種 Lomuto 都退化成 $n(n-1)/2$ 左右，只有 Hoare 不受影響。
\item 隨機輸入上，Hoare 的比較次數比較多：基準沒有歸位，還留在子陣列裡繼續比，兩個指標
交錯時也會多比一兩次。 它的優點在交換少，以及不怕相同的鍵。
\end{itemize}

\section{堆積排序}
\label{sec:so-heap}

第~8~章第~\ref{sec:hp-build}~節最後已經說過堆積排序的想法：先用 \proc{Build-Max-Heap}（虛擬碼~\ref{alg:hp-build}）
把陣列變成最大堆積，最大值就在 $A[1]$；把它和最後一格交換，堆積少一格，再對新的根做
\proc{Sift-Down}（虛擬碼~\ref{alg:hp-sift}）。 每做一次，就有一個最大值放到陣列尾端的正確位
置。 這一節補上完整的虛擬碼、追蹤和分析。

\begin{pseudo}[label=alg:so-heapsort]{堆積排序}
\Procedure{Heapsort}{$A, n$}
    \State \Call{Build-Max-Heap}{$A, n$} \Comment{之後 $A.\mathit{size} = n$}
    \For{$i \gets n$ \DownTo $2$}
        \State 交換 $A[1]$ 和 $A[i]$ \Comment{目前的最大值放到 $A[i]$}
        \State $A.\mathit{size} \gets A.\mathit{size} - 1$
        \State \Call{Sift-Down}{$A, 1$}
    \EndFor
\EndProcedure
\end{pseudo}

\textbf{迴圈不變量}：第 3 行每一輪開始時，$A[1..i]$ 是一個最大堆積，裝著最小的 $i$ 個
數；$A[i+1..n]$ 是最大的 $n - i$ 個數，已經排好。 第 4 行把堆積裡最大的數換到 $A[i]$，
它不大於 $A[i+1..n]$ 的任何一個，所以 $A[i..n]$ 排好了；第 5、6 行把 $A[1..i-1]$ 修回
堆積。 迴圈結束時 $i = 1$，$A[1]$ 是最小的數，$A[2..n]$ 排好，整個陣列就排好了。

第 5 行很重要：堆積的大小要先減一，下沉時才不會把已經排好的 $A[i]$ 又拉回堆積。

\subsection{手算追蹤}

接著第~8~章第~\ref{sec:hp-build}~節的例子，那裡由下而上建好的堆積是
\[
\langle 70, 66, 52, 37, 61, 28, 45, 9, 23, 14 \rangle .
\]
下表每一列是第 3 行的迴圈
做完一輪之後的陣列，「交換」是那一次下沉交換了幾次，灰色的是已經排好、不在堆積裡的
格子。 下沉時兩個子節點相等就換左邊，和第~8~章一樣。

\begin{center}
\begin{tabular}{lc*{10}{c}}
\toprule
& 交換 & 1 & 2 & 3 & 4 & 5 & 6 & 7 & 8 & 9 & 10 \\
\midrule
建好堆積 & & 70 & 66 & 52 & 37 & 61 & 28 & 45 & 9 & 23 & 14 \\
取出 70 & 2 & 66 & 61 & 52 & 37 & 14 & 28 & 45 & 9 & 23 & \sofin{70} \\
取出 66 & 2 & 61 & 37 & 52 & 23 & 14 & 28 & 45 & 9 & \sofin{66} & \sofin{70} \\
取出 61 & 2 & 52 & 37 & 45 & 23 & 14 & 28 & 9 & \sofin{61} & \sofin{66} & \sofin{70} \\
取出 52 & 2 & 45 & 37 & 28 & 23 & 14 & 9 & \sofin{52} & \sofin{61} & \sofin{66} & \sofin{70} \\
取出 45 & 2 & 37 & 23 & 28 & 9 & 14 & \sofin{45} & \sofin{52} & \sofin{61} & \sofin{66} & \sofin{70} \\
取出 37 & 1 & 28 & 23 & 14 & 9 & \sofin{37} & \sofin{45} & \sofin{52} & \sofin{61} & \sofin{66} & \sofin{70} \\
取出 28 & 1 & 23 & 9 & 14 & \sofin{28} & \sofin{37} & \sofin{45} & \sofin{52} & \sofin{61} & \sofin{66} & \sofin{70} \\
取出 23 & 0 & 14 & 9 & \sofin{23} & \sofin{28} & \sofin{37} & \sofin{45} & \sofin{52} & \sofin{61} & \sofin{66} & \sofin{70} \\
取出 14 & 0 & 9 & \sofin{14} & \sofin{23} & \sofin{28} & \sofin{37} & \sofin{45} & \sofin{52} & \sofin{61} & \sofin{66} & \sofin{70} \\
\bottomrule
\end{tabular}
\end{center}

\subsection{分析}

\proc{Build-Max-Heap} 由第~8~章定理~\ref{thm:hp-build} 和它後面的說明，交換少於 $2n$ 次、比較少於 $5n$ 次。 第 3 行的
迴圈做 $n - 1$ 次下沉，堆積大小是 $m$ 時，下沉最多走 $\lfloor \lg m \rfloor$ 層，每層比較
兩次。 所以比較次數最多
\[
5n + \sum_{m=1}^{n-1} 2\lfloor \lg m \rfloor \;\le\; 5n + 2 \lg\bigl((n-1)!\bigr) \;\le\; 2n \lg n + 5n .
\]
堆積排序的最壞情況是 $O(n \lg n)$，額外空間只有幾個變數，下沉也不需要遞迴：它是本章唯
一同時做到「最壞 $O(n \lg n)$」和「原地」的排序。 鍵互不相同時，它的最好情況也是
$\Theta(n \lg n)$（證明比較長，見延伸閱讀）；鍵全部相同時，每次下沉第一步就停，只要
$\Theta(n)$。

堆積排序不穩定。 $\langle 5_a, 5_b, 3 \rangle$ 本來就是最大堆積；第一輪把 $5_a$ 換到
$A[3]$，$3$ 下沉時和 $5_b$ 交換；第二輪把 $5_b$ 換到 $A[2]$，結果是 $\langle 3, 5_b, 5_a
\rangle$。 把最大值換到尾端這個動作，會讓先出現的 $5_a$ 跑到後面。

實務上，堆積排序比快速排序慢一些：下沉時每次跳到兩倍遠的位置，快取的效果差，比較次數
大約 $2n \lg n$，也比快速排序平均的 $1.39\, n \lg n$ 多。 C++ 標準程式庫常用的 introsort
先用快速排序，遞迴太深時改用堆積排序，兼顧平均的速度和最壞的保證。

C 程式在範例~\ref{ex:so-stability}，\texttt{sift\_down} 就是第~8~章的下沉。

\section{各種排序的比較}
\label{sec:so-compare}

\subsection{比較表}

下表整理本章的排序，時間都指鍵的比較次數的級數（也就是執行時間的級數），「平均」是
$n!$ 種排列機會一樣時的期望值，快速排序用隨機基準時也是同一個值。 額外空間包括遞迴堆疊。

\begin{center}
\small
\begin{tabular}{lllllll}
\toprule
排序 & 最好 & 平均 & 最壞 & 額外空間 & 穩定 & 原地 \\
\midrule
氣泡（提早結束） & $\Theta(n)$ & $\Theta(n^2)$ & $\Theta(n^2)$ & $\Theta(1)$ & 是 & 是 \\
選擇 & $\Theta(n^2)$ & $\Theta(n^2)$ & $\Theta(n^2)$ & $\Theta(1)$ & 否 & 是 \\
插入 & $\Theta(n)$ & $\Theta(n^2)$ & $\Theta(n^2)$ & $\Theta(1)$ & 是 & 是 \\
Shell（減半序列） & $\Theta(n \lg n)$ & 見第~\ref{sec:so-shell}~節 & $\Theta(n^2)$ & $\Theta(1)$ & 否 & 是 \\
合併（陣列） & $\Theta(n \lg n)$ & $\Theta(n \lg n)$ & $\Theta(n \lg n)$ & $\Theta(n)$ & 是 & 否 \\
合併（串列） & $\Theta(n \lg n)$ & $\Theta(n \lg n)$ & $\Theta(n \lg n)$ & $\Theta(\lg n)$ & 是 & （不適用） \\
快速 & $\Theta(n \lg n)$ & $\Theta(n \lg n)$ & $\Theta(n^2)$ & $\Theta(\lg n)$～$\Theta(n)$ & 否 & 是 \\
堆積 & $\Theta(n \lg n)$ & $\Theta(n \lg n)$ & $\Theta(n \lg n)$ & $\Theta(1)$ & 否 & 是 \\
\bottomrule
\end{tabular}
\end{center}

幾個要搭配表格記住的條件：

\begin{itemize}
\item 氣泡排序的最好情況 $\Theta(n)$ 要有「沒有交換就結束」那一行；沒有的話三種情況都
是 $\Theta(n^2)$。
\item Shell 排序的最好情況：已經排好時每一趟 $h$ 只比 $n - h$ 次，共 $\Theta(\lg n)$ 趟，
所以是 $\Theta(n \lg n)$ 而不是 $\Theta(n)$。 換成 Knuth 的序列，最壞情況是
$O(n^{3/2})$。
\item 快速排序的空間：遞迴深度最壞是 $n$；用第~\ref{sec:so-quick}~節先遞迴短邊的寫法，
保證是 $O(\lg n)$。 表上的最好、平均是鍵互不相同時的結果，Lomuto 分割遇到大量相同的鍵
會退化。
\item 堆積排序的最好情況 $\Theta(n \lg n)$ 指鍵互不相同；鍵全部相同時是 $\Theta(n)$。
\item 第~15~章會證明任何比較排序最壞都要 $\Omega(n \lg n)$ 次比較，所以合併和堆積排序在
漸近意義下已經最好；不靠比較的計數、基數排序也在第~15~章。
\end{itemize}

怎麼選？ 一般的資料交給標準程式庫；C 標準沒有規定 \texttt{qsort} 用哪
一種演算法，也不保證穩定，需要穩定時把原本的位置加進比較函式當第二個鍵。 資料很少或幾乎
排好，用插入排序；需要穩定或資料在串列上，用合併排序；記憶體很緊又要最壞情況的保證，用
堆積排序；平均最快的通常是快速排序，記得隨機選基準或三數取中。 鍵是範圍不大的整數，或
資料大到放不進記憶體，看第~15~章。

\begin{examnote}
判斷穩定性時，要說「穩定」得說明每一步都不會讓相同的鍵越過對方，例如「只在 $>$ 時交換
相鄰的兩個」；要說「不穩定」只要舉一個反例，用兩個相同的鍵加上下標（$5_a$、$5_b$），寫
出排序後它們的先後。 同一種排序改一個符號，穩定性就可能改變：合併時 $\le$ 改成 $<$、
插入排序的 $>$ 改成 $\ge$ 都會變得不穩定，所以有程式碼的題目要照題目給的寫法判斷。 「原
地」也一樣，題目有定義就照題目的定義。
\end{examnote}

\subsection{實際檢查穩定性}

章首的考卷問題可以直接跑一次。 下面的程式把每筆紀錄存成（分數，座號），輸入時座號由小
到大，用選擇、插入、Shell、快速（Lomuto 和 Hoare）、堆積排序各排一次，印出排好的「分
數:座號」，並檢查同分的座號是不是仍然由小到大。 它也是這幾種排序的 C 實作（合併排序在範
例~\ref{ex:so-merge}），寫法和本章的虛擬碼相同：插入排序就是間隔 $1$ 的 Shell 排序，堆積
排序用第~8~章的下沉；陣列從 $0$ 開始，所以建堆積從 \texttt{n / 2 - 1} 做到 \texttt{0}。 氣
泡排序只是兩層迴圈加一個旗標，請照虛擬碼~\ref{alg:so-bubble} 自己加進去。

\begin{cbookrecap}[第~9~章 9.6、9.9 節]
\texttt{a + k} 是指向 \texttt{a[k]} 的指標，把它當成陣列傳給函式，函式裡的第 $0$ 格就是
\texttt{a[k]}。 快速排序用它遞迴處理右邊那一段，函式只需要長度一個參數。

\texttt{void (*sorter[])(Rec *, int)} 是一個陣列，每一格指向一個「參數是 \texttt{Rec *} 和
\texttt{int}、沒有回傳值」的函式。 函式名稱可以直接放進去，\texttt{sorter[k](a, n)} 就呼叫
第 \texttt{k} 個，同一段測試程式可以輪流套用每一種排序。
\end{cbookrecap}

\begin{example}[label=ex:so-stability]{c}{(sorts.c) 六種排序對同分紀錄的處理}
#include <stdio.h>
#include <string.h>

#define MAXN 100

typedef struct {
    int score;                          // 排序的鍵
    int id;                             // 座號，輸入時由小到大
} Rec;

void swap(Rec *x, Rec *y)
{
    Rec t = *x;
    *x = *y;
    *y = t;
}

void selection_sort(Rec *a, int n)
{
    for (int i = 0; i < n - 1; i++) {
        int m = i;
        for (int j = i + 1; j < n; j++) {
            if (a[j].score < a[m].score) {
                m = j;
            }
        }
        swap(&a[i], &a[m]);
    }
}

void gap_insertion(Rec *a, int n, int h)    // 間隔 h 的插入排序
{
    for (int i = h; i < n; i++) {
        Rec key = a[i];
        int j = i - h;
        while (j >= 0 && a[j].score > key.score) {
            a[j + h] = a[j];
            j -= h;
        }
        a[j + h] = key;
    }
}

void insertion_sort(Rec *a, int n)
{
    gap_insertion(a, n, 1);
}

void shell_sort(Rec *a, int n)          // 增量 n/2, n/4, ..., 1
{
    for (int h = n / 2; h >= 1; h /= 2) {
        gap_insertion(a, n, h);
    }
}

// 快速排序排 a[0..n-1]。遞迴時用指標算術：a + k 是從 a[k] 開始的子陣列。
void quick_lomuto(Rec *a, int n)        // 基準是 a[n-1]
{
    if (n > 1) {
        int i = 0;
        for (int j = 0; j < n - 1; j++) {
            if (a[j].score <= a[n - 1].score) {
                swap(&a[i++], &a[j]);
            }
        }
        swap(&a[i], &a[n - 1]);         // 基準歸位在 a[i]
        quick_lomuto(a, i);             // a[0..i-1]
        quick_lomuto(a + i + 1, n - i - 1);
    }
}

void quick_hoare(Rec *a, int n)         // 基準是 a[0]
{
    if (n > 1) {
        int x = a[0].score, i = 0, j = n - 1;
        for (;;) {
            while (a[i].score < x) {
                i++;
            }
            while (a[j].score > x) {
                j--;
            }
            if (i >= j) {
                break;
            }
            swap(&a[i++], &a[j--]);
        }
        quick_hoare(a, j + 1);          // a[0..j]，包含 a[j]
        quick_hoare(a + j + 1, n - j - 1);
    }
}

void sift_down(Rec *a, int size, int i) // 第 8 章的下沉
{
    for (int j = 2 * i + 1; j < size; i = j, j = 2 * i + 1) {
        if (j + 1 < size && a[j + 1].score > a[j].score) {
            j++;
        }
        if (a[i].score >= a[j].score) {
            return;
        }
        swap(&a[i], &a[j]);
    }
}

void heap_sort(Rec *a, int n)
{
    for (int i = n / 2 - 1; i >= 0; i--) {
        sift_down(a, n, i);
    }
    for (int size = n - 1; size >= 1; size--) {
        swap(&a[0], &a[size]);          // 最大值換到尾端，堆積少一格
        sift_down(a, size, 0);
    }
}

int main(void)
{
    const char *name[] = {"selection", "insertion", "shell",
                          "lomuto", "hoare", "heap"};
    void (*sorter[])(Rec *, int) = {selection_sort, insertion_sort,
                                    shell_sort, quick_lomuto, quick_hoare, heap_sort};
    Rec src[MAXN], a[MAXN];
    int n;
    if (scanf("%d", &n) != 1 || n <= 0 || n > MAXN) {
        return 1;
    }
    for (int i = 0; i < n; i++) {
        src[i].id = i + 1;
        if (scanf("%d", &src[i].score) != 1) {
            return 1;
        }
    }
    for (int k = 0; k < 6; k++) {
        memcpy(a, src, n * sizeof *a);
        sorter[k](a, n);
        int same = 1;                   // 同分的座號是否仍由小到大
        printf("%s:", name[k]);
        for (int i = 0; i < n; i++) {
            printf(" %d:%d", a[i].score, a[i].id);
            if (i > 0 && a[i].score == a[i - 1].score && a[i].id < a[i - 1].id) {
                same = 0;
            }
        }
        printf(" %s\n", same ? "same-order" : "CHANGED");
    }
    return 0;
}
\end{example}
\iosample{8 \\ 90 70 80 80 70 80 70 70}{selection: 70:2 70:5 70:7 70:8 80:6 80:3 80:4 90:1 CHANGED \\ insertion: 70:2 70:5 70:7 70:8 80:3 80:4 80:6 90:1 same-order \\ shell: 70:5 70:2 70:7 70:8 80:3 80:6 80:4 90:1 CHANGED \\ lomuto: 70:2 70:5 70:7 70:8 80:6 80:3 80:4 90:1 CHANGED \\ hoare: 70:5 70:7 70:2 70:8 80:6 80:3 80:4 90:1 CHANGED \\ heap: 70:8 70:7 70:2 70:5 80:6 80:3 80:4 90:1 CHANGED}

插入排序保住了座號的順序，其他五種各自打亂了某些同分的座號，章首的第二種做法（選擇排
序）就在其中。 要注意
\texttt{same-order} 只表示這一組輸入沒有出問題，不能當成穩定的證明；\texttt{CHANGED}
則是一個實際的反例，足以說明不穩定。

\subsection{例題：同一組資料的逐趟追蹤}

考卷上常見的題型是：「用某某排序把下列資料由小到大排好，寫出每一趟（每一輪、每次分
割）之後的結果。」 下面用 $A = \langle 36, 71, 19, 54, 8, 92, 43, 27 \rangle$ 把本章每一種
排序都跑一次，每一趟一列。 這些列是用照本章虛擬碼寫的模擬程式產生的，可以拿來對答案。
「一趟」的定義：氣泡、選擇、插入是外層迴圈的一個 $i$（氣泡排序第 6 趟沒有交換，提早結束）；
Shell 是一個增量（$n = 8$，減半序列是 $4, 2, 1$）；合併排序用由下而上的版本，一趟一列，
$n = 8$ 是 $2$ 的次方，遞迴版每一層合併完的結果和這三列相同；快速排序是一次範圍至少兩個數
的分割，依呼叫的先後；堆積排序是建好堆積一列，之後每取出一個最大值一列。

\begin{center}
\footnotesize
\begin{tabular}{ll*{8}{c}}
\toprule
& & 1 & 2 & 3 & 4 & 5 & 6 & 7 & 8 \\
\midrule
原始陣列 & & 36 & 71 & 19 & 54 & 8 & 92 & 43 & 27 \\
\midrule
\multicolumn{10}{l}{\textbf{氣泡排序}（第二欄：這一趟交換幾次）} \\
第 1 趟 & 5 & 36 & 19 & 54 & 8 & 71 & 43 & 27 & 92 \\
第 2 趟 & 4 & 19 & 36 & 8 & 54 & 43 & 27 & 71 & 92 \\
第 3 趟 & 3 & 19 & 8 & 36 & 43 & 27 & 54 & 71 & 92 \\
第 4 趟 & 2 & 8 & 19 & 36 & 27 & 43 & 54 & 71 & 92 \\
第 5 趟 & 1 & 8 & 19 & 27 & 36 & 43 & 54 & 71 & 92 \\
第 6 趟 & 0 & 8 & 19 & 27 & 36 & 43 & 54 & 71 & 92 \\
\midrule
\multicolumn{10}{l}{\textbf{選擇排序}（第二欄：這一趟的最小值原本在哪一格）} \\
第 1 趟 & 5 & 8 & 71 & 19 & 54 & 36 & 92 & 43 & 27 \\
第 2 趟 & 3 & 8 & 19 & 71 & 54 & 36 & 92 & 43 & 27 \\
第 3 趟 & 8 & 8 & 19 & 27 & 54 & 36 & 92 & 43 & 71 \\
第 4 趟 & 5 & 8 & 19 & 27 & 36 & 54 & 92 & 43 & 71 \\
第 5 趟 & 7 & 8 & 19 & 27 & 36 & 43 & 92 & 54 & 71 \\
第 6 趟 & 7 & 8 & 19 & 27 & 36 & 43 & 54 & 92 & 71 \\
第 7 趟 & 8 & 8 & 19 & 27 & 36 & 43 & 54 & 71 & 92 \\
\midrule
\multicolumn{10}{l}{\textbf{插入排序}（第二欄：$\mathit{key}$）} \\
$i = 2$ & 71 & 36 & 71 & 19 & 54 & 8 & 92 & 43 & 27 \\
$i = 3$ & 19 & 19 & 36 & 71 & 54 & 8 & 92 & 43 & 27 \\
$i = 4$ & 54 & 19 & 36 & 54 & 71 & 8 & 92 & 43 & 27 \\
$i = 5$ & 8 & 8 & 19 & 36 & 54 & 71 & 92 & 43 & 27 \\
$i = 6$ & 92 & 8 & 19 & 36 & 54 & 71 & 92 & 43 & 27 \\
$i = 7$ & 43 & 8 & 19 & 36 & 43 & 54 & 71 & 92 & 27 \\
$i = 8$ & 27 & 8 & 19 & 27 & 36 & 43 & 54 & 71 & 92 \\
\midrule
\multicolumn{10}{l}{\textbf{Shell 排序}，增量 $4, 2, 1$（第二欄：挪動次數）} \\
$h = 4$ & 2 & 8 & 71 & 19 & 27 & 36 & 92 & 43 & 54 \\
$h = 2$ & 3 & 8 & 27 & 19 & 54 & 36 & 71 & 43 & 92 \\
$h = 1$ & 4 & 8 & 19 & 27 & 36 & 43 & 54 & 71 & 92 \\
\bottomrule
\end{tabular}
\end{center}

\begin{center}
\footnotesize
\begin{tabular}{ll*{8}{c}}
\toprule
& & 1 & 2 & 3 & 4 & 5 & 6 & 7 & 8 \\
\midrule
原始陣列 & & 36 & 71 & 19 & 54 & 8 & 92 & 43 & 27 \\
\midrule
\multicolumn{10}{l}{\textbf{合併排序}，由下而上} \\
$w = 1$ & & 36 & 71 & 19 & 54 & 8 & 92 & 27 & 43 \\
$w = 2$ & & 19 & 36 & 54 & 71 & 8 & 27 & 43 & 92 \\
$w = 4$ & & 8 & 19 & 27 & 36 & 43 & 54 & 71 & 92 \\
\midrule
\multicolumn{10}{l}{\textbf{快速排序}，Lomuto（第一欄：範圍；第二欄：基準 $x$；粗體：歸位的基準）} \\
$[1..8]$ & 27 & 19 & 8 & \textbf{27} & 54 & 71 & 92 & 43 & 36 \\
$[1..2]$ & 8 & \textbf{8} & 19 & 27 & 54 & 71 & 92 & 43 & 36 \\
$[4..8]$ & 36 & 8 & 19 & 27 & \textbf{36} & 71 & 92 & 43 & 54 \\
$[5..8]$ & 54 & 8 & 19 & 27 & 36 & 43 & \textbf{54} & 71 & 92 \\
$[7..8]$ & 92 & 8 & 19 & 27 & 36 & 43 & 54 & 71 & \textbf{92} \\
\midrule
\multicolumn{10}{l}{\textbf{快速排序}，Hoare（第二欄：基準 $x$ 和回傳的 $j$）} \\
$[1..8]$ & 36，3 & 27 & 8 & 19 & 54 & 71 & 92 & 43 & 36 \\
$[1..3]$ & 27，2 & 19 & 8 & 27 & 54 & 71 & 92 & 43 & 36 \\
$[1..2]$ & 19，1 & 8 & 19 & 27 & 54 & 71 & 92 & 43 & 36 \\
$[4..8]$ & 54，5 & 8 & 19 & 27 & 36 & 43 & 92 & 71 & 54 \\
$[4..5]$ & 36，4 & 8 & 19 & 27 & 36 & 43 & 92 & 71 & 54 \\
$[6..8]$ & 92，7 & 8 & 19 & 27 & 36 & 43 & 54 & 71 & 92 \\
$[6..7]$ & 54，6 & 8 & 19 & 27 & 36 & 43 & 54 & 71 & 92 \\
\midrule
\multicolumn{10}{l}{\textbf{堆積排序}（第二欄：下沉時交換幾次；灰色：已排好）} \\
建堆積 & & 92 & 71 & 43 & 54 & 8 & 19 & 36 & 27 \\
取出 92 & 2 & 71 & 54 & 43 & 27 & 8 & 19 & 36 & \sofin{92} \\
取出 71 & 1 & 54 & 36 & 43 & 27 & 8 & 19 & \sofin{71} & \sofin{92} \\
取出 54 & 1 & 43 & 36 & 19 & 27 & 8 & \sofin{54} & \sofin{71} & \sofin{92} \\
取出 43 & 2 & 36 & 27 & 19 & 8 & \sofin{43} & \sofin{54} & \sofin{71} & \sofin{92} \\
取出 36 & 1 & 27 & 8 & 19 & \sofin{36} & \sofin{43} & \sofin{54} & \sofin{71} & \sofin{92} \\
取出 27 & 0 & 19 & 8 & \sofin{27} & \sofin{36} & \sofin{43} & \sofin{54} & \sofin{71} & \sofin{92} \\
取出 19 & 0 & 8 & \sofin{19} & \sofin{27} & \sofin{36} & \sofin{43} & \sofin{54} & \sofin{71} & \sofin{92} \\
\bottomrule
\end{tabular}
\end{center}

比較次數：氣泡 27 次（交換 15 次，等於逆序對個數），選擇 28 次，插入 20 次，
Shell 23 次，合併 17 次，快速排序 Lomuto 16 次、Hoare 37 次，堆積 25 次。 合
併排序的 $17$ 次正好是 $n = 8$ 的最壞值 $8 \cdot 3 - 8 + 1$：每一次合併都要比到其中一邊剩最
後一個。

\begin{examnote}
「一趟」的定義各本課本不同：氣泡排序有的從右往左把最小值換到前面，選擇排序有的每趟找
最大值放到後面，快速排序的基準和分割方式有好幾種，合併排序有遞迴和由下而上兩種，建堆積
有逐一插入和由下而上兩種。 題目有指定就照題目做；沒有指定就選一種，寫一句說明，從頭到尾
照它做。
\end{examnote}

\subsection{例題：\texttt{qsort} 比較函式的溢位}

實際寫 C 程式時，排序多半交給標準程式庫的 \texttt{qsort}，自己只寫比較函式。 比較函式
最常見的錯，是用減法決定大小。

\begin{cbookrecap}[第~13~章 13.1、13.2 節，第~18~章 18.6 節]
\texttt{qsort(base, n, size, cmp)} 排序從 \texttt{base} 開始的 \texttt{n} 個元素，每個
\texttt{size} 個位元組。 \texttt{cmp(p, q)} 拿到兩個指向元素的 \texttt{const void *}，第
一個應該排在前面就回傳負數，後面就回傳正數，相等回傳 $0$。 這個比較必須對所有元素定
出一個一致的全序（$a < b$、$b < c$ 就要 $a < c$），否則 C 標準不保證 \texttt{qsort} 會
做什麼。

有號整數的運算結果超出型別範圍是未定義行為；把超出範圍的值轉成有號整數型別，結果則
由編譯器決定（implementation-defined），gcc 的做法是取 $2^{32}$ 的餘數（\texttt{int}
是 32 位元時）。 無號整數的運算一定是取餘數，不會溢位。
\end{cbookrecap}

第一種錯是 \texttt{int} 的比較函式寫成 \texttt{return x - y;}。 $x$ 和 $y$ 差很遠時差會超
出 \texttt{int} 的範圍。 例如 $x = 2{,}000{,}000{,}000$、$y = -1{,}000{,}000{,}000$，真正的差
$3 \times 10^9$ 大於 \texttt{INT\_MAX}（$2{,}147{,}483{,}647$），這是有號整數溢位，屬於未定
義行為：程式可能回傳一個負數，讓 $x$ 排到 $y$ 前面，也可能發生其他事，而且開了
\texttt{-fsanitize=undefined} 就會被抓出來。 本書的程式不能真的執行這種寫法，所以下面
示範第二種同樣常見、但行為確定的錯：鍵是 \texttt{long long}，比較函式卻把差直接當
\texttt{int} 回傳。

情境是一個雲端硬碟依檔案大小（位元組）排序。 大小可以超過 $2^{31}$，所以用
\texttt{long long} 存；\texttt{x - y} 是 \texttt{long long} 的減法，不會溢位，但回傳時被
轉成 \texttt{int}，只留下低 32 位元。 程式先印出錯誤比較函式的三個回傳值，再用它排序。
為了讓結果固定，用錯誤的比較函式排序時用的是自己寫的插入排序，因為把不一致的比較函式
交給 \texttt{qsort} 本身就不受保證。

\begin{example}[label=ex:so-cmp]{c}{(cmp\_trap.c) 比較函式用減法的後果}
#include <stdio.h>
#include <stdlib.h>

// 錯的寫法：long long 的差直接當 int 回傳。差超過 int 的範圍時，
// 轉成 int 的結果由編譯器決定（gcc 取 2^32 的餘數），正負號可能反過來。
int cmp_size_bad(const void *p, const void *q)
{
    long long x = *(const long long *)p;
    long long y = *(const long long *)q;
    return x - y;
}

// 對的寫法：只比較，不相減。
int cmp_size(const void *p, const void *q)
{
    long long x = *(const long long *)p;
    long long y = *(const long long *)q;
    return (x > y) - (x < y);
}

// 用比較函式做插入排序。把錯的比較函式交給 qsort 是未定義行為，
// 所以示範時用自己寫的排序，結果才固定。
void insertion_sort(long long *a, int n, int (*cmp)(const void *, const void *))
{
    for (int i = 1; i < n; i++) {
        long long key = a[i];
        int j = i - 1;
        while (j >= 0 && cmp(&a[j], &key) > 0) {
            a[j + 1] = a[j];
            j--;
        }
        a[j + 1] = key;
    }
}

int main(void)
{
    long long c[3] = {0, 1500000000LL, 3000000000LL};
    for (int i = 0; i < 3; i++) {       // 三個數兩兩比較
        int j = (i + 1) % 3;
        printf("bad(%lld, %lld) = %d\n", c[i], c[j], cmp_size_bad(&c[i], &c[j]));
    }
    long long a[] = {52000000, 3500000000LL, 1200, 2300000000LL, 640000};
    long long b[] = {52000000, 3500000000LL, 1200, 2300000000LL, 640000};
    int n = sizeof a / sizeof a[0];
    insertion_sort(a, n, cmp_size_bad);
    qsort(b, n, sizeof b[0], cmp_size);
    printf("bad:");
    for (int i = 0; i < n; i++) {
        printf(" %lld", a[i]);
    }
    printf("\nqsort:");
    for (int i = 0; i < n; i++) {
        printf(" %lld", b[i]);
    }
    printf("\n");
    return 0;
}
\end{example}
\iosample{}{bad(0, 1500000000) = -1500000000 \\ bad(1500000000, 3000000000) = -1500000000 \\ bad(3000000000, 0) = -1294967296 \\ bad: 2300000000 3500000000 1200 640000 52000000 \\ qsort: 1200 640000 52000000 2300000000 3500000000}

前兩列的差是 $-1.5 \times 10^9$，在 \texttt{int} 的範圍內，答案對。 第三列 $3 \times 10^9 - 0$
超過 $2^{31} - 1$，取 $2^{32}$ 的餘數變成 $3 \times 10^9 - 2^{32} = -1{,}294{,}967{,}296$。 於是
$0$ 比 $1.5 \times 10^9$ 小，$1.5 \times 10^9$ 比 $3 \times 10^9$ 小，$3 \times 10^9$ 卻比 $0$
小，形成一個循環，根本不是全序，任何排序都不可能「照它排對」。 用它排五個檔案的大小，兩
個超過 $2^{31}$ 的大檔案排到了最前面；\texttt{qsort} 配正確的比較函式排出正確的結果。 正確的寫法 \texttt{(x > y) - (x < y)} 不做減法，兩個比較的結果都
是 $0$ 或 $1$，相減只會是 $-1$、$0$、$1$，對 \texttt{int}、\texttt{long long}、
\texttt{double} 都適用。

\section{本章常見錯誤}
\label{sec:so-pitfalls}

\begin{center}
\small
\begin{tabular}{p{0.24\linewidth}p{0.50\linewidth}p{0.16\linewidth}}
\toprule
錯誤 & 例子 & 見 \\
\midrule
最好情況沒看條件 & 說氣泡排序最好是 $\Theta(n)$，但題目的程式沒有提早結束；說堆積排序最好是
$\Theta(n)$，但鍵互不相同 & 第~\ref{sec:so-compare}~節 \\
穩定性只憑印象 & 合併時用 $<$、插入時用 $\ge$，程式已經不穩定，答案卻寫「穩定」 &
第~\ref{sec:so-compare}~節 \\
合併忘了寫回 & 只合併到 $B$ 沒寫回 $A$；由下而上最後一段的右界沒取 $\min$ 而越界 &
第~\ref{sec:so-merge}~節 \\
遞迴版和由下而上混用 & 題目要遞迴版，卻寫出兩兩合併的逐趟結果；$n$ 不是 $2$ 的次方時兩者不同 &
第~\ref{sec:so-merge}~節 \\
分割的遞迴範圍 & Hoare 分割遞迴 $[p..j-1]$；Lomuto 遞迴 $[p..q]$，基準一直留在範圍裡 &
第~\ref{sec:so-quick}~節 \\
Hoare 的基準取最後一個 & $j$ 可能等於 $r$，無窮遞迴 & 第~\ref{sec:so-quick}~節 \\
說快速排序是 $O(n \lg n)$ & 沒說是平均或期望；最壞是 $\Theta(n^2)$，大量相同的鍵用 Lomuto 分割也是 & 第~\ref{sec:so-quick}~節 \\
堆積大小沒有縮小 & 交換以後仍用 $n$ 下沉，把排好的尾端拉回堆積 & 第~\ref{sec:so-heap}~節 \\
比較函式用減法 & \texttt{return x - y;}，差超過 \texttt{int} 的範圍；\texttt{double} 的差被截成
$0$ & 第~\ref{sec:so-compare}~節 \\
\bottomrule
\end{tabular}
\end{center}

\section{觀念題}
\label{sec:so-exercises}

題型和研究所考試、期中期末考的筆試相近，有選擇、計算、追蹤、證明和設計。 沒有特別說
明時，陣列從 $1$ 開始編號，排序方法照本章的虛擬碼。 解答在本章最後。

\begin{enumerate}
\item （選擇，可能不只一個答案）下列哪些敘述正確？
\begin{enumerate}
\item 選擇排序的比較次數和輸入的內容無關。
\item 如果每個數和它排好之後的位置相差不超過 $k$ 格，插入排序的時間是 $O(nk)$。
\item 合併排序（虛擬碼~\ref{alg:so-mergesort}）排已經排好的輸入，只要 $\Theta(n)$ 次比較。
\item 快速排序改成隨機選基準以後，最壞情況是 $\Theta(n \lg n)$。
\item 堆積排序除了輸入陣列以外，只需要 $\Theta(1)$ 的額外空間。
\end{enumerate}

\item （選擇）下列修改，哪些修改之後的排序仍然是穩定的？ 不穩定的請舉反例。
\begin{enumerate}
\item 氣泡排序第 5 行改成 $A[j] \ge A[j+1]$。
\item 插入排序的內層條件改成 $A[j] \ge \mathit{key}$。
\item \proc{Merge} 第 6 行的 $B[i] \le B[j]$ 改成 $B[i] < B[j]$。
\item 選擇排序第 5 行改成 $A[j] \le A[m]$。
\item \proc{Merge-Sort} 第 3 行改成 $q \gets \lceil (p + r)/2 \rceil$。
\end{enumerate}

\item\label{so-q-bubble} （追蹤）用氣泡排序（虛擬碼~\ref{alg:so-bubble}）排
$\langle 45, 12, 78, 33, 9, 61, 27 \rangle$。 寫出每一趟之後的陣列，以及總比較次數和總交換次
數，並驗證交換次數等於逆序對的個數。

\item （計算）令 $W(n)$ 是合併排序（虛擬碼~\ref{alg:so-mergesort}）在長度 $n$ 時最多的比較
次數，$W(1) = 0$，$W(n) = W(\lceil n/2 \rceil) + W(\lfloor n/2 \rfloor) + n - 1$。
\begin{enumerate}
\item 算出 $W(1), W(2), \ldots, W(8)$。
\item 證明 $W(2^k) = k\,2^k - 2^k + 1$。
\item 給一個 $1$ 到 $8$ 的排列，讓合併排序剛好比較 $W(8)$ 次，並說明為什麼。
\end{enumerate}

\item （追蹤）$A = \langle 33, 71, 18, 54, 26, 90, 41 \rangle$。
\begin{enumerate}
\item 用 Lomuto 分割的快速排序（虛擬碼~\ref{alg:so-quick}），寫出每一次分割之後的陣列和基
準的位置，並算出比較次數。
\item 用 Hoare 分割的快速排序（虛擬碼~\ref{alg:so-quick-hoare}），寫出每一次分割的範圍、
回傳的 $j$ 和之後的陣列。
\end{enumerate}

\item （設計）\textbf{三路分割}：設計一個只掃一次的分割程序，把 $A[p..r]$ 排成 $< x$、
$= x$、$> x$ 三段（$x = A[r]$），回傳中間那一段的頭尾位置。 寫出虛擬碼和迴圈不變量，再說
明用它做快速排序時，所有鍵都相同的輸入只要 $\Theta(n)$ 的時間。

\item （說明）Hoare 分割的快速排序。
\begin{enumerate}
\item 如果遞迴寫成 \proc{Quicksort-H}$(A, p, j - 1)$ 和 \proc{Quicksort-H}$(A, j + 1, r)$，
找一個長度 $3$ 的輸入，讓結果沒有排好。
\item 如果 \proc{Hoare-Partition} 第 2 行改成 $x \gets A[r]$，說明 $\langle 1, 2 \rangle$ 會讓
\proc{Quicksort-H} 無窮遞迴。
\end{enumerate}

\item （證明）\proc{Randomized-Quicksort}，鍵互不相同，記號和定理~\ref{thm:so-quick-avg}
一樣。
\begin{enumerate}
\item 證明 $z_i$ 和 $z_{i+1}$ 一定會被比較。
\item 求 $z_1$ 和 $z_n$ 被比較的機率。
\item 證明 $z_k$ 參與的比較次數期望值是 $2(H_k + H_{n-k+1} - 2)$，並用它重新算出
$\mathrm{E}[X]$。
\end{enumerate}

\item （C 程式）下面的比較函式要讓學生依平均成績由低到高排序。 指出它的問題，舉一組會排錯
的成績，並改正。
\begin{code}{c}{觀念題 9}
typedef struct {
    char name[16];
    double gpa;
} Student;

int cmp_gpa(const void *p, const void *q)
{
    const Student *a = p, *b = q;
    return a->gpa - b->gpa;
}
\end{code}
\end{enumerate}

\section{實作練習}
\label{sec:so-practice}

下面是 Aki-Judge 上和本章相關的題目。 幾題只要交排序函式，不寫 \texttt{main}；評測資
料有已經排好、由大到小、全部相同、含 \texttt{INT\_MIN} 和 \texttt{INT\_MAX} 的陣列，交之
前先用這幾種輸入自己測一次。

\begin{center}
\begin{tabular}{llc}
\toprule
小節 & 題目 & 難度 \\
\midrule
比較函式 & \akiref{sorting-integers-with-a-safe-comparator}{整數排序與比較函式溢位} & 暖身 \\
 & \akiref{group-by-distinct-digits}{依不同數字個數分組排序} & 基本 \\
排序的應用 & \akiref{counting-distinct-integers}{計算相異整數的個數} & 暖身 \\
合併排序 & \akiref{merge-sort-function}{合併排序函式} & 基本 \\
 & \akiref{merge-sort-function-dc}{合併排序函式} & 基本 \\
 & \akiref{three-way-merge-sort-by-absolute-value}{依絕對值的三路合併排序} & 進階 \\
快速排序 & \akiref{quicksort-on-adversarial-input}{抗最壞情況的快速排序} & 進階 \\
逆序對 & \akiref{counting-inversions}{計算逆序對的數量} & 進階 \\
 & \akiref{count-inversions-in-an-array}{計算逆序對數量} & 進階 \\
\bottomrule
\end{tabular}
\end{center}

幾題的提示：「整數排序與比較函式溢位」和「依不同數字個數分組排序」用 \texttt{qsort}，比較函
式照第~\ref{sec:so-compare}~節的寫法不要相減，後者依序比三個鍵，數值要用 \texttt{long long}。
「計算相異整數的個數」先排序，再數有幾個位置和前一個不同。 兩題「合併排序函式」要交的函
式一樣，可以一題用遞迴版、一題用由下而上的版本。 「依絕對值的三路合併排序」把比較規則寫成
一個函式，合併三段時呼叫它；輸出的順序就是遞迴的後序。 「抗最壞情況的快速排序」的測資有長
度 $2 \times 10^5$ 的已排好和由大到小的陣列，要用三數取中或隨機基準，並先遞迴短的一邊。 兩
題逆序對都可以在合併時順便數：\proc{Merge} 從右半拿 $B[j]$ 時，左半還沒拿的 $q - i + 1$ 個
數都比它大；相等不算逆序，所以相等時先拿左邊的，答案要用 \texttt{long long}。 第~8~章題單的
「以堆積排序整數陣列」也是第~\ref{sec:so-heap}~節的練習。

\section{讀完你應該會}
\label{sec:so-checklist}

\begin{itemize}
\item 說出每種排序是否穩定、是否原地、額外空間多少；判斷不穩定時能用兩個相同的鍵舉出反例。
\item 給一組資料，寫出氣泡、選擇、插入、Shell、合併（遞迴與由下而上）、快速（Lomuto 與
Hoare）、堆積排序每一趟之後的陣列，並數出比較或交換的次數。
\item 寫出 Lomuto 和 Hoare 分割，說出回傳值的意義和對應的遞迴範圍，用迴圈不變量證明
Lomuto 分割正確。
\item 說出每種排序最好、平均、最壞的時間和發生的條件；用指示變數證明隨機化快速排序
的期望比較次數是 $2(n+1)H_n - 4n$。
\item 寫出不會溢位、能定出全序的 \texttt{qsort} 比較函式，並依資料的特性選擇排序方法。
\end{itemize}

\textbf{延伸閱讀}（書末會整理完整的參考資料）

\begin{itemize}
\item Thomas H. Cormen 等，\textit{Introduction to Algorithms}，第 4 版，第 2.3 節（合併排
序）、第 6 章（堆積排序）、第 7 章（快速排序）。 第 7 章的章末問題有 Hoare 分割的正確性、
相同的鍵、堆疊深度和三數取中的分析。
\item Donald E. Knuth，\textit{The Art of Computer Programming, Vol.~3: Sorting and
Searching}，第 2 版，第 5.2 節：插入與 Shell 排序（5.2.1，含 $h$ 排序好的陣列再做 $k$ 排
序仍是 $h$ 排序好的證明）、交換排序與快速排序（5.2.2）、選擇與堆積排序（5.2.3）、合併
排序（5.2.4），各種比較次數的精確分析都在這裡。
\item Robert Sedgewick、Kevin Wayne，\textit{Algorithms}，第 4 版，第 2 章：實作導向，三
路分割的快速排序講得很清楚。
\item Jon L. Bentley、M. Douglas McIlroy，``Engineering a Sort Function''，
\textit{Software: Practice and Experience} 23(11)，1993：一個實際的 \texttt{qsort} 怎麼
選基準、怎麼處理相同的鍵。
\item Russel Schaffer、Robert Sedgewick，``The Analysis of Heapsort''，\textit{Journal of
Algorithms} 15(1)，1993：堆積排序最好情況也是 $\Theta(n \lg n)$ 的證明。
\item 《C 程式設計基礎》第~18~章 18.6 節：函式指標與 \texttt{qsort}、字串和結構的比較函式。
\end{itemize}

\addcontentsline{toc}{section}{觀念題解答}
\section*{觀念題解答}

\begin{enumerate}
\item (a)、(b)、(e) 正確。 (a)：第 $i$ 趟一定比 $n - i$ 次，共 $n(n-1)/2$ 次。 (b)：設 $x$ 現
在在位置 $i$，排好後在位置 $f$，$|i - f| \le k$。 一個在 $x$ 左邊、比 $x$ 大的數 $y$，排好後的
位置大於 $f$，所以現在的位置至少是 $f + 1 - k$；它又在 $x$ 左邊，位置小於 $i \le f + k$。 這
樣的位置不到 $2k$ 個，所以逆序對 $I < 2kn$，插入排序的時間 $\Theta(n + I) = O(nk)$。 (c)：
排好的輸入每次合併仍要比左半的個數次，共約 $\frac{n}{2}\lg n$ 次。 (d)：隨機選基準只改善
期望值，運氣最差時每次都選到最大或最小，仍是 $\Theta(n^2)$。 (e)：堆積排序是原地的，下沉
也不用遞迴。

\item 只有 (e)。 (a)、(b)、(c)、(d) 都可以用 $\langle 1_a, 1_b \rangle$ 當反例：(a) $A[1] \ge
A[2]$ 成立，交換；(b) $\mathit{key} = 1_b$ 時 $A[1] = 1_a \ge \mathit{key}$，$1_a$ 被挪到右邊；
(c) 合併 $\langle 1_a \rangle$ 和 $\langle 1_b \rangle$ 時 $B[i] < B[j]$ 不成立，先拿右邊的
$1_b$；(d) $j = 2$ 時 $A[2] \le A[1]$ 成立，$m = 2$，交換後是 $\langle 1_b, 1_a \rangle$。 這
四個都變成 $\langle 1_b, 1_a \rangle$。 (e) 只改變切開的位置，左半仍在右半前面，\proc{Merge}
遇到相同的鍵仍先拿左半的，所以仍然穩定。

\item 每一趟之後：
\begin{center}
\begin{tabular}{ll*{7}{c}}
\toprule
& 交換 & 1 & 2 & 3 & 4 & 5 & 6 & 7 \\
\midrule
開始 & & 45 & 12 & 78 & 33 & 9 & 61 & 27 \\
第 1 趟 & 5 & 12 & 45 & 33 & 9 & 61 & 27 & 78 \\
第 2 趟 & 3 & 12 & 33 & 9 & 45 & 27 & 61 & 78 \\
第 3 趟 & 2 & 12 & 9 & 33 & 27 & 45 & 61 & 78 \\
第 4 趟 & 2 & 9 & 12 & 27 & 33 & 45 & 61 & 78 \\
第 5 趟 & 0 & 9 & 12 & 27 & 33 & 45 & 61 & 78 \\
\bottomrule
\end{tabular}
\end{center}
第 5 趟沒有交換，提早結束。 比較 $6 + 5 + 4 + 3 + 2 = 20$ 次，交換 12 次。 逆序對：
$45$ 和 $12, 33, 9, 27$；$12$ 和 $9$；$78$ 和 $33, 9, 61, 27$；$33$ 和 $9, 27$；$61$ 和 $27$，
共 $4 + 1 + 4 + 2 + 1 = 12$ 個。

\item (a) $W(1..8) = 0, 1, 3, 5, 8, 11, 14, 17$。

(b) 對 $k$ 歸納。 $k = 0$：$W(1) = 0 = 0 - 1 + 1$。 $k \ge 1$：
$W(2^k) = 2W(2^{k-1}) + 2^k - 1 = 2\bigl((k-1)2^{k-1} - 2^{k-1} + 1\bigr) + 2^k - 1
= k\,2^k - 2^k + 1$。

(c) 例如 $\langle 1, 5, 3, 7, 2, 6, 4, 8 \rangle$。 一次合併要比滿 $m_1 + m_2 - 1$ 次，條件是
兩邊的數一直交錯，直到其中一邊只剩最後一個。 這個排列把排好的 $1..8$ 依名次的奇偶分成左右
兩半，每一半再照同樣的方法分，所以每一次合併兩邊都完全交錯：四次 $1 + 1$ 的合併各比 $1$
次，兩次 $2 + 2$ 各比 $3$ 次，最後 $4 + 4$ 比 $7$ 次，共 $4 + 6 + 7 = 17$。 第
\ref{sec:so-compare}~節例題的 $\langle 36, 71, 19, 54, 8, 92, 43, 27 \rangle$ 也剛好比 $17$ 次。

\item (a) 粗體是歸位的基準。
\begin{center}
\begin{tabular}{ll*{7}{c}}
\toprule
& $x$ & 1 & 2 & 3 & 4 & 5 & 6 & 7 \\
\midrule
開始 & & 33 & 71 & 18 & 54 & 26 & 90 & 41 \\
$[1..7]$ & 41 & 33 & 18 & 26 & \textbf{41} & 71 & 90 & 54 \\
$[1..3]$ & 26 & 18 & \textbf{26} & 33 & 41 & 71 & 90 & 54 \\
$[5..7]$ & 54 & 18 & 26 & 33 & 41 & \textbf{54} & 90 & 71 \\
$[6..7]$ & 71 & 18 & 26 & 33 & 41 & 54 & \textbf{71} & 90 \\
\bottomrule
\end{tabular}
\end{center}
比較 $6 + 2 + 2 + 1 = 11$ 次。

(b)
\begin{center}
\begin{tabular}{ll*{7}{c}}
\toprule
& $x$，$j$ & 1 & 2 & 3 & 4 & 5 & 6 & 7 \\
\midrule
開始 & & 33 & 71 & 18 & 54 & 26 & 90 & 41 \\
$[1..7]$ & 33，2 & 26 & 18 & 71 & 54 & 33 & 90 & 41 \\
$[1..2]$ & 26，1 & 18 & 26 & 71 & 54 & 33 & 90 & 41 \\
$[3..7]$ & 71，5 & 18 & 26 & 41 & 54 & 33 & 90 & 71 \\
$[3..5]$ & 41，3 & 18 & 26 & 33 & 54 & 41 & 90 & 71 \\
$[4..5]$ & 54，4 & 18 & 26 & 33 & 41 & 54 & 90 & 71 \\
$[6..7]$ & 90，6 & 18 & 26 & 33 & 41 & 54 & 71 & 90 \\
\bottomrule
\end{tabular}
\end{center}

\item 一種寫法（Dijkstra 的「荷蘭國旗」問題）：
\begin{pseudo}{觀念題 6 的解答}
\Function{Three-Way-Partition}{$A, p, r$}
    \State $x \gets A[r]$
    \State $\mathit{lt} \gets p$；$i \gets p$；$\mathit{gt} \gets r$
    \While{$i \le \mathit{gt}$}
        \If{$A[i] < x$}
            \State 交換 $A[\mathit{lt}]$ 和 $A[i]$；$\mathit{lt} \gets \mathit{lt} + 1$；$i \gets i + 1$
        \ElsIf{$A[i] > x$}
            \State 交換 $A[i]$ 和 $A[\mathit{gt}]$；$\mathit{gt} \gets \mathit{gt} - 1$
        \Else
            \State $i \gets i + 1$
        \EndIf
    \EndWhile
    \State \Return $(\mathit{lt}, \mathit{gt})$
\EndFunction
\end{pseudo}
不變量：第 4 行每次檢查時，$A[p..\mathit{lt}-1] < x$，$A[\mathit{lt}..i-1] = x$，
$A[\mathit{gt}+1..r] > x$，$A[i..\mathit{gt}]$ 還沒看。 一開始前三段都是空的。 $A[i] < x$
時，$A[\mathit{lt}]$ 是第二段的第一個（$= x$），或第二段是空的時就是 $A[i]$ 自己，交換後
第一段、第二段各往右一格；$A[i] > x$ 時換到 $A[\mathit{gt}]$，第四段往左長一格，換過來的
數還沒看，所以 $i$ 不動；$A[i] = x$ 時第二段長一格。 每一輪 $i$ 加一或 $\mathit{gt}$ 減一，
$r - p + 1$ 輪內結束，結束時 $i = \mathit{gt} + 1$，中間一段 $A[\mathit{lt}..\mathit{gt}]$
全部等於 $x$，而且至少有 $x$ 自己。 快速排序只遞迴 $[p..\mathit{lt}-1]$ 和
$[\mathit{gt}+1..r]$。 鍵全部相同時，第一次分割每一輪都走 \textbf{else}，結束時
$\mathit{lt} = p$、$\mathit{gt} = r$，兩邊都是空的，總共 $\Theta(n)$。

\item (a) $\langle 3, 1, 2 \rangle$。 $x = 3$：$i$ 停在 $1$，$j$ 停在 $3$（$2 \le 3$），交換得
$\langle 2, 1, 3 \rangle$；接著 $i$ 走到 $3$，$j$ 停在 $2$，回傳 $j = 2$。 正確的遞迴是
$[1..2]$ 和 $[3..3]$；錯的寫法遞迴 $[1..1]$ 和 $[3..3]$，$A[2] = 1$ 沒有參加，結果停在
$\langle 2, 1, 3 \rangle$。

(b) $x = A[2] = 2$。 $i$：$A[1] = 1 < 2$，走到 $i = 2$ 停下。 $j$：$A[2] = 2$ 不大於 $2$，在
$j = 2$ 停下。 $i \ge j$，回傳 $2 = r$。 於是遞迴呼叫 \proc{Quicksort-H}$(A, 1, 2)$，範圍和
原來一模一樣，陣列也沒變，永遠不會結束。

\item (a) $Z_{i,i+1} = \{z_i, z_{i+1}\}$ 只有這兩個數，第一個被選為基準的一定是其中之一，由
引理~\ref{lem:so-pair}，它們一定會被比較。 直觀上，沒有別的數能把它們分開。

(b) $j - i + 1 = n$，機率是 $2/n$：$z_1$ 和 $z_n$ 會被比較，若且唯若第一次分割的基準就是最
小值或最大值。

(c) $z_k$ 參與的比較就是 $\sum_{j \ne k} X_{kj}$，期望值是
\[
\sum_{j \ne k} \frac{2}{|j - k| + 1}
= 2\left(\frac12 + \cdots + \frac1k\right) + 2\left(\frac12 + \cdots + \frac{1}{n-k+1}\right)
= 2(H_k + H_{n-k+1} - 2).
\]
每次比較有兩個數參與，所以把 $k = 1..n$ 加起來再除以 $2$：
$\mathrm{E}[X] = \sum_{k=1}^{n}(H_k + H_{n-k+1} - 2) = 2\sum_{k=1}^{n} H_k - 2n$。 用
$\sum_{k=1}^{n} H_k = (n+1)H_n - n$，得到 $2(n+1)H_n - 4n$，和定理~\ref{thm:so-quick-avg}
一樣。 這和第~9~章第~\ref{sec:bst-height}~節框裡 $j$ 的期望深度 $H_j + H_{n-j+1} - 2$ 只差一個 $2$ 倍：BST 裡
只算祖先，快速排序裡兩個方向都算。

\item \texttt{a->gpa - b->gpa} 是 \texttt{double}，回傳時轉成 \texttt{int} 會把小數部分捨去
（往 $0$ 截斷）。 兩個成績相差不到 $1$ 時差被截成 $0$，比較函式說它們相等。 例如成績
$3.9$ 和 $3.2$，\texttt{qsort} 可以把 $3.9$ 排在 $3.2$ 前面。 更嚴重的是它不是一個一致的全
序：$2.9$ 和 $3.5$「相等」，$3.5$ 和 $4.0$「相等」，$2.9$ 卻比 $4.0$ 小，C 標準不保證
\texttt{qsort} 在這種比較函式下會做什麼。 改成只比較、不相減：最後一行寫成
\texttt{return (a->gpa > b->gpa) - (a->gpa < b->gpa);}。
\end{enumerate}

\chapter{排序的下界、線性時間排序與外部排序}
\label{ch:linear}

% 本章的圖共用的樣式。 敗者樹、勝者樹的圖由模擬程式依虛擬碼逐步執行後產生。
\tikzset{
    dtn/.style={draw=black!65, rounded corners=2pt, minimum width=0.95cm, minimum height=0.55cm, inner sep=1pt, font=\small},
    dtleaf/.style={draw=greenStroke, fill=greenFill, minimum height=0.55cm, inner sep=2pt, font=\small},
    dtbad/.style={draw=black!45, dashed, fill=black!4, minimum height=0.55cm, inner sep=2pt, font=\small, text=black!55},
    dte/.style={draw=black!60},
    dtpath/.style={draw=memStroke, very thick},
    lelab/.style={font=\scriptsize, text=black!65, inner sep=1pt},
    lte/.style={draw=black!60},
    ltn/.style={circle, draw=black!65, minimum size=0.62cm, inner sep=0pt, font=\small},
    ltchg/.style={draw=memStroke, very thick, fill=memFill},
    ltwin/.style={rectangle, draw=coreStroke, thick, fill=coreFill, minimum width=0.7cm, minimum height=0.55cm, font=\small},
    ltleaf/.style={rectangle, draw=black!65, minimum width=0.7cm, minimum height=0.55cm, font=\small},
    ltnew/.style={draw=greenStroke, very thick, fill=greenFill},
    rlab/.style={font=\scriptsize, text=black!60, inner sep=1pt},
    restlab/.style={font=\scriptsize\ttfamily, text=black!55},
    lrunbox/.style={draw=black!60, fill=black!4, minimum height=0.5cm, inner sep=2pt, font=\scriptsize},
}

教務處在加退選結束後，要把全校一百二十萬筆選課紀錄依學號排好，再印出每個人的選
課單。 每筆紀錄 $500$ 位元組，加起來 $600$ MB，排序程式卻只分到 $60$ MB 的記憶
體。 助理想到第~14~章的合併排序，可是有兩件事卡住了他。 第一，只靠「比大小」排
一百二十萬筆，至少要比幾次？ 學號是九位數的整數，能不能利用這一點少做一些事？ 第
二，資料放不進記憶體，合併排序要先切成兩半各自排好，這兩半要放在哪裡？ 往下讀之
前，先想想你會怎麼處理這兩件事。

這一章分成三部分。 第~\ref{sec:lin-lower}~節證明：只靠比較兩個鍵的大小來排序，最
壞情況下一定要 $\Omega(n \lg n)$ 次比較，所以第~14~章的合併排序、堆積排序在這個模
型裡已經是最好的。 第~\ref{sec:lin-linear}~節跳出這個模型：鍵是範圍不大的整數時，
計數排序、基數排序可以在線性時間排好；鍵是均勻分布的實數時，桶子排序的期望時間也
是線性。 第~\ref{sec:lin-external}~節處理放不進記憶體的資料：先切成一段一段排好的
\textbf{串}，再用 $k$ 路合併把它們合起來，中間會講兩個考試常見、但多數演算法教科書
沒有的結構：\textbf{敗者樹}和\textbf{取代選擇}。

這一章用到第~1~章 1.6 節對 $\lg(n!)$ 的估計、第~8~章的堆積，以及第~14~章的合併排
序、快速排序和堆積排序；C 程式用到《C 程式設計基礎》第~14~章的結構和第~19~章的
\texttt{malloc}、\texttt{free}。

\section{比較排序的下界}
\label{sec:lin-lower}

\subsection{比較排序的模型}

第~14~章的排序（插入、合併、快速、堆積、Shell）有一個共同點：它們只靠\textbf{比較
兩個鍵}來得知輸入的資訊。 程式可以搬動元素、可以做任何計算，但關於鍵的值，它唯一
能問的問題是「$a_i \le a_j$ 嗎？」（問 $<$、$>$、$=$ 也一樣，都是一次比較）。 這樣
的排序叫\textbf{比較排序}（comparison sort）。

證明下界時，假設 $n$ 個鍵兩兩不同。 這只是把輸入限制在一部分情況；連這些情況都需
要那麼多次比較，所有情況合起來當然也需要。 鍵兩兩不同時，$a_i \le a_j$ 和
$a_i < a_j$ 是同一個問題，每次比較只有兩種結果。

\subsection{決策樹}

把一個比較排序對固定的 $n$ 所有可能的執行過程畫成一棵二元樹，叫做\textbf{決策樹}
（decision tree）。
\begin{itemize}
\item 每個內部節點標一個比較 $i{:}j$，意思是「比較原本的第 $i$ 個元素 $a_i$ 和第
$j$ 個元素 $a_j$」。 $a_i \le a_j$ 往左，$a_i > a_j$ 往右。
\item 每個葉子標一個排列 $\langle \pi(1), \pi(2), \ldots, \pi(n) \rangle$，意思是
演算法走到這裡就停下來，輸出 $a_{\pi(1)} \le a_{\pi(2)} \le \cdots \le a_{\pi(n)}$。
\item 給定一組輸入，演算法的執行就是從根走到某一個葉子的一條路徑，路徑上的內部節點
數就是這次做的比較次數，也就是那個葉子的\textbf{深度}（從根往下的邊數）。
\end{itemize}

搬動元素、計算索引這些動作都不出現在樹上；決策樹只記錄比較和比較的結果。 下面用一
個只排三個數的小演算法當例子：先比 $a_2$、$a_3$，把後兩個排好；再把 $a_1$ 插進
去，先和較小的那個比，必要時再和較大的那個比。

\begin{center}
\begin{tikzpicture}[x=1cm, y=1cm]
\node[dtn] (r) at (0,0) {$2{:}3$};
\node[dtn] (a) at (-2.6,-1.3) {$1{:}2$};
\node[dtn] (b) at (2.6,-1.3) {$1{:}3$};
\node[dtleaf] (aa) at (-4.0,-2.6) {$\langle 1,2,3 \rangle$};
\node[dtn] (ab) at (-1.3,-2.6) {$1{:}3$};
\node[dtleaf] (ba) at (1.3,-2.6) {$\langle 1,3,2 \rangle$};
\node[dtn] (bb) at (4.0,-2.6) {$1{:}2$};
\node[dtleaf] (aba) at (-2.2,-3.9) {$\langle 2,1,3 \rangle$};
\node[dtleaf] (abb) at (-0.4,-3.9) {$\langle 2,3,1 \rangle$};
\node[dtleaf] (bba) at (3.1,-3.9) {$\langle 3,1,2 \rangle$};
\node[dtleaf] (bbb) at (4.9,-3.9) {$\langle 3,2,1 \rangle$};
\draw[dtpath] (r) -- node[lelab, above left] {$\le$} (a);
\draw[dte] (r) -- node[lelab, above right] {$>$} (b);
\draw[dte] (a) -- node[lelab, above left] {$\le$} (aa);
\draw[dtpath] (a) -- node[lelab, above right] {$>$} (ab);
\draw[dte] (b) -- node[lelab, above left] {$\le$} (ba);
\draw[dte] (b) -- node[lelab, above right] {$>$} (bb);
\draw[dte] (ab) -- node[lelab, left] {$\le$} (aba);
\draw[dtpath] (ab) -- node[lelab, right] {$>$} (abb);
\draw[dte] (bb) -- node[lelab, left] {$\le$} (bba);
\draw[dte] (bb) -- node[lelab, right] {$>$} (bbb);
\end{tikzpicture}
\diagcaption{「先排後兩個，再插入第一個」排三個數的決策樹。 橘色是輸入
$\langle 52, 17, 30 \rangle$ 走的路徑：$17 \le 30$ 往左，$52 > 17$ 往右，$52 > 30$
往右，輸出 $\langle 2, 3, 1 \rangle$，也就是 $17, 30, 52$。}
\label{fig:lin-dtree}
\end{center}

這棵樹有 $6 = 3!$ 個葉子，深度是 $2, 3, 3, 2, 3, 3$。 所以這個演算法最壞要比 $3$
次；六種排列各出現一次的話，平均比 $16/6 = 8/3$ 次。

\subsection{最壞情況的下界}

\begin{theorem}
\label{thm:lin-lower}
任何比較排序在最壞情況下至少要做 $\lceil \lg(n!) \rceil$ 次比較。 因此比較排序的
最壞情況執行時間是 $\Omega(n \lg n)$。
\end{theorem}

\begin{proof}
固定 $n$，看這個演算法的決策樹。

\textbf{第一步：$n!$ 種排列走到 $n!$ 個不同的葉子。} 輸入是 $n$ 個不同的鍵，它們的
大小順序有 $n!$ 種。 假設兩種不同的順序 $\pi \ne \sigma$ 走到同一個葉子，演算法在
兩種輸入上看到完全一樣的比較結果，就會做完全一樣的搬動，輸出同一種重排。 但同一
種重排不可能同時把 $\pi$ 和 $\sigma$ 兩種順序都排好，所以其中一種會排錯。 正確的
演算法因此至少有 $n!$ 個（走得到的）葉子。

\textbf{第二步：深度都不超過 $d$ 的二元樹最多 $2^d$ 個葉子。} 對 $d$ 歸納：$d = 0$
時樹只有根，$1$ 個葉子。 $d \ge 1$ 時，根的左右子樹裡每個葉子的深度都不超過
$d - 1$，各自最多 $2^{d-1}$ 個葉子，合起來最多 $2^d$ 個。

設最壞情況的比較次數是 $d$，也就是最深的葉子深度是 $d$。 由兩步得到 $2^d \ge n!$，
$d \ge \lg(n!)$，$d$ 是整數，所以 $d \ge \lceil \lg(n!) \rceil$。 第~1~章 1.6 節已經
證明 $\lg(n!) \ge \frac{n}{2}\lg\frac{n}{2}$，$n \ge 4$ 時這至少是 $\frac{n}{4}\lg n$，
所以 $d = \Omega(n \lg n)$。 每次比較至少花常數時間，執行時間也是 $\Omega(n \lg n)$。
\end{proof}

\begin{corollary}
\label{cor:lin-optimal}
合併排序和堆積排序最壞情況是 $O(n \lg n)$，所以在比較排序裡是\textbf{漸近最佳}的：
不可能有比較排序在最壞情況下比它們快超過常數倍。
\end{corollary}

快速排序最壞是 $\Theta(n^2)$，不在這個系理裡；它的平均是 $\Theta(n \lg n)$，下面
的平均情況下界說明這個平均也已經是最好的了。

\begin{examnote}
本書的高度數節點（第~7~章），決策樹的「最壞比較次數」是最深葉子的\textbf{深度}，
也就是樹高減 $1$。 英文書常寫「決策樹的高度 $h$ 滿足 $2^h \ge n!$」，那裡的高度數
的是邊。 作答時直接寫「最壞比較次數 $d$ 滿足 $2^d \ge n!$」最不容易出錯；題目問
「決策樹的高度」時，先看題目怎麼定義高度。
\end{examnote}

對小的 $n$，$\lceil \lg(n!) \rceil$ 可以直接算：$n = 4, 5, 8, 12$ 時分別是
$5, 7, 16, 29$。 合併排序最壞的比較次數是 $n\lceil \lg n \rceil - 2^{\lceil \lg n
\rceil} + 1$，同樣這四個 $n$ 是 $5, 8, 17, 33$，和下界只差一點；插入排序最壞是
$n(n-1)/2$，$n = 12$ 時已經是 $66$。 下界不一定達得到：已知 $n \le 11$ 時最少比較次數剛好是
$\lceil \lg(n!) \rceil$（例如五個數最少比 $7$ 次），但 $n = 12$ 時最少要 $30$ 次，比
下界多 $1$（Knuth，TAOCP 第 3 卷 5.3.1 節）。

\subsection{平均情況的下界}

最壞情況的下界只說「有某些輸入很慢」。 會不會有一個比較排序，對大部分輸入都很快，
只有少數輸入很慢？ 下面說明不會：連平均都要 $\lg(n!)$ 次。

\begin{theorem}
\label{thm:lin-avg}
$n!$ 種大小順序出現的機率都相同時，任何比較排序的平均比較次數至少是 $\lg(n!)$。
\end{theorem}

\begin{whybox}[平均深度至少 $\lg L$]
先證一個關於二元樹的不等式：設一棵二元樹的葉子深度是 $d_1, d_2, \ldots, d_L$，則
\[
\sum_{i=1}^{L} 2^{-d_i} \le 1.
\]
對樹的節點數歸納。 只有根時 $L = 1$、$d_1 = 0$，和是 $1$。 否則根的每個子樹裡，葉子
在子樹中的深度比在整棵樹中少 $1$；由歸納假設，每個子樹的葉子在整棵樹裡的
$\sum 2^{-d_i}$ 不超過 $\frac{1}{2}$，最多兩個子樹，合起來不超過 $1$。

現在取決策樹裡 $n!$ 種順序各自走到的 $L = n!$ 個葉子，它們的深度 $d_1, \ldots, d_L$
仍然滿足上面的不等式（少算一些葉子，和只會變小）。 對正數 $2^{-d_i}$ 用算幾不等式：
\[
\frac{1}{L} \ge \frac{1}{L}\sum_{i=1}^{L} 2^{-d_i}
\ge \left( \prod_{i=1}^{L} 2^{-d_i} \right)^{1/L}
= 2^{-\bar{d}}, \qquad \bar{d} = \frac{1}{L}\sum_{i=1}^{L} d_i.
\]
兩邊取 $\lg$，得到 $-\lg L \ge -\bar{d}$，也就是平均深度 $\bar{d} \ge \lg L = \lg(n!)$。
每種順序機率 $1/n!$，平均比較次數就是 $\bar{d}$。
\end{whybox}

圖~\ref{fig:lin-dtree} 的平均是 $8/3 \approx 2.67$，下界是 $\lg 6 \approx 2.58$，差
得不多。 用第~1~章的估計，平均情況一樣是 $\Omega(n \lg n)$。 所以快速排序的平均
$\Theta(n \lg n)$ 也是最佳的。 隨機化也突破不了：把隨機選擇固定下來就是一棵決策樹，
由此可以推出，任何隨機化比較排序都有某個輸入，在它上面的期望比較次數至少 $\lg(n!)$。

\subsection{同樣的方法用在別的問題}

決策樹的論證不只用在排序。 只要一個問題的答案有 $L$ 種可能，每次比較只有兩種結
果，最壞就至少要 $\lceil \lg L \rceil$ 次比較。 這叫\textbf{資訊理論下界}
（information-theoretic lower bound）。 它有時很準，有時很鬆，考題喜歡考兩者的差別。

\textbf{合併兩個排好的串列。} $A$ 有 $m$ 個數、$B$ 有 $n$ 個數，各自排好，所有數兩
兩不同。 合併的結果由「最後的序列裡，哪 $m$ 個位置屬於 $A$」決定，共有
$\binom{m+n}{m}$ 種可能，所以至少要 $\lceil \lg \binom{m+n}{m} \rceil$ 次比較。
\begin{itemize}
\item $m = 1$：至少 $\lceil \lg(n+1) \rceil$ 次。 對 $B$ 做二分搜尋找插入位置剛好
達到，所以這時下界是準的，而一般的合併（逐一往後比）最壞要 $n$ 次，差很多。
\item $m = n$：$\binom{2n}{n} \ge \frac{4^n}{2n+1}$（$\binom{2n}{n}$ 是
$(1+1)^{2n}$ 展開的 $2n+1$ 項裡最大的一項），所以下界至少是
$2n - \lg(2n+1)$。 一般的合併最壞比 $2n - 1$ 次，看起來還差 $\lg n$ 左右。
\end{itemize}

$m = n$ 的情況，其實下界可以提高到剛好 $2n - 1$，用的是另一種論證：
\textbf{對手論證}（adversary argument）。

\begin{whybox}[合併兩個長度 $n$ 的串列至少要比 $2n - 1$ 次]
看這樣的輸入：合併後的順序是 $a_1 < b_1 < a_2 < b_2 < \cdots < a_n < b_n$。 這個序列
裡相鄰的兩個數一共有 $2n - 1$ 對，每一對都來自不同的串列。

說明每一對相鄰的數都一定要直接比過。 假設某個演算法從來沒有比過某一對相鄰的
$x < y$（一個來自 $A$、一個來自 $B$）。 把 $y$ 的值改成介於 $x$ 前一個數和 $x$ 之
間，其他數不動。 $x$ 和 $y$ 交換了先後；$y$ 在自己串列裡的前一個數排在 $x$ 之前，所
以 $y$ 的串列仍然排好；除了 $x$、$y$ 這一對，任何兩個數的大小關係都沒變。 演算法
沒有比過 $x$、$y$，它看到的每次比較結果都和原來一樣，所以輸出也一樣，但這兩種輸
入的正確答案不同，至少有一種錯。 所以 $2n - 1$ 對都要比過，每次比較只能比到一對，
至少 $2n - 1$ 次。
\end{whybox}

例如 $n = 5$：$\lceil \lg \binom{10}{5} \rceil = \lceil \lg 252 \rceil = 8$，對手論
證給出 $9$，一般的合併最壞也是 $9$ 次，所以一般的合併在 $m = n$ 時是最佳的。 觀念
題第 3 題的「找最大值」是另一個資訊理論下界很鬆的例子。


\subsection{下界說了什麼、沒說什麼}

定理~\ref{thm:lin-lower} 只管\textbf{只用比較得知鍵的資訊}的演算法。 如果鍵是整
數，程式可以拿鍵當陣列的索引（例如第~1~章 1.1 節查表找重複學號的做法），一次就得
知「這個鍵是幾」，這不是一次兩種結果的比較，決策樹的論證就不適用。 下一節的計數
排序、基數排序、桶子排序都是這樣做的。 代價是對輸入多了假設：鍵是範圍不大的整
數，或是分布已知的實數。

\section{線性時間排序}
\label{sec:lin-linear}

\subsection{計數排序}
\label{sec:lin-counting}

假設每筆紀錄的鍵是 $0$ 到 $k$ 的整數。 如果知道「鍵小於等於 $v$ 的紀錄有 $c$ 筆」，
那麼鍵是 $v$ 的紀錄裡最後一筆，排好之後就在第 $c$ 格。 \textbf{計數排序}（counting
sort）就是先數出每個鍵出現幾次，再做前綴和，得到每個鍵的最後位置，最後把紀錄一筆
一筆放進輸出陣列。

\begin{pseudo}[label=alg:lin-counting]{計數排序（穩定）}
\Procedure{Counting-Sort}{$A, B, n, k$}
    \State 建立陣列 $C[0..k]$，每一格都設成 $0$
    \For{$j \gets 1$ \To $n$}
        \State $C[A[j].\mathit{key}] \gets C[A[j].\mathit{key}] + 1$
    \EndFor
    \For{$v \gets 1$ \To $k$}
        \State $C[v] \gets C[v] + C[v - 1]$ \Comment{$C[v]$ = 鍵 $\le v$ 的筆數}
    \EndFor
    \For{$j \gets n$ \DownTo $1$}
        \State $v \gets A[j].\mathit{key}$
        \State $B[C[v]] \gets A[j]$
        \State $C[v] \gets C[v] - 1$
    \EndFor
\EndProcedure
\end{pseudo}

\textbf{手算追蹤。} 八筆紀錄，鍵在 $0$ 到 $4$，每筆附一個字母，寫成 $4_a$ 表示「鍵
$4$、字母 $a$」，字母照輸入的順序排，用來看鍵相同的紀錄有沒有保持原來的先後：
\[
A = \langle 4_a, 1_b, 3_c, 4_d, 0_e, 1_f, 4_g, 2_h \rangle .
\]
第 3--4 行數完之後 $C = [1, 2, 1, 1, 3]$（$C[0]$ 到 $C[4]$），第 5--6 行做完前綴和之
後 $C = [1, 3, 4, 5, 8]$。 例如 $C[1] = 3$：鍵 $\le 1$ 的有 $0_e, 1_b, 1_f$ 三筆，所以
最後一個鍵 $1$ 的紀錄要放在第 $3$ 格。 接著第 7--10 行從右往左放：

\begin{center}
\begin{tabular}{cccl l}
\toprule
$j$ & $A[j]$ & 放到 $B$ 的第幾格 & 放完後的 $C[0..4]$ & 放完後的 $B[1..8]$ \\
\midrule
8 & $2_h$ & 4 & 1 3 3 5 8 & -- -- -- $2_h$ -- -- -- -- \\
7 & $4_g$ & 8 & 1 3 3 5 7 & -- -- -- $2_h$ -- -- -- $4_g$ \\
6 & $1_f$ & 3 & 1 2 3 5 7 & -- -- $1_f$ $2_h$ -- -- -- $4_g$ \\
5 & $0_e$ & 1 & 0 2 3 5 7 & $0_e$ -- $1_f$ $2_h$ -- -- -- $4_g$ \\
4 & $4_d$ & 7 & 0 2 3 5 6 & $0_e$ -- $1_f$ $2_h$ -- -- $4_d$ $4_g$ \\
3 & $3_c$ & 5 & 0 2 3 4 6 & $0_e$ -- $1_f$ $2_h$ $3_c$ -- $4_d$ $4_g$ \\
2 & $1_b$ & 2 & 0 1 3 4 6 & $0_e$ $1_b$ $1_f$ $2_h$ $3_c$ -- $4_d$ $4_g$ \\
1 & $4_a$ & 6 & 0 1 3 4 5 & $0_e$ $1_b$ $1_f$ $2_h$ $3_c$ $4_a$ $4_d$ $4_g$ \\
\bottomrule
\end{tabular}
\end{center}

結果裡三個鍵 $4$ 的順序是 $4_a, 4_d, 4_g$，兩個鍵 $1$ 是 $1_b, 1_f$，都和輸入的先後一
樣。 鍵相同的紀錄排完後保持原來的相對順序，這叫\textbf{穩定}（stable，第~5~章合併兩
條串列時已經見過，第~14~章有整理）。 計數排序是穩定的，原因是放的順序：$C[v]$ 一開
始指向「鍵 $v$ 的最後一格」，從右往左掃，先遇到的是鍵 $v$ 裡最後一筆，放在最後一
格，之後每放一筆就往左退一格。 如果第 7 行改成從左往右掃、其他不變，鍵相同的紀錄
會整個顛倒過來（$4_g, 4_d, 4_a$），排序還是對的，但不穩定了。

\textbf{成本。} 第 2 行 $\Theta(k)$，第 3--4 行 $\Theta(n)$，第 5--6 行
$\Theta(k)$，第 7--10 行 $\Theta(n)$，合起來 $\Theta(n + k)$。 $k = O(n)$ 時就是
$\Theta(n)$。 它需要另外 $\Theta(n + k)$ 的空間（陣列 $B$ 和 $C$），不是原地排序。
$k$ 比 $n$ 大很多時就不划算了：排 $1000$ 個九位數學號，$C$ 要開十億格。 下一小節的
基數排序解決這個問題。

\begin{cbookrecap}[第~14~章 結構、第~19~章 19.3 節]
\texttt{struct Rec \{ int key; char tag; \};} 把鍵和附帶的資料包成一筆紀錄。 結構可
以整個指定，\texttt{b[c[v] - 1] = a[j];} 會把兩個欄位一起複製。
\texttt{malloc(n * sizeof *b)} 要一塊放得下 $n$ 筆紀錄的記憶體，用完要
\texttt{free(b)}。 初始化 \texttt{int c[K + 1] = \{0\};} 會把整個陣列設成 $0$。
\end{cbookrecap}

下面的 C 程式就是虛擬碼~\ref{alg:lin-counting}，陣列改成從 $0$ 開始：前綴和之後
$c[v]$ 仍然是「鍵 $\le v$ 的筆數」，所以放進 \texttt{b[c[v] - 1]}。

\begin{example}[label=ex:lin-counting]{c}{(counting\_sort.c) 穩定的計數排序}
#include <stdio.h>
#include <stdlib.h>

#define K 4                       // 鍵的範圍是 0..K

struct Rec {
    int key;
    char tag;                     // 附帶的資料，用來看穩定性
};

// 穩定的計數排序：a[0..n-1] 依 key 排好後放進 b[0..n-1]。
void counting_sort(const struct Rec *a, struct Rec *b, int n)
{
    int c[K + 1] = {0};
    for (int j = 0; j < n; j++)          // 數每個鍵出現幾次
        c[a[j].key]++;
    for (int v = 1; v <= K; v++)         // 前綴和：鍵 <= v 的筆數
        c[v] += c[v - 1];
    for (int j = n - 1; j >= 0; j--) {   // 從右往左放，才會穩定
        int v = a[j].key;
        b[c[v] - 1] = a[j];              // c[v] 是從 1 數的格子
        c[v]--;
    }
}

int main(void)
{
    struct Rec a[] = {{4, 'a'}, {1, 'b'}, {3, 'c'}, {4, 'd'},
                      {0, 'e'}, {1, 'f'}, {4, 'g'}, {2, 'h'}};
    int n = (int)(sizeof a / sizeof a[0]);
    struct Rec *b = malloc(n * sizeof *b);
    if (b == NULL)
        return 1;
    counting_sort(a, b, n);
    printf("result:");
    for (int i = 0; i < n; i++)
        printf(" %d%c", b[i].key, b[i].tag);
    printf("\n");
    free(b);
    return 0;
}
\end{example}
\iosample{}{result: 0e 1b 1f 2h 3c 4a 4d 4g}

\subsection{LSD 基數排序}
\label{sec:lin-lsd}

學號、郵遞區號、日期這類鍵是好幾位數字組成的。 \textbf{基數排序}（radix sort）每次
只看其中一位，用計數排序排一趟，每一趟的計數陣列只有 $10$ 格（十進位）。
\textbf{LSD}（least significant digit first）的做法是從最低位開始：先依個位排，再依
十位排，最後依最高位排。

\begin{pseudo}[label=alg:lin-radix]{LSD 基數排序}
\Procedure{Radix-Sort}{$A, n, d$}
    \For{$i \gets 1$ \To $d$} \Comment{第 $1$ 位是個位}
        \State 用\textbf{穩定}的排序，依第 $i$ 位排 $A[1..n]$
    \EndFor
\EndProcedure
\end{pseudo}

\textbf{手算追蹤。} 十個三位數學號（末三碼）：
$318, 542, 207, 613, 245, 731, 118, 402, 659, 236$。 每一趟一列，畫底線的是這一趟依
據的那一位。

\begin{center}
\small
\setlength{\tabcolsep}{4.5pt}
\begin{tabular}{l*{10}{c}}
\toprule
輸入 & 318 & 542 & 207 & 613 & 245 & 731 & 118 & 402 & 659 & 236 \\
\midrule
第 1 趟（個位） & 73\underline{1} & 54\underline{2} & 40\underline{2} & 61\underline{3} & 24\underline{5} & 23\underline{6} & 20\underline{7} & 31\underline{8} & 11\underline{8} & 65\underline{9} \\
第 2 趟（十位） & 4\underline{0}2 & 2\underline{0}7 & 6\underline{1}3 & 3\underline{1}8 & 1\underline{1}8 & 7\underline{3}1 & 2\underline{3}6 & 5\underline{4}2 & 2\underline{4}5 & 6\underline{5}9 \\
第 3 趟（百位） & \underline{1}18 & \underline{2}07 & \underline{2}36 & \underline{2}45 & \underline{3}18 & \underline{4}02 & \underline{5}42 & \underline{6}13 & \underline{6}59 & \underline{7}31 \\
\bottomrule
\end{tabular}
\end{center}

看第二趟（十位）：$613, 318, 118$ 的十位都是 $1$，它們照第一趟結束時的先後出現；第一
趟已經依個位排好它們（$3 < 8$，而 $318$、$118$ 個位相同，保持輸入的先後），所以第
二趟結束後，每一段十位相同的數都依個位排好了。 第三趟同理：$207, 236, 245$ 百位都
是 $2$，保持第二趟的順序，而第二趟的順序已經依「十位、個位」排好。

\textbf{為什麼內層一定要穩定。} 假設第三趟用不穩定的排序。 $613$ 和 $659$ 的百位都是
$6$，不穩定的排序可能把 $659$ 放在 $613$ 前面，前兩趟依十位、個位排好的成果就沒
了，最後的結果是錯的。 LSD 的每一趟都要\textbf{保住前面幾趟的成果}，靠的就是穩定。

\begin{whybox}[LSD 基數排序的正確性]
對趟數 $i$ 歸納證明：第 $i$ 趟結束後，陣列依「最低的 $i$ 位」組成的數排好了。

$i = 1$：第一趟依個位排，成立。 假設第 $i - 1$ 趟結束後依最低 $i - 1$ 位排好。 第
$i$ 趟依第 $i$ 位排。 任取兩個數 $x$、$y$，看它們最低 $i$ 位的大小：
\begin{itemize}
\item 第 $i$ 位不同：最低 $i$ 位的大小由第 $i$ 位決定，第 $i$ 趟把第 $i$ 位小的放在前
面，順序正確。
\item 第 $i$ 位相同：最低 $i$ 位的大小由最低 $i - 1$ 位決定。 由歸納假設，第 $i$ 趟開
始時它們已經依最低 $i - 1$ 位排好；這一趟的排序穩定，兩者第 $i$ 位相同，所以保持原
來的先後，順序仍然正確。
\end{itemize}
第 $d$ 趟結束後依全部 $d$ 位排好，也就是排好了。 證明只有第二種情況用到穩定。
\end{whybox}

\textbf{成本。} $d$ 趟，每趟用計數排序，每一位有 $b$ 種可能的值（十進位 $b = 10$），
一趟 $\Theta(n + b)$，共 $\Theta(d(n + b))$。 $d$ 是常數、$b = O(n)$ 時是
$\Theta(n)$。

下面的程式排上面那十個學號，每一趟用一次穩定的計數排序（\texttt{sort\_by\_digit}），
\texttt{exp} 依序是 $1, 10, 100$，\texttt{a[j] / exp \% 10} 取出那一位。 計數陣列只有
$10$ 格，直接開在函式裡；暫存的陣列 \texttt{tmp} 和資料一樣長，用 \texttt{malloc}。

\begin{example}[label=ex:lin-radix]{c}{(radix\_lsd.c) LSD 基數排序，印出每一趟}
#include <stdio.h>
#include <stdlib.h>

// 依 a[j] / exp % 10 這一位做一次穩定的計數排序，結果寫回 a。
void sort_by_digit(int *a, int *tmp, int n, int exp)
{
    int c[10] = {0};
    for (int j = 0; j < n; j++)
        c[a[j] / exp % 10]++;
    for (int v = 1; v < 10; v++)
        c[v] += c[v - 1];
    for (int j = n - 1; j >= 0; j--)
        tmp[--c[a[j] / exp % 10]] = a[j];
    for (int j = 0; j < n; j++)
        a[j] = tmp[j];
}

// LSD 基數排序：a[0..n-1] 都是不超過 d 位的非負整數。
void radix_sort(int *a, int n, int d)
{
    int *tmp = malloc(n * sizeof *tmp);
    if (tmp == NULL)
        exit(1);
    int exp = 1;
    for (int pass = 1; pass <= d; pass++, exp *= 10) {
        sort_by_digit(a, tmp, n, exp);
        printf("pass %d:", pass);
        for (int i = 0; i < n; i++)
            printf(" %03d", a[i]);
        printf("\n");
    }
    free(tmp);
}

int main(void)
{
    int a[] = {318, 542, 207, 613, 245, 731, 118, 402, 659, 236};
    radix_sort(a, (int)(sizeof a / sizeof a[0]), 3);
    return 0;
}
\end{example}
\iosample{}{pass 1: 731 542 402 613 245 236 207 318 118 659 \\ pass 2: 402 207 613 318 118 731 236 542 245 659 \\ pass 3: 118 207 236 245 318 402 542 613 659 731}

\subsection{一位要切多大}

「一位」不一定是十進位的一位。 鍵是 $b$ 個位元的非負整數時，可以每 $r$ 個位元當成
一位，共 $\lceil b/r \rceil$ 趟，每趟的計數陣列有 $2^r$ 格，總時間
\[
\Theta\!\left( \frac{b}{r}\,(n + 2^r) \right).
\]
$r$ 小，趟數多；$r$ 大，每趟的計數陣列大。 $2^r$ 不超過 $n$ 時，加大 $r$ 只會讓趟數變
少；$2^r$ 超過 $n$ 很多以後，$2^r$ 那一項會壓過 $n$。 所以 $r$ 取 $\lg n$ 附近最好，
時間是 $\Theta(bn / \lg n)$。

例如 $n = 10^6$ 個 $32$ 位元整數：$r = 8$ 要 $4$ 趟，每趟 $256$ 格；$r = 16$ 要 $2$
趟，每趟 $65536$ 格，仍比 $n$ 小，所以 $r = 16$ 比較快。 實際的程式還要考慮快取：計
數陣列太大時，每次 \texttt{c[v]++} 都可能找不到快取，所以常見的選擇是 $r = 8$ 或
$11$。

另一個常用的結論：$n$ 個鍵都在 $0$ 到 $n^c - 1$（$c$ 是常數）時，以 $n$ 為基底，每個
鍵是 $c$ 位的 $n$ 進位數，基數排序 $c$ 趟、每趟 $\Theta(n + n)$，總共 $\Theta(cn)$。
直接用計數排序要 $\Theta(n^c)$，差很多。

\textbf{基數排序一定比較快嗎？} 不一定。 它不是原地排序，每一趟都要把所有資料搬一
次，每趟還要掃好幾遍陣列；$d \lg n$ 不比 $\lg n$ 小多少的時候，快速排序常常更快。 基
數排序適合鍵短而且固定長度、資料量大的情況。

\subsection{MSD 基數排序}
\label{sec:lin-msd}

\textbf{MSD}（most significant digit first）從最高位開始：先依最高位把資料分成 $10$
個桶子，每個桶子內部再依下一位遞迴地分，桶子只剩一筆（或位數用完）就停。 桶子之間
的先後已經由最高位決定，不會再動，所以不需要後面幾趟去「保住」前面的成果。

用上一小節的十個學號：依百位分成 $\{118\}$、$\{207, 245, 236\}$、$\{318\}$、
$\{402\}$、$\{542\}$、$\{613, 659\}$、$\{731\}$。 只有一筆的桶子不用再分。 百位
是 $2$ 的桶子依十位再分成 $\{207\}$、$\{236\}$、$\{245\}$；百位是 $6$ 的桶子依十位分
成 $\{613\}$、$\{659\}$。 全部只剩一筆，結束。 個位一次都沒看。

MSD 和 LSD 的比較：

\begin{itemize}
\item MSD 可以提早停：桶子小到只剩一筆就不用再看後面的位。 這讓它很適合排\textbf{長
度不一的字串}：依第一個字元分桶，字串在某一位結束時，把它放在「比所有字元都小」的
桶子，短的自然排在前面，這就是字典序。
\item MSD 的遞迴會產生很多小桶子，每個桶子都要開一個計數陣列，小桶子的額外成本很高。
實作時桶子小於某個大小（例如 $16$ 筆）就改用插入排序。
\item LSD 每一趟都掃過全部資料，程式簡單，不用遞迴，但一定要處理完全部 $d$ 位，而且
每一趟都要穩定。 MSD 的每一層分桶如果用計數排序，本身就是穩定的；附帶的資料要保
持原來的先後時才需要在意這一點。
\end{itemize}

\subsection{桶子排序}
\label{sec:lin-bucket}

前面兩種排序假設鍵是整數。 如果鍵是 $[0, 1)$ 裡的實數，而且\textbf{均勻分布}，可以把
$[0, 1)$ 切成 $n$ 個等長的區間當桶子，每個數丟進它所在區間的桶子。 均勻分布時，每個
桶子平均只有一個數，各自排好再依序串起來就好。 這叫\textbf{桶子排序}（bucket
sort）。

\begin{pseudo}[label=alg:lin-bucket]{桶子排序}
\Procedure{Bucket-Sort}{$A, n$}
    \State 建立 $n$ 個空串列 $Q[0..n-1]$
    \For{$i \gets 1$ \To $n$}
        \State 把 $A[i]$ 接到串列 $Q[\lfloor n \cdot A[i] \rfloor]$ 的尾端
    \EndFor
    \For{$j \gets 0$ \To $n - 1$}
        \State 用插入排序排串列 $Q[j]$
    \EndFor
    \State 依序把 $Q[0], Q[1], \ldots, Q[n-1]$ 串起來，寫回 $A$
\EndProcedure
\end{pseudo}

例如 $n = 10$，$A = \langle 0.47, 0.83, 0.11, 0.65, 0.19, 0.58, 0.04, 0.91, 0.43,
0.15 \rangle$。 第 3--4 行之後，非空的桶子是 $Q[0] = \langle 0.04 \rangle$、
$Q[1] = \langle 0.11, 0.19, 0.15 \rangle$、$Q[4] = \langle 0.47, 0.43 \rangle$，以及
$Q[5]$、$Q[6]$、$Q[8]$、$Q[9]$ 各一個數。 $Q[1]$ 排成 $0.11, 0.15, 0.19$，$Q[4]$ 排成 $0.43, 0.47$，再依序串起來就排好了。 桶子
$j$ 裡的數都在 $[j/n, (j+1)/n)$，比桶子 $j + 1$ 的數小，所以串起來的順序是對的。

\textbf{期望時間。} 除了第 6 行，其他都是 $\Theta(n)$。 設桶子 $j$ 最後有 $n_j$ 個
數，插入排序最壞 $O(n_j^2)$，所以總時間是 $\Theta(n) + \sum_j O(n_j^2)$，只要算
$E[n_j^2]$。 每個數獨立地以機率 $1/n$ 落進桶子 $j$，所以 $n_j$ 是二項分布
$B(n, 1/n)$：
\[
E[n_j] = n \cdot \frac{1}{n} = 1, \qquad
\mathrm{Var}[n_j] = n \cdot \frac{1}{n}\left(1 - \frac{1}{n}\right) = 1 - \frac{1}{n},
\]
\[
E[n_j^2] = \mathrm{Var}[n_j] + (E[n_j])^2 = 2 - \frac{1}{n}.
\]
由期望值的線性，總時間的期望值是 $\Theta(n) + n \cdot O(2 - 1/n) = \Theta(n)$。

這裡的「期望」是對輸入的分布取的。 輸入不均勻時（例如全部擠在 $[0, 0.1)$），一個桶
子可能分到 $n$ 個數，最壞是 $\Theta(n^2)$。 期望時間只要求每個桶子的
$E[n_j^2] = O(1)$，桶子數改成 $n$ 的常數倍也可以；桶子數太少就不行了，見觀念題第 6
題。


\section{外部排序}
\label{sec:lin-external}

\subsection{模型與兩個階段}

回到章首的選課紀錄：$600$ MB 的資料、$60$ MB 的記憶體。 資料大部分時間放在磁碟上，
排序的過程中要一塊一塊讀進來、寫出去。 第~11~章 11.3 節說過，磁碟以\textbf{區塊}為
單位讀寫，讀一個區塊比在記憶體裡比較一次慢好幾個數量級，所以\textbf{外部排序}
（external sorting）的成本主要看\textbf{讀寫了幾個區塊}。 這一節用下面的符號：

\begin{itemize}
\item $N$：要排序的紀錄筆數。
\item $M$：記憶體一次放得下的紀錄筆數。
\item $B$：一個區塊放幾筆紀錄。 整個檔案有 $N/B$ 個區塊。
\end{itemize}

外部排序幾乎都用合併排序的想法，分兩個階段。

\begin{enumerate}
\item \textbf{產生初始串}：每次讀進 $M$ 筆，在記憶體裡用任何一種內部排序（第~14~
章）排好，寫回磁碟。 寫出去的每一段排好的資料叫一個\textbf{串}（run）。 這樣得到
$r = \lceil N/M \rceil$ 個串，每個（最後一個除外）長 $M$。
\item \textbf{合併}：每次取 $k$ 個串，合併成一個長串寫回磁碟。 所有的串都這樣合過
一輪叫一\textbf{趟}（pass），一趟之後串的個數大約變成原來的 $1/k$。 重複到只剩一個
串。
\end{enumerate}

合併的時候，記憶體裡只要替每個輸入的串留一個區塊大小的\textbf{緩衝區}，再留一個給
輸出：某個輸入緩衝區用完就從那個串讀下一個區塊，輸出緩衝區滿了就寫出去。 這也說明
$k$ 不能無限大：$k + 1$ 個緩衝區都要放得進記憶體，$k \le M/B - 1$。

\begin{center}
\begin{tikzpicture}[x=1cm, y=1cm]
\foreach \i in {1,...,9} {
    \node[lrunbox, minimum width=0.95cm] (r\i) at ({(\i-1)*1.25}, 0) {串 \i};
}
\foreach \j in {1,2,3} {
    \node[lrunbox, minimum width=2.6cm, fill=coreFill, draw=coreStroke] (s\j) at ({(\j-1)*3.75 + 1.25}, -1.3) {長 $3M$};
}
\node[lrunbox, minimum width=5cm, fill=greenFill, draw=greenStroke] (t) at (5, -2.6) {長 $9M$，排好的檔案};
\foreach \i in {1,2,3} { \draw[lte] (r\i.south) -- (s1.north); }
\foreach \i in {4,5,6} { \draw[lte] (r\i.south) -- (s2.north); }
\foreach \i in {7,8,9} { \draw[lte] (r\i.south) -- (s3.north); }
\foreach \j in {1,2,3} { \draw[lte] (s\j.south) -- (t.north); }
\node[font=\small, anchor=west] at (10.6, 0) {初始串，$r = 9$};
\node[font=\small, anchor=west] at (10.6, -1.3) {第 1 趟};
\node[font=\small, anchor=west] at (10.6, -2.6) {第 2 趟};
\end{tikzpicture}
\diagcaption{$r = 9$ 個初始串用 $3$ 路合併，$\lceil \log_3 9 \rceil = 2$ 趟合併完。}
\label{fig:lin-passes}
\end{center}

\subsection{合併的趟數與 I/O 次數}
\label{sec:lin-passes}

每一趟串數變成 $\lceil \text{串數}/k \rceil$，從 $r$ 個串合到剩一個，需要
\[
\text{合併趟數} = \lceil \log_k r \rceil
\]
趟。 每一趟都把整個檔案讀一次、寫一次，是 $2N/B$ 次區塊 I/O；產生初始串也一樣讀一
次、寫一次。 所以
\[
\text{總 I/O 次數} = \frac{2N}{B}\left(1 + \lceil \log_k r \rceil\right).
\]

\textbf{例 1.} 章首的數字：$N = 1{,}200{,}000$ 筆，每筆 $500$ 位元組；記憶體 $60$ MB，
$M = 120{,}000$ 筆；區塊 $8000$ 位元組，$B = 16$ 筆，檔案共 $N/B = 75{,}000$ 個區
塊。 初始串 $r = 10$ 個。

\begin{center}
\begin{tabular}{cccc}
\toprule
$k$ & 合併趟數 $\lceil \log_k 10 \rceil$ & 含產生初始串的總趟數 & 總 I/O 次數 \\
\midrule
2 & 4 & 5 & $5 \times 150{,}000 = 750{,}000$ \\
3 & 3 & 4 & $600{,}000$ \\
4 & 2 & 3 & $450{,}000$ \\
10 & 1 & 2 & $300{,}000$ \\
\bottomrule
\end{tabular}
\end{center}

$2^3 = 8 < 10$，所以 $2$ 路要 $4$ 趟；$4^2 = 16 \ge 10$，$4$ 路只要 $2$ 趟。 從
$2$ 路換成 $10$ 路，I/O 少了六成。

\begin{examnote}
外部排序的計算題，先寫清楚三件事：(1) 初始串有幾個（$\lceil N/M \rceil$，或取代選擇
的約 $N/(2M)$）；(2) 「趟數」有沒有包含產生初始串的那一趟，題目沒說就兩個都寫出來；
(3) 一趟的 I/O 是讀加寫，每個區塊算兩次，題目只問讀或只問寫時要除以 $2$。 對數一律取
天花板，$\log_k r$ 算出 $1.43$ 就是 $2$ 趟，不能寫 $1.43$ 趟。
\end{examnote}

\textbf{$k$ 越大越好嗎？} 從 I/O 次數看是。 但有兩個代價。 第一，記憶體固定時，$k$ 越
大，每個緩衝區越小，讀同樣多的資料要找更多次磁碟的位置（傳統硬碟找一次位置要幾毫
秒）。 第二，每輸出一筆都要從 $k$ 個候選裡挑最小的；直接一個一個比要 $k - 1$ 次，$k$
大的時候 CPU 會跟不上。 第二個問題可以用下面的勝者樹、敗者樹解決：每輸出一筆只要
$\lceil \lg k \rceil$ 次比較。 這時每筆紀錄在整個合併階段被比較的次數大約是
\[
\lceil \log_k r \rceil \cdot \lceil \lg k \rceil \approx \log_k r \cdot \lg k = \lg r,
\]
和 $k$ 幾乎無關。 加大 $k$ 只減少 I/O，不增加比較，這正是要用樹來挑最小的原因。

\subsection{從 $k$ 個串裡挑最小：勝者樹}

把 $k$ 個串目前排在最前面的鍵想成 $k$ 位選手，兩兩比賽，鍵小的贏。 用一棵 $k$ 個葉
子的完全二元樹排賽程：葉子是選手，內部節點是一場比賽。 編號和第~8~章的堆積一樣，內
部節點是 $1$ 到 $k - 1$，節點 $t$ 的子節點是 $2t$、$2t + 1$；葉子是 $k$ 到 $2k - 1$，
葉子 $k + i - 1$ 代表第 $i$ 個串 $R_i$。 $k$ 不是 $2$ 的次方時，葉子分在最下面兩層，
這沒有關係。

\textbf{勝者樹}（winner tree）在每個內部節點記下那場比賽的\textbf{贏家}，根就是總冠
軍，也就是下一個要輸出的鍵。 輸出之後，冠軍的串補上下一個鍵，然後沿著它的葉子往上
重賽：每一層都要找出\textbf{兄弟節點}的贏家，和自己這邊的贏家比一次，記下新的贏
家。 下圖用四個串：
\[
R_1 = \langle 15, 42, 63 \rangle,\quad
R_2 = \langle 27, 31, 80 \rangle,\quad
R_3 = \langle 8, 56, 71 \rangle,\quad
R_4 = \langle 19, 22, 95 \rangle.
\]
內部節點的數字是贏家的鍵，旁邊小字是它來自哪個串；葉子下面的小字是那個串還沒讀
的部分。

\begin{center}
\begin{tabular}{c@{\hspace{1.2cm}}c}
\begin{tikzpicture}[x=0.70cm,y=1.00cm]
\draw[lte] (3,-1) -- (1,-2);
\draw[lte] (3,-1) -- (5,-2);
\draw[lte] (1,-2) -- (0,-3);
\draw[lte] (1,-2) -- (2,-3);
\draw[lte] (5,-2) -- (4,-3);
\draw[lte] (5,-2) -- (6,-3);
\node[ltn] at (3,-1) {$8$};
\node[rlab, anchor=west] at (3.51429,-0.95) {R3};
\node[ltn] at (1,-2) {$15$};
\node[rlab, anchor=west] at (1.51429,-1.95) {R1};
\node[ltn] at (5,-2) {$8$};
\node[rlab, anchor=west] at (5.51429,-1.95) {R3};
\node[ltleaf] at (0,-3) {$15$};
\node[rlab] at (0,-3.5) {R1};
\node[restlab] at (0,-3.85) {42 63};
\node[ltleaf] at (2,-3) {$27$};
\node[rlab] at (2,-3.5) {R2};
\node[restlab] at (2,-3.85) {31 80};
\node[ltleaf] at (4,-3) {$8$};
\node[rlab] at (4,-3.5) {R3};
\node[restlab] at (4,-3.85) {56 71};
\node[ltleaf] at (6,-3) {$19$};
\node[rlab] at (6,-3.5) {R4};
\node[restlab] at (6,-3.85) {22 95};
\end{tikzpicture}
&
\begin{tikzpicture}[x=0.70cm,y=1.00cm]
\draw[lte] (3,-1) -- (1,-2);
\draw[lte] (3,-1) -- (5,-2);
\draw[lte] (1,-2) -- (0,-3);
\draw[lte] (1,-2) -- (2,-3);
\draw[lte] (5,-2) -- (4,-3);
\draw[lte] (5,-2) -- (6,-3);
\node[ltn, ltchg] at (3,-1) {$15$};
\node[rlab, anchor=west] at (3.51429,-0.95) {R1};
\node[ltn] at (1,-2) {$15$};
\node[rlab, anchor=west] at (1.51429,-1.95) {R1};
\node[ltn, ltchg] at (5,-2) {$19$};
\node[rlab, anchor=west] at (5.51429,-1.95) {R4};
\node[ltleaf] at (0,-3) {$15$};
\node[rlab] at (0,-3.5) {R1};
\node[restlab] at (0,-3.85) {42 63};
\node[ltleaf] at (2,-3) {$27$};
\node[rlab] at (2,-3.5) {R2};
\node[restlab] at (2,-3.85) {31 80};
\node[ltleaf, ltnew] at (4,-3) {$56$};
\node[rlab] at (4,-3.5) {R3};
\node[restlab] at (4,-3.85) {71};
\node[ltleaf] at (6,-3) {$19$};
\node[rlab] at (6,-3.5) {R4};
\node[restlab] at (6,-3.85) {22 95};
\end{tikzpicture}
\\[2pt]
{\small (a) 建好之後，冠軍是 $8$（$R_3$）} & {\small (b) 輸出 $8$，$R_3$ 補上 $56$ 後重賽}
\end{tabular}
\diagcaption{勝者樹。 (b) 中橘色的節點是重賽改過的：節點 3 拿 $56$ 和兄弟 $R_4$ 的
$19$ 比，$19$ 贏；節點 1 拿 $19$ 和兄弟節點 2 的贏家 $15$ 比，$15$ 贏。}
\label{fig:lin-winner}
\end{center}

\subsection{敗者樹}
\label{sec:lin-loser}

勝者樹重賽時，每一層都要去讀兄弟節點，才知道要和誰比。 \textbf{敗者樹}（loser
tree，也叫 tree of losers）換個方式記錄：每個內部節點記下那場比賽的\textbf{輸家}，
另外用一個節點 $0$ 記總冠軍。

為什麼記輸家就夠了？ 冠軍從它的葉子一路贏到根。 在路上的節點 $t$，它打敗的是另一
邊子樹的贏家，這位選手就是節點 $t$ 記的輸家。 冠軍的串補上新的鍵後，新鍵沿同一條路
往上重賽，在節點 $t$ 要面對的對手，正是另一邊子樹的贏家，也就是節點 $t$ 記的那位輸
家。 所以每一層只要和\textbf{自己這個節點}記的選手比：贏家繼續往上，輸家留在這個節
點。 走到根之後，留下來的贏家就是新的冠軍，放進節點 $0$。

下面的虛擬碼用 $\mathit{key}[i]$ 表示串 $R_i$ 目前排在最前面的鍵，串讀完了就是
$\infty$。 \proc{Beats}$(a, b)$ 表示串 $a$ 贏串 $b$：$\mathit{key}[a] <
\mathit{key}[b]$，或者兩者相等而 $a < b$（鍵相同時編號小的串贏，這讓合併是穩定
的）。 $L[0..k-1]$ 是敗者樹，$W[1..k-1]$ 只在建樹時暫存每場比賽的贏家。

\begin{pseudo}[label=alg:lin-loser]{敗者樹：由下往上建樹、沿路徑調整}
\Procedure{Build-Loser-Tree}{$k$}
    \For{$t \gets k - 1$ \DownTo $1$}
        \State $a \gets$ 節點 $2t$ 的贏家 \Comment{葉子 $c$ 的贏家是串 $c - k + 1$；內部節點 $c$ 是 $W[c]$}
        \State $b \gets$ 節點 $2t + 1$ 的贏家
        \If{\Call{Beats}{$a, b$}}
            \State $W[t] \gets a$；$L[t] \gets b$
        \Else
            \State $W[t] \gets b$；$L[t] \gets a$
        \EndIf
    \EndFor
    \State $L[0] \gets W[1]$
\EndProcedure
\Procedure{Adjust}{$s$} \Comment{串 $s$ 剛換了新的鍵}
    \State $t \gets \lfloor (s + k - 1) / 2 \rfloor$ \Comment{串 $s$ 的葉子的父節點}
    \While{$t > 0$}
        \If{\Call{Beats}{$L[t], s$}}
            \State 交換 $s$ 和 $L[t]$ \Comment{原本的輸家贏了，往上走；$s$ 留下}
        \EndIf
        \State $t \gets \lfloor t / 2 \rfloor$
    \EndWhile
    \State $L[0] \gets s$
\EndProcedure
\end{pseudo}

\begin{pseudo}[label=alg:lin-kmerge]{用敗者樹做 $k$ 路合併}
\Procedure{K-Way-Merge}{$R_1, \ldots, R_k$}
    \For{$i \gets 1$ \To $k$}
        \State $\mathit{key}[i] \gets R_i$ 的第一個鍵（空的串是 $\infty$）
    \EndFor
    \State \Call{Build-Loser-Tree}{$k$}
    \While{$\mathit{key}[L[0]] \ne \infty$}
        \State $s \gets L[0]$
        \State 輸出 $\mathit{key}[s]$
        \State $\mathit{key}[s] \gets R_s$ 的下一個鍵（讀完了是 $\infty$）
        \State \Call{Adjust}{$s$}
    \EndWhile
\EndProcedure
\end{pseudo}

\textbf{手算追蹤。} 用勝者樹那四個串。 建樹時由下往上：節點 2 是 $R_1$ 的 $15$ 對
$R_2$ 的 $27$，$15$ 贏，記下輸家 $R_2$；節點 3 是 $8$ 對 $19$，記下輸家 $R_4$；節點 1
是兩邊的贏家 $15$ 對 $8$，記下輸家 $R_1$，冠軍 $R_3$ 放進節點 $0$。 建樹比了 $k - 1 =
3$ 次。 之後每輸出一筆，就沿著那個串的葉子往上調整一次。 下圖畫出前三次輸出之後的
樹：方框是節點 $0$ 的冠軍，圓圈是內部節點記的輸家，橘色是這次調整中換了選手的節點，
綠色是剛補上新鍵的葉子。

\begin{center}
\begin{tabular}{c@{\hspace{1.2cm}}c}
\begin{tikzpicture}[x=0.70cm,y=1.00cm]
\draw[lte] (3,-1) -- (1,-2);
\draw[lte] (3,-1) -- (5,-2);
\draw[lte] (1,-2) -- (0,-3);
\draw[lte] (1,-2) -- (2,-3);
\draw[lte] (5,-2) -- (4,-3);
\draw[lte] (5,-2) -- (6,-3);
\draw[lte] (3,-1) -- (3,0);
\node[ltwin] at (3,0) {$8$};
\node[rlab, anchor=west] at (3.51429,0) {R3};
\node[ltn] at (3,-1) {$15$};
\node[rlab, anchor=west] at (3.51429,-0.95) {R1};
\node[ltn] at (1,-2) {$27$};
\node[rlab, anchor=west] at (1.51429,-1.95) {R2};
\node[ltn] at (5,-2) {$19$};
\node[rlab, anchor=west] at (5.51429,-1.95) {R4};
\node[ltleaf] at (0,-3) {$15$};
\node[rlab] at (0,-3.5) {R1};
\node[restlab] at (0,-3.85) {42 63};
\node[ltleaf] at (2,-3) {$27$};
\node[rlab] at (2,-3.5) {R2};
\node[restlab] at (2,-3.85) {31 80};
\node[ltleaf] at (4,-3) {$8$};
\node[rlab] at (4,-3.5) {R3};
\node[restlab] at (4,-3.85) {56 71};
\node[ltleaf] at (6,-3) {$19$};
\node[rlab] at (6,-3.5) {R4};
\node[restlab] at (6,-3.85) {22 95};
\end{tikzpicture}
&
\begin{tikzpicture}[x=0.70cm,y=1.00cm]
\draw[lte] (3,-1) -- (1,-2);
\draw[lte] (3,-1) -- (5,-2);
\draw[lte] (1,-2) -- (0,-3);
\draw[lte] (1,-2) -- (2,-3);
\draw[lte] (5,-2) -- (4,-3);
\draw[lte] (5,-2) -- (6,-3);
\draw[lte] (3,-1) -- (3,0);
\node[ltwin] at (3,0) {$15$};
\node[rlab, anchor=west] at (3.51429,0) {R1};
\node[ltn, ltchg] at (3,-1) {$19$};
\node[rlab, anchor=west] at (3.51429,-0.95) {R4};
\node[ltn] at (1,-2) {$27$};
\node[rlab, anchor=west] at (1.51429,-1.95) {R2};
\node[ltn, ltchg] at (5,-2) {$56$};
\node[rlab, anchor=west] at (5.51429,-1.95) {R3};
\node[ltleaf] at (0,-3) {$15$};
\node[rlab] at (0,-3.5) {R1};
\node[restlab] at (0,-3.85) {42 63};
\node[ltleaf] at (2,-3) {$27$};
\node[rlab] at (2,-3.5) {R2};
\node[restlab] at (2,-3.85) {31 80};
\node[ltleaf, ltnew] at (4,-3) {$56$};
\node[rlab] at (4,-3.5) {R3};
\node[restlab] at (4,-3.85) {71};
\node[ltleaf] at (6,-3) {$19$};
\node[rlab] at (6,-3.5) {R4};
\node[restlab] at (6,-3.85) {22 95};
\end{tikzpicture}
\\[2pt]
{\small (a) 建好之後} & {\small (b) 輸出 $8$，$R_3$ 補上 $56$} \\[10pt]
\begin{tikzpicture}[x=0.70cm,y=1.00cm]
\draw[lte] (3,-1) -- (1,-2);
\draw[lte] (3,-1) -- (5,-2);
\draw[lte] (1,-2) -- (0,-3);
\draw[lte] (1,-2) -- (2,-3);
\draw[lte] (5,-2) -- (4,-3);
\draw[lte] (5,-2) -- (6,-3);
\draw[lte] (3,-1) -- (3,0);
\node[ltwin] at (3,0) {$19$};
\node[rlab, anchor=west] at (3.51429,0) {R4};
\node[ltn, ltchg] at (3,-1) {$27$};
\node[rlab, anchor=west] at (3.51429,-0.95) {R2};
\node[ltn, ltchg] at (1,-2) {$42$};
\node[rlab, anchor=west] at (1.51429,-1.95) {R1};
\node[ltn] at (5,-2) {$56$};
\node[rlab, anchor=west] at (5.51429,-1.95) {R3};
\node[ltleaf, ltnew] at (0,-3) {$42$};
\node[rlab] at (0,-3.5) {R1};
\node[restlab] at (0,-3.85) {63};
\node[ltleaf] at (2,-3) {$27$};
\node[rlab] at (2,-3.5) {R2};
\node[restlab] at (2,-3.85) {31 80};
\node[ltleaf] at (4,-3) {$56$};
\node[rlab] at (4,-3.5) {R3};
\node[restlab] at (4,-3.85) {71};
\node[ltleaf] at (6,-3) {$19$};
\node[rlab] at (6,-3.5) {R4};
\node[restlab] at (6,-3.85) {22 95};
\end{tikzpicture}
&
\begin{tikzpicture}[x=0.70cm,y=1.00cm]
\draw[lte] (3,-1) -- (1,-2);
\draw[lte] (3,-1) -- (5,-2);
\draw[lte] (1,-2) -- (0,-3);
\draw[lte] (1,-2) -- (2,-3);
\draw[lte] (5,-2) -- (4,-3);
\draw[lte] (5,-2) -- (6,-3);
\draw[lte] (3,-1) -- (3,0);
\node[ltwin] at (3,0) {$22$};
\node[rlab, anchor=west] at (3.51429,0) {R4};
\node[ltn] at (3,-1) {$27$};
\node[rlab, anchor=west] at (3.51429,-0.95) {R2};
\node[ltn] at (1,-2) {$42$};
\node[rlab, anchor=west] at (1.51429,-1.95) {R1};
\node[ltn] at (5,-2) {$56$};
\node[rlab, anchor=west] at (5.51429,-1.95) {R3};
\node[ltleaf] at (0,-3) {$42$};
\node[rlab] at (0,-3.5) {R1};
\node[restlab] at (0,-3.85) {63};
\node[ltleaf] at (2,-3) {$27$};
\node[rlab] at (2,-3.5) {R2};
\node[restlab] at (2,-3.85) {31 80};
\node[ltleaf] at (4,-3) {$56$};
\node[rlab] at (4,-3.5) {R3};
\node[restlab] at (4,-3.85) {71};
\node[ltleaf, ltnew] at (6,-3) {$22$};
\node[rlab] at (6,-3.5) {R4};
\node[restlab] at (6,-3.85) {95};
\end{tikzpicture}
\\[2pt]
{\small (c) 輸出 $15$，$R_1$ 補上 $42$} & {\small (d) 輸出 $19$，$R_4$ 補上 $22$}
\end{tabular}
\diagcaption{敗者樹的前三次輸出。 (b)：$56$ 在節點 3 輸給記在那裡的 $19$，$56$ 留下、
$19$ 往上；在節點 1 輸給 $15$，$19$ 留下、$15$ 成為冠軍。 (d)：$22$ 在節點 3 贏
$56$、在節點 1 贏 $27$，一路上都沒有換人，直接成為冠軍。}
\label{fig:lin-loser}
\end{center}

全部十二次輸出列成下表。 「路徑上的比賽」寫出每個節點那場比賽的結果，例如「3：19
勝 56」是在節點 3，$19$ 贏了 $56$；最後一欄是調整完的 $L[0], L[1], L[2], L[3]$。

\begin{center}
\small
\begin{tabular}{ccclc}
\toprule
輸出 & 來自 & 補上的鍵 & 路徑上的比賽（節點：結果） & $L[0..3]$ \\
\midrule
\multicolumn{4}{l}{建樹} & R3 R1 R2 R4 \\
8 & R3 & 56 & 3：19 勝 56；1：15 勝 19 & R1 R4 R2 R3 \\
15 & R1 & 42 & 2：27 勝 42；1：19 勝 27 & R4 R2 R1 R3 \\
19 & R4 & 22 & 3：22 勝 56；1：22 勝 27 & R4 R2 R1 R3 \\
22 & R4 & 95 & 3：56 勝 95；1：27 勝 56 & R2 R3 R1 R4 \\
27 & R2 & 31 & 2：31 勝 42；1：31 勝 56 & R2 R3 R1 R4 \\
31 & R2 & 80 & 2：42 勝 80；1：42 勝 56 & R1 R3 R2 R4 \\
42 & R1 & 63 & 2：63 勝 80；1：56 勝 63 & R3 R1 R2 R4 \\
56 & R3 & 71 & 3：71 勝 95；1：63 勝 71 & R1 R3 R2 R4 \\
63 & R1 & $\infty$ & 2：80 勝 $\infty$；1：71 勝 80 & R3 R2 R1 R4 \\
71 & R3 & $\infty$ & 3：95 勝 $\infty$；1：80 勝 95 & R2 R4 R1 R3 \\
80 & R2 & $\infty$ & 2：$\infty$ 勝 $\infty$；1：95 勝 $\infty$ & R4 R1 R2 R3 \\
95 & R4 & $\infty$ & 3：$\infty$ 勝 $\infty$；1：$\infty$ 勝 $\infty$ & R1 R3 R2 R4 \\
\bottomrule
\end{tabular}
\end{center}

輸出依序是 $8, 15, 19, 22, 27, 31, 42, 56, 63, 71, 80, 95$。 每次調整比 $2 = \lg 4$
次，建樹 $3$ 次，共 $3 + 12 \times 2 = 27$ 次比較。 串讀完時補上 $\infty$，它比任何鍵
都大，自然一路輸到底；兩邊都是 $\infty$ 時照平手的約定，編號小的贏（表中「$\infty$ 勝
$\infty$」）。 所有串都讀完時冠軍的鍵是 $\infty$，合併結束。

\begin{examnote}
敗者樹的追蹤題，每輸出一筆就畫一次整棵樹，內部節點寫輸家（寫鍵，必要時加上串的編
號），另外畫出冠軍（節點 $0$）；只畫冠軍或只列輸出序列都不完整。 鍵相同時誰贏、串讀
完後用 $\infty$ 補，這兩個約定不同課本不一樣，照題目的規定；題目沒說就自己寫明。
\end{examnote}

\textbf{三種做法的比較。} 從 $k$ 個串挑最小，也可以用第~8~章的最小堆積：堆積裡放每
個串目前的鍵，輸出根以後把根換成同一個串的下一個鍵，再做一次下沉。

\begin{center}
\begin{tabular}{p{0.2\linewidth}p{0.22\linewidth}p{0.22\linewidth}p{0.22\linewidth}}
\toprule
 & 勝者樹 & 敗者樹 & 最小堆積 \\
\midrule
節點記什麼 & 每場比賽的贏家 & 每場比賽的輸家，另記冠軍 & 堆積的元素 \\
每輸出一筆的比較次數 & $\lceil \lg k \rceil$ 以內 & $\lceil \lg k \rceil$ 以內 & 最多約 $2\lfloor \lg k \rfloor$ \\
重賽要讀的節點 & 路徑上的節點和它們的兄弟 & 只有路徑上的節點 & 路徑上的節點和兄弟 \\
建立 & $k - 1$ 次比較 & $k - 1$ 次比較 & $O(k)$（第~8~章） \\
\bottomrule
\end{tabular}
\end{center}

堆積下沉時每一層要比兩次（先比兩個子節點，再和較小的比），所以比較次數大約是兩棵
樹的兩倍。 勝者樹和敗者樹比較次數一樣，差在敗者樹重賽時不用讀兄弟節點，每一層只碰
一個節點，程式也比較短，所以外部排序的 $k$ 路合併多半用敗者樹。

\subsection{C 程式}

下面的程式用陣列模擬磁碟上的四個串，\texttt{struct Run} 記錄陣列、長度和下一個要讀
的位置。 為了讓程式短一點，$k$、\texttt{key}、\texttt{L} 和比較次數寫成全域變數。 C 的串編號從 $0$ 開始：串 $i$ 的葉子是節點 $k + i$，父節點是
$(k + i)/2$，正好是虛擬碼的 $\lfloor (s + k - 1)/2 \rfloor$ 把 $s$ 換成 $i + 1$。 印
的時候把編號加 $1$，和上面的圖、表一致。 讀完的串用 \texttt{INT\_MAX} 當 $\infty$，
所以真正的鍵不能是 \texttt{INT\_MAX}。

\begin{example}[label=ex:lin-loser]{c}{(loser\_merge.c) 用敗者樹合併四個串}
#include <stdio.h>
#include <stdlib.h>
#include <limits.h>

// 用陣列模擬磁碟上的一個串，pos 是下一個要讀的位置。
struct Run {
    const int *a;
    int len, pos;
};

int k;             // 串數，串的編號是 0..k-1
int *key;          // key[i]：串 i 目前最前面的鍵，讀完是 INT_MAX
int *L;            // L[1..k-1]：每場比賽的輸家；L[0]：冠軍
long long cmp;     // 比較次數

// 串 a 贏串 b：鍵小的贏，鍵相同時編號小的贏。
int beats(int a, int b)
{
    cmp++;
    return key[a] < key[b] || (key[a] == key[b] && a < b);
}

int next_key(struct Run *r)
{
    return r->pos < r->len ? r->a[r->pos++] : INT_MAX;
}

// 由下往上建樹；w[x] 暫存節點 x 的贏家。葉子 k + i 代表串 i。
void build(void)
{
    int *w = malloc(k * sizeof *w);
    if (w == NULL)
        exit(1);
    for (int x = k - 1; x >= 1; x--) {
        int a = 2 * x >= k ? 2 * x - k : w[2 * x];
        int b = 2 * x + 1 >= k ? 2 * x + 1 - k : w[2 * x + 1];
        w[x] = beats(a, b) ? a : b;
        L[x] = w[x] == a ? b : a;
    }
    L[0] = k > 1 ? w[1] : 0;
    free(w);
}

// 串 s 換了新的鍵：從它的葉子往上，只和每個節點記的輸家比。
void adjust(int s)
{
    for (int x = (s + k) / 2; x > 0; x /= 2) {
        if (beats(L[x], s)) {          // 記在這裡的輸家贏了：交換
            int t = L[x];
            L[x] = s;
            s = t;
        }
    }
    L[0] = s;
}

int main(void)
{
    static const int r1[] = {15, 42, 63}, r2[] = {27, 31, 80},
                     r3[] = {8, 56, 71}, r4[] = {19, 22, 95};
    struct Run runs[] = {{r1, 3, 0}, {r2, 3, 0},
                         {r3, 3, 0}, {r4, 3, 0}};
    k = 4;
    key = malloc(k * sizeof *key);
    L = malloc(k * sizeof *L);
    if (key == NULL || L == NULL)
        return 1;
    for (int i = 0; i < k; i++)
        key[i] = next_key(&runs[i]);
    build();
    printf("build");
    for (;;) {
        printf(" L =");
        for (int x = 0; x < k; x++)
            printf(" R%d", L[x] + 1);  // 印的時候編號從 1 開始
        printf("\n");
        int s = L[0];
        if (key[s] == INT_MAX)         // 冠軍是無限大：全部讀完了
            break;
        printf("out %d R%d", key[s], s + 1);
        key[s] = next_key(&runs[s]);
        adjust(s);
    }
    printf("comparisons = %lld\n", cmp);
    free(key);
    free(L);
    return 0;
}
\end{example}
\iosample{}{build L = R3 R1 R2 R4 \\ out 8 R3 L = R1 R4 R2 R3 \\ out 15 R1 L = R4 R2 R1 R3 \\ out 19 R4 L = R4 R2 R1 R3 \\ out 22 R4 L = R2 R3 R1 R4 \\ out 27 R2 L = R2 R3 R1 R4 \\ out 31 R2 L = R1 R3 R2 R4 \\ out 42 R1 L = R3 R1 R2 R4 \\ out 56 R3 L = R1 R3 R2 R4 \\ out 63 R1 L = R3 R2 R1 R4 \\ out 71 R3 L = R2 R4 R1 R3 \\ out 80 R2 L = R4 R1 R2 R3 \\ out 95 R4 L = R1 R3 R2 R4 \\ comparisons = 27}

每一列的 $L$ 和上面的表一致，總共 $27$ 次比較。 要合併真正的檔案，只要把
\texttt{next\_key} 改成從檔案（經過緩衝區）讀下一個整數，敗者樹的部分不用動。

\subsection{取代選擇：產生更長的初始串}
\label{sec:lin-replace}

合併趟數是 $\lceil \log_k r \rceil$，初始串越少越好。 用內部排序產生初始串，每個串
剛好長 $M$。 \textbf{取代選擇}（replacement selection）在輸出的同時讀入新的紀錄，
平均可以產生長 $2M$ 的串。

想法是這樣：記憶體裡隨時放 $M$ 筆紀錄。 每次輸出記憶體裡最小的那筆到目前的串，空出
來的位置馬上讀入下一筆。 新讀入的鍵如果不比剛剛輸出的小，它還能接在這個串後面；如
果比較小，這個串已經寫過頭了，它只能留給下一個串，稱為\textbf{凍結}。 記憶體裡全部
都凍結時，這個串結束，解除凍結，開始下一個串。

實作上，把每筆紀錄標上它屬於第幾個串，用（串號, 鍵）當堆積的比較依據：串號小的優
先，串號相同時鍵小的優先。 凍結的紀錄串號多 $1$，自然排在目前的串的所有紀錄後面。

\begin{pseudo}[label=alg:lin-replace]{取代選擇}
\Procedure{Replacement-Selection}{$M$}
    \State 讀入前 $M$ 筆紀錄，每筆的串號都設成 $1$
    \State 依（串號, 鍵）把它們建成最小堆積 $H$
    \While{$H$ 不是空的}
        \State $(r, x) \gets H$ 的根
        \State 把 $x$ 寫到第 $r$ 個串的尾端
        \If{輸入還有紀錄}
            \State $y \gets$ 讀入下一筆
            \If{$y \ge x$}
                \State 把根換成 $(r, y)$
            \Else
                \State 把根換成 $(r + 1, y)$ \Comment{凍結，留給下一個串}
            \EndIf
        \Else
            \State 刪掉根，用堆積的最後一筆補上
        \EndIf
        \State 對根做下沉
    \EndWhile
\EndProcedure
\end{pseudo}

\textbf{手算追蹤。} $M = 3$，輸入依序是
\[
41, 17, 63, 29, 8, 52, 35, 12, 70, 24, 58, 6, 47, 33, 90, 15.
\]
每一列是輸出前記憶體裡的三筆（加括號的是凍結的，屬於下一個串）、這一步輸出什麼、補
進來什麼。

\begin{center}
\small
\begin{tabular}{rlccl}
\toprule
步 & 記憶體（輸出前） & 輸出 & 串 & 補進來 \\
\midrule
1 & 41\ \ 17\ \ 63 & 17 & 1 & 29 \\
2 & 41\ \ 29\ \ 63 & 29 & 1 & 8，比 29 小，凍結 \\
3 & 41\ \ (8)\ \ 63 & 41 & 1 & 52 \\
4 & 52\ \ (8)\ \ 63 & 52 & 1 & 35，凍結 \\
5 & (35)\ \ (8)\ \ 63 & 63 & 1 & 12，凍結；全部凍結，串 1 結束 \\
\midrule
6 & 35\ \ 8\ \ 12 & 8 & 2 & 70 \\
7 & 35\ \ 70\ \ 12 & 12 & 2 & 24 \\
8 & 35\ \ 70\ \ 24 & 24 & 2 & 58 \\
9 & 35\ \ 70\ \ 58 & 35 & 2 & 6，凍結 \\
10 & (6)\ \ 70\ \ 58 & 58 & 2 & 47，凍結 \\
11 & (6)\ \ 70\ \ (47) & 70 & 2 & 33，凍結；全部凍結，串 2 結束 \\
\midrule
12 & 6\ \ 33\ \ 47 & 6 & 3 & 90 \\
13 & 90\ \ 33\ \ 47 & 33 & 3 & 15，凍結 \\
14 & 90\ \ (15)\ \ 47 & 47 & 3 & 輸入讀完 \\
15 & 90\ \ (15) & 90 & 3 & --；串 3 結束 \\
\midrule
16 & 15 & 15 & 4 & -- \\
\bottomrule
\end{tabular}
\end{center}

得到四個串：$\langle 17, 29, 41, 52, 63 \rangle$、$\langle 8, 12, 24, 35, 58, 70
\rangle$、$\langle 6, 33, 47, 90 \rangle$、$\langle 15 \rangle$。 用內部排序每次排
$3$ 筆，$16$ 筆會得到 $\lceil 16/3 \rceil = 6$ 個串；取代選擇得到 $4$ 個，前兩個都
比 $M$ 長。 第 2 步的 $8$ 比剛輸出的 $29$ 小，如果讓它進串 1，串 1 就不是排好的了，
所以凍結。

寫成 C 程式時，把虛擬碼~\ref{alg:lin-replace} 的堆積換成第~8~章的陣列堆積，堆積的元
素是（串號, 鍵）的結構，比較時先比串號、再比鍵，其他照抄第~8~章的下沉就可以了。

\textbf{串有多長。} 幾個容易確認的性質：
\begin{itemize}
\item 輸入已經由小到大排好：讀進來的鍵都不比剛輸出的小，永遠不凍結，整個檔案只有一
個串。
\item 輸入由大到小排好：每個新讀入的鍵都比剛輸出的小，全部凍結，每個串剛好 $M$ 筆
（最後一個可能比較短），和內部排序一樣。
\item 一般情況，除了最後一個串，每個串至少 $M$ 筆（觀念題第 10 題）。
\end{itemize}
輸入是隨機的時候，Knuth 證明串的期望長度趨近 $2M$（TAOCP 第 3 卷 5.4.1 節），所以初
始串大約只有 $N/(2M)$ 個，是內部排序的一半。 這裡不證明。

代價是什麼？ 取代選擇每輸出一筆都要做一次堆積調整，$O(\lg M)$，總共 $O(N \lg M)$，
和用內部排序產生初始串同一級；它的串長度不固定，後面的合併程式要能處理長短不一的
串。 實際上也可以用敗者樹代替堆積，比較次數更少。 串長短不一時，先合併哪幾個串會影響總成本；
這和第~19~章的霍夫曼編碼是同一個問題，每次挑最短的 $k$ 個串合併最省。

\textbf{例 2.} 接例 1：$N = 1{,}200{,}000$，$M = 120{,}000$。 用取代選擇，初始串大
約 $N/(2M) = 5$ 個。 $2$ 路合併要 $\lceil \lg 5 \rceil = 3$ 趟（原本 $4$ 趟），$5$ 路
只要 $1$ 趟（原本 $\lceil \log_5 10 \rceil = 2$ 趟）。 總 I/O 次數：$5$ 路合併時是
$150{,}000 \times (1 + 1) = 300{,}000$，比原本的 $450{,}000$ 少三分之一。

\section{本章常見錯誤}
\label{sec:lin-pitfalls}

\begin{center}
\begin{tabular}{p{0.26\linewidth}p{0.48\linewidth}p{0.16\linewidth}}
\toprule
錯誤 & 例子 & 見 \\
\midrule
下界的證明少了一步 & 只寫「葉子數 $\ge n!$」，沒說明不同的順序不能走到同一個葉子 &
第~\ref{sec:lin-lower}~節 \\
以為下界管所有排序 & 「計數排序是 $\Theta(n + k)$，違反了 $\Omega(n \lg n)$ 下界」 &
第~\ref{sec:lin-lower}~節 \\
計數排序從左往右放 & 前綴和之後從 $j = 1$ 開始放，鍵相同的紀錄整個顛倒，不穩定 &
第~\ref{sec:lin-counting}~節 \\
LSD 的方向或穩定性弄錯 & 先排最高位；每一趟用快速排序 & 第~\ref{sec:lin-lsd}~節 \\
趟數忘了取天花板 & $\log_4 10 \approx 1.66$ 寫成 $1.66$ 趟或 $1$ 趟 &
第~\ref{sec:lin-passes}~節 \\
I/O 次數少算一半 & 一趟只算讀，沒算寫；或漏掉產生初始串那一趟 &
第~\ref{sec:lin-passes}~節 \\
敗者樹記成贏家 & 內部節點寫成每場的贏家，那是勝者樹 & 第~\ref{sec:lin-loser}~節 \\
取代選擇沒有凍結 & 新讀入的鍵比剛輸出的小，仍然放進目前的串，串就沒排好 &
第~\ref{sec:lin-replace}~節 \\
\bottomrule
\end{tabular}
\end{center}

\section{觀念題}
\label{sec:lin-exercises}

題型和研究所考試、期中期末考的筆試相近，有選擇、計算、追蹤和證明。 解答在本章最後。

\begin{enumerate}
\item （選擇，可能不只一個答案）下列哪些敘述正確？
\begin{enumerate}
\item 一個正確的比較排序，排 $n$ 個數的決策樹恰好有 $n!$ 個葉子。
\item 合併排序和堆積排序在比較排序裡是漸近最佳的。
\item $n$ 個鍵都是 $0$ 到 $n^2 - 1$ 的整數時，直接用計數排序（$k = n^2 - 1$）就能在
$O(n)$ 時間排好。
\item LSD 基數排序每一趟改用快速排序，結果仍然一定正確。
\item 計數排序的時間可以低於 $n \lg n$，因為它沒有比較鍵的大小，不受比較排序下界的
限制。
\end{enumerate}

\item （追蹤）用選擇排序的想法排三個數 $A[1..3]$：$m \gets 1$；若 $A[2] < A[m]$ 則
$m \gets 2$；若 $A[3] < A[m]$ 則 $m \gets 3$；交換 $A[1]$ 和 $A[m]$；最後若
$A[3] < A[2]$ 就交換 $A[2]$ 和 $A[3]$。 用原本的編號 $a_1, a_2, a_3$ 畫出它的決策樹（鍵兩兩不同）。 樹有幾個葉子？ 哪些葉子
是任何輸入都走不到的？ 最壞和平均各比幾次？

\item （證明）在 $n$ 個兩兩不同的數裡找最大值。
\begin{enumerate}
\item 用決策樹的論證，可以得到什麼下界？
\item 證明任何只用比較的演算法，最壞至少要 $n - 1$ 次比較。
\item 說明為什麼 (a) 的下界和實際差很多。
\end{enumerate}

\item （追蹤）用虛擬碼~\ref{alg:lin-counting} 排序
$A = \langle 2_p, 0_q, 3_r, 2_s, 1_t, 0_u, 2_v \rangle$（$k = 3$，下標的字母是附帶的
資料）。 寫出第 5--6 行做完之後的 $C$，以及第 7--10 行每一輪放完之後的 $B$。

\item （追蹤）六個四位數的編號 $3172, 5408, 1772, 3105, 5472, 1708$。 用 LSD 基數排
序，但以 $100$ 為基底（每「一位」是兩個十進位數字）。 要幾趟？ 寫出每一趟之後的
順序，並指出哪一對編號的最後順序依賴第二趟的穩定性。

\item （計算）桶子排序改用 $m$ 個桶子（$[0, 1)$ 切成 $m$ 等分），輸入仍是 $n$ 個均勻
分布的數。
\begin{enumerate}
\item 求 $E[n_j^2]$，以及第 6 行（插入排序）的總期望時間，用 $n$、$m$ 表示。
\item $m = n/2$ 和 $m = \sqrt{n}$ 時，整個桶子排序的期望時間各是多少？
\end{enumerate}

\item （計算）一個檔案有 $N = 2{,}000{,}000$ 筆紀錄，記憶體放得下 $M = 50{,}000$ 筆，
一個區塊放 $B = 100$ 筆。
\begin{enumerate}
\item 用內部排序產生初始串，有幾個串？
\item 用 $2$ 路、$5$ 路、$8$ 路合併，各要幾趟合併？
\item $8$ 路合併時，包含產生初始串，總共讀寫幾個區塊？
\item 改用取代選擇，以串的期望長度估計，$5$ 路合併要幾趟？
\end{enumerate}

\item （追蹤）用敗者樹合併五個串：
$R_1 = \langle 23, 48 \rangle$、$R_2 = \langle 11, 60, 74 \rangle$、
$R_3 = \langle 35, 39 \rangle$、$R_4 = \langle 17, 52 \rangle$、
$R_5 = \langle 29, 44, 81 \rangle$。 用本章的編號（葉子 $k + i - 1$ 代表 $R_i$）。
\begin{enumerate}
\item 畫出建好的敗者樹，寫出 $L[0..4]$。
\item 畫出輸出兩筆之後的樹。
\item 合併完總共比較幾次？
\end{enumerate}

\item （追蹤）記憶體放得下 $M = 4$ 筆，用取代選擇處理
\[
55, 31, 72, 18, 64, 9, 40, 83, 27, 50, 12, 66, 35, 71, 5, 44,
\]
寫出產生的每個串。 和用內部排序產生初始串相比，串的個數差多少？

\item （證明）證明取代選擇產生的串，除了最後一個，每個都至少有 $M$ 筆。

\item （證明）第~1~章 1.1 節的「找重複學號」。 假設演算法只能問「$a_i = a_j$ 嗎？」（只
能測相等，不能比大小）。 證明在最壞情況下，判斷 $n$ 個數裡有沒有兩個相同，至少要問
$\binom{n}{2}$ 次。（提示：讓對手每次都回答「不相等」。）
\end{enumerate}

\section{實作練習}
\label{sec:lin-practice}

下面是 Aki-Judge 上和本章相關的題目。 這一章的基數排序、敗者樹、取代選擇目前沒有對
應的題目，請以觀念題和本章的範例程式為主。

\begin{center}
\begin{tabular}{llc}
\toprule
小節 & 題目 & 難度 \\
\midrule
計數排序 & \akiref{counting-sort-function}{計數排序函式} & 基本 \\
MSD 基數排序 & \akiref{sort-words-in-dictionary-order}{將單字依字典序排序} & 基本 \\
$k$ 路合併 & \akiref{merge-sorted-number-files-with-counts}{合併排序數列檔並統計次數} & 進階 \\
\bottomrule
\end{tabular}
\end{center}

「計數排序函式」只要交排序函式；鍵的範圍到 $10^6$，計數陣列要用 \texttt{malloc}，
而且這題只有鍵、沒有附帶資料，可以不用輸出陣列 $B$，數完之後直接依序寫回。 「將單字
依字典序排序」可以用 \texttt{qsort} 加 \texttt{strcmp} 通過，練習時試著用
第~\ref{sec:lin-msd}~節的 MSD 基數排序自己寫：依第一個字母分成 $26$ 個桶子，字串
結束的放在最前面的桶子。 「合併排序數列檔並統計次數」是真正讀檔的 $k$ 路合併，每個
檔案同時只保留一個還沒處理的數；$k \le 20$，每次掃過 $k$ 個候選就夠快，想練習的話可
以把範例~\ref{ex:lin-loser} 的敗者樹搬過去。

\section{讀完你應該會}
\label{sec:lin-checklist}

\begin{itemize}
\item 能畫出一個小比較排序的決策樹，用它算最壞和平均的比較次數；能完整寫出比較排序
最壞 $\Omega(n \lg n)$ 的證明，以及平均情況的下界。
\item 能用資訊理論下界和對手論證估計其他問題（插入、合併、找最大值）的比較次數，並
判斷下界是否達得到。
\item 能逐步追蹤穩定的計數排序（前綴和陣列、從右往左放），說明為什麼從右往左才穩定。
\item 能寫出 LSD 基數排序每一趟的結果，說明為什麼內層要穩定，並依位數切法估計成本；能
說明 MSD 和 LSD 的差別、桶子排序期望線性的理由。
\item 給定 $N$、$M$、$B$、$k$，能算出初始串數、合併趟數和 I/O 次數；能比較用內部排序
和取代選擇產生初始串的差別。
\item 能逐步追蹤敗者樹的建立和每次輸出後的調整，並追蹤取代選擇產生的
串。
\end{itemize}

\textbf{延伸閱讀}（書末會整理完整的參考資料）

\begin{itemize}
\item Thomas H. Cormen 等，\textit{Introduction to Algorithms}，第 4 版，第 8 章：決策樹
下界、計數排序、基數排序、桶子排序；章末問題有合併下界、平均下界等延伸。
\item Donald E. Knuth，\textit{The Art of Computer Programming}，第 3 卷第 2 版：5.3.1 節
最少比較次數的排序，5.4.1 節取代選擇與串的期望長度，5.4 節其他小節有多階段合併等外部
排序的方法。
\item Steven S. Skiena，\textit{The Algorithm Design Manual}，第 3 版，第 4 章：排序的
應用和實務上怎麼選排序方法。
\end{itemize}

\addcontentsline{toc}{section}{觀念題解答}
\section*{觀念題解答}

\begin{enumerate}
\item (b)、(e) 正確。 (a)：葉子至少 $n!$ 個，可以更多，例如第 2 題的決策樹有 $8$ 個葉
子，其中兩個走不到。 (c)：計數排序是 $\Theta(n + k) = \Theta(n^2)$；要 $O(n)$ 得用以
$n$ 為基底的基數排序，每個鍵是兩位的 $n$ 進位數，$2$ 趟，每趟 $\Theta(n)$。 (d)：快
速排序不穩定，會破壞前面幾趟的成果，見第~\ref{sec:lin-lsd}~節 $613$、$659$ 的例子。

\item 根是 $1{:}2$（第一個「若」的比較；$a_1 \le a_2$ 往左表示 $m$ 仍是 $1$）。
\begin{itemize}
\item 左子樹（$m = 1$）：比 $1{:}3$。 $a_1 \le a_3$ 時 $m = 1$，不用換，最後比
$2{:}3$，得到 $\langle 1, 2, 3 \rangle$ 或 $\langle 1, 3, 2 \rangle$。 $a_1 > a_3$ 時
$m = 3$，交換後陣列是 $a_3, a_2, a_1$，最後比 $a_1$ 和 $a_2$：$a_2 \le a_1$ 這邊走不
到（已知 $a_1 < a_2$），另一邊交換，得到 $\langle 3, 1, 2 \rangle$。
\item 右子樹（$m = 2$）：比 $2{:}3$。 $a_2 \le a_3$ 時 $m = 2$，交換後陣列是
$a_2, a_1, a_3$，最後比 $a_1$、$a_3$，得到 $\langle 2, 1, 3 \rangle$ 或
$\langle 2, 3, 1 \rangle$。 $a_2 > a_3$ 時 $m = 3$，交換後是 $a_3, a_2, a_1$，最後比
$a_1$、$a_2$：已知 $a_2 < a_1$，所以只走得到 $\langle 3, 2, 1 \rangle$ 那邊，另一邊走
不到。
\end{itemize}
每條路徑都比 $3$ 次，樹有 $2^3 = 8$ 個葉子，其中 $2$ 個走不到，走得到的剛好 $6 = 3!$
個。 最壞、平均都是 $3$ 次，比圖~\ref{fig:lin-dtree} 的平均 $8/3$ 多：這個演算法第三
次比較有時是在問已經知道答案的問題。

\item (a) 答案有 $n$ 種（最大值是哪一個），下界 $\lceil \lg n \rceil$。

(b) 每次比較恰好讓一個數「輸一次」（比較小的那個）。 假設演算法停下來時回答 $z$，
而另外還有一個數 $w \ne z$ 從來沒輸過。 把 $w$ 改成比所有數都大：$w$ 參加過的比較它
本來就都贏，結果不變；其他比較和 $w$ 無關，結果也不變。 演算法看到一樣的結果，仍然回
答 $z$，但現在最大的是 $w$，答錯了。 所以停下來時，除了答案以外的 $n - 1$ 個數都至少
輸過一次，至少要 $n - 1$ 次比較。 從左到右掃一遍剛好用 $n - 1$ 次，所以這個下界是準
的。

(c) 決策樹的論證只用到「答案有 $n$ 種」，每次比較被當成能把可能的答案砍掉一半。 實際
上一次比較最多淘汰一個候選人，砍不了一半，所以資訊理論下界太樂觀。

\item 數完是 $C = [2, 1, 3, 1]$，前綴和之後 $C = [2, 3, 6, 7]$。 從 $j = 7$ 往回放：

\begin{center}
\small
\begin{tabular}{ccl}
\toprule
$j$ & 放到第幾格 & 放完後的 $B[1..7]$ \\
\midrule
7 & 6 & -- -- -- -- -- $2_v$ -- \\
6 & 2 & -- $0_u$ -- -- -- $2_v$ -- \\
5 & 3 & -- $0_u$ $1_t$ -- -- $2_v$ -- \\
4 & 5 & -- $0_u$ $1_t$ -- $2_s$ $2_v$ -- \\
3 & 7 & -- $0_u$ $1_t$ -- $2_s$ $2_v$ $3_r$ \\
2 & 1 & $0_q$ $0_u$ $1_t$ -- $2_s$ $2_v$ $3_r$ \\
1 & 4 & $0_q$ $0_u$ $1_t$ $2_p$ $2_s$ $2_v$ $3_r$ \\
\bottomrule
\end{tabular}
\end{center}
最後的 $C = [0, 2, 3, 6]$，每一格都退到「鍵 $v$ 的第一格的前一格」。

\item 以 $100$ 為基底，四位數是兩「位」，要 $2$ 趟，每趟的計數陣列有 $100$ 格。
第一趟依末兩位（$72, 08, 72, 05, 72, 08$）：$3105, 5408, 1708, 3172, 1772, 5472$。
第二趟依前兩位：$1708, 1772, 3105, 3172, 5408, 5472$。 $1708$ 和 $1772$ 前兩位都是
$17$，第二趟要保持第一趟「$08$ 在 $72$ 前」的順序，靠的是穩定；$3105$ 與 $3172$、
$5408$ 與 $5472$ 也一樣。

\item (a) $n_j$ 是二項分布 $B(n, 1/m)$：$E[n_j] = n/m$，
$\mathrm{Var}[n_j] = \frac{n}{m}(1 - \frac{1}{m})$，所以
$E[n_j^2] = \frac{n}{m}\left(1 - \frac{1}{m}\right) + \frac{n^2}{m^2}$。 $m$ 個桶子加
起來，插入排序的期望時間是
$O\!\left(n\left(1 - \frac{1}{m}\right) + \frac{n^2}{m}\right) = O\!\left(n + \frac{n^2}{m}\right)$。

(b) 整個桶子排序是 $\Theta(n + m) + O(n + n^2/m)$。 $m = n/2$：$O(n)$，仍是線性。
$m = \sqrt{n}$：$O(n^{1.5})$，每個桶子平均有 $\sqrt{n}$ 個數，插入排序在桶子裡變
成平方。

\item (a) $r = \lceil 2{,}000{,}000 / 50{,}000 \rceil = 40$。

(b) $2$ 路：$2^5 = 32 < 40 \le 64$，$6$ 趟。 $5$ 路：$5^2 = 25 < 40 \le 125$，$3$ 趟。
$8$ 路：$8^2 = 64 \ge 40$，$2$ 趟。

(c) 檔案有 $N/B = 20{,}000$ 個區塊，一趟讀寫 $40{,}000$ 次。 產生初始串 $1$ 趟加合併
$2$ 趟，共 $3 \times 40{,}000 = 120{,}000$ 次。

(d) 串的期望長度約 $2M = 100{,}000$，約 $20$ 個串，$5$ 路合併 $\lceil \log_5 20 \rceil
= 2$ 趟，比 (b) 少一趟。

\item (a) $k = 5$，葉子是節點 $5$ 到 $9$：節點 $5, 6, 7, 8, 9$ 依序代表 $R_1$ 到
$R_5$。 節點 $4$ 的子節點是 $8$、$9$，節點 $2$ 的子節點是 $4$、$5$，節點 $3$ 的子節點
是 $6$、$7$。 由下往上：節點 4 是 $17$ 對 $29$，輸家 $R_5$；節點 3 是 $11$ 對 $35$，
輸家 $R_3$；節點 2 是節點 4 的贏家 $17$ 對 $R_1$ 的 $23$，輸家 $R_1$；節點 1 是 $17$ 對
$11$，輸家 $R_4$。 $L[0..4] = R_2, R_4, R_1, R_3, R_5$。

(b) 輸出 $11$（$R_2$ 補 $60$）：節點 3 記的 $35$ 贏 $60$，交換；節點 1 記的 $17$ 贏
$35$，交換；冠軍 $R_4$。 輸出 $17$（$R_4$ 補 $52$）：節點 4 記的 $29$ 贏 $52$，交換；節
點 2 記的 $23$ 贏 $29$，交換；節點 1 記的 $35$ 輸給 $23$，不換；冠軍 $R_1$。 此時
$L[0..4] = R_1, R_3, R_5, R_2, R_4$。

\begin{center}
\begin{tabular}{c@{\hspace{1cm}}c}
\begin{tikzpicture}[x=0.70cm,y=1.00cm]
\draw[lte] (5,-1) -- (3,-2);
\draw[lte] (5,-1) -- (7,-2);
\draw[lte] (3,-2) -- (1,-3);
\draw[lte] (3,-2) -- (4,-3);
\draw[lte] (7,-2) -- (6,-3);
\draw[lte] (7,-2) -- (8,-3);
\draw[lte] (1,-3) -- (0,-4);
\draw[lte] (1,-3) -- (2,-4);
\draw[lte] (5,-1) -- (5,0);
\node[ltwin] at (5,0) {$11$};
\node[rlab, anchor=west] at (5.51429,0) {R2};
\node[ltn] at (5,-1) {$17$};
\node[rlab, anchor=west] at (5.51429,-0.95) {R4};
\node[ltn] at (3,-2) {$23$};
\node[rlab, anchor=west] at (3.51429,-1.95) {R1};
\node[ltn] at (7,-2) {$35$};
\node[rlab, anchor=west] at (7.51429,-1.95) {R3};
\node[ltn] at (1,-3) {$29$};
\node[rlab, anchor=west] at (1.51429,-2.95) {R5};
\node[ltleaf] at (4,-3) {$23$};
\node[rlab] at (4,-3.5) {R1};
\node[restlab] at (4,-3.85) {48};
\node[ltleaf] at (6,-3) {$11$};
\node[rlab] at (6,-3.5) {R2};
\node[restlab] at (6,-3.85) {60 74};
\node[ltleaf] at (8,-3) {$35$};
\node[rlab] at (8,-3.5) {R3};
\node[restlab] at (8,-3.85) {39};
\node[ltleaf] at (0,-4) {$17$};
\node[rlab] at (0,-4.5) {R4};
\node[restlab] at (0,-4.85) {52};
\node[ltleaf] at (2,-4) {$29$};
\node[rlab] at (2,-4.5) {R5};
\node[restlab] at (2,-4.85) {44 81};
\end{tikzpicture}
&
\begin{tikzpicture}[x=0.70cm,y=1.00cm]
\draw[lte] (5,-1) -- (3,-2);
\draw[lte] (5,-1) -- (7,-2);
\draw[lte] (3,-2) -- (1,-3);
\draw[lte] (3,-2) -- (4,-3);
\draw[lte] (7,-2) -- (6,-3);
\draw[lte] (7,-2) -- (8,-3);
\draw[lte] (1,-3) -- (0,-4);
\draw[lte] (1,-3) -- (2,-4);
\draw[lte] (5,-1) -- (5,0);
\node[ltwin] at (5,0) {$23$};
\node[rlab, anchor=west] at (5.51429,0) {R1};
\node[ltn] at (5,-1) {$35$};
\node[rlab, anchor=west] at (5.51429,-0.95) {R3};
\node[ltn] at (3,-2) {$29$};
\node[rlab, anchor=west] at (3.51429,-1.95) {R5};
\node[ltn] at (7,-2) {$60$};
\node[rlab, anchor=west] at (7.51429,-1.95) {R2};
\node[ltn] at (1,-3) {$52$};
\node[rlab, anchor=west] at (1.51429,-2.95) {R4};
\node[ltleaf] at (4,-3) {$23$};
\node[rlab] at (4,-3.5) {R1};
\node[ltleaf] at (6,-3) {$60$};
\node[rlab] at (6,-3.5) {R2};
\node[ltleaf] at (8,-3) {$35$};
\node[rlab] at (8,-3.5) {R3};
\node[ltleaf] at (0,-4) {$52$};
\node[rlab] at (0,-4.5) {R4};
\node[ltleaf] at (2,-4) {$29$};
\node[rlab] at (2,-4.5) {R5};
\end{tikzpicture}
\\[2pt]
{\small (a) 建好之後} & {\small (b) 輸出 $11$、$17$ 之後}
\end{tabular}
\end{center}

(c) 建樹 $k - 1 = 4$ 次。 $R_1$ 的葉子（節點 5）深度 $2$，$R_2$、$R_3$（節點 6、7）深度
$2$，$R_4$、$R_5$（節點 8、9）深度 $3$。 每輸出一筆，就從那個串的葉子往上比一次每一
層：$R_1$ 輸出 $2$ 筆、$R_2$ 輸出 $3$ 筆、$R_3$ 輸出 $2$ 筆，各比 $2$ 次；$R_4$ 輸出 $2$
筆、$R_5$ 輸出 $3$ 筆，各比 $3$ 次。 共 $4 + (2 + 3 + 2) \times 2 + (2 + 3) \times 3 =
33$ 次。 $k$ 不是 $2$ 的次方時，不同的串每次調整的比較次數可能不同。

\item 三個串：$\langle 18, 31, 55, 64, 72, 83 \rangle$、$\langle 9, 12, 27, 35, 40, 50,
66, 71 \rangle$、$\langle 5, 44 \rangle$。 第 2 步輸出 $31$ 後讀入的 $9$、第 3 步輸出
$55$ 後讀入的 $40$ 都凍結；第 6 步輸出 $83$ 後讀入 $50$，記憶體四筆全部凍結，串 1 結
束。 內部排序每次排 $4$ 筆會得到 $16/4 = 4$ 個串，取代選擇少一個，而且前兩個串長
$6$、$8$，都超過 $M$。

\item $N \le M$ 時只有一個串，不用證。 串 $1$ 開始時記憶體是前 $M$ 筆，全部屬於串
$1$。 串 $r \ge 2$ 開始的那一刻，堆積裡最小的串號是 $r$，也就是記憶體裡的紀錄全部屬
於串 $r$。 如果這時記憶體不滿，表示輸入已經讀完，之後不會再讀入紀錄，也就不會出現
串 $r + 1$，串 $r$ 是最後一個。 所以不是最後一個的串，開始時記憶體是滿的，$M$ 筆都
屬於它。 一個串要到記憶體裡沒有它的紀錄才結束，這 $M$ 筆全部會輸出到這個串，所以
它至少有 $M$ 筆。

\item 對手的策略：不管問哪一對，都回答「不相等」。 假設演算法問完不到 $\binom{n}{2}$
次就停下來，那一定有一對 $(i, j)$ 沒問過。 對手手上有兩組都和所有回答一致的輸入：一組
是 $n$ 個兩兩不同的數；另一組只把 $a_j$ 改成和 $a_i$ 一樣，其他不動，這組裡唯一相等的
一對就是沒問過的 $(i, j)$，所以先前每一個「不相等」都仍然成立。 演算法看到同樣的回答，
對兩組輸入給出同樣的答案，但一組沒有重複、一組有，至少錯一組。 所以最壞要問完全部
$\binom{n}{2}$ 對，第~1~章的兩兩比較已經是這個模型裡最好的。 能比大小的話，先排序
（$O(n \lg n)$）再比相鄰的兩個就好；能拿數值當索引的話（查表打勾），只要 $O(n + M)$。
同一個問題，下界取決於演算法被允許問什麼。
\end{enumerate}

\chapter{選擇問題}
\label{ch:selection}

% 本章的圖共用的樣式
\tikzset{
    sel-outcell/.style={cell, fill=black!7, text=black!40, draw=black!25},
    sel-sc/.style={draw=black!55, minimum width=0.53cm, minimum height=0.56cm, inner sep=0pt, font=\ttfamily\scriptsize},
    sel-scout/.style={sel-sc, fill=black!7, text=black!40, draw=black!25},
    sel-scmed/.style={sel-sc, draw=optStroke, thick, fill=optFill},
    sel-scpiv/.style={sel-sc, draw=memStroke, thick, fill=memFill},
    sel-tt/.style={draw=black!55, rounded corners=2pt, minimum width=0.62cm, minimum height=0.55cm, inner sep=0pt, font=\ttfamily\small, fill=white},
    sel-ttw/.style={sel-tt, draw=memStroke, thick, fill=memFill},
    sel-tto/.style={sel-tt, draw=optStroke, thick, fill=optFill},
    sel-te/.style={draw=black!50},
    sel-colc/.style={draw=black!45, minimum width=0.8cm, minimum height=0.55cm, inner sep=0pt, font=\ttfamily\small},
    sel-colle/.style={sel-colc, fill=coreFill, draw=coreStroke},
    sel-colge/.style={sel-colc, fill=greenFill, draw=greenStroke},
    sel-colx/.style={sel-colc, fill=memFill, draw=memStroke, very thick},
}
% 分割追蹤的一列：\selrow{y}{說明}{數值}{p}{r}{q}；p..r 以外的格子是灰的，q 是橘色
\newcommand{\selrow}[6]{%
    \node[anchor=east, font=\small] at (0.35, #1) {#2};
    \foreach \v [count=\k] in {#3} {
        \ifnum\k=#6\relax
            \node[hlcell, minimum width=0.78cm, minimum height=0.7cm] at (\k*0.82, #1) {\v};
        \else\ifnum\k<#4\relax
            \node[sel-outcell, minimum width=0.78cm, minimum height=0.7cm] at (\k*0.82, #1) {\v};
        \else\ifnum\k>#5\relax
            \node[sel-outcell, minimum width=0.78cm, minimum height=0.7cm] at (\k*0.82, #1) {\v};
        \else
            \node[cell, minimum width=0.78cm, minimum height=0.7cm] at (\k*0.82, #1) {\v};
        \fi\fi\fi
    }
}

期中考改完，老師請助教公布全班三百零一位同學成績的中位數，順便公布最高分和最低分。
助教的第一個念頭是把成績排序，再拿正中間那一個。 只要一個數，卻把三百多個數全部排
好，是不是做了太多事？ 最高分和最低分一起找，最少要比幾次？ 往下讀之前，先估計一
下。

這一章講\textbf{選擇問題}：從 $n$ 個數裡找出第 $i$ 小的那一個。 先看 $i = 1$、
$i = n$ 這兩個最簡單的情形，算出同時找最大和最小、找第二大最少要比幾次；再講兩個在
$O(n)$ 時間內解決一般情形的演算法。 第一個每次隨機挑一個數把陣列分成兩邊，只往其中
一邊找，期望時間是線性；第二個用「中位數的中位數」挑分割的基準，最壞情況也是線性。
兩個演算法都建立在第~14~章快速排序的分割（partition）上。

這一章用到第~1~章的插入排序和期望值、第~2~章的代換法、第~12~章的指示變數；C 程式
用到《C 程式設計基礎》第~15~章的遞迴和第~19~章的 \texttt{malloc}、\texttt{free}。

\section{第~\texorpdfstring{$i$}{i}~小、最大與最小}
\label{sec:sel-minmax}

$n$ 個數由小到大排好以後，排在第 $i$ 個的數叫做第 $i$ 小的數，也叫第 $i$ 個
\textbf{順序統計量}（order statistic）。 第 $1$ 小是最小值，第 $n$ 小是最大值。
\textbf{中位數}（median）是正中間的那一個：$n$ 是奇數時是第 $(n+1)/2$ 小；$n$ 是偶
數時中間有兩個，第 $n/2$ 小叫\textbf{下中位數}，第 $n/2 + 1$ 小叫\textbf{上中位數}。
這本書說「中位數」時一律指下中位數，也就是第 $\lfloor (n+1)/2 \rfloor$ 小，$n$ 是奇
數時兩種說法一樣。 統計課在 $n$ 是偶數時取兩個中間數的平均，那是另一個定義，在演算
法裡不用，因為平均值不一定是輸入裡的數。

\textbf{選擇問題}（selection problem）：輸入 $n$ 個數和一個整數 $i$，$1 \le i \le n$，
輸出第 $i$ 小的數。 為了敘述方便，這一章的分析都假設 $n$ 個數互不相同。 有重複時，
把每個數換成（數值，原本的位置）這一對，先比數值、數值一樣再比位置，就變成互不相同
了；演算法本身大多不用改，第~\ref{sec:sel-random}~節最後會說明哪裡要小心。

最直接的做法是先排序再取第 $i$ 個，要 $\Theta(n \lg n)$ 的時間。 也可以用第~8~章的
堆積：由下而上建最小堆積花 $O(n)$，再取出 $i$ 次最小值，每次 $O(\lg n)$，共
$O(n + i \lg n)$；$i$ 很小時很快，但找中位數時還是 $\Theta(n \lg n)$。 這一章的目標
是對任何 $i$ 都做到 $O(n)$。

這一章跟第~15~章一樣，只數「兩個輸入元素之間的比較」：演算法只能透過比較得知元素的
大小，比較次數就是成本。 迴圈變數、名次之間的比較不算。

\subsection{找最大值：\texorpdfstring{$n - 1$}{n - 1} 次，而且不能更少}

從左到右掃一次，記住目前看過最大的數，每個新的數和它比一次，共比 $n - 1$ 次。 這已
經是最少了。 把每一次比較想成一場比賽，大的贏、小的輸。 如果演算法結束時，除了它回
答的最大值 $m$ 之外，還有一個數 $y$ 一場都沒輸過，那麼把 $y$ 改成比所有數都大，所有
比較的結果都不會變（$y$ 每一場都是贏的，改大了還是贏），演算法卻還是回答 $m$，答錯
了。 所以除了最大值，其他 $n - 1$ 個數每個都至少要輸一場；一次比較只產生一個輸家，
至少要比 $n - 1$ 次。

\subsection{同時找最大和最小}

分開找要 $2(n - 1)$ 次。 可以省下大約四分之一：把數兩個兩個一組，先讓同一組的兩個
數比一次，小的那個才有資格當最小值，大的那個才有資格當最大值，所以小的只和目前的最
小值比，大的只和目前的最大值比。 每兩個數花 $3$ 次比較，而不是 $4$ 次。

開頭要看 $n$ 的奇偶。 $n$ 是奇數時，第一個數同時當作目前的最小和最大，剩下的偶數個
兩兩一組；$n$ 是偶數時，前兩個數先比一次決定最小和最大，剩下的兩兩一組。

\begin{pseudo}[label=alg:sel-minmax]{同時找最小和最大}
\Function{Min-Max}{$A, n$}
    \If{$n$ 是奇數}
        \State $\mathit{lo} \gets A[1]$；$\mathit{hi} \gets A[1]$；$s \gets 2$
    \Else
        \State 比較 $A[1]$、$A[2]$，小的給 $\mathit{lo}$、大的給 $\mathit{hi}$；$s \gets 3$
    \EndIf
    \For{$j \gets s$ \To $n - 1$，每次加 $2$}
        \State 比較 $A[j]$、$A[j + 1]$，小的給 $u$、大的給 $v$
        \If{$u < \mathit{lo}$}
            \State $\mathit{lo} \gets u$
        \EndIf
        \If{$v > \mathit{hi}$}
            \State $\mathit{hi} \gets v$
        \EndIf
    \EndFor
    \State \Return $(\mathit{lo}, \mathit{hi})$
\EndFunction
\end{pseudo}

\begin{theorem}
\label{thm:sel-minmax}
$n \ge 1$ 時，\proc{Min-Max} 剛好比較 $\lceil 3n/2 \rceil - 2$ 次。
\end{theorem}

\begin{proof}
第 6 行的迴圈每一輪處理 $A[j]$、$A[j+1]$ 兩個數，比較 $3$ 次（第 7、8、10 行）。

$n$ 是奇數：第 3 行不比較，迴圈處理 $A[2..n]$ 的 $(n - 1)/2$ 組，共
$3(n-1)/2$ 次。 而 $\lceil 3n/2 \rceil = (3n + 1)/2$，所以
$\lceil 3n/2 \rceil - 2 = (3n - 3)/2 = 3(n-1)/2$，相等。

$n$ 是偶數：第 5 行比較 $1$ 次，迴圈處理 $A[3..n]$ 的 $(n-2)/2$ 組，共
$1 + 3(n-2)/2 = 3n/2 - 2 = \lceil 3n/2 \rceil - 2$ 次。

正確性：迴圈每一輪開始時，$\mathit{lo}$ 和 $\mathit{hi}$ 是 $A[1..j-1]$ 的最小值和最
大值（初始時 $j = s$，第 2--5 行正好處理了 $A[1..s-1]$）。 一輪裡，$A[j]$ 和
$A[j+1]$ 之中較大的 $v$ 不可能是新的最小值（$u$ 比它小），較小的 $u$ 也不可能是新的
最大值，所以只要拿 $u$ 更新最小值、$v$ 更新最大值，不變量就維持到 $j + 2$。 迴圈結
束時 $j = n + 1$，$\mathit{lo}$、$\mathit{hi}$ 是整個陣列的最小和最大。
\end{proof}

這個次數也是最少的：任何只用比較的演算法，在最壞情況下都至少要比
$\lceil 3n/2 \rceil - 2$ 次。 證明用的是第~15~章證明合併下界時用過的對手論證（adversary argument）：想像
有一個對手在演算法執行時才決定輸入，每次演算法要比較兩個數，對手就挑一個對演算法最
不利、但和之前的回答不矛盾的結果。 只要對手能逼演算法比很多次，就證明了在最壞情況
下一定要比這麼多次，因為對手回答完所有比較之後，一定能寫出一組和這些回答一致的實際
數值，這組數值就是那個「最壞的輸入」。

\begin{whybox}[同時找最大最小的下界]
用前面「每一場比賽」的說法。 演算法結束時，除了最大值之外的 $n - 1$ 個數，每個都必
須輸過至少一場；除了最小值之外的 $n - 1$ 個數，每個都必須贏過至少一場（理由和找最
大值時一樣：一個從沒贏過的數可以被改成比所有數都小）。 把「某個數輸過」「某個數贏
過」各算一條\textbf{消息}，演算法總共要收集至少 $2n - 2$ 條消息。

把每個數分成四種狀態：沒比過（N）、只贏過（W）、只輸過（L）、贏過也輸過（B）。 對
手這樣回答：
\begin{itemize}
\item N 對 N：隨便回答，兩個數分別變成 W 和 L，得到 $2$ 條消息。
\item N 對 W：W 贏；N 對 L：L 輸；N 對 B：N 贏。 都只有原本是 N 的那個數得到 $1$ 條
新消息。
\item W 對 W：其中一個輸，變成 B，$1$ 條。 L 對 L：其中一個贏，變成 B，$1$ 條。
\item W 對 L：W 贏；W 或 L 對 B：W 贏、L 輸；B 對 B：照目前定好的數值回答。 都沒有新
消息。
\end{itemize}
這些回答不會互相矛盾：對手隨時可以把只贏過的數改得比所有數都大、把只輸過的數改得比
所有數都小，而不影響之前的結果，上面每一條規則都只用到這兩種修改。

只有 N 對 N 的比較能一次得到 $2$ 條消息，而每個數只會當一次 N，所以這種比較最多
$\lfloor n/2 \rfloor$ 次，其他比較每次最多 $1$ 條消息。 要收集 $2n - 2$ 條，比較次數
至少
\[
\left\lfloor \frac{n}{2} \right\rfloor + \left(2n - 2 - 2\left\lfloor \frac{n}{2} \right\rfloor\right)
= 2n - 2 - \left\lfloor \frac{n}{2} \right\rfloor
= \left\lceil \frac{3n}{2} \right\rceil - 2 .
\]
最後一步用到 $2n - \lfloor n/2 \rfloor = \lceil 3n/2 \rceil$。
\end{whybox}

\subsection{找第二大：錦標賽}

第二大的數只輸給最大的那一個，所以它一定\textbf{直接}和最大值比過：如果它從沒和最大
值比過，又被判定不是最大，它一定輸給了某個別的數，那個數比它大，它就不是第二大了。
（演算法也必須確認它不是最大，所以它一定至少輸過一場。）

這引出\textbf{錦標賽}（tournament）的做法：

\begin{enumerate}
\item 把 $n$ 個數兩兩配對比賽，贏的晉級；人數是奇數時，最後一個人輪空，直接晉級。
一直比到剩下一個人，就是最大值。 每場比賽淘汰一個人，共 $n - 1$ 場。
\item 第二大一定在「直接輸給冠軍的人」裡面。 冠軍每一輪最多比一場，總共
$\lceil \lg n \rceil$ 輪（每一輪人數變成 $\lceil m/2 \rceil$，從 $n$ 變到 $1$ 要
$\lceil \lg n \rceil$ 輪），所以這些人最多 $\lceil \lg n \rceil$ 個，在他們之中找最大，
最多比 $\lceil \lg n \rceil - 1$ 次。
\end{enumerate}

\begin{center}
\begin{tikzpicture}[x=1cm, y=0.95cm]
\node[sel-tt] (L1) at (0.000,0.00) {47};
\node[sel-tt] (L2) at (1.050,0.00) {65};
\node[sel-tto] (L3) at (2.100,0.00) {88};
\node[sel-ttw] (L4) at (3.150,0.00) {91};
\node[sel-tt] (L5) at (4.200,0.00) {38};
\node[sel-tt] (L6) at (5.250,0.00) {72};
\node[sel-tt] (L7) at (6.300,0.00) {56};
\node[sel-tt] (L8) at (7.350,0.00) {83};
\node[sel-tt] (L9) at (8.400,0.00) {60};
\node[sel-tt] (L10) at (9.450,0.00) {79};
\node[sel-tt] (L11) at (10.500,0.00) {74};
\node[sel-tto] (R1_1) at (0.525,1.00) {65};
\node[sel-ttw] (R1_2) at (2.625,1.00) {91};
\node[sel-tt] (R1_3) at (4.725,1.00) {72};
\node[sel-tt] (R1_4) at (6.825,1.00) {83};
\node[sel-tt] (R1_5) at (8.925,1.00) {79};
\node[sel-tt] (R1_6) at (10.500,1.00) {74};
\node[sel-ttw] (R2_1) at (1.575,2.00) {91};
\node[sel-tto] (R2_2) at (5.775,2.00) {83};
\node[sel-tt] (R2_3) at (9.713,2.00) {79};
\node[sel-ttw] (R3_1) at (3.675,3.00) {91};
\node[sel-tto] (R3_2) at (9.713,3.00) {79};
\node[sel-ttw] (R4_1) at (6.694,4.00) {91};
\draw[sel-te] (R1_1) -- (L1);
\draw[sel-te] (R1_1) -- (L2);
\draw[sel-te] (R1_2) -- (L3);
\draw[sel-te] (R1_2) -- (L4);
\draw[sel-te] (R1_3) -- (L5);
\draw[sel-te] (R1_3) -- (L6);
\draw[sel-te] (R1_4) -- (L7);
\draw[sel-te] (R1_4) -- (L8);
\draw[sel-te] (R1_5) -- (L9);
\draw[sel-te] (R1_5) -- (L10);
\draw[sel-te] (R1_6) -- (L11);
\draw[sel-te] (R2_1) -- (R1_1);
\draw[sel-te] (R2_1) -- (R1_2);
\draw[sel-te] (R2_2) -- (R1_3);
\draw[sel-te] (R2_2) -- (R1_4);
\draw[sel-te] (R2_3) -- (R1_5);
\draw[sel-te] (R2_3) -- (R1_6);
\draw[sel-te] (R3_1) -- (R2_1);
\draw[sel-te] (R3_1) -- (R2_2);
\draw[sel-te] (R3_2) -- (R2_3);
\draw[sel-te] (R4_1) -- (R3_1);
\draw[sel-te] (R4_1) -- (R3_2);
\foreach \t/\lab in {0/參賽, 1/第 1 輪, 2/第 2 輪, 3/第 3 輪, 4/第 4 輪} {
    \node[anchor=east, font=\small] at (-0.5, \t) {\lab};
}
\end{tikzpicture}
\diagcaption{$11$ 個數的錦標賽，每個節點寫的是那場比賽的贏家，只有一個子節點的是輪
空。 橘色是冠軍 $91$ 走過的路，粉紅色是直接輸給它的四個數 $88, 65, 83, 79$。 決賽輸
給冠軍的 $79$ 不是第二大，第一輪就輸給冠軍的 $88$ 才是。}
\label{fig:sel-tournament}
\end{center}

\begin{theorem}
\label{thm:sel-second}
$n \ge 2$ 時，錦標賽在最壞情況下用 $n + \lceil \lg n \rceil - 2$ 次比較找出第二大的數。
\end{theorem}

\begin{proof}
第一步 $n - 1$ 次。 第二步：冠軍在每一輪不是比一場就是輪空，最多
$\lceil \lg n \rceil$ 輪，所以直接輸給它的人 $c \le \lceil \lg n \rceil$ 個；在 $c$ 個
數裡找最大要 $c - 1$ 次。 合計不超過 $n - 1 + \lceil \lg n \rceil - 1$。 正確性就是
本小節開頭的論證：第二大在這 $c$ 個人裡面，而這 $c$ 個人都比冠軍小，其中最大的就是
第二大。
\end{proof}

圖~\ref{fig:sel-tournament} 的 $11$ 個數：第一步比 $5 + 3 + 1 + 1 = 10$ 場，冠軍比了
$4$ 場，四個人裡找最大再比 $3$ 次，共 $13 = 11 + \lceil \lg 11 \rceil - 2$ 次。 冠軍
輪空的次數多時，比較會更少（觀念題第~\ref{sel-q-tour}~題）。

這個次數在最壞情況下也是最少的。 下面只寫證明的骨架。

\begin{whybox}[找第二大的下界（骨架）]
任何找第二大的演算法，除了最大和第二大，其他 $n - 2$ 個數都要輸給某個不是最大值的
數（否則它可能是第二大）；再加上第二大本身要輸給最大值。 設最大值總共直接打敗了
$c$ 個數。 這 $c$ 個數裡，除了第二大以外的 $c - 1$ 個，都還要再輸一場給別人，所以
「輸」的總場數至少是 $(n - 1) + (c - 1)$，每場比較只有一個輸家，比較次數至少
$n + c - 2$。

接著用對手論證逼 $c \ge \lceil \lg n \rceil$。 對手給每個數一個權重，一開始都是 $1$。
兩個都沒輸過的數比賽時，權重大的贏，贏家把輸家的權重加到自己身上，輸家的權重變成
$0$；有一方輸過時，讓沒輸過的那一方贏；兩方都輸過時，回答任何和之前一致的結果。 權
重的總和一直是 $n$，最後全部集中在最大值身上。 最大值每贏一個權重為正的對手，自己
的權重最多變成兩倍（對手的權重不比它大），從 $1$ 變成 $n$ 至少要贏
$\lceil \lg n \rceil$ 場，而且每一場的對手都不一樣（輸過的數權重是 $0$）。 所以
$c \ge \lceil \lg n \rceil$，比較次數至少 $n + \lceil \lg n \rceil - 2$。
\end{whybox}

\begin{examnote}
錦標賽的題目，答案要畫出整棵比賽樹（輪空的人照題目的規則處理，沒有說明時讓最後一個
人輪空），標出冠軍一路打敗的人，再在他們之中找最大。 常見的錯是直接拿決賽的輸家當
第二大，圖~\ref{fig:sel-tournament} 就是反例。 比較次數要分開寫「找冠軍」和「在候選
人中找最大」兩部分。
\end{examnote}


\section{隨機化選擇}
\label{sec:sel-random}

一般的 $i$ 要怎麼找？ 第~14~章的快速排序先挑一個基準（pivot），把比它小的
放左邊、比它大的放右邊，再兩邊各自遞迴排序。 分割完之後，基準所在的位置 $q$ 就是它
排序後的位置：左邊有 $q - p$ 個數比它小。 找第 $i$ 小時只需要知道答案在哪一邊，只往
那一邊遞迴，另一邊丟掉不管。 這個想法是 Hoare 在 1961 年提出的，常叫做
\textbf{快速選擇}（quickselect）。

分割直接用第~14~章第~\ref{sec:so-quick}~節的兩個程序。 \proc{Lomuto-Partition}（虛擬
碼~\ref{alg:so-lomuto}）以 $A[r]$ 為基準，結束時 $A[p..q-1]$ 都 $\le A[q]$，$A[q+1..r]$
都 $> A[q]$，回傳 $q$；過程中基準和其他 $r - p$ 個數各比一次。
\proc{Randomized-Partition}（虛擬碼~\ref{alg:so-rand}）先把 $A[p..r]$ 裡隨機挑的一個數
換到 $A[r]$，再做 Lomuto 分割。 \proc{Random}$(p, r)$ 傳回 $p$ 到 $r$ 之間的一個整數，
每個整數的機率都一樣，每次呼叫彼此獨立。

\begin{pseudo}[label=alg:sel-rselect]{隨機化選擇}
\Function{Randomized-Select}{$A, p, r, i$}
    \If{$p = r$}
        \State \Return $A[p]$ \Comment{只剩一個數，$i$ 一定是 $1$}
    \EndIf
    \State $q \gets$ \Call{Randomized-Partition}{$A, p, r$}
    \State $k \gets q - p + 1$ \Comment{基準在 $A[p..r]$ 裡是第 $k$ 小}
    \If{$i = k$}
        \State \Return $A[q]$
    \ElsIf{$i < k$}
        \State \Return \Call{Randomized-Select}{$A, p, q - 1, i$}
    \Else
        \State \Return \Call{Randomized-Select}{$A, q + 1, r, i - k$}
    \EndIf
\EndFunction
\end{pseudo}

呼叫 \proc{Randomized-Select}$(A, 1, n, i)$ 找整個陣列第 $i$ 小的數。 往右邊找時，左
邊的 $k$ 個數（包括基準）都比答案小，所以在右邊要找的是第 $i - k$ 小，這是第 11 行
最容易漏掉的地方。 往左邊找時名次不變。

\textbf{正確性}對 $n = r - p + 1$ 用強歸納法（觀念題第~\ref{sel-q-correct}~題）。 關
鍵是分割之後 $A[q]$ 前面恰好有 $k - 1$ 個數，而且都比它小（互不相同的假設下），後面
的都比它大，所以 $A[q]$ 是第 $k$ 小；$i < k$ 時答案在左邊、名次還是 $i$，$i > k$ 時
答案在右邊、名次是 $i - k$。 遞迴呼叫的範圍不會是空的：$i < k$ 時左邊有
$k - 1 \ge i \ge 1$ 個數，$i > k$ 時右邊有 $n - k \ge i - k \ge 1$ 個數。

\subsection{手算追蹤：成績的中位數}
\label{sec:sel-trace}

考卷上的選擇題通常不用隨機，而是規定基準「固定取最後一個」，答案才唯一。 這相當於
把第 4 行的 \proc{Randomized-Partition} 換成 \proc{Lomuto-Partition}。 下面用這個版本，找
$13$ 位同學成績的中位數，也就是第 $7$ 小。

\begin{center}
\begin{tikzpicture}
\selrow{0.00}{開始}{59,46,47,80,64,90,62,96,73,58,79,44,54}{1}{13}{0}
\selrow{-0.85}{第 1 次}{46,47,44,54,64,90,62,96,73,58,79,59,80}{1}{13}{4}
\selrow{-1.70}{第 2 次}{46,47,44,54,64,62,73,58,79,59,80,96,90}{5}{13}{11}
\selrow{-2.55}{第 3 次}{46,47,44,54,58,59,73,64,79,62,80,96,90}{5}{10}{6}
\selrow{-3.40}{第 4 次}{46,47,44,54,58,59,62,64,79,73,80,96,90}{7}{10}{7}
\foreach \k in {1,...,13} { \node[idxlbl] at (\k*0.82, -4.05) {\k}; }
\end{tikzpicture}
\diagcaption{基準固定取最後一個，找 $13$ 個成績的第 $7$ 小。 每一列是那一次
\proc{Lomuto-Partition} 之後的陣列：白色格子是這一次分割的範圍 $A[p..r]$，橘色是基準最後的位
置 $q$，灰色是已經丟掉、不會再看的部分。}
\label{fig:sel-score-trace}
\end{center}

\begin{center}
\begin{tabular}{cccccccl}
\toprule
次 & $p$ & $r$ & $i$ & 基準 $A[r]$ & $q$ & $k$ & 下一步 \\
\midrule
1 & 1 & 13 & 7 & 54 & 4 & 4 & $i > k$：往右，$i \gets 7 - 4 = 3$ \\
2 & 5 & 13 & 3 & 80 & 11 & 7 & $i < k$：往左，$i$ 不變 \\
3 & 5 & 10 & 3 & 59 & 6 & 2 & $i > k$：往右，$i \gets 3 - 2 = 1$ \\
4 & 7 & 10 & 1 & 62 & 7 & 1 & $i = k$：回傳 $62$ \\
\bottomrule
\end{tabular}
\end{center}

以第 2 次為例：範圍 $A[5..13]$ 的基準是 $80$，比它小的有 $6$ 個，$80$ 落在第 $11$ 格，是
這個範圍的第 $k = 11 - 5 + 1 = 7$ 小；要找的是第 $3$ 小，所以往左邊 $A[5..10]$ 找，名次
不變。 第 1 次往右時，左邊連同基準的 $4$ 個數都比答案小，名次從 $7$ 變成 $3$。

驗算：排好的成績是 $44, 46, 47, 54, 58, 59, 62, 64, 73, 79, 80, 90, 96$，第 $7$ 個是
$62$。 四次分割分別比了 $12, 8, 5, 3$ 次，共 $28$ 次；用插入排序整個排好，這組資料要
比 $49$ 次（可以用第~1~章插入排序的程式驗證）。

\begin{examnote}
選擇的追蹤題，每次分割寫一列：分割之後的整個陣列、這次的範圍 $p..r$、基準、$q$、
$k = q - p + 1$，以及往哪邊走、新的 $i$。 往右走時 $i$ 要減掉 $k$，往左走時不變，閱卷
時這一欄最容易看出有沒有懂。 題目如果規定別的分割方法（例如第~14~章的 Hoare 分割，
或取第一個當基準），陣列的樣子會不同，照題目規定的分割來做。 名次 $i$ 從 $1$ 開始還
是從 $0$ 開始、中位數取上還是取下，也先看題目的定義。
\end{examnote}

\subsection{最壞情況}

每次都挑到範圍裡最大的數當基準，而要找的是最小值，範圍每次只少一個數，比較次數是
$(n-1) + (n-2) + \cdots + 1 = n(n-1)/2$，最壞情況是 $\Theta(n^2)$。 固定取最後一個當
基準時，已經由小到大排好的陣列找最小值就是這種情況。 隨機挑基準時，沒有哪一個輸入
一定會這樣，要連續好幾次都運氣很差才會發生（觀念題第~\ref{sel-q-worst}~題）。 下面證
明平均起來是線性的。

\subsection{期望時間是線性的}

先說清楚「期望」對什麼取。 這裡的輸入是固定的（任何一組 $n$ 個互不相同的數、任何
一個 $i$），隨機性只來自 \proc{Random} 的結果。 所以下面的結論對\textbf{每一個}輸入
都成立，不需要假設輸入是隨機排列的；這和第~1~章插入排序的平均情況不一樣，那裡是對
所有輸入取平均。

\begin{theorem}
\label{thm:sel-rselect}
對 $n$ 個互不相同的數和任何 $i$，\proc{Randomized-Select} 的期望執行時間是
$\Theta(n)$。
\end{theorem}

\begin{proof}
下界：第一次分割就要看過全部 $n$ 個數，所以至少 $\Omega(n)$。

上界。 令 $T(n)$ 是所有大小為 $n$ 的輸入、所有 $i$ 之中，期望執行時間的最大值；再令
$U(n) = \max_{1 \le m \le n} T(m)$，$U(0) = 0$。 $U$ 是不遞減的，而且
$T(n) \le U(n)$。 取常數 $a$ 使得對 $A[p..r]$ 做一次分割、加上遞迴呼叫以外的其他工作，
都不超過 $a n$（$n = r - p + 1 \ge 2$）；$U(1) = T(1)$ 是常數 $b$。

$n \ge 2$ 時，基準是範圍內每一個數的機率都是 $1/n$，所以它在範圍內是第 $k$ 小的機率
也是 $1/n$，$k = 1, \ldots, n$。 已知基準是第 $k$ 小時，接下來最多遞迴一次，範圍大
小是 $k - 1$ 或 $n - k$，之後用到的隨機數是新的，所以這部分的期望時間不超過
$U(\max(k - 1, n - k))$。 對 $k$ 取平均：
\[
T(n) \le a n + \frac{1}{n} \sum_{k=1}^{n} U(\max(k - 1, n - k)).
\]
$k$ 從 $1$ 跑到 $n$ 時，$\max(k - 1, n - k)$ 取到 $\lfloor n/2 \rfloor, \ldots, n-1$
裡的值，每個值最多出現兩次，所以
\[
T(n) \le a n + \frac{2}{n} \sum_{m=\lfloor n/2 \rfloor}^{n-1} U(m).
\]

用代換法（第~2~章第~\ref{sec:rec-subst}~節）證明：取 $c = \max(4a, b)$，所有 $n \ge 1$ 都有
$U(n) \le cn$。 $n = 1$ 時 $U(1) = b \le c$。 設 $n \ge 2$，比 $n$ 小的都成立。 先估
計和：
\[
\sum_{m=\lfloor n/2 \rfloor}^{n-1} m
= \frac{n(n-1)}{2} - \frac{(\lfloor n/2 \rfloor - 1)\lfloor n/2 \rfloor}{2}
\le \frac{n(n-1)}{2} - \frac{(n-3)(n-1)}{8}
= \frac{3(n^2 - 1)}{8} < \frac{3n^2}{8},
\]
中間用到 $\lfloor n/2 \rfloor \ge (n-1)/2 \ge 1/2$，以及 $t(t-1)$ 在 $t \ge 1/2$ 時隨 $t$ 遞增。
代進去：
\[
T(n) \le a n + \frac{2c}{n} \cdot \frac{3n^2}{8} = \frac{3}{4} c n + a n \le c n,
\]
最後一步用到 $c \ge 4a$。 再加上 $U(n - 1) \le c(n-1) \le cn$，得到
$U(n) = \max(U(n-1), T(n)) \le cn$。 所以期望時間 $T(n) \le U(n) = O(n)$。
\end{proof}

證明裡要的是 $U(n) \le cn$ 這個\textbf{同一個} $c$，推出來的是
$\frac{3}{4}cn + an$，多出來的 $an$ 剛好被省下的 $\frac{1}{4}cn$ 吸收。 這就是第~2~章
提醒過的「歸納步驟要推出同一個常數」。 直覺上，基準落在中間一半的機率是 $1/2$，落
在那裡時範圍至少縮成 $3/4$；平均兩次分割就有一次「好的」，總工作量是一個公比小於
$1$ 的等比級數。

\subsection{比較次數算到常數}

上面的證明只說是 $O(n)$。 用第~12~章第~\ref{sec:hash-chain-analysis}~節的指示變數，可以把期望比較次數直接寫
成一個和，再算出常數。

\begin{theorem}
\label{thm:sel-rselect-cmp}
把 $n$ 個互不相同的數由小到大記作 $z_1 < z_2 < \cdots < z_n$。
\proc{Randomized-Select} 找第 $i$ 小時，期望比較次數是
\[
E_n(i) = \sum_{1 \le a < b \le n} \frac{2}{\max(b, i) - \min(a, i) + 1},
\]
而且 $E_n(i) < 4n$。
\end{theorem}

\begin{whybox}[期望比較次數]
比較只發生在分割時，基準和範圍裡的其他數各比一次；當過基準的數之後就不在任何範圍
裡了，所以每一對 $z_a, z_b$ 最多比一次。 令 $X_{ab}$ 是「$z_a$ 和 $z_b$ 比過」的指
示變數，總比較次數是 $\sum_{a<b} X_{ab}$，期望值是 $\sum_{a<b} \Pr[X_{ab} = 1]$。

固定 $a < b$，令 $Z$ 是名次從 $\min(a, i)$ 到 $\max(b, i)$ 的那些數，共
$\max(b, i) - \min(a, i) + 1$ 個；$Z$ 同時包含 $z_a$、$z_b$ 和答案 $z_i$。
\begin{itemize}
\item 還沒有 $Z$ 裡的數被選成基準時，$Z$ 整個都在目前的範圍裡：一個 $Z$ 以外的基準
不是比 $Z$ 全部小，就是比 $Z$ 全部大，而演算法留下的是含有 $z_i$ 的那一邊，也就是
含有整個 $Z$ 的那一邊。
\item 所以第一個落在 $Z$ 裡的基準，在 $Z$ 的每個數上機率都一樣，是 $z_a$ 或 $z_b$ 的
機率是 $2/|Z|$。 是的話，$z_a$、$z_b$ 這時都在範圍裡，比了一次。
\item 如果它是 $Z$ 裡的其他數 $z_c$：$a < c < b$ 時，$z_a$ 和 $z_b$ 被分到兩邊（或
$c = i$，演算法直接結束），以後不會再比；$c$ 在 $a$、$b$ 的範圍外時，$c$ 夾在 $i$ 和
這一對之間（或 $c = i$），留下的那一邊含有 $z_i$ 而不含 $z_a$、$z_b$，也不會再比。
\end{itemize}
所以 $\Pr[X_{ab} = 1] = 2/(\max(b,i) - \min(a,i) + 1)$，加起來就是 $E_n(i)$。

上界分三種 $(a, b)$ 估計：
\begin{itemize}
\item $a \le i \le b$：分母是 $b - a + 1 = L$。 長度為 $L$、包含 $i$ 的區間最多 $L$ 個，
每個貢獻 $2/L$，所以長度 $L$ 的總貢獻不超過 $2$；$L = 2, \ldots, n$，合計不超過
$2(n - 1)$。
\item $i < a < b$：分母是 $b - i + 1$。 固定 $b$，$a$ 有 $b - i - 1$ 種，貢獻
$2(b - i - 1)/(b - i + 1) < 2$；$b$ 最多 $n - i$ 種，合計小於 $2(n - i)$。
\item $a < b < i$：同理，合計小於 $2(i - 1)$。
\end{itemize}
三項加起來小於 $2(n-1) + 2(n-i) + 2(i-1) < 4n$。
\end{whybox}

找最小值（$i = 1$）時，這個和可以算出來是 $2n - 2H_n$（觀念題第~\ref{sel-q-min}~題），
比直接掃一次的 $n - 1$ 多了將近一倍：隨機選擇不是用來找最大最小的。 找中位數時和最
大，$n$ 很大時大約是 $3.39n$。 下面的程式實際跑一次，和公式對照。

\subsection{C 程式}

C 標準函式庫的 \texttt{rand} 在不同的系統上產生的數列不一樣，書上印的輸出就對不起來。
所以程式自己寫一個\textbf{線性同餘產生器}（linear congruential generator，LCG）：
$s \gets s \cdot M + C \bmod 2^{64}$，$M$、$C$ 是固定的常數，種子 $s$ 一開始固定，
每次執行產生的數列都一樣。 LCG 的低位元週期很短，所以取高位元當結果。
\texttt{random\_between(p, r)} 用取餘數把結果縮到 $p..r$，會有一點點不均勻（《C 程式
設計基礎》第~18~章講的模偏差），範圍遠小於 $2^{31}$ 時可以忽略。

\begin{cbookrecap}[第~13~章 13.1、13.5 節，第~18~章 18.8 節]
有號整數溢位是未定義行為；\texttt{unsigned} 型別的運算保證是模 $2^n$，所以 LCG 的乘
法和加法要用 \texttt{unsigned long long}，溢位就是我們要的「模 $2^{64}$」。 寫在函式
外面的 \texttt{static} 變數只在這個檔案裡看得到，而且整個程式執行期間都存在，適合存
產生器的狀態。 \texttt{rand()} 的數列由實作決定，換一台機器就可能不同。
\end{cbookrecap}

程式先用隨機分割找第~\ref{sec:sel-trace}~節那 $13$ 個成績的中位數；再拿 $1001$ 個不同的數，同一
個輸入跑 $2000$ 次，算三個名次的平均比較次數，和定理~\ref{thm:sel-rselect-cmp} 的公式
比較。 \texttt{ncmp} 只在分割的迴圈裡加一，正好是元素之間的比較次數。

\begin{example}[label=ex:sel-rselect]{c}{(rselect.c) 隨機化選擇與期望比較次數}
#include <stdio.h>
#include <stdlib.h>

// 自己寫的線性同餘產生器（LCG），種子固定，每次執行的結果都一樣。
static unsigned long long state = 20261119ULL;

int random_between(int p, int r)            // p..r 之間的一個整數
{
    state = state * 6364136223846793005ULL + 1442695040888963407ULL;
    // 低位元的週期短，取高位元
    unsigned x = (unsigned)(state >> 33);
    return p + (int)(x % (unsigned)(r - p + 1));
}

long long ncmp;                             // 元素之間的比較次數

void swap(int *x, int *y)
{
    int t = *x;
    *x = *y;
    *y = t;
}

// 隨機挑一個數換到 a[r] 當基準，做 Lomuto 分割，回傳基準最後的位置。
int randomized_partition(int *a, int p, int r)
{
    swap(&a[r], &a[random_between(p, r)]);
    int x = a[r], i = p - 1;
    for (int j = p; j < r; j++) {
        ncmp++;
        if (a[j] <= x)
            swap(&a[++i], &a[j]);
    }
    swap(&a[i + 1], &a[r]);
    return i + 1;
}

// 回傳 a[p..r] 裡第 i 小的數，1 <= i <= r - p + 1。
int randomized_select(int *a, int p, int r, int i)
{
    if (p == r)
        return a[p];
    int q = randomized_partition(a, p, r);
    int k = q - p + 1;
    if (i == k)
        return a[q];
    else if (i < k)
        return randomized_select(a, p, q - 1, i);
    else
        return randomized_select(a, q + 1, r, i - k);
}

// 定理的公式：n 個不同的數找第 k 小時的期望比較次數。
double expected_cmp(int n, int k)
{
    double s = 0;
    for (int a = 1; a <= n; a++)
        for (int b = a + 1; b <= n; b++)
            s += 2.0 / ((b > k ? b : k) - (a < k ? a : k) + 1);
    return s;
}

int main(void)
{
    int score[] = {59, 46, 47, 80, 64, 90, 62, 96, 73, 58, 79,
                   44, 54};
    int med = randomized_select(score, 0, 12, 7);
    printf("median = %d, comparisons = %lld\n", med, ncmp);

    // 1001 個不同的數，同一個輸入跑 2000 次，看平均比較次數。
    int N = 1001, T = 2000, ks[] = {1, 251, 501};
    int *base = malloc(N * sizeof *base);
    int *b = malloc(N * sizeof *b);
    if (base == NULL || b == NULL)
        return 1;
    for (int j = 0; j < N; j++)
        base[j] = j + 1;
    // 洗牌，得到 1..N 的一個排列
    for (int j = N - 1; j > 0; j--)
        swap(&base[j], &base[random_between(0, j)]);
    printf("k average formula 4n\n");
    for (int q = 0; q < 3; q++) {
        long long total = 0;
        for (int t = 0; t < T; t++) {
            for (int j = 0; j < N; j++)
                b[j] = base[j];
            ncmp = 0;
            if (randomized_select(b, 0, N - 1, ks[q]) != ks[q])
                printf("WRONG\n");
            total += ncmp;
        }
        printf("%d %.1f %.1f %d\n", ks[q], (double)total / T,
               expected_cmp(N, ks[q]), 4 * N);
    }
    free(base);
    free(b);
    return 0;
}
\end{example}
\iosample{}{median = 62, comparisons = 22 \\ k average formula 4n \\ 1 1980.5 1987.0 4004 \\ 251 3055.8 3079.8 4004 \\ 501 3349.0 3339.7 4004}

同一組成績，隨機挑基準這一次比了 $22$ 次，比固定取最後一個的 $28$ 次少；換一個種子
結果會不同，但答案一定是 $62$。 實驗的部分，$2000$ 次的平均和公式差不到百分之一，三
個名次都小於 $4n$；名次越靠中間，比較越多，中位數大約是 $3.34n$。

\subsection{有重複的數時}

前面的分析假設數互不相同。 \proc{Randomized-Select} 遇到重複的數仍然回傳正確的答
案（分割後左邊都 $\le$ 基準、右邊都 $>$ 基準，論證照樣成立），但時間可能變差。 最極
端的是 $n$ 個數全部一樣：Lomuto 分割每次都把基準放到範圍的最後一格，找中位數要
$\Theta(n^2)$，隨機挑基準也沒有用。 解法是改成第~14~章的三路分割，把範圍分成小於、
等於、大於基準三段，$i$ 落在等於那一段就直接回傳（觀念題第~\ref{sel-q-dup}~題）。

\section{中位數的中位數：最壞情況線性}
\label{sec:sel-mom}

隨機化選擇的期望是線性，最壞還是 $\Theta(n^2)$。 1973 年 Blum、Floyd、Pratt、Rivest、
Tarjan 給了一個最壞情況也是 $O(n)$ 的演算法。 前一節的論證已經指出關鍵：只要基準保
證不太靠兩端，每次都能丟掉固定比例的數，總工作量就是公比小於 $1$ 的等比級數。 問題
是怎麼不靠運氣、又只花線性時間，就挑到這樣的基準。

做法如下（$n \le 5$ 時直接排序）：

\begin{enumerate}
\item 把 $A[p..r]$ 從左到右每 $5$ 個分成一組，最後一組可能不滿 $5$ 個，共
$g = \lceil n/5 \rceil$ 組。
\item 每組用插入排序排好，取出每組的中位數（不滿的那組取下中位數），放進新陣列
$M[1..g]$。
\item \textbf{遞迴}找出 $M$ 的中位數 $x$，這就是「中位數的中位數」。
\item 用 $x$ 當基準分割 $A[p..r]$，接下來和 \proc{Randomized-Select} 一樣，只往一邊遞
迴。
\end{enumerate}

\begin{pseudo}[label=alg:sel-mom, float=!tp]{中位數的中位數}
\Function{Select}{$A, p, r, i$}
    \State $n \gets r - p + 1$
    \If{$n \le 5$}
        \State 用插入排序把 $A[p..r]$ 排好
        \State \Return $A[p + i - 1]$
    \EndIf
    \State $g \gets \lceil n/5 \rceil$
    \State 建立陣列 $M[1..g]$
    \For{$j \gets 1$ \To $g$}
        \State $\ell \gets p + 5(j - 1)$
        \State $h \gets \min(\ell + 4, r)$
        \State 用插入排序把 $A[\ell..h]$ 排好
        \State $M[j] \gets A[\lfloor (\ell + h)/2 \rfloor]$ \Comment{這一組的（下）中位數}
    \EndFor
    \State $x \gets$ \Call{Select}{$M, 1, g, \lceil g/2 \rceil$}
    \State $q \gets$ \Call{Partition-Around}{$A, p, r, x$}
    \State $k \gets q - p + 1$
    \If{$i = k$}
        \State \Return $A[q]$
    \ElsIf{$i < k$}
        \State \Return \Call{Select}{$A, p, q - 1, i$}
    \Else
        \State \Return \Call{Select}{$A, q + 1, r, i - k$}
    \EndIf
\EndFunction
\Statex
\Function{Partition-Around}{$A, p, r, x$}
    \State 找出 $A[p..r]$ 裡等於 $x$ 的位置 $j$
    \State 交換 $A[j]$ 和 $A[r]$
    \State \Return \Call{Lomuto-Partition}{$A, p, r$}
\EndFunction
\end{pseudo}

一組 $h - \ell + 1$ 個數排好之後，下中位數在第 $\lceil (h - \ell + 1)/2 \rceil$ 個，
位置是 $\ell + \lfloor (h - \ell)/2 \rfloor = \lfloor (\ell + h)/2 \rfloor$，這就是第
12 行。 第 13 行對 $M$ 遞迴，$M$ 只有 $\lceil n/5 \rceil$ 個數；第 14 行花 $O(n)$
找到 $x$ 的位置，再交給 Lomuto 分割。 第 4、11 行排的都是不超過 $5$ 個數，每次比較
不超過 $10$ 次。

\subsection{\texorpdfstring{$x$}{x} 保證不靠兩端}

把每一組排好、立起來當成一欄，小的在下、大的在上，再把各欄依照中位數由小到大排好，
如圖~\ref{fig:sel-mom-cols}。 中間那一列是各組的中位數，$x$ 在正中間。 左下方藍色的
數：它們所在的那一欄的中位數 $\le x$，而它們又 $\le$ 那一欄的中位數，所以都 $\le x$。
同理，右上方綠色的數都 $\ge x$。 不管其他數怎麼排，這兩塊都保證存在。

\begin{figure}[tbp]
\centering
\begin{tikzpicture}
\node[sel-colle] at (0.00,0.00) {16};
\node[sel-colle] at (0.00,0.62) {27};
\node[sel-colle] at (0.00,1.24) {31};
\node[sel-colc] at (0.00,1.86) {70};
\node[sel-colc] at (0.00,2.48) {87};
\node[idxlbl] at (0.00,-0.55) {第 4 組};
\node[sel-colle] at (0.95,0.00) {13};
\node[sel-colle] at (0.95,0.62) {30};
\node[sel-colle] at (0.95,1.24) {35};
\node[sel-colc] at (0.95,1.86) {66};
\node[sel-colc] at (0.95,2.48) {85};
\node[idxlbl] at (0.95,-0.55) {第 2 組};
\node[sel-colle] at (1.90,0.00) {15};
\node[sel-colle] at (1.90,0.62) {33};
\node[sel-colx] at (1.90,1.24) {38};
\node[sel-colge] at (1.90,1.86) {88};
\node[sel-colge] at (1.90,2.48) {98};
\node[idxlbl] at (1.90,-0.55) {第 3 組};
\node[sel-colc] at (2.85,0.00) {24};
\node[sel-colc] at (2.85,0.62) {34};
\node[sel-colge] at (2.85,1.24) {50};
\node[sel-colge] at (2.85,1.86) {71};
\node[sel-colge] at (2.85,2.48) {86};
\node[idxlbl] at (2.85,-0.55) {第 5 組};
\node[sel-colc] at (3.80,0.00) {44};
\node[sel-colc] at (3.80,0.62) {48};
\node[sel-colge] at (3.80,1.24) {80};
\node[sel-colge] at (3.80,1.86) {89};
\node[sel-colge] at (3.80,2.48) {91};
\node[idxlbl] at (3.80,-0.55) {第 1 組};
\draw[diagarrow] (-0.85,-0.2) -- node[left, font=\scriptsize, align=center] {由小\\到大} (-0.85,2.7);
\draw[diagarrow] (-0.4,3.25) -- node[above, font=\scriptsize] {中位數由小到大} (4.2,3.25);
\end{tikzpicture}
\diagcaption{第~\ref{sec:sel-mom-trace}~節第一層的 $25$ 個數。 每一欄是排好的一組，欄
的順序依中位數排（演算法本身不搬欄，這只是為了看清楚）。 $x = 38$（橘色）。 藍色的
$8$ 個數一定 $< x$，綠色的 $8$ 個數一定 $> x$；白色的數可能在 $x$ 的任何一邊。}
\label{fig:sel-mom-cols}
\end{figure}

\begin{lemma}
\label{lem:sel-mom-split}
設 $n \ge 6$，$n$ 個數互不相同。 \proc{Select} 第 14 行分割之後，$x$ 的左邊和右邊各有不超
過 $7n/10 + 2$ 個數。 所以第 19 行或第 21 行遞迴的範圍不超過 $7n/10 + 2$。
\end{lemma}

\begin{proof}
$g = \lceil n/5 \rceil$，$x$ 是 $g$ 個組中位數裡第 $\lceil g/2 \rceil$ 小的。 只有最後
一組可能不滿。

數 $\ge x$ 的數。 中位數 $\ge x$ 的組有 $g - \lceil g/2 \rceil + 1 = \lfloor g/2 \rfloor + 1$
組（包括 $x$ 自己那組）。 滿的一組裡，中位數和比它大的兩個數，共 $3$ 個都 $\ge x$。
不滿的那組有 $s$ 個數（$1 \le s \le 4$），$\ge$ 下中位數的有 $s - \lceil s/2 \rceil + 1$
個，至少 $1$ 個，比 $3$ 最多少 $2$。 所以 $\ge x$ 的數至少
$3(\lfloor g/2 \rfloor + 1) - 2 = 3\lfloor g/2 \rfloor + 1$ 個，$x$ 左邊（$< x$）的數
最多
\[
n - 3\left\lfloor \frac{g}{2} \right\rfloor - 1
\le n - 3 \cdot \frac{g - 1}{2} - 1
\le n - \frac{3}{2}\left(\frac{n}{5} - 1\right) - 1
= \frac{7n}{10} + \frac{1}{2}.
\]

數 $\le x$ 的數。 中位數 $\le x$ 的組有 $\lceil g/2 \rceil$ 組，滿的組各貢獻 $3$ 個，
不滿的組至少 $\lceil s/2 \rceil \ge 1$ 個，一樣最多少 $2$。 所以 $\le x$ 的至少
$3\lceil g/2 \rceil - 2$ 個，$x$ 右邊的數最多
\[
n - 3\left\lceil \frac{g}{2} \right\rceil + 2 \le n - \frac{3n}{10} + 2 = \frac{7n}{10} + 2,
\]
用到 $\lceil g/2 \rceil \ge g/2 \ge n/10$。 兩邊都不超過 $7n/10 + 2$。
\end{proof}

圖~\ref{fig:sel-mom-cols} 的例子：$n = 25$，引理保證每一邊不超過 $19$ 個；實際上藍色
和綠色各 $8$ 個，連同 $x$ 自己，每一邊都保證有 $9$ 個數不在那一邊，所以每一邊不超過
$25 - 9 = 16$ 個。

\subsection{最壞情況是線性的}

\begin{theorem}
\label{thm:sel-mom}
對 $n$ 個互不相同的數，\proc{Select} 的最壞執行時間是 $O(n)$。
\end{theorem}

\begin{proof}
令 $T(n)$ 是大小為 $n$ 的所有輸入、所有 $i$ 之中最長的執行時間。 第 13 行的遞迴大小
是 $\lceil n/5 \rceil$；分組、排各組、找 $x$ 的位置、分割，合計不超過 $an$（$a$ 是常
數）；由引理，第 19、21 行最多一次遞迴，大小不超過 $7n/10 + 2$。 所以 $n \ge 6$ 時
\[
T(n) \le T(\lceil n/5 \rceil) + T(m) + an, \qquad m \le \frac{7n}{10} + 2,
\]
沒有遞迴（$i = k$）時 $T(m)$ 那一項不出現。 兩個遞迴的大小都比 $n$ 小（引理證明裡的
兩個上界 $n - 3\lfloor g/2 \rfloor - 1$、$n - 3\lceil g/2 \rceil + 2$ 都小於 $n$），所以
演算法一定會停，每個 $T(n)$ 都是有限的。

用代換法證明存在常數 $c$，所有 $n \ge 1$ 都有 $T(n) \le cn$。 $n \ge 60$ 時對兩個遞
迴用歸納假設：
\[
T(n) \le c\left(\frac{n}{5} + 1\right) + c\left(\frac{7n}{10} + 2\right) + an
= \frac{9}{10}cn + 3c + an.
\]
要這個 $\le cn$，等於要 $\frac{c}{10}(n - 30) \ge an$。 $n \ge 60$ 時
$n - 30 \ge n/2$，左邊至少 $\frac{c}{20} n$，只要 $c \ge 20a$ 就成立。

$1 \le n < 60$ 的有限多個 $n$，$T(n)$ 都不超過某個常數 $B$，取 $c \ge B$ 就有
$T(n) \le B \le cn$。 所以 $c = \max(20a, B)$ 時，對所有 $n$ 都成立，$T(n) = O(n)$。
\end{proof}

證明能過，是因為兩個遞迴的大小加起來是 $\frac{n}{5} + \frac{7n}{10} = \frac{9n}{10}$：
用第~2~章第~\ref{sec:rec-tree}~節的遞迴樹來看，每往下一層總工作量最多是上一層的 $9/10$，各層加起來是
公比 $9/10$ 的等比級數。 門檻 $60$ 只是讓常數好算。

\subsection{為什麼是 5 個一組}

\textbf{3 個一組}：$\ge x$ 的數大約只有 $2 \cdot \frac{g}{2} = \frac{n}{3}$ 個，往下
遞迴的那一邊可能有 $\frac{2n}{3}$ 個，遞迴式變成
\[
T(n) \le T(n/3) + T(2n/3) + an.
\]
$\frac{1}{3} + \frac{2}{3} = 1$，每一層的工作量都不會變少。 代換法試 $T(n) \le cn$
會得到 $cn + an$，永遠多一個 $an$，常數再大也證不過；這個遞迴式的解是
$\Theta(n \lg n)$（第~2~章第~\ref{sec:rec-tree}~節「不平衡的樹」，把 $1/4 : 3/4$ 換成 $1/3 : 2/3$）。
所以用這個分析，$3$ 個一組只能保證 $O(n \lg n)$。

\textbf{7 個一組}：每組貢獻 $4$ 個，$\ge x$ 的數大約 $4 \cdot \frac{g}{2} = \frac{2n}{7}$
個，往下遞迴的一邊大約 $\frac{5n}{7}$，遞迴式是 $T(n/7) + T(5n/7) + an$，
$\frac{1}{7} + \frac{5}{7} = \frac{6}{7} < 1$，一樣是線性。 組越大，找 $x$ 的遞迴越
小，但每組排序的成本越高；$5$ 是能讓兩個比例加起來小於 $1$ 的最小奇數。

\subsection{手算追蹤：25 個數}
\label{sec:sel-mom-trace}

$A = \langle 48, 89, 80, 91, 44, 13, 35, 30, 85, 66, 98, 33, 38, 88, 15, 70, 87, 31,
16, 27, 24, 50, 86, 71, 34 \rangle$，找中位數，也就是第 $13$ 小。 下圖每一列是一個階
段之後的陣列，粗線是組的分界，粉紅色是各組的中位數，橘色是這一層的 $x$（或最後的答
案），灰色是已經丟掉的部分。

\begin{center}
\begin{tikzpicture}
\node[anchor=east, font=\small] at (0.25,0.00) {原始};
\node[sel-sc] at (0.55,0.00) {48};
\node[sel-sc] at (1.10,0.00) {89};
\node[sel-sc] at (1.65,0.00) {80};
\node[sel-sc] at (2.20,0.00) {91};
\node[sel-sc] at (2.75,0.00) {44};
\node[sel-sc] at (3.30,0.00) {13};
\node[sel-sc] at (3.85,0.00) {35};
\node[sel-sc] at (4.40,0.00) {30};
\node[sel-sc] at (4.95,0.00) {85};
\node[sel-sc] at (5.50,0.00) {66};
\node[sel-sc] at (6.05,0.00) {98};
\node[sel-sc] at (6.60,0.00) {33};
\node[sel-sc] at (7.15,0.00) {38};
\node[sel-sc] at (7.70,0.00) {88};
\node[sel-sc] at (8.25,0.00) {15};
\node[sel-sc] at (8.80,0.00) {70};
\node[sel-sc] at (9.35,0.00) {87};
\node[sel-sc] at (9.90,0.00) {31};
\node[sel-sc] at (10.45,0.00) {16};
\node[sel-sc] at (11.00,0.00) {27};
\node[sel-sc] at (11.55,0.00) {24};
\node[sel-sc] at (12.10,0.00) {50};
\node[sel-sc] at (12.65,0.00) {86};
\node[sel-sc] at (13.20,0.00) {71};
\node[sel-sc] at (13.75,0.00) {34};
\draw[very thick, black!80] (3.025,-0.36) -- (3.025,0.36);
\draw[very thick, black!80] (5.775,-0.36) -- (5.775,0.36);
\draw[very thick, black!80] (8.525,-0.36) -- (8.525,0.36);
\draw[very thick, black!80] (11.275,-0.36) -- (11.275,0.36);
\node[anchor=east, font=\small] at (0.25,-0.95) {各組排好};
\node[sel-sc] at (0.55,-0.95) {44};
\node[sel-sc] at (1.10,-0.95) {48};
\node[sel-scmed] at (1.65,-0.95) {80};
\node[sel-sc] at (2.20,-0.95) {89};
\node[sel-sc] at (2.75,-0.95) {91};
\node[sel-sc] at (3.30,-0.95) {13};
\node[sel-sc] at (3.85,-0.95) {30};
\node[sel-scmed] at (4.40,-0.95) {35};
\node[sel-sc] at (4.95,-0.95) {66};
\node[sel-sc] at (5.50,-0.95) {85};
\node[sel-sc] at (6.05,-0.95) {15};
\node[sel-sc] at (6.60,-0.95) {33};
\node[sel-scpiv] at (7.15,-0.95) {38};
\node[sel-sc] at (7.70,-0.95) {88};
\node[sel-sc] at (8.25,-0.95) {98};
\node[sel-sc] at (8.80,-0.95) {16};
\node[sel-sc] at (9.35,-0.95) {27};
\node[sel-scmed] at (9.90,-0.95) {31};
\node[sel-sc] at (10.45,-0.95) {70};
\node[sel-sc] at (11.00,-0.95) {87};
\node[sel-sc] at (11.55,-0.95) {24};
\node[sel-sc] at (12.10,-0.95) {34};
\node[sel-scmed] at (12.65,-0.95) {50};
\node[sel-sc] at (13.20,-0.95) {71};
\node[sel-sc] at (13.75,-0.95) {86};
\draw[very thick, black!80] (3.025,-1.31) -- (3.025,-0.59);
\draw[very thick, black!80] (5.775,-1.31) -- (5.775,-0.59);
\draw[very thick, black!80] (8.525,-1.31) -- (8.525,-0.59);
\draw[very thick, black!80] (11.275,-1.31) -- (11.275,-0.59);
\node[anchor=east, font=\small] at (0.25,-1.90) {對 38 分割};
\node[sel-sc] at (0.55,-1.90) {13};
\node[sel-sc] at (1.10,-1.90) {30};
\node[sel-sc] at (1.65,-1.90) {35};
\node[sel-sc] at (2.20,-1.90) {15};
\node[sel-sc] at (2.75,-1.90) {33};
\node[sel-sc] at (3.30,-1.90) {16};
\node[sel-sc] at (3.85,-1.90) {27};
\node[sel-sc] at (4.40,-1.90) {31};
\node[sel-sc] at (4.95,-1.90) {24};
\node[sel-sc] at (5.50,-1.90) {34};
\node[sel-scpiv] at (6.05,-1.90) {38};
\node[sel-sc] at (6.60,-1.90) {91};
\node[sel-sc] at (7.15,-1.90) {86};
\node[sel-sc] at (7.70,-1.90) {88};
\node[sel-sc] at (8.25,-1.90) {98};
\node[sel-sc] at (8.80,-1.90) {44};
\node[sel-sc] at (9.35,-1.90) {48};
\node[sel-sc] at (9.90,-1.90) {80};
\node[sel-sc] at (10.45,-1.90) {70};
\node[sel-sc] at (11.00,-1.90) {87};
\node[sel-sc] at (11.55,-1.90) {66};
\node[sel-sc] at (12.10,-1.90) {85};
\node[sel-sc] at (12.65,-1.90) {50};
\node[sel-sc] at (13.20,-1.90) {71};
\node[sel-sc] at (13.75,-1.90) {89};
\node[anchor=east, font=\small] at (0.25,-2.85) {各組排好};
\node[sel-scout] at (0.55,-2.85) {13};
\node[sel-scout] at (1.10,-2.85) {30};
\node[sel-scout] at (1.65,-2.85) {35};
\node[sel-scout] at (2.20,-2.85) {15};
\node[sel-scout] at (2.75,-2.85) {33};
\node[sel-scout] at (3.30,-2.85) {16};
\node[sel-scout] at (3.85,-2.85) {27};
\node[sel-scout] at (4.40,-2.85) {31};
\node[sel-scout] at (4.95,-2.85) {24};
\node[sel-scout] at (5.50,-2.85) {34};
\node[sel-scout] at (6.05,-2.85) {38};
\node[sel-sc] at (6.60,-2.85) {44};
\node[sel-sc] at (7.15,-2.85) {86};
\node[sel-scmed] at (7.70,-2.85) {88};
\node[sel-sc] at (8.25,-2.85) {91};
\node[sel-sc] at (8.80,-2.85) {98};
\node[sel-sc] at (9.35,-2.85) {48};
\node[sel-sc] at (9.90,-2.85) {66};
\node[sel-scmed] at (10.45,-2.85) {70};
\node[sel-sc] at (11.00,-2.85) {80};
\node[sel-sc] at (11.55,-2.85) {87};
\node[sel-sc] at (12.10,-2.85) {50};
\node[sel-scpiv] at (12.65,-2.85) {71};
\node[sel-sc] at (13.20,-2.85) {85};
\node[sel-sc] at (13.75,-2.85) {89};
\draw[very thick, black!80] (9.075,-3.21) -- (9.075,-2.49);
\draw[very thick, black!80] (11.825,-3.21) -- (11.825,-2.49);
\node[anchor=east, font=\small] at (0.25,-3.80) {對 71 分割};
\node[sel-scout] at (0.55,-3.80) {13};
\node[sel-scout] at (1.10,-3.80) {30};
\node[sel-scout] at (1.65,-3.80) {35};
\node[sel-scout] at (2.20,-3.80) {15};
\node[sel-scout] at (2.75,-3.80) {33};
\node[sel-scout] at (3.30,-3.80) {16};
\node[sel-scout] at (3.85,-3.80) {27};
\node[sel-scout] at (4.40,-3.80) {31};
\node[sel-scout] at (4.95,-3.80) {24};
\node[sel-scout] at (5.50,-3.80) {34};
\node[sel-scout] at (6.05,-3.80) {38};
\node[sel-sc] at (6.60,-3.80) {44};
\node[sel-sc] at (7.15,-3.80) {48};
\node[sel-sc] at (7.70,-3.80) {66};
\node[sel-sc] at (8.25,-3.80) {70};
\node[sel-sc] at (8.80,-3.80) {50};
\node[sel-scpiv] at (9.35,-3.80) {71};
\node[sel-sc] at (9.90,-3.80) {88};
\node[sel-sc] at (10.45,-3.80) {91};
\node[sel-sc] at (11.00,-3.80) {80};
\node[sel-sc] at (11.55,-3.80) {87};
\node[sel-sc] at (12.10,-3.80) {98};
\node[sel-sc] at (12.65,-3.80) {89};
\node[sel-sc] at (13.20,-3.80) {85};
\node[sel-sc] at (13.75,-3.80) {86};
\node[anchor=east, font=\small] at (0.25,-4.75) {排好 $A[12..16]$};
\node[sel-scout] at (0.55,-4.75) {13};
\node[sel-scout] at (1.10,-4.75) {30};
\node[sel-scout] at (1.65,-4.75) {35};
\node[sel-scout] at (2.20,-4.75) {15};
\node[sel-scout] at (2.75,-4.75) {33};
\node[sel-scout] at (3.30,-4.75) {16};
\node[sel-scout] at (3.85,-4.75) {27};
\node[sel-scout] at (4.40,-4.75) {31};
\node[sel-scout] at (4.95,-4.75) {24};
\node[sel-scout] at (5.50,-4.75) {34};
\node[sel-scout] at (6.05,-4.75) {38};
\node[sel-sc] at (6.60,-4.75) {44};
\node[sel-scpiv] at (7.15,-4.75) {48};
\node[sel-sc] at (7.70,-4.75) {50};
\node[sel-sc] at (8.25,-4.75) {66};
\node[sel-sc] at (8.80,-4.75) {70};
\node[sel-scout] at (9.35,-4.75) {71};
\node[sel-scout] at (9.90,-4.75) {88};
\node[sel-scout] at (10.45,-4.75) {91};
\node[sel-scout] at (11.00,-4.75) {80};
\node[sel-scout] at (11.55,-4.75) {87};
\node[sel-scout] at (12.10,-4.75) {98};
\node[sel-scout] at (12.65,-4.75) {89};
\node[sel-scout] at (13.20,-4.75) {85};
\node[sel-scout] at (13.75,-4.75) {86};
\foreach \k in {1,...,25} { \node[idxlbl] at (\k*0.55, -5.35) {\k}; }
\end{tikzpicture}
\diagcaption{\proc{Select} 找 $25$ 個數的第 $13$ 小。 第一層 $x = 38$，往右；第二層
$A[12..25]$ 的 $x = 71$，往左；第三層只剩 $5$ 個數，直接排序。}
\label{fig:sel-mom-trace}
\end{center}

第一層分成 $5$ 組，組中位數 $M = \langle 80, 35, 38, 31, 50 \rangle$；對 $M$ 找第
$\lceil 5/2 \rceil = 3$ 小，$M$ 只有 $5$ 個數，直接排序得到 $x = 38$。 用 $38$ 分割後
$q = 11$，$k = 11 < 13$，往右邊 $A[12..25]$ 找第 $13 - 11 = 2$ 小。 第二層的 $14$ 個數分
成 $A[12..16]$、$A[17..21]$、$A[22..25]$ 三組，最後一組 $50, 71, 85, 89$ 取下中位數
$71$；$M = \langle 88, 70, 71 \rangle$ 的第 $2$ 小是 $x = 71$。 用 $71$ 分割後
$q = 17$，$k = 6 > 2$，往左邊 $A[12..16]$ 找第 $2$ 小。 第三層只剩 $5$ 個數，直接排成
$44, 48, 50, 66, 70$，答案是 $A[13] = 48$。

驗算：$25$ 個數由小到大的前 $13$ 個是 $13, 15, 16, 24, 27, 30, 31, 33, 34, 35, 38, 44,
48$，第 $13$ 個是 $48$。 第一層丟掉了 $11$ 個，第二層從 $14$ 個留下 $5$ 個；第一層往
右的那一邊有 $14$ 個，不超過上一小節算的 $16$。

\subsection{C 程式}

\texttt{select\_mom} 照虛擬碼翻譯，$M$ 用 \texttt{malloc} 配置，找完 $x$ 就
\texttt{free}；\texttt{partition\_around} 把找位置、交換和 Lomuto 分割寫在一起。 程式
印出第~\ref{sec:sel-mom-trace}~節那 $25$ 個數每一層的追蹤，\texttt{d=} 後面是遞迴到第幾
層的 $M$，\texttt{d=0} 是原陣列。

\begin{example}[label=ex:sel-mom]{c}{(mom\_select.c) 中位數的中位數}
#include <stdio.h>
#include <stdlib.h>

void swap(int *x, int *y)
{
    int t = *x;
    *x = *y;
    *y = t;
}

void insertion_sort(int *a, int l, int h)       // 排好 a[l..h]
{
    for (int i = l + 1; i <= h; i++) {
        int key = a[i], j = i - 1;
        while (j >= l && a[j] > key) {
            a[j + 1] = a[j];
            j--;
        }
        a[j + 1] = key;
    }
}

// 把 x 換到 a[r]，再做 Lomuto 分割，回傳 x 最後的位置。
int partition_around(int *a, int p, int r, int x)
{
    int j = p;
    while (a[j] != x)
        j++;
    swap(&a[j], &a[r]);
    int i = p - 1;
    for (j = p; j < r; j++)
        if (a[j] <= x)
            swap(&a[++i], &a[j]);
    swap(&a[i + 1], &a[r]);
    return i + 1;
}

// 回傳 a[p..r] 裡第 i 小的數，並印出每一層的追蹤。
// d 是遞迴到「中位數陣列」的層數，原陣列是 0。
int select_mom(int *a, int p, int r, int i, int d)
{
    int n = r - p + 1;
    if (n <= 5) {
        insertion_sort(a, p, r);
        printf("d=%d n=%d i=%d base -> %d\n",
               d, n, i, a[p + i - 1]);
        return a[p + i - 1];
    }
    int g = (n + 4) / 5;
    int *m = malloc(g * sizeof *m);
    if (m == NULL)
        exit(1);
    for (int j = 0; j < g; j++) {
        int l = p + 5 * j, h = l + 4 < r ? l + 4 : r;
        insertion_sort(a, l, h);
        m[j] = a[(l + h) / 2];      // 這一組的（下）中位數
    }
    printf("d=%d n=%d i=%d medians", d, n, i);
    for (int j = 0; j < g; j++)
        printf(" %d", m[j]);
    printf("\n");
    int x = select_mom(m, 0, g - 1, (g + 1) / 2, d + 1);
    free(m);
    int k = partition_around(a, p, r, x) - p + 1;
    printf("d=%d x=%d k=%d %s\n", d, x, k,
           i == k ? "found" : (i < k ? "go left" : "go right"));
    if (i == k)
        return x;
    else if (i < k)
        return select_mom(a, p, p + k - 2, i, d);
    else
        return select_mom(a, p + k, r, i - k, d);
}

int main(void)
{
    int a[] = {48, 89, 80, 91, 44, 13, 35, 30, 85, 66, 98, 33, 38,
               88, 15, 70, 87, 31, 16, 27, 24, 50, 86, 71, 34};
    printf("answer %d\n", select_mom(a, 0, 24, 13, 0));
    return 0;
}
\end{example}
\iosample{}{d=0 n=25 i=13 medians 80 35 38 31 50 \\ d=1 n=5 i=3 base -> 38 \\ d=0 x=38 k=11 go right \\ d=0 n=14 i=2 medians 88 70 71 \\ d=1 n=3 i=2 base -> 71 \\ d=0 x=71 k=6 go left \\ d=0 n=5 i=2 base -> 48 \\ answer 48}

追蹤的七行和圖~\ref{fig:sel-mom-trace} 一一對應。 這個演算法每一層都要排好每一組、再遞迴找
$x$，比較次數的常數比範例~\ref{ex:sel-rselect} 的大約 $3.3n$ 大得多。 所以實務上多半用隨機化
選擇，或者先用隨機化選擇、發現遞迴太深時才換成中位數的中位數（Musser 在 1997 年提出的
introselect）。 中位數的中位數的價值在於它證明了
「最壞情況線性」做得到，而它的分析方法（兩個遞迴的比例加起來小於 $1$）在其他問題也
用得上。

\section{本章常見錯誤}
\label{sec:sel-pitfalls}

\begin{center}
\begin{tabular}{p{0.25\linewidth}p{0.50\linewidth}p{0.15\linewidth}}
\toprule
錯誤 & 例子 & 見 \\
\midrule
往右走時名次沒有減 & 遞迴寫成 \proc{Randomized-Select}$(A, q+1, r, i)$，答案錯，或
範圍外存取 & 第~\ref{sec:sel-random}~節 \\
名次和索引差一 & 題目的 $k$ 從 $1$ 開始，C 陣列從 $0$ 開始，直接拿 $k$ 當索引 &
第~\ref{sec:sel-trace}~節 \\
兩邊都遞迴 & 寫成快速排序，時間變成 $\Theta(n \lg n)$ & 第~\ref{sec:sel-random}~節 \\
固定取最後一個當基準 & 已經排好的輸入變成 $\Theta(n^2)$ & 第~\ref{sec:sel-random}~節 \\
忽略重複值 & 全部相等時 Lomuto 分割每次只少一個數，隨機也沒用 &
第~\ref{sec:sel-random}~節 \\
以為 $x$ 是中位數 & 中位數的中位數只保證兩邊各有至少約 $3n/10$ 個數 &
第~\ref{sec:sel-mom}~節 \\
3 個一組也是線性 & 遞迴式 $T(n/3) + T(2n/3) + an$ 的解是 $\Theta(n \lg n)$ &
第~\ref{sec:sel-mom}~節 \\
決賽輸家當第二大 & 第二大可能第一輪就遇到冠軍 & 第~\ref{sec:sel-minmax}~節 \\
最大最小的次數算錯 & 寫成 $3n/2$，沒有分奇偶，或忘了偶數時開頭那一次 &
第~\ref{sec:sel-minmax}~節 \\
\bottomrule
\end{tabular}
\end{center}

\section{觀念題}
\label{sec:sel-exercises}

題型和研究所考試、期中期末考的筆試相近，有選擇、計算、追蹤和證明。 解答在本章最後。

\begin{enumerate}
\item \label{sel-q-tf}（選擇，可能不只一個答案）下列哪些敘述正確？
\begin{enumerate}
\item \proc{Randomized-Select} 的最壞執行時間是 $\Theta(n^2)$。
\item \proc{Select} 第 13 行找到的 $x$ 一定是 $A[p..r]$ 的中位數。
\item 把 \proc{Select} 改成 $3$ 個一組，第~\ref{sec:sel-mom}~節的分析仍然能證明最壞
$O(n)$。
\item 任何只用比較找最大值的演算法，至少要比 $n - 1$ 次。
\item $n = 10$ 時，同時找最大和最小，最壞情況至少要比 $13$ 次。
\end{enumerate}

\item \label{sel-q-tour}（追蹤）
\begin{enumerate}
\item 用錦標賽找 $\langle 93, 41, 67, 28, 85, 52, 76, 39, 60 \rangle$ 的第二大。 相鄰
兩個配對，人數是奇數時最後一個輪空。 畫出比賽樹，寫出直接輸給冠軍的數，以及總比較
次數。
\item 同樣的數改成 $\langle 41, 67, 28, 85, 52, 76, 39, 60, 93 \rangle$，再做一次。 比
較次數為什麼和定理~\ref{thm:sel-second} 不一樣？ 這和下界矛盾嗎？
\end{enumerate}

\item \label{sel-q-trace}（追蹤）用基準固定取最後一個的選擇（把
\proc{Randomized-Partition} 換成 \proc{Lomuto-Partition}），在
$\langle 35, 81, 22, 64, 17, 90, 48, 73, 56 \rangle$ 裡找第 $3$ 小。 每次分割之後寫出
陣列、$p$、$r$、$q$、$k$ 和新的 $i$，並算出總比較次數。

\item \label{sel-q-worst}（計算）
\begin{enumerate}
\item 基準固定取最後一個，在已經由小到大排好的 $n$ 個數裡找最小值，比較幾次？
\item 改用 \proc{Randomized-Select} 找最小值，比較次數和 (a) 一樣多的機率是多少？
\end{enumerate}

\item \label{sel-q-n3}（計算）$n = 3$、$i = 2$，三個數互不相同。 列出
\proc{Randomized-Select} 所有可能的基準選擇，算出期望比較次數，再用
定理~\ref{thm:sel-rselect-cmp} 的公式驗算。

\item \label{sel-q-min}（證明）用定理~\ref{thm:sel-rselect-cmp} 的公式證明：找最小值
（$i = 1$）的期望比較次數是 $2n - 2H_n$，$H_n = 1 + \frac{1}{2} + \cdots + \frac{1}{n}$。
和範例~\ref{ex:sel-rselect} 印出的 $1987.0$ 對照。

\item \label{sel-q-correct}（證明）假設 $A[p..r]$ 的數互不相同，$1 \le i \le r - p + 1$。
對 $n = r - p + 1$ 用強歸納法，證明 \proc{Randomized-Select}$(A, p, r, i)$ 回傳
$A[p..r]$ 裡第 $i$ 小的數。

\item \label{sel-q-dup}（證明）$A[1..n]$ 的 $n$ 個數全部相等，$i = \lfloor (n+1)/2 \rfloor$。
\begin{enumerate}
\item 證明不管 \proc{Random} 傳回什麼，\proc{Randomized-Select} 都要比較
$\Theta(n^2)$ 次。
\item 改寫分割，讓它傳回兩個位置 $q_1 \le q_2$，使 $A[p..q_1 - 1] < x$、
$A[q_1..q_2] = x$、$A[q_2 + 1..r] > x$，再寫出對應的選擇程序（虛擬碼）。 說明為什麼
全部相等時只要比 $O(n)$ 次。
\end{enumerate}

\item \label{sel-q-mom15}（追蹤）用 \proc{Select} 找
$\langle 42, 17, 93, 58, 26, 71, 35, 84, 10, 63, 49, 77, 22, 96, 31 \rangle$ 的中位數。
寫出各組排好的結果、$M$、$x$、分割後的陣列和 $k$。 圖~\ref{fig:sel-mom-cols} 那種
「保證 $\ge x$」的數有幾個？

\item \label{sel-q-quarter}（證明）把 \proc{Select} 第 13 行改成取 $M$ 的第
$\lceil g/4 \rceil$ 小，其他不變。 仿照引理~\ref{lem:sel-mom-split} 估計遞迴範圍的上
界，寫出遞迴式，說明為什麼定理~\ref{thm:sel-mom} 的代換法證不過。

\end{enumerate}

\section{歷屆試題}
\label{sec:sel-pastexam}

以下是研究所入學考試的原題，題目照原卷（語言、記號、配分照錄）。 詳解在本章最後。

\begin{pastexam}[label={px:sel-1}]{中正資工所 113 學年度 第 20 題}
(2\%) Determine if the following statement is true or false: "Find the median of $n$ elements takes $\Omega(nlogn)$ time for any algorithms working in the comparison-based model".
\end{pastexam}

\begin{pastexam}[label={px:sel-2}]{臺大資工所 113 學年度 第 1 題}
（複選題，每個選項單獨計分）

(5 points) Consider the following variation of the Selection Problem: Given a sequence $S=\langle a_1,a_2,\ldots,a_n\rangle$ of $n$ distinct integers and two positive integers $x$ and $y$, where each number in $S$ falls between 1 and $y$, the goal is to report the smallest $x$ numbers in $S$ in ascending order. This scenario assumes the conventional single-processor, random access machine (RAM) model, with a machine word size of at least $\lceil\log y\rceil$ bits. To tackle this problem, five algorithms have been proposed. Please select the correct description(s) below.
\begin{enumerate}[label=\Alph*.]
\item Algorithm A: The algorithm initializes an empty max-heap and sequentially inserts numbers from $S$. During insertion, if the size of the max-heap exceeds $x$, it removes the maximum element from the heap. Finally, the remaining numbers in the heap are sorted by heapsort and reported. If $x\le n/\log n$, Algorithm A runs in $O(n)$ time in the worst case.
\item Algorithm B: The algorithm initially constructs a min-heap and proceeds to extract the minimum element $x$ times. If $x\le n/\log n$, Algorithm B runs in $O(n)$ time in the worst case.
\item Algorithm C: The algorithm begins by performing a base-$n$ radix sort, followed by reporting the smallest $x$ values. If $y\le n^2$, Algorithm C runs in $O(n)$ time in the worst case.
\item Algorithm D: The algorithm performs a counting sort by allocating an array $C$ of length $y$, traversing through $S$ and incrementing $C[a_i]$ for each $a_i$. It finally reports the smallest $x$ values by scanning through the counter array. If $y\le n^2$, Algorithm D runs in $O(n)$ time in the worst case.
\item Algorithm E: The algorithm first applies quickselect, i.e., Randomized-Select, to determine the $x$-th smallest number in $S$. Subsequently, it traverses $S$ for extracting all numbers smaller than or equal to the $x$-th smallest number. Finally, these $x$ numbers are sorted by quicksort and reported. The time complexity of Algorithm E is $O(n+x\log x)$ in expectation.
\end{enumerate}
\end{pastexam}

\begin{pastexam}[label={px:sel-3}]{中央資工所 115 學年度 第 10 題}
（多選題，每一個選項單獨計分，每題 5 分）

Consider the prune-and-search algorithm using the median of medians for finding the $k$-th smallest element in an unsorted array $A$ of $n$ distinct elements, where $1\le k\le n$, as described below:

\begin{quote}
Algorithm Select($A$, $k$):\\
Input: An array $A$ of $n$ elements and an integer $k$, $1\le k\le |A|=n$\\
Output: The $k$-th smallest element in $A$\\
Step 1: If $|A|\le 5$, then sort array $A$ (with any method, e.g., insertion sort) and return the element with rank $k$ (i.e., the $k$-th element in the sorted array).\\
Step 2: Divide the array $A$ into groups of five elements each (ignore the last group if it has fewer than five elements).\\
Step 3: Find the median of each group.\\
Step 4: Recursively compute the median of the found medians, denoted by $m$.\\
Step 5: Partition the original array into the following three sets:\\
$L=\{x\in A\mid x<m\}$\\
$E=\{m\}$\\
$R=\{x\in A\mid x>m\}$\\
Step 6: If $k=|L|+1$ then return $m$\\
\hspace*{3.2em}Else if $k\le|L|$ then return Select$(L,k)$\\
\hspace*{3.2em}Else return Select$(R,k-|L|-1)$.
\end{quote}

Select one or more correct statements.
\begin{enumerate}[label=\Alph*.]
\item After grouping elements into blocks of five and choosing the median of medians as the pivot denoted $m$, at least $3n/10$ elements are guaranteed to be less than or equal to the pivot.
\item The recurrence $\mathrm{T}(n)=\mathrm{T}(n/5)+\mathrm{T}(7n/10)+\mathrm{O}(n)$ implies that the algorithm runs in $\mathrm{O}(n\log n)$ time.
\item Using groups of size 3 instead of 5 still guarantees linear worst-case running time.
\item The correctness of the pruning step relies on the fact that the pivot is an approximate median of the entire array.
\item The median-finding problem is a special case of the $k$-th smallest selection problem, where $k=\lceil n/2\rceil$, and both can be solved in worst-case linear time.
\end{enumerate}
\end{pastexam}

\begin{pastexam}[label={px:sel-4}]{成大資訊聯招 109 學年度 演算法第 3 題}
Show how the time complexity of quicksort can be made to run in $O(n\lg n)$ time. \quad (10\%)
\end{pastexam}

\section{實作練習}
\label{sec:sel-practice}

下面是 Aki-Judge 上和本章相關的題目。 兩題都只交函式（第二題要交標頭檔和 \texttt{.c}
檔），回傳第 $k$ 小的值（$k$ 從 $1$ 開始），不寫 \texttt{main}。

\begin{center}
\begin{tabular}{llc}
\toprule
小節 & 題目 & 難度 \\
\midrule
隨機化選擇 & \akiref{k-th-smallest-element}{第 K 小的元素} & 基本 \\
中位數的中位數 & \akiref{worst-case-linear-k-th-smallest-function}{最壞情況線性時間的第 k 小函式} & 進階 \\
\bottomrule
\end{tabular}
\end{center}

第一題的 $n$ 到 $10^6$。 固定取最後一個當基準，遇到已經排好的輸入會變成
$\Theta(n^2)$；大量重複的值也會讓 Lomuto 分割變慢（第~\ref{sec:sel-random}~節最後一
小節）。 寫的時候用隨機基準，再加上三路分割或第~14~章的 Hoare 分割，程式對任何輸入
都不會退化。 第二題請照
第~\ref{sec:sel-mom}~節寫中位數的中位數；注意題目允許重複值，分割時等於 $x$ 的數不能
都擠到同一邊，否則引理~\ref{lem:sel-mom-split} 的保證就不成立。 交之前先用
範例~\ref{ex:sel-mom} 的資料比對每一層的結果。

\section{讀完你應該會}
\label{sec:sel-checklist}

\begin{itemize}
\item 能寫出成對處理的最大最小演算法，證明它剛好比 $\lceil 3n/2 \rceil - 2$ 次，並說
出對手論證為什麼不能更少。
\item 能畫出錦標賽的比賽樹，找出第二大，算出 $n + \lceil \lg n \rceil - 2$ 次的由
來。
\item 能在給定的分割方法下，逐次寫出選擇演算法每一次分割後的陣列、$q$、$k$ 和新的
名次。
\item 能用 $U(n)$ 的遞迴式和代換法證明隨機化選擇的期望時間是線性，並用指示變數寫出
期望比較次數。
\item 能說明中位數的中位數為什麼保證兩邊都不超過約 $7n/10$，寫出遞迴式並用代換法證
明線性；能說明 $3$ 個一組為什麼不行。
\end{itemize}

\textbf{延伸閱讀}（書末會整理完整的參考資料）

\begin{itemize}
\item Thomas H. Cormen 等，\textit{Introduction to Algorithms}，第 4 版，第 9 章：本章
三節的對應內容。 它證明隨機化選擇的方法和本章不同（數「有幫助的分割」），中位數的
中位數也用了另一種分組方式，可以對照著看。
\item Donald E. Knuth，\textit{The Art of Computer Programming, Vol.~3}，第 2 版，
5.3.3 節：用最少的比較次數找第 $i$ 小，包括第二大的下界的完整證明。
\item M. Blum、R. W. Floyd、V. Pratt、R. L. Rivest、R. E. Tarjan，``Time Bounds for
Selection''，\textit{Journal of Computer and System Sciences} 7(4)，1973：中位數的中
位數的原始論文，也給了選擇問題比較次數的下界。
\item K. Chen、A. Dumitrescu，``Select with Groups of 3 or 4''，WADS 2015：改變挑基準
的方式之後，$3$ 個或 $4$ 個一組也能做到最壞線性時間。
\end{itemize}

\addcontentsline{toc}{section}{觀念題解答}
\section*{觀念題解答}

\begin{enumerate}
\item (a)、(d)、(e) 正確。 (a)：每次都挑到最壞的基準時是 $n(n-1)/2$ 次比較。 (b)：$x$
只保證兩邊各至少約 $3n/10$ 個數，第~\ref{sec:sel-mom-trace}~節的 $x = 38$ 是第 $11$
小，不是第 $13$ 小。 (c)：遞迴式變成 $T(n/3) + T(2n/3) + an$，比例加起來是 $1$，只能
證出 $O(n \lg n)$。 (d)：除了最大值，每個數都要輸一場。 (e)：$\lceil 30/2 \rceil - 2 = 13$。

\item
\begin{enumerate}
\item 第一輪：$93 > 41$、$67 > 28$、$85 > 52$、$76 > 39$，$60$ 輪空。 第二輪：
$93 > 67$、$85 > 76$，$60$ 輪空。 第三輪：$93 > 85$，$60$ 輪空。 第四輪：$93 > 60$。
找冠軍比了 $4 + 2 + 1 + 1 = 8$ 次。 直接輸給 $93$ 的是 $41, 67, 85, 60$，再比 $3$ 次
得到第二大 $85$。 共 $11 = 9 + \lceil \lg 9 \rceil - 2$ 次。
\item 第一輪：$67$、$85$、$76$、$60$ 晉級，$93$ 輪空；第二輪：$85$、$76$ 晉級，$93$
輪空；第三輪：$85$ 晉級，$93$ 輪空；第四輪：$93 > 85$。 找冠軍一樣比 $8$ 次，但冠軍只
比了一場，直接輸給它的只有 $85$，不用再比，共 $8$ 次。 定理~\ref{thm:sel-second} 說
的是最壞情況的次數；這個輸入剛好讓冠軍一路輪空，比較少。 下界也是對最壞情況說的：
任何演算法都\textbf{存在}某個輸入要比 $11$ 次，不是每個輸入都要，所以不矛盾。
\end{enumerate}

\item
\begin{center}
\begin{tabular}{cccccl}
\toprule
$p..r$ & $i$ & 基準 & 分割後的陣列 & $q$, $k$ & 下一步 \\
\midrule
$1..9$ & 3 & 56 & 35 22 17 48 56 90 64 73 81 & 5, 5 & 往左 \\
$1..4$ & 3 & 48 & 35 22 17 48 56 90 64 73 81 & 4, 4 & 往左 \\
$1..3$ & 3 & 17 & 17 22 35 48 56 90 64 73 81 & 1, 1 & 往右，$i \gets 2$ \\
$2..3$ & 2 & 35 & 17 22 35 48 56 90 64 73 81 & 3, 2 & 回傳 $35$ \\
\bottomrule
\end{tabular}
\end{center}
比較次數 $8 + 3 + 2 + 1 = 14$。

\item (a) 每次的基準都是範圍裡最大的數，$k = n$，往左邊少一個數的範圍找。 比較次數
$(n-1) + (n-2) + \cdots + 1 = n(n-1)/2$。 (b) 要比這麼多次，範圍是 $m$ 個數時基準必
須是其中最大的（任何其他基準都會一次丟掉不只一個數），機率 $1/m$。 $m = n, n-1,
\ldots, 2$ 每一步獨立，機率是 $\frac{1}{n} \cdot \frac{1}{n-1} \cdots \frac{1}{2} =
\frac{1}{n!}$。

\item 第一次分割一定比 $2$ 次。 基準是第 $2$ 小（機率 $1/3$）：直接找到，共 $2$ 次。
基準是第 $1$ 小（機率 $1/3$）：往右邊兩個數找第 $1$ 小，再分割一次比 $1$ 次，不管基
準是哪個，下一步不是找到就是只剩一個數，共 $3$ 次。 基準是第 $3$ 小：對稱，$3$ 次。
期望值 $(2 + 3 + 3)/3 = 8/3$。 公式：$(a, b) = (1, 2)$ 的分母是 $2 - 1 + 1 = 2$，
$(1, 3)$ 是 $3$，$(2, 3)$ 是 $2$，和是 $1 + \frac{2}{3} + 1 = \frac{8}{3}$，一樣。

\item $i = 1$ 時 $\min(a, 1) = 1$、$\max(b, 1) = b$，每一對的機率是 $2/b$。 固定 $b$，
$a$ 有 $b - 1$ 種：
\[
E_n(1) = \sum_{b=2}^{n} \frac{2(b-1)}{b} = 2(n - 1) - 2\sum_{b=2}^{n} \frac{1}{b}
= 2(n - 1) - 2(H_n - 1) = 2n - 2H_n.
\]
$n = 1001$：$H_{1001} \approx \ln 1001 + 0.5772 \approx 7.486$，$2002 - 14.97 \approx
1987.0$，和程式印的公式值一樣。

\item $n = 1$：$p = r$，$i$ 只能是 $1$，回傳唯一的數，正確。 $n \ge 2$：假設比 $n$ 小
的範圍都正確。 \proc{Lomuto-Partition} 結束後 $A[p..q-1]$ 都 $< A[q]$、$A[q+1..r]$ 都
$> A[q]$（數互不相同），而且 $A[p..r]$ 仍是原來那些數，所以 $A[q]$ 是第 $k = q - p + 1$
小。 $i = k$ 時回傳 $A[q]$，正確。 $i < k$ 時，第 $i$ 小的數比 $A[q]$ 小，在
$A[p..q-1]$ 裡，而且這一段的數正是全部裡最小的 $k - 1$ 個，所以它在這一段裡也是第 $i$
小；範圍有 $k - 1 \ge i$ 個數、小於 $n$，由歸納假設遞迴回傳正確。 $i > k$ 時，答案在
$A[q+1..r]$，比它小的數裡有 $k$ 個在左邊（含基準），所以它在右邊是第 $i - k$ 小，範圍
有 $n - k \ge i - k \ge 1$ 個數，同樣由歸納假設正確。

\item (a) 所有數相等，隨機交換不改變陣列。 Lomuto 分割裡每個 $A[j] \le x$ 都成立，結
束時 $q = r$、$k = n$。 $i < k$ 時往左邊 $n - 1$ 個數遞迴，名次不變。 範圍從 $n$ 一路
減一，直到範圍大小等於 $i$ 時 $k = i$ 才停，比較次數
$\sum_{m=i}^{n} (m - 1) \ge (n - i + 1)(i - 1)$。 $i = \lfloor (n+1)/2 \rfloor$ 時這是
$\Theta(n^2)$。

(b) 三路分割的一種寫法：用 $\mathit{lt}$、$j$、$\mathit{gt}$ 三個指標，維持
$A[p..\mathit{lt}-1] < x$、$A[\mathit{lt}..j-1] = x$、$A[\mathit{gt}+1..r] > x$，$j$
從左往右掃：$A[j] < x$ 就和 $A[\mathit{lt}]$ 交換、兩者都加一；$A[j] > x$ 就和
$A[\mathit{gt}]$ 交換、$\mathit{gt}$ 減一（$j$ 不動）；相等就只把 $j$ 加一。 掃完傳
回 $q_1 = \mathit{lt}$、$q_2 = \mathit{gt}$。 選擇程序和 \proc{Randomized-Select} 一樣，
只是令 $k_1 = q_1 - p + 1$、$k_2 = q_2 - p + 1$：$i < k_1$ 往左找 $A[p..q_1 - 1]$ 的第
$i$ 小，$i > k_2$ 往右找 $A[q_2 + 1..r]$ 的第 $i - k_2$ 小，否則回傳 $x$。
全部相等時第一次分割就得到 $q_1 = p$、$q_2 = r$，任何 $i$ 都落在中間那段，分割一次就
回傳，比較 $O(n)$ 次。

\item 三組排好：$\langle 17, 26, 42, 58, 93 \rangle$、$\langle 10, 35, 63, 71, 84
\rangle$、$\langle 22, 31, 49, 77, 96 \rangle$，$M = \langle 42, 63, 49 \rangle$，$M$
的第 $2$ 小是 $x = 49$。 用 $49$ 分割後陣列是 $\langle 17, 26, 42, 10, 35, 22, 31, 49,$
$71, 84, 58, 93, 96, 77, 63 \rangle$，$q = 8$，$k = 8 = i$，第一層就找到中位數 $49$。
保證 $\ge 49$ 的：中位數 $\ge 49$ 的兩組（中位數 $63$、$49$）各 $3$ 個，即
$63, 71, 84, 49, 77, 96$，共 $6$ 個；保證 $\le 49$ 的也是 $6$ 個（$17, 26, 42, 22, 31,
49$）。

\item 中位數 $\ge x$ 的組有 $g - \lceil g/4 \rceil + 1$ 組，大約 $3g/4$，保證 $\ge x$
的數大約 $3 \cdot 3g/4 = 9n/20$，左邊最多大約 $11n/20$。 中位數 $\le x$ 的組只有
$\lceil g/4 \rceil$ 組，保證 $\le x$ 的大約 $3g/4 = 3n/20$，右邊最多大約
$17n/20 + O(1)$。 遞迴式是 $T(n) \le T(\lceil n/5 \rceil) + T(17n/20 + O(1)) + an$。
代換法：$c\frac{n}{5} + c\frac{17n}{20} + an = \frac{21}{20}cn + an > cn$，不管 $c$ 多
大都證不出 $T(n) \le cn$。 原因是 $\frac{1}{5} + \frac{17}{20} = \frac{21}{20} > 1$，遞
迴樹每往下一層，總工作量反而變大。 （用第~2~章第~\ref{sec:rec-gf}~節的 Akra--Bazzi 方法，
$(1/5)^p + (17/20)^p = 1$ 的解 $p \approx 1.12$，這個上界只能給出
$O(n^{1.12})$。）

\end{enumerate}

\addcontentsline{toc}{section}{歷屆試題詳解}
\section*{歷屆試題詳解}

\pastsol{px:sel-1}{中正資工所 113 學年度 第 20 題}
答案：False。

\textbf{上界。} 第~\ref{sec:sel-mom}~節的 \proc{Select} 只透過比較元素的大小得知資訊（各組的插入
排序、找 $x$ 的位置、Lomuto 分割都只比較元素），是比較模型裡的演算法；由定理~\ref{thm:sel-mom}，
它找第 $\lfloor (n+1)/2 \rfloor$ 小的最壞時間是 $O(n)$。 既然有比較模型裡的 $O(n)$ 演算法，
「任何」演算法都要 $\Omega(n \lg n)$ 就不對。

\textbf{為什麼排序的下界搬不過來。} 排序的 $\Omega(n \lg n)$ 來自「答案有 $n!$ 種」（第~15~章的
決策樹）。 找中位數只要回答「是哪一個元素」，答案只有 $n$ 種，決策樹只能給出 $\lceil \lg n
\rceil$ 的下界。

\textbf{真正的下界是線性。} 設演算法回答 $m$，而另一個元素 $y$ 一次都沒有參加比較。 $y < m$
時把 $y$ 改成比所有數都大，$y > m$ 時改成比所有數都小；所有比較的結果都不變，演算法仍然回答
$m$，但比 $m$ 小的數的個數變了，$m$ 不再是中位數。 所以除了 $m$，每個元素都至少要參加一次
比較，一次比較只碰兩個元素，至少要 $(n-1)/2$ 次。 找中位數在比較模型裡是 $\Theta(n)$。

\pastsol{px:sel-2}{臺大資工所 113 學年度 第 1 題}
答案：B、C、E。 以下 $\log$ 的底數不影響結論。
\begin{itemize}
\item A 錯誤。 堆積最多 $x + 1$ 個數，每次插入、刪除最大值是 $O(\log x)$，總共
$O(n \log x)$，最後的堆積排序 $O(x \log x)$。 問題在 $O(n \log x)$ 不一定是 $O(n)$。 取
$x = \lfloor n/\log n \rfloor$、$S$ 由小到大排好：堆積裝滿 $x$ 個之後，每個新來的數都比堆積
裡全部的數大，插入時要從新的葉節點一路上浮到根，交換 $\lfloor \log(x+1) \rfloor$ 次。 後面
$n - x \ge n/2$ 個數各花這麼多，共至少 $\frac{n}{2}\lfloor \log(x+1) \rfloor$。 $n \ge 16$ 時
$x + 1 \ge n/\log n \ge \sqrt{n}$，$\log(x+1) \ge \frac{1}{2}\log n$，所以是 $\Omega(n \log n)$，
不是 $O(n)$。
\item B 正確。 由下而上建最小堆積 $O(n)$（第~8~章第~\ref{sec:hp-build}~節），取出最小值 $x$ 次，
每次 $O(\log n)$，依取出的順序就是由小到大。 合計 $O(n + x \log n) \le
O(n + \frac{n}{\log n} \cdot \log n) = O(n)$。
\item C 正確。 每個數 $a_i \le y \le n^2$，以 $n$ 為基底最多 $3$ 位（$n^2$ 本身是 $100$）。
LSD 基數排序做 $3$ 趟，每趟是 $n$ 個數、$n$ 種位值的計數排序，$\Theta(n + n)$
（第~\ref{sec:lin-lsd}~節「一位要切多大」的最後一段）。 字組至少 $\lceil \log y \rceil$ 位元，
$a_i$ 放得進一個字組，取出某一位（除以 $n^t$ 再取餘數）是常數時間。 再輸出前 $x$ 個，
$O(x)$。 合計 $O(n)$。
\item D 錯誤。 光是配置並掃過長度 $y$ 的陣列 $C$ 就要 $\Theta(y)$；條件只說 $y \le n^2$，
$y = n^2$ 時是 $\Theta(n^2)$，不是 $O(n)$。 （$y = O(n)$ 時才是線性，見第~\ref{sec:lin-counting}~節的成本分析。）
\item E 正確。 \proc{Randomized-Select} 的期望時間是 $O(n)$（定理~\ref{thm:sel-rselect}）；
掃一次 $S$ 取出不大於第 $x$ 小的數是 $O(n)$，數兩兩不同，剛好取出 $x$ 個；（隨機化的）快速
排序排 $x$ 個數，期望 $O(x \log x)$（第~14~章定理~\ref{thm:so-quick-avg}）。 由期望值的線性，
總期望時間 $O(n + x \log x)$。
\end{itemize}

\pastsol{px:sel-3}{中央資工所 115 學年度 第 10 題}
答案：A、E。

先說明解讀。 (A) 和 (B) 照字面嚴格讀，結論會和出題的本意相反：(A) 的 $3n/10$ 要扣掉一個
常數才嚴格成立，(B) 的 $O(n \log n)$ 字面上也是對的上界。 這裡照出題的本意判斷（(A) 是課
本常見的近似說法；(B) 問的是這個遞迴式給出的界），下面把嚴格的情況也寫出來。
\begin{itemize}
\item A（照本意）正確。 Step 2 只用滿的組，共 $g = \lfloor n/5 \rfloor$ 組。 中位數不大於 $m$ 的組
至少 $\lceil g/2 \rceil$ 組（包括 $m$ 自己那組），每組的中位數和比它小的兩個數，共 $3$ 個都
$\le m$，所以至少 $3\lceil g/2 \rceil \approx 3n/10$ 個數 $\le m$。 嚴格地說，這個數可能比
$3n/10$ 少一點：$n = 11$，$A = \langle 1, 2, 3, 20, 21, 10, 11, 12, 22, 23, 30 \rangle$，最後的
$30$ 被忽略，兩組的中位數是 $3$、$12$，取下中位數 $m = 3$，$\le m$ 的只有 $3$ 個，比
$3n/10 = 3.3$ 少。 一般來說 $3\lceil g/2 \rceil \ge \frac{3}{2}\lfloor \frac{n}{5} \rfloor
\ge \frac{3(n-4)}{10}$，所以保證的是至少 $\frac{3n}{10} - \frac{6}{5}$ 個。
\item B（照本意）錯誤。 $\frac{1}{5} + \frac{7}{10} = \frac{9}{10} < 1$，這個遞迴式的解是
$O(n)$（定理~\ref{thm:sel-mom} 的代換法，或第~\ref{sec:sel-mom}~節的遞迴樹：每一層的工作量是上
一層的 $9/10$）。 演算法是線性時間，說它「是 $O(n \log n)$」沒有給出它推得出的界。 嚴格地
說，$O(n)$ 的函數也是 $O(n \log n)$，所以字面上這句話不算錯；作答時可以註明這一點。
\item C 錯誤。 $3$ 個一組時，保證不大於 $m$ 的數只有大約 $2 \cdot \frac{g}{2} = \frac{n}{3}$
個，往下遞迴的一邊可能有 $\frac{2n}{3}$ 個，遞迴式是 $T(n) = T(n/3) + T(2n/3) + O(n)$，比例
加起來是 $1$，解是 $\Theta(n \lg n)$（第~\ref{sec:sel-mom}~節「為什麼是 5 個一組」）。 不能保
證線性。
\item D 錯誤。 正確性不需要 $m$ 接近中位數：Step 5 之後 $L$ 的數都比 $m$ 小、$R$ 的都比 $m$
大，$m$ 是第 $|L| + 1$ 小；$k$ 落在哪一段、在那一段裡是第幾小都由此決定，不管 $m$ 是誰都對
（和本章觀念題第~\ref{sel-q-correct}~題的歸納法一樣）。 $m$ 接近中位數只影響\textbf{時間}：它
保證 $L$、$R$ 都不超過大約 $7n/10$。
\item E 正確。 中位數就是第 $\lceil n/2 \rceil$ 小（$n$ 是偶數時是下中位數，和本書的約定相同），
是選擇問題 $k = \lceil n/2 \rceil$ 的特例；\proc{Select} 對任何 $k$ 的最壞時間都是 $O(n)$
（定理~\ref{thm:sel-mom}）。
\end{itemize}

\pastsol{px:sel-4}{成大資訊聯招 109 學年度 演算法第 3 題}
快速排序本來的期望時間就是 $O(n \lg n)$（第~14~章定理~\ref{thm:so-quick-avg}），所以題目要的
是\textbf{最壞情況}也是 $O(n \lg n)$。 最壞情況發生在每次都切得極偏；如果每次都用
\textbf{中位數}當基準，兩邊各剩大約一半，而第~\ref{sec:sel-mom}~節的 \proc{Select} 可以在最壞
$O(n)$ 時間找到中位數。 以下假設鍵互不相同。

\begin{pseudo}{成大 109 學年度演算法第 3 題的解答}
\Procedure{Median-Quicksort}{$A, p, r$}
    \If{$p < r$}
        \State $x \gets$ \Call{Select}{$A, p, r, \lfloor (r - p + 2)/2 \rfloor$} \Comment{$A[p..r]$ 的中位數}
        \State $q \gets$ \Call{Partition-Around}{$A, p, r, x$}
        \State \Call{Median-Quicksort}{$A, p, q - 1$}
        \State \Call{Median-Quicksort}{$A, q + 1, r$}
    \EndIf
\EndProcedure
\end{pseudo}

\proc{Select} 和 \proc{Partition-Around} 都是虛擬碼~\ref{alg:sel-mom} 裡的程序。 範圍有
$n = r - p + 1$ 個數，$\lfloor (r - p + 2)/2 \rfloor = \lfloor (n+1)/2 \rfloor$，就是中位數的名次。

\textbf{正確性。} \proc{Select} 只在 $A[p..r]$ 裡搬動元素，不改變這一段裝的是哪些數。
\proc{Partition-Around} 把 $x$ 換到 $A[r]$ 再做 Lomuto 分割，結束時
$A[p..q-1] < A[q] = x < A[q+1..r]$。 之後和第~14~章的 \proc{Quicksort}（虛擬碼~\ref{alg:so-quick}）
一樣：對 $n$ 歸納，兩邊遞迴的範圍都比 $n$ 小，由歸納假設各自排好，左邊每個數小於 $x$、右邊
每個數大於 $x$，所以整段排好。 快速排序的正確性不依賴基準選誰，換成中位數也一樣。

\textbf{時間。} 設 $T(n)$ 是 $n$ 個數時的最壞時間。 第 3 行由定理~\ref{thm:sel-mom} 是 $O(n)$，
第 4 行找 $x$ 的位置和分割都是 $O(n)$，合計不超過 $cn$。 $x$ 是第 $\lceil n/2 \rceil$ 小，左邊有
$\lceil n/2 \rceil - 1 \le \lfloor n/2 \rfloor$ 個數，右邊有 $n - \lceil n/2 \rceil = \lfloor n/2
\rfloor$ 個數。 所以
\[
T(n) \le 2T(\lfloor n/2 \rfloor) + cn .
\]
用遞迴樹看：第 $d$ 層的每個範圍不超過 $n/2^d$ 個數，同一層的範圍互不重疊，所以每一層的工作量
加起來不超過 $cn$；範圍每往下一層至少減半，$\lfloor \lg n \rfloor + 1$ 層以內就只剩 $0$ 或 $1$
個數。 總時間 $O(n \lg n)$，這也是主定理（第~2~章定理~\ref{thm:rec-master}）$a = b = 2$、
$f(n) = \Theta(n)$ 的情況 2。

\textbf{補充。} 這個做法的常數很大（\proc{Select} 本身就要很多次比較），實務上不會這樣寫。 常見
的做法是隨機選基準，只保證期望 $O(n \lg n)$；要保證最壞情況，C++ 程式庫的 introsort 在遞迴太
深時改用堆積排序（第~14~章第~\ref{sec:so-heap}~節最後）。 作答時寫出上面的演算法和遞迴式就夠了，
可以加一句說明這兩種實務做法。 鍵有重複時，把分割改成三路分割（第~14~章觀念題第 6 題），等於
$x$ 的數全部留在中間，兩邊仍然各不超過一半。


% 第四部：演算法設計
\chapter{分治}
\label{ch:dc}

% 本章的圖共用的樣式（名稱都以 dc 開頭，避免和其他章衝突）
\tikzset{
    dc-cell/.style={draw=black!55, minimum width=0.8cm, minimum height=0.62cm, inner sep=0pt, font=\small},
    dc-run/.style={dc-cell, fill=coreFill, draw=coreStroke},
    dc-ghost/.style={dc-cell, draw=black!30, dashed, text=black!45},
    dc-arr/.style={-{Stealth[length=1.8mm]}, draw=black!65},
    dc-tn/.style={draw=black!55, rounded corners=2pt, inner sep=2.5pt, font=\scriptsize, align=center, fill=white},
    dc-tnw/.style={dc-tn, draw=memStroke, thick, fill=memFill},
    dc-leaf/.style={draw=black!40, rounded corners=2pt, inner sep=2pt, font=\scriptsize, minimum width=0.75cm, align=center},
    dc-edge/.style={draw=black!55},
    dc-pt/.style={circle, fill=black!80, inner sep=0pt, minimum size=4.2pt},
    dc-ptlab/.style={font=\scriptsize, inner sep=1.5pt},
    dc-ray/.style={draw=black!35, dashed},
    dc-rayq/.style={-{Stealth[length=1.8mm]}, dashed, draw=optStroke},
    dc-vec/.style={-{Stealth[length=2mm]}, thick, draw=black!75},
    dc-seg/.style={thick, draw=black!80},
    dc-seg2/.style={thick, draw=coreStroke},
}

夜市一個飲料攤把開張以來每天的盈虧記在本子上，賺錢記正數，賠錢記負數。 老闆想找
出「連續幾天加起來賺最多」的那一段，回頭看那幾天做對了什麼。 一年有三百多天，每
一段都加加看，要試五萬多段。 如果先把這一年從中間切成上半年和下半年，賺最多的那
一段只可能落在哪幾種地方？ 每一種地方各要花多少工夫才找得到？ 往下讀之前，先想
一想。

第~14~章的合併排序把陣列切成兩半，兩半各自排好，再合併起來。 這種做法叫
\textbf{分治法}（divide and conquer），每一層遞迴都做三件事：

\begin{enumerate}
\item \textbf{分}（divide）：把問題切成幾個同類型、規模比較小的子問題。
\item \textbf{治}（conquer）：用同樣的方法遞迴解每個子問題；規模夠小時直接解。
\item \textbf{合}（combine）：把子問題的答案組合成原問題的答案。
\end{enumerate}

切成 $a$ 個規模 $n/b$ 的子問題、分和合一共花 $f(n)$ 時，執行時間是
$T(n) = aT(n/b) + f(n)$，用第~2~章的遞迴樹或主定理（定理~\ref{thm:rec-master}）就能解。 這個式子
也告訴我們，想讓分治法變快，有兩個方向：讓 $f(n)$ 變小（合得更省），或讓 $a$ 變
小（少解幾個子問題）。 這一章的前五節分別示範幾種情形，第六節是計算幾何入門。

\begin{itemize}
\item 第~\ref{sec:dc-bsearch}~節的二分搜尋和快速冪只需要解\textbf{一個}子問題，
$a = 1$。
\item 第~\ref{sec:dc-maxsub}~節的最大子陣列，難的是「合」：最好的一段可能跨過切
開的地方。 這一節最後會用一個線性時間的做法和分治版對照。
\item 第~\ref{sec:dc-mult}~節的 Karatsuba 大數乘法和 Strassen 矩陣乘法，靠代數把
子問題的個數從 $4$ 減成 $3$、從 $8$ 減成 $7$，指數就跟著變小。
\item 第~\ref{sec:dc-closest}~節的平面最近點對，「合」的步驟要靠一個幾何的論證，
才能在線性時間內完成。
\item 第~\ref{sec:dc-inv}~節的逆序數計數，把計數藏在合併排序的合併步驟裡。
\item 第~\ref{sec:dc-geom}~節是計算幾何入門：用外積判斷轉向、線段相交、點在多邊形內，
以及凸包。 凸包的分治做法和最近點對同一個結構；它的 $\Omega(n \lg n)$ 下界則是把排序
歸約到凸包得到的。
\end{itemize}

子問題如果彼此重疊，同一個子問題會被解很多次，分治法就不划算了，這時要改用第~18~
章的動態規劃。

這一章用到第~1~章的迴圈不變量和逆序對、第~2~章的遞迴式和主定理、第~14~章的合併；
最後一節還用到第~3~章的聚合法、第~6~章的堆疊和第~15~章的排序下界。
C 程式用到《C 程式設計基礎》第~9~章的指標參數、第~14~章的結構、第~15~章的遞迴、第~18~章
的 \texttt{qsort} 和第~19~章的 \texttt{malloc}；遞迴在第~2~章複習過，其他的用到時先在框
裡複習。

\section{二分搜尋與快速冪}
\label{sec:dc-bsearch}

分治不一定要把每一份都解。 如果比一次就能判斷答案只可能在哪一份，另一份就整個丟
掉，遞迴式變成 $T(n) = T(n/b) + f(n)$。 有些書把這種只留一個子問題的情形叫做「減
而治之」（decrease and conquer）。 這一節的兩個例子都屬於這一類。

\subsection{找一個值}

輸入是由小到大排好（允許相等）的 $A[1..n]$ 和一個數 $x$。 如果 $x$ 在裡面，回傳
一個位置 $i$ 使 $A[i] = x$；不在就回傳 \NIL。 第~2~章的虛擬碼~\ref{alg:rec-search} 是遞迴版：看中
間那一格，相等就回答，比 $x$ 小就到右半邊找，比 $x$ 大就到左半邊找。 遞迴呼叫是
函式的最後一步（第~6~章的尾遞迴），所以可以直接改成迴圈，用 $\mathit{lo}$、
$\mathit{hi}$ 記住還要找的範圍 $A[\mathit{lo}..\mathit{hi}]$。

\begin{pseudo}[label=alg:dc-bsearch]{二分搜尋（閉區間）}
\Function{Binary-Search}{$A, n, x$}
    \State $\mathit{lo} \gets 1$；$\mathit{hi} \gets n$
    \While{$\mathit{lo} \le \mathit{hi}$}
        \State $\mathit{mid} \gets \lfloor (\mathit{lo} + \mathit{hi}) / 2 \rfloor$
        \If{$A[\mathit{mid}] = x$}
            \State \Return $\mathit{mid}$
        \ElsIf{$A[\mathit{mid}] < x$}
            \State $\mathit{lo} \gets \mathit{mid} + 1$
        \Else
            \State $\mathit{hi} \gets \mathit{mid} - 1$
        \EndIf
    \EndWhile
    \State \Return \NIL
\EndFunction
\end{pseudo}

\textbf{不變量}：每次檢查第 3 行的條件時，$A[1..\mathit{lo}-1]$ 的每個數都小於
$x$，$A[\mathit{hi}+1..n]$ 的每個數都大於 $x$。 照第~\ref{sec:ana-invariant}~節的三步來證。

\begin{itemize}
\item \textbf{初始}：$\mathit{lo} = 1$、$\mathit{hi} = n$，兩段都是空的，敘述成立。
\item \textbf{維持}：迴圈裡 $\mathit{lo} \le \mathit{mid} \le \mathit{hi}$，所以
$A[\mathit{mid}]$ 存在。 如果 $A[\mathit{mid}] = x$，程序直接回傳，不必維持。 如
果 $A[\mathit{mid}] < x$，因為陣列排好了，$A[\mathit{lo}..\mathit{mid}]$ 都不大於
$A[\mathit{mid}]$，也就都小於 $x$；第 8 行把 $\mathit{lo}$ 改成 $\mathit{mid} +
1$ 之後，$A[1..\mathit{lo}-1]$ 仍然都小於 $x$。 $A[\mathit{mid}] > x$ 的情形對稱，
$A[\mathit{mid}..\mathit{hi}]$ 都大於 $x$，第 10 行之後第二段仍然成立。
\item \textbf{結束}：如果是第 6 行回傳，$A[\mathit{mid}] = x$，答案正確。 否則迴
圈停在 $\mathit{lo} = \mathit{hi} + 1$（每一輪 $\mathit{lo}$ 最多跳到
$\mathit{hi} + 1$，$\mathit{hi}$ 最少降到 $\mathit{lo} - 1$），代進不變量：
$A[1..\mathit{hi}]$ 都小於 $x$，$A[\mathit{hi}+1..n]$ 都大於 $x$，整個陣列沒有一
個等於 $x$，回傳 \NIL{} 正確。
\end{itemize}

\textbf{一定會停，而且很快}。 令 $k = \mathit{hi} - \mathit{lo} + 1$ 是還要找的
格數。 中間那格左邊剩 $\mathit{mid} - \mathit{lo} = \lfloor (k-1)/2 \rfloor$ 格，
右邊剩 $\mathit{hi} - \mathit{mid} = \lceil (k-1)/2 \rceil = \lfloor k/2 \rfloor$
格，所以每一輪之後 $k$ 最多剩 $\lfloor k/2 \rfloor$。 $t$ 輪之後
$k \le \lfloor n / 2^t \rfloor$；$t = \lfloor \lg n \rfloor + 1$ 時 $n / 2^t < 1$，
範圍已經是空的。 所以迴圈最多跑 $\lfloor \lg n \rfloor + 1$ 輪，每一輪 $\Theta(1)$，
最壞情況 $\Theta(\lg n)$。 這和第~2~章算出的遞迴版一致。

\subsection{第一個大於等於、第一個大於}

用 $A = \langle 12, 15, 15, 15, 23, 31, 31, 40 \rangle$ 找 $x = 15$：第一輪
$\mathit{mid} = 4$，$A[4] = 15$，馬上回傳 $4$。 $15$ 出現在第 $2$、$3$、$4$ 格，
\proc{Binary-Search} 回傳的是\textbf{其中一個}，不一定是第一個。 很多問題要的是
邊界：第一個 $15$ 在哪裡、$15$ 有幾個、比 $20$ 小的最後一個數在哪裡。 這時改用兩
個「找邊界」的版本：

\begin{itemize}
\item \proc{Lower-Bound}$(A, n, x)$：第一個滿足 $A[i] \ge x$ 的位置 $i$；全部都小
於 $x$ 時回傳 $n + 1$。
\item \proc{Upper-Bound}$(A, n, x)$：第一個滿足 $A[i] > x$ 的位置 $i$；全部都不大
於 $x$ 時回傳 $n + 1$。
\end{itemize}

\proc{Lower-Bound} 的虛擬碼和正確性證明在第~1~章觀念題第~\ref{ana-q-lower-bound}~題。 它用的是
\textbf{半開區間}：答案一定在 $\mathit{lo}..\mathit{hi}$ 之間（可能等於
$\mathit{hi}$），$\mathit{hi}$ 從 $n + 1$ 開始，迴圈條件是 $\mathit{lo} <
\mathit{hi}$，往左縮時寫 $\mathit{hi} \gets \mathit{mid}$。 \proc{Upper-Bound} 只
差一個比較符號：

\begin{pseudo}[label=alg:dc-upper]{第一個大於 $x$ 的位置（半開區間）}
\Function{Upper-Bound}{$A, n, x$}
    \State $\mathit{lo} \gets 1$；$\mathit{hi} \gets n + 1$
    \While{$\mathit{lo} < \mathit{hi}$}
        \State $\mathit{mid} \gets \lfloor (\mathit{lo} + \mathit{hi}) / 2 \rfloor$
        \If{$A[\mathit{mid}] \le x$}
            \State $\mathit{lo} \gets \mathit{mid} + 1$
        \Else
            \State $\mathit{hi} \gets \mathit{mid}$
        \EndIf
    \EndWhile
    \State \Return $\mathit{lo}$
\EndFunction
\end{pseudo}

不變量換成：$A[1..\mathit{lo}-1]$ 都不大於 $x$，$A[\mathit{hi}..n]$ 都大於 $x$。
證明和第~1~章那一題逐字相同，只要把「小於 $x$」換成「不大於 $x$」、「大於等於
$x$」換成「大於 $x$」。 每一輪 $\mathit{hi} - \mathit{lo}$ 從 $k$ 變成
$\lfloor k/2 \rfloor$ 或 $\lceil k/2 \rceil - 1$，都不超過 $\lfloor k/2 \rfloor$，
所以迴圈最多跑 $\lfloor \lg n \rfloor + 1$ 輪。 兩個版本在找不到時都不必另外處
理，這是半開區間寫法的好處。

用上面的陣列手算。 每一列是一輪迴圈開始時的 $\mathit{lo}$、$\mathit{hi}$，以及這
一輪的 $\mathit{mid}$ 和改了哪一邊。

\begin{center}
\begin{tabular}{cccccl}
\toprule
& $\mathit{lo}$ & $\mathit{hi}$ & $\mathit{mid}$ & $A[\mathit{mid}]$ & 動作 \\
\midrule
\proc{Lower-Bound}$(A, 8, 15)$ & 1 & 9 & 5 & 23 & $23 \ge 15$，$\mathit{hi} \gets 5$ \\
& 1 & 5 & 3 & 15 & $15 \ge 15$，$\mathit{hi} \gets 3$ \\
& 1 & 3 & 2 & 15 & $15 \ge 15$，$\mathit{hi} \gets 2$ \\
& 1 & 2 & 1 & 12 & $12 < 15$，$\mathit{lo} \gets 2$ \\
& 2 & 2 & & & 停，回傳 $2$ \\
\midrule
\proc{Upper-Bound}$(A, 8, 15)$ & 1 & 9 & 5 & 23 & $23 > 15$，$\mathit{hi} \gets 5$ \\
& 1 & 5 & 3 & 15 & $15 \le 15$，$\mathit{lo} \gets 4$ \\
& 4 & 5 & 4 & 15 & $15 \le 15$，$\mathit{lo} \gets 5$ \\
& 5 & 5 & & & 停，回傳 $5$ \\
\bottomrule
\end{tabular}
\end{center}

$15$ 佔了第 $2$ 到第 $4$ 格，下界 $2$ 指到第一個 $15$，上界 $5$ 指到最後一個 $15$
的下一格，$5 - 2 = 3$ 正好是 $15$ 的個數。

\begin{center}
\begin{minipage}{\linewidth}
\centering
\begin{tikzpicture}
\foreach \v [count=\k] in {12,15,15,15,23,31,31,40} {
    \ifnum\k>1 \ifnum\k<5
        \node[dc-run] (c\k) at (\k*0.9, 0) {\v};
    \else
        \node[dc-cell] (c\k) at (\k*0.9, 0) {\v};
    \fi \else
        \node[dc-cell] (c\k) at (\k*0.9, 0) {\v};
    \fi
    \node[idxlbl] at (\k*0.9, 0.48) {\k};
}
\node[dc-ghost] (c9) at (8.1, 0) {};
\node[idxlbl] at (8.1, 0.48) {9};
\draw[dc-arr] (0.9, -1.25) node[left, font=\scriptsize] {\proc{Lower-Bound}$(15) = 2$} -- (c2.south);
\draw[dc-arr] (6.0, -1.25) node[right, font=\scriptsize] {\proc{Upper-Bound}$(15) = 5$} -- (c5.south);
\draw[dc-arr] (3.6, 1.5) node[left, font=\scriptsize] {\proc{Lower-Bound}$(31) = 6$} -- ($(c6.north)+(0,0.36)$);
\draw[dc-arr] (8.6, 1.5) node[right, font=\scriptsize] {\proc{Upper-Bound}$(31) = 8$} -- ($(c8.north)+(0,0.36)$);
\end{tikzpicture}
\diagcaption{兩個邊界版本的回傳值。 藍底是 $15$ 佔的範圍，$5 - 2 = 3$ 格；第 $9$ 格是虛線，表示
「超過最後一格」的 $n + 1$。}
\label{fig:dc-bounds}
\end{minipage}
\end{center}

有了這兩個函式，常見的邊界問題都可以直接寫出來：

\begin{center}
\begin{tabular}{ll}
\toprule
要找的 & 寫法 \\
\midrule
$x$ 出現幾次 & \proc{Upper-Bound}$(x) - {}$\proc{Lower-Bound}$(x)$ \\
第一個 $x$ 的位置 & $i = {}$\proc{Lower-Bound}$(x)$，還要檢查 $i \le n$ 且 $A[i] = x$ \\
最後一個 $x$ 的位置 & $i = {}$\proc{Upper-Bound}$(x) - 1$，還要檢查 $i \ge 1$ 且 $A[i] = x$ \\
最後一個小於 $x$ 的位置 & \proc{Lower-Bound}$(x) - 1$（$0$ 表示沒有）\\
最後一個不大於 $x$ 的位置 & \proc{Upper-Bound}$(x) - 1$（$0$ 表示沒有）\\
落在 $[u, v]$ 裡的個數 & \proc{Upper-Bound}$(v) - {}$\proc{Lower-Bound}$(u)$（$u \le v$）\\
\bottomrule
\end{tabular}
\end{center}

兩個版本其實是同一件事。 位置 $1, 2, \ldots, n$ 上有一個條件 $P(i)$，前面一段都不
成立、後面一段都成立（叫做\textbf{單調}的條件），要找第一個成立的位置。
\proc{Lower-Bound} 的 $P(i)$ 是「$A[i] \ge x$」，\proc{Upper-Bound} 的是
「$A[i] > x$」。 只要條件是單調的，即使沒有陣列，例如「容量 $c$ 夠不夠裝下所有
貨物」，同一個迴圈也能找出最小的可行值，每一輪只要能判斷一次 $P(\mathit{mid})$。

\subsection{邊界的三個陷阱}
\label{sec:dc-traps}

二分搜尋的程式很短，卻很容易寫錯。 錯誤幾乎都出在邊界，下面三種最常見。

\textbf{區間的寫法要一致}。 閉區間 $[\mathit{lo}, \mathit{hi}]$ 配
$\mathit{lo} \le \mathit{hi}$ 和 $\mathit{hi} \gets \mathit{mid} - 1$；半開區間
$[\mathit{lo}, \mathit{hi})$ 配 $\mathit{lo} < \mathit{hi}$ 和
$\mathit{hi} \gets \mathit{mid}$。 混用就會出事：如果條件寫成
$\mathit{lo} \le \mathit{hi}$，卻用 $\mathit{hi} \gets \mathit{mid}$，那麼範圍縮到
只剩一格、而且 $A[\mathit{mid}] \ge x$ 時，$\mathit{lo} = \mathit{hi} =
\mathit{mid}$，$\mathit{hi} \gets \mathit{mid}$ 什麼都沒改，迴圈永遠不停。

\textbf{往右縮成 $\mathit{lo} \gets \mathit{mid}$ 時，中點要取天花板}。 例如要找
「最後一個不大於 $x$ 的位置」，有人寫成：$A[\mathit{mid}] \le x$ 就
$\mathit{lo} \gets \mathit{mid}$，否則 $\mathit{hi} \gets \mathit{mid} - 1$。 範圍
只剩兩格（$\mathit{hi} = \mathit{lo} + 1$）時，$\lfloor (\mathit{lo} + \mathit{hi})
/ 2 \rfloor = \mathit{lo}$，如果 $A[\mathit{lo}] \le x$，$\mathit{lo}$ 原地不動，
又是無窮迴圈。 改成 $\mathit{mid} \gets \lceil (\mathit{lo} + \mathit{hi}) / 2
\rceil$ 就保證 $\mathit{mid} > \mathit{lo}$。 更簡單的做法是不要自己寫這一種，直接
用表上的 \proc{Upper-Bound}$(x) - 1$。

\textbf{C 的 $\mathit{lo} + \mathit{hi}$ 會溢位}。 虛擬碼裡的數不會溢位，C 的
\texttt{int} 會。 下標接近 $2^{31}$ 時，\texttt{(lo + hi) / 2} 先算出一個超過
\texttt{INT\_MAX} 的和，這是未定義行為。 改寫成 \texttt{lo + (hi - lo) / 2}，算出
來的值相同，中間值卻不會超過 \texttt{hi}。

檢查一個二分搜尋寫得對不對，最有效的是手算幾個小例子：$n = 1$、$n = 2$；$x$ 比全
部都小、比全部都大；$x$ 等於第一格、等於最後一格；$x$ 出現好幾次。 兩格的例子特
別能抓出無窮迴圈。

\begin{examnote}
二分搜尋的題目先看清楚三件事：位置從 $0$ 還是從 $1$ 開始；找不到時要回傳什麼
（$-1$、\NIL、$n$ 還是 $n + 1$）；有重複時要第一個、最後一個還是任一個。 不同課本
的「上界」「下界」定義不完全一樣，照題目的定義寫。 作答時把迴圈不變量寫在程式旁
邊，再用「只剩兩格」的情況自己代一次，確認迴圈會停。
\end{examnote}

\subsection{C 程式}

C 版只寫兩個邊界版本，對上面的陣列做幾個查詢，印出下界、上界、出現次數、第一個
$x$ 的位置（沒有就印 $-1$）和最後一個不大於 $x$ 的位置，後兩個照上表的寫法算。 C 的
下標從 $0$ 開始，所以 \texttt{lower\_bound} 回傳的是虛擬碼的 \proc{Lower-Bound} 減
一，全部都小於 $x$ 時回傳 \texttt{n}；次數是兩者相減，不受影響。

\begin{example}[label=ex:dc-bsearch]{c}{(bsearch\_bounds.c) 兩個邊界版本的二分搜尋}
#include <stdio.h>

// 半開區間 [lo, hi)：第一個 >= x 的位置，全部都 < x 時回傳 n。
int lower_bound(const int *a, int n, int x)
{
    int lo = 0, hi = n;
    while (lo < hi) {
        int mid = lo + (hi - lo) / 2;      // 不寫 (lo + hi) / 2，避免溢位
        if (a[mid] < x) {
            lo = mid + 1;
        } else {
            hi = mid;
        }
    }
    return lo;
}

// 第一個 > x 的位置，全部都 <= x 時回傳 n。只有比較的方向不同。
int upper_bound(const int *a, int n, int x)
{
    int lo = 0, hi = n;
    while (lo < hi) {
        int mid = lo + (hi - lo) / 2;
        if (a[mid] <= x) {
            lo = mid + 1;
        } else {
            hi = mid;
        }
    }
    return lo;
}

int main(void)
{
    int a[] = {12, 15, 15, 15, 23, 31, 31, 40};
    int q[] = {15, 31, 20, 5, 12, 40, 41};
    printf("x lower upper count first last_le\n");
    for (int k = 0; k < 7; k++) {
        int lb = lower_bound(a, 8, q[k]), ub = upper_bound(a, 8, q[k]);
        int first = (lb < 8 && a[lb] == q[k]) ? lb : -1;    // 第一個 x 的位置
        printf("%d %d %d %d %d %d\n", q[k], lb, ub, ub - lb, first, ub - 1);
    }
    return 0;
}
\end{example}
\iosample{}{x lower upper count first last\_le \\ 15 1 4 3 1 3 \\ 31 5 7 2 5 6 \\ 20 4 4 0 -1 3 \\ 5 0 0 0 -1 -1 \\ 12 0 1 1 0 0 \\ 40 7 8 1 7 7 \\ 41 8 8 0 -1 7}

$20$ 不在陣列裡，下界和上界都是 $4$：$20$ 要插進去的話應該放在 C 的位置 $4$，也就
是 $23$ 的前面，而最後一個不大於 $20$ 的是位置 $3$ 的 $15$。 $5$ 比全部都小，最後一
欄是 $-1$，表示沒有。

\subsection{快速冪}
\label{sec:dc-power}

第二個例子是算 $x^n$，$n$ 是非負整數。 從 $1$ 開始乘 $n$ 次 $x$ 要做 $n$ 次乘法
（第一次乘 $1$ 可以省掉，是 $n - 1$ 次）。 分治的想法是先算一半：
\[
x^n = \begin{cases}
\left(x^{\lfloor n/2 \rfloor}\right)^2, & n \text{ 是偶數}, \\[2pt]
\left(x^{\lfloor n/2 \rfloor}\right)^2 \cdot x, & n \text{ 是奇數}.
\end{cases}
\]
$x^{\lfloor n/2 \rfloor}$ 只要算一次，平方是把它和自己相乘，不必再遞迴第二次。

\begin{pseudo}[label=alg:dc-power]{快速冪（遞迴）}
\Function{Power}{$x, n$}
    \If{$n = 0$}
        \State \Return $1$
    \EndIf
    \If{$n = 1$}
        \State \Return $x$
    \EndIf
    \State $h \gets$ \Call{Power}{$x, \lfloor n/2 \rfloor$}
    \State $r \gets h \cdot h$
    \If{$n$ 是奇數}
        \State $r \gets r \cdot x$
    \EndIf
    \State \Return $r$
\EndFunction
\end{pseudo}

每次呼叫只往下呼叫一次，規模減半，呼叫以外做常數件事，所以
$T(n) = T(\lfloor n/2 \rfloor) + \Theta(1) = \Theta(\lg n)$，就是第~2~章主定理表格
裡 $T(n/2) + 1$ 那一列。 乘法次數可以算得更精確。 把 $n$ 寫成二進位：$n \ge 2$
時，$\lfloor n/2 \rfloor$ 的二進位就是 $n$ 去掉最右邊一位。 遞迴從 $n$ 一路去掉最
低位，直到只剩最高位的 $1$（$n = 1$），途中每去掉一位做一次平方（第 7 行），去掉
的那一位是 $1$ 時再多乘一次 $x$（第 9 行）。 令 $\nu(n)$ 是 $n$ 的二進位裡 $1$ 的
個數，$n$ 有 $\lfloor \lg n \rfloor + 1$ 位，所以
\[
\text{乘法次數} = \lfloor \lg n \rfloor + (\nu(n) - 1) \le 2 \lfloor \lg n \rfloor
\qquad (n \ge 1).
\]
例如 $n = 45 = 101101_2$，$\lfloor \lg 45 \rfloor = 5$，$\nu(45) = 4$，共 $8$ 次乘
法，比連乘的 $44$ 次少很多。

遞迴是從最高位往下處理的。 也可以從最低位往上，寫成迴圈，不用遞迴堆疊：

\begin{pseudo}[label=alg:dc-power-iter]{快速冪（迴圈，從最低位開始）}
\Function{Power-Iter}{$x, n$}
    \State $r \gets 1$；$b \gets x$；$e \gets n$
    \While{$e > 0$}
        \If{$e$ 是奇數}
            \State $r \gets r \cdot b$
        \EndIf
        \State $e \gets \lfloor e/2 \rfloor$
        \If{$e > 0$}
            \State $b \gets b \cdot b$
        \EndIf
    \EndWhile
    \State \Return $r$
\EndFunction
\end{pseudo}

\textbf{不變量}：每次檢查第 3 行的條件時，$r \cdot b^e = x^n$。

\begin{itemize}
\item \textbf{初始}：$r = 1$、$b = x$、$e = n$，$1 \cdot x^n = x^n$。
\item \textbf{維持}：設這一輪開始時 $r b^e = x^n$。 若 $e = 2k + 1$ 是奇數，
$r b^e = (r b) \cdot (b^2)^k$；第 5 行把 $r$ 換成 $rb$，第 6 行把 $e$ 換成 $k$，第
8 行把 $b$ 換成 $b^2$（$k = 0$ 時不換，但這時 $b^k = 1$，換不換都一樣），所以新的
$r b^e$ 仍然是 $x^n$。 若 $e = 2k$ 是偶數，$r b^e = r (b^2)^k$，同樣成立。
\item \textbf{結束}：迴圈停在 $e = 0$，代進不變量得 $r \cdot b^0 = r = x^n$。
\end{itemize}

$e$ 每一輪去掉一個二進位位數，迴圈跑 $\lfloor \lg n \rfloor + 1$ 輪。 第 5 行做
$\nu(n)$ 次，第 8 行做 $\lfloor \lg n \rfloor$ 次，合計比遞迴版多一次，多的那次是
一開始的 $1 \cdot b$。

實際應用時 $x^n$ 大得存不下，通常要的是 $x^n$ 除以某個 $m$ 的餘數。 因為
\[
(u \cdot v) \bmod m = \bigl((u \bmod m)(v \bmod m)\bigr) \bmod m,
\]
每次乘完就取餘數，結果不變，而且所有中間值都小於 $m$。 C 程式要注意的只剩一件事：
兩個小於 $m$ 的數相乘，最大接近 $m^2$，必須放得進 \texttt{long long}，所以 $m$ 不
能超過 $\sqrt{2^{63} - 1} \approx 3.04 \times 10^9$。

用 $x = 7$、$n = 45$、$m = 1000$ 手算 \proc{Power-Iter}，每一列是一輪開始時的值。
$45 = 101101_2$，從最右邊一位讀起。

\begin{center}
\begin{tabular}{ccccl}
\toprule
$e$ & $e$ 的二進位 & $r$ & $b$ & 這一輪做的事（都取除以 $1000$ 的餘數）\\
\midrule
45 & 101101 & 1 & 7 & 奇數：$r = 1 \cdot 7 = 7$；$b = 7^2 = 49$ \\
22 & 10110 & 7 & 49 & 偶數；$b = 49^2 = 2401 \to 401$ \\
11 & 1011 & 7 & 401 & 奇數：$r = 7 \cdot 401 = 2807 \to 807$；$b = 401^2 \to 801$ \\
5 & 101 & 807 & 801 & 奇數：$r = 807 \cdot 801 \to 407$；$b = 801^2 \to 601$ \\
2 & 10 & 407 & 601 & 偶數；$b = 601^2 \to 201$ \\
1 & 1 & 407 & 201 & 奇數：$r = 407 \cdot 201 \to 807$；$e$ 變成 $0$，不再平方 \\
\bottomrule
\end{tabular}
\end{center}

答案是 $7^{45} \bmod 1000 = 807$，共 $\nu(45) + \lfloor \lg 45 \rfloor = 4 + 5 = 9$
次乘法。

快速冪是第~\ref{sec:ana-input-size}~節那個區別的好例子。 輸入的 $n$ 只有 $\lg n$ 個位元。 連乘 $n$
次是位元數的指數；快速冪只要 $\Theta(\lg n)$ 次乘法，是位元數的多項式。

同一個演算法只用到「乘法有結合律」，所以把數換成矩陣也可以。 例如費氏數列
$F_0 = 0$、$F_1 = 1$、$F_{k} = F_{k-1} + F_{k-2}$，對 $n \ge 1$ 歸納可以證明
\[
\begin{pmatrix} 1 & 1 \\ 1 & 0 \end{pmatrix}^{n}
= \begin{pmatrix} F_{n+1} & F_n \\ F_n & F_{n-1} \end{pmatrix},
\]
（$n = 1$ 時兩邊相等；乘上一次左邊的矩陣，右邊的每一格正好變成下一項。）所以
$F_n \bmod m$ 只要 $\Theta(\lg n)$ 次 $2 \times 2$ 矩陣乘法，每次 $8$ 次數的乘
法。 第~25~章的所有點對最短路徑也會用到這種「反覆平方」。

\subsection{C 程式}

\begin{cbookrecap}[第~9~章 9.2 節]
函式只能回傳一個值。 \texttt{power\_mod} 另外用指標參數 \texttt{int *mults} 把乘
法次數帶回去：函式裡的 \texttt{(*mults)++} 改的是呼叫端的變數，呼叫時傳 \texttt{\&c}。
\end{cbookrecap}

C 版只寫迴圈版，每次乘完就取餘數。 測資裡
$1000000007$ 是質數。 數論的\textbf{費馬小定理}說：$p$ 是質數、$a$ 不是 $p$ 的倍數時，
$a^{p-1} \bmod p = 1$。 所以 $5^{1000000006} \bmod 1000000007$ 一定是 $1$，可以拿來檢
查。 最後一組 $m = 1$：任何數除以 $1$ 的餘數都是 $0$，所以 $e = 0$ 時要回傳
\texttt{1 \% m}，不是 \texttt{1}。

\begin{example}[label=ex:dc-power]{c}{(power\_mod.c) 快速冪取餘數}
#include <stdio.h>

// 回傳 b^e mod m。要求 1 <= m <= 3037000499，兩個小於 m 的數相乘才不會溢位。
// 不變量：每一輪開始時，r * b^e 除以 m 的餘數就是答案。*mults 是乘法次數。
long long power_mod(long long b, long long e, long long m, int *mults)
{
    long long r = 1 % m;
    b %= m;
    *mults = 0;
    while (e > 0) {
        if (e % 2 == 1) {
            r = r * b % m;
            (*mults)++;
        }
        e /= 2;
        if (e > 0) {                    // 最後一輪不必再平方
            b = b * b % m;
            (*mults)++;
        }
    }
    return r;
}

int main(void)
{
    long long t[][3] = {{7, 45, 1000}, {5, 1000000006, 1000000007},
                        {123456789, 987654321987LL, 998244353}, {10, 0, 1}};
    for (int k = 0; k < 4; k++) {
        int c;
        long long r = power_mod(t[k][0], t[k][1], t[k][2], &c);
        printf("%lld %lld %lld -> %lld (%d mults)\n", t[k][0], t[k][1], t[k][2], r, c);
    }
    return 0;
}
\end{example}
\iosample{}{7 45 1000 -> 807 (9 mults) \\ 5 1000000006 1000000007 -> 1 (44 mults) \\ 123456789 987654321987 998244353 -> 125393818 (62 mults) \\ 10 0 1 -> 0 (0 mults)}

乘法次數和公式 $\nu(e) + \lfloor \lg e \rfloor$ 一致：$987654321987$ 有 $40$ 個位
元，其中 $23$ 個是 $1$，共 $23 + 39 = 62$ 次；連乘要將近 $10^{12}$ 次。

\section{最大子陣列}
\label{sec:dc-maxsub}

回到章首的飲料攤。 把每天的盈虧依序放進 $A[1..n]$，要找的是\textbf{最大子陣列}
（maximum subarray）：一段連續、至少有一個數的 $A[i..j]$，讓
$A[i] + A[i+1] + \cdots + A[j]$ 最大。 至少要有一個數，是為了全部都賠錢的時候，
答案是賠最少的那一天，而不是「什麼都不選」的 $0$。 下面用 $S(i..j)$ 表示
$A[i..j]$ 的和。

最直接的做法是每一段都算。 固定開頭 $i$，讓結尾 $j$ 從 $i$ 往右走，邊走邊累加，
每一段只要多加一個數，總共 $n(n+1)/2$ 段，$\Theta(n^2)$。 （每一段都從頭加起是
$\Theta(n^3)$。）$n$ 到幾十萬時，$n^2$ 已經太慢。

\subsection{分治：左半、右半、跨過中點}

取中點 $q = \lfloor (p + r)/2 \rfloor$，把 $A[p..r]$ 切成 $A[p..q]$ 和
$A[q+1..r]$。 任何一段 $A[i..j]$ 一定是下面三種之一：

\begin{enumerate}
\item 整段在左半：$p \le i \le j \le q$。
\item 整段在右半：$q + 1 \le i \le j \le r$。
\item 跨過中點：$i \le q < j$，也就是同時包含 $A[q]$ 和 $A[q+1]$。
\end{enumerate}

前兩種是規模減半的同一個問題，遞迴去解。 第三種不是同一個問題（它多了「一定要包
含中點兩側」的限制），但它好解：

\begin{lemma}
\label{lem:dc-crossing}
跨過中點的段裡，和最大的是「結尾在 $q$ 的最佳段」接上「開頭在 $q+1$ 的最佳段」：
\[
\max_{p \le i \le q < j \le r} S(i..j)
= \max_{p \le i \le q} S(i..q) + \max_{q+1 \le j \le r} S(q+1..j).
\]
\end{lemma}

\begin{proof}
跨過中點的段都可以拆成 $S(i..j) = S(i..q) + S(q+1..j)$，左邊只和 $i$ 有關，右邊只
和 $j$ 有關，而 $i$、$j$ 可以各自任選（$i$ 在 $p..q$、$j$ 在 $q+1..r$）。 兩項互不
影響，所以和的最大值等於兩項各自的最大值相加。
\end{proof}

結尾在 $q$ 的最佳段，從 $A[q]$ 往左一格一格加就找得到；開頭在 $q+1$ 的最佳段，從
$A[q+1]$ 往右加。 兩邊合起來只看過 $A[p..r]$ 每一格一次，$\Theta(n)$。

\begin{pseudo}[label=alg:dc-crossing]{跨過中點的最佳段，回傳（開頭、結尾、和）}
\Function{Max-Crossing}{$A, p, q, r$}
    \State $s \gets 0$；$\mathit{left} \gets -\infty$
    \For{$i \gets q$ \DownTo $p$}
        \State $s \gets s + A[i]$
        \If{$s > \mathit{left}$}
            \State $\mathit{left} \gets s$；$i^* \gets i$
        \EndIf
    \EndFor
    \State $s \gets 0$；$\mathit{right} \gets -\infty$
    \For{$j \gets q + 1$ \To $r$}
        \State $s \gets s + A[j]$
        \If{$s > \mathit{right}$}
            \State $\mathit{right} \gets s$；$j^* \gets j$
        \EndIf
    \EndFor
    \State \Return $(i^*, j^*, \mathit{left} + \mathit{right})$
\EndFunction
\end{pseudo}

\begin{pseudo}[label=alg:dc-maxsub]{最大子陣列（分治），回傳（開頭、結尾、和）}
\Function{Max-Subarray}{$A, p, r$}
    \If{$p = r$}
        \State \Return $(p, p, A[p])$
    \EndIf
    \State $q \gets \lfloor (p + r)/2 \rfloor$
    \State $L \gets$ \Call{Max-Subarray}{$A, p, q$}
    \State $R \gets$ \Call{Max-Subarray}{$A, q + 1, r$}
    \State $C \gets$ \Call{Max-Crossing}{$A, p, q, r$}
    \State $\mathit{best} \gets L$
    \If{$C.\mathit{sum} > \mathit{best}.\mathit{sum}$}
        \State $\mathit{best} \gets C$
    \EndIf
    \If{$R.\mathit{sum} > \mathit{best}.\mathit{sum}$}
        \State $\mathit{best} \gets R$
    \EndIf
    \State \Return $\mathit{best}$
\EndFunction
\end{pseudo}

兩個程序都回傳三個值。 虛擬碼把它們當成一個物件的三個欄位：開頭、結尾、和；依第~1~章
$x.\mathit{key}$ 的寫法，第 9 行的 $C.\mathit{sum}$ 就是 $C$ 的「和」這一欄。

\textbf{正確性}：對 $n = r - p + 1$ 用強歸納法。 $n = 1$ 時唯一的一段就是
$A[p..p]$。 $n \ge 2$ 時兩半的規模都比 $n$ 小，由歸納假設 $L$、$R$ 分別是兩半的最
佳段；由引理~\ref{lem:dc-crossing}，$C$ 是跨中點的最佳段。 每一段都屬於三種之一，
所以三者中和最大的就是整個 $A[p..r]$ 的最佳段。

\textbf{執行時間}：兩次遞迴各是一半的規模，\proc{Max-Crossing} 和其他步驟
$\Theta(n)$：
\[
T(n) = 2T(n/2) + \Theta(n) = \Theta(n \lg n),
\]
和合併排序同一個遞迴式（主定理情況 2）。

\textbf{平手的規則}。 \proc{Max-Crossing} 用嚴格的 $>$，所以往左、往右都停在第一
次達到最大值的地方，也就是取最短的那一段。 \proc{Max-Subarray} 的第 9、11 行也用
$>$，三者一樣大時依序取左半、跨中點、右半。 和只有一個，但達到最大和的段可能不只
一段，題目要你寫出「哪一段」時，回答會受這些規則影響。

\begin{examnote}
最大子陣列有兩個地方各書不同：允不允許空的子陣列（允許時全部是負數的答案是
$0$），以及好幾段一樣大時回答哪一段。 作答時先照題目決定這兩件事，寫在答案開頭。
跨中點的段要寫出「往左累加」「往右累加」兩列和它們的最大值，不能只寫結果；這一段
一定同時包含 $A[q]$ 和 $A[q+1]$，往左、往右都至少取一格。
\end{examnote}

\subsection{手算：八天的盈虧}

飲料攤某八天的盈虧（單位：千元）是
\[
A = \langle 4, -6, 3, 5, -2, 7, -9, 6 \rangle .
\]
最上層 $p = 1$、$r = 8$、$q = 4$。 跨中點的部分這樣算：

\begin{center}
\begin{tabular}{l*{8}{c}}
\toprule
第幾天 & 1 & 2 & 3 & 4 & 5 & 6 & 7 & 8 \\
$A$ & 4 & $-6$ & 3 & 5 & $-2$ & 7 & $-9$ & 6 \\
\midrule
從第 4 天往左累加 & 6 & 2 & \textbf{8} & 5 & & & & \\
從第 5 天往右累加 & & & & & $-2$ & \textbf{5} & $-4$ & 2 \\
\bottomrule
\end{tabular}
\end{center}

往左的最大值是 $8$（第 3 到 4 天），往右的最大值是 $5$（第 5 到 6 天），跨中點的最
佳段是第 3 到 6 天，和 $13$。 左半 $\langle 4, -6, 3, 5 \rangle$ 遞迴算出來的最佳
段是第 3 到 4 天（$8$），右半 $\langle -2, 7, -9, 6 \rangle$ 是第 6 天（$7$），三者
取最大，答案是第 3 到 6 天，賺 $13$ 千元。 整棵遞迴樹如下，每個節點寫出它的範圍和
回傳的段，以及這一段來自三種情形的哪一種。

\begin{center}
\begin{minipage}{\linewidth}
\centering
\begin{tikzpicture}[x=1.75cm, y=1.35cm]
\foreach \v [count=\k] in {4,-6,3,5,-2,7,-9,6} {
    \node[dc-leaf] (l\k) at (\k, -3) {$\v$};
}
\node[dc-tn] (m1) at (1.5,-2) {第 1--2 天\\左半 1--1：4};
\node[dc-tn] (m2) at (3.5,-2) {第 3--4 天\\跨中點 3--4：8};
\node[dc-tn] (m3) at (5.5,-2) {第 5--6 天\\右半 6--6：7};
\node[dc-tn] (m4) at (7.5,-2) {第 7--8 天\\右半 8--8：6};
\node[dc-tn] (h1) at (2.5,-1) {第 1--4 天\\右半 3--4：8};
\node[dc-tn] (h2) at (6.5,-1) {第 5--8 天\\左半 6--6：7};
\node[dc-tnw] (t) at (4.5,0) {第 1--8 天\\跨中點 3--6：13};
\foreach \a/\b in {t/h1, t/h2, h1/m1, h1/m2, h2/m3, h2/m4, m1/l1, m1/l2, m2/l3, m2/l4, m3/l5, m3/l6, m4/l7, m4/l8} {
    \draw[dc-edge] (\a) -- (\b);
}
\end{tikzpicture}
\diagcaption{\proc{Max-Subarray} 在八天盈虧上的遞迴樹。 最下面一列是每天的盈虧，
其他節點寫「範圍、最佳段的來源、最佳段：和」。}
\label{fig:dc-maxsub-tree}
\end{minipage}
\end{center}

\subsection{線性的做法}
\label{sec:dc-kadane}

分治版每一層都重新把跨中點的和加一次，其實浪費了。 換一個角度：令
\[
E[i] = \text{結尾\textbf{剛好}在第 } i \text{ 天的段裡，最大的和}.
\]
任何一段都有一個結尾，所以答案是 $\max_{1 \le i \le n} E[i]$。 $E[i]$ 可以從
$E[i-1]$ 算出來：結尾在 $i$ 的段，要嘛只有 $A[i]$ 一天，要嘛是某個結尾在 $i - 1$
的段再接上 $A[i]$。 後一種的和是（結尾在 $i-1$ 的某段的和）$+ A[i]$，最大是
$E[i-1] + A[i]$。 所以
\[
E[1] = A[1], \qquad E[i] = A[i] + \max(E[i-1], 0) \quad (i \ge 2).
\]
白話是：前面那段如果是負的，接上它只會更差，不如從今天重新開始。 這個做法通常稱
為 \textbf{Kadane 演算法}。 每個 $E[i]$ 只用到 $E[i-1]$，所以只要一個變數
$\mathit{cur}$ 記住它。

\begin{pseudo}[label=alg:dc-kadane]{最大子陣列（線性），回傳（開頭、結尾、和）}
\Function{Max-Subarray-Linear}{$A, n$}
    \State $\mathit{cur} \gets A[1]$；$\mathit{start} \gets 1$
    \State $\mathit{best} \gets (1, 1, A[1])$
    \For{$i \gets 2$ \To $n$}
        \If{$\mathit{cur} < 0$}
            \State $\mathit{cur} \gets A[i]$；$\mathit{start} \gets i$
        \Else
            \State $\mathit{cur} \gets \mathit{cur} + A[i]$
        \EndIf
        \If{$\mathit{cur} > \mathit{best}.\mathit{sum}$}
            \State $\mathit{best} \gets (\mathit{start}, i, \mathit{cur})$
        \EndIf
    \EndFor
    \State \Return $\mathit{best}$
\EndFunction
\end{pseudo}

\textbf{不變量}：外層迴圈每一輪開始時（$i$ 是這一輪的值），$\mathit{cur} =
E[i-1] = S(\mathit{start}..i-1)$，而 $\mathit{best}$ 是 $A[1..i-1]$ 裡的最佳段。
初始時 $i = 2$，$\mathit{cur} = A[1] = E[1]$，$A[1..1]$ 只有一段，成立。 維持：第
5--8 行正是 $E[i] = A[i] + \max(E[i-1], 0)$（$E[i-1] = 0$ 時接上或不接一樣，程式選
擇接上），第 9--10 行把 $\mathit{best}$ 更新成 $A[1..i]$ 的最佳段，因為
$A[1..i]$ 裡的段不是在 $A[1..i-1]$ 裡，就是結尾在 $i$，後者最大是 $E[i]$。 結束時
$i = n + 1$，$\mathit{best}$ 是整個陣列的最佳段。 時間 $\Theta(n)$，額外空間
$\Theta(1)$。

同一組八天的數字，一天一列：

\begin{center}
\begin{tabular}{ccccl}
\toprule
$i$ & $A[i]$ & 上一輪的 $\mathit{cur}$ & 這一輪的 $\mathit{cur} = E[i]$ & $\mathit{best}$ \\
\midrule
1 & 4 & & 4 & 1--1：4 \\
2 & $-6$ & 4 & $4 - 6 = -2$ & 1--1：4 \\
3 & 3 & $-2$ & 重新開始：3 & 1--1：4 \\
4 & 5 & 3 & 8 & 3--4：8 \\
5 & $-2$ & 8 & 6 & 3--4：8 \\
6 & 7 & 6 & 13 & 3--6：13 \\
7 & $-9$ & 13 & 4 & 3--6：13 \\
8 & 6 & 4 & 10 & 3--6：13 \\
\bottomrule
\end{tabular}
\end{center}

兩種做法得到同一段。 線性做法其實是第~18~章動態規劃的第一個例子：$E[i]$ 是一個子
問題，它的答案從上一個子問題的答案直接算出來，每個子問題只算一次。 分治版慢在
「合」：每一層都從中點往兩邊重新加一次。 觀念題第~\ref{dc-q-four}~題會讓分治版每
次多回傳幾個值，把「合」變成 $\Theta(1)$，分治版也能做到 $\Theta(n)$。

\subsection{C 程式}

\begin{cbookrecap}[第~14~章 14.2、14.4、14.6 節]
\texttt{typedef struct} 定義的結構可以整個當成回傳值，也可以用 \texttt{=} 整個複
製，\texttt{best = C;} 會把三個欄位一起抄過去。 用結構 \texttt{Range} 把「開頭、結
尾、和」包在一起，就不必用三個指標參數帶回三個值。
\end{cbookrecap}

程式在八天的資料上跑分治版，每處理完一個長度大於一的範圍印一行（後序），再跑線性
版。 天數從 $1$ 開始印。 \texttt{max\_crossing} 把 $\mathit{left}$、
$\mathit{right}$ 的初值設成 \texttt{a[mid]}、\texttt{a[mid + 1]}，而不是 $-\infty$：
C 的整數沒有 $-\infty$，而這兩格本來就一定要取，效果和虛擬碼相同。 金額用
\texttt{long long}。

\begin{example}[label=ex:dc-maxsub]{c}{(max\_subarray.c) 最大子陣列的分治版與線性版}
#include <stdio.h>

typedef struct {
    int lo, hi;                         // 區段 a[lo..hi]
    long long sum;
} Range;

// 跨過中點的最佳段：一定包含 a[mid] 和 a[mid+1]。
Range max_crossing(const long long *a, int lo, int mid, int hi)
{
    Range r = {mid, mid + 1, 0};
    long long s = 0, left = a[mid], right = a[mid + 1];
    for (int i = mid; i >= lo; i--) {   // 結尾固定在 mid 的最佳段
        s += a[i];
        if (s > left) {
            left = s;
            r.lo = i;
        }
    }
    s = 0;
    for (int j = mid + 1; j <= hi; j++) { // 開頭固定在 mid+1 的最佳段
        s += a[j];
        if (s > right) {
            right = s;
            r.hi = j;
        }
    }
    r.sum = left + right;
    return r;
}

// 分治：平手時依序取左半、跨中點、右半。每處理完一段印一行（天數從 1 開始）。
Range max_subarray(const long long *a, int lo, int hi)
{
    if (lo == hi) {
        Range r = {lo, hi, a[lo]};
        return r;
    }
    int mid = lo + (hi - lo) / 2;
    Range L = max_subarray(a, lo, mid);
    Range R = max_subarray(a, mid + 1, hi);
    Range C = max_crossing(a, lo, mid, hi);
    Range best = L;
    if (C.sum > best.sum) {
        best = C;
    }
    if (R.sum > best.sum) {
        best = R;
    }
    printf("days %d-%d: left %d-%d (%lld) cross %d-%d (%lld) right %d-%d (%lld)\n",
           lo + 1, hi + 1, L.lo + 1, L.hi + 1, L.sum, C.lo + 1, C.hi + 1, C.sum,
           R.lo + 1, R.hi + 1, R.sum);
    return best;
}

// 線性做法：cur 是結尾在第 i 天的最佳和。
Range max_subarray_linear(const long long *a, int n)
{
    Range best = {0, 0, a[0]};
    long long cur = a[0];
    int start = 0;
    for (int i = 1; i < n; i++) {
        if (cur < 0) {                  // 前面那段是負的，從今天重新開始
            cur = a[i];
            start = i;
        } else {
            cur += a[i];
        }
        if (cur > best.sum) {
            best.lo = start;
            best.hi = i;
            best.sum = cur;
        }
    }
    return best;
}

int main(void)
{
    long long day[] = {4, -6, 3, 5, -2, 7, -9, 6};
    Range r1 = max_subarray(day, 0, 7);
    Range r2 = max_subarray_linear(day, 8);
    printf("divide: %d-%d (%lld)\n", r1.lo + 1, r1.hi + 1, r1.sum);
    printf("linear: %d-%d (%lld)\n", r2.lo + 1, r2.hi + 1, r2.sum);
    return 0;
}
\end{example}
\iosample{}{days 1-2: left 1-1 (4) cross 1-2 (-2) right 2-2 (-6) \\ days 3-4: left 3-3 (3) cross 3-4 (8) right 4-4 (5) \\ days 1-4: left 1-1 (4) cross 1-4 (6) right 3-4 (8) \\ days 5-6: left 5-5 (-2) cross 5-6 (5) right 6-6 (7) \\ days 7-8: left 7-7 (-9) cross 7-8 (-3) right 8-8 (6) \\ days 5-8: left 6-6 (7) cross 6-8 (4) right 8-8 (6) \\ days 1-8: left 3-4 (8) cross 3-6 (13) right 6-6 (7) \\ divide: 3-6 (13) \\ linear: 3-6 (13)}

第 3 行的 \texttt{cross 1-4 (6)} 說明跨中點的段不一定好：第 1 到 4 天跨過中點最多
只有 $6$，比右半的 $8$ 小。 輸出和圖~\ref{fig:dc-maxsub-tree} 的每個節點一致。

\section{大數乘法與矩陣乘法}
\label{sec:dc-mult}

前兩節的分治法，子問題的個數都是題目本身決定的。 這一節的兩個例子不一樣：直接切
開會得到 $4$ 個（或 $8$ 個）子問題，但用一點代數，可以少算一個。 子問題少一個，
主定理裡的 $\log_b a$ 就變小，整個演算法換了一級。

兩個 $n$ 位數相乘，小學的直式乘法把一個數的每一位乘上另一個數的每一位，共 $n^2$ 次
一位乘一位，再做 $\Theta(n^2)$ 次加法和進位，時間 $\Theta(n^2)$。 這裡的輸入大小
是位數 $n$，不是數值本身（第~\ref{sec:ana-input-size}~節）：密碼學裡的整數動輒上千位，一次乘法不
能當成一步。

設 $x$、$y$ 都是 $n$ 位數，$m = \lfloor n/2 \rfloor$，把每個數切成高位和低位：
\[
x = a \cdot 10^m + b, \qquad y = c \cdot 10^m + d,
\]
$b$、$d$ 是最後 $m$ 位，$a$、$c$ 是前面 $n - m$ 位。 乘開得到
\[
xy = ac \cdot 10^{2m} + (ad + bc) \cdot 10^m + bd .
\]
乘上 $10$ 的次方只是在後面補 $0$，加法是 $\Theta(n)$，所以真正費工的是
$ac$、$ad$、$bc$、$bd$ 四次半長的乘法：
\[
T(n) = 4T(n/2) + \Theta(n).
\]
$n^{\log_2 4} = n^2$ 比 $n$ 大了多項式倍，主定理情況 1 給出 $\Theta(n^2)$，和直式乘
法一樣。 切開本身沒有省到任何東西。

\subsection{Karatsuba：三次就夠}

關鍵在中間那一項只要 $ad + bc$ 的\textbf{和}，不需要 $ad$ 和 $bc$ 各自的值。 注意
\[
(a + b)(c + d) = ac + ad + bc + bd,
\]
所以
\[
ad + bc = (a + b)(c + d) - ac - bd .
\]
$ac$ 和 $bd$ 本來就要算，再多算一次 $(a+b)(c+d)$，就能得到中間項。 三次乘法加上
幾次加減，這是 1960 年代初 Karatsuba 提出的方法。

\begin{pseudo}[label=alg:dc-karatsuba]{Karatsuba 乘法，$x$、$y$ 是非負整數}
\Function{Karatsuba}{$x, y$}
    \If{$x < 10$ \OR $y < 10$}
        \State \Return $x \cdot y$ \Comment{有一個是一位數，直接乘}
    \EndIf
    \State $n \gets$ $x$、$y$ 中比較長的位數；$m \gets \lfloor n/2 \rfloor$
    \State $a \gets \lfloor x / 10^m \rfloor$；$b \gets x \bmod 10^m$
    \State $c \gets \lfloor y / 10^m \rfloor$；$d \gets y \bmod 10^m$
    \State $z_2 \gets$ \Call{Karatsuba}{$a, c$}
    \State $z_0 \gets$ \Call{Karatsuba}{$b, d$}
    \State $z_1 \gets$ \Call{Karatsuba}{$a + b, c + d$} $- z_2 - z_0$
    \State \Return $z_2 \cdot 10^{2m} + z_1 \cdot 10^m + z_0$
\EndFunction
\end{pseudo}

第 2 行處理一位數乘任何數，那是 $\Theta(n)$ 的直式乘法，不再遞迴。 遞迴一定會停：
$n \ge 2$ 時 $m \ge 1$，而 $x - (a + b) = a(10^m - 1)$、$y - (c + d) = c(10^m -
1)$，$a$、$c$ 至少有一個不是 $0$，所以每一次呼叫的 $x + y$ 都比上一層小。

\textbf{執行時間}。 $a$、$c$ 最多 $n - m = \lceil n/2 \rceil$ 位，$b$、$d$ 最多
$m$ 位，$a + b$、$c + d$ 可能進位多一位，最多 $\lceil n/2 \rceil + 1$ 位。 第 5、
6、9、10 行的拆開、移位、加減都是 $\Theta(n)$。 所以
\[
T(n) \le 3T(\lceil n/2 \rceil + 1) + \Theta(n).
\]
多出來的 $+1$ 用第~2~章的換變數處理：令 $S(n) = T(n + 2)$，因為
$\lceil (n+2)/2 \rceil + 1 = \lceil n/2 \rceil + 2$，得到
$S(n) \le 3S(\lceil n/2 \rceil) + \Theta(n)$，這正是第~2~章圖~\ref{fig:rec-tree-3half} 算過的
$3T(n/2) + n$，解是 $\Theta(n^{\lg 3})$。 所以 $T(n) = S(n-2) = O(n^{\lg 3})$，
$\lg 3 \approx 1.585$。 $n = 1000$ 時 $n^2 = 10^6$，$n^{\lg 3}$ 只有大約
$5.7 \times 10^4$。

\subsection{手算：\texorpdfstring{$4726 \times 1853$}{4726 × 1853}}

$n = 4$，$m = 2$，$a = 47$、$b = 26$、$c = 18$、$d = 53$。

\begin{center}
\begin{tabular}{lll}
\toprule
乘積 & 算式 & 值 \\
\midrule
$z_2 = ac$ & $47 \times 18$ & $846$ \\
$z_0 = bd$ & $26 \times 53$ & $1378$ \\
$(a + b)(c + d)$ & $73 \times 71$ & $5183$ \\
$z_1$ & $5183 - 846 - 1378$ & $2959$ \\
\midrule
$xy$ & $846 \times 10^4 + 2959 \times 10^2 + 1378$ & $8757278$ \\
\bottomrule
\end{tabular}
\end{center}

$8460000 + 295900 + 1378 = 8757278$，可以用直式乘法驗算。 三個兩位數的乘法再各自
用同樣的方法拆一次，例如 $73 \times 71$：$a = 7$、$b = 3$、$c = 7$、$d = 1$，
$ac = 49$、$bd = 3$、$(a+b)(c+d) = 10 \times 8 = 80$，中間項 $80 - 49 - 3 = 28$，
$49 \times 100 + 28 \times 10 + 3 = 5183$。 這裡 $a + b = 10$ 進位成兩位數，但另一
個因數 $8$ 只有一位，第 2 行直接乘。

遞迴樹每往下一層，節點數乘以 $3$、規模減半，這個例子最下層是 $9$ 次一位數的乘法；$n = 2^k$ 時最下層有 $3^k = n^{\lg 3}$ 個節點，直式
乘法要 $4^k = n^2$ 次。 $n = 4$ 時是 $9$ 比 $16$，差距不大；位數越多差得越多，而且
每一層還有加減和進位的 $\Theta(n)$ 成本，所以實際的程式在位數少時改用直式乘法，只
有夠長才遞迴。

\subsection{C 程式}

\proc{Karatsuba} 的虛擬碼可以直接翻成 C，只要乘積放得進 \texttt{long long}（$18$ 位
以內）。 真正的大數程式要把數字存成位數的陣列，拆開、相加、移位都在陣列上做，結構
和下面完全相同。 程式先印出 $4726 \times 1853$ 每一次（不是一位數的）呼叫的三個乘
積，順序是遞迴的後序；再對幾組數驗算，並數最下層直接乘的次數。

\begin{example}[label=ex:dc-karatsuba]{c}{(karatsuba.c) Karatsuba 乘法}
#include <stdio.h>

long long leaf_mults;                   // 遞迴停下來時直接做的乘法次數
int trace;                              // 1：每次呼叫做完印一行

int digits(long long v)                 // v 的十進位位數（0 算一位）
{
    int d = 1;
    while (v >= 10) {
        v /= 10;
        d++;
    }
    return d;
}

// x、y 是非負整數，乘積不超過 long long。
long long karatsuba(long long x, long long y)
{
    if (x < 10 || y < 10) {             // 有一個是一位數：直接乘
        leaf_mults++;
        return x * y;
    }
    int n = digits(x) > digits(y) ? digits(x) : digits(y);
    long long p = 1;                    // p = 10^m，m = n / 2
    for (int i = 0; i < n / 2; i++) {
        p *= 10;
    }
    long long a = x / p, b = x % p, c = y / p, d = y % p;
    long long z2 = karatsuba(a, c);
    long long z0 = karatsuba(b, d);
    long long z1 = karatsuba(a + b, c + d) - z2 - z0;
    long long r = z2 * p * p + z1 * p + z0;
    if (trace) {
        printf("%lld * %lld: z2=%lld z1=%lld z0=%lld -> %lld\n", x, y, z2, z1, z0, r);
    }
    return r;
}

int main(void)
{
    trace = 1;
    long long r = karatsuba(4726, 1853);
    printf("leaf mults: %lld (school: 16)\n", leaf_mults);
    trace = 0;
    long long xs[][2] = {{4726, 1853}, {31415926, 27182818}, {99999999, 99999999},
                         {123456789, 987654321}};
    for (int k = 0; k < 4; k++) {
        leaf_mults = 0;
        r = karatsuba(xs[k][0], xs[k][1]);
        printf("%lld * %lld = %lld %s, leaf mults %lld\n", xs[k][0], xs[k][1], r,
               r == xs[k][0] * xs[k][1] ? "ok" : "WRONG", leaf_mults);
    }
    return 0;
}
\end{example}
\iosample{}{47 * 18: z2=4 z1=39 z0=56 -> 846 \\ 26 * 53: z2=10 z1=36 z0=18 -> 1378 \\ 73 * 71: z2=49 z1=28 z0=3 -> 5183 \\ 4726 * 1853: z2=846 z1=2959 z0=1378 -> 8757278 \\ leaf mults: 9 (school: 16) \\ 4726 * 1853 = 8757278 ok, leaf mults 9 \\ 31415926 * 27182818 = 853973398759468 ok, leaf mults 33 \\ 99999999 * 99999999 = 9999999800000001 ok, leaf mults 63 \\ 123456789 * 987654321 = 121932631112635269 ok, leaf mults 45}

前四行是遞迴樹的四個內部節點。 後面三組是八、九位數，最下層
各做 $33$、$63$、$45$ 次乘法，同樣位數的直式乘法要 $64$、$64$、$81$ 次。
$99999999 \times 99999999$ 幾乎沒省到：每一層的 $a + b$ 都進位多一位（$9999 +
9999 = 19998$），子問題比一半大，最下層的乘法也變多。 這就是分析裡那個「$+1$」，漸
近上不影響，常數上卻看得到。 實際的程式會把基礎情況放大（例如幾十位以下就用直式乘
法），並用更大的進位制（例如每格存 $10^9$），減少遞迴的層數。


\subsection{Strassen 矩陣乘法}
\label{sec:dc-strassen}

兩個 $n \times n$ 矩陣相乘，$C = AB$，定義是
$c_{ij} = \sum_{k=1}^{n} a_{ik} b_{kj}$。 $n^2$ 個元素各要 $n$ 次乘法，共 $n^3$ 次，
三層迴圈 $\Theta(n^3)$。

設 $n$ 是 $2$ 的次方，把每個矩陣切成四塊 $\frac{n}{2} \times \frac{n}{2}$ 的子矩
陣：
\[
\begin{pmatrix} C_{11} & C_{12} \\ C_{21} & C_{22} \end{pmatrix}
= \begin{pmatrix} A_{11} & A_{12} \\ A_{21} & A_{22} \end{pmatrix}
\begin{pmatrix} B_{11} & B_{12} \\ B_{21} & B_{22} \end{pmatrix},
\]
\[
\begin{aligned}
C_{11} &= A_{11} B_{11} + A_{12} B_{21}, & C_{12} &= A_{11} B_{12} + A_{12} B_{22}, \\
C_{21} &= A_{21} B_{11} + A_{22} B_{21}, & C_{22} &= A_{21} B_{12} + A_{22} B_{22}.
\end{aligned}
\]
分塊的乘法規則和 $2 \times 2$ 的數字矩陣一模一樣，只是每個元素換成子矩陣。 這要 $8$
次半大的矩陣乘法和 $4$ 次矩陣加法（每次 $\Theta(n^2)$）：
\[
T(n) = 8T(n/2) + \Theta(n^2) = \Theta(n^3),
\]
主定理情況 1（$n^{\log_2 8} = n^3$）。 一樣沒省到。

1969 年 Strassen 發現 $7$ 次乘法就夠。 先算七個乘積：
\[
\begin{aligned}
M_1 &= (A_{11} + A_{22})(B_{11} + B_{22}), &
M_2 &= (A_{21} + A_{22})\,B_{11}, \\
M_3 &= A_{11}\,(B_{12} - B_{22}), &
M_4 &= A_{22}\,(B_{21} - B_{11}), \\
M_5 &= (A_{11} + A_{12})\,B_{22}, &
M_6 &= (A_{21} - A_{11})(B_{11} + B_{12}), \\
M_7 &= (A_{12} - A_{22})(B_{21} + B_{22}),
\end{aligned}
\]
再組合：
\[
C_{11} = M_1 + M_4 - M_5 + M_7, \quad
C_{12} = M_3 + M_5, \quad
C_{21} = M_2 + M_4, \quad
C_{22} = M_1 - M_2 + M_3 + M_6 .
\]
算 $M_1$ 到 $M_7$ 的括號裡有 $10$ 次子矩陣加減，組合 $C$ 又有 $8$ 次，共 $18$ 次，
每次 $\Theta(n^2)$。 所以
\[
T(n) = 7T(n/2) + \Theta(n^2).
\]
$\log_2 7 \approx 2.807$，$n^2 = O(n^{\log_2 7 - 0.8})$，主定理情況 1：
$T(n) = \Theta(n^{\lg 7}) \approx \Theta(n^{2.81})$。 加法從 $4$ 次變成 $18$ 次，
只影響 $f(n)$ 的常數；少掉的那一次乘法卻讓遞迴樹每一層少分出一個分支，改變了葉節
點的個數 $7^{\lg n} = n^{\lg 7}$。

\begin{whybox}[\texorpdfstring{$C_{11} = M_1 + M_4 - M_5 + M_7$}{C11 = M1 + M4 - M5 + M7} 的驗證]
把四個乘積展開（子矩陣的乘法不能交換，但有分配律，展開時保持左右順序）：
\[
\begin{aligned}
M_1 &= A_{11}B_{11} + A_{11}B_{22} + A_{22}B_{11} + A_{22}B_{22}, \\
M_4 &= A_{22}B_{21} - A_{22}B_{11}, \\
-M_5 &= -A_{11}B_{22} - A_{12}B_{22}, \\
M_7 &= A_{12}B_{21} + A_{12}B_{22} - A_{22}B_{21} - A_{22}B_{22}.
\end{aligned}
\]
相加時 $A_{11}B_{22}$、$A_{22}B_{11}$、$A_{22}B_{22}$、$A_{22}B_{21}$、$A_{12}B_{22}$
都一正一負消掉，只剩 $A_{11}B_{11} + A_{12}B_{21}$，正是 $C_{11}$。 其他三塊用同樣
的方法展開驗證，例如 $M_3 + M_5 = A_{11}B_{12} - A_{11}B_{22} + A_{11}B_{22} +
A_{12}B_{22} = C_{12}$。 整個推導只用到加法、減法和分配律，所以子矩陣的元素是整
數、實數或任何可以相減的數都適用。
\end{whybox}

\textbf{手算}。 考卷上常見的是 $2 \times 2$ 的數字矩陣，每一塊都是一個數。 取
\[
A = \begin{pmatrix} 2 & 5 \\ 3 & -1 \end{pmatrix}, \qquad
B = \begin{pmatrix} 4 & 1 \\ -2 & 6 \end{pmatrix}.
\]

\begin{center}
\begin{tabular}{lll}
\toprule
乘積 & 代入 & 值 \\
\midrule
$M_1 = (a_{11} + a_{22})(b_{11} + b_{22})$ & $(2 + (-1))(4 + 6) = 1 \times 10$ & $10$ \\
$M_2 = (a_{21} + a_{22})\,b_{11}$ & $(3 + (-1)) \times 4$ & $8$ \\
$M_3 = a_{11}\,(b_{12} - b_{22})$ & $2 \times (1 - 6)$ & $-10$ \\
$M_4 = a_{22}\,(b_{21} - b_{11})$ & $(-1) \times (-2 - 4)$ & $6$ \\
$M_5 = (a_{11} + a_{12})\,b_{22}$ & $(2 + 5) \times 6$ & $42$ \\
$M_6 = (a_{21} - a_{11})(b_{11} + b_{12})$ & $(3 - 2)(4 + 1)$ & $5$ \\
$M_7 = (a_{12} - a_{22})(b_{21} + b_{22})$ & $(5 - (-1))(-2 + 6) = 6 \times 4$ & $24$ \\
\bottomrule
\end{tabular}
\end{center}

組合起來：
\[
\begin{aligned}
c_{11} &= M_1 + M_4 - M_5 + M_7 = 10 + 6 - 42 + 24 = -2, &
c_{12} &= M_3 + M_5 = -10 + 42 = 32, \\
c_{21} &= M_2 + M_4 = 8 + 6 = 14, &
c_{22} &= M_1 - M_2 + M_3 + M_6 = -3.
\end{aligned}
\]
其中 $c_{22} = 10 - 8 - 10 + 5$。 直接乘驗算：$c_{11} = 2 \cdot 4 + 5 \cdot (-2) = -2$，$c_{12} = 2 \cdot 1 + 5 \cdot 6 =
32$，$c_{21} = 3 \cdot 4 + (-1)(-2) = 14$，$c_{22} = 3 \cdot 1 + (-1) \cdot 6 = -3$，
全部相符。

\begin{examnote}
Strassen 的七個乘積在不同課本裡的編號和寫法不一樣，有的書先定義十個和差
$S_1, \ldots, S_{10}$ 再定義乘積 $P_1, \ldots, P_7$，組合的式子也跟著不同。 題目
給了式子就照題目的式子算；題目沒給，先把你用的七個定義和四個組合式寫出來，再代數
字。 每個乘積寫成「式子、代入、值」三段，最後一定要用一般的乘法驗算 $C$，代錯一
個正負號，驗算馬上看得出來。
\end{examnote}

\textbf{實務上的幾點}。 $n$ 不是 $2$ 的次方時補零補到 $2$ 的次方，時間只差常數倍。
Strassen 的常數比三層迴圈大，浮點數的捨入誤差也比較大，所以實際的程式在子矩陣夠小
時改回一般的乘法。 它還需要「減法」：第~25~章把最短路徑寫成類似矩陣乘法的運算時，
加法換成取最小值，沒有對應的減法，Strassen 就用不上。

\section{最近點對}
\label{sec:dc-closest}

平面上有 $n$ 個點，要找出距離最近的兩個點，距離是一般的直線距離
$\mathit{dist}(p, q) = \sqrt{(p.x - q.x)^2 + (p.y - q.y)^2}$。 兩兩比較要算
$\binom{n}{2}$ 次距離，$\Theta(n^2)$。 這一節用分治做到 $\Theta(n \lg n)$。

一維的情形很簡單：點都在數線上時，排好之後最近的兩點一定相鄰，掃一次相鄰的差就
好。 用分治的眼光看，把排好的點切成左右兩半，跨兩半的點對裡最近的一定是左半最右
邊和右半最左邊那一對。 二維沒有這種「相鄰」，依 $x$ 排好之後，最近的兩點在 $x$ 的順
序上可以相隔很遠。 二維的分治法一樣是切兩半、遞迴、再處理跨兩半的點對，難在最後
一步怎麼不必把左右所有的點兩兩比較。 平面上其他常考的問題（轉向、線段相交、凸包、
最遠點對）在第~\ref{sec:dc-geom}~節。

\subsection{演算法}

先把所有點依 $x$ 排序（$x$ 相同時依 $y$），放在 $P[1..n]$。 遞迴程序
\proc{Closest}$(P, p, r)$ 處理 $P[p..r]$ 這 $n = r - p + 1$ 個點，回傳它們之間的
最近距離 $\delta$。 為了不在每一層重新排序，它還做一件事：進來時 $P[p..r]$ 依
$x$ 排好，回去時改成依 $y$ 排好（$y$ 相同時依 $x$）。 這和合併排序一樣，兩半各自
排好以後合併就好。

\begin{pseudo}[label=alg:dc-closest]{最近點對，$P[p..r]$ 進來時依 $x$ 排好，回去時依 $y$ 排好}
\Function{Closest}{$P, p, r$}
    \If{$r - p + 1 \le 3$}
        \State $\delta \gets P[p..r]$ 兩兩比較的最小距離
        \State 把 $P[p..r]$ 依 $y$ 排好
        \State \Return $\delta$
    \EndIf
    \State $q \gets \lfloor (p + r)/2 \rfloor$
    \State $x_m \gets P[q].x$ \Comment{遞迴會重排 $P$，先記下分界}
    \State $\delta \gets \min($\Call{Closest}{$P, p, q$}$,\ $\Call{Closest}{$P, q + 1, r$}$)$
    \State 依 $y$ 合併 $P[p..q]$ 和 $P[q+1..r]$ \Comment{虛擬碼~\ref{alg:so-merge}，鍵換成 $y$}
    \State $S \gets P[p..r]$ 裡 $|x - x_m| < \delta$ 的點，保持依 $y$ 的順序；$k \gets |S|$
    \For{$i \gets 1$ \To $k - 1$}
        \State $j \gets i + 1$
        \While{$j \le k$ \AND $S[j].y - S[i].y < \delta$}
            \State $\delta \gets \min(\delta, \mathit{dist}(S[i], S[j]))$
            \State $j \gets j + 1$
        \EndWhile
    \EndFor
    \State \Return $\delta$
\EndFunction
\end{pseudo}

第 2--5 行處理三個點以內的範圍，所以遞迴永遠不會碰到只有一個點、沒有距離可言的範
圍。 其他幾行的重點：

\begin{itemize}
\item 第 6--8 行：以第 $q$ 個點的 $x$ 座標當分界。 左半 $L = P[p..q]$ 的點都滿足
$x \le x_m$，右半 $R = P[q+1..r]$ 的點都滿足 $x \ge x_m$（$x$ 相同的點可能分在兩
邊，這不影響後面的論證）。 $x_m$ 一定要在遞迴\textbf{之前}讀：遞迴回來以後
$P[p..q]$ 已經依 $y$ 重排，$P[q]$ 不再是左半 $x$ 最大的點。
\item 第 10 行：第 8 行之後，$\delta$ 是兩點同在一半的最近距離。 如果 $u \in L$、$v \in R$ 而且 $\mathit{dist}(u, v) < \delta$，那麼
$x_m - u.x \le v.x - u.x \le \mathit{dist}(u, v) < \delta$，同理 $v.x - x_m < \delta$，
所以兩點都在以分界線為中心、寬 $2\delta$ 的\textbf{帶狀區域}裡。 帶狀區域以外的點
不可能參與更近的點對。
\item 第 11--15 行：帶狀區域裡的點依 $y$ 排好了。 每個點往後看，直到 $y$ 差不小於
$\delta$ 為止，因為再往後的點 $y$ 差只會更大，距離不可能小於 $\delta$。
\end{itemize}

第 13 行的迴圈看起來可能跑很多次，下面的引理說它最多跑 $7$ 次。

\subsection{帶狀區域裡最多看 7 個點}

\begin{lemma}
\label{lem:dc-seven}
在 $n \ge 4$ 的一次呼叫裡，令 $\delta_0$ 是第 8 行算出的 $\delta$。 對每一個 $i$，
第 13 行的 \textbf{while} 迴圈主體最多執行 $7$ 次。
\end{lemma}

\begin{proof}
$\delta_0 = 0$ 時沒有點滿足 $|x - x_m| < 0$，$S$ 是空的，沒有什麼要證。 以下設
$\delta_0 > 0$。

固定 $i$，令 $y_0 = S[i].y$。 迴圈主體執行時的 $j$ 都滿足 $S[j].y - y_0 < \delta
\le \delta_0$（$\delta$ 在迴圈裡只會變小）。 $S$ 依 $y$ 排好，所以
$S[i], S[i+1], \ldots, S[j]$ 的 $y$ 座標都在 $[y_0, y_0 + \delta_0]$ 裡；它們又在帶
狀區域裡，$x$ 座標在 $(x_m - \delta_0, x_m + \delta_0)$。 因此這些點全都落在長方形
\[
\mathcal{R} = [x_m - \delta_0,\ x_m + \delta_0] \times [y_0,\ y_0 + \delta_0]
\]
裡。 只要證明 $\mathcal{R}$ 裡最多有 $8$ 個 $S$ 的點，就有 $j - i \le 7$，迴圈主體最
多 $7$ 次。

把 $\mathcal{R}$ 的左半 $[x_m - \delta_0, x_m] \times [y_0, y_0 + \delta_0]$ 用四個
邊長 $\delta_0/2$ 的\textbf{閉}正方形蓋滿（圖~\ref{fig:dc-cp-box}）。 $\mathcal{R}$ 裡
屬於 $L$ 的點都滿足 $x \le x_m$，所以都在左半，也就都在某一個正方形裡（落在邊上的
點可能同時屬於兩個正方形，沒有關係）。 同一個閉正方形裡的兩點，距離不超過對角線
$\delta_0 / \sqrt{2} < \delta_0$；但 $L$ 裡任兩點的距離都至少是 $L$ 的最近距離，而
那個距離不小於 $\delta_0 = \min(\delta_L, \delta_R)$。 所以每個正方形最多一個 $L$
的點，$\mathcal{R}$ 裡最多 $4$ 個 $L$ 的點。 右半同理，最多 $4$ 個 $R$ 的點。 合計最
多 $8$ 個，包括 $S[i]$ 自己。
\end{proof}

\begin{center}
\begin{minipage}{\linewidth}
\centering
\begin{tikzpicture}[scale=1]
\def\d{2.6}
\fill[coreFill] (-\d,0) rectangle (\d,\d);
\draw[black!45] (-\d/2,0) -- (-\d/2,\d) (\d/2,0) -- (\d/2,\d) (-\d,\d/2) -- (\d,\d/2);
\draw[thick, black!70] (-\d,0) rectangle (\d,\d);
\draw[very thick, dashed, memStroke] (0,-0.5) -- (0,\d+0.5) node[above, font=\small] {$x = x_m$};
\node[dc-pt, label={[font=\scriptsize]below:$S[i]$}] at (-0.75,0) {};
\node[dc-pt] at (-2.3,0.95) {};
\node[dc-pt] at (-0.25,1.7) {};
\node[dc-pt] at (-1.5,2.45) {};
\node[dc-pt] at (0.35,0.5) {};
\node[dc-pt] at (2.2,0.2) {};
\node[dc-pt] at (1.0,2.05) {};
\node[dc-pt] at (2.45,1.95) {};
\draw[|<->|, black!60] (-\d,-0.75) -- (0,-0.75) node[midway, below, font=\scriptsize] {$\delta_0$};
\draw[|<->|, black!60] (0,-0.75) -- (\d,-0.75) node[midway, below, font=\scriptsize] {$\delta_0$};
\draw[|<->|, black!60] (\d+0.3,0) -- (\d+0.3,\d) node[midway, right, font=\scriptsize] {$\delta_0$};
\draw[|<->|, black!60] (-\d-0.3,0) -- (-\d-0.3,\d/2) node[midway, left, font=\scriptsize] {$\delta_0/2$};
\node[font=\small] at (-\d/2, \d+0.3) {左半 $L$};
\node[font=\small] at (\d/2, \d+0.3) {右半 $R$};
\end{tikzpicture}
\diagcaption{引理~\ref{lem:dc-seven} 的長方形 $\mathcal{R}$：寬 $2\delta_0$、高
$\delta_0$，底邊通過 $S[i]$。 左右各切成四個邊長 $\delta_0/2$ 的正方形，同一邊的兩
點距離至少 $\delta_0$，每個正方形最多一個點，所以 $\mathcal{R}$ 裡最多 $8$ 個點。}
\label{fig:dc-cp-box}
\end{minipage}
\end{center}

\begin{theorem}
\label{thm:dc-closest}
$P[p..r]$ 至少有兩個點、依 $x$ 排好時，\proc{Closest}$(P, p, r)$ 回傳 $P[p..r]$ 的最
近距離，並把 $P[p..r]$ 改成依 $y$ 排好。 先排序再呼叫 \proc{Closest}$(P, 1, n)$，
總時間 $\Theta(n \lg n)$。
\end{theorem}

\begin{proof}
對 $n = r - p + 1$ 用強歸納法。 $n = 2, 3$ 時第 2--5 行直接比完所有點對，也排好了
$y$。

$n \ge 4$ 時，兩半各有 $\lceil n/2 \rceil \ge 2$ 和 $\lfloor n/2 \rfloor \ge 2$ 個
點，由歸納假設，兩次遞迴分別回傳兩半的最近距離 $\delta_L$、$\delta_R$，並把兩半各
自依 $y$ 排好，第 9 行的合併因此得到依 $y$ 排好的 $P[p..r]$。 設 $\delta^*$ 是
$P[p..r]$ 真正的最近距離，$\delta_0 = \min(\delta_L, \delta_R) \ge \delta^*$。 第
14 行只會把 $\delta$ 換成某兩點的實際距離，所以 $\delta$ 永遠 $\ge \delta^*$；只
要證明程序結束時 $\delta \le \delta^*$。

如果最近的一對在同一半，$\delta_0 = \delta^*$，成立。 否則最近的一對是 $u \in L$、
$v \in R$，$\mathit{dist}(u, v) = \delta^* < \delta_0$。 由第 10 行的說明，$u$、$v$
都在 $S$ 裡；設它們在 $S$ 裡的位置是 $a < b$。 考慮外層迴圈 $i = a$ 那一輪。 如果
這一輪開始時 $\delta$ 已經等於 $\delta^*$，就不必再證。 否則 $\delta > \delta^*$，
而對 $a < j \le b$，$S[j].y - S[a].y \le S[b].y - S[a].y \le \mathit{dist}(u, v) =
\delta^* < \delta$，所以第 13 行的條件一路成立到 $j = b$（中途 $\delta$ 如果已經降到 $\delta^*$，也已
經證完），第 14 行算到 $\mathit{dist}(S[a], S[b]) = \delta^*$，$\delta$ 變成
$\delta^*$。

時間：第 9、10 行 $\Theta(n)$；由引理~\ref{lem:dc-seven}，第 11--15 行最多算
$7k \le 7n$ 次距離，也是 $O(n)$。 所以 $T(n) = 2T(n/2) + \Theta(n) = \Theta(n \lg n)$，
加上一開始排序的 $\Theta(n \lg n)$，合計 $\Theta(n \lg n)$。
\end{proof}

如果不讓遞迴順便依 $y$ 排好，改成每一層重新排序帶狀區域，時間會多一個 $\lg n$ 的
因子（觀念題第~\ref{dc-q-tf}~題 (e)）。

\subsection{手算：十個點}

十個點依 $x$ 排好，名字依序是 A 到 J：
\[
\begin{gathered}
A(2, 10),\ B(4, 6),\ C(5, 11),\ D(9, 14),\ E(11, 8), \\
F(12, 5),\ G(13, 9),\ H(14, 14),\ I(16, 1),\ J(18, 5).
\end{gathered}
\]
座標都是整數，所以下面一律寫距離的平方，比較大小時不必開根號：$\delta^2$ 和
$(x \text{ 差})^2$、$(y \text{ 差})^2$ 比。 帶狀區域的條件 $|x - x_m| < \delta$ 也
寫成 $(x - x_m)^2 < \delta^2$。

\begin{center}
\begin{minipage}{\linewidth}
\centering
\begin{tikzpicture}[x=0.62cm, y=0.48cm]
\fill[coreFill] (7.838,0) rectangle (14.162,15.3);
\draw[-{Stealth[length=2mm]}, black!60] (0,0) -- (19.6,0) node[right, font=\small] {$x$};
\draw[-{Stealth[length=2mm]}, black!60] (0,0) -- (0,15.6) node[above, font=\small] {$y$};
\foreach \t in {5,10,15} { \draw[black!60] (\t,0) -- (\t,-0.25) node[below, font=\scriptsize] {\t}; }
\foreach \t in {5,10,15} { \draw[black!60] (0,\t) -- (-0.2,\t) node[left, font=\scriptsize] {\t}; }
\draw[very thick, dashed, memStroke] (11,0) -- (11,15.3) node[above, font=\scriptsize] {$x_m = 11$};
\draw[dashed, black!45] (2,10) -- (5,11);
\draw[dashed, black!45] (12,5) -- (13,9);
\draw[very thick, optStroke] (11,8) -- (13,9);
\foreach \n/\x/\y/\pos in {A/2/10/left, B/4/6/left, C/5/11/above, D/9/14/above, E/11/8/left, F/12/5/below, G/13/9/right, H/14/14/above, I/16/1/right, J/18/5/right} {
    \node[dc-pt, label={[dc-ptlab]\pos:\n}] at (\x,\y) {};
}
\end{tikzpicture}
\diagcaption{十個點。 藍底是最上層的帶狀區域 $|x - 11| < \sqrt{10}$；虛線是左右兩
半各自的最近點對（AC、FG），紅色粗線是答案 EG。}
\label{fig:dc-cp-points}
\end{minipage}
\end{center}

遞迴依後序展開，每一列是一次呼叫做完時的結果。 帶狀區域裡的點依 $y$ 排列；「比
較」一欄列出第 14 行實際算過距離的點對和它們的距離平方。

\begin{center}
\small
\begin{tabular}{lccc>{\raggedright\arraybackslash}p{0.22\linewidth}>{\raggedright\arraybackslash}p{0.27\linewidth}c}
\toprule
範圍 & $x_m$ & $\delta_L^2$ & $\delta_R^2$ & 帶狀區域（依 $y$） & 比較 & 結果 $\delta^2$ \\
\midrule
A--C & & & & & AB 20、AC 10、BC 26 & 10 \\
D--E & & & & & DE 40 & 40 \\
A--E & 5 & 10 & 40 & B, A, C & AC 10 & 10 \\
F--H & & & & & FG 17、FH 85、GH 26 & 17 \\
I--J & & & & & IJ 20 & 20 \\
F--J & 14 & 17 & 20 & I, F, J, G, H & IF 32、IJ 20、FJ 36、FG 17、JG 41 & 17 \\
A--J & 11 & 10 & 17 & F, E, G, D, H & FE 10、EG 5、DH 25 & 5 \\
\bottomrule
\end{tabular}
\end{center}

最上層的帶狀區域依 $y$ 排是 F(5)、E(8)、G(9)、D(14)、H(14)。 FE 的距離平方剛好是
$10$，不比 $\delta^2$ 小，不更新；EG 是 $4 + 1 = 5$，$\delta^2$ 變成 $5$；接著 G 往後
看 D，$y$ 差 $5$，$25 \ge 5$，停；D 和 H 的 $y$ 差是 $0$，算一次 DH。 F 和 J 的 $y$
都是 $5$，依 $x$ 排，F 在前。

答案是 E、G 兩點，距離 $\sqrt{5} \approx 2.24$。 它們分在分界線的兩邊，只看左半
或只看右半都找不到。 整個過程算了 $17$ 次距離，兩兩比較要 $45$ 次。

\subsection{C 程式}

\begin{cbookrecap}[第~18~章 18.6 節、第~19~章 19.3 節]
\texttt{qsort(p, n, sizeof *p, cmp\_x)} 用比較函式 \texttt{cmp\_x} 排序 \texttt{p} 指
到的 \texttt{n} 個元素。 比較函式收到兩個 \texttt{const void *}，先轉回
\texttt{const Point *} 再比；回傳負數、$0$、正數分別表示第一個應該排在前面、一樣、排
在後面。 不要寫成 \texttt{p->x - q->x}，兩個很大的整數相減會溢位，本書第~\ref{sec:so-compare}~節
有例子；這裡用 \texttt{(a > b) - (a < b)}。 暫存陣列 \texttt{tmp} 用
\texttt{malloc} 一次要 $n$ 格，整個遞迴共用，最後 \texttt{free}。
\end{cbookrecap}

程式全部用整數算距離的平方，不用 \texttt{sqrt}，沒有浮點誤差；座標的差先轉成
\texttt{long long} 再平方，不會溢位。 \texttt{best} 的初值 \texttt{LLONG\_MAX} 代表
「還沒找到」，只拿來比大小，不做加減。 點故意以打亂的順序給，先用 \texttt{qsort}
依 $x$ 排好；三個點以內的範圍也用 \texttt{qsort} 依 $y$ 排。 每處理完一個不是基礎
情況的範圍，印出 $x_m$、兩半的結果、帶狀區域的點數和這一層的答案。

\begin{example}[label=ex:dc-closest]{c}{(closest\_pair.c) 平面最近點對}
#include <stdio.h>
#include <stdlib.h>
#include <string.h>
#include <limits.h>

typedef struct {
    int x, y;
} Point;

int dist_calls;                         // 算了幾次距離

long long dist2(Point p, Point q)       // 距離的平方，全部用整數算
{
    long long dx = (long long)p.x - q.x, dy = (long long)p.y - q.y;
    dist_calls++;
    return dx * dx + dy * dy;
}

int cmp_x(const void *pa, const void *pb)   // 先比 x 再比 y，不用相減
{
    const Point *p = pa, *q = pb;
    if (p->x != q->x) {
        return (p->x > q->x) - (p->x < q->x);
    }
    return (p->y > q->y) - (p->y < q->y);
}

int cmp_y(const void *pa, const void *pb)   // 先比 y 再比 x
{
    const Point *p = pa, *q = pb;
    if (p->y != q->y) {
        return (p->y > q->y) - (p->y < q->y);
    }
    return (p->x > q->x) - (p->x < q->x);
}

// p[lo..hi] 進來時依 x 排好，出去時依 y 排好；回傳最近距離的平方。
long long closest(Point *p, Point *tmp, int lo, int hi)
{
    long long best = LLONG_MAX;
    if (hi - lo + 1 <= 3) {             // 三個點以內：兩兩比較，再依 y 排
        for (int i = lo; i <= hi; i++) {
            for (int j = i + 1; j <= hi; j++) {
                long long d = dist2(p[i], p[j]);
                best = d < best ? d : best;
            }
        }
        qsort(p + lo, hi - lo + 1, sizeof *p, cmp_y);
        return best;
    }
    int mid = lo + (hi - lo) / 2;
    int xm = p[mid].x;                  // 遞迴會依 y 重排 p，先記下分界
    long long dl = closest(p, tmp, lo, mid);
    long long dr = closest(p, tmp, mid + 1, hi);
    best = dl < dr ? dl : dr;
    memcpy(tmp + lo, p + lo, (hi - lo + 1) * sizeof *p);   // 依 y 合併兩半
    for (int k = lo, i = lo, j = mid + 1; k <= hi; k++) {
        if (i <= mid && (j > hi || cmp_y(&tmp[i], &tmp[j]) <= 0)) {
            p[k] = tmp[i++];
        } else {
            p[k] = tmp[j++];
        }
    }
    int k = 0;                          // 帶狀區域依 y 的順序放進 s[0..k-1]
    Point *s = tmp + lo;
    for (int t = lo; t <= hi; t++) {
        long long dx = (long long)p[t].x - xm;
        if (dx * dx < best) {
            s[k++] = p[t];
        }
    }
    for (int a = 0; a < k; a++) {
        for (int b = a + 1; b < k; b++) {
            long long dy = (long long)s[b].y - s[a].y;
            if (dy * dy >= best) {      // y 差不小於 delta，後面的更遠
                break;
            }
            long long d = dist2(s[a], s[b]);
            best = d < best ? d : best;
        }
    }
    printf("n=%d xm=%d dL2=%lld dR2=%lld strip=%d d2=%lld\n",
           hi - lo + 1, xm, dl, dr, k, best);
    return best;
}

int main(void)
{
    Point p[] = {{12, 5}, {2, 10}, {16, 1}, {9, 14}, {13, 9},
                 {4, 6}, {18, 5}, {11, 8}, {5, 11}, {14, 14}};
    int n = 10;
    Point *tmp = malloc(n * sizeof *tmp);
    if (tmp == NULL) {
        return 1;
    }
    qsort(p, n, sizeof *p, cmp_x);
    long long d2 = closest(p, tmp, 0, n - 1);
    printf("answer d2=%lld, %d distances (all pairs: %d)\n", d2, dist_calls,
           n * (n - 1) / 2);
    free(tmp);
    return 0;
}
\end{example}
\iosample{}{n=5 xm=5 dL2=10 dR2=40 strip=3 d2=10 \\ n=5 xm=14 dL2=17 dR2=20 strip=5 d2=17 \\ n=10 xm=11 dL2=10 dR2=17 strip=5 d2=5 \\ answer d2=5, 17 distances (all pairs: 45)}

三行追蹤依序是上表的 A--E、F--J、A--J 三列，共算了 $17$ 次距離。 點數多時差距更明
顯：隨機的點上，帶狀區域裡的點通常只往後看一兩個點就停，算距離的次數和 $n$ 同一
級，兩兩比較卻是 $n^2/2$。

\section{逆序數}
\label{sec:dc-inv}

第~1~章定義過\textbf{逆序對}：$i < j$ 而 $A[i] > A[j]$。 逆序對的個數（逆序數）量
的是一個序列離「由小到大」有多遠：排好的序列是 $0$，完全相反的是
$\binom{n}{2}$。 第~1~章證明了插入排序的挪動次數等於逆序數，第~14~章說明氣泡排序
的交換次數也是。 所以這兩種排序都能用來數逆序數，但它們最壞要 $\Theta(n^2)$；兩兩
檢查也是 $\Theta(n^2)$。 這一節用合併排序在 $\Theta(n \lg n)$ 內數完。

\subsection{在合併時順便數}

把 $A[p..r]$ 從 $q$ 切開，逆序對分成三種：兩個數都在左半、都在右半、一個在左一個在
右（叫\textbf{跨兩半}的逆序對）。 前兩種遞迴去數。 第三種只和「左半有哪些數、右半
有哪些數」有關，和兩半內部的順序無關：把左半或右半各自重新排列，跨兩半的數對還是
那些，大小關係也不變。 所以可以先讓遞迴把兩半各自排好，再數跨兩半的，而兩段排好
的序列正是第~14~章 \proc{Merge} 的輸入。

\begin{pseudo}[label=alg:dc-merge-count]{合併 $A[p..q]$ 和 $A[q+1..r]$，回傳跨兩半的逆序對個數}
\Function{Merge-Count}{$A, B, p, q, r$}
    \For{$k \gets p$ \To $r$}
        \State $B[k] \gets A[k]$
    \EndFor
    \State $i \gets p$；$j \gets q + 1$；$c \gets 0$
    \For{$k \gets p$ \To $r$}
        \If{$i \le q$ \AND ($j > r$ \OR $B[i] \le B[j]$)}
            \State $A[k] \gets B[i]$；$i \gets i + 1$
        \Else
            \State $A[k] \gets B[j]$；$j \gets j + 1$
            \State $c \gets c + (q - i + 1)$ \Comment{左半還沒拿的都比它大}
        \EndIf
    \EndFor
    \State \Return $c$
\EndFunction
\end{pseudo}

\begin{pseudo}[label=alg:dc-sort-count]{排序並回傳 $A[p..r]$ 的逆序數}
\Function{Sort-Count}{$A, B, p, r$}
    \If{$p \ge r$}
        \State \Return $0$
    \EndIf
    \State $q \gets \lfloor (p + r)/2 \rfloor$
    \State $a \gets$ \Call{Sort-Count}{$A, B, p, q$}
    \State $b \gets$ \Call{Sort-Count}{$A, B, q + 1, r$}
    \State \Return $a + b + $\Call{Merge-Count}{$A, B, p, q, r$}
\EndFunction
\end{pseudo}

\proc{Merge-Count} 和虛擬碼~\ref{alg:so-merge} 的 \proc{Merge} 只多了計數器 $c$：第 4 行
設成 $0$，第 10 行累加，第 11 行回傳。

\begin{lemma}
\label{lem:dc-cross}
$B[p..q]$、$B[q+1..r]$ 各自由小到大排好時，\proc{Merge-Count} 回傳的 $c$ 等於滿足
$u$ 在左半、$v$ 在右半、$u > v$ 的數對 $(u, v)$ 的個數。
\end{lemma}

\begin{proof}
看右半的某個數 $v = B[j]$ 被第 9 行放進 $A$ 的那一刻。 這時左半已經拿走的是
$B[p..i-1]$，還沒拿的是 $B[i..q]$。
\begin{itemize}
\item 還沒拿的都比 $v$ 大：如果 $i \le q$，走到第 8 行表示第 6 行的條件不成立，而
$i \le q$、$j \le r$，所以 $B[i] > B[j] = v$；左半排好了，$B[i..q]$ 都不小於
$B[i]$，也就都大於 $v$。
\item 已經拿走的都不比 $v$ 大：每個 $B[i']$（$i' < i$）被拿走時，要嘛
$B[i'] \le B[j']$，$j' \le j$ 是當時右半的位置，右半排好了，所以
$B[i'] \le B[j'] \le v$；要嘛當時 $j' > r$，但 $v$ 還沒被拿，不可能。
\end{itemize}
所以左半裡大於 $v$ 的數剛好是 $B[i..q]$，共 $q - i + 1$ 個，第 10 行加的正是以 $v$
為右端的跨兩半逆序對個數。 右半每個數都被拿一次，加起來就是全部跨兩半的逆序對。
\end{proof}

第 6 行的 $B[i] \le B[j]$ 是關鍵：相等時先拿左邊的。 相等的兩個數不算逆序；如果改成
$<$，相等時先拿右邊，左半那個相等的數還沒被拿，就會被多算進去。 這和第~14~章合併
排序要穩定時的規則是同一條。

由引理和前面的討論，對 $n = r - p + 1$ 歸納可得 \proc{Sort-Count} 回傳
$A[p..r]$ 的逆序數，並把它排好。 \proc{Merge-Count} 是 $\Theta(n)$，所以
$T(n) = 2T(n/2) + \Theta(n) = \Theta(n \lg n)$。

\subsection{手算}

用 $A = \langle 6, 2, 9, 4, 4, 1, 8, 3 \rangle$。 下面兩層的六次合併各數到 $1$ 個（例如
$\langle 6 \rangle$ 和 $\langle 2 \rangle$ 合併時數到 $(6, 2)$），左半
$\langle 6, 2, 9, 4 \rangle$ 共 $3$ 個，右半 $\langle 4, 1, 8, 3 \rangle$ 也是
$3$ 個。 最上層合併 $\langle 2, 4, 6, 9 \rangle$ 和 $\langle 1, 3, 4, 8 \rangle$：

\begin{center}
\begin{tabular}{cccccl}
\toprule
$k$ & 比較 & 拿 & 左半還沒拿的 & $c$ 增加 & 說明 \\
\midrule
1 & $2$ 對 $1$ & 右 $1$ & $2, 4, 6, 9$ & 4 & $(2,1)$、$(4,1)$、$(6,1)$、$(9,1)$ \\
2 & $2$ 對 $3$ & 左 $2$ & & & \\
3 & $4$ 對 $3$ & 右 $3$ & $4, 6, 9$ & 3 & \\
4 & $4$ 對 $4$ & 左 $4$ & & & 相等，先拿左邊 \\
5 & $6$ 對 $4$ & 右 $4$ & $6, 9$ & 2 & 左半的 $4$ 已經拿走，不算 \\
6 & $6$ 對 $8$ & 左 $6$ & & & \\
7 & $9$ 對 $8$ & 右 $8$ & $9$ & 1 & \\
8 & 右半用完 & 左 $9$ & & & \\
\midrule
& & & & 10 & \\
\bottomrule
\end{tabular}
\end{center}

總共 $3 + 3 + 10 = 16$ 個逆序對。 第 5 步如果改成相等時先拿右邊，那時左半還沒拿的
是 $4, 6, 9$，會加 $3$ 而不是 $2$，把 $(4, 4)$ 算成逆序，答案多一。

逆序數的一個用途是比較兩份排名：兩位評審替同一群人排名次，依第一位的順序寫下第二位給
的名次，這個序列的逆序數就是兩人意見相反的配對數。

C 程式只要在第~14~章範例 \texttt{merge\_sort.c} 的 \texttt{merge} 裡，從右半拿數的那
一個分支加上 \texttt{cross += mid - i + 1;}，再讓遞迴函式回傳三個計數的和。 $n$ 個數
最多有 $n(n-1)/2$ 個逆序對，$n = 10^5$ 時接近 $5 \times 10^9$，超過 \texttt{int}，計
數要用 \texttt{long long}。


\section{計算幾何入門}
\label{sec:dc-geom}

第~\ref{sec:dc-closest}~節的最近點對是本章第一個平面上的問題。 研究所舊卷還常考另外幾
個平面問題：三個點是左轉還是右轉、兩條線段有沒有相交、一個點在不在多邊形裡面，以
及把一群點用最小的凸多邊形圍起來的\textbf{凸包}。 這一節依序講這幾個問題。 凸包有
兩種 $\Theta(n \lg n)$ 的做法，一種靠排序加堆疊，一種靠分治；最後用第~15~章的排序
下界說明 $n \lg n$ 已經不能再少。

這一節的點都用整數座標，點 $p$ 的座標寫成 $p.x$、$p.y$。 所有判斷都只用加、減、乘，
不用除法、斜率或角度，所以整數座標算出來的答案是精確的。

\subsection{外積與轉向}
\label{sec:dc-orient}

從 $p$ 走到 $q$，接著要走到 $r$，在 $q$ 要左轉還是右轉？ 後面每一個演算法都要問這
個問題。 令 $q - p = (u_x, u_y)$、$r - p = (v_x, v_y)$，定義
\[
\proc{Orient}(p, q, r) = u_x v_y - u_y v_x
= (q.x - p.x)(r.y - p.y) - (q.y - p.y)(r.x - p.x).
\]
它就是二階行列式 $\left|\begin{smallmatrix} u_x & v_x \\ u_y & v_y
\end{smallmatrix}\right|$，也叫兩個向量的\textbf{外積}（cross product；嚴格說是三維
外積的 $z$ 分量）。 高中數學說過，它的絕對值是兩個向量張成的平行四邊形面積，也就是
三角形 $pqr$ 面積的兩倍。 正負號表示方向：

\begin{itemize}
\item $\proc{Orient}(p, q, r) > 0$：$r$ 在有向直線 $p \to q$ 的左邊，$p \to q \to r$
是\textbf{左轉}（逆時針）。
\item $\proc{Orient}(p, q, r) < 0$：$r$ 在右邊，是\textbf{右轉}（順時針）。
\item $\proc{Orient}(p, q, r) = 0$：三點\textbf{共線}。
\end{itemize}

哪個符號是左轉，代最簡單的例子就記得住：$p = (0, 0)$、$q = (1, 0)$、$r = (0, 1)$，
$r$ 在 $x$ 軸正方向的左邊，$\proc{Orient} = 1 \cdot 1 - 0 \cdot 0 = 1 > 0$。

\begin{pseudo}[label=alg:dc-orient]{三點的轉向（外積）}
\Function{Orient}{$p, q, r$}
    \State \Return $(q.x - p.x)(r.y - p.y) - (q.y - p.y)(r.x - p.x)$
\EndFunction
\end{pseudo}

例如 $p = (1, 1)$、$q = (5, 2)$，$q - p = (4, 1)$：

\begin{itemize}
\item $r = (3, 6)$：$4 \cdot 5 - 1 \cdot 2 = 18 > 0$，左轉。
\item $r = (6, 0)$：$4 \cdot (-1) - 1 \cdot 5 = -9 < 0$，右轉。
\item $r = (9, 3)$：$4 \cdot 2 - 1 \cdot 8 = 0$，共線；$r - p = (8, 2)$ 正好是
$q - p$ 的兩倍。
\end{itemize}

\begin{center}
\begin{minipage}{\linewidth}
\centering
\begin{tikzpicture}[x=0.44cm, y=0.44cm]
\begin{scope}
\draw[dc-ray] (1,1) -- (10.2,3.3);
\draw[dc-vec] (1,1) -- (5,2);
\draw[dc-vec] (5,2) -- (3,6);
\node[dc-pt, label={[dc-ptlab]below:$p$}] at (1,1) {};
\node[dc-pt, label={[dc-ptlab]below right:$q$}] at (5,2) {};
\node[dc-pt, label={[dc-ptlab]above:$r$}] at (3,6) {};
\node[font=\scriptsize] at (5,-0.6) {$\proc{Orient} = 18$，左轉};
\end{scope}
\begin{scope}[xshift=5.3cm]
\draw[dc-ray] (1,1) -- (10.2,3.3);
\draw[dc-vec] (1,1) -- (5,2);
\draw[dc-vec] (5,2) -- (6,0);
\node[dc-pt, label={[dc-ptlab]above left:$p$}] at (1,1) {};
\node[dc-pt, label={[dc-ptlab]above:$q$}] at (5,2) {};
\node[dc-pt, label={[dc-ptlab]right:$r$}] at (6,0) {};
\node[font=\scriptsize] at (5,-0.6) {$\proc{Orient} = -9$，右轉};
\end{scope}
\begin{scope}[xshift=10.6cm]
\draw[dc-ray] (1,1) -- (10.2,3.3);
\draw[dc-vec] (1,1) -- (5,2);
\draw[dc-vec] (5,2) -- (9,3);
\node[dc-pt, label={[dc-ptlab]above left:$p$}] at (1,1) {};
\node[dc-pt, label={[dc-ptlab]above:$q$}] at (5,2) {};
\node[dc-pt, label={[dc-ptlab]above:$r$}] at (9,3) {};
\node[font=\scriptsize] at (5,-0.6) {$\proc{Orient} = 0$，共線};
\end{scope}
\end{tikzpicture}
\diagcaption{同一個 $p \to q$，三種 $r$。 灰色虛線是直線 $pq$。}
\label{fig:dc-orient}
\end{minipage}
\end{center}

後面會用到兩個性質，都可以直接展開驗證：把 $q$、$r$ 對調，外積變號；把三個點輪換，
外積不變：
\[
\proc{Orient}(p, q, r) = \proc{Orient}(q, r, p) = \proc{Orient}(r, p, q)
= -\proc{Orient}(p, r, q).
\]
上面的例子 $\proc{Orient}(q, r, p) = (-2)(-1) - 4 \cdot (-4) = 18$，和
$\proc{Orient}(p, q, r)$ 相同。

\textbf{為什麼不用斜率}。 比較斜率要做除法；鉛直線的斜率不存在，要另外處理；浮點數
還有捨入誤差，三點幾乎共線時可能判斷錯。 外積只有乘法和減法，整數座標算出來是精確
的。 C 程式只要注意溢位：座標在 $\pm 10^9$ 以內時，座標差在 $2 \times 10^9$ 以內，
乘積在 $4 \times 10^{18}$ 以內，兩個乘積相減不超過 $8 \times 10^{18}$，放得進
\texttt{long long}（上限約 $9.22 \times 10^{18}$），放不進 \texttt{int}。 座標差要
先轉成 \texttt{long long} 再乘。

\textbf{凸多邊形的判斷}。 一個\textbf{簡單多邊形}是邊和邊除了相鄰兩邊在頂點相接以外
都不相交的多邊形。 把它的頂點依序列成 $v_1, v_2, \ldots, v_n$，它是凸多邊形若且唯若
每三個連續頂點 $v_i, v_{i+1}, v_{i+2}$（下標超過 $n$ 就繞回 $1$）的 \proc{Orient}
全部同號：逆時針列出時全部為正，順時針時全部為負。 掃一次 $\Theta(n)$。 「簡單」這
個條件不能少：把正五邊形的頂點每隔一個連起來，得到五角星，每個轉彎都同一個方向，
但它和自己相交，不是凸多邊形。 有某個外積是 $0$ 表示有三個連續頂點共線，算不算凸照題
目的定義。

\subsection{線段相交}
\label{sec:dc-segint}

給兩條線段 $\overline{p_1 p_2}$、$\overline{p_3 p_4}$，判斷它們有沒有共同點，端點碰到
也算。 先算四個外積：
\[
\begin{aligned}
d_1 &= \proc{Orient}(p_3, p_4, p_1), & d_2 &= \proc{Orient}(p_3, p_4, p_2), \\
d_3 &= \proc{Orient}(p_1, p_2, p_3), & d_4 &= \proc{Orient}(p_1, p_2, p_4).
\end{aligned}
\]
$d_1$、$d_2$ 說明 $p_1$、$p_2$ 各在直線 $p_3 p_4$ 的哪一邊。 $d_1 d_2 < 0$ 表示
$p_1$、$p_2$ 在直線 $p_3 p_4$ 的兩側，這時說線段 $\overline{p_1 p_2}$ \textbf{跨過}
直線 $p_3 p_4$。 $d_3$、$d_4$ 同理。 兩條線段互相跨過對方所在的直線，就一定相交；麻
煩的是有某個 $d$ 等於 $0$ 的情形，要另外檢查那個端點是不是真的落在另一條線段上，而
不只是落在它的延長線上。

\begin{pseudo}[label=alg:dc-onseg]{已知 $r$ 在直線 $pq$ 上，判斷 $r$ 是否在線段 $\overline{pq}$ 上}
\Function{On-Segment}{$p, q, r$}
    \State $a \gets (\min(p.x, q.x) \le r.x \le \max(p.x, q.x))$
    \State $b \gets (\min(p.y, q.y) \le r.y \le \max(p.y, q.y))$
    \State \Return $a$ \AND $b$
\EndFunction
\end{pseudo}

\begin{pseudo}[label=alg:dc-segint]{兩條線段是否相交}
\Function{Segments-Intersect}{$p_1, p_2, p_3, p_4$}
    \State $d_1 \gets$ \Call{Orient}{$p_3, p_4, p_1$}；$d_2 \gets$ \Call{Orient}{$p_3, p_4, p_2$}
    \State $d_3 \gets$ \Call{Orient}{$p_1, p_2, p_3$}；$d_4 \gets$ \Call{Orient}{$p_1, p_2, p_4$}
    \If{$d_1 d_2 < 0$ \AND $d_3 d_4 < 0$}
        \State \Return \TRUE \Comment{互相跨過}
    \EndIf
    \If{$d_1 = 0$ \AND \Call{On-Segment}{$p_3, p_4, p_1$}}
        \State \Return \TRUE
    \EndIf
    \If{$d_2 = 0$ \AND \Call{On-Segment}{$p_3, p_4, p_2$}}
        \State \Return \TRUE
    \EndIf
    \If{$d_3 = 0$ \AND \Call{On-Segment}{$p_1, p_2, p_3$}}
        \State \Return \TRUE
    \EndIf
    \If{$d_4 = 0$ \AND \Call{On-Segment}{$p_1, p_2, p_4$}}
        \State \Return \TRUE
    \EndIf
    \State \Return \FALSE
\EndFunction
\end{pseudo}

\proc{On-Segment} 只檢查 $r$ 在不在以 $p$、$q$ 為對角的長方形裡。 這只在已知 $r$ 在
直線 $pq$ 上時才夠：直線和這個長方形的交集正好是線段 $\overline{pq}$。 所以第 6--13
行每一個都先確認 $d = 0$。

\begin{lemma}
\label{lem:dc-segint}
\proc{Segments-Intersect} 回傳 \TRUE{} 若且唯若兩條線段至少有一個共同點。
\end{lemma}

\begin{proof}
\textbf{回傳 \TRUE{} 時真的相交}。 如果是第 6--13 行回傳，某個端點在另一條線段上，
它就是共同點。 如果是第 5 行回傳，$d_1 d_2 < 0$，$p_1$、$p_2$ 在直線 $p_3 p_4$ 的兩
側，線段 $\overline{p_1 p_2}$ 和直線 $p_3 p_4$ 交於一點 $X$。 同理 $d_3 d_4 < 0$，線
段 $\overline{p_3 p_4}$ 和直線 $p_1 p_2$ 交於一點 $Y$。 兩條直線不重合（$d_1 \ne
0$），也不平行（平行時 $p_1$、$p_2$ 會在直線 $p_3 p_4$ 的同一側），所以只交於一
點，$X$ 和 $Y$ 都是這一點，$X = Y$ 同時在兩條線段上。

\textbf{相交時一定回傳 \TRUE}。 設兩條線段有共同點 $X$。 如果第 6--13 行的四個條件有
一個成立就回傳 \TRUE 了，所以只要證明：四個都不成立時，第 4 行的條件成立。
\begin{itemize}
\item 先證 $d_1 \ne 0$。 假設 $d_1 = 0$，$p_1$ 在直線 $p_3 p_4$ 上，但（第 6 行不成
立）不在線段 $\overline{p_3 p_4}$ 上。 如果 $d_2 \ne 0$，兩條直線不同，只交於 $p_1$；
$X$ 在兩條直線上，所以 $X = p_1$，$p_1$ 在線段 $\overline{p_3 p_4}$ 上，矛盾。 如果
$d_2 = 0$，四個點共線，兩條線段是同一條直線上的兩段，有共同點 $X$。 同一條直線上的兩
段重疊時，起點比較後面的那一段，它的起點一定也在另一段裡，所以某個端點在另一條線段
上，第 6--13 行有一個條件成立，矛盾。 同理 $d_2, d_3, d_4$ 都不是 $0$。
\item 再證 $d_1 d_2 < 0$。 如果 $d_1$、$d_2$ 同號，$p_1$、$p_2$ 在直線 $p_3 p_4$ 的同
一側（嚴格），整條線段 $\overline{p_1 p_2}$ 都在那一側，碰不到直線 $p_3 p_4$ 上的
$X$，矛盾。 同理 $d_3 d_4 < 0$。 \qedhere
\end{itemize}
\end{proof}

圖~\ref{fig:dc-segcases} 是四種情形，每一種的四個外積都在下面的表裡（程式算過）。

\begin{center}
\begin{minipage}{\linewidth}
\centering
\begin{tikzpicture}[x=0.36cm, y=0.36cm]
\begin{scope}
\draw[dc-seg] (1,1) -- (7,5);
\draw[dc-seg2] (1,5) -- (6,1);
\foreach \x/\y/\n/\pos in {1/1/1/below, 7/5/2/above, 1/5/3/above, 6/1/4/below} {
    \node[dc-pt, label={[dc-ptlab]\pos:$p_\n$}] at (\x,\y) {};
}
\node[font=\scriptsize] at (4,-1.6) {(a) 互相跨過};
\end{scope}
\begin{scope}[xshift=3.4cm]
\draw[dc-seg] (1,1) -- (4,3);
\draw[dc-seg2] (5,1) -- (7,6);
\draw[dc-ray] (4,3) -- (6.45,4.64);
\foreach \x/\y/\n/\pos in {1/1/1/below, 4/3/2/above, 5/1/3/below, 7/6/4/above} {
    \node[dc-pt, label={[dc-ptlab]\pos:$p_\n$}] at (\x,\y) {};
}
\node[font=\scriptsize] at (4,-1.6) {(b) 只有一邊跨過};
\end{scope}
\begin{scope}[xshift=6.8cm]
\draw[dc-seg] (0,0) -- (6,3);
\draw[dc-seg2] (4,2) -- (5,6);
\foreach \x/\y/\n/\pos in {0/0/1/below, 6/3/2/below, 4/2/3/below, 5/6/4/above} {
    \node[dc-pt, label={[dc-ptlab]\pos:$p_\n$}] at (\x,\y) {};
}
\node[font=\scriptsize] at (3.5,-1.6) {(c) 端點在另一段上};
\end{scope}
\begin{scope}[xshift=10.2cm]
\draw[dc-seg] (0,0) -- (2,1);
\draw[dc-seg2] (4,2) -- (6,3);
\draw[dc-ray] (-0.6,-0.3) -- (6.8,3.4);
\foreach \x/\y/\n/\pos in {0/0/1/below, 2/1/2/below, 4/2/3/below, 6/3/4/below} {
    \node[dc-pt, label={[dc-ptlab]\pos:$p_\n$}] at (\x,\y) {};
}
\node[font=\scriptsize] at (3,-1.6) {(d) 共線但分開};
\end{scope}
\end{tikzpicture}
\diagcaption{線段相交的四種情形。 黑色是 $\overline{p_1 p_2}$，藍色是
$\overline{p_3 p_4}$，灰色虛線是需要看的延長線。}
\label{fig:dc-segcases}
\end{minipage}
\end{center}

\begin{center}
\small
\begin{tabular}{lcccccl}
\toprule
& 座標 & $d_1$ & $d_2$ & $d_3$ & $d_4$ & 結果 \\
\midrule
(a) & $(1,1)$、$(7,5)$、$(1,5)$、$(6,1)$ & $-20$ & $24$ & $24$ & $-20$ & 第 5 行，\TRUE \\
(b) & $(1,1)$、$(4,3)$、$(5,1)$、$(7,6)$ & $20$ & $9$ & $-8$ & $3$ & 第 14 行，\FALSE \\
(c) & $(0,0)$、$(6,3)$、$(4,2)$、$(5,6)$ & $14$ & $-7$ & $0$ & $21$ & 第 11 行，\TRUE \\
(d) & $(0,0)$、$(2,1)$、$(4,2)$、$(6,3)$ & $0$ & $0$ & $0$ & $0$ & 第 14 行，\FALSE \\
\bottomrule
\end{tabular}
\end{center}

(b) 的 $d_3 d_4 < 0$，$\overline{p_3 p_4}$ 跨過直線 $p_1 p_2$；但 $d_1 d_2 > 0$，
$\overline{p_1 p_2}$ 整條在直線 $p_3 p_4$ 的同一側，兩條直線的交點落在
$\overline{p_1 p_2}$ 的延長線上。 (c) 的 $d_3 = 0$，$p_3 = (4, 2)$ 在直線 $p_1 p_2$
上，而且 $0 \le 4 \le 6$、$0 \le 2 \le 3$，在線段上。 (d) 四個外積都是 $0$，四點共線；
四次 \proc{On-Segment} 都不成立（例如 $p_3$ 的 $x = 4$ 不在 $[0, 2]$ 裡），所以不相
交。 如果只看「有某個 $d = 0$ 就算相交」，(d) 就會答錯。

\proc{Segments-Intersect} 只做常數次乘法和比較，$\Theta(1)$。 C 程式裡不要真的算
$d_1 d_2$：兩個外積都可能接近 $8 \times 10^{18}$，相乘一定溢位。 改成比正負號：
\texttt{(d1 > 0 \&\& d2 < 0) || (d1 < 0 \&\& d2 > 0)}。

$n$ 條線段裡有沒有任兩條相交，兩兩檢查要 $\Theta(n^2)$。 用一條鉛直的\textbf{掃描線}
從左往右掃，同時用平衡二元搜尋樹（第~10~章）記住掃描線碰到的線段由下到上的順序，可以
在 $O(n \lg n)$ 內回答（Shamos 和 Hoey，1976）。 本書不展開，延伸閱讀列了參考資料。

\subsection{點在多邊形內：射線法}
\label{sec:dc-pip}

給一個簡單多邊形，頂點依序是 $V[1..n]$（順時針、逆時針都可以），和一個不在邊界上的點
$q$，判斷 $q$ 在多邊形裡面還是外面。

\textbf{射線法}（ray casting）：從 $q$ 往右畫一條水平射線，數它穿過幾條邊。 奇數在裡
面，偶數在外面。 理由是射線走到夠遠的地方一定在多邊形外面，而每穿過邊界一次，裡外就
交換一次；從外面往回數，穿過奇數次才會回到裡面。

麻煩的是射線剛好通過頂點，或沿著一條水平邊走。 下面的規則一次解決：

\begin{quote}
把高度剛好等於 $q.y$ 的頂點都當成在射線\textbf{下方}。 一條邊 $(a, b)$ 只有在
「$a.y > q.y$」和「$b.y > q.y$」恰好一個成立時，才算和射線所在的水平線相交。
\end{quote}

照這個規則，射線通過頂點 $v$ 時有三種情形：

\begin{itemize}
\item $v$ 的兩個相鄰頂點一上一下（射線真的穿過邊界）：$v$ 算在下方，連到上方鄰居的
那條邊算一次，連到下方鄰居的那條邊不算，合計一次。
\item 兩個鄰居都在上方（射線只是從下面碰到 $v$）：兩條邊都算，合計兩次，奇偶不變。
\item 兩個鄰居都在下方（從上面碰到 $v$）：兩條邊都不算，奇偶也不變。
\end{itemize}

水平邊的兩端都算在下方，不會被計算；它兩端連出去的邊照上面三種情形處理。

接著判斷交點是不是在 $q$ 的右邊。 設邊 $(a, b)$ 和水平線 $y = q.y$ 交於 $X$。 把
$\proc{Orient}(a, b, q)$ 減去 $\proc{Orient}(a, b, X)$：因為 $X.y = q.y$，兩式的第一項
都是 $(b.x - a.x)(q.y - a.y)$，相減消掉；又 $X$ 在直線 $ab$ 上，$\proc{Orient}(a, b, X) = 0$。
所以
\[
\proc{Orient}(a, b, q) = (b.y - a.y)(X.x - q.x).
\]
所以交點在 $q$ 右邊（$X.x > q.x$）若且唯若 $\proc{Orient}(a, b, q)$ 和 $b.y - a.y$
同號。 這兩個數都不會是 $0$：邊跨過水平線，$b.y \ne a.y$；$q$ 不在邊界上，
$X \ne q$。 不必算出 $X$，也不必做除法。

\begin{pseudo}[label=alg:dc-pip]{射線法：$q$ 是否在多邊形 $V[1..n]$ 內（$q$ 不在邊界上）}
\Function{Point-In-Polygon}{$V, n, q$}
    \State $c \gets 0$
    \For{$i \gets 1$ \To $n$}
        \State $a \gets V[i]$；$b \gets V[i \bmod n + 1]$
        \If{$a.y > q.y$ 和 $b.y > q.y$ 恰好一個成立}
            \If{\Call{Orient}{$a, b, q$} $\cdot (b.y - a.y) > 0$} \Comment{交點在 $q$ 的右邊}
                \State $c \gets c + 1$
            \EndIf
        \EndIf
    \EndFor
    \State \Return $c$ 是不是奇數
\EndFunction
\end{pseudo}

每條邊看一次，$\Theta(n)$。 多邊形是凸的時候可以更快：從 $V[1]$ 連到其他頂點，把多
邊形切成 $n - 2$ 個扇形，用二分搜尋找 $q$ 落在哪一個扇形，再用一次外積判斷，
$O(\lg n)$。

\textbf{手算}。 多邊形的八個頂點依逆時針是
\[
v_1(2, 1),\ v_2(10, 2),\ v_3(8, 5),\ v_4(12, 9),\ v_5(6, 10),\ v_6(5, 5),\ v_7(3, 8),\ v_8(1, 4).
\]
它在 $v_3$、$v_6$ 各凹進去一塊。 要判斷的點是 $A(3, 5)$、$B(5, 7)$、$C(2, 8)$。

\begin{center}
\begin{minipage}{\linewidth}
\centering
\begin{tikzpicture}[x=0.55cm, y=0.48cm]
\fill[coreFill] (2,1) -- (10,2) -- (8,5) -- (12,9) -- (6,10) -- (5,5) -- (3,8) -- (1,4) -- cycle;
\draw[thick, black!70] (2,1) -- (10,2) -- (8,5) -- (12,9) -- (6,10) -- (5,5) -- (3,8) -- (1,4) -- cycle;
\draw[dc-rayq] (3,5) -- (13.5,5);
\draw[dc-rayq] (5,7) -- (13.5,7);
\draw[dc-rayq] (2,8) -- (13.5,8);
\foreach \n/\x/\y/\pos in {1/2/1/below, 2/10/2/below right, 3/8/5/left, 4/12/9/right, 5/6/10/above, 6/5/5/below, 7/3/8/above, 8/1/4/left} {
    \node[dc-pt, label={[dc-ptlab]\pos:$v_\n$}] at (\x,\y) {};
}
\foreach \n/\x/\y in {A/3/5, B/5/7, C/2/8} {
    \node[circle, draw=optStroke, fill=white, thick, inner sep=0pt, minimum size=5pt,
          label={[dc-ptlab, text=optStroke]left:$\n$}] at (\x,\y) {};
}
\end{tikzpicture}
\diagcaption{射線法。 藍底是多邊形，紅色空心點是要判斷的點，紅色虛線是往右的射線。
$A$ 的射線通過 $v_6$（從下面碰到）和 $v_3$（穿過），$C$ 的射線從上面碰到 $v_7$。}
\label{fig:dc-pip}
\end{minipage}
\end{center}

$A(3, 5)$ 的射線高度是 $5$，$v_3$、$v_6$ 剛好在這個高度，照規則都算在下方。 只列出
跨過 $y = 5$ 的邊：

\begin{center}
\begin{tabular}{lccccl}
\toprule
邊 & $\proc{Orient}(a, b, A)$ & $b.y - a.y$ & 乘積 & 交點 $X.x$ & 計數 \\
\midrule
$v_3 v_4$ & $20$ & $4$ & $80$ & $8$（$v_3$） & $1$ \\
$v_5 v_6$ & $-10$ & $-5$ & $50$ & $5$（$v_6$） & $2$ \\
$v_6 v_7$ & $6$ & $3$ & $18$ & $5$（$v_6$） & $3$ \\
$v_7 v_8$ & $6$ & $-4$ & $-24$ & $1.5$ & 不算，在左邊 \\
\bottomrule
\end{tabular}
\end{center}

$c = 3$ 是奇數，$A$ 在裡面。 $v_6$ 的兩個鄰居 $v_5$、$v_7$ 都在上方，射線只是從下面碰到
$v_6$，兩條邊各算一次；$v_3$ 的鄰居 $v_2$ 在下方、$v_4$ 在上方，射線真的穿過去，只
有 $v_3 v_4$ 算一次。 交點那一欄只是幫忙看圖，程序本身只看乘積的正負號。

$B(5, 7)$ 在 $v_6$ 那個凹口裡。 跨過 $y = 7$ 的四條邊交點在 $x = 10$、$5.4$、$11/3$、
$2.5$，右邊兩個，$c = 2$，在外面。 $C(2, 8)$ 的射線從上面碰到 $v_7$：$v_6 v_7$、
$v_7 v_8$ 的兩端都不在 $y = 8$ 的上方，都不算；只有 $v_3 v_4$（交點 $x = 11$）和
$v_5 v_6$（交點 $x = 5.6$）算，$c = 2$，在外面。

\begin{examnote}
射線法的答案要寫出三件事：射線的方向、怎麼數（奇數在內），以及射線通過頂點或水平邊
時怎麼處理。 最後一件最常被扣分。 本書的寫法是「高度等於 $q.y$ 的頂點算在下方」；
有的課本算在上方，或把射線稍微往上平移一點點，效果相同，選一種寫清楚就好。 題目如果
允許 $q$ 在邊界上，要先用 \proc{Orient} 和 \proc{On-Segment} 逐邊檢查。
\end{examnote}

\subsection{凸包與 Graham 掃描法}
\label{sec:dc-graham}

平面上有限個點的集合 $Q$，它的\textbf{凸包}（convex hull）$\mathrm{CH}(Q)$ 是包含 $Q$ 的
最小凸多邊形。 想像在每個點釘一根釘子，拿一條橡皮筋撐開套住所有釘子再放手，橡皮筋圍
出來的就是凸包。 凸包的頂點一定是 $Q$ 裡的點；凸包演算法要依\textbf{逆時針}順序輸出
這些頂點。 落在凸包的邊上、但不是轉角的點不算頂點，本書不輸出。

\textbf{Graham 掃描法}（Graham scan，Graham 1972）的步驟：

\begin{enumerate}
\item 找 $y$ 座標最小的點（有好幾個時取 $x$ 最小的），叫它 $p_0$。 它一定是凸包的頂
點。
\item 其他點依「繞 $p_0$ 的極角」由小到大排序。 極角是從 $p_0$ 往右的水平方向逆時
針轉到 $p_0 \to p$ 的角度；極角相同時，離 $p_0$ 近的排前面。
\item 依這個順序逐點處理，用一個堆疊記住目前的凸包。 新的點進來時，只要堆疊頂端兩個
點和新的點不是左轉，就把頂端彈掉；然後把新的點推進去。
\end{enumerate}

\textbf{排序不必算角度}。 $p_0$ 是最低的點（同高時是最左的），其他點都在它上方，或
在它右邊的水平線上，所以極角都在 $[0^\circ, 180^\circ)$ 裡。 兩個極角 $\theta_a$、
$\theta_b$ 都在這個範圍時，$\theta_a < \theta_b$ 若且唯若從 $p_0 \to a$ 轉到
$p_0 \to b$ 是逆時針，也就是 $\proc{Orient}(p_0, a, b) > 0$。 所以比較函式寫成：
$\proc{Orient}(p_0, a, b) > 0$ 時 $a$ 在前；等於 $0$（同一個方向）時近的在前。 距離
比平方就好，不必開根號。

\begin{pseudo}[label=alg:dc-graham]{Graham 掃描法，$Q[1..n]$ 至少兩個不同的點}
\Function{Graham-Scan}{$Q, n$}
    \State 把 $y$ 最小的點（$y$ 相同時取 $x$ 最小的）和 $Q[1]$ 交換
    \State $Q[2..n]$ 依繞 $Q[1]$ 的極角由小到大排序，極角相同時離 $Q[1]$ 近的在前
    \State $S[1] \gets Q[1]$；$S[2] \gets Q[2]$；$t \gets 2$ \Comment{$S[1..t]$ 是堆疊}
    \For{$i \gets 3$ \To $n$}
        \While{$t \ge 2$ \AND \Call{Orient}{$S[t-1], S[t], Q[i]$} $\le 0$}
            \State $t \gets t - 1$ \Comment{彈出頂端}
        \EndWhile
        \State $t \gets t + 1$；$S[t] \gets Q[i]$ \Comment{推入}
    \EndFor
    \State \Return $S[1..t]$
\EndFunction
\end{pseudo}

堆疊用陣列 $S[1..t]$ 表示，$S[t]$ 是頂端，$S[t-1]$ 是頂端下面那一個（第~6~章的陣列堆疊）。
第 6 行的條件用 $\le 0$：共線也要彈掉，這樣落在凸包邊上的點不會留在輸出裡。

\textbf{手算}。 十三個點：
\[
\begin{gathered}
A(6, 1),\ B(11, 3),\ C(9, 5),\ D(13, 8),\ E(10, 6),\ F(7, 11),\ G(4, 7), \\
H(2, 9),\ I(3, 4),\ J(1, 2),\ K(8, 7),\ L(11, 9),\ M(4, 3).
\end{gathered}
\]
$y$ 最小的是 $A$，所以 $p_0 = A$。 依繞 $A$ 的極角排序得到
\[
B,\ D,\ E,\ C,\ L,\ K,\ F,\ G,\ H,\ M,\ I,\ J.
\]
例如 $D - A = (7, 7)$、$E - A = (4, 5)$，$\proc{Orient}(A, D, E) = 7 \cdot 5 - 7 \cdot 4
= 7 > 0$，$D$ 在 $E$ 前面。 $M - A = (-2, 2)$ 和 $I - A = (-3, 3)$ 方向相同，外積是
$0$，$M$ 比較近，排在前面。

\begin{center}
\begin{minipage}{\linewidth}
\centering
\begin{tikzpicture}[x=0.6cm, y=0.5cm]
\foreach \x/\y in {11/3, 9/5, 13/8, 10/6, 7/11, 4/7, 2/9, 3/4, 1/2, 8/7, 11/9, 4/3} {
    \draw[black!18] (6,1) -- (\x,\y);
}
\fill[coreFill, fill opacity=0.6] (6,1) -- (11,3) -- (13,8) -- (7,11) -- (2,9) -- (1,2) -- cycle;
\draw[thick, coreStroke] (6,1) -- (11,3) -- (13,8) -- (7,11) -- (2,9) -- (1,2) -- cycle;
\foreach \n/\x/\y/\pos in {A/6/1/below, B/11/3/right, C/9/5/below right, D/13/8/right,
    E/10/6/right, F/7/11/above, G/4/7/above, H/2/9/above left, I/3/4/left, J/1/2/left,
    K/8/7/above, L/11/9/above right, M/4/3/below} {
    \node[dc-pt, label={[dc-ptlab]\pos:\n}] at (\x,\y) {};
}
\end{tikzpicture}
\diagcaption{Graham 掃描法的例子。 灰線從 $p_0 = A$ 連到其他點，排序就是依這些線的方
向逆時針排。 藍色多邊形是凸包 $A, B, D, F, H, J$；$L$ 落在邊 $DF$ 上，$M$ 和 $I$ 在
同一條灰線上。}
\label{fig:dc-graham}
\end{minipage}
\end{center}

掃描的過程如下。 「彈出」欄括號裡是第 6 行判斷時算的外積，「說明」欄是推入之前最後一
次判斷的外積（大於 $0$，所以停止彈出）。

\begin{center}
\begin{tabular}{clll}
\toprule
處理 & 彈出（外積） & 推入後的堆疊（底 $\to$ 頂） & 說明 \\
\midrule
（開始） & & $A, B$ & \\
$D$ & & $A, B, D$ & $\proc{Orient}(A, B, D) = 21 > 0$ \\
$E$ & & $A, B, D, E$ & $\proc{Orient}(B, D, E) = 11$ \\
$C$ & & $A, B, D, E, C$ & $\proc{Orient}(D, E, C) = 1$ \\
$L$ & $C$（$-2$）、$E$（$-7$） & $A, B, D, L$ & $\proc{Orient}(B, D, L) = 12$ \\
$K$ & & $A, B, D, L, K$ & $\proc{Orient}(D, L, K) = 7$ \\
$F$ & $K$（$-14$）、$L$（$0$） & $A, B, D, F$ & $L$ 在 $DF$ 上，共線也彈 \\
$G$ & & $A, B, D, F, G$ & $\proc{Orient}(D, F, G) = 33$ \\
$H$ & $G$（$-14$） & $A, B, D, F, H$ & $\proc{Orient}(D, F, H) = 27$ \\
$M$ & & $A, B, D, F, H, M$ & $\proc{Orient}(F, H, M) = 34$ \\
$I$ & $M$（$-4$） & $A, B, D, F, H, I$ & $\proc{Orient}(F, H, I) = 27$ \\
$J$ & $I$（$-12$） & $A, B, D, F, H, J$ & $\proc{Orient}(F, H, J) = 33$ \\
\bottomrule
\end{tabular}
\end{center}

凸包是 $A, B, D, F, H, J$（逆時針）。 共推入 $13$ 次（含一開始的兩個），彈出 $7$ 次。
處理 $L$ 時連續彈了兩個點：$E \to C \to L$ 右轉，$C$ 不可能是頂點；彈掉之後
$D \to E \to L$ 還是右轉，$E$ 也彈掉；$B \to D \to L$ 左轉，停。

\begin{whybox}[彈掉的點一定不是頂點]
設處理 $p = Q[i]$ 時，頂端 $s = S[t]$、下面是 $r = S[t-1]$，而 $\proc{Orient}(r, s, p)
\le 0$。 堆疊裡的點是依極角的順序推進去的，所以 $s$ 的極角介於 $r$ 和 $p$ 之間，$s$ 在
以 $p_0$ 為頂點、兩邊是射線 $p_0 r$ 和 $p_0 p$ 的角形區域裡。 另一方面，由輪換和對調
的性質，$\proc{Orient}(r, p, s) = -\proc{Orient}(r, s, p) \ge 0$，$s$ 在有向直線
$r \to p$ 的左邊或線上；$\proc{Orient}(r, p, p_0) = \proc{Orient}(p_0, r, p) > 0$（$r$
的極角比 $p$ 小），$p_0$ 也在左邊。 所以 $s$ 落在三角形 $p_0 r p$ 裡面或邊上。 三角形
的三個頂點都是點集裡的點，$s$ 在它裡面，就不可能是凸包的轉角。 （$r$、$p$ 極角相同時，三點都在從
$p_0$ 出發的同一條射線上，排序讓 $s$ 夾在 $r$、$p$ 之間，同樣不是轉角。）
\end{whybox}

\textbf{正確性}。 不變量：第 5 行每一輪開始時，$S[1..t]$ 依逆時針順序正是
$Q[1..i-1]$ 的凸包頂點。 一開始只有 $Q[1]$、$Q[2]$ 兩點，成立。 處理 $p = Q[i]$ 時，
$p$ 的極角是目前最大的（同方向時是最遠的），所以 $p$ 一定是 $Q[1..i]$ 凸包的頂點；被
彈掉的點由上面的方框都不是頂點；停下來時 $S[t-1] \to S[t] \to p$ 是左轉，堆疊裡原
本的轉彎都還是左轉，$p \to p_0$ 和 $p_0 \to S[2]$ 兩處也是左轉（極角的順序保證），
所以推入 $p$ 後的堆疊圍成一個每個轉角都左轉的凸多邊形。 它包含原來的凸包和 $p$，被彈
掉的點也都在它裡面，頂點又全部是點集裡的點，所以它就是 $Q[1..i]$ 的凸包。 迴圈結束時
$i = n + 1$，回傳的就是全部的凸包。 完整的證明見延伸閱讀。

\textbf{時間}。 第 2 行找最低點 $\Theta(n)$，第 3 行排序 $\Theta(n \lg n)$（用合併排
序或堆積排序，比較函式每次 $\Theta(1)$）。 第 6--7 行的 \textbf{while} 迴圈看起來可能
一次彈很多個，但每個點只推入一次，被彈出最多一次，所以整個 \textbf{for} 迴圈裡第 7 行
合計最多執行 $n$ 次，第 6 行的檢查最多 $(n - 2) + n$ 次（每輪最後一次不成立的檢查加上
成立的次數）。 這就是第~3~章的聚合法：單獨一輪可能很貴，總和是 $O(n)$。 全部合起來
\[
T(n) = \Theta(n) + \Theta(n \lg n) + \Theta(n) = \Theta(n \lg n),
\]
時間花在排序上；掃描那一段不管輸入長什麼樣都是 $O(n)$。 如果點已經依極角排好（例如
直接給一個凸多邊形的頂點順序），整個只要 $\Theta(n)$。

\begin{examnote}
寫 Graham 掃描法的追蹤，先寫出 $p_0$ 和排序後的順序，再一列一列寫「處理哪一點、彈出
哪些點、堆疊變成什麼」。 有三個細節各書不同，作答時先寫你用哪一種：$p_0$ 同高時取最
左還是最右；極角相同時近的在前還是遠的在前（有的書只留最遠的那一個）；共線的點要不要
輸出（本書第 6 行用 $\le 0$，不輸出）。 題目問複雜度時，答 $O(n \lg n)$，並說明排序
是瓶頸、掃描是 $O(n)$（每點最多進出堆疊各一次）。
\end{examnote}

\subsection{C 程式}

程式照虛擬碼~\ref{alg:dc-graham} 寫，點的名字用一個字元。 \texttt{qsort} 的比較函式只
收兩個指標，拿不到 $p_0$，所以把 $p_0$ 放在一個檔案範圍的變數裡，排序前設好。 堆疊直接
用輸出陣列 \texttt{hull}，\texttt{h} 是堆疊裡的點數。 點以打亂的順序給，每處理一個點印
出彈掉的點（括號裡是外積）和堆疊。

\begin{example}[label=ex:dc-graham]{c}{(graham\_scan.c) Graham 掃描法}
#include <stdio.h>
#include <stdlib.h>

typedef struct {
    int x, y;
    char name;
} Point;

// 外積 (q - p) x (r - p)：正是左轉，負是右轉，0 是共線。座標的差先轉成 long long。
long long orient(Point p, Point q, Point r)
{
    return ((long long)q.x - p.x) * ((long long)r.y - p.y)
         - ((long long)q.y - p.y) * ((long long)r.x - p.x);
}

long long dist2(Point p, Point q)
{
    long long dx = (long long)p.x - q.x, dy = (long long)p.y - q.y;
    return dx * dx + dy * dy;
}

Point p0;                               // qsort 的比較函式拿不到參數，基準點放在這裡

// 繞 p0 的極角由小到大；極角相同時近的在前。
int cmp_angle(const void *pa, const void *pb)
{
    const Point *a = pa, *b = pb;
    long long o = orient(p0, *a, *b);
    if (o != 0) {
        return o > 0 ? -1 : 1;          // p0->a 轉到 p0->b 是左轉：a 的極角比較小
    }
    long long da = dist2(p0, *a), db = dist2(p0, *b);
    return (da > db) - (da < db);
}

// 回傳凸包的頂點數，頂點依逆時針順序放在 hull[0..h-1]。要求 n >= 2。
int graham_scan(Point *p, int n, Point *hull)
{
    int low = 0;
    for (int i = 1; i < n; i++) {       // y 最小的點，y 相同取 x 最小
        if (p[i].y < p[low].y || (p[i].y == p[low].y && p[i].x < p[low].x)) {
            low = i;
        }
    }
    Point t = p[0];
    p[0] = p[low];
    p[low] = t;
    p0 = p[0];
    qsort(p + 1, n - 1, sizeof *p, cmp_angle);
    printf("sorted:");
    for (int i = 0; i < n; i++) {
        printf(" %c", p[i].name);
    }
    printf("\n");
    int h = 0;                          // hull[0..h-1] 當堆疊用，hull[h-1] 是頂端
    hull[h++] = p[0];
    hull[h++] = p[1];
    for (int i = 2; i < n; i++) {
        printf("%c:", p[i].name);
        while (h >= 2 && orient(hull[h - 2], hull[h - 1], p[i]) <= 0) {
            printf(" pop %c(%lld)", hull[h - 1].name,
                   orient(hull[h - 2], hull[h - 1], p[i]));
            h--;
        }
        hull[h++] = p[i];
        printf(" ->");
        for (int k = 0; k < h; k++) {
            printf(" %c", hull[k].name);
        }
        printf("\n");
    }
    return h;
}

int main(void)
{
    Point p[] = {{9, 5, 'C'}, {2, 9, 'H'}, {11, 3, 'B'}, {4, 3, 'M'}, {10, 6, 'E'},
                 {6, 1, 'A'}, {8, 7, 'K'}, {1, 2, 'J'}, {13, 8, 'D'}, {4, 7, 'G'},
                 {11, 9, 'L'}, {3, 4, 'I'}, {7, 11, 'F'}};
    int n = 13;
    Point *hull = malloc(n * sizeof *hull);
    if (hull == NULL) {
        return 1;
    }
    int h = graham_scan(p, n, hull);
    printf("hull (%d):", h);
    for (int k = 0; k < h; k++) {
        printf(" %c(%d,%d)", hull[k].name, hull[k].x, hull[k].y);
    }
    printf("\n");
    free(hull);
    return 0;
}
\end{example}
\iosample{}{sorted: A B D E C L K F G H M I J \\ D: -> A B D \\ E: -> A B D E \\ C: -> A B D E C \\ L: pop C(-2) pop E(-7) -> A B D L \\ K: -> A B D L K \\ F: pop K(-14) pop L(0) -> A B D F \\ G: -> A B D F G \\ H: pop G(-14) -> A B D F H \\ M: -> A B D F H M \\ I: pop M(-4) -> A B D F H I \\ J: pop I(-12) -> A B D F H J \\ hull (6): A(6,1) B(11,3) D(13,8) F(7,11) H(2,9) J(1,2)}

輸出和上面的表逐列一致。 \texttt{qsort} 不保證相同的鍵維持原順序，但這裡排序的鍵（極
角，同方向再比距離）不會有兩個點完全相同，所以結果是唯一的。

\subsection{分治的凸包}
\label{sec:dc-hull-dc}

凸包也可以照本章的分治樣板做（Preparata 和 Hong，1977）。 先把點依 $x$ 排序（$x$ 相
同時依 $y$），然後：

\begin{enumerate}
\item \textbf{分}：依排好的順序切成左右兩半 $L$、$R$，各有一半的點，$L$ 的點都在
$R$ 的左邊。
\item \textbf{治}：遞迴求出 $\mathrm{CH}(L)$、$\mathrm{CH}(R)$，各自是逆時針的頂點序
列。 三個點以內直接做。
\item \textbf{合}：找兩條\textbf{橋}（bridge）：上橋是連接 $\mathrm{CH}(L)$ 和
$\mathrm{CH}(R)$ 的線段裡，讓所有點都在它下方（或線上）的那一條；下橋讓所有點都在它
上方。 新的凸包就是：沿 $\mathrm{CH}(L)$ 的外側從上橋的左端走到下橋的左端，走過下
橋，沿 $\mathrm{CH}(R)$ 的外側走到上橋的右端，再走過上橋回來。 兩個小凸包夾在兩條
橋之間、朝內的那幾條邊丟掉。
\end{enumerate}

\begin{center}
\begin{minipage}{\linewidth}
\centering
\begin{tikzpicture}[x=0.55cm, y=0.5cm]
\draw[very thick, dashed, memStroke] (7,0) -- (7,10);
\fill[coreFill, fill opacity=0.6] (3,1) -- (6,3) -- (5,6) -- (2,8) -- (1,4) -- cycle;
\fill[coreFill, fill opacity=0.6] (9,1) -- (13,2) -- (12,5) -- (10,9) -- (8,7) -- cycle;
\draw[thick, coreStroke] (3,1) -- (6,3) -- (5,6) -- (2,8) -- (1,4) -- cycle;
\draw[thick, coreStroke] (9,1) -- (13,2) -- (12,5) -- (10,9) -- (8,7) -- cycle;
\draw[very thick, optStroke] (2,8) -- (10,9) node[midway, above, font=\scriptsize] {上橋};
\draw[very thick, optStroke] (3,1) -- (9,1) node[midway, below, font=\scriptsize] {下橋};
\foreach \x/\y in {3/1, 6/3, 5/6, 2/8, 1/4, 4/4, 9/1, 13/2, 12/5, 10/9, 8/7, 10/4} {
    \node[dc-pt] at (\x,\y) {};
}
\node[font=\small] at (3.5,10.4) {$L$};
\node[font=\small] at (10.5,10.4) {$R$};
\end{tikzpicture}
\diagcaption{分治凸包的「合」。 兩個藍色多邊形是左右兩半各自的凸包，紅色是上橋和下
橋。 合併後的凸包沿左邊凸包的左側、下橋、右邊凸包的右側、上橋繞一圈；兩個凸包朝
內的邊（虛線兩側的那幾條）丟掉。}
\label{fig:dc-hull-merge}
\end{minipage}
\end{center}

找上橋的方法像兩個人往上爬：一開始 $a$ 是 $\mathrm{CH}(L)$ 最右邊的點，$b$ 是
$\mathrm{CH}(R)$ 最左邊的點。 只要 $b$ 在 $\mathrm{CH}(R)$ 上往上走一步（順時針方
向的下一個頂點）會落在有向直線 $a \to b$ 的左邊（上方），就讓 $b$ 走過去；同樣，只要
$a$ 在 $\mathrm{CH}(L)$ 上往上走一步（逆時針方向的下一個頂點）在直線 $a \to b$ 的
上方，就讓 $a$ 走過去。 兩邊都走不動時，$ab$ 就是上橋。 每一步都是一次
\proc{Orient}，$a$、$b$ 各自只往一個方向走，最多走過各自凸包的所有頂點，所以找橋是
$O(n)$；下橋對稱。 圖~\ref{fig:dc-hull-merge} 裡，上橋從 $a = (6, 3)$、$b = (8, 7)$
開始，$a$ 往上走到 $(5, 6)$、$(2, 8)$，$b$ 走到 $(10, 9)$ 停下。 合併時把兩個頂點序
列接起來也是 $O(n)$。 所以
\[
T(n) = 2T(n/2) + \Theta(n) = \Theta(n \lg n),
\]
加上一開始排序的 $\Theta(n \lg n)$，合計 $\Theta(n \lg n)$，和 Graham 掃描法同一級。
它的遞迴結構和第~\ref{sec:dc-closest}~節的最近點對一樣：先依 $x$ 排好，切一半，難處都
在「合」。

\subsection{凸包的下界}
\label{sec:dc-hull-lower}

兩種做法都是 $\Theta(n \lg n)$，而且時間都花在排序。 能不能不排序、做得更快？ 下面用
\textbf{歸約}（reduction）說明不行：如果凸包能做得比 $n \lg n$ 快，排序也就能做得比
$n \lg n$ 快，和第~15~章的下界矛盾。

做法是把要排序的數放到拋物線 $y = x^2$ 上。 給 $n$ 個不同的正數 $a_1, \ldots, a_n$，
做出點 $P_i = (a_i, a_i^2)$。 對三個數 $a < b < c$，
\[
\proc{Orient}\bigl((a, a^2), (b, b^2), (c, c^2)\bigr)
= (b - a)(c^2 - a^2) - (b^2 - a^2)(c - a) = (b - a)(c - a)(c - b) > 0 .
\]
所以依 $x$ 由小到大走，每個轉彎都是左轉；從最右點走回最左點，再走向第二個點，也是左
轉（代進同一個式子可以驗證）。 也就是說，依 $x$ 排好的順序圍成一個凸多邊形，每個點都
是轉角：\textbf{$n$ 個點全部是凸包頂點，而且逆時針順序就是 $x$ 由小到大的順序}（從最
左點開始讀）。

\begin{center}
\begin{minipage}{\linewidth}
\centering
\begin{tikzpicture}[x=0.9cm, y=0.13cm]
\draw[-{Stealth[length=2mm]}, black!60] (0,0) -- (5.8,0) node[right, font=\small] {$x$};
\draw[-{Stealth[length=2mm]}, black!60] (0,0) -- (0,27) node[above, font=\small] {$y$};
\draw[black!30, domain=0:5.15, samples=60] plot (\x, {\x*\x});
\fill[coreFill, fill opacity=0.6] (1,1) -- (2,4) -- (3,9) -- (4,16) -- (5,25) -- cycle;
\draw[thick, coreStroke] (1,1) -- (2,4) -- (3,9) -- (4,16) -- (5,25) -- cycle;
\foreach \a in {1,...,5} {
    \draw[black!60] (\a,0) -- (\a,-0.8) node[below, font=\scriptsize] {\a};
}
\foreach \a/\k in {4/1, 1/2, 3/3, 5/4, 2/5} {
    \node[dc-pt, label={[dc-ptlab]right:$P_{\k}$}] at (\a,\a*\a) {};
}
\end{tikzpicture}
\diagcaption{排序 $\langle 4, 1, 3, 5, 2 \rangle$：做出 $P_i = (a_i, a_i^2)$，五個點都
是凸包頂點。 從最左的 $P_2$ 逆時針讀 $P_2, P_5, P_3, P_1, P_4$，就是 $1, 2, 3, 4, 5$。}
\label{fig:dc-parabola}
\end{minipage}
\end{center}

\begin{theorem}
\label{thm:dc-hull-lower}
設一個凸包演算法對座標只做兩種測試：比較兩個座標的大小，以及判斷 \proc{Orient} 的正
負號（本節的演算法都可以只用這兩種測試寫成）。 那麼它要輸出依逆時針排好的凸包頂點，最壞情況至少要
$\Omega(n \lg n)$ 次測試。
\end{theorem}

\begin{proof}
設這個演算法對 $n$ 個點最多做 $T(n)$ 次測試。 用它做一個排序演算法：輸入 $n$ 個不同
的正數 $a_1, \ldots, a_n$，做出 $P_i = (a_i, a_i^2)$，呼叫凸包演算法，在輸出裡找 $x$
最小的頂點（$n - 1$ 次比較），從它開始依序讀出 $x$ 座標。 由上面的計算，讀出來的就是
由小到大的順序。

這個排序演算法只靠比較鍵就能執行：
\begin{itemize}
\item 比較 $x$ 座標就是比較 $a_i$、$a_j$。
\item 比較 $y$ 座標是比較 $a_i^2$、$a_j^2$；鍵都是正數，和比較 $a_i$、$a_j$ 的結果相同。
\item $\proc{Orient}(P_i, P_j, P_k) = (a_j - a_i)(a_k - a_i)(a_k - a_j)$，它的正負號由
$a_i$、$a_j$、$a_k$ 兩兩比較的三個結果決定。
\end{itemize}
所以它是一個比較排序，最多做 $3T(n) + (n - 1)$ 次比較。 第~15~章的定理~\ref{thm:lin-lower}
說，比較排序排 $n$ 個不同的鍵（限制成正數也一樣，$n!$ 種順序都還在），最壞至少要
$\lceil \lg(n!) \rceil$ 次比較。 因此
\[
3T(n) + n - 1 \ge \lg(n!) = \Omega(n \lg n), \qquad T(n) = \Omega(n \lg n). \qedhere
\]
\end{proof}

定理的前提是一個計算模型，就像第~15~章的下界只管比較排序。 允許更多種運算時（例如可
以判斷任何固定次數多項式的正負號），Ben-Or 在 1983 年證明排序仍然要 $\Omega(n \lg n)$，同
樣的歸約就給出凸包的 $\Omega(n \lg n)$。 只要求找出「哪些點是頂點」、不要求順序時，
下界也還是 $\Omega(n \lg n)$（Yao，1981），證明比較難，本書不談。

\begin{examnote}
「證明凸包至少要 $\Omega(n \lg n)$」的標準寫法：（1）給排序的輸入 $a_1, \ldots, a_n$，
在 $O(n)$ 時間內轉成點 $(a_i, a_i^2)$；（2）說明這些點全部在凸包上，而且凸包的逆時針
順序從最左點開始就是排好的順序，輸出轉回答案也是 $O(n)$；（3）所以凸包如果是
$o(n \lg n)$，排序就是 $o(n \lg n)$，和排序的下界矛盾。 歸約的方向不要寫反：是把\textbf{排
序}變成\textbf{凸包}，證明凸包不比排序容易。 有時間的話補一句計算模型（比較和外積），
並說明轉換本身只花 $O(n)$。
\end{examnote}

各種凸包做法的比較（$h$ 是凸包的頂點數）：

\begin{center}
\begin{tabular}{lll}
\toprule
做法 & 時間 & 想法 \\
\midrule
Graham 掃描法 & $\Theta(n \lg n)$ & 依極角排序，用堆疊掃一次 \\
分治 & $\Theta(n \lg n)$ & 依 $x$ 切兩半，找上下兩條橋 \\
禮物包裝法（Jarvis march） & $\Theta(nh)$ & 每一步掃過所有點，找下一個頂點 \\
任何做法的下界 & $\Omega(n \lg n)$ & 排序可以歸約到凸包 \\
\bottomrule
\end{tabular}
\end{center}

禮物包裝法每找一個頂點要看過所有點，$h$ 很小時比 $n \lg n$ 快；但 $h$ 可以到 $n$（例
如點都在一個圓上），最壞是 $\Theta(n^2)$。 這和下界不衝突：下界用的輸入正是 $h = n$
的情形。

\subsection{最遠點對}
\label{sec:dc-farthest}

和最近點對相對的問題是\textbf{最遠點對}：$n$ 個點裡距離最遠的兩點，也就是點集的直徑。
兩兩比較是 $\Theta(n^2)$。 關鍵的觀察是：\textbf{最遠點對的兩點一定都是凸包的頂點}
（觀念題第~\ref{dc-q-far}~題證明）。 所以先花 $O(n \lg n)$ 求凸包，再只在 $h$ 個頂點
之間找。

$h$ 個頂點兩兩比較要 $\Theta(h^2)$，$h = n$ 時沒有省到。 \textbf{旋轉卡尺}
（rotating calipers）的想法是：沿著凸包逆時針走一圈，對每一條邊，找離這
條邊所在直線最遠的頂點（叫它的\textbf{對踵點}）；邊往前走一條時，對踵點也只會往前走，
不會回頭。 所以用兩個指標一起繞一圈，$O(h)$ 就能看完所有對踵的點對，最遠點對一定是其
中一對。 合計 $O(n \lg n)$。

最近點對不能這樣做：最近的兩點通常在點集的內部，和凸包沒有關係，所以第~\ref{sec:dc-closest}~節
要用分治和帶狀區域的論證。 兩個問題的答案都是 $\Theta(n \lg n)$，但理由完全不同。

\section{本章常見錯誤}
\label{sec:dc-pitfalls}

\begin{longtable}{p{0.22\linewidth}p{0.57\linewidth}p{0.13\linewidth}}
\toprule
錯誤 & 例子 & 見 \\
\midrule
\endhead
區間寫法混用 & 條件寫 $\mathit{lo} \le \mathit{hi}$，卻用 $\mathit{hi} \gets
\mathit{mid}$，只剩一格時停不下來 & 第~\ref{sec:dc-traps}~節 \\
$\mathit{lo} \gets \mathit{mid}$ 配地板 & 剩兩格時 $\mathit{mid} = \mathit{lo}$，
無窮迴圈 & 第~\ref{sec:dc-traps}~節 \\
以為會找到第一個 & 有重複時 \proc{Binary-Search} 回傳任一個；要第一個用
\proc{Lower-Bound} & 第~\ref{sec:dc-bsearch}~節 \\
快速冪溢位 & 沒有先把底數取餘數，或 $m$ 太大，\texttt{b * b} 超過 \texttt{long long} &
第~\ref{sec:dc-power}~節 \\
跨中點只延伸一邊 & 跨中點的段沒有同時包含 $A[q]$ 和 $A[q+1]$ &
第~\ref{sec:dc-maxsub}~節 \\
全部是負數答 $0$ & 題目要求非空的段時，答案是最大的那一個負數 &
第~\ref{sec:dc-maxsub}~節 \\
Karatsuba 還是四次乘法 & 分別算出 $ad$、$bc$ 再相加，沒有用 $(a+b)(c+d)$ &
第~\ref{sec:dc-mult}~節 \\
遞迴之後才讀分界 & 遞迴已把 $P$ 依 $y$ 重排，$P[q].x$ 不是左半最大的 $x$ &
第~\ref{sec:dc-closest}~節 \\
相等時先拿右邊 & 把相等的數對算成逆序 & 第~\ref{sec:dc-inv}~節 \\
逆序數用 \texttt{int} & $n = 10^5$ 時逆序數可達約 $5 \times 10^9$ &
第~\ref{sec:dc-inv}~節 \\
用斜率判斷轉向 & 鉛直線沒有斜率、浮點數有誤差；改用外積，只有乘法和減法 &
第~\ref{sec:dc-orient}~節 \\
外積用 \texttt{int} & 座標到 $10^9$ 時兩個乘積各約 $4 \times 10^{18}$；$d_1 d_2$ 更不能真的乘 &
第~\ref{sec:dc-orient}~節 \\
有 $d = 0$ 就說相交 & 四點共線但分開時四個外積都是 $0$，要再檢查 \proc{On-Segment} &
第~\ref{sec:dc-segint}~節 \\
射線通過頂點算錯次數 & 高度等於 $q.y$ 的頂點要一律歸到同一邊（本書算下方） &
第~\ref{sec:dc-pip}~節 \\
Graham 留下共線的點 & 彈出的條件寫成 $< 0$，邊上的點留在輸出裡 & 第~\ref{sec:dc-graham}~節 \\
下界的歸約方向寫反 & 要把排序變成凸包，才能說凸包至少和排序一樣難 &
第~\ref{sec:dc-hull-lower}~節 \\
\bottomrule
\end{longtable}

\section{觀念題}
\label{sec:dc-exercises}

題型和研究所考試、期中期末考的筆試相近，有選擇、計算、追蹤和證明。 解答在本章最後。

\begin{enumerate}
\item \label{dc-q-tf}（選擇，可能不只一個答案）下列哪些敘述正確？
\begin{enumerate}
\item 把 \proc{Binary-Search} 第 4 行改成 $\mathit{mid} \gets \lfloor (\mathit{lo} +
2\,\mathit{hi})/3 \rfloor$，其他不變，最壞情況仍是 $\Theta(\lg n)$。
\item 把兩個 $n$ 位數各切成三段，如果能只用 $6$ 次三分之一長度的乘法，加上
$\Theta(n)$ 的加減完成乘法，這個方法比 Karatsuba 快。
\item 把矩陣切成四塊、直接做 $8$ 次子矩陣乘法，矩陣加法只有 $4$ 次，比 Strassen 的
$18$ 次少，所以總時間也是 $\Theta(n^{\lg 7})$。
\item 把 \proc{Max-Crossing} 換成 $\Theta(n^2)$ 的做法（兩層迴圈枚舉所有跨中點的段，
邊走邊加），\proc{Max-Subarray} 的時間變成 $\Theta(n^2)$。
\item 最近點對如果不讓遞迴依 $y$ 合併，改成每一層用合併排序把帶狀區域依 $y$ 重排，
最壞情況的時間是 $O(n \lg^2 n)$。
\end{enumerate}

\item \label{dc-q-bounds}（追蹤）$A = \langle 3, 8, 8, 8, 8, 14, 20, 20, 27 \rangle$，
$n = 9$。 寫出 \proc{Lower-Bound}$(A, 9, 8)$、\proc{Upper-Bound}$(A, 9, 8)$、
\proc{Lower-Bound}$(A, 9, 21)$ 每一輪的 $\mathit{lo}$、$\mathit{hi}$、
$\mathit{mid}$、$A[\mathit{mid}]$ 和回傳值。 由結果寫出 $8$ 出現幾次，以及最後一個
不大於 $21$ 的數在第幾格。

\item \label{dc-q-fix}（改錯與證明）下面的程序想回傳最後一個滿足 $A[i] \le x$ 的位置，
$A[1..n]$ 已經由小到大排好，而且保證 $A[1] \le x$。
\begin{pseudo}{最後一個不大於 $x$ 的位置（有錯）}
\Function{Last-Le}{$A, n, x$}
    \State $\mathit{lo} \gets 1$；$\mathit{hi} \gets n$
    \While{$\mathit{lo} < \mathit{hi}$}
        \State $\mathit{mid} \gets \lfloor (\mathit{lo} + \mathit{hi}) / 2 \rfloor$
        \If{$A[\mathit{mid}] \le x$}
            \State $\mathit{lo} \gets \mathit{mid}$
        \Else
            \State $\mathit{hi} \gets \mathit{mid} - 1$
        \EndIf
    \EndWhile
    \State \Return $\mathit{lo}$
\EndFunction
\end{pseudo}
\begin{enumerate}
\item 給一組 $A$ 和 $x$，讓這個程序永遠不停。
\item 只改一行讓它正確，寫出迴圈不變量，用三步證明正確，並說明迴圈一定會停。
\item 用 \proc{Upper-Bound} 寫出同樣的結果。
\end{enumerate}

\item \label{dc-q-pow}（計算）
\begin{enumerate}
\item 用 \proc{Power-Iter} 算 $3^{100}$。 寫出每一輪開始時的 $e$、$r$、$b$（$r$、$b$
寫成 $3$ 的次方即可），以及第 5 行、第 8 行各執行了幾次。
\item \proc{Power}$(x, 15)$ 做幾次乘法？ 找出一種只用 $5$ 次乘法算出 $x^{15}$ 的方
法。
\end{enumerate}

\item \label{dc-q-maxsub}（追蹤）$A = \langle -3, 5, -1, 4, -7, 6, -2 \rangle$。
\begin{enumerate}
\item 對 \proc{Max-Subarray}$(A, 1, 7)$ 的最上層，寫出 $q$、跨中點的兩列累加和它們
的最大值、\proc{Max-Crossing} 的結果，以及左半、右半遞迴回傳的段（可以直接觀察）。
最後的答案是哪一段？
\item 用 \proc{Max-Subarray-Linear} 再算一次，寫出每一輪的 $\mathit{cur}$ 和
$\mathit{best}$。
\end{enumerate}

\item \label{dc-q-four}（設計與計算）讓分治版的最大子陣列每次回傳四個值：整段的和
$T$、最佳前綴和 $P$（開頭固定在最左邊的段的最大和）、最佳後綴和 $S$、最大子陣列的
和 $M$。
\begin{enumerate}
\item 已知左半的 $(T_L, P_L, S_L, M_L)$ 和右半的 $(T_R, P_R, S_R, M_R)$，寫出整段的
四個值，並說明理由。
\item 寫出執行時間的遞迴式並解出來。 和第~\ref{sec:dc-maxsub}~節的分治版比，差在
哪裡？
\end{enumerate}

\item \label{dc-q-strassen}（計算）用 Strassen 的七個乘積（第~\ref{sec:dc-strassen}~節
的式子）計算
\[
\begin{pmatrix} 2 & -3 \\ 1 & 4 \end{pmatrix}
\begin{pmatrix} 2 & 0 \\ 5 & -1 \end{pmatrix},
\]
寫出 $M_1$ 到 $M_7$ 的值，組合出 $C$，再用一般的乘法驗算。

\item \label{dc-q-cp}（追蹤）八個點依 $x$ 排好：
$A(1, 1)$、$B(3, 7)$、$C(7, 7)$、$D(9, 12)$、$E(10, 7)$、$F(11, 11)$、$G(13, 3)$、
$H(14, 9)$。 用 \proc{Closest} 處理全部八個點。 寫出最上層的 $x_m$、
$\delta_L^2$、$\delta_R^2$（兩半的最近距離可以直接兩兩比較求得）、帶狀區域裡的點
（依 $y$ 排列）、第 14 行依序算過的點對和每次之後的 $\delta^2$，以及答案。 落在分界
線上的點屬於哪一半？

\item \label{dc-q-inv}（追蹤）$A = \langle 5, 1, 7, 3, 3, 8, 2, 6 \rangle$。
\begin{enumerate}
\item 用 \proc{Sort-Count} 計算逆序數，寫出每一次 \proc{Merge-Count} 的範圍和回傳
值。
\item 寫出最上層合併的每一步（比較、拿哪一邊、$c$ 增加多少）。
\item 如果把 \proc{Merge-Count} 第 6 行的 $B[i] \le B[j]$ 改成 $B[i] < B[j]$，回傳的
總數是多少？ 多算了哪一對？
\end{enumerate}

\item \label{dc-q-seg}（計算）$P(2, 3)$、$Q(7, 5)$、$R(4, 8)$、$S(12, 7)$、$T(9, 1)$。
\begin{enumerate}
\item 算出 $\proc{Orient}(P, Q, R)$ 和 $\proc{Orient}(P, Q, S)$，說明各是左轉、右轉還是
共線。
\item 用 \proc{Segments-Intersect} 判斷三組線段：$\overline{PS}$ 和 $\overline{QR}$；
$\overline{PR}$ 和 $\overline{QS}$；$\overline{RT}$ 和 $\overline{PS}$（每一組前一條
當 $\overline{p_1 p_2}$）。 寫出 $d_1$ 到 $d_4$，以及程序在第幾行回傳什麼。
\item 如果把 \proc{On-Segment} 拿掉，改成「有某個 $d$ 等於 $0$ 就回傳 \TRUE」，(b) 的
哪一組會答錯？
\end{enumerate}

\item \label{dc-q-pip}（追蹤）多邊形的頂點依序是 $v_1(1, 1)$、$v_2(9, 1)$、$v_3(9, 4)$、
$v_4(6, 4)$、$v_5(5, 7)$、$v_6(3, 4)$，要判斷的點是 $q_1(2, 4)$ 和 $q_2(4, 4)$。 兩條射
線都沿著水平邊 $v_3 v_4$ 走，也都通過 $v_4$、$v_6$。
\begin{enumerate}
\item 對兩個點分別執行 \proc{Point-In-Polygon}，列出跨過 $y = 4$ 的邊、
$\proc{Orient}(a, b, q) \cdot (b.y - a.y)$ 的值和最後的計數，並說明答案。
\item 如果把第 5 行改成「$a.y \ge q.y$ 和 $b.y \ge q.y$ 恰好一個成立」（高度等於
$q.y$ 的頂點改算在上方），兩個點的計數各是多少？ 答案還對嗎？
\end{enumerate}

\item \label{dc-q-graham}（追蹤與分析）
\begin{enumerate}
\item 八個點 $A(7, 0)$、$B(3, 0)$、$C(9, 4)$、$D(6, 3)$、$E(5, 8)$、$F(1, 6)$、
$G(0, 2)$、$H(4, 4)$。 用 \proc{Graham-Scan} 求凸包：寫出 $p_0$、排序後的順序、每處理
一個點時彈出哪些點（附外積）和堆疊，以及最後的凸包。
\item 處理一個點時，第 6--7 行可能連續彈出很多個點。 給一組 $8$ 個點的例子，讓最後
一個點進來時彈出 $5$ 個點，並說明為什麼即使如此，整個掃描仍是 $O(n)$。
\end{enumerate}

\item \label{dc-q-far}（證明）$|u|$ 表示向量 $u$ 的長度。
\begin{enumerate}
\item 設 $p = t a + (1 - t) b$，$0 \le t \le 1$（$p$ 在線段 $\overline{ab}$ 上）。 證明
對任何點 $q$，
\[
|p - q|^2 = t\,|a - q|^2 + (1 - t)\,|b - q|^2 - t(1 - t)\,|a - b|^2 .
\]
\item 由 (a) 證明：$a \ne b$、$0 < t < 1$ 時，$|p - q| < \max(|a - q|, |b - q|)$。
\item 證明：點集裡任何一組最遠點對，兩點都是凸包的頂點。
\item 求凸包之後，在 $h$ 個頂點之間兩兩比較，總時間是多少？ $h$ 最多多大時，這仍是
$O(n \lg n)$？
\end{enumerate}

\end{enumerate}

\section{歷屆試題}
\label{sec:dc-pastexam}

以下是研究所入學考試的原題，題目照原卷（語言、記號、配分照錄）。 詳解在本章最後。

\begin{pastexam}[label={px:dc-1}]{臺大資工所 111 學年度 第 10 題}
(5 points) Consider an integer array $A$ of size $n>0$ with elements $A[1],\ldots,A[n]$,
where the first $p$ elements are positive and the other $(n-p)$ elements are negative. The
following algorithm calculates $p$ correctly by calling \textsc{Compute-P}$(A,n)$. What loop
invariant is maintained for the \textbf{while} loop of the algorithm?

\begin{quote}
\textsc{Compute-P}$(A, len)$\\
1\quad \textbf{if} $A[1]<0$ \textbf{return} 0\\
2\quad $left=1, right=len$\\
3\quad \textbf{while} $left<right$\\
4\quad\quad $m=\text{ceiling}(\frac{left+right}{2})$\\
5\quad\quad \textbf{if} $A[m]>0$\\
6\quad\quad\quad $left=m$\\
7\quad\quad \textbf{else}\\
8\quad\quad\quad $right=m-1$\\
9\quad \textbf{return} $left$
\end{quote}
\begin{enumerate}[label=(\Alph*)]
\item $left+1\le p\le right$
\item $left+1\le p\le right-1$
\item $left\le p\le right$
\item $left\le p\le right-1$
\item none of the other choices
\end{enumerate}
\end{pastexam}

\begin{pastexam}[label={px:dc-2}]{中正資工所 111 學年度 第 1 題}
(2\%) The maximum number of comparisons needed for the binary search of a 2,000 element
array is

(a) 1,999 \qquad (b) 9 \qquad (c) 11 \qquad (d) 13
\end{pastexam}

\begin{pastexam}[label={px:dc-3}]{成大資訊聯招 113 學年度 第 11 題}
(10\%) Consider three $n\times n$ matrices, denoted as $A$, $B$, and $C$. Each of these
matrices is partitioned into four $n/2\times n/2$ submatrices. Assuming that $n$ is an exact
power of 2, we can guarantee that, for $n\ge 2$, the dimension $n/2$ is an integer. The
subdivisions for matrix $A$ yield $A_{11}$, $A_{12}$, $A_{21}$, and $A_{22}$. Similarly,
matrix $B$ is decomposed into $B_{11}$, $B_{12}$, $B_{21}$, and $B_{22}$, while matrix $C$
undergoes division into $C_{11}$, $C_{12}$, $C_{21}$, and $C_{22}$. We have the following
procedure:
\par\vspace{1ex}
{\centering
\begin{tabular}{|l|}
\hline
F($A$, $B$)\\
\hspace*{2em}$n=A$.rows\\
\hspace*{2em}if $n==1$\\
\hspace*{4em}$c_{11}=a_{11}\cdot b_{11}$\\
\hspace*{2em}else\\
\hspace*{4em}$C_{11}=\mathrm{F}(A_{11},B_{11})+\mathrm{F}(A_{12},B_{21})$\\
\hspace*{4em}$C_{12}=\mathrm{F}(A_{11},B_{12})+\mathrm{F}(A_{12},B_{22})$\\
\hspace*{4em}$C_{21}=\mathrm{F}(A_{21},B_{11})+\mathrm{F}(A_{22},B_{21})$\\
\hspace*{4em}$C_{22}=\mathrm{F}(A_{21},B_{12})$\\
\hspace*{2em}return $C$\\
\hline
\end{tabular}\par}
\vspace{1ex}
It is important to note that the procedure for $C_{22}$ is invoked only once and '+' is
matrix addition. Let $T(n)$ represent the time required to process two $n\times n$ matrices
using this particular procedure. Provide an asymptotic tight bound ($\Theta$) for $T(n)$,
assuming that $T(n)$ is a constant for sufficiently small $n$.
\end{pastexam}

\begin{pastexam}[label={px:dc-4}]{政大資科所 107 學年度 第 7 題}
(15\%) A contiguous subsequence of a list $S$ is a subsequence made up of consecutive
elements of $S$. For instance, if $S$ is 5, 15, $-30$, 10, $-5$, 40, 10, then 15, $-30$, 10
is a contiguous subsequence but 5, 15, 40 is not. Design a polynomial-time algorithm for the
following task:\\
Input: A list $S$ of numbers $a_1,a_2,\ldots,a_n$.\\
Output: The contiguous subsequence of $S$ of maximum sum.\\
Note that a subsequence of length zero has sum zero.\\
For the preceding example, the answer would be 10, $-5$, 40, 10, with a sum of 55.
\end{pastexam}

\begin{pastexam}[label={px:dc-5}]{臺大資工所 103 學年度 第 5 題（節錄 (a)(c) 小題）}
(30 points) Computational geometry has applications in a number of fields such as computer
graphics, computer vision, robotics, molecular modeling and statistics. We consider the
following problems of computational geometry in 2D space where the points are specified by
their $x$ and $y$ coordinates, $p_i=(x_i,y_i)$.
\begin{enumerate}[label=(\alph*)]
\item (5 points) Give an algorithm to determine whether two line segments intersect. The
inputs are segments $\overline{p_1p_2}$ and $\overline{p_3p_4}$, and the algorithm returns
TRUE if the two segments intersect and FALSE if they do not.
\item[(c)] (15 points) Give an \textit{efficient} algorithm to determine whether a given
point is inside a given polygon. The vertices of the polygon are in counterclockwise order.
We assume that the boundary of the polygon has no self-intersections, and that the given
point is not on the polygon's boundary. Please describe the key ideas and discuss its time
complexity.

Figure 3 illustrates a polygon. Point A is within the polygon and point B is outside of the
polygon.
\end{enumerate}

\begin{center}
\begin{tikzpicture}[font=\sffamily]
\draw[thick] (3.30,-1.18)--(1.70,-1.90)--(2.48,-3.15)--(2.08,-4.35)--(4.90,-4.78)--(6.52,-2.38)--(3.70,-2.73)--cycle;
\fill (1.73,-3.58) circle (1.3pt) node[above=1pt] {B};
\fill (3.52,-3.70) circle (1.3pt) node[above=1pt] {A};
\end{tikzpicture}\\ Figure 3: A Polygon and two points\end{center}
\end{pastexam}

\section{實作練習}
\label{sec:dc-practice}

下面是 Aki-Judge 上和本章相關的題目。 兩題二分搜尋、「模冪函式」和「耗體力的淘汰
賽對戰表」只交函式（標頭檔或 \texttt{.c} 檔），不寫 \texttt{main}；其他題自己讀輸
入。

\begin{center}
\begin{tabular}{llc}
\toprule
小節 & 題目 & 難度 \\
\midrule
二分搜尋 & \akiref{binary-search-index-function}{二分搜尋索引函式} & 暖身 \\
 & \akiref{lower-and-upper-bound-in-a-sorted-array}{已排序陣列的下界與上界} & 基本 \\
快速冪 & \akiref{modular-exponentiation-function}{模冪函式} & 基本 \\
 & \akiref{fibonacci-number-modulo-1000000007}{費氏數列取模} & 進階 \\
最大子陣列 & \akiref{maximum-contiguous-subarray-sum}{最大連續子陣列和} & 基本 \\
 & \akiref{largest-contiguous-sum}{連續區段的最大和} & 基本 \\
最近點對 & \akiref{smallest-gap-between-points-on-a-line}{數線上最近兩點的距離} & 基本 \\
分治的結構 & \akiref{heavier-half-first}{重的一半排前面} & 基本 \\
 & \akiref{stamina-knockout-bracket}{耗體力的淘汰賽對戰表} & 進階 \\
\bottomrule
\end{tabular}
\end{center}

幾題的提示。 兩題二分搜尋的位置從 $0$ 開始，「二分搜尋索引函式」找不到回傳 $-1$，
「已排序陣列的下界與上界」找不到回傳 $n$；交之前用第~\ref{sec:dc-traps}~節列的小例子
自己測。 「模冪函式」的指數到 $10^{12}$，先把底數取餘數，指數是 $0$ 時回傳 $1$。
「費氏數列取模」用第~\ref{sec:dc-power}~節的矩陣快速冪，每次乘法和加法之後都取餘
數。 兩題最大子陣列都要求非空的段，用 \texttt{long long}；第一題用分治版，第二題的
$n$ 到 $10^6$，用線性版。 「數線上最近兩點的距離」是一維的最近點對，座標可以重複。
最後兩題都是「切成兩半、兩半各自處理完、再由兩半的結果決定這一段」，先寫出遞迴式再
動手；「耗體力的淘汰賽對戰表」要用 \texttt{malloc} 建出整棵遞迴樹，上半區是
$\lceil m/2 \rceil$ 人。 逆序數的兩題在第~14~章的題單，可以用 \proc{Sort-Count} 再
做一次。

\section{讀完你應該會}
\label{sec:dc-checklist}

\begin{itemize}
\item 寫出閉區間和半開區間的二分搜尋、\proc{Lower-Bound}、\proc{Upper-Bound}，用迴圈
不變量證明正確，並用它們算出現次數、最後一個不大於 $x$ 的位置等邊界問題；看得出
混用區間寫法、$\mathit{lo} \gets \mathit{mid}$ 配地板造成的無窮迴圈。
\item 寫出遞迴和迴圈兩種快速冪，算出乘法次數 $\lfloor \lg n \rfloor + \nu(n) - 1$，
在取餘數時不溢位，並能把同樣的方法用在矩陣上。
\item 用分治法和線性做法解最大子陣列，寫出跨中點的兩列累加，說明線性做法的
$E[i]$ 遞迴為什麼對。
\item 手算 Karatsuba 的三個乘積和 Strassen 的七個乘積，寫出它們的遞迴式，用主定理
說明為什麼少一次乘法就改變了指數。
\item 說明最近點對的帶狀區域為什麼寬 $2\delta$、每個點為什麼最多往後看 $7$ 個點，
並能手算每一層的 $x_m$、帶狀區域和比較的點對。
\item 用合併排序在 $\Theta(n \lg n)$ 內數逆序數，說明相等時為什麼要先拿左邊。
\item 用外積判斷左轉、右轉、共線；寫出線段相交的四個外積，說明 \proc{On-Segment} 為什
麼需要；用射線法判斷點在不在多邊形裡，處理射線通過頂點和水平邊的情形。
\item 手算 Graham 掃描法的排序和堆疊的每一步，說明它是 $\Theta(n \lg n)$ 而掃描只要
$O(n)$；說出分治凸包怎麼用兩條橋合併；用排序歸約證明凸包的 $\Omega(n \lg n)$ 下界；
知道最遠點對的兩點都在凸包上。
\end{itemize}

\textbf{延伸閱讀}（書末會整理完整的參考資料）

\begin{itemize}
\item Thomas H. Cormen 等，\textit{Introduction to Algorithms}，第 4 版，4.1、4.2 節：
矩陣乘法的分治版和 Strassen 演算法，Strassen 的七個乘積用另一組式子寫；章末問題
2-4 是逆序數。
\item Steven S. Skiena，\textit{The Algorithm Design Manual}，第 3 版，5.1、5.5、
5.6 節：二分搜尋的各種變化、快速乘法、最大子陣列與最近點對，用很短的篇幅講直覺。
\item Jon Kleinberg、Éva Tardos，\textit{Algorithm Design}，第 5 章：逆序數、最近點
對、整數乘法三節，最近點對的證明用了另一種切格子的方法，可以和本章對照。
\item Jon Bentley，\textit{Programming Pearls}，第 2 版，第 4 欄與第 8 欄：前者用二分
搜尋示範怎麼寫出正確的程式，後者從 $\Theta(n^3)$ 一路改進最大子陣列到線性時間。
\item Donald E. Knuth，\textit{The Art of Computer Programming, Vol.~2: Seminumerical
Algorithms}，第 3 版，4.3.3 節（快速乘法）與 4.6.3 節（冪的計算）：後者討論算 $x^n$
最少要幾次乘法，二進位法不一定最少（觀念題第~\ref{dc-q-pow}~題）。
\item Thomas H. Cormen 等，\textit{Introduction to Algorithms}，第 3 版，第 33 章：外積、
線段相交（含掃描線判斷 $n$ 條線段）、Graham 掃描法和禮物包裝法。 第 4 版刪掉了這一章。
\item Franco P. Preparata、Michael Ian Shamos，\textit{Computational Geometry: An
Introduction}，Springer，1985：凸包的各種做法與下界、點集的直徑（最遠點對）、線段相交
的掃描線演算法，舊卷的計算幾何題大多出自這本書的範圍。
\item Mark de Berg、Otfried Cheong、Marc van Kreveld、Mark Overmars，\textit{Computational
Geometry: Algorithms and Applications}，第 3 版，Springer，2008，第 1、2 章：凸包（包括
共線點和數值誤差怎麼處理）和線段相交的掃描線演算法。
\end{itemize}

\addcontentsline{toc}{section}{觀念題解答}
\section*{觀念題解答}

\begin{enumerate}
\item (a)、(d)、(e) 正確。

(a)：$k = \mathit{hi} - \mathit{lo} + 1$ 格時，$\mathit{mid}$ 左邊約 $2k/3$ 格、右邊
約 $k/3$ 格，最壞每輪剩下不超過約 $2k/3$，輪數不超過 $\log_{3/2} n + O(1)$；每輪
至多縮成三分之一，輪數也至少約 $\log_3 n$（$x$ 不在陣列裡時）。 兩者都是
$\Theta(\lg n)$，只差底數。

(b)：$T(n) = 6T(n/3) + \Theta(n)$，$\log_3 6 \approx 1.631 > \lg 3 \approx 1.585$，
比 Karatsuba 慢。 要比 Karatsuba 快，三段的切法最多只能用 $5$ 次乘法
（$\log_3 5 \approx 1.465$）。

(c)：$8T(n/2) + \Theta(n^2) = \Theta(n^3)$。 加法的次數只影響 $f(n)$ 的常數，決定指
數的是乘法的次數。

(d)：$T(n) = 2T(n/2) + \Theta(n^2)$，$n^2 = \Omega(n^{1 + \varepsilon})$，正則條件
$2(n/2)^2 = \frac{1}{2} n^2$，主定理情況 3 給出 $\Theta(n^2)$。

(e)：帶狀區域最多 $n$ 個點，重排 $O(n \lg n)$，$T(n) = 2T(n/2) + O(n \lg n) =
O(n \lg^2 n)$（定理~\ref{thm:rec-master2}，$k = 1$）。

\item 三次搜尋每一輪的值如下表。
\begin{center}
\begin{tabular}{cccccl}
\toprule
& $\mathit{lo}$ & $\mathit{hi}$ & $\mathit{mid}$ & $A[\mathit{mid}]$ & 動作 \\
\midrule
\proc{Lower-Bound}$(8)$ & 1 & 10 & 5 & 8 & $8 \ge 8$，$\mathit{hi} \gets 5$ \\
& 1 & 5 & 3 & 8 & $\mathit{hi} \gets 3$ \\
& 1 & 3 & 2 & 8 & $\mathit{hi} \gets 2$ \\
& 1 & 2 & 1 & 3 & $3 < 8$，$\mathit{lo} \gets 2$；回傳 $2$ \\
\midrule
\proc{Upper-Bound}$(8)$ & 1 & 10 & 5 & 8 & $8 \le 8$，$\mathit{lo} \gets 6$ \\
& 6 & 10 & 8 & 20 & $20 > 8$，$\mathit{hi} \gets 8$ \\
& 6 & 8 & 7 & 20 & $\mathit{hi} \gets 7$ \\
& 6 & 7 & 6 & 14 & $\mathit{hi} \gets 6$；回傳 $6$ \\
\midrule
\proc{Lower-Bound}$(21)$ & 1 & 10 & 5 & 8 & $8 < 21$，$\mathit{lo} \gets 6$ \\
& 6 & 10 & 8 & 20 & $20 < 21$，$\mathit{lo} \gets 9$ \\
& 9 & 10 & 9 & 27 & $27 \ge 21$，$\mathit{hi} \gets 9$；回傳 $9$ \\
\bottomrule
\end{tabular}
\end{center}
$8$ 出現 $6 - 2 = 4$ 次。 陣列裡沒有 $21$，所以「最後一個不大於 $21$」等於「最後一
個小於 $21$」，是 \proc{Lower-Bound}$(21) - 1 = 8$，$A[8] = 20$。 （用
\proc{Upper-Bound}$(21) - 1$ 也一樣是 $8$。）

\item
\begin{enumerate}
\item $A = \langle 4, 9 \rangle$，$x = 9$。 $\mathit{lo} = 1$、$\mathit{hi} = 2$，
$\mathit{mid} = 1$，$A[1] = 4 \le 9$，$\mathit{lo} \gets 1$，什麼都沒變，永遠重複。 只
要範圍剩兩格、左邊那格不大於 $x$，就會發生。
\item 把第 4 行改成 $\mathit{mid} \gets \lceil (\mathit{lo} + \mathit{hi}) / 2
\rceil$。 不變量：每次檢查第 3 行時，$1 \le \mathit{lo} \le \mathit{hi} \le n$，
$A[\mathit{lo}] \le x$，而且 $A[\mathit{hi}+1..n]$ 都大於 $x$。 初始：
$\mathit{lo} = 1$，題目保證 $A[1] \le x$；$A[n+1..n]$ 是空的。 維持：
$\mathit{lo} < \mathit{hi}$ 時 $\mathit{lo} < \mathit{mid} \le \mathit{hi}$。 若
$A[\mathit{mid}] \le x$，新的 $\mathit{lo} = \mathit{mid}$ 仍滿足 $A[\mathit{lo}] \le
x$；否則 $A[\mathit{mid}..n]$ 都大於 $x$，新的 $\mathit{hi} = \mathit{mid} - 1 \ge
\mathit{lo}$，第二段仍成立。 結束：$\mathit{lo} = \mathit{hi}$，$A[\mathit{lo}] \le
x$ 而 $A[\mathit{lo}+1..n]$ 都大於 $x$，$\mathit{lo}$ 正是最後一個不大於 $x$ 的位
置。 會停：每輪 $\mathit{lo}$ 嚴格變大（$\mathit{mid} > \mathit{lo}$）或
$\mathit{hi}$ 嚴格變小，$\mathit{hi} - \mathit{lo}$ 是非負整數，有限輪後變成 $0$。
\item \proc{Upper-Bound}$(A, n, x) - 1$。 這個寫法不需要「$A[1] \le x$」的保證，
全部都大於 $x$ 時回傳 $0$。
\end{enumerate}

\item
\begin{enumerate}
\item $100 = 1100100_2$。
\begin{center}
\begin{tabular}{cccl}
\toprule
$e$ & $r$ & $b$ & 這一輪 \\
\midrule
100 & 1 & $3$ & 偶數；$b \gets 3^2$ \\
50 & 1 & $3^2$ & 偶數；$b \gets 3^4$ \\
25 & 1 & $3^4$ & 奇數：$r \gets 3^4$；$b \gets 3^8$ \\
12 & $3^4$ & $3^8$ & 偶數；$b \gets 3^{16}$ \\
6 & $3^4$ & $3^{16}$ & 偶數；$b \gets 3^{32}$ \\
3 & $3^4$ & $3^{32}$ & 奇數：$r \gets 3^{36}$；$b \gets 3^{64}$ \\
1 & $3^{36}$ & $3^{64}$ & 奇數：$r \gets 3^{100}$；$e$ 變成 $0$，不再平方 \\
\bottomrule
\end{tabular}
\end{center}
第 5 行執行 $\nu(100) = 3$ 次，第 8 行執行 $\lfloor \lg 100 \rfloor = 6$ 次，共 $9$ 次。
\item $\lfloor \lg 15 \rfloor + \nu(15) - 1 = 3 + 3 = 6$ 次（$x^2$、$x^3$、$x^6$、
$x^7$、$x^{14}$、$x^{15}$）。 $5$ 次的方法：$x^2 = x \cdot x$，$x^3 = x^2 \cdot x$，
$x^6 = x^3 \cdot x^3$，$x^{12} = x^6 \cdot x^6$，$x^{15} = x^{12} \cdot x^3$。 二進位
法不一定用最少的乘法。 最少要幾次沒有簡單的公式，但任何方法至少要 $\lceil \lg n
\rceil$ 次（每次乘法最多讓指數加倍），二進位法不超過 $2\lfloor \lg n \rfloor$ 次，
最多差一倍左右。
\end{enumerate}

\item
\begin{enumerate}
\item $q = 4$。 往左累加（從第 4 個往左）：$4, 3, 8, 5$，最大 $8$（第 2 到 4 個）；往
右累加（從第 5 個往右）：$-7, -1, -3$，最大 $-1$（第 5 到 6 個）。
\proc{Max-Crossing} 回傳第 2 到 6 個，和 $7$。 左半 $\langle -3, 5, -1, 4 \rangle$ 的
最佳段是第 2 到 4 個，和 $8$；右半 $\langle -7, 6, -2 \rangle$ 是第 6 個，和 $6$。
答案是第 2 到 4 個，和 $8$，來自左半。 往右累加的最大值是負的，但跨中點的段一定要
包含第 5 個，所以只能取 $-1$。
\item 每一輪結束時的值如下表，答案同樣是第 2 到 4 個，和 $8$。
\begin{center}
\begin{tabular}{cccl}
\toprule
$i$ & $A[i]$ & $\mathit{cur}$ & $\mathit{best}$ \\
\midrule
1 & $-3$ & $-3$ & 1--1：$-3$ \\
2 & 5 & 重新開始：5 & 2--2：5 \\
3 & $-1$ & 4 & 2--2：5 \\
4 & 4 & 8 & 2--4：8 \\
5 & $-7$ & 1 & 2--4：8 \\
6 & 6 & 7 & 2--4：8 \\
7 & $-2$ & 5 & 2--4：8 \\
\bottomrule
\end{tabular}
\end{center}
\end{enumerate}

\item
\begin{enumerate}
\item $T = T_L + T_R$。 $P = \max(P_L,\ T_L + P_R)$：開頭在最左邊的段，要嘛整段在左
半，要嘛包含整個左半再接右半的一個前綴。 $S = \max(S_R,\ T_R + S_L)$，理由對稱。
$M = \max(M_L,\ M_R,\ S_L + P_R)$：最後一項就是引理~\ref{lem:dc-crossing}，左半的最佳
後綴正是「結尾在 $q$ 的最佳段」。 基礎情況一個數 $a$：四個值都是 $a$。
\item 合併只要常數時間，$T(n) = 2T(n/2) + \Theta(1)$，主定理情況 1，$\Theta(n)$。
第~\ref{sec:dc-maxsub}~節的版本每一層從中點重新往兩邊加，花 $\Theta(n)$；這裡讓子問
題把以後會用到的前綴、後綴資訊一起交回來，合併時就不必再看陣列。
\end{enumerate}

\item 七個乘積：
\[
\begin{aligned}
M_1 &= (2 + 4)(2 + (-1)) = 6, & M_2 &= (1 + 4) \times 2 = 10, \\
M_3 &= 2 \times (0 - (-1)) = 2, & M_4 &= 4 \times (5 - 2) = 12, \\
M_5 &= (2 + (-3)) \times (-1) = 1, & M_6 &= (1 - 2)(2 + 0) = -2, \\
M_7 &= (-3 - 4)(5 + (-1)) = -28.
\end{aligned}
\]
$c_{11} = 6 + 12 - 1 + (-28) = -11$，$c_{12} = 2 + 1 = 3$，$c_{21} = 10 + 12 = 22$，
$c_{22} = 6 - 10 + 2 + (-2) = -4$。 驗算：$c_{11} = 2 \cdot 2 + (-3) \cdot 5 = -11$，
$c_{12} = 2 \cdot 0 + (-3)(-1) = 3$，$c_{21} = 1 \cdot 2 + 4 \cdot 5 = 22$，
$c_{22} = 1 \cdot 0 + 4 \cdot (-1) = -4$。

\item 最上層 $q = 4$，$x_m = 9$（D 的 $x$）。 D 落在分界線上，屬於左半（左半是依
$x$ 排序後的前 $4$ 個點）。 左半 A、B、C、D 的最近距離平方是 BC $= 16$，右半 E、F、
G、H 是 FH $= 9 + 4 = 13$，$\delta^2 = 13$。 帶狀區域 $(x - 9)^2 < 13$：C（$4$）、
D（$0$）、E（$1$）、F（$4$），B（$36$）和 G（$16$）不在。 依 $y$ 排：C(7, 7)、
E(10, 7)、F(11, 11)、D(9, 12)，C 和 E 的 $y$ 相同，依 $x$ 排。 比較：
\begin{itemize}
\item C 往後：CE，$y$ 差 $0$，距離平方 $9 < 13$，$\delta^2 \gets 9$；下一個 F 的
$y$ 差 $4$，$16 \ge 9$，停。
\item E 往後：F 的 $y$ 差 $4$，$16 \ge 9$，停，沒有算距離。
\item F 往後：FD，$y$ 差 $1$，距離平方 $4 + 1 = 5 < 9$，$\delta^2 \gets 5$。
\end{itemize}
答案是 F、D，距離 $\sqrt{5}$。 這一層 $\delta$ 在迴圈裡變小了兩次，第一次（CE）找
到的也是跨兩半的點對，但還不是最近的。

\item
\begin{enumerate}
\item 依遞迴的完成順序：$A[1..2]$：$1$（$(5, 1)$）；$A[3..4]$：$1$（$(7, 3)$）；
$A[1..4]$：$\langle 1, 5 \rangle$ 和 $\langle 3, 7 \rangle$，$1$（$(5, 3)$）；
$A[5..6]$：$0$；$A[7..8]$：$0$；$A[5..8]$：$\langle 3, 8 \rangle$ 和
$\langle 2, 6 \rangle$，$3$（$(3, 2)$、$(8, 2)$、$(8, 6)$）；最上層 $6$。 總數
$1 + 1 + 1 + 0 + 0 + 3 + 6 = 12$。
\item 合併 $\langle 1, 3, 5, 7 \rangle$ 和 $\langle 2, 3, 6, 8 \rangle$：
$1$ 對 $2$，拿左 $1$；$3$ 對 $2$，拿右 $2$，左半還沒拿的 $3, 5, 7$，$+3$；$3$ 對 $3$，
相等，拿左 $3$；$5$ 對 $3$，拿右 $3$，還沒拿的 $5, 7$，$+2$；$5$ 對 $6$，拿左 $5$；
$7$ 對 $6$，拿右 $6$，$+1$；$7$ 對 $8$，拿左 $7$；左半用完，拿右 $8$，$+0$。 共
$6$。
\item 改成 $<$ 之後，$3$ 對 $3$ 時先拿右邊的 $3$，左半還沒拿的是 $3, 5, 7$，加 $3$
而不是 $2$；其他層的合併沒有相等的情形。 總數變成 $13$，多算了左半的 $3$ 和右半的
$3$ 這一對，它們相等，不是逆序。
\end{enumerate}

\item
\begin{enumerate}
\item $Q - P = (5, 2)$。 $R - P = (2, 5)$：$5 \cdot 5 - 2 \cdot 2 = 21 > 0$，左轉。
$S - P = (10, 4)$：$5 \cdot 4 - 2 \cdot 10 = 0$，共線（$S - P$ 正好是 $Q - P$ 的兩倍，
$Q$ 在 $P$、$S$ 之間）。
\item 三組的結果：
\begin{center}
\begin{tabular}{lccccl}
\toprule
線段 & $d_1$ & $d_2$ & $d_3$ & $d_4$ & 結果 \\
\midrule
$\overline{PS}$、$\overline{QR}$ & $21$ & $-21$ & $0$ & $42$ & 第 11 行，\TRUE \\
$\overline{PR}$、$\overline{QS}$ & $0$ & $21$ & $-21$ & $-42$ & 第 14 行，\FALSE \\
$\overline{RT}$、$\overline{PS}$ & $42$ & $-48$ & $-39$ & $51$ & 第 5 行，\TRUE \\
\bottomrule
\end{tabular}
\end{center}
第一組：$d_1 d_2 < 0$ 但 $d_3 = 0$，第 4 行不成立；$d_1$、$d_2$ 不是 $0$；$d_3 = 0$ 且
$Q = (7, 5)$ 在以 $P$、$S$ 為對角的長方形裡（$2 \le 7 \le 12$、$3 \le 5 \le 7$），第 11
行回傳 \TRUE：$\overline{QR}$ 的端點 $Q$ 碰到 $\overline{PS}$。 第二組：$d_1 = 0$，
$P$ 在直線 $QS$ 上，但 $P.x = 2$ 不在 $[7, 12]$ 裡，不在線段上；其他三個 $d$ 都不是
$0$，第 4 行也不成立（$d_1 d_2 = 0$），所以第 14 行回傳 \FALSE。 第三組四個外積都不
是 $0$，$d_1 d_2 < 0$、$d_3 d_4 < 0$，互相跨過。
\item 第二組。 $d_1 = 0$ 只表示 $P$ 在直線 $QS$ 上；$P$ 在 $Q$ 的另一側、線段
$\overline{QS}$ 外面，兩條線段其實沒有碰到。
\end{enumerate}

\item
\begin{enumerate}
\item 高度等於 $4$ 的頂點 $v_3, v_4, v_6$ 都算在下方，所以只有 $v_4 v_5$、$v_5 v_6$ 兩條
邊跨過（它們的上端是 $v_5$），水平邊 $v_3 v_4$ 和 $v_2 v_3$、$v_6 v_1$ 都不算。
\begin{center}
\begin{tabular}{lcccc}
\toprule
& 邊 & $\proc{Orient}(a, b, q)$ & $b.y - a.y$ & 乘積 \\
\midrule
$q_1(2, 4)$ & $v_4 v_5$ & $12$ & $3$ & $36$ \\
& $v_5 v_6$ & $-3$ & $-3$ & $9$ \\
\midrule
$q_2(4, 4)$ & $v_4 v_5$ & $6$ & $3$ & $18$ \\
& $v_5 v_6$ & $3$ & $-3$ & $-9$ \\
\bottomrule
\end{tabular}
\end{center}
$q_1$ 兩個乘積都是正的，$c = 2$，在外面：射線在 $v_6$ 穿進多邊形，在 $v_4$ 穿出去，
之後沿著水平邊 $v_3 v_4$ 走，沒有再穿過邊界。 $q_2$ 只有 $v_4 v_5$ 算，$c = 1$，在裡
面（$v_5 v_6$ 和射線的交點是 $v_6$，在 $q_2$ 左邊）。
\item 改成 $\ge$ 之後，$v_3, v_4, v_6$ 都算在上方，跨過的邊變成 $v_2 v_3$ 和 $v_6 v_1$。
$q_1$：$v_2 v_3$ 的乘積 $21 \cdot 3 = 63 > 0$，$v_6 v_1$ 的乘積 $(-3)(-3) = 9 > 0$，
$c = 2$，在外面。 $q_2$：$v_2 v_3$ 的乘積 $15 \cdot 3 = 45 > 0$，$v_6 v_1$ 的乘積
$3 \cdot (-3) = -9$，不算，$c = 1$，在裡面。 數到的邊不一樣，但奇偶相同，答案仍然正
確。 原因是規則只要求「每個剛好在射線高度的頂點，一律歸到同一邊」，歸到上方或下方都
可以。
\end{enumerate}

\item
\begin{enumerate}
\item $A$、$B$ 的 $y$ 都是 $0$，取 $x$ 小的，$p_0 = B$。 排序：$A$（極角 $0^\circ$）、$C$、
$D$、$H$、$E$、$F$、$G$。 $H - B = (1, 4)$ 和 $E - B = (2, 8)$ 方向相同，$H$ 比較近，排在
前面。
\begin{center}
\begin{tabular}{clll}
\toprule
處理 & 彈出（外積） & 堆疊（底 $\to$ 頂） & 說明 \\
\midrule
（開始） & & $B, A$ & \\
$C$ & & $B, A, C$ & $\proc{Orient}(B, A, C) = 16$ \\
$D$ & & $B, A, C, D$ & $\proc{Orient}(A, C, D) = 10$ \\
$H$ & $D$（$-5$） & $B, A, C, H$ & $\proc{Orient}(A, C, H) = 20$ \\
$E$ & $H$（$-20$） & $B, A, C, E$ & $\proc{Orient}(A, C, E) = 24$ \\
$F$ & & $B, A, C, E, F$ & $\proc{Orient}(C, E, F) = 24$ \\
$G$ & & $B, A, C, E, F, G$ & $\proc{Orient}(E, F, G) = 14$ \\
\bottomrule
\end{tabular}
\end{center}
凸包是 $B, A, C, E, F, G$（逆時針）。 $H$ 在線段 $\overline{BE}$ 上，處理 $E$ 時被彈掉；
$D$ 在凸包內部。
\item 例子：$p_0 = (0, 0)$、$(10, 0)$，再把五個點放在一段向外凸的弧
上：$(9, 4)$、$(8, 6)$、$(6, 8)$、$(4, 9)$、$(1, 10)$，最後一點 $(0, 100)$ 放在很高的地
方。 前七個點依極角處理時每一步都是左轉，堆疊長到七個點；$(0, 100)$ 進來時，那五個點
都在三角形 $(0, 0)$、$(10, 0)$、$(0, 100)$ 裡面，依序被彈出（外積分別是 $-269$、
$-178$、$-172$、$-78$、$-60$），凸包只剩三個點。 即使如此，每個點只推入一次、最多彈
出一次，整個掃描的彈出總數不超過 $n$，第 6 行的檢查總數不超過 $2n$。 某一步很貴，是
因為前面幾步很便宜、累積了很多點在堆疊裡，這正是聚合法的論證。
\end{enumerate}

\item
\begin{enumerate}
\item 令 $u = a - q$、$v = b - q$，則 $p - q = t u + (1 - t) v$。 用內積展開：
\[
|p - q|^2 = t^2 |u|^2 + (1 - t)^2 |v|^2 + 2t(1 - t)\, u \cdot v .
\]
右邊的式子是
\[
t|u|^2 + (1 - t)|v|^2 - t(1 - t)\bigl(|u|^2 + |v|^2 - 2 u \cdot v\bigr),
\]
其中 $|a - b|^2 = |u - v|^2 = |u|^2 + |v|^2 - 2u \cdot v$。 $|u|^2$ 的係數是
$t - t(1 - t) = t^2$，$|v|^2$ 的係數是 $(1 - t) - t(1 - t) = (1 - t)^2$，$u \cdot v$
的係數是 $2t(1 - t)$，兩邊相同。
\item $a \ne b$、$0 < t < 1$ 時最後一項 $t(1 - t)|a - b|^2 > 0$，所以
$|p - q|^2 < t|a - q|^2 + (1 - t)|b - q|^2 \le \max(|a - q|, |b - q|)^2$（兩個數的加權平
均不超過較大的那個），開根號即得。
\item 設 $(p, q)$ 是一組最遠點對，假設 $p$ 不是凸包的頂點。 $p$ 在凸包（含內部）裡，
分兩種情形找一條線段 $\overline{cd}$，讓 $p$ 嚴格在 $c$、$d$ 之間，而 $c$、$d$ 都在凸
包的邊界上：$p$ 在某條邊上（不是端點）時，取那條邊；$p$ 在內部時，過 $p$ 任畫一條直
線，它和凸包的交集是一條線段，兩端在邊界上。 由 (b)，$c$、$d$ 中有一個（設為 $c$）
滿足 $|c - q| > |p - q|$。 如果 $c$ 是頂點，令 $w = c$；否則 $c$ 在某條邊
$\overline{uv}$ 的內部，再用一次 (b)，$u$、$v$ 中有一個（設為 $w$）滿足
$|w - q| > |c - q|$。 不論哪一種，$w$ 都是凸包的頂點，也是點集裡的點，而
$|w - q| > |p - q|$，和 $(p, q)$ 最遠矛盾。 所以 $p$ 是頂點；$q$ 同理。
\item $O(n \lg n + h^2)$。 $h^2 = O(n \lg n)$，也就是 $h = O(\sqrt{n \lg n})$ 時，總時間
仍是 $O(n \lg n)$。 $h$ 可以大到 $n$（點都在一個圓上），這時是 $\Theta(n^2)$，所以需要
旋轉卡尺。
\end{enumerate}

\end{enumerate}

\addcontentsline{toc}{section}{歷屆試題詳解}
\section*{歷屆試題詳解}

\pastsol{px:dc-1}{臺大資工所 111 學年度 第 10 題}
答案是 (C)。

第 1 行保證進入迴圈時 $A[1] > 0$，所以 $p \ge 1$。 正數恰好是前 $p$ 個，因此
$A[m] > 0$ 等價於 $m \le p$，$A[m] < 0$ 等價於 $m \ge p + 1$（題目說每個元素不是正就
是負，沒有 $0$）。 照第~\ref{sec:ana-invariant}~節的三步驗證 (C)：

\begin{itemize}
\item \textbf{初始}：$left = 1 \le p$，$p \le n = right$。
\item \textbf{維持}：迴圈裡 $left < right$，$m = \lceil (left + right)/2 \rceil$ 滿足
$left < m \le right$。 若 $A[m] > 0$，則 $m \le p$，新的 $left = m \le p$，$right$ 不變；
若 $A[m] < 0$，則 $m \ge p + 1$，新的 $right = m - 1 \ge p$，$left$ 不變。 兩種情形
$left \le p \le right$ 都仍成立。
\item \textbf{結束}：每一輪 $left$ 嚴格變大（$m > left$）或 $right$ 嚴格變小，迴圈一定會
停，停在 $left = right$，代入不變量得 $p = left$，第 9 行回傳正確。
\end{itemize}

其他選項：

\begin{itemize}
\item (A) $left + 1 \le p$：$p = 1$ 時（例如 $A = \langle 5, -2 \rangle$），一開始
$left = 1$，$left + 1 = 2 > p$，第一次檢查就不成立。 而且迴圈結束時 $left = p$，這個不
等式一定不成立。
\item (B) 比 (A) 多要求 $p \le right - 1$，同樣在 $p = 1$ 時一開始就不成立；結束時
$left = right$，$left + 1 \le right - 1$ 不可能成立。
\item (D) $p \le right - 1$：全部都是正數時（例如 $A = \langle 3, 1 \rangle$，$p = n = 2$），
一開始 $right = 2$，$p \le 1$ 不成立。 結束時 $left = right$，$left \le p \le left - 1$ 也
不可能成立。
\item (E)：(C) 成立，所以不選。
\end{itemize}

用程式把 $n \le 8$ 的所有情形（每個 $n$、每個 $p$）跑過，每一輪檢查四個不等式，只
有 (C) 每次都成立。 這個程序就是觀念題第~\ref{dc-q-fix}~題改好的「最後一個滿足條件的位
置」：往右縮寫成 $left = m$，所以中點一定要取天花板（第~\ref{sec:dc-traps}~節）。 如果第
4 行改成地板，$left = 1$、$right = 2$、$A[1] > 0$ 時 $m = 1$，$left$ 不動，迴圈不會停。

\pastsol{px:dc-2}{中正資工所 111 學年度 第 1 題}
答案是 (c)。

照考卷的慣例，二分搜尋每一輪把 $x$ 和中間那一格比一次（相等、較小、較大三種結果算一
次比較）。 第~\ref{sec:dc-bsearch}~節證明了：還要找的格數 $k$ 每一輪之後最多剩
$\lfloor k/2 \rfloor$，所以迴圈最多跑 $\lfloor \lg n \rfloor + 1$ 輪。
$2^{10} = 1024 \le 2000 < 2048 = 2^{11}$，$\lfloor \lg 2000 \rfloor = 10$，最多 $11$ 次。

這個上界達得到：$x$ 比全部都大時每一輪都往右，格數依序是
$2000 \to 1000 \to 500 \to 250 \to 125 \to 62 \to 31 \to 15 \to 7 \to 3 \to 1 \to 0$，
每個箭頭是一次比較，共 $11$ 次（用程式對 $2000$ 格的陣列試過所有位置和所有落在兩格之間
的值，最多正好 $11$ 次）。

\begin{itemize}
\item (a) $1999$：這是循序搜尋那一級的次數，二分搜尋不需要。
\item (b) $9$：上面的例子做完 $9$ 次比較後還剩 $3$ 格沒找，無法確定答案。
\item (d) $13$：比最壞情況還多。
\end{itemize}

如果把「是否相等」和「是否較小」算成兩次比較，最壞是 $22$ 次，不在選項裡，所以題目
用的是一輪算一次的算法。

\pastsol{px:dc-3}{成大資訊聯招 113 學年度 第 11 題}
$T(n) = \Theta(n^{\lg 7})$，$\lg 7 \approx 2.807$。

\textbf{寫出遞迴式}。 $n = 1$ 時做一次乘法，$\Theta(1)$。 $n \ge 2$ 時：
\begin{itemize}
\item 遞迴呼叫：$C_{11}$、$C_{12}$、$C_{21}$ 各兩次，$C_{22}$ 一次，共 $7$ 次，每次的
矩陣是 $\frac{n}{2} \times \frac{n}{2}$，合計 $7T(n/2)$。
\item 矩陣加法：$C_{11}$、$C_{12}$、$C_{21}$ 各一次，每次是 $\frac{n}{2} \times \frac{n}{2}$
的 $(n/2)^2$ 次加法，合計 $3n^2/4 = \Theta(n^2)$。
\item 切塊和組合 $C$：用下標計算切塊是 $\Theta(1)$，把子矩陣複製出來是 $\Theta(n^2)$；
組合 $C$ 最多 $\Theta(n^2)$。 不論哪一種實作都是 $O(n^2)$，和加法合起來仍是
$\Theta(n^2)$。
\end{itemize}
所以
\[
T(n) = 7T(n/2) + \Theta(n^2).
\]

\textbf{解遞迴式}。 $a = 7$、$b = 2$，$n^{\log_b a} = n^{\lg 7}$。 $\lg 7 \approx 2.807$，取
$\varepsilon = 0.8$，$n^2 = O(n^{\lg 7 - 0.8})$（$\lg 7 - 0.8 \approx 2.007 > 2$），主定理
情況 1（定理~\ref{thm:rec-master}）給出 $T(n) = \Theta(n^{\lg 7})$。

用遞迴樹檢查：第 $i$ 層有 $7^i$ 個子問題，每個的加法成本是 $c(n/2^i)^2$，整層是
$cn^2 (7/4)^i$，是公比 $7/4 > 1$ 的等比數列，總和由最後一層決定；最後一層有
$7^{\lg n} = n^{\lg 7}$ 個葉節點，每個 $\Theta(1)$。 兩種方法一致。

\textbf{說明}。 這個程序其實算錯了：正確的 $C_{22} = A_{21}B_{12} + A_{22}B_{22}$，它少了
$A_{22}B_{22}$。 但題目只問時間。 它的遞迴式和 Strassen 演算法（第~\ref{sec:dc-strassen}~節）
一樣，都是 $7$ 次半大的乘法加上 $\Theta(n^2)$ 的加減，所以時間也一樣；差別是 Strassen
用 $7$ 次乘法算出\textbf{正確}的四塊。 常見的錯誤是看到 $C_{11}$ 到 $C_{21}$ 各兩次就寫
$8T(n/2)$，要數清楚題目特別提醒的 $C_{22}$ 只呼叫一次。

\pastsol{px:dc-4}{政大資科所 107 學年度 第 7 題}
\textbf{解讀}。 題目允許長度 $0$ 的子序列，和是 $0$。 所以全部都是負數時，答案是空的子
序列（和 $0$），這和第~\ref{sec:dc-maxsub}~節「至少一個數」的約定不同。 輸出用位置
$i..j$ 表示 $a_i, \ldots, a_j$，空的子序列寫成 $j = i - 1$。 多段一樣大時，回傳最先達到
最大和的那一段。

\textbf{想法}。 改寫第~\ref{sec:dc-kadane}~節的線性做法。 令
\[
F[j] = \text{結尾在 } j \text{ 的子序列（可以是空的）的最大和},
\]
$F[0] = 0$。 結尾在 $j$ 的非空子序列，是「結尾在 $j - 1$ 的子序列（可以是空的）」再接上
$a_j$，最大是 $F[j-1] + a_j$；再和空的比，$F[j] = \max(0,\ F[j-1] + a_j)$。 答案是所有
「結尾在 $j$ 的非空最佳和」$F[j-1] + a_j$ 和空序列的 $0$ 之中最大的。

\begin{pseudo}[label=alg:dc-px-maxsub]{最大和的連續子序列（允許空的）}
\Function{Max-Contiguous}{$a, n$}
    \State $\mathit{best} \gets 0$；$\mathit{bi} \gets 1$；$\mathit{bj} \gets 0$ \Comment{先記空的子序列}
    \State $\mathit{cur} \gets 0$；$\mathit{start} \gets 1$
    \For{$j \gets 1$ \To $n$}
        \State $\mathit{cur} \gets \mathit{cur} + a_j$
        \If{$\mathit{cur} > \mathit{best}$}
            \State $\mathit{best} \gets \mathit{cur}$；$\mathit{bi} \gets \mathit{start}$；$\mathit{bj} \gets j$
        \EndIf
        \If{$\mathit{cur} < 0$}
            \State $\mathit{cur} \gets 0$；$\mathit{start} \gets j + 1$ \Comment{接上前面只會更差}
        \EndIf
    \EndFor
    \State \Return $(\mathit{bi}, \mathit{bj}, \mathit{best})$
\EndFunction
\end{pseudo}

\textbf{正確性}。 不變量：第 4 行每一輪開始時，(i) $\mathit{cur} = F[j-1] =
a_{\mathit{start}} + \cdots + a_{j-1}$（$\mathit{start} = j$ 時是空的、和為 $0$）；
(ii) $(\mathit{bi}, \mathit{bj}, \mathit{best})$ 是 $a_1, \ldots, a_{j-1}$ 裡（可以是空
的）和最大的連續子序列。
\begin{itemize}
\item \textbf{初始}：$j = 1$，$\mathit{cur} = 0 = F[0]$，只有空的子序列，(i)(ii) 都成立。
\item \textbf{維持}：第 5 行之後 $\mathit{cur} = F[j-1] + a_j = a_{\mathit{start}} +
\cdots + a_j$，是結尾在 $j$ 的非空子序列的最大和。 $a_1..a_j$ 裡的子序列，不是在
$a_1..a_{j-1}$ 裡，就是結尾在 $j$ 的非空子序列，所以第 6--7 行之後 (ii) 對 $j$ 成立。
第 8--9 行把 $\mathit{cur}$ 變成 $\max(0, \mathit{cur}) = F[j]$：小於 $0$ 時換成空的子
序列，從 $j + 1$ 重新開始。 (i) 對下一輪成立。
\item \textbf{結束}：$j = n + 1$，由 (ii)，回傳的是整個 $S$ 裡和最大的連續子序列。
\end{itemize}

\textbf{時間}。 迴圈 $n$ 輪，每輪常數個動作，$\Theta(n)$；額外空間 $\Theta(1)$。 題目只要
求多項式時間，兩兩枚舉所有 $i \le j$ 邊走邊加的 $\Theta(n^2)$ 或第~\ref{sec:dc-maxsub}~節的
分治 $\Theta(n \lg n)$ 也都符合，但線性做法最好，而且不可能更快：每個數都可能是答案的一
部分，至少要讀過一次。

\textbf{追蹤題目的例子}。 $S = 5, 15, -30, 10, -5, 40, 10$。 表中「第 5 行後」是加上
$a_j$ 之後的 $\mathit{cur}$，「本輪結束」是第 8--9 行之後的值。
\begin{center}
\begin{tabular}{ccccl}
\toprule
$j$ & $a_j$ & 第 5 行後 & 本輪結束 & $(\mathit{bi}, \mathit{bj}, \mathit{best})$ \\
\midrule
1 & $5$ & $5$ & $5$ & $(1, 1, 5)$ \\
2 & $15$ & $20$ & $20$ & $(1, 2, 20)$ \\
3 & $-30$ & $-10$ & $0$，$\mathit{start} = 4$ & $(1, 2, 20)$ \\
4 & $10$ & $10$ & $10$ & $(1, 2, 20)$ \\
5 & $-5$ & $5$ & $5$ & $(1, 2, 20)$ \\
6 & $40$ & $45$ & $45$ & $(4, 6, 45)$ \\
7 & $10$ & $55$ & $55$ & $(4, 7, 55)$ \\
\bottomrule
\end{tabular}
\end{center}
回傳 $(4, 7, 55)$，也就是 $10, -5, 40, 10$，和 $55$，和題目給的答案相同。 也用程式把這個程序和兩兩枚舉在兩萬組隨機的小輸入上比對過，最大和都相同。 全部是負數時，例如
$-3, -1$，$\mathit{best}$ 一直是 $0$，回傳 $(1, 0, 0)$，即空的子序列。

\pastsol{px:dc-5}{臺大資工所 103 學年度 第 5 題（節錄 (a)(c) 小題）}
\textbf{(a)} 用外積判斷轉向，演算法就是第~\ref{sec:dc-segint}~節的
\proc{Segments-Intersect}（虛擬碼~\ref{alg:dc-segint}）：
\begin{enumerate}[label=\arabic*.]
\item 令 $\proc{Orient}(p, q, r) = (q.x - p.x)(r.y - p.y) - (q.y - p.y)(r.x - p.x)$，正、
負、零分別表示 $p \to q \to r$ 左轉、右轉、共線（第~\ref{sec:dc-orient}~節）。
\item 算 $d_1 = \proc{Orient}(p_3, p_4, p_1)$、$d_2 = \proc{Orient}(p_3, p_4, p_2)$、
$d_3 = \proc{Orient}(p_1, p_2, p_3)$、$d_4 = \proc{Orient}(p_1, p_2, p_4)$。
\item 如果 $d_1$、$d_2$ 異號而且 $d_3$、$d_4$ 異號，兩條線段互相跨過對方所在的直線，
回傳 TRUE。
\item 否則，如果某個 $d_k = 0$，而對應的端點落在另一條線段上（座標在另一條線段兩端點
圍成的長方形裡，即 \proc{On-Segment}），回傳 TRUE。 這涵蓋端點碰到另一條線段、以及四
點共線而且有重疊的情形。
\item 其他情形回傳 FALSE。
\end{enumerate}
正確性就是引理~\ref{lem:dc-segint}：第 3 步成立時兩條直線交於一點，而這一點同時在兩
條線段上；反過來，兩條線段有共同點而且沒有端點落在另一條線段上時，四個外積都不是 $0$，
而且兩兩異號。 第 4 步不能省：四點共線但分開時（圖~\ref{fig:dc-segcases} (d)），四個外積
都是 $0$，只有座標範圍的檢查能分辨。 只做 $4$ 次外積和常數次比較，$O(1)$；全部用整數
的乘法和減法，沒有除法，不會遇到斜率不存在或浮點誤差。

\textbf{(c)} 用射線法（第~\ref{sec:dc-pip}~節，虛擬碼~\ref{alg:dc-pip}）。

\textbf{想法}。 從點 $q$ 往右畫一條水平射線，數它和多邊形的邊相交幾次；奇數次在裡面，
偶數次在外面。 射線往右走到夠遠的地方一定在多邊形外面，每穿過一次邊界裡外就交換一次，
而多邊形沒有自我相交，邊界把平面分成裡外兩塊，所以奇偶決定答案。

\textbf{細節}。 逐一檢查每條邊 $(v_i, v_{i+1})$（最後一條是 $(v_n, v_1)$）：
\begin{itemize}
\item 只有兩端點一個嚴格在 $q.y$ 上方、另一個不在上方時才算這條邊跨過水平線
$y = q.y$。 高度剛好等於 $q.y$ 的頂點一律當成在下方，所以射線通過頂點時，穿過去的頂點
只算一次，只碰一下的頂點算 $0$ 次或 $2$ 次，水平邊不算，奇偶都正確。
\item 跨過的邊再判斷交點是否在 $q$ 的右邊：交點 $X$ 和 $q$ 同高，
$\proc{Orient}(v_i, v_{i+1}, q) = (v_{i+1}.y - v_i.y)(X.x - q.x)$，所以只要看
$\proc{Orient}(v_i, v_{i+1}, q)$ 和 $v_{i+1}.y - v_i.y$ 是否同號，不必做除法。 $q$ 不在
邊界上，所以這個乘積不會是 $0$。
\end{itemize}
題目說頂點依逆時針排列，射線法用不到這個條件，順時針也一樣。

\textbf{圖 3 的例子}。 依圖上的位置，從 $A$ 往右的射線只穿過右下方那條往右上的邊
（交點約在 $x = 5.6$），$1$ 次，奇數，$A$ 在裡面。 從 $B$ 往右的射線先穿過左邊往左下的
那條邊（交點約在 $x = 2.3$，在 $B$ 的右邊），再穿過同一條右下方的邊，$2$ 次，偶數，$B$
在外面。 這個多邊形不是凸的（左邊和上方各有一個凹進去的頂點），所以不能用下面凸多邊形
的二分搜尋。

\textbf{時間}。 每條邊做常數次運算，$n$ 個頂點的多邊形是 $\Theta(n)$，額外空間
$\Theta(1)$。 對一般的簡單多邊形，沒有事先處理時這已經是最好的：任何一條沒看過的邊都可
能改變答案，至少要讀過所有的邊。 如果多邊形是凸的，可以從 $v_1$ 把多邊形切成扇形，用二
分搜尋找 $q$ 落在哪個扇形，再做一次外積，$O(\lg n)$；同一個多邊形要回答很多次查詢時，
也可以先花時間建資料結構，讓每次查詢變快，這已經超出題目的範圍。

\chapter{動態規劃}
\label{ch:dp}

% 本章的圖共用的樣式（名稱都加 dp- 前綴，避免和別章撞名）
\tikzset{
    dp-call/.style={rounded corners=2pt, draw=black!60, fill=white, minimum width=0.85cm, minimum height=0.5cm, inner sep=1.5pt, font=\small},
    dp-look/.style={dp-call, draw=memStroke, thick, fill=memFill},
    dp-gone/.style={dp-call, draw=black!25, text=black!35, fill=black!4, dashed},
    dp-base/.style={dp-call, fill=black!6},
    dp-leaf/.style={dp-call, draw=coreStroke, fill=coreFill},
    dp-e/.style={draw=black!55},
    dp-gonee/.style={draw=black!25, dashed},
    dp-arr/.style={-{Stealth[length=1.7mm]}, draw=black!55},
    dp-hd/.style={font=\small, text=black!70},
    dp-c/.style={draw=black!40, minimum width=0.95cm, minimum height=0.62cm, inner sep=0pt, font=\small},
    dp-c0/.style={dp-c, fill=black!5, text=black!55},
    dp-cd0/.style={dp-c},
    dp-cd1/.style={dp-c, fill=coreFill},
    dp-cpath/.style={dp-c, draw=memStroke, thick, fill=memFill},
    dp-chit/.style={dp-c, draw=optStroke, thick, fill=optFill},
    dp-key/.style={rounded corners=2pt, draw=black!65, fill=white, minimum height=0.5cm, inner sep=2pt, font=\small},
    dp-ext/.style={rectangle, draw=black!45, fill=black!6, minimum size=0.42cm, inner sep=1pt, font=\scriptsize},
}

系上的程式作業允許同學在截止前重交。 助教手上有同一位同學交的第一版和第二版，想知
道第二版到底改了多少：兩個版本裡有多少行可以原封不動地對上，對不上的才算改過。 從
頭一行一行往下對很容易對錯位，第二版只要在前面插了一行，後面每一行就都對不上了；可
是把「哪些行要對上」的選法全部試一遍，三百行的程式要試多少種？ 往下讀之前，先想一
想：如果已經知道兩個版本「前面一段」最多能對上幾行，能不能拿這個答案去算更長的一
段？

這一章講\textbf{動態規劃}（dynamic programming，常簡稱 DP）。 它和第~17~章的分治一樣，
把問題拆成較小的子問題；不同的是，分治拆出來的子問題互不相干，動態規劃處理的是子問題
\textbf{會重複出現}的情形。 做法是每個子問題只解一次，把答案記在一張表裡，之後要用時直
接查表。 名稱裡的 programming 是二十世紀中期「規劃、排表」的意思，和寫程式無關。

這一章先用一個計數的例子看出重複的子問題，比較由上而下的記憶化和由下而上的填表，再整
理一個最佳化問題要能用動態規劃的條件。 接著五節是經典問題：矩陣鏈乘、最長共同子序列、
最佳二元搜尋樹、0/1 背包，以及編輯距離與最長遞增子序列，每一個都照考卷的寫法把表格填完、
再把最佳解還原出來。 最後用動態規劃的眼光看所有點對最短路徑，完整的內容在第~25~章。

這一章用到第~\ref{sec:ana-input-size}~節的擬多項式時間、第~2~章的遞迴樹、第~7~章的卡特蘭數
與延伸二元樹、第~9~章的二元搜尋樹，以及第~\ref{sec:dc-bsearch}~節的 \proc{Lower-Bound}。 C 程式用到《C
程式設計基礎》第~9~章把二維表攤平成一維陣列的寫法、第~15~章的記憶化和第~19~章的
\texttt{malloc}、\texttt{free}，用到時會先在框裡複習。

\section{動態規劃的要素}
\label{sec:dp-elements}

\subsection{重疊的子問題：數呼叫次數}
\label{sec:dp-overlap}

從 $n$ 個人裡選 $k$ 個人，方法數是二項式係數 $\binom{n}{k}$。 把第 $n$ 個人分開看：選
他的話，剩下要從 $n - 1$ 人裡選 $k - 1$ 人；不選他的話，要從 $n - 1$ 人裡選 $k$ 人。
這就是巴斯卡公式
\[
\binom{n}{k} = \binom{n-1}{k-1} + \binom{n-1}{k}, \qquad
\binom{n}{0} = \binom{n}{n} = 1 .
\]
照公式直接寫成遞迴：

\begin{pseudo}[label=alg:dp-plain-binom]{照巴斯卡公式直接遞迴}
\Function{Plain-Binom}{$n, k$}
    \If{$k = 0$ \OR $k = n$}
        \State \Return $1$
    \EndIf
    \State \Return $\Call{Plain-Binom}{n - 1, k - 1} + \Call{Plain-Binom}{n - 1, k}$
\EndFunction
\end{pseudo}

下圖是 $\proc{Plain-Binom}(5, 2)$ 的遞迴樹，每個節點寫的是那一次呼叫的參數 $n, k$。
$\binom{5}{2} = 10$，樹上剛好有 $10$ 個回傳 $1$ 的葉節點，因為答案就是這些 $1$ 加起來
的。 同一組參數出現了好幾次：$(3, 1)$ 出現 $2$ 次，$(2, 1)$ 出現 $3$ 次，每一次都從頭
重算一遍。

\begin{center}
\begin{minipage}{\linewidth}
\centering
\begin{tikzpicture}[x=1.32cm, y=1.05cm]
\node[dp-call] (b0) at (3.69,0) {$5,2$};
\node[dp-call] (b1) at (0.88,-1) {$4,1$};
\node[dp-call] (b2) at (0.00,-2) {$3,0$};
\node[dp-call] (b3) at (1.75,-2) {$3,1$};
\node[dp-call] (b4) at (1.00,-3) {$2,0$};
\node[dp-call] (b5) at (2.50,-3) {$2,1$};
\node[dp-call] (b6) at (2.00,-4) {$1,0$};
\node[dp-call] (b7) at (3.00,-4) {$1,1$};
\node[dp-call] (b8) at (6.50,-1) {$4,2$};
\node[dp-look] (b9) at (4.75,-2) {$3,1$};
\node[dp-gone] (b10) at (4.00,-3) {$2,0$};
\node[dp-gone] (b11) at (5.50,-3) {$2,1$};
\node[dp-gone] (b12) at (5.00,-4) {$1,0$};
\node[dp-gone] (b13) at (6.00,-4) {$1,1$};
\node[dp-call] (b14) at (8.25,-2) {$3,2$};
\node[dp-look] (b15) at (7.50,-3) {$2,1$};
\node[dp-gone] (b16) at (7.00,-4) {$1,0$};
\node[dp-gone] (b17) at (8.00,-4) {$1,1$};
\node[dp-call] (b18) at (9.00,-3) {$2,2$};
\draw[dp-e] (b0) -- (b1);
\draw[dp-e] (b1) -- (b2);
\draw[dp-e] (b1) -- (b3);
\draw[dp-e] (b3) -- (b4);
\draw[dp-e] (b3) -- (b5);
\draw[dp-e] (b5) -- (b6);
\draw[dp-e] (b5) -- (b7);
\draw[dp-e] (b0) -- (b8);
\draw[dp-e] (b8) -- (b9);
\draw[dp-gonee] (b9) -- (b10);
\draw[dp-gonee] (b9) -- (b11);
\draw[dp-gonee] (b11) -- (b12);
\draw[dp-gonee] (b11) -- (b13);
\draw[dp-e] (b8) -- (b14);
\draw[dp-e] (b14) -- (b15);
\draw[dp-gonee] (b15) -- (b16);
\draw[dp-gonee] (b15) -- (b17);
\draw[dp-e] (b14) -- (b18);
\end{tikzpicture}
\diagcaption{$\proc{Plain-Binom}(5, 2)$ 的遞迴樹，共 $19$ 次呼叫。 依「先左後右」的順序執
行時，橘色的兩個節點是「這組參數之前算過了」的呼叫；虛線灰色的 $6$ 個節點是它們底下的
重算。 下一小節的記憶化會在橘色節點直接查表，灰色的呼叫都不會發生。}
\label{fig:dp-binom-tree}
\end{minipage}
\end{center}

\begin{theorem}
\label{thm:dp-plain-calls}
$0 \le k \le n$ 時，$\proc{Plain-Binom}(n, k)$ 一共被呼叫 $2\binom{n}{k} - 1$ 次。
\end{theorem}

\begin{proof}
令呼叫次數為 $T(n, k)$，對 $n$ 用歸納法。 $k = 0$ 或 $k = n$ 時只有一次呼叫，
$T = 1 = 2 \cdot 1 - 1$。 其他情況，這一次呼叫加上兩個子呼叫：
\[
T(n, k) = 1 + T(n-1, k-1) + T(n-1, k)
= 1 + \left(2\tbinom{n-1}{k-1} - 1\right) + \left(2\tbinom{n-1}{k} - 1\right)
= 2\tbinom{n}{k} - 1 .
\]
最後一步用了巴斯卡公式。
\end{proof}

這個次數是指數級的：$\binom{2m}{m}$ 是 $\binom{2m}{0}, \binom{2m}{1}, \ldots, \binom{2m}{2m}$
這 $2m + 1$ 個數裡最大的，而它們加起來是 $2^{2m} = 4^m$，所以 $\binom{2m}{m} \ge 4^m / (2m + 1)$。 可是會出現的參數組合 $(i, j)$ 只有大約
$n \cdot k$ 種。 呼叫次數和「不同的子問題個數」差這麼多，表示同一個子問題被算了很多次。
這種情形叫\textbf{重疊的子問題}（overlapping subproblems）。 第~14~章的合併排序不一樣：
每次遞迴處理的都是陣列裡新的一段，沒有一段會被排兩次，所以記表格也省不了什麼。

第~2~章的鋪地磚程式 \texttt{tile(n)} 也是同一種情形，那裡算過呼叫次數是 $\Theta(2^n)$，
而不同的參數只有 $n + 1$ 個。

\subsection{由上而下：記憶化}
\label{sec:dp-memo}

最直接的修法是在遞迴外面加一張表：第一次算出 $\binom{i}{j}$ 時寫進表裡，之後再被問到
就直接回傳。 這叫\textbf{記憶化}（memoization），《C 程式設計基礎》第~15~章 15.9 節用費氏
數列示範過。 表 $\mathit{memo}[0..n, 0..k]$ 一開始每一格都是 \NIL。

\begin{pseudo}[label=alg:dp-memo-binom]{記憶化的二項式係數}
\Function{Memo-Binom}{$n, k$}
    \If{$k = 0$ \OR $k = n$}
        \State \Return $1$
    \EndIf
    \If{$\mathit{memo}[n, k] = \NIL$}
        \State $\mathit{memo}[n, k] \gets \Call{Memo-Binom}{n - 1, k - 1} + \Call{Memo-Binom}{n - 1, k}$
    \EndIf
    \State \Return $\mathit{memo}[n, k]$
\EndFunction
\end{pseudo}

要算幾次？ 從 $(n, k)$ 往下走，每一步不是 $(i, j) \to (i-1, j-1)$（「選」）就是
$(i, j) \to (i-1, j)$（「不選」）。 「選」最多走 $k$ 次、「不選」最多走 $n - k$ 次，而
一個子問題由「已經走了幾次選、幾次不選」決定。 會走到的子問題裡，不是邊界的（$j \ge 1$
而且 $j < i$）正是「選」走了不到 $k$ 次、「不選」走了不到 $n - k$ 次的那些，共
$k(n - k)$ 個。 每個這樣的子問題在第一次被問到時展開一次，發出 $2$ 個呼叫，其他呼叫都在
第 3 行或第 6 行直接回傳，所以
\[
\proc{Memo-Binom}(n, k) \text{ 的呼叫次數} = 1 + 2k(n - k) .
\]
$n = 5$、$k = 2$ 時是 $13$，正好是圖~\ref{fig:dp-binom-tree} 裡白色和橘色節點的個數。
時間從指數變成 $O(nk)$。

把每個不同的子問題畫成一個節點，從每個子問題畫箭頭指向它直接用到的子問題，得到
\textbf{子問題圖}（subproblem graph）。 遞迴樹把同一個子問題畫很多次，子問題圖只畫一
次。 記憶化的呼叫次數等於 $1$ 加上子問題圖的邊數。

\begin{center}
\begin{minipage}{\linewidth}
\centering
\begin{tikzpicture}[x=1.45cm, y=1.0cm]
\node[dp-base] (g10) at (2.00,-4) {$1,0$};
\node[dp-base] (g11) at (3.00,-4) {$1,1$};
\node[dp-base] (g20) at (1.50,-3) {$2,0$};
\node[dp-call] (g21) at (2.50,-3) {$2,1$};
\node[dp-base] (g22) at (3.50,-3) {$2,2$};
\node[dp-base] (g30) at (1.00,-2) {$3,0$};
\node[dp-call] (g31) at (2.00,-2) {$3,1$};
\node[dp-call] (g32) at (3.00,-2) {$3,2$};
\node[dp-call] (g41) at (1.50,-1) {$4,1$};
\node[dp-call] (g42) at (2.50,-1) {$4,2$};
\node[dp-call] (g52) at (2.00,0) {$5,2$};
\draw[dp-arr] (g21) -- (g10);
\draw[dp-arr] (g21) -- (g11);
\draw[dp-arr] (g31) -- (g20);
\draw[dp-arr] (g31) -- (g21);
\draw[dp-arr] (g32) -- (g21);
\draw[dp-arr] (g32) -- (g22);
\draw[dp-arr] (g41) -- (g30);
\draw[dp-arr] (g41) -- (g31);
\draw[dp-arr] (g42) -- (g31);
\draw[dp-arr] (g42) -- (g32);
\draw[dp-arr] (g52) -- (g41);
\draw[dp-arr] (g52) -- (g42);
\end{tikzpicture}
\diagcaption{$\binom{5}{2}$ 的子問題圖：$11$ 個不同的子問題，灰底的是邊界（直接回傳
$1$），$6$ 個白色的各有兩條邊，共 $12$ 條邊。 它沒有環，是一個有向無環圖。}
\label{fig:dp-binom-dag}
\end{minipage}
\end{center}

\subsection{由下而上：填表}
\label{sec:dp-bottomup}

另一種做法完全不遞迴：先想清楚子問題之間誰依賴誰，再從最小的子問題開始，依序把表填滿。
$\binom{i}{j}$ 只用到第 $i - 1$ 列，所以一列一列由上往下填就好。

\begin{pseudo}[label=alg:dp-binom-table]{由下而上填巴斯卡三角形}
\Function{Binom-Table}{$n, k$}
    \For{$i \gets 0$ \To $n$}
        \For{$j \gets 0$ \To $\min(i, k)$}
            \If{$j = 0$ \OR $j = i$}
                \State $C[i, j] \gets 1$
            \Else
                \State $C[i, j] \gets C[i - 1, j - 1] + C[i - 1, j]$
            \EndIf
        \EndFor
    \EndFor
    \State \Return $C[n, k]$
\EndFunction
\end{pseudo}

第 7 行執行時，第 $i - 1$ 列已經填好了，這是由下而上的關鍵：\textbf{填某一格之前，它用到
的格子都要先填好}。 用子問題圖來說，就是箭頭指到的節點要排在前面，也就是把子問題圖做
\textbf{拓樸排序}（第~22~章）再倒過來。 這一章的問題，順序都可以直接看出來：依列號、依
區間長度、依前綴長度。

第 7 行只用到上一列，所以整張表可以換成一列 $\mathit{row}[0..k]$，每一列就地更新。 要注意
更新的方向：新的 $\mathit{row}[j]$ 要用\textbf{舊的} $\mathit{row}[j - 1]$，所以 $j$ 要由
大往小跑，否則 $\mathit{row}[j - 1]$ 已經先被改成新的值了。 第~\ref{sec:dp-knapsack}~節的
0/1 背包會再遇到同一件事。

看實際的數字：$\binom{24}{12} = 2{,}704{,}156$，\proc{Plain-Binom} 要呼叫
$2\binom{24}{12} - 1 = 5{,}408{,}311$ 次，\proc{Memo-Binom} 只要 $1 + 2 \cdot 12 \cdot 12 = 289$
次。 \proc{Binom-Table} 和一列滾動的版本，第 $i$ 列都做 $\min(i - 1, k)$ 次加法，共 $nk - k(k+1)/2 = 210$ 次；它也算了一
些用不到的格子，例如算 $\binom{24}{12}$ 根本用不到 $\binom{20}{1}$，但沒有遞迴，常數很小。

\subsection{最佳子結構}
\label{sec:dp-substructure}

上面是計數問題。 動態規劃更常用在\textbf{最佳化問題}（optimization problem）：每個可行解
有一個值，要找值最小（或最大）的解。 這時除了子問題要重疊，還需要第二個條件。

\begin{definition}[最佳子結構]
\label{def:dp-substructure}
一個問題有\textbf{最佳子結構}（optimal substructure），是說它的最佳解裡面包含的子問題
的解，也都是那些子問題的最佳解。
\end{definition}

有了這個性質，才能寫出「最佳值 $=$ 在幾種選擇中，取（子問題的最佳值 $+$ 這次選擇的成本）
最小的一個」這種遞迴式。 證明一個問題有最佳子結構，幾乎都用同一個手法，叫做
\textbf{剪下貼上}（cut-and-paste）：假設最佳解裡的某個子問題的解不是最佳的，就把它剪掉、
貼上那個子問題真正的最佳解，得到一個更好的整體解，和「原本是最佳解」矛盾。

\textbf{有最佳子結構的例子：最短路徑。} 在一張圖上，如果 $p$ 是從 $u$ 到 $v$ 邊數最少的
路徑，而 $w$ 是 $p$ 上的一個點，那麼 $p$ 從 $u$ 到 $w$ 的那一段 $p_1$ 也是 $u$ 到 $w$ 邊
數最少的路徑。 理由：如果有更短的 $p_1'$，把 $p_1$ 換成 $p_1'$，得到一條從 $u$ 到 $v$
的走法，比 $p$ 短。 這條走法可能經過同一個點兩次，但把中間繞回來的一圈去掉只會更短，
最後得到一條比 $p$ 短的路徑，矛盾。

\textbf{沒有最佳子結構的例子：最長簡單路徑。} 簡單路徑是不重複經過任何點的路徑。 看一個
三角形：三個點 $a, b, c$ 兩兩相連。 從 $a$ 到 $c$ 最長的簡單路徑是 $a \to b \to c$，有
$2$ 條邊。 它從 $a$ 到 $b$ 的那一段只有 $1$ 條邊，可是從 $a$ 到 $b$ 最長的簡單路徑是
$a \to c \to b$，有 $2$ 條邊。 最佳解的一段不是那一段的最佳解。 反過來把兩段的最佳解接起
來，$a \to c \to b$ 接上 $b \to a \to c$，會重複經過 $a$ 和 $c$，根本不是簡單路徑。

差別在於子問題是否\textbf{互相獨立}。 最短路徑的兩段各自用到的點不會互相限制（真的重複
了也能去掉）；最長簡單路徑的兩段卻會搶同一批點，一段用了某個點，另一段就不能用。 剪下
貼上的論證在這裡失效。 最長簡單路徑目前沒有已知的多項式時間演算法，第~28~章會講它和
NP 完備問題的關係。

\subsection{設計動態規劃的步驟}
\label{sec:dp-steps}

一個最佳化問題同時有最佳子結構和重疊的子問題時，可以照下面四步設計演算法：

\begin{enumerate}
\item 描述最佳解的結構：最佳解的「最後一步」或「第一步」有哪幾種選擇，做了某個選擇之後，
剩下的是哪個子問題。
\item 遞迴定義最佳值：寫出子問題的參數、遞迴式和邊界。
\item 計算最佳值：記憶化，或找出填表順序由下而上算。
\item 還原最佳解：填表時順便記下每一格是哪個選擇得到的，最後從答案那一格倒著走回去。
\end{enumerate}

只要最佳值的話，第 4 步可以省略。 執行時間通常是
\[
\text{子問題的個數} \times \text{每個子問題要比較的選擇數},
\]
再乘上每個選擇的成本（通常是常數）。 例如下一節的矩陣鏈乘有 $\Theta(n^2)$ 個子問題，每個
最多 $n - 1$ 種選擇，時間是 $O(n^3)$。

\begin{examnote}
考卷上的動態規劃題，答案通常要包含五件事，順序也照這樣寫：(1) 用一句完整的話定義表格的
一格，例如「$c[i, j]$ 是 $X$ 的前 $i$ 個字元和 $Y$ 的前 $j$ 個字元的最長共同子序列長
度」；(2) 遞迴式，連同邊界條件；(3) 填表的順序；(4) 答案在哪一格；(5) 時間和空間複雜度。
只寫出遞迴式、沒有寫那一格代表什麼，閱卷的人很難判斷你的遞迴式對不對。 填表題要把整張
表寫完，平手時照題目規定的規則選（本章各節都會寫明本書的規則）。
\end{examnote}

\subsection{由上而下還是由下而上}

兩種寫法算的是同一張表，漸近時間通常一樣。 記憶化照遞迴式直接寫，只算真的用得到的子問題，
但有函式呼叫的成本，遞迴也可能很深；由下而上要先想清楚填表順序，通常整張表都填，可是只有迴
圈、常數小，而且常常可以只留一兩列（滾動陣列）來省空間。
用得到的子問題只佔整張表的一小部分時，記憶化比較快，觀念題第~\ref{dp-q-sparse}~題有一個
例子。 遞迴深度可能到 $10^5$ 以上時（例如很長的字串），C 的呼叫堆疊可能不夠，要
寫成由下而上。 考卷上的填表題，幾乎都是由下而上的寫法。

\section{矩陣鏈乘}
\label{sec:dp-matrix}

\subsection{問題}

一個 $p \times q$ 的矩陣乘上一個 $q \times r$ 的矩陣，結果是 $p \times r$；用一般的三層迴圈
算，每一格要做 $q$ 次乘法，共 $pqr$ 次純量乘法。 這一節只數純量乘法的次數。

要算 $n$ 個矩陣的連乘 $A_1 A_2 \cdots A_n$，其中 $A_i$ 是 $p_{i-1} \times p_i$（相鄰兩個的
大小要接得上），所以整條鏈的大小由一個序列 $\langle p_0, p_1, \ldots, p_n \rangle$ 決定。
矩陣乘法有結合律，先乘哪兩個都得到同樣的結果，但次數可能差很多。 這一節的例子是 $6$ 個
矩陣，
\[
\langle p_0, \ldots, p_6 \rangle = \langle 6, 10, 3, 12, 5, 8, 4 \rangle ,
\]
也就是 $A_1$ 是 $6 \times 10$、$A_2$ 是 $10 \times 3$、…、$A_6$ 是 $8 \times 4$。 只看前三個：
\begin{itemize}
\item $(A_1 A_2) A_3$：先算 $6 \cdot 10 \cdot 3 = 180$ 次得到 $6 \times 3$，再乘 $A_3$ 要
$6 \cdot 3 \cdot 12 = 216$ 次，共 $396$ 次。
\item $A_1 (A_2 A_3)$：先算 $10 \cdot 3 \cdot 12 = 360$ 次得到 $10 \times 12$，再算
$6 \cdot 10 \cdot 12 = 720$ 次，共 $1080$ 次。
\end{itemize}
六個全部由左往右乘要 $1188$ 次，最差的加法要 $2472$ 次，等一下會算出最少是 $648$ 次。

\textbf{矩陣鏈乘問題}（matrix-chain multiplication problem）：給定 $\langle p_0, \ldots, p_n
\rangle$，找出一種加滿括號的方式，讓純量乘法的總次數最少。 注意輸出是\textbf{乘的順序}，
不是乘出來的矩陣。

括號的加法有幾種？ 第~7~章算過，$n + 1$ 個數連乘、加滿括號的方法有 $C_n$ 種（卡特蘭數），
所以 $n$ 個矩陣有 $C_{n-1}$ 種。 $C_{n-1} = \frac{1}{n}\binom{2n-2}{n-1}$ 是
$\Omega(4^n / n^{3/2})$，$n = 6$ 已經有 $42$ 種，$n = 20$ 有超過十億種，全部試一遍不可行。

\subsection{最佳子結構與遞迴式}

把 $A_i A_{i+1} \cdots A_j$ 簡寫成 $A_{i..j}$，它是 $p_{i-1} \times p_j$ 的矩陣。 $i < j$ 時，
任何一種加括號的方式，最後一次乘法一定是
\[
A_{i..j} = (A_{i..k}) \cdot (A_{k+1..j}), \qquad i \le k < j ,
\]
最後這一次的成本是 $p_{i-1} p_k p_j$。 如果最佳的乘法在 $k$ 分開，那麼它算 $A_{i..k}$ 的方
式一定也是 $A_{i..k}$ 最省的方式：假如有更省的，把那一段換掉，整體就更省了，這就是剪下
貼上。 右邊 $A_{k+1..j}$ 也一樣。 兩邊用到的矩陣不重疊，互不影響。

令 $m[i, j]$ 為算出 $A_{i..j}$ 最少需要的純量乘法次數。 我們不知道最佳的 $k$ 是哪一個，就每
一個都試：
\[
m[i, j] = \begin{cases}
0, & i = j, \\[2pt]
\displaystyle\min_{i \le k < j} \bigl\{\, m[i, k] + m[k+1, j] + p_{i-1} p_k p_j \,\bigr\}, & i < j .
\end{cases}
\]
另外用 $s[i, j]$ 記下取到最小值的那個 $k$，還原括號時要用。 整個問題的答案是 $m[1, n]$。

子問題是所有的 $1 \le i \le j \le n$，共 $\binom{n}{2} + n = \Theta(n^2)$ 個，而且大量重疊：
例如 $m[3, 5]$ 會被 $m[1, 5]$、$m[2, 5]$、$m[3, 6]$ 都用到。 直接照遞迴式遞迴會是指數時
間（觀念題第~\ref{dp-q-mc-rec}~題）。

\subsection{依鏈長填表}

$m[i, j]$ 用到的 $m[i, k]$ 和 $m[k+1, j]$ 都是比較短的鏈，所以依\textbf{鏈長}
$\ell = j - i + 1$ 由短到長填：先填 $\ell = 2$ 的所有格子，再填 $\ell = 3$，……，最後是
$\ell = n$ 的 $m[1, n]$。 畫成表格的話，$\ell$ 固定的格子排成一條和主對角線平行的斜線，所以
這種填法也叫「一條對角線一條對角線地填」。

\begin{pseudo}[label=alg:dp-mc-order]{矩陣鏈乘的填表}
\Function{Matrix-Chain-Order}{$p, n$}
    \For{$i \gets 1$ \To $n$}
        \State $m[i, i] \gets 0$
    \EndFor
    \For{$\ell \gets 2$ \To $n$} \Comment{鏈長}
        \For{$i \gets 1$ \To $n - \ell + 1$}
            \State $j \gets i + \ell - 1$
            \State $m[i, j] \gets \infty$
            \For{$k \gets i$ \To $j - 1$}
                \State $q \gets m[i, k] + m[k + 1, j] + p_{i-1} p_k p_j$
                \If{$q < m[i, j]$}
                    \State $m[i, j] \gets q$
                    \State $s[i, j] \gets k$
                \EndIf
            \EndFor
        \EndFor
    \EndFor
    \State \Return $m$ 和 $s$
\EndFunction
\end{pseudo}

第 10 行用的是嚴格的 $<$：幾個 $k$ 的成本一樣時，$s[i, j]$ 保留最先試到、也就是最小的
那個 $k$。 這是本書的平手規則。

\textbf{時間}：第 9 行對每一組 $(i, j)$ 執行 $j - i = \ell - 1$ 次，所以總次數是
\[
\sum_{\ell=2}^{n} (n - \ell + 1)(\ell - 1) = \frac{n^3 - n}{6} = \Theta(n^3),
\]
空間是 $m$、$s$ 兩張表，$\Theta(n^2)$。 （和式可以令 $t = \ell - 1$ 寫成
$\sum_{t=1}^{n-1} (n - t)t$ 再用 $\sum t$、$\sum t^2$ 的公式算。）

\subsection{手算追蹤}
\label{sec:dp-mc-trace}

用 $\langle 6, 10, 3, 12, 5, 8, 4 \rangle$ 跑一次。 下表一列是一格，寫出每一個 $k$ 的三項
$m[i, k] + m[k+1, j] + p_{i-1}p_k p_j$，粗體是最小的那個，橫線分開不同的鏈長。

{\small
\begin{center}
\begin{tabular}{c p{0.72\linewidth} r c}
\toprule
$[i, j]$ & 每個 $k$ 的成本 & $m$ & $s$ \\
\midrule
$[1,2]$ & $k=1$：$0+0+6\cdot10\cdot3=\mathbf{180}$ & 180 & 1 \\
$[2,3]$ & $k=2$：$0+0+10\cdot3\cdot12=\mathbf{360}$ & 360 & 2 \\
$[3,4]$ & $k=3$：$0+0+3\cdot12\cdot5=\mathbf{180}$ & 180 & 3 \\
$[4,5]$ & $k=4$：$0+0+12\cdot5\cdot8=\mathbf{480}$ & 480 & 4 \\
$[5,6]$ & $k=5$：$0+0+5\cdot8\cdot4=\mathbf{160}$ & 160 & 5 \\
\midrule
$[1,3]$ & $k=1$：$0+360+6\cdot10\cdot12=1080$；$k=2$：$180+0+6\cdot3\cdot12=\mathbf{396}$ & 396 & 2 \\
$[2,4]$ & $k=2$：$0+180+10\cdot3\cdot5=\mathbf{330}$；$k=3$：$360+0+10\cdot12\cdot5=960$ & 330 & 2 \\
$[3,5]$ & $k=3$：$0+480+3\cdot12\cdot8=768$；$k=4$：$180+0+3\cdot5\cdot8=\mathbf{300}$ & 300 & 4 \\
$[4,6]$ & $k=4$：$0+160+12\cdot5\cdot4=\mathbf{400}$；$k=5$：$480+0+12\cdot8\cdot4=864$ & 400 & 4 \\
\midrule
$[1,4]$ & $k=1$：$0+330+6\cdot10\cdot5=630$；$k=2$：$180+180+6\cdot3\cdot5=\mathbf{450}$；$k=3$：$396+0+6\cdot12\cdot5=756$ & 450 & 2 \\
$[2,5]$ & $k=2$：$0+300+10\cdot3\cdot8=\mathbf{540}$；$k=3$：$360+480+10\cdot12\cdot8=1800$；$k=4$：$330+0+10\cdot5\cdot8=730$ & 540 & 2 \\
$[3,6]$ & $k=3$：$0+400+3\cdot12\cdot4=544$；$k=4$：$180+160+3\cdot5\cdot4=400$；$k=5$：$300+0+3\cdot8\cdot4=\mathbf{396}$ & 396 & 5 \\
\midrule
$[1,5]$ & $k=1$：$0+540+6\cdot10\cdot8=1020$；$k=2$：$180+300+6\cdot3\cdot8=\mathbf{624}$；$k=3$：$396+480+6\cdot12\cdot8=1452$；$k=4$：$450+0+6\cdot5\cdot8=690$ & 624 & 2 \\
$[2,6]$ & $k=2$：$0+396+10\cdot3\cdot4=\mathbf{516}$；$k=3$：$360+400+10\cdot12\cdot4=1240$；$k=4$：$330+160+10\cdot5\cdot4=690$；$k=5$：$540+0+10\cdot8\cdot4=860$ & 516 & 2 \\
\midrule
$[1,6]$ & $k=1$：$0+516+6\cdot10\cdot4=756$；$k=2$：$180+396+6\cdot3\cdot4=\mathbf{648}$；$k=3$：$396+400+6\cdot12\cdot4=1084$；$k=4$：$450+160+6\cdot5\cdot4=730$；$k=5$：$624+0+6\cdot8\cdot4=816$ & 648 & 2 \\
\bottomrule
\end{tabular}
\end{center}
}

整理成兩張表，列是 $i$、欄是 $j$，只用到 $i \le j$ 的上三角。 底色深淺交錯的是不同的對角線
（同一條對角線上的格子鏈長相同）；橘框的是等一下還原括號時會用到的格子。

\begin{center}
\begin{minipage}{\linewidth}
\centering
\begin{tikzpicture}[x=1.05cm, y=0.62cm, baseline=(current bounding box.north)]
\node[font=\small] at (3.50,1.15) {$m[i,j]$};
\node[dp-hd] at (1,0.35) {$j=1$};
\node[dp-hd] at (2,0.35) {$2$};
\node[dp-hd] at (3,0.35) {$3$};
\node[dp-hd] at (4,0.35) {$4$};
\node[dp-hd] at (5,0.35) {$5$};
\node[dp-hd] at (6,0.35) {$6$};
\node[dp-hd, anchor=east] at (0.45,-0.5) {$i=1$};
\node[dp-cd0, minimum width=1.05cm] at (1,-0.5) {0};
\node[dp-cpath, minimum width=1.05cm] at (2,-0.5) {180};
\node[dp-cd0, minimum width=1.05cm] at (3,-0.5) {396};
\node[dp-cd1, minimum width=1.05cm] at (4,-0.5) {450};
\node[dp-cd0, minimum width=1.05cm] at (5,-0.5) {624};
\node[dp-cpath, minimum width=1.05cm] at (6,-0.5) {648};
\node[dp-hd, anchor=east] at (0.45,-1.5) {$2$};
\node[dp-cd0, minimum width=1.05cm] at (2,-1.5) {0};
\node[dp-cd1, minimum width=1.05cm] at (3,-1.5) {360};
\node[dp-cd0, minimum width=1.05cm] at (4,-1.5) {330};
\node[dp-cd1, minimum width=1.05cm] at (5,-1.5) {540};
\node[dp-cd0, minimum width=1.05cm] at (6,-1.5) {516};
\node[dp-hd, anchor=east] at (0.45,-2.5) {$3$};
\node[dp-cd0, minimum width=1.05cm] at (3,-2.5) {0};
\node[dp-cpath, minimum width=1.05cm] at (4,-2.5) {180};
\node[dp-cpath, minimum width=1.05cm] at (5,-2.5) {300};
\node[dp-cpath, minimum width=1.05cm] at (6,-2.5) {396};
\node[dp-hd, anchor=east] at (0.45,-3.5) {$4$};
\node[dp-cd0, minimum width=1.05cm] at (4,-3.5) {0};
\node[dp-cd1, minimum width=1.05cm] at (5,-3.5) {480};
\node[dp-cd0, minimum width=1.05cm] at (6,-3.5) {400};
\node[dp-hd, anchor=east] at (0.45,-4.5) {$5$};
\node[dp-cd0, minimum width=1.05cm] at (5,-4.5) {0};
\node[dp-cd1, minimum width=1.05cm] at (6,-4.5) {160};
\node[dp-hd, anchor=east] at (0.45,-5.5) {$6$};
\node[dp-cd0, minimum width=1.05cm] at (6,-5.5) {0};
\end{tikzpicture}\hspace{0.8cm}\begin{tikzpicture}[x=0.8cm, y=0.62cm, baseline=(current bounding box.north)]
\node[font=\small] at (3.00,1.15) {$s[i,j]$};
\node[dp-hd] at (1,0.35) {$j=2$};
\node[dp-hd] at (2,0.35) {$3$};
\node[dp-hd] at (3,0.35) {$4$};
\node[dp-hd] at (4,0.35) {$5$};
\node[dp-hd] at (5,0.35) {$6$};
\node[dp-hd, anchor=east] at (0.45,-0.5) {$i=1$};
\node[dp-cpath, minimum width=0.8cm] at (1,-0.5) {1};
\node[dp-cd0, minimum width=0.8cm] at (2,-0.5) {2};
\node[dp-cd1, minimum width=0.8cm] at (3,-0.5) {2};
\node[dp-cd0, minimum width=0.8cm] at (4,-0.5) {2};
\node[dp-cpath, minimum width=0.8cm] at (5,-0.5) {2};
\node[dp-hd, anchor=east] at (0.45,-1.5) {$2$};
\node[dp-cd1, minimum width=0.8cm] at (2,-1.5) {2};
\node[dp-cd0, minimum width=0.8cm] at (3,-1.5) {2};
\node[dp-cd1, minimum width=0.8cm] at (4,-1.5) {2};
\node[dp-cd0, minimum width=0.8cm] at (5,-1.5) {2};
\node[dp-hd, anchor=east] at (0.45,-2.5) {$3$};
\node[dp-cpath, minimum width=0.8cm] at (3,-2.5) {3};
\node[dp-cpath, minimum width=0.8cm] at (4,-2.5) {4};
\node[dp-cpath, minimum width=0.8cm] at (5,-2.5) {5};
\node[dp-hd, anchor=east] at (0.45,-3.5) {$4$};
\node[dp-cd1, minimum width=0.8cm] at (4,-3.5) {4};
\node[dp-cd0, minimum width=0.8cm] at (5,-3.5) {4};
\node[dp-hd, anchor=east] at (0.45,-4.5) {$5$};
\node[dp-cd1, minimum width=0.8cm] at (5,-4.5) {5};
\end{tikzpicture}
\diagcaption{$\langle 6, 10, 3, 12, 5, 8, 4 \rangle$ 的 $m$ 表和 $s$ 表。 $m[1, 6] = 648$ 是答
案。 橘框是 \proc{Print-Parens} 走過的格子：$[1,6]$、$[1,2]$、$[3,6]$、$[3,5]$、$[3,4]$。}
\label{fig:dp-mc-tables}
\end{minipage}
\end{center}

從表裡還可以看出幾件事：
\begin{itemize}
\item $m[3, 6]$ 的三個候選是 $544$、$400$、$396$，最後一個只贏了 $4$。 鏈長 $4$ 以上的格子，
候選的差距常常很小，手算要特別小心。
\item $s[1, j] = 2$ 對所有 $j \ge 3$ 都成立：$p_2 = 3$ 是整個序列最小的維度，在它這裡切開，
最後一次乘法的 $p_{i-1}p_k p_j$ 就很小。 這只是這組數字的現象，不能拿來當規則：例如
$\langle 2, 3, 4, 5 \rangle$ 最小的中間維度是 $p_1 = 3$，在那裡切開的 $A_1(A_2 A_3)$ 要
$60 + 30 = 90$ 次，$(A_1 A_2) A_3$ 只要 $24 + 40 = 64$ 次。 子問題還沒算完之前，看不出哪個
$k$ 最好，這正是要把每個 $k$ 都試一遍的原因。
\item 鏈長 $2$ 的格子只有一個候選，所以 $s[i, i+1] = i$。
\end{itemize}

\subsection{還原括號}

$s[i, j] = k$ 表示 $A_{i..j}$ 最後一次乘法在 $A_k$ 和 $A_{k+1}$ 之間，所以從 $s[1, n]$ 開
始遞迴地往兩邊拆：

\begin{pseudo}[label=alg:dp-mc-parens]{印出最佳的括號}
\Procedure{Print-Parens}{$s, i, j$}
    \If{$i = j$}
        \State 印出 $A_i$
    \Else
        \State 印出「(」
        \State \Call{Print-Parens}{$s, i, s[i, j]$}
        \State \Call{Print-Parens}{$s, s[i, j] + 1, j$}
        \State 印出「)」
    \EndIf
\EndProcedure
\end{pseudo}

這組數字：$s[1, 6] = 2$，拆成 $A_{1..2}$ 和 $A_{3..6}$；$s[3, 6] = 5$，拆成 $A_{3..5}$ 和
$A_6$；$s[3, 5] = 4$，拆成 $A_{3..4}$ 和 $A_5$。 印出來是
\[
((A_1 A_2)(((A_3 A_4) A_5) A_6)) .
\]

驗算一次：$A_1 A_2$ 花 $6 \cdot 10 \cdot 3 = 180$ 得到 $6 \times 3$；$A_3 A_4$ 花
$3 \cdot 12 \cdot 5 = 180$ 得到 $3 \times 5$；再乘 $A_5$ 花 $3 \cdot 5 \cdot 8 = 120$ 得到
$3 \times 8$；再乘 $A_6$ 花 $3 \cdot 8 \cdot 4 = 96$ 得到 $3 \times 4$；最後 $6 \times 3$ 乘
$3 \times 4$ 花 $72$。 總共 $180 + 180 + 120 + 96 + 72 = 648$，和 $m[1, 6]$ 一樣。

\subsection{C 程式}

\begin{cbookrecap}[第~9~章 9.7 節、第~19~章 19.3 節]
C 的二維陣列要在編譯時知道列寬。 列寬要等執行時才知道時，本書一律用 \texttt{calloc} 要一
塊 $(n+1)(n+1)$ 格的一維陣列，第 $i$ 列第 $j$ 格放在 \texttt{t[i * (n + 1) + j]}，這就是第~9~
章「多維陣列在記憶體裡其實攤平成一列」的寫法。 \texttt{calloc} 會把整塊設成 $0$，用完要
\texttt{free}。
\end{cbookrecap}

$m$、$s$ 兩張表都這樣攤平，用巨集 \texttt{AT(t, i, j)} 取第 $i$ 列第 $j$ 格；索引和虛擬碼一
樣從 $1$ 開始，第 $0$ 列、第 $0$ 欄空著不用，比較好對照。 $m$ 用 $-1$ 表示「還沒有候選」，
取代虛擬碼的 $\infty$。

\begin{example}[label=ex:dp-mc]{c}{(matrix\_chain.c) 矩陣鏈乘的填表與括號}
#include <stdio.h>
#include <stdlib.h>

#define MAXN 100
#define AT(t, i, j) ((t)[(i) * (n + 1) + (j)])   // (n+1) x (n+1) 的表攤平

// 依 s 表印出 A_i..A_j 的最佳括號
void print_parens(const int *s, int n, int i, int j)
{
    if (i == j) {
        printf("A%d", i);
        return;
    }
    printf("(");
    print_parens(s, n, i, AT(s, i, j));
    print_parens(s, n, AT(s, i, j) + 1, j);
    printf(")");
}

int main(void)
{
    int n;
    long long p[MAXN + 1];
    if (scanf("%d", &n) != 1 || n < 1 || n > MAXN) {
        return 1;
    }
    for (int i = 0; i <= n; i++) {
        if (scanf("%lld", &p[i]) != 1) {
            return 1;
        }
    }
    long long *m = calloc((size_t)(n + 1) * (n + 1), sizeof *m);
    int *s = calloc((size_t)(n + 1) * (n + 1), sizeof *s);
    if (m == NULL || s == NULL) {
        free(m);
        free(s);
        return 1;
    }
    for (int len = 2; len <= n; len++) {           // 鏈長由短到長
        for (int i = 1; i + len - 1 <= n; i++) {
            int j = i + len - 1;
            AT(m, i, j) = -1;                      // -1 代表還沒有候選
            for (int k = i; k < j; k++) {
                long long q = AT(m, i, k) + AT(m, k + 1, j)
                            + p[i - 1] * p[k] * p[j];
                if (AT(m, i, j) < 0 || q < AT(m, i, j)) {
                    AT(m, i, j) = q;               // 平手時保留較小的 k
                    AT(s, i, j) = k;
                }
            }
        }
    }
    printf("m[1][%d] = %lld\n", n, AT(m, 1, n));
    print_parens(s, n, 1, n);
    printf("\n");
    free(m);
    free(s);
    return 0;
}
\end{example}
\iosample{6 \\ 6 10 3 12 5 8 4}{m[1][6] = 648 \\ ((A1A2)(((A3A4)A5)A6))}

如果把外面兩層迴圈改成「$i$ 從 $1$ 到 $n$、
$j$ 從 $i + 1$ 到 $n$」一列一列由上往下填，算 $m[1, 3]$ 時 $m[2, 3]$ 還沒算，答案就錯了。
除了依鏈長，另一種正確的順序是 $i$ 從 $n$ \textbf{往下}到 $1$、$j$ 從 $i + 1$ 往上：這樣
算 $m[i, j]$ 時，第 $i + 1$ 列以下都已經填好，第 $i$ 列左邊的格子也填好了。

\section{最長共同子序列}
\label{sec:dp-lcs}

\subsection{定義}

從序列 $X = \langle x_1, x_2, \ldots, x_m \rangle$ 裡刪掉任意幾個元素（可以不刪），剩下的元素保
持原來的順序，得到的序列叫 $X$ 的\textbf{子序列}（subsequence）。 例如 \texttt{ALE} 是
\texttt{PALETTE} 的子序列（取第 $2, 3, 4$ 個），\texttt{LAP} 不是。 子序列不要求連續，這是
它和「子字串」不同的地方。 $Z$ 同時是 $X$ 和 $Y$ 的子序列時，叫\textbf{共同子序列}；最長
的共同子序列叫 \textbf{LCS}（longest common subsequence）。 LCS 不一定只有一個，但長度是
確定的。

章首的問題就是 LCS：把每一行程式當成一個元素，兩個版本的 LCS 就是最多能原封不動對上的
行。 Unix 的 \texttt{diff}、版本控制系統比較兩個檔案，背後都是這個問題（實際的工具用了更
快的變形）。 生物資訊裡比對兩段 DNA，也常用 LCS 和下一節的編輯距離。

暴力法是列出 $X$ 的 $2^m$ 個子序列，一個一個檢查是不是 $Y$ 的子序列（從左到右掃一次
$Y$，$O(n)$），時間 $\Theta(n 2^m)$。 下面用\textbf{前綴}當子問題：$X_i = \langle x_1,
\ldots, x_i \rangle$ 是 $X$ 的前 $i$ 個元素，$X_0$ 是空序列。

\subsection{最後一個字元的三種情形}

$Z$ 是 $X$ 和 $Y$ 的共同子序列，表示 $Z$ 的每個元素可以依序\textbf{配}到 $X$ 的一個位置和
$Y$ 的一個位置，配到的位置在 $X$、$Y$ 裡都是由左往右遞增的。 只有 $Z$ 的最後一個元素能配到
$x_m$，也只有它能配到 $y_n$。 所以對任何一種配法，下面三種情形一定有一種成立：

\begin{enumerate}
\item $x_m$ 和 $y_n$ 互相配對（這時一定 $x_m = y_n$）；
\item $x_m$ 沒有被配到；
\item $y_n$ 沒有被配到。
\end{enumerate}

令 $c[i, j]$ 為 $X_i$ 和 $Y_j$ 的 LCS 長度。 由這三種情形得到：

\begin{theorem}
\label{thm:dp-lcs}
$i, j \ge 1$ 時，
\[
c[i, j] = \begin{cases}
c[i - 1, j - 1] + 1, & x_i = y_j, \\
\max\bigl(c[i - 1, j],\ c[i, j - 1]\bigr), & x_i \ne y_j,
\end{cases}
\]
而 $c[i, 0] = c[0, j] = 0$。
\end{theorem}

\begin{proof}
先證 $x_i \ne y_j$ 的情形。 $X_{i-1}$ 和 $Y_j$ 的任何共同子序列也是 $X_i$ 和 $Y_j$ 的，所以
$c[i, j] \ge c[i - 1, j]$；同理 $c[i, j] \ge c[i, j - 1]$。 反過來，取 $X_i$、$Y_j$ 的一個
LCS 和它的一種配法。 情形 1 不可能（$x_i \ne y_j$），所以 $x_i$ 或 $y_j$ 沒有被配到；前者
表示它是 $X_{i-1}$ 和 $Y_j$ 的共同子序列，長度不超過 $c[i-1, j]$，後者長度不超過
$c[i, j-1]$。 兩個方向合起來就是等式。

再證 $x_i = y_j$。 在 $X_{i-1}$、$Y_{j-1}$ 的 LCS 後面接上 $x_i$，得到 $X_i$、$Y_j$ 的共同
子序列，所以 $c[i, j] \ge c[i-1, j-1] + 1$。 反過來，取一個 LCS 和它的配法：情形 1 時去掉
最後一個元素，剩下的是 $X_{i-1}$、$Y_{j-1}$ 的共同子序列，長度 $\le c[i-1, j-1]$；情形 2
時它是 $X_{i-1}$、$Y_j$ 的共同子序列，而 $X_{i-1}$ 和 $Y_j$ 的共同子序列去掉最後一個元素
（如果它配到 $y_j$）就是 $X_{i-1}$、$Y_{j-1}$ 的共同子序列，所以長度
$\le c[i-1, j-1] + 1$；情形 3 對稱。 三種情形都不超過 $c[i-1, j-1] + 1$。
\end{proof}

這個定理也說明了最佳子結構：LCS 的前面一段，一定是某一對前綴的 LCS。 子問題有
$(m+1)(n+1)$ 個，每個只看三個鄰居，所以填表是 $\Theta(mn)$。 $x_i = y_j$ 時\textbf{不用}
再和 $c[i-1, j]$、$c[i, j-1]$ 比，配對一定不會比較差，這是證明裡的第二段。

\subsection{填表與箭頭}

除了數字，每一格再記一個箭頭 $b[i, j]$，說明這一格的值是從哪裡來的：$\nwarrow$ 表示
$x_i = y_j$、從左上角加一；$\uparrow$ 表示取上面的 $c[i-1, j]$；$\leftarrow$ 表示取左邊的
$c[i, j-1]$。

\begin{pseudo}[label=alg:dp-lcs-length]{LCS 的填表}
\Function{LCS-Length}{$X, Y, m, n$}
    \For{$i \gets 0$ \To $m$}
        \State $c[i, 0] \gets 0$
    \EndFor
    \For{$j \gets 0$ \To $n$}
        \State $c[0, j] \gets 0$
    \EndFor
    \For{$i \gets 1$ \To $m$}
        \For{$j \gets 1$ \To $n$}
            \If{$x_i = y_j$}
                \State $c[i, j] \gets c[i - 1, j - 1] + 1$
                \State $b[i, j] \gets \text{「}\nwarrow\text{」}$
            \ElsIf{$c[i - 1, j] \ge c[i, j - 1]$}
                \State $c[i, j] \gets c[i - 1, j]$
                \State $b[i, j] \gets \text{「}\uparrow\text{」}$
            \Else
                \State $c[i, j] \gets c[i, j - 1]$
                \State $b[i, j] \gets \text{「}\leftarrow\text{」}$
            \EndIf
        \EndFor
    \EndFor
    \State \Return $c$ 和 $b$
\EndFunction
\end{pseudo}

第 11 行用 $\ge$：上面和左邊一樣大時，箭頭指向上。 這是本書（也是 CLRS）的平手規則。

\begin{pseudo}[label=alg:dp-print-lcs]{沿箭頭印出一個 LCS}
\Procedure{Print-LCS}{$b, X, i, j$}
    \If{$i = 0$ \OR $j = 0$}
        \State \Return
    \EndIf
    \If{$b[i, j] = \text{「}\nwarrow\text{」}$}
        \State \Call{Print-LCS}{$b, X, i - 1, j - 1$}
        \State 印出 $x_i$
    \ElsIf{$b[i, j] = \text{「}\uparrow\text{」}$}
        \State \Call{Print-LCS}{$b, X, i - 1, j$}
    \Else
        \State \Call{Print-LCS}{$b, X, i, j - 1$}
    \EndIf
\EndProcedure
\end{pseudo}

呼叫 $\proc{Print-LCS}(b, X, m, n)$。 每次呼叫 $i$ 或 $j$ 至少少一，所以是 $O(m + n)$。 先遞
迴再印，印出來的順序才是由左到右。

\subsection{手算追蹤}
\label{sec:dp-lcs-trace}

$X = \texttt{PALETTE}$（$m = 7$），$Y = \texttt{PLANET}$（$n = 6$）。 列是 $X$、欄是 $Y$，
一列一列由左往右填。

\begin{center}
\begin{minipage}{\linewidth}
\centering
\begin{tikzpicture}[x=0.98cm, y=0.66cm]
\node[dp-hd] at (0,1) {$j$};
\node[dp-hd] at (-1,0) {$i$};
\node[dp-hd] at (1,1) {0};
\node[dp-hd] at (2,1) {1};
\node[dp-hd] at (2,0.4) {\texttt{P}};
\node[dp-hd] at (3,1) {2};
\node[dp-hd] at (3,0.4) {\texttt{L}};
\node[dp-hd] at (4,1) {3};
\node[dp-hd] at (4,0.4) {\texttt{A}};
\node[dp-hd] at (5,1) {4};
\node[dp-hd] at (5,0.4) {\texttt{N}};
\node[dp-hd] at (6,1) {5};
\node[dp-hd] at (6,0.4) {\texttt{E}};
\node[dp-hd] at (7,1) {6};
\node[dp-hd] at (7,0.4) {\texttt{T}};
\node[dp-hd] at (-1,-0.6) {0};
\node[dp-c0, minimum width=0.98cm] at (1,-0.6) {0};
\node[dp-c0, minimum width=0.98cm] at (2,-0.6) {0};
\node[dp-c0, minimum width=0.98cm] at (3,-0.6) {0};
\node[dp-c0, minimum width=0.98cm] at (4,-0.6) {0};
\node[dp-c0, minimum width=0.98cm] at (5,-0.6) {0};
\node[dp-c0, minimum width=0.98cm] at (6,-0.6) {0};
\node[dp-c0, minimum width=0.98cm] at (7,-0.6) {0};
\node[dp-hd] at (-1,-1.6) {1};
\node[dp-hd] at (-0.3,-1.6) {\texttt{P}};
\node[dp-c0, minimum width=0.98cm] at (1,-1.6) {0};
\node[dp-chit, minimum width=0.98cm] at (2,-1.6) {$\nwarrow$\,1};
\node[dp-cpath, minimum width=0.98cm] at (3,-1.6) {$\leftarrow$\,1};
\node[dp-c, minimum width=0.98cm] at (4,-1.6) {$\leftarrow$\,1};
\node[dp-c, minimum width=0.98cm] at (5,-1.6) {$\leftarrow$\,1};
\node[dp-c, minimum width=0.98cm] at (6,-1.6) {$\leftarrow$\,1};
\node[dp-c, minimum width=0.98cm] at (7,-1.6) {$\leftarrow$\,1};
\node[dp-hd] at (-1,-2.6) {2};
\node[dp-hd] at (-0.3,-2.6) {\texttt{A}};
\node[dp-c0, minimum width=0.98cm] at (1,-2.6) {0};
\node[dp-c, minimum width=0.98cm] at (2,-2.6) {$\uparrow$\,1};
\node[dp-c, minimum width=0.98cm] at (3,-2.6) {$\uparrow$\,1};
\node[dp-chit, minimum width=0.98cm] at (4,-2.6) {$\nwarrow$\,2};
\node[dp-cpath, minimum width=0.98cm] at (5,-2.6) {$\leftarrow$\,2};
\node[dp-c, minimum width=0.98cm] at (6,-2.6) {$\leftarrow$\,2};
\node[dp-c, minimum width=0.98cm] at (7,-2.6) {$\leftarrow$\,2};
\node[dp-hd] at (-1,-3.6) {3};
\node[dp-hd] at (-0.3,-3.6) {\texttt{L}};
\node[dp-c0, minimum width=0.98cm] at (1,-3.6) {0};
\node[dp-c, minimum width=0.98cm] at (2,-3.6) {$\uparrow$\,1};
\node[dp-c, minimum width=0.98cm] at (3,-3.6) {$\nwarrow$\,2};
\node[dp-c, minimum width=0.98cm] at (4,-3.6) {$\uparrow$\,2};
\node[dp-cpath, minimum width=0.98cm] at (5,-3.6) {$\uparrow$\,2};
\node[dp-c, minimum width=0.98cm] at (6,-3.6) {$\uparrow$\,2};
\node[dp-c, minimum width=0.98cm] at (7,-3.6) {$\uparrow$\,2};
\node[dp-hd] at (-1,-4.6) {4};
\node[dp-hd] at (-0.3,-4.6) {\texttt{E}};
\node[dp-c0, minimum width=0.98cm] at (1,-4.6) {0};
\node[dp-c, minimum width=0.98cm] at (2,-4.6) {$\uparrow$\,1};
\node[dp-c, minimum width=0.98cm] at (3,-4.6) {$\uparrow$\,2};
\node[dp-c, minimum width=0.98cm] at (4,-4.6) {$\uparrow$\,2};
\node[dp-c, minimum width=0.98cm] at (5,-4.6) {$\uparrow$\,2};
\node[dp-chit, minimum width=0.98cm] at (6,-4.6) {$\nwarrow$\,3};
\node[dp-c, minimum width=0.98cm] at (7,-4.6) {$\leftarrow$\,3};
\node[dp-hd] at (-1,-5.6) {5};
\node[dp-hd] at (-0.3,-5.6) {\texttt{T}};
\node[dp-c0, minimum width=0.98cm] at (1,-5.6) {0};
\node[dp-c, minimum width=0.98cm] at (2,-5.6) {$\uparrow$\,1};
\node[dp-c, minimum width=0.98cm] at (3,-5.6) {$\uparrow$\,2};
\node[dp-c, minimum width=0.98cm] at (4,-5.6) {$\uparrow$\,2};
\node[dp-c, minimum width=0.98cm] at (5,-5.6) {$\uparrow$\,2};
\node[dp-cpath, minimum width=0.98cm] at (6,-5.6) {$\uparrow$\,3};
\node[dp-c, minimum width=0.98cm] at (7,-5.6) {$\nwarrow$\,4};
\node[dp-hd] at (-1,-6.6) {6};
\node[dp-hd] at (-0.3,-6.6) {\texttt{T}};
\node[dp-c0, minimum width=0.98cm] at (1,-6.6) {0};
\node[dp-c, minimum width=0.98cm] at (2,-6.6) {$\uparrow$\,1};
\node[dp-c, minimum width=0.98cm] at (3,-6.6) {$\uparrow$\,2};
\node[dp-c, minimum width=0.98cm] at (4,-6.6) {$\uparrow$\,2};
\node[dp-c, minimum width=0.98cm] at (5,-6.6) {$\uparrow$\,2};
\node[dp-c, minimum width=0.98cm] at (6,-6.6) {$\uparrow$\,3};
\node[dp-chit, minimum width=0.98cm] at (7,-6.6) {$\nwarrow$\,4};
\node[dp-hd] at (-1,-7.6) {7};
\node[dp-hd] at (-0.3,-7.6) {\texttt{E}};
\node[dp-c0, minimum width=0.98cm] at (1,-7.6) {0};
\node[dp-c, minimum width=0.98cm] at (2,-7.6) {$\uparrow$\,1};
\node[dp-c, minimum width=0.98cm] at (3,-7.6) {$\uparrow$\,2};
\node[dp-c, minimum width=0.98cm] at (4,-7.6) {$\uparrow$\,2};
\node[dp-c, minimum width=0.98cm] at (5,-7.6) {$\uparrow$\,2};
\node[dp-c, minimum width=0.98cm] at (6,-7.6) {$\nwarrow$\,3};
\node[dp-cpath, minimum width=0.98cm] at (7,-7.6) {$\uparrow$\,4};
\end{tikzpicture}
\diagcaption{\texttt{PALETTE} 和 \texttt{PLANET} 的 $c$ 表和箭頭。 第 $0$ 列、第 $0$ 欄是
邊界。 從右下角 $c[7, 6] = 4$ 沿箭頭走回左上，經過的格子加框；粉紅色的是 $\nwarrow$，也就
是被選進 LCS 的字元：\texttt{P}、\texttt{A}、\texttt{E}、\texttt{T}。}
\label{fig:dp-lcs-table}
\end{minipage}
\end{center}

看幾格是怎麼來的：
\begin{itemize}
\item $c[2, 3]$：$x_2 = \texttt{A} = y_3$，從左上 $c[1, 2] = 1$ 加一，得 $2$，箭頭 $\nwarrow$。
\item $c[3, 3]$：$x_3 = \texttt{L} \ne \texttt{A}$，上面 $c[2, 3] = 2$、左邊 $c[3, 2] = 2$，平
手，照規則取上面，箭頭 $\uparrow$。
\item $c[4, 5]$：$x_4 = \texttt{E} = y_5$，從 $c[3, 4] = 2$ 加一，得 $3$。
\item $c[5, 6]$：$x_5 = \texttt{T} = y_6$，從 $c[4, 5] = 3$ 加一，得 $4$。
\item $c[6, 6]$：$x_6$ 也是 \texttt{T}，從 $c[5, 5] = 3$ 加一，也是 $4$，箭頭 $\nwarrow$。
\item $c[7, 6]$：$x_7 = \texttt{E} \ne \texttt{T}$，上面是 $4$、左邊是 $3$，取上面。
\end{itemize}

從 $c[7, 6]$ 往回走：$\uparrow$ 到 $[6, 6]$，$\nwarrow$（選 $x_6 = \texttt{T}$）到 $[5, 5]$，
$\uparrow$ 到 $[4, 5]$，$\nwarrow$（選 \texttt{E}）到 $[3, 4]$，$\uparrow$ 到 $[2, 4]$，
$\leftarrow$ 到 $[2, 3]$，$\nwarrow$（選 \texttt{A}）到 $[1, 2]$，$\leftarrow$ 到
$[1, 1]$，$\nwarrow$（選 \texttt{P}）到 $[0, 0]$。 倒過來讀，LCS 是 \texttt{PAET}。

這兩個字串還有另一個長度 $4$ 的 LCS：\texttt{PLET}（\texttt{PALETTE} 的第 $1, 3, 4, 5$ 個字
元）。 沿著箭頭只能走出一條路，所以 \proc{Print-LCS} 只印出其中一個；印出哪一個由平手規
則決定。

\begin{examnote}
平手時選哪一邊，不同的課本規定不同。 LCS 平手時本書指向上（第 11 行的 $\ge$），有些書指
向左；同一題 \texttt{PALETTE}、\texttt{PLANET}，指向左會還原出 \texttt{PLET}。 矩陣鏈乘
和最佳二元搜尋樹，本書在平手時取最小的 $k$ 或 $r$。 考題有規定時照題目；沒有規定時，在答
案裡寫一句你用的規則（例如「平手時取 $\uparrow$」），長度和最佳值不受影響，只有還原出來
的解可能不同。 題目問「所有的 LCS」時，平手的格子兩個方向都要走。
\end{examnote}

\subsection{省掉箭頭表}

$b$ 表其實可以不存。 還原時站在 $[i, j]$，只要重新判斷一次：$x_i = y_j$ 就往左上；否則比較
$c[i-1, j]$ 和 $c[i, j-1]$，用和填表時一樣的平手規則決定往上還是往左。 每一步 $O(1)$，結
果和用箭頭完全相同，省下一張 $(m+1)(n+1)$ 的表。

\subsection{只留兩列}
\label{sec:dp-lcs-space}

如果只要 LCS 的長度，連 $c$ 表都不用整張留著。 第 $i$ 列只用到第 $i - 1$ 列和自己左邊的格
子，所以只要兩列：$\mathit{prev}$ 是上一列，$\mathit{cur}$ 是正在填的這一列，一列填完後交換
兩者的角色。

\begin{pseudo}[label=alg:dp-lcs-two-rows]{只用兩列算 LCS 長度}
\Function{LCS-Length-Two-Rows}{$X, Y, m, n$}
    \State 建立 $\mathit{prev}[0..n]$ 和 $\mathit{cur}[0..n]$，每一格都設成 $0$
    \For{$i \gets 1$ \To $m$}
        \For{$j \gets 1$ \To $n$}
            \If{$x_i = y_j$}
                \State $\mathit{cur}[j] \gets \mathit{prev}[j - 1] + 1$
            \Else
                \State $\mathit{cur}[j] \gets \max(\mathit{prev}[j], \mathit{cur}[j - 1])$
            \EndIf
        \EndFor
        \State 交換 $\mathit{prev}$ 和 $\mathit{cur}$ \Comment{只交換指標，不複製}
    \EndFor
    \State \Return $\mathit{prev}[n]$
\EndFunction
\end{pseudo}

$\mathit{cur}[0]$ 從來沒被寫過，一直是 $0$，正好是邊界 $c[i, 0]$。 第 9 行換完之後，
$\mathit{prev}$ 存的是剛填好的第 $i$ 列，所以迴圈結束時答案在 $\mathit{prev}[n]$。 空間從
$\Theta(mn)$ 降到 $\Theta(n)$；讓較短的字串當 $Y$，就是 $\Theta(\min(m, n))$。 再省一點，只用一
列也可以：由左往右更新 $\mathit{row}[j]$ 之前，先把舊的 $\mathit{row}[j]$ 存進一個變數，填下一
格時它就是需要的左上角 $c[i-1, j]$。

代價是\textbf{還原不出 LCS 本身}：走回去時需要的那些格子已經被蓋掉了。 要在線性空間裡同時
得到 LCS，可以用 Hirschberg 的方法：用兩列的版本從前面、從後面各算一次，找出 LCS 通過第
$\lfloor m/2 \rfloor$ 列的哪一格，再對左上、右下兩塊分治。 時間仍是 $\Theta(mn)$，這裡不
細講。

\subsection{C 程式}

\texttt{lcs} 填整張表，照前面「省掉箭頭表」的方法直接用 $c$ 表還原。 只要長度的話，把表
換成虛擬碼~\ref{alg:dp-lcs-two-rows} 的兩列即可。 字串的第 $i$ 個字元在 C 裡是 \texttt{x[i - 1]}，表的索引仍然從 $0$ 開始，第 $0$
列、第 $0$ 欄是邊界。

\begin{example}[label=ex:dp-lcs]{c}{(lcs.c) LCS 的填表與還原}
#include <stdio.h>
#include <stdlib.h>
#include <string.h>

#define MAXLEN 1000

// 填完整的 (m+1) x (n+1) 表，再從右下角走回來還原一個 LCS 放進 out。
int lcs(const char *x, const char *y, char *out)
{
    int m = strlen(x), n = strlen(y);
    int *c = calloc((size_t)(m + 1) * (n + 1), sizeof *c);
    if (c == NULL) {
        return -1;
    }
#define C(i, j) c[(i) * (n + 1) + (j)]
    for (int i = 1; i <= m; i++) {
        for (int j = 1; j <= n; j++) {
            if (x[i - 1] == y[j - 1]) {
                C(i, j) = C(i - 1, j - 1) + 1;
            } else if (C(i - 1, j) >= C(i, j - 1)) {
                C(i, j) = C(i - 1, j);          // 平手時往上
            } else {
                C(i, j) = C(i, j - 1);
            }
        }
    }
    int len = C(m, n);
    int i = m, j = n, k = len;
    out[k] = '\0';
    while (i > 0 && j > 0) {                    // 照填表時的規則走回去
        if (x[i - 1] == y[j - 1]) {
            out[--k] = x[i - 1];
            i--;
            j--;
        } else if (C(i - 1, j) >= C(i, j - 1)) {
            i--;
        } else {
            j--;
        }
    }
#undef C
    free(c);
    return len;
}

int main(void)
{
    static char x[MAXLEN + 1], y[MAXLEN + 1], z[MAXLEN + 1];
    while (scanf("%1000s %1000s", x, y) == 2) {
        int len = lcs(x, y, z);
        printf("length %d: %s\n", len, z);
    }
    return 0;
}
\end{example}
\iosample{PALETTE PLANET \\ PLANET PALETTE \\ SUBSEQUENCE CONSEQUENCES}{length 4: PAET \\ length 4: PLET \\ length 8: SEQUENCE}

第二組把兩個字串對調，長度一樣是 $4$，還原出來的卻是 \texttt{PLET}：對調之後列和欄互換，
「平手取上」變成了原本的「平手取左」。 第三組說明 LCS 不必是連續的一段：
\texttt{SUBSEQUENCE} 和 \texttt{CONSEQUENCES} 的 LCS 是 \texttt{SEQUENCE}，長度 $8$。

\section{最佳二元搜尋樹}
\label{sec:dp-obst}

\subsection{問題與成本模型}

編譯器讀程式時，每讀到一個詞，就要判斷它是不是關鍵字。 假設只看五個關鍵字，依字典順序是
\[
k_1 = \texttt{char} < k_2 = \texttt{for} < k_3 = \texttt{goto} < k_4 = \texttt{int}
< k_5 = \texttt{while},
\]
存在一棵二元搜尋樹裡（第~9~章）。 從大量程式統計出來，每 $50$ 次查詢裡，查到各個關鍵字的
次數 $p_i$，以及查不到、落在兩個關鍵字之間的次數 $q_i$ 如下：

\begin{center}
\begin{tabular}{c*{6}{c}}
\toprule
$i$ & 0 & 1 & 2 & 3 & 4 & 5 \\
\midrule
$k_i$ & & \texttt{char} & \texttt{for} & \texttt{goto} & \texttt{int} & \texttt{while} \\
$p_i$ & & 4 & 10 & 3 & 8 & 6 \\
$q_i$ & 3 & 2 & 4 & 1 & 3 & 6 \\
\bottomrule
\end{tabular}
\end{center}

$q_0$ 是比 $k_1$ 小的詞（例如 \texttt{buf}），$q_i$（$1 \le i < n$）是介於 $k_i$ 和
$k_{i+1}$ 之間的詞（例如 $q_3$ 數的是 \texttt{goto} 和 \texttt{int} 之間的 \texttt{i}、
\texttt{idx}），$q_5$ 是比 \texttt{while} 大的詞（例如 \texttt{x}）。 $\sum p_i + \sum q_i
= 31 + 19 = 50$。 有些書用機率，那只是把每個數除以總數，下面的計算都一樣。

把每個查不到的範圍畫成一個\textbf{外部節點} $d_0, d_1, \ldots, d_n$，接在原本是 \NIL 的位置
上，得到第~7~章的延伸二元樹：$n$ 個鍵是內部節點，$n + 1$ 個外部節點是葉節點，中序走訪的
順序是 $d_0, k_1, d_1, k_2, \ldots, k_n, d_n$。 一次查不到的搜尋，一定停在某個 $d_i$。

\textbf{成本模型。} 本書數的是\textbf{和鍵比較的次數}，深度和高度沿用第~7~章的定義（深度數
邊、根的深度是 $0$；高度數節點）：
\begin{itemize}
\item 搜尋 $k_i$，從根往下走到 $k_i$，路上每個鍵都比一次，所以是 $\operatorname{depth}(k_i)
+ 1$ 次比較，也就是 $k_i$ 所在的\textbf{層數}。
\item 查不到、停在 $d_i$，路上比較的是 $d_i$ 上面的每一個鍵，共 $\operatorname{depth}(d_i)$
次；$d_i$ 本身不是鍵，走到它就知道查不到，不再比較。
\end{itemize}
所以一棵樹 $T$ 的總成本是
\[
\operatorname{cost}(T) = \sum_{i=1}^{n} p_i \bigl(\operatorname{depth}_T(k_i) + 1\bigr)
+ \sum_{i=0}^{n} q_i \operatorname{depth}_T(d_i) ,
\]
除以總次數就是平均每次查詢比較幾次。 查到時最多比較的次數等於內部節點組成的那棵樹的
高度，這正是第~7~章「高度數節點」的好處。

\textbf{最佳二元搜尋樹}（optimal binary search tree）是在所有以 $k_1, \ldots, k_n$ 為鍵的
二元搜尋樹中，總成本最小的一棵。 $n$ 個鍵的二元搜尋樹形狀有卡特蘭數 $C_n$ 種（第~7~章），
全部試一遍是指數時間。

\begin{center}
\begin{minipage}{\linewidth}
\centering
\begin{tikzpicture}[x=0.62cm, y=0.95cm, baseline=(current bounding box.north)]
\node[dp-ext] (gd0) at (0,-2) {$d_0$};
\node[dp-key] (gk1) at (1,-1) {\texttt{char}};
\node[dp-ext] (gd1) at (2,-2) {$d_1$};
\node[dp-key] (gk2) at (3,0) {\texttt{for}};
\node[dp-ext] (gd2) at (4,-3) {$d_2$};
\node[dp-key] (gk3) at (5,-2) {\texttt{goto}};
\node[dp-ext] (gd3) at (6,-3) {$d_3$};
\node[dp-key] (gk4) at (7,-1) {\texttt{int}};
\node[dp-ext] (gd4) at (8,-3) {$d_4$};
\node[dp-key] (gk5) at (9,-2) {\texttt{while}};
\node[dp-ext] (gd5) at (10,-3) {$d_5$};
\draw[dp-e] (gk1) -- (gd0);
\draw[dp-e] (gk2) -- (gk1);
\draw[dp-e] (gk1) -- (gd1);
\draw[dp-e] (gk3) -- (gd2);
\draw[dp-e] (gk4) -- (gk3);
\draw[dp-e] (gk3) -- (gd3);
\draw[dp-e] (gk2) -- (gk4);
\draw[dp-e] (gk5) -- (gd4);
\draw[dp-e] (gk4) -- (gk5);
\draw[dp-e] (gk5) -- (gd5);
\node[font=\small] at (5,-3.8) {(a) 每次挑次數最多的當根，成本 $113$};
\end{tikzpicture}\hspace{0.5cm}\begin{tikzpicture}[x=0.62cm, y=0.95cm, baseline=(current bounding box.north)]
\node[dp-ext] (od0) at (0,-3) {$d_0$};
\node[dp-key] (ok1) at (1,-2) {\texttt{char}};
\node[dp-ext] (od1) at (2,-3) {$d_1$};
\node[dp-key] (ok2) at (3,-1) {\texttt{for}};
\node[dp-ext] (od2) at (4,-3) {$d_2$};
\node[dp-key] (ok3) at (5,-2) {\texttt{goto}};
\node[dp-ext] (od3) at (6,-3) {$d_3$};
\node[dp-key] (ok4) at (7,0) {\texttt{int}};
\node[dp-ext] (od4) at (8,-2) {$d_4$};
\node[dp-key] (ok5) at (9,-1) {\texttt{while}};
\node[dp-ext] (od5) at (10,-2) {$d_5$};
\draw[dp-e] (ok1) -- (od0);
\draw[dp-e] (ok2) -- (ok1);
\draw[dp-e] (ok1) -- (od1);
\draw[dp-e] (ok4) -- (ok2);
\draw[dp-e] (ok3) -- (od2);
\draw[dp-e] (ok2) -- (ok3);
\draw[dp-e] (ok3) -- (od3);
\draw[dp-e] (ok5) -- (od4);
\draw[dp-e] (ok4) -- (ok5);
\draw[dp-e] (ok5) -- (od5);
\node[font=\small] at (5,-3.8) {(b) 最佳二元搜尋樹，成本 $109$};
\end{tikzpicture}
\diagcaption{同一組次數的兩棵樹，圓角框是鍵，灰色方塊是外部節點。 (a) 每次挑次數最多的鍵當
根：先挑 \texttt{for}（$10$ 次），右邊再挑 \texttt{int}。 (b) 最佳的樹，根是
\texttt{int}，只出現 $8$ 次。}
\label{fig:dp-obst-trees}
\end{minipage}
\end{center}

直覺的做法是讓最常查的鍵當根。 圖~\ref{fig:dp-obst-trees}(a) 這樣做，成本 $113$；(b) 的
根是次數第二多的 \texttt{int}，成本卻只有 $109$。 和 (b) 相比，(a) 讓 \texttt{for}、\texttt{char}、$d_0$、$d_1$ 各少比一次，省下
$10 + 4 + 3 + 2 = 19$；卻讓 \texttt{int}、\texttt{while}、$d_4$、$d_5$ 各多比一次，多了
$8 + 6 + 3 + 6 = 23$。 下表逐項算 (b) 的成本，
比較次數一欄對鍵是深度加一、對外部節點是深度：

\begin{center}
\begin{tabular}{l*{11}{c}c}
\toprule
節點 & $d_0$ & $k_1$ & $d_1$ & $k_2$ & $d_2$ & $k_3$ & $d_3$ & $k_4$ & $d_4$ & $k_5$ & $d_5$ & 合計 \\
\midrule
深度 & 3 & 2 & 3 & 1 & 3 & 2 & 3 & 0 & 2 & 1 & 2 & \\
比較次數 & 3 & 3 & 3 & 2 & 3 & 3 & 3 & 1 & 2 & 2 & 2 & \\
次數 & 3 & 4 & 2 & 10 & 4 & 3 & 1 & 8 & 3 & 6 & 6 & 50 \\
乘積 & 9 & 12 & 6 & 20 & 12 & 9 & 3 & 8 & 6 & 12 & 12 & 109 \\
\bottomrule
\end{tabular}
\end{center}

平均每次查詢比 $109 / 50 = 2.18$ 次。 這組次數下，$42$ 種形狀裡沒有比 $109$ 更小的，下面
用動態規劃算出這件事。

\subsection{最佳子結構與遞迴式}

二元搜尋樹的任何一棵子樹，裡面的鍵一定是連續的一段 $k_i, \ldots, k_j$，外部節點也一定是對
應的 $d_{i-1}, \ldots, d_j$（中序順序決定的）。 如果最佳的樹 $T$ 有一棵子樹 $T'$ 裝的是
$k_i..k_j$，那麼 $T'$ 一定是這一段的最佳樹：把 $T'$ 換成這一段更好的樹，$T$ 的其他部分不
變，總成本會下降，矛盾。 這又是剪下貼上。

接著看把一棵子樹掛到一個新的根底下會怎樣：子樹裡每個節點的深度都加一，總成本增加的量正好
是子樹裡所有次數的和。 定義
\[
w(i, j) = \sum_{l=i}^{j} p_l + \sum_{l=i-1}^{j} q_l ,
\]
是裝著 $k_i..k_j$ 的子樹的\textbf{總重量}（注意外部節點從 $d_{i-1}$ 數起）。 令 $e[i, j]$
為 $k_i..k_j$ 的最佳樹的成本。 選 $k_r$（$i \le r \le j$）當根時，左子樹是 $k_i..k_{r-1}$，
右子樹是 $k_{r+1}..k_j$，掛上去之後兩邊各多 $w(i, r-1)$ 和 $w(r+1, j)$，根本身比一次，貢獻
$p_r$。 三者加起來剛好是 $w(i, j)$，所以
\[
e[i, j] = \begin{cases}
0, & j = i - 1, \\[2pt]
\displaystyle\min_{i \le r \le j} \bigl\{\, e[i, r-1] + e[r+1, j] \,\bigr\} + w(i, j), & i \le j .
\end{cases}
\]
$j = i - 1$ 是沒有鍵的「空樹」，只有一個外部節點 $d_{i-1}$，走到它不用比較，成本是 $0$。 另
外用 $\mathit{root}[i, j]$ 記下取到最小值的 $r$。 $w$ 也用表來存，避免每次重加：
\[
w[i, i - 1] = q_{i-1}, \qquad w[i, j] = w[i, j - 1] + p_j + q_j .
\]

\begin{pseudo}[label=alg:dp-obst]{最佳二元搜尋樹}
\Function{Optimal-BST}{$p, q, n$}
    \For{$i \gets 1$ \To $n + 1$}
        \State $e[i, i - 1] \gets 0$
        \State $w[i, i - 1] \gets q_{i-1}$
    \EndFor
    \For{$\ell \gets 1$ \To $n$} \Comment{這一段有幾個鍵}
        \For{$i \gets 1$ \To $n - \ell + 1$}
            \State $j \gets i + \ell - 1$
            \State $e[i, j] \gets \infty$
            \State $w[i, j] \gets w[i, j - 1] + p_j + q_j$
            \For{$r \gets i$ \To $j$}
                \State $t \gets e[i, r - 1] + e[r + 1, j] + w[i, j]$
                \If{$t < e[i, j]$}
                    \State $e[i, j] \gets t$
                    \State $\mathit{root}[i, j] \gets r$
                \EndIf
            \EndFor
        \EndFor
    \EndFor
    \State \Return $e$ 和 $\mathit{root}$
\EndFunction
\end{pseudo}

表的大小：$e$ 和 $w$ 的列 $i$ 從 $1$ 到 $n + 1$、欄 $j$ 從 $0$ 到 $n$（要放得下 $e[n+1, n]$
這種空樹），$\mathit{root}$ 只用 $1 \le i \le j \le n$。 結構和矩陣鏈乘幾乎一樣：依區間長度
填，第 10 行的迴圈每格跑 $\ell$ 次，總時間 $\Theta(n^3)$，空間 $\Theta(n^2)$。 第 12 行用
$<$，平手時保留最小的 $r$。

\subsection{手算追蹤}
\label{sec:dp-obst-trace}

$\ell = 1$ 時只有一個鍵，左右都是空樹，$e[i, i] = w[i, i] = q_{i-1} + p_i + q_i$：
$e[1,1] = 3 + 4 + 2 = 9$、$e[2,2] = 16$、$e[3,3] = 8$、$e[4,4] = 12$、$e[5,5] = 15$，根都是
自己。 $\ell \ge 2$ 的每一格寫出每個 $r$ 的 $e[i, r-1] + e[r+1, j] + w[i, j]$：

{\small
\begin{center}
\begin{tabular}{c c p{0.62\linewidth} c c}
\toprule
$(i, j)$ & $w$ & 每個 $r$ 的成本 & $e$ & $\mathit{root}$ \\
\midrule
$(1,2)$ & 23 & $r=1$：$0+16+23=39$；$r=2$：$9+0+23=\mathbf{32}$ & 32 & 2 \\
$(2,3)$ & 20 & $r=2$：$0+8+20=\mathbf{28}$；$r=3$：$16+0+20=36$ & 28 & 2 \\
$(3,4)$ & 19 & $r=3$：$0+12+19=31$；$r=4$：$8+0+19=\mathbf{27}$ & 27 & 4 \\
$(4,5)$ & 24 & $r=4$：$0+15+24=39$；$r=5$：$12+0+24=\mathbf{36}$ & 36 & 5 \\
$(1,3)$ & 27 & $r=1$：$0+28+27=55$；$r=2$：$9+8+27=\mathbf{44}$；$r=3$：$32+0+27=59$ & 44 & 2 \\
$(2,4)$ & 31 & $r=2$：$0+27+31=\mathbf{58}$；$r=3$：$16+12+31=59$；$r=4$：$28+0+31=59$ & 58 & 2 \\
$(3,5)$ & 31 & $r=3$：$0+36+31=67$；$r=4$：$8+15+31=\mathbf{54}$；$r=5$：$27+0+31=58$ & 54 & 4 \\
$(1,4)$ & 38 & $r=1$：$0+58+38=96$；$r=2$：$9+27+38=\mathbf{74}$；$r=3$：$32+12+38=82$；$r=4$：$44+0+38=82$ & 74 & 2 \\
$(2,5)$ & 43 & $r=2$：$0+54+43=97$；$r=3$：$16+36+43=95$；$r=4$：$28+15+43=\mathbf{86}$；$r=5$：$58+0+43=101$ & 86 & 4 \\
$(1,5)$ & 50 & $r=1$：$0+86+50=136$；$r=2$：$9+54+50=113$；$r=3$：$32+36+50=118$；$r=4$：$44+15+50=\mathbf{109}$；$r=5$：$74+0+50=124$ & 109 & 4 \\
\bottomrule
\end{tabular}
\end{center}
}

\begin{center}
\begin{minipage}{\linewidth}
\centering
\begin{tikzpicture}[x=0.85cm, y=0.62cm, baseline=(current bounding box.north)]
\node[font=\small] at (3.50,1.15) {$e[i,j]$};
\node[dp-hd] at (1,0.35) {$j=0$};
\node[dp-hd] at (2,0.35) {$1$};
\node[dp-hd] at (3,0.35) {$2$};
\node[dp-hd] at (4,0.35) {$3$};
\node[dp-hd] at (5,0.35) {$4$};
\node[dp-hd] at (6,0.35) {$5$};
\node[dp-hd, anchor=east] at (0.45,-0.5) {$i=1$};
\node[dp-c0, minimum width=0.85cm] at (1,-0.5) {0};
\node[dp-cd0, minimum width=0.85cm] at (2,-0.5) {9};
\node[dp-cd1, minimum width=0.85cm] at (3,-0.5) {32};
\node[dp-cd0, minimum width=0.85cm] at (4,-0.5) {44};
\node[dp-cd1, minimum width=0.85cm] at (5,-0.5) {74};
\node[dp-cd0, minimum width=0.85cm] at (6,-0.5) {109};
\node[dp-hd, anchor=east] at (0.45,-1.5) {$2$};
\node[dp-c0, minimum width=0.85cm] at (2,-1.5) {0};
\node[dp-cd0, minimum width=0.85cm] at (3,-1.5) {16};
\node[dp-cd1, minimum width=0.85cm] at (4,-1.5) {28};
\node[dp-cd0, minimum width=0.85cm] at (5,-1.5) {58};
\node[dp-cd1, minimum width=0.85cm] at (6,-1.5) {86};
\node[dp-hd, anchor=east] at (0.45,-2.5) {$3$};
\node[dp-c0, minimum width=0.85cm] at (3,-2.5) {0};
\node[dp-cd0, minimum width=0.85cm] at (4,-2.5) {8};
\node[dp-cd1, minimum width=0.85cm] at (5,-2.5) {27};
\node[dp-cd0, minimum width=0.85cm] at (6,-2.5) {54};
\node[dp-hd, anchor=east] at (0.45,-3.5) {$4$};
\node[dp-c0, minimum width=0.85cm] at (4,-3.5) {0};
\node[dp-cd0, minimum width=0.85cm] at (5,-3.5) {12};
\node[dp-cd1, minimum width=0.85cm] at (6,-3.5) {36};
\node[dp-hd, anchor=east] at (0.45,-4.5) {$5$};
\node[dp-c0, minimum width=0.85cm] at (5,-4.5) {0};
\node[dp-cd0, minimum width=0.85cm] at (6,-4.5) {15};
\node[dp-hd, anchor=east] at (0.45,-5.5) {$6$};
\node[dp-c0, minimum width=0.85cm] at (6,-5.5) {0};
\end{tikzpicture}\hspace{0.8cm}\begin{tikzpicture}[x=0.8cm, y=0.62cm, baseline=(current bounding box.north)]
\node[font=\small] at (3.00,1.15) {$\mathit{root}[i,j]$};
\node[dp-hd] at (1,0.35) {$j=1$};
\node[dp-hd] at (2,0.35) {$2$};
\node[dp-hd] at (3,0.35) {$3$};
\node[dp-hd] at (4,0.35) {$4$};
\node[dp-hd] at (5,0.35) {$5$};
\node[dp-hd, anchor=east] at (0.45,-0.5) {$i=1$};
\node[dp-cd0, minimum width=0.8cm] at (1,-0.5) {1};
\node[dp-cd1, minimum width=0.8cm] at (2,-0.5) {2};
\node[dp-cd0, minimum width=0.8cm] at (3,-0.5) {2};
\node[dp-cd1, minimum width=0.8cm] at (4,-0.5) {2};
\node[dp-cd0, minimum width=0.8cm] at (5,-0.5) {4};
\node[dp-hd, anchor=east] at (0.45,-1.5) {$2$};
\node[dp-cd0, minimum width=0.8cm] at (2,-1.5) {2};
\node[dp-cd1, minimum width=0.8cm] at (3,-1.5) {2};
\node[dp-cd0, minimum width=0.8cm] at (4,-1.5) {2};
\node[dp-cd1, minimum width=0.8cm] at (5,-1.5) {4};
\node[dp-hd, anchor=east] at (0.45,-2.5) {$3$};
\node[dp-cd0, minimum width=0.8cm] at (3,-2.5) {3};
\node[dp-cd1, minimum width=0.8cm] at (4,-2.5) {4};
\node[dp-cd0, minimum width=0.8cm] at (5,-2.5) {4};
\node[dp-hd, anchor=east] at (0.45,-3.5) {$4$};
\node[dp-cd0, minimum width=0.8cm] at (4,-3.5) {4};
\node[dp-cd1, minimum width=0.8cm] at (5,-3.5) {5};
\node[dp-hd, anchor=east] at (0.45,-4.5) {$5$};
\node[dp-cd0, minimum width=0.8cm] at (5,-4.5) {5};
\end{tikzpicture}
\diagcaption{$e$ 和 $\mathit{root}$ 兩張表（$w$ 的值見上表第二欄）。 灰底的是 $j = i - 1$ 的空樹，底色深淺交錯的是不
同長度的區間。 答案是 $e[1, 5] = 109$。}
\label{fig:dp-obst-tables}
\end{minipage}
\end{center}

\textbf{由 root 表建樹}：$\mathit{root}[1, 5] = 4$，根是 $k_4 = \texttt{int}$；左邊是
$k_1..k_3$，$\mathit{root}[1, 3] = 2$，根是 \texttt{for}，它的左右各是 $\mathit{root}[1,1] = 1$
的 \texttt{char} 和 $\mathit{root}[3,3] = 3$ 的 \texttt{goto}；右邊是 $k_5..k_5$，就是
\texttt{while}。 這正是圖~\ref{fig:dp-obst-trees}(b)。 和矩陣鏈乘的 \proc{Print-Parens} 一
樣，是一個 $O(n)$ 的遞迴。

表裡有一個現象：每一列的 $\mathit{root}$ 由左往右不會變小，每一欄由上往下也不會變小。 Knuth
在 1971 年證明這對任何次數都成立：
\[
\mathit{root}[i, j - 1] \le \mathit{root}[i, j] \le \mathit{root}[i + 1, j] ,
\]
所以第 10 行的 $r$ 只要從 $\mathit{root}[i, j-1]$ 試到 $\mathit{root}[i+1, j]$，共
$\mathit{root}[i+1, j] - \mathit{root}[i, j-1] + 1$ 個。 固定區間長度 $\ell$，把 $i = 1, 2,
\ldots, n - \ell + 1$ 這些個數加起來（$j = i + \ell - 1$ 跟著 $i$ 變）：第 $i$ 項加的
$\mathit{root}[i+1, j]$，正好是第 $i + 1$ 項減的 $\mathit{root}[(i+1), (j+1) - 1]$，兩兩抵銷，最
後只剩 $\mathit{root}[n - \ell + 2, n] - \mathit{root}[1, \ell - 1]$ 和 $n - \ell + 1$ 個 $1$，不超過
$n + (n - \ell + 1) \le 2n$。 每個 $\ell$ 是 $O(n)$，$n$ 個 $\ell$ 合起來，總時間
降到 $\Theta(n^2)$。 例如 $e[1, 5]$ 只要試
$r = \mathit{root}[1, 4] = 2$ 到 $\mathit{root}[2, 5] = 4$ 這三個。

\begin{examnote}
最佳二元搜尋樹的題目，先看清楚題目的成本怎麼算。 本書數「和鍵比較的次數」：鍵的成本是深度
加一，外部節點的成本是深度，所以空樹 $e[i, i-1] = 0$。 CLRS 把走到外部節點也算一次，外部
節點的成本是深度加一，所以它的 $e[i, i-1] = q_{i-1}$，每一格的 $e$ 都比本書多
$q_{i-1} + q_i + \cdots + q_j$；本節的例子照 CLRS 的成本模型是 $109 + 19 = 128$。 另一種常見
的寫法（例如 Horowitz、Sahni 等人的 \textit{Computer Algorithms}）用 $c(i, j)$ 表示
$a_{i+1}, \ldots, a_j$ 這一段，$c(i, i) = 0$、$w(i, i) = q_i$，數法和本書相同，只是索引差一。 幾種成本模型的 $\mathit{root}$ 表和最佳樹都一樣（相差的是一個和
樹的形狀無關的常數），但 $e$ 的數值不同，答案要照題目的定義寫，必要時寫一句「依題意，查
不到時不計走到外部節點那一步」。
\end{examnote}

\section{0/1 背包}
\label{sec:dp-knapsack}

\subsection{問題}

期末專題還剩 $10$ 天。 組員列出五個還沒做的功能，每個功能要花的天數 $w_i$ 和做完能加的分數
$v_i$ 如下；功能只能整個做完或整個不做，總天數不能超過 $10$ 天。 要做哪幾個，分數最高？

\begin{center}
\begin{tabular}{c*{5}{c}}
\toprule
功能 $i$ & 1 & 2 & 3 & 4 & 5 \\
\midrule
天數 $w_i$ & 3 & 2 & 4 & 5 & 2 \\
分數 $v_i$ & 10 & 7 & 11 & 13 & 6 \\
每天的分數 $v_i / w_i$ & 3.33 & 3.5 & 2.75 & 2.6 & 3 \\
\bottomrule
\end{tabular}
\end{center}

一般地，\textbf{0/1 背包問題}（0/1 knapsack problem）：有 $n$ 件物品，第 $i$ 件重 $w_i$、價值
$v_i$（都是正整數），背包容量 $W$。 每件物品拿或不拿（所以叫 0/1），總重量不超過 $W$，求
總價值最大。

先試一個直覺的做法：照「每天的分數」由高到低挑，放得下就做。 依序是功能 2、1、5（共 $7$
天、$23$ 分），剩下 $3$ 天，功能 3 和 4 都放不下，得到 $23$ 分。 可是功能 1、2、4 剛好
$10$ 天、$30$ 分。 照比值挑的貪心法對 0/1 背包不對；物品可以只拿一部分時（分數背包），它
才是對的。 第~19~章會證明分數背包的貪心法，並把兩者放在一起比較。 全部試一遍要 $2^n$ 種組
合，下面用動態規劃。

\subsection{遞迴式}

看最後一件物品 $n$。 最佳解不是不拿它，就是拿它：
\begin{itemize}
\item 不拿：剩下的就是「前 $n - 1$ 件、容量 $W$」的最佳解。
\item 拿（需要 $w_n \le W$）：其餘的物品要在容量 $W - w_n$ 裡挑，而且一定是「前 $n - 1$ 件、
容量 $W - w_n$」的最佳解，否則換成更好的那組，整體就更好（剪下貼上）。
\end{itemize}
子問題要同時記住「考慮到第幾件」和「剩多少容量」。 令 $K[i, c]$ 為只用前 $i$ 件物品、容量
$c$ 時的最大價值，$0 \le i \le n$，$0 \le c \le W$：
\[
K[i, c] = \begin{cases}
0, & i = 0, \\
K[i - 1, c], & i \ge 1 \text{ 且 } w_i > c, \\
\max\bigl(K[i - 1, c],\ K[i - 1, c - w_i] + v_i\bigr), & i \ge 1 \text{ 且 } w_i \le c .
\end{cases}
\]
答案是 $K[n, W]$。 第 $i$ 列只用到第 $i - 1$ 列，所以一列一列往下填。

\begin{pseudo}[label=alg:dp-knapsack]{0/1 背包的填表}
\Function{Knapsack}{$w, v, n, W$}
    \For{$c \gets 0$ \To $W$}
        \State $K[0, c] \gets 0$
    \EndFor
    \For{$i \gets 1$ \To $n$}
        \For{$c \gets 0$ \To $W$}
            \State $K[i, c] \gets K[i - 1, c]$ \Comment{不拿第 $i$ 件}
            \If{$w_i \le c$ \AND $K[i - 1, c - w_i] + v_i > K[i, c]$}
                \State $K[i, c] \gets K[i - 1, c - w_i] + v_i$ \Comment{拿第 $i$ 件}
            \EndIf
        \EndFor
    \EndFor
    \State \Return $K$
\EndFunction
\end{pseudo}

\textbf{還原}：從 $K[n, W]$ 往上看。 $K[i, c] \ne K[i - 1, c]$ 表示第 $i$ 件一定有拿（不拿的
話值會和上一列一樣），於是容量減掉 $w_i$ 再看上一列；相等就表示不拿第 $i$ 件也能達到，照第 7
行的 $>$，本書這時當作不拿。

\begin{pseudo}[label=alg:dp-knapsack-items]{還原要拿的物品}
\Procedure{Knapsack-Items}{$K, w, n, W$}
    \State $c \gets W$
    \For{$i \gets n$ \DownTo $1$}
        \If{$K[i, c] \ne K[i - 1, c]$}
            \State 印出 $i$ \Comment{第 $i$ 件有拿}
            \State $c \gets c - w_i$
        \EndIf
    \EndFor
\EndProcedure
\end{pseudo}

\subsection{手算追蹤}
\label{sec:dp-knapsack-trace}

\begin{center}
\begin{minipage}{\linewidth}
\centering
\begin{tikzpicture}[x=0.72cm, y=0.6cm]
\node[dp-hd, anchor=east] at (-0.5,0.9) {$c$};
\node[dp-hd] at (0,0.9) {0};
\node[dp-hd] at (1,0.9) {1};
\node[dp-hd] at (2,0.9) {2};
\node[dp-hd] at (3,0.9) {3};
\node[dp-hd] at (4,0.9) {4};
\node[dp-hd] at (5,0.9) {5};
\node[dp-hd] at (6,0.9) {6};
\node[dp-hd] at (7,0.9) {7};
\node[dp-hd] at (8,0.9) {8};
\node[dp-hd] at (9,0.9) {9};
\node[dp-hd] at (10,0.9) {10};
\node[dp-hd, anchor=east] at (-0.45,0) {$i=0$};
\node[dp-cpath, minimum width=0.72cm, minimum height=0.6cm] at (0,0) {0};
\node[dp-c, minimum width=0.72cm, minimum height=0.6cm] at (1,0) {0};
\node[dp-c, minimum width=0.72cm, minimum height=0.6cm] at (2,0) {0};
\node[dp-c, minimum width=0.72cm, minimum height=0.6cm] at (3,0) {0};
\node[dp-c, minimum width=0.72cm, minimum height=0.6cm] at (4,0) {0};
\node[dp-c, minimum width=0.72cm, minimum height=0.6cm] at (5,0) {0};
\node[dp-c, minimum width=0.72cm, minimum height=0.6cm] at (6,0) {0};
\node[dp-c, minimum width=0.72cm, minimum height=0.6cm] at (7,0) {0};
\node[dp-c, minimum width=0.72cm, minimum height=0.6cm] at (8,0) {0};
\node[dp-c, minimum width=0.72cm, minimum height=0.6cm] at (9,0) {0};
\node[dp-c, minimum width=0.72cm, minimum height=0.6cm] at (10,0) {0};
\node[dp-hd, anchor=east] at (-0.45,-1) {$i=1$ (3,10)};
\node[dp-c, minimum width=0.72cm, minimum height=0.6cm] at (0,-1) {0};
\node[dp-c, minimum width=0.72cm, minimum height=0.6cm] at (1,-1) {0};
\node[dp-c, minimum width=0.72cm, minimum height=0.6cm] at (2,-1) {0};
\node[dp-chit, minimum width=0.72cm, minimum height=0.6cm] at (3,-1) {10};
\node[dp-c, minimum width=0.72cm, minimum height=0.6cm] at (4,-1) {10};
\node[dp-c, minimum width=0.72cm, minimum height=0.6cm] at (5,-1) {10};
\node[dp-c, minimum width=0.72cm, minimum height=0.6cm] at (6,-1) {10};
\node[dp-c, minimum width=0.72cm, minimum height=0.6cm] at (7,-1) {10};
\node[dp-c, minimum width=0.72cm, minimum height=0.6cm] at (8,-1) {10};
\node[dp-c, minimum width=0.72cm, minimum height=0.6cm] at (9,-1) {10};
\node[dp-c, minimum width=0.72cm, minimum height=0.6cm] at (10,-1) {10};
\node[dp-hd, anchor=east] at (-0.45,-2) {$i=2$ (2,7)};
\node[dp-c, minimum width=0.72cm, minimum height=0.6cm] at (0,-2) {0};
\node[dp-c, minimum width=0.72cm, minimum height=0.6cm] at (1,-2) {0};
\node[dp-c, minimum width=0.72cm, minimum height=0.6cm] at (2,-2) {7};
\node[dp-c, minimum width=0.72cm, minimum height=0.6cm] at (3,-2) {10};
\node[dp-c, minimum width=0.72cm, minimum height=0.6cm] at (4,-2) {10};
\node[dp-chit, minimum width=0.72cm, minimum height=0.6cm] at (5,-2) {17};
\node[dp-c, minimum width=0.72cm, minimum height=0.6cm] at (6,-2) {17};
\node[dp-c, minimum width=0.72cm, minimum height=0.6cm] at (7,-2) {17};
\node[dp-c, minimum width=0.72cm, minimum height=0.6cm] at (8,-2) {17};
\node[dp-c, minimum width=0.72cm, minimum height=0.6cm] at (9,-2) {17};
\node[dp-c, minimum width=0.72cm, minimum height=0.6cm] at (10,-2) {17};
\node[dp-hd, anchor=east] at (-0.45,-3) {$i=3$ (4,11)};
\node[dp-c, minimum width=0.72cm, minimum height=0.6cm] at (0,-3) {0};
\node[dp-c, minimum width=0.72cm, minimum height=0.6cm] at (1,-3) {0};
\node[dp-c, minimum width=0.72cm, minimum height=0.6cm] at (2,-3) {7};
\node[dp-c, minimum width=0.72cm, minimum height=0.6cm] at (3,-3) {10};
\node[dp-c, minimum width=0.72cm, minimum height=0.6cm] at (4,-3) {11};
\node[dp-cpath, minimum width=0.72cm, minimum height=0.6cm] at (5,-3) {17};
\node[dp-c, minimum width=0.72cm, minimum height=0.6cm] at (6,-3) {18};
\node[dp-c, minimum width=0.72cm, minimum height=0.6cm] at (7,-3) {21};
\node[dp-c, minimum width=0.72cm, minimum height=0.6cm] at (8,-3) {21};
\node[dp-c, minimum width=0.72cm, minimum height=0.6cm] at (9,-3) {28};
\node[dp-c, minimum width=0.72cm, minimum height=0.6cm] at (10,-3) {28};
\node[dp-hd, anchor=east] at (-0.45,-4) {$i=4$ (5,13)};
\node[dp-c, minimum width=0.72cm, minimum height=0.6cm] at (0,-4) {0};
\node[dp-c, minimum width=0.72cm, minimum height=0.6cm] at (1,-4) {0};
\node[dp-c, minimum width=0.72cm, minimum height=0.6cm] at (2,-4) {7};
\node[dp-c, minimum width=0.72cm, minimum height=0.6cm] at (3,-4) {10};
\node[dp-c, minimum width=0.72cm, minimum height=0.6cm] at (4,-4) {11};
\node[dp-c, minimum width=0.72cm, minimum height=0.6cm] at (5,-4) {17};
\node[dp-c, minimum width=0.72cm, minimum height=0.6cm] at (6,-4) {18};
\node[dp-c, minimum width=0.72cm, minimum height=0.6cm] at (7,-4) {21};
\node[dp-c, minimum width=0.72cm, minimum height=0.6cm] at (8,-4) {23};
\node[dp-c, minimum width=0.72cm, minimum height=0.6cm] at (9,-4) {28};
\node[dp-chit, minimum width=0.72cm, minimum height=0.6cm] at (10,-4) {30};
\node[dp-hd, anchor=east] at (-0.45,-5) {$i=5$ (2,6)};
\node[dp-c, minimum width=0.72cm, minimum height=0.6cm] at (0,-5) {0};
\node[dp-c, minimum width=0.72cm, minimum height=0.6cm] at (1,-5) {0};
\node[dp-c, minimum width=0.72cm, minimum height=0.6cm] at (2,-5) {7};
\node[dp-c, minimum width=0.72cm, minimum height=0.6cm] at (3,-5) {10};
\node[dp-c, minimum width=0.72cm, minimum height=0.6cm] at (4,-5) {13};
\node[dp-c, minimum width=0.72cm, minimum height=0.6cm] at (5,-5) {17};
\node[dp-c, minimum width=0.72cm, minimum height=0.6cm] at (6,-5) {18};
\node[dp-c, minimum width=0.72cm, minimum height=0.6cm] at (7,-5) {23};
\node[dp-c, minimum width=0.72cm, minimum height=0.6cm] at (8,-5) {24};
\node[dp-c, minimum width=0.72cm, minimum height=0.6cm] at (9,-5) {28};
\node[dp-cpath, minimum width=0.72cm, minimum height=0.6cm] at (10,-5) {30};
\end{tikzpicture}
\diagcaption{專題例子的 $K$ 表。 列左邊的括號是 $(w_i, v_i)$。 從 $K[5, 10] = 30$ 往上還原，
經過的格子加框；粉紅色的格子和上一列不同，表示拿了那一列的物品：功能 4、2、1。}
\label{fig:dp-knapsack-table}
\end{minipage}
\end{center}

幾格是這樣算的：
\begin{itemize}
\item $K[2, 5]$：功能 2 重 $2$，$\max(K[1, 5], K[1, 3] + 7) = \max(10, 17) = 17$，也就是功能
1、2 都做。
\item $K[3, 7]$：$\max(K[2, 7], K[2, 3] + 11) = \max(17, 21) = 21$，做功能 1、3。
\item $K[4, 10]$：$\max(K[3, 10], K[3, 5] + 13) = \max(28, 30) = 30$。
\item $K[5, 10]$：$\max(K[4, 10], K[4, 8] + 6) = \max(30, 29) = 30$，不做功能 5。
\end{itemize}

還原：$K[5, 10] = K[4, 10]$，不做功能 5；$K[4, 10] = 30 \ne K[3, 10] = 28$，做功能 4，容量剩
$5$；$K[3, 5] = K[2, 5] = 17$，不做功能 3；$K[2, 5] = 17 \ne K[1, 5] = 10$，做功能 2，剩 $3$；
$K[1, 3] = 10 \ne K[0, 3] = 0$，做功能 1，剩 $0$。 答案是功能 1、2、4，$10$ 天、$30$ 分。

\subsection{一維滾動陣列}
\label{sec:dp-knapsack-1d}

第 $i$ 列只用到第 $i - 1$ 列，所以和第~\ref{sec:dp-bottomup}~節的巴斯卡三角形一樣，可以只留一
列 $d[0..W]$，就地從第 $i - 1$ 列改成第 $i$ 列。 更新 $d[c]$ 時要用的是\textbf{舊的}
$d[c - w_i]$（第 $i - 1$ 列的值），而 $c - w_i < c$，所以容量要\textbf{由大到小}跑：改 $d[c]$
的時候，比它小的格子都還沒動過。

\begin{pseudo}[label=alg:dp-knapsack-1d]{0/1 背包，一維陣列}
\Function{Knapsack-1D}{$w, v, n, W$}
    \For{$c \gets 0$ \To $W$}
        \State $d[c] \gets 0$
    \EndFor
    \For{$i \gets 1$ \To $n$}
        \For{$c \gets W$ \DownTo $w_i$}
            \State $d[c] \gets \max(d[c],\ d[c - w_i] + v_i)$
        \EndFor
    \EndFor
    \State \Return $d[W]$
\EndFunction
\end{pseudo}

$c < w_i$ 的格子放不下第 $i$ 件，值和上一列一樣，所以第 5 行的迴圈只跑到 $w_i$。 空間從
$\Theta(nW)$ 降到 $\Theta(W)$，時間不變；但和 LCS 只留兩列一樣，表被蓋掉之後就不能用
\proc{Knapsack-Items} 還原了。

如果第 5 行寫成 $c \gets w_i$ \To $W$ 由小到大，算 $d[c]$ 時 $d[c - w_i]$ 可能已經是這一列的
新值，也就是\textbf{已經拿過第 $i$ 件}的結果，於是同一件物品會被拿好幾次。 只看功能 2：
$d[2] = 7$，接著 $d[4] = d[2] + 7 = 14$、$d[6] = 21$、……、$d[10] = 35$，等於把功能 2 做了
五次。 這個「錯誤」的版本其實正好解另一個問題：每種物品可以拿\textbf{任意多個}的
\textbf{無限背包}（unbounded knapsack），它的遞迴式是
\[
U[c] = \max\Bigl(0,\ \max_{i : w_i \le c} \bigl(U[c - w_i] + v_i\bigr)\Bigr),
\]
由小到大算 $U[c]$ 時，$U[c - w_i]$ 本來就該是已經考慮過所有物品（包括第 $i$ 件）的值。 兩
個問題只差一個迴圈的方向，這一點很容易考，也很容易寫錯。

\subsection{擬多項式時間}
\label{sec:dp-pseudo-poly}

\proc{Knapsack} 的表有 $(n + 1)(W + 1)$ 格，每格 $O(1)$，時間 $\Theta(nW)$。 這看起來是多項式，
但第~\ref{sec:ana-input-size}~節說過，輸入的大小要用寫下輸入所需的位元數來量。 $W$
這個數只要 $b = \lceil \lg(W + 1) \rceil$ 個位元就寫得下，所以 $\Theta(nW) = \Theta(n 2^b)$
是輸入長度的\textbf{指數}函數。 容量從 $10$ 改成 $10^{9}$，輸入只多了幾個字元，表卻大了一億
倍。 這種「對輸入的數值是多項式、對位元數是指數」的時間，就是\textbf{擬多項式時間}。

$W$ 不大時（例如練習題的 $W \le 10^5$），$\Theta(nW)$ 非常好用；$W$ 很大時這張表就開不出來。
0/1 背包的決策版本是 NP 完備問題（第~28~章），目前沒有已知的多項式時間演算法，所以擬多項
式已經是很好的結果。 另一種 DP 改用「達到價值 $v$ 最少要多重」當子問題，時間是
$\Theta(n \sum v_i)$，價值小、重量大時比較好；第~29~章的近似演算法就從這個版本出發。 第~20~
章則用分支限界在不填表的情況下搜尋。

\subsection{C 程式}

\texttt{knap\_2d} 填完整的表並還原；\texttt{knap\_1d} 用一維陣列，參數
\texttt{down} 決定容量的方向。 為了和虛擬碼對照，物品從 $1$ 開始編號，\texttt{w[0]}、
\texttt{v[0]} 不用。

\begin{example}[label=ex:dp-knapsack]{c}{(knapsack.c) 0/1 背包的二維表與一維陣列}
#include <stdio.h>
#include <stdlib.h>

#define MAXN 100

// 二維表 K[i][c]：只用前 i 件、容量 c 的最大價值。take[i] 記錄有沒有拿第 i 件。
long long knap_2d(int n, int cap, const int *w, const int *v, int *take)
{
    long long *K = calloc((size_t)(n + 1) * (cap + 1), sizeof *K);
    if (K == NULL) {
        return -1;
    }
#define T(i, c) K[(i) * (cap + 1) + (c)]
    for (int i = 1; i <= n; i++) {
        for (int c = 0; c <= cap; c++) {
            T(i, c) = T(i - 1, c);                      // 不拿第 i 件
            if (w[i] <= c && T(i - 1, c - w[i]) + v[i] > T(i, c)) {
                T(i, c) = T(i - 1, c - w[i]) + v[i];    // 拿第 i 件
            }
        }
    }
    long long best = T(n, cap);
    for (int i = n, c = cap; i >= 1; i--) {             // 從最後一列往上還原
        take[i] = (T(i, c) != T(i - 1, c));
        if (take[i]) {
            c -= w[i];
        }
    }
#undef T
    free(K);
    return best;
}

// 一維陣列 d[c]。down = 1 時容量由大到小（0/1 背包），0 時由小到大。
long long knap_1d(int n, int cap, const int *w, const int *v, int down)
{
    long long *d = calloc(cap + 1, sizeof *d);
    if (d == NULL) {
        return -1;
    }
    for (int i = 1; i <= n; i++) {
        for (int t = 0; t <= cap - w[i]; t++) {
            int c = down ? cap - t : w[i] + t;
            if (d[c - w[i]] + v[i] > d[c]) {
                d[c] = d[c - w[i]] + v[i];
            }
        }
    }
    long long best = d[cap];
    free(d);
    return best;
}

int main(void)
{
    int n, cap, w[MAXN + 1], v[MAXN + 1], take[MAXN + 1];
    if (scanf("%d %d", &n, &cap) != 2 || n < 1 || n > MAXN || cap < 0) {
        return 1;
    }
    for (int i = 1; i <= n; i++) {
        if (scanf("%d", &w[i]) != 1) {
            return 1;
        }
    }
    for (int i = 1; i <= n; i++) {
        if (scanf("%d", &v[i]) != 1) {
            return 1;
        }
    }
    printf("2D table: best = %lld, items =", knap_2d(n, cap, w, v, take));
    for (int i = 1; i <= n; i++) {
        if (take[i]) {
            printf(" %d", i);
        }
    }
    printf("\n1D, capacity downward: best = %lld\n", knap_1d(n, cap, w, v, 1));
    printf("1D, capacity upward: best = %lld\n", knap_1d(n, cap, w, v, 0));
    return 0;
}
\end{example}
\iosample{5 10 \\ 3 2 4 5 2 \\ 10 7 11 13 6}{2D table: best = 30, items = 1 2 4 \\ 1D, capacity downward: best = 30 \\ 1D, capacity upward: best = 35}

由大到小的一維版本和二維表一樣是 $30$；由小到大的版本得到 $35$，是把功能 2 做五次的無限背
包答案，對 0/1 背包來說是錯的。 一維迴圈為了同時支援兩個方向，用 \texttt{t} 從 $0$ 數到
$W - w_i$，再換算成容量 \texttt{c}；只寫 0/1 背包的話，直接寫
\texttt{for (int c = cap; c >= w[i]; c--)} 就好。

\section{編輯距離與最長遞增子序列}
\label{sec:dp-edit-lis}

這一節的兩個問題，子問題都是「前綴」，和 LCS 是同一類。

\subsection{編輯距離}
\label{sec:dp-edit}

拼字檢查要找「和打錯的字最接近的正確單字」，就要量兩個字串有多接近。 把字串 $X$ 改成 $Y$，
每一步可以做三種操作之一：
\begin{itemize}
\item \textbf{插入}一個字元；
\item \textbf{刪除}一個字元；
\item 把一個字元\textbf{替換}成另一個字元。
\end{itemize}
每種操作的成本都是 $1$，最少的操作次數叫 $X$ 到 $Y$ 的\textbf{編輯距離}（edit distance），也
叫 Levenshtein 距離。 例如 \texttt{PARTNER} 改成 \texttt{PRINTER}：刪掉 \texttt{A}、把
\texttt{T} 換成 \texttt{I}、在 \texttt{N} 後面插入 \texttt{T}，三步。 不可能更少，下面的表會
證明這件事。

把改的過程寫成上下兩列的\textbf{對齊}比較好想：
\begin{center}
\texttt{P A R T N - E R}\\
\texttt{P - R I N T E R}
\end{center}
每一行是三種之一：上下都有字元（相同就不花成本，不同是一次替換）、只有上面（刪除）、只有下面
（插入）。 編輯距離就是最好的對齊裡「不是相同字元」的行數。 看對齊的\textbf{最後一行}：
\begin{itemize}
\item 上面是 $x_m$、下面是 $y_n$：前面是 $X_{m-1}$ 和 $Y_{n-1}$ 的對齊，成本加 $0$ 或 $1$；
\item 上面是 $x_m$、下面空著（刪除 $x_m$）：前面是 $X_{m-1}$ 和 $Y_n$ 的對齊，成本加 $1$；
\item 上面空著、下面是 $y_n$（插入 $y_n$）：前面是 $X_m$ 和 $Y_{n-1}$ 的對齊，成本加 $1$。
\end{itemize}
前面那一段當然要用最好的對齊（剪下貼上）。 令 $D[i, j]$ 為 $X_i$ 改成 $Y_j$ 的編輯距離：
\[
D[i, j] = \min\begin{cases}
D[i - 1, j - 1] + \delta(x_i, y_j), \\
D[i - 1, j] + 1, \\
D[i, j - 1] + 1,
\end{cases}
\qquad
\delta(x_i, y_j) = \begin{cases} 0, & x_i = y_j, \\ 1, & x_i \ne y_j, \end{cases}
\]
邊界是 $D[i, 0] = i$（刪掉 $i$ 個字元）、$D[0, j] = j$（插入 $j$ 個字元）。 和 LCS 不同，
$x_i = y_j$ 時三項還是要一起比；不過可以證明這時取左上角一定不會比較差（觀念題
第~\ref{dp-q-edit-match}~題）。

\begin{pseudo}[label=alg:dp-edit]{編輯距離}
\Function{Edit-Distance}{$X, Y, m, n$}
    \For{$i \gets 0$ \To $m$}
        \State $D[i, 0] \gets i$
    \EndFor
    \For{$j \gets 0$ \To $n$}
        \State $D[0, j] \gets j$
    \EndFor
    \For{$i \gets 1$ \To $m$}
        \For{$j \gets 1$ \To $n$}
            \State $\delta \gets 0$；若 $x_i \ne y_j$，$\delta \gets 1$
            \State $D[i, j] \gets \min(D[i - 1, j - 1] + \delta,\ D[i - 1, j] + 1,\ D[i, j - 1] + 1)$
        \EndFor
    \EndFor
    \State \Return $D[m, n]$
\EndFunction
\end{pseudo}

時間和空間都是 $\Theta(mn)$；只要距離的話，同樣可以只留兩列。

\begin{center}
\begin{minipage}{\linewidth}
\centering
\begin{tikzpicture}[x=0.78cm, y=0.6cm]
\node[dp-hd] at (0,1.55) {0};
\node[dp-hd] at (1,1.55) {1};
\node[dp-hd] at (1,0.85) {\texttt{P}};
\node[dp-hd] at (2,1.55) {2};
\node[dp-hd] at (2,0.85) {\texttt{R}};
\node[dp-hd] at (3,1.55) {3};
\node[dp-hd] at (3,0.85) {\texttt{I}};
\node[dp-hd] at (4,1.55) {4};
\node[dp-hd] at (4,0.85) {\texttt{N}};
\node[dp-hd] at (5,1.55) {5};
\node[dp-hd] at (5,0.85) {\texttt{T}};
\node[dp-hd] at (6,1.55) {6};
\node[dp-hd] at (6,0.85) {\texttt{E}};
\node[dp-hd] at (7,1.55) {7};
\node[dp-hd] at (7,0.85) {\texttt{R}};
\node[dp-hd] at (-1.2,0) {0};
\node[dp-cpath, minimum width=0.78cm, minimum height=0.6cm] at (0,0) {0};
\node[dp-c0, minimum width=0.78cm, minimum height=0.6cm] at (1,0) {1};
\node[dp-c0, minimum width=0.78cm, minimum height=0.6cm] at (2,0) {2};
\node[dp-c0, minimum width=0.78cm, minimum height=0.6cm] at (3,0) {3};
\node[dp-c0, minimum width=0.78cm, minimum height=0.6cm] at (4,0) {4};
\node[dp-c0, minimum width=0.78cm, minimum height=0.6cm] at (5,0) {5};
\node[dp-c0, minimum width=0.78cm, minimum height=0.6cm] at (6,0) {6};
\node[dp-c0, minimum width=0.78cm, minimum height=0.6cm] at (7,0) {7};
\node[dp-hd] at (-1.2,-1) {1};
\node[dp-hd] at (-0.65,-1) {\texttt{P}};
\node[dp-c0, minimum width=0.78cm, minimum height=0.6cm] at (0,-1) {1};
\node[dp-cpath, minimum width=0.78cm, minimum height=0.6cm] at (1,-1) {0};
\node[dp-c, minimum width=0.78cm, minimum height=0.6cm] at (2,-1) {1};
\node[dp-c, minimum width=0.78cm, minimum height=0.6cm] at (3,-1) {2};
\node[dp-c, minimum width=0.78cm, minimum height=0.6cm] at (4,-1) {3};
\node[dp-c, minimum width=0.78cm, minimum height=0.6cm] at (5,-1) {4};
\node[dp-c, minimum width=0.78cm, minimum height=0.6cm] at (6,-1) {5};
\node[dp-c, minimum width=0.78cm, minimum height=0.6cm] at (7,-1) {6};
\node[dp-hd] at (-1.2,-2) {2};
\node[dp-hd] at (-0.65,-2) {\texttt{A}};
\node[dp-c0, minimum width=0.78cm, minimum height=0.6cm] at (0,-2) {2};
\node[dp-cpath, minimum width=0.78cm, minimum height=0.6cm] at (1,-2) {1};
\node[dp-c, minimum width=0.78cm, minimum height=0.6cm] at (2,-2) {1};
\node[dp-c, minimum width=0.78cm, minimum height=0.6cm] at (3,-2) {2};
\node[dp-c, minimum width=0.78cm, minimum height=0.6cm] at (4,-2) {3};
\node[dp-c, minimum width=0.78cm, minimum height=0.6cm] at (5,-2) {4};
\node[dp-c, minimum width=0.78cm, minimum height=0.6cm] at (6,-2) {5};
\node[dp-c, minimum width=0.78cm, minimum height=0.6cm] at (7,-2) {6};
\node[dp-hd] at (-1.2,-3) {3};
\node[dp-hd] at (-0.65,-3) {\texttt{R}};
\node[dp-c0, minimum width=0.78cm, minimum height=0.6cm] at (0,-3) {3};
\node[dp-c, minimum width=0.78cm, minimum height=0.6cm] at (1,-3) {2};
\node[dp-cpath, minimum width=0.78cm, minimum height=0.6cm] at (2,-3) {1};
\node[dp-c, minimum width=0.78cm, minimum height=0.6cm] at (3,-3) {2};
\node[dp-c, minimum width=0.78cm, minimum height=0.6cm] at (4,-3) {3};
\node[dp-c, minimum width=0.78cm, minimum height=0.6cm] at (5,-3) {4};
\node[dp-c, minimum width=0.78cm, minimum height=0.6cm] at (6,-3) {5};
\node[dp-c, minimum width=0.78cm, minimum height=0.6cm] at (7,-3) {5};
\node[dp-hd] at (-1.2,-4) {4};
\node[dp-hd] at (-0.65,-4) {\texttt{T}};
\node[dp-c0, minimum width=0.78cm, minimum height=0.6cm] at (0,-4) {4};
\node[dp-c, minimum width=0.78cm, minimum height=0.6cm] at (1,-4) {3};
\node[dp-c, minimum width=0.78cm, minimum height=0.6cm] at (2,-4) {2};
\node[dp-cpath, minimum width=0.78cm, minimum height=0.6cm] at (3,-4) {2};
\node[dp-c, minimum width=0.78cm, minimum height=0.6cm] at (4,-4) {3};
\node[dp-c, minimum width=0.78cm, minimum height=0.6cm] at (5,-4) {3};
\node[dp-c, minimum width=0.78cm, minimum height=0.6cm] at (6,-4) {4};
\node[dp-c, minimum width=0.78cm, minimum height=0.6cm] at (7,-4) {5};
\node[dp-hd] at (-1.2,-5) {5};
\node[dp-hd] at (-0.65,-5) {\texttt{N}};
\node[dp-c0, minimum width=0.78cm, minimum height=0.6cm] at (0,-5) {5};
\node[dp-c, minimum width=0.78cm, minimum height=0.6cm] at (1,-5) {4};
\node[dp-c, minimum width=0.78cm, minimum height=0.6cm] at (2,-5) {3};
\node[dp-c, minimum width=0.78cm, minimum height=0.6cm] at (3,-5) {3};
\node[dp-cpath, minimum width=0.78cm, minimum height=0.6cm] at (4,-5) {2};
\node[dp-cpath, minimum width=0.78cm, minimum height=0.6cm] at (5,-5) {3};
\node[dp-c, minimum width=0.78cm, minimum height=0.6cm] at (6,-5) {4};
\node[dp-c, minimum width=0.78cm, minimum height=0.6cm] at (7,-5) {5};
\node[dp-hd] at (-1.2,-6) {6};
\node[dp-hd] at (-0.65,-6) {\texttt{E}};
\node[dp-c0, minimum width=0.78cm, minimum height=0.6cm] at (0,-6) {6};
\node[dp-c, minimum width=0.78cm, minimum height=0.6cm] at (1,-6) {5};
\node[dp-c, minimum width=0.78cm, minimum height=0.6cm] at (2,-6) {4};
\node[dp-c, minimum width=0.78cm, minimum height=0.6cm] at (3,-6) {4};
\node[dp-c, minimum width=0.78cm, minimum height=0.6cm] at (4,-6) {3};
\node[dp-c, minimum width=0.78cm, minimum height=0.6cm] at (5,-6) {3};
\node[dp-cpath, minimum width=0.78cm, minimum height=0.6cm] at (6,-6) {3};
\node[dp-c, minimum width=0.78cm, minimum height=0.6cm] at (7,-6) {4};
\node[dp-hd] at (-1.2,-7) {7};
\node[dp-hd] at (-0.65,-7) {\texttt{R}};
\node[dp-c0, minimum width=0.78cm, minimum height=0.6cm] at (0,-7) {7};
\node[dp-c, minimum width=0.78cm, minimum height=0.6cm] at (1,-7) {6};
\node[dp-c, minimum width=0.78cm, minimum height=0.6cm] at (2,-7) {5};
\node[dp-c, minimum width=0.78cm, minimum height=0.6cm] at (3,-7) {5};
\node[dp-c, minimum width=0.78cm, minimum height=0.6cm] at (4,-7) {4};
\node[dp-c, minimum width=0.78cm, minimum height=0.6cm] at (5,-7) {4};
\node[dp-c, minimum width=0.78cm, minimum height=0.6cm] at (6,-7) {4};
\node[dp-cpath, minimum width=0.78cm, minimum height=0.6cm] at (7,-7) {3};
\end{tikzpicture}
\diagcaption{\texttt{PARTNER}（列）改成 \texttt{PRINTER}（欄）的 $D$ 表，答案 $D[7, 7] = 3$。
加框的是還原時走過的格子。}
\label{fig:dp-edit-table}
\end{minipage}
\end{center}

\textbf{還原。} 從 $[m, n]$ 往回走，每一步看這一格的值是由哪一項得到的，依序檢查：左上（相同
或替換）、上面（刪除 $x_i$）、左邊（插入 $y_j$），第一個符合的就走。 這組字串：
$[7,7] \to [6,6] \to [5,5]$ 是兩個相同字元 \texttt{R}、\texttt{E}；$[5,5]$ 的值 $3$ 不等於
$D[4,4] + 1 = 4$（\texttt{N} 和 \texttt{T} 不同），也不等於 $D[4,5] + 1 = 4$，等於
$D[5,4] + 1 = 3$，所以是插入 $y_5 = \texttt{T}$，走到 $[5, 4]$；接著 \texttt{N} 相同到 $[4,3]$；
$[4,3]$ 是 \texttt{T} 換成 \texttt{I}（$D[3,2] + 1 = 2$）；\texttt{R} 相同到 $[2,1]$；$[2,1]$
是刪除 \texttt{A}（$D[1,1] + 1 = 1$）；最後 \texttt{P} 相同到 $[0,0]$。 得到的正是上面的對齊。

\textbf{和 LCS 的關係。} 如果只允許插入和刪除，相同的字元配成一行，不同的字元只能一個刪、一
個插，成本 $2$，所以最好的對齊就是讓相同的行越多越好：
\[
\text{只用插入、刪除的距離} = m + n - 2 \cdot \text{LCS 長度} .
\]
\texttt{PARTNER} 和 \texttt{PRINTER} 的 LCS 長度是 $5$（例如 \texttt{PRTER}），只用插入、刪除
要 $7 + 7 - 10 = 4$ 步，允許替換時是 $3$ 步。

實際應用常把三種操作設成不同的成本（例如鍵盤上相鄰的字母替換得便宜一點），遞迴式只要把
三個 $1$ 換成對應的成本；再加上「相鄰兩個字元對調」也算一步的版本叫 Damerau 距離，要多看
$D[i-2, j-2]$。

\begin{example}[label=ex:dp-edit]{c}{(edit\_distance.c) 編輯距離與一種對齊}
#include <stdio.h>
#include <stdlib.h>
#include <string.h>

#define MAXLEN 1000

int min3(int a, int b, int c)
{
    int m = a < b ? a : b;
    return m < c ? m : c;
}

// 算出 x 改成 y 的編輯距離，並印出一種對齊（倒著存在 ax、ay、op）。
void edit(const char *x, const char *y)
{
    static char ax[2 * MAXLEN], ay[2 * MAXLEN], op[2 * MAXLEN];
    int m = strlen(x), n = strlen(y);
    int *d = malloc((size_t)(m + 1) * (n + 1) * sizeof *d);
    if (d == NULL) {
        return;
    }
#define D(i, j) d[(i) * (n + 1) + (j)]
    for (int i = 0; i <= m; i++) {
        D(i, 0) = i;                        // 刪掉 x 的前 i 個字
    }
    for (int j = 0; j <= n; j++) {
        D(0, j) = j;                        // 插入 y 的前 j 個字
    }
    for (int i = 1; i <= m; i++) {
        for (int j = 1; j <= n; j++) {
            int diff = (x[i - 1] != y[j - 1]);
            D(i, j) = min3(D(i - 1, j - 1) + diff, D(i - 1, j) + 1, D(i, j - 1) + 1);
        }
    }
    int i = m, j = n, k = 0;                // 依序試：左上、上（刪除）、左（插入）
    while (i > 0 || j > 0) {
        if (i > 0 && j > 0 && D(i, j) == D(i - 1, j - 1) + (x[i - 1] != y[j - 1])) {
            ax[k] = x[--i];
            ay[k] = y[--j];
            op[k] = (ax[k] == ay[k]) ? '=' : 'S';
        } else if (i > 0 && D(i, j) == D(i - 1, j) + 1) {
            ax[k] = x[--i];
            ay[k] = '-';
            op[k] = 'D';
        } else {
            ax[k] = '-';
            ay[k] = y[--j];
            op[k] = 'I';
        }
        k++;
    }
    printf("%s -> %s: distance %d\n", x, y, D(m, n));
#undef D
    free(d);
    const char *rows[3] = {ax, ay, op};
    for (int r = 0; r < 3; r++) {
        for (int t = k - 1; t >= 0; t--) {
            printf("%c%c", rows[r][t], t > 0 ? ' ' : '\n');
        }
    }
}

int main(void)
{
    static char x[MAXLEN + 1], y[MAXLEN + 1];
    while (scanf("%1000s %1000s", x, y) == 2) {
        edit(x, y);
    }
    return 0;
}
\end{example}
\iosample{PARTNER PRINTER}{PARTNER -> PRINTER: distance 3 \\ P A R T N - E R \\ P - R I N T E R \\ = D = S = I = =}

這組的最佳對齊只有一種：還原時每一步都只有一個方向符合。 有好幾種時，程式依照「左上、上、
左」的順序只印出其中一種。

\subsection{最長遞增子序列：\texorpdfstring{$O(n^2)$}{O(n²)}}
\label{sec:dp-lis}

給一個數列 $A[1..n]$，\textbf{最長遞增子序列}（longest increasing subsequence，LIS）是最長的
\textbf{嚴格遞增}子序列。 本節的例子是
\[
A = \langle 4, 9, 2, 7, 5, 3, 8, 6, 10, 1, 7 \rangle ,
\]
例如 $\langle 4, 7, 8, 10 \rangle$ 是一個長度 $4$ 的遞增子序列；這個數列沒有長度 $5$ 的。

只用前綴當子問題不夠：知道「$A[1..i]$ 的 LIS 長度」，還是不知道 $A[i+1]$ 接不接得上去，因為
接不接得上要看最後一個數多大。 所以子問題改成：令 $L[i]$ 為\textbf{以 $A[i]$ 結尾}的遞增子
序列的最大長度。 以 $A[i]$ 結尾的遞增子序列，去掉 $A[i]$ 後不是空的，就是以某個 $A[j]$ 結尾
（$j < i$、$A[j] < A[i]$）的遞增子序列，而且那一段要盡量長：
\[
L[i] = 1 + \max\bigl\{\, L[j] : 1 \le j < i,\ A[j] < A[i] \,\bigr\},
\]
大括號裡沒有任何 $j$ 時，最大值當作 $0$。 答案是 $\max_i L[i]$，不是 $L[n]$。

\begin{pseudo}[label=alg:dp-lis-quad]{LIS，$O(n^2)$}
\Function{LIS-Quadratic}{$A, n$}
    \For{$i \gets 1$ \To $n$}
        \State $L[i] \gets 1$
        \State $\mathit{prev}[i] \gets \NIL$
        \For{$j \gets 1$ \To $i - 1$}
            \If{$A[j] < A[i]$ \AND $L[j] + 1 > L[i]$}
                \State $L[i] \gets L[j] + 1$
                \State $\mathit{prev}[i] \gets j$
            \EndIf
        \EndFor
    \EndFor
    \State \Return $\max_{1 \le i \le n} L[i]$
\EndFunction
\end{pseudo}

$\mathit{prev}[i]$ 記下 $A[i]$ 前面接的是誰，平手時保留最小的 $j$（第 6 行用 $>$）。 兩層迴圈，
時間 $\Theta(n^2)$。

\begin{center}
\begin{tabular}{l*{11}{c}}
\toprule
$i$ & 1 & 2 & 3 & 4 & 5 & 6 & 7 & 8 & 9 & 10 & 11 \\
\midrule
$A[i]$ & 4 & 9 & 2 & 7 & 5 & 3 & 8 & 6 & 10 & 1 & 7 \\
$L[i]$ & 1 & 2 & 1 & 2 & 2 & 2 & 3 & 3 & 4 & 1 & 4 \\
$\mathit{prev}[i]$ & -- & 1 & -- & 1 & 1 & 3 & 4 & 5 & 7 & -- & 8 \\
\bottomrule
\end{tabular}
\end{center}

例如 $L[7]$：$A[7] = 8$，前面比 $8$ 小的有 $A[1] = 4$、$A[3] = 2$、$A[4] = 7$、$A[5] = 5$、
$A[6] = 3$，它們的 $L$ 是 $1, 1, 2, 2, 2$，最大的 $2$ 最先出現在 $j = 4$，所以 $L[7] = 3$、
$\mathit{prev}[7] = 4$。 最大的 $L$ 是 $L[9] = 4$（$L[11]$ 也是 $4$，取先出現的），沿著
$\mathit{prev}$ 走回去：$9 \to 7 \to 4 \to 1$，LIS 是 $\langle 4, 7, 8, 10 \rangle$。

另一個看法：LIS 就是 $A$ 和「把 $A$ 排序、去掉重複」之後的序列的 LCS。 這也是 $O(n^2)$。

\subsection{最長遞增子序列：\texorpdfstring{$O(n \lg n)$}{O(n lg n)}}
\label{sec:dp-lis-fast}

$n$ 到 $10^6$ 時 $n^2$ 太慢。 換一個子問題：由左往右掃，對每個長度 $k$，記住
\[
\mathit{tail}[k] = \text{目前看過的數裡，所有長度 $k$ 的遞增子序列中，結尾最小的那一個的位置。}
\]
同樣長度時，結尾越小越好，因為之後越容易接上新的數，所以只記結尾最小的那一個就夠了。

\begin{whybox}[{tail 的值一定是嚴格遞增的}]
設目前有長度 $k$ 和長度 $k+1$ 的遞增子序列。 取一個長度 $k + 1$、結尾最小的，把它的最後一個
數拿掉，得到一個長度 $k$ 的遞增子序列，它的結尾比原本的結尾小。 所以長度 $k$ 的最小結尾，嚴
格小於長度 $k + 1$ 的最小結尾：
\[
A[\mathit{tail}[1]] < A[\mathit{tail}[2]] < \cdots < A[\mathit{tail}[\mathit{len}]] ,
\]
$\mathit{len}$ 是目前最長的長度。 這個陣列是排好序的，所以可以二分搜尋。

處理新的數 $x = A[i]$ 時，找最小的 $k$ 使 $A[\mathit{tail}[k]] \ge x$：
\begin{itemize}
\item $x$ 可以接在長度 $k - 1$、結尾 $A[\mathit{tail}[k-1]] < x$ 的那一條後面，得到長度 $k$、
結尾 $x$ 的遞增子序列，而且 $x \le A[\mathit{tail}[k]]$，所以 $x$ 成為新的長度 $k$ 最小結尾。
\item 長度小於 $k$ 的最小結尾都比 $x$ 小，$x$ 換不掉它們。 長度 $k$ 以上的結尾都 $\ge x$，$x$
接不上去，所以也生不出更長、以 $x$ 結尾的遞增子序列。 其他長度的最小結尾都不變。
\item 都比 $x$ 小時 $k = \mathit{len} + 1$，最長的長度加一。
\end{itemize}
每一步之後 $\mathit{tail}$ 的定義仍然成立，所以掃完時 $\mathit{len}$ 就是 LIS 的長度。
\end{whybox}

「最小的 $k$ 使 $A[\mathit{tail}[k]] \ge x$」正是第~\ref{sec:dc-bsearch}~節 \proc{Lower-Bound}
做的事，每次 $O(\lg n)$，總時間 $O(n \lg n)$。 為了還原，另外記 $\mathit{prev}[i]$：$A[i]$ 放進
第 $k$ 格時，它前面接的是當時的 $\mathit{tail}[k - 1]$。

\begin{pseudo}[label=alg:dp-lis-fast]{LIS，$O(n \lg n)$}
\Function{LIS-Fast}{$A, n$}
    \State $\mathit{len} \gets 0$
    \For{$i \gets 1$ \To $n$}
        \State 二分搜尋最小的 $k$（$1 \le k \le \mathit{len}$）使 $A[\mathit{tail}[k]] \ge A[i]$；沒有就令 $k \gets \mathit{len} + 1$
        \If{$k > 1$}
            \State $\mathit{prev}[i] \gets \mathit{tail}[k - 1]$
        \Else
            \State $\mathit{prev}[i] \gets \NIL$
        \EndIf
        \State $\mathit{tail}[k] \gets i$
        \If{$k = \mathit{len} + 1$}
            \State $\mathit{len} \gets \mathit{len} + 1$
        \EndIf
    \EndFor
    \State \Return $\mathit{len}$
\EndFunction
\end{pseudo}

同樣的數列跑一次。 表中「tail 的值」一欄寫的是 $A[\mathit{tail}[1]], A[\mathit{tail}[2]],
\ldots$，方便看；程式裡存的是位置。

\begin{center}
\begin{tabular}{cclclc}
\toprule
$i$ & $A[i]$ & $k$ & 動作 & tail 的值 & $\mathit{prev}[i]$ \\
\midrule
1 & 4 & 1 & 接在最後 & 4 & -- \\
2 & 9 & 2 & 接在最後 & 4, 9 & 1 \\
3 & 2 & 1 & 取代第 1 格 & 2, 9 & -- \\
4 & 7 & 2 & 取代第 2 格 & 2, 7 & 3 \\
5 & 5 & 2 & 取代第 2 格 & 2, 5 & 3 \\
6 & 3 & 2 & 取代第 2 格 & 2, 3 & 3 \\
7 & 8 & 3 & 接在最後 & 2, 3, 8 & 6 \\
8 & 6 & 3 & 取代第 3 格 & 2, 3, 6 & 6 \\
9 & 10 & 4 & 接在最後 & 2, 3, 6, 10 & 8 \\
10 & 1 & 1 & 取代第 1 格 & 1, 3, 6, 10 & -- \\
11 & 7 & 4 & 取代第 4 格 & 1, 3, 6, 7 & 8 \\
\bottomrule
\end{tabular}
\end{center}

最後 $\mathit{len} = 4$，$\mathit{tail}[4] = 11$。 沿 $\mathit{prev}$ 走回去是位置
$11, 8, 6, 3$，倒過來讀，LIS 是 $A[3]$、$A[6]$、$A[8]$、$A[11]$，也就是 $\langle 2, 3, 6, 7 \rangle$。
它和 $O(n^2)$ 版本找到的 $\langle 4, 7, 8, 10 \rangle$ 不同，但一樣長；遇到好幾個 LIS 時，兩種
方法選的不一樣。

要注意兩件事。 第一，最後的 tail 值 $1, 3, 6, 7$ \textbf{不是}一個子序列：$1$ 在原數列裡排在
$3$、$6$ 後面。 tail 只保證長度對，要還原必須用 $\mathit{prev}$。 第二，相等的數怎麼處理：嚴格遞增時相等的數不能接在一起，所以找的是第一個「$\ge$」的格子，
遇到和某個 tail 值相等的數時，是\textbf{取代}那一格，不會接到它後面。 如果題目要的是
\textbf{非嚴格}遞增（允許相等），要改找第一個「$>$」的格子。

\subsection{C 程式}

\texttt{tail} 存的是位置，比較時用 \texttt{a[tail[mid]]}。 二分搜尋的寫法和範例~\ref{ex:dc-bsearch}
的 \texttt{lower\_bound} 一樣，範圍是左閉右開的 $[0, \mathit{len})$。

\begin{example}[label=ex:dp-lis]{c}{(lis.c) $O(n \lg n)$ 的最長遞增子序列}
#include <stdio.h>
#include <stdlib.h>

#define MAXN 1000000

// 最長嚴格遞增子序列，O(n lg n)。out 放一個 LIS，回傳長度。
// tail[k]：長度 k+1 的遞增子序列裡結尾最小的那一個，結尾的位置。
// pred[i]：以 a[i] 結尾的那一條，前一個數的位置（-1 表示沒有）。
int lis(const int *a, int n, int *out)
{
    int *tail = malloc(n * sizeof *tail);
    int *pred = malloc(n * sizeof *pred);
    if (tail == NULL || pred == NULL) {
        free(tail);
        free(pred);
        return -1;
    }
    int len = 0;
    for (int i = 0; i < n; i++) {
        int lo = 0, hi = len;               // 找第一個 a[tail[k]] >= a[i] 的 k
        while (lo < hi) {
            int mid = lo + (hi - lo) / 2;
            if (a[tail[mid]] < a[i]) {
                lo = mid + 1;
            } else {
                hi = mid;
            }
        }
        pred[i] = (lo > 0) ? tail[lo - 1] : -1;
        tail[lo] = i;                       // 取代，或接在最後面
        if (lo == len) {
            len++;
        }
    }
    for (int k = len - 1, i = tail[len - 1]; k >= 0; k--, i = pred[i]) {
        out[k] = a[i];                      // 沿著 pred 倒著走回去
    }
    free(tail);
    free(pred);
    return len;
}

int main(void)
{
    static int a[MAXN], out[MAXN];
    int n;
    while (scanf("%d", &n) == 1 && n > 0 && n <= MAXN) {
        for (int i = 0; i < n; i++) {
            if (scanf("%d", &a[i]) != 1) {
                return 1;
            }
        }
        int len = lis(a, n, out);
        printf("length %d:", len);
        for (int k = 0; k < len; k++) {
            printf(" %d", out[k]);
        }
        printf("\n");
    }
    return 0;
}
\end{example}
\iosample{11 \\ 4 9 2 7 5 3 8 6 10 1 7 \\ 6 \\ 5 5 5 5 5 5 \\ 7 \\ 9 8 7 6 5 4 3}{length 4: 2 3 6 7 \\ length 1: 5 \\ length 1: 3}

第二組六個 $5$：嚴格遞增時相等的數不能接，LIS 長度是 $1$。 第三組完全遞減，每個數都取代第
一格，長度也是 $1$。

\section{所有點對最短路徑的 DP 觀點}
\label{sec:dp-apsp}

最後用這一章的眼光看一個圖論問題。 有向圖的點編號 $1$ 到 $n$，邊 $(i, j)$ 有權重 $w_{ij}$（可以
是負的，但圖裡沒有總權重為負的環）。 \textbf{所有點對最短路徑}問題要算出每一對 $(i, j)$ 的最
短路徑長度。 第~25~章會完整講這個問題，包括怎麼還原路徑、怎麼偵測負環，以及其他做法；這
一節只看它的遞迴式從哪裡來。

\subsection{子問題：限制中間點}

一條路徑除了起點和終點以外經過的點，叫它的\textbf{中間點}。 定義
\[
d_{ij}^{(k)} = \text{從 } i \text{ 到 } j \text{、中間點都在 } \{1, 2, \ldots, k\} \text{ 裡的最短路徑長度。}
\]
$k = 0$ 時不准有中間點，只能直接走一條邊：$d_{ij}^{(0)} = w_{ij}$，$i = j$ 時是 $0$，沒有邊時
是 $\infty$。 $k = n$ 時沒有限制，$d_{ij}^{(n)}$ 就是答案。

從 $k - 1$ 到 $k$：看中間點都在 $\{1, \ldots, k\}$ 的最短路徑 $p$。 沒有負環時可以假設 $p$
是簡單路徑，所以點 $k$ 在 $p$ 上最多出現一次。
\begin{itemize}
\item $p$ 不經過 $k$：它的中間點都在 $\{1, \ldots, k-1\}$，長度是 $d_{ij}^{(k-1)}$。
\item $p$ 經過 $k$：在 $k$ 切成 $i \leadsto k$ 和 $k \leadsto j$ 兩段，兩段的中間點都在
$\{1, \ldots, k - 1\}$（$k$ 只出現一次，現在是端點），而且各自都要是最短的（剪下貼上），長度是
$d_{ik}^{(k-1)} + d_{kj}^{(k-1)}$。
\end{itemize}
所以
\[
d_{ij}^{(k)} = \min\bigl( d_{ij}^{(k-1)},\ d_{ik}^{(k-1)} + d_{kj}^{(k-1)} \bigr) .
\]
這和 0/1 背包的「第 $k$ 件拿或不拿」是同一種想法：第 $k$ 個點「用或不用」。 和第~\ref{sec:dp-substructure}~節最
長簡單路徑的反例對照：最短路徑的兩段即使經過同一個點，也能把繞的那一圈拿掉而不變長，所以
子問題是獨立的；這正是「沒有負環」這個條件的作用。

子問題有 $n^3$ 個，每個 $O(1)$，時間 $\Theta(n^3)$。 第 $k$ 層只用到第 $k - 1$ 層，而且第 $k$
層的第 $k$ 列、第 $k$ 欄和第 $k - 1$ 層相同：以第 $k$ 欄為例，$d_{kk}^{(k-1)} = 0$，所以
$d_{ik}^{(k)} = \min\bigl(d_{ik}^{(k-1)},\ d_{ik}^{(k-1)} + 0\bigr) = d_{ik}^{(k-1)}$，第 $k$ 列同理。
第 $k$ 層要讀的只有第 $k$ 列、第 $k$ 欄，它們在這一層不會變，所以可以只用一個矩陣就地更新，
空間 $\Theta(n^2)$。 這就是 \textbf{Floyd--Warshall 演算法}：

\begin{pseudo}[label=alg:dp-fw]{Floyd--Warshall（就地更新）}
\Function{Floyd-Warshall}{$W, n$}
    \State $D \gets W$ \Comment{複製一份 $n \times n$ 矩陣}
    \For{$k \gets 1$ \To $n$}
        \For{$i \gets 1$ \To $n$}
            \For{$j \gets 1$ \To $n$}
                \If{$D[i, k] + D[k, j] < D[i, j]$}
                    \State $D[i, j] \gets D[i, k] + D[k, j]$
                \EndIf
            \EndFor
        \EndFor
    \EndFor
    \State \Return $D$
\EndFunction
\end{pseudo}

最外層一定要是 $k$。 把 $k$ 放到最內層，算 $D[i, j]$ 時用到的 $D[i, k]$、$D[k, j]$ 還沒有考慮
過足夠多的中間點，答案會錯，這是寫這個演算法最常見的錯誤。

\subsection{一個例子}

四個點，邊是 $1 \to 2$（$4$）、$1 \to 3$（$9$）、$2 \to 3$（$3$）、$2 \to 4$（$8$）、
$3 \to 1$（$6$）、$3 \to 4$（$-2$）、$4 \to 1$（$5$）。 依序算出 $D^{(0)}$ 到 $D^{(4)}$，粗體是
這一層變小的格子：
\begin{gather*}
D^{(0)} = \begin{pmatrix} 0 & 4 & 9 & \infty \\ \infty & 0 & 3 & 8 \\ 6 & \infty & 0 & -2 \\ 5 & \infty & \infty & 0 \end{pmatrix},\quad D^{(1)} = \begin{pmatrix} 0 & 4 & 9 & \infty \\ \infty & 0 & 3 & 8 \\ 6 & \mathbf{10} & 0 & -2 \\ 5 & \mathbf{9} & \mathbf{14} & 0 \end{pmatrix},\quad D^{(2)} = \begin{pmatrix} 0 & 4 & \mathbf{7} & \mathbf{12} \\ \infty & 0 & 3 & 8 \\ 6 & 10 & 0 & -2 \\ 5 & 9 & \mathbf{12} & 0 \end{pmatrix}, \\
D^{(3)} = \begin{pmatrix} 0 & 4 & 7 & \mathbf{5} \\ \mathbf{9} & 0 & 3 & \mathbf{1} \\ 6 & 10 & 0 & -2 \\ 5 & 9 & 12 & 0 \end{pmatrix},\quad D^{(4)} = \begin{pmatrix} 0 & 4 & 7 & 5 \\ \mathbf{6} & 0 & 3 & 1 \\ \mathbf{3} & \mathbf{7} & 0 & -2 \\ 5 & 9 & 12 & 0 \end{pmatrix}
\end{gather*}

例如 $D^{(3)}$ 的 $d_{14}$ 從 $12$（$1 \to 2 \to 4$）變成 $5$：$d_{13}^{(2)} + d_{34}^{(2)} =
7 + (-2)$，也就是 $1 \to 2 \to 3 \to 4$。 $D^{(4)}$ 的 $d_{31} = 3$ 是 $3 \to 4 \to 1$，比直接
的邊 $6$ 短。 每一層都有格子變小，最後一層就是答案。

子問題也可以改成「最多用 $k$ 條邊」（這裡的 $k$ 是邊數，不是中間點的編號）：
$\ell_{ij}^{(k)} = \min_{t} \bigl(\ell_{it}^{(k-1)} + w_{tj}\bigr)$，看最後一條邊從哪個點 $t$
走進 $j$。 要算到 $k = n - 1$，每層 $\Theta(n^3)$，共 $\Theta(n^4)$；它的形狀和矩陣乘法一樣，
用第~\ref{sec:dc-power}~節快速冪的反覆平方可以降到 $\Theta(n^3 \lg n)$。 同一個問題，
子問題選得不同，時間可以差很多，第~25~章會比較這幾種做法。

\subsection{本章的 DP 整理}

\begin{center}
\small
\begin{tabular}{>{\raggedright\arraybackslash}p{0.2\linewidth}>{\raggedright\arraybackslash}p{0.22\linewidth}cc>{\raggedright\arraybackslash}p{0.22\linewidth}}
\toprule
問題 & 子問題 & 子問題數 & 每個的選擇 & 時間 \\
\midrule
二項式係數 & $\binom{i}{j}$ & $\Theta(nk)$ & $1$ & $\Theta(nk)$ \\
矩陣鏈乘 & 區間 $A_{i..j}$ & $\Theta(n^2)$ & $O(n)$ & $\Theta(n^3)$ \\
LCS & 兩個前綴 $X_i, Y_j$ & $\Theta(mn)$ & $O(1)$ & $\Theta(mn)$ \\
最佳二元搜尋樹 & 區間 $k_i..k_j$ & $\Theta(n^2)$ & $O(n)$ & $\Theta(n^3)$，Knuth 改進後 $\Theta(n^2)$ \\
0/1 背包 & 前 $i$ 件、容量 $c$ & $\Theta(nW)$ & $2$ & $\Theta(nW)$，擬多項式 \\
編輯距離 & 兩個前綴 $X_i, Y_j$ & $\Theta(mn)$ & $3$ & $\Theta(mn)$ \\
LIS & 以 $A[i]$ 結尾 & $n$ & $O(n)$ & $\Theta(n^2)$；改寫後 $O(n \lg n)$ \\
Floyd--Warshall & 中間點限 $\{1..k\}$ & $\Theta(n^3)$ & $2$ & $\Theta(n^3)$ \\
\bottomrule
\end{tabular}
\end{center}

常見的子問題只有幾種形狀：\textbf{前綴}（LCS、編輯距離、背包的前 $i$ 件）、\textbf{區間}（矩陣
鏈乘、最佳二元搜尋樹）、\textbf{以某個元素結尾}（LIS），以及多加一個\textbf{數值參數}（背包的容
量、Floyd--Warshall 的 $k$）。 遇到新問題時，先試這幾種，再看最佳解的「最後一步」有幾種選
擇。 只要值時，前綴型和背包型的表通常都可以只留一兩列；區間型的表依賴的格子散在各列，不容
易這樣省。

\section{本章常見錯誤}
\label{sec:dp-pitfalls}

\begin{longtable}{p{0.25\linewidth}p{0.48\linewidth}>{\raggedright\arraybackslash}p{0.17\linewidth}}
\toprule
錯誤 & 例子 & 見 \\
\midrule
\endhead
沒有定義表格的一格 & 只寫遞迴式，沒寫 $c[i, j]$ 代表什麼，或邊界漏寫 &
第~\ref{sec:dp-steps}~節 \\
記憶化的記號和答案撞在一起 & 用 $0$ 表示「還沒算過」，但子問題的答案本來就可能是 $0$，於
是一直重算 & 第~\ref{sec:dp-memo}~節 \\
填表順序錯 & 矩陣鏈乘一列一列由上往下填，算 $m[1, 3]$ 時 $m[2, 3]$ 還沒算 &
第~\ref{sec:dp-matrix}~節 \\
把一個候選當成最佳 & 矩陣鏈乘一律在最小的維度切開；最佳二元搜尋樹一律讓最常查的鍵當根 &
第~\ref{sec:dp-matrix}、\ref{sec:dp-obst}~節 \\
還原 LCS 的規則不對 & 沒先檢查 $x_i = y_j$ 就比較上、左兩格，漏掉該選的字元；或沒照題目規定的平手方向，還原出另一個 LCS &
第~\ref{sec:dp-lcs}~節 \\
最佳二元搜尋樹的成本模型混用 & 外部節點用深度加一算成本，空樹卻寫 $e[i, i-1] = 0$；
$w(i, j)$ 的 $q$ 從 $q_i$ 開始加 & 第~\ref{sec:dp-obst}~節 \\
一維背包的方向 & 0/1 背包的容量由小到大跑，同一件物品被拿很多次 &
第~\ref{sec:dp-knapsack-1d}~節 \\
說背包 DP 是多項式時間 & $\Theta(nW)$ 是 $W$ 的數值的多項式，不是位元數的 &
第~\ref{sec:dp-pseudo-poly}~節 \\
編輯距離的邊界 & 把 $D[i, 0]$、$D[0, j]$ 都設成 $0$ & 第~\ref{sec:dp-edit}~節 \\
LIS 的答案 & 答案寫成 $L[n]$，不是 $\max_i L[i]$；把最後的 tail 陣列當成一個 LIS &
第~\ref{sec:dp-lis}、\ref{sec:dp-lis-fast}~節 \\
嚴格和非嚴格遞增 & 嚴格遞增卻找第一個 $>$ 的格子，相等的數被接在一起 &
第~\ref{sec:dp-lis-fast}~節 \\
Floyd--Warshall 的迴圈順序 & 把 $k$ 放在最內層 & 第~\ref{sec:dp-apsp}~節 \\
\bottomrule
\end{longtable}

\section{觀念題}
\label{sec:dp-exercises}

題型和研究所考試、期中期末考的筆試相近，有選擇、計算、追蹤、設計和證明。 平手時一律用本章
的規則。 解答在本章最後。

\begin{enumerate}
\item \label{dp-q-tf}（選擇，可能不只一個答案）下列哪些敘述正確？
\begin{enumerate}
\item 同一個遞迴式，記憶化和由下而上填表的漸近時間一定相同。
\item 0/1 背包的 $\Theta(nW)$ 演算法是多項式時間演算法。
\item 矩陣鏈乘只要求 $m[1, n]$ 時，可以像 LCS 一樣只留兩列。
\item Floyd--Warshall 的三層迴圈裡，$k$ 必須在最外層。
\item 一個最佳化問題只要有最佳子結構，就能用動態規劃在多項式時間內解出來。
\item 把 $X$、$Y$ 兩個字串對調，LCS 的長度不變，還原出來的 LCS 也不變。
\end{enumerate}

\item \label{dp-q-mc}（計算）四個矩陣的大小是 $\langle 5, 8, 2, 10, 4 \rangle$。 依鏈長填出
$m$ 表和 $s$ 表（寫出每一格每個 $k$ 的成本），寫出最佳的括號方式，並和由左往右乘
$((A_1 A_2) A_3) A_4$ 的次數比較。

\item \label{dp-q-mc-rec}（證明）不用記憶化，直接照 $m[i, j]$ 的遞迴式寫遞迴函式算
$m[1, n]$。 令 $R(n)$ 為鏈長 $n$ 時函式被呼叫的總次數（$R(1) = 1$）。 寫出 $R(n)$ 的遞迴
式，證明 $R(n) = 3^{n-1}$，並和記憶化時不同子問題的個數比較。

\item \label{dp-q-lcs}（追蹤）$X = \texttt{TEACHER}$（列）、$Y = \texttt{CHEATER}$（欄）。
\begin{enumerate}
\item 用 \proc{LCS-Length} 填出 $c$ 表和箭頭，用 \proc{Print-LCS} 還原一個 LCS。
\item 列出所有的 LCS。
\item 平手時改成指向左，還原出來的是哪一個？
\end{enumerate}

\item \label{dp-q-obst}（計算）三個鍵 $k_1 < k_2 < k_3$，查詢次數
$p = \langle 4, 2, 6 \rangle$，查不到的次數 $q = \langle 1, 3, 2, 2 \rangle$（$q_0$ 到 $q_3$）。
用本章的成本模型填出 $w$、$e$、$\mathit{root}$ 表，畫出最佳二元搜尋樹（含外部節點），算出
平均每次查詢的比較次數。 如果改用「走到外部節點也算一次」的成本模型，$e[1, 3]$ 是多少？

\item \label{dp-q-ks}（追蹤）背包容量 $W = 8$，四件物品的 $(w_i, v_i)$ 是 $(2, 5)$、
$(5, 9)$、$(4, 10)$、$(3, 10)$。
\begin{enumerate}
\item 填出 $K$ 表，用 \proc{Knapsack-Items} 找出要拿的物品。
\item 用一維陣列但容量由小到大跑，$d[0..8]$ 最後是什麼？ 這是哪個問題的答案？
\end{enumerate}

\item \label{dp-q-sparse}（計算）$n = 20$ 件物品，容量 $W = 10^9$，每件的重量都在
$10^7$ 到 $10^8$ 之間。
\begin{enumerate}
\item \proc{Knapsack} 的表有幾格？ 開得出來嗎？
\item 改用記憶化：$\proc{Best}(i, c)$ 回傳「前 $i$ 件、容量 $c$」的最大價值，從
$\proc{Best}(20, W)$ 開始遞迴。 證明會被呼叫到的不同 $(i, c)$ 最多 $2^{21} - 1$ 個。 這時表
要用什麼資料結構存？
\end{enumerate}

\item \label{dp-q-edit}（追蹤）填出 \texttt{SPORT}（列）改成 \texttt{STRONG}（欄）的 $D$ 表，
寫出編輯距離和一種最少步數的改法（用上下兩列的對齊表示）。

\item \label{dp-q-edit-match}（證明）證明 $x_i = y_j$ 時 $D[i, j] = D[i-1, j-1]$。 提示：先
證明 $D[i-1, j-1] \le D[i-1, j] + 1$ 和 $D[i-1, j-1] \le D[i, j-1] + 1$。

\item \label{dp-q-lis}（追蹤）$A = \langle 3, 8, 4, 12, 6, 10, 7, 11, 5 \rangle$。
\begin{enumerate}
\item 用 \proc{LIS-Quadratic} 算出 $L$ 和 $\mathit{prev}$，還原一個 LIS。
\item 用 \proc{LIS-Fast} 寫出每一步之後的 tail 值和 $\mathit{prev}$，還原一個 LIS。
\item 最後的 tail 值是不是 $A$ 的子序列？
\end{enumerate}

\item \label{dp-q-alloc}（設計）期末考前還有 $T = 4$ 小時，要分給三科，每科只能分整數小
時。 第 $i$ 科讀 $h$ 小時預計拿到的分數 $f_i(h)$ 如下表。 定義子問題、寫出遞迴式，填表求出
最高總分和每科分到的時數，並寫出一般 $n$ 科、$T$ 小時的時間複雜度。
\begin{center}
\begin{tabular}{c*{5}{c}}
\toprule
$h$ & 0 & 1 & 2 & 3 & 4 \\
\midrule
$f_1(h)$ & 0 & 5 & 8 & 10 & 11 \\
$f_2(h)$ & 0 & 3 & 7 & 10 & 12 \\
$f_3(h)$ & 0 & 4 & 6 & 7 & 8 \\
\bottomrule
\end{tabular}
\end{center}

\item \label{dp-q-fw}（追蹤）三個點，邊是 $1 \to 2$（$5$）、$1 \to 3$（$4$）、$2 \to 1$
（$2$）、$2 \to 3$（$-3$）、$3 \to 1$（$4$）。 寫出 $D^{(0)}$、$D^{(1)}$、$D^{(2)}$、$D^{(3)}$。

\item \label{dp-q-sub}（選擇並說明）下列哪些問題有最佳子結構？
\begin{enumerate}
\item 有向無環圖上，從 $s$ 到 $t$ 邊數最多的路徑。
\item 一般有向圖上，從 $s$ 到 $t$ 邊數最多的簡單路徑。
\item 帶權有向圖上，從 $s$ 到 $t$、恰好走 $k$ 條邊（點可以重複）的最小權重走法。
\end{enumerate}
\end{enumerate}

\section{實作練習}
\label{sec:dp-practice}

下面是 Aki-Judge 上和本章相關的題目。 矩陣鏈乘和最佳二元搜尋樹目前沒有對應的練習題，請用觀
念題第~\ref{dp-q-mc}、\ref{dp-q-obst}~題練習。

\begin{center}
\begin{tabular}{llc}
\toprule
小節 & 題目 & 難度 \\
\midrule
DP 的要素 & \akiref{cheapest-route-across-a-grid}{網格上最便宜的路線} & 暖身 \\
DP 的要素 & \akiref{cheapest-grid-route-with-blocked-cells}{含封鎖格的最便宜路線} & 基本 \\
DP 的要素 & \akiref{mandatory-checkpoints-on-a-down-right-route}{送貨機器人的必經格} & 進階 \\
DP 的要素 & \akiref{heaviest-set-of-non-adjacent-tree-nodes}{樹上互不相鄰節點的最大權重和} & 進階 \\
LCS & \akiref{length-of-the-longest-common-subsequence}{最長共同子序列長度} & 基本 \\
LCS & \akiref{common-subsequence-without-repeated-neighbors}{相鄰不重複的共同子序列} & 進階 \\
0/1 背包 & \akiref{best-value-within-a-weight-limit}{重量上限內的最大價值} & 基本 \\
0/1 背包 & \akiref{best-value-with-unlimited-copies}{可重複取用的最大價值} & 基本 \\
編輯距離 & \akiref{minimum-edits-between-two-strings}{兩字串之間的最少編輯次數} & 基本 \\
LIS & \akiref{longest-strictly-increasing-subsequence-small-n}{最長嚴格遞增子序列（小範圍）} & 基本 \\
LIS & \akiref{longest-strictly-increasing-subsequence-large-n}{最長嚴格遞增子序列（大範圍）} & 進階 \\
\bottomrule
\end{tabular}
\end{center}

幾點提示：
\begin{itemize}
\item 兩題網格路線：每一步只能往右或往下，所以一格只會從上面或左邊走進來，照列、照欄的順
序填就好。 封鎖格當成成本 $\infty$，走不到終點時輸出題目指定的值。 「必經格」那一題要兩張
表：從起點走不走得到這一格、從這一格走不走得到終點，兩者都成立的格子才在某條路線上。
\item 樹上的題目，子問題是「以 $v$ 為根的子樹，$v$ 選或不選」兩種狀態。 $n$ 到 $2 \times 10^5$，
樹可能是一條長鏈，遞迴太深，請先用 BFS 排出順序，再倒著由葉往根算。
\item 「相鄰不重複的共同子序列」是 LCS 的變形：除了兩個位置，子問題還要記住上一個選進來的
字母（$26$ 種加上「還沒選」），用記憶化最順手。
\item 0/1 背包那題 $n \le 1000$、$W \le 10^5$，二維表太大，要用第~\ref{sec:dp-knapsack-1d}~節的一維
陣列、容量由大到小；可重複取用那題正好相反。
\item 兩題 LCS、編輯距離的字串長度到 $2000$，整張表約四百萬格，可以開；只留兩列更省。
\item LIS 大範圍那題 $n$ 到 $10^6$，要用 $O(n \lg n)$ 的做法，注意是嚴格遞增。
\end{itemize}

\section{讀完你應該會}
\label{sec:dp-checklist}

\begin{itemize}
\item 判斷一個問題有沒有重疊的子問題和最佳子結構，用剪下貼上證明最佳子結構，並舉出最長簡
單路徑這種沒有最佳子結構的例子。
\item 寫動態規劃的答案時，用一句話定義表格的一格，寫出遞迴式、邊界、填表順序、答案位置和
複雜度。
\item 對矩陣鏈乘和最佳二元搜尋樹依區間長度填表，由 $s$ 或 $\mathit{root}$ 表還原括號或樹，
並說清楚最佳二元搜尋樹用的成本模型。
\item 填 LCS 和編輯距離的表，照指定的平手規則還原一個解；只要值時改成只留兩列。
\item 寫出 0/1 背包的二維表和一維版本，說明一維時容量為什麼由大到小，以及 $\Theta(nW)$ 為
什麼是擬多項式時間。
\item 寫出 LIS 的 $O(n^2)$ 和 $O(n \lg n)$ 做法並還原一個 LIS；寫出 Floyd--Warshall 的遞迴式，
說明 $k$ 為什麼在最外層。
\end{itemize}

\textbf{延伸閱讀}（書末會整理完整的參考資料）

\begin{itemize}
\item Thomas H. Cormen 等，\textit{Introduction to Algorithms}，第 4 版，第 14 章：動態規劃的
四個經典例子和「要素」一節的完整討論；所有點對最短路徑在第 23 章。 它的最佳二元搜尋樹把走
到外部節點也算一次比較，和本章差一個常數，對照時要注意。
\item Steven S. Skiena，\textit{The Algorithm Design Manual}，第 3 版，第 10 章：用很多實際例
子講怎麼「想出」子問題，包括編輯距離的各種變形。
\item Antti Laaksonen，\textit{Competitive Programmer's Handbook}，第 7 章：篇幅短，有 LIS、
網格路徑、背包和編輯距離的程式寫法。
\item D. E. Knuth，``Optimum Binary Search Trees''，\textit{Acta Informatica} 1(1)，1971：最佳
二元搜尋樹的 $\Theta(n^2)$ 方法，也就是本章提到的 $\mathit{root}$ 單調性。
\item D. S. Hirschberg，``A Linear Space Algorithm for Computing Maximal Common
Subsequences''，\textit{Communications of the ACM} 18(6)，1975：只用線性空間就能還原 LCS 的
分治方法。
\end{itemize}

\addcontentsline{toc}{section}{觀念題解答}
\section*{觀念題解答}

\begin{enumerate}
\item (d) 正確。
(a) 錯：記憶化只算真的用得到的子問題，可能比整張表少很多（第~\ref{dp-q-sparse}~題）；反過
來說，兩者的上界相同。 (b) 錯：是擬多項式。 (c) 錯：$m[i, j]$ 要用到第 $i$ 列左邊的所有
格子和第 $j$ 欄下面的所有格子，散在很多列，只留兩列不夠。 (d) 對：第 $k$ 輪要用到的
$D[i, k]$、$D[k, j]$ 必須已經是「中間點在 $\{1..k-1\}$」的最短值。 (e) 錯：最佳子結構只保證
遞迴式寫得出來，子問題可能有指數多個，例如旅行推銷員問題（第~20~章：每個城市走一次再回到起點的最短路
線）用動態規劃時，子問題要記住「走過的點的集合」，有 $2^n$ 種。
(f) 錯：長度不變，但平手規則的方向跟著對調，還原出來的可能不同，範例~\ref{ex:dp-lcs} 的前兩
組就是例子。

\item 鏈長 $2$：$m[1,2] = 5 \cdot 8 \cdot 2 = 80$、$m[2,3] = 8 \cdot 2 \cdot 10 = 160$、
$m[3,4] = 2 \cdot 10 \cdot 4 = 80$。
鏈長 $3$：$m[1,3]$：$k = 1$ 得 $0 + 160 + 400 = 560$，$k = 2$ 得 $80 + 0 + 100 = 180$，所以
$180$、$s = 2$。 $m[2,4]$：$k = 2$ 得 $0 + 80 + 64 = 144$，$k = 3$ 得 $160 + 0 + 320 = 480$，所
以 $144$、$s = 2$。
鏈長 $4$：$m[1,4]$：$k = 1$ 得 $0 + 144 + 160 = 304$，$k = 2$ 得 $80 + 80 + 40 = 200$，
$k = 3$ 得 $180 + 0 + 200 = 380$，所以 $200$、$s = 2$。
最佳是 $((A_1 A_2)(A_3 A_4))$，$200$ 次。 由左往右乘正是 $k = 3$ 那個候選，$380$ 次。

\item 每次呼叫自己算一次，再對每個 $k$ 呼叫左右兩邊：
\[
R(n) = 1 + \sum_{k=1}^{n-1} \bigl(R(k) + R(n - k)\bigr) = 1 + 2\sum_{k=1}^{n-1} R(k) .
\]
$n \ge 2$ 時把 $R(n)$ 和 $R(n-1)$ 的式子相減（$R(1) = 1$ 也符合 $R(1) = 1 + 2 \cdot 0$），得
$R(n) - R(n-1) = 2R(n-1)$，即 $R(n) = 3R(n-1)$，所以 $R(n) = 3^{n-1}$。 驗算：
$R(2) = 1 + 2 = 3$，$R(3) = 1 + 2(1 + 3) = 9$。 不同的子問題只有 $n(n+1)/2$ 個，指數次的呼叫
全部是重算。

\item (a) 表如下（列是 $X$，欄是 $Y$）：
{\footnotesize
\begin{center}
\begin{tabular}{cc*{8}{c}}
& &  & \texttt{C} & \texttt{H} & \texttt{E} & \texttt{A} & \texttt{T} & \texttt{E} & \texttt{R} \\
\midrule
0 &  & 0 & 0 & 0 & 0 & 0 & 0 & 0 & 0 \\
1 & \texttt{T} & 0 & $\uparrow$0 & $\uparrow$0 & $\uparrow$0 & $\uparrow$0 & $\nwarrow$1 & $\leftarrow$1 & $\leftarrow$1 \\
2 & \texttt{E} & 0 & $\uparrow$0 & $\uparrow$0 & $\nwarrow$1 & $\leftarrow$1 & $\uparrow$1 & $\nwarrow$2 & $\leftarrow$2 \\
3 & \texttt{A} & 0 & $\uparrow$0 & $\uparrow$0 & $\uparrow$1 & $\nwarrow$2 & $\leftarrow$2 & $\uparrow$2 & $\uparrow$2 \\
4 & \texttt{C} & 0 & $\nwarrow$1 & $\leftarrow$1 & $\uparrow$1 & $\uparrow$2 & $\uparrow$2 & $\uparrow$2 & $\uparrow$2 \\
5 & \texttt{H} & 0 & $\uparrow$1 & $\nwarrow$2 & $\leftarrow$2 & $\uparrow$2 & $\uparrow$2 & $\uparrow$2 & $\uparrow$2 \\
6 & \texttt{E} & 0 & $\uparrow$1 & $\uparrow$2 & $\nwarrow$3 & $\leftarrow$3 & $\leftarrow$3 & $\nwarrow$3 & $\leftarrow$3 \\
7 & \texttt{R} & 0 & $\uparrow$1 & $\uparrow$2 & $\uparrow$3 & $\uparrow$3 & $\uparrow$3 & $\uparrow$3 & $\nwarrow$4 \\
\end{tabular}
\end{center}
}
從 $[7, 7]$：$\nwarrow$ 選 \texttt{R} 到 $[6, 6]$，$\nwarrow$ 選 \texttt{E} 到 $[5, 5]$，
$\uparrow$ 到 $[4, 5]$，$\uparrow$ 到 $[3, 5]$，$\leftarrow$ 到 $[3, 4]$，$\nwarrow$ 選
\texttt{A} 到 $[2, 3]$，$\nwarrow$ 選 \texttt{E} 到 $[1, 2]$，$\uparrow$ 到 $[0, 2]$ 停。 LCS 是
\texttt{EAER}，長度 $4$。
(b) \texttt{CHER} 和 \texttt{EAER} 兩個。
(c) \texttt{CHER}。

\item $w$、$e$、$\mathit{root}$ 依區間長度：
{\footnotesize
\begin{center}
\begin{tabular}{c c p{0.55\linewidth} c c}
\toprule
$(i, j)$ & $w$ & 每個 $r$ 的成本 & $e$ & $\mathit{root}$ \\
\midrule
$(1,1)$ & 8 & $r=1$：$0+0+8=8$ & 8 & 1 \\
$(2,2)$ & 7 & $r=2$：$0+0+7=7$ & 7 & 2 \\
$(3,3)$ & 10 & $r=3$：$0+0+10=10$ & 10 & 3 \\
$(1,2)$ & 12 & $r=1$：$0+7+12=19$；$r=2$：$8+0+12=20$ & 19 & 1 \\
$(2,3)$ & 15 & $r=2$：$0+10+15=25$；$r=3$：$7+0+15=22$ & 22 & 3 \\
$(1,3)$ & 20 & $r=1$：$0+22+20=42$；$r=2$：$8+10+20=38$；$r=3$：$19+0+20=39$ & 38 & 2 \\
\bottomrule
\end{tabular}
\end{center}
}
空樹 $e[i, i-1] = 0$，$w[1,0] = 1$、$w[2,1] = 3$、$w[3,2] = 2$、$w[4,3] = 2$。 最佳樹的根是
$k_2$，左子節點 $k_1$、右子節點 $k_3$，四個外部節點 $d_0, d_1$ 掛在 $k_1$ 下、$d_2, d_3$ 掛在
$k_3$ 下，深度都是 $2$。 成本 $38$，平均 $38 / 20 = 1.9$ 次。 讓次數最多的 $k_3$ 當根的
樹成本是 $39$。 改用「走到外部節點也算一次」時多了 $q_0 + q_1 + q_2 + q_3 = 8$，
$e[1, 3] = 46$，最佳樹不變。

\item (a) $K$ 表（欄是容量 $0$ 到 $8$）：
{\footnotesize
\begin{center}
\begin{tabular}{l*{9}{c}}
$i=0$ & 0 & 0 & 0 & 0 & 0 & 0 & 0 & 0 & 0 \\
$i=1$ & 0 & 0 & 5 & 5 & 5 & 5 & 5 & 5 & 5 \\
$i=2$ & 0 & 0 & 5 & 5 & 5 & 9 & 9 & 14 & 14 \\
$i=3$ & 0 & 0 & 5 & 5 & 10 & 10 & 15 & 15 & 15 \\
$i=4$ & 0 & 0 & 5 & 10 & 10 & 15 & 15 & 20 & 20 \\
\end{tabular}
\end{center}
}
還原：$K[4, 8] = 20 \ne K[3, 8] = 15$，拿第 $4$ 件，剩 $5$；$K[3, 5] = 10 \ne K[2, 5] = 9$，拿
第 $3$ 件，剩 $1$；$K[2, 1] = K[1, 1] = K[0, 1] = 0$，第 $2$、$1$ 件不拿。 拿第 $3$、$4$ 件，重
$7$、價值 $20$。
(b) $d = \langle 0, 0, 5, 10, 10, 15, 20, 20, 25 \rangle$，$d[8] = 25$，例如第 $4$ 件拿兩次再
加第 $1$ 件。 這是無限背包的答案，不是 0/1 背包的。

\item (a) $21 \times (10^9 + 1)$，大約 $2 \times 10^{10}$ 格，就算每格只用 $4$ 個位元組也要
$80$ GB，開不出來。
(b) 從 $(20, W)$ 往下，每一層決定一件物品拿或不拿，所以走到第 $i$ 層時，$c$ 一定是 $W$ 減掉
第 $i+1..20$ 件中某些物品的重量和，最多 $2^{20 - i}$ 種。 全部加起來最多
$\sum_{i=0}^{20} 2^{20-i} = 2^{21} - 1 \approx 2.1 \times 10^6$ 個子問題。 $c$ 的值很大又很分
散，不能拿來當陣列索引，表要用第~12~章的雜湊表，以 $(i, c)$ 為鍵。

\item 表如下，$D[5, 6] = 4$：
{\footnotesize
\begin{center}
\begin{tabular}{cc*{7}{c}}
& &  & \texttt{S} & \texttt{T} & \texttt{R} & \texttt{O} & \texttt{N} & \texttt{G} \\
\midrule
0 &  & 0 & 1 & 2 & 3 & 4 & 5 & 6 \\
1 & \texttt{S} & 1 & 0 & 1 & 2 & 3 & 4 & 5 \\
2 & \texttt{P} & 2 & 1 & 1 & 2 & 3 & 4 & 5 \\
3 & \texttt{O} & 3 & 2 & 2 & 2 & 2 & 3 & 4 \\
4 & \texttt{R} & 4 & 3 & 3 & 2 & 3 & 3 & 4 \\
5 & \texttt{T} & 5 & 4 & 3 & 3 & 3 & 4 & 4 \\
\end{tabular}
\end{center}
}
依「左上、上、左」的順序還原：$[5,6]$ 是 \texttt{T} 換成 \texttt{G}，$[4,5]$ 是 \texttt{R} 換成
\texttt{N}，$[3,4]$ 是 \texttt{O} 相同，$[2,3]$ 是 \texttt{P} 換成 \texttt{R}，$[1,2]$ 是插入
\texttt{T}，$[1,1]$ 是 \texttt{S} 相同。 對齊是
\begin{center}
\texttt{S - P O R T}\\
\texttt{S T R O N G}
\end{center}
一次插入、三次替換，共 $4$ 步。

\item 把 $X_{i-1}$ 改成 $Y_j$ 的任何一種改法，最後再刪掉結尾的 $y_j$，就把 $X_{i-1}$ 改成了
$Y_{j-1}$，所以 $D[i-1, j-1] \le D[i-1, j] + 1$。 同理，先在 $X_{i-1}$ 後面插入 $x_i$ 得到
$X_i$，再用最好的方法把 $X_i$ 改成 $Y_{j-1}$，所以 $D[i-1, j-1] \le 1 + D[i, j-1]$。 $x_i = y_j$
時 $\delta = 0$，遞迴式的第一項 $D[i-1, j-1]$ 不大於另外兩項，最小值就是它。 所以這時可以直接
取左上角，和 LCS 一樣。

\item (a) $L = \langle 1, 2, 2, 3, 3, 4, 4, 5, 3 \rangle$，$\mathit{prev} = \langle \text{--}, 1,
1, 2, 3, 5, 5, 6, 3 \rangle$。 最大的是 $L[8] = 5$，走回去 $8 \to 6 \to 5 \to 3 \to 1$，LIS 是
$\langle 3, 4, 6, 10, 11 \rangle$。
(b) 每一步之後的 tail 值（括號裡是 $\mathit{prev}[i]$）：
$i = 1$：$3$（--）；$i = 2$：$3, 8$（$1$）；$i = 3$：$3, 4$（$1$）；$i = 4$：$3, 4, 12$（$3$）；
$i = 5$：$3, 4, 6$（$3$）；$i = 6$：$3, 4, 6, 10$（$5$）；$i = 7$：$3, 4, 6, 7$（$5$）；
$i = 8$：$3, 4, 6, 7, 11$（$7$）；$i = 9$：$3, 4, 5, 7, 11$（$3$）。
$\mathit{len} = 5$，$\mathit{tail}[5] = 8$，走回去 $8 \to 7 \to 5 \to 3 \to 1$，LIS 是
$\langle 3, 4, 6, 7, 11 \rangle$，和 (a) 不同，一樣長。
(c) 最後的 tail 值是 $3, 4, 5, 7, 11$，其中 $5$ 是 $A[9]$，排在 $7$、$11$ 後面，不是子序列。

\item 令 $G[i, t]$ 為前 $i$ 科共分 $t$ 小時的最高總分。 $G[0, t] = 0$；
$G[i, t] = \max_{0 \le h \le t} \bigl(G[i-1, t-h] + f_i(h)\bigr)$，第 $i$ 科分 $h$ 小時。 答案
$G[3, 4]$。
\begin{center}
\begin{tabular}{c*{5}{c}}
\toprule
$t$ & 0 & 1 & 2 & 3 & 4 \\
\midrule
$G[1, t]$ & 0 & 5 & 8 & 10 & 11 \\
$G[2, t]$ & 0 & 5 & 8 & 12 & 15 \\
$G[3, t]$ & 0 & 5 & 9 & 12 & 16 \\
\bottomrule
\end{tabular}
\end{center}
例如 $G[2, 3] = \max(10 + 0, 8 + 3, 5 + 7, 0 + 10) = 12$（第 $2$ 科分 $2$ 小時）。 $G[3, 4] =
\max(15 + 0, 12 + 4, 8 + 6, 5 + 7, 0 + 8) = 16$，第 $3$ 科分 $1$ 小時；回到 $G[2, 3]$，第 $2$
科分 $2$ 小時；剩 $G[1, 1]$，第 $1$ 科分 $1$ 小時。 最高 $16$ 分（$5 + 7 + 4$）。 一般情況有
$n(T+1)$ 格，每格最多試 $T + 1$ 個 $h$，時間 $\Theta(nT^2)$。 這也是擬多項式（$T$ 是數值）。

\item 四個矩陣如下，粗體是這一層變小的格子。
{\small
\[
D^{(0)} = \begin{pmatrix} 0 & 5 & 4 \\ 2 & 0 & -3 \\ 4 & \infty & 0 \end{pmatrix},\;
D^{(1)} = \begin{pmatrix} 0 & 5 & 4 \\ 2 & 0 & -3 \\ 4 & \mathbf{9} & 0 \end{pmatrix},\;
D^{(2)} = \begin{pmatrix} 0 & 5 & \mathbf{2} \\ 2 & 0 & -3 \\ 4 & 9 & 0 \end{pmatrix},\;
D^{(3)} = \begin{pmatrix} 0 & 5 & 2 \\ \mathbf{1} & 0 & -3 \\ 4 & 9 & 0 \end{pmatrix}.
\]
}
$d_{32}^{(1)} = 4 + 5$（$3 \to 1 \to 2$），$d_{13}^{(2)} = 5 + (-3)$（$1 \to 2 \to 3$），
$d_{21}^{(3)} = -3 + 4$（$2 \to 3 \to 1$）。

\item (a)、(c) 有，(b) 沒有。 (a)：有向無環圖上的路徑一定是簡單的。 若最長路徑經過 $v$，把
$s \leadsto v$ 那段換成更長的一段，接上 $v \leadsto t$ 那段後如果重複經過某個點 $u$，就表示
$u$ 能走到 $v$、$v$ 又能走到 $u$，形成環，矛盾；所以剪下貼上成立，可以依拓樸順序做 DP（第~22~
章）。 (b)：第~\ref{sec:dp-substructure}~節的三角形就是反例。 (c)：恰好 $k$ 條邊的最佳走法，
去掉最後一條邊 $(u, t)$ 後，一定是恰好 $k - 1$ 條邊到 $u$ 的最佳走法；點可以重複，兩段不會互
相限制。 這就是第~\ref{sec:dp-apsp}~節「限制邊數」的子問題。
\end{enumerate}

\chapter{貪心}
\label{ch:greedy}

% 本章的圖共用的樣式（前綴 gr-）
\tikzset{
    gr-leaf/.style={rectangle, draw=black!65, fill=coreFill, align=center,
        font=\scriptsize, inner sep=1.6pt, minimum width=0.55cm},
    gr-inner/.style={circle, draw=black!65, fill=white, minimum size=0.6cm,
        inner sep=0pt, font=\scriptsize},
    gr-new/.style={draw=memStroke, very thick, fill=memFill},
    gr-dummy/.style={rectangle, draw=black!45, dashed, fill=white, align=center,
        font=\scriptsize, inner sep=1.6pt, minimum width=0.55cm},
    gr-edge/.style={draw=black!60},
    gr-bit/.style={font=\tiny, text=black!60, inner sep=1pt},
    gr-plab/.style={font=\small},
    gr-mark/.style={draw=black!45, dashed},
}

系上的小會議室明天收到十二筆借用申請，每一筆寫著開始和結束的時刻，有些時段互相重
疊。 會議室一次只能給一筆申請用，系辦想讓借得到的申請越多越好。 有人說先排最早開始
的，有人說先排用得最短的，也有人說先排最早結束的，排進去之後，和它衝突的就不要了。
這三種做法哪一種一定能排到最多筆？ 往下讀之前，先想一想，再試著各舉一個反例。

上面三種做法都是\textbf{貪心}（greedy）：每一步只看眼前，挑一個「現在看起來最好」
的選擇，做了就不再反悔，然後去處理剩下的問題。 這樣的演算法通常很短、很快，難的是
證明它對。 很多看起來合理的貪心是錯的，所以這一章每一個貪心演算法都會附上證明，最
常用的工具是\textbf{交換論證}（exchange argument）：拿任何一個最佳解，一步一步把它
換成貪心的選擇，證明每一步都不會變差。

第~\ref{sec:gr-activity}~節用活動選擇示範交換論證，第~\ref{sec:gr-knapsack}~節比較分數背包和
0/1 背包，第~\ref{sec:gr-huffman}~節講霍夫曼編碼和外部排序的最佳合併樹，
第~\ref{sec:gr-sched}~節是兩個排程問題（第二個用互斥集合加速），第~\ref{sec:gr-fail}~節
看貪心什麼時候不對，並介紹擬陣。

這一章用到第~8~章的最小堆積、第~13~章的互斥集合、第~14~章的排序和第~15~章的外部排
序。 第~18~章的動態規劃和這一章常常互相對照：貪心解不了的問題，很多要改用動態規劃。
C 程式用到《C 程式設計基礎》第~14~章的結構、第~15~章的遞迴、第~18~章的
\texttt{qsort} 和第~19~章的 \texttt{malloc}，用到時先在框裡複習。

\section{活動選擇}
\label{sec:gr-activity}

\subsection{問題}

有 $n$ 筆申請，第 $i$ 筆要用的時段是半開區間 $[s_i, f_i)$，$s_i < f_i$：從時刻 $s_i$
開始，到時刻 $f_i$ 結束。 兩筆申請的時段沒有交集，就說它們\textbf{相容}
（compatible）。 因為是半開區間，一筆在 $f_i$ 結束，另一筆在同一個時刻 $f_i$ 開始，
不算衝突。 \textbf{活動選擇問題}（activity-selection problem）要挑出最多筆兩兩相容的
申請。

章首的十二筆申請如下（時刻用 $24$ 小時制的整點）：

\begin{center}
\begin{tabular}{c*{12}{c}}
\toprule
編號 & 1 & 2 & 3 & 4 & 5 & 6 & 7 & 8 & 9 & 10 & 11 & 12 \\
\midrule
$s_i$ & 13 & 9 & 17 & 8 & 12 & 10 & 15 & 8 & 11 & 16 & 14 & 19 \\
$f_i$ & 16 & 11 & 20 & 13 & 13 & 12 & 18 & 10 & 14 & 19 & 15 & 21 \\
\bottomrule
\end{tabular}
\end{center}

\subsection{三種貪心，兩種會錯}

章首的三種做法都是「依某個順序看申請，和已選的相容就選」，差別只在順序。 先看兩
個錯的。

\textbf{最早開始的先排。} 申請 $[8, 18)$、$[9, 10)$、$[11, 12)$：最早開始的是
$[8, 18)$，選了它，另外兩筆都和它衝突，只排到 $1$ 筆；不選它可以排 $2$ 筆。 一筆很早
開始、拖很久的申請會把後面整段時間都佔掉。

\textbf{用得最短的先排。} 申請 $[9, 12)$、$[11, 13)$、$[12, 15)$：最短的是
$[11, 13)$，兩小時，可是它同時和另外兩筆衝突。 選它只排到 $1$ 筆；選另外兩筆可以排
$2$ 筆。

\textbf{最早結束的先排。} 這一個是對的。 直覺是：先結束的申請，把會議室還回來得最
早，留給後面的時間最多。 下一小節把這個直覺變成證明。

「和別人衝突最少的先排」也會錯，只是反例要十來筆申請才造得出來。

\subsection{演算法}

先把申請依結束時刻由小到大排好，再從頭看一次：和最後選進來的那一筆相容，就選。
因為選進來的申請結束時刻是遞增的，「和已選的全部相容」只要檢查「開始時刻不早於最
後一筆的結束時刻」。

\begin{pseudo}[label=alg:gr-activity]{活動選擇（申請已依 $f$ 由小到大排好，$n \ge 1$）}
\Function{Activity-Select}{$s, f, n$}
    \State $A \gets \{1\}$
    \State $\mathit{last} \gets f[1]$ \Comment{最後選進來的那筆的結束時刻}
    \For{$i \gets 2$ \To $n$}
        \If{$s[i] \ge \mathit{last}$}
            \State $A \gets A \cup \{i\}$
            \State $\mathit{last} \gets f[i]$
        \EndIf
    \EndFor
    \State \Return $A$
\EndFunction
\end{pseudo}

排序 $\Theta(n \lg n)$，第 4--7 行的迴圈 $\Theta(n)$，總共 $\Theta(n \lg n)$。 輸入已經
依結束時刻排好的話只要 $\Theta(n)$。 結束時刻相同的申請，本書依編號由小到大排；不管
怎麼排，選到的筆數都一樣（觀念題第~\ref{q:gr-act-trace}~題可以自己驗證）。

\subsection{手算追蹤}

章首的十二筆依結束時刻排好之後逐筆檢查，$\mathit{last}$ 是檢查之前的值：

\begin{center}
\begin{tabular}{cccccl}
\toprule
編號 & $[s_i, f_i)$ & $\mathit{last}$ & $s_i \ge \mathit{last}$？ & 結果 & 選進來之後 $\mathit{last}$ \\
\midrule
8 & $[8, 10)$ & -- & -- & 選（第一筆） & 10 \\
2 & $[9, 11)$ & 10 & $9 < 10$ & 跳過 & \\
6 & $[10, 12)$ & 10 & $10 \ge 10$ & 選 & 12 \\
4 & $[8, 13)$ & 12 & $8 < 12$ & 跳過 & \\
5 & $[12, 13)$ & 12 & $12 \ge 12$ & 選 & 13 \\
9 & $[11, 14)$ & 13 & $11 < 13$ & 跳過 & \\
11 & $[14, 15)$ & 13 & $14 \ge 13$ & 選 & 15 \\
1 & $[13, 16)$ & 15 & $13 < 15$ & 跳過 & \\
7 & $[15, 18)$ & 15 & $15 \ge 15$ & 選 & 18 \\
10 & $[16, 19)$ & 18 & $16 < 18$ & 跳過 & \\
3 & $[17, 20)$ & 18 & $17 < 18$ & 跳過 & \\
12 & $[19, 21)$ & 18 & $19 \ge 18$ & 選 & 21 \\
\bottomrule
\end{tabular}
\end{center}

選到 $8, 6, 5, 11, 7, 12$ 共 $6$ 筆。 第 4 筆和第 5 筆都在 $13$ 點結束，依編號第 4 筆
排前面；它在 $8$ 點開始，和第 6 筆衝突，所以輪到它時被跳過。 $10$、$12$、$15$ 這三個時
刻，前一筆剛結束、下一筆立刻開始，在半開區間下不算衝突。

\begin{examnote}
區間題先看清楚端點：題目寫成閉區間 $[s_i, f_i]$，或寫「時刻 $f_i$ 還在使用」，接在後
面的申請就要 $s_j > f_i$ 才相容，虛擬碼~\ref{alg:gr-activity} 第 5 行的 $\ge$ 要改成
$>$，答案可能變少。 作答時寫一句你用的是哪一種，排序時結束時刻相同怎麼排也寫出來。
\end{examnote}

\subsection{為什麼最早結束的一定對}
\label{sec:gr-act-proof}

證明分兩步。 第一步說明「第一個選擇不會錯」，第二步說明「選了之後，剩下的是同一種
問題」。

\begin{theorem}
\label{thm:gr-act-choice}
設 $S$ 是一群申請，$a$ 是 $S$ 裡結束時刻最早的一筆（有好幾筆一樣早時任選一筆）。 那
麼 $S$ 有一個最大的相容子集合包含 $a$。
\end{theorem}

\begin{proof}
任取一個最大的相容子集合 $O$，令 $o$ 是 $O$ 裡結束最早的一筆。 如果 $o = a$，已經得
證。 否則令 $O' = (O - \{o\}) \cup \{a\}$，也就是把 $o$ 換成 $a$。

$O$ 裡其他的申請 $q$ 都和 $o$ 相容，而且 $f_q \ge f_o$。 $q$ 不可能整個在 $o$ 前面
（那樣 $f_q \le s_o < f_o$），所以 $q$ 在 $o$ 後面：$s_q \ge f_o$。 $a$ 是全部申請裡
結束最早的，$f_a \le f_o \le s_q$，所以 $a$ 和每一個 $q$ 都相容。 $O'$ 是相容子集合，
筆數和 $O$ 一樣，也是最大的，而且包含 $a$。
\end{proof}

\textbf{第二步。} 選了 $a$ 之後，還能選的是 $S' = \{q \in S : s_q \ge f_a\}$，其他申
請都和 $a$ 衝突。 如果 $O$ 是 $S$ 的最大相容子集合而且包含 $a$，那麼 $O - \{a\}$ 一
定是 $S'$ 的最大相容子集合：它是 $S'$ 的相容子集合；如果 $S'$ 有更大的相容子集合
$P$，那 $P \cup \{a\}$ 就比 $O$ 大，矛盾。

\textbf{合起來。} 對申請筆數做歸納。 虛擬碼~\ref{alg:gr-activity} 先選了 $S$ 裡結束
最早的 $a$；之後它跳過的申請都在 $f_a$ 之前開始，正好是不在 $S'$ 裡的那些，所以它接
下來做的事，就是對 $S'$ 從頭跑一次同一個演算法。 由歸納假設，它在 $S'$ 選到的筆數就是
$S'$ 的最大筆數 $m'$；由定理~\ref{thm:gr-act-choice} 和第二步，$S$ 的最大筆數是
$1 + m'$，貪心選到的也是 $1 + m'$。

\subsection{貪心法的兩個要素}
\label{sec:gr-elements}

上面的證明拆成兩件事，幾乎每個貪心演算法的證明都是這兩件：

\begin{enumerate}
\item \textbf{貪心選擇性質}（greedy-choice property）：存在一個最佳解，它的第一個選
擇就是貪心的選擇。 通常用交換論證：任取一個最佳解，把它和貪心不同的那個選擇換成貪心
的，說明換完仍然可行、而且不會變差。
\item \textbf{最佳子結構}（optimal substructure）：做了貪心的選擇之後，剩下的是一個同
類型、比較小的問題，而且原問題的最佳解扣掉那個選擇，正好是小問題的最佳解。
\end{enumerate}

第~18~章的動態規劃也需要最佳子結構。 兩者的差別在做選擇的時機：動態規劃先把所有子
問題的答案算出來，再比較各種選擇；貪心在解子問題之前就做了選擇，而且只留下一個子問
題。 所以貪心比較快，但前提是貪心選擇性質成立，這一點每個問題都要另外證明。

寫交換論證時常漏掉兩件事：換完之後解還是\textbf{可行}的（上面的證明花了一半篇幅說
$a$ 和其他申請都相容），以及第一步之後為什麼也對（要靠最佳子結構加歸納，或像觀念題
第~\ref{q:gr-ahead}~題那樣逐步比較）。

\subsection{C 程式}

\begin{cbookrecap}[第~14~章結構、第~18~章 18.6 節]
\texttt{qsort(base, n, size, cmp)} 排序從 \texttt{base} 開始的 \texttt{n} 個元素，每個
\texttt{size} 個位元組。 比較函式拿到兩個 \texttt{const void *}，先轉成指向元素型別的
指標再取欄位；第一個要排前面就回傳負數，後面就回傳正數，相等回傳 $0$。 用
\texttt{x->f - y->f} 回傳差值可能溢位（本書第~14~章 14.7 節的例題），所以下面的程式
一律用比較運算子。
\end{cbookrecap}

程式讀進 $n$ 筆申請，第 $i$ 行是第 $i$ 筆的 $s_i$ 和 $f_i$。 申請存成結構，排序時把編
號一起帶著走，邊掃描邊印出選中的申請。

\begin{example}[label=ex:gr-activity]{c}{(activity.c) 活動選擇：最早結束的先排}
#include <stdio.h>
#include <stdlib.h>

typedef struct {
    int id;         // 申請編號，從 1 開始
    int s, f;       // 借用時段 [s, f)
} Booking;

// 依結束時刻由小到大；結束時刻相同時編號小的在前。 不用減法，避免溢位。
int by_finish(const void *p, const void *q)
{
    const Booking *x = p, *y = q;
    if (x->f != y->f) {
        return x->f < y->f ? -1 : 1;
    }
    return (x->id > y->id) - (x->id < y->id);
}

int main(void)
{
    int n;
    if (scanf("%d", &n) != 1 || n < 0) {
        return 1;
    }
    Booking *b = malloc((n + 1) * sizeof *b);   // 多一格：n = 0 時不 malloc(0)
    if (b == NULL) {
        return 1;
    }
    for (int i = 0; i < n; i++) {
        if (scanf("%d %d", &b[i].s, &b[i].f) != 2 || b[i].s >= b[i].f) {
            free(b);
            return 1;
        }
        b[i].id = i + 1;
    }
    qsort(b, n, sizeof b[0], by_finish);
    int k = 0, last = 0;        // 選了幾筆、最後一筆選中的結束時刻
    for (int i = 0; i < n; i++) {
        if (k == 0 || b[i].s >= last) {         // 半開區間：接著開始不算衝突
            printf("#%d [%d, %d)\n", b[i].id, b[i].s, b[i].f);
            k++;
            last = b[i].f;
        }
    }
    printf("accepted %d of %d\n", k, n);
    free(b);
    return 0;
}
\end{example}
\iosample{12 \\ 13 16 \\ 9 11 \\ 17 20 \\ 8 13 \\ 12 13 \\ 10 12 \\ 15 18 \\ 8 10 \\ 11 14 \\ 16 19 \\ 14 15 \\ 19 21}{\#8 [8, 10) \\ \#6 [10, 12) \\ \#5 [12, 13) \\ \#11 [14, 15) \\ \#7 [15, 18) \\ \#12 [19, 21) \\ accepted 6 of 12}

輸出和手算的表一樣。 \texttt{malloc} 多要一格，是因為 C 標準允許 \texttt{malloc(0)}
回傳 \texttt{NULL}，$n = 0$ 時程式會誤以為要不到記憶體。

\section{分數背包與 0/1 背包}
\label{sec:gr-knapsack}

\subsection{分數背包}

咖啡豆批發行有 $n$ 種豆子，第 $i$ 種的庫存是 $w_i$ 公斤，整批的價值是 $v_i$ 元（$w_i$、
$v_i$ 都是正數）。 袋子最多裝 $W$ 公斤。 每一種豆子都可以只拿一部分：拿第 $i$ 種的比
例 $x_i$（$0 \le x_i \le 1$），得到 $x_i v_i$ 元，佔 $x_i w_i$ 公斤。 在
$\sum x_i w_i \le W$ 的條件下讓 $\sum x_i v_i$ 最大，這叫\textbf{分數背包問題}
（fractional knapsack problem）。

既然可以拆，每一公斤的價值才是重點。 令 $r_i = v_i / w_i$ 是第 $i$ 種的\textbf{單位
重量價值}。 貪心的做法是：依 $r_i$ 由大到小，能整批拿就整批拿，袋子快滿時拿最後一種
的一部分把袋子裝滿。

\begin{pseudo}[label=alg:gr-frac]{分數背包（已依 $v[i] / w[i]$ 由大到小排好）}
\Function{Fractional-Knapsack}{$w, v, n, W$}
    \State $\mathit{room} \gets W$
    \State $\mathit{total} \gets 0$
    \For{$i \gets 1$ \To $n$}
        \State $x_i \gets \min(1, \mathit{room} / w[i])$ \Comment{第 $i$ 種拿的比例}
        \State $\mathit{total} \gets \mathit{total} + x_i \cdot v[i]$
        \State $\mathit{room} \gets \mathit{room} - x_i \cdot w[i]$
    \EndFor
    \State \Return $\mathit{total}$
\EndFunction
\end{pseudo}

袋子滿了以後 $\mathit{room} = 0$，之後每一種的 $x_i$ 都是 $0$。 時間花在排序，
$\Theta(n \lg n)$。

\textbf{例子。} $W = 10$ 公斤，五種豆子：

\begin{center}
\begin{tabular}{cccccc}
\toprule
豆子 & A & B & C & D & E \\
\midrule
庫存 $w_i$（公斤） & 4 & 3 & 5 & 2 & 6 \\
整批價值 $v_i$（元） & 2000 & 1950 & 2250 & 1400 & 2100 \\
每公斤 $r_i$ & 500 & 650 & 450 & 700 & 350 \\
\bottomrule
\end{tabular}
\end{center}

依 $r_i$ 由大到小是 D、B、A、C、E。 D 整批拿（$2$ 公斤，$1400$ 元），B 整批拿（共
$5$ 公斤，$3350$ 元），A 整批拿（共 $9$ 公斤，$5350$ 元）。 剩 $1$ 公斤，拿 C 的
$1/5$，$450$ 元。 總共 $5800$ 元，E 完全沒拿。

\begin{theorem}
\label{thm:gr-frac}
虛擬碼~\ref{alg:gr-frac} 得到分數背包的最大總價值。
\end{theorem}

\begin{proof}
令 $x_1, \ldots, x_n$ 是貪心的解，$y_1, \ldots, y_n$ 是任何一個最佳解。 如果兩者不
同，取第一個不同的位置 $j$。 貪心在第 $j$ 種拿到「能拿的最多」，所以 $y_j < x_j$。
最佳解在第 $j$ 種少拿了 $\delta = (x_j - y_j) w_j$ 公斤。 前 $j - 1$ 種兩者一樣，貪心
拿完第 $j$ 種沒有超重，所以最佳解在第 $j + 1$ 種以後的重量加上袋子的空位，等於
$W - \sum_{i<j} x_i w_i - y_j w_j \ge x_j w_j - y_j w_j = \delta$。 先用袋子的空位，不夠的
部分再從第 $j + 1$ 種以後拿出豆子，湊滿 $\delta$ 公斤，換成第 $j$ 種的 $\delta$ 公斤。
被拿出的豆子每公斤價值都 $\le r_j$，空位的價值是 $0$，所以總價值不會變少，重量也沒有超
過。 換完之後仍是最佳解，而且和貪心解第一個不同的位置往後移了。 重複這個交換，最後最佳
解變成貪心解，所以貪心解也是最佳的。
\end{proof}

\subsection{0/1 背包}

把豆子換成不能拆的東西：每一件只能整件拿或不拿，$x_i \in \{0, 1\}$。 這是
\textbf{0/1 背包問題}（0/1 knapsack problem）。 照單位重量價值拿還對嗎？ 容量 $W = 8$，
三件物品：甲重 $5$、值 $60$（每單位 $12$），乙、丙都是重 $4$、值 $44$（每單位 $11$）。

依單位重量價值，先拿甲（$5$ 公斤，$60$），剩 $3$ 公斤，乙、丙都放不下，總價值 $60$。
不拿甲，拿乙和丙，剛好 $8$ 公斤，總價值 $88$。 貪心錯了。

「價值大的先拿」在這組資料上也是先拿甲；「輕的先拿」在容量 $8$、物品（重 $2$，值
$2$）和（重 $7$，值 $70$）上只得到 $2$。

問題出在定理~\ref{thm:gr-frac} 的證明哪一步？ 交換的時候要從別的物品「拿出 $\delta$
公斤」換成第 $j$ 種，這一步需要把東西切開。 0/1 背包不能切，拿了甲之後剩下的 $3$ 公
斤空間就浪費了，而這個浪費會不會發生，要看後面還有哪些物品，貪心在做第一個選擇時看
不到。 0/1 背包要用第~18~章的動態規劃，時間是 $O(nW)$，這是第~1~章 1.9 節說的擬多
項式時間；第~28~章會看到，一般的 0/1 背包是 NP 完備的，目前不知道有多項式時間的演算
法。

兩個問題都有最佳子結構，差別只在貪心選擇性質，可見「有最佳子結構」不足以讓貪心成立。

另一個用途：分數背包的最佳值一定不小於同一組資料的 0/1 背包最佳值（0/1 的解也是分數
的解），所以它可以當 0/1 背包的\textbf{上界}。 第~20~章的分支限界就用這個上界來剪
枝。

\section{霍夫曼編碼與最佳合併樹}
\label{sec:gr-huffman}

\subsection{定長碼與變長碼}

山上的氣象站每小時記錄一個天氣代碼，一共七種：晴 S、多雲 C、陰 O、雨 R、雷雨 T、霧
F、霾 H。 某一段時間的 $100$ 筆紀錄裡，各代碼出現的次數是：

\begin{center}
\begin{tabular}{c*{7}{c}}
\toprule
代碼 & S & C & O & R & T & F & H \\
\midrule
次數 & 34 & 22 & 15 & 13 & 8 & 5 & 3 \\
\bottomrule
\end{tabular}
\end{center}

要把紀錄存成 $0$、$1$ 組成的位元串。 七種代碼，每種給一個 $3$ 位元的碼（$2^3 = 8 \ge
7$）就夠了，$100$ 筆共 $300$ 位元。 這叫\textbf{定長碼}（fixed-length code）。 晴天
佔了三分之一，如果給 S 比較短的碼、給少見的 H 比較長的碼，總長度可以更短，這叫
\textbf{變長碼}（variable-length code）。

變長碼要能解得開。 如果 S 的碼是 \texttt{1}、C 的碼是 \texttt{10}，讀到 \texttt{10}
時不知道是 C，還是 S 後面接著另一個碼。 最簡單的保證是：\textbf{沒有一個碼是另一個碼
的前綴}（開頭的一段）。 這樣的碼叫\textbf{前綴碼}（prefix code）。 解前綴碼時從頭讀，
一讀到某個完整的碼就可以確定是它，不必往後看。 任何能唯一解碼的碼，都有碼長完全一樣的前綴碼（Kraft--McMillan 不等式），所以只看前綴
碼不會吃虧。

\subsection{碼樹}

前綴碼可以畫成一棵二元樹：從根出發，往左走記 \texttt{0}，往右走記 \texttt{1}，每個
符號放在一個葉節點，根到葉節點的路徑就是它的碼。 符號都在葉節點上，所以沒有一個碼是
另一個的前綴；反過來，任何前綴碼都能這樣畫成一棵樹。 解碼就是從根往下走，走到葉節點
就輸出那個符號，再回到根。

\textbf{碼長和深度。} 本書照第~7~章的定義：節點的\textbf{深度}數的是從根往下走的
\textbf{邊數}，根的深度是 $0$；樹的\textbf{高度}數的是\textbf{節點數}（層數）。 每走
一條邊寫一個位元，所以
\[
\text{符號 $c$ 的碼長} = \text{葉節點 $c$ 的深度 } d_T(c) \text{（邊數）}, \qquad
\text{最長的碼長} = \text{樹的高度} - 1.
\]
符號 $c$ 出現 $c.\mathit{freq}$ 次，整份資料編碼後的總位元數是
\[
B(T) = \sum_{c \in C} c.\mathit{freq} \cdot d_T(c),
\]
這也叫這棵樹的\textbf{加權路徑長}（weighted path length，只算葉節點）。 我們要找
$B(T)$ 最小的樹，對應的碼叫\textbf{最佳前綴碼}。

\textbf{最佳的碼樹是嚴格二元樹。} 如果某個內部節點只有一個孩子，把它刪掉、讓孩子直接
接上去，它底下每個葉節點的深度都少 $1$，$B(T)$ 變小（頻率都是正的）。 所以最佳的碼
樹每個內部節點都有兩個孩子，是第~7~章的嚴格二元樹；$n$ 個符號的碼樹有 $n$ 個葉節點、
$n - 1$ 個內部節點（$n_0 = n_2 + 1$）。 定長碼的樹一般不是嚴格的：七個符號用 $3$ 位
元，有一個碼（例如 \texttt{111}）沒有人用，它的兄弟可以往上提一層。

\subsection{霍夫曼演算法}

霍夫曼（D. A. Huffman，1952）的做法由下往上建樹。 一開始每個符號自己是一棵只有一個
節點的樹，權重是它的頻率。 每一步拿出權重最小的兩棵樹，接到一個新的根底下，新根的權
重是兩者的和，放回去。 每一步樹少一棵，$n - 1$ 步之後只剩一棵，就是碼樹。

「拿出最小的兩個、放回一個」正是第~8~章最小優先佇列的操作。 權重一樣的時候要有固定
的規則，否則同一組資料可以建出不同的樹。 本書的規則：

\begin{itemize}
\item 每個節點有一個編號 $\mathit{id}$。 葉節點照輸入的順序編 $0, 1, \ldots, n - 1$，
內部節點照建立的順序接著編 $n, n + 1, \ldots$。
\item 優先佇列先比權重，權重一樣時比編號，編號小的先出來。 所以權重相同時，葉節點比內
部節點先出來，早建的內部節點比晚建的先出來。
\item 先取出的當左孩子（邊標 \texttt{0}），後取出的當右孩子（邊標 \texttt{1}）。
\end{itemize}

\begin{pseudo}[label=alg:gr-huffman]{霍夫曼編碼}
\Function{Huffman}{$C$} \Comment{$C$：$n \ge 2$ 個符號，各有 $\mathit{freq}$}
    \State 建立最小優先佇列 $Q$，放進 $C$ 的每個符號（鍵是 $(\mathit{freq}, \mathit{id})$）
    \State $k \gets |C|$ \Comment{下一個內部節點的編號}
    \While{$Q$ 裡不只一個節點}
        \State $x \gets \Call{Extract-Min}{Q}$
        \State $y \gets \Call{Extract-Min}{Q}$
        \State 建立節點 $z$：$z.\mathit{left} \gets x$，$z.\mathit{right} \gets y$，
        $z.\mathit{freq} \gets x.\mathit{freq} + y.\mathit{freq}$，$z.\mathit{id} \gets k$
        \State $k \gets k + 1$
        \State \Call{Insert}{$Q, z$}
    \EndWhile
    \State \Return \Call{Extract-Min}{$Q$} \Comment{碼樹的根}
\EndFunction
\end{pseudo}

用二元堆積實作，第 2 行由下而上建堆積是 $O(n)$（第~8~章 8.2 節），迴圈跑 $n - 1$ 次，
每次三個堆積操作，各 $O(\lg n)$，總共 $O(n \lg n)$。

\subsection{逐步建樹}
\label{sec:gr-huff-trace}

用氣象站的資料跑一次。 下面每一張圖是一次合併之後，優先佇列裡的所有樹，由左到右照
$(\mathit{freq}, \mathit{id})$ 排列，所以下一步要取出的永遠是最左邊兩棵。 方框是葉節
點（代碼和次數），圓圈是內部節點（權重），橘色是這一步新建的節點。 圖是用程式照上面
的規則模擬後畫出來的。

\begin{center}
\begin{minipage}{\linewidth}
\centering
\begin{tikzpicture}[baseline=(current bounding box.north)]
\node[gr-leaf] (ha0) at (6.600,0.000) {S\\34};
\node[gr-leaf] (ha1) at (5.500,0.000) {C\\22};
\node[gr-leaf] (ha2) at (4.400,0.000) {O\\15};
\node[gr-leaf] (ha3) at (3.300,0.000) {R\\13};
\node[gr-leaf] (ha4) at (2.200,0.000) {T\\8};
\node[gr-leaf] (ha5) at (1.100,0.000) {F\\5};
\node[gr-leaf] (ha6) at (0.000,0.000) {H\\3};
\node[gr-plab] at (3.300,-1.750) {(a) 開始：七棵單節點的樹};
\end{tikzpicture}\hspace{0.7cm}\begin{tikzpicture}[baseline=(current bounding box.north)]
\node[gr-leaf] (hb0) at (6.320,0.000) {S\\34};
\node[gr-leaf] (hb1) at (5.220,0.000) {C\\22};
\node[gr-leaf] (hb2) at (4.120,0.000) {O\\15};
\node[gr-leaf] (hb3) at (3.020,0.000) {R\\13};
\node[gr-leaf] (hb4) at (0.000,0.000) {T\\8};
\node[gr-leaf] (hb5) at (1.920,-0.900) {F\\5};
\node[gr-leaf] (hb6) at (1.100,-0.900) {H\\3};
\node[gr-inner,gr-new] (hb7) at (1.510,0.000) {8};
\draw[gr-edge] (hb7) -- node[gr-bit, left] {0} (hb6);
\draw[gr-edge] (hb7) -- node[gr-bit, right] {1} (hb5);
\node[gr-plab] at (3.160,-1.750) {(b) 第 1 次：$3 + 5 = 8$};
\end{tikzpicture}

\bigskip
\begin{tikzpicture}[baseline=(current bounding box.north)]
\node[gr-leaf] (hc0) at (6.040,0.000) {S\\34};
\node[gr-leaf] (hc1) at (4.940,0.000) {C\\22};
\node[gr-leaf] (hc2) at (1.100,0.000) {O\\15};
\node[gr-leaf] (hc3) at (0.000,0.000) {R\\13};
\node[gr-leaf] (hc4) at (2.200,-0.900) {T\\8};
\node[gr-leaf] (hc5) at (3.840,-1.800) {F\\5};
\node[gr-leaf] (hc6) at (3.020,-1.800) {H\\3};
\node[gr-inner] (hc7) at (3.430,-0.900) {8};
\node[gr-inner,gr-new] (hc8) at (2.815,0.000) {16};
\draw[gr-edge] (hc7) -- node[gr-bit, left] {0} (hc6);
\draw[gr-edge] (hc7) -- node[gr-bit, right] {1} (hc5);
\draw[gr-edge] (hc8) -- node[gr-bit, left] {0} (hc4);
\draw[gr-edge] (hc8) -- node[gr-bit, right] {1} (hc7);
\node[gr-plab] at (3.020,-2.650) {(c) 第 2 次：$8 + 8 = 16$};
\end{tikzpicture}\hspace{0.7cm}\begin{tikzpicture}[baseline=(current bounding box.north)]
\node[gr-leaf] (hd0) at (5.760,0.000) {S\\34};
\node[gr-leaf] (hd1) at (2.740,0.000) {C\\22};
\node[gr-leaf] (hd2) at (4.660,-0.900) {O\\15};
\node[gr-leaf] (hd3) at (3.840,-0.900) {R\\13};
\node[gr-leaf] (hd4) at (0.000,-0.900) {T\\8};
\node[gr-leaf] (hd5) at (1.640,-1.800) {F\\5};
\node[gr-leaf] (hd6) at (0.820,-1.800) {H\\3};
\node[gr-inner] (hd7) at (1.230,-0.900) {8};
\node[gr-inner] (hd8) at (0.615,0.000) {16};
\node[gr-inner,gr-new] (hd9) at (4.250,0.000) {28};
\draw[gr-edge] (hd7) -- node[gr-bit, left] {0} (hd6);
\draw[gr-edge] (hd7) -- node[gr-bit, right] {1} (hd5);
\draw[gr-edge] (hd8) -- node[gr-bit, left] {0} (hd4);
\draw[gr-edge] (hd8) -- node[gr-bit, right] {1} (hd7);
\draw[gr-edge] (hd9) -- node[gr-bit, left] {0} (hd3);
\draw[gr-edge] (hd9) -- node[gr-bit, right] {1} (hd2);
\node[gr-plab] at (2.880,-2.650) {(d) 第 3 次：$13 + 15 = 28$};
\end{tikzpicture}
\diagcaption{霍夫曼演算法 (a)--(d)：開始和前三次合併。 每一張圖由左到右是優先佇列裡的樹，最左邊兩棵就是下一步要取出的。 邊上的 \texttt{0}、\texttt{1} 是先取出（左）和後取出（右）。}
\label{fig:gr-huff1}
\end{minipage}
\end{center}

第 2 次合併是平手的例子：佇列裡有兩個權重 $8$ 的節點，葉節點 T（編號 $4$）和第 1 次
建的內部節點（編號 $7$）。 照規則葉節點先出來，T 當左孩子。

\begin{center}
\begin{minipage}{\linewidth}
\centering
\begin{tikzpicture}[baseline=(current bounding box.north)]
\node[gr-leaf] (he0) at (1.920,0.000) {S\\34};
\node[gr-leaf] (he1) at (5.480,-0.900) {C\\22};
\node[gr-leaf] (he2) at (0.820,-0.900) {O\\15};
\node[gr-leaf] (he3) at (0.000,-0.900) {R\\13};
\node[gr-leaf] (he4) at (3.020,-1.800) {T\\8};
\node[gr-leaf] (he5) at (4.660,-2.700) {F\\5};
\node[gr-leaf] (he6) at (3.840,-2.700) {H\\3};
\node[gr-inner] (he7) at (4.250,-1.800) {8};
\node[gr-inner] (he8) at (3.635,-0.900) {16};
\node[gr-inner] (he9) at (0.410,0.000) {28};
\node[gr-inner,gr-new] (he10) at (4.558,0.000) {38};
\draw[gr-edge] (he9) -- node[gr-bit, left] {0} (he3);
\draw[gr-edge] (he9) -- node[gr-bit, right] {1} (he2);
\draw[gr-edge] (he7) -- node[gr-bit, left] {0} (he6);
\draw[gr-edge] (he7) -- node[gr-bit, right] {1} (he5);
\draw[gr-edge] (he8) -- node[gr-bit, left] {0} (he4);
\draw[gr-edge] (he8) -- node[gr-bit, right] {1} (he7);
\draw[gr-edge] (he10) -- node[gr-bit, left] {0} (he8);
\draw[gr-edge] (he10) -- node[gr-bit, right] {1} (he1);
\node[gr-plab] at (2.740,-3.550) {(e) 第 4 次：$16 + 22 = 38$};
\end{tikzpicture}\hspace{0.7cm}\begin{tikzpicture}[baseline=(current bounding box.north)]
\node[gr-leaf] (hf0) at (5.200,-0.900) {S\\34};
\node[gr-leaf] (hf1) at (2.460,-0.900) {C\\22};
\node[gr-leaf] (hf2) at (4.380,-1.800) {O\\15};
\node[gr-leaf] (hf3) at (3.560,-1.800) {R\\13};
\node[gr-leaf] (hf4) at (0.000,-1.800) {T\\8};
\node[gr-leaf] (hf5) at (1.640,-2.700) {F\\5};
\node[gr-leaf] (hf6) at (0.820,-2.700) {H\\3};
\node[gr-inner] (hf7) at (1.230,-1.800) {8};
\node[gr-inner] (hf8) at (0.615,-0.900) {16};
\node[gr-inner] (hf9) at (3.970,-0.900) {28};
\node[gr-inner] (hf10) at (1.538,0.000) {38};
\node[gr-inner,gr-new] (hf11) at (4.585,0.000) {62};
\draw[gr-edge] (hf7) -- node[gr-bit, left] {0} (hf6);
\draw[gr-edge] (hf7) -- node[gr-bit, right] {1} (hf5);
\draw[gr-edge] (hf8) -- node[gr-bit, left] {0} (hf4);
\draw[gr-edge] (hf8) -- node[gr-bit, right] {1} (hf7);
\draw[gr-edge] (hf10) -- node[gr-bit, left] {0} (hf8);
\draw[gr-edge] (hf10) -- node[gr-bit, right] {1} (hf1);
\draw[gr-edge] (hf9) -- node[gr-bit, left] {0} (hf3);
\draw[gr-edge] (hf9) -- node[gr-bit, right] {1} (hf2);
\draw[gr-edge] (hf11) -- node[gr-bit, left] {0} (hf9);
\draw[gr-edge] (hf11) -- node[gr-bit, right] {1} (hf0);
\node[gr-plab] at (2.600,-3.550) {(f) 第 5 次：$28 + 34 = 62$};
\end{tikzpicture}

\bigskip
\begin{tikzpicture}[baseline=(current bounding box.north)]
\node[gr-leaf] (hg0) at (4.920,-1.800) {S\\34};
\node[gr-leaf] (hg1) at (2.460,-1.800) {C\\22};
\node[gr-leaf] (hg2) at (4.100,-2.700) {O\\15};
\node[gr-leaf] (hg3) at (3.280,-2.700) {R\\13};
\node[gr-leaf] (hg4) at (0.000,-2.700) {T\\8};
\node[gr-leaf] (hg5) at (1.640,-3.600) {F\\5};
\node[gr-leaf] (hg6) at (0.820,-3.600) {H\\3};
\node[gr-inner] (hg7) at (1.230,-2.700) {8};
\node[gr-inner] (hg8) at (0.615,-1.800) {16};
\node[gr-inner] (hg9) at (3.690,-1.800) {28};
\node[gr-inner] (hg10) at (1.538,-0.900) {38};
\node[gr-inner] (hg11) at (4.305,-0.900) {62};
\node[gr-inner,gr-new] (hg12) at (2.921,0.000) {100};
\draw[gr-edge] (hg7) -- node[gr-bit, left] {0} (hg6);
\draw[gr-edge] (hg7) -- node[gr-bit, right] {1} (hg5);
\draw[gr-edge] (hg8) -- node[gr-bit, left] {0} (hg4);
\draw[gr-edge] (hg8) -- node[gr-bit, right] {1} (hg7);
\draw[gr-edge] (hg10) -- node[gr-bit, left] {0} (hg8);
\draw[gr-edge] (hg10) -- node[gr-bit, right] {1} (hg1);
\draw[gr-edge] (hg9) -- node[gr-bit, left] {0} (hg3);
\draw[gr-edge] (hg9) -- node[gr-bit, right] {1} (hg2);
\draw[gr-edge] (hg11) -- node[gr-bit, left] {0} (hg9);
\draw[gr-edge] (hg11) -- node[gr-bit, right] {1} (hg0);
\draw[gr-edge] (hg12) -- node[gr-bit, left] {0} (hg10);
\draw[gr-edge] (hg12) -- node[gr-bit, right] {1} (hg11);
\node[gr-plab] at (2.460,-4.450) {(g) 第 6 次：$38 + 62 = 100$};
\end{tikzpicture}
\diagcaption{霍夫曼演算法 (e)--(g)：第 4 到第 6 次合併。 (g) 是最後的碼樹，根的權重 $100$ 是紀錄的總筆數。 從根走到葉節點，把邊上的位元串起來就是那個代碼的碼。}
\label{fig:gr-huff2}
\end{minipage}
\end{center}
第 4 次合併時，權重 $16$ 的內部節點比 C 的 $22$ 小，所以它當左孩子，C 當右孩子；所以
C 的碼是 \texttt{01}，不是 \texttt{00}。

從圖 (g) 讀出每個代碼的碼，再和定長碼比較：

\begin{center}
\begin{tabular}{cccccc}
\toprule
代碼 & 次數 & 霍夫曼碼 & 碼長（深度） & 位元數 & 定長碼的位元數 \\
\midrule
S & 34 & \texttt{11} & 2 & 68 & 102 \\
C & 22 & \texttt{01} & 2 & 44 & 66 \\
O & 15 & \texttt{101} & 3 & 45 & 45 \\
R & 13 & \texttt{100} & 3 & 39 & 39 \\
T & 8 & \texttt{000} & 3 & 24 & 24 \\
F & 5 & \texttt{0011} & 4 & 20 & 15 \\
H & 3 & \texttt{0010} & 4 & 12 & 9 \\
\midrule
合計 & 100 & & & 252 & 300 \\
\bottomrule
\end{tabular}
\end{center}

總共 $252$ 位元，平均每筆 $2.52$ 位元，比定長碼省了 $16\%$。 F 和 H 的碼反而比定長碼
長，但它們很少出現。 驗算 $B(T)$ 有另一個方法：它等於所有\textbf{內部節點}的權重和，
$8 + 16 + 28 + 38 + 62 + 100 = 252$。 原因是每個符號被它的每一個祖先（內部節點）各算
一次，祖先的個數正好是它的深度；第~\ref{sec:gr-merge}~節會再用到這個事實。

\textbf{解碼。} 位元串 \texttt{1101100101110010} 從根開始走：\texttt{11} 到 S，回到
根；\texttt{01} 到 C；\texttt{100} 到 R；\texttt{101} 到 O；\texttt{11} 到 S；
\texttt{0010} 到 H。 解出 SCROSH。

\begin{examnote}
霍夫曼樹的追蹤題，權重平手時選哪一個、誰當左孩子，會改變樹的形狀和每個符號的碼，甚至
改變碼長。 例如頻率 $1, 1, 2, 2$：先合併兩個 $1$ 得到 $2$ 以後，三個 $2$ 平手。 合併
兩個原本的 $2$，四個碼長都是 $2$；合併新的 $2$ 和一個原本的 $2$，碼長是 $1, 2, 3, 3$。
兩種樹的總位元數都是 $12$，都是最佳的。 不同課本的平手規則和左右規則不同，作答時照題
目的規定；題目沒有規定，就先寫出你用的規則（例如本書的「先比權重再比編號，先取出的放
左邊」），閱卷的人才能照你的規則核對。 總位元數不受這些規則影響，可以拿來驗算。
\end{examnote}

\subsection{為什麼霍夫曼碼最佳}
\label{sec:gr-huff-proof}

證明的架構和活動選擇一樣：第一個引理是貪心選擇性質，第二個是最佳子結構。 以下 $C$
是符號的集合，頻率都是正數，$|C| \ge 2$。

\begin{lemma}
\label{lem:gr-huff-choice}
設 $x$、$y$ 是 $C$ 裡頻率最小的兩個符號。 存在一棵最佳碼樹，$x$ 和 $y$ 是兄弟，而且
都在最深的那一層。
\end{lemma}

\begin{proof}
任取一棵最佳碼樹 $T$。 它是嚴格二元樹，所以最深的葉節點 $a$ 一定有兄弟 $b$；$b$ 也是
葉節點，否則 $b$ 底下的葉節點比 $a$ 更深。 $d_T(a) = d_T(b)$ 是最大的深度。

想法是把 $x$、$y$ 搬到 $a$、$b$ 的位置。 先看交換兩個葉節點 $u$、$v$ 的效果：只有它們
的深度改變，
\[
B(\text{換後}) - B(\text{換前}) = \bigl(u.\mathit{freq} - v.\mathit{freq}\bigr)
\bigl(d(v) - d(u)\bigr),
\]
所以把頻率比較小的符號換到比較深（或一樣深）的位置，$B$ 不會變大。

如果 $x$ 已經是 $a$ 或 $b$ 其中一個，把名字對調讓 $x = a$；$y$ 也一樣，讓 $y = b$。
剩下沒對好的，逐一交換：$x \ne a$ 時交換 $x$ 和 $a$，$y \ne b$ 時交換 $y$ 和 $b$。 名字
對調過之後，被換上去的 $a$ 或 $b$ 都不是 $x$、$y$，頻率不小於 $y.\mathit{freq} \ge
x.\mathit{freq}$；被換下去的 $x$ 或 $y$ 到了最深的位置。 由上面的式子，每次交換 $B$ 都
不會變大。 換完的樹仍是最佳的，$x$、$y$ 是兄弟，都在最深的一層。
\end{proof}

\begin{lemma}
\label{lem:gr-huff-sub}
設 $x$、$y$ 是 $C$ 裡頻率最小的兩個符號。 令 $C'$ 是把 $x$、$y$ 拿掉、加入一個新符號
$z$，$z.\mathit{freq} = x.\mathit{freq} + y.\mathit{freq}$。 如果 $T'$ 是 $C'$ 的
最佳碼樹，把 $T'$ 的葉節點 $z$ 換成一個內部節點、底下接 $x$ 和 $y$，得到的樹 $T$ 是
$C$ 的最佳碼樹。
\end{lemma}

\begin{proof}
先比較 $T$ 和 $T'$。 除了 $x$、$y$、$z$，每個符號的深度都一樣；$x$、$y$ 的深度是
$d_{T'}(z) + 1$。 所以
\[
B(T) = B(T') - z.\mathit{freq} \cdot d_{T'}(z)
+ (x.\mathit{freq} + y.\mathit{freq})\bigl(d_{T'}(z) + 1\bigr)
= B(T') + x.\mathit{freq} + y.\mathit{freq}.
\]
同樣的計算對任何一棵「$x$、$y$ 是兄弟」的碼樹都成立：把那對兄弟收回成一個葉節點
$z$，$B$ 少 $x.\mathit{freq} + y.\mathit{freq}$。

假設 $T$ 不是最佳的。 由引理~\ref{lem:gr-huff-choice}，有一棵最佳碼樹 $T^*$ 讓 $x$、$y$
是兄弟，而且 $B(T^*) < B(T)$。 把 $T^*$ 的 $x$、$y$ 收回成 $z$，得到 $C'$ 的碼樹
$T^{*\prime}$，
\[
B(T^{*\prime}) = B(T^*) - x.\mathit{freq} - y.\mathit{freq}
< B(T) - x.\mathit{freq} - y.\mathit{freq} = B(T'),
\]
和 $T'$ 是最佳的矛盾。
\end{proof}

\begin{theorem}
\label{thm:gr-huffman}
虛擬碼~\ref{alg:gr-huffman} 建出的碼樹是最佳的。
\end{theorem}

\begin{proof}
對符號個數 $n$ 做歸納。 $n = 2$ 時只有一棵嚴格二元樹（根接兩個葉節點），是最佳的。
$n > 2$ 時，演算法第一次合併取出頻率最小的兩個 $x$、$y$，建出 $z$；之後的每一步，和
「對 $C'$ 從頭跑霍夫曼演算法」做的事完全一樣（$z$ 當成一個葉節點）。 由歸納假設，它
對 $C'$ 建出最佳碼樹 $T'$；最後的樹就是 $T'$ 的 $z$ 展開成 $x$、$y$，由
引理~\ref{lem:gr-huff-sub}，它是 $C$ 的最佳碼樹。
\end{proof}

證明裡沒有用到平手規則：不管平手時怎麼選，取出的都是「頻率最小的兩個」，兩個引理照樣
成立。 這就是上面作答提醒說的，不同規則建出的樹都是最佳的。

\subsection{C 程式}

\begin{cbookrecap}[第~14~章結構、第~15~章遞迴、第~19~章 19.3 節]
\texttt{struct Node} 裡有兩個指向 \texttt{struct Node} 的指標，是建二元樹的寫法。 這裡一次要
一個 $2n - 1$ 格的陣列 \texttt{pool}，不替每個節點各 \texttt{malloc} 一次，前 $n$ 格
放葉節點、後 $n - 1$ 格放內部節點，節點之間仍用指標連起來，用完 \texttt{free} 一次就好。
\texttt{collect} 是遞迴的前序走訪。
\end{cbookrecap}

堆積照第~8~章的陣列堆積寫，改成最小堆積，存的是節點的指標，比較時用
\texttt{before}：先比頻率，再比編號。 碼存在一個 $n \times n$ 的字元陣列，第 $i$ 列
是編號 $i$ 的葉節點的碼；碼長最多 $n - 1$，加上結尾的 \verb|'\0'| 剛好 $n$ 格。 印的
時候照輸入的順序，結果固定。 輸入是符號個數 $n \ge 2$，以及每個符號和它的頻率。

\begin{example}[label=ex:gr-huffman]{c}{(huffman.c) 用最小堆積建霍夫曼樹，照輸入順序印出每個碼}
#include <stdio.h>
#include <stdlib.h>
#include <string.h>

typedef struct Node {
    long long w;                // 頻率（內部節點是兩個孩子的和）
    int id;                     // 葉節點 0..n-1 照輸入順序，內部節點照建立順序
    char sym;                   // 只有葉節點用到
    struct Node *left, *right;  // 葉節點兩個都是 NULL
} Node;

// 平手規則：先比頻率，頻率相同比 id。 id 都不同，所以建出的樹是唯一的。
static int before(const Node *x, const Node *y)
{
    return x->w < y->w || (x->w == y->w && x->id < y->id);
}

// 最小堆積 h[0..*size-1] 存節點的指標。 上浮、下沉同第 8 章，
// 只是把要放的節點先拿在手上，一路把父（子）節點挪過來，最後才放進去。
void push(Node **h, int *size, Node *x)
{
    int i = (*size)++;
    while (i > 0 && before(x, h[(i - 1) / 2])) {
        h[i] = h[(i - 1) / 2];
        i = (i - 1) / 2;
    }
    h[i] = x;
}

Node *pop(Node **h, int *size)          // 呼叫前要確定堆積不是空的
{
    Node *top = h[0], *x = h[--(*size)];
    int i = 0, j;
    while ((j = 2 * i + 1) < *size) {
        if (j + 1 < *size && before(h[j + 1], h[j])) {
            j++;                        // 取較小的孩子
        }
        if (!before(h[j], x)) {
            break;
        }
        h[i] = h[j];
        i = j;
    }
    h[i] = x;
    return top;
}

// 前序走訪，buf[0..depth-1] 是從根到 t 的路徑；到了葉節點就抄進 code 的第 id 列。
void collect(const Node *t, char *buf, int depth, char *code, int n)
{
    if (t->left == NULL) {
        buf[depth] = '\0';
        memcpy(code + (size_t)t->id * n, buf, depth + 1);
        return;
    }
    buf[depth] = '0';
    collect(t->left, buf, depth + 1, code, n);
    buf[depth] = '1';
    collect(t->right, buf, depth + 1, code, n);
}

int main(void)
{
    int n, size = 0;
    if (scanf("%d", &n) != 1 || n < 2) {
        return 1;
    }
    Node *pool = malloc((2 * n - 1) * sizeof *pool);   // n 個葉、n - 1 個內部節點
    Node **heap = malloc(n * sizeof *heap);
    char *code = malloc((size_t)n * n);    // 碼最長 n - 1 位，加 '\0' 共 n 格
    char *buf = malloc(n);
    int ok = pool != NULL && heap != NULL && code != NULL && buf != NULL;
    for (int i = 0; ok && i < n; i++) {
        ok = scanf(" %c %lld", &pool[i].sym, &pool[i].w) == 2 && pool[i].w > 0;
        pool[i].id = i;
        pool[i].left = pool[i].right = NULL;
        if (ok) {
            push(heap, &size, &pool[i]);
        }
    }
    if (ok) {
        for (int k = n; k < 2 * n - 1; k++) {           // n - 1 次合併
            Node *z = &pool[k];
            z->left = pop(heap, &size);                 // 先取出的當左孩子（邊標 0）
            z->right = pop(heap, &size);
            z->w = z->left->w + z->right->w;
            z->id = k;
            push(heap, &size, z);
            printf("merge %d: %lld + %lld = %lld\n",
                   k - n + 1, z->left->w, z->right->w, z->w);
        }
        Node *root = pop(heap, &size);
        collect(root, buf, 0, code, n);
        long long total = 0;
        for (int i = 0; i < n; i++) {                   // 照輸入順序印，結果固定
            const char *c = code + (size_t)i * n;
            printf("%c %lld %s\n", pool[i].sym, pool[i].w, c);
            total += pool[i].w * (long long)strlen(c);
        }
        printf("total bits = %lld\n", total);
    }
    free(pool);
    free(heap);
    free(code);
    free(buf);
    return ok ? 0 : 1;
}
\end{example}
\iosample{7 \\ S 34 \\ C 22 \\ O 15 \\ R 13 \\ T 8 \\ F 5 \\ H 3}{merge 1: 3 + 5 = 8 \\ merge 2: 8 + 8 = 16 \\ merge 3: 13 + 15 = 28 \\ merge 4: 16 + 22 = 38 \\ merge 5: 28 + 34 = 62 \\ merge 6: 38 + 62 = 100 \\ S 34 11 \\ C 22 01 \\ O 15 101 \\ R 13 100 \\ T 8 000 \\ F 5 0011 \\ H 3 0010 \\ total bits = 252}

六次合併和圖 (b)--(g) 一致，每個碼和上面的表相同。 堆積最多放 $n$ 個指標（每次合併
取出兩個、放回一個），所以 \texttt{heap} 只要 $n$ 格。 只有一個符號時碼樹只有根，碼長
是 $0$，沒辦法編碼，所以程式要求 $n \ge 2$。

頻率已經由小到大排好時可以不用堆積：葉節點放一個佇列，新建的內部節點依序放第二個佇
列（它們的權重自動遞增），每次比較兩個佇列的前端，總共 $O(n)$。

\subsection{最佳合併樹}
\label{sec:gr-merge}

第~15~章 15.3 節的外部排序先產生一些排好的串，再一次合併 $k$ 個，直到剩一個。 那裡
假設每個串一樣長，只需要算趟數；用取代選擇產生的串長短不一，「先合併哪幾個」就會影響
總成本。 第~15~章最後留了一句：這和霍夫曼編碼是同一個問題。 這一小節把它說清楚。

先看 $2$ 路合併。 有 $r$ 個排好的串，長度（紀錄筆數）是 $\ell_1, \ldots, \ell_r$。 合
併兩個長度 $a$、$b$ 的串，每筆紀錄讀一次、寫一次，成本算 $a + b$。 每次挑兩個串合併，
$r - 1$ 次之後剩一個。 整個合併過程可以畫成一棵\textbf{合併樹}：每個初始串是一個葉節點，
每次合併是一個內部節點，權重是合併出來的串長；這是一棵有 $r$ 個葉節點的嚴格二元樹。
總成本是所有內部節點的權重和。

\begin{lemma}
\label{lem:gr-merge-cost}
合併樹 $T$ 的總合併成本等於 $\sum_{i=1}^{r} \ell_i \cdot d_T(i)$，$d_T(i)$ 是第 $i$
個串的葉節點的深度（邊數）。
\end{lemma}

\begin{proof}
總成本是每次合併搬動的紀錄數加起來。 換個方向數：第 $i$ 個串的每一筆紀錄，在它的每
一個祖先（內部節點）那一次合併各被搬一次，祖先的個數就是它的深度 $d_T(i)$。 所以總成
本 $= \sum_i \ell_i \, d_T(i)$。
\end{proof}

右邊這個和，把 $\ell_i$ 看成頻率，正是上一小節的 $B(T)$，也叫這棵樹的\textbf{加權
外部路徑長}（weighted external path length）：外部節點（葉節點）的深度乘上權重的和。
所以「怎麼兩兩合併最省」和「怎麼建最佳碼樹」是同一個問題，答案也一樣：\textbf{每次合
併目前最短的兩個串}。 這樣合併出來的樹叫\textbf{最佳合併樹}（optimal merge tree）。

\textbf{例子。} 取代選擇產生了七個串，長度（單位：千筆）如下：

\begin{center}
\begin{tabular}{c*{7}{c}}
\toprule
串 & R1 & R2 & R3 & R4 & R5 & R6 & R7 \\
\midrule
長度 & 12 & 30 & 8 & 25 & 5 & 18 & 40 \\
\bottomrule
\end{tabular}
\end{center}

每次合併最短的兩個（一樣長時初始串先取，和霍夫曼演算法相同）：$5 + 8 = 13$，
$12 + 13 = 25$，$18 + 25 = 43$（R6 和 R4；兩個 $25$ 平手，先取初始串 R4），
$25 + 30 = 55$，$40 + 43 = 83$，$55 + 83 = 138$。 總成本 $357$ 千筆。 畫成合併樹：

\begin{center}
\begin{minipage}{\linewidth}
\centering
\begin{tikzpicture}
\node[gr-leaf] (mt0) at (0.000,-2.700) {R1\\12};
\node[gr-leaf] (mt1) at (2.850,-1.800) {R2\\30};
\node[gr-leaf] (mt2) at (1.900,-3.600) {R3\\8};
\node[gr-leaf] (mt3) at (5.700,-2.700) {R4\\25};
\node[gr-leaf] (mt4) at (0.950,-3.600) {R5\\5};
\node[gr-leaf] (mt5) at (4.750,-2.700) {R6\\18};
\node[gr-leaf] (mt6) at (3.800,-1.800) {R7\\40};
\node[gr-inner] (mt7) at (1.425,-2.700) {13};
\node[gr-inner] (mt8) at (0.712,-1.800) {25};
\node[gr-inner] (mt9) at (5.225,-1.800) {43};
\node[gr-inner] (mt10) at (1.781,-0.900) {55};
\node[gr-inner] (mt11) at (4.512,-0.900) {83};
\node[gr-inner] (mt12) at (3.147,0.000) {138};
\draw[gr-edge] (mt7) -- (mt4);
\draw[gr-edge] (mt7) -- (mt2);
\draw[gr-edge] (mt8) -- (mt0);
\draw[gr-edge] (mt8) -- (mt7);
\draw[gr-edge] (mt10) -- (mt8);
\draw[gr-edge] (mt10) -- (mt1);
\draw[gr-edge] (mt9) -- (mt5);
\draw[gr-edge] (mt9) -- (mt3);
\draw[gr-edge] (mt11) -- (mt6);
\draw[gr-edge] (mt11) -- (mt9);
\draw[gr-edge] (mt12) -- (mt10);
\draw[gr-edge] (mt12) -- (mt11);
\node[idxlbl, anchor=east] at (-0.6,0.000) {深度 0};
\node[idxlbl, anchor=east] at (-0.6,-0.900) {深度 1};
\node[idxlbl, anchor=east] at (-0.6,-1.800) {深度 2};
\node[idxlbl, anchor=east] at (-0.6,-2.700) {深度 3};
\node[idxlbl, anchor=east] at (-0.6,-3.600) {深度 4};
\end{tikzpicture}
\diagcaption{七個串的最佳合併樹。 方框是初始串，圓圈是一次合併，數字是合併出來的串
長。 總成本是圓圈的數字和 $13 + 25 + 43 + 55 + 83 + 138 = 357$，也等於
$\sum \ell_i d_i = 5 \cdot 4 + 8 \cdot 4 + 12 \cdot 3 + 18 \cdot 3 + 25 \cdot 3 +
30 \cdot 2 + 40 \cdot 2 = 357$。}
\label{fig:gr-merge-tree}
\end{minipage}
\end{center}

和第~15~章的做法比較：一趟一趟地把相鄰的兩個串兩兩合併（第 1 趟 R1+R2、R3+R4、
R5+R6，R7 輪空），三趟的成本是 $98 + 138 + 138 = 374$。 那種做法讓每個串的深度差不多，
串一樣長時很好；長短不一時，應該讓長的串少搬幾次。

\begin{examnote}
最佳合併的計算題，先看題目怎麼定義一次合併的成本。 本書用「搬動的紀錄數」$a + b$；
有的題目用最壞的比較次數 $a + b - 1$，$r$ 個串要合併 $r - 1$ 次，總成本差一個常數
$r - 1$，最佳的合併順序不變，但數字要照題目算。 問的是區塊 I/O 時，每次合併讀和寫各
$(a + b)/B$ 個區塊。 用公式 $\sum \ell_i d_i$ 驗算時，$d_i$ 是邊數（深度）；如果用
「根在第 $1$ 層」的層數代入，每一項都要減 $1$，否則多算了 $\sum \ell_i$。
\end{examnote}

\textbf{$k$ 路合併。} 每次合併 $k$ 個串，合併樹的每個內部節點都有 $k$ 個孩子。 引
理~\ref{lem:gr-merge-cost} 的證明完全不用改：總成本仍是 $\sum \ell_i d_i$。 直覺上
應該每次合併最短的 $k$ 個，第~15~章也是這麼說的，但這裡有一個細節。 每次合併讓串的
個數少 $k - 1$，從 $r$ 個變成 $1$ 個，如果每次都恰好合併 $k$ 個，$r - 1$ 必須是
$k - 1$ 的倍數。 不是倍數的時候，一定有某一次合併的串少於 $k$ 個，而這一次放在哪裡會影
響成本。

最好的做法是讓這次「不滿」的合併發生在最底下，處理最短的那幾個串。 實作上先補
\[
d = \bigl(k - 1 - (r - 1) \bmod (k - 1)\bigr) \bmod (k - 1)
\]
個長度 $0$ 的\textbf{虛擬串}，讓 $r + d - 1$ 是 $k - 1$ 的倍數，再每次合併最短的 $k$
個。 虛擬串長度是 $0$，搬它不花成本，它們一定在第一次合併就被取出，效果就是第一次只合
併最短的 $k - d$ 個真正的串。 這個做法的正確性和霍夫曼演算法的證明一樣，只是引
理~\ref{lem:gr-huff-choice} 改成「最小的 $k$ 個是同一個節點的孩子，在最深的一層」；
補虛擬串正是為了讓最深的那一層剛好有 $k$ 個兄弟。

補完之後串的個數 $r + d \equiv 1 \pmod{k - 1}$，每次合併少 $k - 1$ 個，所以每次都取得到
$k$ 個，最後剛好剩一個。 $k = 2$ 時 $d = 0$，就是霍夫曼演算法。

\textbf{例子。} 拿掉上面的 R7，剩下六個串 $12, 30, 8, 25, 5, 18$，用 $3$ 路合併。
$r = 6$，$(r - 1) \bmod 2 = 1$，補 $d = 1$ 個虛擬串。

\begin{center}
\begin{tabular}{lll}
\toprule
 & 補虛擬串 & 不補，直接每次取最短的三個 \\
\midrule
第 1 次 & $0 + 5 + 8 = 13$ & $5 + 8 + 12 = 25$ \\
第 2 次 & $12 + 13 + 18 = 43$ & $18 + 25 + 25 = 68$ \\
第 3 次 & $25 + 30 + 43 = 98$ & $30 + 68 = 98$（只剩兩個） \\
\midrule
總成本 & $154$ & $191$ \\
\bottomrule
\end{tabular}
\end{center}

不補虛擬串時，前兩次都合併了三個，最後只剩兩個串可以合併，「不滿」的那一次落在根。 和
補虛擬串的樹比較，R4 和 R1 的深度各多 $1$，多搬了 $25 + 12 = 37$ 千筆，正好是
$191 - 154$。 補了虛擬串，「不滿」的那一次在最底下，只牽涉最短的 $5$ 和 $8$。 補虛擬串的樹裡，R4、R2 深度 $1$，R1、R6 深度 $2$，R5、R3
深度 $3$，$\sum \ell_i d_i = 55 + 60 + 39 = 154$。

同樣的道理用在 $k$ 元霍夫曼碼（用 $0, 1, \ldots, k - 1$ 這 $k$ 個符號編碼）：先補頻率
$0$ 的虛擬符號，再每次合併頻率最小的 $k$ 個。

\section{排程}
\label{sec:gr-sched}

\subsection{最小化平均完成時間}
\label{sec:gr-spt}

實驗室只有一台 GPU 伺服器，一次只能跑一個訓練工作，一個工作開始了就要跑完。 早上
$9$ 點五位同學同時送來工作，各要跑的小時數是：

\begin{center}
\begin{tabular}{c*{5}{c}}
\toprule
工作 & A & B & C & D & E \\
\midrule
執行時間 $p_i$（小時） & 7 & 2 & 5 & 1 & 4 \\
\bottomrule
\end{tabular}
\end{center}

工作 $i$ 的\textbf{完成時間} $C_i$ 是它跑完的時刻（從 $9$ 點起算）。 每位同學都在等
自己的結果，我們想讓平均完成時間 $\frac{1}{n}\sum C_i$ 最小，也就是讓 $\sum C_i$ 最
小。 第~6~章的印表機佇列最後留了一個問題：先來先服務不一定讓平均等待最短，這一小節回
答它。

照送來的順序 A、B、C、D、E 跑：完成時間 $7, 9, 14, 15, 19$，總和 $64$，平均
$12.8$ 小時。 把短的排前面，D、B、E、C、A：完成時間 $1, 3, 7, 12, 19$，總和 $42$，平
均 $8.4$ 小時。 最後一個工作不管怎麼排都在 $19$ 點完成，差別在前面的工作：排在第 $k$
個的工作，它的執行時間會加進它自己和後面每一個工作的完成時間裡，被算 $n - k + 1$ 次。
寫成式子，順序是 $\pi(1), \pi(2), \ldots, \pi(n)$ 時
\[
\sum_{i=1}^{n} C_i = \sum_{k=1}^{n} (n - k + 1)\, p_{\pi(k)}.
\]
D、B、E、C、A 的順序代進去是 $5 \cdot 1 + 4 \cdot 2 + 3 \cdot 4 + 2 \cdot 5 + 1 \cdot 7 =
42$。 被算最多次的位置放最短的工作，這個貪心叫\textbf{最短工作優先}（shortest
processing time first，SPT）。

\begin{theorem}
\label{thm:gr-spt}
依執行時間由小到大排的順序，$\sum C_i$ 最小。 執行時間一樣的工作之間怎麼排，總和都
一樣。
\end{theorem}

\begin{proof}
任取一個順序。 如果它不是由小到大，一定有\textbf{相鄰}的兩個工作 $i$、$j$，$i$ 排在
$j$ 前面，$p_i > p_j$。 設它們前面的工作在時刻 $t$ 做完。 交換 $i$ 和 $j$ 之前，
$C_i = t + p_i$，$C_j = t + p_i + p_j$；交換之後，$C_j' = t + p_j$，$C_i' = t + p_j + p_i$。
其他工作的完成時間都不變（它們前面的工作總時數沒變）。 所以
\[
(C_i + C_j) - (C_i' + C_j') = p_i - p_j > 0,
\]
交換讓總和嚴格變小。 因此最佳的順序不能有這種相鄰的一對，也就是一定由小到大排好。
執行時間相同的工作交換，總和不變，所以所有由小到大的順序總和都一樣，都是最佳的。
\end{proof}

這個交換論證的寫法和活動選擇不太一樣：它說明任何不是由小到大的順序都還能再改進，沒有
去改造某一個最佳解。 兩種寫法都常見。

\textbf{等待時間。} 工作 $i$ 的等待時間是 $W_i = C_i - p_i$（開始之前等了多久），
$\sum W_i = \sum C_i - \sum p_i$，後面一項和順序無關，所以讓總等待時間最小的也是 SPT。
章首例子的 SPT 順序，總等待時間是 $42 - 19 = 23$ 小時。 演算法就是排序，
$\Theta(n \lg n)$。

工作有輕重之分（例如快要口試的同學比較急）時，每個工作有權重 $w_i$，要讓 $\sum w_i
C_i$ 最小。 同樣的相鄰交換論證可以算出應該依 $p_i / w_i$ 由小到大排，見觀念題
第~\ref{q:gr-wspt}~題。

\subsection{有期限的單位時間工作}
\label{sec:gr-deadline}

系上的雷射切割機開放預約，每個工作剛好用 $1$ 小時。 今天收到 $n$ 件申請，第 $i$ 件
的\textbf{期限}是 $d_i$（正整數）：要在第 $d_i$ 小時結束前做完，申請人才付 $p_i$ 元
（$p_i > 0$）；趕不上就不收費，乾脆不做。 時段編號 $1, 2, 3, \ldots$，時段 $t$ 是第
$t$ 個小時。 工作放在時段 $t$ 能趕上期限，條件是 $t \le d_i$。 要挑一些工作、各放進
不同的時段，讓收到的錢最多。

\begin{center}
\begin{tabular}{c*{9}{c}}
\toprule
工作 & 1 & 2 & 3 & 4 & 5 & 6 & 7 & 8 & 9 \\
\midrule
期限 $d_i$ & 3 & 1 & 4 & 5 & 4 & 1 & 6 & 3 & 2 \\
收費 $p_i$ & 45 & 20 & 60 & 35 & 50 & 25 & 15 & 40 & 30 \\
\bottomrule
\end{tabular}
\end{center}

\textbf{貪心的做法。} 依收費由高到低看每件工作（一樣時編號小的先），如果還找得到一個
不超過期限的空時段，就接下來，放進\textbf{最晚}的那個空時段；找不到就不接。 放最晚的
空時段，是為了把前面的時段留給期限比較早的工作。

期限超過 $n$ 的工作和期限是 $n$ 的沒有差別（$n$ 件工作用不到第 $n$ 個時段以後），可以
先把期限截到 $n$。 最直接的寫法是從期限往回一格一格找空時段，最壞每件工作走 $n$ 步，
總共 $O(n^2)$；第~\ref{sec:gr-dsu}~小節用互斥集合把它降到幾乎線性。

依收費排好是 $3, 5, 1, 8, 4, 9, 6, 2, 7$。 照這個做法，時段 $1$ 到 $6$ 依序放工作
$8, 1, 5, 3, 4, 7$，工作 $9, 6, 2$ 找不到位置，收費 $245$ 元；過程見下一小節的追蹤表。

\subsection{為什麼這樣排最多}

先把「一群工作排得進去」說清楚。 一群工作 $A$ 叫做\textbf{可行}的，是指能把它們放進
不同的時段，每件都不超過自己的期限。 令 $N_t(A)$ 是 $A$ 裡期限 $\le t$ 的工作件數。

\begin{lemma}
\label{lem:gr-feasible}
$A$ 是可行的，若且唯若對每個 $t \ge 1$，$N_t(A) \le t$。
\end{lemma}

\begin{proof}
必要：期限 $\le t$ 的工作都要放在時段 $1, \ldots, t$，而且各佔一個，所以最多 $t$
件。 充分：把 $A$ 依期限由小到大排，第 $k$ 件放時段 $k$。 如果第 $k$ 件的期限
$d < k$，那麼前 $k$ 件的期限都 $\le d$，$N_d(A) \ge k > d$，和條件矛盾。 所以每件都
趕得上。
\end{proof}

\begin{whybox}[貪心是最佳的]
證明分兩部分。

\textbf{(1) 找得到空時段，若且唯若接下來仍然可行。} 設目前接下的工作集合是 $A$（都已
經放好），輪到工作 $j$。 找得到空時段的話，直接得到 $A \cup \{j\}$ 的排法。 反過來，
假設時段 $1, \ldots, d_j$ 全滿了。 令 $u$ 是比 $d_j$ 大的第一個空時段（時段 $n + 1$
以後永遠是空的，所以 $u$ 存在；而且 $u \le n$，否則時段 $1, \ldots, n$ 全滿，$n$ 件工作都已
接下，不會輪到 $j$）。 時段 $1, \ldots, u - 1$ 全滿。 說明這些時段裡每件
工作的期限都 $\le u - 1$：假設某件工作 $q$ 放在時段 $s < u$，期限 $d_q \ge u$。 時段
不會被清空，所以 $q$ 被放進去的時候，時段 $u$ 是空的，而且 $u \le d_q$；演算法放的是
最晚的空時段，應該放在 $u$ 或更晚，不會是 $s$，矛盾。 所以 $A \cup \{j\}$ 裡期限
$\le u - 1$ 的工作至少有 $(u - 1) + 1$ 件（$d_j \le u - 1$），由
引理~\ref{lem:gr-feasible}，$A \cup \{j\}$ 不可行。

\textbf{(2) 交換論證。} 證明下面的不變量：每看完一件工作，已接下的集合 $A$ 都包含於
某個最佳解 $O$（收費最多的可行集合）。 一開始 $A$ 是空的，成立。 輪到工作 $j$ 時：

\begin{itemize}
\item 不接 $j$：由 (1)，$A \cup \{j\}$ 不可行；$O \supseteq A$，所以 $O$ 也不含 $j$
（可行集合的子集合還是可行的），不變量仍然成立。
\item 接下 $j$，而且 $j \in O$：不變量成立。
\item 接下 $j$，但 $j \notin O$：先說明 $O \setminus A$ 裡的工作都還沒輪到。 如果某件
$q \in O \setminus A$ 已經輪過、被拒絕了，那時接下的集合 $A_q \subseteq A \subseteq O$，
$A_q \cup \{q\} \subseteq O$ 是可行的，由 (1) 不會被拒絕，矛盾。 所以它們的收費都
$\le p_j$。 接著，$O \cup \{j\}$ 收費比 $O$ 多，不可能可行；令 $t^*$ 是讓
$N_t(O \cup \{j\}) > t$ 的最小 $t$。 $O$ 可行而多了一件，所以
$N_{t^*}(O \cup \{j\}) = t^* + 1$；$A \cup \{j\}$ 可行，所以 $N_{t^*}(A \cup \{j\})
\le t^*$。 兩者相差的工作都在 $O \setminus A$ 裡，因此有某件 $k \in O \setminus A$，
$d_k \le t^*$。 令 $O' = O \cup \{j\} \setminus \{k\}$。 對 $t < t^*$，
$N_t(O') \le N_t(O \cup \{j\}) \le t$；對 $t \ge t^*$，$k$ 被算在 $N_t$ 裡又被拿掉，
$N_t(O') = N_t(O \cup \{j\}) - 1 \le N_t(O) \le t$。 所以 $O'$ 可行，收費
$\ge$ $O$ 的收費（$p_k \le p_j$），也是最佳解，而且包含 $A \cup \{j\}$。
\end{itemize}

全部看完時，$A \subseteq O$ 而且 $O \setminus A$ 裡沒有工作（都輪過了，由上面的論證
不會被拒絕），所以 $A = O$，貪心的結果是最佳解。
\end{whybox}

\subsection{用互斥集合找最晚的空時段}
\label{sec:gr-dsu}

一格一格往回找太慢。 把時段分組：每一組
是一段連續的時段 $\{e, e + 1, \ldots, e + m\}$，$e$ 是空的（或是 $0$），後面 $m$ 個都
滿了。 時段 $0$ 是哨兵，代表「沒有空時段了」，它永遠不會被填滿。 這樣分組之後，「不
超過 $t$ 的最晚空時段」就是 $t$ 所在那一組的 $e$。 這正是第~13~章的互斥集合：

\begin{itemize}
\item 一開始每個時段 $0, 1, \ldots, D$ 自己一組（\proc{Make-Set}），$D = \min(n,
\max_i d_i)$。
\item 查詢：\proc{Find-Set}$(t)$ 找到 $t$ 那一組的代表，代表上記著這一組的空時段。
\item 時段 $e$ 被填滿以後，它那一組要併進左邊那一組（含 $e - 1$ 的那組）：
\proc{Union}$(e, e - 1)$。 合併後的空時段是左邊那一組原本的空時段。
\end{itemize}

第~13~章的 \proc{Union} 依大小合併，留下來的根不一定是空時段，所以另外用一個陣列
$\mathit{latest}$，只在根上有意義：$\mathit{latest}[r]$ 是根 $r$ 那一組的空時段。

\begin{pseudo}[label=alg:gr-deadline-dsu]{有期限的工作排程，用互斥集合（工作已依 $p$ 由大到小排好）}
\Function{Schedule-Jobs}{$d, p, n$}
    \State $D \gets \min(n, \max_{1 \le i \le n} d[i])$
    \For{$t \gets 0$ \To $D$}
        \State \Call{Make-Set}{$t$}
        \State $\mathit{latest}[t] \gets t$
        \State $\mathit{slot}[t] \gets \NIL$
    \EndFor
    \For{$i \gets 1$ \To $n$}
        \State $e \gets \mathit{latest}[\Call{Find-Set}{\min(d[i], D)}]$
        \If{$e > 0$} \Comment{$e = 0$ 表示沒有空時段，不接}
            \State $\mathit{slot}[e] \gets i$
            \State $\ell \gets \mathit{latest}[\Call{Find-Set}{e - 1}]$ \Comment{左邊那一組的空時段}
            \State \Call{Union}{$e, e - 1$}
            \State $\mathit{latest}[\Call{Find-Set}{e}] \gets \ell$
        \EndIf
    \EndFor
    \State \Return $\mathit{slot}$
\EndFunction
\end{pseudo}

第 11 行要在第 12 行的 \proc{Union} \textbf{之前}讀出左邊那一組的空時段；合併以後根可能換成右邊
那一組的根，就讀不到了。 每件工作做三次 \proc{Find-Set} 和一次 \proc{Union}，用依
大小合併加路徑壓縮，由第~13~章，總時間 $O(n\,\alpha(n))$，整個演算法被排序的
$O(n \lg n)$ 主導。

\textbf{追蹤。} 同樣的九件工作，$D = \min(9, 6) = 6$，時段 $0$ 到 $6$。 \proc{Union}
照第~13~章的規則：依大小合併，一樣大時第二個參數的根接到第一個參數的根底下；
\proc{Find-Set} 做路徑壓縮。 「$p[0..6]$」是做完這件工作之後的父節點陣列，
「$\mathit{latest}$」只列出各個根上的值（根:空時段）。

\begin{center}
\small
\begin{tabular}{cccccl}
\toprule
工作 & $\min(d, D)$ & 根 & 空時段 $e$ & $p[0..6]$ & $\mathit{latest}$（根:空時段） \\
\midrule
開始 & & & & 0 1 2 3 4 5 6 & 每個時段自己 \\
3 & 4 & 4 & 4 & 0 1 2 4 4 5 6 & 0:0, 1:1, 2:2, 4:3, 5:5, 6:6 \\
5 & 4 & 4 & 3 & 0 1 4 4 4 5 6 & 0:0, 1:1, 4:2, 5:5, 6:6 \\
1 & 3 & 4 & 2 & 0 4 4 4 4 5 6 & 0:0, 4:1, 5:5, 6:6 \\
8 & 3 & 4 & 1 & 4 4 4 4 4 5 6 & 4:0, 5:5, 6:6 \\
4 & 5 & 5 & 5 & 4 4 4 4 4 4 6 & 4:0, 6:6 \\
9 & 2 & 4 & 0，不接 & 不變 & \\
6 & 1 & 4 & 0，不接 & 不變 & \\
2 & 1 & 4 & 0，不接 & 不變 & \\
7 & 6 & 6 & 6 & 4 4 4 4 4 4 4 & 4:0 \\
\bottomrule
\end{tabular}
\end{center}

第一件工作 3 填了時段 $4$，\proc{Union}$(4, 3)$ 兩組一樣大，$3$ 接到 $4$ 底下，根是
$4$；但這一組的空時段是 $3$，所以 $\mathit{latest}[4] = 3$。 這就是為什麼不能直接拿根
當空時段。 工作 4 填了時段 $5$ 之後，\proc{Union}$(5, 4)$：$5$ 那一組只有 $1$ 個時段，
$4$ 那一組有 $5$ 個，依大小合併讓 $5$ 接到 $4$ 底下，$\mathit{latest}[4]$ 仍是 $0$。
從此期限 $\le 5$ 的工作一律一次 \proc{Find-Set} 就知道沒有空位。

另一種常見的寫法不依大小合併，填滿時段 $e$ 時直接設 $p[e] \gets \proc{Find-Set}(e - 1)$，
根永遠是空時段，不需要 $\mathit{latest}$；只靠路徑壓縮，總時間 $O(n \lg n)$。

\subsection{C 程式}

\begin{cbookrecap}[第~14~章結構、第~19~章 19.3 節]
\texttt{calloc(count, size)} 和 \texttt{malloc} 一樣要記憶體，但會把整塊設成 $0$；
下面的 \texttt{slot} 用 $0$ 表示「這個時段沒有工作」，工作從 $1$ 編號，所以剛好。 互
斥集合的 \texttt{dsu\_init}、\texttt{find\_iter}、\texttt{dsu\_union} 沿用第~13~章
範例 dsu\_trace.c 的寫法，只拿掉這裡用不到的欄位。
\end{cbookrecap}

程式讀進 $n$ 件工作，每件一列：期限和收費。 每接下一件工作，印出它放的時段和這時的
父節點陣列；最後印出時段 $1$ 到 $D$ 各放了哪件工作（\texttt{-} 是空的）和總收費。

\begin{example}[label=ex:gr-deadline]{c}{(deadline.c) 有期限的工作排程：用互斥集合找最晚的空時段}
#include <stdio.h>
#include <stdlib.h>

// ---- 第 13 章的互斥集合：依大小合併、迴圈版路徑壓縮（只留用得到的欄位）----
typedef struct {
    int *parent;    // parent[x] == x 表示 x 是根
    int *size;      // 只有根的 size 有意義
} DSU;

int dsu_init(DSU *d, int n)
{
    d->parent = malloc(n * sizeof *d->parent);
    d->size = malloc(n * sizeof *d->size);
    if (d->parent == NULL || d->size == NULL) {
        return 0;
    }
    for (int x = 0; x < n; x++) {
        d->parent[x] = x;
        d->size[x] = 1;
    }
    return 1;
}

int find_iter(int *parent, int x)
{
    int root = x;
    while (parent[root] != root) {
        root = parent[root];
    }
    while (parent[x] != root) {
        int next = parent[x];
        parent[x] = root;
        x = next;
    }
    return root;
}

void dsu_union(DSU *d, int x, int y)    // 一樣大時 y 的根接到 x 的根底下
{
    int rx = find_iter(d->parent, x), ry = find_iter(d->parent, y);
    if (rx == ry) {
        return;
    }
    if (d->size[rx] < d->size[ry]) {
        int t = rx;
        rx = ry;
        ry = t;
    }
    d->parent[ry] = rx;
    d->size[rx] += d->size[ry];
}

// ---- 有期限的單位時間工作 ----
typedef struct {
    int id, d, p;   // 編號、期限、收費
} Job;

int by_profit(const void *a, const void *b)    // 收費高的先；一樣時編號小的先
{
    const Job *x = a, *y = b;
    if (x->p != y->p) {
        return x->p > y->p ? -1 : 1;
    }
    return (x->id > y->id) - (x->id < y->id);
}

int main(void)
{
    int n, D = 0;
    if (scanf("%d", &n) != 1 || n <= 0) {
        return 1;
    }
    Job *job = malloc(n * sizeof *job);
    int ok = job != NULL;
    for (int i = 0; ok && i < n; i++) {
        ok = scanf("%d %d", &job[i].d, &job[i].p) == 2 && job[i].d >= 1;
        job[i].id = i + 1;
        if (ok && job[i].d > D) {
            D = job[i].d;
        }
    }
    D = D < n ? D : n;      // 時段 1..D；時段 0 是「沒有空位」的哨兵
    DSU ds = {NULL, NULL};
    int *latest = malloc((D + 1) * sizeof *latest);  // latest[根]：這組的空時段
    int *slot = calloc(D + 1, sizeof *slot);         // slot[t]：放在時段 t 的工作
    ok = ok && latest != NULL && slot != NULL && dsu_init(&ds, D + 1);
    if (ok) {
        for (int t = 0; t <= D; t++) {
            latest[t] = t;
        }
        qsort(job, n, sizeof job[0], by_profit);
        long long total = 0;
        for (int i = 0; i < n; i++) {
            int t = job[i].d < D ? job[i].d : D;
            int e = latest[find_iter(ds.parent, t)];
            printf("job %d (d=%d, p=%d): ", job[i].id, job[i].d, job[i].p);
            if (e == 0) {
                printf("rejected\n");
                continue;
            }
            slot[e] = job[i].id;
            total += job[i].p;
            int left = latest[find_iter(ds.parent, e - 1)];  // 合併前先讀
            dsu_union(&ds, e, e - 1);
            latest[find_iter(ds.parent, e)] = left;
            printf("slot %d, parent =", e);
            for (int x = 0; x <= D; x++) {
                printf(" %d", ds.parent[x]);
            }
            printf("\n");
        }
        printf("schedule:");
        for (int t = 1; t <= D; t++) {
            printf(slot[t] ? " %d" : " -", slot[t]);
        }
        printf("\ntotal profit = %lld\n", total);
    }
    free(job);
    free(latest);
    free(slot);
    free(ds.parent);
    free(ds.size);
    return ok ? 0 : 1;
}
\end{example}
\iosample{9 \\ 3 45 \\ 1 20 \\ 4 60 \\ 5 35 \\ 4 50 \\ 1 25 \\ 6 15 \\ 3 40 \\ 2 30}{job 3 (d=4, p=60): slot 4, parent = 0 1 2 4 4 5 6 \\ job 5 (d=4, p=50): slot 3, parent = 0 1 4 4 4 5 6 \\ job 1 (d=3, p=45): slot 2, parent = 0 4 4 4 4 5 6 \\ job 8 (d=3, p=40): slot 1, parent = 4 4 4 4 4 5 6 \\ job 4 (d=5, p=35): slot 5, parent = 4 4 4 4 4 4 6 \\ job 9 (d=2, p=30): rejected \\ job 6 (d=1, p=25): rejected \\ job 2 (d=1, p=20): rejected \\ job 7 (d=6, p=15): slot 6, parent = 4 4 4 4 4 4 4 \\ schedule: 8 1 5 3 4 7 \\ total profit = 245}

每一列的父節點陣列都和手算的表一樣。 程式的元素直接用時段編號 $0..D$，不必像
第~13~章那樣在輸入輸出時加減一。 \texttt{slot} 和 \texttt{latest} 都要 $D + 1$ 格，
因為時段 $0$ 也是一個元素。

\section{貪心何時不對}
\label{sec:gr-fail}

前面每一個貪心都附了證明。 這一節反過來，看幾個「看起來很合理、其實是錯的」貪心，最
後介紹把貪心整理成一般理論的擬陣。

\subsection{找零錢}

收銀機有幾種面額的硬幣 $c_1 < c_2 < \cdots < c_m$，每種都用不完，要湊出金額 $x$，硬
幣枚數越少越好。 直覺的貪心是：能用大的就用大的。

\begin{pseudo}[label=alg:gr-change]{找零錢的貪心}
\Function{Greedy-Change}{$c, m, x$}
    \State $\mathit{count} \gets 0$
    \For{$i \gets m$ \DownTo $1$}
        \State $\mathit{count} \gets \mathit{count} + \lfloor x / c[i] \rfloor$
        \State $x \gets x \bmod c[i]$
    \EndFor
    \If{$x > 0$}
        \State \Return \NIL \Comment{湊不出來}
    \EndIf
    \State \Return $\mathit{count}$
\EndFunction
\end{pseudo}

新臺幣常用的硬幣 $1, 5, 10, 50$ 元用這個貪心，每個金額都得到最少枚數（觀念題
第~\ref{q:gr-coin}~題會說明怎麼驗證）。 這樣的面額組合叫\textbf{標準的}（canonical）。
換一組面額就不一定了：

\begin{itemize}
\item 面額 $1, 5, 12$，金額 $15$：貪心先拿 $12$，剩 $3$ 只能用三個 $1$，共 $4$ 枚；
三個 $5$ 只要 $3$ 枚。
\item 面額 $4, 10$（沒有 $1$），金額 $12$：貪心先拿 $10$，剩 $2$ 湊不出來，回傳
\NIL；其實 $4 + 4 + 4$ 就可以。 貪心不只不是最少，連有解都找不到。
\end{itemize}

出錯的原因和 0/1 背包一樣：拿了 $12$ 之後剩下的 $3$ 很難湊，貪心在選 $12$ 的時候看不
到這件事。 一般的面額要用動態規劃：令 $M[a]$ 是湊出 $a$ 的最少枚數，$M[0] = 0$，
$M[a] = 1 + \min_{c_i \le a} M[a - c_i]$，時間 $O(mx)$，是擬多項式的（第~18~章）。

一組面額是不是標準的，也有辦法判斷。 Kozen 與 Zaks（1994）證明：如果貪心對某個金額
不是最少枚數，那麼最小的反例一定小於 $c_{m-1} + c_m$。 所以只要用動態規劃算出這個範
圍內每個金額的最少枚數，逐一和貪心比較，就能確定這組面額是不是標準的。

\subsection{其他常見的錯誤貪心}

\begin{itemize}
\item \textbf{0/1 背包}：第~\ref{sec:gr-knapsack}~節。 不管依單位重量價值、價值還是
重量，都有反例。
\item \textbf{加權的活動選擇}：每筆申請有不同的租金，要讓收入最多。 最早結束的先排會
錯：$[8, 10)$ 租金 $1$、$[9, 12)$ 租金 $10$，它先選了前者，後者就排不進去。 這時要用
動態規劃：依結束時刻排好，第 $i$ 筆「選或不選」，選的話接上最後一筆和它相容的申請。
\item \textbf{期限排程放最早的空時段}：第~\ref{sec:gr-deadline}~節如果把工作放在最早
的空時段：工作 A（期限 $2$、收費 $10$）先放進時段 $1$，工作 B（期限 $1$、收費 $8$）就
沒位置了，只收 $10$；放最晚的空時段可以兩件都接，收 $18$。 同一個演算
法，細節換一個，就不對了。
\item 第~25~章的 Dijkstra 演算法也是貪心，邊的權重有負的時候就會錯。
\end{itemize}

遇到一個新的貪心想法，先用小例子找反例，特別是讓兩個選擇平手或差一點點的例子；找不
到，再試著寫交換論證，寫不下去的地方往往就是反例的方向。 也可以另外寫一個窮舉的版本，
在大量隨機的小輸入上和貪心比對答案。 比對一萬組都相同也不是證明，但只要有一組不同，貪
心就是錯的。

\subsection{擬陣：貪心的一般理論（觀念）}
\label{sec:gr-matroid}

很多貪心都長成同一個樣子：有一群元素，每個有正的權重；有一些子集合叫做「獨立的」
（可以同時選）；依權重由大到小看每個元素，加進去之後仍然獨立就加。 第~\ref{sec:gr-deadline}~節的期限排程正是這樣：元素是工作，權重是收費，「獨立」是「可行」。 這種貪心什麼時候對？
有一個漂亮的答案。

一個有限集合 $S$ 和一群「獨立」的子集合 $\mathcal{I}$，如果滿足下面三條，叫做
\textbf{擬陣}（matroid）：

\begin{enumerate}
\item 空集合是獨立的。
\item \textbf{遺傳性}：獨立集合的子集合也是獨立的。
\item \textbf{交換性}：$A$、$B$ 都是獨立的，$|A| < |B|$，那麼 $B$ 裡一定有某個不在
$A$ 裡的元素 $x$，讓 $A \cup \{x\}$ 仍然獨立。
\end{enumerate}

\textbf{定理}（Rado 1957、Edmonds 1971）：在擬陣上，對任何正的權重，上面的貪心都得到
權重和最大的獨立集合。 反過來，如果 $(S, \mathcal{I})$ 滿足前兩條，而且對每一種權重
貪心都對，它一定也滿足第三條，是擬陣。 這裡不證明。

幾個例子：

\begin{itemize}
\item \textbf{圖的擬陣}：元素是無向圖的邊，「獨立」是「不含環」。 權重由小到大的版本就是
第~24~章的 Kruskal 最小生成樹演算法。
\item \textbf{期限排程}：元素是工作，「獨立」是「可行」；交換性可以用
引理~\ref{lem:gr-feasible} 證明。 第~\ref{sec:gr-deadline}~節的演算法就是擬陣上的貪心，「找
得到空時段」就是檢查加進去之後仍然獨立。
\item \textbf{不是擬陣的例子}：兩兩相容的區間。 $A = \{[1, 10)\}$、$B = \{[1, 2), [3, 4)\}$，
$B$ 的兩個都和 $[1, 10)$ 衝突，交換性不成立。 活動選擇的貪心卻是對的，因為它依的是結束
時刻這個特別的順序，不是權重。 擬陣是貪心成立的充分條件，不是必要條件；霍夫曼編碼也不
是擬陣上的問題。
\end{itemize}

CLRS 第 4 版把擬陣的完整內容移到了書的網站，第 3 版的 16.4、16.5 節有證明和用擬陣
處理期限排程的例子；第 4 版第 15 章的章末附注列了原始文獻。

\section{本章常見錯誤}
\label{sec:gr-pitfalls}

\begin{center}
\begin{tabular}{p{0.24\linewidth}p{0.50\linewidth}p{0.16\linewidth}}
\toprule
錯誤 & 例子 & 見 \\
\midrule
排序的鍵選錯 & 活動選擇依開始時刻或長度排序；以為「衝突最少的先排」一定對 &
第~\ref{sec:gr-activity}~節 \\
端點的相容條件弄錯 & 半開區間寫成 $s_i > \mathit{last}$，少選了首尾相接的申請；閉區
間反而寫成 $\ge$ & 第~\ref{sec:gr-activity}~節 \\
交換論證只換一次 & 證明了第一個選擇對，就說整個貪心對，沒有用最佳子結構或歸納接下去；
或換完沒檢查解還可行 & 第~\ref{sec:gr-elements}~節 \\
分數背包的做法用在 0/1 背包 & 依單位重量價值整件拿，袋子剩下的空間浪費掉 &
第~\ref{sec:gr-knapsack}~節 \\
碼長多算 $1$ & 把碼長寫成葉節點所在的層數；碼長是深度（邊數） &
第~\ref{sec:gr-huffman}~節 \\
$k$ 路合併沒補虛擬串 & $r - 1$ 不是 $k - 1$ 的倍數，直接每次取最短的 $k$ 個 &
第~\ref{sec:gr-merge}~節 \\
放到最早的空時段 & 期限排程把工作放進最早的空時段，擋住後面期限早的工作 &
第~\ref{sec:gr-deadline}~節 \\
把根當成空時段 & 互斥集合用了依大小合併，卻直接拿 \proc{Find-Set} 的結果當空時段；
或合併之後才去讀左邊那組的空時段 & 第~\ref{sec:gr-deadline}~節 \\
\bottomrule
\end{tabular}
\end{center}

\section{觀念題}
\label{sec:gr-exercises}

題型和研究所考試、期中期末考的筆試相近，有選擇、計算、追蹤和證明。 解答在本章最後。

\begin{enumerate}
\item \label{q:gr-act-trace}（追蹤）籃球場收到十筆借用申請，時段都是半開區間：
\begin{center}
\begin{tabular}{c*{10}{c}}
\toprule
編號 & 1 & 2 & 3 & 4 & 5 & 6 & 7 & 8 & 9 & 10 \\
\midrule
$s_i$ & 6 & 1 & 2 & 4 & 8 & 5 & 9 & 11 & 3 & 10 \\
$f_i$ & 9 & 3 & 6 & 6 & 10 & 8 & 12 & 13 & 4 & 12 \\
\bottomrule
\end{tabular}
\end{center}
\begin{enumerate}
\item 用虛擬碼~\ref{alg:gr-activity}（結束時刻相同時編號小的先）寫出排序後的順序和選中
的申請。
\item 結束時刻相同時改成編號大的先，選中的申請有沒有變？ 筆數呢？
\item 如果時段改成閉區間（首尾相接也算衝突），最多能選幾筆？
\end{enumerate}

\item （選擇，可能不只一個答案）下列哪些敘述正確？
\begin{enumerate}
\item 分數背包依單位重量價值由大到小拿，得到最大總價值。
\item 0/1 背包依單位重量價值由大到小拿（放得下才拿），得到的總價值至少是最佳值的一半。
\item 霍夫曼演算法建出的碼樹一定是嚴格二元樹。
\item 頻率嚴格最高（比其他每個符號都高）的符號，在任何最佳前綴碼裡，碼長都不比其他
符號長。
\item $2$ 路合併時，每次合併目前最短的兩個串，總合併成本最小。
\item $3$ 路合併時，每次合併目前最短的三個串，總合併成本一定最小。
\item 一個最佳化問題只要有最佳子結構，貪心就一定對。
\end{enumerate}

\item \label{q:gr-huff-trace}（追蹤）八個符號的頻率如下。 用本書的平手規則（先比頻率、再比編號；葉節點照表上
的順序編號 $0$ 到 $7$；先取出的當左孩子，左邊標 \texttt{0}）建霍夫曼樹，寫出每一次合
併、每個符號的碼和總位元數，並和定長碼比較。
\begin{center}
\begin{tabular}{c*{8}{c}}
\toprule
符號 & a & b & c & d & e & f & g & h \\
\midrule
頻率 & 9 & 4 & 15 & 4 & 30 & 7 & 12 & 19 \\
\bottomrule
\end{tabular}
\end{center}

\item （選擇，可能不只一個答案）六個符號，頻率都是正的。 下列哪些碼長組合，可能是某
組頻率的最佳前綴碼的碼長？
\begin{enumerate}
\item $1, 2, 3, 4, 5, 5$
\item $2, 2, 2, 3, 3, 3$
\item $2, 2, 3, 3, 3, 3$
\item $2, 2, 2, 3, 4, 4$
\item $1, 3, 3, 3, 3, 3$
\item $2, 2, 3, 3, 3, 4$
\end{enumerate}

\item \label{q:gr-kmerge-q}（計算）外部排序產生了八個串，長度（千筆）依序是
$14, 6, 22, 9, 30, 11, 3, 17$。 一次合併的成本是合併出來的串長。
\begin{enumerate}
\item 求 $2$ 路合併的最小總成本，寫出每一次合併。
\item 改成 $3$ 路合併：要補幾個虛擬串？ 最小總成本是多少？
\item $3$ 路合併但不補虛擬串，每次直接合併最短的三個（最後不夠三個就有幾個合幾個），
總成本是多少？
\end{enumerate}

\item \label{q:gr-wspt}（證明）同第~\ref{sec:gr-spt}~節，但工作 $i$ 有權重 $w_i > 0$，
要讓 $\sum w_i C_i$ 最小。 證明：依 $p_i / w_i$ 由小到大排的順序是最佳的。 所有
$w_i = 1$ 時，這個結果變成什麼？

\item \label{q:gr-dl-trace}（追蹤）八件單位時間的工作：
\begin{center}
\begin{tabular}{c*{8}{c}}
\toprule
工作 & 1 & 2 & 3 & 4 & 5 & 6 & 7 & 8 \\
\midrule
期限 & 2 & 4 & 1 & 3 & 2 & 4 & 5 & 3 \\
收費 & 40 & 15 & 30 & 55 & 25 & 45 & 10 & 20 \\
\bottomrule
\end{tabular}
\end{center}
用虛擬碼~\ref{alg:gr-deadline-dsu}（\proc{Union} 依大小合併，一樣大時第二個參數的根接
到第一個的根底下；\proc{Find-Set} 做路徑壓縮）排程。 寫出每接下一件工作後的父節點陣
列 $p[0..D]$，以及最後的排程和總收費。

\item \label{q:gr-coin}（計算）
\begin{enumerate}
\item 面額 $1, 4, 9$。 找出貪心（虛擬碼~\ref{alg:gr-change}）不是最少枚數的最小金額，
寫出貪心的結果和最佳的湊法。
\item 用 Kozen 與 Zaks 的結論，要確定 $1, 5, 10, 50$ 是標準的，需要檢查哪些金額？
說明這組面額為什麼是標準的。
\end{enumerate}

\item \label{q:gr-ahead}（證明）另一種證明活動選擇的方法。 設貪心選到的申請依序是
$g_1, g_2, \ldots, g_k$，$o_1, o_2, \ldots, o_m$ 是任何一組兩兩相容的申請，依結束時刻
排好。 證明：對每個 $i \le \min(k, m)$，$f(g_i) \le f(o_i)$。 再用這個結果證明
$k \ge m$。

\item （選擇，可能不只一個答案）下列哪些是擬陣（元素集合／獨立的子集合）？
\begin{enumerate}
\item 一個無向圖的邊／不含環的邊集合。
\item 一群區間／兩兩相容（不重疊）的區間集合。
\item 一群有期限的單位時間工作／能在期限內全部做完的工作集合。
\item 一個無向圖的邊／沒有兩條邊共用端點的邊集合（匹配）。
\item 一個集合 $S$ 的元素／元素個數不超過 $3$ 的子集合。
\end{enumerate}
\end{enumerate}

\section{實作練習}
\label{sec:gr-practice}

下面是 Aki-Judge 上和本章相關的題目。 題目不多，霍夫曼編碼、最佳合併和期限排程目前
沒有對應的實作題，請用觀念題第~\ref{q:gr-huff-trace}--\ref{q:gr-kmerge-q}、\ref{q:gr-dl-trace}~題和範例~\ref{ex:gr-huffman}、
範例~\ref{ex:gr-deadline} 的程式練習。

\begin{center}
\begin{tabular}{llc}
\toprule
小節 & 題目 & 難度 \\
\midrule
活動選擇 & \akiref{most-compatible-bookings}{最多筆不衝突的借用申請} & 基本 \\
分數背包 & \akiref{maximum-value-with-divisible-items}{可分割物品的最大總價值} & 基本 \\
最小化完成時間 & \akiref{smallest-total-waiting-time}{最小的總等待時間} & 暖身 \\
貪心何時不對 & \akiref{fewest-coins-for-an-amount}{湊出金額的最少硬幣數} & 進階 \\
\bottomrule
\end{tabular}
\end{center}

幾題的提示：

\begin{itemize}
\item 「最多筆不衝突的借用申請」：時段是半開區間，和範例~\ref{ex:gr-activity} 一樣；
$n$ 可以是 $0$。 只要印筆數。
\item 「可分割物品的最大總價值」：比值用交叉相乘比較，$v_i w_j$ 對 $v_j w_i$，不要用
浮點數除法；乘積要用 \texttt{long long}。 拿一部分時的價值是整數除法、無條件捨去，
題目也規定了比值相同時的先後。
\item 「最小的總等待時間」問的是等待時間（開始之前等了多久），不是完成時間；總和可能
超過 \texttt{int} 的範圍。
\item 「湊出金額的最少硬幣數」的面額是任意的，貪心會錯，要用第~\ref{sec:gr-fail}~節
的遞迴式做動態規劃（第~18~章）。 先用 $1, 5, 12$ 湊 $15$ 這類例子確認你的程式不是貪
心。
\end{itemize}

\section{讀完你應該會}
\label{sec:gr-checklist}

\begin{itemize}
\item 對一個貪心演算法，寫出「貪心選擇性質」和「最佳子結構」兩部分的證明；會用交換論
證，也會用逐步比較（觀念題第~\ref{q:gr-ahead}~題）的寫法。
\item 手算活動選擇、分數背包，並舉出 0/1 背包、找零錢、加權活動選擇讓貪心失敗的反例。
\item 依指定的平手規則逐步畫出霍夫曼樹，寫出每個符號的碼、總位元數，並解碼一段位元
串；知道碼長是葉節點的深度（邊數）。
\item 算出 $2$ 路和 $k$ 路的最佳合併成本，知道什麼時候要補幾個虛擬串，並用
$\sum \ell_i d_i$ 驗算。
\item 用最短工作優先排程並證明它最佳；用互斥集合實作有期限的工作排程，手算每一步的父
節點陣列。
\item 說出擬陣的定義，判斷一個「獨立集合」的系統是不是擬陣。
\end{itemize}

\textbf{延伸閱讀}（書末會整理完整的參考資料）

\begin{itemize}
\item Thomas H. Cormen 等，\textit{Introduction to Algorithms}，第 4 版，第 15 章：活
動選擇、貪心法的要素、分數與 0/1 背包、霍夫曼碼，證明比本章更形式化。 擬陣和用擬陣
處理期限排程在第 3 版的 16.4、16.5 節（第 4 版移到書的網站）。
\item Jon Kleinberg 與 Éva Tardos，\textit{Algorithm Design}，第 4 章：把貪心的證明分成
「逐步領先」和「交換論證」兩種寫法，例子多，適合想練習寫證明的人。
\item Steven S. Skiena，\textit{The Algorithm Design Manual}，第 3 版，第 1 章 1.2、
1.3 節：用排程問題示範怎麼找貪心的反例。
\item Antti Laaksonen，\textit{Competitive Programmer's Handbook}，第 6 章：找零、排
程、期限、霍夫曼碼，每個主題一兩頁，適合快速複習。
\item Donald E. Knuth，\textit{The Art of Computer Programming}，第 1 卷，第 3 版，
2.3.4.5 節：路徑長與霍夫曼的方法，包括 $t$ 元樹要補權重 $0$ 的節點。
\end{itemize}

\addcontentsline{toc}{section}{觀念題解答}
\section*{觀念題解答}

\begin{enumerate}
\item
\begin{enumerate}
\item 依結束時刻排好（平手編號小的先）：$2, 9, 3, 4, 6, 1, 5, 7, 10, 8$。 選 $2$
$[1, 3)$；$9$ $[3, 4)$ 在 $3$ 開始，選；$3$ $[2, 6)$ 跳過；$4$ $[4, 6)$ 選；$6$
$[5, 8)$ 跳過；$1$ $[6, 9)$ 選；$5$ $[8, 10)$ 跳過；$7$ $[9, 12)$ 選；$10$、$8$ 跳過。
選中 $2, 9, 4, 1, 7$，共 $5$ 筆。
\item 順序變成 $2, 9, 4, 3, 6, 1, 5, 10, 7, 8$。 $4$ 和 $3$ 對調不影響（$3$ 仍然被跳
過）；$10$ $[10, 12)$ 排到 $7$ 前面，它在 $10 \ge 9$ 開始，被選中，$7$ 改被跳過。 選中
$2, 9, 4, 1, 10$，換了一筆，筆數仍是 $5$（定理~\ref{thm:gr-act-choice} 對任何一筆結束最早
的都成立）。
\item 閉區間時相容的條件是 $s_i > \mathit{last}$。 順序同 (a)：選 $2$ $[1, 3]$；$9$
在 $3$ 開始，不行；$3$ 不行；$4$ $[4, 6]$ 選；$6$ 不行；$1$ 在 $6$ 開始，不行；$5$
$[8, 10]$ 選；$7$、$10$ 不行；$8$ $[11, 13]$ 選。 共 $4$ 筆；證明照樣成立，所以最多就是 $4$ 筆。
\end{enumerate}

\item (a)、(c)、(d)、(e) 正確。 (b)：容量 $10$，物品（重 $1$，值 $2$）和（重 $10$，
值 $10$）。 前者比值 $2$ 先拿，剩 $9$ 放不下後者，得 $2$，最佳是 $10$，連一半都不到；
把後者的重量和價值放大，比例可以任意小。 (c)：每次合併都接兩個孩子，沒有只有一個孩
子的節點。 (d)：如果 $x$ 頻率比 $y$ 高卻比 $y$ 深，交換兩者 $B$ 嚴格變小（第~\ref{sec:gr-huff-proof}~節的式子），和最佳矛盾。 (f)：第~\ref{sec:gr-merge}~節
的例子，不補虛擬串得到 $191$，最佳是 $154$。 (g)：0/1 背包有最佳子結構，貪心卻會錯。

\item 編號：a $0$、b $1$、c $2$、d $3$、e $4$、f $5$、g $6$、h $7$，內部節點從 $8$
開始。 七次合併（先取出的寫在前面）：b $4$ + d $4$ = \#8 $8$；f $7$ + \#8 $8$ = \#9 $15$；
a $9$ + g $12$ = \#10 $21$；c $15$ + \#9 $15$ = \#11 $30$（平手，葉節點編號小，先取 c）；
h $19$ + \#10 $21$ = \#12 $40$；e $30$ + \#11 $30$ = \#13 $60$（平手，先取 e）；\#12 + \#13 = $100$。
根的左子樹是 \#12（h 和 \#10），右子樹是 \#13（e 和 \#11）。 碼：h \texttt{00}、
a \texttt{010}、g \texttt{011}、e \texttt{10}、c \texttt{110}、f \texttt{1110}、
b \texttt{11110}、d \texttt{11111}。 總位元數
$9 \cdot 3 + 4 \cdot 5 + 15 \cdot 3 + 4 \cdot 5 + 30 \cdot 2 + 7 \cdot 4 + 12 \cdot 3
+ 19 \cdot 2 = 274$，也等於內部節點的和 $8 + 15 + 21 + 30 + 40 + 60 + 100 = 274$。
定長碼要 $3 \times 100 = 300$ 位元。 把範例~\ref{ex:gr-huffman} 的輸入改成這八個符號，
可以核對每一個碼。

\item (a)、(c)、(d)。 嚴格二元樹的葉節點深度一定滿足 $\sum 2^{-l_i} = 1$（根算 $1$，每個內
部節點把自己的值平分給兩個孩子），反過來滿足這個等式的碼長都排得出一棵嚴格二元樹。 (a)
$\frac12 + \frac14 + \frac18 + \frac1{16} + \frac2{32} = 1$；(b) $\frac34 + \frac38 > 1$，連
前綴碼都不是；(c) $\frac12 + \frac12 = 1$；(d) $\frac34 + \frac18 + \frac2{16} = 1$；(e)
$\frac12 + \frac58 > 1$；(f) $\frac{15}{16} < 1$，碼樹不是嚴格的，不是最佳。 (a)、(c)、(d)
分別是頻率 $16, 8, 4, 2, 1, 1$、全部相同、$4, 4, 4, 2, 1, 1$ 的霍夫曼碼長。

\item
\begin{enumerate}
\item 每次合併最短的兩個：$3 + 6 = 9$；$9 + 9 = 18$（原本的 $9$ 和剛合出的 $9$）；
$11 + 14 = 25$；$17 + 18 = 35$；$22 + 25 = 47$；$30 + 35 = 65$；$47 + 65 = 112$。
總成本 $311$。 （照第~15~章一趟一趟兩兩相鄰合併，三趟共 $336$。）
\item $r = 8$，$(8 - 1) \bmod 2 = 1$，補 $1$ 個虛擬串。 $0 + 3 + 6 = 9$；
$9 + 9 + 11 = 29$；$14 + 17 + 22 = 53$；$29 + 30 + 53 = 112$。 總成本 $203$。
\item $3 + 6 + 9 = 18$；$11 + 14 + 17 = 42$；$18 + 22 + 30 = 70$；最後只剩兩個，
$42 + 70 = 112$。 總成本 $242$，比 (b) 多 $39$。
\end{enumerate}

\item 設相鄰的 $i$、$j$，$i$ 在前，前面的工作在 $t$ 做完。 交換前
$w_i(t + p_i) + w_j(t + p_i + p_j)$，交換後 $w_j(t + p_j) + w_i(t + p_j + p_i)$，其他
工作不變。 相減：
\[
(\text{交換前}) - (\text{交換後}) = w_j p_i - w_i p_j .
\]
若 $p_i / w_i > p_j / w_j$，即 $p_i w_j > p_j w_i$，交換讓總和嚴格變小。 所以最佳順序
不會有相鄰的一對違反「$p / w$ 由小到大」，也就是一定照 $p / w$ 由小到大排；比值相同
的交換總和不變。 $w_i$ 全部是 $1$ 時，$p / w = p$，就是定理~\ref{thm:gr-spt} 的最短工
作優先。

\item $D = \min(8, 5) = 5$。 依收費排好：$4, 6, 1, 3, 5, 8, 2, 7$。
\begin{center}
\begin{tabular}{cccl}
\toprule
工作 & $\min(d, D)$ & 結果 & $p[0..5]$ \\
\midrule
4 & 3 & 時段 3 & 0 1 3 3 4 5 \\
6 & 4 & 時段 4 & 0 1 3 3 3 5 \\
1 & 2 & 時段 2 & 0 3 3 3 3 5 \\
3 & 1 & 時段 1 & 3 3 3 3 3 5 \\
5 & 2 & 不接 & 不變 \\
8 & 3 & 不接 & 不變 \\
2 & 4 & 不接 & 不變 \\
7 & 5 & 時段 5 & 3 3 3 3 3 3 \\
\bottomrule
\end{tabular}
\end{center}
工作 6 填了時段 $4$ 以後，\proc{Union}$(4, 3)$：$4$ 那組只有 $1$ 個，$3$ 那組有 $2$
個，所以 $4$ 接到 $3$ 底下，$\mathit{latest}[3] = 2$。 時段 $1$ 到 $5$ 依序是工作
$3, 1, 4, 6, 7$，總收費 $30 + 40 + 55 + 45 + 10 = 180$。 用範例~\ref{ex:gr-deadline}
跑這組輸入，印出的父節點陣列和上表相同。

\item
\begin{enumerate}
\item 金額 $1$ 到 $11$ 貪心都對（例如 $8 = 4 + 4$、$11 = 9 + 1 + 1$）。 金額 $12$：
貪心拿 $9 + 1 + 1 + 1$，$4$ 枚；$4 + 4 + 4$ 只要 $3$ 枚。 最小的反例是 $12$，也符合
Kozen--Zaks 的界 $12 < 4 + 9 = 13$。
\item 只要檢查 $1$ 到 $59$。 最佳湊法裡 $1$ 元最多 $4$ 枚、$5$ 元最多 $1$ 枚、$10$ 元最多
$4$ 枚（否則可以換成更大的一枚，枚數變少），這三種加起來最多 $49$ 元；所以金額
$\ge 50$ 時最佳湊法一定用到 $50$ 元，扣掉後剩下的金額 $\ge 10$ 時一定用到 $10$ 元，依此類
推，正是貪心的選擇。 範圍內貪心都最少，這組面額是標準的。
\end{enumerate}

\item 對 $i$ 歸納。 $i = 1$：$g_1$ 是所有申請裡結束最早的，$f(g_1) \le f(o_1)$。 假設
$f(g_{i-1}) \le f(o_{i-1})$。 $o_i$ 和 $o_{i-1}$ 相容而且在它之後，
$s(o_i) \ge f(o_{i-1}) \ge f(g_{i-1})$，所以選完 $g_{i-1}$ 時，$o_i$ 仍然可以選。 貪心
接下來選的 $g_i$，是開始時刻 $\ge f(g_{i-1})$ 的申請裡結束最早的，所以
$f(g_i) \le f(o_i)$。

假設 $m > k$。 由上面的結果 $f(g_k) \le f(o_k) \le s(o_{k+1})$，所以 $o_{k+1}$ 在
$g_k$ 之後仍然可以選，而它的結束時刻 $f(o_{k+1}) > s(o_{k+1}) \ge f(g_k)$，排在
$g_k$ 後面。 貪心掃到它的時候 $\mathit{last} = f(g_k) \le s(o_{k+1})$，會把它選進來，
和「貪心只選了 $k$ 筆」矛盾。 所以 $k \ge m$，貪心選的筆數不少於任何一組相容的申請。
這種寫法叫\textbf{逐步領先}（greedy stays ahead）：不改動最佳解，直接證明貪心每一步
都不落後。

\item (a)、(c)、(e)。 (a) 是圖的擬陣（第~\ref{sec:gr-matroid}~節）。 (b)：$A =
\{[0, 5)\}$、$B = \{[0, 1), [2, 3)\}$，$B$ 的區間都和 $[0, 5)$ 重疊，交換性不成立。
(c) 在第~\ref{sec:gr-matroid}~節說明過。 (d)：路徑 $a - b - c - d$，$A = \{bc\}$、
$B = \{ab, cd\}$，交換性不成立。 (e)：$|A| < |B| \le 3$ 時，$B \setminus A$ 的任何元素加進 $A$ 都不超過 $3$
個，交換性成立；這叫均勻擬陣。
\end{enumerate}

\chapter{回溯與分支限界}
\label{ch:backtracking}

% 本章的圖共用的樣式（前綴 bt-）
\tikzset{
    bt-n/.style={circle, draw=black!65, minimum size=0.56cm, inner sep=0pt, font=\scriptsize, fill=white},
    bt-dead/.style={bt-n, draw=black!35, text=black!45, dashed, label={[font=\scriptsize, text=black!55, inner sep=1pt]below:$\times$}},
    bt-ans/.style={bt-n, draw=memStroke, very thick, fill=memFill},
    bt-texp/.style={bt-n, draw=coreStroke, thick, fill=coreFill},
    bt-e/.style={draw=black!55},
    bt-el/.style={font=\tiny, inner sep=0.8pt, fill=white, text=black!75},
    bt-num/.style={font=\tiny, text=black!60, inner sep=0.5pt},
    bt-ord/.style={font=\tiny\bfseries, text=coreStroke, inner sep=0.5pt},
    bt-cost/.style={font=\tiny, text=black!80, inner sep=1pt},
    bt-sbox/.style={rectangle, rounded corners=2pt, draw=black!65, font=\tiny, inner sep=1.5pt, minimum width=0.78cm, minimum height=0.42cm, fill=white},
    bt-sans/.style={bt-sbox, draw=memStroke, very thick, fill=memFill},
    bt-kbox/.style={rectangle, rounded corners=2pt, draw=black!65, font=\scriptsize, align=center, inner sep=1.5pt, minimum width=1.05cm, fill=white},
    bt-kexp/.style={bt-kbox, draw=coreStroke, thick, fill=coreFill},
    bt-kdead/.style={bt-kbox, draw=black!35, text=black!50, dashed},
    bt-kbest/.style={bt-kbox, draw=memStroke, very thick, fill=memFill},
    bt-queen/.style={circle, fill=black!80, minimum size=0.2cm, inner sep=0pt},
    bt-gv/.style={circle, draw=black!70, minimum size=0.6cm, inner sep=0pt, font=\small, fill=white},
    bt-ge/.style={draw=black!60, thick},
    bt-lvl/.style={font=\scriptsize, text=black!55, anchor=east},
}
% \btboard{x}{y}{每一列皇后的欄號，0 表示這列還沒放}：左下角在 (x, y) 的 4×4 小棋盤
\newcommand{\btboard}[3]{%
    \begin{scope}[shift={(#1,#2)}]
    \draw[black!45] (0,0) grid[step=0.36] (1.44,1.44);
    \foreach \btc [count=\btr] in {#3} {
        \ifnum\btc>0
            \node[bt-queen] at ({(\btc-0.5)*0.36},{1.44-(\btr-0.5)*0.36}) {};
        \fi
    }
    \end{scope}%
}

社團週末去爬山，小組要從十二樣裝備裡挑一些放進背包，背包最多只能背十三公斤。 每樣裝備
有重量，也有一個大家討論出來的「有用程度」。 把每一種挑法都列出來是 $4096$ 種，一種一
種檢查還不算慢；裝備變成四十樣，就是一兆多種。 一個省事的做法是一樣一樣決定拿或不拿，
決定到一半發現已經超重，後面那幾樣就不必再試。 還能更省嗎？ 假設目前已經裝進去的，加
上「剩下的全部都拿」，有用程度的總和還比不上已經找到的最好組合，後面還要試嗎？ 往下讀
之前，先想想看這兩種「不必再試」各需要知道什麼。

這一章講兩種有系統地搜尋所有可能解的方法。 \textbf{回溯}（backtracking）把一個解看成一
連串的決定，一次做一個決定，發現走不通就退回上一步改試別的；\textbf{分支限界}（branch
and bound）一樣把所有可能畫成一棵樹，但每次挑「看起來最有希望」的節點往下展開，並用上
下界把不可能更好的子樹整棵丟掉。 兩種方法的最壞情況通常都還是指數時間，它們能做的是讓
實際走過的節點少很多。 研究所考試常考的是手算：畫出狀態空間樹、依序標出節點的編號、寫出
分支限界每一步選了哪個節點、算出界限。 這一章的例子都照這個方式一步一步寫出來。

這一章用到《C 程式設計基礎》第~15~章的遞迴與回溯、第~19~章的 \texttt{malloc}、
\texttt{free}；分支限界的優先佇列用第~8~章的堆積，佇列與堆疊見第~6~章。 0/1 背包的動態規
劃在第~18~章，分數背包的貪心解在第~19~章，這一章會用到後者當作上界。

\section{狀態空間樹與回溯}
\label{sec:bk-tree}

\begin{cbookrecap}[第~15~章 15.15 節]
回溯的程式有固定的形狀：在每一層列出所有可以做的選擇，對每一個選擇先「做」，再遞迴往下
「探索」，回來之後「撤銷」，然後試下一個選擇。 狀態如果用參數傳值（例如還要湊的和
$k$），子呼叫回來時呼叫的人手上的值沒有變，不必另外撤銷；狀態如果是同一塊被就地修改的陣
列（例如棋盤），回來時一定要把改過的格子還原。 那一節寫過三個例子：數有幾個子集合的和是
$k$、數 $n$ 皇后有幾組解、用交換產生所有排列。 這一章不再重講這些程式怎麼寫，重點放在
「搜尋的是哪一棵樹」以及「哪些節點可以不走」。
\end{cbookrecap}

\subsection{解向量、限制與狀態空間樹}

很多問題的答案可以寫成一個向量 $(x_1, x_2, \ldots, x_n)$，每一個 $x_i$ 從一個有限的集合
$S_i$ 裡挑。 例如 $n$ 皇后問題（在 $n \times n$ 的棋盤上放 $n$ 個皇后，任兩個不在同一
列、同一欄或同一條對角線上），可以讓 $x_i$ 表示第 $i$ 列的皇后放在第幾欄。 每一列一定恰
好一個皇后，所以這樣寫不會漏掉任何解。

問題的條件分成兩種：

\begin{itemize}
\item \textbf{顯式限制}（explicit constraint）規定每一個 $x_i$ 能取哪些值，也就是 $S_i$。
例如 $x_i \in \{1, 2, \ldots, n\}$。
\item \textbf{隱式限制}（implicit constraint）規定各個 $x_i$ 之間的關係。 例如任兩個皇后不
同欄（$x_i \ne x_j$），也不在同一條對角線（$|x_i - x_j| \ne |i - j|$）。
\end{itemize}

同一個問題可以把不同的條件放在顯式限制裡。 把「不同欄」也算成顯式限制，$(x_1, \ldots,
x_n)$ 就必須是 $1, \ldots, n$ 的一個排列，要搜尋的可能性從 $n^n$ 種降到 $n!$ 種。

把所有可能的向量組織成一棵樹，叫做\textbf{狀態空間樹}（state space tree）。 根代表什麼都
還沒決定；深度 $k$ 的節點代表已經決定好 $x_1, \ldots, x_k$；一個節點的子節點對應 $x_{k+1}$
的每一個可能值。 樹上的名詞：

\begin{itemize}
\item 每一個節點都是一個\textbf{問題狀態}（problem state）。
\item 從根到某個節點的路徑如果已經是一個滿足顯式限制的完整向量，這個節點是
\textbf{解狀態}（solution state）；向量長度固定為 $n$ 時，解狀態就是深度 $n$ 的葉節點。
\item 解狀態如果也滿足隱式限制，叫做\textbf{答案狀態}（answer state），也叫答案節點。
\end{itemize}

$4$ 皇后的狀態空間樹，如果每一列都允許 $4$ 個欄位，有 $1 + 4 + 4^2 + 4^3 + 4^4 = 341$ 個
節點；用排列樹（同一條路徑上的欄號不重複），只剩 $1 + 4 + 4 \cdot 3 + 4 \cdot 3 \cdot 2 +
4! = 65$ 個節點。 這一章的 $n$ 皇后一律用排列樹。

搜尋時要追蹤每個節點的狀況：

\begin{itemize}
\item \textbf{活節點}（live node）：已經產生出來、但子節點還沒有全部產生的節點。
\item \textbf{E-節點}（E-node，expansion node，擴展節點）：目前正在產生子節點的那一個活節點。
\item \textbf{死節點}（dead node）：被判定不必再展開，或子節點已經全部產生完的節點。
\end{itemize}

判定一個節點不必展開的依據叫做\textbf{界限函數}（bounding function）：給定
$x_1, \ldots, x_k$，如果可以確定它不管怎麼往下補都不會成為答案，界限函數就回答「否」，這
個節點直接變成死節點，整棵子樹都不走。 這個動作叫做\textbf{剪枝}（pruning）。

\textbf{回溯}就是用深度優先的順序走狀態空間樹，邊走邊產生節點，並用界限函數剪枝。 一個
節點一產生出來，只要沒有被剪掉，就馬上變成新的 E-節點，原來的 E-節點要等這個子節點的子樹
走完才繼續產生下一個子節點。 第~\ref{sec:bk-bb}~節的分支限界則是讓一個 E-節點一次把子節點
全部產生完，再從活節點裡挑下一個 E-節點。 兩種方法搜尋的是同一棵樹，差別在節點展開的順
序。

\begin{pseudo}[label=alg:bk-generic, unbreakable]{回溯的骨架}
\Procedure{Backtrack}{$k$}
    \For{$x_k$ 的每一個候選值 $c$，依固定的順序}
        \If{\Call{Promising}{$x_1, \ldots, x_{k-1}, c$}}
            \State $x_k \gets c$
            \If{$(x_1, \ldots, x_k)$ 是答案}
                \State 輸出 $(x_1, \ldots, x_k)$
            \EndIf
            \If{$k < n$}
                \State \Call{Backtrack}{$k + 1$}
            \EndIf
        \EndIf
    \EndFor
\EndProcedure
\end{pseudo}

第 3 行的 \proc{Promising} 就是界限函數，回傳 \TRUE{} 表示這個節點還有希望。 呼叫
\proc{Backtrack}$(1)$ 開始搜尋。 只要一組答案時，在第 6 行找到答案後設一個旗標，讓所有
呼叫都提早結束。 回溯法的時間大約是「產生的節點數 $\times$ 每個節點的成本」；每個節點的
成本主要是第 3 行，節點數則要看界限函數剪掉多少，很難事先算準。

\subsection{\texorpdfstring{$n$}{n} 皇后}
\label{sec:bk-queens}

第 $k$ 列要把皇后放在第 $c$ 欄時，要檢查前面 $k - 1$ 列的皇后。 同一欄的檢查交給排列
樹：已經用過的欄不算子節點。 對角線的檢查有兩種寫法。 直接的寫法是逐一檢查 $i = 1, \ldots,
k - 1$ 是否 $|x_i - c| = k - i$，每個節點花 $O(k)$。 另一種寫法注意到同一條「左下到右上」
的對角線上，列號加欄號是定值；同一條「左上到右下」的對角線上，列號減欄號是定值。 用兩個
布林陣列記錄每一條對角線有沒有皇后，每個節點只要 $O(1)$：

\begin{pseudo}[label=alg:bk-queens]{$n$ 皇后（排列樹，用陣列檢查）}
\Procedure{Queens}{$k$}
    \For{$c \gets 1$ \To $n$}
        \If{$\mathit{col}[c] = \FALSE$} \Comment{排列樹：用過的欄不是子節點}
            \If{$d_1[k + c] = \FALSE$ \AND $d_2[k - c + n] = \FALSE$}
                \State $x_k \gets c$
                \State $\mathit{col}[c] \gets d_1[k + c] \gets d_2[k - c + n] \gets \TRUE$
                \If{$k = n$}
                    \State 輸出 $(x_1, \ldots, x_n)$
                \Else
                    \State \Call{Queens}{$k + 1$}
                \EndIf
                \State $\mathit{col}[c] \gets d_1[k + c] \gets d_2[k - c + n] \gets \FALSE$
            \EndIf
        \EndIf
    \EndFor
\EndProcedure
\end{pseudo}

$k + c$ 的範圍是 $2$ 到 $2n$，$k - c + n$ 的範圍是 $1$ 到 $2n - 1$，所以兩個對角線陣列各
要 $2n - 1$ 格（$d_1$ 從 $2$ 開始編號）。 第 11 行是撤銷：這三個陣列是所有遞迴呼叫共用的，
不還原的話，下一個候選值看到的是上一條分支留下的記號。

\subsubsection{4 皇后的狀態空間樹}

圖~\ref{fig:bk-4queens} 是 \proc{Queens}$(1)$ 在 $n = 4$ 時，找到第一組解之前產生的所有節
點。 節點裡的數字是產生的順序（根是 $1$），邊上的數字是這一列的皇后放在第幾欄。 虛線節點
上打叉的，是第 4 行的對角線檢查沒有通過、產生之後馬上變成死節點的節點；橘色是答案節點。

\begin{center}
\begin{tikzpicture}[x=1.35cm,y=1.05cm]
\node[bt-n] (q1) at (3.681,0.0) {1};
\node[bt-n] (q2) at (1.662,-1.05) {2};
\node[bt-dead] (q3) at (0.0,-2.1) {3};
\node[bt-n] (q4) at (1.425,-2.1) {4};
\node[bt-dead] (q5) at (0.95,-3.15) {5};
\node[bt-dead] (q6) at (1.9,-3.15) {6};
\node[bt-n] (q7) at (3.325,-2.1) {7};
\node[bt-n] (q8) at (2.85,-3.15) {8};
\node[bt-dead] (q9) at (2.85,-4.2) {9};
\node[bt-dead] (q10) at (3.8,-3.15) {10};
\node[bt-n] (q11) at (5.7,-1.05) {11};
\node[bt-dead] (q12) at (4.75,-2.1) {12};
\node[bt-dead] (q13) at (5.7,-2.1) {13};
\node[bt-n] (q14) at (6.65,-2.1) {14};
\node[bt-n] (q15) at (6.65,-3.15) {15};
\node[bt-ans] (q16) at (6.65,-4.2) {16};
\draw[bt-e] (q1) -- node[bt-el] {1} (q2);
\draw[bt-e] (q2) -- node[bt-el] {2} (q3);
\draw[bt-e] (q2) -- node[bt-el] {3} (q4);
\draw[bt-e] (q4) -- node[bt-el] {2} (q5);
\draw[bt-e] (q4) -- node[bt-el] {4} (q6);
\draw[bt-e] (q2) -- node[bt-el] {4} (q7);
\draw[bt-e] (q7) -- node[bt-el] {2} (q8);
\draw[bt-e] (q8) -- node[bt-el] {3} (q9);
\draw[bt-e] (q7) -- node[bt-el] {3} (q10);
\draw[bt-e] (q1) -- node[bt-el] {2} (q11);
\draw[bt-e] (q11) -- node[bt-el] {1} (q12);
\draw[bt-e] (q11) -- node[bt-el] {3} (q13);
\draw[bt-e] (q11) -- node[bt-el] {4} (q14);
\draw[bt-e] (q14) -- node[bt-el] {1} (q15);
\draw[bt-e] (q15) -- node[bt-el] {3} (q16);

\node[bt-lvl] at (-0.4,-1.05) {第 1 列};
\node[bt-lvl] at (-0.4,-2.10) {第 2 列};
\node[bt-lvl] at (-0.4,-3.15) {第 3 列};
\node[bt-lvl] at (-0.4,-4.20) {第 4 列};
\end{tikzpicture}
\diagcaption{4 皇后的回溯，到找到第一組解為止。 節點內是產生的順序，邊上是皇后所在的欄，
左邊標的是這一層決定第幾列。 打叉的節點被對角線檢查剪掉。}
\label{fig:bk-4queens}
\end{center}

照產生的順序讀一遍：

\begin{enumerate}
\item 節點 2：第 1 列放第 1 欄。 第 1 列沒有別的皇后要檢查，一定通過。
\item 節點 3：第 2 列放第 2 欄，和 $(1, 1)$ 在同一條對角線上，剪掉。 節點 4：第 2 列放第
3 欄，通過。
\item 節點 5、6：第 3 列只剩第 2、4 欄可以試。 第 2 欄和 $(2, 3)$ 同一條對角線，第 4 欄
也和 $(2, 3)$ 同一條對角線，兩個都剪掉。 節點 4 的子節點全部是死節點，退回節點 2。
\item 節點 7：第 2 列放第 4 欄。 節點 8：第 3 列放第 2 欄，通過。 節點 9：第 4 列只剩第 3
欄，和 $(3, 2)$ 同一條對角線，剪掉，退回節點 7。 節點 10：第 3 列放第 3 欄，和 $(2, 4)$
同一條對角線，剪掉。 節點 7、節點 2 的子節點都試完了，退回根。
\item 節點 11：第 1 列放第 2 欄。 節點 12（第 1 欄）、13（第 3 欄）都和 $(1, 2)$ 同一條對
角線。 節點 14：第 4 欄，通過。 節點 15：第 3 列放第 1 欄，通過。 節點 16：第 4 列放第 3
欄，通過，得到答案 $(2, 4, 1, 3)$。
\end{enumerate}

一共產生 $16$ 個節點，其中 $7$ 個被剪掉。 排列樹本來有 $65$ 個節點，第一組解在整棵樹裡是
用深度優先順序數過來的第 $31$ 個節點，前面那些子樹大多沒有走到。 圖~\ref{fig:bk-boards}
是節點 4 走不下去和節點 16 找到答案時的棋盤。

\begin{center}
\begin{tikzpicture}
\btboard{0}{0}{1,3,0,0}
\node[font=\footnotesize] at (0.72,-0.35) {節點 4：第 3 列無處可放};
\btboard{4.6}{0}{1,4,2,0}
\node[font=\footnotesize] at (5.32,-0.35) {節點 8：第 4 列無處可放};
\btboard{9.2}{0}{2,4,1,3}
\node[font=\footnotesize] at (9.92,-0.35) {節點 16：答案 $(2, 4, 1, 3)$};
\end{tikzpicture}
\diagcaption{三個時間點的棋盤。 前兩個是死路，最後一個是第一組解。}
\label{fig:bk-boards}
\end{center}

\begin{examnote}
畫狀態空間樹的題目，先看清楚題目用的是哪一種樹（每列 $n$ 個子節點，還是排列樹），以及節點
怎麼編號。 本書用「產生的順序」編號，被剪掉的節點也編號，並畫出來打叉；有些課本改用整棵
排列樹的深度優先編號，同一個答案節點在那種編法裡是第 $31$ 號；也有課本只畫通過檢查的節
點。 欄位一律從小到大試。 題目沒有規定時，寫一句你用的編號方式，畫到題目要求的地方（通常
是第一組解）為止。
\end{examnote}

\subsubsection{C 程式}

下面的程式對輸入的每一個 $n$ 做完整的搜尋，數出解的個數和產生的節點數（包括根和被剪掉的
節點），並記下第一組解是第幾個產生的節點。 搜尋要用到的陣列放在一個結構裡，用指標傳給遞
迴函式，所以程式裡沒有全域變數。

\begin{example}[label=ex:bk-nqueens]{c}{(nqueens.c) $n$ 皇后：計數並印出第一組解}
#include <stdio.h>
#include <stdlib.h>

// 搜尋狀態：第 r 列的皇后在第 col[r] 欄（從 0 起算）
typedef struct {
    int n, *col, *first;
    char *used, *d1, *d2;       // 欄、兩個方向的對角線有沒有皇后
    long long nodes, solutions, first_at;
} Search;

static void place(Search *s, int r)
{
    int n = s->n;
    for (int c = 0; c < n; c++) {
        if (s->used[c]) {
            continue;               // 排列樹裡沒有這個子節點
        }
        s->nodes++;                 // 產生節點 (r, c)
        if (s->d1[r + c] || s->d2[r - c + n - 1]) {
            continue;               // 對角線上已經有皇后：剪掉
        }
        s->col[r] = c;
        s->used[c] = s->d1[r + c] = s->d2[r - c + n - 1] = 1;
        if (r < n - 1) {
            place(s, r + 1);
        } else if (s->solutions++ == 0) {   // 第一組解：記下來
            for (int i = 0; i < n; i++) {
                s->first[i] = s->col[i];
            }
            s->first_at = s->nodes;
        }
        // 撤銷
        s->used[c] = s->d1[r + c] = s->d2[r - c + n - 1] = 0;
    }
}

int main(void)
{
    int n;
    while (scanf("%d", &n) == 1 && n >= 1 && n <= 16) {
        Search s = {n, malloc(n * sizeof(int)),
                    malloc(n * sizeof(int)), calloc(n, 1),
                    calloc(2 * n - 1, 1), calloc(2 * n - 1, 1),
                    1, 0, 0};
        if (s.col && s.first && s.used && s.d1 && s.d2) {
            place(&s, 0);
            printf("n = %d: %lld solutions, %lld nodes\n",
                   n, s.solutions, s.nodes);
            if (s.solutions > 0) {
                printf("first at node %lld:", s.first_at);
                for (int i = 0; i < n; i++) {
                    printf(" %d", s.first[i] + 1);  // 欄號從 1 起算
                }
                printf("\n");
            }
        }
        free(s.col); free(s.first); free(s.used);
        free(s.d1); free(s.d2);
    }
    return 0;
}
\end{example}
\iosample{4 8}{n = 4: 2 solutions, 33 nodes \\ first at node 16: 2 4 1 3 \\ n = 8: 92 solutions, 5509 nodes \\ first at node 337: 1 5 8 6 3 7 2 4}

$n = 4$ 第一組解在第 $16$ 個節點，和圖~\ref{fig:bk-4queens} 一致；搜尋完整棵樹一共產生
$33$ 個節點，找到 $2$ 組解。 下表把 $n = 4$ 到 $8$ 的節點數和兩種狀態空間樹的大小放在一
起（中間幾列用同一支程式輸入 \texttt{5 6 7} 得到）。 界限函數剪掉的比例隨著 $n$ 變大越來
越高：$n = 8$ 時排列樹有十萬多個節點，實際只產生 $5509$ 個。

\begin{center}
\begin{tabular}{crrrr}
\toprule
$n$ & 解的個數 & 回溯產生的節點 & 排列樹的節點 & 每列 $n$ 個子節點的樹 \\
\midrule
4 & 2 & 33 & 65 & 341 \\
5 & 10 & 102 & 326 & 3906 \\
6 & 4 & 357 & 1957 & 55987 \\
7 & 40 & 1346 & 13700 & 960800 \\
8 & 92 & 5509 & 109601 & 19173961 \\
\bottomrule
\end{tabular}
\end{center}

\subsection{子集合和}
\label{sec:bk-subset}

\textbf{子集合和}（sum of subsets）問題：給 $n$ 個正整數 $w_1, \ldots, w_n$ 和目標 $M$，找
出所有和恰好是 $M$ 的子集合。 解向量寫成 $(x_1, \ldots, x_n)$，$x_i = 1$ 表示選 $w_i$，
$x_i = 0$ 表示不選，狀態空間樹是一棵 $n + 1$ 層的滿二元樹，有 $2^{n+1} - 1$ 個節點，左邊的
邊代表選、右邊代表不選。 （也有課本把解寫成「選了哪幾個下標」的遞增序列，長度不固定，樹
的形狀不同，但兩種寫法列舉的是同一群子集合。）

先把 $w$ 排成由小到大。 在決定 $w_k$ 的節點上，記 $s$ 為已經選的數的和，$r = w_k + w_{k+1}
+ \cdots + w_n$ 為還沒決定的數的和。 這個節點要能通往答案，必須同時滿足：

\begin{enumerate}
\item $s + r \ge M$：剩下的全部都選還湊不到 $M$，就不可能湊到。
\item $s + w_k \le M$：剩下的數裡 $w_k$ 最小，它都放不進去的話，剩下的一個也不能選，而
$s$ 又不等於 $M$（等於的話在上一層就已經輸出了）。
\end{enumerate}

第 2 個條件用到了「由小到大排好」；沒有排序時，$w_k$ 放不下不代表後面的放不下。 虛擬碼在
產生子節點之前就先檢查子節點是否滿足這兩個條件，不滿足的子節點不產生。 選 $w_k$ 的左子節點，$s +
r$ 不變，第 1 個條件自動成立，只要檢查 $s + w_k + w_{k+1} \le M$；不選 $w_k$ 的右子節點，兩
個條件是 $s + r - w_k \ge M$ 和 $s + w_{k+1} \le M$。 為了不用另外處理 $k = n$，約定
$w_{n+1} = \infty$。

\begin{pseudo}[label=alg:bk-subset]{子集合和}
\Procedure{Sum-Of-Subsets}{$s, k, r$} \Comment{$w_1 \le \cdots \le w_n$，$w_{n+1} = \infty$}
    \If{$s + w_k = M$}
        \State $x_k \gets 1$；輸出 $x_1, \ldots, x_k$（之後的都是 $0$）
    \ElsIf{$s + w_k + w_{k+1} \le M$}
        \State $x_k \gets 1$
        \State \Call{Sum-Of-Subsets}{$s + w_k, k + 1, r - w_k$}
    \EndIf
    \If{$s + r - w_k \ge M$ \AND $s + w_{k+1} \le M$}
        \State $x_k \gets 0$
        \State \Call{Sum-Of-Subsets}{$s, k + 1, r - w_k$}
    \EndIf
\EndProcedure
\end{pseudo}

第一次呼叫 \proc{Sum-Of-Subsets}$(0, 1, \sum w_i)$，前提是 $w_1 \le M \le \sum w_i$，否則
直接沒有解。 每一次被呼叫的節點都滿足上面兩個條件，所以第 2 行的 $s + w_k = M$ 是唯一可
能在這一層湊滿的情形。 湊滿之後，選 $w_k$ 的那一邊不再往下（後面都是正數，再選就超過），
但第 7 行仍然要試不選 $w_k$ 的那一邊，因為別的組合也可能湊出 $M$。

\subsubsection{手算追蹤}

取 $w = \langle 4, 6, 7, 9, 13, 14 \rangle$、$M = 20$，總和 $53$。 圖~\ref{fig:bk-subset} 畫
出所有被產生的節點，方框裡是 $(s, r)$，左邊的小字是產生的順序，邊上的 $1$、$0$ 是這一層的
$x_k$；橘色方框是答案節點，裡面只寫湊到的和。

\begin{center}
\begin{tikzpicture}[x=1.75cm,y=1.05cm,
    bt-sbox/.append style={font=\scriptsize, minimum width=1.0cm, minimum height=0.5cm},
    bt-sans/.append style={font=\scriptsize},
    bt-num/.append style={font=\scriptsize}, bt-el/.append style={font=\scriptsize}]
\node[bt-sbox, label={[bt-num]left:1}] (s1) at (3.038,0.0) {0,53};
\node[bt-sbox, label={[bt-num]left:2}] (s2) at (1.013,-1.0) {4,49};
\node[bt-sbox, label={[bt-num]left:3}] (s3) at (0.0,-2.0) {10,43};
\node[bt-sbox, label={[bt-num]left:4}] (s4) at (0.0,-3.0) {10,36};
\node[bt-sbox, label={[bt-num]left:5}] (s5) at (2.025,-2.0) {4,43};
\node[bt-sbox, label={[bt-num]left:6}] (s6) at (1.35,-3.0) {11,36};
\node[bt-sans, label={[bt-num]left:7}] (s7) at (1.35,-4.0) {20};
\node[bt-sbox, label={[bt-num]left:8}] (s8) at (2.7,-3.0) {4,36};
\node[bt-sbox, label={[bt-num]left:9}] (s9) at (2.7,-4.0) {4,27};
\node[bt-sbox, label={[bt-num]left:10}] (s10) at (5.062,-1.0) {0,49};
\node[bt-sbox, label={[bt-num]left:11}] (s11) at (4.05,-2.0) {6,43};
\node[bt-sbox, label={[bt-num]left:12}] (s12) at (4.05,-3.0) {6,36};
\node[bt-sbox, label={[bt-num]left:13}] (s13) at (4.05,-4.0) {6,27};
\node[bt-sbox, label={[bt-num]left:14}] (s14) at (4.05,-5.0) {6,14};
\node[bt-sans, label={[bt-num]left:15}] (s15) at (4.05,-6.0) {20};
\node[bt-sbox, label={[bt-num]left:16}] (s16) at (6.075,-2.0) {0,43};
\node[bt-sbox, label={[bt-num]left:17}] (s17) at (5.4,-3.0) {7,36};
\node[bt-sbox, label={[bt-num]left:18}] (s18) at (5.4,-4.0) {7,27};
\node[bt-sans, label={[bt-num]left:19}] (s19) at (5.4,-5.0) {20};
\node[bt-sbox, label={[bt-num]left:20}] (s20) at (6.75,-3.0) {0,36};
\node[bt-sbox, label={[bt-num]left:21}] (s21) at (6.75,-4.0) {0,27};
\draw[bt-e] (s1) -- node[bt-el] {1} (s2);
\draw[bt-e] (s2) -- node[bt-el] {1} (s3);
\draw[bt-e] (s3) -- node[bt-el] {0} (s4);
\draw[bt-e] (s2) -- node[bt-el] {0} (s5);
\draw[bt-e] (s5) -- node[bt-el] {1} (s6);
\draw[bt-e] (s6) -- node[bt-el] {1} (s7);
\draw[bt-e] (s5) -- node[bt-el] {0} (s8);
\draw[bt-e] (s8) -- node[bt-el] {0} (s9);
\draw[bt-e] (s1) -- node[bt-el] {0} (s10);
\draw[bt-e] (s10) -- node[bt-el] {1} (s11);
\draw[bt-e] (s11) -- node[bt-el] {0} (s12);
\draw[bt-e] (s12) -- node[bt-el] {0} (s13);
\draw[bt-e] (s13) -- node[bt-el] {0} (s14);
\draw[bt-e] (s14) -- node[bt-el] {1} (s15);
\draw[bt-e] (s10) -- node[bt-el] {0} (s16);
\draw[bt-e] (s16) -- node[bt-el] {1} (s17);
\draw[bt-e] (s17) -- node[bt-el] {0} (s18);
\draw[bt-e] (s18) -- node[bt-el] {1} (s19);
\draw[bt-e] (s16) -- node[bt-el] {0} (s20);
\draw[bt-e] (s20) -- node[bt-el] {0} (s21);
\end{tikzpicture}
\diagcaption{$w = \langle 4, 6, 7, 9, 13, 14 \rangle$、$M = 20$ 的子集合和回溯。 方框內是
$(s, r)$，邊上是 $x_k$。 第 $d$ 層的節點正要決定 $w_{d+1}$。}
\label{fig:bk-subset}
\end{center}

看幾個節點怎麼被擋下來：

\begin{itemize}
\item 節點 3（已選 $4, 6$，$s = 10$、$r = 43$，正要決定 $7$）：選 $7$ 之後 $s = 17$，再加
下一個 $9$ 是 $26 > 20$，左子節點不產生。 不選 $7$ 的右子節點：$10 + 43 - 7 = 46 \ge 20$，
$10 + 9 = 19 \le 20$，產生節點 4。
\item 節點 4（$s = 10$，正要決定 $9$）：選 $9$ 得 $19$，再加 $13$ 超過；不選 $9$ 的話，$10 +
13 = 23 > 20$。 兩個子節點都不產生，是一條死路。
\item 節點 6（已選 $4, 7$，$s = 11$，正要決定 $9$）：$11 + 9 = 20 = M$，產生答案節點 7，輸
出 $\{4, 7, 9\}$。 不選 $9$ 的右子節點：$11 + 13 = 24 > 20$，不產生。
\item 節點 9（$s = 4$，正要決定 $13$）：選 $13$ 得 $17$，加 $14$ 超過；不選 $13$ 時 $4 + 27 -
13 = 18 < 20$，剩下的全選也不夠。 又是死路。
\end{itemize}

三組答案依序是 $\{4, 7, 9\}$（節點 7）、$\{6, 14\}$（節點 15）、$\{7, 13\}$（節點 19）。 一共
產生 $21$ 個節點，整棵滿二元樹有 $2^7 - 1 = 127$ 個。

\subsubsection{C 程式}

輸入的數要由小到大排好，程式會檢查。 \texttt{nodes} 的數法和
圖~\ref{fig:bk-subset} 一樣：每呼叫一次算一個節點，湊滿時的答案節點另外算一個。

\begin{example}[label=ex:bk-subsets]{c}{(subsets.c) 子集合和：兩個剪枝條件}
#include <stdio.h>
#include <stdlib.h>

// w[0..n-1] 由小到大；x[i] = 1 表示選了 w[i]
typedef struct {
    int n;
    long long m, *w, nodes, found;
    char *x;
} SubsetSearch;

// 正在決定 w[k]；sum 是已選的和，rest = w[k] + ... + w[n-1]
static void sum_of_subsets(SubsetSearch *s, long long sum, int k,
                           long long rest)
{
    long long *w = s->w;
    s->nodes++;
    s->x[k] = 1;
    if (sum + w[k] == s->m) {                   // 答案節點
        s->nodes++;
        s->found++;
        printf("subset:");
        for (int i = 0; i <= k; i++) {
            if (s->x[i]) {
                printf(" %lld", w[i]);
            }
        }
        printf("\n");
    } else if (k + 1 < s->n && sum + w[k] + w[k + 1] <= s->m) {
        sum_of_subsets(s, sum + w[k], k + 1, rest - w[k]);
    }
    s->x[k] = 0;                                // 撤銷
    if (k + 1 < s->n && sum + rest - w[k] >= s->m
        && sum + w[k + 1] <= s->m) {
        sum_of_subsets(s, sum, k + 1, rest - w[k]);
    }
}

int main(void)
{
    int n;
    long long m, total = 0;
    if (scanf("%d %lld", &n, &m) != 2 || n < 1) {
        return 1;
    }
    SubsetSearch s = {n, m, malloc(n * sizeof(long long)), 0, 0,
                      calloc(n, 1)};
    if (s.w == NULL || s.x == NULL) {
        free(s.w); free(s.x);
        return 1;
    }
    for (int i = 0; i < n; i++) {               // 輸入要由小到大
        if (scanf("%lld", &s.w[i]) != 1 || s.w[i] <= 0
            || (i > 0 && s.w[i] < s.w[i - 1])) {
            free(s.w); free(s.x);
            return 1;
        }
        total += s.w[i];
    }
    if (s.w[0] <= m && total >= m) {
        sum_of_subsets(&s, 0, 0, total);
    }
    printf("solutions = %lld, nodes = %lld\n", s.found, s.nodes);
    free(s.w); free(s.x);
    return 0;
}
\end{example}
\iosample{6 20 \\ 4 6 7 9 13 14}{subset: 4 7 9 \\ subset: 6 14 \\ subset: 7 13 \\ solutions = 3, nodes = 21}

C 的下標從 $0$ 開始，虛擬碼的 $w_{n+1} = \infty$ 改成檢查 \texttt{k + 1 < n}：沒有下一個數
時，選了 $w_k$ 之後不能再往下，不選 $w_k$ 的那一邊也不可能再湊滿。 \texttt{x} 是所有呼叫共
用的陣列，所以第二個遞迴之前要先把 \texttt{x[k]} 設回 $0$，這是撤銷。

\subsection{圖著色}
\label{sec:bk-coloring}

給一個無向圖 $G$ 和正整數 $m$，\textbf{$m$ 著色問題}（$m$-coloring）問：能不能用 $m$ 種顏
色替每個頂點上色，使每條邊的兩端顏色不同？ 這是一個\textbf{決策問題}，答案是「能」或「不
能」；能的時候順便輸出一種塗法。 能塗的最小 $m$ 叫做這個圖的\textbf{色數}（chromatic
number），可以依序試 $m = 1, 2, 3, \ldots$，第一個回答「能」的就是色數。 $m \ge 3$ 時這個
決策問題沒有已知的多項式時間演算法，第~28~章會再遇到它。

頂點編號 $1$ 到 $n$，解向量 $(x_1, \ldots, x_n)$，$x_i \in \{1, \ldots, m\}$ 是頂點 $i$ 的顏
色。 狀態空間樹每個內部節點有 $m$ 個子節點，深度 $n$。 界限函數：頂點 $k$ 塗 $c$ 時，它的鄰
居裡編號比 $k$ 小的（已經塗過的）都不能是 $c$。

\begin{pseudo}[label=alg:bk-color]{$m$ 著色}
\Procedure{Color}{$k$}
    \For{$c \gets 1$ \To $m$}
        \If{$k$ 的每個鄰居 $j < k$ 都有 $x_j \ne c$}
            \State $x_k \gets c$
            \If{$k = n$}
                \State 輸出 $(x_1, \ldots, x_n)$
            \Else
                \State \Call{Color}{$k + 1$}
            \EndIf
        \EndIf
    \EndFor
\EndProcedure
\end{pseudo}

第 3 行只讀 $x_1, \ldots, x_{k-1}$，所以 $x_k$ 不必撤銷，下一次寫入會覆蓋它，而且在那之前沒
有人讀它。 狀態空間樹的節點數不超過 $\sum_{i=0}^{n} m^i = O(m^n)$，用鄰接串列檢查第 3
行，每個節點 $O(n)$，最壞 $O(n m^n)$。

\subsubsection{手算追蹤}

圖~\ref{fig:bk-graph} 的圖有 $6$ 個頂點、$9$ 條邊，取 $m = 3$。 下表是 \proc{Color}$(1)$ 找
到第一種塗法之前試過的每一個 $(k, c)$，「部分解」是試之前 $x_1, \ldots, x_{k-1}$ 的顏色。

\begin{center}
\begin{tikzpicture}[x=1cm,y=1cm]
\node[bt-gv] (g1) at (0,1.2) {1};
\node[bt-gv] (g2) at (-1.3,0) {2};
\node[bt-gv] (g3) at (1.3,0) {3};
\node[bt-gv] (g4) at (-1.3,-1.6) {4};
\node[bt-gv] (g5) at (2.6,1.2) {5};
\node[bt-gv] (g6) at (1.3,-1.6) {6};
\draw[bt-ge] (g1) -- (g2);
\draw[bt-ge] (g1) -- (g3);
\draw[bt-ge] (g1) -- (g5);
\draw[bt-ge] (g2) -- (g3);
\draw[bt-ge] (g2) -- (g4);
\draw[bt-ge] (g3) -- (g5);
\draw[bt-ge] (g3) -- (g6);
\draw[bt-ge] (g4) -- (g6);
\draw[bt-ge] (g5) to[bend left=35] (g6);
\end{tikzpicture}
\diagcaption{$m$ 著色的例子：$6$ 個頂點、$9$ 條邊。}
\label{fig:bk-graph}
\end{center}

\begin{center}
\begin{tabular}{cccll}
\toprule
步驟 & $k$ & $c$ & 部分解 $x_1 \ldots x_{k-1}$ & 結果 \\
\midrule
1 & 1 & 1 & （空） & 通過 \\
2 & 2 & 1 & 1 & 和頂點 1 同色 \\
3 & 2 & 2 & 1 & 通過 \\
4 & 3 & 1 & 1 2 & 和頂點 1 同色 \\
5 & 3 & 2 & 1 2 & 和頂點 2 同色 \\
6 & 3 & 3 & 1 2 & 通過 \\
7 & 4 & 1 & 1 2 3 & 通過 \\
8 & 5 & 1 & 1 2 3 1 & 和頂點 1 同色 \\
9 & 5 & 2 & 1 2 3 1 & 通過 \\
10 & 6 & 1 & 1 2 3 1 2 & 和頂點 4 同色 \\
11 & 6 & 2 & 1 2 3 1 2 & 和頂點 5 同色 \\
12 & 6 & 3 & 1 2 3 1 2 & 和頂點 3 同色；頂點 6 三色都不行，退回 \\
13 & 5 & 3 & 1 2 3 1 & 和頂點 3 同色；頂點 5 試完，退回 \\
14 & 4 & 2 & 1 2 3 & 和頂點 2 同色 \\
15 & 4 & 3 & 1 2 3 & 通過 \\
16 & 5 & 1 & 1 2 3 3 & 和頂點 1 同色 \\
17 & 5 & 2 & 1 2 3 3 & 通過 \\
18 & 6 & 1 & 1 2 3 3 2 & 通過，得到塗法 $(1, 2, 3, 3, 2, 1)$ \\
\bottomrule
\end{tabular}
\end{center}

連根一共產生 $19$ 個節點。 步驟 7 把頂點 4 塗成 $1$，當下沒有衝突，但頂點 6 同時和頂點 3
（色 $3$）、4（色 $1$）、5（色 $2$）相鄰，三種顏色都被占走。 回溯一路退到頂點 4，改塗
$3$ 之後才走得通。 這個錯誤要到兩層之後才看得出來，界限函數只檢查已經塗色的鄰居，沒辦法
提早發現。

搜尋完整棵樹會找到 $6$ 種塗法（連根產生 $121$ 個節點），它們其實是同一種分組換顏色名
字：$\{1, 6\}$、$\{2, 5\}$、$\{3, 4\}$ 三組，三種顏色排列有 $3! = 6$ 種。 只想知道「能不能
塗」或想數「本質上不同」的塗法時，可以固定頂點 1 塗顏色 $1$，節點數降到 $41$，找到的 $2$
種塗法乘上 $3$ 就是全部（固定一個頂點只消掉 $3$ 倍；要消掉全部 $3!$ 倍，得再規定顏色第一
次出現的順序）。 $n$ 皇后也能這樣去掉對稱：解左右翻轉還是解，第 1 列只要試前
$\lceil n / 2 \rceil$ 欄，找到的解再翻轉一次（$n$ 是奇數時，第 1 列放正中間的解翻轉後還
是第 1 列放正中間，不能直接乘 $2$）。

\section{分支限界}
\label{sec:bk-bb}

\subsection{E-節點怎麼選：FIFO、LIFO 與最小成本}

回溯一產生子節點就往下走，原來的 E-節點要等整棵子樹走完才繼續。 \textbf{分支限界}
（branch and bound）換一種做法：E-節點一次把所有子節點都產生出來，沒被剪掉的子節點成為活節
點，放進一個容器；然後從容器裡取出下一個 E-節點。 容器怎麼取，決定了搜尋的順序：

\begin{itemize}
\item \textbf{FIFO 分支限界}：容器是佇列，先放進去的先取出，節點依層的順序展開，和廣度優
先搜尋一樣。
\item \textbf{LIFO 分支限界}：容器是堆疊，最後放進去的先取出。 它也往深處走，但每個
E-節點先產生全部的子節點，和回溯的深度優先不完全相同，有的課本叫它 D-search。
\item \textbf{最小成本分支限界}（least-cost，LC）：每個活節點有一個估計值 $\hat c(x)$，代表
「經過 $x$ 的答案，成本大約是多少」，每次取 $\hat c$ 最小的；容器是第~8~章的優先佇列。 求
最大值時改成取上界最大的，這種做法也叫\textbf{最佳優先搜尋}（best-first search）。
\end{itemize}

\textbf{可行性剪枝與界限剪枝。} 第~\ref{sec:bk-tree}~節的界限函數都在回答同一個問題：這個
部分解還有沒有可能補成一個合法的解？ $n$ 皇后檢查對角線、子集合和檢查會不會超過或湊不
到、著色檢查鄰居的顏色，都是這一類，叫做\textbf{可行性剪枝}（feasibility pruning）。 分支
限界通常用來解\textbf{最佳化問題}（要找最大或最小的解），這時還有第二種剪枝：搜尋時記住
目前找到的最好的解，如果能對一個節點算出一個\textbf{界限}，保證這棵子樹裡每一個解都不會比
目前最好的解更好，整棵子樹就可以丟掉，叫做\textbf{界限剪枝}（bounding）。 求最大值時用
\textbf{上界}，求最小值時用\textbf{下界}。 章首的兩種「不必再試」正好是這兩種：超重是可行
性剪枝，「剩下的全部都拿也比不上已經找到的組合」是界限剪枝。

界限要「樂觀」：求最大值的上界寧可估得太大，不能估得太小。 估得太小會把真正的最佳解剪
掉，答案就錯了；估得太大只是剪得少、跑得慢。 在不估錯的前提下，越接近真實值的界限剪得越
多，但通常也算得越久。

界限在分支限界裡有兩個用途：決定先展開哪個節點（只有 LC 用到），以及剪掉不可能比目前最好
的解更好的節點（三種都用）。 以求最小值為例，令 $\mathit{upper}$ 是目前找到的最好答案的
成本，$\hat c(x)$ 是 $x$ 子樹裡所有答案成本的下界；$\hat c(x) \ge \mathit{upper}$ 的節點不必
再展開。

\begin{pseudo}[label=alg:bk-lc]{最小成本分支限界（求最小值）}
\Function{LC-Search}{}
    \State $\mathit{upper} \gets \infty$；把根放進活節點集合 $L$
    \While{$L$ 不是空的}
        \State $E \gets L$ 裡 $\hat c$ 最小的節點（一樣小時取編號小的），從 $L$ 移除
        \If{$\hat c(E) \ge \mathit{upper}$}
            \State \Return $\mathit{upper}$ \Comment{剩下的節點都不可能更好}
        \EndIf
        \For{$E$ 的每一個子節點 $X$，依固定的順序產生並編號}
            \State 算出 $\hat c(X)$
            \If{$X$ 是答案節點}
                \If{$\hat c(X) < \mathit{upper}$}
                    \State $\mathit{upper} \gets \hat c(X)$，記下 $X$
                \EndIf
            \ElsIf{$\hat c(X) < \mathit{upper}$}
                \State 把 $X$ 放進 $L$
            \EndIf
        \EndFor
    \EndWhile
    \State \Return $\mathit{upper}$
\EndFunction
\end{pseudo}

第 4 行改成「取最早放進 $L$ 的節點」就是 FIFO，改成「取最晚放進 $L$ 的節點」就是 LIFO，
其他行都不變。 這個程序要正確，$\hat c$ 必須滿足兩件事：對每個節點，$\hat c(x)$ 不超過
$x$ 子樹裡任何答案的真正成本；在答案節點上，$\hat c$ 就是真正的成本。 這樣第 5 行取出的
節點 $\hat c$ 已經不小於 $\mathit{upper}$ 時，$L$ 裡每個節點的 $\hat c$ 都不小於
$\mathit{upper}$，它們的子樹裡不會有更便宜的答案，$\mathit{upper}$ 就是最佳解的成本。

\subsection{0/1 背包的最佳優先分支限界}
\label{sec:bk-knap}

$n$ 個物品，第 $i$ 個價值 $p_i$、重量 $w_i$，背包容量 $W$，要選一些物品使總重不超過 $W$、
總價值最大。 第~18、19~章把價值寫成 $v_i$；這一章照分支限界課本的習慣寫成 $p_i$（profit），
因為 $v$ 要留給樹上的節點。 狀態空間樹和子集合和一樣是滿二元樹，第 $k$ 層決定第 $k$ 個物品
拿或不拿。 每一個節點本身就是一個合法的解（後面的物品都不拿），所以只要沒超重，就可以拿
它的總價值去更新目前最好的值 $\mathit{best}$。

背包求最大值，要的是上界。 最省事的上界是「目前的總價值加上剩下所有物品的價值」，完全不
管容量，很鬆。 好得多的做法是把剩下的物品當成可以切開：先把物品依 $p_i / w_i$ 由大到小排
好，放得下就整個放，放不下就放進剛好填滿的那一部分。 這就是第~19~章分數背包的貪心解
（虛擬碼~\ref{alg:gr-frac}），加上目前的總價值 $P$ 當作上界：

\begin{pseudo}[label=alg:bk-bound, unbreakable]{分數背包上界}
\Function{Bound}{$k, P, \mathit{Wt}$} \Comment{已決定前 $k$ 個物品，$\mathit{Wt} \le W$}
    \State $b \gets P$；$\mathit{room} \gets W - \mathit{Wt}$；$i \gets k + 1$
    \While{$i \le n$ \AND $w_i \le \mathit{room}$}
        \State $\mathit{room} \gets \mathit{room} - w_i$；$b \gets b + p_i$；$i \gets i + 1$
    \EndWhile
    \If{$i \le n$}
        \State $b \gets b + p_i \cdot \mathit{room} / w_i$ \Comment{第 $i$ 個只放一部分}
    \EndIf
    \State \Return $b$
\EndFunction
\end{pseudo}

這個上界不會估得太小。 這棵子樹裡的任何一個解，都是在剩下的物品 $k + 1, \ldots, n$ 裡選一
個總重不超過 $\mathit{room} = W - \mathit{Wt}$ 的子集合 $T$，總價值是 $P + \sum_{i \in T} p_i$。
把 $T$ 看成「每個物品只拿 $0$ 或 $1$ 份」，它也是同一組物品、同一個容量下分數背包的一個
解；定理~\ref{thm:gr-frac} 說貪心解是分數背包的最佳解，所以 \proc{Bound} 回傳的值不小於
$P + \sum_{i \in T} p_i$。 每次算 \proc{Bound} 要 $O(n)$ 的時間，但它比「剩下的全拿」緊得
多，剪掉的節點也多得多。

分支限界每次展開上界最大的活節點（最佳優先）。 規則整理如下：

\begin{enumerate}
\item 物品依 $p_i / w_i$ 由大到小排好；$\mathit{best} \gets 0$（什麼都不拿也是一個解）。
\item 每個節點記錄層數 $k$（前 $k$ 個物品已決定）、總價值 $P$、總重 $\mathit{Wt}$ 和上界
$b$。 根是節點 1。
\item 從活節點裡取出 $b$ 最大的（一樣大時取編號小的）。 如果 $b \le \mathit{best}$，它是在放
進來之後才變得沒希望的，直接丟掉；否則它是 E-節點。
\item E-節點先產生「拿第 $k + 1$ 個物品」的子節點，再產生「不拿」的子節點，依序編號。 超重的子
節點直接死掉；沒超重的子節點先用 $P$ 更新 $\mathit{best}$，再算 $b$，$k + 1 < n$ 而且 $b >
\mathit{best}$ 才放進活節點。
\end{enumerate}

\begin{pseudo}[label=alg:bk-knap-bb]{0/1 背包的最佳優先分支限界}
\Function{Knapsack-BB}{}
    \State 物品依 $p_i / w_i$ 由大到小排好；$\mathit{best} \gets 0$
    \State 根 $\gets (k = 0, P = 0, \mathit{Wt} = 0)$，$\mathit{root}.b \gets \Call{Bound}{0, 0, 0}$，放進優先佇列 $Q$
    \While{$Q$ 不是空的}
        \State $v \gets Q$ 裡 $b$ 最大的節點（一樣大時取編號小的），從 $Q$ 移除
        \If{$v.b > \mathit{best}$} \Comment{否則 $v$ 已經沒希望，丟掉}
            \For{先「拿」再「不拿」第 $v.k + 1$ 個物品}
                \State 產生子節點 $u$，算出 $u.k$、$u.P$、$u.\mathit{Wt}$，並編號
                \If{$u.\mathit{Wt} \le W$} \Comment{超重的子節點直接死掉}
                    \If{$u.P > \mathit{best}$}
                        \State $\mathit{best} \gets u.P$，記下 $u$ 拿了哪些物品
                    \EndIf
                    \State $u.b \gets \Call{Bound}{u.k, u.P, u.\mathit{Wt}}$
                    \If{$u.k < n$ \AND $u.b > \mathit{best}$}
                        \State 把 $u$ 放進 $Q$
                    \EndIf
                \EndIf
            \EndFor
        \EndIf
    \EndWhile
    \State \Return $\mathit{best}$
\EndFunction
\end{pseudo}

\subsubsection{手算追蹤}

背包容量 $W = 13$，四個物品已經依價值密度排好：

\begin{center}
\begin{tabular}{ccccc}
\toprule
物品 & 1 & 2 & 3 & 4 \\
\midrule
價值 $p_i$ & 27 & 12 & 20 & 15 \\
重量 $w_i$ & 3 & 2 & 4 & 5 \\
$p_i / w_i$ & 9 & 6 & 5 & 3 \\
\bottomrule
\end{tabular}
\end{center}

根的上界：物品 1、2、3 整個放進去，價值 $59$、重量 $9$，剩下的容量 $4$ 放物品 4 的
$4/5$，價值 $15 \cdot 4/5 = 12$，上界 $71$。 下表每一列是一個 E-節點產生的兩個子節點，「之後的
$Q$」列出這一步做完時佇列裡的節點和它們的上界。

\begin{center}
\small
\begin{tabular}{clccccl}
\toprule
E-節點 & 子節點 & $P$ & $\mathit{Wt}$ & $b$ & $\mathit{best}$ & 結果 \\
\midrule
1 & 2：拿 1 & 27 & 3 & 71 & 27 & 活 \\
 & 3：不拿 1 & 0 & 0 & 47 & 27 & 活；$Q$：2(71)、3(47) \\
\midrule
2 & 4：拿 2 & 39 & 5 & 71 & 39 & 活 \\
 & 5：不拿 2 & 27 & 3 & 62 & 39 & 活；$Q$：4(71)、5(62)、3(47) \\
\midrule
4 & 6：拿 3 & 59 & 9 & 71 & 59 & 活 \\
 & 7：不拿 3 & 39 & 5 & 54 & 59 & 死（$54 \le 59$）；$Q$：6(71)、5(62)、3(47) \\
\midrule
6 & 8：拿 4 & 74 & 14 & -- & 59 & 超重 \\
 & 9：不拿 4 & 59 & 9 & 59 & 59 & 死（最後一層）；$Q$：5(62)、3(47) \\
\midrule
5 & 10：拿 3 & 47 & 7 & 62 & 59 & 活 \\
 & 11：不拿 3 & 27 & 3 & 42 & 59 & 死；$Q$：10(62)、3(47) \\
\midrule
10 & 12：拿 4 & 62 & 12 & 62 & 62 & 新的 $\mathit{best}$，最後一層；$Q$：3(47) \\
 & 13：不拿 4 & 47 & 7 & 47 & 62 & 死 \\
\midrule
3 & \multicolumn{6}{l}{取出時 $47 \le 62$，丟掉；$Q$ 空了，結束} \\
\bottomrule
\end{tabular}
\end{center}

\begin{center}
\begin{tikzpicture}[x=1.3cm,y=1.05cm]
\node[bt-kexp, label={[bt-num]left:1}, label={[bt-ord]above right:1}] (k1) at (7.0,0.0) {0,0\\[-1pt]71};
\node[bt-kexp, label={[bt-num]left:2}, label={[bt-ord]above right:2}] (k2) at (4.4,-1.4) {27,3\\[-1pt]71};
\node[bt-kdead, label={[bt-num]left:3}] (k3) at (9.6,-1.4) {0,0\\[-1pt]47};
\node[bt-kexp, label={[bt-num]left:4}, label={[bt-ord]above right:3}] (k4) at (2.0,-2.8) {39,5\\[-1pt]71};
\node[bt-kexp, label={[bt-num]left:5}, label={[bt-ord]above right:5}] (k5) at (6.8,-2.8) {27,3\\[-1pt]62};
\node[bt-kexp, label={[bt-num]left:6}, label={[bt-ord]above right:4}] (k6) at (0.8,-4.2) {59,9\\[-1pt]71};
\node[bt-kdead, label={[bt-num]left:7}] (k7) at (3.2,-4.2) {39,5\\[-1pt]54};
\node[bt-kdead, label={[bt-num]left:8}] (k8) at (0.0,-5.6) {74,14\\[-1pt]$\times$};
\node[bt-kdead, label={[bt-num]left:9}] (k9) at (1.6,-5.6) {59,9\\[-1pt]59};
\node[bt-kexp, label={[bt-num]left:10}, label={[bt-ord]above right:6}] (k10) at (5.6,-4.2) {47,7\\[-1pt]62};
\node[bt-kdead, label={[bt-num]left:11}] (k11) at (8.0,-4.2) {27,3\\[-1pt]42};
\node[bt-kbest, label={[bt-num]left:12}] (k12) at (4.8,-5.6) {62,12\\[-1pt]62};
\node[bt-kdead, label={[bt-num]left:13}] (k13) at (6.4,-5.6) {47,7\\[-1pt]47};
\draw[bt-e] (k1) -- node[bt-el] {1} (k2);
\draw[bt-e] (k1) -- node[bt-el] {0} (k3);
\draw[bt-e] (k2) -- node[bt-el] {1} (k4);
\draw[bt-e] (k2) -- node[bt-el] {0} (k5);
\draw[bt-e] (k4) -- node[bt-el] {1} (k6);
\draw[bt-e] (k4) -- node[bt-el] {0} (k7);
\draw[bt-e] (k6) -- node[bt-el] {1} (k8);
\draw[bt-e] (k6) -- node[bt-el] {0} (k9);
\draw[bt-e] (k5) -- node[bt-el] {1} (k10);
\draw[bt-e] (k5) -- node[bt-el] {0} (k11);
\draw[bt-e] (k10) -- node[bt-el] {1} (k12);
\draw[bt-e] (k10) -- node[bt-el] {0} (k13);
\end{tikzpicture}
\diagcaption{0/1 背包的最佳優先分支限界。 方框上一行是 $P, \mathit{Wt}$，下一行是上界
$b$；左邊的小字是產生的順序，藍色框右上角的粗體數字是它被選為 E-節點的順序。 邊上的 $1$、$0$
表示拿或不拿。 虛線是死節點，橘色是最後的最佳解。}
\label{fig:bk-knap}
\end{center}

這個過程有三個地方要看清楚。 第一，節點 5 比節點 6 早產生，卻比較晚展開：節點 6 的上界 $71$ 比
節點 5 的 $62$ 大，最佳優先搜尋不管節點在哪一層，只看上界。 第二，節點 9 的 $P = 59$ 等於
當時的 $\mathit{best}$，它已經是最後一層，不必放進佇列。 第三，節點 3 在第一步就放進佇
列，上界 $47$ 當時大於 $\mathit{best} = 27$；等到它被取出時 $\mathit{best}$ 已經是 $62$，
才在第 6 行被丟掉。 最佳解是物品 1、3、4，總重 $12$，總價值 $62$，一共產生 $13$ 個節點；整
棵狀態空間樹（滿二元樹）有 $2^5 - 1 = 31$ 個節點。 用第~18~章的動態規劃算這組資料，答案也是 $62$。

\begin{examnote}
背包的分支限界題，先檢查物品有沒有依 $p_i / w_i$ 由大到小排好，沒有的話先排，並寫出排序後
的編號。 本書的約定：先產生「拿」再產生「不拿」；超重的子節點也編號；上界一樣時取編號小
的；$b \le \mathit{best}$ 就剪（等於也剪，因為不可能更好）。 每個節點寫出 $P$、$\mathit{Wt}$、
$b$，並另外列出 E-節點的順序和 $\mathit{best}$ 的變化。 有些課本把背包改寫成求最小值（價值
取負號），每個節點同時算一個下界 $\hat c$ 和一個用貪心整數解得到的上界 $u$，用 $u$ 來更新
$\mathit{upper}$；樹的形狀一樣，但節點被剪掉的時機可能不同，題目指定哪一種就照哪一種寫。
\end{examnote}

\subsubsection{三種順序的比較}

把第 5 行改成 FIFO 或 LIFO，其他規則不變（子節點一樣先拿後不拿，LIFO 因此先展開「不拿」的子
節點），產生的節點數如下。 第二組資料有 $20$ 個物品，第 $i$ 個的重量是
$10 + (37 i \bmod 41)$，價值是重量加上 $(53 i \bmod 29) + 5$，容量取總重的一半 $295$；物品先
依 $p_i / w_i$ 由大到小重新排過，最佳解是 $586$。

\begin{center}
\begin{tabular}{lcccc}
\toprule
 & \multicolumn{2}{c}{4 個物品} & \multicolumn{2}{c}{20 個物品} \\
順序 & 產生的節點 & E-節點 & 產生的節點 & E-節點 \\
\midrule
FIFO & 15 & 7 & 2283 & 1141 \\
LIFO & 23 & 11 & 6175 & 3087 \\
最佳優先 & 13 & 6 & 37 & 18 \\
\bottomrule
\end{tabular}
\end{center}

FIFO 一層一層展開，要到很深的地方才找到好的 $\mathit{best}$，前面幾層的節點很多都剪不掉；
LIFO 在這裡先往「不拿」的方向衝，找到的 $\mathit{best}$ 很差，剪得更少。 最佳優先第一條走
的就是貪心的方向，$\mathit{best}$ 很快就接近最佳解。 代價是記憶體：最佳優先搜尋要把所有活
節點留在佇列裡，最壞情況和節點數同一個量級；回溯只需要一條從根到目前節點的路徑。

\subsubsection{C 程式}

活節點放在一個用 \texttt{realloc} 加倍的陣列裡，每次掃過一遍挑出上界最大的（一樣大時取編號
小的），挑出來的位置用最後一個節點補上。 每次挑選要 $O(\ell)$，$\ell$ 是當時的活節點數；活節點多的時
候改用第~8~章的堆積，挑選就只要 $O(\lg \ell)$。 每個節點用一個 \texttt{unsigned long} 的位元記錄拿了
哪些物品，不必另外存父節點。 輸入的物品要先依 $p_i / w_i$ 由大到小排好；物品資料讀入後不再
改，放在檔案範圍的陣列。 程式印出每一個 E-節點的 $P$、$\mathit{Wt}$、$b$，以及最後被丟掉的節點。

\begin{example}[label=ex:bk-knapsack-bb]{c}{(knapsack\_bb.c) 0/1 背包的最佳優先分支限界}
#include <stdio.h>
#include <stdlib.h>

typedef struct {
    int id, level, profit, weight;  // 編號、已決定幾個、價值、重量
    double bound;                   // 子樹的價值上界
    unsigned long take;             // 第 i 個位元：拿了第 i 個物品
} Node;

// 活節點放在一個可以變長的陣列裡；每次掃過一遍挑出上界最大的
typedef struct {
    Node *a;
    int size, cap;
} List;

static int add(List *q, Node v)
{
    if (q->size == q->cap) {
        int cap = q->cap ? 2 * q->cap : 16;
        Node *a = realloc(q->a, cap * sizeof *a);
        if (a == NULL) {
            return 0;
        }
        q->a = a;
        q->cap = cap;
    }
    q->a[q->size++] = v;
    return 1;
}

// 取出上界最大的節點；上界一樣大時取編號小的
static Node take_best(List *q)
{
    int b = 0;
    for (int i = 1; i < q->size; i++) {
        Node *x = &q->a[i], *y = &q->a[b];
        if (x->bound > y->bound
            || (x->bound == y->bound && x->id < y->id)) {
            b = i;
        }
    }
    Node v = q->a[b];
    q->a[b] = q->a[--q->size];      // 拿最後一個補洞
    return v;
}

static int n, cap, p[30], w[30];    // 物品已依 p/w 由大到小排好

// 分數背包上界
static double bound(int level, int profit, int weight)
{
    double b = profit;
    int room = cap - weight;
    for (int i = level; i < n; i++) {
        if (w[i] > room) {
            return b + (double)(p[i] * room) / w[i];
        }
        room -= w[i];
        b += p[i];
    }
    return b;
}

int main(void)
{
    if (scanf("%d %d", &n, &cap) != 2 || n < 1 || n > 30) {
        return 1;
    }
    for (int i = 0; i < n; i++) {
        if (scanf("%d %d", &p[i], &w[i]) != 2 || w[i] <= 0) {
            return 1;
        }
    }
    List q = {NULL, 0, 0};
    int count = 1, best = 0;
    unsigned long best_take = 0;
    int ok = add(&q, (Node){1, 0, 0, 0, bound(0, 0, 0), 0});
    while (ok && q.size > 0) {
        Node v = take_best(&q);
        if (v.bound <= best) {
            printf("drop %d: b=%g\n", v.id, v.bound);
            continue;
        }
        printf("expand %d: P=%d Wt=%d b=%g\n",
               v.id, v.profit, v.weight, v.bound);
        for (int t = 1; t >= 0 && ok; t--) {    // 先拿，再不拿
            int k = v.level;
            Node u = {++count, k + 1, v.profit + t * p[k],
                      v.weight + t * w[k], 0,
                      v.take | (unsigned long)t << k};
            if (u.weight > cap) {
                continue;                               // 超重
            }
            if (u.profit > best) {
                best = u.profit;
                best_take = u.take;
            }
            u.bound = bound(u.level, u.profit, u.weight);
            if (u.level < n && u.bound > best) {
                ok = add(&q, u);
            }
        }
    }
    printf("best %d, items:", best);
    for (int i = 0; i < n; i++) {
        if (best_take >> i & 1) {
            printf(" %d", i + 1);
        }
    }
    printf("\nnodes generated: %d\n", count);
    free(q.a);
    return ok ? 0 : 1;
}
\end{example}
\iosample{4 13 \\ 27 3 \\ 12 2 \\ 20 4 \\ 15 5}{expand 1: P=0 Wt=0 b=71 \\ expand 2: P=27 Wt=3 b=71 \\ expand 4: P=39 Wt=5 b=71 \\ expand 6: P=59 Wt=9 b=71 \\ expand 5: P=27 Wt=3 b=62 \\ expand 10: P=47 Wt=7 b=62 \\ drop 3: b=47 \\ best 62, items: 1 3 4 \\ nodes generated: 13}

E-節點的順序 1、2、4、6、5、10 和手算表一樣。 上界是 \texttt{double}，用 \texttt{\%g} 印，
整數就印成整數。 \texttt{bound} 裡 \texttt{p * room} 先用整數乘好再除，真正的上界是整數時算出
來的 \texttt{double} 也剛好是那個整數，和整數 \texttt{best} 比較不會因為誤差而判斷錯。

\subsection{旅行推銷員問題：縮減矩陣}
\label{sec:bk-tsp}

\textbf{旅行推銷員問題}（traveling salesperson problem，TSP）：有 $n$ 個城市，$C[i][j]$ 是從
城市 $i$ 到城市 $j$ 的成本（$C[i][j]$ 和 $C[j][i]$ 可以不同），要從城市 1 出發，每個城市恰好
經過一次，最後回到城市 1，總成本最小。 這樣的路線有 $(n - 1)!$ 條。 第~28~章會說明這個問
題是 NP 難的，目前沒有已知的多項式時間演算法；分支限界可以讓小規模的實例不必列舉全部路
線。

狀態空間樹的根是只含城市 1 的路徑；一個節點代表一條從 1 出發的路徑，它的子節點是在路徑後面
接上一個還沒去過的城市，依城市編號由小到大產生。 深度 $n - 1$ 的葉節點已經走過所有城市，接
上回到 1 的邊就是一條完整的路線。 $n = 5$ 時這棵樹有 $1 + 4 + 12 + 24 + 24 = 65$ 個節點。

\subsubsection{縮減矩陣與下界}

把成本矩陣的對角線設成 $\infty$（不能從一個城市走到自己）。 一條路線從每個城市恰好出發
一次、也恰好進入每個城市一次，所以它在矩陣的每一列恰好用到一格，每一行也恰好用到一格。
於是把某一列的每一格都減去同一個數 $m$，每一條路線的成本都恰好少 $m$，誰比較便宜不會
改變。

\textbf{縮減}（reduce）一個矩陣：先對每一列，減去該列的最小值；再對每一行，減去該行（列
縮減之後）的最小值。 全部是 $\infty$ 的列或行不動。 縮減後每一列、每一行都至少有一個
$0$，所有的格子都不是負數。 減掉的總量叫做\textbf{縮減量} $r$。 任何一條路線，原來的成本
等於 $r$ 加上它在縮減後矩陣上的成本，後者不小於 $0$，所以 $r$ 是每一條路線成本的下界，
也就是根的 $\hat c$。

\begin{pseudo}[label=alg:bk-reduce]{縮減矩陣}
\Function{Reduce}{$A$} \Comment{就地修改 $A$；$\infty - m = \infty$}
    \State $r \gets 0$
    \For{$i \gets 1$ \To $n$}
        \State $m \gets$ 第 $i$ 列的最小值
        \If{$0 < m < \infty$}
            \State 第 $i$ 列的每一格減去 $m$；$r \gets r + m$
        \EndIf
    \EndFor
    \For{$j \gets 1$ \To $n$}
        \State $m \gets$ 第 $j$ 行的最小值
        \If{$0 < m < \infty$}
            \State 第 $j$ 行的每一格減去 $m$；$r \gets r + m$
        \EndIf
    \EndFor
    \State \Return $r$
\EndFunction
\end{pseudo}

\textbf{子節點的下界。} 設節點 $S$ 的路徑停在城市 $i$，它的縮減矩陣是 $A$，下界是 $\hat c(S)$。
接上城市 $j$ 得到子節點 $T$：

\begin{pseudo}[label=alg:bk-tsp-child]{TSP 的子節點}
\Function{TSP-Child}{$A, \hat c(S), i, j$}
    \State $B \gets A$ 的複本
    \State $B$ 的第 $i$ 列、第 $j$ 行每一格都設成 $\infty$ \Comment{$i$ 的出邊、$j$ 的入邊已經決定}
    \State $B[j][1] \gets \infty$ \Comment{不能還沒走完就回到 1}
    \State $r \gets \Call{Reduce}{B}$
    \State \Return $B$ 和 $\hat c(T) = \hat c(S) + A[i][j] + r$
\EndFunction
\end{pseudo}

第 3 行：路線已經決定從 $i$ 走到 $j$，$i$ 不會再從別的邊出發，$j$ 也不會再從別的邊進入。
第 4 行：路徑 $1 \to \cdots \to j$ 如果接著走 $j \to 1$，會在還沒經過所有城市時就繞回起點。
$\hat c(T)$ 由三部分組成：父節點已經算進來的下界、這條邊在父節點矩陣上剩下的成本
$A[i][j]$、以及新矩陣再縮減一次的量 $r$。

\begin{whybox}[縮減矩陣給出的是下界，葉節點上剛好是路線成本]
設 $n \ge 3$。 對每個不是葉節點的節點 $X$（路徑 $P_X$ 停在城市 $i$，縮減矩陣 $A_X$），下面的
等式對每一條以 $P_X$ 開頭的路線 $T$ 都成立：
\[
\text{$T$ 的成本} = \hat c(X) + \sum_{e \in T \setminus P_X} A_X[e]. \tag{$*$}
\]
$A_X$ 的每一格都不是負數，所以 $\hat c(X)$ 不超過 $T$ 的成本。

根：路線在每一列、每一行各用一格，縮減時每減掉一個 $m$，每條路線的成本都少 $m$，($*$)
成立。 從 $X$ 到子節點 $Y$（接上 $j$）：以 $P_Y$ 開頭的路線一定用 $(i, j)$，它剩下的邊
$T \setminus P_Y$ 從 $j$ 和還沒去過的城市各出發一次，進入還沒去過的城市和城市 1 各一
次。 第 3、4 行設成 $\infty$ 的格子，這些邊都用不到（$Y$ 不是葉節點時，還有城市沒去，$j \to
1$ 不會出現在 $T$ 裡）。 $B$ 裡還有有限值的列和行，恰好是這些邊用到的列和行，各用一格，
所以縮減掉的 $r$ 對每一條這樣的路線都恰好少算一次，($*$) 對 $Y$ 也成立。

葉節點：設葉節點的父節點 $X$ 的路徑停在 $i$，只剩城市 $j$ 沒去。 $A_X$ 第 $j$ 列除了 $A_X[j][1]$
以外都是 $\infty$，縮減後這一格一定是 $0$；又 $A_X[i][1] = \infty$，所以 $X$ 的路線只有
$i \to j \to 1$ 這一條，($*$) 給出它的成本是 $\hat c(X) + A_X[i][j] + 0$。 葉節點的 $B$ 全部是
$\infty$，$r = 0$，$\hat c(\text{葉節點}) = \hat c(X) + A_X[i][j]$，正好等於這條路線的成本。
\end{whybox}

\subsubsection{手算追蹤：5 個城市}

成本矩陣如下（列是出發的城市，行是到達的城市）。 先做列縮減：五列的最小值是 $4$、$10$、
$3$、$10$、$9$，共 $36$。 減完之後，第 1、2、3 行已經有 $0$，第 4 行的最小值是 $8$，第 5 行
是 $6$，共 $14$。 根的下界 $\hat c = 36 + 14 = 50$，縮減後的矩陣是節點 1 的矩陣。

\begin{center}
\footnotesize
\setlength{\tabcolsep}{4pt}
\begin{tabular}[t]{c|ccccc}
\multicolumn{6}{c}{原始成本矩陣}\\
 & 1 & 2 & 3 & 4 & 5 \\\hline
1 & $\infty$ & 7 & 4 & 21 & 16 \\
2 & 20 & $\infty$ & 10 & 20 & 16 \\
3 & 11 & 3 & $\infty$ & 11 & 16 \\
4 & 10 & 24 & 13 & $\infty$ & 25 \\
5 & 9 & 16 & 24 & 24 & $\infty$ \\
\end{tabular}\qquad
\begin{tabular}[t]{c|ccccc}
\multicolumn{6}{c}{縮減後＝節點 1（$\hat c = 50$）}\\
 & 1 & 2 & 3 & 4 & 5 \\\hline
1 & $\infty$ & 3 & 0 & 9 & 6 \\
2 & 10 & $\infty$ & 0 & 2 & 0 \\
3 & 8 & 0 & $\infty$ & 0 & 7 \\
4 & 0 & 14 & 3 & $\infty$ & 9 \\
5 & 0 & 7 & 15 & 7 & $\infty$ \\
\end{tabular}
\end{center}

LC 搜尋的過程如下。 每一步從活節點裡取 $\hat c$ 最小的當 E-節點，產生它所有的子節點。 下表
列出每一個產生的子節點：$A[i][j]$ 是父節點矩陣上那一格，「再縮減」寫出子節點的矩陣減了哪幾
列、哪幾行。

\begin{center}
\small
\begin{tabular}{cclcclcc}
\toprule
節點 & 父 & 路徑 & $\hat c(\text{父})$ & $A[i][j]$ & 再縮減 & $r$ & $\hat c$ \\
\midrule
2 & 1 & 1$\to$2 & 50 & 3 & 不用減 & 0 & 53 \\
3 & 1 & 1$\to$3 & 50 & 0 & 不用減 & 0 & 50 \\
4 & 1 & 1$\to$4 & 50 & 9 & 列 4 減 3 & 3 & 62 \\
5 & 1 & 1$\to$5 & 50 & 6 & 列 5 減 7 & 7 & 63 \\
\midrule
6 & 3 & 1$\to$3$\to$2 & 50 & 0 & 行 4 減 2 & 2 & 52 \\
7 & 3 & 1$\to$3$\to$4 & 50 & 0 & 列 4 減 9、行 2 減 5 & 14 & 64 \\
8 & 3 & 1$\to$3$\to$5 & 50 & 7 & 列 2 減 2、列 5 減 7 & 9 & 66 \\
\midrule
9 & 6 & 1$\to$3$\to$2$\to$4 & 52 & 0 & 列 4 減 9 & 9 & 61 \\
10 & 6 & 1$\to$3$\to$2$\to$5 & 52 & 0 & 列 5 減 5 & 5 & 57 \\
\midrule
11 & 2 & 1$\to$2$\to$3 & 53 & 0 & 行 5 減 7 & 7 & 60 \\
12 & 2 & 1$\to$2$\to$4 & 53 & 2 & 列 3 減 7、列 4 減 3 & 10 & 65 \\
13 & 2 & 1$\to$2$\to$5 & 53 & 0 & 列 5 減 7、行 3 減 3 & 10 & 63 \\
\midrule
14 & 10 & 1$\to$3$\to$2$\to$5$\to$4 & 57 & 0 & 不用減（葉節點） & 0 & 57 \\
\bottomrule
\end{tabular}
\end{center}

E-節點依序是節點 1（$50$）、3（$50$）、6（$52$）、2（$53$）、10（$57$）。 展開節點 6 之後，活節
點裡最小的是節點 2 的 $53$，比節點 6 的子節點 10（$57$）小，搜尋跳回第 1 層。 節點 2 的三個子
節點都在 $60$ 以上，接著展開節點 10：它唯一的子節點 14 是葉節點，路線
$1 \to 3 \to 2 \to 5 \to 4 \to 1$ 的成本 $4 + 3 + 16 + 24 + 10 = 57$，$\mathit{upper} \gets 57$。 下
一個 $\hat c$ 最小的活節點是 11（$60$），$60 \ge 57$，停止。

最佳路線是 $1 \to 3 \to 2 \to 5 \to 4 \to 1$，成本 $57$。 一共產生 $14$ 個節點、展開 $5$ 個；
把 $24$ 條路線全部算一遍，最小值也是 $57$。 圖~\ref{fig:bk-tsp} 是整個過程的樹，
圖~\ref{fig:bk-tsp-mats} 是被展開的五個節點的縮減矩陣，算子節點的 $A[i][j]$ 和 $r$ 都從這
些矩陣來。

\begin{center}
\begin{tikzpicture}[x=1.4cm,y=1cm]
\node[bt-texp, label={[bt-cost]right:50}, label={[bt-ord]above right:1}] (t1) at (5.175,0.0) {1};
\node[bt-texp, label={[bt-cost]right:53}, label={[bt-ord]above right:4}] (t2) at (1.15,-1.25) {2};
\node[bt-texp, label={[bt-cost]right:50}, label={[bt-ord]above right:2}] (t3) at (5.462,-1.25) {3};
\node[bt-n, label={[bt-cost]right:62}] (t4) at (8.05,-1.25) {4};
\node[bt-n, label={[bt-cost]right:63}] (t5) at (9.2,-1.25) {5};
\node[bt-texp, label={[bt-cost]right:52}, label={[bt-ord]above right:3}] (t6) at (4.025,-2.5) {6};
\node[bt-n, label={[bt-cost]right:64}] (t7) at (5.75,-2.5) {7};
\node[bt-n, label={[bt-cost]right:66}] (t8) at (6.9,-2.5) {8};
\node[bt-n, label={[bt-cost]right:61}] (t9) at (3.45,-3.75) {9};
\node[bt-texp, label={[bt-cost]right:57}, label={[bt-ord]above right:5}] (t10) at (4.6,-3.75) {10};
\node[bt-n, label={[bt-cost]right:60}] (t11) at (0.0,-2.5) {11};
\node[bt-n, label={[bt-cost]right:65}] (t12) at (1.15,-2.5) {12};
\node[bt-n, label={[bt-cost]right:63}] (t13) at (2.3,-2.5) {13};
\node[bt-ans, label={[bt-cost]right:57}] (t14) at (4.6,-5.0) {14};
\draw[bt-e] (t1) -- node[bt-el] {2} (t2);
\draw[bt-e] (t1) -- node[bt-el] {3} (t3);
\draw[bt-e] (t1) -- node[bt-el] {4} (t4);
\draw[bt-e] (t1) -- node[bt-el] {5} (t5);
\draw[bt-e] (t3) -- node[bt-el] {2} (t6);
\draw[bt-e] (t3) -- node[bt-el] {4} (t7);
\draw[bt-e] (t3) -- node[bt-el] {5} (t8);
\draw[bt-e] (t6) -- node[bt-el] {4} (t9);
\draw[bt-e] (t6) -- node[bt-el] {5} (t10);
\draw[bt-e] (t2) -- node[bt-el] {3} (t11);
\draw[bt-e] (t2) -- node[bt-el] {4} (t12);
\draw[bt-e] (t2) -- node[bt-el] {5} (t13);
\draw[bt-e] (t10) -- node[bt-el] {4} (t14);
\end{tikzpicture}
\diagcaption{5 個城市的 LC 分支限界。 節點內是產生的順序，右邊是 $\hat c$，邊上是接上的城
市。 藍色是被展開的節點，右上角的粗體數字是展開的順序；橘色是找到的路線。}
\label{fig:bk-tsp}
\end{center}

\begin{center}
\footnotesize
\setlength{\tabcolsep}{3.5pt}
\begin{tabular}[t]{c|ccccc}
\multicolumn{6}{c}{節點 3（$\hat c = 50$）}\\
 & 1 & 2 & 3 & 4 & 5 \\\hline
1 & $\infty$ & $\infty$ & $\infty$ & $\infty$ & $\infty$ \\
2 & 10 & $\infty$ & $\infty$ & 2 & 0 \\
3 & $\infty$ & 0 & $\infty$ & 0 & 7 \\
4 & 0 & 14 & $\infty$ & $\infty$ & 9 \\
5 & 0 & 7 & $\infty$ & 7 & $\infty$ \\
\end{tabular}\qquad
\begin{tabular}[t]{c|ccccc}
\multicolumn{6}{c}{節點 6（$\hat c = 52$）}\\
 & 1 & 2 & 3 & 4 & 5 \\\hline
1 & $\infty$ & $\infty$ & $\infty$ & $\infty$ & $\infty$ \\
2 & $\infty$ & $\infty$ & $\infty$ & 0 & 0 \\
3 & $\infty$ & $\infty$ & $\infty$ & $\infty$ & $\infty$ \\
4 & 0 & $\infty$ & $\infty$ & $\infty$ & 9 \\
5 & 0 & $\infty$ & $\infty$ & 5 & $\infty$ \\
\end{tabular}

\medskip
\begin{tabular}[t]{c|ccccc}
\multicolumn{6}{c}{節點 2（$\hat c = 53$）}\\
 & 1 & 2 & 3 & 4 & 5 \\\hline
1 & $\infty$ & $\infty$ & $\infty$ & $\infty$ & $\infty$ \\
2 & $\infty$ & $\infty$ & 0 & 2 & 0 \\
3 & 8 & $\infty$ & $\infty$ & 0 & 7 \\
4 & 0 & $\infty$ & 3 & $\infty$ & 9 \\
5 & 0 & $\infty$ & 15 & 7 & $\infty$ \\
\end{tabular}\qquad
\begin{tabular}[t]{c|ccccc}
\multicolumn{6}{c}{節點 10（$\hat c = 57$）}\\
 & 1 & 2 & 3 & 4 & 5 \\\hline
1 & $\infty$ & $\infty$ & $\infty$ & $\infty$ & $\infty$ \\
2 & $\infty$ & $\infty$ & $\infty$ & $\infty$ & $\infty$ \\
3 & $\infty$ & $\infty$ & $\infty$ & $\infty$ & $\infty$ \\
4 & 0 & $\infty$ & $\infty$ & $\infty$ & $\infty$ \\
5 & $\infty$ & $\infty$ & $\infty$ & 0 & $\infty$ \\
\end{tabular}
\diagcaption{被展開的節點的縮減矩陣。 節點 3 是從節點 1 把第 1 列、第 3 行和 $A[3][1]$ 設成
$\infty$ 得到的，剛好不必再縮減。}
\label{fig:bk-tsp-mats}
\end{center}

\begin{examnote}
縮減矩陣的 TSP 題，本書的約定是：先縮減列再縮減行；從 $i$ 走到 $j$ 時把第 $i$ 列、第 $j$
行和 $A[j][1]$ 設成 $\infty$；子節點依城市編號由小到大產生，整棵樹依產生順序編號；選
$\hat c$ 最小的活節點，一樣小時選編號小的；找到葉節點就更新 $\mathit{upper}$，下一個 E-節點的
$\hat c$ 不小於 $\mathit{upper}$ 時停止。 作答時寫出每一個被展開節點的縮減矩陣，每個子節點寫
出 $A[i][j]$ 和 $r$，最後列出 E-節點的順序。 也有課本用別的下界（例如每個城市最便宜的兩
條邊），或先縮減行再縮減列，數字會不一樣，照題目的規定。
\end{examnote}

\section{本章常見錯誤}
\label{sec:bk-pitfalls}

\begin{center}
\begin{tabular}{>{\raggedright\arraybackslash}p{0.26\linewidth}>{\raggedright\arraybackslash}p{0.48\linewidth}>{\raggedright\arraybackslash}p{0.16\linewidth}}
\toprule
錯誤 & 例子 & 見 \\
\midrule
共用的陣列沒有撤銷 & $n$ 皇后的欄、對角線記號在遞迴回來後沒有設回 \FALSE，後面的分支
被看不見的皇后擋住 & 第~\ref{sec:bk-queens}~節 \\
沒排序就用「最小的放不進去」剪枝 & 子集合和的 $w$ 沒有由小到大排好，$s + w_{k+1} > M$
不代表後面的都放不進去，答案被剪掉 & 第~\ref{sec:bk-subset}~節 \\
界限估得太小 & 求最大值時拿貪心的整數解當上界；它是一個可行解的值，可能比子樹裡的最佳解
小，最佳解會被剪掉 & 第~\ref{sec:bk-knap}~節 \\
分數上界沒有先依 $p_i / w_i$ 排序 & 照輸入順序塞物品，算出來的值不是分數背包的最佳值，可
能比真正的最佳解小 & 第~\ref{sec:bk-knap}~節 \\
要找全部最佳解卻在「等於」時剪掉 & 只要一個最佳解時 $b \le \mathit{best}$ 可以剪；要列出所
有最佳解，只能在 $b < \mathit{best}$ 時剪 & 第~\ref{sec:bk-knap}~節 \\
TSP 子節點的 $\hat c$ 少加一項 & 只寫新矩陣的縮減量 $r$，忘了加父節點的 $\hat c$ 和
$A[i][j]$；或 $A[i][j]$ 拿原始矩陣的值，而不是父節點縮減後的值 & 第~\ref{sec:bk-tsp}~節 \\
TSP 忘了設 $\infty$ & 沒有把第 $i$ 列、第 $j$ 行設成 $\infty$，縮減時多減了路線用不到
的格子，算出來的 $\hat c$ 可能超過真正的成本 & 第~\ref{sec:bk-tsp}~節 \\
\bottomrule
\end{tabular}
\end{center}

\section{觀念題}
\label{sec:bk-exercises}

題型和研究所考試、期中期末考的筆試相近，有選擇、計算、追蹤和證明。 解答在本章最後。

\begin{enumerate}
\item （選擇，可能不只一個答案）下列哪些敘述正確？
\begin{enumerate}
\item 用了界限函數的回溯法，在最壞情況下產生的節點數一定比狀態空間樹的節點數少一個數量
級。
\item 回溯是以深度優先的順序走狀態空間樹，邊走邊用界限函數剪枝。
\item FIFO 分支限界用佇列存放活節點，節點展開的順序和廣度優先搜尋相同。
\item LC 分支限界中，只要 $\hat c$ 是下界，第一個被產生出來的答案節點一定是最佳解。
\item 同一個問題，最佳優先搜尋需要的記憶體通常比回溯多。
\end{enumerate}

\item （追蹤）用 \proc{Sum-Of-Subsets}（虛擬碼~\ref{alg:bk-subset}）求
$w = \langle 3, 5, 7, 11, 12, 14 \rangle$、$M = 19$ 的所有答案。 依產生的順序列出每個節點的
$(s, r)$，標出答案節點，並寫出一共產生幾個節點。

\item （證明）證明 \proc{Sum-Of-Subsets} 第 4 行和第 7 行不產生的子節點，子樹裡一定沒有答
案（$w$ 已由小到大排好）。 再舉一個 $w$ 沒有排序的例子，說明同一個程序會漏掉答案。

\item （計算）圖 $C_5$ 是五個頂點連成的一個環：邊是 $\{1,2\}$、$\{2,3\}$、$\{3,4\}$、
$\{4,5\}$、$\{5,1\}$。 取 $m = 3$。
\begin{enumerate}
\item \proc{Color}（虛擬碼~\ref{alg:bk-color}）找到的第一種塗法是什麼？
\item 一共有幾種塗法？ 固定頂點 1 塗顏色 $1$ 時有幾種？
\item 說明為什麼 $m = 2$ 時沒有塗法。
\end{enumerate}

\item （追蹤）用 \proc{Knapsack-BB}（虛擬碼~\ref{alg:bk-knap-bb}）解容量 $W = 16$ 的 0/1 背
包，物品（已依 $p_i / w_i$ 排好）的 $(p_i, w_i)$ 是 $(49, 7)$、$(18, 3)$、$(35, 7)$、
$(20, 5)$。 寫出每一個產生的節點的 $P$、$\mathit{Wt}$、$b$，E-節點的順序，以及最後的答
案。

\item （計算）4 個城市的成本矩陣如下（列是出發、行是到達）。
\[
\begin{array}{c|cccc}
 & 1 & 2 & 3 & 4 \\ \hline
1 & \infty & 13 & 7 & 4 \\
2 & 6 & \infty & 12 & 13 \\
3 & 11 & 11 & \infty & 2 \\
4 & 3 & 14 & 11 & \infty
\end{array}
\]
\begin{enumerate}
\item 縮減這個矩陣，寫出縮減後的矩陣和根的 $\hat c$。
\item 照本節的約定做 LC 分支限界，寫出每一個產生的節點的 $\hat c$、E-節點的順序，以及最佳
路線和它的成本。
\end{enumerate}

\item （追蹤）下面的樹每個節點標著名字和 $\hat c$，子節點由左到右產生。 不做任何剪枝，葉節點被
選到時也算一次 E-節點（它沒有子節點）。 分別寫出 FIFO、LIFO、LC（取 $\hat c$ 最小，一樣小時
取先產生的）三種分支限界的前 $7$ 個 E-節點。
\begin{center}
\begin{tikzpicture}[x=1.6cm,y=1cm,
    bt-sbox/.append style={font=\scriptsize, minimum width=1.0cm, minimum height=0.5cm}]
\node[bt-sbox] (a1) at (2.969,0.0) {A\,10};
\node[bt-sbox] (a2) at (0.475,-1.0) {B\,14};
\node[bt-sbox] (a3) at (3.087,-1.0) {C\,12};
\node[bt-sbox] (a4) at (5.462,-1.0) {D\,11};
\node[bt-sbox] (a5) at (0.0,-2.0) {E\,16};
\node[bt-sbox] (a6) at (0.95,-2.0) {F\,15};
\node[bt-sbox] (a7) at (2.375,-2.0) {G\,13};
\node[bt-sbox] (a8) at (3.8,-2.0) {H\,18};
\node[bt-sbox] (a9) at (4.75,-2.0) {I\,17};
\node[bt-sbox] (a10) at (6.175,-2.0) {J\,12};
\node[bt-sbox] (a11) at (1.9,-3.0) {K\,14};
\node[bt-sbox] (a12) at (2.85,-3.0) {L\,20};
\node[bt-sbox] (a13) at (5.7,-3.0) {M\,19};
\node[bt-sbox] (a14) at (6.65,-3.0) {N\,13};
\draw[bt-e] (a1) -- (a2);
\draw[bt-e] (a1) -- (a3);
\draw[bt-e] (a1) -- (a4);
\draw[bt-e] (a2) -- (a5);
\draw[bt-e] (a2) -- (a6);
\draw[bt-e] (a3) -- (a7);
\draw[bt-e] (a3) -- (a8);
\draw[bt-e] (a4) -- (a9);
\draw[bt-e] (a4) -- (a10);
\draw[bt-e] (a7) -- (a11);
\draw[bt-e] (a7) -- (a12);
\draw[bt-e] (a10) -- (a13);
\draw[bt-e] (a10) -- (a14);
\end{tikzpicture}
\end{center}

\item （設計）\textbf{哈密頓環}問題：給一個無向圖，問有沒有一個環恰好經過每個頂點一次。 用
回溯寫出找出所有哈密頓環的虛擬碼（固定從頂點 1 出發），說明顯式限制、隱式限制和你的界限
函數。 每個環會被找到幾次？ 怎麼讓每個環只出現一次？
\end{enumerate}

\section{歷屆試題}
\label{sec:bk-pastexam}

以下是研究所入學考試的原題，題目照原卷（語言、記號、配分照錄）。 詳解在本章最後。

\begin{pastexam}[label={px:bk-1}]{臺大資工所 112 學年度 第 9 題}
［本卷說明：本題為複選題，5 分，每個選項單獨計分。］

\textbf{DFS advantages:} To find a feasible solution to the eight-queen problem, what is/are the reason(s) that we prefer to use DFS (depth-first search) instead of BFS (breadth-first search)?
\begin{enumerate}[label=(\Alph*)]
\item DFS is more memory efficient.
\item DFS can be implemented using a stack.
\item DFS usually reaches the terminal states of "good" or "bad" faster.
\item DFS is conceptually simpler than BFS.
\item BFS cannot be implemented using a stack.
\end{enumerate}
\end{pastexam}

\begin{pastexam}[label={px:bk-2}]{中山資工所 111 學年度 第 1 題（節錄 (b) 小題）}
What are printed by each of the following C program?
\begin{enumerate}[label=(\alph*),start=2]
\item (10\%)
\begin{minted}[fontsize=\small, style=perldoc]{cpp}
void g(char a[ ], int k, int m)
{
    if (k == m) {
        for (int i = 0; i <= m; i++) cout << a[i] << " ";
            cout << endl;
    }
    else
        for (int i = k; i <= m; i++) {
           swap(a[k], a[i]); // exchange
           g(a, k+1, m);
           swap(a[k], a[i]);
        }
}
int main( )
{
    char a[ ]="abcde";
    g(a,0,2);
}
\end{minted}
\end{enumerate}
\end{pastexam}

\begin{pastexam}[label={px:bk-3}]{臺大生機所 111 學年度 第 2 題}
(15\%) Given a graph, write a pseudocode function that colors the graph's vertices such that no two adjacent vertices share the same color by minimizing the total number of colors used.
\end{pastexam}

\begin{pastexam}[label={px:bk-4}]{中興資工所 115 學年度 第 7 題}
［本卷說明：單選題，每題 2 分。］

Branch and Bound differs from backtracking by
\begin{enumerate}[label=\Alph*.,nosep]
\item Using bounds to prune search space
\item Random selection
\item Depth-only search
\item Greedy choices
\end{enumerate}
\end{pastexam}

\begin{pastexam}[label={px:bk-5}]{中央資工所 100 學年度 第 6 題}
Below are the inputs of an instance of the 0/1 knapsack problem with n=5 items and the knapsack capacity M=38. Note that P$_\mathrm{i}$ is the profit of item i and W$_\mathrm{i}$ is the weight of item i, i=1,\ldots,5. The 0/1 knapsack problem is to maximize the total profit of items that are wholly put into the knapsack under the constraint that their total weight should not exceed the knapsack capacity.
\begin{center}
\begin{tabular}{cccccc}
i & 1 & 2 & 3 & 4 & 5 \\
P$_\mathrm{i}$ & 6 & 10 & 4 & 5 & 6 \\
W$_\mathrm{i}$ & 10 & 19 & 8 & 10 & 12 \\
\end{tabular}\\[2pt]
(P$_\mathrm{i}$/W$_\mathrm{i}\ge$ P$_\mathrm{i+1}$/W$_\mathrm{i+1}$)
\end{center}
The problem can be solved by the branch-and-bound strategy on the basis of the best-first search scheme by deriving the lower bound and upper bound for each node in the search tree.
\begin{enumerate}[itemsep=2pt,topsep=2pt]
\item[6.1] (12\%) Derive the upper bound and the lower bound for the root node and explain how you obtain the bounds.
\item[6.2] (6\%) Explain how the bounds can be used in the branch-and-bound strategy. (6\%)
\end{enumerate}
\end{pastexam}

\section{實作練習}
\label{sec:bk-practice}

下面是 Aki-Judge 上和本章相關的題目。 這幾題原本是給《C 程式設計基礎》遞迴章練習用的，資
料量都很小，重點在把狀態空間樹想清楚、把剪枝的條件寫對，正好用來練這一章的技巧。

\begin{center}
\begin{tabular}{llc}
\toprule
小節 & 題目 & 難度 \\
\midrule
排列樹與限制 & \akiref{presentation-orders-within-time-windows}{限定時段的報告順序} & 暖身 \\
 & \akiref{palindrome-free-letter-arrangements}{不含回文的字母卡排列} & 基本 \\
圖著色 & \akiref{exam-slots-with-room-limits}{有容量上限的考試時段安排} & 基本 \\
組合與剪枝 & \akiref{project-team-with-skill-quotas}{技能有人數門檻的專題小組} & 基本 \\
 & \akiref{kit-combinations-that-use-everything}{剛好用完的零件包組合} & 基本 \\
分組 & \akiref{teams-balanced-on-two-skills}{兩種能力都要顧的分組} & 基本 \\
\bottomrule
\end{tabular}
\end{center}

幾題的提示：「限定時段的報告順序」依時段一個一個決定由哪一組上台，時段不在那一組範圍內的
就不是子節點；組別由小到大試，輸出自然是字典序。 「不含回文的字母卡排列」用每個字母還剩幾
張的計數表挑字母，相同的字母在同一層只試一次，就不會產生重複的排列；只要檢查新放的字母
和前一格、前兩格都不同，因為任何一段回文的正中間一定是長度 $2$ 或 $3$ 的回文。 「有容量上
限的考試時段安排」就是第~\ref{sec:bk-coloring}~節的 $m$ 著色，每個時段多一個容量的檢查；
時段由小到大試，第一個找到的就是字典序最小的。 「技能有人數門檻的專題小組」把人數由小到大
試，固定人數時依編號遞增挑人，合不來的兩人當作可行性剪枝。 「剛好用完的零件包組合」和子
集合和很像，但每種零件包可以拿很多次：規定只能用編號不小於上一包的零件包，同一個組合就不
會算很多次；任何一種零件的需求變成負數就剪掉。 這題規定不能寫迴圈，「試下一種零件包」也
要寫成遞迴。 「兩種能力都要顧的分組」讓每位同學加入已經開的組或開一個新組（一次只能開下
一個編號的組，消掉組的編號互換造成的對稱），剩下的人數不夠把組數補滿時就剪掉。

\section{讀完你應該會}
\label{sec:bk-checklist}

\begin{itemize}
\item 能替一個問題寫出解向量、顯式限制和隱式限制，畫出它的狀態空間樹，說出活節點、E-節點、
死節點的意思。
\item 能依產生的順序畫出 4 皇后、子集合和、圖著色的回溯過程，標出被剪掉的節點和答案節點。
\item 能分辨可行性剪枝和界限剪枝，對 0/1 背包寫出分數背包上界，並說明它為什麼不會估得太小。
\item 能用 FIFO、LIFO、LC 三種順序做分支限界，手算 0/1 背包的最佳優先樹，列出 E-節點的順序。
\item 能縮減 TSP 的成本矩陣，算出每個子節點的 $\hat c$，完成一次 LC 分支限界的手算追蹤。
\end{itemize}

\textbf{延伸閱讀}（書末會整理完整的參考資料）

\begin{itemize}
\item Ellis Horowitz、Sartaj Sahni、Sanguthevar Rajasekaran，\textit{Computer
Algorithms}：回溯與分支限界各一章，狀態空間樹的名詞、子集合和的兩個剪枝條件、縮減矩陣的
TSP 都在這兩章。
\item Richard E. Neapolitan，\textit{Foundations of Algorithms}，第 5 版，第 5、6 章：用很多
小例子講回溯和分支限界，0/1 背包的最佳優先搜尋寫得很仔細。
\item Steven S. Skiena，\textit{The Algorithm Design Manual}，第 3 版，第 9 章：從實作的角度
講回溯和剪枝，例子偏向實務。
\item John D. C. Little、Katta G. Murty、Dura W. Sweeney、Caroline Karel，``An Algorithm for the
Traveling Salesman Problem''，\textit{Operations Research} 11(6)，1963：用縮減矩陣做 TSP 分支
限界的原始論文。
\end{itemize}

\addcontentsline{toc}{section}{觀念題解答}
\section*{觀念題解答}

\begin{enumerate}
\item (b)、(c)、(e) 正確。 (a)：界限函數能剪掉多少要看輸入。 例如 $m$ 著色時，如果圖沒有任何邊，每個節點都通過檢查，產生的節點就是整棵狀態空間樹。 (d)：第一個產生的答案節點只代表找到一個上界；
例如它的成本是 $70$，而另一個活節點的 $\hat c = 60$，底下可能有成本 $62$ 的答案。 要等到下
一個 E-節點的 $\hat c$ 不小於 $\mathit{upper}$ 才能停。 (e)：最佳優先搜尋要把活節點都留在佇
列裡，回溯只需要一條路徑。

\item 總和 $52$。 一共產生 $20$ 個節點，三個答案節點依序是 $\{3, 5, 11\}$（節點 5）、
$\{5, 14\}$（節點 14）、$\{7, 12\}$（節點 18）。 左子樹（選 $3$）：節點 2 $(3, 49)$~$\to$~3 $(8, 44)$
$\to$~不選 $7$ 的 4 $(8, 37)$~$\to$~選 $11$ 得答案節點 5；節點 2 不選 $5$ 的 6 $(3, 44)$~$\to$~7 $(3, 37)$~$\to$~8
$(3, 26)$ 是死路（選 $12$ 後再加 $14$ 超過，不選 $12$ 時 $3 + 26 - 12 = 17 < 19$）。 右子樹（不選
$3$）：9 $(0, 49)$~$\to$~10 $(5, 44)$~$\to$~11 $(5, 37)$~$\to$~12 $(5, 26)$~$\to$~13 $(5, 14)$~$\to$~答案節點 14；9 的
右子節點 15 $(0, 44)$~$\to$~16 $(7, 37)$~$\to$~17 $(7, 26)$~$\to$~答案節點 18；15 的右子節點 19 $(0, 37)$~$\to$~20
$(0, 26)$ 是死路。 用範例~\ref{ex:bk-subsets} 跑輸入 \texttt{6 19} 和 \texttt{3 5 7 11 12 14}，
印出同樣三組答案和 \texttt{solutions = 3, nodes = 20}。

\item 在決定 $w_k$ 的節點，已選的和是 $s$，剩下 $w_k, \ldots, w_n$ 的和是 $r$，而且
$w_k \le w_{k+1} \le \cdots \le w_n$。
\begin{itemize}
\item 第 4 行不產生左子節點，是因為 $s + w_k \ne M$ 且 $s + w_k + w_{k+1} > M$。 左子節點的子
樹裡，每個答案都已經選了 $w_k$，和是 $s + w_k$ 再加上 $w_{k+1}, \ldots, w_n$ 裡選的數。 什
麼都不選，和是 $s + w_k \ne M$；至少選一個，選的數都不小於 $w_{k+1}$，和至少
$s + w_k + w_{k+1} > M$。 兩種都不是答案。
\item 第 7 行不產生右子節點，是因為 $s + r - w_k < M$ 或 $s + w_{k+1} > M$。 右子節點的子樹裡，
每個答案的和是 $s$ 加上 $w_{k+1}, \ldots, w_n$ 裡選的數。 前者表示全選也只有 $s + r - w_k <
M$。 後者表示只要選一個就超過 $M$；一個都不選，和是 $s$，而能走到這個節點代表
$s + w_k \le M$，$s < M$，也不是答案。
\end{itemize}
沒排序的例子：$w = \langle 4, 6, 1 \rangle$、$M = 5$，答案是 $\{4, 1\}$。 第一次呼叫
$(s, k, r) = (0, 1, 11)$：$0 + 4 \ne 5$，$0 + 4 + 6 = 10 > 5$，左子節點不產生；右子節點要
$0 + 6 \le 5$，也不產生。 程序什麼都沒找到。

\item
\begin{enumerate}
\item 頂點 1 塗 $1$；頂點 2 不能是 $1$，塗 $2$；頂點 3 的鄰居只有 2 塗過，塗 $1$；頂點 4 不
能和 3 同色，塗 $2$；頂點 5 的鄰居 4、1 分別是 $2$、$1$，塗 $3$。 第一種塗法是
$(1, 2, 1, 2, 3)$，沒有退回過。
\item 一共 $30$ 種；固定頂點 1 塗 $1$ 時 $10$ 種。 數法：固定頂點 1 塗 $1$，頂點 2、3、4 依
序只要和前一個不同，各有 $2$ 種，共 $8$ 種：$(2,1,2)$、$(2,1,3)$、$(2,3,1)$、$(2,3,2)$、
$(3,1,2)$、$(3,1,3)$、$(3,2,1)$、$(3,2,3)$。 頂點 5 要和頂點 4、頂點 1 都不同：頂點 4 塗 $1$
的 $2$ 種，頂點 5 有 $2$ 種選擇；其他 $6$ 種，頂點 5 只剩 $1$ 種。 共 $2 \cdot 2 + 6 \cdot 1 = 10$
種，頂點 1 有 $3$ 種顏色，全部是 $30$ 種。
\item 兩種顏色時，沿著環走，顏色必須 $1, 2, 1, 2, \ldots$ 交替，第 5 個頂點是 $1$，和頂點 1
相鄰卻同色。 長度是奇數的環都不能用兩種顏色塗。
\end{enumerate}

\item 根的上界：$49 + 18 = 67$（重 $10$），剩 $6$ 放物品 3 的 $6/7$，$35 \cdot 6 / 7 = 30$，
共 $97$。 各節點的 $(P, \mathit{Wt}, b)$：2 $(49, 7, 97)$、3 $(0, 0, 73)$；展開 2 得 4 $(67, 10, 97)$、
5 $(49, 7, 92)$；展開 4 得 6（拿物品 3，重 $17$，超重）、7 $(67, 10, 87)$；展開 5 得 8 $(84, 14,
92)$、9 $(49, 7, 69)$ 死；展開 8 得 10（重 $19$，超重）、11 $(84, 14, 84)$ 死；展開 7 得 12 $(87,
15, 87)$，$\mathit{best} = 87$，13 $(67, 10, 67)$ 死；最後取出節點 3 時 $73 \le 87$，丟掉。 E-節點
的順序是 1、2、4、5、8、7：節點 5、8 的上界 $92$ 都比節點 7 的 $87$ 大，所以先展開。 最佳解
是物品 1、2、4，總重 $15$，價值 $87$，共產生 $13$ 個節點；把範例~\ref{ex:bk-knapsack-bb} 的輸入
換成這組資料，程式印出同樣的 E-節點順序。

\item
\begin{enumerate}
\item 列縮減減去 $4$、$6$、$2$、$3$，共 $15$；之後第 2 行最小 $9$、第 3 行最小 $3$，第 1、4 行
已有 $0$，共 $12$。 根的 $\hat c = 27$，縮減後的矩陣：
\[
\begin{array}{c|cccc}
 & 1 & 2 & 3 & 4 \\ \hline
1 & \infty & 0 & 0 & 0 \\
2 & 0 & \infty & 3 & 7 \\
3 & 9 & 0 & \infty & 0 \\
4 & 0 & 2 & 5 & \infty
\end{array}
\]
\item 節點 1 的子節點：節點 2（$1 \to 2$）：$A[1][2] = 0$，第 2 列減 $3$，$\hat c = 27 + 0 + 3 =
30$；節點 3（$1 \to 3$）：$A[1][3] = 0$，不用減，$\hat c = 27$；節點 4（$1 \to 4$）：
$A[1][4] = 0$，第 4 列減 $2$、第 3 行減 $3$，$\hat c = 32$。 展開節點 3（$27$）：節點 5
（$1 \to 3 \to 2$）：$A[3][2] = 0$，第 2 列減 $7$，$\hat c = 34$；節點 6（$1 \to 3 \to 4$）：
$A[3][4] = 0$，第 4 列減 $2$，$\hat c = 29$。 展開節點 6（$29$）：唯一的子節點 7 是葉節點
$1 \to 3 \to 4 \to 2 \to 1$，$\hat c = 29$，$\mathit{upper} \gets 29$。 下一個最小的活節點是
節點 2（$30$），$30 \ge 29$，停止。 E-節點的順序是 1、3、6，最佳路線
$1 \to 3 \to 4 \to 2 \to 1$，成本 $7 + 2 + 14 + 6 = 29$。 根的下界 $27$ 比答案小 $2$，下界不一定
是緊的。
\end{enumerate}

\item FIFO：A、B、C、D、E、F、G。 LIFO：A、D、J、N、M、I、C（A 產生 B、C、D 後取最後放入
的 D；D 產生 I、J 後取 J；J 的子節點 M、N 取 N，再取 M，再取 I，然後才輪到 C）。 LC：A、D、
C、J、G、N、B。 LC 的過程：A 之後活節點是 B(14)、C(12)、D(11)，取 D；加入 I(17)、J(12)，
C 和 J 都是 $12$，C 先產生，取 C；加入 G(13)、H(18)，取 J(12)；加入 M(19)、N(13)，G 和 N 都是
$13$，取先產生的 G；加入 K(14)、L(20)，取 N(13)；剩下最小的是 B 和 K，都是 $14$，取先產生的
B。

\item 解向量 $(x_1, \ldots, x_n)$ 是頂點的走訪順序，$x_1 = 1$。 顯式限制：$x_k \in \{2, \ldots,
n\}$。 隱式限制：$x_2, \ldots, x_n$ 互不相同；$\{x_{k-1}, x_k\}$ 是邊；$\{x_n, x_1\}$ 也是邊。
界限函數在放 $x_k = v$ 時檢查 $v$ 還沒用過、和 $x_{k-1}$ 相鄰；$k = n$ 時另外檢查 $v$ 和頂點
1 相鄰。
寫成和 \proc{Queens} 一樣的形狀。 先設 $x_1 \gets 1$、把頂點 1 標成用過，再呼叫
\proc{Hamiltonian}$(2)$：
\begin{pseudo}[unbreakable]{觀念題 8 的解答}
\Procedure{Hamiltonian}{$k$}
    \For{$v \gets 2$ \To $n$}
        \If{$\mathit{used}[v] = \FALSE$ \AND $\{x_{k-1}, v\} \in E$}
            \State $x_k \gets v$；$\mathit{used}[v] \gets \TRUE$
            \If{$k < n$}
                \State \Call{Hamiltonian}{$k + 1$}
            \ElsIf{$\{v, 1\} \in E$ \AND $x_2 < x_n$}
                \State 輸出 $(x_1, \ldots, x_n)$
            \EndIf
            \State $\mathit{used}[v] \gets \FALSE$ \Comment{撤銷}
        \EndIf
    \EndFor
\EndProcedure
\end{pseudo}
第 3 行就是界限函數。 固定從頂點 1 出發之後，同一個環還會以兩個方向各被找到一次
（$1 \to a \to \cdots \to b \to 1$ 和 $1 \to b \to \cdots \to a \to 1$）；第 7 行多檢查
$x_2 < x_n$，每個環就只出現一次。 還可以加一個界限：每個還沒用過的頂點 $u$ 在環上需要一個前
一站和一個下一站，前一站只能是還沒用過的頂點或目前的尾端 $x_k$，下一站只能是還沒用過的頂點
或頂點 1；所以 $u$ 在「還沒用過的頂點、$x_k$、頂點 1」之中的鄰居如果少於兩個，就可以剪掉。
\end{enumerate}

\addcontentsline{toc}{section}{歷屆試題詳解}
\section*{歷屆試題詳解}

\pastsol{px:bk-1}{臺大資工所 112 學年度 第 9 題}
本書的答案是 \textbf{(A)、(C)}。 題目問的是「偏好 DFS 的理由」，所以一個選項要算對，必須
是正確的敘述，而且真的是 DFS 在八皇后上勝過 BFS 的原因。

先用數字看兩種走法的差別。 狀態空間樹用第~\ref{sec:bk-queens}~節的排列樹，對角線衝突的節點
一產生就剪掉，節點照產生的順序數（被剪掉的也算）。 八皇后的答案全部在第 $8$ 層（放完 $8$ 個皇
后）。 用程式數出各層通過檢查的節點數：

\begin{center}
\begin{tabular}{l*{8}{r}}
\toprule
層 & 1 & 2 & 3 & 4 & 5 & 6 & 7 & 8 \\
\midrule
通過檢查的節點 & 8 & 42 & 140 & 344 & 568 & 550 & 312 & 92 \\
\bottomrule
\end{tabular}
\end{center}

\begin{itemize}
\item DFS（回溯，範例~\ref{ex:bk-nqueens}）：第一組解是第 $337$ 個產生的節點。 任何時刻只需
要記住從根到目前節點的一條路徑，最多 $8$ 層。
\item BFS（FIFO，同樣的剪枝）：第 $8$ 層的節點要等第 $1$ 到 $6$ 層的節點全部展開、第 $7$ 層
也開始展開之後才會出現。 第一組解是第 $5209$ 個產生的節點，佇列裡同時存放的活節點最多
$573$ 個。
\end{itemize}

逐項看：
\begin{itemize}
\item (A) 對。 DFS 只存一條路徑（深度 $8$ 的遞迴，加上欄和對角線的記號陣列）；BFS 要把整層的
活節點放在佇列裡，上面算出最多 $573$ 個。 一般來說，深度 $d$、每個節點最多 $b$ 個子節點的樹，
DFS 要 $O(d)$ 個節點的空間，BFS 最壞要 $O(b^d)$ 個。
\item (B) 敘述本身是對的（第~\ref{sec:sq-recursion}~節：遞迴就是用堆疊實作），但不是偏好 DFS 的理
由：BFS 也一樣能用佇列實作，用哪一種容器本身沒有好壞。
\item (C) 對。 八皇后的解都在第 $8$ 層，DFS 一路往下走，很快就碰到「放滿 $8$ 個」或「這條路走
不下去」的節點，找到第一組解只產生 $337$ 個節點；BFS 要先把淺層全部走完，產生 $5209$ 個節點
才碰到第一組解。
\item (D) 錯。 「概念比較簡單」是主觀的說法，兩者一樣簡單，不是選 DFS 的理由。
\item (E) 錯。 BFS 需要先進先出的佇列，而佇列可以用兩個堆疊做出來：放入時推進第一個堆疊；取
出時，如果第二個堆疊是空的，先把第一個堆疊的元素全部彈出、依序推進第二個堆疊，再從第二個堆
疊彈出。 所以 BFS 也能用堆疊實作；何況這也不是偏好 DFS 的理由。
\end{itemize}

\pastsol{px:bk-2}{中山資工所 111 學年度 第 1 題（節錄 (b) 小題）}
這是《C 程式設計基礎》第~15~章 15.15 節「用交換產生所有排列」的寫法（C++ 的 \texttt{swap}
交換兩個字元，\texttt{cout} 輸出）。 \texttt{g(a, k, m)} 負責決定 \texttt{a[k]}：依序把
\texttt{a[k..m]} 的每一個字元換到第 \texttt{k} 格，遞迴決定後面的格子，回來後再換回去
（撤銷）。 \texttt{m = 2}，所以只排列前三個字元 \texttt{a}、\texttt{b}、\texttt{c}，
\texttt{d}、\texttt{e} 不會被動到，也不會印出。 另外注意縮排：\texttt{cout << endl;} 不在
\texttt{for} 裡面（那個 \texttt{for} 只有一個敘述），所以每個排列印完才換一次行。

遞迴的過程（第~\ref{sec:bk-tree}~節的狀態空間樹，深度 $k$ 決定 \texttt{a[k]}）：

\begin{center}
\begin{tabular}{lll}
\toprule
呼叫 & 這一層的交換 & 結果 \\
\midrule
\texttt{g(0)}，陣列 \texttt{abc} & $i = 0$：不變，\texttt{abc} & 進入 \texttt{g(1)} \\
\quad \texttt{g(1)}，\texttt{abc} & $i = 1$：不變 & \texttt{g(2)} 印 \texttt{a b c} \\
 & $i = 2$：交換 \texttt{a[1]}、\texttt{a[2]}，\texttt{acb} & \texttt{g(2)} 印 \texttt{a c b}，換回 \texttt{abc} \\
\texttt{g(0)} & $i = 1$：交換 \texttt{a[0]}、\texttt{a[1]}，\texttt{bac} & 進入 \texttt{g(1)} \\
\quad \texttt{g(1)}，\texttt{bac} & $i = 1$：不變 & 印 \texttt{b a c} \\
 & $i = 2$：\texttt{bca} & 印 \texttt{b c a}，換回 \texttt{bac} \\
\texttt{g(0)} & 換回 \texttt{abc}；$i = 2$：交換 \texttt{a[0]}、\texttt{a[2]}，\texttt{cba} & 進入 \texttt{g(1)} \\
\quad \texttt{g(1)}，\texttt{cba} & $i = 1$：不變 & 印 \texttt{c b a} \\
 & $i = 2$：\texttt{cab} & 印 \texttt{c a b}，換回 \texttt{cba} \\
\texttt{g(0)} & 換回 \texttt{abc} & 結束 \\
\bottomrule
\end{tabular}
\end{center}

輸出六行（每個字元後面都有一個空白）：

\noindent\begin{minipage}{\linewidth}
\begin{verbatim}
a b c
a c b
b a c
b c a
c b a
c a b
\end{verbatim}
\end{minipage}
最後兩行是 \texttt{c b a} 在 \texttt{c a b} 前面，不是字典序：第三輪把 \texttt{a[0]} 和
\texttt{a[2]} 交換，得到的是 \texttt{cba}，第 1、2 格的順序也跟著變了。 一共呼叫 \texttt{g}
$1 + 3 + 6 = 10$ 次，正是 $3$ 個元素的排列樹的節點數。

\pastsol{px:bk-3}{臺大生機所 111 學年度 第 2 題}
用的顏色越少越好，就是求圖的色數（第~\ref{sec:bk-coloring}~節）。 做法：依序試
$m = 1, 2, 3, \ldots$，每個 $m$ 用回溯檢查「能不能用 $m$ 種顏色塗」，第一個成功的 $m$ 就是最少
的顏色數。 回溯就是虛擬碼~\ref{alg:bk-color}，改成找到第一種塗法就停。 頂點編號 $1$ 到 $n$
（$n \ge 1$），$x_i$ 是頂點 $i$ 的顏色，$\mathit{found}$ 是全域的旗標。

\begin{pseudo}[unbreakable]{用最少顏色著色}
\Procedure{Try-Color}{$k, m$} \Comment{頂點 $1..k-1$ 已經塗好}
    \For{$c \gets 1$ \To $m$}
        \If{$k$ 的每個鄰居 $j < k$ 都有 $x_j \ne c$}
            \State $x_k \gets c$
            \If{$k = n$}
                \State $\mathit{found} \gets \TRUE$
            \Else
                \State \Call{Try-Color}{$k + 1, m$}
            \EndIf
            \If{$\mathit{found}$}
                \State \Return \Comment{找到一種塗法就不再試別的顏色}
            \EndIf
        \EndIf
    \EndFor
\EndProcedure
\Statex
\Function{Min-Coloring}{$G$}
    \For{$m \gets 1$ \To $n$}
        \State $\mathit{found} \gets \FALSE$
        \State \Call{Try-Color}{$1, m$}
        \If{$\mathit{found}$}
            \State \Return $m$ 和塗法 $(x_1, \ldots, x_n)$
        \EndIf
    \EndFor
\EndFunction
\end{pseudo}

\textbf{正確性。} \proc{Try-Color}$(1, m)$ 走的是 $m$ 著色的狀態空間樹，只剪掉「新塗的頂點和
已經塗過的鄰居同色」的節點，這種節點不管後面怎麼塗都有一條邊兩端同色，子樹裡沒有合法的塗
法。 其他節點都會被走到（除非已經找到），所以 $\mathit{found}$ 變成 \TRUE{} 若且唯若圖能用
$m$ 種顏色塗，而且那時 $x_1, \ldots, x_n$ 就是一種合法的塗法。 $m$ 由小到大試，第一個成功的
$m$ 之前每一個都不能塗，所以它就是最少的顏色數。 $m = n$ 時每個頂點一種顏色一定可以，所以
迴圈一定會回傳。

\textbf{時間。} 對固定的 $m$，狀態空間樹最多 $\sum_{i=0}^{n} m^i = O(m^n)$ 個節點，用鄰接串
列檢查第 3 行，每個節點 $O(n)$，所以一次 \proc{Try-Color}$(1, m)$ 是 $O(n m^n)$。 設最少要
$\chi$ 種顏色，總時間是 $\sum_{m=1}^{\chi} O(n m^n) = O(n \chi^{n+1})$，最壞是指數時間。 這
個問題的決策版本在 $m \ge 3$ 時是 NP 完備的（第~28~章），目前沒有已知的多項式時間演算法；
實務上可以用第~\ref{sec:bk-coloring}~節的「固定頂點 1 塗顏色 $1$」去掉對稱，先塗鄰居多的頂
點，讓搜尋少走一些節點。

\textbf{不要寫成單純的貪心。} 「每個頂點依序塗上鄰居沒用過的最小顏色」只要
$O(n + m_E)$（$m_E$ 是邊數），塗出來一定合法，但顏色不一定最少。 例如四個頂點、邊
$\{1, 3\}$、$\{3, 4\}$、$\{4, 2\}$（一條路徑 $1 - 3 - 4 - 2$）：依編號塗，頂點 1、2 都塗 $1$，
頂點 3 塗 $2$，頂點 4 的鄰居用了 $1$、$2$，只好塗 $3$；其實兩種顏色就夠（$1$、$4$ 一色，
$2$、$3$ 另一色）。 題目要求最少，所以要用上面的搜尋；用 \proc{Min-Coloring} 算這個圖，
$m = 2$ 就成功。

\pastsol{px:bk-4}{中興資工所 115 學年度 第 7 題}
答案是 \textbf{A}。

嚴格說，回溯也會用界限函數剪枝（第~\ref{sec:bk-tree}~節），所以 A 要理解成分支限界的特色：
它通常用來解最佳化問題，對每個節點算出目標值的上界或下界，和目前找到的最好解比較，不可能更
好的子樹整棵丟掉（第~\ref{sec:bk-bb}~節的界限剪枝）；而且 E-節點一次產生全部的子節點，再依
FIFO、LIFO 或最小成本挑下一個 E-節點。 四個選項裡只有 A 說的是分支限界。

\begin{itemize}
\item B 錯：分支限界挑下一個 E-節點的規則是固定的（佇列、堆疊或 $\hat c$ 最小），不是隨機。
\item C 錯：只往深處走是回溯的走法；分支限界可以一層一層走（FIFO），也可以照界限跳來跳去（最
佳優先）。
\item D 錯：貪心一旦做了選擇就不再回頭；分支限界把其他活節點都留著，之後還可能回去展開。 例
如第~\ref{sec:bk-tsp}~節的 TSP，展開節點 6 之後又跳回第 1 層展開節點 2。
\end{itemize}

\pastsol{px:bk-5}{中央資工所 100 學年度 第 6 題}
物品已經依 $P_i / W_i$ 由大到小排好：$0.6, 0.526\ldots, 0.5, 0.5, 0.5$。

\textbf{6.1 上界。} 用第~\ref{sec:bk-knap}~節的分數背包上界（虛擬碼~\ref{alg:bk-bound}）：
把物品照順序放，放得下就整個放，放不下的那一個只放剛好填滿的部分。

\begin{center}
\begin{tabular}{cccl}
\toprule
物品 & 放進去的部分 & 累計重量 & 累計價值 \\
\midrule
1 & 整個 & 10 & 6 \\
2 & 整個 & 29 & 16 \\
3 & 整個 & 37 & 20 \\
4 & 剩下容量 $1$，放 $1/10$ & 38 & $20 + 5 \cdot \frac{1}{10} = 20.5$ \\
\bottomrule
\end{tabular}
\end{center}

上界是 $20.5$。 理由：任何一個 0/1 的解也是同一組資料的分數背包的解，而上面的貪心解是分數背
包的最佳解（定理~\ref{thm:gr-frac}），所以每個 0/1 的解價值都不超過 $20.5$。 價值都是整數，
上界可以再收緊成 $\lfloor 20.5 \rfloor = 20$。

\textbf{6.1 下界。} 找一個實際放得進去的 0/1 解，它的價值就是最佳值的下界。 用同樣的順序做
整數的貪心：物品 1、2、3 放得下（重 $37$），物品 4 放進去會變成 $47 > 38$，跳過；物品 5 放進
去是 $49 > 38$，也跳過。 得到 $\{1, 2, 3\}$，價值 $20$，所以下界是 $20$。

\textbf{6.2 界限怎麼用。} 用最佳優先的分支限界（第~\ref{sec:bk-knap}~節）：
\begin{enumerate}
\item 每個節點代表「前 $k$ 個物品已經決定拿或不拿」，算出它的上界（這棵子樹裡的解最多值多少）
和下界（這棵子樹裡一定做得到的值）。
\item 記住目前所有節點裡最大的下界 $\mathit{best}$，它是某個實際的解。
\item 每次從活節點裡取上界最大的展開，產生「拿」和「不拿」兩個子節點；超重的子節點丟掉。
\item 上界 $\le \mathit{best}$ 的節點不必展開，因為它的子樹裡不會有比 $\mathit{best}$ 更好的
解（界限剪枝）。 活節點裡最大的上界都 $\le \mathit{best}$ 時就停止，$\mathit{best}$ 就是最佳值。
\end{enumerate}

用在這一題：根的上界收緊後是 $20$，等於下界 $20$，所以根就停了，$\{1, 2, 3\}$、價值 $20$ 是
最佳解，不必展開任何節點（把 $32$ 種組合全部列出來檢查，最大值也是 $20$）。 如果上界不取整，
照虛擬碼~\ref{alg:bk-knap-bb} 做，$20.5 > 20$ 的節點會一直被展開：要產生 $15$ 個節點、展開 $7$
個 E-節點（依序是 1、2、4、6、7、9、11）才停下，答案同樣是 $20$。 價值是整數時把上界取整，可
以少走很多節點。

有的課本（第~\ref{sec:bk-knap}~節的作答提醒）把背包改寫成求最小值：價值取負號，根的
$\hat c = -20.5$、上界 $u = -20$，$\mathit{upper}$ 從 $-20$ 開始；數字的正負號相反，用法相同。


% 第五部：圖論
\chapter{圖的表示、BFS 與 DFS}
\label{ch:gtr}

% 本章的圖共用的樣式（前綴 gtr-）
\tikzset{
    gtr-v/.style={circle, draw=black!70, thick, minimum size=0.62cm, inner sep=0pt, font=\small, fill=white},
    gtr-vs/.style={gtr-v, minimum size=0.5cm, font=\scriptsize},
    gtr-vA/.style={gtr-v, draw=coreStroke, fill=coreFill},
    gtr-vB/.style={gtr-v, draw=memStroke, fill=memFill},
    gtr-e/.style={draw=black!65, thick},
    gtr-a/.style={-{Stealth[length=2.2mm]}, draw=black!65, thick},
    gtr-tree/.style={draw=coreStroke, very thick},
    gtr-treea/.style={gtr-tree, -{Stealth[length=2.4mm]}},
    gtr-back/.style={draw=memStroke, thick, dashed, -{Stealth[length=2.2mm]}},
    gtr-fwd/.style={draw=greenStroke, thick, densely dotted, -{Stealth[length=2.2mm]}},
    gtr-cross/.style={draw=optStroke, thick, dash dot, -{Stealth[length=2.2mm]}},
    gtr-non/.style={draw=black!45, thick, dashed},
    gtr-odd/.style={draw=examStroke, very thick},
    gtr-t/.style={font=\scriptsize, text=black!70, inner sep=1pt},
    gtr-k/.style={font=\scriptsize\bfseries, inner sep=1pt, fill=white},
    gtr-lvl/.style={font=\scriptsize, text=black!55, anchor=east},
    gtr-cell/.style={rectangle, draw=black!65, minimum height=0.5cm, minimum width=0.55cm, inner sep=1pt, font=\small},
    gtr-ptr/.style={-{Stealth[length=1.8mm]}, draw=black!70},
    gtr-iv/.style={draw=coreStroke, fill=coreFill, rounded corners=1pt},
}
% 虛擬碼裡的三種顏色
\newcommand{\gtrwhite}{\textsc{white}}
\newcommand{\gtrgray}{\textsc{gray}}
\newcommand{\gtrblack}{\textsc{black}}
% 鄰接串列的一列：\gtradjrow{y}{頂點}{鄰居串列（逗號分開，空的寫 {}）}
\newcommand{\gtradjrow}[3]{%
    \node[gtr-cell, minimum width=0.7cm] (h#2) at (0, #1) {#2};
    \def\gtrprev{h#2}%
    \foreach \gtrv [count=\gtrk] in {#3} {
        \node[gtr-cell] (n#2-\gtrk) at (0.6 + \gtrk*1.05, #1) {\gtrv};
        \draw[gtr-ptr] (\gtrprev.east) -- (n#2-\gtrk.west);
        \xdef\gtrprev{n#2-\gtrk}%
    }
    \node[font=\scriptsize, anchor=west, text=black!60] at (\gtrprev.east) {\,\NIL};
}

登山社把後山的步道畫成一張圖：岔路口和景點是點，兩點之間有一段步道就連一條線。 新
社員從停車場出發，想知道到每個景點最少要走幾段步道；社長想知道地圖上有沒有哪些點從
停車場根本走不到。 社團網站也可以畫成類似的圖：每一頁是一個點，頁面上有連到另一頁的
超連結，就畫一個從這頁指向那頁的箭頭；管理員想知道，一直點連結會不會繞回同一頁。 這
幾件事都要把圖有系統地走一遍。 如果讓你拿著地圖一個點一個點去標記，你會照什麼順序走，
才能保證每個走得到的點都標到、而且不重複？ 往下讀之前，先想一想。

這一章是第五部「圖論」的第一章。 先把圖的名詞和符號定下來，第~22~到~26~章都沿用這一
套；接著講兩種存圖的方法，再講兩種走訪圖的方法。 \textbf{廣度優先搜尋}（breadth-first
search，BFS）從起點一圈一圈往外擴，順便求出到每個點最少要走幾條邊；\textbf{深度優先搜
尋}（depth-first search，DFS）沿著一條路走到底，走不下去才退回上一個岔路口，它留下的
「發現時間」和「完成時間」可以把圖的邊分成四類。 第~22~章的拓樸排序、第~23~章的強連
通元件和割點，都建立在這些時間上。 最後用兩種搜尋解三個問題：連通元件、二分圖、有向
環。

這一章用到《C 程式設計基礎》第~14~章的結構指標、第~15~章的遞迴，以及第~19~章的
\texttt{malloc}、指標的指標和串列。 本書第~5~章的串列插入、第~6~章的佇列和堆疊、
第~7~章的樹和層序走訪、第~13~章的互斥集合都只引用，不再重講。

\section{圖的名詞與符號}
\label{sec:gtr-terms}

\subsection{圖、頂點、邊}

一張\textbf{圖}（graph）寫成 $G = (V, E)$。 $V$ 是有限個\textbf{頂點}（vertex，也
叫節點）組成的集合，$E$ 是\textbf{邊}（edge）組成的集合，每條邊連接兩個頂點。 本書
用
\[
n = |V|, \qquad m = |E|
\]
表示頂點數和邊數，演算法的時間寫成 $n$ 和 $m$ 的函數，例如 $\Theta(n + m)$。

\begin{itemize}
\item \textbf{無向圖}（undirected graph）的邊沒有方向。 連接 $u$、$v$ 的邊是集合
$\{u, v\}$，習慣上也寫成 $(u, v)$，這時 $(u, v)$ 和 $(v, u)$ 是同一條邊。 步道圖是無
向圖：一段步道兩個方向都能走。
\item \textbf{有向圖}（directed graph）的邊有方向。 邊 $(u, v)$ 是有序對，畫成從 $u$
指向 $v$ 的箭頭，也寫成 $u \to v$；$u$ 叫這條邊的\textbf{尾}（tail），$v$ 叫\textbf{頭}
（head）。 $(u, v)$ 和 $(v, u)$ 是兩條不同的邊，可以只有其中一條。 網頁的超連結是有向
圖：$A$ 頁連到 $B$ 頁，不代表 $B$ 頁也連回 $A$。
\end{itemize}

兩端是同一個頂點的邊 $(v, v)$ 叫\textbf{自環}（self-loop）；同一對頂點之間不只一條邊，
叫\textbf{重邊}（multiple edges）。 沒有自環也沒有重邊的圖叫\textbf{簡單圖}（simple
graph）。 \textbf{本書第~21~到~26~章的圖，除非特別說明，都是簡單圖}：無向圖每一對頂點
之間最多一條邊；有向圖沒有自環，從 $u$ 到 $v$ 最多一條邊（$u \to v$ 和 $v \to u$ 可以
同時存在，它們是不同的兩條）。

如果 $(u, v) \in E$，就說 $v$ 和 $u$ \textbf{相鄰}（adjacent），$v$ 是 $u$ 的一個
\textbf{鄰居}（neighbor），而邊 $(u, v)$ 和頂點 $u$、$v$ \textbf{關聯}（incident）。 在
有向圖裡，「$v$ 是 $u$ 的鄰居」只看從 $u$ 出去的邊：$u \to v$ 存在時 $v$ 是 $u$ 的鄰居，
反過來不一定。

圖~\ref{fig:gtr-examples} 是本章一直會用到的兩張圖。 左邊的 $G_1$ 是步道圖，十個頂點用
$1$ 到 $10$ 編號，$1$ 是停車場；$8$、$9$、$10$ 在溪的對岸，和其他點之間沒有步道。 右邊
的 $G_2$ 是社團網站的八個頁面，用大寫字母 $A$ 到 $H$ 命名。 頂點的名字在虛擬碼和手算
裡用 $1..n$ 或字母，在 C 程式裡一律換成 $0..n-1$：頂點 $i$ 存成 \texttt{i - 1}，字母
$X$ 存成 \texttt{X - 'A'}，印出來時再換回去。

\begin{center}
\begin{tikzpicture}[x=1cm, y=1cm]
\begin{scope}
\node[gtr-v] (u1) at (0, 0) {1};
\node[gtr-v] (u2) at (1.5, 1) {2};
\node[gtr-v] (u3) at (1.5, -1) {3};
\node[gtr-v] (u4) at (3, 1) {4};
\node[gtr-v] (u5) at (3, -1) {5};
\node[gtr-v] (u6) at (4.5, 1) {6};
\node[gtr-v] (u7) at (4.5, -1) {7};
\node[gtr-v] (u8) at (6, 1) {8};
\node[gtr-v] (u9) at (6, 0) {9};
\node[gtr-v] (u10) at (6, -1) {10};
\foreach \a/\b in {1/2, 1/3, 2/4, 3/4, 3/5, 4/6, 5/6, 5/7, 6/7, 8/9, 9/10}
    \draw[gtr-e] (u\a) -- (u\b);
\node[font=\small] at (3, -1.9) {$G_1$：步道（無向）};
\end{scope}
\begin{scope}[xshift=7.6cm]
\node[gtr-v] (A) at (0, 0) {$A$};
\node[gtr-v] (B) at (1.6, 1) {$B$};
\node[gtr-v] (C) at (3.2, 1) {$C$};
\node[gtr-v] (E) at (1.6, -1) {$E$};
\node[gtr-v] (D) at (3.2, -1) {$D$};
\node[gtr-v] (G) at (4.8, 1) {$G$};
\node[gtr-v] (F) at (4.8, -1) {$F$};
\node[gtr-v] (H) at (6.2, 0) {$H$};
\draw[gtr-a] (A) -- (B);
\draw[gtr-a] (B) -- (C);
\draw[gtr-a] (C) to[bend right=38] (A);
\draw[gtr-a] (A) -- (E);
\draw[gtr-a] (B) -- (E);
\draw[gtr-a] (E) -- (D);
\draw[gtr-a] (D) -- (C);
\draw[gtr-a] (A) to[bend right=38] (D);
\draw[gtr-a] (F) -- (G);
\draw[gtr-a] (G) -- (C);
\draw[gtr-a] (G) -- (H);
\draw[gtr-a] (F) to[bend left=12] (H);
\draw[gtr-a] (H) to[bend left=12] (F);
\node[font=\small] at (3.1, -2.2) {$G_2$：網站的頁面（有向）};
\end{scope}
\end{tikzpicture}
\diagcaption{本章的兩張例子圖。 $G_1$ 有 $n = 10$、$m = 11$；$G_2$ 有 $n = 8$、$m = 13$。
$G_2$ 的 $F$ 和 $H$ 之間兩個方向各有一條邊。}
\label{fig:gtr-examples}
\end{center}

\subsection{度數}

無向圖裡，和頂點 $v$ 關聯的邊數叫 $v$ 的\textbf{度數}（degree），寫成 $\deg(v)$。 有向
圖分成兩個：從 $v$ 出去的邊數叫\textbf{出度}（out-degree）$\deg^{+}(v)$，指進 $v$ 的邊
數叫\textbf{入度}（in-degree）$\deg^{-}(v)$。 有些課本把度數叫「分支度」，入度、出度
叫「內分支度」「外分支度」；第~7~章樹的「分支度」是子節點個數，和這裡的度數不一樣（樹
當成無向圖時，非根節點的度數是子節點數加 $1$）。

在 $G_1$ 裡，$\deg(3) = \deg(5) = \deg(6) = 3$，$\deg(8) = 1$；在 $G_2$ 裡，
$\deg^{+}(A) = 3$、$\deg^{-}(A) = 1$，$\deg^{+}(C) = 1$、$\deg^{-}(C) = 3$。

\begin{theorem}[握手定理]
\label{thm:gtr-handshake}
無向圖裡 $\displaystyle\sum_{v \in V} \deg(v) = 2m$。 有向圖裡
$\displaystyle\sum_{v \in V} \deg^{+}(v) = \sum_{v \in V} \deg^{-}(v) = m$。
\end{theorem}

\begin{proof}
無向圖的每條邊 $\{u, v\}$ 在 $\deg(u)$ 和 $\deg(v)$ 各算一次，所以每條邊在左邊的和裡剛
好算兩次。 有向圖的每條邊 $(u, v)$ 只在 $\deg^{+}(u)$ 算一次，也只在 $\deg^{-}(v)$ 算一
次。
\end{proof}

一個馬上得到的結果：無向圖裡度數是奇數的頂點有偶數個，因為所有度數加起來是偶數。
$G_1$ 的度數依序是 $2, 2, 3, 3, 3, 3, 2, 1, 2, 1$，和是 $22 = 2 \times 11$，奇數度數的
頂點有 $3, 4, 5, 6, 8, 10$ 共六個。

邊數有上界。 $n$ 個頂點的簡單無向圖，每一對頂點最多一條邊，所以
$m \le \binom{n}{2} = n(n-1)/2$；簡單有向圖每個有序對最多一條邊，$m \le n(n-1)$。 每一
對頂點之間都有邊的無向圖叫\textbf{完全圖}（complete graph），寫成 $K_n$。 所以
$m = O(n^2)$，$\lg m = O(\lg n)$。 $m$ 遠小於 $n^2$ 的圖叫\textbf{稀疏}（sparse）圖，
$m$ 接近 $n^2$ 的叫\textbf{稠密}（dense）圖。 道路、網頁、社群關係這類實際的圖大多很稀
疏，這是下一節選擇存法的主要考量。

\subsection{路徑與環}

\begin{definition}[路徑]
\label{def:gtr-path}
從 $u$ 到 $v$ 的一條\textbf{路徑}（path）是一串頂點 $\langle v_0, v_1, \ldots, v_k
\rangle$，其中 $v_0 = u$、$v_k = v$，而且 $(v_{i-1}, v_i) \in E$，$i = 1, 2, \ldots,
k$。 $k$ 是路徑的\textbf{長度}，也就是經過的邊數；只有一個頂點的 $\langle u \rangle$ 是
長度 $0$ 的路徑。 路徑上的頂點都不相同時，叫\textbf{簡單路徑}（simple path）。
\end{definition}

有向圖的路徑要順著箭頭走。 從 $u$ 到 $v$ 有路徑時，說 $v$ 從 $u$ \textbf{可到達}
（reachable），寫成 $u \leadsto v$。 例如 $G_2$ 裡 $\langle F, G, C, A, B \rangle$ 是長
度 $4$ 的簡單路徑，所以 $F \leadsto B$；但從 $A$ 出發的路徑都到不了 $F$。

路徑可以重複經過頂點，例如 $G_1$ 的 $\langle 1, 2, 4, 3, 1, 2 \rangle$。 如果從 $u$ 到
$v$ 有路徑，就一定有簡單路徑：路徑上某個頂點 $x$ 出現兩次時，把兩次之間繞的那一圈剪
掉，還是從 $u$ 到 $v$ 的路徑，而且比較短；一直剪到沒有重複為止。 所以談「可不可以到
達」「最短路徑」時，只要考慮簡單路徑。

\begin{definition}[環]
\label{def:gtr-cycle}
\textbf{環}（cycle）是一條長度 $k \ge 1$、起點等於終點（$v_0 = v_k$）的路徑，而且
$v_1, v_2, \ldots, v_k$ 互不相同。 無向圖另外要求 $k \ge 3$。 沒有環的圖叫\textbf{無環}
（acyclic）。
\end{definition}

無向圖要求 $k \ge 3$，是為了排除 $\langle u, v, u \rangle$ 這種沿著同一條邊走過去再走回
來的情形。 有向圖的 $\langle u, v, u \rangle$ 用到 $u \to v$ 和 $v \to u$ 兩條不同的邊，
是長度 $2$ 的環，例如 $G_2$ 的 $\langle F, H, F \rangle$。 $G_2$ 的另一個環是
$\langle A, B, C, A \rangle$；$G_1$ 有 $\langle 5, 6, 7, 5 \rangle$、
$\langle 1, 2, 4, 3, 1 \rangle$ 等好幾個環。

沒有環的有向圖叫\textbf{有向無環圖}（directed acyclic graph，DAG）。 先修課程、工程
的先後順序都是 DAG，第~22~章的拓樸排序就在 DAG 上做。

\begin{examnote}
路徑和環的定義，各本書不太一樣。 有的書說的「路徑」就是本書的簡單路徑，另外用
walk 或「通路」稱呼可以重複頂點的走法；有的書把本書的環叫「簡單環」，「環」則允許重
複頂點。 另外，有些課本把有向邊寫成 $\langle u, v \rangle$、無向邊寫成 $(u, v)$ 來區分。
題目問「有幾條路徑」「最長的環」這類會因定義而不同的問題時，先看題目怎麼定義；題目沒
給，就在答案開頭寫一句你用的定義。
\end{examnote}

\subsection{連通}

無向圖裡，如果任兩個頂點之間都有路徑，這張圖叫\textbf{連通}（connected）。 不連通
的圖會分成好幾塊：在無向圖裡，「$u$ 和 $v$ 之間有路徑」滿足自反（長度 $0$ 的路徑）、
對稱（路徑倒過來走）和遞移（兩條路徑接起來），是一個等價關係，它的每一個等價類叫一個
\textbf{連通元件}（connected component，也叫連通分量）。 每個連通元件內部是連通的，
不同元件之間沒有邊。 $G_1$ 有兩個連通元件 $\{1, 2, \ldots, 7\}$ 和 $\{8, 9, 10\}$。
第~13~章的互斥集合處理的就是這種「誰和誰同一塊」的問題，第~\ref{sec:gtr-app}~節會比較
兩種做法。

有向圖的「可到達」沒有對稱性：$G_2$ 裡 $F \leadsto A$，但 $A$ 到不了 $F$。 如果有向圖裡
任兩個頂點 $u$、$v$ 都有 $u \leadsto v$ 而且 $v \leadsto u$，這張圖叫\textbf{強連通}
（strongly connected）。 「互相可達」又是一個等價關係，它的等價類叫\textbf{強連通元
件}（strongly connected component）。 $G_2$ 的強連通元件是 $\{A, B, C, D, E\}$、
$\{F, G, H\}$：$A$ 到 $E$ 這五頁可以互相點過去，$F$、$G$、$H$ 也是，但兩組之間只有從
$G$ 到 $C$ 的單向連結。 強連通元件怎麼有效率地求出來，是第~23~章的主題；把有向邊都當
成無向邊之後連通的有向圖，叫\textbf{弱連通}（weakly connected）。

\subsection{子圖、樹與生成樹}

$G' = (V', E')$ 是 $G = (V, E)$ 的\textbf{子圖}（subgraph），如果 $V' \subseteq V$、
$E' \subseteq E$，而且 $E'$ 的每條邊兩端都在 $V'$ 裡。 $V' = V$ 的子圖叫\textbf{生成子
圖}（spanning subgraph），也就是頂點全留、邊挑一部分。 給一個頂點集合 $V'$，把 $E$ 裡
兩端都在 $V'$ 的邊全部留下來，得到的子圖叫 $V'$ 的\textbf{誘導子圖}（induced
subgraph），寫成 $G[V']$。 例如 $G_1[\{5, 6, 7\}]$ 是一個三角形。

連通而且無環的無向圖叫\textbf{樹}（tree），每個連通元件都是樹的無向圖（也就是無環的無
向圖）叫\textbf{森林}（forest）。 這裡的樹沒有指定根，也叫自由樹（free tree）；挑一個頂
點當根，從根往外看，就得到第~7~章的有根樹。

下面的引理把邊數、連通元件數和環連在一起，後面好幾個地方會用到。

\begin{lemma}
\label{lem:gtr-count}
$n$ 個頂點、$m$ 條邊的無向圖至少有 $n - m$ 個連通元件，而且剛好 $n - m$ 個的充分必要
條件是這張圖沒有環。
\end{lemma}

\begin{proof}
從 $n$ 個孤立的頂點開始（$n$ 個元件），把 $m$ 條邊一條一條加進去。 加入邊 $(u, v)$ 時
有兩種情形：
\begin{itemize}
\item $u$、$v$ 原本就在同一個元件：元件數不變。 原本已經有一條 $u$ 到 $v$ 的簡單路徑，
再加上這條邊就形成一個環。
\item $u$、$v$ 原本在不同元件：兩個元件併成一個，元件數減一。
\end{itemize}
每條邊最多讓元件數減一，所以最後至少剩 $n - m$ 個。 剛好剩 $n - m$ 個，等於說每條邊都
屬於第二種情形。 如果某條邊屬於第一種，最後的圖顯然有環；反過來，如果最後的圖有環，看
這個環上最後加入的那條邊 $(u, v)$，加入它的時候，環上其他的邊已經把 $u$、$v$ 連起來
了，它屬於第一種情形。
\end{proof}

這個證明和第~13~章用互斥集合數連通元件是同一件事：每一次真正的合併讓集合數減一，兩
端已經在同一組的邊不改變集合數，而且正好是會造成環的邊（第~24~章的 Kruskal 演算法就
靠這一點挑邊）。 由引理~\ref{lem:gtr-count} 可以得到樹的幾種等價說法。

\begin{theorem}
\label{thm:gtr-tree}
$G$ 是 $n$ 個頂點、$m$ 條邊的無向圖。 下面四句話互相等價：
\begin{enumerate}
\item $G$ 是樹（連通而且無環）。
\item $G$ 連通，而且 $m = n - 1$。
\item $G$ 無環，而且 $m = n - 1$。
\item $G$ 的任兩個頂點之間剛好有一條簡單路徑。
\end{enumerate}
\end{theorem}

\begin{proof}
連通表示元件數是 $1$。 (1)$\Rightarrow$(2)：無環，由引理元件數等於 $n - m$，又等於
$1$，所以 $m = n - 1$。 (2)$\Rightarrow$(3)：元件數 $1 = n - m$，引理的等號成立，所以
無環。 (3)$\Rightarrow$(1)：無環，元件數 $= n - m = 1$，所以連通。 (1)$\Rightarrow$(4)：
連通保證至少一條簡單路徑；如果 $u$ 到 $v$ 有兩條不同的簡單路徑，從 $u$ 出發，找到兩條第
一次分岔的地方和之後第一次再碰頭的地方，兩段合起來是一個環，矛盾。 (4)$\Rightarrow$
(1)：有路徑所以連通；如果有環 $\langle v_0, \ldots, v_k = v_0 \rangle$，$v_0$ 到 $v_1$ 就
有兩條簡單路徑（直接走邊，或沿著環的另一邊繞過去，$k \ge 3$ 保證這是不同的路徑），矛
盾。
\end{proof}

無向圖 $G$ 的\textbf{生成樹}（spanning tree）是一個生成子圖，而且是一棵樹：$n$ 個頂點
全部留下，挑 $n - 1$ 條邊把它們連起來。 $G$ 有生成樹的充分必要條件是 $G$ 連通。 不連
通的圖，每個連通元件各取一棵生成樹，合起來叫\textbf{生成森林}。 這一章的 BFS 和 DFS 都
會順手產生一棵生成樹（或生成森林）；第~24~章要找的，是邊上有權重時總權重最小的那一棵。

\textbf{加權圖}（weighted graph）是每條邊 $(u, v)$ 另外帶一個數 $w(u, v)$ 的圖，例如步
道的長度、網路連線的費用。 這一章的圖都不帶權重，「最短」指的是邊數最少；第~24~到~26~
章再處理權重。

下表整理本章定下的符號，第~22~到~26~章照用。

\begin{center}
\begin{tabular}{lll}
\toprule
符號 & 意思 & 備註 \\
\midrule
$G = (V, E)$ & 圖、頂點集合、邊集合 & \\
$n = |V|$，$m = |E|$ & 頂點數、邊數 & 時間寫成 $\Theta(n + m)$ \\
$(u, v)$，$u \to v$ & 邊；有向圖是從 $u$ 到 $v$ & 無向圖的 $(u, v)$ 就是 $\{u, v\}$ \\
$\deg(v)$ & 無向圖的度數 & $\sum \deg(v) = 2m$ \\
$\deg^{+}(v)$，$\deg^{-}(v)$ & 有向圖的出度、入度 & 和都是 $m$ \\
$u \leadsto v$ & 從 $u$ 可以到達 $v$ & 有長度 $\ge 0$ 的路徑 \\
$\delta(u, v)$ & 從 $u$ 到 $v$ 最少要走幾條邊 & 到不了是 $\infty$（第~\ref{sec:gtr-bfs}~節） \\
$\mathit{Adj}[u]$ & $u$ 的鄰接串列 & 第~\ref{sec:gtr-repr}~節 \\
$d[v]$，$\pi[v]$ & BFS 的距離、前驅 & 到不了是 $\infty$、\NIL \\
$d[v]$，$f[v]$ & DFS 的發現、完成時間 & 從 $1$ 起，共 $2n$ 個時間 \\
白、灰、黑 & 頂點的狀態（\gtrwhite、\gtrgray、\gtrblack）& 第~\ref{sec:gtr-bfs}、\ref{sec:gtr-dfs}~節 \\
\bottomrule
\end{tabular}
\end{center}

\section{圖的表示法}
\label{sec:gtr-repr}

演算法要問圖的問題，最常見的是兩種：「$(u, v)$ 是不是一條邊？」和「$u$ 的鄰居有哪
些？」。 兩種標準存法各自讓其中一種問題比較快。

\subsection{鄰接矩陣}

頂點編號成 $1..n$ 之後，\textbf{鄰接矩陣}（adjacency matrix）是一個 $n \times n$ 的
矩陣 $A = (a_{uv})$：
\[
a_{uv} = \begin{cases} 1, & (u, v) \in E, \\ 0, & \text{其他}. \end{cases}
\]
下表是 $G_2$ 的鄰接矩陣，第 $u$ 列第 $v$ 行是 $1$ 表示有 $u \to v$ 這條邊。

\begin{center}
\begin{tabular}{c|cccccccc|c}
 & $A$ & $B$ & $C$ & $D$ & $E$ & $F$ & $G$ & $H$ & 列和 $= \deg^{+}$ \\
\hline
$A$ & 0 & 1 & 0 & 1 & 1 & 0 & 0 & 0 & 3 \\
$B$ & 0 & 0 & 1 & 0 & 1 & 0 & 0 & 0 & 2 \\
$C$ & 1 & 0 & 0 & 0 & 0 & 0 & 0 & 0 & 1 \\
$D$ & 0 & 0 & 1 & 0 & 0 & 0 & 0 & 0 & 1 \\
$E$ & 0 & 0 & 0 & 1 & 0 & 0 & 0 & 0 & 1 \\
$F$ & 0 & 0 & 0 & 0 & 0 & 0 & 1 & 1 & 2 \\
$G$ & 0 & 0 & 1 & 0 & 0 & 0 & 0 & 1 & 2 \\
$H$ & 0 & 0 & 0 & 0 & 0 & 1 & 0 & 0 & 1 \\
\hline
行和 $= \deg^{-}$ & 1 & 1 & 3 & 2 & 2 & 1 & 1 & 2 & 13 \\
\end{tabular}
\end{center}

從定義直接看得出幾件事：
\begin{itemize}
\item 第 $u$ 列的和是 $\deg^{+}(u)$，第 $v$ 行的和是 $\deg^{-}(v)$，全部的 $1$ 共有 $m$ 個。
\item 無向圖的鄰接矩陣是對稱的（$a_{uv} = a_{vu}$），$1$ 共有 $2m$ 個；沒有自環時對角線
全是 $0$。 對稱矩陣只要存下三角的 $n(n-1)/2$ 格，位址怎麼算見第~4~章的三角矩陣壓縮。
\item 判斷 $(u, v)$ 是不是邊只要讀一格，$\Theta(1)$；列出 $u$ 的所有鄰居要把第 $u$ 列
整列掃過，$\Theta(n)$，不管 $u$ 實際上有幾個鄰居。
\item 不管邊有多少，都要 $\Theta(n^2)$ 的空間，光是把矩陣設成 $0$ 就要 $\Theta(n^2)$ 的
時間。 $n = 10^5$ 的道路圖，即使每格只用一個位元，也要大約 $1.25 \times 10^9$ 位元
組。
\end{itemize}

加權圖的鄰接矩陣在 $a_{uv}$ 存權重 $w(u, v)$，沒有邊的格子存一個代表「沒有邊」的值，
最短路徑問題通常存 $\infty$。 第~25~章的 Floyd--Warshall 演算法直接在這種矩陣上計算。

\subsection{鄰接串列}

\textbf{鄰接串列}（adjacency list）用一個長度 $n$ 的陣列 $\mathit{Adj}[1..n]$，
$\mathit{Adj}[u]$ 是一條串列，存 $u$ 的所有鄰居。 圖~\ref{fig:gtr-adjlist} 是 $G_1$ 的鄰
接串列。 有向圖的 $\mathit{Adj}[u]$ 只存從 $u$ 出去的邊的另一端；無向圖的每條邊
$\{u, v\}$ 要存兩次，$v$ 在 $\mathit{Adj}[u]$ 裡，$u$ 也在 $\mathit{Adj}[v]$ 裡。

\begin{center}
\begin{tikzpicture}[x=1cm, y=0.62cm]
\gtradjrow{0}{1}{2,3}
\gtradjrow{-1}{2}{1,4}
\gtradjrow{-2}{3}{1,4,5}
\gtradjrow{-3}{4}{2,3,6}
\gtradjrow{-4}{5}{3,6,7}
\gtradjrow{-5}{6}{4,5,7}
\gtradjrow{-6}{7}{5,6}
\gtradjrow{-7}{8}{9}
\gtradjrow{-8}{9}{8,10}
\gtradjrow{-9}{10}{9}
\node[font=\scriptsize, text=black!60] at (0, 0.75) {$\mathit{Adj}$};
\end{tikzpicture}
\diagcaption{$G_1$ 的鄰接串列。 左邊一格是陣列 $\mathit{Adj}$ 的一個元素，存那條串列的
第一個節點；每條串列裡的鄰居由小到大排。 $11$ 條邊，共 $22$ 個節點。}
\label{fig:gtr-adjlist}
\end{center}

串列節點總數是 $m$（有向）或 $2m$（無向），加上陣列的 $n$ 格，空間是 $\Theta(n + m)$。
列出 $u$ 的所有鄰居花 $\Theta(1 + \deg(u))$，只和 $u$ 實際的鄰居數有關；代價是判斷
$(u, v)$ 是不是邊時要在 $\mathit{Adj}[u]$ 裡一個一個找，最壞 $\Theta(\deg(u))$。

同一張圖的鄰接串列不只一種，串列裡的順序可以任意排。 BFS 和 DFS 走的順序就是串列的順
序，所以順序不同，走出來的結果也不同。 \textbf{本書的約定：每條串列裡的鄰居按編號（或
字母）由小到大排}，手算追蹤和 C 程式都照這個順序。

有向圖的鄰接串列只記出去的邊，要知道「誰指向 $v$」得把整個結構掃一遍。 需要常常查入
邊時（例如第~22~章 Kahn 演算法要的入度、第~23~章要的轉置圖），可以另外建一份\textbf{反
鄰接串列}（inverse adjacency list）：$\mathit{Adj}^{-1}[v]$ 存所有 $u \to v$ 的 $u$。
它其實就是把每條邊反過來的圖（轉置圖）的鄰接串列。

\subsection{其他幾種存法}

下面三種存法在資料結構課本和考題裡常見，這裡只講結構，本書後面的程式不會用到。

\textbf{邊串列}（edge list）。 用一個長度 $m$ 的陣列存每條邊的兩端（加權時再存權
重），空間 $\Theta(m)$。 它不能快速列出鄰居，但「把所有邊依權重排序」很方便，第~13~章
的互斥集合和第~24~章的 Kruskal 演算法都直接吃邊串列。

\textbf{循序表示}（也叫壓縮鄰接陣列）。 圖建好之後不再改變時，可以把所有鄰接串列依頂點
順序接成一個長度 $m$（無向圖是 $2m$）的陣列 $\mathit{nbr}$，另外用陣列
$\mathit{start}[1..n+1]$ 記每個頂點的鄰居從哪一格開始：$u$ 的鄰居是
$\mathit{nbr}[\mathit{start}[u] .. \mathit{start}[u+1] - 1]$，$\mathit{start}[n+1]$ 設成
$m + 1$ 當結尾。 $G_2$ 存成：

\begin{center}
\begin{tabular}{l|ccccccccc}
 & $A$ & $B$ & $C$ & $D$ & $E$ & $F$ & $G$ & $H$ & （結尾）\\
\hline
$\mathit{start}$ & 1 & 4 & 6 & 7 & 8 & 9 & 11 & 13 & 14 \\
\end{tabular}

\smallskip
\begin{tabular}{l|ccccccccccccc}
 & 1 & 2 & 3 & 4 & 5 & 6 & 7 & 8 & 9 & 10 & 11 & 12 & 13 \\
\hline
$\mathit{nbr}$ & $B$ & $D$ & $E$ & $C$ & $E$ & $A$ & $C$ & $D$ & $G$ & $H$ & $C$ & $H$ & $F$ \\
\end{tabular}
\end{center}

$\deg^{+}(u) = \mathit{start}[u+1] - \mathit{start}[u]$，例如 $\deg^{+}(F) = 11 - 9 = 2$。
它和鄰接串列一樣是 $\Theta(n + m)$ 的空間，但沒有指標，資料連續存放，走訪時比較快；缺點
是加一條邊要把後面的資料全部往後搬。

\textbf{鄰接多重串列}（adjacency multilist）。 無向圖的鄰接串列把每條邊存兩次，對邊做
標記（例如「這條邊走過了」）時，兩份都要改。 鄰接多重串列讓每條邊只有一個節點，節點裡
存兩端 $u$、$v$，以及兩個指標：$\mathit{link}_u$ 指向下一條和 $u$ 關聯的邊，
$\mathit{link}_v$ 指向下一條和 $v$ 關聯的邊；需要時再加一個標記欄位。 每個頂點的串列頭
指向第一條和它關聯的邊，沿著「自己那一端」的指標走，就能列出它的所有邊。 下面是一張四
個頂點、五條邊的圖，邊依 $e_1$ 到 $e_5$ 的順序串起來：

\begin{center}
\begin{tikzpicture}[x=1cm, y=1cm, baseline=(current bounding box.center)]
\node[gtr-vs] (m1) at (0, 0) {1};
\node[gtr-vs] (m2) at (1.4, 0.8) {2};
\node[gtr-vs] (m3) at (1.4, -0.8) {3};
\node[gtr-vs] (m4) at (2.8, 0) {4};
\draw[gtr-e] (m1) -- node[gtr-k, above left] {$e_1$} (m2);
\draw[gtr-e] (m1) -- node[gtr-k, below left] {$e_2$} (m3);
\draw[gtr-e] (m2) -- node[gtr-k, left] {$e_3$} (m3);
\draw[gtr-e] (m2) -- node[gtr-k, above right] {$e_4$} (m4);
\draw[gtr-e] (m3) -- node[gtr-k, below right] {$e_5$} (m4);
\end{tikzpicture}
\qquad
\begin{tabular}{cccccc}
\toprule
邊 & $u$ & $v$ & $\mathit{link}_u$ & $\mathit{link}_v$ \\
\midrule
$e_1$ & 1 & 2 & $e_2$ & $e_3$ \\
$e_2$ & 1 & 3 & \NIL & $e_3$ \\
$e_3$ & 2 & 3 & $e_4$ & $e_5$ \\
$e_4$ & 2 & 4 & \NIL & $e_5$ \\
$e_5$ & 3 & 4 & \NIL & \NIL \\
\bottomrule
\end{tabular}
\qquad
\begin{tabular}{cc}
\toprule
頂點 & 串列頭 \\
\midrule
1 & $e_1$ \\
2 & $e_1$ \\
3 & $e_2$ \\
4 & $e_4$ \\
\bottomrule
\end{tabular}
\end{center}

列出頂點 $3$ 的邊：從串列頭 $e_2$ 開始，$3$ 是 $e_2$ 的 $v$ 端，下一條看
$\mathit{link}_v = e_3$；$3$ 是 $e_3$ 的 $v$ 端，下一條是 $e_5$；$3$ 是 $e_5$ 的 $u$ 端，
$\mathit{link}_u = \NIL$，結束。 得到 $e_2, e_3, e_5$。 走的時候每到一個節點都要先看自己
是 $u$ 端還是 $v$ 端，再決定跟哪一個指標。 有向圖也有類似的做法，每條邊一個節點，同時
串在尾端的出邊串列和頭端的入邊串列上，叫\textbf{十字串列}（orthogonal list），就是第~5~
章稀疏矩陣的正交串列用在鄰接矩陣上。

\subsection{比較}

\begin{center}
\begin{tabular}{lcc}
\toprule
操作 & 鄰接矩陣 & 鄰接串列 \\
\midrule
空間 & $\Theta(n^2)$ & $\Theta(n + m)$ \\
判斷 $(u, v)$ 是不是邊 & $\Theta(1)$ & $O(\deg(u))$ \\
列出 $u$ 的所有鄰居 & $\Theta(n)$ & $\Theta(1 + \deg(u))$ \\
掃過所有的邊（BFS、DFS）& $\Theta(n^2)$ & $\Theta(n + m)$ \\
算出所有頂點的入度 & $\Theta(n^2)$ & $\Theta(n + m)$ \\
加一條邊 & $\Theta(1)$ & $\Theta(1)$（插在開頭）\\
刪一條邊 $(u, v)$ & $\Theta(1)$ & $O(\deg(u))$（無向：$+\deg(v)$）\\
\bottomrule
\end{tabular}
\end{center}

圖稀疏時（大部分的實際問題），鄰接串列的空間和走訪時間都好得多，本書的圖演算法預設
用鄰接串列。 鄰接矩陣適合三種情形：圖很稠密，$m$ 本來就接近 $n^2$；演算法需要大量
「$(u, v)$ 是不是邊」的查詢；$n$ 很小，寫起來簡單比較重要。 第~25~章的 Floyd--Warshall
和遞移閉包都用矩陣。

\begin{examnote}
題目只給一張圖、問「用鄰接串列表示」時，串列裡的順序會影響之後 BFS、DFS 的答案。
沒有特別說明時，照本書的約定由小到大排，並在答案裡寫一句「串列依編號由小到大」。 題
目已經畫出鄰接串列的，就照它畫的順序走，即使不是由小到大；其他順序走出來的答案會不
同。 問複雜度時記得寫清楚是哪一種表示法：同一個演算法用矩陣是 $\Theta(n^2)$，用串列
是 $\Theta(n + m)$。
\end{examnote}

\subsection{C 的鄰接串列}

C 的鄰接串列照第~5~章的單向串列來寫：每個鄰居是一個用 \texttt{malloc} 要來的節點
（\texttt{struct AdjNode}，存鄰居的編號和 \texttt{next}），\texttt{adj[u]} 指向 \texttt{u}
那條串列的第一個節點，空串列是 \texttt{NULL}。 插入用第~5~章範例~\ref{ex:ll-signup}（\texttt{signup.c}）的
「指標的指標」寫法，找到第一個不小於 \texttt{v} 的位置插進去，每條串列就一直保持由小到大；
遇到相同的 \texttt{v} 就是重邊，不插。 完整的程式放在第~\ref{sec:gtr-c}~節，和 BFS、DFS
寫在一起。

建圖的成本要注意一點。 為了保持由小到大，每次插入要從串列開頭往後找位置，最壞
$\Theta(\deg(u))$，全部插完最壞是 $\Theta\bigl(\sum_u \deg(u)^2\bigr)$，例如一個頂點連到
其他所有頂點、邊又剛好由小到大給的時候是 $\Theta(n^2)$。 不在意順序時，每次插在串列開頭，
建圖是 $\Theta(n + m)$；要排好又要線性時間，可以先把 $m$ 條邊依 $(u, v)$ 用第~15~章的基數
排序排好，再依序接到各串列的尾端。 本章的例子很小，直接排序插入最清楚。

\section{廣度優先搜尋}
\label{sec:gtr-bfs}

\subsection{想法}

回到章首的步道圖：從停車場 $1$ 出發，最少要走幾段步道才能到每一個點？ 一段就能到的是
$1$ 的鄰居；兩段能到、一段到不了的，是那些鄰居的鄰居裡還沒算過的點；依此類推。 BFS 就
是照這個順序做：先把離起點 $s$ 距離 $1$ 的頂點全部找出來，再找距離 $2$ 的，一圈一圈往
外擴。

第~7~章的層序走訪用佇列一層一層拜訪樹的節點，BFS 是同一個想法用在圖上。 圖和樹不一樣
的地方是會有環，同一個頂點可能從好幾個方向被看到（$G_1$ 的 $4$ 是 $2$ 和 $3$ 的共同鄰
居），所以要替頂點做記號，第一次看到時才放進佇列。 本書用三種顏色當記號：

\begin{itemize}
\item \textbf{白色}：還沒被發現。
\item \textbf{灰色}：已經發現、放進了佇列，但它的鄰居還沒看完。 灰色的頂點就是佇列裡的
頂點，也就是目前「已知範圍」的邊界。
\item \textbf{黑色}：已經從佇列拿出來，鄰居都看過了。
\end{itemize}

除了顏色，BFS 替每個頂點 $v$ 記兩個值：$d[v]$ 是從 $s$ 到 $v$ 的距離（目前找到的邊數），
$\pi[v]$ 是 $v$ 是被哪一個頂點發現的，叫 $v$ 的\textbf{前驅}（predecessor）。 還沒發現的
頂點 $d[v] = \infty$、$\pi[v] = \NIL$。

\begin{pseudo}[label=alg:gtr-bfs]{廣度優先搜尋}
\Procedure{BFS}{$G, s$}
    \For{每個頂點 $u \in V$}
        \State $\mathit{color}[u] \gets \gtrwhite$；$d[u] \gets \infty$；$\pi[u] \gets \NIL$
    \EndFor
    \State $\mathit{color}[s] \gets \gtrgray$
    \State $d[s] \gets 0$
    \State $Q \gets$ 空佇列
    \State \Call{Enqueue}{$Q, s$}
    \While{$Q$ 不是空的}
        \State $u \gets \Call{Dequeue}{Q}$
        \For{每個 $v \in \mathit{Adj}[u]$} \Comment{依串列的順序}
            \If{$\mathit{color}[v] = \gtrwhite$}
                \State $\mathit{color}[v] \gets \gtrgray$
                \State $d[v] \gets d[u] + 1$
                \State $\pi[v] \gets u$
                \State \Call{Enqueue}{$Q, v$}
            \EndIf
        \EndFor
        \State $\mathit{color}[u] \gets \gtrblack$
    \EndWhile
\EndProcedure
\end{pseudo}

第 2、3 行把所有頂點設成白色。 第 4 到 7 行把起點塗成灰色、距離設成 $0$、放進佇
列。 主迴圈每一輪從佇列前端拿出一個頂點 $u$（第 9 行），依序檢查它的每個鄰居 $v$：
$v$ 是白色，表示第一次看到它，就塗成灰色，距離比 $u$ 多 $1$，記下是 $u$ 發現它的，放
到佇列尾端（第 11 到 15 行）。 灰色或黑色的鄰居已經發現過了，什麼都不做。 $u$ 的鄰居
全部看完後塗成黑色（第 16 行）。

只有第 11 行會看顏色，而且只分「白」和「不是白」，所以灰和黑其實可以合併成一種「已發
現」。 分成三色是為了說明時方便，也和後面的 DFS 一致；C 程式可以只用 $d[v] = -1$ 代表白
色。

\subsection{手算追蹤}
\label{sec:gtr-bfs-trace}

在 $G_1$ 上從 $s = 1$ 做 BFS，鄰接串列照圖~\ref{fig:gtr-adjlist}。 下表每一列是主迴圈的
一輪：拿出 $u$，依序看 $\mathit{Adj}[u]$ 裡的每個鄰居（括號是看到時的顏色，沒有括號的
是白色），白色的鄰居被發現、放進佇列。 最右欄是這一輪做完之後的佇列，左邊是前端。

\begin{center}
\small
\begin{tabular}{cl>{\raggedright\arraybackslash}p{0.25\linewidth}>{\raggedright\arraybackslash}p{0.27\linewidth}l}
\toprule
輪 & 拿出 $u$ & 依序看的鄰居 & 新發現：$d$，$\pi$ & 之後的佇列 \\
\midrule
0 & --- & --- & $1$：$d = 0$ & $1$ \\
1 & $1$ & $2$，$3$ & $2$：$1$，$1$；$3$：$1$，$1$ & $2\ 3$ \\
2 & $2$ & $1$（黑），$4$ & $4$：$2$，$2$ & $3\ 4$ \\
3 & $3$ & $1$（黑），$4$（灰），$5$ & $5$：$2$，$3$ & $4\ 5$ \\
4 & $4$ & $2$（黑），$3$（黑），$6$ & $6$：$3$，$4$ & $5\ 6$ \\
5 & $5$ & $3$（黑），$6$（灰），$7$ & $7$：$3$，$5$ & $6\ 7$ \\
6 & $6$ & $4$（黑），$5$（黑），$7$（灰） & --- & $7$ \\
7 & $7$ & $5$（黑），$6$（黑） & --- & （空） \\
\bottomrule
\end{tabular}
\end{center}

結束時的 $d$ 和 $\pi$：

\begin{center}
\begin{tabular}{c|cccccccccc}
$v$ & 1 & 2 & 3 & 4 & 5 & 6 & 7 & 8 & 9 & 10 \\
\hline
$d[v]$ & $0$ & $1$ & $1$ & $2$ & $2$ & $3$ & $3$ & $\infty$ & $\infty$ & $\infty$ \\
$\pi[v]$ & \NIL & $1$ & $1$ & $2$ & $3$ & $4$ & $5$ & \NIL & \NIL & \NIL \\
\end{tabular}
\end{center}

這張表有幾個地方要看。 第 3 輪看到 $4$ 時，$4$ 已經是灰色（第 2 輪由 $2$ 發現），所以
$\pi[4] = 2$ 而不是 $3$：誰先發現就算誰的。 第 5 輪的 $6$ 也一樣，$\pi[6] = 4$。 對岸的
$8$、$9$、$10$ 從 $1$ 走不到，BFS 結束時仍是白色，距離是 $\infty$。 佇列裡的頂點，距
離由前到後是 $1, 1$ 或 $1, 2$ 或 $2, 2$ 這種樣子：不會減少，頭尾最多差 $1$。 這是 BFS 正
確的關鍵，下面會證明。

\begin{center}
\begin{tikzpicture}[x=1cm, y=1cm]
\node[gtr-lvl] at (-0.6, 0) {距離 $0$};
\node[gtr-lvl] at (-0.6, -1.2) {距離 $1$};
\node[gtr-lvl] at (-0.6, -2.4) {距離 $2$};
\node[gtr-lvl] at (-0.6, -3.6) {距離 $3$};
\node[gtr-v] (b1) at (1.5, 0) {1};
\node[gtr-v] (b2) at (0.7, -1.2) {2};
\node[gtr-v] (b3) at (2.3, -1.2) {3};
\node[gtr-v] (b4) at (0.7, -2.4) {4};
\node[gtr-v] (b5) at (2.3, -2.4) {5};
\node[gtr-v] (b6) at (0.7, -3.6) {6};
\node[gtr-v] (b7) at (2.3, -3.6) {7};
\foreach \a/\b in {1/2, 1/3, 2/4, 3/5, 4/6, 5/7}
    \draw[gtr-tree] (b\a) -- (b\b);
\draw[gtr-non] (b3) -- (b4);
\draw[gtr-non] (b5) -- (b6);
\draw[gtr-non] (b6) -- (b7);
\node[font=\small, align=left, anchor=west] at (3.6, -1.8)
    {$8$、$9$、$10$：到不了\\$d = \infty$，$\pi = \NIL$};
\end{tikzpicture}
\diagcaption{$G_1$ 從 $1$ 出發的 BFS 樹（粗線，邊 $(\pi[v], v)$）。 虛線是沒用上的邊：
$\{3, 4\}$、$\{5, 6\}$ 連接相鄰的兩個距離，$\{6, 7\}$ 連接同一個距離的兩個頂點。}
\label{fig:gtr-bfs-tree}
\end{center}

有向圖的做法完全一樣，只是只順著箭頭走。 $G_2$ 從 $A$ 出發：

\begin{center}
\small
\begin{tabular}{cl>{\raggedright\arraybackslash}p{0.25\linewidth}>{\raggedright\arraybackslash}p{0.27\linewidth}l}
\toprule
輪 & 拿出 $u$ & 依序看的鄰居 & 新發現：$d$，$\pi$ & 之後的佇列 \\
\midrule
0 & --- & --- & $A$：$d = 0$ & $A$ \\
1 & $A$ & $B$，$D$，$E$ & $B$：$1$，$A$；$D$：$1$，$A$；$E$：$1$，$A$ & $B\ D\ E$ \\
2 & $B$ & $C$，$E$（灰） & $C$：$2$，$B$ & $D\ E\ C$ \\
3 & $D$ & $C$（灰） & --- & $E\ C$ \\
4 & $E$ & $D$（黑） & --- & $C$ \\
5 & $C$ & $A$（黑） & --- & （空） \\
\bottomrule
\end{tabular}
\end{center}

結果 $d[A..E] = 0, 1, 2, 1, 1$，$\pi[B] = \pi[D] = \pi[E] = A$、$\pi[C] = B$；$F$、$G$、$H$
都是 $\infty$。 雖然 $G$ 有邊指向 $C$，但從 $A$ 出發的路都只能順著箭頭走，到不了 $G$。
觀念題會從 $F$ 出發再做一次，那一次每個頂點都到得了。

改變鄰接串列的順序，BFS 拿出頂點的順序和 $\pi$ 可能跟著變（例如 $\mathit{Adj}[1]$ 寫成
$3, 2$，$4$ 就會由 $3$ 發現），但每個頂點的 $d$ 不會變，因為下面會證明 $d[v]$ 一定是最
短距離，和順序無關。

\subsection{時間}

第 2、3 行花 $\Theta(n)$。 頂點只有在白色時才會被放進佇列（第 11 到 15 行），放進去就變
成灰色，不會再變回白色，所以每個頂點最多進佇列一次、出佇列一次，佇列操作共 $O(n)$
次，每次 $\Theta(1)$。 每個頂點出佇列時把它的串列掃一遍，所有串列的長度加起來是 $m$
（有向）或 $2m$（無向），所以第 10 到 15 行總共執行 $O(m)$ 次。 合起來，用鄰接串列時
BFS 的時間是 $O(n + m)$，空間是 $O(n)$（顏色、$d$、$\pi$ 和佇列）。 這是第~3~章聚合法
的典型用法：某一輪的工作量可能很多，但把全部的工作量加起來只有這麼多。

更精確地說，BFS 只會碰到從 $s$ 可到達的頂點和它們的邊，第 2、3 行之外的工作是
$\Theta(n_s + m_s)$，$n_s$、$m_s$ 是可到達的頂點數和它們的出邊數。 用鄰接矩陣存圖時，
第 10 行要掃過一整列，每個出佇列的頂點都花 $\Theta(n)$，最壞是 $\Theta(n^2)$。

\subsection{BFS 求出的是最短距離}
\label{sec:gtr-bfs-proof}

從 $s$ 到 $v$ 的\textbf{最短距離} $\delta(s, v)$ 是所有從 $s$ 到 $v$ 的路徑裡最少的邊數；
到不了時 $\delta(s, v) = \infty$。 距離 $\delta(s, v)$ 的路徑叫\textbf{最短路徑}。 這一章的
「最短」都是指邊數最少，邊上有長度時是第~25~章的問題。

\begin{theorem}
\label{thm:gtr-bfs}
對有向圖或無向圖 $G$ 和起點 $s$ 執行 \proc{BFS}，結束時每個頂點 $v$ 都有
$d[v] = \delta(s, v)$。 而且從 $s$ 可到達的 $v \ne s$ 都有 $\pi[v] \ne \NIL$，從 $s$ 到
$\pi[v]$ 的最短路徑接上邊 $(\pi[v], v)$，就是一條從 $s$ 到 $v$ 的最短路徑。
\end{theorem}

\begin{whybox}[BFS 的距離]
需要三個小事實。

\textbf{(a) 走一條邊，距離最多加一。} 若 $(u, v) \in E$，則 $\delta(s, v) \le \delta(s, u)
+ 1$。 $u$ 到不了時右邊是 $\infty$，顯然成立；否則從 $s$ 到 $u$ 的最短路徑再走 $(u, v)$，
就是一條長度 $\delta(s, u) + 1$ 的路徑到 $v$。

\textbf{(b) $d$ 不會比真正的距離小。} 每個頂點的 $d$ 只在被發現時設定一次。 $d[s] = 0
= \delta(s, s)$。 $v$ 由 $u$ 發現時 $d[v] = d[u] + 1$，而 $u$ 比 $v$ 早被發現，對 $u$ 已經
有 $d[u] \ge \delta(s, u)$，由 (a) 得 $d[v] \ge \delta(s, u) + 1 \ge \delta(s, v)$。 沒被發
現的頂點 $d = \infty$，也成立。

\textbf{(c) 佇列的樣子。} 任何時候，若佇列由前到後是 $\langle v_1, v_2, \ldots, v_r
\rangle$，則
\[
d[v_1] \le d[v_2] \le \cdots \le d[v_r] \le d[v_1] + 1.
\]
一開始只有 $s$，成立。 拿出 $v_1$ 之後，剩下的部分仍然遞增，而且 $d[v_r] \le d[v_1] + 1
\le d[v_2] + 1$，成立。 拿出的 $u = v_1$ 發現 $v$ 而把 $v$ 放到尾端時，$d[v] = d[u] + 1$：
原本的 $d[v_r] \le d[u] + 1 = d[v]$，遞增沒有被破壞；佇列現在的前端（如果還有別人）距
離至少 $d[u]$，所以 $d[v] = d[u] + 1 \le d[\text{前端}] + 1$。 由 (c) 可以知道，頂點進佇
列的順序（也就是被發現的順序）和出佇列的順序，$d$ 值都不會減少。

\textbf{定理的證明。} 對 $k = \delta(s, v)$ 做歸納，證明所有距離是 $k$ 的頂點最後都被發
現、而且 $d[v] = k$。 $k = 0$ 只有 $s$ 自己。 設 $k \ge 1$，距離小於 $k$ 的頂點都已經成
立。 取一條從 $s$ 到 $v$ 的最短路徑，$u$ 是 $v$ 前面那個頂點，$\delta(s, u) = k - 1$。 由
歸納假設，$u$ 被發現、$d[u] = k - 1$，所以 $u$ 會進佇列，之後一定會被拿出來掃串列，看到
$v$。 看到 $v$ 的那一刻：
\begin{itemize}
\item 若 $v$ 是白色，它在這時被發現，$d[v] = d[u] + 1 = k$。
\item 若 $v$ 已經不是白色，它是被某個更早出佇列的頂點 $w$ 發現的，$d[v] = d[w] + 1$。
出佇列的順序 $d$ 值不減，$w$ 比 $u$ 早出來，所以 $d[w] \le d[u] = k - 1$，得 $d[v] \le k$。
再由 (b)，$d[v] \ge k$，所以 $d[v] = k$。
\end{itemize}
到不了的頂點不會被發現（被發現就代表沿著 $\pi$ 有一條路徑），$d = \infty = \delta$。 最
後，$v$ 被 $\pi[v]$ 發現時 $d[v] = d[\pi[v]] + 1$，所以 $\delta(s, \pi[v]) = \delta(s, v) - 1$，
走到 $\pi[v]$ 的最短路徑再接上 $(\pi[v], v)$，長度剛好是 $\delta(s, v)$。
\end{whybox}

(a) 還有一個常用的推論：無向圖的邊 $\{u, v\}$ 兩個方向都可以用 (a)，所以
$|\delta(s, u) - \delta(s, v)| \le 1$，BFS 結束時每條邊兩端的 $d$ 最多差 $1$。 有向圖只
有一個方向：$d[v] \le d[u] + 1$，但 $d[v]$ 可以比 $d[u]$ 小很多（往回指的邊）。

\subsection{BFS 樹與印出路徑}

BFS 結束後，從 $s$ 可到達的頂點，加上所有的邊 $(\pi[v], v)$，構成一棵以 $s$ 為根的樹，叫
\textbf{BFS 樹}（breadth-first tree）。 理由是：除了 $s$，每個可到達的頂點剛好有一個
$\pi$，沿著 $\pi$ 往回走每一步 $d$ 減一，最後一定走到 $s$；所以這些頂點和邊是連通的，
$k$ 個頂點、$k - 1$ 條邊，由定理~\ref{thm:gtr-tree} 是一棵樹。 有向圖的 BFS 樹，邊的方向
都是從 $\pi[v]$ 指向 $v$。 圖~\ref{fig:gtr-bfs-tree} 的粗線就是 $G_1$ 的 BFS 樹。

由定理~\ref{thm:gtr-bfs}，BFS 樹裡從 $s$ 到 $v$ 的那條路徑，長度是 $d[v] = \delta(s,
v)$，就是一條最短路徑。 要把它印出來，從 $v$ 沿著 $\pi$ 往回走到 $s$，但這樣得到的是倒過
來的順序；用遞迴先印前面的部分，再印 $v$，順序就對了。

\begin{pseudo}[label=alg:gtr-path]{沿著前驅印出最短路徑}
\Procedure{Print-Path}{$s, v$}
    \If{$v = s$}
        \State 印出 $s$
    \ElsIf{$\pi[v] = \NIL$}
        \State 印出「從 $s$ 到不了 $v$」
    \Else
        \State \Call{Print-Path}{$s, \pi[v]$}
        \State 印出 $v$
    \EndIf
\EndProcedure
\end{pseudo}

每一層遞迴沿著 $\pi$ 往回一步，共 $d[v] + 1$ 層，時間 $\Theta(d[v] + 1)$。 在 $G_1$ 上，
$\proc{Print-Path}(1, 7)$ 依序遞迴到 $7, 5, 3, 1$，印出 $1, 3, 5, 7$。 遞迴最深 $n$ 層，
圖很大時可以改成先把路徑存進陣列再倒著印。

BFS 樹裡，$s$ 是根，距離 $k$ 的頂點在樹的深度 $k$。 照第~7~章的定義，它們在第 $k + 1$
層，所以本章談 BFS 一律說「距離 $k$」，不說「第幾層」，避免差一的混淆。

\subsection{多個起點}
\label{sec:gtr-multi-source}

有時候要問的是「每個頂點離最近的起點多遠」，起點不只一個，例如每個點到最近的消防站要幾
步。 對每個起點各做一次 BFS 再取最小值，要做好幾次；更好的做法是一開始就把所有起點都設
成 $d = 0$、灰色，全部放進佇列，再照常執行主迴圈。

為什麼對？ 想像加一個新頂點 $s^{*}$，從它連一條邊到每個起點。 從 $s^{*}$ 做 BFS，第一輪
拿出 $s^{*}$，把所有起點以距離 $1$ 放進佇列；之後的過程，和「所有起點以距離 $0$ 開始」
的做法完全一樣，只是每個距離都多 $1$。 由定理~\ref{thm:gtr-bfs}，$\delta(s^{*}, v) - 1$
正是 $v$ 到最近起點的距離。 時間仍然是 $O(n + m)$。

\subsection{C 程式的寫法}

第~\ref{sec:gtr-c}~節的 \texttt{bfs} 函式就是虛擬碼~\ref{alg:gtr-bfs}。 佇列用一個長度 $n$
的陣列 \texttt{q}，\texttt{q[head..tail-1]} 是佇列的內容。 每個頂點最多進佇列一次，
\texttt{tail} 不會超過 $n$，所以不需要第~6~章的環狀佇列，\texttt{head}、\texttt{tail} 只會
往前走，而且 BFS 結束時 \texttt{q[0..tail-1]} 正好是走到的所有頂點，依發現的順序排好。
$d[v] = -1$ 代表 $\infty$，$\pi[v] = -1$ 代表 \NIL。 程式在每拿出一個頂點時印一列，可以和
上面的手算表逐列對照。

\section{深度優先搜尋}
\label{sec:gtr-dfs}

\subsection{想法}

在山裡找遍所有景點，還有另一種走法：每到一個岔路口，挑一條還沒走過的步道一直往前，走
到一個所有相鄰的點都已經去過的地方，才退回上一個岔路口，換另一條還沒走過的路。 這就
是 DFS。 它和第~20~章的回溯是同一種走法：盡量往深處走，走不下去就退回來；和第~7~章的
前序走訪也很像，差別是圖可能有環，要用顏色記住哪些頂點已經去過。

DFS 的三種顏色意思和 BFS 稍有不同：

\begin{itemize}
\item \textbf{白色}：還沒被發現。
\item \textbf{灰色}：已經發現，但還沒完成，也就是 \proc{DFS-Visit} 正在處理它（它的遞迴
呼叫還沒返回）。
\item \textbf{黑色}：已經完成，它的鄰居全部看過，遞迴已經返回。
\end{itemize}

DFS 另外替每個頂點記兩個\textbf{時間戳記}（timestamp）。 用一個全域的計數器
$\mathit{time}$，從 $0$ 開始，每發現一個頂點或完成一個頂點就加一：$d[v]$ 是 $v$ 被發現
（變灰）的時間，$f[v]$ 是 $v$ 完成（變黑）的時間。 每個頂點各被發現、完成一次，所以
$1, 2, \ldots, 2n$ 這 $2n$ 個時間每個剛好用一次，而且 $1 \le d[v] < f[v] \le 2n$。 $\pi[v]$
一樣記是誰發現了 $v$。

\begin{pseudo}[label=alg:gtr-dfs]{深度優先搜尋（遞迴）}
\Procedure{DFS}{$G$}
    \For{每個頂點 $u \in V$}
        \State $\mathit{color}[u] \gets \gtrwhite$；$\pi[u] \gets \NIL$
    \EndFor
    \State $\mathit{time} \gets 0$
    \For{每個頂點 $u \in V$} \Comment{依編號（字母）順序}
        \If{$\mathit{color}[u] = \gtrwhite$}
            \State \Call{DFS-Visit}{$G, u$}
        \EndIf
    \EndFor
\EndProcedure
\Procedure{DFS-Visit}{$G, u$}
    \State $\mathit{time} \gets \mathit{time} + 1$
    \State $d[u] \gets \mathit{time}$
    \State $\mathit{color}[u] \gets \gtrgray$
    \For{每個 $v \in \mathit{Adj}[u]$} \Comment{依串列的順序}
        \If{$\mathit{color}[v] = \gtrwhite$}
            \State $\pi[v] \gets u$
            \State \Call{DFS-Visit}{$G, v$}
        \EndIf
    \EndFor
    \State $\mathit{time} \gets \mathit{time} + 1$
    \State $f[u] \gets \mathit{time}$
    \State $\mathit{color}[u] \gets \gtrblack$
\EndProcedure
\end{pseudo}

\proc{DFS-Visit}$(G, u)$ 從 $u$ 出發，把從 $u$ 經由白色頂點走得到的部分走完。 第 5 到
7 行的外層迴圈依序找還是白色的頂點，從它重新開始；所以 \proc{DFS} 會走遍整張圖，不只是
從某一個起點可到達的部分。 這和 BFS 的習慣不同：BFS 通常是「從 $s$ 出發」，DFS 通常是
「整張圖」，因為 DFS 的應用（判斷有沒有環、拓樸排序、強連通元件）都要看全部的頂點。 只想
從 $s$ 出發時，直接呼叫 $\proc{DFS-Visit}(G, s)$ 就好。

每次在第 7 行呼叫 \proc{DFS-Visit} 都長出一棵樹，樹邊是 $(\pi[v], v)$；所有的樹合起來叫
\textbf{DFS 森林}（depth-first forest）。

\subsection{手算追蹤}
\label{sec:gtr-dfs-trace}

在 $G_2$ 上執行 \proc{DFS}，外層迴圈依字母順序，鄰接串列由小到大（$\mathit{Adj}[A] = \langle B, D, E
\rangle$，其餘見第~\ref{sec:gtr-repr}~節的循序表示）。 考卷上常見的寫法是一個時間一列，同時寫出灰色的頂點有哪些。 灰色的頂點由先
到後排起來，正好是目前的遞迴呼叫堆疊（最右邊是正在處理的頂點）。

\begin{center}
\small
\begin{tabular}{clll}
\toprule
時間 & 事件 & 發生原因 & 灰色頂點（由底到頂） \\
\midrule
1 & 發現 $A$ & 外層迴圈，$A$ 是白色 & $A$ \\
2 & 發現 $B$ & $\mathit{Adj}[A]$ 第一個，白色 & $A\ B$ \\
3 & 發現 $C$ & $\mathit{Adj}[B]$ 第一個，白色 & $A\ B\ C$ \\
4 & 完成 $C$ & $\mathit{Adj}[C] = \langle A \rangle$，$A$ 是灰色 & $A\ B$ \\
5 & 發現 $E$ & 回到 $B$，下一個 $E$ 是白色 & $A\ B\ E$ \\
6 & 發現 $D$ & $\mathit{Adj}[E] = \langle D \rangle$，白色 & $A\ B\ E\ D$ \\
7 & 完成 $D$ & $\mathit{Adj}[D] = \langle C \rangle$，$C$ 是黑色 & $A\ B\ E$ \\
8 & 完成 $E$ & 串列看完 & $A\ B$ \\
9 & 完成 $B$ & 串列看完 & $A$ \\
10 & 完成 $A$ & 回到 $A$，$D$、$E$ 都是黑色 & （空） \\
11 & 發現 $F$ & 外層迴圈：$B$ 到 $E$ 都不是白色，$F$ 是 & $F$ \\
12 & 發現 $G$ & $\mathit{Adj}[F]$ 第一個，白色 & $F\ G$ \\
13 & 發現 $H$ & $\mathit{Adj}[G]$：$C$ 是黑色，$H$ 是白色 & $F\ G\ H$ \\
14 & 完成 $H$ & $\mathit{Adj}[H] = \langle F \rangle$，$F$ 是灰色 & $F\ G$ \\
15 & 完成 $G$ & 串列看完 & $F$ \\
16 & 完成 $F$ & 回到 $F$，$H$ 是黑色 & （空） \\
\bottomrule
\end{tabular}
\end{center}

整理成每個頂點一欄：

\begin{center}
\begin{tabular}{c|cccccccc}
$v$ & $A$ & $B$ & $C$ & $D$ & $E$ & $F$ & $G$ & $H$ \\
\hline
$d[v]$ & 1 & 2 & 3 & 6 & 5 & 11 & 12 & 13 \\
$f[v]$ & 10 & 9 & 4 & 7 & 8 & 16 & 15 & 14 \\
$\pi[v]$ & \NIL & $A$ & $B$ & $E$ & $B$ & \NIL & $F$ & $G$ \\
\end{tabular}
\end{center}

DFS 森林有兩棵樹：$A$ 為根的樹包含 $A$ 到 $E$，$F$ 為根的樹包含 $F$、$G$、$H$
（圖~\ref{fig:gtr-paren} 上半）。 和 BFS 從 $A$ 出發比較：BFS 樹裡 $D$、$E$ 都是 $A$ 的
子節點（距離 $1$），DFS 樹裡 $D$ 卻在深度 $3$，DFS 樹的路徑不是最短路徑。

\subsection{時間}

\proc{DFS-Visit} 只在頂點是白色時被呼叫（第 7、15 行），一進去就把它塗灰（第 11 行），
所以每個頂點剛好被呼叫一次。 每次呼叫把自己的串列掃一遍，第 12 到 15 行總共執行
$\sum_u |\mathit{Adj}[u]| = \Theta(m)$ 次（無向圖是 $2m$）。 加上第 2 到 7 行的
$\Theta(n)$，DFS 的時間是 $\Theta(n + m)$。 和 BFS 一樣是聚合分析；用鄰接矩陣則是
$\Theta(n^2)$。

\subsection{括號定理}

把 $d[v]$ 寫成左括號「$(v$」、$f[v]$ 寫成右括號「$v)$」，依時間排起來，$G_2$ 的結果是
\[
(A\ (B\ (C\ C)\ (E\ (D\ D)\ E)\ B)\ A)\ \ (F\ (G\ (H\ H)\ G)\ F),
\]
一個配對得好好的括號序列。 這不是巧合：遞迴呼叫是後進先出，後發現的頂點一定先完成。

\begin{center}
\begin{tikzpicture}[x=0.62cm, y=1cm]
% DFS 森林
\begin{scope}[yshift=5.6cm]
\node[gtr-v] (fA) at (3.5, 0) {$A$};
\node[gtr-v] (fB) at (3.5, -1) {$B$};
\node[gtr-v] (fC) at (2.2, -2) {$C$};
\node[gtr-v] (fE) at (4.8, -2) {$E$};
\node[gtr-v] (fD) at (4.8, -3) {$D$};
\node[gtr-v] (fF) at (13.5, 0) {$F$};
\node[gtr-v] (fG) at (13.5, -1) {$G$};
\node[gtr-v] (fH) at (13.5, -2) {$H$};
\foreach \a/\b in {A/B, B/C, B/E, E/D, F/G, G/H}
    \draw[gtr-treea] (f\a) -- (f\b);
\foreach \v/\a/\b in {A/1/10, B/2/9, C/3/4, E/5/8, D/6/7, F/11/16, G/12/15, H/13/14}
    \node[gtr-t, anchor=west] at ([xshift=1pt]f\v.east) {\a/\b};
\end{scope}
% 時間軸與區間
\draw[black!50] (0.5, 0) -- (16.5, 0);
\foreach \t in {1,...,16} {
    \draw[black!50] (\t, 0) -- (\t, -0.1);
    \node[font=\scriptsize, text=black!60, below] at (\t, -0.1) {\t};
}
\foreach \v/\a/\b/\r in {A/1/10/1, F/11/16/1, B/2/9/2, G/12/15/2, C/3/4/3, E/5/8/3, H/13/14/3, D/6/7/4} {
    \draw[gtr-iv] (\a - 0.3, 0.45*\r - 0.15) rectangle (\b + 0.3, 0.45*\r + 0.15);
    \node[font=\scriptsize] at ({(\a + \b)/2}, 0.45*\r) {$\v$};
}
\end{tikzpicture}
\diagcaption{上：$G_2$ 的 DFS 森林，頂點旁是 $d/f$。 下：每個頂點的區間 $[d[v], f[v]]$。
任兩個區間不是互相包含就是完全分開，包含關係正好是森林裡的祖先和子孫。}
\label{fig:gtr-paren}
\end{center}

\begin{theorem}[括號定理]
\label{thm:gtr-paren}
在有向圖或無向圖的任何一次 DFS 裡，對任兩個不同的頂點 $u$、$v$，下面三種情形剛好成立一
種：
\begin{enumerate}
\item 區間 $[d[u], f[u]]$ 和 $[d[v], f[v]]$ 完全分開，在 DFS 森林裡 $u$、$v$ 都不是對方的
子孫；
\item $[d[v], f[v]]$ 整個包在 $[d[u], f[u]]$ 裡面，$v$ 是 $u$ 的子孫；
\item $[d[u], f[u]]$ 整個包在 $[d[v], f[v]]$ 裡面，$u$ 是 $v$ 的子孫。
\end{enumerate}
\end{theorem}

\begin{proof}
先看一個事實：任何時刻，灰色的頂點恰好是一條從某棵樹的根往下的樹路徑（就是遞迴堆
疊），新發現的頂點接在這條路徑的最下面，成為所有灰色頂點的子孫。

不妨設 $d[u] < d[v]$（時間不會相同）。 分兩種情形。 若 $d[v] < f[u]$，$v$ 被發現時 $u$
還是灰色，所以 $v$ 是 $u$ 的子孫。 $v$ 的遞迴呼叫在 $u$ 的呼叫裡面，要先返回，$u$ 才會
完成，所以 $f[v] < f[u]$，區間 $[d[v], f[v]]$ 包在 $[d[u], f[u]]$ 裡，這是情形 2。 若
$f[u] < d[v]$，則 $d[u] < f[u] < d[v] < f[v]$，兩個區間分開；$v$ 被發現時 $u$ 已經是黑
色，不在灰色路徑上，所以 $v$ 不是 $u$ 的子孫；$v$ 比 $u$ 晚發現，也不可能是 $u$ 的祖先。
這是情形 1。 $d[u] > d[v]$ 時把兩者對調，得到情形 1 或 3。 兩個區間不可能「部分重疊」，
因為上面兩種情形已經涵蓋所有可能。
\end{proof}

\begin{corollary}
\label{cor:gtr-nest}
$v$ 是 $u$ 的子孫（$v \ne u$）的充分必要條件是 $d[u] < d[v] < f[v] < f[u]$。
\end{corollary}

括號定理把「誰是誰的子孫」變成比大小，之後判斷邊的種類、證明拓樸排序（第~22~章）都用
它。 考卷上常見的題型是給幾個 $d$、$f$ 問「這可能是同一次 DFS 的結果嗎」：只要有兩個區
間部分重疊（例如 $[2, 7]$ 和 $[5, 9]$），或兩個時間相同，就不可能。

\subsection{白路徑定理}

括號定理說的是 DFS 跑完之後的關係。 下面這個定理在 $u$ 剛被發現的那一刻，就能預告
$u$ 的子孫會是哪些頂點。

\begin{theorem}[白路徑定理]
\label{thm:gtr-whitepath}
在 DFS 森林裡，$v$ 是 $u$ 的子孫（$v \ne u$）的充分必要條件是：在時間 $d[u]$，存在一條從
$u$ 到 $v$ 的路徑，路徑上除了 $u$ 之外的頂點全部是白色。
\end{theorem}

\begin{proof}
（$\Rightarrow$）$u$ 到 $v$ 的樹路徑上，除了 $u$ 之外都是 $u$ 的子孫，由系理
\ref{cor:gtr-nest} 它們都比 $u$ 晚發現，所以在時間 $d[u]$ 都是白色。

（$\Leftarrow$）設在時間 $d[u]$ 有一條白色路徑從 $u$ 到 $v$，但 $v$ 不是 $u$ 的子孫。 在
這條路徑上，取第一個不是 $u$ 的子孫的頂點 $w$，它前面的頂點 $x$ 是 $u$ 本身或 $u$ 的子
孫，所以 $f[x] \le f[u]$。 $w$ 在時間 $d[u]$ 是白色，所以 $d[u] < d[w]$。 $x$ 完成之前
一定會看過邊 $(x, w)$，那時 $w$ 若還是白色就會在這時被發現，否則早就被發現了；總之
$d[w] < f[x] \le f[u]$。 於是 $d[u] < d[w] < f[u]$，由括號定理，$w$ 的區間包在 $u$ 的區間
裡，$w$ 是 $u$ 的子孫，矛盾。
\end{proof}

用 $G_2$ 看一次。 時間 $1$ 發現 $A$ 時，所有頂點都是白色，$A$ 經由白色路徑能到的是
$B$、$C$、$D$、$E$，它們正好是 $A$ 那棵樹的全部。 時間 $11$ 發現 $F$ 時，$A$ 到 $E$ 已
經是黑色，$F$ 經由白色頂點能到的只有 $G$、$H$，雖然 $F \leadsto C$，但那條路要經過已經
是黑色的 $C$，所以 $C$ 不會是 $F$ 的子孫。

\subsection{遞迴太深時：改用自己的堆疊}
\label{sec:gtr-dfs-iter}

\begin{cbookrecap}[第~15~章 15.4、15.13、15.14 節]
每一層遞迴呼叫在呼叫堆疊上占一個 frame，放參數、區域變數和返回位置。 呼叫堆疊的大小是
固定的（Linux 上常見的預設是 8 MB），遞迴太深就會用完，程式直接當掉（stack
overflow），而且通常沒有任何錯誤訊息。 任何遞迴都可以改寫成迴圈加上自己管理的堆疊，
自己的堆疊放在 \texttt{malloc} 要來的記憶體裡，大小只受可用記憶體限制。
\end{cbookrecap}

\proc{DFS-Visit} 的遞迴深度等於 DFS 樹裡最長那條樹路徑的頂點數，最壞是 $n$：例如一條鏈
$1 \to 2 \to \cdots \to n$，或任何一張「DFS 一路走到底」的圖。 作者在 Linux 上用 8 MB
的預設堆疊實測，把虛擬碼~\ref{alg:gtr-dfs} 直接翻成 C（不開 sanitizer），在大約
$18$ 萬個頂點的鏈上就當掉了，每層大約用掉 $48$ 個位元組。 圖有十萬、百萬個頂點在實務
上很常見，這時要改用自己的堆疊。

第~\ref{sec:sq-recursion}~節的方法是：堆疊裡放「之後還要做的事」。 \proc{DFS-Visit}$(u)$ 暫停時需要
記住的只有兩件事：正在處理的頂點 $u$，以及 $u$ 的串列看到哪裡了。 每個頂點同一時間最多
在堆疊裡出現一次（只有灰色的頂點在堆疊裡），所以「看到哪裡」可以存在一個以頂點為索引
的陣列 $\mathit{next}[u]$，堆疊裡只放頂點。

\begin{pseudo}[label=alg:gtr-dfs-iter]{DFS-Visit 的迴圈版}
\Procedure{DFS-Visit-Iter}{$G, s$}
    \State $\mathit{time} \gets \mathit{time} + 1$；$d[s] \gets \mathit{time}$；$\mathit{color}[s] \gets \gtrgray$
    \State $\mathit{next}[s] \gets \mathit{Adj}[s]$ 的第一個節點
    \State $S \gets$ 空堆疊；\Call{Push}{$S, s$}
    \While{$S$ 不是空的}
        \State $u \gets \Call{Top}{S}$
        \If{$\mathit{next}[u] \ne \NIL$} \Comment{$u$ 還有鄰居沒看}
            \State $v \gets \mathit{next}[u].\mathit{vertex}$
            \State $\mathit{next}[u] \gets \mathit{next}[u].\mathit{next}$
            \If{$\mathit{color}[v] = \gtrwhite$}
                \State $\pi[v] \gets u$
                \State $\mathit{time} \gets \mathit{time} + 1$；$d[v] \gets \mathit{time}$；$\mathit{color}[v] \gets \gtrgray$
                \State $\mathit{next}[v] \gets \mathit{Adj}[v]$ 的第一個節點
                \State \Call{Push}{$S, v$}
            \EndIf
        \Else \Comment{$u$ 的鄰居都看完了}
            \State \Call{Pop}{$S$}
            \State $\mathit{time} \gets \mathit{time} + 1$；$f[u] \gets \mathit{time}$；$\mathit{color}[u] \gets \gtrblack$
        \EndIf
    \EndWhile
\EndProcedure
\end{pseudo}

把它和虛擬碼~\ref{alg:gtr-dfs} 對照：遞迴版「呼叫 \proc{DFS-Visit}$(G, v)$」變成第 11 到
14 行的「設好 $v$ 再推進堆疊」；遞迴版的 \textbf{for} 迴圈在 $u$ 的串列上前進一步，變成第 8、
9 行把 $\mathit{next}[u]$ 往後移一格；遞迴版返回前的第 16 到 18 行，變成第 16、17 行的 pop
和完成。 每一輪迴圈只看堆疊頂端的那一個頂點，而頂端永遠是遞迴版「正在執行」的那一層，所以兩個
版本發現和完成頂點的順序完全相同，$d$、$f$、$\pi$ 一樣。 堆疊最多 $n$ 個頂點，每一輪要嘛
看掉一條邊、要嘛完成一個頂點，總共 $\Theta(n + m)$ 輪。 外層的 \proc{DFS} 不用改，把第 7
行換成呼叫 \proc{DFS-Visit-Iter} 就好。

網路上和一些課本常見另一種較短的寫法：堆疊裡放「待拜訪的頂點」，拿出一個頂點時如果已
經拜訪過就跳過，否則標記它、把它所有還沒拜訪的鄰居倒著推進堆疊（倒著推，編號小的才會先
出來）。 這種寫法拜訪頂點的順序和遞迴版相同，但同一個頂點可能被推進堆疊好幾次，堆疊最
多要放 $\Theta(m)$ 個元素，而且頂點什麼時候「完成」不容易知道，求不出 $f[v]$。 只要
拜訪順序時可以用它；需要完成時間（第~22、23~章）就用虛擬碼~\ref{alg:gtr-dfs-iter}。

更要小心的是一種\textbf{錯誤}的改法：推進堆疊時就標記為已拜訪，拿出來時不再檢查。 它
看起來只是把標記提早，其實走出來的已經不是 DFS。 在 $G_2$ 上從 $A$ 出發，它一開始就把
$A$ 的三個鄰居 $B$、$D$、$E$ 全部標記並推進堆疊；之後從 $B$ 走到 $C$，退回來時 $D$ 已經
在堆疊裡、被標記過了，於是依序拿出 $D$、$E$，得到拜訪順序 $A, B, C, D, E$。 真正的 DFS
從 $B$ 會先走到 $E$、再從 $E$ 走到 $D$，順序是 $A, B, C, E, D$，而且 $D$ 是 $E$ 的子節點，
不是 $A$ 的子節點。 標記的時機改變了「誰先發現誰」，連 DFS 樹都不一樣了。

第~\ref{sec:gtr-c}~節的程式裡，\texttt{dfs\_visit} 是遞迴版，\texttt{dfs\_iter} 是虛擬碼
\ref{alg:gtr-dfs-iter}，\texttt{stack[0..top-1]} 是灰色的頂點，\texttt{cur[u]} 就是
$\mathit{next}[u]$。 程式用兩個版本各跑一次 $G_2$，時間戳記相同；最後在一條 $10^6$ 個頂點的
鏈上只跑迴圈版，這條鏈若用遞迴版，在一般的堆疊大小下會當掉。

\section{邊的分類}
\label{sec:gtr-class}

\subsection{四種邊}

DFS 跑完之後，相對於得到的 DFS 森林，圖的每一條邊 $(u, v)$ 都可以歸成下面四類之一：

\begin{itemize}
\item \textbf{樹邊}（tree edge）：DFS 森林裡的邊，$v$ 是經由這條邊第一次被發現的，
$\pi[v] = u$。
\item \textbf{後向邊}（back edge）：從 $u$ 指回它的祖先 $v$（有自環的圖，自環也算後向
邊）。
\item \textbf{前向邊}（forward edge）：不是樹邊，從 $u$ 指向它的子孫 $v$。
\item \textbf{交叉邊}（cross edge）：其他所有的邊。 $u$、$v$ 互相都不是祖先和子孫，可
能在同一棵樹的不同分支，也可能在不同的樹。
\end{itemize}

分類不必等 DFS 跑完再對照森林。 \proc{DFS-Visit}$(u)$ 第一次看到邊 $(u, v)$ 時（虛擬碼
\ref{alg:gtr-dfs} 第 13 行），$v$ 的顏色就決定了種類：

\begin{center}
\begin{tabular}{lll}
\toprule
看到 $(u, v)$ 時 $v$ 的顏色 & 種類 & 時間戳記的關係 \\
\midrule
白色 & 樹邊 & $d[u] < d[v] < f[v] < f[u]$ \\
灰色 & 後向邊 & $d[v] < d[u] < f[u] < f[v]$ \\
黑色，而且 $d[u] < d[v]$ & 前向邊 & $d[u] < d[v] < f[v] < f[u]$ \\
黑色，而且 $d[u] > d[v]$ & 交叉邊 & $d[v] < f[v] < d[u] < f[u]$ \\
\bottomrule
\end{tabular}
\end{center}

\begin{proof}[規則為什麼對]
白色：$v$ 在這時被發現，$\pi[v] = u$，是樹邊。 灰色：灰色的頂點是從根到 $u$ 的樹路徑，
$v$ 在上面，是 $u$ 的祖先，所以是後向邊。 黑色時 $v$ 已經完成，$u$ 還沒完成，所以
$f[v] < f[u]$。 若 $d[u] < d[v]$，由括號定理 $v$ 的區間包在 $u$ 的裡面，$v$ 是 $u$ 的子孫；
這條邊不是樹邊（$v$ 不是經由它被發現的），所以是前向邊。 若 $d[v] < d[u]$，$v$ 的區間不
可能包住 $u$ 的（$f[v] < f[u]$），只能完全分開，$f[v] < d[u]$，兩者互不為祖先子孫，是交叉
邊。 最右欄就是由這些推論加上括號定理得到的。
\end{proof}

從最右欄可以讀出一個常用的事實：樹邊、前向邊、交叉邊都有 $f[v] < f[u]$，\textbf{只有後
向邊}有 $f[v] > f[u]$。 也就是說，沿著一條邊走，完成時間會變小，除非那是後向邊。 第~22
章的「依完成時間由大到小排列就是拓樸順序」就是從這裡來的。

交叉邊的方向也有規律：$d[v] < f[v] < d[u]$，交叉邊總是從比較晚發現的頂點，指向比較早就
已經完成的頂點。 畫 DFS 森林時，如果每棵樹的子節點由左到右照發現的順序排、先長出的樹畫在左
邊，交叉邊一律從右指向左。

\subsection{$G_2$ 的分類}

照第~\ref{sec:gtr-dfs-trace}~節那次 DFS，依照邊被第一次看到的順序：

\begin{center}
\begin{tabular}{clll}
\toprule
順序 & 邊 & 看到時 & 種類 \\
\midrule
1 & $(A, B)$ & $B$ 是白色 & 樹邊 \\
2 & $(B, C)$ & $C$ 是白色 & 樹邊 \\
3 & $(C, A)$ & $A$ 是灰色 & 後向邊 \\
4 & $(B, E)$ & $E$ 是白色 & 樹邊 \\
5 & $(E, D)$ & $D$ 是白色 & 樹邊 \\
6 & $(D, C)$ & $C$ 是黑色，$d[D] = 6 > d[C] = 3$ & 交叉邊 \\
7 & $(A, D)$ & $D$ 是黑色，$d[A] = 1 < d[D] = 6$ & 前向邊 \\
8 & $(A, E)$ & $E$ 是黑色，$d[A] = 1 < d[E] = 5$ & 前向邊 \\
9 & $(F, G)$ & $G$ 是白色 & 樹邊 \\
10 & $(G, C)$ & $C$ 是黑色，$d[G] = 12 > d[C] = 3$ & 交叉邊 \\
11 & $(G, H)$ & $H$ 是白色 & 樹邊 \\
12 & $(H, F)$ & $F$ 是灰色 & 後向邊 \\
13 & $(F, H)$ & $H$ 是黑色，$d[F] = 11 < d[H] = 13$ & 前向邊 \\
\bottomrule
\end{tabular}
\end{center}

共 $6$ 條樹邊、$2$ 條後向邊、$3$ 條前向邊、$2$ 條交叉邊，和第~\ref{sec:gtr-c}~節程式印
出的一樣。 樹邊有 $n - 2 = 6$ 條，因為森林有兩棵樹。 $(D, C)$ 是同一棵樹裡不同分支之間的交叉邊，
$(G, C)$ 是從第二棵樹指向第一棵樹的交叉邊，兩條都從晚發現的頂點指向早已完成的頂點。

\begin{center}
\begin{tikzpicture}[x=1.5cm, y=1.35cm]
\node[gtr-v] (A) at (0, 0) {$A$};
\node[gtr-v] (B) at (1.6, 1) {$B$};
\node[gtr-v] (C) at (3.2, 1) {$C$};
\node[gtr-v] (E) at (1.6, -1) {$E$};
\node[gtr-v] (D) at (3.2, -1) {$D$};
\node[gtr-v] (G) at (4.8, 1) {$G$};
\node[gtr-v] (F) at (4.8, -1) {$F$};
\node[gtr-v] (H) at (6.2, 0) {$H$};
\draw[gtr-treea] (A) -- node[gtr-k, pos=0.45, above left] {T} (B);
\draw[gtr-treea] (B) -- node[gtr-k, below] {T} (C);
\draw[gtr-back] (C) to[bend right=40] node[gtr-k, above] {B} (A);
\draw[gtr-fwd] (A) -- node[gtr-k, pos=0.45, above right] {F} (E);
\draw[gtr-treea] (B) -- node[gtr-k, right] {T} (E);
\draw[gtr-treea] (E) -- node[gtr-k, above] {T} (D);
\draw[gtr-cross] (D) -- node[gtr-k, right] {C} (C);
\draw[gtr-fwd] (A) to[bend right=40] node[gtr-k, below] {F} (D);
\draw[gtr-treea] (F) -- node[gtr-k, left] {T} (G);
\draw[gtr-cross] (G) -- node[gtr-k, above] {C} (C);
\draw[gtr-treea] (G) -- node[gtr-k, above right] {T} (H);
\draw[gtr-fwd] (F) to[bend left=14] node[gtr-k, above left] {F} (H);
\draw[gtr-back] (H) to[bend left=14] node[gtr-k, below right] {B} (F);
\node[gtr-t, left=2pt of A] {1/10};
\node[gtr-t, left=2pt of B] {2/9};
\node[gtr-t, above right=0pt and 0pt of C] {3/4};
\node[gtr-t, left=2pt of E] {5/8};
\node[gtr-t, below right=1pt and 0pt of D] {6/7};
\node[gtr-t, above=2pt of G] {12/15};
\node[gtr-t, below=2pt of F] {11/16};
\node[gtr-t, right=2pt of H] {13/14};
\end{tikzpicture}
\diagcaption{$G_2$ 的邊分類，頂點旁是 $d/f$。 粗實線 T 是樹邊，虛線 B 是後向邊，點線 F 是
前向邊，點虛線 C 是交叉邊。}
\label{fig:gtr-class}
\end{center}

\subsection{無向圖只有樹邊和後向邊}

無向圖的邊 $\{u, v\}$ 在 $\mathit{Adj}[u]$ 和 $\mathit{Adj}[v]$ 各出現一次，DFS 會從兩端各看
到一次。 分類時\textbf{以第一次看到的方向為準}：先從 $u$ 那一端看到，就當成 $(u, v)$ 來
分類，之後從 $v$ 那一端再看到時不再算。

\begin{theorem}
\label{thm:gtr-undirected}
對無向圖做 DFS，每一條邊不是樹邊就是後向邊。
\end{theorem}

\begin{proof}
設 $\{u, v\}$ 是一條邊，不妨設 $d[u] < d[v]$。 在時間 $d[u]$，$v$ 是白色，邊 $\{u, v\}$ 本身
就是一條白色路徑，由白路徑定理 $v$ 是 $u$ 的子孫，所以 $v$ 的整個處理過程（從 $d[v]$ 到
$f[v]$）發生在 $u$ 還是灰色的時候。 如果這條邊第一次是從 $u$ 那端看到的，那時 $v$ 一定還
是白色：$v$ 若已經被發現，在 $v$ 完成之前就會掃到 $\mathit{Adj}[v]$ 裡的 $u$，第一次就是
從 $v$ 那端了。 所以它是樹邊。 如果第一次是從 $v$ 那端看到的，那時 $u$ 是灰色（$v$ 還在處
理中，$u$ 是它的祖先），所以是後向邊。
\end{proof}

實作時要注意：從 $v$ 看到 $\mathit{Adj}[v]$ 裡的父節點 $\pi[v]$ 時，$\pi[v]$ 是灰色，但這是
那條樹邊的第二次，不是後向邊；從 $u$ 看到已經是黑色的子孫時，那是後向邊的第二次。 第~\ref{sec:gtr-c}~節
的 \texttt{dfs\_visit} 在無向圖遇到這兩種情形就跳過（因為本書的圖沒有重
邊，「$v$ 是 $\pi[u]$」只會是樹邊的另一端）。

$G_1$ 的 DFS（外層迴圈依編號，串列由小到大）畫成圖~\ref{fig:gtr-dfs-g1}，頂點旁是 $d/f$。

\begin{center}
\begin{tikzpicture}[x=1cm, y=0.85cm]
\foreach \v/\y/\a/\b in {1/0/1/14, 2/-1/2/13, 4/-2/3/12, 3/-3/4/11, 5/-4/5/10, 6/-5/6/9, 7/-6/7/8} {
    \node[gtr-v] (w\v) at (0, \y) {\v};
    \node[gtr-t, anchor=east] at ([xshift=-2pt]w\v.west) {\a/\b};
}
\foreach \a/\b in {1/2, 2/4, 4/3, 3/5, 5/6, 6/7}
    \draw[gtr-tree] (w\a) -- (w\b);
\draw[gtr-back] (w3) to[out=20, in=-20, looseness=1.1] (w1);
\draw[gtr-back] (w6) to[out=15, in=-15, looseness=1.3] (w4);
\draw[gtr-back] (w7) to[out=30, in=-30, looseness=1.4] (w5);
\begin{scope}[xshift=3.2cm]
\foreach \v/\y/\a/\b in {8/0/15/20, 9/-1/16/19, 10/-2/17/18} {
    \node[gtr-v] (w\v) at (0, \y) {\v};
    \node[gtr-t, anchor=east] at ([xshift=-2pt]w\v.west) {\a/\b};
}
\draw[gtr-tree] (w8) -- (w9);
\draw[gtr-tree] (w9) -- (w10);
\end{scope}
\end{tikzpicture}
\diagcaption{$G_1$ 的 DFS 森林。 左邊的樹剛好是一條路徑 $1, 2, 4, 3, 5, 6, 7$；虛線是三條
後向邊 $(3, 1)$、$(6, 4)$、$(7, 5)$，都從子孫指回祖先。 沒有前向邊和交叉邊。}
\label{fig:gtr-dfs-g1}
\end{center}

$11$ 條邊裡 $8$ 條是樹邊（$n$ 減去樹的個數 $2$），$3$ 條是後向邊。 從 $4$ 出發時，
$\mathit{Adj}[4] = \langle 2, 3, 6 \rangle$，$2$ 是灰色的父節點，跳過；$3$ 是白色，所以先
走到 $3$，而不是 $6$。 這和 BFS 樹完全不同：在 BFS 樹裡 $3$ 和 $2$ 同樣離 $1$ 一步，在 DFS
樹裡 $3$ 卻在深度 $3$。

\begin{examnote}
邊的分類題，答案寫成「邊、種類」的表，每條邊只出現一次；無向圖要先說明每條邊依第一次
被看到的方向分類。 時間戳記照本書從 $1$ 開始、發現和完成都加一；有的課本只記發現的
順序（有時從 $0$ 開始），題目給了定義就照題目。 外層迴圈從哪個頂點開始、串列的順序，題
目沒說就寫出你的假設（例如「依字母順序，串列由小到大」），因為換一個順序，樹邊、前向
邊、交叉邊的歸類都可能改變；但「有沒有後向邊」不會變，下一節會證明它等於「有沒有環」。
\end{examnote}

\section{應用：連通元件、二分圖、有向環}
\label{sec:gtr-app}

\subsection{連通元件}

無向圖從頂點 $s$ 做一次 DFS 或 BFS，走到的頂點正好是 $s$ 所在的連通元件：可到達的都走
得到（白路徑定理，或 BFS 的定理~\ref{thm:gtr-bfs}），走不到的不在同一個元件。 所以和 \proc{DFS} 的外層迴圈一樣，
依序檢查每個頂點，遇到還沒走過的頂點 $u$，就把計數 $k$ 加一，從 $u$ 做一次
\proc{DFS-Visit}（或 BFS），把走到的每個頂點 $v$ 記成 $\mathit{comp}[v] = k$。
時間和 DFS 一樣是 $\Theta(n + m)$，$k$ 最後就是元件個數。 之後問「$u$、
$v$ 是否連通」只要比較 $\mathit{comp}[u] = \mathit{comp}[v]$，$\Theta(1)$。 $G_1$ 的結果是
$\mathit{comp}[1..7] = 1$、$\mathit{comp}[8..10] = 2$，$k = 2$。 用 DFS 或 BFS 都可以，第~\ref{sec:gtr-c}~節的程式就是用 BFS 做的；兩者結果相同，BFS 不會有遞迴太深的問題。

第~13~章用互斥集合也能回答連通的問題，兩種做法各有適合的場合：

\begin{center}
\small
\begin{tabular}{>{\raggedright\arraybackslash}p{0.25\linewidth}>{\raggedright\arraybackslash}p{0.32\linewidth}>{\raggedright\arraybackslash}p{0.32\linewidth}}
\toprule
 & DFS／BFS & 互斥集合（第~13~章） \\
\midrule
需要先有什麼 & 整張圖的鄰接串列 & 只要邊，一條一條讀進來就好 \\
全部算完的時間 & $\Theta(n + m)$ & $O(n + m\,\alpha(n))$ \\
額外的記憶體 & 鄰接串列 $\Theta(n + m)$ & 父節點和大小陣列 $\Theta(n)$ \\
邊陸續加入，中間隨時要回答 & 每次重算 $\Theta(n + m)$ & 每條邊、每個詢問平攤 $O(\alpha(n))$ \\
有刪邊 & 刪掉之後重算 & 不支援；事先知道順序時倒著做（範例~\ref{ex:ds-lab-cut}）\\
順便得到 & 元件的成員、路徑、距離、生成樹 & 只有「是否同組」和組的大小 \\
有向圖 & 可以再發展成強連通元件（第~23~章）& 不適用 \\
\bottomrule
\end{tabular}
\end{center}

圖一次給齊、還要做別的事（找路徑、算距離）時用 DFS 或 BFS；邊一條一條來、要隨時回答，
或圖大到不想存鄰接串列時，用互斥集合。 兩者的關係在引理~\ref{lem:gtr-count} 已經看過：
互斥集合每做一次真正的合併，就相當於 DFS 森林裡長出一條樹邊。

\subsection{二分圖}

學期末要把幾門課的考試排進兩個時段，同一個學生修的兩門課不能排在同一個時段。 把課當
頂點、「有人同時修」當邊，問題就是：能不能把頂點分成兩組，讓每條邊的兩端在不同組？

\begin{definition}[二分圖]
\label{def:gtr-bipartite}
無向圖 $G = (V, E)$ 是\textbf{二分圖}（bipartite graph），如果 $V$ 可以分成兩個不相交的
集合 $L$、$R$，使得每條邊都是一端在 $L$、一端在 $R$。 等價的說法是：可以用兩種顏色替
頂點著色，讓每條邊兩端的顏色不同。
\end{definition}

沿著一個環走，每走一條邊就換一邊，回到起點時必須換了偶數次，所以二分圖不可能有長度是奇
數的環（\textbf{奇環}）。 反過來也對，而且 BFS 能直接判斷：

\begin{theorem}
\label{thm:gtr-bipartite}
無向圖是二分圖的充分必要條件是它沒有奇環。 對每個連通元件從任一頂點做 BFS，圖是二分圖
的充分必要條件是：沒有任何一條邊的兩端 $d$ 值相同。 這時把 $d$ 是偶數的頂點放 $L$、奇
數的放 $R$，就是一種分法。
\end{theorem}

\begin{proof}
二分圖沒有奇環，上面已經說明。 接著在每個連通元件做 BFS，設起點是 $s$。 由
第~\ref{sec:gtr-bfs-proof}~節的推論，無向圖每條邊兩端的 $d$ 最多差 $1$。

若沒有一條邊兩端的 $d$ 相同，每條邊兩端的 $d$ 就剛好差 $1$，一奇一偶，依奇偶分成 $L$、$R$
就符合定義，圖是二分圖。

若某條邊 $\{u, v\}$ 有 $d[u] = d[v]$，在 BFS 樹裡從 $u$ 和 $v$ 同時沿著 $\pi$ 往上走，每
走一步兩邊的 $d$ 都減一、仍然相等，第一次走到同一個頂點 $w$ 時停下。 樹路徑 $u \leadsto
w$ 和 $v \leadsto w$ 除了 $w$ 之外沒有共同的頂點，長度都是 $d[u] - d[w]$，再加上邊
$\{u, v\}$，得到一個長度 $2(d[u] - d[w]) + 1$ 的環，是奇環，所以圖不是二分圖。 兩種情形
合起來，也證明了「沒有奇環就是二分圖」。
\end{proof}

\begin{pseudo}[label=alg:gtr-bip]{判斷二分圖（BFS 著色）}
\Function{Is-Bipartite}{$G$}
    \For{每個頂點 $u \in V$}
        \State $\mathit{side}[u] \gets \NIL$
    \EndFor
    \For{每個頂點 $s \in V$}
        \If{$\mathit{side}[s] = \NIL$} \Comment{$s$ 開始一個新的連通元件}
            \State $\mathit{side}[s] \gets 0$
            \State $Q \gets$ 空佇列；\Call{Enqueue}{$Q, s$}
            \While{$Q$ 不是空的}
                \State $u \gets \Call{Dequeue}{Q}$
                \For{每個 $v \in \mathit{Adj}[u]$}
                    \If{$\mathit{side}[v] = \NIL$}
                        \State $\mathit{side}[v] \gets 1 - \mathit{side}[u]$
                        \State \Call{Enqueue}{$Q, v$}
                    \ElsIf{$\mathit{side}[v] = \mathit{side}[u]$}
                        \State \Return \FALSE \Comment{同色的鄰居：有奇環}
                    \EndIf
                \EndFor
            \EndWhile
        \EndIf
    \EndFor
    \State \Return \TRUE
\EndFunction
\end{pseudo}

$\mathit{side}[v]$ 就是 $d[v]$ 的奇偶，所以第 14 行的條件「同色」等於「$d$ 相同」（$d$ 最
多差一）。 時間 $\Theta(n + m)$。 這裡用 DFS 也可以，只要讓子節點的顏色和父節點相反；同色
的鄰居一樣代表奇環。

用 $G_1$ 試一次。 第~\ref{sec:gtr-bfs-trace}~節的 BFS 裡，$6$ 和 $7$ 的距離都是 $3$，邊
$\{6, 7\}$ 讓圖不是二分圖。 往上找共同的祖先：$6 \to 4 \to 2 \to 1$、$7 \to 5 \to 3 \to 1$，
得到奇環 $\langle 6, 4, 2, 1, 3, 5, 7, 6 \rangle$，長度 $7$。 圖裡其實還有更短的奇環
$\langle 5, 6, 7, 5 \rangle$；BFS 保證找到的是奇環，不保證是最短的。 另一個元件
$\{8, 9, 10\}$ 自己是二分圖：$L = \{8, 10\}$、$R = \{9\}$。

第~\ref{sec:gtr-c}~節的程式把兩件事合在同一次 BFS 裡：外層迴圈替元件編號，同時記住第一
條兩端 $d$ 相同的邊。

\subsection{有向環}

章首的網站管理員想知道「一直點連結會不會繞回同一頁」，就是問有向圖有沒有環。 第~22~
章的拓樸排序、課程先修關係的檢查，也都要先確定圖是 DAG。

\begin{theorem}
\label{thm:gtr-cycle}
有向圖 $G$ 有環的充分必要條件是：對 $G$ 做 DFS 會出現後向邊。 不管外層迴圈和鄰接串列的順
序怎麼排，結論都一樣。
\end{theorem}

\begin{proof}
（$\Leftarrow$）後向邊 $(u, v)$ 的 $v$ 是 $u$ 的祖先（或 $u$ 自己），樹路徑 $v \leadsto u$
再接上 $(u, v)$ 就是一個環。

（$\Rightarrow$）設 $c$ 是一個環。 取 $c$ 上最早被發現的頂點 $v$，$u$ 是環上 $v$ 的前一個
頂點，$(u, v) \in E$。 在時間 $d[v]$，環上其他頂點都還是白色，沿著環從 $v$ 走到 $u$ 是一
條白色路徑，由白路徑定理 $u$ 是 $v$ 的子孫。 $u$ 處理 $(u, v)$ 時，$v$ 還沒完成（$u$ 的區
間在 $v$ 的裡面），是灰色，所以 $(u, v)$ 是後向邊。
\end{proof}

偵測時不需要分出四種邊，只要看「有沒有遇到灰色的鄰居」，遇到就可以停：

\begin{pseudo}[label=alg:gtr-cycle]{判斷有向圖有沒有環}
\Function{Has-Cycle}{$G$}
    \For{每個頂點 $u \in V$}
        \State $\mathit{color}[u] \gets \gtrwhite$
    \EndFor
    \For{每個頂點 $u \in V$}
        \If{$\mathit{color}[u] = \gtrwhite$ \AND \Call{Visit}{$G, u$}}
            \State \Return \TRUE
        \EndIf
    \EndFor
    \State \Return \FALSE
\EndFunction
\Function{Visit}{$G, u$} \Comment{從 $u$ 出發找得到環就回傳 \TRUE}
    \State $\mathit{color}[u] \gets \gtrgray$
    \For{每個 $v \in \mathit{Adj}[u]$}
        \If{$\mathit{color}[v] = \gtrgray$}
            \State \Return \TRUE \Comment{後向邊 $(u, v)$}
        \EndIf
        \If{$\mathit{color}[v] = \gtrwhite$ \AND \Call{Visit}{$G, v$}}
            \State \Return \TRUE
        \EndIf
    \EndFor
    \State $\mathit{color}[u] \gets \gtrblack$
    \State \Return \FALSE
\EndFunction
\end{pseudo}

時間 $O(n + m)$。 $G_2$ 的 DFS 在時間 $4$ 之前就遇到後向邊 $(C, A)$，對應的環是樹路徑
$A \to B \to C$ 加上 $C \to A$；另一條後向邊 $(H, F)$ 對應環 $F \to G \to H \to F$。 第~\ref{sec:gtr-c}~節的程式印出的 \texttt{cycle: yes} 就是數到了後向邊。

為什麼一定要分灰和黑？ 黑色的鄰居不代表有環。 $G_2$ 的 $D \to C$ 碰到的 $C$ 是黑色，那
是交叉邊；把 $C \to A$ 和 $H \to F$ 刪掉之後圖沒有環，DFS 照樣會從 $D$ 看到黑色的 $C$。
同樣的道理，BFS 在有向圖裡看到「已經發現過的頂點」也不能斷定有環：只有 $A \to B$、
$A \to C$、$B \to C$ 三條邊的圖沒有環，但從 $A$ 做 BFS，拿出 $B$ 時會看到灰色的 $C$。 用
BFS 的想法判斷有向環，是第~22~章 Kahn 演算法的做法，要用到入度，而不是顏色。

無向圖的環比較單純。 由定理~\ref{thm:gtr-undirected}，DFS 只有樹邊和後向邊，每條後向
邊都和樹路徑圍出一個環，所以無向圖有環的充分必要條件也是有後向邊；實作時記得排除「看到
父節點」那一次。 更簡單的判斷是數邊：由引理~\ref{lem:gtr-count}，有 $c$ 個連通元件的無
向圖沒有環，若且唯若 $m = n - c$。 $G_1$ 有 $m = 11 > 10 - 2$，所以有環，而且恰好有
$11 - 8 = 3$ 條後向邊。

\subsection{方格也是圖}

很多問題的圖沒有明白給出來。 一張 $r$ 列 $c$ 行的方格地圖，每一格是一個頂點，上下左右
相鄰的兩格之間有邊（有的題目連斜角也算，每格最多八個鄰居），牆壁的格子不算頂點。 這種
圖不必建鄰接串列：格子 $(i, j)$ 的鄰居用兩個位移陣列 $\mathit{di} = \langle -1, 1, 0, 0
\rangle$、$\mathit{dj} = \langle 0, 0, -1, 1 \rangle$ 當場算出來，檢查沒有出界、不是牆壁
就好。 顏色、距離存在 $r \times c$ 的二維陣列，或用第~4~章的列優先公式攤成一維。

這時 $n = rc$，四方向的邊數最多大約 $2rc$，BFS、DFS 都是 $\Theta(rc)$。 常見的問題都是前
面講過的東西換個樣子：從一格出發能走到幾格，是一次 DFS 或 BFS；地圖上有幾塊互不相連的
區域、最大的一塊有多大，是數連通元件；每一格離最近的出口多遠，是第~\ref{sec:gtr-multi-source}~節
的多起點 BFS。 要注意遞迴深度：$1000 \times 1000$ 的方格，DFS 可能一路彎彎曲曲走過大部分
的格子，遞迴深度到幾十萬層，這時用 BFS 或第~\ref{sec:gtr-dfs-iter}~節的迴圈版。

\section{C 實作}
\label{sec:gtr-c}

這一節把本章的資料結構和演算法寫成一支完整的 C 程式：鄰接串列（第~\ref{sec:gtr-repr}~節）、
BFS 和印出路徑（第~\ref{sec:gtr-bfs}~節）、遞迴和迴圈兩種 DFS 與邊的分類（第
\ref{sec:gtr-dfs}、\ref{sec:gtr-class}~節）、連通元件和二分圖檢查（第~\ref{sec:gtr-app}~節）。
每一部分前面的註解寫了對應的小節。

\begin{cbookrecap}[第~14~章 14.5 節、第~19~章 19.3、19.9--19.11 節]
\texttt{AdjNode **pp} 是「指向某個 \texttt{next} 欄位（或指向 \texttt{adj[u]}）」的指
標。 \texttt{*pp} 是那個欄位目前指的節點，\texttt{pp = \&(*pp)->next} 往下走一格；
插入時 \texttt{x->next = *pp; *pp = x;}，不管插在開頭還是中間都是這兩行。
\texttt{calloc(n, sizeof *g->adj)} 配置 $n$ 個指標並全部設成 $0$，也就是 $n$ 條空串列。
釋放時每條串列要一個節點一個節點 \texttt{free}，先記下 \texttt{next} 再釋放目前的節點，
最後才釋放 \texttt{adj} 陣列本身。
\end{cbookrecap}

輸入依序是兩張圖。 每張圖的第一行是 $n$、$m$ 和一個旗標（$1$ 是有向圖，$0$ 是無向圖），
接著 $m$ 行各一條邊；頂點寫成大寫字母或從 $1$ 開始的數字，\texttt{read\_vertex} 把它換成
$0..n-1$，\texttt{put} 印的時候再換回來。 第一張圖後面多一行，是 BFS 的起點和要印路徑的終
點。 為了讓程式短一點，搜尋用的四個陣列 \texttt{color}、\texttt{d}、\texttt{f}、\texttt{pi}
和 DFS 的時間 \texttt{now} 放成全域變數，每次搜尋前由 \texttt{setup} 配置一塊記憶體切成四段、設成
白色和 \NIL，用完 \texttt{free(color)} 一次釋放。 下面用 $G_1$（從 $1$ 到 $7$）和 $G_2$ 當輸入。

\begin{example}[label=ex:gtr-graph]{c}{(graph.c) 鄰接串列、BFS、DFS、邊的分類、連通元件與二分圖}
#include <ctype.h>
#include <stdio.h>
#include <stdlib.h>

enum { WHITE, GRAY, BLACK };

typedef struct AdjNode {
    int v;                      // 鄰居的編號 0..n-1
    struct AdjNode *next;
} AdjNode;

typedef struct {
    int n, directed;            // directed：1 是有向圖，0 是無向圖
    AdjNode **adj;              // adj[u]：u 的鄰接串列，由小到大
} Graph;

// ---- 鄰接串列（21.2 節）----

// 把 v 插進 *pp 開頭的串列，保持由小到大。
// 回傳 1：插入了；0：已經有 v（重邊）；-1：malloc 失敗。
int list_insert(AdjNode **pp, int v)
{
    while (*pp != NULL && (*pp)->v < v)
        pp = &(*pp)->next;
    if (*pp != NULL && (*pp)->v == v)
        return 0;
    AdjNode *x = malloc(sizeof *x);
    if (x == NULL)
        return -1;
    x->v = v;
    x->next = *pp;
    *pp = x;
    return 1;
}

Graph *graph_new(int n, int directed)
{
    Graph *g = malloc(sizeof *g);
    if (g != NULL && (g->adj = calloc(n, sizeof *g->adj)) == NULL) {
        free(g);
        g = NULL;
    }
    if (g != NULL) {
        g->n = n;
        g->directed = directed;
    }
    return g;
}

// 加一條邊 (u, v)；自環和重邊不加。 失敗回傳 0。
int graph_add_edge(Graph *g, int u, int v)
{
    if (u == v)
        return 1;
    int r = list_insert(&g->adj[u], v);
    if (r == 1 && !g->directed)
        r = list_insert(&g->adj[v], u);
    return r >= 0;
}

void graph_free(Graph *g)
{
    for (int u = 0; u < g->n; u++) {
        AdjNode *p = g->adj[u];
        while (p != NULL) {
            AdjNode *next = p->next;
            free(p);
            p = next;
        }
    }
    free(g->adj);
    free(g);
}

// 頂點的名字：大寫字母 A, B, ... 或數字 1, 2, ...；程式裡一律是 0..n-1。
int letters = 0;

int read_vertex(int n, int *v)
{
    char s[16];
    if (scanf("%15s", s) != 1)
        return 0;
    letters = isalpha((unsigned char)s[0]) != 0;
    *v = letters ? s[0] - 'A' : atoi(s) - 1;
    return *v >= 0 && *v < n;
}

void put(int v)                 // 印一個空白和頂點的名字；-1 印成 NIL
{
    if (v < 0)
        printf(" NIL");
    else if (letters)
        printf(" %c", 'A' + v);
    else
        printf(" %d", v + 1);
}

// 輸入格式：n m directed，接著 m 行，每行一條邊 u v。
Graph *read_graph(void)
{
    int n, m, directed;
    if (scanf("%d %d %d", &n, &m, &directed) != 3 || n <= 0)
        return NULL;
    Graph *g = graph_new(n, directed);
    for (int i = 0; g != NULL && i < m; i++) {
        int u, v;
        if (!read_vertex(n, &u) || !read_vertex(n, &v)
            || !graph_add_edge(g, u, v)) {
            graph_free(g);
            g = NULL;
        }
    }
    return g;
}

// ---- 搜尋共用的陣列：color、d、f、pi（-1 代表 NIL）----
int *color, *d, *f, *pi, now;   // now：DFS 的全域時間

// ---- BFS（21.3 節）：d[v] = -1 代表無限大 ----
// 佇列是長度 n 的陣列 q[head..tail-1]：每個頂點最多進佇列一次。
// trace 是 1 時，每拿出一個頂點就印出新發現的頂點和之後的佇列。
// 第一條「兩端距離相同」的邊記在 *bu、*bv（二分圖檢查用，21.6 節）。
int bfs(const Graph *g, int s, int *q, int trace, int *bu, int *bv)
{
    int head = 0, tail = 0;
    color[s] = GRAY;
    d[s] = 0;
    pi[s] = -1;
    q[tail++] = s;
    while (head < tail) {
        int u = q[head++];
        if (trace) {
            printf("dequeue");
            put(u);
            printf(" ->");
        }
        for (AdjNode *p = g->adj[u]; p != NULL; p = p->next) {
            int v = p->v;
            if (color[v] == WHITE) {
                color[v] = GRAY;
                d[v] = d[u] + 1;
                pi[v] = u;
                q[tail++] = v;
                if (trace)
                    put(v);
            } else if (d[v] == d[u] && *bu < 0) {
                *bu = u;
                *bv = v;
            }
        }
        color[u] = BLACK;
        if (trace) {
            printf(" | queue");
            for (int i = head; i < tail; i++)
                put(q[i]);
            printf(head < tail ? "\n" : " (empty)\n");
        }
    }
    return tail;                // 走到幾個頂點；它們在 q[0..tail-1]
}

void print_path(int s, int v)   // 虛擬碼的 Print-Path
{
    if (v == s)
        put(s);
    else if (pi[v] < 0)
        printf(" no path");
    else {
        print_path(s, pi[v]);
        put(v);
    }
}

// ---- DFS（21.4、21.5 節）----
int back_edges;

// 遞迴版，順便替每條邊分類。
void dfs_visit(const Graph *g, int u)
{
    d[u] = ++now;
    color[u] = GRAY;
    for (AdjNode *p = g->adj[u]; p != NULL; p = p->next) {
        int v = p->v;
        const char *kind;
        if (color[v] == WHITE)
            kind = "tree";
        else if (!g->directed && (v == pi[u] || color[v] == BLACK))
            continue;           // 無向圖：這條邊從另一端分類過了
        else if (color[v] == GRAY)
            kind = "back";
        else
            kind = d[u] < d[v] ? "forward" : "cross";
        printf("edge");
        put(u);
        put(v);
        printf(" %s\n", kind);
        back_edges += color[v] == GRAY;
        if (color[v] == WHITE) {
            pi[v] = u;
            dfs_visit(g, v);
        }
    }
    color[u] = BLACK;
    f[u] = ++now;
}

// 迴圈版（虛擬碼 21.4）：stack[0..top-1] 是灰色的頂點，
// cur[u] 是 u 的串列裡下一個要看的節點。 回傳堆疊最深幾層。
int dfs_iter(const Graph *g, int s, int *stack, AdjNode **cur)
{
    int top = 0, deepest = 1;
    d[s] = ++now;
    color[s] = GRAY;
    cur[s] = g->adj[s];
    stack[top++] = s;
    while (top > 0) {
        int u = stack[top - 1];
        if (cur[u] != NULL) {               // 還有鄰居沒看
            int v = cur[u]->v;
            cur[u] = cur[u]->next;
            if (color[v] == WHITE) {
                pi[v] = u;
                d[v] = ++now;
                color[v] = GRAY;
                cur[v] = g->adj[v];
                stack[top++] = v;
                if (top > deepest)
                    deepest = top;
            }
        } else {                            // u 的鄰居都看完了
            top--;
            color[u] = BLACK;
            f[u] = ++now;
        }
    }
    return deepest;
}

// 配置搜尋用的四個陣列（同一塊記憶體切成四段），全部設成白色、NIL。
int setup(int n)
{
    if ((color = malloc(4 * (size_t)n * sizeof *color)) == NULL)
        return 0;
    d = color + n;
    f = d + n;
    pi = f + n;
    for (int v = 0; v < n; v++) {
        color[v] = WHITE;
        d[v] = pi[v] = -1;
    }
    now = 0;
    return 1;
}

// 印出一個陣列，-1 印成 inf（距離）或 NIL（前驅）。
void print_row(const char *name, const int *a, int n, int is_vertex)
{
    printf("%s:", name);
    for (int v = 0; v < n; v++) {
        if (is_vertex)
            put(a[v]);
        else if (a[v] < 0)
            printf(" inf");
        else
            printf(" %d", a[v]);
    }
    printf("\n");
}

int main(void)
{
    // 第一張圖（G1）：從 s 做 BFS、印出到 t 的路徑；再找連通元件。
    Graph *g = read_graph();
    int s, t, bu = -1, bv = -1;
    if (g == NULL || !read_vertex(g->n, &s) || !read_vertex(g->n, &t)
        || !setup(g->n))
        return 1;
    int *q = malloc(g->n * sizeof *q);
    if (q == NULL)
        return 1;
    bfs(g, s, q, 1, &bu, &bv);
    print_row("d", d, g->n, 0);
    print_row("pi", pi, g->n, 1);
    printf("path:");
    print_path(s, t);
    printf("\n");
    free(color);
    if (!setup(g->n))
        return 1;
    bu = bv = -1;
    for (int k = 0, v = 0; v < g->n; v++) {
        if (color[v] == WHITE) {            // v 開始一個新的連通元件
            int cnt = bfs(g, v, q, 0, &bu, &bv);
            printf("component %d:", ++k);
            for (int i = 0; i < cnt; i++)
                put(q[i]);
            printf("\n");
        }
    }
    if (bu < 0) {
        printf("bipartite: yes\n");
    } else {
        printf("bipartite: no, edge");
        put(bu);
        put(bv);
        printf(" joins two vertices at distance %d\n", d[bu]);
    }
    free(color);
    free(q);
    graph_free(g);

    // 第二張圖（G2）：DFS 並分類邊，再用迴圈版跑一次。
    if ((g = read_graph()) == NULL || !setup(g->n))
        return 1;
    for (int u = 0; u < g->n; u++)          // 虛擬碼 21.3 的 DFS
        if (color[u] == WHITE)
            dfs_visit(g, u);
    print_row("d", d, g->n, 0);
    print_row("f", f, g->n, 0);
    print_row("pi", pi, g->n, 1);
    printf("cycle: %s\n", back_edges > 0 ? "yes" : "no");
    free(color);
    int *stack = malloc(g->n * sizeof *stack);
    AdjNode **cur = malloc(g->n * sizeof *cur);
    if (!setup(g->n) || stack == NULL || cur == NULL)
        return 1;
    for (int u = 0; u < g->n; u++)
        if (color[u] == WHITE)
            dfs_iter(g, u, stack, cur);
    print_row("iterative d", d, g->n, 0);
    print_row("iterative f", f, g->n, 0);
    free(stack);
    free(cur);
    free(color);
    graph_free(g);

    // 一條 10^6 個頂點的鏈 1 -> 2 -> ... -> n，只用迴圈版。
    int n = 1000000;
    g = graph_new(n, 1);
    stack = malloc(n * sizeof *stack);
    cur = malloc(n * sizeof *cur);
    if (g == NULL || stack == NULL || cur == NULL || !setup(n))
        return 1;
    for (int i = 0; i + 1 < n; i++)
        if (!graph_add_edge(g, i, i + 1))
            return 1;
    int deepest = dfs_iter(g, 0, stack, cur);
    printf("chain: n = %d, f[1] = %d, deepest stack = %d\n", n, f[0], deepest);
    free(stack);
    free(cur);
    free(color);
    graph_free(g);
    return 0;
}
\end{example}
\iosample{10 11 0 \\ 1 2 \\ 1 3 \\ 2 4 \\ 3 4 \\ 3 5 \\ 4 6 \\ 5 6 \\ 5 7 \\ 6 7 \\ 8 9 \\ 9 10 \\ 1 7 \\ 8 13 1 \\ A B \\ A D \\ A E \\ B C \\ B E \\ C A \\ D C \\ E D \\ F G \\ F H \\ G C \\ G H \\ H F}{dequeue 1 -> 2 3 | queue 2 3 \\ dequeue 2 -> 4 | queue 3 4 \\ dequeue 3 -> 5 | queue 4 5 \\ dequeue 4 -> 6 | queue 5 6 \\ dequeue 5 -> 7 | queue 6 7 \\ dequeue 6 -> | queue 7 \\ dequeue 7 -> | queue (empty) \\ d: 0 1 1 2 2 3 3 inf inf inf \\ pi: NIL 1 1 2 3 4 5 NIL NIL NIL \\ path: 1 3 5 7 \\ component 1: 1 2 3 4 5 6 7 \\ component 2: 8 9 10 \\ bipartite: no, edge 6 7 joins two vertices at distance 3 \\ edge A B tree \\ edge B C tree \\ edge C A back \\ edge B E tree \\ edge E D tree \\ edge D C cross \\ edge A D forward \\ edge A E forward \\ edge F G tree \\ edge G C cross \\ edge G H tree \\ edge H F back \\ edge F H forward \\ d: 1 2 3 6 5 11 12 13 \\ f: 10 9 4 7 8 16 15 14 \\ pi: NIL A B E B NIL F G \\ cycle: yes \\ iterative d: 1 2 3 6 5 11 12 13 \\ iterative f: 10 9 4 7 8 16 15 14 \\ chain: n = 1000000, f[1] = 2000000, deepest stack = 1000000}

對照前面的手算：\texttt{dequeue} 那七列就是第~\ref{sec:gtr-bfs-trace}~節的佇列表，\texttt{d}、
\texttt{pi} 兩列是那裡的結果表；\texttt{component} 兩列是 $G_1$ 的兩個連通元件，每一列依 BFS
發現的順序列出；\texttt{bipartite} 指出的邊 $\{6, 7\}$ 和第~\ref{sec:gtr-app}~節找到的相同。
\texttt{edge} 那十三列和第~\ref{sec:gtr-class}~節的分類表逐列相同，接著的 \texttt{d}、\texttt{f}、
\texttt{pi} 是第~\ref{sec:gtr-dfs-trace}~節的表，迴圈版印出的兩列和遞迴版一樣。

最後一列是 $10^6$ 個頂點的鏈 $1 \to 2 \to \cdots \to 10^6$。 迴圈版的堆疊最深 $10^6$ 層，用
的是 \texttt{malloc} 要來的 $4 \times 10^6$ 位元組；起點最晚完成，$f = 2n$。 同一條鏈交給
\texttt{dfs\_visit}，程式會在遞迴到十幾萬層時當掉。

幾個寫法上的細節：
\begin{itemize}
\item \texttt{bfs} 多帶兩個參數 \texttt{bu}、\texttt{bv}，記下第一條兩端距離相同的邊，求連通元件
時順便檢查二分圖。 只記第一條，BFS 照樣跑完，元件才會完整。
\item \texttt{dfs\_visit} 在無向圖看到父節點、或看到黑色的頂點時直接跳過：這兩種都是同一條邊
從另一端第二次被看到（第~\ref{sec:gtr-class}~節）。 本書的圖沒有重邊，所以「$v$ 是 $\pi[u]$」
只會是樹邊的另一端；有重邊的圖要改成記住「從哪一條邊進來」。
\item 主程式在每次配置失敗時直接結束。 正常執行時每一塊記憶體都有釋放，sanitizer 檢查不到
任何洩漏。
\end{itemize}

\section{本章常見錯誤}
\label{sec:gtr-pitfalls}

\begin{center}
\small
\begin{tabular}{p{0.27\linewidth}p{0.48\linewidth}p{0.15\linewidth}}
\toprule
錯誤 & 例子 & 見 \\
\midrule
沒照題目給的串列順序走 & 題目畫的串列不是由小到大，卻自己排序後再做 BFS、DFS &
第~\ref{sec:gtr-repr}~節 \\
BFS 拿出來才標記 & 放進佇列時沒有改顏色，同一個頂點進佇列好幾次，$d$ 被後來的值蓋掉，
佇列也超過 $n$ 格 & 第~\ref{sec:gtr-bfs}~節 \\
把 DFS 樹的路徑當最短路徑 & $G_2$ 的 DFS 樹裡 $D$ 在深度 $3$，其實 $\delta(A, D) = 1$ &
第~\ref{sec:gtr-dfs-trace}~節 \\
時間戳記算錯 & 只在發現時加一；完成時沒加；換一棵樹時把時間歸零 &
第~\ref{sec:gtr-dfs}~節 \\
DFS 改成迴圈時推進就標記 & 走出 $A, B, C, D, E$，已經不是 DFS 的順序 &
第~\ref{sec:gtr-dfs-iter}~節 \\
大圖用遞迴 DFS & 一條十幾萬個頂點的鏈就讓程式當掉 & 第~\ref{sec:gtr-dfs-iter}~節 \\
無向圖把父節點當成後向邊 & 每條樹邊都被當成環，連樹都被判成有環 &
第~\ref{sec:gtr-class}~節 \\
有向圖看到「走過的頂點」就說有環 & 黑色的鄰居是前向邊或交叉邊，BFS 看到灰色的頂點也
不代表有環 & 第~\ref{sec:gtr-app}~節 \\
只從一個起點做，就下整張圖的結論 & 從 $1$ 做一次 BFS 就說 $G_1$ 是二分圖、或只數到一個
元件；忘了外層迴圈 & 第~\ref{sec:gtr-app}~節 \\
用鄰接矩陣卻寫 $\Theta(n + m)$ & 矩陣上列出鄰居要掃整列，BFS、DFS 是 $\Theta(n^2)$ &
第~\ref{sec:gtr-repr}~節 \\
\bottomrule
\end{tabular}
\end{center}

\section{觀念題}
\label{sec:gtr-exercises}

題型和研究所考試、期中期末考的筆試相近。 沒有特別說明時，鄰接串列依編號（字母）由小到
大，DFS 的外層迴圈也依這個順序。 解答在本章最後。

\begin{enumerate}
\item \label{q:gtr-facts}（選擇，可能不只一個答案）下列敘述哪些正確？
\begin{enumerate}
\item $n$ 個頂點的簡單無向圖最多有 $n(n-1)/2$ 條邊。
\item 有向圖所有頂點的入度和等於出度和。
\item $n$ 個頂點、$n - 1$ 條邊的無向圖一定是樹。
\item 用鄰接矩陣存圖時，BFS 的時間是 $\Theta(n + m)$。
\item 無向圖裡度數是奇數的頂點一定有偶數個。
\item 對無向圖做 DFS，不會產生交叉邊。
\end{enumerate}

\item \label{q:gtr-degree}（計算）以下都是簡單無向圖。
\begin{enumerate}
\item 一張圖有 $12$ 條邊，其中 $3$ 個頂點的度數是 $4$，其他頂點的度數都是 $3$。 這張圖
有幾個頂點？
\item 有沒有一張圖，五個頂點的度數分別是 $4, 3, 3, 2, 1$？
\item $10$ 個頂點的連通圖最少有幾條邊？ $10$ 個頂點、\textbf{不連通}的圖最多有幾條邊？
\end{enumerate}

\item \label{q:gtr-repr}（表示）有向圖 $P$ 的頂點是 $1$ 到 $6$，鄰接矩陣如下（第 $u$ 列第 $v$
行是 $1$ 表示有 $u \to v$）。
\[
\begin{pmatrix}
0 & 1 & 0 & 1 & 0 & 0 \\
0 & 0 & 1 & 0 & 0 & 1 \\
1 & 0 & 0 & 0 & 1 & 0 \\
0 & 0 & 0 & 0 & 1 & 0 \\
0 & 0 & 0 & 0 & 0 & 1 \\
0 & 0 & 0 & 1 & 0 & 0
\end{pmatrix}
\]
(a) 寫出鄰接串列，以及每個頂點的入度和出度。 (b) 寫出循序表示的 $\mathit{start}[1..7]$ 和
$\mathit{nbr}$。 (c) 寫出反鄰接串列。

\item \label{q:gtr-bfs-f}（追蹤）在 $G_2$（圖~\ref{fig:gtr-examples}）上從 $F$ 做 BFS。 照
第~\ref{sec:gtr-bfs-trace}~節的格式寫出每一輪拿出的頂點和之後的佇列，以及最後的 $d$、
$\pi$。 從 $F$ 到 $E$ 的最短路徑是什麼？

\item \label{q:gtr-dfs-q}（追蹤）有向圖 $Q$ 的頂點是 $1$ 到 $7$，鄰接串列如下：
\[
1\colon 2, 4 \qquad 2\colon 3, 5 \qquad 3\colon 1 \qquad 4\colon 3 \qquad
5\colon 4, 6 \qquad 6\colon 7 \qquad 7\colon 5, 6
\]
執行 \proc{DFS}，寫出每個頂點的 $d$、$f$、$\pi$，把每條邊分類，並判斷 $Q$ 有沒有環。

\item \label{q:gtr-interval}（選擇）$u$、$v$ 是同一張有向圖的兩個頂點。 下列哪些可能是同一次
DFS 得到的 $[d[u], f[u]]$ 和 $[d[v], f[v]]$？ 可能的，說明 $u$、$v$ 在 DFS 森林裡的關係。
\begin{enumerate}
\item $[1, 6]$ 和 $[2, 5]$
\item $[2, 7]$ 和 $[5, 10]$
\item $[3, 4]$ 和 $[5, 8]$
\item $[4, 9]$ 和 $[1, 10]$
\item $[2, 5]$ 和 $[5, 8]$
\end{enumerate}

\item \label{q:gtr-orders}（選擇）無向圖 $X$ 的頂點是 $1$ 到 $6$，邊是 $\{1, 2\}$、$\{1, 3\}$、
$\{2, 4\}$、$\{3, 4\}$、$\{3, 5\}$、$\{4, 6\}$、$\{5, 6\}$。 鄰接串列的順序可以任意安排，
從 $1$ 出發。 下列各個頂點順序，哪些可能是 BFS 拿出頂點的順序？ 哪些可能是 DFS 發現頂點的
順序？
\begin{enumerate}
\item $1, 3, 2, 5, 4, 6$
\item $1, 2, 3, 5, 4, 6$
\item $1, 2, 4, 3, 5, 6$
\item $1, 2, 4, 6, 3, 5$
\item $1, 3, 4, 2, 6, 5$
\end{enumerate}

\item \label{q:gtr-acyclic}（證明）對有向圖 $G$ 做一次 DFS，把所有後向邊刪掉。 (a) 證明剩下的圖沒有
環。 (b) 刪掉的邊數一定是「讓 $G$ 沒有環最少要刪幾條邊」嗎？ 證明或舉反例。

\item \label{q:gtr-tf}（證明或反例）下列敘述對的請證明，錯的請給反例。
\begin{enumerate}
\item 對有向圖做 DFS，交叉邊的兩端一定在不同的 DFS 樹裡。
\item 對無向圖做 DFS，若 $\{u, v\}$ 是後向邊，則 $u$、$v$ 在同一個環上。
\item 對有向圖做 DFS，若 $u$、$v$ 在同一個強連通元件裡，則它們在同一棵 DFS 樹裡。
\end{enumerate}

\item \label{q:gtr-count-paths}（設計）不帶權重的圖 $G$ 和起點 $s$，設計一個 $O(n + m)$ 的
演算法，求出從 $s$ 到每個頂點的\textbf{最短路徑有幾條}，並說明為什麼對。 用你的方法算出
$G_1$ 從 $1$ 出發的結果。

\item \label{q:gtr-on-path}（設計）有向圖 $G$ 用鄰接串列存，給兩個頂點 $s$、$t$，而且 $t$ 從 $s$
可到達。 設計一個 $O(n + m)$ 的演算法，找出所有「至少在一條從 $s$ 到 $t$ 的最短路徑上」的
頂點，並說明為什麼對。 用你的方法找出 $G_1$ 從 $1$ 到 $7$ 的答案（$G_1$ 是無向圖，把每條邊看
成兩個方向的有向邊）。

\item \label{q:gtr-wrong-dfs}（追蹤）第~\ref{sec:gtr-dfs-iter}~節提到一種錯誤的迴圈版 DFS：
推進堆疊時就標記，鄰居倒著推入。 在 $G_1$ 上從 $1$ 執行它，寫出拜訪順序，指出它從哪一
步開始和真正的 DFS 不同，原因是什麼。

\item \label{q:gtr-make-bip}（計算）$G_1$ 刪掉哪一條邊之後會變成二分圖？ 找出所有答案，並
各寫出一種分法 $L$、$R$。

\item \label{q:gtr-leaf}（證明）對連通的無向圖做 DFS，$u$ 是 DFS 樹裡的葉節點（沒有子節點），而
且不是根。 證明 $u$ 的每一個鄰居都是 $u$ 的祖先。
\end{enumerate}

\section{實作練習}
\label{sec:gtr-practice}

下面是 Aki-Judge 上和本章相關的題目。 題目裡的「鄰接表」就是本章的鄰接串列，「拜訪」就
是發現。 頂點編號從 $1$ 開始的題目，程式裡記得整段換成 $0..n-1$，或直接多開一格。

\begin{center}
\begin{tabular}{llc}
\toprule
小節 & 題目 & 難度 \\
\midrule
圖的表示 & \akiref{sorted-adjacency-list-graph}{維持排序的鄰接表圖} & 基本 \\
BFS & \akiref{fewest-edges-between-two-vertices}{兩點間最少經過的邊數} & 暖身 \\
 & \akiref{fewest-edges-from-a-source}{從起點出發的最少邊數} & 基本 \\
DFS 的迴圈版 & \akiref{stack-based-depth-first-visit-order}{用堆疊實作的深度優先走訪順序} & 基本 \\
連通元件 & \akiref{counting-connected-groups}{計算連通區塊數} & 暖身 \\
二分圖 & \akiref{two-coloring-a-graph}{圖能否二分著色} & 基本 \\
有向環 & \akiref{detecting-a-directed-cycle}{判斷有向圖是否有環} & 基本 \\
方格 & \akiref{reachable-cells-in-a-grid-maze}{迷宮格中可到達的格數} & 暖身 \\
 & \akiref{largest-connected-region-of-ones}{最大的 1 連通區域} & 基本 \\
 & \akiref{ink-blots-on-paper}{紙上的墨漬} & 進階 \\
多個起點 & \akiref{distance-to-the-nearest-source-on-a-grid}{格子上到最近起點的距離} & 進階 \\
延伸 & \akiref{cheapest-route-with-zero-and-one-costs}{只有 0 與 1 成本的最短路} & 進階 \\
\bottomrule
\end{tabular}
\end{center}

幾題的提示：

\begin{itemize}
\item 「維持排序的鄰接表圖」只交標頭檔和函式。 插入就是第~\ref{sec:gtr-c}~節的
\texttt{list\_insert}：無向邊兩端各插一次，重邊和自環都不加，邊數只在真的加進去時才加一。
第 $i$ 個鄰居從串列開頭往後走 $i$ 步。
\item 兩題 BFS 的頂點多達 $2 \times 10^5$，而且有長鏈測資。 BFS 本來就沒有遞迴，佇列用第~\ref{sec:gtr-c}~節
的長度 $n$ 陣列就夠；到不了的輸出 $-1$。
\item 「用堆疊實作的深度優先走訪順序」指定了第~\ref{sec:gtr-dfs-iter}~節的第二種寫法：鄰居
由大到小推入、拿出來時才檢查是否拜訪過。 同一個頂點可能被推入好幾次，堆疊的大小要以邊
數估計，最多 $2m + 1$ 個，不是 $n$。
\item 「計算連通區塊數」沒有邊的頂點自己算一個元件，外層迴圈每開始一次就加一。
\item 「圖能否二分著色」的圖可能不連通，每個元件都要檢查，就是虛擬碼~\ref{alg:gtr-bip}。
\item 「判斷有向圖是否有環」用虛擬碼~\ref{alg:gtr-cycle}；只分「走過、沒走過」兩種狀態
會把交叉邊誤判成環。
\item 方格的三題都不必建鄰接串列，用位移陣列算鄰居。 「紙上的墨漬」是八個方向；一塊墨漬
的「起點」就是依列優先順序掃描時，第一個碰到的格子，也就是外層迴圈從它開始走這一塊的那
一格。 紙最大 $200 \times 200$，遞迴深度最多四萬層，遞迴寫法還可以，用 BFS 更保險。
\item 「格子上到最近起點的距離」是第~\ref{sec:gtr-multi-source}~節的多起點 BFS；哪些格子要
輸出 $-1$ 照題目的說明。
\item 「只有 0 與 1 成本的最短路」超出本章一點：邊的成本不全相同，佇列要換成第~6~章的雙
端佇列。 成本 $0$ 的邊，鄰居放到前端；成本 $1$ 的，放到尾端；一個頂點的距離變小時才放
進去。 這樣拿出來的頂點距離仍然不減，道理和 BFS 的佇列性質一樣，第~25~章的 Dijkstra 演算
法是它的一般形式。
\end{itemize}

\section{讀完你應該會}
\label{sec:gtr-checklist}

\begin{itemize}
\item 用本章的符號描述一張圖（$n$、$m$、度數、路徑、環、連通元件、強連通、DAG、生成樹），
用握手定理和引理~\ref{lem:gtr-count} 做邊數、度數的計算。
\item 把一張圖寫成鄰接矩陣、鄰接串列、循序表示和鄰接多重串列，說出各操作在矩陣和串列上的
成本。
\item 照指定的串列順序手算 BFS 的每一輪佇列、$d$、$\pi$，畫出 BFS 樹，印出最短路徑；說明
為什麼 $d[v] = \delta(s, v)$。
\item 手算 DFS 的時間戳記和森林，用括號定理和白路徑定理判斷祖先、子孫關係，並把每條邊分成
樹邊、後向邊、前向邊、交叉邊。
\item 用 C 寫出鄰接串列、BFS（陣列佇列）、遞迴和迴圈兩種 DFS，知道遞迴深度的限制。
\item 用 BFS 或 DFS 求連通元件、判斷二分圖（找出奇環）、判斷有向圖有沒有環，並說出它們和
互斥集合各自適合的場合。
\end{itemize}

\textbf{延伸閱讀}（書末會整理完整的參考資料）

\begin{itemize}
\item Thomas H. Cormen 等，\textit{Introduction to Algorithms}，第 4 版，20.1--20.3 節：
圖的兩種表示法、BFS 的正確性證明、DFS 的括號定理、白路徑定理和邊的分類，證明寫得最完
整。 20.4 節的拓樸排序接本書第~22~章。
\item Steven S. Skiena，\textit{The Algorithm Design Manual}，第 3 版，第 7 章：從應用的角度
講 BFS 和 DFS，包括二分著色、找環、方格上的走訪，適合想看「什麼問題可以化成走訪」的讀者。
\item Antti Laaksonen，\textit{Competitive Programmer's Handbook}，第 11、12 章：圖的名詞、
存法和兩種走訪，篇幅很短，適合快速複習。
\item Ellis Horowitz、Sartaj Sahni、Susan Anderson-Freed，\textit{Fundamentals of Data
Structures in C}，第 6 章：鄰接多重串列、十字串列、循序表示等資料結構課常考的存法。
\item Robert E. Tarjan，``Depth-First Search and Linear Graph Algorithms''，\textit{SIAM
Journal on Computing} 1(2)，1972：用 DFS 在線性時間解決強連通元件和雙連通元件的經典論
文，第~23~章會再用到。
\end{itemize}

\addcontentsline{toc}{section}{觀念題解答}
\section*{觀念題解答}

\begin{enumerate}
\item (a)、(b)、(e)、(f) 正確。 (a)：每一對頂點最多一條邊。 (b)：握手定理，兩者都是 $m$。
(c)：錯，例如一個三角形加上一個孤立頂點，$n = 4$、$m = 3$，有環也不連通。 (d)：錯，列出
每個頂點的鄰居要掃一整列，是 $\Theta(n^2)$。 (e)：度數和是偶數 $2m$。 (f)：定理
\ref{thm:gtr-undirected}。

\item
\begin{enumerate}
\item 由握手定理 $3 \cdot 4 + 3(n - 3) = 2 \cdot 12$，得 $n = 7$。 這樣的圖確實存在，例如邊
$\{1,2\}$、$\{1,3\}$、$\{1,4\}$、$\{1,5\}$、$\{2,3\}$、$\{2,4\}$、$\{2,5\}$、$\{3,6\}$、
$\{3,7\}$、$\{4,6\}$、$\{5,7\}$、$\{6,7\}$，頂點 $1, 2, 3$ 的度數是 $4$，其他是 $3$。
\item 沒有。 度數和 $13$ 是奇數，違反握手定理。
\item 連通表示只有一個元件，由引理~\ref{lem:gtr-count}，$1 \ge 10 - m$，$m \ge 9$；一條路
徑剛好 $9$ 條邊。 不連通時，頂點可以分成沒有邊相連的兩群，$k$ 個和 $10 - k$ 個
（$1 \le k \le 9$），邊數最多 $\binom{k}{2} + \binom{10-k}{2}$。 這是 $k$ 的二次函數、開口
向上，最大值在兩端：$k = 1$ 或 $9$ 時是 $\binom{9}{2} = 36$。 例子是 $K_9$ 加一個孤立頂
點。
\end{enumerate}

\item (a) 鄰接串列和度數：

\begin{center}
\begin{tabular}{cccc}
\toprule
$u$ & $\mathit{Adj}[u]$ & $\deg^{+}(u)$ & $\deg^{-}(u)$ \\
\midrule
1 & $2, 4$ & 2 & 1 \\
2 & $3, 6$ & 2 & 1 \\
3 & $1, 5$ & 2 & 1 \\
4 & $5$ & 1 & 2 \\
5 & $6$ & 1 & 2 \\
6 & $4$ & 1 & 2 \\
\bottomrule
\end{tabular}
\end{center}

(b) $\mathit{start}[1..7] = \langle 1, 3, 5, 7, 8, 9, 10 \rangle$，$\mathit{nbr}[1..9] =
\langle 2, 4, 3, 6, 1, 5, 5, 6, 4 \rangle$。 (c) 反鄰接串列：$1\colon 3$；$2\colon 1$；
$3\colon 2$；$4\colon 1, 6$；$5\colon 3, 4$；$6\colon 2, 5$。 它就是矩陣每一行的 $1$ 所在的列。

\item 從 $F$ 出發：

\begin{center}
\small
\begin{tabular}{cl>{\raggedright\arraybackslash}p{0.25\linewidth}>{\raggedright\arraybackslash}p{0.27\linewidth}l}
\toprule
輪 & 拿出 $u$ & 依序看的鄰居 & 新發現：$d$，$\pi$ & 之後的佇列 \\
\midrule
0 & --- & --- & $F$：$d = 0$ & $F$ \\
1 & $F$ & $G$，$H$ & $G$：$1$，$F$；$H$：$1$，$F$ & $G\ H$ \\
2 & $G$ & $C$，$H$（灰） & $C$：$2$，$G$ & $H\ C$ \\
3 & $H$ & $F$（黑） & --- & $C$ \\
4 & $C$ & $A$ & $A$：$3$，$C$ & $A$ \\
5 & $A$ & $B$，$D$，$E$ & $B$：$4$，$A$；$D$：$4$，$A$；$E$：$4$，$A$ & $B\ D\ E$ \\
6 & $B$ & $C$（黑），$E$（灰） & --- & $D\ E$ \\
7 & $D$ & $C$（黑） & --- & $E$ \\
8 & $E$ & $D$（黑） & --- & （空） \\
\bottomrule
\end{tabular}
\end{center}

$d[A..H] = 3, 4, 2, 4, 4, 0, 1, 1$；$\pi[A] = C$，$\pi[B] = \pi[D] = \pi[E] = A$，
$\pi[C] = G$，$\pi[G] = \pi[H] = F$，$\pi[F] = \NIL$。 \proc{Print-Path}$(F, E)$ 印出
$F, G, C, A, E$，長度 $4$。 把範例~\ref{ex:gtr-graph} 輸入的第一張圖換成 $G_2$、
起點終點那一行改成 \texttt{F E}，可以驗證。

\item 事件依序是：發現 $1$、$2$、$3$，完成 $3$（$\mathit{Adj}[3] = \langle 1 \rangle$，$1$ 是
灰色），發現 $5$、$4$，完成 $4$（$3$ 是黑色），發現 $6$、$7$，完成 $7$、$6$、$5$、$2$、$1$。

\begin{center}
\begin{tabular}{c|ccccccc}
$v$ & 1 & 2 & 3 & 4 & 5 & 6 & 7 \\
\hline
$d[v]$ & 1 & 2 & 3 & 6 & 5 & 8 & 9 \\
$f[v]$ & 14 & 13 & 4 & 7 & 12 & 11 & 10 \\
$\pi[v]$ & \NIL & $1$ & $2$ & $5$ & $2$ & $5$ & $6$ \\
\end{tabular}
\end{center}

依被看到的順序分類：$(1, 2)$ 樹邊，$(2, 3)$ 樹邊，$(3, 1)$ 後向邊，$(2, 5)$ 樹邊，
$(5, 4)$ 樹邊，$(4, 3)$ 交叉邊（$3$ 是黑色，$d[4] = 6 > d[3] = 3$），$(5, 6)$ 樹邊，
$(6, 7)$ 樹邊，$(7, 5)$ 後向邊，$(7, 6)$ 後向邊，$(1, 4)$ 前向邊（$4$ 是黑色，$d[1] = 1 <
d[4] = 6$）。 共 $6$ 條樹邊、$3$ 條後向邊、$1$ 條前向邊、$1$ 條交叉邊。 有後向邊，所以
有環，例如 $\langle 1, 2, 3, 1 \rangle$、$\langle 5, 6, 7, 5 \rangle$、$\langle 6, 7, 6
\rangle$。 只有一棵樹，所以樹邊是 $n - 1 = 6$ 條。

\item (a) 可能，$v$ 的區間包在 $u$ 的裡面，$v$ 是 $u$ 的子孫。 (b) 不可能，兩個區間部分重
疊，違反括號定理。 (c) 可能，區間分開，兩者互不為祖先子孫；$u$ 在 $v$ 被發現前就完成了
（可能在同一棵樹的不同分支，也可能在不同的樹）。 (d) 可能，$u$ 是 $v$ 的子孫。 (e) 不可能，
每個時間只會用一次，$f[u]$ 和 $d[v]$ 不會都是 $5$。

\item BFS 可能的只有 (a)；DFS 可能的是 (c)、(e)。 (a) 是 BFS：取 $\mathit{Adj}[1] = \langle 3, 2
\rangle$、$\mathit{Adj}[3] = \langle 5, 4 \rangle$；不是 DFS，因為發現 $3$ 之後它還有白色的鄰居
$4$、$5$，下一個不能是 $2$。 (b)：BFS 拿出 $2$ 時放進 $4$，拿出 $3$ 時才放進 $5$，$4$ 一定
在 $5$ 前面；DFS 發現 $2$ 之後下一個必須是 $4$。 (c)：DFS 走 $1 \to 2 \to 4 \to 3 \to 5 \to 6$；
BFS 拿出 $1$ 之後接下來必須是 $2$、$3$。 (d)：DFS 發現 $6$ 時，$5$ 還是白色，下一個必須是
$5$。 (e)：DFS 走 $1 \to 3 \to 4 \to 2$，$2$ 沒有白色鄰居，退回 $4$，再 $4 \to 6 \to 5$；不是
BFS，理由同 (c)。

\item (a) 由第~\ref{sec:gtr-class}~節的表，樹邊、前向邊、交叉邊都有 $f[v] < f[u]$。 刪掉後向
邊之後，每條邊 $(u, v)$ 都有 $f[v] < f[u]$。 若還有環 $\langle v_0, v_1, \ldots, v_k = v_0
\rangle$，沿著環走 $f$ 嚴格遞減，得到 $f[v_0] > f[v_k] = f[v_0]$，矛盾。 (b) 不一定。 反例：
邊 $1 \to 2$、$2 \to 3$、$3 \to 1$、$3 \to 2$。 從 $1$ 做 DFS，$1 \to 2 \to 3$ 都是樹邊，$3 \to
1$ 和 $3 \to 2$ 都是後向邊，刪兩條；但只刪 $2 \to 3$ 一條，兩個環 $\langle 1, 2, 3, 1
\rangle$、$\langle 2, 3, 2 \rangle$ 就都斷了。 $G_2$ 剛好兩者相等：兩條後向邊各在一個
強連通元件裡，兩個元件各至少要刪一條。 （「最少刪幾條邊變成 DAG」在一般的圖上沒有已知
的多項式時間演算法，這類問題見第~28~章。）

\item
\begin{enumerate}
\item 錯。 $G_2$ 的 $(D, C)$ 是交叉邊，$D$、$C$ 都在 $A$ 那棵樹裡，只是在不同的分支。
\item 對。 後向邊 $\{u, v\}$ 中 $v$ 是 $u$ 的祖先（或反過來），而且 $v$ 不是 $u$ 的父節點
（那是樹邊的另一端）。 所以樹路徑 $v \leadsto u$ 至少有兩條邊，再加上 $\{u, v\}$，是一個長
度至少 $3$、頂點不重複的環。
\item 對。 設 $S$ 是那個強連通元件，$r$ 是 $S$ 裡最早被發現的頂點。 對 $S$ 裡任一個頂點
$v$，取一條從 $r$ 到 $v$ 的路徑；路徑上的每個頂點 $x$ 都有 $r \leadsto x$ 和 $x \leadsto v
\leadsto r$，所以 $x \in S$。 在時間 $d[r]$，$S$ 裡除了 $r$ 都是白色，這是一條白色路徑，
由白路徑定理 $v$ 是 $r$ 的子孫。 所以整個 $S$ 都在 $r$ 那棵樹裡。 第~23~章求強連通元件
就從這一點出發。
\end{enumerate}

\item 在 BFS 裡多記 $\mathit{cnt}[v]$，一開始 $\mathit{cnt}[s] = 1$、其他是 $0$。 拿出 $u$
看鄰居 $v$ 時：$v$ 是白色，照常發現，並設 $\mathit{cnt}[v] \gets \mathit{cnt}[u]$；$v$ 不
是白色但 $d[v] = d[u] + 1$，就 $\mathit{cnt}[v] \gets \mathit{cnt}[v] + \mathit{cnt}[u]$。
時間仍是 $O(n + m)$。

為什麼對：一條到 $v$ 的最短路徑，最後一條邊 $(u, v)$ 一定滿足 $\delta(s, u) = \delta(s, v)
- 1$，前面的部分是到 $u$ 的最短路徑；反過來每條到這種 $u$ 的最短路徑接上 $(u, v)$ 都是到
$v$ 的最短路徑。 所以 $v$ 的條數是所有這種 $u$ 的條數的和，上面的做法正好把每個這種 $u$
加一次（在 $u$ 出佇列時）。 還要確定 $u$ 出佇列時 $\mathit{cnt}[u]$ 已經算完：貢獻給
$\mathit{cnt}[u]$ 的都是距離 $d[u] - 1$ 的頂點，由佇列的性質，它們全部在 $u$ 之前出佇列。

$G_1$ 從 $1$ 出發：$\mathit{cnt}[1..7] = 1, 1, 1, 2, 1, 3, 1$，$8$、$9$、$10$ 是 $0$。 例如
到 $6$ 的三條是 $1, 2, 4, 6$、$1, 3, 4, 6$、$1, 3, 5, 6$；到 $7$ 只有 $1, 3, 5, 7$，因為
$6$ 和 $7$ 距離相同，邊 $\{6, 7\}$ 不在任何最短路徑上。 條數可能大到超過 \texttt{long long}
（例如一串菱形接起來，每多一個菱形條數加倍），實作時題目通常要求取餘數。

\item 從 $s$ 做一次 BFS 得到 $\delta(s, v)$；把每條邊反過來（用第~\ref{sec:gtr-repr}~節的反鄰
接串列，建它要 $O(n + m)$），從 $t$ 做一次 BFS，得到每個 $v$ 到 $t$ 的距離 $\delta(v, t)$。
答案是所有滿足 $\delta(s, v) + \delta(v, t) = \delta(s, t)$ 的 $v$。 理由：若 $v$ 在某條最短路
徑上，這條路徑在 $v$ 切成兩段，兩段各自至少是 $\delta(s, v)$、$\delta(v, t)$，加起來又等於
$\delta(s, t)$，所以兩段都剛好是最短距離，等式成立；反過來，等式成立時，把 $s$ 到 $v$ 的最短
路徑和 $v$ 到 $t$ 的最短路徑接起來，長度是 $\delta(s, t)$，就是一條經過 $v$ 的最短路徑（這條
走法不會重複頂點，因為重複的話剪掉一圈會比 $\delta(s, t)$ 還短）。 $G_1$ 從 $1$ 到 $7$：
$\delta(1, \cdot)$ 依序是 $0, 1, 1, 2, 2, 3, 3$，從 $7$ 做 BFS 得 $\delta(\cdot, 7)$ 依序是
$3, 3, 2, 2, 1, 1, 0$，$\delta(1, 7) = 3$，和是 $3$ 的頂點是 $1, 3, 5, 7$。 這也和第~\ref{q:gtr-count-paths}~題
「到 $7$ 的最短路徑只有一條」一致。

\item 拜訪順序是 $1, 2, 4, 6, 5, 7, 3$。 過程（堆疊左邊是底）：拿出 $1$，把 $3$、$2$ 標記並
推入，堆疊 $\langle 3, 2 \rangle$；拿出 $2$，推入 $4$，$\langle 3, 4 \rangle$；拿出 $4$，$3$
已經標記過，只推入 $6$，$\langle 3, 6 \rangle$；拿出 $6$，推入 $7$、$5$，
$\langle 3, 7, 5 \rangle$；拿出 $5$（鄰居都標記過）、$7$、$3$。 真正的 DFS 是 $1, 2, 4, 3,
5, 6, 7$，第四個頂點開始不同：在 $4$ 的時候，$3$ 還沒被發現，DFS 應該走過去；但錯誤的版本
在處理 $1$ 時就把 $3$ 標記了，於是跳過它，改走 $6$。 它記號的意思變成「已經排進待辦」，
而不是「已經發現」，這兩件事在 DFS 裡不一樣。

\item 刪 $\{5, 7\}$ 或刪 $\{6, 7\}$，只有這兩個答案。 $G_1$ 裡有奇環 $\langle 5, 6, 7, 5
\rangle$ 和 $\langle 3, 4, 6, 7, 5, 3 \rangle$（長度 $5$），刪掉的邊必須同時在這兩個環上，
否則留下的那個奇環讓圖仍不是二分圖；兩個環共同的邊只有 $\{6, 7\}$、$\{5, 7\}$。 再驗證這兩
個確實可以：刪 $\{6, 7\}$ 後，$L = \{1, 4, 5, 8, 10\}$、$R = \{2, 3, 6, 7, 9\}$；刪
$\{5, 7\}$ 後，$L = \{1, 4, 5, 7, 8, 10\}$、$R = \{2, 3, 6, 9\}$。 兩種分法都可以逐條邊檢
查，每條邊兩端都在不同組。 也可以把範例~\ref{ex:gtr-graph} 輸入的 $G_1$ 刪掉一條邊（$m$ 改成 $10$），看程式是否
印出 \texttt{bipartite: yes}。

\item 由定理~\ref{thm:gtr-undirected}，和 $u$ 關聯的邊 $\{u, w\}$ 不是樹邊就是後向邊。 樹邊
連接父子：$u$ 沒有子節點，所以 $w$ 只能是 $u$ 的父節點，是祖先。 後向邊連接一個頂點和它的祖
先：若 $u$ 是 $w$ 的祖先，$w$ 就是 $u$ 的子孫，但葉節點沒有子孫，所以只能是 $w$ 為 $u$ 的祖先。
第~23~章的割點判斷，靠的就是無向圖 DFS 的這個性質：不是樹邊的邊只會連到祖先。
\end{enumerate}

\chapter{拓樸排序、AOV 與 AOE 網路}
\label{ch:topological}

% 本章的圖共用的樣式（前綴 tps-）
\tikzset{
    tps-v/.style={circle, draw=black!70, minimum size=0.62cm, inner sep=0pt, font=\small, fill=white},
    tps-src/.style={tps-v, draw=coreStroke, thick, fill=coreFill},
    tps-e/.style={-{Stealth[length=2.2mm]}, draw=black!65, thick},
    tps-add/.style={tps-e, draw=examStroke, dashed},
    tps-ev/.style={circle, draw=black!70, minimum size=0.66cm, inner sep=0pt, font=\small, fill=white},
    tps-cev/.style={tps-ev, draw=examStroke, very thick},
    tps-a/.style={-{Stealth[length=2.2mm]}, draw=black!60, thick},
    tps-crit/.style={-{Stealth[length=2.6mm]}, draw=examStroke, line width=1.5pt},
    tps-dummy/.style={-{Stealth[length=2.2mm]}, draw=black!60, thick, dashed},
    tps-cdummy/.style={-{Stealth[length=2.6mm]}, draw=examStroke, line width=1.5pt, dashed},
    tps-al/.style={font=\scriptsize, inner sep=1pt, fill=white, text=black!85},
    tps-time/.style={font=\scriptsize, text=coreStroke, inner sep=1pt},
    tps-root/.style={rectangle, rounded corners=2pt, draw=black!65, font=\scriptsize, inner sep=2pt, fill=black!6},
    tps-tn/.style={circle, draw=black!65, minimum size=0.52cm, inner sep=0pt, font=\small, fill=white},
    tps-tleaf/.style={tps-tn, draw=memStroke, thick, fill=memFill},
    tps-te/.style={draw=black!55},
    tps-tord/.style={font=\scriptsize\ttfamily, text=black!80},
    tps-lvl/.style={font=\scriptsize, text=black!55, anchor=east},
    tps-gone/.style={text=black!35},
}
% 手算表裡已經取出的頂點
\newcommand{\tpsgone}{\textcolor{black!35}{--}}

家裡的老公寓要重新裝修。 工班給了一張工作清單：拆除舊裝潢、水電配管、浴室防水和貼
磚、天花板木作、油漆、訂做並安裝系統櫃、裝燈具、最後清潔。 有些工作要等別的做完才
能開始，油漆一定在木作之後，裝櫃子要等油漆乾；有些工作互不相干，可以同時做。 屋主
想知道三件事：這些工作可以排成什麼先後順序？ 人手足夠的話，最快幾天能完工？ 哪幾項
工作只要晚一天，整個工程就跟著晚一天？ 往下讀之前，先想想看要回答這三個問題，各需
要知道哪些資料。

第一個問題是\textbf{拓樸排序}（topological sort）：把有先後限制的工作排成一列，每一
項都排在它必須等待的工作後面。 第~\ref{sec:tps-kahn}~節的 Kahn 演算法一直挑「現在
沒有人擋住的」工作；第~\ref{sec:tps-dfs}~節用第~21~章的深度優先搜尋，依完成時間倒過
來排。 兩種方法都是線性時間，但排出來的順序常常不一樣，因為拓樸順序通常不只一種，考
試很喜歡問「有幾種」「全部列出來」，第~\ref{sec:tps-count}~節處理這件事。 後兩個問題
要加上每項工作的天數，這是第~\ref{sec:tps-aoe}~節 AOE 網路的\textbf{關鍵路徑}
（critical path）。 中間的第~\ref{sec:tps-aov}~節先講 AOV 網路，它是課本上描述「工作
之間的先後關係」的標準說法。

這一章假設你讀過第~21~章：有向圖、鄰接串列、深度優先搜尋的白灰黑三色，以及發現時間
$d[v]$ 和完成時間 $f[v]$（都從 $1$ 開始，每次發現或完成一個頂點，全域的時間加 $1$）。
C 程式用到《C 程式設計基礎》第~19~章的 \texttt{malloc}、\texttt{free} 和串列，佇列與
堆積見本書第~6~章和第~8~章。 AOE 網路的最早時間就是有向無環圖上的最長路徑，計算方式是
第~18~章的動態規劃，這裡會直接引用。

\section{拓樸排序與 Kahn 演算法}
\label{sec:tps-kahn}

\subsection{有向無環圖與拓樸順序}

先重述幾個名詞。 有向圖 $G$ 有 $n$ 個頂點、$m$ 條邊。 頂點 $v$ 的\textbf{入度}
（in-degree）$\deg^{-}(v)$ 是指向 $v$ 的邊數，\textbf{出度}（out-degree）$\deg^{+}(v)$
是從 $v$ 出去的邊數。 入度 $0$ 的頂點叫\textbf{源點}（source），出度 $0$ 的叫\textbf{匯點}（sink）。 沒有任何
有向環的有向圖叫做\textbf{有向無環圖}（directed acyclic graph，簡稱 DAG）。 邊
$(u, v)$ 寫成 $u \to v$，讀成「$u$ 要在 $v$ 之前」。

\begin{definition}[拓樸順序]
\label{def:tps-order}
有向圖 $G$ 的一個\textbf{拓樸順序}（topological order）是把全部 $n$ 個頂點排成一列
$v_1, v_2, \ldots, v_n$，使得每一條邊 $(u, v)$ 的 $u$ 都排在 $v$ 前面。 找出一個拓
樸順序叫做\textbf{拓樸排序}。
\end{definition}

把頂點依拓樸順序由左到右擺在一條線上，所有的邊都往右指。 這和第~14~章的排序不一樣：
那裡任兩個鍵都能比大小，這裡沒有邊（也沒有路徑）相連的兩個頂點，誰先誰後都可以。

圖~\ref{fig:tps-g1} 是本章前兩節共用的例子，$8$ 個頂點、$11$ 條邊。 頂點的編號是隨
便給的，不是拓樸順序，例如有 $2 \to 1$ 這條邊。

\begin{center}
\begin{tikzpicture}[x=1.15cm, y=0.9cm]
\node[tps-v] (2) at (0, 1) {2};
\node[tps-v] (3) at (0, -1) {3};
\node[tps-v] (5) at (2, 2) {5};
\node[tps-v] (1) at (2, 0) {1};
\node[tps-v] (6) at (2, -2) {6};
\node[tps-v] (4) at (4, 1) {4};
\node[tps-v] (8) at (4, -1) {8};
\node[tps-v] (7) at (6, 0) {7};
\foreach \u/\v in {2/1, 3/1, 2/5, 3/6, 1/4, 5/4, 6/4, 5/8, 6/8, 4/7, 8/7} {
    \draw[tps-e] (\u) -- (\v);
}
\node[font=\small, anchor=west, align=left] at (6.9, 0) {
$\mathit{Adj}[1]$: 4\\ $\mathit{Adj}[2]$: 1, 5\\ $\mathit{Adj}[3]$: 1, 6\\
$\mathit{Adj}[4]$: 7\\ $\mathit{Adj}[5]$: 4, 8\\ $\mathit{Adj}[6]$: 4, 8\\
$\mathit{Adj}[7]$: （空）\\ $\mathit{Adj}[8]$: 7};
\end{tikzpicture}
\diagcaption{本章的例子 $G_1$：$8$ 個頂點、$11$ 條邊的有向無環圖。 右邊是鄰接串列，每個頂點的鄰居依編號由小到大排。}
\label{fig:tps-g1}
\end{center}

什麼樣的圖才有拓樸順序？ 有環一定不行：環 $v_1 \to v_2 \to \cdots \to v_k \to v_1$
上的每條邊都要求前一個排在後一個前面，繞一圈回來變成 $v_1$ 排在 $v_1$ 前面。 反過來，
沒有環就一定排得出來，關鍵是下面這個引理。

\begin{lemma}
\label{lem:tps-source}
至少有一個頂點的有向無環圖，一定有入度 $0$ 的頂點，也一定有出度 $0$ 的頂點。
\end{lemma}

\begin{proof}
假設每個頂點的入度都至少是 $1$。 從任何一個頂點 $x_0$ 出發，它有一條進來的邊
$x_1 \to x_0$；$x_1$ 也有進來的邊 $x_2 \to x_1$……一直往回走 $n$ 步，得到
$x_0, x_1, \ldots, x_n$ 共 $n + 1$ 個頂點，一定有兩個相同，設 $x_i = x_j$
（$i < j$）。 那麼 $x_j \to x_{j-1} \to \cdots \to x_i$ 是一個有向環，矛盾。 出度的
情形一樣，改成順著出邊往前走。
\end{proof}

\begin{theorem}
\label{thm:tps-exists}
有向圖 $G$ 有拓樸順序，若且唯若 $G$ 是有向無環圖。
\end{theorem}

\begin{proof}
有環就沒有拓樸順序，上面已經說明。 反方向對 $n$ 做歸納：$n = 1$ 時顯然成立。 $n > 1$
時，由引理~\ref{lem:tps-source}，$G$ 有一個入度 $0$ 的頂點 $s$。 把 $s$ 和它的出邊拿
掉，剩下的圖仍然沒有環，由歸納假設有拓樸順序；把 $s$ 放在最前面，$s$ 的邊都往後指，
其他邊照舊，所以整列是 $G$ 的拓樸順序。
\end{proof}

這個證明本身就是一個演算法：一直拿掉入度 $0$ 的頂點。 下一小節把它寫成有效率的程序。

\subsection{Kahn 演算法}

直接照證明做，每一輪都要找一個入度 $0$ 的頂點，掃一遍是 $O(n)$，總共 $O(n^2)$。
Kahn 演算法的改進是記帳：
\begin{itemize}
\item 用陣列 $\mathit{in}[v]$ 記錄 $v$ 還有幾個前驅沒有排出去。
\item 把入度已經是 $0$、還沒排出去的頂點放在一個容器 $Q$ 裡。
\item 每次從 $Q$ 取出一個頂點 $u$ 排到答案的尾端，再把 $u$ 的每條出邊 $u \to v$
「拿掉」，也就是 $\mathit{in}[v]$ 減一。 減到 $0$ 的 $v$ 放進 $Q$。
\end{itemize}
入度要減到 $0$ 只會發生在減一的那一刻，所以不必重新掃描所有頂點。

$Q$ 裡同時有好幾個頂點時，取哪一個都對，但答案會不一樣。 本書的規則是：\textbf{$Q$
用佇列（先進先出）}；一開始依編號由小到大把入度 $0$ 的頂點放進去，處理 $u$ 的出邊時
照鄰接串列的順序（編號由小到大）。 想要「每次取編號最小的」，就把佇列換成最小優先佇
列，第~\ref{sec:tps-kahn-trace}~節會比較兩種結果。

\begin{pseudo}[label=alg:tps-kahn]{Kahn 演算法}
\Function{Kahn-Topo-Sort}{$G$}
    \For{$v \gets 1$ \To $n$}
        \State $\mathit{in}[v] \gets 0$
    \EndFor
    \For{$u \gets 1$ \To $n$}
        \For{$\mathit{Adj}[u]$ 裡的每一個 $v$}
            \State $\mathit{in}[v] \gets \mathit{in}[v] + 1$
        \EndFor
    \EndFor
    \State $Q \gets$ 空的佇列；$L \gets$ 空的串列
    \For{$v \gets 1$ \To $n$} \Comment{編號由小到大}
        \If{$\mathit{in}[v] = 0$}
            \State \Call{Enqueue}{$Q, v$}
        \EndIf
    \EndFor
    \While{$Q$ 不是空的}
        \State $u \gets$ \Call{Dequeue}{$Q$}
        \State 把 $u$ 接到 $L$ 的尾端
        \For{$\mathit{Adj}[u]$ 裡的每一個 $v$} \Comment{編號由小到大}
            \State $\mathit{in}[v] \gets \mathit{in}[v] - 1$
            \If{$\mathit{in}[v] = 0$}
                \State \Call{Enqueue}{$Q, v$}
            \EndIf
        \EndFor
    \EndWhile
    \If{$L$ 的頂點少於 $n$ 個}
        \State \Return 「有環，沒有拓樸順序」
    \EndIf
    \State \Return $L$
\EndFunction
\end{pseudo}

第 2--6 行算入度，第 7--10 行把一開始的源點放進佇列，第 11--17 行是主迴圈。 第 18
行的檢查不能省：有環的時候主迴圈也會正常結束，只是排出來的頂點不夠 $n$ 個。

\subsection{手算追蹤}
\label{sec:tps-kahn-trace}

用 $G_1$（圖~\ref{fig:tps-g1}）跑一次。 一開始的入度是 $\mathit{in}[1..8] = 2, 0, 0,
3, 1, 1, 2, 2$，入度 $0$ 的是 $2$ 和 $3$。 下表每一列是取出一個頂點、處理完它的出邊
之後的狀態，「--」表示這個頂點已經排出去了；佇列的前端寫在左邊。

\begin{center}
\small
\begin{tabular}{cc*{8}{c}ll}
\toprule
步驟 & 取出 & \multicolumn{8}{c}{$\mathit{in}[1..8]$} & 佇列 $Q$ & $L$ \\
\cmidrule(lr){3-10}
 & & 1 & 2 & 3 & 4 & 5 & 6 & 7 & 8 & & \\
\midrule
開始 & & 2 & 0 & 0 & 3 & 1 & 1 & 2 & 2 & 2 3 & \\
1 & 2 & 1 & \tpsgone & 0 & 3 & 0 & 1 & 2 & 2 & 3 5 & 2 \\
2 & 3 & 0 & \tpsgone & \tpsgone & 3 & 0 & 0 & 2 & 2 & 5 1 6 & 2 3 \\
3 & 5 & 0 & \tpsgone & \tpsgone & 2 & \tpsgone & 0 & 2 & 1 & 1 6 & 2 3 5 \\
4 & 1 & \tpsgone & \tpsgone & \tpsgone & 1 & \tpsgone & 0 & 2 & 1 & 6 & 2 3 5 1 \\
5 & 6 & \tpsgone & \tpsgone & \tpsgone & 0 & \tpsgone & \tpsgone & 2 & 0 & 4 8 & 2 3 5 1 6 \\
6 & 4 & \tpsgone & \tpsgone & \tpsgone & \tpsgone & \tpsgone & \tpsgone & 1 & 0 & 8 & 2 3 5 1 6 4 \\
7 & 8 & \tpsgone & \tpsgone & \tpsgone & \tpsgone & \tpsgone & \tpsgone & 0 & \tpsgone & 7 & 2 3 5 1 6 4 8 \\
8 & 7 & \tpsgone & \tpsgone & \tpsgone & \tpsgone & \tpsgone & \tpsgone & \tpsgone & \tpsgone & （空）& 2 3 5 1 6 4 8 7 \\
\bottomrule
\end{tabular}
\end{center}

逐步看幾個細節。 第 1 步取出 $2$，$\mathit{Adj}[2] = 1, 5$：$\mathit{in}[1]$ 從
$2$ 變 $1$，還不能進佇列；$\mathit{in}[5]$ 從 $1$ 變 $0$，放進佇列。 第 2 步取出 $3$
之後，$1$ 和 $6$ 依鄰接串列的順序先後進佇列，排在第 1 步就進來的 $5$ 後面。 第 5 步取
出 $6$，$4$ 的最後一個前驅排出去了，$4$ 和 $8$ 同時變成 $0$，依 $\mathit{Adj}[6] =
4, 8$ 的順序進佇列。 結果是
\[
L = 2, 3, 5, 1, 6, 4, 8, 7.
\]
逐條檢查 $11$ 條邊，起點都排在終點前面。

\textbf{改成每次取編號最小的。} 把佇列換成最小優先佇列（第~\ref{sec:hp-pq}~節），其他不變。
前兩步一樣取出 $2$、$3$；第 2 步之後 $Q$ 裡有 $1, 5, 6$，佇列會取最先進來的 $5$，最
小優先佇列取 $1$。 之後每一步的 $Q$ 是 $\{5, 6\}$、$\{6\}$、$\{4, 8\}$、$\{8\}$、
$\{7\}$，結果是
\[
L' = 2, 3, 1, 5, 6, 4, 8, 7.
\]
這個版本排出的是所有拓樸順序裡\textbf{字典序最小}的一個：每一步都在「現在可以排的頂點」
裡挑最小的，而任何拓樸順序在同一個位置都只能從這些頂點裡挑。 代價是每次進出優先佇列
花 $O(\lg n)$，總時間變成 $O(m + n \lg n)$。

\begin{examnote}
Kahn 演算法的追蹤題，答案取決於三個約定：(1) 入度 $0$ 的頂點用佇列、堆疊還是「每次挑
最小的」來存（有的課本用堆疊，排出來又不一樣）；(2) 一開始依什麼順序放進去；(3) 處理
出邊時照什麼順序。 本書用佇列、編號由小到大，題目有規定就照題目；題目只給圖、沒有說
時，在答案開頭寫一句你用的約定。 DFS 也一樣，外層的起點順序和鄰接串列的順序不同，
$d$、$f$ 和排出來的順序都會變，照題目給的鄰接串列走。
\end{examnote}

\subsection{正確性與執行時間}

\begin{theorem}
\label{thm:tps-kahn}
若 $G$ 是有向無環圖，\proc{Kahn-Topo-Sort} 回傳 $G$ 的一個拓樸順序；若 $G$ 有環，它
回傳「有環」。
\end{theorem}

\begin{proof}
主迴圈有一個不變量：每次檢查第 11 行的條件時，對每個還沒排出去的頂點 $v$，
$\mathit{in}[v]$ 等於「還沒排出去的前驅」的個數；$Q$ 裡恰好是那些 $\mathit{in}[v] = 0$
而且還沒排出去的頂點。 一開始沒有頂點排出去，$\mathit{in}$ 就是入度，第 7--10 行把入
度 $0$ 的都放進 $Q$。 取出 $u$ 時，$u$ 從「還沒排出去」變成「排出去了」，第 14--17 行
正好把 $u$ 的每個後繼減一，減到 $0$ 的放進 $Q$，不變量維持。

\textbf{排出來的順序合乎要求。} 頂點 $v$ 進 $Q$ 時 $\mathit{in}[v] = 0$，表示它的前驅
全都已經排出去了，所以每條邊 $u \to v$ 的 $u$ 都排在 $v$ 前面。

\textbf{沒有環時會排完。} 假設主迴圈結束時還有頂點沒排出去。 這些頂點連同它們之間的邊
形成 $G$ 的一個子圖，沒有環，由引理~\ref{lem:tps-source} 其中有一個頂點 $v$ 在子圖裡
入度 $0$，也就是它的前驅都已經排出去了，$\mathit{in}[v] = 0$。 由不變量，$v$ 在 $Q$
裡，$Q$ 不是空的，和迴圈已經結束矛盾。

\textbf{有環時排不完。} 設 $C$ 是一個環。 如果 $C$ 上有頂點被排出去，看其中最早排出
去的一個 $v$，它在 $C$ 上的前一個頂點 $u$ 是 $v$ 的前驅，必須更早排出去，矛盾。 所以
$C$ 上的頂點一個都排不出去，$L$ 不到 $n$ 個。
\end{proof}

注意「有環時排不完」只說環上的頂點排不出去。 卡住時剩下的頂點不一定都在環上：從環走
得到的頂點也會被擋住，第~\ref{sec:tps-aov}~節有一個例子。

\textbf{執行時間。} 用鄰接串列時，第 2--6 行是 $O(n + m)$；每個頂點進出 $Q$ 各一次，
每條邊在第 15 行被處理一次，主迴圈也是 $O(n + m)$。 總共 $\Theta(n + m)$。 用鄰接矩陣
時，算入度和找每個頂點的出邊都要掃一整列，總共 $\Theta(n^2)$。

\section{用 DFS 的完成時間排序}
\label{sec:tps-dfs}

第二種方法不數入度，直接對整張圖做一次深度優先搜尋，記下每個頂點的完成時間 $f[v]$，
再\textbf{依 $f$ 由大到小}排列。 實作時不必真的排序：每完成一個頂點就把它插到串列的最
前面，最後完成的自然在最前面。

\subsection{演算法}

\begin{pseudo}[label=alg:tps-dfs]{用 DFS 做拓樸排序}
\Function{DFS-Topo-Sort}{$G$}
    \For{$u \gets 1$ \To $n$}
        \State $\mathit{color}[u] \gets$ 白
    \EndFor
    \State $\mathit{time} \gets 0$；$L \gets$ 空的串列
    \For{$u \gets 1$ \To $n$}
        \If{$\mathit{color}[u] =$ 白}
            \State \Call{Topo-Visit}{$G, u$}
        \EndIf
    \EndFor
    \State \Return $L$
\EndFunction
\Procedure{Topo-Visit}{$G, u$}
    \State $\mathit{time} \gets \mathit{time} + 1$；$d[u] \gets \mathit{time}$
    \State $\mathit{color}[u] \gets$ 灰
    \For{$\mathit{Adj}[u]$ 裡的每一個 $v$} \Comment{編號由小到大}
        \If{$\mathit{color}[v] =$ 白}
            \State \Call{Topo-Visit}{$G, v$}
        \ElsIf{$\mathit{color}[v] =$ 灰}
            \State 回報「找到後向邊 $(u, v)$，有環」並停止
        \EndIf
    \EndFor
    \State $\mathit{color}[u] \gets$ 黑
    \State $\mathit{time} \gets \mathit{time} + 1$；$f[u] \gets \mathit{time}$
    \State 把 $u$ 插到 $L$ 的最前面
\EndProcedure
\end{pseudo}

$\mathit{time}$ 和 $L$ 是兩個程序共用的變數。 這就是第~21~章的 DFS 多了兩件事：第 16
行遇到灰色的頂點就停下來報告有環，第 19 行把剛完成的頂點插到 $L$ 最前面。 DFS 是
$\Theta(n + m)$，每次插到串列最前面是 $O(1)$，所以整個程序也是 $\Theta(n + m)$。

\subsection{為什麼一定對}

要證明的是：在有向無環圖上，每一條邊 $(u, v)$ 都滿足 $f[v] < f[u]$，這樣 $u$ 比 $v$ 晚
完成、在 $L$ 裡排在 $v$ 前面。 證明只用到第~21~章的\textbf{括號定理}（定理~\ref{thm:gtr-paren}）：任兩個頂點的時
間區間 $[d[u], f[u]]$ 和 $[d[v], f[v]]$ 不是互不相交，就是一個完全包含另一個；而
$v$ 的區間包含在 $u$ 的區間裡，正好表示 $v$ 是 $u$ 在 DFS 森林裡的子孫。

\begin{theorem}
\label{thm:tps-dfs}
對有向無環圖 $G$ 做 DFS，每一條邊 $(u, v)$ 都有 $f[v] < f[u]$。 因此依 $f$ 由大到小
排列（也就是 \proc{DFS-Topo-Sort} 回傳的 $L$）是一個拓樸順序。
\end{theorem}

\begin{proof}
看 $u$、$v$ 誰先被發現。

\textbf{情形一：$d[u] < d[v]$。} 在時間 $d[u]$，$v$ 還是白的。 $u$ 完成之前一定會檢查
到邊 $(u, v)$，那時 $v$ 不是還白著（於是現在被發現），就是在 $d[u]$ 之後、檢查這條邊之
前已經被發現了。 兩種情況都有 $d[u] < d[v] < f[u]$。 兩個區間有交集，由括號定理一個包
含另一個；$d[u] < d[v]$ 表示是 $u$ 的區間包含 $v$ 的，所以 $f[v] < f[u]$。

\textbf{情形二：$d[v] < d[u]$。} 如果 $u$ 在 $v$ 還沒完成時被發現，也就是 $d[v] <
d[u] < f[v]$，由括號定理 $u$ 的區間在 $v$ 的區間裡，$u$ 是 $v$ 的子孫，DFS 樹上有一條
從 $v$ 到 $u$ 的路徑；再加上邊 $u \to v$ 就是一個環，和 $G$ 沒有環矛盾。 所以 $u$ 是在
$v$ 完成之後才被發現的：$f[v] < d[u] < f[u]$。
\end{proof}

情形二也說明了第 15--16 行為什麼能偵測環：檢查邊 $(u, v)$ 時 $v$ 是灰的，表示 $u$ 是
$v$ 的子孫，這條邊是第~21~章說的\textbf{後向邊}，圖上一定有環。 第~21~章的定理~\ref{thm:gtr-cycle} 也證
明了反方向：有環的圖做 DFS 一定會出現後向邊。 所以 \proc{DFS-Topo-Sort} 不會漏掉環。

只看情形一、二的結論，可以把邊依種類整理：樹邊和前向邊是情形一（$v$ 是 $u$ 的子孫）；
交叉邊是情形二（$v$ 早就完成了）；後向邊在有向無環圖裡不會出現。 每一種都有
$f[v] < f[u]$。

\subsection{手算追蹤}

在 $G_1$ 上照虛擬碼~\ref{alg:tps-dfs} 走一次。 外層依 $1, 2, \ldots, 8$ 找白色的頂點，
鄰接串列照圖~\ref{fig:tps-g1}。

\begin{enumerate}
\item 從 $1$ 開始：$d[1] = 1$。 $\mathit{Adj}[1] = 4$，發現 $4$（$d = 2$）；
$\mathit{Adj}[4] = 7$，發現 $7$（$d = 3$）。 $7$ 沒有出邊，完成 $f[7] = 4$；接著 $4$
完成（$5$）、$1$ 完成（$6$）。 $L = 1, 4, 7$。
\item 外層到 $2$：$d[2] = 7$。 $\mathit{Adj}[2] = 1, 5$；$1$ 已經是黑的（交叉邊），跳過；
發現 $5$（$d = 8$）。 $\mathit{Adj}[5] = 4, 8$；$4$ 是黑的，發現 $8$（$d = 9$）。
$\mathit{Adj}[8] = 7$，黑的，$8$ 完成（$10$），$5$ 完成（$11$），$2$ 完成（$12$）。
$L = 2, 5, 8, 1, 4, 7$。
\item 外層的 $3$ 還是白的：$d[3] = 13$。 $1$ 是黑的；發現 $6$（$d = 14$），$\mathit{Adj}[6]
= 4, 8$ 都是黑的，$6$ 完成（$15$），$3$ 完成（$16$）。 $L = 3, 6, 2, 5, 8, 1, 4, 7$。
\item 外層的 $4..8$ 都已經是黑的，結束。
\end{enumerate}

\begin{center}
\begin{tabular}{c*{8}{c}}
\toprule
頂點 $v$ & 1 & 2 & 3 & 4 & 5 & 6 & 7 & 8 \\
\midrule
$d[v]$ & 1 & 7 & 13 & 2 & 8 & 14 & 3 & 9 \\
$f[v]$ & 6 & 12 & 16 & 5 & 11 & 15 & 4 & 10 \\
\bottomrule
\end{tabular}
\end{center}

依 $f$ 由大到小：$f[3] = 16$、$f[6] = 15$、$f[2] = 12$、$f[5] = 11$、$f[8] = 10$、
$f[1] = 6$、$f[4] = 5$、$f[7] = 4$，得到
\[
L = 3, 6, 2, 5, 8, 1, 4, 7.
\]
和 Kahn 演算法的 $2, 3, 5, 1, 6, 4, 8, 7$ 不同，但一樣合乎要求。 圖~\ref{fig:tps-line}
把頂點照這個順序排在一條線上，$11$ 條邊全部往右指。 可以順便驗證定理：例如交叉邊
$2 \to 1$ 有 $f[1] = 6 < 12 = f[2]$，$1$ 在 $2$ 被發現之前就完成了。

\begin{center}
\begin{tikzpicture}[x=1.3cm]
\node[tps-v] (L3) at (0,0) {3};
\node[tps-v] (L6) at (1,0) {6};
\node[tps-v] (L2) at (2,0) {2};
\node[tps-v] (L5) at (3,0) {5};
\node[tps-v] (L8) at (4,0) {8};
\node[tps-v] (L1) at (5,0) {1};
\node[tps-v] (L4) at (6,0) {4};
\node[tps-v] (L7) at (7,0) {7};
\foreach \v/\t in {3/16, 6/15, 2/12, 5/11, 8/10, 1/6, 4/5, 7/4} {
    \node[tps-time] at ($(L\v)+(0,-0.48)$) {$f = \t$};
}
\draw[tps-e] (L3) -- (L6);
\draw[tps-e] (L2) -- (L5);
\draw[tps-e] (L5) -- (L8);
\draw[tps-e] (L1) -- (L4);
\draw[tps-e] (L4) -- (L7);
\draw[tps-e] (L3) to[bend left=40] (L1);
\draw[tps-e] (L6) to[bend left=34] (L4);
\draw[tps-e] (L5) to[bend left=38] (L4);
\draw[tps-e] (L6) to[bend right=40] (L8);
\draw[tps-e] (L2) to[bend right=40] (L1);
\draw[tps-e] (L8) to[bend right=40] (L7);
\end{tikzpicture}
\diagcaption{$G_1$ 的頂點依 DFS 完成時間由大到小排成一列，每條邊都往右指。}
\label{fig:tps-line}
\end{center}

\subsection{C 程式}

下面的程式把兩種方法寫在一起：讀進好幾張圖，每張圖先跑 Kahn 演算法（佇列），再跑 DFS。
輸入的頂點是 $1..n$，程式裡一律減一變成 $0..n-1$，印出來時再加一。 鄰接串列插入新邊時找
到正確的位置，讓每個頂點的鄰居保持由小到大，這樣不管邊的輸入順序如何，結果都和手算一致。

\begin{cbookrecap}[第~14~章 14.5 節、第~19~章 19.3、19.10 節]
\texttt{Node **p = \&g->head[u];} 讓 \texttt{p} 指向「要改的那個指標」：一開始是串列
頭，往後走時變成前一個節點的 \texttt{next} 欄位。 找到位置之後，\texttt{x->next = *p;
*p = x;} 兩行就把新節點接進去，插在最前面和插在中間不必分開寫。 \texttt{calloc} 配
置的指標陣列每一格都是 \texttt{NULL}，代表空串列。
\end{cbookrecap}

Kahn 演算法的佇列直接用答案陣列 \texttt{q}：每個頂點最多進佇列一次，排出去的順序就是進佇
列的順序，所以 \texttt{q[0..head-1]} 是已經取出的，\texttt{q[head..cnt-1]} 是還在佇列裡
的。 \texttt{kahn} 把入度複製到 \texttt{deg} 再減，不改圖本身。 DFS 把「插到串列最前面」
換成陣列：\texttt{order} 有 $n$ 格，\texttt{pos} 從 $n$ 開始，每完成一個頂點就放進
\texttt{order[--pos]}。 找到後向邊時記下來，讓所有呼叫盡快結束，只印出那條邊。

\begin{example}[label=ex:tps-topo]{c}{(topo.c) Kahn 演算法與 DFS 完成時間法}
#include <stdio.h>
#include <stdlib.h>

enum { WHITE, GRAY, BLACK };

typedef struct node {           // 鄰接串列的一個節點：一條邊 u -> to
    int to;
    struct node *next;
} Node;

typedef struct {
    int n;
    Node **head;                // head[u]：u 的出邊，終點由小到大
    int *indeg;                 // 入度
    int *color, *d, *f, time;   // DFS 用
    int *order, pos;            // DFS 完成的頂點從 order 的最後一格往前放
    int back_u, back_v;         // 找到的後向邊；沒有時 back_u = -1
} Graph;

int graph_init(Graph *g, int n)
{
    g->n = n;
    g->head = calloc(n, sizeof *g->head);   // 每一格都是 NULL：空串列
    g->indeg = calloc(n, sizeof(int));
    g->color = calloc(n, sizeof(int));      // 0 就是 WHITE
    g->d = calloc(n, sizeof(int));
    g->f = calloc(n, sizeof(int));
    g->order = calloc(n, sizeof(int));
    g->time = 0;
    g->pos = n;
    g->back_u = g->back_v = -1;
    return g->head && g->indeg && g->color && g->d && g->f && g->order;
}

void graph_free(Graph *g)
{
    for (int u = 0; u < g->n && g->head != NULL; u++) {
        Node *x = g->head[u];
        while (x != NULL) {
            Node *nx = x->next;
            free(x);
            x = nx;
        }
    }
    free(g->head);
    free(g->indeg);
    free(g->color);
    free(g->d);
    free(g->f);
    free(g->order);
}

// 加入邊 u -> v，插在串列裡讓終點保持由小到大
int add_edge(Graph *g, int u, int v)
{
    Node *x = malloc(sizeof *x);
    if (x == NULL) {
        return 0;
    }
    Node **p = &g->head[u];
    while (*p != NULL && (*p)->to < v) {
        p = &(*p)->next;
    }
    x->to = v;
    x->next = *p;
    *p = x;
    g->indeg[v]++;
    return 1;
}

// Kahn 演算法（佇列）。拓樸順序寫進 q，回傳排出幾個頂點。
// q 本身就是佇列：q[head..cnt-1] 是還在佇列裡的頂點。
int kahn(const Graph *g, int *q)
{
    int *deg = malloc(g->n * sizeof *deg), cnt = 0;
    if (deg == NULL) {
        return -1;
    }
    for (int v = 0; v < g->n; v++) {
        deg[v] = g->indeg[v];
        if (deg[v] == 0) {
            q[cnt++] = v;
        }
    }
    for (int head = 0; head < cnt; head++) {
        for (Node *x = g->head[q[head]]; x != NULL; x = x->next) {
            if (--deg[x->to] == 0) {
                q[cnt++] = x->to;
            }
        }
    }
    free(deg);
    return cnt;
}

void visit(Graph *g, int u)
{
    g->d[u] = ++g->time;
    g->color[u] = GRAY;
    for (Node *x = g->head[u]; x != NULL && g->back_u < 0; x = x->next) {
        if (g->color[x->to] == WHITE) {
            visit(g, x->to);
        } else if (g->color[x->to] == GRAY) {   // 後向邊：有環
            g->back_u = u;
            g->back_v = x->to;
        }
    }
    g->color[u] = BLACK;
    g->f[u] = ++g->time;
    g->order[--g->pos] = u;                     // 插到串列最前面
}

void print_list(const char *name, const int *a, int k)
{
    printf("%s:", name);
    for (int i = 0; i < k; i++) {
        printf(" %d", a[i] + 1);                // 印回 1 起算的編號
    }
}

int main(void)
{
    int n, m;
    while (scanf("%d %d", &n, &m) == 2 && n > 0) {
        Graph g;
        int ok = graph_init(&g, n);
        for (int i = 0; i < m && ok; i++) {
            int u, v;
            ok = scanf("%d %d", &u, &v) == 2 && u >= 1 && u <= n && v >= 1
                 && v <= n && add_edge(&g, u - 1, v - 1);
        }
        int *q = ok ? malloc(n * sizeof *q) : NULL;
        int cnt = q != NULL ? kahn(&g, q) : -1;
        if (cnt < 0) {
            free(q);
            graph_free(&g);
            return 1;
        }
        print_list("kahn", q, cnt);
        if (cnt < n) {
            printf(" (stuck, %d left: cycle)", n - cnt);
        }
        printf("\n");
        for (int u = 0; u < n && g.back_u < 0; u++) {
            if (g.color[u] == WHITE) {
                visit(&g, u);
            }
        }
        if (g.back_u >= 0) {
            printf("dfs: back edge %d -> %d (cycle)\n", g.back_u + 1, g.back_v + 1);
        } else {
            print_list("dfs", g.order, n);
            printf("\nd f:");
            for (int v = 0; v < n; v++) {
                printf(" %d/%d", g.d[v], g.f[v]);
            }
            printf("\n");
        }
        free(q);
        graph_free(&g);
    }
    return 0;
}
\end{example}
\iosample{8 11 \\ 2 1 3 1 1 4 2 5 3 6 5 4 \\ 6 4 5 8 6 8 4 7 8 7 \\ 5 6 \\ 1 2 2 3 3 4 4 2 1 5 5 3}{kahn: 2 3 5 1 6 4 8 7 \\ dfs: 3 6 2 5 8 1 4 7 \\ d f: 1/6 7/12 13/16 2/5 8/11 14/15 3/4 9/10 \\ kahn: 1 5 (stuck, 3 left: cycle) \\ dfs: back edge 4 -> 2 (cycle)}

第一張圖是 $G_1$，兩種方法的順序和 $d$、$f$（\texttt{d f:} 那一行依頂點 $1..8$ 列出
$d/f$）都和手算相同。 第二張圖有環 $2 \to 3 \to 4 \to 2$：Kahn 演算法取出 $1$、$5$ 之
後，$2$、$3$、$4$ 的入度都還是 $1$，佇列空了；DFS 從 $1$ 出發依序發現 $2$、$3$、$4$，在
$4$ 檢查邊 $4 \to 2$ 時 $2$ 還是灰的。 想要字典序最小的版本，把佇列換成第~8~章的最小堆
積就好，其他不變。 遞迴深度最多是 $n$；$n$ 很大（例如 $10^6$ 個頂點排成一條鏈）時，C 的
呼叫堆疊可能不夠，要改用自己的堆疊模擬遞迴（第~\ref{sec:gtr-dfs-iter}~節），或乾脆用 Kahn 演算法。

\subsection{兩種方法的比較}

\begin{center}
\begin{tabular}{p{0.2\linewidth}p{0.34\linewidth}p{0.34\linewidth}}
\toprule
 & Kahn 演算法 & DFS 完成時間法 \\
\midrule
想法 & 一直取出入度 $0$ 的頂點 & 依完成時間由大到小 \\
產生順序 & 由前往後，第一個先確定 & 由後往前，最後一個先確定 \\
需要的資料 & 入度陣列、一個佇列 & 顏色、時間戳記、遞迴（或堆疊） \\
偵測環 & 排出的頂點少於 $n$ 個 & 遇到灰色的頂點（後向邊） \\
指定平手規則 & 換容器就好，例如最小優先佇列得到字典序最小 & 不容易控制 \\
時間 & $\Theta(n + m)$ & $\Theta(n + m)$ \\
\bottomrule
\end{tabular}
\end{center}

「指定平手規則」那一列的意思是：想要字典序最小的拓樸順序，Kahn 演算法只要換容器；DFS
的順序由搜尋路徑決定，沒有簡單的改法。 反過來，DFS 不必先算入度，而且同一次搜尋還能
拿來做別的事，例如第~23~章的強連通元件。

\section{拓樸順序不唯一}
\label{sec:tps-count}

$G_1$ 已經出現了三個不同的拓樸順序。 它其實有 $38$ 個（下面的程式可以算出來）。 這一
節回答三個考試常問的問題：什麼時候只有一個？ 怎麼不重複、不遺漏地全部列出來？ 有沒有
快一點的數法？

\subsection{什麼時候唯一}

\begin{theorem}
\label{thm:tps-unique}
設 $v_1, v_2, \ldots, v_n$ 是有向無環圖 $G$ 的一個拓樸順序。 $G$ 只有這一個拓樸順序，
若且唯若每一對相鄰的 $v_i, v_{i+1}$ 之間都有邊 $v_i \to v_{i+1}$；也就是 $G$ 有一條經
過所有頂點的路徑。
\end{theorem}

\begin{proof}
若每對相鄰的頂點都有邊，路徑 $v_1 \to v_2 \to \cdots \to v_n$ 規定了每一對相鄰頂點的
先後，順序只能是這一個。 反過來，若某一對 $v_i, v_{i+1}$ 之間沒有邊，把這兩個對調。 被
影響的只有一端是 $v_i$ 或 $v_{i+1}$ 的邊；這兩個頂點之間沒有邊，而它們和其他頂點的前
後關係都沒變，所以每條邊的起點仍然在終點前面，得到第二個拓樸順序。
\end{proof}

用 Kahn 演算法來看，拓樸順序唯一，等於每一步佇列裡都恰好只有一個頂點；佇列裡出現兩個
以上，取哪一個都會得到不同的順序。 所以判斷「拓樸順序是否唯一」只要跑一次 Kahn 演算
法，$\Theta(n + m)$。

\subsection{全部列出來}

列舉所有拓樸順序，就是第~20~章的回溯：一個位置一個位置決定，每個位置可以放的是「還
沒放、而且前驅都放好了」的頂點，也就是目前入度 $0$ 的頂點。 試的時候依字母（或編號）
由小到大，列出來的順序自然是字典序，不會重複也不會漏掉。

用一個小例子 $H$：$5$ 個頂點 A--E，邊是 $A \to C$、$B \to C$、$B \to D$、$C \to E$。
圖~\ref{fig:tps-enum} 是它的列舉樹：第 $k$ 層的節點是第 $k$ 個位置放的頂點，同一個節點
的子節點就是那時候入度 $0$ 的頂點。 例如放了 A 之後，C 還在等 B、D 也在等 B，所以只有
B 一個子節點；放了 B、A 之後，C 和 D 都可以放。

\begin{center}
\begin{tikzpicture}[x=1cm, y=0.95cm]
\begin{scope}[shift={(-3.6,-1.6)}, x=0.95cm, y=0.95cm]
\node[tps-v] (hA) at (0, 1) {A};
\node[tps-v] (hB) at (0, -1) {B};
\node[tps-v] (hC) at (1.4, 0.3) {C};
\node[tps-v] (hD) at (1.4, -1.4) {D};
\node[tps-v] (hE) at (2.6, 0.3) {E};
\draw[tps-e] (hA) -- (hC);
\draw[tps-e] (hB) -- (hC);
\draw[tps-e] (hB) -- (hD);
\draw[tps-e] (hC) -- (hE);
\node[font=\small] at (1.3, -2.3) {圖 $H$};
\end{scope}
\node[tps-root] (t0) at (4.94,0.00) {開始};
\node[tps-tn] (t1) at (1.94,-1.00) {A};
\draw[tps-te] (t0) -- (t1);
\node[tps-tn] (t2) at (1.94,-2.00) {B};
\draw[tps-te] (t1) -- (t2);
\node[tps-tn] (t3) at (0.78,-3.00) {C};
\draw[tps-te] (t2) -- (t3);
\node[tps-tn] (t4) at (0.00,-4.00) {D};
\draw[tps-te] (t3) -- (t4);
\node[tps-tleaf] (t5) at (0.00,-5.00) {E};
\draw[tps-te] (t4) -- (t5);
\node[tps-tord] at (0.00,-5.62) {ABCDE};
\node[tps-tn] (t6) at (1.55,-4.00) {E};
\draw[tps-te] (t3) -- (t6);
\node[tps-tleaf] (t7) at (1.55,-5.00) {D};
\draw[tps-te] (t6) -- (t7);
\node[tps-tord] at (1.55,-5.62) {ABCED};
\node[tps-tn] (t8) at (3.10,-3.00) {D};
\draw[tps-te] (t2) -- (t8);
\node[tps-tn] (t9) at (3.10,-4.00) {C};
\draw[tps-te] (t8) -- (t9);
\node[tps-tleaf] (t10) at (3.10,-5.00) {E};
\draw[tps-te] (t9) -- (t10);
\node[tps-tord] at (3.10,-5.62) {ABDCE};
\node[tps-tn] (t11) at (7.94,-1.00) {B};
\draw[tps-te] (t0) -- (t11);
\node[tps-tn] (t12) at (6.59,-2.00) {A};
\draw[tps-te] (t11) -- (t12);
\node[tps-tn] (t13) at (5.42,-3.00) {C};
\draw[tps-te] (t12) -- (t13);
\node[tps-tn] (t14) at (4.65,-4.00) {D};
\draw[tps-te] (t13) -- (t14);
\node[tps-tleaf] (t15) at (4.65,-5.00) {E};
\draw[tps-te] (t14) -- (t15);
\node[tps-tord] at (4.65,-5.62) {BACDE};
\node[tps-tn] (t16) at (6.20,-4.00) {E};
\draw[tps-te] (t13) -- (t16);
\node[tps-tleaf] (t17) at (6.20,-5.00) {D};
\draw[tps-te] (t16) -- (t17);
\node[tps-tord] at (6.20,-5.62) {BACED};
\node[tps-tn] (t18) at (7.75,-3.00) {D};
\draw[tps-te] (t12) -- (t18);
\node[tps-tn] (t19) at (7.75,-4.00) {C};
\draw[tps-te] (t18) -- (t19);
\node[tps-tleaf] (t20) at (7.75,-5.00) {E};
\draw[tps-te] (t19) -- (t20);
\node[tps-tord] at (7.75,-5.62) {BADCE};
\node[tps-tn] (t21) at (9.30,-2.00) {D};
\draw[tps-te] (t11) -- (t21);
\node[tps-tn] (t22) at (9.30,-3.00) {A};
\draw[tps-te] (t21) -- (t22);
\node[tps-tn] (t23) at (9.30,-4.00) {C};
\draw[tps-te] (t22) -- (t23);
\node[tps-tleaf] (t24) at (9.30,-5.00) {E};
\draw[tps-te] (t23) -- (t24);
\node[tps-tord] at (9.30,-5.62) {BDACE};
\end{tikzpicture}
\diagcaption{左：圖 $H$。 右：列舉 $H$ 所有拓樸順序的樹，每個節點的子節點是那時入度 $0$ 的頂點，依字母排。 $7$ 個葉節點由左到右就是字典序的 $7$ 個拓樸順序。}
\label{fig:tps-enum}
\end{center}

所以 $H$ 共有 $7$ 個拓樸順序：
\[
\texttt{ABCDE},\ \texttt{ABCED},\ \texttt{ABDCE},\ \texttt{BACDE},\ \texttt{BACED},\
\texttt{BADCE},\ \texttt{BDACE}.
\]

\begin{examnote}
「列出所有拓樸順序」的題目，照列舉樹的方式寫：第一個位置先試最小的源點，把它底下的全
部列完，再換下一個。 答案依字典序排好，最後寫出總數；閱卷的人一眼就能看出有沒有漏。
寫完可以抽一兩個檢查每條邊的起點都在終點前面。 題目只問「有幾種」時，能用下一小節的
公式就用公式，再用列舉樹驗算小的部分。
\end{examnote}

下面的程式照列舉樹的方式回溯。 $n$ 很小，所以用鄰接矩陣：放下 $v$ 時從 \texttt{indeg} 減
掉 $v$ 那一列（拿掉它的出邊），回來時再加回去，這是第~20~章的「做、探索、撤
銷」。

\begin{example}[label=ex:tps-all]{c}{(all\_topo.c) 列出所有拓樸順序}
#include <stdio.h>
#include <stdlib.h>

#define MAXN 26                 // 頂點用大寫字母 A..Z

int n;
int adj[MAXN][MAXN];            // adj[u][v] = 1 表示有邊 u -> v
int indeg[MAXN];                // 還沒放進順序的前驅有幾個
int used[MAXN];
char seq[MAXN + 1];
long long total;

// 已經排好 k 個頂點，試第 k + 1 個位置放誰
void extend(int k)
{
    if (k == n) {
        total++;
        printf("%lld: %s\n", total, seq);
        return;
    }
    for (int v = 0; v < n; v++) {           // 字母由小到大試
        if (used[v] || indeg[v] != 0) {
            continue;
        }
        used[v] = 1;                        // 做：放 v，拿掉它的出邊
        seq[k] = (char)('A' + v);
        for (int w = 0; w < n; w++) {
            indeg[w] -= adj[v][w];
        }
        extend(k + 1);                      // 探索
        for (int w = 0; w < n; w++) {       // 撤銷
            indeg[w] += adj[v][w];
        }
        used[v] = 0;
    }
}

int main(void)
{
    int m;
    if (scanf("%d %d", &n, &m) != 2 || n < 1 || n > MAXN) {
        return 1;
    }
    for (int i = 0; i < m; i++) {
        char a, b;
        if (scanf(" %c %c", &a, &b) != 2 || a < 'A' || a >= 'A' + n
            || b < 'A' || b >= 'A' + n) {
            return 1;
        }
        int u = a - 'A', v = b - 'A';
        if (!adj[u][v]) {                   // 重複的邊只算一次
            adj[u][v] = 1;
            indeg[v]++;
        }
    }
    seq[n] = '\0';
    extend(0);
    printf("total = %lld\n", total);
    return 0;
}
\end{example}
\iosample{5 4 \\ A C \\ B C \\ B D \\ C E}{1: ABCDE \\ 2: ABCED \\ 3: ABDCE \\ 4: BACDE \\ 5: BACED \\ 6: BADCE \\ 7: BDACE \\ total = 7}

把 $G_1$ 的頂點 $1..8$ 換成 A--H 輸入，程式會列出 $38$ 行，第一行就是字典序最小的
\texttt{BCAEFDHG}，也就是 $2, 3, 1, 5, 6, 4, 8, 7$，和最小優先佇列版的 Kahn 演算法一致。
有環的圖會一個都列不出來，印出 \texttt{total = 0}。

\subsection{數拓樸順序}

列舉的時間和拓樸順序的個數成正比，而個數可以多到 $n!$（沒有邊的圖）。 一般的圖沒有
簡單的公式，但考試常見的幾種結構可以直接算。

\textbf{互不相連的幾塊。} 圖分成沒有邊相連的兩塊，各有 $n_1$、$n_2$ 個頂點、$c_1$、
$c_2$ 個拓樸順序。 整張圖的拓樸順序，等於先決定哪 $n_1$ 個位置給第一塊（其餘給第二
塊），再在這些位置裡各自照一個拓樸順序填。 所以共有
\[
\binom{n_1 + n_2}{n_1} \, c_1 \, c_2
\]
個。 特例：兩條鏈 $a_1 \to \cdots \to a_p$ 和 $b_1 \to \cdots \to b_q$，每條鏈自己只有
$1$ 種順序，合起來有 $\binom{p + q}{p}$ 種，就是「把兩疊牌洗在一起、各自保持順序」的
方法數。 推廣到 $k$ 塊，是多項式係數 $\frac{(n_1 + \cdots + n_k)!}{n_1! \cdots n_k!}
\, c_1 \cdots c_k$。

\textbf{必經的頂點。} 若頂點 $x$ 和其他每一個頂點之間都有路徑（不是 $x$ 走得到它，就
是它走得到 $x$），$x$ 在每個拓樸順序裡的位置都一樣，前面是所有走得到 $x$ 的頂點，後面
是 $x$ 走得到的頂點。 個數等於「前面那部分的個數」乘上「後面那部分的個數」。

用 $H$ 驗算：C 不是必經的（D 和 C 之間沒有路徑），沒辦法直接切。 換個角度：E 一定在 C
後面，而 C 在 A、B 後面。 先不管 D，A、B、C、E 的順序只有 $\texttt{ABCE}$、
$\texttt{BACE}$ 兩種；D 只要求在 B 後面，插進去的位置：$\texttt{ABCE}$ 裡 B 後面有
$3$ 個空位，$\texttt{BACE}$ 裡有 $4$ 個，合計 $3 + 4 = 7$，和列舉樹一樣。 一般的圖要
這樣分情況，頂點多了就容易漏，用列舉或子集合上的動態規劃（狀態是「已經放好的頂點集
合」，$2^n$ 個狀態）比較可靠。

\section{AOV 網路}
\label{sec:tps-aov}

資料結構的課本描述「工作之間的先後關係」時，用的是 AOV 網路。

\begin{definition}[AOV 網路]
\label{def:tps-aov}
\textbf{AOV 網路}（activity on vertex network）是一個有向圖：每個頂點代表一項\textbf{活
動}（activity），邊 $\langle i, j \rangle$ 表示活動 $i$ 必須在活動 $j$ 開始之前完成。
這時稱 $i$ 是 $j$ 的\textbf{直接前置活動}（immediate predecessor），$j$ 是 $i$ 的
\textbf{直接後繼活動}（immediate successor）。 若從 $i$ 有一條路徑到 $j$，稱 $i$ 是 $j$
的\textbf{前置活動}（predecessor）。
\end{definition}

\subsection{前置關係與偏序}

把「$i$ 是 $j$ 的前置活動」記成 $i \prec j$。 這個關係有兩個性質：
\begin{itemize}
\item \textbf{遞移}（transitive）：$i \prec j$ 且 $j \prec k$，則 $i \prec k$。 路徑接起
來還是路徑。
\item \textbf{非自反}（irreflexive）：沒有活動滿足 $i \prec i$，自己不能是自己的前置活
動。
\end{itemize}
同時滿足這兩個性質的關係叫做（嚴格的）\textbf{偏序}（partial order），兩者合起來還推得
出\textbf{反對稱}：$i \prec j$ 和 $j \prec i$ 不會同時成立，否則由遞移得到 $i \prec i$。
「偏」的意思是不必任兩個活動都分得出先後。 拓樸順序就是把偏序擴充成一個全部排好的順序
（線性順序），而且不違反原本的任何一對先後關係。

$i \prec i$ 成立，正好表示從 $i$ 有一條路徑回到 $i$，也就是有環。 所以\textbf{前置關係
是偏序，若且唯若 AOV 網路沒有環}；一個有環的 AOV 網路代表某些活動要等自己完成才能開
始，這個計畫根本做不完。 檢查一份計畫可不可行，就是做一次拓樸排序（定理~\ref{thm:tps-kahn}）。

畫 AOV 網路時，通常只畫直接前置的邊。 若已經有 $i \to j \to k$，再畫 $i \to k$ 並不會
多出任何限制（拓樸順序的集合不變），這種邊叫\textbf{遞移邊}（transitive edge），可以
省略。 題目給的表格如果包含遞移邊，畫圖時照畫也沒關係，拓樸排序的結果不受影響，只是
Kahn 演算法多處理幾條邊。

\subsection{例子：系學會迎新}

系學會準備迎新活動，工作和直接前置活動如下表。 圖~\ref{fig:tps-aov} 是它的 AOV 網路。

\begin{center}
\begin{tabular}{clc}
\toprule
活動 & 內容 & 直接前置活動 \\
\midrule
A & 確定日期與預算 & -- \\
B & 借教室 & A \\
C & 設計海報與報名表 & A \\
D & 開放報名 & B, C \\
E & 準備闖關道具 & B \\
F & 訂餐 & D \\
G & 工作人員彩排 & E \\
H & 迎新當天 & F, G \\
\bottomrule
\end{tabular}
\end{center}

\begin{center}
\begin{tikzpicture}[x=1.45cm, y=1cm]
\node[tps-v] (A) at (0, 0) {A};
\node[tps-v] (B) at (1.5, 1) {B};
\node[tps-v] (C) at (1.5, -1) {C};
\node[tps-v] (D) at (3, -1) {D};
\node[tps-v] (E) at (3, 1) {E};
\node[tps-v] (F) at (4.5, -1) {F};
\node[tps-v] (G) at (4.5, 1) {G};
\node[tps-v] (H) at (6, 0) {H};
\foreach \u/\v in {A/B, A/C, B/D, C/D, B/E, D/F, E/G, F/H, G/H} {
    \draw[tps-e] (\u) -- (\v);
}
\draw[tps-add] (D) to[bend right=35] node[right, font=\scriptsize, text=examStroke] {加上 $D \to B$ 就有環} (B);
\end{tikzpicture}
\diagcaption{迎新活動的 AOV 網路（實線）。 紅色虛線是後面討論的不合理要求。}
\label{fig:tps-aov}
\end{center}

用 Kahn 演算法（佇列，字母由小到大）排一次：A 先出來，B、C 進佇列；取出 B，E 變成入度 $0$，
D 還在等 C；取出 C，D 進佇列……得到
\[
\text{A, B, C, E, D, G, F, H}.
\]
如果只有一個人做事、一次做一項，這就是一個可行的工作順序。 這張圖共有 $16$ 個拓樸順
序，任何一個都可以。

假設有人提議「報名人數確定後才去借教室，才知道要借多大的」，也就是加上 $D \to B$。 可
是報名表上要寫活動地點，原本就有 $B \to D$。 這時 Kahn 演算法取出 A，C 進佇列（B 的入度
還有 $1$，來自 D）；取出 C 之後 D 的入度還有 $1$（來自 B），佇列空了，只排出 A、C 兩個。 剩
下的六項裡，只有 B、D 在環 $B \to D \to B$ 上；E、F、G、H 是被環擋住的下游活動。 這就是
前面說的「卡住時剩下的頂點不一定都在環上」。

\subsection{活動有工期的時候}

如果每項活動 $v$ 還有所需時間 $t(v)$，而且人手足夠、任何可以開始的活動都立刻開始，就
能問「每項活動最早什麼時候完成」。 把最早完成時間記成 $\mathit{EF}(v)$：
\[
\mathit{EF}(v) = t(v) + \max_{u \text{ 是 } v \text{ 的直接前置活動}} \mathit{EF}(u),
\]
沒有前置活動時 $\max$ 那一項當成 $0$。 依拓樸順序算，每個 $\mathit{EF}(u)$ 用到時都已
經算好了；也可以寫成記憶化的遞迴，從要問的活動往前置活動找。 整個計畫的最短工期是最大
的 $\mathit{EF}$，它等於一條路徑上所有頂點的工期總和的最大值，也就是「權重在頂點上」
的最長路徑。

這件事在 AOV 網路上也算得出來，但課本和考卷處理工期，習慣用下一節的 AOE 網路：活動改
放在邊上，頂點變成「某些活動都完成了」的時間點。 兩種網路的差別整理如下。

\begin{center}
\begin{tabular}{p{0.2\linewidth}p{0.34\linewidth}p{0.34\linewidth}}
\toprule
 & AOV 網路 & AOE 網路 \\
\midrule
頂點 & 活動 & 事件（時間點） \\
邊 & 先後關係 & 活動，權重是工期 \\
主要問題 & 可不可行、排出一個順序 & 最短工期、關鍵活動、寬裕時間 \\
用到的演算法 & 拓樸排序 & 拓樸排序 $+$ 最長路徑（往前、往回各一次） \\
\bottomrule
\end{tabular}
\end{center}

\section{AOE 網路與關鍵路徑}
\label{sec:tps-aoe}

\subsection{事件與活動}

\begin{definition}[AOE 網路]
\label{def:tps-aoe}
\textbf{AOE 網路}（activity on edge network）是一個邊上有非負權重的有向無環圖。 每條邊
代表一項活動，權重 $\mathit{dur}(a)$ 是活動 $a$ 需要的時間；每個頂點代表一個\textbf{事
件}（event），意思是「所有指向它的活動都已經完成」，從它出發的活動在這個時刻之後才能開
始。 入度 $0$ 的事件叫\textbf{開始事件}（源點），出度 $0$ 的叫\textbf{結束事件}（匯點）。
\end{definition}

本節假設 AOE 網路只有一個開始事件、一個結束事件，而且每個事件都在某條從開始到結束的路
徑上。 有好幾個開始事件時，加一個新的開始事件，用工期 $0$ 的邊連到原本的每一個；結束事
件也一樣處理。 也假設人手足夠：一項活動只要它的起點事件發生了就能開工，不必和別的活動
搶人。

有時候兩項活動之間只有先後關係，沒有實際的工作，就用一條工期 $0$ 的\textbf{虛活動}
（dummy activity）表示，圖上畫成虛線。 下面的例子會看到為什麼需要它。

\textbf{例子：裝修工程。} 章首的裝修工程整理成下面兩張表：第一張是事件的意思，第二張是活動。
\begin{center}
\small
\begin{tabular}{cl@{\hspace{2em}}cl}
1 & 開工 & 6 & 可以油漆（泥作、木作都完成） \\
2 & 拆除完成 & 7 & 油漆完成，系統櫃也運到了 \\
3 & 系統櫃設計定案 & 8 & 系統櫃、衛浴設備都裝好 \\
4 & 水電配管完成 & 9 & 完工 \\
5 & 窗框與浴室泥作完成 & & \\
\end{tabular}
\end{center}

\begin{center}
\begin{tabular}{clcc}
\toprule
活動 & 內容 & 邊 & 工期（天） \\
\midrule
$a_1$ & 拆除舊裝潢 & $\langle 1, 2 \rangle$ & 4 \\
$a_2$ & 丈量並設計系統櫃 & $\langle 1, 3 \rangle$ & 3 \\
$a_3$ & 水電配管 & $\langle 2, 4 \rangle$ & 5 \\
$a_4$ & 更換窗框 & $\langle 2, 5 \rangle$ & 3 \\
$a_5$ & 工廠製作系統櫃 & $\langle 3, 7 \rangle$ & 14 \\
$a_6$ & 浴室防水與貼磚 & $\langle 4, 5 \rangle$ & 6 \\
$a_7$ & 天花板木作 & $\langle 4, 6 \rangle$ & 6 \\
$a_8$ & （虛活動） & $\langle 5, 6 \rangle$ & 0 \\
$a_9$ & 安裝衛浴設備 & $\langle 5, 8 \rangle$ & 2 \\
$a_{10}$ & 油漆 & $\langle 6, 7 \rangle$ & 5 \\
$a_{11}$ & 安裝系統櫃 & $\langle 7, 8 \rangle$ & 3 \\
$a_{12}$ & 裝燈具與清潔 & $\langle 8, 9 \rangle$ & 2 \\
\bottomrule
\end{tabular}
\end{center}

\begin{center}
\begin{tikzpicture}[x=1.22cm, y=1.1cm]
\node[tps-ev] (e1) at (0, 1.2) {1};
\node[tps-ev] (e2) at (1.9, 1.2) {2};
\node[tps-ev] (e3) at (1.9, 2.8) {3};
\node[tps-ev] (e4) at (3.9, 1.2) {4};
\node[tps-ev] (e5) at (5.5, -0.2) {5};
\node[tps-ev] (e6) at (7.1, 1.2) {6};
\node[tps-ev] (e7) at (9.1, 1.2) {7};
\node[tps-ev] (e8) at (10.5, -0.2) {8};
\node[tps-ev] (e9) at (12.1, -0.2) {9};
\draw[tps-a] (e1) -- node[tps-al, below] {$a_1 = 4$} (e2);
\draw[tps-a] (e1) -- node[tps-al, above left] {$a_2 = 3$} (e3);
\draw[tps-a] (e2) -- node[tps-al, above] {$a_3 = 5$} (e4);
\draw[tps-a] (e2) -- node[tps-al, below left] {$a_4 = 3$} (e5);
\draw[tps-a] (e3) to[out=0, in=110] node[tps-al, above, pos=0.45] {$a_5 = 14$} (e7);
\draw[tps-a] (e4) -- node[tps-al, right=1pt] {$a_6 = 6$} (e5);
\draw[tps-a] (e4) -- node[tps-al, above] {$a_7 = 6$} (e6);
\draw[tps-dummy] (e5) -- node[tps-al, right=1pt] {$a_8 = 0$} (e6);
\draw[tps-a] (e5) -- node[tps-al, below] {$a_9 = 2$} (e8);
\draw[tps-a] (e6) -- node[tps-al, above] {$a_{10} = 5$} (e7);
\draw[tps-a] (e7) -- node[tps-al, right=1pt] {$a_{11} = 3$} (e8);
\draw[tps-a] (e8) -- node[tps-al, above] {$a_{12} = 2$} (e9);
\end{tikzpicture}
\diagcaption{裝修工程的 AOE 網路。 圓圈是事件，邊是活動，數字是工期（天）；虛線是虛活動。}
\label{fig:tps-aoe}
\end{center}

虛活動 $a_8$ 的用處：油漆要等泥作（事件 5）和木作（$a_7$）都完成，衛浴設備卻只要等泥
作，不必等木作。 如果讓 $a_7$ 也指向事件 5，衛浴設備就被迫等木作；如果讓 $a_4$、$a_6$
直接指向事件 6，衛浴設備（從事件 6 出發）一樣要等木作。 用兩個事件加一條虛活動，兩種
先後關係才都對。

\subsection{事件的最早時間：往前算}

\textbf{事件 $j$ 的最早發生時間} $\mathit{ee}[j]$，是從開工算起，事件 $j$ 最早在第幾天
發生。 開始事件是 $0$。 事件 $j$ 要等每一項指向它的活動 $\langle i, j \rangle$ 都完成，
而這項活動最早在 $\mathit{ee}[i]$ 開工、$\mathit{ee}[i] + \mathit{dur}$ 完工，所以
\[
\mathit{ee}[\text{開始}] = 0, \qquad
\mathit{ee}[j] = \max_{\langle i, j \rangle} \bigl( \mathit{ee}[i] + \mathit{dur}(\langle i, j \rangle) \bigr).
\]
這就是第~18~章的動態規劃。 $\mathit{ee}[j]$ 是從開始事件到 $j$ 的\textbf{最長路徑}長
度：子問題是「到每個事件的最長路徑」，最後一條邊有幾種選擇就比幾次。 有向無環圖上的最長
路徑有最佳子結構，第~18~章觀念題第~\ref{dp-q-sub}~題 (a) 用剪下貼上證明過。 填表時，填一
格之前它用到的格子都要先填好（第~\ref{sec:dp-bottomup}~節），在圖上就是\textbf{依拓樸順
序}算。

沒有環是必要的。 圖上有環時，最長路徑可以一直繞圈變長，這個遞迴式就失效了；一般圖上的最
長簡單路徑也沒有已知的多項式時間解法（第~\ref{sec:dp-substructure}~節）。 把 $\max$ 換成
$\min$ 就是有向無環圖上的最短路徑，第~25~章會再用到。

裝修工程的事件編號 $1..9$ 剛好就是一個拓樸順序，依序算：

\begin{center}
\begin{tabular}{cclc}
\toprule
事件 $j$ & 指向 $j$ 的活動 & 計算 & $\mathit{ee}[j]$ \\
\midrule
1 & -- & 開始事件 & 0 \\
2 & $a_1$ & $0 + 4$ & 4 \\
3 & $a_2$ & $0 + 3$ & 3 \\
4 & $a_3$ & $4 + 5$ & 9 \\
5 & $a_4$、$a_6$ & $\max(4 + 3,\ 9 + 6)$ & 15 \\
6 & $a_7$、$a_8$ & $\max(9 + 6,\ 15 + 0)$ & 15 \\
7 & $a_5$、$a_{10}$ & $\max(3 + 14,\ 15 + 5)$ & 20 \\
8 & $a_9$、$a_{11}$ & $\max(15 + 2,\ 20 + 3)$ & 23 \\
9 & $a_{12}$ & $23 + 2$ & 25 \\
\bottomrule
\end{tabular}
\end{center}

結束事件的 $\mathit{ee}[9] = 25$，就是\textbf{最短工期}：人手再多，也要 $25$ 天才能完
工。 記成 $T = 25$。

\subsection{事件的最晚時間：往回算}

\textbf{事件 $i$ 的最晚發生時間} $\mathit{le}[i]$，是在不拖延工期 $T$ 的前提下，事件
$i$ 最晚可以在第幾天發生。 結束事件最晚就是 $T$。 從事件 $i$ 出發的每一項活動
$\langle i, j \rangle$，都要在 $\mathit{le}[j]$ 之前完成，所以最晚要在
$\mathit{le}[j] - \mathit{dur}$ 開工；$i$ 必須趕上其中最早的那一個：
\[
\mathit{le}[\text{結束}] = \mathit{ee}[\text{結束}] = T, \qquad
\mathit{le}[i] = \min_{\langle i, j \rangle} \bigl( \mathit{le}[j] - \mathit{dur}(\langle i, j \rangle) \bigr).
\]
$\mathit{le}[i]$ 用到的是 $i$ 的後繼，所以\textbf{依拓樸順序倒過來}算。

\begin{center}
\begin{tabular}{cclc}
\toprule
事件 $i$ & 從 $i$ 出發的活動 & 計算 & $\mathit{le}[i]$ \\
\midrule
9 & -- & 結束事件，$= \mathit{ee}[9]$ & 25 \\
8 & $a_{12}$ & $25 - 2$ & 23 \\
7 & $a_{11}$ & $23 - 3$ & 20 \\
6 & $a_{10}$ & $20 - 5$ & 15 \\
5 & $a_8$、$a_9$ & $\min(15 - 0,\ 23 - 2)$ & 15 \\
4 & $a_6$、$a_7$ & $\min(15 - 6,\ 15 - 6)$ & 9 \\
3 & $a_5$ & $20 - 14$ & 6 \\
2 & $a_3$、$a_4$ & $\min(9 - 5,\ 15 - 3)$ & 4 \\
1 & $a_1$、$a_2$ & $\min(4 - 4,\ 6 - 3)$ & 0 \\
\bottomrule
\end{tabular}
\end{center}

開始事件的 $\mathit{le}$ 一定算回 $0$，可以拿來驗算：如果不是 $0$，前面某一格算錯了。
$\mathit{ee}[j] = \mathit{le}[j]$ 的事件叫\textbf{關鍵事件}，這裡除了事件 3（$\mathit{ee}
= 3$、$\mathit{le} = 6$）以外都是。

\subsection{活動的最早、最晚開始時間與寬裕時間}

對活動 $a = \langle i, j \rangle$：
\begin{itemize}
\item \textbf{最早開始時間} $e(a) = \mathit{ee}[i]$：起點事件一發生就能開工。
\item \textbf{最晚開始時間} $l(a) = \mathit{le}[j] - \mathit{dur}(a)$：再晚開工，終點事件
就趕不上它的最晚時間。
\item \textbf{寬裕時間}（slack）$l(a) - e(a)$：這項活動可以延後開工（或拖長工期）多少天而
不影響總工期。
\item 寬裕時間為 $0$ 的活動叫\textbf{關鍵活動}（critical activity）。
\end{itemize}
注意 $l(a)$ 用的是\textbf{終點}的 $\mathit{le}$，不是起點的 $\mathit{le}[i]$。 起點的
$\mathit{le}[i]$ 是「從 $i$ 出發的所有活動裡，最急的那一項」的最晚開工時間，其他活動可
以更晚。 把 $l(a)$ 寫成 $\mathit{le}[i]$ 是最常見的錯誤。

\begin{center}
\begin{tabular}{clccccc}
\toprule
活動 & 內容 & 工期 & $e$ & $l$ & $l - e$ & 關鍵 \\
\midrule
$a_1$ & 拆除舊裝潢 & 4 & 0 & $4 - 4 = 0$ & 0 & $\checkmark$ \\
$a_2$ & 丈量並設計系統櫃 & 3 & 0 & $6 - 3 = 3$ & 3 & \\
$a_3$ & 水電配管 & 5 & 4 & $9 - 5 = 4$ & 0 & $\checkmark$ \\
$a_4$ & 更換窗框 & 3 & 4 & $15 - 3 = 12$ & 8 & \\
$a_5$ & 工廠製作系統櫃 & 14 & 3 & $20 - 14 = 6$ & 3 & \\
$a_6$ & 浴室防水與貼磚 & 6 & 9 & $15 - 6 = 9$ & 0 & $\checkmark$ \\
$a_7$ & 天花板木作 & 6 & 9 & $15 - 6 = 9$ & 0 & $\checkmark$ \\
$a_8$ & （虛活動） & 0 & 15 & $15 - 0 = 15$ & 0 & $\checkmark$ \\
$a_9$ & 安裝衛浴設備 & 2 & 15 & $23 - 2 = 21$ & 6 & \\
$a_{10}$ & 油漆 & 5 & 15 & $20 - 5 = 15$ & 0 & $\checkmark$ \\
$a_{11}$ & 安裝系統櫃 & 3 & 20 & $23 - 3 = 20$ & 0 & $\checkmark$ \\
$a_{12}$ & 裝燈具與清潔 & 2 & 23 & $25 - 2 = 23$ & 0 & $\checkmark$ \\
\bottomrule
\end{tabular}
\end{center}

\begin{examnote}
AOE 的記號各本書不同，作答前先對照題目用的是哪一套。 本書用 $\mathit{ee}$／$\mathit{le}$
表示事件的最早、最晚時間，$e$／$l$ 表示活動的最早、最晚開始時間。 常見的其他寫法有：
事件用 $\mathit{ve}(j)$／$\mathit{vl}(j)$、$\mathit{TE}$／$\mathit{TL}$ 或
$\mathit{early}$／$\mathit{late}$；活動用 $\mathit{ES}$、$\mathit{EF}$、$\mathit{LS}$、
$\mathit{LF}$（最早開始、最早完成、最晚開始、最晚完成，專案管理的要徑法常用），這時寬裕
時間是 $\mathit{LS} - \mathit{ES}$，也等於 $\mathit{LF} - \mathit{EF}$。 寬裕時間也叫
浮時（float）。 表格一律畫完整：事件表有 $\mathit{ee}$、$\mathit{le}$ 兩欄，活動表有
$e$、$l$、$l - e$，最後寫出關鍵活動和關鍵路徑。
\end{examnote}

\subsection{關鍵路徑}

\begin{definition}[關鍵路徑]
從開始事件到結束事件、全部由關鍵活動組成的路徑，叫做\textbf{關鍵路徑}（critical path）。
\end{definition}

\begin{theorem}
\label{thm:tps-critical}
在 AOE 網路裡，活動 $a$ 是關鍵活動，若且唯若 $a$ 在某一條從開始事件到結束事件的最長路
徑上。 因此關鍵路徑就是最長路徑，長度都等於最短工期 $T$，而且至少有一條。
\end{theorem}

證明的關鍵是看懂 $\mathit{le}$ 的意思。 $\mathit{ee}[i]$ 是「開始到 $i$」的最長路徑；
往回算的 $\mathit{le}$ 對稱地和「$j$ 到結束」的最長路徑有關。

\begin{whybox}[定理~\ref{thm:tps-critical}]
記 $\mathit{tail}[j]$ 為從事件 $j$ 到結束事件的最長路徑長度。 先證明
$\mathit{le}[j] = T - \mathit{tail}[j]$。 依拓樸順序倒過來歸納：結束事件
$\mathit{tail} = 0$，$\mathit{le} = T$，成立。 其他事件 $i$ 的最長路徑一定先走某條邊
$\langle i, j \rangle$，所以 $\mathit{tail}[i] = \max_{\langle i, j \rangle}
(\mathit{dur} + \mathit{tail}[j])$；由歸納假設，
\[
\mathit{le}[i] = \min_{\langle i, j \rangle} \bigl( T - \mathit{tail}[j] - \mathit{dur} \bigr)
= T - \max_{\langle i, j \rangle} \bigl( \mathit{dur} + \mathit{tail}[j] \bigr) = T - \mathit{tail}[i].
\]

對活動 $a = \langle i, j \rangle$，經過 $a$ 的最長路徑長度是
$\mathit{ee}[i] + \mathit{dur}(a) + \mathit{tail}[j]$（前段、後段各取最長，有向無環圖
上接起來不會重複經過頂點）。 寬裕時間
\[
l(a) - e(a) = \mathit{le}[j] - \mathit{dur}(a) - \mathit{ee}[i]
= T - \bigl( \mathit{ee}[i] + \mathit{dur}(a) + \mathit{tail}[j] \bigr).
\]
所以寬裕時間就是「整張圖的最長路徑」減去「經過 $a$ 的最長路徑」，一定 $\ge 0$；等於
$0$ 若且唯若 $a$ 在某條長度 $T$ 的路徑上。 從開始到結束的最長路徑一定存在（路徑只有有
限多條），它上面的每一項活動都是關鍵活動，所以至少有一條關鍵路徑。
\end{whybox}

證明中間那個式子很好用：\textbf{寬裕時間 $=$ 工期 $-$ 經過這項活動的最長路徑}。 想知
道一項活動能拖幾天，看經過它的最長路徑比 $T$ 短多少就好。

裝修工程的關鍵活動是 $a_1, a_3, a_6, a_7, a_8, a_{10}, a_{11}, a_{12}$，它們組成\textbf{兩
條}關鍵路徑（圖~\ref{fig:tps-aoe-crit}）：
\begin{center}
\small
\begin{tabular}{lll}
\toprule
路徑 & 活動 & 長度 \\
\midrule
$1 \to 2 \to 4 \to 5 \to 6 \to 7 \to 8 \to 9$ & $a_1, a_3, a_6, a_8, a_{10}, a_{11}, a_{12}$ & $4 + 5 + 6 + 0 + 5 + 3 + 2 = 25$ \\
$1 \to 2 \to 4 \to 6 \to 7 \to 8 \to 9$ & $a_1, a_3, a_7, a_{10}, a_{11}, a_{12}$ & $4 + 5 + 6 + 5 + 3 + 2 = 25$ \\
\bottomrule
\end{tabular}
\end{center}
浴室貼磚和天花板木作同時開工、一樣要 $6$ 天，兩條路一樣長。

\begin{center}
\begin{tikzpicture}[x=1.22cm, y=1.1cm]
\node[tps-cev] (e1) at (0, 1.2) {1};
\node[tps-cev] (e2) at (1.9, 1.2) {2};
\node[tps-ev] (e3) at (1.9, 2.8) {3};
\node[tps-cev] (e4) at (3.9, 1.2) {4};
\node[tps-cev] (e5) at (5.5, -0.2) {5};
\node[tps-cev] (e6) at (7.1, 1.2) {6};
\node[tps-cev] (e7) at (9.1, 1.2) {7};
\node[tps-cev] (e8) at (10.5, -0.2) {8};
\node[tps-cev] (e9) at (12.1, -0.2) {9};
\foreach \v/\ee/\le/\pos in {1/0/0/left, 2/4/4/above, 3/3/6/above, 4/9/9/above, 5/15/15/below, 6/15/15/above, 7/20/20/right, 8/23/23/below, 9/25/25/below} {
    \node[tps-time, \pos=0.33cm] at (e\v) {$[\ee, \le]$};
}
\draw[tps-crit] (e1) -- node[tps-al, above] {$a_1$} (e2);
\draw[tps-a] (e1) -- node[tps-al, above left] {$a_2$} (e3);
\draw[tps-crit] (e2) -- node[tps-al, above] {$a_3$} (e4);
\draw[tps-a] (e2) -- node[tps-al, below left] {$a_4$} (e5);
\draw[tps-a] (e3) to[out=0, in=110] node[tps-al, above, pos=0.45] {$a_5$} (e7);
\draw[tps-crit] (e4) -- node[tps-al, right=1pt] {$a_6$} (e5);
\draw[tps-crit] (e4) -- node[tps-al, above] {$a_7$} (e6);
\draw[tps-cdummy] (e5) -- node[tps-al, right=1pt] {$a_8$} (e6);
\draw[tps-a] (e5) -- node[tps-al, below] {$a_9$} (e8);
\draw[tps-crit] (e6) -- node[tps-al, above] {$a_{10}$} (e7);
\draw[tps-crit] (e7) -- node[tps-al, right=1pt] {$a_{11}$} (e8);
\draw[tps-crit] (e8) -- node[tps-al, above] {$a_{12}$} (e9);
\end{tikzpicture}
\diagcaption{裝修工程算完之後：方括號是 $[\mathit{ee}, \mathit{le}]$，紅色粗框是關鍵事件，紅色粗線是關鍵活動。 兩條關鍵路徑在事件 4 分開、在事件 6 會合。}
\label{fig:tps-aoe-crit}
\end{center}

\subsection{考試常問的觀念}

下面幾點都用裝修工程說明，數字可以對照上面的表。

\textbf{(1) 兩端都是關鍵事件的活動，不一定是關鍵活動。} 事件 2、5 都是關鍵事件，
$a_4$（更換窗框）卻有 $8$ 天的寬裕；事件 5、8 也都是關鍵事件，$a_9$ 有 $6$ 天。 關鍵
事件只說明「這個時間點沒有彈性」，連接它們的活動要另外檢查 $\mathit{ee}[i] +
\mathit{dur} = \mathit{le}[j]$ 才算關鍵。 所以\textbf{找關鍵路徑要看活動表，不能只把關
鍵事件連起來}。

\textbf{(2) 縮短非關鍵活動，工期不會變。} 經過 $a_5$ 的最長路徑是 $3 + 14 + 3 + 2 = 22$
天，就算系統櫃一天做好，另外兩條 $25$ 天的路徑還在。

\textbf{(3) 縮短一項關鍵活動，工期不一定變短。} 有多條關鍵路徑時，只縮短其中一條上的活
動沒有用：$a_6$（貼磚）縮成 $4$ 天，經過 $a_7$ 的路徑仍然是 $25$ 天，工期不變。 要縮短
工期，得縮短\textbf{每一條}關鍵路徑，例如縮短各條共用的 $a_1$、$a_3$、$a_{10}$、
$a_{11}$、$a_{12}$，或同時縮短 $a_6$ 和 $a_7$。

\textbf{(4) 縮短 $k$ 天，工期不一定少 $k$ 天。} 原本不是關鍵的路徑可能變成新的關鍵路徑。
$a_3$（水電配管）縮短 $4$ 天變成 $1$ 天，兩條關鍵路徑都變成 $21$ 天，但經過 $a_2$、$a_5$
的路徑是 $22$ 天，新的工期是 $22$ 天，只少了 $3$ 天。 縮短後一定要重算，或至少把每條路徑
的長度重新比一次。

\textbf{(5) 延後不超過寬裕時間，工期不變；超過多少，工期就延後多少。} $a_9$（衛浴設備）
晚 $6$ 天開工都沒關係，晚 $8$ 天就讓工期變成 $27$ 天。 這是因為寬裕時間等於 $T$ 減去經過
它的最長路徑（定理~\ref{thm:tps-critical} 的證明）。

\textbf{(6) 同一條路徑上的寬裕時間不能相加。} $a_2$、$a_5$ 的寬裕都是 $3$ 天，但它們在同
一條路徑 $1 \to 3 \to 7$ 上，兩個合起來只能延後 $3$ 天：$a_2$ 晚 $2$ 天、$a_5$ 再多拖 $2$
天，經過它們的路徑變成 $26$ 天。 寬裕時間是「其他活動都照最早時間做」時，這一項單獨能用
的餘裕。

\subsection{C 程式}

程式讀入事件數 $n$、活動數 $m$ 和每項活動的起點、終點、工期（活動依輸入順序編號
$a_1, a_2, \ldots$）。 先用 Kahn 演算法排出拓樸順序，順便檢查有沒有環；接著依拓樸順序
往前算 $\mathit{ee}$，倒過來算 $\mathit{le}$，再印出事件表、活動表和所有的關鍵路徑。

這支程式的出邊串列不用 \texttt{malloc}，改用陣列下標當指標：\texttt{first[u]} 是 $u$ 的
第一項活動，\texttt{nxt[k]} 是同一個起點的下一項，$-1$ 代表 \texttt{NULL}。 新的活動接
在尾端（\texttt{last[u]} 記住尾巴），串列保持輸入的順序。 往回算時，每個事件的
$\mathit{le}$ 先設成 $T$ 再取最小值，結束事件沒有出邊，就留在 $T$，有好幾個結束事件也不
用另外處理。 列出關鍵路徑是一個小回溯：從開始事件出發，只沿著寬裕時間 $0$ 的活動往下
走，走到結束事件就印出一條。

\begin{example}[label=ex:tps-aoe]{c}{(aoe.c) AOE 網路：事件時間、寬裕時間與關鍵路徑}
#include <stdio.h>

#define MAXN 100                // 事件數上限
#define MAXM 1000               // 活動數上限

int n, m;                       // 事件數、活動數
int from[MAXM], to[MAXM], dur[MAXM];    // 活動 k：事件 from[k] -> to[k]
int first[MAXN], last[MAXN], nxt[MAXM]; // 出邊串列，用陣列下標當指標，-1 是結尾
int indeg[MAXN], order[MAXN], ee[MAXN], le[MAXN], path[MAXN];

// Kahn 演算法：order 本身當佇列用，回傳排出幾個事件
int topo_order(void)
{
    int deg[MAXN], cnt = 0;
    for (int v = 0; v < n; v++) {
        deg[v] = indeg[v];
        if (deg[v] == 0) {
            order[cnt++] = v;
        }
    }
    for (int h = 0; h < cnt; h++) {
        for (int k = first[order[h]]; k != -1; k = nxt[k]) {
            if (--deg[to[k]] == 0) {
                order[cnt++] = to[k];
            }
        }
    }
    return cnt;
}

// 從事件 u 只沿著寬裕時間 0 的活動往下走，走到結束事件就印出一條關鍵路徑
void print_critical(int u, int len)
{
    path[len++] = u;
    if (first[u] == -1) {
        printf("critical path:");
        for (int i = 0; i < len; i++) {
            printf(" %d", path[i] + 1);
        }
        printf("\n");
        return;
    }
    for (int k = first[u]; k != -1; k = nxt[k]) {
        if (le[to[k]] - dur[k] == ee[u]) {
            print_critical(to[k], len);
        }
    }
}

int main(void)
{
    if (scanf("%d %d", &n, &m) != 2 || n < 1 || n > MAXN || m < 0 || m > MAXM) {
        return 1;
    }
    for (int v = 0; v < n; v++) {
        first[v] = last[v] = -1;
    }
    for (int k = 0; k < m; k++) {
        if (scanf("%d %d %d", &from[k], &to[k], &dur[k]) != 3 || from[k] < 1
            || from[k] > n || to[k] < 1 || to[k] > n || dur[k] < 0) {
            return 1;
        }
        int u = --from[k];              // 輸入的事件從 1 編號
        to[k]--;
        nxt[k] = -1;                    // 接在 u 的串列尾端，保持輸入的順序
        if (first[u] == -1) {
            first[u] = k;
        } else {
            nxt[last[u]] = k;
        }
        last[u] = k;
        indeg[to[k]]++;
    }
    if (topo_order() < n) {
        printf("not a DAG\n");
        return 0;
    }
    for (int i = 0; i < n; i++) {           // 往前算：依拓樸順序
        int u = order[i];
        for (int k = first[u]; k != -1; k = nxt[k]) {
            if (ee[u] + dur[k] > ee[to[k]]) {
                ee[to[k]] = ee[u] + dur[k];
            }
        }
    }
    int T = 0;
    for (int v = 0; v < n; v++) {
        if (ee[v] > T) {
            T = ee[v];
        }
    }
    for (int i = n - 1; i >= 0; i--) {      // 往回算：拓樸順序倒過來
        int u = order[i];
        le[u] = T;                          // 結束事件沒有出邊，留在 T
        for (int k = first[u]; k != -1; k = nxt[k]) {
            if (le[to[k]] - dur[k] < le[u]) {
                le[u] = le[to[k]] - dur[k];
            }
        }
    }
    printf("event ee le\n");
    for (int v = 0; v < n; v++) {
        printf("%d %d %d\n", v + 1, ee[v], le[v]);
    }
    printf("act edge dur e l slack\n");
    for (int k = 0; k < m; k++) {
        int e = ee[from[k]], l = le[to[k]] - dur[k];
        printf("a%d %d->%d %d %d %d %d%s\n", k + 1, from[k] + 1, to[k] + 1,
               dur[k], e, l, l - e, l == e ? " *" : "");
    }
    printf("project length = %d\n", T);
    for (int v = 0; v < n; v++) {
        if (indeg[v] == 0 && le[v] == 0) {  // 從開始事件出發
            print_critical(v, 0);
        }
    }
    return 0;
}
\end{example}
\iosample{9 12 \\ 1 2 4 \\ 1 3 3 \\ 2 4 5 \\ 2 5 3 \\ 3 7 14 \\ 4 5 6 \\ 4 6 6 \\ 5 6 0 \\ 5 8 2 \\ 6 7 5 \\ 7 8 3 \\ 8 9 2}{event ee le \\ 1 0 0 \\ 2 4 4 \\ 3 3 6 \\ 4 9 9 \\ 5 15 15 \\ 6 15 15 \\ 7 20 20 \\ 8 23 23 \\ 9 25 25 \\ act edge dur e l slack \\ a1 1->2 4 0 0 0 * \\ a2 1->3 3 0 3 3 \\ a3 2->4 5 4 4 0 * \\ a4 2->5 3 4 12 8 \\ a5 3->7 14 3 6 3 \\ a6 4->5 6 9 9 0 * \\ a7 4->6 6 9 9 0 * \\ a8 5->6 0 15 15 0 * \\ a9 5->8 2 15 21 6 \\ a10 6->7 5 15 15 0 * \\ a11 7->8 3 20 20 0 * \\ a12 8->9 2 23 23 0 * \\ project length = 25 \\ critical path: 1 2 4 5 6 7 8 9 \\ critical path: 1 2 4 6 7 8 9}

輸出和手算的兩張表一致，星號是關鍵活動。 要試第 (4) 點，把輸入裡的 \texttt{2 4 5} 改成
\texttt{2 4 1}，程式會印出 \texttt{project length = 22}，關鍵路徑只剩
\texttt{1 3 7 8 9}。

\section{本章常見錯誤}
\label{sec:tps-pitfalls}

\begin{center}
\begin{tabular}{p{0.26\linewidth}p{0.48\linewidth}p{0.16\linewidth}}
\toprule
錯誤 & 例子 & 見 \\
\midrule
一開始只放一個源點 & Kahn 演算法只把第一個入度 $0$ 的頂點放進佇列；圖分成好幾塊，或
有好幾個源點時，後面的永遠排不出來 & 第~\ref{sec:tps-kahn}~節 \\
沒有檢查排出幾個 & 有環時主迴圈照樣結束，把不完整的 $L$ 當成答案 &
第~\ref{sec:tps-kahn}~節 \\
以為剩下的都在環上 & Kahn 卡住時剩下的頂點，包含被環擋住的下游頂點 &
第~\ref{sec:tps-aov}~節 \\
DFS 依發現時間排 & 進入頂點時就加進串列，或依 $d$ 排序；要在完成時插到最前面，等於
依 $f$ 由大到小 & 第~\ref{sec:tps-dfs}~節 \\
追蹤時換了平手規則 & 題目規定佇列卻用了「挑最小的」，或鄰接串列沒照題目的順序走 &
第~\ref{sec:tps-kahn-trace}~節 \\
列舉時漏掉 & 想到哪寫到哪，沒有照列舉樹一個位置一個位置展開 &
第~\ref{sec:tps-count}~節 \\
$l(a)$ 用了起點的 $\mathit{le}$ & 活動 $\langle i, j \rangle$ 的最晚開始時間寫成
$\mathit{le}[i]$；應該是 $\mathit{le}[j] - \mathit{dur}$ & 第~\ref{sec:tps-aoe}~節 \\
$\mathit{le}$ 的初值設錯 & 結束事件的 $\mathit{le}$ 設成 $0$ 或 $\infty$；應該等於
$\mathit{ee}[\text{結束}]$ & 第~\ref{sec:tps-aoe}~節 \\
把關鍵事件連起來當關鍵路徑 & 裝修工程的 $a_4$、$a_9$ 兩端都是關鍵事件，卻有寬裕時間 &
第~\ref{sec:tps-aoe}~節 \\
縮短 $k$ 天就少 $k$ 天 & 忽略其他關鍵路徑，或忽略次長的路徑變成新的關鍵路徑 &
第~\ref{sec:tps-aoe}~節 \\
寬裕時間相加 & 同一條路徑上兩項活動各有 $3$ 天寬裕，以為合起來可以拖 $6$ 天 &
第~\ref{sec:tps-aoe}~節 \\
\bottomrule
\end{tabular}
\end{center}

\section{觀念題}
\label{sec:tps-exercises}

題型和研究所考試、期中期末考的筆試相近，有選擇、計算、追蹤、證明和設計。 第~\ref{tps-q-kahn}--\ref{tps-q-list}~題用
下面的圖 $G_2$：$7$ 個頂點，邊是
\[
1 \to 3,\ 4 \to 1,\ 4 \to 2,\ 2 \to 5,\ 3 \to 6,\ 5 \to 3,\ 5 \to 6,\ 7 \to 2,\ 7 \to 6,
\]
鄰接串列依編號由小到大。 解答在本章最後。

\begin{enumerate}
\item \label{tps-q-tf}（選擇，可能不只一個答案）下列哪些敘述正確？
\begin{enumerate}
\item 每個至少有一個頂點的有向無環圖，都同時有入度 $0$ 和出度 $0$ 的頂點。
\item 對有環的有向圖做 Kahn 演算法，一個頂點都排不出來。
\item 把 Kahn 演算法的佇列換成堆疊，在有向無環圖上排出來的仍然是拓樸順序。
\item 對有向無環圖做 DFS，依發現時間 $d$ 由小到大排列，是一個拓樸順序。
\item $n$ 個頂點的有向無環圖最多有 $n(n-1)/2$ 條邊。
\end{enumerate}

\item \label{tps-q-kahn}（追蹤）對 $G_2$ 做 Kahn 演算法。
\begin{enumerate}
\item 用佇列（本書的約定），寫出每一步取出的頂點、$\mathit{in}[1..7]$ 和佇列的內容。
\item 改用最小優先佇列，寫出結果。
\item 改用堆疊：一開始依編號由小到大推入，處理出邊時也依編號由小到大推入，每次取堆疊
頂端。 寫出結果。
\end{enumerate}

\item \label{tps-q-dfs}（追蹤）對 $G_2$ 做 \proc{DFS-Topo-Sort}（虛擬碼~\ref{alg:tps-dfs}），
寫出每個頂點的 $d$、$f$ 和回傳的 $L$。

\item \label{tps-q-list}（列舉）列出 $G_2$ 的所有拓樸順序，依字典序排列，並寫出總數。

\item \label{tps-q-count}（計算）
\begin{enumerate}
\item 圖有 $6$ 個頂點 A--F，邊只有 $A \to B$、$B \to C$、$D \to E$，F 沒有任何邊。 共有幾個
拓樸順序？
\item 在 (a) 的圖再加一條邊 $C \to E$，共有幾個？
\end{enumerate}

\item \label{tps-q-upper}（證明）證明：有向圖 $G$ 沒有環，若且唯若可以把頂點重新編號成
$1..n$，使每一條邊都從小編號指向大編號。 並說明這時 $G$ 的鄰接矩陣（依新編號排列）長什
麼樣子。

\item \label{tps-q-dfs-tf}（證明或反例）$G$ 是有向無環圖，對它做一次 DFS。 下列敘述對的請
證明，錯的請給反例。
\begin{enumerate}
\item 依 $f$ 由小到大排列，是把 $G$ 每條邊反向之後那張圖的拓樸順序。
\item 依 $d$ 由大到小排列，是 $G$ 的拓樸順序。
\item 若有一條從 $u$ 到 $v$ 的路徑（$u \ne v$），則 $f[v] < f[u]$。
\item 若 $f[v] < f[u]$，則有一條從 $u$ 到 $v$ 的路徑。
\end{enumerate}

\item \label{tps-q-cycle}（設計）Kahn 演算法排不完時只知道「有環」。 設計一個
$O(n + m)$ 的方法，接在 Kahn 演算法後面，找出一個環並依邊的方向印出它的頂點。 說明為什麼
一定找得到。

\item \label{tps-q-matrix}（選擇）圖用鄰接矩陣存，Kahn 演算法的時間是多少？ \quad
(A) $\Theta(n + m)$ \quad (B) $\Theta(n^2)$ \quad (C) $\Theta(nm)$ \quad
(D) $\Theta(m \lg n)$

\item \label{tps-q-aov}（計算）第~\ref{sec:tps-aov}~節的迎新活動，工期（天）是 A $1$、B $2$、
C $3$、D $5$、E $4$、F $1$、G $1$、H $1$，人手足夠。
\begin{enumerate}
\item 依拓樸順序算出每項活動的最早完成時間 $\mathit{EF}$，整個計畫最快幾天完成？ 是哪一串
活動決定的？
\item 不拖延整個計畫的前提下，G 最晚要在第幾天結束？ B 呢？
\end{enumerate}

\item \label{tps-q-aoe}（計算）AOE 網路有事件 $1..7$，活動如下（工期單位是天）：
\begin{center}
\begin{tabular}{*{10}{c}}
\toprule
$a_1$ & $a_2$ & $a_3$ & $a_4$ & $a_5$ & $a_6$ & $a_7$ & $a_8$ & $a_9$ & $a_{10}$ \\
\midrule
$\langle 1, 2 \rangle$ & $\langle 1, 3 \rangle$ & $\langle 1, 4 \rangle$ & $\langle 2, 5 \rangle$ &
$\langle 2, 6 \rangle$ & $\langle 3, 5 \rangle$ & $\langle 3, 6 \rangle$ & $\langle 4, 6 \rangle$ &
$\langle 5, 7 \rangle$ & $\langle 6, 7 \rangle$ \\
5 & 4 & 2 & 3 & 3 & 6 & 5 & 4 & 4 & 6 \\
\bottomrule
\end{tabular}
\end{center}
算出每個事件的 $\mathit{ee}$、$\mathit{le}$，每項活動的 $e$、$l$、寬裕時間，並寫出最短工期
和所有關鍵路徑。

\item \label{tps-q-aoe2}（計算）承第~\ref{tps-q-aoe}~題，各小題互相獨立。
\begin{enumerate}
\item $a_6$ 延後 $2$ 天開工，工期變成幾天？
\item $a_7$ 縮短 $2$ 天，工期變成幾天？
\item 每項活動縮短 $1$ 天的成本都一樣。 要把工期縮成 $13$ 天，合計最少要縮短幾天？ 給一種
做法並說明為什麼不能更少。
\end{enumerate}

\item \label{tps-q-aoe-tf}（是非）下列關於 AOE 網路的敘述，對的寫「對」，錯的給反例或說明。
\begin{enumerate}
\item 事件 $i$、$j$ 都是關鍵事件，則活動 $\langle i, j \rangle$ 一定是關鍵活動。
\item 所有關鍵活動的工期加起來等於最短工期。
\item 關鍵路徑可能不只一條。
\item 活動 $a$ 有 $s > 0$ 天寬裕時間。 把 $a$ 延後 $s$ 天開工後，其他活動的寬裕時間都不變。
\item 只有一個開始事件時，它的 $\mathit{le}$ 一定是 $0$。
\end{enumerate}

\item \label{tps-q-ee}（證明）AOE 網路只有一個開始事件 $s$。 證明：依拓樸順序用
$\mathit{ee}[j] = \max_{\langle i, j \rangle} (\mathit{ee}[i] + \mathit{dur})$ 算出來的
$\mathit{ee}[j]$，等於從 $s$ 到 $j$ 的最長路徑長度。 （提示：先說明每個事件都走得到。）
\end{enumerate}

\section{實作練習}
\label{sec:tps-practice}

下面是 Aki-Judge 上和本章相關的題目。

\begin{center}
\begin{tabular}{llc}
\toprule
小節 & 題目 & 難度 \\
\midrule
Kahn 演算法 & \akiref{smallest-first-topological-order}{字典序最小的拓撲排序} & 基本 \\
DFS 完成時間法 & \akiref{topological-order-by-dfs-finish-time}{以 DFS 完成順序求拓撲排序} & 基本 \\
依拓樸順序做 DP & \akiref{longest-path-in-a-dag}{有向無環圖的最長路徑} & 基本 \\
依拓樸順序做 DP & \akiref{shortest-path-in-a-directed-acyclic-graph}{有向無環圖上的最短路徑} & 基本 \\
AOV 網路的工期 & \akiref{earliest-finish-of-dependent-steps}{工序的最早完成時間} & 基本 \\
AOE 網路 & \akiref{longest-path-in-a-dag-weighted}{有向無環圖的最長路徑} & 進階 \\
\bottomrule
\end{tabular}
\end{center}

幾題的提示：「字典序最小的拓撲排序」就是最小優先佇列版的 Kahn 演算法；有環時的輸出格
式看清楚題目的說明。 「以 DFS 完成順序求拓撲排序」的走訪順序和本書不同：外層一樣依編號，
但每個頂點的出邊照\textbf{輸入的順序由後往前}試，這正好練習作答提醒框說的「照題目給的順
序走」。 表裡有兩題同名的「有向無環圖的最長路徑」：第一題數邊數，答案是所有頂點的 DP 值
取最大，包含不連通的各塊；第二題邊上有權重，只算從頂點 $1$ 出發到頂點 $n$ 的路線，相當
於 AOE 網路的 $\mathit{ee}[n]$。 但它的圖可能有別的源點，從頂點 $1$ 走不到的頂點要設成
$-\infty$，不能和 $1$ 一樣從 $0$ 開始。 「有向無環圖上的最短路徑」把 $\max$ 換成 $\min$，
走不到的設成 $\infty$，第~25~章會再講一次。 「工序的最早完成時間」是第~\ref{sec:tps-aov}~
節的 $\mathit{EF}$，前置步驟的編號可能比自己大，所以要依拓樸順序算，或寫成記憶化遞迴；
每個步驟記下「完成時間最大的前置步驟」（平手取編號小的），一路往回走就是一條讓它沒辦法更
早完成的關鍵鏈。

\section{讀完你應該會}
\label{sec:tps-checklist}

\begin{itemize}
\item 能用 Kahn 演算法和 DFS 完成時間法對一張圖做拓樸排序，照指定的平手規則寫出入度表、
佇列內容或 $d$、$f$ 表，並判斷圖有沒有環。
\item 能證明有向圖有拓樸順序若且唯若沒有環，並用括號定理證明 DFS 完成時間法的正確性。
\item 能判斷拓樸順序是否唯一，用列舉樹列出所有拓樸順序，並用 $\binom{p + q}{p}$ 一類的公式
算出常見結構的個數。
\item 能把一張工作清單畫成 AOV 網路，說明前置關係為什麼必須是偏序（沒有環）。
\item 能對 AOE 網路往前算 $\mathit{ee}$、往回算 $\mathit{le}$，填好活動的 $e$、$l$、寬裕時
間，找出所有關鍵路徑。
\item 能回答縮短或延後某項活動之後工期怎麼變，並指出「關鍵事件之間的活動不一定關鍵」這類
陷阱。
\end{itemize}

\textbf{延伸閱讀}（書末會整理完整的參考資料）

\begin{itemize}
\item Thomas H. Cormen 等，\textit{Introduction to Algorithms}，第 4 版，20.4 節與 22.2 節：
DFS 完成時間法的拓樸排序，以及有向無環圖上的最短路徑；22.2 節後半講怎麼用它找關鍵路徑。
\item Ellis Horowitz、Sartaj Sahni、Susan Anderson-Freed，\textit{Fundamentals of Data
Structures in C}，第 2 版，第 6 章：AOV、AOE 網路的名詞和考卷上常見的表格寫法都來自這一系
列的資料結構課本。
\item Donald E. Knuth，\textit{The Art of Computer Programming, Vol.~1}，第 3 版，2.2.3 節：
用串列實作的線性時間拓樸排序，以及偏序的討論。
\item A. B. Kahn，``Topological Sorting of Large Networks''，\textit{Communications of the
ACM} 5(11)，1962：反覆取出沒有前驅的頂點的做法，本章的 Kahn 演算法以他命名。
\item James E. Kelley Jr.、Morgan R. Walker，``Critical-Path Planning and Scheduling''，
\textit{Proceedings of the Eastern Joint Computer Conference}，1959：要徑法（CPM）的早期論
文，$\mathit{ES}$、$\mathit{LS}$ 一類記號的來源。
\end{itemize}

\addcontentsline{toc}{section}{觀念題解答}
\section*{觀念題解答}

\begin{enumerate}
\item (a)、(c)、(e) 正確。 (a)：引理~\ref{lem:tps-source}。 (b)：錯，範例~\ref{ex:tps-topo}
的第二張圖排出了 $1$、$5$ 才卡住；只有環上和被環擋住的頂點排不出來。 (c)：對，
定理~\ref{thm:tps-kahn} 的證明只用到「$Q$ 裡恰好是還沒排出去、入度已經是 $0$ 的頂點」，
和取出的順序無關。 (d)：錯，兩個頂點、一條邊 $2 \to 1$，DFS 先從 $1$ 開始，$d[1] = 1 <
d[2] = 3$，依 $d$ 排是 $1, 2$，違反這條邊。 (e)：對，依拓樸順序，每條邊都從前面的頂點指
向後面的，不同的邊對應不同的「一前一後」的頂點對，最多 $\binom{n}{2}$ 對；讓每一對 $i < j$
都有邊 $i \to j$ 就達到。

\item 一開始 $\mathit{in}[1..7] = 1, 2, 2, 0, 1, 3, 0$，$\mathit{Adj}[4] = 1, 2$，
$\mathit{Adj}[7] = 2, 6$。
\begin{enumerate}
\item
\begin{center}
\small
\begin{tabular}{cc*{7}{c}l}
\toprule
步驟 & 取出 & 1 & 2 & 3 & 4 & 5 & 6 & 7 & 佇列 \\
\midrule
開始 & & 1 & 2 & 2 & 0 & 1 & 3 & 0 & 4 7 \\
1 & 4 & 0 & 1 & 2 & \tpsgone & 1 & 3 & 0 & 7 1 \\
2 & 7 & 0 & 0 & 2 & \tpsgone & 1 & 2 & \tpsgone & 1 2 \\
3 & 1 & \tpsgone & 0 & 1 & \tpsgone & 1 & 2 & \tpsgone & 2 \\
4 & 2 & \tpsgone & \tpsgone & 1 & \tpsgone & 0 & 2 & \tpsgone & 5 \\
5 & 5 & \tpsgone & \tpsgone & 0 & \tpsgone & \tpsgone & 1 & \tpsgone & 3 \\
6 & 3 & \tpsgone & \tpsgone & \tpsgone & \tpsgone & \tpsgone & 0 & \tpsgone & 6 \\
7 & 6 & \tpsgone & \tpsgone & \tpsgone & \tpsgone & \tpsgone & \tpsgone & \tpsgone & （空） \\
\bottomrule
\end{tabular}
\end{center}
結果 $4, 7, 1, 2, 5, 3, 6$。
\item 第 1 步之後 $Q = \{1, 7\}$，取 $1$；之後 $7$、$2$、$5$、$3$、$6$ 都是唯一的選擇。
結果 $4, 1, 7, 2, 5, 3, 6$。
\item 堆疊一開始是 $4, 7$（$7$ 在頂端）。 取出 $7$：$\mathit{in}[2] = 1$、$\mathit{in}[6] = 2$。
取出 $4$：$1$、$2$ 依序推入，頂端是 $2$。 取出 $2$，推入 $5$；取出 $5$，$\mathit{in}[3] = 1$、
$\mathit{in}[6] = 1$；取出 $1$，推入 $3$；取出 $3$，推入 $6$；取出 $6$。 結果
$7, 4, 2, 5, 1, 3, 6$。 三種容器得到三個不同的答案，這就是作答提醒框要你先寫約定的原因。
\end{enumerate}

\item 從 $1$：$1 \to 3 \to 6$，$d = 1, 2, 3$，$f[6] = 4$、$f[3] = 5$、$f[1] = 6$。 外層到
$2$：$d[2] = 7$，發現 $5$（$d = 8$），$\mathit{Adj}[5] = 3, 6$ 都是黑的，$f[5] = 9$、
$f[2] = 10$。 外層到 $4$：$d[4] = 11$，$1$、$2$ 都是黑的，$f[4] = 12$。 外層到 $7$：
$d[7] = 13$，$f[7] = 14$。
\begin{center}
\begin{tabular}{c*{7}{c}}
\toprule
$v$ & 1 & 2 & 3 & 4 & 5 & 6 & 7 \\
\midrule
$d$ & 1 & 7 & 2 & 11 & 8 & 3 & 13 \\
$f$ & 6 & 10 & 5 & 12 & 9 & 4 & 14 \\
\bottomrule
\end{tabular}
\end{center}
依 $f$ 由大到小：$L = 7, 4, 2, 5, 1, 3, 6$。 （和第~\ref{tps-q-kahn}~題 (c) 剛好相同，這是巧
合，一般不會一樣。）

\item 源點是 $4$、$7$。 $1$ 只等 $4$；$2$ 等 $4$、$7$；$5$ 等 $2$；$3$ 等 $1$、$5$；$6$ 最後。
照列舉樹展開：
\begin{center}
\begin{tabular}{ll}
$4, 1, 7, 2, 5, 3, 6$ & $7, 4, 1, 2, 5, 3, 6$ \\
$4, 7, 1, 2, 5, 3, 6$ & $7, 4, 2, 1, 5, 3, 6$ \\
$4, 7, 2, 1, 5, 3, 6$ & $7, 4, 2, 5, 1, 3, 6$ \\
$4, 7, 2, 5, 1, 3, 6$ & \\
\end{tabular}
\end{center}
共 $7$ 個（左欄由上往下，再接右欄，是字典序）。 檢查：$3$、$6$ 一定是最後兩個；前五個
裡 $1$ 在 $4$ 後面、$2$ 在 $4$ 和 $7$ 後面、$5$ 在 $2$ 後面。 以 $4$ 開頭時，$7, 2, 5$ 必須
依序，$1$ 插進它們前後的 $4$ 個位置之一，有 $4$ 個；以 $7$ 開頭時，$4$ 一定第二，剩下的
$1$ 插進 $2, 5$ 前後的 $3$ 個位置之一，有 $3$ 個。 合計 $7$。

\item (a) 三塊：鏈 ABC（$1$ 種）、鏈 DE（$1$ 種）、F（$1$ 種），共
$\frac{6!}{3!\,2!\,1!} = 60$ 個。 (b) 加上 $C \to E$ 後，A--E 連成一塊：E 一定是這五個裡
最後的，A、B、C 依序，D 可以放在 A、B、C 之間或前後 $4$ 個位置之一，所以這塊有 $4$ 種；
再和 F 合起來，$\binom{6}{1} \cdot 4 = 24$ 個。

\item 若沒有環，由定理~\ref{thm:tps-exists} 有拓樸順序 $v_1, \ldots, v_n$，把 $v_k$ 編成
$k$，每條邊都從前面指到後面，也就是從小編號指向大編號。 反過來，如果每條邊都從小編號指向
大編號，沿著任何路徑編號嚴格變大，不可能回到起點，所以沒有環。 依新編號排列的鄰接矩陣，
$A[i][j] = 1$ 只會出現在 $i < j$，對角線和對角線以下全是 $0$：是嚴格上三角矩陣。

\item
\begin{enumerate}
\item 對。 定理~\ref{thm:tps-dfs}：每條邊 $u \to v$ 有 $f[v] < f[u]$；反向後這條邊是 $v \to u$，
依 $f$ 由小到大時 $v$ 排在 $u$ 前面，合乎要求。
\item 錯。 三個頂點，邊 $1 \to 2$、$1 \to 3$、$3 \to 2$。 從 $1$ 開始：$d[1] = 1$，發現 $2$
（$d = 2$、$f = 3$），再發現 $3$（$d = 4$，$2$ 已經是黑的，$f = 5$），$f[1] = 6$。 依 $d$
由大到小是 $3, 2, 1$，違反 $1 \to 2$。
\item 對。 路徑 $u = x_0 \to x_1 \to \cdots \to x_k = v$ 的每條邊都有 $f[x_{i+1}] < f[x_i]$，
串起來得到 $f[v] < f[u]$。
\item 錯。 兩個頂點、沒有邊：$f[1] = 2 < 4 = f[2]$，但 $2$ 走不到 $1$。
\end{enumerate}

\item 設 $R$ 是 Kahn 演算法結束時還沒排出去的頂點，非空。 由不變量，$R$ 裡每個頂點 $v$ 的
$\mathit{in}[v] \ge 1$，也就是至少有一個前驅也在 $R$ 裡。 掃一遍所有的邊，對每條兩端都在
$R$ 的邊 $u \to v$ 設 $\mathit{pre}[v] \gets u$（$O(n + m)$）。 從 $R$ 裡任取一個頂點 $x$，
一直走 $x \gets \mathit{pre}[x]$ 並標記走過的頂點；$R$ 是有限集合，最多 $|R| + 1$ 步一定走到
一個已經標記過的頂點 $y$。 $y$ 沿著 $\mathit{pre}$ 走一圈會回到 $y$，把這一圈記下來再倒過來
印，就是依邊的方向的一個環。 例如範例~\ref{ex:tps-topo} 的第二張圖，$R = \{2, 3, 4\}$，
$\mathit{pre}[2] = 4$、$\mathit{pre}[3] = 2$、$\mathit{pre}[4] = 3$（邊 $5 \to 3$ 的 $5$ 已
經排出去，不算），從 $2$ 走 $2, 4, 3, 2$，倒過來印 $2 \to 3 \to 4 \to 2$。

\item (B)。 算入度要看矩陣的每一格；每取出一個頂點，要掃它那一整列找出邊。 兩件事都是
$n^2$ 格，不管 $m$ 多小。

\item (a) 依拓樸順序 A, B, C, E, D, G, F, H：
$\mathit{EF}(A) = 1$、$\mathit{EF}(B) = 1 + 2 = 3$、$\mathit{EF}(C) = 1 + 3 = 4$、
$\mathit{EF}(E) = 3 + 4 = 7$、$\mathit{EF}(D) = \max(3, 4) + 5 = 9$、
$\mathit{EF}(G) = 7 + 1 = 8$、$\mathit{EF}(F) = 9 + 1 = 10$、
$\mathit{EF}(H) = \max(10, 8) + 1 = 11$。 最快 $11$ 天，由 A, C, D, F, H 這一串決定
（$1 + 3 + 5 + 1 + 1 = 11$）。
(b) H 最晚第 $11 - 1 = 10$ 天開始，所以 G 最晚第 $10$ 天結束（比最早的 $8$ 晚 $2$ 天）。
F 也最晚第 $10$ 天結束，D 最晚第 $10 - 1 = 9$ 天結束、第 $9 - 5 = 4$ 天開始；E 最晚在 G
開始的第 $10 - 1 = 9$ 天結束、第 $9 - 4 = 5$ 天開始。 B 要趕上 D 和 E 中較早的開始時間，
最晚第 $\min(4, 5) = 4$ 天結束（比最早的 $3$ 晚 $1$ 天）。 這就是往回算的 $\mathit{le}$，
只是放在活動上。

\item 編號 $1..7$ 已經是拓樸順序。
\begin{center}
\begin{tabular}{cll}
\toprule
事件 & $\mathit{ee}$ & $\mathit{le}$ \\
\midrule
1 & 0 & $\min(6 - 5,\ 4 - 4,\ 5 - 2) = 0$ \\
2 & $0 + 5 = 5$ & $\min(11 - 3,\ 9 - 3) = 6$ \\
3 & $0 + 4 = 4$ & $\min(11 - 6,\ 9 - 5) = 4$ \\
4 & $0 + 2 = 2$ & $9 - 4 = 5$ \\
5 & $\max(5 + 3,\ 4 + 6) = 10$ & $15 - 4 = 11$ \\
6 & $\max(5 + 3,\ 4 + 5,\ 2 + 4) = 9$ & $15 - 6 = 9$ \\
7 & $\max(10 + 4,\ 9 + 6) = 15$ & $15$ \\
\bottomrule
\end{tabular}
\qquad
\begin{tabular}{cccc}
\toprule
活動 & $e$ & $l$ & $l - e$ \\
\midrule
$a_1$ & 0 & 1 & 1 \\
$a_2$ & 0 & 0 & 0 \\
$a_3$ & 0 & 3 & 3 \\
$a_4$ & 5 & 8 & 3 \\
$a_5$ & 5 & 6 & 1 \\
$a_6$ & 4 & 5 & 1 \\
$a_7$ & 4 & 4 & 0 \\
$a_8$ & 2 & 5 & 3 \\
$a_9$ & 10 & 11 & 1 \\
$a_{10}$ & 9 & 9 & 0 \\
\bottomrule
\end{tabular}
\end{center}
最短工期 $15$ 天。 關鍵活動 $a_2, a_7, a_{10}$，只有一條關鍵路徑
$1 \to 3 \to 6 \to 7$（$4 + 5 + 6 = 15$）。 把輸入改成這張圖，範例~\ref{ex:tps-aoe} 會印出
同樣的表。

\item 五條從 $1$ 到 $7$ 的路徑長度：$1\!-\!2\!-\!5\!-\!7$ 是 $12$、$1\!-\!2\!-\!6\!-\!7$ 是
$14$、$1\!-\!3\!-\!5\!-\!7$ 是 $14$、$1\!-\!3\!-\!6\!-\!7$ 是 $15$、$1\!-\!4\!-\!6\!-\!7$ 是
$12$。
(a) $a_6$ 只有 $1$ 天寬裕，延後 $2$ 天讓 $1\!-\!3\!-\!5\!-\!7$ 變成 $16$，工期 $16$ 天。
(b) $1\!-\!3\!-\!6\!-\!7$ 變成 $13$，但還有兩條 $14$ 天的路徑，工期是 $14$ 天，只少 $1$ 天。
(c) 最少 $2$ 天：$a_2$、$a_{10}$ 各縮短 $1$ 天。 這時 $1\!-\!3\!-\!6\!-\!7$ 是 $13$，
$1\!-\!3\!-\!5\!-\!7$ 少了 $a_2$ 的 $1$ 天變成 $13$，$1\!-\!2\!-\!6\!-\!7$ 少了 $a_{10}$ 的
$1$ 天變成 $13$，$1\!-\!2\!-\!5\!-\!7$ 仍是 $12$，$1\!-\!4\!-\!6\!-\!7$ 變成 $11$。 不能更少：$1\!-\!3\!-\!6\!-\!7$ 本身就要少 $2$ 天，它
的活動合計至少縮短 $2$ 天。

\item (a) 錯，裝修工程的 $a_4$：事件 2、5 都是關鍵事件，$a_4$ 有 $8$ 天寬裕。 (b) 錯，
裝修工程的關鍵活動工期加起來是 $4 + 5 + 6 + 6 + 0 + 5 + 3 + 2 = 31 \ne 25$，因為有兩條
關鍵路徑，$a_6$ 和 $a_7$ 是並行的。 只有一條關鍵路徑時才相等。 (c) 對，裝修工程就有兩條。
(d) 錯，裝修工程的 $a_2$ 有 $3$ 天寬裕，延後 $3$ 天之後，同一條路徑上的 $a_5$ 寬裕從 $3$ 變
成 $0$。 (e) 對，$\mathit{le}[s] = T - \mathit{tail}[s]$（定理~\ref{thm:tps-critical} 的
證明），而所有事件都在從 $s$ 出發的路徑上，從 $s$ 出發的最長路徑就是整張圖的最長路徑，
$\mathit{tail}[s] = T$。

\item 先說明每個事件都走得到。 對事件 $j \ne s$，$j$ 的入度至少 $1$（只有 $s$ 入度 $0$），
從 $j$ 沿著進來的邊往回走；圖沒有環，不會重複，有限步內一定停在入度 $0$ 的頂點，也就是
$s$，所以 $s$ 走得到 $j$。 記 $L(j)$ 為從 $s$ 到 $j$ 的最長路徑長度。 依拓樸順序歸納證明
$\mathit{ee}[j] = L(j)$。 $j = s$ 時沒有邊指向 $s$，唯一的路徑是只有 $s$ 的空路徑，
$L(s) = 0 = \mathit{ee}[s]$。 $j \ne s$ 時，任何從 $s$ 到 $j$ 的路徑最後一條邊是某個
$\langle i, j \rangle$，前面那段是 $s$ 到 $i$ 的路徑，長度不超過 $L(i) = \mathit{ee}[i]$
（$i$ 在拓樸順序裡比 $j$ 前面，用歸納假設），所以整條不超過 $\mathit{ee}[i] + \mathit{dur}
\le \mathit{ee}[j]$，得到 $L(j) \le \mathit{ee}[j]$。 反過來，取使 $\max$ 成立的 $i$，把
$s$ 到 $i$ 的最長路徑接上 $\langle i, j \rangle$，是一條長度 $\mathit{ee}[j]$ 的路徑（有向
無環圖上不會重複經過 $j$），所以 $L(j) \ge \mathit{ee}[j]$。
\end{enumerate}

\chapter{連通性：強連通元件、割點、橋與雙連通元件}
\label{ch:connectivity}

% 本章的圖共用的樣式（前綴 con-）。 追蹤表和圖上的 d、f、low 都由模擬程式照虛擬碼
% 逐步執行後產生，再和本章的 C 程式輸出對過。
\tikzset{
    con-v/.style={circle, draw=black!70, fill=white, minimum size=0.62cm, inner sep=0pt, font=\small},
    con-cut/.style={draw=memStroke, very thick, fill=memFill},
    con-arc/.style={-{Stealth[length=1.8mm]}, draw=black!60, shorten >=0.5pt},
    con-tree/.style={draw=black, very thick},
    con-nt/.style={densely dashed},
    con-elab/.style={font=\scriptsize, fill=white, inner sep=0.8pt, text=black!75},
    con-lab/.style={font=\scriptsize, text=black!80},
    con-scc/.style={rounded corners=7pt, fill=coreFill, draw=coreStroke!50, inner sep=5pt},
    con-sccname/.style={font=\scriptsize, text=coreStroke},
    con-cv/.style={rounded corners=3pt, draw=coreStroke!70, fill=coreFill, font=\small, inner sep=4pt},
    con-e/.style={draw=black!60, thick},
    con-bridge/.style={draw=optStroke, line width=1.8pt},
    con-back/.style={draw=black!55, densely dashed, thick},
}

市政府把老城區的十幾條路改成單行道。 改完以後，交通局想先確認兩件事：從任何一個路
口出發，是不是都還能開到其他路口；如果不能，哪些路口之間可以來回，哪些只能單向進
出。 另一邊，學校要整修校園裡的大樓和走廊，一次只封一棟或一條。 總務處想知道，封
掉哪一棟大樓、哪一條走廊，會讓其他大樓之間走不通。 這兩個問題都可以一個一個試：每
封一處就從頭走一遍圖，看還連不連通。 往下讀之前，先想想有沒有辦法只走一遍圖，就把
答案全部找出來。

這一章用第~21~章的深度優先搜尋（DFS）回答這兩類問題。 有向圖的部分是\textbf{強連
通元件}：把「彼此可以來回」的頂點分成一群一群。 我們會講兩個 $\Theta(n + m)$ 的演
算法：Kosaraju 演算法做兩次 DFS，第二次在把邊全部反向的圖上做；Tarjan 演算法只做
一次 DFS，用一個叫 $\mathit{low}$ 的值和一個堆疊。 每一群縮成一個點以後，得到的圖一定
沒有環，第~22~章的拓樸排序就能用上。 無向圖的部分是\textbf{割點}、\textbf{橋}和
\textbf{雙連通元件}，也用 $\mathit{low}$ 值判斷，寫法和 Tarjan 演算法很像，但細節不
同，考試最常考的就是這張 $\mathit{low}$ 值的手算表。

這一章用到第~21~章的這些東西：鄰接串列、DFS 的白灰黑三色和時間戳記 $d[v]$、$f[v]$、
DFS 森林、括號定理（定理~\ref{thm:gtr-paren}）、白路徑定理（定理~\ref{thm:gtr-whitepath}），以及有向圖的四種邊（樹邊、後向邊、前向邊、交叉
邊）。 第~22~章「依完成時間由大到小排列」的拓樸排序法也會用到。 程式方面用到《C 程
式設計基礎》第~14~章的結構、第~15~章的遞迴和第~19~章的 \texttt{malloc}、鏈結串
列。

本章的約定和第~21~章相同：圖有 $n$ 個頂點、$m$ 條邊，頂點編號 $1$ 到 $n$（觀念題有
時用英文字母）；鄰接串列裡每個頂點的鄰居由小到大排，DFS 照這個順序走，主迴圈也從
編號小的頂點開始。 C 程式裡的頂點從 $0$ 編號，讀入時減一、印出時加一，所以輸出和手
算用的是同一套編號。

\section{強連通元件}
\label{sec:con-scc}

記號照第~21~章：有向圖裡有一條從 $u$ 到 $v$ 的路徑，就說 $v$ 可以從 $u$ 到達，記成
$u \leadsto v$。 長度 $0$ 的路徑也算，所以每個頂點都可以到達自己。 第~21~章已經簡單介紹過強連
通元件，這裡寫成正式的定義。

\begin{definition}
有向圖 $G = (V, E)$ 裡，如果 $u \leadsto v$ 而且 $v \leadsto u$，就說 $u$ 和 $v$
\textbf{互相可達}。 \textbf{強連通元件}（strongly connected component，簡稱 SCC）是
一個最大的頂點集合 $C \subseteq V$，其中任兩個頂點都互相可達。「最大」的意思是再
加進任何一個 $C$ 以外的頂點，這個性質就不成立。 整張圖只有一個強連通元件時，叫做
\textbf{強連通}（strongly connected）的圖。
\end{definition}

「互相可達」是一個\textbf{等價關係}：每個頂點和自己互相可達（長度 $0$ 的路徑）；
$u$、$v$ 互相可達時 $v$、$u$ 也是；$u$、$v$ 互相可達、$v$、$w$ 互相可達，把路徑接起
來就得到 $u$、$w$ 互相可達。 等價關係把頂點分成互不重疊的類別，每一類就是一個強連通
元件。 所以每個頂點恰好屬於一個強連通元件，元件可以只有一個頂點。

\begin{center}
\begin{tikzpicture}[x=1.15cm, y=1.15cm]
\node[con-v] (v1) at (1.5,1.8) {1};
\node[con-v] (v2) at (3.3,1.8) {2};
\node[con-v] (v3) at (2.6,-0.9) {3};
\node[con-v] (v4) at (2.4,0.45) {4};
\node[con-v] (v5) at (5.1,-0.3) {5};
\node[con-v] (v6) at (5.1,1.8) {6};
\node[con-v] (v7) at (0,0.9) {7};
\node[con-v] (v8) at (0.9,-0.9) {8};
\node[con-v] (v9) at (6.5,0.75) {9};
\begin{scope}[on background layer]
\node[con-scc, fit=(v1)(v2)(v4)] (s2) {};
\node[con-scc, fit=(v3)(v8)] (s3) {};
\node[con-scc, fit=(v5)(v6)(v9)] (s4) {};
\node[con-scc, fit=(v7)] (s1) {};
\end{scope}
\node[con-sccname, anchor=south] at (s1.north) {$C_1$};
\node[con-sccname, anchor=south] at (s2.north) {$C_2$};
\node[con-sccname, anchor=north] at (s3.south) {$C_3$};
\node[con-sccname, anchor=south] at (s4.north) {$C_4$};
\draw[con-arc] (v1) to[bend left=14] (v2);
\draw[con-arc] (v2) to[bend left=14] (v1);
\draw[con-arc] (v1) -- (v4);
\draw[con-arc] (v2) -- (v6);
\draw[con-arc] (v3) -- (v5);
\draw[con-arc] (v3) to[bend left=14] (v8);
\draw[con-arc] (v8) to[bend left=14] (v3);
\draw[con-arc] (v4) -- (v2);
\draw[con-arc] (v4) -- (v3);
\draw[con-arc] (v5) -- (v6);
\draw[con-arc] (v6) -- (v9);
\draw[con-arc] (v9) -- (v5);
\draw[con-arc] (v7) -- (v1);
\draw[con-arc] (v7) -- (v8);
\end{tikzpicture}
\diagcaption{老城區的單行道路網：$9$ 個路口、$14$ 條單行道。 藍底的四塊是四個強連通
元件 $C_1 = \{7\}$、$C_2 = \{1, 2, 4\}$、$C_3 = \{3, 8\}$、$C_4 = \{5, 6, 9\}$。}
\label{fig:con-roads}
\end{center}

圖~\ref{fig:con-roads} 是章首的路網，本章前四節都用它。 $1 \to 4 \to 2 \to 1$ 是一
個環，所以 $1, 2, 4$ 互相可達；$3$ 和 $8$ 之間兩個方向都有路；$6 \to 9 \to 5 \to 6$
又是一個環。 路口 $7$ 沒有任何路開進來，除了自己以外沒有誰能到 $7$，所以它自己一個元
件。 元件之間的路都只有一個方向：從 $C_2$ 可以經過 $2 \to 6$ 進入 $C_4$，但 $C_4$
裡沒有任何一條路開出去，一開進 $C_4$ 就回不到其他元件。

\subsection{轉置圖和元件圖}

把 $G$ 的每一條邊都反向，得到\textbf{轉置圖} $G^{\mathrm{T}} = (V, E^{\mathrm{T}})$，
其中 $E^{\mathrm{T}} = \{(v, u) : (u, v) \in E\}$。 $G$ 裡從 $u$ 到 $v$ 的路徑反過來，
就是 $G^{\mathrm{T}}$ 裡從 $v$ 到 $u$ 的路徑，所以「互相可達」在兩張圖裡完全一樣：
\textbf{$G$ 和 $G^{\mathrm{T}}$ 的強連通元件相同}。 用鄰接串列表示時，掃一遍所有的
邊、把每條 $(u, v)$ 加進 $G^{\mathrm{T}}$ 的 $v$ 的串列，花 $\Theta(n + m)$ 時間。

把每個強連通元件縮成一個點，得到的圖叫\textbf{元件圖}（component graph），記成
$G^{\mathrm{SCC}}$：它的頂點是 $G$ 的強連通元件 $C_1, C_2, \ldots, C_k$；只要 $G$ 裡
有一條邊從 $C_i$ 的某個頂點指到 $C_j$ 的某個頂點（$i \ne j$），就加一條邊 $(C_i,
C_j)$，同一對元件之間最多一條。 元件內部的邊縮掉了，不留自環。

\begin{lemma}
\label{lem:con-dag}
任何有向圖的元件圖都沒有環，也就是說 $G^{\mathrm{SCC}}$ 是一個 DAG。
\end{lemma}

\begin{proof}
假設元件圖有一個環 $C_1 \to C_2 \to \cdots \to C_j \to C_1$，$j \ge 2$，而且這些元
件互不相同。 邊 $(C_i, C_{i+1})$ 表示有 $x \in C_i$、$y \in C_{i+1}$ 使得 $(x, y)
\in E$。 元件裡的頂點互相可達，所以 $C_i$ 的任何頂點都能先走到 $x$，經過 $(x, y)$，
再走到 $C_{i+1}$ 的任何頂點。 沿著環走一圈，$C_1$ 的頂點能到 $C_2$ 的頂點，$C_2$ 的頂點
也能繞回 $C_1$ 的頂點，兩者互相可達，應該在同一個元件裡，和 $C_1 \ne C_2$ 矛盾。
\end{proof}

最直接的找法是從每個頂點各做一次 DFS，求出它能到哪些頂點，再把互相可達的配對起
來。 這要做 $n$ 次 DFS，總共 $O(n(n + m))$。 接下來兩節的演算法都只要 $\Theta(n + m)$。

\section{Kosaraju 演算法}
\label{sec:con-kosaraju}

先想一個比較容易的情況。 假設我們知道某個元件 $C$ 在元件圖裡是\textbf{匯點}（沒有
邊出去），從 $C$ 裡的任何一個頂點開始做 DFS，走得到的頂點都在 $C$ 裡，而且 $C$ 的每
個頂點都走得到，這一棵 DFS 樹剛好就是 $C$。 把 $C$ 拿掉以後，剩下的元件圖還是 DAG，
一定還有匯點，再從那裡做一次……每次都從匯點開始，就能把元件一個一個剝下來。

問題是我們不知道匯點在哪裡。 Kosaraju 演算法的想法是：先對 $G$ 做一次完整的 DFS，
\textbf{完成時間最大的頂點一定落在元件圖的源點}（沒有邊進來的元件）。 $G$ 的源點元件
在 $G^{\mathrm{T}}$ 裡正好變成匯點元件，所以第二次 DFS 改在 $G^{\mathrm{T}}$ 上做，
依完成時間由大到小挑起點，每一棵樹就是一個強連通元件。

\begin{pseudo}[label=alg:con-kosaraju]{Kosaraju 演算法}
\Procedure{Kosaraju}{$G$}
    \State 對 $G$ 做第~21~章的 DFS，求出每個頂點的完成時間 $f[u]$
    \State 建立轉置圖 $G^{\mathrm{T}}$
    \State 把每個 $\mathit{comp}[u]$ 設成 $0$；$k \gets 0$
    \For{依 $f[u]$ 由大到小的每個頂點 $u$}
        \If{$\mathit{comp}[u] = 0$}
            \State $k \gets k + 1$
            \State \Call{Assign}{$G^{\mathrm{T}}, u, k$}
        \EndIf
    \EndFor
\EndProcedure
\Procedure{Assign}{$H, u, k$}
    \State $\mathit{comp}[u] \gets k$ \Comment{$u$ 屬於第 $k$ 個元件}
    \For{每個 $v \in H.\mathit{Adj}[u]$}
        \If{$\mathit{comp}[v] = 0$}
            \State \Call{Assign}{$H, v, k$}
        \EndIf
    \EndFor
\EndProcedure
\end{pseudo}

第 5 行要的「依 $f$ 由大到小」的順序不用另外排序：和第~22~章的拓樸排序一樣，每個
頂點在第 2 行的 DFS 完成時就把它插到一個串列的最前面，DFS 做完時，串列就是完成時間由大到小
的順序。 \proc{Assign} 是只做「發現」的 DFS，用 $\mathit{comp}[v] = 0$ 代表白色。 兩
次 DFS 和建立 $G^{\mathrm{T}}$ 各花 $\Theta(n + m)$，總時間是 $\Theta(n + m)$。

\subsection{為什麼對}

下面的 $d[u]$、$f[u]$ 都指第一次 DFS（在 $G$ 上）的時間戳記。 對一個頂點集合 $C$，令
$f(C) = \max\{f[u] : u \in C\}$，也就是 $C$ 裡最晚完成的時間。

\begin{lemma}
\label{lem:con-finish}
設 $C$、$C'$ 是 $G$ 的兩個不同的強連通元件，而且有一條邊 $(x, y) \in E$，$x \in C$、
$y \in C'$。 那麼 $f(C) > f(C')$。
\end{lemma}

\begin{proof}
令 $z$ 是 $C \cup C'$ 裡最先被發現的頂點。 在時間 $d[z]$，$C \cup C'$ 的頂點全是白
色。 分兩種情況。

\textbf{$z \in C$}。 $C$ 裡任何頂點都能從 $z$ 在 $C$ 裡面走到；$C'$ 的任何頂點 $v$
也能從 $z$ 走到：先在 $C$ 裡走到 $x$，經過 $(x, y)$，再在 $C'$ 裡走到 $v$。 這些路徑
只經過 $C \cup C'$ 的頂點，在時間 $d[z]$ 全是白色，由白路徑定理，$C \cup C'$ 的其他
頂點都是 $z$ 在 DFS 樹裡的子孫。 由括號定理，子孫比 $z$ 先完成，所以 $f(C) = f[z] >
f(C')$。

\textbf{$z \in C'$}。 同樣的理由，$C'$ 的其他頂點都是 $z$ 的子孫，$f(C') = f[z]$。
$z$ 走不到 $C$ 的任何頂點：如果走得到，加上邊 $(x, y)$ 和 $C'$ 內部的路，$C$ 和 $C'$
就互相可達了。 在 $d[z]$ 到 $f[z]$ 這段時間裡被發現的頂點都是 $z$ 的子孫（括號定理），
也都從 $z$ 走得到，所以 $C$ 的頂點在時間 $f[z]$ 仍然是白色，之後才被發現、完成。 因此
$C$ 的每個頂點的完成時間都大於 $f[z]$，$f(C) > f[z] = f(C')$。
\end{proof}

換句話說，元件圖裡每一條邊 $(C, C')$ 都從 $f$ 較大的元件指向 $f$ 較小的元件。 完成時
間最大的頂點所在的元件 $f$ 值最大，不可能有邊指進來，是一個源點。

\begin{theorem}
\label{thm:con-kosaraju}
虛擬碼~\ref{alg:con-kosaraju} 第二次 DFS 的每一棵樹（同一個 $k$ 的頂點）恰好是一個強
連通元件。
\end{theorem}

\begin{proof}
對樹的個數做歸納。 假設前 $k - 1$ 棵樹各是一個完整的強連通元件，看第 $k$ 棵樹，它的
根 $r$ 是目前 $\mathit{comp} = 0$ 的頂點裡 $f$ 最大的一個。 令 $C$ 是 $r$ 所在的元件。
由歸納假設，已經分好的頂點都屬於完整的元件，$C$ 不是其中之一，所以 $C$ 的頂點全都還
沒分。 於是 $f[r] = f(C)$，而且對每個還沒分的元件 $C''$ 都有 $f(C) \ge f(C'')$。

(1) $C$ 的頂點都在這棵樹裡。 $C$ 在 $G^{\mathrm{T}}$ 裡也是強連通的，開始 $\proc{Assign}
(G^{\mathrm{T}}, r, k)$ 時 $C$ 的頂點全是白色（$\mathit{comp} = 0$），由白路徑定理，它
們都會成為 $r$ 的子孫。

(2) 這棵樹沒有 $C$ 以外的頂點。 假設有，看第一個被加進來、不在 $C$ 裡的頂點 $v$：它
是從樹裡某個頂點 $u$ 沿著 $G^{\mathrm{T}}$ 的邊 $(u, v)$ 走到的，而 $u$ 在 $C$ 裡。
$G^{\mathrm{T}}$ 的邊 $(u, v)$ 就是 $G$ 的邊 $(v, u)$，從 $v$ 的元件 $C_v$ 指向 $C$，
由引理~\ref{lem:con-finish}，$f(C_v) > f(C)$。 但 $v$ 的 $\mathit{comp}$ 是 $0$，
$C_v$ 還沒分，這和 $f(C)$ 是還沒分的元件裡最大的矛盾。
\end{proof}

證明的 (2) 也說明了第二次 DFS 為什麼一定要在 $G^{\mathrm{T}}$ 上做。 如果還在 $G$
上，從源點元件出發會順著邊走進下游的元件，一棵樹就吃下好幾個元件。

再看元件產生的順序。 第 $k$ 棵樹的根是當時 $f$ 最大的頂點，所以 $f(C_1) > f(C_2) >
\cdots$。 由引理~\ref{lem:con-finish}，元件圖的邊 $(C_i, C_j)$ 一定滿足 $f(C_i) >
f(C_j)$，也就是 $i < j$。 所以 \textbf{Kosaraju 演算法依元件圖的拓樸順序產生元件}，
第~\ref{sec:con-dag}~節會用到這一點。

\subsection{手算追蹤}

在圖~\ref{fig:con-roads} 的路網上跑一次。 第一次 DFS 在 $G$ 上，從 $1$ 開始，鄰居由小到大：
$1 \to 2$，$2$ 的鄰居 $1$ 是灰色（後向邊），接著 $2 \to 6 \to 9 \to 5$；$5 \to 6$ 是
後向邊，$5$、$9$、$6$、$2$ 依序完成。 回到 $1$，走 $1 \to 4$：$4 \to 2$ 時 $2$ 已經是
黑色，而且 $d[2] < d[4]$，是交叉邊；$4 \to 3$，$3 \to 5$ 又是交叉邊，$3 \to 8$，$8
\to 3$ 是後向邊。 最後主迴圈走到 $7$，它的兩條邊都指向黑色的頂點，是交叉邊。

\begin{center}
\begin{tikzpicture}[x=1.15cm, y=1.15cm]
\node[con-v] (v1) at (1.5,1.8) {1};
\node[con-v] (v2) at (3.3,1.8) {2};
\node[con-v] (v3) at (2.6,-0.9) {3};
\node[con-v] (v4) at (2.4,0.45) {4};
\node[con-v] (v5) at (5.1,-0.3) {5};
\node[con-v] (v6) at (5.1,1.8) {6};
\node[con-v] (v7) at (0,0.9) {7};
\node[con-v] (v8) at (0.9,-0.9) {8};
\node[con-v] (v9) at (6.5,0.75) {9};
\node[con-lab, above=0.5mm of v1] {1/16};
\node[con-lab, above=0.5mm of v2] {2/9};
\node[con-lab, below=0.5mm of v3] {11/14};
\node[con-lab, right=0.5mm of v4] {10/15};
\node[con-lab, below=0.5mm of v5] {5/6};
\node[con-lab, above=0.5mm of v6] {3/8};
\node[con-lab, above=0.5mm of v7] {17/18};
\node[con-lab, below=0.5mm of v8] {12/13};
\node[con-lab, right=0.5mm of v9] {4/7};
\draw[con-arc, con-tree] (v1) to[bend left=14] (v2);
\draw[con-arc, con-nt] (v2) to[bend left=14] node[con-elab, below] {B} (v1);
\draw[con-arc, con-tree] (v1) -- (v4);
\draw[con-arc, con-tree] (v2) -- (v6);
\draw[con-arc, con-nt] (v3) -- node[con-elab, pos=0.55] {C} (v5);
\draw[con-arc, con-tree] (v3) to[bend left=14] (v8);
\draw[con-arc, con-nt] (v8) to[bend left=14] node[con-elab, above] {B} (v3);
\draw[con-arc, con-nt] (v4) -- node[con-elab] {C} (v2);
\draw[con-arc, con-tree] (v4) -- (v3);
\draw[con-arc, con-nt] (v5) -- node[con-elab] {B} (v6);
\draw[con-arc, con-tree] (v6) -- (v9);
\draw[con-arc, con-tree] (v9) -- (v5);
\draw[con-arc, con-nt] (v7) -- node[con-elab] {C} (v1);
\draw[con-arc, con-nt] (v7) -- node[con-elab] {C} (v8);
\end{tikzpicture}
\diagcaption{第一次 DFS（在 $G$ 上）。 每個頂點旁邊是 $d/f$，粗線是樹邊，虛線是其他
邊：B 是後向邊，C 是交叉邊。 這次沒有前向邊。}
\label{fig:con-dfs1}
\end{center}

依 $f$ 由大到小排，順序是
\[
7\,(18),\ 1\,(16),\ 4\,(15),\ 3\,(14),\ 8\,(13),\ 2\,(9),\ 6\,(8),\ 9\,(7),\ 5\,(6).
\]
第二次 DFS 在 $G^{\mathrm{T}}$ 上做。 $G^{\mathrm{T}}$ 的鄰接串列（一樣由小到大）是
\[
\begin{array}{lllll}
1: 2, 7 & 2: 1, 4 & 3: 4, 8 & 4: 1 & 5: 3, 9 \\
6: 2, 5 & 7: \text{（空）} & 8: 3, 7 & 9: 6 &
\end{array}
\]
照上面的順序挑起點：

\begin{center}
\small
\begin{tabular}{clc}
\toprule
起點 & 這棵樹的走法（在 $G^{\mathrm{T}}$ 上） & 元件 \\
\midrule
$7$ & $7$ 沒有出邊 & $C_1 = \{7\}$ \\
$1$ & $1 \to 2$，$2 \to 4$（$2$ 的 $1$ 已走過；$4$ 的 $1$ 已走過）；$1$ 的 $7$ 已屬
於 $C_1$ & $C_2 = \{1, 2, 4\}$ \\
$4$ & 已分好，跳過 & \\
$3$ & $3$ 的 $4$ 已屬於 $C_2$；$3 \to 8$，$8$ 的 $3$ 走過、$7$ 屬於 $C_1$ & $C_3 =
\{3, 8\}$ \\
$8$、$2$ & 已分好，跳過 & \\
$6$ & $6$ 的 $2$ 屬於 $C_2$；$6 \to 5$，$5$ 的 $3$ 屬於 $C_3$；$5 \to 9$，$9$ 的 $6$
走過 & $C_4 = \{5, 6, 9\}$ \\
$9$、$5$ & 已分好，跳過 & \\
\bottomrule
\end{tabular}
\end{center}

第二次 DFS 碰到的「已分好」的頂點，全都屬於\textbf{之前}產生的元件，這正是定理的
(2)。 例如從 $3$ 出發時看到 $4$：$G^{\mathrm{T}}$ 的邊 $3 \to 4$ 是 $G$ 的邊 $4 \to
3$，從 $C_2$ 指向 $C_3$，$C_2$ 早就分好了，不會被吃進來。 得到的四個元件和
圖~\ref{fig:con-roads} 一致，產生的順序 $C_1, C_2, C_3, C_4$ 也是元件圖的一個拓樸順
序。

\begin{examnote}
Kosaraju 的追蹤題通常要寫三樣東西：第一次 DFS 的 $d/f$（或只寫 $f$）、依 $f$ 由大到
小的順序、第二次 DFS 的每一棵樹。 兩次 DFS 的鄰居順序和主迴圈的順序都會影響中間的數
字（元件本身不受影響），照題目給的鄰接串列走；題目沒給時寫明「鄰居由小到大」。 第
二次 DFS 是在 $G^{\mathrm{T}}$ 上照 $f$ \textbf{由大到小}挑起點。 改在 $G$ 上做，
不管照 $f$ 由大到小還是由小到大挑起點，都可能把兩個以上的元件併成一棵樹（觀念
題第~\ref{q:con-choice} 題）。
\end{examnote}

\section{Tarjan 演算法}
\label{sec:con-tarjan}

Kosaraju 演算法要做兩次 DFS，還要多存一張轉置圖。 Tarjan 演算法只做一次 DFS，邊走
邊把元件切出來。 它靠兩個觀察。

第一，每個強連通元件 $C$ 裡最先被發現的頂點，叫它 $C$ 的\textbf{頭}（head）。 發現頭
的時候，$C$ 的其他頂點都是白色，而且都能從頭在 $C$ 裡面走到，由白路徑定理，它們全是
頭的子孫。 所以每個元件在 DFS 森林裡都「掛」在它的頭底下。

第二，DFS 發現一個頂點，就把它推進一個堆疊。 一個頭 $u$ 完成的時候，比 $u$ 晚進堆疊、
還留在堆疊上的頂點都是 $u$ 的子孫；如果在它之前完成的頭，都已經把自己的元件從堆疊上拿
走了，那麼 $u$ 上面剩下的就剛好是 $u$ 的元件。 所以只要能在 $u$ 完成時判斷「$u$ 是不
是頭」，是頭就把堆疊彈到 $u$ 為止，彈出來的就是一個元件。

判斷頭的工具是 $\mathit{low}$ 值。

\subsection{編號和 low 值}

本章的 Tarjan 演算法和第~\ref{sec:con-cut}~節的割點演算法裡，$d[u]$ 只在\textbf{發
現}時加一，等於頂點被發現的順序 $1, 2, \ldots, n$，完成時不加。 這和第~21~章同時記
$d$ 和 $f$ 的時間戳記數字不同，但兩種編號的大小順序完全一樣，而下面所有的判斷都只比
大小，所以結果相同。 用發現順序是因為 $\mathit{low}$ 的手算表比較好填。

\begin{definition}
在 Tarjan 演算法裡，頂點 $u$ 的 $\mathit{low}[u]$ 是下面這些數的最小值：
\begin{itemize}
\item $d[u]$；
\item $\mathit{low}[w]$，$w$ 是 $u$ 在 DFS 樹裡的子節點；
\item $d[w]$，$(u, w)$ 是一條邊，而且檢查這條邊的時候 $w$ 已經被發現、\textbf{還在
堆疊裡}。
\end{itemize}
\end{definition}

把子節點的部分一層一層展開，$\mathit{low}[u]$ 就是「$u$ 自己的 $d$，以及從 $u$ 的子
樹裡任何頂點出發、一條邊走到堆疊上的頂點 $w$ 時，那些 $d[w]$」的最小值。 直覺上，
$\mathit{low}[u] < d[u]$ 表示 $u$ 的子樹裡有路能回到比 $u$ 更早發現、而且和 $u$ 在同
一個元件的頂點，$u$ 就不會是頭。

\begin{pseudo}[label=alg:con-tarjan]{Tarjan 強連通元件演算法}
\Procedure{SCC-Tarjan}{$G$}
    \State $\mathit{time} \gets 0$；$k \gets 0$；$S \gets$ 空堆疊
    \State 把每個 $d[u]$ 設成 $0$、$\mathit{onstack}[u]$ 設成 \FALSE
    \For{每個頂點 $u$（由小到大）}
        \If{$d[u] = 0$}
            \State \Call{Tarjan-Visit}{$G, u$}
        \EndIf
    \EndFor
\EndProcedure
\Procedure{Tarjan-Visit}{$G, u$}
    \State $\mathit{time} \gets \mathit{time} + 1$；$d[u] \gets \mathit{time}$；$\mathit{low}[u] \gets \mathit{time}$
    \State \Call{Push}{$S, u$}；$\mathit{onstack}[u] \gets \TRUE$
    \For{每個 $w \in \mathit{Adj}[u]$}
        \If{$d[w] = 0$} \Comment{樹邊}
            \State \Call{Tarjan-Visit}{$G, w$}
            \State $\mathit{low}[u] \gets \min(\mathit{low}[u], \mathit{low}[w])$
        \ElsIf{$\mathit{onstack}[w]$} \Comment{$w$ 還沒分進任何元件}
            \State $\mathit{low}[u] \gets \min(\mathit{low}[u], d[w])$
        \EndIf
    \EndFor
    \If{$\mathit{low}[u] = d[u]$} \Comment{$u$ 是頭}
        \State $k \gets k + 1$
        \Repeat
            \State $x \gets$ \Call{Pop}{$S$}；$\mathit{onstack}[x] \gets \FALSE$；$\mathit{comp}[x] \gets k$
        \Until{$x = u$}
    \EndIf
\EndProcedure
\end{pseudo}

$d[w] = 0$ 就是白色。 第 14 行的條件是整個演算法最容易寫錯的地方：要問的是
\textbf{$w$ 還在不在堆疊上}，不是 $w$ 是不是灰色。

\begin{itemize}
\item $w$ 已經不在堆疊上，表示它的元件已經切出去了。 $w$ 走不回 $u$（否則它們會在同一
個元件），這條邊對 $u$ 的元件沒有影響，要跳過。 這種邊一定是交叉邊，漏了這個檢查，
$\mathit{low}[u]$ 會被一個不相干的小 $d$ 值拉低，$u$ 明明是頭卻判斷成不是。
\item $w$ 還在堆疊上，但可能已經是黑色：它完成了，只是它的頭還沒完成。 這時邊 $(u, w)$
是交叉邊或前向邊，一樣要算進去。 只看灰色的話會漏掉這種情況。
\end{itemize}

每個頂點進出堆疊各一次，每條邊在第 10 行檢查一次，總時間是 $\Theta(n + m)$。 C 程式放
在第~\ref{sec:con-dag}~節，和縮點一起寫。

\subsection{手算追蹤}

一樣用圖~\ref{fig:con-roads} 的路網。 下表一列是一個動作；$d$、$\mathit{low}$ 寫在
「改變」欄，堆疊由左（底）往右（頂）寫，是做完這一列之後的內容。

\begin{center}
\footnotesize
\setlength{\tabcolsep}{3.5pt}
\begin{tabular}{r>{\raggedright\arraybackslash}p{0.27\linewidth}>{\raggedright\arraybackslash}p{0.43\linewidth}lc}
\toprule
步 & 動作 & 改變 & 堆疊 & 輸出 \\
\midrule
1 & 主迴圈發現 $1$ & $d[1]=\mathit{low}[1]=1$ & 1 &  \\
2 & 樹邊 $1\to2$，發現 $2$ & $d[2]=\mathit{low}[2]=2$ & 1 2 &  \\
3 & $2\to1$，$1$ 在堆疊裡 & $\mathit{low}[2]=\min(2,d[1])=1$ & 1 2 &  \\
4 & 樹邊 $2\to6$，發現 $6$ & $d[6]=\mathit{low}[6]=3$ & 1 2 6 &  \\
5 & 樹邊 $6\to9$，發現 $9$ & $d[9]=\mathit{low}[9]=4$ & 1 2 6 9 &  \\
6 & 樹邊 $9\to5$，發現 $5$ & $d[5]=\mathit{low}[5]=5$ & 1 2 6 9 5 &  \\
7 & $5\to6$，$6$ 在堆疊裡；$5$ 完成 & $\mathit{low}[5]=\min(5,d[6])=3$；留在堆疊 & 1 2 6 9 5 &  \\
8 & 回到 $9$；$9$ 完成 & $\mathit{low}[9]=\min(4,\mathit{low}[5])=3$；留在堆疊 & 1 2 6 9 5 &  \\
9 & 回到 $6$；$6$ 完成 & $\mathit{low}[6]=\min(3,\mathit{low}[9])=3=d[6]$，彈出 & 1 2 & $\{5,9,6\}$ \\
10 & 回到 $2$；$2$ 完成 & $\mathit{low}[2]=\min(1,\mathit{low}[6])=1$；留在堆疊 & 1 2 &  \\
11 & 回到 $1$ & $\mathit{low}[1]=\min(1,\mathit{low}[2])=1$ & 1 2 &  \\
12 & 樹邊 $1\to4$，發現 $4$ & $d[4]=\mathit{low}[4]=6$ & 1 2 4 &  \\
13 & $4\to2$，$2$ 在堆疊裡 & $\mathit{low}[4]=\min(6,d[2])=2$ & 1 2 4 &  \\
14 & 樹邊 $4\to3$，發現 $3$ & $d[3]=\mathit{low}[3]=7$ & 1 2 4 3 &  \\
15 & $3\to5$，$5$ 不在堆疊 & 略過 & 1 2 4 3 &  \\
16 & 樹邊 $3\to8$，發現 $8$ & $d[8]=\mathit{low}[8]=8$ & 1 2 4 3 8 &  \\
17 & $8\to3$，$3$ 在堆疊裡；$8$ 完成 & $\mathit{low}[8]=\min(8,d[3])=7$；留在堆疊 & 1 2 4 3 8 &  \\
18 & 回到 $3$；$3$ 完成 & $\mathit{low}[3]=\min(7,\mathit{low}[8])=7=d[3]$，彈出 & 1 2 4 & $\{8,3\}$ \\
19 & 回到 $4$；$4$ 完成 & $\mathit{low}[4]=\min(2,\mathit{low}[3])=2$；留在堆疊 & 1 2 4 &  \\
20 & 回到 $1$；$1$ 完成 & $\mathit{low}[1]=\min(1,\mathit{low}[4])=1=d[1]$，彈出 & （空） & $\{4,2,1\}$ \\
21 & 主迴圈發現 $7$ & $d[7]=\mathit{low}[7]=9$ & 7 &  \\
22 & $7\to1$，$1$ 不在堆疊 & 略過 & 7 &  \\
23 & $7\to8$，$8$ 不在堆疊；$7$ 完成 & $\mathit{low}[7]=d[7]$，彈出 & （空） & $\{7\}$ \\
\bottomrule
\end{tabular}
\end{center}

最後的值整理成一列，四個頭 $6, 3, 1, 7$ 正是 $\mathit{low} = d$ 的頂點：
\[
\begin{array}{c|ccccccccc}
v & 1 & 2 & 3 & 4 & 5 & 6 & 7 & 8 & 9 \\
\hline
d[v] & 1 & 2 & 7 & 6 & 5 & 3 & 9 & 8 & 4 \\
\mathit{low}[v] & 1 & 1 & 7 & 2 & 3 & 3 & 9 & 7 & 3
\end{array}
\]

幾個要注意的地方：

\begin{itemize}
\item 第 13 步：$2$ 在第 10 步就完成了，是黑色，但它的頭 $1$ 還沒完成，所以 $2$ 還在
堆疊上，$4 \to 2$ 要算。 這一步讓 $\mathit{low}[4] = 2 < 6$，$4$ 不是頭，最後和 $1$、
$2$ 一起被彈出。
\item 第 15 步：$3 \to 5$ 指向一個已經在第 9 步切出去的頂點。 如果沒有檢查
$\mathit{onstack}$，$\mathit{low}[3]$ 會變成 $d[5] = 5$，第 18 步就不會彈出，$3$、$8$
會被錯誤地併進 $\{1, 2, 4\}$。
\item 元件輸出的順序是 $\{5,9,6\}$、$\{8,3\}$、$\{4,2,1\}$、$\{7\}$，正好是
Kosaraju 順序 $C_1, C_2, C_3, C_4$ 倒過來，第~\ref{sec:con-dag}~節會證明這不是巧合。
\item $\mathit{low}[u]$ 不一定等於 $u$ 所在元件的頭的 $d$ 值：$\mathit{low}[4] = 2$，
但 $4$ 的頭 $1$ 的 $d$ 是 $1$。 $\mathit{low}$ 只用來判斷「是不是頭」，不要拿來當元件編號。
\end{itemize}

\subsection{為什麼對}

\begin{theorem}
\label{thm:con-tarjan}
在虛擬碼~\ref{alg:con-tarjan} 裡，頂點 $u$ 完成時（第 16 行），$\mathit{low}[u] =
d[u]$ 若且唯若 $u$ 是它所在元件的頭；而且這時第 17--20 行彈出的頂點恰好是 $u$ 的元件。
\end{theorem}

\begin{whybox}[定理~\ref{thm:con-tarjan}]
令 $T(u)$ 是 $u$ 在 DFS 森林裡的子樹（$u$ 和它所有的子孫），$h(x)$ 是 $x$ 所在元件的
頭。 對頂點完成的先後做歸納：假設比 $u$ 早完成的頂點都符合定理。 這個假設推得出一個
性質：在 $u$ 完成之前的任何時刻，一個已發現的頂點 $x$ 在堆疊上，若且唯若 $h(x)$ 還
沒完成。（$x$ 一被發現就進堆疊；由歸納假設，只有頭完成時會彈出，而且彈出的恰好是那
個頭的元件。）

另外，把第 13 行的子節點一層一層展開，$u$ 完成時
\[
\mathit{low}[u] = \min\bigl(\{d[u]\} \cup \{d[w] : x \in T(u)，\text{第 15 行處理邊 } (x, w) \text{ 時 } w \text{ 在堆疊上}\}\bigr).
\]

\textbf{$u$ 是頭 $\Rightarrow$ $\mathit{low}[u] = d[u]$。} 看上式裡任何一個 $d[w]$。 處
理邊 $(x, w)$ 時 $w$ 在堆疊上，所以 $h(w)$ 還沒完成。 假設 $d[w] < d[u]$，那麼 $d[h(w)]
\le d[w] < d[u]$：$h(w)$ 比 $u$ 早發現，而這時 DFS 還在 $u$ 的子樹裡，$h(w)$ 沒有完成，
一定是灰色，也就是 $x$ 的祖先，而且是 $u$ 的真祖先。 於是 $h(w) \leadsto u$（樹上的
路）、$u \leadsto x \to w$、$w \leadsto h(w)$（同一個元件），$u$ 和 $h(w)$ 互相可達、
在同一個元件，$h(w)$ 又比 $u$ 早發現，和「$u$ 是頭」矛盾。 所以每個 $d[w] \ge d[u]$，
$\mathit{low}[u] = d[u]$。

\textbf{$u$ 不是頭 $\Rightarrow$ $\mathit{low}[u] < d[u]$。} 令 $h = h(u)$，它比 $u$ 早
發現，而且由白路徑定理，$u$ 是 $h$ 的子孫，所以 $h$ 是 $u$ 的真祖先，$h \notin T(u)$。
取一條從 $u$ 到 $h$ 的路徑，路徑上的頂點都和 $u$ 在同一個元件（$u$ 能到它們，它們能到
$h$）。 路徑從 $T(u)$ 裡出發，終點在 $T(u)$ 外面，所以有第一條邊 $(x, w)$ 滿足 $x \in
T(u)$、$w \notin T(u)$。 $x$ 處理這條邊時，$w$ 不可能是白色（否則 $w$ 會變成 $x$ 的子
孫，落在 $T(u)$ 裡）。 在 $u$ 的子樹還沒做完時被發現的頂點都在 $T(u)$ 裡，所以 $w$ 比
$u$ 早發現，$d[w] < d[u]$。 又 $h(w) = h$ 是 $u$ 的祖先，還沒完成，所以 $w$ 在堆疊上，
第 15 行會執行，$\mathit{low}[u] \le d[w] < d[u]$。

\textbf{彈出的恰好是 $u$ 的元件。} 設 $u$ 是頭。 堆疊上 $u$ 以上的頂點都比 $u$ 晚發現、
在 $u$ 完成前就發現了，所以都是 $u$ 的子孫，而且它們的頭都還沒完成。 如果其中某個
$y$ 的頭 $h(y) \ne u$：$h(y)$ 是 $y$ 的祖先，要嘛是 $u$ 的真子孫（那它已經完成，$y$
應該已經被彈出），要嘛是 $u$ 的真祖先（那 $h(y) \leadsto u \leadsto y \leadsto h(y)$，
$u$ 和 $h(y)$ 同一個元件，$u$ 就不是頭），兩種都不可能。 所以 $u$ 上面的頂點都屬於
$u$ 的元件。 反過來，$u$ 的元件裡每個頂點都是 $u$ 的子孫，已經發現，頭 $u$ 還沒完成，
所以都還在堆疊上。 彈到 $u$ 為止，拿到的就是整個元件。
\end{whybox}


\section{縮點成 DAG}
\label{sec:con-dag}

找出強連通元件以後，常見的下一步是把每個元件縮成一個點，在元件圖 $G^{\mathrm{SCC}}$
上繼續做。 引理~\ref{lem:con-dag} 保證元件圖沒有環，第~22~章的拓樸排序、DAG 上的動
態規劃都能直接用。 這一步叫\textbf{縮點}（condensation）。

\begin{center}
\begin{tikzpicture}[x=1.25cm, y=1.1cm]
\node[con-cv] (c1) at (0,0) {$C_1 = \{7\}$};
\node[con-cv] (c2) at (2.6,0.95) {$C_2 = \{1, 2, 4\}$};
\node[con-cv] (c3) at (2.6,-0.95) {$C_3 = \{3, 8\}$};
\node[con-cv] (c4) at (5.4,0) {$C_4 = \{5, 6, 9\}$};
\draw[con-arc] (c1) -- node[con-elab, above left] {$7 \to 1$} (c2);
\draw[con-arc] (c1) -- node[con-elab, below left] {$7 \to 8$} (c3);
\draw[con-arc] (c2) -- node[con-elab, right] {$4 \to 3$} (c3);
\draw[con-arc] (c2) -- node[con-elab, above right] {$2 \to 6$} (c4);
\draw[con-arc] (c3) -- node[con-elab, below right] {$3 \to 5$} (c4);
\end{tikzpicture}
\diagcaption{圖~\ref{fig:con-roads} 的元件圖。 每條邊旁邊是原圖裡造成這條邊的一條路；
$C_1, C_2, C_3, C_4$ 是一個拓樸順序。}
\label{fig:con-dag}
\end{center}

\subsection{建立元件圖}

有了每個頂點的元件編號 $\mathit{comp}[u]$，掃過原圖的每一條邊 $(u, v)$：
$\mathit{comp}[u] \ne \mathit{comp}[v]$ 時，元件圖要有 $(\mathit{comp}[u],
\mathit{comp}[v])$ 這條邊。 同一對元件之間可能有好幾條原圖的邊，例如 $C_2$ 到 $C_4$ 只
有 $2 \to 6$ 一條，但如果再加一條 $4 \to 5$，就會出現第二次。 要去掉重複，可以照元件
一個一個處理：處理元件 $c$ 的所有出邊時，用一個陣列 $\mathit{mark}$ 記下「$c$ 已經連到
哪些元件」，$\mathit{mark}[c'] = c$ 表示這一輪已經加過 $(c, c')$。 每條原圖的邊只看一次，
總時間 $\Theta(n + m)$，$\mathit{mark}$ 也不用每輪清空。

\subsection{元件產生的順序}

兩個演算法產生元件的順序剛好相反。 第~\ref{sec:con-kosaraju}~節已經證明 Kosaraju 演算
法依拓樸順序產生元件。 Tarjan 演算法則是\textbf{反向的拓樸順序}：元件圖裡的匯點最先
被彈出。

\begin{theorem}
\label{thm:con-tarjan-order}
若元件圖有邊 $(C, C')$，Tarjan 演算法先輸出 $C'$，後輸出 $C$。
\end{theorem}

\begin{proof}
Tarjan 演算法在元件的頭完成時輸出這個元件。 元件的其他頂點都是頭的子孫，比頭先完
成。 用第~21~章的時間戳記記下這次 DFS 的完成時間 $f$（Tarjan 演算法沒有存 $f$，但
DFS 的走法已經決定了它），頭的完成時間就是 $f(C)$，所以元件依 $f(C)$ 由小到大輸出。
引理~\ref{lem:con-finish} 對任何一次 DFS 都成立，邊 $(C, C')$ 表示 $f(C) > f(C')$，所
以 $C'$ 先輸出。
\end{proof}

換句話說，Tarjan 演算法的輸出倒過來就是元件圖的拓樸順序。 第~22~章用 DFS 做拓樸排序時
是「完成的先後倒過來」；在 DAG 上跑 Tarjan 演算法，每個元件都只有一個頂點，輸出順序就
是完成的順序，兩件事是一致的。

\subsection{應用：從哪裡出發可以到每個路口}

縮點最常見的用途，是把「有環的圖」上的問題換成「DAG」上的問題。 下面是一個例子。

\begin{theorem}
\label{thm:con-source}
有向圖裡存在一個頂點可以到達所有頂點，若且唯若元件圖恰好有一個源點（入度為 $0$ 的
元件）。 這時能到達所有頂點的，正好是那個源點元件裡的頂點。 一般來說，至少要選幾個
起點才能到達所有頂點，答案是元件圖的源點個數。
\end{theorem}

\begin{proof}
元件圖是 DAG，至少有一個源點。 在 DAG 裡從任何一個點出發、一直逆著邊往回走，因為沒有
環，最後一定停在某個源點；所以每個元件都能從某個源點元件到達。 源點元件沒有邊進來，
外面的頂點都到不了它，所以每個源點元件裡至少要選一個起點；每個源點元件各選一個頂點
又已經足夠（元件內部互相可達），最少的起點數就是源點個數。 只有一個源點 $S$ 時，$S$
到得了每個元件，$S$ 的每個頂點都能到達所有頂點；其他元件的頂點到不了 $S$。 有兩個以
上的源點時，任何頂點都到不了「另一個」源點元件，所以沒有頂點能到達所有頂點。
\end{proof}

章首的路網只有一個源點 $C_1 = \{7\}$，所以從路口 $7$ 出發可以開到每個路口，而且只有
$7$ 可以。 交通局要讓每個路口都能互相到達，至少得再加幾條路，這也是在元件圖上算源點和
匯點的問題（結論是 $\max(\text{源點數}, \text{匯點數})$，元件多於一個時成立），證明比
較長，有興趣可以看 Eswaran 和 Tarjan 的論文（延伸閱讀）。

其他常見的用法：

\begin{itemize}
\item 頂點有權重（例如每個路口有一間店，可以收集點數），問一條路徑最多收集多少。 原
圖有環，同一個元件裡的頂點可以全部走到再離開，所以把元件的權重加總，在元件圖上做
「DAG 最長路徑」的動態規劃（第~18~章的想法、第~22~章的拓樸順序）。
\item 判斷可達性：同一個元件裡的頂點能到達的頂點集合完全相同，只要對每個元件算一次。
\item 2-SAT（每個子句只有兩個變數的布林式能不能滿足）：把每個變數的真、假各當一個頂
點、每個子句換成兩條邊，布林式可以滿足若且唯若沒有變數的真、假落在同一個強連通元件
裡（Aspvall、Plass、Tarjan 1979）。 這是延伸題，本書不展開。
\end{itemize}

\subsection{C 程式}

\begin{cbookrecap}[第~14~章結構、第~15~章遞迴、第~19~章鏈結串列]
鄰接串列的每個節點是一個 \texttt{Node}，存鄰居編號和指向下一個節點的指標，
\texttt{adj[u]} 是 \texttt{u} 的串列的開頭。 \texttt{add\_arc} 要把新節點插在「由小
到大」的位置，用的是本書第~5~章「依序插入」的寫法：\texttt{Node **p} 指向「要改的那
個指標」，一開始是 \texttt{\&g->adj[u]}，往後走時變成前一個節點的
\texttt{\&(*p)->next}，最後 \texttt{*p = x} 把新節點接上，插在開頭和插在中間寫成同一
段。

\texttt{visit} 是遞迴函式，每深一層就在呼叫堆疊上多一個框。 DFS 樹是一條很長的鏈時
（例如 $10^6$ 個頂點排成一列），遞迴會用光呼叫堆疊，要改用第~\ref{sec:gtr-dfs-iter}~節的寫法；本章的例子很小，直接用遞迴。
\end{cbookrecap}

下面的程式讀入有向圖，照虛擬碼~\ref{alg:con-tarjan} 找強連通元件，同時建立元件圖。 元
件依彈出的順序編號 $1, 2, \ldots$；每彈出一個元件，就印出它的頂點（彈出的順序），再
在 \texttt{->} 後面印出它有邊連到的元件。 由定理~\ref{thm:con-tarjan-order}，剛彈出的元件的
出邊只會連到自己或更早彈出的元件，這些頂點的 \texttt{comp} 都已經填好，所以可以當場建
邊；去掉重複用的就是前面說的 \texttt{mark}。 彈出的頂點雖然離開了堆疊，數值還留在
\texttt{stack[top..end-1]}，直接拿來掃。 最後印出每個頂點的 $d$、$\mathit{low}$，以及
元件圖的源點；源點只有一個時，再印出能到達所有頂點的起點。 演算法用到的變數都包在
\texttt{Tarjan} 結構裡，用指標傳給 \texttt{visit}。 輸入是圖~\ref{fig:con-roads} 的路
網：第一列是 $n$、$m$，接著 $m$ 列各是一條邊 $u\ v$。

\begin{example}[label=ex:con-tarjan]{c}{(scc\_tarjan.c) Tarjan 演算法找強連通元件並建立元件圖}
#include <stdio.h>
#include <stdlib.h>

typedef struct Node {
    int v;
    struct Node *next;
} Node;

typedef struct {
    int n;
    Node **adj;     // adj[u]：u 的出邊，終點由小到大
} Graph;

// 加入有向邊 u -> v，插在串列裡由小到大的位置
int add_arc(Graph *g, int u, int v)
{
    Node *x = malloc(sizeof *x);
    if (x == NULL) {
        return 0;
    }
    Node **p = &g->adj[u];
    while (*p != NULL && (*p)->v < v) {
        p = &(*p)->next;
    }
    x->v = v;
    x->next = *p;
    *p = x;
    return 1;
}

void free_graph(Graph *g)
{
    for (int u = 0; g->adj != NULL && u < g->n; u++) {
        Node *x = g->adj[u];
        while (x != NULL) {
            Node *next = x->next;
            free(x);
            x = next;
        }
    }
    free(g->adj);
}

typedef struct {
    const Graph *g;
    int *d, *low;       // d[u] = 0 表示還沒發現
    int *on_stack;
    int *stack, top;
    int time, count;    // count：已經找到幾個 SCC
    int *comp;          // comp[v]：v 屬於第幾個 SCC，0 表示還沒分
    int *mark, *indeg;  // 建元件圖用，以元件編號當索引
} Tarjan;

void visit(Tarjan *t, int u)
{
    t->d[u] = t->low[u] = ++t->time;
    t->stack[t->top++] = u;
    t->on_stack[u] = 1;
    for (Node *x = t->g->adj[u]; x != NULL; x = x->next) {
        int w = x->v;
        if (t->d[w] == 0) {                     // 樹邊
            visit(t, w);
            if (t->low[w] < t->low[u]) {
                t->low[u] = t->low[w];
            }
        } else if (t->on_stack[w] && t->d[w] < t->low[u]) {
            t->low[u] = t->d[w];                // w 還沒分進任何 SCC
        }
    }
    if (t->low[u] == t->d[u]) {                 // u 是這個 SCC 的頭
        int c = ++t->count, end = t->top, w;
        printf("SCC %d:", c);
        do {
            w = t->stack[--t->top];
            t->on_stack[w] = 0;
            t->comp[w] = c;
            printf(" %d", w + 1);
        } while (w != u);
        printf(" ->");
        // 剛彈出的頂點還留在 stack[top..end-1]。 它們的出邊只會連到
        // 自己的元件或更早彈出的元件，所以 comp 都已經填好了。
        for (int i = t->top; i < end; i++) {
            for (Node *x = t->g->adj[t->stack[i]]; x != NULL;
                 x = x->next) {
                int c2 = t->comp[x->v];
                if (c2 != c && t->mark[c2] != c) {  // 每對只記一次
                    t->mark[c2] = c;
                    t->indeg[c2]++;
                    printf(" %d", c2);
                }
            }
        }
        printf("\n");
    }
}

int main(void)
{
    int n, m;
    if (scanf("%d %d", &n, &m) != 2 || n <= 0 || m < 0) {
        return 1;
    }
    Graph g = {n, calloc(n, sizeof(Node *))};
    Tarjan t = {&g, calloc(n, sizeof(int)), calloc(n, sizeof(int)),
                calloc(n, sizeof(int)), malloc(n * sizeof(int)),
                0, 0, 0, calloc(n, sizeof(int)),
                calloc(n + 1, sizeof(int)),
                calloc(n + 1, sizeof(int))};
    int ok = g.adj != NULL && t.d != NULL && t.low != NULL
             && t.on_stack != NULL && t.stack != NULL
             && t.comp != NULL && t.mark != NULL && t.indeg != NULL;
    for (int i = 0; ok && i < m; i++) {
        int u, v;
        ok = scanf("%d %d", &u, &v) == 2 && u >= 1 && u <= n
             && v >= 1 && v <= n && add_arc(&g, u - 1, v - 1);
    }
    if (ok) {
        for (int u = 0; u < n; u++) {
            if (t.d[u] == 0) {
                visit(&t, u);
            }
        }
        printf("v:");
        for (int u = 0; u < n; u++) {
            printf(" %d", u + 1);
        }
        printf("\nd:");
        for (int u = 0; u < n; u++) {
            printf(" %d", t.d[u]);
        }
        printf("\nlow:");
        for (int u = 0; u < n; u++) {
            printf(" %d", t.low[u]);
        }
        printf("\nsources:");
        int sources = 0, src = 0;
        for (int c = 1; c <= t.count; c++) {
            if (t.indeg[c] == 0) {
                printf(" %d", c);
                sources++;
                src = c;
            }
        }
        printf("\n");
        if (sources == 1) {
            printf("reach every vertex from:");
            for (int v = 0; v < n; v++) {
                if (t.comp[v] == src) {
                    printf(" %d", v + 1);
                }
            }
            printf("\n");
        }
    }
    free_graph(&g);
    free(t.d);
    free(t.low);
    free(t.on_stack);
    free(t.stack);
    free(t.comp);
    free(t.mark);
    free(t.indeg);
    return ok ? 0 : 1;
}
\end{example}
\iosample{9 14 \\ 1 2 \\ 1 4 \\ 2 1 \\ 2 6 \\ 3 5 \\ 3 8 \\ 4 2 \\ 4 3 \\ 5 6 \\ 6 9 \\ 7 1 \\ 7 8 \\ 8 3 \\ 9 5}{SCC 1: 5 9 6 -> \\ SCC 2: 8 3 -> 1 \\ SCC 3: 4 2 1 -> 1 2 \\ SCC 4: 7 -> 3 2 \\ v: 1 2 3 4 5 6 7 8 9 \\ d: 1 2 7 6 5 3 9 8 4 \\ low: 1 1 7 2 3 3 9 7 3 \\ sources: 4 \\ reach every vertex from: 7}

元件、$d$、$\mathit{low}$ 都和手算表一致。 程式的元件編號是彈出的順序，和
圖~\ref{fig:con-dag} 的拓樸順序相反：程式的 $1, 2, 3, 4$ 依序是 $C_4, C_3, C_2, C_1$，元件
圖的每條邊都從大編號指到小編號。 元件 $3$（$\{1, 2, 4\}$）印出 \texttt{-> 1 2}，因為先掃
到頂點 $4$ 的 $4 \to 3$，再掃到 $2 \to 6$。 唯一的源點是元件 $4$，也就是 $\{7\}$，和
定理~\ref{thm:con-source} 的結論相同。 \texttt{mark}、\texttt{indeg} 用元件編號當索
引，編號最大是 $n$，所以開 $n + 1$ 格。 建鄰接串列時每次插入要走到正確的位置，最壞花
$O(\text{度數})$；這只是為了讓 DFS 的順序和手算相同，找元件和建元件圖本身是 $\Theta(n +
m)$。

\section{割點與橋}
\label{sec:con-cut}

接下來兩節講無向圖。 本章的無向圖都假設是\textbf{簡單圖}：沒有自環，兩個頂點之間最
多一條邊。 有平行邊時哪裡要改，在本節最後說明。

\begin{definition}
無向圖 $G$ 裡，如果拿掉頂點 $u$（連同所有接到 $u$ 的邊）之後，連通元件的個數變多，就
說 $u$ 是\textbf{割點}（cut vertex，也叫關節點 articulation point）。 如果拿掉邊 $e$
之後連通元件的個數變多，就說 $e$ 是\textbf{橋}（bridge）。
\end{definition}

定義用「元件個數變多」，所以原本就不連通的圖也適用：割點、橋只會影響它所在的那個連
通元件。 孤立的頂點拿掉後元件個數變少，不是割點。

\begin{center}
\begin{tikzpicture}[x=1.05cm, y=1.05cm]
\node[con-v] (u10) at (0,1.2) {10};
\node[con-v] (u1) at (1.4,1.2) {1};
\node[con-v] (u5) at (2.3,2.3) {5};
\node[con-v] (u3) at (2.6,0.5) {3};
\node[con-v] (u6) at (4.1,0.5) {6};
\node[con-v] (u4) at (5.0,1.7) {4};
\node[con-v] (u2) at (4.6,3.0) {2};
\node[con-v] (u8) at (5.9,0.5) {8};
\node[con-v] (u7) at (7.3,1.3) {7};
\node[con-v] (u9) at (7.3,-0.3) {9};
\draw[con-e] (u1) -- (u3);
\draw[con-e] (u1) -- (u5);
\draw[con-e] (u1) -- (u10);
\draw[con-e] (u2) -- (u4);
\draw[con-e] (u3) -- (u5);
\draw[con-e] (u3) -- (u6);
\draw[con-e] (u4) -- (u6);
\draw[con-e] (u4) -- (u8);
\draw[con-e] (u6) -- (u8);
\draw[con-e] (u7) -- (u8);
\draw[con-e] (u7) -- (u9);
\draw[con-e] (u8) -- (u9);
\end{tikzpicture}
\diagcaption{校園的大樓和走廊：$10$ 棟大樓、$12$ 條走廊。}
\label{fig:con-corridor}
\end{center}

圖~\ref{fig:con-corridor} 是章首的校園。 用眼睛看：封掉大樓 $3$，$\{1, 5, 10\}$ 和右邊
就分開了，$3$ 是割點；走廊 $3$--$6$ 是左右兩半之間唯一的一條，是橋。 大樓 $5$ 不是割
點，封掉它，$1$ 和 $3$ 之間還有走廊 $1$--$3$。 一個一個試的話，每個頂點拿掉後做一次
DFS，要 $O(n(n + m))$；每條邊試一次要 $O(m(n + m))$。 下面用一次 DFS 全部找出來。

\subsection{無向圖的 DFS 樹}

第~21~章的定理~\ref{thm:gtr-undirected} 證明過，無向圖的 DFS 只會產生\textbf{樹邊}和\textbf{後向邊}，沒有前向邊和
交叉邊。 換句話說，每一條邊的兩端，一定有一個是另一個在 DFS 樹裡的祖先。 這是本節所
有證明的基礎。 實作時有兩個細節：

\begin{itemize}
\item 無向邊在鄰接串列裡出現兩次。 從子節點 $v$ 看到父節點 $u$ 時，那是樹邊 $\{u, v\}$
本身，不是後向邊，要跳過。 簡單圖裡用「$w$ 是不是父節點」判斷就夠了。
\item 後向邊 $\{x, w\}$（$w$ 是 $x$ 的祖先）也會被看到兩次：從 $x$ 看到 $w$ 時 $w$ 是灰
色；從 $w$ 看到 $x$ 時，$x$ 已經完成。 後者 $d[x] > d[w]$，拿來更新 $\mathit{low}[w]$ 不
會改變任何值，可以不必特別排除。
\end{itemize}

令 $T(v)$ 表示 $v$ 在 DFS 樹裡的子樹。 一條邊如果一端在 $T(v)$ 裡、另一端在外面，因為
兩端有祖先關係，外面那一端一定是 $v$ 的真祖先。 所以：

\begin{lemma}
\label{lem:con-leave}
設 $v$ 是 $u$ 在 DFS 樹裡的子節點。 連接 $T(v)$ 和 $T(v)$ 以外頂點的邊只有兩種：樹邊
$\{u, v\}$，以及從 $T(v)$ 裡的頂點連到 $u$ 或 $u$ 的祖先的後向邊。
\end{lemma}

\subsection{low 值}

\begin{definition}
在無向圖的 DFS 裡，頂點 $u$ 的 $\mathit{low}[u]$ 是下面這些數的最小值：
\begin{itemize}
\item $d[u]$；
\item $\mathit{low}[v]$，$v$ 是 $u$ 的子節點；
\item $d[w]$，$\{u, w\}$ 是後向邊，$w$ 是 $u$ 的祖先（不含走下來的那條樹邊的父節點）。
\end{itemize}
\end{definition}

$d$ 的編號和第~\ref{sec:con-tarjan}~節一樣是發現順序 $1, 2, \ldots, n$。 把子節點一層
一層展開，就得到另一種寫法：
\[
\mathit{low}[u] = \min\bigl(\{d[u]\} \cup \{d[w] : x \in T(u)，\{x, w\} \text{ 是後向邊、} w \text{ 是 } x \text{ 的祖先}\}\bigr),
\]
也就是「從 $u$ 的子樹出發，往下走任意條樹邊，最後\textbf{最多走一條後向邊}，能到達的
最早發現的頂點」的 $d$ 值。 有些課本用遞迴的寫法定義，有些用「子樹裡的後向邊」定義，
兩者算出來的數字相同。 子樹裡的頂點 $d$ 都比 $u$ 大，所以 $\mathit{low}[u] < d[u]$ 表
示子樹裡有後向邊回到 $u$ 的真祖先；$\mathit{low}[u] = d[u]$ 表示沒有。

和 Tarjan 演算法比，這裡沒有堆疊：無向圖沒有交叉邊，不會碰到「已經切出去的元件」，
後向邊另一端的祖先一定還是灰色。

由引理~\ref{lem:con-leave}，對 $u$ 的子節點 $v$，後向邊離開 $T(v)$ 時只能到 $u$（$d$
值是 $d[u]$）或 $u$ 的真祖先（$d$ 值比 $d[u]$ 小），而 $T(v)$ 裡的頂點 $d$ 值都比
$d[u]$ 大。 所以：
\begin{itemize}
\item $\mathit{low}[v] \ge d[u]$ 若且唯若 $T(v)$ 沒有後向邊連到 $u$ 的\textbf{真祖先}；
\item $\mathit{low}[v] > d[u]$ 若且唯若 $T(v)$ 沒有後向邊連到 $u$ \textbf{或} $u$ 的真
祖先，也就是除了樹邊 $\{u, v\}$ 之外，沒有任何邊離開 $T(v)$。
\end{itemize}

\subsection{三個判斷規則}

\begin{theorem}[根]
\label{thm:con-root}
DFS 樹的根 $r$ 是割點，若且唯若 $r$ 在 DFS 樹裡至少有兩個子節點。
\end{theorem}

\begin{proof}
設 $r$ 有兩個子節點 $v_1$、$v_2$。 由引理~\ref{lem:con-leave}，離開 $T(v_1)$ 的邊只能
連到 $r$（$r$ 沒有祖先），所以拿掉 $r$ 後，$T(v_1)$ 和 $T(v_2)$ 之間沒有邊，分成不同的
元件，個數變多，$r$ 是割點。 反過來，$r$ 沒有子節點時 $r$ 是孤立點，拿掉後元件少一
個；$r$ 只有一個子節點 $v$ 時，$r$ 的元件拿掉 $r$ 後剩下 $T(v)$，靠樹邊就連通，元件個
數不變。 這兩種情況 $r$ 都不是割點。
\end{proof}

\begin{theorem}[非根]
\label{thm:con-nonroot}
非根的頂點 $u$ 是割點，若且唯若 $u$ 有某個子節點 $v$ 滿足 $\mathit{low}[v] \ge d[u]$。
\end{theorem}

\begin{proof}
設 $p$ 是 $u$ 的父節點，$K$ 是 $u$ 所在的連通元件（也就是這棵 DFS 樹的頂點）。

若某個子節點 $v$ 滿足 $\mathit{low}[v] \ge d[u]$：$T(v)$ 沒有後向邊連到 $u$ 的真祖先，由
引理~\ref{lem:con-leave}，離開 $T(v)$ 的邊全都接在 $u$ 上。 拿掉 $u$ 以後，$T(v)$ 和
$p$ 不再相連（$p$ 不在 $T(v)$ 裡），$K$ 至少分成兩塊，$u$ 是割點。

若每個子節點 $v$ 都滿足 $\mathit{low}[v] < d[u]$：拿掉 $u$ 後，把 $K$ 剩下的頂點分成
$R = K - T(u)$ 和各個 $T(v)$。 樹拿掉子樹 $T(u)$ 後仍然是一棵樹，所以 $R$ 靠樹邊連通，
而且 $p \in R$。 每個 $T(v)$ 自己靠樹邊連通，又有一條後向邊連到 $u$ 的真祖先，那個祖先
在 $R$ 裡。 所以 $K - u$ 仍然連通，元件個數不變，$u$ 不是割點。
\end{proof}

\begin{theorem}[橋]
\label{thm:con-bridge}
後向邊都不是橋。 樹邊 $\{u, v\}$（$u$ 是 $v$ 的父節點）是橋，若且唯若 $\mathit{low}[v]
> d[u]$。
\end{theorem}

\begin{proof}
後向邊 $\{x, w\}$ 加上樹上從 $w$ 到 $x$ 的路，構成一個環。 一條邊在環上，拿掉它以後兩
端還能沿著環的另一半相連，元件個數不變，所以不是橋。

樹邊 $\{u, v\}$：若 $\mathit{low}[v] > d[u]$，除了 $\{u, v\}$ 本身，沒有任何邊離開
$T(v)$，拿掉它，$T(v)$ 就和 $u$ 分開了，是橋。 若 $\mathit{low}[v] \le d[u]$，$T(v)$ 裡
有某個 $x$ 經後向邊連到 $w$，$w$ 是 $u$ 或 $u$ 的真祖先。 這條後向邊不是 $\{u, v\}$ 本
身（那是樹邊），所以 $w \leadsto u \to v \leadsto x \to w$ 是一個包含 $\{u, v\}$ 的環，
$\{u, v\}$ 不是橋。
\end{proof}

橋的判斷是嚴格的 $>$，非根割點的判斷是 $\ge$。 差別在「後向邊剛好回到 $u$」：這時拿掉
$u$ 會把 $T(v)$ 切開，但拿掉 $\{u, v\}$ 不會，因為 $T(v)$ 還能經過那條後向邊回到 $u$。

\begin{pseudo}[label=alg:con-cut]{割點與橋}
\Procedure{Cut-Bridge}{$G$}
    \State $\mathit{time} \gets 0$
    \State 把每個 $d[u]$ 設成 $0$、$\mathit{cut}[u]$ 設成 \FALSE
    \For{每個頂點 $u$（由小到大）}
        \If{$d[u] = 0$}
            \State \Call{Cut-Visit}{$G, u, \NIL$}
        \EndIf
    \EndFor
\EndProcedure
\Procedure{Cut-Visit}{$G, u, p$} \Comment{$p$ 是 $u$ 的父節點}
    \State $\mathit{time} \gets \mathit{time} + 1$；$d[u] \gets \mathit{time}$；$\mathit{low}[u] \gets \mathit{time}$
    \State $c \gets 0$ \Comment{子節點個數}
    \For{每個 $w \in \mathit{Adj}[u]$}
        \If{$d[w] = 0$} \Comment{樹邊}
            \State $c \gets c + 1$
            \State \Call{Cut-Visit}{$G, w, u$}
            \State $\mathit{low}[u] \gets \min(\mathit{low}[u], \mathit{low}[w])$
            \If{$p \ne \NIL$ \AND $\mathit{low}[w] \ge d[u]$}
                \State $\mathit{cut}[u] \gets \TRUE$
            \EndIf
            \If{$\mathit{low}[w] > d[u]$}
                \State 輸出橋 $\{u, w\}$
            \EndIf
        \ElsIf{$w \ne p$} \Comment{後向邊}
            \State $\mathit{low}[u] \gets \min(\mathit{low}[u], d[w])$
        \EndIf
    \EndFor
    \If{$p = \NIL$ \AND $c \ge 2$}
        \State $\mathit{cut}[u] \gets \TRUE$
    \EndIf
\EndProcedure
\end{pseudo}

第 15 行可能對同一個 $u$ 成立好幾次（$u$ 有好幾個子節點都滿足），所以用
$\mathit{cut}[u]$ 記下來，不要每次成立就把割點個數加一。 第 19--20 行的 $w$ 可能是已經完成
的子孫（後向邊從祖先那一端看到），$d[w] > d[u]$，取 $\min$ 不會改變任何值。 每條邊在
第 10 行被看兩次，總時間 $\Theta(n + m)$。

\subsection{手算追蹤}

在圖~\ref{fig:con-corridor} 上從 $1$ 開始，鄰居由小到大。 $1$ 的鄰居是 $3, 5, 10$，先走
到 $3$；$3$ 的鄰居 $1$ 是父節點，跳過，接著走到 $5$；$5$ 的鄰居 $1$ 是祖先，$\{5, 1\}$
是後向邊。 $5$ 完成，回到 $3$ 走 $6$，$6 \to 4 \to 2$；$2$ 只有父節點這個鄰居，完成；
$4$ 的下一個鄰居 $6$ 是父節點，再來 $8$；$8$ 的鄰居 $4$ 是父節點，$6$ 是祖先（後向邊
$\{8, 6\}$），接著 $8 \to 7 \to 9$，$9$ 的鄰居 $8$ 是祖先（後向邊 $\{9, 8\}$）。 一路完
成回到 $1$ 以後，$1$ 的鄰居 $5$ 已經完成（$d[5] > d[1]$，略過），最後走到 $10$。 得到
的 DFS 樹如下。

\begin{center}
\begin{tikzpicture}[x=1.0cm, y=0.95cm]
\node[con-v] (t1) at (0,0) {1};
\node[con-v] (t3) at (-1.3,-1) {3};
\node[con-v] (t10) at (1.3,-1) {10};
\node[con-v] (t5) at (-2.4,-2) {5};
\node[con-v] (t6) at (-0.4,-2) {6};
\node[con-v] (t4) at (-0.4,-3) {4};
\node[con-v] (t2) at (-1.4,-4) {2};
\node[con-v] (t8) at (0.6,-4) {8};
\node[con-v] (t7) at (0.6,-5) {7};
\node[con-v] (t9) at (0.6,-6) {9};
\draw[con-tree] (t1) -- (t3);
\draw[con-tree] (t1) -- (t10);
\draw[con-tree] (t3) -- (t5);
\draw[con-tree] (t3) -- (t6);
\draw[con-tree] (t6) -- (t4);
\draw[con-tree] (t4) -- (t2);
\draw[con-tree] (t4) -- (t8);
\draw[con-tree] (t8) -- (t7);
\draw[con-tree] (t7) -- (t9);
\draw[con-back] (t5) to[bend left=65] (t1);
\draw[con-back] (t8) to[bend right=45] (t6);
\draw[con-back] (t9) to[bend right=50] (t8);
\node[con-lab, right=1mm of t1] {1/1};
\node[con-lab, above left=0.3mm of t3] {2/1};
\node[con-lab, right=1mm of t10] {10/10};
\node[con-lab, left=1mm of t5] {3/1};
\node[con-lab, left=1mm of t6] {4/4};
\node[con-lab, left=1mm of t4] {5/4};
\node[con-lab, left=1mm of t2] {6/6};
\node[con-lab, left=1mm of t8] {7/4};
\node[con-lab, left=1mm of t7] {8/7};
\node[con-lab, left=1mm of t9] {9/7};
\end{tikzpicture}
\diagcaption{圖~\ref{fig:con-corridor} 從 $1$ 開始的 DFS 樹。 粗線是樹邊，虛線是後向
邊；每個頂點旁邊是 $d/\mathit{low}$。}
\label{fig:con-corridor-tree}
\end{center}

$\mathit{low}$ 要從葉節點往上算：子節點的 $\mathit{low}$ 算好，父節點才能算。 下表依 $d$
排列，最後一欄是每條樹邊 $\{u, v\}$ 的判斷（$v$ 是子節點）。

\begin{center}
\small
\setlength{\tabcolsep}{4pt}
\begin{tabular}{ccccc>{\raggedright\arraybackslash}p{0.44\linewidth}}
\toprule
$u$ & $d$ & 子節點 & 後向邊到 & $\mathit{low}$ & 判斷 \\
\midrule
1 & 1 & 3, 10 & — & $\min(1, 1, 10) = 1$ & 根，兩個子節點：割點 \\
3 & 2 & 5, 6 & — & $\min(2, 1, 4) = 1$ & $\mathit{low}[6] = 4 \ge 2$：割點；$4 > 2$：$\{3, 6\}$ 是橋 \\
5 & 3 & — & 1 & $\min(3, d[1]) = 1$ & 葉節點 \\
6 & 4 & 4 & — & $\min(4, 4) = 4$ & $\mathit{low}[4] = 4 \ge 4$：割點；$4 \not> 4$：不是橋 \\
4 & 5 & 2, 8 & — & $\min(5, 6, 4) = 4$ & $\mathit{low}[2] = 6 \ge 5$：割點；$6 > 5$：$\{4, 2\}$ 是橋 \\
2 & 6 & — & — & $6$ & 葉節點 \\
8 & 7 & 7 & 6 & $\min(7, d[6], 7) = 4$ & $\mathit{low}[7] = 7 \ge 7$：割點；不是橋 \\
7 & 8 & 9 & — & $\min(8, 7) = 7$ & $\mathit{low}[9] = 7 < 8$ \\
9 & 9 & — & 8 & $\min(9, d[8]) = 7$ & 葉節點 \\
10 & 10 & — & — & $10$ & 葉節點；$\mathit{low}[10] = 10 > 1$：$\{1, 10\}$ 是橋 \\
\bottomrule
\end{tabular}
\end{center}

結果：割點是 $1, 3, 4, 6, 8$，橋是 $\{1, 10\}$、$\{3, 6\}$、$\{2, 4\}$。 幾個要注意的地方：

\begin{itemize}
\item $3$ 有兩個子節點，$5$ 那邊有後向邊回到 $1$（$\mathit{low}[5] = 1 < 2$），$6$ 那邊沒
有。 只要\textbf{有一個}子節點滿足條件就是割點。
\item $6$ 和 $8$ 都是「$\mathit{low}[v] = d[u]$」的情況：子樹的後向邊剛好回到 $u$，$u$
是割點，但邊 $\{u, v\}$ 不是橋。
\item $4$ 和 $6$ 都是割點，但 $\{4, 6\}$ 不是橋。 兩端都是割點的邊不一定是橋；反過來，
橋的兩端也不一定都是割點，例如 $\{1, 10\}$ 的 $10$ 是葉節點。 DFS 樹的葉節點沒有子節點，定理~\ref{thm:con-nonroot} 的條件不會成立，所以一定不是割點。
\item 根 $1$ 不能用非根的規則：對子節點 $3$，$\mathit{low}[3] = 1 \ge d[1]$ 永遠成立，根
要另外數子節點。
\end{itemize}

\begin{center}
\begin{tikzpicture}[x=1.05cm, y=1.05cm]
\node[con-v] (u10) at (0,1.2) {10};
\node[con-v, con-cut] (u1) at (1.4,1.2) {1};
\node[con-v] (u5) at (2.3,2.3) {5};
\node[con-v, con-cut] (u3) at (2.6,0.5) {3};
\node[con-v, con-cut] (u6) at (4.1,0.5) {6};
\node[con-v, con-cut] (u4) at (5.0,1.7) {4};
\node[con-v] (u2) at (4.6,3.0) {2};
\node[con-v, con-cut] (u8) at (5.9,0.5) {8};
\node[con-v] (u7) at (7.3,1.3) {7};
\node[con-v] (u9) at (7.3,-0.3) {9};
\draw[con-e] (u1) -- (u3);
\draw[con-e] (u1) -- (u5);
\draw[con-bridge] (u1) -- (u10);
\draw[con-bridge] (u2) -- (u4);
\draw[con-e] (u3) -- (u5);
\draw[con-bridge] (u3) -- (u6);
\draw[con-e] (u4) -- (u6);
\draw[con-e] (u4) -- (u8);
\draw[con-e] (u6) -- (u8);
\draw[con-e] (u7) -- (u8);
\draw[con-e] (u7) -- (u9);
\draw[con-e] (u8) -- (u9);
\end{tikzpicture}
\diagcaption{校園的割點（橘色）和橋（紅色粗線）。}
\label{fig:con-corridor-cut}
\end{center}

\begin{examnote}
$\mathit{low}$ 的手算表各課本寫法不同，作答前先看題目的約定。 (1) $d$ 的編號有的從
$0$ 開始，有的用第~21~章同時記 $d$、$f$ 的時間戳記，數字會不同，但判斷只比大小，割
點、橋的結果一樣。 (2) 遞迴的定義和「子樹裡的後向邊」的定義算出來的數字相同。 (3) 有
的程式在後向邊寫 $\mathit{low}[u] \gets \min(\mathit{low}[u], \mathit{low}[w])$，用
$\mathit{low}[w]$ 代替 $d[w]$。 祖先 $w$ 這時的 $\mathit{low}$ 可能已經比 $d[w]$ 小，算
出的數字就會變小：上例 $8$ 先處理後向邊 $\{8, 6\}$，$\mathit{low}[8]$ 已經是 $4$，之後
$9$ 的後向邊 $\{9, 8\}$ 會讓 $\mathit{low}[9] = \mathit{low}[7] = 4$，$\mathit{low}[7] = 4
< d[8] = 7$，\textbf{割點 $8$ 就找不到了}。 這種寫法找橋的結果仍然正確，找割點會錯，
作答時一律用 $d[w]$。（Tarjan 強連通元件演算法做同樣的替換，數字也會變，但找到的元件不
變，見觀念題第~\ref{q:con-variant-proof}~題。）
\end{examnote}

\subsection{平行邊}

如果兩棟大樓之間有兩條走廊，拿掉其中一條還有另一條，它們都不是橋。 但虛擬碼第 19 行
用「$w \ne p$」判斷，會把兩條都當成走下來的樹邊跳過，$\mathit{low}[v]$ 算大了，就把
$\{u, v\}$ 誤判成橋。 有平行邊時，要幫每條邊編號，遞迴時傳「走下來的那條邊的編號」，
只跳過那一條，其他連回父節點的邊照樣當後向邊。 割點不受影響：回到父節點 $u$ 的後向邊
最多讓 $\mathit{low}[v]$ 降到 $d[u]$，$\mathit{low}[v] \ge d[u]$ 的判斷結果不變。

\section{雙連通元件}
\label{sec:con-bcc}

割點把圖分成幾塊，每一塊內部「封掉任何一棟大樓都還走得通」。 這些塊就是雙連通元件。

\begin{definition}
無向圖裡兩條邊 $e_1$、$e_2$，如果 $e_1 = e_2$，或者有一個\textbf{簡單環}（頂點不重複
的環）同時經過 $e_1$ 和 $e_2$，就說它們屬於同一塊。 這是一個等價關係，它把所有的邊分
成若干類，每一類叫做一個\textbf{雙連通元件}（biconnected component），也叫\textbf{區
塊}（block）。
\end{definition}

自反（$e_1 = e_1$）和對稱（同一個環也同時經過 $e_2$、$e_1$）直接從定義看得出來；遞移性（$e_1, e_2$ 在同一個環上、$e_2, e_3$ 在同一個環上，則
$e_1, e_3$ 也在某個環上）需要一點功夫，這裡不證，見延伸閱讀。 幾個常用的性質：

\begin{itemize}
\item 橋不在任何環上，所以每一座橋自己是一個區塊。
\item 兩個不同的區塊最多共用一個頂點，共用的頂點一定是割點（觀念題第~\ref{q:con-block-share}~題）。
\item 一個頂點是割點，若且唯若它屬於兩個以上的區塊。
\end{itemize}

有的課本改用頂點定義：沒有割點的連通圖叫\textbf{雙連通圖}（只有一條邊、兩個頂點的
圖也算），雙連通元件是最大的雙連通子圖。 兩種定義對有邊的部分得到同樣的分法；差別只在孤立
頂點，頂點定義把它算成一個元件，邊定義不算。

圖~\ref{fig:con-corridor-cut} 的校園有 $6$ 個區塊：
\[
\{2\text{--}4\},\ \{7\text{--}8, 7\text{--}9, 8\text{--}9\},\ \{4\text{--}6, 4\text{--}8, 6\text{--}8\},\
\{3\text{--}6\},\ \{1\text{--}3, 1\text{--}5, 3\text{--}5\},\ \{1\text{--}10\}.
\]
每個割點都
剛好落在兩個區塊的交界：$4$ 在第一、三塊，$8$ 在第二、三塊，$6$ 在第三、四塊，$3$ 在第
四、五塊，$1$ 在第五、六塊。

\subsection{用邊的堆疊找區塊}

在割點演算法上加一個\textbf{邊的堆疊}就能把區塊切出來，想法和 Tarjan 強連通元件演算
法一樣：DFS 走過的邊依序推進堆疊；當子節點 $w$ 做完、發現 $\mathit{low}[w] \ge d[u]$
時，$u$ 把 $T(w)$ 這一側和圖的其他部分隔開，堆疊上從樹邊 $(u, w)$ 以上、還沒被彈出的
邊，就是一個區塊。

\begin{pseudo}[label=alg:con-bcc]{雙連通元件（主程序和虛擬碼~\ref{alg:con-cut} 相同，另外設 $S \gets$ 空堆疊、$k \gets 0$）}
\Procedure{BCC-Visit}{$G, u, p$}
    \State $\mathit{time} \gets \mathit{time} + 1$；$d[u] \gets \mathit{time}$；$\mathit{low}[u] \gets \mathit{time}$
    \For{每個 $w \in \mathit{Adj}[u]$}
        \If{$d[w] = 0$} \Comment{樹邊}
            \State \Call{Push}{$S, (u, w)$}
            \State \Call{BCC-Visit}{$G, w, u$}
            \State $\mathit{low}[u] \gets \min(\mathit{low}[u], \mathit{low}[w])$
            \If{$\mathit{low}[w] \ge d[u]$}
                \State $k \gets k + 1$ \Comment{新的區塊}
                \Repeat
                    \State $(x, y) \gets$ \Call{Pop}{$S$}；把邊 $\{x, y\}$ 標成第 $k$ 個區塊
                \Until{$(x, y) = (u, w)$}
            \EndIf
        \ElsIf{$w \ne p$ \AND $d[w] < d[u]$} \Comment{往祖先的後向邊}
            \State \Call{Push}{$S, (u, w)$}
            \State $\mathit{low}[u] \gets \min(\mathit{low}[u], d[w])$
        \EndIf
    \EndFor
\EndProcedure
\end{pseudo}

和虛擬碼~\ref{alg:con-cut} 比，有三個差別。 第一，第 13 行多了 $d[w] < d[u]$：後向邊從
祖先那一端看到時不能再推一次，否則同一條邊會進堆疊兩次。 第二，第 8 行的判斷\textbf{包
含根}：根的每個子節點都會切出區塊，根有兩個子節點時就切出兩組，這正好對應「根是割點」。
第三，每條邊恰好進出堆疊各一次，所以總時間仍然是 $\Theta(n + m)$，堆疊最多放 $m$ 條邊。

為什麼一次彈出的剛好是一個區塊？ 完整的證明比較長，這裡只說想法。 對彈出的先後做歸
納：$T(w)$ 裡比較深的地方，凡是被某個割點隔開的部分，都在更早的時候被彈走了；留下
來的邊，每一條都能透過子樹裡的後向邊和 $(u, w)$ 繞成環，所以都在同一個區塊。 另一方
面，$\mathit{low}[w] \ge d[u]$ 表示 $T(w)$ 沒有邊越過 $u$ 往上，$u$ 以上的邊和這些邊不可
能在同一個環上，所以這一塊不會再變大。

\subsection{手算追蹤}

DFS 的走法和圖~\ref{fig:con-corridor-tree} 相同。 邊寫成 $u$--$w$ 表示從 $u$ 看到 $w$
時推進堆疊（樹邊由父到子，後向邊由子孫到祖先）。 只有判斷成立的六個時刻會彈出：

\begin{center}
\small
\setlength{\tabcolsep}{4pt}
\begin{tabular}{ll>{\raggedright\arraybackslash}p{0.42\linewidth}l}
\toprule
時刻 & 判斷 & 彈出前的堆疊（底 $\to$ 頂） & 彈出的區塊 \\
\midrule
從 $2$ 回到 $4$ & $\mathit{low}[2] = 6 \ge 5$ & \mbox{1--3}, \mbox{3--5}, \mbox{5--1}, \mbox{3--6}, \mbox{6--4}, \mbox{4--2} & \mbox{4--2} \\
從 $7$ 回到 $8$ & $\mathit{low}[7] = 7 \ge 7$ & \mbox{1--3}, \mbox{3--5}, \mbox{5--1}, \mbox{3--6}, \mbox{6--4}, \mbox{4--8}, \mbox{8--6}, \mbox{8--7}, \mbox{7--9}, \mbox{9--8} & \mbox{9--8}, \mbox{7--9}, \mbox{8--7} \\
從 $4$ 回到 $6$ & $\mathit{low}[4] = 4 \ge 4$ & \mbox{1--3}, \mbox{3--5}, \mbox{5--1}, \mbox{3--6}, \mbox{6--4}, \mbox{4--8}, \mbox{8--6} & \mbox{8--6}, \mbox{4--8}, \mbox{6--4} \\
從 $6$ 回到 $3$ & $\mathit{low}[6] = 4 \ge 2$ & \mbox{1--3}, \mbox{3--5}, \mbox{5--1}, \mbox{3--6} & \mbox{3--6} \\
從 $3$ 回到 $1$ & $\mathit{low}[3] = 1 \ge 1$ & \mbox{1--3}, \mbox{3--5}, \mbox{5--1} & \mbox{5--1}, \mbox{3--5}, \mbox{1--3} \\
從 $10$ 回到 $1$ & $\mathit{low}[10] = 10 \ge 1$ & \mbox{1--10} & \mbox{1--10} \\
\bottomrule
\end{tabular}
\end{center}

其他幾次回到父節點時判斷都不成立：從 $5$ 回到 $3$（$\mathit{low}[5] = 1 < 2$）、從 $8$
回到 $4$（$4 < 5$）、從 $9$ 回到 $7$（$7 < 8$）。 這些邊留在堆疊上，等上面的頂點一起
彈出。 注意第二列：$4$--$8$ 比 $8$--$6$ 早進堆疊，但它不屬於第二個區塊，彈到 $8$--$7$
就停了，$4$--$8$ 和 $8$--$6$ 留給第三列。

\begin{examnote}
雙連通元件的答案可以寫成邊的集合，也可以寫成頂點的集合（例如第三塊是 $\{4, 6, 8\}$）；
寫頂點時，割點會同時出現在好幾塊裡，這是正常的。 題目用頂點定義、圖裡又有孤立頂點時，
每個孤立頂點自己算一塊，邊的堆疊不會產生它，要另外補上。 有的題目只問區塊個數，就是
第 8 行成立的次數。
\end{examnote}

\subsection{C 程式}

下面的程式把本章無向圖的三件事一次做完：照虛擬碼~\ref{alg:con-bcc} 的走法，在判斷成
立時印出區塊（如果同時 $\mathit{low}[w] > d[u]$，先印出橋），最後印出每個頂點的 $d$、
$\mathit{low}$ 和所有割點。 每條無向邊在鄰接串列裡加兩次。 邊的堆疊是一個
\texttt{Edge} 陣列，最多放 $m$ 條邊；\texttt{(Edge)\{u, w\}} 是 C99 起的複合字面值，直接
造出一個結構值。 輸入是圖~\ref{fig:con-corridor} 的校園。

\begin{example}[label=ex:con-cut]{c}{(cut\_bridge.c) 割點、橋與雙連通元件}
#include <stdio.h>
#include <stdlib.h>

typedef struct Node {
    int v;
    struct Node *next;
} Node;

typedef struct {
    int n;
    Node **adj;     // adj[u]：u 的鄰居，由小到大
} Graph;

int add_arc(Graph *g, int u, int v)
{
    Node *x = malloc(sizeof *x);
    if (x == NULL) {
        return 0;
    }
    Node **p = &g->adj[u];
    while (*p != NULL && (*p)->v < v) {
        p = &(*p)->next;
    }
    x->v = v;
    x->next = *p;
    *p = x;
    return 1;
}

void free_graph(Graph *g)
{
    for (int u = 0; g->adj != NULL && u < g->n; u++) {
        Node *x = g->adj[u];
        while (x != NULL) {
            Node *next = x->next;
            free(x);
            x = next;
        }
    }
    free(g->adj);
}

typedef struct {
    int u, w;
} Edge;

typedef struct {
    const Graph *g;
    int *d, *low;       // d[u] = 0 表示還沒發現
    int *is_cut;
    Edge *es;           // 邊的堆疊
    int top, time, blocks;
} Search;

void visit(Search *s, int u, int parent)
{
    s->d[u] = s->low[u] = ++s->time;
    int children = 0;
    for (Node *x = s->g->adj[u]; x != NULL; x = x->next) {
        int w = x->v;
        if (s->d[w] == 0) {                         // 樹邊
            children++;
            s->es[s->top++] = (Edge){u, w};
            visit(s, w, u);
            if (s->low[w] < s->low[u]) {
                s->low[u] = s->low[w];
            }
            if (s->low[w] >= s->d[u]) {     // w 的子樹到不了 u 上面
                if (parent != -1) {
                    s->is_cut[u] = 1;
                }
                if (s->low[w] > s->d[u]) {
                    printf("bridge %d-%d\n", u + 1, w + 1);
                }
                printf("block %d:", ++s->blocks);
                Edge e;
                do {
                    e = s->es[--s->top];
                    printf(" %d-%d", e.u + 1, e.w + 1);
                } while (e.u != u || e.w != w);
                printf("\n");
            }
        } else if (w != parent && s->d[w] < s->d[u]) {  // 後向邊
            s->es[s->top++] = (Edge){u, w};
            if (s->d[w] < s->low[u]) {
                s->low[u] = s->d[w];
            }
        }
    }
    if (parent == -1 && children >= 2) {
        s->is_cut[u] = 1;
    }
}

int main(void)
{
    int n, m;
    if (scanf("%d %d", &n, &m) != 2 || n <= 0 || m < 0) {
        return 1;
    }
    Graph g = {n, calloc(n, sizeof(Node *))};
    Search s = {&g, calloc(n, sizeof(int)), calloc(n, sizeof(int)),
                calloc(n, sizeof(int)),
                malloc((m + 1) * sizeof(Edge)), 0, 0, 0};
    int ok = g.adj != NULL && s.d != NULL && s.low != NULL
             && s.is_cut != NULL && s.es != NULL;
    for (int i = 0; ok && i < m; i++) {
        int u, v;
        ok = scanf("%d %d", &u, &v) == 2 && u >= 1 && u <= n
             && v >= 1 && v <= n && u != v
             && add_arc(&g, u - 1, v - 1)
             && add_arc(&g, v - 1, u - 1);
    }
    if (ok) {
        for (int u = 0; u < n; u++) {
            if (s.d[u] == 0) {
                visit(&s, u, -1);
            }
        }
        printf("v:");
        for (int u = 0; u < n; u++) {
            printf(" %d", u + 1);
        }
        printf("\nd:");
        for (int u = 0; u < n; u++) {
            printf(" %d", s.d[u]);
        }
        printf("\nlow:");
        for (int u = 0; u < n; u++) {
            printf(" %d", s.low[u]);
        }
        printf("\ncut vertices:");
        for (int u = 0; u < n; u++) {
            if (s.is_cut[u]) {
                printf(" %d", u + 1);
            }
        }
        printf("\n");
    }
    free_graph(&g);
    free(s.d);
    free(s.low);
    free(s.is_cut);
    free(s.es);
    return ok ? 0 : 1;
}
\end{example}
\iosample{10 12 \\ 1 3 \\ 1 5 \\ 1 10 \\ 2 4 \\ 3 5 \\ 3 6 \\ 4 6 \\ 4 8 \\ 6 8 \\ 7 8 \\ 7 9 \\ 8 9}{bridge 4-2 \\ block 1: 4-2 \\ block 2: 9-8 7-9 8-7 \\ block 3: 8-6 4-8 6-4 \\ bridge 3-6 \\ block 4: 3-6 \\ block 5: 5-1 3-5 1-3 \\ bridge 1-10 \\ block 6: 1-10 \\ v: 1 2 3 4 5 6 7 8 9 10 \\ d: 1 6 2 5 3 4 8 7 9 10 \\ low: 1 6 1 4 1 4 7 4 7 10 \\ cut vertices: 1 3 4 6 8}

輸出和手算一致。 程式只在 $d[w] < d[u]$ 時處理後向邊，所以「從祖先那一端看到已完成的子
孫」的情況直接略過，連 $\mathit{low}$ 都不用更新。 \texttt{is\_cut} 是旗標陣列，同一個
頂點被好幾個子節點判斷成割點也只記一次。 程式假設輸入是簡單圖（讀入時擋掉了自環，平
行邊沒有檢查）。

把所有的橋拿掉，剩下的每個連通元件叫做\textbf{邊雙連通元件}（2-edge-connected
component）：裡面拿掉任何一條邊都還連通。 校園的例子拿掉三座橋後剩下 $\{1, 3, 5\}$、
$\{4, 6, 7, 8, 9\}$、$\{2\}$、$\{10\}$ 四塊。 注意 $\{4, 6, 7, 8, 9\}$ 包含兩個區塊：
邊雙連通只管邊，不管割點。

\section{本章常見錯誤}
\label{sec:con-pitfalls}

\begin{center}
\small
\begin{tabular}{>{\raggedright\arraybackslash}p{0.24\linewidth}>{\raggedright\arraybackslash}p{0.52\linewidth}>{\raggedright\arraybackslash}p{0.14\linewidth}}
\toprule
錯誤 & 例子 & 見 \\
\midrule
第二次 DFS 在 $G$ 上做 & Kosaraju 的第二次 DFS 沒有轉置，從 $7$ 出發就走遍全圖，四個元件併成一個 &
第~\ref{sec:con-kosaraju}~節 \\
挑起點的順序錯 & 第二次 DFS 照 $f$ 由小到大，或照 $d$、照編號挑起點 & 第~\ref{sec:con-kosaraju}~節 \\
沒檢查 $\mathit{onstack}$ & Tarjan 處理 $3 \to 5$ 時，$5$ 已經彈出還拿 $d[5]$ 更新，$\{3, 8\}$ 沒有被切出來 &
第~\ref{sec:con-tarjan}~節 \\
用灰色代替「在堆疊上」 & $4 \to 2$ 的 $2$ 已是黑色但還在堆疊上，只看灰色會讓 $\mathit{low}[4]$ 留在
$6$，$4$ 被誤判成頭 & 第~\ref{sec:con-tarjan}~節 \\
拿 $\mathit{low}$ 當元件編號 & $\mathit{low}[4] = 2$，但 $4$ 的元件的頭是 $1$ & 第~\ref{sec:con-tarjan}~節 \\
父節點的邊當成後向邊 & 無向圖沒有跳過父節點：每個 $\mathit{low}[v] \le d[u]$，一座橋都找不到（割點
不受影響） & 第~\ref{sec:con-cut}~節 \\
根用非根的規則 & 對根來說 $\mathit{low}[v] \ge d[r]$ 永遠成立，每個有子節點的根都被當成割點 &
定理~\ref{thm:con-root} \\
橋寫成 $\ge$ & 校園的 $\{6, 4\}$、$\{8, 7\}$ 被誤判成橋 & 定理~\ref{thm:con-bridge} \\
割點重複計數 & $u$ 有兩個子節點都滿足 $\mathit{low}[v] \ge d[u]$，割點個數加兩次 &
第~\ref{sec:con-cut}~節 \\
後向邊用 $\mathit{low}[w]$ & 校園的割點 $8$ 找不到 & 第~\ref{sec:con-cut}~節作答提醒 \\
有平行邊還用 $w \ne p$ & 兩條走廊連同一對大樓，被誤判成橋 & 第~\ref{sec:con-cut}~節 \\
後向邊推兩次 & 雙連通元件的邊堆疊沒有限制 $d[w] < d[u]$，同一條邊出現在兩個區塊 &
第~\ref{sec:con-bcc}~節 \\
遞迴太深 & 頂點排成很長的鏈時，遞迴的 DFS 用光呼叫堆疊 & 第~\ref{sec:gtr-dfs-iter}~節 \\
\bottomrule
\end{tabular}
\end{center}

\section{觀念題}
\label{sec:con-exercises}

題型和研究所考試、期中期末考的筆試相近。 沒有特別說明時，鄰接串列的鄰居依字母（或編號）
由小到大，主迴圈也從最小的頂點開始；Tarjan 演算法和割點演算法的 $d$ 是發現順序
$1, 2, \ldots$；無向圖是簡單圖。 解答在本章最後。

第~\ref{q:con-kosaraju}--\ref{q:con-dag-h} 題用同一張有向圖 $H$，頂點 $a$ 到 $h$，$14$ 條邊：
\[
\begin{array}{lllllll}
a \to d & a \to e & a \to h & b \to d & b \to e & c \to e & c \to f \\
d \to e & e \to a & e \to b & e \to g & f \to c & h \to d & h \to g
\end{array}
\]

\begin{enumerate}
\item \label{q:con-choice}（選擇，可能不只一個答案）下列哪些敘述正確？
\begin{enumerate}
\item Kosaraju 演算法的第二次 DFS 改在原圖 $G$ 上做，仍然依 $f$ 由大到小挑起點，結果一樣
正確。
\item Tarjan 演算法裡，頂點 $u$ 完成時若 $\mathit{low}[u] = d[u]$，則 $u$ 是它的強連通元件
裡最先被發現的頂點。
\item 元件圖可能有環，但不會有自環。
\item 連通的無向圖裡，兩端都是割點的邊一定是橋。
\item 至少有 $3$ 個頂點的連通無向圖如果有橋，就一定有割點。
\item 無向圖的 DFS 樹裡，沒有子節點的頂點一定不是割點。
\item 非根頂點 $u$ 有子節點 $v$ 滿足 $\mathit{low}[v] = d[u]$，則 $\{u, v\}$ 是橋。
\end{enumerate}

\item \label{q:con-kosaraju}（追蹤）對 $H$ 執行 Kosaraju 演算法。 (a) 寫出第一次 DFS 每個
頂點的 $d/f$（第~21~章的時間戳記），並把每條邊分類。 (b) 寫出依 $f$ 由大到小的順序和
$H^{\mathrm{T}}$ 的鄰接串列。 (c) 寫出第二次 DFS 的每一棵樹。

\item \label{q:con-tarjan-h}（追蹤）對 $H$ 執行 Tarjan 演算法（虛擬碼~\ref{alg:con-tarjan}）。
(a) 寫出每個頂點的 $d$、$\mathit{low}$，以及元件被彈出的順序。 (b) 如果第 15 行改成
$\mathit{low}[u] \gets \min(\mathit{low}[u], \mathit{low}[w])$，哪些頂點的 $\mathit{low}$
會改變？ 輸出的元件有沒有改變？（第~\ref{q:con-variant-proof} 題證明一般的情形。）

\item \label{q:con-dag-h}（計算）(a) 畫出 $H$ 的元件圖，寫出一個拓樸順序。 (b) 從哪些頂點出
發可以到達所有頂點？ (c) 把邊 $c \to e$ 刪掉，至少要選幾個起點才能到達所有頂點？ 舉一組。

\item \label{q:con-variant-proof}（證明）證明：把 Tarjan 演算法第 15 行的 $d[w]$ 換成
$\mathit{low}[w]$，找到的強連通元件和彈出的順序都不變。（提示：令新的值是
$\mathit{low}'$。 證明 $\mathit{low}'[u] \le \mathit{low}[u]$，以及 $\mathit{low}'[u] \ge
d[h(u)]$，其中 $h(u)$ 是 $u$ 的元件的頭。）

\item （計算）(a) $n$ 個頂點的有向圖有 $k$ 個強連通元件，元件圖最多有幾條邊？ (b) $n \ge 2$
個頂點的強連通有向圖最少有幾條邊？ 證明你的答案。 (c) 對一個 DAG 執行 Tarjan 演算法，
每個元件有幾個頂點？ 元件被彈出的順序和第~22~章的拓樸排序有什麼關係？

\item \label{q:con-cut-k}（追蹤）無向圖 $K$ 有頂點 $a$ 到 $i$ 和下面 $12$ 條邊：
\[
\begin{array}{llllll}
a\text{--}e & a\text{--}g & a\text{--}i & b\text{--}f & b\text{--}h & c\text{--}f \\
d\text{--}h & d\text{--}i & e\text{--}g & f\text{--}h & g\text{--}i & h\text{--}i
\end{array}
\]
從 $a$ 開始做 DFS。 (a) 畫出 DFS 樹和後向邊，寫出每個頂點的 $d$ 和 $\mathit{low}$。 (b) 找
出所有割點和橋。 (c) 用虛擬碼~\ref{alg:con-bcc} 依序寫出彈出的區塊。

\item （追蹤）承上題，如果後向邊改用 $\mathit{low}[u] \gets \min(\mathit{low}[u],
\mathit{low}[w])$，哪些頂點的 $\mathit{low}$ 會改變？ 用非根規則會找到哪些割點？ 橋呢？

\item \label{q:con-block-share}（證明）證明：簡單無向圖裡，兩個不同的區塊最多共用一個頂點，而且共用的頂點一定是割點。

\item （追蹤）四棟大樓 $1, 2, 3, 4$ 之間有五條走廊：$1$--$2$ 有兩條，另外是 $2$--$3$、
$3$--$4$、$2$--$4$。 (a) 正確的割點和橋是哪些？ (b) 直接用虛擬碼~\ref{alg:con-cut}
（第 19 行判斷 $w \ne p$），從 $1$ 開始會輸出什麼？ (c) 怎麼改才對？

\item （證明）(a) 一個連通的簡單無向圖有割點、但沒有橋，最少有幾個頂點？ 畫出一個例子，
並證明不能更少。 (b) 舉一個有橋、但沒有割點的連通圖。
\end{enumerate}

\section{實作練習}
\label{sec:con-practice}

下面是 Aki-Judge 上和本章相關的題目。 強連通元件本身目前沒有題目，可以用範
例~\ref{ex:con-tarjan} 自己出幾張圖練習，再和手算對照。

\begin{center}
\begin{tabular}{llc}
\toprule
小節 & 題目 & 難度 \\
\midrule
割點 & \akiref{count-cut-vertices}{割點個數} & 基本 \\
橋 & \akiref{count-bridge-edges}{橋邊個數} & 基本 \\
雙連通元件 & \akiref{count-biconnected-blocks}{雙連通區塊個數} & 進階 \\
縮點 & \akiref{count-reachable-pairs}{可達點對個數} & 進階 \\
\bottomrule
\end{tabular}
\end{center}

幾題的提示：

\begin{itemize}
\item 「割點個數」的圖不一定連通，主迴圈要從每個還沒發現的頂點各做一次 DFS，每棵樹的根
都用根的規則。 同一個頂點可能被好幾個子節點判斷成割點，用旗標陣列記，最後再數。
\item 「橋邊個數」照定理~\ref{thm:con-bridge}，用 $\mathit{low}[v] > d[u]$。 題目保證沒有平
行邊，跳過父節點就好。
\item 「雙連通區塊個數」只問個數，不用真的存邊的堆疊：數虛擬碼~\ref{alg:con-bcc} 第 8 行成
立幾次（包含根）。 沒有邊的孤立頂點不算。
\item 「可達點對個數」的 $n$ 很小，從每個頂點各做一次 DFS 數走得到的頂點就夠快，總共
$O(n(n + m))$。 用本章的想法可以再省：同一個強連通元件裡的頂點能到達的集合一樣，每個元
件只要算一次，再乘上元件的大小。 第~25~章的遞移閉包是另一種做法。
\end{itemize}

\section{讀完你應該會}
\label{sec:con-checklist}

\begin{itemize}
\item 說出強連通元件、元件圖的定義，證明元件圖沒有環，知道 $G$ 和 $G^{\mathrm{T}}$ 的元件
相同。
\item 手算 Kosaraju 演算法：第一次 DFS 的 $d/f$、依 $f$ 由大到小的順序、第二次在
$G^{\mathrm{T}}$ 上的每一棵樹，並用完成時間的引理說明它為什麼對。
\item 手算 Tarjan 演算法的 $d$、$\mathit{low}$ 和堆疊，說明為什麼要檢查「在不在堆疊上」，
以及為什麼 $\mathit{low}[u] = d[u]$ 就是元件的頭。
\item 把元件縮成 DAG，知道 Kosaraju 依拓樸順序、Tarjan 依反向拓樸順序產生元件，用源點
個數回答「從哪裡出發能到每個頂點」。
\item 在無向圖上填 $d$、$\mathit{low}$ 表，用根的規則、$\mathit{low}[v] \ge d[u]$、
$\mathit{low}[v] > d[u]$ 找割點和橋，並證明這三個規則。
\item 用邊的堆疊切出雙連通元件，知道不同課本 $\mathit{low}$ 的寫法哪些等價、哪些會算錯割點。
\end{itemize}

\textbf{延伸閱讀}（書末會整理完整的參考資料）

\begin{itemize}
\item Thomas H. Cormen 等，\textit{Introduction to Algorithms}，第 4 版，20.5 節：用兩次
DFS 求強連通元件，附元件圖和完成時間的引理；第 20 章章末的問題 20-2 用一連串小題帶出割點、
橋和雙連通元件。
\item Robert E. Tarjan，``Depth-First Search and Linear Graph Algorithms''，\textit{SIAM
Journal on Computing} 1(2)，1972：一次 DFS 求強連通元件和雙連通元件的原始論文。
\item M. Sharir，``A Strong-Connectivity Algorithm and Its Applications in Data Flow
Analysis''，\textit{Computers and Mathematics with Applications} 7(1)，1981：兩次 DFS 的做
法；Aho、Hopcroft、Ullman 的 \textit{Data Structures and Algorithms}（1983）把它歸功於
Kosaraju 未發表的筆記，所以常叫 Kosaraju 演算法。
\item John E. Hopcroft、Robert E. Tarjan，``Algorithm 447: Efficient Algorithms for Graph
Manipulation''，\textit{Communications of the ACM} 16(6)，1973：雙連通元件的線性時間
演算法。
\item K. P. Eswaran、R. E. Tarjan，``Augmentation Problems''，\textit{SIAM Journal on
Computing} 5(4)，1976：最少加幾條邊讓圖變成強連通或雙連通。
\item Bengt Aspvall、Michael F. Plass、Robert E. Tarjan，``A Linear-Time Algorithm for
Testing the Truth of Certain Quantified Boolean Formulas''，\textit{Information
Processing Letters} 8(3)，1979：用強連通元件解 2-SAT。
\item Steven S. Skiena，\textit{The Algorithm Design Manual}，第 3 版，7.9、7.10 節：割點和強
連通元件的另一種講法，偏重怎麼想到要用 DFS。
\end{itemize}

\addcontentsline{toc}{section}{觀念題解答}
\section*{觀念題解答}

\begin{enumerate}
\item (b)、(e)、(f) 正確。
(a)：圖~\ref{fig:con-roads} 裡 $f$ 最大的是 $7$，在 $G$ 上從 $7$ 出發會走遍全部 $9$ 個頂
點，四個元件併成一棵樹。 在 $G$ 上改照 $f$ 由小到大挑起點也不對：三個頂點、邊 $1 \to 2$、
$2 \to 1$、$1 \to 3$，第一次 DFS 得到 $f[2] = 3$、$f[3] = 5$、$f[1] = 6$；從 $f$ 最小的 $2$
出發，在 $G$ 上走得到 $1$ 和 $3$，把 $\{1, 2\}$、$\{3\}$ 兩個元件併成一棵樹。 (b)：定理~\ref{thm:con-tarjan}。 (c)：引理~\ref{lem:con-dag}，元
件圖沒有環。 (d)：校園的 $4$、$6$ 都是割點，$\{4, 6\}$ 在環 $4, 6, 8$ 上，不是橋。 (e)：設
$\{u, v\}$ 是橋，拿掉它後分成含 $u$ 的一塊和含 $v$ 的一塊。 頂點至少 $3$ 個，其中一塊至
少有 $2$ 個頂點，設是 $u$ 那塊；這塊連通，所以 $u$ 在塊裡有鄰居 $x$。 拿掉 $u$ 後 $x$ 和
$v$ 不再相連，$u$ 是割點。 (f)：非根而沒有子節點，定理~\ref{thm:con-nonroot} 的條件不可
能成立；根沒有子節點就是孤立點，也不是割點。 (g)：校園的 $\mathit{low}[4] = d[6] = 4$，
但 $\{6, 4\}$ 不是橋；橋要 $\mathit{low}[v] > d[u]$。

\item (a) 第一次 DFS：$a \to d \to e$；$e \to a$ 是後向邊；$e \to b$，$b \to d$、$b \to e$
都是後向邊，$b$ 完成；$e \to g$，$g$ 沒有出邊；$e$、$d$ 完成。 回到 $a$，$a \to e$ 的
$e$ 已是黑色、$d[e] > d[a]$，是前向邊；$a \to h$，$h \to d$、$h \to g$ 是交叉邊。 主迴
圈接著到 $c$：$c \to e$ 是交叉邊，$c \to f$，$f \to c$ 是後向邊。
\[
\begin{array}{c|cccccccc}
v & a & b & c & d & e & f & g & h \\
\hline
d/f & 1/12 & 4/5 & 13/16 & 2/9 & 3/8 & 14/15 & 6/7 & 10/11
\end{array}
\]
樹邊：$a \to d$、$d \to e$、$e \to b$、$e \to g$、$a \to h$、$c \to f$；後向邊：$e \to a$、
$b \to d$、$b \to e$、$f \to c$；前向邊：$a \to e$；交叉邊：$h \to d$、$h \to g$、$c \to e$。

(b) 順序是 $c, f, a, h, d, e, g, b$。 $H^{\mathrm{T}}$：$a: e$；$b: e$；$c: f$；$d: a, b,
h$；$e: a, b, c, d$；$f: c$；$g: e, h$；$h: a$。

(c) 從 $c$：$c \to f$，得 $\{c, f\}$。 $f$ 已分好。 從 $a$：$a \to e$；$e$ 的 $a$ 走過，
$e \to b$（$b$ 的 $e$ 走過）；$e$ 的 $c$ 屬於第一棵；$e \to d$，$d$ 的 $a$、$b$ 走過，$d
\to h$，$h$ 的 $a$ 走過。 得 $\{a, b, d, e, h\}$。 $h$、$d$、$e$ 已分好。 從 $g$：$e$、
$h$ 都已分好，得 $\{g\}$。 $b$ 已分好。 三個元件依序是 $\{c, f\}$、$\{a, b, d, e, h\}$、
$\{g\}$。

\item (a) 發現順序是 $a, d, e, b, g, h, c, f$。
\[
\begin{array}{c|cccccccc}
v & a & b & c & d & e & f & g & h \\
\hline
d & 1 & 4 & 7 & 2 & 3 & 8 & 5 & 6 \\
\mathit{low} & 1 & 2 & 7 & 1 & 1 & 7 & 5 & 2
\end{array}
\]
關鍵幾步：$e \to a$ 讓 $\mathit{low}[e] = 1$。 $b \to d$ 讓 $\mathit{low}[b] = 2$，$b \to e$
是 $\min(2, 3)$，不變。 $g$ 沒有出邊，$\mathit{low}[g] = d[g]$，第一個彈出 $\{g\}$。 回到
$a$，$a \to e$ 是前向邊，$e$ 在堆疊上，$\min(1, 3)$ 不變；$h \to d$ 讓 $\mathit{low}[h]
= 2$，$h \to g$ 的 $g$ 已經不在堆疊上，略過。 $a$ 完成時彈出 $h, b, e, d, a$。 最後 $c \to
e$ 略過，$f \to c$ 讓 $\mathit{low}[f] = 7$，$c$ 完成時彈出 $f, c$。 彈出順序：$\{g\}$、
$\{h, b, e, d, a\}$、$\{f, c\}$，正好是第~\ref{q:con-kosaraju}~題的順序倒過來。

(b) 改用 $\mathit{low}[w]$ 後，$b \to e$ 時 $\mathit{low}[e]$ 已經是 $1$，$\mathit{low}[b]$ 變成
$1$；$h \to d$ 時 $\mathit{low}[d] = 1$，$\mathit{low}[h]$ 變成 $1$。 其他頂點不變。 $b$、$h$
本來就不是頭，改變後仍然 $\mathit{low} < d$，輸出的元件和順序都不變。

\item (a) 元件圖有三個點：$Y = \{c, f\}$、$X = \{a, b, d, e, h\}$、$Z = \{g\}$，邊是 $Y \to
X$（來自 $c \to e$）和 $X \to Z$（$e \to g$、$h \to g$ 只算一條）。 拓樸順序 $Y, X, Z$，和
第~\ref{q:con-kosaraju}~題 Kosaraju 產生的順序相同。 (b) 只有一個源點 $Y$，所以從 $c$ 或 $f$ 出發可以到達所有
頂點，其他頂點都不行。 (c) 刪掉 $c \to e$ 後，$Y$ 和其他元件之間沒有邊，$Y$、$X$ 都變成源
點，至少要 $2$ 個起點，例如 $c$ 和 $a$。

\item 令 $\mathit{low}'$ 是改過的值。 Tarjan 演算法的每一步只取決於「$\mathit{low} = d$ 是
否成立」，所以只要證明對每個頂點 $u$，$\mathit{low}'[u] = d[u]$ 若且唯若 $\mathit{low}[u]
= d[u]$，兩種寫法的堆疊變化就完全相同。 對完成的先後做歸納。

$\mathit{low}'[u] \le \mathit{low}[u]$：三種候選值裡，$d[u]$ 一樣；子節點的
$\mathit{low}'[w] \le \mathit{low}[w]$（歸納假設）；堆疊上的 $w$ 用 $\mathit{low}'[w] \le
d[w]$。 所以 $u$ 不是頭時 $\mathit{low}'[u] \le \mathit{low}[u] < d[u]$。

$\mathit{low}'[u] \ge d[h(u)]$：(1) $d[u] \ge d[h(u)]$。 (2) 子節點 $w$ 和 $u$ 同一個元件
時，$\mathit{low}'[w] \ge d[h(w)] = d[h(u)]$；不同元件時，$h(w)$ 是 $w$ 的祖先，若它是 $u$ 或
$u$ 的祖先，則 $h(w) \leadsto u \to w \leadsto h(w)$，$u$ 和 $w$ 同一個元件，矛盾；所以
$h(w) = w$，$\mathit{low}'[w] \ge d[w] > d[u]$。 (3) 邊 $(u, w)$ 而 $w$ 在堆疊上：$h(w)$
還沒完成，又比 $w$ 早發現，所以是灰色，是 $u$ 的祖先或 $u$ 本身；$u \to w \leadsto h(w)
\leadsto u$，$w$ 和 $u$ 同一個元件，$\mathit{low}'[w] \ge d[h(u)]$。 所以 $u$ 是頭時
$d[u] = d[h(u)] \le \mathit{low}'[u] \le d[u]$，$\mathit{low}'[u] = d[u]$。

\item (a) $k(k - 1)/2$。 元件圖是 DAG，排成拓樸順序後每條邊都從前面指向後面，每一對元件最
多一條邊。 這個數做得到：取 $k - 1$ 個單一頂點和一個由其餘頂點繞成環的元件，排成一列，每個
元件都連一條邊到後面所有的元件。 (b) $n$ 條。 強連通時每個頂點都要能到達別的頂點，出度至
少 $1$，所以至少 $n$ 條邊；把所有頂點繞成一個有向環正好 $n$ 條。 (c) DAG 沒有環，任兩個不
同的頂點不會互相可達，每個元件只有一個頂點，每個頂點都是頭，在自己完成時彈出。 所以彈出順
序就是完成的順序，倒過來就是第~22~章用 DFS 完成時間得到的拓樸順序。

\item (a) $a$ 的鄰居 $e, g, i$，先走 $e$；之後依序走到 $g$、$i$、$d$、$h$、$b$、$f$、$c$，
DFS 樹剛好是一條鏈 $a - e - g - i - d - h - b - f - c$。 後向邊有四條：$g$--$a$、$i$--$a$、
$f$--$h$、$h$--$i$（$h$ 的鄰居依序是 $b, d, f, i$，所以 $h$--$i$ 在 $b$ 的子樹做完後才看到）。
\[
\begin{array}{c|ccccccccc}
v & a & b & c & d & e & f & g & h & i \\
\hline
d & 1 & 7 & 9 & 5 & 2 & 8 & 3 & 6 & 4 \\
\mathit{low} & 1 & 6 & 9 & 4 & 1 & 6 & 1 & 4 & 1
\end{array}
\]
(b) 根 $a$ 只有一個子節點 $e$，不是割點。 沿著鏈檢查每條樹邊 $\{u, v\}$：$\{f, c\}$ 有
$\mathit{low}[c] = 9 \ge 8$，$f$ 是割點，而且 $9 > 8$，$\{c, f\}$ 是橋；$\{h, b\}$ 有
$\mathit{low}[b] = 6 \ge 6$，$h$ 是割點；$\{i, d\}$ 有 $\mathit{low}[d] = 4 \ge 4$，$i$ 是
割點；其他樹邊都是 $\mathit{low}[v] < d[u]$。 割點是 $f$、$h$、$i$，橋只有 $\{c, f\}$。
(c) 依序彈出四個區塊：$\{f\text{--}c\}$；$\{f\text{--}h, b\text{--}f, h\text{--}b\}$；
$\{h\text{--}i, d\text{--}h, i\text{--}d\}$；$\{i\text{--}a, g\text{--}i, g\text{--}a,
e\text{--}g, a\text{--}e\}$。

\item 改變的是 $d$ 和 $h$：$i$ 先處理後向邊 $i$--$a$，$\mathit{low}[i]$ 已經是 $1$，之後 $h$
的後向邊 $h$--$i$ 讓 $\mathit{low}[h] = 1$，再傳給 $\mathit{low}[d] = 1$。（$f$ 的後向邊
$f$--$h$ 發生時 $\mathit{low}[h]$ 還是 $6$，所以 $\mathit{low}[f]$ 不變。） 這時 $\mathit{low}[d]
= 1 < d[i] = 4$，割點 $i$ 找不到，只找到 $f$、$h$。 橋仍然只有 $\{c, f\}$，和正確答案相同。

\item 設兩個不同的區塊 $B_1$、$B_2$ 共用兩個頂點 $x \ne y$。 區塊裡任兩條邊在同一個環上，
所以區塊的邊構成一個連通的子圖，可以在 $B_1$ 裡找一條 $x$ 到 $y$ 的路徑 $P_1$，在 $B_2$ 裡找
一條 $P_2$，兩者沒有共同的邊。 沿著 $P_1$ 從 $x$ 出發，令 $z$ 是第一個也在 $P_2$ 上的頂點
（$y$ 就在 $P_2$ 上，所以 $z$ 存在）。 $P_1$ 從 $x$ 到 $z$ 的一段和 $P_2$ 從 $x$ 到 $z$ 的一段
只在兩端相交；兩段都至少一條邊，而且不可能都只有一條邊（那會是 $x$、$z$ 之間的兩條邊，
簡單圖不行），所以合起來是一個簡單環。 這個環同時有 $B_1$ 和 $B_2$ 的邊，它們應該在同一個區
塊，矛盾。

再設 $v$ 是共用的頂點，$\{v, a\} \in B_1$、$\{v, b\} \in B_2$，$a \ne b$（簡單圖）。 如果拿
掉 $v$ 後 $a$、$b$ 仍有一條路徑 $Q$ 相連，$v, a, Q, b, v$ 就是經過這兩條邊的簡單環，矛
盾。 所以拿掉 $v$ 後 $a$、$b$ 分開了，而原本它們經過 $v$ 相連，$v$ 所在的元件被切成兩塊以
上，$v$ 是割點。

\item (a) 割點只有 $2$（拿掉它，$1$ 就孤立了）；沒有橋：$1$--$2$ 有兩條，拿掉一條還有另一
條，$2, 3, 4$ 是三角形。 (b) DFS：$d[1] = 1$、$d[2] = 2$、$d[3] = 3$、$d[4] = 4$。 $2$ 的
鄰居裡兩個 $1$ 都等於父節點，都被跳過；$4$ 的後向邊 $4$--$2$ 讓 $\mathit{low}[4] =
\mathit{low}[3] = 2$，$\mathit{low}[2] = 2$。 回到 $2$ 時 $\mathit{low}[3] = 2 \ge d[2]$，$2$
是割點（對）；回到 $1$ 時 $\mathit{low}[2] = 2 > d[1] = 1$，輸出橋 $\{1, 2\}$（錯）。 根 $1$
只有一個子節點，不是割點（對）。 (c) 每條邊編號，\proc{Cut-Visit} 改傳走下來的邊的編號
$e$，第 19 行改成「這條邊不是 $e$」。 第二條 $1$--$2$ 就會被當成後向邊，$\mathit{low}[2] =
d[1] = 1$，不再滿足 $> d[1]$。

\item (a) $5$ 個。 例子：兩個三角形 $\{1, 2, 3\}$、$\{3, 4, 5\}$ 共用頂點 $3$；$3$ 是割點，每條
邊都在三角形上，沒有橋。 不能更少：設 $v$ 是割點，$G - v$ 至少有兩個連通元件。 如果某個元件
只有一個頂點 $x$，因為 $G$ 連通，$x$ 只能接到 $v$，而且簡單圖只有一條邊 $\{x, v\}$，這條邊
就是橋。 所以每個元件至少 $2$ 個頂點，總共至少 $1 + 2 + 2 = 5$ 個。 (b) 兩個頂點、一條邊：
這條邊是橋；拿掉任一個頂點，剩下一個頂點，元件個數仍是 $1$，所以沒有割點。
\end{enumerate}

\chapter{最小生成樹}
\label{ch:mst}

% 本章的圖共用的樣式（前綴 mst-）。 追蹤表、圖上的粗線和程式輸出都由模擬程式照虛擬碼
% 逐步執行後產生，再和本章的 C 程式輸出對過。
\tikzset{
    mst-v/.style={circle, draw=black!70, fill=white, minimum size=0.62cm, inner sep=0pt, font=\small},
    mst-e/.style={draw=black!50, thick},
    mst-t/.style={draw=black, line width=1.6pt},
    mst-x/.style={draw=black!35, densely dashed},
    mst-w/.style={font=\scriptsize, fill=white, inner sep=1pt},
    mst-cross/.style={draw=memStroke, line width=1.4pt},
    mst-light/.style={draw=optStroke, line width=2.2pt},
    mst-new/.style={draw=optStroke, line width=1.6pt},
    mst-side/.style={rounded corners=8pt, fill=coreFill, draw=coreStroke!50},
    mst-hull/.style={line width=0.95cm, line cap=round, line join=round, draw=coreFill, fill=coreFill},
    mst-lab/.style={font=\scriptsize, text=black!80},
    mst-pick/.style={-{Stealth[length=2mm]}, draw=optStroke, thick, shorten >=1pt},
}
% Prim 追蹤表：框起來的是這一回合取出的頂點，底色是這一回合改過的 key。
\newcommand{\mstsel}[1]{\fbox{#1}}
\newcommand{\mstchg}[1]{\colorbox{memFill}{#1}}

山上的露營區有八個營地，管理處打算拉電線，讓每個營地都有電。 電只要能從一個營地經過
別的營地送過來就好，不需要每兩個營地之間都有線。 工程行估了價：有些營地之間地形平
緩，拉線便宜；有些要跨過溪谷，貴很多；還有些營地之間根本拉不過去。 管理處想用最少
的錢，讓八個營地全部連在一起。 往下讀之前，先想想：從最便宜的那段開始拉，一段一段
加上去，最後的總價會是最少的嗎？ 如果不是每段都要，哪些段可以放心跳過？

這一章處理的就是這個問題：在一張每條邊都有權重的連通無向圖裡，挑出一些邊把所有頂點
連起來，而且總權重最小。 挑出來的一定是一棵生成樹，叫做\textbf{最小生成樹}。 我們
先證明兩個性質：\textbf{切割性質}告訴我們哪些邊一定可以放心選，\textbf{環性質}告訴我們
哪些邊可以放心不選。 三個經典演算法都是切割性質的應用，差別只在每次看哪一個切割：
Kruskal 演算法把邊由輕到重排好，一條一條試；Prim 演算法從一個頂點開始，把樹往外長；
Sollin 演算法（也叫 Borůvka 演算法）讓每個元件同時挑一條最輕的邊。 最後談幾個考試常
問的觀念：權重相同時怎麼處理、最小生成樹什麼時候唯一、次小生成樹，以及最小生成樹和最
短路徑樹為什麼不一樣。

這一章用到前面的這些東西：第~21~章的生成樹、樹的等價性質（定理~\ref{thm:gtr-tree}：$n$ 個頂點的樹恰好有
$n - 1$ 條邊、任兩點之間恰好一條簡單路徑）和鄰接串列；第~13~章的互斥集合；第~8~章的帶索
引最小堆積和 \proc{Decrease-Key}；第~19~章的貪心法和交換論證。 程式方面用到《C 程式
設計基礎》第~14~章的結構、第~19~章的 \texttt{malloc} 和鏈結串列，以及標準函式庫的
\texttt{qsort}。

本章的約定和第~21~章相同：圖是\textbf{連通的簡單無向圖}，有 $n$ 個頂點、$m$ 條邊，頂點
編號 $1$ 到 $n$（觀念題有時用英文字母）；每條邊 $(u, v)$ 有一個權重 $w(u, v)$，可以是
任何實數，包括 $0$ 和負數。 鄰接串列裡的鄰居由小到大排。 C 程式裡的頂點從 $0$ 編號，讀
入時減一、印出時加一。

\section{最小生成樹與兩個性質}
\label{sec:mst-props}

\subsection{問題}

對一個邊的集合 $A$，令 $w(A) = \sum_{(u, v) \in A} w(u, v)$，叫做 $A$ 的\textbf{總權
重}。

\begin{definition}
$G = (V, E)$ 是連通的無向圖，每條邊有權重 $w$。 $G$ 的生成樹 $T$ 裡，總權重 $w(T)$ 最
小的叫做 $G$ 的\textbf{最小生成樹}（minimum spanning tree，簡稱 MST）。 $G$ 不連通時，
每個連通元件各取一棵最小生成樹，合起來叫\textbf{最小生成森林}。
\end{definition}

由第~21~章樹的等價性質（定理~\ref{thm:gtr-tree}），每一棵生成樹都恰好有 $n - 1$ 條邊，所以「總權重最小」比的是同
樣多條邊的總和。 最小生成樹不一定只有一棵：權重相同的邊可能互相替換，總權重一樣。 「最
小」指的是總權重，不是邊數，也不是最重的那條邊。

\begin{center}
\begin{tikzpicture}[x=1.2cm, y=1.2cm]
\node[mst-v] (v1) at (0,2) {$1$};
\node[mst-v] (v2) at (2,2) {$2$};
\node[mst-v] (v3) at (4,2) {$3$};
\node[mst-v] (v4) at (6,2) {$4$};
\node[mst-v] (v5) at (0,0) {$5$};
\node[mst-v] (v6) at (2,0) {$6$};
\node[mst-v] (v7) at (4,0) {$7$};
\node[mst-v] (v8) at (6,0) {$8$};
\draw[mst-e] (v1) -- node[mst-w] {7} (v2);
\draw[mst-e] (v2) -- node[mst-w] {12} (v3);
\draw[mst-e] (v3) -- node[mst-w] {14} (v4);
\draw[mst-e] (v5) -- node[mst-w] {9} (v6);
\draw[mst-e] (v6) -- node[mst-w] {4} (v7);
\draw[mst-e] (v7) -- node[mst-w] {13} (v8);
\draw[mst-e] (v1) -- node[mst-w] {1} (v5);
\draw[mst-e] (v2) -- node[mst-w] {6} (v6);
\draw[mst-e] (v3) -- node[mst-w] {11} (v7);
\draw[mst-e] (v4) -- node[mst-w] {3} (v8);
\draw[mst-e] (v1) -- node[mst-w] {10} (v6);
\draw[mst-e] (v2) -- node[mst-w] {5} (v7);
\draw[mst-e] (v3) -- node[mst-w] {8} (v8);
\end{tikzpicture}
\diagcaption{露營區的八個營地。 每條邊是一段拉得過去的電線，數字是費用（萬元）。 本章的
Kruskal、Prim、Sollin 三個演算法都在這張圖上追蹤。}
\label{fig:mst-camp}
\end{center}

圖~\ref{fig:mst-camp} 是章首的露營區，記成 $G$。 它有 $13$ 條邊，每棵生成樹要挑其中 $7$
條；用程式列舉，一共有 $377$ 棵生成樹。 一棵一棵算總權重太慢，需要一套挑邊的規則，而且
要能證明它不會挑錯。

最直接的想法是貪心：從最便宜的邊開始拉，會形成環的就跳過。 $1$--$5$（$1$）、$4$--$8$
（$3$）、$6$--$7$（$4$）、$2$--$7$（$5$）都可以拉；接下來 $2$--$6$（$6$）的兩端已經經過
$7$ 連在一起，拉了只會多花錢……這就是第~\ref{sec:mst-kruskal}~節的 Kruskal 演算法。 它
看起來合理，但第~19~章說過，貪心法每一步都要證明不會錯。 證明的工具是下面的切割性質。

\subsection{切割}

\begin{definition}
把頂點分成兩個都不是空的集合 $S$ 和 $V - S$，叫做一個\textbf{切割}（cut），寫成 $(S, V -
S)$。 一條邊的兩端一個在 $S$、一個在 $V - S$，就說這條邊\textbf{跨越}這個切割。 跨越
切割的邊裡，權重最小的叫做這個切割的\textbf{輕邊}（light edge）；權重最小的邊有好幾條
時，每一條都是輕邊。 如果邊集合 $A$ 裡沒有任何一條邊跨越切割，就說這個切割\textbf{尊
重} $A$。
\end{definition}

\begin{center}
\begin{tikzpicture}[x=1.2cm, y=1.2cm]
\node[mst-v] (v1) at (0,2) {$1$};
\node[mst-v] (v2) at (2,2) {$2$};
\node[mst-v] (v3) at (4,2) {$3$};
\node[mst-v] (v4) at (6,2) {$4$};
\node[mst-v] (v5) at (0,0) {$5$};
\node[mst-v] (v6) at (2,0) {$6$};
\node[mst-v] (v7) at (4,0) {$7$};
\node[mst-v] (v8) at (6,0) {$8$};
\begin{scope}[on background layer]
\node[mst-side, fit=(v1)(v2)(v5)(v6), inner sep=7pt] (S) {};
\end{scope}
\draw[mst-e] (v1) -- node[mst-w] {7} (v2);
\draw[mst-cross] (v2) -- node[mst-w] {12} (v3);
\draw[mst-e] (v3) -- node[mst-w] {14} (v4);
\draw[mst-e] (v5) -- node[mst-w] {9} (v6);
\draw[mst-light] (v6) -- node[mst-w] {4} (v7);
\draw[mst-e] (v7) -- node[mst-w] {13} (v8);
\draw[mst-e] (v1) -- node[mst-w] {1} (v5);
\draw[mst-e] (v2) -- node[mst-w] {6} (v6);
\draw[mst-e] (v3) -- node[mst-w] {11} (v7);
\draw[mst-e] (v4) -- node[mst-w] {3} (v8);
\draw[mst-e] (v1) -- node[mst-w] {10} (v6);
\draw[mst-cross] (v2) -- node[mst-w] {5} (v7);
\draw[mst-e] (v3) -- node[mst-w] {8} (v8);
\node[mst-lab, anchor=south] at (S.north) {$S$};
\node[mst-lab, anchor=south] at ($(v3)!0.5!(v4)+(0,0.45)$) {$V - S$};
\end{tikzpicture}
\diagcaption{切割 $(S, V - S)$，$S = \{1, 2, 5, 6\}$。 跨越它的邊有三條（橘色）：
$6$--$7$（$4$）、$2$--$7$（$5$）、$2$--$3$（$12$）。 輕邊是 $6$--$7$（紅色）。}
\label{fig:mst-cut}
\end{center}

圖~\ref{fig:mst-cut} 是一個例子。 不管最後的生成樹長什麼樣子，它一定要用到至少一條跨越
這個切割的邊，否則 $S$ 和 $V - S$ 就不相連了。 直覺上，既然一定要過去，就挑最便宜的那條
過去。 下面的定理說這個直覺是對的，而且可以在「已經選了一些邊」的情況下繼續用。

\subsection{切割性質}

在證明之前先定一個說法。 $A$ 是一些邊的集合，而且 $A$ 包含在某一棵最小生成樹裡。 如果
再加一條邊 $(u, v)$，$A \cup \{(u, v)\}$ 仍然包含在某一棵最小生成樹裡，就說 $(u, v)$
對 $A$ 是\textbf{安全的}（safe）。 從 $A = \emptyset$ 開始，每次加一條安全的邊，加了 $n
- 1$ 條之後，$A$ 本身就是一棵最小生成樹。

\begin{theorem}[切割性質]
\label{thm:mst-cut}
$G$ 是連通的加權無向圖，邊集合 $A$ 包含在 $G$ 的某一棵最小生成樹裡。 $(S, V - S)$ 是尊
重 $A$ 的任何一個切割，$(u, v)$ 是這個切割的一條輕邊。 那麼 $(u, v)$ 對 $A$ 是安全的。
\end{theorem}

\begin{proof}
令 $T$ 是包含 $A$ 的一棵最小生成樹。 如果 $(u, v) \in T$，$T$ 本身就包含 $A \cup \{(u,
v)\}$，得證。

否則把 $(u, v)$ 加進 $T$。 $T$ 裡 $u$ 到 $v$ 本來就有一條簡單路徑 $P$，加上 $(u, v)$ 就
形成一個環。 $u$ 和 $v$ 一個在 $S$、一個在 $V - S$，所以 $P$ 從 $S$ 出發、走到 $V - S$，
途中至少有一條邊 $(x, y)$ 跨越這個切割。 這條邊有三個性質：
\begin{itemize}
\item $(x, y) \notin A$，因為切割尊重 $A$，$A$ 裡沒有跨越切割的邊；
\item $w(u, v) \le w(x, y)$，因為 $(u, v)$ 是輕邊，而 $(x, y)$ 也跨越切割；
\item $(x, y) \ne (u, v)$，因為 $(x, y)$ 在 $T$ 裡，$(u, v)$ 不在。
\end{itemize}
令 $T' = T - \{(x, y)\} \cup \{(u, v)\}$。 從 $T$ 拿掉 $(x, y)$，樹分成兩塊，一塊含 $x$、
一塊含 $y$；$u$、$v$ 分別在這兩塊（路徑 $P$ 被切斷的兩邊），加上 $(u, v)$ 又連成一棵
樹，所以 $T'$ 是生成樹。 它的總權重
\[
w(T') = w(T) - w(x, y) + w(u, v) \le w(T),
\]
而 $T$ 已經是最小的，所以 $w(T') = w(T)$，$T'$ 也是最小生成樹。 $A \subseteq T$ 而且
$(x, y) \notin A$，所以 $A \subseteq T'$，再加上 $(u, v) \in T'$，得證。
\end{proof}

證明裡的動作「加一條邊形成環，再從環上拿掉另一條」叫做\textbf{換邊}，本章會一直用到。
它是第~19~章交換論證在生成樹上的版本：換完之後要說明兩件事，結果仍然是生成樹，而且總權
重沒有變大。

取 $A = \emptyset$，任何切割都尊重它，得到一個常用的特例。

\begin{corollary}
\label{cor:mst-cut-strict}
任何一個切割的輕邊，都屬於某一棵最小生成樹。 如果一個切割的輕邊只有一條（權重嚴格最
小），它屬於\textbf{每一棵}最小生成樹。
\end{corollary}

\begin{proof}
第一句是定理~\ref{thm:mst-cut} 取 $A = \emptyset$。 第二句：假設某棵最小生成樹 $T$ 不含
這條唯一的輕邊 $(u, v)$，照上面的證明換邊，換掉的 $(x, y)$ 跨越切割、又不是 $(u, v)$，
所以 $w(x, y) > w(u, v)$，$w(T') < w(T)$，和 $T$ 最小矛盾。
\end{proof}

在圖~\ref{fig:mst-cut} 裡，$6$--$7$ 是唯一的輕邊，所以每一棵最小生成樹都有 $6$--$7$。 把
$S$ 換成 $\{v\}$（一個頂點自己一邊），就得到：\textbf{每個頂點連出去的邊裡最輕的那一條，
一定在某一棵最小生成樹裡}。 第~\ref{sec:mst-sollin}~節的 Sollin 演算法第一回合就是靠這
件事。

\subsection{環性質}

切割性質說哪些邊可以選，環性質說哪些邊可以不選。

\begin{theorem}[環性質]
\label{thm:mst-cycle}
$C$ 是 $G$ 裡的一個環，$e$ 是 $C$ 上權重最大的一條邊（$C$ 上其他邊的權重都 $\le w(e)$）。
那麼 $G$ 有一棵最小生成樹不含 $e$。 如果 $e$ 是 $C$ 上\textbf{唯一}權重最大的邊，$e$ 不在
任何一棵最小生成樹裡。
\end{theorem}

\begin{proof}
令 $T$ 是一棵最小生成樹，$e = (u, v)$。 如果 $e \notin T$，第一句已經成立。 設 $e \in
T$。 從 $T$ 拿掉 $e$，樹分成兩塊，令 $S$ 是含 $u$ 的那一塊的頂點，$v \in V - S$。 環 $C$
拿掉 $e$ 之後是一條從 $u$ 到 $v$ 的路徑，這條路徑從 $S$ 走到 $V - S$，一定有一條邊 $f$ 跨越
$(S, V - S)$，而且 $f \ne e$。 $T$ 裡跨越這個切割的邊只有 $e$（$T - \{e\}$ 的兩塊之間沒
有別的邊），所以 $f \notin T$。

$T' = T - \{e\} \cup \{f\}$ 又把兩塊連起來，是一棵生成樹，$w(T') = w(T) - w(e) + w(f)$。
$f$ 也在環 $C$ 上，$w(f) \le w(e)$，所以 $w(T') \le w(T)$，$T'$ 是不含 $e$ 的最小生成樹。
如果 $e$ 是唯一最大的，$w(f) < w(e)$，$w(T') < w(T)$，和 $T$ 最小矛盾，所以 $T$ 根本不能
含 $e$。
\end{proof}

在圖~\ref{fig:mst-camp} 裡，環 $2 \to 6 \to 7 \to 2$ 的三條邊是 $6$、$4$、$5$，$2$--$6$
（$6$）唯一最重，所以 $2$--$6$ 不在任何最小生成樹裡。 環性質也可以反過來用：依權重由大
到小看每條邊，拿掉之後圖還連通（也就是它在某個環上，而且是那個環上最重的）就拿掉。 這
個做法叫\textbf{反向刪除}（reverse-delete），結果一樣是最小生成樹，但每次都要檢查連通，
比本章的三個演算法慢。 有些中文教材叫它「破圈法」。

\subsection{通用的做法}

三個演算法都套用同一個骨架：維護一個邊的集合 $A$，保持「$A$ 包含在某一棵最小生成樹
裡」，每次加一條安全的邊。

\begin{pseudo}[label=alg:mst-generic]{通用的最小生成樹做法}
\Procedure{Generic-MST}{$G, w$}
    \State $A \gets \emptyset$
    \While{$A$ 還不是一棵生成樹}
        \State 找一條對 $A$ 安全的邊 $(u, v)$
        \State $A \gets A \cup \{(u, v)\}$
    \EndWhile
    \State \Return $A$
\EndProcedure
\end{pseudo}

$A$ 在過程中一直是一個森林（沒有環），它的每個連通元件是一棵樹。 第 4 行怎麼找，就是
三個演算法的差別：

\begin{center}
\small
\begin{tabular}{lll}
\toprule
演算法 & $A$ 的樣子 & 用哪個切割 \\
\midrule
Kruskal & 好幾棵樹組成的森林 & 下一條要加的邊的一端所在的那棵樹 vs. 其他 \\
Prim & 從起點長出來的一棵樹 & 這棵樹 vs. 其他頂點 \\
Sollin & 好幾棵樹組成的森林 & 每棵樹各自 vs. 其他，同時進行 \\
\bottomrule
\end{tabular}
\end{center}

第~\ref{sec:gr-matroid}~節介紹過擬陣。 圖的邊、以「不含環」為獨立的條件，構成一個擬陣（圖擬陣），所以依權
重由小到大的貪心一定對，這是 Kruskal 演算法正確的另一種說法。 本章用切割性質證明，因為
它同時涵蓋三個演算法。

\subsection{權重互不相同時，最小生成樹唯一}

\begin{theorem}
\label{thm:mst-unique}
如果 $G$ 每條邊的權重都不相同，$G$ 的最小生成樹只有一棵。
\end{theorem}

\begin{proof}
假設 $T_1 \ne T_2$ 都是最小生成樹。 兩者都有 $n - 1$ 條邊，所以兩邊各有對方沒有的邊。 在
所有「恰好屬於其中一棵」的邊裡，取權重最小的一條 $e$；權重互不相同，$e$ 是唯一的。 不
妨設 $e \in T_1$、$e \notin T_2$。

把 $e$ 加進 $T_2$ 形成一個環。 這個環上一定有一條邊 $f$ 不在 $T_1$ 裡，否則整個環都在
$T_1$ 裡，和 $T_1$ 是樹矛盾。 $f$ 在 $T_2$ 裡而不在 $T_1$ 裡，也是「恰好屬於其中一棵」的
邊，又 $f \ne e$，由 $e$ 的取法，$w(f) > w(e)$。 和切割性質的證明一樣，$T_2 - \{f\} \cup
\{e\}$ 是一棵生成樹，總權重 $w(T_2) - w(f) + w(e) < w(T_2)$，和 $T_2$ 最小矛盾。
\end{proof}

反過來不成立：權重有相同的，最小生成樹仍然可能唯一。 例如三角形的三條邊權重是 $1$、$1$、
$2$，最小生成樹只有 $\{1, 1\}$ 那一棵（權重 $2$ 的邊是環上唯一最重的，由環性質一定不選）。
更極端的例子是 $G$ 本身就是一棵樹，不管權重是多少，最小生成樹都是它自己。
第~\ref{sec:mst-more}~節會給一個「最小生成樹唯一」的充分必要條件。

另外，最小生成樹不唯一時，每一棵的總權重當然相同。 更強的事也成立（這裡不證）：把每一
棵的邊權重由小到大排成一列，得到的數列都一樣。 例如某張圖的一棵最小生成樹用了兩段 $3$ 萬、
一段 $5$ 萬，另一棵最小生成樹也一定是兩段 $3$ 萬、一段 $5$ 萬，只是選的位置不同。

\subsection{權重相同時：平手規則}
\label{sec:mst-ties}

考試和程式裡的圖常常有權重相同的邊。 這時最小生成樹可能不只一棵，演算法選哪一棵取決於
平手時怎麼排。 本書的約定是：每條邊寫成 $(u, v)$ 且 $u < v$，兩條邊比較時

\begin{enumerate}
\item 權重小的比較「輕」；
\item 權重相同時，$u$（小的端點）小的比較輕；
\item 權重和 $u$ 都相同時，$v$（大的端點）小的比較輕。
\end{enumerate}

用字母當頂點時，照字母順序。 這樣任兩條不同的邊都分得出輕重，寫成 $e \prec f$（$e$ 比 $f$
輕）。 Kruskal 演算法排序和 Sollin 演算法挑邊都照 $\prec$。 Prim 演算法比的是頂點的
$\mathit{key}$，它的平手規則另外在第~\ref{sec:mst-prim}~節說明。

\begin{lemma}
\label{lem:mst-ties}
照 $\prec$ 比輕重，就好像所有權重都不相同：在 $\prec$ 之下最輕的生成樹 $T_{\prec}$ 只
有一棵，而且 $T_{\prec}$ 是 $G$（原本的權重）的一棵最小生成樹。
\end{lemma}

\begin{whybox}[引理~\ref{lem:mst-ties}]
把邊照 $\prec$ 由輕到重排成 $e_1, e_2, \ldots, e_m$，令 $r(e_i) = i$。 取一個很小的正
數 $\varepsilon$，定義新的權重
\[
w'(e) = w(e) + \varepsilon \cdot r(e).
\]
先看權重都是整數的情況（考試的圖幾乎都是），取 $\varepsilon = 1/(mn)$。
\begin{itemize}
\item \textbf{$w'$ 互不相同，而且大小順序就是 $\prec$}：$w(e) < w(f)$ 時 $w(f) - w(e) \ge
1$，而 $\varepsilon\,|r(e) - r(f)| < \varepsilon m < 1$，所以 $w'(e) < w'(f)$；$w(e) = w(f)$
時由 $r$ 決定。 所以照 $\prec$ 比較，等於照 $w'$ 比較。
\item \textbf{$w'$ 的最小生成樹也是 $w$ 的最小生成樹}：任何一棵生成樹 $T$ 有 $n - 1$ 條
邊，$0 < w'(T) - w(T) = \varepsilon \sum_{e \in T} r(e) \le \varepsilon \cdot (n - 1) m < 1$。
設 $T'$ 是 $w'$ 的最小生成樹，$T$ 是任何一棵生成樹。 如果 $w(T) < w(T')$，因為都是整數，
$w(T) \le w(T') - 1$，於是 $w'(T) < w(T) + 1 \le w(T') < w'(T')$，和 $T'$ 對 $w'$ 最小矛
盾。 所以 $w(T') \le w(T)$。
\end{itemize}
$w'$ 的權重互不相同，由定理~\ref{thm:mst-unique}，它的最小生成樹只有一棵，就是 $T_{\prec}$。

權重是一般實數時，把上面的「$1$」換成「兩條邊權重的差」和「兩棵生成樹總權重的差」裡最
小的正數（有限多個，一定有最小值），$\varepsilon$ 取得比這個數除以 $mn$ 還小，論證完全
相同。
\end{whybox}

所以只要所有比較都照 $\prec$ 做，演算法就像在一張權重互不相同的圖上執行，結果一定是
$T_{\prec}$ 這一棵。 第~\ref{sec:mst-kruskal}~節的 Kruskal 演算法和第~\ref{sec:mst-sollin}~節
的 Sollin 演算法都是這樣，對同一張圖會選出同一棵樹。

\begin{examnote}
權重相同時，「最小生成樹是哪一棵」的答案不只一個，所以追蹤題通常會規定平手時怎麼排（例
如「權重相同時先選頂點編號小的邊」「權重相同時依字母順序」「Prim 演算法權重相同時先加入
編號小的頂點」）。 不同課本的約定不同，照題目寫；題目沒規定時，在答案開頭寫明自己用的規
則，例如本書的「先比權重，再比小端點，再比大端點」。 問總權重時答案只有一個，不受平手規
則影響。
\end{examnote}

\section{Kruskal 演算法}
\label{sec:mst-kruskal}

Kruskal 演算法就是第~\ref{sec:mst-props}~節開頭那個貪心：把邊依權重由小到大排好，一條一
條看；兩端已經連通的邊會形成環，跳過，否則選進來。 「兩端是不是已經連通」用第~13~章的互
斥集合判斷：每個集合是目前森林裡一棵樹的頂點，\proc{Find-Set} 找代表，\proc{Union} 把兩
棵樹併起來。

\begin{pseudo}[label=alg:mst-kruskal]{Kruskal 演算法}
\Procedure{MST-Kruskal}{$G, w$}
    \State $A \gets \emptyset$
    \For{每個頂點 $v \in G.V$}
        \State \Call{Make-Set}{$v$}
    \EndFor
    \State 把 $G.E$ 的邊照 $\prec$ 由輕到重排序 \Comment{第~\ref{sec:mst-ties}~節的平手規則}
    \For{依排好的順序，每條邊 $(u, v)$}
        \If{$\Call{Find-Set}{u} \ne \Call{Find-Set}{v}$}
            \State $A \gets A \cup \{(u, v)\}$
            \State \Call{Union}{$u, v$}
        \EndIf
    \EndFor
    \State \Return $A$
\EndProcedure
\end{pseudo}

第 7 行的判斷就是第~21~章引理~\ref{lem:gtr-count}（邊數和連通元件數）的兩種情形：兩端已經在同一個集合，
加進來會形成環；兩端在不同集合，加進來會把兩棵樹接成一棵，集合數減一。 選滿 $n - 1$ 條
邊時已經是一棵生成樹，後面的邊一定都會跳過，可以提早結束迴圈。

\begin{theorem}
\label{thm:mst-kruskal}
$G$ 連通時，\proc{MST-Kruskal} 回傳一棵最小生成樹。 而且只要排序時平手照 $\prec$，回傳
的就是引理~\ref{lem:mst-ties} 的 $T_{\prec}$。
\end{theorem}

\begin{proof}
證明第 8 行每次加進來的邊對 $A$ 都是安全的。 設第 8 行要加 $(u, v)$，這時 $u$ 和 $v$ 在 $A$
的不同的樹裡。 令 $S$ 是 $u$ 所在那棵樹的頂點，$v \in V - S$。

\begin{itemize}
\item 切割 $(S, V - S)$ 尊重 $A$：$S$ 是 $A$ 的一個連通元件，$A$ 的邊不會一端在裡面、一
端在外面。
\item $(u, v)$ 跨越這個切割，而且是輕邊。 比 $(u, v)$ 早處理的每條邊，處理時要嘛被加進
$A$，要嘛兩端已經在同一棵樹裡；森林只會越併越大，所以現在它的兩端仍在同一棵樹裡，也就是
不跨越 $(S, V - S)$。 因此跨越切割的邊都排在 $(u, v)$ 或它後面，權重都 $\ge w(u, v)$。
\end{itemize}

由定理~\ref{thm:mst-cut}，$(u, v)$ 是安全的。 一開始 $A = \emptyset$ 包含在任何一棵最小生成
樹裡，歸納下去，$A$ 始終包含在某一棵最小生成樹裡。 最後 $A$ 是一棵生成樹：如果 $A$ 還有兩
棵以上的樹，$G$ 連通，一定有一條邊連接兩棵不同的樹，它被處理時兩端在不同集合，會被選進
來，矛盾。 一棵生成樹包含在某一棵最小生成樹裡，邊數又一樣，兩者相等。

照 $\prec$ 排序時，演算法的每個判斷都和「在 $w'$ 上執行」相同，而 $w'$ 的最小生成樹只有
$T_{\prec}$ 一棵（引理~\ref{lem:mst-ties}），所以回傳的是 $T_{\prec}$。
\end{proof}

圖不連通時，演算法照樣跑完，得到的是最小生成森林：證明裡的每一步都沒有用到連通，只有
最後「只剩一棵樹」那一步要用。 選到的邊數是 $n$ 減去連通元件數。

\subsection{手算追蹤}

在圖~\ref{fig:mst-camp} 的露營區上跑一次。 $13$ 條邊的權重都不同，排序不用平手規則。
下表一列是一條邊；最後一欄是處理完這條邊之後，互斥集合裡有兩個頂點以上的集合（其他頂點
各自一組）。

\begin{center}
\small
\begin{tabular}{rcrll}
\toprule
步 & 邊 & 權重 & 結果 & 處理完的集合 \\
\midrule
1 & $1$--$5$ & 1 & 選 & $\{1, 5\}$ \\
2 & $4$--$8$ & 3 & 選 & $\{1, 5\}$、$\{4, 8\}$ \\
3 & $6$--$7$ & 4 & 選 & $\{1, 5\}$、$\{4, 8\}$、$\{6, 7\}$ \\
4 & $2$--$7$ & 5 & 選 & $\{1, 5\}$、$\{2, 6, 7\}$、$\{4, 8\}$ \\
5 & $2$--$6$ & 6 & 略過（成環） & $\{1, 5\}$、$\{2, 6, 7\}$、$\{4, 8\}$ \\
6 & $1$--$2$ & 7 & 選 & $\{1, 2, 5, 6, 7\}$、$\{4, 8\}$ \\
7 & $3$--$8$ & 8 & 選 & $\{1, 2, 5, 6, 7\}$、$\{3, 4, 8\}$ \\
8 & $5$--$6$ & 9 & 略過（成環） & $\{1, 2, 5, 6, 7\}$、$\{3, 4, 8\}$ \\
9 & $1$--$6$ & 10 & 略過（成環） & $\{1, 2, 5, 6, 7\}$、$\{3, 4, 8\}$ \\
10 & $3$--$7$ & 11 & 選 & $\{1, 2, 3, 4, 5, 6, 7, 8\}$ \\
11--13 & \multicolumn{4}{l}{$2$--$3$（$12$）、$7$--$8$（$13$）、$3$--$4$（$14$）：已經選滿 $7$ 條，不用再看} \\
\bottomrule
\end{tabular}
\end{center}

\begin{center}
\begin{tikzpicture}[x=1.2cm, y=1.2cm]
\node[mst-v] (v1) at (0,2) {$1$};
\node[mst-v] (v2) at (2,2) {$2$};
\node[mst-v] (v3) at (4,2) {$3$};
\node[mst-v] (v4) at (6,2) {$4$};
\node[mst-v] (v5) at (0,0) {$5$};
\node[mst-v] (v6) at (2,0) {$6$};
\node[mst-v] (v7) at (4,0) {$7$};
\node[mst-v] (v8) at (6,0) {$8$};
\draw[mst-t] (v1) -- node[mst-w] {7} (v2);
\draw[mst-x] (v2) -- node[mst-w] {12} (v3);
\draw[mst-x] (v3) -- node[mst-w] {14} (v4);
\draw[mst-x] (v5) -- node[mst-w] {9} (v6);
\draw[mst-t] (v6) -- node[mst-w] {4} (v7);
\draw[mst-x] (v7) -- node[mst-w] {13} (v8);
\draw[mst-t] (v1) -- node[mst-w] {1} (v5);
\draw[mst-x] (v2) -- node[mst-w] {6} (v6);
\draw[mst-t] (v3) -- node[mst-w] {11} (v7);
\draw[mst-t] (v4) -- node[mst-w] {3} (v8);
\draw[mst-x] (v1) -- node[mst-w] {10} (v6);
\draw[mst-t] (v2) -- node[mst-w] {5} (v7);
\draw[mst-t] (v3) -- node[mst-w] {8} (v8);
\end{tikzpicture}
\diagcaption{露營區的最小生成樹（粗線），總權重 $1 + 3 + 4 + 5 + 7 + 8 + 11 = 39$。 虛線是沒有
選的邊。}
\label{fig:mst-tree}
\end{center}

幾個要注意的地方：

\begin{itemize}
\item 第 5 步的 $2$--$6$ 雖然只有 $6$ 萬，但 $2$、$6$ 已經經過 $7$ 連通了，拉了只是多一
條備用線。 這正是環性質：環 $2, 6, 7$ 上它最重。
\item 第 10 步的 $3$--$7$ 要 $11$ 萬，比被跳過的 $5$--$6$、$1$--$6$ 都貴，卻一定要選：它是
切割 $(\{1, 2, 5, 6, 7\}, \{3, 4, 8\})$ 的輕邊，跨越這個切割的另外兩條 $2$--$3$、$7$--$8$
更貴。 最小生成樹不是「挑最便宜的 $n - 1$ 條邊」。
\item 結果和第~\ref{sec:mst-props}~節開頭的貪心一模一樣，總共 $39$ 萬。 權重互不相同，由
定理~\ref{thm:mst-unique}，這是唯一的一棵最小生成樹。
\end{itemize}

\subsection{時間}

第 3--4 行做 $n$ 次 \proc{Make-Set}。 第 5 行排序 $m$ 條邊，用第~14~章的合併排序或堆積排
序是 $O(m \lg m)$。 第 6--9 行做 $2m$ 次 \proc{Find-Set} 和 $n - 1$ 次 \proc{Union}，用依
大小（或依秩）合併加路徑壓縮，總共 $O(m\,\alpha(n))$（第~13~章）。 簡單圖 $m < n^2$，所
以 $\lg m < 2 \lg n$，$O(m \lg m) = O(m \lg n)$；連通圖 $m \ge n - 1$。 合起來
\[
O(m \lg m) + O(m\,\alpha(n)) = O(m \lg n),
\]
瓶頸在排序。 如果邊已經依權重排好，或者權重是範圍很小的整數、可以用第~15~章的計數排序，
Kruskal 演算法就幾乎是線性時間。

\subsection{C 程式}

\begin{cbookrecap}[第~14~章結構、第~18~章 18.6 節、第~19~章 \texttt{malloc}]
邊存在一個 \texttt{Edge} 結構的陣列裡，用 \texttt{malloc} 依讀到的 $m$ 配置。
\texttt{qsort(e, m, sizeof *e, cmp\_edge)} 把整個陣列排好；比較函式拿到兩個
\texttt{const void *}，先轉成 \texttt{const Edge *} 再比欄位，第一個要排前面就回傳負數。
這裡依平手規則依序比權重、小端點、大端點，所以讀入時先把每條邊的兩端整理成 \texttt{u <
v}。 寫 \texttt{return x->w - y->w;} 雖然短，但兩個權重相差很大時減法會溢位，所以本章一律
用比較的寫法。
\end{cbookrecap}

下面的程式讀入無向圖（第一列 $n$、$m$，接著 $m$ 列各是 $u\ v\ w$），照
虛擬碼~\ref{alg:mst-kruskal} 執行，每看一條邊印一列，選滿 $n - 1$ 條就停。 互斥集合用第~13~
章的依大小合併和迴圈版路徑壓縮。 自環的兩端相同，\texttt{union\_set} 會回傳 $0$，自然被跳
過。 總權重用 \texttt{long long} 累加，因為 $n - 1$ 條邊的權重加起來可能超過 \texttt{int}。

\begin{example}[label=ex:mst-kruskal]{c}{(kruskal.c) Kruskal 演算法：排序加互斥集合}
#include <stdio.h>
#include <stdlib.h>

typedef struct {
    int u, v, w;    // 讀入時讓 u < v
} Edge;

// 平手規則：權重小的先；權重相同時比較小的端點，再比較大的端點
int cmp_edge(const void *a, const void *b)
{
    const Edge *x = a, *y = b;
    if (x->w != y->w) {
        return x->w < y->w ? -1 : 1;
    }
    if (x->u != y->u) {
        return x->u < y->u ? -1 : 1;
    }
    return (x->v > y->v) - (x->v < y->v);
}

// 第 13 章的互斥集合：依大小合併、路徑壓縮
typedef struct {
    int *p, *size;
} DSU;

int find_set(DSU *s, int x)
{
    int r = x;
    while (s->p[r] != r) {
        r = s->p[r];
    }
    while (s->p[x] != r) {
        int next = s->p[x];
        s->p[x] = r;
        x = next;
    }
    return r;
}

// 已經同一組時回傳 0，真的合併時回傳 1
int union_set(DSU *s, int x, int y)
{
    int rx = find_set(s, x), ry = find_set(s, y);
    if (rx == ry) {
        return 0;
    }
    if (s->size[rx] < s->size[ry]) {
        int t = rx;
        rx = ry;
        ry = t;
    }
    s->p[ry] = rx;
    s->size[rx] += s->size[ry];
    return 1;
}

int main(void)
{
    int n, m;
    if (scanf("%d %d", &n, &m) != 2 || n <= 0 || m < 0) {
        return 1;
    }
    Edge *e = malloc((m > 0 ? m : 1) * sizeof *e);
    DSU s = {malloc(n * sizeof(int)), malloc(n * sizeof(int))};
    int ok = e != NULL && s.p != NULL && s.size != NULL;
    for (int i = 0; ok && i < m; i++) {
        int u, v, w;
        ok = scanf("%d %d %d", &u, &v, &w) == 3 && u >= 1 && u <= n
             && v >= 1 && v <= n;
        if (ok) {
            e[i].u = (u < v ? u : v) - 1;
            e[i].v = (u < v ? v : u) - 1;
            e[i].w = w;
        }
    }
    if (ok) {
        for (int x = 0; x < n; x++) {
            s.p[x] = x;
            s.size[x] = 1;
        }
        qsort(e, m, sizeof *e, cmp_edge);
        long long total = 0;
        int taken = 0;
        for (int i = 0; i < m && taken < n - 1; i++) {
            int take = union_set(&s, e[i].u, e[i].v);
            printf("%d-%d (%d) %s\n", e[i].u + 1, e[i].v + 1, e[i].w,
                   take ? "take" : "skip");
            if (take) {
                total += e[i].w;
                taken++;
            }
        }
        if (taken == n - 1) {
            printf("MST weight %lld, %d edges\n", total, taken);
        } else {
            printf("not connected: forest weight %lld, %d edges\n", total, taken);
        }
    }
    free(e);
    free(s.p);
    free(s.size);
    return ok ? 0 : 1;
}
\end{example}
\iosample{8 13 \\ 1 2 7 \\ 2 3 12 \\ 3 4 14 \\ 5 6 9 \\ 6 7 4 \\ 7 8 13 \\ 1 5 1 \\ 2 6 6 \\ 3 7 11 \\ 4 8 3 \\ 1 6 10 \\ 2 7 5 \\ 3 8 8}{1-5 (1) take \\ 4-8 (3) take \\ 6-7 (4) take \\ 2-7 (5) take \\ 2-6 (6) skip \\ 1-2 (7) take \\ 3-8 (8) take \\ 5-6 (9) skip \\ 1-6 (10) skip \\ 3-7 (11) take \\ MST weight 39, 7 edges}

輸出的十列和手算表的前十步一致。 程式把「找代表、比較、合併」寫成一個
\texttt{union\_set}，回傳值同時當作第 7 行的判斷，所以一條邊只找兩次代表。 \texttt{malloc}
的大小寫成 \texttt{m > 0 ? m : 1}，是為了 $m = 0$（只有一個頂點）時不配置 $0$ 個位元組。
圖不連通時，迴圈把邊全部看完，選到的邊數少於 $n - 1$，最後一列改印最小生成森林的權重。

\section{Prim 演算法}
\label{sec:mst-prim}

Prim 演算法只維護一棵樹。 從指定的起點 $r$ 開始，每一步在「一端在樹裡、一端在樹外」的
邊裡挑最輕的一條，把樹外那一端拉進來。 用切割的話說，切割永遠是（樹裡的頂點，樹外的頂
點），每次加的就是這個切割的輕邊。

要快速找到這條邊，對每個樹外的頂點 $v$ 記兩個值：
\begin{itemize}
\item $\mathit{key}[v]$：$v$ 連到樹裡的邊之中最輕的權重；連不到時是 $\infty$。
\item $\pi[v]$：那條最輕的邊在樹裡的那一端；連不到時是 \NIL。
\end{itemize}
下一個加進樹的，就是 $\mathit{key}$ 最小的樹外頂點 $u$，加的邊是 $(\pi[u], u)$。 $u$ 進了
樹以後，只有 $u$ 的鄰居的 $\mathit{key}$ 可能變小，掃一遍 $u$ 的鄰接串列更新就好。 「找
$\mathit{key}$ 最小的」和「把某個 $\mathit{key}$ 改小」正好是第~8~章最小優先佇列的
\proc{Extract-Min} 和 \proc{Decrease-Key}。

\begin{pseudo}[label=alg:mst-prim]{Prim 演算法（帶索引的最小堆積）}
\Procedure{MST-Prim}{$G, w, r$}
    \For{每個頂點 $v \in G.V$}
        \State $\mathit{key}[v] \gets \infty$；$\pi[v] \gets \NIL$；$\mathit{inT}[v] \gets \FALSE$
    \EndFor
    \State $\mathit{key}[r] \gets 0$
    \State $Q \gets$ 含全部頂點、以 $\mathit{key}$ 為鍵的帶索引最小堆積
    \While{$Q$ 不是空的}
        \State $u \gets \Call{Extract-Min}{Q}$ \Comment{$\mathit{key}$ 相同時取編號小的}
        \State $\mathit{inT}[u] \gets \TRUE$
        \For{每個 $v \in G.\mathit{Adj}[u]$}
            \If{$\mathit{inT}[v] = \FALSE$ \AND $w(u, v) < \mathit{key}[v]$}
                \State $\pi[v] \gets u$
                \State \Call{Decrease-Key}{$Q, v, w(u, v)$}
            \EndIf
        \EndFor
    \EndWhile
    \State \Return $\{(\pi[v], v) : v \ne r\}$
\EndProcedure
\end{pseudo}

$\mathit{inT}[v]$ 表示 $v$ 已經在樹裡，也就是已經離開 $Q$。 第 10 行的比較用嚴格的 $<$：
另一條一樣輕的邊不會換掉原來的 $\pi[v]$。 這和第 7 行「$\mathit{key}$ 相同取編號小的」合
起來，是本書 Prim 演算法的平手規則。

\begin{theorem}
\label{thm:mst-prim}
$G$ 連通時，\proc{MST-Prim} 回傳一棵最小生成樹。
\end{theorem}

\begin{proof}
令 $A = \{(\pi[v], v) : \mathit{inT}[v] = \TRUE,\ v \ne r\}$，也就是目前這棵樹的邊。 證明
每次執行第 7 行之前，下面兩件事都成立：
\begin{enumerate}
\item $A$ 包含在某一棵最小生成樹裡；
\item 對每個 $\mathit{inT}[v] = \FALSE$ 的頂點，$\mathit{key}[v]$ 是 $v$ 連到樹裡頂點的
邊之中最輕的權重（沒有這種邊時是 $\infty$），$\pi[v]$ 是那條邊在樹裡的一端。
\end{enumerate}
第一次執行時樹是空的，$A = \emptyset$，(1) 成立；(2) 裡沒有樹裡的頂點，所有 $\mathit{key}$
都是 $\infty$（$r$ 是 $0$，它會第一個被取出，沒有 $\pi$），成立。

設第 7 行取出 $u \ne r$。 令 $S$ 是目前樹裡的頂點。 切割 $(S, V - S)$ 尊重 $A$，因為 $A$ 的
邊兩端都在樹裡。 由 (2)，每個樹外頂點的 $\mathit{key}$ 就是它跨越這個切割的最輕的邊；$u$
的 $\mathit{key}$ 是所有樹外頂點裡最小的，所以 $(\pi[u], u)$ 是這個切割的輕邊（有平手時
它是輕邊之一）。 由定理~\ref{thm:mst-cut}，它對 $A$ 是安全的，加進來之後 (1) 仍成立。 $u$
進樹以後，樹外頂點 $v$「連到樹裡的邊」多了 $(u, v)$ 這一條，第 9--12 行正好拿它和原本的
$\mathit{key}[v]$ 比，(2) 也仍成立。

$G$ 連通，所以樹外還有頂點時，至少有一個頂點的 $\mathit{key}$ 不是 $\infty$，每次取出的
頂點都真的連得進樹。 迴圈結束時 $n$ 個頂點都在樹裡，$A$ 是一棵包含在某一棵最小生成樹裡
的生成樹，所以就是最小生成樹。
\end{proof}

如果取出的頂點 $\mathit{key}$ 是 $\infty$，表示剩下的頂點都連不到樹，圖不連通。 這時可以停
下來，或者再從剩下的頂點挑一個當新的起點，得到最小生成森林。

\subsection{手算追蹤}

在露營區上從頂點 $1$ 開始。 下表每一列是一回合：先取出 $\mathit{key}$ 最小的頂點（框起
來的格子），再更新它的鄰居。 格子寫成 $\mathit{key}/\pi$，是這一回合做完之後的值；有底
色的是這一回合改過的；空白的頂點已經在樹裡。

\begin{center}
\footnotesize
\setlength{\tabcolsep}{3.2pt}
\setlength{\fboxsep}{1.2pt}
\begin{tabular}{cccrrrrrrrr}
\toprule
回合 & 取出 & 加入的邊 & $1$ & $2$ & $3$ & $4$ & $5$ & $6$ & $7$ & $8$ \\
\midrule
1 & $1$ & — & \mstsel{$0$} & \mstchg{$7/1$} & $\infty$ & $\infty$ & \mstchg{$1/1$} & \mstchg{$10/1$} & $\infty$ & $\infty$ \\
2 & $5$ & $1$--$5$ &  & $7/1$ & $\infty$ & $\infty$ & \mstsel{$1/1$} & \mstchg{$9/5$} & $\infty$ & $\infty$ \\
3 & $2$ & $1$--$2$ &  & \mstsel{$7/1$} & \mstchg{$12/2$} & $\infty$ &  & \mstchg{$6/2$} & \mstchg{$5/2$} & $\infty$ \\
4 & $7$ & $2$--$7$ &  &  & \mstchg{$11/7$} & $\infty$ &  & \mstchg{$4/7$} & \mstsel{$5/2$} & \mstchg{$13/7$} \\
5 & $6$ & $7$--$6$ &  &  & $11/7$ & $\infty$ &  & \mstsel{$4/7$} &  & $13/7$ \\
6 & $3$ & $7$--$3$ &  &  & \mstsel{$11/7$} & \mstchg{$14/3$} &  &  &  & \mstchg{$8/3$} \\
7 & $8$ & $3$--$8$ &  &  &  & \mstchg{$3/8$} &  &  &  & \mstsel{$8/3$} \\
8 & $4$ & $8$--$4$ &  &  &  & \mstsel{$3/8$} &  &  &  &  \\
\bottomrule
\end{tabular}
\end{center}

加入的邊依序是 $1$--$5$、$1$--$2$、$2$--$7$、$7$--$6$、$7$--$3$、$3$--$8$、$8$--$4$，和
Kruskal 演算法選的是同一組邊（權重互不相同，最小生成樹只有一棵），但\textbf{順序不同}：
權重依序是 $1, 7, 5, 4, 11, 8, 3$，不是由小到大。 Prim 演算法每次挑的是「目前這棵樹」的
切割的輕邊，第 3 回合時權重 $4$ 的 $6$--$7$ 兩端都還不在樹裡，輪不到它。

再看 $\mathit{key}[6]$ 的變化：第 1 回合是 $10$（經過 $1$），第 2 回合變 $9$（經過 $5$），
第 3 回合變 $6$（經過 $2$），第 4 回合變 $4$（經過 $7$）。 一個頂點的 $\mathit{key}$ 可以
被改小好幾次，每次都是一次 \proc{Decrease-Key}。 這張圖一共做了 $13$ 次（每條邊都讓某個
$\mathit{key}$ 變小過一次），剛好碰到下面時間分析裡的上界 $m$。

\begin{examnote}
Kruskal 和 Prim 的追蹤題最常要的是「依序選了哪些邊」。 Kruskal 照權重（和平手規則）由
小到大，跳過的邊要寫原因（兩端已經連通）；Prim 要先看清楚題目指定的起點，選邊的順序取
決於樹長到哪裡，不是由小到大。 Prim 的過程用上面這種 $\mathit{key}$ 表最不容易漏，每回合
寫出取出的頂點和改過的 $\mathit{key}$。 兩者選到的邊集合在權重互不相同時一定相同，可以
拿來互相檢查；總權重任何時候都相同。
\end{examnote}

\subsection{時間}

用第~8~章的帶索引二元堆積：第 5 行把 $n$ 個頂點放進堆積，$\mathit{key}$ 只有 $r$ 是 $0$、
其他都是 $\infty$，直接排好就是堆積，$O(n)$。 第 7 行執行 $n$ 次，每次 $O(\lg n)$。 第 12
行最多執行 $m$ 次：邊 $(u, v)$ 在 $u$、$v$ 先進樹的那一端被看到時，另一端還在樹外，可能更
新一次；後進樹的那一端看到它時，另一端已經在樹裡，第 10 行不成立。 每次
\proc{Decrease-Key} 是 $O(\lg n)$。 第 9 行的迴圈總共掃過每個鄰接串列一次，$\Theta(m)$。
合起來
\[
O(n \lg n + m \lg n) = O(m \lg n) \qquad (\text{連通圖 } m \ge n - 1).
\]
第~\ref{sec:hp-fib}~節的費氏堆積把 \proc{Decrease-Key} 的平攤成本降到 $O(1)$，
\proc{Extract-Min} 仍是 $O(\lg n)$，總時間變成 $O(m + n \lg n)$。

\subsection{不做 decrease-key 的寫法}

第~8~章說過，帶索引的堆積要多維護一個位置陣列，寫起來比較長。 另一種常見寫法完全不
用 \proc{Decrease-Key}：堆積裡放的是（權重，樹外頂點，樹裡頂點）這樣的三元組，一個頂點的
$\mathit{key}$ 變小時，不去改舊的那筆，直接再放一筆新的。 同一個頂點可能在堆積裡有好幾
筆，最輕的那筆一定最先出來；後面出來的都是\textbf{過期}的，這時頂點已經在樹裡，丟掉就好。

\begin{pseudo}[label=alg:mst-prim-lazy]{Prim 演算法（允許重複放入）}
\Procedure{MST-Prim-Lazy}{$G, w, r$}
    \For{每個頂點 $v \in G.V$}
        \State $\mathit{inT}[v] \gets \FALSE$
    \EndFor
    \State $A \gets \emptyset$；$Q \gets$ 空的最小堆積
    \State \Call{Insert}{$Q, (0, r, \NIL)$}
    \While{$Q$ 不是空的}
        \State $(k, u, p) \gets \Call{Extract-Min}{Q}$
        \If{$\mathit{inT}[u] = \FALSE$} \Comment{否則是過期的資料，丟掉}
            \State $\mathit{inT}[u] \gets \TRUE$
            \If{$p \ne \NIL$}
                \State $A \gets A \cup \{(p, u)\}$
            \EndIf
            \For{每個 $v \in G.\mathit{Adj}[u]$}
                \If{$\mathit{inT}[v] = \FALSE$}
                    \State \Call{Insert}{$Q, (w(u, v), v, u)$}
                \EndIf
            \EndFor
        \EndIf
    \EndWhile
    \State \Return $A$
\EndProcedure
\end{pseudo}

三元組先比權重，權重相同時比樹外頂點的編號，這和虛擬碼~\ref{alg:mst-prim} 的平手規則一
致。 在露營區上從 $1$ 開始，堆積的變化如下（只寫三元組的前兩項）：

\begin{center}
\footnotesize
\setlength{\tabcolsep}{4pt}
\begin{tabular}{rcllc}
\toprule
步 & 取出 & 動作 & 放入 & 之後堆積裡有幾筆 \\
\midrule
1 & $(0, 1)$ & 起點 $1$ & $(7, 2)$、$(1, 5)$、$(10, 6)$ & 3 \\
2 & $(1, 5)$ & 加入 $1$--$5$ & $(9, 6)$ & 3 \\
3 & $(7, 2)$ & 加入 $1$--$2$ & $(12, 3)$、$(6, 6)$、$(5, 7)$ & 5 \\
4 & $(5, 7)$ & 加入 $2$--$7$ & $(11, 3)$、$(4, 6)$、$(13, 8)$ & 7 \\
5 & $(4, 6)$ & 加入 $7$--$6$ &  & 6 \\
6 & $(6, 6)$ & 過期，丟掉 & & 5 \\
7 & $(9, 6)$ & 過期，丟掉 & & 4 \\
8 & $(10, 6)$ & 過期，丟掉 & & 3 \\
9 & $(11, 3)$ & 加入 $7$--$3$ & $(14, 4)$、$(8, 8)$ & 4 \\
10 & $(8, 8)$ & 加入 $3$--$8$ & $(3, 4)$ & 4 \\
11 & $(3, 4)$ & 加入 $8$--$4$ &  & 3 \\
12--14 & \multicolumn{4}{l}{$(12, 3)$、$(13, 8)$、$(14, 4)$ 依序取出，都過期，丟掉} \\
\bottomrule
\end{tabular}
\end{center}

頂點 $6$ 一共被放進去四次（權重 $10, 9, 6, 4$），只有最輕的那筆有用。 加入的邊和順序都和
前一個表相同。 第 11 步之後 $8$ 個頂點都進了樹，其實可以直接停下來，剩下三筆不用取出。

成本：每條邊 $(u, v)$ 只在 $u$、$v$ 先進樹的那一端被放進去一次（後進樹的那一端看到它時，
另一端已經在樹裡），加上起點那一筆，最多 $m + 1$ 次 \proc{Insert}，\proc{Extract-Min}
次數也一樣。 堆積最多有 $O(m)$ 筆，每次操作 $O(\lg m) = O(\lg n)$，總時間 $O(m \lg n)$，
和帶索引的版本同級；代價是堆積最多佔 $\Theta(m)$ 的空間（帶索引的版本是 $\Theta(n)$），而
且比較次數多一些。 好處是程式短、不用位置陣列。 這個例子放了 $14$ 筆，帶索引的版本做了
$13$ 次 \proc{Decrease-Key}，兩者差不多。

\subsection{陣列版：\texorpdfstring{$O(n^2)$}{O(n²)}}

如果不用堆積，每回合直接掃過所有頂點找 $\mathit{key}$ 最小的，就是 $\Theta(n)$。 圖用鄰
接矩陣存，$W[u][v]$ 是邊的權重，沒有邊時是 $\infty$。

\begin{pseudo}[label=alg:mst-prim-array]{Prim 演算法（陣列版）}
\Procedure{MST-Prim-Array}{$W, n, r$}
    \For{$v \gets 1$ \To $n$}
        \State $\mathit{key}[v] \gets \infty$；$\pi[v] \gets \NIL$；$\mathit{inT}[v] \gets \FALSE$
    \EndFor
    \State $\mathit{key}[r] \gets 0$
    \For{$i \gets 1$ \To $n$}
        \State $u \gets \NIL$
        \For{$v \gets 1$ \To $n$} \Comment{找 $\mathit{key}$ 最小的；平手時編號小的}
            \If{$\mathit{inT}[v] = \FALSE$ \AND ($u = \NIL$ \OR $\mathit{key}[v] < \mathit{key}[u]$)}
                \State $u \gets v$
            \EndIf
        \EndFor
        \State $\mathit{inT}[u] \gets \TRUE$
        \For{$v \gets 1$ \To $n$}
            \If{$\mathit{inT}[v] = \FALSE$ \AND $W[u][v] < \mathit{key}[v]$}
                \State $\mathit{key}[v] \gets W[u][v]$；$\pi[v] \gets u$
            \EndIf
        \EndFor
    \EndFor
    \State \Return $\pi$
\EndProcedure
\end{pseudo}

第 8 行用嚴格的 $<$，$v$ 又由小到大，所以 $\mathit{key}$ 相同時留下的是編號小的，和堆積
版的平手規則相同，兩個版本每一回合取出的頂點完全一樣。 第 5 行的迴圈做 $n$ 回合，每回合
第 7、11 行各掃 $n$ 個頂點，總時間 $\Theta(n^2)$，和邊數無關。 鄰接矩陣要 $\Theta(n^2)$
的空間；如果圖用鄰接串列存，第 11--13 行改成掃 $u$ 的鄰接串列，總時間是 $\Theta(n^2 + m)
= \Theta(n^2)$，空間只要 $\Theta(n + m)$。

\subsection{稠密的圖選哪一個}

把本章的寫法放在一起比（$G$ 連通，所以 $n - 1 \le m < n^2/2$）：

\begin{center}
\small
\begin{tabular}{llll}
\toprule
寫法 & 時間 & 額外空間 & 適合 \\
\midrule
Kruskal & $O(m \lg n)$ & 邊陣列 $\Theta(m)$ & 稀疏的圖、邊已排好 \\
Prim，帶索引堆積 & $O(m \lg n)$ & $\Theta(n)$ & 稀疏的圖 \\
Prim，允許重複放入 & $O(m \lg n)$ & $\Theta(m)$ & 稀疏的圖，程式最短 \\
Prim，陣列 & $\Theta(n^2)$ & $\Theta(n)$ & 稠密的圖 \\
Prim，費氏堆積 & $O(m + n \lg n)$ & $\Theta(n)$ & 理論上最好，實作複雜 \\
Sollin（第~\ref{sec:mst-sollin}~節） & $O(m \lg n)$ & $\Theta(n + m)$ & 可以平行做 \\
\bottomrule
\end{tabular}
\end{center}

\textbf{稠密}的圖指 $m = \Theta(n^2)$，例如任兩點之間都能拉線的完全圖。 這時堆積版是
$O(n^2 \lg n)$，陣列版是 $\Theta(n^2)$，陣列版勝出；而且不用堆積、不用邊陣列，常數也小。
\textbf{稀疏}的圖指 $m = O(n)$，例如道路網，每個頂點只連到少數幾個鄰居。 這時堆積版和
Kruskal 都是 $O(n \lg n)$，陣列版仍是 $\Theta(n^2)$，堆積版勝出。 兩者打平的地方大約是 $m
\lg n \approx n^2$，也就是 $m \approx n^2 / \lg n$。

用數字感受一下。 $n = 2000$ 的完全圖有 $m = 1\,999\,000$ 條邊：陣列版兩層迴圈大約 $2n^2 =
8 \times 10^6$ 步；堆積版最壞要做接近 $m$ 次 \proc{Decrease-Key}，每次最多上浮 $\lg n
\approx 11$ 層，大約 $2 \times 10^7$ 次交換；Kruskal 光是排序 $2 \times 10^6$ 條邊就要大約
$m \lg m \approx 4 \times 10^7$ 次比較，還要先把邊全部存下來。 反過來，$n = 10^5$、$m = 3
\times 10^5$ 的道路網，陣列版要 $10^{10}$ 步，完全不能用；堆積版大約 $(n + m) \lg n \approx
7 \times 10^6$。

\subsection{權重相同時，Kruskal 和 Prim 可能選出不同的樹}

下圖左邊的圖有四條權重 $3$ 的邊，最小生成樹不只一棵（總權重都是 $9$）。

\begin{center}
\begin{tabular}{ccc}
\begin{tikzpicture}[x=0.9cm, y=0.9cm]
\node[mst-v] (v1) at (0,1) {$1$};
\node[mst-v] (v2) at (1.6,2) {$2$};
\node[mst-v] (v3) at (1.6,0) {$3$};
\node[mst-v] (v4) at (3.2,2) {$4$};
\node[mst-v] (v5) at (3.2,0) {$5$};
\draw[mst-e] (v1) -- node[mst-w] {4} (v2);
\draw[mst-e] (v1) -- node[mst-w] {3} (v3);
\draw[mst-e] (v2) -- node[mst-w] {1} (v3);
\draw[mst-e] (v2) -- node[mst-w] {3} (v4);
\draw[mst-e] (v3) -- node[mst-w] {3} (v4);
\draw[mst-e] (v3) -- node[mst-w] {2} (v5);
\draw[mst-e] (v4) -- node[mst-w] {3} (v5);
\end{tikzpicture}
&
\begin{tikzpicture}[x=0.9cm, y=0.9cm]
\node[mst-v] (v1) at (0,1) {$1$};
\node[mst-v] (v2) at (1.6,2) {$2$};
\node[mst-v] (v3) at (1.6,0) {$3$};
\node[mst-v] (v4) at (3.2,2) {$4$};
\node[mst-v] (v5) at (3.2,0) {$5$};
\draw[mst-x] (v1) -- node[mst-w] {4} (v2);
\draw[mst-t] (v1) -- node[mst-w] {3} (v3);
\draw[mst-t] (v2) -- node[mst-w] {1} (v3);
\draw[mst-t] (v2) -- node[mst-w] {3} (v4);
\draw[mst-x] (v3) -- node[mst-w] {3} (v4);
\draw[mst-t] (v3) -- node[mst-w] {2} (v5);
\draw[mst-x] (v4) -- node[mst-w] {3} (v5);
\end{tikzpicture}
&
\begin{tikzpicture}[x=0.9cm, y=0.9cm]
\node[mst-v] (v1) at (0,1) {$1$};
\node[mst-v] (v2) at (1.6,2) {$2$};
\node[mst-v] (v3) at (1.6,0) {$3$};
\node[mst-v] (v4) at (3.2,2) {$4$};
\node[mst-v] (v5) at (3.2,0) {$5$};
\draw[mst-x] (v1) -- node[mst-w] {4} (v2);
\draw[mst-t] (v1) -- node[mst-w] {3} (v3);
\draw[mst-t] (v2) -- node[mst-w] {1} (v3);
\draw[mst-x] (v2) -- node[mst-w] {3} (v4);
\draw[mst-t] (v3) -- node[mst-w] {3} (v4);
\draw[mst-t] (v3) -- node[mst-w] {2} (v5);
\draw[mst-x] (v4) -- node[mst-w] {3} (v5);
\end{tikzpicture}
\\
\small 原圖 & \small Kruskal（照 $\prec$） & \small Prim（從 $1$ 開始）
\end{tabular}
\diagcaption{平手時兩個演算法選出不同的最小生成樹。 頂點 $4$ 一定要從三條權重 $3$ 的邊
$2$--$4$、$3$--$4$、$4$--$5$ 裡挑一條接進來，三種挑法都是最小生成樹。}
\label{fig:mst-ties}
\end{center}

Kruskal 照 $\prec$ 排序：$2$--$3$（$1$）、$3$--$5$（$2$），接著四條權重 $3$ 的邊依小端點、
大端點排成 $1$--$3$、$2$--$4$、$3$--$4$、$4$--$5$。 $1$--$3$ 和 $2$--$4$ 都選，這時已經
選滿 $4$ 條。 Prim 從 $1$ 開始：先取出 $3$（$\mathit{key} = 3$），$3$ 把 $\mathit{key}[2]$
改成 $1$、$\mathit{key}[4]$ 改成 $3$（$\pi[4] = 3$）、$\mathit{key}[5]$ 改成 $2$；接著取出 $2$，
邊 $2$--$4$ 的權重 $3$ 沒有\textbf{小於} $\mathit{key}[4] = 3$，不更新；再取出 $5$，$4$--$5$ 也一
樣；最後 $4$ 經由 $\pi[4] = 3$ 接進來。 兩棵樹差在頂點 $4$ 接在哪裡，都對。

這說明了兩件事。 第一，Prim 的平手規則比的是頂點，不是邊，所以它的結果不一定是引
理~\ref{lem:mst-ties} 的 $T_{\prec}$；它仍然是最小生成樹，因為定理~\ref{thm:mst-prim} 的證明
只用到「輕邊之一」。 第二，第 10 行如果改成 $\le$，$\pi[4]$ 會被後來一樣輕的邊換掉，最後變成
$4$--$5$，又是另一棵。 所以答案不只一個時，平手規則（包括更新時用 $<$ 還是 $\le$）一定要
照題目。

\subsection{C 程式}

下面兩個程式都從頂點 $1$ 開始，印出每回合加入的頂點和邊。 第一個照
虛擬碼~\ref{alg:mst-prim}，用鄰接串列和第~8~章的帶索引最小堆積，另外數
\proc{Decrease-Key} 做了幾次；第二個照虛擬碼~\ref{alg:mst-prim-array}，用鄰接矩陣。 兩
個程式的平手規則相同，對同一張圖印出的過程應該完全相同。

\begin{cbookrecap}[第~14~章結構、第~19~章 \texttt{malloc} 與鏈結串列]
堆積的三個陣列包在 \texttt{MinHeap} 結構裡：\texttt{heap[i]} 是第 \texttt{i} 格放的頂點，
\texttt{pos[v]} 是頂點 \texttt{v} 在哪一格，\texttt{key[v]} 是它的鍵。 交換兩格時兩個頂點的
\texttt{pos} 都要改，寫成 \texttt{swap\_at}；這和第~8~章急診室候診名單的程式是同一個寫
法。 函式收到的是 \texttt{MinHeap *h}，用 \texttt{h->key[v]} 存取，改的是呼叫端那一份。
\end{cbookrecap}

\begin{example}[label=ex:mst-prim-heap]{c}{(prim\_heap.c) Prim 演算法：鄰接串列加帶索引最小堆積}
#include <stdio.h>
#include <stdlib.h>
#include <limits.h>

typedef struct Node {
    int v, w;
    struct Node *next;
} Node;

typedef struct {
    int n;
    Node **adj;     // adj[u]：u 的鄰居，編號由小到大
} Graph;

// 把 v 插進 adj[u]，保持由小到大
int add_arc(Graph *g, int u, int v, int w)
{
    Node *x = malloc(sizeof *x);
    if (x == NULL) {
        return 0;
    }
    Node **p = &g->adj[u];
    while (*p != NULL && (*p)->v < v) {
        p = &(*p)->next;
    }
    x->v = v;
    x->w = w;
    x->next = *p;
    *p = x;
    return 1;
}

void free_graph(Graph *g)
{
    for (int u = 0; g->adj != NULL && u < g->n; u++) {
        Node *x = g->adj[u];
        while (x != NULL) {
            Node *next = x->next;
            free(x);
            x = next;
        }
    }
    free(g->adj);
}

// 第 8 章的帶索引最小堆積。 比較的鍵是 (key, 頂點編號)，
// 所以 key 相同時編號小的先出來。
typedef struct {
    int *heap;      // heap[0..size-1]：頂點編號
    int *pos;       // pos[v]：v 在 heap 的哪一格，-1 表示已經取出
    long long *key;
    int size;
} MinHeap;

static int before(const MinHeap *h, int u, int v)
{
    return h->key[u] < h->key[v] || (h->key[u] == h->key[v] && u < v);
}

static void swap_at(MinHeap *h, int i, int j)
{
    int u = h->heap[i], v = h->heap[j];
    h->heap[i] = v;
    h->heap[j] = u;
    h->pos[v] = i;
    h->pos[u] = j;
}

static void sift_up(MinHeap *h, int i)
{
    while (i > 0 && before(h, h->heap[i], h->heap[(i - 1) / 2])) {
        swap_at(h, i, (i - 1) / 2);
        i = (i - 1) / 2;
    }
}

static void sift_down(MinHeap *h, int i)
{
    for (;;) {
        int j = 2 * i + 1;
        if (j >= h->size) {
            return;
        }
        if (j + 1 < h->size && before(h, h->heap[j + 1], h->heap[j])) {
            j++;
        }
        if (!before(h, h->heap[j], h->heap[i])) {
            return;
        }
        swap_at(h, i, j);
        i = j;
    }
}

int extract_min(MinHeap *h)
{
    int u = h->heap[0];
    swap_at(h, 0, h->size - 1);
    h->size--;
    h->pos[u] = -1;
    sift_down(h, 0);
    return u;
}

void decrease_key(MinHeap *h, int v, long long k)
{
    h->key[v] = k;
    sift_up(h, h->pos[v]);
}

int main(void)
{
    int n, m;
    if (scanf("%d %d", &n, &m) != 2 || n <= 0 || m < 0) {
        return 1;
    }
    Graph g = {n, calloc(n, sizeof(Node *))};
    MinHeap h = {malloc(n * sizeof(int)), malloc(n * sizeof(int)),
                 malloc(n * sizeof(long long)), n};
    int *pi = malloc(n * sizeof(int));
    int ok = g.adj != NULL && h.heap != NULL && h.pos != NULL
             && h.key != NULL && pi != NULL;
    for (int i = 0; ok && i < m; i++) {
        int u, v, w;
        ok = scanf("%d %d %d", &u, &v, &w) == 3 && u >= 1 && u <= n
             && v >= 1 && v <= n;
        if (ok && u != v) {     // 自環不會在生成樹裡，直接丟掉
            ok = add_arc(&g, u - 1, v - 1, w) && add_arc(&g, v - 1, u - 1, w);
        }
    }
    if (ok) {
        // 所有頂點一開始都在堆積裡，key 是無限大；起點是頂點 1
        for (int v = 0; v < n; v++) {
            h.heap[v] = v;
            h.pos[v] = v;
            h.key[v] = LLONG_MAX;
            pi[v] = -1;
        }
        h.key[0] = 0;
        long long total = 0;
        int calls = 0, added = 0;
        while (h.size > 0) {
            int u = extract_min(&h);
            if (h.key[u] == LLONG_MAX) {
                break;                      // 剩下的頂點都連不到
            }
            added++;
            if (pi[u] < 0) {
                printf("add %d\n", u + 1);
            } else {
                printf("add %d via %d-%d (%lld)\n", u + 1, pi[u] + 1, u + 1, h.key[u]);
                total += h.key[u];
            }
            for (Node *x = g.adj[u]; x != NULL; x = x->next) {
                int v = x->v;
                if (h.pos[v] >= 0 && x->w < h.key[v]) {
                    pi[v] = u;
                    decrease_key(&h, v, x->w);
                    calls++;
                }
            }
        }
        printf("decrease-key calls %d\n", calls);
        if (added == n) {
            printf("MST weight %lld\n", total);
        } else {
            printf("not connected: reached %d of %d vertices\n", added, n);
        }
    }
    free_graph(&g);
    free(h.heap);
    free(h.pos);
    free(h.key);
    free(pi);
    return ok ? 0 : 1;
}
\end{example}
\iosample{8 13 \\ 1 2 7 \\ 2 3 12 \\ 3 4 14 \\ 5 6 9 \\ 6 7 4 \\ 7 8 13 \\ 1 5 1 \\ 2 6 6 \\ 3 7 11 \\ 4 8 3 \\ 1 6 10 \\ 2 7 5 \\ 3 8 8}{add 1 \\ add 5 via 1-5 (1) \\ add 2 via 1-2 (7) \\ add 7 via 2-7 (5) \\ add 6 via 7-6 (4) \\ add 3 via 7-3 (11) \\ add 8 via 3-8 (8) \\ add 4 via 8-4 (3) \\ decrease-key calls 13 \\ MST weight 39}

一開始所有 $\mathit{key}$ 都是 \texttt{LLONG\_MAX}，只有頂點 $0$（書上的 $1$）是 $0$，陣列
\texttt{heap} 照編號排就已經是一個合法的堆積，不用另外建。 \texttt{pos[v] >= 0} 就是虛擬
碼的「$v$ 還在 $Q$ 裡」。 自環在讀入時就丟掉；平行邊不用特別處理，比較輕的那條自然會讓
$\mathit{key}$ 變得比較小。

\begin{example}[label=ex:mst-prim-array]{c}{(prim\_array.c) Prim 演算法：鄰接矩陣、$O(n^2)$}
#include <stdio.h>
#include <stdlib.h>
#include <limits.h>

#define NO_EDGE INT_MAX     // 鄰接矩陣裡「沒有邊」，也當作 key 的無限大

int main(void)
{
    int n, m;
    if (scanf("%d %d", &n, &m) != 2 || n <= 0 || m < 0) {
        return 1;
    }
    int *W = malloc((size_t)n * n * sizeof(int));  // W[u * n + v]
    int *key = malloc(n * sizeof(int));
    int *pi = malloc(n * sizeof(int));
    int *in_tree = calloc(n, sizeof(int));
    int ok = W != NULL && key != NULL && pi != NULL && in_tree != NULL;
    for (int i = 0; ok && i < n * n; i++) {
        W[i] = NO_EDGE;
    }
    for (int i = 0; ok && i < m; i++) {
        int u, v, w;
        ok = scanf("%d %d %d", &u, &v, &w) == 3 && u >= 1 && u <= n
             && v >= 1 && v <= n && w < NO_EDGE;
        if (ok && u != v && w < W[(u - 1) * n + (v - 1)]) {
            W[(u - 1) * n + (v - 1)] = w;   // 平行邊只留最輕的
            W[(v - 1) * n + (u - 1)] = w;
        }
    }
    if (ok) {
        for (int v = 0; v < n; v++) {
            key[v] = NO_EDGE;
            pi[v] = -1;
        }
        key[0] = 0;
        long long total = 0;
        int added = 0;
        for (int i = 0; i < n; i++) {
            int u = -1;
            for (int v = 0; v < n; v++) {   // 找 key 最小的；平手時編號小的
                if (!in_tree[v] && (u < 0 || key[v] < key[u])) {
                    u = v;
                }
            }
            if (key[u] == NO_EDGE) {
                break;                      // 剩下的頂點都連不到
            }
            in_tree[u] = 1;
            added++;
            if (pi[u] < 0) {
                printf("add %d\n", u + 1);
            } else {
                printf("add %d via %d-%d (%d)\n", u + 1, pi[u] + 1, u + 1, key[u]);
                total += key[u];
            }
            for (int v = 0; v < n; v++) {
                if (!in_tree[v] && W[u * n + v] < key[v]) {
                    key[v] = W[u * n + v];
                    pi[v] = u;
                }
            }
        }
        if (added == n) {
            printf("MST weight %lld\n", total);
        } else {
            printf("not connected: reached %d of %d vertices\n", added, n);
        }
    }
    free(W);
    free(key);
    free(pi);
    free(in_tree);
    return ok ? 0 : 1;
}
\end{example}
\iosample{8 13 \\ 1 2 7 \\ 2 3 12 \\ 3 4 14 \\ 5 6 9 \\ 6 7 4 \\ 7 8 13 \\ 1 5 1 \\ 2 6 6 \\ 3 7 11 \\ 4 8 3 \\ 1 6 10 \\ 2 7 5 \\ 3 8 8}{add 1 \\ add 5 via 1-5 (1) \\ add 2 via 1-2 (7) \\ add 7 via 2-7 (5) \\ add 6 via 7-6 (4) \\ add 3 via 7-3 (11) \\ add 8 via 3-8 (8) \\ add 4 via 8-4 (3) \\ MST weight 39}

鄰接矩陣用一維陣列存，第 $u$ 列第 $v$ 行放在 \texttt{W[u * n + v]}，只要一次
\texttt{malloc}。 配置大小先轉成 \texttt{size\_t} 再乘，避免 $n$ 很大時 \texttt{int} 溢位。
\texttt{INT\_MAX} 同時代表「沒有邊」和 $\infty$：沒有邊的格子永遠不會小於 $\mathit{key}$，
更新 $\mathit{key}$ 的迴圈自然跳過它。 這個程式不需要堆積，整個程式不到八十行，$n$ 幾千、
邊很密的時候反而是最快的寫法。

\section{Sollin 演算法（Borůvka 演算法）}
\label{sec:mst-sollin}

Kruskal 和 Prim 每一步只加一條邊。 Sollin 演算法一次加很多條：把目前森林裡的每一棵樹看
成一個\textbf{元件}，每個元件同時挑一條「連到別的元件的最輕的邊」，全部加進來，元件就兩
兩併起來。 重複幾回合，直到只剩一個元件。 這個方法最早是 Borůvka 在 1926 年為了規劃電網
提出的，後來又被好幾個人重新發現，不少資料結構課本叫它 Sollin 演算法。

每個元件 $C$ 挑的邊，是切割 $(C, V - C)$ 的輕邊。 切割尊重目前的森林（森林的邊不會一端在
$C$ 裡、一端在外面），所以由切割性質，\textbf{單獨看}每一條都是安全的。 麻煩在「同時」：
定理~\ref{thm:mst-cut} 一次只保證一條。 下面會看到，權重相同時如果每個元件各挑各的，一次加
進來的邊可能形成環。

\begin{pseudo}[label=alg:mst-sollin]{Sollin 演算法}
\Procedure{MST-Sollin}{$G, w$}
    \State $A \gets \emptyset$
    \For{每個頂點 $v \in G.V$}
        \State \Call{Make-Set}{$v$}
    \EndFor
    \While{$|A| < n - 1$}
        \For{每個頂點 $v \in G.V$}
            \State $\mathit{cheap}[v] \gets \NIL$
        \EndFor
        \For{每條邊 $e = (u, v) \in G.E$}
            \State $a \gets \Call{Find-Set}{u}$；$b \gets \Call{Find-Set}{v}$
            \If{$a \ne b$}
                \If{$\mathit{cheap}[a] = \NIL$ \OR $e \prec \mathit{cheap}[a]$}
                    \State $\mathit{cheap}[a] \gets e$
                \EndIf
                \If{$\mathit{cheap}[b] = \NIL$ \OR $e \prec \mathit{cheap}[b]$}
                    \State $\mathit{cheap}[b] \gets e$
                \EndIf
            \EndIf
        \EndFor
        \For{每個 $\mathit{cheap}[c] \ne \NIL$ 的 $c$}
            \State $(u, v) \gets \mathit{cheap}[c]$
            \If{$\Call{Find-Set}{u} \ne \Call{Find-Set}{v}$} \Comment{兩個元件可能挑到同一條}
                \State $A \gets A \cup \{(u, v)\}$；\Call{Union}{$u, v$}
            \EndIf
        \EndFor
    \EndWhile
    \State \Return $A$
\EndProcedure
\end{pseudo}

$\mathit{cheap}[c]$ 用元件的代表 $c$ 當索引，記這個元件目前找到最輕的往外的邊。 第 8--14
行掃一遍所有的邊，每條跨越兩個元件的邊，兩端的元件各拿它比一次；比較一律用
第~\ref{sec:mst-ties}~節的 $\prec$。 第 15--18 行把挑到的邊加進來。 同一條邊常常被兩端的
元件同時挑中（例如兩個元件互相是對方最近的），第二次處理時兩端已經合併，第 17 行把它跳
過。 圖不連通時第 5 行的條件永遠不會成立，實作時要改成「這一回合沒有加任何邊就停」，結果
是最小生成森林。

\subsection{回合數}

\begin{theorem}
\label{thm:mst-rounds}
$G$ 連通時，Sollin 演算法每一回合至少讓元件數減半，所以最多 $\lfloor \lg n \rfloor$ 回合。
\end{theorem}

\begin{proof}
設某回合開始時有 $k \ge 2$ 個元件。 $G$ 連通，每個元件都有往外的邊，都會挑到一條，這條邊
把它和另一個元件連起來。 所以這一回合結束時，每個新的元件都至少包含兩個舊的元件，新的
元件數 $\le \lfloor k/2 \rfloor$。 從 $n$ 個元件開始，$r$ 回合後最多剩 $\lfloor n / 2^r
\rfloor$ 個（每回合都無條件捨去），$r = \lfloor \lg n \rfloor$ 時這個數是 $1$。
\end{proof}

每一回合第 6--7 行是 $\Theta(n)$，第 8--14 行對每條邊做兩次 \proc{Find-Set}，第 15--18 行
最多 $n$ 次，用第~13~章的互斥集合是 $O((n + m)\,\alpha(n))$。 也可以每回合開始時用第~21~章
的 DFS 把森林的元件重新編號，不用互斥集合，每回合 $\Theta(n + m)$。 總時間 $O(m \lg n)$，
和另外兩個演算法同級。 它的好處是每一回合裡各個元件的工作互不相干，很適合分給很多處理器
同時做；而且回合數少，每回合結束時把元件縮成一個點、重複的邊只留最輕的，圖會越縮越小。
Karger、Klein、Tarjan 的線性期望時間演算法（見延伸閱讀）就用 Sollin 的回合當成基本步驟。

\subsection{平手時要用一致的規則}

看一個三角形：三個頂點 $1, 2, 3$，三條邊的權重都是 $5$。 第一回合每個頂點自己是一個元件，
每個都要挑一條最輕的往外的邊，而每個頂點的兩條邊一樣輕。 如果各自隨便挑，可能出現
「$1$ 挑 $1$--$2$、$2$ 挑 $2$--$3$、$3$ 挑 $3$--$1$」：每一條單獨看都是輕邊，但三條一起加
進來是一個環。

\begin{center}
\begin{tabular}{cc}
\begin{tikzpicture}[x=1.2cm, y=1.2cm]
\node[mst-v] (a) at (0,0) {$1$};
\node[mst-v] (b) at (2,0) {$2$};
\node[mst-v] (c) at (1,1.6) {$3$};
\draw[mst-new] (a) -- node[mst-w] (ab) {5} (b);
\draw[mst-new] (b) -- node[mst-w] (bc) {5} (c);
\draw[mst-new] (c) -- node[mst-w] (ca) {5} (a);
\draw[mst-pick] (a) to[bend right=25] (ab);
\draw[mst-pick] (b) to[bend right=25] (bc);
\draw[mst-pick] (c) to[bend right=25] (ca);
\end{tikzpicture}
&
\begin{tikzpicture}[x=1.2cm, y=1.2cm]
\node[mst-v] (a) at (0,0) {$1$};
\node[mst-v] (b) at (2,0) {$2$};
\node[mst-v] (c) at (1,1.6) {$3$};
\draw[mst-t] (a) -- node[mst-w] (ab) {5} (b);
\draw[mst-x] (b) -- node[mst-w] (bc) {5} (c);
\draw[mst-t] (c) -- node[mst-w] (ca) {5} (a);
\draw[mst-pick] (a) to[bend right=25] (ab);
\draw[mst-pick] (b) to[bend left=25] (ab);
\draw[mst-pick] (c) to[bend left=25] (ca);
\end{tikzpicture}
\\
\small 各挑各的：三條邊形成環 & \small 照 $\prec$：$1$、$2$ 都挑 $1$--$2$，$3$ 挑 $1$--$3$
\end{tabular}
\diagcaption{三條邊權重相同的三角形。 箭頭從頂點指向它在第一回合挑的邊。}
\label{fig:mst-sollin-cycle}
\end{center}

照 $\prec$ 比較時，頂點 $1$ 的兩條邊 $(5, 1, 2) \prec (5, 1, 3)$，挑 $1$--$2$；頂點 $2$ 的兩
條邊 $(5, 1, 2) \prec (5, 2, 3)$，也挑 $1$--$2$；頂點 $3$ 的兩條邊 $(5, 1, 3) \prec (5, 2,
3)$，挑 $1$--$3$。 加進來的是 $1$--$2$ 和 $1$--$3$，沒有環。 一般的情形是下面的定理：

\begin{theorem}
\label{thm:mst-sollin}
如果所有元件挑邊時都照同一個嚴格的全序 $\prec$（第~\ref{sec:mst-ties}~節）比較，Sollin 演
算法每一回合挑到的邊加起來不會形成環，最後得到的是引理~\ref{lem:mst-ties} 的 $T_{\prec}$，也
就是一棵最小生成樹。
\end{theorem}

\begin{proof}
由引理~\ref{lem:mst-ties}，照 $\prec$ 比較等於在權重互不相同的 $w'$ 上比較，$w'$ 的最小生
成樹只有 $T_{\prec}$ 一棵。 對回合做歸納，假設回合開始時 $A \subseteq T_{\prec}$（一開始
$A = \emptyset$）。 這一回合元件 $C$ 挑的邊 $e_C$，是切割 $(C, V - C)$ 在 $w'$ 之下唯一的輕
邊，由系理~\ref{cor:mst-cut-strict}，$e_C$ 在 $w'$ 的每一棵最小生成樹裡，也就是 $e_C \in
T_{\prec}$。 所以 $A$ 加上這一回合所有挑到的邊，仍然包含在 $T_{\prec}$ 裡，而 $T_{\prec}$
是樹，這些邊不會形成環。 最後 $A$ 有 $n - 1$ 條邊，就是 $T_{\prec}$。
\end{proof}

證明的關鍵是「每個元件挑的邊都在\textbf{同一棵}樹 $T_{\prec}$ 裡」。 各挑各的時候，每條邊
雖然各自屬於某一棵最小生成樹，卻可能屬於不同的那一棵，合起來就不保證了。

虛擬碼~\ref{alg:mst-sollin} 第 17 行的檢查，在程式裡其實也會把形成環的邊擋掉。 但手算時
習慣把整回合挑到的邊一次畫上去，平手規則不一致就會畫出環；更重要的是，上面的正確性證明
要靠一致的 $\prec$。 所以寫 Sollin 演算法時，不管手算還是程式，平手規則都要先定好、全部
元件共用。

\subsection{手算追蹤}

露營區的權重互不相同，不會遇到平手。 第一回合每個頂點自己是一個元件：

\begin{center}
\small
\begin{tabular}{cccccccc}
\toprule
$\{1\}$ & $\{2\}$ & $\{3\}$ & $\{4\}$ & $\{5\}$ & $\{6\}$ & $\{7\}$ & $\{8\}$ \\
\midrule
$1$--$5$（$1$） & $2$--$7$（$5$） & $3$--$8$（$8$） & $4$--$8$（$3$） & $1$--$5$（$1$） &
$6$--$7$（$4$） & $6$--$7$（$4$） & $4$--$8$（$3$） \\
\bottomrule
\end{tabular}
\end{center}

八個元件挑了八次，但 $1$--$5$、$6$--$7$、$4$--$8$ 各被挑了兩次，去掉重複後加入五條邊：
$1$--$5$、$2$--$7$、$3$--$8$、$4$--$8$、$6$--$7$。 元件剩下 $\{1, 5\}$、$\{2, 6, 7\}$、
$\{3, 4, 8\}$ 三個。

\begin{center}
\begin{tikzpicture}[x=1.2cm, y=1.2cm]
\node[mst-v] (v1) at (0,2) {$1$};
\node[mst-v] (v2) at (2,2) {$2$};
\node[mst-v] (v3) at (4,2) {$3$};
\node[mst-v] (v4) at (6,2) {$4$};
\node[mst-v] (v5) at (0,0) {$5$};
\node[mst-v] (v6) at (2,0) {$6$};
\node[mst-v] (v7) at (4,0) {$7$};
\node[mst-v] (v8) at (6,0) {$8$};
\begin{scope}[on background layer]
\draw[mst-hull] (v1.center) -- (v5.center);
\filldraw[mst-hull] (v2.center) -- (v7.center) -- (v6.center) -- cycle;
\filldraw[mst-hull] (v3.center) -- (v8.center) -- (v4.center) -- cycle;
\end{scope}
\draw[mst-new] (v1) -- node[mst-w] {7} (v2);
\draw[mst-x] (v2) -- node[mst-w] {12} (v3);
\draw[mst-x] (v3) -- node[mst-w] {14} (v4);
\draw[mst-x] (v5) -- node[mst-w] {9} (v6);
\draw[mst-t] (v6) -- node[mst-w] {4} (v7);
\draw[mst-x] (v7) -- node[mst-w] {13} (v8);
\draw[mst-t] (v1) -- node[mst-w] {1} (v5);
\draw[mst-x] (v2) -- node[mst-w] {6} (v6);
\draw[mst-new] (v3) -- node[mst-w] {11} (v7);
\draw[mst-t] (v4) -- node[mst-w] {3} (v8);
\draw[mst-x] (v1) -- node[mst-w] {10} (v6);
\draw[mst-t] (v2) -- node[mst-w] {5} (v7);
\draw[mst-t] (v3) -- node[mst-w] {8} (v8);
\end{tikzpicture}
\diagcaption{第一回合後的三個元件（藍底），粗黑線是第一回合加入的五條邊。 紅線是第二回合
加入的兩條，虛線是始終沒選的邊。}
\label{fig:mst-sollin}
\end{center}

第二回合，三個元件往外的邊：

\begin{center}
\small
\begin{tabular}{lll}
\toprule
元件 & 往外的邊（權重） & 挑 \\
\midrule
$\{1, 5\}$ & $1$--$2$（$7$）、$5$--$6$（$9$）、$1$--$6$（$10$） & $1$--$2$（$7$） \\
$\{2, 6, 7\}$ & $1$--$2$（$7$）、$5$--$6$（$9$）、$1$--$6$（$10$）、$3$--$7$（$11$）、$2$--$3$（$12$）、$7$--$8$（$13$） & $1$--$2$（$7$） \\
$\{3, 4, 8\}$ & $3$--$7$（$11$）、$2$--$3$（$12$）、$7$--$8$（$13$） & $3$--$7$（$11$） \\
\bottomrule
\end{tabular}
\end{center}

$\{1, 5\}$ 和 $\{2, 6, 7\}$ 都挑 $1$--$2$，加入 $1$--$2$、$3$--$7$ 兩條，只剩一個元件，結束。 兩
回合，總共 $7$ 條邊、總權重 $39$，和前兩節的結果相同。 $n = 8$ 時定理~\ref{thm:mst-rounds}
的上限是 $3$ 回合，這張圖兩回合就做完了。

\subsection{C 程式}

下面的程式照虛擬碼~\ref{alg:mst-sollin} 執行，每回合先依元件裡最小的頂點編號，印出每個元
件和它挑的邊，再印出這一回合真的加入的邊。 為了印出元件的成員，程式每次都重新掃一遍所
有頂點，$O(n^2)$，只是為了輸出好看；演算法本身不需要。 互斥集合只用路徑壓縮，因為這裡
的重點是挑邊。

\begin{example}[label=ex:mst-sollin]{c}{(sollin.c) Sollin 演算法：每個元件同時挑最輕的往外的邊}
#include <stdio.h>
#include <stdlib.h>

typedef struct {
    int u, v, w;    // 讀入時讓 u < v
} Edge;

// 一致的平手規則：先比權重，再比小端點，再比大端點
int lighter(const Edge *x, const Edge *y)
{
    if (x->w != y->w) {
        return x->w < y->w;
    }
    if (x->u != y->u) {
        return x->u < y->u;
    }
    return x->v < y->v;
}

int find_set(int *p, int x)
{
    int r = x;
    while (p[r] != r) {
        r = p[r];
    }
    while (p[x] != r) {
        int next = p[x];
        p[x] = r;
        x = next;
    }
    return r;
}

int main(void)
{
    int n, m;
    if (scanf("%d %d", &n, &m) != 2 || n <= 0 || m < 0) {
        return 1;
    }
    Edge *e = malloc((m > 0 ? m : 1) * sizeof *e);
    int *p = malloc(n * sizeof(int));
    int *cheap = malloc(n * sizeof(int));  // cheap[r]：根 r 的元件挑到的邊（索引）
    int ok = e != NULL && p != NULL && cheap != NULL;
    for (int i = 0; ok && i < m; i++) {
        int u, v, w;
        ok = scanf("%d %d %d", &u, &v, &w) == 3 && u >= 1 && u <= n
             && v >= 1 && v <= n;
        if (ok) {
            e[i].u = (u < v ? u : v) - 1;
            e[i].v = (u < v ? v : u) - 1;
            e[i].w = w;
        }
    }
    if (ok) {
        for (int x = 0; x < n; x++) {
            p[x] = x;
        }
        long long total = 0;
        int taken = 0, round = 0, progress = 1;
        while (taken < n - 1 && progress) {
            for (int x = 0; x < n; x++) {
                cheap[x] = -1;
            }
            // 每條邊兩端的元件各自比一次
            for (int i = 0; i < m; i++) {
                int a = find_set(p, e[i].u), b = find_set(p, e[i].v);
                if (a == b) {
                    continue;
                }
                if (cheap[a] < 0 || lighter(&e[i], &e[cheap[a]])) {
                    cheap[a] = i;
                }
                if (cheap[b] < 0 || lighter(&e[i], &e[cheap[b]])) {
                    cheap[b] = i;
                }
            }
            printf("round %d\n", ++round);
            // 依元件裡最小的頂點編號，印出每個元件和它挑的邊
            for (int x = 0; x < n; x++) {
                int r = find_set(p, x);
                int first = 1;
                for (int y = 0; y < x && first; y++) {
                    first = find_set(p, y) != r;
                }
                if (!first || cheap[r] < 0) {
                    continue;
                }
                printf("{");
                for (int y = x, sep = 0; y < n; y++) {
                    if (find_set(p, y) == r) {
                        printf(sep++ ? ",%d" : "%d", y + 1);
                    }
                }
                const Edge *c = &e[cheap[r]];
                printf("} picks %d-%d (%d)\n", c->u + 1, c->v + 1, c->w);
            }
            // 把挑到的邊加進來；兩個元件挑到同一條時，第二次會被略過
            progress = 0;
            printf("added:");
            for (int r = 0; r < n; r++) {
                if (cheap[r] < 0) {
                    continue;
                }
                const Edge *c = &e[cheap[r]];
                int a = find_set(p, c->u), b = find_set(p, c->v);
                if (a != b) {
                    p[b] = a;
                    total += c->w;
                    taken++;
                    progress = 1;
                    printf(" %d-%d", c->u + 1, c->v + 1);
                }
            }
            printf("\n");
        }
        if (taken == n - 1) {
            printf("MST weight %lld, %d edges\n", total, taken);
        } else {
            printf("not connected: forest weight %lld, %d edges\n", total, taken);
        }
    }
    free(e);
    free(p);
    free(cheap);
    return ok ? 0 : 1;
}
\end{example}
\iosample{8 13 \\ 1 2 7 \\ 2 3 12 \\ 3 4 14 \\ 5 6 9 \\ 6 7 4 \\ 7 8 13 \\ 1 5 1 \\ 2 6 6 \\ 3 7 11 \\ 4 8 3 \\ 1 6 10 \\ 2 7 5 \\ 3 8 8}{round 1 \\ \{1\} picks 1-5 (1) \\ \{2\} picks 2-7 (5) \\ \{3\} picks 3-8 (8) \\ \{4\} picks 4-8 (3) \\ \{5\} picks 1-5 (1) \\ \{6\} picks 6-7 (4) \\ \{7\} picks 6-7 (4) \\ \{8\} picks 4-8 (3) \\ added: 1-5 2-7 3-8 4-8 6-7 \\ round 2 \\ \{1,5\} picks 1-2 (7) \\ \{2,6,7\} picks 1-2 (7) \\ \{3,4,8\} picks 3-7 (11) \\ added: 1-2 3-7 \\ MST weight 39, 7 edges}

\texttt{cheap} 以每回合開始時的根當索引；加邊時元件會合併、根會改變，但 \texttt{cheap} 記
的是邊的索引，加邊前重新 \texttt{find\_set} 檢查兩端，所以不受影響。 \texttt{progress} 處
理圖不連通的情形：某一回合沒有任何元件挑得到邊，就停下來，印出最小生成森林的權重。
如果比較只看權重，每個元件留下的是它最先看到的那條，挑到哪一條會隨邊的輸入順序改變，手
算的答案就可能對不上。

\section{判定、次小生成樹、瓶頸與最短路徑樹}
\label{sec:mst-more}

這一節整理幾個常和最小生成樹一起考的觀念。 下面 $T$ 是一棵生成樹；對任兩個頂點 $x$、$y$，
$T$ 裡 $x$ 到 $y$ 恰好有一條簡單路徑，記成 $P_T(x, y)$，路徑上最重的邊的權重記成
$\mathit{max}_T(x, y)$。

\subsection{怎麼判斷一棵生成樹是不是最小的}

\begin{theorem}
\label{thm:mst-check}
生成樹 $T$ 是最小生成樹，若且唯若每一條不在 $T$ 裡的邊 $(x, y)$ 都滿足
\[
w(x, y) \ge \mathit{max}_T(x, y),
\]
也就是說，把任何一條非樹邊加進 $T$ 形成的環上，這條非樹邊是最重的（可以平手）。
\end{theorem}

\begin{proof}[$\Rightarrow$ 的證明]
假設某條非樹邊 $(x, y)$ 的權重小於 $P_T(x, y)$ 上某條邊 $f$。 把 $(x, y)$ 加進 $T$、拿掉 $f$：
$f$ 在 $x$ 到 $y$ 的路徑上，拿掉後 $x$、$y$ 分在兩塊，$(x, y)$ 把兩塊接回來，得到一棵生成
樹，總權重比 $w(T)$ 小，$T$ 就不是最小的。
\end{proof}

\begin{whybox}[定理~\ref{thm:mst-check} 的 $\Leftarrow$]
設 $T$ 滿足條件。 在所有最小生成樹裡，取一棵和 $T$ 共同的邊最多的，叫它 $T^*$。 如果 $T^* =
T$ 就證完了。 否則兩棵樹邊數一樣，$T^*$ 有一條不在 $T$ 裡的邊 $f = (x, y)$。

從 $T^*$ 拿掉 $f$，分成兩塊，$x$、$y$ 各在一塊。 $T$ 的路徑 $P_T(x, y)$ 從 $x$ 走到 $y$，一
定有一條邊 $e$ 跨越這兩塊。 $e$ 不是 $f$（$f$ 不在 $T$ 裡），而 $T^*$ 跨越這兩塊的邊只有 $f$，
所以 $e \notin T^*$。 由條件，$w(f) \ge \mathit{max}_T(x, y) \ge w(e)$。

令 $T'' = T^* - \{f\} \cup \{e\}$，它是生成樹，$w(T'') = w(T^*) - w(f) + w(e) \le w(T^*)$，所
以也是最小生成樹。 但 $T''$ 比 $T^*$ 多了一條和 $T$ 共同的邊 $e$（$f$ 本來就不是共同的），
和 $T^*$ 的取法矛盾。 所以 $T^* = T$，$T$ 是最小生成樹。
\end{whybox}

用這個定理檢查一棵樹，要知道每一對 $(x, y)$ 的 $\mathit{max}_T(x, y)$。 從每個頂點在 $T$ 上
做一次 BFS 或 DFS，沿路記下「到目前為止最重的邊」，一次 $\Theta(n)$，$n$ 次就把所有配對算
完，$\Theta(n^2)$；再掃一遍所有非樹邊，總共 $\Theta(n^2 + m)$。

在露營區驗證一下。 六條非樹邊各自形成的環：

\begin{center}
\small
\begin{tabular}{clcc}
\toprule
非樹邊 & $T$ 上的路徑 & 路徑上最重的邊 & 換進來的總權重 \\
\midrule
$2$--$6$（$6$） & $2, 7, 6$ & $2$--$7$（$5$） & $39 + 6 - 5 = 40$ \\
$5$--$6$（$9$） & $5, 1, 2, 7, 6$ & $1$--$2$（$7$） & $39 + 9 - 7 = 41$ \\
$1$--$6$（$10$） & $1, 2, 7, 6$ & $1$--$2$（$7$） & $39 + 10 - 7 = 42$ \\
$2$--$3$（$12$） & $2, 7, 3$ & $3$--$7$（$11$） & $39 + 12 - 11 = 40$ \\
$7$--$8$（$13$） & $7, 3, 8$ & $3$--$7$（$11$） & $39 + 13 - 11 = 41$ \\
$3$--$4$（$14$） & $3, 8, 4$ & $3$--$8$（$8$） & $39 + 14 - 8 = 45$ \\
\bottomrule
\end{tabular}
\end{center}

每一條非樹邊都比它的環上其他的邊重，所以 $T$ 是最小生成樹；而且都是\textbf{嚴格}比較重。
這也可以用來判斷最小生成樹是不是唯一（觀念題第~\ref{q:mst-unique-iff}~題）：

\begin{corollary}
\label{cor:mst-unique-iff}
最小生成樹 $T$ 是唯一的最小生成樹，若且唯若每一條非樹邊 $(x, y)$ 都滿足 $w(x, y) >
\mathit{max}_T(x, y)$。
\end{corollary}

\subsection{次小生成樹（觀念）}
\label{sec:mst-second}

\textbf{次小生成樹}（second-best minimum spanning tree）是除了最小生成樹以外，總權重最小
的生成樹。 上表最後一欄就是「換一條邊」得到的生成樹的總權重：把非樹邊 $(x, y)$ 換進來，同
時拿掉它的環上最重的樹邊，比拿掉環上其他樹邊都划算。

\begin{theorem}
\label{thm:mst-second}
如果 $G$ 的權重互不相同，$T$ 是它唯一的最小生成樹，那麼有一棵次小生成樹和 $T$ 只差一條
邊：它是 $T - \{f\} \cup \{(x, y)\}$，其中 $(x, y)$ 是某條非樹邊，$f$ 是 $P_T(x, y)$ 上最重的
邊。
\end{theorem}

證明用的也是換邊，這裡不證（CLRS 第 21 章章末的問題有完整的推導）。 由這個定理，求次小生
成樹的做法是：

\begin{enumerate}
\item 求出最小生成樹 $T$，$O(m \lg n)$。
\item 算出所有的 $\mathit{max}_T(x, y)$，$\Theta(n^2)$。
\item 對每一條非樹邊 $(x, y)$，算 $w(T) + w(x, y) - \mathit{max}_T(x, y)$，取最小的，$\Theta(m)$。
\end{enumerate}

露營區的次小生成樹總權重是 $40$，而且有兩棵：換進 $2$--$6$、拿掉 $2$--$7$，或換進
$2$--$3$、拿掉 $3$--$7$。 最小生成樹唯一，次小生成樹卻不一定唯一。

\begin{examnote}
「次小生成樹」有兩種定義，作答前先看題目用哪一種。 一種是「和最小生成樹不同的生成樹裡最
輕的」，最小生成樹不唯一時，它的權重就等於最小生成樹的權重；另一種是「總權重\textbf{嚴格
大於}最小生成樹的生成樹裡最輕的」。 用換邊法算第二種時，如果非樹邊和環上最重的樹邊一樣
重，要改拿環上嚴格比它輕的最重的邊。 另外，題目問「最小生成樹」和「最短路徑」的差別時，
常見的答法是舉一個小反例（下面的三角形），比長篇說明更有說服力。
\end{examnote}

\subsection{瓶頸：最小生成樹上的路徑讓最重的一段最輕}
\label{sec:mst-bottleneck}

一條路徑上最重的邊叫做它的\textbf{瓶頸}。 例如一輛貨車要從營地 $1$ 開到營地 $4$，每段路
有載重限制，它能載多重取決於路上限制最嚴的那一段。 把權重看成「需要的門檻」，問題就是：
在所有 $x$ 到 $y$ 的路徑裡，瓶頸最小是多少？ 最小生成樹直接給出答案。

\begin{theorem}
\label{thm:mst-bottleneck}
$T$ 是 $G$ 的一棵最小生成樹。 對任兩個頂點 $x \ne y$ 和 $G$ 裡任何一條 $x$ 到 $y$ 的路徑 $Q$，
$Q$ 上最重的邊至少是 $\mathit{max}_T(x, y)$。 也就是說，$T$ 上的路徑是瓶頸最小的路徑。
\end{theorem}

\begin{proof}
令 $e$ 是 $P_T(x, y)$ 上最重的邊。 從 $T$ 拿掉 $e$，分成兩塊，$x$、$y$ 各在一塊，令 $S$ 是 $x$
那一塊。 $e$ 一定是切割 $(S, V - S)$ 的輕邊：如果有一條跨越的邊 $f$ 比 $e$ 輕，$T - \{e\}
\cup \{f\}$ 就是更輕的生成樹。 路徑 $Q$ 從 $S$ 走到 $V - S$，至少經過一條跨越的邊 $f$，$w(f)
\ge w(e)$，所以 $Q$ 上最重的邊 $\ge w(e) = \mathit{max}_T(x, y)$。
\end{proof}

露營區裡 $1$ 到 $4$ 的樹上路徑是 $1, 2, 7, 3, 8, 4$，瓶頸是 $3$--$7$（$11$）。 任何一條從 $1$
到 $4$ 的路都要跨越切割 $(\{1, 2, 5, 6, 7\}, \{3, 4, 8\})$，跨越的邊是 $3$--$7$（$11$）、
$2$--$3$（$12$）、$7$--$8$（$13$），最輕也要 $11$。 實作時不一定要建出整棵樹：照 Kruskal 演
算法由輕到重加邊，$x$ 和 $y$ 第一次連通時加進來的那條邊，它的權重就是答案。

\subsection{最小生成樹不是最短路徑樹}
\label{sec:mst-spt}

從頂點 $r$ 出發，到每個頂點都走最短路徑（總權重最小的路徑），這些路徑合起來也是一棵生成
樹，叫做以 $r$ 為根的\textbf{最短路徑樹}，第~25~章的 Dijkstra 演算法會求它。 兩種樹都「讓
某個總和最小」，很容易混在一起，但它們最小化的東西不同：

\begin{itemize}
\item 最小生成樹讓\textbf{整棵樹的邊權重總和}最小，和根無關；
\item 最短路徑樹讓\textbf{從 $r$ 到每個頂點的路徑長度}各自最小，換一個根就可能換一棵樹。
\end{itemize}

\begin{center}
\begin{tabular}{ccc}
\begin{tikzpicture}[x=1.1cm, y=1.1cm]
\node[mst-v] (a) at (0,0) {$1$};
\node[mst-v] (b) at (2,0) {$2$};
\node[mst-v] (c) at (1,1.6) {$3$};
\draw[mst-e] (a) -- node[mst-w] {4} (b);
\draw[mst-e] (b) -- node[mst-w] {3} (c);
\draw[mst-e] (c) -- node[mst-w] {5} (a);
\end{tikzpicture}
&
\begin{tikzpicture}[x=1.1cm, y=1.1cm]
\node[mst-v] (a) at (0,0) {$1$};
\node[mst-v] (b) at (2,0) {$2$};
\node[mst-v] (c) at (1,1.6) {$3$};
\draw[mst-t] (a) -- node[mst-w] {4} (b);
\draw[mst-t] (b) -- node[mst-w] {3} (c);
\draw[mst-x] (c) -- node[mst-w] {5} (a);
\end{tikzpicture}
&
\begin{tikzpicture}[x=1.1cm, y=1.1cm]
\node[mst-v] (a) at (0,0) {$1$};
\node[mst-v] (b) at (2,0) {$2$};
\node[mst-v] (c) at (1,1.6) {$3$};
\draw[mst-t] (a) -- node[mst-w] {4} (b);
\draw[mst-x] (b) -- node[mst-w] {3} (c);
\draw[mst-t] (c) -- node[mst-w] {5} (a);
\end{tikzpicture}
\\
\small 原圖 & \small 最小生成樹，總權重 $7$ & \small 以 $1$ 為根的最短路徑樹，總權重 $9$
\end{tabular}
\diagcaption{最小生成樹和最短路徑樹不同的最小例子。 最小生成樹上 $1$ 到 $3$ 要走 $4 + 3 =
7$，最短路徑是直接走 $5$。}
\label{fig:mst-spt}
\end{center}

圖~\ref{fig:mst-spt} 是最小的反例：三角形的邊是 $3$、$4$、$5$。 最小生成樹丟掉最重的 $5$，
總權重 $7$；但從 $1$ 到 $3$，樹上要繞經 $2$ 走 $7$，直接走只要 $5$。 以 $1$ 為根的最短路徑
樹用 $1$--$2$ 和 $1$--$3$，總權重 $9$。 露營區也一樣：以 $1$ 為根的最短路徑樹有兩棵：$1$--$2$、
$1$--$5$、$2$--$3$、$2$--$7$、$7$--$8$、$8$--$4$ 這六條是共同的，頂點 $6$ 可以經由 $1$--$6$
或 $5$--$6$ 接進來（兩條路都是 $10$），總權重 $51$ 或 $50$，都比最小生成樹的 $39$ 重；反過來，最小生成樹上 $1$ 到 $4$ 要走 $7 + 5 + 11 + 8 + 3 = 34$，最短路徑只要 $28$。

幾個相關的事實，都可以直接從定義看出來：

\begin{itemize}
\item 最短路徑樹也是生成樹，所以它的總權重一定 $\ge$ 最小生成樹的總權重。
\item 每條邊的權重都加上同一個常數 $c$：每棵生成樹都恰好 $n - 1$ 條邊，總權重都加 $(n - 1)c$，
最小生成樹不變。 最短路徑則可能改變，因為邊數多的路徑加得多。
\item 把權重換成任何一個嚴格遞增的函數（例如權重都是正數時取平方）：Kruskal 演算法只看邊
的大小順序，順序不變，最小生成樹就不變。 最短路徑會變。
\item 有負權重的邊時，本章三個演算法都照樣正確（證明沒有用到權重的正負）；最短路徑遇到負
權重則要小心，第~25~章會講。
\item 要找總權重\textbf{最大}的生成樹，把權重全部變號再求最小生成樹，或者 Kruskal 改成由重
到輕排序。
\end{itemize}

\section{本章常見錯誤}
\label{sec:mst-pitfalls}

\begin{center}
\small
\begin{tabular}{p{0.25\linewidth}p{0.51\linewidth}p{0.14\linewidth}}
\toprule
錯誤 & 例子 & 見 \\
\midrule
挑最輕的 $n - 1$ 條邊 & 露營區最輕的七條包含成環的 $2$--$6$，卻沒有必選的 $3$--$7$（$11$） &
第~\ref{sec:mst-kruskal}~節 \\
用「兩端都碰過」判斷成環 & Kruskal 第 6 步的 $1$--$2$ 兩端都已經有邊，但在不同的樹裡，應該選 &
第~\ref{sec:mst-kruskal}~節 \\
Prim 的選邊順序寫成由小到大 & 露營區從 $1$ 開始，權重依序是 $1, 7, 5, 4, 11, 8, 3$，不是 $1, 3, 4,
\ldots$ & 第~\ref{sec:mst-prim}~節 \\
Prim 沒檢查 $v$ 還在不在 $Q$ 裡 & 已經在樹裡的頂點的 $\mathit{key}$、$\pi$ 被改掉，回傳的邊錯 &
虛擬碼~\ref{alg:mst-prim} \\
允許重複放入卻沒丟掉過期資料 & 頂點 $6$ 被加入四次，多出三條邊 & 虛擬碼~\ref{alg:mst-prim-lazy} \\
平手規則沒照題目 & 更新 $\mathit{key}$ 用 $\le$ 而題目用 $<$，頂點 $4$ 接到 $5$ 而不是 $3$ &
圖~\ref{fig:mst-ties} \\
Sollin 平手各挑各的 & 三條權重相同的三角形，三個頂點挑的三條邊形成環 &
圖~\ref{fig:mst-sollin-cycle} \\
Sollin 同一條邊算兩次 & 露營區第一回合 $1$--$5$ 被兩端都挑到，總權重多算 $1$ & 第~\ref{sec:mst-sollin}~節 \\
比較函式用減法 & \texttt{return x->w - y->w;} 權重相差超過 \texttt{int} 範圍時溢位 &
範例~\ref{ex:mst-kruskal} \\
總權重用 \texttt{int} & $10^5$ 條邊、每條 $10^6$，總和超過 $2^{31}$ & 範例~\ref{ex:mst-kruskal} \\
大的稀疏圖用陣列版 Prim & $n = 10^5$ 時要 $10^{10}$ 步 & 第~\ref{sec:mst-prim}~節 \\
把最小生成樹當最短路徑樹 & 三角形 $3, 4, 5$ 裡，樹上 $1$ 到 $3$ 要走 $7$，最短是 $5$ & 圖~\ref{fig:mst-spt} \\
\bottomrule
\end{tabular}
\end{center}

\section{觀念題}
\label{sec:mst-exercises}

題型和研究所考試、期中期末考的筆試相近。 沒有特別說明時，圖是連通的簡單無向圖；平手時照
第~\ref{sec:mst-ties}~節的 $\prec$（先比權重，再比小端點，再比大端點，字母照字母順序）；Prim
演算法照虛擬碼~\ref{alg:mst-prim}。 解答在本章最後。

第~\ref{q:mst-kruskal-h}--\ref{q:mst-spt-h} 題用同一張圖 $H$：七個頂點 $a$ 到 $g$、$12$ 條邊。

\begin{center}
\begin{minipage}{\linewidth}
\centering
\begin{tikzpicture}[x=1.1cm, y=1.1cm]
\node[mst-v] (v1) at (-2,0) {$a$};
\node[mst-v] (v2) at (-1,1.7) {$b$};
\node[mst-v] (v3) at (1,1.7) {$c$};
\node[mst-v] (v4) at (2,0) {$d$};
\node[mst-v] (v6) at (1,-1.7) {$f$};
\node[mst-v] (v5) at (-1,-1.7) {$e$};
\node[mst-v] (v7) at (0,0) {$g$};
\draw[mst-e] (v1) -- node[mst-w] {3} (v2);
\draw[mst-e] (v2) -- node[mst-w] {9} (v3);
\draw[mst-e] (v3) -- node[mst-w] {5} (v4);
\draw[mst-e] (v4) -- node[mst-w] {2} (v6);
\draw[mst-e] (v5) -- node[mst-w] {5} (v6);
\draw[mst-e] (v1) -- node[mst-w] {3} (v5);
\draw[mst-e] (v1) -- node[mst-w] {5} (v7);
\draw[mst-e] (v2) -- node[mst-w] {4} (v7);
\draw[mst-e] (v3) -- node[mst-w] {1} (v7);
\draw[mst-e] (v4) -- node[mst-w] {8} (v7);
\draw[mst-e] (v5) -- node[mst-w] {1} (v7);
\draw[mst-e] (v6) -- node[mst-w] {6} (v7);
\end{tikzpicture}
\diagcaption{觀念題用的圖 $H$。}
\label{fig:mst-h}
\end{minipage}
\end{center}

\begin{enumerate}
\item \label{q:mst-choice}（選擇，可能不只一個答案）$G$ 是連通的加權無向圖，至少有三條邊。 下
列哪些敘述正確？
\begin{enumerate}
\item 權重互不相同時，權重第二小的邊一定在最小生成樹裡。
\item 權重互不相同時，權重第三小的邊一定在最小生成樹裡。
\item 每個頂點連出去的邊裡最重的那一條，一定不在任何最小生成樹裡。
\item 每條邊的權重都加上同一個常數 $c$（$c$ 可以是負數），最小生成樹的邊集合不變。
\item 權重有正有負時，把每個權重換成它的平方，最小生成樹的邊集合不變。
\item Prim 演算法在有負權重的圖上仍然正確。
\item 最小生成樹上任兩點之間的路徑，就是兩點之間的最短路徑。
\item 最小生成樹只有一棵，則權重一定互不相同。
\item Sollin 演算法某一回合開始時有 $k \ge 2$ 個元件，這一回合至少加入 $\lceil k/2 \rceil$ 條邊。
\end{enumerate}

\item \label{q:mst-kruskal-h}（追蹤）對 $H$ 執行 Kruskal 演算法。 依序寫出看過的每一條邊、選或
略過（略過的寫原因），以及最小生成樹的總權重。

\item \label{q:mst-prim-h}（追蹤）對 $H$ 從 $a$ 開始執行 Prim 演算法。 (a) 照本章的格式寫出每一
回合的 $\mathit{key}/\pi$ 表和加入的邊。 (b) 有哪幾個回合遇到 $\mathit{key}$ 平手？ 平手規則
決定了什麼？ (c) \proc{Decrease-Key} 一共做了幾次？

\item \label{q:mst-lazy-h}（計算）對 $H$ 從 $a$ 開始執行虛擬碼~\ref{alg:mst-prim-lazy}（三元組先
比權重、再比頂點字母）。 \proc{Insert} 一共做了幾次？ 取出的資料有幾筆是過期的？ 堆積裡最
多同時有幾筆？ 和上一題的 \proc{Decrease-Key} 次數比較。

\item \label{q:mst-sollin-h}（追蹤）對 $H$ 執行 Sollin 演算法。 (a) 寫出每一回合每個元件挑的邊和
這一回合加入的邊。 (b) 第一回合頂點 $g$ 的兩條權重 $1$ 的邊平手。 如果 $g$ 改挑另一條，這一
回合加入的邊會不會不同？

\item \label{q:mst-count-h}（計算）(a) $H$ 有幾棵最小生成樹？ 全部列出。 (b) Kruskal 演算法在權
重相同時改成「照 $\prec$ 倒過來的順序」，會得到哪一棵？

\item \label{q:mst-second-h}（計算）令 $T$ 是第~\ref{q:mst-kruskal-h}~題的最小生成樹。 (a) 對每
一條非樹邊，寫出它和 $T$ 形成的環上最重的樹邊，以及換進來之後的總權重。 (b) 依「和 $T$ 不同
的生成樹裡最輕的」的定義，次小生成樹的總權重是多少？ (c) 依「總權重嚴格大於最小生成樹」的定
義呢？

\item \label{q:mst-spt-h}（計算）(a) 寫出 $H$ 裡從 $a$ 到每個頂點的最短距離 $\delta(a, v)$。 (b)
以 $a$ 為根的最短路徑樹有幾棵？ 它們的總權重各是多少？ (c) 其中有沒有哪一棵也是最小生成樹？
(d) 在第~\ref{q:mst-kruskal-h}~題的最小生成樹上，$a$ 到 $f$ 的路徑長是多少？ 和 $\delta(a, f)$
比較。

\item \label{q:mst-two-light}（證明）$G$ 是權重互不相同的連通圖，$n \ge 3$。 (a) 證明權重最小和
第二小的兩條邊都在最小生成樹裡。 (b) 舉一個例子，權重第三小的邊不在最小生成樹裡。 (c) 對每個頂
點，取它連出去的邊裡最輕的一條。 證明這些邊都在最小生成樹裡，而且其中不同的邊至少有 $\lceil
n/2 \rceil$ 條。

\item \label{q:mst-unique-iff}（證明）證明系理~\ref{cor:mst-unique-iff}：最小生成樹 $T$ 是唯一
的，若且唯若每一條非樹邊 $(x, y)$ 都滿足 $w(x, y) > \mathit{max}_T(x, y)$。（提示：$\Leftarrow$
可以模仿定理~\ref{thm:mst-check} 的證明，或用環性質。）

\item \label{q:mst-max}（證明）(a) Kruskal 演算法改成依權重由大到小排序，其他不變。 證明它得
到總權重最大的生成樹。 (b) 一條路徑上最輕的邊叫它的「寬度」（例如管線的最小管徑）。 證明：
在總權重最大的生成樹上，任兩點之間的路徑，是兩點之間所有路徑裡寬度最大的。

\item \label{q:mst-dense}（選擇與計算）下面三種情況，各用本章哪一種寫法最合適？ 估計大約要幾
步。 (a) $n = 3000$ 的完全圖。 (b) $n = 2 \times 10^5$、$m = 4 \times 10^5$ 的道路網。 (c) $n =
10^6$、$m = 3 \times 10^6$，邊已經依權重由小到大排好。

\item \label{q:mst-rounds}（計算）(a) $n$ 個頂點排成一條路徑 $1, 2, \ldots, n$，邊 $(i, i + 1)$
的權重是 $i$。 Sollin 演算法要幾回合？ (b) 畫一張 $8$ 個頂點的連通圖，讓 Sollin 演算法恰好要
$3$ 回合（定理~\ref{thm:mst-rounds} 的上限），並寫出每回合加入的邊。
\end{enumerate}

\section{實作練習}
\label{sec:mst-practice}

下面是 Aki-Judge 上和本章相關的題目。 Sollin 演算法目前沒有題目，可以把
範例~\ref{ex:mst-sollin} 的輸出拿掉，改成只印總權重，拿前兩題的測資驗證。

\begin{center}
\begin{tabular}{llc}
\toprule
小節 & 題目 & 難度 \\
\midrule
Kruskal & \akiref{total-weight-of-a-minimum-spanning-tree}{最小生成樹的總權重} & 基本 \\
Kruskal & \akiref{mst-total-weight-on-a-connected-graph}{連通圖的最小生成樹總權重} & 基本 \\
Kruskal & \akiref{compare-two-kruskal-runs}{兩次 Kruskal 的總權重比較} & 基本 \\
Prim & \akiref{minimum-spanning-tree-grown-from-vertex-1}{從頂點 1 擴張的最小生成樹} & 基本 \\
瓶頸 & \akiref{smallest-threshold-connecting-first-and-last-vertex}{使頂點 1 與 n 連通的最小門檻} & 進階 \\
\bottomrule
\end{tabular}
\end{center}

幾題的提示：

\begin{itemize}
\item 「最小生成樹的總權重」和「連通圖的最小生成樹總權重」都是範例~\ref{ex:mst-kruskal} 拿掉
逐邊輸出；前一題的 $n$、$m$ 比較大，排序和互斥集合都要用本章的寫法。 總和用 \texttt{long
long}。
\item 「兩次 Kruskal 的總權重比較」要把 Kruskal 寫成一個函式，每次呼叫前重新初始化互斥集合。
用同一份邊陣列呼叫兩次時，第二次拿到的是已經排好的陣列，結果一樣。
\item 「從頂點 1 擴張的最小生成樹」的 $n$ 最多 $2000$，範例~\ref{ex:mst-prim-heap} 和
範例~\ref{ex:mst-prim-array} 都夠快；每回合重新掃過所有的邊找最輕的寫法是 $O(nm)$，才會太
慢。 輸入有自環，兩個範例都已經處理。
\item 「使頂點 1 與 n 連通的最小門檻」問的是瓶頸，不是最短路徑。 照
定理~\ref{thm:mst-bottleneck} 後面的說明，Kruskal 由輕到重加邊，$1$ 和 $n$ 第一次連通時加進來
的那條邊的權重就是答案；$n = 1$ 時答案是 $0$。
\end{itemize}

\section{讀完你應該會}
\label{sec:mst-checklist}

\begin{itemize}
\item 說出切割、輕邊、安全的邊的定義，用換邊論證證明切割性質和環性質。
\item 手算 Kruskal 演算法的選邊表（跳過的邊寫出原因），以及 Prim 演算法的 $\mathit{key}/\pi$
表，說明兩者選邊的順序為什麼不同。
\item 比較 Prim 演算法三種寫法（帶索引堆積、允許重複放入、陣列）的時間和空間，判斷稠密圖、
稀疏圖各該用哪一種。
\item 手算 Sollin 演算法每回合每個元件挑的邊，說明平手規則為什麼要一致，以及回合數為什麼不
超過 $\lfloor \lg n \rfloor$。
\item 證明權重互不相同時最小生成樹唯一，用非樹邊的條件判斷一棵生成樹是不是（唯一的）最小生
成樹，用換邊求次小生成樹。
\item 舉反例說明最小生成樹不是最短路徑樹，知道最小生成樹上的路徑是瓶頸最小的路徑。
\end{itemize}

\textbf{延伸閱讀}（書末會整理完整的參考資料）

\begin{itemize}
\item Thomas H. Cormen 等，\textit{Introduction to Algorithms}，第 4 版，第 21 章：用安全的邊
統一說明 Kruskal 和 Prim，章末的問題討論次小生成樹、瓶頸生成樹和稀疏圖的加速。
\item Joseph B. Kruskal，``On the Shortest Spanning Subtree of a Graph and the Traveling
Salesman Problem''，\textit{Proceedings of the American Mathematical Society} 7(1)，1956：
Kruskal 演算法的原始論文，只有三頁。
\item Robert C. Prim，``Shortest Connection Networks and Some Generalizations''，\textit{Bell
System Technical Journal} 36(6)，1957。 同一個方法更早由 Vojtěch Jarník 在 1930 年發表，所以
也叫 Jarník 演算法。
\item Otakar Borůvka，``O jistém problému minimálním''，\textit{Práce moravské
přírodovědecké společnosti} 3，1926：Sollin 演算法的最早版本（捷克文）。
\item R. L. Graham、Pavol Hell，``On the History of the Minimum Spanning Tree Problem''，
\textit{Annals of the History of Computing} 7(1)，1985：三個演算法誰先誰後的完整考證。
\item Robert E. Tarjan，\textit{Data Structures and Network Algorithms}，SIAM，1983，第 6 章：
用「藍色規則」（切割性質）和「紅色規則」（環性質）統一所有最小生成樹演算法。
\item David R. Karger、Philip N. Klein、Robert E. Tarjan，``A Randomized Linear-Time
Algorithm to Find Minimum Spanning Trees''，\textit{Journal of the ACM} 42(2)，1995：期望時間
$O(m)$ 的演算法，用到 Sollin 的回合和第~\ref{sec:mst-more}~節的判定定理。
\item Steven S. Skiena，\textit{The Algorithm Design Manual}，第 3 版，8.1 節：最小生成樹的
應用和變形，偏重怎麼把實際問題換成最小生成樹。
\end{itemize}

\addcontentsline{toc}{section}{觀念題解答}
\section*{觀念題解答}

\begin{enumerate}
\item (a)、(d)、(f)、(i) 正確。 (a)：最輕的兩條邊不會形成環（環至少三條邊），Kruskal 演算法
一定選它們；權重互不相同時最小生成樹唯一，所以它們在最小生成樹裡。 (b)：三角形的三條邊權重
$1, 2, 3$，第三小的是環上最重的，不選。 (c)：一個度數 $1$ 的頂點只有一條邊，它既是最重也是最
輕，一定要選。 (d)：每棵生成樹都恰好 $n - 1$ 條邊，總權重都加 $(n - 1)c$，大小順序不變。 (e)：
三角形權重 $-3, 1, 2$，原本選 $-3, 1$；平方後是 $9, 1, 4$，改選 $1, 4$，也就是原本的 $1, 2$。
(f)：定理~\ref{thm:mst-prim} 的證明沒有用到權重的正負。 (g)：圖~\ref{fig:mst-spt}。 (h)：
第~\ref{sec:mst-props}~節權重 $1, 1, 2$ 的三角形。 (i)：$k$ 個元件各挑一條邊，每條邊最多被兩
個元件挑到，所以不同的邊至少 $\lceil k/2 \rceil$ 條；由定理~\ref{thm:mst-sollin}，它們都會被
加入。

\item 照 $\prec$ 排序：$c$--$g$（$1$）、$e$--$g$（$1$）、$d$--$f$（$2$）、$a$--$b$（$3$）、
$a$--$e$（$3$）、$b$--$g$（$4$）、$a$--$g$（$5$）、$c$--$d$（$5$）、$e$--$f$（$5$）、
$f$--$g$（$6$）、$d$--$g$（$8$）、$b$--$c$（$9$）。

\begin{center}
\small
\begin{tabular}{rcrl}
\toprule
步 & 邊 & 權重 & 結果 \\
\midrule
1 & $c$--$g$ & 1 & 選 \\
2 & $e$--$g$ & 1 & 選 \\
3 & $d$--$f$ & 2 & 選 \\
4 & $a$--$b$ & 3 & 選 \\
5 & $a$--$e$ & 3 & 選（$\{a, b\}$ 和 $\{c, e, g\}$ 合併） \\
6 & $b$--$g$ & 4 & 略過：$b$、$g$ 已經經過 $a, e$ 連通 \\
7 & $a$--$g$ & 5 & 略過：$a$、$g$ 已經經過 $e$ 連通 \\
8 & $c$--$d$ & 5 & 選，已滿 $6$ 條 \\
\bottomrule
\end{tabular}
\end{center}
後面四條不用看。 最小生成樹是 $\{c\text{--}g, e\text{--}g, d\text{--}f, a\text{--}b, a\text{--}e,
c\text{--}d\}$，總權重 $1 + 1 + 2 + 3 + 3 + 5 = 15$。

\item (a)
\begin{center}
\footnotesize
\setlength{\tabcolsep}{3.2pt}
\setlength{\fboxsep}{1.2pt}
\begin{tabular}{cccrrrrrrr}
\toprule
回合 & 取出 & 加入的邊 & $a$ & $b$ & $c$ & $d$ & $e$ & $f$ & $g$ \\
\midrule
1 & $a$ & — & \mstsel{$0$} & \mstchg{$3/a$} & $\infty$ & $\infty$ & \mstchg{$3/a$} & $\infty$ & \mstchg{$5/a$} \\
2 & $b$ & $a$--$b$ &  & \mstsel{$3/a$} & \mstchg{$9/b$} & $\infty$ & $3/a$ & $\infty$ & \mstchg{$4/b$} \\
3 & $e$ & $a$--$e$ &  &  & $9/b$ & $\infty$ & \mstsel{$3/a$} & \mstchg{$5/e$} & \mstchg{$1/e$} \\
4 & $g$ & $e$--$g$ &  &  & \mstchg{$1/g$} & \mstchg{$8/g$} &  & $5/e$ & \mstsel{$1/e$} \\
5 & $c$ & $g$--$c$ &  &  & \mstsel{$1/g$} & \mstchg{$5/c$} &  & $5/e$ &  \\
6 & $d$ & $c$--$d$ &  &  &  & \mstsel{$5/c$} &  & \mstchg{$2/d$} &  \\
7 & $f$ & $d$--$f$ &  &  &  &  &  & \mstsel{$2/d$} &  \\
\bottomrule
\end{tabular}
\end{center}
加入的邊依序是 $a$--$b$、$a$--$e$、$e$--$g$、$g$--$c$、$c$--$d$、$d$--$f$，總權重 $15$，和上一題
同一棵樹。 (b) 第 2 回合 $b$、$e$ 的 $\mathit{key}$ 都是 $3$，取字母小的 $b$；這只影響順序，兩
條邊最後都會選。 第 6 回合 $d$、$f$ 的 $\mathit{key}$ 都是 $5$，取 $d$，於是加入 $c$--$d$，之後
$f$ 改經由 $d$--$f$（$2$）接進來；如果先取 $f$，加入的是 $e$--$f$，$d$ 再經由 $d$--$f$ 接進
來，得到另一棵最小生成樹（第~\ref{q:mst-count-h}~題）。 (c) 表裡有底色的格子共 $11$ 個，也就
是 $11$ 次；其中第一次把 $\infty$ 改小的有 $6$ 次，真正把有限的值再改小的有 $5$ 次（$g$ 兩次，
$c$、$d$、$f$ 各一次）。

\item \proc{Insert} 共 $13$ 次（$12$ 條邊各一次，加上起點）。 取出的順序是 $(0, a)$、$(3, b)$、
$(3, e)$、$(1, g)$、$(1, c)$、$(4, g)$ 過期、$(5, d)$、$(2, f)$，之後 $(5, f)$、$(5, g)$、$(6, f)$、
$(8, d)$、$(9, c)$ 都過期，共 $6$ 筆過期（加入第 $7$ 個頂點後就可以停，不用取出最後五筆）。 堆
積最多同時有 $7$ 筆，發生在取出 $g$、放入 $(1, c)$、$(8, d)$、$(6, f)$ 之後。 帶索引的版本做
了 $11$ 次 \proc{Decrease-Key}，堆積最多 $7$ 個頂點；這裡多做了兩次插入，堆積大小差不多。 圖
大的時候差別在空間：一個是 $\Theta(m)$，一個是 $\Theta(n)$。

\item (a) 第一回合：
\begin{center}
\small
\begin{tabular}{ccccccc}
\toprule
$\{a\}$ & $\{b\}$ & $\{c\}$ & $\{d\}$ & $\{e\}$ & $\{f\}$ & $\{g\}$ \\
\midrule
$a$--$b$（$3$） & $a$--$b$（$3$） & $c$--$g$（$1$） & $d$--$f$（$2$） & $e$--$g$（$1$） & $d$--$f$（$2$） &
$c$--$g$（$1$） \\
\bottomrule
\end{tabular}
\end{center}
$a$ 的 $a$--$b$、$a$--$e$ 平手，$(3, a, b) \prec (3, a, e)$；$g$ 的 $c$--$g$、$e$--$g$ 平手，取
$c$--$g$。 加入 $a$--$b$、$c$--$g$、$d$--$f$、$e$--$g$，剩三個元件 $\{a, b\}$、$\{c, e, g\}$、
$\{d, f\}$。 第二回合：$\{a, b\}$ 挑 $a$--$e$（$3$）；$\{c, e, g\}$ 也挑 $a$--$e$（$3$）；
$\{d, f\}$ 的 $c$--$d$、$e$--$f$ 都是 $5$，$(5, c, d) \prec (5, e, f)$，挑 $c$--$d$。 加入
$a$--$e$、$c$--$d$，結束，總權重 $15$，和第~\ref{q:mst-kruskal-h}~題同一棵（定理~\ref{thm:mst-sollin}）。
(b) 不會。 $c$ 自己挑 $c$--$g$、$e$ 自己挑 $e$--$g$，兩條本來就都會加入，$g$ 挑哪一條都只是重
複。

\item (a) 兩棵。 $c$--$g$、$e$--$g$、$d$--$f$、$a$--$b$、$a$--$e$ 是兩棵共同的：前三條是各自
頂點唯一最輕的邊；$a$--$b$ 是 $b$ 和其他頂點之間的切割裡唯一最輕的（$b$ 的邊是 $3, 4, 9$）；
$a$--$e$ 是切割 $(\{a, b\}, \{c, \ldots, g\})$ 唯一最輕的（跨越的邊是 $3, 4, 5, 9$）。 剩下要把
$\{d, f\}$ 接上 $\{a, b, c, e, g\}$，跨越的邊是 $c$--$d$（$5$）、$e$--$f$（$5$）、$f$--$g$（$6$）、
$d$--$g$（$8$），選 $c$--$d$ 或 $e$--$f$ 都是總權重 $15$。 兩棵分別是第~\ref{q:mst-kruskal-h}~題
的樹，以及把其中的 $c$--$d$ 換成 $e$--$f$。 (b) 權重 $5$ 的三條邊順序倒過來，變成 $e$--$f$、
$c$--$d$、$a$--$g$，先看到 $e$--$f$ 就選它，得到第二棵。

\item (a)
\begin{center}
\small
\begin{tabular}{clcc}
\toprule
非樹邊 & $T$ 上的路徑 & 環上最重的樹邊 & 換進來的總權重 \\
\midrule
$b$--$g$（$4$） & $b, a, e, g$ & $a$--$b$ 或 $a$--$e$（$3$） & $16$ \\
$a$--$g$（$5$） & $a, e, g$ & $a$--$e$（$3$） & $17$ \\
$e$--$f$（$5$） & $e, g, c, d, f$ & $c$--$d$（$5$） & $15$ \\
$f$--$g$（$6$） & $f, d, c, g$ & $c$--$d$（$5$） & $16$ \\
$d$--$g$（$8$） & $d, c, g$ & $c$--$d$（$5$） & $18$ \\
$b$--$c$（$9$） & $b, a, e, g, c$ & $a$--$b$ 或 $a$--$e$（$3$） & $21$ \\
\bottomrule
\end{tabular}
\end{center}
(b) $15$：換進 $e$--$f$、拿掉 $c$--$d$，得到另一棵最小生成樹。 (c) $16$。 $e$--$f$ 和環上最重的
$c$--$d$ 一樣重，照作答提醒改拿環上嚴格比 $5$ 輕的最重的邊 $d$--$f$（$2$），得 $15 + 5 - 2 =
18$；其他各列不變。 最小的是 $16$，從 $T$ 換一條邊得到的有三棵：換進 $b$--$g$ 拿掉 $a$--$b$ 或 $a$--$e$，或換進
$f$--$g$ 拿掉 $c$--$d$。

\item (a) $\delta(a, b) = 3$、$\delta(a, e) = 3$、$\delta(a, g) = 4$（$a, e, g$）、$\delta(a, c) =
5$（$a, e, g, c$）、$\delta(a, f) = 8$（$a, e, f$）、$\delta(a, d) = 10$。 (b) 兩棵。 $b, e, g, c,
f$ 的最短路徑都只有一條；$d$ 可以經過 $c$（$5 + 5$）或經過 $f$（$8 + 2$）。 共同的邊是
$a$--$b$、$a$--$e$、$e$--$g$、$g$--$c$、$e$--$f$，總權重 $13$，再加 $c$--$d$ 是 $18$，加
$d$--$f$ 是 $15$。 (c) 有。 加 $d$--$f$ 的那一棵總權重 $15$，正是第~\ref{q:mst-count-h}~題的第
二棵最小生成樹。 這是巧合，不是規則：圖~\ref{fig:mst-spt} 的最短路徑樹就比最小生成樹重。 (d)
第~\ref{q:mst-kruskal-h}~題的樹上 $a$ 到 $f$ 走 $a, e, g, c, d, f$，長 $3 + 1 + 1 + 5 + 2 = 12$，
比 $\delta(a, f) = 8$ 長。

\item (a) 最輕的邊 $e_1$ 自己不會形成環，Kruskal 演算法第一個選它。 第二輕的 $e_2$ 和 $e_1$ 最多
共用一個端點，兩條邊不會形成環（簡單圖的環至少三條邊），所以第二個也選。 權重互不相同，最小
生成樹唯一（定理~\ref{thm:mst-unique}），就是 Kruskal 選的那一棵。 (b) 三角形，權重 $1, 2, 3$；
或任何一個三條最輕的邊剛好圍成三角形的圖。 (c) 頂點 $v$ 最輕的邊是切割 $(\{v\}, V - \{v\})$ 唯
一的輕邊，由系理~\ref{cor:mst-cut-strict}，它在最小生成樹裡。 $n$ 個頂點各取一條，一條邊最多
被它的兩個端點取到，所以不同的邊至少 $\lceil n/2 \rceil$ 條。

\item $\Rightarrow$：假設某條非樹邊 $(x, y)$ 有 $w(x, y) \le \mathit{max}_T(x, y)$。 由
定理~\ref{thm:mst-check}，$T$ 是最小生成樹，所以 $w(x, y) \ge \mathit{max}_T(x, y)$，兩者相等。
令 $f$ 是 $P_T(x, y)$ 上權重等於 $w(x, y)$ 的樹邊，$T - \{f\} \cup \{(x, y)\}$ 是另一棵生成樹，
總權重相同，也是最小生成樹，和唯一矛盾。

$\Leftarrow$：設每條非樹邊都嚴格比 $\mathit{max}_T$ 重，而 $T' \ne T$ 也是最小生成樹。 取一
條 $f = (x, y) \in T' - T$。 從 $T'$ 拿掉 $f$ 分成兩塊，$P_T(x, y)$ 有一條樹邊 $e$ 跨越這兩塊，
$e \notin T'$。 由假設 $w(f) > \mathit{max}_T(x, y) \ge w(e)$，$T' - \{f\} \cup \{e\}$ 比 $T'$ 輕，
矛盾。 所以 $T$ 是唯一的最小生成樹。

\item (a) 令 $w''(e) = -w(e)$。 由重到輕排 $w$ 就是由輕到重排 $w''$，演算法的每一步都和「在
$w''$ 上執行原本的 Kruskal」相同，所以得到 $w''$ 的最小生成樹 $T$。 對任何生成樹 $T_1$，
$w''(T) \le w''(T_1)$，也就是 $w(T) \ge w(T_1)$。 本章的證明沒有用到權重的正負，所以換成負的
權重也成立。 (b) 由 (a)，總權重最大的生成樹 $T$ 是 $w''$ 的最小生成樹。 對 $w''$ 用
定理~\ref{thm:mst-bottleneck}：任何 $x$ 到 $y$ 的路徑 $Q$ 上 $w''$ 最大的邊，至少是 $T$ 的路徑上
$w''$ 最大的邊。 $w''$ 最大就是 $w$ 最小，所以 $Q$ 上 $w$ 最小的邊，至多是 $T$ 的路徑上 $w$ 最小
的邊，也就是 $Q$ 的寬度不超過 $T$ 上路徑的寬度。

\item (a) 完全圖 $m = n(n - 1)/2 \approx 4.5 \times 10^6$，是稠密圖，用陣列版 Prim，大約 $2n^2 =
1.8 \times 10^7$ 步，不用存邊。 堆積版最壞大約 $m \lg n \approx 5 \times 10^7$；Kruskal 光排序
就要 $m \lg m \approx 10^8$ 次比較，還要先存下 $4.5 \times 10^6$ 條邊。 (b) 稀疏圖，用堆積版
Prim 或 Kruskal：前者大約 $(n + m) \lg n \approx 6 \times 10^5 \times 18 \approx 10^7$，後者排
序大約 $m \lg m \approx 7.4 \times 10^6$；陣列版要 $n^2 = 4 \times 10^{10}$ 步，不能用。 (c) 用
Kruskal，邊已經排好就省掉排序，只剩 $m$ 次左右的互斥集合操作，$O(m\,\alpha(n))$，大約幾百萬
步；Prim 至少要 $(n + m) \lg n \approx 8 \times 10^7$。

\item (a) $1$ 回合。 頂點 $1$ 只有邊 $(1, 2)$；頂點 $i \ge 2$ 的兩條邊權重是 $i - 1$ 和 $i$，挑
$(i - 1, i)$。 所以每條邊都被挑到，第一回合就加入全部 $n - 1$ 條邊。 這說明定理~\ref{thm:mst-rounds}
只是上限，回合數可以遠少於 $\lg n$。 (b) 一個例子：路徑 $1, 2, \ldots, 8$，七條邊依序的權重是
$1, 10, 2, 20, 3, 11, 4$（$(1, 2)$ 是 $1$、$(2, 3)$ 是 $10$、$(3, 4)$ 是 $2$，依此類推）。 第一回合每
個頂點挑它在 $\{1, 2\}$、$\{3, 4\}$、$\{5, 6\}$、$\{7, 8\}$ 裡那條邊，加入 $1$--$2$、$3$--$4$、
$5$--$6$、$7$--$8$。 第二回合 $\{1, 2\}$ 和 $\{3, 4\}$ 都挑 $2$--$3$（$10$，比 $4$--$5$ 的 $20$
輕），$\{5, 6\}$ 和 $\{7, 8\}$ 都挑 $6$--$7$（$11$），加入兩條，剩 $\{1, 2, 3, 4\}$、$\{5, 6, 7,
8\}$。 第三回合兩邊都挑 $4$--$5$。 共 $3$ 回合。
\end{enumerate}

\chapter{最短路徑}
\label{ch:spt}

% 本章的圖共用的樣式（前綴 spt-）。 追蹤表、快照圖和矩陣序列都由模擬程式照虛擬碼
% 逐步執行後產生，再和本章的 C 程式輸出對過。
\tikzset{
    spt-v/.style={circle, draw=black!70, fill=white, minimum size=0.62cm, inner sep=0pt, font=\small},
    spt-vs/.style={spt-v, minimum size=0.5cm, font=\scriptsize},
    spt-vsdone/.style={spt-vs, draw=coreStroke, fill=coreFill},
    spt-vsnew/.style={spt-vs, draw=memStroke, fill=memFill, very thick},
    spt-a/.style={-{Stealth[length=2.2mm]}, draw=black!60, thick},
    spt-neg/.style={spt-a, draw=optStroke},
    spt-sa/.style={-{Stealth[length=1.6mm]}, draw=black!35},
    spt-ta/.style={-{Stealth[length=2mm]}, draw=coreStroke, very thick},
    spt-tree/.style={-{Stealth[length=2.4mm]}, draw=coreStroke, line width=1.5pt},
    spt-w/.style={font=\scriptsize, inner sep=1pt, fill=white, text=black!85},
    spt-wn/.style={spt-w, text=optStroke},
    spt-d/.style={font=\scriptsize, inner sep=1.5pt, text=coreStroke},
    spt-h/.style={font=\scriptsize, inner sep=1.5pt, text=memStroke},
}
% 追蹤表裡的記號：這一輪取出的頂點、這一輪變小的值、已經確定的值
\newcommand{\sptpick}[1]{\textbf{#1}}
\newcommand{\sptnew}[1]{\underline{#1}}
\newcommand{\sptold}[1]{\textcolor{black!45}{#1}}

外送員從店裡出發，要把餐點送到城裡好幾個地方。 地圖上每一段路都標了騎車要幾分鐘，
有些路是單行道。 從店裡到每個地方，最快要幾分鐘？ 把路線一條一條列出來比較行不通，
路口一多，路線的數目會多到列不完。 換個角度想：離店裡最近的那個地方，它的最快走法
一定很短，很容易確定；確定了近的地方，能不能一步一步推到遠的地方？ 如果有些路段
走了反而能領到補貼，可以把它想成「負的分鐘數」，剛才的想法還成立嗎？ 往下讀之前，先
想一想。

這一章處理邊上有權重的圖。 第~21~章的 BFS 求的是邊數最少的路徑，相當於每條邊的權重都
是 $1$；權重不一樣時，要找的是權重總和最小的路徑。 先定下問題和一個所有演算法共用的
動作「鬆弛」，證明幾個鬆弛的性質，之後的每個演算法都用這些性質證明。 接著是三個單一
起點的演算法：權重不是負數時用 \textbf{Dijkstra 演算法}，它是第~19~章的貪心法，搭配
第~8~章的優先佇列；有負的權重時用 \textbf{Bellman--Ford 演算法}，它還能發現負環；圖沒有
環時，依第~22~章的拓樸順序鬆弛一遍就好。 最後是所有點對的最短路徑：第~18~章已經用動
態規劃推出 Floyd--Warshall 演算法的遞迴式，這一章補上路徑還原、負環偵測和手算的矩陣序
列，再講同一個架構求的遞移閉包，以及在稀疏圖上比較快的 Johnson 演算法。

這一章用到的前面內容：第~21~章的圖的名詞、鄰接串列和 \proc{Print-Path}，第~22~章的拓
樸排序，第~\ref{sec:dp-substructure}~節的最佳子結構與第~\ref{sec:dp-apsp}~節的 Floyd--Warshall
遞迴式，第~\ref{sec:hp-pq}~節的 \proc{Decrease-Key}。 程式方面用到《C 程式設計基礎》第~13~章的 \texttt{long long}、
第~19~章的 \texttt{malloc} 和鏈結串列。

本章的約定和第~21~章相同：圖有 $n$ 個頂點、$m$ 條邊，頂點在手算時編號 $1$ 到 $n$，鄰
接串列裡的鄰居由小到大排；C 程式裡的頂點從 $0$ 編號，讀入時減一、印出時加一。

\section{最短路徑問題與鬆弛}
\label{sec:spt-relax}

\subsection{路徑的權重與最短距離}

\textbf{加權有向圖}是一張有向圖 $G = (V, E)$，加上一個\textbf{權重函數} $w$，每條邊
$(u, v)$ 有一個實數權重 $w(u, v)$。 一條路徑 $p = \langle v_0, v_1, \ldots, v_k
\rangle$ 的\textbf{權重}是它經過的邊的權重總和：
\[
w(p) = \sum_{i=1}^{k} w(v_{i-1}, v_i).
\]
長度 $0$ 的路徑 $\langle u \rangle$ 權重是 $0$。 權重可以代表時間、距離、費用，名稱不重
要；這一章說「最短」都是指權重總和最小，不是邊數最少。

\begin{definition}[最短距離]
\label{def:spt-delta}
從 $u$ 到 $v$ 的\textbf{最短距離}是
\[
\delta(u, v) = \begin{cases}
\min\{\, w(p) : p \text{ 是從 } u \text{ 到 } v \text{ 的路徑} \,\}, & u \leadsto v, \\
\infty, & \text{從 } u \text{ 到不了 } v.
\end{cases}
\]
權重等於 $\delta(u, v)$ 的路徑叫從 $u$ 到 $v$ 的一條\textbf{最短路徑}。
\end{definition}

這和第~21~章的 $\delta(u, v)$ 是同一個符號：每條邊的權重都是 $1$ 時，路徑的權重就是邊數，
兩個定義一致，BFS（定理~\ref{thm:gtr-bfs}）就是這個特例的解法。 上面的 $\min$ 有一個
問題：路徑可以重複經過頂點，條數可能有無限多，最小值不一定存在。 下一小節處理這件事。

常見的問題有四種：
\begin{itemize}
\item \textbf{單一起點}（single-source）：給定起點 $s$，求 $s$ 到每個頂點的最短距離和路
徑。 本章第~\ref{sec:spt-dijkstra}～\ref{sec:spt-dag}~節都是這一種。
\item \textbf{單一終點}：求每個頂點到 $t$ 的最短距離。 把每條邊反向，就變成從 $t$ 出發的單
一起點問題。
\item \textbf{單一點對}：只要 $s$ 到 $t$。 目前還沒有已知的方法在最壞情況下比單一起點更快，
通常就跑單一起點的演算法，算到 $t$ 確定時提早停下。
\item \textbf{所有點對}（all-pairs）：每一對 $(u, v)$ 都要。 第~\ref{sec:spt-fw}、
\ref{sec:spt-johnson}~節。
\end{itemize}

無向圖的一條邊 $\{u, v\}$，看成兩條方向相反、權重相同的有向邊 $(u, v)$ 和 $(v, u)$，本章
的演算法就都能用。 但要小心：權重是負的無向邊，換成兩條有向邊後，$u \to v \to u$ 本身就
是一個權重為負的環，下一小節會看到這代表最短距離沒有意義。 所以本章說的「負邊」都是在
有向圖上。

\subsection{負邊與負環}

權重可以是負的。 外送平台在某些路段發補貼，或者財務上的「這一步會收到錢」，都可以寫成
負的權重。 負邊本身沒有問題，有問題的是負環。

\begin{definition}[負環]
\label{def:spt-negcycle}
權重總和小於 $0$ 的環叫\textbf{負環}（negative-weight cycle）。
\end{definition}

如果從 $s$ 走得到一個負環，而且從這個負環走得到 $v$，那麼到 $v$ 的路徑可以在負環上多繞
一圈，權重就再小一些，沒有最小值。 這時規定 $\delta(s, v) = -\infty$。
圖~\ref{fig:spt-negcycle} 是一個例子：從 $1$ 出發，$2 \to 3 \to 2$ 是權重 $-3 + 1 = -2$
的負環，$2$、$3$ 和從它們走得到的 $4$ 都是 $-\infty$；$5$ 從 $1$ 到不了，是 $\infty$；$1$
自己是 $0$。

\begin{center}
\begin{tikzpicture}[x=1cm, y=1cm]
\node[spt-v] (n5) at (0, 0) {5};
\node[spt-v] (n1) at (2, 0) {1};
\node[spt-v] (n2) at (4, 0) {2};
\node[spt-v] (n3) at (6, 0) {3};
\node[spt-v] (n4) at (8, 0) {4};
\draw[spt-a] (n5) -- node[spt-w, above] {$1$} (n1);
\draw[spt-a] (n1) -- node[spt-w, above] {$4$} (n2);
\draw[spt-neg] (n2) to[bend left=25] node[spt-wn, above] {$-3$} (n3);
\draw[spt-a] (n3) to[bend left=25] node[spt-w, below] {$1$} (n2);
\draw[spt-a] (n3) -- node[spt-w, above] {$2$} (n4);
\foreach \nd/\lab in {n5/\infty, n1/0, n2/-\infty, n3/-\infty, n4/-\infty}
    \node[spt-d, anchor=north] at ([yshift=-6pt]\nd.south) {$\lab$};
\end{tikzpicture}\par\nopagebreak
\diagcaption{從 $1$ 出發的最短距離（藍字）。 $2 \to 3 \to 2$ 是負環，負的邊畫成紅色。}
\label{fig:spt-negcycle}
\end{center}

\begin{lemma}
\label{lem:spt-simple}
如果從 $s$ 走得到的地方沒有負環，那麼對每個 $s$ 走得到的 $v$，$\delta(s, v)$ 是一個有限的
數，而且有一條從 $s$ 到 $v$ 的最短路徑是簡單路徑，最多 $n - 1$ 條邊。
\end{lemma}

\begin{proof}
任取一條從 $s$ 到 $v$ 的路徑 $p$。 如果 $p$ 經過某個頂點 $x$ 兩次，兩次之間繞的那一段是
一個環，這個環上的頂點都從 $s$ 走得到，依假設它的權重 $\ge 0$。 把這一段剪掉，還是一條從
$s$ 到 $v$ 的路徑，權重不會變大。 一直剪到沒有重複頂點為止，得到一條簡單路徑 $p'$，
$w(p') \le w(p)$。 所以只要在簡單路徑裡找最小值就好；簡單路徑最多 $n - 1$ 條邊，條數有限，
最小值存在。
\end{proof}

權重剛好是 $0$ 的環沒有關係：剪不剪掉權重都一樣，引理照樣成立。 本章的單一起點演算法都
只保證在「從 $s$ 走得到的地方沒有負環」時算出正確的 $\delta$；Bellman--Ford 演算法另外會
回報負環。

\subsection{最短路徑的子路徑還是最短路徑}

\begin{lemma}
\label{lem:spt-subpath}
設 $p = \langle v_0, v_1, \ldots, v_k \rangle$ 是從 $v_0$ 到 $v_k$ 的最短路徑。 對任何
$0 \le i \le j \le k$，$p$ 從 $v_i$ 到 $v_j$ 的那一段 $p_{ij}$ 是從 $v_i$ 到 $v_j$ 的最短
路徑。
\end{lemma}

\begin{proof}
這是第~\ref{sec:dp-substructure}~節的剪下貼上。 把 $p$ 拆成 $v_0 \leadsto v_i$、$p_{ij}$、$v_j \leadsto
v_k$ 三段，$w(p)$ 是三段的和。 如果有一條從 $v_i$ 到 $v_j$ 的路徑 $p'$ 比 $p_{ij}$ 輕，把
$p_{ij}$ 換成 $p'$，得到一條從 $v_0$ 到 $v_k$、比 $p$ 輕的路徑，和 $p$ 是最短路徑矛盾。
\end{proof}

第~18~章用邊數證過同一件事，那時候要把重複的頂點剪掉；這裡換來的路徑可以重複頂點，不必
剪，因為 $\delta$ 的定義本來就允許重複。 這個引理是本章所有演算法能「一段一段算」的原因。
另一個天天用到的事實是三角不等式：

\begin{lemma}[三角不等式]
\label{lem:spt-triangle}
對每條邊 $(u, v)$，$\delta(s, v) \le \delta(s, u) + w(u, v)$。
\end{lemma}

\begin{proof}
$s$ 到不了 $u$ 時右邊是 $\infty$。 否則任取一條從 $s$ 到 $u$ 的路徑，接上邊 $(u, v)$ 就是到
$v$ 的路徑，權重是那條路徑的權重加 $w(u, v)$；對所有到 $u$ 的路徑取最小值，就得到不等式
（$\delta(s, u) = -\infty$ 時兩邊都是 $-\infty$）。
\end{proof}

\subsection{鬆弛}

本章的單一起點演算法都替每個頂點 $v$ 記兩個值，和第~21~章 BFS 的 $d$、$\pi$ 用法一樣：
\begin{itemize}
\item $d[v]$：目前找到的、從 $s$ 到 $v$ 的最短路徑的權重，叫做\textbf{距離估計}。 它永遠是
$\delta(s, v)$ 的上界，演算法結束時等於 $\delta(s, v)$。
\item $\pi[v]$：目前那條路徑上 $v$ 的前一個頂點（\textbf{前驅}），沒有時是 \NIL。
\end{itemize}

\begin{pseudo}[label=alg:spt-init]{初始化單一起點}
\Procedure{Initialize-Single-Source}{$G, s$}
    \For{每個頂點 $v \in V$}
        \State $d[v] \gets \infty$；$\pi[v] \gets \NIL$
    \EndFor
    \State $d[s] \gets 0$
\EndProcedure
\end{pseudo}

\textbf{鬆弛}（relax）一條邊 $(u, v)$，是問：先走到 $u$ 再走這條邊，會不會比目前到 $v$ 的走
法更好？ 會的話就更新 $d[v]$ 和 $\pi[v]$。

\begin{pseudo}[label=alg:spt-relax]{鬆弛一條邊}
\Procedure{Relax}{$u, v, w$}
    \If{$d[u] + w(u, v) < d[v]$}
        \State $d[v] \gets d[u] + w(u, v)$
        \State $\pi[v] \gets u$
    \EndIf
\EndProcedure
\end{pseudo}

例如 $d[u] = 5$、$d[v] = 9$、$w(u, v) = 3$，因為 $5 + 3 < 9$，鬆弛之後 $d[v] = 8$、
$\pi[v] = u$。 若 $w(u, v) = 4$，$5 + 4 = 9$ 不小於 $9$，什麼都不改；比較用的是嚴格的
$<$，相等不更新，這會影響 $\pi$ 記下哪一條路徑（第~\ref{sec:spt-dijkstra}~節的追蹤會遇到）。

$\infty$ 的運算規則約定成：$\infty + x = \infty$（$x$ 是任何有限數），$\infty < \infty$ 不成立。
所以 $d[u] = \infty$ 時鬆弛什麼都不做。 這在 C 程式裡要明寫：用一個很大的數代表 $\infty$
時，$w(u, v)$ 若是負的，「很大的數加上負數」會比「很大的數」小，程式就會從一個根本還沒走
到的 $u$ 去更新 $v$。 本章的程式都先檢查 $d[u]$ 是不是 $\infty$。

鬆弛只會讓 $d[v]$ 變小或不變，不會變大。 本章每一個單一起點演算法都是先呼叫一次
\proc{Initialize-Single-Source}，之後只用鬆弛修改 $d$ 和 $\pi$；它們的差別只在\textbf{鬆
弛哪些邊、照什麼順序、鬆弛幾次}。 Dijkstra 演算法和 DAG 的演算法每條邊各鬆弛一次，
Bellman--Ford 演算法每條邊鬆弛 $n - 1$ 次。

\subsection{鬆弛的性質}

下面的性質只假設「先初始化，之後只做鬆弛」，不管鬆弛的順序。 證明 Dijkstra 和
Bellman--Ford 正確時，只需要這幾條。

\begin{theorem}[鬆弛的性質]
\label{thm:spt-props}
對 $G$ 和起點 $s$ 呼叫 \proc{Initialize-Single-Source}，之後做任意一串 \proc{Relax}。
\begin{enumerate}
\item \textbf{上界性質}：任何時候都有 $d[v] \ge \delta(s, v)$；而且一旦 $d[v] = \delta(s,
v)$，之後就不再改變。
\item \textbf{無路徑性質}：如果 $s$ 到不了 $v$，$d[v]$ 一直是 $\infty$。
\item \textbf{收斂性質}：如果 $s \leadsto u \to v$ 是到 $v$ 的一條最短路徑，而在鬆弛
$(u, v)$ 之前的某個時候已經有 $d[u] = \delta(s, u)$，那麼鬆弛 $(u, v)$ 之後
$d[v] = \delta(s, v)$，以後也一直如此。
\item \textbf{路徑鬆弛性質}：設 $p = \langle v_0, v_1, \ldots, v_k \rangle$ 是從
$s = v_0$ 到 $v_k$ 的最短路徑。 如果邊 $(v_0, v_1), (v_1, v_2), \ldots, (v_{k-1}, v_k)$
依這個先後順序都被鬆弛過（中間可以夾著任何其他的鬆弛，同一條邊也可以多鬆弛幾次），那麼
從此 $d[v_k] = \delta(s, v_k)$。
\end{enumerate}
\end{theorem}

\begin{whybox}[鬆弛的性質]
\textbf{上界性質。} 對鬆弛的次數做歸納。 初始化之後 $d[s] = 0 \ge \delta(s, s)$（$\delta(s,
s)$ 不是 $0$ 就是 $-\infty$，長度 $0$ 的路徑權重是 $0$），其他頂點是 $\infty$，成立。 鬆弛
$(u, v)$ 時若改了 $d[v]$，新的值是
\[
d[u] + w(u, v) \ge \delta(s, u) + w(u, v) \ge \delta(s, v),
\]
第一個 $\ge$ 是歸納假設，第二個是三角不等式（引理~\ref{lem:spt-triangle}）。 只有 $d[v]$
被改過，其他頂點不變，所以成立。 $d[v]$ 達到 $\delta(s, v)$ 以後，鬆弛不會讓它變大，也不
可能讓它變小（會比 $\delta(s, v)$ 還小），所以不再改變。

\textbf{無路徑性質。} $\infty \ge d[v] \ge \delta(s, v) = \infty$。

\textbf{收斂性質。} 由上界性質，$d[u]$ 一旦等於 $\delta(s, u)$ 就不再變，所以鬆弛 $(u, v)$
的時候 $d[u] = \delta(s, u)$。 鬆弛之後
\[
d[v] \le d[u] + w(u, v) = \delta(s, u) + w(u, v) = \delta(s, v).
\]
最後一個等號：$s \leadsto u \to v$ 是最短路徑，由引理~\ref{lem:spt-subpath} 它到 $u$ 的那一
段是到 $u$ 的最短路徑，權重 $\delta(s, u)$，整條的權重就是 $\delta(s, u) + w(u, v)$。 再和
上界性質合起來，$d[v] = \delta(s, v)$，之後不再變。

\textbf{路徑鬆弛性質。} 對 $i$ 歸納：鬆弛完 $(v_{i-1}, v_i)$ 之後 $d[v_i] = \delta(s,
v_i)$。 $i = 0$：$p$ 是最短路徑，$\delta(s, s) = 0$（否則 $s$ 在負環上，到 $v_k$ 就沒有最
短路徑），初始化後 $d[s] = 0$，由上界性質不再改。 $i \ge 1$：由引理~\ref{lem:spt-subpath}，
$\langle v_0, \ldots, v_i \rangle$ 是到 $v_i$ 的最短路徑；鬆弛 $(v_{i-1}, v_i)$ 之前已經
有 $d[v_{i-1}] = \delta(s, v_{i-1})$（歸納假設），由收斂性質成立。
\end{whybox}

路徑鬆弛性質是本章的核心：只要能保證「最短路徑上的邊依序被鬆弛過」，演算法就對。 三個
單一起點演算法各用一種辦法做到這件事。 DAG 的演算法依拓樸順序鬆弛，最短路徑上的邊一定
照順序出現；Bellman--Ford 演算法把所有邊鬆弛 $n - 1$ 輪，第 $i$ 輪一定輪到最短路徑上的第
$i$ 條邊；Dijkstra 演算法的保證最微妙，要用到權重不是負的。

\subsection{最短路徑樹}

演算法結束時，除了 $d$，還要能印出路徑。 和 BFS 樹一樣，沿著 $\pi$ 往回走就是路徑。

\begin{theorem}
\label{thm:spt-tree}
設從 $s$ 走得到的地方沒有負環。 在任意一串鬆弛之後，如果每個頂點都有 $d[v] = \delta(s,
v)$，那麼對每個 $s$ 走得到的 $v \ne s$，$\pi[v] \ne \NIL$；從 $v$ 沿著 $\pi$ 往回走會走到
$s$，走過的路徑倒過來就是一條從 $s$ 到 $v$ 的最短路徑。
\end{theorem}

這些邊 $(\pi[v], v)$ 和 $s$ 走得到的頂點構成一棵以 $s$ 為根的樹，叫\textbf{最短路徑樹}
（shortest-paths tree）。 印出路徑用第~21~章的 \proc{Print-Path}（虛擬碼~\ref{alg:gtr-path}）
就好，不必另寫。 最短路徑可能不只一條，最短路徑樹也就不唯一；演算法記下的是哪一棵，取決
於鬆弛的順序和「相等不更新」的規則。

\begin{whybox}[定理~\ref{thm:spt-tree}]
\textbf{(i) $\pi$ 邊的一個不等式。} 任何時候，只要 $\pi[v] = u \ne \NIL$，就有 $d[v] \ge d[u]
+ w(u, v)$。 設定 $\pi[v] \gets u$ 的那一刻等號成立；之後 $d[u]$ 只會變小，$d[v]$ 要變就
連 $\pi[v]$ 一起改。

\textbf{(ii) $\pi$ 不會形成環。} 假設某一次鬆弛把 $\pi[v_k]$ 設成 $v_{k-1}$，讓
$\langle v_0, v_1, \ldots, v_k = v_0 \rangle$ 第一次成為由 $\pi$ 邊組成的環。 這一刻之前，
這條邊被鬆弛成功，所以 $d[v_k] > d[v_{k-1}] + w(v_{k-1}, v_k)$（這裡用的是更新前的
$d[v_k]$）；環上其他的邊由 (i) 有 $d[v_i] \ge d[v_{i-1}] + w(v_{i-1}, v_i)$。 把這 $k$ 個
不等式加起來，兩邊的 $d$ 剛好各出現一次而消掉，得到 $0 > \sum_{i=1}^{k} w(v_{i-1}, v_i)$：
這是一個負環。 環上的頂點都有 $\pi \ne \NIL$，$d$ 是有限的，都從 $s$ 走得到，和假設矛盾。

\textbf{(iii) 沿著 $\pi$ 走是最短路徑。} 設 $v \ne s$ 走得到，$d[v] = \delta(s, v) < \infty$，
所以 $d[v]$ 被鬆弛改過，$\pi[v] \ne \NIL$。 沒有負環時 $d[s] = 0$ 從沒被改過，$\pi[s] = \NIL$。
從 $v$ 沿著 $\pi$ 往回走，由 (ii) 不會繞圈，只能停在 $\pi = \NIL$ 的頂點，而走過的頂點 $d$
都有限，所以停在 $s$。 對路上每條邊 $(u, x)$（$u = \pi[x]$），由 (i) 和三角不等式，
\[
\delta(s, x) = d[x] \ge d[u] + w(u, x) = \delta(s, u) + w(u, x) \ge \delta(s, x),
\]
所以 $\delta(s, x) = \delta(s, u) + w(u, x)$。 從 $s$ 一路加到 $v$，這條路徑的權重剛好是
$\delta(s, v)$。
\end{whybox}

\section{Dijkstra 演算法}
\label{sec:spt-dijkstra}

\subsection{想法}

先看所有權重都不是負數的情形。 BFS 的佇列有一個性質：頂點出佇列的順序，距離不會變小
（第~\ref{sec:gtr-bfs-proof}~節）。 也就是說，BFS 是\textbf{由近到遠}一個一個確定頂點的距離。 權重
不一樣時，「下一個最近的」不一定是佇列前端的那個，但想法可以照搬：

\begin{itemize}
\item 維護一個集合 $S$，裡面的頂點距離已經確定，$d[v] = \delta(s, v)$。
\item 每一步從 $S$ 外面挑 $d$ 最小的頂點 $u$，宣布它的距離確定了，放進 $S$。
\item 鬆弛 $u$ 的每條出邊，讓 $u$ 的鄰居有機會經過 $u$ 得到更小的 $d$。
\end{itemize}

為什麼 $S$ 外面 $d$ 最小的頂點可以直接確定？ 直覺是這樣：任何一條到 $u$ 的其他走法，
一定在某個地方第一次離開 $S$，走到 $S$ 外的某個頂點 $y$；到 $y$ 的花費已經不小於
$d[u]$，後面再走的邊權重都不是負的，只會更大。 這是第~19~章的貪心法：每一步做眼前
最好的選擇，而且永遠不回頭修改。 權重不是負數這個條件在這個論證裡是必要的，
第~\ref{sec:spt-dijkstra-neg}~節會看到少了它會怎樣。

\subsection{演算法}

$S$ 外面的頂點放在一個最小優先佇列 $Q$ 裡，鍵是 $d$。 用到第~\ref{sec:hp-pq}~節的三個操作：
\proc{Insert}、\proc{Extract-Min}、\proc{Decrease-Key}。 鬆弛讓 $d[v]$ 變小時，用
\proc{Decrease-Key} 通知佇列。

\begin{pseudo}[label=alg:spt-dijkstra]{Dijkstra 演算法}
\Procedure{Dijkstra}{$G, w, s$}
    \State \Call{Initialize-Single-Source}{$G, s$}
    \State $S \gets \emptyset$
    \State $Q \gets$ 空的最小優先佇列
    \For{每個頂點 $u \in V$}
        \State \Call{Insert}{$Q, u, d[u]$}
    \EndFor
    \While{$Q$ 不是空的}
        \State $u \gets$ \Call{Extract-Min}{$Q$} \Comment{$d$ 最小；平手取編號小的}
        \State $S \gets S \cup \{u\}$
        \For{每個 $v \in \mathit{Adj}[u]$} \Comment{依串列的順序}
            \If{$v \notin S$ \AND $d[u] + w(u, v) < d[v]$}
                \State $d[v] \gets d[u] + w(u, v)$
                \State $\pi[v] \gets u$
                \State \Call{Decrease-Key}{$Q, v, d[v]$}
            \EndIf
        \EndFor
    \EndWhile
\EndProcedure
\end{pseudo}

第 2--6 行初始化，所有頂點都放進 $Q$。 第 8 行取出 $d$ 最小的頂點 $u$，第 9 行把它放
進 $S$。 第 11--14 行就是 \proc{Relax}$(u, v, w)$，只多了「$v$ 還不在 $S$ 裡」這個條件，
外加通知佇列。 權重不是負數時，這個多出來的條件其實不會擋掉任何一次真的會成功的鬆弛
（下面的定理證完就知道：$v \in S$ 時 $d[v] = \delta(s, v) \le d[u] + w(u, v)$），寫出來
是為了讓 \proc{Decrease-Key} 只作用在佇列裡的元素上，很多資料結構課本的寫法也是這樣。

\textbf{本書的約定}：第 8 行如果有好幾個頂點的 $d$ 一樣小，取\textbf{編號最小}的；第 10
行照鄰接串列的順序，也就是鄰居編號由小到大。 這兩條規則不影響最後的 $d$，但會影響取出
的順序和 $\pi$，手算時要照著做。

\subsection{手算追蹤}
\label{sec:spt-dijkstra-trace}

圖~\ref{fig:spt-g1}~(a) 是外送的路網 $G_1$：$7$ 個路口、$14$ 條單行道，權重是騎車的分鐘
數，店在路口 $1$。

\begin{center}
\begin{tikzpicture}[x=1cm, y=1cm]
\begin{scope}
\node[spt-v] (a1) at (0, 0) {1};
\node[spt-v] (a2) at (2.4, 1.5) {2};
\node[spt-v] (a3) at (2.4, 0) {3};
\node[spt-v] (a4) at (2.4, -1.5) {4};
\node[spt-v] (a5) at (4.8, 0) {5};
\node[spt-v] (a6) at (4.8, 1.5) {6};
\node[spt-v] (a7) at (7.2, 0) {7};
\draw[spt-a] (a1) -- node[spt-w, above left] {$7$} (a2);
\draw[spt-a] (a1) -- node[spt-w, above] {$2$} (a3);
\draw[spt-a] (a1) -- node[spt-w, below left] {$6$} (a4);
\draw[spt-a] (a3) -- node[spt-w, right] {$4$} (a2);
\draw[spt-a] (a3) -- node[spt-w, right] {$4$} (a4);
\draw[spt-a] (a3) -- node[spt-w, above, pos=0.4] {$8$} (a5);
\draw[spt-a] (a2) -- node[spt-w, pos=0.4] {$2$} (a5);
\draw[spt-a] (a2) -- node[spt-w, below, pos=0.5] {$9$} (a6);
\draw[spt-a] (a6) to[bend right=30] node[spt-w, above] {$1$} (a2);
\draw[spt-a] (a4) -- node[spt-w, pos=0.45] {$4$} (a5);
\draw[spt-a] (a4) to[bend right=18] node[spt-w, below] {$10$} (a7);
\draw[spt-a] (a5) -- node[spt-w, right] {$3$} (a6);
\draw[spt-a] (a5) -- node[spt-w, above] {$6$} (a7);
\draw[spt-a] (a6) -- node[spt-w, above right] {$1$} (a7);
\node[font=\small] at (3.6, -2.8) {(a) $G_1$};
\end{scope}
\begin{scope}[xshift=8.6cm]
\node[spt-v] (b1) at (0, 0) {1};
\node[spt-v] (b2) at (2.4, 1.5) {2};
\node[spt-v] (b3) at (2.4, 0) {3};
\node[spt-v] (b4) at (2.4, -1.5) {4};
\node[spt-v] (b5) at (4.8, 0) {5};
\node[spt-v] (b6) at (4.8, 1.5) {6};
\node[spt-v] (b7) at (6.6, 0) {7};
\draw[spt-tree] (b1) -- (b3);
\draw[spt-tree] (b3) -- (b2);
\draw[spt-tree] (b1) -- (b4);
\draw[spt-tree] (b2) -- (b5);
\draw[spt-tree] (b5) -- (b6);
\draw[spt-tree] (b6) -- (b7);
\node[spt-d, left=1pt of b1] {$0$};
\node[spt-d, above=1pt of b2] {$6$};
\node[spt-d, right=1pt of b3] {$2$};
\node[spt-d, below=1pt of b4] {$6$};
\node[spt-d, below=1pt of b5] {$8$};
\node[spt-d, above=1pt of b6] {$11$};
\node[spt-d, below=1pt of b7] {$12$};
\node[font=\small] at (3.3, -2.8) {(b) 最短路徑樹，藍字是 $\delta(1, v)$};
\end{scope}
\end{tikzpicture}\par\nopagebreak
\diagcaption{外送路網 $G_1$ 和 Dijkstra 演算法從 $1$ 出發得到的最短路徑樹。}
\label{fig:spt-g1}
\end{center}

鄰接串列依鄰居編號排好是
\begin{gather*}
\mathit{Adj}[1] = 2, 3, 4;\quad \mathit{Adj}[2] = 5, 6;\quad \mathit{Adj}[3] = 2, 4, 5;\quad \mathit{Adj}[4] = 5, 7; \\
\mathit{Adj}[5] = 6, 7;\quad \mathit{Adj}[6] = 2, 7;\quad \mathit{Adj}[7] = \text{空}.
\end{gather*}
下表每一列是一輪主迴圈：取出 $u$、鬆弛完它的出邊之後的 $d$。 \sptpick{粗體}是這一輪取出
的頂點，\sptnew{底線}是這一輪變小的值，\sptold{灰色}是之前已經取出、不會再變的值。 到這
一列為止取出過的頂點就是 $S$。

\begin{center}
\begin{tabular}{cc*{7}{c}}
\toprule
輪 & 取出 $u$ & $d[1]$ & $d[2]$ & $d[3]$ & $d[4]$ & $d[5]$ & $d[6]$ & $d[7]$ \\
\midrule
初始 & & $0$ & $\infty$ & $\infty$ & $\infty$ & $\infty$ & $\infty$ & $\infty$ \\
1 & $1$ & \sptpick{0} & \sptnew{7} & \sptnew{2} & \sptnew{6} & $\infty$ & $\infty$ & $\infty$ \\
2 & $3$ & \sptold{0} & \sptnew{6} & \sptpick{2} & $6$ & \sptnew{10} & $\infty$ & $\infty$ \\
3 & $2$ & \sptold{0} & \sptpick{6} & \sptold{2} & $6$ & \sptnew{8} & \sptnew{15} & $\infty$ \\
4 & $4$ & \sptold{0} & \sptold{6} & \sptold{2} & \sptpick{6} & $8$ & $15$ & \sptnew{16} \\
5 & $5$ & \sptold{0} & \sptold{6} & \sptold{2} & \sptold{6} & \sptpick{8} & \sptnew{11} & \sptnew{14} \\
6 & $6$ & \sptold{0} & \sptold{6} & \sptold{2} & \sptold{6} & \sptold{8} & \sptpick{11} & \sptnew{12} \\
7 & $7$ & \sptold{0} & \sptold{6} & \sptold{2} & \sptold{6} & \sptold{8} & \sptold{11} & \sptpick{12} \\
\midrule
$\pi$ & & \NIL & $3$ & $1$ & $1$ & $2$ & $5$ & $6$ \\
\bottomrule
\end{tabular}
\end{center}

幾個要注意的地方：
\begin{itemize}
\item 第 2 輪取出 $3$（$d = 2$），鬆弛 $3 \to 2$ 得到 $2 + 4 = 6 < 7$，$d[2]$ 變小，
$\pi[2] = 3$。 鬆弛 $3 \to 4$ 得到 $2 + 4 = 6$，和原本的 $d[4] = 6$ 一樣，\textbf{相等不
更新}，$\pi[4]$ 仍然是 $1$。 到 $4$ 其實有兩條最短路徑，$1 \to 4$ 和 $1 \to 3 \to 4$，
記下哪一條由這條規則決定。
\item 第 3 輪，$2$ 和 $4$ 的 $d$ 都是 $6$，平手取編號小的 $2$。 取 $4$ 也對，只是後面的列
順序不同。
\item 第 4 輪取出 $4$，$4 \to 5$ 得到 $6 + 4 = 10$，不比 $8$ 小，不更新。 $d[7]$ 先是 $16$，
第 5 輪變成 $14$，第 6 輪再變成 $12$：同一個頂點的 $d$ 可以被改好幾次，每次都是一次
\proc{Decrease-Key}。
\item 第 6 輪取出 $6$，邊 $6 \to 2$ 的 $2$ 已經在 $S$ 裡，第 11 行直接跳過（就算去比，
$11 + 1 = 12$ 也不比 $6$ 小）。
\end{itemize}

最後 $d = (0, 6, 2, 6, 8, 11, 12)$，$\pi$ 構成圖~\ref{fig:spt-g1}~(b) 的最短路徑樹。 到 $7$
的路徑沿著 $\pi$ 往回是 $7, 6, 5, 2, 3, 1$，倒過來是 $1 \to 3 \to 2 \to 5 \to 6 \to 7$，
權重 $2 + 4 + 2 + 3 + 1 = 12$。 圖~\ref{fig:spt-snap} 畫出前六輪結束時的樣子：橘色是這一輪
取出的頂點，藍色是之前取出的，藍色粗箭頭是目前的 $\pi$ 邊。

\begin{center}
\begin{tikzpicture}
\begin{scope}[xshift=0.0cm, yshift=0.0cm, x=0.55cm, y=0.6cm]
\node[spt-vsnew] (p1) at (0, 0) {1};
\node[spt-d, left=0pt of p1] {$0$};
\node[spt-vs] (p2) at (2.4, 1.5) {2};
\node[spt-d, above left=0pt of p2] {$7$};
\node[spt-vs] (p3) at (2.4, 0) {3};
\node[spt-d, below left=0pt of p3] {$2$};
\node[spt-vs] (p4) at (2.4, -1.5) {4};
\node[spt-d, below=0pt of p4] {$6$};
\node[spt-vs] (p5) at (4.8, 0) {5};
\node[spt-d, below right=0pt of p5] {$\infty$};
\node[spt-vs] (p6) at (4.8, 1.5) {6};
\node[spt-d, above right=0pt of p6] {$\infty$};
\node[spt-vs] (p7) at (7.2, 0) {7};
\node[spt-d, right=0pt of p7] {$\infty$};
\draw[spt-ta] (p1) -- (p2);
\draw[spt-ta] (p1) -- (p3);
\draw[spt-ta] (p1) -- (p4);
\draw[spt-sa] (p3) -- (p2);
\draw[spt-sa] (p3) -- (p4);
\draw[spt-sa] (p3) -- (p5);
\draw[spt-sa] (p2) -- (p5);
\draw[spt-sa] (p2) -- (p6);
\draw[spt-sa] (p4) -- (p5);
\draw[spt-sa] (p4) to[bend right=18] (p7);
\draw[spt-sa] (p5) -- (p6);
\draw[spt-sa] (p5) -- (p7);
\draw[spt-sa] (p6) -- (p7);
\draw[spt-sa] (p6) to[bend right=30] (p2);
\node[font=\small] at (3.6, -3.1) {第 1 輪：取出 $1$};
\end{scope}
\begin{scope}[xshift=5.45cm, yshift=0.0cm, x=0.55cm, y=0.6cm]
\node[spt-vsdone] (p1) at (0, 0) {1};
\node[spt-d, left=0pt of p1] {$0$};
\node[spt-vs] (p2) at (2.4, 1.5) {2};
\node[spt-d, above left=0pt of p2] {$6$};
\node[spt-vsnew] (p3) at (2.4, 0) {3};
\node[spt-d, below left=0pt of p3] {$2$};
\node[spt-vs] (p4) at (2.4, -1.5) {4};
\node[spt-d, below=0pt of p4] {$6$};
\node[spt-vs] (p5) at (4.8, 0) {5};
\node[spt-d, below right=0pt of p5] {$10$};
\node[spt-vs] (p6) at (4.8, 1.5) {6};
\node[spt-d, above right=0pt of p6] {$\infty$};
\node[spt-vs] (p7) at (7.2, 0) {7};
\node[spt-d, right=0pt of p7] {$\infty$};
\draw[spt-sa] (p1) -- (p2);
\draw[spt-ta] (p1) -- (p3);
\draw[spt-ta] (p1) -- (p4);
\draw[spt-ta] (p3) -- (p2);
\draw[spt-sa] (p3) -- (p4);
\draw[spt-ta] (p3) -- (p5);
\draw[spt-sa] (p2) -- (p5);
\draw[spt-sa] (p2) -- (p6);
\draw[spt-sa] (p4) -- (p5);
\draw[spt-sa] (p4) to[bend right=18] (p7);
\draw[spt-sa] (p5) -- (p6);
\draw[spt-sa] (p5) -- (p7);
\draw[spt-sa] (p6) -- (p7);
\draw[spt-sa] (p6) to[bend right=30] (p2);
\node[font=\small] at (3.6, -3.1) {第 2 輪：取出 $3$};
\end{scope}
\begin{scope}[xshift=10.9cm, yshift=0.0cm, x=0.55cm, y=0.6cm]
\node[spt-vsdone] (p1) at (0, 0) {1};
\node[spt-d, left=0pt of p1] {$0$};
\node[spt-vsnew] (p2) at (2.4, 1.5) {2};
\node[spt-d, above left=0pt of p2] {$6$};
\node[spt-vsdone] (p3) at (2.4, 0) {3};
\node[spt-d, below left=0pt of p3] {$2$};
\node[spt-vs] (p4) at (2.4, -1.5) {4};
\node[spt-d, below=0pt of p4] {$6$};
\node[spt-vs] (p5) at (4.8, 0) {5};
\node[spt-d, below right=0pt of p5] {$8$};
\node[spt-vs] (p6) at (4.8, 1.5) {6};
\node[spt-d, above right=0pt of p6] {$15$};
\node[spt-vs] (p7) at (7.2, 0) {7};
\node[spt-d, right=0pt of p7] {$\infty$};
\draw[spt-sa] (p1) -- (p2);
\draw[spt-ta] (p1) -- (p3);
\draw[spt-ta] (p1) -- (p4);
\draw[spt-ta] (p3) -- (p2);
\draw[spt-sa] (p3) -- (p4);
\draw[spt-sa] (p3) -- (p5);
\draw[spt-ta] (p2) -- (p5);
\draw[spt-ta] (p2) -- (p6);
\draw[spt-sa] (p4) -- (p5);
\draw[spt-sa] (p4) to[bend right=18] (p7);
\draw[spt-sa] (p5) -- (p6);
\draw[spt-sa] (p5) -- (p7);
\draw[spt-sa] (p6) -- (p7);
\draw[spt-sa] (p6) to[bend right=30] (p2);
\node[font=\small] at (3.6, -3.1) {第 3 輪：取出 $2$};
\end{scope}
\begin{scope}[xshift=0.0cm, yshift=-3.7cm, x=0.55cm, y=0.6cm]
\node[spt-vsdone] (p1) at (0, 0) {1};
\node[spt-d, left=0pt of p1] {$0$};
\node[spt-vsdone] (p2) at (2.4, 1.5) {2};
\node[spt-d, above left=0pt of p2] {$6$};
\node[spt-vsdone] (p3) at (2.4, 0) {3};
\node[spt-d, below left=0pt of p3] {$2$};
\node[spt-vsnew] (p4) at (2.4, -1.5) {4};
\node[spt-d, below=0pt of p4] {$6$};
\node[spt-vs] (p5) at (4.8, 0) {5};
\node[spt-d, below right=0pt of p5] {$8$};
\node[spt-vs] (p6) at (4.8, 1.5) {6};
\node[spt-d, above right=0pt of p6] {$15$};
\node[spt-vs] (p7) at (7.2, 0) {7};
\node[spt-d, right=0pt of p7] {$16$};
\draw[spt-sa] (p1) -- (p2);
\draw[spt-ta] (p1) -- (p3);
\draw[spt-ta] (p1) -- (p4);
\draw[spt-ta] (p3) -- (p2);
\draw[spt-sa] (p3) -- (p4);
\draw[spt-sa] (p3) -- (p5);
\draw[spt-ta] (p2) -- (p5);
\draw[spt-ta] (p2) -- (p6);
\draw[spt-sa] (p4) -- (p5);
\draw[spt-ta] (p4) to[bend right=18] (p7);
\draw[spt-sa] (p5) -- (p6);
\draw[spt-sa] (p5) -- (p7);
\draw[spt-sa] (p6) -- (p7);
\draw[spt-sa] (p6) to[bend right=30] (p2);
\node[font=\small] at (3.6, -3.1) {第 4 輪：取出 $4$};
\end{scope}
\begin{scope}[xshift=5.45cm, yshift=-3.7cm, x=0.55cm, y=0.6cm]
\node[spt-vsdone] (p1) at (0, 0) {1};
\node[spt-d, left=0pt of p1] {$0$};
\node[spt-vsdone] (p2) at (2.4, 1.5) {2};
\node[spt-d, above left=0pt of p2] {$6$};
\node[spt-vsdone] (p3) at (2.4, 0) {3};
\node[spt-d, below left=0pt of p3] {$2$};
\node[spt-vsdone] (p4) at (2.4, -1.5) {4};
\node[spt-d, below=0pt of p4] {$6$};
\node[spt-vsnew] (p5) at (4.8, 0) {5};
\node[spt-d, below right=0pt of p5] {$8$};
\node[spt-vs] (p6) at (4.8, 1.5) {6};
\node[spt-d, above right=0pt of p6] {$11$};
\node[spt-vs] (p7) at (7.2, 0) {7};
\node[spt-d, right=0pt of p7] {$14$};
\draw[spt-sa] (p1) -- (p2);
\draw[spt-ta] (p1) -- (p3);
\draw[spt-ta] (p1) -- (p4);
\draw[spt-ta] (p3) -- (p2);
\draw[spt-sa] (p3) -- (p4);
\draw[spt-sa] (p3) -- (p5);
\draw[spt-ta] (p2) -- (p5);
\draw[spt-sa] (p2) -- (p6);
\draw[spt-sa] (p4) -- (p5);
\draw[spt-sa] (p4) to[bend right=18] (p7);
\draw[spt-ta] (p5) -- (p6);
\draw[spt-ta] (p5) -- (p7);
\draw[spt-sa] (p6) -- (p7);
\draw[spt-sa] (p6) to[bend right=30] (p2);
\node[font=\small] at (3.6, -3.1) {第 5 輪：取出 $5$};
\end{scope}
\begin{scope}[xshift=10.9cm, yshift=-3.7cm, x=0.55cm, y=0.6cm]
\node[spt-vsdone] (p1) at (0, 0) {1};
\node[spt-d, left=0pt of p1] {$0$};
\node[spt-vsdone] (p2) at (2.4, 1.5) {2};
\node[spt-d, above left=0pt of p2] {$6$};
\node[spt-vsdone] (p3) at (2.4, 0) {3};
\node[spt-d, below left=0pt of p3] {$2$};
\node[spt-vsdone] (p4) at (2.4, -1.5) {4};
\node[spt-d, below=0pt of p4] {$6$};
\node[spt-vsdone] (p5) at (4.8, 0) {5};
\node[spt-d, below right=0pt of p5] {$8$};
\node[spt-vsnew] (p6) at (4.8, 1.5) {6};
\node[spt-d, above right=0pt of p6] {$11$};
\node[spt-vs] (p7) at (7.2, 0) {7};
\node[spt-d, right=0pt of p7] {$12$};
\draw[spt-sa] (p1) -- (p2);
\draw[spt-ta] (p1) -- (p3);
\draw[spt-ta] (p1) -- (p4);
\draw[spt-ta] (p3) -- (p2);
\draw[spt-sa] (p3) -- (p4);
\draw[spt-sa] (p3) -- (p5);
\draw[spt-ta] (p2) -- (p5);
\draw[spt-sa] (p2) -- (p6);
\draw[spt-sa] (p4) -- (p5);
\draw[spt-sa] (p4) to[bend right=18] (p7);
\draw[spt-ta] (p5) -- (p6);
\draw[spt-sa] (p5) -- (p7);
\draw[spt-ta] (p6) -- (p7);
\draw[spt-sa] (p6) to[bend right=30] (p2);
\node[font=\small] at (3.6, -3.1) {第 6 輪：取出 $6$};
\end{scope}
\end{tikzpicture}\par\nopagebreak
\diagcaption{Dijkstra 演算法在 $G_1$ 上前六輪結束時的 $d$ 和 $\pi$。 第 7 輪取出 $7$，
不再改變任何值。}
\label{fig:spt-snap}
\end{center}

\begin{examnote}
Dijkstra 的追蹤題，答案會被三個約定影響：(1) 好幾個頂點的 $d$ 一樣小時取哪一個；(2) 鬆弛
時「相等」要不要更新 $\pi$；(3) 表格的格式。 本書取編號最小的、相等不更新，表格每一列寫
取出的頂點和全部的 $d$。 有的課本每一列只寫集合 $S$ 和距離陣列，有的每格寫成
「$d/\pi$」，也有的題目只問「第幾個被確定的頂點是誰」。 題目有規定就照題目；沒有規定
時，在答案開頭寫一句你用的平手規則，表格裡每一輪都把取出的頂點標出來。 無向圖的題目，
每條邊兩個方向都要鬆弛。
\end{examnote}

\subsection{為什麼一定對}

\begin{theorem}
\label{thm:spt-dijkstra}
設 $G$ 每條邊的權重都 $\ge 0$。 對 $G$ 和起點 $s$ 執行 \proc{Dijkstra}，每個頂點 $u$ 在第
8 行被取出時都有 $d[u] = \delta(s, u)$；所以演算法結束時，每個頂點都有 $d[v] = \delta(s,
v)$，而 $\pi$ 構成一棵最短路徑樹。
\end{theorem}

\begin{proof}
權重不是負數，沒有負環，所以 $\delta(s, s) = 0$。 第 11--13 行只是在 $v \notin S$ 時做
\proc{Relax}，少做幾次鬆弛不影響定理~\ref{thm:spt-props}，那幾條性質都可以用。

假設定理不對，設 $u$ 是\textbf{第一個}被取出時 $d[u] \ne \delta(s, u)$ 的頂點。 第一個取出
的是 $s$，$d[s] = 0 = \delta(s, s)$，所以 $u \ne s$。 由上界性質 $d[u] > \delta(s, u)$，所以
$\delta(s, u) < \infty$，$s$ 走得到 $u$；取一條從 $s$ 到 $u$ 的最短路徑 $p$。 看 $u$ 被取出
的那一刻（第 9 行之前）：$s$ 在 $S$ 裡，$u$ 不在。 沿著 $p$ 從 $s$ 往前走，設 $y$ 是第一個
不在 $S$ 裡的頂點，$x$ 是 $p$ 上 $y$ 的前一個頂點，$x \in S$。 （$y$ 可能就是 $u$。）

$x$ 比 $u$ 早取出，依 $u$ 的選法，$x$ 取出時 $d[x] = \delta(s, x)$。 取出 $x$ 之後第 10
行看過邊 $(x, y)$，那時 $y$ 不在 $S$ 裡，所以這條邊被鬆弛過。 由引理~\ref{lem:spt-subpath}，
$p$ 到 $y$ 的那一段 $s \leadsto x \to y$ 是到 $y$ 的最短路徑，由收斂性質
$d[y] = \delta(s, y)$。

$y$ 在 $p$ 上排在 $u$ 前面（或者就是 $u$），從 $y$ 走到 $u$ 的那一段權重 $\ge 0$，所以
\[
d[y] = \delta(s, y) \le \delta(s, u) < d[u].
\]
可是 $y$ 和 $u$ 都在 $Q$ 裡，第 8 行選了 $u$，表示 $d[u] \le d[y]$，矛盾。 （$y = u$ 時，上式
直接寫成 $d[u] = \delta(s, u)$，一樣矛盾。）

所以每個頂點取出時 $d = \delta$，由上界性質之後不再改變；$s$ 到不了的頂點由無路徑性質一
直是 $\infty$。 最後所有 $d[v] = \delta(s, v)$，由定理~\ref{thm:spt-tree}，$\pi$ 構成最短路
徑樹。
\end{proof}

證明裡唯一用到「權重 $\ge 0$」的地方是 $\delta(s, y) \le \delta(s, u)$：最短路徑上後面的
頂點距離不會比前面的小。 順便得到一個常用的事實：權重 $\ge 0$ 時，Dijkstra 演算法取出頂點
的順序，$d$ 值不會變小（觀念題第~\ref{spt-q-mono}~題），和 BFS 的佇列一樣是由近到遠。

\subsection{有負邊時會錯}
\label{sec:spt-dijkstra-neg}

圖~\ref{fig:spt-neg} 的 $N_1$ 只有一條負邊 $3 \to 2$，沒有負環。 真正的最短距離是
$\delta(1, 2) = 5 - 4 = 1$（$1 \to 3 \to 2$）、$\delta(1, 4) = 1 + 2 = 3$（$1 \to 3 \to 2
\to 4$）。

\begin{center}
\begin{tikzpicture}[x=1cm, y=1cm]
\node[spt-v] (m1) at (0, 0) {1};
\node[spt-v] (m2) at (2.4, 1) {2};
\node[spt-v] (m3) at (2.4, -1) {3};
\node[spt-v] (m4) at (4.8, 1) {4};
\draw[spt-a] (m1) -- node[spt-w, above left] {$3$} (m2);
\draw[spt-a] (m1) -- node[spt-w, below left] {$5$} (m3);
\draw[spt-neg] (m3) -- node[spt-wn, right] {$-4$} (m2);
\draw[spt-a] (m2) -- node[spt-w, above] {$2$} (m4);
\end{tikzpicture}\par\nopagebreak
\diagcaption{$N_1$：一條負邊就足以讓 Dijkstra 演算法出錯。}
\label{fig:spt-neg}
\end{center}

照虛擬碼~\ref{alg:spt-dijkstra} 做：

\begin{center}
\begin{tabular}{cc*{4}{c}l}
\toprule
輪 & 取出 $u$ & $d[1]$ & $d[2]$ & $d[3]$ & $d[4]$ & 說明 \\
\midrule
初始 & & $0$ & $\infty$ & $\infty$ & $\infty$ & \\
1 & $1$ & \sptpick{0} & \sptnew{3} & \sptnew{5} & $\infty$ & \\
2 & $2$ & \sptold{0} & \sptpick{3} & $5$ & \sptnew{5} & $2$ 被宣布確定，$d[2] = 3$ \\
3 & $3$ & \sptold{0} & \sptold{3} & \sptpick{5} & $5$ & $3 \to 2$：$2 \in S$，跳過 \\
4 & $4$ & \sptold{0} & \sptold{3} & \sptold{5} & \sptpick{5} & \\
\bottomrule
\end{tabular}
\end{center}

第 3 輪 $d[3] = d[4] = 5$ 平手，取編號小的 $3$。 結果 $d[2] = 3$、$d[4] = 5$，都錯了。 問
題出在第 2 輪：$2$ 的 $d = 3$ 是當時 $Q$ 裡最小的，演算法就宣布它確定了；證明裡的
「$\delta(s, y) \le \delta(s, u)$」在這裡不成立，因為經過 $3$ 再走一條負邊，可以比 $3$ 本身
還近。 就算第 3 輪不管 $2 \in S$、硬把 $d[2]$ 改成 $1$，$2$ 的出邊已經在第 2 輪處理過，不
會再處理一次，$d[4]$ 還是錯的 $5$。 這正是第~\ref{sec:gr-fail}~節說的「貪心何時不對」：一旦選了
就不回頭。

\begin{itemize}
\item 有負邊不代表一定算錯，例如把 $N_1$ 的 $1 \to 2$ 改成 $0$，結果剛好全對。 重點是\textbf{不
保證}對，所以有負邊就不能用。
\item 有些程式不維護 $S$，只要某個頂點的 $d$ 變小就再把它放進佇列，連已經取出過的也一樣。
這樣在有負邊、沒有負環的圖上答案會對，但最壞情況下同一個頂點會被取出非常多次，時間
可以是指數級，已經不是 Dijkstra 演算法了。 有負邊時應該用下一節的 Bellman--Ford。
\item 「把每條邊都加上同一個數，讓權重全變成非負」也不行：邊數不同的路徑加的量不同，最
短路徑會跟著變（觀念題第~\ref{spt-q-const}~題）。 第~\ref{sec:spt-johnson}~節有正確的改
權重方法。
\end{itemize}

\subsection{執行時間}

每個頂點進出 $Q$ 各一次，每條邊在第 10 行被看一次，所以 \proc{Insert} 和
\proc{Extract-Min} 各做 $n$ 次，\proc{Decrease-Key} 最多做 $m$ 次。 總時間取決於優先佇列
怎麼做：

\begin{center}
\begin{tabular}{lccc}
\toprule
$Q$ 的實作 & \proc{Extract-Min} & \proc{Decrease-Key} & 總時間 \\
\midrule
陣列（以頂點編號為索引） & $\Theta(n)$ & $\Theta(1)$ & $\Theta(n^2 + m) = \Theta(n^2)$ \\
二元堆積（第~\ref{sec:hp-pq}~節） & $O(\lg n)$ & $O(\lg n)$ & $O((n + m) \lg n)$ \\
$d$ 元堆積（第~\ref{sec:hp-dary}~節） & $O(d \log_d n)$ & $O(\log_d n)$ & $O((nd + m) \log_d n)$ \\
費氏堆積（第~\ref{sec:hp-fib}~節，平攤） & $O(\lg n)$ & $O(1)$ & $O(n \lg n + m)$ \\
\bottomrule
\end{tabular}
\end{center}

陣列版是最直接的寫法：$d$ 本身就是陣列，再用一個布林陣列標記誰在 $S$ 裡；
\proc{Extract-Min} 掃過所有不在 $S$ 裡的頂點找最小的，\proc{Decrease-Key} 就是直接改
$d[v]$。 圖用鄰接矩陣存時，第 10 行掃一整列也是 $\Theta(n)$，總時間一樣是 $\Theta(n^2)$。

哪一種比較好看圖的疏密。 稠密圖 $m = \Theta(n^2)$：陣列版 $\Theta(n^2)$，二元堆積是
$\Theta(n^2 \lg n)$，陣列版反而快。 稀疏圖 $m = O(n)$：陣列版還是 $\Theta(n^2)$，二元堆積
是 $O(n \lg n)$，快很多。 實際的道路圖、網路圖大多稀疏，所以程式裡最常見的是堆積版。
$d$ 元堆積取 $d \approx \max(2, m/n)$ 時，總時間是 $O(m \log_d n)$，邊越多的圖比二元堆積
快得越多。 費氏堆積兩邊都不輸，但常數大，主要用在理論上的時間界（第~\ref{sec:hp-fib}~節）。

\subsection{C 程式}

下面的程式用二元堆積，但不做 \proc{Decrease-Key}：$d[v]$ 變小時，直接把新的一筆
$(d[v], v)$ 放進堆積，舊的那筆留在裡面；取出時發現頂點已經處理過，就丟掉。 這是第~\ref{sec:hp-pq}~節
提到的寫法，不必記錄每個頂點在堆積的哪一格。 每次成功的鬆弛放進一筆，堆積最
多 $m + 1$ 筆，時間 $O(m \lg m) = O(m \lg n)$（因為 $m \le n^2$，$\lg m \le 2 \lg n$）。

堆積比較的是 $(d, v)$ 這一對：先比 $d$，平手比頂點編號。 這樣取出的順序和虛擬碼的平手規
則完全一樣。 理由：頂點 $v$ 還沒處理時，堆積裡一定有 $(d[v], v)$ 這一筆（最後一次變小時放
進去的），舊的幾筆 $(k, v)$ 都有 $k > d[v]$。 所以堆積裡最小的一筆如果屬於還沒處理的頂點，
它就是所有還沒處理的頂點裡 $(d, \text{編號})$ 最小的；過期的筆數不影響誰先出來。

\begin{cbookrecap}[第~13~章 13.1 節、第~19~章 19.10 節]
距離用 \texttt{long long}，邊數多、權重大時總和才不會溢位。 $\infty$ 用
\texttt{LLONG\_MAX / 4}：比任何真實的距離都大，再加上一條邊的權重也不會超出
\texttt{long long} 的範圍。 \texttt{Node **p = \&head[u];} 讓 \texttt{p} 指向「要改的那個
指標」，插入新節點時串列頭和中間不必分開寫，和第~22~章的程式一樣。
\end{cbookrecap}

輸入有好幾張圖，每張圖第一行是 $n$、$m$、起點 $s$ 和要印出路徑的終點 $t$，接著是 $m$ 條邊
$u$、$v$、$w$。 程式印出取出的順序（括號是當時的 $d$）、最後的 $d$ 和 $\pi$，以及到 $t$ 的
路徑。

\begin{example}[label=ex:spt-dijkstra]{c}{(dijkstra.c) Dijkstra 演算法：二元堆積、不做 decrease-key}
#include <stdio.h>
#include <stdlib.h>
#include <limits.h>

#define INF (LLONG_MAX / 4)     // 代表 ∞，加上一條邊也不會溢位

typedef struct node {           // 邊 u -> to，權重 w
    int to;
    long long w;
    struct node *next;
} Node;

typedef struct {                // 堆積的一筆：頂點 v，放進去時的 d 值
    long long key;
    int v;
} Item;

// 加入邊 u -> v，插在串列裡讓終點保持由小到大
int add_edge(Node **head, int u, int v, long long w)
{
    Node *x = malloc(sizeof *x);
    if (x == NULL) {
        return 0;
    }
    Node **p = &head[u];
    while (*p != NULL && (*p)->to < v) {
        p = &(*p)->next;
    }
    x->to = v;
    x->w = w;
    x->next = *p;
    *p = x;
    return 1;
}

int less(Item x, Item y)        // 先比 key，平手比頂點編號
{
    return x.key < y.key || (x.key == y.key && x.v < y.v);
}

void push(Item *a, int *size, long long key, int v)
{
    int i = (*size)++;
    Item x = { key, v };
    while (i > 0 && less(x, a[(i - 1) / 2])) {          // 上浮
        a[i] = a[(i - 1) / 2];
        i = (i - 1) / 2;
    }
    a[i] = x;
}

Item pop(Item *a, int *size)
{
    Item top = a[0], x = a[--*size];
    int i = 0, c;
    while ((c = 2 * i + 1) < *size) {                   // 下沉
        if (c + 1 < *size && less(a[c + 1], a[c])) {
            c++;
        }
        if (!less(a[c], x)) {
            break;
        }
        a[i] = a[c];
        i = c;
    }
    a[i] = x;
    return top;
}

// Dijkstra：頂點 0..n-1，起點 s；順便印出取出的順序
void dijkstra(Node **head, int n, int m, int s, long long *d, int *pi)
{
    char *done = calloc(n, 1);
    Item *heap = malloc((m + 1) * sizeof *heap);
    int size = 0;
    if (done != NULL && heap != NULL) {
        for (int v = 0; v < n; v++) {
            d[v] = INF;
            pi[v] = -1;
        }
        d[s] = 0;
        push(heap, &size, 0, s);
        printf("order:");
        while (size > 0) {
            int u = pop(heap, &size).v;
            if (done[u]) {                      // 過期的一筆
                continue;
            }
            done[u] = 1;
            printf(" %d(%lld)", u + 1, d[u]);
            for (Node *x = head[u]; x != NULL; x = x->next) {
                if (!done[x->to] && d[u] + x->w < d[x->to]) {
                    d[x->to] = d[u] + x->w;
                    pi[x->to] = u;
                    push(heap, &size, d[x->to], x->to);
                }
            }
        }
        printf("\n");
    }
    free(done);
    free(heap);
}

void print_path(const int *pi, int v)
{
    if (pi[v] != -1) {
        print_path(pi, pi[v]);
        printf(" ");
    }
    printf("%d", v + 1);
}

int main(void)
{
    int n, m, s, t;
    while (scanf("%d %d %d %d", &n, &m, &s, &t) == 4) {
        Node **head = calloc(n, sizeof *head);
        long long *d = malloc(n * sizeof *d);
        int *pi = malloc(n * sizeof *pi), u, v;
        long long w;
        if (head == NULL || d == NULL || pi == NULL) {
            return 1;
        }
        for (int i = 0; i < m; i++) {
            if (scanf("%d %d %lld", &u, &v, &w) != 3
                || !add_edge(head, u - 1, v - 1, w)) {
                return 1;
            }
        }
        dijkstra(head, n, m, s - 1, d, pi);
        printf("d:");
        for (v = 0; v < n; v++) {
            if (d[v] == INF) {
                printf(" inf");
            } else {
                printf(" %lld", d[v]);
            }
        }
        printf("\npi:");
        for (v = 0; v < n; v++) {
            if (pi[v] == -1) {
                printf(" -");
            } else {
                printf(" %d", pi[v] + 1);
            }
        }
        if (d[t - 1] == INF) {
            printf("\n%d cannot reach %d\n", s, t);
        } else {
            printf("\npath to %d: ", t);
            print_path(pi, t - 1);
            printf(" (length %lld)\n", d[t - 1]);
        }
        for (u = 0; u < n; u++) {
            while (head[u] != NULL) {
                Node *nx = head[u]->next;
                free(head[u]);
                head[u] = nx;
            }
        }
        free(head);
        free(d);
        free(pi);
    }
    return 0;
}
\end{example}
\iosample{7 14 1 7 \\ 1 2 7 1 3 2 1 4 6 3 2 4 3 4 4 \\ 3 5 8 2 5 2 2 6 9 4 5 4 4 7 10 \\ 5 6 3 5 7 6 6 7 1 6 2 1 \\ 4 4 1 4 \\ 1 2 3 1 3 5 3 2 -4 2 4 2 \\ 3 1 1 3 \\ 1 2 4}{order: 1(0) 3(2) 2(6) 4(6) 5(8) 6(11) 7(12) \\ d: 0 6 2 6 8 11 12 \\ pi: - 3 1 1 2 5 6 \\ path to 7: 1 3 2 5 6 7 (length 12) \\ order: 1(0) 2(3) 3(5) 4(5) \\ d: 0 3 5 5 \\ pi: - 1 1 2 \\ path to 4: 1 2 4 (length 5) \\ order: 1(0) 2(4) \\ d: 0 4 inf \\ pi: - 1 - \\ 1 cannot reach 3}

第一張圖是 $G_1$，取出順序 $1, 3, 2, 4, 5, 6, 7$ 和每個頂點的 $d$、$\pi$ 都和手算的表相同。
第二張圖是 $N_1$，程式忠實地重現了錯誤的答案 $d[2] = 3$、$d[4] = 5$。 第三張圖頂點 $3$ 沒
有任何邊連進來，$d[3]$ 保持 $\infty$，$\pi[3]$ 保持 \NIL。 印路徑的 \texttt{print\_path} 和
\proc{Print-Path} 一樣用遞迴，先印前面的部分再印自己。

\section{Bellman--Ford 演算法與負環}
\label{sec:spt-bf}

\subsection{演算法}

有負邊時，沒辦法像 Dijkstra 那樣一個一個「確定」頂點。 Bellman--Ford 演算法換一個保守
的做法：每一輪把\textbf{所有的邊}都鬆弛一次，做 $n - 1$ 輪。 最後再掃一次所有的邊，如果還
有邊能鬆弛，就表示從 $s$ 走得到一個負環。

\begin{pseudo}[label=alg:spt-bf]{Bellman--Ford 演算法}
\Function{Bellman-Ford}{$G, w, s$}
    \State \Call{Initialize-Single-Source}{$G, s$}
    \For{$i \gets 1$ \To $n - 1$}
        \For{每條邊 $(u, v) \in E$} \Comment{依固定的邊順序}
            \State \Call{Relax}{$u, v, w$}
        \EndFor
    \EndFor
    \For{每條邊 $(u, v) \in E$}
        \If{$d[u] + w(u, v) < d[v]$}
            \State \Return \FALSE \Comment{從 $s$ 走得到負環}
        \EndIf
    \EndFor
    \State \Return \TRUE
\EndFunction
\end{pseudo}

第 3--5 行是 $n - 1$ 輪鬆弛，第 6--8 行是檢查。 邊用什麼順序掃，最後的答案都一樣（下面的
定理不管順序），但\textbf{每一輪結束時的 $d$ 會不一樣}，所以手算時一定要照指定的順序。
\textbf{本書的約定}：每一輪都照同一個順序掃邊，題目或輸入給了邊的順序就照它；沒有給時，
照起點編號由小到大、同一個起點再照終點由小到大，也就是依序掃 $\mathit{Adj}[1]$、
$\mathit{Adj}[2]$、\ldots。

\subsection{手算追蹤}
\label{sec:spt-bf-trace}

圖~\ref{fig:spt-bf} 的 $B$ 有兩條負邊 $3 \to 2$、$4 \to 5$，沒有負環（唯一的環
$3 \to 2 \to 4 \to 5 \to 3$ 權重 $-3 + 2 - 1 + 3 = 1$）。 起點是 $1$，邊依約定的順序是
$(1, 2)$、$(1, 3)$、$(2, 4)$、$(2, 5)$、$(3, 2)$、$(3, 4)$、$(4, 5)$、$(5, 3)$。

\begin{center}
\begin{tikzpicture}[x=1cm, y=0.8cm]
\node[spt-v] (c1) at (0, 0) {1};
\node[spt-v] (c2) at (2.4, 1.2) {2};
\node[spt-v] (c3) at (2.4, -1.2) {3};
\node[spt-v] (c4) at (4.8, -1.2) {4};
\node[spt-v] (c5) at (7.2, 0) {5};
\draw[spt-a] (c1) -- node[spt-w, above left] {$5$} (c2);
\draw[spt-a] (c1) -- node[spt-w, below left] {$3$} (c3);
\draw[spt-neg] (c3) -- node[spt-wn, left] {$-3$} (c2);
\draw[spt-a] (c2) -- node[spt-w, pos=0.4, above right] {$2$} (c4);
\draw[spt-a] (c2) -- node[spt-w, above] {$6$} (c5);
\draw[spt-a] (c3) -- node[spt-w, above] {$4$} (c4);
\draw[spt-neg] (c4) -- node[spt-wn, below right] {$-1$} (c5);
\begin{pgfinterruptboundingbox}
\draw[spt-a] (c5) to[out=-100, in=-60, looseness=1.0] node[spt-w, below] {$3$} (c3);
\end{pgfinterruptboundingbox}
\path (0, -2.8);   % 只把弧線真正畫到的地方算進圖的範圍
\end{tikzpicture}\par\nopagebreak
\diagcaption{圖 $B$。 紅色是負邊。 把 $5 \to 3$ 的權重從 $3$ 改成 $1$ 就得到第~\ref{sec:spt-negcycle}~節的 $B'$。}
\label{fig:spt-bf}
\end{center}

先把第 1 輪一條邊一條邊寫出來：

\begin{center}
\small
\begin{tabular}{clcl}
\toprule
邊 & 計算 & 成功？ & 之後的 $d[1..5]$ \\
\midrule
$(1, 2)$ & $0 + 5 = 5 < \infty$ & 是 & $0, 5, \infty, \infty, \infty$ \\
$(1, 3)$ & $0 + 3 = 3 < \infty$ & 是 & $0, 5, 3, \infty, \infty$ \\
$(2, 4)$ & $5 + 2 = 7 < \infty$ & 是 & $0, 5, 3, 7, \infty$ \\
$(2, 5)$ & $5 + 6 = 11 < \infty$ & 是 & $0, 5, 3, 7, 11$ \\
$(3, 2)$ & $3 - 3 = 0 < 5$ & 是 & $0, 0, 3, 7, 11$ \\
$(3, 4)$ & $3 + 4 = 7 \not< 7$ & 否 & \\
$(4, 5)$ & $7 - 1 = 6 < 11$ & 是 & $0, 0, 3, 7, 6$ \\
$(5, 3)$ & $6 + 3 = 9 \not< 3$ & 否 & \\
\bottomrule
\end{tabular}
\end{center}

注意同一輪裡，前面剛改過的 $d$ 馬上就被後面的邊用到：$(2, 4)$ 用的是這一輪剛算出的
$d[2] = 5$，$(4, 5)$ 又用了剛算出的 $d[4] = 7$。 第 1 輪結束時 $d[5] = 6$，對應的是
$1 \to 2 \to 4 \to 5$，已經用了三條邊。 之後每一輪只寫結果：

\begin{center}
\begin{tabular}{c*{5}{c}l}
\toprule
輪 & $d[1]$ & $d[2]$ & $d[3]$ & $d[4]$ & $d[5]$ & 這一輪成功鬆弛的邊 \\
\midrule
初始 & $0$ & $\infty$ & $\infty$ & $\infty$ & $\infty$ & \\
1 & $0$ & $0$ & $3$ & $7$ & $6$ & $(1,2)$ $(1,3)$ $(2,4)$ $(2,5)$ $(3,2)$ $(4,5)$ \\
2 & $0$ & $0$ & $3$ & \sptnew{2} & \sptnew{1} & $(2,4)$ $(4,5)$ \\
3 & $0$ & $0$ & $3$ & $2$ & $1$ & 沒有 \\
4 & $0$ & $0$ & $3$ & $2$ & $1$ & 沒有 \\
\midrule
$\pi$ & \NIL & $3$ & $1$ & $2$ & $4$ & 檢查：沒有邊能鬆弛，回傳 \TRUE \\
\bottomrule
\end{tabular}
\end{center}

第 2 輪 $(2, 4)$ 用到第 1 輪掃過 $(2, 4)$ 之後才改好的 $d[2] = 0$，得到 $d[4] = 2$；同一輪的 $(4, 5)$ 再
得到 $d[5] = 1$。 最短路徑樹是 $1 \to 3 \to 2 \to 4 \to 5$ 這一條鏈，到 $5$ 的最短路徑用了
四條邊、兩條負邊，權重 $3 - 3 + 2 - 1 = 1$。

\textbf{換一個順序。} 如果邊的順序剛好倒過來，$(5,3), (4,5), (3,4), (3,2), (2,5), (2,4),
(1,3), (1,2)$，每一輪的結果是：

\begin{center}
\begin{tabular}{c*{5}{c}}
\toprule
輪 & $d[1]$ & $d[2]$ & $d[3]$ & $d[4]$ & $d[5]$ \\
\midrule
1 & $0$ & $5$ & $3$ & $\infty$ & $\infty$ \\
2 & $0$ & $0$ & $3$ & $2$ & $6$ \\
3 & $0$ & $0$ & $3$ & $2$ & $1$ \\
4 & $0$ & $0$ & $3$ & $2$ & $1$ \\
\bottomrule
\end{tabular}
\end{center}

第 1 輪只有從 $1$ 出發的兩條邊有用，因為其他邊在 $d[2]$、$d[3]$ 還是 $\infty$ 的時候就掃過
了；要到第 3 輪才收斂。 最後的 $d$ 和 $\pi$ 一樣，過程不一樣。

\begin{examnote}
Bellman--Ford 的追蹤題，先看題目要的是哪一種表。 (1) 本書（以及 CLRS）的寫法是\textbf{就
地更新}：每一輪照固定的邊順序鬆弛，同一輪裡改過的 $d$ 立刻被後面的邊用到，表格寫每一輪
結束時的 $d$。 (2) 有的資料結構課本寫成 $\mathit{dist}^{k}[v]$，意思是「最多用 $k$ 條邊」的
最短距離，第 $k$ 欄只用第 $k - 1$ 欄的值算，同一輪不互相影響。 圖 $B$ 照 (2) 算出來是
$\mathit{dist}^1 = (0, 5, 3, \infty, \infty)$、$\mathit{dist}^2 = (0, 0, 3, 7, 11)$、
$\mathit{dist}^3 = (0, 0, 3, 2, 6)$、$\mathit{dist}^4 = (0, 0, 3, 2, 1)$，和上面的表不同，
但最後一欄相同。 題目給了邊的順序或指定了寫法就照題目；都沒有時，寫明你用哪一種、邊照
什麼順序。
\end{examnote}

\subsection{為什麼一定對}

\begin{theorem}
\label{thm:spt-bf}
對 $G$ 和起點 $s$ 執行 \proc{Bellman-Ford}。
\begin{enumerate}
\item 如果從 $s$ 走得到的地方沒有負環，它回傳 \TRUE，而且每個頂點都有 $d[v] = \delta(s,
v)$，$\pi$ 構成最短路徑樹。
\item 如果從 $s$ 走得到一個負環，它回傳 \FALSE。
\end{enumerate}
\end{theorem}

\begin{proof}
(1) 設 $v$ 從 $s$ 走得到。 由引理~\ref{lem:spt-simple}，有一條從 $s$ 到 $v$ 的最短路徑
$p = \langle v_0, v_1, \ldots, v_k \rangle$ 是簡單路徑，$k \le n - 1$。 每一輪都會鬆弛所有
的邊，所以第 $i$ 輪一定鬆弛過 $(v_{i-1}, v_i)$，這 $k$ 條邊在第 $1, 2, \ldots, k$ 輪依序被
鬆弛過。 由路徑鬆弛性質，第 $k$ 輪結束後 $d[v] = \delta(s, v)$，之後不變。 走不到的頂點由
無路徑性質是 $\infty = \delta$。 第 6--8 行檢查時，對每條邊由三角不等式，
\[
d[v] = \delta(s, v) \le \delta(s, u) + w(u, v) = d[u] + w(u, v),
\]
所以第 7 行的條件都不成立，回傳 \TRUE。 $\pi$ 由定理~\ref{thm:spt-tree} 構成最短路徑樹。

(2) 設負環 $c = \langle v_0, v_1, \ldots, v_k = v_0 \rangle$ 從 $s$ 走得到，$w(c) < 0$。 先看
環上的 $d$ 都是有限的：對 $i$ 歸納可知第 $i$ 輪結束後，凡是能用 $\le i$ 條邊從 $s$ 走到的頂
點，$d$ 都有限；環上每個頂點都能用一條 $\le n - 1$ 條邊的簡單路徑走到。 假設演算法回傳
\TRUE，那麼第 7 行對環上每條邊都不成立：
\[
d[v_i] \le d[v_{i-1}] + w(v_{i-1}, v_i), \qquad i = 1, 2, \ldots, k.
\]
把 $k$ 個不等式加起來。 因為 $v_0 = v_k$，左邊的 $\sum d[v_i]$ 和右邊的 $\sum d[v_{i-1}]$
是同一群數的和，又都有限，消掉之後得到 $0 \le w(c)$，矛盾。
\end{proof}

注意 (2) 只管「從 $s$ 走得到」的負環。 走不到的負環，上面的 $d$ 都是 $\infty$，第 7 行永
遠不成立，演算法照樣回傳 \TRUE；這時從 $s$ 出發的答案本來就不受影響，回傳 \TRUE 沒有錯。
要偵測\textbf{整張圖}有沒有負環，常見的做法是加一個新頂點 $q$，從 $q$ 連一條權重 $0$ 的
邊到每個頂點，再從 $q$ 跑 Bellman--Ford；等價的寫法是一開始把所有頂點的 $d$ 都設成 $0$。

\textbf{執行時間。} 初始化 $\Theta(n)$，第 3--5 行 $(n - 1) \cdot m$ 次鬆弛，檢查 $m$ 次，共
$\Theta(nm)$。 圖用鄰接矩陣存時，「掃所有的邊」要看 $n^2$ 格，變成 $\Theta(n^3)$。 和
Dijkstra 演算法的 $O(m \lg n)$ 比，慢了大約 $n / \lg n$ 倍，這是處理負邊的代價。

\textbf{提早結束。} 每一輪做什麼只由當時的 $d$ 決定。 如果某一輪沒有任何 $d$ 改變，下一輪
看到的 $d$ 一模一樣，也不會有任何改變，以後都是這樣；所以可以直接停。 停下來的那一輪本身就
是「沒有邊能鬆弛」，等於第 6--8 行的檢查已經通過，回傳 \TRUE。 圖 $B$ 照本書的順序在第 3 輪
就能停。 另一種常見的加速只鬆弛「上一輪 $d$ 有變過的頂點」的出邊，用一個佇列記這些頂點；
實務上常快很多，但最壞情況仍是 $O(nm)$。

\textbf{和第~18~章的關係。} 第~\ref{sec:dp-apsp}~節最後「限制邊數」的 DP，限制成單一起點，
就是作答提醒框裡的 $\mathit{dist}^{k}$。 Bellman--Ford 是它就地更新的版本：第 $i$ 輪結束後，
$d[v]$ \textbf{不大於}任何一條最多 $i$ 條邊、從 $s$ 到 $v$ 的路徑的權重（同一輪用到新的值，
只會更小），所以收斂得一樣快或更快。

\subsection{找出負環}
\label{sec:spt-negcycle}

把圖 $B$ 的 $5 \to 3$ 的權重改成 $1$，叫它 $B'$。 環 $3 \to 2 \to 4 \to 5 \to 3$ 的權重變成
$-3 + 2 - 1 + 1 = -1$，是負環。 照同樣的邊順序：

\begin{center}
\begin{tabular}{c*{5}{c}l}
\toprule
輪 & $d[1]$ & $d[2]$ & $d[3]$ & $d[4]$ & $d[5]$ & 這一輪成功鬆弛的邊 \\
\midrule
1 & $0$ & $0$ & $3$ & $7$ & $6$ & $(1,2)$ $(1,3)$ $(2,4)$ $(2,5)$ $(3,2)$ $(4,5)$ \\
2 & $0$ & $0$ & \sptnew{2} & \sptnew{2} & \sptnew{1} & $(2,4)$ $(4,5)$ $(5,3)$ \\
3 & $0$ & \sptnew{$-1$} & $2$ & $2$ & $1$ & $(3,2)$ \\
4 & $0$ & $-1$ & \sptnew{1} & \sptnew{1} & \sptnew{0} & $(2,4)$ $(4,5)$ $(5,3)$ \\
\midrule
檢查 & \multicolumn{5}{c}{$(3, 2)$：$1 - 3 = -2 < -1$} & 回傳 \FALSE \\
\bottomrule
\end{tabular}
\end{center}

環上的 $d$ 每繞一圈就少 $1$，永遠停不下來；做完 $n - 1 = 4$ 輪，檢查時 $(3, 2)$ 還能鬆弛。

只知道「有負環」常常不夠，還要把它找出來。 方法是把檢查那一輪也照樣鬆弛，記下最後一個 $d$
變小的頂點 $x$，然後從 $x$ 沿著 $\pi$ 往回走 $n$ 步。 走完之後所在的頂點一定在一個由 $\pi$
邊組成的環上，再沿著 $\pi$ 繞一圈回到自己，就是負環（倒過來就是邊的方向）。 $B'$ 的檢查
那一輪把 $d[2]$ 改成 $-2$，$x = 2$；$\pi$ 是 $\pi[2] = 3$、$\pi[3] = 5$、$\pi[5] = 4$、
$\pi[4] = 2$，從 $2$ 往回走 $5$ 步是 $2, 3, 5, 4, 2, 3$，停在 $3$；從 $3$ 繞一圈是 $3, 5, 4,
2, 3$，倒過來得到負環 $2 \to 4 \to 5 \to 3 \to 2$。

\begin{whybox}[為什麼往回走 $n$ 步會落在負環上]
先證明一個不變量：第 $i$ 輪結束後，$d[v]$ 不大於任何一條最多 $i$ 條邊、從 $s$ 到 $v$ 的路
徑的權重。 對 $i$ 歸納。 $i = 0$ 時只有 $s$ 有長度 $0$ 的路徑，$d[s] = 0$。 $i \ge 1$ 時，設 $p$ 是
一條最多 $i$ 條邊、從 $s$ 到 $v$ 的路徑；長度 $0$ 的情形同上，否則設最後一條邊是 $(u, v)$，前面那段
$p'$ 最多 $i - 1$ 條邊。 由歸納假設，第 $i - 1$ 輪結束時 $d[u] \le w(p')$；$d$ 只會變小，所以第
$i$ 輪鬆弛 $(u, v)$ 的時候仍有 $d[u] \le w(p')$，鬆弛之後 $d[v] \le d[u] + w(u, v) \le w(p)$，之後
也只會更小。

設第 $n$ 輪（檢查那一輪）把 $d[x]$ 變小。 從 $x$ 沿著 $\pi$ 往回走，假設會走到 \NIL。
$d$ 有限而 $\pi = \NIL$ 的頂點只有 $s$，而且這時 $d[s]$ 從沒被改過，還是 $0$。 這條
$\pi$ 鏈如果有重複頂點就已經是環了，所以它是一條簡單路徑 $q$，最多 $n - 1$ 條邊。 由定理
~\ref{thm:spt-tree} 證明裡的 (i)，鏈上每條邊都有 $d[\text{頭}] \ge d[\text{尾}] + w$，加起來
得到 $d[x] \ge d[s] + w(q) = w(q)$。 可是第 $n - 1$ 輪結束時已經有 $d[x] \le w(q)$（不變量，
$q$ 最多 $n - 1$ 條邊），第 $n$ 輪又讓 $d[x]$ 嚴格變小，所以現在 $d[x] < w(q)$，矛盾。

所以從 $x$ 往回走永遠不會碰到 \NIL，在 $n$ 個頂點裡一直走，$n$ 步之內一定進入一個環，走
滿 $n$ 步時已經在環上。 由定理~\ref{thm:spt-tree} 證明裡的 (ii)，由 $\pi$ 邊組成的環一定
是負環。
\end{whybox}

\textbf{哪些頂點是 $-\infty$。} 從 $s$ 走得到的負環，以及從這些負環走得到的頂點，$\delta$ 都
是 $-\infty$。 做法：在檢查那一輪，把每條還能鬆弛的邊 $(u, v)$ 的 $v$ 記下來，再從這些頂點做
一次 BFS 或 DFS，走到的都是 $-\infty$；其他頂點的 $d$ 就是正確的 $\delta$。 理由有兩半：
從 $s$ 走得到的每個負環上，至少有一條邊在檢查時還能鬆弛（定理~\ref{thm:spt-bf} (2) 的證
明），所以每個負環都會被找到；反過來，如果 $v$ 不能從任何這種負環走到，到 $v$ 的路徑上沒有
負環，引理~\ref{lem:spt-simple} 的論證對 $v$ 照樣成立，做完 $n - 1$ 輪時 $d[v] = \delta(s, v)$，
指向 $v$ 的邊都鬆弛不動。 $B'$ 裡 $2, 3, 4, 5$ 都是 $-\infty$，只有 $d[1] = 0$ 是對的。

\subsection{C 程式}

圖存成邊的陣列，照輸入的順序掃，這樣「邊的順序」就是輸入的順序，和手算一致。 程式每一輪
都印出 $d$；檢查那一輪也照樣鬆弛，記下最後一個變小的頂點，有的話就照上面的方法印出負環，
沒有的話印出 $\pi$。 \texttt{relax\_all} 在 \texttt{d[u]} 是 $\infty$ 時跳過，原因在
第~\ref{sec:spt-relax}~節說過：負的權重加上代表 $\infty$ 的大數，會比那個大數小。

\begin{example}[label=ex:spt-bf]{c}{(bellman\_ford.c) Bellman--Ford 演算法與找出負環}
#include <stdio.h>
#include <stdlib.h>
#include <limits.h>

#define INF (LLONG_MAX / 4)

typedef struct {                // 一條邊 u -> v，權重 w
    int u, v;
    long long w;
} Edge;

// 照 e[0..m-1] 的順序把每條邊鬆弛一次
// 回傳最後一個 d 變小的頂點，沒有就回傳 -1
int relax_all(const Edge *e, int m, long long *d, int *pi)
{
    int last = -1;
    for (int i = 0; i < m; i++) {
        if (d[e[i].u] != INF && d[e[i].u] + e[i].w < d[e[i].v]) {
            d[e[i].v] = d[e[i].u] + e[i].w;
            pi[e[i].v] = e[i].u;
            last = e[i].v;
        }
    }
    return last;
}

void print_row(const char *label, const long long *d, int n)
{
    printf("%s", label);
    for (int v = 0; v < n; v++) {
        if (d[v] == INF) {
            printf(" inf");
        } else {
            printf(" %lld", d[v]);
        }
    }
    printf("\n");
}

// 從 x 沿著 pi 往回走 n 步，一定落在負環上；再繞一圈，倒過來印出環
void print_cycle(const int *pi, int n, int x)
{
    int *cyc = malloc((n + 1) * sizeof *cyc), len = 0;
    if (cyc == NULL) {
        return;
    }
    for (int i = 0; i < n; i++) {
        x = pi[x];
    }
    int y = x;
    do {
        cyc[len++] = y;
        y = pi[y];
    } while (y != x);
    printf("negative cycle:");
    for (int i = len - 1; i >= 0; i--) {
        printf(" %d", cyc[i] + 1);
    }
    printf(" %d\n", cyc[len - 1] + 1);
    free(cyc);
}

int main(void)
{
    int n, m, s;
    while (scanf("%d %d %d", &n, &m, &s) == 3) {
        Edge *e = malloc(m * sizeof *e);
        long long *d = malloc(n * sizeof *d);
        int *pi = malloc(n * sizeof *pi);
        if (e == NULL || d == NULL || pi == NULL) {
            return 1;
        }
        for (int i = 0; i < m; i++) {
            if (scanf("%d %d %lld", &e[i].u, &e[i].v,
                      &e[i].w) != 3) {
                return 1;
            }
            e[i].u--;
            e[i].v--;
        }
        for (int v = 0; v < n; v++) {
            d[v] = INF;
            pi[v] = -1;
        }
        d[s - 1] = 0;
        for (int r = 1; r <= n - 1; r++) {
            char label[32];
            relax_all(e, m, d, pi);
            snprintf(label, sizeof label, "round %d:", r);
            print_row(label, d, n);
        }
        // 第 n 輪：還有邊能鬆弛就是有負環
        int x = relax_all(e, m, d, pi);
        if (x == -1) {
            printf("pi:");
            for (int v = 0; v < n; v++) {
                if (pi[v] == -1) {
                    printf(" -");
                } else {
                    printf(" %d", pi[v] + 1);
                }
            }
            printf("\n");
        } else {
            print_cycle(pi, n, x);
        }
        free(e);
        free(d);
        free(pi);
    }
    return 0;
}
\end{example}
\iosample{5 8 1 \\ 1 2 5 1 3 3 2 4 2 2 5 6 \\ 3 2 -3 3 4 4 4 5 -1 5 3 3 \\ 5 8 1 \\ 5 3 3 4 5 -1 3 4 4 3 2 -3 \\ 2 5 6 2 4 2 1 3 3 1 2 5 \\ 5 8 1 \\ 1 2 5 1 3 3 2 4 2 2 5 6 \\ 3 2 -3 3 4 4 4 5 -1 5 3 1}{round 1: 0 0 3 7 6 \\ round 2: 0 0 3 2 1 \\ round 3: 0 0 3 2 1 \\ round 4: 0 0 3 2 1 \\ pi: - 3 1 2 4 \\ round 1: 0 5 3 inf inf \\ round 2: 0 0 3 2 6 \\ round 3: 0 0 3 2 1 \\ round 4: 0 0 3 2 1 \\ pi: - 3 1 2 4 \\ round 1: 0 0 3 7 6 \\ round 2: 0 0 2 2 1 \\ round 3: 0 -1 2 2 1 \\ round 4: 0 -1 1 1 0 \\ negative cycle: 2 4 5 3 2}

三張圖依序是照本書順序的 $B$、邊倒過來的 $B$、以及 $B'$，每一輪的 $d$ 都和上面三張表相
同。 程式沒有做提早結束，為的是把 $n - 1$ 輪都印出來和表對照；實際使用時，某一輪
\texttt{relax\_all} 回傳 $-1$ 就可以停了。

\section{有向無環圖上的最短路徑}
\label{sec:spt-dag}

\subsection{依拓樸順序鬆弛}

圖沒有環時，有一個比 Bellman--Ford 快得多、又允許負邊的做法。 第~22~章的拓樸順序有一個
性質：每條邊都從前面的頂點指向後面的頂點。 所以任何一條路徑上的頂點，在拓樸順序裡也是
從前到後排的。 只要依拓樸順序處理頂點，每處理一個頂點就鬆弛它所有的出邊，任何一條路徑
的邊都會照路徑上的先後被鬆弛到，路徑鬆弛性質直接給出答案。

\begin{pseudo}[label=alg:spt-dag, unbreakable]{有向無環圖上的單一起點最短路徑}
\Procedure{DAG-Shortest-Paths}{$G, w, s$}
    \State 把 $G$ 的頂點依拓樸順序排好 \Comment{虛擬碼~\ref{alg:tps-kahn} 或 \ref{alg:tps-dfs}}
    \State \Call{Initialize-Single-Source}{$G, s$}
    \For{依拓樸順序的每個頂點 $u$}
        \For{每個 $v \in \mathit{Adj}[u]$}
            \State \Call{Relax}{$u, v, w$}
        \EndFor
    \EndFor
\EndProcedure
\end{pseudo}

\begin{theorem}
\label{thm:spt-dag}
$G$ 是有向無環圖（權重可以是負的）。 \proc{DAG-Shortest-Paths} 結束時每個頂點都有
$d[v] = \delta(s, v)$，$\pi$ 構成最短路徑樹。 時間是 $\Theta(n + m)$。
\end{theorem}

\begin{proof}
沒有環就沒有負環，引理~\ref{lem:spt-simple} 成立。 設 $s$ 走得到 $v$，取一條從 $s$ 到 $v$ 的
最短路徑 $\langle v_0, v_1, \ldots, v_k \rangle$，$v_0 = s$、$v_k = v$。 拓樸順序裡 $v_0, v_1, \ldots, v_k$ 依序出現，
第 4 行依序處理它們，處理 $v_{i-1}$ 時鬆弛了 $(v_{i-1}, v_i)$；這 $k$ 條邊依路徑的順序被鬆弛
過，由路徑鬆弛性質 $d[v] = \delta(s, v)$。 走不到的由無路徑性質是 $\infty$。 $\pi$ 由
定理~\ref{thm:spt-tree}。 拓樸排序 $\Theta(n + m)$，第 4--6 行每個頂點一次、每條邊一次，也是
$\Theta(n + m)$。
\end{proof}

注意拓樸順序裡排在 $s$ 前面的頂點，$s$ 一定走不到（走得到就要排在 $s$ 後面），處理它們時
$d[u] = \infty$，鬆弛什麼都不做。 所以也可以從 $s$ 在順序裡的位置才開始處理。

\subsection{手算追蹤}

圖~\ref{fig:spt-dag} 的 $D$ 是一張有向無環圖，頂點已經照一個拓樸順序 $1, 2, \ldots, 6$ 從左
排到右（用第~22~章的 Kahn 演算法排出來也是這個順序）。 起點是 $s = 2$。

\begin{center}
\begin{tikzpicture}[x=1cm, y=1cm]
\foreach \i in {1, ..., 6}
    \node[spt-v] (g\i) at ({(\i - 1) * 2.2}, 0) {\i};
\draw[spt-a] (g1) -- node[spt-w, above] {$3$} (g2);
\draw[spt-a] (g2) -- node[spt-w, above] {$2$} (g3);
\draw[spt-neg] (g3) -- node[spt-wn, above] {$-3$} (g4);
\draw[spt-a] (g4) -- node[spt-w, above] {$1$} (g5);
\draw[spt-neg] (g5) -- node[spt-wn, above] {$-2$} (g6);
\draw[spt-a] (g1) to[bend left=40] node[spt-w, above] {$2$} (g3);
\draw[spt-a] (g3) to[bend left=40] node[spt-w, above] {$4$} (g5);
\draw[spt-a] (g2) to[bend right=40] node[spt-w, below] {$6$} (g4);
\draw[spt-a] (g4) to[bend right=40] node[spt-w, below] {$5$} (g6);
\end{tikzpicture}\par\nopagebreak
\diagcaption{有向無環圖 $D$，頂點依拓樸順序由左到右。}
\label{fig:spt-dag}
\end{center}

每一列是處理完一個頂點（鬆弛完它所有的出邊）之後的 $d$：

\begin{center}
\begin{tabular}{c*{6}{c}l}
\toprule
處理 $u$ & $d[1]$ & $d[2]$ & $d[3]$ & $d[4]$ & $d[5]$ & $d[6]$ & 說明 \\
\midrule
初始 & $\infty$ & $0$ & $\infty$ & $\infty$ & $\infty$ & $\infty$ & \\
$1$ & $\infty$ & $0$ & $\infty$ & $\infty$ & $\infty$ & $\infty$ & $d[1] = \infty$，不鬆弛 \\
$2$ & $\infty$ & $0$ & \sptnew{2} & \sptnew{6} & $\infty$ & $\infty$ & \\
$3$ & $\infty$ & $0$ & $2$ & \sptnew{$-1$} & \sptnew{6} & $\infty$ & $2 - 3 < 6$ \\
$4$ & $\infty$ & $0$ & $2$ & $-1$ & \sptnew{0} & \sptnew{4} & \\
$5$ & $\infty$ & $0$ & $2$ & $-1$ & $0$ & \sptnew{$-2$} & \\
$6$ & $\infty$ & $0$ & $2$ & $-1$ & $0$ & $-2$ & 沒有出邊 \\
\midrule
$\pi$ & \NIL & \NIL & $2$ & $3$ & $4$ & $5$ & \\
\bottomrule
\end{tabular}
\end{center}

最短路徑樹是 $2 \to 3 \to 4 \to 5 \to 6$ 這條鏈。 和 Bellman--Ford 不同，這裡每條邊只鬆弛一
次，處理到 $u$ 時 $d[u]$ 已經是最後的答案（$u$ 的前驅都在它前面處理完了）。

\subsection{最長路徑與關鍵路徑}

一般的圖上，最長簡單路徑沒有已知的多項式時間演算法（第~\ref{sec:dp-substructure}~節）。 有向無環圖沒有
環，路徑一定是簡單的，最長路徑就容易了：把每條邊的權重變號，求最短路徑，答案再變號回來；
或者直接把 \proc{Relax} 的 $<$ 改成 $>$、初值的 $\infty$ 改成 $-\infty$。 第~22~章
第~\ref{sec:tps-aoe}~節 AOE 網路的最早時間 $\mathit{ee}[j]$，正是從開始事件出發的最長路徑，
當時依拓樸順序取 $\max$ 的做法，就是這一節的演算法把 $\min$ 換成 $\max$。 $D$ 從 $2$ 出發
的最長距離是 $d[3..6] = 2, 6, 7, 11$，路徑是 $2 \to 4 \to 6$（$6 + 5 = 11$）。

同一個架構還能算很多東西，只要答案能「依拓樸順序從前驅推過來」。 例如從 $s$ 到每個頂點的
路徑條數：$c[s] = 1$，其他先設成 $0$，依拓樸順序處理 $u$ 時，對每條出邊 $(u, v)$ 做
$c[v] \gets c[v] + c[u]$。 只數最短路徑的條數，就只在 $d[u] + w(u, v) = d[v]$ 的邊上累加
（先把 $d$ 算完，再掃一遍）。

\textbf{C 程式。} 這一節不另外寫完整的程式：把第~22~章範例~\ref{ex:tps-topo} 的 Kahn 演算法排出
的順序 \texttt{q[0..n-1]} 依序走一遍，對每個 $d$ 不是 $\infty$ 的頂點，照範例~\ref{ex:spt-bf} 的
\texttt{relax\_all} 那樣鬆弛它的出邊就好。 要求最長路徑，把比較的 \texttt{<} 改成 \texttt{>}、
初值改成代表 $-\infty$ 的 \texttt{-INF}，記得一樣要先跳過 $d$ 是 $-\infty$ 的頂點。

\section{所有點對最短路徑：Floyd--Warshall 與遞移閉包}
\label{sec:spt-fw}

\subsection{輸入與輸出}

所有點對的問題通常用矩陣表示（第~\ref{sec:gtr-repr}~節的加權鄰接矩陣）。 頂點編號 $1..n$，輸入是
$n \times n$ 的\textbf{權重矩陣} $W = (w_{ij})$：
\[
w_{ij} = \begin{cases}
0, & i = j, \\
w(i, j), & i \ne j \text{ 而且 } (i, j) \in E, \\
\infty, & i \ne j \text{ 而且 } (i, j) \notin E.
\end{cases}
\]
輸出是距離矩陣 $D = (d_{ij})$，$d_{ij} = \delta(i, j)$，以及\textbf{前驅矩陣} $\Pi =
(\pi_{ij})$：$\pi_{ij}$ 是某一條從 $i$ 到 $j$ 的最短路徑上 $j$ 的前一個頂點；$i = j$ 或到不了
時是 \NIL。 $\Pi$ 的第 $i$ 列，就是從 $i$ 出發的最短路徑樹的 $\pi$。 這一節先假設圖沒有負環，
最後再看有負環時會怎樣。

最直接的做法是對每個頂點各跑一次單一起點的演算法。 權重不是負數時跑 $n$ 次 Dijkstra：陣
列版共 $\Theta(n^3)$，二元堆積版 $O(nm \lg n)$。 有負邊時跑 $n$ 次 Bellman--Ford，共
$\Theta(n^2 m)$，稠密圖上是 $\Theta(n^4)$。 下面三個方法都比後者好。

\subsection{用邊數當子問題：像矩陣乘法的做法}

第~\ref{sec:dp-apsp}~節最後提過另一種子問題：$\ell_{ij}^{(k)}$ 是從 $i$ 到 $j$、\textbf{最多用 $k$
條邊}的路徑的最小權重。 $k = 1$ 時就是 $W$（$w_{ii} = 0$ 讓「不動」也算一種走法）。 看最後
一條邊從哪個頂點 $t$ 走進 $j$：
\[
\ell_{ij}^{(k)} = \min_{1 \le t \le n} \bigl( \ell_{it}^{(k-1)} + w_{tj} \bigr).
\]
（$t = j$ 那一項是 $\ell_{ij}^{(k-1)} + 0$，對應「不用第 $k$ 條邊」。） 沒有負環時，最短路徑
可以取簡單路徑，最多 $n - 1$ 條邊（引理~\ref{lem:spt-simple}），所以 $\delta(i, j) =
\ell_{ij}^{(n-1)} = \ell_{ij}^{(n)} = \cdots$。

\begin{pseudo}[label=alg:spt-extend]{把最短路徑多延伸一條邊}
\Function{Extend-Paths}{$L, W, n$}
    \State $L' \gets$ 新的 $n \times n$ 矩陣
    \For{$i \gets 1$ \To $n$}
        \For{$j \gets 1$ \To $n$}
            \State $L'[i, j] \gets \infty$
            \For{$t \gets 1$ \To $n$}
                \State $L'[i, j] \gets \min(L'[i, j],\ L[i, t] + W[t, j])$
            \EndFor
        \EndFor
    \EndFor
    \State \Return $L'$
\EndFunction
\end{pseudo}

把第 5 行的 $\infty$ 換成 $0$、第 7 行的 $\min$ 換成 $+$、$+$ 換成 $\times$，就是一般的矩陣乘
法 $C = A \cdot B$。 所以把 \proc{Extend-Paths} 看成一種「乘法」，記成 $L \otimes W$，
$L^{(k)} = L^{(k-1)} \otimes W$，答案是 $W$ 自己「乘」$n - 2$ 次，每次 $\Theta(n^3)$，共
$\Theta(n^4)$。

這種乘法和一般矩陣乘法一樣有結合律。 一般矩陣乘法的結合律靠的是分配律 $a(b + c) = ab + ac$；
這裡 $\min$ 扮演 $+$、$+$ 扮演 $\times$，對應的分配律 $a + \min(b, c) = \min(a + b, a + c)$ 也成
立，同樣的證明照搬過來就好。 所以可以用第~\ref{sec:dc-power}~節快速冪的反覆平方：
\[
L^{(1)} = W,\quad L^{(2)} = W \otimes W,\quad L^{(4)} = L^{(2)} \otimes L^{(2)},\quad \ldots
\]
算到 $L^{(2^r)}$ 而 $2^r \ge n - 1$ 為止，因為超過 $n - 1$ 條邊也不會更短。 只要
$\lceil \lg(n - 1) \rceil$ 次乘法，共 $\Theta(n^3 \lg n)$。 不能再用 Strassen 加速：它需要減
法，而 $\min$ 沒有對應的「減法」（第~\ref{sec:dc-strassen}~節）。

以圖~\ref{fig:spt-fw} 的 $F$ 為例（$n = 4$，只要算到 $L^{(4)}$，兩次乘法）：
\[
L^{(1)} = W = \begin{pmatrix} 0 & 3 & \infty & 8 \\ \infty & 0 & -2 & 6 \\ 4 & \infty & 0 & 2 \\ 2 & 1 & \infty & 0 \end{pmatrix},\quad
L^{(2)} = \begin{pmatrix} 0 & 3 & 1 & 8 \\ 2 & 0 & -2 & 0 \\ 4 & 3 & 0 & 2 \\ 2 & 1 & -1 & 0 \end{pmatrix},\quad
L^{(4)} = \begin{pmatrix} 0 & 3 & 1 & 3 \\ 2 & 0 & -2 & 0 \\ 4 & 3 & 0 & 2 \\ 2 & 1 & -1 & 0 \end{pmatrix}.
\]
例如 $\ell_{14}^{(2)} = 8$：兩條邊以內，$1 \to 2 \to 4$ 是 $9$，不如直接的 $8$；
$\ell_{14}^{(4)} = 3$ 是三條邊的 $1 \to 2 \to 3 \to 4$。 $L^{(4)}$ 就是答案，和下面
Floyd--Warshall 算出的 $D^{(4)}$ 相同。

\subsection{Floyd--Warshall：中間點當子問題}

第~\ref{sec:dp-apsp}~節已經推過遞迴式：$d_{ij}^{(k)}$ 是從 $i$ 到 $j$、中間點都在 $\{1, \ldots, k\}$
裡的最短路徑權重，
\[
d_{ij}^{(0)} = w_{ij}, \qquad
d_{ij}^{(k)} = \min\bigl( d_{ij}^{(k-1)},\ d_{ik}^{(k-1)} + d_{kj}^{(k-1)} \bigr),
\]
而且可以只用一個矩陣就地更新（虛擬碼~\ref{alg:dp-fw}），時間 $\Theta(n^3)$、空
間 $\Theta(n^2)$。 這裡補上 $\Pi$。 $k = 0$ 時，有邊 $(i, j)$ 就 $\pi_{ij}^{(0)} = i$，否則是
\NIL。 第 $k$ 層如果改用經過 $k$ 的路徑 $i \leadsto k \leadsto j$，$j$ 的前一個頂點就是
$k \leadsto j$ 那一段裡 $j$ 的前一個，也就是 $\pi_{kj}^{(k-1)}$：
\[
\pi_{ij}^{(k)} = \begin{cases}
\pi_{kj}^{(k-1)}, & d_{ik}^{(k-1)} + d_{kj}^{(k-1)} < d_{ij}^{(k-1)}, \\
\pi_{ij}^{(k-1)}, & \text{其他}.
\end{cases}
\]

\begin{pseudo}[label=alg:spt-fw]{Floyd--Warshall 演算法，含前驅矩陣}
\Function{Floyd-Warshall-Path}{$W, n$}
    \For{$i \gets 1$ \To $n$}
        \For{$j \gets 1$ \To $n$}
            \State $D[i, j] \gets W[i, j]$；$\Pi[i, j] \gets \NIL$
            \If{$i \ne j$ \AND $W[i, j] < \infty$}
                \State $\Pi[i, j] \gets i$
            \EndIf
        \EndFor
    \EndFor
    \For{$k \gets 1$ \To $n$}
        \For{$i \gets 1$ \To $n$}
            \For{$j \gets 1$ \To $n$}
                \If{$D[i, k] + D[k, j] < D[i, j]$}
                    \State $D[i, j] \gets D[i, k] + D[k, j]$
                    \State $\Pi[i, j] \gets \Pi[k, j]$
                \EndIf
            \EndFor
        \EndFor
    \EndFor
    \State \Return $D, \Pi$
\EndFunction
\end{pseudo}

第 2--6 行初始化，第 7--12 行是三層迴圈，$k$ 一定在最外層（第~18~章說過放錯位置會錯，觀念題
第~\ref{spt-q-order}~題要你找一個反例）。 就地更新沒有問題：第 $k$ 層裡第 $k$ 列和第 $k$ 行
不會被改（$d_{kk} = 0$，$d_{ik} + d_{kk} = d_{ik}$ 不會更小），所以第 10--12 行讀到的
$D[i, k]$、$D[k, j]$、$\Pi[k, j]$ 都還是第 $k - 1$ 層的值。

\subsection{手算：$D^{(k)}$ 和 $\Pi^{(k)}$ 的序列}
\label{sec:spt-fw-trace}

考試最常見的寫法是把 $D^{(0)}, D^{(1)}, \ldots, D^{(n)}$（常常連同 $\Pi^{(k)}$）一個一個寫出
來。 用圖~\ref{fig:spt-fw} 的 $F$：$4$ 個頂點、$8$ 條邊，$2 \to 3$ 是負邊，沒有負環。

\begin{center}
\begin{tikzpicture}[x=1cm, y=1cm]
\node[spt-v] (f1) at (0, 0) {1};
\node[spt-v] (f2) at (2.4, 1.8) {2};
\node[spt-v] (f3) at (4.8, 0) {3};
\node[spt-v] (f4) at (2.4, -1.8) {4};
\draw[spt-a] (f1) -- node[spt-w, above left] {$3$} (f2);
\draw[spt-neg] (f2) -- node[spt-wn, above right] {$-2$} (f3);
\draw[spt-a] (f3) -- node[spt-w, below right] {$2$} (f4);
\draw[spt-a] (f1) to[bend left=25] node[spt-w, pos=0.5] {$8$} (f4);
\draw[spt-a] (f4) to[bend left=25] node[spt-w, pos=0.5] {$2$} (f1);
\draw[spt-a] (f2) to[bend left=25] node[spt-w] {$6$} (f4);
\draw[spt-a] (f4) to[bend left=25] node[spt-w] {$1$} (f2);
\draw[spt-a] (f3) to[out=-90, in=-90, looseness=1.6] node[spt-w, below] {$4$} (f1);
\end{tikzpicture}\par\nopagebreak
\diagcaption{圖 $F$。 $1$、$4$ 之間和 $2$、$4$ 之間兩個方向各有一條邊。}
\label{fig:spt-fw}
\end{center}

每一層\textbf{粗體}是這一層變小的格子，$\Pi$ 裡的 \NIL 寫成「$-$」。
\begin{gather*}
D^{(0)} = \begin{pmatrix} 0 & 3 & \infty & 8 \\ \infty & 0 & -2 & 6 \\ 4 & \infty & 0 & 2 \\ 2 & 1 & \infty & 0 \end{pmatrix},\quad
\Pi^{(0)} = \begin{pmatrix} - & 1 & - & 1 \\ - & - & 2 & 2 \\ 3 & - & - & 3 \\ 4 & 4 & - & - \end{pmatrix}; \\
D^{(1)} = \begin{pmatrix} 0 & 3 & \infty & 8 \\ \infty & 0 & -2 & 6 \\ 4 & \mathbf{7} & 0 & 2 \\ 2 & 1 & \infty & 0 \end{pmatrix},\quad
\Pi^{(1)} = \begin{pmatrix} - & 1 & - & 1 \\ - & - & 2 & 2 \\ 3 & \mathbf{1} & - & 3 \\ 4 & 4 & - & - \end{pmatrix}; \\
D^{(2)} = \begin{pmatrix} 0 & 3 & \mathbf{1} & 8 \\ \infty & 0 & -2 & 6 \\ 4 & 7 & 0 & 2 \\ 2 & 1 & \mathbf{-1} & 0 \end{pmatrix},\quad
\Pi^{(2)} = \begin{pmatrix} - & 1 & \mathbf{2} & 1 \\ - & - & 2 & 2 \\ 3 & 1 & - & 3 \\ 4 & 4 & \mathbf{2} & - \end{pmatrix}; \\
D^{(3)} = \begin{pmatrix} 0 & 3 & 1 & \mathbf{3} \\ \mathbf{2} & 0 & -2 & \mathbf{0} \\ 4 & 7 & 0 & 2 \\ 2 & 1 & -1 & 0 \end{pmatrix},\quad
\Pi^{(3)} = \begin{pmatrix} - & 1 & 2 & \mathbf{3} \\ \mathbf{3} & - & 2 & \mathbf{3} \\ 3 & 1 & - & 3 \\ 4 & 4 & 2 & - \end{pmatrix}; \\
D^{(4)} = \begin{pmatrix} 0 & 3 & 1 & 3 \\ 2 & 0 & -2 & 0 \\ 4 & \mathbf{3} & 0 & 2 \\ 2 & 1 & -1 & 0 \end{pmatrix},\quad
\Pi^{(4)} = \begin{pmatrix} - & 1 & 2 & 3 \\ 3 & - & 2 & 3 \\ 3 & \mathbf{4} & - & 3 \\ 4 & 4 & 2 & - \end{pmatrix}.
\end{gather*}

逐層看：
\begin{itemize}
\item $k = 1$：第 $1$ 行裡只有 $d_{31} = 4$、$d_{41} = 2$ 是有限的，第 $1$ 列有限的是 $d_{12}
= 3$、$d_{14} = 8$。 候選是 $d_{32}$：$4 + 3 = 7 < \infty$，改；$d_{34}$：$4 + 8 = 12$，不比
$2$ 小；$d_{42}$：$2 + 3 = 5$，不比 $1$ 小；$d_{44}$：不會變。 $\pi_{32} \gets \pi_{12} = 1$。
\item $k = 2$：$d_{13} = 3 + (-2) = 1$，$d_{43} = 1 + (-2) = -1$。 $d_{14}$ 經過 $2$ 是 $3 + 6
= 9$，不比 $8$ 小。
\item $k = 3$：$d_{14} = 1 + 2 = 3$（$1 \to 2 \to 3 \to 4$），$d_{21} = -2 + 4 = 2$，$d_{24} =
-2 + 2 = 0$。 $\pi_{14} \gets \pi_{34} = 3$，$\pi_{21} \gets \pi_{31} = 3$。
\item $k = 4$：$d_{32} = 2 + 1 = 3$（$3 \to 4 \to 2$），比 $k = 1$ 時的 $7$ 好；$\pi_{32} \gets
\pi_{42} = 4$。
\end{itemize}

手算時有兩個省力的地方。 第 $k$ 層的第 $k$ 列和第 $k$ 行照抄。 第 $k$ 行裡是 $\infty$ 的那些
列、第 $k$ 列裡是 $\infty$ 的那些行，這一層也都不會變。 例如 $k = 1$ 時第 $1$ 行的 $d_{21} =
\infty$，第 $2$ 列整列不必看；第 $1$ 列的 $d_{13} = \infty$，第 $3$ 行整行不必看，剩下第 $3$、$4$
列和第 $2$、$4$ 行交叉的四格，就是上面列出的那四個候選。

\subsection{還原路徑}

從 $\Pi$ 印出 $i$ 到 $j$ 的路徑，和 \proc{Print-Path} 一樣，只是 $\pi$ 換成 $\Pi$ 的第 $i$
列：

\begin{pseudo}[label=alg:spt-fwpath]{用前驅矩陣印出路徑}
\Procedure{Print-All-Pairs-Path}{$\Pi, i, j$}
    \If{$i = j$}
        \State 印出 $i$
    \ElsIf{$\Pi[i, j] = \NIL$}
        \State 印出「從 $i$ 到不了 $j$」
    \Else
        \State \Call{Print-All-Pairs-Path}{$\Pi, i, \Pi[i, j]$}
        \State 印出 $j$
    \EndIf
\EndProcedure
\end{pseudo}

$F$ 裡 $1$ 到 $4$：$\pi_{14} = 3$，$\pi_{13} = 2$，$\pi_{12} = 1$，路徑 $1 \to 2 \to 3 \to 4$，權
重 $3 - 2 + 2 = 3 = d_{14}$。 $3$ 到 $2$：$\pi_{32} = 4$，$\pi_{34} = 3$，路徑 $3 \to 4 \to 2$，
權重 $2 + 1 = 3$。

另一種常見的記法是記\textbf{中繼點}：$p_{ij}$ 記下讓 $d_{ij}$ 最後一次變小的那個 $k$，沒有變
過就是 $0$（直接走邊）。 印路徑時遞迴印 $i \leadsto p_{ij}$ 和 $p_{ij} \leadsto j$ 兩段。 兩種
記法存的東西不一樣，題目要哪一種就填哪一種，見下面的作答提醒框。

\subsection{有負環時}

Floyd--Warshall 不會在負環上停下來，它照樣跑完三層迴圈，但有負環時有些 $d_{ij}$ 不是
$\delta(i, j)$（$\delta$ 是 $-\infty$）。 判斷的方法很簡單：\textbf{跑完之後，若某個
$d_{ii} < 0$，圖就有負環，而且 $i$ 在某個負環上}；反過來，有負環時，負環上一定有某個頂點
$i$ 的 $d_{ii} < 0$。

理由：$D[i, j]$ 任何時候都等於某一條從 $i$ 到 $j$ 的走法（可以重複頂點）的權重，所以
$d_{ii} < 0$ 表示有一條從 $i$ 回到 $i$ 的走法權重是負的，它裡面一定有一個負環。 反過來，設
$c$ 是一個簡單的負環，$k$ 是 $c$ 上編號最大的頂點，$i$ 是 $c$ 上另一個頂點。 環上從 $i$ 到 $k$、
從 $k$ 到 $i$ 的兩段，中間點都小於 $k$；第 $k - 1$ 層結束時 $D[i, k]$、$D[k, i]$ 不大於這兩段
的權重（就地更新只會讓值更小），所以第 $k$ 層時 $D[i, i] \le D[i, k] + D[k, i] \le w(c) < 0$。

\subsection{C 程式}

矩陣直接用二維陣列，頂點數上限 \texttt{MAXN} 寫死（第~22~章的 \texttt{aoe.c} 也是這樣）。 讀邊時
遇到重複的邊只留權重最小的。 第 $k$ 層開始時如果 \texttt{d[i][k]} 是 $\infty$，這一列整列跳過，
這就是上面說的省力技巧，也順便避免拿 $\infty$ 去相加。 程式印出每一層改了哪些格子，最後印出
$D$、$\Pi$ 和幾條路徑；有負環時只印出在負環上的頂點。

\begin{example}[label=ex:spt-fw]{c}{(floyd.c) Floyd--Warshall 演算法：前驅矩陣、路徑與負環}
#include <stdio.h>
#include <limits.h>

#define MAXN 100
#define INF (LLONG_MAX / 4)

int n;
long long d[MAXN][MAXN];        // 距離矩陣 D
int p[MAXN][MAXN];              // 前驅矩陣 Π：-1 代表 NIL

void floyd_warshall(void)
{
    for (int k = 0; k < n; k++) {
        printf("k=%d:", k + 1);
        for (int i = 0; i < n; i++) {
            if (d[i][k] == INF) {               // 這一列不會變
                continue;
            }
            for (int j = 0; j < n; j++) {
                if (d[k][j] != INF && d[i][k] + d[k][j] < d[i][j]) {
                    d[i][j] = d[i][k] + d[k][j];
                    p[i][j] = p[k][j];
                    printf(" (%d,%d)", i + 1, j + 1);
                }
            }
        }
        printf("\n");
    }
}

void print_path(int i, int j)
{
    if (i != j) {
        print_path(i, p[i][j]);
        printf(" ");
    }
    printf("%d", j + 1);
}

int main(void)
{
    int m, q, u, v, neg;
    long long w;
    while (scanf("%d %d", &n, &m) == 2 && n <= MAXN) {
        for (int i = 0; i < n; i++) {
            for (int j = 0; j < n; j++) {
                d[i][j] = (i == j) ? 0 : INF;
                p[i][j] = -1;
            }
        }
        for (int e = 0; e < m; e++) {
            if (scanf("%d %d %lld", &u, &v, &w) != 3) {
                return 1;
            }
            if (w < d[u - 1][v - 1]) {          // 重複的邊只留最小的
                d[u - 1][v - 1] = w;
                p[u - 1][v - 1] = u - 1;
            }
        }
        floyd_warshall();
        neg = 0;
        for (int i = 0; i < n; i++) {
            if (d[i][i] < 0) {
                printf(neg ? " %d" : "on a negative cycle: %d",
                       i + 1);
                neg = 1;
            }
        }
        if (neg) {
            printf("\n");
        }
        for (int i = 0; i < n && !neg; i++) {
            printf("D%d:", i + 1);
            for (int j = 0; j < n; j++) {
                if (d[i][j] == INF) {
                    printf(" inf");
                } else {
                    printf(" %lld", d[i][j]);
                }
            }
            printf(" P%d:", i + 1);
            for (int j = 0; j < n; j++) {
                if (p[i][j] == -1) {
                    printf(" -");
                } else {
                    printf(" %d", p[i][j] + 1);
                }
            }
            printf("\n");
        }
        if (scanf("%d", &q) != 1) {
            q = 0;
        }
        for (int t = 0; t < q && scanf("%d %d", &u, &v) == 2; t++) {
            if (neg) {
                continue;
            }
            if (d[u - 1][v - 1] == INF) {
                printf("%d -> %d: no path\n", u, v);
            } else {
                printf("%d -> %d: ", u, v);
                print_path(u - 1, v - 1);
                printf(" (length %lld)\n", d[u - 1][v - 1]);
            }
        }
    }
    return 0;
}
\end{example}
\iosample{4 8 \\ 1 2 3 1 4 8 2 3 -2 2 4 6 \\ 3 1 4 3 4 2 4 2 1 4 1 2 \\ 2 3 2 1 4 \\ 2 2 \\ 1 2 1 2 1 -3 \\ 0 \\ 2 1 \\ 1 2 5 \\ 1 2 1}{k=1: (3,2) \\ k=2: (1,3) (4,3) \\ k=3: (1,4) (2,1) (2,4) \\ k=4: (3,2) \\ D1: 0 3 1 3 P1: - 1 2 3 \\ D2: 2 0 -2 0 P2: 3 - 2 3 \\ D3: 4 3 0 2 P3: 3 4 - 3 \\ D4: 2 1 -1 0 P4: 4 4 2 - \\ 3 -> 2: 3 4 2 (length 3) \\ 1 -> 4: 1 2 3 4 (length 3) \\ k=1: (2,2) \\ k=2: (1,1) (1,2) (2,1) (2,2) \\ on a negative cycle: 1 2 \\ k=1: \\ k=2: \\ D1: 0 5 P1: - 1 \\ D2: inf 0 P2: - - \\ 2 -> 1: no path}

第一張圖是 $F$，每一層改的格子就是上面粗體的位置，最後的 \texttt{D}、\texttt{P} 兩列一組，
分別是 $D^{(4)}$ 和 $\Pi^{(4)}$ 的每一列。 第二張圖 $1 \to 2 \to 1$ 的權重是 $1 - 3 = -2$，兩個頂點的
$d_{ii}$ 都變成負的。 第三張圖只有兩個頂點和一條邊 $1 \to 2$，$2$ 到 $1$ 沒有路。

\subsection{遞移閉包}

只想知道「從 $i$ 能不能走到 $j$」時，權重不重要。 有向圖 $G$ 的\textbf{遞移閉包}
（transitive closure）是一個 $n \times n$ 的 $0/1$ 矩陣，記錄每一對頂點之間有沒有路徑。 這裡
有兩種定義，差在對角線：
\begin{itemize}
\item $A^{+}$：$a^{+}_{ij} = 1$ 若且唯若從 $i$ 到 $j$ 有一條\textbf{長度至少 $1$} 的路徑。 所
以 $a^{+}_{ii} = 1$ 表示 $i$ 在某個環上。
\item $A^{*}$（自反遞移閉包）：長度可以是 $0$，所以對角線全是 $1$，其他格子和 $A^{+}$ 相同。
\end{itemize}

把 Floyd--Warshall 的 $\min$ 換成「或」、$+$ 換成「且」，就是 \textbf{Warshall 演算法}：
$t_{ij}^{(k)} = 1$ 表示從 $i$ 到 $j$ 有一條長度至少 $1$、中間點都在 $\{1, \ldots, k\}$ 的路徑，
\[
t_{ij}^{(0)} = a_{ij}, \qquad t_{ij}^{(k)} = t_{ij}^{(k-1)} \lor \bigl( t_{ik}^{(k-1)} \land
t_{kj}^{(k-1)} \bigr).
\]
從鄰接矩陣 $A$（對角線是 $0$）開始，得到的是 $A^{+}$；要 $A^{*}$，最後把對角線設成 $1$，或
一開始就把對角線設成 $1$。

\begin{pseudo}[label=alg:spt-tc]{Warshall 演算法：遞移閉包}
\Function{Transitive-Closure}{$A, n$}
    \State $T \gets A$ 的複本
    \For{$k \gets 1$ \To $n$}
        \For{$i \gets 1$ \To $n$}
            \For{$j \gets 1$ \To $n$}
                \State $T[i, j] \gets T[i, j] \lor \bigl( T[i, k] \land T[k, j] \bigr)$
            \EndFor
        \EndFor
    \EndFor
    \State \Return $T$
\EndFunction
\end{pseudo}

就地更新的理由和 Floyd--Warshall 相同：第 $k$ 層裡 $T[i, k]$ 和 $T[k, j]$ 不會因為這一層而改
變（$T[i, k] \lor (T[i, k] \land T[k, k])$ 還是 $T[i, k]$）。 用位元運算時常數比 Floyd--Warshall
小很多，所以另外寫出來。

\textbf{手算。} 圖~\ref{fig:spt-tc} 的 $C$ 有一個三角形的環 $1 \to 2 \to 3 \to 1$，從 $3$ 走到
$4$，$4$ 和 $5$ 互相連。

\begin{center}
\begin{tikzpicture}[x=1cm, y=1cm]
\node[spt-v] (t1) at (0, 0.9) {1};
\node[spt-v] (t2) at (1.8, 0.9) {2};
\node[spt-v] (t3) at (0.9, -0.5) {3};
\node[spt-v] (t4) at (3.0, -0.5) {4};
\node[spt-v] (t5) at (4.9, -0.5) {5};
\draw[spt-a] (t1) -- (t2);
\draw[spt-a] (t2) -- (t3);
\draw[spt-a] (t3) -- (t1);
\draw[spt-a] (t3) -- (t4);
\draw[spt-a] (t4) to[bend left=20] (t5);
\draw[spt-a] (t5) to[bend left=20] (t4);
\end{tikzpicture}\par\nopagebreak
\diagcaption{圖 $C$：求遞移閉包用的有向圖。}
\label{fig:spt-tc}
\end{center}

粗體是這一層新變成 $1$ 的格子：
\begin{gather*}
T^{(0)} = \begin{pmatrix} 0&1&0&0&0 \\ 0&0&1&0&0 \\ 1&0&0&1&0 \\ 0&0&0&0&1 \\ 0&0&0&1&0 \end{pmatrix},\quad
T^{(1)} = \begin{pmatrix} 0&1&0&0&0 \\ 0&0&1&0&0 \\ 1&\mathbf{1}&0&1&0 \\ 0&0&0&0&1 \\ 0&0&0&1&0 \end{pmatrix},\quad
T^{(2)} = \begin{pmatrix} 0&1&\mathbf{1}&0&0 \\ 0&0&1&0&0 \\ 1&1&\mathbf{1}&1&0 \\ 0&0&0&0&1 \\ 0&0&0&1&0 \end{pmatrix}, \\
T^{(3)} = \begin{pmatrix} \mathbf{1}&1&1&\mathbf{1}&0 \\ \mathbf{1}&\mathbf{1}&1&\mathbf{1}&0 \\ 1&1&1&1&0 \\ 0&0&0&0&1 \\ 0&0&0&1&0 \end{pmatrix},\quad
T^{(4)} = \begin{pmatrix} 1&1&1&1&\mathbf{1} \\ 1&1&1&1&\mathbf{1} \\ 1&1&1&1&\mathbf{1} \\ 0&0&0&0&1 \\ 0&0&0&1&\mathbf{1} \end{pmatrix},\quad
T^{(5)} = \begin{pmatrix} 1&1&1&1&1 \\ 1&1&1&1&1 \\ 1&1&1&1&1 \\ 0&0&0&\mathbf{1}&1 \\ 0&0&0&1&1 \end{pmatrix}.
\end{gather*}

$T^{(5)} = A^{+}$。 對角線全是 $1$，表示五個頂點都在某個環上；把對角線忽略，$1$、$2$、$3$
各走得到另外四個頂點，$4$、$5$ 只走得到彼此，共 $3 \times 4 + 2 \times 1 = 14$ 個有序點對
$(i, j)$，$i \ne j$。 $t_{ij} = t_{ji} = 1$ 的 $i$、$j$ 互相可達，所以 $\{1, 2, 3\}$、$\{4, 5\}$
是兩個強連通元件（第~23~章）。

\begin{examnote}
所有點對的題目，先確認題目要的是哪一種矩陣。 (1) 路徑：有的課本記前驅 $\pi_{ij}$（本書），
有的記最後一次讓 $d_{ij}$ 變小的中繼點 $k$；同一張圖兩種矩陣的內容完全不同，還原路徑的方法也
不同。 (2) 遞移閉包：$A^{+}$ 的對角線要看有沒有環，$A^{*}$ 的對角線全是 $1$；題目只說「遞移
閉包」時，看它的定義或題目給的例子，沒有給就寫明你算的是哪一個。 (3) 寫 $D^{(k)}$ 序列時，
每一層都寫完整的矩陣，並標出改過的格子，閱卷的人才看得出你是一層一層算的。
\end{examnote}

\textbf{時間與 C 實作。} Warshall 演算法是 $\Theta(n^3)$。 另一種做法是從每個頂點各做一次 BFS
或 DFS，共 $\Theta(n(n + m))$，稀疏圖上比較快。 用位元可以把 Warshall 加速：把 $T$ 的每一列存
成一串位元，每 $64$ 個頂點一組放進一個 \texttt{unsigned long long}。 內層迴圈「對每個 $j$，
$T[i, j] \gets T[i, j] \lor T[k, j]$（當 $T[i, k] = 1$）」就變成把第 $k$ 列整列 OR 進第 $i$
列，一次處理 $64$ 個 $j$，時間降到大約 $n^3 / 64$。

\begin{cbookrecap}[第~12~章 12.3、12.6 節，第~13~章 13.2 節]
\texttt{(row[j / 64] >> (j \% 64)) \& 1} 取出第 $j$ 個位元：先找到它在第幾個字，再把它移到最
低位。 \texttt{row[j / 64] |= (Word)1 << (j \% 64)} 把它設成 $1$。 移位的對象要先轉成
\texttt{unsigned long long}，寫成 \texttt{1 << 40} 的話，\texttt{1} 是 \texttt{int}，移太多位
是未定義行為。
\end{cbookrecap}

\begin{example}[label=ex:spt-tc]{c}{(closure.c) Warshall 演算法：用位元存每一列}
#include <stdio.h>
#include <stdlib.h>

typedef unsigned long long Word;    // 一個字存 64 個位元
#define BITS 64

// 第 i 列存成 w 個字：row[i * w + j / BITS] 的第 j % BITS 個位元是 T[i][j]
int get_bit(const Word *row, int j)
{
    return (int)((row[j / BITS] >> (j % BITS)) & 1u);
}

void set_bit(Word *row, int j)
{
    row[j / BITS] |= (Word)1 << (j % BITS);
}

// Warshall：T[i][k] = 1 時，第 k 列整列 OR 進第 i 列
void warshall(Word *t, int n, int w)
{
    for (int k = 0; k < n; k++) {
        for (int i = 0; i < n; i++) {
            if (get_bit(&t[i * w], k)) {
                for (int x = 0; x < w; x++) {
                    t[i * w + x] |= t[k * w + x];
                }
            }
        }
    }
}

int main(void)
{
    int n, m;
    while (scanf("%d %d", &n, &m) == 2) {
        int w = (n + BITS - 1) / BITS;
        Word *t = calloc((size_t)n * w, sizeof *t);
        if (t == NULL) {
            return 1;
        }
        for (int e = 0; e < m; e++) {
            int u, v;
            if (scanf("%d %d", &u, &v) != 2) {
                return 1;
            }
            set_bit(&t[(u - 1) * w], v - 1);
        }
        warshall(t, n, w);
        long long pairs = 0;
        for (int i = 0; i < n; i++) {
            printf("T%d: ", i + 1);
            for (int j = 0; j < n; j++) {
                int b = get_bit(&t[i * w], j);
                printf("%d", b);
                if (b && i != j) {
                    pairs++;
                }
            }
            printf("\n");
        }
        printf("reachable pairs (i != j): %lld\n", pairs);
        printf("on a cycle:");
        int any = 0;
        for (int i = 0; i < n; i++) {
            // T[i][i] = 1：i 走得回自己
            if (get_bit(&t[i * w], i)) {
                printf(" %d", i + 1);
                any = 1;
            }
        }
        printf("%s\n", any ? "" : " none");
        free(t);
    }
    return 0;
}
\end{example}
\iosample{5 6 \\ 1 2 2 3 3 1 3 4 4 5 5 4 \\ 3 2 \\ 1 2 2 3}{T1: 11111 \\ T2: 11111 \\ T3: 11111 \\ T4: 00011 \\ T5: 00011 \\ reachable pairs (i != j): 14 \\ on a cycle: 1 2 3 4 5 \\ T1: 011 \\ T2: 001 \\ T3: 000 \\ reachable pairs (i != j): 3 \\ on a cycle: none}

第一張圖是 $C$，結果和 $T^{(5)}$ 相同。 第二張圖是一條鏈 $1 \to 2 \to 3$，沒有環，對角線都
是 $0$。 第~23~章實作練習裡的「可達點對個數」，用這支程式數 \texttt{reachable pairs} 就是一
種做法。

\section{Johnson 演算法}
\label{sec:spt-johnson}

Floyd--Warshall 不管圖的疏密都是 $\Theta(n^3)$。 稀疏圖上，對每個頂點跑一次 Dijkstra 只要
$O(nm \lg n)$，$m = O(n)$ 時是 $O(n^2 \lg n)$，快得多；可惜有負邊時 Dijkstra 不能用。
\textbf{Johnson 演算法}的想法是：先把權重改成非負的，而且改完之後最短路徑還是同一條，再跑
$n$ 次 Dijkstra。 這一節只講觀念和一個小例子。

\subsection{改權重}

第~\ref{sec:spt-dijkstra-neg}~節說過，每條邊加同一個數不行。 正確的改法是替每個頂點選一個
數 $h(v)$，令
\[
\hat{w}(u, v) = w(u, v) + h(u) - h(v).
\]

\begin{lemma}
\label{lem:spt-reweight}
對任何一條路徑 $p = \langle v_0, v_1, \ldots, v_k \rangle$，$\hat{w}(p) = w(p) + h(v_0) -
h(v_k)$。 所以 $p$ 是 $w$ 之下從 $v_0$ 到 $v_k$ 的最短路徑，若且唯若它是 $\hat{w}$ 之下的最短
路徑；$G$ 在 $w$ 之下有負環，若且唯若在 $\hat{w}$ 之下有負環。
\end{lemma}

\begin{proof}
$\hat{w}(p) = \sum_{i=1}^{k} \bigl( w(v_{i-1}, v_i) + h(v_{i-1}) - h(v_i) \bigr)$，中間的 $h$
一加一減全部消掉，只剩 $h(v_0) - h(v_k)$。 從 $v_0$ 到 $v_k$ 的每一條路徑都加上同一個數
$h(v_0) - h(v_k)$，大小順序不變。 環的起點等於終點，$\hat{w}(c) = w(c)$。
\end{proof}

加常數失敗，是因為每條路徑加的量和它的邊數有關；這裡加的量只和兩個端點有關，所以不會改變
誰比較短。 剩下的問題是怎麼選 $h$ 讓每個 $\hat{w}(u, v) \ge 0$，也就是 $h(v) \le h(u) +
w(u, v)$。 這正是三角不等式的形狀：加一個新頂點 $q$，從 $q$ 連一條權重 $0$ 的邊到每個頂點，
得到 $G'$；令 $h(v) = \delta(q, v)$。 由引理~\ref{lem:spt-triangle}，
\[
h(v) = \delta(q, v) \le \delta(q, u) + w(u, v) = h(u) + w(u, v).
\]
$q$ 走得到每個頂點，所以 $h(v)$ 有限（而且 $\le 0$）；$q$ 沒有入邊，$G'$ 的環都是 $G$ 原本的環。
用 Bellman--Ford 從 $q$ 出發算 $h$，順便檢查 $G$ 有沒有負環。

\begin{pseudo}[label=alg:spt-johnson]{Johnson 演算法}
\Function{Johnson}{$G, w$}
    \State $G' \gets G$ 加上新頂點 $q$，以及權重 $0$ 的邊 $(q, v)$，$v \in V$
    \If{\Call{Bellman-Ford}{$G', w, q$} $=$ \FALSE}
        \State \Return 「圖有負環」
    \EndIf
    \For{每個頂點 $v \in V$}
        \State $h(v) \gets d[v]$ \Comment{也就是 $\delta(q, v)$}
    \EndFor
    \For{每條邊 $(u, v) \in E$}
        \State $\hat{w}(u, v) \gets w(u, v) + h(u) - h(v)$
    \EndFor
    \For{每個頂點 $u \in V$}
        \State 用 \Call{Dijkstra}{$G, \hat{w}, u$} 求出 $\hat{d}[v] = \hat{\delta}(u, v)$
        \For{每個頂點 $v \in V$}
            \State $D[u, v] \gets \hat{d}[v] - h(u) + h(v)$ \Comment{引理~\ref{lem:spt-reweight}}
        \EndFor
    \EndFor
    \State \Return $D$
\EndFunction
\end{pseudo}

\textbf{時間。} Bellman--Ford 在 $n + 1$ 個頂點、$n + m$ 條邊的 $G'$ 上是 $O(nm)$（假設每個頂點
都至少有一條邊，$m \ge n/2$），改權重 $O(m)$，$n$ 次 Dijkstra 用二元堆積共 $O(nm \lg n)$，用費
氏堆積共 $O(n^2 \lg n + nm)$。 稀疏圖上明顯比 $\Theta(n^3)$ 好；稠密圖 $m = \Theta(n^2)$ 時用
費氏堆積也只是 $\Theta(n^3)$，和 Floyd--Warshall 打平，而 Floyd--Warshall 的程式短得多、常數
小，所以稠密圖通常用 Floyd--Warshall。

\subsection{一個例子}

圖~\ref{fig:spt-johnson}~(a) 的 $J$ 有兩條負邊。 加上 $q$ 之後跑 Bellman--Ford，$h(1) = 0$，
$h(2) = -2$（$q \to 1 \to 2$），$h(3) = -1$（$q \to 1 \to 2 \to 3$），$h(4) = -4$（$q \to 1 \to 2
\to 3 \to 4$）。 改完的權重在 (b)，全部 $\ge 0$。 例如 $\hat{w}(4, 2) = 5 + (-4) - (-2) = 3$，
$\hat{w}(3, 4) = -3 + (-1) - (-4) = 0$。

\begin{center}
\begin{tikzpicture}[x=1cm, y=1cm]
\begin{scope}
\node[spt-v] (j1) at (0, 0) {1};
\node[spt-v] (j2) at (2.2, 1.2) {2};
\node[spt-v] (j3) at (2.2, -1.2) {3};
\node[spt-v] (j4) at (4.4, 0) {4};
\draw[spt-neg] (j1) -- node[spt-wn, above left] {$-2$} (j2);
\draw[spt-a] (j1) -- node[spt-w, below left] {$4$} (j3);
\draw[spt-a] (j2) -- node[spt-w, left] {$1$} (j3);
\draw[spt-neg] (j3) -- node[spt-wn, below right] {$-3$} (j4);
\draw[spt-a] (j2) to[bend left=15] node[spt-w, above right] {$3$} (j4);
\draw[spt-a] (j4) to[bend left=15] node[spt-w, below left, pos=0.45] {$5$} (j2);
\node[font=\small] at (2.2, -2.3) {(a) 原本的權重 $w$};
\end{scope}
\begin{scope}[xshift=7.2cm]
\node[spt-v] (k1) at (0, 0) {1};
\node[spt-v] (k2) at (2.2, 1.2) {2};
\node[spt-v] (k3) at (2.2, -1.2) {3};
\node[spt-v] (k4) at (4.4, 0) {4};
\draw[spt-a] (k1) -- node[spt-w, above left] {$0$} (k2);
\draw[spt-a] (k1) -- node[spt-w, below left] {$5$} (k3);
\draw[spt-a] (k2) -- node[spt-w, left] {$0$} (k3);
\draw[spt-a] (k3) -- node[spt-w, below right] {$0$} (k4);
\draw[spt-a] (k2) to[bend left=15] node[spt-w, above right] {$5$} (k4);
\draw[spt-a] (k4) to[bend left=15] node[spt-w, below left, pos=0.45] {$3$} (k2);
\node[spt-h, left=2pt of k1] {$h = 0$};
\node[spt-h, above=2pt of k2] {$h = -2$};
\node[spt-h, below=2pt of k3] {$h = -1$};
\node[spt-h, right=2pt of k4] {$h = -4$};
\node[font=\small] at (2.2, -2.3) {(b) 改過的權重 $\hat{w}$，橘字是 $h$};
\end{scope}
\end{tikzpicture}\par\nopagebreak
\diagcaption{Johnson 演算法的改權重。 最短路徑上的邊改完之後權重常常變成 $0$。}
\label{fig:spt-johnson}
\end{center}

在 $\hat{w}$ 上從 $1$ 跑 Dijkstra，$\hat{d}[2] = \hat{d}[3] = \hat{d}[4] = 0$（$1 \to 2 \to 3 \to 4$
三條邊都是 $0$），換回來 $\delta(1, 4) = 0 - h(1) + h(4) = -4$，正是 $-2 + 1 - 3$。 從 $4$ 跑，
$\hat{d}[3] = 3 + 0 = 3$（$4 \to 2 \to 3$），$\delta(4, 3) = 3 - (-4) + (-1) = 6 = 5 + 1$。 注意
$\hat{w}$ 上最短路徑的邊很多是 $0$：$h(v) = \delta(q, v)$ 時，由 $q$ 出發的最短路徑樹上每條
邊都滿足三角不等式的等號，改完剛好是 $0$。

\subsection{本章演算法整理}

\begin{center}
\small
\begin{tabular}{>{\raggedright\arraybackslash}p{0.17\linewidth}>{\raggedright\arraybackslash}p{0.10\linewidth}>{\raggedright\arraybackslash}p{0.14\linewidth}>{\raggedright\arraybackslash}p{0.27\linewidth}>{\raggedright\arraybackslash}p{0.18\linewidth}}
\toprule
演算法 & 問題 & 權重 & 時間 & 負環 \\
\midrule
BFS（第~21~章） & 單一起點 & 全部相同 & $\Theta(n + m)$ & 不會有 \\
Dijkstra & 單一起點 & $\ge 0$ & 陣列 $\Theta(n^2)$；二元堆積 $O((n + m) \lg n)$；費氏堆積 $O(m + n \lg n)$ & 不會有 \\
Bellman--Ford & 單一起點 & 任意 & $\Theta(nm)$ & 偵測從 $s$ 走得到的 \\
DAG 最短路徑 & 單一起點 & 任意，無環 & $\Theta(n + m)$ & 不會有 \\
反覆平方 & 所有點對 & 任意 & $\Theta(n^3 \lg n)$ & 對角線變負 \\
Floyd--Warshall & 所有點對 & 任意 & $\Theta(n^3)$ & $d_{ii} < 0$ \\
Johnson & 所有點對 & 任意 & $O(nm \lg n)$ & 由 Bellman--Ford 偵測 \\
Warshall & 可達性 & 不看權重 & $\Theta(n^3)$，用位元約 $n^3/64$ & 不適用 \\
\bottomrule
\end{tabular}
\end{center}

選擇的順序大致是：圖沒有環就用 DAG 的演算法；權重全部相同用 BFS；沒有負邊用 Dijkstra；有負
邊才用 Bellman--Ford。 所有點對時，$n$ 不大（幾百）或圖稠密用 Floyd--Warshall，又大又稀疏用
Johnson（沒有負邊時直接跑 $n$ 次 Dijkstra）。

\section{本章常見錯誤}
\label{sec:spt-pitfalls}

\begin{center}
\begin{tabular}{>{\raggedright\arraybackslash}p{0.26\linewidth}>{\raggedright\arraybackslash}p{0.48\linewidth}>{\raggedright\arraybackslash}p{0.16\linewidth}}
\toprule
錯誤 & 例子 & 見 \\
\midrule
有負邊還用 Dijkstra & $N_1$ 只有一條負邊，算出 $d[4] = 5$，正確是 $3$ &
第~\ref{sec:spt-dijkstra-neg}~節 \\
用大數代表 $\infty$ 卻不先檢查 & \texttt{INF + (-3) < INF}，從一個還沒走到的頂點去更新別人 &
第~\ref{sec:spt-relax}~節 \\
相等也更新 & $G_1$ 第 2 輪 $3 \to 4$ 得到相等的 $6$，更新的話 $\pi[4]$ 變成 $3$，和題目要的
樹不同 & 第~\ref{sec:spt-dijkstra-trace}~節 \\
平手規則、邊的順序沒照題目 & Dijkstra 平手時取了別的頂點；Bellman--Ford 沒照指定的邊順序，
每一輪的表都不一樣 & 第~\ref{sec:spt-dijkstra-trace}、\ref{sec:spt-bf-trace}~節 \\
就地更新和 $\mathit{dist}^{k}$ 混用 & 同一張表前幾欄用了同一輪剛改的值，後幾欄又沒用 &
第~\ref{sec:spt-bf-trace}~節 \\
Bellman--Ford 沒做檢查 & 有負環時 $n - 1$ 輪照樣做完，把錯的 $d$ 當答案 &
第~\ref{sec:spt-bf}~節 \\
以為 \TRUE 代表整張圖沒負環 & 從 $s$ 走不到的負環偵測不到；要偵測整張圖，所有 $d$ 從 $0$
開始 & 第~\ref{sec:spt-bf}~節 \\
每條邊加同一個數 & 邊數多的路徑被加得比較多，最短路徑跟著變 &
第~\ref{sec:spt-johnson}~節 \\
無向圖有負邊 & 一條負的無向邊本身就是負環，本章的演算法都不能直接用 &
第~\ref{sec:spt-relax}~節 \\
最長路徑的初值寫 $0$ & DAG 上從 $s$ 走不到的頂點應該是 $-\infty$，寫 $0$ 會把它們當成起點 &
第~\ref{sec:spt-dag}~節 \\
Floyd--Warshall 的 $k$ 不在最外層 & 讀到還沒算好的 $D[i, k]$ 或 $D[k, j]$，答案偏大 &
第~\ref{sec:spt-fw}~節 \\
$\Pi$ 更新寫成 $\Pi[i, j] \gets k$ & 把前驅矩陣和中繼點矩陣混在一起，還原出錯的路徑 &
第~\ref{sec:spt-fw}~節 \\
$W$ 的對角線沒有設 $0$ & 對角線用 $\infty$ 開始，算出的 $d_{ii}$ 是經過 $i$ 的最短環，不是 $0$ &
第~\ref{sec:spt-fw}~節 \\
遞移閉包的對角線 & $A^{+}$ 的 $t_{ii}$ 要看 $i$ 在不在環上，照抄 $1$ 或 $0$ 都可能錯 &
第~\ref{sec:spt-fw}~節 \\
\bottomrule
\end{tabular}
\end{center}

\section{觀念題}
\label{sec:spt-exercises}

題型和研究所考試、期中期末考的筆試相近，有選擇、計算、追蹤、證明和設計。 沒有特別說明時，
照本書的約定：Dijkstra 平手取編號小的、相等不更新；鄰接串列的鄰居由小到大；Bellman--Ford 照
題目給的邊順序。 解答在本章最後。

\begin{enumerate}
\item \label{spt-q-tf}（選擇，可能不只一個答案）下列哪些敘述正確？
\begin{enumerate}
\item 每條邊的權重都是相異的正整數時，從 $s$ 到 $v$ 的最短路徑是唯一的。
\item 每條邊的權重都乘上同一個正數 $c$，任兩點之間的最短路徑不變。
\item 每條邊的權重都加上同一個正數 $c$，任兩點之間的最短路徑不變。
\item 無向圖裡有一條權重為負的邊；把每條邊換成兩個方向的有向邊之後，圖一定有負環。
\item \proc{Bellman-Ford} 回傳 \TRUE，表示整張圖沒有負環。
\item 有向無環圖即使有負邊，\proc{DAG-Shortest-Paths} 也能求出正確的最短距離。
\end{enumerate}

\item \label{spt-q-dij}（追蹤）無向圖 $H$ 有 $6$ 個頂點，邊和權重是
\begin{gather*}
\{1,2\}\!:4,\ \{1,3\}\!:1,\ \{2,3\}\!:2,\ \{2,4\}\!:5,\ \{3,4\}\!:8, \\
\{3,5\}\!:10,\ \{4,5\}\!:2,\ \{4,6\}\!:5,\ \{5,6\}\!:3.
\end{gather*}
\begin{enumerate}
\item 從 $s = 1$ 執行 Dijkstra 演算法，照本章的格式寫出每一輪取出的頂點和 $d[1..6]$。
\item 寫出最後的 $\pi$，畫出最短路徑樹，寫出 $1$ 到 $6$ 的路徑。
\item 從 $1$ 到 $6$ 的最短路徑有幾條？ $\pi$ 記下的是哪一條，為什麼？ 如果鬆弛時「相等也
更新」，$\pi$ 哪裡會不同？
\end{enumerate}

\item \label{spt-q-dk}（計算）對圖~\ref{fig:spt-g1} 的 $G_1$ 從 $1$ 執行 Dijkstra 演算法。
\begin{enumerate}
\item \proc{Decrease-Key} 一共被呼叫幾次？ 和 $m$ 比較。
\item 用陣列版，每次 \proc{Extract-Min} 要看過所有還在 $Q$ 裡的頂點，七次一共看過幾個頂點？
\item 一般來說，$m = \Theta(n^{3/2})$ 的圖，陣列版和二元堆積版哪一個的時間界比較好？
\end{enumerate}

\item \label{spt-q-const}（證明與反例）
\begin{enumerate}
\item 證明：每條邊的權重都乘上同一個正數 $c$，任兩點之間的最短路徑不變，最短距離變成 $c$ 倍。
\item 有人想讓 Dijkstra 演算法處理負邊：找出最小的權重 $w_{\min} < 0$，每條邊都加上
$-w_{\min}$。 舉一個反例說明這樣算出的最短路徑可能不對。
\item 說明 Johnson 演算法的 $\hat{w}$ 為什麼沒有 (b) 的問題。
\end{enumerate}

\item \label{spt-q-bf}（追蹤）有向圖有 $5$ 個頂點，起點 $s = 1$，邊依下面的順序掃描：
\[
(3, 5)\!:-2,\ (4, 3)\!:-3,\ (2, 3)\!:4,\ (1, 2)\!:6,\ (2, 5)\!:5,\ (1, 4)\!:7,\ (4, 5)\!:9,\ (5, 1)\!:2.
\]
\begin{enumerate}
\item 執行 \proc{Bellman-Ford}，寫出每一輪結束時的 $d[1..5]$ 和最後的 $\pi$。 第幾輪之後 $d$
就不再改變？
\item 重新排列這 $8$ 條邊的順序，讓第 $1$ 輪結束時 $d$ 就全部正確。
\item 證明：只要從 $s$ 走得到的地方沒有負環，一定存在一種邊的順序，讓 Bellman--Ford 第 $1$
輪結束時 $d$ 就全部等於 $\delta$。
\end{enumerate}

\item \label{spt-q-neg}（計算）有向圖有 $7$ 個頂點，起點 $s = 1$，邊依序是
\[
(1, 2)\!:2,\ (2, 3)\!:1,\ (3, 4)\!:-4,\ (4, 2)\!:1,\ (4, 5)\!:3,\ (1, 6)\!:5,\ (6, 5)\!:1,\ (7, 1)\!:2.
\]
\begin{enumerate}
\item 寫出 \proc{Bellman-Ford} 前兩輪結束時的 $d[1..7]$，並說明之後每一輪 $d$ 怎麼變、第 $6$
輪結束時的 $d$ 是多少、回傳什麼。
\item 照第~\ref{sec:spt-negcycle}~節的方法，從 $\pi$ 找出負環。
\item 寫出每個頂點的 $\delta(1, v)$（有限值、$-\infty$ 或 $\infty$）。
\end{enumerate}

\item \label{spt-q-dag}（追蹤）有向無環圖有 $7$ 個頂點，邊是
\begin{gather*}
1 \to 3\!:2,\ 2 \to 3\!:-1,\ 2 \to 6\!:1,\ 3 \to 4\!:3,\ 3 \to 5\!:6, \\
4 \to 5\!:-2,\ 4 \to 6\!:4,\ 5 \to 6\!:1,\ 5 \to 7\!:5,\ 6 \to 7\!:-3.
\end{gather*}
\begin{enumerate}
\item 用 Kahn 演算法（佇列，第~22~章的約定）排出拓樸順序。
\item 從 $s = 3$ 執行 \proc{DAG-Shortest-Paths}，寫出每處理一個頂點之後的 $d$，以及最後的
$\pi$。
\item 求從 $3$ 出發到每個頂點的最長路徑長度和一條最長路徑。
\end{enumerate}

\item \label{spt-q-count}（設計）$G$ 每條邊的權重都是正的。 設計一個 $O(m \lg n)$ 的演算法，求
出從 $s$ 到每個頂點的\textbf{最短路徑條數}，說明為什麼對，並算出 $G_1$ 從 $1$ 出發的結果。
為什麼要求權重是正的，而不只是非負？ （第~21~章觀念題第~\ref{q:gtr-count-paths}~題是權重
都是 $1$ 的版本。）

\item \label{spt-q-fw}（追蹤）有向圖有 $4$ 個頂點，邊是
\[
1 \to 2\!:5,\ 1 \to 3\!:2,\ 2 \to 1\!:-1,\ 2 \to 4\!:4,\ 3 \to 2\!:1,\ 3 \to 4\!:7,\ 4 \to 1\!:3,\
4 \to 3\!:-2.
\]
\begin{enumerate}
\item 寫出 $D^{(0)}$ 到 $D^{(4)}$ 和 $\Pi^{(0)}$ 到 $\Pi^{(4)}$，標出每一層改過的格子。
\item 用 $\Pi^{(4)}$ 還原 $4$ 到 $2$、$1$ 到 $4$ 的最短路徑，驗算權重。
\item 改用 \proc{Extend-Paths} 和反覆平方，這張圖要做幾次乘法？ $n = 100$ 的圖呢？
\end{enumerate}

\item \label{spt-q-order}（反例）有人把 Floyd--Warshall 的三層迴圈寫成 $i$ 最外層、$j$ 中間、$k$
最內層，其他不變。
\begin{enumerate}
\item 找一張最多 $4$ 個頂點、權重都是正數的圖，讓這個版本算錯，指出錯的是哪一格。
\item 改成 $i$ 最外層、$k$ 中間、$j$ 最內層呢？ 用同一張圖說明。
\end{enumerate}

\item \label{spt-q-tc}（計算）有向圖有 $5$ 個頂點，邊是 $1 \to 2$、$2 \to 4$、$4 \to 1$、$3 \to 2$、
$3 \to 5$、$5 \to 3$。
\begin{enumerate}
\item 用 Warshall 演算法寫出 $T^{(1)}$ 到 $T^{(5)}$。
\item 寫出 $A^{+}$ 和 $A^{*}$。
\item 有幾個有序點對 $(i, j)$，$i \ne j$，從 $i$ 走得到 $j$？ 哪些頂點在環上？ 強連通元件有哪
些？
\end{enumerate}

\item \label{spt-q-johnson}（計算）對圖~\ref{fig:spt-fw} 的 $F$ 做 Johnson 演算法。
\begin{enumerate}
\item 求出 $h(1..4)$。
\item 寫出每條邊的 $\hat{w}$，確認都 $\ge 0$。
\item 在 $\hat{w}$ 上從 $3$ 執行 Dijkstra，換回原本的權重，和 $D^{(4)}$ 的第 $3$ 列比較。
\end{enumerate}

\item \label{spt-q-mono}（證明與反例）證明：所有權重都 $\ge 0$ 時，Dijkstra 演算法依序取出的頂點，
取出當時的 $d$ 值不會變小。 再舉一個有負邊的小例子，說明這件事不一定成立。

\item \label{spt-q-choose}（配對）下列每個情況，選本章最適合的演算法，寫出時間。
\begin{enumerate}
\item $n = 10^5$、$m = 3 \times 10^5$，權重是 $1$ 到 $100$ 的整數，求一個起點到所有頂點的距離。
\item $n = 300$，有負邊、沒有負環，要回答任兩點之間的最短距離。
\item 工程的先後關係圖（沒有環），每項工作有一個可能是負的「成本」，求從開工事件到各事件的
最小成本。
\item $n = 2000$、$m = 5000$，有負邊，要知道從某個起點出發會不會遇到負環。
\item 每條邊的權重都相同，求一個起點到所有頂點的距離。
\item $n = 5000$、$m = 2 \times 10^4$，有負邊，要任兩點之間的最短距離。
\end{enumerate}

\item \label{spt-q-via}（設計）$G$ 是有向圖，權重都 $\ge 0$，給定 $s$、$t$。
\begin{enumerate}
\item 用 $O(m \lg n)$ 時間求出\textbf{每個頂點到 $t$} 的最短距離。
\item 求從 $s$ 出發、\textbf{一定要經過}頂點 $x$、到達 $t$ 的最短路徑長度。
\item 用 $O(m \lg n)$ 時間找出所有「至少在一條 $s$ 到 $t$ 的最短路徑上」的邊，並證明你的判斷
條件。
\end{enumerate}
\end{enumerate}

\section{實作練習}
\label{sec:spt-practice}

下面是 Aki-Judge 上和本章相關的題目。

\begin{center}
\begin{tabular}{llc}
\toprule
小節 & 題目 & 難度 \\
\midrule
Dijkstra & \akiref{shortest-distance-in-an-undirected-weighted-graph}{無向加權圖的最短距離} & 基本 \\
 & \akiref{lexicographically-smallest-shortest-path}{最短路徑的字典序最小還原} & 進階 \\
 & \akiref{minimum-cost-across-a-grid}{方格地圖上的最小通行成本} & 基本 \\
Bellman--Ford & \akiref{shortest-distance-in-a-directed-graph-with-negative-edges}{含負權重有向圖的最短距離} & 基本 \\
 & \akiref{does-a-directed-graph-contain-a-negative-cycle}{判斷有向圖是否有負環} & 基本 \\
DAG 上的 DP & \akiref{relay-groups-and-delivery-paths}{中繼群組的送貨路徑} & 進階 \\
Floyd--Warshall & \akiref{all-pairs-shortest-distance-in-a-directed-graph}{有向圖上任兩點的最短距離} & 基本 \\
 & \akiref{patrol-route-with-deadlines}{限時巡查路線} & 進階 \\
\bottomrule
\end{tabular}
\end{center}

幾題的提示：

\begin{itemize}
\item 「無向加權圖的最短距離」的 $n$ 到 $5 \times 10^4$，陣列版 $\Theta(n^2)$ 太慢，要用堆積；範
例~\ref{ex:spt-dijkstra} 的寫法可以直接用。 平行邊不必特別處理，兩條都放進鄰接串列，鬆弛時自
然取小的。 無向邊要放兩個方向。 距離用 \texttt{long long}。
\item 「最短路徑的字典序最小還原」不能只靠 $\pi$：$\pi$ 是每個頂點記「前一個」，平手規則決定的
是路徑的尾巴，不是開頭。 先從 $t$ 跑一次 Dijkstra（無向圖不必反向），得到每個頂點到 $t$ 的距離
$\delta(v, t)$；再從 $s$ 出發，每一步走到編號最小、而且 $w(u, v) + \delta(v, t) = \delta(u, t)$
的鄰居 $v$。 權重是正的，每走一步剩下的距離嚴格變小，不會繞圈。
\item 「方格地圖上的最小通行成本」的成本在格子上而不是邊上：走進格子 $v$ 的邊，權重就是 $v$
的成本，起點格的成本另外加。 鄰居用上下左右的位移算，不必建鄰接串列。
\item 「含負權重有向圖的最短距離」是範例~\ref{ex:spt-bf} 的前半；一定要先檢查 $d[u]$ 是不是
$\infty$ 再鬆弛。 「判斷有向圖是否有負環」不限起點，所有 $d$ 從 $0$ 開始（第~\ref{sec:spt-bf}~
節），做完 $n - 1$ 輪再檢查一次。
\item 「中繼群組的送貨路徑」是第~\ref{sec:spt-dag}~節「依拓樸順序累加路徑條數」的應用，群組
的編號本身就是拓樸順序。 同一個群組裡的站點完全對稱，以群組為單位算，每條規則 $u \to v$ 讓
到 $v$ 的條數加上「到 $u$ 的條數 $\times$ $u$ 的站點數」。 不在任何路徑上的站點，是從起點走不
到或走不到終點的：往前、往後各做一次可達性。
\item 「有向圖上任兩點的最短距離」的 $n \le 400$，$\Theta(n^3)$ 可以；重複的邊讀入時取最小，
$k$ 放最外層。
\item 「限時巡查路線」要列出所有經過每個檢查站恰好一次的走法，是第~20~章的回溯，不是最短路
徑問題。 但最短路徑可以拿來剪枝：先用 Floyd--Warshall 算出任兩站之間的最短距離 $\delta$，走到
$u$、時間是 $t$ 時，任何還沒去的站 $v$ 最早也要到 $t + \delta(u, v)$ 才能抵達（就算中間經過其他
還沒去的站，也不會比最短路徑快），超過期限就整支剪掉。
\end{itemize}

\section{讀完你應該會}
\label{sec:spt-checklist}

\begin{itemize}
\item 說出最短距離、負環、$-\infty$ 的定義，用剪下貼上證明最短路徑的子路徑還是最短路徑，並證
明鬆弛的上界、收斂、路徑鬆弛三個性質。
\item 照指定的平手規則手算 Dijkstra 的逐輪表，畫出最短路徑樹；證明它在權重非負時正確，並給
出有負邊時的反例。
\item 照指定的邊順序手算 Bellman--Ford 的逐輪表，分辨就地更新和 $\mathit{dist}^{k}$ 兩種寫法；
證明它能偵測從 $s$ 走得到的負環，並從 $\pi$ 找出負環。
\item 依拓樸順序求有向無環圖的最短路徑和最長路徑，說明它和第~22~章關鍵路徑的關係。
\item 寫出 Floyd--Warshall 的 $D^{(k)}$、$\Pi^{(k)}$ 序列並還原路徑，用 $d_{ii} < 0$ 判斷負環；
用 Warshall 演算法求 $A^{+}$、$A^{*}$。
\item 說明 Johnson 演算法為什麼改權重之後最短路徑不變、權重一定非負，並依圖的大小、疏密和
權重的正負選擇演算法。
\end{itemize}

\textbf{延伸閱讀}（書末會整理完整的參考資料）

\begin{itemize}
\item Thomas H. Cormen 等，\textit{Introduction to Algorithms}，第 4 版，第 22、23 章：本章的主要
參考。 22.4 節的差分約束（把一組 $x_j - x_i \le b$ 的不等式化成最短路徑）本章沒有講，是
Bellman--Ford 的一個應用。
\item Ellis Horowitz、Sartaj Sahni、Susan Anderson-Freed，\textit{Fundamentals of Data Structures
in C}，第 2 版，第 6 章：$\mathit{dist}^{k}$ 形式的 Bellman--Ford、$A^{+}$ 與 $A^{*}$ 等寫法，都是
這一系列資料結構課本常見的記號。
\item E. W. Dijkstra，``A Note on Two Problems in Connexion with Graphs''，\textit{Numerische
Mathematik} 1，1959：三頁的原始論文，同時給了最小生成樹（第~24~章）和最短路徑的做法。
\item R. Bellman，``On a Routing Problem''，\textit{Quarterly of Applied Mathematics} 16(1)，1958；
L. R. Ford Jr.，\textit{Network Flow Theory}，RAND Paper P-923，1956：Bellman--Ford 演算法名稱的
由來。
\item R. W. Floyd，``Algorithm 97: Shortest Path''，\textit{Communications of the ACM} 5(6)，1962；
S. Warshall，``A Theorem on Boolean Matrices''，\textit{Journal of the ACM} 9(1)，1962。
\item D. B. Johnson，``Efficient Algorithms for Shortest Paths in Sparse Networks''，\textit{Journal of
the ACM} 24(1)，1977：改權重的方法。
\item M. L. Fredman、R. E. Tarjan，``Fibonacci Heaps and Their Uses in Improved Network
Optimization Algorithms''，\textit{Journal of the ACM} 34(3)，1987：費氏堆積，以及 Dijkstra 演算法
因此得到的 $O(m + n \lg n)$。
\end{itemize}

\addcontentsline{toc}{section}{觀念題解答}
\section*{觀念題解答}

\begin{enumerate}
\item (b)、(d)、(f) 正確。 (a)：錯，$1 \to 2$ 權重 $1$、$2 \to 3$ 權重 $4$、$1 \to 3$ 權重 $5$，三個
權重相異，$1$ 到 $3$ 卻有兩條最短路徑，權重都是 $5$。 (b)：對，見第~\ref{spt-q-const}~題。
(c)：錯，見第~\ref{spt-q-const}~題 (b)，加正數也一樣，邊數多的路徑吃虧。 (d)：對，負邊 $\{u, v\}$
變成 $u \to v \to u$，權重是兩倍的負數。 (e)：錯，從 $s$ 走不到的負環偵測不到
（第~\ref{sec:spt-bf}~節）。 (f)：對，定理~\ref{thm:spt-dag}。

\item 鄰接串列：$\mathit{Adj}[1] = 2, 3$；$\mathit{Adj}[2] = 1, 3, 4$；$\mathit{Adj}[3] = 1, 2, 4, 5$；
$\mathit{Adj}[4] = 2, 3, 5, 6$；$\mathit{Adj}[5] = 3, 4, 6$；$\mathit{Adj}[6] = 4, 5$。
\begin{enumerate}
\item 每一輪取出的頂點和鬆弛之後的 $d$：
\begin{center}
\begin{tabular}{cc*{6}{c}}
\toprule
輪 & 取出 $u$ & $d[1]$ & $d[2]$ & $d[3]$ & $d[4]$ & $d[5]$ & $d[6]$ \\
\midrule
初始 & & $0$ & $\infty$ & $\infty$ & $\infty$ & $\infty$ & $\infty$ \\
1 & $1$ & \sptpick{0} & \sptnew{4} & \sptnew{1} & $\infty$ & $\infty$ & $\infty$ \\
2 & $3$ & \sptold{0} & \sptnew{3} & \sptpick{1} & \sptnew{9} & \sptnew{11} & $\infty$ \\
3 & $2$ & \sptold{0} & \sptpick{3} & \sptold{1} & \sptnew{8} & $11$ & $\infty$ \\
4 & $4$ & \sptold{0} & \sptold{3} & \sptold{1} & \sptpick{8} & \sptnew{10} & \sptnew{13} \\
5 & $5$ & \sptold{0} & \sptold{3} & \sptold{1} & \sptold{8} & \sptpick{10} & $13$ \\
6 & $6$ & \sptold{0} & \sptold{3} & \sptold{1} & \sptold{8} & \sptold{10} & \sptpick{13} \\
\bottomrule
\end{tabular}
\end{center}
\item $\pi[2..6] = 3, 1, 2, 4, 4$。 樹的邊是 $\{1,3\}$、$\{3,2\}$、$\{2,4\}$、$\{4,5\}$、$\{4,6\}$。
$1$ 到 $6$：$1, 3, 2, 4, 6$，權重 $1 + 2 + 5 + 5 = 13$。
\item 兩條：$1, 3, 2, 4, 6$ 和 $1, 3, 2, 4, 5, 6$，都是 $13$。 第 4 輪從 $4$ 鬆弛得到 $d[6] = 13$、
$\pi[6] = 4$；第 5 輪從 $5$ 得到 $10 + 3 = 13$，相等不更新，所以 $\pi$ 記下的是經過 $4 \to 6$ 那一
條。 相等也更新的話，$\pi[6] = 5$，其他不變（整個過程只有這一次相等）。
\end{enumerate}

\item (a) 每一輪成功鬆弛的次數是 $3, 2, 2, 1, 2, 1, 0$，共 $11$ 次，比 $m = 14$ 少：有三條邊鬆弛不
成功（$3 \to 4$ 相等、$4 \to 5$、$6 \to 2$）。 (b) $7 + 6 + \cdots + 1 = 28$ 個。 (c) 二元堆積
$O((n + m) \lg n) = O(n^{3/2} \lg n)$，陣列版 $\Theta(n^2)$；$n^{3/2} \lg n = o(n^2)$，二元堆積好。

\item (a) 任何一條路徑 $p$，新的權重是 $\sum c \cdot w(e) = c \cdot w(p)$。 $c > 0$，所以兩條路徑誰
比較輕不會改變，同一對端點之間的最短路徑還是那些，最短距離乘上 $c$（$\pm\infty$ 不變）。
(b) $1 \to 2$：$-2$，$2 \to 3$：$-2$，$1 \to 3$：$-3$。 原本 $1$ 到 $3$ 的最短路徑是 $1 \to 2 \to 3$，
權重 $-4$。 每條邊加 $3$ 之後變成 $1, 1, 0$，$1 \to 2 \to 3$ 是 $2$，直接走是 $0$，最短路徑變了。
$k$ 條邊的路徑被加了 $3k$，邊數不同的路徑加的量不同。 (c) $\hat{w}$ 加在一條路徑上的總量是
$h(v_0) - h(v_k)$（引理~\ref{lem:spt-reweight}），只和兩端有關，同一對端點之間的每條路徑都加
同一個數。

\item (a) 第 1 輪掃 $(3, 5)$、$(4, 3)$、$(2, 3)$ 時起點都還是 $\infty$；第 2 輪才用 $d[4] = 7$
得到 $d[3] = 4$，第 3 輪才用 $d[3]$ 得到 $d[5] = 2$，之後不再改變。 每一輪結束時的 $d$：
\begin{center}
\begin{tabular}{c*{5}{c}}
\toprule
輪 & $d[1]$ & $d[2]$ & $d[3]$ & $d[4]$ & $d[5]$ \\
\midrule
1 & $0$ & $6$ & $\infty$ & $7$ & $11$ \\
2 & $0$ & $6$ & $4$ & $7$ & $11$ \\
3 & $0$ & $6$ & $4$ & $7$ & $2$ \\
4 & $0$ & $6$ & $4$ & $7$ & $2$ \\
\bottomrule
\end{tabular}
\end{center}
最後的 $\pi[2..5] = 1, 4, 1, 3$。
(b) 例如 $(1,4), (4,3), (3,5), (1,2), (2,3), (2,5), (4,5), (5,1)$：第 1 輪依序得到 $d[4] = 7$、
$d[3] = 4$、$d[5] = 2$、$d[2] = 6$，之後的邊都鬆弛不動，已經全部正確。
(c) 依定理~\ref{thm:spt-bf} 和定理~\ref{thm:spt-tree}，存在一棵以 $s$ 為根的最短路徑樹 $T$（例
如先跑一次 Bellman--Ford 得到的 $\pi$）。 把 $T$ 的邊依頭的深度由淺到深排在最前面（同深度隨
意），其他的邊排在後面。 對每個 $s$ 走得到的 $v$，$T$ 裡從 $s$ 到 $v$ 的路徑是一條最短路徑，它
的邊深度遞增，在第 $1$ 輪裡照路徑的先後被鬆弛，由路徑鬆弛性質第 $1$ 輪結束時 $d[v] =
\delta(s, v)$；走不到的頂點本來就是 $\infty$。 照深度排只是一種保證「每條樹上路徑的邊由根往
下出現」的排法；(b) 的順序沒有照深度，但也滿足這一點。

\item (a) 第 1 輪：$d = (0, 0, 3, -1, 2, 5, \infty)$（$(4, 2)$ 用到這一輪剛算出的 $d[4] = -1$，把
$d[2]$ 從 $2$ 改成 $0$）。 第 2 輪：$d = (0, -2, 1, -3, 0, 5, \infty)$。 之後每一輪 $d[2]$、$d[3]$、
$d[4]$、$d[5]$ 都再少 $2$（繞一圈負環 $-2$），第 6 輪結束時 $d = (0, -10, -7, -11, -8, 5, \infty)$。
檢查時 $(2, 3)$：$-10 + 1 = -9 < -7$，回傳 \FALSE。
(b) $\pi[2] = 4$、$\pi[3] = 2$、$\pi[4] = 3$、$\pi[5] = 4$、$\pi[6] = 1$。 檢查那一輪照樣鬆弛，最後一
個變小的是 $5$（$(2,3)$、$(3,4)$、$(4,2)$、$(4,5)$ 依序成功）。 從 $5$ 往回走 $7$ 步：$4, 3, 2, 4,
3, 2, 4$，停在 $4$；繞一圈 $4, 3, 2, 4$，倒過來是負環 $2 \to 3 \to 4 \to 2$，權重 $1 - 4 + 1 = -2$。
(c) $\delta(1, 1) = 0$，$\delta(1, 6) = 5$；$2, 3, 4$ 在負環上，$5$ 從負環走得到，都是 $-\infty$；
$7$ 沒有邊進來，是 $\infty$。 注意 $d[6] = 5$ 是對的：$6$ 和負環在同一張圖上，但從負環走不到 $6$。

\item (a) 入度 $\mathit{in}[1..7] = 0, 0, 2, 1, 2, 3, 2$，佇列一開始是 $1, 2$。 取出 $1$：$\mathit{in}[3]
= 1$。 取出 $2$：$\mathit{in}[3] = 0$ 進佇列，$\mathit{in}[6] = 2$。 之後依序取出 $3, 4, 5, 6, 7$，每次
剛好放進下一個。 拓樸順序 $1, 2, 3, 4, 5, 6, 7$。
(b) 每處理一個頂點之後的 $d$：
\begin{center}
\begin{tabular}{c*{7}{c}}
\toprule
處理 $u$ & $d[1]$ & $d[2]$ & $d[3]$ & $d[4]$ & $d[5]$ & $d[6]$ & $d[7]$ \\
\midrule
$1$、$2$ & $\infty$ & $\infty$ & $0$ & $\infty$ & $\infty$ & $\infty$ & $\infty$ \\
$3$ & $\infty$ & $\infty$ & $0$ & \sptnew{3} & \sptnew{6} & $\infty$ & $\infty$ \\
$4$ & $\infty$ & $\infty$ & $0$ & $3$ & \sptnew{1} & \sptnew{7} & $\infty$ \\
$5$ & $\infty$ & $\infty$ & $0$ & $3$ & $1$ & \sptnew{2} & \sptnew{6} \\
$6$ & $\infty$ & $\infty$ & $0$ & $3$ & $1$ & $2$ & \sptnew{$-1$} \\
$7$ & $\infty$ & $\infty$ & $0$ & $3$ & $1$ & $2$ & $-1$ \\
\bottomrule
\end{tabular}
\end{center}
$1$、$2$ 排在 $3$ 前面，$3$ 走不到它們，處理時 $d = \infty$ 不鬆弛。 $\pi[4..7] = 3, 4, 5, 6$，到 $7$
的路徑 $3 \to 4 \to 5 \to 6 \to 7$，權重 $3 - 2 + 1 - 3 = -1$。
(c) 最長距離 $d[4..7] = 3, 6, 7, 11$：$3 \to 4$；$3 \to 5$（$6 > 3 - 2$）；到 $6$ 有兩條都是 $7$
（$3 \to 4 \to 6$ 和 $3 \to 5 \to 6$）；$3 \to 5 \to 7$ 是 $11$。

\item 先跑 Dijkstra 得到 $d = \delta(s, \cdot)$ 和取出的順序。 令 $c[s] = 1$，其他是 $0$；依取出的順序
處理每個頂點 $u$，對每條出邊 $(u, v)$，若 $d[u] + w(u, v) = d[v]$ 就 $c[v] \gets c[v] + c[u]$。 時間
是 Dijkstra 的 $O(m \lg n)$ 再加 $O(n + m)$。

為什麼對：到 $v \ne s$ 的每一條最短路徑，最後一條邊 $(u, v)$ 都滿足 $\delta(s, u) + w(u, v) = \delta(s,
v)$，前面那段是到 $u$ 的最短路徑（引理~\ref{lem:spt-subpath}）；反過來，到這種 $u$ 的任何一條最
短路徑接上 $(u, v)$，都是到 $v$ 的最短路徑。 所以 $c[v] = \sum c[u]$，和是對所有這種 $u$ 取。 權
重是正的，這種 $u$ 都有 $d[u] < d[v]$，比 $v$ 早取出，加到 $c[v]$ 時 $c[u]$ 已經算完。

$G_1$：$c[1..7] = 1, 1, 1, 2, 1, 1, 1$。 只有 $4$ 有兩條：$1 \to 4$ 和 $1 \to 3 \to 4$，都是 $6$。
權重可以是 $0$ 時，$(u, v)$ 權重 $0$ 會讓 $d[u] = d[v]$，平手規則可能先處理 $v$，$c[v]$ 就少算了
$c[u]$；如果還有權重 $0$ 的環，最短路徑（可以重複頂點）有無限多條。 沒有權重 $0$ 的環時，同一
個 $d$ 值的頂點改依權重 $0$ 的邊的拓樸順序處理就能修正。

\item (a) 粗體是這一層改過的格子，$\Pi$ 裡的「$-$」是 \NIL。
\begin{gather*}
D^{(0)} = \begin{pmatrix} 0 & 5 & 2 & \infty \\ -1 & 0 & \infty & 4 \\ \infty & 1 & 0 & 7 \\ 3 & \infty & -2 & 0 \end{pmatrix},\quad
\Pi^{(0)} = \begin{pmatrix} - & 1 & 1 & - \\ 2 & - & - & 2 \\ - & 3 & - & 3 \\ 4 & - & 4 & - \end{pmatrix}; \\
D^{(1)} = \begin{pmatrix} 0 & 5 & 2 & \infty \\ -1 & 0 & \mathbf{1} & 4 \\ \infty & 1 & 0 & 7 \\ 3 & \mathbf{8} & -2 & 0 \end{pmatrix},\quad
\Pi^{(1)} = \begin{pmatrix} - & 1 & 1 & - \\ 2 & - & \mathbf{1} & 2 \\ - & 3 & - & 3 \\ 4 & \mathbf{1} & 4 & - \end{pmatrix}; \\
D^{(2)} = \begin{pmatrix} 0 & 5 & 2 & \mathbf{9} \\ -1 & 0 & 1 & 4 \\ \mathbf{0} & 1 & 0 & \mathbf{5} \\ 3 & 8 & -2 & 0 \end{pmatrix},\quad
\Pi^{(2)} = \begin{pmatrix} - & 1 & 1 & \mathbf{2} \\ 2 & - & 1 & 2 \\ \mathbf{2} & 3 & - & \mathbf{2} \\ 4 & 1 & 4 & - \end{pmatrix}; \\
D^{(3)} = \begin{pmatrix} 0 & \mathbf{3} & 2 & \mathbf{7} \\ -1 & 0 & 1 & 4 \\ 0 & 1 & 0 & 5 \\ \mathbf{-2} & \mathbf{-1} & -2 & 0 \end{pmatrix},\quad
\Pi^{(3)} = \begin{pmatrix} - & \mathbf{3} & 1 & 2 \\ 2 & - & 1 & 2 \\ 2 & 3 & - & 2 \\ \mathbf{2} & \mathbf{3} & 4 & - \end{pmatrix}.
\end{gather*}
$D^{(4)} = D^{(3)}$、$\Pi^{(4)} = \Pi^{(3)}$，第 $4$ 層沒有格子變小。 幾個要注意的地方：$d_{14}$ 在
第 $2$ 層變成 $9$（$1 \to 2 \to 4$），第 $3$ 層又變成 $7$（$1 \to 3 \to 2 \to 4$），$\pi_{14}$ 都是
$2$，因為兩條路徑的最後一條邊都是 $2 \to 4$；$\pi_{12}$ 在第 $3$ 層從 $1$ 變成 $3$。
(b) $4$ 到 $2$：$\pi_{42} = 3$、$\pi_{43} = 4$，路徑 $4 \to 3 \to 2$，權重 $-2 + 1 = -1 = d_{42}$。
$1$ 到 $4$：$\pi_{14} = 2$、$\pi_{12} = 3$、$\pi_{13} = 1$，路徑 $1 \to 3 \to 2 \to 4$，權重 $2 + 1 + 4 =
7 = d_{14}$。
(c) $n - 1 = 3$，算到 $L^{(4)}$，兩次乘法。 $n = 100$ 時 $n - 1 = 99$，$2^7 = 128 \ge 99$，七次。

\item (a) 三條邊 $1 \to 4$、$4 \to 3$、$3 \to 2$，權重都是 $1$。 錯的版本算 $D[1, 2]$ 時（$i = 1$、
$j = 2$），試 $k = 1..4$：$D[1, 3]$ 要到 $j = 3$ 才會算出來，現在還是 $\infty$；$D[1, 4] = 1$ 但
$D[4, 2]$ 要到 $i = 4$ 才會算出來，現在也是 $\infty$。 所以 $D[1, 2]$ 停在 $\infty$，正確答案是
$3$（$1 \to 4 \to 3 \to 2$）。
(b) 也錯。 $i = 1$ 時，$k = 3$ 那一輪 $D[1, 3]$ 還是 $\infty$（要等 $k = 4$ 才變成 $2$），用不上
$3 \to 2$；$k = 4$ 那一輪 $D[4, 2]$ 還沒算（第 $4$ 列要到 $i = 4$ 才處理），也是 $\infty$。 $D[1, 2]$
同樣停在 $\infty$。 $k$ 必須在最外層，才能保證第 $k$ 層用到的 $D[i, k]$、$D[k, j]$ 已經考慮過所有
$\{1, \ldots, k - 1\}$ 的中間點。

\item (a) 只寫每一層新變成 $1$ 的格子：$T^{(1)}$：$(4, 2)$。 $T^{(2)}$：$(1, 4)$、$(3, 4)$、$(4, 4)$。
$T^{(3)}$：$(5, 2)$、$(5, 4)$、$(5, 5)$。 $T^{(4)}$：$(1, 1)$、$(2, 1)$、$(2, 2)$、$(3, 1)$、$(5, 1)$。
$T^{(5)}$：$(3, 3)$。
(b)
\[
A^{+} = T^{(5)} = \begin{pmatrix} 1&1&0&1&0 \\ 1&1&0&1&0 \\ 1&1&1&1&1 \\ 1&1&0&1&0 \\ 1&1&1&1&1 \end{pmatrix},
\]
對角線已經全是 $1$，所以 $A^{*} = A^{+}$。
(c) 不看對角線，各列的 $1$ 有 $2, 2, 4, 2, 4$ 個，共 $14$ 個點對。 五個頂點都在環上
（$1 \to 2 \to 4 \to 1$、$3 \to 5 \to 3$）。 強連通元件 $\{1, 2, 4\}$、$\{3, 5\}$；從 $\{3, 5\}$ 走得到
$\{1, 2, 4\}$，反過來不行。

\item (a) 從 $q$ 跑 Bellman--Ford：$h(1) = 0$、$h(2) = 0$、$h(3) = -2$（$q \to 2 \to 3$）、$h(4) = 0$。
(b) $\hat{w}$：$1 \to 2$：$3$，$1 \to 4$：$8$，$2 \to 3$：$-2 + 0 + 2 = 0$，$2 \to 4$：$6$，$3 \to 1$：
$4 - 2 - 0 = 2$，$3 \to 4$：$2 - 2 - 0 = 0$，$4 \to 1$：$2$，$4 \to 2$：$1$，全部 $\ge 0$。
(c) 從 $3$：取出 $3$，$\hat{d}[1] = 2$、$\hat{d}[4] = 0$；取出 $4$，$\hat{d}[2] = 1$（$4 \to 1$ 得到 $2$，
相等不更新）；再取出 $2$、$1$，沒有變化。 $\hat{d}[1..4] = 2, 1, 0, 0$。 換回來 $D[3, v] = \hat{d}[v] -
h(3) + h(v) = \hat{d}[v] + 2 + h(v)$，得 $4, 3, 0, 2$，和 $D^{(4)}$ 的第 $3$ 列相同。

\item 對取出的次數歸納，證明「每次取出 $u$ 之後，$Q$ 裡每個頂點的 $d$ 都 $\ge d[u]$」。 取出 $u$ 的
那一刻，$u$ 是 $Q$ 裡最小的，其他頂點的 $d \ge d[u]$。 接下來鬆弛 $u$ 的出邊，被改的 $d[v]$ 變成
$d[u] + w(u, v) \ge d[u]$（權重 $\ge 0$），沒被改的本來就 $\ge d[u]$。 所以下一個取出的頂點 $u'$
有 $d[u'] \ge d[u]$，而 $d[u']$ 在取出之後不再改變（定理~\ref{thm:spt-dijkstra}），取出當時的值
就是這個。 反例：$1 \to 2$ 權重 $2$、$2 \to 3$ 權重 $-5$。 依序取出 $1$（$0$）、$2$（$2$）、$3$
（$-3$），$d$ 變小了。

\item (a) Dijkstra 加二元堆積，$O(m \lg n)$，約 $3 \times 10^5 \times 17$。 (b) Floyd--Warshall，
$\Theta(n^3) = 2.7 \times 10^7$。 (c) DAG 最短路徑，$\Theta(n + m)$。 (d) Bellman--Ford，$\Theta(nm)
= 10^7$。 (e) BFS，$\Theta(n + m)$。 (f) Johnson，$O(nm \lg n) \approx 5000 \times 2 \times 10^4 \times 13
\approx 1.3 \times 10^9$；Floyd--Warshall 要 $1.25 \times 10^{11}$。

\item (a) 把每條邊反向得到 $G^{\mathrm{T}}$，從 $t$ 跑 Dijkstra：$G$ 裡從 $v$ 到 $t$ 的路徑，正好是
$G^{\mathrm{T}}$ 裡從 $t$ 到 $v$ 的路徑倒過來，權重相同。 (b) 答案是 $\delta(s, x) + \delta(x, t)$：從
$s$ 跑一次 Dijkstra，從 $x$ 再跑一次（或在 $G^{\mathrm{T}}$ 上從 $t$ 跑）。 經過 $x$ 的路徑在 $x$ 切成兩
段，兩段各自最短時總和最小；接起來可能重複頂點，這在路徑的定義裡是允許的
（定義~\ref{def:gtr-path}）。 (c) 從 $s$ 在 $G$ 上、從 $t$ 在 $G^{\mathrm{T}}$ 上各跑一次 Dijkstra，得到
$d_s(v) = \delta(s, v)$、$d_t(v) = \delta(v, t)$。 設 $\delta(s, t) < \infty$。 邊 $(u, v)$ 在某一條 $s$
到 $t$ 的最短路徑上，若且唯若
\[
d_s(u) + w(u, v) + d_t(v) = \delta(s, t).
\]
必要：最短路徑 $p$ 經過 $(u, v)$，$p$ 到 $u$ 的那段權重 $\ge d_s(u)$、$v$ 之後那段 $\ge d_t(v)$，所以
$\delta(s, t) = w(p) \ge d_s(u) + w(u, v) + d_t(v)$；而左邊這三段接起來是一條 $s$ 到 $t$ 的路徑，權重
$\ge \delta(s, t)$，等號成立。 充分：把 $s$ 到 $u$ 的最短路徑、邊 $(u, v)$、$v$ 到 $t$ 的最短路徑接
起來，權重正好是 $\delta(s, t)$，是一條經過 $(u, v)$ 的最短路徑。 每條邊檢查一次，總時間
$O(m \lg n)$。 第~21~章觀念題第~\ref{q:gtr-on-path}~題是權重都是 $1$ 的版本。
\end{enumerate}

\chapter{網路流與二分匹配}
\label{ch:flw}

% 本章的圖共用的樣式（前綴 flw-）。 Edmonds--Karp 每一輪的剩餘網路、距離表和最後的
% 流，都由模擬程式照虛擬碼逐步執行後產生，再和本章 C 程式的輸出對過。
\tikzset{
    flw-v/.style={circle, draw=black!70, thick, minimum size=0.62cm, inner sep=0pt, font=\small, fill=white},
    flw-vS/.style={draw=coreStroke, fill=coreFill},
    flw-a/.style={-{Stealth[length=2.2mm]}, draw=black!65, thick},
    flw-e/.style={draw=black!55, thick},
    flw-rev/.style={densely dashed},
    flw-aug/.style={draw=examStroke, very thick},
    flw-cut/.style={draw=memStroke, very thick},
    flw-m/.style={draw=examStroke, line width=1.8pt},
    flw-c/.style={font=\scriptsize, fill=white, inner sep=1pt},
    flw-cA/.style={text=examStroke, font=\scriptsize\bfseries},
    flw-t/.style={font=\small},
}

一家賣生鮮蔬果的網路商店，每天晚上要把貨從郊區的集貨場運到市區的配送中心。 中間有五
個轉運站，每一段路線由不同的車隊負責，一個晚上最多只能載固定的板數。 貨可以在轉運站換
車、分開走不同的路線，但轉運站沒有冷藏庫，當晚進站多少板，就得當晚全部送出去。 物流經
理想知道一個晚上最多能送多少板到配送中心；如果不夠，應該先替哪一段路線加車。 另一邊，
系上的四位助教要分配到五堂實習課，每位助教只帶一堂，每堂課只要一位助教，而每位助教只有
某幾堂的時間可以。 最多能讓幾堂課有助教？ 兩個問題都可以先隨便排一種送法、一種分配，再
慢慢改。 往下讀之前，先想一想：排到一半發現卡住時，要怎麼把先前的決定改掉，又不會越改
越亂？

這一章處理\textbf{最大流}（maximum flow）問題：在每條邊有容量上限的有向圖上，從起點送
最多的量到終點。 主要的方法是 Ford--Fulkerson：一再找一條「還能多送一點」的路徑，沿著它
送。 關鍵是\textbf{剩餘網路}（也叫殘餘網路），它除了記錄每條邊還能多送多少，也允許把先前送出去的量退回
來，所以一開始選錯路不會讓演算法卡住。 Edmonds--Karp 演算法規定每次用 BFS 找邊數最少的
路徑，時間是 $O(nm^2)$，和容量的大小無關。 證明這一切正確的是\textbf{最大流最小割定
理}：能送的最大量，剛好等於要把起點和終點完全隔開、最少得切掉多少容量。 最後把二分圖
的最大匹配化成最大流，助教的分配就是一個例子。

這一章用到第~21~章的有向圖、BFS（虛擬碼~\ref{alg:gtr-bfs}、定理~\ref{thm:gtr-bfs}）和
二分圖（定義~\ref{def:gtr-bipartite}），名詞和符號都沿用那一章：$n = |V|$、$m = |E|$，
$u \leadsto v$ 表示有路徑，$\delta(u, v)$ 是邊數最少的路徑長度。 第~\ref{sec:ana-input-size}~節的
「擬多項式」也會出現一次。 C 程式用到《C 程式設計基礎》第~6~章的二維陣列、第~19~章的
\texttt{malloc}。

本章的例子圖頂點用字母，$s$ 是起點、$t$ 是終點。 手算時，「依序看鄰居」的順序寫在每個
例子裡；C 程式的頂點一律是 $0..n-1$，$s$ 是 $0$、$t$ 是 $n - 1$，印出時再換回字母。

\section{流網路與流}
\label{sec:flw-network}

\subsection{定義}

\begin{definition}[流網路]
\label{def:flw-network}
\textbf{流網路}（flow network）是一張有向圖 $G = (V, E)$，加上兩個指定的頂點：
\textbf{源點}（source）$s$ 和\textbf{匯點}（sink）$t$，以及每條邊 $(u, v) \in E$ 的
\textbf{容量}（capacity）$c(u, v) \ge 0$。 不是邊的頂點對，規定 $c(u, v) = 0$。
\end{definition}

本章另外做兩個假設，都不會讓問題變窄，第~\ref{sec:flw-model}~節會說明怎麼把一般的情
形轉成這樣。

\begin{enumerate}
\item \textbf{沒有自環，也沒有反平行邊}：$(u, v) \in E$ 時，$(v, u) \notin E$。 第~21~章的
簡單有向圖允許 $u \to v$ 和 $v \to u$ 同時存在，這裡不允許。 這個假設讓下一節剩餘網路的
定義只有一種情況要考慮。
\item \textbf{每個頂點都在某條 $s \leadsto t$ 的路徑上}。 走不到、或走了也到不了 $t$ 的頂
點，對送貨沒有幫助，可以先刪掉。 這樣 $s$ 以外的每個頂點都至少有一條進來的邊，所以
$m \ge n - 1$，$O(n + m)$ 可以寫成 $O(m)$。
\end{enumerate}

\begin{definition}[流]
\label{def:flw-flow}
流網路 $G$ 上的\textbf{流}（flow）是一個函數 $f : V \times V \to \mathbb{R}$，滿足
\begin{itemize}
\item \textbf{容量限制}：所有 $u, v \in V$，$0 \le f(u, v) \le c(u, v)$。
\item \textbf{流量守恆}：所有 $u \in V - \{s, t\}$，
\[
\sum_{v \in V} f(v, u) = \sum_{v \in V} f(u, v),
\]
也就是流進 $u$ 的總量等於流出 $u$ 的總量。
\end{itemize}
$f(u, v)$ 叫做邊 $(u, v)$ 上的流量。 流的\textbf{值}（value）是從源點淨流出的量：
\[
|f| = \sum_{v \in V} f(s, v) - \sum_{v \in V} f(v, s).
\]
\textbf{最大流問題}：給定流網路，找一個值最大的流。
\end{definition}

不是邊的頂點對，容量是 $0$，所以流量也只能是 $0$；守恆式裡的和實際上只跑過真正的邊。
$|f|$ 的寫法看起來像絕對值，其實只是一個記號。 本章的例子裡沒有邊指進 $s$，第二個和
是 $0$；定義保留它，是因為下面的證明要對一般的網路成立。

圖~\ref{fig:flw-net}(a) 是章首的物流網路。 集貨場是 $s$，配送中心是 $t$，五個轉運站是
$a$ 到 $e$；邊上的數字是每晚最多能載的板數。 圖 (b) 是其中一個流 $f_1$，每條邊標成
「流量$/$容量」，斜線只是分隔，不是除法。

\begin{center}
\begin{tikzpicture}[x=0.95cm, y=0.95cm]
\begin{scope}[xshift=0cm]
\node[flw-v] (Ns) at (0, 0) {$s$};
\node[flw-v] (Na) at (2.2, 1.6) {$a$};
\node[flw-v] (Nb) at (2.2, 0) {$b$};
\node[flw-v] (Nc) at (2.2, -1.6) {$c$};
\node[flw-v] (Nd) at (4.7, 0.95) {$d$};
\node[flw-v] (Ne) at (4.7, -0.95) {$e$};
\node[flw-v] (Nt) at (6.9, 0) {$t$};
\draw[flw-a] (Ns) -- node[flw-c] {5} (Na);
\draw[flw-a] (Ns) -- node[flw-c] {11} (Nb);
\draw[flw-a] (Ns) -- node[flw-c] {7} (Nc);
\draw[flw-a] (Nb) -- node[flw-c] {2} (Na);
\draw[flw-a] (Na) -- node[flw-c] {12} (Nd);
\draw[flw-a] (Nb) -- node[flw-c] {5} (Nd);
\draw[flw-a] (Nb) -- node[flw-c] {6} (Ne);
\draw[flw-a] (Nc) -- node[flw-c] {7} (Ne);
\draw[flw-a] (Nd) -- node[flw-c] {2} (Ne);
\draw[flw-a] (Nd) -- node[flw-c] {12} (Nt);
\draw[flw-a] (Ne) -- node[flw-c] {9} (Nt);
\node[flw-t] at (3.45, -2.45) {(a) 容量 $c$};
\end{scope}
\begin{scope}[xshift=8.4cm]
\node[flw-v] (Ms) at (0, 0) {$s$};
\node[flw-v] (Ma) at (2.2, 1.6) {$a$};
\node[flw-v] (Mb) at (2.2, 0) {$b$};
\node[flw-v] (Mc) at (2.2, -1.6) {$c$};
\node[flw-v] (Md) at (4.7, 0.95) {$d$};
\node[flw-v] (Me) at (4.7, -0.95) {$e$};
\node[flw-v] (Mt) at (6.9, 0) {$t$};
\draw[flw-a] (Ms) -- node[flw-c] {5/5} (Ma);
\draw[flw-a] (Ms) -- node[flw-c] {11/11} (Mb);
\draw[flw-a] (Ms) -- node[flw-c] {0/7} (Mc);
\draw[flw-a] (Mb) -- node[flw-c] {0/2} (Ma);
\draw[flw-a] (Ma) -- node[flw-c] {5/12} (Md);
\draw[flw-a] (Mb) -- node[flw-c] {5/5} (Md);
\draw[flw-a] (Mb) -- node[flw-c] {6/6} (Me);
\draw[flw-a] (Mc) -- node[flw-c] {0/7} (Me);
\draw[flw-a] (Md) -- node[flw-c] {0/2} (Me);
\draw[flw-a] (Md) -- node[flw-c] {10/12} (Mt);
\draw[flw-a] (Me) -- node[flw-c] {6/9} (Mt);
\node[flw-t] at (3.45, -2.45) {(b) 一個流 $f_1$，$|f_1| = 16$};
\end{scope}
\end{tikzpicture}\par\nopagebreak
\diagcaption{章首的物流網路（單位：板）。 (a) 每條邊的容量。 (b) 一個流 $f_1$，每條邊標
「流量$/$容量」。 第~\ref{sec:flw-ek-trace}~節會看到，它就是 Edmonds--Karp 做完三輪之後
的流。}
\label{fig:flw-net}
\end{center}

檢查 $f_1$ 是一個流：每條邊的流量都沒有超過容量；在 $b$，流進 $11$，流出
$f_1(b, a) + f_1(b, d) + f_1(b, e) = 0 + 5 + 6 = 11$；在 $d$，流進 $5 + 5 = 10$，流出
$0 + 10 = 10$；其他轉運站也一樣。 值是 $|f_1| = 5 + 11 + 0 = 16$，剛好等於流進 $t$ 的
$10 + 6$。 第~\ref{sec:flw-mfmc}~節會證明這一定成立：流的值等於流進匯點的淨量。 這
個流還不是最大的：第~\ref{sec:flw-ek}~節會算出這張網路的最大流是 $21$。

\begin{examnote}
有些課本（例如 CLRS 第 2 版）用另一種定義：$f(u, v)$ 代表「淨流量」，要求
$f(u, v) = -f(v, u)$（斜對稱），守恆寫成 $\sum_{v} f(u, v) = 0$。 兩種定義的最大流值相同，
剩餘容量也相同，差在寫法：淨流量的定義下，剩餘容量只用一個式子 $c_f(u, v) = c(u, v) - f(u, v)$，
流量表裡反方向的格子會出現負數。 作答前先看題目用的是哪一種，照題目的寫法填表；
題目沒有說明時，用本書的寫法（流量都不是負數），並在答案開頭寫一句「$f(u, v)$ 表示邊
$(u, v)$ 上的流量」。
\end{examnote}

\subsection{幾種常見的轉換}
\label{sec:flw-model}

實際的問題常常不直接長成上面的樣子，下面三種轉換最常用，都只多出 $O(n + m)$ 個頂點和
邊。

\textbf{反平行邊。} 如果 $u \to v$ 和 $v \to u$ 都有容量，挑其中一條，例如 $(v, u)$，
加一個新頂點 $x$，把它換成 $v \to x$ 和 $x \to u$ 兩條邊，容量都是 $c(v, u)$。 原本走
$v \to u$ 的流改走 $v \to x \to u$，反過來也一樣，所以兩張網路的流一一對應、值相同。
本章第~\ref{sec:flw-c}~節的程式用矩陣存剩餘容量，連這一步都可以省掉，到時再說明。

\textbf{多個源點、多個匯點。} 好幾個集貨場 $s_1, \ldots, s_k$ 和好幾個配送中心
$t_1, \ldots, t_\ell$ 時，加一個\textbf{超級源點} $s^{*}$，從它連到每個 $s_i$；再加一
個\textbf{超級匯點} $t^{*}$，從每個 $t_j$ 連過去。 新邊的容量設成「夠大」就好，例如
$s_i$ 所有出邊的容量和。 這和第~\ref{sec:gtr-multi-source}~節多起點 BFS 加一個虛擬起點是同一個想法。

\textbf{頂點的容量。} 如果轉運站 $v$ 一晚最多只能處理 $k$ 板，把 $v$ 拆成 $v_{\mathrm{in}}$
和 $v_{\mathrm{out}}$ 兩個頂點，中間連一條容量 $k$ 的邊 $v_{\mathrm{in}} \to v_{\mathrm{out}}$；
原本指進 $v$ 的邊改指進 $v_{\mathrm{in}}$，原本從 $v$ 出去的邊改從 $v_{\mathrm{out}}$ 出
去。 經過 $v$ 的流都得走那條新邊，所以總量不會超過 $k$。

\begin{center}
\begin{tikzpicture}[x=1cm, y=1cm]
\begin{scope}
\node[flw-v] (u) at (0, 0) {$u$};
\node[flw-v] (v) at (2.4, 0) {$v$};
\draw[flw-a] (u) to[bend left=18] node[flw-c] {6} (v);
\draw[flw-a] (v) to[bend left=18] node[flw-c] {4} (u);
\node at (3.35, 0) {$\Rightarrow$};
\node[flw-v] (u2) at (4.3, 0) {$u$};
\node[flw-v] (v2) at (6.7, 0) {$v$};
\node[flw-v] (x2) at (5.5, -1.1) {$x$};
\draw[flw-a] (u2) -- node[flw-c] {6} (v2);
\draw[flw-a] (v2) -- node[flw-c] {4} (x2);
\draw[flw-a] (x2) -- node[flw-c] {4} (u2);
\node[flw-t] at (3.35, -1.9) {(a) 拆掉反平行邊};
\end{scope}
\begin{scope}[xshift=8.6cm]
\node[flw-v] (p) at (0, 0.7) {$p$};
\node[flw-v] (q) at (0, -0.7) {$q$};
\node[flw-v, minimum size=0.8cm] (bi) at (1.9, 0) {$b_{\mathrm{in}}$};
\node[flw-v, minimum size=0.8cm] (bo) at (3.9, 0) {$b_{\mathrm{out}}$};
\node[flw-v] (r) at (5.8, 0.7) {$r$};
\node[flw-v] (w) at (5.8, -0.7) {$w$};
\draw[flw-a] (p) -- (bi);
\draw[flw-a] (q) -- (bi);
\draw[flw-a, flw-cut] (bi) -- node[flw-c] {$k$} (bo);
\draw[flw-a] (bo) -- (r);
\draw[flw-a] (bo) -- (w);
\node[flw-t] at (2.9, -1.9) {(b) 頂點 $b$ 的容量 $k$};
\end{scope}
\end{tikzpicture}\par\nopagebreak
\diagcaption{兩種轉換。 (a) $u$、$v$ 之間兩個方向都有邊，把 $v \to u$ 換成經過新頂點
$x$ 的兩條邊。 (b) 把頂點拆成兩個，中間的邊（粗線）就是頂點的容量。}
\label{fig:flw-model}
\end{center}

\section{剩餘網路與增廣路徑}
\label{sec:flw-residual}

\subsection{只會往前加，會卡住}

最直覺的做法是：找一條從 $s$ 到 $t$、每條邊都還沒滿的路徑，送到路上最緊的那條邊滿為
止，一直重複到找不到這種路徑。 這個做法會停在不是最大的地方。

圖~\ref{fig:flw-greedy}(a) 的小網路有 $s, x, y, t$ 四個頂點。 假設第一次挑了
$s \to x \to y \to t$，三條邊還能送的量是 $2, 2, 2$，送 $2$。 之後 $s \to x$ 滿了，
$s \to y$ 還能送 $1$，但 $y \to t$ 也滿了；$x \to t$ 還能送 $1$，可是到 $x$ 的唯一入口已
經滿了。 每條從 $s$ 到 $t$ 的路徑都至少有一條滿的邊，只能停在 $2$。 但這張網路其實可以
送 $3$：$s \to x$ 送 $2$，其中 $1$ 走 $x \to t$、$1$ 走 $x \to y \to t$；另外 $s \to y$ 送
$1$ 走 $y \to t$。

差別在 $x \to y$ 上的那兩板。 要從 $2$ 變成 $3$，得把其中一板「退回」$x$，讓它改走
$x \to t$，而 $y$ 少掉的那一板由 $s \to y$ 補上。 只會往前加的做法沒有「退回」這個動作，
所以第一步選錯就回不來。 下面的剩餘網路把「退回」也當成一條可以走的邊。

\begin{center}
\begin{tikzpicture}[x=1.0cm, y=1.0cm]
\begin{scope}[xshift=0cm]
\node[flw-v] (gs) at (0, 0) {$s$};
\node[flw-v] (gx) at (2, 1.1) {$x$};
\node[flw-v] (gy) at (2, -1.1) {$y$};
\node[flw-v] (gt) at (4, 0) {$t$};
\draw[flw-a] (gs) -- node[flw-c] {2/2} (gx);
\draw[flw-a] (gs) -- node[flw-c] {0/1} (gy);
\draw[flw-a] (gx) -- node[flw-c] {2/2} (gy);
\draw[flw-a] (gx) -- node[flw-c] {0/1} (gt);
\draw[flw-a] (gy) -- node[flw-c] {2/2} (gt);
\node[flw-t] at (2, -1.85) {(a) 先走 $s \to x \to y \to t$};
\end{scope}
\begin{scope}[xshift=5.2cm]
\node[flw-v] (hs) at (0, 0) {$s$};
\node[flw-v] (hx) at (2, 1.1) {$x$};
\node[flw-v] (hy) at (2, -1.1) {$y$};
\node[flw-v] (ht) at (4, 0) {$t$};
\draw[flw-a, flw-rev] (hx) -- node[flw-c] {2} (hs);
\draw[flw-a, flw-aug] (hs) -- node[flw-c, flw-cA] {1} (hy);
\draw[flw-a, flw-rev, flw-aug] (hy) -- node[flw-c, flw-cA] {2} (hx);
\draw[flw-a, flw-aug] (hx) -- node[flw-c, flw-cA] {1} (ht);
\draw[flw-a, flw-rev] (ht) -- node[flw-c] {2} (hy);
\node[flw-t] at (2, -1.85) {(b) 剩餘網路 $G_f$};
\end{scope}
\begin{scope}[xshift=10.4cm]
\node[flw-v] (ks) at (0, 0) {$s$};
\node[flw-v] (kx) at (2, 1.1) {$x$};
\node[flw-v] (ky) at (2, -1.1) {$y$};
\node[flw-v] (kt) at (4, 0) {$t$};
\draw[flw-a] (ks) -- node[flw-c] {2/2} (kx);
\draw[flw-a] (ks) -- node[flw-c] {1/1} (ky);
\draw[flw-a] (kx) -- node[flw-c] {1/2} (ky);
\draw[flw-a] (kx) -- node[flw-c] {1/1} (kt);
\draw[flw-a] (ky) -- node[flw-c] {2/2} (kt);
\node[flw-t] at (2, -1.85) {(c) 增廣之後，$|f| = 3$};
\end{scope}
\end{tikzpicture}\par\nopagebreak
\diagcaption{只會往前加會卡住的例子。 (a) 先沿 $s \to x \to y \to t$ 送 $2$，之後每條路
徑都有一條滿的邊。 (b) 這個流的剩餘網路，虛線是反向邊；紅色的路徑經過反向邊
$y \to x$。 (c) 沿紅色路徑送 $1$ 之後的流。}
\label{fig:flw-greedy}
\end{center}

\subsection{剩餘容量與剩餘網路}

\begin{definition}[剩餘網路]
\label{def:flw-residual}
$f$ 是流網路 $G$ 上的一個流。 對每一對頂點 $u, v$，定義\textbf{剩餘容量}（residual
capacity）
\[
c_f(u, v) =
\begin{cases}
c(u, v) - f(u, v) & \text{若 } (u, v) \in E, \\
f(v, u) & \text{若 } (v, u) \in E, \\
0 & \text{其他。}
\end{cases}
\]
\textbf{剩餘網路}（residual network）是 $G_f = (V, E_f)$，其中
$E_f = \{(u, v) : c_f(u, v) > 0\}$，每條邊的容量是 $c_f(u, v)$。
\end{definition}

因為沒有反平行邊，前兩種情況不會同時成立，$c_f(u, v)$ 的值是確定的。 原本的每條邊
$(u, v)$ 在剩餘網路裡最多變成兩條：
\begin{itemize}
\item \textbf{順向邊} $(u, v)$，容量 $c(u, v) - f(u, v)$：這條路還能多送多少。 邊滿了就沒
有這條。
\item \textbf{反向邊} $(v, u)$，容量 $f(u, v)$：最多能把多少已經送出去的量退回來。 邊上
沒有流就沒有這條。
\end{itemize}
所以 $|E_f| \le 2m$。 本章的圖把反向邊畫成虛線，方便分辨；在剩餘網路裡，兩種邊的地位完
全一樣。

圖~\ref{fig:flw-greedy}(b) 是送了 $2$ 之後的剩餘網路。 $s \to x$、$y \to t$ 滿了，只剩反
向邊 $x \to s$、$t \to y$；$x \to y$ 全部用掉，換成容量 $2$ 的反向邊 $y \to x$。 這時有
一條從 $s$ 到 $t$ 的路徑 $s \to y \to x \to t$（紅色），中間那一步就是「把 $x \to y$ 上的
量退回 $x$」。

\subsection{增廣路徑}

\begin{definition}[增廣路徑]
\label{def:flw-augpath}
剩餘網路 $G_f$ 裡從 $s$ 到 $t$ 的簡單路徑叫\textbf{增廣路徑}（augmenting path），也叫
擴增路徑。 路徑 $p$ 的\textbf{剩餘容量}是路上最小的剩餘容量
\[
c_f(p) = \min\{c_f(u, v) : (u, v) \text{ 在 } p \text{ 上}\},
\]
也叫 $p$ 的\textbf{瓶頸}（bottleneck）。 剩餘容量剛好等於 $c_f(p)$ 的邊叫 $p$ 的
\textbf{關鍵邊}（critical edge），一條路徑至少有一條。
\end{definition}

沿著增廣路徑 $p$ \textbf{增廣}（augment），就是對 $p$ 上的每條剩餘邊 $(u, v)$ 做一件事：
\begin{itemize}
\item 若 $(u, v) \in E$（順向邊），$f(u, v) \gets f(u, v) + c_f(p)$；
\item 若 $(v, u) \in E$（反向邊），$f(v, u) \gets f(v, u) - c_f(p)$。
\end{itemize}

\begin{lemma}
\label{lem:flw-augment}
$p$ 是 $G_f$ 的增廣路徑。 沿 $p$ 增廣之後得到的 $f'$ 仍然是一個流，而且
$|f'| = |f| + c_f(p)$。
\end{lemma}

\begin{proof}
記 $b = c_f(p) > 0$。 $p$ 是簡單路徑，不會同時經過 $(u, v)$ 和 $(v, u)$，所以每條原本的
邊最多被改一次。

\textbf{容量限制。} 順向邊：$f'(u, v) = f(u, v) + b \le f(u, v) + c_f(u, v) = c(u, v)$，
而且顯然不是負的。 反向邊：$f'(v, u) = f(v, u) - b \ge f(v, u) - c_f(u, v) = 0$，而且沒有變
大，不會超過 $c(v, u)$。

\textbf{流量守恆。} 不在 $p$ 上的頂點，相關的流量都沒變。 $p$ 經過的頂點 $w \ne s, t$ 剛
好有一條剩餘邊進來、一條剩餘邊出去。 進來的那條如果是順向邊 $(x, w)$，流進 $w$ 的量多
$b$；如果是反向邊（原本的邊是 $(w, x)$），流出 $w$ 的量少 $b$。 兩種情況下，「流進減流
出」都多了 $b$。 同理，出去的那條讓「流進減流出」少了 $b$。 合起來沒有變，守恆仍成立。

\textbf{值。} $p$ 從 $s$ 出發，不會再回到 $s$，所以 $s$ 只有一條出去的剩餘邊
$(s, y)$。 它是順向邊時 $f(s, y)$ 多 $b$；是反向邊時，原本的邊 $(y, s)$ 上的流少 $b$。
兩種情況都讓 $|f|$ 多 $b$。
\end{proof}

回到圖~\ref{fig:flw-greedy}：路徑 $s \to y \to x \to t$ 的剩餘容量是 $1, 2, 1$，瓶頸是
$1$。 $s \to y$、$x \to t$ 是順向邊，流量各加 $1$；$y \to x$ 是反向邊，$x \to y$ 的流量從
$2$ 減成 $1$。 得到圖 (c)，值是 $3$。

\section{Ford--Fulkerson 方法}
\label{sec:flw-ff}

\subsection{方法}

把上一節的動作一直重複，就是 Ford 和 Fulkerson 在 1956 年提出的方法：從全部是 $0$ 的流開
始，只要剩餘網路裡還有增廣路徑，就沿著它增廣。

\begin{pseudo}[label=alg:flw-ff]{Ford--Fulkerson 方法}
\Procedure{Ford-Fulkerson}{$G, s, t$}
    \For{每條邊 $(u, v) \in E$}
        \State $f(u, v) \gets 0$
    \EndFor
    \While{$G_f$ 裡有從 $s$ 到 $t$ 的簡單路徑 $p$}
        \State $b \gets \min\{c_f(u, v) : (u, v) \text{ 在 } p \text{ 上}\}$
        \For{$p$ 上的每條邊 $(u, v)$}
            \If{$(u, v) \in E$}
                \State $f(u, v) \gets f(u, v) + b$ \Comment{順向：多送}
            \Else
                \State $f(v, u) \gets f(v, u) - b$ \Comment{反向：退回}
            \EndIf
        \EndFor
    \EndWhile
    \State \Return $f$
\EndProcedure
\end{pseudo}

第 4 行沒有規定要找哪一條路徑，用 DFS、BFS 或任何方法找都可以，習慣上叫它 Ford--Fulkerson
「方法」：不同的找法，跑的步數可能差很多（見下面），但只要停下來，答案都是最大流。
「停下來時一定是最大流」要用第~\ref{sec:flw-mfmc}~節的最大流最小割定理證明；這裡先看它
會不會停、要跑多久。

\subsection{會不會停、跑多久}

\textbf{容量都是整數時一定會停。} 用歸納法：一開始流量都是 $0$；如果目前的流量都是整
數，剩餘容量也都是整數，瓶頸 $b$ 是正整數，增廣後流量仍是整數。 所以每一輪 $|f|$ 至少
增加 $1$。 另一方面 $|f|$ 不會超過從 $s$ 出去的容量總和，所以迴圈最多執行 $|f^{*}|$ 輪，
$f^{*}$ 是最大流。

每一輪要建出（或直接在原圖上模擬）剩餘網路、找一條路徑、沿路更新，都是 $O(n + m)$，在
本章的假設下就是 $O(m)$。 所以容量是整數時，Ford--Fulkerson 的時間是 $O(m\,|f^{*}|)$。

\begin{theorem}[整數性]
\label{thm:flw-integral}
流網路的容量都是整數時，Ford--Fulkerson 方法求出的最大流，每條邊的流量都是整數。 特別
是，一定存在一個流量全是整數的最大流。
\end{theorem}

上面的歸納已經證明流量全是整數；「它是最大流」等第~\ref{sec:flw-mfmc}~節再補上。 這個定
理看起來平凡，第~\ref{sec:flw-match}~節把匹配化成流的時候，靠的就是它：流量只能是 $0$
或 $1$，才能解讀成「配」或「不配」。

\subsection{選錯路徑的代價}
\label{sec:flw-bad}

$O(m\,|f^{*}|)$ 和容量的大小有關。 圖~\ref{fig:flw-bad} 的網路最大流是 $200$，兩次增廣
就夠：$s \to u \to t$、$s \to v \to t$ 各送 $100$。 但如果第 4 行很不巧，第一次選
$s \to u \to v \to t$，瓶頸是中間那條容量 $1$ 的邊，只送 $1$；第二次選
$s \to v \to u \to t$，其中 $v \to u$ 是反向邊，把剛剛送的那 $1$ 退回，又只送 $1$。 兩種
路徑輪流選，每次只增加 $1$，要 $200$ 輪才結束。

\begin{center}
\begin{tikzpicture}[x=1.1cm, y=1.1cm]
\node[flw-v] (s) at (0, 0) {$s$};
\node[flw-v] (u) at (2.2, 1.1) {$u$};
\node[flw-v] (v) at (2.2, -1.1) {$v$};
\node[flw-v] (t) at (4.4, 0) {$t$};
\draw[flw-a] (s) -- node[flw-c] {100} (u);
\draw[flw-a] (s) -- node[flw-c] {100} (v);
\draw[flw-a] (u) -- node[flw-c] {1} (v);
\draw[flw-a] (u) -- node[flw-c] {100} (t);
\draw[flw-a] (v) -- node[flw-c] {100} (t);
\end{tikzpicture}\par\nopagebreak
\diagcaption{兩條路徑輪流經過中間那條容量 $1$ 的邊，Ford--Fulkerson 要增廣 $200$ 次；BFS
找最短的路徑只要 $2$ 次。}
\label{fig:flw-bad}
\end{center}

把 $100$ 換成 $10^9$，就要二十億輪。 輸入只是五個數字，每個用三十個位元就寫得下，執行
時間卻和數值本身成正比，這正是第~\ref{sec:ana-input-size}~節說的\textbf{擬多項式}時間。 容量不是整數時
情況更糟：可以構造出容量是無理數的網路，讓 Ford--Fulkerson 永遠不停，流的值甚至收斂到比
最大流小的數。 下一節的 Edmonds--Karp 只改了第 4 行「怎麼找路徑」，就把這兩個問題一起解
決。

\section{Edmonds--Karp 演算法}
\label{sec:flw-ek}

\subsection{用 BFS 找增廣路徑}

Edmonds 和 Karp 在 1972 年發表的做法很簡單：每一輪都在剩餘網路上從 $s$ 做 BFS，找\textbf{邊
數最少}的增廣路徑。 由第~21~章的定理~\ref{thm:gtr-bfs}，BFS 樹上從 $s$ 到 $t$ 的路徑就是
$G_f$ 裡最短的路徑，長度是 $\delta_f(s, t)$。 這裡的 $\delta_f$ 是在剩餘網路 $G_f$ 上算的
距離，每條邊都算長度 $1$，和容量多大無關。

寫成程式時，不必真的把 $G_f$ 建出來。 只要記住每一對頂點的剩餘容量 $c_f(u, v)$：一開始
$c_f(u, v) = c(u, v)$；沿路徑送 $b$ 時，路上每條邊 $(u, v)$ 的 $c_f(u, v)$ 減 $b$，反方向的
$c_f(v, u)$ 加 $b$。 這一個規則同時處理了順向邊和反向邊：送過去的量變成反方向可以退回的
量，退回的量變成順方向可以再送的量。

\begin{pseudo}[label=alg:flw-ek]{Edmonds--Karp（用剩餘容量寫）}
\Function{Edmonds-Karp}{$G, s, t$}
    \For{每一對頂點 $u, v$}
        \State $c_f(u, v) \gets c(u, v)$
    \EndFor
    \State $\mathit{val} \gets 0$
    \While{\TRUE}
        \State \Call{BFS}{$G_f, s$}，只走 $c_f(u, v) > 0$ 的邊，得到 $\pi$
        \If{$t$ 沒有被發現}
            \State \Return $\mathit{val}$
        \EndIf
        \State $b \gets \infty$；$v \gets t$
        \While{$v \ne s$} \Comment{沿 $\pi$ 往回找瓶頸}
            \State $b \gets \min(b, c_f(\pi[v], v))$
            \State $v \gets \pi[v]$
        \EndWhile
        \State $v \gets t$
        \While{$v \ne s$} \Comment{沿 $\pi$ 往回增廣}
            \State $c_f(\pi[v], v) \gets c_f(\pi[v], v) - b$
            \State $c_f(v, \pi[v]) \gets c_f(v, \pi[v]) + b$
            \State $v \gets \pi[v]$
        \EndWhile
        \State $\mathit{val} \gets \mathit{val} + b$
    \EndWhile
\EndFunction
\end{pseudo}

第 6 行的 BFS 就是虛擬碼~\ref{alg:gtr-bfs}，只是「$v$ 是 $u$ 的鄰居」換成「$c_f(u, v) > 0$」。
第 9--12 行沿著前驅 $\pi$ 從 $t$ 走回 $s$，求瓶頸；第 13--17 行再走一次，更新剩餘容
量。 結束後，每條邊的流量是 $f(u, v) = c(u, v) - c_f(u, v)$，$(u, v) \in E$。

\textbf{鄰居的順序。} 剩餘網路裡可能有好幾條一樣短的增廣路徑，BFS 找到哪一條，要看鄰居
依什麼順序檢查。 本章的手算一律照頂點的固定順序（物流網路是 $s, a, b, c, d, e, t$）檢查
每個頂點的剩餘邊，不分順向、反向；C 程式也照這個順序，輸出可以逐列對照。 不同的順序可
能走出不同的路徑和不同的最大流，但最大流的\textbf{值}一定相同。

\subsection{手算追蹤}
\label{sec:flw-ek-trace}

在圖~\ref{fig:flw-net}(a) 的物流網路上執行 Edmonds--Karp。 下表每一列是一輪：先在目前的剩
餘網路上從 $s$ 做 BFS，列出每個頂點的距離 $\delta_f(s, v)$，再寫出 BFS 樹上的增廣路徑、
瓶頸、關鍵邊，以及增廣之後的 $|f|$。

\begin{center}
\small
\begin{tabular}{c|ccccccc|lccc}
\toprule
輪 & $s$ & $a$ & $b$ & $c$ & $d$ & $e$ & $t$ & 增廣路徑 & 瓶頸 & 關鍵邊 & $|f|$ \\
\midrule
1 & 0 & 1 & 1 & 1 & 2 & 2 & 3 & $s \to a \to d \to t$ & $5$ & $(s, a)$ & $5$ \\
2 & 0 & 2 & 1 & 1 & 2 & 2 & 3 & $s \to b \to d \to t$ & $5$ & $(b, d)$ & $10$ \\
3 & 0 & 2 & 1 & 1 & 3 & 2 & 3 & $s \to b \to e \to t$ & $6$ & $(s, b)$、$(b, e)$ & $16$ \\
4 & 0 & 4 & 3 & 1 & 4 & 2 & 3 & $s \to c \to e \to t$ & $3$ & $(e, t)$ & $19$ \\
5 & 0 & 4 & 3 & 1 & 5 & 2 & 6 & $s \to c \to e \to b \to a \to d \to t$ & $2$ & $(b, a)$、$(d, t)$ & $21$ \\
6 & 0 & $\infty$ & 3 & 1 & $\infty$ & 2 & $\infty$ & 走不到 $t$，結束 & & & \\
\bottomrule
\end{tabular}
\end{center}

圖~\ref{fig:flw-ek-rounds} 是每一輪開始時的剩餘網路，紅色是那一輪的增廣路徑，紅色的數字
是路徑上的剩餘容量。

\begin{center}
\begin{tikzpicture}[x=0.95cm, y=0.95cm]
\begin{scope}[xshift=0.0cm]
\node[flw-v] (R0s) at (0, 0) {$s$};
\node[flw-v] (R0a) at (2.2, 1.6) {$a$};
\node[flw-v] (R0b) at (2.2, 0) {$b$};
\node[flw-v] (R0c) at (2.2, -1.6) {$c$};
\node[flw-v] (R0d) at (4.7, 0.95) {$d$};
\node[flw-v] (R0e) at (4.7, -0.95) {$e$};
\node[flw-v] (R0t) at (6.9, 0) {$t$};
\draw[flw-a, flw-aug] (R0s) -- node[flw-c, flw-cA] {5} (R0a);
\draw[flw-a] (R0s) -- node[flw-c] {11} (R0b);
\draw[flw-a] (R0s) -- node[flw-c] {7} (R0c);
\draw[flw-a] (R0b) -- node[flw-c] {2} (R0a);
\draw[flw-a, flw-aug] (R0a) -- node[flw-c, flw-cA] {12} (R0d);
\draw[flw-a] (R0b) -- node[flw-c] {5} (R0d);
\draw[flw-a] (R0b) -- node[flw-c] {6} (R0e);
\draw[flw-a] (R0c) -- node[flw-c] {7} (R0e);
\draw[flw-a] (R0d) -- node[flw-c] {2} (R0e);
\draw[flw-a, flw-aug] (R0d) -- node[flw-c, flw-cA] {12} (R0t);
\draw[flw-a] (R0e) -- node[flw-c] {9} (R0t);
\node[flw-t] at (3.45, -2.45) {(a) 第 1 輪：$s \to a \to d \to t$，瓶頸 $5$};
\end{scope}
\begin{scope}[xshift=8.4cm]
\node[flw-v] (R1s) at (0, 0) {$s$};
\node[flw-v] (R1a) at (2.2, 1.6) {$a$};
\node[flw-v] (R1b) at (2.2, 0) {$b$};
\node[flw-v] (R1c) at (2.2, -1.6) {$c$};
\node[flw-v] (R1d) at (4.7, 0.95) {$d$};
\node[flw-v] (R1e) at (4.7, -0.95) {$e$};
\node[flw-v] (R1t) at (6.9, 0) {$t$};
\draw[flw-a, flw-rev] (R1a) -- node[flw-c] {5} (R1s);
\draw[flw-a, flw-aug] (R1s) -- node[flw-c, flw-cA] {11} (R1b);
\draw[flw-a] (R1s) -- node[flw-c] {7} (R1c);
\draw[flw-a] (R1b) -- node[flw-c] {2} (R1a);
\draw[flw-a] (R1a) to[bend left=17] node[flw-c] {7} (R1d);
\draw[flw-a, flw-rev] (R1d) to[bend left=17] node[flw-c] {5} (R1a);
\draw[flw-a, flw-aug] (R1b) -- node[flw-c, flw-cA] {5} (R1d);
\draw[flw-a] (R1b) -- node[flw-c] {6} (R1e);
\draw[flw-a] (R1c) -- node[flw-c] {7} (R1e);
\draw[flw-a] (R1d) -- node[flw-c] {2} (R1e);
\draw[flw-a, flw-aug] (R1d) to[bend left=17] node[flw-c, flw-cA] {7} (R1t);
\draw[flw-a, flw-rev] (R1t) to[bend left=17] node[flw-c] {5} (R1d);
\draw[flw-a] (R1e) -- node[flw-c] {9} (R1t);
\node[flw-t] at (3.45, -2.45) {(b) 第 2 輪：$s \to b \to d \to t$，瓶頸 $5$};
\end{scope}
\end{tikzpicture}

\medskip
\begin{tikzpicture}[x=0.95cm, y=0.95cm]
\begin{scope}[xshift=0.0cm]
\node[flw-v] (R2s) at (0, 0) {$s$};
\node[flw-v] (R2a) at (2.2, 1.6) {$a$};
\node[flw-v] (R2b) at (2.2, 0) {$b$};
\node[flw-v] (R2c) at (2.2, -1.6) {$c$};
\node[flw-v] (R2d) at (4.7, 0.95) {$d$};
\node[flw-v] (R2e) at (4.7, -0.95) {$e$};
\node[flw-v] (R2t) at (6.9, 0) {$t$};
\draw[flw-a, flw-rev] (R2a) -- node[flw-c] {5} (R2s);
\draw[flw-a, flw-aug] (R2s) to[bend left=17] node[flw-c, flw-cA] {6} (R2b);
\draw[flw-a, flw-rev] (R2b) to[bend left=17] node[flw-c] {5} (R2s);
\draw[flw-a] (R2s) -- node[flw-c] {7} (R2c);
\draw[flw-a] (R2b) -- node[flw-c] {2} (R2a);
\draw[flw-a] (R2a) to[bend left=17] node[flw-c] {7} (R2d);
\draw[flw-a, flw-rev] (R2d) to[bend left=17] node[flw-c] {5} (R2a);
\draw[flw-a, flw-rev] (R2d) -- node[flw-c] {5} (R2b);
\draw[flw-a, flw-aug] (R2b) -- node[flw-c, flw-cA] {6} (R2e);
\draw[flw-a] (R2c) -- node[flw-c] {7} (R2e);
\draw[flw-a] (R2d) -- node[flw-c] {2} (R2e);
\draw[flw-a] (R2d) to[bend left=17] node[flw-c] {2} (R2t);
\draw[flw-a, flw-rev] (R2t) to[bend left=17] node[flw-c] {10} (R2d);
\draw[flw-a, flw-aug] (R2e) -- node[flw-c, flw-cA] {9} (R2t);
\node[flw-t] at (3.45, -2.45) {(c) 第 3 輪：$s \to b \to e \to t$，瓶頸 $6$};
\end{scope}
\begin{scope}[xshift=8.4cm]
\node[flw-v] (R3s) at (0, 0) {$s$};
\node[flw-v] (R3a) at (2.2, 1.6) {$a$};
\node[flw-v] (R3b) at (2.2, 0) {$b$};
\node[flw-v] (R3c) at (2.2, -1.6) {$c$};
\node[flw-v] (R3d) at (4.7, 0.95) {$d$};
\node[flw-v] (R3e) at (4.7, -0.95) {$e$};
\node[flw-v] (R3t) at (6.9, 0) {$t$};
\draw[flw-a, flw-rev] (R3a) -- node[flw-c] {5} (R3s);
\draw[flw-a, flw-rev] (R3b) -- node[flw-c] {11} (R3s);
\draw[flw-a, flw-aug] (R3s) -- node[flw-c, flw-cA] {7} (R3c);
\draw[flw-a] (R3b) -- node[flw-c] {2} (R3a);
\draw[flw-a] (R3a) to[bend left=17] node[flw-c] {7} (R3d);
\draw[flw-a, flw-rev] (R3d) to[bend left=17] node[flw-c] {5} (R3a);
\draw[flw-a, flw-rev] (R3d) -- node[flw-c] {5} (R3b);
\draw[flw-a, flw-rev] (R3e) -- node[flw-c] {6} (R3b);
\draw[flw-a, flw-aug] (R3c) -- node[flw-c, flw-cA] {7} (R3e);
\draw[flw-a] (R3d) -- node[flw-c] {2} (R3e);
\draw[flw-a] (R3d) to[bend left=17] node[flw-c] {2} (R3t);
\draw[flw-a, flw-rev] (R3t) to[bend left=17] node[flw-c] {10} (R3d);
\draw[flw-a, flw-aug] (R3e) to[bend left=17] node[flw-c, flw-cA] {3} (R3t);
\draw[flw-a, flw-rev] (R3t) to[bend left=17] node[flw-c] {6} (R3e);
\node[flw-t] at (3.45, -2.45) {(d) 第 4 輪：$s \to c \to e \to t$，瓶頸 $3$};
\end{scope}
\end{tikzpicture}

\medskip
\begin{tikzpicture}[x=0.95cm, y=0.95cm]
\begin{scope}[xshift=0.0cm]
\node[flw-v] (R4s) at (0, 0) {$s$};
\node[flw-v] (R4a) at (2.2, 1.6) {$a$};
\node[flw-v] (R4b) at (2.2, 0) {$b$};
\node[flw-v] (R4c) at (2.2, -1.6) {$c$};
\node[flw-v] (R4d) at (4.7, 0.95) {$d$};
\node[flw-v] (R4e) at (4.7, -0.95) {$e$};
\node[flw-v] (R4t) at (6.9, 0) {$t$};
\draw[flw-a, flw-rev] (R4a) -- node[flw-c] {5} (R4s);
\draw[flw-a, flw-rev] (R4b) -- node[flw-c] {11} (R4s);
\draw[flw-a, flw-aug] (R4s) to[bend left=17] node[flw-c, flw-cA] {4} (R4c);
\draw[flw-a, flw-rev] (R4c) to[bend left=17] node[flw-c] {3} (R4s);
\draw[flw-a, flw-aug] (R4b) -- node[flw-c, flw-cA] {2} (R4a);
\draw[flw-a, flw-aug] (R4a) to[bend left=17] node[flw-c, flw-cA] {7} (R4d);
\draw[flw-a, flw-rev] (R4d) to[bend left=17] node[flw-c] {5} (R4a);
\draw[flw-a, flw-rev] (R4d) -- node[flw-c] {5} (R4b);
\draw[flw-a, flw-rev, flw-aug] (R4e) -- node[flw-c, flw-cA] {6} (R4b);
\draw[flw-a, flw-aug] (R4c) to[bend left=17] node[flw-c, flw-cA] {4} (R4e);
\draw[flw-a, flw-rev] (R4e) to[bend left=17] node[flw-c] {3} (R4c);
\draw[flw-a] (R4d) -- node[flw-c] {2} (R4e);
\draw[flw-a, flw-aug] (R4d) to[bend left=17] node[flw-c, flw-cA] {2} (R4t);
\draw[flw-a, flw-rev] (R4t) to[bend left=17] node[flw-c] {10} (R4d);
\draw[flw-a, flw-rev] (R4t) -- node[flw-c] {9} (R4e);
\node[flw-t] at (3.45, -2.45) {(e) 第 5 輪：$s \to c \to e \to b \to a \to d \to t$，瓶頸 $2$};
\end{scope}
\begin{scope}[xshift=8.4cm]
\node[flw-v, flw-vS] (R5s) at (0, 0) {$s$};
\node[flw-v] (R5a) at (2.2, 1.6) {$a$};
\node[flw-v, flw-vS] (R5b) at (2.2, 0) {$b$};
\node[flw-v, flw-vS] (R5c) at (2.2, -1.6) {$c$};
\node[flw-v] (R5d) at (4.7, 0.95) {$d$};
\node[flw-v, flw-vS] (R5e) at (4.7, -0.95) {$e$};
\node[flw-v] (R5t) at (6.9, 0) {$t$};
\draw[flw-a, flw-rev] (R5a) -- node[flw-c] {5} (R5s);
\draw[flw-a, flw-rev] (R5b) -- node[flw-c] {11} (R5s);
\draw[flw-a] (R5s) to[bend left=17] node[flw-c] {2} (R5c);
\draw[flw-a, flw-rev] (R5c) to[bend left=17] node[flw-c] {5} (R5s);
\draw[flw-a, flw-rev] (R5a) -- node[flw-c] {2} (R5b);
\draw[flw-a] (R5a) to[bend left=17] node[flw-c] {5} (R5d);
\draw[flw-a, flw-rev] (R5d) to[bend left=17] node[flw-c] {7} (R5a);
\draw[flw-a, flw-rev] (R5d) -- node[flw-c] {5} (R5b);
\draw[flw-a] (R5b) to[bend left=17] node[flw-c] {2} (R5e);
\draw[flw-a, flw-rev] (R5e) to[bend left=17] node[flw-c] {4} (R5b);
\draw[flw-a] (R5c) to[bend left=17] node[flw-c] {2} (R5e);
\draw[flw-a, flw-rev] (R5e) to[bend left=17] node[flw-c] {5} (R5c);
\draw[flw-a] (R5d) -- node[flw-c] {2} (R5e);
\draw[flw-a, flw-rev] (R5t) -- node[flw-c] {12} (R5d);
\draw[flw-a, flw-rev] (R5t) -- node[flw-c] {9} (R5e);
\node[flw-t] at (3.45, -2.45) {(f) 結束：走不到 $t$};
\end{scope}
\end{tikzpicture}\par\nopagebreak
\diagcaption{物流網路上 Edmonds--Karp 每一輪開始時的剩餘網路 $G_f$。 實線是順向邊，虛線是
反向邊，數字是剩餘容量；兩個方向都有剩餘容量時畫成兩條弧線。 (f) 的藍色頂點是從 $s$ 還走得
到的頂點。}
\label{fig:flw-ek-rounds}
\end{center}

逐輪看幾個地方：

\begin{itemize}
\item \textbf{第 1 輪。} 一開始 $G_f$ 就是原圖。 BFS 依序發現 $a, b, c$（距離 $1$），再從
$a$ 發現 $d$、從 $b$ 發現 $e$（距離 $2$），最後從 $d$ 發現 $t$。 $s \to b \to d \to t$ 也
一樣短，但 $a$ 比 $b$ 早被拿出佇列，$d$ 是 $a$ 發現的，所以路徑經過 $a$。
\item \textbf{第 2 輪。} $s \to a$ 滿了，$a$ 只能經由 $b \to a$ 到達，距離變成 $2$。 同時
多了反向邊 $a \to s$、$d \to a$、$t \to d$，容量都是 $5$。
\item \textbf{第 3 輪。} $b \to d$ 滿了，$d$ 改由 $a$ 發現，距離變成 $3$。 這一輪的瓶頸 $6$
同時出現在 $s \to b$ 和 $b \to e$，兩條都是關鍵邊。 做完之後的流就是圖~\ref{fig:flw-net}(b)
的 $f_1$。
\item \textbf{第 4 輪。} $s$ 出去的邊只剩 $s \to c$。 $b$ 要經過 $c \to e$ 和反向邊
$e \to b$ 才到得了，距離跳到 $3$，$a$ 跳到 $4$。 最短的增廣路徑是 $s \to c \to e \to t$，
瓶頸是 $e \to t$ 剩下的 $3$。
\item \textbf{第 5 輪。} $e \to t$ 也滿了，$t$ 的距離從 $3$ 跳到 $6$。 增廣路徑
$s \to c \to e \to b \to a \to d \to t$ 經過反向邊 $e \to b$：$b \to e$ 上原本有 $6$ 板，
其中 $2$ 板「退回」$b$，改走 $b \to a \to d \to t$；$e$ 少掉的那 $2$ 板，由
$s \to c \to e$ 送來的補上。 這和第~\ref{sec:flw-residual}~節的小例子是同一件事。
\item \textbf{第 6 輪。} BFS 只走得到 $s, b, c, e$，找不到增廣路徑，最大流是 $21$。
\end{itemize}

表裡每一欄由上往下看，距離都沒有變小過：$a$ 是 $1, 2, 2, 4, 4, \infty$，$t$ 是
$3, 3, 3, 3, 6, \infty$。 下一小節會證明這件事一定成立，它就是 Edmonds--Karp 跑得快的原因。

圖~\ref{fig:flw-final} 是最後的流。 藍色的頂點和橘色的邊是下一節要講的最小割。

\begin{center}
\begin{tikzpicture}[x=1.1cm, y=1.1cm]
\node[flw-v, flw-vS] (Zs) at (0, 0) {$s$};
\node[flw-v] (Za) at (2.2, 1.6) {$a$};
\node[flw-v, flw-vS] (Zb) at (2.2, 0) {$b$};
\node[flw-v, flw-vS] (Zc) at (2.2, -1.6) {$c$};
\node[flw-v] (Zd) at (4.7, 0.95) {$d$};
\node[flw-v, flw-vS] (Ze) at (4.7, -0.95) {$e$};
\node[flw-v] (Zt) at (6.9, 0) {$t$};
\draw[flw-a, flw-cut] (Zs) -- node[flw-c] {5/5} (Za);
\draw[flw-a] (Zs) -- node[flw-c] {11/11} (Zb);
\draw[flw-a] (Zs) -- node[flw-c] {5/7} (Zc);
\draw[flw-a, flw-cut] (Zb) -- node[flw-c] {2/2} (Za);
\draw[flw-a] (Za) -- node[flw-c] {7/12} (Zd);
\draw[flw-a, flw-cut] (Zb) -- node[flw-c] {5/5} (Zd);
\draw[flw-a] (Zb) -- node[flw-c] {4/6} (Ze);
\draw[flw-a] (Zc) -- node[flw-c] {5/7} (Ze);
\draw[flw-a] (Zd) -- node[flw-c] {0/2} (Ze);
\draw[flw-a] (Zd) -- node[flw-c] {12/12} (Zt);
\draw[flw-a, flw-cut] (Ze) -- node[flw-c] {9/9} (Zt);
\end{tikzpicture}\par\nopagebreak
\diagcaption{Edmonds--Karp 求出的最大流，$|f| = 21$。 藍色的 $S = \{s, b, c, e\}$ 是最後的
剩餘網路裡從 $s$ 走得到的頂點，橘色的四條邊從 $S$ 指向 $T = \{a, d, t\}$，全部是滿的。}
\label{fig:flw-final}
\end{center}

\begin{examnote}
手算增廣的題目，每一輪寫出：目前的剩餘網路（或列出所有 $c_f > 0$ 的邊和剩餘容量）、增
廣路徑、瓶頸、增廣後的 $|f|$。 最常漏的是反向邊：每送一次，路上每條邊都要多一條反方向、
同樣大小的剩餘容量。 題目沒有規定找路徑的順序時，在答案開頭寫明「BFS，鄰居依字母順
序」，因為不同的順序會走出不同的路徑。 最後再寫出最大流的值和一個最小割，兩者相等，就是
最好的驗算。
\end{examnote}

\subsection{增廣次數是 $O(nm)$}
\label{sec:flw-ek-proof}

以下 $f$ 是某一輪增廣之前的流，$f'$ 是增廣之後的流，$\delta_f(s, v)$、$\delta_{f'}(s, v)$
分別是 $G_f$、$G_{f'}$ 裡從 $s$ 到 $v$ 的邊數距離。

\begin{lemma}
\label{lem:flw-mono}
執行 Edmonds--Karp 時，對每個頂點 $v \ne s$，$\delta_f(s, v)$ 在每一次增廣之後都不會變小。
\end{lemma}

\begin{proof}
用反證法。 假設某次增廣（從 $f$ 變成 $f'$）讓某些頂點的距離變小了，在這些頂點裡取
$\delta_{f'}(s, v)$ 最小的一個 $v$。 $v$ 在 $G_{f'}$ 裡走得到（距離變小，不會是 $\infty$），
取 $G_{f'}$ 裡從 $s$ 到 $v$ 的一條最短路徑，$u$ 是 $v$ 前面那個頂點，所以
\[
\delta_{f'}(s, u) = \delta_{f'}(s, v) - 1 .
\]
$u$ 的距離在 $G_{f'}$ 裡比 $v$ 小，由 $v$ 的取法，$u$ 的距離沒有變小（$u = s$ 時距離本來
就是 $0$）：$\delta_f(s, u) \le \delta_{f'}(s, u)$。

先看 $(u, v)$ 是不是 $G_f$ 的邊。 如果是，由第~\ref{sec:gtr-bfs-proof}~節 BFS 證明裡的事實 (a)（走一條邊距離最
多加一），
\[
\delta_f(s, v) \le \delta_f(s, u) + 1 \le \delta_{f'}(s, u) + 1 = \delta_{f'}(s, v),
\]
和「$v$ 的距離變小」矛盾。 所以 $(u, v)$ 不在 $G_f$ 裡，卻在 $G_{f'}$ 裡。 剩餘容量從 $0$
變成正的，只可能是這次的增廣路徑經過了反方向的 $(v, u)$。 增廣路徑是 $G_f$ 裡的最短路
徑，它經過 $(v, u)$，表示 $\delta_f(s, u) = \delta_f(s, v) + 1$。 於是
\[
\delta_f(s, v) = \delta_f(s, u) - 1 \le \delta_{f'}(s, u) - 1 = \delta_{f'}(s, v) - 2,
\]
距離反而至少增加了 $2$，又矛盾。
\end{proof}

\begin{theorem}
\label{thm:flw-ek}
Edmonds--Karp 演算法最多增廣 $nm$ 次，時間是 $O(nm^2)$。
\end{theorem}

\begin{proof}
看一條可能出現在剩餘網路裡的邊 $(u, v)$，也就是 $(u, v) \in E$ 或 $(v, u) \in E$，這種邊最
多 $2m$ 條。 每次增廣，路徑上至少有一條關鍵邊，增廣之後它的剩餘容量變成 $0$，從剩餘網路
消失。

假設 $(u, v)$ 在某一輪（流是 $f$）當了關鍵邊。 它在 BFS 求出的最短路徑上，所以
$\delta_f(s, v) = \delta_f(s, u) + 1$。 增廣之後它消失了；要再出現，之後某一輪（流是
$f''$）的增廣路徑必須經過 $(v, u)$，那時 $\delta_{f''}(s, u) = \delta_{f''}(s, v) + 1$。 由
引理~\ref{lem:flw-mono}，$\delta_{f''}(s, v) \ge \delta_f(s, v)$，所以
\[
\delta_{f''}(s, u) = \delta_{f''}(s, v) + 1 \ge \delta_f(s, v) + 1 = \delta_f(s, u) + 2 .
\]
也就是說，$(u, v)$ 每多當一次關鍵邊，$u$ 的距離至少多 $2$。 $u$ 在增廣路徑上又不是 $t$，
距離最多是 $n - 2$；從 $0$ 開始每次至少加 $2$，$(u, v)$ 當關鍵邊最多
$(n - 2)/2 + 1 = n/2$ 次。

每次增廣至少有一條關鍵邊，所以增廣次數不超過 $2m \cdot n/2 = nm$。 每一輪的 BFS 和更新都
是 $O(n + m) = O(m)$，總時間 $O(nm^2)$。
\end{proof}

這個上界只和 $n$、$m$ 有關，跟容量多大、甚至是不是整數都無關：無理數的容量也一定在
$nm$ 輪內停下來。 圖~\ref{fig:flw-bad} 的網路，Edmonds--Karp 第一輪就找到長度 $2$ 的
$s \to u \to t$，兩輪就結束。 證明裡「每次至少加 $2$」的估計並不緊，實際上增廣次數通常
遠少於 $nm$；物流網路 $n = 7$、$m = 11$，只增廣了 $5$ 次。

最大流還有更快的演算法。 Dinic（1970）在每個階段用 BFS 把頂點分層，一次找出很多條同樣
長度的增廣路徑，時間 $O(n^2 m)$；推送--重標（push-relabel）一類的方法不找完整的路徑，
讓頂點暫時多收一些流，再往前推，可以做到 $O(n^3)$。 這些都超出本書的範圍，延伸閱讀列了
參考資料。

\section{最大流最小割定理}
\label{sec:flw-mfmc}

\subsection{割}

章首的物流經理還問了一句：要增加送貨量，該先替哪一段路線加車？ 直覺上，如果某幾條路線
一斷，集貨場就完全送不到配送中心，那這幾條路線的容量加起來，就是送貨量的上限。 這一節
把這個直覺寫成定理。

\begin{definition}[割]
\label{def:flw-cut}
流網路的一個\textbf{割}（cut）$(S, T)$ 是把 $V$ 分成兩部分 $S$ 和 $T = V - S$，而且
$s \in S$、$t \in T$。 割的\textbf{容量}是從 $S$ 指向 $T$ 的邊的容量和：
\[
c(S, T) = \sum_{u \in S} \sum_{v \in T} c(u, v) .
\]
給定流 $f$，通過割的\textbf{淨流量}是
\[
f(S, T) = \sum_{u \in S} \sum_{v \in T} f(u, v) - \sum_{u \in S} \sum_{v \in T} f(v, u) .
\]
容量最小的割叫\textbf{最小割}（minimum cut）。
\end{definition}

兩個定義不對稱，要特別注意：容量只算從 $S$ 到 $T$ 的邊；淨流量則是從 $S$ 到 $T$ 的流
量，減掉從 $T$ 回到 $S$ 的流量。

以物流網路和圖~\ref{fig:flw-net}(b) 的流 $f_1$ 為例，取 $S = \{s, a, d\}$，
$T = \{b, c, e, t\}$。 從 $S$ 到 $T$ 的邊是 $s \to b$、$s \to c$、$d \to e$、$d \to t$，
從 $T$ 到 $S$ 的邊是 $b \to a$、$b \to d$。 所以
\begin{align*}
c(S, T) &= 11 + 7 + 2 + 12 = 32, \\
f_1(S, T) &= (11 + 0 + 0 + 10) - (0 + 5) = 16 .
\end{align*}
淨流量剛好是 $|f_1| = 16$。 再取 $S = \{s, b, c, e\}$：從 $S$ 到 $T$ 的邊是 $s \to a$、
$b \to a$、$b \to d$、$e \to t$，$c(S, T) = 5 + 2 + 5 + 9 = 21$；從 $T$ 到 $S$ 的邊只有
$d \to e$，它的容量不算進 $c(S, T)$。

\begin{lemma}
\label{lem:flw-cutflow}
$f$ 是流，$(S, T)$ 是任意一個割，則 $f(S, T) = |f|$。
\end{lemma}

\begin{proof}
對每個頂點 $u$，記 $\mathrm{out}(u) = \sum_{v} f(u, v)$、$\mathrm{in}(u) = \sum_{v} f(v, u)$。
由定義，$\mathrm{out}(s) - \mathrm{in}(s) = |f|$；由流量守恆，$S$ 裡其他的頂點（$t$ 不在
$S$ 裡）都是 $\mathrm{out}(u) - \mathrm{in}(u) = 0$。 加起來，
\[
\sum_{u \in S} \bigl(\mathrm{out}(u) - \mathrm{in}(u)\bigr) = |f| .
\]
把左邊每一項 $f(x, y)$ 依另一端在 $S$ 還是 $T$ 分開。 兩端都在 $S$ 的 $f(x, y)$，在
$\mathrm{out}(x)$ 裡加一次、在 $\mathrm{in}(y)$ 裡減一次，互相抵消。 剩下的，是 $u \in S$
流向 $T$ 的量（加）和 $T$ 流進 $u \in S$ 的量（減），合起來正是 $f(S, T)$。
\end{proof}

取 $S = V - \{t\}$，就得到 $|f|$ 等於流進 $t$ 的淨量，這是前面 $f_1$ 的例子看到的事。

\begin{corollary}
\label{cor:flw-weak}
任何流 $f$ 和任何割 $(S, T)$，都有 $|f| \le c(S, T)$。
\end{corollary}

\begin{proof}
$|f| = f(S, T) \le \sum_{u \in S} \sum_{v \in T} f(u, v) \le \sum_{u \in S} \sum_{v \in T}
c(u, v) = c(S, T)$。 第一個不等號丟掉了減掉的那一項（它不是負的），第二個用容量限制。
\end{proof}

所以\textbf{任何}一個割的容量都是\textbf{任何}一個流的上限，最大流的值不超過最小割的容
量。 物流網路的割 $\{s, b, c, e\}$ 容量 $21$，Edmonds--Karp 也找到了值 $21$ 的流，兩邊碰
在一起，就同時證明了「流 $21$ 是最大的」和「割 $21$ 是最小的」。 下面的定理說，兩邊一定
會碰在一起。

\subsection{定理與證明}

\begin{theorem}[最大流最小割定理]
\label{thm:flw-mfmc}
$f$ 是流網路 $G$ 上的一個流。 下面三件事等價：
\begin{enumerate}
\item $f$ 是最大流。
\item 剩餘網路 $G_f$ 裡沒有增廣路徑。
\item 有某個割 $(S, T)$ 使得 $|f| = c(S, T)$。
\end{enumerate}
\end{theorem}

\begin{proof}
\textbf{(1) $\Rightarrow$ (2)。} 如果 $G_f$ 有增廣路徑 $p$，由引理~\ref{lem:flw-augment}，
沿 $p$ 增廣得到的流值是 $|f| + c_f(p) > |f|$，$f$ 就不是最大流。

\textbf{(2) $\Rightarrow$ (3)。} 令 $S$ 是 $G_f$ 裡從 $s$ 走得到的頂點，$T = V - S$。
$s \in S$；$G_f$ 沒有增廣路徑，所以 $t \notin S$，$(S, T)$ 是一個割。 取 $u \in S$、
$v \in T$：
\begin{itemize}
\item 若 $(u, v) \in E$，則 $f(u, v) = c(u, v)$。 否則 $c_f(u, v) = c(u, v) - f(u, v) > 0$，
$(u, v)$ 是 $G_f$ 的邊，$v$ 也走得到，和 $v \in T$ 矛盾。
\item 若 $(v, u) \in E$，則 $f(v, u) = 0$。 否則 $c_f(u, v) = f(v, u) > 0$，同樣矛盾。
\end{itemize}
不是邊的頂點對流量都是 $0$。 所以
\[
f(S, T) = \sum_{u \in S} \sum_{v \in T} c(u, v) - 0 = c(S, T),
\]
再由引理~\ref{lem:flw-cutflow}，$|f| = f(S, T) = c(S, T)$。

\textbf{(3) $\Rightarrow$ (1)。} 由系理~\ref{cor:flw-weak}，任何流 $g$ 都有
$|g| \le c(S, T) = |f|$，所以 $f$ 是最大流。
\end{proof}

這個定理一次給了三個結論。 第一，\textbf{最大流的值等於最小割的容量}：由 (1) 推到 (3)，
最大流的值等於某個割的容量，而系理~\ref{cor:flw-weak} 說沒有割比它小。 第二，
\textbf{Ford--Fulkerson 是對的}：虛擬碼~\ref{alg:flw-ff} 停下來時 $G_f$ 沒有增廣路徑，由
(2) 推到 (1)，手上的流就是最大流。 這也補完了定理~\ref{thm:flw-integral}：容量是整數時，
Ford--Fulkerson 會停，停下來的流量全是整數，又是最大流。 第三，\textbf{最小割可以順便求
出來}，就是證明裡的 $S$。

\subsection{手找最小割}
\label{sec:flw-mincut}

求出最大流之後，在最後的剩餘網路上從 $s$ 做一次 BFS 或 DFS，走得到的頂點是 $S$，其他是
$T$，多花 $O(n + m)$。 Edmonds--Karp 最後那一輪 BFS 已經做了這件事：表裡第 6 列距離不是
$\infty$ 的頂點，就是 $S = \{s, b, c, e\}$。

驗算的方法是看證明的 (2) $\Rightarrow$ (3)：從 $S$ 到 $T$ 的邊都要滿，從 $T$ 到 $S$ 的邊
都要沒有流。 圖~\ref{fig:flw-final} 裡，$s \to a$（$5/5$）、$b \to a$（$2/2$）、
$b \to d$（$5/5$）、$e \to t$（$9/9$）都滿了，唯一從 $T$ 回到 $S$ 的 $d \to e$ 流量是
$0$；容量和是 $21$，等於 $|f|$。

最小割不一定只有一個。 物流網路的 $S' = \{s, a, b, c, d, e\}$（只有 $t$ 在另一邊）也是最
小割：$d \to t$ 和 $e \to t$ 的容量和是 $12 + 9 = 21$。 用剩餘網路找到的那一個，$S$ 是所有
最小割裡最小的（觀念題第~\ref{q:flw-smallest}~題）。

回到物流經理的問題。 加大某條邊的容量，最大流要變大，所有的最小割都得變大，所以這條邊
必須出現在\textbf{每一個}最小割裡；只在其中一個最小割裡的邊，加了也沒用，因為另一個最小
割還是 $21$。 物流網路的兩個最小割共有的邊只有 $e \to t$。 實際算一下：$s \to a$ 加 $3$
板，最大流還是 $21$；$e \to t$ 加 $3$ 板，最大流變成 $23$（不是 $24$，因為又有別的割變成
最緊的）。 只靠「加在最小割的邊上」不保證有用，但不在每個最小割裡的邊一定沒用。

\begin{examnote}
題目問「最小割」時，寫出 $S$、$T$、從 $S$ 到 $T$ 的每條邊和它們的容量和，並說明它等於
最大流的值。 容量只加從 $S$ 指向 $T$ 的邊；從 $T$ 指回 $S$ 的邊，不論容量多大都不算。 題
目問「所有的最小割」時，不能只交剩餘網路找到的那一個，要逐一檢查其他的分法（可以用觀念
題第~\ref{q:flw-allcuts}~題的方法縮小範圍）。
\end{examnote}

\section{二分匹配}
\label{sec:flw-match}

\subsection{匹配}

\begin{definition}[匹配]
\label{def:flw-matching}
無向圖 $G = (V, E)$ 的一個\textbf{匹配}（matching）是一個邊的子集合 $M \subseteq E$，其中
任兩條邊都沒有共同的端點。 $M$ 裡某條邊的端點叫\textbf{已匹配}（matched）的頂點，其他叫
\textbf{未匹配}（free，也叫自由）的頂點。 邊數最多的匹配叫\textbf{最大匹配}（maximum
matching）；每個頂點都已匹配的叫\textbf{完美匹配}（perfect matching）。
\end{definition}

這一節只處理二分圖。 照第~21~章的定義~\ref{def:gtr-bipartite}，二分圖的頂點分成 $L$ 和 $R$
兩邊，每條邊一端在 $L$、一端在 $R$；寫成 $G = (L \cup R, E)$。 章首的助教分配就是二分圖：
$L = \{x_1, x_2, x_3, x_4\}$ 是四位助教，$R = \{y_1, \ldots, y_5\}$ 是五堂實習課，助教
$x_i$ 有空帶 $y_j$ 就連一條邊。

\begin{center}
\small
\begin{tabular}{cl}
\toprule
助教 & 有空的實習課 \\
\midrule
$x_1$ & $y_3$、$y_5$ \\
$x_2$ & $y_1$、$y_2$、$y_4$、$y_5$ \\
$x_3$ & $y_1$、$y_2$ \\
$x_4$ & $y_2$ \\
\bottomrule
\end{tabular}
\end{center}

一個匹配就是一種分配：每位助教最多帶一堂，每堂課最多一位助教。 最直接的做法是依序讓每
位助教挑第一堂還沒人帶的課：$x_1$ 拿 $y_3$、$x_2$ 拿 $y_1$、$x_3$ 拿 $y_2$，輪到 $x_4$
時，他唯一能帶的 $y_2$ 已經被拿走了（圖~\ref{fig:flw-ta}(a)）。 這三條邊已經不能再加任何
一條邊，但最大匹配有四條（圖 (b)）。

\begin{center}
\begin{tikzpicture}[x=1cm,y=0.95cm]
\begin{scope}[xshift=0cm]
\node[flw-v] (Ax1) at (0, 1.5) {$x_1$};
\node[flw-v] (Ax2) at (0, 0.5) {$x_2$};
\node[flw-v] (Ax3) at (0, -0.5) {$x_3$};
\node[flw-v] (Ax4) at (0, -1.5) {$x_4$};
\node[flw-v] (Ay1) at (2.4, 2) {$y_1$};
\node[flw-v] (Ay2) at (2.4, 1) {$y_2$};
\node[flw-v] (Ay3) at (2.4, 0) {$y_3$};
\node[flw-v] (Ay4) at (2.4, -1) {$y_4$};
\node[flw-v] (Ay5) at (2.4, -2) {$y_5$};
\draw[flw-m] (Ax1) -- (Ay3);
\draw[flw-e] (Ax1) -- (Ay5);
\draw[flw-m] (Ax2) -- (Ay1);
\draw[flw-e] (Ax2) -- (Ay2);
\draw[flw-e] (Ax2) -- (Ay4);
\draw[flw-e] (Ax2) -- (Ay5);
\draw[flw-e] (Ax3) -- (Ay1);
\draw[flw-m] (Ax3) -- (Ay2);
\draw[flw-e] (Ax4) -- (Ay2);
\node[flw-t] at (1.2, -2.75) {(a) 依序配給第一個有空的課};
\end{scope}
\begin{scope}[xshift=5.2cm]
\node[flw-v] (Bx1) at (0, 1.5) {$x_1$};
\node[flw-v] (Bx2) at (0, 0.5) {$x_2$};
\node[flw-v] (Bx3) at (0, -0.5) {$x_3$};
\node[flw-v] (Bx4) at (0, -1.5) {$x_4$};
\node[flw-v] (By1) at (2.4, 2) {$y_1$};
\node[flw-v] (By2) at (2.4, 1) {$y_2$};
\node[flw-v] (By3) at (2.4, 0) {$y_3$};
\node[flw-v] (By4) at (2.4, -1) {$y_4$};
\node[flw-v] (By5) at (2.4, -2) {$y_5$};
\draw[flw-m] (Bx1) -- (By3);
\draw[flw-e] (Bx1) -- (By5);
\draw[flw-e] (Bx2) -- (By1);
\draw[flw-e] (Bx2) -- (By2);
\draw[flw-m] (Bx2) -- (By4);
\draw[flw-e] (Bx2) -- (By5);
\draw[flw-m] (Bx3) -- (By1);
\draw[flw-e] (Bx3) -- (By2);
\draw[flw-m] (Bx4) -- (By2);
\node[flw-t] at (1.2, -2.75) {(b) 最大匹配};
\end{scope}
\end{tikzpicture}\par\nopagebreak
\diagcaption{助教分配的二分圖，粗紅線是匹配。 (a) 依序讓 $x_1, x_2, x_3$ 挑第一堂沒人帶的課，
$x_4$ 沒課可帶；這個匹配不能再加邊。 (b) 四條邊的最大匹配。}
\label{fig:flw-ta}
\end{center}

\begin{examnote}
「不能再加邊」的匹配叫\textbf{極大匹配}（maximal matching），「邊數最多」的叫
\textbf{最大匹配}（maximum matching）。 最大匹配一定是極大匹配，反過來不一定，圖~
\ref{fig:flw-ta}(a) 就是反例。 中文的「極大」「最大」只差一個字，英文也只差兩個字母；作
答時遇到這兩個詞，先確認題目要的是哪一個，必要時把英文寫在旁邊。
\end{examnote}

\subsection{轉成最大流}
\label{sec:flw-reduce}

把二分圖 $G = (L \cup R, E)$ 改成流網路 $G'$：
\begin{itemize}
\item 頂點是 $L \cup R$，再加上新的源點 $s$ 和匯點 $t$。
\item 從 $s$ 到每個 $x \in L$ 連一條邊；從每個 $y \in R$ 到 $t$ 連一條邊。
\item $G$ 的每條邊 $\{x, y\}$（$x \in L$、$y \in R$）改成有向邊 $x \to y$，方向一律從 $L$
指向 $R$。
\item 所有邊的容量都是 $1$。
\end{itemize}
$G'$ 有 $n + 2$ 個頂點、$m + |L| + |R| = m + n$ 條邊，圖~\ref{fig:flw-tanet}(a) 是助教例子
的 $G'$。 想成：$s$ 給每位助教一張「派課單」，派課單只能沿著「有空」的邊交給一堂課，每堂
課最多收一張，再交給 $t$。

\begin{lemma}
\label{lem:flw-match-flow}
(1) $G$ 的每個匹配 $M$ 對應 $G'$ 的一個整數流 $f$，$|f| = |M|$。 (2) 反過來，$G'$ 的每個
整數流 $f$ 對應 $G$ 的一個匹配 $M$，$|M| = |f|$。
\end{lemma}

\begin{proof}
(1) 對 $M$ 的每條邊 $\{x, y\}$，在 $s \to x \to y \to t$ 三條邊上各放流量 $1$，其他邊都是
$0$。 $M$ 的邊沒有共同端點，所以每條 $s \to x$、$y \to t$ 最多被放一次，不超過容量；每個
$x$、$y$ 流進和流出都是 $0$ 或 $1$，守恆成立。 值是從 $s$ 出去的量，等於 $|M|$。

(2) 容量都是 $1$，整數流每條邊的流量不是 $0$ 就是 $1$。 令
$M = \{\{x, y\} : x \in L,\ y \in R,\ f(x, y) = 1\}$。 每個 $x \in L$ 只有一條進來的邊
$s \to x$，流進最多 $1$；由守恆，$x$ 出去的邊最多一條有流。 同理每個 $y \in R$ 只有一條出
去的邊 $y \to t$，進來的邊最多一條有流。 所以 $M$ 的邊沒有共同端點，是一個匹配。 取割
$S = \{s\} \cup L$：$s \to x$ 兩端都在 $S$ 裡，所以從 $S$ 到 $T$ 的邊只有 $L$ 到 $R$ 的
邊；從 $T$ 回到 $S$ 的邊一條也沒有（$G'$ 沒有從 $R$ 指回 $L$ 的邊，也沒有指向 $s$ 的邊）。
由引理~\ref{lem:flw-cutflow}，
$|f| = f(S, T) = \sum_{x \in L} \sum_{y \in R} f(x, y) = |M|$。
\end{proof}

\begin{theorem}
\label{thm:flw-match}
二分圖 $G$ 的最大匹配的邊數，等於 $G'$ 的最大流的值。 在 $G'$ 上執行 Ford--Fulkerson，把
有流的 $x \to y$ 邊挑出來，就是一個最大匹配，時間 $O(n(n + m))$。
\end{theorem}

\begin{proof}
由定理~\ref{thm:flw-integral}，Ford--Fulkerson 求出的最大流 $f$ 是整數流，由引理~\ref{lem:flw-match-flow}(2)
它對應一個 $|f|$ 條邊的匹配。 如果有更大的匹配 $M$，由引理~\ref{lem:flw-match-flow}(1)，
就有值 $|M| > |f|$ 的流，和 $f$ 是最大流矛盾。 時間：最大匹配最多 $\min(|L|, |R|) \le n/2$
條邊，所以 $|f^{*}| \le n/2$，迴圈最多 $n/2$ 輪；每一輪在 $G'$ 上找路徑是
$O(n + m)$。
\end{proof}

如果 $G$ 沒有孤立的頂點，$n \le 2m$，時間可以寫成 $O(nm)$。 這裡的瓶頸永遠是 $1$，所以用
DFS 還是 BFS 找路徑，輪數都一樣是 $|f^{*}|$；Edmonds--Karp 的 $O(nm^2)$ 在這裡反而是比較
鬆的估計。

\subsection{手算：分配助教}
\label{sec:flw-match-trace}

在 $G'$ 上執行 Edmonds--Karp，鄰居依 $s, x_1, \ldots, x_4, y_1, \ldots, y_5, t$ 的順序檢查。
前三輪的剩餘網路裡，$s \to x_i \to y_j \to t$ 這種長度 $3$ 的路徑還在，BFS 依序找到

\begin{center}
\small
\begin{tabular}{clcl}
\toprule
輪 & 增廣路徑 & 長度 & 匹配的變化 \\
\midrule
1 & $s \to x_1 \to y_3 \to t$ & 3 & 加入 $\{x_1, y_3\}$ \\
2 & $s \to x_2 \to y_1 \to t$ & 3 & 加入 $\{x_2, y_1\}$ \\
3 & $s \to x_3 \to y_2 \to t$ & 3 & 加入 $\{x_3, y_2\}$ \\
4 & $s \to x_4 \to y_2 \to x_3 \to y_1 \to x_2 \to y_4 \to t$ & 7 &
加入 $\{x_4, y_2\}$、$\{x_3, y_1\}$、$\{x_2, y_4\}$；\\
& & & 拿掉 $\{x_3, y_2\}$、$\{x_2, y_1\}$ \\
\bottomrule
\end{tabular}
\end{center}

前三輪和「依序挑第一堂沒人帶的課」完全一樣。 第 4 輪開始時（圖~\ref{fig:flw-tanet}(b)），
$s$ 只剩 $s \to x_4$ 一條順向邊；已匹配的邊 $x_i \to y_j$ 都滿了，換成反向邊
$y_j \to x_i$；已匹配的課 $y_1, y_2, y_3$ 到 $t$ 的邊也滿了，只有 $y_4$、$y_5$ 還連得到
$t$。 BFS 從 $x_4$ 只能走到 $y_2$，再沿反向邊回到 $y_2$ 現在的助教 $x_3$，$x_3$ 走到他另
一堂有空的 $y_1$，再回到 $y_1$ 的助教 $x_2$，$x_2$ 走到沒人帶的 $y_4$，最後到 $t$。

\begin{center}
\begin{tikzpicture}[x=0.95cm,y=0.95cm]
\begin{scope}[xshift=0cm]
\node[flw-v] (Ps) at (-2, 0) {$s$};
\node[flw-v] (Pt) at (4.4, 0) {$t$};
\node[flw-v] (Px1) at (0, 1.5) {$x_1$};
\node[flw-v] (Px2) at (0, 0.5) {$x_2$};
\node[flw-v] (Px3) at (0, -0.5) {$x_3$};
\node[flw-v] (Px4) at (0, -1.5) {$x_4$};
\node[flw-v] (Py1) at (2.4, 2) {$y_1$};
\node[flw-v] (Py2) at (2.4, 1) {$y_2$};
\node[flw-v] (Py3) at (2.4, 0) {$y_3$};
\node[flw-v] (Py4) at (2.4, -1) {$y_4$};
\node[flw-v] (Py5) at (2.4, -2) {$y_5$};
\draw[flw-a] (Ps) -- (Px1);
\draw[flw-a] (Ps) -- (Px2);
\draw[flw-a] (Ps) -- (Px3);
\draw[flw-a] (Ps) -- (Px4);
\draw[flw-a] (Px1) -- (Py3);
\draw[flw-a] (Px1) -- (Py5);
\draw[flw-a] (Px2) -- (Py1);
\draw[flw-a] (Px2) -- (Py2);
\draw[flw-a] (Px2) -- (Py4);
\draw[flw-a] (Px2) -- (Py5);
\draw[flw-a] (Px3) -- (Py1);
\draw[flw-a] (Px3) -- (Py2);
\draw[flw-a] (Px4) -- (Py2);
\draw[flw-a] (Py1) -- (Pt);
\draw[flw-a] (Py2) -- (Pt);
\draw[flw-a] (Py3) -- (Pt);
\draw[flw-a] (Py4) -- (Pt);
\draw[flw-a] (Py5) -- (Pt);
\node[flw-t] at (1.2, -2.75) {(a) $G'$：每條邊容量都是 $1$};
\end{scope}
\begin{scope}[xshift=8.0cm]
\node[flw-v] (Qs) at (-2, 0) {$s$};
\node[flw-v] (Qt) at (4.4, 0) {$t$};
\node[flw-v] (Qx1) at (0, 1.5) {$x_1$};
\node[flw-v] (Qx2) at (0, 0.5) {$x_2$};
\node[flw-v] (Qx3) at (0, -0.5) {$x_3$};
\node[flw-v] (Qx4) at (0, -1.5) {$x_4$};
\node[flw-v] (Qy1) at (2.4, 2) {$y_1$};
\node[flw-v] (Qy2) at (2.4, 1) {$y_2$};
\node[flw-v] (Qy3) at (2.4, 0) {$y_3$};
\node[flw-v] (Qy4) at (2.4, -1) {$y_4$};
\node[flw-v] (Qy5) at (2.4, -2) {$y_5$};
\draw[flw-a, flw-rev] (Qx1) -- (Qs);
\draw[flw-a, flw-rev] (Qx2) -- (Qs);
\draw[flw-a, flw-rev] (Qx3) -- (Qs);
\draw[flw-a, flw-aug] (Qs) -- (Qx4);
\draw[flw-a, flw-rev] (Qy3) -- (Qx1);
\draw[flw-a] (Qx1) -- (Qy5);
\draw[flw-a, flw-rev, flw-aug] (Qy1) -- (Qx2);
\draw[flw-a] (Qx2) -- (Qy2);
\draw[flw-a, flw-aug] (Qx2) -- (Qy4);
\draw[flw-a] (Qx2) -- (Qy5);
\draw[flw-a, flw-aug] (Qx3) -- (Qy1);
\draw[flw-a, flw-rev, flw-aug] (Qy2) -- (Qx3);
\draw[flw-a, flw-aug] (Qx4) -- (Qy2);
\draw[flw-a, flw-rev] (Qt) -- (Qy1);
\draw[flw-a, flw-rev] (Qt) -- (Qy2);
\draw[flw-a, flw-rev] (Qt) -- (Qy3);
\draw[flw-a, flw-aug] (Qy4) -- (Qt);
\draw[flw-a] (Qy5) -- (Qt);
\node[flw-t] at (1.2, -2.75) {(b) 第 4 輪的剩餘網路};
\end{scope}
\end{tikzpicture}\par\nopagebreak
\diagcaption{(a) 助教分配轉成的流網路 $G'$，邊的方向一律是 $s \to L \to R \to t$。 (b) 前三
輪做完之後的剩餘網路，虛線是反向邊（已匹配的邊和用掉的 $s \to x$、$y \to t$）；紅色是第 4
輪的增廣路徑，在 $L$、$R$ 之間來回，交替走未匹配的邊和已匹配的邊。}
\label{fig:flw-tanet}
\end{center}

用分配的話說：$x_4$ 要 $y_2$，於是 $x_3$ 讓出 $y_2$ 改帶 $y_1$，$x_2$ 讓出 $y_1$ 改帶
$y_4$。 每位已經有課的助教還是有課，只是換了一堂，多出來的是 $x_4$。 最後的匹配是
$\{x_1, y_3\}$、$\{x_2, y_4\}$、$\{x_3, y_1\}$、$\{x_4, y_2\}$，四位助教都有課，$y_5$ 沒有助
教。 $|L| = 4$，不可能更多，這就是最大匹配。

\subsection{交錯路徑}
\label{sec:flw-alt}

第 4 輪的路徑去掉頭尾的 $s$、$t$，是 $G$ 裡的路徑 $x_4, y_2, x_3, y_1, x_2, y_4$：第一、
三、五條邊不在匹配裡，第二、四條在。 這種「不在 $M$、在 $M$」交替出現的路徑叫做
\textbf{交錯路徑}（alternating path）；兩端都是未匹配頂點的交錯路徑，叫做（對 $M$ 的）
\textbf{增廣路徑}。

匹配的增廣路徑和 $G'$ 的增廣路徑是同一件事。 $M$ 對應的流，剩餘網路長這樣：未匹配的 $x$
才有 $s \to x$；不在 $M$ 的邊是順向邊 $x \to y$，在 $M$ 的邊只剩反向邊 $y \to x$；未匹配
的 $y$ 才有 $y \to t$。 所以 $G'_f$ 裡從 $s$ 到 $t$ 的路徑，一定是從一個未匹配的 $x$ 出
發，從 $L$ 往 $R$ 走不在 $M$ 的邊，從 $R$ 往 $L$ 走在 $M$ 的邊，最後停在一個未匹配的 $y$。
沿它增廣，等於把路徑上的邊「在 $M$ 的拿掉、不在 $M$ 的加進來」，路徑上不在 $M$ 的邊比在
$M$ 的多一條，匹配就多一條邊。

再加上最大流最小割定理（$G'_f$ 沒有增廣路徑 $\Leftrightarrow$ 流是最大的），得到：

\begin{corollary}
\label{cor:flw-berge}
二分圖的匹配 $M$ 是最大匹配，若且唯若 $G$ 裡沒有對 $M$ 的增廣路徑。
\end{corollary}

這個結論對一般的圖也成立（Berge，1957），只是一般圖的增廣路徑不好找，超出本書範圍。

既然增廣路徑可以直接在 $G$ 上找，就不必真的建出 $G'$。 虛擬碼~\ref{alg:flw-alt} 對每個
$x \in L$ 用 DFS 找一條從 $x$ 出發的增廣路徑：$\mathit{mate}[y]$ 是目前和 $y$ 配對的頂點，
沒有配對時是 \NIL；$\mathit{seen}[y]$ 記錄這次 DFS 已經試過的 $y$，避免繞圈。

\begin{pseudo}[label=alg:flw-alt]{用交錯路徑求二分圖的最大匹配}
\Procedure{Try-Augment}{$x$} \Comment{從 $x \in L$ 找增廣路徑}
    \For{每個 $y \in \mathit{Adj}[x]$}
        \If{$\mathit{seen}[y] = \FALSE$}
            \State $\mathit{seen}[y] \gets \TRUE$
            \If{$\mathit{mate}[y] = \NIL$ \OR \Call{Try-Augment}{$\mathit{mate}[y]$}}
                \State $\mathit{mate}[y] \gets x$
                \State \Return \TRUE
            \EndIf
        \EndIf
    \EndFor
    \State \Return \FALSE
\EndProcedure
\Statex
\Procedure{Bipartite-Matching}{$G$}
    \For{每個 $y \in R$}
        \State $\mathit{mate}[y] \gets \NIL$
    \EndFor
    \State $\mathit{size} \gets 0$
    \For{每個 $x \in L$}
        \For{每個 $y \in R$}
            \State $\mathit{seen}[y] \gets \FALSE$
        \EndFor
        \If{\Call{Try-Augment}{$x$}}
            \State $\mathit{size} \gets \mathit{size} + 1$
        \EndIf
    \EndFor
    \State \Return $\mathit{size}$
\EndProcedure
\end{pseudo}

第 5 行是整個演算法的核心：$y$ 沒人配，路徑就在 $y$ 結束；$y$ 已經配給別人，就遞迴問「那
個人能不能改配別堂」，問得到，路徑就經過 $y$ 和 $\mathit{mate}[y]$ 繼續延伸。 找到之後，
第 6 行在遞迴回來的途中，把路徑上的每個 $y$ 改配給前一個 $x$，也就是把交錯路徑翻轉。

每次呼叫 \proc{Try-Augment} 從一個 $x$ 出發，是一次 DFS，每個 $y$ 最多試一次，花
$O(n + m)$；外層對每個 $x$ 做一次，總時間 $O(|L|(n + m)) = O(nm)$（沒有孤立頂點時）。 一
個 $x$ 這次找不到增廣路徑，之後也不必再試它（觀念題第~\ref{q:flw-once}~題），所以外層迴
圈每個 $x$ 只跑一次就夠。 第~\ref{sec:flw-c}~節最後用一張隨機的大圖比較
過，兩種做法求出的匹配大小相同。

把二分匹配做得更快的是 Hopcroft--Karp 演算法（1973）：每個階段用 BFS 分層，一次找出很多
條互不相交的最短增廣路徑，時間 $O(m\sqrt{n})$。 邊上有權重、要讓總成本最小的指派問題
（assignment problem），可以用匈牙利演算法（Hungarian algorithm）在 $O(n^3)$ 時間解決。
這兩個演算法本書都不展開，延伸閱讀列了參考資料；觀念題第~\ref{q:flw-konig}~題用最大流最
小割證明二分圖的另一個重要性質（König 定理）。

\section{C 實作}
\label{sec:flw-c}

\subsection{鄰接矩陣的 Edmonds--Karp}

本節的程式用一個 $n \times n$ 的陣列 \texttt{res} 存剩餘容量 $c_f(u, v)$，另一個陣列
\texttt{cap} 存原本的容量，就是虛擬碼~\ref{alg:flw-ek} 的直接翻譯。 選鄰接矩陣有三個理由：

\begin{itemize}
\item \textbf{反向邊不用另外存。} $(u, v)$ 的反向邊就是 \texttt{res[v][u]}，增廣時一格減、
一格加。
\item \textbf{反平行邊不用轉換。} 如果 $u \to v$、$v \to u$ 兩條邊都有，\texttt{res[u][v]}
一開始是 $c(u, v)$，之後自動變成「$c(u, v)$ 減掉 $u$ 往 $v$ 送的、加上 $v$ 往 $u$ 送的」，
兩個方向合併成一格。 第~\ref{sec:flw-model}~節加新頂點的轉換在這裡可以省掉。 印流量時取
$c(u, v) - c_f(u, v)$，只印正的：兩個方向互相抵消後，只有一個方向留下流量。
\item \textbf{鄰居順序和手算一樣。} BFS 掃一整列 $v = 0, 1, \ldots, n - 1$，頂點編號照輸入
的順序，和手算「依固定順序檢查鄰居」的規則相同，輸出可以逐列對照。
\end{itemize}

代價是空間 $\Theta(n^2)$，而且每次 BFS 要掃整個矩陣，花 $\Theta(n^2)$ 而不是 $O(n + m)$，
總時間變成 $O(nm) \cdot \Theta(n^2) = O(n^3 m)$。 幾百個頂點以內沒有問題，更大的稀疏圖要
換成下一小節的邊串列。

\begin{cbookrecap}[第~6~章 6.11 節、第~19~章 19.3 節]
二維陣列在記憶體裡是一列接一列排的，第 $i$ 列第 $j$ 個元素的位移是 \texttt{i * C + j}。
大小執行時才知道的二維表，用 \texttt{malloc} 配置 \texttt{n * n} 個元素的一塊，自己算位移；
這裡用巨集 \texttt{RES(g, u, v)} 包起來，寫起來和 \texttt{res[u][v]} 一樣。 \texttt{calloc}
配置後全部設成 $0$，正好當成「沒有邊」；\texttt{memcpy} 一次把整張容量表複製成剩餘容量
表。 每一塊 \texttt{malloc}、\texttt{calloc} 來的記憶體都要 \texttt{free}。
\end{cbookrecap}

輸入有好幾張網路，一直讀到檔案結束。 每張網路第一行是 $n$、$m$，第二行是 $n$ 個頂點的名
字（最多三個字元，第一個是 $s$、最後一個是 $t$），接著是 $m$ 條邊，每條寫成兩個頂點的名字
和容量，一行可以寫好幾條；同一對頂點重複出現時容量相加。 每一輪印一列：BFS 求出的距離（依頂點的順序，走不到印 \texttt{-}）、
增廣路徑、瓶頸和增廣之後的值。 最後印出有流的邊和最小割。 輸入的三張網路依序是物流網
路、一張有反平行邊的小網路，以及助教分配的 $G'$。

\begin{example}[label=ex:flw-maxflow]{c}{(maxflow.c) 鄰接矩陣的 Edmonds--Karp、最小割與二分匹配}
#include <stdio.h>
#include <stdlib.h>
#include <string.h>

#define MAXN 64

typedef struct {
    int n;
    char name[MAXN][4];         // 頂點的名字，0 是 s，n-1 是 t
    int *cap;                   // cap[u*n+v]：原本的容量 c(u, v)
    int *res;                   // res[u*n+v]：剩餘容量 c_f(u, v)
} Network;

#define CAP(g, u, v) ((g)->cap[(u) * (g)->n + (v)])
#define RES(g, u, v) ((g)->res[(u) * (g)->n + (v)])

int find_vertex(const Network *g, const char *s)
{
    for (int i = 0; i < g->n; i++)
        if (strcmp(g->name[i], s) == 0)
            return i;
    return -1;
}

// 在剩餘網路上從 s 做 BFS（虛擬碼 21.1），鄰居照編號由小到大看。
// dist[v] = -1 表示走不到。回傳 1：走得到 t。
int bfs(const Network *g, int *dist, int *pi, int *queue)
{
    int n = g->n, head = 0, tail = 0;
    for (int v = 0; v < n; v++)
        dist[v] = pi[v] = -1;
    dist[0] = 0;
    queue[tail++] = 0;
    while (head < tail) {
        int u = queue[head++];
        for (int v = 0; v < n; v++)
            if (dist[v] == -1 && RES(g, u, v) > 0) {
                dist[v] = dist[u] + 1;
                pi[v] = u;
                queue[tail++] = v;
            }
    }
    return dist[n - 1] != -1;
}

// Edmonds-Karp（虛擬碼 26.2）。每一輪印出距離、增廣路徑、瓶頸和 |f|。
// 結束時 dist 是最後一次 BFS 的結果，走得到的頂點就是最小割的 S。
long edmonds_karp(Network *g, int *dist, int *pi, int *queue)
{
    int n = g->n, t = n - 1;
    long value = 0;
    memcpy(g->res, g->cap, (size_t)n * n * sizeof *g->res);
    for (;;) {
        int found = bfs(g, dist, pi, queue);
        for (int v = 0; v < n; v++) {
            if (v > 0)
                putchar(' ');
            if (dist[v] == -1)
                putchar('-');
            else
                printf("%d", dist[v]);
        }
        if (!found) {
            printf(" | no path\n");
            return value;
        }
        int b = RES(g, pi[t], t);                 // 瓶頸 c_f(p)
        for (int v = t; v != 0; v = pi[v])
            if (RES(g, pi[v], v) < b)
                b = RES(g, pi[v], v);
        int len = 0;
        for (int v = t; v != 0; v = pi[v]) {      // 沿路徑增廣
            RES(g, pi[v], v) -= b;
            RES(g, v, pi[v]) += b;
            queue[len++] = v;                     // 倒著存路徑
        }
        value += b;
        printf(" | s");
        while (len > 0)
            printf(" %s", g->name[queue[--len]]);
        printf(" | +%d=%ld\n", b, value);
    }
}

// 印出有流的邊（f/c），再印出最小割。
void report(const Network *g, const int *dist)
{
    int n = g->n, k = 0;
    for (int u = 0; u < n; u++)
        for (int v = 0; v < n; v++) {
            int f = CAP(g, u, v) - RES(g, u, v);  // 反平行邊可能為負
            if (CAP(g, u, v) == 0 || f <= 0)
                continue;
            printf("%s%s->%s %d/%d", k % 3 == 0 ? "" : ", ",
                   g->name[u], g->name[v], f, CAP(g, u, v));
            if (++k % 3 == 0)
                printf("\n");
        }
    if (k % 3 != 0)
        printf("\n");
    long c = 0;                                   // 割的容量 c(S, T)
    for (int u = 0; u < n; u++)
        for (int v = 0; v < n; v++)
            if (dist[u] != -1 && dist[v] == -1)   // u 在 S、v 在 T
                c += CAP(g, u, v);
    printf("min cut S = {s");
    for (int v = 1; v < n; v++)
        if (dist[v] != -1)
            printf(" %s", g->name[v]);
    printf("}, capacity %ld\ncut edges:", c);
    for (int u = 0; u < n; u++)
        for (int v = 0; v < n; v++)
            if (dist[u] != -1 && dist[v] == -1 && CAP(g, u, v) > 0)
                printf(" %s->%s", g->name[u], g->name[v]);
    printf("\n");
}

int main(void)
{
    int n, m, id = 0;
    while (scanf("%d %d", &n, &m) == 2) {
        if (n < 2 || n > MAXN)
            return 1;
        Network g = {.n = n};
        for (int i = 0; i < n; i++)
            if (scanf("%3s", g.name[i]) != 1)
                return 1;
        g.cap = calloc((size_t)n * n, sizeof *g.cap);
        g.res = malloc((size_t)n * n * sizeof *g.res);
        int *buf = malloc(3 * (size_t)n * sizeof *buf);
        if (g.cap == NULL || g.res == NULL || buf == NULL)
            return 1;
        for (int i = 0; i < m; i++) {
            char a[4], b[4];
            int c;
            if (scanf("%3s %3s %d", a, b, &c) != 3)
                return 1;
            int u = find_vertex(&g, a), v = find_vertex(&g, b);
            if (u < 0 || v < 0 || u == v || c < 0)
                return 1;
            CAP(&g, u, v) += c;                   // 重邊的容量相加
        }
        printf("network %d\n", ++id);
        long value = edmonds_karp(&g, buf, buf + n, buf + 2 * n);
        printf("max flow %ld\n", value);
        report(&g, buf);
        free(buf);
        free(g.cap);
        free(g.res);
    }
    return 0;
}
\end{example}
\iosample{7 11 \\ s a b c d e t \\ s a 5 s b 11 s c 7 \\ b a 2 a d 12 b d 5 \\ b e 6 c e 7 d e 2 \\ d t 12 e t 9 \\ 4 6 \\ s u v t \\ s u 3 s v 2 u v 6 \\ v u 4 u t 2 v t 3 \\ 11 18 \\ s x1 x2 x3 x4 y1 y2 y3 y4 y5 t \\ s x1 1 s x2 1 s x3 1 \\ s x4 1 x1 y3 1 x1 y5 1 \\ x2 y1 1 x2 y2 1 x2 y4 1 \\ x2 y5 1 x3 y1 1 x3 y2 1 \\ x4 y2 1 y1 t 1 y2 t 1 \\ y3 t 1 y4 t 1 y5 t 1}{network 1 \\ 0 1 1 1 2 2 3 | s a d t | +5=5 \\ 0 2 1 1 2 2 3 | s b d t | +5=10 \\ 0 2 1 1 3 2 3 | s b e t | +6=16 \\ 0 4 3 1 4 2 3 | s c e t | +3=19 \\ 0 4 3 1 5 2 6 | s c e b a d t | +2=21 \\ 0 - 3 1 - 2 - | no path \\ max flow 21 \\ s->a 5/5, s->b 11/11, s->c 5/7 \\ a->d 7/12, b->a 2/2, b->d 5/5 \\ b->e 4/6, c->e 5/7, d->t 12/12 \\ e->t 9/9 \\ min cut S = \{s b c e\}, capacity 21 \\ cut edges: s->a b->a b->d e->t \\ network 2 \\ 0 1 1 2 | s u t | +2=2 \\ 0 1 1 2 | s v t | +2=4 \\ 0 1 2 3 | s u v t | +1=5 \\ 0 - - - | no path \\ max flow 5 \\ s->u 3/3, s->v 2/2, u->v 1/6 \\ u->t 2/2, v->t 3/3 \\ min cut S = \{s\}, capacity 5 \\ cut edges: s->u s->v \\ network 3 \\ 0 1 1 1 1 2 2 2 2 2 3 | s x1 y3 t | +1=1 \\ 0 5 1 1 1 2 2 4 2 2 3 | s x2 y1 t | +1=2 \\ 0 5 3 1 1 2 2 4 4 4 3 | s x3 y2 t | +1=3 \\ 0 9 5 3 1 4 2 8 6 6 7 | s x4 y2 x3 y1 x2 y4 t | +1=4 \\ 0 - - - - - - - - - - | no path \\ max flow 4 \\ s->x1 1/1, s->x2 1/1, s->x3 1/1 \\ s->x4 1/1, x1->y3 1/1, x2->y4 1/1 \\ x3->y1 1/1, x4->y2 1/1, y1->t 1/1 \\ y2->t 1/1, y3->t 1/1, y4->t 1/1 \\ min cut S = \{s\}, capacity 4 \\ cut edges: s->x1 s->x2 s->x3 s->x4}

第一張網路的前六列和第~\ref{sec:flw-ek-trace}~節的表逐列相同，最後一列距離是 \texttt{-}
的頂點就是 $T$；流量那四列是圖~\ref{fig:flw-final}（沒有流的 $d \to e$ 不印），最後兩列是
第~\ref{sec:flw-mincut}~節的最小割。 第二張網路的 $u \to v$、$v \to u$ 是一對反平行邊，程
式不需要任何轉換就算出最大流 $5$，流量只留在 $u \to v$ 一個方向。 第三張是第~\ref{sec:flw-match-trace}~節
的四輪增廣，$x \to y$ 有流的四條邊就是最大匹配。 注意第三張的 BFS 找到 $t$ 之後沒有停，
繼續沿反向邊從 $t$ 走回已匹配的 $y$，所以距離列裡有些 $y$ 的距離比 $t$ 還大；這不影響找到
的路徑，只是為了印出完整的距離。 這張的最小割是 $S = \{s\}$：四位助教都有課，$s$ 出去的
邊全滿。

\subsection{邊串列：大而稀疏的圖}

頂點有幾萬個時，矩陣太浪費。 這時把每條邊存成一對\textbf{弧}（arc）：第 $2i$ 條弧是原本
的邊 $u \to v$，剩餘容量從 $c(u, v)$ 開始；第 $2i + 1$ 條是它的反向弧 $v \to u$，從 $0$ 開
始。 兩條弧的編號只差最低位元，\texttt{k \^{} 1} 就是弧 \texttt{k} 的另一半，增廣時
\texttt{res[k]} 減、\texttt{res[k \^{} 1]} 加。 每個頂點的弧用陣列串成串列
（\texttt{first}、\texttt{next}，和第~21~章的鄰接串列是同一個結構，只是指標換成陣列下標），
BFS 只掃真正存在的弧，一輪 $O(n + m)$，總時間就是定理~\ref{thm:flw-ek} 的 $O(nm^2)$。 BFS
記下走進每個頂點的弧 \texttt{pe[v]}；弧的起點是另一半的終點 \texttt{to[pe[v] \^{} 1]}，沿著
它就能從 $t$ 走回 $s$。

\begin{code}{c}{邊串列的 Edmonds--Karp（片段）}
typedef struct {
    int n, arcs;
    int *first, *last;          // 頂點的第一條、最後一條弧
    int *to, *res, *next;       // 弧的終點、剩餘容量、下一條
} FlowNet;

// 把弧 u -> v 接在 u 的串列尾端；-1 表示沒有弧
void push_arc(FlowNet *g, int u, int v, int c)
{
    int k = g->arcs++;
    g->to[k] = v;
    g->res[k] = c;
    g->next[k] = -1;
    if (g->last[u] == -1)
        g->first[u] = k;
    else
        g->next[g->last[u]] = k;
    g->last[u] = k;
}

void add_edge(FlowNet *g, int u, int v, int c)
{
    push_arc(g, u, v, c);       // 弧 2i
    push_arc(g, v, u, 0);       // 弧 2i+1
}

// max_flow 的主迴圈；pe、queue 是長度 n 的陣列
for (;;) {
    for (int v = 0; v < n; v++)
        pe[v] = -2;                         // -2：還沒發現
    pe[s] = -1;
    int head = 0, tail = 0;
    queue[tail++] = s;
    while (head < tail && pe[t] == -2) {    // 找到 t 就可以停
        int u = queue[head++];
        for (int k = g->first[u]; k != -1; k = g->next[k])
            if (g->res[k] > 0 && pe[g->to[k]] == -2) {
                pe[g->to[k]] = k;
                queue[tail++] = g->to[k];
            }
    }
    if (pe[t] == -2)
        break;
    int b = g->res[pe[t]];
    for (int v = t; v != s; v = g->to[pe[v] ^ 1])
        if (g->res[pe[v]] < b)
            b = g->res[pe[v]];
    for (int v = t; v != s; v = g->to[pe[v] ^ 1]) {
        g->res[pe[v]] -= b;
        g->res[pe[v] ^ 1] += b;
    }
    flow += b;
}
\end{code}

用這個結構求二分匹配時，照第~\ref{sec:flw-reduce}~節建 $G'$，先加 $s \to x$、再加
$x \to y$、最後加 $y \to t$。 弧接在串列尾端，輸入的邊依 $(x, y)$ 排好時，每個頂點的弧剛好
依對方的編號由小到大，BFS 的順序和手算一樣。 結束後，原本的邊 $x \to y$（偶數號的弧）剩
餘容量是 $0$，表示上面有流，這條邊就在匹配裡。 拿左右各 $2000$ 個頂點、$8000$ 條邊的隨機
二分圖測試，$G'$ 有 $4002$ 個頂點：用上一小節的矩陣，兩張表要 $2 \times 4002^2$ 個
\texttt{int}，大約 $128$ MB，每次 BFS 要掃一千六百萬格；用邊串列只要幾萬個弧，和虛擬碼~\ref{alg:flw-alt}
的寫法算出的匹配大小相同。

\section{本章常見錯誤}
\label{sec:flw-pitfalls}

\begin{center}
\small
\begin{tabular}{>{\raggedright\arraybackslash}p{0.27\linewidth}>{\raggedright\arraybackslash}p{0.48\linewidth}>{\raggedright\arraybackslash}p{0.15\linewidth}}
\toprule
錯誤 & 例子 & 見 \\
\midrule
剩餘網路漏了反向邊 & 只畫「還能送多少」，沒畫「能退回多少」；物流網路第 5 輪的路徑
$s \to c \to e \to b \to \cdots$ 就找不到，停在 $19$ & 第~\ref{sec:flw-residual}~節 \\
反向邊的容量寫錯 & 把 $(v, u)$ 的剩餘容量寫成 $c(u, v)$ 或 $c(u, v) - f(u, v)$；正確是
$f(u, v)$ & 定義~\ref{def:flw-residual} \\
沿反向邊增廣時把流量加上去 & 路徑經過反向邊 $e \to b$，應該是 $f(b, e)$ 減 $2$，卻寫成
$f(e, b) = 2$ & 引理~\ref{lem:flw-augment} \\
割的容量算錯方向 & 把從 $T$ 指回 $S$ 的 $d \to e$ 也加進去（或減掉） &
第~\ref{sec:flw-mfmc}~節 \\
淨流量忘了減回流 & $S = \{s, a, d\}$ 時沒減 $f_1(b, d) = 5$，算出 $21$ 而不是 $16$ &
引理~\ref{lem:flw-cutflow} \\
以為最小割、最大流都只有一個 & 只交剩餘網路找到的 $\{s, b, c, e\}$，漏了 $V - \{t\}$ &
第~\ref{sec:flw-mincut}~節 \\
找不到「每條邊都沒滿的路徑」就停 & 圖~\ref{fig:flw-greedy} 停在 $2$；要在剩餘網路上找才
對 & 第~\ref{sec:flw-residual}~節 \\
把 Ford--Fulkerson 寫成多項式時間 & $O(m\,|f^{*}|)$ 和容量的數值成正比，是擬多項式 &
第~\ref{sec:flw-bad}~節 \\
把增廣次數當成時間 & Edmonds--Karp 增廣 $O(nm)$ 次，每次 $O(m)$，共 $O(nm^2)$；用矩陣
做 BFS 是 $O(n^3 m)$ & 第~\ref{sec:flw-ek-proof}~節 \\
二分圖的邊畫成雙向 & $x \to y$ 和 $y \to x$ 都放容量 $1$，流可以在 $L$、$R$ 之間來回，對
不回匹配 & 第~\ref{sec:flw-reduce}~節 \\
極大匹配當成最大匹配 & 依序貪心配到不能再加，就說是最大匹配 & 第~\ref{sec:flw-match}~節 \\
\bottomrule
\end{tabular}
\end{center}

\section{觀念題}
\label{sec:flw-exercises}

題型和研究所考試、期中期末考的筆試相近。 沒有特別說明時，Edmonds--Karp 的 BFS 依題目列出
頂點的順序檢查鄰居。 「物流網路」指圖~\ref{fig:flw-net}(a)。 解答在本章最後。

\begin{enumerate}
\item \label{q:flw-tf}（選擇，可能不只一個答案）下列敘述哪些正確？
\begin{enumerate}
\item 任何一個流的值，都不超過任何一個割的容量。
\item 割 $(S, T)$ 的容量，是從 $S$ 到 $T$ 的邊的容量和，減掉從 $T$ 到 $S$ 的邊的容量和。
\item 若剩餘網路 $G_f$ 裡沒有從 $s$ 到 $t$ 的路徑，$f$ 是最大流。
\item 一張流網路的最大流只有一個。
\item 容量都是整數時，每一個最大流的每條邊流量都是整數。
\item Edmonds--Karp 每一輪找到的增廣路徑，邊數不會比前一輪少。
\item 增廣之後，這條路徑的關鍵邊不會出現在新的剩餘網路裡。
\end{enumerate}

\item \label{q:flw-check}（計算）下表是物流網路上的三個函數，每一欄是一條邊的值。
\begin{center}
\small
\begin{tabular}{c|ccccccccccc}
 & $sa$ & $sb$ & $sc$ & $ba$ & $ad$ & $bd$ & $be$ & $ce$ & $de$ & $dt$ & $et$ \\
\hline
$f_A$ & 5 & 8 & 6 & 2 & 7 & 3 & 3 & 6 & 1 & 9 & 10 \\
$f_B$ & 5 & 8 & 4 & 2 & 7 & 3 & 3 & 4 & 1 & 9 & 8 \\
$f_C$ & 4 & 10 & 7 & 2 & 6 & 5 & 4 & 7 & 0 & 11 & 9 \\
\end{tabular}
\end{center}
(a) 哪些是流？ 不是流的，指出違反了哪一條規則、在哪裡。 (b) 對是流的那一個，寫出它的值，
列出剩餘網路的所有邊和剩餘容量，並用 BFS（順序 $s, a, b, c, d, e, t$）找出一條最短的增廣
路徑和瓶頸。

\item \label{q:flw-trace}（追蹤）流網路的頂點依序是 $s, 1, 2, 3, 4, t$，邊和容量是
$s \to 1$：$5$、$s \to 2$：$6$、$1 \to 2$：$2$、$1 \to 3$：$4$、$1 \to 4$：$9$、
$2 \to 3$：$2$、$2 \to 4$：$2$、$4 \to 3$：$3$、$3 \to t$：$4$、$4 \to t$：$8$。 執行
Edmonds--Karp，每一輪寫出各頂點的距離、增廣路徑、瓶頸和 $|f|$；最後寫出每條邊的流量和一
個最小割。

\item \label{q:flw-allcuts}（計算與證明）證明：$f$ 是最大流時，割 $(S, T)$ 是最小割，若且
唯若在 $G_f$ 裡沒有從 $S$ 指向 $T$ 的邊。 用這個事實找出物流網路所有的最小割。

\item \label{q:flw-smallest}（證明）$f$ 是最大流，$R$ 是 $G_f$ 裡從 $s$ 走得到的頂點。 證
明對每一個最小割 $(S, T)$，都有 $R \subseteq S$。 也就是說，第~\ref{sec:flw-mincut}~節
找到的最小割，源點那一側是所有最小割裡最小的。

\item \label{q:flw-plus1}（證明或反例）$(S, T)$ 是 $G$ 的一個最小割。 下列敘述對的請證
明，錯的請給反例。
\begin{enumerate}
\item 每條邊的容量都乘以同一個正數 $k$，$(S, T)$ 仍是最小割。
\item 每條邊的容量都加 $1$，$(S, T)$ 仍是最小割。
\item 每條邊的容量都加 $1$，最大流的值至少增加 $1$。
\end{enumerate}

\item \label{q:flw-bad}（計算）把圖~\ref{fig:flw-bad} 的四個 $100$ 換成正整數 $K$。 (a) Ford--Fulkerson
在最壞的路徑選擇下要增廣幾次？ 寫出造成最壞情況的選法。 (b) Edmonds--Karp（順序
$s, u, v, t$）增廣幾次？ 每一次的路徑是什麼？ (c) $K = 10^9$ 時，(a) 的次數和輸入的長度（位
元數）是什麼關係？

\item \label{q:flw-time}（計算）用 $n$、$m$ 和最大流的值 $F$ 寫出下列做法的時間上界（容量都
是整數，圖滿足本章的兩個假設）。
\begin{enumerate}
\item Ford--Fulkerson，用 DFS 在鄰接串列上找增廣路徑。
\item Edmonds--Karp，用鄰接串列（或第~\ref{sec:flw-c}~節的邊串列）。
\item Edmonds--Karp，用鄰接矩陣，每次 BFS 掃整個矩陣。
\item 二分圖最大匹配，在 $G'$ 上用 Ford--Fulkerson。
\end{enumerate}

\item \label{q:flw-vertex}（計算）物流網路的轉運站 $b$ 改建，一晚最多只能處理 $8$ 板（進
出 $b$ 的量都不超過 $8$），其他不變。 畫出轉換後的網路，求新的最大流的值，並給一個最小割
證明它是最大的。

\item \label{q:flw-disjoint}（設計與證明）$G$ 是有向圖，$s \ne t$ 是兩個頂點。 兩條路徑如
果沒有共用的邊，就說它們\textbf{邊互斥}（edge-disjoint）。 把每條邊的容量都設成 $1$，證
明：(a) 從 $s$ 到 $t$ 最多有幾條兩兩邊互斥的路徑，就等於這個網路的最大流的值；(b) 這個值也
等於「最少要刪掉幾條邊，才能讓 $s$ 走不到 $t$」。

\item \label{q:flw-mtrace}（追蹤）二分圖 $L = \{x_1, \ldots, x_5\}$、
$R = \{y_1, \ldots, y_5\}$，邊是：$x_1$ 連 $y_2, y_3$；$x_2$ 連 $y_5$；$x_3$ 連 $y_2$；$x_4$
連 $y_2, y_3, y_4, y_5$；$x_5$ 連 $y_2, y_5$。 (a) 在 $G'$ 上執行 Edmonds--Karp（順序
$s, x_1, \ldots, x_5, y_1, \ldots, y_5, t$），寫出每一輪的增廣路徑和最後的匹配。 (b) 改用虛擬碼~\ref{alg:flw-alt}
（$x$、$y$ 都依編號），寫出每次呼叫 \proc{Try-Augment} 的結果。 (c) 說明為什麼不可能有 $5$
條邊的匹配。

\item \label{q:flw-konig}（證明）無向圖的\textbf{頂點覆蓋}（vertex cover）是一個頂點的集合
$C$，每條邊至少有一個端點在 $C$ 裡。 證明 König 定理：二分圖裡，最大匹配的邊數等於最小頂
點覆蓋的頂點數。 （提示：把 $G'$ 中間那些邊的容量改成 $n + 1$，看最小割。） 用你的證明，
找出第~\ref{q:flw-mtrace}~題那張圖的一個最小頂點覆蓋。

\item \label{q:flw-once}（證明）$M$ 是二分圖的匹配，$x \in L$ 未匹配，而且 $G$ 裡沒有從
$x$ 出發、對 $M$ 的增廣路徑。 $P$ 是對 $M$ 的另一條增廣路徑，$M'$ 是沿 $P$ 翻轉之後的匹配。
證明對 $M'$ 也沒有從 $x$ 出發的增廣路徑。 （這說明虛擬碼~\ref{alg:flw-alt} 每個 $x$ 只試
一次就夠。）

\item \label{q:flw-maximal}（證明）$M$ 是任意一張圖（不一定是二分圖）的極大匹配，$M^{*}$
是最大匹配。 證明 $|M^{*}| \le 2|M|$，並給一個等號成立的例子。
\end{enumerate}

\section{實作練習}
\label{sec:flw-practice}

Aki-Judge 目前還沒有和本章對應的題目，請用觀念題和第~\ref{sec:flw-c}~節的程式練習。
例如把觀念題第~\ref{q:flw-trace}~題的網路改寫成範例~\ref{ex:flw-maxflow} 的輸入格式，和自
己的手算對照；或把第~\ref{q:flw-mtrace}~題的二分圖照第~\ref{sec:flw-reduce}~節改成 $G'$，交給同一支程
式。 想練邊串列的話，把第~\ref{sec:flw-c}~節的片段補成完整的程式，用同樣的輸入對答案。 也可
以替範例程式加上頂點容量（第~\ref{sec:flw-model}~節的拆點），用第~\ref{q:flw-vertex}~題
檢查。

\section{讀完你應該會}
\label{sec:flw-checklist}

\begin{itemize}
\item 寫出流網路、流、流的值、剩餘容量、剩餘網路、增廣路徑、割的定義，並把反平行邊、多
個源點匯點、頂點容量轉成標準的流網路。
\item 檢查一個函數是不是流；給定流，畫出剩餘網路（順向邊和反向邊都要有），沿增廣路徑更
新流量。
\item 照指定的順序手算 Edmonds--Karp：每一輪的距離、增廣路徑、瓶頸、$|f|$，以及最後的
流。
\item 證明最大流最小割定理，用最後的剩餘網路找出最小割，並用「$S$ 到 $T$ 全滿、$T$ 到 $S$
全空」驗算。
\item 說明 Ford--Fulkerson 為什麼是擬多項式、Edmonds--Karp 為什麼是 $O(nm^2)$（距離不減、
每條邊當關鍵邊最多 $n/2$ 次）。
\item 把二分圖最大匹配轉成最大流，用交錯路徑解釋每一次增廣，寫出 C 程式（矩陣或邊串列）。
\end{itemize}

\textbf{延伸閱讀}（書末會整理完整的參考資料）

\begin{itemize}
\item Thomas H. Cormen 等，\textit{Introduction to Algorithms}，第 4 版，第 24 章：流網路、
Ford--Fulkerson、Edmonds--Karp 的完整證明和二分匹配。 第 25 章接著講 Hopcroft--Karp 演算
法、穩定婚姻問題和匈牙利演算法。
\item Jon Kleinberg、Éva Tardos，\textit{Algorithm Design}，第 7 章：網路流的應用講得最多，
包括邊互斥路徑、專案選擇、影像分割等「怎麼把問題建成流」的例子。
\item Steven S. Skiena，\textit{The Algorithm Design Manual}，第 3 版，8.5 節：從應用的角度
講二分匹配和網路流，適合想知道實務上什麼問題會用到流的讀者。
\item Antti Laaksonen，\textit{Competitive Programmer's Handbook}，第 20 章：流與割、互斥
路徑、最大匹配和路徑覆蓋，篇幅很短，程式寫法偏向競賽。
\item L. R. Ford, Jr.、D. R. Fulkerson，``Maximal Flow through a Network''，\textit{Canadian
Journal of Mathematics} 8，1956：Ford--Fulkerson 方法和最大流最小割定理的原始論文。
\item Jack Edmonds、Richard M. Karp，``Theoretical Improvements in Algorithmic Efficiency for
Network Flow Problems''，\textit{Journal of the ACM} 19(2)，1972：證明用最短的增廣路徑時，
增廣次數和容量無關。
\item John E. Hopcroft、Richard M. Karp，``An $n^{5/2}$ Algorithm for Maximum Matchings in
Bipartite Graphs''，\textit{SIAM Journal on Computing} 2(4)，1973：一次找很多條最短增廣路
徑的二分匹配演算法。
\end{itemize}

\addcontentsline{toc}{section}{觀念題解答}
\section*{觀念題解答}

\begin{enumerate}
\item (a)、(c)、(f)、(g) 正確。 (a)：系理~\ref{cor:flw-weak}。 (b)：錯，容量只算從 $S$ 到
$T$ 的邊，從 $T$ 到 $S$ 的邊不加也不減。 (c)：定理~\ref{thm:flw-mfmc} 的 (2) $\Rightarrow$
(1)。 (d)：錯。 例如 $s \to a$、$a \to t$、$a \to b$、$b \to t$ 四條邊容量都是 $1$，最大流的
值是 $1$，那一單位可以走 $a \to t$，也可以走 $a \to b \to t$，是兩個不同的最大流。 (e)：
錯，同一個例子，兩條路各走 $1/2$ 也是最大流。 定理~\ref{thm:flw-integral} 只保證「存在」
整數的最大流，而且 Ford--Fulkerson 會找到它。 (f)：由引理~\ref{lem:flw-mono}，
$\delta_f(s, t)$ 不會變小，而增廣路徑的邊數就是 $\delta_f(s, t)$。 (g)：關鍵邊的剩餘容量等
於瓶頸，增廣後變成 $0$。

\item (a) $f_A$ 不是流：$e \to t$ 的流量 $10$ 超過容量 $9$（各頂點都守恆）。 $f_C$ 不是流：
$b$ 流進 $10$、流出 $2 + 5 + 4 = 11$；$e$ 流進 $4 + 7 + 0 = 11$、流出 $9$，兩處不守恆（容
量都沒超過）。 $f_B$ 是流：容量都沒超過；$a$ 流進 $5 + 2$、流出 $7$；$b$ 流進 $8$、流出
$2 + 3 + 3$；$c$ 進出都是 $4$；$d$ 流進 $7 + 3$、流出 $1 + 9$；$e$ 流進 $3 + 4 + 1$、流出
$8$。

(b) $|f_B| = 5 + 8 + 4 = 17$（也等於流進 $t$ 的 $9 + 8$）。 剩餘網路的順向邊：$s \to b$：
$3$、$s \to c$：$3$、$a \to d$：$5$、$b \to d$：$2$、$b \to e$：$3$、$c \to e$：$3$、
$d \to e$：$1$、$d \to t$：$3$、$e \to t$：$1$（$s \to a$、$b \to a$ 已滿）。 反向邊：
$a \to s$：$5$、$b \to s$：$8$、$c \to s$：$4$、$a \to b$：$2$、$d \to a$：$7$、$d \to b$：
$3$、$e \to b$：$3$、$e \to c$：$4$、$e \to d$：$1$、$t \to d$：$9$、$t \to e$：$8$。 BFS 從
$s$ 發現 $b$、$c$（距離 $1$），從 $b$ 發現 $d$、$e$（距離 $2$），從 $d$ 發現 $t$（距離
$3$）；$a$ 要經過 $d \to a$，距離 $3$。 增廣路徑 $s \to b \to d \to t$，剩餘容量
$3, 2, 3$，瓶頸 $2$。 所以 $f_B$ 不是最大流；繼續做下去，會再經過
$s \to b \to e \to t$（$1$）和 $s \to c \to e \to d \to t$（$1$）到達 $21$。

\item 距離依 $s, 1, 2, 3, 4, t$ 的順序：

\begin{center}
\small
\begin{tabular}{c|cccccc|lcc}
\toprule
輪 & $s$ & $1$ & $2$ & $3$ & $4$ & $t$ & 增廣路徑 & 瓶頸 & $|f|$ \\
\midrule
1 & 0 & 1 & 1 & 2 & 2 & 3 & $s \to 1 \to 3 \to t$ & $4$ & $4$ \\
2 & 0 & 1 & 1 & 2 & 2 & 3 & $s \to 1 \to 4 \to t$ & $1$ & $5$ \\
3 & 0 & 3 & 1 & 2 & 2 & 3 & $s \to 2 \to 4 \to t$ & $2$ & $7$ \\
4 & 0 & 3 & 1 & 2 & 4 & 5 & $s \to 2 \to 3 \to 1 \to 4 \to t$ & $2$ & $9$ \\
5 & 0 & $\infty$ & 1 & $\infty$ & $\infty$ & $\infty$ & 走不到 $t$ & & \\
\bottomrule
\end{tabular}
\end{center}

第 1 輪 $1 \to 3$、$3 \to t$ 同時滿了（瓶頸 $4$）；第 2 輪 $s \to 1$ 滿了，之後 $1$ 只能經
過 $2 \to 3$ 和反向邊 $3 \to 1$ 到達，距離變成 $3$。 第 4 輪的路徑經過反向邊 $3 \to 1$：
$1 \to 3$ 上的 $4$ 單位退回 $2$ 單位，改走 $1 \to 4 \to t$，$3$ 少掉的由 $2 \to 3$ 補上。
最後的流量：$s \to 1$：$5$、$s \to 2$：$4$、$1 \to 2$：$0$、$1 \to 3$：$2$、$1 \to 4$：
$3$、$2 \to 3$：$2$、$2 \to 4$：$2$、$4 \to 3$：$0$、$3 \to t$：$4$、$4 \to t$：$5$，值 $9$。
最小割 $S = \{s, 2\}$、$T = \{1, 3, 4, t\}$，割邊 $s \to 1$、$2 \to 3$、$2 \to 4$，容量
$5 + 2 + 2 = 9$；從 $T$ 指回 $S$ 的 $1 \to 2$ 不算，它的流量是 $0$。 這張網路的最小割只有這
一個（可以用第~\ref{q:flw-allcuts}~題的方法驗證）。

\item 先證明。 由引理~\ref{lem:flw-cutflow}，
$|f| = f(S, T) = \sum_{S \to T} f(u, v) - \sum_{T \to S} f(v, u) \le c(S, T)$，等號成立若且
唯若從 $S$ 到 $T$ 的邊全滿、從 $T$ 到 $S$ 的邊全空。 $f$ 是最大流，由最大流最小割定理
$|f|$ 等於最小割的容量，所以 $(S, T)$ 是最小割 $\Leftrightarrow$ $c(S, T) = |f|$
$\Leftrightarrow$ 上面的等號成立。 而「$S \to T$ 的邊全滿、$T \to S$ 的邊全空」正好是說：
沒有從 $S$ 到 $T$ 的順向剩餘邊（全滿），也沒有從 $S$ 到 $T$ 的反向剩餘邊（反向邊來自 $T$
到 $S$ 的流）。

再找物流網路的最小割。 用圖~\ref{fig:flw-final} 的最大流。 $G_f$ 裡從 $s$ 走得到的是
$R = \{s, b, c, e\}$，由第~\ref{q:flw-smallest}~題，每個最小割的 $S$ 都包含 $R$，只剩 $a$、
$d$ 要決定放哪一邊。 $G_f$ 裡有 $a \to d$（剩餘 $5$）和 $d \to a$（反向，$7$），所以 $a$、
$d$ 必須同一邊，否則就有一條從 $S$ 到 $T$ 的剩餘邊。 兩種選法都檢查：$S = R$ 是第~\ref{sec:flw-mincut}~節
的最小割；$S = R \cup \{a, d\} = V - \{t\}$ 時，從 $S$ 出去的剩餘邊只可能指向 $t$，但
$d \to t$、$e \to t$ 都滿了，所以也是最小割，容量 $12 + 9 = 21$。 物流網路剛好有這兩個最
小割。

\item 取任意一個最小割 $(S, T)$。 由第~\ref{q:flw-allcuts}~題，$G_f$ 裡沒有從 $S$ 指向 $T$
的邊。 從 $s \in S$ 出發，沿著 $G_f$ 的邊走，永遠出不了 $S$，所以走得到的頂點都在 $S$ 裡，
$R \subseteq S$。 由定理~\ref{thm:flw-mfmc} 的證明，$(R, V - R)$ 本身也是最小割，所以它是
源點一側最小的那一個。

\item
\begin{enumerate}
\item 對。 每個割的容量都乘以 $k$，大小順序不變，原本最小的還是最小。
\item 錯。 反例：$s \to x$ 容量 $4$；$x \to y_1$、$x \to y_2$、$x \to y_3$ 容量 $1$；
$y_1 \to t$、$y_2 \to t$、$y_3 \to t$ 容量 $5$。 最小割是 $S = \{s, x\}$，容量 $3$（$S = \{s\}$
是 $4$，其他的割至少 $5$）。 每條邊加 $1$ 之後，$S = \{s, x\}$ 變成 $6$，$S = \{s\}$ 只有
$5$，原本的最小割不再是最小。 割邊多的割，加的也多。
\item 對。 由本章的假設，$t$ 在某條從 $s$ 出發的路徑上，這條路徑一定有一條邊從 $S$ 跨到 $T$，
所以每個割至少有一條 $S \to T$ 的邊，容量至少增加 $1$。 新的最小割容量至少是舊的加 $1$，
由最大流最小割定理，最大流的值也至少加 $1$。
\end{enumerate}

\item (a) $2K$ 次。 奇數輪選 $s \to u \to v \to t$，偶數輪選 $s \to v \to u \to t$。 奇數輪
後 $u \to v$ 滿了，留下容量 $1$ 的反向邊 $v \to u$；偶數輪經過它，把那一單位退回，$u \to v$
又空了。 每一輪的瓶頸都是中間那條容量 $1$ 的邊（或它的反向邊），每輪只增加 $1$，最大流是
$2K$，所以要 $2K$ 輪。 這也是 $O(m\,|f^{*}|)$ 這個上界在這張圖上會被達到的情形。 (b) $2$ 次。
第 1 輪 BFS 依序發現 $u$、$v$，再從 $u$ 發現 $t$，路徑 $s \to u \to t$，瓶頸 $K$。 第 2 輪
$s \to u$ 滿了，路徑 $s \to v \to t$，瓶頸 $K$。 (c) 輸入的長度大約是 $4 \lg K$ 個位元
（四個 $K$，再加上常數），次數 $2K$ 是輸入長度的指數，$K = 10^9$ 時要二十億輪；而 Edmonds--Karp 的次數
和 $K$ 無關。

\item (a) 每一輪 $O(n + m) = O(m)$，最多 $F$ 輪：$O(mF)$。 (b) $O(nm)$ 輪，每輪 $O(m)$：
$O(nm^2)$。 (c) $O(nm)$ 輪，每輪 BFS $\Theta(n^2)$：$O(n^3 m)$；矩陣在 $m$ 接近 $n^2$ 的
稠密圖上才划算。 (d) $F \le n/2$，$G'$ 有 $n + 2$ 個頂點、$m + n$ 條邊，每輪 $O(n + m)$：
$O(n(n + m))$。

\item 把 $b$ 拆成 $b_{\mathrm{in}}$、$b_{\mathrm{out}}$，加邊 $b_{\mathrm{in}} \to b_{\mathrm{out}}$，
容量 $8$；$s \to b$ 改成 $s \to b_{\mathrm{in}}$，$b \to a$、$b \to d$、$b \to e$ 改從
$b_{\mathrm{out}}$ 出發。 最大流是 $20$。 一個值 $20$ 的流：$s \to a$：$5$、$s \to c$：$7$、
$s \to b_{\mathrm{in}} \to b_{\mathrm{out}}$：$8$，其中 $b_{\mathrm{out}} \to d$：$5$、
$b_{\mathrm{out}} \to e$：$2$、$b_{\mathrm{out}} \to a$：$1$；$a \to d$：$6$、$c \to e$：
$7$、$d \to e$：$0$、$d \to t$：$11$、$e \to t$：$9$（可以逐點檢查守恆）。 割
$S = \{s, b_{\mathrm{in}}\}$ 的割邊是 $s \to a$、$s \to c$、$b_{\mathrm{in}} \to b_{\mathrm{out}}$，
容量 $5 + 7 + 8 = 20$，所以 $20$ 是最大的。

\item (a) 設最大流的值是 $F$，最多有 $k$ 條兩兩邊互斥的路徑。 \textbf{$k \le F$}：在 $k$ 條路徑
經過的每條邊上放流量 $1$。 路徑沒有共用邊，每條邊最多 $1$；路徑經過中間的頂點時一進一出，
守恆成立（同一條路徑重複經過某個頂點也一樣）；$s$ 每條路徑出去一次，值是 $k$。
\textbf{$F \le k$}：由定理~\ref{thm:flw-integral}，有一個流量全是 $0$ 或 $1$ 的最大流。 從 $s$
出發，每次沿一條流量是 $1$ 的邊往前走。 如果走回已經走過的頂點，就繞出了一個環：把環上
的流量都改成 $0$，每個頂點流進、流出各少 $1$，仍是流，值也不變；丟掉這個環，從環的起點
繼續走。 走到 $v \ne s, t$ 時，$v$ 有一條流進的有流邊，由守恆就有一條流出的有流邊，可以
繼續；回到 $s$ 時，$|f| \ge 1$ 表示 $s$ 流出的比流進的多，也還有出邊。 邊數有限，所以最後
會走到 $t$，得到一條簡單路徑。 把這條路徑上的流量都改成 $0$，剩下的仍是流，值少 $1$。 重
複 $F$ 次，得到 $F$ 條邊互斥的路徑。 (b) 刪掉一個最小割的所有割邊，$s$ 就走不到
$t$，所以最少的刪邊數不超過最小割的容量 $= F$。 反過來，刪掉某個邊集合 $D$ 之後 $s$ 走不到
$t$，令 $S$ 是刪邊後 $s$ 走得到的頂點，則原圖裡所有從 $S$ 指向 $V - S$ 的邊都在 $D$ 裡，
$|D| \ge c(S, V - S) \ge F$。 這個結果叫 Menger 定理（邊的版本）。

\item (a) 依序是 $s \to x_1 \to y_2 \to t$、$s \to x_2 \to y_5 \to t$、$s \to x_4 \to y_3 \to t$，
第 4 輪 $s \to x_3 \to y_2 \to x_1 \to y_3 \to x_4 \to y_4 \to t$。 第 3 輪輪到 $x_3$ 時，它唯一
的 $y_2$ 已經配給 $x_1$，從 $x_3$ 出發的路徑比較長；BFS 先找到從 $x_4$ 出發、長度 $3$ 的
那一條。 第 4 輪：$x_3$ 要 $y_2$，$x_1$ 改配 $y_3$，$x_4$ 改配 $y_4$。 最後的匹配
$\{x_1, y_3\}$、$\{x_2, y_5\}$、$\{x_3, y_2\}$、$\{x_4, y_4\}$，$x_5$ 和 $y_1$ 沒配到。 這時
剩餘網路裡 $s$ 只走得到 $x_5$、$y_2$、$y_5$、$x_3$、$x_2$，到不了 $t$。

(b) $x_1$：$y_2$ 沒人配，配給 $x_1$。 $x_2$：$y_5$ 沒人配，配給 $x_2$。 $x_3$：$y_2$ 配給了
$x_1$，遞迴 \proc{Try-Augment}$(x_1)$：$y_2$ 這次已看過，$y_3$ 沒人配，配給 $x_1$；回來後
$y_2$ 改配 $x_3$。 $x_4$：$y_2$ 屬於 $x_3$，$x_3$ 沒有別的課，失敗；$y_3$ 屬於 $x_1$，$x_1$ 的
$y_2$、$y_3$ 都看過了，失敗；$y_4$ 沒人配，配給 $x_4$。 $x_5$：$y_2$ 屬於 $x_3$，失敗；$y_5$
屬於 $x_2$，$x_2$ 沒有別的課，失敗，傳回 \FALSE。 結果和 (a) 相同，大小 $4$。

(c) $x_2$、$x_3$、$x_5$ 這三位能帶的課都在 $\{y_2, y_5\}$ 裡，三個人兩堂課，至少一個人配不到，
所以最多 $4$ 條。 換個說法：$\{x_1, x_4, y_2, y_5\}$ 碰到每一條邊，匹配的每條邊都要用掉其中
一個不同的頂點，匹配最多 $4$ 條（這就是第~\ref{q:flw-konig}~題的頂點覆蓋）。

\item 匹配 $M$ 和頂點覆蓋 $C$：$M$ 的每條邊至少有一個端點在 $C$ 裡，$M$ 的邊沒有共同端點，
這些端點都不同，所以 $|M| \le |C|$，最大匹配 $\le$ 最小覆蓋。

反方向：把 $G'$ 中間的邊 $x \to y$ 容量改成 $n + 1$，$s \to x$、$y \to t$ 仍是 $1$。 每個 $x$
流進最多 $1$，所以中間邊的流量其實不會超過 $1$；引理~\ref{lem:flw-match-flow} 的證明照樣成
立，最大流的值仍等於最大匹配的邊數 $\nu$。 取一個最小割 $(S, T)$。 割 $\{s\}$ 的容量是 $|L|
\le n$，所以最小割的容量不超過 $n$，不可能包含容量 $n + 1$ 的中間邊：\textbf{沒有一條邊從
$L \cap S$ 指到 $R \cap T$}。 令 $C = (L \cap T) \cup (R \cap S)$。 對任一條邊 $\{x, y\}$，
若 $x \notin C$ 則 $x \in S$，由上一句 $y \in S$，所以 $y \in C$；$C$ 是頂點覆蓋。 割邊只剩
$s \to x$（$x \in L \cap T$）和 $y \to t$（$y \in R \cap S$），容量剛好是 $|C|$。 所以
$|C| = $ 最小割 $= \nu$，最小覆蓋 $\le$ 最大匹配。

第~\ref{q:flw-mtrace}~題那張圖：最後的剩餘網路裡 $s$ 走得到 $S = \{s, x_2, x_3, x_5, y_2, y_5\}$
（中間邊容量改大，只多出從已匹配的 $x$ 走向它伴侶的順向邊；已匹配的 $x$ 只能從伴侶那
裡走到，伴侶本來就在 $S$ 裡，所以走得到的頂點不變），$C = \{x_1, x_4\} \cup \{y_2, y_5\}$，四個頂點，等於最
大匹配的大小。

\item 令 $Z$ 是從 $x$ 出發、沿對 $M$ 的交錯路徑走得到的頂點（第一步走不在 $M$ 的邊）。 從 $x$
沒有增廣路徑，所以 $Z$ 裡 $R$ 的頂點都已匹配，而且它們在 $M$ 裡的伴侶也在 $Z$ 裡。 先證明
$P$ 和 $Z$ 沒有共同的頂點。 $P$ 的兩端是未匹配的 $x' \in L$（$x' \ne x$，否則 $P$ 就是從 $x$
出發的增廣路徑）和未匹配的 $y' \in R$（$y' \notin Z$）。 假設 $P$ 碰到 $Z$，在 $P$ 上取最靠
近 $y'$ 的 $Z$ 頂點 $w$。 若 $w \in L$，$P$ 從 $w$ 往 $y'$ 走的下一條邊不在 $M$，接在到 $w$ 的
交錯路徑後面，下一個頂點也在 $Z$ 裡；若 $w \in R$，$P$ 往 $y'$ 走的下一條邊是 $w$ 在 $M$ 的
邊，它的伴侶也在 $Z$ 裡。 兩種情況都和「$w$ 最靠近 $y'$」矛盾（$w \ne y'$，因為 $y' \notin Z$）。
所以 $P$ 和 $Z$ 不相交，翻轉 $P$ 不會改變任何碰到 $Z$ 的邊在不在匹配裡。 從 $x$ 沿對 $M'$ 的
交錯路徑走，每一步看的都是碰到 $Z$ 的邊，走到的還是 $Z$，$Z$ 裡 $R$ 的頂點在 $M'$ 中仍已匹
配，所以對 $M'$ 也沒有從 $x$ 出發的增廣路徑。

\item $M^{*}$ 的每條邊 $\{u, v\}$ 至少有一個端點已被 $M$ 匹配，否則可以把它加進 $M$，和 $M$
是極大匹配矛盾。 把 $M^{*}$ 的每條邊對應到它一個被 $M$ 匹配的端點；$M^{*}$ 的邊沒有共同端
點，所以這些端點都不同。 被 $M$ 匹配的頂點只有 $2|M|$ 個，所以 $|M^{*}| \le 2|M|$。 等號的例
子：一條路徑 $a - b - c - d$，$M = \{\{b, c\}\}$ 是極大匹配（另外兩條邊都碰到 $b$ 或 $c$），
最大匹配 $\{\{a, b\}, \{c, d\}\}$ 有兩條邊。
\end{enumerate}


% 第六部：其他
\chapter{字串比對}
\label{ch:stringmatching}

% 本章的圖共用的樣式（前綴 stm-）。 比對過程的對齊圖、失敗函數表、KMP 追蹤表和
% 自動機轉移表都由模擬程式照虛擬碼逐步執行後產生，再和本章的 C 程式輸出對過。
\tikzset{
    stm-c/.style={draw=black!55, minimum width=0.5cm, minimum height=0.5cm, inner sep=0pt, font=\ttfamily\small},
    stm-ok/.style={stm-c, fill=coreFill, draw=coreStroke},
    stm-bad/.style={stm-c, fill=optFill, draw=optStroke, thick},
    stm-known/.style={stm-c, fill=greenFill, draw=greenStroke},
    stm-idle/.style={stm-c, draw=black!25, text=black!40},
    stm-lab/.style={anchor=east, font=\small},
    stm-note/.style={font=\scriptsize, text=black!70},
}

生物系的實驗室定序出一條很長的 DNA，幾百萬個鹼基，每個鹼基是 A、C、G、T 其中一
個。 研究助理手上有一段十個鹼基的片段，想知道它在長序列裡出現在哪些位置，重疊的
也要算。 最直接的做法是從長序列的每一個位置開始，把片段一個一個對上去，對不上就換
下一個位置重來。 片段裡有很多重複的鹼基時，這個做法會把同一段長序列看上好幾次。
往下讀之前，先想想：前一個位置已經對上的那幾個鹼基，能不能告訴我們下一個位置不必
再從頭比？

這一章講在一段文字裡找出一個字串所有出現位置的方法。 第~\ref{sec:stm-naive}~節先
寫清楚問題，再分析最直接的樸素比對，它的最壞情況是 $\Theta(nm)$ 這一級。
第~\ref{sec:stm-rk}~節的 Rabin--Karp 演算法把字串看成數字，用第~12~章的雜湊先篩
掉大部分位置，平均很快，但最壞情況沒有比較好。 第~\ref{sec:stm-kmp}~節的 KMP 演算
法事先算好模式自己和自己重疊的情形（失敗函數），比對時文字的每個字元只往前讀一次，
總時間 $\Theta(n + m)$；證明它是線性時間要用第~3~章的平攤分析。
第~\ref{sec:stm-dfa}~節換一個角度，把 KMP 看成一台有限自動機。

C 程式用到《C 程式設計基礎》第~11~章的字串（以 \texttt{'\textbackslash 0'} 結尾的
字元陣列、\texttt{strlen}、\texttt{strncmp}）和第~19~章的 \texttt{malloc}，用到時會
先在框裡複習。

\section{問題與樸素比對}
\label{sec:stm-naive}

\subsection{符號}

\textbf{字母表}（alphabet）$\Sigma$ 是一個有限的字元集合，例如 DNA 的
$\Sigma = \{\mathtt{A}, \mathtt{C}, \mathtt{G}, \mathtt{T}\}$，或是 $26$ 個小寫英文
字母。 \textbf{字串}（string）是由 $\Sigma$ 裡的字元排成的有限序列，長度記成
$|x|$；長度 $0$ 的字串叫\textbf{空字串}，記成 $\varepsilon$。 這一章的字串和陣列一
樣從 $1$ 編號，$x[i..j]$ 是第 $i$ 到第 $j$ 個字元組成的子字串，$i > j$ 時是空字串。

\begin{definition}
\label{def:stm-problem}
給\textbf{文字}（text）$T[1..n]$ 和\textbf{模式}（pattern）$P[1..m]$，$1 \le m \le
n$。 整數 $s$ 滿足 $0 \le s \le n - m$ 而且
\[
T[s + 1 .. s + m] = P[1..m]
\]
時，說 $P$ 以\textbf{位移}（shift）$s$ 出現在 $T$ 裡，$s$ 是一個\textbf{有效位移}
（valid shift）。 \textbf{字串比對}（string matching）問題要找出所有有效位移。
\end{definition}

位移 $s$ 的意思是把 $P$ 往右移 $s$ 格，和 $T$ 對齊；$P[1]$ 對到 $T[s + 1]$。 C 程式
裡字串從 $0$ 編號，位移 $s$ 剛好就是出現位置的起點 \texttt{t[s]}，兩邊的數字一樣，
不用換算。 出現的位置可以重疊：$P = \mathtt{AA}$ 在 $T = \mathtt{AAAA}$ 裡的有效位
移是 $0, 1, 2$。

後面講 KMP 時還要用到前綴和後綴。 $x[1..k]$（$0 \le k \le |x|$）叫 $x$ 的
\textbf{前綴}（prefix），$x[|x| - k + 1 .. |x|]$ 叫 $x$ 的\textbf{後綴}（suffix）；
$k < |x|$ 時加上「真」字，叫\textbf{真前綴}、\textbf{真後綴}。 空字串是每個字串的
前綴，也是後綴。 例如 \texttt{AAC} 是 \texttt{AACAA} 的前綴，\texttt{CAA} 是它的後
綴。

\subsection{樸素比對}

最直接的做法是把每一個可能的位移都試一次：$s$ 從 $0$ 到 $n - m$，從左往右比
$P[j]$ 和 $T[s + j]$，全部相同就是有效位移，遇到第一個不同的字元就換下一個 $s$。

\begin{pseudo}[label=alg:stm-naive]{樸素比對}
\Procedure{Naive-Match}{$T, P, n, m$}
    \For{$s \gets 0$ \To $n - m$}
        \State $j \gets 1$
        \While{$j \le m$ \AND $P[j] = T[s + j]$}
            \State $j \gets j + 1$
        \EndWhile
        \If{$j = m + 1$}
            \State 輸出 $s$ \Comment{$P[1..m]$ 全部對上}
        \EndIf
    \EndFor
\EndProcedure
\end{pseudo}

用 $T = \mathtt{AACAACAACAAAC}$（$n = 13$）、$P = \mathtt{AACAACAAAC}$（$m = 10$）跑
一次。 下圖每一列是一個位移，藍底是比對成功的字元，粉紅是第一個不同的字元，灰字是
沒有比到的字元；右邊是這個位移比了幾次（第 4 行的 $P[j] = T[s + j]$ 執行的次數，
$j = m + 1$ 那次因為 \textbf{and} 左邊不成立，不算）。

\begin{center}\begin{minipage}{\linewidth}\centering
\begin{tikzpicture}
\node[stm-lab] at (-0.20,0.00) {$T$};
\node[stm-c] at (0.00,0.00) {A};
\node[stm-c] at (0.50,0.00) {A};
\node[stm-c] at (1.00,0.00) {C};
\node[stm-c] at (1.50,0.00) {A};
\node[stm-c] at (2.00,0.00) {A};
\node[stm-c] at (2.50,0.00) {C};
\node[stm-c] at (3.00,0.00) {A};
\node[stm-c] at (3.50,0.00) {A};
\node[stm-c] at (4.00,0.00) {C};
\node[stm-c] at (4.50,0.00) {A};
\node[stm-c] at (5.00,0.00) {A};
\node[stm-c] at (5.50,0.00) {A};
\node[stm-c] at (6.00,0.00) {C};
\node[idxlbl] at (0.00,0.42) {1};
\node[idxlbl] at (0.50,0.42) {2};
\node[idxlbl] at (1.00,0.42) {3};
\node[idxlbl] at (1.50,0.42) {4};
\node[idxlbl] at (2.00,0.42) {5};
\node[idxlbl] at (2.50,0.42) {6};
\node[idxlbl] at (3.00,0.42) {7};
\node[idxlbl] at (3.50,0.42) {8};
\node[idxlbl] at (4.00,0.42) {9};
\node[idxlbl] at (4.50,0.42) {10};
\node[idxlbl] at (5.00,0.42) {11};
\node[idxlbl] at (5.50,0.42) {12};
\node[idxlbl] at (6.00,0.42) {13};
\node[stm-lab] at (-0.20,-0.75) {$s=0$};
\node[stm-ok] at (0.00,-0.75) {A};
\node[stm-ok] at (0.50,-0.75) {A};
\node[stm-ok] at (1.00,-0.75) {C};
\node[stm-ok] at (1.50,-0.75) {A};
\node[stm-ok] at (2.00,-0.75) {A};
\node[stm-ok] at (2.50,-0.75) {C};
\node[stm-ok] at (3.00,-0.75) {A};
\node[stm-ok] at (3.50,-0.75) {A};
\node[stm-bad] at (4.00,-0.75) {A};
\node[stm-idle] at (4.50,-0.75) {C};
\node[stm-note, anchor=west] at (4.95,-0.75) {9 次};
\node[stm-lab] at (-0.20,-1.40) {$s=1$};
\node[stm-ok] at (0.50,-1.40) {A};
\node[stm-bad] at (1.00,-1.40) {A};
\node[stm-idle] at (1.50,-1.40) {C};
\node[stm-idle] at (2.00,-1.40) {A};
\node[stm-idle] at (2.50,-1.40) {A};
\node[stm-idle] at (3.00,-1.40) {C};
\node[stm-idle] at (3.50,-1.40) {A};
\node[stm-idle] at (4.00,-1.40) {A};
\node[stm-idle] at (4.50,-1.40) {A};
\node[stm-idle] at (5.00,-1.40) {C};
\node[stm-note, anchor=west] at (5.45,-1.40) {2 次};
\node[stm-lab] at (-0.20,-2.05) {$s=2$};
\node[stm-bad] at (1.00,-2.05) {A};
\node[stm-idle] at (1.50,-2.05) {A};
\node[stm-idle] at (2.00,-2.05) {C};
\node[stm-idle] at (2.50,-2.05) {A};
\node[stm-idle] at (3.00,-2.05) {A};
\node[stm-idle] at (3.50,-2.05) {C};
\node[stm-idle] at (4.00,-2.05) {A};
\node[stm-idle] at (4.50,-2.05) {A};
\node[stm-idle] at (5.00,-2.05) {A};
\node[stm-idle] at (5.50,-2.05) {C};
\node[stm-note, anchor=west] at (5.95,-2.05) {1 次};
\node[stm-lab] at (-0.20,-2.70) {$s=3$};
\node[stm-ok] at (1.50,-2.70) {A};
\node[stm-ok] at (2.00,-2.70) {A};
\node[stm-ok] at (2.50,-2.70) {C};
\node[stm-ok] at (3.00,-2.70) {A};
\node[stm-ok] at (3.50,-2.70) {A};
\node[stm-ok] at (4.00,-2.70) {C};
\node[stm-ok] at (4.50,-2.70) {A};
\node[stm-ok] at (5.00,-2.70) {A};
\node[stm-ok] at (5.50,-2.70) {A};
\node[stm-ok] at (6.00,-2.70) {C};
\node[stm-note, anchor=west] at (6.45,-2.70) {10 次};
\end{tikzpicture}
\diagcaption{樸素比對試過的四個位移。 $s = 3$ 是唯一的有效位移，總共比較了
$9 + 2 + 1 + 10 = 22$ 次。}
\label{fig:stm-naive}
\end{minipage}\end{center}

\subsection{分析}

外層迴圈跑 $n - m + 1$ 次，每次內層最多比 $m$ 次，所以時間是
$O((n - m + 1)m)$。 這個上界會達到：$T = \mathtt{A}^n$（$n$ 個 \texttt{A}）、
$P = \mathtt{A}^{m-1}\mathtt{C}$ 時，每個位移都要比到最後一個字元才失敗，剛好比
$(n - m + 1)m$ 次。 $P = \mathtt{A}^m$ 時每個位移都是有效位移，也一樣多。 所以樸素
比對的最壞情況是 $\Theta((n - m + 1)m)$；$m = \lfloor n/2 \rfloor$ 時是
$\Theta(n^2)$。

最好的情況是每個位移的第一個字元就不同，共比 $n - m + 1$ 次。 一般的文字介於兩者
之間。 如果 $T$ 和 $P$ 的字元都像是從 $|\Sigma|$ 個字元裡隨機挑的，每個位移第一個
字元相同的機率是 $1/|\Sigma|$，前兩個都相同是 $1/|\Sigma|^2$，所以每個位移平均比不
到 $1 + 1/|\Sigma| + 1/|\Sigma|^2 + \cdots = |\Sigma| / (|\Sigma| - 1)$ 次；DNA 是
$4/3$。 真實的資料不是隨機的，DNA 和文章裡常有大段的重複，最壞情況就會出現。

浪費在哪裡，從圖~\ref{fig:stm-naive} 就看得出來。 $s = 0$ 時已經知道 $T[1..8] =
\mathtt{AACAACAA}$，$s = 1$ 和 $s = 2$ 卻又從頭去讀 $T[2]$、$T[3]$。 其實光憑已經對
上的這 $8$ 個字元，就能算出 $s = 1$、$s = 2$ 不可能成功，而 $s = 3$ 的前 $5$ 個字元
一定對得上。 第~\ref{sec:stm-kmp}~節的 KMP 演算法就是把這個推理事先算好。

\subsection{C 程式}

\begin{cbookrecap}[第~11~章 11.1--11.3、11.5 節，第~19~章 19.3 節]
C 的字串是一個字元陣列，最後放一個 \texttt{'\textbackslash 0'} 表示結束；
\texttt{strlen(s)} 從頭數到 \texttt{'\textbackslash 0'} 之前有幾個字元，每呼叫一次
就走一遍，花 $\Theta(|s|)$ 的時間，所以不要寫在迴圈條件裡。 \texttt{scanf("\%100000s",
t)} 讀一個不含空白的字，最多 $100000$ 個字元，陣列至少要開 $100001$ 格放
\texttt{'\textbackslash 0'}；不寫寬度的 \texttt{\%s} 可能寫出陣列外。
\texttt{strncmp(a, b, k)} 比較前 $k$ 個字元，相同時回傳 $0$。 標準函式庫的
\texttt{strstr(t, p)} 也能找子字串，但 C 標準沒有規定它用什麼方法、花多少時間。
陣列用 \texttt{malloc} 向系統要，用完 \texttt{free}。
\end{cbookrecap}

下面的程式讀進好幾組 $T$、$P$，印出所有有效位移，以及比較的次數和
$(n - m + 1)m$。 第一組是圖~\ref{fig:stm-naive} 的例子，第二組是一段普通的 DNA，最
後兩組是上面說的最壞情況。

\begin{example}[label=ex:stm-naive]{c}{(naive\_match.c) 樸素比對並計算比較次數}
#include <stdio.h>
#include <stdlib.h>
#include <string.h>

#define MAXLEN 100000           // T、P 最長的長度

// 樸素比對：印出 P 在 T 裡的每一個有效位移，回傳字元比較的次數。
long long naive_match(const char *t, int n, const char *p, int m)
{
    long long cmp = 0;
    int found = 0;
    printf("shifts:");
    for (int s = 0; s <= n - m; s++) {
        int j = 0;
        while (j < m) {          // 從左往右比，第一個不同就停
            cmp++;
            if (p[j] != t[s + j]) {
                break;
            }
            j++;
        }
        if (j == m) {
            printf(" %d", s);
            found++;
        }
    }
    printf("%s\n", found ? "" : " none");
    return cmp;
}

int main(void)
{
    char *t = malloc(MAXLEN + 1);
    char *p = malloc(MAXLEN + 1);
    int cases;
    if (t == NULL || p == NULL || scanf("%d", &cases) != 1) {
        free(t);
        free(p);
        return 1;
    }
    while (cases--) {
        if (scanf("%100000s %100000s", t, p) != 2) {
            break;
        }
        int n = (int)strlen(t), m = (int)strlen(p);
        long long cmp = naive_match(t, n, p, m);
        long long bound = (n >= m) ? (long long)(n - m + 1) * m : 0;
        printf("comparisons = %lld, (n-m+1)m = %lld\n", cmp, bound);
    }
    free(t);
    free(p);
    return 0;
}
\end{example}
\iosample{4 \\ AACAACAACAAAC AACAACAAAC \\ TGCCTCAGTCATCC CAGT \\ AAAAAAAAAAAA AAAAC \\ AAAAAAAAAAAA AAAA}{shifts: 3 \\ comparisons = 22, (n-m+1)m = 40 \\ shifts: 5 \\ comparisons = 18, (n-m+1)m = 44 \\ shifts: none \\ comparisons = 40, (n-m+1)m = 40 \\ shifts: 0 1 2 3 4 5 6 7 8 \\ comparisons = 36, (n-m+1)m = 36}

普通的 DNA 那一組只比了 $18$ 次，比 $n = 14$ 多一點；最壞的兩組都剛好比了
$(n - m + 1)m$ 次。

\section{Rabin--Karp 演算法}
\label{sec:stm-rk}

\subsection{把字串當成數字}

先假設字母表只有十個數字 $0$ 到 $9$。 長度 $m$ 的字串可以看成一個 $m$ 位數，兩個字
串相同若且唯若兩個數相同。 如果能很快算出 $T$ 每一個長度 $m$ 的視窗代表的數，比對
就變成比較兩個整數，一次 $O(1)$。

一般的字母表也一樣。 把每個字元對應到 $0$ 到 $d - 1$ 的一個數字（$d = |\Sigma|$），
字串就是一個 $d$ 進位的數。 這一節的 DNA 例子用 $d = 4$，
$\mathtt{A} = 0$、$\mathtt{C} = 1$、$\mathtt{G} = 2$、$\mathtt{T} = 3$。 令
\[
t_s = \sum_{j=1}^{m} T[s + j] \cdot d^{m-j}
\]
是第 $s$ 個視窗 $T[s + 1 .. s + m]$ 的值（這裡 $T[i]$ 指字元對應的數字），$p$ 是
$P$ 的值。 $p$ 和 $t_0$ 用第~4~章的 Horner 法（虛擬碼~\ref{alg:arr-horner}）各花
$\Theta(m)$ 算出來。 之後的視窗不用重算：從 $t_s$ 到 $t_{s+1}$，最高位的
$T[s + 1]$ 離開，其他位數往左移一位，$T[s + m + 1]$ 從最低位進來：
\begin{equation}
t_{s+1} = d \bigl( t_s - T[s + 1] \cdot d^{m-1} \bigr) + T[s + m + 1].
\label{eq:stm-roll}
\end{equation}
這叫\textbf{滾動雜湊}（rolling hash），每次只要常數個算術運算。

問題是 $t_s$ 有 $m$ 位，$m$ 一大就放不進一個字組，算術運算不能再當成一步
（第~\ref{sec:ana-input-size}~節）。 解法是全部對一個質數 $q$ 取餘數，也就是第~12~章字串鍵的雜湊函數
（第~\ref{sec:hash-func}~節）：
\[
p \bmod q, \qquad t_s \bmod q, \qquad
t_{s+1} \equiv d \bigl( t_s - T[s + 1] \cdot h \bigr) + T[s + m + 1] \pmod q,
\]
其中 $h = d^{m-1} \bmod q$ 事先算好。 取餘數之後，$t_s \not\equiv p$ 一定表示視窗和
$P$ 不同，可以直接跳過；但 $t_s \equiv p$ 不保證相同，可能只是兩個不同的字串雜湊到
同一個值，也就是第~12~章的\textbf{碰撞}。 雜湊值相等叫\textbf{命中}（hit），命中了
但字串其實不同，叫\textbf{偽命中}（spurious hit）。 所以命中之後還要逐字確認一次。

\begin{pseudo}[label=alg:stm-rk]{Rabin--Karp 比對，$\bmod$ 的結果一律在 $0$ 到 $q - 1$}
\Procedure{Rabin-Karp}{$T, P, n, m, d, q$}
    \State $h \gets d^{m-1} \bmod q$
    \State $p \gets 0$
    \State $t \gets 0$
    \For{$i \gets 1$ \To $m$} \Comment{Horner 法}
        \State $p \gets (d \cdot p + P[i]) \bmod q$
        \State $t \gets (d \cdot t + T[i]) \bmod q$
    \EndFor
    \For{$s \gets 0$ \To $n - m$}
        \If{$p = t$} \Comment{命中}
            \If{$P[1..m] = T[s + 1 .. s + m]$} \Comment{逐字確認}
                \State 輸出 $s$
            \EndIf
        \EndIf
        \If{$s < n - m$}
            \State $t \gets \bigl( d \, (t - T[s + 1] \cdot h) + T[s + m + 1] \bigr) \bmod q$
        \EndIf
    \EndFor
\EndProcedure
\end{pseudo}

第 8 行的迴圈每次開始時，$t = t_s \bmod q$；第 13 行用式~\eqref{eq:stm-roll} 把它換
成 $t_{s+1} \bmod q$，所以這個性質一直成立。

\subsection{手算}

$T = \mathtt{TGCCTCAGTCATCC}$（$n = 14$），$P = \mathtt{CAGT}$（$m = 4$），$d = 4$，
$q = 11$。 先算 $h = 4^3 \bmod 11 = 64 \bmod 11 = 9$，
$p = 1 \cdot 64 + 0 \cdot 16 + 2 \cdot 4 + 3 = 75$，$75 \bmod 11 = 9$。 下表列出每
個視窗；考試時「視窗的值」那一欄可以不寫，直接用滾動的式子算餘數，這裡列出來是為了
驗算。

\begin{center}
\begin{tabular}{cccccl}
\toprule
$s$ & 視窗 & 視窗的值 $t_s$ & $t_s \bmod 11$ & 命中？ & 結果 \\
\midrule
0 & \texttt{TGCC} & 229 & 9 & 是 & 偽命中 \\
1 & \texttt{GCCT} & 151 & 8 &  &  \\
2 & \texttt{CCTC} & 93 & 5 &  &  \\
3 & \texttt{CTCA} & 116 & 6 &  &  \\
4 & \texttt{TCAG} & 210 & 1 &  &  \\
5 & \texttt{CAGT} & 75 & 9 & 是 & 有效位移 \\
6 & \texttt{AGTC} & 45 & 1 &  &  \\
7 & \texttt{GTCA} & 180 & 4 &  &  \\
8 & \texttt{TCAT} & 211 & 2 &  &  \\
9 & \texttt{CATC} & 77 & 0 &  &  \\
10 & \texttt{ATCC} & 53 & 9 & 是 & 偽命中 \\
\bottomrule
\end{tabular}
\end{center}

三次命中裡只有 $s = 5$ 是真的，$s = 0$ 和 $s = 10$ 是偽命中，各要多做一次逐字比較
才知道。 從 $s = 0$ 滾到 $s = 1$ 的計算：離開的是 $T[1] = \mathtt{T} = 3$，進來的是
$T[5] = \mathtt{T} = 3$，
\[
t_1 \equiv 4 \, (9 - 3 \cdot 9) + 3 = 4 \cdot (-18) + 3 = -69 \equiv -69 + 77 = 8
\pmod{11},
\]
和表上 $151 \bmod 11 = 8$ 一樣。 中間出現負數是正常的，取餘數時加上 $q$ 的倍數變回
$0$ 到 $q - 1$ 就好。

\begin{examnote}
Rabin--Karp 的計算題先看清楚三件事：字元怎麼對應到數字（題目可能用
$\mathtt{a} = 1$，或直接用 ASCII 碼）、$d$ 和 $q$ 是多少、字串的第一個字元是最高位
還是最低位。 先寫出 $h = d^{m-1} \bmod q$ 和 $p \bmod q$，再一個視窗一列，列出
$t_s \bmod q$、有沒有命中、是有效位移還是偽命中。 題目問「偽命中幾次」時，有效位移
不算在內。
\end{examnote}

\subsection{時間}

前處理（第 2--7 行）花 $\Theta(m)$；第 2 行的 $d^{m-1} \bmod q$ 連乘 $m - 1$ 次就
夠。 比對的部分，每個位移花 $O(1)$ 算雜湊、比雜湊值，命中時再加 $O(m)$ 確認。 設有
$v$ 個有效位移、$x$ 次偽命中，比對時間是
\[
\Theta(n - m + 1) + O\bigl(m (v + x)\bigr).
\]
最壞情況沒有比較好：$T = \mathtt{A}^n$、$P = \mathtt{A}^m$ 時每個位移都要逐字確
認，又是 $\Theta((n - m + 1)m)$。 偽命中也可能很多，例如 $q$ 選得不好時（觀念題
第~\ref{q:stm-power-mod}~題）。

偽命中有多少，要看雜湊值像不像隨機。 照第~12~章的習慣先做一個假設：每個不是有效位
移的視窗，雜湊值等於 $p \bmod q$ 的機率大約是 $1/q$。 那麼偽命中的期望次數大約是
$(n - m + 1)/q$，比對時間的期望值是
\[
O\Bigl(n + m \Bigl(v + \frac{n}{q}\Bigr)\Bigr),
\]
有效位移不多、$q \ge m$ 時就是 $O(n + m)$。 上面的例子有 $11$ 個視窗、$q = 11$，期
望大約 $1$ 次偽命中，實際上有 $2$ 次。

這個假設不能證明，因為 $d$ 和 $q$ 是固定的，有心人可以專挑會碰撞的文字。 要得到真
正的保證，可以像第~12~章的通用雜湊（第~\ref{sec:hash-dynamic}~節）一樣，把隨機性
放進演算法裡：質數 $q$ 固定，底數不用 $d$，改成每次執行時從 $0$ 到 $q - 1$ 隨機挑
一個 $x$，字串 $a_1 a_2 \cdots a_m$ 的雜湊值是
$H_x(a) = \bigl( \sum_{j=1}^{m} a_j \, x^{m-j} \bigr) \bmod q$。 滾動更新的式子不
變，只是把 $d$ 換成 $x$。

\begin{theorem}
\label{thm:stm-rk-random}
$q$ 是質數，字元對應的數字都小於 $q$，$x$ 從 $\{0, 1, \ldots, q - 1\}$ 均勻隨機挑
選。 對任兩個長度 $m$ 的不同字串 $a$、$b$，
\[
\Pr\bigl[ H_x(a) = H_x(b) \bigr] \le \frac{m - 1}{q}.
\]
所以偽命中的期望次數不超過 $(n - m + 1)(m - 1)/q$。
\end{theorem}

\begin{proof}
令 $D(x) = \sum_{j=1}^{m} (a_j - b_j) \, x^{m-j}$，這是次數不超過 $m - 1$ 的多項
式。 $a \ne b$，所以有某個 $j$ 使 $a_j \ne b_j$；兩個數都在 $0$ 到 $q - 1$ 之間，差
不是 $q$ 的倍數，所以 $D$ 在模 $q$ 下不是零多項式。 數論的一個基本事實是：模質數
時，次數 $k \ge 0$ 的非零多項式最多有 $k$ 個根。 $H_x(a) = H_x(b)$ 正是
$D(x) \equiv 0 \pmod q$，在 $q$ 個可能的 $x$ 裡最多 $m - 1$ 個會讓它成立。 第二句：
每個不是有效位移的視窗都是一對不同的字串，偽命中的機率不超過 $(m - 1)/q$，對最多
$n - m + 1$ 個視窗用期望值的線性加起來。
\end{proof}

選 $q > nm$ 的質數，偽命中的期望次數就小於 $1$。 這種「結果一定對、時間是隨機的」
演算法叫 Las Vegas 型；如果為了省時間，命中就直接回報、不逐字確認，結果就有一點機會
錯，那叫 Monte Carlo 型。 Rabin--Karp 的好處是容易推廣：同時找很多個長度相同的模
式時，把它們的雜湊值放進一張雜湊表，每個視窗查一次表就好，KMP 做不到這麼簡單。

\subsection{C 程式}

程式把 \texttt{A}、\texttt{C}、\texttt{G}、\texttt{T} 換成 $0$ 到 $3$，固定 $d = 4$，
$q$ 由輸入給。 C 的 \texttt{\%} 對負數會得到負的餘數（\texttt{-69 \% 11} 是
\texttt{-3}），所以減掉最高位時先加 $q$，讓結果保持在 $0$ 到 $q - 1$。 $q$ 不超過
$10^9$ 時，\texttt{ht * D} 和 \texttt{code(t[s]) * h} 都放得進 \texttt{long long}。

\begin{example}[label=ex:stm-rk]{c}{(rabin\_karp.c) Rabin--Karp 比對，列出每一次命中}
#include <stdio.h>
#include <stdlib.h>
#include <string.h>

#define MAXLEN 100000
#define D 4                     // 字母表大小：A C G T

// 把鹼基換成 0..3。題目保證只有 A、C、G、T。
int code(char c)
{
    switch (c) {
    case 'A': return 0;
    case 'C': return 1;
    case 'G': return 2;
    default:  return 3;         // 'T'
    }
}

// Rabin-Karp：模數 q。印出每一次命中，回傳偽命中的次數。
int rabin_karp(const char *t, int n, const char *p, int m,
               long long q)
{
    long long h = 1;            // h = D^(m-1) mod q
    for (int i = 1; i < m; i++) {
        h = h * D % q;
    }
    long long hp = 0, ht = 0;   // P 和第一個視窗的雜湊值
    for (int i = 0; i < m; i++) {
        hp = (hp * D + code(p[i])) % q;
        ht = (ht * D + code(t[i])) % q;
    }
    int spurious = 0;
    for (int s = 0; s <= n - m; s++) {
        if (hp == ht) {                       // 命中：再逐字確認
            if (strncmp(t + s, p, m) == 0) {
                printf("s = %d: hit, match\n", s);
            } else {
                printf("s = %d: hit, spurious\n", s);
                spurious++;
            }
        }
        if (s < n - m) {                      // 滾到下一個視窗
            // 先減掉最高位；加 q 讓結果不是負數
            ht = (ht - code(t[s]) * h % q + q) % q;
            ht = (ht * D + code(t[s + m])) % q;
        }
    }
    return spurious;
}

int main(void)
{
    char *t = malloc(MAXLEN + 1);
    char *p = malloc(MAXLEN + 1);
    int cases;
    if (t == NULL || p == NULL || scanf("%d", &cases) != 1) {
        free(t);
        free(p);
        return 1;
    }
    while (cases--) {
        long long q;
        if (scanf("%lld %100000s %100000s", &q, t, p) != 3) {
            break;
        }
        int n = (int)strlen(t), m = (int)strlen(p);
        if (m == 0 || m > n) {
            continue;
        }
        printf("q = %lld\n", q);
        int sp = rabin_karp(t, n, p, m, q);
        printf("spurious hits = %d\n", sp);
    }
    free(t);
    free(p);
    return 0;
}
\end{example}
\iosample{2 \\ 11 TGCCTCAGTCATCC CAGT \\ 1000003 TGCCTCAGTCATCC CAGT}{q = 11 \\ s = 0: hit, spurious \\ s = 5: hit, match \\ s = 10: hit, spurious \\ spurious hits = 2 \\ q = 1000003 \\ s = 5: hit, match \\ spurious hits = 0}

$q = 11$ 的結果和手算的表一樣。 $q$ 換成質數 $1000003$ 以後，同一段文字只剩真正的
那次命中。

\section{KMP 演算法}
\label{sec:stm-kmp}

Knuth、Morris 和 Pratt 在 1970 年代提出這個演算法，所以叫 KMP。 它的比對時間是
$\Theta(n)$，而且不靠任何機率。

\subsection{對上的字元告訴我們什麼}

回到圖~\ref{fig:stm-naive}。 位移 $s$ 已經對上了 $q$ 個字元，第 $q + 1$ 個不同。 這
時我們知道 $T[s + 1 .. s + q] = P[1..q]$，這段文字是什麼，完全由 $P$ 決定，不用再
讀 $T$。 下一個可能成功的位移 $s' > s$ 要滿足：把 $P$ 移到 $s'$ 以後，和已知的這段
文字重疊的部分要對得上。 重疊的部分是 $P$ 的一個前綴 $P[1..k]$，對著
$P[1..q]$ 的後綴，其中 $k = s + q - s' < q$。 所以要找的是
\[
\text{同時是 } P[1..q] \text{ 的真前綴和真後綴的字串，越長越好。}
\]
越長，$s'$ 越小，才不會跳過任何一個可能的位移；而且新位移的前 $k$ 個字元已知相同，
直接從 $P[k + 1]$ 接著比。

圖~\ref{fig:stm-naive} 裡 $q = 8$，$P[1..8] = \mathtt{AACAACAA}$。 它的前綴和後綴相
同的有 \texttt{AACAA}（長度 $5$）、\texttt{AA}、\texttt{A}，最長的是 $5$。 所以下一個
要試的位移是 $s' = 0 + 8 - 5 = 3$，而且前 $5$ 個字元不必比，從 $P[6]$ 開始。 樸素
比對在 $s = 1$、$s = 2$ 花的 $3$ 次比較全部省掉，前 $5$ 個字元也不用重比。 這個
「最長的長度」只和 $P$ 有關，可以事先對每個 $q$ 算好。

\subsection{邊界與失敗函數}

\begin{definition}
\label{def:stm-border}
字串 $x$ 的\textbf{邊界}（border）是同時為 $x$ 的真前綴和真後綴的字串。 空字串是
每個非空字串的邊界。 模式 $P[1..m]$ 的\textbf{失敗函數}（failure function）$\pi$ 定
義在 $q = 1, 2, \ldots, m$ 上：
\[
\pi[q] = P[1..q] \text{ 最長的邊界的長度}.
\]
\end{definition}

邊界是真前綴，所以 $0 \le \pi[q] < q$，特別是 $\pi[1] = 0$。 失敗函數也叫\textbf{前
綴函數}（prefix function），名字來自它記錄的是 $P$ 的前綴之間的關係；Aki-Judge 的題
目用的是這個名字。 本章的例子 $P = \mathtt{AACAACAAAC}$ 的失敗函數如下，下一小節會
示範怎麼算。

\begin{center}
\begin{tabular}{c*{10}{c}}
\toprule
$q$ & 1 & 2 & 3 & 4 & 5 & 6 & 7 & 8 & 9 & 10 \\
$P[q]$ & \texttt{A} & \texttt{A} & \texttt{C} & \texttt{A} & \texttt{A} & \texttt{C} & \texttt{A} & \texttt{A} & \texttt{A} & \texttt{C} \\
$\pi[q]$ & 0 & 1 & 0 & 1 & 2 & 3 & 4 & 5 & 2 & 3 \\
\bottomrule
\end{tabular}
\end{center}

最長的邊界有了，其他的邊界怎麼找？ 下面的引理說：邊界的邊界還是邊界，而且所有的邊
界都可以這樣一層一層找到。

\begin{lemma}[邊界的鏈]
\label{lem:stm-chain}
對 $1 \le q \le m$，$P[1..q]$ 所有邊界的長度，由大到小恰好是
\[
\pi[q], \quad \pi[\pi[q]], \quad \pi[\pi[\pi[q]]], \quad \ldots
\]
一直到 $0$ 為止。
\end{lemma}

\begin{proof}
鏈上的每一個都是邊界：$\pi[q]$ 是 $P[1..q]$ 的邊界。 如果 $u$ 是 $v$ 的邊界、$v$ 是
$x$ 的邊界，$u$ 是 $v$ 的前綴，$v$ 是 $x$ 的前綴，所以 $u$ 是 $x$ 的前綴；後綴同理，
而且 $|u| < |v| < |x|$，所以 $u$ 也是 $x$ 的邊界。 鏈上每一項嚴格變小，一定會走到
$0$。

沒有漏掉：假設 $P[1..b]$ 是 $P[1..q]$ 的邊界，但 $b$ 不在鏈上。 $b \le \pi[q]$（因
為 $\pi[q]$ 最大），又 $b \ne \pi[q]$，所以鏈上有比 $b$ 大的數；令 $c$ 是其中最小
的一個。 $P[1..b]$ 和 $P[1..c]$ 都是 $P[1..q]$ 的後綴，$b < c$，所以 $P[1..b]$ 是
$P[1..c]$ 的後綴，它本來就是 $P[1..c]$ 的前綴，因此是 $P[1..c]$ 的邊界，
$b \le \pi[c]$。 $\pi[c]$ 是鏈上 $c$ 的下一項，比 $c$ 小；由 $c$ 的選法，它不會比
$b$ 大，所以 $\pi[c] = b$，$b$ 在鏈上，矛盾。
\end{proof}

下圖是 $q = 8$ 的鏈。 每一列把 $P$ 的前綴對齊在 $P[1..8]$ 的尾巴下面，對得上的部分
就是一個邊界。

\begin{center}\begin{minipage}{\linewidth}\centering
\begin{tikzpicture}
\node[stm-lab] at (-0.20,0.00) {$P[1..8]$};
\node[stm-c] at (0.00,0.00) {A};
\node[stm-c] at (0.50,0.00) {A};
\node[stm-c] at (1.00,0.00) {C};
\node[stm-c] at (1.50,0.00) {A};
\node[stm-c] at (2.00,0.00) {A};
\node[stm-c] at (2.50,0.00) {C};
\node[stm-c] at (3.00,0.00) {A};
\node[stm-c] at (3.50,0.00) {A};
\node[idxlbl] at (0.00,0.42) {1};
\node[idxlbl] at (0.50,0.42) {2};
\node[idxlbl] at (1.00,0.42) {3};
\node[idxlbl] at (1.50,0.42) {4};
\node[idxlbl] at (2.00,0.42) {5};
\node[idxlbl] at (2.50,0.42) {6};
\node[idxlbl] at (3.00,0.42) {7};
\node[idxlbl] at (3.50,0.42) {8};
\node[stm-lab] at (-0.20,-0.75) {$\pi[8]=5$};
\node[stm-ok] at (1.50,-0.75) {A};
\node[stm-ok] at (2.00,-0.75) {A};
\node[stm-ok] at (2.50,-0.75) {C};
\node[stm-ok] at (3.00,-0.75) {A};
\node[stm-ok] at (3.50,-0.75) {A};
\node[stm-lab] at (-0.20,-1.40) {$\pi[5]=2$};
\node[stm-ok] at (3.00,-1.40) {A};
\node[stm-ok] at (3.50,-1.40) {A};
\node[stm-lab] at (-0.20,-2.05) {$\pi[2]=1$};
\node[stm-ok] at (3.50,-2.05) {A};
\node[stm-lab] at (-0.20,-2.70) {$\pi[1]=0$};
\node[stm-note, anchor=west] at (0.00,-2.70) {空字串};
\end{tikzpicture}
\diagcaption{$P[1..8] = \mathtt{AACAACAA}$ 的邊界：長度 $5, 2, 1, 0$，依序是
$\pi[8]$、$\pi[5]$、$\pi[2]$、$\pi[1]$。}
\label{fig:stm-borders}
\end{minipage}\end{center}

\subsection{計算失敗函數}

$\pi[q]$ 可以從前面的值推出來。 設 $P[1..r]$ 是 $P[1..q]$ 的邊界，$r \ge 1$。 拿掉
兩邊的最後一個字元，$P[1..r - 1]$ 是 $P[1..q - 1]$ 的邊界，而且 $P[r] = P[q]$。 反
過來，$P[1..k]$ 是 $P[1..q - 1]$ 的邊界、$P[k + 1] = P[q]$ 時，$P[1..k + 1]$ 是
$P[1..q]$ 的邊界。 所以
\[
\pi[q] = 1 + \max\bigl\{ k : k \text{ 是 } P[1..q-1] \text{ 的邊界長度，且 }
P[k + 1] = P[q] \bigr\},
\]
集合是空的時候 $\pi[q] = 0$。 由引理~\ref{lem:stm-chain}，$P[1..q - 1]$ 的邊界長度
就是從 $\pi[q - 1]$ 開始的鏈，從大往小試，第一個 $P[k + 1] = P[q]$ 的就是最大的。
這也順便說明了 $\pi[q] \le \pi[q - 1] + 1$：失敗函數每往右一格，最多加一。

\begin{pseudo}[label=alg:stm-failure]{計算失敗函數}
\Function{Compute-Failure}{$P, m$}
    \State $\pi[1] \gets 0$
    \State $k \gets 0$ \Comment{$k = \pi[q - 1]$}
    \For{$q \gets 2$ \To $m$}
        \While{$k > 0$ \AND $P[k + 1] \ne P[q]$}
            \State $k \gets \pi[k]$ \Comment{換成下一個較短的邊界}
        \EndWhile
        \If{$P[k + 1] = P[q]$}
            \State $k \gets k + 1$
        \EndIf
        \State $\pi[q] \gets k$
    \EndFor
    \State \Return $\pi$
\EndFunction
\end{pseudo}

第 5 行的迴圈從 $k = \pi[q - 1]$ 開始沿著鏈往下走，走到 $P[k + 1] = P[q]$ 或
$k = 0$ 為止。 停在 $k = 0$ 時，第 7 行還要看 $P[1]$ 和 $P[q]$ 是否相同，因為空字串
接上一個字元也可能是邊界。

下表是 $P = \mathtt{AACAACAAAC}$ 的計算過程。 「比較」一欄依序寫出每一次比較
$P[k + 1]$ 和 $P[q]$ 的結果；不同時 $k$ 退到 $\pi[k]$ 再比。

\begin{center}
\small
\begin{tabular}{ccc>{\raggedright\arraybackslash}p{0.56\linewidth}c}
\toprule
$q$ & $P[q]$ & 開始的 $k$ & 比較 & $\pi[q]$ \\
\midrule
1 & \texttt{A} & & 定義 & 0 \\
2 & \texttt{A} & 0 & $P[1]=\mathtt{A}$ 相同，$k\gets1$ & 1 \\
3 & \texttt{C} & 1 & $P[2]=\mathtt{A}\ne\mathtt{C}$，$k\gets\pi[1]=0$；$P[1]=\mathtt{A}\ne\mathtt{C}$，$k=0$ & 0 \\
4 & \texttt{A} & 0 & $P[1]=\mathtt{A}$ 相同，$k\gets1$ & 1 \\
5 & \texttt{A} & 1 & $P[2]=\mathtt{A}$ 相同，$k\gets2$ & 2 \\
6 & \texttt{C} & 2 & $P[3]=\mathtt{C}$ 相同，$k\gets3$ & 3 \\
7 & \texttt{A} & 3 & $P[4]=\mathtt{A}$ 相同，$k\gets4$ & 4 \\
8 & \texttt{A} & 4 & $P[5]=\mathtt{A}$ 相同，$k\gets5$ & 5 \\
9 & \texttt{A} & 5 & $P[6]=\mathtt{C}\ne\mathtt{A}$，$k\gets\pi[5]=2$；$P[3]=\mathtt{C}\ne\mathtt{A}$，$k\gets\pi[2]=1$；$P[2]=\mathtt{A}$ 相同，$k\gets2$ & 2 \\
10 & \texttt{C} & 2 & $P[3]=\mathtt{C}$ 相同，$k\gets3$ & 3 \\
\bottomrule
\end{tabular}
\end{center}

$q = 9$ 那一列退了兩次：$P[1..8]$ 的最長邊界 $5$ 接不上 \texttt{A}，退到 $\pi[5] = 2$
也接不上，再退到 $\pi[2] = 1$，$P[2] = \mathtt{A}$ 接上了，得到 $\pi[9] = 2$。 這三個
$k$ 正是圖~\ref{fig:stm-borders} 的鏈。

\subsection{KMP 比對}

比對文字的做法和計算失敗函數幾乎一樣，只是拿 $T$ 去對 $P$。 變數 $q$ 記錄目前對上
了幾個字元。 讀進 $T[i]$ 時，如果 $P[q + 1] = T[i]$，就多對上一個；不同的話，把 $q$
換成 $\pi[q]$（相當於把 $P$ 往右移 $q - \pi[q]$ 格），再試，直到對上或 $q = 0$。
$q = m$ 時找到一次出現，接著把 $q$ 設成 $\pi[m]$，繼續找可能和它重疊的下一次。

\begin{pseudo}[label=alg:stm-kmp]{KMP 比對}
\Procedure{KMP-Match}{$T, P, n, m$}
    \State $\pi \gets$ \Call{Compute-Failure}{$P, m$}
    \State $q \gets 0$ \Comment{已經對上的字元數}
    \For{$i \gets 1$ \To $n$}
        \While{$q > 0$ \AND $P[q + 1] \ne T[i]$}
            \State $q \gets \pi[q]$ \Comment{失配，退到較短的邊界}
        \EndWhile
        \If{$P[q + 1] = T[i]$}
            \State $q \gets q + 1$
        \EndIf
        \If{$q = m$}
            \State 輸出 $i - m$ \Comment{$T[i - m + 1 .. i] = P$}
            \State $q \gets \pi[q]$
        \EndIf
    \EndFor
\EndProcedure
\end{pseudo}

\begin{theorem}
\label{thm:stm-kmp-correct}
\proc{KMP-Match} 輸出的恰好是所有有效位移。
\end{theorem}

\begin{whybox}[定理~\ref{thm:stm-kmp-correct}]
令 $f(i)$ 是「$P$ 的前綴之中，是 $T[1..i]$ 的後綴的最長長度」，$g(i)$ 是同樣的東西
但只考慮長度小於 $m$ 的前綴；$f(0) = g(0) = 0$。 證明這個不變量：第 4 行的迴圈每
一輪開始時，$q = g(i - 1)$。 有了它，第 5--8 行之後 $q = f(i)$（下面證明），而
$f(i) = m$ 正好表示 $T[i - m + 1 .. i] = P$，也就是位移 $i - m$ 有效，所以第 10 行
輸出的恰好是所有有效位移。

初始：$i = 1$ 時 $q = 0 = g(0)$。

第 5--8 行算出 $f(i)$：如果 $P[1..k]$ 是 $T[1..i]$ 的後綴而且 $k \ge 1$，拿掉最後一
個字元，$P[1..k - 1]$ 是 $T[1..i - 1]$ 的後綴，而且 $P[k] = T[i]$。 所以 $f(i)$ 是
$j + 1$ 的最大值，$j$ 跑遍「$P[1..j]$ 是 $T[1..i - 1]$ 的後綴、$j < m$、
$P[j + 1] = T[i]$」的所有 $j$；沒有這樣的 $j$ 時 $f(i) = 0$。 這些 $j$ 都不超過
$g(i - 1) = q$；$P[1..j]$ 和 $P[1..q]$ 都是 $T[1..i - 1]$ 的後綴，所以 $j = q$ 或
$P[1..j]$ 是 $P[1..q]$ 的邊界。 由引理~\ref{lem:stm-chain}，候選的 $j$ 恰好是
$q, \pi[q], \pi[\pi[q]], \ldots, 0$。 第 5--6 行由大到小試，停在第一個
$P[j + 1] = T[i]$ 的 $j$（或 $0$），第 7--8 行再加一，得到的就是 $f(i)$。

維持：若 $f(i) < m$，$g(i) = f(i)$，下一輪開始時 $q = g(i)$。 若 $f(i) = m$，
$T[1..i]$ 以 $P$ 結尾；長度 $k < m$ 的前綴 $P[1..k]$ 是 $T[1..i]$ 的後綴，等於說它
是 $P$ 的後綴，也就是 $P$ 的邊界，所以 $g(i) = \pi[m]$，正是第 11 行設的值。
\end{whybox}

\subsection{手算追蹤}

$T = \mathtt{AACAACAACAAACAACAAAGT}$（$n = 21$），$P$ 同上。 下表一個字元一列：讀
$T[i]$ 之前的 $q$、失配時依序退到的 $q$（第 6 行），以及讀完之後的 $q$。

\begin{center}
\small
\begin{tabular}{ccccc>{\raggedright\arraybackslash}p{0.38\linewidth}}
\toprule
$i$ & $T[i]$ & 開始的 $q$ & 退到 & 結束的 $q$ & 說明 \\
\midrule
1 & \texttt{A} & 0 & -- & 1 &  \\
2 & \texttt{A} & 1 & -- & 2 &  \\
3 & \texttt{C} & 2 & -- & 3 &  \\
4 & \texttt{A} & 3 & -- & 4 &  \\
5 & \texttt{A} & 4 & -- & 5 &  \\
6 & \texttt{C} & 5 & -- & 6 &  \\
7 & \texttt{A} & 6 & -- & 7 &  \\
8 & \texttt{A} & 7 & -- & 8 &  \\
9 & \texttt{C} & 8 & 5 & 6 &  \\
10 & \texttt{A} & 6 & -- & 7 &  \\
11 & \texttt{A} & 7 & -- & 8 &  \\
12 & \texttt{A} & 8 & -- & 9 &  \\
13 & \texttt{C} & 9 & -- & 10 & 找到，位移 $s=3$；$q\gets\pi[10]=3$ \\
14 & \texttt{A} & 3 & -- & 4 &  \\
15 & \texttt{A} & 4 & -- & 5 &  \\
16 & \texttt{C} & 5 & -- & 6 &  \\
17 & \texttt{A} & 6 & -- & 7 &  \\
18 & \texttt{A} & 7 & -- & 8 &  \\
19 & \texttt{A} & 8 & -- & 9 &  \\
20 & \texttt{G} & 9 & 2 $\to$ 1 $\to$ 0 & 0 &  \\
21 & \texttt{T} & 0 & -- & 0 &  \\
\bottomrule
\end{tabular}
\end{center}

同一個過程畫成對齊圖：每次 $q$ 退到 $\pi[q]$，就等於把 $P$ 移到一個新的位移。 綠底
是不用比就知道相同的字元（前一個位移留下來的邊界），藍底是這個位移比對成功的字元，
粉紅是失配的字元。 超出 $T$ 的部分不畫。

\begin{center}\begin{minipage}{\linewidth}\centering
\begin{tikzpicture}
\node[stm-lab] at (-0.20,0.00) {$T$};
\node[stm-c] at (0.00,0.00) {A};
\node[stm-c] at (0.50,0.00) {A};
\node[stm-c] at (1.00,0.00) {C};
\node[stm-c] at (1.50,0.00) {A};
\node[stm-c] at (2.00,0.00) {A};
\node[stm-c] at (2.50,0.00) {C};
\node[stm-c] at (3.00,0.00) {A};
\node[stm-c] at (3.50,0.00) {A};
\node[stm-c] at (4.00,0.00) {C};
\node[stm-c] at (4.50,0.00) {A};
\node[stm-c] at (5.00,0.00) {A};
\node[stm-c] at (5.50,0.00) {A};
\node[stm-c] at (6.00,0.00) {C};
\node[stm-c] at (6.50,0.00) {A};
\node[stm-c] at (7.00,0.00) {A};
\node[stm-c] at (7.50,0.00) {C};
\node[stm-c] at (8.00,0.00) {A};
\node[stm-c] at (8.50,0.00) {A};
\node[stm-c] at (9.00,0.00) {A};
\node[stm-c] at (9.50,0.00) {G};
\node[stm-c] at (10.00,0.00) {T};
\node[idxlbl] at (0.00,0.42) {1};
\node[idxlbl] at (0.50,0.42) {2};
\node[idxlbl] at (1.00,0.42) {3};
\node[idxlbl] at (1.50,0.42) {4};
\node[idxlbl] at (2.00,0.42) {5};
\node[idxlbl] at (2.50,0.42) {6};
\node[idxlbl] at (3.00,0.42) {7};
\node[idxlbl] at (3.50,0.42) {8};
\node[idxlbl] at (4.00,0.42) {9};
\node[idxlbl] at (4.50,0.42) {10};
\node[idxlbl] at (5.00,0.42) {11};
\node[idxlbl] at (5.50,0.42) {12};
\node[idxlbl] at (6.00,0.42) {13};
\node[idxlbl] at (6.50,0.42) {14};
\node[idxlbl] at (7.00,0.42) {15};
\node[idxlbl] at (7.50,0.42) {16};
\node[idxlbl] at (8.00,0.42) {17};
\node[idxlbl] at (8.50,0.42) {18};
\node[idxlbl] at (9.00,0.42) {19};
\node[idxlbl] at (9.50,0.42) {20};
\node[idxlbl] at (10.00,0.42) {21};
\node[stm-lab] at (-0.20,-0.75) {$s=0$};
\node[stm-ok] at (0.00,-0.75) {A};
\node[stm-ok] at (0.50,-0.75) {A};
\node[stm-ok] at (1.00,-0.75) {C};
\node[stm-ok] at (1.50,-0.75) {A};
\node[stm-ok] at (2.00,-0.75) {A};
\node[stm-ok] at (2.50,-0.75) {C};
\node[stm-ok] at (3.00,-0.75) {A};
\node[stm-ok] at (3.50,-0.75) {A};
\node[stm-bad] at (4.00,-0.75) {A};
\node[stm-idle] at (4.50,-0.75) {C};
\node[stm-note, anchor=west] at (4.95,-0.75) {比 9 次};
\node[stm-lab] at (-0.20,-1.40) {$s=3$};
\node[stm-known] at (1.50,-1.40) {A};
\node[stm-known] at (2.00,-1.40) {A};
\node[stm-known] at (2.50,-1.40) {C};
\node[stm-known] at (3.00,-1.40) {A};
\node[stm-known] at (3.50,-1.40) {A};
\node[stm-ok] at (4.00,-1.40) {C};
\node[stm-ok] at (4.50,-1.40) {A};
\node[stm-ok] at (5.00,-1.40) {A};
\node[stm-ok] at (5.50,-1.40) {A};
\node[stm-ok] at (6.00,-1.40) {C};
\node[stm-note, anchor=west] at (6.45,-1.40) {已知 5，比 5 次};
\node[stm-lab] at (-0.20,-2.05) {$s=10$};
\node[stm-known] at (5.00,-2.05) {A};
\node[stm-known] at (5.50,-2.05) {A};
\node[stm-known] at (6.00,-2.05) {C};
\node[stm-ok] at (6.50,-2.05) {A};
\node[stm-ok] at (7.00,-2.05) {A};
\node[stm-ok] at (7.50,-2.05) {C};
\node[stm-ok] at (8.00,-2.05) {A};
\node[stm-ok] at (8.50,-2.05) {A};
\node[stm-ok] at (9.00,-2.05) {A};
\node[stm-bad] at (9.50,-2.05) {C};
\node[stm-note, anchor=west] at (9.95,-2.05) {已知 3，比 7 次};
\node[stm-lab] at (-0.20,-2.70) {$s=17$};
\node[stm-known] at (8.50,-2.70) {A};
\node[stm-known] at (9.00,-2.70) {A};
\node[stm-bad] at (9.50,-2.70) {C};
\node[stm-idle] at (10.00,-2.70) {A};
\node[stm-note, anchor=west] at (10.55,-2.70) {已知 2，比 1 次};
\node[stm-lab] at (-0.20,-3.35) {$s=18$};
\node[stm-known] at (9.00,-3.35) {A};
\node[stm-bad] at (9.50,-3.35) {A};
\node[stm-idle] at (10.00,-3.35) {C};
\node[stm-note, anchor=west] at (10.55,-3.35) {已知 1，比 1 次};
\node[stm-lab] at (-0.20,-4.00) {$s=19$};
\node[stm-bad] at (9.50,-4.00) {A};
\node[stm-idle] at (10.00,-4.00) {A};
\node[stm-note, anchor=west] at (10.55,-4.00) {比 1 次};
\node[stm-lab] at (-0.20,-4.65) {$s=20$};
\node[stm-bad] at (10.00,-4.65) {A};
\node[stm-note, anchor=west] at (10.55,-4.65) {比 1 次};
\end{tikzpicture}
\diagcaption{KMP 在 $T = \mathtt{AACAACAACAAACAACAAAGT}$ 上試過的位移。 $s = 1, 2$、$s = 4$ 到 $9$、$s = 11$ 到 $16$ 都直接跳過，$T$ 的每個字元只有失配時才會再比一次。}
\label{fig:stm-kmp-align}
\end{minipage}\end{center}

$i = 20$ 那一列：已經對上 $9$ 個字元，$T[20] = \mathtt{G}$ 卻不是 $P[10] =
\mathtt{C}$。 $q$ 沿著鏈退成 $\pi[9] = 2$、$\pi[2] = 1$、$\pi[1] = 0$，每一次都接不上
\texttt{G}，最後 $q = 0$。 字母 \texttt{G} 根本不在 $P$ 裡，任何位移只要蓋到
$T[20]$ 都不可能成功，所以結果一定是 $0$；但演算法不知道這件事，還是一層一層退下來。

\begin{examnote}
KMP 的追蹤題，常見的兩種寫法是：像上表一個字元一列寫出 $q$ 的變化，或像
圖~\ref{fig:stm-kmp-align} 每換一個位移畫一列對齊。 兩種都要寫出失配時退到哪裡，
不能只寫最後的 $q$。 找到一次出現之後，記得 $q$ 是設成 $\pi[m]$，不是 $0$，否則會
漏掉重疊的出現。 題目問比較次數時，先說明你怎麼數（例如第 5 行和第 7 行比的是同一
對字元時只算一次），再照那個數法算。
\end{examnote}

\subsection{時間：平攤分析}

第 5 行的 \textbf{while} 迴圈在某些字元上可能跑很多次。 例如
$P = \mathtt{A}^{m-1}\mathtt{C}$、$T = \mathtt{A}^{m-1}\mathtt{G}\cdots$，讀到
\texttt{G} 時 $q$ 從 $m - 1$ 一路退到 $0$，退了 $m - 1$ 次。 所以單看一個字元，最
壞是 $\Theta(m)$；照第~1~章的方法把每個字元的最壞情況乘起來，只能得到
$O(nm)$。 這和第~3~章的 \proc{Multipop}（虛擬碼~\ref{alg:multipop}）是同一種情形：要退很多次，之前一定先往上加過
很多次。

\begin{theorem}
\label{thm:stm-kmp-time}
\proc{Compute-Failure} 的時間是 $\Theta(m)$，\proc{KMP-Match} 除了第 2 行以外的時
間是 $\Theta(n)$。 因此 KMP 的總時間是 $\Theta(n + m)$。
\end{theorem}

\begin{proof}
先看 \proc{KMP-Match}。 用第~\ref{sec:amo-potential}~節的位能法，位能取
$\Phi = q$。 一開始 $q = 0$，之後 $q$ 永遠不是負的，所以
$\Phi(D_i) \ge \Phi(D_0) = 0$，定理~\ref{thm:potential} 可以用。

第 $i$ 輪（讀 $T[i]$）的實際成本記成 $1 + w_i$，$w_i$ 是第 6 行執行的次數，其他的
行都是常數。 第 6 行每執行一次，$q$ 變成 $\pi[q] < q$，至少減少 $1$；第 8 行最多加
$1$；第 11 行只會讓 $q$ 變小。 所以這一輪的位能變化 $\Delta\Phi \le 1 - w_i$，平攤成
本
\[
\hat{c}_i = (1 + w_i) + \Delta\Phi \le (1 + w_i) + (1 - w_i) = 2.
\]
$n$ 輪的總成本不超過 $2n$，是 $O(n)$；迴圈至少跑 $n$ 輪，所以是 $\Theta(n)$。

\proc{Compute-Failure} 完全一樣，位能取 $\Phi = k$：$k$ 從 $0$ 開始、不會變成負的，
每執行一次第 6 行至少減 $1$，每一輪第 8 行最多加 $1$，所以 $m - 1$ 輪的總成本不超過
$2(m - 1)$，是 $\Theta(m)$。
\end{proof}

用第~\ref{sec:amo-aggregate}~節的聚合法也可以講：$q$ 每一輪最多加 $1$，總共最多加
$n$；每執行一次第 6 行至少減 $1$，而 $q$ 不會變成負的，所以總共減的量不超過加的
量，第 6 行全部加起來最多執行 $n$ 次。 位能 $\Phi = q$ 就是記帳法裡「已經對上的字
元身上存的錢」：每對上一個字元多存 $1$，失配退回時拿來付。

也可以直接數字元比較的次數。 把第 5 行和第 7 行比的同一對字元只算一次（C 程式就是
這樣寫），讀 $T[i]$ 時比較的次數是 $1 + w_i$，總數是 $n + \sum w_i \le 2n$。 和樸素
比對的 $(n - m + 1)m$ 比，差別在 $T$ 的指標 $i$ 從不往回走：每次比較不是讓 $i$ 前
進，就是讓 $P$ 往右移至少一格。

\subsection{C 程式}

C 的陣列從 $0$ 編號，失敗函數也跟著整段改：\texttt{pi[j]} 是 \texttt{p[0..j]} 最長
邊界的長度，也就是虛擬碼的 $\pi[j + 1]$。 $q$ 的意思不變（已經對上的字元數），所以
下一個要比的字元是 \texttt{p[q]}，失配時退到 \texttt{pi[q - 1]}。 為了數比較次數，
把 \textbf{while} 和後面的 \textbf{if} 合成一個迴圈：比一次，相同就前進並結束；不同
而且 $q = 0$ 也結束；否則退一格再比。 最後一行印出建失敗函數和比對各比了幾次，以及
$2m + 2n$。

\begin{example}[label=ex:stm-kmp]{c}{(kmp.c) 失敗函數與 KMP 比對，並計算比較次數}
#include <stdio.h>
#include <stdlib.h>
#include <string.h>

#define MAXLEN 100000

// 失敗函數。pi[j] 是 p[0..j] 最長邊界的長度，等於虛擬碼的第 j+1 格。
// 回傳字元比較的次數。
long long compute_failure(const char *p, int m, int *pi)
{
    long long cmp = 0;
    int k = 0;                  // 目前的邊界長度
    pi[0] = 0;
    for (int q = 1; q < m; q++) {
        for (;;) {
            cmp++;
            if (p[k] == p[q]) { // 邊界可以往後接一個字元
                k++;
                break;
            }
            if (k == 0) {
                break;
            }
            k = pi[k - 1];      // 換成次長的邊界再試
        }
        pi[q] = k;
    }
    return cmp;
}

// KMP 比對：印出每一個有效位移，回傳字元比較的次數。
long long kmp_match(const char *t, int n, const char *p, int m,
                    const int *pi)
{
    long long cmp = 0;
    int q = 0;                  // 已經對上的字元數
    int found = 0;
    printf("shifts:");
    for (int i = 0; i < n; i++) {
        for (;;) {
            cmp++;
            if (p[q] == t[i]) {
                q++;
                break;
            }
            if (q == 0) {
                break;
            }
            q = pi[q - 1];
        }
        if (q == m) {
            printf(" %d", i - m + 1);
            found++;
            q = pi[q - 1];      // 接著找可能重疊的下一次
        }
    }
    printf("%s\n", found ? "" : " none");
    return cmp;
}

int main(void)
{
    char *t = malloc(MAXLEN + 1);
    char *p = malloc(MAXLEN + 1);
    int *pi = malloc(MAXLEN * sizeof *pi);
    int cases;
    if (t == NULL || p == NULL || pi == NULL
        || scanf("%d", &cases) != 1) {
        free(t);
        free(p);
        free(pi);
        return 1;
    }
    while (cases--) {
        if (scanf("%100000s %100000s", t, p) != 2) {
            break;
        }
        int n = (int)strlen(t), m = (int)strlen(p);
        long long c1 = compute_failure(p, m, pi);
        printf("pi:");
        for (int j = 0; j < m; j++) {
            printf(" %d", pi[j]);
        }
        printf("\n");
        long long c2 = kmp_match(t, n, p, m, pi);
        printf("comparisons: %lld + %lld <= %d\n",
               c1, c2, 2 * (m + n));
    }
    free(t);
    free(p);
    free(pi);
    return 0;
}
\end{example}
\iosample{4 \\ AACAACAACAAACAACAAAGT AACAACAAAC \\ AACAACAAACAACAAAC AACAACAAAC \\ AAAAAAAAAAAA AAAAC \\ AAAAAAAAAAAA AAAA}{pi: 0 1 0 1 2 3 4 5 2 3 \\ shifts: 3 \\ comparisons: 12 + 25 <= 62 \\ pi: 0 1 0 1 2 3 4 5 2 3 \\ shifts: 0 7 \\ comparisons: 12 + 17 <= 54 \\ pi: 0 1 2 3 0 \\ shifts: none \\ comparisons: 7 + 20 <= 34 \\ pi: 0 1 2 3 \\ shifts: 0 1 2 3 4 5 6 7 8 \\ comparisons: 3 + 12 <= 32}

第一組是上面的手算例子，比對用了 $25$ 次比較，建失敗函數用了 $12$ 次。 第二組的兩次出現重疊了 $3$ 個字
元（$\pi[10] = 3$），找到第一次之後 $q$ 設成 $3$ 才找得到第二次。 後兩組是樸素比對的
最壞情況：範例~\ref{ex:stm-naive} 分別比了 $40$ 次和 $36$ 次，這裡只要 $20$ 次和
$12$ 次。 $n$ 越大差得越多：樸素比對是 $(n - m + 1)m$，KMP 不超過 $2n$。 每一組都在定理~\ref{thm:stm-kmp-time} 的證明給的 $2(m - 1) + 2n$ 以內，程式印的
$2m + 2n$ 是更寬一點的上界。

\subsection{不同課本的失敗函數}

失敗函數的寫法，各本書不一樣。 有的從 $0$ 編號，有的值從 $-1$ 開始，有的記「下一個
要比的位置」而不是邊界的長度。 意思都一樣，換算的時候整排一起換。

\begin{examnote}
寫失敗函數之前先看題目的定義，下表是常見的四種，用本章的 $P = \mathtt{AACAACAAAC}$
對照。
\begin{center}
\small
\begin{tabular}{p{0.17\linewidth}p{0.42\linewidth}l}
\toprule
寫法 & 定義 & 值 \\
\midrule
本書（CLRS） & $\pi[q]$，$q = 1..m$：$P[1..q]$ 最長邊界的長度 & $0\ 1\ 0\ 1\ 2\ 3\ 4\ 5\ 2\ 3$ \\
C 程式、競賽常用 & \texttt{pi[j]}，$j = 0..m-1$：$= \pi[j + 1]$ & $0\ 1\ 0\ 1\ 2\ 3\ 4\ 5\ 2\ 3$ \\
Horowitz 系列 & $f(j)$，$j = 0..m-1$：$= \pi[j + 1] - 1$，沒有邊界時 $-1$ & $-1\ 0\ -1\ 0\ 1\ 2\ 3\ 4\ 1\ 2$ \\
$\mathit{next}$ 陣列 & $\mathit{next}[1] = 0$，$j \ge 2$ 時 $= \pi[j - 1] + 1$ & $0\ 1\ 2\ 1\ 2\ 3\ 4\ 5\ 6\ 3$ \\
\bottomrule
\end{tabular}
\end{center}
前兩種只是編號不同，值完全一樣。 Horowitz 等人的資料結構課本（C 版）用 $f$，比對時
在 $p_j$ 失配就改比 $p_{f(j-1)+1}$，等於本書的 $q \gets \pi[q]$。 $\mathit{next}[j]$
記的是「在 $P[j]$ 失配後，下一個拿來比的是 $P$ 的第幾個字元」，$0$ 表示 $T$ 往前一
格、從 $P[1]$ 重來。 有些書還有改良版（常叫 $\mathit{nextval}$）：如果
$P[\mathit{next}[j]] = P[j]$，跳到那裡一定又失配，就直接沿用
$\mathit{nextval}[\mathit{next}[j]]$；這個版本的值和上面四種都不同，題目沒有指定就
不要用。
\end{examnote}

\section{比對自動機}
\label{sec:stm-dfa}

KMP 失配時要沿著鏈退好幾步，才知道下一個 $q$ 是多少。 如果事先把「目前對上 $q$ 個、
讀到字元 $a$ 之後變成幾個」全部算好存成表，每個字元就只要查一次表。 這張表就是一台
\textbf{有限自動機}（finite automaton）：狀態是 $0, 1, \ldots, m$，代表目前對上了幾個
字元；從狀態 $q$ 讀到字元 $a$，移到狀態
\[
\delta(q, a) = \text{$P$ 的前綴之中，是 $P[1..q]\,a$ 的後綴的最長長度}.
\]
照定理~\ref{thm:stm-kmp-correct} 證明裡同樣的推理（把 $q$ 換成 $f(i - 1)$）可以得到
$f(i) = \delta(f(i - 1), T[i])$，所以讀完 $T[i]$ 之後的狀態就是 $f(i)$。 進入狀態 $m$ 就表示找到一次出現。

\begin{pseudo}[label=alg:stm-dfa-match]{用自動機比對}
\Procedure{Automaton-Match}{$T, n, \delta, m$}
    \State $q \gets 0$
    \For{$i \gets 1$ \To $n$}
        \State $q \gets \delta(q, T[i])$
        \If{$q = m$}
            \State 輸出 $i - m$
        \EndIf
    \EndFor
\EndProcedure
\end{pseudo}

這裡不需要特別處理找到之後的情形，$\delta(m, a)$ 本來就照定義算。 比對時間是
$\Theta(n)$，而且每個字元剛好查一次表，比 KMP 的平攤 $O(1)$ 更整齊。 代價是轉移表有
$(m + 1)|\Sigma|$ 格。 用失敗函數可以一格 $O(1)$ 填完：

\begin{pseudo}[label=alg:stm-dfa-build]{用失敗函數建轉移表}
\Function{Build-Automaton}{$P, m, \Sigma$}
    \State $\pi \gets$ \Call{Compute-Failure}{$P, m$}
    \For{每個 $a \in \Sigma$}
        \State $\delta(0, a) \gets 0$
    \EndFor
    \State $\delta(0, P[1]) \gets 1$
    \For{$q \gets 1$ \To $m$}
        \For{每個 $a \in \Sigma$}
            \If{$q < m$ \AND $a = P[q + 1]$}
                \State $\delta(q, a) \gets q + 1$
            \Else
                \State $\delta(q, a) \gets \delta(\pi[q], a)$ \Comment{$\pi[q] < q$，已經填好}
            \EndIf
        \EndFor
    \EndFor
    \State \Return $\delta$
\EndFunction
\end{pseudo}

第 11 行的理由：$a$ 接不上 $P[q + 1]$ 時，新的前綴一定是 $P[1..q]$ 某個邊界接上 $a$，
這和 KMP 沿著鏈退的情形一樣；而從 $\pi[q]$ 開始找，正是 $\delta(\pi[q], a)$ 的定義。
（$q = m$ 時也一樣，因為長度 $m + 1$ 的前綴不存在。） 前處理的時間是
$\Theta(m |\Sigma|)$。

$P = \mathtt{AACAACAAAC}$ 的轉移表如下。 \texttt{G} 和 \texttt{T} 不在 $P$ 裡，那
兩欄全部是 $0$，表上省略。 對照第~\ref{sec:stm-kmp}~節的追蹤：狀態 $9$ 讀到
\texttt{G} 直接到 $0$，KMP 要退三次；狀態 $8$ 讀到 \texttt{C} 到 $6$，就是 $i = 9$
那一列退到 $\pi[8] = 5$ 再加一。

\begin{center}
\small
\begin{tabular}{clcc}
\toprule
狀態 $q$ & 已對上的前綴 & $\delta(q, \mathtt{A})$ & $\delta(q, \mathtt{C})$ \\
\midrule
0 & $\varepsilon$ & 1 & 0 \\
1 & \texttt{A} & 2 & 0 \\
2 & \texttt{AA} & 2 & 3 \\
3 & \texttt{AAC} & 4 & 0 \\
4 & \texttt{AACA} & 5 & 0 \\
5 & \texttt{AACAA} & 2 & 6 \\
6 & \texttt{AACAAC} & 7 & 0 \\
7 & \texttt{AACAACA} & 8 & 0 \\
8 & \texttt{AACAACAA} & 9 & 6 \\
9 & \texttt{AACAACAAA} & 2 & 10 \\
10 & \texttt{AACAACAAAC} & 4 & 0 \\
\bottomrule
\end{tabular}
\end{center}

KMP 可以看成「不存整張表，需要時才沿著鏈現算」的自動機：前處理只要 $\Theta(m)$，和
$|\Sigma|$ 無關，比對時間仍是 $\Theta(n)$。 字母表很大時（例如 Unicode），KMP 比較
划算；字母表很小、同一個模式要在很多段文字裡找時，自動機每個字元固定一次查表，常數
比較小。

本章的四種方法整理如下（$v$ 是有效位移的個數）。

\begin{center}
\small
\begin{tabular}{lllll}
\toprule
方法 & 前處理 & 比對（最壞） & 比對（期望或典型） & 額外空間 \\
\midrule
樸素比對 & 無 & $\Theta((n - m + 1)m)$ & 隨機文字約 $O(n)$ & $O(1)$ \\
Rabin--Karp & $\Theta(m)$ & $\Theta((n - m + 1)m)$ & $O(n + m(v + n/q))$ & $O(1)$ \\
比對自動機 & $\Theta(m|\Sigma|)$ & $\Theta(n)$ & $\Theta(n)$ & $\Theta(m|\Sigma|)$ \\
KMP & $\Theta(m)$ & $\Theta(n)$ & $\Theta(n)$ & $\Theta(m)$ \\
\bottomrule
\end{tabular}
\end{center}

\section{本章常見錯誤}
\label{sec:stm-pitfalls}

\begin{center}
\small
\begin{tabular}{p{0.26\linewidth}p{0.50\linewidth}p{0.14\linewidth}}
\toprule
錯誤 & 例子 & 見 \\
\midrule
失敗函數的定義混用 & 題目要 Horowitz 的 $f$（從 $-1$ 開始），卻寫了 $\pi$；或把
$\pi$ 的編號從 $0$ 開始，整排差一格 & 第~\ref{sec:stm-kmp}~節 \\
失配時退錯格 & 虛擬碼寫成 $q \gets \pi[q - 1]$；C 從 $0$ 編號卻寫
\texttt{q = pi[q]} & 第~\ref{sec:stm-kmp}~節 \\
\textbf{while} 寫成 \textbf{if} & 只退一次就去比 $P[q + 1]$，$i = 20$ 那一步會停在
$q = 2$ 而不是 $0$，之後的結果全錯 & 圖~\ref{fig:stm-kmp-align} \\
找到之後把 $q$ 設成 $0$ & $P = \mathtt{AACAACAAAC}$ 在 $\mathtt{AACAACAAACAACAAAC}$
裡只找到位移 $0$，漏掉位移 $7$ & 範例~\ref{ex:stm-kmp} \\
以為 KMP 每個字元都是 $O(1)$ & 單一字元可能退 $\Theta(m)$ 次；$O(1)$ 是平攤，證明要
寫出位能或總和 & 定理~\ref{thm:stm-kmp-time} \\
Rabin--Karp 命中就輸出 & 沒有逐字確認，把偽命中當成出現位置 &
第~\ref{sec:stm-rk}~節 \\
滾動時出現負數 & C 的 \texttt{\%} 對負數得到負數，之後的雜湊值全錯；要先加 $q$ &
範例~\ref{ex:stm-rk} \\
$q$ 選成 $d$ 的次方 & $d = 4$、$q = 256$ 時，雜湊值只看最後 $4$ 個字元 &
觀念題~\ref{q:stm-power-mod} \\
把 \texttt{strlen} 寫在迴圈條件 & \texttt{for (i = 0; i < strlen(t); i++)} 每輪都
重數一次，樸素比對憑空多出 $\Theta(n^2)$ & 第~\ref{sec:stm-naive}~節 \\
\bottomrule
\end{tabular}
\end{center}

\section{觀念題}
\label{sec:stm-exercises}

題型和研究所考試、期中期末考的筆試相近，有選擇、計算、追蹤、證明和設計。 除非題目
另外說明，失敗函數用本章的 $\pi$，Rabin--Karp 的字元對應是 $\mathtt{A} = 0$、
$\mathtt{C} = 1$、$\mathtt{G} = 2$、$\mathtt{T} = 3$，第一個字元是最高位。 解答在本章
最後。

\begin{enumerate}
\item （選擇，可能不只一個答案）下列哪些敘述正確？
\begin{enumerate}
\item 樸素比對的最壞情況是 $\Theta((n - m + 1)m)$。
\item Rabin--Karp 的最壞情況是 $\Theta(n + m)$。
\item 對任何 $P$ 和 $1 \le q \le m$，$\pi[q] < q$。
\item $P$ 的字元兩兩不同時，$\pi[q] = 0$ 對所有 $q$ 成立。
\item KMP 讀進某一個 $T[i]$ 時，第 6 行最多執行 $O(1)$ 次。
\end{enumerate}

\item （計算）用樸素比對在 $T = \mathtt{GTAGTAGGTAGT}$ 裡找 $P = \mathtt{GTAGG}$。 寫
出每個位移比較了幾次、總比較次數，以及有效位移。

\item （計算）Rabin--Karp，$d = 4$、$q = 7$，$T = \mathtt{CTAGCTACGATAC}$，
$P = \mathtt{TAC}$。
\begin{enumerate}
\item 求 $h$ 和 $p \bmod 7$。
\item 用滾動的式子從 $t_0 \bmod 7$ 算出 $t_1 \bmod 7$，寫出中間的負數怎麼處理。
\item 列出所有命中的位移，指出哪些是偽命中。
\end{enumerate}

\item （計算）$P = \mathtt{GTAGTGTAG}$。
\begin{enumerate}
\item 用 \proc{Compute-Failure} 求 $\pi[1..9]$，寫出 $q = 6$ 那一步的比較過程。
\item 換算成 Horowitz 的 $f(0..8)$ 和 $\mathit{next}[1..9]$。
\end{enumerate}

\item （追蹤）用第 4 題的 $P$，以 \proc{KMP-Match} 掃描
$T = \mathtt{GTAGTAGTGTAGTGTAG}$。 寫出每次失配時 $q$ 退到哪裡、所有有效位移，以及
（照範例~\ref{ex:stm-kmp} 的數法）比對階段的字元比較次數。

\item （選擇）下列陣列哪些可能是某個字串的失敗函數 $\pi[1..]$？ 說明理由。
(A) $0, 0, 1, 2, 3, 1$ \quad (B) $0, 1, 0, 2$ \quad (C) $0, 1, 2, 3, 4$ \quad
(D) $0, 0, 1, 0, 2$

\item （證明）證明對 $2 \le q \le m$，$\pi[q] \le \pi[q - 1] + 1$。 等號何時成立？

\item （證明）用位能法證明 \proc{Compute-Failure} 的第 6 行總共最多執行 $m - 1$ 次。
寫出位能函數、它為什麼不會小於初始值，以及每一輪的平攤成本。

\item （構造）給一個長度 $m$ 的模式和一段文字，使得 KMP 讀到某一個字元時，第 6 行連
續執行 $m - 1$ 次。 再說明為什麼這不影響定理~\ref{thm:stm-kmp-time}。

\item （證明與計算）令 $p = m - \pi[m]$。
\begin{enumerate}
\item 證明 $P[i] = P[i + p]$ 對 $1 \le i \le m - p$ 都成立（$p$ 是 $P$ 的一個週期）。
\item 證明沒有比 $p$ 更小的正整數週期。
\item $P = \mathtt{CAGCAGCAGCAG}$，求 $\pi[12]$ 和 $p$。 $p$ 整除 $m$ 時，$P$ 是
$P[1..p]$ 重複 $m/p$ 次；這個 $P$ 是哪個字串重複幾次？
\end{enumerate}

\item （設計）改寫 \proc{KMP-Match}，改成數「彼此不重疊」的出現最多有幾次（例如
$P = \mathtt{AA}$、$T = \mathtt{AAAAA}$ 時答案是 $2$）。 只改一行，並說明為什麼由左往
右、能找到就算的做法會得到最多的個數。

\item （計算）$P = \mathtt{ACAG}$，字母表 $\{\mathtt{A}, \mathtt{C}, \mathtt{G},
\mathtt{T}\}$。
\begin{enumerate}
\item 求 $\pi$，再用 \proc{Build-Automaton} 填出完整的轉移表 $\delta(q, a)$，
$q = 0..4$。
\item 用自動機掃描 $T = \mathtt{CACACAGACAG}$，寫出每個字元之後的狀態和所有有效位
移。
\end{enumerate}

\item \label{q:stm-power-mod}（說明）Rabin--Karp 用 $d = 4$，有人為了用位元運算算得
快，取 $q = 256$。 證明這時 $t_s \bmod q$ 只和視窗的最後 $4$ 個字元有關。 $m = 10$ 時
這會造成什麼問題？ 和第~12~章引理~\ref{lem:hash-gcd} 前後討論的「除法的表格大小不要選
$2$ 的次方」比較。

\item （是非）「Rabin--Karp 的 $q$ 選一個比 $nm$ 大的質數，底數 $x$ 從 $0$ 到 $q - 1$
隨機挑，就保證不會有偽命中。」 這句話對嗎？ 如果不對，正確的敘述是什麼？
\end{enumerate}

\section{實作練習}
\label{sec:stm-practice}

下面是 Aki-Judge 上和本章相關的題目。 Rabin--Karp 目前沒有題目，可以用
範例~\ref{ex:stm-rk} 搭配範例~\ref{ex:stm-naive}，對同一段文字比較兩者的輸出，自己驗
證偽命中有沒有被濾掉。

\begin{center}
\begin{tabular}{llc}
\toprule
小節 & 題目 & 難度 \\
\midrule
樸素比對 & \akiref{brute-force-first-match}{暴力比對：第一次出現位置} & 暖身 \\
樸素比對 & \akiref{count-overlapping-occurrences-brute-force}{暴力比對：計算出現次數} & 暖身 \\
樸素比對 & \akiref{two-fragments-in-order}{先後出現的兩段字串} & 基本 \\
KMP & \akiref{first-match-with-prefix-function}{以前綴函數找第一次出現位置} & 基本 \\
KMP & \akiref{count-occurrences-with-prefix-function}{以前綴函數計算出現次數} & 基本 \\
\bottomrule
\end{tabular}
\end{center}

幾題的提示：

\begin{itemize}
\item 前兩題把範例~\ref{ex:stm-naive} 的 \texttt{naive\_match} 改成找到第一個就回傳，或
改成計數。 輸出的位置從 $0$ 開始，就是本章的位移 $s$。
\item 「先後出現的兩段字串」要在每個單字裡找 $A$ 最早的出現，再從它結束的位置往後
找 $B$。 把樸素比對改成「從某個起點開始找第一次出現」的函式，呼叫兩次就好。 為什麼
只要看 $A$ 最早的那次，題目的說明已經講了。
\item 後兩題照範例~\ref{ex:stm-kmp}。 計數那題的出現可以重疊，找到之後 $q$ 要設成
\texttt{pi[m - 1]}，不是 $0$。
\end{itemize}

\section{讀完你應該會}
\label{sec:stm-checklist}

\begin{itemize}
\item 用位移的語言寫出字串比對問題，算出樸素比對在給定輸入上的比較次數，並舉出達到
$\Theta((n - m + 1)m)$ 的輸入。
\item 用 Rabin--Karp 一個視窗一列算出 $t_s \bmod q$，處理滾動時的負數，分出有效位移
和偽命中；說明期望時間的假設，以及隨機底數時偽命中機率 $(m - 1)/q$ 的證明。
\item 用 \proc{Compute-Failure} 手算失敗函數，寫出每一步沿著邊界的鏈退到哪裡，並在
$\pi$、C 的 \texttt{pi}、Horowitz 的 $f$、$\mathit{next}$ 之間換算。
\item 手算 KMP 的比對過程（$q$ 的變化表或對齊圖），找出所有重疊的出現。
\item 用位能法或聚合法證明 KMP 是線性時間，說明為什麼單一字元的最壞成本是
$\Theta(m)$ 但平攤是 $O(1)$。
\item 用失敗函數建出比對自動機的轉移表，比較自動機和 KMP 的前處理成本。
\end{itemize}

\textbf{延伸閱讀}（書末會整理完整的參考資料）

\begin{itemize}
\item Thomas H. Cormen 等，\textit{Introduction to Algorithms}，第 4 版，第 32 章
（32.1--32.4 節）：樸素比對、Rabin--Karp、比對自動機和 KMP，自動機和失敗函數的正確
性證明寫得比本章完整；32.5 節的後綴陣列是本章沒有講的另一條路。
\item Donald E. Knuth、James H. Morris Jr.、Vaughan R. Pratt，``Fast Pattern Matching in
Strings''，\textit{SIAM Journal on Computing} 6(2)，1977：KMP 的原始論文，也講了
$\mathit{next}$ 的改良版。
\item Richard M. Karp、Michael O. Rabin，``Efficient Randomized Pattern-Matching
Algorithms''，\textit{IBM Journal of Research and Development} 31(2)，1987：
Rabin--Karp 的原始論文，用隨機挑選的質數當模數。
\item Robert S. Boyer、J Strother Moore，``A Fast String Searching Algorithm''，
\textit{Communications of the ACM} 20(10)，1977：從模式的右邊往左比，實務上常比 KMP
快，文字編輯器的搜尋常用它的變形。
\item Steven S. Skiena，\textit{The Algorithm Design Manual}，第 3 版，6.7 節與 21.3
節：Rabin--Karp 的直觀講法，以及實務上該選哪一種比對方法。
\item Antti Laaksonen，\textit{Competitive Programmer's Handbook}，第 26 章：字串雜湊
和 Z 演算法，後者是和 KMP 等價、但想法不同的線性時間方法。
\item Dan Gusfield，\textit{Algorithms on Strings, Trees, and Sequences}，Cambridge
University Press，1997：字串演算法的專書，從 Z 演算法一路講到後綴樹和生物序列的應用。
\end{itemize}

\addcontentsline{toc}{section}{觀念題解答}
\section*{觀念題解答}

\begin{enumerate}
\item (a)、(c)、(d) 正確。 (a)：第~\ref{sec:stm-naive}~節的 $\mathtt{A}^n$ 和
$\mathtt{A}^{m-1}\mathtt{C}$。 (b)：錯，$T = \mathtt{A}^n$、$P = \mathtt{A}^m$ 時每
個位移都要逐字確認，是 $\Theta((n - m + 1)m)$。 (c)：邊界是真前綴。 (d)：長度 $r \ge 1$
的邊界要有 $P[1] = P[q - r + 1]$，字元兩兩不同時只能 $q - r + 1 = 1$，即 $r = q$，不
是真前綴。 (e)：錯，單一字元最多可以退 $m - 1$ 次（第 9 題）；$O(1)$
是平攤的結果。

\item 位移 $0$ 到 $7$ 的比較次數依序是 $5$、$1$、$1$、$5$、$1$、$1$、$2$、$5$，共 $21$ 次。 有效位移只有
$s = 3$（$T[4..8] = \mathtt{GTAGG}$）。 $s = 0$ 對上前 $4$ 個字元，第 $5$ 個
$T[5] = \mathtt{T} \ne \mathtt{G}$；$s = 6$ 對上 \texttt{G} 之後 $T[8] = \mathtt{G}
\ne \mathtt{T}$，比 $2$ 次；$s = 7$ 是 \texttt{GTAGT}，前 $4$ 個相同，比 $5$ 次。

\item
\begin{enumerate}
\item $h = 4^2 \bmod 7 = 16 \bmod 7 = 2$。 $p = 3 \cdot 16 + 0 \cdot 4 + 1 = 49$，
$49 \bmod 7 = 0$。
\item $t_0 = \mathtt{CTA} = 1 \cdot 16 + 3 \cdot 4 + 0 = 28$，$28 \bmod 7 = 0$。 離開的
是 $T[1] = \mathtt{C} = 1$，進來的是 $T[4] = \mathtt{G} = 2$：
$t_1 \equiv 4\,(0 - 1 \cdot 2) + 2 = -6 \equiv -6 + 7 = 1 \pmod 7$。 驗算：
$\mathtt{TAG} = 48 + 0 + 2 = 50$，$50 \bmod 7 = 1$。
\item 各位移的 $t_s \bmod 7$：
\begin{center}
\small
\begin{tabular}{c*{11}{c}}
\toprule
$s$ & 0 & 1 & 2 & 3 & 4 & 5 & 6 & 7 & 8 & 9 & 10 \\
視窗 & \texttt{CTA} & \texttt{TAG} & \texttt{AGC} & \texttt{GCT} & \texttt{CTA} & \texttt{TAC} & \texttt{ACG} & \texttt{CGA} & \texttt{GAT} & \texttt{ATA} & \texttt{TAC} \\
$t_s \bmod 7$ & \textbf{0} & 1 & 2 & 4 & \textbf{0} & \textbf{0} & 6 & 3 & \textbf{0} & 5 & \textbf{0} \\
\bottomrule
\end{tabular}
\end{center}
命中的是 $s = 0, 4, 5, 8, 10$。 $s = 5$ 和 $s = 10$ 是 \texttt{TAC}，有效；$s = 0$、
$s = 4$（\texttt{CTA}）和 $s = 8$（\texttt{GAT}）是偽命中，共 $3$ 次。
\end{enumerate}

\item
\begin{enumerate}
\item $\pi = 0,\ 0,\ 0,\ 1,\ 2,\ 1,\ 2,\ 3,\ 4$。 $q = 6$：開始時 $k = \pi[5] = 2$，$P[3] = \mathtt{A} \ne
P[6] = \mathtt{G}$，退到 $k = \pi[2] = 0$；$P[1] = \mathtt{G} = P[6]$，$k = 1$，所以
$\pi[6] = 1$。 其他各步都沒有退：$q = 2, 3$ 時 $P[1] = \mathtt{G}$ 不同，是 $0$；
$q = 4, 5$ 接上 \texttt{G}、\texttt{T}；$q = 7, 8, 9$ 從 $k = 1$ 接上
\texttt{T}、\texttt{A}、\texttt{G}。
\item $f = -1,\ -1,\ -1,\ 0,\ 1,\ 0,\ 1,\ 2,\ 3$（每一項減 $1$）。 $\mathit{next} = 0,\ 1,\ 1,\ 1,\ 2,\ 3,\ 2,\ 3,\ 4$（$\mathit{next}[1] =
0$，之後是前一格的 $\pi$ 加 $1$）。
\end{enumerate}

\item $q$ 的變化：$i = 1$ 到 $5$ 一路對上，$q = 5$。 $i = 6$ 讀到 \texttt{A}，
$P[6] = \mathtt{G}$ 不同，退到 $\pi[5] = 2$，$P[3] = \mathtt{A}$ 相同，$q = 3$。 之後
$i = 7$ 到 $12$ 全部對上，$i = 12$ 時 $q = 9$，輸出位移 $12 - 9 = 3$，$q \gets \pi[9]
= 4$。 $i = 13$ 到 $17$ 又全部對上，$i = 17$ 時 $q = 9$，輸出位移 $8$，$q \gets 4$。
有效位移是 $3$ 和 $8$，兩次出現重疊了 $4$ 個字元（\texttt{GTAG}）。 比較次數：每個字元
比 $1$ 次，加上失配退一格的 $1$ 次，共 $17 + 1 = 18$ 次。 可以用範例~\ref{ex:stm-kmp}
驗證，輸入 \texttt{1}、\texttt{GTAGTAGTGTAGTGTAG GTAGTGTAG}，程式印出
\texttt{comparisons: 9 + 18 <= 52}。

\item (A)、(C) 可能，(B)、(D) 不可能。 (A) 例如 \texttt{CGCGCC}：\texttt{CGCGC} 的最長邊界 \texttt{CGC} 後面是
\texttt{G}，接不上 \texttt{C}；退到邊界 \texttt{C}，後面也是 \texttt{G}；退到空字
串，$P[1] = \mathtt{C}$ 接上，得 $1$。 (C) 是 \texttt{AAAAA}。 (B) 的
$\pi[4] = 2 > \pi[3] + 1 = 1$，違反第 7 題。 (D) 的 $\pi[5] = 2 > \pi[4] + 1 = 1$，也
違反。

\item 令 $r = \pi[q]$。 $r = 0$ 時顯然成立。 $r \ge 1$ 時，$P[1..r]$ 是 $P[1..q]$ 的邊
界，拿掉兩邊的最後一個字元，$P[1..r - 1]$ 是 $P[1..q - 1]$ 的前綴也是後綴，而且
$r - 1 < q - 1$，所以是 $P[1..q - 1]$ 的邊界，$r - 1 \le \pi[q - 1]$。 等號成立若且唯
若 $P[\pi[q - 1] + 1] = P[q]$，也就是最長的邊界直接接得上，\proc{Compute-Failure} 第
5 行一次都不用退。

\item 位能取 $\Phi = k$。 一開始 $k = 0$；第 6 行把 $k$ 換成 $\pi[k] \ge 0$，第 8 行加
一，$k$ 永遠不小於 $0$，所以 $\Phi \ge \Phi_0 = 0$。 第 $q$ 輪（$q = 2..m$）第 6 行執行
$w_q$ 次，每次 $k$ 至少減 $1$，第 8 行最多加 $1$，所以 $\Delta\Phi \le 1 - w_q$。 把
「第 6 行執行一次」的成本記成 $1$、其他不計，這一輪的平攤成本是
$w_q + \Delta\Phi \le 1$。 由定理~\ref{thm:potential}，$m - 1$ 輪的總成本
$\sum w_q \le m - 1$。

\item $P = \mathtt{A}^{m-1}\mathtt{C}$，$\pi = 0, 1, 2, \ldots, m - 2, 0$；
$T = \mathtt{A}^{m-1}\mathtt{G}$。 讀完前 $m - 1$ 個 \texttt{A} 時 $q = m - 1$；讀到
\texttt{G}，$P[m] = \mathtt{C}$ 不同，退到 $\pi[m - 1] = m - 2$，$P[m - 1] =
\mathtt{A}$ 也不同，再退……一直退到 $0$，共 $m - 1$ 次。 但要退 $m - 1$ 次，$q$ 之前
得先加到 $m - 1$，這要讀 $m - 1$ 個字元，每個字元只付 $1$ 次比較。 平攤下來每個字元
仍是 $O(1)$，定理~\ref{thm:stm-kmp-time} 的位能論證本來就容許某一輪很貴。

\item
\begin{enumerate}
\item $k = \pi[m]$，$P[1..k]$ 是 $P$ 的後綴，所以 $P[i] = P[m - k + i] = P[i + p]$
對 $1 \le i \le k = m - p$ 成立。
\item 若 $p' < p$ 是週期，$P[i] = P[i + p']$ 對 $1 \le i \le m - p'$ 成立，等於說
$P[1..m - p']$ 和 $P[p' + 1 .. m]$ 相同，也就是 $P$ 有長度 $m - p'$ 的邊界
（$p' \ge 1$，所以是真前綴）。 但 $m - p' > m - p = \pi[m]$，和 $\pi[m]$ 是最長的矛盾。
\item $\pi = 0,\ 0,\ 0,\ 1,\ 2,\ 3,\ 4,\ 5,\ 6,\ 7,\ 8,\ 9$，$\pi[12] = 9$，$p = 3$。 $3$ 整除 $12$，$P$ 是
\texttt{CAG} 重複 $4$ 次。
\end{enumerate}

\item 把第 11 行的 $q \gets \pi[q]$ 改成 $q \gets 0$：找到一次以後，下一次出現只能從
$i + 1$ 開始，前面的字元一個都不能共用。 由左往右、能找到就算是最好的，理由是第~19~章的
交換論證：設某個最佳解選的第一次出現結束得比貪心選的晚（或一樣），把它換成貪心選的那
次，結束位置只會更早，後面的出現都還能用，個數不變。 一直換下去，貪心的個數不少於最
佳解。 $T = \mathtt{AAAAA}$、$P = \mathtt{AA}$ 時，貪心選位移 $0$、$2$，答案是 $2$。

\item
\begin{enumerate}
\item $\pi = 0,\ 0,\ 1,\ 0$。 轉移表（列是 $q$，欄是讀到的字元）：
\begin{center}
\small
\begin{tabular}{ccccc}
\toprule
$q$ & \texttt{A} & \texttt{C} & \texttt{G} & \texttt{T} \\
\midrule
0 & 1 & 0 & 0 & 0 \\
1 & 1 & 2 & 0 & 0 \\
2 & 3 & 0 & 0 & 0 \\
3 & 1 & 2 & 4 & 0 \\
4 & 1 & 0 & 0 & 0 \\
\bottomrule
\end{tabular}
\end{center}
例如 $\delta(3, \mathtt{C})$：$P[4] = \mathtt{G}$ 接不上，查 $\delta(\pi[3], \mathtt{C})
= \delta(1, \mathtt{C}) = 2$。 $\delta(4, \cdot)$ 全部照抄 $\delta(\pi[4], \cdot) =
\delta(0, \cdot)$。
\item 狀態依序是 $0$、$1$、$2$、$3$、$2$、$3$、$4$、$1$、$2$、$3$、$4$。 到達狀態 $4$ 的是 $i = 7$ 和 $i = 11$，有效位移
$3$ 和 $7$。
\end{enumerate}

\item $q = 256 = 4^4$。 視窗的值 $t_s = \sum_{j=1}^{m} T[s + j]\, 4^{m-j}$ 裡，
$j \le m - 4$ 的項都含 $4^4$ 這個因數，除以 $256$ 的餘數是 $0$，所以
$t_s \bmod 256$ 只由最後 $4$ 項，也就是最後 $4$ 個字元決定。 $m = 10$ 時，任何最後
$4$ 個字元和 $P$ 相同的視窗都會命中，前 $6$ 個字元完全沒被檢查。 DNA 這類重複很多的
資料，偽命中會大量出現，每次都要花 $O(m)$ 確認，時間退化成接近樸素比對。 這和第~12~章
除法的表格大小取 $2^p$ 時只用到鍵的最低 $p$ 個位元是同一件事：模數整除底數的某個次方
時（這裡 $256 = 4^4$），高位的資訊全部丟掉。 所以 $q$ 要選不整除 $d$ 的質數，實際上就是
選一個比 $d$ 大的質數。

\item 不對。 定理~\ref{thm:stm-rk-random} 給的是機率的上界：每個不同的視窗偽命中的機
率不超過 $(m - 1)/q < 1/n$，所以偽命中的\textbf{期望}次數小於 $(n - m + 1)/n \le 1$。
運氣不好時還是可能發生，只是很少；正因為如此，第 10 行的逐字確認不能拿掉。 要拿掉確
認，就得接受結果偶爾會錯（Monte Carlo 型）。
\end{enumerate}

\chapter{NP 完備}
\label{ch:npc}

% 本章的圖共用的樣式（前綴 npc-）
\tikzset{
    npc-v/.style={circle, draw=black!70, thick, minimum size=0.62cm, inner sep=0pt, font=\small, fill=white},
    npc-vA/.style={npc-v, draw=coreStroke, fill=coreFill, very thick},
    npc-vB/.style={npc-v, draw=memStroke, fill=memFill, very thick},
    npc-l/.style={circle, draw=black!70, thick, minimum size=0.8cm, inner sep=0pt, font=\small, fill=white},
    npc-lsel/.style={npc-l, draw=coreStroke, fill=coreFill, very thick},
    npc-e/.style={draw=black!65, thick},
    npc-x/.style={draw=memStroke, thick, dashed},
    npc-ham/.style={draw=coreStroke, line width=1.8pt},
    npc-t/.style={font=\small},
    npc-box/.style={draw=black!70, thick, rounded corners=2pt, minimum height=0.8cm, inner sep=4pt, align=center, font=\small, fill=white},
    npc-pbox/.style={npc-box, font=\footnotesize\sffamily, minimum height=0.6cm, inner sep=3pt},
    npc-arr/.style={-{Stealth[length=2.2mm]}, draw=black!70, thick},
    npc-darr/.style={npc-arr, dashed},
}
% 問題的名稱（決策問題用無襯線字）
\newcommand{\npcP}[1]{\textsf{#1}}

教務處要替二十門課排期末考。 同一個學生修的兩門課不能排在同一個時段，可是考場只借到三個時段。
助理排了一個晚上還排不出來，開始懷疑根本排不下。 如果有人拿一份排好的表給你，你只要一門課一門
課核對有沒有衝突，幾分鐘就能確定它對不對；要自己從頭排，或是要確定「真的排不下」，好像就只能一
種一種試。 檢查一個答案和找出一個答案，難度真的差這麼多嗎？ 往下讀之前，先想一想。

前面二十幾章的問題幾乎都有多項式時間的演算法：排序 $\Theta(n \lg n)$、最短路徑、最小生成樹、強
連通元件都是 $n$ 和 $m$ 的低次多項式。 也有幾個問題一直沒解決：第~18~章的最長簡單路徑、第~20~
章的圖著色和旅行推銷員問題，我們只能用回溯或分支限界去搜，最壞情況是指數時間。 這一章說明這些問
題為什麼一直找不到快的演算法，以及怎麼\textbf{證明}一個新問題至少和它們一樣難。

這一章以觀念和證明為主。 先定義 P 和 NP，再定義多項式時間歸約，接著用一串歸約證明十幾個常見
問題都是 NP 完備的；每個歸約都會證兩個方向，並拿一個自己出的小公式或小圖實際做一次。 最後回到
第~\ref{sec:ana-input-size}~節的擬多項式時間，說明為什麼 0/1 背包有 $\Theta(nW)$ 的動態規劃，卻
仍然是 NP 完備。 C 程式只有兩支，用暴力法驗證歸約的兩邊答案一致，用到《C 程式設計基礎》第~12~章
的位元運算、第~14~章的結構和第~19~章的 \texttt{malloc}。

\section{決策問題與 P}
\label{sec:npc-decision}

\subsection{最佳化問題和決策問題}

本書大部分的問題都在求「最好」的答案：最短的路徑、總重最小的生成樹、價值最大的背包。 這種問題
叫\textbf{最佳化問題}（optimization problem）。 這一章改看\textbf{決策問題}（decision problem）：
答案只有「是」或「否」。 每個最佳化問題都可以加一個門檻變成決策問題：

\begin{center}
\begin{tabular}{p{0.24\linewidth}p{0.66\linewidth}}
\toprule
最佳化問題 & 決策版本 \\
\midrule
最短路徑（第~25~章）& 給圖 $G$、兩點 $s$、$t$ 和整數 $L$，有沒有從 $s$ 到 $t$、長度不超過 $L$ 的
路徑？ \\
最大團 & 給圖 $G$ 和整數 $k$，有沒有 $k$ 個兩兩相鄰的頂點？ \\
0/1 背包（第~18~章）& 給重量 $w_i$、價值 $v_i$、容量 $W$ 和目標 $V$，能不能挑一些物品，總重不超
過 $W$、總價值至少 $V$？ \\
旅行推銷員（第~20~章）& 給成本矩陣和整數 $B$，有沒有一條經過每個城市恰好一次再回到起點、總成本
不超過 $B$ 的路線？ \\
\bottomrule
\end{tabular}
\end{center}

決策版本不會比最佳化版本難：最佳值算出來之後，和門檻比一下就知道答案。 所以只要證明某個決策版本
「很難」，它的最佳化版本也一樣難。 反過來，決策版本會解的話，最佳化版本通常也只多花多項式倍的時
間，第~\ref{sec:npc-search}~節會示範。 只看決策問題的好處是答案的形式固定，下一節「驗證答案」
和第~\ref{sec:npc-reduce}~節「把一個問題轉成另一個問題」都比較好說清楚。

一個決策問題的每一個輸入叫一個\textbf{實例}（instance）。 答案是「是」的實例叫\textbf{是實例}，
答案是「否」的叫\textbf{否實例}。 例如「圖~$G$、$k = 3$」是團問題的一個實例，$G$ 裡有三個兩兩相
鄰的頂點時，它是是實例。

\subsection{輸入大小與 P}

第~\ref{sec:ana-input-size}~節說過，輸入大小要用寫下輸入所需的位元數來量。 這一章所有的「多項
式時間」都是指\textbf{輸入的位元數}的多項式。 寫下輸入的方式叫\textbf{編碼}（encoding），只要編
碼合理，換一種寫法只會讓長度差多項式倍，不影響「是不是多項式」：

\begin{itemize}
\item 整數用二進位（或十進位）寫。 數值 $x$ 只占 $\Theta(\lg x)$ 個位元。 用一進位（寫 $x$ 個
$1$）不算合理的編碼，它會把長度灌大成 $x$，第~\ref{sec:npc-pseudo}~節會看到這個差別有多重要。
\item 圖用鄰接串列或鄰接矩陣都可以。 前者長度約 $(n + m)\lg n$ 位元，後者 $n^2$ 位元，彼此是多項
式的關係。
\item 布林公式照字面寫下來，長度和變數、子句的個數成正比（再乘上寫一個變數編號的位元數）。
\end{itemize}

\begin{definition}[P]
\label{def:npc-p}
\textbf{P} 是所有能在多項式時間內解決的決策問題所成的集合：問題 $A$ 屬於 P，如果有一個演算法和
一個常數 $c$，對 $A$ 的每一個長度為 $N$ 的實例，演算法都在 $O(N^c)$ 步內答出正確的「是」或「否」。
\end{definition}

本書前面的例子有很多在 P 裡：第~21~章判斷一個圖是不是二分圖（也就是能不能用兩種顏色著色）、第~23~
章的 2-SAT、第~24~章的最小生成樹、第~25~章的最短路徑，它們的決策版本都在 P 裡。

為什麼用「多項式」當分界？ 主要是三個理由。 第一，多項式在組合之後還是多項式：多項式時間的程序
呼叫多項式次多項式時間的副程式，總時間仍是多項式，下一節的歸約就靠這一點。 第二，各種合理的計算
模型（RAM 模型、圖靈機）之間的模擬只差多項式倍，P 不會因為換一種模型而改變。 第三，實際上遇到的
多項式演算法次數都不大。 反過來看，指數時間的演算法在 $N$ 稍大時就完全不能用，這個分界很少錯判。
當然 $\Theta(N^{100})$ 也算多項式，「在 P 裡」不代表「實際上很快」；這只是理論上的分界。

\subsection{決策版本夠用嗎}
\label{sec:npc-search}

假設有一個程序 \proc{Clique}$(G, k)$ 能在多項式時間內回答「$G$ 有沒有大小 $k$ 的團」。 用它可以求
出最大團的大小，甚至把團找出來：

\begin{pseudo}[label=alg:npc-find-clique]{用決策版本找出最大團}
\Function{Find-Clique}{$G$}
    \State $k \gets n$
    \While{\proc{Clique}$(G, k)$ 回答「否」}
        \State $k \gets k - 1$
    \EndWhile
    \For{$v \gets 1$ \To $n$}
        \If{\proc{Clique}$(G - v, k)$ 回答「是」}
            \State $G \gets G - v$ \Comment{不用 $v$ 也有大小 $k$ 的團}
        \EndIf
    \EndFor
    \State \Return $G$ 剩下的頂點
\EndFunction
\end{pseudo}

$G - v$ 表示刪掉頂點 $v$ 和它關聯的邊。 第 3、4 行從 $n$ 往下試，停下時 $k$ 就是最大團的大小（至
少有一個頂點時，$k = 1$ 一定回答「是」）。 第 5 到 7 行一個一個試著刪頂點，刪掉之後還有大小 $k$
的團就真的刪掉。

\begin{whybox}[剩下的頂點恰好是一個團]
令 $G'$ 是迴圈結束後的圖。 每次刪除都保證刪完的圖還有大小 $k$ 的團，所以 $G'$ 有一個大小 $k$ 的
團 $Q$。 假設 $G'$ 還有 $Q$ 以外的頂點 $u$。 迴圈檢查 $u$ 的時候，當時的圖包含 $G'$ 的全部頂點
（之後只會再刪），也就包含 $Q$，而 $u \notin Q$，所以刪掉 $u$ 之後 $Q$ 還在，第 6 行應該回答
「是」並刪掉 $u$，矛盾。 因此 $G'$ 的頂點恰好是 $Q$。
\end{whybox}

\proc{Clique} 最多被呼叫 $2n$ 次，每次的輸入都不比原來的圖大，所以只要 \proc{Clique} 是多項式時間，
整個程序也是。 這種「用決策版本把答案一塊一塊確定下來」的做法叫\textbf{自我歸約}（self-reduction），
大部分 NP 完備問題都可以這樣做。

答案是一個數值的時候，要改用二分搜尋。 旅行推銷員問題的成本都是不超過 $C$ 的非負整數時，最佳成
本在 $0$ 到 $nC$ 之間，二分搜尋只要問決策版本 $\lceil \lg(nC + 1) \rceil$ 次，這是輸入位元數的多
項式（$\lg C$ 不超過寫下一個成本的位元數）。 如果從 $B = 0, 1, 2, \ldots$ 一個一個問，最多要問
$nC + 1$ 次，這是 $C$ 的數值的多項式，卻是位元數的指數，第~\ref{sec:npc-pseudo}~節會再遇到這個
差別。

\section{NP：答案可以快速驗證的問題}
\label{sec:npc-np}

\subsection{憑證與驗證者}

回到章首的排考。 「三個時段排得下嗎？」是一個決策問題。 答案是「是」的時候，有一個很短的證據可以
說服別人：一份排好的時段表。 拿到這份表，逐一檢查每一對有共同學生的課，時間和課程對數成正比。
這份證據叫\textbf{憑證}（certificate），檢查憑證的演算法叫\textbf{驗證者}（verifier）。

\begin{definition}[NP]
\label{def:npc-np}
決策問題 $A$ 屬於 \textbf{NP}，如果有一個多項式時間的演算法 $V$（驗證者）和一個多項式 $p$，使得
對 $A$ 的每一個實例 $x$：
\[
x \text{ 是是實例} \iff \text{存在一個字串 } y\text{，} |y| \le p(|x|)\text{，而且 } V(x, y)
\text{ 回答「是」。}
\]
\end{definition}

這個定義有三個要點：

\begin{itemize}
\item \textbf{憑證要短}：長度是 $|x|$ 的多項式。 不能拿「把所有可能都列出來」當憑證。
\item \textbf{驗證要快}：$V$ 在 $|x|$ 的多項式時間內做完（$y$ 也只有多項式長，所以兩者一樣）。
\item \textbf{兩個方向都要成立}：是實例至少有一個憑證能讓 $V$ 回答「是」；否實例則不管給什麼 $y$，
$V$ 都回答「否」，不會被騙。
\end{itemize}

定義只要求「存在」憑證，沒有說要怎麼找到它。 NP 描述的是「答案是『是』的時候，能不能被快速確
認」，和找答案要花多久是兩回事。

\begin{center}
\begin{tabular}{p{0.2\linewidth}p{0.3\linewidth}p{0.4\linewidth}}
\toprule
問題 & 憑證 & 驗證 \\
\midrule
\npcP{SAT}：布林公式能不能滿足 & 每個變數的真假 & 代進去算每個子句，和公式長度成正比 \\
\npcP{CLIQUE}：有沒有大小 $k$ 的團 & $k$ 個頂點 & 檢查 $\binom{k}{2}$ 對都相鄰，$O(k^2)$（用鄰接
矩陣）\\
\npcP{HAM-CYCLE}：有沒有哈密頓環 & 頂點的一個排列 & 虛擬碼~\ref{alg:npc-verify-ham}，$\Theta(n)$ \\
\npcP{SUBSET-SUM}：能不能湊出 $t$ & 挑出的那些數 & 加起來和 $t$ 比較，位元數的多項式 \\
\npcP{3-COLOR}：能不能三著色 & 每個頂點的顏色 & 檢查每條邊兩端不同色，$O(n + m)$ \\
\npcP{COMPOSITE}：整數 $x$ 是不是合數 & 一個因數 $d$，$1 < d < x$ & 做一次除法，$\lg x$ 的多項式 \\
\bottomrule
\end{tabular}
\end{center}

\textbf{哈密頓環}（Hamiltonian cycle）是恰好經過每個頂點一次的環。 給無向圖 $G$（鄰接矩陣）和一
個頂點序列 $c = \langle c_1, \ldots, c_n \rangle$ 當憑證，驗證者這樣檢查：

\begin{pseudo}[label=alg:npc-verify-ham]{哈密頓環的驗證者}
\Function{Verify-Ham-Cycle}{$G, c$}
    \If{$c$ 的長度不是 $n$}
        \State \Return \FALSE
    \EndIf
    \State 建立陣列 $\mathit{used}[1..n]$，每一格都設成 \FALSE
    \For{$i \gets 1$ \To $n$}
        \If{$c_i < 1$ \OR $c_i > n$ \OR $\mathit{used}[c_i]$}
            \State \Return \FALSE \Comment{不是頂點，或重複了}
        \EndIf
        \State $\mathit{used}[c_i] \gets \TRUE$
    \EndFor
    \For{$i \gets 1$ \To $n$}
        \State $j \gets i \bmod n + 1$ \Comment{$c_n$ 的下一個是 $c_1$}
        \If{$(c_i, c_j) \notin E$}
            \State \Return \FALSE
        \EndIf
    \EndFor
    \State \Return \TRUE
\EndFunction
\end{pseudo}

第 5 到 8 行確認 $c$ 是 $1..n$ 的一個排列，第 9 到 12 行確認相鄰兩個頂點（含最後一個接回第一個）
之間都有邊。 兩個迴圈各 $n$ 輪，用鄰接矩陣查邊是 $O(1)$，總共 $\Theta(n)$。 $G$ 有哈密頓環時，把
環上的頂點照順序列出來就是一個會通過的憑證；$G$ 沒有哈密頓環時，任何通過第 5 到 12 行的序列都會
是一個哈密頓環，不存在，所以不管給什麼都回答 \FALSE。

\subsection{P 和 NP 的關係}

\begin{theorem}
\label{thm:npc-p-in-np}
$\mathrm{P} \subseteq \mathrm{NP}$。
\end{theorem}

\begin{proof}
設 $A \in \mathrm{P}$，演算法 $M$ 在多項式時間內解 $A$。 令驗證者 $V(x, y)$ 不看 $y$，直接執行
$M(x)$ 並回答一樣的答案。 $V$ 是多項式時間；$x$ 是是實例時，空字串 $y$ 就讓 $V$ 回答「是」；$x$
是否實例時，任何 $y$ 都讓 $V$ 回答「否」。 所以 $A \in \mathrm{NP}$。
\end{proof}

反過來，NP 的問題至少可以在指數時間內解決：把所有長度不超過 $p(|x|)$ 的字串都當憑證試一次，最多
$2^{p(|x|) + 1}$ 個，每個用多項式時間檢查。 第~20~章的回溯法其實就是在有系統地列舉憑證，只是靠剪枝
少試一些。

NP 這個名字是 nondeterministic polynomial time（非決定性多項式時間）的縮寫。 另一個等價的定義是：
有一台「每一步可以同時試好幾種選擇」的機器，只要其中一條路在多項式步內回答「是」就算「是」；這台
機器「猜」的那一串選擇，正好就是憑證。 \textbf{NP 不是 non-polynomial 的縮寫}，也不是「不能在多項式
時間內解決的問題」：P 裡的每個問題都在 NP 裡。

\textbf{P 等不等於 NP}，也就是「能快速驗證的問題，是不是都能快速解決」，是到現在還沒有答案的問
題。 大多數人相信 $\mathrm{P} \ne \mathrm{NP}$，但沒有人證明出來。 這一章所有「某問題沒有多項式
時間演算法」的說法，都要加上「除非 $\mathrm{P} = \mathrm{NP}$」。

定義~\ref{def:npc-np} 對是實例和否實例不對稱。 「公式可以滿足」有短憑證（一組指派）；「公式不能滿
足」要說服別人，目前不知道有沒有短的憑證，總不能把 $2^n$ 組指派全部列出來。 補問題（是、否對調）
在 NP 裡的問題所成的集合叫 \textbf{co-NP}。 $\mathrm{NP} = \text{co-NP}$ 是否成立也還不知道，所以
「$A$ 在 NP 裡，所以 $A$ 的補問題也在 NP 裡」是不能用的推論。

表裡的 \npcP{COMPOSITE} 說明了另一件事。 第~\ref{sec:ana-input-size}~節的試除法要 $\sqrt{x}$ 步，
是位元數的指數；但給一個因數當憑證，驗證只要一次除法。 2002 年 Agrawal、Kayal、Saxena 證明了判斷
質數（也就能判斷合數）其實在 P 裡；而真的把因數找出來（質因數分解）目前仍不知道有沒有多項式時間的演
算法。 判斷「有沒有因數」和「把因數找出來」，難度可以差很多。

\section{多項式時間歸約}
\label{sec:npc-reduce}

\subsection{定義}

要說「問題 $B$ 至少和問題 $A$ 一樣難」，做法是把 $A$ 的每個實例改寫成 $B$ 的一個實例，答案不變，
而且改寫很快。 這樣只要 $B$ 會解，$A$ 也就會解。

\begin{definition}[多項式時間歸約]
\label{def:npc-reduce}
決策問題 $A$ \textbf{多項式時間歸約}到決策問題 $B$，記作 $A \le_{\mathrm{p}} B$，如果有一個能在多
項式時間內算出的函數 $f$，把 $A$ 的實例變成 $B$ 的實例，而且對 $A$ 的\textbf{每一個}實例 $x$：
\[
x \text{ 是 } A \text{ 的是實例} \iff f(x) \text{ 是 } B \text{ 的是實例}.
\]
\end{definition}

\begin{center}
\begin{tikzpicture}[x=1cm, y=1cm]
\node[npc-box] (f) at (2.4, 0) {$f$\\[-1pt]{\footnotesize 多項式時間}};
\node[npc-box] (b) at (6.4, 0) {解 $B$ 的\\[-1pt]演算法};
\node[npc-t] (x) at (0, 0) {$x$};
\node[npc-t] (yes) at (9.3, 0.45) {是};
\node[npc-t] (no) at (9.3, -0.45) {否};
\draw[npc-arr] (x) -- (f);
\draw[npc-arr] (f) -- node[above, font=\small] {$f(x)$} (b);
\draw[npc-arr] (b.east) -- (yes);
\draw[npc-arr] (b.east) -- (no);
\draw[black!50, dashed, rounded corners=4pt] (1.1, -1.0) rectangle (8.0, 1.05);
\node[font=\small, text=black!60, anchor=north west] at (1.15, -1.05) {解 $A$ 的演算法};
\end{tikzpicture}
\diagcaption{歸約 $A \le_{\mathrm{p}} B$：先用 $f$ 把 $x$ 改寫成 $B$ 的實例，再交給解 $B$ 的演算
法，它的答案就是 $x$ 的答案。}
\label{fig:npc-reduce}
\end{center}

定義裡的「$\iff$」有兩個方向，證明一個歸約時兩個都要證：

\begin{enumerate}
\item $x$ 是是實例 $\Rightarrow$ $f(x)$ 是是實例；
\item $f(x)$ 是是實例 $\Rightarrow$ $x$ 是是實例（等價地：$x$ 是否實例 $\Rightarrow$ $f(x)$ 是否實
例）。
\end{enumerate}

只證第一個方向，$f$ 可以把每個實例都變成同一個是實例，什麼都沒證到。 另外，$f$ 必須對
\textbf{所有}實例都有定義、都在多項式時間內算完，而且算的時候\textbf{不知道} $x$ 的答案。 「如果
$x$ 是是實例就輸出某某，否則輸出某某」這種寫法不是歸約，除非你能在多項式時間內判斷 $x$ 的答案。

\begin{lemma}
\label{lem:npc-reduce-p}
若 $A \le_{\mathrm{p}} B$ 且 $B \in \mathrm{P}$，則 $A \in \mathrm{P}$。
\end{lemma}

\begin{proof}
設 $f$ 在 $O(N^a)$ 時間內算完，解 $B$ 的演算法在 $O(N^b)$ 時間內算完。 長度 $N$ 的 $x$，算 $f(x)$
最多走 $O(N^a)$ 步，每步最多寫一個符號，所以 $|f(x)| = O(N^a)$。 照圖~\ref{fig:npc-reduce} 先算
$f(x)$ 再解 $B$，總時間 $O(N^a) + O((N^a)^b) = O(N^{ab})$，是多項式。 答案正確是因為 $x$ 和 $f(x)$
的答案相同。
\end{proof}

換一個方向讀：若 $A \le_{\mathrm{p}} B$ 且 $A$ 不在 P 裡，則 $B$ 也不在 P 裡。 所以 $A \le_{\mathrm{p}}
B$ 讀作「$B$ 至少和 $A$ 一樣難」。 記號 $\le$ 的方向和難度一致：左邊不比右邊難。

\begin{lemma}
\label{lem:npc-trans}
若 $A \le_{\mathrm{p}} B$ 且 $B \le_{\mathrm{p}} C$，則 $A \le_{\mathrm{p}} C$。
\end{lemma}

\begin{proof}
令 $f$、$g$ 是兩個歸約，取 $h(x) = g(f(x))$。 和引理~\ref{lem:npc-reduce-p} 的證明一樣，$|f(x)|$ 是
$|x|$ 的多項式，所以 $h$ 是多項式時間。 $x$ 是 $A$ 的是實例 $\iff$ $f(x)$ 是 $B$ 的是實例 $\iff$
$g(f(x))$ 是 $C$ 的是實例。
\end{proof}

\begin{examnote}
寫歸約時先寫出「$A \le_{\mathrm{p}} B$」，並說明哪一個是已知的難問題、哪一個是要證的問題：要證 $B$
很難，就把\textbf{已知很難的 $A$ 轉成 $B$}，方向是 $A \le_{\mathrm{p}} B$。 中文說「把 $A$ 歸約到
$B$」也是 $A \le_{\mathrm{p}} B$。 寫完構造之後，用兩個小標題分開寫兩個方向，最後一句交代 $f$ 的計
算時間。
\end{examnote}

\subsection{第一組例子：獨立集、團、頂點覆蓋}
\label{sec:npc-is-vc}

先看三個關係很近的圖論問題。 $G = (V, E)$ 是無向圖，$S \subseteq V$。

\begin{itemize}
\item $S$ 是\textbf{獨立集}（independent set）：$S$ 裡任兩個頂點都不相鄰。 \npcP{IND-SET}：給 $G$ 和
$k$，有沒有大小 $k$ 的獨立集？
\item $S$ 是\textbf{團}（clique）：$S$ 裡任兩個頂點都相鄰。 \npcP{CLIQUE}：給 $G$ 和 $k$，有沒有大
小 $k$ 的團？
\item $S$ 是\textbf{頂點覆蓋}（vertex cover）：每條邊至少有一端在 $S$ 裡。 \npcP{VERTEX-COVER}：給
$G$ 和 $k$，有沒有大小 $k$ 的頂點覆蓋？
\end{itemize}

（問「恰好 $k$ 個」或「至少 $k$ 個」對獨立集和團是一樣的，因為獨立集、團的子集合仍然是獨立集、團；
頂點覆蓋則是「恰好 $k$」和「至多 $k$」一樣，因為頂點覆蓋加頂點仍是頂點覆蓋。 以下都在 $k \le n$ 的
範圍討論。）

$G$ 的\textbf{補圖} $\overline{G}$ 和 $G$ 有相同的頂點，兩個不同的頂點在 $\overline{G}$ 裡相鄰，若且
唯若它們在 $G$ 裡不相鄰。

\begin{lemma}
\label{lem:npc-complement}
對 $G = (V, E)$ 和 $S \subseteq V$，下面三件事等價：
(1) $S$ 是 $G$ 的獨立集；(2) $V - S$ 是 $G$ 的頂點覆蓋；(3) $S$ 是 $\overline{G}$ 的團。
\end{lemma}

\begin{proof}
(1)$\iff$(3)：$S$ 裡兩個頂點在 $G$ 裡不相鄰，就是在 $\overline{G}$ 裡相鄰，照定義直接成立。

(1)$\Rightarrow$(2)：任取一條邊 $\{u, v\}$。 $S$ 是獨立集，$u$、$v$ 不會都在 $S$ 裡，所以至少一端
在 $V - S$ 裡。

(2)$\Rightarrow$(1)：任取 $u, v \in S$。 如果 $\{u, v\} \in E$，這條邊的兩端都不在 $V - S$ 裡，和
$V - S$ 是頂點覆蓋矛盾。 所以 $S$ 裡任兩點都不相鄰。
\end{proof}

圖~\ref{fig:npc-is-vc} 的 $G$ 有 $6$ 個頂點、$7$ 條邊。 $S = \{1, 4, 6\}$ 是獨立集（藍色），
$V - S = \{2, 3, 5\}$ 是頂點覆蓋（橘色）：七條邊每一條都碰到 $2$、$3$、$5$ 之一。 右邊的
$\overline{G}$ 有 $\binom{6}{2} - 7 = 8$ 條邊，$\{1, 4, 6\}$ 在裡面兩兩相鄰，是一個團。

\begin{center}
\begin{tikzpicture}[x=1cm, y=1cm]
\begin{scope}
\node[npc-vA] (a1) at (0, 0.8) {1};
\node[npc-vB] (a2) at (1.3, 1.6) {2};
\node[npc-vB] (a3) at (1.3, 0) {3};
\node[npc-vA] (a4) at (2.8, 1.6) {4};
\node[npc-vB] (a5) at (2.8, 0) {5};
\node[npc-vA] (a6) at (4.1, 0.8) {6};
\foreach \a/\b in {1/2, 1/3, 2/3, 2/4, 3/5, 4/5, 5/6}
    \draw[npc-e] (a\a) -- (a\b);
\node[npc-t] at (2.05, -0.8) {$G$};
\end{scope}
\begin{scope}[xshift=6.6cm]
\node[npc-vA] (b1) at (0, 0.6) {1};
\node[npc-vA] (b4) at (1.5, 1.7) {4};
\node[npc-vA] (b6) at (3.0, 0.6) {6};
\node[npc-v] (b3) at (3.2, 1.9) {3};
\node[npc-v] (b5) at (0.8, -0.4) {5};
\node[npc-v] (b2) at (2.2, -0.4) {2};
\draw[npc-ham] (b1) -- (b4) -- (b6) -- (b1);
\foreach \a/\b in {1/5, 5/2, 2/6, 3/4, 3/6}
    \draw[npc-e] (b\a) -- (b\b);
\node[npc-t] at (1.6, -1.15) {$\overline{G}$};
\end{scope}
\end{tikzpicture}
\diagcaption{左：$G$ 的獨立集 $\{1, 4, 6\}$（藍）和頂點覆蓋 $\{2, 3, 5\}$（橘）。 右：補圖
$\overline{G}$（頂點位置重新排過，比較好看出結構），粗線是團 $\{1, 4, 6\}$。}
\label{fig:npc-is-vc}
\end{center}

由引理~\ref{lem:npc-complement} 立刻得到兩個歸約：

\begin{itemize}
\item $\npcP{IND-SET} \le_{\mathrm{p}} \npcP{CLIQUE}$：$f(G, k) = (\overline{G}, k)$。 $G$ 有大小 $k$
的獨立集 $S$ $\Rightarrow$ $S$ 是 $\overline{G}$ 的大小 $k$ 的團；反過來 $\overline{G}$ 有大小 $k$ 的
團 $S$ $\Rightarrow$ $S$ 是 $G$ 的獨立集。 建 $\overline{G}$ 的鄰接矩陣要 $\Theta(n^2)$。
\item $\npcP{IND-SET} \le_{\mathrm{p}} \npcP{VERTEX-COVER}$：$f(G, k) = (G, n - k)$。 $G$ 有大小 $k$ 的
獨立集 $S$ $\Rightarrow$ $V - S$ 是大小 $n - k$ 的頂點覆蓋；反過來大小 $n - k$ 的頂點覆蓋 $C$
$\Rightarrow$ $V - C$ 是大小 $k$ 的獨立集。 只要把 $k$ 換成 $n - k$，$O(N)$ 時間。
\end{itemize}

同樣的論證也給出反方向的歸約（$\npcP{CLIQUE} \le_{\mathrm{p}} \npcP{IND-SET}$ 等），三個問題彼此可以
互相歸約，難度一樣。 到這裡我們只知道它們「一樣難」，還不知道「很難」；第~\ref{sec:npc-sat-is}~節
會把一個已知很難的問題歸約到獨立集。

注意「兩個方向」和「互相歸約」是兩件事。 證 $A \le_{\mathrm{p}} B$ 時要證的兩個方向，是同一個 $f$
下「$x$ 是是實例 $\Rightarrow$ $f(x)$ 是是實例」和它的反方向；並不需要另外證 $B \le_{\mathrm{p}} A$。

\section{NP 完備與 Cook--Levin 定理}
\label{sec:npc-npc}

\subsection{NP 難與 NP 完備}

\begin{definition}[NP 難、NP 完備]
\label{def:npc-npc}
決策問題 $B$ 是 \textbf{NP 難}（NP-hard）的，如果 NP 裡的\textbf{每一個}問題 $A$ 都滿足
$A \le_{\mathrm{p}} B$。 如果 $B$ 是 NP 難的，而且 $B \in \mathrm{NP}$，就說 $B$ 是 \textbf{NP 完
備}（NP-complete）的。 NP 完備問題所成的集合記作 NPC。
\end{definition}

NP 完備問題是 NP 裡「最難」的問題：NP 裡任何問題都能轉成它。

\begin{theorem}
\label{thm:npc-pnp}
如果有任何一個 NP 完備問題在 P 裡，則 $\mathrm{P} = \mathrm{NP}$。 等價地，如果
$\mathrm{P} \ne \mathrm{NP}$，則沒有任何 NP 完備問題在 P 裡。
\end{theorem}

\begin{proof}
設 $B$ 是 NP 完備的且 $B \in \mathrm{P}$。 任取 $A \in \mathrm{NP}$，由定義 $A \le_{\mathrm{p}} B$，
由引理~\ref{lem:npc-reduce-p}，$A \in \mathrm{P}$。 所以 $\mathrm{NP} \subseteq \mathrm{P}$，加上定
理~\ref{thm:npc-p-in-np} 得 $\mathrm{P} = \mathrm{NP}$。
\end{proof}

所以證明一個問題是 NP 完備的，等於說「除非 $\mathrm{P} = \mathrm{NP}$，它沒有多項式時間演算法」。 幾
十年來，成千上萬個 NP 完備問題都沒有人找到多項式時間演算法，這是大家相信 $\mathrm{P} \ne \mathrm{NP}$
的主要理由。

\begin{center}
\begin{tikzpicture}[x=1cm, y=1cm]
\begin{scope}
\clip (0, 0) ellipse (2.9 and 1.55);
\fill[memFill] (1.35, -2) rectangle (6, 2);
\end{scope}
\fill[coreFill] (-1.25, 0) ellipse (1.1 and 0.8);
\draw[thick, coreStroke] (-1.25, 0) ellipse (1.1 and 0.8);
\draw[thick] (0, 0) ellipse (2.9 and 1.55);
\draw[thick, memStroke, dashed, rounded corners=6pt] (1.35, -1.95) rectangle (5.6, 1.95);
\node[font=\small] at (-1.25, 0) {P};
\node[font=\small] at (-0.1, 1.15) {NP};
\node[font=\small, align=center] at (2.05, 0) {NP\\完備};
\node[font=\small, align=center] at (4.35, 0) {NP 難\\（不在 NP 裡\\的部分）};
\end{tikzpicture}
\diagcaption{假設 $\mathrm{P} \ne \mathrm{NP}$ 時的關係。 NP 完備是 NP 和 NP 難的交集；NP 難的問題不
一定在 NP 裡。}
\label{fig:npc-classes}
\end{center}

圖~\ref{fig:npc-classes} 是一般相信的樣子。 如果 $\mathrm{P} = \mathrm{NP}$，圖會整個塌掉：P、NP 和
NP 完備幾乎是同一個集合（只差「永遠回答是」和「永遠回答否」這兩個沒有否實例或沒有是實例的問題，
別的問題無法歸約到它們）。 $\mathrm{P} \ne \mathrm{NP}$ 時，Ladner 在 1975 年證明了 NP 裡一定還有既
不在 P、也不是 NP 完備的問題，質因數分解的決策版本常被當成可能的例子，但目前沒有證明。

\textbf{NP 難的問題不一定在 NP 裡}。 兩個常見的例子：
\begin{itemize}
\item 最佳化問題。 「求最短的旅行推銷員路線」不是決策問題，談不上在不在 NP 裡；但它會解的話決策版
本就會解，所以習慣上也說它是 NP 難的。 嚴格的說法是：它若有多項式時間演算法，則
$\mathrm{P} = \mathrm{NP}$。
\item 比 NP 更難的決策問題。 例如停機問題（給一支程式和輸入，問它會不會停）：任何公式 $\varphi$ 都
可以改寫成一支程式「依序試所有指派，找到滿足的就停，全部試完都沒有就進入無窮迴圈」，它會停若且唯若
$\varphi$ 可以滿足，所以 $\npcP{SAT} \le_{\mathrm{p}}$ 停機問題，停機問題是 NP 難的（下一小節會說明
\npcP{SAT} 是 NP 完備的）。 但停機問題根本沒有演算法能解，當然不在 NP 裡。
\end{itemize}

\subsection{怎麼證明一個問題是 NP 完備}

直接用定義要對 NP 裡「每一個」問題都給一個歸約，這做不到。 有了第一個 NP 完備問題之後，遞移性
（引理~\ref{lem:npc-trans}）讓證明只需要一個歸約：

\begin{lemma}
\label{lem:npc-recipe}
若 $A$ 是 NP 完備的，$A \le_{\mathrm{p}} B$，而且 $B \in \mathrm{NP}$，則 $B$ 是 NP 完備的。
\end{lemma}

\begin{proof}
任取 $C \in \mathrm{NP}$。 $A$ 是 NP 難的，所以 $C \le_{\mathrm{p}} A$；再由 $A \le_{\mathrm{p}} B$ 和
引理~\ref{lem:npc-trans}，$C \le_{\mathrm{p}} B$。 所以 $B$ 是 NP 難的，又 $B \in \mathrm{NP}$，$B$ 是
NP 完備的。
\end{proof}

\begin{examnote}
證明「$B$ 是 NP 完備的」照四步寫，每一步一個小標題：
(1)~$B \in \mathrm{NP}$：說出憑證是什麼、驗證者做什麼、為什麼是多項式時間。
(2)~挑一個已知 NP 完備的 $A$，寫出「證明 $A \le_{\mathrm{p}} B$」。
(3)~寫出構造 $f$，並說明 $f$ 是多項式時間（輸出大小也是多項式）。
(4)~證兩個方向：$x \in A \Rightarrow f(x) \in B$，以及 $f(x) \in B \Rightarrow x \in A$。
只證 (2)～(4) 得到的是「NP 難」；少了 (1) 就不能說「NP 完備」。
\end{examnote}

\subsection{布林公式與 Cook--Levin 定理}

\textbf{布林變數} $x_1, \ldots, x_n$ 的值是真（$1$）或假（$0$）。 一個變數 $x_i$ 或它的否定
$\neg x_i$ 叫一個\textbf{文字}（literal）；幾個文字用「或」（$\vee$）接起來叫一個\textbf{子句}（clause）；
幾個子句用「且」（$\wedge$）接起來，就是\textbf{合取範式}（conjunctive normal form，CNF）。 例如
\[
(x_1 \vee \neg x_2) \wedge (\neg x_1 \vee x_2 \vee x_3) \wedge (\neg x_3)
\]
是有三個子句的 CNF 公式。 給每個變數一個真假值叫一組\textbf{指派}（assignment）；一組指派讓每個子句
都至少有一個文字為真，就說它\textbf{滿足}這個公式。 上面的公式可以用 $x_1 = x_2 = 1$、$x_3 = 0$ 滿
足。

\begin{itemize}
\item \npcP{SAT}：給一個 CNF 公式，它可以被滿足嗎？
\item \npcP{3-SAT}：同樣的問題，但公式是 \textbf{3-CNF}：每個子句恰好有三個文字，而且這三個文字的變
數互不相同。
\end{itemize}

（有的課本的 3-SAT 允許子句少於三個文字，或允許同一個變數在一個子句裡出現兩次。 這些版本都一樣是
NP 完備的，題目有定義時照題目。）

\npcP{SAT} 在 NP 裡：憑證是一組指派，驗證者把它代進每個子句，時間和公式長度成正比。

\begin{theorem}[Cook--Levin]
\label{thm:npc-cook}
\npcP{SAT} 是 NP 完備的。
\end{theorem}

這是第一個被證明的 NP 完備問題（Cook 1971 年；Levin 獨立得到相同的結果）。 證明要對 NP 裡「每一個」
問題直接給歸約，本書不寫，只說想法。

\begin{whybox}[Cook--Levin 定理的想法]
任取 $A \in \mathrm{NP}$，它有一個驗證者 $V$，在 $p(N)$ 步內做完。 把 $V$ 在輸入 $(x, y)$ 上執行的
過程畫成一張表：第 $t$ 列是第 $t$ 步時整台機器的狀態（記憶體內容、目前執行到哪裡），共 $p(N)$ 列，
每列 $O(p(N))$ 格。 每一格的內容只由上一列附近固定幾格決定，這種「局部」的規則可以寫成常數大小的
子句。 再加上「第一列是 $x$ 和某個 $y$」「最後一列回答是」的子句，得到一個多項式大小的 CNF 公式
$\varphi_x$。 $y$ 的位元是公式裡的自由變數：有某個 $y$ 讓 $V$ 回答是，若且唯若 $\varphi_x$ 可以滿
足。 所以 $A \le_{\mathrm{p}} \npcP{SAT}$。 細節見 Sipser 第 7.4 節或 CLRS 第 34.3 節。
\end{whybox}

有了 \npcP{SAT} 之後，引理~\ref{lem:npc-recipe} 讓其他問題只要從一個已知的 NP 完備問題歸約過來。 下一
節就照這個方式，一步一步把 NP 完備的問題串起來。

\section{歸約鏈}
\label{sec:npc-chain}

圖~\ref{fig:npc-chain} 是本節要證明的歸約。 實線箭頭在本章證明，虛線箭頭的結論本書只敘述，不給證明。

\begin{center}
\begin{tikzpicture}[x=1cm, y=1cm]
\node[npc-pbox] (sat) at (0, 0.6) {SAT};
\node[npc-pbox] (s3) at (2.0, 0.6) {3-SAT};
\node[npc-pbox] (is) at (4.4, 2.2) {IND-SET};
\node[npc-pbox] (ss) at (4.4, 0.6) {SUBSET-SUM};
\node[npc-pbox] (hc) at (4.4, -0.9) {HAM-CYCLE};
\node[npc-pbox] (c3) at (4.4, -2.2) {3-COLOR};
\node[npc-pbox] (cl) at (7.6, 2.8) {CLIQUE};
\node[npc-pbox] (vc) at (7.6, 1.8) {VERTEX-COVER};
\node[npc-pbox] (kn) at (7.6, 0.6) {KNAPSACK};
\node[npc-pbox] (tsp) at (7.6, -0.4) {TSP};
\node[npc-pbox] (hp) at (7.6, -1.4) {HAM-PATH};
\node[npc-pbox] (lp) at (10.6, -1.4) {LONGEST-PATH};
\draw[npc-arr] (sat) -- (s3);
\draw[npc-arr] (s3) -- (is);
\draw[npc-arr] (s3) -- (ss);
\draw[npc-darr] (s3) -- (hc);
\draw[npc-darr] (s3) -- (c3);
\draw[npc-arr] (is) -- (cl);
\draw[npc-arr] (is) -- (vc);
\draw[npc-arr] (ss) -- (kn);
\draw[npc-arr] (hc) -- (tsp);
\draw[npc-darr] (hc) -- (hp);
\draw[npc-arr] (hp) -- (lp);
\node[font=\scriptsize, text=black!60, align=center] at (0, -0.2) {Cook--Levin};
\end{tikzpicture}
\diagcaption{本章的歸約鏈。 箭頭 $A \to B$ 表示 $A \le_{\mathrm{p}} B$；實線在本章證明，虛線只敘述結論。}
\label{fig:npc-chain}
\end{center}

\subsection{SAT 歸約到 3-SAT}
\label{sec:npc-sat-3sat}

\begin{theorem}
\label{thm:npc-3sat}
\npcP{3-SAT} 是 NP 完備的。
\end{theorem}

\npcP{3-SAT} 是 \npcP{SAT} 的特殊情形，驗證方法相同，所以在 NP 裡。 下面證明 $\npcP{SAT} \le_{\mathrm{p}}
\npcP{3-SAT}$：把每個子句換成一組恰好三個文字的子句，換的時候加入新的變數。 先做一個簡單的整理：一
個子句裡重複的文字只留一個；同時含有 $x$ 和 $\neg x$ 的子句一定為真，直接刪掉（全部刪光的話，輸出
一個固定可以滿足的 3-CNF 公式）。 然後對每個子句 $C = (\ell_1 \vee \cdots \vee \ell_r)$，依長度 $r$
換成：

\begin{itemize}
\item $r = 3$：不變。
\item $r = 2$：加一個新變數 $b$，換成 $(\ell_1 \vee \ell_2 \vee b) \wedge (\ell_1 \vee \ell_2 \vee \neg b)$。
\item $r = 1$：加兩個新變數 $c$、$d$，換成
$(\ell_1 \vee c \vee d) \wedge (\ell_1 \vee c \vee \neg d) \wedge (\ell_1 \vee \neg c \vee d) \wedge
(\ell_1 \vee \neg c \vee \neg d)$。
\item $r \ge 4$：加 $r - 3$ 個新變數 $a_1, \ldots, a_{r-3}$，換成 $r - 2$ 個子句串成的鏈
\[
(\ell_1 \vee \ell_2 \vee a_1) \wedge (\neg a_1 \vee \ell_3 \vee a_2) \wedge \cdots \wedge (\neg a_{r-4}
\vee \ell_{r-2} \vee a_{r-3}) \wedge (\neg a_{r-3} \vee \ell_{r-1} \vee \ell_r).
\]
第一個子句放 $\ell_1$、$\ell_2$，最後一個放 $\ell_{r-1}$、$\ell_r$，中間第 $s$ 個子句（$2 \le s \le r - 3$）
是 $(\neg a_{s-1} \vee \ell_{s+1} \vee a_s)$。 例如 $r = 4$ 時只有兩個子句 $(\ell_1 \vee \ell_2 \vee a_1)
\wedge (\neg a_1 \vee \ell_3 \vee \ell_4)$。
\end{itemize}

每個子句都用自己的新變數，不同子句的新變數不共用。 一個長度 $r$ 的子句最多換成 $\max(4, r - 2)$
個子句，所以新公式的長度是原公式的常數倍，$f$ 是線性時間。

\begin{proof}[證明（兩個方向）]
只要對每一種長度證明：\textbf{原來的子句 $C$ 被某組指派滿足，若且唯若這組指派能擴充到 $C$ 的新變
數上，使換出來的子句全部為真。} 每個子句的新變數只出現在它自己換出來的子句裡，可以各自擴充，所以
對整個公式也成立。

$r = 2$。 ($\Rightarrow$) $\ell_1 \vee \ell_2$ 為真，兩個新子句不管 $b$ 是什麼都為真。 ($\Leftarrow$)
兩個新子句都為真。 $b$ 和 $\neg b$ 有一個為假，含那個文字的子句只能靠 $\ell_1 \vee \ell_2$，所以
$\ell_1 \vee \ell_2$ 為真。

$r = 1$。 同樣的理由：$(c, d)$ 的四種組合中，恰好有一個子句的 $c$、$d$ 兩個文字都為假，它只能靠
$\ell_1$。

$r \ge 4$。 ($\Rightarrow$) 設 $\ell_i$ 為真。 令 $a_1, \ldots, a_{i-2}$ 為真，$a_{i-1}, \ldots,
a_{r-3}$ 為假（超出 $1..r-3$ 的範圍就不管）。 含 $\ell_i$ 的子句因 $\ell_i$ 為真；它左邊的子句各
含一個正的 $a_j$，$j \le i - 2$，為真；右邊的子句各含一個 $\neg a_j$，$j \ge i - 1$，也為真。
($\Leftarrow$) 用反證。 設 $\ell_1, \ldots, \ell_r$ 全部為假，而換出來的子句全部為真。 第一個子句
$(\ell_1 \vee \ell_2 \vee a_1)$ 只能靠 $a_1$，所以 $a_1 = 1$；第二個子句的 $\neg a_1$ 為假、$\ell_3$
為假，只能靠 $a_2$，所以 $a_2 = 1$；依此類推 $a_1 = \cdots = a_{r-3} = 1$。 最後一個子句
$(\neg a_{r-3} \vee \ell_{r-1} \vee \ell_r)$ 三個文字都為假，矛盾。
\end{proof}

\textbf{例子}。 公式
\[
\alpha = (x_1 \vee \neg x_2 \vee x_3 \vee x_4 \vee \neg x_5) \wedge (\neg x_1 \vee x_5) \wedge (x_2)
\]
換成 3-CNF：

\begin{center}
\begin{tabular}{p{0.3\linewidth}p{0.5\linewidth}l}
\toprule
原子句 & 換成 & 新變數 \\
\midrule
$(x_1 \vee \neg x_2 \vee x_3 \vee x_4 \vee \neg x_5)$ & $(x_1 \vee \neg x_2 \vee a_1) \wedge (\neg a_1 \vee
x_3 \vee a_2) \wedge (\neg a_2 \vee x_4 \vee \neg x_5)$ & $a_1, a_2$ \\
$(\neg x_1 \vee x_5)$ & $(\neg x_1 \vee x_5 \vee b) \wedge (\neg x_1 \vee x_5 \vee \neg b)$ & $b$ \\
$(x_2)$ & $(x_2 \vee c \vee d) \wedge (x_2 \vee c \vee \neg d) \wedge (x_2 \vee \neg c \vee d) \wedge
(x_2 \vee \neg c \vee \neg d)$ & $c, d$ \\
\bottomrule
\end{tabular}
\end{center}

共 $9$ 個子句、$10$ 個變數。 指派 $x = (0, 1, 0, 1, 0)$ 滿足 $\alpha$：第一個子句只靠 $x_4$（$\neg x_5$
也為真，但 $x_4$ 就夠了）。 照證明的做法，$\ell_4 = x_4$ 為真，令 $a_1 = a_2 = 1$：$(x_1 \vee \neg x_2
\vee a_1)$ 靠 $a_1$，$(\neg a_1 \vee x_3 \vee a_2)$ 靠 $a_2$，$(\neg a_2 \vee x_4 \vee \neg x_5)$ 靠
$x_4$。 $b$、$c$、$d$ 隨便給。 反過來，如果把 $x_4$ 改成 $0$、$x_5$ 改成 $1$（第二個子句仍為真），第
一個子句的五個文字全為假，鏈上 $a_1$、$a_2$ 被迫為真，最後一個子句 $(\neg a_2 \vee x_4 \vee \neg x_5)$
就不能滿足。

\subsection{3-SAT 歸約到獨立集}
\label{sec:npc-sat-is}

\begin{theorem}
\label{thm:npc-is}
\npcP{IND-SET}、\npcP{CLIQUE}、\npcP{VERTEX-COVER} 都是 NP 完備的。
\end{theorem}

三個問題都在 NP 裡：憑證是那 $k$ 個頂點，檢查兩兩不相鄰、兩兩相鄰或每條邊都被碰到，都是多項式時間。
第~\ref{sec:npc-is-vc}~節已經有 $\npcP{IND-SET} \le_{\mathrm{p}} \npcP{CLIQUE}$ 和 $\npcP{IND-SET}
\le_{\mathrm{p}} \npcP{VERTEX-COVER}$，所以只要證明 $\npcP{3-SAT} \le_{\mathrm{p}} \npcP{IND-SET}$。

想法：獨立集要在每個子句裡挑一個「負責讓它為真」的文字，而且挑出來的文字不能互相矛盾。 每個子句做
成一個三角形（同一個三角形最多挑一個），矛盾的兩個文字之間連一條邊（不能同時挑）。

\begin{pseudo}[label=alg:npc-sat-is]{3-SAT 歸約到獨立集}
\Function{Sat-To-IS}{$\varphi$}
    \State $V \gets \emptyset$，$E \gets \emptyset$
    \For{$j \gets 1$ \To $k$}
        \State 加入頂點 $v_{j,1}, v_{j,2}, v_{j,3}$，$v_{j,p}$ 標上文字 $\ell_{j,p}$
        \State 加入邊 $\{v_{j,1}, v_{j,2}\}$、$\{v_{j,1}, v_{j,3}\}$、$\{v_{j,2}, v_{j,3}\}$
    \EndFor
    \For{每一對頂點 $v_{j,p}$、$v_{j',p'}$，$j < j'$}
        \If{$\ell_{j,p}$ 和 $\ell_{j',p'}$ 是同一個變數的正、反兩面}
            \State 加入邊 $\{v_{j,p}, v_{j',p'}\}$ \Comment{矛盾邊}
        \EndIf
    \EndFor
    \State \Return $(G = (V, E),\ k)$
\EndFunction
\end{pseudo}

這裡 $\varphi = C_1 \wedge \cdots \wedge C_k$，$C_j = (\ell_{j,1} \vee \ell_{j,2} \vee \ell_{j,3})$。 圖
有 $3k$ 個頂點；第 6 行檢查 $O(k^2)$ 對頂點，所以 $f$ 是 $O(k^2)$ 時間。 輸出的門檻就是子句數 $k$。

\begin{proof}[證明：$\varphi$ 可以滿足 $\iff$ $G$ 有大小 $k$ 的獨立集]
($\Rightarrow$) 取一組滿足 $\varphi$ 的指派。 每個子句至少有一個文字為真，在每個子句挑一個為真的文
字，令 $S$ 是對應的 $k$ 個頂點。 $S$ 的頂點來自不同的子句，沒有三角形的邊；矛盾邊的兩端是 $x$ 和
$\neg x$，不可能同時為真，所以 $S$ 裡也沒有矛盾邊。 $S$ 是大小 $k$ 的獨立集。

($\Leftarrow$) 設 $S$ 是大小 $k$ 的獨立集。 同一個三角形的頂點兩兩相鄰，$S$ 在每個三角形最多一個頂
點；共有 $k$ 個三角形，所以\textbf{恰好}各一個。 把 $S$ 的頂點標的文字都設成真：$S$ 裡沒有矛盾邊，
不會同時要求 $x = 1$ 和 $x = 0$，這組要求是一致的；沒被提到的變數隨便給。 每個子句都有一個文字在
$S$ 裡而為真，所以 $\varphi$ 被滿足。
\end{proof}

證明的 ($\Leftarrow$) 用了兩個構造上的細節：三角形保證「每個子句恰好挑一個」，矛盾邊保證「挑出來的
文字不互相矛盾」。 少了任何一種邊，($\Leftarrow$) 就會失敗：沒有三角形的邊，獨立集可以在同一個子句挑
兩個、在別的子句一個都不挑；沒有矛盾邊，可以同時挑 $x_1$ 和 $\neg x_1$。

\textbf{例子}。 公式
\[
\varphi = (x_1 \vee \neg x_2 \vee x_5) \wedge (\neg x_1 \vee x_2 \vee \neg x_3) \wedge (x_1 \vee x_3 \vee
x_4) \wedge (\neg x_1 \vee \neg x_4 \vee \neg x_5)
\]
有 $5$ 個變數、$k = 4$ 個子句，得到的圖如圖~\ref{fig:npc-sat-is}：$12$ 個頂點、$12$ 條三角形的邊、
$8$ 條矛盾邊（$x_1$ 出現兩次、$\neg x_1$ 出現兩次，在外圈連成四條；$x_2$、$x_3$、$x_4$、$x_5$ 各一
條，在內圈）。

\begin{center}
\begin{tikzpicture}[x=0.95cm, y=0.95cm]
\node[npc-lsel] (s00) at (-0.6,4.6) {$x_1$};
\node[npc-l] (s01) at (0.8,4.1) {$\neg x_2$};
\node[npc-l] (s02) at (-0.1,3.2) {$x_5$};
\node[npc-l] (s10) at (6.6,4.6) {$\neg x_1$};
\node[npc-lsel] (s11) at (5.2,4.1) {$x_2$};
\node[npc-l] (s12) at (6.1,3.2) {$\neg x_3$};
\node[npc-lsel] (s20) at (6.6,-0.6) {$x_1$};
\node[npc-l] (s21) at (6.1,0.8) {$x_3$};
\node[npc-l] (s22) at (5.2,-0.1) {$x_4$};
\node[npc-l] (s30) at (-0.6,-0.6) {$\neg x_1$};
\node[npc-lsel] (s31) at (0.8,-0.1) {$\neg x_4$};
\node[npc-l] (s32) at (-0.1,0.8) {$\neg x_5$};
\draw[npc-e] (s00) -- (s01) -- (s02) -- (s00);
\draw[npc-e] (s10) -- (s11) -- (s12) -- (s10);
\draw[npc-e] (s20) -- (s21) -- (s22) -- (s20);
\draw[npc-e] (s30) -- (s31) -- (s32) -- (s30);
\draw[npc-x] (s00) to[bend left=22] (s10);
\draw[npc-x] (s00) to[bend right=22] (s30);
\draw[npc-x] (s01) -- (s11);
\draw[npc-x] (s02) -- (s32);
\draw[npc-x] (s10) to[bend left=22] (s20);
\draw[npc-x] (s12) -- (s21);
\draw[npc-x] (s20) to[bend left=22] (s30);
\draw[npc-x] (s22) -- (s31);
\node[npc-t] at (-1.5,4.6) {$C_1$};
\node[npc-t] at (7.5,4.6) {$C_2$};
\node[npc-t] at (7.5,-0.6) {$C_3$};
\node[npc-t] at (-1.5,-0.6) {$C_4$};
\end{tikzpicture}
\diagcaption{$\varphi$ 歸約出來的圖。 實線是三角形的邊，橘色虛線是矛盾邊；藍色的四個頂點是對應指
派 $x = (1, 1, 0, 0, 0)$ 的獨立集。}
\label{fig:npc-sat-is}
\end{center}

($\Rightarrow$) 方向：指派 $x = (1, 1, 0, 0, 0)$ 滿足 $\varphi$。 每個子句挑第一個為真的文字：$C_1$
挑 $x_1$，$C_2$ 的 $\neg x_1$ 為假、挑 $x_2$，$C_3$ 挑 $x_1$，$C_4$ 的 $\neg x_1$ 為假、挑 $\neg x_4$。
這四個頂點（圖中藍色）兩兩不相鄰：$C_1$ 和 $C_3$ 的兩個 $x_1$ 是同一個文字，之間沒有邊；$x_2$ 的矛
盾邊連到 $C_1$ 的 $\neg x_2$，$\neg x_4$ 的矛盾邊連到 $C_3$ 的 $x_4$，都沒被挑。

($\Leftarrow$) 方向：反過來從這個獨立集讀出 $x_1 = 1$、$x_2 = 1$、$x_4 = 0$；$x_3$、$x_5$ 沒被提到，
給 $0$。 得到的正是 $(1, 1, 0, 0, 0)$。 如果獨立集挑的是別的頂點，例如 $C_3$ 改挑 $x_3$（它的矛盾邊連
到 $C_2$ 的 $\neg x_3$，沒被挑），讀出來的指派變成 $(1, 1, 1, 0, 0)$，一樣滿足 $\varphi$。

\subsection{C 程式：兩邊各用暴力法算一次}

下面的程式讀入 3-CNF 公式，照虛擬碼~\ref{alg:npc-sat-is} 建出圖，再用兩個互不相干的暴力法分別回答
兩邊的問題：\texttt{brute\_sat} 試全部 $2^{\mathit{nv}}$ 組指派；\texttt{max\_is} 用第~20~章的回溯
求最大獨立集，完全不知道圖是怎麼造出來的。 兩邊的答案應該一致：公式可以滿足若且唯若最大獨立集的大小
是 $k$。 接著把證明的兩個方向實際做一次：從指派挑出獨立集，從獨立集讀回指派，再用兩個驗證者
\texttt{satisfies}、\texttt{independent} 檢查。

第一組輸入是上面的 $\varphi$；第二組是三個變數的八個子句，每種正負號組合各一個，任何指派都會讓其中
一個子句的三個文字全為假，不能滿足。 文字用整數表示：$+i$ 是 $x_i$，$-i$ 是 $\neg x_i$，印出來時寫成 \texttt{x2}、\texttt{!x2}。 獨立集
（\texttt{IS}）依頂點編號印出，每個子句恰好一個，依序是 $C_1, C_2, \ldots$ 挑的文字。

\begin{cbookrecap}[第~12~章位元運算、第~14~章結構、第~19~章 \texttt{malloc}]
\texttt{(m >> s) \& 1} 取出整數 \texttt{m} 從右數第 \texttt{s} 位（從 $0$ 數起）。 \texttt{brute\_sat}
讓 \texttt{m} 從 $0$ 跑到 $2^{\mathit{nv}} - 1$，第 $i$ 個變數取第 $\mathit{nv} - i$ 位，於是指派依
$x_1 x_2 \cdots x_{\mathit{nv}}$ 的字典序產生。 \texttt{1LL << nv} 用 \texttt{long long} 做位移，避免
\texttt{nv} 稍大時溢位。

結構 \texttt{Formula} 包一個用 \texttt{malloc} 要來的文字陣列，函式收結構的指標，用
\texttt{f->lit[i]} 存取。 鄰接矩陣用 \texttt{calloc} 要一塊 $n \times n$ 個 \texttt{char} 的一維陣列，
$u$ 列 $v$ 行在 \texttt{adj[u * n + v]}。
\end{cbookrecap}

\begin{example}[label=ex:npc-sat-is]{c}{(npc\_sat\_is.c) 3-SAT 歸約到獨立集，兩邊用暴力法比對}
#include <stdio.h>
#include <stdlib.h>

// 3-CNF 公式：nv 個變數 x1..xnv，k 個子句。 lit[3*j+p] 是
// 第 j 個子句的第 p 個文字：+i 表示 xi，-i 表示 !xi。
typedef struct {
    int nv, k;
    int *lit;
} Formula;

// 文字 l 在指派 x[1..nv] 之下是否為真
int lit_true(int l, const int *x)
{
    return l > 0 ? x[l] == 1 : x[-l] == 0;
}

// 驗證者一：指派 x 是否滿足每一個子句
int satisfies(const Formula *f, const int *x)
{
    for (int j = 0; j < f->k; j++) {
        int ok = 0;
        for (int p = 0; p < 3; p++) {
            ok = ok || lit_true(f->lit[3 * j + p], x);
        }
        if (!ok) {
            return 0;
        }
    }
    return 1;
}

// 歸約：頂點 3*j+p 對應第 j 個子句的第 p 個文字。 同一個子句
// 的三個頂點連成三角形，xi 和 !xi 之間連矛盾邊。 回傳鄰接矩陣。
char *reduce(const Formula *f, int *edges)
{
    int n = 3 * f->k;
    char *adj = calloc((size_t)n * n, 1);
    *edges = 0;
    for (int u = 0; u < n; u++) {
        for (int v = u + 1; v < n; v++) {
            if (u / 3 == v / 3 || f->lit[u] == -f->lit[v]) {
                adj[u * n + v] = adj[v * n + u] = 1;
                (*edges)++;
            }
        }
    }
    return adj;
}

// 驗證者二：頂點 s[0..cnt-1] 是否兩兩不相鄰
int independent(const char *adj, int n, const int *s, int cnt)
{
    for (int a = 0; a < cnt; a++) {
        for (int b = a + 1; b < cnt; b++) {
            if (adj[s[a] * n + s[b]]) {
                return 0;
            }
        }
    }
    return 1;
}

// 暴力法一：依 x1 x2 ... xnv 的字典序試全部 2^nv 組指派
int brute_sat(const Formula *f, int *x)
{
    for (long long m = 0; m < (1LL << f->nv); m++) {
        for (int i = 1; i <= f->nv; i++) {
            x[i] = (int)((m >> (f->nv - i)) & 1);
        }
        if (satisfies(f, x)) {
            return 1;
        }
    }
    return 0;
}

// 暴力法二：回溯求最大獨立集，不知道圖是怎麼造出來的。
// cur[0..cnt-1] 是目前選的頂點，接下來決定頂點 v。
void max_is(const char *adj, int n, int v, int *cur, int cnt,
            int *best, int *bcnt)
{
    if (cnt + (n - v) <= *bcnt) {     // 剩下的全選也贏不了
        return;
    }
    if (v == n) {
        *bcnt = cnt;
        for (int i = 0; i < cnt; i++) {
            best[i] = cur[i];
        }
        return;
    }
    int ok = 1;
    for (int i = 0; i < cnt; i++) {
        if (adj[cur[i] * n + v]) {
            ok = 0;
        }
    }
    if (ok) {                         // 先試「選 v」
        cur[cnt] = v;
        max_is(adj, n, v + 1, cur, cnt + 1, best, bcnt);
    }
    max_is(adj, n, v + 1, cur, cnt, best, bcnt);
}

// 印出頂點 v 標的文字（獨立集依頂點編號印，正好一個子句一個）
void print_lit(const Formula *f, int v)
{
    int l = f->lit[v];
    printf(" %sx%d", l < 0 ? "!" : "", abs(l));
}

void print_x(const char *label, const int *x, int nv)
{
    printf("%s", label);
    for (int i = 1; i <= nv; i++) {
        printf(" %d", x[i]);
    }
    printf("\n");
}

int main(void)
{
    int t;
    if (scanf("%d", &t) != 1) {
        return 1;
    }
    for (int c = 1; c <= t; c++) {
        Formula f;
        if (scanf("%d %d", &f.nv, &f.k) != 2) {
            return 1;
        }
        int n = 3 * f.k, edges;
        f.lit = malloc(n * sizeof *f.lit);
        for (int i = 0; i < n; i++) {
            if (scanf("%d", &f.lit[i]) != 1) {
                free(f.lit);
                return 1;
            }
        }
        char *adj = reduce(&f, &edges);
        int *x = calloc(f.nv + 1, sizeof *x);
        int *cur = malloc(n * sizeof *cur);
        int *best = malloc(n * sizeof *best);
        int sat = brute_sat(&f, x), bcnt = 0;
        max_is(adj, n, 0, cur, 0, best, &bcnt);
        printf("formula %d: %d vertices, %d edges\n", c, n, edges);
        printf("sat = %s, max IS = %d (k = %d)\n",
               sat ? "yes" : "no", bcnt, f.k);
        if (sat) {           // 方向一：指派 -> 每個子句挑一個真文字
            print_x("x =", x, f.nv);
            printf("-> IS:");
            for (int j = 0; j < f.k; j++) {
                int p = 0;
                while (!lit_true(f.lit[3 * j + p], x)) {
                    p++;
                }
                cur[j] = 3 * j + p;
                print_lit(&f, cur[j]);
            }
            printf("\n-> independent? %s\n",
                   independent(adj, n, cur, f.k) ? "yes" : "no");
        }
        if (bcnt == f.k) {   // 方向二：獨立集 -> 指派
            printf("IS:");
            for (int i = 1; i <= f.nv; i++) {
                x[i] = 0;    // 沒被提到的變數給 0
            }
            for (int i = 0; i < bcnt; i++) {
                x[abs(f.lit[best[i]])] = (f.lit[best[i]] > 0);
                print_lit(&f, best[i]);
            }
            printf("\n");
            print_x("-> x =", x, f.nv);
            printf("-> satisfies? %s\n",
                   satisfies(&f, x) ? "yes" : "no");
        }
        free(x);
        free(cur);
        free(best);
        free(adj);
        free(f.lit);
    }
    return 0;
}
\end{example}
\iosample{2 \\ 5 4 \\ 1 -2 5 \\ -1 2 -3 \\ 1 3 4 \\ -1 -4 -5 \\ 3 8 \\ 1 2 3 \\ 1 2 -3 \\ 1 -2 3 \\ 1 -2 -3 \\ -1 2 3 \\ -1 2 -3 \\ -1 -2 3 \\ -1 -2 -3}{formula 1: 12 vertices, 20 edges \\ sat = yes, max IS = 4 (k = 4) \\ x = 0 0 0 1 0 \\ -> IS: !x2 !x1 x4 !x1 \\ -> independent? yes \\ IS: x1 x2 x1 !x4 \\ -> x = 1 1 0 0 0 \\ -> satisfies? yes \\ formula 2: 24 vertices, 72 edges \\ sat = no, max IS = 7 (k = 8)}

第一組：$12$ 個頂點、$20$ 條邊（$12$ 條三角形的邊加 $8$ 條矛盾邊），和上面算的一樣。 \texttt{brute\_sat} 依字典序找到的第一組指派是
$(0, 0, 0, 1, 0)$，每個子句挑第一個為真的文字，得到另一個大小 $4$ 的獨立集。 \texttt{max\_is} 依頂點
編號先試「選」，找到的正是圖~\ref{fig:npc-sat-is} 的四個藍色頂點，讀回 $(1, 1, 0, 0, 0)$。 第二組：
$24$ 個頂點，$8$ 個三角形 $24$ 條邊，每個文字和另外四個子句裡它的否定各連一條，共 $24 \times 4 / 2
= 48$ 條矛盾邊；公式不能滿足，最大獨立集只有 $7 < 8$，兩邊都回答「否」。

兩個暴力法都是指數時間，只能跑小例子；它們的用途是檢查歸約做對了。 歸約本身（\texttt{reduce}）是
$O(k^2)$。

\subsection{3-SAT 歸約到子集合和}
\label{sec:npc-subset}

第~20~章的\textbf{子集合和}問題列出所有和為 $M$ 的子集合；這裡用它的決策版本。 \npcP{SUBSET-SUM}：給
正整數 $s_1, \ldots, s_m$ 和目標 $t$，能不能挑一些數（每個最多一次），和恰好是 $t$？

\begin{theorem}
\label{thm:npc-subset}
\npcP{SUBSET-SUM} 是 NP 完備的。
\end{theorem}

憑證是挑出來的那些數，加起來和 $t$ 比較，所以在 NP 裡。 下面證明 $\npcP{3-SAT} \le_{\mathrm{p}}
\npcP{SUBSET-SUM}$。 3-CNF 公式有 $n$ 個變數、$k$ 個子句。 造出來的每個數都寫成 $n + k$ 位的
\textbf{十進位}數，左邊 $n$ 位各對應一個變數，右邊 $k$ 位各對應一個子句：

\begin{itemize}
\item 每個變數 $x_i$ 造兩個數。 $y_i$ 代表「$x_i$ 為真」：變數位 $i$ 是 $1$，子句位 $j$ 在 $C_j$ 含
$x_i$ 時是 $1$，其他位是 $0$。 $z_i$ 代表「$x_i$ 為假」：變數位 $i$ 是 $1$，子句位 $j$ 在 $C_j$ 含
$\neg x_i$ 時是 $1$。
\item 每個子句 $C_j$ 造兩個一樣的\textbf{補數} $g_j = h_j$：只有子句位 $j$ 是 $1$。
\item 目標 $t$：每個變數位是 $1$，每個子句位是 $3$。
\end{itemize}

一共 $2n + 2k$ 個數，每個 $n + k$ 位，寫下它們要 $O((n + k)^2)$ 個數字，是公式長度的多項式。 注意
這些數的\textbf{數值}大約是 $10^{n+k}$，是公式長度的指數；第~\ref{sec:npc-pseudo}~節會用到這一點。

\begin{proof}[證明（兩個方向）]
先看一件事：\textbf{加任何一些數都不會進位}。 每一位上最多有幾個數是 $1$？ 變數位 $i$ 只有 $y_i$、
$z_i$ 兩個；子句位 $j$ 有三個文字對應的數（每個文字的變數不同，各貢獻一個）加上 $g_j$、$h_j$，最多
$5$。 都小於 $10$，所以挑出來的數相加時，每一位可以分開算。

($\Rightarrow$) 取一組滿足的指派。 $x_i$ 為真就挑 $y_i$，為假就挑 $z_i$，每個變數位剛好是 $1$。 子句
位 $j$ 等於 $C_j$ 裡為真的文字個數，介於 $1$ 和 $3$ 之間（至少 $1$ 因為指派滿足 $C_j$）；差幾個到
$3$，就從 $g_j$、$h_j$ 裡挑幾個（$0$、$1$ 或 $2$ 個）。 每一位都和 $t$ 相同，和恰好是 $t$。

($\Leftarrow$) 設挑出來的數和是 $t$。 不會進位，所以每一位分別等於 $t$ 的那一位。 變數位 $i$ 是 $1$，
表示 $y_i$、$z_i$ 恰好挑了一個：挑 $y_i$ 就令 $x_i = 1$，挑 $z_i$ 就令 $x_i = 0$，得到一組指派。 子句
位 $j$ 是 $3$，而 $g_j$、$h_j$ 最多貢獻 $2$，所以至少有一個挑出來的 $y_i$ 或 $z_i$ 在這一位是 $1$，
也就是 $C_j$ 至少有一個文字在這組指派下為真。 每個子句都被滿足。
\end{proof}

\textbf{例子}。 公式 $\psi = (x_1 \vee x_2 \vee \neg x_3) \wedge (x_1 \vee \neg x_2 \vee x_3)$，$n = 3$、
$k = 2$，每個數 $5$ 位：

\begin{center}
\begin{tabular}{c|ccc|cc}
\toprule
數 & $x_1$ & $x_2$ & $x_3$ & $C_1$ & $C_2$ \\
\midrule
$y_1$ & 1 & 0 & 0 & 1 & 1 \\
$z_1$ & 1 & 0 & 0 & 0 & 0 \\
$y_2$ & 0 & 1 & 0 & 1 & 0 \\
$z_2$ & 0 & 1 & 0 & 0 & 1 \\
$y_3$ & 0 & 0 & 1 & 0 & 1 \\
$z_3$ & 0 & 0 & 1 & 1 & 0 \\
$g_1 = h_1$ & 0 & 0 & 0 & 1 & 0 \\
$g_2 = h_2$ & 0 & 0 & 0 & 0 & 1 \\
\midrule
$t$ & 1 & 1 & 1 & 3 & 3 \\
\bottomrule
\end{tabular}
\end{center}

$z_1$ 的子句位全是 $0$：$\neg x_1$ 沒有出現在任何子句裡。 指派 $x = (1, 1, 0)$ 滿足 $\psi$：$C_1$ 的
三個文字 $x_1$、$x_2$、$\neg x_3$ 全為真，$C_2$ 只有 $x_1$ 為真。 挑 $y_1 + y_2 + z_3 = 10011 + 01010
+ 00110 = 11131$，$C_1$ 位已經是 $3$，$C_2$ 位只有 $1$，再挑 $g_2$、$h_2$ 補 $2$，得到 $11133 = t$。
另一方面，指派 $(0, 0, 1)$ 讓 $C_1$ 的三個文字全為假：挑 $z_1$、$z_2$、$y_3$ 之後 $C_1$ 位是 $0$，補數
最多再加 $2$，湊不到 $3$。

下面的程式照這個構造產生數字，再用和第~18~章 0/1 背包同一種的填表（只記「湊不湊得出」）判斷能不能
湊出 $t$，並從表中還原挑了哪些數、讀回指派。 動態規劃的表有 $t + 1$ 格，程式在
超過一千萬格時只印 $t$，不做動態規劃。 第一組輸入是 $\psi$，第二組是上一小節的 $\varphi$。

\begin{example}[label=ex:npc-subset]{c}{(npc\_subset\_sum.c) 3-SAT 歸約到子集合和，再用動態規劃判斷}
#include <stdio.h>
#include <stdlib.h>

#define LIMIT 10000000LL   // DP 表超過一千萬格就不做

// 第 i 個數的名字：y1 z1 y2 z2 ... g1 h1 g2 h2 ...
void print_name(int i, int nv)
{
    if (i < 2 * nv) {
        printf("%c%d", i % 2 ? 'z' : 'y', i / 2 + 1);
    } else {
        printf("%c%d", i % 2 ? 'h' : 'g', (i - 2 * nv) / 2 + 1);
    }
}

// 3-SAT 歸約到子集合和，再用 DP 判斷。 每個數有 w = nv + k 位
// （十進位）：左邊 nv 位對應變數，右邊 k 位對應子句。
void solve(int nv, int k, const int *lit)
{
    int w = nv + k, cnt = 2 * w;
    long long pos[18], p = 1, t = 0;      // pos[d]：由左數第 d 位
    for (int d = w - 1; d >= 0; d--, p *= 10) {
        pos[d] = p;
    }
    // a[2i-2] = y_i，a[2i-1] = z_i；之後是 g1 h1 g2 h2 ...
    long long *a = calloc(cnt, sizeof *a);
    for (int i = 1; i <= nv; i++) {
        a[2 * i - 2] = a[2 * i - 1] = pos[i - 1];
        t += pos[i - 1];
    }
    for (int j = 0; j < k; j++) {
        for (int q = 0; q < 3; q++) {   // xi 加到 y_i，!xi 加到 z_i
            int l = lit[3 * j + q];
            a[l > 0 ? 2 * l - 2 : -2 * l - 1] += pos[nv + j];
        }
        a[2 * nv + 2 * j] = a[2 * nv + 2 * j + 1] = pos[nv + j];
        t += 3 * pos[nv + j];
    }
    printf("t = %0*lld\n", w, t);
    if (t + 1 > LIMIT) {                  // 表太大：只印 t 就結束
        printf("DP: %lld cells, skipped\n", t + 1);
        free(a);
        return;
    }
    for (int i = 0; i < cnt; i += 2) {
        print_name(i, nv);
        printf(" %0*lld | ", w, a[i]);
        print_name(i + 1, nv);
        printf(" %0*lld\n", w, a[i + 1]);
    }

    // reach[s]：湊得出 s；from[s]：第一次湊出 s 用的數（編號+1）
    char *reach = calloc(t + 1, 1), *from = calloc(t + 1, 1);
    reach[0] = 1;
    for (int i = 0; i < cnt; i++) {
        for (long long s = t; s >= a[i]; s--) {   // 每個數只用一次
            if (reach[s - a[i]] && !reach[s]) {
                reach[s] = 1;
                from[s] = (char)(i + 1);
            }
        }
    }
    printf("DP: %lld cells, %s\n", t + 1, reach[t] ? "yes" : "no");
    if (reach[t]) {
        int *x = calloc(nv + 1, sizeof *x);
        printf("subset:");
        for (long long s = t; s > 0; s -= a[from[s] - 1]) {
            int i = from[s] - 1;
            printf(" ");
            print_name(i, nv);
            if (i < 2 * nv && i % 2 == 0) {   // 挑了 y_i：xi 為真
                x[i / 2 + 1] = 1;
            }
        }
        printf("\n-> x =");
        for (int i = 1; i <= nv; i++) {
            printf(" %d", x[i]);
        }
        printf("\n");
        free(x);
    }
    free(reach);
    free(from);
    free(a);
}

int main(void)
{
    int tc;
    if (scanf("%d", &tc) != 1) {
        return 1;
    }
    for (int c = 1; c <= tc; c++) {
        int nv, k;
        if (scanf("%d %d", &nv, &k) != 2 || nv + k > 18) {
            return 1;
        }
        int *lit = malloc(3 * k * sizeof *lit);
        for (int i = 0; i < 3 * k; i++) {
            if (scanf("%d", &lit[i]) != 1) {
                free(lit);
                return 1;
            }
        }
        printf("formula %d: %d numbers, %d digits\n",
               c, 2 * (nv + k), nv + k);
        solve(nv, k, lit);
        free(lit);
    }
    return 0;
}
\end{example}
\iosample{2 \\ 3 2 \\ 1 2 -3 \\ 1 -2 3 \\ 5 4 \\ 1 -2 5 \\ -1 2 -3 \\ 1 3 4 \\ -1 -4 -5}{formula 1: 10 numbers, 5 digits \\ t = 11133 \\ y1 10011 | z1 10000 \\ y2 01010 | z2 01001 \\ y3 00101 | z3 00110 \\ g1 00010 | h1 00010 \\ g2 00001 | h2 00001 \\ DP: 11134 cells, yes \\ subset: h1 g1 y3 z2 y1 \\ -> x = 1 0 1 \\ formula 2: 18 numbers, 9 digits \\ t = 111113333 \\ DP: 111113334 cells, skipped}

$\psi$ 的表只有 $11134$ 格，動態規劃找到 $h_1 + g_1 + y_3 + z_2 + y_1$，讀回指派 $(1, 0, 1)$：$C_1$ 只
有 $x_1$ 為真，所以補了 $g_1$、$h_1$；$C_2$ 的三個文字全為真，不用補。 $\varphi$ 只多了兩個變數、兩個子句，數字從 $5$ 位變成 $9$ 位，表
就要一億一千多萬格。 動態規劃的時間和 $t$ 的\textbf{數值}成正比，而歸約造出來的 $t$ 大約是
$10^{n+k}$，每多一個變數或子句，表就大十倍。

\textbf{0/1 背包}。 第~18~章 0/1 背包的決策版本 \npcP{KNAPSACK}：給 $w_i$、$v_i$、$W$、$V$，能不能挑一些物
品，總重不超過 $W$、總價值至少 $V$？ 它在 NP 裡（憑證是挑的物品）。 證明 $\npcP{SUBSET-SUM}
\le_{\mathrm{p}} \npcP{KNAPSACK}$：令每個物品的重量和價值都等於 $s_i$，$W = V = t$。

\begin{itemize}
\item ($\Rightarrow$) 挑出的數和恰好是 $t$，同一批物品總重 $t \le W$、總價值 $t \ge V$。
\item ($\Leftarrow$) 挑出的物品總重不超過 $t$、總價值至少 $t$；兩者是同一個和，所以恰好是 $t$。
\end{itemize}

所以 \npcP{KNAPSACK} 是 NP 完備的，這就是第~18~章和第~19~章說「0/1 背包目前沒有已知的多項式時間演算
法」的原因。 在子集合和的數之外再加兩個適當的數，也可以證明\textbf{分割問題}（\npcP{PARTITION}：能不能把一堆正整數分成和相等的
兩堆）是 NP 完備的，本書不展開。

\subsection{哈密頓環歸約到旅行推銷員問題}
\label{sec:npc-tsp}

\begin{theorem}[Karp 1972]
\label{thm:npc-ham}
\npcP{HAM-CYCLE}（無向圖有沒有哈密頓環）是 NP 完備的。
\end{theorem}

它在 NP 裡（虛擬碼~\ref{alg:npc-verify-ham}）。 NP 難的證明要從 \npcP{3-SAT} 或 \npcP{VERTEX-COVER}
歸約過來，需要設計比較複雜的「零件」（gadget）子圖來模擬變數和子句，本書只敘述結論，有興趣可以看
CLRS 第 34.5.3 節（從頂點覆蓋歸約）或 Sipser 第 7.5 節（從 3-SAT 先歸約到有向圖的哈密頓路徑，再轉成無向圖）。

旅行推銷員問題的決策版本 \npcP{TSP}：給 $n$ 個城市，兩兩之間的成本 $c(i, j) = c(j, i)$ 是非負整數，以
及門檻 $B$，有沒有一條經過每個城市恰好一次、回到起點、總成本不超過 $B$ 的路線（巡迴）？

\begin{theorem}
\label{thm:npc-tsp}
\npcP{TSP} 是 NP 完備的。
\end{theorem}

\begin{proof}
\textbf{在 NP 裡}：憑證是城市的排列，檢查它是排列、算總成本和 $B$ 比，$\Theta(n)$。

\textbf{歸約 $\npcP{HAM-CYCLE} \le_{\mathrm{p}} \npcP{TSP}$}。 給無向圖 $G = (V, E)$，$n = |V|$，造 $n$
個城市（一個頂點一個），成本
\[
c(i, j) = \begin{cases} 1, & \{i, j\} \in E, \\ 2, & \{i, j\} \notin E, \end{cases}
\]
門檻 $B = n$。 填 $n^2$ 格成本，$\Theta(n^2)$ 時間。

($\Rightarrow$) $G$ 有哈密頓環。 照環的順序走，就是一條經過每個城市一次的巡迴，它的 $n$ 段都是 $G$ 的
邊，成本 $n \le B$。

($\Leftarrow$) 有一條成本不超過 $n$ 的巡迴。 巡迴恰好有 $n$ 段，每段成本至少 $1$，總和不超過 $n$ 只
可能每段都是 $1$，也就是每段都是 $G$ 的邊。 這條巡迴照順序就是 $G$ 的一個哈密頓環。
\end{proof}

成本只用了 $1$ 和 $2$，而 $1 + 1 \ge 2$，所以造出來的實例滿足三角不等式 $c(i, j) \le c(i, k) + c(k,
j)$。 也就是說，即使限制成本滿足三角不等式（第~29~章的「度量 TSP」），\npcP{TSP} 仍是 NP 完備的。 許
多課本改用 $0$ 和 $1$（邊的成本 $0$、非邊 $1$、門檻 $0$），證明完全一樣，但那樣造出來的實例不滿足三
角不等式。 第~20~章的 TSP 允許 $c(i, j) \ne c(j, i)$，對稱的實例是它的特殊情形，所以那個版本的決策
問題也是 NP 完備的，求最佳路線是 NP 難的。

\textbf{例子}。 圖~\ref{fig:npc-ham} 左邊的 $H_1$ 有哈密頓環 $1 \to 2 \to 5 \to 3 \to 4 \to 1$（粗線）。
歸約出的成本矩陣如下，這條巡迴的成本是 $1 + 1 + 1 + 1 + 1 = 5 = B$，回答「是」。

\begin{center}
\begin{tikzpicture}[x=1cm, y=1cm]
\begin{scope}
\node[npc-v] (h1) at (90:1.35) {1};
\node[npc-v] (h2) at (18:1.35) {2};
\node[npc-v] (h5) at (-54:1.35) {5};
\node[npc-v] (h3) at (-126:1.35) {3};
\node[npc-v] (h4) at (162:1.35) {4};
\draw[npc-ham] (h1) -- (h2) -- (h5) -- (h3) -- (h4) -- (h1);
\draw[npc-e] (h2) -- (h3);
\draw[npc-e] (h4) -- (h5);
\node[npc-t] at (0, -1.9) {$H_1$};
\end{scope}
\begin{scope}[xshift=5.2cm]
\node[npc-v] (k1) at (-1.6, 0.8) {1};
\node[npc-v] (k2) at (-1.6, -0.8) {2};
\node[npc-v] (k3) at (0, 0) {3};
\node[npc-v] (k4) at (1.6, 0.8) {4};
\node[npc-v] (k5) at (1.6, -0.8) {5};
\foreach \a/\b in {1/2, 1/3, 2/3, 3/4, 3/5, 4/5}
    \draw[npc-e] (k\a) -- (k\b);
\node[npc-t] at (0, -1.9) {$H_2$};
\end{scope}
\end{tikzpicture}
\diagcaption{$H_1$ 有哈密頓環（粗線）；$H_2$ 沒有，因為頂點 $3$ 是割點。}
\label{fig:npc-ham}
\end{center}

\begin{center}
\begin{tabular}{c|ccccc}
\toprule
$c(i, j)$ & 1 & 2 & 3 & 4 & 5 \\
\midrule
1 & -- & 1 & 2 & 1 & 2 \\
2 & 1 & -- & 1 & 2 & 1 \\
3 & 2 & 1 & -- & 1 & 1 \\
4 & 1 & 2 & 1 & -- & 1 \\
5 & 2 & 1 & 1 & 1 & -- \\
\bottomrule
\end{tabular}
\end{center}

右邊的 $H_2$ 是兩個共用頂點 $3$ 的三角形。 哈密頓環刪掉一個頂點後還是連通的（剩下一條路徑），但
$H_2$ 刪掉 $3$ 就斷成兩塊（$3$ 是第~23~章的割點），所以 $H_2$ 沒有哈密頓環。 對應的 \npcP{TSP} 實例
每條巡迴都至少要用一段成本 $2$，例如 $1 \to 2 \to 3 \to 4 \to 5 \to 1$ 的成本是 $1 + 1 + 1 + 1 + 2 =
6 > 5$，回答「否」，和 $H_2$ 一致。

\textbf{哈密頓路徑與最長路徑}。 \npcP{HAM-PATH}（有沒有一條經過每個頂點恰好一次的路徑）也是 NP 完備
的（同樣只敘述）。 第~18~章說過最長簡單路徑沒有最佳子結構；它的決策版本 \npcP{LONGEST-PATH}（給 $G$
和 $L$，有沒有一條至少 $L$ 條邊的簡單路徑）是 NP 完備的：憑證是那條路徑；而 $\npcP{HAM-PATH}
\le_{\mathrm{p}} \npcP{LONGEST-PATH}$ 只要取 $L = n - 1$，因為恰好 $n - 1$ 條邊的簡單路徑經過全部 $n$
個頂點，正是哈密頓路徑。 第~25~章的最短路徑在 P 裡，換成「最長」就變成 NP 完備；在 DAG 上的最長路
徑則因為沒有環，可以用第~22~章的拓樸順序加動態規劃在線性時間內解決。

\subsection{其他常見的 NP 完備問題}

下面這些問題也是 NP 完備的，本書只敘述結論。 右欄是外表很像、卻在 P 裡的問題，研究所考試很常拿來對
照。

\begin{center}
\begin{tabular}{p{0.44\linewidth}p{0.46\linewidth}}
\toprule
NP 完備 & 相近但在 P 裡 \\
\midrule
\npcP{3-SAT}：每個子句三個文字 & 2-SAT：每個子句兩個文字（第~23~章，強連通元件） \\
\npcP{3-COLOR}：三種顏色著色（第~20~章的 $m$ 著色，$m \ge 3$）& 兩種顏色著色，就是判斷二分圖（第~21~章）\\
\npcP{HAM-CYCLE}：每個\textbf{頂點}恰好經過一次 & 歐拉環：每條\textbf{邊}恰好經過一次（連通且每個頂點
度數都是偶數） \\
\npcP{LONGEST-PATH}：最長簡單路徑 & 最短路徑（第~25~章）；DAG 上的最長路徑（第~22~章） \\
\npcP{IND-SET}：一般圖的最大獨立集 & 樹上的最大獨立集（樹形動態規劃）\\
\npcP{SUBSET-SUM}、\npcP{KNAPSACK} & 分數背包（第~19~章，貪心）\\
最大割：把頂點分兩邊，跨兩邊的邊最多 & 最小割（第~26~章，最大流）\\
\npcP{VERTEX-COVER}：一般圖的最小頂點覆蓋 & 二分圖的最小頂點覆蓋（大小等於最大匹配，可用第~26~章的最大流求）\\
最小回饋邊集：最少刪幾條邊讓有向圖沒有環 & 判斷有向圖有沒有環（第~21~章的 DFS）\\
\bottomrule
\end{tabular}
\end{center}

這張表要記住的是：問題的外表只差一點，難度可以從多項式跳到 NP 完備。 判斷一個新問題時，不要只看它
「像」哪一個問題。

\section{擬多項式時間與強 NP 完備}
\label{sec:npc-pseudo}

\subsection{\texorpdfstring{$\Theta(mt)$}{Θ(mt)} 的動態規劃並沒有證明 P = NP}

\npcP{SUBSET-SUM} 有一個 $\Theta(mt)$ 的動態規劃（$m$ 個數、目標 $t$，範例~\ref{ex:npc-subset} 用的就是
它），\npcP{KNAPSACK} 有第~18~章 $\Theta(nW)$ 的動態規劃。 它們都是 NP 完備問題，看起來卻有多項式
時間的演算法，難道 $\mathrm{P} = \mathrm{NP}$？

答案在第~\ref{sec:ana-input-size}~節：輸入大小是位元數。 $t$ 只占 $b = \lceil \lg(t + 1) \rceil$ 個
位元，$\Theta(mt) = \Theta(m \cdot 2^b)$ 是輸入長度的指數。 範例~\ref{ex:npc-subset} 的第二組輸入就是
實例：輸入只是 $18$ 個九位數，表卻有一億多格。 這種演算法叫擬多項式時間。

\begin{definition}[擬多項式時間]
\label{def:npc-pseudo}
一個處理整數的演算法是\textbf{擬多項式時間}（pseudo-polynomial time）的，如果它的時間是輸入長度 $N$
和輸入中最大的數值 $M$ 這兩者的多項式。
\end{definition}

$M$ 最大可以到 $2^N$ 左右，所以擬多項式時間不一定是多項式時間。 但如果實例裡的數都不大（例如 $M \le
N^2$），擬多項式時間的演算法在這些實例上就是多項式時間。 第~\ref{sec:npc-search}~節「從 $B = 0$ 一
個一個問」的做法也是同一回事：問的次數是 $C$ 的數值的多項式，是擬多項式。

用一進位寫數字（$t$ 寫成 $t$ 個 $1$）時，輸入長度至少是 $t$，$\Theta(mt)$ 就成了輸入長度的多項式。
所以「一進位編碼的子集合和」在 P 裡。 這不矛盾：一進位把輸入灌大了指數倍，演算法相對於這個膨脹過
的長度當然顯得快。 這也是第~\ref{sec:npc-decision}~節說一進位不算合理編碼的原因。

\subsection{強 NP 完備與弱 NP 完備}

子集合和的 NP 完備性完全靠「數很大」：第~\ref{sec:npc-subset}~節的歸約造出約 $10^{n+k}$ 的數，如果數
值被限制在輸入長度的多項式以內，動態規劃就能在多項式時間內解掉。 有些問題不是這樣，數再小也一樣難。

\begin{definition}[強 NP 完備]
\label{def:npc-strong}
一個 NP 完備問題是\textbf{強 NP 完備}（strongly NP-complete）的，如果存在多項式 $q$，使得它限制在
「所有數都不超過 $q(N)$」的實例上仍是 NP 完備的。 有擬多項式時間演算法的 NP 完備問題叫\textbf{弱 NP
完備}（weakly NP-complete）。
\end{definition}

\begin{theorem}
\label{thm:npc-strong}
如果 $\mathrm{P} \ne \mathrm{NP}$，強 NP 完備的問題沒有擬多項式時間演算法。
\end{theorem}

\begin{proof}
設強 NP 完備的問題 $A$ 有時間 $O((NM)^c)$ 的演算法。 把它用在「所有數都不超過 $q(N)$」的實例上，時
間是 $O((N q(N))^c)$，是多項式。 這個受限的版本是 NP 完備的，由定理~\ref{thm:npc-pnp}，
$\mathrm{P} = \mathrm{NP}$。
\end{proof}

\begin{itemize}
\item \npcP{TSP} 是強 NP 完備的：定理~\ref{thm:npc-tsp} 的歸約只用了成本 $1$ 和 $2$，限制成本不超過
$2$ 的 \npcP{TSP} 仍是 NP 完備。 所以不要期待 \npcP{TSP} 有「依成本大小」的動態規劃。
\item 輸入裡根本沒有數值的問題（\npcP{SAT}、\npcP{CLIQUE}、\npcP{HAM-CYCLE}）照定義都是強 NP 完備的。
\item \npcP{SUBSET-SUM}、\npcP{KNAPSACK}、\npcP{PARTITION} 都是弱 NP 完備的：有擬多項式時間的動態規劃。
\item 裝箱問題（bin packing）和 3-分割問題（把 $3q$ 個數分成 $q$ 組、每組三個數、和都相同）是強 NP
完備的，雖然它們看起來和分割問題很像（只敘述）。
\end{itemize}

\begin{examnote}
題目問某個演算法「是不是多項式時間」時，先寫出輸入長度：$n$ 個不超過 $M$ 的整數，長度約
$n \lg M$ 位元。 再把時間寫成 $n$ 和 $\lg M$ 的函數。 $\Theta(nM)$、$\Theta(n\sqrt{M})$ 都含有 $M$ 本身，
是 $\lg M$ 的指數，屬於擬多項式；$\Theta(n^2 \lg M)$ 才是多項式。
\end{examnote}

\section{遇到 NP 完備問題怎麼辦}
\label{sec:npc-cope}

證明一個問題是 NP 完備的，意思是「不要再找對所有輸入都又快又正確的演算法」，但問題還是要解。 常見
的方向有幾種，本書前後幾章都已經或將會用到：

\begin{itemize}
\item \textbf{看輸入有沒有特別的結構}。 圖是樹、是 DAG、是二分圖，或子句只有兩個文字時，很多 NP 完備
問題會變成多項式時間（上一節的表）。 NP 完備說的是\textbf{最壞情況}，不是每個實例都難。
\item \textbf{數不大時用擬多項式時間}。 背包、子集合和在容量、目標不大時，第~18~章的動態規劃非常實用。
\item \textbf{聰明一點的指數時間}。 第~20~章的回溯和分支限界靠剪枝讓實際試的組合少很多；\npcP{TSP} 用
「子集合 $\times$ 終點」當狀態的動態規劃是 $O(n^2 2^n)$，比 $(n - 1)!$ 條路線少得多。
\item \textbf{參數小時}。 頂點覆蓋問 $k$ 個頂點夠不夠：任取一條還沒被覆蓋的邊 $\{u, v\}$，$u$、$v$ 至少
要選一個，分兩支遞迴，每支 $k$ 減 $1$，搜尋樹最多 $2^k$ 個葉節點，時間 $O(2^k (n + m))$。 $k$ 很小時，
即使 $n$ 很大也跑得動。
\item \textbf{近似}。 不求最佳解，只求保證離最佳解不遠的解，例如頂點覆蓋的 2-近似，見第~29~章。
\end{itemize}

\section{本章常見錯誤}
\label{sec:npc-pitfalls}

\begin{center}
\begin{tabular}{p{0.24\linewidth}p{0.50\linewidth}p{0.16\linewidth}}
\toprule
寫歸約時的錯誤 & 例子 & 見 \\
\midrule
歸約方向反了 & 要證新問題 $B$ 是 NP 難的，卻證了 $B \le_{\mathrm{p}} \npcP{3-SAT}$；這只說明 $B$ 不比
\npcP{3-SAT} 難 & 第~\ref{sec:npc-reduce}~節 \\
只證一個方向 & 只證「$\varphi$ 可以滿足 $\Rightarrow$ 圖有大小 $k$ 的獨立集」；把每個實例都送到同一個
是實例的 $f$ 也滿足這一半 & 第~\ref{sec:npc-reduce}~節 \\
構造時偷看答案 & 「$\varphi$ 可以滿足就輸出一個三角形，否則輸出一條邊」：$f$ 要先判斷 $\varphi$，這本身
就不知道能不能在多項式時間做到 & 第~\ref{sec:npc-reduce}~節 \\
構造的大小不是多項式 & 把 $\varphi$ 的 $2^n$ 組指派各做成一個頂點；$f$ 的時間（和輸出長度）必須是多項式
& 第~\ref{sec:npc-reduce}~節 \\
\bottomrule
\end{tabular}
\end{center}

\begin{center}
\begin{tabular}{p{0.24\linewidth}p{0.50\linewidth}p{0.16\linewidth}}
\toprule
觀念上的錯誤 & 例子 & 見 \\
\midrule
把 NP 讀成「非多項式」 & 「NP 的問題都沒有多項式時間演算法」；其實 $\mathrm{P} \subseteq
\mathrm{NP}$ & 第~\ref{sec:npc-np}~節 \\
把猜想當成定理 & 「已經證明 \npcP{TSP} 沒有多項式時間演算法」；要說「除非 $\mathrm{P} = \mathrm{NP}$」
& 第~\ref{sec:npc-np}~節 \\
以為 NP 難的都在 NP 裡 & 說「求最短巡迴是 NP 完備的」；最佳化問題不是決策問題，只能說 NP 難 &
第~\ref{sec:npc-npc}~節 \\
忘了證在 NP 裡 & 只給了歸約就說「NP 完備」，得到的其實只有 NP 難 & 第~\ref{sec:npc-npc}~節 \\
特殊情形和一般化搞反 & 「2-SAT 是 3-SAT 的特殊情形，所以 2-SAT 也是 NP 完備」；特殊情形可能比較簡單
& 第~\ref{sec:npc-chain}~節 \\
用數值當輸入大小 & 「子集合和有 $\Theta(mt)$ 的動態規劃，所以 $\mathrm{P} = \mathrm{NP}$」 &
第~\ref{sec:npc-pseudo}~節 \\
以為每個實例都難 & 「獨立集是 NP 完備的，所以樹上的最大獨立集也沒有多項式時間演算法」 &
第~\ref{sec:npc-cope}~節 \\
\bottomrule
\end{tabular}
\end{center}

\section{觀念題}
\label{sec:npc-exercises}

題型和研究所考試、期中期末考的筆試相近，有選擇、計算、追蹤和證明。 除非題目另外說明，都可以使用本章
證明或敘述過的 NP 完備問題（\npcP{SAT}、\npcP{3-SAT}、\npcP{IND-SET}、\npcP{CLIQUE}、
\npcP{VERTEX-COVER}、\npcP{SUBSET-SUM}、\npcP{KNAPSACK}、\npcP{HAM-CYCLE}、\npcP{HAM-PATH}、\npcP{TSP}、
\npcP{3-COLOR}）。 解答在本章最後。

\begin{enumerate}
\item （選擇，可能不只一個答案）下列哪些敘述正確？
\begin{enumerate}
\item $\mathrm{P} \subseteq \mathrm{NP}$。
\item NP 是所有「不能在多項式時間內解決」的決策問題所成的集合。
\item 只要有一個 NP 完備問題有多項式時間演算法，就有 $\mathrm{P} = \mathrm{NP}$。
\item 每一個 NP 難的問題都在 NP 裡。
\item 若決策問題 $A$ 在 NP 裡，則 $A$ 的補問題（是、否對調）也在 NP 裡。
\item 若 $\mathrm{P} \ne \mathrm{NP}$，則 NP 完備問題的每一個實例都要指數時間才能解。
\end{enumerate}

\item （選擇，可能不只一個答案）已知 $A \le_{\mathrm{p}} B$。 下列哪些\textbf{一定}成立？
\begin{enumerate}
\item 若 $B \in \mathrm{P}$，則 $A \in \mathrm{P}$。
\item 若 $A \in \mathrm{P}$，則 $B \in \mathrm{P}$。
\item 若 $A$ 是 NP 完備的且 $B \in \mathrm{NP}$，則 $B$ 是 NP 完備的。
\item 若 $B$ 是 NP 完備的，則 $A$ 是 NP 完備的。
\item 若 $A \notin \mathrm{P}$，則 $B \notin \mathrm{P}$。
\item $B \le_{\mathrm{p}} A$。
\end{enumerate}

\item （是非，說明理由）「子集合和問題有 $\Theta(mt)$ 的動態規劃，$m$ 是數的個數、$t$ 是目標；而子集合
和是 NP 完備的，所以 $\mathrm{P} = \mathrm{NP}$。」

\item （分類）輸入是 $n$ 個正整數，每個都不超過 $M$。 下列執行時間各是多項式時間、擬多項式時間（但不
是多項式時間），還是兩者都不是？ 說明理由。
\begin{gather*}
\text{(a) } \Theta(n^3 \lg M) \qquad \text{(b) } \Theta(nM) \qquad \text{(c) } \Theta(n\sqrt{M}) \\
\text{(d) } \Theta(2^n \lg M) \qquad \text{(e) } \Theta(n^2 (\lg M)^3) \qquad \text{(f) } \Theta(n^{\lg M})
\end{gather*}

\item （追蹤）用第~\ref{sec:npc-sat-3sat}~節的方法把
\[
\beta = (x_1 \vee \neg x_2 \vee x_3 \vee \neg x_4 \vee x_5 \vee x_6) \wedge (\neg x_3 \vee x_4) \wedge
(\neg x_1) \wedge (x_2 \vee \neg x_5 \vee \neg x_6)
\]
換成 3-CNF。
\begin{enumerate}
\item 寫出換出來的每一個子句，新公式有幾個子句、幾個變數？
\item 指派 $x = (0, 1, 1, 1, 0, 0)$ 滿足 $\beta$。 照證明的做法替新變數決定值，並逐一檢查第一個子句換出
來的子句都為真。
\item 一般地，長度 $r \ge 4$ 的子句換出幾個子句、幾個新變數？ 為什麼整個歸約是線性時間？
\end{enumerate}

\item （追蹤）對公式 $\rho = (x_1 \vee x_2 \vee \neg x_3) \wedge (\neg x_1 \vee x_3 \vee x_4) \wedge
(\neg x_2 \vee \neg x_4 \vee x_3)$ 做第~\ref{sec:npc-sat-is}~節的歸約。
\begin{enumerate}
\item 圖有幾個頂點、幾條三角形的邊、幾條矛盾邊？ 列出每一條矛盾邊。
\item 找一個大小 $3$ 的獨立集，並讀出對應的指派，驗證它滿足 $\rho$。
\item 找一個包含 $C_1$ 的 $\neg x_3$ 的大小 $3$ 的獨立集，或說明不存在。
\end{enumerate}

\item （計算）用第~\ref{sec:npc-subset}~節的公式 $\psi$ 和它的數字表。
\begin{enumerate}
\item 指派 $(1, 0, 0)$ 滿足 $\psi$。 寫出它對應的子集合（含補數），並驗算和是 $t$。
\item 一共有幾個子集合的和恰好是 $t$？（$g_j$、$h_j$ 數值相同，但算成兩個不同的數。）
\item 如果把每個數改用\textbf{二進位}解讀（同樣的 $0$、$1$ 位，但每一位代表 $2$ 的次方；$t$ 的子句位
仍是 $3$），證明裡哪一句話不再成立？ 找一個和是 $t$、卻不對應任何指派的子集合。
\end{enumerate}

\item （證明）$k$-\npcP{COLOR}：給圖 $G$，能不能用 $k$ 種顏色替頂點著色，使每條邊兩端不同色？ 已知
\npcP{3-COLOR} 是 NP 完備的。 證明 \npcP{4-COLOR} 是 NP 完備的，並說明怎麼推廣到每個 $k \ge 3$。
（提示：期末考如果多了一個時段，可以多排一門「和每門課都衝突」的課。）

\item （證明）設 $G$ 至少有兩個頂點。 令 $G'$ 是在 $G$ 上加一個新頂點 $w$，並把 $w$ 和 $G$ 的每個頂點
相連。
\begin{enumerate}
\item 證明：$G$ 有哈密頓路徑若且唯若 $G'$ 有哈密頓環。
\item 這給出哪一個方向的歸約？ 它能不能用來證明 \npcP{HAM-PATH} 是 NP 難的？ 如果不能，它證明了什麼？
\end{enumerate}

\item （證明）社團要選 $k$ 位幹部，名單上有些同學兩兩合不來，不能同時當選；另外有一群同學 $R$ 是指定
一定要選的。 \npcP{TEAM}：給合不來關係的圖 $G$、整數 $k$ 和集合 $R$，有沒有 $k$ 位兩兩合得來、而且
包含 $R$ 全部的人選？
\begin{enumerate}
\item 證明 \npcP{TEAM} 是 NP 完備的。
\item 反過來，說明怎麼用一次解 \npcP{IND-SET} 的程序回答 \npcP{TEAM}，所以兩者一樣難。
\end{enumerate}

\item （是非，說明理由）
\begin{enumerate}
\item 「2-SAT 是 \npcP{3-SAT} 的特殊情形，而 \npcP{3-SAT} 是 NP 完備的，所以 2-SAT 也是 NP 完備的。」
\item 「\npcP{SAT} 包含 \npcP{3-SAT} 當特殊情形，\npcP{3-SAT} 是 NP 難的，所以 \npcP{SAT} 也是 NP 難的。」
\item 「要證明新問題 $B$ 是 NP 難的，可以寫一個把 $B$ 的實例轉成 3-CNF 公式的多項式時間歸約。」
\item 「\npcP{3-SAT} 是 \npcP{SAT} 的特殊情形，所以只要說明這一點，就證明了 \npcP{3-SAT} 是 NP 難的。」
\end{enumerate}

\item （證明）$C$ 表示 \npcP{TSP} 實例裡最大的成本。
\begin{enumerate}
\item 證明：如果 \npcP{TSP} 有時間 $O((nC)^5)$ 的演算法，則 $\mathrm{P} = \mathrm{NP}$。
\item 0/1 背包有 $\Theta(nW)$ 的演算法，形式和 (a) 的假設很像，為什麼不能得到同樣的結論？
\end{enumerate}

\item （計算）完全二分圖 $K_{2,3}$：頂點 $1, 2$ 在一邊、$3, 4, 5$ 在另一邊，兩邊之間每一對都有邊，同一
邊內沒有邊。
\begin{enumerate}
\item 寫出定理~\ref{thm:npc-tsp} 的歸約得到的成本矩陣和門檻 $B$。
\item 求最便宜的巡迴成本，並證明不可能更小。
\item 由 (b) 判斷 $K_{2,3}$ 有沒有哈密頓環。
\end{enumerate}
\end{enumerate}

\section{實作練習}
\label{sec:npc-practice}

Aki-Judge 目前沒有和本章對應的題目；這一章的練習以觀念題為主。 範例~\ref{ex:npc-sat-is} 和
範例~\ref{ex:npc-subset} 可以拿來檢查追蹤題：把觀念題 6 的 $\rho$ 改成輸入，程式會印出圖的頂點數、邊
數和一組大小 $3$ 的獨立集。 第二支程式的第一組輸入就是觀念題 7 的 $\psi$，可以把它改成只輸出「和是
$t$ 的子集合有幾個」，對照觀念題 7(b) 的答案。

\section{讀完你應該會}
\label{sec:npc-checklist}

\begin{itemize}
\item 能把最佳化問題寫成決策版本，並用決策版本的演算法（自我歸約或二分搜尋）求出最佳解，判斷呼叫次數
是不是輸入長度的多項式。
\item 能替一個問題說出憑證和驗證者，判斷它在不在 NP 裡；能說明 $\mathrm{P} \subseteq \mathrm{NP}$，不把
NP 當成「非多項式」。
\item 能寫出多項式時間歸約的定義，判斷題目裡歸約的方向，並推出「若 $B \in \mathrm{P}$ 則 $A \in
\mathrm{P}$」這類結論。
\item 能照四步證明一個問題是 NP 完備的；對 \npcP{SAT} $\to$ \npcP{3-SAT}、\npcP{3-SAT} $\to$
\npcP{IND-SET}、\npcP{IND-SET} $\to$ \npcP{CLIQUE}／\npcP{VERTEX-COVER}、\npcP{3-SAT} $\to$
\npcP{SUBSET-SUM}、\npcP{HAM-CYCLE} $\to$ \npcP{TSP}，能寫出構造、證明兩個方向，並在小例子上實際做一次。
\item 能分辨多項式、擬多項式和指數時間，說明 0/1 背包的 $\Theta(nW)$ 為什麼沒有證明 $\mathrm{P} =
\mathrm{NP}$，以及強、弱 NP 完備的差別。
\end{itemize}

\textbf{延伸閱讀}（書末會整理完整的參考資料）

\begin{itemize}
\item Thomas H. Cormen 等，\textit{Introduction to Algorithms}，第 4 版，第 34 章：用電路可滿足性證明
Cook--Levin 定理，並完整證明哈密頓環、子集合和等問題的 NP 完備性。
\item Michael Sipser，\textit{Introduction to the Theory of Computation}，第 3 版，第 7 章：從圖靈機的角度
定義 P 和 NP，第 7.4、7.5 節是 Cook--Levin 定理和一串歸約，零件的畫法很清楚。
\item Michael R. Garey、David S. Johnson，\textit{Computers and Intractability: A Guide to the Theory of
NP-Completeness}，1979：附錄列了三百多個 NP 完備問題，第 4 章講強 NP 完備；遇到新問題時可以先查這本。
\item Jon Kleinberg、Éva Tardos，\textit{Algorithm Design}，第 8 章：把 NP 完備問題分成幾大類（打包、覆
蓋、排列、分割、數值），對「從哪一個問題歸約過來」很有幫助。
\item Steven S. Skiena，\textit{The Algorithm Design Manual}，第 3 版，第 11 章：從實務的角度談怎麼證
明、怎麼面對 NP 完備問題。
\item Stephen A. Cook，``The Complexity of Theorem-Proving Procedures''，STOC 1971；Richard M. Karp，
``Reducibility Among Combinatorial Problems''，1972：第一個 NP 完備問題，以及最早的二十一個 NP 完備問題。
\end{itemize}

\addcontentsline{toc}{section}{觀念題解答}
\section*{觀念題解答}

\begin{enumerate}
\item (a)、(c) 正確。 (b)：P 裡的問題都在 NP 裡，NP 是 nondeterministic polynomial 的縮寫。 (d)：停機問題
是 NP 難的，但沒有演算法能解，不在 NP 裡；求最短巡迴這種最佳化問題也不在 NP 裡。 (e)：這是
$\mathrm{NP} = \text{co-NP}$ 的問題，目前不知道；例如「公式不能滿足」沒有已知的短憑證。 (f)：NP 完備說
的是最壞情況。 很多實例（例如沒有矛盾的公式、樹狀的圖）很容易；而且 $\mathrm{P} \ne \mathrm{NP}$ 只保證
沒有多項式時間演算法，不代表一定要指數時間（例如 $n^{\lg n}$ 介於兩者之間）。

\item (a)、(c)、(e) 一定成立。 (a) 是引理~\ref{lem:npc-reduce-p}，(e) 是 (a) 的換質位。 (c) 是引
理~\ref{lem:npc-recipe}。 (b)：$A$ 簡單不代表 $B$ 簡單；P 裡的任何問題都能歸約到 \npcP{3-SAT}（直接在
多項式時間內解出答案，再輸出一個固定的是實例或否實例），但不能因此說 \npcP{3-SAT} 在 P 裡。 (d)：同一
個例子，$A$ 在 P 裡、$B = \npcP{3-SAT}$；除非 $\mathrm{P} = \mathrm{NP}$，$A$ 不是 NP 完備的。 (f)：歸
約不必反過來成立。

\item 錯。 輸入大小是寫下 $m$ 個數和 $t$ 所需的位元數，$t$ 只占約 $\lg t$ 個位元，$\Theta(mt) = \Theta(m 2^{\lg t})$ 是位元數的指數，屬於擬多項式時間，不是多
項式時間。 第~\ref{sec:npc-subset}~節的歸約造出來的 $t$ 大約是 $10^{n+k}$，這時動態規劃要指數時間。

\item 輸入長度約 $n \lg M$。
(a) 多項式。 (b) 擬多項式：$M = 2^{\lg M}$ 是 $\lg M$ 的指數。 (c) 擬多項式：$\sqrt{M} = 2^{(\lg M)/2}$，
仍是 $\lg M$ 的指數，但它是 $n$ 和 $M$ 的多項式。 (d) 兩者都不是：$2^n$ 連 $n$ 的多項式都不是。 (e) 多項
式。 (f) 兩者都不是：$n^{\lg M} = M^{\lg n}$，$M$ 的次方隨 $n$ 變大，不是 $n$ 和 $M$ 的多項式（固定
$M = 2$ 時是 $n$，但 $M = n$ 時是 $n^{\lg n}$）。

\item
\begin{enumerate}
\item 第一個子句 $r = 6$，用 $a_1, a_2, a_3$：
$(x_1 \vee \neg x_2 \vee a_1) \wedge (\neg a_1 \vee x_3 \vee a_2) \wedge (\neg a_2 \vee \neg x_4 \vee a_3)
\wedge (\neg a_3 \vee x_5 \vee x_6)$。 第二個子句 $r = 2$，用 $b$：$(\neg x_3 \vee x_4 \vee b) \wedge
(\neg x_3 \vee x_4 \vee \neg b)$。 第三個子句 $r = 1$，用 $c$、$d$：$(\neg x_1 \vee c \vee d) \wedge
(\neg x_1 \vee c \vee \neg d) \wedge (\neg x_1 \vee \neg c \vee d) \wedge (\neg x_1 \vee \neg c \vee
\neg d)$。 第四個子句不變。 共 $4 + 2 + 4 + 1 = 11$ 個子句，$6 + 3 + 1 + 2 = 12$ 個變數。
\item 檢查 $\beta$：第一個子句靠 $x_3$，第二個靠 $x_4$，第三個靠 $\neg x_1$，第四個靠 $x_2$。 第一個子句
裡為真的是 $\ell_3 = x_3$，令 $a_1 = 1$、$a_2 = a_3 = 0$：$(x_1 \vee \neg x_2 \vee a_1)$ 靠 $a_1$；
$(\neg a_1 \vee x_3 \vee a_2)$ 靠 $x_3$；$(\neg a_2 \vee \neg x_4 \vee a_3)$ 靠 $\neg a_2$；
$(\neg a_3 \vee x_5 \vee x_6)$ 靠 $\neg a_3$。 $b$、$c$、$d$ 任意。
\item $r - 2$ 個子句、$r - 3$ 個新變數。 每個子句換出的文字個數最多是原來長度的常數倍（$r = 1$ 時 $12$
個文字、$r \ge 4$ 時 $3(r - 2) < 3r$），所以新公式長度是原公式的常數倍，掃一遍就能寫出來。
\end{enumerate}

\item
\begin{enumerate}
\item $9$ 個頂點、$9$ 條三角形的邊、$5$ 條矛盾邊：$x_1$($C_1$)--$\neg x_1$($C_2$)、$x_2$($C_1$)--$\neg x_2$
($C_3$)、$\neg x_3$($C_1$)--$x_3$($C_2$)、$\neg x_3$($C_1$)--$x_3$($C_3$)、$x_4$($C_2$)--$\neg x_4$($C_3$)。
\item 例如 $\{x_2(C_1), x_3(C_2), x_3(C_3)\}$：兩個 $x_3$ 是同一個文字，沒有邊；$x_2$ 的矛盾邊連到
$\neg x_2(C_3)$，沒被選。 讀出 $x_2 = x_3 = 1$，$x_1$、$x_4$ 任意（例如都給 $0$）。 驗證：$C_1$ 靠
$x_2$，$C_2$、$C_3$ 靠 $x_3$。
\item 選了 $\neg x_3(C_1)$，兩個 $x_3$ 都不能選。 $C_2$ 只能選 $\neg x_1$ 或 $x_4$，$C_3$ 只能選
$\neg x_2$ 或 $\neg x_4$，而 $x_4$、$\neg x_4$ 不能同時選。 例如 $\{\neg x_3(C_1), \neg x_1(C_2),
\neg x_2(C_3)\}$，對應 $x_1 = x_2 = x_3 = 0$；或 $\{\neg x_3(C_1), x_4(C_2), \neg x_2(C_3)\}$，對應
$x_2 = x_3 = 0$、$x_4 = 1$。 都滿足 $\rho$。
\end{enumerate}

\item
\begin{enumerate}
\item $x_1 = 1$ 挑 $y_1 = 10011$，$x_2 = 0$ 挑 $z_2 = 01001$，$x_3 = 0$ 挑 $z_3 = 00110$，和是
$11122$：$C_1$ 有 $x_1$、$\neg x_3$ 兩個文字為真，$C_2$ 有 $x_1$、$\neg x_2$ 兩個。 兩個子句各差 $1$，
挑 $g_1 = 00010$ 和 $g_2 = 00001$（或換成 $h_1$、$h_2$），總和 $11133 = t$。
\item 每組滿足的指派決定了 $y$、$z$ 怎麼挑；子句 $j$ 有 $c_j$ 個文字為真時，要從 $g_j$、$h_j$ 挑
$3 - c_j$ 個，有 $\binom{2}{3 - c_j}$ 種（$c_j = 2$ 時 $2$ 種，其他 $1$ 種）。 $8$ 組指派中 $(0, 0, 1)$
不滿足 $C_1$、$(0, 1, 0)$ 不滿足 $C_2$，其餘六組的 $(c_1, c_2)$ 是：$(0, 0, 0)$：$(1, 1)$；
$(0, 1, 1)$：$(1, 1)$；$(1, 0, 0)$：$(2, 2)$；$(1, 0, 1)$：$(1, 3)$；$(1, 1, 0)$：$(3, 1)$；
$(1, 1, 1)$：$(2, 2)$。 子集合個數 $1 + 1 + 4 + 1 + 1 + 4 = 12$。
\item 「加任何一些數都不會進位」不再成立：二進位每一位只要有兩個 $1$ 就進位。 例如挑 $y_1$、$z_1$、
$g_1$：二進位下 $y_1 = 10011_2 = 19$、$z_1 = 10000_2 = 16$、$g_1 = 00010_2 = 2$，和是 $37$；而
$t = 16 + 8 + 4 + 3 \cdot 2 + 3 \cdot 1 = 37$。 這個子集合同時挑了 $y_1$ 和 $z_1$，又沒有挑任何 $x_2$、
$x_3$ 的數，不對應任何指派：$x_1$ 位的 $16 + 16$ 多出來的 $16$，剛好補上缺少的 $x_2$、$x_3$ 兩位
（$8 + 4$）和子句位少掉的 $4$。 用十進位（或任何大於 $5$ 的進位制）就不會發生。
\end{enumerate}

\item \textbf{在 NP 裡}：憑證是每個頂點的顏色，檢查每條邊兩端不同色，$O(n + m)$。 \textbf{歸約
$\npcP{3-COLOR} \le_{\mathrm{p}} \npcP{4-COLOR}$}：$G'$ 是 $G$ 加一個新頂點 $w$，$w$ 和 $G$ 的每個頂點相
連，$O(n + m)$ 時間。 ($\Rightarrow$) $G$ 有三著色，$w$ 塗第四種顏色，$w$ 的鄰居都不是這個顏色，得到
$G'$ 的四著色。 ($\Leftarrow$) $G'$ 有四著色。 $w$ 和所有其他頂點相鄰，它們都不是 $w$ 的顏色，只用了另外
三種，限制到 $G$ 上就是三著色。 推廣：同樣的構造給出 $k$-\npcP{COLOR} $\le_{\mathrm{p}}$
$(k+1)$-\npcP{COLOR}，由 $k = 3$ 開始歸納，每個 $k \ge 3$ 都是 NP 完備的。 （$k = 2$ 在 P 裡，就是判斷
二分圖，所以不能從 $k = 2$ 開始。）

\item
\begin{enumerate}
\item ($\Rightarrow$) $G$ 有哈密頓路徑 $v_1, v_2, \ldots, v_n$。 $w$ 和 $v_1$、$v_n$ 都相鄰，
$w, v_1, \ldots, v_n, w$ 是 $G'$ 的哈密頓環（$n \ge 2$，所以至少三個頂點，是一個環）。 ($\Leftarrow$)
$G'$ 有哈密頓環，它經過 $w$ 恰好一次。 從環上刪掉 $w$，剩下一條依序經過 $G$ 的全部頂點的路徑；這條路徑
上的邊都不碰 $w$，所以都是 $G$ 原本的邊。 它是 $G$ 的哈密頓路徑。
\item 這是 $\npcP{HAM-PATH} \le_{\mathrm{p}} \npcP{HAM-CYCLE}$（$G \mapsto G'$，$O(n + m)$）。 它說明
\npcP{HAM-CYCLE} 至少和 \npcP{HAM-PATH} 一樣難，不能用來證明 \npcP{HAM-PATH} 是 NP 難的：那需要把一個已
知的 NP 完備問題歸約\textbf{到} \npcP{HAM-PATH}。 它真正給出的是：若 \npcP{HAM-CYCLE} 在 P 裡，則
\npcP{HAM-PATH} 也在 P 裡。 （\npcP{HAM-CYCLE} 已知是 NP 完備的，所以這個結論本來就由「\npcP{HAM-PATH}
在 NP 裡」得到了，這個歸約沒有提供新的難度資訊。）
\end{enumerate}

\item
\begin{enumerate}
\item \textbf{在 NP 裡}：憑證是 $k$ 位人選，檢查人數、兩兩之間沒有邊、$R$ 的每個人都在裡面，多項式時
間。 \textbf{歸約 $\npcP{IND-SET} \le_{\mathrm{p}} \npcP{TEAM}$}：$f(G, k) = (G, k, \emptyset)$。
($\Rightarrow$) 大小 $k$ 的獨立集就是包含空集合的人選。 ($\Leftarrow$) 包含 $\emptyset$ 的人選就是大小
$k$ 的獨立集。 所以 \npcP{TEAM} 是 NP 完備的。 一般化（多加一個條件可以設成「沒有限制」）的問題至少
和原問題一樣難。
\item 先檢查 $R$ 本身兩兩合得來、$|R| \le k$，不是就回答否。 再刪掉 $R$ 和所有和 $R$ 裡某人合不來的人，
在剩下的圖 $G''$ 上問有沒有大小 $k - |R|$ 的獨立集。 剩下的人和 $R$ 都合得來，所以 $G''$ 的獨立集加上
$R$ 就是合法的人選；反過來，合法人選扣掉 $R$ 一定落在 $G''$ 裡。 這是
$\npcP{TEAM} \le_{\mathrm{p}} \npcP{IND-SET}$，兩者一樣難。
\end{enumerate}

\item
\begin{enumerate}
\item 錯。 一個 NP 完備問題的特殊情形可能很簡單；2-SAT 在 P 裡（第~23~章）。 證明 NP 難要把已知的難問題
歸約到 2-SAT，而這做不到（除非 $\mathrm{P} = \mathrm{NP}$）。
\item 對。 取 $f$ 為恆等函數，3-CNF 公式本身就是 CNF 公式，答案相同，這是 $\npcP{3-SAT} \le_{\mathrm{p}}
\npcP{SAT}$。 一般化的問題至少和特殊情形一樣難。
\item 錯。 那是 $B \le_{\mathrm{p}} \npcP{3-SAT}$，只說明 $B$ 不比 \npcP{3-SAT} 難（由 $\npcP{3-SAT} \in
\mathrm{NP}$，這個歸約可以推出 $B \in \mathrm{NP}$，但推不出 $B$ 是 NP 難的）。 要證 $B$ 是 NP 難的，應該把 3-CNF 公式轉成 $B$ 的實例。
\item 錯。 理由同 (a)：特殊情形可能比較簡單，必須另外證 $\npcP{SAT} \le_{\mathrm{p}} \npcP{3-SAT}$
（第~\ref{sec:npc-sat-3sat}~節）。
\end{enumerate}

\item
\begin{enumerate}
\item 給 \npcP{HAM-CYCLE} 的實例 $G$，用定理~\ref{thm:npc-tsp} 的歸約造出 \npcP{TSP} 實例，它的成本只有 $1$
和 $2$，$C = 2$。 在它上面跑假設的演算法，時間 $O((2n)^5) = O(n^5)$，加上歸約的 $\Theta(n^2)$，得到
\npcP{HAM-CYCLE} 的多項式時間演算法。 \npcP{HAM-CYCLE} 是 NP 完備的，由定理~\ref{thm:npc-pnp}，
$\mathrm{P} = \mathrm{NP}$。 （這正是「\npcP{TSP} 是強 NP 完備」的意思。）
\item (a) 的論證需要一群數值都很小、卻仍然 NP 完備的實例（\npcP{TSP} 有：成本只有 $1$、$2$ 的實
例）。 0/1 背包是弱 NP 完備的，它的難度來自很大的數：從子集合和歸約過來時，$W = t$ 可以是輸入長度的指
數，$\Theta(nW)$ 在這些實例上是指數時間。 而在 $W$ 不超過輸入長度的多項式的實例上，$\Theta(nW)$ 已經
是多項式時間，這樣的實例不可能是 NP 完備的（除非 $\mathrm{P} = \mathrm{NP}$），所以推不出結論。
\end{enumerate}

\item
\begin{enumerate}
\item 兩邊之間的成本是 $1$，同一邊之內是 $2$，$B = 5$：
\begin{center}
\begin{tabular}{c|ccccc}
$c(i, j)$ & 1 & 2 & 3 & 4 & 5 \\
\hline
1 & -- & 2 & 1 & 1 & 1 \\
2 & 2 & -- & 1 & 1 & 1 \\
3 & 1 & 1 & -- & 2 & 2 \\
4 & 1 & 1 & 2 & -- & 2 \\
5 & 1 & 1 & 2 & 2 & -- \\
\end{tabular}
\end{center}
\item 最小成本是 $6$，例如 $1 \to 3 \to 2 \to 4 \to 5 \to 1$：$1 + 1 + 1 + 2 + 1$。 成本 $5$ 的巡迴每一
段都要是邊，也就是每一段都跨兩邊。 巡迴上 $3$、$4$、$5$ 各有兩個相鄰的城市，共 $6$ 個「相鄰位置」都要
是 $1$ 或 $2$；但 $1$、$2$ 在巡迴上各只有兩個鄰居，最多提供 $4$ 個。 不可能，所以至少一段成本是 $2$，
總成本至少 $6$。
\item 最小成本 $6 > B = 5$，\npcP{TSP} 實例回答否，由歸約的兩個方向，$K_{2,3}$ 沒有哈密頓環。 一般地，
兩邊頂點數不同的二分圖都沒有哈密頓環：環上兩邊的頂點必須交錯出現。
\end{enumerate}
\end{enumerate}

\chapter{近似演算法}
\label{ch:approximation}

% 本章的圖共用的樣式（前綴 apx-）
\tikzset{
    apx-v/.style={circle, draw=black!70, thick, minimum size=0.62cm, inner sep=0pt, font=\small, fill=white},
    apx-vC/.style={apx-v, draw=coreStroke, fill=coreFill},
    apx-vO/.style={apx-v, draw=greenStroke, fill=greenFill},
    apx-e/.style={draw=black!55, thick},
    apx-m/.style={draw=memStroke, line width=2.2pt},
    apx-pt/.style={circle, draw=black!70, thick, minimum size=0.5cm, inner sep=0pt, font=\scriptsize, fill=white},
    apx-grid/.style={draw=black!12, thin},
    apx-tree/.style={draw=coreStroke, very thick},
    apx-tour/.style={draw=memStroke, thick, -{Stealth[length=2mm]}},
    apx-opt/.style={draw=greenStroke, thick, -{Stealth[length=2mm]}},
    apx-w/.style={font=\scriptsize, text=black!70, inner sep=1pt, fill=white},
    apx-dot/.style={circle, fill=black!70, minimum size=0.16cm, inner sep=0pt},
    apx-row/.style={draw=greenStroke, thick, rounded corners=3pt},
    apx-blk/.style={draw=memStroke, thick, rounded corners=3pt, dashed},
    apx-cap/.style={font=\small},
}

校園裡有七個收件點，工友每天下午推著推車從收發室出發，每個點都要去一次，最後回到收發室。
路是棋盤狀的，兩點之間的距離就是橫向加縱向走的格數。 第~20~章的分支限界能算出最短的路線，
以「走過的點的集合」當子問題的動態規劃（第~18~章提過）也可以，可是點數一多，它們要花的時間
都會爆炸；第~28~章更說明了，這類問題很可能
根本沒有快的精確解法。 換個要求：如果有一個很快的做法，保證排出來的路線\textbf{不會超過最
短路線的兩倍}，你願意用嗎？ 再想一件事：最短路線有多長你根本不知道，要怎麼保證「不超過它的兩
倍」？ 往下讀之前，先想一想。

第~28~章說，NP 完備問題目前都沒有已知的多項式時間演算法。 這些問題的最佳化版本還是每天都要
解，常見的出路有三條：輸入很小的時候用回溯或分支限界（第~20~章）；輸入有特殊結構的時候找專
用的演算法（例如樹上的問題常常可以用動態規劃）；或者放棄「一定最佳」，改求\textbf{保證離最佳
不遠}的解。 這一章講第三條路。

上面那個問題的答案，是這一章所有證明共用的想法：最佳解的值不知道，但可以找一個算得出來的量，
證明最佳解一定不比它好，再證明演算法的答案不超過這個量的某個倍數。 第~\ref{sec:apx-ratio}~節
定義近似比；第~\ref{sec:apx-vc}~節用頂點覆蓋示範最簡單的 $2$-近似；第~\ref{sec:apx-tsp}~節用最
小生成樹做度量 TSP 的 $2$-近似，並證明沒有三角不等式的一般 TSP 連常數倍都保證不了；
第~\ref{sec:apx-sc}~節的集合覆蓋用貪心，近似比是 $H_d \le \ln d + 1$；
第~\ref{sec:apx-knap}~節回到第~18~章的 0/1 背包，看怎麼用動態規劃做到「誤差要多小都可以」。
每個近似比都附證明；前三個問題還各附一個讓比值逼近上界的例子，說明分析已經不能再改進。

這一章用到第~19~章的貪心與交換論證、第~21~章的圖論名詞、第~24~章的最小生成樹和第~28~章的 NP
完備，都只引用，不重講。 C 程式用到《C 程式設計基礎》第~12~章的位元運算、第~14~章的結構和第~19~章的
\texttt{malloc}。

\section{近似比}
\label{sec:apx-ratio}

\subsection{最佳化問題與近似演算法}

這一章的問題都是\textbf{最佳化問題}：每個輸入有很多個\textbf{可行解}（feasible solution），每
個可行解有一個\textbf{成本}（或價值），要找成本最小（或價值最大）的那一個。 最佳解的值寫成
$C^*$，演算法找到的解的值寫成 $C$。 本章假設所有可行解的值都是正的，比值才有意義。

\begin{definition}[近似比]
\label{def:apx-ratio}
一個演算法對大小 $n$ 的每一個輸入都滿足
\[
\max\left(\frac{C}{C^*},\ \frac{C^*}{C}\right) \le \rho(n),
\]
就說它的\textbf{近似比}（approximation ratio）是 $\rho(n)$，這個演算法叫\textbf{$\rho(n)$-近似演算
法}（$\rho(n)$-approximation algorithm）。 本章說的近似演算法都要在多項式時間內跑完。
\end{definition}

最小化問題 $C \ge C^*$，所以上式就是 $C \le \rho(n) \cdot C^*$；最大化問題 $C \le C^*$，上式是
$C \ge C^* / \rho(n)$。 兩種情形的 $\rho(n)$ 都 $\ge 1$，越接近 $1$ 越好，$\rho(n) = 1$ 就是精確
的演算法。 $\rho(n)$ 可以是常數（例如 $2$），也可以隨 $n$ 變大（例如 $\ln n + 1$）。

\begin{examnote}
近似比的寫法各本書不同。 有的書對最大化問題寫 $C \ge \alpha C^*$，$\alpha \le 1$，於是同一個
演算法在那裡叫「$\frac12$-近似」，在本書叫「$2$-近似」；有的書把「近似比 $\rho$」寫成「效能
保證」或「$\rho$ 倍近似」。 看清楚題目用哪一種，答案寫成「$C \le 2C^*$」或「$C \ge C^*/2$」這種
不等式，不會被誤解。
\end{examnote}

\subsection{不知道最佳解，怎麼證明}

證明 $C \le \rho C^*$ 最難的地方，是 $C^*$ 算不出來（算得出來就不必近似了）。 本章每個證明都
照同一個模式（以最小化為例）：

\begin{enumerate}
\item 找一個多項式時間算得出來、或至少能在證明裡談論的量 $L$，證明 $L \le C^*$。 $L$ 叫
\textbf{下界}（lower bound）。
\item 證明演算法的答案 $C \le \rho L$。
\item 合起來，$C \le \rho L \le \rho C^*$。
\end{enumerate}

章首的問題也是這樣回答：任何一條巡迴路線拿掉一段之後都是一棵生成樹，所以最短路線不會比最小
生成樹短；只要排出來的路線不超過最小生成樹的兩倍，就不超過最短路線的兩倍。 本章用到的下界
整理如下，後面各節再一一證明。

\begin{center}
\begin{tabular}{lll}
\toprule
問題 & 下界 $L$ & 近似比 \\
\midrule
頂點覆蓋 & 一組兩兩不共用端點的邊的條數 & $2$ \\
度量 TSP & 最小生成樹的權重 & $2$ \\
集合覆蓋 & 把成本分攤到元素上的「價格」 & $H_d$ \\
0/1 背包（最大化） & 四捨五入前後的價值差 & $1/(1 - \varepsilon)$ \\
\bottomrule
\end{tabular}
\end{center}

\subsection{近似比是緊的}

證明了 $C \le 2C^*$，只說明這個演算法\textbf{最差不會超過} $2$ 倍，沒說它真的會差到 $2$ 倍。 也
許更仔細的分析能證明 $1.5$ 倍。 如果能找到一族輸入，讓 $C / C^*$ 等於 $\rho$，或是任意接近
$\rho$，就說這個近似比的分析是\textbf{緊的}（tight）：對這個演算法，$\rho$ 已經不能再改小。

「緊」說的是\textbf{這個演算法}的分析，不是說這個問題不可能有更好的演算法。 例如
第~\ref{sec:apx-tsp}~節的 TSP 演算法，近似比 $2$ 是緊的，但另有一個演算法可以做到 $1.5$。 舉
緊的例子時，演算法的答案 $C$ 和最佳解 $C^*$ 都要算出來，而且 $C^*$ 要附理由（為什麼不可能更
好），只算演算法這一邊不夠。

\section{頂點覆蓋}
\label{sec:apx-vc}

\subsection{問題}

社區的巷弄畫成一張無向圖：路口是頂點，巷子是邊。 晚上要排志工在路口站崗，站在一個路口的志工
看得到從這個路口出去的每一條巷子；每條巷子至少要有一端的路口有人。 最少要排幾個人？

\begin{definition}[頂點覆蓋]
\label{def:apx-vc}
無向圖 $G = (V, E)$ 的\textbf{頂點覆蓋}（vertex cover）是一個頂點集合 $C \subseteq V$，使得每
一條邊 $(u, v) \in E$ 都有 $u \in C$ 或 $v \in C$（也可以兩端都在）。 \textbf{最小頂點覆蓋}問
題要找頂點數最少的頂點覆蓋。
\end{definition}

整個 $V$ 一定是頂點覆蓋，難的是找最小的。 第~28~章證明了頂點覆蓋的決策版本（「有沒有大小
$\le k$ 的頂點覆蓋」）是 NP 完備的，所以除非 P $=$ NP，最小頂點覆蓋沒有多項式時間的演算法。

\subsection{演算法}

想法很簡單：看到一條還沒被蓋住的邊 $(u, v)$，不知道 $u$、$v$ 哪一個在最佳解裡，就兩個都拿。
這樣浪費最多一半，下面會證明。

\begin{pseudo}[label=alg:apx-vc]{頂點覆蓋的 $2$-近似}
\Function{Approx-Vertex-Cover}{$G$}
    \State $C \gets \emptyset$
    \For{$E$ 裡的每一條邊 $(u, v)$} \Comment{依邊表的順序}
        \If{$u \notin C$ \AND $v \notin C$}
            \State $C \gets C \cup \{u, v\}$
        \EndIf
    \EndFor
    \State \Return $C$
\EndFunction
\end{pseudo}

用一個長度 $n$ 的布林陣列記錄「在不在 $C$ 裡」，第 4、5 行都是 $O(1)$，總時間 $\Theta(n + m)$。
第 5 行把兩端都加進來的那些邊，叫做\textbf{被選中的邊}，集合寫成 $M$。

\subsection{手算追蹤}

圖~\ref{fig:apx-vc} 左邊的圖有 $n = 8$ 個路口、$m = 10$ 條巷子。 邊表依字典順序排好：
$(1,2), (1,3), (2,3), (2,4), (3,4), (4,5), (4,6), (5,7), (6,7), (7,8)$。

\begin{center}
\begin{tabular}{clc}
\toprule
邊 & 動作 & $C$ \\
\midrule
$(1, 2)$ & 兩端都不在 $C$，選中 & $\{1, 2\}$ \\
$(1, 3)$ & $1 \in C$，已蓋住 & \\
$(2, 3)$ & 已蓋住 & \\
$(2, 4)$ & 已蓋住 & \\
$(3, 4)$ & 兩端都不在 $C$，選中 & $\{1, 2, 3, 4\}$ \\
$(4, 5)$、$(4, 6)$ & 已蓋住 & \\
$(5, 7)$ & 兩端都不在 $C$，選中 & $\{1, 2, 3, 4, 5, 7\}$ \\
$(6, 7)$、$(7, 8)$ & 已蓋住 & \\
\bottomrule
\end{tabular}
\end{center}

演算法排了 $6$ 個人。 最少要幾個？ 三角形 $1, 2, 3$ 的三條邊，一個頂點最多蓋兩條，所以
$\{1, 2, 3\}$ 裡至少要選兩個；$4, 5, 7, 6$ 圍成一個四邊形，四條邊也至少要兩個頂點才蓋得住。
兩組頂點沒有重疊，所以 $C^* \ge 4$；$\{1, 2, 4, 7\}$ 是大小 $4$ 的頂點覆蓋（右圖），$C^* = 4$。
這個例子 $C / C^* = 1.5$。

\begin{center}
\begin{tikzpicture}[x=1cm, y=1cm]
\begin{scope}
\node[apx-vC] (a1) at (0, 0) {1};
\node[apx-vC] (a2) at (1.2, 0.85) {2};
\node[apx-vC] (a3) at (1.2, -0.85) {3};
\node[apx-vC] (a4) at (2.4, 0) {4};
\node[apx-vC] (a5) at (3.6, 0.85) {5};
\node[apx-v] (a6) at (3.6, -0.85) {6};
\node[apx-vC] (a7) at (4.8, 0) {7};
\node[apx-v] (a8) at (6.0, 0) {8};
\foreach \a/\b in {1/3, 2/3, 2/4, 4/5, 4/6, 6/7, 7/8}
    \draw[apx-e] (a\a) -- (a\b);
\foreach \a/\b in {1/2, 3/4, 5/7}
    \draw[apx-m] (a\a) -- (a\b);
\node[apx-cap] at (3, -1.6) {演算法：$C = \{1, 2, 3, 4, 5, 7\}$};
\end{scope}
\begin{scope}[xshift=7.4cm]
\node[apx-vO] (b1) at (0, 0) {1};
\node[apx-vO] (b2) at (1.2, 0.85) {2};
\node[apx-v] (b3) at (1.2, -0.85) {3};
\node[apx-vO] (b4) at (2.4, 0) {4};
\node[apx-v] (b5) at (3.6, 0.85) {5};
\node[apx-v] (b6) at (3.6, -0.85) {6};
\node[apx-vO] (b7) at (4.8, 0) {7};
\node[apx-v] (b8) at (6.0, 0) {8};
\foreach \a/\b in {1/2, 1/3, 2/3, 2/4, 3/4, 4/5, 4/6, 5/7, 6/7, 7/8}
    \draw[apx-e] (b\a) -- (b\b);
\node[apx-cap] at (3, -1.6) {最佳解：$C^* = \{1, 2, 4, 7\}$};
\end{scope}
\end{tikzpicture}
\diagcaption{左：虛擬碼~\ref{alg:apx-vc} 的結果，粗的橘色邊是被選中的三條邊，藍色頂點在 $C$ 裡。
右：最小頂點覆蓋（綠色）。}
\label{fig:apx-vc}
\end{center}

\subsection{為什麼是 2-近似}

\begin{theorem}
\label{thm:apx-vc}
\proc{Approx-Vertex-Cover} 傳回的 $C$ 是頂點覆蓋，而且 $|C| \le 2\,|C^*|$。
\end{theorem}

\begin{proof}
\textbf{是覆蓋。} 看任何一條邊 $(u, v)$。 第 3 行輪到它的時候，要不是 $u$ 或 $v$ 已經在 $C$
裡，就是它被選中、兩端都加進 $C$。 $C$ 只會變大，所以結束時 $(u, v)$ 一定被蓋住。

\textbf{下界。} 被選中的邊兩兩沒有共用端點：一條邊被選中時，它的兩端都還不在 $C$ 裡，所以沒有
碰到之前選中的任何一條邊；選中之後兩端都進了 $C$，之後碰到這兩點的邊都不會再被選中。 任何
頂點覆蓋，包括最佳的 $C^*$，都要蓋住 $M$ 的每一條邊，每條至少用掉一個端點；$M$ 的邊沒有共用
端點，用掉的端點兩兩不同。 所以 $|C^*| \ge |M|$。

\textbf{上界。} 每選中一條邊，$C$ 剛好多兩個頂點（兩端原本都不在 $C$ 裡），所以 $|C| = 2|M|$。
合起來，$|C| = 2|M| \le 2|C^*|$。
\end{proof}

證明裡的 $M$ 在圖論裡叫做\textbf{匹配}（matching）：兩兩沒有共用端點的一群邊。 演算法結束時，
每條邊都至少碰到 $M$ 的一個端點，所以 $M$ 不能再多加任何一條邊，這叫\textbf{極大匹配}
（maximal matching）。 「極大」是加不進去，不是最多；邊數最多的匹配叫\textbf{最大匹配}
（maximum matching），第~26~章用網路流求二分圖的最大匹配。 圖~\ref{fig:apx-vc} 的 $M$ 有三條
邊，於是 $C^* \ge 3$；實際上 $C^* = 4$，這個下界不一定剛好等於 $C^*$。

\subsection{近似比是緊的}

讓 $C / C^*$ 剛好是 $2$ 的例子很多。 最簡單的是 $k$ 條互不相連的邊：每一條都被選中，
$|C| = 2k$，而每條邊各要一個端點，$C^* = k$。 連通的例子也有：完全二分圖 $K_{k,k}$，左邊是
$1..k$，右邊是 $k+1..2k$，每個左邊的點和每個右邊的點都有邊，邊表依字典順序。 掃到 $(i, k+i)$
時，前面選中的邊是 $(1, k+1), \ldots, (i-1, k+i-1)$，都碰不到 $i$ 和 $k+i$，所以它被選中；其他的
邊 $(i, k+j)$，$j \ne i$，一定有一端已經在 $C$ 裡（$j < i$ 時 $k+j$ 已經在，$j > i$ 時 $i$ 已經
在）。 結果 $C$ 是全部 $2k$ 個點；只取左邊 $k$ 個點就蓋住了所有邊，而 $(1, k+1), \ldots,
(k, 2k)$ 這 $k$ 條邊沒有共用端點，至少要 $k$ 個點，所以 $C^* = k$。

還有一個問題：證明裡用的下界 $|M|$ 本身會不會太鬆？ 會。 $2k+1$ 個點的完全圖 $K_{2k+1}$，任何
匹配最多 $k$ 條邊，但頂點覆蓋至少要 $2k$ 個點（漏掉兩個點，它們之間的邊就沒蓋到）。 所以只靠
「$C^* \ge |M|$」這個下界，不可能證明比 $2$ 更好的近似比。 目前沒有人知道比 $2$ 好一個常數的
多項式時間演算法；這是近似演算法裡很有名的未解問題。

\subsection{兩個看起來更聰明的做法}

\textbf{只拿一端。} 每條沒蓋住的邊只拿一個端點，例如編號小的那個，好像比較省。 星狀圖（中心
$n$ 連到 $1, 2, \ldots, n-1$）就會讓它拿了 $n - 1$ 個葉節點，最佳解只要中心一個點。 不知道該拿哪
一端時，兩端都拿，反而能保證比值。

\textbf{先拿度數最大的。} 每次拿目前蓋住最多條還沒蓋住的邊的頂點，直覺上很合理，第~19~章也
提醒過這種貪心要先找反例。 一隻三條腿的「蜘蛛」：中心 $7$ 連到 $1, 2, 3$，再接上 $(1,4)$、
$(2,5)$、$(3,6)$。 中心的度數是 $3$，其他點都 $\le 2$，貪心先拿 $7$；剩下的三條邊 $(1,4), (2,5), (3,6)$ 互不相連，還要再拿三個點，共
$4$ 個。 最佳解是 $\{1, 2, 3\}$，$3$ 個。

這個貪心其實是第~\ref{sec:apx-sc}~節集合覆蓋貪心的特例，由定理~\ref{thm:apx-sc} 它的近似比不超過
$H_\Delta \le \ln \Delta + 1$，$\Delta$ 是最大度數；已經知道有些圖會讓它差到 $\Theta(\log n)$ 倍，
所以它\textbf{沒有}常數的近似比，雖然在一般的圖上它常常比虛擬碼~\ref{alg:apx-vc} 的結果小。

\subsection{C 程式}

\texttt{approx\_vc} 照虛擬碼~\ref{alg:apx-vc} 掃描邊表；\texttt{optimal\_vc} 試遍 $2^n$ 個頂點子集
合（窮舉所有子集合，用 \texttt{unsigned} 的位元表示子集合），求出最小頂點覆蓋的大小，只能用在很小
的圖上。 \texttt{main} 先跑圖~\ref{fig:apx-vc} 的例子，再產生 $2000$ 張 $12$ 個點的隨機圖，統計比
值的分布。 書上的頂點是 $1..n$，程式裡一律減一，用 $0..n-1$。

\begin{cbookrecap}[第~12~章位元運算、第~19~章 19.3 節]
\texttt{(s >> x) \& 1u} 取出 \texttt{s} 的第 \texttt{x} 個位元：子集合 \texttt{s} 含頂點
\texttt{x} 時是 $1$。 \texttt{1u << n} 是 $2^n$，\texttt{u} 表示無號整數，$n \le 20$ 時不會溢位。
\texttt{malloc} 要到的記憶體用完要 \texttt{free}；本書的程式都用 \texttt{-fsanitize=address} 檢
查，漏掉一個 \texttt{free} 也會被抓到。
\end{cbookrecap}

\begin{example}[label=ex:apx-vc]{c}{(vc\_approx.c) 頂點覆蓋的 2-近似和暴力最佳解}
#include <stdio.h>
#include <stdlib.h>

#define N 12                    // 隨機圖的頂點數

struct edge {
    int u, v;                   // 兩個端點，0..n-1
};

// 2-近似：依邊的順序掃描，兩端都還沒選的邊就把兩端都選進來。
// in_c[x] 記錄 x 有沒有被選，回傳選了幾個點。
int approx_vc(int n, int m, const struct edge *e, char *in_c)
{
    int size = 0;
    for (int x = 0; x < n; x++) {
        in_c[x] = 0;
    }
    for (int i = 0; i < m; i++) {
        if (!in_c[e[i].u] && !in_c[e[i].v]) {
            in_c[e[i].u] = in_c[e[i].v] = 1;
            size += 2;
        }
    }
    return size;
}

// 暴力：試遍 2^n 個子集合 s，回傳最小頂點覆蓋的大小。
int optimal_vc(int n, int m, const struct edge *e)
{
    int opt = n;                // 全部的頂點一定是覆蓋
    for (unsigned s = 0; s < (1u << n); s++) {
        int cnt = 0, ok = 1;
        for (int x = 0; x < n; x++) {
            cnt += (s >> x) & 1u;
        }
        for (int i = 0; i < m && ok; i++) {
            ok = ((s >> e[i].u) & 1u) || ((s >> e[i].v) & 1u);
        }
        if (ok && cnt < opt) {
            opt = cnt;
        }
    }
    return opt;
}

// 線性同餘亂數，每次執行都產生一樣的圖。
static unsigned long long seed = 2029;
static unsigned next_rand(void)
{
    seed = seed * 6364136223846793005ULL + 1442695040888963407ULL;
    return (unsigned)(seed >> 33);
}

int main(void)
{
    // 上面手算的例子：8 個點，邊依字典順序，頂點寫成 1..8。
    const int ex[10][2] = {{1, 2}, {1, 3}, {2, 3}, {2, 4}, {3, 4},
                           {4, 5}, {4, 6}, {5, 7}, {6, 7}, {7, 8}};
    struct edge *e = malloc(N * (N - 1) / 2 * sizeof *e);
    char *in_c = malloc(N);
    if (e == NULL || in_c == NULL) {
        free(e);
        free(in_c);
        return 1;
    }
    for (int i = 0; i < 10; i++) {
        e[i].u = ex[i][0] - 1;  // 程式裡用 0..n-1
        e[i].v = ex[i][1] - 1;
    }
    printf("approx (%d):", approx_vc(8, 10, e, in_c));
    for (int x = 0; x < 8; x++) {
        if (in_c[x]) {
            printf(" %d", x + 1);
        }
    }
    printf("\noptimal: %d\n", optimal_vc(8, 10, e));

    // 2000 張隨機圖，每一對頂點以 1/4 的機率有邊。
    int cnt[3] = {0, 0, 0};     // 比值是 1、在 1 和 2 之間、是 2
    for (int t = 0; t < 2000; t++) {
        int m = 0;
        for (int u = 0; u < N; u++) {
            for (int v = u + 1; v < N; v++) {
                if (next_rand() % 4 == 0) {
                    e[m].u = u;
                    e[m].v = v;
                    m++;
                }
            }
        }
        int a = approx_vc(N, m, e, in_c);
        int o = optimal_vc(N, m, e);
        cnt[a == o ? 0 : (a == 2 * o ? 2 : 1)]++;
    }
    printf("ratio 1: %d, between: %d, ratio 2: %d\n", cnt[0], cnt[1], cnt[2]);
    free(e);
    free(in_c);
    return 0;
}
\end{example}
\iosample{}{approx (6): 1 2 3 4 5 7 \\ optimal: 4 \\ ratio 1: 7, between: 1576, ratio 2: 417}

前兩行和上面的手算一致。 隨機圖的部分，$2000$ 張圖沒有一張超過 $2$ 倍，符合
定理~\ref{thm:apx-vc}；比值剛好是 $1$ 的只有 $7$ 張，剛好是 $2$ 的反而有 $417$ 張。 比值是 $2$
表示 $|C^*| = |M|$，也就是最小覆蓋在每一條被選中的邊上只用了一個端點；這種很稀疏的圖常常是這
樣。 近似比是最壞情況的保證，實際的比值要看輸入。

\section{度量 TSP}
\label{sec:apx-tsp}

\subsection{問題與三角不等式}

第~20~章用分支限界解過\textbf{旅行推銷員問題}（TSP）。 這一節用對稱的版本：$G = (V, E)$ 是
完全無向圖（每一對頂點都有邊），每條邊有成本 $c(u, v) = c(v, u) \ge 0$。 一條\textbf{巡迴路線}
（tour）從某個頂點出發，每個頂點恰好經過一次，再回到出發點，也就是一個經過所有頂點的環
（\textbf{哈密頓環}，Hamiltonian cycle）；它的成本是環上 $n$ 條邊的成本和。 要找成本最小的巡迴
路線，最小成本寫成 $c(H^*)$。 第~28~章說明了 TSP 是 NP 難的。

章首的工友走的是棋盤狀的路，成本有一個特別的性質：從 $u$ 直接走到 $w$，不會比繞經 $v$ 遠。

\begin{definition}[三角不等式]
\label{def:apx-triangle}
成本函數 $c$ 滿足\textbf{三角不等式}（triangle inequality），如果對任意三個頂點 $u, v, w$，
\[
c(u, w) \le c(u, v) + c(v, w).
\]
滿足三角不等式的 TSP 叫\textbf{度量 TSP}（metric TSP）。
\end{definition}

平面上的直線距離、棋盤狀街道的曼哈頓距離 $|x_1 - x_2| + |y_1 - y_2|$，還有「在一張道路圖上
的最短路徑長」（第~25~章），都滿足三角不等式。 把三角不等式用很多次，可以得到：沿著任何一條
路徑 $\langle v_0, v_1, \ldots, v_k \rangle$ 走，成本不會比直接從 $v_0$ 走到 $v_k$ 少，
\[
c(v_0, v_k) \le c(v_0, v_1) + c(v_1, v_2) + \cdots + c(v_{k-1}, v_k).
\]
這個性質叫「抄捷徑不會變遠」，下面的證明就靠它。

\subsection{下界：最小生成樹}

\begin{lemma}
\label{lem:apx-mst-lb}
$T$ 是 $G$ 的最小生成樹，$H^*$ 是最佳巡迴路線，則 $c(T) \le c(H^*)$。 這個引理不需要三角不等
式，只需要成本 $\ge 0$。
\end{lemma}

\begin{proof}
從 $H^*$ 拿掉任何一條邊，剩下一條經過所有頂點的路徑，它連通、沒有環，是一棵生成樹，成本
$\le c(H^*)$（拿掉的邊成本 $\ge 0$）。 最小生成樹的成本不超過任何一棵生成樹，所以
$c(T) \le c(H^*)$。
\end{proof}

\subsection{演算法：最小生成樹的前序走訪}

從根出發，沿著最小生成樹 $T$ 做深度優先的走訪，每條樹邊走下去一次、走回來一次，最後回到根，
這樣的\textbf{完整走訪}（full walk）經過每個頂點，成本剛好 $2c(T)$。 但它會重複經過頂點，不是巡
迴路線。 把重複的頂點刪掉，只留每個頂點第一次出現的位置，剩下的順序就是 $T$ 的\textbf{前序走
訪}（第~7~章）；刪掉一個重複的頂點，等於從它前一個點直接走到它後一個點，由三角不等式成本不會
變大。

\begin{pseudo}[label=alg:apx-tsp]{度量 TSP 的 $2$-近似}
\Function{Approx-TSP-Tour}{$G, c$}
    \State $r \gets 1$ \Comment{任選一個頂點當根}
    \State $T \gets$ 以 $r$ 為根的最小生成樹 \Comment{第~24~章的 Prim}
    \State $L \gets$ 空串列
    \State \Call{Preorder}{$T, r, L$}
    \State \Return 依 $L$ 的順序走過每個頂點，最後回到 $r$
\EndFunction
\Statex
\Function{Preorder}{$T, u, L$}
    \State 把 $u$ 接在 $L$ 的尾端
    \For{$u$ 在 $T$ 裡的每一個子節點 $v$} \Comment{編號由小到大}
        \State \Call{Preorder}{$T, v, L$}
    \EndFor
\EndFunction
\end{pseudo}

完全圖有 $m = n(n-1)/2$ 條邊，第 3 行用陣列版的 Prim（不用堆積）是 $\Theta(n^2)$，前序走訪
$\Theta(n)$，總時間 $\Theta(n^2)$，和輸入（成本矩陣）的大小同階。

\begin{theorem}
\label{thm:apx-tsp}
成本滿足三角不等式時，\proc{Approx-TSP-Tour} 傳回的巡迴路線 $H$ 滿足 $c(H) \le 2\,c(H^*)$。
\end{theorem}

\begin{proof}
令 $W$ 是 $T$ 從 $r$ 開始的完整走訪（子節點的順序和 \proc{Preorder} 一樣）。 每條樹邊在 $W$ 裡剛
好走兩次，一次往下、一次往上，所以 $c(W) = 2c(T)$。

$W$ 裡每個頂點第一次出現的順序，就是 \proc{Preorder} 產生的 $L$：\proc{Preorder} 第一次碰到 $u$ 時
把它接到 $L$，完整走訪也是在這時第一次走到 $u$。 從 $W$ 刪掉不是第一次出現的頂點（最後回到 $r$
的那一次保留），得到的就是 $H$。 $H$ 的每一條邊 $(x, y)$ 都對應到 $W$ 裡從 $x$ 到 $y$ 的一段
路，中間經過的都是被刪掉的頂點；不同的邊對應的段落不重疊，合起來正好是整個 $W$。 由三角不等
式，$c(x, y)$ 不超過那一段路的成本，加總得到 $c(H) \le c(W)$。

合起來，$c(H) \le c(W) = 2c(T) \le 2c(H^*)$，最後一步是引理~\ref{lem:apx-mst-lb}。
\end{proof}

\subsection{手算追蹤}

章首的七個收件點，座標（單位是一個街區）如下，$1$ 是收發室；成本是曼哈頓距離，例如
$c(1, 5) = |0 - 5| + |2 - 3| = 6$。
\[
1\,(0, 2),\quad 2\,(0, 5),\quad 3\,(2, 1),\quad 4\,(2, 2),\quad 5\,(5, 3),\quad 6\,(5, 5),\quad 7\,(7, 0)
\]

\textbf{第一步：最小生成樹。} 用第~24~章的 Prim 從 $1$ 開始，頂點加入的順序是
$1, 4, 3, 2, 5, 6, 7$，樹邊和成本是
\[
(1,4)\ 2,\quad (4,3)\ 1,\quad (1,2)\ 3,\quad (4,5)\ 4,\quad (5,6)\ 2,\quad (5,7)\ 5,
\]
$c(T) = 17$。 這張圖的最小生成樹只有這一棵（Kruskal 也得到同樣的邊），所以平手規則不影響結果。
以 $1$ 為根，$1$ 的子節點是 $2, 4$；$4$ 的子節點是 $3, 5$；$5$ 的子節點是 $6, 7$。

\textbf{第二步：完整走訪與前序。} 子節點依編號由小到大，完整走訪是
\[
1,\ 2,\ 1,\ 4,\ 3,\ 4,\ 5,\ 6,\ 5,\ 7,\ 5,\ 4,\ 1,
\]
成本 $2 \times 17 = 34$。 只留第一次出現的頂點，得到前序 $1, 2, 4, 3, 5, 6, 7$，巡迴路線是
$1 \to 2 \to 4 \to 3 \to 5 \to 6 \to 7 \to 1$。

\textbf{第三步：抄捷徑省了多少。} 下表列出每一段捷徑和它取代的那段完整走訪：

\begin{center}
\begin{tabular}{ccccc}
\toprule
捷徑 & 取代的路 & 原本成本 & 捷徑成本 & 省下 \\
\midrule
$2 \to 4$ & $2, 1, 4$ & $3 + 2 = 5$ & $5$ & $0$ \\
$3 \to 5$ & $3, 4, 5$ & $1 + 4 = 5$ & $5$ & $0$ \\
$6 \to 7$ & $6, 5, 7$ & $2 + 5 = 7$ & $7$ & $0$ \\
$7 \to 1$ & $7, 5, 4, 1$ & $5 + 4 + 2 = 11$ & $9$ & $2$ \\
\bottomrule
\end{tabular}
\end{center}

其他三段 $1 \to 2$、$4 \to 3$、$5 \to 6$ 本來就是樹邊。 巡迴路線的成本是
$3 + 5 + 1 + 5 + 2 + 7 + 9 = 32 = 34 - 2$。 曼哈頓距離下，只要繞的路沒有「折回頭」，捷徑就省不
到東西，所以前三段都沒省。

最佳路線是 $1 \to 2 \to 6 \to 5 \to 7 \to 3 \to 4 \to 1$，成本 $3 + 5 + 2 + 5 + 6 + 1 + 2 = 24$。
為什麼不可能更短？ 七個點的 $x$ 座標從 $0$ 到 $7$，任何巡迴路線都要從最左邊走到最右邊、再走回
來，橫向至少走 $2 \times 7 = 14$；同理縱向至少 $2 \times 5 = 10$。 曼哈頓距離是橫向加縱向，所以
每條路線至少 $24$。 三個數的關係是
\[
c(T) = 17 \;\le\; c(H^*) = 24 \;\le\; c(H) = 32 \;\le\; 2c(T) = 34 ,
\]
比值 $32 / 24 \approx 1.33$。 注意證明用的是 $c(H) \le 2c(T)$，不需要知道 $24$ 這個數。

\begin{center}
\begin{tikzpicture}[x=0.55cm, y=0.55cm]
\foreach \k/\ttl in {0/最小生成樹 $T$（$17$）, 1/近似路線 $H$（$32$）, 2/最佳路線 $H^*$（$24$）} {
    \begin{scope}[xshift=\k*5.06cm]
    \draw[apx-grid] (0, 0) grid (7, 5);
    \node[apx-cap] at (3.5, -1) {\ttl};
    \node[apx-pt] (p1-\k) at (0, 2) {1};
    \node[apx-pt] (p2-\k) at (0, 5) {2};
    \node[apx-pt] (p3-\k) at (2, 1) {3};
    \node[apx-pt] (p4-\k) at (2, 2) {4};
    \node[apx-pt] (p5-\k) at (5, 3) {5};
    \node[apx-pt] (p6-\k) at (5, 5) {6};
    \node[apx-pt] (p7-\k) at (7, 0) {7};
    \end{scope}
}
\begin{scope}[on background layer]
\foreach \a/\b in {1/4, 4/3, 1/2, 4/5, 5/6, 5/7}
    \draw[apx-tree] (p\a-0) -- (p\b-0);
\end{scope}
\foreach \a/\b in {1/2, 2/4, 4/3, 3/5, 5/6, 6/7, 7/1}
    \draw[apx-tour] (p\a-1) -- (p\b-1);
\foreach \a/\b in {1/2, 2/6, 6/5, 5/7, 7/3, 3/4, 4/1}
    \draw[apx-opt] (p\a-2) -- (p\b-2);
\end{tikzpicture}
\diagcaption{七個收件點。 圖上畫的是直線，實際的成本是曼哈頓距離（沿著格線走）。}
\label{fig:apx-tsp}
\end{center}

\begin{examnote}
追蹤題寫出三樣東西：最小生成樹的邊（和用哪個演算法、從哪個點開始）、前序走訪的順序（子節點依什麼
順序）、路線的總成本。 總成本要包含\textbf{最後回到起點}的那一段，這一段最常漏。 題目沒規定子節點
的順序時，寫一句你的約定；換一個順序，路線和成本都可能不同，但都不超過 $2c(T)$。
\end{examnote}

\subsection{近似比是緊的}

要讓 $c(H)$ 接近 $2c(H^*)$，需要兩件事同時發生：抄捷徑幾乎沒省到（$c(H) \approx 2c(T)$），而且
最佳路線幾乎和最小生成樹一樣便宜（$c(H^*) \approx c(T)$）。 下面的例子成本只有 $1$ 和 $2$ 兩種，
任何這樣的成本都滿足三角不等式（$2 \le 1 + 1$）。

取 $n = k + 1$ 個頂點，$k \ge 2$：頂點 $1$ 是中心，$2, 3, \ldots, k + 1$ 是「葉」。 成本規定為
\[
c(u, v) =
\begin{cases}
1 & u, v \text{ 有一個是 } 1\text{，或 } u, v \text{ 都是葉而且 } |u - v| = 2, \\
2 & \text{其他。}
\end{cases}
\]
也就是說，中心到每片葉都是 $1$；葉之間只有編號差 $2$ 的是 $1$，編號差 $1$ 的都是 $2$。
圖~\ref{fig:apx-tsp-tight} 是 $k = 6$ 的情形。

\textbf{演算法的結果。} Prim 從 $1$ 開始，第一步就把每片葉的 key 設成 $c(1, v) = 1$，這已經是最小
的成本，之後不會再變，所以每片葉的父節點都是 $1$，$T$ 是以 $1$ 為中心的星狀樹，$c(T) = k$。 前序
走訪是 $1, 2, 3, \ldots, k + 1$，相鄰兩片葉的編號差 $1$，成本都是 $2$：
\[
c(H) = \underbrace{1}_{1 \to 2} + \underbrace{2(k - 1)}_{\text{葉到葉}} + \underbrace{1}_{k+1 \to 1}
= 2k = 2c(T).
\]

\textbf{最佳路線。} 路線 $1$，所有偶數葉由小到大，所有奇數葉由大到小，再回到 $1$，只有「最大的偶數
葉到最大的奇數葉」這一段成本是 $2$（兩者編號差 $1$），其他都是 $1$，總成本 $n + 1 = k + 2$。 不
可能更便宜：路線有 $n$ 段，每段至少 $1$；如果全部是 $1$，拿掉中心之後，剩下的是一條只用成本 $1$
的邊、經過所有葉的路徑，但成本 $1$ 的葉間邊只連編號同奇偶的葉，偶數葉和奇數葉之間沒有這種邊，
走不通。 所以 $c(H^*) = k + 2$，
\[
\frac{c(H)}{c(H^*)} = \frac{2k}{k + 2} \;\longrightarrow\; 2 \qquad (k \to \infty).
\]

\begin{center}
\begin{tikzpicture}[x=1cm, y=1cm]
\foreach \k/\ttl in {0/$T$ 和成本 $1$ 的葉間邊, 1/近似路線（$12$）, 2/最佳路線（$8$）} {
    \begin{scope}[xshift=\k*4.6cm]
    \node[apx-v] (c-\k) at (0, 0) {1};
    \foreach \v [count=\i] in {2, 3, 4, 5, 6, 7}
        \node[apx-v] (l\v-\k) at ({150 - 60 * (\i - 1)}:1.6) {\v};
    \node[apx-cap] at (0, -2.3) {\ttl};
    \end{scope}
}
\foreach \v in {2, 3, 4, 5, 6, 7}
    \draw[apx-tree] (c-0) -- (l\v-0);
\foreach \a/\b in {2/4, 4/6, 3/5, 5/7}
    \draw[apx-e] (l\a-0) -- (l\b-0);
\draw[apx-tour] (c-1) -- (l2-1);
\foreach \a/\b in {2/3, 3/4, 4/5, 5/6, 6/7}
    \draw[apx-tour] (l\a-1) -- (l\b-1);
\draw[apx-tour] (l7-1) -- (c-1);
\draw[apx-opt] (c-2) -- (l2-2);
\foreach \a/\b in {2/4, 4/6, 6/7, 7/5, 5/3}
    \draw[apx-opt] (l\a-2) -- (l\b-2);
\draw[apx-opt] (l3-2) -- (c-2);
\end{tikzpicture}
\diagcaption{$k = 6$ 的緊例子。 左圖畫出的邊成本是 $1$，沒畫出的是 $2$。 近似路線有五段葉到葉、
成本都是 $2$；最佳路線只有 $6 \to 7$ 一段成本 $2$。}
\label{fig:apx-tsp-tight}
\end{center}

$2$ 不是度量 TSP 最好的近似比。 Christofides（1976）把演算法改成：在 $T$ 裡度數是奇數的頂點之間
加一組成本最小的完美匹配，讓每個頂點的度數都變成偶數，再走一條經過每條邊一次的路線（歐拉迴路），
最後抄捷徑。 他證明這樣的成本不超過 $1.5\,c(H^*)$。 這個 $1.5$ 保持了四十多年，2021 年 Karlin、
Klein 與 Oveis Gharan 才用隨機的方法把它改進了一個極小的量。 本書不證明這些結果。

\subsection{一般的 TSP 不能近似}
\label{sec:apx-tsp-general}

拿掉三角不等式，情況完全不同：連「不超過最佳的一百萬倍」都做不到。

\begin{theorem}
\label{thm:apx-tsp-hard}
如果 P $\ne$ NP，那麼對任何常數 $\rho \ge 1$，一般的（成本不一定滿足三角不等式的）TSP 都沒有多
項式時間的 $\rho$-近似演算法。
\end{theorem}

\begin{proof}
用反證法。 假設某個多項式時間的演算法 $A$ 是 $\rho$-近似的。 用它在多項式時間內判斷任何無向圖
$G = (V, E)$ 有沒有哈密頓環；第~28~章說明了這個決策問題是 NP 完備的，所以這會推出 P $=$ NP。

給 $G$，$n = |V|$，在同一組頂點上建一個完全圖，成本定為
\[
c(u, v) =
\begin{cases}
1 & (u, v) \in E, \\
\rho n + 1 & (u, v) \notin E .
\end{cases}
\]
$\rho$ 是常數，這個成本矩陣有 $n^2$ 格，每一格的數只要 $O(\lg n)$ 個位元，建表是多項式時間。

\begin{itemize}
\item 如果 $G$ 有哈密頓環，那條環的每一段都是 $G$ 的邊，成本 $n$，所以 $c(H^*) = n$（每段至少
$1$）。 $A$ 傳回的路線成本 $\le \rho n$。
\item 如果 $G$ 沒有哈密頓環，任何巡迴路線都至少有一段不是 $G$ 的邊，成本至少
$(\rho n + 1) + (n - 1) > \rho n$。 $A$ 傳回的路線成本也一樣 $> \rho n$。
\end{itemize}

所以「$A$ 傳回的成本 $\le \rho n$」若且唯若「$G$ 有哈密頓環」。 建表、執行 $A$、比較一次，都是多
項式時間，這就在多項式時間內回答了哈密頓環問題。
\end{proof}

這個證明裡，$c(u, w) = \rho n + 1$ 卻有 $c(u, v) + c(v, w) = 2$ 的情形，三角不等式被嚴重破壞；
定理~\ref{thm:apx-tsp} 的捷徑論證在這裡完全失效。 這個證明的做法是：把一個 NP 完備的決策問題
變成最佳化問題，讓「是」的實例最佳值很小、「否」的實例最佳值超過 $\rho$ 倍，兩者之間留一個近似
演算法跨不過去的\textbf{間隙}（gap）。

\section{集合覆蓋}
\label{sec:apx-sc}

\subsection{問題}

營隊有 $12$ 項工作（編號 $1$ 到 $12$），六位志工報名，每位會做其中幾項：

\begin{center}
\begin{tabular}{cl}
\toprule
志工 & 會做的工作 \\
\midrule
$S_1$ & $2, 5, 7$ \\
$S_2$ & $4, 5, 12$ \\
$S_3$ & $7, 9, 11, 12$ \\
$S_4$ & $4, 6, 8, 10$ \\
$S_5$ & $3, 4, 5, 6, 7, 10$ \\
$S_6$ & $1, 2, 3, 5$ \\
\bottomrule
\end{tabular}
\end{center}

最少要找幾位，才能讓每項工作都有人會做？

\begin{definition}[集合覆蓋]
\label{def:apx-sc}
給一個有限集合 $X$（$|X| = n$）和一族子集合 $\mathcal{F} = \{S_1, S_2, \ldots, S_k\}$，每個
$S_i \subseteq X$，而且 $S_1 \cup \cdots \cup S_k = X$。 $\mathcal{C} \subseteq \mathcal{F}$ 是一
個\textbf{集合覆蓋}（set cover），如果 $\mathcal{C}$ 裡的集合聯集起來是 $X$。 \textbf{最小集合覆
蓋}問題要找集合個數最少的覆蓋。 \textbf{加權}的版本每個集合有成本 $c(S) > 0$，要讓總成本最小；
所有 $c(S) = 1$ 就是不加權的版本。 另外記 $d = \max_{S \in \mathcal{F}} |S|$，最大的集合有幾個元素。
\end{definition}

頂點覆蓋是集合覆蓋的特例：令 $X = E$，每個頂點 $v$ 對應一個集合 $S_v$，內容是和 $v$ 關聯的邊。
選一群頂點蓋住每條邊，就是選一群 $S_v$ 把 $E$ 聯集出來；這時 $d$ 是最大度數 $\Delta$。 所以集合
覆蓋至少和頂點覆蓋一樣難，決策版本也是 NP 完備的。

\subsection{貪心演算法}

不加權時，每次選\textbf{涵蓋最多個還沒被涵蓋的元素}的集合。 加權時，比較的是「每多涵蓋一個新
元素，平均要花多少」，也就是 $c(S) / |S \cap U|$，$U$ 是還沒被涵蓋的元素；選平均花費最少的。
不加權時 $c(S) = 1$，平均花費最少就是新元素最多。

\begin{pseudo}[label=alg:apx-sc]{集合覆蓋的貪心}
\Function{Greedy-Set-Cover}{$X, \mathcal{F}, c$}
    \State $U \gets X$ \Comment{還沒被涵蓋的元素}
    \State $\mathcal{C} \gets \emptyset$
    \While{$U \ne \emptyset$}
        \State 選 $|S \cap U| > 0$ 而且 $c(S) / |S \cap U|$ 最小的 $S \in \mathcal{F}$ \Comment{平手取編號小的}
        \For{$S \cap U$ 裡的每一個元素 $e$}
            \State $\mathit{price}[e] \gets c(S) / |S \cap U|$ \Comment{只為了分析}
        \EndFor
        \State $U \gets U - S$
        \State $\mathcal{C} \gets \mathcal{C} \cup \{S\}$
    \EndWhile
    \State \Return $\mathcal{C}$
\EndFunction
\end{pseudo}

第 6、7 行不影響結果，只是把這一輪花的成本 $c(S)$ 平均分給這一輪新涵蓋的元素，叫做元素的
\textbf{價格}。 每一輪新涵蓋的元素價格加起來剛好是 $c(S)$，所以
\begin{equation}
c(\mathcal{C}) = \sum_{e \in X} \mathit{price}[e].
\label{eq:apx-price}
\end{equation}
用一個長度 $n$ 的布林陣列標記哪些元素還在 $U$ 裡，第 5 行一次掃過所有集合的所有元素，花
$O\bigl(\sum_{S \in \mathcal{F}} |S|\bigr)$；每一輪至少涵蓋一個新元素、也至少用掉一個集合，最多
$\min(n, k)$ 輪，總時間 $O\bigl(\min(n, k) \sum_{S} |S|\bigr)$，是輸入大小的多項式。

\subsection{手算追蹤}

不加權，平手取編號小的。 表中每一格是 $|S_i \cap U|$，方框是這一輪選中的。

\begin{center}
\begin{tabular}{cl*{6}{c}l}
\toprule
輪 & $U$ & $S_1$ & $S_2$ & $S_3$ & $S_4$ & $S_5$ & $S_6$ & 新元素的價格 \\
\midrule
1 & $1..12$ & 3 & 3 & 4 & 4 & \fbox{6} & 4 & $3, 4, 5, 6, 7, 10$ 各 $\frac16$ \\
2 & $1, 2, 8, 9, 11, 12$ & 1 & 1 & \fbox{3} & 1 & 0 & 2 & $9, 11, 12$ 各 $\frac13$ \\
3 & $1, 2, 8$ & 1 & 0 & -- & 1 & -- & \fbox{2} & $1, 2$ 各 $\frac12$ \\
4 & $8$ & 0 & 0 & -- & \fbox{1} & -- & -- & $8$ 是 $1$ \\
\bottomrule
\end{tabular}
\end{center}

貪心選了 $S_5, S_3, S_6, S_4$，$4$ 位。 但 $S_3, S_4, S_6$ 三位就夠了：三個集合各 $4$ 項，沒有重疊，
合起來正好是 $1..12$。 兩位不夠，因為最大的兩個集合加起來只有 $6 + 4 = 10 < 12$ 項。 所以
$|\mathcal{C}^*| = 3$，貪心多用了一位。 錯在第一步：$S_5$ 最大，可是它的工作和其他人重疊很多，
選了它之後，剩下的工作分散在好幾個人身上。

價格加起來是 $6 \cdot \frac16 + 3 \cdot \frac13 + 2 \cdot \frac12 + 1 = 4$，正好是貪心選的人數，
這就是式~\eqref{eq:apx-price}。

\begin{examnote}
集合覆蓋的追蹤題，每一輪寫出還沒被涵蓋的元素和每個集合「還能涵蓋幾個」，不要寫集合原本的大
小；平手時照題目的規則，題目沒說就寫明你的規則（例如編號小的先）。 加權的版本比的是
$c(S) / |S \cap U|$，不是 $c(S)$，也不是 $|S \cap U|$。
\end{examnote}

\subsection{為什麼是 \texorpdfstring{$H_d$}{H\_d}-近似}

記 $H_t = 1 + \frac12 + \frac13 + \cdots + \frac1t$，第 $t$ 個\textbf{調和數}。 由第~1~章的積分
估計，$\ln(t + 1) \le H_t \le \ln t + 1$。

\begin{theorem}
\label{thm:apx-sc}
\proc{Greedy-Set-Cover} 傳回的 $\mathcal{C}$ 滿足 $c(\mathcal{C}) \le H_d \cdot c(\mathcal{C}^*)$，
$\mathcal{C}^*$ 是成本最小的集合覆蓋，$d$ 是最大集合的大小。 因為 $d \le n$，也有
$c(\mathcal{C}) \le (\ln n + 1)\, c(\mathcal{C}^*)$。
\end{theorem}

\begin{proof}
關鍵是下面這個性質：\textbf{任何一個集合 $T \in \mathcal{F}$，它的元素的價格總和不超過
$H_{|T|} \cdot c(T)$。}

把 $T$ 的元素依被涵蓋的先後排成 $t_1, t_2, \ldots, t_s$（$s = |T|$；同一輪被涵蓋的，順序任意）。
看 $t_j$ 被涵蓋的那一輪，演算法選了集合 $S$。 這一輪開始時，$t_j, t_{j+1}, \ldots, t_s$ 都還沒被
涵蓋（它們不比 $t_j$ 早），所以 $|T \cap U| \ge s - j + 1$，選 $T$ 的平均花費
\[
\frac{c(T)}{|T \cap U|} \le \frac{c(T)}{s - j + 1}.
\]
演算法選的 $S$ 平均花費最小，所以 $\mathit{price}[t_j] = c(S) / |S \cap U| \le c(T) / (s - j + 1)$。
對 $j$ 加總，
\[
\sum_{j=1}^{s} \mathit{price}[t_j] \le c(T) \left(\frac1s + \frac1{s-1} + \cdots + \frac11\right)
= H_s\, c(T) \le H_d\, c(T).
\]

接著用式~\eqref{eq:apx-price}。 $\mathcal{C}^*$ 是覆蓋，每個元素至少屬於 $\mathcal{C}^*$ 的一個集合，
價格又都是正的，所以
\[
c(\mathcal{C}) = \sum_{e \in X} \mathit{price}[e]
\le \sum_{T \in \mathcal{C}^*} \sum_{e \in T} \mathit{price}[e]
\le \sum_{T \in \mathcal{C}^*} H_d\, c(T) = H_d\, c(\mathcal{C}^*). \qedhere
\]
\end{proof}

這個證明的下界是「價格」：它沒有直接算出 $c(\mathcal{C}^*)$，而是說明每個最佳解裡的集合，最多只
「吸收」得了 $H_d$ 倍自己成本的價格。 手算例子裡 $\mathcal{C}^* = \{S_3, S_4, S_6\}$，三個集合的價
格和分別是 $\frac16 + 3 \cdot \frac13 = \frac76$、$3 \cdot \frac16 + 1 = \frac32$、
$2 \cdot \frac16 + 2 \cdot \frac12 = \frac43$，都不超過 $H_4 = \frac{25}{12}$，加起來是 $4$。

套到頂點覆蓋：第~\ref{sec:apx-vc}~節「先拿度數最大的」正是不加權的集合覆蓋貪心，所以它是
$H_\Delta$-近似；第~\ref{sec:apx-vc}~節的蜘蛛 $\Delta = 3$，$4 \le H_3 \cdot 3 = 5.5$。 $\Delta$
很大時 $H_\Delta$ 比 $2$ 大，但\textbf{加權}的頂點覆蓋，虛擬碼~\ref{alg:apx-vc} 完全不能用（觀念
題第~\ref{q:apx-wvc}~題），這個貪心還是 $H_\Delta$-近似。

\subsection{近似比是緊的}

\textbf{加權：剛好 $H_n$。} 一個串流平台有 $n$ 個頻道 $e_1, \ldots, e_n$。 可以買全餐方案
$X$，價錢 $1 + \delta$（$\delta > 0$ 很小）；也可以單買頻道 $e_i$，價錢 $1/i$。 輪到還沒買的是
$e_1, \ldots, e_j$ 時，全餐的平均花費是 $(1 + \delta)/j$，單買 $e_i$ 是 $1/i$，最小的是單買 $e_j$
的 $1/j < (1 + \delta)/j$。 所以貪心依序單買 $e_n, e_{n-1}, \ldots, e_1$，總共
$\frac1n + \frac1{n-1} + \cdots + 1 = H_n$；最佳是直接買全餐，$1 + \delta$（$n \ge 2$ 時，不買全餐
就得單買全部，花 $H_n > 1 + \delta$）。 比值 $H_n/(1 + \delta)$，$\delta \to 0$ 時趨近
$H_n$，而這裡 $d = n$，所以定理~\ref{thm:apx-sc} 的 $H_d$ 不能再改小。

\textbf{不加權。} 不加權的貪心也會差到 $\Theta(\log n)$ 倍。 一種造法：元素排成三列，最佳解是
三個「整列」的集合；另外把各行從左到右切成寬度大約每次乘 $\frac32$ 的幾段，每一段的三列合起來是
一個集合，而且每一段都比它左邊剩下的一整列略大。 貪心每一輪都會選目前最右邊的那一段，一共選了
$\Theta(\log n)$ 段（觀念題第~\ref{q:apx-rows}~題）。 所以不加權時，定理~\ref{thm:apx-sc} 的界也只
差常數倍。

集合覆蓋的近似比幾乎不可能再改進：Dinur 與 Steurer（2014）證明，如果 P $\ne$ NP，對任何常數
$\varepsilon > 0$，都沒有多項式時間的 $(1 - \varepsilon)\ln n$-近似演算法。 這個簡單的貪心，已經
差不多是能做到的最好了。

\section{0/1 背包的近似方案}
\label{sec:apx-knap}

前面三個問題的近似比都是固定的數（或隨 $n$ 變大）。 有些問題可以做得更好：使用者自己指定誤差
$\varepsilon$，演算法保證答案在最佳解的 $(1 + \varepsilon)$ 倍以內，$\varepsilon$ 越小跑得越久。

\begin{definition}[近似方案]
\label{def:apx-scheme}
一個演算法多吃一個參數 $\varepsilon > 0$，對每個固定的 $\varepsilon$ 都是 $(1 + \varepsilon)$-近似演算
法，叫\textbf{近似方案}（approximation scheme）。 對每個固定的 $\varepsilon$，時間都是 $n$ 的多項式
（例如 $O(n^{1/\varepsilon})$），叫\textbf{多項式時間近似方案}（PTAS）；時間是 $n$ 和
$1/\varepsilon$ 兩者的多項式（例如 $O(n^3/\varepsilon)$），叫\textbf{完全多項式時間近似方案}
（FPTAS）。
\end{definition}

這一節把第~18~章的 0/1 背包做成 FPTAS：$n$ 件物品，重量 $w_i$、價值 $v_i$ 都是正整數，背包容量
$W$，要讓裝進去的總價值最大。 第~\ref{sec:dp-pseudo-poly}~節說過，以容量為索引的動態規劃要
$\Theta(nW)$，是擬多項式時間。

\subsection{以價值為索引的動態規劃}

第~\ref{sec:dp-pseudo-poly}~節的最後提到另一種表：令 $M[i, p]$ 是「只用前 $i$ 件、總價值剛好是 $p$」的最小總重量，做不
到時是 $\infty$。 第 $i$ 件拿或不拿，和以容量為索引的表一樣分兩種情形：
\[
M[i, p] = \min\bigl(M[i-1, p],\ w_i + M[i-1, p - v_i]\bigr),
\]
第二項只在 $p \ge v_i$ 時才有。 答案是使 $M[n, p] \le W$ 的最大 $p$。

\begin{pseudo}[label=alg:apx-knap-value]{0/1 背包，以價值為索引}
\Function{Knapsack-By-Value}{$w, v, n, W$}
    \State $V \gets v_1 + v_2 + \cdots + v_n$
    \State $M[0, 0] \gets 0$
    \For{$p \gets 1$ \To $V$}
        \State $M[0, p] \gets \infty$
    \EndFor
    \For{$i \gets 1$ \To $n$}
        \For{$p \gets 0$ \To $V$}
            \State $M[i, p] \gets M[i-1, p]$
            \If{$p \ge v_i$ \AND $w_i + M[i-1, p - v_i] < M[i, p]$}
                \State $M[i, p] \gets w_i + M[i-1, p - v_i]$
            \EndIf
        \EndFor
    \EndFor
    \State \Return 使 $M[n, p] \le W$ 的最大 $p$
\EndFunction
\end{pseudo}

表有 $(n + 1)(V + 1)$ 格，時間 $\Theta(nV)$；還原拿了哪些物品的方法和第~18~章的
\proc{Knapsack-Items} 一樣，從 $M[n, p]$ 往回看 $M[i, p]$ 和 $M[i-1, p]$ 相不相等。 這也是擬多項
式的，但時間只和\textbf{價值的總和}有關，和容量、重量無關。 想法於是很自然：價值太大，就把價值
等比例縮小再取整數，讓 $V$ 變小；縮小造成的誤差，用 $\varepsilon$ 控制。

\subsection{縮小價值}

\begin{pseudo}[label=alg:apx-fptas]{0/1 背包的 FPTAS}
\Function{Knapsack-FPTAS}{$w, v, n, W, \varepsilon$}
    \State 刪掉 $w_i > W$ 的物品（一定放不下），$n$ 改成剩下的件數
    \State $K \gets \varepsilon \cdot v_{\max} / n$ \Comment{$v_{\max}$ 是最大的 $v_i$}
    \For{$i \gets 1$ \To $n$}
        \State $v'_i \gets \lfloor v_i / K \rfloor$
    \EndFor
    \State 用 \Call{Knapsack-By-Value}{$w, v', n, W$} 找出 $v'$ 總和最大的可行集合 $S'$
    \State \Return $S'$
\EndFunction
\end{pseudo}

\begin{theorem}
\label{thm:apx-fptas}
\proc{Knapsack-FPTAS} 傳回的 $S'$ 放得進背包，而且 $v(S') \ge (1 - \varepsilon)\, v(S^*)$，$S^*$ 是
最佳解；時間 $O(n^3 / \varepsilon)$。 所以它是 FPTAS（$1/(1-\varepsilon) \le 1 + 2\varepsilon$
當 $\varepsilon \le \frac12$，要 $(1 + \varepsilon')$-近似就取 $\varepsilon = \varepsilon'/2$）。
\end{theorem}

\begin{proof}
第 2 行刪掉的物品不在任何可行解裡，不影響最佳解。 刪完之後，價值最大的那一件自己就放得進背包，
所以 $v_{\max} \le v(S^*)$。

由取整數的定義，$K v'_i \le v_i < K(v'_i + 1)$，也就是每件物品縮小再放大之後，損失不到 $K$。 $S'$ 在
價值 $v'$ 下最大，$S^*$ 也是可行解，所以 $v'(S') \ge v'(S^*)$。 於是
\[
v(S') \ge K\, v'(S') \ge K\, v'(S^*) > \sum_{i \in S^*} (v_i - K) \ge v(S^*) - nK
= v(S^*) - \varepsilon v_{\max} \ge (1 - \varepsilon)\, v(S^*).
\]

時間：每個 $v'_i \le v_i / K \le v_{\max}/K = n/\varepsilon$，所以 $V' = \sum v'_i \le n^2/\varepsilon$，
\proc{Knapsack-By-Value} 花 $\Theta(nV') = O(n^3/\varepsilon)$，其他步驟 $O(n)$。
\end{proof}

證明的下界和前幾節不一樣：最大化問題要的是「最佳解不會比演算法好太多」，這裡用的量是
$K\,v'(S^*)$，它算是最佳解在縮小後的價值，和最佳解只差每件不到 $K$。 誤差的總量 $nK$ 被設計成
$\varepsilon v_{\max}$，而 $v_{\max}$ 又不超過最佳解，這就是 $K$ 取 $\varepsilon v_{\max}/n$ 的原因。

\subsection{例子}

容量 $W = 10$，五件物品，$\varepsilon = \frac12$：

\begin{center}
\begin{tabular}{c*{5}{c}}
\toprule
物品 $i$ & 1 & 2 & 3 & 4 & 5 \\
\midrule
重量 $w_i$ & 6 & 8 & 4 & 9 & 4 \\
價值 $v_i$ & 35 & 60 & 15 & 54 & 29 \\
$v'_i = \lfloor v_i / 6 \rfloor$ & 5 & 10 & 2 & 9 & 4 \\
\bottomrule
\end{tabular}
\end{center}

每件都放得下，$v_{\max} = 60$，$K = \frac12 \cdot 60 / 5 = 6$。 重量不超過 $10$ 的組合只有單獨一件，以
及 $\{1, 3\}$（重 $10$）、$\{1, 5\}$（重 $10$）、$\{3, 5\}$（重 $8$）三對，任何三件都超重。 用 $v'$
比，$\{2\}$ 的 $10$ 最大（$\{1, 5\}$ 和 $\{4\}$ 都是 $9$），演算法傳回 $\{2\}$，價值 $60$。 真正的最佳
解是 $\{1, 5\}$，價值 $64$。 物品 1、5 取整數時各損失 $5$（$35 \to 30$、$29 \to 24$），物品 2 剛好
是 $6$ 的倍數，沒有損失，所以縮小之後 $\{1, 5\}$ 看起來比 $\{2\}$ 差。

結果 $60 \ge (1 - \frac12) \cdot 64 = 32$，符合定理，而且實際的比值 $60/64 \approx 0.94$ 比保證的
$0.5$ 好得多。 證明裡的 $nK$ 假設最佳解的每一件都損失將近 $K$，而且最佳解只比 $v_{\max}$ 大一點，
這兩件事很難同時發生，所以這個界通常不會碰到。 對近似方案來說，重點不在某一個 $\varepsilon$ 的界
緊不緊，而在 $\varepsilon$ 可以任意小：要 $1\%$ 以內，就取 $\varepsilon = 0.01$，代價是時間多
$100$ 倍。

不是每個 NP 難問題都有 FPTAS。 0/1 背包有，是因為它有擬多項式時間的演算法，縮小數值就能把時間
壓下來。 第~28~章提到的\textbf{強 NP 完備}問題（數值很小時仍然是 NP 完備的），只要目標值是整數、
而且不超過輸入大小和數值的某個多項式，除非 P $=$ NP 都沒有 FPTAS；度量 TSP 也已經知道連 PTAS 都沒有（除非 P $=$ NP）。

\section{本章常見錯誤}
\label{sec:apx-pitfalls}

\begin{center}
\begin{tabular}{p{0.24\linewidth}p{0.50\linewidth}p{0.16\linewidth}}
\toprule
錯誤 & 例子 & 見 \\
\midrule
證明裡直接用 $C^*$ 的值 & 寫「設最佳解有 $k$ 個點，演算法選了 $2k$ 個」，卻沒說明演算法的答案和 $k$ 有
什麼關係；應該先找一個下界，再拿演算法的答案和下界比 & 第~\ref{sec:apx-ratio}~節 \\
最大化問題的比值寫反 & 背包寫成 $C \le \rho C^*$；或把 $\frac12$-近似和 $2$-近似當成兩件事 &
第~\ref{sec:apx-ratio}~節 \\
緊的例子只算一邊 & 算出演算法得到 $2k$，就說比值是 $2$，沒有證明最佳解是 $k$ &
第~\ref{sec:apx-ratio}~節 \\
頂點覆蓋只拿一端 & 每條沒蓋住的邊只拿一個端點，星狀圖會差 $n - 1$ 倍 & 第~\ref{sec:apx-vc}~節 \\
極大和最大匹配搞混 & 以為被選中的邊是最大匹配；它只是加不進去的極大匹配 & 第~\ref{sec:apx-vc}~節 \\
路線成本漏了最後一段 & 前序走訪的成本只加到最後一個頂點，沒加回到起點的那段 &
第~\ref{sec:apx-tsp}~節 \\
沒有三角不等式也套 $2$ 倍 & 成本矩陣裡有 $c(u, w) > c(u, v) + c(v, w)$，還說 MST 走訪是 $2$-近似 &
第~\ref{sec:apx-tsp-general}~節 \\
間隙的權重取太小 & 不可近似的證明裡，非邊的成本取 $\rho$ 或 $\rho n$ 以下，「否」的實例也可能落在
$\rho n$ 以內 & 第~\ref{sec:apx-tsp-general}~節 \\
貪心比的量弄錯 & 加權集合覆蓋比 $c(S)$ 或 $|S|$，不是 $c(S)/|S \cap U|$；或用集合原本的大小，不是
還沒被涵蓋的個數 & 第~\ref{sec:apx-sc}~節 \\
$H_d$ 的 $d$ 弄錯 & 把 $d$ 寫成集合個數；$d$ 是最大集合的元素個數 & 第~\ref{sec:apx-sc}~節 \\
沒先刪掉放不下的物品 & FPTAS 的 $v_{\max}$ 來自一件放不進背包的物品，$v_{\max} \le C^*$ 不成立，
誤差界失效 & 第~\ref{sec:apx-knap}~節 \\
\bottomrule
\end{tabular}
\end{center}

\section{觀念題}
\label{sec:apx-exercises}

題型和研究所考試、期中期末考的筆試相近，有選擇、計算、追蹤和證明。 解答在本章最後。

\begin{enumerate}
\item \label{q:apx-mc}（選擇，可能不只一個答案）下列哪些敘述正確？
\begin{enumerate}
\item $2$-近似演算法對每一個輸入，傳回的答案都比最佳解差。
\item 某個最小化問題的最佳值一定 $\ge 10$，而某個演算法對每個輸入都保證 $C \le C^* + 5$，那麼它是
$1.5$-近似演算法。
\item 定理~\ref{thm:apx-tsp-hard} 說明，除非 P $=$ NP，度量 TSP 沒有多項式時間的常數近似演算法。
\item 虛擬碼~\ref{alg:apx-vc} 傳回的頂點個數一定是偶數。
\item 如果 P $=$ NP，最小頂點覆蓋就有多項式時間的 $1$-近似演算法。
\item 每個集合都最多兩個元素時，集合覆蓋的貪心演算法是 $1.5$-近似。
\item 一個演算法的近似比是緊的，表示這個問題沒有任何多項式時間演算法能做得更好。
\end{enumerate}

\item \label{q:apx-vc-trace}（追蹤）$3 \times 3$ 的格子圖：頂點
\[
\begin{matrix} 1 & 2 & 3 \\ 4 & 5 & 6 \\ 7 & 8 & 9 \end{matrix}
\]
左右相鄰、上下相鄰的兩點之間有邊，共 $12$ 條。
\begin{enumerate}
\item 邊表的順序是 $(2,5), (4,7), (6,9)$，接著其他邊依字典順序。 寫出虛擬碼~\ref{alg:apx-vc} 選中的
邊和 $C$。
\item 邊表全部依字典順序，結果呢？
\item 求最小頂點覆蓋的大小，並證明不能更小。
\item 證明：不論邊表的順序如何，$|C| \ge 6$。
\end{enumerate}

\item \label{q:apx-matching}（證明）$M$ 是無向圖 $G$ 的任何一個極大匹配，$M^*$ 是最大匹配。 證明
$|M| \ge |M^*| / 2$，並舉一個 $|M| = |M^*|/2$ 的例子。

\item \label{q:apx-wvc}（證明）\textbf{加權頂點覆蓋}：每個頂點有權重 $w(v) > 0$，要讓覆蓋的總權重最小。
\begin{enumerate}
\item 舉例說明虛擬碼~\ref{alg:apx-vc} 的總權重可以是最佳解的任意多倍。
\item 改成「沒蓋住的邊只拿權重比較小的那一端」，舉例說明它也可以差任意多倍。
\item 說明怎麼用加權的 \proc{Greedy-Set-Cover} 得到 $H_\Delta$-近似，$\Delta$ 是最大度數。
\end{enumerate}

\item \label{q:apx-tsp-trace}（追蹤）六個點的座標是 $1\,(1, 0)$、$2\,(2, 2)$、$3\,(4, 1)$、$4\,(5, 4)$、
$5\,(6, 2)$、$6\,(6, 4)$，成本是曼哈頓距離。
\begin{enumerate}
\item 用 Prim 從 $1$ 開始求最小生成樹（key 一樣時取編號小的），寫出頂點加入的順序和樹邊。
\item 照虛擬碼~\ref{alg:apx-tsp}，以 $1$ 為根，寫出巡迴路線和成本。
\item 改用 $6$ 當根（子節點仍依編號由小到大），路線和成本是多少？
\item 最佳路線的成本是 $18$。 驗證 (b)、(c) 都符合定理~\ref{thm:apx-tsp}。
\end{enumerate}

\item \label{q:apx-tsp-better}（證明）$n \ge 3$，成本滿足三角不等式。
\begin{enumerate}
\item 證明 $c(T) \le (1 - \frac1n)\, c(H^*)$。
\item 由此證明 \proc{Approx-TSP-Tour} 的路線滿足 $c(H) \le 2(1 - \frac1n)\, c(H^*)$。
\item 和第~\ref{sec:apx-tsp}~節的緊例子比較：$n = k + 1$ 時，(b) 的界和實際的比值各是多少？
\end{enumerate}

\item \label{q:apx-gap}（計算）定理~\ref{thm:apx-tsp-hard} 的證明裡，非邊的成本取 $\rho n + 1$。
\begin{enumerate}
\item 如果改成取整數 $x$，$x$ 至少要多大，證明才成立？ 用 $\rho$、$n$ 表示。
\item $\rho = 1$ 時，(a) 的答案是多少？ 這時「近似」的意思是什麼？
\end{enumerate}

\item \label{q:apx-sc-trace}（追蹤）加權集合覆蓋，$X = \{1, \ldots, 8\}$：
\begin{center}
\begin{tabular}{clc}
\toprule
集合 & 元素 & 成本 \\
\midrule
$S_1$ & $1, 3, 4, 5, 6$ & 3 \\
$S_2$ & $6, 8$ & 7 \\
$S_3$ & $1, 3, 6, 7, 8$ & 4 \\
$S_4$ & $2, 4, 5$ & 5 \\
$S_5$ & $1, 8$ & 7 \\
\bottomrule
\end{tabular}
\end{center}
\begin{enumerate}
\item 用虛擬碼~\ref{alg:apx-sc} 求覆蓋，寫出每一輪各集合的 $c(S)/|S \cap U|$ 和每個元素的價格。
\item 求最佳覆蓋，並檢查定理~\ref{thm:apx-sc} 證明裡的性質：最佳解的每個集合 $T$，價格和
$\le H_{|T|}\, c(T)$。
\end{enumerate}

\item \label{q:apx-rows}（證明）$3L$ 個元素排成 $3$ 列、$L$ 行。 $R_1, R_2, R_3$ 各是一整列。 把 $L$ 行從左到
右切成寬度 $w_1, \ldots, w_k$ 的段，$B_j$ 是第 $j$ 段的 $3w_j$ 個元素；寬度取 $w_1 = 1$，
$w_j = \lfloor s_{j-1}/2 \rfloor + 1$，$s_j = w_1 + \cdots + w_j$，$L = s_k$。 集合的編號順序是
$R_1, R_2, R_3, B_1, \ldots, B_k$，不加權。
\begin{enumerate}
\item 寫出 $k = 4$ 的寬度，並追蹤貪心。
\item 證明 $k \ge 3$ 時最佳解是 $3$ 個集合，貪心選了 $k$ 個。
\item 證明 $k > \log_{3/2}(n/6)$，$n = 3L$；所以貪心和最佳的比值是 $\Omega(\log n)$。
\end{enumerate}

\item \label{q:apx-knap-half}（證明）0/1 背包，每件物品都放得進背包（$w_i \le W$）。 把物品依
$v_i / w_i$ 由大到小排好，從第一件開始依序放，遇到第一件放不下的就停，得到集合 $P$。 令
$A = \max\bigl(v(P),\ v_{\max}\bigr)$。
\begin{enumerate}
\item 證明 $A \ge \frac12 C^*$。（提示：第~\ref{sec:gr-knapsack}~節的分數背包。）
\item 舉一族例子，讓 $C^* / A$ 任意接近 $2$。
\item 只用 $v(P)$、不和 $v_{\max}$ 比，為什麼不行？
\end{enumerate}

\item \label{q:apx-additive}（證明）0/1 背包的價值都是正整數。 證明：如果 P $\ne$ NP，就沒有多項式時
間演算法能保證 $C^* - C \le 1$。 （提示：把每件的價值乘以 $2$。 0/1 背包的決策版本是 NP 完備的。）

\end{enumerate}

\section{實作練習}
\label{sec:apx-practice}

Aki-Judge 目前沒有和本章對應的題目。 請用觀念題和範例~\ref{ex:apx-vc} 練習；例如把
\texttt{approx\_vc} 換成「先拿度數最大的」，在同樣的隨機圖上比較兩種做法。

\section{讀完你應該會}
\label{sec:apx-checklist}

\begin{itemize}
\item 寫出近似比的定義，分清楚最小化和最大化問題的不等式方向；說明「近似比是緊的」要證明哪兩件事。
\item 手算頂點覆蓋的 $2$-近似，用被選中的邊（極大匹配）當下界證明 $2$ 倍，舉出比值剛好是 $2$ 的例子。
\item 手算度量 TSP 的 MST 前序走訪，寫出完整走訪、抄捷徑和總成本；證明 $2$ 倍並說明三角不等式用在
哪裡。
\item 用哈密頓環的歸約證明一般 TSP 沒有常數近似（除非 P $=$ NP），說出「間隙」的取法。
\item 手算加權與不加權的集合覆蓋貪心和每個元素的價格，證明 $H_d$ 倍，並用串流方案的例子說明它
是緊的。
\item 說出 PTAS 和 FPTAS 的差別，照縮小價值的做法求背包的近似解，並證明 $(1 - \varepsilon)$ 的保證。
\end{itemize}

\textbf{延伸閱讀}（書末會整理完整的參考資料）

\begin{itemize}
\item Thomas H. Cormen 等，\textit{Introduction to Algorithms}，第 4 版，第 35 章：頂點覆蓋、TSP、
集合覆蓋，另外有用線性規劃做加權頂點覆蓋（35.4 節）和子集合和的 FPTAS（35.5 節）。
\item Steven S. Skiena，\textit{The Algorithm Design Manual}，第 3 版，第 12 章 12.1--12.6 節：先講直
覺再講證明，也介紹了區域搜尋、模擬退火這些沒有保證但很實用的啟發式方法。
\item Jon Kleinberg、Éva Tardos，\textit{Algorithm Design}，第 11 章：負載平衡、中心選址、定價法
（pricing method）做加權頂點覆蓋和背包的近似方案，證明寫得很細。
\item Vijay V. Vazirani，\textit{Approximation Algorithms}，Springer，2001，第 1--3、8 章：集合覆蓋、
頂點覆蓋、TSP 和背包的近似演算法，是這個領域的入門專書。
\item David P. Williamson、David B. Shmoys，\textit{The Design of Approximation Algorithms}，
Cambridge University Press，2011（作者網站提供免費電子版），第 1--3 章：用集合覆蓋示範好幾種設計
方法，第 2 章有 Christofides 演算法的完整證明。
\end{itemize}

\addcontentsline{toc}{section}{觀念題解答}
\section*{觀念題解答}

\begin{enumerate}
\item (b)、(d)、(e)、(f) 正確。 (a)：$2$-近似只保證不超過 $2$ 倍，也可能剛好是最佳解。 例如三角形：虛擬碼~\ref{alg:apx-vc}
選中一條邊、拿到兩個點，最小頂點覆蓋也是兩個點。 (b)：$C \le C^* + 5
\le C^* + \frac12 C^* = 1.5\, C^*$。 (c)：定理說的是\textbf{一般的} TSP；度量 TSP 有
定理~\ref{thm:apx-tsp} 的 $2$-近似。 (d)：每次加兩個點，而且兩點原本都不在 $C$ 裡。 (e)：P $=$ NP
時，最小頂點覆蓋可以在多項式時間求出（用決策版本對 $k = 0, 1, \ldots, n$ 一一問，再用第~\ref{sec:npc-search}~節
的自我歸約逐點確定覆蓋），比值是 $1$。 (f)：$d \le 2$，$H_2 = 1.5$。 (g)：緊只說這個演算法的分析不能再改
進，例如度量 TSP 的 MST 走訪是緊的 $2$，Christofides 卻是 $1.5$。

\item
\begin{enumerate}
\item 選中 $(2,5)$、$(4,7)$、$(6,9)$，之後依字典順序：$(1,2)$ 有 $2$、$(1,4)$ 有 $4$、$(2,3)$ 有 $2$、
$(3,6)$ 有 $6$、$(4,5)$、$(5,6)$、$(5,8)$ 有 $5$、$(7,8)$ 有 $7$、$(8,9)$ 有 $9$，都已蓋住。
$C = \{2, 4, 5, 6, 7, 9\}$，$6$ 個點。
\item 選中 $(1,2)$、$(3,6)$、$(4,5)$、$(7,8)$，$C = \{1, \ldots, 8\}$，$8$ 個點。
\item $\{2, 4, 6, 8\}$ 是頂點覆蓋：每條邊都連著一個「邊中點」（$2, 4, 6, 8$）。 $(1,2)$、$(3,6)$、
$(4,7)$、$(5,8)$ 沒有共用端點，至少要 $4$ 個點。 最小是 $4$。
\item 被選中的邊 $M$ 是極大匹配，$|C| = 2|M|$，只要說明 $|M| \ge 3$。 假設 $|M| \le 2$，至少有
$5$ 個點不是 $M$ 的端點；它們兩兩不相鄰（否則它們之間的邊還沒被蓋住，演算法會選它），是一組
至少 $5$ 個點的\textbf{獨立集}（兩兩不相鄰的頂點，第~28~章）。 把頂點分成 $\{1,2\}$、$\{3,6\}$、
$\{9,8\}$、$\{7,4\}$、$\{5\}$ 五組，前四組各是一條邊，獨立集每組最多取一個，所以最多 $5$ 個；
剛好 $5$ 個時一定取了 $5$，而 $5$ 和 $2, 4, 6, 8$ 都相鄰，於是只能是 $\{1, 3, 5, 7, 9\}$。 這樣
$M$ 的 $4$ 個端點就是 $2, 4, 6, 8$，但這四點兩兩不相鄰，$M$ 裡的邊無從連起，矛盾。
\end{enumerate}

\item $M^*$ 的每一條邊 $e$ 都和 $M$ 的某條邊共用端點：$e \in M$ 時顯然；否則 $e$ 加不進 $M$，表示
它的某一端已經被 $M$ 用掉。 讓 $e$ 對應到「它碰到的 $M$ 的某個端點」。 $M^*$ 的邊兩兩不共用端點，
所以不同的 $e$ 對應到不同的端點；$M$ 的端點一共 $2|M|$ 個，所以 $|M^*| \le 2|M|$。 例子：路徑
$a - b - c - d$，$M = \{(b, c)\}$ 是極大匹配，$M^* = \{(a, b), (c, d)\}$。 這和
定理~\ref{thm:apx-vc} 是同一個論證：極大匹配的端點同時是頂點覆蓋。

\item
\begin{enumerate}
\item 一條邊 $(1, 2)$，$w(1) = 1$、$w(2) = N$。 演算法兩點都拿，總權重 $N + 1$；只拿 $1$ 就夠，
比值 $N + 1$，$N$ 可以任意大。
\item 星狀圖：中心權重 $2$，$N$ 片葉各權重 $1$。 每條邊比較輕的是葉，於是 $N$ 片葉全拿，總權重
$N$；只拿中心是 $2$。 比值 $N/2$。
\item 照第~\ref{sec:apx-sc}~節的對應：$X = E$，頂點 $v$ 對應集合 $S_v$（和 $v$ 關聯的邊），成本
$w(v)$。 頂點覆蓋就是集合覆蓋，最大的集合有 $\Delta$ 個元素，由定理~\ref{thm:apx-sc} 得到
$H_\Delta$ 倍。 每一輪選的是 $w(v)$ 除以「$v$ 還沒蓋住的邊數」最小的頂點。
\end{enumerate}

\item
\begin{enumerate}
\item 加入順序 $1, 2, 3, 5, 6, 4$（key 依序是 $0, 3, 3, 3, 2, 1$）。 樹邊 $(1,2)\ 3$、$(2,3)\ 3$、
$(3,5)\ 3$、$(5,6)\ 2$、$(6,4)\ 1$，$c(T) = 12$；這棵樹是一條路徑 $1 - 2 - 3 - 5 - 6 - 4$。
\item 前序 $1, 2, 3, 5, 6, 4$，路線 $1 \to 2 \to 3 \to 5 \to 6 \to 4 \to 1$，成本
$3 + 3 + 3 + 2 + 1 + 8 = 20$。
\item 最小生成樹只有這一棵，換根只改變父子方向。 以 $6$ 為根，$6$ 的子節點是 $4, 5$；$5$ 的子節點是 $3$，$3$ 的子節點是 $2$，$2$ 的子節點是 $1$。 前序
$6, 4, 5, 3, 2, 1$，成本 $c(6,4) + c(4,5) + c(5,3) + c(3,2) + c(2,1) + c(1,6) = 1 + 3 + 3 + 3 + 3 + 9
= 22$。
\item $20 \le 22 \le 2 \cdot 12 = 24 \le 2 \cdot 18 = 36$。 同一棵樹，根不同，路線成本可以不同，
但都不超過 $2c(T)$。 最佳成本 $18$ 也可以用正文的方法驗證下界：橫向至少 $2 \times 5$，縱向至少
$2 \times 4$。
\end{enumerate}

\item
\begin{enumerate}
\item $H^*$ 有 $n$ 條邊，最貴的一條 $e$ 成本至少是平均值 $c(H^*)/n$。 拿掉 $e$ 得到一棵生成樹，所以
$c(T) \le c(H^*) - c(H^*)/n$。
\item $c(H) \le 2c(T) \le 2(1 - \frac1n)\, c(H^*)$。
\item 界是 $2(1 - \frac1{k+1}) = \frac{2k}{k+1}$，實際比值 $\frac{2k}{k+2}$，兩者只差一點點，$k$ 變大
時都趨近 $2$。
\end{enumerate}

\item
\begin{enumerate}
\item 「是」的實例最佳值仍然是 $n$，$A$ 傳回 $\le \rho n$。 「否」的實例每條路線至少一段成本 $x$，
其他 $n - 1$ 段至少 $1$，成本 $\ge x + n - 1$。 要讓它 $> \rho n$，需要 $x > (\rho - 1)n + 1$，最小的整數是
$x = \lfloor (\rho - 1)n \rfloor + 2$。 證明裡的 $\rho n + 1$ 比這個大，當然成立；取得比較大只是
為了好記。
\item $\rho = 1$ 時 $x = 2$：非邊成本 $2$、邊成本 $1$，最佳值 $= n$ 若且唯若有哈密頓環。 $1$-近似
就是求出最佳解，這說明一般 TSP（甚至只有成本 $1$、$2$ 的 TSP）求最佳解是 NP 難的。 注意成本
只有 $1$、$2$ 時滿足三角不等式，所以這也說明度量 TSP 求最佳解是 NP 難的；它只是沒辦法放大到
$\rho > 1$。
\end{enumerate}

\item
\begin{enumerate}
\item 第 1 輪：$S_1\ 3/5$、$S_2\ 7/2$、$S_3\ 4/5$、$S_4\ 5/3$、$S_5\ 7/2$，選 $S_1$，元素
$1, 3, 4, 5, 6$ 價格各 $\frac35$。 剩 $\{2, 7, 8\}$。 第 2 輪：$S_2\ 7/1$、$S_3\ 4/2$、$S_4\ 5/1$、
$S_5\ 7/1$，選 $S_3$，$7, 8$ 各 $2$。 剩 $\{2\}$，只有 $S_4$ 含它，選 $S_4$，$2$ 的價格 $5$。 覆蓋
$\{S_1, S_3, S_4\}$，成本 $3 + 4 + 5 = 12$，價格和 $5 \cdot \frac35 + 2 \cdot 2 + 5 = 12$。
\item 元素 $2$ 只在 $S_4$ 裡，元素 $7$ 只在 $S_3$ 裡，所以 $S_3$、$S_4$ 一定要選，成本 $9$，而
$S_3 \cup S_4 = X$。 最佳成本 $9$，貪心多花了 $3$。 檢查：$S_3$ 的價格和
$3 \cdot \frac35 + 2 + 2 = 5.8 \le H_5 \cdot 4 \approx 9.13$；$S_4$ 的價格和
$5 + \frac35 + \frac35 = 6.2 \le H_3 \cdot 5 \approx 9.17$。 兩者加起來 $12$，正好是貪心的成本。
\end{enumerate}

\item
\begin{enumerate}
\item 寬度 $1, 1, 2, 3$，$L = 7$，$n = 21$。 第 1 輪：每列還有 $7$ 個，$B_1, B_2, B_3, B_4$ 各
$3, 3, 6, 9$ 個，選 $B_4$。 第 2 輪：每列 $4$ 個，$B_3$ 有 $6$ 個，選 $B_3$。 第 3 輪：每列 $2$ 個，
$B_1$、$B_2$ 各 $3$ 個，平手選 $B_1$。 第 4 輪：每列 $1$ 個，選 $B_2$。 貪心 $4$ 個，最佳 $3$ 個。
\item 三個列集合是覆蓋。 一個覆蓋如果少用了某一列 $R_r$，第 $r$ 列的元素只能靠段集合蓋，每一行都
屬於某一段，所以 $k$ 個段集合全都要用，至少 $k \ge 3$ 個。 所以最佳是 $3$。 貪心：寬度的規則保證
$2w_j > s_{j-1}$，也就是 $3w_j > s_j$，而且寬度不會變小。 用歸納法，假設已經選了
$B_k, \ldots, B_{j+1}$，還沒涵蓋的是左邊 $s_j$ 行：每個列集合還能涵蓋 $s_j$ 個，$B_j$ 能涵蓋
$3w_j > s_j$ 個，比 $B_j$ 左邊的段都不少。 所以貪心選 $B_j$（只有 $w_1 = w_2$ 時會先選 $B_1$，結果
一樣），從來不選列集合，一共 $k$ 個。
\item $w_j \le s_{j-1}/2 + 1$，所以 $s_j \le \frac32 s_{j-1} + 1$。 對 $j$ 歸納可得
$s_j \le 2\bigl((\frac32)^j - 1\bigr)$（$j = 1$ 時 $1 \le 1$；歸納步驟
$\frac32 \cdot 2\bigl((\frac32)^{j-1} - 1\bigr) + 1 = 2(\frac32)^j - 2$）。 所以
$L = s_k < 2(\frac32)^k$，$k > \log_{3/2}(L/2) = \log_{3/2}(n/6)$。 比值 $k/3 = \Omega(\log n)$。
\end{enumerate}

\item
\begin{enumerate}
\item 設第一件放不下的是排序後的第 $j$ 件（沒有這一件的話 $P$ 就是全部物品，$v(P) = C^*$）。 由
定理~\ref{thm:gr-frac}，分數背包的最佳值是：前 $j - 1$ 件全拿，第 $j$ 件拿一部分，所以分數背包的
最佳值 $< v(P) + v_j$。 0/1 背包的解也是分數背包的解，$C^* \le$ 分數背包的最佳值
$< v(P) + v_j \le v(P) + v_{\max} \le 2A$。
\item 容量 $2N$：物品 $a$ 重 $1$、值 $2$；物品 $b$、$c$ 各重 $N$、值 $N$。 依比值，$a$ 排第一（$2$），
$b$、$c$ 都是 $1$。 $P = \{a, b\}$（重 $N + 1$），放 $c$ 會超重，停。 $A = \max(N + 2, N) = N + 2$，
最佳是 $\{b, c\}$，$2N$。 比值 $\frac{2N}{N + 2} \to 2$。
\item 容量 $N$：物品 $a$ 重 $1$、值 $2$；物品 $b$ 重 $N$、值 $N$。 $P = \{a\}$，$v(P) = 2$，最佳是
$N$，比值 $N/2$，任意大。 第~19~章觀念題的反例就是這種情形。
\end{enumerate}

\item 假設演算法 $A$ 保證 $C^* - C \le 1$。 給一個背包實例，把每件的價值乘以 $2$，其他不變。 新實例
的可行解和舊的一樣，每個解的價值都是偶數，最佳值是 $2C^*_{\text{舊}}$。 $A$ 在新實例上的答案 $C$ 是偶數，
$2C^*_{\text{舊}} - C \le 1$ 又是偶數，所以等於 $0$：$A$ 找到了最佳解，價值除以 $2$ 就是原實例的最
佳解。 這樣在多項式時間解出了 0/1 背包，它的決策版本是 NP 完備的，所以 P $=$ NP。 乘上 $2$ 這一招
說明「加法的誤差」對有數值的問題沒有意義，近似演算法都用比值。

\end{enumerate}


% 書末（不編章號）
% % 書末：參考資料與延伸閱讀（不編章號）。全書寫完時再整理，格式照 C++ 演算法解題 2。
\chapter*{參考資料與延伸閱讀}
\addcontentsline{toc}{chapter}{參考資料與延伸閱讀}


\end{document}
